\documentclass[11pt,a4paper,openany]{book}
\usepackage[margin=25mm]{geometry}
\usepackage{fontspec}
\setmainfont{OLMarathiSerif-Regular.ttf}[Path=../../fonts/,Script=Devanagari,Language=Marathi,BoldFont=OLMarathiSerif-Bold.ttf,ItalicFont=OLMarathiSerif-Regular.ttf]
\setsansfont{OLMarathiSerif-Regular.ttf}[Path=../../fonts/,Script=Devanagari,Language=Marathi,BoldFont=OLMarathiSerif-Bold.ttf]
\usepackage{amsmath,amssymb,amsthm,nicefrac,graphicx,tikz}
\usepackage{../../upstream/sty/ptolemaicastronomy}
\usepackage{multirow}
\usepackage{stmaryrd}
\usepackage{mathdots}
\usepackage{../../upstream/sty/bussproofs-extra}
\usepackage[tableaux]{prooftrees}
\usepackage{flafter}
\usepackage{rotating}
\usepackage[unicode,colorlinks=true,linkcolor=blue,urlcolor=blue]{hyperref}
\usepackage{microtype}
\setlength{\parindent}{0pt}
\setlength{\parskip}{0.45em}
\linespread{1.18}
\setlength{\emergencystretch}{3em}
\renewcommand{\contentsname}{अनुक्रमणिका}
\renewcommand{\chaptername}{प्रकरण}
\renewcommand{\partname}{भाग}
\renewcommand{\thepart}{\arabic{part}}
\renewcommand{\figurename}{आकृती}
\renewcommand{\proofname}{सिद्धता}
\theoremstyle{definition}
\newtheorem{defn}{व्याख्या}[section]
\newtheorem{ex}[defn]{उदाहरण}
\newtheorem{prop}[defn]{विधान}
\newtheorem{thm}[defn]{प्रमेय}
\newtheorem{prob}{सराव}[chapter]
\newenvironment{explain}{}{}
\newenvironment{digress}{\par\smallskip}{\par\smallskip}
\newcommand{\Setabs}[2]{\{#1:#2\}}
\newcommand{\Pow}[1]{\wp(#1)}
\newcommand{\tuple}[1]{\langle #1\rangle}
\newcommand{\Nat}{\mathbb{N}}
\newcommand{\unitline}{\text{L}}
\newcommand{\unitsquare}{\text{S}}
\newcommand{\Int}{\mathbb{Z}}
\newcommand{\Rat}{\mathbb{Q}}
\newcommand{\Real}{\mathbb{R}}
\newcommand{\PosInt}{\mathbb{Z}^{+}}
\newcommand{\Bin}{\mathbb{B}}
\newcommand{\len}[1]{\mathrm{len}(#1)}
\newcommand{\lif}{\mathbin{\to}}
\colorlet{oldiagcolorA}{black}
\colorlet{oldiagcolorB}{gray}
\definecolor{oldiagcolorC}{HTML}{a81c21}
\newcommand{\olasset}[2][0.52\textwidth]{\centering\resizebox{#1}{!}{\input{../../upstream/#2}}}
\hypersetup{pdftitle={मुक्त तर्कशास्त्र: संपूर्ण मराठी आवृत्ती},pdfauthor={Open Logic Project; OpenAI Codex: GPT-5.6 Sol, GPT-6 Sol; corrective review GPT-6.1 Sol; Ultra effort}}
\newenvironment{intro}{}{}
\newcommand{\Id}[1]{\mathord{\mathrm{Id}_{#1}}}
\newcommand{\emptyseq}{\Lambda}
\newcommand{\liff}{\mathbin{\leftrightarrow}}
\newcommand{\equivrep}[2]{[#1]_{#2}}
\newcommand{\equivclass}[2]{#1/_{\!{#2}}}
\newcommand{\funrestrictionto}[2]{#1\mathord{\restriction}_{#2}}
\newcommand{\funimage}[2]{#1[#2]}
\newcommand{\citeyear}[1]{\hyperref[bib:#1]{1965}}
\newcommand{\dom}[1]{\mathrm{dom}(#1)}
\newtheorem{cor}[defn]{अनुप्रमेय}
\newcommand{\ran}[1]{\mathrm{ran}(#1)}
\newcommand{\comp}[2]{#2\circ #1}
\newcommand{\pto}{\mathrel{\ooalign{\hfil$\mapstochar\mkern 5mu$\hfil\cr$\to$}}}
\newcommand{\fdefined}{\downarrow}
\newcommand{\fundefined}{\uparrow}
\newcommand{\card}[1]{\left|#1\right|}
\newcommand{\cardle}[2]{#1 \preceq #2}
\newcommand{\cardless}[2]{#1 \prec #2}
\newcommand{\cardeq}[2]{#1 \approx #2}
\newcommand{\cardneq}[2]{#1 \not\approx #2}
\newcommand{\Intequiv}{\sim}
\newcommand{\Ratequiv}{\backsim}
\newcommand{\Realequiv}{\Bumpeq}
\newcommand{\defis}{=}
\newcommand{\closureofunder}[2]{\mathrm{clo}_{#1}(#2)}
\newcommand{\Closureofunder}[2]{\mathrm{Clo}_{#1}(#2)}
\newtheorem{lem}[defn]{पूर्वप्रमेय}
\newenvironment{editorial}{\begin{quote}\small\itshape}{\end{quote}}
\definecolor{oldiagcolorD}{HTML}{1973ba}
\definecolor{oldiagcolorE}{HTML}{4d6f39}
\newlength{\olphotowidth}\setlength{\olphotowidth}{0.35\textwidth}

\newcommand{\applytofirst}[2]{{\expandafter#1#2}}
\newcommand{\Struct}[1]{\applytofirst{\mathfrak}{#1}}
\newcommand{\Lang}[1]{\applytofirst{\mathcal}{#1}}
\newcommand{\Obj}[1]{\mathsf{#1}}
\newcommand{\PVar}{\mathrm{At}_0}
\NewDocumentCommand{\Frm}{o}{\IfNoValueTF{#1}{\mathrm{Frm}}{\mathrm{Frm}(\Lang{#1})}}
\newcommand{\True}{\ensuremath{\mathbb{T}}}
\newcommand{\False}{\ensuremath{\mathbb{F}}}
\newcommand{\lfalse}{\bot}
\newcommand{\ltrue}{\top}
\newcommand{\Ptemp}{\mathord{\mathsf{P}}}
\newcommand{\Htemp}{\mathord{\mathsf{H}}}
\newcommand{\Ftemp}{\mathord{\mathsf{F}}}
\newcommand{\Gtemp}{\mathord{\mathsf{G}}}
\newcommand{\Since}{\mathord{\mathsf{S}}}
\newcommand{\Until}{\mathord{\mathsf{U}}}
\newcommand{\Knows}{\mathord{\mathsf{K}}}
\newcommand{\EKnows}{\mathord{\mathsf{E}}}
\newcommand{\CKnows}{\mathord{\mathsf{C}}}
\DeclareSymbolFont{symbolsC}{U}{ntxsyc}{m}{n}
\DeclareMathSymbol{\fishhookright}{\mathbin}{symbolsC}{74}
\newcommand{\ident}{\equiv}
\newcommand{\subst}[2]{#1/#2}
\newcommand{\SSubst}[2]{#1[#2]}
\newcommand{\Subst}[3]{#1[\subst{#2}{#3}]}
\NewDocumentCommand{\Entails}{t{/} o}{%
  \IfBooleanTF{#1}{\IfNoValueTF{#2}{\nvDash}{\nvDash_{#2}}}{\IfNoValueTF{#2}{\vDash}{\vDash_{#2}}}}
\newcommand{\pAssign}[1]{\applytofirst{\mathfrak}{#1}}
\NewDocumentCommand{\pValue}{m o d() o}{%
  \overline{\pAssign{#1}}%
  \IfNoValueTF{#2}{\IfNoValueF{#4}{_{#4}}}{_{#2}}%
  \IfNoValueF{#3}{(#3)}}
\NewDocumentCommand{\pSat}{t{/} m m o}{%
  \pAssign{#2}\IfBooleanTF{#1}{\nvDash}{\vDash}%
  \IfNoValueF{#4}{_{#4}}#3}
\NewDocumentCommand{\tf}{m o}{%
  \widetilde{#1}\IfNoValueF{#2}{_{#2}}}
\newcommand{\indfrm}{}
\NewDocumentCommand{\indcase}{m m +m}{\renewcommand{\indfrm}{#1}$#1 \ident #2$: #3}
\def\startycommalist{\def\ycomma{\def\ycomma{, }}}

\newcommand{\Sequent}{\Rightarrow}
\newcommand{\nSequent}{\mid}
\NewDocumentCommand{\iR}{m m o}{\ensuremath{{#1\IfNoValueTF{#3}{}{_{#3}}}{#2}}}
\renewcommand{\fCenter}{\ensuremath{\,\Sequent\,}}
\NewDocumentCommand{\LeftR}{m}{\ensuremath{{#1}\mathrm{L}}}
\NewDocumentCommand{\RightR}{m}{\ensuremath{{#1}\mathrm{R}}}
\newcommand{\Weakening}{\ensuremath{\mathrm{W}}}
\NewDocumentCommand{\Intro}{m o}{\ensuremath{{#1}\mathrm{Intro}\IfNoValueTF{#2}{}{_{#2}}}}
\NewDocumentCommand{\Elim}{m o}{\ensuremath{{#1}\mathrm{Elim}\IfNoValueTF{#2}{}{_{#2}}}}
\newcommand{\FalseInt}{\ensuremath{\lfalse_I}}
\NewDocumentCommand{\Discharge}{m m}{[#1]^{#2}}
\NewDocumentCommand{\DischargeRule}{m m}{\RightLabel{#1}\LeftLabel{\scriptsize $#2$}}
\NewDocumentCommand{\sFmla}{m m o}{\ensuremath{\IfNoValueTF{#3}{}{#3\,}\hbox to.8em{\ensuremath{#1}\hfil} #2}}
\NewDocumentCommand{\pFmla}{m m m}{\ensuremath{\hskip 3em{\llap{$#3$}\,}\hbox to1.3em{\ensuremath{#1}\hfil} #2}}
\NewDocumentCommand{\TRule}{m m o}{\ensuremath{{#2}{#1}\IfNoValueTF{#3}{}{\, #3}}}
\newcommand{\TAss}{\text{गृहीतक}}
\NewDocumentCommand{\Proves}{t{/} o}{%
  \IfBooleanTF{#1}{\IfNoValueTF{#2}{\nvdash}{\nvdash_{#2}}}{\IfNoValueTF{#2}{\vdash}{\vdash_{#2}}}}
\newcommand{\formula}[1]{#1}
\renewenvironment{prooftree}{\begin{center}\bottomAlignProof}{\DisplayProof\end{center}}
\newenvironment{oltableau}{\center\tableau{}}{\endtableau\endcenter}
\newenvironment{derivation}{~\begin{trivlist}\item\begin{tabular}[b]{@{}rll@{}}}{\end{tabular}\end{trivlist}}

\NewDocumentCommand{\lexists}{t{!} o o}{%
  \exists\IfBooleanT{#1}{\mathord{!}}%
  \IfNoValueF{#2}{#2}\IfNoValueF{#3}{\,#3}}
\NewDocumentCommand{\lforall}{o o}{%
  \IfNoValueTF{#1}{\forall}{\forall #1}\IfNoValueF{#2}{\,#2}}
\NewDocumentCommand{\eq}{t{/} o o}{%
  \IfNoValueTF{#3}{\IfBooleanTF{#1}{\neq}{=}}{%
    \IfBooleanTF{#1}{#2\neq#3}{#2=#3}}}
\NewDocumentCommand{\Atom}{m m}{\mathord{#1}(#2)}
\NewDocumentCommand{\Assign}{m m}{\mathord{#1^{\Struct{#2}}}}
\NewDocumentCommand{\Sat}{t{/} m m o}{%
  \IfBooleanTF{#1}{%
    \IfNoValueTF{#4}{\Struct{#2}\nvDash#3}{\Struct{#2},#4\nvDash#3}}{%
    \IfNoValueTF{#4}{\Struct{#2}\vDash#3}{\Struct{#2},#4\vDash#3}}}
\NewDocumentCommand{\varAssign}{m m m o}{%
  \IfNoValueTF{#4}{#1\sim_{#3}#2}{#1=#2[^{#4}/#3]}}
\NewDocumentCommand{\Value}{m m o}{%
  \IfNoValueTF{#3}{\mathrm{Val}^{\Struct{#2}}(#1)}{%
    \mathrm{Val}^{\Struct{#2}}_{#3}(#1)}}
\NewDocumentCommand{\Log}{m o}{%
  \ensuremath{\mathbf{#1}\IfNoValueF{#2}{_{#2}}}}
\newcommand{\LogCL}{\Log{C}}
\newcommand{\Undef}{\ensuremath{\mathbb{U}}}
\newcommand{\LogLuk}{\Log{\textbf{\L}}}
\newcommand{\LogGod}{\Log{G}}
\newcommand{\LogKs}{\Log{Ks}}
\newcommand{\LogKw}{\Log{Kw}}
\newcommand{\LogLP}{\Log{LP}}
\newcommand{\LogRM}{\Log{RM}}
\newcommand{\LogHal}{\Log{Hal}}
\newcommand{\Contraction}{\ensuremath{\mathrm{C}}}
\newcommand{\Exchange}{\ensuremath{\mathrm{X}}}
\newcommand{\Cut}{\ensuremath{\mathrm{Cut}}}
\newcommand{\FalseCl}{\ensuremath{\lfalse_C}}
\newenvironment{defish}{}{}
\newcommand{\readerexternalref}[1]{\textit{मूळ ग्रंथातील संबंधित विभाग}}
\makeatletter
\renewcommand{\@pnumwidth}{2em}
\renewcommand*{\l@section}{\@dottedtocline{1}{1.5em}{2.8em}}
\makeatother

\newtheorem{rem}[defn]{टीप}
\newenvironment{history}{\begin{quote}\small}{\end{quote}}
\newcommand{\MP}{\textsc{mp}}
\newcommand{\QR}{\textsc{qr}}
\newcommand{\Hyp}{\textsc{Hyp}}
\newcommand{\Nec}{\textsc{nec}}
\newcommand{\Var}{\mathrm{Var}}
\NewDocumentCommand{\Trm}{o}{\IfNoValueTF{#1}{\mathrm{Trm}}{\mathrm{Trm}(\Lang{#1})}}
\NewDocumentCommand{\Sent}{o}{\IfNoValueTF{#1}{\mathrm{Sent}}{\mathrm{Sent}(\Lang{#1})}}
\newcommand{\Th}[1]{\mathbf{#1}}
\newcommand{\Theory}[1]{\mathrm{Th}(\Struct{#1})}
\newcommand{\substruct}{\subseteq}
\NewDocumentCommand{\elemequiv}{t{/} o}{%
  \IfBooleanTF{#1}{\IfNoValueTF{#2}{\not\equiv}{\not\equiv_{#2}}}{%
  \IfNoValueTF{#2}{\equiv}{\equiv_{#2}}}}
\NewDocumentCommand{\iso}{t{/} o}{%
  \IfBooleanTF{#1}{\IfNoValueTF{#2}{\not\simeq}{\not\simeq_{#2}}}{%
  \IfNoValueTF{#2}{\simeq}{\simeq_{#2}}}}
\newcommand{\QuantRank}[1]{\mathrm{qr}(#1)}
\newcommand{\Domain}[1]{\left|\Struct{#1}\right|}
\newcommand{\Expan}[2]{(\Struct{#1},#2)}
\newcommand{\nszero}{\mathbf{z}}
\newcommand{\nssucc}{*}
\newcommand{\nsplus}{\oplus}
\newcommand{\nstimes}{\otimes}
\newcommand{\nsless}{\varolessthan}
\newcommand{\concat}{\frown}
\NewDocumentCommand{\lambd}{o o}{%
  \IfNoValueTF{#1}{\lambda}{\lambda #1}\IfNoValueF{#2}{.\,#2}}
\newcommand{\num}[1]{\overline{#1}}

\newtheorem{conv}[defn]{संकेतनरूढी}
\NewDocumentCommand{\redone}{o}{%
  \IfNoValueTF{#1}{\xrightarrow{}}{\xrightarrow{#1}}}
\NewDocumentCommand{\xrightarrowdbl}{o m}{%
  \IfNoValueTF{#1}
  {\xrightarrow{#2}\mathrel{\mkern-14mu}\rightarrow}
  {\xrightarrow[#1]{#2}\mathrel{\mkern-14mu}\rightarrow}}
\NewDocumentCommand{\red}{o}{%
  \IfNoValueTF{#1}{\xrightarrowdbl{}}{\xrightarrowdbl{#1}}}
\newcommand{\bredone}{\redone[\beta]}
\NewDocumentCommand{\equal}{o}{%
  \IfNoValueTF{#1}{\eq}{\stackrel{#1}{\eq}}}
\newcommand{\eqs}{\equiv}

\newcommand{\FV}[1]{\mathrm{FV}(#1)}
\NewDocumentCommand{\rep}{m o}{%
  \IfNoValueTF{#2}{\underline{#1}}{{\underline{#1}}_{#2}}}
\newcommand{\aconvone}{\redone[\alpha]}
\newcommand{\aconv}{\red[\alpha]}
\newcommand{\aeq}{\equal[\alpha]}
\newcommand{\bred}{\red[\beta]}
\newcommand{\eredone}{\redone[\eta]}
\newcommand{\bered}{\red[\beta\eta]}
\newcommand{\ext}{\ensuremath{\mathit{ext}}}

\newcommand{\xredone}{\redone[X]}
\newcommand{\xred}{\red[X]}
\newcommand{\beredone}{\redone[\beta\eta]}
\NewDocumentCommand{\redpar}{o}{%
  \IfNoValueTF{#1}{\Longrightarrow}{\mathrel{\stackrel{#1}{\Longrightarrow}}}}
\newcommand{\bredpar}{\redpar[\beta]}
\newcommand{\beredpar}{\redpar[\beta\eta]}
\newcommand{\bcd}[1]{{#1}^{*{\beta}}}
\newcommand{\becd}[1]{{#1}^{*{\beta\eta}}}
\newcommand{\mModel}[1]{\applytofirst{\mathfrak}{#1}}
\newcommand{\Top}[1]{\mathcal{#1}}
\newcommand{\Interior}[1]{\mathrm{Int}(#1)}
\newcommand{\Prop}[2]{{[\!\![} #2 {]\!\!]_{\mModel{#1}}}}
\NewDocumentCommand{\mSat}{t{/} m m o}{%
  \IfBooleanTF{#1}{\IfNoValueTF{#4}{\mModel{#2}\nVdash#3}{\mModel{#2},#4\nVdash#3}}{%
  \IfNoValueTF{#4}{\mModel{#2}\Vdash#3}{\mModel{#2},#4\Vdash#3}}}
\RenewDocumentCommand{\indcase}{s t{!} m m +m}{%
  \renewcommand{\indfrm}{#3}%
  \def\indfrmp{#3}%
  \def\indcomplex{#4}%
  \IfBooleanTF{#1}{$#3$ हे आण्विक सूत्र आहे: }{$#3 \ident #4$: }%
  \IfBooleanTF{#2}{सराव.}{#5}}
\NewDocumentCommand{\OPrf}{o}{\mathsf{Prf}\IfNoValueF{#1}{_{#1}}}
\NewDocumentCommand{\OCon}{o}{\mathsf{Con}\IfNoValueF{#1}{_{#1}}}
\newcommand{\PAx}{\mathrm{Ax}_0}
\newcommand{\PIso}[1]{\mathcal{#1}}
\newcommand{\Part}[2]{\Atom{\Obj P}{#1,#2}}
\newcommand{\gn}[1]{\ulcorner #1\urcorner}

\usetikzlibrary{arrows,positioning,intersections}
\NewDocumentCommand{\Mod}{o d() m}{%
  \IfNoValueTF{#2}{%
    \IfNoValueTF{#1}{\mathrm{Mod}(#3)}{\mathrm{Mod}^{\Lang{#1}}(#3)}}{%
    \IfNoValueTF{#1}{\mathrm{Mod}_{#2}(#3)}{\mathrm{Mod}_{#2}^{\Lang{#1}}(#3)}}}

\usetikzlibrary{automata}
\usetikzlibrary{lindenmayersystems}
\pgfdeclarelindenmayersystem{Hilbert curve}{
  \rule{L -> +RF-LFL-FR+}
  \rule{R -> -LF+RFR+FL-}}

\newcommand{\TMendtape}{\triangleright}
\newcommand{\TMblank}{0}
\newcommand{\TMstroke}{1}
\newcommand{\TMright}{R}
\newcommand{\TMleft}{L}
\newcommand{\TMstay}{N}
\NewDocumentCommand{\TMtrans}{m m m}{\ensuremath{#1, #2, #3}}
\newcommand{\fn}[1]{\mathrm{#1}}

\newcommand{\Zero}{\fn{zero}}
\newcommand{\Succ}{\fn{succ}}
\newcommand{\Add}{\fn{add}}
\newcommand{\Mult}{\fn{mult}}
\NewDocumentCommand{\Proj}{m m}{P^{#1}_{#2}}
\newcommand{\tsub}{\mathbin{\dot-}}
\newcommand{\defiff}{\Leftrightarrow}
\NewDocumentCommand{\umin}{m m}{\mu #1\;#2}
\NewDocumentCommand{\bmin}{m m}{(\fn{min}\;#1)\,#2}
\NewDocumentCommand{\bexists}{m m}{(\exists #1)\;#2}
\NewDocumentCommand{\bforall}{m m}{(\forall #1)\;#2}
\NewDocumentCommand{\cfind}{m o}{\IfNoValueTF{#2}{\varphi_{#1}}{\varphi_{#1}^{#2}}}
\NewDocumentCommand{\Complement}{m}{\overline{#1}}
\NewDocumentCommand{\fact}{m}{#1\,!}
\newcommand{\Char}[1]{\chi_{#1}}
\NewDocumentCommand{\Prov}{o}{\mathrm{Prov}\IfNoValueF{#1}{_{#1}}}
\NewDocumentCommand{\OProv}{o}{\mathsf{Prov}\IfNoValueF{#1}{_{#1}}}

\NewDocumentCommand{\scode}{m}{\fn{c}_{#1}}
\NewDocumentCommand{\Gn}{m}{{^{\reflectbox{\tiny\#}}}{#1}{^{\mbox{\tiny\#}}}}
\NewDocumentCommand{\Prf}{o}{\mathrm{Prf}\IfNoValueF{#1}{_{#1}}}
\newcommand{\openTuple}{\langle}
\newcommand{\closeTuple}{\rangle}
\DeclareMathOperator\lcm{\mathrm{lcm}}

\DeclareDocumentCommand{\Refut}{o}{\mathrm{Ref}\IfNoValueF{#1}{_{#1}}}
\DeclareDocumentCommand{\ORefut}{o}{\mathsf{Ref}\IfNoValueF{#1}{_{#1}}}
\DeclareDocumentCommand{\ORProv}{o}{\mathsf{RProv}\IfNoValueF{#1}{_{#1}}}

\DeclareDocumentCommand{\TrmSOL}{o}{
  \IfNoValueTF{#1}{\mathrm{Trm}^2}{\mathrm{Trm}^2({\Lang #1})}}
\DeclareDocumentCommand{\FrmSOL}{o}{
  \IfNoValueTF{#1}{\mathrm{Frm}^2}{\mathrm{Frm}^2({\Lang #1})}}

\newcommand{\mClass}[1]{\mathcal{#1}}
\newcommand{\mTrue}[1]{\ensuremath{#1}}
\newcommand{\mFalse}[1]{\ensuremath{\lnot #1}}
\NewDocumentCommand{\Ax}{m}{\ensuremath{\mathrm{#1}}}
\newcommand{\Dual}{\textsc{dual}}
\tikzset{
  modal/.style={>=stealth',shorten >=1pt,shorten <=1pt,auto,
    node distance=1.5cm,label distance=2pt,semithick},
  every label/.style={phantom,align=left},
  world/.style={circle,draw,minimum size=0.5cm,fill=gray!15},
  modal every node/.style={world},
  point/.style={circle,draw,inner sep=0.5mm,fill=black},
  phantom/.style={rectangle,inner sep=0pt,draw=none,fill=none},
  reflexive above/.style={->,loop,looseness=7,in=60,out=120},
  reflexive below/.style={->,loop,looseness=7,in=240,out=300},
  reflexive left/.style={->,loop,looseness=7,in=150,out=210},
  reflexive right/.style={->,loop,looseness=7,in=30,out=330}
}
\newcommand{\ST}{\mathord{\mathrm{ST}}}
\newcommand{\strictif}{\fishhookright}
\DeclareMathSymbol{\boxright}{\mathbin}{symbolsC}{128}
\newcommand{\cif}{\boxright}
\newtheorem{axiom}[defn]{स्वयंसिद्धक}
\newcommand{\stagesord}{\emph{टप्प्यांचा क्रम}}
\newcommand{\stageshier}{\emph{टप्पे मूलभूत}}
\newcommand{\stagesacc}{\emph{टप्प्यांवरील संचय}}
\newcommand{\stagessucc}{\emph{टप्पे सुरूच राहतात}}
\newcommand{\stagesinf}{\emph{अनंत टप्पा}}
\newcommand{\Z}{\Th{Z}}
\newcommand{\Zminus}{\Z^{-}}
\newcommand{\ZF}{\Th{ZF}}
\newcommand{\ZFminus}{\ZF^{-}}
\newcommand{\LT}{\Th{LT}}
\newcommand{\Zr}{\Th{Zr}}
\newcommand{\limofsize}{\emph{आकारमानाच्या मर्यादेचे तत्त्व}}
\newcommand{\stagesinex}{\emph{टप्पे निरपेक्षपणे अनंत आहेत}}
\newcommand{\stagescofin}{\emph{टप्पे संचांपलीकडे जातात}}
\newcommand{\ordplus}{+}
\newcommand{\ordtimes}{\cdot}
\newcommand{\ordexpo}[2]{#1^{(#2)}}
\newcommand{\disjointsum}{\sqcup}
\newcommand{\rlexless}{\mathrel{\sphericalangle}}
\newcommand{\ZFC}{\Th{ZFC}}
\newcommand{\cardfont}[1]{\mathfrak{#1}}
\newcommand{\cardsucc}[1]{#1^{\oplus}}
\newcommand{\canonord}{\lhd}
\newcommand{\cardplus}{\oplus}
\newcommand{\cardtimes}{\otimes}
\newcommand{\cardexpo}[2]{#1^{#2}}
\newcommand{\funfromto}[2]{{}^{#1}{#2}}
\newcommand{\cardnless}[2]{#1 \npreceq #2}
\newcommand{\onesphere}{\mathbf{S}}
\newcommand{\rotationsgroup}{R}
\NewDocumentCommand{\fregenum}{m m}{\# #1\,#2}
\newcommand{\ordtype}[1]{\mathrm{ord}(#1)}
\newcommand{\ordsucc}[1]{#1^{+}}
\newcommand{\ordeq}[2]{#1 \cong #2}
\newcommand{\ordneq}[2]{#1 \ncong #2}
\newcommand{\isomorphic}{\cong}
\DeclareMathOperator*{\supstrict}{\mathrm{lsub}}
\newcommand{\trcl}[1]{\mathrm{trcl}(#1)}
\newcommand{\setrank}[1]{\mathrm{rank}(#1)}
\newcommand{\fregeext}[2]{\epsilon #1\, #2}
\newcommand{\Taut}{\textsc{taut}}
\newcommand{\PL}{\textsc{pl}}
\newcommand{\RK}{\textsc{rk}}
\usepackage{longtable}
\usepackage{cleveref}
\usepackage{xurl}
\tikzset{initial text=सुरुवात}
\renewcommand{\tablename}{तक्ता}
\providecommand{\gitissue}[1]{\href{https://github.com/OpenLogicProject/OpenLogic/issues/#1}{मूळ स्रोताची नोंद #1}}
\crefname{lem}{पूर्वप्रमेय}{पूर्वप्रमेये}
\crefname{table}{तक्ता}{तक्ते}
\Crefname{table}{तक्ता}{तक्ते}
\makeatletter\renewcommand*\l@subsection{\@dottedtocline{2}{3.8em}{4.2em}}\makeatother
\newcommand{\lnand}{\mathbin{\uparrow}}
\newcommand{\lnor}{\mathbin{\downarrow}}
\providecommand{\CutCS}{\ensuremath{\mathrm{Cut}_{\mathrm{CS}}}}
\providecommand{\maeh}[2]{{#1\mathrel{;}#2}}
\providecommand{\pheight}[1]{\fn{ht}(#1)}
\providecommand{\depth}[1]{\fn{dp}(#1)}
\providecommand{\cheight}[1]{\fn{ch}(#1)}
\providecommand{\cutr}[1]{\fn{cr}(#1)}
\providecommand{\cutrank}[1]{\fn{cr}(#1)}
\NewDocumentCommand{\typeof}{m m}{#1^{#2}}
\NewDocumentCommand{\andi}{m m}{\tuple{#1, #2}}
\NewDocumentCommand{\ande}{m m}{\fn{p}_{#1}(#2)}
\NewDocumentCommand{\ori}{m m m}{\fn{in}_{#1}^{#2}(#3)}
\NewDocumentCommand{\ore}{m m m m m}{\fn{case}(#1, #2.#3, #4.#5)}
\NewDocumentCommand{\pair}{m m}{\tuple{#1,#2}}
\NewDocumentCommand{\proj}{m m}{\pi_{#1}(#2)}
\NewDocumentCommand{\dcase}{m m m m m}{\delta\, #1\, #2.#3\, #4.#5}
\NewDocumentCommand{\inj}{o m m}{\IfNoValueTF{#1}{\iota_{#2}(#3)}{\iota_{#2}^{#1}(#3)}}
\NewDocumentCommand{\abort}{m m}{\varepsilon^{#1}(#2)}
\NewDocumentCommand{\maxrank}{m}{\fn{mr}(#1)}
\begin{document}
\begin{titlepage}
\vspace*{2cm}
{\Huge\bfseries मुक्त तर्कशास्त्र\par}
\vspace{1cm}
{\Large ओपन लॉजिक प्रकल्पाची संपूर्ण मराठी आवृत्ती\par}
\vspace{1cm}
मूळ इंग्रजी मजकूर: Open Logic Project.\par
सर्व 722 स्रोत-एककांचे भाषांतर केले आहे. मुख्य अध्यायांबरोबर
पर्यायी आणि मूळ मार्गाबाहेरील विभाग पुढील पूरक भागात दिले आहेत.\par
यंत्रानुवाद, दुरुस्ती आणि तपासणी: OpenAI Codex —
GPT-5.6 Sol आणि GPT-6 Sol, दोन्ही Ultra effort.\par
दुरुस्ती-पुनरावलोकन: GPT-6.1 Sol, Ultra effort.\par
स्रोताशी तुलना, शब्दनिर्णय आणि यांत्रिक तपासण्या केल्या आहेत.
स्वतंत्र मानवी संपादनाचा दावा नाही.
काही तांत्रिक संज्ञा तज्ज्ञ-पुनरावलोकनासाठी खुल्या आहेत.\par
सर्व 722 एककांच्या 6,644 जुळवलेल्या मजकूरखंडांची नोंद आहे.
अधिक तपशीलवार शब्दनिर्णय-पुनरावलोकन 281 एककांपुरते आहे;
हा तपशीलवार पुनरावलोकनाचा दावा संपूर्ण ग्रंथासाठी केलेला नाही.\par
मूळ मजकूर आणि हे रूपांतर: Creative Commons Attribution 4.0.
मूळ घटकांचे स्वतंत्र परवाने लागू राहतात.\par
\url{https://github.com/OpenLogicProject/OpenLogic}\par
\url{https://github.com/KokunoYumeto/OpenLogic-translations}
\end{titlepage}
\chapter*{ओपन लॉजिक प्रकल्पाविषयी}
\addcontentsline{toc}{chapter}{ओपन लॉजिक प्रकल्पाविषयी}

\textit{Open Logic Text} हे औपचारिक अधितर्कशास्त्र आणि
औपचारिक पद्धती यांवरील मुक्त-स्रोत, सहयोगाने लिहिलेले पाठ्यपुस्तक
आहे. ते मध्यम स्तरापासून सुरू होते (म्हणजे औपचारिक तर्कशास्त्राचा
प्राथमिक अभ्यासक्रम पूर्ण झाल्यानंतर). गणितेतर वाचकांना, विशेषतः
तत्त्वज्ञान आणि संगणकशास्त्राच्या विद्यार्थ्यांना, उद्देशून लिहिलेले
असले तरी ते काटेकोर आहे.

सध्या समाविष्ट असलेल्या काही विषयांचा ऊहापोह अजून अपूर्ण
असू शकतो आणि अनेक विभागांत मोठी सुधारणा आवश्यक आहे. भविष्यात
अधिक विषय समाविष्ट करण्याचा आमचा विचार आहे. शब्दसूची, पुढील
वाचनाची यादी, ऐतिहासिक नोंदी, चित्रे, अधिक चांगली स्पष्टीकरणे,
तत्त्वज्ञान, संगणकशास्त्र आणि गणित यांमध्ये निष्कर्षांचे महत्त्व
उलगडणारे विभाग, तसेच अधिक प्रश्न आणि उदाहरणे यांचीही भर
घालण्याचा आमचा विचार आहे. तुम्हाला चूक आढळली किंवा काही सूचना
असली, तर \href{https://github.com/OpenLogicProject/OpenLogic/wiki/Contributing}{कृपया प्रकल्पाच्या चमूला कळवा}.

हा प्रकल्प मुक्त-स्रोताच्या भावनेने चालतो. मजकूर विनामूल्य
उपलब्ध तर आहेच; त्याबरोबर त्याचा LaTeX स्रोत क्रिएटिव्ह कॉमन्स
अट्रिब्यूशन परवान्याखाली उपलब्ध करून दिला आहे. त्यामुळे योग्य
श्रेय दिल्यास कोणालाही आमचे काम डाउनलोड करणे, वापरणे, बदलणे,
त्याची मांडणी बदलणे, दुसऱ्या स्वरूपात रूपांतर करणे आणि पुन्हा
वितरित करणे यांचा अधिकार आहे. अधिक माहितीसाठी ओपन लॉजिक
प्रकल्पाचे \href{http://openlogicproject.org/}{openlogicproject.org} संकेतस्थळ पहा.

\tableofcontents
\part{अनौपचारिक संचसिद्धान्त}\label{sfr:part}\label{sfr:::part}
\begin{editorial}
  या भागातील सामग्री ही मूलभूत अनौपचारिक संचसिद्धान्ताची ओळख
  आहे. टिम बटन यांच्या Open Set Theory चा समावेश केल्यामुळे
  संख्याप्रणालींची रचना आणि अनंततेवरील चर्चाही येथे येते;
  ओपन लॉजिक प्रकल्पाच्या तार्किक भागांसाठी त्या आवश्यक नाहीत.
\end{editorial}
\chapter{संच}








\section{विस्तारात्मकता}\label{sfr:set:bas:sec}

वस्तूंचा समूह एकच वस्तू {म्हणून} विचारात घेतला, की त्याला \emph{संच} म्हणतात.
संच बनवणाऱ्या वस्तूंना त्या संचाचे \emph{घटक} किंवा \emph{सदस्य} म्हणतात.
$x$ हा $A$ या संचाचा घटक असेल, तर आपण $x \in A$ असे लिहितो;
नसेल, तर $x \notin A$ असे लिहितो. ज्या संचात एकही घटक नसतो,
त्याला \emph{रिक्त} संच म्हणतात आणि तो ``$\emptyset$'' या चिन्हाने दर्शवतात.

\begin{explain}
आपण संच \emph{कसा निर्दिष्ट करतो}, त्यातील घटक \emph{कोणत्या क्रमाने}
मांडतो किंवा तेच घटक \emph{किती वेळा} मोजतो, याला महत्त्व नसते.
त्यात कोणते घटक आहेत, एवढेच महत्त्वाचे असते. ही गोष्ट आपण पुढील
तत्त्वात नेमकेपणाने मांडतो.
\end{explain}

\begin{defn}[विस्तारात्मकता]
  $A$ आणि $B$ हे संच असतील, तर $A = B$ तेव्हा आणि केवळ तेव्हाच, जेव्हा
  $A$ मधील प्रत्येक घटक हा $B$ चाही घटक असतो, आणि उलटही.
\end{defn}

विस्तारात्मकतेमुळे काही चिन्हपद्धती वापरणे योग्य ठरते. सर्वसाधारणपणे,
$a_{1}$, \dots, $a_{n}$ या वस्तू दिलेल्या असतील, तर
$\{a_{1}, \dots, a_{n}\}$ म्हणजे ज्यातील घटक हे
$a_1, \ldots, a_n$ आहेत \emph{तो विशिष्ट} संच होय. येथे ``\emph{तो विशिष्ट}''
या शब्दांवर भर आहे, कारण विस्तारात्मकतेनुसार असा \emph{एकच} संच असू शकतो.
विस्तारात्मकतेमुळे पुढील समानताही योग्य ठरते:
  \[    \{a, a, b\} = \{a, b\} = \{b,a\}.
  \]
म्हणजेच, संचांचा विचार करताना त्यातील घटक कोणत्या क्रमाने किंवा
किती वेळा निर्दिष्ट केले आहेत, याला महत्त्व नसते.


\begin{ex}
काही वस्तू दिलेल्या असतील, तर त्या एकत्र करून त्यांचा संच बनवता येतो.
उदाहरणार्थ, रिचर्डच्या भावंडांच्या संचात एक व्यक्ती आहे, आणि तो
$S=\{\textrm{Ruth}\}$ असा लिहिता येईल.
$4$ पेक्षा लहान धन पूर्णांकांचा संच $\{1, 2, 3\}$ आहे; पण तो
$\{3, 2, 1\}$ किंवा अगदी $\{1, 2, 1, 2, 3\}$ असाही लिहिता येतो.
विस्तारात्मकतेनुसार हे सर्व एकच संच आहेत. कारण $\{1, 2, 3\}$ मधील
प्रत्येक घटक हा $\{3, 2, 1\}$ चा (आणि $\{1,
2, 1, 2, 3\}$ चाही) घटक आहे, आणि उलटही.
\end{ex}


अनेकदा संचातील सर्व घटक यांना समान असलेला एखादा गुणधर्म सांगून आपण
संच निर्दिष्ट करू. त्यासाठी पुढील संक्षिप्त चिन्हपद्धती वापरू:
$\Setabs{x}{\phi(x)}$. येथे $\phi(x)$ हा असा गुणधर्म दर्शवतो की,
$x$ ची संचातील घटक मध्ये गणना होण्यासाठी त्यात तो गुणधर्म असला पाहिजे.


\begin{ex}
आपल्या उदाहरणातील $S$ हा संच पुढीलप्रमाणेही निर्दिष्ट करता आला असता:
\[S = \Setabs{x}{x \text{ हे रिचर्डचे भावंड आहे}}.
\]
\end{ex}



\begin{ex}
एखादी संख्या तिच्या स्वतःखेरीज इतर सर्व धन विभाजकांच्या बेरजेइतकी असेल तेव्हा
आणि केवळ तेव्हाच तिला \emph{परिपूर्ण} म्हणतात (हे विभाजक म्हणजे त्या संख्येला
निःशेष भाग जाणाऱ्या, पण तिच्याशी समान नसलेल्या संख्या).
उदाहरणार्थ, $6$ ही परिपूर्ण संख्या आहे, कारण तिचे असे विभाजक
$1$, $2$ आणि~$3$ आहेत, आणि $6 = 1 + 2 + 3$.
खरे तर, $10$ पेक्षा लहान धन पूर्णांकांपैकी $6$ हीच एकमेव परिपूर्ण संख्या आहे.
म्हणून विस्तारात्मकतेचा उपयोग करून आपण असे म्हणू शकतो:
\[	\{6\} = \Setabs{x}{x\text{ ही परिपूर्ण संख्या आहे आणि }0 \leq x \leq 10}
\]
उजवीकडील चिन्हलेखनाचे वाचन असे करतो: ``$x$ ही परिपूर्ण संख्या आहे आणि
$0 \leq x \leq 10$ अशा सर्व $x$ चा संच.'' संच कसा निर्दिष्ट केला आहे,
याला संचाचा विचार करताना महत्त्व नसते, ही गोष्ट वरील समानता स्पष्ट करते.
अधिक सर्वसाधारणपणे, $\phi(x)$ असणाऱ्या सर्व $x$ चा संच नेहमी एकच असतो,
याची विस्तारात्मकता खात्री देते.
त्यामुळे $\Setabs{x}{\phi(x)}$ याला $\phi(x)$ असणाऱ्या सर्व $x$ चा
\emph{तो विशिष्ट} संच म्हणणे विस्तारात्मकतेमुळे योग्य ठरते.
\end{ex}


संच समान असल्याचे दाखवण्याची पद्धत विस्तारात्मकतेमुळे मिळते:
$A = B$ हे दाखवण्यासाठी, $x \in A$ असेल तेव्हा $x \in B$ देखील असते,
आणि $y \in B$ असेल तेव्हा $y \in A$ देखील असते, हे दाखवा.

\begin{prob}
जास्तीत जास्त एकच रिक्त संच असतो हे सिद्ध करा. म्हणजे, $A$ आणि $B$ या
संचांत एकही घटक नसेल, तर $A = B$ हे दाखवा.
\end{prob}











\section{उपसंच आणि घातसंच}\label{sfr:set:sub:sec}

\begin{explain}
आपल्याला अनेकदा संचांची तुलना करावी लागेल. तुलना करण्याचा एक सरळ प्रकार
असा आहे: \emph{एका संचातील प्रत्येक वस्तू दुसऱ्या संचातही असते}.
ही स्थिती पुरेशी महत्त्वाची असल्यामुळे तिच्यासाठी आपण नवीन चिन्हपद्धतीची
ओळख करून घेऊ.
\end{explain}

\begin{defn}[उपसंच]
$A$ या संचातील प्रत्येक घटक हा $B$ चाही घटक असेल,
तर $A$ हा $B$ चा \emph{उपसंच} आहे असे म्हणतो, आणि $A \subseteq B$ असे लिहितो.
$A$ हा $B$ चा उपसंच नसेल, तर $A \not\subseteq B$ असे लिहितो.
$A \subseteq B$ पण $A \neq B$ असेल, तर $A \subsetneq B$ असे लिहितो
आणि $A$ हा $B$ चा \emph{उचित उपसंच} आहे असे म्हणतो.
\end{defn}

\begin{ex}
प्रत्येक संच स्वतःचाच उपसंच असतो, आणि $\emptyset$ हा प्रत्येक संचाचा उपसंच असतो.
सम नैसर्गिक संख्यांचा संच हा नैसर्गिक संख्यांच्या संचाचा उपसंच आहे.
तसेच, $\{ a, b \} \subseteq \{ a, b, c \}$.
पण $\{ a, b, e
\}$ हा $\{ a, b, c \}$ चा उपसंच नाही.
\end{ex}

\begin{ex}
$2$ ही संख्या पूर्णांकांच्या संचाचा घटक आहे, तर सम संख्यांचा संच
हा पूर्णांकांच्या संचाचा उपसंच आहे. मात्र, एखादा संच दुसऱ्या संचाचा
घटक आणि उपसंच \emph{दोन्हीही} असू शकतो. उदाहरणार्थ,
$\{0\} \in \{0, \{0\}\}$ आणि $\{0\} \subseteq \{0,
\{0\}\}$ देखील.
\end{ex}

विस्तारात्मकतेमुळे संचांच्या समानतेचा निकष मिळतो: $A = B$ तेव्हा आणि केवळ तेव्हाच,
जेव्हा $A$ मधील प्रत्येक घटक हा $B$ चाही घटक असतो, आणि उलटही.
``उपसंच'' या व्याख्येत $A \subseteq B$ याचा अर्थ या निकषाचा नेमका पहिला
भाग आहे: $A$ मधील प्रत्येक घटक हा $B$ चा घटक असतो.
अर्थात, $A$ आणि $B$ यांची अदलाबदल केली तरी ही व्याख्या लागू होते:
$B \subseteq A$ तेव्हा आणि केवळ तेव्हाच, जेव्हा $B$ मधील प्रत्येक घटक
हा $A$ चा घटक असतो. हाच विस्तारात्मकतेतील ``आणि उलटही'' हा भाग होय.
दुसऱ्या शब्दांत, संच एकमेकांचे उपसंच असतील तेव्हा आणि केवळ तेव्हाच ते समान
असतात, हे विस्तारात्मकतेवरून निष्पन्न होते.

\begin{prop}
$A = B$ तेव्हा आणि केवळ तेव्हाच, जेव्हा $A \subseteq B$ आणि $B \subseteq A$ दोन्ही असतात.
\end{prop}

आणखी काही उपयुक्त चिन्हपद्धतींची ओळख करून घेण्याची ही योग्य वेळ आहे.
$A$ हा $B$ चा उपसंच केव्हा असतो याची व्याख्या करताना आपण
``$A$ मधील प्रत्येक घटक हा \dots'' असे म्हटले, आणि त्या
``$\dots$'' च्या जागी ``$B$ चा घटक असतो'' असे घातले.
परंतु असा \emph{आकार} असणाऱ्या अभिव्यक्ती इतक्या नेहमी वापरल्या जातात की,
त्यांच्यासाठी आकारिक चिन्हपद्धती उपयोगी पडेल.

\begin{defn}\label{sfr:set:sub:forallxina}
$(\forall x \in A)\phi$ हे $\forall x(x \in A \lif
\phi)$ याचे संक्षिप्त रूप आहे. त्याचप्रमाणे, $(\exists x \in A)\phi$ हे
$\exists x(x
\in A \land \phi)$ याचे संक्षिप्त रूप आहे.
\end{defn}

या चिन्हपद्धतीचा उपयोग करून, $A \subseteq B$ तेव्हा आणि केवळ तेव्हाच, जेव्हा
$(\forall
x \in A)x \in B$, असे म्हणता येते.

आता आपण एका विशिष्ट प्रकारच्या संचाचा विचार करू: दिलेल्या संचाच्या सर्व
उपसंचांचा संच.

\begin{defn}[घातसंच]
$A$ या संचाच्या सर्व उपसंचांचा मिळून जो संच बनतो, त्याला
$A$ चा \emph{घातसंच} म्हणतात आणि तो $\Pow{A}$ असा लिहितात.
  \[    \Pow{A} = \Setabs{B}{B \subseteq A}
  \]
\end{defn}

\begin{ex}
$\{ a, b, c \}$ चे सर्व शक्य उपसंच कोणते? ते असे आहेत:
$\emptyset$, $\{a \}$, $\{b\}$, $\{c\}$, $\{a, b\}$, $\{a, c\}$, $\{b,
c\}$, $\{a, b, c\}$. या सर्व उपसंचांचा संच म्हणजे
$\Pow{\{a,b,c\}}$:
\[\Pow{\{ a, b, c \}} = \{\emptyset, \{a \}, \{b\}, \{c\}, \{a, b\},
\{b, c\}, \{a, c\}, \{a, b, c\}\}
\]
\end{ex}

\begin{prob}
$\{a, b, c, d\}$ चे सर्व उपसंच लिहा.
\end{prob}

\begin{prob}
$A$ मध्ये $n$ घटक असतील, तर $\Pow{A}$ मध्ये $2^n$
घटक असतात हे दाखवा.
\end{prob}











\section{काही महत्त्वाचे संच}\label{sfr:set:imp:sec}

\begin{ex}
आपण प्रामुख्याने ज्यातील घटक या गणिती वस्तू आहेत अशा संचांचा
विचार करणार आहोत. असे चार संच इतके महत्त्वाचे आहेत की त्यांना
विशिष्ट नावे दिलेली आहेत:
\begin{multline*}
    \Nat = \{0, 1, 2, 3, \ldots\} \\
    \shoveright{\text{नैसर्गिक संख्यांचा संच}}\\
    \shoveleft{\Int = \{\ldots, -2, -1, 0, 1, 2, \ldots\}} \\
    \shoveright{\text{पूर्णांकांचा संच}}\\
    \shoveleft{\Rat = \Setabs{\nicefrac{m}{n}}{m, n \in \Int\text{ आणि }n \neq 0}}\\
    \shoveright{\text{परिमेय संख्यांचा संच}}\\
    \shoveleft{\Real = (-\infty, \infty)}\\
    \text{वास्तव संख्यांचा संच (सांतत्य)}
\end{multline*}
हे सर्व \emph{अनंत} संच आहेत; म्हणजेच, प्रत्येकात अनंत घटक आहेत.

या संचांतून पुढे जाताना आपण आपल्या संग्रहात \emph{आणखी} संख्या समाविष्ट
करत असतो. खरे तर, $\Nat \subseteq \Int
\subseteq \Rat \subseteq \Real$ हे स्पष्ट असावे: कारण प्रत्येक नैसर्गिक संख्या
पूर्णांक असते; प्रत्येक पूर्णांक परिमेय असतो; आणि प्रत्येक परिमेय संख्या वास्तव असते.
त्याचप्रमाणे, $\Nat \subsetneq \Int \subsetneq
\Rat$ हेही स्पष्ट असावे, कारण $-1$ हा पूर्णांक आहे पण नैसर्गिक संख्या नाही,
आणि $\nicefrac{1}{2}$ ही परिमेय संख्या आहे पण पूर्णांक नाही.
$\Rat \subsetneq \Real$, म्हणजे परिमेय नसलेल्या काही वास्तव संख्या आहेत,
ही गोष्ट तितकी सहज स्पष्ट होत नाही.

कधीकधी आपण धन पूर्णांकांचा संच $\PosInt = \{1,
2, 3, \dots\}$ आणि फक्त पहिल्या दोन नैसर्गिक संख्या असणारा संच
$\Bin = \{0, 1\}$ यांचाही उपयोग करू.
\end{ex}


\begin{ex}[चिन्हमाला]
आणखी एक लक्षवेधी उदाहरण म्हणजे $A$ या वर्णमालेवरील \emph{सांत चिन्हमालांचा}
संच $A^{*}$: $A$ मधील घटकांचा कोणताही सांत अनुक्रम हा $A$ वरील चिन्हमाला
असतो. प्रत्येक $A$ वर्णमालेसाठी, $A$ वरील चिन्हमालांमध्ये आपण
\emph{रिक्त चिन्हमाला $\Lambda$} हिचाही समावेश करतो. उदाहरणार्थ,
\begin{multline*}
\Bin^*
=\{\Lambda,0,1,00,01,10,11,\\
000,001,010,011,100,101,110,111,0000,\ldots\}.
\end{multline*}
$x=x_{1}\ldots x_{n}\in A^{*}$ ही $A$ मधील $n$ ``अक्षरांनी'' बनलेली
चिन्हमाला असेल, तर तिची \emph{लांबी} $n$ आहे असे म्हणतो आणि $\len{x}=n$ असे लिहितो.
\end{ex}


\begin{ex}[अनंत अनुक्रम]
कोणत्याही $A$ संचासाठी आपण $A$ मधील घटक यांच्या अनंत अनुक्रमांचा
संच $A^\omega$ याचाही विचार करू शकतो. अनंत अनुक्रम
$a_1a_2a_3a_4\dots$ म्हणजे वस्तूंची एका दिशेने अनंत लांब जाणारी यादी;
यादीतील प्रत्येक वस्तू $A$ चा घटक असते.
\end{ex}











\section{संयोग आणि छेद}\label{sfr:set:uni:sec}

\begin{explain}
\ref{sfr:set:bas:sec} मध्ये आपण गुणधर्मानुसार संचांची व्याख्या करणे,
म्हणजे $\Setabs{x}{\phi(x)}$ या आकाराची व्याख्या करणे, पाहिले.
येथे आपण एखादा $\phi$ गुणधर्म वापरतो; या गुणधर्मात आधीच व्याख्या केलेल्या
संचांचा उल्लेख असू शकतो. उदाहरणार्थ, $A$ आणि $B$ हे संच असतील,
तर $\Setabs{x}{x \in A \lor x \in B}$ या संचात $A$ किंवा $B$ मधील
घटक असणाऱ्या सर्व वस्तू येतात. म्हणजेच, $A$ आणि $B$ मधील
घटक एकत्र करणारा हा संच आहे. \ref{sfr:set:uni:fig:union} मध्ये
दाखवल्याप्रमाणे त्याचे चित्र काढता येते; ठळक केलेला भाग $A$ आणि $B$
या दोन संचांतील घटक एकत्र दर्शवतो.

\begin{figure}
  \olasset{assets/diagrams/union.tikz}
  \caption{दोन संचांचा संयोग $A \cup B$ म्हणजे $A$ मधील आणि $B$ मधील
    घटक एकत्र घेऊन बनणारा संच.}
  \label{sfr:set:uni:fig:union}
\end{figure}

संच एकत्र करणारी ही क्रिया अतिशय उपयुक्त आहे आणि वारंवार वापरली जाते.
म्हणून तिला आपण आकारिक नाव आणि चिन्ह देतो.
\end{explain}

\begin{defn}[संयोग]
$A$ आणि $B$ या दोन संचांचा \emph{संयोग} $A \cup B$ असा लिहितात.
तो $A$, $B$ किंवा दोन्ही संचांतील घटक असणाऱ्या सर्व वस्तूंचा संच आहे.
\[A \cup B = \Setabs{x}{x \in A \lor x \in B}
\]
\end{defn}

\begin{ex}
तेच घटक किती वेळा आले आहेत याला महत्त्व नसल्याने, दोन संचांत
सामाईक घटक असेल, तर त्यांच्या संयोगात तो घटक एकदाच येतो.
उदाहरणार्थ, $\{ a, b, c\} \cup \{ a, 0, 1\} = \{a, b, c, 0, 1\}$.

संच आणि त्याचा एखादा उपसंच यांचा संयोग म्हणजे मोठा संचच:
$\{a,
b, c \} \cup \{a \} = \{a, b, c\}$.

संचाचा रिक्त संचाशी संयोग हा मूळ संचाशी समान असतो:
$\{a,
b, c \} \cup \emptyset = \{a, b, c \}$.
\end{ex}

\begin{prob}
$A \subseteq B$ असेल, तर $A \cup B = B$ हे सिद्ध करा.
\end{prob}

\begin{explain}
संयोगाच्या ``द्वैत'' क्रियेचाही विचार करता येतो. या क्रियेत असे सर्व
घटक घेतले जातात की जे $A$ मधील घटक आणि
$B$ मधीलही घटक आहेत, आणि त्यांचा संच बनवला जातो.
या क्रियेला \emph{छेद} म्हणतात. तिचे चित्र \ref{sfr:set:uni:fig:intersection} प्रमाणे काढता येते.
\begin{figure}
  \olasset{assets/diagrams/intersection.tikz}
  \caption{दोन संचांचा छेद $A \cap B$ म्हणजे त्यांतील सामाईक
    घटक यांचा संच.}
  \label{sfr:set:uni:fig:intersection}
\end{figure}
\end{explain}

\begin{defn}[छेद]
$A$ आणि $B$ या दोन संचांचा \emph{छेद} $A \cap B$ असा लिहितात.
तो $A$ आणि $B$ या दोन्ही संचांतील घटक असणाऱ्या सर्व वस्तूंचा संच आहे.
\[A \cap B = \Setabs{x}{x \in A \land x \in B}
\]
दोन संचांचा छेद रिक्त असेल, तर त्यांना \emph{विभक्त} म्हणतात.
याचा अर्थ, त्यांच्यात सामाईक घटक नसतात.
\end{defn}

\begin{ex}
दोन संचांत सामाईक घटक नसतील, तर त्यांचा छेद रिक्त असतो:
$\{ a, b, c\} \cap \{ 0, 1\} = \emptyset$.

दोन संचांत सामाईक घटक असतील, तर त्यांचा छेद म्हणजे त्या सर्वांचा संच:
$\{a, b, c \} \cap \{a, b, d \} = \{a, b\}$.

संच आणि त्याचा एखादा उपसंच यांचा छेद म्हणजे लहान संचच:
$\{a, b, c\} \cap \{a, b\} = \{a, b\}$.

कोणत्याही संचाचा रिक्त संचाशी छेद रिक्त असतो:
$\{a, b, c \}
\cap \emptyset = \emptyset$.
\end{ex}

\begin{prob}
$A \subseteq B$ असेल, तर $A \cap B = A$ हे काटेकोरपणे सिद्ध करा.
\end{prob}

\begin{explain}
दोनपेक्षा अधिक संचांचा संयोग किंवा छेदही बनवता येतो.
हे सर्वसाधारणपणे हाताळण्याची एक सुटसुटीत पद्धत पुढीलप्रमाणे आहे:
ज्या सर्व संचांचा संयोग (किंवा छेद) करायचा आहे, ते सर्व एका संचात एकत्र केले आहेत असे समजा.
मग या संचाच्या किमान एका घटक मध्ये असणाऱ्या सर्व वस्तूंचा संच
म्हणजे मूळ सर्व संचांचा संयोग अशी व्याख्या करता येते.
तसेच, या संचाच्या प्रत्येक घटक मध्ये असणाऱ्या सर्व वस्तूंचा संच
म्हणजे त्यांचा छेद अशी व्याख्या करता येते.
\end{explain}

\begin{defn}
$A$ हा संचांचा संच असेल, तर $\bigcup A$ म्हणजे $A$ मधील
घटक यांत असणाऱ्या घटक यांचा संच:
\begin{align*}
\bigcup A & = \Setabs{x}{x \text{ ही पुढील संचाच्या एका घटकात असते: } A},
\text{ म्हणजे,}\\
& = \Setabs{x}{\text{असा } B \in A
  \text{ असतो की } x \in B}
\end{align*}
\end{defn}

\begin{defn}
$A$ हा संचांचा संच असेल, तर $\bigcap A$ म्हणजे $A$ मधील सर्व घटकांत
सामाईक असणाऱ्या वस्तूंचा संच:
\begin{align*}
\bigcap A & = \Setabs{x}{x \text{ ही पुढील संचाच्या प्रत्येक घटकात असते: } A},
\text{ म्हणजे,}\\
 & = \Setabs{x}{\text{प्रत्येक } B \in A, x \in B}
\end{align*}
\end{defn}

\begin{ex}
$A = \{ \{ a, b \}, \{ a, d, e \}, \{ a, d \} \}$ असे समजा.
मग $\bigcup A = \{ a, b, d, e \}$ आणि $\bigcap A = \{ a \}$.
\end{ex}
\begin{prob}
	$A$ हा संच असेल आणि $A \in B$ असेल, तर $A \subseteq \bigcup B$ हे दाखवा.
\end{prob}

$A_1$, $A_2$, \dots या संचांच्या अनुक्रमासाठीही हेच करता येईल:
\begin{align*}
\bigcup_i A_i & = \Setabs{x}{x \text{ ही पुढीलपैकी एखाद्या संचात असते: } A_i}\\
\bigcap_i A_i & = \Setabs{x}{x \text{ ही पुढील प्रत्येक संचात असते: } A_i}.
\end{align*}

संचांसाठी \emph{निर्देशांक संच} दिलेला असेल, म्हणजे एखादा $I$ संच असा असेल की,
प्रत्येक $i \in I$ साठी आपण $A_i$ चा विचार करत असू, तर पुढील संक्षिप्त
रूपेही वापरता येतात:
\begin{align*}
	\bigcup_{i \in I} A_i & = \bigcup \Setabs{A_i }{i \in I}\\
	\bigcap_{i \in I} A_i & = \bigcap\Setabs{A_i}{i \in I}
\end{align*}

शेवटी, $A$ मध्ये असणारे पण $B$ मध्ये नसणारे सर्व घटक यांच्या
संचाचा विचार करायचा असू शकतो. त्याचे चित्र \ref{sfr:set:uni:difference} प्रमाणे काढता येते.

\begin{figure}
  \olasset{assets/diagrams/difference.tikz}
  \caption{दोन संचांचा फरक $A \setminus B$ म्हणजे $A$ मधील असे
    घटक की जे $B$ मधील घटक नाहीत, यांचा संच.}
  \label{sfr:set:uni:difference}
\end{figure}

\begin{defn}[फरक]
\emph{संचांचा फरक} $A \setminus B$ म्हणजे $A$ मधील असे सर्व
घटक की जे $B$ मधील घटक नाहीत, यांचा संच. म्हणजेच,
\[A\setminus B = \Setabs{x}{x\in A \text{ आणि } x \notin B}.
\]
\end{defn}

\begin{prob}
	$A \subsetneq B$ असेल, तर $B \setminus A \neq \emptyset$ हे सिद्ध करा.
\end{prob}











\section{जोड्या, क्रमित घटकसमूह आणि कार्तीय गुणाकार}\label{sfr:set:pai:sec}

\begin{explain}
संचांतील घटकांना क्रम नसतो, हे विस्तारात्मकतेवरून निष्पन्न होते.
म्हणून क्रम दर्शवायचा असेल, तर आपण \emph{क्रमित जोड्या} $\tuple{x, y}$ वापरतो.
अक्रमित जोडी $\{x, y\}$ मध्ये क्रमाला महत्त्व नसते:
$\{x, y\} = \{y, x\}$. क्रमित जोडीमध्ये मात्र क्रम महत्त्वाचा असतो:
$x
\neq y$ असेल, तर $\tuple{x, y} \neq \tuple{y, x}$.

संच सिद्धांतात क्रमित जोड्यांचा विचार कसा करावा? दोन क्रमित जोड्यांचे
पहिले घटक समान असतील आणि दुसरे घटकही समान असतील तेव्हा आणि केवळ तेव्हाच
त्या जोड्या समान असतात, ही कल्पना जपणे आवश्यक आहे. म्हणजेच:
\[  \tuple{a, b}= \tuple{c, d}\text{ तेव्हा आणि केवळ तेव्हाच, जेव्हा }a = c \text{ आणि }b=d.
\]
वीनर--कुराटोव्स्की यांच्या व्याख्येने संच सिद्धांतात क्रमित जोड्यांची व्याख्या करता येते.
\end{explain}

\begin{defn}[क्रमित जोडी]\label{sfr:set:pai:wienerkuratowski}
	$\tuple{a, b} = \{\{a\}, \{a, b\}\}$.
\end{defn}

\begin{prob}
	\ref{sfr:set:pai:wienerkuratowski} वापरून, $\tuple{a,
	b}= \tuple{c, d}$ तेव्हा आणि केवळ तेव्हाच, जेव्हा $a = c$ आणि $b=d$
दोन्ही असतात, हे सिद्ध करा.
\end{prob}

\begin{explain}
क्रमित जोडीची व्याख्या निश्चित केल्यावर तिच्या साहाय्याने आणखी संचांची
व्याख्या करता येते. उदाहरणार्थ, कधीकधी दोनपेक्षा अधिक वस्तूंचे क्रमित
अनुक्रम हवे असतात: \emph{त्रिके} $\tuple{x, y, z}$,
\emph{चतुष्के} $\tuple{x, y, z, u}$ इत्यादी.
त्रिक म्हणजे ज्याचा पहिला घटक स्वतःच क्रमित जोडी आहे अशी विशिष्ट क्रमित
जोडी असे मानता येते: $\tuple{x, y, z}$ म्हणजे $\tuple{\tuple{x, y},z}$.
चतुष्कांसाठीही हेच लागू होते: $\tuple{x,y,z,u}$ म्हणजे
$\tuple{\tuple{\tuple{x,y},z},u}$, आणि असेच पुढे.
सर्वसाधारणपणे आपण \emph{क्रमित $n$-घटकसमूह} $\tuple{x_1, \dots, x_n}$ यांचा उल्लेख करतो.

क्रमित जोड्यांचे, किंवा इतर क्रमित $n$-घटकसमूहांचे, काही संच उपयुक्त ठरतील.
\end{explain}

\begin{defn}[कार्तीय गुणाकार]
$A$ आणि $B$ हे संच दिलेले असतील, तर त्यांचा \emph{कार्तीय गुणाकार}
$A \times B$ याची व्याख्या अशी करतात:
\[  A \times B = \Setabs{\tuple{x, y}}{x \in A \text{ आणि } y \in B}.
\]
\end{defn}

\begin{ex}
$A = \{0, 1\}$ आणि $B = \{1, a, b\}$ असतील, तर त्यांचा गुणाकार असा आहे:
\[A \times B = \{ \tuple{0, 1}, \tuple{0, a}, \tuple{0, b},
    \tuple{1, 1}, \tuple{1, a}, \tuple{1, b} \}.
\]
\end{ex}

\begin{ex}
$A$ हा संच असेल, तर $A$ चा स्वतःशी गुणाकार $A \times A$ हा $A^2$
असाही लिहितात. तो $x, y \in A$ असणाऱ्या \emph{सर्व} $\tuple{x, y}$ जोड्यांचा
संच आहे. सर्व $\tuple{x, y, z}$ त्रिकांचा संच $A^3$ असतो, आणि असेच पुढे.
पुढील पुनरावर्ती व्याख्या देता येते:
\begin{align*}
  A^1 & = A\\
  A^{k+1} & = A^k \times A
\end{align*}
\end{ex}

\begin{prob}
$\{1, 2, 3\}^3$ मधील सर्व घटक लिहा.
\end{prob}

\begin{prop}\label{sfr:set:pai:cardnmprod}
$A$ मध्ये $n$ घटक आणि $B$ मध्ये $m$ घटक असतील,
तर $A
\times B$ मध्ये $n\cdot m$ घटक असतात.
\end{prop}

\begin{proof}
$A$ मधील प्रत्येक घटक $x$ साठी, $\tuple{x, y} \in A \times B$
या आकाराचे $m$ घटक असतात.
$B_x = \Setabs{\tuple{x, y}}{y
  \in B}$ असे मानू. $x_1 \neq x_2$ असेल तेव्हा $\tuple{x_1, y} \neq
\tuple{x_2, y}$ असल्यामुळे $B_{x_1} \cap B_{x_2} = \emptyset$.
पण $A = \{x_1,
\dots, x_n\}$ असेल, तर $A \times B = B_{x_1} \cup \dots \cup B_{x_n}$;
म्हणून त्यात $n\cdot m$ घटक असतात.

हे चित्ररूपाने पाहण्यासाठी $A \times B$ मधील घटक जाळीच्या स्वरूपात मांडा:
\[\begin{array}{rcccc}
  B_{x_1} = & \{\tuple{x_1, y_1} & \tuple{x_1, y_2} & \dots & \tuple{x_1, y_m}\}\\
  B_{x_2} = & \{\tuple{x_2, y_1} & \tuple{x_2, y_2} & \dots & \tuple{x_2, y_m}\}\\
  \vdots & & \vdots\\
  B_{x_n} = & \{\tuple{x_n, y_1} & \tuple{x_n, y_2} & \dots & \tuple{x_n, y_m}\}
\end{array}
\]
सर्व $x_i$ वेगवेगळे आहेत आणि सर्व $y_j$ वेगवेगळे आहेत. त्यामुळे या जाळीतील
कोणत्याही दोन जोड्या समान नाहीत, आणि त्यांची संख्या $n\cdot m$ आहे.
\end{proof}

\begin{prob}
$k$ वरील गणिती विगमनाने असे दाखवा की, प्रत्येक $k \ge 1$ साठी,
$A$ मध्ये $n$ घटक असतील, तर $A^k$ मध्ये $n^k$ घटक असतात.
\end{prob}

\begin{ex}
$A$ हा संच असेल, तर $A$ वरील \emph{शब्द} म्हणजे $A$ मधील
घटक यांचा कोणताही अनुक्रम.
असा अनुक्रम म्हणजे $A$ मधील घटक यांचा $n$-घटकसमूह असे मानता येते.
उदाहरणार्थ, $A = \{a, b, c\}$ असेल, तर ``$bac$'' या अनुक्रमाचा
$\tuple{b, a, c}$ हे त्रिक म्हणून विचार करता येतो.
शब्दांना, म्हणजे चिन्हांच्या अनुक्रमांना, संगणकशास्त्रात अत्यंत महत्त्व आहे.
संकेतानुसार, $A$ मधील घटक यांना आपण $1$ लांबीचे अनुक्रम मानतो,
आणि $\emptyset$ याला $0$ लांबीचा अनुक्रम मानतो.
मग $A$ वरील \emph{सर्व} शब्दांचा संच असा आहे:
\[A^* = \{\emptyset\} \cup A \cup A^2 \cup A^3 \cup \dots
\]
\end{ex}












\section{रसेलची विरोधापत्ती}\label{sfr:set:rus:sec}

विस्तारात्मकतेमुळे $\Setabs{x}{\phi(x)}$ ही चिन्हपद्धती $\phi(x)$ असणाऱ्या
सर्व $x$ च्या \emph{त्या विशिष्ट} संचासाठी वापरणे योग्य ठरते.
मात्र, विस्तारात्मकतेमुळे \emph{प्रत्यक्षात} फक्त पुढील विचार योग्य ठरतो.
ज्याचे सदस्य सर्व आणि केवळ $\phi$ असणाऱ्या वस्तू आहेत असा संच \emph{असेल},
\emph{तर} असा एकच संच असतो. दुसऱ्या शब्दांत: एखादा $\phi$ निश्चित केल्यावर,
$\Setabs{x}{\phi(x)}$ हा संच \emph{अस्तित्वात असल्यास} तो एकमेव असतो.

पण ही अट महत्त्वाची आहे!{} प्रत्येक गुणधर्मापासून \emph{गुणधर्माधारित संचरचना}
करता येतेच असे नाही. म्हणजेच, काही गुणधर्म संचांची व्याख्या करत \emph{नाहीत}.
सर्वच गुणधर्मांनी संचांची व्याख्या होत असती, तर आपल्याला थेट व्याघातांना
सामोरे जावे लागले असते. याचे सर्वांत प्रसिद्ध उदाहरण म्हणजे रसेलची विरोधापत्ती.

संच हे इतर संचांचे घटक असू शकतात; उदाहरणार्थ, $A$ या संचाचा घातसंच
हा संचांनी बनलेला असतो. म्हणून एखादा संच दुसऱ्या संचाचा घटक आहे का,
असा प्रश्न विचारणे किंवा त्याचा शोध घेणे अर्थपूर्ण आहे. संच स्वतःचाच सदस्य
असू शकतो का? संचाच्या कल्पनेत असे काही दिसत नाही की ज्यामुळे ही शक्यता
नाकारली जाईल. उदाहरणार्थ, \emph{सर्व} संच मिळून वस्तूंचा समूह बनत असेल,
तर त्या सर्वांना एका संचात एकत्र करता येईल असे वाटू शकते; तो म्हणजे सर्व
संचांचा संच. तो स्वतः एक संच असल्यामुळे सर्व संचांच्या संचाचा घटक असेल.

स्वतःला घटक म्हणून न सामावणाऱ्या, म्हणजे \emph{स्वतःचे सदस्य नसणाऱ्या}
संचांचा गुणधर्म विचारात घेतल्यावर रसेलची विरोधापत्ती उद्भवते.
स्वतःला घटक म्हणून न सामावणाऱ्या सर्व संचांचा एक संच आहे असे आपण
मानले, तर काय होईल? असा
\[R = \Setabs{x}{x \notin x}
\]
संच अस्तित्वात आहे का? तो अस्तित्वात नसतो हे आपण सिद्ध करू शकतो.

\begin{thm}[रसेलची विरोधापत्ती]\label{sfr:set:rus:thm:russells-paradox}
	$R = \Setabs{x}{x \notin x}$ असा कोणताही संच नाही.
\end{thm}

\begin{proof}
$R = \Setabs{x}{x \notin x}$ अस्तित्वात असेल, तर
$R \in R$ तेव्हा आणि केवळ तेव्हाच, जेव्हा $R \notin R$; हा व्याघात आहे.
\end{proof}


\begin{explain}
ही सिद्धता अधिक सावकाश पाहू. $R$ अस्तित्वात असेल, तर $R \in
R$ आहे की नाही हा प्रश्न अर्थपूर्ण आहे. खरोखर $R \in R$ आहे असे समजा.
आता, स्वतःचे घटक नसणाऱ्या सर्व संचांचा संच अशी $R$ ची व्याख्या केली होती.
त्यामुळे $R \in R$ असेल, तर $R$ मध्ये स्वतःच $R$ चा व्याख्यक गुणधर्म नसतो.
पण हा गुणधर्म असणारेच संच $R$ मध्ये असतात. म्हणून $R$ हा $R$ चा
घटक असू शकत नाही; म्हणजेच $R \notin R$.
परंतु $R$ हा $R$ चा घटक आहे आणि नाही, असे दोन्ही होऊ शकत नाही;
म्हणून व्याघात मिळतो.

$R \in R$ या गृहीतकामुळे व्याघात येत असल्याने $R \notin R$ मिळते.
पण यामुळेही व्याघात येतो!{} कारण $R
\notin R$ असेल, तर $R$ मध्ये स्वतःच $R$ चा व्याख्यक गुणधर्म असतो;
म्हणून स्वतःचे सदस्य नसणाऱ्या इतर सर्व संचांप्रमाणे $R$ हा $R$ चा
घटक असेल. आणि पुन्हा, $R$ चा घटक नसणे आणि असणे
या दोन्ही गोष्टी एकत्र शक्य नाहीत.
\end{explain}


\begin{digress}
रसेलची विरोधापत्ती टाळणारा, म्हणजे $R = \Setabs{x}{x \notin x}$
अस्तित्वात आहे हा \emph{विसंगत} दावा न करणारा संच सिद्धांत कसा मांडायचा?
त्यासाठी संच केव्हा अस्तित्वात असतात (आणि केव्हा नसतात) हे सांगणाऱ्या
अगदी नेमक्या अटी देणारी स्वयंसिद्धके मांडावी लागतील.

या प्रकरणात आराखडा दिलेला संच सिद्धांत असे करत नाही. तो
\emph{खरोखरच अनौपचारिक} आहे. त्यातून एवढेच सांगितले जाते की संच
विस्तारात्मकतेचे पालन करतात आणि काही संच दिलेले असतील, तर त्यांचा संयोग,
छेद इत्यादी बनवता येतात. संच सिद्धांत यापेक्षा अधिक काटेकोरपणे विकसित करणे
शक्य आहे.
\end{digress}




\chapter{संबंध}\label{sfr:rel::chap}








\section{संच म्हणून संबंध}\label{sfr:rel:set:sec}

\begin{explain}
\ref{sfr:set:imp:sec} मध्ये आपण काही महत्त्वाच्या संचांचा उल्लेख केला:
$\Nat$, $\Int$, $\Rat$, $\Real$. यांपैकी काही संचांतील घटक
यांच्यातील काही लक्षवेधी संबंध तुम्हाला आठवत असतील. उदाहरणार्थ, या प्रत्येक
संचावर एक सर्वमान्य \emph{क्रमसंबंध} असतो. तसेच, प्रत्येक वस्तूचा फक्त
स्वतःशी आणि इतर कशाशीही नसणारा \emph{एकरूप असण्याचा} संबंध असतो. पुढे
आपल्याला आणखी अनेक लक्षवेधी संबंध भेटतील, आणि शक्य संबंध तर त्याहूनही
अधिक आहेत. त्यांचा आढावा घेण्यापूर्वी मात्र, संबंधांकडे एका विशिष्ट प्रकारचे
संच म्हणून पाहता येते, हे आपण दाखवू.

यासाठी \ref{sfr:set:pai:sec} मधील दोन गोष्टी आठवा. प्रथम,
\emph{क्रमित जोडी} ही संकल्पना: $a$ आणि $b$ दिलेले असतील, तर
$\tuple{a, b}$ बनवता येते. येथे घटकांचा क्रम महत्त्वाचा \emph{असतो}.
म्हणून $a \neq b$ असेल, तर $\tuple{a, b} \neq \tuple{b, a}$.
(याची अक्रमित जोड्यांशी, म्हणजे $2$-घटक संचांशी, तुलना करा: तेथे
$\{a, b\}=\{b, a\}$.) दुसरे म्हणजे, \emph{कार्तीय गुणाकार} ही संकल्पना
आठवा: $A$ आणि $B$ हे संच असतील, तर $A \times B$ बनवता येतो. तो
$x \in A$ आणि $y \in B$ असणाऱ्या सर्व $\tuple{x, y}$ जोड्यांचा संच
आहे. विशेषतः, $A^{2}= A \times A$ हा $A$ मधून घेतलेल्या सर्व क्रमित
जोड्यांचा संच आहे.

आता एका संचावरील विशिष्ट संबंध पाहू: नैसर्गिक संख्यांच्या $\Nat$ या
संचावरील $<$-संबंध. $n<m$ असणाऱ्या संख्यांच्या सर्व $\tuple{n, m}$
जोड्यांचा संच विचारात घ्या, म्हणजे,
\[  R=\Setabs{\tuple{n, m}}{n, m \in \Nat \text{ आणि } n<m}.
\]
$n$ ही संख्या $m$ पेक्षा लहान असणे आणि $\tuple{n, m}$ ही जोडी
$R$ ची सदस्य असणे यांचा निकटचा संबंध आहे, तो असा:
\[      n<m\text{ तेव्हा आणि केवळ तेव्हाच }\tuple{n, m} \in R.
\]
खरे तर, कोणतीही माहिती न गमावता $R$ हा संच \emph{म्हणजेच} $\Nat$
वरील $<$-संबंध आहे, असे आपण मानू शकतो.

त्याचप्रमाणे, संख्यांमधील कोणत्याही संबंधासाठी $\Nat^{2}$ चा एक उपसंच
तयार करता येतो. उलट, संख्यांच्या जोड्यांचा कोणताही $S \subseteq \Nat^{2}$
संच दिला असेल, तर त्याला अनुरूप संख्यांमधील एक संबंध असतो: $n$ चा $m$
शी तो संबंध तेव्हा आणि केवळ तेव्हाच असतो, जेव्हा $\tuple{n, m} \in S$.
यावरून पुढील व्याख्या समर्थित होते:
\end{explain}

\begin{defn}[द्विपदी संबंध]
$A$ या संचावरील \emph{द्विपदी संबंध} म्हणजे $A^{2}$ चा एक उपसंच.
$R \subseteq A^{2}$ हा $A$ वरील द्विपदी संबंध असेल आणि $x, y \in A$
असतील, तर $\tuple{x, y} \in R$ यासाठी आपण कधी कधी $Rxy$ (किंवा $xRy$)
असे लिहितो.
\end{defn}

\begin{ex}
  \label{sfr:rel:set:relations}
नैसर्गिक संख्यांच्या जोड्यांचा $\Nat^{2}$ हा संच अशा द्विमितीय सारणीत
मांडता येतो:
\[  \begin{array}{ccccc}
  \mathbf{\tuple{ 0,0 }} & \tuple{ 0,1 } &
    \tuple{ 0,2 } & \tuple{ 0,3 } & \ldots\\
  \tuple{ 1,0 } & \mathbf{\tuple{ 1,1 }} &
    \tuple{ 1,2 } & \tuple{ 1,3 } & \ldots\\
  \tuple{ 2,0 } & \tuple{ 2,1 } &
    \mathbf{\tuple{ 2,2 }} & \tuple{ 2,3 } & \ldots\\
  \tuple{ 3,0 } & \tuple{ 3,1 } & \tuple{ 3,2 } &
    \mathbf{\tuple{ 3,3 }} & \ldots\\
  \vdots & \vdots & \vdots & \vdots & \mathbf{\ddots}
  \end{array}
\]
येथे कर्णावरील नोंदी ठळक केल्या आहेत, कारण त्या जोड्यांचा बनलेला
$\Nat^2$ चा उपसंच, म्हणजे,
\[  \{\tuple{0,0 }, \tuple{ 1,1 }, \tuple{ 2,2 }, \dots\},
  \]
हा \emph{$\Nat$ वरील एकरूपता संबंध} आहे. (एकरूपता संबंध वारंवार वापरला
जातो, म्हणून कोणत्याही $A$ संचासाठी $\Id{A}=\Setabs{\tuple{ x,x }}{x \in
A}$ अशी व्याख्या करू.) कर्णाच्या वर असलेल्या सर्व जोड्यांचा उपसंच, म्हणजे,
\[  L = \{\tuple{ 0,1 },\tuple{ 0,2 },\ldots,\tuple{ 1,2 },
  \tuple{ 1,3 }, \dots, \tuple{ 2,3 },\tuple{ 2,4 },\ldots\},
\]
हा \emph{लहान असण्याचा} संबंध आहे, म्हणजे $Lnm$ तेव्हा आणि केवळ तेव्हाच,
जेव्हा $n<m$. कर्णाच्या खाली असलेल्या जोड्यांचा उपसंच, म्हणजे,
\[  G=\{ \tuple{ 1,0 },\tuple{ 2,0 },\tuple{
    2,1 }, \tuple{ 3,0 },\tuple{ 3,1 },\tuple{ 3,2 }, \dots\},
\]
हा \emph{मोठे असण्याचा} संबंध आहे, म्हणजे $Gnm$ तेव्हा आणि केवळ तेव्हाच,
जेव्हा $n>m$. $L$ आणि $I$ यांचा संयोग $K=L\cup I$ असा लिहू; तो
\emph{लहान किंवा समान असण्याचा} संबंध आहे: $Knm$ तेव्हा आणि केवळ तेव्हाच,
जेव्हा $n \le m$. त्याचप्रमाणे, $H=G \cup I$ हा \emph{मोठे किंवा समान
असण्याचा संबंध} आहे. $L$, $G$, $K$ आणि $H$ हे \emph{क्रम} नावाच्या
विशिष्ट प्रकारचे संबंध आहेत. $L$ आणि $G$ यांचा गुणधर्म असा: कोणत्याही संख्येचा स्वतःशी $L$ किंवा $G$
संबंध नसतो (म्हणजे, सर्व $n$ साठी $Lnn$ आणि $Gnn$ दोन्ही असत्य असतात).
हा गुणधर्म असणाऱ्या संबंधांना \emph{अपरावर्ती} म्हणतात; आणि ते क्रमही
असतील, तर त्यांना \emph{काटेकोर क्रम} म्हणतात.
\end{ex}

\begin{explain}
क्रम आणि एकरूपता हे महत्त्वाचे व स्वाभाविक संबंध असले, तरी आपल्या
व्याख्येनुसार $A^{2}$ चा \emph{कोणताही} उपसंच हा $A$ वरील संबंध असतो,
तो कितीही कृत्रिम किंवा ओढूनताणून बनवलेला वाटला तरीही. विशेषतः,
$\emptyset$ हा कोणत्याही संचावरील संबंध असतो (हा \emph{रिक्त संबंध};
कोणत्याही घटकजोडीमध्ये तो नसतो), आणि $A^{2}$ हाही $A$ वरील संबंध असतो
(प्रत्येक जोडीमध्ये असणारा); त्याला \emph{सार्वत्रिक संबंध} म्हणतात.
तसेच, $E=\Setabs{\tuple{n, m}}{n>5 \text{ किंवा } m \times n \ge 34}$
यासारखी गोष्टही संबंध ठरते.
\end{explain}

\begin{prob}
  $\Pow{\{a, b, c\}}$ या संचावरील $\subseteq$ संबंधाचे घटक लिहा.
\end{prob}











\section{तात्त्विक चिंतन}\label{sfr:rel:ref:sec}

\ref{sfr:rel:set:sec} मध्ये आपण संबंधांची व्याख्या काही विशिष्ट संच म्हणून
केली. येथे थांबून एक छोटासा तात्त्विक प्रश्न विचारू: अशी व्याख्या नेमके
काय \emph{करते}? काही सत्तामीमांसक एकरूपतेची तथ्ये आपण \emph{शोधून काढली},
असा दावा आपल्याला करायचा आहे का, याबद्दल खूपच शंका आहे. उदाहरणार्थ,
$\Nat$ वरील क्रमसंबंध हा \ref{sfr:rel:set:sec} मध्ये व्याख्यिलेला
$R=  \Setabs{\tuple{n,m}}{n, m \in \Nat\text{ आणि }
n < m}$ संचच \emph{निघाला}, असे म्हणायचे का? अशी शंका घेण्याची तीन
कारणे पुढीलप्रमाणे आहेत.

पहिले: \ref{sfr:set:pai:wienerkuratowski} मध्ये आपण
$\tuple{a, b} = \{\{a\}, \{a, b\}\}$ अशी व्याख्या केली. त्याऐवजी
$\lVert a,b\rVert = \{\{b\}, \{a, b\}\} = \tuple{b,a}$ ही व्याख्या
विचारात घ्या. $a \neq b$ असेल, तर $\tuple{a, b} \neq \lVert a,b\rVert$.
पण क्रमित जोडीची व्याख्या $\tuple{a,b}$ ऐवजी $\lVert a,b\rVert$ अशीही
तितक्याच योग्य रीतीने करता आली असती. दोन्ही व्याख्या तितक्याच उपयोगी
ठरल्या असत्या. त्यामुळे नैसर्गिक संख्यांवरील क्रमसंबंध ``असण्यासाठी''
आता दोन तितकेच योग्य उमेदवार आहेत:
\begin{align*}
		R &= \Setabs{\tuple{n,m}}{n, m \in \Nat \text{ आणि }n < m}\\
		S &= \Setabs{\lVert n,m\rVert}{n, m \in \Nat \text{ आणि }n < m}.
\end{align*}
विस्तारात्मकतेनुसार $R \neq S$ असल्याने ते \emph{दोन्ही} $\Nat$ वरील
क्रमसंबंधाशी एकरूप असू शकत नाहीत. पण प्रत्यक्षात $S$ ऐवजी $R$ (किंवा
उलट) \emph{हाच} क्रमसंबंध आहे, असा दावा करणे मनमानीचे आणि म्हणून काहीसे
अडचणीचे ठरेल. (संचसैद्धान्तिक न्यूनीकरणवादाविरुद्ध बेनासेराफ यांनी
\citeyear{Benacerraf1965} मध्ये प्रसिद्ध केलेल्या युक्तिवादाचे हे अतिशय
सोपे उदाहरण आहे. आपण त्याकडे आणखी काही वेळा परत येऊ.)

दुसरे: \emph{प्रत्येक} संबंध एखाद्या संचाशी एकरूप मानला पाहिजे असे आपल्याला
वाटत असेल, तर संचसदस्यत्वाचा $\in$ हा संबंधसुद्धा एक विशिष्ट संच असला
पाहिजे. तो $\Setabs{\tuple{x,y}}{x \in y}$ हा संच असावा लागेल. पण हा
संच अस्तित्वात आहे का? रसेलची विरोधापत्ती लक्षात घेता, असा संच अस्तित्वात
आहे हा दावा सहज स्वीकारता येत नाही. खरे तर,
संच सिद्धांत
काटेकोरपणे स्वयंसिद्धकांवर आधारित सिद्धांत म्हणून विकसित करता येतो,
आणि तो सिद्धांत या संचाचे अस्तित्व नाकारतो.
म्हणून काही संबंधांना संच म्हणून वागवता येत असले, तरी संचसदस्यत्वाचा
संबंध अपवादात्मक प्रकरण म्हणून घ्यावा लागेल.

तिसरे: संबंध संचांशी ``एकरूप'' मानताना, $\tuple{x,y} \in R$ यासाठी
$Rxy$ असे लिहिण्याची मुभा आपण घेतली. सदस्यत्वाचा ``$\in$'' हा संबंध
विधेय \emph{म्हणून} वापरला, तर हे योग्य आहे. पण ``$\in$'' हे एका
विशिष्ट प्रकारच्या संचाचे नाव आहे असे मानले, तर ``$\tuple{x,y} \in R$''
या व्यंजकात संचांचा निर्देश करणारी फक्त तीन एकवाची पदे राहतात:
``$\tuple{x,y}$'', ``$\in$'' आणि ``$R$''. अशा नावांच्या यादीतून विधान
व्यक्त होऊ शकत नाही; ``तो कप लेखणीधारक ते टेबल'' या निरर्थक शब्दमालेइतकीच
ती निष्फळ आहे. म्हणून पुन्हा, काही संबंधांना संच म्हणून वागवता येत असले,
तरी संचसदस्यत्वाचा संबंध अपवादात्मक प्रकरण असला पाहिजे. (यात फ्रेगे यांच्या
\emph{घोडा} संकल्पनेच्या विरोधापत्तीचे एक सोपे रूप आणि विट्गेन्स्टाइन यांनी
रसेल यांच्याविरुद्ध मांडलेला एक प्रसिद्ध आक्षेप एकत्र आले आहेत.)

मग यातून काय निष्पन्न होते? संख्यांवरील संबंध हे संच आहेत असे म्हणण्यात
काही \emph{चूक} नाही. ते म्हणण्यामागचा हेतू मात्र नीट समजून घ्यावा लागेल.
आपण सत्तामीमांसक एकरूपतेचे तथ्य सांगत नाही. फक्त एवढेच नोंदवत आहोत की,
काही संदर्भांत काही संबंधांना काही विशिष्ट संच \emph{म्हणून वागवता}
येते (आणि आपण तसे वागवणार आहोत).











\section{संबंधांचे विशेष गुणधर्म}\label{sfr:rel:prp:sec}

\begin{intro}
काही प्रकारचे संबंध इतके वारंवार आढळतात की त्यांना विशेष नावे दिली
आहेत. उदाहरणार्थ, $\le$ आणि $\subseteq$ हे आपापल्या प्रांतांतील घटकांना
सारख्याच प्रकारे जोडतात (उदा., $\le$ साठी $\Nat$ आणि $\subseteq$ साठी
$\Pow{A}$). त्यांच्यातील साम्य आणि भेद नेमके समजण्यासाठी संबंधांचे काही
विशेष गुणधर्मांनुसार वर्गीकरण करू. यांपैकी काही गुणधर्मांची संयोजने विशेष
महत्त्वाची ठरतात: क्रम आणि तुल्यता संबंध.
\end{intro}

\begin{defn}[परावर्तिता]
$R \subseteq A^2$ हा संबंध \emph{परावर्ती} तेव्हा आणि केवळ तेव्हाच असतो,
जेव्हा प्रत्येक $x \in A$ साठी $Rxx$ असते.
\end{defn}

\begin{defn}[संक्रमकता]
$R \subseteq A^2$ हा संबंध \emph{संक्रमक} तेव्हा आणि केवळ तेव्हाच असतो,
जेव्हा $Rxy$ आणि $Ryz$ असले, की $Rxz$ सुद्धा असते.
\end{defn}

\begin{defn}[सममिती]
$R \subseteq A^2$ हा संबंध \emph{सममित} तेव्हा आणि केवळ तेव्हाच असतो,
जेव्हा $Rxy$ असले, की $Ryx$ सुद्धा असते.
\end{defn}

\begin{defn}[प्रतिसममिती]
$R \subseteq A^2$ हा संबंध \emph{प्रति\-सम\-मित} तेव्हा आणि केवळ तेव्हाच
असतो, जेव्हा $Rxy$ आणि $Ryx$ दोन्ही असले, की $x=y$ असते (दुसऱ्या शब्दांत:
$x\neq y$ असेल, तर $\lnot Rxy$ किंवा $\lnot Ryx$ असते).
\end{defn}

\begin{explain}
सममित संबंधात $Rxy$ आणि $Ryx$ नेहमी एकत्र सत्य असतात किंवा दोन्ही असत्य
असतात. प्रतिसममित संबंधात $Rxy$ आणि $Ryx$ दोन्ही सत्य असण्यासाठी $x = y$
असणे आवश्यक असते. पण $x = y$ असताना $Rxy$ आणि $Ryx$ सत्य \emph{असलेच
पाहिजेत}, अशी अट यात नाही; फक्त त्यांना मनाई नाही. म्हणून प्रतिसममित
संबंध परावर्ती असू शकतो, पण प्रत्येक प्रतिसममित संबंध परावर्ती असतोच
असे नाही. तसेच, प्रतिसममित असणे आणि केवळ सममित नसणे या वेगळ्या अटी आहेत.
खरे तर, एकच संबंध एकाच वेळी सममित आणि प्रतिसममितही असू शकतो
(उदा., एकरूपता संबंध).
\end{explain}

\begin{defn}[संयोजितता]
$R \subseteq A^2$ या संबंधाला \emph{संयोजित} म्हणतात, जर सर्व $x,y\in A$
साठी $x \neq y$ असल्यास $Rxy$ किंवा $Ryx$ असते.
\end{defn}

\begin{prob}
पुढील गुणधर्म असणाऱ्या संबंधांची उदाहरणे द्या: (a) परावर्ती आणि सममित,
पण संक्रमक नाही; (b) परावर्ती आणि प्रतिसममित; (c) प्रतिसममित व संक्रमक,
पण परावर्ती नाही; आणि (d) परावर्ती, सममित व संक्रमक. संख्या किंवा संच
यांवरील संबंध वापरू नका.
\end{prob}

\begin{defn}[अपरावर्तिता]
$R \subseteq A^2$ या संबंधाला \emph{अपरावर्ती} म्हणतात, जर सर्व $x \in A$
साठी $Rxx$ असत्य असते.
\end{defn}

\begin{defn}[असममिती]
$R \subseteq A^2$ या संबंधाला \emph{असममित} म्हणतात, जर $x,y\in A$ अशा
कोणत्याही जोडीसाठी $Rxy$ आणि $Ryx$ दोन्ही सत्य नसतात.
\end{defn}

लक्षात घ्या: $A \neq \emptyset$ असेल, तर $A$ वरील कोणताही अपरावर्ती
संबंध परावर्ती नसतो, आणि $A$ वरील प्रत्येक असममित संबंध प्रतिसममितही
असतो. तरीही, काही $R \subseteq A^2$ संबंध परावर्तीही नसतात आणि
अपरावर्तीही नसतात; तसेच काही प्रतिसममित संबंध असममित नसतात.












\section{तुल्यता संबंध}\label{sfr:rel:eqv:sec}

एखाद्या संचावरील एकरूपता संबंध परावर्ती, सममित आणि संक्रमक असतो.
हे तिन्ही गुणधर्म असणारे $R$ संबंध वारंवार आढळतात.

\begin{defn}[तुल्यता संबंध]
परावर्ती, सममित आणि संक्रमक असणाऱ्या $R \subseteq A^2$ या संबंधाला
\emph{तुल्यता संबंध} म्हणतात. $A$ मधील घटक $x$ आणि $y$ यांना
\emph{$R$-तुल्य} म्हणतात, जर $Rxy$ असते.
\end{defn}

तुल्यता संबंधांमुळे \emph{तुल्यतावर्ग} ही संकल्पना मिळते. तुल्यता संबंध
प्रांताचे वेगवेगळ्या भागांत ``तुकडे'' करतो. प्रत्येक भागात सर्व वस्तू
एकमेकींशी संबंधित असतात; भिन्न भागांतील कोणत्याही वस्तू एकमेकींशी संबंधित
नसतात. काही वेळा या भागांबद्दल \emph{थेट} बोलणे उपयोगी ठरते. त्यासाठी
पुढील व्याख्या करू:

\begin{defn}\label{sfr:rel:eqv:def:equivalenceclass}
$R \subseteq A^2$ हा तुल्यता संबंध असू द्या. प्रत्येक $x \in A$ साठी,
$A$ मधील $x$ चा \emph{तुल्यतावर्ग} म्हणजे
$\equivrep{x}{R} = \Setabs{y \in A}{Rxy}$ हा संच. $R$ नुसार $A$ चा
\emph{भागाकारसंच} म्हणजे $\equivclass{A}{R} = \Setabs{\equivrep{x}{R}}{x \in A}$,
अर्थात या तुल्यतावर्गांचा संच.
\end{defn}

पुढील निष्कर्ष तुल्यतावर्गाच्या व्याख्येला आधार देतो: तुल्यतावर्ग हे खरोखर
$A$ चे विभाजन करणारे भाग असतात, हे तो सिद्ध करतो.

\begin{prop}
$R \subseteq A^2$ हा तुल्यता संबंध असेल, तर $Rxy$ तेव्हा आणि केवळ
तेव्हाच, जेव्हा $\equivrep{x}{R} = \equivrep{y}{R}$.
\end{prop}

\begin{proof}
डावीकडून उजवीकडे जाणाऱ्या दिशेसाठी, $Rxy$ असे समजा आणि
$z \in \equivrep{x}{R}$ असू द्या. व्याख्येनुसार $Rxz$. $R$ हा तुल्यता
संबंध असल्याने $Ryz$. (सविस्तर सांगायचे तर: $Rxy$ असल्याने आणि $R$
सममित असल्याने $Ryx$; तसेच $Rxz$ असल्याने आणि $R$ संक्रमक असल्याने
$Ryz$.) म्हणून $z \in \equivrep{y}{R}$. हे कोणत्याही अशा घटकासाठी
लागू असल्याने $\equivrep{x}{R} \subseteq \equivrep{y}{R}$. अगदी
त्याचप्रमाणे $\equivrep{y}{R} \subseteq \equivrep{x}{R}$. म्हणून
विस्तारात्मकतेनुसार $\equivrep{x}{R} = \equivrep{y}{R}$.

उजवीकडून डावीकडे जाणाऱ्या दिशेसाठी,
$\equivrep{x}{R} = \equivrep{y}{R}$ असे समजा. $R$ परावर्ती असल्याने
$Ryy$, म्हणून $y \in \equivrep{y}{R}$. मग
$\equivrep{x}{R} = \equivrep{y}{R}$ या गृहीतकामुळे
$y \in \equivrep{x}{R}$ सुद्धा. म्हणून $Rxy$.
\end{proof}

\begin{ex}
तुल्यता संबंधांचे एक चांगले उदाहरण शेषांक अंकगणितातून मिळते. कोणत्याही
$a$, $b$ आणि $n \in \PosInt$ साठी, $a$ ला $n$ ने भागल्यावर येणारी
बाकी आणि $b$ ला $n$ ने भागल्यावर येणारी बाकी सारखी असेल, तेव्हा आणि
केवळ तेव्हाच $a \equiv_n b$ असे म्हणू. (चिन्हांत मांडायचे तर:
$a \equiv_n b$ तेव्हा आणि केवळ तेव्हाच, जेव्हा एखाद्या $k \in \Int$
साठी $a - b = kn$.) मग कोणत्याही $n$ साठी $\equiv_n$ हा तुल्यता
संबंध असतो. $\equiv_n$ मुळे नेमके $n$ भिन्न तुल्यतावर्ग मिळतात; म्हणजे
$\equivclass{\Nat}{\equiv_n}$ मध्ये $n$ घटक असतात. ते असे:
$n$ ने निःशेष भाग जाणाऱ्या संख्यांचा संच, म्हणजे $\equivrep{0}{\equiv_n}$;
$n$ ने भागल्यावर बाकी $1$ येणाऱ्या संख्यांचा संच, म्हणजे
$\equivrep{1}{\equiv_n}$; \ldots; आणि $n$ ने भागल्यावर बाकी $n-1$
येणाऱ्या संख्यांचा संच, म्हणजे $\equivrep{n-1}{\equiv_n}$.
\end{ex}

\begin{prob}
कोणत्याही $n \in \PosInt$ साठी $\equiv_n$ हा तुल्यता संबंध आहे आणि
$\equivclass{\Nat}{\equiv_n}$ मध्ये नेमके $n$ सदस्य आहेत, हे दाखवा.
\end{prob}











\section{क्रम}\label{sfr:rel:ord:sec}

\begin{explain}
अनेक तुलनांमध्ये आपण काही वस्तू इतर वस्तूंपेक्षा एखाद्या दृष्टीने
``लहान'', ``समान'' किंवा ``मोठ्या'' आहेत असे म्हणतो. यात \emph{क्रम}
संबंध येतात. पण क्रमसंबंधांचे वेगवेगळे प्रकार आहेत. उदाहरणार्थ, काहींमध्ये
कोणत्याही दोन वस्तूंची तुलना करता आली पाहिजे, तर इतरांमध्ये तशी अट नसते.
काहींमध्ये एकरूपतेचा समावेश असतो (जसे $\le$), तर काहींमध्ये नसतो
(जसे $<$). येथे या प्रकारांचे वर्गीकरण उपयोगी पडेल.
\end{explain}

\begin{defn}[पूर्वक्रम]
परावर्ती आणि संक्रमक असणाऱ्या संबंधाला \emph{पूर्वक्रम} म्हणतात.
\end{defn}

\begin{defn}[अंशतः क्रम]
प्रतिसममितही असणाऱ्या पूर्वक्रमाला \emph{अंशतः क्रम} म्हणतात.
\end{defn}

\begin{defn}[रेषीय क्रम]\label{sfr:rel:ord:def:linearorder}
संयोजितही असणाऱ्या अंशतः क्रमाला \emph{पूर्ण क्रम} किंवा \emph{रेषीय क्रम}
म्हणतात.
\end{defn}

\begin{ex}
प्रत्येक रेषीय क्रम अंशतः क्रमही असतो आणि प्रत्येक अंशतः क्रम पूर्वक्रमही
असतो; पण यांची उलट विधाने सत्य नाहीत. $A$ वरील सार्वत्रिक संबंध परावर्ती
आणि संक्रमक असल्याने तो पूर्वक्रम असतो. पण $A$ मध्ये एकाहून अधिक
घटक असतील, तर सार्वत्रिक संबंध प्रतिसममित नसतो, म्हणून तो
अंशतः क्रम नसतो.
\end{ex}

\begin{ex}
$\Bin^*$ वरील \emph{अधिक लांब नसण्याचा} $\preccurlyeq$ हा संबंध पाहा:
$x \preccurlyeq y$ तेव्हा आणि केवळ तेव्हाच, जेव्हा $\len{x} \le \len{y}$.
हा पूर्वक्रम (परावर्ती आणि संक्रमक) आहे, शिवाय संयोजितही आहे; पण
प्रतिसममित नसल्याने अंशतः क्रम नाही. उदाहरणार्थ,
$01 \preccurlyeq 10$ आणि $10 \preccurlyeq 01$, पण $01 \neq 10$.
\end{ex}

\begin{ex}
संचांच्या संचावरील $\subseteq$ हा एक महत्त्वाचा अंशतः क्रम आहे.
सर्वसाधारणपणे तो रेषीय क्रम नसतो. कारण $a \neq b$ असेल आणि
$\Pow{\{a, b\}} = \{\emptyset, \{a\}, \{b\}, \{a,b\}\}$ विचारात घेतला,
तर $\{a\} \nsubseteq \{b\}$, $\{a\} \neq \{b\}$ आणि
$\{b\} \nsubseteq \{a\}$ असे दिसते.
\end{ex}

\begin{ex}
\emph{निःशेष विभाज्यतेचा} संबंध हा रेषीय नसलेला अंशतः क्रम देतो.
$n$ आणि $m$ या पूर्णांकांसाठी, $n \mid m$ म्हणजे $n$ ने $m$ ला निःशेष
भाग जातो, अर्थात एखादा पूर्णांक $k$ असा असतो की $m = kn$.
$\Nat$ वर हा अंशतः क्रम आहे, पण रेषीय क्रम नाही: उदाहरणार्थ,
$2 \nmid 3$ आणि $3 \nmid 2$. $\Int$ वरील संबंध म्हणून पाहिल्यास
विभाज्यता केवळ पूर्वक्रम आहे, कारण ती प्रतिसममित नाही:
$1 \mid -1$ आणि $-1 \mid 1$, पण $1 \neq -1$.
\end{ex}

\begin{ex}
अनुक्रमांच्या $A^*$ या संचावरील \emph{विस्तार} संबंध असा आहे:
$s \sqsubseteq s'$ तेव्हा आणि केवळ तेव्हाच, जेव्हा $s = \emptyseq$
(रिक्त अनुक्रम), $s = s'$, किंवा $s = \tuple{s_1, \dots, s_n}$ आणि
$s' = \tuple{s_1, \dots, s_n, s_{n+1}, \dots, s_m}$.
$s \sqsubseteq s'$ असल्यास $s$ हा $s'$ चा \emph{आरंभीचा खंड} आहे असेही
म्हणतो. $A^*$ वरील विस्तार संबंध अंशतः क्रम आहे, पण रेषीय क्रम नाही;
उदा., $a \neq b$ असेल, तर $ab \not\sqsubseteq ba$ आणि
$ba \not\sqsubseteq ab$.
\end{ex}

\begin{defn}[काटेकोर क्रम]
\emph{काटेकोर क्रम} म्हणजे अपरावर्ती, असममित आणि संक्रमक असणारा संबंध.
\end{defn}

\begin{defn}[काटेकोर रेषीय क्रम]\label{sfr:rel:ord:def:strictlinearorder}
संयोजितही असणाऱ्या काटेकोर क्रमाला \emph{काटेकोर पूर्ण क्रम} किंवा
\emph{काटेकोर रेषीय क्रम} म्हणतात.
\end{defn}

\begin{ex}
$\le$ हा काटेकोर रेषीय क्रम $<$ याला अनुरूप रेषीय क्रम आहे.
$\subseteq$ हा काटेकोर क्रम $\subsetneq$ याला अनुरूप अंशतः क्रम आहे.
\end{ex}

$A$ वरील कोणत्याही काटेकोर क्रम $R$ मध्ये कर्ण $\Id{A}$, म्हणजे सर्व
$\tuple{x, x}$ जोड्या, मिळवल्यास अंशतः क्रम तयार होतो. (याला $R$ चे
\emph{परावर्ती संवरण} म्हणतात.) उलट, अंशतः क्रमातून $\Id{A}$ काढून
टाकल्यास काटेकोर क्रम मिळतो. पुढील दोन निष्कर्ष हे नेमके स्पष्ट करतात.

\begin{prop}\label{sfr:rel:ord:prop:stricttopartial}
$R$ हा $A$ वरील काटेकोर क्रम असेल, तर $R^+ = R \cup \Id{A}$ हा अंशतः
क्रम असतो. शिवाय, $R$ हा काटेकोर रेषीय क्रम असेल, तर $R^+$ हा रेषीय
क्रम असतो.
\end{prop}

\begin{proof}
$R$ हा काटेकोर क्रम आहे असे समजा, म्हणजे $R \subseteq A^2$ आणि $R$
अपरावर्ती, असममित व संक्रमक आहे. $R^+ = R \cup \Id{A}$ असू द्या.
$R^+$ परावर्ती, प्रतिसममित आणि संक्रमक आहे हे दाखवायचे आहे.

$R^+$ स्पष्टपणे परावर्ती आहे, कारण सर्व $x \in A$ साठी
$\tuple{x, x} \in \Id{A} \subseteq R^+$.

$R^+$ प्रतिसममित आहे हे दाखवण्यासाठी व्याघात काढू. $R^+xy$ आणि $R^+yx$
असूनही $x \neq y$ आहे असे समजा. $\tuple{x,y} \in R \cup \Id{A}$,
पण $\tuple{x, y} \notin \Id{A}$ असल्याने $\tuple{x, y} \in R$, म्हणजे
$Rxy$ असले पाहिजे. त्याचप्रमाणे $Ryx$. पण $R$ असममित आहे या
गृहीतकाशी हे विसंगत आहे.

संक्रमकता दाखवण्यासाठी $R^+xy$ आणि $R^+yz$ असे समजा.
$\tuple{x, y} \in R$ आणि $\tuple{y,z} \in R$ दोन्ही असतील, तर $R$
संक्रमक असल्याने $\tuple{x, z} \in R$. अन्यथा $\tuple{x, y} \in \Id{A}$,
म्हणजे $x = y$, किंवा $\tuple{y, z} \in \Id{A}$, म्हणजे $y = z$.
पहिल्या प्रकरणात गृहीतकानुसार $R^+yz$ आणि $x = y$, म्हणून $R^+xz$.
दुसऱ्या प्रकरणातही त्याचप्रमाणे. प्रत्येक प्रकरणात $R^+xz$, म्हणून
$R^+$ संक्रमकही आहे.

``शिवाय'' या भागासाठी, $R$ संयोजितही आहे असे समजा. म्हणून सर्व
$x \neq y$ साठी $Rxy$ किंवा $Ryx$, म्हणजे $\tuple{x, y} \in R$ किंवा
$\tuple{y, x} \in R$. $R \subseteq R^+$ असल्याने $R^+$ बाबतही हे सत्य
राहते, म्हणून $R^+$ संयोजितही आहे.
\end{proof}

\begin{prop}\label{sfr:rel:ord:prop:partialtostrict}
$R$ हा $A$ वरील अंशतः क्रम असेल, तर $R^- = R \setminus \Id{A}$ हा
काटेकोर क्रम असतो. शिवाय, $R$ हा रेषीय क्रम असेल, तर $R^-$ हा काटेकोर
रेषीय क्रम असतो.
\end{prop}

\begin{proof}
हे सरावासाठी ठेवले आहे.
\end{proof}

\begin{prob}
\ref{sfr:rel:ord:prop:partialtostrict} याची सिद्धता द्या.
\end{prob}

काटेकोर रेषीय क्रम विस्तारात्मकतेसारखा एक गुणधर्म पाळतात, हे पुढील
सोप्या निष्कर्षातून दिसते:

\begin{prop}\label{sfr:rel:ord:prop:extensionality-strictlinearorders}
$<$ हा $A$ वरील काटेकोर रेषीय क्रम असेल, तर:
\[  (\forall a, b \in A)((\forall x \in A)(x < a \liff x < b) \lif a = b).
\]
\end{prop}

\begin{proof}
$(\forall x \in A)(x < a \liff x < b)$ असे समजा. $a < b$ असेल, तर
$a < a$; पण $<$ अपरावर्ती आहे याच्याशी हे विसंगत आहे. म्हणून $a \nless b$.
अगदी त्याचप्रमाणे $b \nless a$. $<$ संयोजित असल्याने $a = b$.
\end{proof}











\section{आलेख}\label{sfr:rel:grp:sec}

\emph{आलेख} म्हणजे अशी आकृती की जिच्यात ``गाठी'' किंवा ``शिखरे''
(``शिखर'' या शब्दाचे अनेकवचन) म्हटले जाणारे बिंदू कडांनी जोडलेले असतात. विविक्त गणित आणि संगणकशास्त्रात
आलेख सर्वत्र वापरले जातात. विविध जाळ्यांसारख्या ठोस गोष्टींपासून ते
निर्णयांच्या शक्य परिणामांसारख्या अमूर्त संरचनांपर्यंत अनेक संबंध आणि
संरचना दर्शवण्यासाठी व त्यांचे दृश्य रूप देण्यासाठी ते अतिशय उपयोगी
असतात. ग्रंथसाहित्यात आलेखांचे अनेक प्रकार आढळतात. उदाहरणार्थ, कडांना
दिशा आहे की नाही, नावे आहेत की नाहीत, एखाद्या शिखरापासून त्याच शिखरापर्यंत
कड असू शकते की नाही, त्याच शिखरांमध्ये अनेक कडा असू शकतात की नाहीत,
इत्यादी बाबतींत हे प्रकार वेगळे असतात. \emph{दिशित आलेखांचा} संबंधांशी
विशेष संबंध आहे.

\begin{defn}[दिशित आलेख]
\emph{दिशित आलेख} $G = \tuple{V, E}$ म्हणजे \emph{शिखरांचा} $V$ हा संच
आणि \emph{कडांचा} $E \subseteq V^2$ हा संच.
\end{defn}

\begin{explain}
आपल्या व्याख्येनुसार आलेख म्हणजे एक संच आणि त्या संचावरील एक संबंध.
अर्थात, आलेखांबद्दल बोलताना त्यांचे चित्र काढले जावे ही अपेक्षा स्वाभाविक
आहे: $v_1$ आणि $v_2$ ही शिखरे बाणाने तेव्हा आणि केवळ तेव्हाच जोडून
आलेख काढता येतो, जेव्हा $\tuple{v_1, v_2} \in E$.
नुसता संबंध आणि आलेख यांतील फरक एवढाच की आलेखात शिखरांचा संचही स्पष्ट
केलेला असतो, म्हणजे आलेखात एकाकी शिखरे असू शकतात. महत्त्वाचा मुद्दा असा:
$X$ संचावरील प्रत्येक $R$ संबंधाकडे $\tuple{X, R}$ हा दिशित आलेख म्हणून
पाहता येते; आणि उलट, $\tuple{V, E}$ या दिशित आलेखाकडे, $V$ हा संच स्पष्ट
केलेला असताना, $E \subseteq V^2$ हा संबंध म्हणून पाहता येते.
\end{explain}

\begin{ex}
$V = \{1, 2, 3, 4\}$ आणि $E =
\{\tuple{1,1}, \allowbreak \tuple{1, 2}, \allowbreak \tuple{1, 3},
\allowbreak \tuple{2, 3}\}$ असलेला $\tuple{V, E}$ हा आलेख असा दिसतो:
\begin{align*}
& \begin{tikzpicture}[->,node distance=2cm]
  \node[draw,circle] (A) {$1$};
  \node[draw,circle] (B) [right of=A] {$2$};
  \node[draw,circle] (C) [below of=B] {$3$};
  \node[draw,circle] (D) [right of=B] {$4$};
  \draw (A) to [loop above]  (A);
  \draw (A) to  (B);
  \draw (A) to  (C);
  \draw (B) to  (C);
  \end{tikzpicture}
\intertext{हा आलेख $V' = \{1, 2, 3\}$ असणाऱ्या $\tuple{V', E}$ या
  आलेखापेक्षा वेगळा आहे; तो असा दिसतो:}
& \begin{tikzpicture}[->,node distance=2cm]
  \node[draw,circle] (A) {$1$};
  \node[draw,circle] (B) [right of=A] {$2$};
  \node[draw,circle] (C) [below of=B] {$3$};
  \draw (A) to [loop above]  (A);
  \draw (A) to  (B);
  \draw (A) to  (C);
  \draw (B) to  (C);
  \end{tikzpicture}
\end{align*}
\end{ex}

\begin{prob}
  $\{1, 2, 3, 4\}$ या संचावरील लहान-किंवा-समान असण्याचा $\le$ संबंध
  आलेख म्हणून पाहा आणि त्याची आकृती काढा.
\end{prob}











\section{वृक्ष}\label{sfr:rel:tre:sec}

तर्कशास्त्राच्या सर्व शाखांत महत्त्वाची भूमिका बजावणारा अंशतः क्रमाचा
एक विशिष्ट प्रकार म्हणजे \emph{वृक्ष}. तर्कशास्त्राच्या प्राथमिक भागांत
सांत वृक्ष येतात: उदाहरणार्थ, सूत्रे विन्यासवृक्षात विभागून
समजून घेता येतात, आणि अनेक निष्पत्ती पद्धतींतील
निष्पत्ती सुद्धा सांत वृक्षांच्या स्वरूपात असतात.

विधानीय आणि प्रथमस्तरीय तर्कशास्त्राच्या परिपूर्णता प्रमेयांच्या सिद्धतेतच
अनंत वृक्ष येऊ लागतात; पुढे संपूर्ण गणिती तर्कशास्त्रात त्यांचा वापर होतो.

वृक्षाची संचसैद्धान्तिक संकल्पना आणि आलेख सिद्धांतातील वृक्षाची संकल्पना
यांचा निकटचा संबंध आहे. एका सांत वृक्षाचे चित्र असे आहे:

\begin{center}
\begin{tikzpicture}[nodes={draw, circle}, -]
\node{$r$} [grow'=up]
    child { node {$a$}
        child { node {$c$} }
        child { node {$d$} }
        child { node {$e$} }
    }
    child { node {$b$} };
\end{tikzpicture}
\end{center}

सर्वांत खालचे $r$ हे शिखर मूळ आहे. $r$ व्यतिरिक्त प्रत्येक शिखराच्या
लगेच खाली नेमके एक पालक शिखर आहे. $x$ पासून कडांवरून वर जात $y$ पर्यंत
पोहोचता येत असेल, तर $x$ चा $y$ शी असणारा संबंध म्हणजे $x$ हा $y$ चा
\emph{पूर्वज} असणे, असे समजू शकतो.

वृक्षातील पूर्वज संबंध हा काटेकोर अंशतः क्रम असतो. यावरून संचसैद्धान्तिक
व्याख्या सुचते. ती मांडण्यासाठी दोन संकल्पना लागतात. $\le$ ने अंशतः
क्रमित केलेल्या $A$ या संचातील \emph{लघुतम घटक} म्हणजे असा
घटक $x \in A$ की सर्व $y \in A$ साठी $x \le y$ असते.
एखाद्या संचाच्या प्रत्येक अरिक्त उपसंचात लघुतम घटक असेल, तर तो संच
$\le$ ने \emph{सुक्रमित} आहे असे म्हणतात.

\begin{defn}[वृक्ष]
\emph{वृक्ष} म्हणजे $T = \tuple{A, \le}$ अशी जोडी की $A$ हा संच असतो,
$\le$ हा $A$ वरील अंशतः क्रम असतो, त्याचा $r \in A$ हा एकमेव लघुतम
घटक असतो (त्याला \emph{मूळ} म्हणतात), आणि सर्व $x \in A$ साठी
$\Setabs{y}{y \le x}$ हा संच $\le$ ने सुक्रमित असतो.
\end{defn}

\begin{defn}[उत्तरवर्ती]
$T = \tuple{A, \le}$ हा वृक्ष आहे असे समजा.
$x,y \in A$, $x < y$ असतील आणि $x < z < y$ असा कोणताही $z \in A$
नसेल, तर $y$ हा $x$ चा \emph{उत्तरवर्ती} आहे असे म्हणतो.
\end{defn}

$x \in A$ च्या उत्तरवर्तींना त्याची \emph{अपत्ये} असेही म्हणतात.
$y$ हा $x$ चा उत्तरवर्ती असेल, तर $x$ ला $y$ चा \emph{पूर्ववर्ती}
किंवा \emph{पालक} म्हणतो.

\begin{prop}
$\tuple{A,\le}$ हा वृक्ष असेल, तर मूळ वगळता प्रत्येक $x \in A$ ला
जास्तीत जास्त एक पूर्ववर्ती असतो.
\end{prop}

\begin{proof}
  $y_1 < x$, $y_2 < x$ आणि $y_1 \neq y_2$ असे समजा. मग
  $\{y_1, y_2\} \subseteq \Setabs{z}{z<x}$. $\Setabs{z}{z<x}$ हा संच
  $\le$ ने सुक्रमित असल्याने त्याच्या $\{y_1, y_2\}$ या उपसंचाला लघुतम
  घटक असतो. तो स्पष्टपणे $y_1$ किंवा $y_2$ असला पाहिजे. म्हणून
  $y_1 \le y_2$ किंवा $y_2 \le y_1$. आपण $y_1 \neq y_2$ गृहीत धरले,
  म्हणून प्रत्यक्षात $y_1 < y_2$ किंवा $y_2 < y_1$. शिवाय,
  $y_1 < x$ आणि $y_2 < x$ या गृहीतकांमुळे $y_1 < y_2 < x$ किंवा
  $y_2 < y_1 < x$. म्हणून $y_1$ आणि $y_2$ हे दोघेही $x$ चे पूर्ववर्ती
  असू शकत नाहीत.
\end{proof}

\begin{defn}
$T = \tuple{A, \le}$ या वृक्षातील $A$ हा अनंत संच असेल, तर वृक्षाला
\emph{अनंत}, आणि अन्यथा \emph{सांत} म्हणतात. $T$ मधील प्रत्येक
$x \in A$ ला फक्त सांत संख्येने उत्तरवर्ती असतील, तर $T$ ला
\emph{सांत शाखनाचा} वृक्ष म्हणतात.
\end{defn}

\begin{defn}[शाखा]
$T = \tuple{A, \le}$ हा वृक्ष दिला असता, $T$ ची \emph{शाखा} म्हणजे
$T$ मधील अधिकतम शृंखला; म्हणजे $B \subseteq A$ असा संच की कोणत्याही
$x, y \in B$ साठी $x \le y$ किंवा $y \le x$ असते, आणि प्रत्येक
$z \in X \setminus B$ साठी असा $u \in B$ असतो की $z \le u$ आणि
$u \le z$ दोन्ही असत्य असतात.

$T$ च्या सर्व शाखांचा संच दर्शवण्यासाठी आपण $[T]$ वापरतो.
\end{defn}

\begin{ex}
सांत शाखनाच्या वृक्षाचे एक प्रसिद्ध उदाहरण म्हणजे $0$ आणि $1$ यांच्या
सांत अनुक्रमांचा \emph{अनंत द्विमान वृक्ष}. त्याला कधी $\{0,1\}^*$ किंवा
$\Bin^*$ असे लिहितात आणि विस्तार संबंध $\sqsubseteq$ ने क्रमित करतात
(उदा., $101 \sqsubseteq 101101$). कोणत्याही द्विमान चिन्हमालेच्या शेवटी
$0$ किंवा $1$ जोडून तिचा विस्तार नेहमी करता येतो. म्हणून या वृक्षात
अनंत घटक असतात: प्रत्येक $s$ या घटकाला नेमके दोन उत्तरवर्ती, $s0$ आणि
$s1$, असतात. त्याचे मूळ म्हणजे $\emptyseq$ हा रिक्त अनुक्रम.
\end{ex}

\begin{ex}
थोडे अधिक सर्वसाधारण उदाहरण म्हणजे नैसर्गिक संख्यांच्या सांत अनुक्रमांचा
$\Nat^*$ हा संच आणि त्यावरील विस्तार संबंध $\sqsubseteq$; हाही वृक्ष
असतो. तो स्पष्टपणे सांत शाखनाचा नाही: प्रत्येक $s \in \Nat^*$ ला अनंत
उत्तरवर्ती $sn$ असतात, प्रत्येक $n \in \Nat$ साठी एक. $\sqsubseteq$
नुसार बंद असलेला प्रत्येक $A \subseteq \Nat^*$ हा $\Nat^*$ चा
\emph{उपवृक्ष} असतो. (म्हणजे $A$ असा आहे की $s \in A$ आणि
$s' \sqsubseteq s$ असतील, तर $s' \in A$ सुद्धा.) सर्व सांत वृक्ष
$\Nat^*$ चे सांत उपवृक्ष म्हणून दर्शवता येतात.
\end{ex}

\begin{prop}[K\H{o}nig यांचे उपप्रमेय]
$T = \tuple{A,\le}$ हा सांत शाखनाचा अनंत वृक्ष असेल, तर $T$ ला
एक अनंत शाखा असते.
\end{prop}

संगणनीयता सिद्धांतात वारंवार वापरले जाणारे K\H{o}nig यांच्या उपप्रमेयाचे
एक विशेष रूप, ज्याला \emph{दुर्बल K\H{o}nig उपप्रमेय} म्हणतात, असे आहे:
$\{0,1\}^*$ च्या प्रत्येक अनंत उपवृक्षाला एक अनंत शाखा असते.











\section{संबंधांवरील क्रिया}\label{sfr:rel:ops:sec}

संबंधांत बदल करणे किंवा त्यांना एकत्र करणे अनेकदा उपयोगी पडते.
\ref{sfr:rel:ord:prop:stricttopartial} मध्ये आपण संबंधांचा
\emph{संयोग} पाहिला; तो म्हणजे दोन संबंधांकडे जोड्यांचे संच म्हणून
पाहिल्यावर त्यांचा संयोग. त्याचप्रमाणे,
\ref{sfr:rel:ord:prop:partialtostrict} मध्ये संबंधांचा सापेक्ष फरक
पाहिला. संबंधांवर करता येणाऱ्या आणखी काही क्रिया पुढीलप्रमाणे आहेत.

\begin{defn}\label{sfr:rel:ops:relationoperations}
$R$, $S$ हे संबंध आणि $A$ हा कोणताही संच असू द्या.

$R$ चा \emph{व्युत्क्रम} म्हणजे $R^{-1} = \Setabs{\tuple{y, x}}{\tuple{x,
    y} \in R}$.

$R$ आणि $S$ यांचा \emph{सापेक्ष गुणाकार} म्हणजे $(R \mid S) =
\{\tuple{x, z} : \exists y(Rxy \land Syz)\}$.

$R$ चे $A$ वरील \emph{मर्यादन} म्हणजे $\funrestrictionto{R}{A}= R
\cap A^2$.

$R$ चा $A$ वर \emph{प्रयोग} म्हणजे $\funimage{R}{A} = \{y :
(\exists x \in A)Rxy\}$
\end{defn}

\begin{ex}
$S \subseteq \Int^2$ हा $\Int$ वरील उत्तरवर्ती संबंध असू द्या, म्हणजे
$S = \Setabs{\tuple{x, y} \in \Int^2}{x + 1 = y}$; त्यामुळे $Sxy$
तेव्हा आणि केवळ तेव्हाच, जेव्हा $x + 1 = y$.

$S^{-1}$ हा $\Int$ वरील पूर्ववर्ती संबंध आहे, म्हणजे
$\Setabs{\tuple{x,y}\in\Int^2}{x -1 =y}$.

$S\mid S$ म्हणजे
$ \Setabs{\tuple{x,y}\in\Int^2}{x + 2 =y}$

$\funrestrictionto{S}{\Nat}$ हा $\Nat$ वरील उत्तरवर्ती संबंध आहे.

$\funimage{S}{\{1,2,3\}}$ म्हणजे $\{2, 3, 4\}$.
\end{ex}

\begin{defn}[संक्रमक संवरण] $R \subseteq A^2$ हा द्विपदी संबंध असू द्या.

$R$ चे \emph{संक्रमक संवरण} म्हणजे $R^+ = \bigcup_{0 < n \in
\Nat} R^n$, जेथे $R^1 = R$ आणि $R^{n+1} = R^n \mid R$ अशी पुनरावर्ती
व्याख्या करतो.

$R$ चे \emph{परावर्ती संक्रमक संवरण} म्हणजे $R^* = R^+ \cup \Id{A}$.
\end{defn}

\begin{ex}
$S \subseteq \Int^2$ हा उत्तरवर्ती संबंध घ्या. $S^2xy$ तेव्हा आणि केवळ
तेव्हाच, जेव्हा $x + 2 = y$; $S^3xy$ तेव्हा आणि केवळ तेव्हाच, जेव्हा
$x + 3 = y$; इत्यादी. म्हणून $S^+xy$ तेव्हा आणि केवळ तेव्हाच, जेव्हा
एखाद्या $n \geq 1$ साठी $x + n = y$. दुसऱ्या शब्दांत, $S^+xy$ तेव्हा
आणि केवळ तेव्हाच, जेव्हा $x < y$; आणि $S^*xy$ तेव्हा आणि केवळ तेव्हाच,
जेव्हा $x \le y$.
\end{ex}

\begin{prob}
$R$ चे संक्रमक संवरण खरोखर संक्रमक आहे, हे दाखवा.
\end{prob}


\chapter{फलने}\label{sfr:fun::chap}





\section{मूलभूत संकल्पना}\label{sfr:fun:bas:sec}

\begin{explain}
दिलेल्या संचातील प्रत्येक घटक दुसऱ्या एखाद्या दिलेल्या संचातील
एका ठरावीक घटक शी जोडणारी संगती म्हणजे \emph{फलन}.
उदाहरणार्थ, $1$ मिळवण्याच्या क्रियेने एक फलन ठरते: प्रत्येक संख्या~$n$
एका एकमेव संख्या~$n+1$ शी जोडली जाते.

अधिक व्यापकपणे पाहता, फलनांना आदान म्हणून जोड्या, त्रिके इत्यादी देता
येतात आणि त्यांच्याकडून कोणत्या तरी प्रकारचे प्रदान मिळते. प्राथमिक
अंकगणितातून अनेक फलने आपल्याला परिचित आहेत. उदाहरणार्थ, बेरीज आणि
गुणाकार ही फलने आहेत. ती दोन संख्या आदान म्हणून घेतात आणि तिसरी संख्या देतात.

या गणिती, अमूर्त अर्थाने फलन म्हणजे एक \emph{काळी पेटी}: कोणत्या आदानाशी
कोणते प्रदान जोडलेले आहे एवढेच महत्त्वाचे असते; प्रदानाची गणना करण्याची
पद्धत महत्त्वाची नसते.
\end{explain}

\begin{defn}[फलन]
\emph{फलन} $f \colon A \to B$ ही~$A$ मधील प्रत्येक घटक ची~$B$
मधील एका घटक शी लावलेली संगती असते.

$A$ ला~$f$ चा \emph{प्रांत} आणि $B$ ला~$f$ चा \emph{सहप्रांत} म्हणतो.
~$A$ मधील घटक ना~$f$ ची आदाने किंवा \emph{फलसाधक} म्हणतात.
फलसाधक~$x$ शी~$f$ ने जोडलेल्या~$B$ मधील घटक ला फलसाधक~$x$
साठीचे \emph{$f$ चे मूल्य} म्हणतात आणि ते~$f(x)$ असे लिहितात.

~$f$ ची \emph{व्याप्ती} $\ran{f}$ म्हणजे सहप्रांताचा असा उपसंच की
त्याचे घटक कोणत्या तरी फलसाधकासाठीची~$f$ ची मूल्ये असतात;
$\ran{f} = \Setabs{f(x)}{x \in A}$.
\end{defn}

फलनांचा विचार करण्यासाठी \ref{sfr:fun:bas:fig:function} मधील आकृती उपयोगी पडू शकते.
डावीकडील लंबवर्तुळ फलनाचा \emph{प्रांत} दर्शवते; उजवीकडील लंबवर्तुळ
त्याचा \emph{सहप्रांत} दर्शवते; आणि बाण प्रांतातील \emph{फलसाधकापासून}
सहप्रांतातील त्याच्या संगत \emph{मूल्याकडे} जातो.

\begin{figure}
  \olasset{assets/diagrams/function.tikz}
  \caption{फलन एका संचातील प्रत्येक घटक ची दुसऱ्या संचातील
    घटक शी संगती लावते. बाण प्रांतातील फलसाधकापासून
    सहप्रांतातील त्याच्या संगत मूल्याकडे जातो.}
  \label{sfr:fun:bas:fig:function}
\end{figure}

\begin{ex}
गुणाकार नैसर्गिक संख्यांच्या जोड्या आदान म्हणून घेतो आणि त्यांना प्रदान
म्हणून नैसर्गिक संख्या जोडतो. म्हणून तो $\Nat \times \Nat$ या प्रांताकडून
$\Nat$ या सहप्रांताकडे जातो. त्याची व्याप्तीही $\Nat$ च असते, कारण
प्रत्येक $n \in \Nat$ हा $n \times 1$ असतो.
\end{ex}

\begin{ex}
गुणाकार हे फलन आहे, कारण तो प्रत्येक आदानाशी---नैसर्गिक संख्यांच्या
प्रत्येक जोडीशी---एकच प्रदान जोडतो: $\times \colon \Nat^2 \to \Nat$.
याउलट, नैसर्गिक संख्या~$n$ शी $x^2 = n$ असणारी वास्तव संख्या~$x$
जोडण्याची संगती फलन होत नाही, कारण प्रत्येक धन पूर्णांक~$n$ ची दोन
वर्गमुळे असतात: $\sqrt{n}$ आणि~$-\sqrt{n}$. फक्त धन वर्गमूळ देऊन
आपण तिचे फलन करू शकतो: फक्त ऋणेतर (प्रधान) वर्गमूळ परत द्यायचे,
$\sqrt{\phantom{X}} \colon \Nat \to \Real$.
\end{ex}
\begin{quote}\small\textbf{स्रोतदुरुस्ती.} OLFUN-002 — गोठवलेल्या इंग्रजी स्रोतामध्ये सर्व नैसर्गिक संख्यांसाठी positive square root म्हटले आहे. शून्याचे प्रधान वर्गमूळ शून्य असल्यामुळे येथे ऋणेतर (प्रधान) असे दुरुस्त केले आहे; धन पूर्णांकांना दोन वास्तव वर्गमुळे असतात हे आधीचे विधान योग्य आहे.\end{quote}

\begin{ex}
वर्गातील प्रत्येक विद्यार्थ्याशी त्याची अंतिम श्रेणी जोडणारा संबंध हे
फलन आहे---एकाच वर्गात कोणत्याही विद्यार्थ्याला दोन वेगवेगळ्या अंतिम
श्रेणी मिळू शकत नाहीत. वर्गातील प्रत्येक विद्यार्थ्याशी त्याचे पालक
जोडणारा संबंध फलन नाही: विद्यार्थ्यांचे पालक शून्य, दोन किंवा अधिक असू शकतात.
\end{ex}

\begin{explain}
प्रत्येक शक्य फलसाधकासाठी फलनाचे मूल्य काय असते हे नेमक्या पद्धतीने
सांगून फलनाची व्याख्या करता येते. सूत्र देणे, मूल्याची गणना करण्याची
पद्धत वर्णन करणे किंवा प्रत्येक फलसाधकासाठीची मूल्ये यादीत देणे या
त्याच्या वेगवेगळ्या पद्धती आहेत. फलनाची व्याख्या कोणत्याही पद्धतीने
केली, तरी प्रत्येक फलसाधकासाठी एक आणि केवळ एकच मूल्य दिलेले असेल
याची खात्री केली पाहिजे.
\end{explain}


\begin{ex}
$f \colon \Nat \to \Nat$ ची व्याख्या $f(x) = x+1$ अशी करू.
या व्याख्येत $f$ हे नैसर्गिक संख्या आदान म्हणून घेऊन नैसर्गिक संख्या
प्रदान करणारे फलन आहे असे सांगितले आहे. नैसर्गिक संख्या~$x$ दिली,
की $f$ तिची उत्तरवर्ती संख्या~$x+1$ देईल असे ती सांगते.
या उदाहरणात सहप्रांत $\Nat$ हा~$f$ ची व्याप्ती नाही, कारण नैसर्गिक
संख्या~$0$ ही कोणत्याही नैसर्गिक संख्येची उत्तरवर्ती संख्या नाही.
~$f$ ची व्याप्ती सर्व धन पूर्णांकांचा संच $\Int^{+}$ आहे.
\end{ex}

\begin{ex}\label{sfr:fun:bas:examplefunext}
$g \colon \Nat \to \Nat$ ची व्याख्या $g(x) = x+2-1$ अशी करू.
यावरून $g$ हे नैसर्गिक संख्या आदान म्हणून घेऊन नैसर्गिक संख्या प्रदान
करणारे फलन आहे हे कळते. नैसर्गिक संख्या x दिली, की $g$ हा~$x$ च्या
उत्तरवर्ती संख्येच्या उत्तरवर्ती संख्येची पूर्ववर्ती संख्या,
म्हणजेच $x+1$, देईल.
\end{ex}
\begin{quote}\small\textbf{स्रोतदुरुस्ती.} OLFUN-003 — गोठवलेल्या इंग्रजी स्रोताच्या गद्यामध्ये याच उदाहरणात आदानासाठी आधी $n$ आणि लगेच x वापरले आहेत. सूत्र न बदलता मराठी गद्याने x हे एकच नाव सातत्याने वापरले आहे.\end{quote}

\begin{explain}
आपण नुकतीच वेगवेगळ्या \emph{व्याख्या} असलेली $f$ आणि $g$ ही दोन
फलने पाहिली. पण ती \emph{एकच फलन} आहेत. कारण कोणत्याही नैसर्गिक
संख्या~$n$ साठी $f(n) = n+1 = n+2-1 = g(n)$ असते. दुसऱ्या शब्दांत,
$f$ आणि~$g$ यांच्या व्याख्या वेगवेगळ्या समीकरणांनी एकच संगती सांगतात.
म्हणून आपण अप्रत्यक्षपणे फलनांसाठीच्या विस्तारात्मकतेच्या तत्त्वावर
अवलंबून आहोत:
\[  \text{जर }\forall x\, f(x) = g(x)\text{, तर }f = g
\]
मात्र $f$ आणि~$g$ यांचा प्रांत आणि सहप्रांत एकच असला पाहिजे.
\end{explain}

\begin{ex}
प्रकरणे पाडूनही फलनाची व्याख्या करता येते. उदाहरणार्थ,
$h \colon \Nat \to \Nat$ ची व्याख्या अशी करता येईल:
\[h(x) =
\begin{cases}
  \frac{x}{2} & \text{जर $x$ सम असेल} \\
  \frac{x+1}{2} & \text{जर $x$ विषम असेल.}
\end{cases}
\]
प्रत्येक नैसर्गिक संख्या सम किंवा विषम असते, म्हणून या फलनाचे प्रदान
नेहमी नैसर्गिक संख्या असेल. प्रकरणे पाडून फलनाची व्याख्या करताना
प्रत्येक शक्य आदान नेमक्या एका प्रकरणात बसले पाहिजे, हे लक्षात ठेवा.
काही वेळा सर्व शक्यता त्या प्रकरणांत येतात आणि कोणतीही शक्यता दोन
प्रकरणांत येत नाही याची सिद्धता करावी लागेल.
\end{ex}








\section{फलनांचे प्रकार}\label{sfr:fun:kin:sec}

\begin{explain}
आपल्याला वारंवार भेटणाऱ्या फलनांच्या काही प्रकारांचे वर्गीकरण करून
घेणे उपयोगी ठरेल.

सुरुवातीला, सहप्रांतातील प्रत्येक सदस्य फलनाचे मूल्य असतो असा गुणधर्म
असलेल्या फलनांचा विचार करू. अशा फलनांना आच्छादक म्हणतात.
त्यांची आकृती \ref{sfr:fun:kin:fig:surjective} प्रमाणे काढता येते.

\begin{figure}
  \olasset{assets/diagrams/surjective.tikz}
  \caption{आच्छादक फलनात सहप्रांतातील प्रत्येक घटक
    हे फलनाचे मूल्य असते.}
  \label{sfr:fun:kin:fig:surjective}
\end{figure}
\end{explain}

\begin{defn}[आच्छादक फलन]
$f \colon A \rightarrow B$ हे फलन \emph{आच्छादक} असते तेव्हा
आणि केवळ तेव्हाच, जेव्हा $B$ हा~$f$ ची व्याप्तीही असतो. म्हणजे,
प्रत्येक $y \in B$ साठी~$f(x) = y$ असणारा किमान एक $x \in A$
असतो. चिन्हांत:
\[  (\forall y \in B)(\exists x \in A)f(x) = y.
\]
अशा फलनाला $A$ पासून $B$ पर्यंतचे आच्छादन म्हणतो.
\end{defn}

\begin{explain}
$f$ हे आच्छादन आहे असे दाखवायचे असल्यास, $f$ च्या सहप्रांतातील
प्रत्येक वस्तू कोणत्या तरी आदान $x$ साठी $f(x)$ चे मूल्य आहे असे दाखवावे लागते.

कोणत्याही फलनापासून एक आच्छादन \emph{निर्माण होते}, हे लक्षात घ्या.
कारण $f \colon A \to B$ हे फलन दिले असल्यास, $f' \colon A \to \ran{f}$
ची व्याख्या $f'(x) = f(x)$ अशी करू. $\ran{f}$ ची \emph{व्याख्याच}
$\Setabs{f(x) \in B}{x \in A}$ अशी असल्यामुळे हे फलन $f'$ निश्चितच
आच्छादन असते.
\end{explain}

\begin{explain}
कोणतेही फलन प्रत्येक शक्य आदानाशी एकमेव प्रदान जोडते. पण वेगवेगळ्या
आदानांशी कधीच एकच प्रदान न जोडणारी फलनेही असतात. अशा फलनांना
एकास-एक म्हणतात. त्यांची आकृती \ref{sfr:fun:kin:fig:injective} प्रमाणे काढता येते.
\begin{figure}
  \olasset{assets/diagrams/injective.tikz}
  \caption{एकास-एक फलन कधीच दोन वेगवेगळ्या फलसाधकांशी
    एकच मूल्य जोडत नाही.}
  \label{sfr:fun:kin:fig:injective}
\end{figure}
\end{explain}

\begin{defn}[एकास-एक फलन]
$f \colon A \rightarrow B$ हे फलन \emph{एकास-एक} असते तेव्हा
आणि केवळ तेव्हाच, जेव्हा प्रत्येक $y \in B$ साठी~$f(x) = y$ असणारा
जास्तीत जास्त एक $x \in A$ असतो. अशा फलनाला $A$ पासून~$B$ पर्यंतचे
एकास-एक फलन म्हणतो.
\end{defn}

\begin{explain}
$f$ हे एकास-एक फलन आहे असे दाखवायचे असल्यास, $f$ च्या प्रांतातील
कोणत्याही घटक $x$ आणि $y$ साठी $f(x)=f(y)$ असल्यास $x=y$
असते असे दाखवावे लागते.
\end{explain}

\begin{ex}
$f(x) = 1$ ने दिलेले स्थिर फलन $f\colon \Nat \to \Nat$ हे
एकास-एक ही नाही आणि आच्छादक ही नाही.

$f(x) = x$ ने दिलेले एकरूपता फलन $f\colon \Nat \to \Nat$ हे
एकास-एक आणि आच्छादक दोन्ही आहे.

$f(x) = x+1$ ने दिलेले उत्तरवर्ती फलन $f \colon \Nat \to \Nat$
हे एकास-एक आहे, पण आच्छादक नाही.

पुढीलप्रमाणे व्याख्या केलेले फलन $f \colon \Nat \to \Nat$:
\[  f(x) =
  \begin{cases}
    \frac{x}{2} & \text{जर $x$ सम असेल} \\
    \frac{x+1}{2} & \text{जर $x$ विषम असेल.}
  \end{cases}
\]
हे आच्छादक आहे, पण एकास-एक नाही.
\end{ex}

\begin{explain}
अनेकदा एकास-एक आणि आच्छादक हे दोन्ही गुणधर्म असलेल्या
फलनांचा विचार करायचा असतो. अशा फलनांना एकास-एक व आच्छादक म्हणतो.
ती \ref{sfr:fun:kin:fig:bijective} मधील फलनासारखी दिसतात. एकास-एक आच्छादने ना
कधी कधी \emph{एकास-एक संगती} असेही म्हणतात, कारण ती सहप्रांतातील
घटकांची प्रांतातील घटकांशी एकमेव जोडी लावतात.
\begin{figure}
  \olasset{assets/diagrams/bijective.tikz}
  \caption{एकास-एक व आच्छादक फलन सहप्रांतातील घटकांची प्रांतातील
    घटकांशी एकमेव जोडी लावते.}
  \label{sfr:fun:kin:fig:bijective}
\end{figure}
\end{explain}

\begin{defn}[एकास-एक आच्छादन]
$f \colon A \to B$ हे फलन \emph{एकास-एक व आच्छादक} असते तेव्हा आणि
केवळ तेव्हाच, जेव्हा ते आच्छादक आणि एकास-एक दोन्ही असते.
अशा फलनाला $A$ पासून~$B$ पर्यंतचे (किंवा $A$ आणि~$B$ यांमधील)
एकास-एक आच्छादन म्हणतो.
\end{defn}










\section{संबंध म्हणून फलने}\label{sfr:fun:rel:sec}

\begin{explain}
~$A$ मधील घटक ची~$B$ मधील घटक शी संगती लावणारे
फलन $A$ आणि~$B$ यांच्यामधील एक संबंध ठरवते: $f(x) = y$ असेल
तेव्हा आणि केवळ तेव्हाच $x$ आणि $y$ यांच्यात हा संबंध असतो.
खरे तर, उदाहरणार्थ गणिताच्या पायाभूत घटकांची संख्या कमी करायची असेल,
तर आपण फलन~$f$ आणि हा संबंध, म्हणजे जोड्यांचा एक संच, \emph{एकरूप}
मानू शकतो. मग असा प्रश्न उभा राहतो: कोणते संबंध अशा रीतीने फलने ठरवतात?
\end{explain}

\begin{defn}[फलनाचा आलेख] $f\colon A \to B$ हे फलन असू द्या.
~$f$ चा \emph{आलेख} म्हणजे पुढील व्याख्येने मिळणारा संबंध
$R_f \subseteq A \times B$:
\[R_f = \Setabs{\tuple{x,y}}{f(x) = y}.
\]
\end{defn}

\begin{explain}
विस्तारात्मकतेमुळे फलनाचा आलेख एकमेव ठरतो. शिवाय, संचांसाठीच्या
विस्तारात्मकतेमुळे फलनांसाठीचे अप्रत्यक्ष विस्तारात्मकता तत्त्व लगेच
समर्थित होते: $f$ आणि~$g$ यांचा प्रांत व सहप्रांत एकच असेल आणि
सर्व मूल्यांवर ती जुळत असतील, तर ती एकच फलन असतात.

तसेच, एखादा संबंध ``फलनरूप'' असेल, तर तो एका फलनाचा आलेख असतो.
\end{explain}

\begin{prop}\label{sfr:fun:rel:prop:graph-function}
$R \subseteq A \times B$ हा संबंध पुढील अटी पूर्ण करतो असे समजा:
\begin{enumerate}
\item $Rxy$ आणि $Rxz$ असल्यास $y = z$ असते; आणि
\item प्रत्येक $x \in A$ साठी $\tuple{x,
y} \in R$ असणारा काही $y \in B$ असतो.
\end{enumerate}
मग $f(x) = y$ तेव्हा आणि केवळ तेव्हाच $Rxy$, अशा व्याख्येने मिळणाऱ्या
फलन $f\colon A \to B$ चा आलेख $R$ असतो.
\end{prop}

\begin{proof}
$Rxy$ असणारा एक $y$ आहे असे समजा. $Rxz$ असणारा दुसरा $z \neq y$
असता, तर~$R$ वरील अट मोडली गेली असती. म्हणून $Rxy$ असणारा एक $y$
असेल, तर हा $y$ एकमेव असतो. त्यामुळे $f$ सुस्पष्टपणे परिभाषित आहे.
उघडच, $R_f = R$.
\end{proof}

\begin{explain}
प्रत्येक फलन $f\colon A \to B$ ला एक आलेख असतो, म्हणजे $f(x) = y$
ने परिभाषित होणारा A आणि B यांच्यातील, $A \times B$ मध्ये
समाविष्ट असलेला एक संबंध असतो. उलट,
\ref{sfr:fun:rel:prop:graph-function} मधील गुणधर्म असलेला प्रत्येक संबंध~$R
\subseteq A \times B$ हा एका फलन~$f \colon A \to B$ चा आलेख असतो.
फलने आणि त्यांचे आलेख यांचा हा निकटचा संबंध असल्यामुळे फलन म्हणजे
त्याचा आलेखच असे आपण मानू शकतो. दुसऱ्या शब्दांत, फलने विशिष्ट
संबंधांशी, म्हणजे विशिष्ट क्रमित घटकसमूहांच्या संचांशी, एकरूप मानता
येतात. पण या ``एकरूप मानण्याचा'' आशय
\ref{sfr:rel:ref:sec} मधल्यासारखाच आहे, हे लक्षात घ्या: तो
फलनांच्या सत्तामीमांसेविषयीचा दावा नाही, तर फलनांना विशिष्ट संच
\emph{मानणे} सोयीचे आहे हे निरीक्षण आहे. ही सोय होण्याचे एक कारण असे की,
संबंधांवर जशा क्रिया केल्या तशाच क्रिया फलनांवर करण्याचा आता आपण
विचार करू शकतो (\ref{sfr:rel:ops:sec} पाहा). विशेषतः:
\end{explain}
\begin{quote}\small\textbf{स्रोतदुरुस्ती.} OLFUN-004 — इंग्रजी स्रोत येथे graph ला relation on A times B म्हणतो; पण त्याच ग्रंथाच्या व्याख्येनुसार त्याचा शब्दशः अर्थ (A times B) च्या वर्गाचा उपसंच होईल. मराठी मजकुरात अचूक अर्थ A आणि B यांच्यातील, A times B मध्ये समाविष्ट संबंध असा दिला आहे.\end{quote}

\begin{defn}\label{sfr:fun:rel:defn:funimage}
$f \colon A \to B$ हे फलन असू द्या आणि $C\subseteq A$ असू द्या.

~$f$ चे~$C$ पर्यंतचे \emph{मर्यादन} म्हणजे सर्व $x \in C$ साठी
$(\funrestrictionto{f}{C})(x) = f(x)$ या व्याख्येने मिळणारे
फलन~$\funrestrictionto{f}{C}\colon C \to B$. दुसऱ्या शब्दांत,
$\funrestrictionto{f}{C} = \Setabs{\tuple{x, y} \in R_f}{x \in C}$.

~$f$ चे~$C$ वरील \emph{उपयोजन} म्हणजे $\funimage{f}{C} =
\Setabs{f(x)}{x \in C}$. यालाच~$f$ अंतर्गत~$C$ ची \emph{प्रतिमा}
असेही म्हणतो.
\end{defn}

\begin{explain}
या व्याख्यांवरून कोणत्याही फलन~$f$ साठी $\ran{f} =
\funimage{f}{\dom{f}}$ असते.
\ref{sfr:rel:ops:sec} मधील
व्याख्या आणि फलनांना संबंधांशी एकरूप मानण्याची आपली पद्धत लक्षात
घेतल्यास या संकल्पना अनुरूप आहेत. मात्र फलनाचे C पर्यंतचे मर्यादन
केवळ आदान निर्देशांक C मध्ये ठेवते; संबंधाचे C वरील मर्यादन
दोन्ही निर्देशांक C मध्ये ठेवते. त्यामुळे त्या अगदी एकच क्रिया नाहीत.
व्युत्क्रम आणि सापेक्ष गुणाकार या इतर दोन क्रियांना थोडे अधिक
स्पष्टीकरण लागते. ते आपण
\ref{sfr:fun:inv:sec} आणि \ref{sfr:fun:cmp:sec} मध्ये देऊ.
\end{explain}
\begin{quote}\small\textbf{स्रोतदुरुस्ती.} OLFUN-005 — फलनाच्या मर्यादनाचे स्पष्ट सूत्र योग्य आहे: आलेखात पहिला निर्देशांक C मध्ये असतो. आधीचे संबंधावरील मर्यादन दोन्ही निर्देशांक C मध्ये ठेवते. इंग्रजी स्रोताचा exactly असा दावा येथे अनुरूप पण भिन्न असा स्पष्ट केला आहे.\end{quote}








\section{फलनांचे व्युत्क्रम}\label{sfr:fun:inv:sec}

\begin{explain}
फलनांचा विचार आपण संगती म्हणून करतो. मग ही संगती ``उलटवता'' येते का,
हा प्रश्न स्वाभाविकपणे पडतो. उदाहरणार्थ, उत्तरवर्ती फलन $f(x) = x + 1$
उलटवता येते, या अर्थाने की $g(y) = y - 1$ हे फलन $f$ ने केलेली
क्रिया ``पूर्ववत करते''.

पण येथे काळजी घेतली पाहिजे. ~$g$ च्या व्याख्येने $\Int \to \Int$
असे फलन ठरत असले, तरी $\Nat \to \Nat$ असे \emph{फलन} ठरत नाही,
कारण $g(0) \notin \Nat$. म्हणून साध्या उदाहरणांतसुद्धा फलन उलटवता
येते का हे तितकेसे उघड नसते; ते प्रांत आणि सहप्रांतावर अवलंबून असू शकते.

फलनाच्या व्युत्क्रमाच्या संकल्पनेने हा विचार अधिक नेमका करता येतो.
\end{explain}

\begin{defn}
सर्व $x \in A$ आणि $y \in B$ साठी $f(g(y)) = y$ आणि $g(f(x)) = x$
असेल, तर $g \colon B \to A$ हे फलन $f \colon A \to B$ या फलनाचा
\emph{व्युत्क्रम} असते.
\end{defn}

$f$ ला व्युत्क्रम~$g$ असेल, तर आपण अनेकदा~$g$ ऐवजी $f^{-1}$ असे लिहितो.

\begin{explain}
आता कोणत्या फलनांना व्युत्क्रम असतात हे ठरवू. $f\colon A \to B$
च्या व्युत्क्रमासाठी एक संभाव्य फलन म्हणजे पुढीलप्रमाणे ``व्याख्या केलेले''
$g\colon B \to A$:
\[g(y) = \text{$f(x) = y$ असणारा ``तो'' $x$.}
\]
पण ``व्याख्या केलेले'' आणि ``तो'' यांभोवतीची अवतरणचिन्हे सुचवतात की
ही खरे तर व्याख्या नाही. किमान पूर्ण व्यापकतेने ती नेहमी लागू पडणार
नाही. कारण या व्याख्येने फलन ठरवायचे असेल, तर $f(x) = y$ असणारा
एक आणि केवळ एकच~$x$ असला पाहिजे---~$g$ चे प्रदान एकमेव ठरले पाहिजे.
शिवाय ते प्रत्येक $y \in B$ साठी ठरले पाहिजे. $x_1$ आणि $x_2 \in A$
हे $x_1 \neq x_2$ असूनही $f(x_1) = f(x_2)$ असतील, तर
$y = f(x_1) = f(x_2)$ साठी $g(y)$ एकमेव ठरणार नाही.
आणि $f(x) = y$ असणारा~$x$ मुळीच नसेल, तर $g(y)$ मुळीच ठरणार नाही.
दुसऱ्या शब्दांत, $g$ परिभाषित होण्यासाठी $f$ हे एकास-एक आणि
आच्छादक दोन्ही असले पाहिजे.

हे टप्प्याटप्प्याने पाहू. प्रश्नाचे दोन भाग करू: फलन~$f\colon A \to B$
दिले असता, $g(f(x)) = x$ असणारे फलन $g\colon B \to A$ कधी असते?
असे $g$ हे $f$ ने केलेली क्रिया ``पूर्ववत करते'', आणि त्याला~$f$ चा
\emph{डावा व्युत्क्रम} म्हणतात. दुसरे म्हणजे, $f(h(y)) = y$ असणारे
फलन $h\colon B \to A$ कधी असते? अशा $h$ ला~$f$ चा
\emph{उजवा व्युत्क्रम} म्हणतात---$f$ हे~$h$ ने केलेली क्रिया ``पूर्ववत करते''.
\end{explain}

\begin{prop}
A अरिक्त असेल आणि $f\colon A \to B$ हे एकास-एक असेल,
तर~$f$ चा एक
\emph{डावा व्युत्क्रम}~$g\colon B \to A$ असतो, आणि सर्व $x \in A$
साठी $g(f(x)) = x$ असते.
\end{prop}
\begin{quote}\small\textbf{स्रोतदुरुस्ती.} OLFUN-001 — गोठवलेल्या इंग्रजी प्रमेयात A अरिक्त असण्याची अट सुटली आहे. रिक्त A पासून अरिक्त B कडेचे रिक्त फलन एकास-एक असूनही B पासून A कडे फलन नसते. नेमकी सामान्य अट: एकास-एक फलनाला डावा व्युत्क्रम तेव्हा आणि केवळ तेव्हाच असतो, जेव्हा A अरिक्त असतो किंवा B रिक्त असतो. मराठी प्रमेयात साधी पुरेशी A अरिक्त ही अट जोडली आहे.\end{quote}

\begin{proof}
A अरिक्त आहे आणि $f\colon A \to B$ हे एकास-एक आहे असे समजा.
$y \in B$ विचारात घ्या.
$y \in \ran{f}$ असेल, तर $f(x) = y$ असणारा एक $x \in A$ असतो.
$f$ हे एकास-एक असल्यामुळे असा~$x \in A$ एकच असतो. म्हणून
$g(y) = x$ अशी व्याख्या करता येते; म्हणजे $g(y)$ हा $f(x) = y$
असणारा ``तो'' $x \in A$ असतो. $y \notin \ran{f}$ असेल, तर त्याची
संगती कोणत्याही~$a \in A$ शी लावता येते. म्हणून एक $a \in A$
निवडून $g\colon B \to A$ ची व्याख्या अशी करता येते:
\[g(y) = \begin{cases}
    x & \text{जर $f(x) = y$ असेल}\\
    a & \text{जर $y \notin \ran{f}$ असेल.}
\end{cases}
\]
हे सर्व $y \in B$ साठी परिभाषित आहे, कारण अशा प्रत्येक $y \in \ran{f}$
साठी $f(x) = y$ असणारा नेमका एक $x \in A$ असतो. व्याख्येनुसार,
$y = f(x)$ असल्यास $g(y) = x$, म्हणजे $g(f(x)) = x$.
\end{proof}

\begin{prob}
$f\colon A \to B$ ला डावा व्युत्क्रम~$g$ असल्यास $f$ हे
एकास-एक असते, हे दाखवा.
\end{prob}

\begin{prop}
    $f\colon A \to B$ हे आच्छादक असेल, तर~$f$ चा एक
    \emph{उजवा व्युत्क्रम}~$h\colon B \to A$ असतो, आणि सर्व~$y \in B$
    साठी $f(h(y)) = y$ असते.
\end{prop}

\begin{proof}
$f\colon A \to B$ हे आच्छादक आहे असे समजा. $y \in B$ विचारात घ्या.
~$f$ हे आच्छादक असल्यामुळे $f(x_y) = y$ असणारा एक $x_y \in A$
असतो. म्हणून $h(y) = x_y$ अशी व्याख्या करता येते: प्रत्येक $y \in B$
साठी $f(x) = y$ असणारा काही $x \in A$ आपण निवडतो. ~$f$ हे
आच्छादक असल्यामुळे निवडण्यासाठी नेहमी किमान एक घटक असतो.\footnote{
$f$ हे आच्छादक असल्यामुळे प्रत्येक~$y \in B$ साठी
$\Setabs{x}{f(x) = y}$ हा संच अरिक्त असतो. ~$h$ ची आपली व्याख्या
या प्रत्येक संचातून एकच $x$ निवडण्याची मागणी करते. हे नेहमी करता येते
ही गोष्ट खरे तर उघड नाही---या निवडी करता येतात हे एक स्वयंसिद्धक म्हणून
गृहीत धरले जाते. दुसऱ्या शब्दांत, ही प्रतिज्ञा तथाकथित निवडीचे स्वयंसिद्धक
गृहीत धरते. या प्रश्नाकडे आपण सध्या दुर्लक्ष करू.
मात्र अनेक विशिष्ट उदाहरणांत, जसे $A = \Nat$ असेल किंवा सांत असेल,
किंवा $f$ हे एकास-एक व आच्छादक असेल, तेव्हा निवडीचे स्वयंसिद्धक लागत नाही.
($f$ हे एकास-एक व आच्छादक असण्याच्या विशिष्ट उदाहरणात प्रत्येक $y \in B$
साठी $\Setabs{x}{f(x) = y}$ या संचात नेमका एक घटक असतो,
म्हणून निवड करण्याचा प्रश्नच उरत नाही.)}
व्याख्येनुसार, $x = h(y)$ असल्यास $f(x) = y$ असते; म्हणजे कोणत्याही
$y \in B$ साठी $f(h(y)) = y$.
\end{proof}

\begin{prob}
$f\colon A \to B$ ला उजवा व्युत्क्रम~$h$ असल्यास $f$ हे
आच्छादक असते, हे दाखवा.
\end{prob}

\begin{explain}
  मागील सिद्धतेतील कल्पना एकत्र केल्यावर प्रत्येक एकास-एक आच्छादन ला
  व्युत्क्रम असतो हे मिळते. म्हणजे~$f$ चा डावा आणि उजवा व्युत्क्रम
  असे दोन्ही असणारे एकच फलन असते.
\end{explain}

\begin{prop}\label{sfr:fun:inv:prop:bijection-inverse}
$f\colon A \to B$ हे एकास-एक व आच्छादक असेल, तर एक
फलन~$f^{-1}\colon B \to A$ असते, आणि सर्व $x \in A$ साठी
$f^{-1}(f(x)) = x$ तसेच सर्व $y \in B$ साठी $f(f^{-1}(y)) = y$ असते.
\end{prop}

\begin{proof}
सरावासाठी.
\end{proof}

\begin{prob}
\ref{sfr:fun:inv:prop:bijection-inverse} सिद्ध करा. ~$f^{-1}$ ची
व्याख्या करा, ते फलन आहे हे दाखवा आणि ते~$f$ चा व्युत्क्रम आहे हे
दाखवा. म्हणजे, सर्व $x \in A$ आणि $y \in B$ साठी
$f^{-1}(f(x)) = x$ आणि $f(f^{-1}(y)) = y$ हे दाखवा.
\end{prob}

\begin{explain}
व्युत्क्रम मिळवण्याची थोडी अधिक व्यापक पद्धत आहे. \ref{sfr:fun:kin:sec}
मध्ये आपण पाहिले की प्रत्येक फलन $f$ पासून सर्व $x \in A$ साठी
$f'(x) = f(x)$ असे ठेवून आच्छादन $f' \colon A \to \ran{f}$
मिळते. उघडच, ~$f$ हे एकास-एक असेल, तर~$f'$ हे एकास-एक व आच्छादक
असते. म्हणून \ref{sfr:fun:inv:prop:bijection-inverse} नुसार त्याला एकमेव
व्युत्क्रम असतो. संकेतलेखनात अगदी थोडी सवलत घेऊन आपण कधी कधी
$f'$ च्या व्युत्क्रमाला फक्त ``~$f$ चा व्युत्क्रम'' म्हणतो.
\end{explain}

\begin{prop}\label{sfr:fun:inv:prop:left-right}%
  $f\colon A \to B$ ला डावा व्युत्क्रम~$g$ आणि उजवा व्युत्क्रम~$h$
  असल्यास $h = g$ असते, हे दाखवा.
\end{prop}

\begin{proof}
  सरावासाठी.
\end{proof}

\begin{prob}
  \ref{sfr:fun:inv:prop:left-right} सिद्ध करा.
\end{prob}

\begin{prop}\label{sfr:fun:inv:prop:inverse-unique}
प्रत्येक फलन~$f$ ला जास्तीत जास्त एक व्युत्क्रम असतो.
\end{prop}

\begin{proof}
  $g$ आणि $h$ हे दोन्ही~$f$ चे व्युत्क्रम आहेत असे समजा. मग विशेषतः
  ~$g$ हा~$f$ चा डावा व्युत्क्रम आणि~$h$ हा उजवा व्युत्क्रम असतो.
  \ref{sfr:fun:inv:prop:left-right} नुसार $g = h$.
\end{proof}








\section{फलनांचे संयोजन}\label{sfr:fun:cmp:sec}

\begin{explain}
\ref{sfr:fun:inv:sec} मध्ये आपण पाहिले की
एकास-एक आच्छादन~$f$ चा व्युत्क्रम~$f^{-1}$ हाही एक फलन असतो.
फलनांवरील आणखी एक क्रिया म्हणजे संयोजन: $f$ आणि~$g$ या दोन
फलनांचे संयोजन करून, म्हणजे आधी $f$ आणि मग~$g$ उपयोजून,
एक नवे फलन परिभाषित करता येते. अर्थात, व्याप्ती आणि प्रांत यांचा मेळ
असेल तेव्हाच हे शक्य असते; म्हणजे~$f$ ची व्याप्ती ही~$g$ च्या प्रांताचा
उपसंच असली पाहिजे. फलनांवरील ही क्रिया
\ref{sfr:rel:ops:sec} मधील संबंधांवरील सापेक्ष गुणाकाराच्या क्रियेशी
समांतर आहे.

संयोजनाची कल्पना स्पष्ट करण्यासाठी आकृती उपयोगी पडू शकते.
\ref{sfr:fun:cmp:fig:composition} मध्ये $f \colon A \to B$ आणि $g \colon B \to C$
ही दोन फलने आणि त्यांचे संयोजन~$(\comp{f}{g})$ दाखवले आहे.
$(\comp{f}{g}) \colon A \to C$ हे फलन~$A$ मधील प्रत्येक घटक
शी~$C$ मधील घटक जोडते. $A$ मधील घटक शी~$C$
मधील कोणता घटक जोडायचा हे असे ठरवतो: आदान $x \in A$
दिले असता, प्रथम~$x$ वर फलन $f$ उपयोजा. त्यामुळे काही
$f(x) = y \in B$ हे प्रदान मिळेल. मग~$y$ वर फलन $g$ उपयोजा.
त्यामुळे काही $g(f(x)) = g(y) = z \in C$ हे प्रदान मिळेल.
\begin{figure}
  \olasset[2\olphotowidth]{assets/diagrams/composition.tikz}
  \caption{$f$ आणि~$g$ या दोन फलनांचे संयोजन $g \circ f$.}
  \label{sfr:fun:cmp:fig:composition}
\end{figure}
\end{explain}

\begin{defn}[संयोजन]
$f\colon A \to B$ आणि $g\colon B \to C$ ही फलने असू द्या.
$f$ चे~$g$ बरोबरचे \emph{संयोजन} म्हणजे $\comp{f}{g} \colon A \to C$,
जेथे $(\comp{f}{g})(x) = g(f(x))$.
\end{defn}

\begin{ex}
$f(x) = x + 1$ आणि $g(x) = 2x$ ही फलने विचारात घ्या.
$(\comp{f}{g})(x) = g(f(x))$ असल्यामुळे प्रत्येक आदान~$x$ साठी
प्रथम त्याची उत्तरवर्ती संख्या घेऊन नंतर तिला दोनने गुणले पाहिजे.
म्हणून त्यांचे संयोजन $(\comp{f}{g})(x) = 2(x+1)$ असे मिळते.
\end{ex}

\begin{prob}
$f \colon A \to B$ आणि $g \colon B \to C$ ही दोन्ही फलने
एकास-एक असल्यास $\comp{f}{g}\colon A \to C$ हे
एकास-एक असते, हे दाखवा.
\end{prob}

\begin{prob}
$f \colon A \to B$ आणि $g \colon B \to C$ ही दोन्ही फलने
आच्छादक असल्यास $\comp{f}{g}\colon A \to C$ हे
आच्छादक असते, हे दाखवा.
\end{prob}

\begin{prob}
$f \colon A \to B$ आणि $g \colon B \to C$ असे समजा.
$\comp{f}{g}$ चा आलेख $R_f \mid R_g$ आहे हे दाखवा.
\end{prob}










\section{अंशतः फलने}\label{sfr:fun:par:sec}

\begin{explain}
कधी कधी फलनाच्या व्याख्येत थोडी सवलत देणे उपयोगी ठरते: सर्व शक्य
आदानांसाठी फलनाचे प्रदान परिभाषित असलेच पाहिजे ही अट काढून टाकता
येते. अशा संगतींना \emph{अंशतः फलने} म्हणतात.
\end{explain}

\begin{defn}
\emph{अंशतः फलन} $f \colon A \pto B$ ही अशी संगती असते की ती~$A$
मधील प्रत्येक घटक शी~$B$ मधील जास्तीत जास्त एक घटक
जोडते. $f$ ने $x \in A$ शी~$B$ मधील एखादा घटक जोडला असेल, तर
$f(x)$ \emph{परिभाषित} आहे असे म्हणतो; अन्यथा ते \emph{अपरिभाषित}
असते. $f(x)$ परिभाषित असल्यास $f(x) \fdefined$ असे लिहितो;
अन्यथा $f(x) \fundefined$ असे लिहितो. अंशतः फलन~$f$ चा \emph{प्रांत}
म्हणजे ते परिभाषित असलेल्या~$A$ मधील घटकांचा उपसंच:
$\dom{f} = \Setabs{x \in A}{f(x) \fdefined}$.
\end{defn}

\begin{ex}
प्रत्येक फलन $f\colon A \to B$ हे अंशतः फलनही असते. ~$A$ वरील
सर्व ठिकाणी परिभाषित असलेल्या अंशतः फलनांना---म्हणजे ज्यांना आपण
आतापर्यंत फक्त फलन म्हणत आलो, त्यांना---\emph{सर्वत्र परिभाषित}
फलने असेही म्हणतात.
\end{ex}

\begin{ex}
$f(x) = 1/x$ ने दिलेले अंशतः फलन $f \colon \Real \pto \Real$
हे $x = 0$ साठी अपरिभाषित असते आणि इतर सर्व ठिकाणी परिभाषित असते.
\end{ex}

\begin{prob}
$f\colon A \pto B$ दिले असता, अंशतः फलन $g\colon B \pto A$ ची
व्याख्या अशी करा: कोणत्याही $y \in B$ साठी $f(x) = y$ असणारा एकमेव
$x \in A$ असेल, तर $g(y) = x$; अन्यथा $g(y) \fundefined$.
$f$ हे एकास-एक असल्यास सर्व $x \in \dom{f}$ साठी $g(f(x)) = x$
आणि सर्व $y \in \ran{f}$ साठी $f(g(y)) = y$ असते, हे दाखवा.
\end{prob}

\begin{defn}[अंशतः फलनाचा आलेख]
$f\colon A \pto B$ हे अंशतः फलन असू द्या. ~$f$ चा \emph{आलेख}
म्हणजे पुढील व्याख्येने मिळणारा संबंध $R_f \subseteq A \times B$:
\[R_f = \Setabs{\tuple{x,y}}{f(x) = y}.
\]
\end{defn}

\begin{prop}
$R \subseteq A \times B$ या संबंधाचा असा गुणधर्म आहे असे समजा की
$Rxy$ आणि $Rxy'$ असल्यास $y = y'$ असते. मग $R$ हा पुढीलप्रमाणे
व्याख्या केलेल्या अंशतः फलन $f\colon A \pto B$ चा आलेख असतो:
$Rxy$ असणारा एक $y$ असेल, तर $f(x) = y$; अन्यथा $f(x) \fundefined$.
शिवाय $R$ हा \emph{प्रत्येक आदानाशी किमान एक प्रदान जोडणारा} असेल,
म्हणजे प्रत्येक $x \in A$ साठी $Rxy$ असणारा एक $y \in B$ असेल,
तर $f$ हे सर्वत्र परिभाषित असते.
\end{prop}

\begin{proof}
$Rxy$ असणारा एक $y$ आहे असे समजा. $Rxy'$ असणारा दुसरा
$y' \neq y$ असता, तर $R$ वरील अट मोडली गेली असती. म्हणून
$Rxy$ असणारा एक $y$ असेल, तर तो $y$ एकमेव असतो. त्यामुळे $f$
सुस्पष्टपणे परिभाषित आहे. उघडच $R_f = R$, आणि~$R$ प्रत्येक आदानाशी
किमान एक प्रदान जोडत असेल, तर $f$ हे सर्वत्र परिभाषित असते.
\end{proof}



\chapter{संचांचे आकारमान}\label{sfr:siz::chap}






\section{परिचय}\label{sfr:siz:int:sec}

1870 च्या दशकात गेओर्क कँटर यांनी संच सिद्धांत
विकसित केला, तेव्हा अनंत समूहाची कल्पना---मध्ययुगीन विचारवंतांनी
म्हटल्याप्रमाणे प्रत्यक्ष अस्तित्वातील अनंतता---स्वीकारार्ह करणे हे
त्यांच्या उद्दिष्टांपैकी एक होते. वेगवेगळ्या संचांच्या \emph{आकारमानाचा}
त्यांनी केलेला विचार हा त्यातील महत्त्वाचा भाग होता. $a$, $b$ आणि $c$
हे सर्व भिन्न असतील, तर $\{a, b, c\}$ हा संच $\{a, b\}$ पेक्षा
\emph{मोठा} आहे असे सहज वाटते. पण अनंत संचांचे काय? ते सर्व
एकमेकांएवढेच मोठे असतात का? तसे नसते, हे दिसून येते.

येथील पहिली महत्त्वाची कल्पना म्हणजे प्रगणन. कोणत्याही सांत संचातील
सर्व घटक लिहून त्याची यादी देता येते. काही अनंत संचांतील
सर्व घटक चीही यादी देता येते, जर यादी स्वतः अनंत असण्याची
मुभा दिली, तर. अशा संचांना गणनीय म्हणतात. काही अनंत संच
गणनीय नसतात, हा कँटर यांचा आश्चर्यकारक निष्कर्ष होता.
या प्रकरणाच्या अखेरीस तो आपल्याला पूर्णपणे समजेल.









\section{प्रगणने आणि गणनीय संच}\label{sfr:siz:enm:sec}

\begin{editorial}
  या विभागात संचांच्या प्रगणनांची चर्चा आहे. त्यांची व्याख्या
  $\PosInt$ पासूनच्या आच्छादनांनी केली आहे. गणिताची फारशी पार्श्वभूमी
  नसलेल्या वाचकांसाठी विषय हळूहळू मांडला आहे. पर्यायी, अधिक संक्षिप्त
  आवृत्ती \ref{sfr:siz:enm-alt:sec} मध्ये दिली आहे. तेथे प्रगणनांची वेगळी
  व्याख्या केली आहे: $\Nat$ किंवा त्याच्या आरंभीच्या खंडाशी
  एकास-एक आच्छादन असणे.
\end{editorial}

\begin{explain}
संचांचे घटक यादीत देऊन आपण आधीच संचांची उदाहरणे दिली आहेत.
संच अनंत असला, तरी त्याच्या घटक ची यादी कशी आणि कधी
देता येते, याची आता अधिक व्यापकपणे चर्चा करू.
\end{explain}

\begin{defn}[प्रगणनाची अनौपचारिक व्याख्या]
अनौपचारिकपणे, संच~$A$ चे \emph{प्रगणन} म्हणजे~$A$ मधील
घटक ची अशी यादी (कदाचित अनंत) की $A$ मधील प्रत्येक
घटक त्या यादीत कोणत्या तरी सांत क्रमांकाच्या स्थानी येतो.
$A$ चे प्रगणन असेल, तर $A$ हा \emph{गणनीय} आहे असे म्हणतात.
\end{defn}

\begin{explain}
प्रगणनांविषयी काही मुद्दे:
\begin{enumerate}
\item ज्या यादीला सुरुवात असते आणि जिच्यात पहिल्या घटकाव्यतिरिक्त
  प्रत्येक घटक च्या लगेच आधी एकच घटक असतो, अशीच यादी
  आपण प्रगणन मानतो. दुसऱ्या शब्दांत, यादीच्या पहिल्या घटक
  आणि इतर कोणत्याही घटक यांच्यामध्ये केवळ सांत संख्येने घटक
  असतात. विशेषतः, प्रगणनातील प्रत्येक घटक चे स्थान सांत
  क्रमांकाचे असते: पहिल्या घटक चे स्थान~$1$, दुसऱ्याचे~$2$, इत्यादी.
\item एकाच संच~$A$ ची वेगवेगळी प्रगणने असू शकतात, ज्यांत
  घटक येण्याचा क्रम वेगळा असतो: $4$, $1$, $25$, $16$,~$9$
  हे पहिल्या पाच वर्गसंख्यांच्या संचाचे प्रगणन आहे, तसेच
  $1$, $4$, $9$, $16$,~$25$ हेही आहे.
\item पुनरुक्ती असलेली प्रगणनेही प्रगणनेच असतात: $1$, $1$, $2$,
  $2$, $3$, $3$,~\dots{} हे त्याच संचाचे प्रगणन आहे, ज्याचे
  प्रगणन $1$, $2$, $3$,~\dots{} हे आहे.
\item विशिष्ट प्रगणन सांगताना क्रम आणि पुनरुक्ती यांना महत्त्व
  \emph{असते}: धन पूर्णांकांचे प्रगणन $1$, $2$, $3$, $1$, \dots{}
  अशा सुरुवातीने करता येते; पण नेहमीच्या $1$, $2$, $3$, $4$,~\dots
  या पद्धतीने प्रगणन केले, तर त्यातील नमुना अधिक सहज दिसतो.
\item प्रगणनाला सुरुवात असली पाहिजे: \dots, $3$, $2$, $1$ हे
  धन पूर्णांकांचे प्रगणन नाही, कारण त्याला पहिला घटक नाही.
  अनौपचारिक व्याख्येवरून हे कसे येते ते पाहण्यासाठी स्वतःला विचारा:
  ``या यादीत 76 ही संख्या कोणत्या स्थानी येते?''
\item पुढील यादी धन पूर्णांकांचे प्रगणन नाही:
  $1$, $3$, $5$, \dots, $2$, $4$, $6$, \dots\@ अडचण अशी की
  सम संख्या सांत क्रमांकांच्या स्थानांवर येण्याऐवजी
  $\infty + 1$, $\infty + 2$, $\infty + 3$ या स्थानांवर येतात.
\item रिक्त संच गणनीय आहे: रिक्त यादी हे त्याचे प्रगणन आहे!{}
\end{enumerate}
\end{explain}

\begin{prop}
  $A$ चे प्रगणन असेल, तर त्याचे पुनरुक्ती नसलेले प्रगणनही असते.
\end{prop}

\begin{proof}
  $A$ चे $x_1$, $x_2$, \dots{} असे प्रगणन आहे, आणि प्रत्येक
  $x_i$ हा~$A$ चा घटक आहे असे समजा. प्रगणनातून पुन्हा
  आलेले घटक काढून टाकून त्यातील पुनरुक्ती काढता येते.
  उदाहरणार्थ, पुढीलप्रमाणे नवे प्रगणन मिळवता येते: $x_i$ हा~$A$
  मधील घटक असून $x_1$, \dots, $x_{i-1}$ यांमध्ये नसेल,
  तर त्याला यादीत लिहायचे; आणि $x_i$ हा आधीच $x_1$, \dots,~$x_{i-1}$
  यांमध्ये असेल, तर त्याला यादीतून काढून टाकायचे.
\end{proof}

मागील युक्तिवाद दाखवतो की प्रगणने आणि गणनीय संच नीट
समजून घेण्यासाठी, तसेच त्यांविषयी गोष्टी सिद्ध करण्यासाठी,
अधिक नेमकी व्याख्या हवी. ती पुढे दिली आहे.

\begin{defn}[प्रगणनाची औपचारिक व्याख्या]
$A \neq \emptyset$ या संचाचे \emph{प्रगणन} म्हणजे कोणतेही
आच्छादक फलन $f \colon \PosInt \to A$.
\end{defn}

\begin{explain}
औपचारिक व्याख्या आणि कदाचित अनंत असणाऱ्या यादीचा वापर करणारी
अनौपचारिक व्याख्या समतुल्य आहेत याची खात्री करू. प्रथम, $\PosInt$
पासून संच~$A$ पर्यंतचे कोणतेही आच्छादक फलन~$A$ चे प्रगणन करते.
अशा फलनाने वरील अनौपचारिक अर्थाचे प्रगणन ठरते: $f(1)$, $f(2)$,
$f(3)$, \dots{} ही यादी. $f$ हे आच्छादक असल्यामुळे~$A$
मधील प्रत्येक घटक हा कोणत्या तरी~$n \in \PosInt$ साठीचे
~$f(n)$ चे मूल्य असतो. म्हणून $A$ मधील प्रत्येक घटक यादीत
सांत क्रमांकाच्या स्थानी येतो. फलन एकास-एक असेलच असे नसल्याने
यादीत पुनरुक्ती असू शकते; पण वर म्हटल्याप्रमाणे ती चालते.

उलट,~$A$ मधील सर्व घटक चे प्रगणन करणारी यादी दिली असल्यास
आच्छादक फलन $f\colon \PosInt \to A$ असे परिभाषित करता येते:
$f(n)$ म्हणजे यादीतील $n$ व्या स्थानी असलेला घटक; आणि
$n$ व्या स्थानी घटक नसेल, तर यादीचा शेवटचा घटक.
$A$ रिक्त असतो आणि म्हणून यादीही रिक्त असते, फक्त या एकाच बाबतीत
या पद्धतीने आच्छादक फलन मिळत नाही. म्हणून प्रत्येक अरिक्त
यादी एक आच्छादक फलन $f\colon \PosInt \to A$ ठरवते.
\end{explain}

\begin{defn}
  \label{sfr:siz:enm:defn:enumerable}
  संच~$A$ हा गणनीय असतो तेव्हा आणि केवळ तेव्हाच,
  जेव्हा तो रिक्त असतो किंवा त्याचे प्रगणन असते.
\end{defn}

\begin{ex}
धन पूर्णांकांचे ($\PosInt$) प्रगणन करणारे फलन म्हणजे फक्त
$f(n) = n$ ने दिलेले एकरूपता फलन. नैसर्गिक संख्या $\Nat$ यांचे
प्रगणन करणारे फलन म्हणजे $g(n) = n - 1$.
\end{ex}

\begin{ex}
$f\colon \PosInt \to \PosInt$ आणि $g \colon \PosInt \to \PosInt$
ही फलने अशी दिली आहेत:
\begin{align*}
f(n) & = 2n \text{ आणि}\\
g(n) & = 2n - 1
\end{align*}
ती अनुक्रमे सम धन पूर्णांकांचे आणि विषम धन पूर्णांकांचे प्रगणन करतात.
मात्र यांपैकी कोणतेही फलन $\PosInt$ चे प्रगणन नाही, कारण कोणतेही
आच्छादक नाही.
\end{ex}

\begin{prob}
धन वर्गसंख्या $1$, $4$, $9$, $16$, \dots{} यांचे प्रगणन परिभाषित करा.
\end{prob}

\begin{ex}
$f(n) = (-1)^{n} \lceil \frac{(n-1)}{2}\rceil$ हे फलन पूर्णांकांचा
संच~$\Int$ प्रगणित करते. येथे $\lceil x \rceil$ हे
\emph{ऊर्ध्व पूर्णांक फलन} दर्शवते; ते $x$ ला त्याच्यापेक्षा कमी
नसलेल्या सर्वांत जवळच्या पूर्णांकाकडे नेते. धन आणि ऋण पूर्णांकांमध्ये
आलटूनपालटून ``उड्या मारून'' $f$ हे $\Int$ ची मूल्ये कशी निर्माण
करते ते पाहा:
\[\begin{array}{c c c c c c c c}
f(1) & f(2) & f(3) & f(4) & f(5) & f(6) & f(7) & \dots \\ \\
- \lceil \tfrac{0}{2} \rceil & \lceil \tfrac{1}{2}\rceil & - \lceil \tfrac{2}{2} \rceil & \lceil \tfrac{3}{2} \rceil & - \lceil \tfrac{4}{2} \rceil  & \lceil \tfrac{5}{2}
\rceil & - \lceil \tfrac{6}{2} \rceil & \dots \\ \\
0 & 1 & -1 & 2 & -2 & 3 & -3 & \dots
\end{array}
\]
\begin{quote}\small\textbf{स्रोतदुरुस्ती.} OLSIZ-001: गोठवलेल्या इंग्रजी तक्त्याच्या शेवटच्या ओळीत f(7) साठी अपेक्षित \ensuremath{-}3 हे मूल्य सुटले आहे; सूत्र आणि वरील दोन ओळी ते मूल्य निश्चित करतात. मराठी तक्त्यात \ensuremath{-}3 समाविष्ट करून ही उणीव दुरुस्त केली आहे.\end{quote}
$f$ ची व्याख्या पुढीलप्रमाणे प्रकरणे पाडून केली आहे असेही मानता येते:
\[f(n) = \begin{cases}
  0 & \text{जर $n = 1$ असेल}\\
  n/2 & \text{जर $n$ सम असेल}\\
  -(n-1)/2 & \text{जर $n$ विषम आणि $>1$ असेल}
  \end{cases}
\]
\end{ex}

\begin{prob}
  $A$ आणि $B$ हे गणनीय असतील, तर $A \cup B$ हाही
  गणनीय असतो, हे दाखवा. यासाठी आच्छादक फलने
  $f\colon \PosInt \to A$ आणि $g\colon \PosInt \to B$ आहेत असे
  समजा. एक आच्छादक फलन~$h\colon \PosInt \to A \cup B$
  परिभाषित करा आणि ते आच्छादक आहे हे सिद्ध करा. $A$ किंवा
  ~$B = \emptyset$ असण्याच्याही बाबतींचा विचार करा.
\end{prob}

\begin{prob}
  $B \subseteq A$ आणि $A$ हा गणनीय असेल, तर~$B$ हाही
  गणनीय असतो, हे दाखवा. यासाठी एक आच्छादक फलन
  $f\colon \PosInt \to A$ आहे असे समजा. एक आच्छादक
  फलन~$g\colon \PosInt \to B$ परिभाषित करा आणि ते
  आच्छादक आहे हे सिद्ध करा. $B = \emptyset$ असल्यास काय होते?
\end{prob}

\begin{prob}
  $A_1$, $A_2$, \dots, $A_n$ हे सर्व गणनीय असतील, तर
  $A_1 \cup \dots \cup A_n$ हाही गणनीय असतो, हे
  $n$ वर विगमन करून दाखवा. $A$ आणि~$B$ हे दोन संच
  गणनीय असल्यास~$A \cup B$ हाही गणनीय असतो
  हे तथ्य तुम्ही गृहीत धरू शकता.
\end{prob}

संचातील घटक ची यादी करताना $1$ व्या घटक पासून
मोजायला सुरुवात करणे कदाचित अधिक स्वाभाविक असले, तरी गणितज्ञांना
मोजणीसाठी नैसर्गिक संख्या~$\Nat$ वापरायला आवडतात. ते यादीतील
$0$ व्या, $1$ व्या, $2$ व्या, अशा घटक विषयी बोलतात.
त्याला अनुसरून प्रगणनाची व्याख्या $\Nat$ पासून~$A$ पर्यंतचे
आच्छादक फलन अशी करता येते. अर्थात, या दोन व्याख्या समतुल्य आहेत.

\begin{prop}\label{sfr:siz:enm:prop:enum-shift}
  आच्छादन $f\colon \PosInt \to A$ असते तेव्हा आणि
  केवळ तेव्हाच, जेव्हा आच्छादन $g\colon \Nat \to A$ असते.
\end{prop}

\begin{proof}
  आच्छादन $f\colon \PosInt \to A$ दिले असता, सर्व
  $n \in \Nat$ साठी $g(n) = f(n+1)$ अशी व्याख्या करता येते.
  $g\colon \Nat \to A$ हे आच्छादक आहे हे सहज दिसते.
  उलट, आच्छादन $g\colon \Nat \to A$ दिले असता,
  $f(n) = g(n-1)$ अशी व्याख्या करा.
\end{proof}

यावरून पुढील निष्कर्ष मिळतो:

\begin{cor}\label{sfr:siz:enm:cor:enum-nat}
संच $A$ हा गणनीय असतो तेव्हा आणि केवळ तेव्हाच, जेव्हा
तो रिक्त असतो किंवा आच्छादक फलन $f\colon \Nat \to A$ असते.
\end{cor}

संच~$A$ मधील घटक च्या यादीतून पुनरुक्ती नसलेली यादी
मिळवता येते याची वर चर्चा केली. प्रगणनांबाबतही हे खरे आहे,
पण ते काटेकोरपणे मांडणे आणि सिद्ध करणे थोडे कठीण आहे. कोणतेही
फलन $f\colon \PosInt \to A$ हे सर्व $n \in \PosInt$ साठी
परिभाषित असले पाहिजे. ~$A$ मध्ये फक्त सांत संख्येने घटक
असतील, तर~$\PosInt$ मधील अनंत संख्येने असलेल्या घटक
वर परिभाषित असणारे आणि~$A$ मधील सर्व घटक मूल्ये म्हणून
घेणारे, पण एकच मूल्य दोनदा कधीच न घेणारे फलन असू शकत नाही.
अशा बाबतीत, म्हणजे पुनरुक्ती नसलेली यादी सांत असते तेव्हा,~$f$
साठी वेगळा प्रांत निवडला पाहिजे; त्यात केवळ सांत संख्येने
घटक असले पाहिजेत. पुनरुक्ती नसणे म्हणजे $f$ हे
एकास-एक असले पाहिजे. ते आच्छादक ही असल्यामुळे
आपण कोणता तरी सांत संच $\{1, \dots, n\}$ किंवा $\PosInt$
आणि~$A$ यांच्यामधील एकास-एक आच्छादन शोधत आहोत.

\begin{prop}\label{sfr:siz:enm:prop:enum-bij}
$f\colon \PosInt \to A$ हे आच्छादक असेल (म्हणजे~$A$
चे प्रगणन असेल), तर एकास-एक आच्छादन $g\colon Z \to A$ असते,
जेथे $Z$ हा~$\PosInt$ किंवा कोणत्या तरी~$n \in \PosInt$
साठी $\{1, \dots, n\}$ असतो.
\end{prop}

\begin{proof}
  फलन $g$ ची व्याख्या आपण पुनरावर्ती रीतीने करतो: $g(1) = f(1)$
  असे ठेवा. $g(i)$ आधीच परिभाषित झाले असेल, तर $f(1)$, $f(2)$,
  \dots{} यांपैकी $g(1)$, \dots, $g(i)$ यांत आधीच न आलेले
  पहिले मूल्य, असे मूल्य असेल तर, $g(i+1)$ म्हणून घ्या.
  $A$ मध्ये नेमके $n$ घटक असतील, तर $g(1)$, \dots, $g(n)$
  ही सर्व मूल्ये परिभाषित असतात, आणि आपल्याला
  $g\colon \{1, \dots, n\} \to A$ असे फलन मिळते.
  $A$ मध्ये अनंत संख्येने घटक असतील, तर कोणत्याही $i$
  साठी $f(1)$, $f(2)$, \dots{} या प्रगणनात~$A$ मधील असा
  घटक असलाच पाहिजे, जो $g(1)$, \dots, $g(i)$ यांत
  आधीच आलेला नाही. अशा बाबतीत आपण $g\colon \PosInt \to A$
  हे फलन परिभाषित केले आहे.

  फलन $g$ हे आच्छादक आहे, कारण~$A$ मधील कोणताही घटक
  $f(1)$, $f(2)$, \dots{} यांमध्ये येतो ($f$ हे आच्छादक
  असल्यामुळे), आणि म्हणून कधी ना कधी तो कोणत्या तरी~$i$
  साठीचे~$g(i)$ चे मूल्य होईल. ते एकास-एक ही आहे:
  $g(j) = g(i)$ असणारे $j < i$ असते, तर $g(i)$ हे मूल्य
  आधीच $g(1)$, \dots, $g(i-1)$ यांमध्ये आलेले असते;
  पण हे~$g$ च्या आपल्या व्याख्येविरुद्ध आहे.
\end{proof}

\begin{cor}\label{sfr:siz:enm:cor:enum-nat-bij}
संच $A$ हा गणनीय असतो तेव्हा आणि केवळ तेव्हाच, जेव्हा
तो रिक्त असतो किंवा एकास-एक आच्छादन $f\colon N \to A$ असते,
जेथे $N = \Nat$ किंवा कोणत्या तरी $n \in \Nat$ साठी
$N = \{0, \dots, n\}$ असते.
\end{cor}

\begin{proof}
$A$ हा गणनीय असतो तेव्हा आणि केवळ तेव्हाच, जेव्हा $A$
रिक्त असतो किंवा आच्छादक $f\colon \PosInt \to A$ असते.
\ref{sfr:siz:enm:prop:enum-bij} नुसार दुसरी अट तेव्हा आणि केवळ तेव्हाच
खरी असते, जेव्हा एकास-एक व आच्छादक फलन~$f\colon Z \to A$ असते,
जेथे $Z = \PosInt$ किंवा कोणत्या तरी $n \in \PosInt$ साठी
$Z = \{1, \dots, n\}$ असते. \ref{sfr:siz:enm:prop:enum-shift} च्या
सिद्धतेतील युक्तिवादानुसार, हे तेव्हा आणि केवळ तेव्हाच शक्य होते,
जेव्हा एकास-एक आच्छादन $g\colon N \to A$ असते, जेथे
$N = \Nat$ किंवा $N = \{0, \dots, n-1\}$ असते.
\begin{quote}\small\textbf{स्रोतदुरुस्ती.} MRSIZ-001: विधानातील सांत प्रांत 0 ते n आणि सिद्धतेतील 0 ते n\ensuremath{-}1 हे वेगवेगळ्या चलनामांनी त्याच सांत आरंभीच्या खंडांच्या कुटुंबाचे वर्णन करतात; येथे स्रोताचे दोन्ही निर्देशांक जसेच्या तसे ठेवले आहेत.\end{quote}
\end{proof}

\begin{prob}
  \ref{sfr:siz:enm:defn:enumerable} नुसार संच $A$ हा गणनीय
  असतो तेव्हा आणि केवळ तेव्हाच, जेव्हा $A = \emptyset$ किंवा
  आच्छादक $f\colon \PosInt \to A$ असते.
  ``गणनीय संच'' याची नेमकी व्याख्या पुढीलप्रमाणेही करता येते:
  संच गणनीय असतो तेव्हा आणि केवळ तेव्हाच, जेव्हा एकास-एक
  फलन $g\colon A \to \PosInt$ असते. या व्याख्या समतुल्य आहेत
  हे दाखवा. म्हणजे, एकास-एक फलन $g\colon A \to \PosInt$
  असते तेव्हा आणि केवळ तेव्हाच, जेव्हा $A = \emptyset$ किंवा
  आच्छादक $f\colon \PosInt \to A$ असते, हे दाखवा.
\end{prob}









\section{कँटर यांची नागमोडी पद्धत}\label{sfr:siz:zigzag:sec}

\begin{explain}
काही ``सोप्या'' प्रगणनांचा आपण आधीच विचार केला आहे. आता जरा कठीण
उदाहरण पाहू. नैसर्गिक संख्यांच्या जोड्यांचा संच विचारात घ्या
; त्याची व्याख्या आपण
\ref{sfr:set:pai:sec} मध्ये पुढीलप्रमाणे केली:
\[\Nat \times \Nat = \Setabs{\tuple{n,m}}{n,m \in \Nat}
\]
या क्रमित जोड्या आपण पुढीलप्रमाणे एका \emph{सरणीत} मांडू शकतो:
\[\begin{array}{ c | c | c | c | c | c}
& \mathbf 0 & \mathbf 1 & \mathbf 2 & \mathbf 3 & \dots \\
\hline
\mathbf 0 & \tuple{0,0} & \tuple{0,1} & \tuple{0,2} & \tuple{0,3} & \dots \\
\hline
\mathbf 1 & \tuple{1,0} & \tuple{1,1} & \tuple{1,2} & \tuple{1,3} & \dots \\
\hline
\mathbf 2 & \tuple{2,0} & \tuple{2,1} & \tuple{2,2} & \tuple{2,3} & \dots \\
\hline
\mathbf 3 & \tuple{3,0} & \tuple{3,1} & \tuple{3,2} & \tuple{3,3} & \dots \\
\hline
\vdots & \vdots & \vdots & \vdots & \vdots & \ddots\\
\end{array}
\]
$\Nat \times \Nat$ मधील प्रत्येक क्रमित जोडी या सरणीत नेमकी एकदाच
येईल, हे स्पष्ट आहे. विशेषतः, $\tuple{n,m}$ ही जोडी $n$ व्या ओळीत
आणि $m$ व्या स्तंभात येईल. पण अशा सरणीतील घटकांची ``एकमितीय'' यादी
कशी करायची? पुढील सरणीतील नमुना ती करण्याचा एक मार्ग दाखवतो
(अर्थातच इतर अनेक मार्गही आहेत):
\[\begin{array}{ c | c | c | c | c | c | c}
& \mathbf 0 & \mathbf 1 & \mathbf 2 & \mathbf 3 & \mathbf 4 &\dots \\
\hline
\mathbf 0 & 0  & 1& 3 & 6& 10 &\ldots \\
\hline
\mathbf 1 &2 & 4& 7 & 11 & \dots &\ldots \\
\hline
\mathbf 2 & 5 & 8 & 12 & \ldots & \dots&\ldots \\
\hline
\mathbf 3 & 9 & 13 & \ldots & \ldots & \dots & \ldots \\
\hline
\mathbf 4 & 14 & \ldots & \ldots & \ldots & \dots & \ldots \\
\hline
\vdots & \vdots & \vdots & \vdots & \vdots&\ldots & \ddots\\
\end{array}
\]\noindent
या नमुन्याला \emph{कँटर यांची नागमोडी पद्धत} म्हणतात. तो
$\Nat \times \Nat$ चे पुढीलप्रमाणे प्रगणन करतो:
\[\tuple{0,0}, \tuple{0,1}, \tuple{1,0}, \tuple{0,2}, \tuple{1,1},
\tuple{2,0}, \tuple{0,3}, \tuple{1,2}, \tuple{2,1}, \tuple{3,0}, \dots
\]
यावरून पुढील विधान सिद्ध होते:
\end{explain}

\begin{prop}\label{sfr:siz:zigzag:natsquaredenumerable}
$\Nat \times \Nat$ हा गणनीय आहे.
\end{prop}

\begin{proof}
प्रत्येक $k \in \Nat$ ला कँटर यांच्या नागमोडी सरणीतील ज्या $n$ व्या
ओळीच्या आणि $m$ व्या स्तंभाच्या जागेचे मूल्य $k$ आहे, त्या
$\tuple{n,m} \in \Nat \times \Nat$ या क्रमित घटकाशी जोडणारे
$f \colon \Nat \to \Nat\times\Nat$ असे फलन घ्या.
\end{proof}

\begin{explain}
ही पद्धत अधिक व्यापक बाबतीतही सहज लागू होते. उदाहरणार्थ, नैसर्गिक
संख्यांच्या क्रमित तीन-घटकसमूहांचा पुढील संच प्रगणित करण्यासाठी ती
वापरता येते:
\[\Nat \times \Nat \times \Nat = \Setabs{\tuple{n,m,k}}{n,m,k \in \Nat}
\]
$\Nat \times \Nat \times \Nat$ हा $\Nat \times \Nat$ आणि $\Nat$
यांचा कार्तीय गुणाकार आहे असे आपण मानतो; म्हणजे,
\[\Nat^3 = (\Nat \times \Nat) \times \Nat =
\Setabs{\tuple{\tuple{n,m},k}}{n, m, k
  \in \Nat }
\]
म्हणून एका अक्षाला $\Nat$ चे प्रगणन आणि दुसऱ्या अक्षाला $\Nat^2$
चे प्रगणन अशी नावे देऊन, एका सरणीच्या साहाय्याने $\Nat^3$ चे
प्रगणन करता येते:
\[\begin{array}{ c | c | c | c | c | c}
& \mathbf 0 & \mathbf 1 & \mathbf 2 & \mathbf 3 & \dots \\
\hline
\mathbf{\tuple{0,0}} & \tuple{0,0,0} & \tuple{0,0,1} & \tuple{0,0,2} & \tuple{0,0,3} & \dots \\
\hline
\mathbf{\tuple{0,1}} & \tuple{0,1,0} & \tuple{0,1,1} & \tuple{0,1,2} & \tuple{0,1,3} & \dots \\
\hline
\mathbf{\tuple{1,0}} & \tuple{1,0,0} & \tuple{1,0,1} & \tuple{1,0,2} & \tuple{1,0,3} & \dots \\
\hline
\mathbf{\tuple{0,2}} & \tuple{0,2,0} & \tuple{0,2,1} & \tuple{0,2,2} & \tuple{0,2,3} & \dots\\
\hline
\vdots & \vdots & \vdots & \vdots & \vdots & \ddots \\
\end{array}
\]
अशा रीतीने कँटर यांच्या नागमोडी पद्धतीसारखी पद्धत वापरून $\Nat^3$
चे प्रगणन मिळते. ही प्रक्रिया पुढे चालू ठेवून, प्रत्येक नैसर्गिक संख्या
$n$ साठी $\Nat^n$ चे प्रगणन मिळवता येते. म्हणून पुढील विधान मिळते:
\end{explain}

\begin{prop}
प्रत्येक $n \in \Nat$ साठी $\Nat^n$ हा गणनीय आहे.
\end{prop}

\begin{prob}\label{sfr:siz:zigzag:prob:posint-n}
प्रत्येक $n \in \Nat$ साठी $(\PosInt)^n$ हा गणनीय आहे, हे दाखवा.
\end{prob}

\begin{prob}\label{sfr:siz:zigzag:prob:posint-star}
  $(\PosInt)^*$ हा गणनीय आहे, हे दाखवा. तुम्ही \ref{sfr:siz:zigzag:prob:posint-n} गृहीत धरू शकता.
\end{prob}








\section{जोडीकरण फलने आणि संकेतांक}\label{sfr:siz:pai:sec}

\begin{explain}
कँटर यांच्या नागमोडी पद्धतीमुळे $\Nat^n$ ची गणनीयता चित्रातून
स्पष्ट दिसते. पण $\Nat^2$ दाखवणाऱ्या आपल्या सरणीवर लक्ष केंद्रित करू.
सरणीतील नागमोडी रेषेचा मागोवा घेऊन स्थाने मोजल्यास,
$\tuple{1,2}$ ही जोडी~$7$ या संख्येशी जोडली आहे हे तपासता येते.
पण ही संख्या अधिक थेट मोजता आली, तर सोयीचे होईल. म्हणजे, नागमोडी
प्रगणनाचे \emph{व्युत्क्रम} $g\colon \Nat^2 \to \Nat$ हाताशी असावे,
आणि त्यासाठी
\[g(\tuple{0,0}) = 0, \;
g(\tuple{0,1}) = 1, \;
g(\tuple{1,0}) = 2, \; \dots,
g(\tuple{1,2}) = 7, \; \dots
\]
असे असावे. त्यामुळे $\tuple{n, m}$ ही जोडी आपल्या प्रगणनात नेमकी
कोणत्या स्थानी येईल ते मोजता येईल.

खरे तर दोन निरीक्षणांच्या साहाय्याने $g$ ची थेट व्याख्या करता येते.
पहिले निरीक्षण: $n$ व्या ओळीतील आणि $m$ व्या स्तंभातील मूल्य~$v$
असेल, तर $(n+1)$ व्या ओळीतील आणि $(m-1)$ व्या स्तंभातील मूल्य
$v + 1$ असते. दुसरे निरीक्षण: आपल्या प्रगणनाची पहिली ओळ $0$, $1$,
$3$, $6$, इत्यादींपासून सुरू होणाऱ्या त्रिकोणी संख्यांची आहे.
$k$ वी त्रिकोणी संख्या म्हणजे $< k$ ही अट पूर्ण करणाऱ्या नैसर्गिक संख्यांची
बेरीज; ती $k(k+1)/2$ अशी मोजता येते. ही दोन निरीक्षणे एकत्र करून
पुढील फलन विचारात घ्या:
\[  g(n,m) = \frac{(n+m+1)(n+m)}{2} + n
\]
$g(\tuple{n, m})$ ऐवजी आपण बहुधा फक्त $g(n, m)$ असे लिहितो, कारण
ते पाहायला सोपे असते. या सूत्रानुसार आधी $(n+m)^\text{वी}$ त्रिकोणी
संख्या काढायची आणि नंतर तिच्यात $n$ मिळवायचा. त्यामुळे सरणी नेमकी
आपल्याला हवी तशी भरते. विशेषतः, $\tuple{1, 2}$ ही जोडी
$\frac{4 \times 3}{2} + 1 = 7$ कडे पाठवली जाते.

हे फलन $g$ म्हणजे जोड्यांच्या संचाच्या एका प्रगणनाचे \emph{व्युत्क्रम}
आहे. अशा फलनांना \emph{जोडीकरण फलने} म्हणतात.
\end{explain}

\begin{defn}[जोडीकरण फलन]
  $f\colon A \times B \to \Nat$ यातील $f$ हे एकास-एक असेल, तर ते
  \emph{अंकगणिती जोडीकरण फलन} असते. $f$ हे $A \times B$ चे
  \emph{सांकेतीकरण करते} आणि $f(x,y)$ हा $\tuple{x,y}$ चा
  \emph{संकेतांक} असतो असेही आपण म्हणतो.
\end{defn}

\begin{explain}
जोडीकरण फलनांनी, उदाहरणार्थ, नैसर्गिक संख्यांच्या जोड्यांचे
सांकेतीकरण करता येते; म्हणजेच घटकांची प्रत्येक \emph{जोडी} एका
\emph{एकाच} संख्येने दर्शवता येते. जोडीकरण फलनाचे व्युत्क्रम वापरून
संख्येचे \emph{विसांकेतीकरण} करता येते, म्हणजे ती संख्या कोणती जोडी
दर्शवते हे शोधता येते.
\end{explain}

\begin{prob}
सर्व ऋणेतर परिमेय संख्यांच्या संचाचे एक प्रगणन द्या.
\end{prob}

\begin{prob}
$\Rat$ हा गणनीय आहे, हे दाखवा. कोणतीही परिमेय संख्या
$z/m$ या अपूर्णांकाच्या रूपात लिहिता येते, जेथे $z \in \Int$ आणि
$m \in \Nat^+$ असतात, हे आठवा.
\end{prob}

\begin{prob}
$\Bin^*$ चे एक प्रगणन परिभाषित करा.
\end{prob}

\begin{prob}
प्रास्ताविक तर्कशास्त्राच्या अभ्यासक्रमातून आठवा की प्रत्येक संभाव्य
सत्यताकोष्टक एक सत्यता-फलन व्यक्त करतो. दुसऱ्या शब्दांत, सत्यता-फलने
म्हणजे कोणत्या तरी~$k$ साठी $\Bin^k \to \Bin$ अशी सर्व फलने आहेत.
सर्व सत्यता-फलनांचा संच गणनीय आहे, हे सिद्ध करा.
\end{prob}

\begin{prob}
कोणत्याही अनंत गणनीय संचाच्या सर्व सांत उपसंचांचा संच
गणनीय आहे, हे दाखवा.
\end{prob}

\begin{prob}
$\Nat$ च्या एखाद्या सांत उपसंचाचा पूरक असलेल्या $\Nat$ च्या उपसंचाला
\emph{सांत-पूरक} म्हणतात. म्हणजे, $A \subseteq \Nat$ हा सांत-पूरक
असतो तेव्हा आणि केवळ तेव्हाच, जेव्हा $\Nat\setminus A$ हा सांत असतो.
\begin{quote}\small\textbf{स्रोतदुरुस्ती.} OLSIZ-002: गोठवलेल्या इंग्रजी वाक्यात “complement of a finite set Nat” अशी संबंधदर्शक शब्दांशिवाय अपूर्ण रचना आहे. लगेचचे सूत्र अभिप्रेत अर्थ ठरवते: Nat मधील एखाद्या सांत उपसंचाचा Nat-सापेक्ष पूरक. मराठी मजकुरात हा संबंध स्पष्ट केला आहे.\end{quote}
$\Nat$ चे सांत आणि सांत-पूरक उपसंच हेच ज्याचे घटक आहेत,
असा संच $I$ घ्या. $I$ हा गणनीय आहे, हे दाखवा.
\end{prob}

\begin{prob}
गणनीय इतक्या गणनीय संचांचा संयोग गणनीय असतो,
हे दाखवा. म्हणजे, $A_1$, $A_2$, \dots{} हे संच असतील आणि प्रत्येक
$A_i$ हा गणनीय असेल, तर त्या सर्वांचा संयोग
$\bigcup_{i=1}^\infty A_i$ हाही गणनीय असतो. [टीप: हे कठीण आहे!]
\end{prob}

\begin{prob}
$f \colon A \times B \to \Nat$ हे कोणतेही जोडीकरण फलन असू द्या.
हे $f$ आच्छादकही असेल, तर त्याचे व्युत्क्रम $A \times B$ चे प्रगणन
आहे, हे दाखवा. सर्वसाधारण बाबतीत व्युत्क्रम केवळ फलनाच्या व्याप्तीवर
परिभाषित असते; यावरून जोड्यांचा हा संच गणनीय आहे, हे सिद्ध करा.
\begin{quote}\small\textbf{स्रोतदुरुस्ती.} MRSIZ-002: गोठवलेल्या इंग्रजी सरावात कोणत्याही एकास-एक जोडीकरण फलनाचे व्युत्क्रम थेट प्रगणन असल्याचा दावा आहे. फलन आच्छादक नसेल, तर त्याचे व्युत्क्रम सर्व नैसर्गिक संख्यांवर परिभाषित नसते. मराठी सरावात आच्छादक बाब आणि सामान्य अंशतः-व्युत्क्रम बाब वेगळ्या केल्या आहेत.\end{quote}
\end{prob}

\begin{prob}
$\Nat^3$ चे सांकेतीकरण करणारे फलन स्पष्टपणे द्या.
\end{prob}









\section{एक पर्यायी जोडीकरण फलन}\label{sfr:siz:pai-alt:sec}

\begin{explain}
$\Nat^2$ ची अशी इतर प्रगणनेही आहेत की त्यांची व्युत्क्रमे शोधणे अधिक
सोपे होते. त्यांपैकी एक पुढे दिले आहे. प्रगणन सरणीत दाखवण्याऐवजी,
सुरुवातीला रिकाम्या असलेल्या जागांशी जोडलेल्या धन पूर्णांकांच्या
यादीपासून सुरुवात करा. या जागा पुढीलप्रमाणे एकामागून एक
$\tuple{n,m}$ जोड्यांनी भरत आहोत, अशी कल्पना करा. पहिल्या स्थानी~$0$
असलेल्या जोड्यांपासून (म्हणजे $\tuple{0,m}$ जोड्यांपासून) सुरुवात करा.
पहिली जोडी (म्हणजे $\tuple{0,0}$) पहिल्या रिकाम्या जागेत ठेवा; मग एक
रिकामी जागा सोडून दुसरी जोडी पुढील रिकाम्या जागेत ठेवा आणि पुन्हा एक
जागा सोडा; ही प्रक्रिया अशीच पुढे चालू ठेवा. गोठवलेल्या इंग्रजी
मजकुरात दुसरी जोडी $\tuple{0,2}$ अशी दिली आहे; परंतु खालील सारणीनुसार
येथे $\langle 0,1\rangle$ अभिप्रेत आहे. आता आपल्या प्रगणनाची (अपूर्ण) सुरुवात अशी
दिसते:
\begin{quote}\small\textbf{स्रोतदुरुस्ती.} MRSIZ-003: गोठवलेल्या इंग्रजी मजकुरात एक रिकामी जागा सोडल्यानंतरची दुसरी जोडी \ensuremath{\langle}0,2\ensuremath{\rangle} अशी दिली आहे, पण लगेचची सारणी आणि वर्णनाची क्रमवार रचना ती \ensuremath{\langle}0,1\ensuremath{\rangle} असावी हे दाखवतात. सूत्र-संरचना जशीच्या तशी ठेवून मराठी मजकुरात अपेक्षित जोडी स्पष्ट केली आहे.\end{quote}
\[\begin{array}{@{}c c c c c c c c c c c@{}}
\mathbf 1 & \mathbf 2 & \mathbf 3 & \mathbf 4 & \mathbf 5 & \mathbf 6 & \mathbf 7 & \mathbf 8 & \mathbf 9 & \mathbf{10} & \dots \\ \\
\tuple{0,0} &  & \tuple{0,1} &  & \tuple{0,2} &  & \tuple{0,3} & & \tuple{0,4} &  & \dots \\
\end{array}
\]
अजून रिकाम्या राहिलेल्या जागांसाठी $\tuple{1,m}$ जोड्यांबाबत हीच
प्रक्रिया पुन्हा करा; पुन्हा प्रत्येक दुसरी रिकामी जागा वगळा:
\[\begin{array}{@{}c c c c c c c c c c c@{}}
\mathbf 1 & \mathbf 2 & \mathbf 3 & \mathbf 4 & \mathbf 5 & \mathbf 6 & \mathbf 7 & \mathbf 8 & \mathbf 9 & \mathbf{10} & \dots \\ \\
\tuple{0,0} & \tuple{1,0} & \tuple{0,1} &  & \tuple{0,2} & \tuple{1,1} &
\tuple{0,3} & & \tuple{0,4} &  \tuple{1,2} & \dots \\
\end{array}
\]
यानंतर $\tuple{2,m}$ जोड्या आणि मग $\tuple{3,m}$ जोड्या त्याच पद्धतीने
भरा; त्यापुढेही पहिला घटक अनुक्रमे 4, 5, इत्यादी करत पुढे जायचे आहे.
आपल्या पूर्ण प्रगणनाची सुरुवात म्हणून अशी दिसते:
\begin{quote}\small\textbf{स्रोतदुरुस्ती.} OLSIZ-003: गोठवलेल्या इंग्रजी मजकुरात “pairs \ensuremath{\langle}2,m\ensuremath{\rangle}, \ensuremath{\langle}2,m\ensuremath{\rangle}, etc.” अशी पुनरुक्ती आहे. पूर्ण सारणीतील \ensuremath{\langle}3,0\ensuremath{\rangle} आणि पुढील सर्व ओळींवरून दुसऱ्या ठिकाणी \ensuremath{\langle}3,m\ensuremath{\rangle} अभिप्रेत असल्याचे निश्चित होते; मराठी मजकुरात दुसरी जोडी त्यानुसार दुरुस्त केली आहे.\end{quote}
\[\begin{array}{@{}cc c c c c c c c c c@{}}
\mathbf 1 & \mathbf 2 & \mathbf 3 & \mathbf 4 & \mathbf 5 & \mathbf 6 & \mathbf 7 & \mathbf 8 & \mathbf 9 & \mathbf{10} & \dots \\ \\
\tuple{0,0} & \tuple{1,0} & \tuple{0,1} & \tuple{2,0}  & \tuple{0,2} &
\tuple{1,1} & \tuple{0,3} & \tuple{3,0}  & \tuple{0,4} &  \tuple{1,2} & \dots \\
\end{array}
\]
वरील सरणीतील चौकटींना या प्रगणनानुसार क्रमांक दिले, तर आपल्याला
नीट नागमोडी रेषा दिसणार नाही; त्याऐवजी पुढील मांडणी दिसेल:
\[\begin{array}{ c | c | c | c | c | c | c | c }
& \mathbf 0 & \mathbf 1 & \mathbf 2 & \mathbf 3 & \mathbf 4 & \mathbf 5 & \dots \\
\hline
\mathbf 0 & 1 & 3 & 5 & 7 & 9 & 11 & \dots \\
\hline
\mathbf 1 & 2 & 6 & 10 & 14 & 18 & \dots & \dots \\
\hline
\mathbf 2 & 4 & 12 & 20 & 28 & \dots & \dots & \dots \\
\hline
\mathbf 3 & 8 & 24 & 40 & \dots & \dots & \dots & \dots \\
\hline
\mathbf 4 & 16 & 48 & \dots & \dots & \dots & \dots & \dots \\
\hline
\mathbf 5 & 32 & \dots & \dots & \dots & \dots & \dots & \dots \\
\hline
\vdots & \vdots & \vdots & \vdots & \vdots & \vdots & \vdots & \ddots\\
\end{array}
\]
आपल्या प्रगणनात ओळ~$0$ मधील जोड्या विषम क्रमांकांच्या स्थानी आहेत,
हे दिसते; म्हणजे $\tuple{0,m}$ ही जोडी $2m+1$ या स्थानी आहे. दुसऱ्या
ओळीतील $\tuple{1,m}$ जोड्या विषम संख्येच्या दुप्पट क्रमांक असलेल्या
स्थानी आहेत; विशेषतः त्या $2 \cdot (2m+1)$ या स्थानी आहेत. तिसऱ्या
ओळीतील $\tuple{2,m}$ जोड्या विषम संख्येच्या चौपट क्रमांक असलेल्या
स्थानी, म्हणजे $4 \cdot (2m+1)$ या स्थानी आहेत; आणि ही रचना अशीच
पुढे चालते. प्रत्येक ओळीत $(2m+1)$ च्या गुणकांचे मूल्य $1$, $2$,
$4$, $8$, \dots{} असे आहे; हे नेमके~$2$ चे घात आहेत:
$1= 2^0$, $2 = 2^1$, $4 = 2^2$, $8 = 2^3$, \dots\@ विचाराधीन
जोडीचा पहिला घटक हाच नेहमी संबंधित घातांक असतो. म्हणून
$\tuple{n,m}$ या जोडीसाठी गुणक $2^n$ आहे. यातून आपल्याला
$2^n \cdot (2m+1)$ हे सर्वसाधारण सूत्र मिळते. पण या फलनाने जोड्या
\emph{धन} पूर्णांकांकडे पाठवल्या जातात; म्हणजे $\tuple{0,0}$ चे
स्थान~$1$ आहे. स्थानांची सुरुवात~$0$ पासून करायची असेल, तर उत्तरातून
$1$ वजा करावा लागतो. त्यामुळे पुढील फलन मिळते:
\end{explain}

\begin{ex}
$h\colon \Nat^2 \to \Nat$ हे पुढीलप्रमाणे दिलेले फलन घ्या:
\[h(n,m) = 2^n (2m+1) - 1
\]
हे नैसर्गिक संख्यांच्या जोड्यांचा संच~$\Nat^2$ यासाठीचे जोडीकरण फलन
आहे.
\end{ex}

\begin{explain}
त्यानुसार, $\Nat^2$ च्या आपल्या दुसऱ्या प्रगणनात $\tuple{0,0}$ या
जोडीचा संकेतांक $h(0,0) = 2^0(2\cdot 0+1) - 1 = 0$ आहे;
$\tuple{1,2}$ हिचा संकेतांक $2^{1} \cdot (2 \cdot 2 + 1) - 1 = 2
\cdot 5 - 1 = 9$ आहे; आणि $\tuple{2,6}$ हिचा संकेतांक $2^{2} \cdot (2
\cdot 6 + 1) - 1 = 51$ आहे.
\end{explain}

कधी कधी नैसर्गिक संख्यांच्या जोड्यांचे~$\Nat^2$ सांकेतीकरण करताना ते
आच्छादक असण्याची आवश्यकता नसते. अशा सांकेतीकरणांची व्युत्क्रमे केवळ
अंशतः फलने असतात.

\begin{ex}
$j\colon \Nat^2 \to \Nat^+$ हे पुढीलप्रमाणे दिलेले फलन घ्या:
\[j(n,m) = 2^n3^m
\]
हे $\Nat^2 \to \Nat$ असे एकास-एक फलन आहे.
\end{ex}









\section{अगणनीय संच}\label{sfr:siz:nen:sec}

\begin{editorial}
  या विभागात \ref{sfr:siz:enm:sec} मधील व्याख्या वापरून $\Bin^\omega$ आणि
  $\Pow{\PosInt}$ यांची अगणनीयता सिद्ध केली आहे. ही मांडणी
  \ref{sfr:siz:enm-alt:sec} मधील आवृत्तीपेक्षा थोडी अधिक प्राथमिक आणि
  थोडी अधिक तपशीलवार ठेवली आहे.
\end{editorial}

धन पूर्णांकांचा संच~$\PosInt$ यासारखे काही संच अनंत असतात. आतापर्यंत
पाहिलेल्या अनंत संचांची सर्व उदाहरणे गणनीय होती. पण हा गुणधर्म
नसलेले अनंत संचही असतात. अशा संचांना \emph{अगणनीय} म्हणतात.

सुरुवातीलाच, अगणनीय संच असतात ही गोष्ट कदाचित आश्चर्यकारक
वाटेल. कोणत्याही गणनीय संच~$A$ साठी $f \colon \PosInt \to A$
असे आच्छादक फलन असते. संच अगणनीय असेल, तर असे फलन
असू शकत नाही. म्हणजे, $\PosInt$ मधील अनंत घटक ना~$A$ कडे
पाठवणारे कोणतेही फलन~$A$ मधील सर्व घटक व्यापू शकत नाही. या अर्थाने,
अनंत असलेल्या धन पूर्णांकांपेक्षा~$A$ मध्ये ``अधिक'' घटक असतात.

एखादा संच अगणनीय आहे हे कसे सिद्ध करायचे? तसे कोणतेही
आच्छादक फलन अस्तित्वात असू शकत नाही, हे दाखवावे लागते. समतुल्यपणे,
$A$ मधील घटकांचे एकदिश अनंत यादीत प्रगणन करता येत नाही, हे दाखवावे
लागते. हे करण्याचा सर्वोत्तम मार्ग म्हणजे~$A$ मधील घटक च्या
प्रत्येक यादीतून किमान एक घटक सुटतो, किंवा $f\colon \PosInt \to A$
असे कोणतेही फलन आच्छादक असू शकत नाही, हे दाखवणे. यासाठी कँटर
यांची \emph{कर्णरेषीय पद्धत} वापरता येते. $A$ मधील घटक ची
एखादी यादी, म्हणजे $x_1$, $x_2$, \dots, दिली असता, आपण~$A$ मधील असा
दुसरा घटक रचतो की त्याच्या रचनेमुळे तो त्या यादीत असणे अशक्य ठरते.

आपले पहिले उदाहरण म्हणजे $0$ आणि $1$ यांच्या सर्व अनंत व मध्ये खंड
नसलेल्या अनुक्रमांचा संच~$\Bin^\omega$.

\begin{thm}
\label{sfr:siz:nen:thm:nonenum-bin-omega}
$\Bin^\omega$ हा अगणनीय आहे.
\end{thm}

\begin{proof}
व्याघाताने सिद्ध करण्यासाठी असे समजा की $\Bin^\omega$ हा गणनीय आहे; म्हणजे,
$\Bin^\omega$ मधील सर्व घटक ची $s_{1}$, $s_{2}$, $s_{3}$,
$s_{4}$, \dots{} अशी यादी आहे, असे समजा. यांपैकी प्रत्येक $s_i$ हा
स्वतः $0$ आणि~$1$ यांचा अनंत अनुक्रम आहे. या यादीतील $i$ व्या अनुक्रमाच्या
$j$ व्या घटकाला $s_i(j)$ म्हणू. मग $i$ वा अनुक्रम~$s_i$ असा आहे:
\[s_i(1), s_i(2), s_i(3), \dots
\]

ही यादी आणि तिच्यातील प्रत्येक अनुक्रम~$s_i$ चे घटक आपण पुढील सरणीत
मांडू शकतो:
\[\begin{array}{c|c|c|c|c|c}
& 1 & 2 & 3 & 4 & \dots \\\hline
1 & \mathbf{s_{1}(1)} & s_{1}(2) & s_{1}(3) & s_1(4) & \dots \\\hline
2 & s_{2}(1)& \mathbf{s_{2}(2)} & s_2(3) & s_2(4) & \dots \\\hline
3 & s_{3}(1)& s_{3}(2) & \mathbf{s_3(3)} & s_3(4) & \dots \\\hline
4 & s_{4}(1)& s_{4}(2) & s_4(3) & \mathbf{s_4(4)} & \dots \\\hline
\vdots & \vdots & \vdots & \vdots & \vdots & \mathbf{\ddots}
\end{array}
\]
डाव्या बाजूकडील खुणा यादीतील अनुक्रमांना $s_1$, $s_2$, \dots{} असे
क्रमांक देतात; वरच्या बाजूकडील संख्या प्रत्येक अनुक्रमातील घटक
ना क्रमांक देतात. उदाहरणार्थ, $s_{1}(1)$ हे अनुक्रम~$s_{1}$ मधील
पहिला घटक असलेल्या संख्येचे नाव आहे; ती संख्या $0$ किंवा~$1$
असू शकते. पुढेही याचप्रमाणे अर्थ घ्यायचा.

आता आपण $0$ आणि $1$ यांचा असा अनंत अनुक्रम~$\overline{s}$ रचू, जो
या यादीत असणे अशक्य आहे. $\overline{s}$ ची व्याख्या $s_1$, $s_2$,
\dots{} या यादीवर अवलंबून असेल. $0$ आणि $1$ यांच्या अनंत अनुक्रमांची
कोणतीही अनंत यादी असा अनंत अनुक्रम~$\overline{s}$ निर्माण करते, जो
त्या यादीत नसेल याची खात्री असते.

$\overline{s}$ ची व्याख्या करण्यासाठी आपण त्याचे सर्व घटक
ठरवतो; म्हणजे, प्रत्येक $n \in \PosInt$ साठी $\overline{s}(n)$
ठरवतो. वरील सरणीच्या कर्णरेषेवरून खाली वाचून (म्हणूनच ``कर्णरेषीय
पद्धत'' हे नाव) प्रत्येक $1$ चे $0$ मध्ये आणि प्रत्येक $0$ चे~$1$
मध्ये रूपांतर करून आपण हे करतो. अधिक अमूर्तपणे, कर्णरेषेवरील $n$ वा
घटक, म्हणजे $s_n(n)$, हा $1$ आहे की $0$ यानुसार
$\overline{s}(n)$ हे $0$ किंवा $1$ असे परिभाषित करतो:
\[\overline{s}(n) =
\begin{cases}
1 & \text{जर $s_{n}(n) = 0$ असेल}\\
0 & \text{जर $s_{n}(n) = 1$ असेल}.
\end{cases}
\]
प्रकरणांनुसार व्याख्येपेक्षा सूत्रे अधिक आवडत असतील, तर
$\overline{s}(n) = 1 - s_n(n)$ अशीही व्याख्या करता येईल.

$\overline{s}$ हा स्पष्टपणे $0$ आणि $1$ यांचा अनंत अनुक्रम आहे, कारण
तो आपल्या सरणीच्या कर्णरेषेवर दिसणाऱ्या $0$ आणि $1$ यांच्या अनुक्रमाचा
पूरक अनुक्रम आहे. म्हणून $\overline{s}$ हा $\Bin^\omega$ चा
घटक आहे. पण तो $s_1$, $s_2$, \dots{} या यादीत असू शकत नाही.
असे का?

तो यादीतील पहिला अनुक्रम~$s_1$ असू शकत नाही, कारण त्याचा पहिला
घटक हा $s_1$ च्या पहिल्या घटकापेक्षा वेगळा आहे. $s_1(1)$
काहीही असो, $\overline{s}(1)$ त्याच्या विरुद्ध असेल अशी आपण व्याख्या
केली. तो यादीतील दुसरा अनुक्रमही असू शकत नाही, कारण $\overline{s}$
चा दुसरा घटक $s_2$ च्या दुसऱ्या घटकापेक्षा वेगळा आहे: $s_2(2)$ हा
$0$ असेल, तर $\overline{s}(2)$ हा $1$ आहे, आणि उलट. पुढेही असेच.

अधिक नेमकेपणाने सांगायचे तर, $\overline{s}$ यादीत असता, तर कोणत्या
तरी $k$ साठी $\overline{s} = s_{k}$ असते. दोन अनुक्रम प्रत्येक स्थानी
जुळतात तेव्हा आणि केवळ तेव्हाच ते एकरूप असतात; म्हणजे, कोणत्याही~$n$
साठी $\overline{s}(n) = s_{k}(n)$ असते. म्हणून, विशेष प्रकरण म्हणून
$n = k$ घेतल्यास $\overline{s}(k) = s_{k}(k)$ असणे आवश्यक ठरते.
$s_k(k)$ हे $0$ किंवा~$1$ आहे. ते $0$ असेल, तर $\overline{s}(k)$ हे~$1$
असले पाहिजे---कारण आपण $\overline{s}$ ची तशीच व्याख्या केली आहे. पण
$s_k(k) = 1$ असेल, तर पुन्हा $\overline{s}$ च्या व्याख्येमुळे
$\overline{s}(k) = 0$ होते. दोन्ही बाबतीत
$\overline{s}(k) \neq s_{k}(k)$.

सुरुवातीला आपण $\Bin^\omega$ मधील घटक ची $s_1$, $s_2$,
\dots{} अशी यादी आहे असे गृहीत धरले. या यादीपासून आपण असा
अनुक्रम~$\overline{s}$ रचला की तो यादीत असू शकत नाही, हे सिद्ध केले.
पण सर्व $s_i$ हे $0$ आणि $1$ यांचे अनुक्रम असतील, तर तो नक्कीच $0$
आणि $1$ यांचा अनुक्रम असतो; म्हणजे, $\overline{s} \in
\Bin^\omega$. यावरून विशेषतः $\Bin^\omega$ मधील \emph{सर्व}
घटक ची यादी असू शकत नाही. कारण अशी कोणतीही यादी दिली, तरी
तिच्यात नसेल असा अनुक्रम~$\overline{s}$ आपण रचू शकतो. म्हणून
$\Bin^\omega$ मधील सर्व अनुक्रमांची यादी आहे हे गृहीतक व्याघात निर्माण
करते.
\end{proof}

\begin{explain}
या सिद्धता-पद्धतीला ``कर्णीकरण'' म्हणतात, कारण तिच्यात सरणीच्या
कर्णरेषेचा वापर करून~$\overline{s}$ ची व्याख्या केली जाते. कर्णीकरणात
सरणी असलीच पाहिजे असे नाही: प्रत्यक्ष सरणी किंवा कर्णरेषा नसतानाही
हीच कल्पना वापरून संच गणनीय नाहीत, हे दाखवता येते.
\end{explain}

\begin{thm}
\label{sfr:siz:nen:thm:nonenum-pownat}
$\Pow{\PosInt}$ हा गणनीय नाही.
\end{thm}

\begin{proof}
आपण त्याच पद्धतीने पुढे जाऊ: $\PosInt$ च्या उपसंचांची प्रत्येक यादी
दिली असता त्या यादीत नसलेला $\PosInt$ चा एक उपसंच असतो, हे दाखवू.
$\PosInt$ च्या उपसंचांची पुढील यादी दिली आहे असे समजा:
\[Z_{1}, Z_{2}, Z_{3}, \dots
\]
आता आपण $\overline{Z}$ हा असा संच परिभाषित करतो की प्रत्येक
$n \in \PosInt$ साठी $n \in \overline{Z}$ तेव्हा आणि केवळ तेव्हाच,
जेव्हा $n \notin Z_{n}$:
\[\overline{Z} = \Setabs{n \in \PosInt}{n \notin Z_n}
\]
$\overline{Z}$ हा स्पष्टपणे धन पूर्णांकांचा संच आहे, कारण गृहीतकानुसार
प्रत्येक~$Z_n$ तसा आहे; म्हणून $\overline{Z} \in
\Pow{\PosInt}$. पण $\overline{Z}$ यादीत असू शकत नाही. हे दाखवण्यासाठी
प्रत्येक $k \in \PosInt$ साठी $\overline{Z} \neq Z_k$, हे सिद्ध करू.

आता $k \in \PosInt$ स्वैर असू द्या. प्रत्येक $n \in \PosInt$ साठी
$n \in \overline{Z}$ तेव्हा आणि केवळ तेव्हाच, जेव्हा $n \notin Z_n$,
अशी $\overline{Z}$ ची व्याख्या केली आहे. विशेषतः $n=k$ घेतल्यास,
$k \in \overline{Z}$ तेव्हा आणि केवळ तेव्हाच, जेव्हा $k \notin Z_k$.
पण यावरून $\overline{Z} \neq Z_k$ हे दिसते, कारण $k$ हा एकाचा
घटक आहे आणि दुसऱ्याचा नाही; म्हणून $\overline{Z}$ आणि $Z_k$
यांचे घटक वेगवेगळे आहेत. $k$ स्वैर असल्याने $\overline{Z}$ हा
$Z_1$, $Z_2$, \dots{} या यादीत नाही.
\end{proof}

\begin{explain}
मागील सिद्धतेत कर्णरेषेचा उल्लेख झाला नाही; पण पुढील चित्र मनात आणले,
तर तिच्यात कर्णरेषा आहे असे समजता येते. $Z_1$, $Z_2$, \dots{} हे संच
एका सरणीत लिहिले आहेत अशी कल्पना करा; प्रत्येक घटक~$j \in Z_i$
हा $j$ व्या स्तंभात लिहिला आहे. समजा, त्या यादीतील पहिले चार संच
$\{1,2,3,\dots\}$, $\{2, 4, 6, \dots\}$, $\{1,2,5\}$ आणि
$\{3,4,5,\dots\}$ आहेत. मग सरणीची सुरुवात अशी होईल:
\[\begin{array}{r@{}rrrrrrr}
  Z_1 = \{ & \mathbf{1}, & 2, & 3, & 4, & 5, & 6, & \dots\}\\
  Z_2 = \{ &  & \mathbf{2}, &  & 4, &  & 6, & \dots\}\\
  Z_3 = \{ & 1, & 2, &  &  & 5\phantom{,} &  & \}\\
  Z_4 = \{ &  &  & 3, & \mathbf{4}, & 5, & 6, & \dots\}\\
  \vdots & & & & & \ddots
\end{array}
\]
यानंतर कर्णरेषेवरून खाली जाताना कर्णरेषेवर दिसणाऱ्या संख्या वगळून,
आणि $j$ व्या ओळीच्या त्याच क्रमांकाच्या स्तंभातील जागा रिकामी असेल अशा~$j$
चा समावेश करून, $\overline{Z}$ हा संच मिळतो. वरील उदाहरणात
आपण $1$ आणि $2$ वगळू, $3$ चा समावेश करू, $4$ वगळू, इत्यादी.
\end{explain}

\begin{prob}
कर्णरेषीय युक्तिवादाने $\Pow{\Nat}$ हा अगणनीय आहे, हे दाखवा.
\end{prob}

\begin{prob}\label{sfr:siz:nen:prob:f-posint}
$f \colon \PosInt \to \PosInt$ अशा फलनांचा संच अगणनीय
आहे, हे स्पष्ट कर्णरेषीय युक्तिवादाने दाखवा. म्हणजे, $f_1$, $f_2$,
\dots{} ही फलनांची यादी असेल आणि प्रत्येक $f_i\colon
\PosInt \to \PosInt$ असेल, तर या यादीत नसलेले कोणते तरी
$\overline{f}\colon \PosInt \to \PosInt$ असते, हे दाखवा.
\end{prob}









\section{न्यूनीकरण}\label{sfr:siz:red:sec}

\begin{editorial}
  या विभागात \ref{sfr:siz:nen:sec} मधील निष्कर्षांशी जुळणारी अगणनीयता
  न्यूनीकरणाने सिद्ध केली आहे. \ref{sfr:siz:nen-alt:sec} मधील निष्कर्षांशी
  जुळणारी, किंचित अधिक संक्षिप्त पर्यायी आवृत्ती \ref{sfr:siz:red-alt:sec}
  मध्ये दिली आहे.
\end{editorial}

कर्णीकरणाच्या युक्तिवादाने $\Pow{\PosInt}$ हा अगणनीय आहे,
हे आपण दाखवले. सर्व अनंत $0$--$1$ अनुक्रमांचा संच~$\Bin^\omega$ हा
अगणनीय आहे, याची सिद्धता आपल्याकडे आधीच होती.
$\Pow{\PosInt}$ हा अगणनीय आहे हे सिद्ध करण्याचा आणखी एक मार्ग
असा: \emph{$\Pow{\PosInt}$ हा गणनीय असेल, तर $\Bin^\omega$
  हाही गणनीय असतो} हे दाखवा. $\Bin^\omega$ हा गणनीय
नाही, हे माहीत असल्याने $\Pow{\PosInt}$ हाही तसा असू शकत नाही. यालाच
एक समस्या दुसऱ्या समस्येकडे \emph{न्यूनीत करणे} म्हणतात---या बाबतीत
$\Bin^\omega$ चे प्रगणन करण्याची समस्या आपण $\Pow{\PosInt}$ चे प्रगणन
करण्याच्या समस्येकडे न्यूनीत करतो. नंतरच्या समस्येचे उत्तर---म्हणजे
$\Pow{\PosInt}$ चे प्रगणन---पहिल्या समस्येचे उत्तर---म्हणजे
$\Bin^\omega$ चे प्रगणन---देईल.

संच~$B$ च्या प्रगणनाची समस्या संच~$A$ च्या प्रगणनाच्या समस्येकडे कशी
न्यूनीत करायची? $A$ चे प्रगणन $B$ च्या प्रगणनात रूपांतरित करण्याची
पद्धत आपण देतो. त्यासाठीचा सर्वांत सोपा मार्ग म्हणजे
$f\colon A \to B$ असे आच्छादक फलन परिभाषित करणे. $x_1$, $x_2$,
\dots{} हे~$A$ चे प्रगणन असेल, तर $f(x_1)$, $f(x_2)$, \dots{} हे~$B$
चे प्रगणन करील. आपल्या बाबतीत आपण
$f\colon \Pow{\PosInt} \to \Bin^\omega$ असे आच्छादक फलन शोधत आहोत.

\begin{prob}
$g\colon B \to A$ असे एकास-एक फलन असेल आणि $B$ हा
अगणनीय असेल, तर $A$ हाही अगणनीय असतो, हे दाखवा.
$g$ वापरून~$A$ चे प्रगणन~$B$ च्या प्रगणनात कसे रूपांतरित करता येते,
हे दाखवून सिद्धता द्या.
\end{prob}

\begin{proof}[न्यूनीकरणाने {\ref{sfr:siz:nen:thm:nonenum-pownat}} ची सिद्धता]
$\Pow{\PosInt}$ हा गणनीय आहे, आणि म्हणून त्याचे $Z_{1}$,
$Z_{2}$, $Z_{3}$, \dots{} असे प्रगणन आहे, असे समजा.

$f \colon \Pow{\PosInt} \to \Bin^\omega$ हे फलन पुढीलप्रमाणे परिभाषित
करा: $f(Z)$ हा असा अनुक्रम~$s$ असू द्या की $n \in Z$ असेल तेव्हा
आणि केवळ तेव्हाच $s(n) = 1$, आणि अन्यथा $s(n) = 0$. याने
फलन स्पष्टपणे परिभाषित होते, कारण $Z \subseteq \PosInt$ असताना कोणताही
$n \in \PosInt$ हा $Z$ चा घटक असतो किंवा नसतो. उदाहरणार्थ, धन सम संख्यांचा संच
$2\PosInt = \{2, 4, 6, \dots\}$ हा $010101\dots$ या अनुक्रमाकडे,
रिक्त संच $0000\dots$ कडे आणि $\PosInt$ हा संच स्वतः $1111\dots$ कडे
पाठवला जातो.
\begin{quote}\small\textbf{स्रोतदुरुस्ती.} OLSIZ-004: गोठवलेल्या इंग्रजी व्याख्येत f(Z) ला s\_k असे म्हटले आहे, पण k चा Z शी संबंध किंवा निवड दिलेली नाही. अभिप्रेत अनुक्रम Z नेच ठरतो; मराठी व्याख्येत एकच बद्ध निर्गत नाव s आणि त्याचे घटक s(n) सातत्याने वापरले आहेत.\end{quote}

हे फलन आच्छादक सुद्धा आहे: $0$ आणि $1$ यांचा प्रत्येक अनुक्रम धन
पूर्णांकांच्या कोणत्या तरी संचाशी जुळतो; तो संच म्हणजे अनुक्रमात ज्या
स्थानी~$1$ आहे त्या स्थानांना अनुरूप पूर्णांकांचा संच. अधिक नेमकेपणाने,
$s \in \Bin^\omega$ असे समजा. पुढीलप्रमाणे $Z \subseteq \PosInt$
परिभाषित करा:
\[Z = \Setabs{n \in \PosInt}{s(n) = 1}
\]
मग $f$ च्या व्याख्येवरून तपासता येते की $f(Z) = s$.

आता पुढील यादी विचारात घ्या:
\[f(Z_1), f(Z_2), f(Z_3), \dots
\]
$f$ हे आच्छादक असल्याने $\Bin^\omega$ मधील प्रत्येक सदस्य हा
$f$ च्या कोणत्या तरी फलसाधकावरील मूल्य असला पाहिजे, आणि म्हणून तो
यादीत आला पाहिजे. त्यामुळे ही यादी $\Bin^\omega$ मधील सर्व घटकांचे
प्रगणन करते.

म्हणून $\Pow{\PosInt}$ हा गणनीय असता, तर $\Bin^\omega$ हाही
गणनीय असता. पण $\Bin^\omega$ हा अगणनीय आहे
(\ref{sfr:siz:nen:thm:nonenum-bin-omega}). त्यामुळे $\Pow{\PosInt}$ हा
अगणनीय आहे.
\end{proof}

\begin{explain}
न्यूनीकरण कोणत्या दिशेने केले जाते याबद्दल गोंधळ होणे सोपे आहे.
उदाहरणार्थ, $g \colon \Bin^\omega \to B$ असे आच्छादक फलन
$B$ हा अगणनीय आहे, हे \emph{सिद्ध करत नाही}. ($g
\colon \Bin^\omega \to \Bin$ हे $g(s) = s(1)$ ने परिभाषित केलेले,
$0$ आणि $1$ यांच्या अनुक्रमाला त्याच्या पहिल्या घटक कडे
पाठवणारे फलन विचारात घ्या. काही अनुक्रम $0$ ने आणि काही $1$ ने सुरू
होत असल्याने ते आच्छादक आहे. पण $\Bin$ हा सांत संच आहे.)
तसेच फलन~$f$ हे आच्छादक असलेच पाहिजे; अन्यथा युक्तिवाद लागू
होत नाही: $f(x_1)$, $f(x_2)$, \dots{} यात~$B$ मधील सर्व घटक
असतील याची खात्री उरत नाही. उदाहरणार्थ,
\[h(n) = \underbrace{000\dots0}_{\text{$n$ वेळा $0$}}111\dots
\]
याने $h\colon \PosInt \to
\Bin^\omega$ असे फलन परिभाषित होते: प्रत्येक मूल्याची सुरुवात n शून्यांनी
होते आणि त्यानंतर कायम एक येतात. या फलनाच्या व्याप्तीत अखेरीस कायम एक
न होणारे अनुक्रम येत नाहीत, म्हणून ते आच्छादक नाही; पण $\PosInt$ हा
गणनीय आहे.
\begin{quote}\small\textbf{स्रोतदुरुस्ती.} OLSIZ-005: गोठवलेल्या इंग्रजी प्रदर्शनात h(n) चे मूल्य फक्त n शून्यांची सांत चिन्हमाला आहे, जी अनंत अनुक्रमांच्या Bin-omega संचात येत नाही. मराठी प्रदर्शनात त्यानंतर अनंत एकांची शेपटी जोडून वैध, अनाच्छादक फलन दिले आहे.\end{quote}
\end{explain}

\begin{prob}
धन पूर्णांकांच्या जोड्यांच्या \emph{सर्व संचांचा} संच न्यूनीकरणाच्या
युक्तिवादाने अगणनीय आहे, हे दाखवा.
\end{prob}

\begin{prob}\label{sfr:siz:red:prob:nat-nat}
  $f\colon \Nat \to \Nat$ अशा सर्व फलनांचा संच~$X$ हा न्यूनीकरणाच्या
  युक्तिवादाने अगणनीय आहे, हे दाखवा. (सूचना: $X$ पासून
  $\Bin^\omega$ कडे जाणारे आच्छादक फलन द्या.)
\end{prob}

\begin{prob}
नैसर्गिक संख्यांच्या सर्व अनंत अनुक्रमांचा संच~$\Nat^\omega$ हा
न्यूनीकरणाच्या युक्तिवादाने अगणनीय आहे, हे दाखवा.
\end{prob}

\begin{prob}
$P$ हा धन पूर्णांकांच्या संचापासून $\{0\}$ या संचाकडे जाणाऱ्या फलनांचा
संच असू द्या; आणि $Q$ हा धन पूर्णांकांच्या संचापासून $\{0\}$ या
संचाकडे जाणाऱ्या \emph{अंशतः} फलनांचा संच असू द्या. $P$ हा
गणनीय आहे आणि $Q$ नाही, हे दाखवा. (सूचना: $\Bin^\omega$ चे
प्रगणन करण्याची समस्या~$Q$ चे प्रगणन करण्याच्या समस्येकडे न्यूनीत करा.)
\end{prob}

\begin{prob}
$S$ हा धन पूर्णांकांच्या संचापासून \{0,1\} या संचाकडे जाणाऱ्या सर्व
आच्छादक फलनांचा संच असू द्या; म्हणजे, $S$ मध्ये सर्व
आच्छादक~$f\colon \PosInt \to \Bin$ आहेत. $S$ हा
अगणनीय आहे, हे दाखवा.
\end{prob}

\begin{prob}
सर्व वास्तव संख्यांचा संच~$\Real$ हा अगणनीय आहे, हे दाखवा.
\end{prob}









\section{तुल्यबलता}\label{sfr:siz:equ:sec}

संचांच्या ``आकारमाना''ची एक अंतःप्रज्ञ कल्पना आपल्याकडे आहे आणि ती
सांत संचांसाठी व्यवस्थित चालते. पण अनंत संचांचे काय? \emph{कोणत्याही}
आकारमानाच्या दोन संचांच्या आकारमानांची औपचारिक रीतीने तुलना करायची
असेल, तर दोन संच एकाच आकारमानाचे कधी असतात याची व्याख्या प्रथम करणे
योग्य ठरते. फ्रेगे यांचे पुढील उदाहरण पाहा:
\begin{quote}
  उपाहारगृहातील कर्मचाऱ्याला मेजावर ताटांइतक्याच सुऱ्या मांडल्या आहेत
  याची खात्री करायची असेल, तर प्रत्येक ताटाच्या उजवीकडे एक सुरी ठेवून
  आणि मेजावरील प्रत्येक सुरी कोणत्या तरी ताटाच्या उजवीकडे येईल असे
  केल्यास त्याला ताटे किंवा सुऱ्या मोजण्याची गरज नसते. अशा प्रकारे
  ताटे आणि सुऱ्या एकास-एक संगतीने परस्परांशी जोडली जातात, आणि तीही
  त्याच अवकाशीय संबंधाने. (Frege, 1884, \S70)
\end{quote}
या उताऱ्यातील अंतर्दृष्टी पुढील औपचारिक व्याख्येत मांडता येते:

\begin{defn}\label{sfr:siz:equ:comparisondef}
  $A$ हा $B$ शी \emph{तुल्यबल} असतो, आणि ते $\cardeq{A}{B}$ असे लिहितो,
  तेव्हा आणि केवळ तेव्हाच, जेव्हा $f \colon A \to B$ असे
  एकास-एक आच्छादन असते.
\end{defn}

\begin{prop}\label{sfr:siz:equ:equinumerosityisequi}
तुल्यबलता हा तुल्यता संबंध आहे.
\end{prop}

\begin{proof}
तुल्यबलता परावर्ती, सममित आणि संक्रमक आहे, हे दाखवावे लागेल. $A, B$
आणि $C$ हे संच असू द्या.

\emph{परावर्तिता.} प्रत्येक $x \in A$ साठी $\Id{A} (x) = x$ असे
असलेले एकरूपता फलन $\Id{A} \colon A \to A$ हे एकास-एक आच्छादन आहे.
म्हणून $\cardeq{A}{A}$.

\emph{सममितता.} $\cardeq{A}{B}$ आहे, म्हणजे $f\colon A \to B$ असे
एकास-एक आच्छादन आहे, असे समजा. $f$ हे एकास-एक व आच्छादक असल्याने त्याचे
व्युत्क्रम $f^{-1}$ अस्तित्वात आहे आणि तेही एकास-एक व आच्छादक आहे. म्हणून
$f^{-1}\colon B \to A$ हे एकास-एक आच्छादन आहे; त्यामुळे
$\cardeq{B}{A}$.

\emph{संक्रमकता.} $\cardeq{A}{B}$ आणि $\cardeq{B}{C}$ आहे, म्हणजे
$f\colon A \to B$ आणि $g\colon B \to
C$ अशी एकास-एक आच्छादने आहेत, असे समजा. मग संयोजन
$\comp{f}{g}\colon A \to C$ हे एकास-एक व आच्छादक आहे; म्हणून
$\cardeq{A}{C}$.
\end{proof}

\begin{prop}
$\cardeq{A}{B}$ असेल, तर $A$ हा गणनीय असतो तेव्हा आणि केवळ
तेव्हाच $B$ तसा असतो.
\end{prop}

\begin{editorial}
पुढील सिद्धतेत \ref{sfr:siz:enm:sec} समाविष्ट असेल, तर
\ref{sfr:siz:enm:defn:enumerable} वापरली जाते; अन्यथा
\ref{sfr:siz:enm-alt:defn:enumerable} वापरली जाते.
\end{editorial}

\begin{proof}
$\cardeq{A}{B}$ आहे, म्हणून $f \colon A
\to B$ असे काही एकास-एक आच्छादन आहे, आणि $A$ हा गणनीय आहे,
असे समजा.

  मग एकतर $A = \emptyset$ असतो, किंवा $g\colon \PosInt \to A$ असे
  आच्छादक फलन असते. $A = \emptyset$ असेल, तर $B = \emptyset$
  सुद्धा असतो (अन्यथा एखादा घटक~$y \in B$ असेल,
  पण $x \in A$ असा कोणताही घटक नसल्याने $f(x) = y$ अशी पूर्वप्रतिमा
  मिळणार नाही). दुसरीकडे, $g\colon \PosInt \to A$ हे आच्छादक
  असेल, तर $\comp{g}{f} \colon \PosInt \to B$ हे आच्छादक आहे.
  हे पाहण्यासाठी $y \in B$ असू द्या. $f$ हे आच्छादक असल्याने
  $f(x) = y$ असा एखादा $x \in A$ घटक असतो. $g$ हे आच्छादक
  असल्याने $g(n) =
  x$ अशी एखादी $n \in \PosInt$ संख्या असते. त्यामुळे
\[(\comp{g}{f})(n) = f(g(n)) = f(x) = y
\]
  आणि म्हणून $\comp{g}{f}$ हे आच्छादक आहे. $\comp{g}{f}$ हे~$B$
  चे प्रगणन आहे; त्यामुळे $B$ हा गणनीय आहे.
\begin{quote}\small\textbf{स्रोतदुरुस्ती.} OLSIZ-006: गोठवलेल्या इंग्रजी सिद्धतेच्या दोन्ही पर्यायी शाखांतील रिक्त-संच प्रकरणात g(x)=y असे लिहिले आहे. त्या शाखेत g अजून अस्तित्वात नसते; A पासून B पर्यंतचे आधी दिलेले द्विएक फलन f वापरूनच B रिक्त असल्याचा निष्कर्ष मिळतो. मराठी मजकुरात दोन्ही ठिकाणी f(x)=y केले आहे.\end{quote}
\begin{quote}\small\textbf{स्रोतदुरुस्ती.} MRSIZ-004: गोठवलेल्या पर्यायी शाखेत एकाच प्रांताला आधी “initial sequence of natural numbers” आणि नंतर “initial segment” म्हटले आहे. आधीच्या औपचारिक व्याख्येशी सुसंगत म्हणून मराठीत दोन्हीकडे “आरंभीचा खंड” वापरला आहे.\end{quote}

$B$ हा गणनीय असेल, तर~$f$ ऐवजी
एकास-एक आच्छादन $f^{-1}\colon B \to A$ घेऊन तोच युक्तिवाद पुन्हा केल्यास
$A$ हा गणनीय आहे, हे मिळते.
\end{proof}

\begin{prob}
$\cardeq{A}{C}$ आणि $\cardeq{B}{D}$ असेल, आणि $A \cap B =
C \cap D = \emptyset$ असेल, तर $\cardeq{A \cup B}{C \cup D}$ आहे,
हे दाखवा.
\end{prob}

\begin{prob}
$A$ हा अनंत आणि गणनीय असेल, तर $\cardeq{A}{\Nat}$ आहे, हे
दाखवा.
\end{prob}









\section{वेगवेगळ्या आकारमानांचे संच आणि कँटर यांचे प्रमेय}\label{sfr:siz:car:sec}

\begin{explain}
दोन संचांचे आकारमान समान असते या कल्पनेचे नेमके विधान आपण दिले आहे.
एक संच दुसऱ्यापेक्षा लहान असतो या कल्पनेचेही नेमके विधान करता येते.
``लहान (किंवा तुल्यबल) असणे'' याच्या व्याख्येत संचांदरम्यानचे
एकास-एक आच्छादन न घेता पहिल्या संचापासून दुसऱ्या संचाकडे जाणारे
एकास-एक फलन घ्यावे लागेल. असे फलन अस्तित्वात असेल, तर पहिल्या
संचाचे आकारमान दुसऱ्या संचाच्या आकारमानापेक्षा लहान किंवा त्याच्याइतके
असते. अंतःप्रज्ञेने, एका संचापासून दुसऱ्या संचाकडे जाणारे
एकास-एक फलन हे फलनाच्या व्याप्तीत प्रांताइतके तरी घटक आहेत
याची खात्री देते, कारण प्रांतातील कोणतेही दोन घटक व्याप्तीतील
एकाच घटक कडे पाठवले जात नाहीत.
\end{explain}

\begin{defn}
$A$ हा $B$ पेक्षा \emph{आकारमानाने मोठा नाही}, हे $\cardle{A}{B}$ असे
लिहितो, तेव्हा आणि केवळ तेव्हाच, जेव्हा $f \colon A \to B$ असे
एकास-एक फलन असते.
\end{defn}

हा संबंध परावर्ती आणि संक्रमक आहे, पण सममित नाही, हे स्पष्ट आहे (हे
सरावासाठी सोडले आहे). एक संच दुसऱ्या संचापेक्षा (काटेकोरपणे) लहान आहे
असे सांगणारी कल्पनाही मांडता येते.

\begin{defn}
$A$ हा $B$ पेक्षा \emph{आकारमानाने लहान} आहे, हे $\cardless{A}{B}$ असे
लिहितो, तेव्हा आणि केवळ तेव्हाच, जेव्हा $f\colon A \to B$ असे
एकास-एक फलन असते, पण $g\colon A
\to B$ असे कोणतेही एकास-एक आच्छादन नसते; म्हणजे, $\cardle{A}{B}$ आणि
$\cardneq{A}{B}$.
\end{defn}

हा संबंध अपरावर्ती आणि संक्रमक आहे, हे स्पष्ट आहे. (हे सरावासाठी
सोडले आहे.) या संकेतनाने संच $A$ हा गणनीय असतो तेव्हा आणि
केवळ तेव्हाच $\cardle{A}{\Nat}$, आणि $A$ हा अगणनीय असतो
तेव्हा आणि केवळ तेव्हाच $\cardless{\Nat}{A}$, असे म्हणता येते. त्यामुळे
%
\ref{sfr:siz:nen-alt:thm:nonenum-pownat}
हे प्रमेय $\cardless{\Nat}{\Pow{\Nat}}$ असे पुन्हा मांडता येते.
खरे तर, हा शेवटचा मुद्दा \emph{पूर्णपणे व्यापक} आहे, हे
Cantor (1892) यांनी सिद्ध केले:

\begin{thm}[कँटर]\label{sfr:siz:car:thm:cantor}
कोणत्याही संच $A$ साठी $\cardless{A}{\Pow{A}}$.
\end{thm}

\begin{proof}
$f(x) = \{x\}$ हे $f \colon A \to \Pow{A}$ असे एकास-एक फलन आहे;
कारण $x \neq y$ असेल, तर विस्तारात्मकतेमुळे $\{x\} \neq \{y\}$ सुद्धा
असते, आणि म्हणून $f(x) \neq f(y)$. त्यामुळे $\cardle{A}{\Pow{A}}$.

\begin{editorial}
\ref{sfr:siz:nen:sec} समाविष्ट असेल, तर आपण संथ सिद्धता देतो; अन्यथा
\ref{sfr:siz:nen-alt:sec} शी जुळणारी अधिक जलद सिद्धता देतो.
\end{editorial}

%
  आता $g\colon A \to
  \Pow{A}$ असे आच्छादक फलन, आणि त्यामुळे एकास-एक व आच्छादक फलनही,
  असू शकत नाही हे दाखवू; म्हणून $\cardneq{A}{\Pow{A}}$. त्यासाठी
  $g\colon A \to \Pow{A}$ असे समजा. $g$ हे सर्वत्र परिभाषित असल्याने
  प्रत्येक $x \in A$ हा $g(x)
  \subseteq A$ अशा उपसंचाकडे पाठवला जातो. $g$ हे आच्छादक असू शकत नाही,
  हे दाखवता येते. त्यासाठी व्याख्येमुळेच~$g$ च्या व्याप्तीत येऊ न
  शकणारा $\overline{A} \subseteq A$ असा उपसंच परिभाषित करू. पुढील संच घ्या:
  \[  \overline{A} = \Setabs{x \in A}{x \notin g(x)}.
  \]
  प्रत्येक $x \in A$ साठी $g(x)$ परिभाषित असल्याने $\overline{A}$ हा
  स्पष्टपणे~$A$ चा सुयोग्य रीतीने परिभाषित उपसंच आहे. पण तो~$g$ च्या
  व्याप्तीत असू शकत नाही. $x \in A$ स्वैर असू द्या; आपण
  $\overline{A} \neq g(x)$ दाखवू. $x \in g(x)$ असेल, तर तो
  $x \notin g(x)$ हा गुणधर्म पूर्ण करत नाही, आणि म्हणून~$\overline{A}$
  च्या व्याख्येवरून $x \notin
  \overline{A}$. $x \in \overline{A}$ असेल, तर त्याने~$\overline{A}$
  चा व्याख्यात्मक गुणधर्म पूर्ण केला पाहिजे; म्हणजे $x \in A$ आणि
  $x \notin
  g(x)$. $x$ स्वैर असल्याने प्रत्येक $x \in A$ साठी $x \in g(x)$ तेव्हा
  आणि केवळ तेव्हाच $x \notin \overline{A}$; म्हणून
  $g(x) \neq \overline{A}$.
  दुसऱ्या शब्दांत, $\overline{A}$ हा~$g$ च्या व्याप्तीत येऊ शकत नाही;
  त्यामुळे~$g$ आच्छादक आहे या गृहीतकाला व्याघात होतो.{}
\begin{quote}\small\textbf{स्रोतदुरुस्ती.} OLSIZ-007: गोठवलेल्या संथ सिद्धतेत x स्वैरपणे A मधून घेतल्यानंतर निष्कर्षाचा आवाका चुकून “प्रत्येक x जो A-bar मध्ये आहे” असा दिला आहे. g(x) आणि A-bar वेगळे आहेत हे प्रत्येक x\ensuremath{\in}A साठी दाखवणे आवश्यक आहे; मराठी सिद्धतेत तो आवाका दुरुस्त केला आहे.\end{quote}
\end{proof}

\begin{explain}

  \ref{sfr:siz:car:thm:cantor} ची सिद्धता आणि
  \ref{sfr:siz:nen:thm:nonenum-pownat} ची सिद्धता यांची तुलना करणे उद्बोधक
  आहे. तेथे~$\PosInt$ च्या उपसंचांची कोणतीही $Z_1$, $Z_2$, \dots{}
  अशी यादी दिली असता, यादीत नसेल असा संख्यांचा संच~$\overline{Z}$
  रचता येतो, हे दाखवले. तो यादीत नसतो याची खात्री यामुळे मिळते की
  प्रत्येक $n \in
  \PosInt$ साठी $n \in Z_n$ तेव्हा आणि केवळ तेव्हाच
  $n \notin \overline{Z}$. त्यामुळे $Z_n$ आणि $\overline{Z}$ यांपैकी
  एकाचा घटक असलेली, पण दुसऱ्याचा नसलेली कोणती तरी संख्या
  नेहमी असते. येथेही तीच कल्पना वापरतो; फरक एवढाच की निर्देशांक~$n$
  आता~$\PosInt$ चे नसून~$A$ चे घटक आहेत. $\overline{A}$ ची
  व्याख्या अशी केली आहे की प्रत्येक $x \in A$ साठी तो~$g(x)$ पेक्षा
  वेगळा असतो, कारण $x \in g(x)$ तेव्हा आणि केवळ तेव्हाच
  $x \notin \overline{A}$. पुन्हा, $g(x)$ आणि $\overline{A}$ यांपैकी
  एकाचा घटक असलेला, पण दुसऱ्याचा नसलेला~$A$ चा कोणता तरी
  घटक नेहमी असतो. म्हणूनच जसा $\overline{Z}$ हा $Z_1$,
  $Z_2$, \dots{} या यादीत येऊ शकत नाही, तसा $\overline{A}$ हा~$g$ च्या
  व्याप्तीत येऊ शकत नाही.


  \ref{sfr:siz:car:thm:cantor} ची सिद्धता आणि
  \ref{sfr:siz:nen-alt:thm:nonenum-pownat} ची सिद्धता यांची तुलना करणे
  उद्बोधक आहे. तेथे~$\Nat$ च्या उपसंचांची कोणतीही $N_0$, $N_1$,
  $N_2$, \dots{} अशी यादी दिली असता, यादीत नसेल असा संख्यांचा संच~$D$
  रचता येतो, हे दाखवले. प्रत्येक $n \in \Nat$ साठी $n \in N_n$ तेव्हा
  आणि केवळ तेव्हाच $n \notin
  D$ असल्यामुळे तो यादीत नसतो याची खात्री मिळते. येथेही तीच कल्पना
  वापरतो; फरक एवढाच की निर्देशांक~$n$ आता~$\Nat$ चे नसून~$A$ चे
  घटक आहेत. $D$ ची व्याख्या अशी केली आहे की प्रत्येक $x
  \in A$ साठी तो~$g(x)$ पेक्षा वेगळा असतो, कारण $x \in g(x)$ तेव्हा
  आणि केवळ तेव्हाच $x \notin D$.

या सिद्धतेची रसेल यांच्या विरोधापत्तीच्या सिद्धतेशीही तुलना करणे उपयुक्त
आहे: \ref{sfr:set:rus:thm:russells-paradox}. खरे तर, कँटर यांच्या
प्रमेयातूनच रसेल यांना स्वतःची विरोधापत्ती सुचली.
\end{explain}

\begin{prob}
  कोणत्याही संच~$A$ साठी $g\colon \Pow{A} \to
  A$ असे एकास-एक फलन असू शकत नाही, हे दाखवा. सूचना: $g\colon
  \Pow{A} \to A$ हे एकास-एक आहे, असे समजा. पुढील संच विचारात घ्या:
  $D = \Setabs{g(B)}{B \subseteq A \text{ आणि
  } g(B) \notin B}$. $x = g(D)$ असे घ्या. $g$ हे एकास-एक आहे हा
  हे तथ्य वापरून व्याघात मिळवा.
\end{prob}









\section{आकारमानाची कल्पना आणि श्रेडर--बर्नस्टाइन प्रमेय}\label{sfr:siz:sb:sec}

\begin{explain}
पुढील विचार अंतर्ज्ञानाला पटणारा आहे: $A$ हा $B$ पेक्षा आकारमानाने मोठा नाही
आणि $B$ हा $A$ पेक्षा आकारमानाने मोठा नाही, तर $A$ आणि $B$ तुल्यबल
आहेत. खरे सांगायचे तर, हा विचार \emph{चुकीचा} असता, तर तुल्यबलतेच्या
आपल्या व्याख्येचा संचांच्या ``आकारमानां''च्या तुलनेशी काही संबंध आहे,
असे म्हणण्याला आपल्याजवळ जवळजवळ आधारच उरला नसता!{} सुदैवाने हा
अंतःप्रज्ञ विचार बरोबर आहे. श्रेडर--बर्नस्टाइन प्रमेय त्याचे समर्थन
करते.
\end{explain}

\begin{thm}[श्रेडर--बर्नस्टाइन]
	\label{sfr:siz:sb:thm:schroder-bernstein}
	$\cardle{A}{B}$ आणि $\cardle{B}{A}$ असेल,
	तर $\cardeq{A}{B}$.
\end{thm}

\begin{explain}
दुसऱ्या शब्दांत, $A$ पासून~$B$ कडे जाणारे एकास-एक फलन आणि $B$
पासून~$A$ कडे जाणारे एकास-एक फलन असेल, तर $A$ पासून~$B$ कडे
जाणारे एकास-एक आच्छादन असते.

पण या निष्कर्षाची सिद्धता खरोखरच बरीच \emph{कठीण} आहे. कँटर यांनी
निष्कर्ष मांडला असला, तरी इतरांनी तो सिद्ध केला.
\footnote{इतिहासाच्या अधिक माहितीसाठी, उदाहरणार्थ,
Potter (2004, pp.~165--6) पाहा.}
%
आत्तासाठी हा निष्कर्ष विश्वासाने स्वीकारावा लागेल.

सुदैवाने श्रेडर--बर्नस्टाइन प्रमेय \emph{बरोबर} आहे आणि म्हणून आपण
परिभाषित केलेले $\cardeq{A}{B}$ आणि $\cardle{A}{B}$ हे संबंध
``आकारमाना''शी निगडित आहेत, असे मानणे समर्थित होते. शिवाय
श्रेडर--बर्नस्टाइन प्रमेय फार \emph{उपयुक्त} आहे. दोन तुल्यबल संचांमध्ये
एकास-एक आच्छादन शोधणे कठीण असू शकते. या प्रमेयामुळे तुलना दोन दिशांत
विभागता येते: दोन्ही दिशांतील एकास-एक फलन स्वतंत्रपणे शोधावी
लागतात.
\end{explain}








\section{प्रगणने आणि गणनीय संच}\label{sfr:siz:enm-alt:sec}

\begin{editorial}
  या विभागात संचसिद्धान्तात रूढ असलेल्या पद्धतीने, $\Nat$ च्या
  (आरंभीच्या खंडां)शी असलेली एकास-एक आच्छादने म्हणून प्रगणनांची
  व्याख्या केली आहे. त्यामुळे \ref{sfr:siz:enm:sec} मधील व्याख्यांशी तिचा
  किंचित विरोध होतो, आणि तेथील सर्व उदाहरणे येथे पुन्हा येतात. हा
  विभाग त्या विभागापेक्षा थोडा अधिक संक्षिप्तही आहे.
\end{editorial}

सांत संच त्यातील घटक केवळ क्रमाने मांडून सांगता येतो. पुढीलप्रमाणे
संचाची व्याख्या करताना आपण हेच करतो:
\[  A = \{a_1, a_2, \ldots, a_n\}.
\]
घटक $a_1$, \dots, $a_n$ हे सर्व भिन्न आहेत असे गृहीत धरल्यास,
यातून $A$ आणि पहिल्या $n$ नैसर्गिक संख्या $0$, \dots, $n-1$ यांच्यात
एकास-एक आच्छादन मिळते. उलट, प्रत्येक सांत संचात सांत संख्येनेच
घटक असल्यामुळे प्रत्येक सांत संच अशी संगती लावून मांडता येतो.
दुसऱ्या शब्दांत, $A$ सांत असेल, तर $A$ आणि $\{0, \dots, n-1\}$
यांच्यात एकास-एक आच्छादन असते; येथे~$A$ मध्ये $n$ इतके घटक
असतात.

विशिष्ट प्रकारचे अनंत संच मान्य केल्यास, काही अनंत संचांचेही प्रगणन
करता येईल. अनंत संच आणि संपूर्ण~$\Nat$ यांच्यातील एकास-एक आच्छादन त्या
संचाचे प्रगणन करते, असे म्हणून हे नेमकेपणाने मांडता येते.

\begin{defn}[प्रगणनाची संचसैद्धान्तिक व्याख्या]
संच $A$ चे \emph{प्रगणन} म्हणजे असे एकास-एक आच्छादन की ज्याची व्याप्ती
$A$ असते आणि ज्याचा प्रांत नैसर्गिक संख्यांचा एखादा आरंभीचा खंड $\{0,
1, \ldots, n\}$ {किंवा} नैसर्गिक संख्यांचा संपूर्ण संच~$\Nat$ असतो.
\end{defn}

\begin{explain}
\emph{प्रगणन} या शब्दाच्या अशा वापरामागे सहज समजणारा आधार आहे.
संच $A$ चे प्रगणन केले आहे असे म्हणणे म्हणजे संच~$A$ मधील घटक मोजत
जाण्यास अनुमती देणारे एकास-एक आच्छादन $f$ आहे असे म्हणणे. $0$ वा घटक
$f(0)$, पहिला $f(1)$, \ldots आणि $n$ वा $f(n)$,~\ldots असतो.
\footnote{होय, आपण $0$ पासून मोजतो. अर्थात~$1$ पासूनही सुरुवात करता
येईल. त्यामुळे मोठा फरक पडणार नाही; फक्त~$\Nat$ ऐवजी~$\PosInt$
घ्यावा लागेल.} पुढील व्याख्या जोडल्यावर या मागचे कारण आणखी स्पष्ट होते:
\end{explain}

\begin{defn}
  \label{sfr:siz:enm-alt:defn:enumerable}
  संच~$A$ हा गणनीय असतो तेव्हा आणि केवळ तेव्हाच, जेव्हा
  $A = \emptyset$ असते किंवा~$A$ चे प्रगणन असते. $A$ हा
  अगणनीय आहे असे म्हणतो तेव्हा आणि केवळ तेव्हाच, जेव्हा
  $A$ हा गणनीय नसतो.
\end{defn}

\begin{explain}
म्हणून संच रिक्त असेल किंवा प्रगणन वापरून त्यातील घटक मोजता
येत असतील तेव्हा आणि केवळ तेव्हाच तो गणनीय असतो.
\end{explain}

\begin{ex}
नैसर्गिक संख्यांचे प्रगणन करणारे फलन म्हणजे फक्त $\Id{\Nat} \colon
\Nat \to \Nat$ हे $\Id{\Nat}(n) = n$ ने दिलेले एकरूपता फलन.
\emph{धन} नैसर्गिक संख्या $\Nat^+ = \Nat \setminus \{0\}$ यांचे
प्रगणन करणारे फलन $g(n) = n + 1$ आहे; म्हणजेच ते उत्तरवर्ती फलन आहे.
\end{ex}

\begin{prob}
संच $A$ हा गणनीय असतो तेव्हा आणि केवळ तेव्हाच $A =
\emptyset$ असते किंवा $f\colon \Nat \to A$ असे आच्छादन
असते, हे दाखवा. $A$ हा गणनीय असतो तेव्हा आणि केवळ तेव्हाच
$g\colon A \to \Nat$ असे एकास-एक फलन असते, हेही दाखवा.
\end{prob}

\begin{ex}
पुढीलप्रमाणे दिलेली $f\colon \Nat \to \Nat$ आणि $g \colon \Nat
\to \Nat$ ही फलने
\begin{align*}
  f(n) & = 2n \text{ आणि}\\
  g(n) & = 2n+1
\end{align*}
अनुक्रमे सम नैसर्गिक संख्यांचे आणि विषम नैसर्गिक संख्यांचे प्रगणन
करतात. पण यांपैकी कोणतेही फलन आच्छादक नाही; म्हणून कोणतेही
$\Nat$ चे प्रगणन नाही.
\end{ex}

\begin{prob}
वर्गसंख्या $1$, $4$, $9$, $16$, \dots{} यांचे प्रगणन परिभाषित करा.
\end{prob}

\begin{ex}
$\lceil x \rceil$ हे $x$ ला त्याच्यापेक्षा कमी नसलेल्या सर्वांत जवळच्या
पूर्णांकाकडे नेणारे \emph{ऊर्ध्व पूर्णांक} फलन असू द्या. मग पुढीलप्रमाणे
दिलेले फलन $f \colon \Nat \to \Int$:
\[  f(n) = (-1)^{n} \left\lceil\tfrac{n}{2}\right\rceil
\]
पूर्णांकांचा संच~$\Int$ पुढीलप्रमाणे प्रगणित करते:
\[\begin{array}{c c c c c c c c}
f(0) & f(1) & f(2) & f(3) & f(4) & f(5) & f(6) & \dots \\ \\
\big\lceil \tfrac{0}{2} \big\rceil & -\big\lceil \tfrac{1}{2}\big\rceil &  \big\lceil \tfrac{2}{2} \big\rceil & -\big\lceil \tfrac{3}{2} \big\rceil & \big\lceil \tfrac{4}{2} \big\rceil  & -\big\lceil \tfrac{5}{2}\big\rceil & \big\lceil \tfrac{6}{2} \big\rceil & \dots \\ \\
0 & -1 & 1 & -2 & 2 & -3 & 3& \dots
\end{array}
\]
धन आणि ऋण पूर्णांकांमध्ये आलटूनपालटून ``उड्या मारून'' $f$ हे $\Int$
ची मूल्ये कशी निर्माण करते ते पाहा. $f$ ची व्याख्या पुढीलप्रमाणे
प्रकरणे पाडून केली आहे असेही मानता येते:
\[f(n) = \begin{cases}
  \frac{n}{2} & \text{जर $n$ सम असेल}\\
  -\frac{n+1}{2} & \text{जर $n$ विषम असेल}
  \end{cases}
\]
\end{ex}

\begin{prob}
$A$ आणि $B$ हे गणनीय असतील, तर $A \cup B$ सुद्धा गणनीय
आहे, हे दाखवा.
\end{prob}

\begin{prob}
$A_1$, $A_2$, \dots, $A_n$ हे सर्व गणनीय असतील, तर
$A_1 \cup \dots \cup A_n$ सुद्धा गणनीय आहे, हे $n$ वरील गणिती विगमनाने
दाखवा.
\end{prob}









\section{अगणनीय संच}\label{sfr:siz:nen-alt:sec}

\begin{editorial}
  या विभागात \ref{sfr:siz:enm-alt:sec} मधील व्याख्या वापरून
  $\Bin^\omega$ आणि $\Pow{\Nat}$ यांची अगणनीयता सिद्ध केली आहे;
  म्हणजे $\PosInt$ पासून आच्छादन मागण्याऐवजी~$\Nat$ शी एकास-एक
  आच्छादन आवश्यक मानले आहे.
\end{editorial}

\begin{explain}
नैसर्गिक संख्यांचा संच $\Nat$ हा अनंत आहे. तो सहजपणे
गणनीय देखील आहे. परंतु
\emph{अगणनीय} संच, म्हणजे गणनीय नसलेले संच,
अस्तित्वात आहेत ही लक्षणीय गोष्ट आहे
(पाहा \ref{sfr:siz:enm-alt:defn:enumerable}).

हे आश्चर्यकारक वाटू शकते. कारण $A$ हा अगणनीय आहे असे
म्हणणे म्हणजे \emph{कोणतेही} एकास-एक आच्छादन $f \colon \Nat \to A$
नाही असे म्हणणे; म्हणजे~$\Nat$ मधील अनंतपणे अनेक घटक~$A$
मध्ये पाठवणारे कोणतेही फलन संपूर्ण~$A$ व्यापत नाही. म्हणून $A$ हा
अगणनीय असेल, तर नैसर्गिक संख्यांपेक्षा~$A$ मध्ये ``अधिक''
घटक असतात.

संच अगणनीय आहे हे सिद्ध करण्यासाठी अनुरूप एकास-एक आच्छादन
अस्तित्वात असू शकत नाही हे दाखवावे लागते. याचा सर्वोत्तम मार्ग म्हणजे
~$A$ मधील घटक प्रगणित करण्याचा प्रत्येक प्रयत्न किमान एक
घटक वगळतो हे दाखवणे; यावरून कोणतेही फलन $f\colon \Nat \to A$
हे आच्छादक नाही असे दिसते. हे प्रस्थापित करण्याची एक सर्वसाधारण
युक्ती म्हणजे कँटर यांची \emph{कर्णरेषीय पद्धत} वापरणे. $A$ मधील
घटक $x_1$, $x_2$, \dots{} असे मांडले असता, रचनेमुळे त्या
यादीत असूच शकणार नाही असा~$A$ चा आणखी एक घटक आपण बनवतो.

हे सर्व उदाहरणाने उत्तम समजते. म्हणून आपले पहिले उदाहरण म्हणजे $0$
आणि $1$ यांच्या सर्व अनंत चिन्हमालांचा संच~$\Bin^\omega$. (`$\Bin$'
हे द्विमान दर्शवते आणि त्याचा अर्थ फक्त दोन घटकांचा संच $\{0,1\}$
असा घेता येतो.)
तरी सध्या या किंचित शिथिल मांडणीमुळे गोंधळ होऊ नये.
\end{explain}

\begin{thm}
\label{sfr:siz:nen-alt:thm:nonenum-bin-omega}
$\Bin^\omega$~हा अगणनीय आहे.
\end{thm}

\begin{proof}
$\Bin^\omega$ च्या एखाद्या उपसंचाचे कोणतेही प्रगणन घ्या. म्हणजे
$s_{0}$, $s_{1}$, $s_{2}$, \dots{} अशी यादी आपल्याकडे आहे, जेथे
प्रत्येक $s_n$ ही $0$ आणि~$1$ यांची अनंत चिन्हमाला आहे. या यादीतील
$n$ व्या चिन्हमालेतील $m$ वा अंक $s_n(m)$ असू द्या.
\begin{quote}\small\textbf{स्रोतदुरुस्ती.} OLSIZ-008: गोठवलेल्या इंग्रजी वर्णनात s\_n(m) ला m व्या चिन्हमालेतील n वा अंक म्हटले आहे; लगेचचे वाक्य आणि सारणी मात्र तो n व्या ओळीतील, म्हणजे n व्या चिन्हमालेतील, m वा अंक असल्याचे निश्चित करतात. मराठी मजकुरात ओळ-स्तंभाशी सुसंगत निर्देशांकक्रम वापरला आहे.\end{quote}
मग यादीकडे सरणी म्हणून पाहता येते, ज्यात $s_n(m)$ हा $n$ व्या ओळीत
आणि $m$ व्या स्तंभात ठेवलेला आहे:
\[\begin{array}{c|c|c|c|c|c}
& 0 & 1 & 2 & 3 & \dots \\\hline
0 & \mathbf{s_{0}(0)} & s_{0}(1) & s_{0}(2) & s_0(3) & \dots \\\hline
1 & s_{1}(0)& \mathbf{s_{1}(1)} & s_1(2) & s_1(3) & \dots \\\hline
2 & s_{2}(0)& s_{2}(1) & \mathbf{s_2(2)} & s_2(3) & \dots \\\hline
3 & s_{3}(0)& s_{3}(1) & s_3(2) & \mathbf{s_3(3)} & \dots \\\hline
\vdots & \vdots & \vdots & \vdots & \vdots & \mathbf{\ddots}
\end{array}
\]
आता या यादीत नसलेली $0$ आणि $1$ यांची $d$ ही अनंत चिन्हमाला आपण
बनवू. तिची प्रत्येक नोंद सांगून, म्हणजे सर्व $n \in \Nat$ साठी
$d(n)$ सांगून, हे करू. सहज कल्पना अशी: वरील सरणीचा कर्ण खाली वाचायचा
(त्यामुळे ``कर्णरेषीय पद्धत'' हे नाव) आणि प्रत्येक $1$ चा $0$, तर
प्रत्येक $0$ चा~$1$ करायचा.
\begin{quote}\small\textbf{स्रोतदुरुस्ती.} OLSIZ-009: गोठवलेल्या इंग्रजी सूचनेत “प्रत्येक 1 चा 0” हेच परिवर्तन दोनदा आले आहे. लगेचचे प्रकरणनिहाय सूत्र परस्परपूरक दोन्ही बदल निश्चित करते: 1 चा 0 आणि 0 चा 1. मराठी मजकुरात हे दोन्ही बदल स्पष्ट केले आहेत.\end{quote}
अधिक अमूर्तपणे, कर्णावरील $n$ वा घटक, $s_n(n)$, हा $1$ किंवा
$0$ आहे त्यानुसार $d(n)$ हा $0$ किंवा $1$ असावा अशी व्याख्या करतो;
म्हणजे:
\[d(n) =
\begin{cases}
1 & \text{जर $s_{n}(n) = 0$ असेल}\\
0 & \text{जर $s_{n}(n) = 1$ असेल}
\end{cases}
\]
स्पष्टच, $d \in \Bin^\omega$, कारण ती $0$ आणि $1$ यांची अनंत
चिन्हमाला आहे. पण आपण $d$ अशी बनवली की कोणत्याही $n \in \Nat$
साठी $d(n) \neq s_n(n)$. म्हणजे $d$ ही $s_n$ पेक्षा तिच्या $n$ व्या
नोंदीत वेगळी आहे. म्हणून कोणत्याही $n\in \Nat$ साठी $d \neq s_n$.
त्यामुळे $d$ ही $s_0$, $s_1$, $s_2$, \dots{} या यादीत असू शकत नाही.

$\Bin^\omega$ च्या एखाद्या उपसंचाचे स्वैर प्रगणन दिले असता ते
$\Bin^\omega$ चा कोणता तरी घटक वगळेल हे आपण दाखवले आहे.
म्हणून $\Bin^\omega$ या संचाचे प्रगणन नाही; म्हणजे $\Bin^\omega$ हा
अगणनीय आहे.
\end{proof}

\begin{explain}
या सिद्धतापद्धतीला ``कर्णीकरण'' म्हणतात, कारण ती~$d$ ची व्याख्या
करण्यासाठी सरणीचा कर्ण वापरते. मात्र कर्णीकरणात सरणी असणे आवश्यक नाही.
प्रत्यक्षात, सरणी किंवा प्रत्यक्ष कर्ण नसतानाही हिच्याशी मिळतीजुळती कल्पना
वापरून एखादा संच अगणनीय आहे हे दाखवता येते. पुढील फलित हे
कसे करता येते ते दाखवते.
\end{explain}

\begin{thm}
\label{sfr:siz:nen-alt:thm:nonenum-pownat}
$\Pow{\Nat}$ हा गणनीय नाही.
\end{thm}

\begin{proof}
~$\Nat$ च्या उपसंचांची प्रत्येक यादी $\Nat$ चा कोणता तरी उपसंच वगळते
हे दाखवून आपण त्याच पद्धतीने पुढे जाऊ. म्हणून $\Nat$ च्या उपसंचांची
$N_0, N_1, N_2, \ldots$ अशी एखादी यादी आहे असे समजा. $D$ या संचाची
व्याख्या पुढीलप्रमाणे करतो: $n \in D$ तेव्हा आणि केवळ तेव्हाच, जेव्हा
$n \notin N_{n}$:
\[D = \Setabs{n \in \Nat}{n \notin N_n}
\]
स्पष्टच $D\subseteq \Nat$. पण $D$ हा यादीत असू शकत नाही. कारण
रचनेनुसार $n \in D$ तेव्हा आणि केवळ तेव्हाच, जेव्हा $n\notin N_n$;
म्हणून कोणत्याही $n \in \Nat$ साठी $D \neq N_n$.
\end{proof}

\begin{explain}
आधीच्या सिद्धतेत कर्णाचा उल्लेख नव्हता. तरी पुढीलप्रमाणे चित्र डोळ्यांसमोर
आणल्यास त्यात कर्ण आहे असे मानता येते: $N_0$, $N_1$, \dots{} हे संच
एका सरणीत लिहिल्याची कल्पना करा; $m \in N_n$ असेल तेव्हा आणि केवळ
तेव्हाच $m$ व्या स्तंभात $m$ लिहून $N_n$ हा $n$ व्या ओळीत लिहायचा.
उदाहरणार्थ, त्या यादीतील पहिले चार संच $\{0,1,2,\dots\}$,
$\{1, 3, 5, \dots\}$, $\{0,1,4\}$ आणि $\{2,3,4,\dots\}$ आहेत असे
समजा; मग आपल्या सरणीची सुरुवात अशी असेल:
\[\begin{array}{r@{}rrrrrrr}
  N_0 = \{ & \mathbf{0}, & 1, & 2, & & & & \dots\}\\
  N_1 = \{ &  & \mathbf{1}, &  & 3, &  & 5, & \dots\}\\
  N_2 = \{ & 0, & 1, &  &  & 4\phantom{,} &  & \}\\
  N_3 = \{ &  &  & 2, & \mathbf{3}, & 4, & & \dots\}\\
  &\vdots & & & & & & \ddots\phantom{\}}
  \end{array}
\]
मग कर्णाच्या खाली जात, $n$ कर्णावर \emph{नसेल} तेव्हा आणि केवळ
तेव्हाच $n \in D$ असे ठेवून $D$ हा संच मिळतो. म्हणून वरील उदाहरणात
आपण $0$ आणि $1$ वगळू,~$2$ समाविष्ट करू,~$3$ वगळू, इत्यादी.
\end{explain}

\begin{prob}
सर्व फलने $f \colon \Nat \to \Nat$ यांचा संच स्पष्ट कर्णरेषीय
युक्तिवादाने अगणनीय आहे हे दाखवा. म्हणजे $f_1$, $f_2$,
\dots{} ही फलनांची यादी असून प्रत्येक $f_i\colon \Nat \to \Nat$
असेल, तर या यादीत नसलेले कोणते तरी $g \colon \Nat \to \Nat$ असते,
हे दाखवा.
\end{prob}









\section{न्यूनीकरण}\label{sfr:siz:red-alt:sec}

\begin{editorial}
  या विभागात \ref{sfr:siz:nen-alt:sec} मधील फलितांशी जुळणारी अगणनीयता
  न्यूनीकरणाने सिद्ध केली आहे. \ref{sfr:siz:nen:sec} मधील फलितांशी जुळणारी
  किंचित अधिक विस्तृत पर्यायी आवृत्ती \ref{sfr:siz:red:sec} मध्ये दिली आहे.
\end{editorial}

कर्णीकरणाच्या युक्तिवादाने $\Bin^\omega$ हा अगणनीय आहे,
हे आपण सिद्ध केले. तसाच कर्णीकरणाचा युक्तिवाद वापरून $\Pow{\Nat}$ हा
अगणनीय आहे हेही दाखवले. पण $\Pow{\Nat}$ हा
अगणनीय असल्याचे सिद्ध करण्याचा आणखी एक मार्ग आहे:
\emph{$\Pow{\Nat}$ हा गणनीय असेल, तर $\Bin^\omega$ हाही
गणनीय असतो} हे दाखवा. $\Bin^\omega$ हा अगणनीय
आहे हे आपल्याला माहीत असल्यामुळे $\Pow{\Nat}$ हाही तसाच आहे असा
निष्कर्ष निघेल.

याला एक समस्या दुसऱ्या समस्येकडे \emph{न्यूनीत करणे} म्हणतात. या
बाबतीत $\Bin^\omega$ चे प्रगणन करण्याची समस्या आपण $\Pow{\Nat}$ चे
प्रगणन करण्याच्या समस्येकडे न्यूनीत करतो. नंतरच्या समस्येचे उत्तर---
$\Pow{\Nat}$ चे प्रगणन---पहिल्या समस्येचे उत्तर---$\Bin^\omega$ चे
प्रगणन---देईल.

संच~$B$ च्या प्रगणनाची समस्या संच~$A$ च्या प्रगणनाच्या समस्येकडे
न्यूनीत करण्यासाठी~$A$ चे प्रगणन~$B$ च्या प्रगणनात रूपांतरित करण्याची
पद्धत आपण देतो. हे करण्याचा सर्वांत सोपा मार्ग म्हणजे
$f\colon A \to B$ असे आच्छादन परिभाषित करणे. $x_1$, $x_2$,
\dots{} हे~$A$ चे प्रगणन असेल, तर $f(x_1)$, $f(x_2)$, \dots{} हे~$B$
चे प्रगणन करील. आपल्या बाबतीत आपण
$f\colon \Pow{\Nat} \to \Bin^\omega$ असे आच्छादन शोधत आहोत.

\begin{prob}
$g\colon B \to A$ असे एकास-एक फलन असेल आणि $B$ हा
अगणनीय असेल, तर $A$ हाही तसाच असतो, हे दाखवा. $g$ वापरून
~$A$ चे प्रगणन~$B$ च्या प्रगणनात कसे रूपांतरित करता येते हे दाखवून
सिद्धता द्या.
\end{prob}

\begin{proof}[न्यूनीकरणाने {\ref{sfr:siz:nen-alt:thm:nonenum-pownat}} ची सिद्धता]
न्यूनीकरणासाठी, $\Pow{\Nat}$ हा गणनीय आहे आणि त्यामुळे त्याचे
$N_{1}$, $N_{2}$, $N_{3}$, \dots{} असे प्रगणन आहे, असे समजा.

$f \colon \Pow{\Nat} \to \Bin^\omega$ हे फलन पुढीलप्रमाणे परिभाषित
करा: $f(N)$ ही अशी चिन्हमाला~$s$ असू द्या की $s(n) = 1$ तेव्हा आणि
केवळ तेव्हाच, जेव्हा $n \in N$; अन्यथा $s(n) = 0$.
\begin{quote}\small\textbf{स्रोतदुरुस्ती.} OLSIZ-010: गोठवलेल्या इंग्रजी व्याख्येत f(N) ला s\_k असे म्हटले आहे, पण k बांधलेला किंवा N वरून ठरवलेला नाही. घटकनिहाय समीकरण N ची अभिलक्षण-चिन्हमाला एकमेव ठरवते; मराठी व्याख्येत एकच बद्ध निर्गत नाव s आणि त्याचे घटक s(n) सातत्याने वापरले आहेत.\end{quote}

याने फलन स्पष्टपणे परिभाषित होते, कारण $N \subseteq \Nat$ असताना
कोणताही $n \in \Nat$ हा $N$ चा घटक असतो किंवा नसतो.
उदाहरणार्थ, सम नैसर्गिक संख्यांचा संच
$2\Nat = \Setabs{2n}{n \in \Nat} = \{0,2, 4, 6, \dots\}$ हा
$1010101\dots$ या चिन्हमालेकडे; $\emptyset$ हा $0000\dots$ कडे; आणि
$\Nat$ हा $1111\dots$ कडे पाठवला जातो.

हे फलन आच्छादक देखील आहे: $0$ आणि $1$ यांची प्रत्येक चिन्हमाला
नैसर्गिक संख्यांच्या कोणत्या तरी संचाशी अनुरूप असते; तो संच म्हणजे
चिन्हमालेत ज्या स्थानांवर~$1$ आहे त्यांना अनुरूप नैसर्गिक संख्या
सदस्य म्हणून असलेला संच. अधिक नेमकेपणाने, $s \in \Bin^\omega$ असेल,
तर पुढीलप्रमाणे $N \subseteq \Nat$ परिभाषित करा:
\[N = \Setabs{n \in \Nat}{s(n) = 1}
\]
मग~$f$ ची व्याख्या पाहून $f(N) = s$ हे तपासता येते.

आता पुढील यादी विचारात घ्या:
\[f(N_1), f(N_2), f(N_3), \dots
\]
$f$ हे आच्छादक असल्यामुळे $\Bin^\omega$ चा प्रत्येक सदस्य
हा~$f$ चे कोणत्या तरी फलसाधकावरील मूल्य असला पाहिजे आणि म्हणून यादीत
आला पाहिजे. त्यामुळे ही यादी संपूर्ण~$\Bin^\omega$ चे प्रगणन करते.

म्हणून $\Pow{\Nat}$ हा गणनीय असता, तर $\Bin^\omega$ हाही
गणनीय असता. पण $\Bin^\omega$ हा अगणनीय आहे
(\ref{sfr:siz:nen-alt:thm:nonenum-bin-omega}). त्यामुळे $\Pow{\Nat}$ हा
अगणनीय आहे.
\end{proof}


\begin{prob}\label{sfr:siz:red:prob:nat-nat-alt}
  $f\colon \Nat \to \Nat$ अशा सर्व फलनांचा संच~$X$ हा
  न्यूनीकरणाच्या युक्तिवादाने अगणनीय आहे, हे दाखवा.
  (सूचना: $X$ पासून~$\Bin^\omega$ कडे जाणारे आच्छादक फलन द्या.)
\end{prob}

\begin{prob}
नैसर्गिक संख्यांच्या जोड्यांच्या \emph{सर्व संचांचा} संच, म्हणजे
$\Pow{\Nat \times \Nat}$, न्यूनीकरणाच्या युक्तिवादाने
अगणनीय आहे, हे दाखवा.
\end{prob}

\begin{prob}
नैसर्गिक संख्यांच्या अनंत अनुक्रमांचा संच $\Nat^\omega$ हा
न्यूनीकरणाच्या युक्तिवादाने अगणनीय आहे, हे दाखवा.
\end{prob}


\begin{prob}
$S$ हा $\Nat$ पासून $\{0,1\}$ या संचाकडे जाणारी सर्व
आच्छादक फलने असलेला संच असू द्या; म्हणजे $S$ मध्ये
$f \colon \Nat \to \Bin$ अशी सर्व आच्छादक फलने आहेत. $S$ हा
अगणनीय आहे, हे दाखवा.
\end{prob}

\begin{prob}
सर्व वास्तव संख्यांचा संच~$\Real$ हा अगणनीय आहे, हे दाखवा.
\end{prob}



\chapter{अंकगणितीकरण}\label{sfr:arith::chap}
\begin{editorial}
या प्रकरणातील सामग्री अनौपचारिक संचसिद्धान्तात संख्याप्रणालींची रचना
मांडते. ती टिम बटन यांच्या Open Set Theory या ग्रंथातून घेतली आहे.
\end{editorial}






\section{$\Nat$ पासून $\Int$ कडे}\label{sfr:arith:int:sec}

येथे दोन मूलभूत निरीक्षणे आहेत:
	\begin{enumerate}
		\item प्रत्येक पूर्णांक $n - m$ या रूपात लिहिता येतो, जेथे $n, m \in\Nat$.
		\item $n - m$ या व्यंजकात सांकेतिक केलेली माहिती $\tuple{n, m}$ या क्रमित जोडीद्वारेही तितकीच सांकेतिक करता येते.
	\end{enumerate}
नैसर्गिक संख्यांच्या क्रमित जोड्या $\Nat^2$ चे घटक आहेत हे आपल्याला आधीच माहीत आहे. आणि आपण $\Nat$ समजतो असे गृहीत धरत आहोत. त्यामुळे या दोन निरीक्षणांवर आधारलेली एक भोळी सूचना पुढे करता येईल: \emph{नैसर्गिक संख्यांच्या क्रमित जोड्यांना पूर्णांक म्हणून हाताळू या}.

प्रत्यक्षात ही सूचना अतिशय भोळी आहे. अर्थात, $0- 2 = 4 - 6$ असले पाहिजे असे आपल्याला हवे आहे. पण स्पष्टपणे $\tuple{0, 2 }\neq \tuple{4, 6}$. त्यामुळे $\Nat^2$ हाच पूर्णांकांचा संच आहे असे सरळ म्हणता येत नाही.

मागील अडचणीचे सामान्यीकरण केल्यास आपल्याला पुढील अट हवी आहे:
	$$a - b = c - d \text{ तेव्हा आणि केवळ तेव्हाच }a + d = c + b$$
(पूर्णांकांचे वर्तन \emph{असेच अभिप्रेत} आहे हे स्पष्ट असावे: दोन्ही बाजूंना फक्त $b$ आणि $d$ मिळवा.) हे वर्तन निश्चित करण्याचा सोपा मार्ग म्हणजे क्रमित जोड्यांमधील $\Intequiv$ हा तुल्यता संबंध पुढीलप्रमाणे व्याख्यायित करणे:
$$\tuple{ a, b } \Intequiv \tuple{c, d}\text{ तेव्हा आणि केवळ तेव्हाच }a + d = c + b$$
आता हा तुल्यता संबंध आहे हे दाखवावे लागेल.
\begin{prop} $\Intequiv$ हा तुल्यता संबंध आहे.
\end{prop}
	\begin{proof}
	$\Intequiv$ परावर्ती, सममित आणि संक्रमक आहे हे दाखवले पाहिजे.

	\emph{परावर्तिता:} $a + b = b + a$ असल्यामुळे स्पष्टपणे $\tuple{a, b} \Intequiv \tuple{a, b}$.

	\emph{सममितता:} $\tuple{a, b} \Intequiv \tuple{c, d}$ असे समजा; म्हणजे $a + d = c + b$. मग $c + b = a + d$, त्यामुळे $\tuple{c, d} \Intequiv \tuple{a, b}$.

	\emph{संक्रमकता:} $\tuple{a, b} \Intequiv \tuple{c, d}\Intequiv \tuple{m, n}$ असे समजा. त्यामुळे $a + d = c + b$ आणि $c + n = m + d$. म्हणून $a + d + c + n = c + b + m + d$, व म्हणून $a + n = m + b$. परिणामी $\tuple{a, b } \Intequiv \tuple{m, n}$.
\end{proof}

आता या तुल्यता संबंधाचा वापर करून तुल्यतावर्ग घेता येतात:

\begin{defn}
		नैसर्गिक संख्यांच्या क्रमित जोड्यांचे $\Intequiv$ अंतर्गत तुल्यतावर्ग म्हणजे पूर्णांक; म्हणजेच $\Int = \equivclass{\Nat^2}{\Intequiv}$.
\end{defn}

या संकेतजनक व्याख्येबद्दल अनेक निरनिराळ्या \emph{तात्त्विक} प्रतिक्रिया असू शकतात. त्या प्रतिक्रियांचा विचार करण्यापूर्वी मात्र काही तांत्रिक तपशील पुढे नेणे उपयुक्त ठरेल.

पूर्णांक म्हणजे काय हे ठरवल्यावर त्यांच्यावरील मूलभूत फलने आणि संबंध व्याख्यायित करावे लागतील. ज्या $\Intequiv$ अंतर्गत तुल्यतावर्गात $\tuple{m, n}$ हा घटक आहे, त्या वर्गासाठी $\equivrep{m, n}{\Intequiv}$ असे लिहू या.\footnote{टीप: \ref{sfr:rel:eqv:def:equivalenceclass} मध्ये दिलेले संकेतन वापरल्यास त्याच वर्गासाठी आपण $\equivrep{\tuple{m,n}}{\Intequiv}$ असे लिहिले असते. पण ते वाचायला जरा कठीण आहे.} म्हणजे:
	$$\equivrep{m, n}{\Intequiv} = \Setabs{\tuple{a, b} \in \Nat^2}{\tuple{a, b}\Intequiv \tuple{m, n}}$$
आता काही व्याख्या देऊ:
	\begin{align*}
		\equivrep{a, b}{\Intequiv} + \equivrep{c, d}{\Intequiv} &= \equivrep{a + c, b + d}{\Intequiv}\\
		\equivrep{a, b}{\Intequiv} \times \equivrep{c, d}{\Intequiv} &= \equivrep{a c + b  d, a  d + b c}{\Intequiv}\\
		\equivrep{a, b}{\Intequiv} \leq \equivrep{c, d}{\Intequiv} &\text{ तेव्हा आणि केवळ तेव्हाच }a + d \leq b + c
	\end{align*}
(प्रथेप्रमाणे, स्वयंसिद्धे वाचायला सोपी व्हावीत म्हणून मी `$(a \times b)$' या अर्थी `$ab$' लिहित आहे.) आता या व्याख्या \emph{अपेक्षेप्रमाणे} वागतात याची खात्री केली पाहिजे. याचा नेमका अर्थ उलगडून संपूर्ण पडताळणी करणे बरेच कष्टाचे आहे; तपशील आपण \ref{sfr:arith:check:sec} कडे सोपवतो. पण थोडक्यात: सर्व काही व्यवस्थित चालते!

आता फक्त एक गोष्ट उरते. आपण नैसर्गिक संख्यांच्या साहाय्याने
पूर्णांकांची रचना केली आहे. पण त्यामुळे नैसर्गिक संख्या \emph{स्वतः
पूर्णांक नसतील}. याचे तात्त्विक महत्त्व आपण \ref{sfr:arith:ref:sec} मध्ये
पुन्हा पाहू. मात्र पूर्णपणे तांत्रिक दृष्टीने, नैसर्गिक संख्यांना
पूर्णांक \emph{म्हणून} हाताळण्याची काही पद्धत आपल्याला हवी. कल्पना अगदी
सोपी आहे: प्रत्येक $n \in \Nat$ साठी $n_\Int = \equivrep{n, 0}{\Intequiv}$
असे आपण ठरवतो. ही व्याख्या सुयोग्य रीतीने वागते याची, म्हणजे कोणत्याही
$m, n \in \Nat$ साठी पुढील अटी पूर्ण होतात याची खात्री केली पाहिजे:
\begin{align*}
	(m + n)_\Int &= m_\Int + n_\Int\\
	(m \times n)_\Int &= m_\Int \times n_\Int\\
	m \leq n &\liff m_\Int \leq n_\Int
\end{align*}
ही सर्व पडताळणी अगदी सरळ आहे. उदाहरणार्थ, दुसरी अट खरी आहे हे
दाखवताना नैसर्गिक संख्यांच्या परिचित वर्तनाचा थेट वापर करून पुढीलप्रमाणे
युक्तिवाद करता येतो:
	\begin{align*}
		(m \times n)_\Int  &= \equivrep{m \times n, 0}{\Intequiv} \\
		&= \equivrep{m \times n + 0 \times 0, m \times 0 + 0 \times n}{\Intequiv} \\
		&= \equivrep{m, 0}{\Intequiv} \times \equivrep{n, 0}{\Intequiv} \\
		&= m_\Int \times n_\Int
	\end{align*}
इतर दोन अटी पडताळण्याचे काम सरावासाठी सोडतो.
\begin{prob}
	कोणत्याही $m, n \in \Nat$ साठी $(m + n)_\Int = m_\Int + n_\Int$ आणि $m \leq n \liff m_\Int \leq n_\Int$ हे दाखवा.
\end{prob}







\section{$\Int$ पासून $\Rat$ कडे}\label{sfr:arith:rat:sec}

अनौपचारिक संचसिद्धान्त वापरून नैसर्गिक संख्यांपासून पूर्णांकांची रचना कशी
करायची हे आपण नुकतेच पाहिले. आता अगदी त्याच प्रकारे पूर्णांकांपासून
परिमेय संख्यांची रचना कशी करायची ते पाहू. सुरुवातीची निरीक्षणे अशी:
\begin{enumerate}
	\item प्रत्येक परिमेय संख्या $\nicefrac{i}{j}$ या रूपात लिहिता येते,
	जेथे $i$ आणि $j$ दोन्ही पूर्णांक असतात, पण $j$ शून्येतर असतो.
	\item $\nicefrac{i}{j}$ या व्यंजकात सांकेतिक केलेली माहिती
	$\tuple{ i, j}$ या क्रमित जोडीद्वारेही तितकीच सांकेतिक करता येते.
\end{enumerate}
यावरून परिमेय संख्यांकडे $\Int \times (\Int \setminus \{0_\Int\})$
मधून घेतलेल्या क्रमित जोड्या \emph{म्हणून} पाहणे हा स्वाभाविक मार्ग
वाटेल. मात्र आधीप्रमाणे ही कल्पना जरा अतिभोळी ठरेल, कारण आपल्याला
$\nicefrac{3}{2} = \nicefrac{6}{4}$ हवे आहे, पण $\tuple{ 3, 2}\neq \tuple{ 6,
4}$. अधिक सर्वसाधारणपणे आपल्याला पुढील अट हवी:
\[	\nicefrac{a}{b} = \nicefrac{c}{d} \text{ तेव्हा आणि केवळ तेव्हाच } a \times d = b \times c
\]
ही अट मिळवण्यासाठी $\Int \times (\Int \setminus \{0_\Int\})$ वरील
{तुल्यता संबंध} आपण पुढीलप्रमाणे व्याख्यायित करतो:
\[	\tuple{ a, b }\Ratequiv \tuple{ c, d} \text{ तेव्हा आणि केवळ तेव्हाच }a \times d = b \times c
\]
हा तुल्यता संबंध आहे हे तपासले पाहिजे. ही तपासणी बरीचशी $\Intequiv$
च्या प्रकरणासारखीच आहे आणि ती आपण सरावासाठी सोडू.
\begin{prob}
$\Ratequiv$ हा तुल्यता संबंध आहे हे दाखवा.
\end{prob}
त्यामुळे पुढील व्याख्या देता येते:
\begin{defn}
दुसरा घटक शून्येतर असलेल्या पूर्णांकांच्या जोड्यांचे $\Ratequiv$ अंतर्गत
तुल्यतावर्ग म्हणजे परिमेय संख्या. म्हणजेच $\Rat =
\equivclass{(\Int \times (\Int\setminus \{0_\Int\}))}{\Ratequiv}$.
\end{defn}

पूर्णांकांप्रमाणे येथेही काही मूलभूत क्रिया व्याख्यायित करायच्या आहेत.
$\tuple{i, j}$ हा घटक असलेला $\Ratequiv$ अंतर्गत तुल्यतावर्ग
$\equivrep{i,j}{\Ratequiv}$ असेल, तर आपण पुढील व्याख्या करतो:
\begin{align*}
	\equivrep{a, b}{\Ratequiv} + \equivrep{c, d}{\Ratequiv} &= \equivrep{ad + bc,  bd}{\Ratequiv}\\
	\equivrep{a, b}{\Ratequiv} \times \equivrep{c, d}{\Ratequiv} &= \equivrep{a  c, b d}{\Ratequiv}.
\intertext{या परिमेय संख्यांवर $r \leq s$ ची व्याख्या करण्यासाठी आपण पुढील तथ्य वापरतो: $r \le s$ तेव्हा आणि केवळ तेव्हाच $s - r$ हा फरक ऋण नसतो; म्हणजे $s - r$ हा $\nicefrac{i}{j}$ या रूपात लिहिता येतो, जेथे $i$ ऋणेतर आणि $j$~धन असतो:}
	\equivrep{a, b}{\Ratequiv} \leq \equivrep{c, d}{\Ratequiv} &\text{ तेव्हा आणि केवळ तेव्हाच }
	\equivrep{c, d}{\Ratequiv} - \equivrep{a, b}{\Ratequiv} =
	\equivrep{i_\Int, j_\Int}{\Ratequiv}
\end{align*}
\begin{quote}\small\textbf{स्रोतदुरुस्ती.} MRARITH-001: गोठवलेल्या इंग्रजी स्पष्टीकरणात “s-r is not negative” नंतर चुकून r-s ला ऋणेतर अपूर्णांकाच्या रूपात लिहिले आहे. लगेचची प्रदर्शित व्याख्या पुन्हा s-r वापरते. मराठी मजकुरात या दोन्ही नियंत्रक ठिकाणांशी सुसंगत s-r वापरले आहे.\end{quote}
असे काही $i \in \Nat$ आणि $0 \neq j \in \Nat$ असले पाहिजेत.

त्यानंतर या व्याख्या \emph{अपेक्षेप्रमाणे} वागतात हे तपासावे लागेल;
ही तपासणी आपण \ref{sfr:arith:check:sec} कडे सोपवतो. आणि त्या खरोखर तशाच
वागतात!{} शेवटी, पूर्णांकांना परिमेय संख्या \emph{म्हणून} हाताळण्याची
काही पद्धत हवी; म्हणून प्रत्येक $i \in \Int$ साठी $i_\Rat =
\equivrep{i, 1_\Int}{\Ratequiv}$ असे आपण ठरवतो. हे सर्व सुयोग्य रीतीने
वागते याची तपासणीही आपण \ref{sfr:arith:check:sec} मध्ये करतो.

\begin{prob}
कोणत्याही $i, j \in \Int$ साठी $(i + j)_\Rat = i_\Rat+ j_\Rat$ आणि $(i \times j)_\Rat =
i_\Rat \times j_\Rat$ आणि $i \leq j \liff i_\Rat \leq j_\Rat$ हे दाखवा.
\end{prob}








\section{वास्तव संख्या-रेषा}\label{sfr:arith:real:sec}

पुढील पायरी म्हणजे परिमेय संख्यांपासून वास्तव संख्यांची रचना कशी करायची
हे दाखवणे. त्याआधी वास्तव संख्यांचे \emph{वैशिष्ट्य} नेमके काय आहे हे
समजून घ्यावे लागेल.

वास्तव संख्या परिमेय संख्यांसारख्याच बऱ्याच अंशी वागतात. (तांत्रिक
भाषेत, दोन्ही \emph{क्रमित क्षेत्रांची} उदाहरणे आहेत; त्याची व्याख्या
\ref{sfr:arith:check:orderedfield} मध्ये पाहा.) आता, तुम्ही \ref{sfr:siz::chap}
मधील सराव सोडवले असतील, तर वास्तव संख्या परिमेय संख्यांपेक्षा काटेकोरपणे
अधिक आहेत, म्हणजे $\cardless{\Rat}{\Real}$, हे तुम्हाला माहीत असेल.
हे कँटर यांनी प्रथम सिद्ध केले. परंतु अपरिमेय संख्या, म्हणजे परिमेय
नसलेल्या वास्तव संख्या, अस्तित्वात आहेत हे सुमारे अडीच सहस्रकांपासून
माहीत आहे. खरे तर:
\begin{quote}\small\textbf{स्रोतदुरुस्ती.} MRARITH-002: गोठवलेल्या इंग्रजीतील निष्क्रिय केलेली तळटीप \ensuremath{\sqrt{2}} हे धन आणि ऋण अशी दोन्ही मुळे दर्शवते असे म्हणते. पारंपरिक मूलचिन्ह मात्र मुख्य, ऋणेतर वर्गमूळ दर्शवते; म्हणून मराठीतील निष्क्रिय तळटीप दुरुस्त केली आहे. सक्रिय वाचनीय मजकुरावर याचा परिणाम नाही.\end{quote}
\begin{thm}\label{sfr:arith:real:root2irrational}
$\sqrt{2}$ ही परिमेय संख्या नाही; म्हणजे $\sqrt{2} \notin \Rat$.
\end{thm}

\begin{proof}
व्याघातासाठी $\sqrt{2}$ ही परिमेय संख्या आहे असे समजा. मग काही नैसर्गिक
संख्या $m$ आणि $n$ यांसाठी $\sqrt{2} = \nicefrac{m}{n}$. हा अपूर्णांक
यापुढे संक्षिप्त करता येणार नाही अशा रीतीने आपण $m$ आणि $n$ निवडू शकतो.
पुनर्मांडणी केल्यास $m^{2} = 2n^{2}$. येथून सिद्धता दोन प्रकारे पूर्ण
करता येते:

\emph{प्रथम, भूमितीय रीतीने} (टेनेनबाउम यांच्या पद्धतीने).\footnote{ही
सिद्धता Conway (2006) यांनी नोंदवली आहे.} पुढील चौरस पाहा:
\begin{center}
	\begin{tikzpicture}
		\draw[thick] (0,0) rectangle (3,3);
		\draw[thick, fill=red!50] (0,0) rectangle (2.3,2.3);
		\draw[thick, fill=yellow!50] (0.7,0.7) rectangle (3, 3);
		\draw[thick, fill=orange!50] (0.7,0.7) rectangle (2.3, 2.3);
		\draw[<->] (4, 0.7)--(4, 3);
		\node at (4.25, 1.85) (n) {$n$};
		\draw[<->] (5, 0)--(5, 3);
		\node at (5.25, 1.5) (m) {$m$};
	\end{tikzpicture}
\end{center}
$m^2 = 2n^2$ असल्यामुळे $n$ बाजू असलेले दोन चौरस ज्या प्रदेशात
एकमेकांवर येतात त्याचे क्षेत्रफळ, या दोन्ही चौरसांपैकी एकानेही न झाकलेल्या
प्रदेशाच्या क्षेत्रफळाइतके आहे; म्हणजे नारिंगी चौरसाचे क्षेत्रफळ हे दोन
न रंगवलेल्या चौरसांच्या क्षेत्रफळांच्या बेरजेइतके आहे. नारिंगी चौरसाची
बाजू $p$ आणि प्रत्येक न रंगवलेल्या चौरसाची बाजू $q$ असेल, तर $p^2 =
2q^2$. पण आता $p < m$, $q < n$ आणि $p, q \in \Nat$ असून
$\sqrt{2} = \nicefrac{p}{q}$. हे $m$ आणि $n$ शक्य तितके लहान निवडले
होते या तथ्याशी विसंगत आहे.

\emph{दुसरी, आकारिक रीतीने.} $m^{2} = 2n^{2}$ असल्यामुळे $m$ सम आहे.
($x$ विषम असेल, तर $x^2$ विषम असतो हे सहज दाखवता येते.) म्हणून काही
$r \in \Nat$ साठी $m = 2r$. पुनर्मांडणी केल्यास $2r^2 = n^2$,
त्यामुळे $n$ देखील सम आहे. म्हणून $m$ आणि $n$ दोन्ही सम आहेत, आणि
त्यामुळे $\nicefrac{m}{n}$ हा अपूर्णांक पुढे \emph{संक्षिप्त करता येतो}.
व्याघात!
\end{proof}

जाताजाता, या आकृतीवरील सिद्धतेमुळे पर्यायी विभाग
मधील सामग्रीचा पुन्हा विचार करता येतो. टेनेनबाउम (1927--2005) हे पूर्णपणे
आधुनिक गणितज्ञ होते; तरी ही सिद्धता निर्विवादपणे सुंदर, पूर्णतः तर्कशुद्ध
आहे आणि भूमितीय अंतःप्रज्ञेला आवाहन करते!
\begin{quote}\small\textbf{स्रोतदुरुस्ती.} MRARITH-003: गोठवलेल्या इंग्रजी मजकुरात निधनवर्ष 2006 दिले आहे. टेनेनबाउम यांच्या स्मृतिस्थळावरील चरित्र नोंद निधनाची तारीख 4 मे 2005 देते; मराठी मजकुरात 1927--2005 वापरले आहे.\end{quote}

कसेही असो, वास्तव संख्या परिमेय संख्यांपेक्षा ``अधिक विस्तृत'' आहेत.
एका अर्थाने परिमेय संख्यांमध्ये ``पोकळ्या'' आहेत आणि वास्तव संख्या त्या
भरून काढतात. वायश्ट्रास यांच्या लक्षात आले की या वर्णनातून वास्तव
संख्यांचा एकच गुणधर्म व्यक्त होतो, जो त्यांना परिमेय संख्यांपासून वेगळे
करतो; तो म्हणजे संपूर्णता गुणधर्म: \emph{ऊर्ध्वबंध असलेल्या वास्तव
संख्यांच्या प्रत्येक अरिक्त संचाला लघुतम ऊर्ध्वबंध असतो.}

परिमेय संख्यांना संपूर्णता गुणधर्म नाही हे सहज दिसते. उदाहरणार्थ,
$\sqrt{2}$ पेक्षा लहान परिमेय संख्यांचा पुढील संच पाहा:
\[	\Setabs{p \in \Rat}{p^2 < 2 \text{ किंवा }p < 0}
\]
या संचाला परिमेय संख्यांमध्ये ऊर्ध्वबंध आहे; उदाहरणार्थ, त्याचे सर्व
घटक< $3$ पेक्षा लहान आहेत. पण त्याचा लघुतम ऊर्ध्वबंध कोणता?
आपल्याला ` $\sqrt{2}$ ' असे म्हणायचे आहे; पण $\sqrt{2}$ ही \emph{परिमेय
संख्या नाही} हे आपण नुकतेच पाहिले. आणि $\sqrt{2}$ पेक्षा मोठी अशी कोणतीही
\emph{लघुतम} परिमेय संख्या नाही. म्हणून या संचाला ऊर्ध्वबंध आहे, पण लघुतम
ऊर्ध्वबंध नाही. त्यामुळे परिमेय संख्यांमध्ये संपूर्णता गुणधर्माचा अभाव आहे.

याउलट, सांतत्याला संपूर्णता गुणधर्म स्वभावतः असायला हवा. $\sqrt{2}$ ही
केवळ वास्तव संख्या असावी एवढेच आपल्याला नको; परिमेय संख्या-रेषेतील सर्व
``पोकळ्या'' भरायच्या आहेत. खरे तर सांतत्यात स्वतःमध्येच कोणतीही ``पोकळी''
नसावी असे आपल्याला हवे आहे. संपूर्णतेद्वारे नेमके हेच मिळेल.







\section{$\Rat$ पासून $\Real$ कडे}\label{sfr:arith:cuts:sec}

मूलतः, संपूर्णता गुणधर्म असे दाखवतो की वास्तव संख्या-रेषेवरील कोणताही
$\alpha$ बिंदू त्या रेषेचे अचूक दोन भाग करतो: खालच्या भागाचा $\alpha$
हा लघुतम ऊर्ध्वबंध असतो आणि वरच्या भागाचा $\alpha$ हा महत्तम निम्नबंध
असतो. परिमेय संख्यांपासून वास्तव संख्यांची \emph{रचना} करण्यासाठी
डेडेकिंड यांनी वास्तव संख्यांना परिमेय संख्यांचे विभाजन करणारे
\emph{छेद} मानण्याची सूचना केली. म्हणजेच, $< \sqrt{2}$ असलेल्या परिमेय
संख्यांना $> \sqrt{2}$ असलेल्या परिमेय संख्यांपासून वेगळे करणाऱ्या
\emph{छेदाशी} आपण $\sqrt{2}$ ची एकरूपता मानतो.

ही मांडणी नीट नेमकी करू या. परिमेय संख्यांचे दोन भाग केले, तर आपण केलेले
विभाजन केवळ त्याचा \emph{खालचा} भाग विचारात घेऊन अनन्यपणे ओळखता येते.
म्हणून नेमकी पुढील व्याख्या देऊ:

\begin{defn}[छेद]
\emph{छेद} $\alpha$ म्हणजे परिमेय संख्यांचा असा कोणताही अरिक्त उचित
आरंभीचा खंड ज्याला महत्तम घटक नाही. म्हणजे $\alpha$ हा छेद तेव्हा आणि
केवळ तेव्हाच असतो, जेव्हा:
\begin{enumerate}
	\item \emph{अरिक्त, उचित}: $\emptyset \neq \alpha \subsetneq \Rat$
	\item \emph{आरंभीचा}: सर्व $p,q \in \Rat$ साठी: $p < q \in \alpha$ असेल, तर $p \in \alpha$
	\item \emph{महत्तम घटक नाही}: प्रत्येक $p \in \alpha$ साठी $p < q$ असे काही $q \in \alpha$ असते
\end{enumerate}
मग $\Real$ हा छेदांचा संच आहे.
\end{defn}

आता $\sqrt{2} = \Setabs{p \in \Rat}{p^2 < 2\text{
किंवा }p < 0}$ असे म्हणता येते. अर्थात, हा खरोखर \emph{छेद} आहे हे
तपासले पाहिजे; ती तपासणी आपण \ref{sfr:arith:check:sec} कडे सोपवतो.

आधीप्रमाणेच, काही वस्तूंची व्याख्या केल्यानंतर त्यांच्यावरील मूलभूत
फलने आणि संबंध व्याख्यायित करावे लागतात. एका सोप्या व्याख्येपासून सुरुवात करू:
\begin{align*}
	\alpha \leq \beta \text{ तेव्हा आणि केवळ तेव्हाच }\alpha \subseteq \beta
\end{align*}
क्रमाची ही व्याख्या छेदांच्या संचाला संपूर्णता गुणधर्म आहे हा मध्यवर्ती
निष्कर्ष \emph{मांडू} देते. पूर्णपणे उलगडल्यास विधानाचे रूप असे आहे.
$S$ हा ऊर्ध्वबंध असलेला छेदांचा अरिक्त संच असेल, तर $S$ ला लघुतम
ऊर्ध्वबंध असतो. अधिक तपशीलाने: $S$ चा ऊर्ध्वबंध असणारा $\lambda$ हा छेद
असतो, म्हणजे $(\forall \alpha \in S)\alpha \subseteq \lambda$, आणि
$\lambda$ हा अशा छेदांपैकी लघुतम असतो, म्हणजे $(\forall \beta \in \Real)((\forall \alpha \in S)\alpha \subseteq \beta \lif \lambda \subseteq \beta)$. आता या निष्कर्षाची सिद्धता पाहा:

\begin{thm}\label{sfr:arith:cuts:realcompleteness}
छेदांच्या संचाला संपूर्णता गुणधर्म आहे.
\end{thm}

\begin{proof}
$S$ हा ऊर्ध्वबंध असलेला छेदांचा कोणताही अरिक्त संच असू द्या. $\lambda
= \bigcup S$ असे ठरवू.
प्रथम $\lambda$ हा छेद आहे असा दावा करू:
\begin{enumerate}
\item $S$ अरिक्त असल्यामुळे $S$ मध्ये किमान एक छेद आहे, म्हणून
$\emptyset \neq \lambda$. $S$ हा छेदांचा संच असल्यामुळे $\lambda
\subseteq \Rat$. $S$ ला ऊर्ध्वबंध असल्यामुळे प्रत्येक $\alpha \in S$
मधून अनुपस्थित असलेली अशी काही $p \in \Rat$ आहे. म्हणून $p\notin \lambda$,
आणि परिणामी $\lambda \subsetneq \Rat$.
\begin{quote}\small\textbf{स्रोतदुरुस्ती.} MRARITH-004: गोठवलेल्या इंग्रजी सिद्धतेत \ensuremath{\lambda} अरिक्त असल्याचे पहिले कारण चुकून “S has an upper bound” असे दिले आहे. आवश्यक आणि आधीच गृहीत असलेले कारण S अरिक्त असणे हे आहे; मराठी सिद्धतेत तेच वापरले आहे. पुढील उचित-उपसंच युक्तिवादासाठी मात्र ऊर्ध्वबंधाचे गृहीतक योग्य आहे.\end{quote}
\item $p < q \in \lambda$ असे समजा. मग $q \in \alpha$ असे काही
$\alpha \in S$ आहे. $\alpha$ हा छेद असल्यामुळे $p \in \alpha$.
म्हणून $p \in \lambda$.
\item $p \in \lambda$ असे समजा. मग $p \in \alpha$ असे काही
$\alpha \in S$ आहे. $\alpha$ हा छेद असल्यामुळे $p < q$ असे काही
$q \in \alpha$ आहे. म्हणून $q \in \lambda$.
\end{enumerate}
यावरून दावा सिद्ध होतो. शिवाय, स्पष्टपणे $(\forall \alpha \in S)\alpha
\subseteq \bigcup S = \lambda$, म्हणजे $\lambda$ हा $S$ चा ऊर्ध्वबंध
आहे. आता $\beta \in \mathbb{R}$ हाही ऊर्ध्वबंध आहे असे समजा, म्हणजे
$(\forall \alpha \in S)\alpha \subseteq \beta$. कोणत्याही $p \in \Rat$
साठी, $p \in \lambda$ असेल तर $p \in \alpha$ असे काही $\alpha \in S$
आहे, आणि त्यामुळे $p \in \beta$. सार्वत्रिकीकरण केल्यास $\lambda
\subseteq \beta$. म्हणून $\lambda$ हा $S$ चा \emph{लघुतम} ऊर्ध्वबंध आहे.
\end{proof}

म्हणून आपल्याकडे संपूर्णता गुणधर्म पूर्ण करणाऱ्या अनेक वस्तू आहेत.
हे सांगण्याचा एक मार्ग असा: आपल्या छेदांमध्ये ``पोकळ्या'' नाहीत.
(त्यामुळे परिमेय संख्यांऐवजी वास्तव संख्यांचे आणखी ``छेद'' घेतल्यास
काही रोचक नवी वस्तू मिळणार नाही.)

पुढे वास्तव संख्यांवरील काही क्रिया व्याख्यायित केल्या पाहिजेत. प्रत्येक
$p \in \Rat$ साठी $p_\Real = \Setabs{q \in \Rat}{q < p}$ असे ठरवून
परिमेय संख्यांना वास्तव संख्यांत अंतःस्थापित करण्यापासून सुरुवात करू.
मग पुढील व्याख्या करतो:
\begin{align*}
	\alpha + \beta &= \Setabs{p + q}{p \in \alpha \land q \in \beta}\\
	\alpha \times \beta &=
	\Setabs{p \times q}{0 \leq p \in \alpha \land 0 \leq q \in \beta} \cup 0_\Real & \text{जर }\alpha, \beta \geq 0_\Real
\end{align*}
\begin{quote}\small\textbf{स्रोतदुरुस्ती.} MRARITH-005: गोठवलेल्या इंग्रजी गुणाकार-व्याख्येत शून्याचा छेद एकदाच 0\textasciicircum{}R असा विसंगत ऊर्ध्वनिर्देशांक वापरून लिहिला आहे; त्याच परिच्छेदातील सर्व नियंत्रक संकेतन 0\_R वापरते. मराठी सूत्रात 0\_R वापरले आहे.\end{quote}
गुणाकाराची इतर प्रकरणे हाताळण्यासाठी प्रथम $0_\Real \times \alpha = 0_\Real = \alpha \times 0_\Real$ असे ठरवू आणि पुढील व्याख्या करू:
\begin{quote}\small\textbf{स्रोतदुरुस्ती.} MRARITH-006: गोठवलेल्या इंग्रजी स्रोताने शून्यासोबतच्या गुणाकाराची ही आवश्यक व्याख्या त्याच ओळीवर निष्क्रिय केली आहे; त्यामुळे एक घटक ऋण व दुसरा शून्य असलेली प्रकरणे पुढील प्रकरणनिहाय व्याख्येत येत नाहीत. स्रोताच्याच निष्क्रिय सूत्राला मराठीत सक्रिय केले आहे.\end{quote}
\begin{align*}
	-\alpha &= \Setabs{p - q}{p < 0 \land q \notin \alpha}
\end{align*}
आणि मग पुढीलप्रमाणे ठरवू:
\begin{align*}
	\alpha \times \beta &\defis
	\begin{cases}
		\mathord{-}\alpha \times \mathord{-}\beta &\text{जर }\alpha < 0_\Real\text{ आणि }\beta < 0_\Real\\
		\mathord{-}(\mathord{-}\alpha \times \beta) &\text{जर }\alpha < 0_\Real \text{ आणि }\beta > 0_\Real\\
		\mathord{-}(\alpha \times \mathord{-}\beta) &\text{जर }\alpha > 0_\Real \text{ आणि }\beta < 0_\Real
	\end{cases}
\end{align*}
यानंतर यांपैकी प्रत्येक व्याख्येतून नेहमी छेदच मिळतो हे तपासावे लागेल.
आणि शेवटी, अशा प्रकारे व्याख्यायित केलेले छेद अपेक्षेप्रमाणेच वागतात
हे दाखवणारी सोपी (पण लांबलचक) सिद्धता पूर्ण केली पाहिजे. मात्र हे सर्व
आपण \ref{sfr:arith:check:sec} कडे सोपवतो.







\section{काही तात्त्विक चिंतन}\label{sfr:arith:ref:sec}

तांत्रिक तपशील एवढेच. पण त्यातून साधले तरी काय?

एक तर, अगदी निर्विवादपणे, त्यातून आपल्याला काही {सुंदर} शुद्ध गणित
मिळाले. शिवाय काही खोल संकल्पनात्मक कामगिरीही झाली. वास्तव संख्या आणि
परिमेय संख्या यांतील निर्णायक फरक संपूर्णता गुणधर्म व्यक्त करतो, हे
ओळखणे ही एक सखोल अंतर्दृष्टी होती. तसेच, डेडेकिंड छेद म्हणून वास्तव
संख्यांची स्पष्ट रचना केल्यामुळे विश्लेषणाच्या विषयाला भक्कम आधार मिळतो.
छेदांपासून नेमके असेच क्षेत्र बनत असल्यामुळे \emph{संपूर्ण क्रमित क्षेत्र}
ही संकल्पना सुसंगत आहे, हे आपल्याला माहीत आहे.

एवढे असूनही, या कामगिरीबद्दलच्या काही शंका उघडपणे मांडल्या पाहिजेत.

पहिली शंका अशी की, वास्तव संख्यांचा परिचित (कदाचित अनंत) दशांश विस्तार
विचारात घेण्यापेक्षा छेदांच्या रूपात त्यांचा विचार करणे \emph{अधिक}
काटेकोर आहे, हे स्पष्ट नाही. दशांश विस्तारांनी वास्तव संख्यांची केलेली
ही ``रचना'' कोशी क्रमिकांद्वारे केलेल्या रचनेस काहीशी मिळतीजुळती आहे;
परंतु प्रत्यक्षात ती सतराव्या शतकाच्या प्रारंभापासूनच गणितज्ञांना मूलतः
माहीत होती (पाहा \ref{sfr:arith:cauchy:sec}). वास्तव संख्यांना संपूर्णता गुणधर्म
आहे याची जाणीव झाल्यामुळे काटेकोरपणात खरी वाढ झाली; वास्तव संख्यांची काही
विशिष्ट संच म्हणून रचना करता येणे हे कदाचित स्वतःपुरते फारसे रोचक नाही.

पूर्णांकांचे किंवा परिमेय संख्यांचे (त्याहून खूप सोपे) अंकगणितीकरण त्या
क्षेत्रांतील काटेकोरपणा वाढवते, हे तर आणखी कमी स्पष्ट आहे. येथे एक साधे
निरीक्षण उपयुक्त ठरेल. नैसर्गिक संख्यांच्या क्रमित जोड्यांचे तुल्यतावर्ग
म्हणून पूर्णांकांची \emph{रचना} केल्यानंतर, मग पूर्णांकांच्या क्रमित
जोड्यांचे तुल्यतावर्ग म्हणून परिमेय संख्यांची रचना केल्यानंतर, आणि मग
परिमेय संख्यांचे संच म्हणून वास्तव संख्यांची रचना केल्यानंतर आपण त्या
रचना लगेचच \emph{विसरून जातो}. विशेषतः, कोणालाही गणिती सिद्धतेत या
रचनांचा \emph{उपयोग} करावासा वाटणार नाही (त्या रचना अपेक्षेप्रमाणेच
वागतात हे सिद्ध करताना मात्र अर्थातच त्यांचा उपयोग करावा लागेल).
नैसर्गिक संख्यांच्या संचांच्या संचांच्या संचांच्या संचांच्या संचांच्या
संचांच्या एखाद्या संचाबद्दल बोलण्यापेक्षा वास्तव संख्येबद्दल थेट बोलणे
खूपच सोपे आहे.

या व्याख्या आपल्याला पूर्णांक, परिमेय संख्या किंवा वास्तव संख्या
\emph{नेमक्या काय आहेत}, \emph{सत्तामीमांसक दृष्ट्या सांगायचे तर}, हे
सांगतात, याबद्दल तर
सर्वाधिक शंका आहे. म्हणजे, वास्तव संख्या (उदाहरणार्थ) काही विशिष्ट संच
(संचांच्या संचांचे संच\ldots) \emph{आहेत}, हे संशयास्पद आहे. अशा
दृष्टिकोनासमोरील मुख्य अडथळा हा की ही रचना अनेक वेगवेगळ्या पद्धतींनी करता
आली असती. वास्तव संख्यांच्या बाबतीत खरोखर रोचक रीतीने भिन्न असलेल्या काही
रचना आहेत (पाहा \ref{sfr:arith:cauchy:sec}). पण काही वेगळ्या रचना मिळवण्याची
एक अगदी क्षुल्लक पद्धत अशी: \ref{sfr:rel:ref:sec} प्रमाणे आपण
क्रमित जोडीची व्याख्या किंचित वेगळी करू शकलो असतो; क्रमित जोडीची ही
पर्यायी संकल्पना वापरली असती, तर आपल्या रचना तितक्याच यशस्वी झाल्या
असत्या, पण शेवटी मिळालेल्या वस्तू वेगळ्या असत्या. त्यामुळे पूर्णांक,
परिमेय संख्या आणि वास्तव संख्या यांच्या अनेक प्रतिस्पर्धी संचसैद्धान्तिक
रचना आहेत. आता पूर्णांक (उदाहरणार्थ) \emph{त्या} संचांऐवजी \emph{हेच}
संच आहेत, असा दावा करणे केवळ मनमानीचे (आणि अडचणीचे) ठरेल.
(\ref{sfr:rel:ref:sec} प्रमाणे, Benacerraf (1965) यांनी
प्रसिद्ध केलेल्या युक्तिवादाचे हे एक उदाहरण आहे.)

आणखी एक मुद्दा उपस्थित करायला हवा: आपल्या रचनांमध्ये काहीतरी फारच
\emph{विचित्र} आहे. आपण नैसर्गिक संख्यांपासून सुरुवात केली. मग पूर्णांकांची
रचना केली आणि ``पूर्णांकांचा $0$'', म्हणजे $ \equivrep{0,0}{\Intequiv}$,
याची रचना केली. पण $0 \neq \equivrep{0,0}{\Intequiv}$. खरे तर, आपल्या
रचनांनुसार \emph{एकही} नैसर्गिक संख्या पूर्णांक नाही. पण हे सहजज्ञानाच्या
अगदी विरुद्ध वाटते. प्रत्यक्षात \ref{sfr:set:imp:sec} मध्ये आपण
फारसा युक्तिवाद न करता $\Nat \subseteq \Rat$ असा दावा केला होता. संख्या
नेमक्या \emph{काय आहेत} हे या रचना सांगत असतील, तर तो दावा उघडपणे असत्य
होता.

मग थोडे मागे हटून पाहिल्यास आपण कुठवर पोहोचलो? अनौपचारिक संचसिद्धान्तात
काम करताना आणि नैसर्गिक संख्या गृहीत धरताना, पूर्णांक, परिमेय संख्या आणि
वास्तव संख्या यांना काही विशिष्ट संच \emph{म्हणून वागवणे} आपल्याला शक्य
होते. त्या अर्थाने, या वस्तूंचे सिद्धान्त संचसिद्धान्तात
\emph{अंतःस्थापित} करता येतात. पण या अंतःस्थापनाचा तात्त्विक आशय अजिबात
सरळसोट नाही.

अर्थात, याने या विषयावरील शेवटचा शब्द सांगितलेला नाही!{} मुद्दा फक्त
एवढाच आहे. वास्तव संख्यांच्या अंकगणितीकरणाला खोल तात्त्विक महत्त्व
\emph{आहे} हे दाखवण्यासाठी आणखी काही \emph{तात्त्विक} युक्तिवाद आवश्यक असेल.









\section{क्रमित वलये आणि क्षेत्रे}\label{sfr:arith:check:sec}

या प्रकरणभर काही व्याख्या ``जशा वागायला हव्यात'' तशा वागतात, असे
आपण म्हटले. या तांत्रिक परिशिष्टात त्याचा अर्थ स्पष्ट करू आणि त्या
व्याख्या ``योग्य रीतीने'' वागतात हे कसे दाखवायचे याचा आराखडा देऊ.

\ref{sfr:arith:int:sec} मध्ये आपण $\Int$ वरील बेरीज आणि गुणाकार व्याख्यायित
केले. व्याख्येप्रमाणे या क्रिया $\Int$ ला आपण ``इच्छित'' असलेली संरचना
देतात, हे दाखवायचे आहे. विशेषतः, येथे अभिप्रेत संरचना क्रमनिरपेक्ष
वलयाची आहे.

\begin{defn}
	\emph{क्रमनिरपेक्ष वलय} म्हणजे विशिष्ट $0$ व $1$ घटक आणि $+$ व $\times$
	क्रिया असलेला असा संच $S$, जो पुढील आठ सूत्रे पूर्ण करतो:
	\begin{align*}
		\emph{सहचारिता}&&a + (b+ c) & = (a + b) + c \\
		&& (a \times b) \times c & = a \times (b\times c)\\
		\emph{क्रमनिरपेक्षता}&&a + b &= b+ a  \\
		&&  a \times b&= b\times a\\
		\emph{अविकारक}&&a + 0 &= a \\
		&& a \times 1 &= a\\
		\emph{बेरीज व्यस्त}&&(\exists b\in S)0&=a + b\\
		\emph{वितरणता}&&a \times (b+ c ) &= (a \times b) + (a \times c)
	\end{align*}
	येथे ही सर्व सूत्रे $S$ पुरते मर्यादित सार्वत्रिक परिमाणकांनी बद्ध आहेत,
	हे अध्याहृत आहे. तसेच येथील $0$ आणि~$1$ हे घटक याच नावाच्या नैसर्गिक
	संख्या असतीलच असे नाही.
\end{defn}

म्हणून पूर्णांक क्रमनिरपेक्ष वलय बनवतात हे तपासण्यासाठी या आठ अटी पूर्ण
होतात एवढेच तपासायचे आहे. यांपैकी कोणतीही अट सिद्ध करणे {अवघड} नाही,
पण ते थोडे कष्टाचे आहे. उदाहरणार्थ, बेरजेच्या बाबतीत
\emph{सहचारिता} अशी सिद्ध करता येते:

\begin{proof}
$i, j, k \in \Int$ निश्चित करा. मग $i = \equivrep{a_1, b_1}{}$ आणि $j =
\equivrep{a_2,b_2}{}$ आणि $k = \equivrep{a_3, b_3}{}$ असे होतील अशा
$a_1, b_1, a_2, b_2, a_3, b_3 \in \Nat$ संख्या आहेत. (वाचनीयतेसाठी
``$\equivrep{x, y}{\Intequiv}$'' ऐवजी ``$\equivrep{x, y}{}$'' असे लिहितो;
या संपूर्ण विभागात आपण हेच करू.) आता:
\begin{align*}
	i + (j + k) &= \equivrep{a_1, b_1}{}+(\equivrep{a_2, b_2}{} + \equivrep{a_3, b_3}{}) \\
	&= \equivrep{a_1,  b_1}{} + \equivrep{a_2+a_3, b_2+b_3}{}\\
	&= \equivrep{a_1 + (a_2 + a_3), b_1 + (b_2 + b_3)}{}\\
	&= \equivrep{(a_1 + a_2) + a_3, (b_1 + b_2) + b_3}{}\\
	&= \equivrep{a_1 + a_2, b_1 + b_2}{} + \equivrep{a_3, b_3}{}\\
	&= (\equivrep{a_1, b_1}{} + \equivrep{a_2, b_2}{}) + \equivrep{a_3, b_3}{}\\
	&= (i+j) + k
\end{align*}
येथे $\Nat$ वरील बेरजेच्या वर्तनाचा आपण मुक्तपणे उपयोग केला आहे.
\end{proof}

त्याचप्रमाणे, \emph{बेरीज व्यस्त} असा सिद्ध करता येतो:

\begin{proof}
$i \in \Int$ निश्चित करा; म्हणजे काही $a,b \in \Nat$ साठी $i =
\equivrep{a,b}{}$. $j = \equivrep{b,a}{} \in \Int$ असे ठरवू. नैसर्गिक
संख्यांच्या वर्तनाचा उपयोग केल्यास $(a+b) + 0 = 0 + (a+b)$; म्हणून
व्याख्येनुसार $\tuple{a+b, b+a} \sim_\Int \tuple{0,0}$, आणि त्यामुळे
$\equivrep{a+b, b+a}{} = \equivrep{0, 0}{} = 0_\Int$. आता $i + j =
\equivrep{a,b}{}+\equivrep{b,a}{}=\equivrep{a+b, b+a}{}= \equivrep{0,
0}{} = 0_\Int$.
\end{proof}

आणि \emph{वितरणतेची} सिद्धता अशी आहे:

\begin{proof}
वरीलप्रमाणे, $i = \equivrep{a_1, b_1}{}$ आणि $j =
\equivrep{a_2,b_2}{}$ आणि $k = \equivrep{a_3, b_3}{}$ निश्चित करा. आता:
\begin{align*}
	i \times (j + k)
	&= \equivrep{a_1, b_1}{} \times (\equivrep{a_2,b_2}{} + \equivrep{a_3, b_3}{})\\
	&= \equivrep{a_1, b_1}{} \times \equivrep{a_2 + a_3,b_2+b_3}{}\\
	&= \equivrep{a_1  (a_2 + a_3) + b_1  (b_2+b_3), a_1  (b_2 + b_3) + b_1 (a_2 + a_3)}{}\\
	&= \equivrep{a_1 a_2 + a_1a_3 + b_1 b_2+b_1b_3, a_1 b_2 + a_1b_3 + a_2b_1 + a_3b_1}{}\\
	&= \equivrep{a_1a_2 + b_1b_2, a_1b_2 + a_2b_1}{} + \equivrep{a_1a_3 + b_1b_3, a_1b_3 + a_3b_1}{}\\
	&= (\equivrep{a_1, b_1}{} \times \equivrep{a_2,b_2}{}) + (\equivrep{a_1, b_1}{} \times  \equivrep{a_3, b_3}{})\\
	&= (i \times j) + (i \times k)
\end{align*}
\end{proof}

उरलेल्या पाच अटी सिद्ध करण्याचे काम सरावासाठी सोडतो. ते पूर्ण केल्यावर
$\Int$ हे क्रमनिरपेक्ष वलय आहे, म्हणजे व्याख्यायित केलेली बेरीज आणि
गुणाकार अपेक्षेप्रमाणे वागतात, हे आपण दाखवलेले असेल.

\begin{prob}
$\Int$ हे क्रमनिरपेक्ष वलय आहे हे सिद्ध करा.
\end{prob}

पण आपले काम अजून संपलेले नाही. $\Int$ वरील बेरीज आणि गुणाकार यांच्याबरोबरच
आपण $\leq$ हा क्रमसंबंध व्याख्यायित केला होता; तो अपेक्षेप्रमाणे वागतो हेही
तपासले पाहिजे. अधिक नेमकेपणाने, $\Int$ हे \emph{क्रमित} वलय आहे हे दाखवले
पाहिजे.\footnote{\ref{sfr:rel:ord:def:linearorder} मधील व्याख्या आठवा:
	पूर्ण क्रम म्हणजे परावर्ती, संक्रमक, प्रतिसममित आणि संयोजित संबंध.
	क्रमसंबंधांच्या संदर्भात संयोजिततेला कधीकधी \emph{त्रिपर्याय} म्हणतात,
	कारण कोणत्याही $a$ आणि $b$ साठी $a \leq b \lor a
	= b \lor a \geq b$.}

\begin{defn}
\emph{क्रमित वलय} म्हणजे $\leq$ या पूर्ण क्रमसंबंधानेही सज्ज असलेले
असे क्रमनिरपेक्ष वलय की:
\begin{align*}
	a \leq b &\lif a + c \leq b + c\\
	(a \leq b \land 0 \leq c) &\lif a \times c \leq b \times c
\end{align*}
\end{defn}

\begin{prob}
$\Int$ हे क्रमित वलय आहे हे सिद्ध करा.
\end{prob}

पूर्वीप्रमाणेच, रचलेला $\Int$ हा क्रमित वलय आहे हे दाखवणे कष्टाचे पण
नित्याचे काम आहे. ते आम्ही तुमच्यासाठी सोडतो.

यामुळे पूर्णांकांचे काम संपले. आता परिमेय संख्यांबद्दल अगदी तशाच गोष्टी
दाखवायच्या आहेत. विशेषतः, $+$, $\times$ आणि $\leq$ यांच्या आपल्या
व्याख्यांनुसार परिमेय संख्या क्रमित \emph{क्षेत्र} बनवतात हे दाखवायचे आहे:
\begin{defn}\label{sfr:arith:check:orderedfield}
\emph{क्रमित क्षेत्र} म्हणजे पुढील अटही पूर्ण करणारे क्रमित वलय:
\begin{align*}
	\emph{गुणाकार व्यस्त}& & (\forall a \in S \setminus \{0\})(\exists b \in S) a\times b& = 1
\end{align*}
\end{defn}

$\Int$ हे क्रमित वलय आहे हे दाखवल्यानंतर $\Rat$ हे क्रमित क्षेत्र आहे
हे दाखवणे सोपे पण कष्टाचे आहे.

\begin{prob}
$\Rat$ हे क्रमित क्षेत्र आहे हे सिद्ध करा.
\end{prob}

पूर्णांक आणि परिमेय संख्या हाताळल्यानंतर केवळ वास्तव संख्या उरतात.
विशेषतः, $\Real$ हे \emph{संपूर्ण} क्रमित क्षेत्र, म्हणजे संपूर्णता गुणधर्म
असलेले क्रमित क्षेत्र आहे, हे दाखवले पाहिजे. \ref{sfr:arith:cuts:realcompleteness}
मध्ये $\Real$ ला संपूर्णता गुणधर्म आहे हे सिद्ध केले. मात्र $\Real$ हे
क्रमित क्षेत्र आहे याची (कंटाळवाणी) तपासणी अजून करायची आहे.
\begin{quote}\small\textbf{स्रोतदुरुस्ती.} MRARITH-007: गोठवलेल्या इंग्रजी वाक्यात “run through the (tedious) of checking” येथे आवश्यक नाम गाळले आहे. मराठीत अभिप्रेत “कंटाळवाणी तपासणी” स्पष्टपणे दिली आहे.\end{quote}

\emph{त्या} कष्टाच्या सरावाकडे वळण्यापूर्वी आणखी काही ``तात्काळ'' गोष्टी
तपासल्या पाहिजेत. उदाहरणार्थ, कोणत्याही $\alpha$ आणि $\beta$ छेदांसाठी
व्याख्यायित केलेला $\alpha + \beta$ खरोखरच \emph{छेद} आहे याची खात्री
हवी. त्याची सिद्धता अशी:

\begin{proof}
$\alpha$ आणि $\beta$ हे दोन्ही छेद असल्यामुळे $\alpha + \beta = \Setabs{p
+ q}{p \in \alpha \land q \in \beta}$ हा $\Rat$ चा अरिक्त उचित उपसंच
आहे. आता काही $p \in \alpha$ आणि $q \in \beta$ साठी $x < p + q$ असे
समजा. मग $x - p < q$, म्हणून $x - p \in \beta$, आणि $x = p + (x - p)
\in \alpha + \beta$. म्हणून $\alpha + \beta$ हा $\Rat$ चा आरंभीचा खंड
आहे. शेवटी, कोणत्याही $p + q \in \alpha + \beta$ साठी, $\alpha$ आणि
$\beta$ हे दोन्ही छेद असल्यामुळे $p < p_1$ आणि $q < q_1$ असे काही
$p_1 \in \alpha$ आणि $q_1 \in \beta$ असतात; म्हणून $p + q < p_1 + q_1
\in \alpha + \beta$; त्यामुळे $\alpha + \beta$ ला महत्तम घटक नाही.
\end{proof}

अशाच प्रयत्नांनी $\alpha - \beta$, $\alpha \times \beta$ आणि $\alpha
\div \beta$ हे छेद आहेत हे तुम्हाला तपासता येईल (शेवटच्या प्रकरणात
$\beta$ हा शून्य-छेद असलेले प्रकरण वगळून). पुन्हा, हे काम आम्ही
तुमच्यासाठीच सोडतो.
\begin{quote}\small\textbf{स्रोतदुरुस्ती.} MRARITH-008: गोठवलेल्या इंग्रजी स्रोताच्या आधीच्या cuts विभागात वजाबाकी आणि भागाकाराची सूत्रे निष्क्रिय टिप्पण्यांतच आहेत; त्यामुळे हे वाक्य सक्रियपणे पूर्ण व्याख्या न दिलेल्या क्रियांना गृहीत धरते. मराठी मजकूर दावा जतन करतो आणि ही स्रोत-अपूर्णता स्पष्ट करतो.\end{quote}

\begin{prob}
$\Real$ हे क्रमित क्षेत्र आहे हे सिद्ध करा.
\end{prob}

पण एक लहानसा अपूर्ण मुद्दा आवरायचा आहे. \ref{sfr:arith:cuts:sec} मध्ये आपण
$\sqrt{2} = \Setabs{p \in \Rat}{p < 0 \text{ किंवा }p^2 < 2}$ असे घेता
येईल असे सुचवले. मात्र हा संच \emph{छेद} आहे हे दाखवणे आवश्यक आहे.
त्याची सिद्धता अशी:

\begin{proof}
हा परिमेय संख्यांचा अरिक्त उचित आरंभीचा खंड आहे हे स्पष्ट आहे; म्हणून
त्याला महत्तम घटक नाही एवढे दाखवणे पुरेसे आहे. विशेषतः, $p$ ही $p^2 < 2$
असलेली धन परिमेय संख्या आणि $q = \frac{2p+2}{p+2}$ असेल, तर $p < q$
आणि $q^2 < 2$ हे दोन्ही दाखवणे पुरेसे आहे. $p < q$ हे पाहण्यासाठी फक्त
पुढील नोंद घ्या:
\begin{align*}
	p^2 &< 2\\
	p^2 + 2p &< 2 + 2p\\
	p(p + 2) &< 2 + 2p\\
	p &< \tfrac{2+2p}{p+2} = q
\end{align*}
$q^2 < 2$ हे पाहण्यासाठी फक्त पुढील नोंद घ्या:
\begin{align*}
	p^2 &< 2\\
	2p^2 + 4p + 2 &< p^2 + 4p+ 4\\
	4p^2 + 8p + 4 &< 2(p^2 + 4p + 4)\\
	(2p+2)^2 & <2(p+2)^2\\
	\tfrac{(2p+2)^2}{(p+2)^2} &< 2\\
	q^2 &< 2
\end{align*}
\end{proof}








\section{परिशिष्ट: कोशी क्रमिकांच्या रूपातील वास्तव संख्या}\label{sfr:arith:cauchy:sec}

\ref{sfr:arith:cuts:sec} मध्ये आपण डेडेकिंड छेद म्हणून वास्तव संख्यांची रचना
केली. या विभागात आपण एक पर्यायी रचना स्पष्ट करू. आज आपण जिला कोशी क्रमिका
म्हणतो तिची कोशी यांनी दिलेली व्याख्या या रचनेचा आधार आहे; पण त्या
व्याख्येचा उपयोग करून वास्तव संख्यांची \emph{रचना} करण्याचे श्रेय
एकोणिसाव्या शतकातील इतर लेखकांना, विशेषतः वायश्ट्रास, हाइने, मेरे आणि
कँटर यांना जाते. (या इतिहासाच्या छानशा आढाव्यासाठी
O'Connor आणि Robertson (2005) यांचा लेख पाहा.)

एकोणिसाव्या शतकाकडे वळण्यापूर्वी सायमन स्टेव्हिन (1548--1620) यांचा
विचार करणे उपयुक्त ठरेल. थोडक्यात, प्रत्येक वास्तव संख्येचा विचार तिच्या
दशांश विस्ताराच्या रूपात करता येतो हे स्टेव्हिन यांच्या लक्षात आले.
त्यामुळे $\sqrt{2}$ सारख्या अपरिमेय संख्येलाही एक व्यवस्थित दशांश विस्तार
असतो; त्याची सुरुवात अशी:
\[	1.41421356237\ldots
\]
संचसिद्धान्तात दशांश विस्तारांचे प्रतिमानीकरण करणे अगदी सोपे आहे: त्यांना
$d \colon \Nat \to \Nat$ अशी फलने माना, जिथे $d(n)$ हे आपल्याला हवे
असलेले $n$ वे दशांश स्थान आहे. मग दशांश चिन्हाच्या आधी येणारा वास्तव
संख्येचा भाग (येथे फक्त $1$) हाताळण्यासाठी या मांडणीत थोडा बदल करावा
लागेल. तसेच, उदाहरणार्थ $0.999\ldots = 1$ याची खात्री करण्यासाठी आणखी
एक बदल (तुल्यता संबंध) करावा लागेल. पण संचसिद्धान्ताच्या चौकटीत
स्टेव्हिन यांच्या पद्धतीने वास्तव संख्यांची पूर्णतः काटेकोर रचना देणे
कठीण नाही.

स्टेव्हिन हा आपला मुख्य विषय नाही. (स्टेव्हिन यांच्याविषयी अधिक माहितीसाठी
Katz आणि Katz (2012) पाहा.) पण त्याच्याशी जवळून संबंधित असा एक विचार
करू. $\sqrt{2}$ चा दशांश विस्तार थेट हाताळण्याऐवजी अधिकाधिक अचूक होत
जाणारे विस्तार घेऊन आपण $\sqrt{2}$ च्या अधिकाधिक अचूक परिमेय निकटनांची
\emph{क्रमिका} विचारात घेऊ शकतो:
\[1, 1.4, 1.414, 1.4142, 1.41421,\ldots
\]
वास्तव संख्यांचा विचार “अधिकाधिक चांगल्या निकटनांद्वारे” करता येतो या
कल्पनेतून अंतर्दृष्टींची आणखी एक मालिका मिळते (जशी आपण $\Int$ पासून
$\Rat$ किंवा $\Nat$ पासून $\Int$ रचताना वापरली होती). त्या मूलभूत
अंतर्दृष्टी अशा:
\begin{enumerate}
	\item प्रत्येक वास्तव संख्या एका (कदाचित अनंत) दशांश विस्ताराने
	लिहिता येते.
	\item एका (कदाचित अनंत) दशांश विस्तारात सांकेतिक केलेली माहिती
	परिमेय संख्यांच्या क्रमिकेनेही सांकेतिक करता येते.
	\item परिमेय संख्यांची क्रमिका ही $\Nat$ पासून $\Rat$ कडे जाणारे
	फलन मानता येते; क्रमिकेतील $n$ वी परिमेय संख्या $f(n)$ माना.
\end{enumerate}
अर्थात, $\Nat$ पासून $\Rat$ कडे जाणारे \emph{कोणतेही} फलन आपल्याला वास्तव
संख्या देईल असे नाही. उदाहरणार्थ, हे फलन पाहा:
\[	f(n) =\begin{cases}
			1 & \text{जर }n\text{ विषम असेल}\\
			0 &\text{जर }n\text{ सम असेल}
		\end{cases}
\]
येथील मुख्य अडचण अशी की $0,1,0,1,0,1,0,\ldots$ ही क्रमिका कोणत्याही
वास्तव संख्येकडे जवळ जाताना दिसत नाही. म्हणून एखाद्या वास्तव संख्येकडे
जवळ जाणाऱ्या क्रमिकांचाच विचार करण्यासाठी आपले लक्ष काही \emph{सीमा}
असलेल्या क्रमिकांपुरते मर्यादित केले पाहिजे.

सीमेची कल्पना पर्यायी विभाग मध्ये आपण यापूर्वी पाहिली
आहे. पण तेथे वापरलेली \emph{अगदी तीच} व्याख्या येथे वापरता येणार नाही.
तेथील “$(\forall \epsilon>0)$” या पदावलीत वास्तव संख्यांवरील
परिमाणन गृहीत धरले होते; आणि आपण वास्तव संख्यांवरील फलनांच्या सीमा
विचारात घेत होतो. म्हणून ती व्याख्या येथे वापरणे म्हणजे वास्तव संख्या
आधीच उपलब्ध मानणे होईल; पण नेमकी त्यांचीच आपल्याला \emph{रचना} करायची
आहे. सुदैवाने, सीमेच्या अगदी जवळच्या कल्पनेने आपण काम करू शकतो.
\begin{defn}\label{sfr:arith:cauchy:def:CauchySequence}
  $f: \Nat \to \Rat$ हे फलन \emph{कोशी क्रमिका} आहे तेव्हाच आणि
  तेव्हाच, जेव्हा प्रत्येक धन $\epsilon \in \Rat$ साठी $(\exists \ell
  \in \Nat)(\forall m, n > \ell)|f(m) - f(n)| < \epsilon$.
\end{defn}

सीमेची व्यापक कल्पना पूर्वीसारखीच आहे: यथार्थतेची एखादी पातळी
(जिचे मापन~$\epsilon$ ने केले आहे) हवी असेल, तर पाहण्यासाठी एक
“प्रदेश” असतो ($\ell$ पेक्षा मोठे कोणतेही आदान). आपल्या $1$, $1.4$,
$1.414$, $1.4142$, $1.41421$\ldots या क्रमिकेला सीमा आहे हेही सहज
दिसते: $\sqrt{2}$ चे निकटन $\nicefrac{1}{10^{n}}$ इतक्या त्रुटीच्या
आत हवे असेल, तर $n$ व्या पदानंतरचे कोणतेही पद पाहा.

मग वास्तव संख्या ही कोणतीही कोशी क्रमिका \emph{आहेच} असे म्हणण्याचा
विचार साहजिक वाटेल. पण $\Int$ आणि $\Rat$ यांच्या रचनांप्रमाणेच तो फार
भोळा ठरेल: कोणत्याही एका वास्तव संख्येचा निर्देश अनेक भिन्न कोशी क्रमिका
करतात. हे पाहण्याची एक सोपी पद्धत अशी. $f$ ही कोशी क्रमिका दिली असता,
$g(0)\neq f(0)$ एवढा अपवाद सोडून $g$ हे $f$ सारखेच फलन ठरवा. पहिल्या
संख्येनंतर दोन्ही क्रमिका सर्वत्र जुळत असल्यामुळे
\ref{sfr:arith:cauchy:def:CauchySequence} मध्ये वापरलेल्या अर्थाने त्यांची सीमा समान
आहे, आणि म्हणून त्या एकाच वास्तव संख्येची “व्याख्या” करतात, असे
आपल्याला शेवटी म्हणायचे असेल. त्यामुळे या कोशी क्रमिका एकच वास्तव
संख्या आहेत असे आपण खरे तर मानले पाहिजे.

म्हणून कोशी क्रमिकांवर पुन्हा एक तुल्यता संबंध परिभाषित करून वास्तव
संख्यांची ओळख त्या संबंधाखालील तुल्यतावर्गांशी करावी लागेल.
\begin{quote}\small\textbf{स्रोतदुरुस्ती.} MRARITH-009: गोठवलेल्या इंग्रजी मजकुरात वास्तव संख्यांची ओळख equivalence relations शी करावी असे म्हटले आहे; पुढील वाक्ये आणि रचना स्पष्टपणे equivalence classes वापरतात. मराठी मजकुरात अपेक्षित तुल्यतावर्ग वापरले आहेत.\end{quote}
सर्वप्रथम ज्याची सीमा $0$ आहे अशा फलनाची कल्पना हवी. कोणत्याही $h :
\Nat \to \Rat$ या फलनाविषयी असे म्हणा की \emph{$h$ ची सीमा $0$ आहे}
तेव्हाच आणि तेव्हाच, जेव्हा प्रत्येक धन $\epsilon \in \Rat$ साठी
$(\exists \ell \in \Nat)(\forall n > \ell)|h(n)| < \epsilon$.
\footnote{याची तुलना पर्यायी विभाग मधील $\lim_{x
\mathord{\rightarrow}\infty}f(x) = 0$ या व्याख्येशी करा.}
शिवाय $f$ आणि $g$ ही $\Nat \to \Rat$ अशी फलने असतील, तर
$(f-g)(n) = f(n) - g(n)$ असे माना. आता अशी व्याख्या करा:
\[	f \Realequiv g \text{ तेव्हाच आणि तेव्हाच जेव्हा $(f-g)$ ची सीमा $0$ आहे}.
\]
$\Realequiv$ हा तुल्यता संबंध आहे हे तपासावे लागेल; आणि तो तसा आहे.
मग, हवे असल्यास, $\Nat \to \Rat$ अशा सर्व कोशी क्रमिकांचे
$\Realequiv$ खालील तुल्यतावर्ग म्हणजे वास्तव संख्या, अशी व्याख्या करता
येईल.

\begin{prob}
प्रत्येक $n$ साठी $f(n) = 0$ असे माना. $g(n) =
\frac{1}{(n+1)^2}$ असे माना. ही दोन्ही फलने कोशी क्रमिका आहेत, आणि
दोन्ही फलनांची सीमा $0$ आहे, हे दाखवा; त्यामुळे $f \sim_\Real g$
हेही दाखवा.
\end{prob}

हे केल्यानंतर, नेहमीप्रमाणे, $f$ हा घटक असलेल्या तुल्यतावर्गासाठी
आपण $\equivrep{f}{\Realequiv}$ असे लिहू. मात्र मजकूर वाचनीय ठेवण्यासाठी
पुढे अधोलिखित भाग वगळून फक्त $\equivrep{f}{}$ असे लिहू. प्रत्येक $q
\in \Rat$ साठी $q_{\Real} = \equivrep{c_{q}}{}$ असेही आपण ठरवतो,
जिथे प्रत्येक $n \in \Nat$ साठी $c_q(n) = q$ असे स्थिर फलन $c_{q}$
आहे. मग वास्तव संख्यांवरील मूलभूत संबंध आणि क्रिया परिभाषित करतो;
उदाहरणार्थ:
\begin{align*}
	\equivrep{f}{} + \equivrep{g}{} &= 	\equivrep{(f + g)}{} \\
	\equivrep{f}{} \times 	\equivrep{g}{} &= \equivrep{(f \times g)}{}
\end{align*}
जिथे $(f + g)(n) = f(n) + g(n)$ आणि $(f \times g)(n) = f(n) \times
g(n)$. $f$ आणि $g$ कोशी क्रमिका असतील तेव्हा $(f + g)$, $(f-g)$ आणि
$(f\times g)$ यांपैकी प्रत्येकही कोशी क्रमिका आहे हे अर्थात तपासावे
लागेल; तसे ते आहेत, आणि ही सिद्धता तुमच्यासाठी सोडली आहे.
\begin{quote}\small\textbf{स्रोतदुरुस्ती.} MRARITH-012: गोठवलेला मजकूर येथे बेरीज, वजाबाकी आणि गुणाकाराने कोशी क्रमिका मिळतात एवढीच संवृति तपासायला सांगतो. तुल्यतावर्गांवरील क्रिया सुव्याख्यायित होण्यासाठी प्रतिनिधी बदलल्यावर वर्ग बदलत नाही हेही तपासणे आवश्यक आहे; तो स्वतंत्र तपशील स्रोतामध्ये दिलेला नाही.\end{quote}

शेवटी क्रमाची संकल्पना परिभाषित करू. $\equivrep{f}{}$ हा वर्ग
\emph{धन} आहे तेव्हाच आणि तेव्हाच, जेव्हा $\equivrep{f}{}\neq 0_\Real$
आणि $(\exists \ell \in \Nat)(\forall n > \ell)0 < f(n)$
या दोन्ही अटी पूर्ण होतात. मग $\equivrep{f}{} <
\equivrep{g}{}$ तेव्हाच आणि तेव्हाच, जेव्हा $\equivrep{(g - f)}{}$
धन आहे, असे म्हणा.
\begin{quote}\small\textbf{स्रोतदुरुस्ती.} MRARITH-010: गोठवलेल्या इंग्रजीतील धनतेच्या व्याख्येत वास्तव तुल्यतावर्गाची तुलना 0\_Rat शी केली आहे. येथे त्याच रचनेतील वास्तव शून्य 0\_Real वापरले आहे.\end{quote}
ही व्याख्या सुव्याख्यायित आहे, म्हणजे ती तुल्यतावर्गातून निवडलेल्या
“प्रतिनिधी” फलनावर अवलंबून नाही, हे तपासावे लागेल.
हे केल्यानंतर या व्याख्यांमुळे योग्य बैजिक गुणधर्म मिळतात हे दाखवणे
बरेच सोपे आहे; म्हणजे:
\begin{thm}\label{sfr:arith:cauchy:thm:cauchyorderedfield}
कोशी क्रमिकांचे वरील तुल्यतावर्ग एक क्रमित क्षेत्र बनवतात.
\end{thm}
\begin{quote}\small\textbf{स्रोतदुरुस्ती.} MRARITH-011: गोठवलेल्या इंग्रजीतील प्रमेय, पुढील प्रश्न आणि संपूर्णतेचे विधान Cauchy sequences हेच क्षेत्र बनवतात अशी संक्षिप्त भाषा वापरतात. प्रत्यक्ष क्षेत्राचे घटक Realequiv संबंधाखालील तुल्यतावर्ग आहेत; मराठी मजकुरात ते स्पष्ट केले आहे. पुढील सिद्धतेत जेव्हा क्रमिकेवर वास्तव-क्रम लावला जातो तेव्हा तिचा तुल्यतावर्ग अभिप्रेत आहे.\end{quote}

\begin{proof}
सराव.
\end{proof}

\begin{prob}
कोशी क्रमिकांचे तुल्यतावर्ग एक क्रमित क्षेत्र बनवतात हे सिद्ध करा.
\end{prob}

अशा रीतीने रचलेल्या वास्तव संख्यांना संपूर्णता गुणधर्म आहे हे सिद्ध
करणे अधिक कठीण आहे; म्हणून त्याची सिद्धता आपण देऊ.

\begin{thm}
ऊर्ध्वबंध असलेल्या कोशी क्रमिकांच्या तुल्यतावर्गांच्या प्रत्येक अरिक्त
संचाला लघुतम ऊर्ध्वबंध असतो.
\end{thm}

\begin{proof}[सिद्धतेचा आराखडा]
$S$ हा ऊर्ध्वबंध असलेल्या कोशी क्रमिकांचा कोणताही अरिक्त संच असू द्या.
म्हणून काही $p \in \Rat$ असे आहे की $p_{\Real}$ हा $S$ चा ऊर्ध्वबंध
आहे. $r \in S$ असू द्या; मग काही $q \in \Rat$ असे आहे की
$q_{\Real} < r$. त्यामुळे $S$ ला लघुतम ऊर्ध्वबंध असेल, तर तो
$q_\Real$ आणि $p_\Real$ यांच्या दरम्यान (दोन्ही टोकांसह) असेल.

लघुतम ऊर्ध्वबंधाकडे खालून आणि वरून एकाच वेळी जवळ जात आपण तो शोधू.
विशेषतः, $f, g \colon \Nat \to \Rat$ अशी दोन फलने परिभाषित करू:
$f$ ने वरून आणि $g$ ने खालून लघुतम ऊर्ध्वबंधाकडे जवळ जावे हा उद्देश
आहे. सुरुवात अशा व्याख्यांनी करू:
\begin{align*}
	f(0) &= p \\
	g(0) &= q
\end{align*}
मग $a_n = \frac{f(n) + g(n)}{2}$ असेल, तर पुढील व्याख्या
करा:\footnote{ही पुनरावर्ती व्याख्या आहे. पण पुनरावर्ती व्याख्या ग्राह्य
आहेत असे मानण्याचे कोणतेही कारण आपण \emph{अजून} दिलेले नाही.}
\begin{align*}
	f(n+1) &=
	\begin{cases}
		a_n &\text{जर }(\forall h \in S)\equivrep{h}{} \leq (a_n)_\Real\\
		f(n)&\text{अन्यथा}
	\end{cases}\\
	g(n+1) &=
	\begin{cases}
		a_n &\text{जर }(\exists h \in S)\equivrep{h}{} \geq (a_n)_\Real\\
	 	g(n) &\text{अन्यथा}
	\end{cases}
\end{align*}
$f$ आणि $g$ या दोन्ही कोशी क्रमिका आहेत. (हे बऱ्यापैकी सहज तपासता
येते, पण तो सराव तुमच्यासाठी सोडला आहे.) प्रत्येक पायरीवर $f$ आणि
$g$ मधील फरक निम्मा होत असल्यामुळे $(f-g)$ या फलनाची सीमा $0$ आहे.
त्यामुळे $\equivrep{f}{} = \equivrep{g}{}$.

सर्वप्रथम $\equivrep{f}{}$ हा $S$ चा ऊर्ध्वबंध आहे, म्हणजे
$(\forall h \in S)\equivrep{h}{} \leq \equivrep{f}{}$, हे दाखवू.
(वाटेत आपण \ref{sfr:arith:cauchy:thm:cauchyorderedfield} वापरू.) $h \in S$ असू द्या
आणि व्याघातासाठी $\equivrep{f}{} < \equivrep{h}{}$ असे समजा; म्हणजे
$0_\Real < \equivrep{(h-f)}{}$. $f$ ही एकस्वनिक घटती कोशी क्रमिका
असल्यामुळे काही $n \in \Nat$ असे आहे की
$\equivrep{(c_{f(n)} - f)}{} < \equivrep{(h-f)}{}$. म्हणून:
\[	(f(n))_\Real = \equivrep{c_{f(n)}}{} < \equivrep{f}{} + \equivrep{(h-f)}{} = \equivrep{h}{},
\]
पण रचनेनुसार $\equivrep{h}{} \leq (f(n))_\Real$ असल्यामुळे हा
व्याघात आहे.

आता $\equivrep{f}{} = \equivrep{g}{}$ हाच $S$ चा \emph{लघुतम}
ऊर्ध्वबंध आहे हे दाखवू. $j$ ही कोणतीही कोशी क्रमिका असू द्या आणि
$\equivrep{j}{} < \equivrep{g}{}$ असे समजा. वरच्याप्रमाणेच युक्तिवाद
केल्यास ($g$ हे \emph{वाढते} आहे या तथ्याचा उपयोग करून) काही $n \in
\Nat$ असे आहे की $\equivrep{j}{} < (g(n))_\Real$. पण रचनेनुसार काही
$h \in S$ असे आहे की $(g(n))_\Real \leq \equivrep{h}{}$; म्हणून
$\equivrep{j}{} < \equivrep{h}{}$. त्यामुळे $\equivrep{j}{}$ हा
$S$ चा ऊर्ध्वबंध नाही.
\end{proof}



\chapter{अनंत संच}\label{sfr:infinite::chap}
\begin{editorial}
अनंत संचांवरील हे प्रकरण टिम बटन यांच्या \emph{Open Set Theory}
या ग्रंथातून घेतले आहे.
\end{editorial}




\section{हिल्बर्ट यांचे हॉटेल}\label{sfr:infinite:hilbert:sec}

नैसर्गिक संख्यांचा संच स्पष्टपणे अनंत आहे. म्हणून, नैसर्गिक संख्या
स्वतःच \emph{आयत्या गृहीत धरायच्या} नसतील, तर नैसर्गिक संख्यांचा उल्लेख
न करता अनंत संचाचे वैशिष्ट्य सांगणे ही आपली पहिली पायरी असली पाहिजे.
सन 1924 मधील एका व्याख्यानात हिल्बर्ट यांनी मांडलेला पुढील दृष्टिकोन
सुबक आहे. ते आपल्याला अशी कल्पना करायला सांगतात:
\begin{quote}
[\ldots] सांत संख्येने खोल्या असलेले एक हॉटेल. या सर्व खोल्यांपैकी
प्रत्येक खोलीत नेमका एक पाहुणा असावा. आता पाहुण्यांनी कशाही प्रकारे
आपापल्या खोल्या बदलल्या, [पण] प्रत्येक खोलीत अजूनही एकापेक्षा अधिक
व्यक्ती राहणार नाही अशा रीतीने, तर एकही खोली मोकळी होणार नाही आणि
हॉटेलमालकाला या मार्गाने नव्याने येणाऱ्या पाहुण्यासाठी नवी जागा निर्माण
करता येणार नाही [\ldots
\textparagraph \ldots]

आता हॉटेलमध्ये $1$, $2$, $3$, $4$, $5$, \dots अशा क्रमांकांच्या अनंत
खोल्या असतील आणि प्रत्येक खोलीत नेमका एक पाहुणा असेल, असे आपण ठरवू.
नवा पाहुणा येताच मालकाने प्रत्येक जुन्या पाहुण्याला त्याच्या खोलीच्या
क्रमांकापेक्षा एकाने मोठा क्रमांक असलेल्या खोलीत हलवायचे असते; मग नव्याने
आलेल्या पाहुण्यासाठी खोली~$1$ मोकळी होईल.
	\begin{center}
		\begin{tikzpicture}[scale = .75]
		\foreach \x in {1, 2, 3, 4, 5, 6, 7, 8, 9}
		{
			\node (\x a) at (\x, 1) {\small{\x}};
			\node (\x b) at (\x, 2) {\small{\x}};
		}
		\node (dotsa) at (10, 1) {\small{\ldots}};
		\node (dotsb) at (10, 2) {\small{\ldots}};
		\draw[->] (1b)--(2a);
		\draw[->] (2b)--(3a);
		\draw[->] (3b)--(4a);
		\draw[->] (4b)--(5a);
		\draw[->] (5b)--(6a);
		\draw[->] (6b)--(7a);
		\draw[->] (7b)--(8a);
		\draw[->] (8b)--(9a);
		\draw[->] (9b)--(dotsa);
		\draw (1,1) circle (.4);
		\end{tikzpicture}
	\end{center}
	(Ewald आणि Sieg (2013, p.~730) मध्ये प्रकाशित; अनुवाद आमचा)
\end{quote}
निर्णायक मुद्दा असा की हिल्बर्ट यांच्या हॉटेलमध्ये अनंत खोल्या आहेत;
आणि असे म्हणण्याचा अर्थ काय, याची व्याख्या म्हणून आपण त्यांचे स्पष्टीकरण
घेऊ शकतो. खरे तर हाच डेडेकिंड यांचा दृष्टिकोन होता (अर्थातच येथे फार
मोठा कालविपर्यास करून तो मांडला आहे; डेडेकिंड यांची व्याख्या
1888 मधील आहे):

\begin{defn}\label{sfr:infinite:hilbert:defn:DedekindInfinite}
संच $A$ हा \emph{डेडेकिंड-अनंत} असतो जर आणि तरच $A$ पासून~$A$ च्या
एखाद्या उचित उपसंचाकडे एकास-एक फलन असते. म्हणजे, असा काही $o \in A$
आणि एकास-एक फलन $f \colon A \to A$ असतात की $o \notin \ran{f}$.
\end{defn}







\section{डेडेकिंड बीजसंरचना}\label{sfr:infinite:dedekind:sec}

नैसर्गिक संख्या केवळ अनंत असाव्यात असेच आपल्याला वाटत नाही; त्यांना काही
(बीजगणितीय) गुणधर्म असावेत असेही आपल्याला वाटते: बेरीज, गुणाकार इत्यादी
क्रियांखाली त्यांनी योग्य रीतीने वागले पाहिजे.

\emph{उत्तरवर्ती फलनाची} कल्पना मूलभूत मानून, पुढील गुणधर्म असलेल्या
वस्तू म्हणजे संख्या असे त्यांचे वैशिष्ट्य सांगावे, ही डेडेकिंड यांची
कल्पना होती:
\begin{enumerate}
	\item अशी $0$ ही संख्या आहे जी कोणत्याही संख्येची उत्तरवर्ती नाही
	\\म्हणजे, $0 \notin \ran{s}$
	\\म्हणजे, $\forall x\ s(x) \neq 0$
	\item भिन्न संख्यांच्या उत्तरवर्ती संख्या भिन्न असतात
	\\म्हणजे, $s$ हे एकास-एक फलन आहे
	\\म्हणजे, $\forall x \forall y (s(x) = s(y) \lif x = y)$
	\item\label{sfr:infinite:dedekind:repeatedapplication} प्रत्येक संख्या $0$ पासून
	उत्तरवर्ती फलन वारंवार लावून मिळते.
\end{enumerate}
पहिल्या दोन अटी प्रथम-क्रम तर्कशास्त्र वापरून सहज हाताळता येतात (वर
पाहा). पण केवळ प्रथम-क्रम तर्कशास्त्र वापरून \ref{sfr:infinite:dedekind:repeatedapplication}
ही अट हाताळता येत नाही. \ref{sfr:infinite:dedekind:repeatedapplication} ही अट पुढीलप्रमाणे
संचसैद्धान्तिक रूपात पुन्हा मांडणे ही डेडेकिंड यांची निर्णायक कामगिरी होती:
\begin{enumerate}
	\item[3$'$.] नैसर्गिक संख्यांचा संच हा \emph{उत्तरवर्ती फलनाखाली
	बंद} असलेला लघुतम संच आहे: म्हणजे, संचातील कोणत्याही घटक वर
	$s$ लावल्यास संचातीलच दुसरा घटक मिळतो.
\end{enumerate}
पण हे आपल्याला सावकाश आणि स्पष्टपणे मांडावे लागेल.

\begin{defn}\label{sfr:infinite:dedekind:Closure}
	कोणत्याही फलन $f$ साठी, संच $X$ हा $f$-\emph{बंद} असतो जर आणि तरच
	$(\forall x \in X)f(x) \in X$. आता कोणत्याही $o$ साठी अशी व्याख्या करू:
	$$\closureofunder{f}{o} = \bigcap\Setabs{X}{o \in X\text{ आणि }X
\text{ हा $f$-बंद आहे}}$$
\end{defn}

म्हणून $\closureofunder{f}{o}$ हा $o$ हा घटक असलेल्या सर्व
$f$-बंद संचांचा छेद आहे. त्यामुळे सहजज्ञानानुसार,
$\closureofunder{f}{o}$ हा $o$ हा घटक असलेला \emph{लघुतम}
$f$-बंद संच आहे. पुढील निष्कर्ष हे सहजज्ञान नेमकेपणाने मांडतो:
\begin{lem}\label{sfr:infinite:dedekind:closureproperties}
	कोणतेही फलन $f$ आणि कोणतेही $o \in A$ यांसाठी:
	\begin{enumerate}
		\item\label{sfr:infinite:dedekind:closurehaselem} $o \in \closureofunder{f}{o}$; आणि
		\item\label{sfr:infinite:dedekind:closureclosed} $\closureofunder{f}{o}$ हा $f$-बंद आहे; आणि
		\item\label{sfr:infinite:dedekind:closuresmallest} $X$ हा $f$-बंद असून $o \in
		X$ असेल, तर $\closureofunder{f}{o} \subseteq X$
	\end{enumerate}
\end{lem}

\begin{proof}
$o$ हा घटक असलेला किमान एक $f$-बंद संच आहे, हे लक्षात घ्या;
तो म्हणजे $\ran{f}\cup \{o\}$. म्हणून अशा \emph{सर्व} संचांचा छेद
$\closureofunder{f}{o}$ अस्तित्वात आहे. आता आपल्याला
\ref{sfr:infinite:dedekind:closurehaselem}--\ref{sfr:infinite:dedekind:closuresmallest} तपासावे लागतील.

\ref{sfr:infinite:dedekind:closurehaselem} बाबत: छेदातील सर्व संचांत $o$ हा घटक
असल्यामुळे $o \in \closureofunder{f}{o}$.

\ref{sfr:infinite:dedekind:closureclosed} बाबत: $x \in \closureofunder{f}{o}$ असे समजा.
म्हणून $o \in X$ आणि $X$ हा $f$-बंद असेल, तर $x \in X$; आणि $X$ हा
$f$-बंद असल्यामुळे $f(x) \in X$. म्हणून
$f(x) \in \closureofunder{f}{o}$.

\ref{sfr:infinite:dedekind:closuresmallest} बाबत: सर्वसाधारणपणे, $X \in C$ असेल, तर
$\bigcap C \subseteq X$.
\end{proof}

याचा उपयोग करून आपण असे म्हणू शकतो:

\begin{defn}
संच $A$, फलन $f \colon A \to A$ आणि असा काही $o \in A$ यांनी मिळून
\emph{डेडेकिंड बीजसंरचना} बनते, जर:
	\begin{enumerate}
		\item \label{sfr:infinite:dedekind:ded:proper} $o \notin \ran{f}$
		\item \label{sfr:infinite:dedekind:ded:injection} $f$ हे एकास-एक फलन आहे
		\item \label{sfr:infinite:dedekind:ded:closure} $A = \closureofunder{f}{o}$
	\end{enumerate}
\end{defn}

$A = \closureofunder{f}{o}$ असल्यामुळे, आधीच्या निष्कर्षावरून $A$ हा
$o$ हा घटक असलेला लघुतम $f$-बंद संच आहे. डेडेकिंड बीजसंरचना
स्पष्टपणे डेडेकिंड-अनंत असते; व्याख्येतील \ref{sfr:infinite:dedekind:ded:proper} आणि
\ref{sfr:infinite:dedekind:ded:injection} ही कलमे पाहा. पण कोणत्याही डेडेकिंड-अनंत संचाचे
डेडेकिंड बीजसंरचनेत रूपांतर करता येते, ही अधिक रोचक गोष्ट आहे.

\begin{thm}\label{sfr:infinite:dedekind:thm:DedekindInfiniteAlgebra}
डेडेकिंड-अनंत संच अस्तित्वात असेल, तर डेडेकिंड बीजसंरचना अस्तित्वात असते.
\end{thm}

\begin{proof}
$D$ हा डेडेकिंड-अनंत असू द्या. मग एकास-एक फलन $g \colon D \to D$
आणि घटक $o  \in D \setminus \ran{g}$ अस्तित्वात आहेत. आता $A =
\closureofunder{g}{o}$ असे मानू; \ref{sfr:infinite:dedekind:closureproperties} नुसार $A$
अस्तित्वात आहे आणि $o \in A$. $f = \funrestrictionto{g}{A}$ असे मानू.
$A, f, o$ हे मिळून डेडेकिंड बीजसंरचना बनतात, हे आपण दाखवू.

\ref{sfr:infinite:dedekind:ded:proper} बाबत: $o \notin \ran{g}$ आणि $\ran{f}
\subseteq \ran{g}$, म्हणून $o\notin \ran{f}$.

\ref{sfr:infinite:dedekind:ded:injection} बाबत: $g$ हे $D$ वरील एकास-एक फलन आहे; म्हणून
$f \subseteq g$ हेही एकास-एक फलन असले पाहिजे.

\ref{sfr:infinite:dedekind:ded:closure} बाबत: \ref{sfr:infinite:dedekind:closureproperties} नुसार $A$ हा
$g$-बंद आहे; त्यामुळे विशेषतः $A$ हा $f$-बंद आहे. म्हणून
\ref{sfr:infinite:dedekind:closureproperties} नुसार $\closureofunder{f}{o} \subseteq A$.
तसेच $\closureofunder{f}{o}$ हा $f$-बंद आहे आणि
$f = \funrestrictionto{g}{A}$ असल्यामुळे $\closureofunder{f}{o}$ हा
$g$-बंद आहे. म्हणून \ref{sfr:infinite:dedekind:closureproperties} नुसार
$A \subseteq \closureofunder{f}{o}$.
\end{proof}








\section{डेडेकिंड बीजसंरचना आणि अंकगणितीय विगमन}\label{sfr:infinite:induction:sec}

आता निर्णायक बाब अशी की एखादी डेडेकिंड बीजसंरचना---खरे तर,
\emph{कोणतीही} डेडेकिंड बीजसंरचना---नैसर्गिक संख्यांचे पर्यायी प्रतिरूप
म्हणून उपयोगी पडेल. हे पुढील सहज निष्पत्तीमुळे घडते:

\begin{thm}[अंकगणितीय विगमन]\label{sfr:infinite:induction:thm:dedinfiniteinduction}
$N, s, o$ यांनी मिळून डेडेकिंड बीजसंरचना बनते असे समजा. मग कोणत्याही
संच $X$ साठी:
	\begin{center}
		जर $o \in X$ आणि $(\forall {n} \in N \cap X){s}({n}) \in X$, {तर} $N \subseteq X$.
	\end{center}
\end{thm}

\begin{proof}
डेडेकिंड बीजसंरचनेच्या व्याख्येनुसार $N = \closureofunder{s}{o}$.
आता ${o} \in X$ आणि $(\forall {n} \in N)(n \in X \lif
{s}({n}) \in X)$ या दोन्ही अटी पूर्ण होत असतील, तर
$N = \closureofunder{s}{o} \subseteq X$.
\end{proof}

विगमन हा नैसर्गिक संख्यांचा वैशिष्ट्यपूर्ण गुणधर्म असल्यामुळे मुद्दा असा
आहे: कोणताही डेडेकिंड-अनंत संच दिला असता, आपण डेडेकिंड बीजसंरचना बनवू
शकतो आणि ती बीजसंरचना नैसर्गिक संख्यांचे पर्यायी प्रतिरूप म्हणून वापरू
शकतो.

हे मान्यच की \ref{sfr:infinite:induction:thm:dedinfiniteinduction} मध्ये विगमन
\emph{संचसैद्धान्तिक} भाषेत मांडले आहे. पण हे तत्त्व अधिक परिचित वाटेल
अशा भाषेत आपण सहज मांडू शकतो:

\begin{cor}\label{sfr:infinite:induction:natinductionschema}
$N, s, o$ यांनी मिळून डेडेकिंड बीजसंरचना बनते असे समजा. मग प्राचले असू
शकतील अशा कोणत्याही सूत्र $\phi(x)$ साठी:
\begin{center}
	जर $\phi(o)$ आणि $(\forall {n} \in N)(\phi(n)\lif
	\phi({s}({n})))$, {तर} $(\forall n \in N)\phi(n)$
\end{center}
\end{cor}

\begin{proof}
$X = \Setabs{n \in N}{\phi(n)}$ असे मानून आता
\ref{sfr:infinite:induction:thm:dedinfiniteinduction} वापरा.
\end{proof}

या निष्कर्षात आपण सूत्राला ``प्राचले असण्याबद्दल'' बोललो. याचा ढोबळ
अर्थ असा की कोणत्याही वस्तू $c_1, \ldots, c_k$ साठी आपण
$\phi(x, c_1, \ldots, c_k)$ बरोबर काम करू शकतो. अधिक नेमकेपणाने,
``प्राचले'' असा उल्लेख न करता निष्कर्ष पुढीलप्रमाणे मांडता येतो.
ज्याची सर्व मुक्त चरे दाखवली आहेत अशा कोणत्याही सूत्र
$\phi(x, v_1, \ldots, v_k)$ साठी:
	\begin{align*}
			\forall v_1 \ldots \forall v_k((&\phi(o, v_1,\ldots, v_k) \land {}\\
			&	(\forall x \in N)(\phi(x,v_1, \ldots, v_k) \lif \phi(s(x), v_1,\ldots, v_k))) \lif {}\\
			&\hspace{3em} (\forall x \in N)\phi(x, v_1,\ldots, v_k))
	\end{align*}
``प्राचले असणे'' असे म्हणाल्याने मजकूर वाचणे फारच सोपे होते, हे उघड आहे.
(पर्यायी विभाग मध्ये आपण ही युक्ती वारंवार वापरू.)

पुन्हा डेडेकिंड बीजसंरचनांकडे वळू: कोणतीही डेडेकिंड बीजसंरचना दिली
असता, बेरीज, गुणाकार आणि घातांक क्रिया यांची नेहमीची अंकगणितीय फलनेही
आपण परिभाषित करू शकतो. मात्र हे क्षुल्लक नाही आणि त्यात
\emph{पुनरावर्ती व्याख्येचे} तंत्र वापरावे लागते. हे तंत्र आपण बरेच नंतर,
त्यापेक्षा अधिक व्यापक संदर्भात मांडू आणि समर्थित करू.
(उत्सुक वाचकांनी पर्यायी विभाग नंतर याकडे पुन्हा
वळावे किंवा Potter (2004, pp.~95--8) यांसारखा पर्यायी विवेचनग्रंथ
वाचावा.) पण $N, s, o$ यांनी डेडेकिंड बीजसंरचना बनत असेल, तर शेवटी
आपल्याला पुढीलप्रमाणे अर्थ ठरवता येईल:
\begin{align*}
	{a} + {o} &= {a} & & & {a} \times {o} &= {o} & & & {a}^{o} &= s(o)\\
{a} + {s}({b}) &= {s}({a}+{b}) &&& {a} \times {s}({b}) &= ({a}\times {b}) + {a}  & & & {a}^{{s}({b})} &= {a}^{b} \times {a}
\end{align*}
आणि ही फलने आपल्या अपेक्षेप्रमाणे वागतात, हे दाखवता येईल.









\section{अनंत संचाच्या अस्तित्वाची
डेडेकिंड यांची ``सिद्धता''}\label{sfr:infinite:dedekindsproof:sec}

या प्रकरणात आपण डेडेकिंड बीजसंरचनांच्या साहाय्याने नैसर्गिक संख्यांचे
संचसैद्धान्तिक विवेचन दिले. \ref{sfr:arith:ref:sec} मध्ये आपण इतर
गोष्टींबरोबर विश्लेषणाच्या अंकगणितीकरणाच्या तात्त्विक महत्त्वाचा विचार
केला. आता येथे आपण जे साध्य केले आहे त्याच्या महत्त्वाचा विचार करायला हवा.

संपूर्ण \ref{sfr:arith::chap} मध्ये आपण नैसर्गिक संख्या आयत्या
मानल्या आणि त्यांचा वापर करून पूर्णांक, परिमेय संख्या आणि वास्तव संख्या
यांची स्पष्ट रचना केली. या प्रकरणात आपण नैसर्गिक संख्यांची स्पष्ट रचना
दिलेली नाही. आपण एवढेच दाखवले आहे की \emph{कोणताही डेडेकिंड-अनंत संच
दिला असता}, आपण असा संच परिभाषित करू शकतो जो आपल्याला~$\Nat$ ने जसे
वागावेसे वाटते अगदी तसाच वागेल.

त्यामुळे \emph{कोणत्या वस्तू नैसर्गिक संख्या आहेत} अशा सत्तामीमांसक
प्रश्नाचे उत्तर आपण दिले आहे, असा दावा अर्थातच करता येणार नाही. पण ही
चांगलीच गोष्ट आहे. शेवटी, \ref{sfr:arith:ref:sec} मध्ये आपण यावर
भर दिला होता की $\Real$ म्हणजे डेडेकिंड छेदांचा संच ही व्याख्या सोयीची
संकेतजनक व्याख्या न मानता \emph{शोध} मानणे चुकीचे ठरेल. डेडेकिंड छेद
वास्तव संख्यांच्या प्रमुख गणितीय गुणधर्मांचे उदाहरण देतात, हे निर्णायक
निरीक्षण आहे. येथेही तसेच: \emph{कोणतीही} डेडेकिंड बीजसंरचना नैसर्गिक
संख्यांच्या प्रमुख गणितीय गुणधर्मांचे उदाहरण देते, हे निर्णायक निरीक्षण
आहे. (खरे तर सर्व डेडेकिंड बीजसंरचना परस्पर \emph{समरूपी} आहेत, हे
सिद्ध करून डेडेकिंड यांनी हा मुद्दा ठामपणे दाखवला
(1888, प्रमेये 132--3). त्यामुळे आजचे अनेक
``संरचनावादी'' डेडेकिंड यांना पूर्वसूरी मानतात, यात आश्चर्य नाही.)

 शिवाय आपण नैसर्गिक संख्यांचा सिद्धान्त अशा अनौपचारिक साध्या
 संचसिद्धान्तात कसा अंतःस्थापित करायचा हे दाखवले आहे, जो स्वतः अजूनही
 बराच अनौपचारिक आहे, पण नैसर्गिक संख्या आयत्या गृहीत धरत नाही असे
 (वरवर तरी) दिसते. त्यामुळे डेडेकिंड यांचा स्वतःचा महत्त्वाकांक्षी
 प्रकल्प प्रत्यक्षात आणण्याच्या मार्गावर आपण असू शकतो. त्यांनी तो असा
 स्पष्ट केला:
\begin{quote}
	विज्ञानात सिद्ध करण्याजोगी कोणतीही गोष्ट सिद्धतेशिवाय मान्य करू नये.
	ही मागणी रास्त वाटत असली, तरी अगदी साध्या विज्ञानाचा पाया घालण्याच्या
	सर्वांत अलीकडील पद्धतींमध्येही ती पूर्ण झाली आहे असे मला मानता येत
	नाही; म्हणजे, संख्यासिद्धान्ताशी संबंधित तर्कशास्त्राच्या त्या भागातही.
	अंकगणित (बीजगणित, विश्लेषण) हे तर्कशास्त्राचाच केवळ एक भाग आहे असे
	म्हणताना माझा आशय असा की संख्या ही संकल्पना अवकाश आणि काळ यांच्या
	कल्पना किंवा अंतःप्रज्ञा यांपासून पूर्णपणे स्वतंत्र आहे---ती विचारांच्या
	शुद्ध नियमांचे थेट फलित आहे असेच मी मानतो.
	(Dedekind, 1888, प्रस्तावना)
\end{quote}
डेडेकिंड यांची धाडसी कल्पना अशी आहे. आपण नुकतेच पाहिले की केवळ
(अनौपचारिक) संचसिद्धान्त वापरून नैसर्गिक संख्या कशा घडवायच्या.
\ref{sfr:arith::chap} मध्ये आपण पाहिले की नैसर्गिक संख्या आणि काही
संचसिद्धान्त दिला असता वास्तव संख्यांची रचना कशी करायची. त्यामुळे कदाचित
``अंकगणित (बीजगणित, विश्लेषण)'' हे ``तर्कशास्त्राचाच केवळ एक भाग'' ठरते
(डेडेकिंड यांच्या ``तर्कशास्त्र'' या शब्दाच्या विस्तारित अर्थाने).

ही झाली कल्पना. पण क्षणभर थांबा. डेडेकिंड बीजसंरचनेची आपली रचना
(म्हणजे नैसर्गिक संख्यांचे आपले पर्यायी प्रतिरूप) डेडेकिंड-अनंत संचाच्या
अस्तित्वावर अवलंबून आहे. (मागे \ref{sfr:infinite:dedekind:thm:DedekindInfiniteAlgebra}
पाहा.) डेडेकिंड-अनंत संचाचे अस्तित्व ``तर्कशास्त्र'' किंवा ``विचारांचे
शुद्ध नियम'' यांच्या साहाय्याने प्रस्थापित करता आले नाही, तर हा प्रकल्प
अडून पडतो.

मग डेडेकिंड-अनंत संचाचे अस्तित्व ``विचारांच्या शुद्ध नियमांनी''
प्रस्थापित \emph{करता येते का}? डेडेकिंड यांचा प्रयत्न असा होता:
\begin{quote}
  माझे स्वतःचे विचारविश्व, म्हणजे माझ्या विचारांचे विषय होऊ शकणाऱ्या सर्व
  वस्तूंची समग्रता $S$, अनंत आहे. कारण $s$ हा~$S$ चा घटक दर्शवत असेल,
  तर $s$ हा माझ्या विचाराचा विषय होऊ शकतो हा विचार $s'$ स्वतः~$S$ चा
  घटक आहे. याला घटक~$s$ ची प्रतिमा $\phi(s)$ मानले, तर
  \dots~$S$ हा [डेडेकिंडच्या अर्थाने] अनंत आहे, हेच सिद्ध करायचे होते.
	(Dedekind, 1888, \S66)
\end{quote}
मुख्यतः अत्यंत काटेकोर गणितीय सिद्धतांनी भरलेल्या ग्रंथाच्या मध्यभागी
अशी गोष्ट सापडणे खरोखरच चकित करणारे आहे. येथे दोन टिप्पण्या कराव्याशा
वाटतात.

पहिली: या ``सिद्धतेला'' आज आपण ``गणितीय'' स्वरूप म्हणून ओळखू असे स्वरूप
जवळजवळ नाहीच. ती मानसशास्त्रीय वस्तूंबद्दल (विचारांबद्दल) बोलते आणि
त्यातही केवळ \emph{संभाव्य} विचारांबद्दल बोलते.

दुसरी: निदान आपण डेडेकिंड बीजसंरचना ज्या रीतीने मांडल्या आहेत त्या
रीतीनुसार, या ``सिद्धतेत'' एक सरळ तांत्रिक उणीव आहे. डेडेकिंड यांचा
युक्तिवाद यशस्वी झाला, तरी तो केवळ अनंत वस्तू (विशेषतः अनंत विचार)
अस्तित्वात आहेत एवढेच प्रस्थापित करतो. पण~$S$ ला अनेक घटक असलेला
एकच \emph{संच} का मानावे,~$S$ ला केवळ \emph{काही वस्तू} (अनेकवचनात) का मानू
नये, याचे कारणही डेडेकिंड यांनी द्यायला हवे.

डेडेकिंड यांना येथे उणीव दिसली नाही, यावरून कदाचित त्यांनी वापरलेला
``समग्रता'' हा शब्द \emph{आपण} वापरत असलेल्या ``संच'' या शब्दाशी
तंतोतंत जुळत नाही असे सूचित होते.\footnote{तसे तो जुळत नव्हता असे
मानण्याची इतर कारणेही आपल्याकडे आहेत; पाहा Potter (2004, p.~23).}
पण हे फार आश्चर्यकारक ठरणार नाही. मागील दोन प्रकरणांत आपण राबवलेला
प्रकल्प---नैसर्गिक संख्यांची ``रचना'' आणि त्यांच्यापासून पूर्णांक,
वास्तव संख्या व परिमेय संख्या यांची ``रचना''---संपूर्णपणे अनौपचारिक
रीतीने राबवला आहे. नेमके कोणते संच कधी अस्तित्वात असतात याविषयी अनेक
सामान्य तत्त्वे न मांडता, गरज पडेल तेव्हा आपण हा संच किंवा तो संच आयता
मानला. पण काही सामान्य तत्त्वांची आपल्याला \emph{गरज} आहे हे माहीत आहे;
नाहीतर आपल्याला रसेलच्या विरोधापत्तीला सामोरे जावे लागेल.

 आता अनौपचारिकतेच्या या मर्यादा ओलांडण्याची वेळ आली आहे.








\section{परिशिष्ट: श्रेडर--बर्नस्टाइनची सिद्धता}\label{sfr:infinite:card-sb:sec}

अनौपचारिक संचसिद्धान्ताचा निरोप घेण्यापूर्वी, आणखी एका अनौपचारिक
(पण प्रगत!) सिद्धतेचा विचार करायचा आहे. ही श्रेडर--बर्नस्टाइनची सिद्धता
आहे (\ref{sfr:siz:sb:thm:schroder-bernstein}): जर
$\cardle{A}{B}$ आणि $\cardle{B}{A}$, तर $\cardeq{A}{B}$; म्हणजे,
एकास-एक फलने $f \colon A \to B$ आणि $g \colon B \to A$ दिली असता
एकास-एक आच्छादन $h \colon A \to B$ असते.

या प्रकरणात आपण डेडेकिंड यांच्या \emph{संवरणांच्या} संकल्पनेचा वापर
केला. किंबहुना याच संकल्पनेचा उपयोग करून डेडेकिंड यांनी
श्रेडर--बर्नस्टाइनची एक सुंदर सिद्धता दिली; ती येथे मांडू. ही सिद्धता
Potter (2004, pp.~157--8) यांच्या विवेचनाला अगदी जवळून अनुसरते;
तेथील मांडणी किंचित वेगळी पण मूलतः समान आहे. इंटरनेटवर थोडा शोध
घेतल्यास असेही दिसेल की---$\sqrt{2}$ च्या अपरिमेयतेप्रमाणेच---या
प्रमेयाच्या \emph{अनेक} रोचक आणि वेगवेगळ्या सिद्धता आहेत.

\ref{sfr:infinite:dedekind:Closure} मधील संकेतनासारखेच संकेतन
वापरून, प्रत्येक संच $B$ आणि फलन~$f$ यांसाठी
\[\Closureofunder{f}{B} = \bigcap \Setabs{X}{B \subseteq X
\text{ आणि $X$ हा $f$-बंद आहे}}
\]
असे मानू. अशा रीतीने परिभाषित केलेला $\Closureofunder{f}{B}$ हा~$B$
चा समावेश असलेला लघुतम $f$-बंद संच आहे; म्हणजे:

\begin{lem}\label{sfr:infinite:card-sb:Closureprops}
	कोणतेही फलन $f$ आणि कोणताही $B$ यांसाठी:
	\begin{enumerate}
		\item\label{sfr:infinite:card-sb:Closurehaselem} $B \subseteq \Closureofunder{f}{B}$; आणि
		\item\label{sfr:infinite:card-sb:Closureclosed} $\Closureofunder{f}{B}$ हा $f$-बंद आहे; आणि
		\item\label{sfr:infinite:card-sb:Closuresmallest} $X$ हा $f$-बंद असून $B
		\subseteq X$ असेल, तर $\Closureofunder{f}{B} \subseteq X$.
	\end{enumerate}
\end{lem}

\begin{proof}
\ref{sfr:infinite:dedekind:closureproperties} मधील सिद्धतेप्रमाणेच.
\end{proof}

श्रेडर--बर्नस्टाइनपर्यंत पोहोचण्यासाठी आपल्याला आणखी एका तथ्याची गरज आहे:

\begin{prop}\label{sfr:infinite:card-sb:sbhelper}
जर $A \subseteq B \subseteq C$ आणि $A \approx C$, तर $\cardeq{B}{C}$.
\end{prop}
\begin{quote}\small\textbf{स्रोतदुरुस्ती.} MRINF-003 पात्रता: गोठवलेल्या इंग्रजीतील अंतःस्थ cardeq रचना $A\approx B\approx C$ अशी वैध रीतीने विस्तारते; ती स्रोतदोष नाही. मराठीतील स्पष्ट $B\approx C$ निष्कर्ष त्याच तुल्यबलता-साखळीतील, पुढील उपयोगासाठी केंद्रित केलेली सममूल्य मांडणी आहे.\end{quote}

\begin{proof}
एकास-एक आच्छादन  $f \colon C \to A$ दिले असता, $F =
\Closureofunder{f}{C \setminus B}$ असे मानू आणि प्रांत $C$ असलेले फलन
$g$ पुढीलप्रमाणे परिभाषित करू:
\[	g(x) =
	\begin{cases}
			f(x) &\text{जर $x \in F$}\\
			x & \text{अन्यथा}
	\end{cases}
\]
$g$ हे $C \to B$ असे एकास-एक आच्छादन आहे, हे आपण दाखवू. त्यावरून
$\comp{f^{-1}}{g} \colon A \to B$ हे एकास-एक आच्छादन आहे आणि सिद्धता
पूर्ण होते.

प्रथम आपला दावा असा आहे की $x \in F$ पण $y\notin F$ असेल, तर $g(x) \neq
g(y)$. विरोधाभासाने सिद्ध करण्यासाठी तसे नाही असे समजा; मग $y = g(y) =
g(x) = f(x)$. $x \in F$ आहे आणि \ref{sfr:infinite:card-sb:Closureprops} नुसार $F$ हा
$f$-बंद आहे, म्हणून $y = f(x) \in F$; हा विरोधाभास आहे.

आता $g(x) = g(y)$ असे समजा. वरील दाव्यानुसार, $x \in F$ तेव्हा आणि
केवळ तेव्हाच, जेव्हा $y \in F$. जर $x, y \in F$, तर $f(x) = g (x) =
g(y) = f(y)$; आणि $f$ हे एकास-एक आच्छादन असल्यामुळे $x = y$. जर
$x, y \notin F$, तर $x = g(x) = g(y) = y$. म्हणून $g$ हे
एकास-एक फलन आहे.

$\ran{g} = B$ हे दाखवणे उरते. म्हणून $x \in B \subseteq C$ स्थिर करू.
जर $x \notin F$, तर $g(x) = x$. जर $x \in F$, तर कोणत्यातरी
$y \in F$ साठी $x = f(y)$; कारण तसे नसते, तर $F \setminus \{x\}$ हा
$f$-बंद संच असता आणि त्यात $C\setminus B$ चा समावेश असता; पण
\ref{sfr:infinite:card-sb:Closureprops} नुसार हे अशक्य आहे. आता $g(y) = f(y) = x$.
\end{proof}

शेवटी, मुख्य निष्कर्षाची सिद्धता पुढीलप्रमाणे. फलन $h$ आणि संच $D$ दिले
असता आपण $\funimage{h}{D} = \Setabs{h(x)}{x \in D}$ अशी व्याख्या करतो,
हे आठवा.

\begin{proof}[श्रेडर--बर्नस्टाइनची सिद्धता] $f \colon A \to B$
आणि $g \colon B \to A$ ही एकास-एक फलने असू द्या.
$\funimage{f}{A} \subseteq B$ असल्यामुळे
$\funimage{g}{\funimage{f}{A}} \subseteq g[B] \subseteq A$.
तसेच, $\comp{f}{g} \colon A \to
\funimage{g}{\funimage{f}{A}}$ हे एकास-एक फलन आहे, कारण $g$ आणि
$f$ दोन्ही तशी आहेत; आणि त्याचा सहप्रांत ज्या रीतीने परिभाषित केला आहे,
त्यावरून $\comp{f}{g}$ हे खरे तर एकास-एक आच्छादन आहे. म्हणून
$\cardeq{\funimage{g}{\funimage{f}{A}}}{A}$; आणि त्यामुळे
\ref{sfr:infinite:card-sb:sbhelper} नुसार $h \colon A \to \funimage{g}{B}$ असे
एकास-एक आच्छादन आहे. शिवाय, $g^{-1}$ हे
$\funimage{g}{B} \to B$ असे एकास-एक आच्छादन आहे. म्हणून
$\comp{h}{g^{-1}} \colon A \to B$ हे एकास-एक आच्छादन आहे.
\end{proof}



\chapter{विधानीय तर्कशास्त्र: विन्यासमीमांसा आणि चिन्हार्थमीमांसा}\label{pl:syn::chap}
\begin{editorial}
  या भागात अभिजात विधानीय तर्कशास्त्रावरील सामग्री आहे. पहिले
  प्रकरण तुलनेने प्राथमिक स्वरूपाचे आहे आणि त्यात केवळ व्याख्या
  व निष्पत्त्या दिल्या आहेत; अनेक सिद्धता पूर्ण केलेल्या नसून त्या
  सरावासाठी सोडल्या आहेत. सिद्धता-पद्धतींवरील आणि संपूर्णता
  प्रमेयावरील सामग्री प्रथम-क्रम तर्कशास्त्राच्या भागातून, ``FOL''
  हा टॅग false असा ठेवून, समाविष्ट केली आहे. त्यामुळे विधेये,
  पदे आणि संख्यापक यांच्याशी संबंधित सर्व मजकूर वगळला जातो;
  तसेच संरचना~$\Struct{M}$ या संरचनांविषयीच्या चर्चेऐवजी
  सत्य-मूल्यांकने~$\pAssign{v}$ या सत्य-मूल्यांकनांविषयी चर्चा येते.

  या भागात अधिक तपशील समाविष्ट करण्यासाठी त्याचा विस्तार करणे,
  आणि सत्यता-फलनात्मक संपूर्णतेसारखे पुढील विषय व निष्पत्त्या जोडणे,
  असे नियोजन आहे.
\end{editorial}
\begin{editorial}
  हा केवळ व्याख्यांचा अतिशय संक्षिप्त आढावा आहे. सूत्रांवरील
  विगमनाने सिद्धता करण्याचा सुलभ परिचय आणि आणखी बरीच उदाहरणे
  देऊन त्याचा विस्तार करायला हवा.
\end{editorial}






\section{परिचय}\label{pl:syn:int:sec}

विधानीय तर्कशास्त्रात विधानीय संयोजके
$\lnot$, $\land$, $\lor$, $\lif$ आणि $\liff$ वापरून
विधानीय चलांपासून बनवलेल्या सूत्रांचा विचार
केला जातो. अंतःप्रेरणेने पाहता, विधानीय चल~$p$
हे सत्य किंवा असत्य असलेल्या एखाद्या वाक्याच्या किंवा विधानाच्या
जागी येते. सूत्र मधील विधानीय चलाचे
``सत्यतामूल्य'' ठरले की, विधानीय संयोजके वापरून त्यांच्यापासून
बनवलेल्या कोणत्याही सूत्रांचे सत्यतामूल्यही ठरते. विधानीय
तर्कशास्त्राला आपण \emph{सत्यता-फलनात्मक} म्हणतो, कारण त्याची
चिन्हार्थमीमांसा सत्यतामूल्यांच्या फलनांनी दिली जाते. विशेषतः,
विधानीय तर्कशास्त्रात सत्य आणि असत्य यांचे पुढील प्रकारचे निर्धारण
आपण विचारात घेत नाही: उदाहरणार्थ, एखादी गोष्ट केवळ परिस्थितीवश
सत्य असण्याऐवजी अनिवार्यपणे सत्य आहे का, ती सत्य आहे हे ज्ञात आहे
का, किंवा ती आत्ता सत्य आहे की पूर्वी सत्य होती अथवा पुढे सत्य
असेल. आपण फक्त सत्य ($\True$) आणि असत्य ($\False$) ही दोन
सत्यतामूल्ये विचारात घेतो. म्हणून एखादे विधान सत्यही नसते व
असत्यही नसते, किंवा ते केवळ अर्धसत्य असते, ही शक्यता चर्चेतून
वगळतो. तसेच, ज्यांच्यापासून बनवलेल्या सूत्राचे सत्यतामूल्य
त्याच्या भागांच्या सत्यतामूल्यांनी पूर्णपणे ठरते (त्याच्या अर्थाने,
उदाहरणार्थ, ठरत नाही), अशाच संयोजकांवर आपण लक्ष केंद्रित करतो.
विशेषतः, इंग्रजीतील सशर्त विधानांचे सत्यतामूल्य या अर्थाने
सत्यता-फलनात्मक असते की नाही, हा वादग्रस्त प्रश्न आहे. वास्तविक
अभिव्यंजन~$\lif$ सत्यता-फलनात्मक आहे; सत्यता-फलनात्मक नसलेल्या सशर्त
विधानांचा विचार इतर तर्कपद्धती करतात.

सत्यता-फलनात्मक विधानीय तर्कशास्त्राची उपपत्ती आणि अधिउपपत्ती
विकसित करण्यासाठी, आपण प्रथम त्याच्या व्यंजकांची विन्यासमीमांसा
आणि चिन्हार्थमीमांसा परिभाषित केली पाहिजे. संयोजके वापरून
विधानीय चलांपासून सूत्रे तयार करण्याची एक
पद्धत आपण वर्णन करू. पर्यायी व्याख्या शक्य आहेत. इतर पद्धती वेगळी
चिन्हे निवडतील, संयोजकांचे वेगळे संच मूलभूत म्हणून निवडतील आणि
कंसांचा वेगळ्या रीतीने वापर करतील (किंवा तथाकथित पोलिश
संकेतपद्धतीप्रमाणे कंस अजिबात वापरणार नाहीत). तरी सर्व दृष्टिकोनांत
सामाईक बाब ही की रचनानियम सूत्रांचा संच \emph{विगमनाने}
परिभाषित करतात. हे योग्य रीतीने केल्यास, रचनानियमांनुसार प्रत्येक
व्यंजक मूलतः एकाच रीतीने तयार होऊ शकते. त्यामुळे तयार होणारी
व्यंजके \emph{एकमेव रीतीने वाचनीय} असतात. परिणामी, त्यांची ही
विगमनाधारित व्याख्या आपल्याला त्याच पद्धतीने---विगमनाधारित
व्याख्येने---त्या व्यंजकांना अर्थ देता येतो, याची खात्री देते.

व्यंजकांना अर्थ देणे हा चिन्हार्थमीमांसेचा विषय आहे. विधानीय
तर्कशास्त्राच्या चिन्हार्थमीमांसेतील मध्यवर्ती संकल्पना म्हणजे
सत्य-मूल्यांकन मधील पूर्ति. सत्य-मूल्यांकन~$\pAssign{v}$ हे
विधानीय चलांना $\True$, $\False$ ही सत्यतामूल्ये
नेमून देते. कोणतेही सत्य-मूल्यांकन कोणत्याही सूत्र~$\varphi$
साठी $\pValue{v}(\varphi)$ हे सत्यतामूल्य ठरवते. सूत्राची
सत्य-मूल्यांकन~$\pAssign{v}$ मध्ये पूर्ति तेव्हा आणि केवळ तेव्हाच
होते जेव्हा $\pValue{v}(\varphi) = \True$; हे आपण $\pSat{v}{\varphi}$ असे
लिहितो. तार्किक संयोजकांची सत्यता-फलने वापरून, उदाहरणार्थ
$\varphi \land \psi$ ची पूर्ति $\varphi$ आणि~$\psi$ यांची पूर्ति होते (किंवा
होत नाही) यावरून परिभाषित करून, हा संबंधही~$\varphi$ च्या रचनेवरील
विगमनाने परिभाषित करता येतो.

वाक्यांसाठी असलेल्या $\pSat{v}{\varphi}$ या पूर्ति-संबंधाच्या आधारे
आपण सर्वतः सत्यता, तार्किक निष्पन्नता आणि पूर्ततायोग्यता या मूलभूत
चिन्हार्थविषयक संकल्पना परिभाषित करू शकतो. सूत्र हे सर्वतः
सत्य सूत्र, $\Entails \varphi$, असते, जर प्रत्येक सत्य-मूल्यांकन त्याची
पूर्ति करत असेल, म्हणजे कोणत्याही~$\pAssign{v}$ साठी
$\pValue{v}(\varphi) = \True$ असेल. सूत्रांच्या एखाद्या संचापासून,
$\Gamma \Entails \varphi$, ते तार्किकरीत्या निष्पन्न होते, जर~$\Gamma$
मधील सर्व सूत्रांची पूर्ति करणारे प्रत्येक सत्य-मूल्यांकन
~$\varphi$ चीही पूर्ति करत असेल. तसेच, सूत्रांचा एखादा संच
पूर्ततायोग्य असतो, जर एखादे सत्य-मूल्यांकन त्यातील सर्व सूत्रांची एकाच वेळी पूर्ति करत असेल. सूत्रे विगमनाने
परिभाषित केलेली असल्यामुळे, आणि पूर्तिही सूत्रांच्या रचनेवरील
विगमनाने परिभाषित केलेली असल्यामुळे, आपल्या चिन्हार्थमीमांसेचे
गुणधर्म सिद्ध करण्यासाठी आणि परिभाषित केलेल्या चिन्हार्थविषयक
संकल्पनांचा परस्परसंबंध दाखवण्यासाठी आपण विगमन वापरू शकतो.









\section{विधानीय सूत्र}\label{pl:syn:fml:sec}

विधानीय तर्कशास्त्रातील सूत्रे ही
\emph{विधानीय चले}, विधानीय
  स्थिर~$\lfalse$
    आणि विधानीय स्थिर~$\ltrue$ यांपासून \emph{तार्किक
  संयोजके} वापरून तयार होतात.

\begin{enumerate}
\item प्रगणनीय संच~$\PVar$: त्यातील विधानीय चले
  $\Obj p_0$, $\Obj p_1$, \dots
\item असत्यता साठीचे विधानीय स्थिर~$\lfalse$.
\item सत्यता साठीचे विधानीय स्थिर~$\ltrue$.
\item तार्किक संयोजके:
  \startycommalist
  \ycomma $\lnot$ (नकरण)%
  \ycomma $\land$ (संधियोग)%
  \ycomma $\lor$ (विकल्पयोग)%
  \ycomma $\lif$ (शर्त)%
  \ycomma $\liff$ (द्विशर्त)%
\item विरामचिन्हे: (, ), आणि स्वल्पविराम.
\end{enumerate}

विधानीय तर्कशास्त्राची ही भाषा आपण $\Lang L_0$ ने दर्शवतो.

.




\begin{intro}
वर आपण वापरलेल्या संज्ञांपेक्षा आणि चिन्हांपेक्षा वेगळ्या संज्ञा व
चिन्हे तुम्हाला परिचित असू शकतात. तर्कशास्त्राच्या ग्रंथांत (आणि
अध्यापनात) ``नकरण'' यासाठी सामान्यतः $\sim$, $\neg$, आणि~!\ यांपैकी
एखादे, तर ``संधियोग'' यासाठी $\wedge$, $\cdot$, आणि $\&$ यांपैकी
एखादे चिन्ह वापरतात. ``सशर्त (संकेतार्थक)'' किंवा ``अभिव्यंजन'' यांसाठी
$\rightarrow$, $\Rightarrow$, आणि $\supset$ ही चिन्हे सामान्यतः
वापरली जातात.
``द्विसशर्त'', ``द्वि-अभिव्यंजन'' किंवा
  ``(वास्तविक) सममूल्यता'' यांसाठी $\leftrightarrow$,
  $\Leftrightarrow$, आणि $\equiv$ ही चिन्हे वापरतात.
$\lfalse$ या चिन्हाला वेगवेगळ्या ठिकाणी
  ``असत्यता'', ``फाल्सम'', ``विसंगती'' किंवा ``तळ'' म्हणतात.
$\ltrue$ या चिन्हाला वेगवेगळ्या ठिकाणी
  ``सत्यता'', ``वेरम'' किंवा ``शिखर'' म्हणतात.
\end{intro}

\begin{defn}[सूत्र]
\label{pl:syn:fml:defn:formulas}
विधानीय तर्कशास्त्रातील \emph{सूत्रे} चा संच~$\Frm[L_0]$
पुढीलप्रमाणे विगमनाने परिभाषित केला जातो:
\begin{enumerate}
\item $\lfalse$ हे आण्विक सूत्र आहे.

\item $\ltrue$ हे आण्विक सूत्र आहे.

\item प्रत्येक विधानीय चल~$\Obj p_i$ हे आण्विक
  सूत्र आहे.

\item $\varphi$ हे सूत्र असेल, तर $\lnot \varphi$ हे
  सूत्र आहे.

\item $\varphi$ आणि $\psi$ ही सूत्रे असतील, तर $(\varphi \land
  \psi)$ हे सूत्र आहे.

\item $\varphi$ आणि $\psi$ ही सूत्रे असतील, तर $(\varphi \lor \psi)$
  हे सूत्र आहे.

\item $\varphi$ आणि $\psi$ ही सूत्रे असतील, तर $(\varphi \lif \psi)$
  हे सूत्र आहे.

\item $\varphi$ आणि $\psi$ ही सूत्रे असतील, तर $(\varphi \liff \psi)$
  हे सूत्र आहे.

\item इतर काहीही सूत्र नाही.
\end{enumerate}
\end{defn}

\begin{explain}
सूत्रांची ही व्याख्या \emph{विगमनाधारित व्याख्या} आहे. मूलतः,
आपण सूत्रांचा संच अनंतपणे अनेक टप्प्यांत तयार करतो. आरंभीच्या
टप्प्यात आपण सर्व आण्विक सूत्रांना सूत्रे घोषित करतो; हे व्याख्येतील
पहिल्या काही बाबींशी, म्हणजे
$\ltrue$, %
$\lfalse$, %
$\Obj p_i$ यांच्या बाबींशी जुळते. म्हणून ``आण्विक सूत्र''
म्हणजे या रूपातील कोणतेही सूत्र.

व्याख्येतील इतर बाबी आधीच तयार केलेल्या सूत्रांपासून नवी
सूत्रे तयार करण्याचे नियम देतात. दुसऱ्या टप्प्यात या नियमांनी
आण्विक सूत्रांपासून सूत्रे तयार करता येतात. तिसऱ्या
टप्प्यात आण्विक सूत्रे आणि दुसऱ्या टप्प्यात मिळालेली सूत्रे यांपासून
नवी सूत्रे तयार करतो, आणि ही प्रक्रिया पुढे चालू राहते. अशा एखाद्या
टप्प्यात शेवटी तयार होणारी प्रत्येक गोष्ट सूत्र असते; इतर
काहीही सूत्र नसते.
\end{explain}

द्विस्थानी संयोजक~$\ast$ वापरून $\psi$, $\chi$ पासून तयार केलेले
सूत्र $(\psi \ast \chi)$ लिहिताना आपण अनेकदा सर्वांत बाहेरची कंसांची
जोडी वगळून केवळ~$\psi \ast \chi$ असे लिहू.



\begin{defn}[विन्यासात्मक अनन्यता]
$\ident$ हे चिन्ह चिन्हमालांमधील विन्यासात्मक अनन्यता व्यक्त करते;
म्हणजे $\varphi \ident \psi$ तेव्हा आणि केवळ तेव्हाच, जेव्हा $\varphi$ आणि $\psi$
या समान लांबीच्या चिन्हमाला असतात आणि प्रत्येक स्थानी त्यांच्यात तेच
चिन्ह असते.
\end{defn}

$\ident$ या चिन्हाच्या दोन्ही बाजूंना जोडणीने मिळालेल्या चिन्हमाला
येऊ शकतात. उदाहरणार्थ, $\varphi \ident (\psi \lor \chi)$ याचा अर्थ असा:
आरंभीचा कंस, चिन्हमाला~$\psi$, $\lor$ हे चिन्ह, चिन्हमाला~$\chi$ आणि
शेवटचा कंस या क्रमाने जोडून मिळणारी चिन्हमाला आणि चिन्हमाला~$\varphi$
या एकच आहेत. असे असेल, तर $\varphi$ चे पहिले चिन्ह आरंभीचा कंस आहे;
$\varphi$ मध्ये दुसऱ्या चिन्हापासून सुरू होणारी $\psi$ ही उपचिन्हमाला आहे;
तिच्यानंतर $\lor$ येते; इत्यादी बाबी आपल्याला माहीत होतात.









\section{पूर्वतयारी}\label{pl:syn:pre:sec}

\begin{thm}[\emph{सूत्रांवरील विगमनाचे तत्त्व}]
  \label{pl:syn:pre:thm:induction}
  एखादा गुणधर्म~$P$ सर्व आण्विक सूत्रांसाठी लागू होत असेल आणि
  पुढील अटी पूर्ण करत असेल:
  \begin{enumerate}
    \item तो~$\varphi$ साठी लागू असताना $\lnot \varphi$ साठीही
      लागू होतो;
    \item तो~$\varphi$ आणि~$\psi$ यांच्यासाठी लागू असताना
      $(\varphi \land \psi)$ साठीही लागू होतो;
    \item तो~$\varphi$ आणि~$\psi$ यांच्यासाठी लागू असताना
      $(\varphi \lor \psi)$ साठीही लागू होतो;
    \item तो~$\varphi$ आणि~$\psi$ यांच्यासाठी लागू असताना
      $(\varphi \lif \psi)$ साठीही लागू होतो;
    \item तो~$\varphi$ आणि~$\psi$ यांच्यासाठी लागू असताना
      $(\varphi \liff \psi)$ साठीही लागू होतो;
  \end{enumerate}
  तर $P$ सर्व सूत्रांसाठी लागू होतो.
\end{thm}

\begin{proof}
  गुणधर्म~$P$ असलेल्या सर्व सूत्रांचा संग्रह~$S$ असू द्या.
  स्पष्टपणे $S \subseteq \Frm[L_0]$. \ref{pl:syn:fml:defn:formulas} मधील
  सर्व अटी $S$~पूर्ण करतो: त्यात सर्व आण्विक सूत्रे आहेत आणि
  तो संकारकांखाली बंद आहे.  $\Frm[L_0]$ हा असा सर्वांत लहान
  वर्ग आहे; म्हणून $\Frm[L_0] \subseteq S$. त्यामुळे $\Frm[L_0] = S$,
  आणि प्रत्येक सूत्राला गुणधर्म~$P$ आहे.
\end{proof}

\begin{prop}\label{pl:syn:pre:prop:balanced}
   $\Frm[L_0]$ मधील प्रत्येक सूत्र \emph{संतुलित} आहे; म्हणजे
   त्यात जेवढे डावे कंस आहेत तेवढेच उजवे कंस आहेत.
\end{prop}

\begin{prob}
\ref{pl:syn:pre:prop:balanced} सिद्ध करा.
\end{prob}

\begin{prop} \label{pl:syn:pre:prop:noinit}
कोणत्याही सूत्राचा कोणताही उचित आरंभीचा खंड सूत्र नसतो.
\end{prop}

\begin{prob}
  \ref{pl:syn:pre:prop:noinit} सिद्ध करा.
\end{prob}

\begin{prop}[एकमेव वाचनीयता]
$ \Frm[L_0]$ मधील कोणत्याही सूत्र~$\varphi$ चे पुढीलपैकी नेमक्या एका
रूपात रचनाविश्लेषण करता येते:
\begin{enumerate}
\item $\lfalse$.

\item $\ltrue$.

\item एखाद्या $\Obj p_n \in  \PVar$ साठी $\Obj p_n$.

\item एखाद्या सूत्र~$\psi$ साठी $\lnot \psi$.

\item काही सूत्रे $\psi$ आणि~$\chi$ यांच्यासाठी
  $(\psi \land \chi)$.

\item काही सूत्रे $\psi$ आणि~$\chi$ यांच्यासाठी
  $(\psi \lor \chi)$.

\item काही सूत्रे $\psi$ आणि~$\chi$ यांच्यासाठी
  $(\psi \lif \chi)$.

\item काही सूत्रे $\psi$ आणि~$\chi$ यांच्यासाठी
  $(\psi \liff \chi)$.
\end{enumerate}
शिवाय, हे रचनाविश्लेषण \emph{एकमेव} असते.
\end{prop}

\begin{proof}
$\varphi$ वरील विगमनाने सिद्ध करू. उदाहरणार्थ, $\varphi$ चे $(\psi \lif \chi)$ आणि
$(\psi' \lif \chi')$ अशी दोन भिन्न रचनाविश्लेषणे आहेत असे समजा. मग $\psi$
आणि $\psi'$ एकच असले पाहिजेत (अन्यथा त्यांपैकी एक हा दुसऱ्याचा उचित
आरंभीचा खंड असेल); म्हणून $\varphi$ ची ही दोन रचनाविश्लेषणे भिन्न असतील,
तर त्याचे कारण असेच असले पाहिजे की $\chi$ आणि $\chi'$ ही एकाच
चिन्हमालेची भिन्न रचनाविश्लेषणे आहेत; पण विगमन गृहीतकानुसार ते अशक्य
आहे.
\end{proof}


\begin{defn}[एकरूप आदेशन]
$\varphi$ आणि $\psi$ ही सूत्रे असतील आणि $\Obj p_i$ हे विधानीय चल असेल, तर $\varphi$ मधील $\Obj p_i$ ची प्रत्येक
आस्थिति $\psi$ च्या एका आस्थितीने बदलल्यावर मिळणारा परिणाम
$\Subst{\varphi}{\psi}{\Obj p_i}$ ने दर्शवला जातो; त्याचप्रमाणे,
$\Obj p_1$, \dots,~$\Obj p_n$ यांची सूत्रे $\psi_1$,
\dots,~$\psi_n$ यांनी एकाच वेळी केलेले आदेशन
$\SSubst{\varphi}{\subst{\psi_1}{\Obj p_1},\dots,\subst{\psi_n}{\Obj p_n}}$
ने दर्शवले जाते.
\end{defn}

\begin{prob} खालील पाच सूत्रांपैकी प्रत्येकासाठी, ते
 सूत्र \( \Subst{\varphi}{\psi}{\Obj p_i} \) या आदेशनाच्या रूपात व्यक्त
 करता येते का ते ठरवा; येथे \( \varphi \) हे अनुक्रमे (i) \( \Obj p_0 \);
 (ii) \( ( \lnot \Obj p_0 \land \Obj p_1) \); आणि (iii)
 \( ( ( \lnot \Obj p_0 \lif \Obj p_1 ) \land \Obj p_2 ) \) आहे.
 प्रत्येक बाबतीत संबद्ध आदेशन स्पष्ट करा.
  \begin{enumerate}
    \item \( \Obj p_1 \)
    \item \( ( \lnot \Obj p_0 \land \Obj p_0 ) \)
    \item \( ( ( \Obj p_0 \lor \Obj p_1 ) \land \Obj p_2 )  \)
    \item \( \lnot ( ( \Obj p_0 \lif \Obj p_1 ) \land \Obj p_2 ) \)
    \item \( (( \lnot ( \Obj p_0 \lif \Obj p_1 ) \lif ( \Obj p_0 \lor \Obj p_1 )) \land \lnot ( \Obj p_0 \land \Obj p_1 )) \)
  \end{enumerate}
\end{prob}

\begin{prob}
$\Subst{\varphi}{\psi}{p}$ ची विगमनाधारित, गणितीयदृष्ट्या काटेकोर व्याख्या
द्या.
\end{prob}









\section{रचनाक्रमिका}\label{pl:syn:fseq:sec}

सूत्रांची विगमनाधारित व्याख्या देणे आणि त्याला पूरक तंत्र म्हणून
सूत्रांचे गुणधर्म विगमनाने सिद्ध करणे हा सुबक व कार्यक्षम
दृष्टिकोन आहे. तथापि, सूत्रांच्या रचनेचा तळापासून वर जाणारा,
टप्प्याटप्प्याचा दृष्टिकोनही उपयुक्त ठरू शकतो; येथे आपण तो
\emph{रचनाक्रमिका} या संकल्पनेच्या साहाय्याने मांडतो.

\begin{defn}[सूत्रांसाठी रचनाक्रमिका]
\label{pl:syn:fseq:defn:fseq-frm}
भाषा~$\Lang L_0$ मधील चिन्हांपासून बनलेल्या चिन्हमालांची सांत क्रमिका
$\tuple{\varphi_0,\dotsc,\varphi_n}$ ही $\varphi$ साठी \emph{रचनाक्रमिका} असते, जर
$\varphi \ident \varphi_n$ आणि प्रत्येक $i \leq n$ साठी एकतर $\varphi_i$ हे आण्विक
सूत्र असेल किंवा पुढीलपैकी एखादी अट पूर्ण करणारे $j,k < i$ अस्तित्वात
असतील:
\begin{enumerate}
    \item $\varphi_i \ident \lnot \varphi_j$.%
    \item $\varphi_i \ident (\varphi_j \land \varphi_k)$.%
    \item $\varphi_i \ident (\varphi_j \lor \varphi_k)$.%
    \item $\varphi_i \ident (\varphi_j \lif \varphi_k)$.%
    \item $\varphi_i \ident (\varphi_j \liff \varphi_k)$.%
\end{enumerate}
\end{defn}

\begin{ex}
\[\tuple{
    \Obj p_0,
    \Obj p_1,
    (\Obj p_1 \land \Obj p_0),
    \lnot (\Obj p_1 \land \Obj p_0)
}
\]
ही $\lnot (\Obj p_1 \land \Obj p_0)$ ची एक रचनाक्रमिका आहे; तसेच
पुढील क्रमिकाही तिची रचनाक्रमिका आहे:
\[\tuple{
    \Obj p_0,
    \Obj p_1,
    \Obj p_0,
    (\Obj p_1 \land \Obj p_0),
    (\Obj p_0 \lif \Obj p_1),
    \lnot (\Obj p_1 \land \Obj p_0)
}.
\]
दुसऱ्या उदाहरणावरून दिसते की रचनाक्रमिकेत `अडगळ' असू शकते: अनावश्यक
असलेली किंवा रचनेत काहीही हातभार न लावणारी सूत्रे.
\end{ex}

\begin{prop}
\label{pl:syn:fseq:prop:formed}
$\Frm[L_0]$ मधील प्रत्येक सूत्र~$\varphi$ ला एक रचनाक्रमिका असते.
\end{prop}

\begin{proof}
$\varphi$ आण्विक आहे असे समजा. मग क्रमिका~$\tuple{\varphi}$ ही $\varphi$ साठी
रचनाक्रमिका आहे.
आता $\psi$ आणि~$\chi$ यांना अनुक्रमे $\tuple{\psi_0,\dotsc,\psi_n}$ आणि
$\tuple{\chi_0,\dotsc,\chi_m}$ या रचनाक्रमिका आहेत असे समजा.
\begin{enumerate}
  \item जर $\varphi \ident \lnot \psi$, तर
      $\tuple{\psi_0,\dotsc,\psi_n,\lnot \psi_n}$ ही $\varphi$ साठी
      रचनाक्रमिका आहे.
  \item जर $\varphi \ident (\psi \land \chi)$, तर
      $\tuple{\psi_0,\dotsc,\psi_n,\chi_0,\dotsc,\chi_m,(\psi_n \land \chi_m)}$
      ही $\varphi$ साठी रचनाक्रमिका आहे.
  \item जर $\varphi \ident (\psi \lor \chi)$, तर
      $\tuple{\psi_0,\dotsc,\psi_n,\chi_0,\dotsc,\chi_m,(\psi_n \lor \chi_m)}$
      ही $\varphi$ साठी रचनाक्रमिका आहे.
  \item जर $\varphi \ident (\psi \lif \chi)$, तर
      $\tuple{\psi_0,\dotsc,\psi_n,\chi_0,\dotsc,\chi_m,(\psi_n \lif \chi_m)}$
      ही $\varphi$ साठी रचनाक्रमिका आहे.
  \item जर $\varphi \ident (\psi \liff \chi)$, तर
      $\tuple{\psi_0,\dotsc,\psi_n,\chi_0,\dotsc,\chi_m,(\psi_n \liff \chi_m)}$
      ही $\varphi$ साठी रचनाक्रमिका आहे.
\end{enumerate}
सूत्रांवरील विगमनाच्या तत्त्वानुसार प्रत्येक सूत्राला एक
रचनाक्रमिका असते.
\end{proof}

याचे उलट विधानही आपण सिद्ध करू शकतो. ते महत्त्वाचे आहे, कारण त्यामुळे
सूत्रे परिभाषित करण्याच्या आपल्या दोन्ही पद्धती सममूल्य आहेत, म्हणजे
त्या समान परिणाम देतात, हे दिसते. तसेच रचनाक्रमिकांच्या लांबीवर साधे
विगमन वापरून सूत्रांविषयीची प्रमेये सिद्ध करता येतात.

\begin{lem}
\label{pl:syn:fseq:lem:fseq-init}
$\tuple{\varphi_0,\dotsc,\varphi_n}$ ही $\varphi_n$ साठी रचनाक्रमिका आहे आणि
$k \leq n$ असे समजा. मग $\tuple{\varphi_0,\dotsc,\varphi_k}$ ही $\varphi_k$ साठी
रचनाक्रमिका आहे.
\end{lem}

\begin{proof}
अभ्यासप्रश्न.
\end{proof}

\begin{thm}
\label{pl:syn:fseq:thm:fseq-frm-equiv}
$\Frm[L_0]$ हा भाषा~$\Lang L_0$ मधील ज्या सर्व चिन्हमालांना
रचनाक्रमिका आहे, त्यांचा संच आहे.
\end{thm}

\begin{proof}
भाषा~$\Lang L_0$ मधील ज्या सर्व चिन्हमालांना रचनाक्रमिका आहे, त्यांचा
संच~$F$ असू द्या. $\Frm[L_0] \subseteq F$ हे
\ref{pl:syn:fseq:prop:formed} मध्ये आपण पाहिले आहे; म्हणून आता उलट
समावेश सिद्ध करू.

$\varphi$ ला $\tuple{\varphi_0,\dotsc,\varphi_n}$ ही रचनाक्रमिका आहे असे समजा.
$n$ वरील प्रबल विगमनाने $\varphi \in \Frm[L_0]$ हे सिद्ध करू. लांबी
$m < n$ असलेली रचनाक्रमिका ज्या प्रत्येक चिन्हमालेला आहे, ती
$\Frm[L_0]$ मध्ये आहे, हे आपले विगमन गृहीतक आहे. रचनाक्रमिकेच्या
व्याख्येनुसार एकतर $\varphi_n$ आण्विक आहे किंवा पुढीलपैकी एखादी अट पूर्ण
करणारे $j,k < n$ अस्तित्वात असले पाहिजेत:
\begin{enumerate}
    \item $\varphi_n \ident \lnot \varphi_j$.%
    \item $\varphi_n \ident (\varphi_j \land \varphi_k)$.%
    \item $\varphi_n \ident (\varphi_j \lor \varphi_k)$.%
    \item $\varphi_n \ident (\varphi_j \lif \varphi_k)$.%
    \item $\varphi_n \ident (\varphi_j \liff \varphi_k)$.%
\end{enumerate}
आता प्रत्येक शक्यतेचा स्वतंत्र विचार करू. $\varphi_n$ आण्विक असेल, तर
$\varphi_n \in \Frm[L_0]$. त्याऐवजी $\varphi \ident
(\varphi_j \land \varphi_k)$ असे समजा.
\begin{quote}\small\textbf{स्रोतदुरुस्ती.} MRPL-002: गोठवलेल्या इंग्रजीत या एकाच शक्यतेत विन्यासात्मक अनन्यता-चिन्हाऐवजी पर्यायी सममूल्यता-चिन्ह आले आहे. आधीची स्पष्ट व्याख्या आणि इतर सर्व रचना-बाबी यांच्याशी जुळण्यासाठी येथे विन्यासात्मक अनन्यता वापरली आहे.\end{quote}
\ref{pl:syn:fseq:lem:fseq-init} नुसार
$\tuple{\varphi_0,\dotsc,\varphi_j}$ आणि $\tuple{\varphi_0,\dotsc,\varphi_k}$ या अनुक्रमे
$\varphi_j$ आणि~$\varphi_k$ साठी रचनाक्रमिका आहेत. या $\varphi$ च्या
रचनाक्रमिकेच्या उचित आरंभीच्या उपक्रमिका असल्यामुळे दोन्हींची लांबी
$n$ पेक्षा कमी आहे. म्हणून विगमन गृहीतकानुसार $\varphi_j$ आणि~$\varphi_k$ ही
$\Frm[L_0]$ मध्ये आहेत; त्यामुळे सूत्राच्या व्याख्येनुसार
$(\varphi_j \land \varphi_k)$ हेही त्यात आहे. इतर शक्यता समांतर युक्तिवादाने
सिद्ध होतात.
\end{proof}









\section{सत्य-मूल्यांकन आणि पूर्ति}\label{pl:syn:val:sec}

\begin{defn}[सत्य-मूल्यांकने]
``सत्य'' आणि ``असत्य'' ही दोन सत्यतामूल्ये असलेला संच
$\{\True, \False\}$ असू द्या. $\Lang{L_0}$ साठीचे
\emph{सत्य-मूल्यांकन} म्हणजे भाषेतील विधानीय चलांना
$\True$ किंवा $\False$ यांपैकी एक मूल्य देणारे फलन~$\pAssign{v}$;
म्हणजेच $\pAssign{v} \colon \PVar \to \{\True, \False \}$.
\end{defn}

\begin{defn}
  दिलेले सत्य-मूल्यांकन~$\pAssign{v}$ असताना, मूल्यन फलन
  $\pValue{v} \colon \Frm[L_0] \to \{\True, \False \}$ पुढीलप्रमाणे
  विगमनाने परिभाषित करू:
  \begin{align*}
    \pValue{v}(\lfalse) & = \False; \\
    \pValue{v}(\ltrue) & = \True; \\
    \pValue{v}(\Obj p_n) & = \pAssign{v}(\Obj p_n); \\
    \pValue{v}(\lnot \varphi) & = \begin{cases}
        \True & \text{जर } \pValue{v}(\varphi) = \False;\\
        \False & \text{अन्यथा.}
      \end{cases} \\
    \pValue{v}(\varphi \land \psi) & = \begin{cases}
        \True &
        \text{जर $\pValue{v}(\varphi) = \True$ आणि $\pValue{v}(\psi) = \True$;}\\
        \False &
        \text{जर $\pValue{v}(\varphi) = \False$ किंवा $\pValue{v}(\psi) = \False$}.
      \end{cases}\\
    \pValue{v}(\varphi \lor \psi) & = \begin{cases}
        \True &
        \text{जर $\pValue{v}(\varphi) = \True$ किंवा $\pValue{v}(\psi) = \True$;}\\
        \False &
        \text{जर $\pValue{v}(\varphi) = \False$ आणि $\pValue{v}(\psi) = \False$}.
      \end{cases}\\
    \pValue{v}(\varphi \lif \psi) & = \begin{cases}
        \True &
        \text{जर $\pValue{v}(\varphi) = \False$ किंवा $\pValue{v}(\psi) = \True$;}\\
        \False &
        \text{जर $\pValue{v}(\varphi) = \True$ आणि $\pValue{v}(\psi) = \False$}.
      \end{cases}\\
    \pValue{v}(\varphi \liff \psi) & = \begin{cases}
        \True &
        \text{जर $\pValue{v}(\varphi) = \pValue{v}(\psi)$;}\\
        \False &
        \text{जर $\pValue{v}(\varphi) \neq \pValue{v}(\psi)$}.
      \end{cases}
\end{align*}
\end{defn}

\begin{explain}
या बाबींना पुढील सत्यताकोष्टके अनुरूप आहेत:
\begin{center}

  \begin{tabular}{|c||c|} \hline
    $\varphi$ & $ \lnot \varphi$ \\
    \hline \hline
    $\True$ & $\False$ \\
    $\False$ & $\True$ \\
    \hline
  \end{tabular}


  \begin{tabular}{|cc||c|} \hline
    $\varphi$ & $\psi$ & $\varphi \land \psi$ \\
    \hline \hline
    $\True$ & $\True$ & $\True$ \\
    $\True$ & $\False$ & $\False$ \\
    $\False$ & $\True$ & $\False$ \\
    $\False$ & $\False$ & $\False$ \\
    \hline
  \end{tabular}


  \begin{tabular}{|cc||c|} \hline
    $\varphi$ & $\psi$ & $\varphi \lor \psi$ \\
    \hline \hline
    $\True$ & $\True$ & $\True$ \\
    $\True$ & $\False$ & $\True$ \\
    $\False$ & $\True$ & $\True$ \\
    $\False$ & $\False$ & $\False$ \\
    \hline
  \end{tabular}
%
\\[1em]
  \begin{tabular}{|cc||c|} \hline
    $\varphi$ & $\psi$ & $\varphi \lif \psi$ \\
    \hline \hline
    $\True$ & $\True$ & $\True$ \\
    $\True$ & $\False$ & $\False$ \\
    $\False$ & $\True$ & $\True$ \\
    $\False$ & $\False$ & $\True$ \\
    \hline
  \end{tabular}
  \begin{tabular}{|cc||c|} \hline
    $\varphi$ & $\psi$ & $\varphi \liff \psi$ \\
    \hline \hline
    $\True$ & $\True$ & $\True$ \\
    $\True$ & $\False$ & $\False$ \\
    $\False$ & $\True$ & $\False$ \\
    $\False$ & $\False$ & $\True$ \\
    \hline
  \end{tabular}
\end{center}
\end{explain}
\begin{prob}
$\Lang{L_0}$ मध्ये $\diamondsuit$ हा त्रिस्थानी संयोजक वाढवण्याचा विचार
करा; त्याचे मूल्यन पुढीलप्रमाणे दिले आहे:
\begin{gather*}
  \pValue{v}(\diamondsuit( \varphi , \psi , \chi ) ) = \begin{cases}
    \pValue{v}( \psi ) &
    \text{जर $\pValue{v}(\varphi) = \True$;}\\
    \pValue{v}( \chi ) &
    \text{जर $\pValue{v}(\varphi) = \False $}.
  \end{cases}
\end{gather*}
या संयोजकाचे सत्यताकोष्टक लिहा.
\end{prob}
\begin{thm}[स्थानिक निर्धारण]
\label{pl:syn:val:thm:LocalDetermination} $\pAssign{v_1}$ आणि
$\pAssign{v_2}$ ही अशी सत्य-मूल्यांकने आहेत की $\varphi$ मध्ये आस्थित असलेल्या
विधानीय चलांना ती समान मूल्ये देतात; म्हणजे, एखादे $\Obj p_n$ हे
सूत्र~$\varphi$ मध्ये आस्थित असेल तेव्हा $\pAssign{v_1}(\Obj p_n) =
\pAssign{v_2}(\Obj p_n)$. मग $\pValue{v_1}$ आणि $\pValue{v_2}$ हीही
$\varphi$ वर एकमत असतात; म्हणजे $\pValue{v_1}(\varphi) = \pValue{v_2}(\varphi)$.
\end{thm}
\begin{proof}
$\varphi$ वरील विगमनाने.
\end{proof}
\begin{defn}[पूर्ति]
\label{pl:syn:val:defn:satisfaction} आपण
  \emph{सत्य-मूल्यांकन~$\pAssign{v}$ कडून सूत्र~$\varphi$ ची पूर्ति},
  $\pSat{v}{\varphi}$, ही संकल्पना पुढीलप्रमाणे विगमनाने परिभाषित करू शकतो.
  (``$\pSat{v}{\varphi}$ नाही'' हे दर्शवण्यासाठी आपण $\pSat/{v}{\varphi}$ लिहितो.)
\begin{enumerate}
\item %
  \indcase{\varphi}{\lfalse}{$\pSat/{v}{\indfrm}$.}
\item %
  \indcase{\varphi}{\ltrue}{$\pSat{v}{\indfrm}$.}
\item \indcase{\varphi}{\Obj p_i}{$\pSat{v}{\indfrm}$ तेव्हा आणि केवळ
  तेव्हाच, जेव्हा $\pAssign{v}(\Obj p_i) = \True$.}
\item %
  \indcase{\varphi}{\lnot \psi}{$\pSat{v}{\indfrm}$ तेव्हा आणि केवळ
    तेव्हाच, जेव्हा $\pSat/{v}{\psi}$.}
\item %
  \indcase{\varphi}{(\psi \land \chi)}{$\pSat{v}{\indfrm}$ तेव्हा आणि केवळ
    तेव्हाच, जेव्हा $\pSat{v}{\psi}$ आणि $\pSat{v}{\chi}$.}
\item %
  \indcase{\varphi}{(\psi \lor \chi)}{$\pSat{v}{\indfrm}$ तेव्हा आणि केवळ
    तेव्हाच, जेव्हा $\pSat{v}{\psi}$ किंवा $\pSat{v}{\chi}$ (किंवा दोन्ही).}
\item %
  \indcase{\varphi}{(\psi \lif \chi)}{$\pSat{v}{\indfrm}$ तेव्हा आणि केवळ
    तेव्हाच, जेव्हा $\pSat/{v}{\psi}$ किंवा $\pSat{v}{\chi}$ (किंवा दोन्ही).}
\item %
  \indcase{\varphi}{(\psi \liff \chi)}{$\pSat{v}{\indfrm}$ तेव्हा आणि केवळ
    तेव्हाच, जेव्हा $\pSat{v}{\psi}$ आणि $\pSat{v}{\chi}$ ही दोन्ही, किंवा
    $\pSat{v}{\psi}$ आणि $\pSat{v}{\chi}$ यांपैकी कोणतेही नाही.}
\end{enumerate}
$\Gamma$ हा सूत्रांचा संच असेल, तर $\pSat{v}{\Gamma}$ तेव्हा आणि
केवळ तेव्हाच, जेव्हा प्रत्येक~$\varphi \in \Gamma$ साठी $\pSat{v}{\varphi}$.
\end{defn}
\begin{prop}\label{pl:syn:val:prop:sat-value}
  $\pSat{v}{\varphi}$ तेव्हा आणि केवळ तेव्हाच, जेव्हा
  $\pValue{v}(\varphi) = \True$.
\end{prop}
\begin{proof}
  $\varphi$ वरील विगमनाने.
\end{proof}
\begin{prob}
  \ref{pl:syn:val:prop:sat-value} सिद्ध करा.
\end{prob}
\section{चिन्हार्थविषयक संकल्पना}\label{pl:syn:sem:sec}
पुढील चिन्हार्थविषयक संकल्पना परिभाषित करू:
\begin{defn}
\begin{enumerate}
\item सूत्र~$\varphi$ हे \emph{पूर्ततायोग्य} असते, जर एखाद्या
  $\pAssign{v}$ साठी $\pSat{v}{\varphi}$; आणि कोणत्याही $\pAssign{v}$ साठी
  $\pSat{v}{\varphi}$ नसल्यास ते \emph{अपूर्ततायोग्य} असते;
\item सर्व सत्य-मूल्यांकने~$\pAssign{v}$ साठी $\pSat{v}{\varphi}$ असेल, तर
  सूत्र~$\varphi$ हे \emph{सर्वतः सत्य सूत्र} असते;
\item सूत्र~$\varphi$ हे पूर्ततायोग्य असेल, पण सर्वतः सत्य सूत्र नसेल,
  तर ते \emph{परिस्थितीवश} असते;
\item $\Gamma$ हा सूत्रांचा संच असेल, तर $\Gamma \Entails \varphi$
  (``$\Gamma$ पासून $\varphi$ तार्किकरीत्या निष्पन्न होते'') तेव्हा आणि केवळ
  तेव्हाच, जेव्हा $\pSat{v}{\Gamma}$ असलेल्या प्रत्येक
  सत्य-मूल्यांकन~$\pAssign{v}$ साठी $\pSat{v}{\varphi}$.
\item $\Gamma$ हा सूत्रांचा संच असेल, तर एखाद्या
  सत्य-मूल्यांकन~$\pAssign{v}$ साठी $\pSat{v}{\Gamma}$ असल्यास $\Gamma$
  हा \emph{पूर्ततायोग्य} असतो; अन्यथा $\Gamma$ हा
  \emph{अपूर्ततायोग्य} असतो.
\end{enumerate}
\end{defn}
\begin{prob}
 पुढील चार सूत्रांपैकी प्रत्येक सूत्र (a)~पूर्ततायोग्य,
 (b)~सर्वतः सत्य सूत्र आणि (c)~परिस्थितीवश आहे की नाही ते ठरवा.
\begin{enumerate}
  \item \( ( \Obj p_0 \lif ( \lnot \Obj p_1 \lif \lnot \Obj p_0 ) ) \).
  \item \( ( ( \Obj p_0 \land \lnot \Obj p_1 ) \lif ( \lnot \Obj p_0 \land \Obj p_2 )) \liff ( ( \Obj p_2 \lif \Obj p_0 ) \lif ( \Obj p_0 \lif \Obj p_1 )) \).
  \item \( ( \Obj p_0 \liff \Obj p_1 ) \lif ( \Obj p_2 \liff \lnot \Obj p_1 ) \).
  \item \( (( \Obj p_0 \liff ( \lnot \Obj p_1 \land \Obj p_2 )) \lor ( \Obj p_2 \lif ( \Obj p_0 \liff \Obj p_1 ))) \).
\end{enumerate}
\end{prob}
\begin{prop}
\label{pl:syn:sem:prop:semanticalfacts}
\begin{enumerate}
\item $\varphi$ हे सर्वतः सत्य सूत्र तेव्हा आणि केवळ तेव्हाच असते, जेव्हा
  $\emptyset \Entails \varphi$;
\item $\Gamma \Entails \varphi$ आणि $\Gamma \Entails \varphi \lif \psi$ असल्यास
  $\Gamma \Entails \psi$;
\item $\Gamma$ पूर्ततायोग्य असेल, तर $\Gamma$ चा प्रत्येक सांत उपसंचही
  पूर्ततायोग्य असतो;
\item \label{pl:syn:sem:def:monotonicity} एकस्वनिकता: $\Gamma \subseteq \Delta$
  आणि $\Gamma \Entails \varphi$ असल्यास $\Delta \Entails \varphi$ हेही असते;
\item \label{pl:syn:sem:def:Cut} संक्रमकता: $\Gamma \Entails \varphi$ आणि
  $\Delta \cup \{ \varphi\} \Entails \psi$ असल्यास $\Gamma \cup \Delta \Entails
  \psi$.
\end{enumerate}
\end{prop}
\begin{proof}
सरावासाठी.
\end{proof}
\begin{prob}
\ref{pl:syn:sem:prop:semanticalfacts} सिद्ध करा.
\end{prob}
\begin{prop}\label{pl:syn:sem:prop:entails-unsat}
  $\Gamma \Entails \varphi$ तेव्हा आणि केवळ तेव्हाच, जेव्हा
  $\Gamma \cup \{\lnot \varphi\}$ अपूर्ततायोग्य असतो.
\end{prop}
\begin{proof}
सरावासाठी.
\end{proof}
\begin{prob}
\ref{pl:syn:sem:prop:entails-unsat} सिद्ध करा.
\end{prob}
\begin{thm}[चिन्हार्थविषयक निगमन प्रमेय]
  \label{pl:syn:sem:thm:sem-deduction} $\Gamma \Entails \varphi \lif \psi$ तेव्हा आणि केवळ
  तेव्हाच, जेव्हा $\Gamma \cup \{\varphi\} \Entails \psi$.
\end{thm}
\begin{proof}
सरावासाठी.
\end{proof}
\begin{prob}
\ref{pl:syn:sem:thm:sem-deduction} सिद्ध करा.
\end{prob}
\chapter{सिद्धता-पद्धती}\label{fol:prf::chap}
\begin{editorial}
या प्रकरणात निष्पत्ती-पद्धतींविषयीची सर्वसाधारण सामग्री
एकत्रित केली आहे. एखादी विशिष्ट पद्धत वापरणारे पाठ्यपुस्तक,
ज्या प्रकरणात त्या पद्धतीचा परिचय करून दिला आहे त्या प्रकरणाच्या
सुरुवातीस प्रस्तावना-विभाग आणि संबंधित आढावा-विभाग समाविष्ट करू शकते.
\end{editorial}
\section{परिचय}\label{fol:prf:int:sec}
तर्कप्रणालींमध्ये सामान्यतः चिन्हार्थमीमांसा आणि निष्पत्ती-पद्धती हे दोन्ही घटक असतात. चिन्हार्थमीमांसेत सत्यता, पूर्ततायोग्यता,
वैधता आणि तार्किक निष्पन्नता अशा संकल्पनांचा विचार केला जातो.
तार्किक निष्पन्नता आणि वैधता प्रस्थापित करण्यासाठी निव्वळ विन्यासात्मक
रीत उपलब्ध करून देणे हा निष्पत्ती-पद्धतींचा उद्देश असतो. त्या
निव्वळ विन्यासात्मक असतात याचा अर्थ असा की अशा पद्धतीतील
निष्पत्ती ही एक सांत विन्यासात्मक वस्तू असते; सहसा ती
वाक्ये किंवा सूत्रे यांचा अनुक्रम (किंवा इतर एखादी सांत
मांडणी) असते. वाक्ये किंवा सूत्रे यांचा कोणताही दिलेला
अनुक्रम किंवा मांडणी ``योग्य'' आहे, हे यांत्रिक रीतीने पडताळता येते,
हा चांगल्या निष्पत्ती-पद्धतींचा गुणधर्म असतो.
प्रथम-क्रम तर्कशास्त्रासाठीच्या सर्वांत सोप्या (आणि इतिहासात प्रथम
आलेल्या) निष्पत्ती-पद्धती \emph{स्वयंसिद्धकीय} होत्या. अशा
पद्धतीत सूत्रांचा एखादा अनुक्रम तेव्हाच निष्पत्ती मानला
जातो, जेव्हा त्यातील प्रत्येक स्वतंत्र सूत्र एकतर
``स्वयंसिद्धकांच्या'' एका निश्चित संचात असते किंवा त्या अनुक्रमात
त्याच्या आधी आलेल्या सूत्रांपासून ठरावीक संख्येतील ``निगमन
नियमां''पैकी एखाद्या नियमाने निष्पन्न होते. शिवाय, सूत्र
स्वयंसिद्धक आहे का आणि ते इतर सूत्रांपासून एखाद्या निगमन
नियमानुसार योग्य रीतीने निष्पन्न होते का, हे यांत्रिक रीतीने पडताळता
येते. स्वयंसिद्धकीय निष्पत्ती-पद्धतींचे वर्णन करणे सोपे असते आणि
अधिउपपत्तीच्या पातळीवरही त्या हाताळणे सोपे असते; पण त्यांतील
निष्पत्त्या वाचणे आणि समजून घेणे कठीण असते, तसेच त्या तयार करणेही
कठीण असते.
निष्पत्त्या तयार करणे सोपे व्हावे किंवा पूर्ण झालेल्या
निष्पत्त्या समजून घेणे सोपे व्हावे, या उद्देशाने इतर
निष्पत्ती-पद्धती विकसित करण्यात आल्या आहेत. नैसर्गिक निगमन,
सत्यता-वृक्ष---ज्यांना टॅब्लो सिद्धता असेही म्हणतात---आणि क्रमवर्ती कलन
ही त्यांची उदाहरणे आहेत. काही निष्पत्ती-पद्धती विशेषतः
यांत्रिकीकरण डोळ्यांसमोर ठेवून रचल्या जातात. उदाहरणार्थ, निराकरण पद्धत
संगणकीय प्रणालीत कार्यान्वित करणे सोपे असते (पण तिच्यातील
निष्पत्त्या समजून घेणे जवळजवळ अशक्य असते). यांपैकी बहुतेक इतर
निष्पत्ती-पद्धती निष्पत्त्यांना अनुक्रमांऐवजी सूत्रांचे
वृक्ष म्हणून दर्शवतात. त्यामुळे निष्पत्तीचा कोणता भाग इतर
कोणत्या भागांवर अवलंबून आहे, हे पाहणे सोपे होते.
म्हणून प्रथम-क्रम तर्कशास्त्रासारखी एखादी तर्कप्रणाली घेतल्यास,
वाक्य हे \emph{प्रमेय} असणे म्हणजे काय आणि वाक्य हे
इतर काही वाक्यांपासून निष्पन्न करता येण्याजोगे असणे म्हणजे काय, यांची
वेगवेगळ्या निष्पत्ती-पद्धती वेगवेगळी नेमकी स्पष्टीकरणे देतील. हे
कसेही केले (स्वयंसिद्धकीय निष्पत्त्या, नैसर्गिक निगमने, क्रमवर्ती
निष्पत्त्या, सत्यता-वृक्ष किंवा निराकरण-खंडने यांच्या साहाय्याने),
तरी हे संबंध वैधता आणि तार्किक निष्पन्नता या चिन्हार्थविषयक संकल्पनांशी
जुळावेत, अशी आपली अपेक्षा असते. ``$\varphi$~हे प्रमेय आहे'' यासाठी
$\Proves \varphi$ आणि ``$\varphi$~हे~$\Gamma$ पासून निष्पन्न करता येण्याजोगे आहे'' यासाठी
``$\Gamma \Proves \varphi$'' लिहू या. $\Proves$~ची व्याख्या कोणतीही असो,
ती $\Entails$~शी जुळावी, म्हणजे:
\begin{enumerate}
\item $\Proves \varphi$ जर आणि फक्त जर $\Entails \varphi$
\item $\Gamma \Proves \varphi$ जर आणि फक्त जर $\Gamma \Entails \varphi$
\end{enumerate}
वरील विधानांतील ``फक्त जर'' दिशेला \emph{निर्दोषता} म्हणतात.
निष्पत्ती-पद्धत तेव्हा निर्दोष असते, जेव्हा निष्पन्नतेमुळे
तार्किक निष्पन्नतेची (किंवा वैधतेची) हमी मिळते. प्रत्येक समाधानकारक
निष्पत्ती-पद्धत निर्दोष असलीच पाहिजे; सदोष निष्पत्ती-पद्धती
मुळीच उपयोगी नसतात. कारण निष्पत्तीचा संपूर्ण उद्देशच वैधतेची
किंवा तार्किक निष्पन्नतेची विन्यासात्मक हमी देणे हा असतो. आपण मांडत
असलेल्या निष्पत्ती-पद्धतींची निर्दोषता पुढे सिद्ध करू.
उलट ``जर'' दिशेलाही महत्त्व आहे: तिला \emph{संपूर्णता} म्हणतात.
एखादी संपूर्ण निष्पत्ती-पद्धत इतकी समर्थ असते की $\varphi$~वैध असेल
तेव्हा $\varphi$~हे प्रमेय आहे, तसेच $\Gamma \Entails \varphi$ असेल तेव्हा
$\Gamma \Proves \varphi$, हे ती दाखवू शकते. संपूर्णता प्रस्थापित करणे अधिक
कठीण असते आणि काही तर्कप्रणालींसाठी संपूर्ण निष्पत्ती-पद्धतीच
नसतात. प्रथम-क्रम तर्कशास्त्रासाठी मात्र त्या आहेत. कुर्ट गोडेल यांनी
आपल्या १९२९ सालच्या प्रबंधात प्रथम-क्रम तर्कशास्त्राच्या
निष्पत्ती-पद्धतीची संपूर्णता सर्वप्रथम सिद्ध केली.
निष्पत्ती-पद्धतींशी जोडलेली आणखी एक संकल्पना म्हणजे
\emph{सुसंगतता}. वाक्यांच्या एखाद्या संचापासून कोणतीही गोष्ट
निष्पन्न करता येत असेल तर तो संच विसंगत म्हणतात; अन्यथा तो सुसंगत
असतो. विसंगतता ही अपूर्ततायोग्यतेची विन्यासात्मक समकक्ष संकल्पना आहे:
जसे अपूर्ततायोग्य संच चांगल्या उपपत्ती ठरत नाहीत, तसेच वाक्यांचे
विसंगत संचही चांगल्या उपपत्ती ठरत नाहीत; त्यांच्यात मूलभूत स्वरूपाचा
दोष असतो. वाक्यांचे सुसंगत संच सत्य किंवा उपयोगी असतीलच असे
नाही, पण ते निदान तार्किक उपयोगितेचा किमान उंबरठा ओलांडतात. वेगवेगळ्या
निष्पत्ती-पद्धतींसाठी वाक्यांच्या संचांच्या सुसंगततेची नेमकी
व्याख्या वेगळी असू शकते; पण~$\Proves$ प्रमाणेच सुसंगतताही तिच्या
चिन्हार्थविषयक समकक्ष संकल्पनेशी, म्हणजे पूर्ततायोग्यतेशी, जुळली पाहिजे.
$\Gamma$ सुसंगत आहे जर आणि फक्त जर तो पूर्ततायोग्य आहे, हे नेहमी खरे
असावे अशी आपली अपेक्षा आहे. येथे ``फक्त जर'' दिशा संपूर्णतेइतकीच ठरते
(सुसंगततेमुळे पूर्ततायोग्यतेची हमी मिळते), आणि ``जर'' दिशा
निर्दोषतेइतकीच ठरते (पूर्ततायोग्यतेमुळे सुसंगततेची हमी मिळते).
प्रत्यक्षात, अभिजात प्रथम-क्रम तर्कशास्त्रासाठी निर्दोषता आणि संपूर्णता
यांची ही दोन रूपे सममूल्य आहेत.
\section{क्रमवर्ती कलन}\label{fol:prf:seq:sec}
अनेक निष्पत्ती-पद्धती वाक्यांच्या मांडण्यांवर कार्य करतात,
तर क्रमवर्ती कलन \emph{क्रमवर्तींवर} कार्य करते. क्रमवर्ती ही पुढील रूपाची
अभिव्यक्ती असते:
\[
\varphi_1, \dots, \varphi_m \Sequent \psi_1, \dots, \psi_m,
\]
म्हणजे वाक्यांच्या अशा दोन अनुक्रमांची जोडी, ज्यांना क्रमवर्ती
चिन्ह~$\Sequent$ वेगळे करते. यांपैकी कोणताही अनुक्रम रिकामा असू शकतो.
क्रमवर्ती कलनातील निष्पत्ती म्हणजे क्रमवर्तींचा वृक्ष होय. त्यातील
सर्वांत वरच्या क्रमवर्ती विशेष रूपाच्या असतात (त्यांना ``आरंभीच्या
क्रमवर्ती'' किंवा ``स्वयंसिद्धके'' म्हणतात), आणि इतर प्रत्येक क्रमवर्ती
तिच्या लगत वर असलेल्या क्रमवर्तींपासून निगमन नियमांपैकी एखाद्या नियमाने
निष्पन्न होते. निगमन नियम एकतर क्रमवर्तींतील वाक्ये हाताळतात
(डाव्या किंवा उजव्या बाजूला त्यांची भर घालतात, त्यांना काढतात किंवा त्यांचा क्रम
बदलतात), किंवा नियमाच्या निष्कर्षात एखादे जटिल सूत्र आणतात.
उदाहरणार्थ, $\LeftR{\land}$ नियमामुळे $\varphi, \Gamma \Sequent \Delta$ पासून
$\varphi \land \psi, \Gamma \Sequent \Delta$ हा निष्कर्ष काढता येतो; आणि
$\RightR{\lif}$ नियमामुळे $\varphi, \Gamma \Sequent \Delta, \psi$ पासून
$\Gamma \Sequent \Delta, \varphi \lif \psi$ हा निष्कर्ष काढता येतो. हे कोणत्याही
$\Gamma$, $\Delta$, $\varphi$ आणि~$\psi$ साठी लागू आहे. (विशेषतः, $\Gamma$
आणि~$\Delta$ रिकामे असू शकतात.)

क्रमवर्ती कलनावर आधारित $\Proves$ संबंधाची व्याख्या पुढीलप्रमाणे केली
जाते: $\Gamma \Proves \varphi$ तेव्हा आणि केवळ तेव्हाच, जेव्हा असा एखादा
अनुक्रम~$\Gamma_0$ अस्तित्वात असतो की $\Gamma_0$ मधील प्रत्येक $\varphi$
हा~$\Gamma$ मध्ये असतो आणि ज्याच्या मुळाशी $\Gamma_0 \Sequent \varphi$ ही
क्रमवर्ती आहे अशी निष्पत्ती अस्तित्वात असते. $\Sequent \varphi$ या क्रमवर्तीची
निष्पत्ती असेल तर क्रमवर्ती कलनात $\varphi$ हे प्रमेय असते. उदाहरणार्थ,
$\Proves (\varphi \land \psi) \lif \varphi$ हे दाखवणारी निष्पत्ती पुढे दिली आहे:
\begin{prooftree}
  \Axiom$\varphi \fCenter \varphi$
  \RightLabel{\LeftR{\land}}
  \UnaryInf$\varphi \land \psi \fCenter \varphi$
  \RightLabel{\RightR{\lif}}
  \UnaryInf$\fCenter (\varphi \land \psi) \lif \varphi$
\end{prooftree}

प्रत्येक $\varphi \in \Gamma_0$ हे~$\Gamma$ मध्ये असेल अशा एखाद्या
$\Gamma_0 \Sequent$ ची निष्पत्ती अस्तित्वात असेल (म्हणजे
क्रमवर्तीची उजवी बाजू रिकामी असेल), तर क्रमवर्ती कलनात संच~$\Gamma$
विसंगत असतो. \RightR{\Weakening} नियम वापरून विसंगत संचापासून कोणतेही
वाक्य निष्पन्न करता येते.

क्रमवर्ती कलनाची निर्मिती १९३० च्या दशकात गरहार्ड गेंटझेन यांनी केली.
त्याची रचना पद्धतशीर आणि सममित असल्यामुळे निष्पत्त्यांची उपपत्ती
विकसित करण्यासाठी ते अतिशय उपयुक्त आकारिक साधन आहे. क्रमवर्ती कलनात
निष्पत्त्या शोधणे तुलनेने सोपे असते; पण या निष्पत्त्या अनेकदा
वाचायला कठीण असतात आणि त्यांचा सिद्धतांशी असलेला संबंधही काही वेळा
सहज दिसत नाही. तरी निष्पत्ती-पद्धतींकडे पाहण्याची ही अतिशय सुबक
रीत ठरली आहे, आणि अनेक तर्कप्रणालींसाठी क्रमवर्ती-कलन पद्धती उपलब्ध आहेत.









\section{नैसर्गिक निगमन}\label{fol:prf:ntd:sec}

नैसर्गिक निगमन ही प्रत्यक्ष तर्कविचाराचे अनुकरण करण्यासाठी घडवलेली
निष्पत्ती-पद्धत आहे (विशेषतः गणितज्ञ वापरतात त्या नियमबद्ध
तर्कविचाराचे). प्रत्यक्ष तर्कविचार अनेक ``नैसर्गिक'' पद्धतींनी पुढे जातो.
उदाहरणार्थ, प्रकरणांनुसार सिद्धतेमुळे विकल्पयोगी आधारविधानाच्या आधारे
निष्कर्ष प्रस्थापित करता येतो: विकल्पयोगाच्या प्रत्येक घटकापासून
तो निष्कर्ष निष्पन्न होतो, हे स्वतंत्रपणे दाखवले जाते. अप्रत्यक्ष सिद्धतेत
निष्कर्षाचे नकरण व्याघाताकडे नेते हे दाखवून निष्कर्ष प्रस्थापित करता येतो.
सोपाधिक सिद्धतेत पूर्वांगापासून उत्तरांग निष्पन्न होते हे दाखवून
``जर \dots तर \dots'' असा सोपाधिक दावा प्रस्थापित केला जातो. नैसर्गिक
निगमन म्हणजे अशा नैसर्गिक अनुमानांपैकी काहींचे आकारिकीकरण होय. प्रत्येक
तार्किक संयोजक आणि संख्यापक यांना दोन नियम असतात—एक प्रवेशन नियम आणि एक
विलोपन नियम—आणि हे नियम प्रत्येकी अशाच एका नैसर्गिक अनुमान-पद्धतीशी
जुळतात. उदाहरणार्थ, $\Intro{\lif}$ सोपाधिक सिद्धतेशी, तर $\Elim{\lor}$
प्रकरणांनुसार सिद्धतेशी जुळतो. $\Elim{\land}$ हा विशेषतः सोपा नियम आहे;
त्यामुळे $\varphi \land \psi$ पासून~$\varphi$ (किंवा $\psi$) असा निष्कर्ष काढता येतो.

नैसर्गिक निगमनाला इतर निष्पत्ती-पद्धतींपासून वेगळे करणारे एक
वैशिष्ट्य म्हणजे त्यात गृहीतकांचा होणारा वापर. नैसर्गिक निगमनातील निष्पत्ती म्हणजे सूत्रांचा वृक्ष होय. या सूत्रांच्या
वृक्षाच्या मुळाशी एकच सूत्र असते, आणि वृक्षाची ``पाने'' ही अशी
सूत्रे असतात की ज्यांच्यापासून निष्कर्ष काढला जातो. नैसर्गिक
निगमनात पानांवरील काही सूत्रे ही निष्पत्तीमध्ये भूमिका बजावतात;
पण निष्पत्ती निष्कर्षापर्यंत पोहोचेपर्यंत ती ``वापरून संपवलेली'' असतात.
प्रत्यक्ष तर्कविचारात फक्त थोड्या काळापुरती लागू असणारी गृहीतके मांडण्याच्या
प्रथेशी हे जुळते. उदाहरणार्थ, प्रकरणांनुसार सिद्धतेत आपण विकल्पयोगाच्या
प्रत्येक घटकाची सत्यता गृहीत धरतो; सोपाधिक सिद्धतेत पूर्वांगाची सत्यता
गृहीत धरतो; आणि अप्रत्यक्ष सिद्धतेत निष्कर्षाच्या नकरणाची सत्यता गृहीत
धरतो. अशी तात्पुरती गृहीतके मांडणे आणि एखादी मधली पायरी प्रस्थापित करण्यासाठी
ती नंतर बाजूला काढणे, हे नैसर्गिक निगमनाचे खास वैशिष्ट्य आहे. नैसर्गिक
निगमनातील निष्पत्तीच्या पानांवरील सूत्रांना गृहीतके म्हणतात,
आणि काही निगमन नियम त्यांना ``मुक्त'' करू शकतात. उदाहरणार्थ, काही
गृहीतकांपासून~$\psi$ ची अशी निष्पत्ती असेल आणि त्या गृहीतकांत~$\varphi$
असेल, तर $\Intro{\lif}$ नियमामुळे~$\varphi \lif \psi$ असा निष्कर्ष काढता येतो
आणि~$\varphi$ या रूपातील कोणतेही गृहीतक मुक्त करता येते. (कोणती गृहीतके
कोणत्या अनुमानाच्या वेळी मुक्त केली जातात याची नोंद ठेवण्यासाठी, त्या
अनुमानाला आणि त्याने मुक्त केलेल्या गृहीतकांना
एक क्रमांक देतो.) निष्पत्तीच्या शेवटी जी गृहीतके मुक्त न केलेली
राहतात, ती निष्कर्षाच्या सत्यतेसाठी एकत्रितपणे पुरेशी असतात; म्हणून
निष्पत्ती ही प्रस्थापित करते की तिची मुक्त न केलेली गृहीतके तिचा निष्कर्ष तार्किकरीत्या निष्पन्न करतात.

नैसर्गिक निगमनावर आधारित $\Gamma \Proves \varphi$ हा संबंध तेव्हा आणि केवळ
तेव्हाच लागू होतो, जेव्हा अशी निष्पत्ती असते की त्यात वृक्षातील
शेवटचे वाक्य हे~$\varphi$ असते आणि मुक्त न केलेले प्रत्येक पान
हे~$\Gamma$ मध्ये असते. नैसर्गिक निगमनात $\varphi$ हे तेव्हा आणि केवळ तेव्हाच
प्रमेय असते, जेव्हा अशी निष्पत्ती असते की त्यातील शेवटचे
वाक्य हे~$\varphi$ असते आणि सर्व गृहीतके मुक्त झालेली असतात.
उदाहरणार्थ, $\Proves (\varphi \land \psi) \lif \varphi$ हे दाखवणारी निष्पत्ती
पुढे दिली आहे:
\begin{prooftree}
  \AxiomC{$\Discharge{\varphi \land \psi}{1}$}
  \RightLabel{\Elim{\land}}
  \UnaryInfC{$\varphi$}
  \DischargeRule{\Intro{\lif}}{1}
  \UnaryInfC{$(\varphi \land \psi) \lif \varphi$}
\end{prooftree}
$1$ ही खूण दाखवते की $\varphi \land \psi$ हे गृहीतक \Intro{\lif} अनुमानाच्या
वेळी मुक्त केले जाते.

नैसर्गिक निगमनात संच~$\Gamma$ तेव्हा आणि केवळ तेव्हाच विसंगत असतो,
जेव्हा $\Gamma \Proves \lfalse$. \FalseInt{} नियमामुळे विसंगत संचापासून
कोणतेही वाक्य निष्पन्न करता येते.

नैसर्गिक निगमन पद्धती १९३० च्या दशकात गरहार्ड गेंटझेन आणि स्टॅनिस्लॉ
यास्कोव्हस्की (Stanis\l{}aw Ja\'skowski) यांनी विकसित केल्या; पुढे डॅग
प्रावित्झ आणि फ्रेडरिक फिच यांनी त्यांचा आणखी विकास केला. त्यांतील अनुमाने
सिद्धतेच्या नैसर्गिक पद्धतींचे अनुकरण करत असल्यामुळे तत्त्वज्ञ नैसर्गिक
निगमनाला पसंती देतात.
फिच यांनी विकसित केलेल्या आवृत्त्या तर्कशास्त्राच्या प्रास्ताविक
पाठ्यपुस्तकांमध्ये अनेकदा वापरल्या जातात. तर्कशास्त्राच्या तत्त्वज्ञानात
नैसर्गिक निगमनाचे नियम तार्किक कारकांचे अर्थ देतात, असे काही वेळा मानले
गेले आहे (``सिद्धता-उपपत्तीय चिन्हार्थमीमांसा'').









\section{टॅब्लो}\label{fol:prf:tab:sec}

अनेक निष्पत्ती-पद्धती वाक्यांच्या मांडणीवर काम करतात, तर
टॅब्लो चिन्हांकित सूत्रांवर काम करतात. चिन्हांकित सूत्र म्हणजे
सत्यतामूल्याचे चिन्ह ($\True$ किंवा $\False$) आणि वाक्य यांची जोडी:
\[\sFmla{\True}{\varphi} \text{ किंवा } \sFmla{\False}{\varphi}.
\]
टॅब्लोमध्ये चिन्हांकित सूत्रे खालीच्या दिशेने शाखा फुटणाऱ्या
वृक्षात मांडलेली असतात. त्याची सुरुवात काही \emph{गृहीतकांपासून} होते आणि
त्यांच्या वर असलेल्या चिन्हांकित सूत्रांपैकी एखाद्याला निगमन नियमांपैकी एक
नियम लावल्यामुळे मिळणाऱ्या चिन्हांकित सूत्रांनी तो पुढे चालू राहतो.
प्रत्येक नियमामुळे एखाद्या शाखेच्या शेवटी एक किंवा अधिक चिन्हांकित सूत्रे
जोडता येतात, किंवा दोन चिन्हांकित सूत्रे शेजारी शेजारी जोडता येतात---दुसऱ्या
स्थितीत शाखा दोन शाखांत फुटते आणि नव्याने जोडलेली दोन चिन्हांकित सूत्रे त्या दोन शाखांची टोके बनतात.

संमिश्र चिन्हांकित सूत्राला नियम लावल्यावर तात्काळ उप-सूत्रे असलेली
चिन्हांकित सूत्रे जोडली जातात. हे नियम जोड्यांनी येतात—दोन्ही चिन्हांसाठी
प्रत्येकी एक नियम. उदाहरणार्थ, $\TRule{\True}{\land}$ हा नियम
$\sFmla{\True}{\varphi \land \psi}$ ला लागू होतो आणि
$\sFmla{\True}{\varphi \land \psi}$ असलेल्या कोणत्याही शाखेच्या शेवटी
$\sFmla{\True}{\varphi}$ आणि~$\sFmla{\True}{\psi}$ ही दोन्ही चिन्हांकित सूत्रे
जोडू देतो; तसेच $\TRule{\False}{\varphi \land \psi}$ हा नियम
$\sFmla{\False}{\varphi}$ आणि $\sFmla{\False}{\psi}$ शेजारी शेजारी जोडून शाखा
फोडू देतो. टॅब्लोच्या प्रत्येक शाखेत $\sFmla{\True}{\varphi}$ आणि
$\sFmla{\False}{\varphi}$ ही चिन्हांकित सूत्रांची जुळणारी जोडी असेल, तर तो
बंद असतो.

टॅब्लोवर आधारित $\Proves$ संबंधाची व्याख्या पुढीलप्रमाणे केली जाते:
$\Gamma \Proves \varphi$ तेव्हा आणि केवळ तेव्हाच, जेव्हा असा काही सांत संच
~$\Gamma_0 = \{\psi_1, \dots, \psi_n\} \subseteq \Gamma$ असतो की पुढील
गृहीतकांसाठी बंद टॅब्लो असतो:
\[\{\sFmla{\False}{\varphi}, \sFmla{\True}{\psi_1}, \dots, \sFmla{\True}{\psi_n}\}
\]
उदाहरणार्थ, $\Proves (\varphi \land \psi) \lif \varphi$ हे दाखवणारा बंद टॅब्लो
पुढे दिला आहे:
\begin{oltableau}
  [\sFmla{\False}{(\varphi \land \psi) \lif \varphi}, just = \TAss
    [\sFmla{\True}{\varphi \land \psi}, just = {\TRule{\False}{\lif}[1]}
      [\sFmla{\False}{\varphi}, just = {\TRule{\False}{\lif}[1]}
        [\sFmla{\True}{\varphi}, just = {\TRule{\True}{\land}[2]}
          [\sFmla{\True}{\psi}, just = {\TRule{\True}{\land}[2]}, close]
        ]
      ]
    ]
  ]
\end{oltableau}

काही $\psi_i \in \Gamma$ साठी पुढील गृहीतकांचा बंद टॅब्लो असेल, तर
संच $\Gamma$ हा टॅब्लो कलनात विसंगत असतो:
\[\{\sFmla{\True}{\psi_1}, \dots, \sFmla{\True}{\psi_n}\}
\]

टॅब्लोचा शोध १९५० च्या दशकात एव्हर्ट बेथ आणि याक्को हिंटिक्का यांनी
स्वतंत्रपणे लावला; रेमंड स्मलियन यांनी त्यांचे सुलभीकरण करून ते लोकप्रिय
केले. टॅब्लो तयार करणे ही अत्यंत पद्धतशीर प्रक्रिया असल्यामुळे त्यांचा
वापर अगदी सोपा आहे. टॅब्लोच्या पद्धतशीर स्वरूपामुळे त्यांची
संगणकीय अंमलबजावणीही सुलभ होते. तथापि, टॅब्लो वाचणे अनेकदा कठीण
असते आणि त्यांचा सिद्धतांशी असलेला संबंध काही वेळा सहज दिसत नाही. ही
दृष्टीही बरीच व्यापक आहे आणि अनेक भिन्न तर्कशास्त्रांसाठी टॅब्लो
पद्धती उपलब्ध आहेत. दिलेली वाक्ये (किंवा त्यांचे संच) पूर्ण करणाऱ्या
संरचना शोधण्यासही टॅब्लो मदत करतात: संच पूर्ततायोग्य असेल,
तर त्याचा बंद टॅब्लो असणार नाही; म्हणजेच कोणत्याही टॅब्लोमध्ये
एक खुली शाखा असेल. त्या शाखेवर लागू करता येणारा प्रत्येक नियम लागू केलेला
असेल, तर खुल्या शाखेवरून पूर्ति करणारी संरचना ``वाचून काढता'' येते.
क्रमवर्ती कलनाशीही याचा अत्यंत निकटचा संबंध आहे: मूलतः, बंद टॅब्लो
म्हणजे उलटी लिहिलेली क्रमवर्ती कलनातील संक्षेपित निष्पत्ती होय.









\section{स्वयंसिद्धकीय निष्पत्ती}\label{fol:prf:axd:sec}

स्वयंसिद्धकीय निष्पत्त्या या सर्वांत जुन्या आणि सोप्या तार्किक
निष्पत्ती-पद्धती आहेत. त्यांतील निष्पत्त्या म्हणजे केवळ
वाक्यांचे अनुक्रम असतात. वाक्यांचा अनुक्रम तेव्हा योग्य
निष्पत्ती मानला जातो, जेव्हा त्यातील प्रत्येक वाक्य~$\varphi$ पुढील
अटींपैकी एक अट पूर्ण करते:
\begin{enumerate}
\item $\varphi$ हे स्वयंसिद्धक असते, किंवा
\item $\varphi$ हे दिलेल्या वाक्यांच्या संच~$\Gamma$ चा घटक असते, किंवा
\item $\varphi$ ला निगमन नियमामुळे समर्थन मिळते.
\end{enumerate}
$\varphi$ हे स्वयंसिद्धक असण्यासाठी त्याचे रूप काही ठरावीक वाक्य
रूपबंधांपैकी एक असावे लागते. प्रथम-क्रम तर्कशास्त्रासाठी समाधानकारक
(निर्दोष आणि संपूर्ण) निष्पत्ती-पद्धत देणारे स्वयंसिद्धक-रूपबंधांचे
अनेक संच आहेत. त्यांपैकी काही संच, ते ज्या संयोजकांना नियंत्रित करतात
त्यांनुसार मांडलेले असतात; उदाहरणार्थ, पुढील रूपबंध
\[\varphi \lif (\psi \lif \varphi) \qquad \psi \lif (\psi \lor \chi) \qquad (\psi \land \chi) \lif \psi
\]
हे $\lif$, $\lor$ आणि~$\land$ यांना नियंत्रित करणारे सर्वसाधारण स्वयंसिद्धक
आहेत. काही स्वयंसिद्धकीय पद्धतींमध्ये स्वयंसिद्धकांची संख्या कमीत कमी
ठेवण्याचे उद्दिष्ट असते. कोणते संयोजक आदिम म्हणून स्वीकारले आहेत यानुसार,
एकाच स्वयंसिद्धकाची पद्धत मिळणेही शक्य असते.

निगमन नियम म्हणजे असे सोपाधिक विधान की जे निष्पत्तीमधील
वाक्य समर्थित असण्यासाठी पुरेशी अट देते. विध्यात्मक अनुमान हा असा
एक अतिशय प्रचलित नियम आहे: $\varphi$ आणि $\varphi \lif \psi$ यांना आधीच समर्थन मिळाले
असेल, तर $\psi$ लाही समर्थन मिळते, असे तो सांगतो. म्हणजे, निष्पत्तीमध्ये वाक्य~$\psi$ असलेल्या ओळीला समर्थन मिळते, मात्र $\varphi$ आणि
$\varphi \lif \psi$ ही दोन्ही सूत्रे (काही वाक्य~$\varphi$ साठी) त्या
निष्पत्तीमध्ये~$\psi$ च्या आधी आलेली असली पाहिजेत.

स्वयंसिद्धकीय निष्पत्त्यांवर आधारित $\Proves$ संबंधाची व्याख्या
पुढीलप्रमाणे केली जाते: $\Gamma \Proves \varphi$ तेव्हा आणि केवळ तेव्हाच, जेव्हा
अशी निष्पत्ती असते की तिच्या शेवटच्या सूत्राच्या जागी
वाक्य~$\varphi$ असते (आणि वरच्या~(2) मुळे समर्थन मिळालेल्या त्या
निष्पत्तीमधील वाक्यांचा संच म्हणजे $\Gamma$ असे घेतले जाते).
~$\varphi$ ची अशी निष्पत्ती असेल की तिच्यासाठी~$\Gamma$ रिक्त असेल, म्हणजेच
त्या निष्पत्तीमधील प्रत्येक वाक्याला (1) किंवा~(3) मुळे समर्थन
मिळाले असेल, तर~$\varphi$ हे प्रमेय असते. उदाहरणार्थ, $\Proves
\varphi  \lif (\psi \lif (\psi \lor \varphi))$ हे दाखवणारी निष्पत्ती पुढे दिली आहे:
\begin{derivation}
  1. & $\psi \lif (\psi \lor \varphi)$ \\
  2. & $(\psi \lif (\psi \lor \varphi)) \lif (\varphi  \lif (\psi \lif (\psi \lor \varphi)))$\\
  3. & $\varphi  \lif (\psi \lif (\psi \lor \varphi))$
\end{derivation}
ओळ~1 वरील वाक्य हे $\varphi \lif (\varphi \lor \psi)$ या स्वयंसिद्धकाच्या
रूपाचे आहे ($\varphi$ आणि $\psi$ यांच्या भूमिका उलट आहेत). ओळ~2 वरील वाक्य
$\varphi \lif (\psi \lif \varphi)$ या स्वयंसिद्धकाच्या रूपाचे आहे. त्यामुळे दोन्ही
ओळींना समर्थन मिळते. ओळ~3 ला विध्यात्मक अनुमानामुळे समर्थन मिळते: त्याचा
संक्षेप $\theta$ असा केल्यास, ओळ~2 चे रूप $\chi \lif \theta$ असे आहे, आणि त्यातील
$\chi$ म्हणजे $\psi \lif (\psi \lor \varphi)$, म्हणजेच ओळ~1 होय.

$\Gamma \Proves \lfalse$ असेल, तर संच $\Gamma$ विसंगत असतो. संपूर्ण
स्वयंसिद्धकीय पद्धत कोणत्याही~$\varphi$ साठी $\lfalse \lif \varphi$ हेही सिद्ध करेल;
त्यामुळे $\Gamma$ विसंगत असेल, तर कोणत्याही~$\varphi$ साठी $\Gamma \Proves \varphi$.

तर्कशास्त्रासाठी स्वयंसिद्धकीय निष्पत्त्यांच्या पद्धती सर्वप्रथम गोटलोप
फ्रेग यांनी त्यांच्या १८७९ मधील \emph{Begriffsschrift} या ग्रंथात दिल्या;
म्हणून हा ग्रंथ आधुनिक तर्कशास्त्रातील पहिला ग्रंथ मानला जातो. ॲल्फ्रेड
नॉर्थ व्हाइटहेड आणि बर्ट्रंड रसेल यांच्या \emph{Principia Mathematica} मध्ये,
तसेच १९२० च्या दशकात डाव्हीट हिल्बर्ट आणि त्यांच्या विद्यार्थ्यांकडून, या
पद्धतींना परिपूर्ण रूप मिळाले. म्हणून त्यांना अनेकदा ``फ्रेग पद्धती'' किंवा
``हिल्बर्ट पद्धती'' म्हणतात. एखाद्या तर्कशास्त्रासाठी स्वयंसिद्धकीय पद्धत
शोधणे बहुधा सोपे असल्यामुळे या पद्धती अतिशय बहुपयोगी आहेत. निष्पत्त्यांची रचना अत्यंत सोपी असून त्यांत फक्त एक किंवा दोन निगमन नियम असतात;
त्यामुळे \emph{त्यांच्याविषयीच्या} गोष्टी सिद्ध करणेही तुलनेने सोपे असते. तथापि,
प्रत्यक्ष व्यवहारात त्यांचा वापर फार कठीण असतो; म्हणजेच सिद्धता शोधणे आणि
लिहिणे कठीण असते.



\chapter{क्रमवर्ती कलन}\label{fol:seq::chap}
\begin{editorial}
  या प्रकरणात अभिजात प्रथम-क्रम तर्कशास्त्रासाठी गेंटझेन यांचे मानक
  क्रमवर्ती कलन LK मांडले आहे. यात आणखी उदाहरणे आणि सराव समाविष्ट करणे
  उपयुक्त ठरेल. सिद्धता-पद्धती म्हणून क्रमवर्ती कलनाशी संबंधित मजकूर
  समाविष्ट करण्यासाठी किंवा वगळण्यासाठी ``prfLK'' हा टॅग वापरा.
\end{editorial}






\section{नियम आणि निष्पत्ती}\label{fol:seq:rul:sec}

पुढील मजकुरात, $\Gamma, \Delta, \Pi, \Lambda$ या वाक्यांच्या सांत
क्रमिका दर्शवतील, असे मानू.

\begin{defn}[क्रमवर्ती]
\emph{क्रमवर्ती} म्हणजे पुढील रूपाची अभिव्यक्ती होय:
\[\Gamma \Sequent \Delta
\]
येथे $\Gamma$ आणि $\Delta$ या $\Lang L$ भाषेतील वाक्यांच्या
सांत (संभवतः रिक्त) क्रमिका आहेत. $\Gamma$ ला \emph{पूर्वांग}, तर
$\Delta$ ला \emph{उत्तरांग} म्हणतात.
\end{defn}

\begin{explain}
क्रमवर्तीमागील अंतःप्रेरित कल्पना अशी आहे: पूर्वांगातील सर्व वाक्ये
सत्य असतील, तर उत्तरांगातील वाक्यांपैकी किमान एक सत्य असते. म्हणजे,
$\Gamma = \tuple{\varphi_1, \dots, \varphi_m}$ आणि
$\Delta = \tuple{\psi_1, \dots, \psi_n}$ असतील, तर $\Gamma \Sequent \Delta$
तेव्हा आणि केवळ तेव्हाच सत्य असते, जेव्हा
\[(\varphi_1 \land \cdots \land \varphi_m) \lif (\psi_1 \lor \cdots \lor
\psi_n)
\]
सत्य असते. याची दोन विशेष प्रकरणे आहेत: $\Gamma$ रिक्त असण्याचे आणि
$\Delta$~रिक्त असण्याचे. $\Gamma$ रिक्त असेल, म्हणजेच $m = 0$ असेल, तर
$\quad \Sequent \Delta$ तेव्हा आणि केवळ तेव्हाच सत्य असते, जेव्हा
$\psi_1 \lor \dots \lor \psi_n$ सत्य असते. $\Delta$ रिक्त असेल, म्हणजेच
$n = 0$ असेल, तर $\Gamma \Sequent \quad$ तेव्हा आणि केवळ तेव्हाच सत्य
असते, जेव्हा $\lnot(\varphi_1 \land \dots \land \varphi_m)$ सत्य असते. एखाद्या
क्रमवर्तीला तेव्हा आणि केवळ तेव्हाच वैध म्हणतो, जेव्हा तिच्याशी संबंधित
वाक्य वैध असते.
\end{explain}

$\Gamma$ ही वाक्यांची क्रमिका असेल, तर $\Gamma$ च्या उजव्या टोकाला
$\varphi$ जोडून मिळणाऱ्या क्रमिकेसाठी आपण $\Gamma, \varphi$ लिहितो (आणि $\Gamma$
च्या डाव्या टोकाला $\varphi$ जोडून मिळणाऱ्या क्रमिकेसाठी $\varphi, \Gamma$
लिहितो). $\Delta$ हीदेखील वाक्यांची क्रमिका असेल, तर $\Gamma,
\Delta$ ही त्या दोन क्रमिकांची जोडणी असते.

\begin{defn}[आरंभीची क्रमवर्ती]
\emph{आरंभीची क्रमवर्ती} म्हणजे, भाषेतील कोणत्याही वाक्य~$\varphi$ साठी,
पुढीलपैकी एका रूपाची क्रमवर्ती होय:
  \begin{enumerate}
    \item $\varphi \Sequent \varphi$
    \item $\quad \Sequent \ltrue$
    \item $\lfalse \Sequent \quad$
  \end{enumerate}.
\end{defn}

क्रमवर्ती कलनातील निष्पत्त्या हे क्रमवर्तींचे विशिष्ट वृक्ष असतात.
त्यांतील सर्वांत वरच्या क्रमवर्ती आरंभीच्या क्रमवर्ती असतात; आणि एखादी
क्रमवर्ती इतर एक किंवा दोन क्रमवर्तींच्या खाली असेल, तर ती निगमनाच्या
एखाद्या नियमानुसार त्यांच्यापासून योग्य रीतीने निष्पन्न झाली पाहिजे.
$\Log{LK}$ चे नियम दोन मुख्य प्रकारांत विभागले आहेत: \emph{तार्किक} नियम
आणि \emph{संरचनात्मक} नियम. खालच्या क्रमवर्तीतील $\varphi$ आणि/किंवा $\psi$
असलेल्या वाक्याच्या मुख्य संकारक वरून तार्किक नियमांना नावे दिली
जातात. प्रत्येक नियमाच्या दोन आवृत्त्या असतात: संकारक असलेले
वाक्य डावीकडे असणारी क्रमवर्ती निष्पन्न करण्यासाठी एक, आणि ते
वाक्य उजवीकडे असणारी क्रमवर्ती निष्पन्न करण्यासाठी दुसरी.









\section{विधानीय नियम}\label{fol:seq:prl:sec}

\subsection{$\lnot$ साठीचे नियम}

\begin{defish}
\Axiom$ \Gamma \fCenter \Delta, \varphi $
\RightLabel{\LeftR{\lnot}}
\UnaryInf$ \lnot \varphi, \Gamma \fCenter \Delta$
\DisplayProof
\hfill
\Axiom$\varphi, \Gamma \fCenter \Delta$
\RightLabel{\RightR{\lnot}}
\UnaryInf$ \Gamma \fCenter \Delta, \lnot \varphi $
\DisplayProof
\end{defish}

\subsection{$\land$ साठीचे नियम}

\begin{defish}\noindent
\begin{tabular}{l}
\Axiom$ \varphi, \Gamma \fCenter \Delta$
\RightLabel{\LeftR{\land}}
\UnaryInf$ \varphi \land \psi, \Gamma \fCenter \Delta$
\DisplayProof
\\[3ex]
\Axiom$\psi, \Gamma \fCenter \Delta$
\RightLabel{\LeftR{\land}}
\UnaryInf$\varphi \land \psi, \Gamma \fCenter \Delta$
\DisplayProof
\end{tabular}
\hfill
\Axiom$\Gamma \fCenter \Delta, \varphi$
\Axiom$ \Gamma \fCenter \Delta, \psi$
\RightLabel{\RightR{\land}}
\BinaryInf$ \Gamma \fCenter \Delta, \varphi \land \psi $
\DisplayProof
\end{defish}
\subsection{$\lor$ साठीचे नियम}
\begin{defish}
\Axiom$\varphi, \Gamma \fCenter \Delta$
\Axiom$\psi, \Gamma \fCenter \Delta$
\RightLabel{\LeftR{\lor}}
\BinaryInf$\varphi \lor \psi, \Gamma \fCenter \Delta$
\DisplayProof
\hfill
\begin{tabular}{r}
\Axiom$\Gamma \fCenter \Delta, \varphi$
\RightLabel{\RightR{\lor}}
\UnaryInf$ \Gamma \fCenter \Delta, \varphi \lor \psi$
\DisplayProof
\\[3ex]
\Axiom$ \Gamma \fCenter \Delta, \psi$
\RightLabel{\RightR{\lor}}
\UnaryInf$ \Gamma \fCenter \Delta, \varphi \lor \psi$
\DisplayProof
\end{tabular}
\end{defish}
\subsection{$\lif$ साठीचे नियम}
\begin{defish}
\Axiom$ \Gamma \fCenter \Delta, \varphi$
\Axiom$ \psi, \Pi \fCenter \Lambda$
\RightLabel{\LeftR{\lif}}
\BinaryInf$ \varphi \lif \psi, \Gamma, \Pi \fCenter \Delta, \Lambda$
\DisplayProof
\hfill
\Axiom$ \varphi, \Gamma \fCenter \Delta, \psi$
\RightLabel{\RightR{\lif}}
\UnaryInf$ \Gamma \fCenter \Delta, \varphi \lif \psi $
\DisplayProof
\end{defish}
\section{संख्यापकांचे नियम}\label{fol:seq:qrl:sec}
\subsection{$\lforall$ साठीचे नियम}
\begin{defish}
\Axiom$ \varphi(t), \Gamma \fCenter \Delta$
\RightLabel{\LeftR{\lforall}}
\UnaryInf$ \lforall[x][\varphi(x)],\Gamma \fCenter \Delta$
\DisplayProof
\hfill
\Axiom$ \Gamma \fCenter \Delta, \varphi(a) $
\RightLabel{\RightR{\lforall}}
\UnaryInf$ \Gamma \fCenter \Delta, \lforall[x][\varphi(x)]$
\DisplayProof
\end{defish}
\LeftR{\lforall} मध्ये, $t$ हे बंद पद आहे (म्हणजे, चरे नसलेले पद).
\RightR{\lforall} मध्ये, $a$~हा असा स्थिरांक आहे की तो
\RightR{\lforall} नियमाच्या खालच्या क्रमवर्तीमध्ये कोठेही येता कामा नये.
\RightR{\forall} अनुमानाचा $a$~हा \emph{आयगेन चर} आहे, असे आपण
म्हणतो.\footnote{वरील नियमात $a$ हा स्थिरांक असला तरी आपण
``आयगेन चर'' ही संज्ञा वापरतो. यामागे ऐतिहासिक कारणे आहेत.}
\enlargethispage{1pt}
\subsection{$\lexists$ साठीचे नियम}
\begin{defish}
\Axiom$ \varphi(a), \Gamma \fCenter \Delta $
\RightLabel{\LeftR{\lexists}}
\UnaryInf$ \lexists[x][\varphi(x)], \Gamma \fCenter \Delta$
\DisplayProof
\hfill
\Axiom$ \Gamma \fCenter \Delta, \varphi(t) $
\RightLabel{\RightR{\lexists}}
\UnaryInf$ \Gamma \fCenter \Delta, \lexists[x][\varphi(x)]$
\DisplayProof
\end{defish}
पुन्हा, $t$~हे बंद पद आहे आणि $a$~हा असा स्थिरांक आहे की तो
\LeftR{\lexists} नियमाच्या खालच्या क्रमवर्तीमध्ये येत नाही.
\LeftR{\lexists} अनुमानाचा $a$~हा \emph{आयगेन चर} आहे, असे आपण म्हणतो.
\RightR{\lforall} किंवा \LeftR{\lexists} अनुमानाच्या खालच्या
क्रमवर्तीमध्ये आयगेन चर येत नसावा, या अटीला \emph{आयगेन चराची अट}
म्हणतात.
\begin{explain}
ही प्रथा आठवा: $\varphi$ हे असे सूत्र असेल की त्यात
चल~$x$ मुक्त असेल, तर हे आपण~$\varphi(x)$ असे लिहून दर्शवतो.
त्याच संदर्भात, मग $\varphi(t)$ हा~$\Subst{\varphi}{t}{x}$ चा संक्षेप असतो.
त्यामुळे \RightR{\lexists} नियम पुढीलप्रमाणेही लिहिता येईल:
\begin{prooftree}
  \Axiom$\Gamma \fCenter \Delta, \Subst{\varphi}{t}{x}$
  \RightLabel{\RightR{\lexists}}
  \UnaryInf$\Gamma \fCenter \Delta, \lexists[x][\varphi]$
\end{prooftree}
लक्षात घ्या की~$\varphi$ मध्ये $t$ आधीच येऊ शकते; उदा., $\varphi$ हे
~$\Atom{\Obj P}{t,x}$ असू शकते. म्हणून,~$\Gamma \Sequent \Delta,
\Atom{\Obj P}{t,t}$ पासून $\Gamma \Sequent \Delta,
\lexists[x][\Atom{\Obj P}{t,x}]$ निष्पन्न करणे हा
\RightR{\lexists} चा योग्य उपयोग आहे---आधारविधानात~$t$ ज्या ठिकाणी
आले आहे, त्यांपैकी एक किंवा अधिक ठिकाणी (म्हणजे, त्या पदाच्या एक किंवा अधिक
आस्थितींच्या जागी) बद्ध चल~$x$ घालता येतो; सर्वच आस्थिती बदलणे आवश्यक नाही. परंतु
\RightR{\lforall} आणि~\LeftR{\lexists} मधील आयगेन चराच्या अटींनुसार
स्थिरांक~$a$ हा~$\varphi$ मध्ये येता कामा नये. म्हणून,
$\RightR{\lforall}$ वापरून $\Gamma \Sequent \Delta, \Atom{\Obj P}{a,a}$
पासून $\Gamma \Sequent \Delta, \lforall[x][\Atom{\Obj P}{a,x}]$
योग्य रीतीने निष्पन्न करता येत नाही.
\end{explain}
\begin{explain}
\RightR{\lexists} आणि \LeftR{\lforall} मध्ये पद~$t$ बंद असण्याव्यतिरिक्त
त्यावर आणखी कोणतेही निर्बंध नाहीत. याउलट, \LeftR{\lexists} आणि \RightR{\lforall}
नियमांमध्ये आयगेन चराच्या अटीनुसार वरच्या क्रमवर्तीमध्ये~$\varphi(a)$ च्या
बाहेर स्थिरांक~$a$ कोठेही येता कामा नये. ही पद्धत निर्दोष आहे,
म्हणजे ती केवळ वैध क्रमवर्तीच निष्पन्न करते, याची खात्री करण्यासाठी ही
अट आवश्यक आहे. ही अट नसती, तर पुढील निष्पत्त्या अनुज्ञात असत्या:
\begin{prooftree}
  \Axiom$\varphi(a) \fCenter \varphi(a)$
  \RightLabel{*\LeftR{\lexists}}
  \UnaryInf$\lexists[x][\varphi(x)] \fCenter \varphi(a)$
  \RightLabel{\RightR{\lforall}}
  \UnaryInf$\lexists[x][\varphi(x)] \fCenter \lforall[x][\varphi(x)]$
  \DisplayProof\bottomAlignProof
  \qquad
  \Axiom$\varphi(a) \fCenter \varphi(a)$
  \RightLabel{*\RightR{\lforall}}
  \UnaryInf$\varphi(a) \fCenter \lforall[x][\varphi(x)]$
  \RightLabel{\LeftR{\lexists}}
  \UnaryInf$\lexists[x][\varphi(x)] \fCenter \lforall[x][\varphi(x)]$
\end{prooftree}
परंतु $\lexists[x][\varphi(x)] \Sequent \lforall[x][\varphi(x)]$ ही क्रमवर्ती
वैध नाही.
\end{explain}
\section{संरचनात्मक नियम}\label{fol:seq:srl:sec}
क्रमवर्तीच्या डाव्या आणि उजव्या बाजूंवरील वाक्यांची मांडणी बदलता येईल,
असे आणखी काही नियम आपल्याला आवश्यक आहेत. तार्किक नियम ज्या वाक्ये वर
लागू होतो, ती आधारविधानात अगदी डाव्या किंवा अगदी उजव्या टोकाला असणे आवश्यक
असते. त्यामुळे वाक्ये योग्य स्थानी हलवता यावीत म्हणून आपल्याला
``अदलाबदल'' नियमाची गरज आहे. तसेच कधी कधी दोन एकसारखी वाक्ये एकात
एकत्र करता येणे आणि दोन्हींपैकी कोणत्याही बाजूला वाक्य जोडता येणे
महत्त्वाचे असते.
\subsection{दुर्बलीकरण}
\begin{defish}
\Axiom$ \Gamma \fCenter \Delta $
\RightLabel{\LeftR{\Weakening}}
\UnaryInf$ \varphi, \Gamma \fCenter \Delta$
\DisplayProof
\hfill
\Axiom$ \Gamma \fCenter \Delta$
\RightLabel{\RightR{\Weakening}}
\UnaryInf$ \Gamma \fCenter \Delta, \varphi$
\DisplayProof
\end{defish}
\subsection{आकुंचन}
\begin{defish}
\Axiom$ \varphi, \varphi, \Gamma \fCenter \Delta $
\RightLabel{\LeftR{\Contraction}}
\UnaryInf$ \varphi, \Gamma \fCenter \Delta$
\DisplayProof
\hfill
\Axiom$ \Gamma \fCenter \Delta, \varphi, \varphi$
\RightLabel{\RightR{\Contraction}}
\UnaryInf$ \Gamma \fCenter \Delta, \varphi$
\DisplayProof
\end{defish}
\subsection{अदलाबदल}
\begin{defish}
\Axiom$ \Gamma, \varphi, \psi, \Pi \fCenter \Delta $
\RightLabel{\LeftR{\Exchange}}
\UnaryInf$ \Gamma, \psi, \varphi, \Pi \fCenter \Delta$
\DisplayProof
\hfill
\Axiom$ \Gamma \fCenter \Delta, \varphi, \psi, \Lambda$
\RightLabel{\RightR{\Exchange}}
\UnaryInf$ \Gamma \fCenter \Delta, \psi, \varphi, \Lambda$
\DisplayProof
\end{defish}
दुर्बलीकरण, आकुंचन आणि अदलाबदल या अनुमानांची मालिका अनेकदा दुहेरी
अनुमानरेषांनी दर्शवली जाते.
पुढील ``कट'' नावाचा नियम काटेकोर अर्थाने आवश्यक नाही; मात्र त्यामुळे
निष्पत्त्या पुन्हा वापरणे आणि एकत्र करणे पुष्कळ सोपे होते.
\begin{defish}
\[
\Axiom$ \Gamma \fCenter \Delta, \varphi$
\Axiom$ \varphi, \Pi \fCenter \Lambda $
\RightLabel{\Cut}
\BinaryInf$ \Gamma, \Pi \fCenter \Delta, \Lambda$
\DisplayProof
\]
\end{defish}









\section{निष्पत्ती}\label{fol:seq:der:sec}

\begin{explain}
आरंभीची क्रमवर्ती कशी असते हे आपण सांगितले आहे आणि निगमन नियम दिले आहेत.
क्रमवर्ती कलनातील निष्पत्त्या यांपासून विगमनाधारित रीतीने निर्माण केल्या
जातात: प्रत्येक निष्पत्ती ही एकतर स्वतःच आरंभीची क्रमवर्ती असते, किंवा
एका अनुमानाच्या वर असलेल्या एक किंवा दोन निष्पत्त्या ने बनलेली असते.
\end{explain}

\begin{defn}[$\Log{LK}$ निष्पत्ती]
क्रमवर्ती~$S$ ची \emph{$\Log{LK}$-निष्पत्ती} म्हणजे पुढील अटी पूर्ण
करणारा क्रमवर्तींचा सांत वृक्ष होय:
\begin{enumerate}
\item वृक्षातील सर्वांत वरच्या क्रमवर्ती आरंभीच्या क्रमवर्ती असतात.
\item वृक्षातील सर्वांत खालची क्रमवर्ती~$S$ असते.
\item वृक्षातील $S$ व्यतिरिक्त प्रत्येक क्रमवर्ती ही निगमन नियमाच्या योग्य
  उपयोजनाचे आधारविधान असते आणि त्या उपयोजनाचा निष्कर्ष वृक्षात त्या
  क्रमवर्तीच्या थेट खाली असतो.
\end{enumerate}
यानंतर $S$ ला त्या निष्पत्तीची \emph{अंतिम क्रमवर्ती} म्हणतो आणि
$S$ ही \emph{$\Log{LK}$ मध्ये निष्पन्न करता येण्याजोगे} आहे (किंवा ती
$\Log{LK}$-निष्पन्न करता येण्याजोगे आहे) असे म्हणतो.
\end{defn}

\begin{ex}
प्रत्येक आरंभीची क्रमवर्ती, उदा., $\chi \Sequent \chi$, ही निष्पत्ती असते.
उदाहरणार्थ, $\LeftR{\Weakening}$ नियम लागू करून आपण तिच्यापासून नवी
निष्पत्ती मिळवू शकतो:
\begin{prooftree}
\Axiom$ \Gamma \fCenter \Delta $
\RightLabel{\LeftR{\Weakening}}
\UnaryInf$ \varphi, \Gamma \fCenter \Delta$
\end{prooftree}
मात्र हा नियम सर्वसाधारण स्वरूपाचा आहे: नियमातील $\varphi$ च्या जागी कोणतेही
वाक्य, उदा., $\theta$ सुद्धा, घेता येते. आधारविधान आपल्या आरंभीच्या
क्रमवर्ती $\chi \Sequent \chi$ शी जुळत असेल, तर $\Gamma$ आणि $\Delta$ ही दोन्ही
केवळ~$\chi$ असतात आणि निष्कर्ष $\theta, \chi \Sequent \chi$ असा असेल. म्हणून पुढील
ही निष्पत्ती आहे:
\begin{prooftree}
\Axiom$ \chi \fCenter \chi $
\RightLabel{\LeftR{\Weakening}}
\UnaryInf$ \theta, \chi \fCenter \chi$
\end{prooftree}
आता आपण दुसरा नियम, समजा $\LeftR{\Exchange}$, लागू करू शकतो; त्यामुळे
डावीकडील दोन वाक्यांची अदलाबदल करता येते. म्हणून पुढीलही योग्य
निष्पत्ती आहे:
\begin{prooftree}
\Axiom$ \chi \fCenter \chi $
\RightLabel{\LeftR{\Weakening}}
\UnaryInf$ \theta, \chi \fCenter \chi$
\RightLabel{\LeftR{\Exchange}}
\UnaryInf$ \chi, \theta \fCenter \chi$
\end{prooftree}
नियमाच्या या उपयोजनात—मूळ नियम पुढीलप्रमाणे दिला होता—
\begin{prooftree}
\Axiom$ \Gamma, \varphi, \psi, \Pi \fCenter \Delta $
\RightLabel{\LeftR{\Exchange}}
\UnaryInf$ \Gamma, \psi, \varphi, \Pi \fCenter \Delta,$
\end{prooftree}
$\Gamma$ आणि $\Pi$ दोन्ही रिक्त ठेवले होते, $\Delta$ हे $\chi$ होते, आणि
$\varphi$ व $\psi$ यांच्या भूमिका अनुक्रमे $\theta$ आणि~$\chi$ यांनी बजावल्या. अगदी
याच प्रकारे पुढीलही निष्पत्ती आहे हे आपल्याला दिसते:
\begin{prooftree}
\Axiom$ \theta \fCenter \theta $
\RightLabel{\LeftR{\Weakening}}
\UnaryInf$ \chi, \theta \fCenter \theta$
\end{prooftree}
आता या दोन निष्पत्त्या घेऊन $\RightR{\land}$ च्या साहाय्याने त्यांना
एकत्र करता येते. तो नियम असा होता:
\begin{prooftree}
\Axiom$\Gamma \fCenter \Delta, \varphi$
\Axiom$ \Gamma \fCenter \Delta, \psi$
\RightLabel{\RightR{\land}}
\BinaryInf$ \Gamma \fCenter \Delta, \varphi \land \psi $
\end{prooftree}
आपल्या उदाहरणात आधारविधाने त्यांवर संपणाऱ्या निष्पत्त्यांच्या अंतिम
क्रमवर्तींशी जुळली पाहिजेत. याचा अर्थ $\Gamma$ हे $\chi, \theta$ आहे, $\Delta$
रिक्त आहे, $\varphi$ हे $\chi$ आहे आणि $\psi$ हे $\theta$ आहे. म्हणून अनुमान योग्य
असेल, तर निष्कर्ष $\chi, \theta \Sequent \chi \land \theta$ हा आहे.
\begin{prooftree}
\Axiom$ \chi \fCenter \chi $
\RightLabel{\LeftR{\Weakening}}
\UnaryInf$ \theta, \chi \fCenter \chi$
\RightLabel{\LeftR{\Exchange}}
\UnaryInf$ \chi, \theta \fCenter \chi$
\Axiom$ \theta \fCenter \theta $
\RightLabel{\LeftR{\Weakening}}
\UnaryInf$ \chi, \theta \fCenter \theta$
\RightLabel{\RightR{\land}}
\BinaryInf$ \chi, \theta \fCenter \chi \land \theta $
\end{prooftree}
अर्थात, आधारविधानांचा क्रम उलटाही करता येतो; मग $\varphi$ हे $\theta$ आणि $\psi$
हे~$\chi$ असेल.
\begin{prooftree}
\Axiom$ \theta \fCenter \theta $
\RightLabel{\LeftR{\Weakening}}
\UnaryInf$ \chi, \theta \fCenter \theta$
\Axiom$ \chi \fCenter \chi $
\RightLabel{\LeftR{\Weakening}}
\UnaryInf$ \theta, \chi \fCenter \chi$
\RightLabel{\LeftR{\Exchange}}
\UnaryInf$ \chi, \theta \fCenter \chi$
\RightLabel{\RightR{\land}}
\BinaryInf$ \chi, \theta \fCenter \theta \land \chi $
\end{prooftree}
\end{ex}









\section{निष्पत्ती: उदाहरणे}\label{fol:seq:pro:sec}

\begin{ex}
क्रमवर्ती $\varphi \land \psi \Sequent \varphi$ ची एक $\Log{LK}$-निष्पत्ती द्या.

हवी असलेली अंतिम क्रमवर्ती निष्पत्तीच्या तळाशी लिहून आपण सुरुवात करतो.
\begin{prooftree}
\AxiomC{}
\UnaryInf$\varphi\land \psi \fCenter \varphi$
\end{prooftree}
यानंतर, या स्वरूपाची खालची क्रमवर्ती कोणत्या प्रकारच्या अनुमानामुळे मिळू शकते
हे शोधावे लागते. ते एखाद्या संरचनात्मक नियमाचे उपयोजन असू शकते; पण आधी
तार्किक नियम शोधणे योग्य ठरते. खालच्या क्रमवर्तीत येणारा एकमेव तार्किक
संयोजक $\land$ आहे; म्हणून आपण $\land$ चा नियम शोधत आहोत. $\land$ हे चिन्ह
पूर्वांगात येत असल्यामुळे आपण \LeftR{\land} नियम पाहतो.
\begin{prooftree}
\AxiomC{}
\RightLabel{\LeftR{\land}}
\UnaryInf$\varphi\land \psi \fCenter \varphi$
\end{prooftree}
\LeftR{\land} अनुमानाची वरची क्रमवर्ती काय असू शकते यासाठी दोन पर्याय आहेत:
वरची क्रमवर्ती $\varphi \Sequent \varphi$ किंवा $\psi \Sequent \varphi$ असू शकते. स्पष्टपणे,
$\varphi \Sequent \varphi$ ही आरंभीची क्रमवर्ती आहे (ही आपल्या दृष्टीने चांगली बाब
आहे), तर $\psi \Sequent \varphi$ ही सर्वसाधारणपणे निष्पन्न करता येत नाही. आपण वरची
क्रमवर्ती भरू:
\begin{prooftree}
\Axiom$\varphi \fCenter \varphi$
\RightLabel{\LeftR{\land}}
\UnaryInf$\varphi\land \psi \fCenter \varphi$
\end{prooftree}
आता आपल्याकडे क्रमवर्ती $\varphi \land \psi \Sequent \varphi$ ची योग्य
$\Log{LK}$-निष्पत्ती आहे.
\end{ex}

\begin{ex}
क्रमवर्ती $\lnot \varphi \lor \psi \Sequent \varphi \lif \psi$ ची एक
$\Log{LK}$-निष्पत्ती द्या.

हवी असलेली अंतिम क्रमवर्ती निष्पत्तीच्या तळाशी लिहून सुरुवात करा.
\begin{prooftree}
\AxiomC{}
\UnaryInf$\lnot \varphi \lor \psi \fCenter \varphi \lif \psi$
\end{prooftree}
ही अंतिम क्रमवर्ती देऊ शकेल असा तार्किक नियम शोधण्यासाठी तिच्यातील तार्किक
संयोजक पाहू: $\lnot$, $\lor$ आणि $\lif$. सध्या आपल्याला फक्त $\lor$ आणि
$\lif$ यांचा विचार करायचा आहे, कारण ती अंतिम क्रमवर्तीतील वाक्यांची
मुख्य संकारक आहेत; $\lnot$ दुसऱ्या संयोजकाच्या व्याप्तीच्या आत आहे,
म्हणून त्याचा विचार आपण नंतर करू. शेवटच्या अनुमानासाठी आपल्याकडे
\LeftR{\lor} नियम आणि \RightR{\lif} नियम हे तार्किक नियमांचे पर्याय आहेत.
खरे तर कोणताही नियम निवडता येईल; पण आपण \RightR{\lif} नियम निवडू (दुसरे
काही कारण नसले तरी त्यामुळे दोन शाखांत विभागणे काही काळ पुढे ढकलता येते).
ही खालची क्रमवर्ती देऊ शकणाऱ्या \RightR{\lif} अनुमानांच्या स्वरूपानुसार ते
पुढीलप्रमाणे असले पाहिजे:
\begin{prooftree}
\AxiomC{}
\UnaryInf$ \varphi, \lnot \varphi \lor \psi \fCenter \psi $
\RightLabel{\RightR{\lif}} \UnaryInf$ \lnot \varphi \lor \psi \fCenter \varphi \lif \psi $
\end{prooftree}

$\lnot \varphi \lor \psi$ हे पूर्वांगाच्या बाहेरील टोकाला हलवले, तर आपण
\LeftR{\lor} नियम लागू करू शकतो. रूपबंधानुसार त्याचे पुढीलप्रमाणे दोन वरच्या
क्रमवर्तींत विभाजन झाले पाहिजे:
\begin{prooftree}
\AxiomC{}
\UnaryInf$\lnot \varphi, \varphi \fCenter \psi$
\AxiomC{}
\UnaryInf$\psi, \varphi \fCenter \psi$
\RightLabel{\LeftR{\lor}}
\BinaryInf$ \lnot \varphi \lor \psi, \varphi \fCenter \psi $
\RightLabel{\LeftR{\Exchange}}
\UnaryInf$ \varphi, \lnot \varphi \lor \psi \fCenter \psi $
\RightLabel{\RightR{\lif}}
\UnaryInf$ \lnot \varphi \lor \psi \fCenter \varphi \lif \psi $
\end{prooftree}
आपण वरच्या दिशेने आरंभीच्या क्रमवर्तींपर्यंत वाट काढण्याचा प्रयत्न करीत आहोत
हे लक्षात ठेवा; आपण बरेच जवळ आलो आहोत असे दिसते!{} उजवी शाखा आरंभीच्या
क्रमवर्तीपासून केवळ एक दुर्बलीकरण आणि एक अदलाबदल दूर आहे; ते केल्यावर ती
पूर्ण होते:
\begin{prooftree}
\AxiomC{}
\UnaryInf$\lnot \varphi, \varphi \fCenter \psi$
\Axiom$\psi \fCenter \psi$
\RightLabel{\LeftR{\Weakening}}
\UnaryInf$\varphi, \psi \fCenter \psi$
\RightLabel{\LeftR{\Exchange}}
\UnaryInf$\psi, \varphi \fCenter \psi$
\RightLabel{\LeftR{\lor}}
\BinaryInf$\lnot \varphi \lor \psi, \varphi \fCenter \psi $
\RightLabel{\LeftR{\Exchange}}
\UnaryInf$ \varphi, \lnot \varphi \lor \psi \fCenter \psi $
\RightLabel{\RightR{\lif}}
\UnaryInf$ \lnot \varphi \lor \psi \fCenter \varphi \lif \psi $
\end{prooftree}

आता डावी शाखा पाहू. कोणत्याही वाक्य मधील एकमेव तार्किक संयोजक हे
पूर्वांगातील वाक्ये मधील $\lnot$ चिन्ह आहे; म्हणून आपण
\LeftR{\lnot} नियमाचे एक उपयोजन शोधत आहोत.
\begin{prooftree}
\AxiomC{}
\UnaryInf$ \varphi \fCenter \psi, \varphi$
\RightLabel{\LeftR{\lnot}}
\UnaryInf$\lnot \varphi, \varphi \fCenter \psi$
\Axiom$\psi \fCenter \psi$
\RightLabel{\LeftR{\Weakening}}
\UnaryInf$\varphi, \psi \fCenter \psi$
\RightLabel{\LeftR{\Exchange}}
\UnaryInf$\psi, \varphi \fCenter \psi$
\RightLabel{\LeftR{\lor}}
\BinaryInf$\lnot \varphi \lor \psi, \varphi \fCenter \psi $
\RightLabel{\LeftR{\Exchange}}
\UnaryInf$ \varphi, \lnot \varphi \lor \psi \fCenter \psi $
\RightLabel{\RightR{\lif}}
\UnaryInf$ \lnot \varphi \lor \psi \fCenter \varphi \lif \psi $
\end{prooftree}
उजवी शाखा जशी पूर्ण केली, त्याचप्रमाणे ही डावी शाखासुद्धा पूर्ण होण्यासाठी
आता केवळ एक दुर्बलीकरण आणि एक अदलाबदल करायची आहे.
\begin{prooftree}
\Axiom$\varphi \fCenter \varphi$
\RightLabel{\RightR{\Weakening}}
\UnaryInf$ \varphi \fCenter \varphi, \psi$
\RightLabel{\RightR{\Exchange}}
\UnaryInf$ \varphi \fCenter \psi, \varphi$
\RightLabel{\LeftR{\lnot}}
\UnaryInf$\lnot \varphi, \varphi \fCenter \psi$
\Axiom$\psi \fCenter \psi$
\RightLabel{\LeftR{\Weakening}}
\UnaryInf$\varphi, \psi \fCenter \psi$
\RightLabel{\LeftR{\Exchange}}
\UnaryInf$\psi, \varphi \fCenter \psi$
\RightLabel{\LeftR{\lor}}
\BinaryInf$\lnot \varphi \lor \psi, \varphi \fCenter \psi $
\RightLabel{\LeftR{\Exchange}}
\UnaryInf$ \varphi, \lnot \varphi \lor \psi \fCenter \psi $
\RightLabel{\RightR{\lif}}
\UnaryInf$ \lnot \varphi \lor \psi \fCenter \varphi \lif \psi $
\end{prooftree}
\end{ex}

\begin{ex}
क्रमवर्ती $\lnot \varphi \lor \lnot \psi \Sequent \lnot (\varphi \land \psi)$ ची एक
$\Log{LK}$-निष्पत्ती द्या.

वरील तंत्रे वापरून, हवी असलेली अंतिम क्रमवर्ती तळाशी लिहून आपण सुरुवात करतो.
\begin{prooftree}
\AxiomC{}
\UnaryInf$ \lnot \varphi \lor \lnot \psi \fCenter \lnot (\varphi \land \psi) $
\end{prooftree}
अंतिम क्रमवर्तीतील वाक्यांची उपलब्ध मुख्य संक्रियके $\lor$ चिन्ह आणि
$\lnot$ चिन्ह ही आहेत. येथे \LeftR{\lor} किंवा \RightR{\lnot} यांपैकी कोणताही
नियम लागू केला तरी चालेल; पण दोन शाखांत विभागणे क्षणभर टाळता यावे म्हणून आपण
\RightR{\lnot} नियमापासून सुरुवात करू:
\begin{prooftree}
\AxiomC{}
\UnaryInf$\varphi \land \psi, \lnot \varphi \lor \lnot \psi \fCenter $
\RightLabel{\RightR{\lnot}}
\UnaryInf$\lnot \varphi \lor \lnot \psi \fCenter \lnot (\varphi \land \psi)$
\end{prooftree}
आता \LeftR{\land} नियम पाहायचा की \LeftR{\lor} नियम, असा पर्याय आहे.
\LeftR{\land} नियम लागू केल्यावर काय होते ते पाहू: $\varphi, \lnot \varphi \lor \lnot \psi
\Sequent \quad$ या क्रमवर्तीपासून किंवा $\psi, \lnot \varphi \lor \lnot \psi \Sequent
\quad$ या क्रमवर्तीपासून सुरुवात करण्याचा पर्याय आपल्याकडे आहे. ही
निष्पत्ती $\varphi$ आणि $\psi$ यांच्या दृष्टीने सममित असल्यामुळे आपण पहिला पर्याय
निवडू:
\begin{prooftree}
\AxiomC{}
\UnaryInf$\varphi, \lnot \varphi \lor \lnot \psi \fCenter $
\RightLabel{\LeftR{\land}}
\UnaryInf$\varphi \land \psi, \lnot \varphi \lor \lnot \psi \fCenter $
\RightLabel{\RightR{\lnot}}
\UnaryInf$\lnot \varphi \lor \lnot \psi \fCenter \lnot (\varphi \land \psi)$
\end{prooftree}
निष्पत्ती पुढे भरत राहिल्यावर आपल्यासमोर एक अडचण येते:
\begin{prooftree}
\Axiom$\varphi \fCenter \varphi$
\RightLabel{\LeftR{\lnot}}
\UnaryInf$ \lnot \varphi, \varphi \fCenter$
\AxiomC{}
\RightLabel{?}
\UnaryInf$\varphi \fCenter \psi$
\RightLabel{\LeftR{\lnot}}
\UnaryInf$ \lnot \psi, \varphi \fCenter$
\RightLabel{\LeftR{\lor}}
\BinaryInf$\lnot \varphi \lor \lnot \psi, \varphi \fCenter $
\RightLabel{\LeftR{\Exchange}}
\UnaryInf$\varphi, \lnot \varphi \lor \lnot \psi \fCenter $
\RightLabel{\LeftR{\land}}
\UnaryInf$\varphi \land \psi, \lnot \varphi \lor \lnot \psi \fCenter $
\RightLabel{\RightR{\lnot}}
\UnaryInf$\lnot \varphi \lor \lnot \psi \fCenter \lnot (\varphi \land \psi)$
\end{prooftree}
उजव्या शाखेच्या वरच्या टोकाचे आणखी लघुकरण करता येत नाही आणि संरचनात्मक
अनुमानांच्या साहाय्याने तिला आरंभीची क्रमवर्ती बनवताही येत नाही; म्हणून हा
योग्य मार्ग नाही. त्यामुळे वर \LeftR{\land} नियम लागू करणे ही स्पष्टपणे चूक
होती. मागील अवस्थेकडे परत जाऊन त्याऐवजी \LeftR{\lor} नियम लागू केल्यावर
आपल्याला पुढील निष्पत्ती मिळते:
\begin{prooftree}
\AxiomC{}
\UnaryInf$\lnot \varphi, \varphi \land \psi \fCenter $

\AxiomC{}
\UnaryInf$\lnot \psi, \varphi \land \psi \fCenter $

\RightLabel{\LeftR{\lor}}
\BinaryInf$\lnot \varphi \lor \lnot \psi, \varphi \land \psi \fCenter $
\RightLabel{\LeftR{\Exchange}}
\UnaryInf$\varphi \land \psi, \lnot \varphi \lor \lnot \psi \fCenter $
\RightLabel{\RightR{\lnot}}
\UnaryInf$\lnot \varphi \lor \lnot \psi \fCenter \lnot (\varphi \land \psi)$
\end{prooftree}
पूर्वी केल्याप्रमाणे प्रत्येक शाखा पूर्ण केल्यावर आपल्याला पुढील निष्पत्ती
मिळते:
\begin{prooftree}
\Axiom$ \varphi \fCenter\varphi$
\RightLabel{\LeftR{\land}}
\UnaryInf$\varphi \land \psi \fCenter \varphi$
\RightLabel{\LeftR{\lnot}}
\UnaryInf$\lnot \varphi, \varphi \land \psi \fCenter $

\Axiom$ \psi \fCenter \psi$
\RightLabel{\LeftR{\land}}
\UnaryInf$\varphi \land \psi \fCenter \psi$
\RightLabel{\LeftR{\lnot}}
\UnaryInf$\lnot \psi, \varphi \land \psi \fCenter $

\RightLabel{\LeftR{\lor}}
\BinaryInf$\lnot \varphi \lor \lnot \psi, \varphi \land \psi \fCenter $
\RightLabel{\LeftR{\Exchange}}
\UnaryInf$\varphi \land \psi, \lnot \varphi \lor \lnot \psi \fCenter $
\RightLabel{\RightR{\lnot}}
\UnaryInf$\lnot \varphi \lor \lnot \psi \fCenter \lnot (\varphi \land \psi)$
\end{prooftree}
(या पायऱ्यांत $\lnot$ नियमांपेक्षा खाली $\land$ नियम लागू केले असते, तरीही
आपल्याला योग्य निष्पत्ती मिळाली असती).
\end{ex}

\begin{ex}
आतापर्यंत आपण आकुंचन नियम वापरलेला नाही; परंतु कधी कधी तो आवश्यक असतो. असे
घडणारे एक उदाहरण पाहू. आपल्याला $\quad \Sequent \varphi \lor \lnot \varphi$ सिद्ध
करायचे आहे असे समजा. $\RightR{\lor}$ उलट्या दिशेने लागू केल्यास पुढील दोन
निष्पत्त्या पैकी एक मिळेल:
\begin{prooftree}
\AxiomC{}
\UnaryInf$ \fCenter \varphi$
\RightLabel{\RightR{\lor}}
\UnaryInf$ \fCenter \varphi \lor \lnot \varphi$
\DisplayProof\qquad\bottomAlignProof
\AxiomC{}
\UnaryInf$\varphi \fCenter $
\RightLabel{\RightR{\lnot}}
\UnaryInf$ \fCenter \lnot \varphi$
\RightLabel{\RightR{\lor}}
\UnaryInf$ \fCenter \varphi \lor \lnot \varphi$
\end{prooftree}
यांपैकी कोणतीही अर्थातच आरंभीच्या क्रमवर्तीवर संपत नाही. यासाठीची युक्ती अशी
की आकुंचन नियम आपल्याला वाक्याच्या दोन प्रती एकात एकत्र करू देतो—आणि
सिद्धता शोधताना, म्हणजे तळापासून वर जाताना, आपण आधारविधानात $\varphi \lor \lnot
\varphi$ ची एक प्रत राखून ठेवू शकतो, उदा.,
\begin{prooftree}
\AxiomC{}
\UnaryInf$ \fCenter \varphi \lor \lnot \varphi, \varphi$
\RightLabel{\RightR{\lor}}
\UnaryInf$ \fCenter \varphi \lor \lnot \varphi, \varphi \lor \lnot \varphi$
\RightLabel{\RightR{\Contraction}}
\UnaryInf$ \fCenter \varphi \lor \lnot \varphi$
\end{prooftree}
आता आपण $\RightR{\lor}$ दुसऱ्यांदा लागू करून~$\lnot \varphi$ सुद्धा मिळवू शकतो;
त्यातून एक पूर्ण निष्पत्ती मिळते.
\begin{prooftree}
\Axiom$\varphi \fCenter \varphi$
\RightLabel{\RightR{\lnot}}
\UnaryInf$\fCenter \varphi, \lnot \varphi$
\RightLabel{\RightR{\lor}}
\UnaryInf$\fCenter \varphi, \varphi \lor \lnot \varphi$
\RightLabel{\RightR{\Exchange}}
\UnaryInf$ \fCenter \varphi \lor \lnot \varphi, \varphi$
\RightLabel{\RightR{\lor}}
\UnaryInf$ \fCenter \varphi \lor \lnot \varphi, \varphi \lor \lnot \varphi$
\RightLabel{\RightR{\Contraction}}
\UnaryInf$ \fCenter \varphi \lor \lnot \varphi$
\end{prooftree}
\end{ex}

\begin{prob}
पुढील क्रमवर्तींसाठी निष्पत्त्या तयार करा:
\begin{enumerate}
\item $\varphi \land (\psi \land \chi) \Sequent (\varphi \land \psi) \land \chi$.
\item $\varphi \lor (\psi \lor \chi) \Sequent (\varphi \lor \psi) \lor \chi$.
\item $\varphi \lif (\psi \lif \chi) \Sequent \psi \lif (\varphi \lif \chi)$.
\item $\varphi \Sequent \lnot\lnot \varphi$.
\end{enumerate}
\end{prob}

\begin{prob}
पुढील क्रमवर्तींसाठी निष्पत्त्या तयार करा:
\begin{enumerate}
\item $(\varphi \lor \psi) \lif \chi \Sequent \varphi \lif \chi$.
\item $(\varphi \lif \chi) \land (\psi \lif \chi) \Sequent (\varphi \lor \psi) \lif \chi$.
\item $\Sequent \lnot(\varphi \land \lnot \varphi)$.
\item $\psi \lif \varphi \Sequent \lnot \varphi \lif \lnot \psi$.
\item $\Sequent (\varphi \lif \lnot \varphi) \lif \lnot \varphi$.
\item $\Sequent \lnot(\varphi \lif \psi) \lif \lnot \psi$.
\item $\varphi \lif \chi \Sequent \lnot (\varphi \land \lnot \chi)$.
\item $\varphi \land \lnot \chi \Sequent \lnot (\varphi \lif \chi)$.
\item $\varphi \lor \psi, \lnot \psi \Sequent \varphi$.
\item $\lnot \varphi \lor \lnot \psi \Sequent \lnot(\varphi \land \psi)$.
\item $\Sequent (\lnot \varphi \land \lnot \psi) \lif\lnot(\varphi \lor \psi)$.
\item $\Sequent \lnot(\varphi \lor \psi) \lif (\lnot \varphi \land \lnot \psi)$.
\end{enumerate}
\end{prob}

\begin{prob}
पुढील क्रमवर्तींसाठी निष्पत्त्या तयार करा:
\begin{enumerate}
\item $\lnot(\varphi \lif \psi) \Sequent \varphi$.
\item $\lnot(\varphi \land \psi) \Sequent \lnot \varphi \lor \lnot \psi$.
\item $\varphi \lif \psi \Sequent \lnot \varphi \lor \psi$.
\item $\Sequent \lnot \lnot \varphi \lif \varphi$.
\item $\varphi \lif \psi, \lnot \varphi \lif \psi \Sequent \psi$.
\item $(\varphi \land \psi) \lif \chi \Sequent (\varphi \lif \chi) \lor (\psi \lif \chi)$.
\item $(\varphi \lif \psi) \lif \varphi \Sequent \varphi$.
\item $\Sequent (\varphi \lif \psi) \lor (\psi \lif \chi)$.
\end{enumerate}
(या सर्वांसाठी $\RightR{\Contraction}$~नियम आवश्यक आहे.)
\end{prob}








\section{संख्यापकांसह निष्पत्ती}\label{fol:seq:prq:sec}

\begin{ex}
क्रमवर्ती $\lexists[x][\lnot \varphi(x)]
\Sequent \lnot \lforall[x][\varphi(x)]$ ची एक $\Log{LK}$-निष्पत्ती द्या.

संख्यापक हाताळताना आयगेन चराच्या अटीचा भंग होणार नाही याची खात्री करावी
लागते; काही वेळा त्यासाठी ठरावीक अनुमाने कोणत्या क्रमाने करायची, यात फेरबदल
करावा लागतो. सामान्यतः आयगेन चराची अट लागू असलेले नियम आधी हाताळण्याचा
प्रयत्न केल्यास मदत होते (पूर्ण झालेल्या सिद्धतेत ते खालच्या बाजूस असतील).
तसेच पुढे काय येईल याचा विचार करून आरंभीची क्रमवर्ती कशी असेल याचा अंदाज
लावणे उपयुक्त ठरते. आपल्या बाबतीत ती $\varphi(a) \Sequent \varphi(a)$ यासारखी असावी
लागेल. याचा अर्थ असा की संख्यापकांचे नियम ``उलट्या दिशेने'' लागू करताना
$\lforall$ नियम आणि $\lexists$ नियम या दोन्हींसाठी एकच पद---ज्याला आपण $a$
म्हणू---निवडावे लागेल. प्रत्येक नियमासाठी वेगवेगळी पदे निवडली, तर शेवटी
$\varphi(a) \Sequent \varphi(b)$ यासारखी क्रमवर्ती मिळेल; ती अर्थातच निष्पन्न करता
येत नाही.

नेहमीप्रमाणे आपण पुढील क्रमवर्ती लिहून सुरुवात करतो:
\begin{prooftree}
\AxiomC{}
\UnaryInf$\lexists[x][\lnot \varphi(x)] \fCenter \lnot \lforall[x][\varphi(x)]$
\end{prooftree}
आपण \LeftR{\exists} नियम किंवा \RightR{\lnot} नियम यांपैकी कोणताही लागू करू
शकतो. \LeftR{\exists} नियमाला आयगेन चराची अट लागू असल्यामुळे तो उशिरा
हाताळण्यापेक्षा लवकर हाताळणे योग्य ठरेल; म्हणून तोच नियम आधी लागू करू.
\begin{prooftree}
\AxiomC{}
\UnaryInf$ \lnot \varphi(a) \fCenter \lnot \lforall[x][\varphi(x)]$
\RightLabel{\LeftR{\lexists}}
\UnaryInf$ \lexists[x][\lnot \varphi(x)] \fCenter \lnot \lforall[x][\varphi(x)]$
\end{prooftree}
\LeftR{\lnot} आणि \RightR{\lnot} नियम उलट्या दिशेने लागू केल्यावर पुढील
निष्पत्ती मिळते:
\begin{prooftree}
\AxiomC{}
\UnaryInf$\lforall[x][\varphi(x)] \fCenter \varphi(a)$
\RightLabel{\LeftR{\lnot}}
\UnaryInf$\lnot \varphi(a), \lforall[x][\varphi(x)] \fCenter $
\RightLabel{\LeftR{\Exchange}}
\UnaryInf$\lforall[x][\varphi(x)], \lnot \varphi(a) \fCenter $
\RightLabel{\RightR{\lnot}}
\UnaryInf$ \lnot \varphi(a) \fCenter \lnot \lforall[x] \varphi(x)$
\RightLabel{\LeftR{\lexists}}
\UnaryInf$ \lexists[x] \lnot \varphi(x) \fCenter \lnot \lforall[x] \varphi(x)$
\end{prooftree}
या टप्प्यावर \LeftR{\forall} नियम लागू करणे हाच आपल्यापुढील एकमेव पर्याय
आहे. या नियमाला आयगेन चराचे बंधन लागू नसल्यामुळे आपला मार्ग मोकळा आहे.
आपल्याला आरंभीची क्रमवर्ती (जिचे स्वरूप $\varphi(a) \Sequent \varphi(a)$ असे आहे)
मिळवण्याचा प्रयत्न करायचा आहे, हे आठवा; म्हणून नियम लागू करताना $\varphi$ मध्ये
घालायचे पद म्हणून $a$ निवडले पाहिजे.
\begin{prooftree}
\Axiom$\varphi(a) \fCenter \varphi(a)$
\RightLabel{\LeftR{\lforall}}
\UnaryInf$\lforall[x][\varphi(x)] \fCenter \varphi(a)$
\RightLabel{\LeftR{\lnot}}
\UnaryInf$\lnot \varphi(a), \lforall[x][\varphi(x)] \fCenter $
\RightLabel{\LeftR{\Exchange}}
\UnaryInf$\lforall[x][\varphi(x)], \lnot \varphi(a) \fCenter $
\RightLabel{\RightR{\lnot}}
\UnaryInf$ \lnot \varphi(a) \fCenter \lnot \lforall[x][\varphi(x)]$
\RightLabel{\LeftR{\lexists}}
\UnaryInf$ \lexists[x][ \lnot \varphi(x)] \fCenter \lnot \lforall[x][\varphi(x)]$
\end{prooftree}
विशेषतः संख्यापक हाताळताना, या टप्प्यावर आयगेन चराच्या अटीचा भंग झालेला नाही
ना, हे पुन्हा तपासणे महत्त्वाचे असते. आयगेन चराची अट लागू असलेला आपण वापरलेला
एकमेव नियम \LeftR{\exists} होता आणि आयगेन चर~$a$ त्याच्या खालच्या क्रमवर्तीत
(म्हणजे अंतिम क्रमवर्तीत) येत नाही; म्हणून ही योग्य निष्पत्ती आहे.
\end{ex}

\begin{prob}
पुढील क्रमवर्तींसाठी निष्पत्त्या तयार करा:
\begin{enumerate}
\item $\Sequent (\lforall[x][\varphi(x)] \land \lforall[y][\psi(y)]) \lif
\lforall[z][(\varphi(z) \land \psi(z))]$.
\item $\Sequent (\lexists[x][\varphi(x)] \lor \lexists[y][\psi(y)]) \lif
\lexists[z][(\varphi(z) \lor \psi(z))]$.
\item $\lforall[x][(\varphi(x) \lif \psi)] \Sequent \lexists[y][\varphi(y)] \lif \psi$.
\item $\lforall[x][\lnot \varphi(x)] \Sequent \lnot\lexists[x][\varphi(x)]$.
\item $\Sequent \lnot\lexists[x][\varphi(x)] \lif \lforall[x][\lnot \varphi(x)]$.
\item $\Sequent \lnot\lexists[x][\lforall[y][((\varphi(x,y) \lif \lnot
\varphi(y,y)) \land (\lnot \varphi(y,y) \lif \varphi(x,y)))]]$.
\end{enumerate}
\end{prob}

\begin{prob}
पुढील क्रमवर्तींसाठी निष्पत्त्या तयार करा:
\begin{enumerate}
\item $\Sequent \lnot\lforall[x][\varphi(x)] \lif \lexists[x][\lnot\varphi(x)]$.
\item $(\lforall[x][\varphi(x)] \lif \psi) \Sequent \lexists[y][(\varphi(y) \lif \psi)]$.
\item $\Sequent \lexists[x][(\varphi(x) \lif \lforall[y][\varphi(y)])]$.
\end{enumerate}
(या सर्वांसाठी $\RightR{\Contraction}$~नियम आवश्यक आहे.)
\end{prob}







\begin{editorial}
  या विभागात नैसर्गिक निगमनासाठी सिद्धतायोग्यता संबंध आणि सुसंगतता यांच्या
  व्याख्या एकत्रित केल्या आहेत.
\end{editorial}



\section{सिद्धता-उपपत्तीय संकल्पना}\label{fol:seq:ptn:sec}

\begin{explain}
जशा आपण अनेक महत्त्वाच्या चिन्हार्थविषयक संकल्पना (वैधता, तार्किक निष्पन्नता,
पूर्ततायोग्यता) परिभाषित केल्या आहेत, तशाच त्यांना अनुरूप
\emph{सिद्धता-उपपत्तीय संकल्पना} आता परिभाषित करू. संरचना मध्ये
वाक्यांची पूर्ति होते की नाही, याचा आधार घेऊन या संकल्पना परिभाषित
केल्या जात नाहीत; त्याऐवजी ठरावीक क्रमवर्तींची निष्पन्नता किंवा
अनिष्पन्नता यांचा आधार घेतला जातो. या दोन प्रकारच्या संकल्पना
एकमेकांशी जुळतात, हा एक महत्त्वाचा शोध होता. त्या जुळतात, हाच
\emph{निर्दोषता प्रमेय} आणि \emph{संपूर्णता प्रमेय} यांचा आशय आहे.
\end{explain}

\begin{defn}[प्रमेये]
वाक्य~$\varphi$ हे \emph{प्रमेय} असते, जर $\Log{LK}$ मध्ये
$\quad \Sequent \varphi$ या क्रमवर्तीची निष्पत्ती असेल. $\varphi$ हे प्रमेय
असेल तर आपण $\Proves \varphi$ लिहितो आणि ते प्रमेय नसेल तर $\Proves/ \varphi$ लिहितो.
\end{defn}

\begin{defn}[निष्पन्नता]
वाक्य~$\varphi$ हे वाक्यांच्या संच~$\Gamma$ \emph{पासून
निष्पन्न करता येण्याजोगे} असते, म्हणजे $\Gamma \Proves \varphi$, तेव्हा आणि केवळ तेव्हाच,
जेव्हा एखादा सांत उपसंच~$\Gamma_0 \subseteq \Gamma$ आणि $\Gamma_0$ मधील
वाक्यांची एखादी क्रमिका $\Gamma_0'$ अशी असते की $\Log{LK}$
$\Gamma_0' \Sequent \varphi$ हे निष्पन्न करते. $\varphi$ हे $\Gamma$ पासून
निष्पन्न करता येण्याजोगे नसेल, तर आपण $\Gamma \Proves/ \varphi$ लिहितो.
\end{defn}

आकुंचन, दुर्बलीकरण आणि अदलाबदल नियमांमुळे $\Gamma_0'$ मधील
वाक्यांचा क्रम आणि त्यांची पुनरावृत्ती कितीदा होते हे महत्त्वाचे नसते:
$\Gamma_0' \Sequent \varphi$ ही क्रमवर्ती निष्पन्न करता येण्याजोगे असेल, तर
$\Gamma_0'$ मध्ये असलेली तीच वाक्ये असणाऱ्या कोणत्याही
$\Gamma_0''$ साठी $\Gamma_0'' \Sequent \varphi$ ही क्रमवर्तीसुद्धा निष्पन्न करता
येते. उदाहरणार्थ, $\Gamma_0 = \{\psi, \chi\}$ असेल, तर
$\Gamma_0' = \tuple{\psi, \psi, \chi}$ आणि $\Gamma_0'' =
\tuple{\chi, \chi, \psi}$ या दोन्ही क्रमिकांमध्ये $\Gamma_0$ मधीलच
वाक्ये आहेत. यांपैकी एक क्रमिका असणारी क्रमवर्ती निष्पन्न करता येण्याजोगे असेल,
तर दुसरी असणारी क्रमवर्तीसुद्धा निष्पन्न करता येते; उदा.:
\begin{prooftree}
  \AxiomC{}
  \Deduce$\psi, \psi, \chi \fCenter \varphi$
  \RightLabel{\LeftR{\Contraction}}
  \UnaryInf$\psi, \chi \fCenter \varphi$
  \RightLabel{\LeftR{\Exchange}}
  \UnaryInf$\chi, \psi \fCenter \varphi$
  \RightLabel{\LeftR{\Weakening}}
  \UnaryInf$\chi, \chi, \psi \fCenter \varphi$
\end{prooftree}
यापुढे $\Gamma_0$ हा वाक्यांचा सांत संच असेल, तर
$\Gamma_0 \Sequent \varphi$ याचा अर्थ असा कोणताही क्रमवर्ती घेऊ की ज्याच्या
पूर्वांगात $\Gamma_0$ मधील वाक्यांची क्रमिका आहे; तसेच आवश्यक असल्यास
आकुंचन, अदलाबदल आणि दुर्बलीकरणे गृहीत धरू.

\begin{defn}[सुसंगतता]
वाक्यांचा संच~$\Gamma$ हा \emph{विसंगत} असतो तेव्हा आणि केवळ तेव्हाच, जेव्हा
एखादा सांत उपसंच~$\Gamma_0 \subseteq \Gamma$ असा असतो की $\Log{LK}$
$\Gamma_0 \Sequent \quad$ हे निष्पन्न करते. $\Gamma$ विसंगत नसेल,
म्हणजे प्रत्येक सांत $\Gamma_0 \subseteq \Gamma$ साठी $\Log{LK}$
$\Gamma_0 \Sequent \quad$ हे निष्पन्न करत नसेल, तर त्याला
\emph{सुसंगत} म्हणतो.
\end{defn}

\begin{prop}[परावर्तिता]
\label{fol:seq:ptn:prop:reflexivity}
$\varphi \in \Gamma$ असेल, तर $\Gamma \Proves \varphi$.
\end{prop}

\begin{proof}
$\varphi \Sequent \varphi$ ही आरंभीची क्रमवर्ती निष्पन्न करता येण्याजोगे आहे आणि $\{\varphi\}
\subseteq \Gamma$.
\end{proof}

\begin{prop}[एकस्वनिकता]
\label{fol:seq:ptn:prop:monotonicity}
$\Gamma \subseteq \Delta$ आणि $\Gamma \Proves \varphi$ असेल, तर $\Delta
\Proves \varphi$.
\end{prop}

\begin{proof}
$\Gamma \Proves \varphi$ असे समजा; म्हणजे, एखादा सांत $\Gamma_0
\subseteq \Gamma$ असा आहे की $\Gamma_0 \Sequent \varphi$ ही क्रमवर्ती
निष्पन्न करता येण्याजोगे आहे. $\Gamma \subseteq \Delta$ असल्यामुळे $\Gamma_0$ हा
$\Delta$ चासुद्धा सांत उपसंच आहे. म्हणून $\Gamma_0 \Sequent \varphi$ ची
निष्पत्ती ही $\Delta \Proves \varphi$ हेसुद्धा दाखवते.
\end{proof}

\begin{prop}[संक्रमकता]
\label{fol:seq:ptn:prop:transitivity}
$\Gamma \Proves \varphi$ आणि $\{\varphi\} \cup \Delta \Proves
\psi$ असेल, तर $\Gamma \cup \Delta \Proves \psi$.
\end{prop}

\begin{proof}
$\Gamma \Proves \varphi$ असेल, तर एखादा सांत $\Gamma_0 \subseteq \Gamma$
आणि $\Gamma_0 \Sequent \varphi$ ची निष्पत्ती~$\pi_0$ असते.
$\{\varphi\} \cup \Delta \Proves \psi$ असेल, तर एखाद्या सांत उपसंचासाठी
$\Delta_0 \subseteq \Delta$, $\varphi, \Delta_0 \Sequent \psi$ ची
निष्पत्ती~$\pi_1$ असते. पुढील निष्पत्ती विचारात घ्या:
\begin{prooftree}
  \AxiomC{}
  \RightLabel{$\pi_0$}
  \Deduce$\Gamma_0 \fCenter \varphi$
  \AxiomC{}
  \RightLabel{$\pi_1$}
  \Deduce$\varphi, \Delta_0 \fCenter \psi$
  \RightLabel{\Cut}
  \BinaryInf$\Gamma_0, \Delta_0 \fCenter \psi$
\end{prooftree}
$\Gamma_0 \cup \Delta_0 \subseteq \Gamma \cup \Delta$ असल्यामुळे यावरून
$\Gamma \cup \Delta \Proves \psi$ हे दिसते.
\end{proof}

विशेषतः, याचा अर्थ असा की $\Gamma \Proves \varphi$ आणि $\varphi
\Proves \psi$ असेल, तर $\Gamma \Proves \psi$. तसेच $\varphi_1,
\dots, \varphi_n \Proves \psi$ असेल आणि प्रत्येक~$i$ साठी $\Gamma \Proves \varphi_i$
असेल, तर $\Gamma \Proves \psi$, हेसुद्धा यावरून मिळते.

\begin{prop}
\label{fol:seq:ptn:prop:incons}
$\Gamma$ विसंगत असतो तेव्हा आणि केवळ तेव्हाच, जेव्हा प्रत्येक
  वाक्य~$\varphi$ साठी $\Gamma \Proves {\varphi}$.
\end{prop}

\begin{proof}
सराव.
\end{proof}

\begin{prob}
\ref{fol:seq:ptn:prop:incons} सिद्ध करा.
\end{prob}

\begin{prop}[संहतता]
  \label{fol:seq:ptn:prop:proves-compact}
  \begin{enumerate}
  \item $\Gamma \Proves \varphi$ असेल, तर एखादा सांत उपसंच $\Gamma_0
    \subseteq \Gamma$ असा आहे की $\Gamma_0 \Proves \varphi$.
  \item $\Gamma$ चा प्रत्येक सांत उपसंच सुसंगत असेल, तर $\Gamma$ सुसंगत
    असतो.
  \end{enumerate}
\end{prop}

\begin{proof}
  \begin{enumerate}
    \item $\Gamma \Proves \varphi$ असेल, तर एखादा सांत उपसंच
      $\Gamma_0 \subseteq \Gamma$ असा आहे की $\Gamma_0
      \Sequent \varphi$ या क्रमवर्तीची निष्पत्ती आहे. परिणामी,
      $\Gamma_0 \Proves \varphi$.
    \item $\Gamma$ विसंगत असेल, तर एखादा सांत
      उपसंच~$\Gamma_0 \subseteq \Gamma$ असा आहे की $\Log{LK}$
      $\Gamma_0 \Sequent \quad$ हे निष्पन्न करते. पण मग $\Gamma_0$ हा
      $\Gamma$ चा विसंगत सांत उपसंच आहे.
  \end{enumerate}
\end{proof}









\section{निष्पन्नता आणि सुसंगतता}\label{fol:seq:prv:sec}

आता आपण निष्पन्नता संबंधाचे अनेक गुणधर्म सिद्ध करू. ते स्वतंत्रपणेही
महत्त्वाचे आहेत, पण संपूर्णता प्रमेयाच्या सिद्धतेत प्रत्येकाची भूमिका असेल.

\begin{prop}\label{fol:seq:prv:prop:provability-contr}
  जर $\Gamma \Proves \varphi$ आणि $\Gamma \cup \{\varphi\}$ विसंगत असेल, तर
  $\Gamma$ विसंगत आहे.
\end{prop}

\begin{proof}
अशी सांत $\Gamma_0$ आणि $\Gamma_1 \subseteq \Gamma$ असतात की
$\Log{LK}$ निष्पन्न $\Gamma_0 \Sequent \varphi$ आणि $\varphi, \Gamma_1
\Sequent \quad$. $\Gamma_0 \Sequent \varphi$ ची $\Log{LK}$-निष्पत्ती
~$\pi_0$ आणि $\Gamma_1, \varphi \Sequent \quad$ ची $\Log{LK}$-निष्पत्ती
~$\pi_1$ अशी नावे देऊ. मग आपण निष्पन्न करू शकतो
\begin{prooftree}
\AxiomC{}
\RightLabel{$\pi_0$}
\Deduce$ \Gamma_0 \fCenter \varphi $
\AxiomC{}
\RightLabel{$\pi_1$}
\Deduce$\varphi, \Gamma_1 \fCenter $
\RightLabel{\Cut}
\BinaryInf$ \Gamma_0,\Gamma_1 \fCenter $
\end{prooftree}

कारण $\Gamma_0 \subseteq \Gamma$ आणि $\Gamma_1 \subseteq \Gamma$,
$\Gamma_0 \cup \Gamma_1 \subseteq \Gamma$; म्हणून $\Gamma$ विसंगत आहे.
\end{proof}

\begin{prop}
\label{fol:seq:prv:prop:prov-incons}
$\Gamma \Proves \varphi$ तेव्हा आणि केवळ तेव्हाच, जेव्हा $\Gamma \cup \{\lnot \varphi\}$
विसंगत असते.
\end{prop}

\begin{proof}
प्रथम $\Gamma \Proves \varphi$ असे गृहीत धरा; म्हणजे $\Gamma \Sequent \varphi$ ची
निष्पत्ती~$\pi_0$ आहे. $\LeftR{\lnot}$ नियम जोडल्यावर आपल्याला
$\lnot \varphi, \Gamma
\Sequent \quad$ ची निष्पत्ती मिळते; म्हणजे $\Gamma \cup \{\lnot \varphi\}$
विसंगत आहे.

जर $\Gamma \cup \{\lnot \varphi\}$ विसंगत असेल, तर $\lnot \varphi, \Gamma
\Sequent \quad$ ची निष्पत्ती~$\pi_1$ असते. खालील $\Gamma \Sequent \varphi$
ची निष्पत्ती आहे:
\begin{prooftree}
  \Axiom$\varphi \fCenter \varphi$
  \RightLabel{\RightR{\lnot}}
  \UnaryInf$\fCenter \varphi, \lnot \varphi$
  \AxiomC{}
  \RightLabel{$\pi_1$}
  \Deduce$\lnot \varphi, \Gamma \fCenter$
  \RightLabel{\Cut}
  \BinaryInf$\Gamma \fCenter \varphi$
\end{prooftree}
\end{proof}

\begin{prob}
$\Gamma \Proves \lnot \varphi$ तेव्हा आणि केवळ तेव्हाच, जेव्हा $\Gamma \cup \{\varphi\}$
विसंगत असते, हे सिद्ध करा.
\end{prob}

\begin{prop}\label{fol:seq:prv:prop:explicit-inc}
  जर $\Gamma \Proves \varphi$ आणि $\lnot \varphi \in \Gamma$ असेल, तर $\Gamma$
  विसंगत आहे.
\end{prop}

\begin{proof}
  समजा $\Gamma \Proves \varphi$ आणि $\lnot \varphi \in \Gamma$. मग $\Gamma_0 \Sequent
  \varphi$ ची निष्पत्ती~$\pi$ असते. $\lnot \varphi, \Gamma_0 \Sequent \quad$
  ही क्रमवर्तीही निष्पन्न करता येण्याजोगे आहे:
  \begin{prooftree}
    \AxiomC{}
    \RightLabel{$\pi$}
    \Deduce$\Gamma_0 \fCenter \varphi$
    \Axiom$\varphi \fCenter \varphi$
    \RightLabel{\LeftR{\lnot}}
    \UnaryInf$\lnot \varphi, \varphi \fCenter$
    \RightLabel{\LeftR{\Exchange}}
    \UnaryInf$\varphi, \lnot \varphi \fCenter$
    \RightLabel{\Cut}
    \BinaryInf$\Gamma_0, \lnot \varphi \fCenter$
  \end{prooftree}
  कारण $\lnot \varphi \in \Gamma$ आणि $\Gamma_0 \subseteq \Gamma$, यावरून
  $\Gamma$ विसंगत असल्याचे दिसते.
\end{proof}

\begin{prop}\label{fol:seq:prv:prop:provability-exhaustive}
  जर $\Gamma \cup \{\varphi\}$ आणि $\Gamma \cup \{\lnot \varphi\}$ दोन्ही
  विसंगत असतील, तर $\Gamma$ विसंगत आहे.
\end{prop}

\begin{proof}
अशी सांत संच $\Gamma_0 \subseteq \Gamma$ आणि $\Gamma_1
\subseteq \Gamma$, तसेच $\Log{LK}$-निष्पत्त्या $\pi_0$ आणि $\pi_1$
$\varphi, \Gamma_0 \Sequent \quad$ आणि $\lnot \varphi, \Gamma_1 \Sequent
\quad$ यांचे अनुक्रमे, असतात. मग आपण निष्पन्न करू शकतो
\begin{prooftree}
\AxiomC{}
\RightLabel{$\pi_0$}
\Deduce$ \varphi, \Gamma_0 \fCenter $
\RightLabel{\RightR{\lnot}}
\UnaryInf$ \Gamma_0 \fCenter \lnot \varphi$
\AxiomC{}
\RightLabel{$\pi_1$}
\Deduce$\lnot \varphi, \Gamma_1 \fCenter  $
\RightLabel{\Cut}
\BinaryInf$ \Gamma_0, \Gamma_1 \fCenter $
\end{prooftree}
कारण $\Gamma_0 \subseteq \Gamma$ आणि $\Gamma_1 \subseteq \Gamma$,
$\Gamma_0 \cup \Gamma_1 \subseteq \Gamma$. म्हणून $\Gamma$ विसंगत आहे.
\end{proof}










\section{निष्पन्नता आणि विधानीय संयोजक}\label{fol:seq:ppr:sec}

\begin{explain}
  क्रमवर्ती कलनातील निष्पन्नता संबंध~$\Proves$ हा विधानीय संयोजकांशी
  संबंधित काही मूलभूत तथ्ये सिद्ध करण्यासाठी पुरेसा सामर्थ्यवान आहे; उदाहरणार्थ,
  $\varphi \land \psi
  \Proves \varphi$ आणि $\varphi, \varphi \lif \psi \Proves \psi$ (विध्यात्मक अनुमान). ही
  तथ्ये संपूर्णता प्रमेयाच्या सिद्धतेसाठी आवश्यक आहेत.
\end{explain}

\begin{prop}\label{fol:seq:ppr:prop:provability-land}
  \begin{enumerate}
  \item \label{fol:seq:ppr:prop:provability-land-left} $\varphi \land \psi \Proves
    \varphi$ आणि $\varphi \land \psi \Proves \psi$ ही दोन्ही तथ्ये सत्य आहेत.
  \item \label{fol:seq:ppr:prop:provability-land-right} $\varphi, \psi \Proves \varphi \land
    \psi$.
  \end{enumerate}
\end{prop}

\begin{proof}
  \begin{enumerate}
  \item $\varphi \land \psi \Sequent \varphi$ आणि $\varphi \land \psi \Sequent
    \psi$ या दोन्ही क्रमवर्ती निष्पन्न करता येण्याजोगे आहेत:
    \begin{prooftree}
      \Axiom$\varphi \fCenter \varphi$
      \RightLabel{\LeftR{\land}}
      \UnaryInf$\varphi \land \psi \fCenter \varphi$
      \DisplayProof\qquad\bottomAlignProof
      \Axiom$\psi \fCenter \psi$
      \RightLabel{\LeftR{\land}}
      \UnaryInf$\varphi \land \psi \fCenter \psi$
    \end{prooftree}
    \item $\varphi, \psi \Sequent \varphi \land \psi$ या क्रमवर्तीची निष्पत्ती येथे आहे:
    \begin{prooftree}
      \Axiom$\varphi \fCenter \varphi$
      \Axiom$\psi \fCenter \psi$
      \RightLabel{\RightR{\land}}
      \BinaryInf$\varphi, \psi \fCenter \varphi \land \psi$
    \end{prooftree}
  \end{enumerate}
\end{proof}

\begin{prop}\label{fol:seq:ppr:prop:provability-lor}
  \begin{enumerate}
  \item $\varphi \lor \psi, \lnot \varphi, \lnot \psi$ विसंगत आहे.
  \item $\varphi \Proves \varphi \lor \psi$ आणि $\psi \Proves \varphi \lor \psi$ ही दोन्ही तथ्ये आहेत.
  \end{enumerate}
\end{prop}

\begin{proof}
  \begin{enumerate}
  \item $\varphi \lor \psi, \lnot \varphi,
    \lnot \psi \Sequent$ या क्रमवर्तीची निष्पत्ती देऊ:
    \begin{prooftree}
      \Axiom$\varphi \fCenter \varphi$
      \RightLabel{\LeftR{\lnot}}
      \UnaryInf$\lnot \varphi, \varphi \fCenter$
      \doubleLine
      \UnaryInf$\varphi, \lnot \varphi, \lnot \psi \fCenter$
      \Axiom$\psi \fCenter \psi$
      \RightLabel{\LeftR{\lnot}}
      \UnaryInf$\lnot \psi, \psi \fCenter$
      \doubleLine
      \UnaryInf$\psi, \lnot \varphi, \lnot \psi \fCenter$
      \RightLabel{\LeftR{\lor}}
      \BinaryInf$ \varphi \lor \psi, \lnot \varphi, \lnot \psi \fCenter $
    \end{prooftree}
    (दुहेरी अनुमान-रेषा दुर्बलीकरण, आकुंचन आणि अदलाबदल यांच्या अनेक अनुमानांना
    दर्शवतात, हे लक्षात ठेवा.)
  \item   $\varphi \Sequent \varphi \lor \psi$ आणि $\psi \Sequent \varphi
    \lor \psi$ या दोन्ही क्रमवर्तींना निष्पत्त्या आहेत:
    \begin{prooftree}
      \Axiom$\varphi \fCenter \varphi$
      \RightLabel{\RightR{\lor}}
      \UnaryInf$\varphi \fCenter \varphi \lor \psi$
      \DisplayProof\qquad\bottomAlignProof
      \Axiom$\psi \fCenter \psi$
      \RightLabel{\RightR{\lor}}
      \UnaryInf$\psi \fCenter \varphi \lor \psi$
    \end{prooftree}
  \end{enumerate}
\end{proof}

\begin{prop}\label{fol:seq:ppr:prop:provability-lif}
  \begin{enumerate}
  \item \label{fol:seq:ppr:prop:provability-lif-left}  $\varphi, \varphi \lif \psi \Proves \psi$.
  \item \label{fol:seq:ppr:prop:provability-lif-right}
    $\lnot \varphi \Proves \varphi \lif \psi$ आणि $\psi \Proves \varphi \lif \psi$ ही दोन्ही तथ्ये आहेत.
  \end{enumerate}
\end{prop}

\begin{proof}
  \begin{enumerate}
    \item $\varphi \lif \psi, \varphi \Sequent \psi$ ही क्रमवर्ती निष्पन्न करता येण्याजोगे आहे:
      \begin{prooftree}
        \Axiom$\varphi \fCenter \varphi$
        \Axiom$\psi \fCenter \psi$
        \RightLabel{\LeftR{\lif}}
        \BinaryInf$\varphi \lif \psi, \varphi  \fCenter \psi$
      \end{prooftree}
    \item $\lnot \varphi \Sequent \varphi \lif \psi$ आणि $\psi
      \Sequent \varphi \lif \psi$ या दोन्ही क्रमवर्ती निष्पन्न करता येण्याजोगे आहेत:
      \begin{prooftree}
        \Axiom$\varphi \fCenter \varphi$
        \RightLabel{\LeftR{\lnot}}
        \UnaryInf$\lnot \varphi, \varphi \fCenter$
        \RightLabel{\LeftR{\Exchange}}
        \UnaryInf$\varphi, \lnot \varphi \fCenter$
        \RightLabel{\RightR{\Weakening}}
        \UnaryInf$\varphi, \lnot \varphi \fCenter \psi$
        \RightLabel{\RightR{\lif}}
        \UnaryInf$\lnot \varphi \fCenter \varphi \lif \psi$
        \DisplayProof\qquad\bottomAlignProof
        \Axiom$\psi \fCenter \psi$
        \RightLabel{\LeftR{\Weakening}}
        \UnaryInf$\varphi, \psi \fCenter \psi$
        \RightLabel{\RightR{\lif}}
        \UnaryInf$\psi \fCenter \varphi \lif \psi$
      \end{prooftree}
  \end{enumerate}
\end{proof}








\section{निष्पन्नता आणि संख्यापक}\label{fol:seq:qpr:sec}

\begin{explain}
  संपूर्णता प्रमेयासाठी या विभागात स्थापित केलेली~$\Proves$ संबंधी तथ्ये
  क्रमवर्ती कलनातील नियमांनी निष्पन्न होतात, हेही आवश्यक आहे.
\end{explain}

\begin{thm}
\label{fol:seq:qpr:thm:strong-generalization} जर $c$ हे $\Gamma$ किंवा $\varphi(x)$ मध्ये
न येणारे स्थिर असेल आणि $\Gamma \Proves \varphi(c)$, तर $\Gamma
\Proves \lforall[x][\varphi(x)]$.
\end{thm}

\begin{proof}
$\pi_0$ ही $\Gamma_0 \Sequent \varphi(c)$ या क्रमवर्तीची $\Log{LK}$-निष्पत्ती
असलेली एखादी सांत $\Gamma_0 \subseteq \Gamma$ घ्या.  $\RightR{\lforall}$
अनुमान जोडल्यावर आपल्याला $\Gamma_0 \Sequent
\lforall[x][\varphi(x)]$ ची निष्पत्ती मिळते, कारण $c$ हे $\Gamma$ किंवा
$\varphi(x)$ मध्ये येत नाही आणि म्हणून आयगेन चराची अट पूर्ण होते.
\end{proof}

\begin{prop}
\label{fol:seq:qpr:prop:provability-quantifiers}
\begin{enumerate}
\item $\varphi(t) \Proves \lexists[x][\varphi(x)]$.
\item $\lforall[x][\varphi(x)] \Proves \varphi(t)$.
\end{enumerate}
\end{prop}

\begin{proof}
\begin{enumerate}
\item $\varphi(t) \Sequent \lexists[x][\varphi(x)]$ ही क्रमवर्ती
  निष्पन्न करता येण्याजोगे आहे:
  \begin{prooftree}
    \Axiom$\varphi(t) \fCenter \varphi(t)$
    \RightLabel{\RightR{\lexists}}
    \UnaryInf$\varphi(t) \fCenter \lexists[x][\varphi(x)]$
    \end{prooftree}
\item $\lforall[x][\varphi(x)] \Sequent \varphi(t)$ ही क्रमवर्ती
  निष्पन्न करता येण्याजोगे आहे:
  \begin{prooftree}
    \Axiom$\varphi(t) \fCenter \varphi(t)$
    \RightLabel{\LeftR{\lforall}}
    \UnaryInf$\lforall[x][\varphi(x)] \fCenter \varphi(t)$
    \end{prooftree}
\end{enumerate}
\end{proof}









\section{निर्दोषता}\label{fol:seq:sou:sec}

\begin{explain}
क्रमवर्ती कलनासारखी निष्पत्ती-पद्धत \emph{निर्दोष} असते, जर ती
प्रत्यक्षात धारण न करणाऱ्या गोष्टी निष्पन्न करू शकत नसेल. त्यामुळे
निर्दोषता हा निष्पत्ती-पद्धतींसाठी हमी दिलेल्या सुरक्षिततेचा एक प्रकार
आहे. कोणता सिद्धता-उपपत्तीय गुणधर्म विचारात आहे यावर अवलंबून, उदाहरणार्थ,
आपल्याला पुढील गोष्टी हव्या असतात:
\begin{enumerate}
\item प्रत्येक निष्पन्न करता येण्याजोगे~$\varphi$ हे वैध असते;
\item एखादे वाक्य इतर काही वाक्यांपासून निष्पन्न करता येण्याजोगे असेल, तर ते
  त्यांचे तार्किक परिणामही असते;
\item वाक्यांचा संच विसंगत असेल, तर तो अपूर्ततायोग्य असतो.
\end{enumerate}
हे निष्पत्ती-पद्धतीचे महत्त्वाचे गुणधर्म आहेत. यांपैकी एखादा गुणधर्म
खरा नसेल, तर निष्पत्ती-पद्धतीत दोष आहे—ती खूप जास्त निष्पत्ती
निष्पन्न करेल. म्हणून निष्पत्ती-पद्धतीची निर्दोषता सिद्ध करणे अत्यंत
महत्त्वाचे आहे.

वरील सर्व सिद्धता-उपपत्तीय गुणधर्म क्रमवर्ती कलनातील ठरावीक क्रमवर्तींच्या
निष्पन्नता द्वारे परिभाषित केलेले असल्याने, (1)--(3) सिद्ध करण्यासाठी
निष्पन्न करता येण्याजोगे क्रमवर्तींच्या चिन्हार्थविषयक गुणधर्मांबद्दल काहीतरी सिद्ध करावे
लागते. प्रथम क्रमवर्ती \emph{वैध} असणे म्हणजे काय ते परिभाषित करू आणि मग
प्रत्येक निष्पन्न करता येण्याजोगे क्रमवर्ती वैध असते हे दाखवू. त्यानंतर (1)--(3) या
परिणामाच्या उपपरिणामांप्रमाणे मिळतील.
\end{explain}

\begin{defn}
संरचना~$\Struct
  M$ क्रमवर्ती
$\Gamma \Sequent \Delta$ \emph{पूर्ण करते} तेव्हा आणि केवळ तेव्हाच, जेव्हा
$\Sat/{M}{\varphi}$ हे एखाद्या $\varphi \in \Gamma$ साठी किंवा
$\Sat{M}{\varphi}$ हे एखाद्या $\varphi \in \Delta$ साठी असते.

क्रमवर्ती \emph{वैध} असते तेव्हा आणि केवळ तेव्हाच, जेव्हा प्रत्येक
संरचना~$\Struct
  M$ ते पूर्ण करते.
\end{defn}

\begin{thm}[निर्दोषता]
\label{fol:seq:sou:thm:sequent-soundness} जर $\Log{LK}$ निष्पन्न $\Theta
\Sequent \Xi$, तर $\Theta \Sequent \Xi$ वैध आहे.
\end{thm}

\begin{proof}
  $\Theta \Sequent \Xi$ ची निष्पत्ती~$\pi$ असू द्या. $\pi$ मधील
अनुमानांची संख्या~$n$ यावर विगमनाने पुढे जाऊ.

अनुमानांची संख्या~$0$ असेल, तर $\pi$ मध्ये फक्त आरंभीची क्रमवर्ती असते.
प्रत्येक आरंभीची क्रमवर्ती $\varphi \Sequent \varphi$ स्पष्टपणे वैध असते, कारण प्रत्येक
$\Struct M$ साठी एकतर
  $\Sat/{M}{\varphi}$ किंवा
$\Sat{M}{\varphi}$.

  अनुमानांची संख्या~0 पेक्षा जास्त असेल, तर सर्वांत खालच्या अनुमानाच्या
प्रकारानुसार प्रकरणे वेगळी करू. विगमन गृहीतकानुसार त्या अनुमानाची आधारविधाने
वैध आहेत असे गृहीत धरू, कारण कोणत्याही आधारविधानाच्या निष्पत्तीमधील
अनुमानांची संख्या~$n$ पेक्षा कमी असते.

प्रथम, एकच आधारविधान असलेली संभाव्य अनुमाने पाहू.
\begin{enumerate}
\item शेवटचे अनुमान दुर्बलीकरण आहे. मग $\Theta \Sequent \Xi$ हे एकतर
  $\varphi, \Gamma \Sequent \Delta$ (शेवटचे अनुमान
  \LeftR{\Weakening} असेल) किंवा $\Gamma \Sequent \Delta, \varphi$ (ते
  \RightR{\Weakening} असेल), आणि निष्पत्ती पुढीलपैकी एका रूपात संपते:
  \begin{prooftree}
    \AxiomC{}
    \Deduce$\Gamma \fCenter \Delta$
    \RightLabel{\LeftR{\Weakening}}
    \UnaryInf$\varphi, \Gamma \fCenter \Delta$
    \DisplayProof\qquad\bottomAlignProof
    \AxiomC{}
    \Deduce$\Gamma \fCenter \Delta$
    \RightLabel{\RightR{\Weakening}}
    \UnaryInf$\Gamma \fCenter \Delta, \varphi$
  \end{prooftree}
  विगमन गृहीतकानुसार $\Gamma \Sequent \Delta$ वैध आहे; म्हणजे प्रत्येक
  संरचना~$\Struct{M}$,
  एकतर एखादे $\chi \in \Gamma$ असे आहे की
  $\Sat/{M}{\chi}$, किंवा एखादे $\chi \in
  \Delta$ असे आहे की $\Sat{M}{\chi}$.

  जर $\Sat/{M}{\chi}$ हे एखाद्या $\chi \in \Gamma$ साठी
  असेल, तर $\Theta = \varphi, \Gamma$ असल्याने $\chi \in \Theta$ देखील आहे आणि
  म्हणून $\Sat/{M}{\chi}$ हे एखाद्या $\chi \in \Theta$ साठी
  आहे. त्याचप्रमाणे, जर $\Sat{M}{\chi}$ हे एखाद्या
  $\chi \in \Delta$ साठी असेल, तर $\chi \in \Xi$ आणि
  $\Sat{M}{\chi}$ हे एखाद्या $\chi \in \Xi$ साठी आहे.
  म्हणून $\Theta \Sequent \Xi$ वैध आहे.
\item शेवटचे अनुमान \LeftR{\lnot} आहे. मग त्याचे आधारविधान
  $\Gamma \Sequent \Delta, \varphi$ आणि निष्कर्ष $\lnot \varphi, \Gamma \Sequent
  \Delta$ असतो; म्हणजे निष्पत्ती पुढील रूपात संपते:
  \begin{prooftree}
    \AxiomC{}
    \Deduce$\Gamma \fCenter \Delta, \varphi$
    \RightLabel{\LeftR{\lnot}}
    \UnaryInf$\lnot \varphi, \Gamma \fCenter \Delta$
  \end{prooftree}
  आणि $\Theta = \lnot \varphi, \Gamma$, तर $\Xi = \Delta$.

  विगमन गृहीतकानुसार $\Gamma \Sequent \Delta, \varphi$ वैध आहे; म्हणजे प्रत्येक
  $\Struct{M}$ साठी एकतर (a) एखाद्या
  $\chi \in \Gamma$ साठी
  $\Sat/{M}{\chi}$, किंवा (b) एखाद्या $\chi \in
  \Delta$ साठी, $\Sat{M}{\chi}$, किंवा (c)
  $\Sat{M}{\varphi}$. $\Theta
  \Sequent \Xi$ देखील वैध आहे हे दाखवायचे आहे. $\Struct{M}$
    ही संरचना घ्या. (a) असल्यास,
  एखादे $\chi \in \Gamma$ असे आहे की
  $\Sat/{M}{\chi}$, आणि $\chi \in \Theta$ मध्ये
  देखील आहे. (b) असल्यास, एखादे $\chi \in \Delta$ असे आहे की
  $\Sat{M}{\chi}$, आणि $\chi \in \Xi$ मध्ये
  देखील आहे. शेवटी, जर $\Sat{M}{\varphi}$ असेल, तर
  $\Sat/{M}{\lnot \varphi}$. कारण $\lnot \varphi \in
  \Theta$ मध्ये असल्याने, एखादे $\chi \in \Theta$ असे आहे की
  $\Sat/{M}{\chi}$. म्हणून $\Theta
  \Sequent \Xi$ वैध आहे.
\item शेवटचे अनुमान \RightR{\lnot} आहे: सरावासाठी.
\item शेवटचे अनुमान \LeftR{\land} आहे. याचे दोन प्रकार आहेत: आधारविधानाच्या
  डाव्या बाजूला असलेल्या $\varphi$ किंवा $\psi$ पासून $\varphi
  \land \psi$ डावीकडे निष्पन्न केले जाऊ शकते. पहिल्या प्रकारात $\pi$ पुढील रूपात
  संपते:
  \begin{prooftree}
    \AxiomC{}
    \Deduce$\varphi, \Gamma \fCenter \Delta$
    \RightLabel{\LeftR{\land}}
    \UnaryInf$\varphi \land \psi, \Gamma \fCenter \Delta$
  \end{prooftree}
  आणि $\Theta = \varphi \land \psi, \Gamma$, तर $\Xi = \Delta$. आता
  संरचना~$\Struct
  M$ घ्या. विगमन गृहीतकानुसार
  $\varphi, \Gamma \Sequent \Delta$ वैध असल्याने, (a)
  $\Sat/{M}{\varphi}$, (b)
  $\Sat/{M}{\chi}$ हे एखाद्या $\chi \in \Gamma$ साठी, किंवा
  (c)~$\Sat{M}{\chi}$ हे एखाद्या $\chi \in \Delta$ साठी.
  (a) प्रकरणात $\Sat/{M}{\varphi \land \psi}$, म्हणून
  $\chi \in \Theta$ (म्हणजे $\varphi \land \psi$) असे आहे की
  $\Sat/{M}{\chi}$. (b) प्रकरणात एखादे $\chi
  \in \Gamma$ असे आहे की $\Sat/{M}{\chi}$ आणि
  $\chi \in \Theta$ देखील आहे. (c) प्रकरणात एखादे $\chi \in \Delta$ असे आहे
  की $\Sat{M}{\chi}$ आणि $\Xi = \Delta$ असल्याने
  $\chi \in \Xi$ देखील आहे. म्हणून प्रत्येक प्रकरणात
  $\Struct
    M$ $\varphi \land \psi, \Gamma \Sequent
  \Delta$ पूर्ण करते. $\Struct{M}$ स्वेच्छ
  असल्याने $\Gamma \Sequent \Delta$ वैध आहे. $\varphi \land \psi$ हे $\psi$ पासून
  निष्पन्न होण्याचे प्रकरणही असेच हाताळता येते; फक्त $\varphi$ ऐवजी $\psi$ घ्या.
\item शेवटचे अनुमान \RightR{\lor} आहे. याचे दोन प्रकार आहेत: आधारविधानाच्या
  उजव्या बाजूला असलेल्या $\varphi$ किंवा $\psi$ पासून $\varphi
  \lor \psi$ उजवीकडे निष्पन्न केले जाऊ शकते. पहिल्या प्रकारात $\pi$ पुढील रूपात
  संपते:
  \begin{prooftree}
    \AxiomC{}
    \Deduce$\Gamma \fCenter \Delta, \varphi$
    \RightLabel{\RightR{\lor}}
    \UnaryInf$\Gamma \fCenter \Delta, \varphi \lor \psi$
  \end{prooftree}
  येथे $\Theta = \Gamma$ आणि $\Xi = \Delta, \varphi \lor \psi$. आता
  संरचना~$\Struct{M}$ घ्या.
  $\Gamma \Sequent \Delta, \varphi$ वैध असल्याने, (a)
  $\Sat{M}{\varphi}$, (b)
  $\Sat/{M}{\chi}$ हे एखाद्या $\chi \in \Gamma$ साठी, किंवा
  (c)~$\Sat{M}{\chi}$ हे एखाद्या $\chi \in \Delta$ साठी.
  (a) प्रकरणात $\Sat{M}{\varphi \lor \psi}$. (b) प्रकरणात
  एखादे $\chi \in \Gamma$ असे आहे की
  $\Sat/{M}{\chi}$. (c) प्रकरणात एखादे $\chi
  \in \Delta$ असे आहे की $\Sat{M}{\chi}$. म्हणून
  प्रत्येक प्रकरणात $\Struct{M}$
  $\Gamma \Sequent \Delta, \varphi \lor \psi$, म्हणजेच $\Theta \Sequent \Xi$ पूर्ण
  करते. $\Struct{M}$ स्वेच्छ असल्याने
  $\Theta \Sequent \Xi$ वैध आहे. $\varphi \lor \psi$ हे $\psi$ पासून निष्पन्न होण्याचे
  प्रकरणही असेच हाताळता येते; फक्त $\varphi$ ऐवजी $\psi$ घ्या.
\item शेवटचे अनुमान \RightR{\lif} आहे. मग $\pi$ पुढील रूपात संपते:
  \begin{prooftree}
    \AxiomC{}
    \Deduce$\varphi, \Gamma \fCenter \Delta, \psi$
    \RightLabel{\RightR{\lif}}
    \UnaryInf$\Gamma \fCenter \Delta, \varphi \lif \psi$
  \end{prooftree}
  पुन्हा, विगमन गृहीतकानुसार आधारविधान वैध आहे; निष्कर्षही वैध आहे हे
  दाखवायचे आहे. $\Struct{M}$ स्वेच्छ घ्या.
  $\varphi, \Gamma \Sequent \Delta, \psi$ वैध असल्याने पुढीलपैकी किमान एक प्रकरण
  लागू होते: (a) $\Sat/{M}{\varphi}$, (b)
  $\Sat{M}{\psi}$, (c)
  $\Sat/{M}{\chi}$ हे एखाद्या~$~\chi \in \Gamma$ साठी, किंवा
  (d) $\Sat{M}{\chi}$ हे एखाद्या $\chi \in \Delta$ साठी.
  (a) आणि (b) प्रकरणांत $\Sat{M}{\varphi \lif \psi}$
  आणि म्हणून $\chi \in \Delta, \varphi \lif \psi$ असे आहे की
  $\Sat{M}{\chi}$. (c) प्रकरणात एखाद्या $\chi \in
  \Gamma$ साठी $\Sat/{M}{\chi}$. (d) प्रकरणात
  एखाद्या $\chi \in \Delta$ साठी $\Sat{M}{\chi}$. प्रत्येक
  प्रकरणात $\Struct{M}$ $\Gamma
  \Sequent \Delta, \varphi \lif \psi$ पूर्ण करते. $\Struct{M}$
  स्वेच्छ असल्याने $\Gamma \Sequent \Delta, \varphi \lif \psi$ वैध आहे.
\item शेवटचे अनुमान \LeftR{\lforall} आहे. मग सूत्र~$\varphi(x)$ आणि
  बंद पद~$t$ असे आहेत की $\pi$ पुढील रूपात संपते:
  \begin{prooftree}
    \AxiomC{}
    \Deduce$\varphi(t), \Gamma \fCenter \Delta$
    \RightLabel{\LeftR{\lforall}}
    \UnaryInf$\lforall[x][\varphi(x)], \Gamma \fCenter \Delta$
  \end{prooftree}
  निष्कर्ष $\lforall[x][\varphi(x)], \Gamma
  \Sequent \Delta$ वैध आहे हे दाखवायचे आहे. संरचना~$\Struct M$ घ्या.
  आधारविधान $\varphi(t), \Gamma \Sequent \Delta$ वैध असल्याने, (a)
  $\Sat/{M}{\varphi(t)}$, (b) $\Sat/{M}{\chi}$ हे एखाद्या $\chi \in \Gamma$ साठी, किंवा
  (c)~$\Sat{M}{\chi}$ हे एखाद्या $\chi \in \Delta$ साठी असते. (a) प्रकरणात,
  \ref{fol:syn:sem:prop:quant-terms} नुसार, जर
  $\Sat{M}{\lforall[x][\varphi(x)]}$ असेल, तर $\Sat{M}{\varphi(t)}$. म्हणून
  $\Sat/{M}{\varphi(t)}$ असल्याने $\Sat/{M}{\lforall[x][\varphi(x)]}$. (b) आणि
  (c) प्रकरणांत $\Struct{M}$ देखील $\lforall[x][\varphi(x)], \Gamma
  \Sequent \Delta$ पूर्ण करते. $\Struct M$ स्वेच्छ असल्याने,
  $\lforall[x][\varphi(x)], \Gamma \Sequent \Delta$ वैध आहे.
\item शेवटचे अनुमान \RightR{\lexists} आहे: सरावासाठी.
\item शेवटचे अनुमान \RightR{\lforall} आहे. मग
  सूत्र~$\varphi(x)$ आणि स्थिरांक~$a$ असे आहेत की $\pi$ पुढील रूपात
  संपते:
  \begin{prooftree}
    \AxiomC{}
    \Deduce$\Gamma \fCenter \Delta, \varphi(a)$
    \RightLabel{\RightR{\lforall}}
    \UnaryInf$\Gamma \fCenter \Delta, \lforall[x][\varphi(x)]$
  \end{prooftree}
  येथे आयगेन चराची अट पूर्ण आहे, म्हणजे $a$ हे $\varphi(x)$, $\Gamma$ किंवा
  $\Delta$ मध्ये येत नाही. विगमन गृहीतकानुसार शेवटच्या अनुमानाचे आधारविधान
  वैध आहे. निष्कर्षही वैध आहे हे दाखवायचे आहे; म्हणजे कोणत्याही
  संरचना~$\Struct M$ साठी (a) $\Sat{M}{\lforall[x][\varphi(x)]}$, (b)
  एखाद्या $\chi \in \Gamma$ साठी $\Sat/{M}{\chi}$, किंवा (c) एखाद्या
  $\chi \in \Delta$ साठी $\Sat{M}{\chi}$.

  $\Struct{M}$ ही स्वेच्छ संरचना घ्या. (b) किंवा (c) लागू असेल तर
  काम झाले; म्हणून दोन्ही लागू नाहीत असे गृहीत धरा: प्रत्येक $\chi \in
  \Gamma$ साठी $\Sat{M}{\chi}$ आणि प्रत्येक $\chi \in \Delta$ साठी
  $\Sat/{M}{\chi}$. आता (a), म्हणजेच $\Sat{M}{\lforall[x][\varphi(x)]}$, हे दाखवायचे
  आहे. \ref{fol:syn:ass:prop:sat-quant} नुसार सर्व चल-विन्यास~$s$ साठी
  $\Sat{M}{\varphi(x)}[s]$ दाखवणे पुरेसे आहे. म्हणून $s$ हा स्वेच्छ चल-विन्यास घ्या.
  $\Struct{M'}$ ही $\Struct{M}$ सारखीच रचना घ्या, फक्त
  $\Assign{a}{M'} = s(x)$ असेल. \ref{fol:syn:ext:cor:extensionality-sent}
  नुसार प्रत्येक $\chi \in \Gamma$ साठी $a$ हे $\Gamma$ मध्ये येत नसल्यामुळे
  $\Sat{M'}{\chi}$, आणि प्रत्येक $\chi \in \Delta$ साठी $\Sat/{M'}{\chi}$.
  पण आधारविधान वैध असल्याने $\Sat{M'}{\varphi(a)}$. \ref{fol:syn:ass:prop:sentence-sat-true}
  नुसार $\varphi(a)$ हे वाक्य असल्यामुळे $\Sat{M'}{\varphi(a)}[s]$.
  आता $\varAssign{s}{s}{x}$ आणि $s(x) = \Value{a}{M'}[s]$, कारण
  $\Struct{M'}$ याच प्रकारे परिभाषित केले आहे. म्हणून
  \ref{fol:syn:ext:prop:ext-formulas} लागू होते आणि $\Sat{M'}{\varphi(x)}[s]$
  मिळते. $a$ हे $\varphi(x)$ मध्ये येत नसल्याने
  \ref{fol:syn:ext:prop:extensionality} नुसार $\Sat{M}{\varphi(x)}[s]$.
  $s$ स्वेच्छ असल्याने सर्व चल-विन्यासांसाठी $\Sat{M}{\varphi(x)}[s]$ सिद्ध झाले.
\item शेवटचे अनुमान \LeftR{\lexists} आहे: सरावासाठी.

\end{enumerate}
आता दोन आधारविधाने असलेली संभाव्य अनुमाने पाहू.
\begin{enumerate}
\item शेवटचे अनुमान कट आहे; म्हणून $\pi$ पुढील रूपात संपते:
  \begin{prooftree}
    \AxiomC{}
    \Deduce$\Gamma \fCenter \Delta, \varphi$
    \AxiomC{}
    \Deduce$\varphi, \Pi \fCenter \Lambda$
    \RightLabel{\Cut}
    \BinaryInf$\Gamma, \Pi \fCenter \Delta, \Lambda$
  \end{prooftree}
  $\Struct{M}$ ही संरचना घ्या. विगमन गृहीतकानुसार आधारविधाने वैध असल्याने
  $\Struct{M}$ दोन्ही आधारविधाने पूर्ण करते.
  दोन प्रकरणे पाहू: (a) $\Sat/{M}{\varphi}$ आणि (b)
  $\Sat{M}{\varphi}$. (a) प्रकरणात डावे आधारविधान पूर्ण
  करण्यासाठी $\Struct{M}$ ने
  $\Gamma \Sequent \Delta$ पूर्ण केले पाहिजे; मग ते निष्कर्षही पूर्ण करते.
  (b) प्रकरणात उजवे आधारविधान पूर्ण करण्यासाठी
  $\Struct{M}$ ने $\Pi \setminus \Lambda$ पूर्ण
  केले पाहिजे. पुन्हा, $\Struct{M}$ निष्कर्ष पूर्ण करते.
\item शेवटचे अनुमान \RightR{\land} आहे. मग $\pi$ पुढील रूपात संपते:
  \begin{prooftree}
    \AxiomC{}
    \Deduce$\Gamma \fCenter \Delta, \varphi$
    \AxiomC{}
    \Deduce$\Gamma \fCenter \Delta, \psi$
    \RightLabel{\RightR{\land}}
    \BinaryInf$\Gamma \fCenter \Delta, \varphi \land \psi$
  \end{prooftree}
  संरचना~$\Struct
    M$ घ्या. $\Struct{M}$
  ने $\Gamma \Sequent \Delta$ पूर्ण केले तर काम झाले. तसे नसल्याचे गृहीत धरा.
  विगमन गृहीतकानुसार $\Gamma \fCenter
  \Delta, \varphi$ वैध असल्याने,
  $\Sat{M}{\varphi}$. त्याचप्रमाणे, कारण $\Gamma
  \Sequent \Delta, \psi$ वैध असल्याने,
  $\Sat{M}{\psi}$. म्हणून
  $\Sat{M}{\varphi \land \psi}$.
\item शेवटचे अनुमान \LeftR{\lor} आहे: सरावासाठी.
\item शेवटचे अनुमान \LeftR{\lif} आहे. मग $\pi$ पुढील रूपात संपते:
  \begin{prooftree}
    \AxiomC{}
    \Deduce$\Gamma \fCenter \Delta, \varphi$
    \AxiomC{}
    \Deduce$\psi, \Pi \fCenter \Lambda$
    \RightLabel{\LeftR{\lif}}
    \BinaryInf$\varphi \lif \psi, \Gamma, \Pi \fCenter \Delta, \Lambda$
  \end{prooftree}
  पुन्हा,
  संरचना~$\Struct{M}$
  $\Struct{M}$ ने $\Gamma, \Pi \Sequent
  \Delta, \Lambda$ पूर्ण केलेले नाही असे गृहीत धरा. आपल्याला
  $\Sat/{M}{\varphi \lif \psi}$ दाखवायचे आहे.
  $\Struct{M}$ ने $\Gamma,
  \Pi \Sequent \Delta, \Lambda$ पूर्ण केलेले नसल्याने ते $\Gamma \Sequent
  \Delta$ किंवा $\Pi \Sequent \Lambda$ यांपैकी कोणतेही पूर्ण करत नाही. $\Gamma
  \Sequent \Delta, \varphi$ वैध असल्याने $\Sat{M}{\varphi}$.
  $\psi, \Pi \Sequent \Lambda$ वैध असल्याने
  $\Sat/{M}{\psi}$. म्हणून
  $\Sat/{M}{\varphi \lif \psi}$, जे दाखवायचे होते तेच.
\end{enumerate}
\end{proof}


\begin{prob}
\ref{fol:seq:sou:thm:sequent-soundness} ची सिद्धता पूर्ण करा.
\end{prob}




\begin{cor}
\label{fol:seq:sou:cor:weak-soundness}
जर $\Proves \varphi$, तर $\varphi$ हे वैध आहे.
\end{cor}

\begin{cor}
\label{fol:seq:sou:cor:entailment-soundness}
जर $\Gamma \Proves \varphi$, तर $\Gamma \Entails \varphi$.
\end{cor}

\begin{proof}
जर $\Gamma \Proves \varphi$ असेल, तर $\Gamma_0 \subseteq \Gamma$ असा काही
मर्यादित संच आणि $\Gamma_0 \Sequent \varphi$ ची निष्पत्ती असते. \ref{fol:seq:sou:thm:sequent-soundness}
नुसार, प्रत्येक
संरचना~$\Struct{M}$
किंवा $\psi \in \Gamma_0$ असत्य करते, किंवा $\varphi$ सत्य करते. म्हणून,
जर $\Sat{M}{\Gamma}$ असेल, तर
$\Sat{M}{\varphi}$ देखील असेल.
\end{proof}

\begin{cor}
\label{fol:seq:sou:cor:consistency-soundness}
जर $\Gamma$ पूर्ण करण्यायोग्य असेल, तर तो सुसंगत आहे.
\end{cor}

\begin{proof}
प्रतिलोम विधान सिद्ध करू. $\Gamma$ सुसंगत नाही असे गृहीत धरा. मग
$\Gamma_0 \subseteq \Gamma$ असा काही मर्यादित संच आणि
$\Gamma_0 \Sequent \quad$ ची निष्पत्ती असते. \ref{fol:seq:sou:thm:sequent-soundness}
नुसार $\Gamma_0 \Sequent \quad$ वैध आहे. दुसऱ्या शब्दांत, प्रत्येक
संरचना~$\Struct{M}$ साठी
$\chi \in \Gamma_0$ असे आहे की
$\Sat/{M}{\chi}$; आणि $\Gamma_0 \subseteq \Gamma$
असल्याने ते $\chi$ $\Gamma$ मध्येही आहे. त्यामुळे कोणतेही
$\Struct{M}$ $\Gamma$ पूर्ण करत नाही, आणि
$\Gamma$ पूर्ण करण्यायोग्य नाही.
\end{proof}









\section{निष्पत्ती सह एकरूपता}\label{fol:seq:ide:sec}

निष्पत्त्या आणि एकरूपता साठी अतिरिक्त आरंभीच्या क्रमवर्ती
आणि अनुमाननियमांची आवश्यकता असते.

\begin{defn}[$\eq$ साठी आरंभीच्या क्रमवर्ती]
जर $t$ हे बंद पद असेल, तर ${} \Sequent \eq[t][t]$ ही आरंभीची क्रमवर्ती असते.
\end{defn}

$\eq$ साठीचे नियम पुढीलप्रमाणे आहेत ($t_1$ आणि $t_2$ ही बंद पदे आहेत):

\begin{defish}
\Axiom$ \eq[t_1][t_2], \Gamma \fCenter \Delta, \varphi(t_1) $
\RightLabel{$\eq$}
\UnaryInf$\eq[t_1][t_2], \Gamma \fCenter \Delta, \varphi(t_2)$
\DisplayProof
\hfill
\Axiom$\eq[t_1][t_2], \Gamma \fCenter \Delta, \varphi(t_2) $
\RightLabel{$\eq$}
\UnaryInf$\eq[t_1][t_2], \Gamma  \fCenter \Delta, \varphi(t_1)$
\DisplayProof
\end{defish}

\begin{ex}
जर $s$ आणि $t$ ही बंद पदे असतील, तर $\eq[s][t], \varphi(s)
\Proves \varphi(t)$:
\begin{prooftree}
\Axiom$ \varphi(s) \fCenter \varphi(s)$
\RightLabel{\LeftR{\Weakening}}
\UnaryInf$\eq[s][t], \varphi(s)  \fCenter \varphi(s)$
\RightLabel{$\eq$}
\UnaryInf$\eq[s][t], \varphi(s)  \fCenter \varphi(t)$
\end{prooftree}
याला एकरूप वस्तूंच्या प्रतिस्थापनाचे तत्त्व किंवा लाइब्निझचा नियम
म्हणून ओळखले जाऊ शकते.

$\Log{LK}$ सिद्ध करते की $\eq$ सममित आणि संक्रमक आहे:
\begin{prooftree}
\Axiom$ \fCenter \eq[t_1][t_1] $
\RightLabel{\LeftR{\Weakening}}
\UnaryInf$ \eq[t_1][t_2] \fCenter \eq[t_1][t_1] $
\RightLabel{$\eq$}
\UnaryInf$ \eq[t_1][t_2] \fCenter \eq[t_2][t_1]$
\DisplayProof\qquad\bottomAlignProof
\Axiom$ \eq[t_1][t_2] \fCenter \eq[t_1][t_2] $
\RightLabel{\LeftR{\Weakening}}
\UnaryInf$\eq[t_2][t_3], \eq[t_1][t_2]  \fCenter \eq[t_1][t_2] $
\RightLabel{$\eq$}
\UnaryInf$\eq[t_2][t_3], \eq[t_1][t_2]  \fCenter \eq[t_1][t_3]$
\RightLabel{\LeftR{\Exchange}}
\UnaryInf$\eq[t_1][t_2], \eq[t_2][t_3]  \fCenter \eq[t_1][t_3]$
\end{prooftree}
डावीकडील निष्पत्तीमध्ये सूत्र~$\eq[x][t_1]$ हे आपले
$\varphi(x)$ आहे. उजवीकडे $\varphi(x)$ हे~$\eq[t_1][x]$ असे घेतो.
\end{ex}

\begin{prob}
पुढील क्रमवर्तींच्या निष्पत्त्या द्या:
\begin{enumerate}
\item $\Sequent \lforall[x][\lforall[y][((x = y \land \varphi(x)) \lif \varphi(y))]]$
\item $\lexists[x][\varphi(x)] \land \lforall[y][\lforall[z][((\varphi(y) \land
    \varphi(z)) \lif y = z)]] \Sequent
\lexists[x][(\varphi(x) \land \lforall[y][(\varphi(y) \lif y = x)])]$
\end{enumerate}
\end{prob}









\section{एकरूपता सह निर्दोषता}\label{fol:seq:sid:sec}

\begin{prop}
$\Log{LK}$ मधील आरंभीच्या क्रमवर्ती आणि एकरूपतेच्या नियमांसहची पद्धत निर्दोष आहे.
\end{prop}

\begin{proof}
${} \Sequent \eq[t][t]$ या रूपातील आरंभीच्या क्रमवर्ती वैध आहेत, कारण
प्रत्येक संरचना~$\Struct M$ साठी $\Sat{M}{\eq[t][t]}$ असते. (लक्षात घ्या की
$t$ हे बंद पद आहे असे आपण गृहीत धरतो, म्हणजे त्यात कोणतेही चल नाही;
म्हणून चल-विन्यास अप्रासंगिक असतात.)

निष्पत्तीमधील शेवटचे अनुमान $=$ आहे असे गृहीत धरा. मग आधारविधान
$\eq[t_1][t_2], \Gamma \Sequent \Delta, \varphi(t_1)$ आणि निष्कर्ष
$\eq[t_1][t_2], \Gamma \Sequent \Delta, \varphi(t_2)$ असतो. संरचना~$\Struct M$
घ्या. निष्कर्ष वैध आहे हे दाखवायचे आहे; म्हणजे $\Sat{M}{\eq[t_1][t_2]}$ आणि
$\Sat{M}{\Gamma}$ असल्यास, एकतर $\Sat{M}{\chi}$ हे एखाद्या $\chi \in \Delta$ साठी
असते किंवा $\Sat{M}{\varphi(t_2)}$ असते.

विगमन गृहीतकानुसार आधारविधान वैध आहे. म्हणजे $\Sat{M}{\eq[t_1][t_2]}$ आणि
$\Sat{M}{\Gamma}$ असल्यास, एकतर (a) एखाद्या $\chi \in \Delta$ साठी
$\Sat{M}{\chi}$ किंवा (b) $\Sat{M}{\varphi(t_1)}$. (a) प्रकरणात काम झाले.
(b) प्रकरणाचा विचार करू. $s$ हा चल-विन्यास घ्या, ज्यासाठी
$s(x) = \Value{t_1}{M}$. \ref{fol:syn:ass:prop:sentence-sat-true} नुसार,
$\Sat{M}{\varphi(t_1)}[s]$. $\varAssign{s}{s}{x}$ असल्याने,
\ref{fol:syn:ext:prop:ext-formulas} नुसार $\Sat{M}{\varphi(x)}[s]$. कारण
$\Sat{M}{\eq[t_1][t_2]}$, आपल्याला $\Value{t_1}{M} = \Value{t_2}{M}$,
आणि म्हणून $s(x) = \Value{t_2}{M}$ मिळते. \ref{fol:syn:ext:prop:ext-formulas}
पुन्हा लागू केल्याने $\Sat{M}{\varphi(t_2)}[s]$ देखील मिळते.
\ref{fol:syn:ass:prop:sentence-sat-true} नुसार,
$\Sat{M}{\varphi(t_2)}$.
\end{proof}



\chapter{नैसर्गिक निगमन}\label{fol:ntd::chap}
\begin{editorial}
हे प्रकरण Gentzen/Prawitz शैलीतील नैसर्गिक निगमन पद्धत मांडते.

नैसर्गिक निगमनाशी सिद्धता-पद्धती म्हणून संबंधित सामग्री समाविष्ट किंवा
वगळण्यासाठी ``prfND'' हा टॅग वापरा.
\end{editorial}






\section{नियम आणि निष्पत्ती}\label{fol:ntd:rul:sec}

\begin{explain}
नैसर्गिक निगमन पद्धतींचा उद्देश गणितीय सिद्धतेत वापरल्या जाणाऱ्या अनौपचारिक
तर्कपद्धतीशी जवळचे साम्य राखणे हा आहे (म्हणूनच ती काहीशी ``नैसर्गिक'' आहे).
नैसर्गिक निगमनातील सिद्धता गृहीतकांपासून सुरू होतात.
त्यानंतर अनुमाननियम लागू केले जातात. \Intro{\lnot}, \Intro{\lif},
\Elim{\lor} आणि \Elim{\lexists} अनुमाननियमांद्वारे
गृहीतके ``मुक्त'' केली जातात; स्पष्टतेसाठी मुक्त केलेल्या
गृहीतकाचा लेबल त्या अनुमानाजवळ ठेवला जातो.
\end{explain}

\begin{defn}[गृहीतक]
\emph{गृहीतक} म्हणजे कोणत्याही शाखेच्या सर्वांत वरच्या स्थानावर असलेले
कोणतेही वाक्य होय.
\end{defn}

नैसर्गिक निगमनातील निष्पत्त्या ही वाक्यांची विशिष्ट वृक्षरचना असते;
त्यातील सर्वांत वरच्या वाक्ये गृहीतके असतात. एखादे वाक्य एक,
दोन किंवा तीन इतर क्रमवर्तींच्या खाली असल्यास, ते अनुमाननियमाने योग्य
प्रकारे आलेले असले पाहिजे. अनुमानाच्या शीर्षस्थानी असलेल्या वाक्यांना
\emph{आधारविधाने} म्हणतात आणि खालील वाक्याला अनुमानाचा
\emph{निष्कर्ष} म्हणतात. प्रत्येक संकारक साठी नियम जोड्यांमध्ये येतात:
एक प्रवेशन नियम आणि एक विलोपन नियम. ते निष्कर्षात संकारक आणतात
किंवा नियमाच्या आधारविधानातून संकारक काढतात. काही नियम विशिष्ट
प्रकारचे गृहीतक \emph{मुक्त} करण्याची परवानगी देतात. कोणते
गृहीतक कोणत्या अनुमानाने मुक्त केले हे दर्शवण्यासाठी गृहीतक आणि
अनुमान दोन्हींना लेबल देतो. हे गृहीतक ``$\Discharge{\varphi}{n}$.` असे लिहून दर्शवले जाते.


सर्व संकारक $\land$, $\lor$, $\lif$, $\lnot$ आणि $\lfalse$ यांच्यासाठीचे नियम विचारात घेणे प्रचलित आहे, जरी त्यांपैकी काही परिभाषित केलेले असले तरी.











\section{विधानीय नियम}\label{fol:ntd:prl:sec}


\subsection{$\land$ साठीचे नियम}

\begin{defish}
\AxiomC{$\varphi$}
\AxiomC{$\psi$}
\RightLabel{\Intro{\land}}
\BinaryInfC{$\varphi \land \psi$}
\DisplayProof
\hfill
\begin{tabular}{r}
\AxiomC{$\varphi \land \psi$}
\RightLabel{\Elim{\land}}
\UnaryInfC{$\varphi$}
\DisplayProof
\\[3ex]
\AxiomC{$\varphi \land \psi$}
\RightLabel{\Elim{\land}}
\UnaryInfC{$\psi$}
\DisplayProof
\end{tabular}
\end{defish}
\subsection{$\lor$ साठीचे नियम}
\begin{defish}
\begin{tabular}{r}
\AxiomC{$\varphi$}
\RightLabel{\Intro{\lor}}
\UnaryInfC{$\varphi \lor \psi$}
\DisplayProof
\\[3ex]
\AxiomC{$\psi$}
\RightLabel{\Intro{\lor}}
\UnaryInfC{$\varphi \lor \psi$}
\DisplayProof
\end{tabular}
\hfill
\AxiomC{$\varphi \lor \psi$}
\AxiomC{$\Discharge{\varphi}{n}$}
\DeduceC{$\chi$}
\AxiomC{$\Discharge{\psi}{n}$}
\DeduceC{$\chi$}
\DischargeRule{\Elim{\lor}}{n}
\TrinaryInfC{$\chi$}
\DisplayProof
\end{defish}
\subsection{$\lif$ साठीचे नियम}
\begin{defish}
\AxiomC{$\Discharge{\varphi}{n}$}
\DeduceC{$\psi$}
\DischargeRule{\Intro{\lif}}{n}
\UnaryInfC{$\varphi \lif \psi$}
\DisplayProof
\hfill
\AxiomC{$\varphi \lif \psi$}
\AxiomC{$\varphi$}
\RightLabel{\Elim{\lif}}
\BinaryInfC{$\psi$}
\DisplayProof
\end{defish}
\subsection{$\lnot$ साठीचे नियम}
\begin{defish}
\AxiomC{$\Discharge{\varphi}{n}$}
\noLine
\DeduceC{$\lfalse$}
\DischargeRule{\Intro{\lnot}}{n}
\UnaryInfC{$\lnot \varphi$}
\DisplayProof
\hfill
\AxiomC{$\lnot \varphi$}
\AxiomC{$\varphi$}
\RightLabel{\Elim{\lnot}}
\BinaryInfC{$\lfalse$}
\DisplayProof
\end{defish}
\subsection{$\lfalse$ साठीचे नियम}
\begin{defish}
\AxiomC{$\lfalse$}
\RightLabel{\FalseInt}
\UnaryInfC{$\varphi$}
\DisplayProof
\hfill
\AxiomC{$\Discharge{\lnot \varphi}{n}$}
\DeduceC{$\lfalse$}
\DischargeRule{\FalseCl}{n}
\UnaryInfC{$\varphi$}
\DisplayProof
\end{defish}
लक्षात घ्या की $\Intro{\lnot}$ आणि $\FalseCl$ हे अगदी सारखे आहेत: फरक
असा की $\Intro{\lnot}$ हा नकारयुक्त वाक्य~$\lnot \varphi$ निष्पन्न
करतो, तर $\FalseCl$ हा सकारात्मक वाक्य~$\varphi$ निष्पन्न करतो.
जेव्हा एखादा नियम एखादे गृहीतक मुक्त करता येते असे दर्शवतो, तेव्हा आपण ती
परवानगी मानतो, आवश्यकता नाही. उदा., $\Intro{\lif}$ नियमात आधारविधान~$\psi$
च्या निष्पत्तीमधील $\varphi$ या रूपाची कितीही गृहीतके—शून्यसुद्धा—आपण
मुक्त करू शकतो.
\section{संख्यापकांचे नियम}\label{fol:ntd:qrl:sec}
\subsection{$\lforall$ साठीचे नियम}
\begin{defish}
\AxiomC{$\varphi(a)$}
\RightLabel{\Intro{\lforall}}
\UnaryInfC{$\lforall[x][\Atom{\varphi}{x}]$}
\DisplayProof
\hfill
\AxiomC{$\lforall[x][\Atom{\varphi}{x}]$}
\RightLabel{\Elim{\lforall}}
\UnaryInfC{$\varphi(t)$}
\DisplayProof
\end{defish}
$\lforall$ च्या नियमांमध्ये, $t$ हे बंद पद आहे (म्हणजे, कोणतेही चर
नसलेले पद), आणि $a$~हा असा स्थिरांक आहे की तो
निष्कर्ष~$\lforall[x][\varphi(x)]$ मध्ये किंवा आधारविधान~$\varphi(a)$ ने समाप्त
होणाऱ्या निष्पत्तीमध्ये मुक्त न केलेली असलेल्या कोणत्याही
गृहीतकात येत नाही. आपण $a$ ला \Intro{\lforall} अनुमानाचा
\emph{आयगेन चर} म्हणतो.\footnote{वरील नियमात $a$ हा अचर असला तरी आपण
``आयगेन चर'' ही संज्ञा वापरतो. यामागे ऐतिहासिक कारणे आहेत.}
\subsection{$\lexists$ साठीचे नियम}
\begin{defish}
\AxiomC{$\Atom{\varphi}{t}$}
\RightLabel{\Intro{\lexists}}
\UnaryInfC{$\lexists[x][\Atom{\varphi}{x}]$}
\DisplayProof
\hfill
\AxiomC{$\lexists[x][\Atom{\varphi}{x}]$}
\AxiomC{[$\Atom{\varphi}{a}$]$^n$}
\DeduceC{$\chi$}
\DischargeRule{\Elim{\lexists}}{n}
\BinaryInfC{$\chi$}
\DisplayProof
\end{defish}
पुन्हा, $t$ हे बंद पद आहे आणि $a$ हा असा स्थिरांक आहे की तो
आधारविधान $\lexists[x][\varphi(x)]$ मध्ये, निष्कर्ष~$\chi$ मध्ये किंवा दोन
आधारविधानांनी समाप्त होणाऱ्या निष्पत्त्या मधील कोणत्याही
मुक्त न केलेली गृहीतकात ($\varphi(a)$ ही गृहीतके वगळता) येत नाही. आपण
$a$ ला \Elim{\lexists} अनुमानाचा \emph{आयगेन चर} म्हणतो.
\Intro{\lforall} आणि \Elim{\lexists} अनुमानांसाठी आयगेन चरावर लागू होणारे
निर्बंध वर स्वतंत्रपणे सांगितले आहेत. त्यांत संबंधित आधारविधान व निष्कर्ष,
तसेच आधारविधानांकडे नेणाऱ्या निष्पत्त्या मधील मुक्त न केलेली गृहीतके,
यांबाबत वेगवेगळे निर्बंध आहेत. या निर्बंधांना \emph{आयगेन चराची अट} म्हणतात.
\begin{explain}
ही प्रथा आठवा: $\varphi$ हे असे सूत्र असेल की त्यात
चल~$x$ मुक्त असेल, तर हे आपण~$\varphi(x)$ असे लिहून दर्शवतो.
त्याच संदर्भात, मग $\varphi(t)$ हा~$\Subst{\varphi}{t}{x}$ चा संक्षेप असतो.
त्यामुळे $\Intro\lexists$ नियम पुढीलप्रमाणेही लिहिता येईल:
\begin{prooftree}
\AxiomC{$\Subst{\varphi}{t}{x}$}
\RightLabel{\Intro{\lexists}}
\UnaryInfC{$\lexists[x][\varphi]$}
\end{prooftree}
लक्षात घ्या की~$\varphi$ मध्ये $t$ आधीपासून येऊ शकते; उदा., $\varphi$ हे
$\Atom{\Obj P}{t,x}$ असू शकते. म्हणून,~$\Atom{\Obj P}{t,t}$ पासून
$\lexists[x][\Atom{\Obj P}{t,x}]$ निष्पन्न करणे हे~$\Intro\lexists$
चे योग्य उपयोजन आहे---आधारविधानात~$t$ जेथे येते त्यांपैकी एक किंवा
अधिक ठिकाणी त्याऐवजी बद्ध चल~$x$ घालता येतो; सर्वच ठिकाणी
असा ``बदल'' करणे आवश्यक नाही. परंतु $\Intro\lforall$ आणि
$\Elim\lexists$ मधील आयगेन चराच्या अटीनुसार स्थिरांक~$a$ हा~$\varphi$
मध्ये येता कामा नये. म्हणून, $\Intro\lforall$ वापरून
$\Atom{\Obj P}{a,a}$ पासून $\lforall[x][\Atom{\Obj P}{a,x}]$ योग्य
रीतीने निष्पन्न करता येत नाही.
\end{explain}
\begin{explain}
\Intro{\lexists} आणि \Elim{\lforall} मध्ये पद~$t$ बंद
असण्याव्यतिरिक्त त्यावर आणखी कोणतेही निर्बंध नाहीत; म्हणून इतर कोणत्याही
अटीची काळजी करावी लागत नाही. याउलट, \Elim{\lexists} आणि
\Intro{\lforall} नियमांमध्ये आयगेन चराच्या अटीनुसार स्थिरांक~$a$ हा
निष्कर्षात किंवा कोणत्याही मुक्त न केलेली गृहीतकात कोठेही येता कामा नये.
ही पद्धत निर्दोष आहे, म्हणजे ती केवळ अशा मुक्त न केलेली गृहीतकांपासूनच
वाक्ये निष्पन्न करते की ज्यांच्यापासून ती तार्किकरीत्या निष्पन्न
होतात, याची खात्री करण्यासाठी ही अट आवश्यक आहे. ही अट नसती, तर पुढील
निष्पत्तीला परवानगी मिळाली असती:
\begin{prooftree}
  \AxiomC{$\lexists[x][\varphi(x)]$}
  \AxiomC{$\Discharge{\varphi(a)}{1}$}
  \RightLabel{*\Intro{\lforall}}
  \UnaryInfC{$\lforall[x][\varphi(x)]$}
  \RightLabel{\Elim{\lexists}}
  \BinaryInfC{$\lforall[x][\varphi(x)]$}
\end{prooftree}
परंतु, $\lexists[x][\varphi(x)] \Entails/ \lforall[x][\varphi(x)]$.
संख्यापकांच्या विलोपन नियमांमध्ये चलांच्या जागी केवळ बंद पदे
ठेवण्याची परवानगी असल्यामुळे, वाक्यांच्या संचापासून निष्पन्न करता
येणारे कोणतेही सूत्र हे स्वतः वाक्य असते, हे निष्पन्न होते.
\end{explain}
\section{निष्पत्ती}\label{fol:ntd:der:sec}
\begin{explain}
गृहीतक म्हणजे काय हे आपण सांगितले आहे आणि निगमन नियम दिले आहेत.
नैसर्गिक निगमनातील निष्पत्त्या यांपासून आगमनाने निर्माण होतात:
प्रत्येक निष्पत्ती ही एकतर स्वतंत्र गृहीतक असते किंवा एक, दोन वा तीन
निष्पत्त्या नंतर केलेले योग्य अनुमान असते.
\end{explain}
\begin{defn}[निष्पत्ती]
\emph{निष्पत्ती} म्हणजे गृहीतके~$\Gamma$
यांपासून वाक्य~$\varphi$ ची पुढील अटी पूर्ण करणारी वाक्यांची
सांत वृक्षरचना होय:
\begin{enumerate}
\item वृक्षातील सर्वांत वरची वाक्ये ही एकतर $\Gamma$ मध्ये असतात
  किंवा वृक्षातील एखाद्या अनुमानाने मुक्त केलेली असतात.
\item वृक्षातील सर्वांत खालचे वाक्य हे~$\varphi$ असते.
\item वृक्षाच्या तळाशी असलेले वाक्य~$\varphi$ वगळता, वृक्षातील प्रत्येक
  वाक्य हे एखाद्या निगमन नियमाच्या योग्य उपयोजनाचे आधारविधान असते
  आणि त्या उपयोजनाचा निष्कर्ष वृक्षात त्या वाक्याच्या थेट खाली येतो.
\end{enumerate}
मग $\varphi$ हा त्या निष्पत्तीचा \emph{निष्कर्ष} आणि $\Gamma$ ही तिची
मुक्त न केलेली गृहीतके आहेत, असे आपण म्हणतो.
जर~$\Gamma$ या गृहीतकांपासून $\varphi$ ची निष्पत्ती अस्तित्वात असेल, तर
$\varphi$ हे~$\Gamma$ पासून \emph{निष्पन्न करता येण्याजोगे} आहे असे म्हणतो; चिन्हांत:
$\Gamma \Proves \varphi$. प्रत्येक गृहीतक मुक्त केलेले असणारी $\varphi$
ची निष्पत्ती असेल, तर आपण~$\Proves \varphi$ असे लिहितो.
\end{defn}
\begin{ex}
प्रत्येक गृहीतक स्वतःच निष्पत्ती असते. म्हणून, उदा., एकटे $\varphi$ ही
निष्पत्ती आहे आणि एकटे $\psi$ सुद्धा. यांवर, समजा,
$\Intro{\land}$ नियम लागू करून नवी निष्पत्ती मिळवता येते:
\begin{prooftree}
\AxiomC{$\varphi$}
\AxiomC{$\psi$}
\RightLabel{\Intro{\land}}
\BinaryInfC{$\varphi \land \psi$}
\end{prooftree}
हे नियम सर्वसाधारण असतात: त्यांतील $\varphi$ आणि~$\psi$ यांच्या जागी कोणतीही
वाक्ये, उदा., $\chi$ आणि~$\theta$, ठेवता येतात. मग निष्कर्ष
$\chi \land \theta$ असेल, आणि म्हणून
\begin{prooftree}
\AxiomC{$\chi$}
\AxiomC{$\theta$}
\RightLabel{\Intro{\land}}
\BinaryInfC{$\chi \land \theta$}
\end{prooftree}
ही योग्य निष्पत्ती आहे. अर्थात, गृहीतकांची अदलाबदल करून $\theta$ ला
$\varphi$ ची आणि $\chi$ ला~$\psi$ ची भूमिका देता येते. त्यामुळे,
\begin{prooftree}
\AxiomC{$\theta$}
\AxiomC{$\chi$}
\RightLabel{\Intro{\land}}
\BinaryInfC{$\theta \land \chi$}
\end{prooftree}
हीसुद्धा योग्य निष्पत्ती आहे.
आता आपण दुसरा नियम, समजा, $\Intro{\lif}$ लागू करू शकतो. त्यामुळे सोपाधिक
विधानाचा निष्कर्ष काढता येतो आणि त्या विधानाच्या पूर्वांगाशी एकरूप असलेले
कोणतेही गृहीतक मुक्त करता येते. म्हणून पुढील दोन्या निष्पत्त्या योग्य आहेत:
\begin{prooftree}
\AxiomC{$\Discharge{\chi}{1}$}
\AxiomC{$\theta$}
\RightLabel{\Intro{\land}}
\BinaryInfC{$\chi \land \theta$}
\DischargeRule{\Intro{\lif}}{1}
\UnaryInfC{$\chi \lif (\chi \land \theta)$}
\DisplayProof\bottomAlignProof
\AxiomC{$\chi$}
\AxiomC{$\Discharge{\theta}{1}$}
\RightLabel{\Intro{\land}}
\BinaryInfC{$\chi \land \theta$}
\DischargeRule{\Intro{\lif}}{1}
\UnaryInfC{$\theta \lif (\chi \land \theta)$}
\end{prooftree}
त्या अनुक्रमे $\theta \Proves \chi \lif (\chi \land \theta)$ आणि
$\chi \Proves \theta \lif (\chi \land \theta)$ हे दाखवतात.
गृहीतके मुक्त करणे ही परवानगी आहे, आवश्यकता नाही, हे आठवा: गृहीतके मुक्त
करणे बंधनकारक नसते. विशेषतः, निष्पत्तीमध्ये ती गृहीतके नसली तरीही
नियम लागू करता येतो. उदाहरणार्थ, मुक्त करण्यासाठी~$\varphi$ हे गृहीतक
नसले तरी पुढील उपयोजन अनुज्ञात आहे:
\begin{prooftree}
  \AxiomC{$\psi$}
  \DischargeRule{\Intro{\lif}}{1}
  \UnaryInfC{$\varphi \lif \psi$}
\end{prooftree}
\end{ex}
\section{निष्पत्ती: उदाहरणे}\label{fol:ntd:pro:sec}
\begin{ex}
वाक्य $(\varphi \land \psi) \lif \varphi$ ची निष्पत्ती देऊ.
हवा असलेला निष्कर्ष निष्पत्तीच्या तळाशी लिहून आपण सुरुवात करतो.
\begin{prooftree}
\AxiomC{}
\UnaryInfC{$(\varphi\land \psi) \lif \varphi$}
\end{prooftree}
पुढे, या रूपाचे वाक्य कोणत्या प्रकारच्या अनुमानामुळे मिळू शकते हे
शोधावे लागते. निष्कर्षाचा मुख्य संकारक हा $\lif$ आहे; म्हणून
\Intro{\lif} नियम वापरून निष्कर्षापर्यंत पोहोचण्याचा प्रयत्न करू. काम पुढे
जाताना संबंधित गृहीतके लिहून ठेवणे आणि निगमन नियमांना लेबले देणे उत्तम;
त्यामुळे सिद्धतेच्या शेवटी सर्व गृहीतके मुक्त झाली आहेत का हे सहज
दिसते.
\begin{prooftree}
\AxiomC{$\Discharge{\varphi \land \psi}{1}$}
\DeduceC{$\varphi$}
\DischargeRule{\Intro{\lif}}{1}
\UnaryInfC{$(\varphi\land \psi) \lif \varphi$}
\end{prooftree}
आता गृहीतक $\varphi \land \psi$ पासून $\varphi$ पर्यंतच्या पायऱ्या भराव्या लागतील.
येथे हाताळण्यासाठी फक्त एकच संयोजक, $\land$, असल्यामुळे $\land$ चा विलोपन
नियम वापरलाच पाहिजे. त्यामुळे पुढील सिद्धता मिळते:
\begin{prooftree}
\AxiomC{$\Discharge{\varphi \land \psi}{1}$}
\RightLabel{\Elim{\land}}
\UnaryInfC{$\varphi$}
\DischargeRule{\Intro{\lif}}{1}
\UnaryInfC{$(\varphi\land \psi) \lif \varphi$}
\end{prooftree}
आता आपल्याकडे $(\varphi \land \psi) \lif \varphi$ ची योग्य निष्पत्ती आहे.
\end{ex}
\begin{ex}
आता $(\lnot \varphi \lor \psi) \lif (\varphi \lif \psi)$ ची निष्पत्ती देऊ.
हवा असलेला निष्कर्ष निष्पत्तीच्या तळाशी लिहून आपण सुरुवात करतो.
\begin{prooftree}
\AxiomC{}
\UnaryInfC{$(\lnot \varphi \lor \psi) \lif (\varphi \lif \psi)$}
\end{prooftree}
हा निष्कर्ष देऊ शकणारा तार्किक नियम शोधण्यासाठी निष्कर्षातील $\lnot$,
$\lor$ आणि $\lif$ हे तार्किक संयोजक पाहतो. सध्या $\lif$ ची पहिली आस्थितीच
महत्त्वाची आहे, कारण अंतिम क्रमवर्तीतील वाक्याचा तो मुख्य संकारक
आहे; $\lnot$, $\lor$ आणि $\lif$ ची दुसरी आस्थिती दुसऱ्या संयोजकाच्या
व्याप्तीत असल्यामुळे त्यांचा विचार नंतर करू. म्हणून आपण \Intro{\lif}
नियमाने सुरुवात करतो. त्याचे योग्य उपयोजन पुढीलप्रमाणे असले पाहिजे:
\begin{prooftree}
\AxiomC{$\Discharge{\lnot \varphi \lor \psi}{1}$}
\DeduceC{$\varphi \lif \psi$}
\DischargeRule{\Intro{\lif}}{1}
\UnaryInfC{$(\lnot \varphi \lor \psi) \lif (\varphi \lif \psi)$}
\end{prooftree}
पुढे जाण्यासाठी आता दोन पर्याय आहेत. एकतर तळापासून वरच्या दिशेने काम
चालू ठेवून \Intro{\lif} नियमाचे आणखी एक उपयोजन शोधता येईल, किंवा वरून
खाली काम करून \Elim{\lor} नियम लागू करता येईल. दुसरा पर्याय घेऊ.
\Elim{\lor} चे सर्वांत डावे आधारविधान म्हणून $\lnot \varphi \lor \psi$ हे
गृहीतक वापरू. \Elim{\lor} च्या योग्य उपयोजनासाठी इतर दोन्ही आधारविधाने
$\varphi \lif \psi$ या निष्कर्षाशी एकरूप असली पाहिजेत; मात्र प्रत्येकाला
अनुक्रमे आणखी एका गृहीतकापासून, म्हणजे $\lnot \varphi \lor \psi$ च्या दोन
विकल्पघटकांपैकी एकापासून, निष्पन्न करता येते. त्यामुळे आपली
निष्पत्ती पुढीलप्रमाणे दिसेल:
\begin{prooftree}
\AxiomC{$\Discharge{\lnot \varphi \lor \psi}{1}$}
\AxiomC{$\Discharge{\lnot \varphi}{2}$}
\DeduceC{$\varphi \lif \psi$}
\AxiomC{$\Discharge{\psi}{2}$}
\DeduceC{$\varphi \lif \psi$}
\DischargeRule{\Elim{\lor}}{2}
\TrinaryInfC{$\varphi \lif \psi$}
\DischargeRule{\Intro{\lif}}{1}
\UnaryInfC{$(\lnot \varphi \lor \psi) \lif (\varphi \lif \psi)$}
\end{prooftree}
उजवीकडील दोन्ही शाखांत आपल्याला $\varphi \lif \psi$ हे निष्पन्न करायचे आहे;
त्यासाठी \Intro{\lif} वापरणे सर्वांत योग्य आहे.
\begin{prooftree}
\AxiomC{$\Discharge{\lnot \varphi \lor \psi}{1}$}
\AxiomC{$\Discharge{\lnot \varphi}{2}, \Discharge{\varphi}{3}$}
\DeduceC{$\psi$}
\DischargeRule{\Intro{\lif}}{3}
\UnaryInfC{$\varphi \lif \psi$}
\AxiomC{$\Discharge{\psi}{2}, \Discharge{\varphi}{4}$}
\DeduceC{$\psi$}
\DischargeRule{\Intro{\lif}}{4}
\UnaryInfC{$\varphi \lif \psi$}
\DischargeRule{\Elim{\lor}}{2}
\TrinaryInfC{$\varphi \lif \psi$}
\DischargeRule{\Intro{\lif}}{1}
\UnaryInfC{$(\lnot \varphi \lor \psi) \lif (\varphi \lif \psi)$}
\end{prooftree}
निष्पत्तीच्या उरलेल्या दोन भागांसाठी, मधल्या भागात $\lnot \varphi$ आणि
$\varphi$ यांपासून, तर उजवीकडील भागात $\varphi$ आणि $\psi$ यांपासून, $\psi$ मिळवणाऱ्या
निष्पत्त्या हवी आहेत. यांपैकी पहिली आधी घेऊ. $\lnot \varphi$ आणि $\varphi$ ही
\Elim{\lnot} ची दोन आधारविधाने आहेत:
\begin{prooftree}
\AxiomC{$\Discharge{\lnot \varphi}{2}$}
\AxiomC{$\Discharge{\varphi}{3}$}
\RightLabel{\Elim{\lnot}}
\BinaryInfC{$\lfalse$}
\DeduceC{$\psi$}
\end{prooftree}
\FalseInt वापरून $\psi$ हा निष्कर्ष मिळवता येतो आणि शाखा पूर्ण करता येते.
\begin{prooftree}
\AxiomC{$\Discharge{\lnot \varphi \lor \psi}{1}$}
\AxiomC{$\Discharge{\lnot \varphi}{2}$}
\AxiomC{$\Discharge{\varphi}{3}$}
\RightLabel{\Intro{\lfalse}}
\BinaryInfC{$\lfalse$}
\RightLabel{\FalseInt}
\UnaryInfC{$\psi$}
\DischargeRule{\Intro{\lif}}{3}
\UnaryInfC{$\varphi \lif \psi$}
\AxiomC{$\Discharge{\psi}{2}, \Discharge{\varphi}{4}$}
\DeduceC{$\psi$}
\DischargeRule{\Intro{\lif}}{4}
\UnaryInfC{$\varphi \lif \psi$}
\DischargeRule{\Elim{\lor}}{2}
\TrinaryInfC{$\varphi \lif \psi$}
\DischargeRule{\Intro{\lif}}{1}
\UnaryInfC{$(\lnot \varphi \lor \psi) \lif (\varphi \lif \psi)$}
\end{prooftree}
आता सर्वांत उजवी शाखा पाहू. येथे हे लक्षात घेणे महत्त्वाचे आहे की
निष्पत्तीची व्याख्या
\emph{गृहीतके मुक्त करण्याची परवानगी देते}, पण
\emph{तसे करणे आवश्यक ठरवत नाही}. दुसऱ्या शब्दांत, $\varphi$ आणि $\psi$ या
गृहीतकांपैकी एकाचा उपयोग न करता दुसऱ्यापासून $\psi$ निष्पन्न करता येत असेल,
तर ते योग्य आहे. $\psi$ पासून $\psi$ हे निष्पन्न करणे तर क्षुल्लक आहे:
एकटे $\psi$ ही अशी निष्पत्ती असून कोणत्याही अनुमानाची गरज नाही. म्हणून
गृहीतक~$\varphi$ सरळ वगळता येते.
\begin{prooftree}
\AxiomC{$\Discharge{\lnot \varphi \lor \psi}{1}$}
\AxiomC{$\Discharge{\lnot \varphi}{2}$}
\AxiomC{$\Discharge{\varphi}{3}$}
\RightLabel{\Elim{\lnot}}
\BinaryInfC{$\lfalse$}
\RightLabel{\FalseInt}
\UnaryInfC{$\psi$}
\DischargeRule{\Intro{\lif}}{3}
\UnaryInfC{$\varphi \lif \psi$}
\AxiomC{$\Discharge{\psi}{2}$}
\RightLabel{\Intro{\lif}}
\UnaryInfC{$\varphi \lif \psi$}
\DischargeRule{\Elim{\lor}}{2}
\TrinaryInfC{$\varphi \lif \psi$}
\DischargeRule{\Intro{\lif}}{1}
\UnaryInfC{$(\lnot \varphi \lor \psi) \lif (\varphi \lif \psi)$}
\end{prooftree}
लक्षात घ्या की पूर्ण झालेल्या निष्पत्तीमध्ये सर्वांत उजवीकडील
\Intro{\lif} अनुमान प्रत्यक्षात कोणतेही गृहीतक मुक्त करत नाही.
\end{ex}
\begin{ex}
आतापर्यंत आपल्याला \FalseCl{} नियमाची गरज पडलेली नाही. हा नियम विशेष आहे,
कारण नियमाच्या निष्कर्षाचे उप-सूत्र नसलेले गृहीतक तो मुक्त करू देतो.
त्याचा \FalseInt{} नियमाशी जवळचा संबंध आहे. खरे तर \FalseInt{} नियम हे
\FalseCl{} नियमाचे विशेष प्रकरण आहे---``अंतःप्रज्ञावादी तर्कशास्त्र''
नावाचे एक तर्कशास्त्र आहे आणि त्यात केवळ \FalseInt{} अनुज्ञात आहे. इतर
कोणताही मार्ग चालत नसताना \FalseCl{} नियम हा अखेरचा उपाय असतो. उदाहरणार्थ,
आपल्याला $\varphi \lor \lnot \varphi$ हे निष्पन्न करायचे आहे असे समजा. नेहमीच्या
पद्धतीनुसार $\Intro{\lor}$ वापरून $\varphi \lor \lnot \varphi$ हे निष्पन्न
करण्याचा प्रयत्न करू. पण त्यासाठी कोणत्याही गृहीतकांशिवाय $\varphi$ किंवा
$\lnot \varphi$ यांपैकी एक निष्पन्न करावे लागेल, आणि ते शक्य नाही. अशा वेळी
\FalseCl{} उपयोगी पडतो!
\begin{prooftree}
  \AxiomC{$\Discharge{\lnot(\varphi \lor \lnot \varphi)}{1}$}
  \DeduceC{$\lfalse$}
  \DischargeRule{\FalseCl}{1}
  \UnaryInfC{$\varphi \lor \lnot \varphi$}
\end{prooftree}
आता $\lnot(\varphi \lor \lnot \varphi)$ पासून $\lfalse$ ची निष्पत्ती शोधायची
आहे. $\lfalse$ हा $\Elim{\lnot}$ चा निष्कर्ष असल्यामुळे तो नियम वापरून
पाहू:
\begin{prooftree}
  \AxiomC{$\Discharge{\lnot(\varphi \lor \lnot \varphi)}{1}$}
  \DeduceC{$\lnot \varphi$}
  \AxiomC{$\Discharge{\lnot(\varphi \lor \lnot \varphi)}{1}$}
  \DeduceC{$\varphi$}
  \RightLabel{\Elim{\lnot}}
  \BinaryInfC{$\lfalse$}
  \DischargeRule{\FalseCl}{1}
  \UnaryInfC{$\varphi \lor \lnot \varphi$}
\end{prooftree}
$\lnot \varphi$ ची निष्पत्ती शोधण्याच्या आपल्या पद्धतीनुसार
$\Intro{\lnot}$ चे उपयोजन करावे लागेल:
\begin{prooftree}
  \AxiomC{$\Discharge{\lnot(\varphi \lor \lnot \varphi)}{1}, \Discharge{\varphi}{2}$}
  \DeduceC{$\lfalse$}
  \DischargeRule{\Intro{\lnot}}{2}
  \UnaryInfC{$\lnot \varphi$}
  \AxiomC{$\Discharge{\lnot(\varphi \lor \lnot \varphi)}{1}$}
  \DeduceC{$\varphi$}
  \RightLabel{\Elim{\lnot}}
  \BinaryInfC{$\lfalse$}
  \DischargeRule{\FalseCl}{1}
  \UnaryInfC{$\varphi \lor \lnot \varphi$}
\end{prooftree}
येथे $\lnot(\varphi \lor \lnot \varphi)$ हे गृहीतक आणि आपल्या नव्या $\varphi$
गृहीतकापासून~$\Intro{\lor}$ ने निष्पन्न होणारे $\varphi \lor \lnot \varphi$
यांना $\Elim{\lnot}$ लावून $\lfalse$ सहज मिळवता येते:
\begin{prooftree}
  \AxiomC{$\Discharge{\lnot(\varphi \lor \lnot \varphi)}{1}$}
  \AxiomC{$\Discharge{\varphi}{2}$}
  \RightLabel{\Intro{\lor}}
  \UnaryInfC{$\varphi \lor \lnot \varphi$}
  \RightLabel{\Elim{\lnot}}
  \BinaryInfC{$\lfalse$}
  \DischargeRule{\Intro{\lnot}}{2}
  \UnaryInfC{$\lnot \varphi$}
  \AxiomC{$\Discharge{\lnot(\varphi \lor \lnot \varphi)}{1}$}
  \DeduceC{$\varphi$}
  \RightLabel{\Elim{\lnot}}
  \BinaryInfC{$\lfalse$}
  \DischargeRule{\FalseCl}{1}
  \UnaryInfC{$\varphi \lor \lnot \varphi$}
\end{prooftree}
उजव्या बाजूला हीच पद्धत वापरतो; फक्त तेथे $\varphi$ हे~\FalseCl ने मिळवतो:
\begin{prooftree}
  \AxiomC{$\Discharge{\lnot(\varphi \lor \lnot \varphi)}{1}$}
  \AxiomC{$\Discharge{\varphi}{2}$}
  \RightLabel{\Intro{\lor}}
  \UnaryInfC{$\varphi \lor \lnot \varphi$}
  \RightLabel{\Elim{\lnot}}
  \BinaryInfC{$\lfalse$}
  \DischargeRule{\Intro{\lnot}}{2}
  \UnaryInfC{$\lnot \varphi$}
  \AxiomC{$\Discharge{\lnot(\varphi \lor \lnot \varphi)}{1}$}
  \AxiomC{$\Discharge{\lnot \varphi}{3}$}
  \RightLabel{\Intro{\lor}}
  \UnaryInfC{$\varphi \lor \lnot \varphi$}
  \RightLabel{\Elim{\lnot}}
  \BinaryInfC{$\lfalse$}
  \DischargeRule{\FalseCl}{3}
  \UnaryInfC{$\varphi$}
  \RightLabel{\Elim{\lnot}}
  \BinaryInfC{$\lfalse$}
  \DischargeRule{\FalseCl}{1}
  \UnaryInfC{$\varphi \lor \lnot \varphi$}
\end{prooftree}
\end{ex}
\begin{prob}
पुढील गोष्टी दाखवणाऱ्या निष्पत्त्या द्या:
\begin{enumerate}
\item $\varphi \land (\psi \land \chi) \Proves (\varphi \land \psi) \land \chi$.
\item $\varphi \lor (\psi \lor \chi) \Proves (\varphi \lor \psi) \lor \chi$.
\item $\varphi \lif (\psi \lif \chi) \Proves \psi \lif (\varphi \lif \chi)$.
\item $\varphi \Proves \lnot\lnot \varphi$.
\end{enumerate}
\end{prob}
\begin{prob}
पुढील गोष्टी दाखवणाऱ्या निष्पत्त्या द्या:
\begin{enumerate}
\item $(\varphi \lor \psi) \lif \chi \Proves \varphi \lif \chi$.
\item $(\varphi \lif \chi) \land (\psi \lif \chi) \Proves (\varphi \lor \psi) \lif \chi$.
\item $\Proves \lnot(\varphi \land \lnot \varphi)$.
\item $\psi \lif \varphi \Proves \lnot \varphi \lif \lnot \psi$.
\item $\Proves (\varphi \lif \lnot \varphi) \lif \lnot \varphi$.
\item $\Proves \lnot(\varphi \lif \psi) \lif \lnot \psi$.
\item $\varphi \lif \chi \Proves \lnot (\varphi \land \lnot \chi)$.
\item $\varphi \land \lnot \chi \Proves \lnot (\varphi \lif \chi)$.
\item $\varphi \lor \psi, \lnot \psi \Proves \varphi$.
\item $\lnot \varphi \lor \lnot \psi \Proves \lnot(\varphi \land \psi)$.
\item $\Proves (\lnot \varphi \land \lnot \psi) \lif\lnot(\varphi \lor \psi)$.
\item $\Proves \lnot(\varphi \lor \psi) \lif (\lnot \varphi \land \lnot \psi)$.
\end{enumerate}
\end{prob}
\begin{prob}
पुढील गोष्टी दाखवणाऱ्या निष्पत्त्या द्या:
\begin{enumerate}
\item $\lnot(\varphi \lif \psi) \Proves \varphi$.
\item $\lnot(\varphi \land \psi) \Proves \lnot \varphi \lor \lnot \psi$.
\item $\varphi \lif \psi \Proves \lnot \varphi \lor \psi$.
\item $\Proves \lnot \lnot \varphi \lif \varphi$.
\item $\varphi \lif \psi, \lnot \varphi \lif \psi \Proves \psi$.
\item $(\varphi \land \psi) \lif \chi \Proves (\varphi \lif \chi) \lor (\psi \lif \chi)$.
\item $(\varphi \lif \psi) \lif \varphi \Proves \varphi$.
\item $\Proves (\varphi \lif \psi) \lor (\psi \lif \chi)$.
\end{enumerate}
(या सर्वांसाठी $\FalseCl$~नियम आवश्यक आहे.)
\end{prob}
\section{संख्यापकांसह निष्पत्ती}\label{fol:ntd:prq:sec}
\begin{ex}
संख्यापक हाताळताना आयगेन चराच्या अटीचा भंग होणार नाही याची खात्री करावी
लागते; त्यासाठी कधीकधी काही अनुमाने करण्याचा क्रम बदलून पाहावा लागतो.
सामान्यतः आयगेन चराच्या अटीच्या अधीन असलेले नियम आधी हाताळण्याचा प्रयत्न
उपयोगी ठरतो (पूर्ण सिद्धतेत ते अधिक खाली असतील).
सूत्र $\lexists[x][\lnot \varphi(x)] \lif \lnot \lforall[x][\varphi(x)]$
ची निष्पत्ती कशी द्यायची ते पाहू. नेहमीप्रमाणे सुरुवात करून लिहू:
\begin{prooftree}
\AxiomC{}
\UnaryInfC{$\lexists[x][\lnot \varphi(x)]\lif \lnot \lforall[x][\varphi(x)]$}
\end{prooftree}
ही शेवटची पायरी \Intro{\lif} नियमाने समर्थित होण्यासाठी काय आवश्यक आहे ते
प्रथम लिहू.
\begin{prooftree}
\AxiomC{$\Discharge{\lexists[x][\lnot \varphi(x)]}{1}$}
\DeduceC{$\lnot \lforall[x][\varphi(x)]$}
\DischargeRule{\Intro{\lif}}{1}
\UnaryInfC{$\lexists[x][\lnot \varphi(x)]\lif \lnot \lforall[x][\varphi(x)]$}
\end{prooftree}
$\lnot \lforall[x][\varphi(x)]$ वर लागू करता येईल असा कोणताही स्पष्ट नियम
नसल्यामुळे, \Elim{\lexists} नियम वापरता येईल अशा रीतीने निष्पत्ती
मांडू. येथे आयगेन चराच्या अटीकडे लक्ष देऊन, प्रमुख आधारविधान
$\lexists[x][\lnot \varphi(x)]$ किंवा निष्कर्ष यांत, तसेच त्यांच्याकडे नेणाऱ्या
निष्पत्तीतील इतर कोणत्याही मुक्त न केलेल्या गृहीतकात न येणारा अचर निवडावा
लागतो. (परंतु येथे कोणतेही स्थिरांक येत नसल्यामुळे कोणताही अचर चालेल.)
\begin{prooftree}
\AxiomC{$\Discharge{\lexists[x][\lnot \varphi(x)]}{1}$}
\AxiomC{$\Discharge{\lnot \varphi(a)}{2}$}
\DeduceC{$\lnot \lforall[x][\varphi(x)]$}
\DischargeRule{\Elim{\lexists}}{2}
\BinaryInfC{$\lnot \lforall[x][\varphi(x)]$}
\DischargeRule{\Intro{\lif}}{1}
\UnaryInfC{$\lexists[x][\lnot \varphi(x)] \lif \lnot \lforall[x][\varphi(x)]$}
\end{prooftree}
$\lnot \lforall[x][\varphi(x)]$ हे निष्पन्न करण्यासाठी \Intro{\lnot} नियम
वापरून पाहू: यासाठी, आवश्यक असल्यास $\lforall[x][\varphi(x)]$ हे अतिरिक्त गृहीतक वापरून,
व्याघात निष्पन्न करावा लागेल. अर्थात या व्याघातात $\lnot \varphi(a)$ हे गृहीतक
वापरले जाऊ शकते; ते \Elim{\lexists} अनुमानाने मुक्त केले जाईल. ही रचना
पुढीलप्रमाणे करता येते:
\begin{prooftree}
\AxiomC{$\Discharge{\lexists[x][\lnot \varphi(x)]}{1}$}
\AxiomC{$\Discharge{\lnot \varphi(a)}{2}, \Discharge{\lforall[x][\varphi(x)]}{3}$}
\DeduceC{$\lfalse$}
\DischargeRule{\Intro{\lnot}}{3}
\UnaryInfC{$\lnot \lforall[x][\varphi(x)]$}
\DischargeRule{\Elim{\lexists}}{2}
\BinaryInfC{$\lnot \lforall[x][\varphi(x)]$}
\DischargeRule{\Intro{\lif}}{1}
\UnaryInfC{$\lexists[x][\lnot \varphi(x)]\lif \lnot \lforall[x][\varphi(x)]$}
\end{prooftree}
व्याघात मिळण्याच्या आपण अगदी जवळ आलो आहोत. आयगेन चराची अट नसलेला
\Elim{\lforall} नियम येथे लागू करणे सर्वांत सोपे आहे. सार्वत्रिक
संख्यापकाने बद्ध असलेल्या~$x$ च्या जागी हवे ते कोणतेही बंद पद ठेवता येते;
व्याघातापर्यंत पोहोचण्यासाठी~$a$ चाच पुढे वापर करणे सर्वांत सयुक्तिक आहे.
\begin{prooftree}
\AxiomC{$\Discharge{\lexists[x][\lnot \varphi(x)]}{1}$}
\AxiomC{$\Discharge{\lnot \varphi(a)}{2}$}
\AxiomC{$\Discharge{\lforall[x][\varphi(x)]}{3}$}
\RightLabel{\Elim{\lforall}}
\UnaryInfC{$\varphi(a)$}
\RightLabel{\Elim{\lnot}}
\BinaryInfC{$\lfalse$}
\DischargeRule{\Intro{\lnot}}{3}
\UnaryInfC{$\lnot \lforall[x][\varphi(x)]$}
\DischargeRule{\Elim{\lexists}}{2}
\BinaryInfC{$\lnot \lforall[x][\varphi(x)]$}
\DischargeRule{\Intro{\lif}}{1}
\UnaryInfC{$\lexists[x][\lnot \varphi(x)]\lif \lnot \lforall[x][\varphi(x)]$}
\end{prooftree}
विशेषतः संख्यापक हाताळताना, आयगेन चराच्या अटीचा भंग झालेला नाही याची या
टप्प्यावर पुन्हा तपासणी करणे महत्त्वाचे आहे. आपण लावलेल्या नियमांपैकी फक्त
\Elim{\exists} हा नियम त्या अटीच्या अधीन होता. त्याचा आयगेन चर~$a$ प्रमुख
आधारविधानात किंवा निष्कर्षात येत नाही आणि नियमाने मुक्त केलेले तात्पुरते
गृहीतक वगळता संबंधित कोणत्याही मुक्त न केलेल्या गृहीतकातही येत नाही. म्हणून
ही योग्य निष्पत्ती आहे.
\end{ex}
\begin{ex}
कधीकधी इतर सूत्रांपासून सूत्र निष्पन्न करायचे असते. अशा वेळी
काही गृहीतके मुक्त न केलेली राहू शकतात. अंतिम ध्येयाबरोबरच गृहीतकांचीही
नोंद ठेवणे महत्त्वाचे आहे.
गृहीतके $\lexists[x][(\varphi(x) \land \psi(x))]$ आणि
$\lforall[x][(\psi(x) \lif \chi(x,b))]$ यांपासून सूत्र
$\lexists[x][\chi(x,b)]$ ची निष्पत्ती कशी द्यायची ते पाहू.
नेहमीप्रमाणे सुरुवात करून निष्कर्ष तळाशी लिहू.
\begin{prooftree}
\AxiomC{}
\UnaryInfC{$\lexists[x][\chi(x,b)]$}
\end{prooftree}
आपल्याकडे वापरण्यासाठी दोन आधारविधाने आहेत. पहिले वापरण्यासाठी, म्हणजे
$\lexists[x][(\varphi(x) \land \psi(x))]$ पासून
$\lexists[x][\chi(x, b)]$ ची निष्पत्ती शोधण्याचा प्रयत्न करण्यासाठी,
\Elim{\lexists} नियम वापरू. त्याला आयगेन चराची अट असल्यामुळे तो नियम आधी
लावू. मग पुढील मिळते:
\begin{prooftree}
\AxiomC{$\lexists[x][(\varphi(x) \land \psi(x))]$}
\AxiomC{$\Discharge{\varphi(a) \land \psi(a)}{1}$}
\DeduceC{$\lexists[x][\chi(x,b)]$}
\DischargeRule{\Elim{\lexists}}{1}
\BinaryInfC{$\lexists[x][\chi(x,b)]$}
\end{prooftree}
आपण वापरत असलेल्या दोन्ही गृहीतकांत~$\psi$ समान आहे. या टप्प्यावर
\Elim{\land} लावून~$\psi(a)$ वेगळे काढणे उपयोगी ठरेल.
\begin{prooftree}
\AxiomC{$\lexists[x][(\varphi(x) \land \psi(x)])$}
\AxiomC{$\Discharge{\varphi(a)
\land \psi(a)}{1}$}
\RightLabel{\Elim{\land}}
\UnaryInfC{$\psi(a)$}
\DeduceC{$\lexists[x][\chi(x,b)]$}
\DischargeRule{\Elim{\lexists}}{1}
\BinaryInfC{$\lexists[x][\chi(x,b)]$}
\end{prooftree}
वापरण्यासाठी असलेले दुसरे गृहीतक म्हणजे~$\lforall[x][(\psi(x) \lif
  \chi(x,b))]$. येथे आयगेन चराची अट नसल्यामुळे \Elim{\lforall} वापरून
स्थिरांक~$a$ हे $x$ च्या जागी ठेवता येते आणि
$\psi(a) \lif \chi(a, b)$ मिळते. आता आपल्याकडे $\psi(a) \lif \chi(a,b)$ आणि
$\psi(a)$ दोन्ही आहेत. यानंतर \Elim{\lif} नियमाचे सरळ उपयोजन करावे.
\begin{prooftree}
\AxiomC{$\lexists[x][(\varphi(x) \land \psi(x))]$}
\AxiomC{$\lforall[x][(\psi(x) \lif \chi(x,b))]$}
\RightLabel{\Elim{\lforall}}
\UnaryInfC{$\psi(a) \lif \chi(a,b)$}
\AxiomC{$\Discharge{\varphi(a)
\land \psi(a)}{1}$}
\RightLabel{\Elim{\land}}
\UnaryInfC{$\psi(a)$}
\RightLabel{\Elim{\lif}}
\BinaryInfC{$\chi(a,b)$}
\DeduceC{$\lexists[x][\chi(x,b)]$}
\DischargeRule{\Elim{\lexists}}{1}
\insertBetweenHyps{\hspace{-5em}}
\BinaryInfC{$\lexists[x][\chi(x,b)]$}
\end{prooftree}
आपण ध्येयाच्या अगदी जवळ आलो आहोत!{} \Intro{\lexists} चे एक उपयोजन केले
की ध्येय गाठले.
\begin{prooftree}
\AxiomC{$\lexists[x][(\varphi(x) \land \psi(x))]$}
\AxiomC{$\lforall[x][(\psi(x) \lif \chi(x,b))]$}
\RightLabel{\Elim{\lforall}}
\UnaryInfC{$\psi(a) \lif \chi(a,b)$}
\AxiomC{$\Discharge{\varphi(a)
\land \psi(a)}{1}$}
\RightLabel{\Elim{\land}}
\UnaryInfC{$\psi(a)$}
\RightLabel{\Elim{\lif}}
\BinaryInfC{$\chi(a,b)$}
\RightLabel{\Intro{\lexists}}
\UnaryInfC{$\lexists[x][\chi(x,b)]$}
\DischargeRule{\Elim{\lexists}}{1}
\insertBetweenHyps{\hspace{-5em}}
\BinaryInfC{$\lexists[x][\chi(x,b)]$}
\end{prooftree}
प्रत्येक पायरीवर आयगेन चराच्या अटींचा भंग झालेला नाही याची खात्री केल्यामुळे
ही योग्य निष्पत्ती आहे असा विश्वास ठेवता येतो.
\end{ex}
\begin{ex}
गृहीतके $\lforall[x][\varphi(x)] \lif \lexists[y][\psi(y)]$ आणि
$\lnot\lexists[y][\psi(y)]$ यांपासून सूत्र
$\lnot\lforall[x][\varphi(x)]$ ची निष्पत्ती द्या.
नेहमीप्रमाणे सुरुवात करून अपेक्षित सूत्र तळाशी लिहू.
\begin{prooftree}
\AxiomC{}
\UnaryInfC{$\lnot\lforall[x][\varphi(x)]$}
\end{prooftree}
निष्पत्तीची शेवटची ओळ नकरण आहे; म्हणून \Intro{\lnot} वापरून पाहू.
त्यासाठी व्याघात कसा निष्पन्न करायचा ते शोधावे लागेल.
\begin{prooftree}
\AxiomC{$\Discharge{\lforall[x][\varphi(x)]}{1}$}
\DeduceC{$\lfalse$}
\DischargeRule{\Intro{\lnot}}{1}
\UnaryInfC{$\lnot\lforall[x][\varphi(x)]$}
\end{prooftree}
येथपर्यंत सर्व ठीक आहे. \Elim{\lforall} वापरता येईल, पण त्याने ध्येयापर्यंत
पोहोचण्यास मदत होईल का हे स्पष्ट नाही. त्याऐवजी आपले एक गृहीतक वापरू.
$\lforall[x][\varphi(x)] \lif \lexists[y][\psi(y)]$ आणि
$\lforall[x][\varphi(x)]$ यांच्या आधारे \Elim{\lif} नियम वापरता येईल.
\begin{prooftree}
\AxiomC{$\lforall[x][\varphi(x)] \lif \lexists[y][\psi(y)]$}
\AxiomC{$\Discharge{\lforall[x][\varphi(x)]}{1}$}
\RightLabel{\Elim{\lif}}
\BinaryInfC{$\lexists[y][\psi(y)]$}
\DeduceC{$\lfalse$}
\DischargeRule{\Intro{\lnot}}{1}
\UnaryInfC{$\lnot\lforall[x][\varphi(x)]$}
\end{prooftree}
आता वापरण्यासाठी एक अखेरचे गृहीतक उरले आहे. \Elim{\lnot} वापरून
व्याघातापर्यंत पोहोचण्यास ते मदत करेल असे दिसते.
\begin{prooftree}
\AxiomC{$\lnot\lexists[y][\psi(y)]$}
\AxiomC{$\lforall[x][\varphi(x)] \lif \lexists[y][\psi(y)]$}
\AxiomC{$\Discharge{\lforall[x][\varphi(x)]}{1}$}
\RightLabel{\Elim{\lif}}
\BinaryInfC{$\lexists[y][\psi(y)]$}
\RightLabel{\Elim{\lnot}}
\BinaryInfC{$\lfalse$}
\DischargeRule{\Intro{\lnot}}{1}
\UnaryInfC{$\lnot\lforall[x][\varphi(x)]$}
\end{prooftree}
\end{ex}
\begin{prob}
पुढील गोष्टी दाखवणाऱ्या निष्पत्त्या द्या:
\begin{enumerate}
\item $\Proves (\lforall[x][\varphi(x)] \land \lforall[y][\psi(y)]) \lif
\lforall[z][(\varphi(z) \land \psi(z))]$.
\item $\Proves (\lexists[x][\varphi(x)] \lor \lexists[y][\psi(y)]) \lif
\lexists[z][(\varphi(z) \lor \psi(z))]$.
\item $\lforall[x][(\varphi(x) \lif \psi)] \Proves \lexists[y][\varphi(y)] \lif \psi$.
\item $\lforall[x][\lnot \varphi(x)] \Proves \lnot\lexists[x][\varphi(x)]$.
\item $\Proves \lnot\lexists[x][\varphi(x)] \lif \lforall[x][\lnot \varphi(x)]$.
\item $\Proves \lnot\lexists[x][\lforall[y][((\varphi(x,y) \lif \lnot
\varphi(y,y)) \land (\lnot \varphi(y,y) \lif \varphi(x,y)))]]$.
\end{enumerate}
\end{prob}
\begin{prob}
पुढील गोष्टी दाखवणाऱ्या निष्पत्त्या द्या:
\begin{enumerate}
\item $\Proves \lnot\lforall[x][\varphi(x)] \lif \lexists[x][\lnot\varphi(x)]$.
\item $(\lforall[x][\varphi(x)] \lif \psi) \Proves \lexists[y][(\varphi(y) \lif \psi)]$.
\item $\Proves \lexists[x][(\varphi(x) \lif \lforall[y][\varphi(y)])]$.
\end{enumerate}
(या सर्वांसाठी $\FalseCl$~नियम आवश्यक आहे.)
\end{prob}
\section{सिद्धता-उपपत्तीय संकल्पना}\label{fol:ntd:ptn:sec}
\begin{editorial}
या विभागात नैसर्गिक निगमनासाठी सिद्धतायोग्यता संबंध आणि सुसंगतता यांच्या
व्याख्या एकत्रित केल्या आहेत.
\end{editorial}
\begin{explain}
जशा आपण अनेक महत्त्वाच्या चिन्हार्थविषयक संकल्पना (वैधता, तार्किक निष्पन्नता,
पूर्ततायोग्यता) परिभाषित केल्या आहेत, तशाच आता त्यांना अनुरूप
\emph{सिद्धता-उपपत्तीय संकल्पना} परिभाषित करू. या संकल्पना संरचना मध्ये
वाक्यांची पूर्ति होते की नाही याच्या आधारे परिभाषित केलेल्या नसून,
काही वाक्यांची इतरांपासूनची निष्पन्नता किंवा
अनिष्पन्नता यांचा आधार घेऊन परिभाषित केल्या आहेत. या दोन प्रकारच्या
संकल्पना एकमेकांशी जुळतात, हा एक महत्त्वाचा शोध होता. त्या जुळतात, हाच
\emph{निर्दोषता प्रमेय} आणि \emph{संपूर्णता प्रमेय} यांचा आशय आहे.
\end{explain}
\begin{defn}[प्रमेये]
एखादे वाक्य~$\varphi$ हे \emph{प्रमेय} असते, जर नैसर्गिक निगमनात $\varphi$ ची
अशी निष्पत्ती असेल की तिच्यातील सर्व गृहीतके मुक्त आहेत.
$\varphi$ हे प्रमेय असेल तर आपण $\Proves \varphi$ लिहितो आणि ते प्रमेय नसेल तर
$\Proves/ \varphi$ लिहितो.
\end{defn}
\begin{defn}[निष्पन्नता]
वाक्य $\varphi$ हे वाक्यांच्या संच~$\Gamma$ \emph{पासून
निष्पन्न करता येण्याजोगे} असते, म्हणजे $\Gamma \Proves \varphi$, तेव्हा आणि केवळ तेव्हाच,
जेव्हा निष्कर्ष~$\varphi$ असणारी अशी निष्पत्ती असेल की तिच्यातील प्रत्येक
गृहीतक एकतर मुक्त आहे किंवा $\Gamma$ मध्ये आहे. $\varphi$ हे $\Gamma$
पासून निष्पन्न करता येण्याजोगे नसेल, तर आपण $\Gamma \Proves/ \varphi$ लिहितो.
\end{defn}
\begin{defn}[सुसंगतता]
वाक्यांचा संच~$\Gamma$ हा \emph{विसंगत} असतो तेव्हा आणि केवळ तेव्हाच,
जेव्हा $\Gamma \Proves \lfalse$. $\Gamma$ विसंगत नसेल, म्हणजे
$\Gamma \Proves/ \lfalse$ असेल, तर त्याला \emph{सुसंगत} म्हणतो.
\end{defn}
\begin{prop}[परावर्तिता]
\label{fol:ntd:ptn:prop:reflexivity}
$\varphi \in \Gamma$ असेल, तर $\Gamma \Proves \varphi$.
\end{prop}
\begin{proof}
केवळ $\varphi$ हे गृहीतक स्वतःच $\varphi$ ची अशी निष्पत्ती आहे की तिच्यातील
प्रत्येक मुक्त न केलेली गृहीतक (म्हणजे $\varphi$) $\Gamma$ मध्ये आहे.
\end{proof}
\begin{prop}[एकस्वनिकता]
\label{fol:ntd:ptn:prop:monotonicity}
$\Gamma \subseteq \Delta$ आणि $\Gamma \Proves \varphi$ असेल, तर $\Delta
\Proves \varphi$.
\end{prop}
\begin{proof}
$\Gamma$ पासून $\varphi$ ची कोणतीही निष्पत्ती ही $\Delta$ पासून $\varphi$ चीही
निष्पत्ती असते.
\end{proof}
\begin{prop}[संक्रमकता]
\label{fol:ntd:ptn:prop:transitivity}
$\Gamma \Proves \varphi$ आणि $\{\varphi\} \cup \Delta \Proves
\psi$ असेल, तर $\Gamma \cup \Delta \Proves \psi$.
\end{prop}
\begin{proof}
$\Gamma \Proves \varphi$ असेल, तर $\varphi$ ची अशी निष्पत्ती~$\delta_0$ असते
की तिच्यातील सर्व मुक्त न केलेली गृहीतके $\Gamma$ मध्ये असतात. $\{\varphi\}
\cup \Delta \Proves \psi$ असेल, तर $\psi$ ची अशी निष्पत्ती~$\delta_1$
असते की तिच्यातील सर्व मुक्त न केलेली गृहीतके $\{\varphi\} \cup \Delta$ मध्ये
असतात. आता पुढील रचना विचारात घ्या:
\begin{prooftree}
  \AxiomC{$\Delta, \Discharge{\varphi}{1}$}
  \RightLabel{$\delta_1$}
  \DeduceC{$\psi$}
  \DischargeRule{\Intro{\lif}}{1}
  \UnaryInfC{$\varphi \lif \psi$}
  \AxiomC{$\Gamma$}
  \RightLabel{$\delta_0$}
  \DeduceC{$\varphi$}
  \RightLabel{\Elim{\lif}}
  \BinaryInfC{$\psi$}
\end{prooftree}
आता सर्व मुक्त न केलेली गृहीतके $\Gamma \cup \Delta$ मध्ये आहेत; म्हणून
यावरून $\Gamma \cup \Delta \Proves \psi$ हे दिसते.
\end{proof}
$\Gamma = \{\varphi_1, \varphi_2, \ldots, \varphi_k\}$ हा सांत संच असेल, तेव्हा
$\Gamma \Proves \psi$ ऐवजी आपण $\varphi_1,\varphi_2,\ldots,\varphi_k \Proves \psi$ हे
सरलीकृत संकेतलेखन वापरू शकतो. विशेषतः, $\varphi \Proves \psi$ याचा अर्थ
$\{\varphi\} \Proves \psi$ असा होतो.
लक्षात घ्या की $\Gamma \Proves \varphi$ आणि $\varphi
\Proves \psi$ असेल, तर $\Gamma \Proves \psi$. तसेच $\varphi_1,
\dots, \varphi_n \Proves \psi$ असेल आणि प्रत्येक~$i$ साठी $\Gamma \Proves \varphi_i$
असेल, तर $\Gamma \Proves \psi$, हेसुद्धा यावरून मिळते.
\begin{prop}
\label{fol:ntd:ptn:prop:incons}
पुढील विधाने सममूल्य आहेत.
    \begin{enumerate}
      \item \( \Gamma \) विसंगत आहे.
      \item प्रत्येक वाक्य~\( {\varphi} \) साठी \( \Gamma \Proves {\varphi} \).
      \item एखाद्या वाक्य~\( {\varphi} \) साठी \( \Gamma \Proves {\varphi} \) आणि \( \Gamma \Proves \lnot {\varphi} \).
    \end{enumerate}
\end{prop}
\begin{proof}
सराव.
\end{proof}
\begin{prob}
\ref{fol:ntd:ptn:prop:incons} सिद्ध करा.
\end{prob}
\begin{prop}[संहतता]
  \label{fol:ntd:ptn:prop:proves-compact}
  \begin{enumerate}
  \item $\Gamma \Proves \varphi$ असेल, तर एखादा सांत उपसंच $\Gamma_0
    \subseteq \Gamma$ असा आहे की $\Gamma_0 \Proves \varphi$.
  \item $\Gamma$ चा प्रत्येक सांत उपसंच सुसंगत असेल, तर $\Gamma$ सुसंगत
    असतो.
  \end{enumerate}
\end{prop}
\begin{proof}
  \begin{enumerate}
    \item $\Gamma \Proves \varphi$ असेल, तर $\Gamma$ पासून $\varphi$ ची
      निष्पत्ती~$\delta$ असते. $\Gamma_0$ हा $\delta$ मधील
      मुक्त न केलेली गृहीतकांचा संच असू द्या. प्रत्येक निष्पत्ती सांत
      असल्यामुळे $\Gamma_0$ मधील वाक्यांची संख्या सांतच असते.
      म्हणून $\delta$ ही सांत~$\Gamma_0 \subseteq \Gamma$ पासून $\varphi$ ची
      निष्पत्ती आहे.
    \item (1) मध्ये $\varphi
      \ident \lfalse$ हे विशेष प्रकरण घेतल्यावर मिळणारे हे निषेधात्मक उलट
      विधान आहे.
  \end{enumerate}
\end{proof}
\section{निष्पन्नता आणि सुसंगतता}\label{fol:ntd:prv:sec}
आता आपण निष्पन्नता संबंधाचे अनेक गुणधर्म सिद्ध करू. ते स्वतंत्रपणेही
महत्त्वाचे आहेत, पण संपूर्णता प्रमेयाच्या सिद्धतेत प्रत्येकाची भूमिका असेल.
\begin{prop}\label{fol:ntd:prv:prop:provability-contr}
  जर $\Gamma \Proves \varphi$ आणि $\Gamma \cup \{\varphi\}$ विसंगत असेल, तर
  $\Gamma$ विसंगत आहे.
\end{prop}
\begin{proof}
$\Gamma$ पासून $\varphi$ ची निष्पत्ती~$\delta_1$ आणि $\Gamma \cup \{\varphi\}$
पासून $\lfalse$ ची निष्पत्ती~$\delta_2$ असू द्या. मग आपण पुढीलप्रमाणे
निष्पन्न करू शकतो:
\begin{prooftree}
\AxiomC{$\Gamma, \Discharge{\varphi}{1}$}
\RightLabel{$\delta_2$}
\DeduceC{$\lfalse$}
\DischargeRule{\Intro{\lnot}}{1}
\UnaryInfC{$\lnot \varphi$}
\AxiomC{$\Gamma$}
\RightLabel{$\delta_1$}
\DeduceC{$\varphi$}
\RightLabel{\Elim{\lnot}}
\BinaryInfC{$\lfalse$}
\end{prooftree}
नव्या निष्पत्तीमध्ये गृहीतक~$\varphi$ हे मुक्त आहे; म्हणून ही
$\Gamma$ पासूनची निष्पत्ती आहे.
\end{proof}
\begin{prop}
\label{fol:ntd:prv:prop:prov-incons}
$\Gamma \Proves \varphi$ तेव्हा आणि केवळ तेव्हाच, जेव्हा $\Gamma \cup \{\lnot \varphi\}$
विसंगत असते.
\end{prop}
\begin{proof}
प्रथम $\Gamma \Proves \varphi$ असे समजा; म्हणजे $\Gamma$ मधील
मुक्त न केलेली गृहीतकांपासून $\varphi$ ची निष्पत्ती~$\delta_0$ असते.
पुढीलप्रमाणे $\Gamma \cup \{\lnot \varphi\}$ पासून $\lfalse$ ची
निष्पत्ती मिळते:
\begin{prooftree}
  \AxiomC{$\lnot \varphi$}
  \AxiomC{$\Gamma$}
  \RightLabel{$\delta_0$}
  \DeduceC{$\varphi$}
  \RightLabel{\Elim{\lnot}}
  \BinaryInfC{$\lfalse$}
\end{prooftree}
आता $\Gamma \cup \{\lnot \varphi\}$ विसंगत आहे असे समजा आणि $\Gamma \cup
\{\lnot \varphi\}$ मधील मुक्त न केलेली गृहीतकांपासून $\lfalse$ ची संबंधित
निष्पत्ती~$\delta_1$ असू द्या. $\FalseCl$ वापरून पुढीलप्रमाणे केवळ
$\Gamma$ पासून $\varphi$ ची निष्पत्ती मिळते:
\begin{prooftree}
  \AxiomC{$\Gamma, \Discharge{\lnot \varphi}{1}$}
  \RightLabel{$\delta_1$}
  \DeduceC{$\lfalse$}
  \RightLabel{\FalseCl}
  \DischargeRule{\FalseCl}{1}
  \UnaryInfC{$\varphi$}
\end{prooftree}
\end{proof}
\begin{prob}
$\Gamma \Proves \lnot \varphi$ तेव्हा आणि केवळ तेव्हाच, जेव्हा $\Gamma \cup \{\varphi\}$
विसंगत असते, हे सिद्ध करा.
\end{prob}
\begin{prop}\label{fol:ntd:prv:prop:explicit-inc}
  जर $\Gamma \Proves \varphi$ आणि $\lnot \varphi \in \Gamma$ असेल, तर $\Gamma$
  विसंगत आहे.
\end{prop}
\begin{proof}
  समजा $\Gamma \Proves \varphi$ आणि $\lnot \varphi \in \Gamma$. मग $\Gamma$ पासून
  $\varphi$ ची निष्पत्ती~$\delta$ असते. $\Elim{\lnot}$ नियमाचा पुढील सोपा
  उपयोग विचारात घ्या:
  \begin{prooftree}
    \AxiomC{$\lnot \varphi$}
    \AxiomC{$\Gamma$}
    \RightLabel{$\delta$}
    \DeduceC{$\varphi$}
    \RightLabel{\Elim{\lnot}}
    \BinaryInfC{$\lfalse$}
  \end{prooftree}
  $\lnot \varphi \in \Gamma$ असल्यामुळे सर्व मुक्त न केलेली गृहीतके
  $\Gamma$ मध्ये आहेत. म्हणून यावरून $\Gamma \Proves \lfalse$ हे दिसते.
\end{proof}
\begin{prop}\label{fol:ntd:prv:prop:provability-exhaustive}
  $\Gamma \cup \{\varphi\}$ आणि $\Gamma \cup \{\lnot \varphi\}$ हे दोन्ही संच
  विसंगत असतील, तर $\Gamma$ विसंगत आहे.
\end{prop}
\begin{proof}
$\delta_1$ आणि $\delta_2$ या अनुक्रमे $\Gamma \cup \{ \varphi \}$ पासून
$\lfalse$ ची आणि $\Gamma \cup \{ \lnot \varphi \}$ पासून $\lfalse$ ची
निष्पत्त्या असू द्या. मग आपण पुढीलप्रमाणे निष्पन्न करू शकतो:
\begin{prooftree}
\AxiomC{$\Gamma, \Discharge{\lnot \varphi}{2}$}
\RightLabel{$\delta_2$}
\DeduceC{$\lfalse$}
\DischargeRule{\Intro{\lnot}}{2}
\UnaryInfC{$\lnot \lnot \varphi$}
\AxiomC{$\Gamma, \Discharge{\varphi}{1}$}
\RightLabel{$\delta_1$}
\DeduceC{$\lfalse$}
\DischargeRule{\Intro{\lnot}}{1}
\UnaryInfC{$\lnot \varphi$}
\RightLabel{\Elim{\lnot}}
\BinaryInfC{$\lfalse$}
\end{prooftree}
$\varphi$ आणि $\lnot \varphi$ ही गृहीतके मुक्त असल्यामुळे ही केवळ
$\Gamma$ पासून $\lfalse$ ची निष्पत्ती आहे. म्हणून $\Gamma$ विसंगत
आहे.
\end{proof}
\section{निष्पन्नता आणि विधानीय संयोजक}\label{fol:ntd:ppr:sec}
\begin{explain}
  नैसर्गिक निगमनाचा निष्पन्नता संबंध~$\Proves$ हा विधानीय संयोजकांशी
  संबंधित काही मूलभूत तथ्ये सिद्ध करण्यासाठी पुरेसा सामर्थ्यवान आहे; उदाहरणार्थ,
  $\varphi \land \psi
  \Proves \varphi$ आणि $\varphi, \varphi \lif \psi \Proves \psi$ (विध्यात्मक अनुमान). ही
  तथ्ये संपूर्णता प्रमेयाच्या सिद्धतेसाठी आवश्यक आहेत.
\end{explain}
\begin{prop}\label{fol:ntd:ppr:prop:provability-land}
  \begin{enumerate}
  \item \label{fol:ntd:ppr:prop:provability-land-left} $\varphi \land \psi \Proves
    \varphi$ आणि $\varphi \land \psi \Proves \psi$ ही दोन्ही तथ्ये सत्य आहेत.
  \item \label{fol:ntd:ppr:prop:provability-land-right} $\varphi, \psi \Proves \varphi \land \psi$.
  \end{enumerate}
\end{prop}
\begin{proof}
  \begin{enumerate}
  \item आपण पुढील दोन्ही निष्पन्न करू शकतो:
    \begin{prooftree}
      \AxiomC{$\varphi \land \psi$}
      \RightLabel{\Elim{\land}}
      \UnaryInfC{$\varphi$}
      \DisplayProof\qquad\bottomAlignProof
      \AxiomC{$\varphi \land \psi$}
      \RightLabel{\Elim{\land}}
      \UnaryInfC{$\psi$}
    \end{prooftree}
  \item आपण पुढील निष्पन्न करू शकतो:
    \begin{prooftree}
      \AxiomC{$\varphi$}
      \AxiomC{$\psi$}
      \RightLabel{\Intro{\land}}
      \BinaryInfC{$\varphi \land \psi$}
    \end{prooftree}
  \end{enumerate}
\end{proof}
\begin{prop}\label{fol:ntd:ppr:prop:provability-lor}
  \begin{enumerate}
  \item $\varphi \lor \psi, \lnot \varphi, \lnot \psi$ विसंगत आहे.
  \item $\varphi \Proves \varphi \lor \psi$ आणि $\psi \Proves \varphi \lor \psi$ ही दोन्ही
    तथ्ये सत्य आहेत.
  \end{enumerate}
\end{prop}
\begin{proof}
  \begin{enumerate}
  \item पुढील निष्पत्ती विचारात घ्या:
    \begin{prooftree}
      \AxiomC{$\varphi \lor \psi$}
      \AxiomC{$\lnot \varphi$}
      \AxiomC{$\Discharge{\varphi}{1}$}
      \RightLabel{\Elim{\lnot}}
      \BinaryInfC{$\lfalse$}
      \AxiomC{$\lnot \psi$}
      \AxiomC{$\Discharge{\psi}{1}$}
      \RightLabel{\Elim{\lnot}}
      \BinaryInfC{$\lfalse$}
      \DischargeRule{\Elim{\lor}}{1}
      \TrinaryInfC{$\lfalse$}
    \end{prooftree}
    ही $\varphi \lor \psi$, $\lnot \varphi$ आणि $\lnot \psi$ या मुक्त न केलेली
    गृहीतकांपासून $\lfalse$ ची निष्पत्ती आहे.
  \item आपण पुढील दोन्ही निष्पन्न करू शकतो:
    \begin{prooftree}
      \AxiomC{$\varphi$}
      \RightLabel{\Intro{\lor}}
      \UnaryInfC{$\varphi \lor \psi$}
      \DisplayProof\qquad\bottomAlignProof
      \AxiomC{$\psi$}
      \RightLabel{\Intro{\lor}}
      \UnaryInfC{$\varphi \lor \psi$}
    \end{prooftree}
  \end{enumerate}
\end{proof}
\begin{prop}\label{fol:ntd:ppr:prop:provability-lif}
  \begin{enumerate}
  \item \label{fol:ntd:ppr:prop:provability-lif-left} $\varphi, \varphi \lif \psi \Proves \psi$.
  \item \label{fol:ntd:ppr:prop:provability-lif-right}
    $\lnot \varphi \Proves \varphi \lif \psi$ आणि $\psi \Proves \varphi \lif \psi$ ही दोन्ही
    तथ्ये सत्य आहेत.
  \end{enumerate}
\end{prop}
\begin{proof}
  \begin{enumerate}
  \item आपण पुढील निष्पन्न करू शकतो:
    \begin{prooftree}
      \AxiomC{$\varphi \lif \psi$}
      \AxiomC{$\varphi$}
      \RightLabel{\Elim{\lif}}
      \BinaryInfC{$\psi$}
    \end{prooftree}
  \item हे पुढील दोन निष्पत्त्या दाखवतात:
    \begin{prooftree}
      \AxiomC{$\lnot \varphi$}
      \AxiomC{$\Discharge{\varphi}{1}$}
      \RightLabel{\Elim{\lnot}}
      \BinaryInfC{$\lfalse$}
      \RightLabel{\FalseInt}
      \UnaryInfC{$\psi$}
      \DischargeRule{\Intro{\lif}}{1}
      \UnaryInfC{$\varphi \lif \psi$}
      \DisplayProof\qquad\bottomAlignProof
      \AxiomC{$\psi$}
      \RightLabel{\Intro{\lif}}
      \UnaryInfC{$\varphi \lif \psi$}
    \end{prooftree}
    लक्षात घ्या की $\Intro{\lif}$ गृहीतक~$\varphi$ हे मुक्त करू शकतो;
    पण ते करणे आवश्यक नाही.
  \end{enumerate}
\end{proof}
\section{निष्पन्नता आणि संख्यापक}\label{fol:ntd:qpr:sec}
\begin{explain}
  संपूर्णता प्रमेयासाठी या विभागात स्थापित केलेली~$\Proves$ संबंधी तथ्ये
  नैसर्गिक निगमनाच्या नियमांनी निष्पन्न होतात, हेही आवश्यक आहे.
\end{explain}
\begin{thm}
\label{fol:ntd:qpr:thm:strong-generalization} जर $c$ हे $\Gamma$ किंवा $\varphi(x)$ मध्ये
न येणारे स्थिर असेल आणि $\Gamma \Proves \varphi(c)$, तर $\Gamma
\Proves \lforall[x][\varphi(x)]$.
\end{thm}
\begin{proof}
$\delta$ ही $\Gamma$ पासून $\varphi(c)$ ची निष्पत्ती असू द्या.
\Intro{\lforall} अनुमान जोडल्यावर आपल्याला $\lforall[x][\varphi(x)]$ ची
निष्पत्ती मिळते. $c$ हे $\Gamma$ किंवा $\varphi(x)$ मध्ये येत नसल्यामुळे
आयगेन चराची अट पूर्ण होते.
\end{proof}
\begin{prop}
\label{fol:ntd:qpr:prop:provability-quantifiers}
\begin{enumerate}
\item $\varphi(t) \Proves \lexists[x][\varphi(x)]$.
\item $\lforall[x][\varphi(x)] \Proves \varphi(t)$.
\end{enumerate}
\end{prop}
\begin{proof}
\begin{enumerate}
\item पुढील ही~$\lexists[x][\varphi(x)]$ ची $\varphi(t)$ पासूनची
  निष्पत्ती आहे:
\begin{prooftree}
\AxiomC{$\varphi(t)$}
\RightLabel{\Intro{\lexists}}
\UnaryInfC{$\lexists[x][\varphi(x)]$}
\end{prooftree}
\item पुढील ही~$\varphi(t)$ ची $\lforall[x][\varphi(x)]$ पासूनची
  निष्पत्ती आहे:
\begin{prooftree}
\AxiomC{$\lforall[x][\varphi(x)]$}
\RightLabel{\Elim{\lforall}}
\UnaryInfC{$\varphi(t)$}
\end{prooftree}
\end{enumerate}
\end{proof}
\section{निर्दोषता}\label{fol:ntd:sou:sec}
\begin{explain}
नैसर्गिक निगमनासारखी निष्पत्ती-पद्धत \emph{निर्दोष} असते, जर ती
प्रत्यक्षात निष्पन्न न होणाऱ्या गोष्टी निष्पन्न करू शकत नसेल. त्यामुळे
निर्दोषता हा निष्पत्ती-पद्धतींसाठी हमी दिलेल्या सुरक्षिततेचा एक प्रकार
आहे. कोणता सिद्धता-उपपत्तीय गुणधर्म विचारात आहे यावर अवलंबून, उदाहरणार्थ,
आपल्याला पुढील गोष्टी हव्या असतात:
\begin{enumerate}
\item प्रत्येक निष्पन्न करता येण्याजोगे वाक्य हे वैध असते;
\item एखादे वाक्य इतर काही वाक्यांपासून निष्पन्न करता येण्याजोगे असेल, तर ते
  त्यांचा तार्किक परिणामही असते;
\item वाक्यांचा संच विसंगत असेल, तर तो अपूर्ततायोग्य असतो.
\end{enumerate}
हे निष्पत्ती-पद्धतीचे महत्त्वाचे गुणधर्म आहेत. यांपैकी एखादा गुणधर्म
खरा नसेल, तर निष्पत्ती-पद्धतीत दोष आहे---ती खूप जास्त गोष्टी
निष्पन्न करेल. म्हणून निष्पत्ती-पद्धतीची निर्दोषता सिद्ध करणे अत्यंत
महत्त्वाचे आहे.
\end{explain}
\begin{thm}[निर्दोषता]
\label{fol:ntd:sou:thm:soundness}
$\varphi$ हे मुक्त न केलेली गृहीतके~$\Gamma$ यांच्यापासून निष्पन्न करता येण्याजोगे असेल, तर
$\Gamma \Entails \varphi$.
\end{thm}
\begin{proof}
$\delta$ ही $\varphi$ ची निष्पत्ती असू द्या. $\delta$ मधील अनुमानांच्या
संख्येवर विगमनाने पुढे जाऊ.
विगमनाच्या आधारपायरीसाठी अनुमानांची संख्या~$0$ असेल तेव्हा दावा सिद्ध करू.
या प्रकरणात $\delta$ मध्ये फक्त एकच वाक्य~$\varphi$, म्हणजे एक गृहीतक,
असते. ते गृहीतक मुक्त न केलेली असते, कारण गृहीतके फक्त अनुमानांद्वारे
मुक्त होऊ शकतात आणि येथे कोणतेही अनुमान नाही. म्हणून सिद्धतेतील
सर्व मुक्त न केलेली गृहीतके पूर्ण करणारी कोणतीही
संरचना~$\Struct{M}$
$\varphi$ हेदेखील पूर्ण करते.
आता विगमन पायरी पाहू. $\delta$ मध्ये $n$ अनुमाने आहेत असे समजा. सर्वांत
खालच्या अनुमानाची आधारविधाने अशा उप-निष्पत्त्या वापरून निष्पन्न केली
आहेत की प्रत्येकामध्ये $n$ पेक्षा कमी अनुमाने आहेत. विगमन गृहीतक असे धरू:
सर्वांत खालच्या अनुमानाची आधारविधाने, त्या आधारविधानांवर संपणाऱ्या
उप-निष्पत्त्यांच्या मुक्त न केलेली गृहीतकांपासून निष्पन्न होतात. संपूर्ण
सिद्धतेच्या मुक्त न केलेली गृहीतकांपासून निष्कर्ष~$\varphi$ निष्पन्न होतो, हे
आपल्याला दाखवायचे आहे.
सर्वांत खालच्या अनुमानाच्या प्रकारानुसार प्रकरणे वेगळी करू. प्रथम एकच
आधारविधान असलेली संभाव्य अनुमाने पाहू.
\begin{enumerate}
\item शेवटचे अनुमान \Intro{\lnot} आहे असे समजा: निष्पत्तीचे रूप
  पुढीलप्रमाणे आहे:
  \begin{prooftree}
    \AxiomC{$\Gamma, \Discharge{\varphi}{n}$}
    \RightLabel{$\delta_1$}
    \DeduceC{$\lfalse$}
    \DischargeRule{\Intro{\lnot}}{n}
    \UnaryInfC{$\lnot \varphi$}
  \end{prooftree}
  विगमन गृहीतकानुसार $\lfalse$ हे $\delta_1$ च्या मुक्त न केलेली
  गृहीतकांपासून~$\Gamma \cup \{\varphi\}$ निष्पन्न होते. एखादी
  संरचना~$\Struct{M}$
  विचारात घ्या. जर $\Sat{M}{\Gamma}$, तर
  $\Sat{M}{\lnot \varphi}$, हे दाखवायचे आहे.
  विसंगतीद्वारे सिद्धतेसाठी असे समजा की
  $\Sat{M}{\Gamma}$, पण
  $\Sat/{M}{\lnot \varphi}$, म्हणजे
  $\Sat{M}{\varphi}$. याचा अर्थ
  $\Sat{M}{\Gamma \cup \{\varphi\}}$ असा होईल. हे आपल्या
  विगमन गृहीतकाच्या विरुद्ध आहे. म्हणून
  $\Sat{M}{\lnot \varphi}$.
\item शेवटचे अनुमान \Elim{\land} आहे. याचे दोन प्रकार आहेत: $\varphi \land
  \psi$ या आधारविधानापासून $\varphi$ किंवा $\psi$ निष्पन्न केले जाऊ शकते. पहिले
  प्रकरण पाहू. निष्पत्ती~$\delta$ पुढीलप्रमाणे दिसते:
  \begin{prooftree}
    \AxiomC{$\Gamma$}
    \RightLabel{$\delta_1$}
    \DeduceC{$\varphi \land \psi$}
    \RightLabel{\Elim{\land}}
    \UnaryInfC{$\varphi$}
  \end{prooftree}
  विगमन गृहीतकानुसार $\varphi \land \psi$ हे $\delta_1$ च्या
  मुक्त न केलेली गृहीतकांपासून~$\Gamma$ निष्पन्न होते. एखादी
  संरचना~$\Struct{M}$ विचारात घ्या.
  जर $\Sat{M}{\Gamma}$, तर
  $\Sat{M}{\varphi}$, हे दाखवायचे आहे.
  $\Sat{M}{\Gamma}$ असे समजा. आपल्या विगमन
  गृहीतकानुसार ($\Gamma \Entails \varphi \land \psi$),
  $\Sat{M}{\varphi \land \psi}$ हे आपल्याला माहीत आहे.
  व्याख्येनुसार, $\Sat{M}{\varphi \land \psi}$ तेव्हा आणि
  केवळ तेव्हाच, जेव्हा $\Sat{M}{\varphi}$ आणि
  $\Sat{M}{\psi}$. ($\varphi \land \psi$ पासून $\psi$
  निष्पन्न करण्याचे प्रकरणही असेच हाताळता येते.)
\item शेवटचे अनुमान \Intro{\lor} आहे. याचे दोन प्रकार आहेत: $\varphi$ या
  आधारविधानापासून किंवा $\psi$ या आधारविधानापासून $\varphi \lor \psi$ निष्पन्न
  केले जाऊ शकते. पहिले प्रकरण पाहू. निष्पत्तीचे रूप पुढीलप्रमाणे आहे:
  \begin{prooftree}
    \AxiomC{$\Gamma$}
    \RightLabel{$\delta_1$}
    \DeduceC{$\varphi$}
    \RightLabel{\Intro{\lor}}
    \UnaryInfC{$\varphi \lor \psi$}
  \end{prooftree}
  विगमन गृहीतकानुसार $\varphi$ हे $\delta_1$ च्या मुक्त न केलेली
  गृहीतकांपासून~$\Gamma$ निष्पन्न होते. एखादी
  संरचना~$\Struct{M}$
  विचारात घ्या. जर $\Sat{M}{\Gamma}$, तर
  $\Sat{M}{\varphi \lor \psi}$, हे दाखवायचे आहे.
  $\Sat{M}{\Gamma}$ असे समजा; मग
  $\Gamma \Entails \varphi$ (विगमन गृहीतक) असल्यामुळे
  $\Sat{M}{\varphi}$. म्हणून
  $\Sat{M}{\varphi \lor \psi}$ हेदेखील असलेच पाहिजे.
  ($\psi$ पासून $\varphi \lor \psi$ निष्पन्न करण्याचे प्रकरणही असेच हाताळता येते.)
\item शेवटचे अनुमान \Intro{\lif} आहे. $\varphi$ हे गृहीतक आणि $\psi$ हा निष्कर्ष
  असलेल्या उपसिद्धतेपासून $\varphi \lif \psi$ निष्पन्न केले जाते, म्हणजे:
  \begin{prooftree}
    \AxiomC{$\Gamma, \Discharge{\varphi}{n}$}
    \RightLabel{$\delta_1$}
    \DeduceC{$\psi$}
    \DischargeRule{\Intro{\lif}}{n}
    \UnaryInfC{$\varphi \lif \psi$}
  \end{prooftree}
  विगमन गृहीतकानुसार $\psi$ हे $\delta_1$ च्या मुक्त न केलेली
  गृहीतकांपासून निष्पन्न होते, म्हणजे $\Gamma \cup \{\varphi\} \Entails \psi$.
  एखादी
  संरचना~$\Struct{M}$
  विचारात घ्या. शेवटच्या अनुमानात $\varphi$ चे विसर्जन झाल्यामुळे $\delta$
  ची मुक्त न केलेली गृहीतके फक्त $\Gamma$ आहेत. म्हणून
  $\Gamma \Entails \varphi \lif \psi$ हे दाखवायचे आहे. विसंगतीद्वारे सिद्धतेसाठी
  असे समजा की एखाद्या
  संरचना~$\Struct{M}$
  साठी $\Sat{M}{\Gamma}$, पण
  $\Sat/{M}{\varphi \lif \psi}$. मग
  $\Sat{M}{\varphi}$ आणि
  $\Sat/{M}{\psi}$. पण गृहीतकानुसार $\psi$ हा
  $\Gamma \cup \{\varphi\}$ चा तार्किक परिणाम आहे, म्हणजे
  $\Sat{M}{\psi}$; ही विसंगती आहे. म्हणून
  $\Gamma \Entails \varphi \lif \psi$.
\item शेवटचे अनुमान \FalseInt आहे. येथे $\delta$ पुढीलप्रमाणे संपते:
  \begin{prooftree}
    \AxiomC{$\Gamma$}
    \RightLabel{$\delta_1$}
    \DeduceC{$\lfalse$}
    \RightLabel{\FalseInt}
    \UnaryInfC{$\varphi$}
  \end{prooftree}
  विगमन गृहीतकानुसार $\Gamma \Entails \lfalse$. आपल्याला
  $\Gamma \Entails \varphi$ हे दाखवायचे आहे. तसे नसल्याचे समजा; मग एखाद्या
  $\Struct{M}$ साठी
    $\Sat{M}{\Gamma}$ आणि
    $\Sat/{M}{\varphi}$. पण नेहमीच
    $\Sat/{M}{\lfalse}$ असल्यामुळे याचा अर्थ
    $\Gamma \Entails/ \lfalse$ असा होईल, जो विगमन गृहीतकाच्या विरुद्ध आहे.
\item शेवटचे अनुमान \FalseCl आहे: सराव.
\item शेवटचे अनुमान \Intro{\lforall} आहे. मग $\delta$ चे रूप पुढीलप्रमाणे आहे:
  \begin{prooftree}
    \AxiomC{$\Gamma$}
    \RightLabel{$\delta_1$}
    \DeduceC{$\varphi(a)$}
    \RightLabel{\Intro{\lforall}}
    \UnaryInfC{$\lforall[x][\varphi(x)]$}
  \end{prooftree}
  विगमन गृहीतकानुसार आधारविधान~$\varphi(a)$ हे मुक्त न केलेली
  गृहीतके~$\Gamma$ यांचा तार्किक परिणाम आहे.
  $\Sat{M}{\Gamma}$ अशी एखादी रचना~$\Struct{M}$ विचारात घ्या. आपल्याला
  $\Sat{M}{\lforall[x][\varphi(x)]}$ हे दाखवायचे आहे.
  $\lforall[x][\varphi(x)]$ हे वाक्य असल्यामुळे याचा अर्थ प्रत्येक
  चल-विन्यास~$s$ साठी $\Sat{M}{\varphi(x)}[s]$ दाखवायचे आहे
  (\ref{fol:syn:ass:prop:sat-quant}). $\Gamma$ मध्ये केवळ वाक्ये असल्यामुळे
  \ref{fol:syn:sat:defn:satisfaction} नुसार प्रत्येक $\psi \in \Gamma$ साठी
  $\Sat{M}{\psi}[s]$. $\Struct{M'}$ ही $\Struct{M}$ सारखीच रचना असू द्या,
  फक्त $\Assign{a}{M'} = s(x)$. $a$ हे $\Gamma$ मध्ये येत नसल्यामुळे
  \ref{fol:syn:ext:cor:extensionality-sent} नुसार
  $\Sat{M'}{\Gamma}$. $\Gamma \Entails \varphi(a)$ असल्यामुळे
  $\Sat{M'}{\varphi(a)}$. $\varphi(a)$ हे वाक्य असल्यामुळे
  \ref{fol:syn:ass:prop:sentence-sat-true} नुसार
  $\Sat{M'}{\varphi(a)}[s]$. \ref{fol:syn:ext:prop:ext-formulas} नुसार
  $\Sat{M'}{\varphi(x)}[s]$ तेव्हा आणि केवळ तेव्हाच, जेव्हा
  $\Sat{M'}{\varphi(a)}$ (लक्षात घ्या की $\varphi(a)$ म्हणजे
  $\Subst{\varphi(x)}{a}{x}$). म्हणून $\Sat{M'}{\varphi(x)}[s]$. $a$ हे $\varphi(x)$
  मध्ये येत नसल्यामुळे \ref{fol:syn:ext:prop:extensionality} नुसार
  $\Sat{M}{\varphi(x)}[s]$. पण $s$ हा स्वेच्छ चल-विन्यास होता, म्हणून
  $\Sat{M}{\lforall[x][\varphi(x)]}$.
\item शेवटचे अनुमान \Intro{\lexists} आहे: सराव.
\item शेवटचे अनुमान \Elim{\forall} आहे: सराव.
\end{enumerate}
आता अनेक आधारविधाने असलेली संभाव्य अनुमाने पाहू:
\Elim{\lor}, \Intro{\land}, \Elim{\lif} आणि
  \Elim{\lexists}.
\begin{enumerate}
\item शेवटचे अनुमान \Intro{\land} आहे. $\varphi$ आणि $\psi$ या आधारविधानांपासून
  $\varphi \land \psi$ निष्पन्न केले जाते आणि $\delta$ चे रूप पुढीलप्रमाणे आहे:
  \begin{prooftree}
    \AxiomC{$\Gamma_1$}
    \RightLabel{$\delta_1$}
    \DeduceC{$\varphi$}
    \AxiomC{$\Gamma_2$}
    \RightLabel{$\delta_2$}
    \DeduceC{$\psi$}
    \RightLabel{\Intro{\land}}
    \BinaryInfC{$\varphi \land \psi$}
  \end{prooftree}
  विगमन गृहीतकानुसार $\varphi$ हे $\delta_1$ च्या मुक्त न केलेली
  गृहीतकांपासून~$\Gamma_1$ निष्पन्न होते आणि $\psi$ हे $\delta_2$ च्या
  मुक्त न केलेली गृहीतकांपासून~$\Gamma_2$ निष्पन्न होते. $\delta$ ची
  मुक्त न केलेली गृहीतके $\Gamma_1 \cup \Gamma_2$ आहेत; म्हणून
  $\Gamma_1 \cup \Gamma_2 \Entails \varphi \land \psi$ हे दाखवायचे आहे.
  $\Sat{M}{\Gamma_1 \cup \Gamma_2}$ असलेली एखादी
  संरचना~$\Struct{M}$
  विचारात घ्या. $\Sat{M}{\Gamma_1}$ आणि
  $\Gamma_1 \Entails \varphi$ असल्यामुळे
  $\Sat{M}{\varphi}$ असलेच पाहिजे; तसेच
  $\Sat{M}{\Gamma_2}$ आणि
  $\Gamma_2 \Entails \psi$ असल्यामुळे
  $\Sat{M}{\psi}$. या दोन्हींवरून
  $\Sat{M}{\varphi \land \psi}$.
\item शेवटचे अनुमान \Elim{\lor} आहे: सराव.
\item शेवटचे अनुमान \Elim{\lif} आहे. $\varphi \lif \psi$ आणि $\varphi$ या
  आधारविधानांपासून $\psi$ निष्पन्न केले जाते. निष्पत्ती~$\delta$
  पुढीलप्रमाणे दिसते:
  \begin{prooftree}
    \AxiomC{$\Gamma_1$}
    \RightLabel{$\delta_1$}
    \DeduceC{$\varphi \lif \psi$}
    \AxiomC{$\Gamma_2$}
    \RightLabel{$\delta_2$}
    \DeduceC{$\varphi$}
    \RightLabel{\Elim{\lif}}
    \BinaryInfC{$\psi$}
  \end{prooftree}
  विगमन गृहीतकानुसार $\varphi \lif \psi$ हे $\delta_1$ च्या
  मुक्त न केलेली गृहीतकांपासून~$\Gamma_1$ निष्पन्न होते आणि $\varphi$ हे
  $\delta_2$ च्या मुक्त न केलेली गृहीतकांपासून~$\Gamma_2$ निष्पन्न होते.
  एखादी
  संरचना~$\Struct{M}$
  विचारात घ्या. जर $\Sat{M}{\Gamma_1 \cup
    \Gamma_2}$, तर $\Sat{M}{\psi}$, हे दाखवायचे आहे.
  $\Sat{M}{\Gamma_1 \cup \Gamma_2}$ असे समजा.
  $\Gamma_1 \Entails \varphi \lif \psi$ असल्यामुळे
  $\Sat{M}{\varphi \lif \psi}$. तसेच
  $\Gamma_2 \Entails \varphi$ असल्यामुळे
  $\Sat{M}{\varphi}$. याचा अर्थ
  $\Sat{M}{\psi}$ असा होतो (कारण
  $\Sat/{M}{\psi}$ असेल, तर
  $\Sat{M}{\varphi}$ असल्यामुळे
  $\Sat/{M}{\varphi \lif \psi}$ असे होईल, जे
  $\Sat{M}{\varphi \lif \psi}$ च्या विरुद्ध आहे).
\item शेवटचे अनुमान \Elim{\lnot} आहे: सराव.
\item शेवटचे अनुमान \Elim{\lexists} आहे: सराव.
\end{enumerate}
\end{proof}
\begin{prob}
\ref{fol:ntd:sou:thm:soundness} ची सिद्धता पूर्ण करा.
\end{prob}
\begin{cor}
\label{fol:ntd:sou:cor:weak-soundness}
जर $\Proves \varphi$, तर $\varphi$ हे वैध आहे.
\end{cor}
\begin{cor}
\label{fol:ntd:sou:cor:consistency-soundness}
जर $\Gamma$ पूर्ततायोग्य असेल, तर तो सुसंगत आहे.
\end{cor}
\begin{proof}
प्रतिलोम विधान सिद्ध करू. $\Gamma$ सुसंगत नाही असे समजा. मग
$\Gamma \Proves \lfalse$, म्हणजे $\Gamma$ मधील मुक्त न केलेली
गृहीतकांपासून $\lfalse$ ची निष्पत्ती आहे. \ref{fol:ntd:sou:thm:soundness}
नुसार $\Gamma$ पूर्ण करणारी कोणतीही
संरचना~$\Struct{M}$
$\lfalse$ देखील पूर्ण करते. प्रत्येक
संरचना~$\Struct{M}$
साठी $\Sat/{M}{\lfalse}$ असल्यामुळे कोणतीही
$\Struct{M}$ $\Gamma$ पूर्ण करू शकत नाही,
म्हणजे $\Gamma$ पूर्ततायोग्य नाही.
\end{proof}
\section{निष्पत्ती सह एकरूपता}\label{fol:ntd:ide:sec}
एकरूपता असलेल्या निष्पत्त्या साठी अतिरिक्त अनुमाननियम आवश्यक असतात.
\par\medskip\noindent
\begin{defish}
\AxiomC{}
\RightLabel{\Intro{\eq}}
\UnaryInfC{$\eq[t][t]$}
\DisplayProof
\par\medskip\noindent
\begin{tabular}{r}
\AxiomC{$\eq[t_1][t_2]$}
\AxiomC{$\varphi(t_1)$}
\RightLabel{\Elim{\eq}}
\BinaryInfC{$\varphi(t_2)$}
\DisplayProof
\\[3ex]
\AxiomC{$\eq[t_1][t_2]$}
\AxiomC{$\varphi(t_2)$}
\RightLabel{\Elim{\eq}}
\BinaryInfC{$\varphi(t_1)$}
\DisplayProof
\end{tabular}
\end{defish}
वरील नियमांमध्ये $t$, $t_1$ आणि $t_2$ ही बंद पदे आहेत. \Intro{\eq} नियम
आपल्याला कोणत्याही गृहीतकांशिवाय, थेट $\eq[t][t]$ या रूपाचे कोणतेही
एकरूपता-विधान निष्पन्न करू देतो.
\begin{ex}
$s$ आणि $t$ ही बंद पदे असतील, तर $\varphi(s), \eq[s][t] \Proves \varphi(t)$:
\begin{prooftree}
\AxiomC{$\eq[s][t]$}
\AxiomC{$\varphi(s)$}
\RightLabel{$\Elim{\eq}$}
\BinaryInfC{$\varphi(t)$}
\end{prooftree}
याला “एकरूप वस्तूंच्या प्रतिस्थापनाचे तत्त्व” किंवा लाइब्निझचा नियम
म्हणून ओळखले जाऊ शकते.
\end{ex}
\begin{prob}
$=$ हे सममित आणि संक्रमक दोन्ही आहे हे सिद्ध करा; म्हणजे
$\lforall[x][\lforall[y][(\eq[x][y] \lif
    \eq[y][x])]]$ आणि $\lforall[x][\lforall[y][\lforall[z]((\eq[x][y]
    \land \eq[y][z]) \lif \eq[x][z])]]$ यांच्या
निष्पत्त्या द्या.
\end{prob}
\begin{ex}
आपण पुढील वाक्य निष्पन्न करतो:
\begin{align*}
& \lforall[x][\lforall[y][((\varphi(x) \land \varphi(y)) \lif \eq[x][y])]]
\intertext{या वाक्य पासून:}
& \lexists[x][\lforall[y][(\varphi(y) \lif \eq[y][x])]]
\end{align*}
निष्पत्ती उलट्या दिशेने विकसित करू:
\begin{prooftree}
\AxiomC{$\lexists[x][\lforall[y][(\varphi(y) \lif \eq[y][x])]]
\quad \Discharge{\varphi(a) \land \varphi(b)}{1}$}
\DeduceC{$\eq[a][b]$}
\DischargeRule{\Intro{\lif}}{1}
\UnaryInfC{$((\varphi(a) \land \varphi(b)) \lif \eq[a][b])$}
\RightLabel{\Intro{\lforall}}
\UnaryInfC{$\lforall[y][((\varphi(a) \land \varphi(y)) \lif \eq[a][y])]$}
\RightLabel{\Intro{\lforall}}
\UnaryInfC{$\lforall[x][\lforall[y][((\varphi(x) \land \varphi(y)) \lif \eq[x][y])]]$}
\end{prooftree}
आता मुख्य गृहीतक वापरावे लागेल: ते अस्तित्ववाची सूत्र असल्यामुळे
\Elim{\lexists} वापरून मध्यवर्ती निष्कर्ष~$\eq[a][b]$ निष्पन्न करू.
\begin{prooftree}
\AxiomC{$\lexists[x][\lforall[y][(\varphi(y) \lif \eq[y][x])]]$}
\AxiomC{$\Discharge{\lforall[y][(\varphi(y) \lif \eq[y][c]])}{2}$}
\noLine
\UnaryInfC{$\Discharge{\varphi(a) \land \varphi(b)}{1}$}
\DeduceC{$\eq[a][b]$}
\DischargeRule{\Elim{\lexists}}{2}
\BinaryInfC{$\eq[a][b]$}
\DischargeRule{\Intro{\lif}}{1}
\UnaryInfC{$((\varphi(a) \land \varphi(b)) \lif \eq[a][b])$}
\RightLabel{\Intro{\lforall}}
\UnaryInfC{$\lforall[y][((\varphi(a) \land \varphi(y)) \lif \eq[a][y])]$}
\RightLabel{\Intro{\lforall}}
\UnaryInfC{$\lforall[x][\lforall[y][((\varphi(x) \land \varphi(y)) \lif \eq[x][y])]]$}
\end{prooftree}
वरच्या उजव्या बाजूची उप-निष्पत्ती तिच्या गृहीतकांचा वापर करून
$\eq[a][c]$ आणि $\eq[b][c]$ दाखवल्यावर पूर्ण होते. यासाठी दोन स्वतंत्र
निष्पत्त्या आवश्यक आहेत. $\eq[a][c]$ साठीची निष्पत्ती
पुढीलप्रमाणे आहे:
\begin{prooftree}
\AxiomC{$\Discharge{\lforall[y][(\varphi(y) \lif \eq[y][c]])}{2}$}
\RightLabel{\Elim{\lforall}}
\UnaryInfC{$\varphi(a) \lif \eq[a][c]$}
\AxiomC{$\Discharge{\varphi(a) \land \varphi(b)}{1}$}
\RightLabel{\Elim{\land}}
\UnaryInfC{$\varphi(a)$}
\RightLabel{\Elim{\lif}}
\BinaryInfC{$\eq[a][c]$}
\end{prooftree}
$\eq[a][c]$ आणि $\eq[b][c]$ यांच्यापासून \Elim{\eq} वापरून
$\eq[a][b]$ निष्पन्न करतो.
\end{ex}
\begin{prob}
पुढील सूत्रांच्या निष्पत्त्या द्या:
\begin{enumerate}
\item $\lforall[x][\lforall[y][((\eq[x][y] \land \varphi(x)) \lif \varphi(y))]]$
\item $\lexists[x][\varphi(x)] \land \lforall[y][\lforall[z][((\varphi(y) \land
    \varphi(z)) \lif \eq[y][z])]] \lif \lexists[x][(\varphi(x) \land
  \lforall[y][(\varphi(y) \lif \eq[y][x])])]$
\end{enumerate}
\end{prob}
\section{एकरूपता सह निर्दोषता}\label{fol:ntd:sid:sec}
\begin{prop}
$\eq$ साठीच्या नियमांसह नैसर्गिक निगमन निर्दोष आहे.
\end{prop}
\begin{proof}
$\eq[t][t]$ या रूपातील कोणतेही सूत्र वैध आहे, कारण प्रत्येक
संरचना~$\Struct M$ साठी $\Sat{M}{\eq[t][t]}$. (लक्षात घ्या की
$t$ हे बंद पद आहे असे आपण गृहीत धरतो, म्हणजे त्यात कोणतेही चल नाही;
म्हणून चल-विन्यास अप्रासंगिक असतात.)
निष्पत्तीमधील शेवटचे अनुमान \Elim{\eq} आहे असे समजा, म्हणजे
निष्पत्तीचे रूप पुढीलप्रमाणे आहे:
\begin{prooftree}
  \AxiomC{$\Gamma_1$}
  \RightLabel{$\delta_1$}
  \DeduceC{$\eq[t_1][t_2]$}
  \AxiomC{$\Gamma_2$}
  \RightLabel{$\delta_2$}
  \DeduceC{$\varphi(t_1)$}
  \RightLabel{\Elim{\eq}}
  \BinaryInfC{$\varphi(t_2)$}
\end{prooftree}
आधारविधाने~$\eq[t_1][t_2]$ आणि $\varphi(t_1)$ ही अनुक्रमे
मुक्त न केलेली गृहीतके~$\Gamma_1$ आणि $\Gamma_2$ यांच्यापासून
निष्पन्न केली आहेत. $\varphi(t_2)$ हे $\Gamma_1 \cup \Gamma_2$ पासून
निष्पन्न होते, हे दाखवायचे आहे. $\Sat{M}{\Gamma_1 \cup
  \Gamma_2}$ असलेली संरचना~$\Struct{M}$ विचारात घ्या. विगमन
गृहीतकानुसार $\Sat{M}{\varphi(t_1)}$ आणि $\Sat{M}{\eq[t_1][t_2]}$.
त्यामुळे $\Value{t_1}{M} = \Value{t_2}{M}$. $s$ हा कोणताही चल-विन्यास
आणि $m = \Value{t_1}{M} = \Value{t_2}{M}$ असू द्या.
\ref{fol:syn:ext:prop:ext-formulas} नुसार
$\Sat{M}{\varphi(t_1)}[s]$ तेव्हा आणि केवळ तेव्हाच, जेव्हा
$\Sat{M}{\varphi(x)}[\Subst{s}{m}{x}]$, आणि हे तेव्हा आणि केवळ तेव्हाच, जेव्हा
$\Sat{M}{\varphi(t_2)}[s]$. $\Sat{M}{\varphi(t_1)}$ असल्यामुळे
$\Sat{M}{\varphi(t_2)}$.
\end{proof}
\chapter{टॅब्लो}\label{fol:tab::chap}\begingroup\setlength{\emergencystretch}{3em}\NewDocumentCommand{\sFmlaWide}{m m o}{\ensuremath{\IfNoValueTF{#3}{}{#3\,}\hbox to 1.2em{\ensuremath{#1}\hfil} #2}}
\begin{editorial}
हे प्रकरण चिन्हांकित सूत्रांवर आधारलेली विश्लेषणात्मक टॅब्लो पद्धत मांडते.
टॅब्लोशी सिद्धता-पद्धती म्हणून संबंधित सामग्री समाविष्ट किंवा
वगळण्यासाठी ``prfTab'' टॅग वापरा.
\end{editorial}
\section{नियम आणि टॅब्लो}\label{fol:tab:rul:sec}
टॅब्लो म्हणजे वाक्य संरचना मध्ये ज्या संभाव्य प्रकारे
सत्य अथवा असत्य असू शकते त्यांचा पद्धतशीर आढावा होय. टॅब्लोचे मूलभूत
घटक चिन्हांकित सूत्रे असतात: वाक्ये सोबत सत्यतामूल्याचे
``चिन्ह,'' म्हणजे $\True$ किंवा~$\False$. ही चिन्हांकित सूत्रे
(खाली वाढणाऱ्या) वृक्षात मांडली जातात.
\begin{defn}
  एक \emph{चिन्हांकित सूत्र} म्हणजे सत्यतामूल्य आणि वाक्य यांची जोडी,
  म्हणजे पुढीलपैकी एक:
  \[
  \sFmla{\True}{\varphi} \text{ किंवा } \sFmla{\False}{\varphi}.
  \]
\end{defn}

सहज अर्थानुसार, $\sFmla{\True}{\varphi}$ चे वाचन आपण ``$\varphi$ सत्य असू
शकते'' आणि $\sFmla{\False}{\varphi}$ चे वाचन ``$\varphi$ असत्य असू शकते''
(एखाद्या संरचना मध्ये) असे करू शकतो.

वृक्षातील प्रत्येक चिन्हांकित सूत्र एकतर \emph{गृहीतक} असते (ती
वृक्षाच्या अगदी वरच्या टोकाला नोंदवलेली असतात), किंवा तिच्या वर असलेल्या
चिन्हांकित सूत्राला अनेक निगमन नियमांपैकी एखादा नियम लावून ती मिळते.
पूर्ववर्ती सूत्राच्या प्रत्येक संभाव्य मुख्य संकारक साठी दोन नियम
असतात: चिन्ह~$\True$ असेल त्या स्थितीसाठी एक आणि चिन्ह~$\False$ असेल त्या
स्थितीसाठी एक. काही नियम वृक्षाची शाखा फोडू देतात, तर काही नियम शाखेत फक्त
चिन्हांकित सूत्रे जोडतात. एखादा नियम फक्त अगदी आधी येणाऱ्या
चिन्हांकित सूत्र वरच लागू करावा असे नाही; मुळापासून नियम लागू करण्याच्या
स्थानापर्यंतच्या शाखेतील कोणत्याही चिन्हांकित सूत्र वर तो लागू करता येतो
(आणि अनेकदा आधीच्या एखाद्यावर लावावाच लागतो).

एखाद्या शाखेत $\sFmla{\True}{\varphi}$ आणि $\sFmla{\False}{\varphi}$ ही दोन्ही
असतील तेव्हा ती शाखा \emph{बंद} असते. ज्याची प्रत्येक शाखा बंद असते तो
टॅब्लो बंद असतो. सहज अर्थानुसार, कोणतीही शाखा एक संयुक्त शक्यता
वर्णन करते; परंतु $\sFmla{\True}{\varphi}$ आणि $\sFmla{\False}{\varphi}$ एकाच वेळी
शक्य नाहीत. दुसऱ्या शब्दांत, एखादी शाखा बंद असेल तर तिने वर्णन केलेली शक्यता
बाद होते. विशेषतः याचा अर्थ असा की, बंद टॅब्लो हा
$\sFmla{\True}{\varphi}$ स्वरूपाचे प्रत्येक गृहीतक सत्य आणि
$\sFmla{\False}{\varphi}$ स्वरूपाचे प्रत्येक गृहीतक असत्य करणाऱ्या सर्व शक्यता
बाद करतो.

बंद टॅब्लो \emph{$\varphi$ साठीचा} म्हणजे ज्याचे
मूळ~$\sFmla{\False}{\varphi}$ आहे असा बंद टॅब्लो. असा बंद टॅब्लो
अस्तित्वात असल्यास, $\varphi$ असत्य असण्याच्या सर्व शक्यता बाद होतात; म्हणजेच
$\varphi$ प्रत्येक संरचना मध्ये सत्य असले पाहिजे.












\section{विधानीय नियम}\label{fol:tab:prl:sec}

\subsection{$\lnot$ साठीचे नियम}

\begin{defish}
\AxiomC{\sFmla{\True}{\lnot \varphi}}
\RightLabel{\TRule{\True}{\lnot}}
\UnaryInfC{\sFmla{\False}{\varphi}}
\DisplayProof
\hfill
\AxiomC{\sFmla{\False}{\lnot \varphi}}
\RightLabel{\TRule{\False}{\lnot}}
\UnaryInfC{\sFmla{\True}{\varphi}}
\DisplayProof
\end{defish}

\subsection{$\land$ साठीचे नियम}

\begin{defish}\noindent
\AxiomC{\sFmla{\True}{\varphi \land \psi}}
\RightLabel{\TRule{\True}{\land}}
\UnaryInfC{\sFmla{\True}{\varphi}}
\noLine
\UnaryInfC{\sFmla{\True}{\psi}}
\DisplayProof
\hfill
\AxiomC{\sFmla{\False}{\varphi \land \psi}}
\RightLabel{\TRule{\False}{\land}}
\UnaryInfC{$\sFmla{\False}{\varphi} \quad \mid \quad \sFmla{\False}{\psi}$}
\DisplayProof
\end{defish}

\subsection{$\lor$ साठीचे नियम}

\begin{defish}
\AxiomC{\sFmla{\True}{\varphi \lor \psi}}
\RightLabel{\TRule{\True}{\lor}}
\UnaryInfC{$\sFmla{\True}{\varphi} \quad \mid \quad \sFmla{\True}{\psi}$}
\DisplayProof
\hfill
\AxiomC{\sFmla{\False}{\varphi \lor \psi}}
\RightLabel{\TRule{\False}{\lor}}
\UnaryInfC{\sFmla{\False}{\varphi}}
\noLine
\UnaryInfC{\sFmla{\False}{\psi}}
\DisplayProof
\end{defish}

\subsection{$\lif$ साठीचे नियम}

\begin{defish}
\AxiomC{\sFmla{\True}{\varphi \lif \psi}}
\RightLabel{\TRule{\True}{\lif}}
\UnaryInfC{$\sFmla{\False}{\varphi} \quad \mid \quad \sFmla{\True}{\psi}$}
\DisplayProof
\hfill
\AxiomC{\sFmla{\False}{\varphi \lif \psi}}
\RightLabel{\TRule{\False}{\lif}}
\UnaryInfC{\sFmla{\True}{\varphi}}
\noLine
\UnaryInfC{\sFmla{\False}{\psi}}
\DisplayProof
\end{defish}

\subsection{कट नियम}

\begin{defish}
\AxiomC{}
\RightLabel{\Cut}
\UnaryInfC{$\sFmla{\True}{\varphi} \quad \mid \quad \sFmla{\False}{\varphi}$}
\DisplayProof
\end{defish}

\Cut{} नियम पूर्वीच्या चिन्हांकित सूत्र वर ``लावला'' जात नाही; त्याऐवजी
तो टॅब्लो मधील प्रत्येक शाखा दोन शाखांत फोडू देतो: एका शाखेत
$\sFmla{\True}{\varphi}$ आणि दुसरीत~$\sFmla{\False}{\varphi}$ असते. हा नियम आवश्यक
नाही---ज्या चिन्हांकित सूत्रांच्या संचासाठी बंद टॅब्लो आहे, त्या
संचासाठी \Cut{} न वापरणाराही एक बंद टॅब्लो असतो---परंतु त्यामुळे
टॅब्लो सोयीस्कर रीतीने एकत्र जोडता येतात.












\section{संख्यापकांचे नियम}\label{fol:tab:qrl:sec}

\subsection{$\lforall$ साठीचे नियम}

\begin{defish}
\AxiomC{\sFmla{\True}{\lforall[x][\varphi(x)]}}
\RightLabel{\TRule{\True}{\forall}}
\UnaryInfC{\sFmla{\True}{\varphi(t)}}
\DisplayProof
\hfill
\AxiomC{\sFmla{\False}{\lforall[x][\varphi(x)]}}
\RightLabel{\TRule{\False}{\lforall}}
\UnaryInfC{\sFmla{\False}{\varphi(a)}}
\DisplayProof
\end{defish}

\TRule{\True}{\lforall} मध्ये, $t$ हे बंद पद आहे (म्हणजे, चरे नसलेले
पद). \TRule{\False}{\lforall} मध्ये, $a$~हा असा स्थिरांक आहे की तो
\TRule{\False}{\lforall} नियमाच्या वरच्या शाखेत कोठेही येता कामा नये.
\TRule{\False}{\forall} अनुमानाचा $a$ हा \emph{आयगेन चर} आहे, असे आपण
म्हणतो.\footnote{वरील नियमात $a$ हा स्थिरांक असला तरी आपण
``आयगेन चर'' ही संज्ञा वापरतो. यामागे ऐतिहासिक कारणे आहेत.}

\subsection{$\lexists$ साठीचे नियम}

\begin{defish}
\AxiomC{\sFmla{\True}{\lexists[x][\varphi(x)]}}
\RightLabel{\TRule{\True}{\lexists}}
\UnaryInfC{\sFmla{\True}{\varphi(a)}}
\DisplayProof
\hfill
\AxiomC{\sFmla{\False}{\lexists[x][\varphi(x)]}}
\RightLabel{\TRule{\False}{\lexists}}
\UnaryInfC{\sFmla{\False}{\varphi(t)}}
\DisplayProof
\end{defish}

पुन्हा, $t$~हे बंद पद आहे आणि $a$~हा असा स्थिरांक आहे की तो
\TRule{\True}{\lexists} नियमाच्या वरच्या शाखेत येत नाही. आपण $a$ ला
\TRule{\True}{\lexists} अनुमानाचा \emph{आयगेन चर} म्हणतो.

\TRule{\False}{\lforall} किंवा \TRule{\True}{\lexists} अनुमानाच्या वरच्या
शाखेत आयगेन चर येत नसावा, या अटीला \emph{आयगेन चराची अट} म्हणतात.

\begin{explain}
ही प्रथा आठवा: $\varphi$ हे असे सूत्र असेल की त्यात
चल~$x$ मुक्त असेल, तर हे आपण~$\varphi(x)$ असे लिहून दर्शवतो.
त्याच संदर्भात, मग $\varphi(t)$ हा~$\Subst{\varphi}{t}{x}$ चा संक्षेप असतो.
त्यामुळे \TRule{\False}{\lexists} नियम पुढीलप्रमाणेही लिहिता येईल:
\begin{prooftree}
  \AxiomC{\sFmla{\False}{\lexists[x][\varphi]}}
  \RightLabel{\TRule{\False}{\lexists}}
  \UnaryInfC{\sFmla{\False}{\Subst{\varphi}{t}{x}}}
\end{prooftree}
लक्षात घ्या की~$\varphi$ मध्ये $t$ आधीच येऊ शकते; उदा., $\varphi$ हे
~$\Atom{\Obj P}{t,x}$ असू शकते. म्हणून,
$\sFmla{\False}{\lexists[x][\Atom{\Obj P}{t,x}]}$ पासून
$\sFmla{\False}{\Atom{\Obj P}{t,t}}$ निष्पन्न करणे हा
\TRule{\False}{\lexists} चा योग्य उपयोग आहे. परंतु
\TRule{\False}{\lforall} आणि~\TRule{\True}{\lexists} मधील आयगेन चराच्या
अटींनुसार स्थिरांक~$a$ हा~$\varphi$ मध्ये येता कामा नये. म्हणून,
$\sFmla{\False}{\lforall[x][\Atom{\Obj P}{a,x}]}$ पासून
$\sFmla{\False}{\Atom{\Obj P}{a,a}}$ हे
~$\TRule{\False}{\lforall}$ वापरून योग्य रीतीने निष्पन्न करता येत नाही.
\end{explain}

\begin{explain}
\TRule{\True}{\lforall} आणि \TRule{\False}{\lexists} मध्ये पद~$t$ बंद
असण्याव्यतिरिक्त त्यावर आणखी कोणतेही निर्बंध नाहीत. याउलट,
\TRule{\True}{\lexists} आणि \TRule{\False}{\lforall} नियमांमध्ये आयगेन
चराच्या अटीनुसार संबंधित अनुमानांच्या वरच्या शाखांमध्ये स्थिरांक~$a$
कोठेही येता कामा नये. ही पद्धत निर्दोष राहील याची खात्री करण्यासाठी ही अट
आवश्यक आहे. ही अट नसती, तर पुढील आकृती
$\lexists[x][\varphi(x)] \lif \lforall[x][\varphi(x)]$ साठीचा
बंद टॅब्लो ठरली असती:
\begin{center}
\begin{tableau}{}
  [\sFmla{\False}{\lexists[x][\varphi(x)] \lif \lforall[x][\varphi(x)]}, just=\TAss
    [\sFmla{\True}{\lexists[x][\varphi(x)]},
      just={\TRule{\False}{\lif}[1]}
      [\sFmla{\False}{\lforall[x][\varphi(x)]},
        just={\TRule{\False}{\lif}[1]}
        [\sFmla{\True}{\varphi(a)},
          just={\TRule{\True}{\lexists}[2]}
          [\sFmla{\False}{\varphi(a)},
            just={\TRule{\False}{\lforall}[3]}, close]
        ]
      ]
    ]
  ]
\end{tableau}
\end{center}
परंतु $\lexists[x][\varphi(x)] \lif
\lforall[x][\varphi(x)]$ हे वैध नाही.
\end{explain}












\section{टॅब्लो}\label{fol:tab:der:sec}

\begin{explain}
गृहीतक म्हणजे काय हे आपण सांगितले आहे आणि निगमन नियम दिले आहेत.
टॅब्लो यांपासून विगमनाधारित रीतीने निर्माण होतात: प्रत्येक
टॅब्लो ही एकतर एक किंवा अधिक गृहीतकांनी बनलेली एकच शाखा असते, किंवा
टॅब्लोच्या एखाद्या शाखेवर निगमन नियमांपैकी एक नियम लावल्याने ती मिळते.
\end{explain}

\begin{defn}[टॅब्लो]
$\sFmlaWide{S_1}{\varphi_1}$, \dots,\linebreak $\sFmlaWide{S_n}{\varphi_n}$ या गृहीतकांसाठीचा
टॅब्लो (येथे प्रत्येक $S_i$ हे $\True$ किंवा~$\False$ यांपैकी एक
आहे) म्हणजे पुढील अटी पूर्ण करणारा चिन्हांकित सूत्रांचा सांत वृक्ष होय:
\begin{enumerate}
\item वृक्षातील सर्वांत वरची $n$ चिन्हांकित सूत्रे म्हणजे एकाखाली एक
  असलेली\linebreak $\sFmlaWide{S_i}{\varphi_i}$ ही सूत्रे होत.
\item वृक्षातील गृहीतक नसलेली प्रत्येक चिन्हांकित सूत्र तिच्या वरच्या
  शाखेतील चिन्हांकित सूत्राला निगमन नियम योग्य रीतीने लावल्याने मिळते.
\end{enumerate}
टॅब्लोची एखादी शाखा $\sFmla{\True}{\varphi}$ आणि
~$\sFmla{\False}{\varphi}$ ही दोन्ही सूत्रे धारण करत असेल तर आणि तरच ती
\emph{बंद} असते; अन्यथा ती \emph{खुली} असते. ज्याची प्रत्येक शाखा बंद
असते तो टॅब्लो (त्याच्या गृहीतकांच्या संचासाठीचा)
\emph{बंद टॅब्लो} असतो. टॅब्लो बंद नसेल, म्हणजे त्यात किमान
एक खुली शाखा असेल, तर तो \emph{खुला} असतो.
\end{defn}

\begin{ex}
  गृहीतकांचा प्रत्येक संच स्वतःच टॅब्लो असतो, परंतु साधारणतः तो बंद
  नसतो. (स्पष्टपणे, तो बंद असेल तर गृहीतकांत $\sFmla{\True}{\varphi}$ आणि
  ~$\sFmla{\False}{\varphi}$ ही चिन्हांकित सूत्रांची जोडी आधीच असली पाहिजे.)

  एखाद्या टॅब्लो मधील चिन्हांकित सूत्र~$\varphi$ ला निगमन नियमांपैकी
  एक नियम लावून त्या टॅब्लोपासून (तो खुला असो वा बंद) नवा आणि अधिक मोठा
  टॅब्लो मिळवता येतो. नियम लागू केलेली~$\varphi$ ची आस्थिती ज्या ज्या शाखेत
  आहे, त्या प्रत्येक शाखेच्या शेवटी तो नियम एक किंवा अधिक
  चिन्हांकित सूत्रे जोडतो.

  उदाहरणार्थ, $\sFmla{\True}{\varphi \land \lnot \varphi}$ हे गृहीतक विचारात घ्या.
  केवळ त्या गृहीतकाने बनलेला (खुला) टॅब्लो पुढीलप्रमाणे आहे:
  \begin{center}
    \begin{tableau}{}
      [\sFmla{\True}{\varphi \land \lnot \varphi}, just=\TAss]
    \end{tableau}{}
  \end{center}
  त्या गृहीतकाला $\TRule{\True}{\land}$ नियम लावून त्यापासून नवा
  टॅब्लो मिळतो. हा नियम टॅब्लोमध्ये $\sFmla{\True}{\varphi}$ आणि
  $\sFmla{\True}{\lnot \varphi}$ या दोन नव्या ओळी जोडू देतो:
  \begin{center}
    \begin{tableau}{}
      [\sFmla{\True}{\varphi \land \lnot \varphi}, just=\TAss
        [\sFmla{\True}{\varphi}, just={\TRule{\True}{\land}[1]},
          [\sFmla{\True}{\lnot \varphi}, just={\TRule{\True}{\land}[1]}
          ]
        ]
      ]
    \end{tableau}{}
  \end{center}
  टॅब्लो लिहिताना आपण लावलेले नियम उजवीकडे नोंदवतो (उदा.,
  $\TRule{\True}{\land} 1$ याचा अर्थ त्या ओळीवरील चिन्हांकित सूत्र ही
  ओळ~$1$ वरील चिन्हांकित सूत्राला $\TRule{\True}{\land}$ नियम लावल्याने
  मिळाली आहे). या नव्या टॅब्लोमध्ये आता अधिक चिन्हांकित सूत्रे
  आहेत, परंतु त्यांपैकी फक्त एकीला ($\sFmla{\True}{\lnot \varphi}$) नियम लावता
  येतो (या स्थितीत $\TRule{\True}{\lnot}$ नियम). त्यामुळे पुढील बंद
  टॅब्लो मिळतो:
  \begin{center}
    \begin{tableau}{}
      [\sFmla{\True}{\varphi \land \lnot \varphi}, just=\TAss
        [\sFmla{\True}{\varphi}, just={\TRule{\True}{\land}[1]}
          [\sFmla{\True}{\lnot \varphi}, just={\TRule{\True}{\land}[1]}
            [\sFmla{\False}{\varphi}, just={\TRule{\True}{\lnot}[3]}, close]
          ]
        ]
      ]
    \end{tableau}{}
  \end{center}
\end{ex}












\section{टॅब्लोची उदाहरणे}\label{fol:tab:pro:sec}

\begin{ex}
$(\varphi \land \psi) \lif \varphi$ या वाक्य साठी बंद टॅब्लो शोधू या.

संबंधित गृहीतक टॅब्लोच्या सर्वांत वर लिहून आपण सुरुवात करतो.
\begin{oltableau}
  [\sFmla{\False}{(\varphi \land \psi) \lif \varphi},
    just = \TAss]
\end{oltableau}

एकच गृहीतक असल्यामुळे ज्याला नियम लावता येईल अशी एकच चिन्हांकित सूत्र
आहे. (प्रत्येक चिन्हांकित सूत्राला नेहमी जास्तीत जास्त एकच नियम लागू
करता येतो: तो म्हणजे वाक्याच्या संबंधित चिन्हाचा आणि
मुख्य संकारक चा नियम.) या स्थितीत याचा अर्थ असा की आपण
$\TRule{\False}{\lif}$ लावला पाहिजे.
\begin{oltableau}
  [\sFmla{\False}{(\varphi \land \psi) \lif \varphi},
    checked, just = \TAss
    [\sFmla{\True}{\varphi \land \psi},
      just={\TRule{\False}{\lif}[1]}
      [\sFmla{\False}{\varphi}, just={\TRule{\False}{\lif}[1]}]
    ]
  ]
\end{oltableau}
कोणत्या चिन्हांकित सूत्रे ना त्यांचे संबंधित नियम लावले आहेत याची नोंद
ठेवण्यासाठी आपण त्या वाक्याशेजारी खूण करतो. मात्र नियम सर्व खुल्या शाखांवर
लावला असेल \emph{तरच} खूण करा. एखाद्या चिन्हांकित सूत्राला प्रत्येक
खुल्या शाखेत संबंधित नियम लावल्यानंतर तिच्याकडे परत जाऊन तो नियम पुन्हा
लावावा लागणार नाही. येथे एकच शाखा असल्यामुळे नियम एकदाच लावावा लागतो.
(या खुणा केवळ टॅब्लो तयार करण्याची सोय आहेत आणि टॅब्लोच्या विन्यासाचा
अधिकृत भाग नाहीत, हे लक्षात घ्या.)

नियम लावता येईल अशी एक नवी चिन्हांकित सूत्र आहे: ओळ~$2$ वरील
$\sFmla{\True}{\varphi \land \psi}$. $\TRule{\True}{\land}$ नियम लावल्यावर
पुढील परिणाम मिळतो:
\begin{oltableau}
  [\sFmla{\False}{(\varphi \land \psi) \lif \varphi},
    checked, just = \TAss
    [\sFmla{\True}{\varphi \land \psi},
      just={\TRule{\False}{\lif}[1]}, checked
      [\sFmla{\False}{\varphi}, just={\TRule{\False}{\lif}[1]}
        [\sFmla{\True}{\varphi}, just={\TRule{\True}{\land}[2]}
          [\sFmla{\True}{\psi}, just={\TRule{\True}{\land}[2]}, close
          ]
        ]
      ]
    ]
  ]
\end{oltableau}
शाखेत आता $\sFmla{\True}{\varphi}$ (ओळ~$4$ वर) आणि
$\sFmla{\False}{\varphi}$ (ओळ~$3$ वर) ही दोन्ही असल्यामुळे शाखा बंद आहे.
ती एकमेव शाखा असल्यामुळे टॅब्लो बंद आहे. आपण
~$(\varphi \land \psi) \lif \varphi$ साठी बंद टॅब्लो शोधला आहे.
\end{ex}

\begin{ex}
आता $(\lnot \varphi \lor \psi) \lif (\varphi \lif \psi)$ साठी बंद टॅब्लो शोधू या.

संबंधित गृहीतकाने आपण सुरुवात करतो:
\begin{oltableau}
  [\sFmla{\False}{(\lnot \varphi \lor \psi) \lif
      (\varphi \lif \psi)}, just=\TAss]
\end{oltableau}
या टॅब्लो मधील एकमेव चिन्हांकित सूत्राचा मुख्य संकारक~$\lif$
आणि चिन्ह~$\False$ आहे; म्हणून तिला $\TRule{\False}{\lif}$ नियम लावल्यावर
पुढील परिणाम मिळतो:
\begin{oltableau}
  [\sFmla{\False}{(\lnot \varphi \lor \psi)
      \lif (\varphi \lif \psi)}, just=\TAss, checked
    [\sFmla{\True}{\lnot \varphi \lor \psi},
      just={\TRule{\False}{\lif}[1]}
      [\sFmla{\False}{(\varphi \lif \psi)},
        just={\TRule{\False}{\lif}[1]}
      ]
    ]
  ]
\end{oltableau}
आता ओळ~$2$ ला~$\TRule{\True}{\lor}$ लावायचा की ओळ~$3$ ला
$\TRule{\False}{\lif}$ लावायचा, असा पर्याय आपल्यापुढे आहे. प्रत्येक
चिन्हांकित सूत्राला तिचा संबंधित नियम प्रत्येक शाखेत लावला, तर आपण कोणता
क्रम निवडतो याने प्रत्यक्षात फरक पडत नाही. म्हणून पहिला पर्याय निवडू या.
$\TRule{\True}{\lor}$ नियमामुळे टॅब्लोची शाखा फुटू शकते आणि नियमाचे
दोन निष्कर्ष त्या दोन नव्या शाखांना जोडलेल्या नव्या चिन्हांकित सूत्रे
असतील. त्यामुळे पुढील परिणाम मिळतो:
\begin{oltableau}
  [\sFmla{\False}{(\lnot \varphi \lor \psi) \lif
      (\varphi \lif \psi)}, just=\TAss, checked
    [\sFmla{\True}{\lnot \varphi \lor \psi},
      just={\TRule{\False}{\lif}[1]}, checked
      [\sFmla{\False}{(\varphi \lif \psi)},
        just={\TRule{\False}{\lif}[1]}
        [\sFmla{\True}{\lnot \varphi}, just={\TRule{\True}{\lor}[2]}]
        [\sFmla{\True}{\psi}, just={\TRule{\True}{\lor}[2]}]
      ]
    ]
  ]
\end{oltableau}
ओळ~$3$ ला आपण अद्याप $\TRule{\False}{\lif}$ नियम लावलेला नाही; तो आता
लावू या. वेळ वाचवण्यासाठी तो दोन्ही शाखांवर लावू. प्रत्येक खुल्या शाखेत
संबंधित नियम लावला असेल तरच चिन्हांकित सूत्र शेजारी खूण करतो, हे आठवा.
म्हणून नियम ज्याला लागू होतो ती चिन्हांकित सूत्र असलेल्या प्रत्येक
शाखेच्या शेवटी तो नियम लावणे उचित ठरते. त्यामुळे विविध शाखांत खाली जाऊन त्या
चिन्हांकित सूत्र कडे पुन्हा परतावे लागत नाही.
\begin{oltableau}
  [\sFmla{\False}{(\lnot \varphi \lor \psi) \lif
      (\varphi \lif \psi)}, just=\TAss, checked
    [\sFmla{\True}{\lnot \varphi \lor \psi},
      just={\TRule{\False}{\lif}[1]}, checked
      [\sFmla{\False}{(\varphi \lif \psi)},
        just={\TRule{\False}{\lif}[1]}, checked
        [\sFmla{\True}{\lnot \varphi}, just={\TRule{\True}{\lor}[2]}
          [\sFmla{\True}{\varphi}, just={\TRule{\False}{\lif}[3]}
            [\sFmla{\False}{\psi}, just={\TRule{\False}{\lif}[3]}]
          ]
        ]
        [\sFmla{\True}{\psi}, just={\TRule{\True}{\lor}[2]}
          [\sFmla{\True}{\varphi}, just={\TRule{\False}{\lif}[3]}
            [\sFmla{\False}{\psi}, just={\TRule{\False}{\lif}[3]}, close]
          ]
        ]
      ]
    ]
  ]
\end{oltableau}
उजवी शाखा आता बंद आहे. डाव्या शाखेत आपण अजूनही ओळ~$4$ ला
$\TRule{\True}{\lnot}$ नियम लावू शकतो. त्यामुळे
~$\sFmla{\False}{\varphi}$ मिळते आणि डावी शाखा बंद होते:
\begin{oltableau}
  [\sFmla{\False}{(\lnot \varphi \lor \psi) \lif
      (\varphi \lif \psi)}, just=\TAss, checked
    [\sFmla{\True}{\lnot \varphi \lor \psi},
      just={\TRule{\False}{\lif}[1]}, checked
      [\sFmla{\False}{(\varphi \lif \psi)},
        just={\TRule{\False}{\lif}[1]}, checked
        [\sFmla{\True}{\lnot \varphi}, just={\TRule{\True}{\lor}[2]}
          [\sFmla{\True}{\varphi}, just={\TRule{\False}{\lif}[3]}
            [\sFmla{\False}{\psi}, just={\TRule{\False}{\lif}[3]}
              [\sFmla{\False}{\varphi},
                just={\TRule{\True}{\lnot}[4]}, close]
            ]
          ]
        ]
        [\sFmla{\True}{\psi}, just={\TRule{\True}{\lor}[2]}
          [\sFmla{\True}{\varphi}, just={\TRule{\False}{\lif}[3]}
            [\sFmla{\False}{\psi},
              just={\TRule{\False}{\lif}[3]}, close]
          ]
        ]
      ]
    ]
  ]
\end{oltableau}
\end{ex}

\begin{ex}
कितीही चिन्हांकित सूत्रे गृहीतके म्हणून घेऊन त्यांच्यासाठी टॅब्लो
देता येतात. शाखा फोडू देणारे एकाहून अधिक नियम लावणेही अनेकदा आवश्यक असते;
आणि सर्वसाधारणपणे टॅब्लोमध्ये कितीही शाखा असू शकतात. उदाहरणार्थ,
पुढील सूत्रांसाठीचा टॅब्लो विचारात घ्या:
$\{\sFmla{\True}{\varphi \lor (\psi \land \chi)}, \sFmla{\False}{(\varphi \lor \psi) \land (\varphi
  \lor \chi)}\}$. पहिल्या गृहीतकाला $\TRule{\True}{\lor}$ लावून आपण
सुरुवात करतो:
\begin{oltableau}
  [\sFmla{\True}{\varphi \lor (\psi \land \formula{C})},
    just=\TAss, checked
    [\sFmla{\False}{(\varphi \lor \psi) \land
        (\varphi \lor \formula{C})}, just=\TAss
      [\sFmla{\True}{\varphi}, just={\TRule{\True}{\lor}[1]}]
      [\sFmla{\True}{\psi \land \formula{C}},
        just={\TRule{\True}{\lor}[1]}]
    ]
  ]
\end{oltableau}
आता ओळ~$2$ ला $\TRule{\False}{\land}$ नियम लावता येतो. दोन्ही शाखांवर
तो एकाच वेळी लावल्यामुळे आपण ओळ~$2$ वर खूण करू शकतो:
\begin{oltableau}
  [\sFmla{\True}{\varphi \lor (\psi \land \formula{C})},
    just=\TAss, checked
    [\sFmla{\False}{(\varphi \lor \psi) \land
        (\varphi \lor \formula{C})}, just=\TAss, checked
      [\sFmla{\True}{\varphi}, just={\TRule{\True}{\lor}[1]}
        [\sFmla{\False}{\varphi \lor \psi},
          just={\TRule{\False}{\land}[2]}]
        [\sFmla{\False}{\varphi \lor \formula{C}},
          just={\TRule{\False}{\land}[2]}]
      ]
      [\sFmla{\True}{\psi \land \formula{C}},
        just={\TRule{\True}{\lor}[1]}
        [\sFmla{\False}{\varphi \lor \psi},
          just={\TRule{\False}{\land}[2]}]
        [\sFmla{\False}{\varphi \lor \formula{C}},
          just={\TRule{\False}{\land}[2]}]
      ]
    ]
  ]
\end{oltableau}
आता $\varphi \lor \psi$ असलेल्या सर्व शाखांवर $\TRule{\False}{\lor}$ लावता
येतो:
\begin{oltableau}
  [\sFmla{\True}{\varphi \lor (\psi \land \formula{C})},
    just=\TAss, checked
    [\sFmla{\False}{(\varphi \lor \psi) \land
        (\varphi \lor \formula{C})}, just=\TAss, checked
      [\sFmla{\True}{\varphi}, just={\TRule{\True}{\lor}[1]}
        [\sFmla{\False}{\varphi \lor \psi},
          just={\TRule{\False}{\land}[2]}, checked
          [\sFmla{\False}{\varphi}, just={\TRule{\False}{\lor}[4]}
            [\sFmla{\False}{\psi},
              just={\TRule{\False}{\lor}[4]}, close]
          ]
        ]
        [\sFmla{\False}{\varphi \lor \formula{C}},
          just={\TRule{\False}{\land}[2]}]
      ]
      [\sFmla{\True}{\psi \land \formula{C}},
        just={\TRule{\True}{\lor}[1]}
        [\sFmla{\False}{\varphi \lor \psi},
          just={\TRule{\False}{\land}[2]}, checked
          [\sFmla{\False}{\varphi}, just={\TRule{\False}{\lor}[4]}
            [\sFmla{\False}{\psi},
              just={\TRule{\False}{\lor}[4]}]
          ]
        ]
        [\sFmla{\False}{\varphi \lor \formula{C}},
          just={\TRule{\False}{\land}[2]}]
      ]
    ]
  ]
\end{oltableau}
सर्वांत डावी शाखा आता बंद आहे. आता $\varphi \lor \chi$ ला
$\TRule{\False}{\lor}$ लावू या:
\begin{oltableau}
  [\sFmla{\True}{\varphi \lor (\psi \land \formula{C})},
    just=\TAss, checked
    [\sFmla{\False}{(\varphi \lor \psi) \land
        (\varphi \lor \formula{C})}, just=\TAss, checked
      [\sFmla{\True}{\varphi}, just={\TRule{\True}{\lor}[1]}
        [\sFmla{\False}{\varphi \lor \psi},
          just={\TRule{\False}{\land}[2]}, checked
          [\sFmla{\False}{\varphi},
            just={\TRule{\False}{\lor}[4]}
            [\sFmla{\False}{\psi},
              just={\TRule{\False}{\lor}[4]}, close]
          ]
        ]
        [\sFmla{\False}{\varphi \lor \formula{C}},
          just={\TRule{\False}{\land}[2]}, checked
          [\sFmla{\False}{\varphi},
            just={\TRule{\False}{\lor}[4]}, move by=2
            [\sFmla{\False}{\formula{C}},
              just={\TRule{\False}{\lor}[4]}, close]
          ]
        ]
      ]
      [\sFmla{\True}{\psi \land \formula{C}},
        just={\TRule{\True}{\lor}[1]}
        [\sFmla{\False}{\varphi \lor \psi},
          just={\TRule{\False}{\land}[2]}, checked
          [\sFmla{\False}{\varphi},
            just={\TRule{\False}{\lor}[4]}
            [\sFmla{\False}{\psi},
              just={\TRule{\False}{\lor}[4]}
            ]
          ]
        ]
        [\sFmla{\False}{\varphi \lor \formula{C}},
          just={\TRule{\False}{\land}[2]}, checked
          [\sFmla{\False}{\varphi},
            just={\TRule{\False}{\lor}[4]}, move by=2
            [\sFmla{\False}{\formula{C}},
              just={\TRule{\False}{\lor}[4]}]
          ]
        ]
      ]
    ]
  ]
\end{oltableau}
स्पष्टतेसाठी $\TRule{\False}{\lor}$ दुसऱ्यांदा लावल्याचा परिणाम आपण खाली
हलवला आहे, हे लक्षात घ्या. येथे समर्थन समान असते, त्यामुळे तसे करण्याची
गरज नव्हती.

दोन शाखा खुल्या राहतात आणि ओळ~$3$ वरील $\sFmla{\True}{\psi \land \chi}$
वर खूण झालेली नाही. तिला $\TRule{\True}{\land}$ लावून पुढील बंद
टॅब्लो मिळतो:
\begin{oltableau}
  [\sFmla{\True}{\varphi \lor (\psi \land \formula{C})},
    just=\TAss, checked
    [\sFmla{\False}{(\varphi \lor \psi) \land
        (\varphi \lor \formula{C})}, just=\TAss, checked
      [\sFmla{\True}{\varphi}, just={\TRule{\True}{\lor}[1]}
        [\sFmla{\False}{\varphi \lor \psi},
          just={\TRule{\False}{\land}[2]}, checked
          [\sFmla{\False}{\varphi},
            just={\TRule{\False}{\lor}[4]}
            [\sFmla{\False}{\psi},
              just={\TRule{\False}{\lor}[4]}, close]
          ]
        ]
        [\sFmla{\False}{\varphi \lor \formula{C}},
          just={\TRule{\False}{\land}[2]}, checked
          [\sFmla{\False}{\varphi},
            just={\TRule{\False}{\lor}[4]}
            [\sFmla{\False}{\formula{C}},
              just={\TRule{\False}{\lor}[4]}, close]
          ]
        ]
      ]
      [\sFmla{\True}{\psi \land \formula{C}},
        just={\TRule{\True}{\lor}[1]}, checked
        [\sFmla{\False}{\varphi \lor \psi},
          just={\TRule{\False}{\land}[2]}, checked
          [\sFmla{\False}{\varphi},
            just={\TRule{\False}{\lor}[4]}
            [\sFmla{\False}{\psi},
              just={\TRule{\False}{\lor}[4]}
              [\sFmla{\True}{\psi},
                just={\TRule{\True}{\land}[3]}
                [\sFmla{\True}{\formula{C}},
                  just={\TRule{\True}{\land}[3]},close
                ]
              ]
            ]
          ]
        ]
        [\sFmla{\False}{\varphi \lor \formula{C}},
          just={\TRule{\False}{\land}[2]}, checked
          [\sFmla{\False}{\varphi},
            just={\TRule{\False}{\lor}[4]}
            [\sFmla{\False}{\formula{C}},
              just={\TRule{\False}{\lor}[4]}
              [\sFmla{\True}{\psi},
                just={\TRule{\True}{\land}[3]}
                [\sFmla{\True}{\formula{C}},
                  just={\TRule{\True}{\land}[3]},close
                ]
              ]
            ]
          ]
        ]
      ]
    ]
  ]
\end{oltableau}

तुलनेसाठी, त्याच गृहीतकांच्या संचासाठी नियम वेगळ्या क्रमाने लावलेला बंद
टॅब्लो पुढे दिला आहे:
\begin{oltableau}
  [\sFmla{\True}{\varphi \lor (\psi \land \formula{C})},
    just=\TAss, checked
    [\sFmla{\False}{(\varphi \lor \psi) \land
        (\varphi \lor \formula{C})}, just=\TAss, checked
      [\sFmla{\False}{\varphi \lor \psi},
        just={\TRule{\False}{\land}[2]}, checked
        [\sFmla{\False}{\varphi},
          just={\TRule{\False}{\lor}[3]}
          [\sFmla{\False}{\psi},
            just={\TRule{\False}{\lor}[3]}
            [\sFmla{\True}{\varphi},
              just={\TRule{\True}{\lor}[1]},close
            ]
            [\sFmla{\True}{\psi \land \formula{C}},
              just={\TRule{\True}{\lor}[1]}, checked
              [\sFmla{\True}{\psi},
                just={\TRule{\True}{\land}[6]}
                [\sFmla{\True}{\formula{C}},
                  just={\TRule{\True}{\land}[6]},close
                ]
              ]
            ]
          ]
        ]
      ]
      [\sFmla{\False}{\varphi \lor \formula{C}},
        just={\TRule{\False}{\land}[2]}, checked
        [\sFmla{\False}{\varphi},
          just={\TRule{\False}{\lor}[3]}
          [\sFmla{\False}{\formula{C}},
            just={\TRule{\False}{\lor}[3]}
            [\sFmla{\True}{\varphi},
              just={\TRule{\True}{\lor}[1]},close
            ]
            [\sFmla{\True}{\psi \land \formula{C}},
              just={\TRule{\True}{\lor}[1]}, checked
              [\sFmla{\True}{\psi},
                just={\TRule{\True}{\land}[6]}
                [\sFmla{\True}{\formula{C}},
                  just={\TRule{\True}{\land}[6]},close
                ]
              ]
            ]
          ]
        ]
      ]
    ]
  ]
\end{oltableau}
\end{ex}

\begin{prob}
पुढील सूत्रांचे बंद टॅब्लो द्या:
\begin{enumerate}
\item $\sFmla{\True}{\varphi \land (\psi \land \chi)}, \sFmla{\False}{(\varphi \land \psi) \land \chi}$.
\item $\sFmla{\True}{\varphi \lor (\psi \lor \chi)}, \sFmla{\False}{(\varphi \lor \psi) \lor \chi}$.
\item $\sFmla{\True}{\varphi \lif (\psi \lif \chi)}, \sFmla{\False}{\psi \lif (\varphi \lif \chi)}$.
\item $\sFmla{\True}{\varphi}, \sFmla{\False}{\lnot\lnot \varphi}$.
\end{enumerate}
\end{prob}

\begin{prob}
पुढील सूत्रांचे बंद टॅब्लो द्या:
\begin{enumerate}
\item $\sFmla{\True}{(\varphi \lor \psi) \lif \chi}, \sFmla{\False}{\varphi \lif \chi}$.
\item $\sFmla{\True}{(\varphi \lif \chi) \land (\psi \lif \chi)}, \sFmla{\False}{(\varphi \lor \psi) \lif \chi}$.
\item $\sFmla{\False}{\lnot(\varphi \land \lnot \varphi)}$.
\item $\sFmla{\True}{\psi \lif \varphi}, \sFmla{\False}{\lnot \varphi \lif \lnot \psi}$.
\item $\sFmla{\False}{(\varphi \lif \lnot \varphi) \lif \lnot \varphi}$.
\item $\sFmla{\False}{\lnot(\varphi \lif \psi) \lif \lnot \psi}$.
\item $\sFmla{\True}{\varphi \lif \chi}, \sFmla{\False}{\lnot (\varphi \land \lnot \chi)}$.
\item $\sFmla{\True}{\varphi \land \lnot \chi}, \sFmla{\False}{\lnot (\varphi \lif \chi)}$.
\item $\sFmla{\True}{\varphi \lor \psi, \lnot \psi}, \sFmla{\False}{\varphi}$.
\item $\sFmla{\True}{\lnot \varphi \lor \lnot \psi}, \sFmla{\False}{\lnot(\varphi \land \psi)}$.
\item $\sFmla{\False}{(\lnot \varphi \land \lnot \psi) \lif\lnot(\varphi \lor \psi)}$.
\item $\sFmla{\False}{\lnot(\varphi \lor \psi) \lif (\lnot \varphi \land \lnot \psi)}$.
\end{enumerate}
\end{prob}

\begin{prob}
पुढील सूत्रांचे बंद टॅब्लो द्या:
\begin{enumerate}
\item $\sFmla{\True}{\lnot(\varphi \lif \psi)}, \sFmla{\False}{\varphi}$.
\item $\sFmla{\True}{\lnot(\varphi \land \psi)}, \sFmla{\False}{\lnot \varphi \lor \lnot \psi}$.
\item $\sFmla{\True}{\varphi \lif \psi}, \sFmla{\False}{\lnot \varphi \lor \psi}$.
\item $\sFmla{\False}{\lnot \lnot \varphi \lif \varphi}$.
\item $\sFmla{\True}{\varphi \lif \psi}, \sFmla{\True}{\lnot \varphi \lif \psi}, \sFmla{\False}{\psi}$.
\item $\sFmla{\True}{(\varphi \land \psi) \lif \chi}, \sFmla{\False}{(\varphi \lif \chi) \lor (\psi \lif \chi)}$.
\item $\sFmla{\True}{(\varphi \lif \psi) \lif \varphi}, \sFmla{\False}{\varphi}$.
\item $\sFmla{\False}{(\varphi \lif \psi) \lor (\psi \lif \chi)}$.
\end{enumerate}
\end{prob}












\section{संख्यापकांसह टॅब्लो}\label{fol:tab:prq:sec}

\begin{ex}
संख्यापक हाताळताना आयगेन चराच्या अटीचा भंग होणार नाही याची खात्री करावी
लागते; त्यासाठी कधीकधी काही अनुमाने करण्याचा क्रम बदलून पाहावा लागतो.
सामान्यतः आयगेन चराच्या अटीच्या अधीन असलेले नियम आधी हाताळण्याचा प्रयत्न
उपयोगी ठरतो (पूर्ण टॅब्लोमध्ये ते अधिक वर असतील).

$\lexists[x][\lnot \varphi(x)] \lif \lnot \lforall[x][\varphi(x)]$ या
वाक्य साठी टॅब्लो कसा द्यायचा ते पाहू. नेहमीप्रमाणे संबंधित
गृहीतक नोंदवून आपण सुरुवात करतो:
\begin{oltableau}
[\sFmla{\False}{\lexists[x][\lnot \varphi(x)]\lif \lnot
    \lforall[x][\varphi(x)]}, just=\TAss]
\end{oltableau}
मुख्य संकारक $\lif$ असल्यामुळे आपण $\TRule{\False}{\lif}$ लावतो:
\begin{oltableau}
[\sFmla{\False}{\lexists[x][\lnot \varphi(x)]\lif \lnot
    \lforall[x][\varphi(x)]}, just=\TAss, checked
  [\sFmla{\True}{\lexists[x][\lnot \varphi(x)]},
    just={\TRule{\False}{\lif}[1]}
    [\sFmla{\False}{\lnot\lforall[x][\varphi(x)]},
      just={\TRule{\False}{\lif}[1]}]
  ]
]
\end{oltableau}
यानंतर ओळ~$2$ हाताळायची आहे. आपण $\TRule{\True}{\lexists}$ वापरतो.
त्यासाठी नवा स्थिरांक आवश्यक आहे; अद्याप कोणतेही स्थिरांक आलेले
नसल्यामुळे कोणताही एक, समजा~$a$, निवडता येतो.
\begin{oltableau}
[\sFmla{\False}{\lexists[x][\lnot \varphi(x)]\lif \lnot
    \lforall[x][\varphi(x)]}, just=\TAss, checked
  [\sFmla{\True}{\lexists[x][\lnot \varphi(x)]},
    just={\TRule{\False}{\lif}[1]}, checked
    [\sFmla{\False}{\lnot\lforall[x][\varphi(x)]},
      just={\TRule{\False}{\lif}[1]}
      [\sFmla{\True}{\lnot\varphi(a)}, just={\TRule{\True}{\lexists}[2]}]
    ]
  ]
]
\end{oltableau}
आता ओळ~$3$ ला $\TRule{\False}{\lnot}$ लावतो:
\begin{oltableau}
[\sFmla{\False}{\lexists[x][\lnot \varphi(x)]\lif \lnot
    \lforall[x][\varphi(x)]}, just=\TAss, checked
  [\sFmla{\True}{\lexists[x][\lnot \varphi(x)]},
    just={\TRule{\False}{\lif}[1]}, checked
    [\sFmla{\False}{\lnot\lforall[x][\varphi(x)]},
      just={\TRule{\False}{\lif}[1]}, checked
      [\sFmla{\True}{\lnot\varphi(a)}, just={\TRule{\True}{\lexists}[2]}
        [\sFmla{\True}{\lforall[x][\varphi(x)]},
          just = {\TRule{\False}{\lnot}[3]}
        ]
      ]
    ]
  ]
]
\end{oltableau}
ओळ~$4$ ला $\TRule{\True}{\lnot}$ आणि त्यानंतर ओळ~$5$ ला
$\TRule{\True}{\lforall}$ लावल्यावर बंद टॅब्लो मिळतो.
\begin{oltableau}
[\sFmla{\False}{\lexists[x][\lnot \varphi(x)]\lif \lnot
    \lforall[x][\varphi(x)]}, just=\TAss, checked
  [\sFmla{\True}{\lexists[x][\lnot \varphi(x)]},
    just={\TRule{\False}{\lif}[1]}, checked
    [\sFmla{\False}{\lnot\lforall[x][\varphi(x)]},
      just={\TRule{\False}{\lif}[1]}, checked
      [\sFmla{\True}{\lnot\varphi(a)}, just={\TRule{\True}{\lexists}[2]}
        [\sFmla{\True}{\lforall[x][\varphi(x)]},
          just = {\TRule{\False}{\lnot}[3]}
          [\sFmla{\False}{\varphi(a)}, just={\TRule{\True}{\lnot}[4]}
            [\sFmla{\True}{\varphi(a)},
              just={\TRule{\True}{\lforall}[5]}, close
            ]
          ]
        ]
      ]
    ]
  ]
]
\end{oltableau}
\end{ex}

\begin{ex}
पुढील संचासाठी टॅब्लो कसा द्यायचा ते पाहू:
\[\sFmla{\False}{\lexists[x][\chi(x,b)]},
\sFmla{\True}{\lexists[x][(\varphi(x) \land \psi(x))]},
\sFmla{\True}{\lforall[x][(\psi(x) \lif \chi(x,b))]}.
\]
नेहमीप्रमाणे गृहीतकांनी सुरुवात करतो:
\begin{oltableau}
  [\sFmla{\False}{\lexists[x][\formula{C}(x,b)]}, just=\TAss
    [\sFmla{\True}{\lexists[x][(\varphi(x) \land \psi(x))]},
      just=\TAss
      [\sFmla{\True}{\lforall[x][(\psi(x) \lif \formula{C}(x,b))]},
        just=\TAss]
    ]
  ]
\end{oltableau}
आयगेन चराची अट असलेला नियम आपण नेहमी आधी लावला पाहिजे; येथे तो म्हणजे
ओळ~$2$ ला $\TRule{\True}{\lexists}$. गृहीतकांत स्थिरांक~$b$ असल्यामुळे
वेगळा अचर वापरावा लागतो; पुन्हा~$a$ निवडू या.
\begin{oltableau}
  [\sFmla{\False}{\lexists[x][\formula{C}(x,b)]}, just=\TAss
    [\sFmla{\True}{\lexists[x][(\varphi(x) \land \psi(x))]},
      just=\TAss, checked
      [\sFmla{\True}{\lforall[x][(\psi(x) \lif \formula{C}(x,b))]},
        just=\TAss
        [\sFmla{\True}{\varphi(a) \land \psi(a)},
          just={\TRule{\True}{\lexists}[2]}]
      ]
    ]
  ]
\end{oltableau}
आता ओळ~$1$ ला $\TRule{\False}{\lexists}$ किंवा ओळ~$3$ ला
$\TRule{\True}{\lforall}$ लावल्यास, $x$ च्या जागी कोणते बंद पद~$t$ ठेवायचे
ते ठरवावे लागते. या नियमांना आयगेन चराची अट नसल्यामुळे हवे ते कोणतेही बंद
पद निवडता येते. काही स्थितींमध्ये वेगवेगळी~$t$ वापरून नियम अनेकदा लावावा
लागू शकतो. मात्र सर्वसाधारणपणे टॅब्लोमध्ये आधीच आलेल्या पदांपैकी
एक---येथे $a$ किंवा~$b$---निवडणे फायदेशीर ठरते; आणि येथे $a$ मुळे बंद
शाखा मिळण्याची शक्यता अधिक आहे असा अंदाज करता येतो.
\begin{oltableau}
  [\sFmla{\False}{\lexists[x][\formula{C}(x,b)]}, just=\TAss
    [\sFmla{\True}{\lexists[x][(\varphi(x) \land \psi(x))]},
      just=\TAss, checked
      [\sFmla{\True}{\lforall[x][(\psi(x) \lif \formula{C}(x,b))]}, just=\TAss
        [\sFmla{\True}{\varphi(a) \land \psi(a)}, just={\TRule{\True}{\lexists}[2]}
          [\sFmla{\False}{\formula{C}(a,b)}, just={\TRule{\False}{\lexists}[1]}
            [\sFmla{\True}{\psi(a) \lif \formula{C}(a,b)},
              just={\TRule{\True}{\lforall}[3]}
            ]
          ]
        ]
      ]
    ]
  ]
\end{oltableau}
ओळी $1$ आणि~$3$ वरील चिन्हांकित सूत्रे पुन्हा वापराव्या लागू शकतात,
म्हणून त्यांवर खूण करत नाही. आता ओळ~$4$ ला $\TRule{\True}{\land}$ लावा:
\begin{oltableau}
  [\sFmla{\False}{\lexists[x][\formula{C}(x,b)]}, just=\TAss
    [\sFmla{\True}{\lexists[x][(\varphi(x) \land \psi(x))]},
      just=\TAss, checked
      [\sFmla{\True}{\lforall[x][(\psi(x) \lif \formula{C}(x,b))]}, just=\TAss
        [\sFmla{\True}{\varphi(a) \land \psi(a)},
          just={\TRule{\True}{\lexists}[2]}, checked
          [\sFmla{\False}{\formula{C}(a,b)}, just={\TRule{\False}{\lexists}[1]}
            [\sFmla{\True}{\psi(a) \lif \formula{C}(a,b)},
              just={\TRule{\True}{\lforall}[3]}
              [\sFmla{\True}{\varphi(a)}, just={\TRule{\True}{\land}[4]}
                [\sFmla{\True}{\psi(a)}, just={\TRule{\True}{\land}[4]}
                ]
              ]
            ]
          ]
        ]
      ]
    ]
  ]
\end{oltableau}
आता ओळ~$6$ ला $\TRule{\True}{\lif}$ लावल्यास टॅब्लो बंद होतो:
\begin{oltableau}
  [\sFmla{\False}{\lexists[x][\formula{C}(x,b)]}, just=\TAss
    [\sFmla{\True}{\lexists[x][(\varphi(x) \land \psi(x))]},
      just=\TAss, checked
      [\sFmla{\True}{\lforall[x][(\psi(x) \lif \formula{C}(x,b))]},
        just=\TAss
        [\sFmla{\True}{\varphi(a) \land \psi(a)},
          just={\TRule{\True}{\lexists}[2]}, checked
          [\sFmla{\False}{\formula{C}(a,b)},
            just={\TRule{\False}{\lexists}[1]}
            [\sFmla{\True}{\psi(a) \lif \formula{C}(a,b)},
              just={\TRule{\True}{\lforall}[3]}, checked
              [\sFmla{\True}{\varphi(a)},
                just={\TRule{\True}{\land}[4]}
                [\sFmla{\True}{\psi(a)},
                  just={\TRule{\True}{\land}[4]}
                  [\sFmla{\False}{\psi(a)},
                    just={\TRule{\True}{\lif}[6]},close
                  ]
                  [\sFmla{\True}{\formula{C}(a,b)},
                    just={\TRule{\True}{\lif}[6]},close
                  ]
                ]
              ]
            ]
          ]
        ]
      ]
    ]
  ]
\end{oltableau}


\end{ex}

\begin{ex}
पुढील संचासाठी आपण टॅब्लो तयार करतो:
\[\sFmla{\True}{\lforall[x][\varphi(x)]}, \sFmla{\True}{\lforall[x][\varphi(x)]
  \lif \lexists[y][\psi(y)]}, \sFmla{\True}{\lnot\lexists[y][\psi(y)]}.
\]
नेहमीप्रमाणे गृहीतके लिहून सुरुवात करतो:
\begin{oltableau}
  [\sFmla{\True}{\lforall[x][\varphi(x)]}, just=\TAss
    [\sFmla{\True}{\lforall[x][\varphi(x)] \lif
        \lexists[y][\psi(y)]}, just=\TAss
      [\sFmla{\True}{\lnot\lexists[y][\psi(y)]}, just=\TAss
      ]
    ]
  ]
\end{oltableau}
ओळ~$3$ ला $\TRule{\True}{\lnot}$ नियम लावून सुरुवात करतो. ``आयगेन
चराची अट असलेले नियम नेहमी आधी लावा'' या नियमाचा एक परिणाम असा की
``आयगेन चराची अट नसलेले संख्यापक-नियम आवश्यक होईपर्यंत पुढे ढकला.''
शाखा फोडणारे नियमही पुढे ढकला.
\begin{oltableau}
  [\sFmla{\True}{\lforall[x][\varphi(x)]}, just=\TAss
    [\sFmla{\True}{\lforall[x][\varphi(x)] \lif
        \lexists[y][\psi(y)]}, just=\TAss
      [\sFmla{\True}{\lnot\lexists[y][\psi(y)]}, just=\TAss, checked
        [\sFmla{\False}{\lexists[y][\psi(y)]}, just={\TRule{\True}{\lnot}[3]}]
      ]
    ]
  ]
\end{oltableau}
नव्या ओळ~$4$ साठी $\TRule{\False}{\lexists}$ आवश्यक आहे आणि हा आयगेन
चराची अट नसलेला संख्यापक-नियम आहे. म्हणून तो पुढे ढकलून त्याऐवजी ओळ~$2$
वर $\TRule{\True}{\lif}$ वापरतो.
\begin{oltableau}
  [\sFmla{\True}{\lforall[x][\varphi(x)]}, just=\TAss
    [\sFmla{\True}{\lforall[x][\varphi(x)] \lif
        \lexists[y][\psi(y)]}, just=\TAss, checked
      [\sFmla{\True}{\lnot\lexists[y][\psi(y)]}, just=\TAss, checked
        [\sFmla{\False}{\lexists[y][\psi(y)]},
          just={\TRule{\True}{\lnot}[3]},
          [\sFmla{\False}{\lforall[x][\varphi(x)]}, just={\TRule{\True}{\lif}[2]}]
          [\sFmla{\True}{\lexists[y][\psi(y)]}, just={\TRule{\True}{\lif}[2]}]
        ]
      ]
    ]
  ]
\end{oltableau}
दोन्ही नव्या चिन्हांकित सूत्रे ना आयगेन चराच्या अटी असलेले नियम आवश्यक
आहेत; म्हणून ते नियम यानंतर लावावेत:
\begin{oltableau}
  [\sFmla{\True}{\lforall[x][\varphi(x)]}, just=\TAss
    [\sFmla{\True}{\lforall[x][\varphi(x)] \lif
        \lexists[y][\psi(y)]}, just=\TAss, checked
      [\sFmla{\True}{\lnot\lexists[y][\psi(y)]}, just=\TAss, checked
        [\sFmla{\False}{\lexists[y][\psi(y)]},
          just={\TRule{\True}{\lnot}[3]}
          [\sFmla{\False}{\lforall[x][\varphi(x)]}, just={\TRule{\True}{\lif}[2]},checked
            [\sFmla{\False}{\varphi(b)}, just={\TRule{\False}{\lforall}[5]}]
          ]
          [\sFmla{\True}{\lexists[y][\psi(y)]}, just={\TRule{\True}{\lif}[2]},checked
            [\sFmla{\True}{\psi(c)}, just={\TRule{\True}{\lexists}[5]}]
          ]
        ]
      ]
    ]
  ]
\end{oltableau}
शाखा बंद करण्यासाठी ओळी $1$ आणि~$4$ वरील चिन्हांकित सूत्रे वापराव्या
लागतात. त्यांच्या संबंधित नियमांना (\TRule{\True}{\lforall} आणि
\TRule{\False}{\lexists}) आयगेन चराची अट नाही; म्हणून योग्य असलेली कोणतीही
बंद पदे निवडण्याचे स्वातंत्र्य आहे. येथे ती अनुक्रमे $b$ आणि~$c$ आहेत.
\begin{oltableau}
  [\sFmla{\True}{\lforall[x][\varphi(x)]}, just=\TAss
    [\sFmla{\True}{\lforall[x][\varphi(x)] \lif
        \lexists[y][\psi(y)]}, just=\TAss, checked
      [\sFmla{\True}{\lnot\lexists[y][\psi(y)]}, just=\TAss, checked
        [\sFmla{\False}{\lexists[y][\psi(y)]},
          just={\TRule{\True}{\lnot}[3]}
          [\sFmla{\False}{\lforall[x][\varphi(x)]},
            just={\TRule{\True}{\lif}[2]},checked
            [\sFmla{\False}{\varphi(b)},
              just={\TRule{\False}{\lforall}[5]}
              [\sFmla{\True}{\varphi(b)},
                just={\TRule{\True}{\lforall}[1]},close
              ]
            ]
          ]
          [\sFmla{\True}{\lexists[y][\psi(y)]},
            just={\TRule{\True}{\lif}[2]},checked
            [\sFmla{\True}{\psi(c)},
              just={\TRule{\True}{\lexists}[5]}
              [\sFmla{\False}{\psi(c)},
                just={\TRule{\False}{\lexists}[4]},close
              ]
            ]
          ]
        ]
      ]
    ]
  ]
\end{oltableau}
\end{ex}

\begin{prob}
पुढील सूत्रांचे बंद टॅब्लो द्या:
\begin{enumerate}
\item $\sFmla{\False}{(\lforall[x][\varphi(x)] \land \lforall[y][\psi(y)])
\lif \lforall[z][(\varphi(z) \land \psi(z))]}$.
\item $\sFmla{\False}{(\lexists[x][\varphi(x)] \lor \lexists[y][\psi(y)])
\lif \lexists[z][(\varphi(z) \lor \psi(z))]}$.
\item $\sFmla{\True}{\lforall[x][(\varphi(x) \lif \psi)]},
\sFmla{\False}{\lexists[y][\varphi(y)] \lif \psi}$.
\item $\sFmla{\True}{\lforall[x][\lnot \varphi(x)]},
\sFmla{\False}{\lnot\lexists[x][\varphi(x)]}$.
\item $\sFmla{\False}{\lnot\lexists[x][\varphi(x)] \lif \lforall[x][\lnot
\varphi(x)]}$.
\item $\sFmla{\False}{\lnot\lexists[x][\lforall[y][((\varphi(x,y) \lif
\lnot \varphi(y,y)) \land (\lnot \varphi(y,y) \lif \varphi(x,y)))]]}$.
\end{enumerate}
\end{prob}

\begin{prob}
पुढील सूत्रांचे बंद टॅब्लो द्या:
\begin{enumerate}
\item $\sFmla{\False}{\lnot\lforall[x][\varphi(x)] \lif \lexists[x][\lnot\varphi(x)]}$.
\item $\sFmla{\True}{(\lforall[x][\varphi(x)] \lif \psi)}, \sFmla{\False}{\lexists[y][(\varphi(y) \lif \psi)]}$.
\item $\sFmla{\False}{\lexists[x][(\varphi(x) \lif \lforall[y][\varphi(y)])]}$.
\end{enumerate}
\end{prob}












\section{सिद्धता-उपपत्तीय संकल्पना}\label{fol:tab:ptn:sec}

\begin{editorial}
या विभागात टॅब्लोसाठी सिद्धतायोग्यता संबंध आणि सुसंगतता यांच्या व्याख्या
एकत्रित केल्या आहेत.
\end{editorial}

\begin{explain}
जशा आपण अनेक महत्त्वाच्या चिन्हार्थविषयक संकल्पना (वैधता, तार्किक निष्पन्नता,
पूर्ततायोग्यता) परिभाषित केल्या आहेत, तशाच त्यांना अनुरूप
\emph{सिद्धता-उपपत्तीय संकल्पना} आता परिभाषित करू. संरचना मध्ये
वाक्यांची पूर्ति होते की नाही याचा आधार घेऊन या संकल्पना परिभाषित
केल्या जात नाहीत; त्याऐवजी ठरावीक बंद टॅब्लोx च्या अस्तित्वाचा आधार
घेतला जातो. या दोन प्रकारच्या संकल्पना एकमेकांशी जुळतात, हा एक महत्त्वाचा
शोध होता. त्या जुळतात, हाच \emph{निर्दोषता प्रमेय} आणि
\emph{संपूर्णता प्रमेय} यांचा आशय आहे.
\end{explain}


\begin{defn}[प्रमेये]
एखादे वाक्य~$\varphi$ हे \emph{प्रमेय} असते, जर
$\sFmla{\False}{\varphi}$ साठी बंद टॅब्लो असेल. $\varphi$ हे प्रमेय असेल
तर आपण $\Proves \varphi$ लिहितो आणि ते प्रमेय नसेल तर $\Proves/ \varphi$ लिहितो.
\end{defn}

\begin{defn}[निष्पन्नता]
वाक्य $\varphi$ हे वाक्यांच्या संच~$\Gamma$ \emph{पासून
निष्पन्न करता येण्याजोगे} असते, म्हणजे $\Gamma \Proves \varphi$, तेव्हा आणि केवळ तेव्हाच,
जेव्हा एखादा सांत संच $\{\psi_1, \dots, \psi_n\} \subseteq \Gamma$ असा असतो
आणि पुढील संचासाठी बंद टॅब्लो असतो:
\[\{
  \sFmla{\False}{\varphi},
  \sFmla{\True}{\psi_1}, \dots,
  \sFmla{\True}{\psi_n}
  \}.
\]
$\varphi$ हे $\Gamma$ पासून निष्पन्न करता येण्याजोगे नसेल, तर आपण $\Gamma \Proves/
\varphi$ लिहितो.
\end{defn}

\begin{defn}[सुसंगतता]
वाक्यांचा संच~$\Gamma$ हा \emph{विसंगत} असतो तेव्हा आणि केवळ तेव्हाच,
जेव्हा एखादा सांत संच $\{\psi_1, \dots, \psi_n\} \subseteq \Gamma$ असा असतो
आणि पुढील संचासाठी बंद टॅब्लो असतो:
\[\{
  \sFmla{\True}{\psi_1}, \dots,
  \sFmla{\True}{\psi_n}
  \}.
\]
$\Gamma$ विसंगत नसेल, तर त्याला \emph{सुसंगत} म्हणतो.
\end{defn}

\begin{prop}[परावर्तिता]
\label{fol:tab:ptn:prop:reflexivity}
$\varphi \in \Gamma$ असेल, तर $\Gamma \Proves \varphi$.
\end{prop}

\begin{proof}
$\varphi \in \Gamma$ असेल, तर $\{\varphi\}$ हा $\Gamma$ चा सांत उपसंच आहे आणि पुढील टॅब्लो
\begin{oltableau}
  [\sFmla{\False}{\varphi}, just = \TAss
    [\sFmla{\True}{\varphi}, just = \TAss,close]
  ]
\end{oltableau}
बंद आहे.
\end{proof}


\begin{prop}[एकस्वनिकता]
\label{fol:tab:ptn:prop:monotonicity}
$\Gamma \subseteq \Delta$ आणि $\Gamma \Proves \varphi$ असेल, तर $\Delta
\Proves \varphi$.
\end{prop}

\begin{proof}
$\Gamma$ चा कोणताही सांत उपसंच हा $\Delta$ चासुद्धा सांत उपसंच असतो.
\end{proof}

\begin{prop}[संक्रमकता]
\label{fol:tab:ptn:prop:transitivity}
$\Gamma \Proves \varphi$ आणि $\{\varphi\} \cup
\Delta \Proves \psi$ असेल, तर $\Gamma \cup \Delta \Proves \psi$.
\end{prop}

\begin{proof}
$\{\varphi\} \cup \Delta \Proves \psi$ असेल, तर $\Delta_0 =
\{\chi_1, \dots, \chi_n\} \subseteq \Delta$ असा सांत उपसंच असतो आणि पुढील संचासाठी
\begin{align*}
\{\sFmla{\False}{\psi}, & \sFmla{\True}{\varphi}, \sFmla{\True}{\chi_1},
\dots, \sFmla{\True}{\chi_n}\}
\intertext{बंद टॅब्लो असतो. जर $\Gamma \Proves \varphi$ असेल, तर
  $\{\theta_1, \dots, \theta_m\} \subseteq \Gamma$ असा संच असतो आणि पुढील संचासाठीही}
\{\sFmla{\False}{\varphi}, & \sFmla{\True}{\theta_1},
\dots, \sFmla{\True}{\theta_m}\}
\end{align*}
बंद टॅब्लो असतो.

आता पुढील गृहीतके असलेला टॅब्लो विचारात घ्या:
\[\sFmla{\False}{\psi},
\sFmla{\True}{\chi_1}, \dots, \sFmla{\True}{\chi_n},
\sFmla{\True}{\theta_1}, \dots, \sFmla{\True}{\theta_m}.
\]
त्यात~$\varphi$ वर \Cut{} नियम लावा. यामुळे दोन शाखा तयार होतात: एका शाखेत
$\sFmla{\True}{\varphi}$, तर दुसऱ्या शाखेत $\sFmla{\False}{\varphi}$ असते. त्यामुळे
पहिल्या शाखेत पुढील सर्व गृहीतके उपलब्ध असतात:
\[\{\sFmla{\False}{\psi}, \sFmla{\True}{\varphi},
\sFmla{\True}{\chi_1}, \dots, \sFmla{\True}{\chi_n}\}
\]
या गृहीतकांसाठी बंद टॅब्लो असल्यामुळे तो त्या शाखेला जोडता येतो;
$\sFmla{\True}{\varphi}$ मधून जाणारी प्रत्येक शाखा बंद होते. दुसऱ्या शाखेत
पुढील सर्व गृहीतके उपलब्ध असतात:
\[\{\sFmla{\False}{\varphi}, \sFmla{\True}{\theta_1}, \dots,
\sFmla{\True}{\theta_m}\}
\]
म्हणून दुसरी शाखासुद्धा पूर्ण करून बंद टॅब्लो मिळवता येतो. यावरून
$\Gamma \cup \Delta \Proves \psi$ हे दिसते.
\end{proof}

विशेषतः, याचा अर्थ असा की $\Gamma \Proves \varphi$ आणि $\varphi
\Proves \psi$ असेल, तर $\Gamma \Proves \psi$. तसेच $\varphi_1,
\dots, \varphi_n \Proves \psi$ असेल आणि प्रत्येक~$i$ साठी $\Gamma \Proves \varphi_i$
असेल, तर $\Gamma \Proves \psi$, हेसुद्धा यावरून मिळते.

\begin{prop}
\label{fol:tab:ptn:prop:incons}
$\Gamma$ विसंगत असतो तेव्हा आणि केवळ तेव्हाच, जेव्हा प्रत्येक
  वाक्य~$\varphi$ साठी $\Gamma \Proves \varphi$.
\end{prop}

\begin{proof}
सराव.
\end{proof}


\begin{prob}
\ref{fol:tab:ptn:prop:incons} सिद्ध करा.
\end{prob}




\begin{prop}[संहतता]
  \label{fol:tab:ptn:prop:proves-compact}
  \begin{enumerate}
  \item $\Gamma \Proves \varphi$ असेल, तर एखादा सांत उपसंच $\Gamma_0
    \subseteq \Gamma$ असा आहे की $\Gamma_0 \Proves \varphi$.
  \item $\Gamma$ चा प्रत्येक सांत उपसंच सुसंगत असेल, तर $\Gamma$ सुसंगत
    असतो.
  \end{enumerate}
\end{prop}

\begin{proof}
  \begin{enumerate}
    \item $\Gamma \Proves \varphi$ असेल, तर $\Gamma_0 =
      \{\psi_1, \dots, \psi_n\}$ असा त्याचा सांत उपसंच आणि पुढील संचासाठी बंद टॅब्लो असतो:
      \[      \{\sFmla{\False}{\varphi}, \sFmla{\True}{\psi_1}, \dots, \sFmla{\True}{\psi_n}\}
      \]
      हा टॅब्लो $\Gamma_0 \Proves \varphi$ हेसुद्धा दाखवतो.
    \item $\Gamma$ विसंगत असेल, तर
      $\Gamma_0 = \{\psi_1, \dots, \psi_n\}$ असा त्याचा एखादा सांत उपसंच असून पुढील
      संचाचा बंद टॅब्लो असतो:
      \[      \{\sFmla{\True}{\psi_1}, \dots, \sFmla{\True}{\psi_n}\}
      \]
      हा बंद टॅब्लो $\Gamma_0$ विसंगत असल्याचे दाखवतो.
  \end{enumerate}
\end{proof}












\section{निष्पन्नता आणि सुसंगतता}\label{fol:tab:prv:sec}

आता आपण निष्पन्नता संबंधाचे अनेक गुणधर्म सिद्ध करू. ते स्वतंत्रपणेही
महत्त्वाचे आहेत, पण संपूर्णता प्रमेयाच्या सिद्धतेत प्रत्येकाची भूमिका असेल.

\begin{prop}\label{fol:tab:prv:prop:provability-contr}
  जर $\Gamma \Proves \varphi$ आणि $\Gamma \cup \{\varphi\}$ विसंगत असेल, तर
  $\Gamma$ विसंगत आहे.
\end{prop}

\begin{proof}
$\Gamma_0 = \{\psi_1, \dots, \psi_n\}$ आणि $\Gamma_1
  =\{\chi_1, \dots, \chi_m\} \subseteq \Gamma$ असे सांत उपसंच असतात की
  \begin{align*}
    \{\sFmla{\False}{\varphi}, &
    \sFmla{\True}{\psi_1}, \dots, \sFmla{\True}{\psi_n}\} \\
    \{\sFmla{\True}{\varphi}, &
    \sFmla{\True}{\chi_1}, \dots, \sFmla{\True}{\chi_m}\}
  \end{align*}
  यांसाठी बंद टॅब्लो असतात. $\varphi$ वर \Cut{} नियम वापरून हे दोन्ही
  एकाच बंद टॅब्लोमध्ये एकत्र करता येतात; त्यावरून $\Gamma_0
  \cup \Gamma_1$ विसंगत असल्याचे दिसते. $\Gamma_0
  \subseteq \Gamma$ आणि $\Gamma_1 \subseteq \Gamma$ असल्यामुळे $\Gamma_0 \cup
  \Gamma_1 \subseteq \Gamma$; म्हणून $\Gamma$ विसंगत आहे.
\end{proof}

\begin{prop}
\label{fol:tab:prv:prop:prov-incons}
$\Gamma \Proves \varphi$ तेव्हा आणि केवळ तेव्हाच, जेव्हा $\Gamma \cup \{\lnot \varphi\}$
विसंगत असते.
\end{prop}

\begin{proof}
प्रथम $\Gamma \Proves \varphi$ असे समजा; म्हणजे पुढील संचासाठी बंद
टॅब्लो असतो:
\[\{\sFmla{\False}{\varphi},
\sFmla{\True}{\psi_1}, \dots, \sFmla{\True}{\psi_n}\}
\]
$\TRule{\True}{\lnot}$ नियम वापरून त्याचे पुढील संचासाठीच्या बंद
टॅब्लोमध्ये रूपांतर करता येते:
\[\{\sFmla{\True}{\lnot \varphi},
\sFmla{\True}{\psi_1}, \dots, \sFmla{\True}{\psi_n}\}.
\]

दुसरीकडे, दुसऱ्या संचासाठी बंद टॅब्लो असेल, तर त्यातील पहिले
गृहीतक~$\sFmla{\True}{\lnot \varphi}$ आणि त्या गृहीतकाला
\TRule{\True}{\lnot} लावून मिळणारी सर्व सूत्रे काढून, तसेच
$\sFmla{\False}{\varphi}$ हे गृहीतक जोडून, त्याचे पहिल्या संचासाठीच्या बंद
टॅब्लोमध्ये रूपांतर करता येते. कारण एखादी शाखा पूर्वी
$\sFmla{\True}{\lnot \varphi}$ ला \TRule{\True}{\lnot} लावून मिळणारा
निष्कर्ष, म्हणजे $\sFmla{\False}{\varphi}$, असल्यामुळे बंद झाली असेल, तर नव्या
टॅब्लो मधील संबंधित शाखाही बंद असते. जुन्या टॅब्लोतील एखादी शाखा
$\sFmla{\True}{\lnot \varphi}$ हे गृहीतक आणि $\sFmla{\False}{\lnot \varphi}$
ही दोन्ही सूत्रे असल्यामुळे बंद झाली असेल, तर $\sFmla{\False}{\lnot \varphi}$
ला $\TRule{\False}{\lnot}$ लावून $\sFmla{\True}{\varphi}$ मिळवता येते आणि
ती शाखा बंद करता येते. आपण $\sFmla{\False}{\varphi}$ हे गृहीतक जोडलेले
असल्यामुळे ही शाखा बंद होते.
\end{proof}

\begin{prob}
$\Gamma \Proves \lnot \varphi$ तेव्हा आणि केवळ तेव्हाच, जेव्हा $\Gamma \cup \{\varphi\}$
विसंगत असते, हे सिद्ध करा.
\end{prob}

\begin{prop}\label{fol:tab:prv:prop:explicit-inc}
  जर $\Gamma \Proves \varphi$ आणि $\lnot \varphi \in \Gamma$ असेल, तर $\Gamma$
  विसंगत आहे.
\end{prop}

\begin{proof}
  समजा $\Gamma \Proves \varphi$ आणि $\lnot \varphi \in \Gamma$. मग
  $\psi_1$, \dots, $\psi_n \in \Gamma$ अशी वाक्ये असतात की
  \[  \{\sFmla{\False}{\varphi}, \sFmla{\True}{\psi_1}, \dots, \sFmla{\True}{\psi_n}\}
  \]
  यासाठी बंद टॅब्लो असतो. \sFmla{\False}{\varphi} हे गृहीतक
  \sFmla{\True}{\lnot \varphi} ने बदला आणि \sFmla{\True}{\lnot \varphi} ला
  \TRule{\True}{\lnot} लावून मिळणारा निष्कर्ष गृहीतकांनंतर घाला. जुन्या
  टॅब्लोमध्ये ओळ~$1$ च्या आधारे समर्थित कोणतेही वाक्य आता
  ओळ~$n+2$ च्या आधारे समर्थित होते. म्हणून जुना टॅब्लो बंद असेल, तर
  नवा टॅब्लोसुद्धा बंद असतो. सर्व गृहीतके~$\Gamma$ मध्ये असल्यामुळे हा टॅब्लो
  $\Gamma$ विसंगत असल्याचे दाखवतो.
\end{proof}

\begin{prop}\label{fol:tab:prv:prop:provability-exhaustive}
  $\Gamma \cup \{\varphi\}$ आणि $\Gamma \cup \{\lnot \varphi\}$ हे दोन्ही संच
  विसंगत असतील, तर $\Gamma$ विसंगत आहे.
\end{prop}

\begin{proof}
  $\psi_1$, \dots, $\psi_n \in \Gamma$ आणि $\chi_1$, \dots,
  $\chi_m \in \Gamma$ अशी वाक्ये असतील आणि
  \begin{align*}
    \{\sFmla{\True}{\varphi}, &
    \sFmla{\True}{\psi_1}, \dots, \sFmla{\True}{\psi_n}\} \text{ आणि}\\
    \{\sFmla{\True}{\lnot \varphi}, &
    \sFmla{\True}{\chi_1}, \dots, \sFmla{\True}{\chi_m}\}
  \end{align*}
  या दोन्हींसाठी बंद टॅब्लो असतील, तर
  $\sFmla{\True}{\psi_1}$, \dots, $\sFmla{\True}{\psi_n}$ आणि
  $\sFmla{\True}{\chi_1}$, \dots, $\sFmla{\True}{\chi_m}$ ही गृहीतके घेऊन,
  त्यानंतर \Cut{} नियम लावून, $\Gamma$ विसंगत असल्याचे दाखवणारा एकच
  संयुक्त टॅब्लो तयार करता येतो. त्यामुळे दोन शाखा मिळतात: एक
  $\sFmla{\True}{\varphi}$ ने, तर दुसरी $\sFmla{\False}{\varphi}$ ने सुरू होते.

  डाव्या बाजूस पहिल्या टॅब्लो मधील त्याच्या गृहीतकांखालचा भाग जोडा.
  येथे प्रत्येक नियमाचा उपयोग अजूनही योग्य आहे, कारण
  $\sFmla{\True}{\varphi}$ सहित पहिल्या टॅब्लो मधील प्रत्येक गृहीतक
  उपलब्ध आहे. त्यामुळे $\sFmla{\True}{\varphi}$ च्या खालील प्रत्येक शाखा बंद होते.

  उजव्या बाजूस दुसऱ्या टॅब्लो मधील त्याच्या गृहीतकांखालचा भाग जोडा;
  मात्र $\sFmla{\True}{\lnot \varphi}$ ला~$\TRule{\True}{\lnot}$ लावून मिळणारे
  सर्व परिणाम काढून टाका. $\sFmla{\True}{\lnot \varphi}$ ला
  $\TRule{\True}{\lnot}$ लावल्याचा निष्कर्ष $\sFmla{\False}{\varphi}$ हा आहे;
  तरीही तो उपलब्ध असतो, कारण तो संयुक्त टॅब्लोच्या उजव्या बाजूवरील
  \Cut{} नियमाचा निष्कर्ष आहे.

  दुसऱ्या टॅब्लोतील एखादी शाखा $\sFmla{\True}{\lnot \varphi}$ हे गृहीतक
  (जे आता संयुक्त टॅब्लोमध्ये गृहीतक म्हणून येत नाही) आणि
  $\sFmla{\False}{\lnot \varphi}$ ही दोन्ही सूत्रे असल्यामुळे बंद झाली असेल,
  तर $\sFmla{\False}{\lnot \varphi}$ ला $\TRule{\False}{\lnot}$ लावून
  $\sFmla{\True}{\varphi}$ मिळवता येते. मग संयुक्त टॅब्लो मधील संबंधित
  शाखासुद्धा बंद होते, कारण तिच्यात \Cut{} नियमाचा उजवीकडील निष्कर्ष
  $\sFmla{\False}{\varphi}$ असतो. दुसऱ्या टॅब्लो मधील एखादी शाखा अन्य
  कोणत्याही कारणाने बंद झाली असेल, तर संयुक्त टॅब्लो मधील संबंधित
  शाखासुद्धा बंद होते; कारण जुन्या दुसऱ्या टॅब्लो मधील त्या शाखेवर
  येणारी $\sFmla{\True}{\lnot \varphi}$ व्यतिरिक्त इतर सर्व चिन्हांकित सूत्रे
  संयुक्त टॅब्लो मधील संबंधित शाखेवरही येतात.
  \end{proof}















\section{निष्पन्नता आणि विधानीय संयोजक}\label{fol:tab:ppr:sec}

\begin{explain}
  टॅब्लोचा निष्पन्नता संबंध~$\Proves$ हा विधानीय संयोजकांशी संबंधित
  काही मूलभूत तथ्ये सिद्ध करण्यासाठी पुरेसा सामर्थ्यवान आहे; उदाहरणार्थ,
  $\varphi \land \psi \Proves \varphi$ आणि $\varphi, \varphi \lif \psi \Proves \psi$
  (विध्यात्मक अनुमान). ही तथ्ये संपूर्णता प्रमेयाच्या सिद्धतेसाठी आवश्यक आहेत.
\end{explain}

\begin{prop}\label{fol:tab:ppr:prop:provability-land}
  \begin{enumerate}
  \item \label{fol:tab:ppr:prop:provability-land-left} $\varphi \land \psi \Proves
    \varphi$ आणि $\varphi \land \psi \Proves \psi$ ही दोन्ही तथ्ये सत्य आहेत.
  \item \label{fol:tab:ppr:prop:provability-land-right} $\varphi, \psi \Proves \varphi \land
    \psi$.
  \end{enumerate}
\end{prop}

\begin{proof}
  \begin{enumerate}
  \item $\{\sFmla{\False}{\varphi}, \sFmla{\True}{\varphi \land \psi}\}$ आणि
    $\{\sFmla{\False}{\psi}, \sFmla{\True}{\varphi \land \psi}\}$ या दोन्हींसाठी
    बंद टॅब्लो आहेत:
    \begin{oltableau}{}
      [\sFmla{\False}{\varphi}, just=\TAss
        [\sFmla{\True}{\varphi \land \psi}, just=\TAss
          [\sFmla{\True}{\varphi},just={\TRule{\True}{\land}[2]}
            [\sFmla{\True}{\psi},just={\TRule{\True}{\land}[2]}, close
            ]
          ]
        ]
      ]
    \end{oltableau}
    \begin{oltableau}{}
      [\sFmla{\False}{\psi}, just=\TAss
        [\sFmla{\True}{\varphi \land \psi}, just=\TAss
          [\sFmla{\True}{\varphi},just={\TRule{\True}{\land}[2]}
            [\sFmla{\True}{\psi},just={\TRule{\True}{\land}[2]}, close
            ]
          ]
        ]
      ]
    \end{oltableau}
    \item $\{\sFmla{\True}{\varphi}, \sFmla{\True}{\psi},
      \sFmla{\False}{\varphi \land \psi}\}$ साठीचा बंद टॅब्लो पुढीलप्रमाणे आहे:
      \begin{oltableau}
        [\sFmla{\False}{\varphi \land \psi}, just = \TAss
          [\sFmla{\True}{\varphi}, just = \TAss
            [\sFmla{\True}{\psi}, just=\TAss
              [\sFmla{\False}{\varphi}, just = {\TRule{\False}{\land}[1]}, close]
              [\sFmla{\False}{\psi}, just = {\TRule{\False}{\land}[1]}, close]
            ]
          ]
        ]
      \end{oltableau}
  \end{enumerate}
\end{proof}

\begin{prop}\label{fol:tab:ppr:prop:provability-lor}
  \begin{enumerate}
  \item $\{\varphi \lor \psi, \lnot \varphi, \lnot \psi\}$ विसंगत आहे.
  \item $\varphi \Proves \varphi \lor \psi$ आणि $\psi \Proves \varphi \lor \psi$ ही दोन्ही
    तथ्ये सत्य आहेत.
  \end{enumerate}
\end{prop}

\begin{proof}
  \begin{enumerate}
  \item $\{\sFmla{\True}{\varphi \lor \psi}, \sFmla{\True}{\lnot \varphi},
    \sFmla{\True}{\lnot \psi}\}$ साठी बंद टॅब्लो देऊ:
    \begin{oltableau}
      [\sFmla{\True}{\varphi \lor \psi}, just = \TAss
        [\sFmla{\True}{\lnot \varphi}, just = \TAss
          [\sFmla{\True}{\lnot \psi}, just = \TAss
            [\sFmla{\False}{\varphi}, just = {\TRule{\True}{\lnot}[2]}
              [\sFmla{\False}{\psi}, just = {\TRule{\True}{\lnot}[3]}
                [\sFmla{\True}{\varphi}, just = {\TRule{\True}{\lor}[1]}, close]
                [\sFmla{\True}{\psi}, just = {\TRule{\True}{\lor}[1]}, close]
              ]
            ]
          ]
        ]
      ]
    \end{oltableau}
  \item $\{\sFmla{\False}{\varphi \lor \psi}, \sFmla{\True}{\varphi}\}$ आणि
    $\{\sFmla{\False}{\varphi \lor \psi}, \sFmla{\True}{\psi}\}$ या दोन्हींसाठी
    बंद टॅब्लो आहेत:
    \begin{oltableau}{}
      [\sFmla{\False}{\varphi \lor \psi}, just=\TAss
        [\sFmla{\True}{\varphi}, just=\TAss
          [\sFmla{\False}{\varphi},just={\TRule{\False}{\lor}[1]}
            [\sFmla{\False}{\psi},just={\TRule{\False}{\lor}[1]}, close
            ]
          ]
        ]
      ]
    \end{oltableau}
    \begin{oltableau}{}
      [\sFmla{\False}{\varphi \lor \psi}, just=\TAss
        [\sFmla{\True}{\psi}, just=\TAss
          [\sFmla{\False}{\varphi},just={\TRule{\False}{\lor}[1]}
            [\sFmla{\False}{\psi},just={\TRule{\False}{\lor}[1]}, close
            ]
          ]
        ]
      ]
    \end{oltableau}
  \end{enumerate}
\end{proof}

\begin{prop}\label{fol:tab:ppr:prop:provability-lif}
  \begin{enumerate}
  \item \label{fol:tab:ppr:prop:provability-lif-left}  $\varphi, \varphi \lif \psi \Proves \psi$.
  \item \label{fol:tab:ppr:prop:provability-lif-right}
    $\lnot \varphi \Proves \varphi \lif \psi$ आणि $\psi \Proves \varphi \lif \psi$ ही दोन्ही
    तथ्ये सत्य आहेत.
  \end{enumerate}
\end{prop}

\begin{proof}
  \begin{enumerate}
    \item $\{\sFmla{\False}{\psi}, \sFmla{\True}{\varphi \lif \psi},
      \sFmla{\True}{\varphi}\}$ साठी बंद टॅब्लो आहे:
      \begin{oltableau}
        [\sFmla{\False}{\psi}, just=\TAss
          [\sFmla{\True}{\varphi \lif \psi}, just=\TAss
            [\sFmla{\True}{\varphi}, just=\TAss
              [\sFmla{\False}{\varphi}, just = {\TRule{\True}{\lif}[2]}, close]
              [\sFmla{\True}{\psi}, just = {\TRule{\True}{\lif}[2]}, close]
            ]
          ]
        ]
      \end{oltableau}
    \item $\{\sFmla{\False}{\varphi \lif \psi}, \sFmla{\True}{\lnot \varphi}\}$ आणि
      $\{\sFmla{\False}{\varphi \lif \psi}, \sFmla{\True}{\psi}\}$ या दोन्हींसाठी
      बंद टॅब्लो आहेत:
      \begin{oltableau}
        [\sFmla{\False}{\varphi \lif \psi}, just = \TAss
          [\sFmla{\True}{\lnot \varphi}, just = \TAss
            [\sFmla{\True}{\varphi}, just = {\TRule{\False}{\lif}[1]}
              [\sFmla{\False}{\psi}, just = {\TRule{\False}{\lif}[1]}
                [\sFmla{\False}{\varphi}, just = {\TRule{\True}{\lnot}[2]}, close]
              ]
            ]
          ]
        ]
      \end{oltableau}
      \begin{oltableau}
        [\sFmla{\False}{\varphi \lif \psi}, just = \TAss
          [\sFmla{\True}{\psi}, just = \TAss
            [\sFmla{\True}{\varphi}, just = {\TRule{\False}{\lif}[1]}
              [\sFmla{\False}{\psi}, just = {\TRule{\False}{\lif}[1]},close
              ]
            ]
          ]
        ]
      \end{oltableau}
  \end{enumerate}
\end{proof}











\section{निष्पन्नता आणि संख्यापक}\label{fol:tab:qpr:sec}

\begin{explain}
  संपूर्णता प्रमेयासाठी या विभागात स्थापित केलेली~$\Proves$ संबंधी तथ्ये
  टॅब्लोच्या नियमांनी निष्पन्न होतात, हेही आवश्यक आहे.
\end{explain}

\begin{thm}
\label{fol:tab:qpr:thm:strong-generalization} जर $c$ हे $\Gamma$ किंवा~$\varphi(x)$ मध्ये
न येणारे अचर असेल आणि $\Gamma \Proves \varphi(c)$, तर $\Gamma
\Proves \lforall[x][\varphi(x)]$.
\end{thm}

\begin{proof}
$\Gamma \Proves \varphi(c)$ असे समजा; म्हणजे $\psi_1$, \dots, $\psi_n
\in \Gamma$ अशी वाक्ये आणि पुढील संचासाठी बंद टॅब्लो असतो:
\begin{align*}
\{ \sFmla{\False}{\varphi(c)}, &
\sFmla{\True}{\psi_1}, \dots, \sFmla{\True}{\psi_n} \}.
\intertext{पुढील संचासाठीही बंद टॅब्लो असतो, हे दाखवायचे आहे:}
\{ \sFmla{\False}{\lforall[x][\varphi(x)]}, &
\sFmla{\True}{\psi_1}, \dots, \sFmla{\True}{\psi_n} \}.
\end{align*}
बंद टॅब्लो घेऊन त्यातील पहिले गृहीतक
$\sFmla{\False}{\lforall[x][\varphi(x)]}$ ने बदला आणि गृहीतकांनंतर
$\sFmla{\False}{\varphi(c)}$ घाला.
\begin{center}
\begin{tableau}{not line numbering}
  [\sFmla{\False}{\varphi(c)}
    [\sFmla{\True}{\psi_1}
      [\vdots
      [\sFmla{\True}{\psi_n},
        [][][]]]]]
\end{tableau}\qquad
\begin{tableau}{not line numbering}
  [\sFmla{\False}{\lforall[x][\varphi(x)]}
    [\sFmla{\True}{\psi_1}
      [\vdots
        [\sFmla{\True}{\psi_n},
          [\sFmla{\False}{\varphi(c)}
        [][][]]]]]]
\end{tableau}
\end{center}
हा टॅब्लो अजूनही बंद आहे, कारण पूर्वी गृहीतके म्हणून उपलब्ध असलेली सर्व
वाक्ये नव्या टॅब्लोच्या वरच्या भागात अजूनही उपलब्ध आहेत. घातलेली
ओळ ही $\TRule{\False}{\lforall}$ च्या योग्य उपयोगाचा परिणाम आहे, कारण
स्थिरांक~$c$ हे $\psi_1$, \dots, $\psi_n$ किंवा
$\lforall[x][\varphi(x)]$ मध्ये येत नाही; म्हणजे ते नव्या टॅब्लोमध्ये घातलेल्या
ओळीच्या वर येत नाही.
\end{proof}

\begin{prop}
\label{fol:tab:qpr:prop:provability-quantifiers}
\begin{enumerate}
\item $\varphi(t) \Proves \lexists[x][\varphi(x)]$.
\item $\lforall[x][\varphi(x)] \Proves \varphi(t)$.
\end{enumerate}
\end{prop}

\begin{proof}
\begin{enumerate}
\item
  $\sFmla{\False}{\lexists[x][\varphi(x)]}, \sFmla{\True}{\varphi(t)}$ साठीचा बंद
  टॅब्लो पुढीलप्रमाणे आहे:
  \begin{oltableau}
    [\sFmla{\False}{\lexists[x][\varphi(x)]}, just = \TAss
      [\sFmla{\True}{\varphi(t)}, just = \TAss
        [\sFmla{\False}{\varphi(t)}, just = {\TRule{\False}{\lexists}[1]}, close]]]
    \end{oltableau}
\item
  $\sFmla{\False}{\varphi(t)}, \sFmla{\True}{\lforall[x][\varphi(x)]}, $
  साठीचा बंद टॅब्लो पुढीलप्रमाणे आहे:
  \begin{oltableau}
    [\sFmla{\False}{\varphi(t)}, just = \TAss
      [\sFmla{\True}{\lforall[x][\varphi(x)]}, just = \TAss
        [\sFmla{\True}{\varphi(t)}, just = {\TRule{\True}{\lforall}[2]}, close]]]
    \end{oltableau}
\end{enumerate}
\end{proof}












\section{निर्दोषता}\label{fol:tab:sou:sec}

\begin{explain}
टॅब्लोसारखी निष्पत्ती-पद्धत \emph{निर्दोष} असते, जर ती प्रत्यक्षात
धारण न करणाऱ्या गोष्टी निष्पन्न करू शकत नसेल. त्यामुळे निर्दोषता हा
निष्पत्ती-पद्धतींसाठी हमी दिलेल्या सुरक्षिततेचा एक प्रकार आहे. कोणता
सिद्धता-उपपत्तीय गुणधर्म विचारात आहे यावर अवलंबून, उदाहरणार्थ, आपल्याला
पुढील गोष्टी हव्या असतात:
\begin{enumerate}
\item प्रत्येक निष्पन्न करता येण्याजोगे~$\varphi$ हे वैध असते;
\item एखादे वाक्य इतर काही वाक्यांपासून निष्पन्न करता येण्याजोगे असेल, तर ते
  त्यांचा तार्किक परिणामही असते;
\item वाक्यांचा संच विसंगत असेल, तर तो अपूर्ततायोग्य असतो.
\end{enumerate}
हे निष्पत्ती-पद्धतीचे महत्त्वाचे गुणधर्म आहेत. यांपैकी एखादा गुणधर्म
खरा नसेल, तर निष्पत्ती-पद्धतीत दोष आहे---ती खूप जास्त गोष्टी
निष्पन्न करेल. म्हणून निष्पत्ती-पद्धतीची निर्दोषता सिद्ध करणे अत्यंत
महत्त्वाचे आहे.

हे सर्व सिद्धता-उपपत्तीय गुणधर्म कोणत्या ना कोणत्या प्रकारच्या बंद
टॅब्लो द्वारे परिभाषित केलेले असल्यामुळे, वरील (1)--(3) सिद्ध करण्यासाठी
बंद टॅब्लोच्या चिन्हार्थविषयक गुणधर्मांबद्दल काहीतरी सिद्ध करावे लागते.
प्रथम एखाद्या चिन्हांकित सूत्राची संरचनेत पूर्ती होणे म्हणजे काय ते
परिभाषित करू; मग टॅब्लो बंद असेल, तर कोणतीही संरचना त्याची सर्व गृहीतके
पूर्ण करत नाही, हे दाखवू. त्यानंतर (1)--(3) या परिणामाचे उपपरिणाम म्हणून
मिळतील.
\end{explain}

\begin{defn}
संरचना~$\Struct{M}$
चिन्हांकित सूत्र $\sFmla{\True}{\varphi}$ \emph{पूर्ण करते} तेव्हा आणि केवळ तेव्हाच, जेव्हा
$\Sat{M}{\varphi}$; आणि ते
$\sFmla{\False}{\varphi}$ पूर्ण करते तेव्हा आणि केवळ तेव्हाच, जेव्हा
$\Sat/{M}{\varphi}$. $\Struct{M}$
हा चिन्हांकित सूत्रांचा संच~$\Gamma$ पूर्ण करतो तेव्हा आणि केवळ तेव्हाच,
जेव्हा तो प्रत्येक $\sFmla{S}{\varphi} \in \Gamma$ पूर्ण करतो. $\Gamma$ पूर्ण
करणारी एखादी संरचना असेल, तर तो
\emph{पूर्ततायोग्य} आणि अन्यथा \emph{अपूर्ततायोग्य} असतो.
\end{defn}

\begin{thm}[निर्दोषता]
  \label{fol:tab:sou:thm:tableau-soundness}
  $\Gamma$ साठी बंद टॅब्लो असेल, तर $\Gamma$ अपूर्ततायोग्य असतो.
\end{thm}

\begin{proof}
टॅब्लोच्या शाखेवरील चिन्हांकित सूत्रांचा संच पूर्ततायोग्य असेल, तर
आणि केवळ तेव्हाच त्या शाखेला पूर्ततायोग्य म्हणू; तसेच टॅब्लोमध्ये
किमान एक पूर्ततायोग्य शाखा असेल, तर त्याला पूर्ततायोग्य म्हणू.

पुढील गोष्ट दाखवू: एखाद्या पूर्ततायोग्य टॅब्लोचा अनुमानाच्या कोणत्याही
एका नियमाने विस्तार केल्यास मिळणारा टॅब्लो नेहमी पूर्ततायोग्य असतो.
यावरून प्रमेय सिद्ध होईल: $\Gamma$ मधील केवळ गृहीतके असलेल्या टॅब्लो ला
अनुमानाचे नियम लावून कोणताही बंद टॅब्लो मिळतो. म्हणून $\Gamma$
पूर्ततायोग्य असेल, तर त्याच्यासाठीचा कोणताही टॅब्लो पूर्ततायोग्य असेल.
पण बंद टॅब्लो स्पष्टपणे पूर्ततायोग्य नसतो: प्रत्येक शाखेत
$\sFmla{\True}{\varphi}$ आणि $\sFmla{\False}{\varphi}$ ही दोन्ही सूत्रे असतात, आणि
कोणतीही संरचना $\varphi$ ची पूर्ती आणि अपूर्ती या दोन्ही करू शकत नाही.

आपल्याकडे पूर्ततायोग्य टॅब्लो, म्हणजे किमान एक पूर्ततायोग्य शाखा असलेला
टॅब्लो, आहे असे समजा. अनुमानाचा नियम लावल्यावर एकतर शाखेत
चिन्हांकित सूत्रे जोडली जातात, किंवा शाखा दोन भागांत विभागली जाते. त्या
नियमाच्या उपयोगाने विस्तारली न गेलेली एखादी पूर्ततायोग्य शाखा टॅब्लोमध्ये
असेल, तर विस्तारित टॅब्लोमध्येही ती पूर्ततायोग्य शाखा राहते; म्हणून
विस्तारित टॅब्लोसुद्धा पूर्ततायोग्य असतो. त्यामुळे पूर्ततायोग्य शाखेलाच नियम
लावला जातो ते प्रकरण पाहणे पुरेसे आहे.

त्या शाखेवरील चिन्हांकित सूत्रांचा संच~$\Gamma$ असू द्या, आणि ज्या
चिन्हांकित सूत्राला नियम लावला जातो ते $\sFmla{S}{\varphi} \in \Gamma$ असू
द्या. नियम लावल्यावर शाखा विभागली जात नसेल, तर विस्तारित शाखा, म्हणजे
$\Gamma$ आणि नियमाचे निष्कर्ष मिळून बनलेला संच, अजूनही पूर्ततायोग्य आहे हे
दाखवावे लागेल. नियम लावल्यावर शाखा विभागली जात असेल, तर मिळणाऱ्या दोन
शाखांपैकी किमान एक पूर्ततायोग्य आहे हे दाखवावे लागेल.

प्रथम, शाखा न विभागणारी संभाव्य अनुमाने पाहू.
\begin{enumerate}
\item $\sFmla{\True}{\lnot \psi} \in \Gamma$ ला
  $\TRule{\True}{\lnot}$ लावून शाखेचा विस्तार केला आहे असे समजा. मग
  विस्तारित शाखेवरील चिन्हांकित सूत्रांचा संच $\Gamma \cup
  \{\sFmla{\False}{\psi}\}$ असतो. समजा
  $\Sat{M}{\Gamma}$. विशेषतः,
  $\Sat{M}{\lnot \psi}$. त्यामुळे
  $\Sat/{M}{\psi}$, म्हणजे
  $\Struct{M}$ हे
  $\sFmla{\False}{\psi}$ पूर्ण करते.
\item $\sFmla{\False}{\lnot \psi} \in \Gamma$ ला
  $\TRule{\False}{\lnot}$ लावून शाखेचा विस्तार केला आहे: सराव.
\item $\sFmla{\True}{\psi \land \chi} \in \Gamma$ ला
  $\TRule{\True}{\land}$ लावून शाखेचा विस्तार केला आहे. यामुळे शाखेवर
  $\sFmla{\True}{\psi}$ आणि $\sFmla{\True}{\chi}$ ही दोन नवी
  चिन्हांकित सूत्रे मिळतात. समजा
  $\Sat{M}{\Gamma}$; विशेषतः,
  $\Sat{M}{\psi \land \chi}$. मग
  $\Sat{M}{\psi}$ आणि
  $\Sat{M}{\chi}$. म्हणजे
  $\Struct{M}$ हे
  $\sFmla{\True}{\psi}$ आणि $\sFmla{\True}{\chi}$ ही दोन्ही सूत्रे पूर्ण करते.
\item $\sFmla{\False}{\psi \lor \chi} \in \Gamma$ ला
  $\TRule{\False}{\lor}$ लावून शाखेचा विस्तार केला आहे: सराव.
\item $\sFmla{\False}{\psi \lif \chi} \in \Gamma$ ला
  $\TRule{\False}{\lif}$ लावून शाखेचा विस्तार केला आहे. यामुळे शाखेवर
  $\sFmla{\True}{\psi}$ आणि $\sFmla{\False}{\chi}$ ही दोन नवी
  चिन्हांकित सूत्रे मिळतात. समजा
  $\Sat{M}{\Gamma}$; विशेषतः,
  $\Sat/{M}{\psi \lif \chi}$. मग
  $\Sat{M}{\psi}$ आणि
  $\Sat/{M}{\chi}$. म्हणजे
  $\Struct{M}$ हे
  $\sFmla{\True}{\psi}$ आणि $\sFmla{\False}{\chi}$ ही दोन्ही सूत्रे पूर्ण करते.


\item $\sFmla{\True}{\lforall[x][\varphi(x)]} \in \Gamma$ ला
  $\TRule{\True}{\lforall}$ लावून शाखेचा विस्तार केला आहे. यामुळे शाखेवर
  नवीन चिन्हांकित सूत्र~$\sFmla{\True}{\varphi(t)}$ मिळते.
  समजा $\Sat{M}{\Gamma}$; विशेषतः,
  $\Sat{M}{\lforall[x][\varphi(x)]}$.
  \ref{fol:syn:sem:prop:quant-terms} नुसार $\Sat{M}{\varphi(t)}$. परिणामी,
  $\Struct{M}$ हे $\sFmla{\True}{\varphi(t)}$ पूर्ण करते.

\item $\sFmla{\False}{\lforall[x][\varphi(x)]} \in \Gamma$ ला
  $\TRule{\False}{\lforall}$ लावून शाखेचा विस्तार केला आहे. यामुळे शाखेवर
  नवीन चिन्हांकित सूत्र~$\sFmla{\False}{\varphi(a)}$ मिळते, जेथे $a$ हे
  $\Gamma$ मध्ये न येणारे स्थिरांक आहे. $\Gamma$ पूर्ततायोग्य असल्यामुळे
  $\Sat{M}{\Gamma}$ अशी एखादी $\Struct{M}$ असते; विशेषतः,
  $\Sat/{M}{\lforall[x][\varphi(x)]}$. आता
  $\Gamma \cup \{\sFmla{\False}{\varphi(a)}\}$ पूर्ततायोग्य आहे हे दाखवायचे आहे.
  त्यासाठी पुढीलप्रमाणे योग्य~$\Struct{M'}$ परिभाषित करू.

  \ref{fol:syn:ass:prop:sat-quant} नुसार,
  $\Sat/{M}{\lforall[x][\varphi(x)]}$ तेव्हा आणि केवळ तेव्हाच, जेव्हा काही $s$
  साठी $\Sat/{M}{\varphi(x)}[s]$. आता $\Struct{M'}$ ही $\Struct{M}$ सारखीच
  असू द्या, फक्त $\Assign{a}{M'} = s(x)$. $a$ हे $\Gamma$ मध्ये येत
  नसल्यामुळे, \ref{fol:syn:ext:cor:extensionality-sent} नुसार कोणत्याही
  $\sFmla{\True}{\chi} \in \Gamma$ साठी $\Sat{M'}{\chi}$, आणि कोणत्याही
  $\sFmla{\False}{\chi} \in \Gamma$ साठी $\Sat/{M'}{\chi}$.

  \ref{fol:syn:ext:prop:extensionality} नुसार $\Sat/{M'}{\varphi(x)}[s]$.
  \ref{fol:syn:ext:prop:ext-formulas} नुसार $\Sat/{M'}{\varphi(a)}[s]$.
  $\varphi(a)$ हे वाक्य असल्यामुळे,
  \ref{fol:syn:ass:prop:sentence-sat-true} नुसार $\Sat/{M'}{\varphi(a)}$;
  म्हणजे $\Struct{M'}$ हे $\sFmla{\False}{\varphi(a)}$ पूर्ण करते.
\item $\sFmla{\True}{\lexists[x][\psi(x)]} \in \Gamma$ ला
  $\TRule{\True}{\lexists}$ लावून शाखेचा विस्तार केला आहे: सराव.
\item $\sFmla{\False}{\lexists[x][\psi(x)]} \in \Gamma$ ला
  $\TRule{\False}{\lexists}$ लावून शाखेचा विस्तार केला आहे: सराव.
\end{enumerate}
आता शाखा विभागणारी संभाव्य अनुमाने पाहू.
\begin{enumerate}
\item $\sFmla{\False}{\psi \land \chi} \in \Gamma$ ला
  $\TRule{\False}{\land}$ लावून शाखेचा विस्तार केला आहे. यामुळे दोन शाखा
  मिळतात: डावीकडील शाखा $\sFmla{\False}{\psi}$ मधून आणि उजवीकडील शाखा
  $\sFmla{\False}{\chi}$ मधून पुढे जाते. समजा
  $\Sat{M}{\Gamma}$; विशेषतः,
  $\Sat/{M}{\psi \land \chi}$. मग एकतर
  $\Sat/{M}{\psi}$ किंवा
  $\Sat/{M}{\chi}$. पहिल्या प्रकरणात
  $\Struct{M}$ हे
  $\sFmla{\False}{\psi}$ पूर्ण करते; म्हणजे
  $\Struct{M}$ डावीकडील शाखेवरील सूत्रे पूर्ण करते.
  दुसऱ्या प्रकरणात $\Struct{M}$ हे
  $\sFmla{\False}{\chi}$ पूर्ण करते; म्हणजे
  $\Struct{M}$ उजवीकडील शाखेवरील सूत्रे पूर्ण करते.
\item $\sFmla{\True}{\psi \lor \chi} \in \Gamma$ ला
  $\TRule{\True}{\lor}$ लावून शाखेचा विस्तार केला आहे: सराव.
\item $\sFmla{\True}{\psi \lif \chi} \in \Gamma$ ला
  $\TRule{\True}{\lif}$ लावून शाखेचा विस्तार केला आहे: सराव.
\item \Cut{} लावून शाखेचा विस्तार केला आहे. यामुळे दोन शाखा मिळतात: एका
  शाखेत $\sFmla{\True}{\psi}$ आणि दुसऱ्या शाखेत $\sFmla{\False}{\psi}$ असते.
  $\Sat{M}{\Gamma}$ आणि एकतर
  $\Sat{M}{\psi}$ किंवा
  $\Sat/{M}{\psi}$ असल्यामुळे,
  $\Struct{M}$ हे डावीकडील किंवा उजवीकडील शाखा
  पूर्ण करते.
\end{enumerate}
\end{proof}


\begin{prob}
\ref{fol:tab:sou:thm:tableau-soundness} ची सिद्धता पूर्ण करा.
\end{prob}




\begin{cor}
\label{fol:tab:sou:cor:weak-soundness}
जर $\Proves \varphi$, तर $\varphi$ हे वैध आहे.
\end{cor}

\begin{cor}
\label{fol:tab:sou:cor:entailment-soundness}
जर $\Gamma \Proves \varphi$, तर $\Gamma \Entails \varphi$.
\end{cor}

\begin{proof}
  जर $\Gamma \Proves \varphi$, तर काही $\psi_1$, \dots, $\psi_n \in
  \Gamma$ साठी $\{\sFmla{\False}{\varphi}, \sFmla{\True}{\psi_1}, \dots,
  \sFmla{\True}{\psi_n}\}$ चा बंद टॅब्लो असतो.
  \ref{fol:tab:sou:thm:tableau-soundness} नुसार, प्रत्येक
  संरचना~$\Struct{M}$
  काही $\psi_i$ असत्य करते किंवा $\varphi$ सत्य करते. म्हणून
  $\Sat{M}{\Gamma}$ असेल, तर
  $\Sat{M}{\varphi}$ सुद्धा असेल.
\end{proof}

\begin{cor}
\label{fol:tab:sou:cor:consistency-soundness}
जर $\Gamma$ पूर्ततायोग्य असेल, तर तो सुसंगत आहे.
\end{cor}

\begin{proof}
प्रतिलोम विधान सिद्ध करू. $\Gamma$ सुसंगत नाही असे समजा. मग
$\psi_1$, \dots, $\psi_n \in \Gamma$ अशी वाक्ये आणि
$\{\sFmla{\True}{\psi_1}, \dots, \sFmla{\True}{\psi_n}\}$ चा बंद
टॅब्लो असतो. \ref{fol:tab:sou:thm:tableau-soundness} नुसार, प्रत्येक
$i=1$, \dots,~$n$ साठी $\Sat{M}{\psi_i}$ अशी
कोणतीही $\Struct{M}$ नसते. पण मग $\Gamma$
पूर्ततायोग्य नसतो.
\end{proof}












\section{टॅब्लो सह एकरूपता}\label{fol:tab:ide:sec}

टॅब्लोमध्ये एकरूपता असल्यास अतिरिक्त निगमन नियम आवश्यक असतात.
$\eq$ चे नियम पुढीलप्रमाणे आहेत ($t$, $t_1$ आणि $t_2$ ही बंद पदे आहेत):

\begin{defish}
\AxiomC{}
\RightLabel{$\eq$}
\UnaryInfC{\sFmla{\True}{\eq[t][t]}}
\DisplayProof
\hfill
\AxiomC{\sFmla{\True}{\eq[t_1][t_2]}}
\noLine
\UnaryInfC{\sFmla{\True}{\varphi(t_1)}}
\RightLabel{$\TRule{\True}{\eq}$}
\UnaryInfC{\sFmla{\True}{\varphi(t_2)}}
\DisplayProof
\hfill
\AxiomC{\sFmla{\True}{\eq[t_1][t_2]}}
\noLine
\UnaryInfC{\sFmla{\False}{\varphi(t_1)}}
\RightLabel{$\TRule{\False}{\eq}$}
\UnaryInfC{\sFmla{\False}{\varphi(t_2)}}
\DisplayProof
\end{defish}
लक्षात घ्या की इतर सर्व नियमांपेक्षा वेगळे, $\TRule{\True}{\eq}$ आणि
$\TRule{\False}{\eq}$ यांना शाखेवर आधीपासून \emph{दोन} चिन्हांकित
सूत्रे असणे आवश्यक आहे, म्हणजे $\sFmla{\True}{\eq[t_1][t_2]}$
आणि $\sFmla{S}{\varphi(t_1)}$ ही दोन्ही.

\begin{ex}
जर $s$ आणि $t$ ही बंद पदे असतील, तर $\eq[s][t], \varphi(s)
\Proves \varphi(t)$:
\begin{oltableau}
  [\sFmla{\False}{\varphi(t)}, just = \TAss
    [\sFmla{\True}{\eq[s][t]}, just = \TAss
      [\sFmla{\True}{\varphi(s)}, just = \TAss
        [\sFmla{\True}{\varphi(t)}, just={\TRule{\True}{\eq}[2, 3]}, close]
      ]
    ]
  ]
\end{oltableau}
याला एकरूप वस्तूंच्या प्रतिस्थापनाचे तत्त्व किंवा लाइब्निझचा नियम म्हणून
ओळखले जाऊ शकते.

टॅब्लो द्वारे $\eq$ सममित आहे, म्हणजे $\eq[s_1][s_2]
\Proves \eq[s_2][s_1]$, हे सिद्ध होते:
\begin{oltableau}
  [\sFmla{\False}{\eq[s_2][s_1]}, just = \TAss
    [\sFmla{\True}{\eq[s_1][s_2]}, just = \TAss
      [\sFmla{\True}{\eq[s_1][s_1]}, just = {$\eq$}
        [\sFmla{\True}{\eq[s_2][s_1]}, just = {\TRule{\True}{\eq}[2, 3]}, close]
      ]
    ]
  ]
\end{oltableau}
येथे ओळ~$2$ ही $\TRule{\True}{\eq}$ साठीचे पहिले आधारभूत
सूत्र, $\sFmla{\True}{\eq[s_1][s_2]}$, आहे. ओळ~$3$ हे
$\sFmla{\True}{\varphi(s_1)}$ या रूपातील दुसरे आधारविधान आहे---$\varphi(x)$ हे
$\eq[x][s_1]$ आहे असे विचार करा; मग $\varphi(s_1)$ म्हणजे
$\eq[s_1][s_1]$ आणि $\varphi(s_2)$ म्हणजे $\eq[s_2][s_1]$.

तसेच $\eq$ संक्रमक आहे, म्हणजे $\eq[s_1][s_2],
\eq[s_2][s_3] \Proves \eq[s_1][s_3]$, हेही ते सिद्ध करतात:
\begin{oltableau}
  [\sFmla{\False}{\eq[s_1][s_3]}, just = \TAss
    [\sFmla{\True}{\eq[s_1][s_2]}, just = \TAss
      [\sFmla{\True}{\eq[s_2][s_3]}, just = \TAss
        [\sFmla{\True}{\eq[s_1][s_3]}, just = {\TRule{\True}{\eq}[3, 2]}, close]
      ]
    ]
  ]
\end{oltableau}
या टॅब्लोमध्ये $\TRule{\True}{\eq}$ चे पहिले आधारभूत सूत्र
ओळ~$3$ वरील $\sFmla{\True}{\eq[s_2][s_3]}$ आहे ($s_2$ ही~$t_1$ ची,
आणि $s_3$ ही~$t_2$ ची भूमिका बजावतात). $\sFmla{\True}{\varphi(s_2)}$
या रूपातील दुसरे आधारविधान ओळ~$2$ आहे. येथे $\varphi(x)$ हे
$\eq[s_1][x]$ आहे असे विचार करा; त्यामुळे $\varphi(s_2)$ हे
$\eq[s_1][s_2]$ (म्हणजे ओळ~$2$) आणि निष्कर्षातील
सूत्र~$\eq[s_1][s_3]$ हे $\varphi(s_3)$ होते.
\end{ex}

\begin{prob}
पुढील गोष्टींसाठी बंद टॅब्लो द्या:
\begin{enumerate}
\item $\sFmla{\False}{\lforall[x][\lforall[y][((x = y \land \varphi(x))
      \lif \varphi(y))]]}$
\item $\sFmla{\False}{\lexists[x][(\varphi(x) \land
    \lforall[y][(\varphi(y) \lif y = x)])]}$,\\
  $\sFmla{\True}{\lexists[x][\varphi(x)] \land
    \lforall[y][\lforall[z][((\varphi(y) \land \varphi(z)) \lif y = z)]]}$
\end{enumerate}
\end{prob}












\section{एकरूपता सह निर्दोषता}\label{fol:tab:sid:sec}

\begin{prop}
एकरूपतेचे नियम असलेले टॅब्लो निर्दोष आहेत: कोणताही बंद टॅब्लो
  पूर्ततायोग्य नसतो.
\end{prop}

\begin{proof}
मागीलप्रमाणे एवढेच दाखवायचे आहे की टॅब्लोची एखादी शाखा पूर्ततायोग्य
असेल, तर तिला~$\eq$ चा कोणताही नियम लावून मिळणारी शाखासुद्धा पूर्ततायोग्य
असते. शाखेवरील चिन्हांकित सूत्रांचा संच~$\Gamma$ असू द्या आणि
$\Struct{M}$ ही $\Gamma$ पूर्ण करणारी संरचना असू द्या.

शाखेचा~$\eq$ वापरून विस्तार केला आहे, म्हणजे चिन्हांकित
सूत्र~$\sFmla{\True}{\eq[t][t]}$ जोडले आहे असे समजा. स्पष्टपणे
$\Sat{M}{\eq[t][t]}$, म्हणून $\Struct{M}$ हे
$\Gamma \cup \{\sFmla{\True}{\eq[t][t]}\}$ सुद्धा पूर्ण करते.

$\TRule{\True}{\eq}$ वापरून शाखेचा विस्तार केला असेल, तर
सूत्र~$\sFmla{S}{\varphi(t_2)}$ जोडतो; परंतु $\Gamma$ मध्ये
$\sFmla{\True}{\eq[t_1][t_2]}$ आणि $\sFmla{\True}{\varphi(t_1)}$ ही दोन्ही
असतात. म्हणून $\Sat{M}{\eq[t_1][t_2]}$ आणि $\Sat{M}{\varphi(t_1)}$.
$s$ हा $s(x) = \Value{t_1}{M}$ असलेला चल-विन्यास असू द्या.
\ref{fol:syn:ass:prop:sentence-sat-true} नुसार
$\Sat{M}{\varphi(t_1)}[s]$. $\varAssign{s}{s}{x}$ असल्यामुळे,
\ref{fol:syn:ext:prop:ext-formulas} नुसार $\Sat{M}{\varphi(x)}[s]$.
$\Sat{M}{\eq[t_1][t_2]}$ असल्यामुळे
$\Value{t_1}{M} = \Value{t_2}{M}$, आणि म्हणून
$s(x) = \Value{t_2}{M}$. \ref{fol:syn:ext:prop:ext-formulas} पुन्हा
लावल्याने $\Sat{M}{\varphi(t_2)}[s]$ सुद्धा मिळते.
\ref{fol:syn:ass:prop:sentence-sat-true} नुसार
$\Sat{M}{\varphi(t_2)}$. $\TRule{\False}{\eq}$ चे प्रकरण याचप्रमाणे आहे.
\end{proof}



\endgroup
\chapter{स्वयंसिद्धकीय निष्पत्ती}\label{fol:axd::chap}
\begin{editorial}
  या प्रकरणातील मजकूर कोणते संयोजक आणि संख्यापक आदिम किंवा परिभाषित
  आहेत हे दर्शवणाऱ्या विविध टॅग्जचे पालन करतो याची खात्री करण्याचा अद्याप
  कोणताही प्रयत्न केलेला नाही: यांपैकी $\liff$ परिभाषित आहे, तर बाकी
  सर्व आदिम आहेत असे गृहीत धरले आहे. FOL टॅग सत्य
  असल्यास आपण संख्यापकांसह आवृत्ती तयार करतो; अन्यथा संख्यापकांविना
  आवृत्ती तयार करतो.
\end{editorial}









\section{नियम आणि निष्पत्ती}\label{fol:axd:rul:sec}

\begin{explain}
  स्वयंसिद्धकीय निष्पत्त्या ही तर्कशास्त्रासाठी कदाचित सर्वांत सोपी
  निष्पत्ती-पद्धत आहे. निष्पत्ती म्हणजे केवळ सूत्रांचा
  अनुक्रम. निष्पत्ती म्हणून गणले जाण्यासाठी अनुक्रमातील प्रत्येक
  सूत्र एकतर स्वयंसिद्धकाचे उदाहरण असले पाहिजे, किंवा अनुक्रमात
  त्याच्या आधी येणाऱ्या एक वा अधिक सूत्रांपासून निगमन नियमाने
  निष्पन्न झाले पाहिजे. निष्पत्ती तिचे शेवटचे सूत्र
  निष्पन्न करते.
\end{explain}

\begin{defn}[निष्पन्नता]
जर $\Gamma$ हा $\Lang L$ मधील सूत्रांचा संच असेल, तर
$\Gamma$ पासूनची \emph{निष्पत्ती} म्हणजे सूत्रांचा सांत अनुक्रम
$\varphi_1$, \dots,~$\varphi_n$ असा की प्रत्येक $i \le n$ साठी पुढीलपैकी एक अट
पूर्ण होते:
\begin{enumerate}
\item $\varphi_i \in \Gamma$; किंवा
\item $\varphi_i$ हे स्वयंसिद्धक आहे; किंवा
\item $j < i$ (आणि $k < i$) असलेल्या एखाद्या $\varphi_j$ (आणि $\varphi_k$)
  पासून निगमन नियमाने $\varphi_i$ निष्पन्न होते.
\end{enumerate}
\end{defn}

योग्य निष्पत्ती म्हणून कशाची गणना होते हे आपण कोणते निगमन नियम
मान्य करतो (आणि अर्थातच कोणती सूत्रे स्वयंसिद्धके मानतो) यावर अवलंबून
असते. निगमन नियम म्हणजे असे जर-तर विधान जे विशिष्ट अटींमध्ये
निष्पत्तीमधील पायरी~$A_i$ ही निगमनाची योग्य पायरी आहे असे सांगते.

\begin{defn}[निगमन नियम]
\emph{निगमन नियम} हा $\Gamma$ पासूनच्या निष्पत्तीमध्ये निगमनाची
योग्य पायरी म्हणून कशाची गणना होते यासाठी पुरेशी अट देतो.
\end{defn}

उदाहरणार्थ, $\varphi \in \Gamma$ असलेला कोणताही एक-घटकी अनुक्रम $\varphi$ हा
आपोआपच निष्पत्ती म्हणून गणला जात असल्यामुळे, पुढील नियम हा
अतिशय सोपा निगमन नियम असू शकतो:
\begin{quote}
  जर $\varphi \in \Gamma$, तर $\Gamma$ पासूनच्या कोणत्याही निष्पत्तीमध्ये $\varphi$ ही नेहमीच निगमनाची योग्य पायरी असते.
\end{quote}
त्याचप्रमाणे, $\varphi$ हे स्वयंसिद्धकांपैकी एक असेल, तर एकटे $\varphi$ ही
निष्पत्ती असते; म्हणून पुढील विधानही निगमन नियम आहे:
\begin{quote}
  जर $\varphi$ हे स्वयंसिद्धक असेल, तर $\varphi$ ही निगमनाची योग्य पायरी असते.
\end{quote}
विचाराधीन पायरीच्या आधी दिसणाऱ्या सूत्रांचा निगमन नियम आधार घेतो
तेव्हा गोष्ट अधिक रोचक होते. पुढील नियमाला \emph{विध्यात्मक अनुमान}
म्हणतात:
\begin{quote}
  जर $\psi \lif \varphi$ आणि $\psi$ हे निष्पत्तीमध्ये यापूर्वी आढळत असतील,
  तर~$\varphi$ ही निगमनाची योग्य पायरी असते.
\end{quote}
हा एकमेव निगमन नियम असेल, तर वर दिलेल्या निष्पत्तीच्या व्याख्येचा
अर्थ असा होतो: प्रत्येक $i \le n$ साठी पुढीलपैकी एक अट पूर्ण होत असेल
तेव्हा आणि केवळ तेव्हाच $\varphi_1$, \dots,~$\varphi_n$ ही निष्पत्ती असते:
\begin{enumerate}
\item $\varphi_i \in \Gamma$; किंवा
\item $\varphi_i$ हे स्वयंसिद्धक आहे; किंवा
\item एखाद्या $j < i$ साठी $\varphi_j$ हे $\psi \lif \varphi_i$ आहे, आणि एखाद्या
  $k < i$ साठी $\varphi_k$ हे~$\psi$ आहे.
\end{enumerate}
शेवटचे कलम असे सांगते की~$\varphi_j$ ($\psi \lif \varphi_i$) आणि $\varphi_k$ ($\psi$)
यांपासून विध्यात्मक अनुमानाने $\varphi_i$ निष्पन्न होते. आपण $1$ पासून~$n$
पर्यंत जाऊ शकलो आणि प्रत्येक वेळी~$\Gamma$ मध्ये असलेले, स्वयंसिद्धक
असलेले किंवा निगमन नियमाने योग्य निगमन-पायरी ठरणारे सूत्र~$\varphi_i$
मिळाले, तर संपूर्ण अनुक्रम योग्य निष्पत्ती म्हणून गणला जातो.

\begin{defn}[निष्पन्नता]
जर $\Gamma$ पासून सुरू होऊन $\varphi$ वर संपणारी निष्पत्ती असेल, तर
सूत्र~$\varphi$ हे $\Gamma$ पासून \emph{निष्पन्न करता येण्याजोगे} असते; हे आपण
$\Gamma \Proves \varphi$ असे लिहितो.
\end{defn}

\begin{defn}[प्रमेये]
रिक्त संचापासून $\varphi$ ची निष्पत्ती असेल, तर सूत्र~$\varphi$ हे
\emph{प्रमेय} असते. $\varphi$ हे प्रमेय असेल तर आपण $\Proves \varphi$ लिहितो आणि
ते प्रमेय नसेल तर $\Proves/ \varphi$ लिहितो.
\end{defn}













\section{विधानीय संयोजकांसाठी स्वयंसिद्धके आणि नियम}\label{fol:axd:prp:sec}

\begin{defn}[स्वयंसिद्धके]
विधानीय संयोजकांसाठीच्या \emph{स्वयंसिद्धकांचा} संच $\PAx$ हा पुढील रूपांच्या
सर्व सूत्रांचा बनलेला आहे:
\begin{align}
  & (\varphi \land \psi) \lif \varphi \label{fol:axd:prp:ax:land1}\\
  & (\varphi \land \psi) \lif \psi \label{fol:axd:prp:ax:land2}\\
  & \varphi \lif (\psi \lif (\varphi \land \psi)) \label{fol:axd:prp:ax:land3}\\
  & \varphi \lif (\varphi \lor \psi) \label{fol:axd:prp:ax:lor1}\\
  & \varphi \lif (\psi \lor \varphi) \label{fol:axd:prp:ax:lor2}\\
  & (\varphi \lif \chi) \lif ((\psi \lif \chi) \lif ((\varphi \lor \psi) \lif \chi)) \label{fol:axd:prp:ax:lor3}\\
  & \varphi \lif (\psi \lif \varphi) \label{fol:axd:prp:ax:lif1}\\
  & (\varphi \lif (\psi \lif \chi)) \lif ((\varphi \lif \psi) \lif (\varphi \lif \chi)) \label{fol:axd:prp:ax:lif2}\\
  & (\varphi \lif \psi) \lif ((\varphi \lif \lnot \psi) \lif \lnot \varphi) \label{fol:axd:prp:ax:lnot1}\\
  & \lnot \varphi \lif (\varphi \lif \psi) \label{fol:axd:prp:ax:lnot2}\\
  & \ltrue \label{fol:axd:prp:ax:ltrue}\\
  & \lfalse \lif \varphi \label{fol:axd:prp:ax:lfalse1}\\
  & (\varphi \lif \lfalse) \lif \lnot \varphi \label{fol:axd:prp:ax:lfalse2}\\
  & \lnot\lnot \varphi \lif \varphi \label{fol:axd:prp:ax:dne}
\end{align}
\end{defn}

\begin{defn}[विध्यात्मक अनुमान]
  जर $\psi$ आणि $\psi \lif \varphi$ हे निष्पत्तीमध्ये आधीच आलेले असतील,
  तर $\varphi$ ही निगमनाची योग्य पायरी असते.
\end{defn}

विध्यात्मक अनुमान या नियमाचा संक्षेप आपण ``\MP'' असा करू.












\section{संख्यापकांसाठी स्वयंसिद्धके आणि नियम}\label{fol:axd:qua:sec}

\begin{defn}[संख्यापकांसाठी स्वयंसिद्धके]
संख्यापकांचे नियमन करणारी \emph{स्वयंसिद्धके} म्हणजे पुढील रूपांचे सर्व
दृष्टांत होत:
\begin{align}
\label{fol:axd:qua:ax:q1} & \lforall[x][\psi] \lif \psi(t), \\
\label{fol:axd:qua:ax:q2} & \psi(t) \lif \lexists[x][\psi].
\end{align}
कोणत्याही बंद पद~$t$ साठी.
\end{defn}

\begin{defn}[संख्यापकांसाठी नियम]
\item जर $\psi \lif \varphi(a)$ हे निष्पत्तीमध्ये आधीच आलेले असेल आणि
  $a$ ची~$\Gamma$ किंवा~$\psi$ मध्ये आस्थिति नसेल, तर $\psi \lif
  \lforall[x][\varphi(x)]$ ही निगमनाची योग्य पायरी असते.
\item जर $\varphi(a) \lif \psi$ हे निष्पत्तीमध्ये आधीच आलेले असेल आणि
  $a$ ची~$\Gamma$ किंवा~$\psi$ मध्ये आस्थिति नसेल, तर
  $\lexists[x][\varphi(x)] \lif \psi$ ही निगमनाची योग्य पायरी असते.
\end{defn}

यांपैकी कोणत्याही नियमाचा संक्षेप आपण ``\QR'' असा करू.












\section{निष्पत्ती: उदाहरणे}\label{fol:axd:pro:sec}


\begin{ex}
आपल्याला $(\lnot \theta \lor \varphi_{E}) \lif (\theta \lif
\varphi_{E})$ हे सिद्ध करायचे आहे असे समजा. स्पष्टपणे, हे आपल्या कोणत्याही
स्वयंसिद्धकाचे उदाहरण नाही; म्हणून ते निष्पन्न करण्यासाठी \MP{} नियम
वापरावा लागेल. आपला एकमेव नियम~MP आहे; $\varphi$ आणि $\varphi \lif \psi$ दिले असता
या नियमाने~$\psi$ ला समर्थित करता येते. एक युक्ती अशी असेल: $\varphi$ म्हणून
$\lnot \theta$, $\psi$ म्हणून $\varphi_{E}$ आणि $\chi$ म्हणून $\theta \lif \varphi_{E}$ घेऊन
\ref{fol:axd:prp:ax:lor3} वापरणे, म्हणजे पुढील उदाहरण वापरणे:
\[(\lnot \theta \lif (\theta \lif \varphi_{E})) \lif ((\varphi_{E} \lif (\theta \lif \varphi_{E})) \lif ((\lnot
\theta \lor \varphi_{E}) \lif (\theta \lif \varphi_{E}))).
\]
का? MP चे दोन उपयोग केल्यावर शेवटचा भाग मिळतो, आणि आपल्याला नेमके तेच
हवे आहे. तसेच $\lnot \theta \lif (\theta \lif \varphi_{E})$ हे
\ref{fol:axd:prp:ax:lnot2} चे उदाहरण आणि $\varphi_{E} \lif (\theta \lif \varphi_{E})$ हे
\ref{fol:axd:prp:ax:lif1} चे उदाहरण आहे, हे सहज दिसते. म्हणून आपली
निष्पत्ती अशी आहे:
\begin{derivation}
  1. &  $\lnot \theta \lif (\theta \lif \varphi_{E})$ & \ref{fol:axd:prp:ax:lnot2} \\
  2. & $(\lnot \theta \lif (\theta \lif \varphi_{E})) \lif {}$\\
  &\qquad $((\varphi_{E} \lif (\theta \lif \varphi_{E})) \lif ((\lnot \theta \lor \varphi_{E}) \lif (\theta \lif \varphi_{E})))$ & \ref{fol:axd:prp:ax:lor3}\\
  3. & $(\varphi_{E} \lif (\theta \lif \varphi_{E})) \lif ((\lnot \theta \lor \varphi_{E}) \lif (\theta \lif \varphi_{E}))$ & 1, 2, \MP\\
  4. & $\varphi_{E} \lif (\theta \lif \varphi_{E})$ & \ref{fol:axd:prp:ax:lif1}\\
  5. & $(\lnot \theta \lor \varphi_{E}) \lif (\theta \lif \varphi_{E})$ & 3, 4, \MP
\end{derivation}
\end{ex}

\begin{ex}\label{fol:axd:pro:ex:identity}
$\theta \lif \theta$ ची निष्पत्ती शोधून पाहू. हे स्वयंसिद्धकाचे उदाहरण
नाही; म्हणून ते निष्पन्न करण्यासाठी आपल्याला \MP{} वापरावा लागेल.
\ref{fol:axd:prp:ax:lif1} हे~$\varphi \lif \psi$ या रूपाचे स्वयंसिद्धक आहे आणि त्याला
आपण~\MP लावू शकतो. त्याचा उपयोग व्हायचा असेल, तर अर्थातच या वेळी
\MP{} ने योग्य पायरी म्हणून समर्थित होणारे $\psi$ हे~$\theta \lif \theta$ असले
पाहिजे, कारण आपल्याला तेच निष्पन्न करायचे आहे. त्यामुळे $\varphi$ सुद्धा
$\theta$ असले पाहिजे; म्हणजे आपण \ref{fol:axd:prp:ax:lif1} चे पुढील उदाहरण पाहू शकतो:
\[\theta \lif (\theta \lif \theta)
\]
\MP लावण्यासाठी आपल्याला संबंधित दुसरे आधारविधान, म्हणजे~$\varphi$, सुद्धा
समर्थित करावे लागेल. पण आपल्या प्रकरणात ते~$\theta$ असेल, आणि एकटे~$\theta$
निष्पन्न करता येणार नाही. म्हणून आपल्याला वेगळी युक्ती हवी.

केवळ~$\lif$ असलेले दुसरे स्वयंसिद्धक म्हणजे \ref{fol:axd:prp:ax:lif2}, म्हणजे
\[(\varphi \lif (\psi \lif \chi)) \lif ((\varphi \lif \psi) \lif (\varphi \lif \chi))
\]
\MP{} दोनदा लावून आपण शेवटच्या अंतःस्थ सशर्त विधानापर्यंत पोहोचू शकतो.
पुन्हा, याचा अर्थ असा की \ref{fol:axd:prp:ax:lif2} चे असे उदाहरण आपल्याला हवे
ज्यात $\varphi \lif \chi$ हे $\theta \lif \theta$, म्हणजे आपण साध्य करू पाहत असलेले
सूत्र, असेल. मग अर्थातच $\varphi$ आणि $\chi$ ही दोन्ही~$\theta$ असतात.
$\varphi \lif (\psi \lif \chi)$ आणि $\varphi \lif \psi$, म्हणजे आपल्या प्रकरणात
$\theta \lif (\psi \lif \theta)$ आणि $\theta \lif \psi$, ही दोन्हीही निष्पन्न करता येण्याजोगे
होतील अशा रीतीने~$\psi$ कसे निवडावे? आपण $\psi$ काहीही निवडले, तरी यांपैकी
पहिले हे आधीच \ref{fol:axd:prp:ax:lif1} चे उदाहरण आहे. आणि $\psi$ हे
$(\theta \lif \theta)$ असेल, तर $\theta \lif \psi$ हे \ref{fol:axd:prp:ax:lif1} चे आणखी
एक उदाहरण असेल. म्हणून आपली निष्पत्ती अशी आहे:
\begin{derivation}
  1. & $\theta \lif ((\theta \lif \theta) \lif \theta)$ & \ref{fol:axd:prp:ax:lif1}\\
  2. & $(\theta \lif ((\theta \lif \theta) \lif \theta)) \lif {}$\\
  & \qquad $((\theta \lif (\theta \lif \theta)) \lif (\theta \lif \theta))$ & \ref{fol:axd:prp:ax:lif2}\\
  3. & $(\theta \lif (\theta \lif \theta)) \lif (\theta \lif \theta)$ & 1, 2, \MP\\
  4. & $\theta \lif (\theta \lif \theta)$ & \ref{fol:axd:prp:ax:lif1}\\
  5. & $\theta \lif \theta$ & 3, 4, \MP
\end{derivation}
\end{ex}

\begin{ex}\label{fol:axd:pro:ex:chain}
कधी कधी काही इतर सूत्रे~$\Gamma$ पासून एखाद्या सूत्राची
निष्पत्ती आहे हे आपल्याला दाखवायचे असते. उदाहरणार्थ,
$\Gamma = \{\varphi \lif \psi, \psi \lif \chi\}$ पासून $\varphi \lif \chi$ हे
निष्पन्न करता येते, हे दाखवू.
\begin{derivation}
  1. & $\varphi \lif \psi$ & \Hyp\\
  2. & $\psi \lif \chi$ & \Hyp\\
  3. & $(\psi \lif \chi) \lif (\varphi \lif (\psi \lif \chi))$ & \ref{fol:axd:prp:ax:lif1} \\
  4. & $\varphi \lif (\psi \lif \chi)$ & 2, 3, \MP\\
  5. & $(\varphi \lif (\psi \lif \chi)) \lif {}$\\
  & \qquad $((\varphi \lif \psi) \lif (\varphi \lif \chi))$ & \ref{fol:axd:prp:ax:lif2}\\
  6. & $((\varphi \lif \psi) \lif (\varphi \lif \chi))$ & 4, 5, \MP\\
  7. & $\varphi \lif \chi$ & 1, 6, \MP
\end{derivation}
``\Hyp'' (``hypothesis'', म्हणजे ``गृहीतक'') असे लेबल असलेल्या ओळी
त्या ओळीवरील सूत्र हा~$\Gamma$ चा घटक असल्याचे दर्शवतात.
\end{ex}

\begin{prop}
  \label{fol:axd:pro:prop:chain} जर $\Gamma \Proves \varphi \lif \psi$ आणि $\Gamma
  \Proves \psi \lif \chi$, तर $\Gamma \Proves \varphi \lif \chi$
\end{prop}

\begin{proof}
  $\Gamma \Proves \varphi \lif \psi$ आणि $\Gamma \Proves \psi \lif
  \chi$ असे समजा. मग~$\Gamma$ पासून $\varphi \lif \psi$ ची निष्पत्ती
  असते; तसेच~$\Gamma$ पासून $\psi \lif \chi$ चीही निष्पत्ती असते.
  त्यांची जोडणी करून त्या एकाच निष्पत्तीमध्ये एकत्र करा. आता मागील
  उदाहरणातील निष्पत्तीच्या ओळी 3--7 जोडा. ही~$\Gamma$ पासूनची
  निष्पत्ती आहे, आणि नव्या निष्पत्तीची शेवटची ओळ
  $\varphi \lif \chi$ आहे. लक्षात घ्या की ओळी 4 आणि~7 साठी दिलेली
  समर्थने पुढील बदलांनंतरही वैध राहतात: ओळ क्रमांक~1 चा निर्देश $\varphi
  \lif \psi$ च्या निष्पत्तीच्या शेवटच्या ओळीच्या निर्देशाने आणि ओळ
  क्रमांक~2 चा निर्देश $\psi \lif \chi$ च्या निष्पत्तीच्या शेवटच्या
  ओळीच्या निर्देशाने बदला.
\end{proof}

\begin{prob}
स्वयंसिद्धकांपासूनच्या निष्पत्त्या मांडून पुढील विधाने सिद्ध करा:
\begin{enumerate}
\item $(\varphi \land \psi) \lif (\psi \land \varphi)$
\item $((\varphi \land \psi) \lif \chi) \lif (\varphi \lif (\psi \lif \chi))$
\item $\lnot(\varphi \lor \psi) \lif \lnot \varphi$
\end{enumerate}
\end{prob}














\section{संख्यापकांसह निष्पत्ती}\label{fol:axd:prq:sec}

\begin{ex}
$(\lforall[x][\varphi(x)] \land
\lforall[y][\psi(y)]) \lif \lforall[x][(\varphi(x) \land \psi(x))]$ ची
निष्पत्ती देऊ.

प्रथम, लक्षात घ्या की
\begin{align*}
  (\lforall[x][\varphi(x)] \land \lforall[y][\psi(y)]) & \lif \lforall[x][\varphi(x)]\\
  \intertext{हे \ref{fol:axd:prp:ax:land1} चे उदाहरण आहे, आणि}
  \lforall[x][\varphi(x)] & \lif \varphi(a) \\
  \intertext{हे \ref{fol:axd:qua:ax:q1} चे उदाहरण आहे. म्हणून \ref{fol:axd:pro:prop:chain} वरून आपल्याला माहीत आहे की}
  (\lforall[x][\varphi(x)] \land \lforall[y][\psi(y)]) & \lif \varphi(a) \\
  \intertext{हे निष्पन्न करता येण्याजोगे आहे. त्याचप्रमाणे,}
  (\lforall[x][\varphi(x)] \land \lforall[y][\psi(y)]) & \lif \lforall[y][\psi(y)] \qquad\text{आणि}\\
    \lforall[y][\psi(y)] & \lif \psi(a)\\
    \intertext{ही अनुक्रमे \ref{fol:axd:prp:ax:land2} आणि \ref{fol:axd:qua:ax:q1} यांची उदाहरणे असल्यामुळे,}
    (\lforall[x][\varphi(x)] \land \lforall[y][\psi(y)]) & \lif \psi(a)\\
    \intertext{हे \ref{fol:axd:pro:prop:chain} मुळे सिद्धतायोग्य आहे. \ref{fol:axd:prp:ax:land3} चे योग्य उदाहरण आणि~\MP चे दोन उपयोग करून आपल्याला दिसते की}
    (\lforall[x][\varphi(x)] \land \lforall[y][\psi(y)]) & \lif (\varphi(a) \land \psi(a))\\
    \intertext{हे सिद्धतायोग्य आहे. आता \QR{} लावून पुढील सूत्र मिळते:}
(\lforall[x][\varphi(x)] \land \lforall[y][\psi(y)]) & \lif \lforall[x][(\varphi(x) \land \psi(x))].
\end{align*}
\end{ex}












\section{सिद्धता-उपपत्तीय संकल्पना}\label{fol:axd:ptn:sec}

\begin{explain}
जशा आपण अनेक महत्त्वाच्या चिन्हार्थविषयक संकल्पना
(वैधता, तार्किक निष्पन्नता, पूर्ततायोग्यता)
परिभाषित केल्या आहेत, तशाच आता त्यांना अनुरूप
\emph{सिद्धता-उपपत्तीय संकल्पना} परिभाषित करू. या संकल्पना संरचना
मध्ये वाक्यांची पूर्ति होते की नाही याचा आधार घेऊन परिभाषित केल्या
जात नाहीत; त्याऐवजी विशिष्ट सूत्रांची निष्पन्नता किंवा
अनिष्पन्नता यांचा आधार घेतला जातो. या दोन प्रकारच्या संकल्पना
एकमेकांशी जुळतात, हा एक महत्त्वाचा शोध होता. त्या जुळतात, हाच
\emph{निर्दोषता प्रमेय} आणि \emph{संपूर्णता प्रमेय} यांचा आशय आहे.
\end{explain}

\begin{defn}[निष्पन्नता]
सूत्र~$\varphi$ हे $\Gamma$ \emph{पासून निष्पन्न करता येण्याजोगे} असते, म्हणजे
$\Gamma \Proves \varphi$, जर~$\Gamma$ पासून सुरू होऊन~$\varphi$ वर संपणारी
निष्पत्ती असेल.
\end{defn}

\begin{defn}[प्रमेये]
रिक्त संचापासून $\varphi$ ची निष्पत्ती असेल, तर सूत्र~$\varphi$ हे
\emph{प्रमेय} असते. $\varphi$ हे प्रमेय असेल तर आपण $\Proves \varphi$ लिहितो आणि
ते प्रमेय नसेल तर $\Proves/ \varphi$ लिहितो.
\end{defn}

\begin{defn}[सुसंगतता]
सूत्रांचा संच~$\Gamma$ हा $\Gamma\Proves/ \lfalse$ असेल तेव्हा आणि
केवळ तेव्हाच \emph{सुसंगत} असतो; अन्यथा तो \emph{विसंगत} असतो.
\end{defn}

\begin{prop}[परावर्तिता]
\label{fol:axd:ptn:prop:reflexivity}
$\varphi \in \Gamma$ असेल, तर $\Gamma \Proves \varphi$.
\end{prop}

\begin{proof}
  एकटे सूत्र~$\varphi$ ही~$\Gamma$ पासून $\varphi$ ची निष्पत्ती आहे.
\end{proof}

\begin{prop}[एकस्वनिकता]
\label{fol:axd:ptn:prop:monotonicity}
$\Gamma \subseteq \Delta$ आणि $\Gamma \Proves \varphi$ असेल, तर $\Delta
\Proves \varphi$.
\end{prop}

\begin{proof}
$\Gamma$ पासून $\varphi$ ची कोणतीही निष्पत्ती ही~$\Delta$ पासून $\varphi$ चीही
निष्पत्ती असते.
\end{proof}

\begin{prop}[संक्रमकता]
\label{fol:axd:ptn:prop:transitivity}
$\Gamma \Proves \varphi$ आणि $\{\varphi\} \cup \Delta \Proves
\psi$ असेल, तर $\Gamma \cup \Delta \Proves \psi$.
\end{prop}

\begin{proof}
  $\{\varphi\} \cup \Delta \Proves \psi$ असे समजा. मग~$\{\varphi\} \cup
  \Delta$ पासून $\psi_1$, \dots, $\psi_l = \psi$ ही निष्पत्ती असते.
  त्या निष्पत्तीमधील काही पायऱ्या योग्य असतात, कारण एखादा नियम
  त्यांपूर्वीच्या~$\psi_i = \varphi$ या ओळीचा निर्देश करतो. गृहीतकानुसार,
  $\Gamma$ पासून~$\varphi$ ची निष्पत्ती असते, म्हणजे निष्पत्ती
  $\varphi_1$, \dots, $\varphi_k = \varphi$ अशी असते की प्रत्येक $\varphi_i$ हे स्वयंसिद्धक,
  $\Gamma$ चा घटक, किंवा निगमन नियमाने योग्य ठरलेले असते. आता
  पुढील अनुक्रम विचारात घ्या:
  \[  \varphi_1, \dots, \varphi_k = \varphi, \psi_1, \dots, \psi_l = \psi.
  \]
  ही $\Gamma \cup \Delta$ पासून~$\psi$ ची योग्य निष्पत्ती आहे, कारण
  प्रत्येक $B_i = \varphi$ ला आता~$\varphi_k = \varphi$ ला समर्थित करणाऱ्या त्याच नियमाने
  समर्थन मिळते.
\end{proof}

लक्षात घ्या की विशेषतः याचा अर्थ असा होतो: $\Gamma \Proves \varphi$ आणि $\varphi
\Proves \psi$ असेल, तर $\Gamma \Proves \psi$. तसेच $\varphi_1,
\dots, \varphi_n \Proves \psi$ असेल आणि प्रत्येक~$i$ साठी $\Gamma \Proves \varphi_i$
असेल, तर $\Gamma \Proves \psi$, हेसुद्धा यावरून मिळते.

\begin{prop}
\label{fol:axd:ptn:prop:incons}
$\Gamma$ विसंगत असतो तेव्हा आणि केवळ तेव्हाच, जेव्हा प्रत्येक~$\varphi$ साठी
$\Gamma \Proves \varphi$.
\end{prop}

\begin{proof}
सराव.
\end{proof}


\begin{prob}
\ref{fol:axd:ptn:prop:incons} सिद्ध करा.
\end{prob}




\begin{prop}[संहतता]
\label{fol:axd:ptn:prop:proves-compact}
  \begin{enumerate}
  \item $\Gamma \Proves \varphi$ असेल, तर एखादा सांत उपसंच $\Gamma_0
    \subseteq \Gamma$ असा असतो की $\Gamma_0 \Proves \varphi$.
  \item $\Gamma$ चा प्रत्येक सांत उपसंच सुसंगत असेल, तर $\Gamma$ सुसंगत
    असतो.
  \end{enumerate}
\end{prop}

\begin{proof}
  \begin{enumerate}
    \item $\Gamma \Proves \varphi$ असेल, तर सूत्रांचा सांत अनुक्रम
      $\varphi_1$, \dots,~$\varphi_n$ असा असतो की $\varphi \ident \varphi_n$ आणि प्रत्येक
      $\varphi_i$ हे एकतर तर्कशास्त्रीय स्वयंसिद्धक, $\Gamma$ चा घटक,
      किंवा आधीच्या सूत्रांपासून विध्यात्मक अनुमानाने निष्पन्न झालेले
      असते. $\Gamma$ मध्ये असलेल्या त्या $\varphi_i$ चा संच~$\Gamma_0$ म्हणून
      घ्या. मग ही निष्पत्ती त्याचप्रमाणे~$\Gamma_0$ पासूनची
      निष्पत्ती असते, आणि म्हणून $\Gamma_0 \Proves \varphi$.
    \item (1) मध्ये $\varphi
      \ident \lfalse$ हे विशेष प्रकरण घेतल्यावर मिळणारे हे निषेधात्मक उलट
      विधान आहे.
\end{enumerate}
\end{proof}















\section{निगमन प्रमेय}\label{fol:axd:ded:sec}

आपण पाहिल्याप्रमाणे, स्वयंसिद्धकीय पद्धतीत निष्पत्त्या देणे किचकट
असते आणि निष्पत्त्या शोधणे कठीण असू शकते. सूत्रांच्या लांब
याद्या प्रत्यक्ष लिहिण्याऐवजी काही सोपे निष्कर्ष वापरून अशा
निष्पत्त्या अस्तित्वात आहेत असा युक्तिवाद करणे सहसा सोपे असते. असे तीन
निष्कर्ष आपण आधीच प्रस्थापित केले आहेत: $\varphi \in \Gamma$ हे माहीत असेल,
तर $\Gamma \Proves \varphi$ असे नेहमी म्हणता येते, हे
\ref{fol:axd:ptn:prop:reflexivity} सांगतो. $\Gamma \Proves \varphi$ असेल, तर
$\Gamma \cup \{\psi\} \Proves \varphi$ सुद्धा असते, हे
\ref{fol:axd:ptn:prop:monotonicity} सांगतो. आणि $\Gamma \Proves \varphi$ व
$\varphi \Proves \psi$ असेल, तर $\Gamma \Proves \psi$, असे
\ref{fol:axd:ptn:prop:transitivity} वरून मिळते. पुढे आणखी एक सोपा निष्कर्ष,
विध्यात्मक अनुमानाची ``अधि''-आवृत्ती, दिली आहे:

\begin{prop}
  \label{fol:axd:ded:prop:mp} जर $\Gamma \Proves \varphi$ आणि $\Gamma \Proves \varphi \lif
  \psi$, तर $\Gamma \Proves \psi$.
\end{prop}

\begin{proof}
  $\{\varphi, \varphi \lif \psi\} \Proves \psi$ हे आपल्याकडे आहे:
  \begin{derivation}
    1. & $\varphi$ & Hyp.\\
    2. & $\varphi \lif \psi$ & Hyp.\\
    3. & $\psi$ & 1, 2, MP
  \end{derivation}
  \ref{fol:axd:ptn:prop:transitivity} वरून $\Gamma \Proves \psi$.
\end{proof}

या संदर्भात आपण वापरणार असलेला सर्वांत महत्त्वाचा निष्कर्ष म्हणजे निगमन
प्रमेय:

\begin{thm}[निगमन प्रमेय]
\label{fol:axd:ded:thm:deduction-thm} $\Gamma \cup \{\varphi\} \Proves \psi$ तेव्हा आणि
केवळ तेव्हाच, जेव्हा $\Gamma \Proves \varphi \lif \psi$.
\end{thm}

\begin{proof}
``जर'' ही दिशा तात्काळ मिळते. $\Gamma \Proves \varphi \lif \psi$ असेल, तर
\ref{fol:axd:ptn:prop:monotonicity} वरून $\Gamma \cup \{\varphi\} \Proves \varphi \lif \psi$
सुद्धा असते. तसेच \ref{fol:axd:ptn:prop:reflexivity} वरून
$\Gamma \cup \{\varphi\} \Proves \varphi$. म्हणून \ref{fol:axd:ded:prop:mp} वरून
$\Gamma \cup \{\varphi\} \Proves \psi$.

``केवळ जर'' या दिशेसाठी, $\Gamma \cup \{\varphi\}$ पासून $\psi$ च्या
निष्पत्तीच्या लांबीवर विगमन करू.

विगमनाच्या आधारपायरीसाठी, लांबी~$1$ असलेल्या प्रत्येक निष्पत्ती साठी
दावा सिद्ध करू. $\Gamma \cup \{\varphi\}$ पासून~$\psi$ ची लांबी~$1$ असलेली
निष्पत्ती म्हणजे एकटे $\psi$; ती योग्य असेल, तर $\psi$ हे एकतर
$\in \Gamma \cup \{\varphi\}$ असते किंवा स्वयंसिद्धक असते. $\psi \in \Gamma$
असेल किंवा ते स्वयंसिद्धक असेल, तर $\Gamma \Proves \psi$. तसेच
\ref{fol:axd:prp:ax:lif1} वरून $\Gamma \Proves \psi \lif (\varphi \lif \psi)$, आणि
\ref{fol:axd:ded:prop:mp} वरून $\Gamma \Proves \varphi \lif \psi$. $\psi \in \{ \varphi\}$
असेल, तर $\Gamma \Proves \varphi \lif \psi$, कारण शेवटचे वाक्य~$\varphi
\lif \psi$ हे $\varphi \lif \varphi$ सारखेच आहे आणि ते आपण
\ref{fol:axd:pro:ex:identity} मध्ये निष्पन्न केले आहे.

विगमन पायरीसाठी, $\Gamma \cup \{\varphi\}$ पासून~$\psi$ ची एखादी
निष्पत्ती विध्यात्मक अनुमानाने समर्थित असलेल्या~$\psi$ या पायरीने
संपते असे समजा. (तिला विध्यात्मक अनुमानाने समर्थन मिळत नसेल, तर
$\psi \in \Gamma$, $\psi \ident \varphi$, किंवा $\psi$ हे स्वयंसिद्धक असते, आणि
विगमनाच्या आधारपायरीतील युक्तिवादच लागू होतो.) मग एखाद्या
सूत्र~$\chi$ साठी त्या निष्पत्तीमधील आधीच्या काही पायऱ्या
$\chi \lif \psi$ आणि $\chi$ असतात; म्हणजे $\Gamma \cup \{\varphi\} \Proves \chi \lif \psi$
आणि $\Gamma \cup \{\varphi\} \Proves \chi$. त्यांच्याशी संबंधित
निष्पत्त्या लहान असल्यामुळे त्यांना विगमन गृहीतक लागू होते. म्हणून
आपल्याकडे पुढील दोन्ही आहेत:
\begin{align*}
  & \Gamma \Proves \varphi \lif (\chi \lif \psi); \\
  & \Gamma \Proves \varphi \lif \chi.
\end{align*}
पण \ref{fol:axd:prp:ax:lif2} वरून पुढीलही मिळते:
\[\Gamma \Proves (\varphi \lif (\chi \lif \psi)) \lif
((\varphi\lif \chi)  \lif (\varphi \lif \psi)),
\]
आणि \ref{fol:axd:ded:prop:mp} चे दोन उपयोग केल्यावर आवश्यकतेनुसार
$\Gamma \Proves \varphi \lif \psi$ मिळते.
\end{proof}

निगमन प्रमेय खरे ठरावे यासाठीच \ref{fol:axd:prp:ax:lif1} आणि
\ref{fol:axd:prp:ax:lif2} यांची नेमकी अशी निवड केली होती, हे लक्षात घ्या.

निष्पन्नता विषयीची पुढील काही उपयुक्त तथ्ये सरावासाठी सोडली आहेत.

\begin{prop}
  \label{fol:axd:ded:prop:derivfacts}
\begin{enumerate}
\item $\Proves (\varphi \lif \psi) \lif ((\psi \lif \chi)
  \lif (\varphi \lif \chi))$; \label{fol:axd:ded:derivfacts:a}

\item $\Gamma \cup \{ \lnot \varphi\}
  \Proves \lnot \psi$ असेल, तर $\Gamma \cup \{ \psi\} \Proves
  \varphi$ (प्रतिपरिवर्तन); \label{fol:axd:ded:derivfacts:b}
\item  $\{ \varphi, \lnot\varphi\} \Proves
    \psi$ (Ex Falso Quodlibet, म्हणजे असत्यापासून काहीही; स्फोट); \label{fol:axd:ded:derivfacts:c}
\item  $\{ \lnot\lnot\varphi\} \Proves
  \varphi$ (द्विनकरण विलोपन);\label{fol:axd:ded:derivfacts:d}
\item $\Gamma \Proves \lnot\lnot\varphi$ असेल, तर $\Gamma \Proves
  \varphi$;\label{fol:axd:ded:derivfacts:e}
\end{enumerate}
\end{prop}


\begin{prob}
  \ref{fol:axd:ded:prop:derivfacts} सिद्ध करा.
\end{prob}















\section{संख्यापकांसह निगमन प्रमेय}\label{fol:axd:ddq:sec}

\begin{thm}[निगमन प्रमेय]
\label{fol:axd:ddq:thm:deduction-thm-q} जर $\Gamma \cup \{\varphi\} \Proves \psi$ असेल,
तर $\Gamma \Proves \varphi \lif \psi$.
\end{thm}

\begin{proof}
$\Gamma \cup \{\varphi\}$ पासून $\psi$ च्या निष्पत्तीच्या लांबीवर
आपण पुन्हा विगमन करू.

विगमनाच्या आधारपायरीचा पुरावा~\ref{fol:axd:ded:thm:deduction-thm} च्या
पुराव्यातील आधारपायरीच्या पुराव्यासारखाच आहे.

विगमन पायरीसाठी, $\Gamma \cup \{\varphi\}$ पासून $\psi$ ची निष्पत्ती
एखाद्या निगमन नियमाने समर्थित असलेल्या~$\psi$ या पायरीने संपते असे पुन्हा
समजा. निगमन नियम विध्यात्मक अनुमान असेल, तर आपण
\ref{fol:axd:ded:thm:deduction-thm} च्या पुराव्याप्रमाणेच पुढे जाऊ. निगमन
नियम \QR असेल, तर $\psi \ident \chi \lif \lforall[x][\theta(x)]$ आहे आणि
$\chi \lif \theta(a)$ या रूपाचे सूत्र त्या निष्पत्तीमध्ये आधी
येते, जेथे $a$ हे~$\chi$, $\varphi$, किंवा $\Gamma$ मध्ये येत नाही. म्हणून
आपल्याकडे पुढील आहे:
\begin{align*}
  \Gamma \cup \{\varphi\} & \Proves \chi \lif \theta(a),\\
  \intertext{आणि विगमन गृहीतक लागू होते; म्हणजे पुढील मिळते:}
    \Gamma & \Proves \varphi \lif (\chi \lif \theta(a)).\\
  \intertext{पुढील सूत्रावरून}
  & \Proves (\varphi \lif (\chi \lif \theta(a))) \lif ((\varphi \land \chi) \lif \theta(a))\\
  \intertext{आणि विध्यात्मक अनुमानामुळे पुढील मिळते:}
  \Gamma & \Proves (\varphi \land \chi) \lif \theta(a).\\
  \intertext{आयगेन चराची अट अजूनही लागू असल्यामुळे, या निष्पत्तीमध्ये \QR ने समर्थित एक पायरी जोडून पुढील मिळवता येते:}
    \Gamma & \Proves (\varphi \land \chi) \lif \lforall[x][\theta(x)].\\
    \intertext{आपल्याकडे पुढीलही आहे:}
    & \Proves ((\varphi \land \chi) \lif \lforall[x][\theta(x)]) \lif (\varphi \lif (\chi \lif \lforall[x][\theta(x)])),\\
    \intertext{म्हणून विध्यात्मक अनुमानाने,}
    \Gamma & \Proves \varphi \lif (\chi \lif \lforall[x][\theta(x)]),
\end{align*}

म्हणजे $\Gamma \Proves \varphi \lif \psi$.


$\psi$ ला \QR या नियमाने समर्थन मिळते, पण त्याचे रूप
$\lexists[x][\theta(x)] \lif \chi$ आहे, हे प्रकरण आपण सरावासाठी सोडतो.
\end{proof}

\begin{prob}
  \ref{fol:axd:ddq:thm:deduction-thm-q} चा पुरावा पूर्ण करा.
\end{prob}












\section{निष्पन्नता आणि सुसंगतता}\label{fol:axd:prv:sec}

आता आपण निष्पन्नता संबंधाचे अनेक गुणधर्म सिद्ध करू. ते स्वतंत्रपणेही
महत्त्वाचे आहेत, पण संपूर्णता प्रमेयाच्या सिद्धतेत प्रत्येकाची भूमिका असेल.

\begin{prop}\label{fol:axd:prv:prop:provability-contr}
  जर $\Gamma \Proves \varphi$ आणि $\Gamma \cup \{\varphi\}$ विसंगत असेल, तर
  $\Gamma$ विसंगत आहे.
\end{prop}

\begin{proof}
  जर $\Gamma \cup \{\varphi\}$ विसंगत असेल, तर $\Gamma \cup \{\varphi\}
  \Proves \lfalse$. \ref{fol:axd:ptn:prop:reflexivity} वरून प्रत्येक
  $\psi \in \Gamma$ साठी $\Gamma \Proves \psi$. गृहीतकानुसार
  $\Gamma \Proves \varphi$ सुद्धा असल्यामुळे प्रत्येक $\psi \in \Gamma \cup
  \{\varphi\}$ साठी $\Gamma \Proves \psi$. \ref{fol:axd:ptn:prop:transitivity}
  वरून $\Gamma \Proves \lfalse$, म्हणजे $\Gamma$ विसंगत आहे.
\end{proof}

\begin{prop}
\label{fol:axd:prv:prop:prov-incons}
$\Gamma \Proves \varphi$ तेव्हा आणि केवळ तेव्हाच, जेव्हा $\Gamma \cup \{\lnot \varphi\}$
विसंगत असते.
\end{prop}

\begin{proof}
प्रथम $\Gamma \Proves \varphi$ असे समजा. मग \ref{fol:axd:ptn:prop:monotonicity}
वरून $\Gamma \cup \{\lnot \varphi\} \Proves \varphi$. \ref{fol:axd:ptn:prop:reflexivity}
वरून $\Gamma \cup \{\lnot \varphi\} \Proves \lnot \varphi$. तसेच
\ref{fol:axd:prp:ax:lnot2} वरून $\Proves \lnot \varphi \lif (\varphi \lif \lfalse)$.
म्हणून \ref{fol:axd:ded:prop:mp} चे दोन उपयोग करून $\Gamma \cup \{\lnot
\varphi\} \Proves \lfalse$ मिळते.

आता $\Gamma \cup \{\lnot \varphi\}$ विसंगत आहे असे गृहीत धरा, म्हणजे
$\Gamma \cup \{\lnot \varphi\} \Proves \lfalse$. निगमन प्रमेयावरून $\Gamma
\Proves \lnot \varphi \lif \lfalse$. \ref{fol:axd:prp:ax:lfalse2} वरून $\Gamma
\Proves (\lnot \varphi \lif \lfalse) \lif \lnot\lnot \varphi$, म्हणून
\ref{fol:axd:ded:prop:mp} वरून $\Gamma \Proves \lnot\lnot \varphi$. आपल्याकडे
$\Gamma \Proves \lnot\lnot \varphi \lif \varphi$ हेही असल्यामुळे
(\ref{fol:axd:prp:ax:dne}), \ref{fol:axd:ded:prop:mp} च्या आणखी एका उपयोगाने
$\Gamma \Proves \varphi$ मिळते.
\end{proof}

\begin{prob}
$\Gamma \Proves \lnot \varphi$ तेव्हा आणि केवळ तेव्हाच, जेव्हा $\Gamma \cup \{\varphi\}$
विसंगत असते, हे सिद्ध करा.
\end{prob}

\begin{prop}\label{fol:axd:prv:prop:explicit-inc}
  जर $\Gamma \Proves \varphi$ आणि $\lnot \varphi \in \Gamma$ असेल, तर $\Gamma$
  विसंगत आहे.
\end{prop}

\begin{proof}
  \ref{fol:axd:prp:ax:lnot2} वरून $\Gamma \Proves \lnot \varphi \lif (\varphi \lif \lfalse)$.
  \ref{fol:axd:ded:prop:mp} चे दोन उपयोग करून $\Gamma \Proves \lfalse$ मिळते.
\end{proof}


\begin{prop}\label{fol:axd:prv:prop:provability-exhaustive}
  $\Gamma \cup \{\varphi\}$ आणि $\Gamma \cup \{\lnot \varphi\}$ हे दोन्ही संच
  विसंगत असतील, तर $\Gamma$ विसंगत आहे.
\end{prop}

\begin{proof}
सराव.
\end{proof}


\begin{prob}
  \ref{fol:axd:prv:prop:provability-exhaustive} सिद्ध करा
\end{prob}
















\section{निष्पन्नता आणि विधानीय संयोजक}\label{fol:axd:ppr:sec}

\begin{explain}
  स्वयंसिद्धकीय निगमनाचा निष्पन्नता संबंध~$\Proves$ हा विधानीय
  संयोजकांशी संबंधित काही मूलभूत तथ्ये सिद्ध करण्यासाठी पुरेसा सामर्थ्यवान
  आहे; उदाहरणार्थ, $\varphi \land \psi \Proves \varphi$ आणि $\varphi, \varphi \lif \psi
  \Proves \psi$ (विध्यात्मक अनुमान). ही तथ्ये संपूर्णता प्रमेयाच्या
  सिद्धतेसाठी आवश्यक आहेत.
\end{explain}

\begin{prop}\label{fol:axd:ppr:prop:provability-land}
  \begin{enumerate}
  \item \label{fol:axd:ppr:prop:provability-land-left} $\varphi \land \psi \Proves
    \varphi$ आणि $\varphi \land \psi \Proves \psi$ ही दोन्ही तथ्ये सत्य आहेत.
  \item \label{fol:axd:ppr:prop:provability-land-right} $\varphi, \psi \Proves \varphi \land \psi$.
  \end{enumerate}
\end{prop}

\begin{proof}
  \begin{enumerate}
    \item \ref{fol:axd:prp:ax:land1} आणि \ref{fol:axd:prp:ax:land2} यांपासून,
      अनुक्रमे विध्यात्मक अनुमान लावून.

  \item \ref{fol:axd:prp:ax:land3} पासून विध्यात्मक अनुमान दोनदा लावून.
  \end{enumerate}
\end{proof}


\begin{prop}\label{fol:axd:ppr:prop:provability-lor}
  \begin{enumerate}
  \item $\varphi \lor \psi, \lnot \varphi, \lnot \psi$ विसंगत आहे.
  \item $\varphi \Proves \varphi \lor \psi$ आणि $\psi \Proves \varphi \lor \psi$ ही दोन्ही
    तथ्ये सत्य आहेत.
  \end{enumerate}
\end{prop}

\begin{proof}
  \begin{enumerate}
  \item \ref{fol:axd:prp:ax:lnot2} पासून $\Proves \lnot \varphi \lif (\varphi
    \lif \lfalse)$ आणि $\Proves \lnot \psi \lif (\psi \lif \lfalse)$
    मिळतात. म्हणून निगमन प्रमेयाने $\{\lnot \varphi\} \Proves \varphi \lif
    \lfalse$ आणि $\{\lnot \psi\} \Proves \psi \lif \lfalse$ मिळतात.
    \ref{fol:axd:prp:ax:lor3} पासून $\{\lnot \varphi, \lnot \psi\} \Proves (\varphi
    \lor \psi) \lif \lfalse$ मिळते. निगमन प्रमेयाने $\{\varphi \lor \psi,
    \lnot \varphi, \lnot \psi\} \Proves \lfalse$.

  \item \ref{fol:axd:prp:ax:lor1} आणि \ref{fol:axd:prp:ax:lor2} यांपासून विध्यात्मक
    अनुमानाने.
  \end{enumerate}
\end{proof}

\begin{prop}\label{fol:axd:ppr:prop:provability-lif}
  \begin{enumerate}
  \item \label{fol:axd:ppr:prop:provability-lif-left} $\varphi, \varphi \lif \psi \Proves \psi$.
  \item \label{fol:axd:ppr:prop:provability-lif-right}
    $\lnot \varphi \Proves \varphi \lif \psi$ आणि $\psi \Proves \varphi \lif \psi$ ही दोन्ही
    तथ्ये सत्य आहेत.
  \end{enumerate}
\end{prop}

\begin{proof}
  \begin{enumerate}
  \item आपण पुढील निष्पन्न करू शकतो:
    \begin{derivation}
      1. & $\varphi$ & \Hyp\\
      2. & $\varphi \lif \psi$ & \Hyp\\
      3. & $\psi$ & 1, 2, \MP
    \end{derivation}

  \item अनुक्रमे \ref{fol:axd:prp:ax:lnot2}, \ref{fol:axd:prp:ax:lif1} आणि निगमन
    प्रमेय वापरून.
  \end{enumerate}
\end{proof}















\section{निष्पन्नता आणि संख्यापक}\label{fol:axd:qpr:sec}

\begin{explain}
  संपूर्णता प्रमेयासाठी या विभागात स्थापित केलेली~$\Proves$ संबंधी तथ्ये
  स्वयंसिद्धकीय निगमनातून निष्पन्न होतात, हेही आवश्यक आहे.
\end{explain}

\begin{thm}
\label{fol:axd:qpr:thm:strong-generalization} जर $c$ हे $\Gamma$ किंवा $\varphi(x)$ मध्ये
न येणारे स्थिरांक असेल आणि $\Gamma \Proves \varphi(c)$, तर $\Gamma
\Proves \lforall[x][\varphi(x)]$.
\end{thm}

\begin{proof}
निगमन प्रमेयाने $\Gamma \Proves \ltrue \lif \varphi(c)$. $c$ हे
$\Gamma$ किंवा~$\top$ मध्ये येत नसल्यामुळे $\Gamma \Proves
\ltrue \lif \lforall[x][\varphi(x)]$ मिळते. त्यानंतर सत्य स्थिरांकाचे
स्वयंसिद्धक आणि विध्यात्मक अनुमान वापरून $\Gamma \Proves
\lforall[x][\varphi(x)]$ मिळते.

\end{proof}

\begin{prop}
\label{fol:axd:qpr:prop:provability-quantifiers}
\begin{enumerate}
\item $\varphi(t) \Proves \lexists[x][\varphi(x)]$.

\item $\lforall[x][\varphi(x)] \Proves \varphi(t)$.
\end{enumerate}
\end{prop}

\begin{proof}
\begin{enumerate}
\item \ref{fol:axd:qua:ax:q2} आणि निगमन प्रमेय वापरून.
\item \ref{fol:axd:qua:ax:q1} आणि निगमन प्रमेय वापरून.
\end{enumerate}
\end{proof}












\section{निर्दोषता}\label{fol:axd:sou:sec}

\begin{explain}
स्वयंसिद्धकीय निगमनासारखी निष्पत्ती-पद्धत
\emph{निर्दोष} असते, जर ती प्रत्यक्षात निष्पन्न न होणाऱ्या गोष्टी
निष्पन्न करू शकत नसेल. त्यामुळे निर्दोषता हा निष्पत्ती-पद्धतींसाठी
हमी दिलेल्या सुरक्षिततेचा एक प्रकार आहे. कोणता सिद्धता-उपपत्तीय गुणधर्म
विचारात आहे यावर अवलंबून, उदाहरणार्थ, आपल्याला पुढील गोष्टी हव्या असतात:
\begin{enumerate}
\item प्रत्येक निष्पन्न करता येण्याजोगे~$\varphi$ वैध असते;
\item $\varphi$ हे इतर काही सूत्रे~$\Gamma$ यांच्यापासून निष्पन्न करता येण्याजोगे असेल,
  तर ते त्यांचा तार्किक परिणामही असते;
\item सूत्रांचा संच~$\Gamma$ विसंगत असेल, तर तो अपूर्ततायोग्य असतो.
\end{enumerate}
हे निष्पत्ती-पद्धतीचे महत्त्वाचे गुणधर्म आहेत. यांपैकी एखादा गुणधर्म
खरा नसेल, तर निष्पत्ती-पद्धतीत दोष आहे---ती खूप जास्त गोष्टी
निष्पन्न करेल. म्हणून निष्पत्ती-पद्धतीची निर्दोषता सिद्ध करणे अत्यंत
महत्त्वाचे आहे.
\end{explain}

\begin{prop}
  जर $\varphi$ हे स्वयंसिद्धक असेल, तर
  प्रत्येक संरचना~$\Struct{M}$ आणि चल-विन्यास~$s$ साठी $\Sat{M}{\varphi}[s]$.
\end{prop}

\begin{proof}
सर्व स्वयंसिद्धके वैध आहेत हे तपासावे लागते. उदाहरणार्थ,
  \ref{fol:axd:qua:ax:q1} चे प्रकरण पाहू: $t$ हे $\varphi$ मध्ये $x$ साठी
  आदेशनासाठी मुक्त आहे असे समजा आणि
  $\Sat{M}{\lforall[x][\varphi]}[s]$ असे गृहीत धरा. मग पूर्ततेच्या व्याख्येनुसार
  प्रत्येक $\varAssign{s'}{s}{x}$ साठी $\Sat{M}{\varphi}[s']$, आणि विशेषतः
  $s'(x) = \Value{t}{M}[s]$ असताना हे खरे आहे.
  \ref{fol:syn:ext:prop:ext-formulas} नुसार
  $\Sat{M}{\Subst{\varphi}{t}{x}}[s]$. यावरून
  $\Sat{M}{(\lforall[x][\varphi] \lif \Subst{\varphi}{t}{x})}[s]$ हे सिद्ध होते.
\end{proof}

\begin{thm}[निर्दोषता]
\label{fol:axd:sou:thm:soundness}
जर $\Gamma \Proves \varphi$, तर $\Gamma \Entails \varphi$.
\end{thm}

\begin{proof}
$\Gamma$ पासून $\varphi$ च्या निष्पत्तीच्या लांबीवर विगमन करू. अनुमानांनी
समर्थित कोणत्याही पायऱ्या नसतील, तर निष्पत्तीमधील सर्व सूत्रे
ही स्वयंसिद्धकांची दृष्टांते असतात किंवा~$\Gamma$ मध्ये असतात. आधीच्या
विधानानुसार सर्व स्वयंसिद्धके वैध आहेत;
म्हणून $\varphi$ हे स्वयंसिद्धक असेल, तर $\Gamma \Entails \varphi$. जर
$\varphi \in \Gamma$, तर स्पष्टपणे $\Gamma \Entails \varphi$.

जर~$\varphi$ च्या निष्पत्तीमधील शेवटची पायरी विध्यात्मक अनुमानाने समर्थित
असेल, तर त्या निष्पत्तीमध्ये $\psi$ आणि $\psi \lif \varphi$ ही सूत्रे
असतात; आणि त्या सूत्रे वर संपणाऱ्या निष्पत्तीच्या भागांना विगमन
गृहीतक लागू होते (कारण त्यांच्यात अनुमानाने समर्थित किमान एक पायरी कमी
असते). म्हणून विगमन गृहीतकानुसार $\Gamma \Entails \psi$ आणि
$\Gamma \Entails \psi \lif \varphi$. मग
\ref{fol:syn:sem:thm:sem-deduction}
नुसार $\Gamma \Entails \varphi$.

आता शेवटची पायरी \QR ने समर्थित आहे असे समजा. मग त्या पायरीचे
रूप $\chi \lif \lforall[x][\psi(x)]$ असते आणि तिच्याआधीची पायरी
$\chi \lif \psi(c)$ असते; येथे $c$ हे $\Gamma$, $\chi$ किंवा
$\lforall[x][B(x)]$ यांमध्ये येत नाही. विगमन गृहीतकानुसार
$\Gamma \Entails \chi \lif \psi(c)$.
\ref{fol:syn:sem:thm:sem-deduction} नुसार $\Gamma \cup \{\chi\} \Entails
\psi(c)$.


$\Sat{M}{\Gamma \cup
  \{\chi\}}$ अशी एखादी रचना~$\Struct{M}$ विचारात घ्या. आपल्याला
$\Sat{M}{\lforall[x][\psi(x)]}$ हे दाखवायचे आहे.
$\lforall[x][\psi(x)]$ हे वाक्य असल्यामुळे याचा अर्थ प्रत्येक
चल-विन्यास~$s$ साठी $\Sat{M}{\psi(x)}[s]$ दाखवायचे आहे
(\ref{fol:syn:ass:prop:sat-quant}). $\Gamma \cup \{\chi\}$ मध्ये केवळ
वाक्ये असल्यामुळे \ref{fol:syn:sat:defn:satisfaction} नुसार प्रत्येक
$\theta \in \Gamma$ साठी $\Sat{M}{\theta}[s]$. $\Struct{M'}$ ही
$\Struct{M}$ सारखीच रचना असू द्या, फक्त $\Assign{c}{M'} = s(x)$.
$c$ हे~$\Gamma$ किंवा~$\chi$ मध्ये येत नसल्यामुळे
\ref{fol:syn:ext:cor:extensionality-sent} नुसार
$\Sat{M'}{\Gamma \cup \{\chi\}}$. $\Gamma \cup \{\chi\}
\Entails \psi(c)$ असल्यामुळे $\Sat{M'}{\psi(c)}$. $\psi(c)$ हे वाक्य
असल्यामुळे \ref{fol:syn:ass:prop:sentence-sat-true} नुसार
$\Sat{M'}{\psi(c)}[s]$. \ref{fol:syn:ext:prop:ext-formulas} नुसार
$\Sat{M'}{\psi(x)}[s]$ तेव्हा आणि केवळ तेव्हाच, जेव्हा
$\Sat{M'}{\psi(c)}[s]$ (लक्षात घ्या की $\psi(c)$ म्हणजे
$\Subst{\psi(x)}{c}{x}$). म्हणून $\Sat{M'}{\psi(x)}[s]$. $c$ हे~$\psi(x)$
मध्ये येत नसल्यामुळे \ref{fol:syn:ext:prop:extensionality} नुसार
$\Sat{M}{\psi(x)}[s]$. पण $s$ हा स्वेच्छ चल-विन्यास होता, म्हणून
$\Sat{M}{\lforall[x][\psi(x)]}$. त्यामुळे $\Gamma \cup \{\chi\} \Entails
\lforall[x][\psi(x)]$. \ref{fol:syn:sem:thm:sem-deduction} नुसार
$\Gamma \Entails \chi \lif \lforall[x][\psi(x)]$.


$\varphi$ हे \QR ने समर्थित पण $\lexists[x][\psi(x)] \lif \chi$ या रूपाचे असलेले
प्रकरण सरावासाठी सोडले आहे.
\end{proof}


\begin{prob}
  \ref{fol:axd:sou:thm:soundness} ची सिद्धता पूर्ण करा.
\end{prob}


\begin{cor}
\label{fol:axd:sou:cor:weak-soundness}
जर $\Proves \varphi$, तर $\varphi$ हे वैध आहे.
\end{cor}

\begin{cor}
\label{fol:axd:sou:cor:consistency-soundness}
जर $\Gamma$ पूर्ततायोग्य असेल, तर तो सुसंगत आहे.
\end{cor}

\begin{proof}
प्रतिलोम विधान सिद्ध करू. $\Gamma$ सुसंगत नाही असे समजा. मग
$\Gamma \Proves \lfalse$, म्हणजे~$\Gamma$ पासून $\lfalse$ ची
निष्पत्ती असते. \ref{fol:axd:sou:thm:soundness} नुसार $\Gamma$ पूर्ण करणाऱ्या
प्रत्येक
संरचना~$\Struct{M}$
ने~$\lfalse$ पूर्ण केले पाहिजे. प्रत्येक
संरचना~$\Struct{M}$ साठी
  $\Sat/{M}{\lfalse}$ असल्यामुळे कोणतीही
$\Struct{M}$ $\Gamma$ पूर्ण करू शकत नाही;
म्हणजे $\Gamma$ अपूर्ततायोग्य आहे.
\end{proof}














\section{निष्पत्ती सह एकरूपता}\label{fol:axd:ide:sec}

निष्पत्त्यांमध्ये $\eq$ समाविष्ट करण्यासाठी आपण फक्त नवीन
स्वयंसिद्धक-आकार जोडतो. निष्पत्ती आणि $\Proves$ यांच्या व्याख्या
तशाच राहतात; आता आपण नवीन स्वयंसिद्धकेही मान्य करतो.

\begin{defn}[एकरूपता साठीची स्वयंसिद्धके]
\begin{align}
\label{fol:axd:ide:ax:id1} & \eq[t][t], \\
\label{fol:axd:ide:ax:id2} & \eq[t_1][t_2] \lif (\psi(t_1) \lif
  \psi(t_2)),
\end{align}
येथे $t$, $t_1$, $t_2$ ही कोणतीही बंद पदे आहेत.
\end{defn}

\begin{prop}\label{fol:axd:ide:prop:sound}
  \ref{fol:axd:ide:ax:id1} आणि \ref{fol:axd:ide:ax:id2} ही स्वयंसिद्धके वैध आहेत.
\end{prop}

\begin{proof}
  सराव.
\end{proof}

\begin{prob}
\ref{fol:axd:ide:prop:sound} सिद्ध करा.
\end{prob}

\begin{prop}\label{fol:axd:ide:prop:iden1}
 कोणतेही पद $t$ आणि कोणताही संच~$\Gamma$ यांसाठी
 $\Gamma \Proves \eq[t][t]$.
\end{prop}

\begin{prop}\label{fol:axd:ide:prop:iden2}
  जर $\Gamma \Proves \varphi(t_1)$ आणि $\Gamma \Proves
  \eq[t_1][t_2]$, तर $\Gamma \Proves \varphi(t_2)$.
\end{prop}

\begin{proof}
पुढील सूत्र
\[(\eq[t_1][t_2] \lif (\varphi(t_1) \lif \varphi(t_2)))
\]
हा~\ref{fol:axd:ide:ax:id2} चा दृष्टांत आहे. निष्कर्ष~\MP चे दोन उपयोग केल्यावर
मिळतो.
\end{proof}



\chapter{संपूर्णता प्रमेय}\label{fol:com::chap}









\section{परिचय}\label{fol:com:int:sec}

संपूर्णता प्रमेय हा तर्कशास्त्राविषयीच्या सर्वांत मूलभूत निष्कर्षांपैकी एक
आहे. त्याच्या दोन मांडण्या आहेत; त्यांची सममूल्यता आपण सिद्ध करू. त्याची
पहिली मांडणी चिन्हार्थी परिणाम आणि आपल्या निष्पत्ती-पद्धतीतील संबंधाबद्दल
एक मूलभूत विधान करते: जर वाक्य~$\varphi$ हे काही वाक्ये
$\Gamma$ यांच्यापासून निष्पन्न होत असेल, तर $\Gamma \Proves \varphi$ हे
प्रस्थापित करणारी निष्पत्ती सुद्धा असते. त्यामुळे प्रत्यक्षात निष्पन्न
न होणाऱ्या गोष्टी सिद्ध न करता निष्पत्ती-पद्धत जितकी सामर्थ्यवान असू
शकते तितकी ती असते.

दुसऱ्या मांडणीत ते प्रतिमानाच्या अस्तित्वाविषयीचा निष्कर्ष म्हणून सांगता येते:
वाक्यांचा प्रत्येक सुसंगत संच पूर्ततायोग्य असतो. सुसंगतता ही
सिद्धता-उपपत्तीय संकल्पना आहे: आपली निष्पत्ती-पद्धत विशिष्ट
निष्पत्त्या निर्माण करू शकत नाही, असे ती सांगते. पण~$\Gamma$ पासून
विशिष्ट प्रकारच्या निष्पत्त्यांनाहीत एवढ्यावरूनच
संरचना~$\Struct{M}$
अस्तित्वात असून $\Sat{M}{\Gamma}$ आहे, याची हमी
आहे, असे कोण म्हणेल? संपूर्णता प्रमेय प्रथम सिद्ध होण्यापूर्वी---
खरे तर आज आपल्याकडे असलेल्या निष्पत्ती-पद्धती निर्माण होण्यापूर्वी---
महान जर्मन गणितज्ञ डाव्हीट हिल्बर्ट यांचे मत होते की गणिती सिद्धांतांची
सुसंगतता त्यांच्या विषयातील वस्तूंच्या अस्तित्वाची हमी देते. गोटलोप फ्रेग
यांना लिहिलेल्या पत्रात त्यांनी ते पुढीलप्रमाणे मांडले:
\begin{quote}
  स्वेच्छपणे दिलेली स्वयंसिद्धके आणि त्यांचे सर्व परिणाम परस्परविरोधी नसतील,
  तर ती सत्य असतात आणि स्वयंसिद्धकांनी परिभाषित केलेल्या वस्तू अस्तित्वात
  असतात. माझ्यासाठी सत्य आणि अस्तित्वाचा निकष हाच आहे.
\end{quote}
फ्रेग यांनी याला तीव्र विरोध केला. संपूर्णता प्रमेयाची दुसरी मांडणी दाखवते
की निदान पुढील अर्थाने हिल्बर्ट बरोबर होते: स्वयंसिद्धके सुसंगत असतील, तर
त्या सर्वांना सत्य करणारी \emph{एखादी}
संरचना अस्तित्वात असते.

संपूर्णता प्रमेय---किंवा अधिक नेमके सांगायचे तर त्याची सिद्धता---महत्त्वाची
असण्याची हीच काही एकमेव कारणे नाहीत. त्याचे अनेक महत्त्वाचे परिणाम आहेत;
त्यांपैकी काहींची स्वतंत्र चर्चा करू. उदाहरणार्थ, $\Gamma \Proves \varphi$ हे
दाखवणारी कोणतीही निष्पत्ती सांत असल्यामुळे ती~$\Gamma$ मधील सांतच
वाक्ये वापरू शकते. म्हणून संपूर्णता प्रमेयावरून असे निष्पन्न होते की
$\varphi$ हा~$\Gamma$ चा परिणाम असेल, तर तो~$\Gamma$ च्या एखाद्या सांत उपसंचाचाही
परिणाम असतो. याला \emph{संहतता} म्हणतात. सममूल्य रीतीने, $\Gamma$ चा प्रत्येक
सांत उपसंच सुसंगत असेल, तर $\Gamma$ स्वतः सुसंगत असला पाहिजे.

संहतता प्रमेय हे संपूर्णता प्रमेयापासून निष्पत्त्या मधून जाणाऱ्या वळणाने
मिळत असले, तरी संपूर्णता प्रमेयाची \emph{सिद्धता} वापरून ते थेट प्रस्थापित
करणेही शक्य आहे. कारण त्या सिद्धतेत विशिष्ट गुणधर्म---सुसंगतता---असलेला
वाक्यांचा संच घेऊन त्यातून विशिष्ट गुणधर्म असलेली संरचना
रचली जाते (या प्रकरणात ती त्या संचाची पूर्ति करते). सुसंगत संचांऐवजी
``सांततः पूर्ततायोग्य'' वाक्यांच्या संचांपासून सुरुवात केल्यास जवळजवळ
हीच रचना वापरून संहतता थेट प्रस्थापित करता येते.
या रचनेतून इतर परिणामही मिळतात; उदाहरणार्थ, वाक्यांच्या
  कोणत्याही पूर्ततायोग्य संचाला सांत किंवा गणनीय अनंत असे प्रतिमान असते.
  (या निष्कर्षाला ल्योव्हेनहाइम--स्कोलेम प्रमेय म्हणतात.) सर्वसाधारणपणे,
  वाक्यांच्या संचांपासून संरचना रचण्याची पद्धत तर्कशास्त्रात
  वारंवार आणि कधी कधी तत्त्वज्ञानातही वापरली जाते.












\section{सिद्धतेचा आराखडा}\label{fol:com:out:sec}

संपूर्णता प्रमेयाची सिद्धता काहीशी गुंतागुंतीची आहे आणि ती प्रथम वाचताना
वाट चुकणे सोपे आहे. म्हणून सिद्धतेचा आराखडा पाहू. पहिली पायरी म्हणजे
दृष्टिकोनातील बदल; त्यामुळे सिद्धतेकडे जाणारा मार्ग दिसू लागतो. संपूर्णतेचा
अर्थ “जेव्हा जेव्हा $\Gamma \Entails \varphi$, तेव्हा $\Gamma \Proves \varphi$”
असा घेतल्यास कल्पनाही सुचणे कठीण जाऊ शकते: कारण $\Gamma \Proves \varphi$ हे
दाखवण्यासाठी आपल्याला निष्पत्ती शोधावी लागते, आणि $\Gamma \Entails
\varphi$ हे गृहीतक त्यासाठी कोणत्याही प्रकारे मदत करते असे दिसत नाही. काही
सिद्धता-पद्धतींसाठी निष्पत्ती थेट रचता येते; पण आपण थोडा वेगळा मार्ग
घेऊ. आवश्यक दृष्टिकोन-बदल असा आहे: संपूर्णता पुढीलप्रमाणेही मांडता येते:
“$\Gamma$ सुसंगत असेल, तर तो पूर्ततायोग्य असतो.” कदाचित~$\Gamma$ मधील
माहिती आणि तो सुसंगत असल्याचे गृहीतक वापरून आपण~$\Gamma$ मधील प्रत्येक
वाक्य पूर्ण करणारी
संरचना रचू शकू. शेवटी, आपल्याला
कोणत्या प्रकारची संरचना हवी आहे हे
माहीत आहे:~$\Gamma$ जिचे वर्णन करतो अशीच!{}

$\Gamma$ मध्ये फक्त आण्विक
  वाक्ये असतील, तर त्याच्यासाठी
प्रतिमान रचणे सोपे आहे. सर्व आण्विक वाक्ये ही
  $\Atom{P}{a_1,\dots,a_n}$ या रूपाची आहेत असे समजा; येथे $a_i$ ही
  स्थिरांक आहेत. आपल्याला फक्त

  प्रांत~$\Domain{M}$ आणि~$P$ साठी असे निर्देशन ठरवायचे आहे की
  $\Sat{M}{\Atom{P}{a_1,\ldots,a_n}}$. पण ते फार कठीण नाही:
  $\Domain{M} = \Nat$, $\Assign{\Obj c_i}{M} = i$ असे ठेवा आणि प्रत्येक
  $\Atom{P}{a_1,\ldots,a_n} \in \Gamma$ साठी $\tuple{k_1,
    \dots, k_n}$ हे क्रमित बहुपद $\Assign{P}{M}$ मध्ये घाला; येथे $k_i$
  हा स्थिर चिन्ह~$a_i$ चा निर्देशांक आहे (म्हणजे $a_i \ident \Obj
  c_{k_i}$).

आता $\Gamma$ मध्ये~$\lnot \psi$ हे एखादे सूत्र असून $\psi$ आण्विक आहे
असे समजा. $\Struct{M}$ ची रचना $\lnot \psi$ सत्य
करण्याच्या शक्यतेत अडथळा आणेल, अशी काळजी वाटू शकते. पण येथेच~$\Gamma$ ची
सुसंगतता उपयोगी पडते: $\lnot \psi \in \Gamma$ असेल, तर $\psi \notin \Gamma$;
अन्यथा $\Gamma$ विसंगत असेल. आणि $\psi \notin \Gamma$ असेल, तर आपल्या
$\Struct{M}$ च्या रचनेनुसार
$\Sat/{M}{\psi}$; म्हणून
$\Sat{M}{\lnot \psi}$. आतापर्यंत सर्व ठीक आहे.

$\Gamma$ मध्ये गुंतागुंतीची, अनाण्विक सूत्रे असतील तर काय? उदाहरणार्थ,
त्यात $\varphi \land \psi$ आहे असे समजा. ते सत्य करण्यासाठी $\varphi$ आणि $\psi$ ही
दोन्ही~$\Gamma$ मध्ये असल्याप्रमाणे आपण पुढे गेले पाहिजे. आणि $\varphi \lor
\psi \in \Gamma$ असेल, तर त्यांपैकी किमान एक सत्य करावे लागेल; म्हणजे
त्यांपैकी एक~$\Gamma$ मध्ये असल्याप्रमाणे पुढे जावे लागेल.

यावरून पुढील कल्पना सुचते: आपण~$\Gamma$ मध्ये अतिरिक्त सूत्रे
अशा रीतीने भर घालतो की (a)~मिळणारा संच सुसंगत राहील, (b)~प्रत्येक
संभाव्य आण्विक वाक्य~$\varphi$ साठी मिळणाऱ्या संचात $\varphi$ किंवा
$\lnot \varphi$ यांपैकी एक असेल, आणि (c)~संचात $\varphi \land \psi$ असेल तेव्हा
$\varphi$ व $\psi$ ही दोन्ही असतील, तर $\varphi \lor \psi$ असेल तेव्हा $\varphi$ किंवा
$\psi$ यांपैकी किमान एकही असेल, इत्यादी. हे आपण करत राहतो (कदाचित
अनंतकाळपर्यंत). अशा रीतीने भर घातलेल्या सर्व सूत्रांच्या संचाला
$\Gamma^*$ म्हणू. मग वरील रचनेतून आपल्याला अशी
संरचना~$\Struct{M}$
मिळेल की ती~$\Gamma^*$ मधील सर्व वाक्ये पूर्ण करते, हे विगमनाने सिद्ध
करता येईल. $\Gamma \subseteq \Gamma^*$ असल्यामुळे ती~$\Gamma$ मधीलही
सर्व वाक्ये पूर्ण करेल. प्रत्यक्षात (a) आणि~(b) यांची हमी देणे पुरेसे
ठरते. ज्या वाक्यसंचासाठी (b) खरे असते त्याला \emph{संपूर्ण} म्हणतात.
त्यामुळे सुसंगत संच~$\Gamma$ चा विस्तार करून सुसंगत आणि संपूर्ण संच
$\Gamma^*$ मिळवणे हे आपले उद्दिष्ट असेल.


या योजनेत एक अडचण आहे: $\lexists[x][\varphi(x)] \in \Gamma$ असेल, तर या
प्रक्रियेत एखादे स्थिरांक~$c$ निवडून $\varphi(c)$ ची भर घालता यावी अशी
आपली अपेक्षा आहे. पण ते नेहमी करता येईल हे कसे कळणार? कदाचित आपल्या
भाषेत मोजकीच स्थिरांक असतील आणि त्यांपैकी प्रत्येकासाठी
$\lnot \varphi(c) \in \Gamma$ असेल. आपण $\varphi(c)$ चीही भर घालू शकत नाही,
कारण त्यामुळे संच विसंगत होईल; आणि $\Struct{M}$ ने $\varphi(c)$ सत्य करावे की
$\lnot \varphi(c)$ हे आपल्याला कळणार नाही. शिवाय, $\Gamma$ मध्ये अशा भाषेतील
वाक्येच असू शकतात जिच्यात एकही स्थिर चिन्ह नाही (उदा., संचसिद्धान्ताची
भाषा).

या समस्येचा उपाय म्हणजे सुरुवातीलाच अनंत स्थिरे आणि त्यांना संख्यापकांशी
योग्य प्रकारे जोडणारी वाक्ये यांची भर घालणे. (अर्थात, त्यामुळे विसंगती
निर्माण होऊ शकत नाही, हे आपल्याला तपासावे लागेल.)

आण्विक वाक्यांमध्ये फक्त स्थिरांक असतील, तर आपली मूळ रचना चांगली
चालते. पण भाषेत फलने सुद्धा असू शकतात. अशा वेळी सर्व काही योग्य
रीतीने चालण्यासाठी या फलने ना नेमून द्यायची~$\Nat$ वरील योग्य
फलने शोधणे अवघड ठरू शकते. म्हणून आणखी एक युक्ती वापरू: $\Obj c_i$ चा
अर्थ लावण्यासाठी $i$ वापरण्याऐवजी स्थिरांकाचा संचच प्रांत म्हणून
घेऊ. मग $\Struct M$ प्रत्येक स्थिरांक ला तेच स्थिर नेमून देऊ शकते:
$\Assign{\Obj c_i}{M} = \Obj c_i$. पण पूर्ण मार्गच का घेऊ नये:
$\Domain{M}$ हा भाषेतील सर्व \emph{पदांचा} संच असू द्या!{} असे केल्यास
फलने ना फलनांचे स्पष्ट निर्देशन मिळते (या फलनांची प्रचले पदे असतात
आणि त्यांची मूल्येही पदे असतात): फलन~$\Obj f^n_i$ ला आपण असे
फलन नेमून देतो की $n$ पदे $t_1$, \dots,~$t_n$ आदान म्हणून दिल्यावर ते
$\Obj f^n_i(t_1, \dots, t_n)$ हे पद मूल्य म्हणून देते.

कोड्याचा शेवटचा भाग म्हणजे~$\eq$ चे काय करायचे. विधेय~$\eq$ चा
अर्थ स्थिर आहे: $\Sat{M}{\eq[t][t']}$ तेव्हा आणि केवळ तेव्हाच, जेव्हा
$\Value{t}{M} = \Value{t'}{M}$. आता पद~$t$ चे मूल्य हे $t$ स्वतःच
असेल अशी मांडणी केली, तर $t$ आणि $t'$ हे अगदी एकच पद असल्याखेरीज ही
संरचना $\eq[t][t']$ या रूपातील \emph{एकही} वाक्य सत्य करणार नाही.
आणि ही अर्थातच समस्या आहे; कारण फलने असलेल्या भाषेतील जवळजवळ
प्रत्येक रोचक उपपत्तीमध्ये $t$ व $t'$ ही वेगळी पदे असूनही $\eq[t][t']$
या रूपाची वाक्ये प्रमेये असतील (उदा., अंकगणिताच्या उपपत्तींमध्ये
$\eq[(\Obj 0+ \Obj 0)][\Obj 0]$). ही समस्या सोडवण्यासाठी आपण
$\Struct M$ चा प्रांत बदलतो: $\Domain{M}$ मधील वस्तू म्हणून पदे वापरण्याऐवजी
पदांचे संच वापरतो; आणि प्रत्येक संचात~$\Gamma$ मधील वाक्यांनी समान असणे
आवश्यक ठरवलेली सर्व पदे असतात. उदाहरणार्थ, $\Gamma$ ही अंकगणिताची उपपत्ती
असेल, तर अशा संचांपैकी एकात $\Obj 0$, $(\Obj 0 + \Obj 0)$, $(\Obj
0 \times \Obj 0)$, इत्यादी असतील. हाच संच आपण $\Obj 0$ ला नेमून देऊ;
आणि तो त्यातील सर्व पदांचे मूल्यही आहे, उदा., $(\Obj 0 + \Obj 0)$ चेही,
हे निष्पन्न होईल. म्हणून सुधारित संरचनांमध्ये
$\eq[(\Obj 0+ \Obj 0)][\Obj 0]$ हे वाक्य सत्य असेल.

मग आपण पुढीलप्रमाणे काम करू. प्रथम संपूर्ण सुसंगत संचांच्या गुणधर्मांचा
अभ्यास करू. विशेषतः, संपूर्ण सुसंगत संचात $\varphi \land \psi$ तेव्हा आणि
केवळ तेव्हाच असते, जेव्हा त्यात $\varphi$ व~$\psi$ ही दोन्ही असतात; $\varphi \lor \psi$
तेव्हा आणि केवळ तेव्हाच असते, जेव्हा त्यांपैकी किमान एक असते; इत्यादी,
हे सिद्ध करू (\ref{fol:com:ccs:prop:ccs}). त्यानंतर वाक्यांच्या
  “संतृप्त” संचांची व्याख्या करून त्यांचा अभ्यास करू. प्रत्येक संख्यकीकृत
  वाक्याला तिच्या दृष्टांतांशी जोडणाऱ्या सशर्तांचा समावेश असलेला
  संच म्हणजे संतृप्त संच (\ref{fol:com:hen:defn:henkin-exp}). कोणत्याही सुसंगत
  संच~$\Gamma$ चा विस्तार करून संतृप्त संच~$\Gamma'$ नेहमी मिळवता येतो,
  हे आपण दाखवू (\ref{fol:com:hen:lem:henkin}). एखादा संच सुसंगत, संतृप्त आणि
  संपूर्ण असेल, तर त्यात $\lexists[x][\varphi(x)]$ तेव्हा
    आणि केवळ तेव्हाच असते, जेव्हा एखाद्या बंद पद~$t$ साठी त्यात $\varphi(t)$
    असते आणि
  $\lforall[x][\varphi(x)]$ तेव्हा आणि केवळ तेव्हाच असते,
    जेव्हा प्रत्येक बंद पद~$t$ साठी त्यात $\varphi(t)$ असते
  (\ref{fol:com:hen:prop:saturated-instances}), हाही गुणधर्म त्याला असतो.
त्यानंतर संतृप्त सुसंगत संच
$\Gamma'$ घेऊन त्याचा विस्तार करून
संतृप्त, सुसंगत आणि संपूर्ण संच~$\Gamma^*$
मिळवता येतो, हे दाखवू (\ref{fol:com:lin:lem:lindenbaum}). याच~$\Gamma^*$ चा
वापर आपण आपले पद-प्रतिमान~$\Struct M(\Gamma^*)$ परिभाषित करण्यासाठी करू.
पद-प्रतिमानाचा प्रांत हा बंद पदांचा संच असतो आणि त्यातील
  विधेयांचा अर्थ आण्विक वाक्ये यांच्या~$\Gamma^*$ मधील सदस्यत्वाने ठरते
(\ref{fol:com:mod:defn:termmodel}). संतृप्त, संपूर्ण सुसंगत
संचांचे गुणधर्म वापरून
$\Sat{M(\Gamma^*)}{\varphi}$ तेव्हा आणि
केवळ तेव्हाच, जेव्हा $\varphi \in \Gamma^*$, हे दाखवू
(\ref{fol:com:mod:lem:truth}); म्हणून विशेषतः
$\Sat{M(\Gamma^*)}{\Gamma}$.
शेवटी, $\Gamma$ मध्ये~$\eq$ सुद्धा असेल तर पद-प्रतिमानाची
  व्याख्या कशी करायची याचा विचार करू
  (\ref{fol:com:ide:defn:term-model-factor}) आणि ते~$\Gamma^*$ ची पूर्ति करते,
  हे दाखवू (\ref{fol:com:ide:lem:truth}).
















\section{वाक्याचे संपूर्ण सुसंगत संच}\label{fol:com:ccs:sec}

\begin{defn}[संपूर्ण संच]
\label{fol:com:ccs:def:complete-set} वाक्यांचा संच~$\Gamma$
\emph{संपूर्ण} असतो तेव्हा आणि केवळ तेव्हाच, जेव्हा कोणत्याही
वाक्य~$\varphi$ साठी $\varphi \in \Gamma$ किंवा $\lnot \varphi \in \Gamma$.
\end{defn}

\begin{explain}
संपूर्ण वाक्यसंच कोणताही प्रश्न अनुत्तरित ठेवत नाहीत. कोणत्याही
वाक्य~$\varphi$ साठी $\varphi$ सत्य आहे की असत्य हे~$\Gamma$ “सांगतो.”
संपूर्ण संचांचे महत्त्व संपूर्णता प्रमेयाच्या सिद्धतेपलीकडेही आहे.
उदाहरणार्थ, जी उपपत्ती संपूर्ण असून स्वयंसिद्धकीय रीतीने मांडता
येते, तिच्याविषयी नेहमी निर्णय घेता येतो.
\end{explain}

\begin{explain}
संपूर्णतेच्या सिद्धतेत संपूर्ण सुसंगत संच महत्त्वाचे असतात, कारण
वाक्यांचा प्रत्येक सुसंगत संच~$\Gamma$ हा एखाद्या संपूर्ण
सुसंगत संच~$\Gamma^*$ मध्ये समाविष्ट असतो, याची हमी आपण देऊ शकतो.
संपूर्ण सुसंगत संचात प्रत्येक वाक्य~$\varphi$ साठी $\varphi$ किंवा
त्याचे नकरण $\lnot \varphi$ यांपैकी एक असते, पण दोन्ही नसतात. हे विशेषतः
आण्विक वाक्ये यांच्यासाठी खरे असते.
म्हणून स्थिरांकाची योग्य भर घालून विस्तारित केलेल्या
भाषेतील संपूर्ण सुसंगत संचापासून आपण अशी
संरचना रचू शकतो की तिच्यातील
विधेयांचा अर्थ
हा~$\Gamma^*$ मध्ये कोणती आण्विक
वाक्ये आहेत त्यानुसार परिभाषित
केला जातो. त्यानंतर ही संरचना
$\Gamma^*$ मधील सर्व वाक्ये (आणि म्हणून~$\Gamma$ मधीलही सर्व)
सत्य करते, हे दाखवता येते. या शेवटच्या तथ्याच्या सिद्धतेसाठी
$\lnot \varphi \in \Gamma^*$ तेव्हा आणि केवळ तेव्हाच, जेव्हा $\varphi \notin
\Gamma^*$; $(\varphi \lor \psi) \in \Gamma^*$ तेव्हा आणि केवळ तेव्हाच, जेव्हा
$\varphi \in \Gamma^*$ किंवा $\psi \in \Gamma^*$; इत्यादी तथ्ये आवश्यक असतात.
\end{explain}

पुढे आपण $\Proves$ ची परावर्तिता, एकस्वनिकता आणि संक्रमकता हे गुणधर्म
अनेकदा अध्याहृतपणे वापरू (पाहा

  \ref{fol:seq:ptn:sec}, \ref{fol:ntd:ptn:sec}, \ref{fol:axd:ptn:sec}, \ref{fol:tab:ptn:sec}).

\begin{prop}
\label{fol:com:ccs:prop:ccs}
$\Gamma$ हा संपूर्ण आणि सुसंगत आहे असे समजा. मग:
\begin{enumerate}
\item \label{fol:com:ccs:prop:ccs-prov-in} जर $\Gamma \Proves \varphi$, तर $\varphi \in
  \Gamma$.

\item \label{fol:com:ccs:prop:ccs-and} $\varphi \land \psi \in \Gamma$
  तेव्हा आणि केवळ तेव्हाच, जेव्हा $\varphi \in \Gamma$ आणि $\psi \in \Gamma$
  ही दोन्ही खरे आहेत.

\item \label{fol:com:ccs:prop:ccs-or} $\varphi \lor \psi \in \Gamma$ तेव्हा
  आणि केवळ तेव्हाच, जेव्हा $\varphi \in \Gamma$ किंवा $\psi \in \Gamma$.

\item \label{fol:com:ccs:prop:ccs-if} $\varphi \lif \psi \in \Gamma$ तेव्हा
  आणि केवळ तेव्हाच, जेव्हा $\varphi \notin \Gamma$ किंवा $\psi \in \Gamma$.
\end{enumerate}
\end{prop}

\begin{proof}
पुढील सर्व प्रकरणांत $\Gamma$ हा संपूर्ण आणि सुसंगत आहे असे समजू.
\begin{enumerate}
\item जर $\Gamma \Proves \varphi$, तर $\varphi \in \Gamma$.

$\Gamma \Proves \varphi$ असे समजा. उलटपक्षी $\varphi \notin \Gamma$ असे गृहीत
धरा. $\Gamma$ हा संपूर्ण असल्यामुळे $\lnot \varphi \in \Gamma$. पुढील
संदर्भांनुसार

  \ref{fol:seq:prv:prop:explicit-inc}, \ref{fol:ntd:prv:prop:explicit-inc}, \ref{fol:axd:prv:prop:explicit-inc}, \ref{fol:tab:prv:prop:explicit-inc},
$\Gamma$ विसंगत आहे. हे $\Gamma$ सुसंगत असल्याच्या गृहीतकाशी विसंगत आहे.
म्हणून $\varphi \notin \Gamma$ असणे शक्य नाही; त्यामुळे $\varphi \in \Gamma$.

\item

$\varphi \land \psi \in \Gamma$ तेव्हा आणि केवळ तेव्हाच, जेव्हा $\varphi \in
\Gamma$ आणि $\psi \in \Gamma$ ही दोन्ही खरे आहेत:

पुढे जाणाऱ्या दिशेसाठी $\varphi \land \psi \in \Gamma$ असे समजा. मग
पुढील संदर्भांतील

  \ref{fol:seq:ppr:prop:provability-land}, \ref{fol:ntd:ppr:prop:provability-land}, \ref{fol:axd:ppr:prop:provability-land}, \ref{fol:tab:ppr:prop:provability-land}, बाब~(1) नुसार
$\Gamma \Proves \varphi$ आणि $\Gamma \Proves \psi$. \ref{fol:com:ccs:prop:ccs-prov-in}
नुसार $\varphi \in \Gamma$ आणि $\psi \in \Gamma$, अपेक्षेप्रमाणे.

उलट दिशेसाठी $\varphi \in \Gamma$ आणि $\psi \in \Gamma$ असे समजा. पुढील
संदर्भांतील

  \ref{fol:seq:ppr:prop:provability-land}, \ref{fol:ntd:ppr:prop:provability-land}, \ref{fol:axd:ppr:prop:provability-land}, \ref{fol:tab:ppr:prop:provability-land}, बाब~(2) नुसार
$\Gamma \Proves \varphi \land \psi$. \ref{fol:com:ccs:prop:ccs-prov-in} नुसार $\varphi \land
\psi \in \Gamma$.

\item

प्रथम $\varphi \lor \psi \in \Gamma$ असेल, तर $\varphi \in \Gamma$ किंवा $\psi \in
\Gamma$, हे दाखवू. $\varphi \lor \psi \in \Gamma$, पण $\varphi \notin \Gamma$ आणि
$\psi \notin \Gamma$ असे समजा. $\Gamma$ हा संपूर्ण असल्यामुळे
$\lnot \varphi \in \Gamma$ आणि $\lnot \psi \in \Gamma$. पुढील संदर्भांतील

  \ref{fol:seq:ppr:prop:provability-lor}, \ref{fol:ntd:ppr:prop:provability-lor}, \ref{fol:axd:ppr:prop:provability-lor}, \ref{fol:tab:ppr:prop:provability-lor},
बाब~(1) नुसार $\Gamma$ विसंगत आहे; हा विरोधाभास आहे. म्हणून $\varphi \in
\Gamma$ किंवा $\psi \in \Gamma$.

उलट दिशेसाठी $\varphi \in \Gamma$ किंवा $\psi \in \Gamma$ असे समजा. पुढील
संदर्भांतील

  \ref{fol:seq:ppr:prop:provability-lor}, \ref{fol:ntd:ppr:prop:provability-lor}, \ref{fol:axd:ppr:prop:provability-lor}, \ref{fol:tab:ppr:prop:provability-lor}, बाब~(2) नुसार
$\Gamma \Proves \varphi \lor \psi$. \ref{fol:com:ccs:prop:ccs-prov-in} नुसार $\varphi \lor
\psi \in \Gamma$, अपेक्षेप्रमाणे.

\item

पुढे जाणाऱ्या दिशेसाठी $\varphi \lif \psi \in \Gamma$ असे समजा आणि उलटपक्षी
$\varphi \in \Gamma$ व $\psi \notin \Gamma$ असे गृहीत धरा. या गृहीतकांनुसार
$\varphi \lif \psi \in \Gamma$ आणि $\varphi \in \Gamma$. पुढील संदर्भांतील

  \ref{fol:seq:ppr:prop:provability-lif}, \ref{fol:ntd:ppr:prop:provability-lif}, \ref{fol:axd:ppr:prop:provability-lif}, \ref{fol:tab:ppr:prop:provability-lif}, बाब~(1) नुसार
$\Gamma \Proves \psi$. पण मग \ref{fol:com:ccs:prop:ccs-prov-in} नुसार $\psi \in
\Gamma$; हे $\psi \notin \Gamma$ या गृहीतकाशी विसंगत आहे.

उलट दिशेसाठी प्रथम $\varphi \notin \Gamma$ हे प्रकरण पाहू. $\Gamma$ हा
संपूर्ण असल्यामुळे $\lnot \varphi \in \Gamma$. पुढील संदर्भांतील

  \ref{fol:seq:ppr:prop:provability-lif}, \ref{fol:ntd:ppr:prop:provability-lif}, \ref{fol:axd:ppr:prop:provability-lif}, \ref{fol:tab:ppr:prop:provability-lif}, बाब~(2) नुसार
$\Gamma \Proves \varphi \lif \psi$. पुन्हा \ref{fol:com:ccs:prop:ccs-prov-in} नुसार
$\varphi \lif \psi \in \Gamma$, अपेक्षेप्रमाणे.

आता $\psi \in \Gamma$ हे प्रकरण पाहू. पुढील संदर्भांतील

  \ref{fol:seq:ppr:prop:provability-lif}, \ref{fol:ntd:ppr:prop:provability-lif}, \ref{fol:tab:ppr:prop:provability-lif}, \ref{fol:axd:ppr:prop:provability-lif},
पुन्हा बाब~(2) नुसार $\Gamma \Proves \varphi \lif \psi$.
\ref{fol:com:ccs:prop:ccs-prov-in} नुसार $\varphi \lif \psi \in \Gamma$.
\end{enumerate}
\end{proof}


\begin{prob}
\ref{fol:com:ccs:prop:ccs} ची सिद्धता पूर्ण करा.
\end{prob}















\section{हेंकिन विस्तार}\label{fol:com:hen:sec}

\begin{explain}
संपूर्णता प्रमेय सिद्ध करण्यातील एक आव्हान असे आहे की संपूर्ण सुसंगत
संच~$\Gamma$ पासून आपण रचलेल्या प्रतिमानाने~$\Gamma$ मधील सर्व संख्यकीकृत
सूत्रे सत्य केली पाहिजेत. याची हमी देण्यासाठी लिऑन हेंकिन यांची एक
युक्ती वापरतो. थोडक्यात, भाषेत अनंत स्थिरांकाची भर घालून तिचा विस्तार
करणे आणि एक मुक्त चल असलेल्या प्रत्येक सूत्र~$\varphi(x)$ साठी
पुढील रूपाचे सूत्र जोडणे, ही ती युक्ती:
$\lexists[x][\varphi(x)] \lif \varphi(c)$,
येथे $c$ हे नवीन स्थिरांक पैकी एक आहे. $\Gamma$ पूर्ण करणारी
संरचना आपण रचू तेव्हा, यामुळे प्रत्येक
सत्य अस्तित्ववाची वाक्याचा साक्षी
नवीन स्थिरांमध्ये असेल, याची हमी मिळेल.
\end{explain}

\begin{prop}
\label{fol:com:hen:prop:lang-exp}
$\Gamma$ हा $\Lang L$ मध्ये सुसंगत असेल आणि $\Lang L$ मध्ये नवीन
स्थिरांकाचा गणनीय अनंत संच $\Obj d_0$, $\Obj d_1$, \dots
मिळवून $\Lang L'$ तयार केली असेल, तर $\Gamma$ हा~$\Lang L'$ मध्ये
सुसंगत असतो.
\end{prop}

\begin{defn}[संतृप्त संच]
भाषा $\Lang {L}$ मधील सूत्रांचा संच $\Gamma$
\emph{संतृप्त} असतो तेव्हा आणि केवळ तेव्हाच, जेव्हा एक मुक्त
चल~$x$ असलेल्या प्रत्येक सूत्र~$\varphi(x) \in \Frm[L]$
साठी असे स्थिरांक~$c \in \Lang{L}$ असते की
$\lexists[x][\varphi(x)] \lif \varphi(c) \in \Gamma$.
\end{defn}

पुढील प्रमेयाच्या सिद्धतेत खालील व्याख्या वापरली जाईल.

\begin{defn}
\label{fol:com:hen:defn:henkin-exp}
$\Lang L'$ ही \ref{fol:com:hen:prop:lang-exp} मधीलप्रमाणे असू द्या. ज्या
$\Lang L'$ मधील सूत्रे~$\varphi_i(x_i)$ मध्ये एक चर ($x_i$) मुक्तपणे
येतो, त्या सर्वांचे $\varphi_0(x_0)$, $\varphi_1(x_1)$, \dots असे प्रगणन निश्चित
करा. आपण~$n$ वर विगमनाने वाक्ये~$\theta_n$ परिभाषित करतो.

$\Lang{L}$ मध्ये आपण भर घातलेल्या $\Obj d_i$ पैकी~$\varphi_0(x_0)$ मध्ये
न येणारे पहिले स्थिरांक म्हणजे $c_0$ असू द्या. $\theta_0$,
\dots,~$\theta_{n-1}$ आधीच परिभाषित केली आहेत असे गृहीत धरून, नवीन
स्थिरांक~$\Obj d_i$ पैकी~$\theta_0$, \dots,~$\theta_{n-1}$ यांमध्ये किंवा
$\varphi_n(x_n)$ मध्ये न येणारे पहिले स्थिर म्हणजे $c_n$ असू द्या.

आता $\theta_{n}$ हे पुढील सूत्र असू द्या:
$\lexists[x_{n}][\varphi_{n}(x_{n})] \lif
  \varphi_{n}(c_{n})$.
\end{defn}

\begin{lem}
\label{fol:com:hen:lem:henkin}
प्रत्येक सुसंगत संच~$\Gamma$ चा विस्तार करून संतृप्त सुसंगत
संच~$\Gamma'$ मिळवता येतो.
\end{lem}

\begin{proof}
भाषा~$\Lang{L}$ मधील वाक्यांचा सुसंगत संच~$\Gamma$ दिला असता,
नवीन स्थिरांकाचा गणनीय अनंत संच जोडून भाषेचा विस्तार
$\Lang{L'}$ पर्यंत करा. \ref{fol:com:hen:prop:lang-exp} नुसार $\Gamma$ या अधिक
समृद्ध भाषेतही सुसंगत असतो. तसेच $\theta_i$ ही
\ref{fol:com:hen:defn:henkin-exp} मधीलप्रमाणे असू द्या. पुढील व्याख्या करा:
\begin{align*}
\Gamma_0 & = \Gamma \\
\Gamma_{n+1} & = \Gamma_n \cup \{\theta_n \}
\end{align*}
म्हणजे $\Gamma_{n+1} = \Gamma \cup \{ \theta_0, \dots, \theta_n \}$; आणि
$\Gamma' = \bigcup_{n} \Gamma_n$ असू द्या. $\Gamma'$ स्पष्टपणे संतृप्त आहे.

$\Gamma'$ विसंगत असेल, तर कोणत्यातरी $n$ साठी $\Gamma_n$ विसंगत असेल
(सराव: का ते स्पष्ट करा). म्हणून $\Gamma'$ सुसंगत आहे हे दाखवण्यासाठी
प्रत्येक संच~$\Gamma_n$ सुसंगत आहे हे~$n$ वर विगमनाने दाखवणे पुरेसे आहे.

विगमनाचा आधार म्हणजे $\Gamma_0 = \Gamma$ सुसंगत आहे हे विधानच; ते
प्रमेयाचे गृहीतक आहे. विगमन-पायरीसाठी $\Gamma_{n}$ सुसंगत आहे, पण
$\Gamma_{n+1} = \Gamma_n \cup \{\theta_n\}$ विसंगत आहे असे समजा. लक्षात
घ्या की $\theta_n$ हे
$\lexists[x_{n}][\varphi_{n}(x_n)] \lif \varphi_{n}(c_{n})$
आहे; येथे $\varphi_n(x_n)$ हे फक्त~$x_n$ चर मुक्त असलेले $\Lang{L'}$ चे
सूत्र आहे. $c_n$ निवडण्याच्या पद्धतीनुसार
(\ref{fol:com:hen:defn:henkin-exp} पाहा), $c_n$ हे~$\varphi_n(x_n)$ किंवा~$\Gamma_n$
मध्ये येत नाही.

$\Gamma_n \cup \{\theta_n\}$ विसंगत असेल, तर $\Gamma_n
\Proves \lnot \theta_n$; म्हणून पुढील दोन्ही विधाने खरी आहेत:

\[\Gamma_n \Proves \lexists[x_n][\varphi_n(x_n)]
\qquad
\Gamma_n \Proves \lnot \varphi_n(c_n)
\]
$c_n$ हे $\Gamma_n$ किंवा~$\varphi_n(x_n)$ मध्ये येत नसल्यामुळे
\ref{fol:axd:qpr:thm:strong-generalization}, \ref{fol:seq:qpr:thm:strong-generalization}, \ref{fol:ntd:qpr:thm:strong-generalization}, \ref{fol:tab:qpr:thm:strong-generalization}
लागू होते.
$\Gamma_n \Proves \lnot \varphi_n(c_n)$
यापासून अनुक्रमे
$\Gamma_n \Proves \lforall[x_n][\lnot \varphi_n(x_n)]$
मिळते. अशा रीतीने आपल्याकडे पुढील दोन्ही विधाने आहेत:
$\Gamma_n \Proves \lexists[x_n][\varphi_n(x_n)]$ आणि
$\Gamma_n \Proves \lforall[x_n][\lnot \varphi_n(x_n)]$;
म्हणून $\Gamma_n$ स्वतः विसंगत आहे.
(लक्षात घ्या की
$\lforall[x_n][\lnot \varphi_n(x_n)] \Proves
  \lnot\lexists[x_n][\varphi_n(x_n)]$.)
हा विरोधाभास आहे: $\Gamma_n$ सुसंगत असल्याचे गृहीत होते. म्हणून
$\Gamma_n \cup \{ \theta_n\}$ सुसंगत आहे.

\end{proof}

\begin{explain}
आता आपण दाखवू की संतृप्त असलेल्या \emph{संपूर्ण}, सुसंगत संचांना पुढील
गुणधर्म असतो: त्यात सर्ववाची संख्यकीकृत वाक्य
  तेव्हा आणि केवळ तेव्हाच असते, जेव्हा त्याचे सर्व दृष्टांत त्यात असतात
   आणि त्यात अस्तित्ववाची
  संख्यकीकृत वाक्य तेव्हा आणि केवळ तेव्हाच असते, जेव्हा त्याचा
  किमान एक दृष्टांत त्यात असतो. संपूर्ण, सुसंगत, संतृप्त संचापासून
आपण निर्माण करू ती संरचना त्यातील सर्व संख्यकीकृत वाक्ये सत्य करते,
हे दाखवण्यासाठी हा गुणधर्म वापरू.
\end{explain}

\begin{prop}\label{fol:com:hen:prop:saturated-instances}
$\Gamma$ हा संपूर्ण, सुसंगत आणि संतृप्त आहे असे समजा.
\begin{enumerate}
\item $\lexists[x][\varphi(x)] \in \Gamma$ तेव्हा आणि केवळ तेव्हाच,
  जेव्हा किमान एका बंद पद~$t$ साठी $\varphi(t) \in \Gamma$.
\item $\lforall[x][\varphi(x)] \in \Gamma$ तेव्हा आणि केवळ तेव्हाच,
  जेव्हा सर्व बंद पदे~$t$ यांच्यासाठी $\varphi(t) \in \Gamma$.
\end{enumerate}
\end{prop}

\begin{proof}
  \begin{enumerate}
  \item
    प्रथम $\lexists[x][\varphi(x)]
      \in \Gamma$ असे समजा. $\Gamma$ संतृप्त असल्यामुळे कोणत्यातरी
      स्थिरांक~$c$ साठी
      $(\lexists[x][\varphi(x)] \lif \varphi(c)) \in \Gamma$. पुढील संदर्भांतील
      \ref{fol:axd:ppr:prop:provability-lif}, \ref{fol:seq:ppr:prop:provability-lif}, \ref{fol:ntd:ppr:prop:provability-lif}, \ref{fol:tab:ppr:prop:provability-lif},
      बाब~(1), आणि
      \ref{fol:com:ccs:prop:ccs}\ref{fol:com:ccs:prop:ccs-prov-in} नुसार $\varphi(c)
      \in \Gamma$.

    उलट दिशेसाठी संतृप्तता आवश्यक नाही: $\varphi(t) \in \Gamma$ असे समजा.
    मग पुढील संदर्भांतील
      \ref{fol:axd:qpr:prop:provability-quantifiers}, \ref{fol:seq:qpr:prop:provability-quantifiers}, \ref{fol:ntd:qpr:prop:provability-quantifiers}, \ref{fol:tab:qpr:prop:provability-quantifiers},
      बाब~(1) नुसार $\Gamma \Proves \lexists[x][\varphi(x)]$.
    \ref{fol:com:ccs:prop:ccs}\ref{fol:com:ccs:prop:ccs-prov-in} नुसार
    $\lexists[x][\varphi(x)] \in \Gamma$.
  \item
    सर्व बंद पदे~$t$ यांच्यासाठी
      $\varphi(t) \in \Gamma$ असे समजा. विरोधासाठी
      $\lforall[x][\varphi(x)] \notin \Gamma$ असे गृहीत धरा. $\Gamma$ संपूर्ण
      असल्यामुळे $\lnot\lforall[x][\varphi(x)] \in \Gamma$. संतृप्ततेनुसार
      कोणत्यातरी स्थिरांक~$c$ साठी
      $(\lexists[x][\lnot \varphi(x)] \lif \lnot \varphi(c)) \in
        \Gamma$. $c$ हे बंद पद असल्यामुळे गृहीतकानुसार $\varphi(c) \in
      \Gamma$. पण यामुळे $\Gamma$ विसंगत होईल. (सराव: पुढील संच विसंगत
      असल्याचे दाखवणारी निष्पत्ती द्या:
      \[      \lnot \lforall[x][\varphi(x)],
      \lexists[x][\lnot \varphi(x)] \lif \lnot
        \varphi(c), \varphi(c)
      \]
      .)

      उलट दिशेसाठी संतृप्तता आवश्यक नाही:
      $\lforall[x][\varphi(x)] \in \Gamma$ असे समजा. मग पुढील संदर्भांतील
      \ref{fol:axd:qpr:prop:provability-quantifiers}, \ref{fol:seq:qpr:prop:provability-quantifiers}, \ref{fol:ntd:qpr:prop:provability-quantifiers}, \ref{fol:tab:qpr:prop:provability-quantifiers},
      बाब~(2) नुसार $\Gamma \Proves \varphi(t)$. \ref{fol:com:ccs:prop:ccs}
      नुसार $\varphi(t) \in \Gamma$.
  \end{enumerate}
\end{proof}












\section{लिंडेनबाउमचे पूर्वप्रमेय}\label{fol:com:lin:sec}

\begin{explain}
आता आपण असे पूर्वप्रमेय सिद्ध करू की वाक्यांचा कोणताही सुसंगत
संच केवळ सुसंगतच नव्हे, तर संपूर्ण असलेल्या एखाद्या वाक्यसंचात
समाविष्ट असतो. या सिद्धतेत प्रत्येक वेळी एक वाक्य जोडले जाते
आणि प्रत्येक पायरीवर संच सुसंगत राहील याची हमी दिली जाते. हे अशा रीतीने
करतो की प्रत्येक $\varphi$ साठी कोणत्यातरी टप्प्यावर $\varphi$ किंवा $\lnot
\varphi$ यांपैकी एक जोडले जाईल. मग त्या रचनेतील सर्व टप्प्यांचा संघ $\varphi$
किंवा त्याचे नकरण~$\lnot \varphi$ यांपैकी एक समाविष्ट करतो आणि म्हणून तो
संपूर्ण असतो. प्रत्येक टप्प्यावर विसंगती येणार नाही याची खात्री केल्यामुळे
तो सुसंगतही असतो.
\end{explain}

\begin{lem}[लिंडेनबाउमचे पूर्वप्रमेय]
\label{fol:com:lin:lem:lindenbaum} भाषा~$\Lang{L}$ मधील प्रत्येक सुसंगत
संच~$\Gamma$ चा विस्तार करून संपूर्ण आणि सुसंगत असा
संच~$\Gamma^*$ मिळवता येतो.
\end{lem}

\begin{proof}
$\Gamma$ सुसंगत असू द्या. $\Lang L$ मधील सर्व वाक्यांचे
$\varphi_0$, $\varphi_1$, \dots{} असे प्रगणन असू द्या. $\Gamma_0 = \Gamma$
अशी व्याख्या करा, आणि
\[\Gamma_{n+1} =
\begin{cases}
\Gamma_n \cup \{ \varphi_n \} & \textrm{जर $\Gamma_n \cup \{\varphi_n\}$
  सुसंगत असेल;} \\
\Gamma_n \cup \{ \lnot \varphi_n \} & \textrm{अन्यथा.}
\end{cases}
\]
$\Gamma^* = \bigcup_{n \geq 0} \Gamma_n$ असू द्या.

प्रत्येक $\Gamma_n$ सुसंगत आहे: व्याख्येनुसार $\Gamma_0$ सुसंगत आहे.
$\Gamma_{n+1} = \Gamma_n \cup \{\varphi_n\}$ असेल, तर याचे कारण हा
दुसरा संच सुसंगत असणे हेच आहे. तो सुसंगत नसेल, तर $\Gamma_{n+1} =
\Gamma_n \cup \{\lnot \varphi_n\}$. $\Gamma_n \cup \{\lnot \varphi_n\}$
सुसंगत आहे हे आपल्याला पडताळावे लागेल. तो सुसंगत नाही असे समजा. मग
$\Gamma_n \cup \{\varphi_n\}$ आणि $\Gamma_n \cup \{\lnot \varphi_n\}$ हे
\emph{दोन्ही} विसंगत आहेत. याचा अर्थ

  \ref{fol:axd:prv:prop:provability-exhaustive}, \ref{fol:seq:prv:prop:provability-exhaustive}, \ref{fol:ntd:prv:prop:provability-exhaustive}, \ref{fol:tab:prv:prop:provability-exhaustive}
नुसार $\Gamma_n$ विसंगत असेल; पण हे विगमन-गृहीतकाच्या विरुद्ध आहे.

प्रत्येक~$n$ आणि प्रत्येक $i < n$ साठी $\Gamma_i \subseteq \Gamma_n$.
हे~$n$ वरील सोप्या विगमनाने मिळते. $n=0$ असेल, तर $i < 0$ असा
कोणताही निर्देशांक नसतो; म्हणून विधान आपोआप लागू होते. विगमन-पायरीसाठी ते~$n$
साठी खरे आहे असे समजा. $i < n+1$ असेल, तर $\Gamma_i \subseteq
\Gamma_{n+1}$ हे दाखवू. रचनेनुसार $\Gamma_{n+1} = \Gamma_n \cup
\{\varphi_n\}$ किंवा $= \Gamma_n \cup \{\lnot \varphi_n\}$. म्हणून $\Gamma_n
\subseteq \Gamma_{n+1}$. $i < n+1$ असेल, तर विगमन-गृहीतकानुसार
$\Gamma_i \subseteq \Gamma_n$ ($i < n$ असल्यास), किंवा $\Gamma_n
\subseteq \Gamma_n$ हे क्षुल्लक तथ्य लागू होते ($i = n$ असल्यास).
म्हणून~$\subseteq$ च्या संक्रमणशीलतेने $\Gamma_i \subseteq
\Gamma_{n+1}$ मिळते.

यावरून $\Gamma^*$ सुसंगत आहे हे मिळते. ते कसे ते पाहू: $\Gamma'
\subseteq \Gamma^*$ हा सांत संच असू द्या. प्रत्येक $\psi \in \Gamma'$
हे कोणत्यातरी~$i$ साठी~$\Gamma_i$ मध्येही असते. या निर्देशांकांतील
सर्वांत मोठा निर्देशांक $n$ असू द्या. $i \le n$ असेल, तर $\Gamma_i
\subseteq \Gamma_n$ असल्यामुळे प्रत्येक $\psi \in \Gamma'$ हे
$\in \Gamma_n$ देखील आहे, म्हणजे $\Gamma' \subseteq \Gamma_n$, आणि
$\Gamma_n$ सुसंगत आहे. म्हणून प्रत्येक सांत उपसंच $\Gamma' \subseteq
\Gamma^*$ सुसंगत आहे.

  \ref{fol:axd:ptn:prop:proves-compact}, \ref{fol:seq:ptn:prop:proves-compact}, \ref{fol:ntd:ptn:prop:proves-compact}, \ref{fol:tab:ptn:prop:proves-compact} नुसार $\Gamma^*$
सुसंगत आहे.

$\Frm[L]$ मधील प्रत्येक वाक्य हा~$\Gamma^*$ ची व्याख्या
करण्यासाठी वापरलेल्या यादीत येतो. $\varphi_n \notin \Gamma^*$ असेल, तर
$\Gamma_n \cup \{\varphi_n\}$ विसंगत होते हेच त्याचे कारण आहे. पण मग
$\lnot \varphi_n \in \Gamma^*$; म्हणून $\Gamma^*$ हा संपूर्ण आहे.
\end{proof}












\section{प्रतिमानाची रचना}\label{fol:com:mod:sec}

\begin{explain}
सध्या आपण~$\eq$ चा विचार करत नाही; म्हणजे~$\eq$ नसलेला
  वाक्यांचा सुसंगत संच~$\Gamma$ पूर्ततायोग्य आहे एवढेच आपल्याला
  दाखवायचे आहे. प्रथम~$\Gamma$ चा विस्तार करून सुसंगत, संपूर्ण
  आणि संतृप्त असा संच~$\Gamma^*$ मिळवतो. या बाबतीत
  प्रतिमान~$\Struct{M(\Gamma^*)}$ ची व्याख्या सोपी आहे: आपण~$\Lang{L'}$
  मधील बंद पदांचा संच प्रांत म्हणून घेतो. प्रत्येक स्थिरांक ला तेच
  स्थिर नेमतो आणि, अधिक व्यापकपणे, प्रत्येक बंद पद~$t$ साठी
  $\Value{t}{M(\Gamma^*)} = t$ याची खात्री करतो. विधेये ना
  असे विस्तार नेमतो की एखादे आण्विक वाक्य हे
  $\Struct{M(\Gamma^*)}$ मध्ये तेव्हा आणि केवळ तेव्हाच सत्य असते, जेव्हा
  ते~$\Gamma^*$ मध्ये असते. त्यामुळे~$\Gamma^*$ मधील सर्व आण्विक
  वाक्ये हे~$\Struct{M(\Gamma^*)}$ मध्ये सत्य होतील, हे उघड आहे.
  आपण सुरुवात केलेला~$\Gamma^*$ सुसंगत, संपूर्ण आणि संतृप्त असेल, तर
  उरलेली वाक्येही सत्य असतात.
\end{explain}


\begin{defn}[पद-प्रतिमान]
\label{fol:com:mod:defn:termmodel}
$\Gamma^*$ हा भाषा~$\Lang L$ मधील वाक्यांचा संपूर्ण,
सुसंगत आणि संतृप्त संच असू द्या. $\Gamma^*$ चे
\emph{पद-प्रतिमान}~$\Struct M(\Gamma^*)$ ही पुढीलप्रमाणे परिभाषित केलेली
संरचना आहे:
\begin{enumerate}
\item प्रांत~$\Domain{M(\Gamma^*)}$ हा~$\Lang L$ मधील सर्व बंद
  पदांचा संच आहे.
\item स्थिरांक~$c$ चा निर्वचनार्थ स्वतः~$c$ आहे:
  $\Assign{c}{M(\Gamma^*)} = c$.
\item फलन~$f$ ला असे फलन नेमले जाते की बंद पदे $t_1$, \dots,
  $t_n$ ही वितर्के दिली असता त्याचे मूल्य बंद पद $f(t_1, \dots, t_n)$
  असते:
\[\Assign{f}{M(\Gamma^*)}(t_1, \dots, t_n) = f(t_1,\dots, t_n)
\]
\item $R$ हे $n$-स्थानी विधेय असेल, तर
  \[  \tuple{t_1, \dots,
  t_n} \in \Assign{R}{M(\Gamma^*)} \text{ तेव्हा आणि केवळ तेव्हाच } \Atom{R}{t_1, \dots,
    t_n} \in \Gamma^*.
  \]
\end{enumerate}
\end{defn}



आपल्याकडे खरोखरच $\Value{t}{M(\Gamma^*)} = t$ आहे, हे आता पडताळू.

\begin{lem}
\label{fol:com:mod:lem:val-in-termmodel} \ref{fol:com:mod:defn:termmodel} मधील
$\Struct M(\Gamma^*)$ हे पद-प्रतिमान असू द्या; मग
$\Value{t}{M(\Gamma^*)} = t$.
\end{lem}
\begin{proof}
 ही सिद्धता~$t$ वरील विगमनाने मिळते. $t$ हे स्थिरांक असण्याचा
 आधार-प्रसंग पद-प्रतिमानाच्या व्याख्येवरून थेट मिळतो. विगमन-पायरीसाठी
 $t_1, \ldots, t_n$ ही अशी बंद पदे आहेत की
 $\Value{t_i}{M(\Gamma^*)} = t_i$, आणि $f$ हे $n$-स्थानी
 फलन आहे, असे समजा. मग
\begin{align*}
\Value{f(t_1,\ldots,t_n)}{M(\Gamma^*)} &= \Assign{f}{M(\Gamma^*)}(\Value{t_1}
 {M(\Gamma^*)},
\ldots, \Value{t_n}{M(\Gamma^*)}) \\
&= \Assign{f}{M(\Gamma^*)}(t_1, \dots, t_n) \\
&= f(t_1,\dots, t_n),
\end{align*}
आणि म्हणून विगमनाने हे प्रत्येक बंद पद~$t$ साठी लागू होते.
\end{proof}



\begin{explain}
  एखादी संरचना~$\Struct{M}$ अस्तित्ववाची संख्यकीकृत
    वाक्य~$\lexists[x][\varphi(x)]$ सत्य करू शकते, पण ती सत्य करते
    असा कोणताही दृष्टांत~$\varphi(t)$ नसेल.
  एखादी संरचना~$\Struct{M}$ सर्ववाची संख्यकीकृत
    वाक्य~$\lforall[x][\varphi(x)]$ चे सर्व दृष्टांत~$\varphi(t)$ सत्य करू
    शकते, पण~$\lforall[x][\varphi(x)]$ सत्य करणार नाही. याचे कारण सामान्यतः
  प्रांत~$\Domain{M}$ मधील प्रत्येक घटक हे एखाद्या बंद पदाचे
  मूल्य नसते ($\Struct{M}$ बंद पदांनी आच्छादित नसेल). म्हणूनच
  पूर्ति-संबंधाची व्याख्या चर-मूल्यांकनांच्या साहाय्याने केली जाते. पण आपल्या
  पद-प्रतिमान~$\Struct{M(\Gamma^*)}$ साठी त्याची आवश्यकता नसती---कारण
  ते बंद पदांनी आच्छादित आहे. पुढील निष्पत्तीचा आशय हाच आहे.
\end{explain}

\begin{prop}
\label{fol:com:mod:prop:quant-termmodel}
\ref{fol:com:mod:defn:termmodel} मधील $\Struct M(\Gamma^*)$ हे पद-प्रतिमान असू द्या.
\begin{enumerate}
\item $\Sat{M(\Gamma^*)}{\lexists[x][\varphi(x)]}$ तेव्हा आणि केवळ तेव्हाच,
  जेव्हा किमान एका बंद पद~$t$ साठी $\Sat{M(\Gamma^*)}{\varphi(t)}$.
\item $\Sat{M(\Gamma^*)}{\lforall[x][\varphi(x)]}$ तेव्हा आणि केवळ तेव्हाच,
  जेव्हा सर्व बंद पदे~$t$ यांच्यासाठी $\Sat{M(\Gamma^*)}{\varphi(t)}$.
\end{enumerate}
\end{prop}

\begin{proof}
\begin{enumerate}
\item
  \ref{fol:syn:ass:prop:sat-quant} नुसार
    $\Sat{M(\Gamma^*)}{\lexists[x][\varphi(x)]}$ तेव्हा आणि केवळ तेव्हाच,
    जेव्हा किमान एका चर-मूल्यांकन~$s$ साठी
    $\Sat{M(\Gamma^*)}{\varphi(x)}[s]$. $\Domain{M(\Gamma^*)}$ मध्ये
    $\Lang{L}$ मधील बंद पदे असल्यामुळे, हे तेव्हा आणि केवळ तेव्हाच लागू
    होते, जेव्हा असे किमान एक बंद पद~$t$ असते की $s(x) = t$ आणि
    $\Sat{M(\Gamma^*)}{\varphi(x)}[s]$. \ref{fol:syn:ext:prop:ext-formulas}
    नुसार, $s(x) = t$ असेल तेथे $\Sat{M(\Gamma^*)}{\varphi(x)}[s]$ तेव्हा आणि
    केवळ तेव्हाच, जेव्हा $\Sat{M(\Gamma^*)}{\varphi(t)}[s]$. तसेच
    \ref{fol:syn:ass:prop:sentence-sat-true} नुसार,
    $\varphi(t)$ हे वाक्य असल्यामुळे $\Sat{M(\Gamma^*)}{\varphi(t)}[s]$ तेव्हा आणि
    केवळ तेव्हाच, जेव्हा $\Sat{M(\Gamma^*)}{\varphi(t)}$.
\item
  \ref{fol:syn:ass:prop:sat-quant} नुसार
    $\Sat{M(\Gamma^*)}{\lforall[x][\varphi(x)]}$ तेव्हा आणि केवळ तेव्हाच,
    जेव्हा प्रत्येक चर-मूल्यांकन $s$ साठी
    $\Sat{M(\Gamma^*)}{\varphi(x)}[s]$. लक्षात घ्या की
    $\Domain{M(\Gamma^*)}$ मध्ये~$\Lang{L}$ मधील बंद पदे असतात. त्यामुळे
    प्रत्येक बंद पद~$t$ साठी $s(x) = t$ असे चर-मूल्यांकन मिळते, आणि
    कोणत्याही चर-मूल्यांकनासाठी $s(x)$ हे एखादे बंद पद~$t$ असते.
    \ref{fol:syn:ext:prop:ext-formulas} नुसार, $s(x) = t$ असेल तेथे
    $\Sat{M(\Gamma^*)}{\varphi(x)}[s]$ तेव्हा आणि केवळ तेव्हाच, जेव्हा
    $\Sat{M(\Gamma^*)}{\varphi(t)}[s]$. तसेच
    \ref{fol:syn:ass:prop:sentence-sat-true} नुसार,
    $\varphi(t)$ हे वाक्य असल्यामुळे $\Sat{M(\Gamma^*)}{\varphi(t)}[s]$ तेव्हा आणि
    केवळ तेव्हाच, जेव्हा $\Sat{M(\Gamma^*)}{\varphi(t)}$.
\end{enumerate}
\end{proof}




\begin{lem}[सत्यता पूर्वप्रमेय]
\label{fol:com:mod:lem:truth} $\varphi$ मध्ये~$\eq$ येत नाही असे समजा. मग
$\Sat{M(\Gamma^*)}{\varphi}$ तेव्हा आणि केवळ तेव्हाच, जेव्हा $\varphi \in \Gamma^*$.
\end{lem}

\begin{proof}
दोन्ही दिशा एकाच वेळी, $\varphi$ वरील विगमनाने सिद्ध करू.
\begin{enumerate}
\item \indcase{\varphi}{\lfalse}
  {$\Sat/{M(\Gamma^*)}{\lfalse}$
    हे पूर्तिच्या व्याख्येनुसार मिळते. दुसरीकडे, $\Gamma^*$ सुसंगत
    असल्यामुळे $\lfalse \notin \Gamma^*$.}

\item \indcase{\varphi}{\ltrue}
  {$\Sat{M(\Gamma^*)}{\ltrue}$
    हे पूर्तिच्या व्याख्येनुसार मिळते. दुसरीकडे, $\Gamma^*$ सुसंगत व
    संपूर्ण असल्यामुळे आणि $\Gamma^* \Proves \ltrue$ असल्यामुळे
    $\ltrue \in \Gamma^*$.}

\item \indcase{\varphi}{R(t_1, \dots, t_n)}
  {$\Sat{M(\Gamma^*)}{\Atom{R}{t_1, \dots, t_n}}$ तेव्हा आणि केवळ तेव्हाच,
    जेव्हा $\tuple{t_1, \dots, t_n} \in \Assign{R}{M(\Gamma^*)}$
    (पूर्तिच्या व्याख्येनुसार); आणि हे तेव्हा आणि केवळ तेव्हाच, जेव्हा
    $R(t_1, \dots, t_n) \in \Gamma^*$ ($\Struct M(\Gamma^*)$ च्या
    रचनेनुसार).}

\item
  \indcase{\varphi}{\lnot \psi}
    {$\Sat{M(\Gamma^*)}{\indfrm}$ तेव्हा
    आणि केवळ तेव्हाच, जेव्हा
    $\Sat/{M(\Gamma^*)}{\psi}$
    (पूर्तिच्या व्याख्येनुसार). विगमन-गृहीतकानुसार,
    $\Sat/{M(\Gamma^*)}{\psi}$ तेव्हा
    आणि केवळ तेव्हाच, जेव्हा $\psi \notin \Gamma^*$. $\Gamma^*$ सुसंगत
    आणि संपूर्ण असल्यामुळे $\psi \notin \Gamma^*$ तेव्हा आणि केवळ
    तेव्हाच, जेव्हा $\lnot \psi \in \Gamma^*$.}

\item

  \indcase{\varphi}{\psi \land
    \chi}{$\Sat{M(\Gamma^*)}{\indfrm}$
    तेव्हा आणि केवळ तेव्हाच, जेव्हा
    $\Sat{M(\Gamma^*)}{\psi}$ आणि
    $\Sat{M(\Gamma^*)}{\chi}$ हे दोन्ही
    लागू होतात (पूर्तिच्या व्याख्येनुसार); आणि हे तेव्हा आणि केवळ तेव्हाच,
    जेव्हा $\psi \in \Gamma^*$ आणि $\chi \in \Gamma^*$ हे दोन्ही लागू
    होतात (विगमन-गृहीतकानुसार). तसेच
    \ref{fol:com:ccs:prop:ccs}\ref{fol:com:ccs:prop:ccs-and} नुसार, हे तेव्हा आणि
    केवळ तेव्हाच लागू होते, जेव्हा $(\psi \land \chi) \in \Gamma^*$.}

\item

  \indcase{\varphi}{\psi \lor \chi}
    {$\Sat{M(\Gamma^*)}{\indfrm}$
    तेव्हा आणि केवळ तेव्हाच, जेव्हा
    $\Sat{M(\Gamma^*)}{\psi}$ किंवा
    $\Sat{M(\Gamma^*)}{\chi}$
    (पूर्तिच्या व्याख्येनुसार); आणि हे तेव्हा आणि केवळ तेव्हाच, जेव्हा
    $\psi \in \Gamma^*$ किंवा $\chi \in \Gamma^*$ (विगमन-गृहीतकानुसार).
    \ref{fol:com:ccs:prop:ccs}\ref{fol:com:ccs:prop:ccs-or} नुसार, हे तेव्हा आणि
    केवळ तेव्हाच लागू होते, जेव्हा $(\psi \lor \chi) \in \Gamma^*$.}

\item

  \indcase{\varphi}{\psi \lif \chi}
    {$\Sat{M(\Gamma^*)}{\indfrm}$
    तेव्हा आणि केवळ तेव्हाच, जेव्हा
    $\Sat/{M(\Gamma^*)}{\psi}$ किंवा $\Sat{M (\Gamma^*)}{\chi}$
    (पूर्तिच्या व्याख्येनुसार); आणि हे तेव्हा आणि केवळ तेव्हाच, जेव्हा
    $\psi \notin \Gamma^*$ किंवा $\chi \in \Gamma^*$ (विगमन-गृहीतकानुसार).
    \ref{fol:com:ccs:prop:ccs}\ref{fol:com:ccs:prop:ccs-if} नुसार, हे तेव्हा आणि
    केवळ तेव्हाच लागू होते, जेव्हा $(\psi \lif \chi) \in \Gamma^*$.}


\item

    \indcase{\varphi}{\lforall[x][\psi(x)]}{$\Sat{M(\Gamma^*)}{\indfrm}$ तेव्हा
      आणि केवळ तेव्हाच, जेव्हा सर्व बंद पदे~$t$ यांच्यासाठी
      $\Sat{M(\Gamma^*)}{\psi(t)}$ (\ref{fol:com:mod:prop:quant-termmodel}).
      विगमन-गृहीतकानुसार हे तेव्हा आणि केवळ तेव्हाच लागू होते, जेव्हा सर्व
      बंद पदे~$t$ यांच्यासाठी $\psi(t) \in \Gamma^*$; आणि
      \ref{fol:com:hen:prop:saturated-instances} नुसार हे पुन्हा तेव्हा आणि
      केवळ तेव्हाच लागू होते, जेव्हा $\lforall[x][\psi(x)] \in \Gamma^*$.}


\item

    \indcase{\varphi}{\lexists[x][\psi(x)]}{$\Sat{M(\Gamma^*)}{\indfrm}$ तेव्हा
      आणि केवळ तेव्हाच, जेव्हा किमान एका बंद पद~$t$ साठी
      $\Sat{M(\Gamma^*)}{\psi(t)}$ (\ref{fol:com:mod:prop:quant-termmodel}).
      विगमन-गृहीतकानुसार हे तेव्हा आणि केवळ तेव्हाच लागू होते, जेव्हा किमान
      एका बंद पद~$t$ साठी $\psi(t) \in \Gamma^*$. तसेच
      \ref{fol:com:hen:prop:saturated-instances} नुसार हे पुन्हा तेव्हा आणि
      केवळ तेव्हाच लागू होते, जेव्हा $\lexists[x][\psi(x)] \in \Gamma^*$.}


\end{enumerate}
\end{proof}
















\section{एकरूपता}\label{fol:com:ide:sec}

\begin{explain}
मागील विभागात दिलेली पद-प्रतिमानाची रचना ही~$\eq$ नसलेल्या
संच~$\Gamma$ साठी प्रथम-कोटी तर्कशास्त्राची संपूर्णता प्रस्थापित करण्यास
पुरेशी आहे. पद-प्रतिमान हे~$\eq$ नसलेले प्रत्येक $\varphi \in \Gamma^*$
पूर्त करते (आणि म्हणून सर्व $\varphi \in \Gamma$ पूर्त करते). पण~$\eq$
उपस्थित असेल, तर ही रचना चालत नाही. कारण मग~$\Gamma^*$ मध्ये
वाक्य~$\eq[t][t']$ असू शकते, पण पद-प्रतिमानात कोणत्याही पदाचे
मूल्य ते पद स्वतःच असते. म्हणून $t$ आणि $t'$ ही भिन्न पदे असतील, तर
पद-प्रतिमानातील त्यांची मूल्ये---म्हणजे अनुक्रमे $t$ आणि $t'$---भिन्न
असतात; त्यामुळे $\eq[t][t']$ असत्य असते. मात्र ``भागाकार घेणे'' म्हणून
ओळखली जाणारी रचना वापरून हे दुरुस्त करता येते.
\end{explain}

\begin{defn}
  $\Gamma^*$ हा~$\Lang L$ मधील वाक्यांचा सुसंगत आणि संपूर्ण संच
  असू द्या. $\Lang L$ मधील बंद पदांच्या संचावर आपण~$\approx$ हा संबंध
  पुढीलप्रमाणे परिभाषित करतो:
  \[  t \approx t' \text{\quad तेव्हा आणि केवळ तेव्हाच \quad} \eq[t][t'] \in \Gamma^*
  \]
\end{defn}

\begin{prop}
\label{fol:com:ide:prop:approx-equiv}
$\approx$ या संबंधाला पुढील गुणधर्म आहेत:
\begin{enumerate}
\item $\approx$ परावर्ती आहे.
\item $\approx$ सममित आहे.
\item $\approx$ संक्रमक आहे.
\item $t \approx t'$, $f$ हे फलन, आणि $t_1$, \dots,
  $t_{i-1}$, $t_{i+1}$, \dots, $t_n$ ही बंद पदे असतील, तर
\[\Atom{f}{t_1,\dots, t_{i-1}, t, t_{i+1}, \dots, t_n} \approx
\Atom{f}{t_1,\dots, t_{i-1}, t', t_{i+1}, \dots, t_n}.
\]
\item $t \approx t'$, $R$ हे विधेय, आणि $t_1$, \dots,
  $t_{i-1}$, $t_{i+1}$, \dots, $t_n$ ही बंद पदे असतील, तर
\begin{multline*}
\Atom{R}{t_1,\dots, t_{i-1}, t, t_{i+1}, \dots, t_n} \in \Gamma^* \text{ तेव्हा आणि केवळ तेव्हाच } \\
\Atom{R}{t_1,\dots, t_{i-1}, t', t_{i+1}, \dots, t_n} \in \Gamma^*.
\end{multline*}
\end{enumerate}
\end{prop}

\begin{proof}
$\Gamma^*$ सुसंगत आणि संपूर्ण असल्यामुळे $\eq[t][t'] \in
\Gamma^*$ तेव्हा आणि केवळ तेव्हाच, जेव्हा $\Gamma^* \Proves
\eq[t][t']$. म्हणून पुढील विधाने दाखवणे पुरेसे आहे:
\begin{enumerate}
\item प्रत्येक बंद पद~$t$ साठी $\Gamma^* \Proves \eq[t][t]$.
\item $\Gamma^* \Proves \eq[t][t']$ असेल, तर $\Gamma^* \Proves \eq[t'][t]$.
\item $\Gamma^* \Proves \eq[t][t']$ आणि $\Gamma^* \Proves
  \eq[t'][t'']$ असेल, तर $\Gamma^* \Proves \eq[t][t'']$.
\item $\Gamma^* \Proves \eq[t][t']$ असेल, तर
\[\Gamma^* \Proves
\eq[\Atom{f}{t_1,\dots,t_{i-1},t,t_{i+1},\dots,t_n}][\Atom{f}{t_1,\dots,t_{i-1},t',t_{i+1},\dots,t_n}]
\]
प्रत्येक $n$-स्थानी फलन~$f$ आणि बंद पदे $t_1$, \dots,
$t_{i-1}$, $t_{i+1}$, \dots,~$t_n$ यांच्यासाठी.
\item $\Gamma^* \Proves \eq[t][t']$ आणि
$\Gamma^* \Proves
\Atom{R}{t_1,\dots,t_{i-1},t,t_{i+1},\dots,t_n}$ असेल, तर
\[\Gamma^* \Proves
\Atom{R}{t_1,\dots,t_{i-1},t',t_{i+1},\dots,t_n}
\]
प्रत्येक $n$-स्थानी विधेय~$R$ आणि बंद पदे $t_1$, \dots,
$t_{i-1}$, $t_{i+1}$, \dots,~$t_n$ यांच्यासाठी.
\end{enumerate}

\end{proof}

\begin{prob}
\ref{fol:com:ide:prop:approx-equiv} ची सिद्धता पूर्ण करा.
\end{prob}

\begin{defn}
$\Gamma^*$ हा भाषा~$\Lang L$ मधील सुसंगत आणि संपूर्ण संच आहे,
$t$ हे बंद पद आहे, आणि~$\approx$ हा मागील व्याख्येतील संबंध आहे असे
समजा. मग:
\[\equivrep{t}{\approx} = \Setabs{t'}{t'\in \Trm[L], t \approx t'}
\]
आणि $\equivclass{\Trm[L]}{\approx} = \Setabs{\equivrep{t}{\approx}}{t \in \Trm[L]}$.
\end{defn}

\begin{defn}
\label{fol:com:ide:defn:term-model-factor}
$\Struct M = \Struct M(\Gamma^*)$ हे \ref{fol:com:mod:defn:termmodel} मधील
$\Gamma^*$ चे पद-प्रतिमान असू द्या. मग
$\Struct{\equivclass{M}{\approx}}$ ही पुढील संरचना आहे:
\begin{enumerate}
\item $\Domain{\equivclass{M}{\approx}} = \equivclass{\Trm[L]}{\approx}$.
\item $\Assign{c}{\equivclass{M}{\approx}} = \equivrep{c}{\approx}$
\item $\Assign{f}{\equivclass{M}{\approx}}(\equivrep{t_1}{\approx}, \dots,
  \equivrep{t_n}{\approx}) = \equivrep{\Atom{f}{t_1,\dots, t_n}}{\approx}$
\item $\tuple{\equivrep{t_1}{\approx}, \dots, \equivrep{t_n}{\approx}} \in
  \Assign{R}{\equivclass{M}{\approx}}$ तेव्हा आणि केवळ तेव्हाच, जेव्हा
  $\Sat{M}{\Atom{R}{t_1,\dots, t_n}}$; म्हणजे तेव्हा आणि केवळ तेव्हाच,
  जेव्हा $\Atom{R}{t_1,\dots, t_n} \in \Gamma^*$.
\end{enumerate}
\end{defn}

\begin{explain}
लक्षात घ्या की~$\equivclass{\Trm[L]}{\approx}$ च्या घटकांना
$\equivrep{t}{\approx}$ असे संबोधून, म्हणजेच
\emph{प्रतिनिधी}~$t \in \equivrep{t}{\approx}$ यांच्या साहाय्याने,
आपण त्यांच्यासाठी $\Assign{f}{\equivclass{M}{\approx}}$ आणि
$\Assign{R}{\equivclass{M}{\approx}}$ परिभाषित केली आहेत. या व्याख्या
प्रतिनिधींच्या निवडीवर अवलंबून नाहीत याची आपल्याला खात्री केली पाहिजे.
म्हणजे, तेच तुल्यतावर्ग निश्चित करणारी दुसरी~$t'$ निवडली
($\equivrep{t}{\approx} = \equivrep{t'}{\approx}$), तरी व्याख्यांतून तेच
फल मिळाले पाहिजे. उदाहरणार्थ, $R$ हे एकस्थानी विधेय असेल, तर
व्याख्येच्या शेवटच्या बाबीनुसार $\equivrep{t}{\approx} \in
\Assign{R}{\equivclass{M}{\approx}}$ तेव्हा आणि केवळ तेव्हाच, जेव्हा
$\Sat{M}{\Atom{R}{t}}$. आता $t \approx t'$ असलेल्या दुसऱ्या
पद~$t'$ साठी $\Sat/{M}{\Atom{R}{t'}}$ असेल, तर व्याख्येनुसार
$\equivrep{t'}{\approx} \notin \Assign{R}{\equivclass{M}{\approx}}$ असावे
लागेल. $t \approx t'$ असेल, तर $\equivrep{t}{\approx} =
\equivrep{t'}{\approx}$; पण $\equivrep{t}{\approx} \in
\Assign{R}{\equivclass{M}{\approx}}$ आणि $\equivrep{t}{\approx}
\notin \Assign{R}{\equivclass{M}{\approx}}$ ही दोन्ही विधाने लागू होऊ
शकत नाहीत. मात्र \ref{fol:com:ide:prop:approx-equiv} मुळे असे होऊ शकत नाही, याची
हमी मिळते.

\end{explain}

\begin{prop}
$\Struct{\equivclass{M}{\approx}}$ सुव्याख्यायित आहे; म्हणजे $t_1$,
  \dots, $t_n$, $t_1'$, \dots, $t_n'$ ही बंद पदे असतील आणि
  $t_i \approx t_i'$ असेल, तर
\begin{enumerate}
\item $\equivrep{\Atom{f}{t_1,\dots, t_n}}{\approx} =
    \equivrep{\Atom{f}{t_1',\dots, t_n'}}{\approx}$, म्हणजे
  \[  \Atom{f}{t_1,\dots, t_n} \approx \Atom{f}{t_1',\dots, t_n'}
  \]
  आणि
\item $\Sat{M}{\Atom{R}{t_1,\dots, t_n}}$ तेव्हा आणि केवळ तेव्हाच,
  जेव्हा $\Sat{M}{\Atom{R}{t_1',\dots, t_n'}}$; म्हणजे
  \[    \Atom{R}{t_1,\dots, t_n} \in \Gamma^* \text{ तेव्हा आणि केवळ तेव्हाच }
    \Atom{R}{t_1',\dots, t_n'} \in \Gamma^*.
  \]
\end{enumerate}
\end{prop}

\begin{proof}
\ref{fol:com:ide:prop:approx-equiv} वरून~$n$ वरील विगमनाने मिळते.
\end{proof}

पद-प्रतिमानाच्या बाबतीत होते त्याप्रमाणेच, सत्यता पूर्वप्रमेय सिद्ध
करण्यापूर्वी आपल्याला पुढील पूर्वप्रमेयाची गरज आहे.

\begin{lem}
\label{fol:com:ide:lem:val-in-termmodel-factored} $\Struct M = \Struct M
 (\Gamma^*)$ असू द्या; मग $\Value{t}{\equivclass{M}{\approx}} = \equivrep{t}
 {\approx}$.
\end{lem}
\begin{proof}
ही सिद्धता \ref{fol:com:mod:lem:val-in-termmodel} च्या सिद्धतेसारखीच आहे.
\end{proof}

\begin{prob}
\ref{fol:com:ide:lem:val-in-termmodel-factored} ची सिद्धता पूर्ण करा.
\end{prob}

\begin{lem}
\label{fol:com:ide:lem:truth}
सर्व वाक्ये~$\varphi$ यांच्यासाठी $\Sat{\equivclass{M}{\approx}}{\varphi}$ तेव्हा आणि
केवळ तेव्हाच, जेव्हा $\varphi \in \Gamma^*$.
\end{lem}

\begin{proof}
\ref{fol:com:mod:lem:truth} च्या सिद्धतेप्रमाणेच~$\varphi$ वरील विगमनाने सिद्धता
मिळते. $\varphi \ident \eq[t][t']$ असण्याच्या प्रसंगाकडेच फक्त अधिक लक्ष
देण्याची गरज आहे.
\begin{align*}
\Sat{\equivclass{M}{\approx}}{\eq[t][t']} & \text{ तेव्हा आणि केवळ तेव्हाच } \equivrep{t}{\approx} = \equivrep{t'}{\approx}
\text{ ($\Struct{\equivclass{M}{\approx}}$ च्या व्याख्येनुसार)}\\
& \text{ तेव्हा आणि केवळ तेव्हाच } t \approx t' \text{ ($\equivrep{t}{\approx}$ च्या व्याख्येनुसार)}\\
& \text{ तेव्हा आणि केवळ तेव्हाच } \eq[t][t'] \in \Gamma^* \text{ ($\approx$ च्या व्याख्येनुसार).}
\end{align*}
\end{proof}

\begin{digress}
लक्षात घ्या की $\Struct{M(\Gamma^*)}$ हे नेहमी गणनीय आणि अनंत
असले, तरी $\Struct{\equivclass{M}{\approx}}$ सांत असू शकते; कारण
$\equivrep{t}{\approx}$ असे फक्त सांत अनेक वर्ग असू शकतात. हे अपेक्षितच
आहे, कारण~$\Gamma$ मध्ये अशी वाक्ये असू शकतात की ती सत्य असणारी
कोणतीही संरचना सांत असणे आवश्यक असते. उदाहरणार्थ,
$\lforall[x][\lforall[y][x = y]]$ हे सुसंगत वाक्य आहे, पण ज्याच्या
प्रांत मध्ये नेमका एकच घटक आहे अशाच संरचनांमध्ये
त्याची पूर्ति होते.
\end{digress}












\section{संपूर्णता प्रमेय}\label{fol:com:cth:sec}

\begin{explain}
आपल्या निष्पत्ती एकत्र करू: आपल्याला संपूर्णता प्रमेय मिळते.
\end{explain}

\begin{thm}[संपूर्णता प्रमेय]
\label{fol:com:cth:thm:completeness}
$\Gamma$ हा वाक्यांचा संच असू द्या. $\Gamma$ सुसंगत असेल, तर
तो पूर्ततायोग्य आहे.
\end{thm}

\begin{proof}
$\Gamma$ सुसंगत आहे असे समजा.
  \ref{fol:com:hen:lem:henkin} नुसार असा संतृप्त सुसंगत संच
  $\Gamma' \supseteq \Gamma$ असतो. \ref{fol:com:lin:lem:lindenbaum} नुसार
असा $\Gamma^* \supseteq \Gamma'$ असतो की तो
सुसंगत आणि संपूर्ण असतो.

  $\Gamma' \subseteq \Gamma^*$ असल्यामुळे प्रत्येक
  सूत्र~$\varphi(x)$ साठी~$\Gamma^*$ मध्ये पुढील रूपाचे वाक्य
  असते:
  $\lexists[x][\varphi(x)] \lif \varphi(c)$.
  म्हणून~$\Gamma^*$ संतृप्त आहे. $\Gamma$ मध्ये~$\eq$ येत नसेल, तर
\ref{fol:com:mod:lem:truth} नुसार
$\Sat{M(\Gamma^*)}{\varphi}$ तेव्हा आणि केवळ तेव्हाच, जेव्हा $\varphi \in
\Gamma^*$. यावरून विशेषतः प्रत्येक $\varphi \in \Gamma$ साठी
$\Sat{M(\Gamma^*)}{\varphi}$; म्हणून
$\Gamma$ पूर्ततायोग्य आहे.
  $\Gamma$ मध्ये~$\eq$ येत असेल, तर \ref{fol:com:ide:lem:truth} नुसार
  सर्व वाक्ये~$\varphi$ यांच्यासाठी
  $\Sat{\equivclass{M}{\approx}}{\varphi}$ तेव्हा आणि केवळ तेव्हाच, जेव्हा
  $\varphi \in \Gamma^*$. विशेषतः, प्रत्येक $\varphi \in \Gamma$ साठी
  $\Sat{\equivclass{M}{\approx}}{\varphi}$; म्हणून $\Gamma$ पूर्ततायोग्य आहे.
\end{proof}

\begin{cor}[संपूर्णता प्रमेय, दुसरी आवृत्ती]
\label{fol:com:cth:cor:completeness}
प्रत्येक $\Gamma$ आणि प्रत्येक वाक्ये~$\varphi$ साठी: जर
$\Gamma \Entails \varphi$, तर $\Gamma \Proves \varphi$.
\end{cor}

\begin{proof}
लक्षात घ्या की \ref{fol:com:cth:cor:completeness} आणि
\ref{fol:com:cth:thm:completeness} यांमधील $\Gamma$ वर सर्ववाची संख्यकीकरण केलेले
आहे. आपला गोंधळ होऊ नये म्हणून \ref{fol:com:cth:thm:completeness} हे वेगळा चर
वापरून पुन्हा मांडू: वाक्यांचा कोणताही संच~$\Delta$ दिला असता,
जर $\Delta$ सुसंगत असेल, तर तो पूर्ततायोग्य आहे. प्रतिपरिवर्तनाने,
$\Delta$ पूर्ततायोग्य नसेल, तर $\Delta$ विसंगत आहे. अनुप्रमेय सिद्ध
करण्यासाठी याचा वापर करू.

$\Gamma \Entails \varphi$ असे समजा. मग \ref{fol:syn:sem:prop:entails-unsat}
नुसार $\Gamma \cup \{\lnot \varphi\}$ अपूर्ततायोग्य आहे. $\Gamma \cup
\{\lnot \varphi\}$ हाच आपला $\Delta$ घेतल्यास,
\ref{fol:com:cth:thm:completeness} च्या मागील आवृत्तीनुसार $\Gamma \cup
\{\lnot \varphi\}$ विसंगत आहे. पुढील संदर्भांनुसार,

  \ref{fol:axd:prv:prop:prov-incons}, \ref{fol:seq:prv:prop:prov-incons}, \ref{fol:ntd:prv:prop:prov-incons}, \ref{fol:tab:prv:prop:prov-incons},
$\Gamma \Proves \varphi$.
\end{proof}


\begin{prob}
\ref{fol:com:cth:cor:completeness} वापरून
\ref{fol:com:cth:thm:completeness} सिद्ध करा आणि अशा रीतीने
संपूर्णता प्रमेयाच्या दोन्ही मांडण्या समतुल्य आहेत हे दाखवा.
\end{prob}





\begin{prob}
एखादी निष्पत्ती प्रणाली संपूर्ण असण्यासाठी तिचे नियम इतके समर्थ
असले पाहिजेत की प्रत्येक अपूर्ततायोग्य संच विसंगत असल्याचे ते सिद्ध करू
शकतील. संपूर्णता सिद्ध करण्यासाठी निष्पत्तीचे कोणते नियम आवश्यक
होते? यांपैकी कोणताही नियम सिद्धतेत कोठेच वापरलेला नाही का? या प्रश्नांची
उत्तरे देण्यासाठी, \ref{fol:com:cth:thm:completeness} च्या सिद्धतेपर्यंत
नेणाऱ्या कोणत्या निष्पत्तींमध्ये निष्पत्तीचे कोणते नियम वापरले, हे
दाखवणारी यादी किंवा आकृती तयार करा. या सिद्धतांमधील नियमांच्या कोणत्याही
अध्याहृत वापरांची नोंद अवश्य करा.
\end{prob}















\section{संहतता प्रमेय}\label{fol:com:com:sec}

संपूर्णता प्रमेयाची एक महत्त्वाची निष्पत्ती म्हणजे संहतता प्रमेय.
वाक्यांच्या एखाद्या संचाचा प्रत्येक \emph{सांत} उपसंच पूर्ततायोग्य असेल,
तर संपूर्ण संच पूर्ततायोग्य असतो---तो संच स्वतः अनंत असला तरीही, असे
संहतता प्रमेय सांगते. हे मुळीच उघड नाही. प्रथमदर्शनी तरी अशी शक्यता
नाकारली जाईल असे काही दिसत नाही की वाक्यांचे अनंत संच व्याघाती
असतील, पण व्याघात, जणू काही, अनंत संख्येमुळेच उद्भवत असेल. अशी परिस्थिती
नाकारता येते, असे संहतता प्रमेय सांगते: ज्यांचा प्रत्येक सांत उपसंच
पूर्ततायोग्य आहे असे वाक्यांचे अपूर्ततायोग्य अनंत संच नसतात.
संपूर्णता प्रमेयाप्रमाणे याची तार्किक निष्पन्नतेशी संबंधित आवृत्तीही आहे:
वाक्यांच्या अनंत संचापासून काही निष्पन्न होत असेल, तर ते आधीच
एखाद्या सांत उपसंचापासून निष्पन्न होते.

\begin{defn}
  सूत्रांचा संच~$\Gamma$ \emph{सांततः पूर्ततायोग्य} असतो तेव्हा आणि
  केवळ तेव्हाच, जेव्हा प्रत्येक सांत $\Gamma_0 \subseteq \Gamma$
  पूर्ततायोग्य असतो.
\end{defn}

\begin{thm}[संहतता प्रमेय]
\label{fol:com:com:thm:compactness}
वाक्यांचा कोणताही संच~$\Gamma$ आणि कोणतेही वाक्य~$\varphi$ यांच्यासाठी पुढील
विधाने लागू होतात:
\begin{enumerate}
  \item $\Gamma \Entails \varphi$ तेव्हा आणि केवळ तेव्हाच, जेव्हा असा सांत
    $\Gamma_0 \subseteq \Gamma$ असतो की $\Gamma_0 \Entails \varphi$.
  \item $\Gamma$ पूर्ततायोग्य असतो तेव्हा आणि केवळ तेव्हाच, जेव्हा तो
    सांततः पूर्ततायोग्य असतो.
\end{enumerate}
\end{thm}

\begin{proof}
आपण (2) सिद्ध करू. $\Gamma$ पूर्ततायोग्य असेल, तर अशी
संरचना~$\Struct{M}$
असते की प्रत्येक $\varphi \in \Gamma$ साठी
$\Sat{M}{\varphi}$. अर्थात, ही
$\Struct{M}$ देखील~$\Gamma$ चा प्रत्येक सांत
उपसंच पूर्त करते; म्हणून~$\Gamma$ सांततः पूर्ततायोग्य आहे.

आता $\Gamma$ सांततः पूर्ततायोग्य आहे असे समजा. मग प्रत्येक सांत
उपसंच~$\Gamma_0 \subseteq \Gamma$ पूर्ततायोग्य आहे. निर्दोषतेनुसार
(
  \ref{fol:axd:sou:cor:consistency-soundness}, \ref{fol:seq:sou:cor:consistency-soundness}, \ref{fol:ntd:sou:cor:consistency-soundness}, \ref{fol:tab:sou:cor:consistency-soundness}),
प्रत्येक सांत उपसंच सुसंगत आहे. मग पुढील संदर्भांनुसार $\Gamma$ स्वतः
सुसंगत असला पाहिजे:

  \ref{fol:axd:ptn:prop:proves-compact}, \ref{fol:seq:ptn:prop:proves-compact}, \ref{fol:ntd:ptn:prop:proves-compact}, \ref{fol:tab:ptn:prop:proves-compact}.
संपूर्णतेनुसार (\ref{fol:com:cth:thm:completeness}), $\Gamma$ सुसंगत
असल्यामुळे तो पूर्ततायोग्य आहे.
\end{proof}


\begin{prob}
\ref{fol:com:com:thm:compactness} मधील (1) सिद्ध करा.
\end{prob}





\begin{ex}
उपपत्ती~$\Gamma$ च्या प्रत्येक प्रतिमान~$\Struct{M}$ मध्ये प्रत्येक
पद~$t$ हे अर्थातच~$\Domain{M}$ मधील घटक निर्देशित करते.
$\Domain{M}$ मधील प्रत्येक घटक देखील कोणत्यातरी पदाने निर्देशित
केला जातो, याची हमी देता येते का? दुसऱ्या शब्दांत, ज्यांची सर्व प्रतिमाने
बंद पदांनी आच्छादित आहेत अशा उपपत्ती~$\Gamma$ असतात का? $\Gamma$ ला
अनंत प्रतिमाने असतील, तर असे नसते, हे संहतता प्रमेय दाखवते. ते पुढीलप्रमाणे
दिसते: $\Struct{M}$ हे~$\Gamma$ चे अनंत प्रतिमान असू द्या आणि $c$ हे
$\Gamma$ च्या भाषेत नसलेले स्थिरांक असू द्या. $\Gamma$ ची
भाषा~$\Lang{L}$ मधील प्रत्येक पद~$t$ साठी असलेल्या सर्व
$\eq/[c][t]$ या वाक्यांचा~$\Delta$ हा संच असू द्या, म्हणजे
  \[  \Delta = \Setabs{\eq/[c][t]}{t \in \Trm[L]}.
  \]
$\Gamma \cup \Delta$ चा सांत उपसंच $\Gamma' \cup \Delta'$ असा लिहिता
येतो; येथे $\Gamma' \subseteq \Gamma$ आणि $\Delta' \subseteq \Delta$.
$\Delta'$ सांत असल्यामुळे त्यात फक्त सांत अनेक पदे असू शकतात.
त्यांपैकी कोणत्याही पदाने निर्देशित न केलेला $\Domain{M}$ मधील
घटक $a \in \Domain{M}$ असू द्या आणि $\Struct{M'}$ ही
$\Struct{M}$ सारखीच, पण $\Assign{c}{M'} = a$ देखील असलेली संरचना
असू द्या. $\Delta'$ मध्ये येणाऱ्या प्रत्येक~$t$ साठी
$a \neq \Value{t}{M}$ असल्यामुळे $\Sat{M'}{\Delta'}$. तसेच
$\Sat{M}{\Gamma}$, $\Gamma' \subseteq \Gamma$, आणि~$\Gamma$ मध्ये $c$
येत नसल्यामुळे $\Sat{M'}{\Gamma'}$ देखील. म्हणून $\Gamma \cup \Delta$
च्या प्रत्येक सांत उपसंच $\Gamma' \cup \Delta'$ साठी
$\Sat{M'}{\Gamma' \cup \Delta'}$. त्यामुळे $\Gamma \cup \Delta$ चा
प्रत्येक सांत उपसंच पूर्ततायोग्य आहे. संहततेनुसार $\Gamma \cup \Delta$
स्वतः पूर्ततायोग्य आहे. म्हणून $\Sat{M}{\Gamma \cup \Delta}$ अशी
प्रतिमाने असतात. असे प्रत्येक $\Struct{M}$ हे~$\Gamma$ चे प्रतिमान आहे,
पण ते बंद पदांनी आच्छादित नाही; कारण~$\Lang{L}$ मधील सर्व पदे~$t$
यांच्यासाठी $\Value{c}{M} \neq \Value{t}{M}$.
\end{ex}

\begin{ex}
विधेय~$<$, स्थिरांक~$\Obj{0}$, $\Obj{1}$, आणि
फलने~$+$, $\times$, व~$-$ असलेल्या भाषा~$\Lang{L}$ चा
विचार करा. प्रांत~$\Rat$ आणि उघड निर्वचनार्थ असलेल्या
संरचना~$\Struct{Q}$ मध्ये सत्य असलेल्या या भाषा मधील सर्व
वाक्यांचा संच~$\Gamma$ असू द्या. $\Gamma$ म्हणजे परिमेय
संख्यांविषयी सत्य असलेल्या~$\Lang{L}$ मधील सर्व वाक्यांचा संच.
अर्थात~$\Rat$ मध्ये (आणि~$\Real$ मध्येही) अशा संख्या~$r$ नसतात की त्या
$0$ पेक्षा मोठ्या, पण प्रत्येक $k \in \PosInt$ साठी $1/k$ पेक्षा लहान
असतात. अशी संख्या अस्तित्वात असती, तर ती \emph{अतिसूक्ष्म संख्या} असती:
शून्येतर, पण अनंतपटीने लहान. $\Gamma$ ची अतिसूक्ष्म संख्या असलेली प्रतिमाने
अस्तित्वात आहेत, हे दाखवण्यासाठी संहतता प्रमेय वापरता येते. आपल्या भाषेत
भागाकाराचे फलन नाही (शून्याने भागाकार अपरिभाषित असतो आणि
फलनांचे निर्वचन सर्वत्र परिभाषित फलनांनी केले पाहिजे). तरीही
$r < 1/k$ हे व्यक्त करता येते; कारण हे तेव्हा आणि केवळ तेव्हाच लागू होते,
जेव्हा $r \cdot k < 1$. आता $c$ हे नवीन स्थिरांक असू द्या आणि
$\Delta$ पुढीलप्रमाणे असू द्या:
\[\{\Obj{0} < c\}
\cup \Setabs{ c \times \num k < \Obj{1} }{k \in \PosInt}
\]
(येथे $\num{k} = (\Obj{1} + (\Obj{1} + \dots + (\Obj{1} +
\Obj{1})\dots))$ मध्ये $k$ इतके~$\Obj{1}$ आहेत). $\Delta$ च्या कोणत्याही
सांत उपसंच~$\Delta_0$ साठी असा~$K$ असतो की~$\Delta_0$ मधील सर्व
वाक्ये $c \times \num{k} < \Obj{1}$ यांमध्ये $k < K$ असते.
$\Assign{c}{Q'} = 1/K$ असे ठेवून $\Struct{Q}$ चा विस्तार
$\Struct{Q'}$ पर्यंत केला, तर कोणत्याही सांत $\Gamma_0 \subseteq \Gamma$
साठी $\Sat{Q'}{\Gamma_0 \cup \Delta_0}$, आणि म्हणून $\Gamma \cup \Delta$
सांततः पूर्ततायोग्य आहे (सराव: हे सविस्तर सिद्ध करा). संहततेनुसार
$\Gamma \cup \Delta$ पूर्ततायोग्य आहे. $\Gamma \cup \Delta$ च्या
कोणत्याही प्रतिमान~$\Struct{S}$ मध्ये एक अतिसूक्ष्म संख्या असते, ती म्हणजे
$\Assign{c}{S}$.
\end{ex}



\begin{prob}
अंकगणिताच्या प्रमाण प्रतिमान~$\Struct{N}$ मध्ये प्रत्येक सूत्र~$\num{n}
< x$ ची पूर्ति करणारा कोणताही घटक~$k \in \Domain{N}$ नाही
(येथे $\num{n}$ हे $n$ इतकी~$\prime$ चिन्हे असलेले
$\Obj{0}^{\prime\dots\prime}$ आहे). अंकगणिताच्या प्रमाण
प्रतिमान~$\Struct{N}$ मध्ये सत्य असलेली अंकगणिताच्या भाषेतील
वाक्ये ही प्रत्येक सूत्र~$\num{n} < x$ ची पूर्ति \emph{करणारा}
घटक असलेल्या संरचना~$\Struct{N'}$ मध्येही सत्य आहेत,
हे दाखवण्यासाठी संहतता प्रमेय वापरा.
\end{prob}



\begin{ex}
एकरूपता असलेले प्रथम-कोटी तर्कशास्त्र हे प्रांताच्या आकाराला
एखादी किमान मर्यादा असली पाहिजे, असे व्यक्त करू शकते हे आपल्याला माहीत
आहे: वाक्य~$\varphi_{\ge n}$ (जे ``किमान $n$ भिन्न वस्तू आहेत'' असे सांगते)
फक्त अशा रचनांमध्ये सत्य असते ज्यांत $\Domain{M}$ मध्ये किमान~$n$ वस्तू
असतात. म्हणून
  \[  \Delta = \Setabs{\varphi_{\ge n}}{n \ge 1}
  \]
असे घेतले, तर~$\Delta$ चे कोणतेही प्रतिमान अनंत असले पाहिजे. म्हणून
एखाद्या उपपत्तीला फक्त अनंत प्रतिमानेच आहेत याची हमी तिच्यात~$\Delta$
मिळवून देता येते: $\Gamma \cup \Delta$ ची प्रतिमाने म्हणजे~$\Gamma$ ची
सर्व आणि फक्त अनंत प्रतिमाने.

अशा रीतीने प्रथम-कोटी तर्कशास्त्र अनंतता व्यक्त करू शकते. पण ते सांतता
व्यक्त करू शकत नाही, हे संहतता प्रमेय दाखवते. कारण वाक्यांचा
एखादा संच~$\Lambda$ सर्व आणि फक्त सांत संरचनांमध्ये पूर्त होतो
असे समजा. मग $\Delta \cup \Lambda$ सांततः पूर्ततायोग्य आहे. का? सांत
$\Delta' \cup \Lambda' \subseteq \Delta \cup \Lambda$ असा संच असू
द्या; येथे $\Delta' \subseteq \Delta$ आणि $\Lambda' \subseteq \Lambda$.
$\varphi_{\ge n} \in \Delta'$ असेल अशा संख्यांपैकी सर्वांत मोठी संख्या~$n$
असू द्या. $\Lambda$ ची सर्व सांत रचनांमध्ये पूर्ति होत असल्यामुळे,
त्याला सांत अनेक, पण $\ge n$ इतके घटक असलेले
प्रतिमान~$\Struct{M}$ असते. पण मग $\Sat{M}{\Delta' \cup \Lambda'}$.
संहततेनुसार $\Delta \cup \Lambda$ ला अनंत प्रतिमान असते; पण हे
$\Lambda$ ची पूर्ति फक्त सांत संरचनांमध्ये होते या गृहीतकाच्या
विरुद्ध आहे.
\end{ex}













\section{संहतता प्रमेयाची प्रत्यक्ष सिद्धता}\label{fol:com:cpd:sec}

संपूर्णता प्रमेयाच्या सिद्धतेतील कल्पनाच वापरून, संपूर्णता प्रमेयाचा आधार
न घेता संहतता प्रमेय प्रत्यक्ष सिद्ध करता येते. संपूर्णता प्रमेयाच्या
सिद्धतेत आपण वाक्यांच्या सुसंगत संच~$\Gamma$ पासून सुरुवात केली,
त्याचा विस्तार करून वाक्यांचा सुसंगत, संतृप्त आणि
संपूर्ण संच~$\Gamma^*$ मिळवला; आणि मग~$\Gamma^*$ पासून रचलेल्या
पद-प्रतिमान~$\Struct{M(\Gamma^*)}$
मध्ये~$\Gamma$ मधील सर्व वाक्ये सत्य आहेत, म्हणजे~$\Gamma$
पूर्ततायोग्य आहे, हे दाखवले.

सांततः पूर्ततायोग्य वाक्यसंच पूर्ततायोग्य असतो हे दाखवण्यासाठी हीच पद्धत
वापरता येते. फक्त सत्यता पूर्वप्रमेयापर्यंत नेणाऱ्या निष्पत्तींच्या अशा
समांतर आवृत्त्या सिद्ध कराव्या लागतील की त्यांत ``सुसंगत'' ऐवजी ``सांततः
पूर्ततायोग्य'' वापरलेले असेल.

\begin{prop}
\label{fol:com:cpd:prop:fsat-ccs}
$\Gamma$ हा संपूर्ण आणि सांततः पूर्ततायोग्य आहे असे समजा. मग:
\begin{enumerate}
\item $(\varphi \land \psi) \in \Gamma$
  तेव्हा आणि केवळ तेव्हाच, जेव्हा $\varphi \in \Gamma$ आणि $\psi \in \Gamma$
  ही दोन्ही विधाने लागू होतात.

\item $(\varphi \lor \psi) \in \Gamma$ तेव्हा आणि केवळ तेव्हाच,
  जेव्हा $\varphi \in \Gamma$ किंवा $\psi \in \Gamma$ यांपैकी एक लागू होते.

\item $(\varphi \lif \psi) \in \Gamma$ तेव्हा आणि केवळ तेव्हाच,
  जेव्हा $\varphi \notin \Gamma$ किंवा $\psi \in \Gamma$ यांपैकी एक लागू होते.
\end{enumerate}
\end{prop}


\begin{prob}
\ref{fol:com:cpd:prop:fsat-ccs} सिद्ध करा. $\Proves$ चा वापर टाळा.
\end{prob}





\begin{lem}
\label{fol:com:cpd:lem:fsat-henkin} प्रत्येक सांततः पूर्ततायोग्य संच~$\Gamma$ चा
विस्तार करून संतृप्त आणि सांततः पूर्ततायोग्य संच~$\Gamma'$ मिळवता येतो.
\end{lem}



\begin{prob}
\ref{fol:com:cpd:lem:fsat-henkin} सिद्ध करा. (सूचना: महत्त्वाची
पायरी म्हणजे निष्पत्त्या किंवा सुसंगततेचा अजिबात आधार न घेता,
$\Gamma_n$ सांततः पूर्ततायोग्य असेल, तर $\Gamma_n \cup \{\theta_n\}$
देखील सांततः पूर्ततायोग्य आहे हे दाखवणे.)
\end{prob}



\begin{prop}\label{fol:com:cpd:prop:fsat-instances}
$\Gamma$ हा संपूर्ण, सांततः पूर्ततायोग्य आणि संतृप्त आहे असे समजा.
\begin{enumerate}
\item $\lexists[x][\varphi(x)] \in \Gamma$ तेव्हा आणि केवळ तेव्हाच,
  जेव्हा किमान एका बंद पद~$t$ साठी $\varphi(t) \in \Gamma$.
\item $\lforall[x][\varphi(x)] \in \Gamma$ तेव्हा आणि केवळ तेव्हाच,
  जेव्हा सर्व बंद पदे~$t$ यांच्यासाठी $\varphi(t) \in \Gamma$.
\end{enumerate}
\end{prop}



\begin{prob}
\ref{fol:com:cpd:prop:fsat-instances} सिद्ध करा.
\end{prob}


\begin{lem}
\label{fol:com:cpd:lem:fsat-lindenbaum} प्रत्येक सांततः पूर्ततायोग्य संच~$\Gamma$
चा विस्तार करून संपूर्ण आणि सांततः पूर्ततायोग्य संच~$\Gamma^*$
मिळवता येतो.
\end{lem}


\begin{prob}
\ref{fol:com:cpd:lem:fsat-lindenbaum} सिद्ध करा. (सूचना: महत्त्वाची
पायरी म्हणजे $\Gamma_n$ सांततः पूर्ततायोग्य असेल, तर $\Gamma_n \cup
\{\varphi_n\}$ किंवा $\Gamma_n \cup \{\lnot \varphi_n\}$ यांपैकी एक संच सांततः
पूर्ततायोग्य असतो हे दाखवणे.)
\end{prob}




\begin{thm}[संहतता]
\label{fol:com:cpd:thm:compactness-direct} $\Gamma$ पूर्ततायोग्य असतो तेव्हा आणि
केवळ तेव्हाच, जेव्हा तो सांततः पूर्ततायोग्य असतो.
\end{thm}

\begin{proof}
$\Gamma$ पूर्ततायोग्य असेल, तर अशी
संरचना~$\Struct{M}$
असते की प्रत्येक $\varphi \in \Gamma$ साठी
$\Sat{M}{\varphi}$. अर्थात, ही
$\Struct{M}$ देखील~$\Gamma$ चा प्रत्येक
सांत उपसंच पूर्त करते; म्हणून~$\Gamma$ सांततः पूर्ततायोग्य आहे.

आता $\Gamma$ सांततः पूर्ततायोग्य आहे असे समजा.
  \ref{fol:com:cpd:lem:fsat-henkin} नुसार असा सांततः पूर्ततायोग्य, संतृप्त संच
  $\Gamma' \supseteq \Gamma$ असतो. \ref{fol:com:cpd:lem:fsat-lindenbaum} नुसार
$\Gamma'$ चा विस्तार करून संपूर्ण आणि सांततः
पूर्ततायोग्य संच~$\Gamma^*$ मिळवता येतो, आणि $\Gamma^*$
  अजूनही संतृप्त असतो. \ref{fol:com:mod:defn:termmodel} मधीलप्रमाणे
पद-प्रतिमान~$\Struct{M(\Gamma^*)}$
रचा. लक्षात घ्या की \ref{fol:com:mod:prop:quant-termmodel} हे
$\Gamma^*$ सुसंगत आहे (किंवा त्या दृष्टीने संपूर्ण वा संतृप्त आहे)
या तथ्यावर अवलंबून नव्हते; ते फक्त $\Struct{M(\Gamma^*)}$ हे बंद पदांनी
आच्छादित आहे या तथ्यावर अवलंबून होते.
सत्यता पूर्वप्रमेयाची सिद्धता (\ref{fol:com:mod:lem:truth}) तशीच लागू होते,
जर \ref{fol:com:ccs:prop:ccs} च्या संदर्भाऐवजी
\ref{fol:com:cpd:prop:fsat-ccs} चा संदर्भ वापरला आणि
\ref{fol:com:hen:prop:saturated-instances} च्या संदर्भाऐवजी
\ref{fol:com:cpd:prop:fsat-instances} चा संदर्भ वापरला.

\end{proof}


\begin{prob}
\ref{fol:com:cpd:thm:compactness-direct} च्या सिद्धतेला आवश्यक असलेली
सत्यता पूर्वप्रमेयाची (\ref{fol:com:mod:lem:truth}) संपूर्ण सिद्धता
लिहा.
\end{prob}














\section{लोव्हेनहाइम--स्कोलेम प्रमेय}\label{fol:com:dls:sec}

एखाद्या उपपत्तीला अनंत प्रतिमान असेल, तर तिला फार तर गणनीय अनंत
असलेले प्रतिमानही असते, असे लोव्हेनहाइम--स्कोलेम प्रमेय सांगते. या
तथ्याची तात्काळ निष्पत्ती अशी की एखाद्या संरचनेचा आकार
अगणनीय आहे हे प्रथम-कोटी तर्कशास्त्र व्यक्त करू शकत नाही: सर्व
अगणनीय संरचनांमध्ये पूर्त होणारे कोणतेही वाक्य
किंवा वाक्यांचा कोणताही संच हा एखाद्या गणनीय रचनेतही
पूर्त होतो.

\begin{thm}
\label{fol:com:dls:thm:downward-ls} $\Gamma$ सुसंगत असेल, तर त्याला
गणनीय प्रतिमान असते; म्हणजे ज्याचा प्रांत सांत किंवा
गणनीय अनंत आहे अशा संरचनांमध्ये तो पूर्ततायोग्य असतो.
\end{thm}

\begin{proof}
$\Gamma$ सुसंगत असेल, तर संपूर्णता प्रमेयाच्या सिद्धतेतून मिळणाऱ्या
संरचना~$\Struct M$ चा प्रांत~$\Domain{M}$ हा भाषा~$\Lang L$ मधील
पदांच्या संचापेक्षा मोठा नसतो. म्हणून $\Struct M$ फार तर
गणनीय अनंत आहे.
\end{proof}

\begin{thm}
\label{fol:com:dls:noidentity-ls} $\Gamma$ हा एकरूपतेविना प्रथम-कोटी तर्कशास्त्राच्या
भाषेतील वाक्यांचा सुसंगत संच असेल, तर त्याला गणनीय अनंत
प्रतिमान असते; म्हणजे ज्याचा प्रांत अनंत आणि गणनीय आहे अशा
संरचनांमध्ये तो पूर्ततायोग्य असतो.
\end{thm}

\begin{proof}
$\Gamma$ सुसंगत असेल आणि त्यात एकरूपता येणारी वाक्ये नसतील, तर संपूर्णता
प्रमेयाच्या सिद्धतेतून मिळणाऱ्या संरचना~$\Struct M$ चा
प्रांत~$\Domain{M}$ हा भाषा~$\Lang L'$ मधील पदांच्या संचाशी एकरूप असतो.
म्हणून $\Trm[L']$ हे गणनीय अनंत असल्यामुळे $\Struct{M}$ देखील
गणनीय अनंत आहे.
\end{proof}

\begin{ex}[स्कोलेमची विरोधापत्ती]
त्सेर्मेलो--फ्रेंकेल संचसिद्धान्त~$\Log{ZFC}$ ही एक अतिशय समर्थ चौकट
आहे; संचांच्या आकारांविषयीच्या तथ्यांसह जवळजवळ सर्व गणिती विधाने तिच्यात
व्यक्त करता येतात. म्हणून, उदाहरणार्थ, वास्तव संख्यांचा संच~$\Real$ हा
अगणनीय आहे हे~$\Log{ZFC}$ सिद्ध करू शकते; कोणत्याही संचाचा
घातसंच त्या संचापेक्षा मोठा असतो हे कँटर यांचे प्रमेयही ती सिद्ध करू शकते;
आणि असेच पुढे. $\Log{ZFC}$ सुसंगत असेल, तर तिची सर्व प्रतिमाने अनंत
असतात. शिवाय, त्या सर्व प्रतिमानांत असे घटक असतात की उपपत्तीच्या
मते ते अगणनीय असतात; उदाहरणार्थ, नैसर्गिक संख्यांचा घातसंच
अस्तित्वात आहे हे~$\Log{ZFC}$ चे प्रमेय सत्य करणारा घटक. लोव्हेनहाइम--स्कोलेम
प्रमेयानुसार $\Log{ZFC}$ ला गणनीय प्रतिमानेही असतात---ज्यांत
``अगणनीय'' संच असतात, पण जी स्वतः गणनीय असतात.
\end{ex}



\chapter{प्रथम-क्रम तर्कशास्त्राचा परिचय}\label{fol:int::chap}









\section{प्रथम-क्रम तर्कशास्त्र}\label{fol:int:fol:sec}

आकारिक तर्कशास्त्राशी तुमची पहिली ओळख झाली तेव्हापासून प्रथम-क्रम
तर्कशास्त्र बहुधा तुम्हाला परिचित असेल.\footnote{खरे तर, तुम्ही त्याच्याशी
परिचित आहात असे आम्ही जवळजवळ गृहीतच धरतो!{} तसे नसल्यास,
\emph{forall x} (Magnus आणि सहकारी, 2021) यासारख्या अधिक प्राथमिक
पाठ्यपुस्तकाची उजळणी करू शकता.} ते तुम्हाला ``संख्यापक तर्कशास्त्र''
किंवा ``विधेय तर्कशास्त्र'' या नावांनी माहीत असू शकते. सर्वप्रथम,
प्रथम-क्रम तर्कशास्त्र ही एक आकारिक भाषा आहे. म्हणजे तिचा एक विशिष्ट
शब्दसंग्रह असतो आणि तिच्या अभिव्यक्ती या त्या शब्दसंग्रहातील चिन्हमाला
असतात. पण प्रत्येक चिन्हमाला अनुमत नसते. अनुमत अभिव्यक्तींचे विविध
प्रकार आहेत: पदे, सूत्रे आणि वाक्ये. आपल्याला मुख्यतः
प्रथम-क्रम तर्कशास्त्रातील वाक्ये मध्ये रस आहे: ती आपल्याला
इंग्रजीतील वाक्यांचे आकारिक प्रतिरूप देतात आणि त्यांच्याविषयी
तर्कशास्त्रज्ञाला नेहमी रस असणारे प्रश्न आपण विचारू शकतो. उदाहरणार्थ:
\begin{itemize}
    \item $\psi$ हे~$\varphi$ पासून तार्किकरीत्या निष्पन्न होते का?
    \item $\varphi$ तार्किकरीत्या सत्य, तार्किकरीत्या असत्य, की
परिस्थितीवश आहे?
    \item $\varphi$ आणि $\psi$ सममूल्य आहेत का?
\end{itemize}

हे प्रश्न मुख्यतः प्रथम-क्रम तर्कशास्त्रातील वाक्यांच्या
``अर्था''विषयीचे आहेत. उदाहरणार्थ, $\psi$ हे~$\varphi$ पासून तार्किकरीत्या
निष्पन्न होते का या प्रश्नाचे तत्त्वज्ञ पुढीलप्रमाणे विश्लेषण करेल:
$\varphi$ सत्य पण~$\psi$ असत्य असणारा काही प्रसंग आहे का ($\psi$ हे~$\varphi$
पासून निष्पन्न होत नाही), की $\varphi$ सत्य करणारा प्रत्येक प्रसंग~$\psi$ ही
सत्य करतो ($\psi$ हे~$\varphi$ पासून निष्पन्न होते)? पण ``प्रसंग'' म्हणजे काय
हे अजून सांगितलेले नाही---ते \emph{चिन्हार्थमीमांसेचे} काम आहे.
प्रथम-क्रम तर्कशास्त्राची चिन्हार्थमीमांसा तत्त्वज्ञाच्या ``प्रसंग'' या
अंतःप्रज्ञ कल्पनेचे गणितीयदृष्ट्या काटेकोर प्रतिमान देते आणि---हे
महत्त्वाचे आहे---एखाद्या प्रसंगात वाक्य~$\varphi$ \emph{सत्य असणे}
म्हणजे काय याचेही प्रतिमान देते. आपण विकसित करणार असलेल्या या
गणितीयदृष्ट्या काटेकोर प्रतिमानाला संरचना म्हणतो. ``यात सत्य''
याला काटेकोर अर्थ देणाऱ्या संबंधाला \emph{पूर्ति}-संबंध म्हणतात. म्हणून
वाक्ये~$\varphi$ आणि संरचना~$\Struct{M}$ यांच्यासाठी आपण
``$\varphi$ ची~$\Struct{M}$ मध्ये पूर्ति होते'' (चिन्हांत:
$\Sat{M}{\varphi}$) हे परिभाषित करणार आहोत. हे झाल्यावर ``पासून निष्पन्न
होते'' किंवा ``तार्किकरीत्या सत्य आहे'' यांसारख्या इतर चिन्हार्थविषयक
संज्ञांच्याही काटेकोर व्याख्या देता येतील. उदाहरणार्थ,
$\lforall[x][(\varphi(x) \lif \psi(x))],
\lexists[x][\varphi(x)] \Entails \lexists[x][\psi(x)]$ आहे की नाही हे पुन्हा
गणितीय काटेकोरपणे ठरवणे या व्याख्यांमुळे शक्य होईल. अर्थातच उत्तर
``होय'' असेल. इंग्रजी वाक्यांचे प्रथम-क्रम तर्कशास्त्रात चिन्हांकन
करण्याचे प्रशिक्षण तुम्ही आधीच घेतले असेल, तर ही, उदाहरणार्थ, ``सर्व
मुंग्या कीटक आहेत, मुंग्या अस्तित्वात आहेत, म्हणून कीटक अस्तित्वात
आहेत'' यांची चिन्हांकने आहेत हे तुम्ही ओळखाल. हा उघडपणे वैध युक्तिवाद
आहे आणि म्हणून आपल्या आकारिक भाषेतील ``पासून निष्पन्न होते'' याच्या
गणितीय प्रतिमानानेही तेच उत्तर दिले पाहिजे.


आकारिक तर्कशास्त्राशी तुमची पहिली ओळख झाली तेव्हापासून आठवत असणारा
आणखी एक विषय म्हणजे \emph{निष्पत्त्या}. तुम्ही आकारिक
तर्कशास्त्राचा पहिला अभ्यासक्रम घेतला असेल, तर तुमच्या शिक्षकांनी अशा
निष्पत्त्या शोधण्याचा सराव तुमच्याकडून करून घेतला असेल; कदाचित
वरील तार्किक निष्पन्नता लागू असल्याचे दाखवणारी निष्पत्ती देखील.
निष्पत्त्या देण्याच्या अनेक पद्धती आहेत: तुम्ही ``नैसर्गिक निगमन''
किंवा ``सत्यता-वृक्ष'' नावाची पद्धत केली असेल; पण इतर अनेक पद्धतीही
आहेत. वरील तर्कशास्त्रीय प्रश्नांची उत्तरे देता येतील अशी साधने पुरवणे
हे निष्पत्ती-पद्धतींचे उद्दिष्ट आहे. उदाहरणार्थ, ज्या नैसर्गिक
निगमनातील निष्पत्तीमध्ये $\lforall[x][(\varphi(x) \lif \psi(x))]$ आणि
$\lexists[x][\varphi(x)]$ ही आधारविधाने आणि $\lexists[x][\psi(x)]$ हा
निष्कर्ष (शेवटची ओळ) आहे, ती $\lexists[x][\psi(x)]$ हे
$\lforall[x][(\varphi(x) \lif \psi(x))]$ आणि $\lexists[x][\varphi(x)]$ यांपासून
तार्किकरीत्या निष्पन्न होते याची \emph{पुष्टी करते}.

पण असे का? वरवर पाहता, निष्पत्ती-पद्धतींचा चिन्हार्थमीमांसेशी
काहीच संबंध नाही: आकारिक निष्पत्ती देणे म्हणजे केवळ काही
नियमांनी नियंत्रित रीतीने चिन्हांची मांडणी करणे; त्यात ``प्रसंगां''चा
किंवा ``यात सत्य''चा मुळीच उल्लेख नसतो. निष्पत्ती-पद्धती आणि
चिन्हार्थमीमांसा यांच्यातील संबंध अधितार्किक अन्वेषणाने प्रस्थापित
करावा लागतो. उदाहरणार्थ, आधारविधाने
$\lforall[x][(\varphi(x) \lif \psi(x))]$ आणि $\lexists[x][\varphi(x)]$ यांपासून
$\lexists[x][\psi(x)]$ ची आकारिक निष्पत्ती शक्य असते तेव्हा आणि
केवळ तेव्हाच $\lforall[x][(\varphi(x) \lif \psi(x))]$ आणि
$\lexists[x][\varphi(x)]$ ही दोन्ही मिळून
$\lexists[x][\psi(x)]$ तार्किकरीत्या निष्पन्न करतात, याची गणितीय
सिद्धता आवश्यक आहे. मात्र, हे करता येण्यापूर्वी विन्यासमीमांसा आणि
चिन्हार्थमीमांसा यांच्या व्याख्या अचूक करण्याचे बरेच कष्टप्रद काम
करावे लागते.












\section{विन्यासमीमांसा}\label{fol:int:syn:sec}

प्रथम-क्रम तर्कशास्त्रातील कोणत्या चिन्हमाला वाक्ये म्हणून गणल्या
जातात हे प्रथम काटेकोरपणे सांगावे लागेल. हे आपण नंतर करू; सध्या
उदाहरणांच्या साहाय्याने पुढे जाऊ. प्रथम-क्रम तर्कशास्त्राचे मूलभूत
रचनाघटक---म्हणजे त्याचा शब्दसंग्रह---दोन भागांत विभागलेले आहेत. पहिल्या
भागात विशिष्ट गोष्टींबद्दल काही सांगण्यासाठी किंवा विशिष्ट गोष्टी
निर्देशित करण्यासाठी आपण वापरतो ती चिन्हे येतात. गोष्टी निर्देशित
करण्यासाठी आपण स्थिरांक वापरतो आणि निर्देशित केलेल्या गोष्टींबद्दल
काही सांगण्यासाठी विधेये वापरतो. उदाहरणार्थ, एखादी एक गोष्ट
निर्देशित करण्यासाठी $\Obj a$ हे स्थिरांक म्हणून वापरू शकतो आणि मग
वाक्य~$\Atom{\Obj P}{\Obj a}$ वापरून तिच्याबद्दल काही सांगू शकतो.
``$\Obj a$'' आणि ``$\Obj P$'' यांचे काही अर्थ तुम्ही मनात धरले असतील,
तर $\Atom{\Obj P}{\Obj a}$ हे इंग्रजीतील वाक्य म्हणून वाचू शकता (आणि
आकारिक तर्कशास्त्र प्रथम शिकताना तुम्ही बहुधा तसे केलेही असेल).
प्रथम-क्रम तर्कशास्त्रातील अशी साधी वाक्ये मिळाल्यानंतर,
शब्दसंग्रहातील दुसरा भाग---तार्किक चिन्हे (संयोजक आणि संख्यापक)---वापरून
त्यांपासून अधिक जटिल वाक्ये रचता येतात. म्हणून, उदाहरणार्थ,
$(\Atom{\Obj P}{\Obj a} \land \Atom{\Obj Q}{\Obj b})$ किंवा
~$\lexists[\Obj x][\Atom{\Obj P}{\Obj x}]$ यांसारख्या अभिव्यक्ती
रचता येतात.

काटेकोर अधितार्किक अभ्यासासाठी आवश्यक असलेल्या चिन्हार्थमीमांसेच्या
नेमक्या व्याख्या आणि आपल्या निष्पत्ती-पद्धतींचे नियम देता यावेत,
यासाठी प्रथम-क्रम तर्कशास्त्रातील वाक्य म्हणून नेमके काय गणले
जाते याची काटेकोर व्याख्या सर्वप्रथम द्यावी लागेल. मूलभूत कल्पना समजणे
पुरेसे सोपे आहे: फक्त विधेये आणि स्थिरांक यांपासून आपण
$\Atom{\Obj P}{\Obj a}$ यासारखी काही साधी वाक्ये रचू शकतो. मग
संयोजक आणि संख्यापक वापरून त्यांपासून अधिक जटिल वाक्ये रचतो. पण अधिक
जटिल वाक्ये रचण्याची अनुमती देणारे नियम नेमके कोणते? ते स्पष्ट
करावेच लागतील; अन्यथा ``प्रथम-क्रम तर्कशास्त्रातील वाक्य'' याची
आपण पुरेशी काटेकोर व्याख्या केलेली नाही. येथे काही अडचणी आहेत. पहिली
अडचण म्हणजे योग्य चिन्हमालाच वाक्ये म्हणून गणल्या जातील याची
खात्री करणे. दुसरी म्हणजे हे अशा रीतीने करणे की
\emph{सर्व}~वाक्ये विषयी गणितीय सिद्धता देता येतील. शेवटी,
``$\varphi$ मधील प्रत्येक $\Obj x$ च्या जागी~$\Obj b$ ठेवा'' यासारख्या
वाक्ये वरील काही प्राथमिक क्रियांच्याही काटेकोर व्याख्या द्याव्या
लागतील. प्रथम-क्रम तर्कशास्त्रातील संख्यापक आणि चले यांमुळे हे
कसे करावे हे पूर्णपणे स्पष्ट नसते, ही अडचण आहे. उदाहरणार्थ,
$\lexists[\Obj x][\Atom{\Obj P}{\Obj a}]$ हे वाक्य म्हणून
गणले जावे का? $\lexists[\Obj x][\lexists[\Obj x][\Atom{\Obj P}{\Obj
x}]]$ चे काय? आणि ``$(\Atom{\Obj P}{\Obj x} \land \lexists[\Obj
x][\Atom{\Obj P}{\Obj x}])$ मध्ये $\Obj x$ च्या जागी~$\Obj b$ ठेवा''
याचा परिणाम काय असावा?












\section{सूत्र}\label{fol:int:fml:sec}

प्रथम-क्रम तर्कशास्त्रातील वाक्ये काटेकोरपणे निर्दिष्ट करण्यासाठी
आणि चलांच्या वापरामुळे निर्माण होणाऱ्या अडचणी हाताळण्यासाठी
आपण पुढील पद्धत वापरणार आहोत. प्रथम आपण अभिव्यक्तींचा एक
\emph{वेगळा} संच परिभाषित करतो: सूत्रे. ते केल्यानंतर त्यांत
चले कोणती भूमिका बजावतात याचा विचार करता येतो---आणि
``मुक्त'' व ``बद्ध'' चले या काही इतर कल्पनांच्या आधारे
वाक्य म्हणजे काय हे परिभाषित करता येते (म्हणजे, मुक्त
चले नसलेले सूत्र). यामुळे ``वाक्य'' ची व्याख्या
अधिक हाताळण्याजोगी होते एवढ्याचसाठी आपण हे करत नाही; पूर्ति ही
चिन्हार्थविषयक संकल्पना ज्या रीतीने परिभाषित करू त्यासाठीही ते
निर्णायक ठरेल.

फक्त एक विधेय~$\Obj P$, एक स्थिरांक~$\Obj a$ आणि
$\lnot$, $\land$ व~$\lexists$ ही तार्किक चिन्हे असलेल्या एका साध्या
प्रथम-क्रम भाषेसाठी ``सूत्र'' परिभाषित करू या. आपल्या संपूर्ण
व्याख्या यापेक्षा खूप व्यापक असतील: आपण अनंतपणे अनेक विधेये
आणि स्थिरांक स्वीकारू. किंबहुना, स्थिरांक आणि चले
यांच्याशी जोडून ``पदे'' रचता येतील अशा फलनांचाही विचार करू.
सध्या $\Obj a$ आणि चरे एवढीच आपली पदे असतील. आपल्याला अनंतपणे अनेक
चले मात्र आवश्यक आहेत. अधिकृतपणे आपण $\Obj v_0$, $\Obj
v_1$, \dots, ही चिन्हे चरे म्हणून वापरू.

\begin{defn}
\emph{सूत्रे} चा संच~$\Frm$ पुढीलप्रमाणे परिभाषित केला आहे:
\begin{enumerate}
\item\label{fol:int:fml:fmls-atom} $\Atom{\Obj P}{\Obj a}$ आणि $\Atom{\Obj
  P}{\Obj v_i}$ ही सूत्रे आहेत ($i \in \Nat$).

\item \label{fol:int:fml:fmls-not}$\varphi$ हे सूत्र असेल, तर $\lnot \varphi$ हे
  सूत्र आहे.

\item $\varphi$ आणि $\psi$ ही सूत्रे असतील, तर $(\varphi \land
  \psi)$ हे सूत्र आहे.

\item \label{fol:int:fml:fmls-ex}$\varphi$ हे सूत्र आणि $x$ हे चल
  असेल, तर $\lexists[x][\varphi]$ हे सूत्र आहे.

\item \label{fol:int:fml:fmls-limit}दुसरे काहीही सूत्र नाही.
\end{enumerate}
\end{defn}

कोणत्याही $i \in \Nat$ साठी $\Atom{\Obj P}{\Obj a}$ आणि
$\Atom{\Obj P}{\Obj v_i}$ ही सूत्रे आहेत, असे
\ref{fol:int:fml:fmls-atom} सांगतो. त्यांना तथाकथित \emph{आण्विक} सूत्रे
म्हणतात. त्यांपासून सुरुवात करता येते. इतर उपवाक्ये आपण आधीच
रचलेल्यांपासून नवी सूत्रे रचण्याचे मार्ग देतात. म्हणून,
उदाहरणार्थ, $\Atom{\Obj P}{\Obj v_2}$ हे \ref{fol:int:fml:fmls-atom} नुसार आधीच
सूत्र असल्यामुळे, $\lnot \Atom{\Obj P}{\Obj v_2}$ हे
\ref{fol:int:fml:fmls-not} नुसार सूत्र आहे. मग
\ref{fol:int:fml:fmls-ex} नुसार $\lexists[\Obj v_2][\lnot \Atom{\Obj P}{\Obj
v_2}]$ हे आणखी एक सूत्र आहे; आणि असेच पुढे जाता येते. \emph{फक्त} याच
रीतीने रचता येणाऱ्या चिन्हमाला सूत्रे म्हणून गणल्या जातात, असे
\ref{fol:int:fml:fmls-limit} सांगतो. विशेषतः,
$\lexists[\Obj v_0][\Atom{\Obj P}{\Obj a}]$ आणि $\lexists[\Obj
  v_0][\lexists[\Obj v_0][\Atom{\Obj P}{\Obj a}]]$ ही
\emph{खरोखर} सूत्रे म्हणून गणली जातात; परंतु अनावश्यक बाहेरील
कंसांमुळे $(\lnot \Atom{\Obj P}{\Obj a})$ हे गणले जात नाही.

सूत्रे परिभाषित करण्याच्या या पद्धतीला \emph{विगमनाधारित
व्याख्या} म्हणतात आणि त्यामुळे \emph{संरचनात्मक विगमन} म्हणतात त्या
विगमनाने सिद्धता करण्याच्या प्रकाराचा वापर करून सूत्रे विषयी
गोष्टी सिद्ध करता येतात. यांची सर्वसाधारण चर्चा
\ref{mth:ind:idf:sec} आणि \ref{mth:ind:sti:sec} येथे केली
आहे; पुढील सिद्धतांचा सखोल अभ्यास करण्यापूर्वी तिची उजळणी करावी. मूलतः,
सर्व सूत्रे साठी एखादी गोष्ट सत्य आहे याची सिद्धता द्यायची असेल,
तर ती आण्विक सूत्रे साठी सत्य आहे हे प्रथम दाखवता आणि नंतर ती
कोणत्याही सूत्र~$\varphi$ साठी (आणि~$\psi$ साठी) सत्य \emph{असेल}, तर
ती $\lnot \varphi$, $(\varphi \land \psi)$ आणि $\lexists[x][\varphi]$ यांच्यासाठीही
\emph{सत्य असते}, हे दाखवता. उदाहरणार्थ, यावरून ती
$\lexists[\Obj v_2][\lnot \Atom{\Obj P}{\Obj v_2}]$ साठी सत्य आहे हे
सिद्ध होते: पहिल्या भागावरून ती आण्विक सूत्र~$\Atom{\Obj
P}{\Obj v_2}$ साठी सत्य आहे हे माहीत असते. मग दुसऱ्या भागावरून ती
$\lnot \Atom{\Obj P}{\Obj v_2}$ साठी सत्य आहे आणि त्याच भागावरून पुन्हा
$\lexists[\Obj v_2][\lnot \Atom{\Obj P}{\Obj v_2}]$ याच्यासाठी सत्य
आहे, हे मिळते. सर्व सूत्रे आण्विक सूत्रांपासून
विगमनाधारित रीतीने निर्माण केली जात असल्यामुळे, ही पद्धत त्यांपैकी
कोणत्याही सूत्रासाठी चालते.












\section{पूर्ति}\label{fol:int:sat:sec}

सूत्रे म्हणजे काय हे माहीत झाल्यावर आपण लगेच प्रथम-क्रम
तर्कशास्त्राच्या चिन्हार्थमीमांसेकडे वळू शकतो: येथे मूलभूत व्याख्या
संरचनेची आहे. आपल्या साध्या भाषेसाठी संरचना~$\Struct M$
चे फक्त तीन घटक असतात: \emph{प्रांत} म्हणतात तो अरिक्त संच
$\Domain{M}$, $\Struct M$ मध्ये $\Obj a$ ने निर्देशित केलेली गोष्ट, आणि
$\Struct M$ मध्ये $\Obj P$ ज्या गोष्टींबाबत सत्य आहे त्या गोष्टी.
$\Obj a$ ने निर्देशित केलेली वस्तू~$\Assign{\Obj a}{M}$ ने, तर
$\Obj P$ ज्या गोष्टींबाबत सत्य आहे त्यांचा संच~$\Assign{\Obj P}{M}$ ने
दर्शवतात. संरचना~$\Struct{M}$ मध्ये याच तीन गोष्टी असतात:
$\Domain{M}$, $\Assign{\Obj a}{M} \in \Domain{M}$ आणि
$\Assign{\Obj P}{M} \subseteq \Domain{M}$. सर्वसाधारण स्थिती अधिक
गुंतागुंतीची असेल, कारण त्यात अनेक विधेये आणि स्थिरांक
असतील, विधेये ना एकाहून अधिक स्थाने असू शकतील, आणि
फलने ही असतील.


ज्या सूत्रे मध्ये चले नाहीत त्यांच्यासाठी पूर्तिची
व्याख्या देण्यास हे पुरेसे आहे. आपण सूत्रे ज्या रीतीने परिभाषित
केली तिचे प्रतिबिंब दाखवणारी विगमनाधारित व्याख्या देण्याची कल्पना आहे.
एखाद्या आण्विक सूत्राची~$\Struct{M}$ मध्ये पूर्ति केव्हा होते हे आपण
सांगतो आणि मग, उदाहरणार्थ, $\varphi$ ची~$\Struct{M}$ मध्ये पूर्ति होते की
नाही याच्या आधारे $\lnot \varphi$ ची~$\Struct{M}$ मध्ये पूर्ति केव्हा होते
हे सांगतो. उदाहरणार्थ, पुढीलप्रमाणे व्याख्या करता येईल:
\begin{enumerate}
  \item $\Atom{\Obj P}{\Obj a}$ ची~$\Struct{M}$ मध्ये पूर्ति तेव्हा आणि
  केवळ तेव्हाच होते, जेव्हा $\Assign{\Obj a}{M} \in \Assign{\Obj P}{M}$.
  \item $\lnot \varphi$ ची~$\Struct{M}$ मध्ये पूर्ति तेव्हा आणि केवळ तेव्हाच
  होते, जेव्हा $\varphi$ ची~$\Struct{M}$ मध्ये पूर्ति होत नाही.
  \item $(\varphi \land \psi)$ ची~$\Struct{M}$ मध्ये पूर्ति तेव्हा आणि केवळ
  तेव्हाच होते, जेव्हा $\varphi$ ची~$\Struct{M}$ मध्ये पूर्ति होते आणि
  $\psi$ ची~$\Struct{M}$ मध्येही पूर्ति होते.
\end{enumerate}
$\Domain{M} = \{0, 1, 2\}$, $\Assign{\Obj a}{M} = 1$ आणि
$\Assign{\Obj P}{M} = \{1, 2\}$ असे समजू. ही व्याख्या आपल्याला
$\Atom{\Obj P}{\Obj a}$ ची~$\Struct{M}$ मध्ये पूर्ति होते असे सांगेल
(कारण $\Assign{\Obj a}{M} =  1 \in \{1,2\} = \Assign{\Obj P}{M}$).
ती पुढे असे सांगते की $\lnot \Atom{\Obj P}{\Obj a}$ ची~$\Struct{M}$
मध्ये पूर्ति होत नाही; म्हणून $\lnot\lnot \Atom{\Obj P}{\Obj a}$ ची
पूर्ति होते आणि $(\lnot \Atom{\Obj P}{\Obj a} \land \Atom{\Obj P}{\Obj a})$
ची पूर्ति होत नाही; आणि असेच पुढे जाता येते.

संख्यापकांसाठी व्याख्या द्यायची झाल्यावर अडचण निर्माण होते: आपल्याला
``$\lexists[\Obj v_0][\Atom{\Obj P}{\Obj v_0}]$ ची पूर्ति तेव्हा आणि
केवळ तेव्हाच होते, जेव्हा $\Atom{\Obj P}{\Obj v_0}$ ची पूर्ति होते''
असे काहीतरी म्हणायचे आहे. पण संरचना~$\Struct{M}$ हे
चले विषयी काय करायचे ते सांगत नाही. आपल्याला प्रत्यक्षात
$\Atom{\Obj P}{\Obj v_0}$ ची \emph{$\Obj v_0$ च्या एखाद्या मूल्यासाठी}
पूर्ति होते असे म्हणायचे आहे. हे काटेकोर करण्यासाठी फक्त $\Obj a$ लाच
नव्हे, तर~$\Obj v_0$ लाही~$\Domain{M}$ मधील घटक नेमून देण्याचा
मार्ग हवा. यासाठी आपण चल ची \emph{मूल्यांकने} आणतो.
चल चे मूल्यांकन म्हणजे चले ना~$\Domain{M}$ मधील
घटक कडे पाठवणारे फलन~$s$ (आपल्या उदाहरणात, $0$, $1$ किंवा~$2$
यांपैकी एका घटकाकडे). एखाद्या सूत्रामध्ये कोणती चले
येतील हे आपल्याला आधी माहीत नसल्यामुळे, $s$ कोणत्या चले ना
मूल्ये नेमून देते यावर मर्यादा घालता येत नाही. सोपा उपाय म्हणजे $s$ ने
\emph{सर्व} चले~$\Obj v_0$, $\Obj v_1$, \dots\@ यांना मूल्ये
नेमून द्यावीत अशी अट ठेवणे. आपल्याला लागतील तीच मूल्ये आपण वापरू.


सूत्रांची पूर्ति फक्त संरचनेच्या सापेक्ष परिभाषित
करण्याऐवजी, आपण ती संरचना~$\Struct{M}$ \emph{आणि}
चल चे मूल्यांकन~$s$ या दोन्हींच्या सापेक्ष परिभाषित करू आणि
संक्षेपाने $\Sat{M}{\varphi}[s]$ असे लिहू. चले असलेली आण्विक
सूत्रे हाताळण्यासाठी आपल्या व्याख्येत आता एक अतिरिक्त उपवाक्य
असेल:
\begin{enumerate}
  \item $\Sat{M}{\Atom{\Obj P}{\Obj a}}[s]$ तेव्हा आणि केवळ तेव्हाच,
  जेव्हा $\Assign{\Obj a}{M} \in \Assign{\Obj P}{M}$.
  \item $\Sat{M}{\Atom{\Obj P}{\Obj v_i}}[s]$ तेव्हा आणि केवळ तेव्हाच,
  जेव्हा $s(\Obj v_i) \in \Assign{\Obj P}{M}$.
  \item $\Sat{M}{\lnot \varphi}[s]$ तेव्हा आणि केवळ तेव्हाच, जेव्हा
  $\Sat{M}{\varphi}[s]$ नाही.
  \item $\Sat{M}{(\varphi \land \psi)}[s]$ तेव्हा आणि केवळ तेव्हाच, जेव्हा
  $\Sat{M}{\varphi}[s]$ आणि $\Sat{M}{\psi}[s]$ दोन्ही आहेत.
\end{enumerate}
ठीक आहे, यामुळे एक अडचण सुटते: $s(\Obj v_0)$ या मूल्यासाठी
$\Struct{M}$ हे $\Atom{\Obj P}{\Obj v_0}$ ची पूर्ति केव्हा करते हे आता
सांगता येते. $\lexists[\Obj v_0][\Atom{\Obj P}{\Obj v_0}]$ साठी योग्य
व्याख्या मिळवण्यासाठी आणखी एक गोष्ट करावी लागेल:
$\Sat{M}{\lexists[\Obj v_0][\Atom{\Obj P}{\Obj v_0}]}[s]$ तेव्हा आणि
केवळ तेव्हाच असावे, जेव्हा~$\Obj v_0$ ला मूल्य नेमून देणाऱ्या
\emph{एखाद्या} पद्धती~$s'$ साठी
$\Sat{M}{\Atom{\Obj P}{\Obj v_0}}[s']$ असते. पण~$\Obj v_0$ ला नेमलेले
मूल्य $s(\Obj v_0)$ ने निर्देशित केलेले मूल्यच असणे आवश्यक नाही.
त्यासाठी आपण एक संकेतपद्धती आणू: $m \in \Domain{M}$ असेल, तर
$\Subst{s}{m}{\Obj v_0}$ हे~$s$ सारखेच असणारे (फक्त~$\Obj v_0$ सोडून
इतर सर्व चले साठी) पण~$\Obj v_0$ ला~$m$ नेमून देणारे
मूल्यांकन मानू. आता आपली व्याख्या अशी देता येईल:
\begin{enumerate}\setcounter{enumi}{4}
  \item $\Sat{M}{\lexists[\Obj v_i][\varphi]}[s]$ तेव्हा आणि केवळ तेव्हाच,
  जेव्हा एखाद्या~$m \in \Domain{M}$ साठी
  $\Sat{M}{\varphi}[\Subst{s}{m}{\Obj v_i}]$.
\end{enumerate}
हे योग्य रीतीने चालते का? प्रत्येक~$i \in \Nat$ साठी
$s(\Obj v_i) = 0$ ठेवू या. एखादा $m \in \Domain{M}$ असा असेल की
$\Sat{M}{\Atom{\Obj P}{\Obj v_0}}[\Subst{s}{m}{\Obj v_0}]$, तेव्हा आणि
केवळ तेव्हाच $\Sat{M}{\lexists[\Obj v_0][\Atom{\Obj P}{\Obj v_0}]}[s]$.
आणि असा घटक आहे: आपण $m = 1$ किंवा $m = 2$ निवडू शकतो. $s$ नेच
~$\Obj v_0$ ला नेमून दिलेले $s(\Obj v_0)$ हे मूल्य---या उदाहरणात,
$0$---हे काम करत नसले तरी हे सत्य आहे, हे लक्षात घ्या. आपल्याकडे
$\Sat{M}{\Atom{\Obj P}{\Obj v_0}}[\Subst{s}{1}{\Obj v_0}]$ आहे, पण
$\Sat{M}{\Atom{\Obj P}{\Obj v_0}}[s]$ नाही.

हे गोंधळात टाकणारे आणि अवजड वाटत असेल, तर ते तसे आहे. पण सर्व
सूत्रे साठी पूर्तिची काटेकोर, विगमनाधारित व्याख्या देण्यासाठी ही
अतिरिक्त गुंतागुंत आवश्यक आहे आणि चिन्हार्थविषयक संकल्पना काटेकोरपणे
परिभाषित करण्यासाठी आपल्याला अशाच काहीतरी पद्धतीची गरज आहे. हे
करण्याचे इतरही मार्ग आहेत, पण ते सर्व तितकेच (अन)सुबक आहेत.












\section{वाक्य}\label{fol:int:snt:sec}

ठीक आहे, आता आपल्याकडे संरचना आणि सूत्रे यांच्यासाठी
पूर्तिच्या (``यात सत्य'') व्याख्येचा आराखडा आहे. पण तिला हा अतिरिक्त
भाग---चल चे मूल्यांकन---आवश्यक आहे आणि आपल्याला हवी होती ती
वाक्यांची व्याख्या. मूल्यांकने कशी काढून टाकायची आणि
वाक्ये म्हणजे काय?

आकारिक तर्कशास्त्राशी तुमची पहिली ओळख झाली तेव्हा चले आणि
संख्यापक यांच्यातील संबंधाची केलेली चर्चा तुम्हाला बहुधा आठवत असेल.
संख्यापकानंतर नेहमी चल येते आणि मग त्या संख्यापकाचा
वाक्य मधील ज्या भागाला उपयोग होतो (त्याची ``व्याप्ती''), त्या
भागात ते चल त्या संख्यापकाने ``बद्ध'' केले आहे असे आपण
समजतो. सूत्रे मध्ये प्रत्येक चल ला जुळणारा संख्यापक
असणे आवश्यक नव्हते आणि जुळणारे संख्यापक नसलेली चले ही
``मुक्त'' किंवा ``अबद्ध'' असतात. मुक्त चले नसलेली सर्व
सूत्रे म्हणजे वाक्ये असे आपण मानू.

सूत्र~$\varphi$ मधील चल ची एखादी आस्थिती केव्हा बद्ध असते,
तिला कोणता संख्यापक बद्ध करतो आणि ती केव्हा मुक्त असते, याची अंतःप्रज्ञ
कल्पना समजणे पुन्हा अवघड नाही. हे तपासण्याची, कदाचित कंस मोजण्याची,
एखादी पद्धत तुम्ही शिकला असाल. पण आपण काटेकोर व्याख्येचा आग्रह धरला
पाहिजे---आणि सूत्रांची व्याख्या आपण विगमनाने केली असल्यामुळे,
सूत्र~$\varphi$ मधील चल~$x$ च्या मुक्त आणि बद्ध आस्थितींची
व्याख्याही विगमनाने देता येते. उदाहरणार्थ, आपल्या सोप्या भाषेसाठी ती
अशी दिसू शकेल:
\begin{enumerate}
  \item $\varphi$ हे आण्विक असेल, तर त्यातील $x$ च्या सर्व आस्थिती मुक्त
  असतात (म्हणजे, $\Atom{\Obj P}{x}$ मधील $x$ ची आस्थिती मुक्त असते).
  \item $\varphi$ हे $\lnot \psi$ या रूपाचे असेल, तर $\lnot \psi$ मधील~$x$ ची
  आस्थिती तेव्हा आणि केवळ तेव्हाच मुक्त असते, जेव्हा~$\psi$ मधील~$x$ ची
  संबंधित आस्थिती मुक्त असते (म्हणजे,~$\psi$ मधील चरांच्या मुक्त
  आस्थिती म्हणजे नेमक्या $\lnot \psi$ मधील संबंधित आस्थिती होत).
  \item $\varphi$ हे $(\psi \land \chi)$ या रूपाचे असेल, तर $(\psi \land \chi)$
  मधील~$x$ ची आस्थिती तेव्हा आणि केवळ तेव्हाच मुक्त असते, जेव्हा~$\psi$
  किंवा~$\chi$ मधील~$x$ ची संबंधित आस्थिती मुक्त असते.
  \item $\varphi$ हे $\lexists[x][\psi]$ या रूपाचे असेल, तर $\varphi$ मधील $x$ ची
  कोणतीही आस्थिती मुक्त नसते; ते $\lexists[y][\psi]$ या रूपाचे असेल आणि
  $y$ हे~$x$ पेक्षा वेगळे चल असेल, तर
  $\lexists[y][\psi]$ मधील~$x$ ची आस्थिती तेव्हा आणि केवळ तेव्हाच मुक्त
  असते, जेव्हा~$\psi$ मधील~$x$ ची संबंधित आस्थिती मुक्त असते.
\end{enumerate}

चरांच्या मुक्त आणि बद्ध आस्थितींची काटेकोर व्याख्या मिळाल्यानंतर आपण
सरळ असे म्हणू शकतो: मुक्त चलांची आस्थिती नसलेले कोणतेही
सूत्र म्हणजे वाक्य होय.












\section{चिन्हार्थविषयक संकल्पना}\label{fol:int:sem:sec}

$\Sat{M}{\varphi}[s]$ लागू आहे का याचा विचार करताना, सोयीसाठी आपण $s$ कडून
सर्व चले ना मूल्ये नेमून घेतो, पण~$\varphi$ मधील चले ना
त्याने नेमून दिलेली मूल्येच फक्त वापरली जातात, असे आपण वर म्हटले.
खरे तर, $\varphi$ मधील \emph{मुक्त} चरांची मूल्येच महत्त्वाची असतात. आपण
काटेकोर असल्यामुळे अर्थातच हे तथ्य सिद्ध करणार आहोत. वाक्ये मध्ये
मुक्त चरे नसल्यामुळे, संरचनांमध्ये त्यांची पूर्ति होते की नाही
यासाठी $s$ चा काहीही फरक पडत नाही. म्हणून, $\varphi$ हे वाक्य असेल,
तेव्हा ``प्रत्येक~$s$ साठी $\Sat{M}{\varphi}[s]$'' असा $\Sat{M}{\varphi}$ चा अर्थ
परिभाषित करता येतो; आणि तो योगायोगाने किमान एका~$s$ साठी
$\Sat{M}{\varphi}[s]$ असते तेव्हा आणि केवळ तेव्हाच सत्य असतो. सूत्रे
साठी पूर्तिची कार्यक्षम व्याख्या मिळवण्यासाठी आपल्याला चल ची
मूल्यांकने आणावी लागतात, पण वाक्ये साठी पूर्ति ही चल
च्या मूल्यांकनांपासून स्वतंत्र असते.

``$\Sat{M}{\varphi}$'' ची व्याख्या मिळाल्यानंतर प्रथम-क्रम तर्कशास्त्रातील
वाक्यांच्या संदर्भात ``प्रसंग'' आणि ``यात सत्य'' यांचा अर्थ काय
हे आपल्याला माहीत होते. मग वाक्ये साठीच्या $\Sat{M}{\varphi}$ च्या
व्याख्येच्या आधारे वैधता, तार्किक निष्पन्नता आणि पूर्ततायोग्यता या मूलभूत
चिन्हार्थविषयक संकल्पना परिभाषित करता येतात. प्रत्येक संरचना
एखाद्या वाक्याची पूर्ति करत असेल, तर ते वाक्य वैध, $\Entails \varphi$, असते.
~$\Gamma$ मधील सर्व वाक्यांची पूर्ति करणारे प्रत्येक
संरचना हे~$\varphi$ चीही पूर्ति करत असेल, तर वाक्यांच्या संचातून
ते तार्किकरीत्या निष्पन्न होते, $\Gamma \Entails \varphi$. आणि एखादे
संरचना एखाद्या वाक्यांच्या संचातील सर्व वाक्यांची एकाच
वेळी पूर्ति करत असेल, तर तो संच पूर्ततायोग्य असतो.

सूत्रे विगमनाधारित रीतीने परिभाषित केली आहेत आणि पूर्ति हीही
सूत्रांच्या संरचनेवरील विगमनाने परिभाषित केली आहे; म्हणून आपल्या
चिन्हार्थमीमांसेचे गुणधर्म सिद्ध करण्यासाठी आणि परिभाषित
चिन्हार्थविषयक संकल्पनांचा परस्परसंबंध लावण्यासाठी विगमन वापरता येते.
यांपैकी काही गुणधर्म आपण एकत्र करून सिद्ध करू; काही अंशी प्रत्येक
गुणधर्म स्वतंत्रपणे रोचक असल्यामुळे, पण मुख्यतः चिन्हार्थमीमांसा आणि
निष्पत्ती-पद्धती यांच्यातील संबंधाचा अभ्यास पुढे करताना त्यांपैकी
अनेक उपयोगी पडणार असल्यामुळे. हे करण्यासाठी आपण अजून उल्लेख न केलेल्या
काही विन्यासविषयक संकल्पना आणि क्रियाही (काटेकोरपणे, म्हणजे विगमनाने)
परिभाषित कराव्या लागतील.












\section{आदेशन}\label{fol:int:sub:sec}

सर्व गोष्टींचा परस्परसंबंध कसा आहे आणि विन्यासमीमांसा व
चिन्हार्थमीमांसा यांचा विकास पुढील अधिक प्रगत अभ्यासाचा पाया कसा घालतो,
हे दाखवण्यासाठी एका उदाहरणाची चर्चा करू. आपल्या निष्पत्ती-पद्धतींनी
आपल्याला $\lforall[\Obj v_0][\Atom{\Obj P}{\Obj v_0}]$ पासून
$\Atom{\Obj P}{\Obj a}$ निष्पन्न करू दिले पाहिजे. कदाचित हे निगमन नियम
म्हणूनही मांडायचे असेल. पण त्यासाठी ते शक्य तितक्या सर्वसाधारण शब्दांत
मांडता आले पाहिजे: फक्त $\Obj P$, $\Obj a$ आणि $\Obj v_0$ यांच्यासाठी
नव्हे, तर कोणत्याही सूत्र~$\varphi$, पद~$t$ आणि चल~$x$
यांच्यासाठी. (स्थिरांक ही पदे असतात हे आठवा; पण स्थिरांक आणि
फलने यांच्यापासून रचलेल्या अधिक जटिल पदांचाही आपण विचार करू.)
म्हणून आपल्याला असे काहीतरी म्हणता आले पाहिजे: ``जेव्हा तुम्ही
$\lforall[x][\varphi(x)]$ निष्पन्न केले असेल, तेव्हा~$\varphi(t)$ चे अनुमान
करणे समर्थित असते---$\lforall[x]$ काढून आणि~$x$ च्या जागी~$t$ ठेवून
मिळणारे हे सूत्र आहे.'' पण ``$x$ च्या जागी~$t$ ठेवणे'' म्हणजे नेमके
काय? $\varphi(x)$ आणि~$\varphi(t)$ यांच्यातील संबंध कोणता? हे नेहमी चालते का?


हे काटेकोर करण्यासाठी आपण \emph{आदेशन} ही क्रिया परिभाषित करतो.
आदेशन प्रत्यक्षात गुंतागुंतीचे आहे, कारण~$\varphi$ मधील सर्व~$x$ फक्त~$t$ ने
बदलता येत नाहीत आणि प्रत्येक~$t$ चे कोणत्याही~$x$ साठी आदेशन करता येत
नाही. आपण हे पुन्हा विगमनाधारित व्याख्या वापरून हाताळू. पण हे
केल्यानंतर ``$\lforall[x][\varphi(x)]$ पासून $\varphi(t)$ चे अनुमान करा'' असा
निगमन नियम निर्दिष्ट करणे ही काटेकोर व्याख्या बनते. शिवाय,
$\lforall[x][\varphi(x)]$ पासून~$\varphi(t)$ तार्किकरीत्या निष्पन्न होते या अर्थाने
हा चांगला निगमन नियम आहे, हे दाखवता येईल. पण हे सिद्ध करण्यासाठी प्रथम
दर्शनी ``आपण हे का करत आहोत?'' असा प्रश्न पडेल अशी आणखी एक गोष्ट पुन्हा
सिद्ध करावी लागेल. $\lforall[x][\varphi(x)]$ पासून~$\varphi(t)$ तार्किकरीत्या
निष्पन्न होते हे पुढील तथ्यावर अवलंबून आहे: $\Sat{M}{\varphi(t)}$ लागू आहे
की नाही हे फक्त पद~$t$ च्या मूल्यावर अवलंबून असते. म्हणजे, $t$ ने
निर्देशित होणारा~$\Domain{M}$ मधील कोणताही घटक आपण $m$ मानला,
तर $\Sat{M}{\varphi(t)}[s]$ तेव्हा आणि केवळ तेव्हाच, जेव्हा
$\Sat{M}{\varphi(x)}[\Subst{s}{m}{x}]$. $t$ मध्ये चले असली तरीही
हे लागू असते; पण हा परिणाम नेमका कसा मांडायचा याबाबत काळजी घ्यावी
लागेल.












\section{प्रतिमाने आणि उपपत्ती}\label{fol:int:mod:sec}

प्रथम-क्रम तर्कशास्त्राची विन्यासमीमांसा आणि चिन्हार्थमीमांसा परिभाषित
केल्यानंतर संरचनेचे गुणधर्म आणि चिन्हार्थविषयक संकल्पना यांचा
अभ्यास सुरू करता येतो. निष्पत्ती-पद्धतीही परिभाषित करून त्यांचा
अभ्यास करता येतो. वाक्यांचा एखादा संच घेतल्यावर आपण विचारू
शकतो: त्या संचातील सर्व वाक्ये कोणती संरचना सत्य करतात?
वाक्यांचा संच~$\Gamma$ दिला असता, त्यांची पूर्ति करणाऱ्या
संरचना~$\Struct{M}$ ला \emph{$\Gamma$ चे प्रतिमान} म्हणतात.
~$\Gamma$ पासून सुरुवात करून आपण त्याची प्रतिमाने शोधू शकतो---ती कशी
दिसतात? ती किती मोठी किंवा लहान असावी लागतात? पण एखादे एक
संरचना किंवा संरचनेचा संग्रह घेऊनही विचारता येते: त्यांत
कोणती वाक्ये सत्य आहेत? या संरचनांमध्ये ती आणि फक्त तीच
सत्य असतात या अर्थाने त्यांचे \emph{अभिलक्षण} करणारी वाक्ये आहेत
का? अशा प्रकारचे प्रश्न \emph{प्रतिमान उपपत्ती} या क्षेत्रातील आहेत.
त्यांवरच \emph{स्वयंसिद्धकीय पद्धत} ही आधारित असते: एखाद्या
संरचनेच्या संग्रहाचे उपपत्तीच्या स्वयंसिद्धकांचा, म्हणजे
वाक्यांच्या संचाचा, वापर करून वर्णन करणे. स्वयंसिद्धकांपासून
प्रथम-क्रम तर्कशास्त्रात तार्किकरीत्या निष्पन्न होणारी नेमकी तीच
वाक्ये स्वयंसिद्धकांच्या सर्व प्रतिमानांत सत्य असतात, या
निरीक्षणामुळे हे शक्य होते.

अतिशय सोपे उदाहरण म्हणून पूर्वक्रमांचा विचार करा. एखाद्या संच~$A$ वरील
परावर्ती आणि संक्रमक असा संबंध~$R$ म्हणजे पूर्वक्रम होय. द्विस्थानी
संबंध $R \subseteq A \times A$ असलेला संच~$A$ म्हणजे एकच द्विस्थानी
संबंधचिन्ह~$\Obj P$ असलेल्या प्रथम-क्रम भाषेला संरचना देण्यासाठी
नेमके जे आवश्यक आहे ते: आपण $\Domain{M} = A$ आणि
$\Assign{\Obj P}{M} = R$ ठेवू. $R$ हा पूर्वक्रम असल्यामुळे तो परावर्ती
आणि संक्रमक आहे आणि हे सांगणारा प्रथम-क्रम तर्कशास्त्रातील
वाक्यांचा संच~$\Gamma$ आपल्याला मिळू शकतो:
\begin{align*}
  & \lforall[\Obj v_0][\Atom{\Obj P}{\Obj v_0,\Obj v_0}]\\
  & \lforall[\Obj v_0][\lforall[\Obj v_1][\lforall[\Obj v_2][((\Atom{\Obj P}{\Obj v_0,\Obj v_1} \land \Atom{\Obj P}{\Obj v_1,\Obj v_2}) \lif \Atom{\Obj P}{\Obj v_0,\Obj v_2})]]]
\end{align*}
ही वाक्ये म्हणजे फक्त ``कोणत्याही~$x$ साठी $Rxx$'' ($R$ हा
परावर्ती आहे) आणि ``जेव्हा $Rxy$ आणि $Ryz$, तेव्हा $Rxz$ ही'' ($R$ हा
संक्रमक आहे) यांची चिन्हांकने आहेत. संरचना~$\Struct{M}$ हे या
दोन वाक्ये~$\Gamma$ चे प्रतिमान तेव्हा आणि केवळ तेव्हाच असते,
जेव्हा $R$ (म्हणजे $\Assign{\Obj P}{M}$) हा~$A$ (म्हणजे
$\Domain{M}$) वरील पूर्वक्रम असतो. दुसऱ्या शब्दांत, $\Gamma$ ची प्रतिमाने
म्हणजे नेमके पूर्वक्रम होत. फक्त~$\Obj P$ हे विधेय असलेल्या
प्रथम-क्रम भाषेत व्यक्त करता येणारा सर्व पूर्वक्रमांचा कोणताही गुणधर्म
(वरील परावर्तिता आणि संक्रमणता यांसारखा) हा~$\Gamma$ मधील दोन
वाक्यांपासून तार्किकरीत्या निष्पन्न होतो आणि उलटही तसेच होते.
म्हणून~$\Gamma$ च्या प्रतिमानांविषयी आपण जे काही सिद्ध करू शकतो, ते सर्व
पूर्वक्रमांविषयी सिद्ध केलेले असते.

कोणत्याही विशिष्ट उपपत्तीसाठी आणि प्रतिमानांच्या वर्गासाठी (जसे
$\Gamma$ आणि सर्व पूर्वक्रम) संबंधित प्रथम-क्रम भाषेत काय व्यक्त करता
येते आणि काय व्यक्त करता येत नाही याविषयी रोचक प्रश्न असतील. सर्व भाषा
आणि प्रतिमानवर्गांसाठी रोचक असणारे संरचनेचे काही गुणधर्म आहेत;
विशेषतः प्रांताच्या आकारमानाशी संबंधित गुणधर्म. उदाहरणार्थ,
कोणत्याही~$n \in \PosInt$ साठी प्रांत मध्ये नेमके
$n$~घटक आहेत असे नेहमी व्यक्त करता येते. अनंतपणे अनेक
वाक्यांचा संच वापरून प्रांत अनंत आहे असेही व्यक्त करता येते.
पण प्रांत सांत आहे किंवा प्रांत अगणनीय आहे असे व्यक्त करता
येत नाही. प्रथम-क्रम भाषांच्या मर्यादांविषयीचे हे परिणाम संहतता आणि
लोव्हेनहाइम--स्कोलेम प्रमेयांचे परिणाम आहेत.












\section{निर्दोषता आणि संपूर्णता}\label{fol:int:scp:sec}

प्रथम-क्रम तर्कशास्त्रासाठीच्या निष्पत्ती-पद्धतींचाही आपण परिचय करून
घेऊ. तर्कशास्त्रज्ञांनी अनेक निष्पत्ती-पद्धती विकसित केल्या आहेत,
पण त्या सर्व वाक्ये दरम्यान तोच निष्पन्नता संबंध परिभाषित
करतात. विशिष्ट आणि काटेकोरपणे परिभाषित प्रकारची निष्पत्ती असेल,
तर~$\Gamma$ हे~$\varphi$ \emph{निष्पन्न} करते, $\Gamma \Proves \varphi$, असे
आपण म्हणतो. निष्पत्त्या म्हणजे नेहमी चिन्हांच्या सांत मांडण्या
असतात---कदाचित वाक्यांची यादी किंवा त्याहून गुंतागुंतीची काही
संरचना. एखाद्या संच~$\Gamma$ पासून वाक्य तार्किकरीत्या निष्पन्न
होते का हे ठरवण्याचे साधन पुरवणे हे निष्पत्ती-पद्धतींचे उद्दिष्ट
आहे. ते उद्दिष्ट साध्य करण्यासाठी $\Gamma \Entails \varphi$ तेव्हा आणि केवळ
तेव्हाच $\Gamma \Proves \varphi$ असले पाहिजे.

$\Gamma \Proves \varphi$ पण $\Gamma \Entails \varphi$ नसेल, तर आपली
निष्पत्ती-पद्धत अतिप्रबल असेल, म्हणजे ती गरजेपेक्षा जास्त सिद्ध
करेल. $\Gamma \Proves \varphi$ असेल तर $\Gamma \Entails \varphi$ असते या
गुणधर्माला \emph{निर्दोषता} म्हणतात आणि कोणत्याही चांगल्या
निष्पत्ती-पद्धतीसाठी ही किमान आवश्यकता आहे. दुसरीकडे,
$\Gamma \Entails \varphi$ पण $\Gamma \Proves \varphi$ नसेल, तर आपली
निष्पत्ती-पद्धत अतिदुर्बल असेल, म्हणजे ती पुरेसे सिद्ध करणार नाही.
$\Gamma \Entails \varphi$ असेल तर $\Gamma \Proves \varphi$ असते या गुणधर्माला
\emph{संपूर्णता} म्हणतात. निर्दोषता सिद्ध करणे सहसा तुलनेने सोपे असते
(विगमनाधारित रीतीने परिभाषित केलेल्या निष्पत्त्यांच्या संरचनेवर
विगमन करून). संपूर्णता सिद्ध करणे अधिक कठीण असते.

निर्दोषता आणि संपूर्णता यांचे अनेक महत्त्वाचे परिणाम आहेत.
वाक्यांचा एखादा संच~$\Gamma$ हा व्याघात (जसे
$\varphi \land \lnot \varphi$) निष्पन्न करत असेल, तर त्याला \emph{विसंगत}
म्हणतात. विसंगत~$\Gamma$ ला कोणतेही प्रतिमान असू शकत नाही; तो
अपूर्ततायोग्य असतो. संपूर्णतेवरून याचा उलट निष्कर्ष मिळतो: विसंगत
नसलेल्या---किंवा, आपण म्हणू त्याप्रमाणे, \emph{सुसंगत}---प्रत्येक
$\Gamma$ ला प्रतिमान असते. खरे तर, हे संपूर्णतेशी सममूल्य आहे आणि
संपूर्णतेचे हेच रूप आपण प्रत्यक्षात सिद्ध करू. हा सखोल आणि कदाचित
आश्चर्यकारक परिणाम आहे: $\Gamma$ पासून $\varphi \land \lnot \varphi$ सिद्ध करता
येत नाही एवढ्यामुळेच~$\Gamma$ वर्णन करते तसे संरचना अस्तित्वात
असण्याची हमी मिळते. म्हणून संपूर्णता पुढील प्रश्नाचे उत्तर देते:
वाक्यांच्या कोणत्या संचांना प्रतिमाने असतात? उत्तर: सर्व आणि फक्त
सुसंगत संचांना.

निर्दोषता आणि संपूर्णता प्रमेयांचे दोन महत्त्वाचे परिणाम आहेत: संहतता
प्रमेय आणि लोव्हेनहाइम--स्कोलेम प्रमेय. प्रतिमानांच्या उपपत्तीतील हे
महत्त्वाचे परिणाम आहेत आणि अनेक रोचक निष्कर्ष प्रस्थापित करण्यासाठी
त्यांचा उपयोग करता येतो. आपण आधीच दोन निष्कर्षांचा उल्लेख केला आहे:
प्रथम-क्रम तर्कशास्त्रात संरचनेचा प्रांत सांत आहे किंवा तो
अगणनीय आहे असे व्यक्त करता येत नाही.

ऐतिहासिकदृष्ट्या हे सर्व---प्रथम-क्रम तर्कशास्त्राची विन्यासमीमांसा आणि
चिन्हार्थमीमांसा कशी परिभाषित करायची, चांगल्या निष्पत्ती-पद्धती
कशा परिभाषित करायच्या, त्या निर्दोष आणि संपूर्ण आहेत हे कसे सिद्ध
करायचे, प्रथम-क्रम भाषांत काय व्यक्त करता येते आणि काय व्यक्त करता येत
नाही हे स्पष्ट करणे---शोधून योग्य रीतीने ठरवण्यास बराच काळ लागला. हे
कसे करायचे हे आता माहीत असले, तरी सर्व तपशील समजून घेणे अजूनही
गोंधळात टाकणारे आणि कंटाळवाणे ठरू शकते. पण ते महत्त्वाचेही आहे, कारण
प्रथम-क्रम तर्कशास्त्राच्या आकारिक भाषेसाठी येथे विकसित केलेल्या पद्धती
तर्कशास्त्र, संगणकशास्त्र आणि भाषाविज्ञान यांत सर्वत्र वापरल्या जातात.
म्हणून सर्व तपशीलांचा अभ्यास करण्याचे दीर्घकाळात फळ मिळते.



\chapter{प्रथम-क्रम तर्कशास्त्राची विन्यासमीमांसा}\label{fol:syn::chap}









\section{परिचय}\label{fol:syn:itx:sec}

प्रथम-क्रम तर्कशास्त्राची उपपत्ती आणि अधिउपपत्ती विकसित करण्यासाठी,
आपण प्रथम त्याच्या व्यंजकांची विन्यासमीमांसा आणि चिन्हार्थमीमांसा
परिभाषित केली पाहिजे. प्रथम-क्रम तर्कशास्त्राची व्यंजके म्हणजे पदे आणि
सूत्रे. पदे चले, स्थिरांक आणि फलने
यांपासून तयार होतात. सूत्रे मात्र विधेये आणि पदे
यांपासून तयार होतात (यांतून सर्वांत लहान, ``आण्विक'' सूत्रे
तयार होतात); मग तार्किक संयोजके आणि संख्यापक वापरून आण्विक
सूत्रांपासून अधिक जटिल सूत्रे तयार करता येतात. रचनानियम
मांडण्याचे अनेक निरनिराळे मार्ग आहेत; आपण त्यांपैकी केवळ एक शक्य
मार्ग देतो. इतर पद्धती वेगळी चिन्हे निवडतील, संयोजकांचे वेगळे संच
आदिम म्हणून निवडतील आणि कंसांचा वेगळ्या रीतीने वापर करतील (किंवा
तथाकथित पोलिश संकेतपद्धतीप्रमाणे कंस अजिबात वापरणार नाहीत). तरी सर्व
दृष्टिकोनांत सामाईक बाब ही की रचनानियम पदांचा आणि सूत्रांचा
संच \emph{विगमनाने} परिभाषित करतात. हे योग्य रीतीने केल्यास,
रचनानियमांनुसार प्रत्येक व्यंजक मूलतः एकाच रीतीने तयार होऊ शकते.
त्यामुळे तयार होणारी व्यंजके \emph{एकमेव रीतीने वाचनीय} असतात. परिणामी,
त्यांची ही विगमनाधारित व्याख्या आपल्याला त्याच पद्धतीने---विगमनाधारित
व्याख्येने---त्या व्यंजकांना अर्थ देता येतो, याची खात्री देते.












\section{प्रथम-क्रम भाषा}\label{fol:syn:fol:sec}


प्रथम-क्रम तर्कशास्त्राची व्यंजके अशा मूलभूत शब्दसंग्रहापासून तयार
केली जातात ज्यात \emph{चले}, \emph{स्थिरांक},
\emph{विधेये} आणि कधीकधी \emph{फलने} असतात. त्यांपासून,
तार्किक संयोजके, संख्यापक आणि कंस व स्वल्पविराम यांसारखी विरामचिन्हे
यांच्या साहाय्याने \emph{पदे} आणि \emph{सूत्रे} तयार केली जातात.

\begin{explain}
अनौपचारिकपणे, विधेये ही गुणधर्मांची आणि संबंधांची नावे,
स्थिरांक ही स्वतंत्र वस्तूंची नावे आणि फलने ही
संगतींची नावे आहेत. एकरूपता~$\eq$ सोडल्यास हीच \emph{अतार्किक
चिन्हे} होत आणि ती मिळून एक भाषा बनते. कोणतीही प्रथम-क्रम
भाषा~$\Lang L$ तिच्या अतार्किक चिन्हांनी निर्धारित होते. सर्वांत
सामान्य बाबतीत $\Lang L$ मध्ये प्रत्येक प्रकारची अनंत चिन्हे असतात.
\end{explain}

सामान्य बाबतीत प्रथम-क्रम तर्कशास्त्रात आपण पुढील चिन्हे वापरतो:

\begin{enumerate}
\item तार्किक चिन्हे
\begin{enumerate}
\item तार्किक संयोजके:
  \startycommalist
  \ycomma $\lnot$ (नकरण)
  \ycomma $\land$ (संधियोग)
  \ycomma $\lor$ (विकल्पयोग)
  \ycomma $\lif$ (शर्त)
  \ycomma $\liff$ (द्विशर्त)
  \ycomma $\lforall$ (वैश्विक संख्यापक)
  \ycomma $\lexists$ (अस्तित्ववाची संख्यापक).
\item असत्यता~$\lfalse$ साठीचे विधानीय स्थिर.
\item सत्यता~$\ltrue$ साठीचे विधानीय स्थिर.
\item द्विस्थानी एकरूपता~$\eq$.
\item चलांचा एक गणनीय अनंत संच: $\Obj v_0$, $\Obj v_1$, $\Obj
  v_2$, \dots
\end{enumerate}
\item प्रथम-क्रम तर्कशास्त्राची \emph{मानक भाषा} बनवणारी अतार्किक चिन्हे
\begin{enumerate}
\item प्रत्येक $n>0$ साठी $n$-स्थानी विधेयांचा एक गणनीय अनंत संच: $\Obj
  A^n_0$, $\Obj A^n_1$, $\Obj A^n_2$, \dots
\item स्थिरांकाचा एक गणनीय अनंत संच: $\Obj c_0$, $\Obj c_1$, $\Obj
  c_2$, \dots.
\item प्रत्येक $n>0$ साठी $n$-स्थानी फलनांचा एक गणनीय अनंत संच:
  $\Obj f^n_0$, $\Obj f^n_1$, $\Obj f^n_2$, \dots
\end{enumerate}
\item विरामचिन्हे: (, ) आणि स्वल्पविराम.
\end{enumerate}

आपल्या बहुतेक व्याख्या आणि निष्कर्ष प्रथम-क्रम तर्कशास्त्राच्या संपूर्ण
मानक भाषेसाठी मांडले जातील. मात्र, उपयोगानुसार आपण भाषा केवळ काही
विधेये, स्थिरांक आणि फलने पुरती मर्यादितही करू शकतो.

\begin{ex}
अंकगणिताची भाषा~$\Lang L_A$ हिच्यात एक द्विस्थानी विधेय~$<$,
एक स्थिरांक~$\Obj 0$, एक एकस्थानी फलन~$\prime$ आणि दोन
द्विस्थानी फलने~$+$ व~$\times$ असतात.
\end{ex}

\begin{ex}
संचसिद्धान्ताची भाषा~$\Lang L_Z$ हिच्यात केवळ एक द्विस्थानी
विधेय~$\in$ असते.
\end{ex}

\begin{ex}
क्रमांची भाषा~$\Lang L_\le$ हिच्यात केवळ द्विस्थानी
विधेय~$\le$ असते.
\end{ex}

पुन्हा, हे संकेत आहेत: औपचारिकदृष्ट्या ही केवळ पर्यायी नावे आहेत;
उदा., $<$, $\in$ आणि $\le$ ही $\Obj A^2_0$ ची, $\Obj 0$ हे
$\Obj c_0$ चे, $\prime$ हे $\Obj f^1_0$ चे, $+$ हे $\Obj f^2_0$ चे आणि $\times$ हे
$\Obj f^2_1$ चे पर्यायी नाव आहे.

.





\begin{intro}
वर आपण वापरलेल्या संज्ञांपेक्षा आणि चिन्हांपेक्षा वेगळ्या संज्ञा व
चिन्हे तुम्हाला परिचित असू शकतात. तर्कशास्त्राच्या ग्रंथांत (आणि
अध्यापनात) ``नकरण'' यासाठी सामान्यतः $\sim$, $\neg$ किंवा~!\ यांपैकी
एखादे, तर ``संधियोग'' यासाठी $\wedge$, $\cdot$ किंवा $\&$ यांपैकी
एखादे चिन्ह वापरतात. ``सशर्त (संकेतार्थक)'' किंवा ``अभिव्यंजन'' यांसाठी
$\rightarrow$, $\Rightarrow$ आणि $\supset$ ही चिन्हे सामान्यतः वापरली जातात.
``द्विसशर्त,'' ``द्वि-अभिव्यंजन'' किंवा
  ``(वास्तविक) सममूल्यता'' यांसाठीची चिन्हे $\leftrightarrow$,
  $\Leftrightarrow$ आणि $\equiv$ ही आहेत.
$\lfalse$ या चिन्हाला निरनिराळ्या ठिकाणी
  ``असत्यता,'' ``फाल्सम,'' ``विसंगती'' किंवा ``तळ'' म्हणतात.
$\ltrue$ या चिन्हाला निरनिराळ्या ठिकाणी
  ``सत्यता,'' ``वेरम'' किंवा ``शिखर'' म्हणतात.

स्थिरांक (ज्यांना कधीकधी नावे म्हणतात) यांच्यासाठी लॅटिन
वर्णमालेच्या सुरुवातीची लघु अक्षरे (उदा., $a$, $b$, $c$) आणि
चले साठी तिच्या शेवटची लघु अक्षरे (उदा., $x$, $y$, $z$)
वापरण्याचा संकेत आहे. संख्यापक चले सोबत जोडले जातात, उदा.,
$x$; वैश्विक संख्यापकासाठी $\forall x$, $(\forall x)$, $(x)$, $\Pi x$,
$\bigwedge_x$ आणि अस्तित्ववाची संख्यापकासाठी $\exists x$, $(\exists
x)$, $(Ex)$, $\Sigma x$, $\bigvee_x$ ही संकेतलेखनातील रूपांतरे आहेत.
\end{intro}

\begin{explain}
आपण सर्व विधानीय कारक आणि दोन्ही संख्यापक यांना भाषेची आदिम चिन्हे
मानू शकतो. त्याऐवजी आदिम चिन्हांचा लहान संच निवडून उरलेले
संकारक परिभाषित मानू शकतो. बूलीय कारकांच्या
``सत्यता-फलनात्मकदृष्ट्या संपूर्ण'' संचांत $\{ \lnot, \lor \}$,
$\{ \lnot, \land \}$ आणि $\{ \lnot,
\lif\}$ यांचा समावेश होतो---यांपैकी प्रत्येक संच कोणत्याही एका
संख्यापकासोबत जोडल्यास अभिव्यक्तीदृष्ट्या संपूर्ण प्रथम-क्रम भाषा मिळते.

इतर दोन संकारक तुम्हाला परिचित असू शकतात: शेफर स्ट्रोक~$|$
(हे नाव हेन्री शेफर यांच्यावरून दिले आहे) आणि पिअर्स बाण~$\downarrow$,
ज्याला क्वाइनचा डॅगर असेही म्हणतात. अनुक्रमे ``नँड'' आणि ``नॉर'' हे
त्यांचे नेहमीचे वाचन दिल्यास, हे कारक स्वतःच सत्यता-फलनात्मकदृष्ट्या
संपूर्ण असतात.
\end{explain}












\section{पदे आणि सूत्र}\label{fol:syn:frm:sec}

प्रथम-क्रम भाषा~$\Lang L$ दिली की, $\Lang L$ च्या मूलभूत
शब्दसंग्रहापासून तयार होणारी व्यंजके आपण परिभाषित करू शकतो. त्यांत
विशेषतः \emph{पदे} आणि \emph{सूत्रे} यांचा समावेश होतो.

\begin{defn}[पदे]
\label{fol:syn:frm:defn:terms}
$\Lang L$ मधील \emph{पदांचा} संच~$\Trm[L]$ पुढीलप्रमाणे विगमनाने
परिभाषित केला जातो:
\begin{enumerate}
\item प्रत्येक चल हे पद आहे.
\item $\Lang L$ मधील प्रत्येक स्थिरांक हे पद आहे.
\item $f$ हे $n$-स्थानी फलन आणि $t_1$, \dots, $t_n$ ही
  पदे असतील, तर $\Atom{f}{t_1, \ldots, t_n}$ हे पद आहे.
\item
  इतर काहीही पद नाही.
\end{enumerate}
ज्या पदात एकही चले नाही ते \emph{बंद पद} होय.
\end{defn}

\begin{explain}
भाषेचे विनिर्देशन करताना स्थिरांक ही पदांपासून वेगळी चिन्हांची
श्रेणी म्हणून दिसतात; पण त्याऐवजी ती शून्यस्थानी फलने मध्ये
समाविष्ट करता आली असती. मग पदांच्या व्याख्येतील दुसरी बाब आवश्यक
राहिली नसती. $n = 0$ असेल तेव्हा $\Atom{f}{t_1, \ldots, t_n}$ याचा
अर्थ केवळ स्वतंत्र $f$ असा घ्यावा लागेल.
\end{explain}

\begin{defn}[सूत्रे]
\label{fol:syn:frm:defn:formulas}
$\Lang L$ या भाषेतील \emph{सूत्रे} चा संच~$\Frm[L]$
पुढीलप्रमाणे विगमनाने परिभाषित केला जातो:
\begin{enumerate}
\item $\lfalse$ हे आण्विक सूत्र आहे.

\item $\ltrue$ हे आण्विक सूत्र आहे.

\item $R$ हे $\Lang L$ मधील $n$-स्थानी विधेय आणि $t_1$, \dots,
  $t_n$ ही $\Lang L$ मधील पदे असतील, तर $\Atom{R}{t_1,\ldots, t_n}$ हे
  आण्विक सूत्र आहे.

\item $t_1$ आणि $t_2$ ही $\Lang L$ मधील पदे असतील, तर $\Atom{\eq}{t_1, t_2}$
  हे आण्विक सूत्र आहे.

\item $\varphi$ हे सूत्र असेल, तर $\lnot \varphi$ हे
  सूत्र आहे.

\item $\varphi$ आणि $\psi$ ही सूत्रे असतील, तर $(\varphi \land
  \psi)$ हे सूत्र आहे.

\item $\varphi$ आणि $\psi$ ही सूत्रे असतील, तर $(\varphi \lor \psi)$
  हे सूत्र आहे.

\item $\varphi$ आणि $\psi$ ही सूत्रे असतील, तर $(\varphi \lif \psi)$
  हे सूत्र आहे.

\item $\varphi$ आणि $\psi$ ही सूत्रे असतील, तर $(\varphi \liff \psi)$
  हे सूत्र आहे.

\item $\varphi$ हे सूत्र आणि $x$ हे चल असेल,
  तर $\lforall[x][\varphi]$ हे सूत्र आहे.

\item $\varphi$ हे सूत्र आणि $x$ हे चल असेल,
  तर $\lexists[x][\varphi]$ हे सूत्र आहे.

\item इतर काहीही सूत्र नाही.
\end{enumerate}
\end{defn}

\begin{explain}
पदांच्या संचाची आणि सूत्रांच्या संचाची व्याख्या या
\emph{विगमनाधारित व्याख्या} आहेत. मूलतः, आपण सूत्रांचा संच
अनंतपणे अनेक टप्प्यांत तयार करतो. आरंभीच्या टप्प्यात सर्व आण्विक सूत्रे
ही सूत्रे आहेत असे घोषित करतो; हे व्याख्येतील पहिल्या काही बाबींशी,
म्हणजे
$\ltrue$,
$\lfalse$,
$\Atom{R}{t_1,\dots,t_n}$ आणि $\Atom{\eq}{t_1,t_2}$ यांच्या बाबींशी जुळते.
म्हणून ``आण्विक सूत्र'' म्हणजे या रूपांपैकी कोणतेही सूत्र.

व्याख्येतील इतर बाबी आधीच तयार केलेल्या सूत्रांपासून नवी
सूत्रे तयार करण्याचे नियम देतात. दुसऱ्या टप्प्यात या नियमांनी
आण्विक सूत्रांपासून सूत्रे तयार करता येतात. तिसऱ्या
टप्प्यात आण्विक सूत्रे आणि दुसऱ्या टप्प्यात मिळालेली सूत्रे यांपासून
नवी सूत्रे तयार करतो, आणि ही प्रक्रिया पुढे चालू राहते. अशा एखाद्या
टप्प्यात शेवटी तयार होणारी प्रत्येक गोष्ट सूत्र असते; इतर
काहीही सूत्र नसते.
\end{explain}

संकेताप्रमाणे, आपण $\eq$ हे चिन्ह त्याच्या युक्तिवादांच्या मध्ये लिहून
कंस वगळतो: $\eq[t_1][t_2]$ हा $\Atom{\eq}{t_1,t_2}$ चा संक्षेप आहे.
तसेच, $\lnot \Atom{\eq}{t_1,t_2}$ चा संक्षेप $\eq/[t_1][t_2]$ असा
लिहितो. द्विस्थानी संयोजक~$\ast$ वापरून $\psi$, $\chi$ पासून तयार केलेले
सूत्र $(\psi \ast \chi)$ लिहिताना आपण अनेकदा सर्वांत बाहेरची कंसांची
जोडी वगळून केवळ~$\psi \ast \chi$ असे लिहू.

\begin{intro}
काही तर्कशास्त्राच्या ग्रंथांत
$\lexists[x][\varphi]$
आणि
$\lforall[x][\varphi]$
यांची गणना
सूत्रे म्हणून करण्यासाठी
चल~$x$ हे~$\varphi$ मध्ये आलेले असणे आवश्यक मानतात. ही अट न
घातल्याने कोणतीही अडचण येत नाही आणि काम अधिक सोपे होते.
\end{intro}



विशिष्ट उपयोगासाठीच्या भाषेत काम करताना आपण अनेकदा द्विस्थानी
विधेये आणि फलने त्यांच्या संबंधित पदांच्या मध्ये
लिहितो; उदा., अंकगणिताच्या भाषेत $t_1 < t_2$ आणि $(t_1 + t_2)$, तर
संचसिद्धान्ताच्या भाषेत $t_1 \in t_2$. अंकगणिताच्या भाषेतील उत्तराधिकारी
फलन तर संकेताप्रमाणे त्याच्या युक्तिवादाच्या \emph{नंतर} लिहिले
जाते:~$t'$. मात्र, औपचारिकदृष्ट्या हे अनुक्रमे
$\Atom{\Obj A^2_0}{t_1, t_2}$, $\Obj f^2_0(t_1, t_2)$,
$\Atom{\Obj A^2_0}{t_1, t_2}$ आणि $f^1_0(t)$ यांचे केवळ रूढ संक्षेप आहेत.

\begin{defn}[विन्यासात्मक अनन्यता]
$\ident$ हे चिन्ह चिन्हमालांमधील विन्यासात्मक अनन्यता व्यक्त करते;
म्हणजे $\varphi \ident \psi$ तेव्हा आणि केवळ तेव्हाच, जेव्हा $\varphi$ आणि $\psi$
या समान लांबीच्या चिन्हमाला असतात आणि प्रत्येक स्थानी त्यांच्यात तेच
चिन्ह असते.
\end{defn}

$\ident$ या चिन्हाच्या दोन्ही बाजूंना जोडणीने मिळालेल्या चिन्हमाला
येऊ शकतात. उदाहरणार्थ, $\varphi \ident (\psi \lor \chi)$ याचा अर्थ असा:
आरंभीचा कंस, चिन्हमाला~$\psi$, $\lor$ हे चिन्ह, चिन्हमाला~$\chi$ आणि
शेवटचा कंस या क्रमाने जोडून मिळणारी चिन्हमाला आणि चिन्हमाला~$\varphi$
या एकच आहेत. असे असेल, तर $\varphi$ चे पहिले चिन्ह आरंभीचा कंस आहे;
$\varphi$ मध्ये दुसऱ्या चिन्हापासून सुरू होणारी $\psi$ ही उपचिन्हमाला आहे;
तिच्यानंतर $\lor$ येते; इत्यादी बाबी आपल्याला माहीत होतात.

पदे आणि सूत्रे ही मूलभूत घटकांपासून विगमनाधारित व्याख्यांनी तयार
होत असल्यामुळे, त्यांच्याविषयीच्या गोष्टी सिद्ध करण्यासाठी पुढील विगमन
तत्त्वे वापरता येतात.

\begin{lem}[\emph{पदांवरील विगमनाचे तत्त्व}]
    \label{fol:syn:frm:lem:trmind}
    $\Lang L$ ही प्रथम-क्रम भाषा असू द्या.
    एखादा गुणधर्म~$P$ असा असेल की

    \begin{enumerate}
        \item तो प्रत्येक चल~$v$ ला लागू होतो,

        \item तो $\Lang L$ मधील प्रत्येक स्थिरांक~$a$ ला लागू होतो, आणि

        \item तो $t_1$,~\dots, $t_n$ यांना लागू होत असेल आणि $f$ हे
        $\Lang L$ मधील $n$-स्थानी फलन असेल, तर तो
        $f(t_1,\dotsc,t_n)$ लाही लागू होतो
    \end{enumerate}
    (येथे $t_1$,~\dots, $t_n$ ही $\Lang{L}$ मधील पदे आहेत असे गृहीत धरले आहे),
    तर $P$ हा $\Trm[L]$ मधील प्रत्येक पदाला लागू होतो.
\end{lem}

\begin{prob}
    \ref{fol:syn:frm:lem:trmind} सिद्ध करा.
\end{prob}

\begin{lem}[\emph{सूत्रे वरील विगमनाचे तत्त्व}]
    \label{fol:syn:frm:thm:frmind}
    $\Lang L$ ही प्रथम-क्रम भाषा असू द्या.
    एखादा गुणधर्म~$P$ सर्व आण्विक सूत्रे ला लागू होत असेल आणि तो असा असेल की

    \begin{enumerate}
        \item तो $\varphi$ ला लागू होत असेल, तर तो $\lnot \varphi$ ला
            लागू होतो;
        \item तो $\varphi$ आणि~$\psi$ ला लागू होत असेल, तर तो
            $(\varphi \land \psi)$ ला लागू होतो;
        \item तो $\varphi$ आणि~$\psi$ ला लागू होत असेल, तर तो
            $(\varphi \lor \psi)$ ला लागू होतो;
        \item तो $\varphi$ आणि~$\psi$ ला लागू होत असेल, तर तो
            $(\varphi \lif \psi)$ ला लागू होतो;
        \item तो $\varphi$ आणि~$\psi$ ला लागू होत असेल, तर तो
            $(\varphi \liff \psi)$ ला लागू होतो;
        \item तो $\varphi$ ला लागू होत असेल, तर तो
            $\lexists[x][\varphi]$ ला लागू होतो;
        \item तो $\varphi$ ला लागू होत असेल, तर तो
            $\lforall[x][\varphi]$ ला लागू होतो;
  \end{enumerate}
  (येथे $\varphi$ आणि $\psi$ ही $\Lang{L}$ मधील सूत्रे आहेत असे गृहीत धरले आहे),
  तर $P$ हा $\Frm[L]$ मधील सर्व सूत्रांना लागू होतो.
\end{lem}












\section{एकमेव वाचनीयता}\label{fol:syn:unq:sec}

\begin{explain}
आपण सूत्रांची ज्या रीतीने व्याख्या केली आहे त्यामुळे प्रत्येक
सूत्राचे \emph{एकमेव वाचन} असते; म्हणजे सूत्रे साठीच्या
आपल्या रचनानियमांनुसार ते तयार करण्याचा मूलतः एकच मार्ग आणि त्याची
``अर्थलावणी'' करण्याचा एकच मार्ग असतो. तसे नसते, तर आपल्याला संदिग्ध
सूत्रे मिळाली असती, म्हणजे एकाहून अधिक वाचने किंवा अर्थलावण्या
असलेली सूत्रे---आणि हे आपल्याला टाळायचे आहे हे स्पष्ट आहे. पण
त्याहून महत्त्वाचे म्हणजे, या गुणधर्माविना आपण पुढे देणार असलेल्या
बहुतेक व्याख्या आणि सिद्धता चालणार नाहीत.

एकमेव वाचनीयतेची खात्री न देणारे सूत्रे तयार करण्याचे सदोष नियम
दिले असते तर काय झाले असते हे पाहणे, हा मुद्दा स्पष्ट करण्याचा बहुधा
सर्वोत्तम मार्ग आहे. उदाहरणार्थ, संयोजकांच्या रचनानियमांत आपण कंस
विसरलो असतो; उदा., पुढील गोष्टीला परवानगी दिली असती:
\begin{quote}
$\varphi$ आणि $\psi$ ही सूत्रे असतील, तर $\varphi \lif \psi$ हेही तसेच आहे.
\end{quote}
$\theta$ या आण्विक सूत्रापासून सुरुवात करून यामुळे $\theta \lif \theta$ तयार करता
आले असते. यापासून, $\theta$ सोबत, $\theta \lif \theta \lif \theta$ मिळाले असते. पण
हे करण्याचे दोन मार्ग आहेत:
\begin{enumerate}
\item $\theta$ हे $\varphi$ आणि $\theta \lif \theta$ हे $\psi$ मानणे.
\item $\varphi$ हे $\theta \lif \theta$ आणि $\psi$ हे~$\theta$ मानणे.
\end{enumerate}
त्याप्रमाणे, सूत्र~$\theta \lif \theta \lif \theta$ ``वाचण्याचे'' दोन मार्ग
आहेत. ते $\psi \lif \chi$ या रूपात असून $\psi$ हे $\theta$ आणि $\chi$ हे $\theta
\lif \theta$ आहे; पण ते $\psi \lif \chi$ या रूपात \emph{देखील आहे}, जेथे $\psi$
हे $\theta \lif \theta$ आणि $\chi$ हे~$\theta$ आहे.

असे झाले, तर आपल्या व्याख्या नेहमी चालणार नाहीत. उदाहरणार्थ, एखाद्या
सूत्राचा मुख्य संकारक परिभाषित करताना आपण म्हणतो: $\psi \lif \chi$
या रूपातील सूत्रात~$\lif$ ची दर्शवलेली आस्थिति हाच मुख्य संकारक
आहे. पण $\theta \lif \theta \lif \theta$ या सूत्राची वर सांगितलेल्या दोन वेगळ्या
मार्गांनी $\psi \lif \chi$ शी जुळणी करता आली, तर एका बाबतीत $\lif$ ची
पहिली आस्थिति मुख्य संकारक मिळते आणि दुसऱ्या बाबतीत दुसरी
आस्थिति मिळते. पण मुख्य संकारक हे सूत्राचे \emph{फलन} असावे
असा आपला उद्देश आहे; म्हणजे प्रत्येक सूत्रामध्ये मुख्य संकारक ची नेमकी एक आस्थिति असली पाहिजे.
\end{explain}

\begin{lem}
सूत्र~$\varphi$ मधील डाव्या आणि उजव्या कंसांची संख्या समान असते.
\end{lem}

\begin{proof}
$\varphi$ ज्या रीतीने तयार होते तिच्यावर विगमन करून हे सिद्ध करू. यासाठी
दोन गोष्टी आवश्यक आहेत: (a) प्रथम, सर्व आण्विक सूत्रांना संबंधित
गुणधर्म लागू होतो हे सिद्ध करावे लागते (विगमनाचा आधार). (b) नंतर,
दिलेल्या सूत्रांपासून नवी सूत्रे तयार केली आणि जुन्या सूत्रांना तो
गुणधर्म लागू होत असेल, तर नव्या सूत्रांनाही तो लागू होतो हे सिद्ध
करावे लागते.

$l(\varphi)$ ही $\varphi$ मधील डाव्या कंसांची संख्या आणि $r(\varphi)$ ही उजव्या
कंसांची संख्या असू द्या; त्याचप्रमाणे $l(t)$ आणि $r(t)$ या पद~$t$
मधील अनुक्रमे डाव्या आणि उजव्या कंसांच्या संख्या असू द्या.

\begin{prob}
    कोणत्याही पद~$t$ साठी $l(t) = r(t)$ हे सिद्ध करा.
\end{prob}

\begin{enumerate}
\item \indcase{\varphi}{\lfalse}{$\indfrm$ मध्ये $0$ डावे आणि $0$
    उजवे कंस आहेत.}

\item \indcase{\varphi}{\ltrue}{$\indfrm$ मध्ये $0$ डावे आणि $0$
    उजवे कंस आहेत.}

\item \indcase{\varphi}{\Atom{R}{t_1,\dots,t_n}}{$l(\indfrm) = 1 + l(t_1) +
  \dots + l(t_n) = 1 + r(t_1) + \dots + r(t_n) = r(\indfrm)$. येथे
  कोणत्याही पद~$t$ साठी $l(t) = r(t)$ ही, सरावासाठी सोडलेली, बाब
  वापरली आहे.}

\item \indcase{\varphi}{\eq[t_1][t_2]}{$l(\indfrm) = l(t_1) + l(t_2) =
  r(t_1) + r(t_2) = r(\indfrm)$.}

\item \indcase{\varphi}{\lnot \psi}{विगमन गृहीतकानुसार,
    $l(\psi) = r(\psi)$. म्हणून $l(\indfrm) = l(\psi) = r(\psi) =
    r(\indfrm)$.}

\item \indcase{\varphi}{(\psi \ast \chi)}{विगमन गृहीतकानुसार, $l(\psi) =
  r(\psi)$ आणि $l(\chi) = r(\chi)$. म्हणून $l(\indfrm) = 1 + l(\psi) + l(\chi) =
  1 + r(\psi) + r(\chi) = r(\indfrm)$.}

\item \indcase{\varphi}{\lforall[x][\psi]}{विगमन
    गृहीतकानुसार, $l(\psi) = r(\psi)$. म्हणून $l(\indfrm) = l(\psi) = r(\psi) =
    r(\indfrm)$.}





\item \indcase{\varphi}{\lexists[x][\psi]}{त्याचप्रमाणे.}
\end{enumerate}
\end{proof}

\begin{defn}[उचित पूर्वखंड]
$\psi$ ही चिन्हमाला आणि एक रिकामी नसलेली चिन्हमाला यांची जोडणी केल्यावर
चिन्हमाला~$\varphi$ मिळत असेल, तर चिन्हमाला $\psi$ ही चिन्हमाला~$\varphi$ चा
\emph{उचित पूर्वखंड} आहे.
\end{defn}

\begin{lem}\label{fol:syn:unq:lem:no-prefix}
$\varphi$ हे सूत्र आणि $\psi$ हा $\varphi$ चा उचित पूर्वखंड असेल, तर $\psi$
हे सूत्र नाही.
\end{lem}

\begin{proof}
सराव.
\end{proof}

\begin{prob}
\ref{fol:syn:unq:lem:no-prefix} सिद्ध करा.
\end{prob}

\begin{prop}
\label{fol:syn:unq:prop:unique-atomic}
$\varphi$ हे आण्विक सूत्र असेल, तर पुढीलपैकी एक आणि फक्त एकच अट त्याला
लागू होते.
\begin{enumerate}
\item $\varphi \ident \lfalse$.
\item $\varphi \ident \ltrue$.
\item $\varphi \ident \Atom{R}{t_1,\dots,t_n}$, जेथे $R$ हे $n$-स्थानी
  विधेय, $t_1$, \dots, $t_n$ ही पदे आणि $R$, $t_1$, \dots,
  $t_n$ यांपैकी प्रत्येक एकमेव रीतीने निर्धारित आहे.
\item $\varphi \ident \eq[t_1][t_2]$, जेथे $t_1$ आणि $t_2$ ही एकमेव रीतीने
  निर्धारित पदे आहेत.
\end{enumerate}
\end{prop}

\begin{proof}
सराव.
\end{proof}

\begin{prob}
\ref{fol:syn:unq:prop:unique-atomic} सिद्ध करा (सूचना: पदांसाठी
\ref{fol:syn:unq:lem:no-prefix} ची एक आवृत्ती मांडा आणि सिद्ध करा.)
\end{prob}

\begin{prop}[एकमेव वाचनीयता]
प्रत्येक सूत्राला पुढीलपैकी एक आणि फक्त एकच अट लागू होते.
\begin{enumerate}
\item $\varphi$ हे आण्विक आहे.

\item $\varphi$ हे $\lnot \psi$ या रूपात आहे.

\item $\varphi$ हे $(\psi \land \chi)$ या रूपात आहे.

\item $\varphi$ हे $(\psi \lor \chi)$ या रूपात आहे.

\item $\varphi$ हे $(\psi \lif \chi)$ या रूपात आहे.

\item $\varphi$ हे $(\psi \liff \chi)$ या रूपात आहे.

\item $\varphi$ हे $\lforall[x][\psi]$ या रूपात आहे.

\item $\varphi$ हे $\lexists[x][\psi]$ या रूपात आहे.
\end{enumerate}
शिवाय, प्रत्येक बाबतीत $\psi$, किंवा $\psi$ आणि $\chi$, एकमेव रीतीने
निर्धारित असतात. म्हणजे, उदाहरणार्थ, $\psi$, $\chi$ आणि $\psi'$, $\chi'$ अशा
दोन वेगळ्या जोड्या अस्तित्वात नाहीत की $\varphi$ हे एकाच वेळी
$(\psi \lif \chi)$ आणि $(\psi' \lif \chi')$ या दोन्ही रूपांत असेल.
\end{prop}

\begin{proof}
रचनानियमांनुसार, सूत्र आण्विक नसेल, तर त्याची सुरुवात आरंभीच्या
कंसाने~(, $\lnot$ ने, किंवा एखाद्या संख्यापकाने झाली
पाहिजे. दुसरीकडे, पुढीलपैकी एखाद्या चिन्हाने सुरू होणारे प्रत्येक
सूत्र आण्विक असले पाहिजे: विधेय, फलन,
स्थिरांक, $\lfalse$, $\ltrue$.

म्हणून आपल्याला प्रत्यक्षात एवढेच दाखवायचे आहे: $\varphi$ हे $(\psi \ast
\chi)$ या रूपात आणि $(\psi' \mathbin{\ast'} \chi')$ या रूपातही असेल, तर
$\psi \ident \psi'$, $\chi \ident \chi'$ आणि $\ast = {\ast'}$.

म्हणून $\varphi \ident (\psi \ast \chi)$ आणि $\varphi \ident (\psi'
\mathbin{\ast'} \chi')$ हे दोन्ही गृहीत धरा. मग एकतर $\psi \ident \psi'$
किंवा तसे नाही. तसे असेल, तर स्पष्टपणे $\ast = {\ast'}$ आणि $\chi
\ident \chi'$, कारण त्या स्थितीत त्या $\varphi$ च्या एकाच स्थानी सुरू होणाऱ्या
आणि समान लांबीच्या उपचिन्हमाला आहेत. उरलेली बाब म्हणजे $\psi
\not\ident \psi'$. $\psi$ आणि $\psi'$ या दोन्ही $\varphi$ च्या एकाच स्थानी सुरू
होणाऱ्या उपचिन्हमाला असल्यामुळे, त्यांपैकी एक दुसरीचा उचित पूर्वखंड
असला पाहिजे. पण \ref{fol:syn:unq:lem:no-prefix} नुसार हे अशक्य आहे.
\end{proof}













\section{सूत्राचा मुख्य संकारक}\label{fol:syn:mai:sec}

\begin{explain}
सूत्र~$\varphi$ तयार करताना सर्वांत शेवटी वापरलेल्या कारकाविषयी
बोलणे अनेकदा उपयुक्त असते. या कारकाला $\varphi$ चा \emph{मुख्य कारक}
म्हणतात. अंतःप्रेरणेने तो $\varphi$ चा ``सर्वांत बाहेरचा'' कारक होय.
उदाहरणार्थ, $\lnot \varphi$ चा मुख्य कारक $\lnot$, $(\varphi \lor \psi)$ चा
मुख्य कारक $\lor$, इत्यादी.
\end{explain}


\begin{defn}[मुख्य संकारक]
\label{fol:syn:mai:def:main-op}
सूत्र~$\varphi$ चा \emph{मुख्य संकारक} पुढीलप्रमाणे परिभाषित केला जातो:
\begin{enumerate}
\item \indcase*{\varphi}{\varphi}{$\indfrm$ ला मुख्य संकारक नाही.}

\item \indcase{\varphi}{\lnot \psi}{$\indfrm$ चा मुख्य संकारक
  हा~$\lnot$ आहे.}

\item \indcase{\varphi}{(\psi \land \chi)}{$\indfrm$ चा
  मुख्य संकारक हा~$\land$ आहे.}

\item \indcase{\varphi}{(\psi \lor \chi)}{$\indfrm$ चा
  मुख्य संकारक हा~$\lor$ आहे.}

\item \indcase{\varphi}{(\psi \lif \chi)}{$\indfrm$ चा
  मुख्य संकारक हा~$\lif$ आहे.}

\item \indcase{\varphi}{(\psi \liff \chi)}{$\indfrm$ चा
  मुख्य संकारक हा~$\liff$ आहे.}

\item \indcase{\varphi}{\lforall[x][\psi]}{$\indfrm$ चा
  मुख्य संकारक हा~$\lforall$ आहे.}

\item \indcase{\varphi}{\lexists[x][\psi]}{$\indfrm$ चा
  मुख्य संकारक हा~$\lexists$ आहे.}
\end{enumerate}
\end{defn}

प्रत्येक बाबतीत सूत्रातील मुख्य संकारक ची नेमकी दर्शवलेली
\emph{आस्थिति} आपल्याला अभिप्रेत आहे. उदाहरणार्थ,
$((\theta \lif \varphi_{E}) \lif (\varphi_{E} \lif \theta))$ हे सूत्र $(\psi \lif \chi)$ या रूपात
असून $\psi$ हे $(\theta \lif \varphi_{E})$ आणि $\chi$ हे $(\varphi_{E} \lif \theta)$ असल्यामुळे,
$\lif$ ची दुसरी आस्थिति मुख्य संकारक आहे.

\begin{explain}
ही अशा फलनाची \emph{पुनरावर्ती} व्याख्या आहे जे सर्व अनाण्विक
सूत्रे ना त्यांच्या मुख्य संकारक च्या आस्थितीकडे पाठवते.
सूत्रांची व्याख्या विगमनाने केलेली असल्यामुळे प्रत्येक
सूत्र~$\varphi$ ला \ref{fol:syn:mai:def:main-op} मधील एखादी बाब लागू होते.
त्यामुळे प्रत्येक अनाण्विक सूत्र~$\varphi$ साठी मुख्य संकारक
अस्तित्वात आहे याची खात्री मिळते. प्रत्येक सूत्राला यांपैकी
फक्त एकच अट लागू होते आणि प्रत्येक बाबतीत $\varphi$ ज्यांपासून तयार केले
आहे ती लहान सूत्रे एकमेव रीतीने निर्धारित असतात. त्यामुळे $\varphi$
मधील मुख्य संकारक ची आस्थिति एकमेव असते आणि म्हणून आपण फलन
परिभाषित केले आहे.
\end{explain}

सूत्रांचा मुख्य संकारक कोणते चिन्ह आहे यानुसार आपण त्यांना
\ref{fol:syn:mai:tab:main-op} मधील नावे देतो.

\begin{table}[!h]
\centering
\begin{tabular}{c | c | c}
मुख्य संकारक & सूत्राचा प्रकार & उदाहरण\\
\hline
काही नाही & आण्विक (सूत्र) &
$\lfalse$,
$\ltrue$,
$\Atom{R}{t_1, \dots, t_n}$,
$\eq[t_1][t_2]$\\
$\lnot$ & नकरण & $\lnot \varphi$ \\
$\land$ & संधियोग & $(\varphi \land \psi$) \\
$\lor$ & विकल्पयोग & $(\varphi \lor \psi$) \\
$\lif$ & शर्त & $(\varphi \lif \psi$) \\
$\liff$ & द्विशर्त & $(\varphi \liff \psi)$ \\
$\lforall[][]$ & वैश्विक (सूत्र)& $\lforall[x][\varphi]$ \\
$\lexists[][]$ & अस्तित्ववाची (सूत्र)& $\lexists[x][\varphi]$
\end{tabular}
\caption{सूत्रांचे मुख्य कारक आणि नावे}
\label{fol:syn:mai:tab:main-op}
\end{table}












\section{उपसूत्र}\label{fol:syn:sbf:sec}

\begin{explain}
दिलेल्या सूत्राचे ``घटक'' असलेल्या सूत्रे विषयी बोलणे अनेकदा
उपयुक्त असते. त्यांना आपण त्याची \emph{उपसूत्रे} म्हणतो. कोणतेही
सूत्र हे स्वतःचे उपसूत्र मानले जाते; $\varphi$ स्वतः सोडून
$\varphi$ चे इतर उपसूत्र म्हणजे \emph{उचित उपसूत्र} होय.
\end{explain}

\begin{defn}[तात्काळ उपसूत्र]
$\varphi$ हे सूत्र असेल, तर $\varphi$ ची \emph{तात्काळ उपसूत्रे}
पुढीलप्रमाणे विगमनाने परिभाषित केली जातात:
\begin{enumerate}
\item आण्विक सूत्रे ना तात्काळ उपसूत्रे नसतात.

\item \indcase{\varphi}{\lnot \psi}{$\indfrm$ चे एकमेव तात्काळ
    उपसूत्र हे~$\psi$ आहे.}

\item \indcase{\varphi}{(\psi \ast \chi)}{$\indfrm$ ची तात्काळ उपसूत्रे
  $\psi$ आणि $\chi$ आहेत ($\ast$ हे कोणतेही द्विस्थानी संयोजक आहे).}

\item \indcase{\varphi}{\lforall[x][\psi]}{$\indfrm$ चे एकमेव तात्काळ
    उपसूत्र हे~$\psi$ आहे.}

\item \indcase{\varphi}{\lexists[x][\psi]}{$\indfrm$ चे एकमेव तात्काळ
    उपसूत्र हे~$\psi$ आहे.}
\end{enumerate}
\end{defn}

\begin{defn}[उचित उपसूत्र]
$\varphi$ हे सूत्र असेल, तर $\varphi$ ची \emph{उचित उपसूत्रे}
पुढीलप्रमाणे पुनरावर्ती रीतीने परिभाषित केली जातात:
\begin{enumerate}
\item आण्विक सूत्रे ना उचित उपसूत्रे नसतात.

\item \indcase{\varphi}{\lnot \psi}{$\indfrm$ ची उचित
    उपसूत्रे म्हणजे~$\psi$ आणि $\psi$ ची सर्व उचित उपसूत्रे.}

\item \indcase{\varphi}{(\psi \ast \chi)}{$\indfrm$ ची उचित उपसूत्रे
  म्हणजे $\psi$, $\chi$, $\psi$ ची सर्व उचित उपसूत्रे आणि~$\chi$ ची
  सर्व उचित उपसूत्रे.}

\item \indcase{\varphi}{\lforall[x][\psi]}{$\indfrm$ ची उचित
    उपसूत्रे म्हणजे~$\psi$ आणि $\psi$ ची सर्व उचित
    उपसूत्रे.}

\item \indcase{\varphi}{\lexists[x][\psi]}{$\indfrm$ ची उचित
    उपसूत्रे म्हणजे~$\psi$ आणि $\psi$ ची सर्व उचित
    उपसूत्रे.}
\end{enumerate}
\end{defn}

\begin{defn}[उपसूत्र]
$\varphi$ ची उपसूत्रे म्हणजे स्वतः $\varphi$ आणि त्याची सर्व उचित
उपसूत्रे.
\end{defn}

\begin{explain}
तात्काळ उपसूत्रे आणि उचित उपसूत्रे यांच्या व्याख्यांतील
सूक्ष्म फरक लक्षात घ्या. पहिल्या बाबतीत, $\varphi$ च्या प्रत्येक शक्य
रूपासाठी आपण सूत्र~$\varphi$ ची तात्काळ उपसूत्रे थेट परिभाषित केली
आहेत. ही बाबींनुसार केलेली स्पष्ट व्याख्या आहे आणि त्या बाबी
सूत्रांच्या संचाच्या विगमनाधारित व्याख्येला समांतर आहेत. दुसऱ्या
बाबतीतही सर्व सूत्रांचा संच ज्या रीतीने परिभाषित केला आहे तिला
आपण समांतर गेलो आहोत; पण प्रत्येक बाबतीत लहान सूत्रे $\psi$, $\chi$
यांची उचित उपसूत्रे या सूत्रे ना स्वतःलाही समाविष्ट करून
घेतली आहेत.
त्यामुळे व्याख्या \emph{पुनरावर्ती} होते. सर्वसाधारणपणे, विगमनाने
परिभाषित केलेल्या संचावरील फलनाची व्याख्या (आपल्या बाबतीत हा संच म्हणजे
सूत्रे) पुनरावर्ती असते, जर फलनाच्या व्याख्येतील बाबी स्वतः त्या
फलनाचा वापर करत असतील. मात्र ती सुनिर्धारित असण्यासाठी, आपण परिभाषित
करत असलेल्या युक्तिवादाच्या ``आधी'' येणाऱ्या युक्तिवादांसाठीचीच फलनाची
मूल्ये वापरतो याची खात्री केली पाहिजे---आपल्या बाबतीत $(\psi \ast \chi)$
ची ``उचित उपसूत्र'' परिभाषित करताना आपण केवळ ``आधीच्या''
सूत्रे $\psi$ आणि $\chi$ यांची उचित उपसूत्रे वापरतो.
\end{explain}

\begin{prop}
\label{fol:syn:sbf:prop:subfrm-trans}
$\psi$ हे $\varphi$ चे उपसूत्र आणि $\chi$ हे $\psi$ चे उपसूत्र आहे असे गृहीत धरा.
मग $\chi$ हे $\varphi$ चे उपसूत्र आहे. दुसऱ्या शब्दांत, उपसूत्र-संबंध
संक्रमक आहे.
\end{prop}

\begin{prob}
\ref{fol:syn:sbf:prop:subfrm-trans} सिद्ध करा.
\end{prob}

\begin{prop}
\label{fol:syn:sbf:prop:count-subfrms}
$\varphi$ हे $n$ संयोजके आणि संख्यापक असलेले सूत्र आहे असे गृहीत धरा.
मग $\varphi$ ची जास्तीत जास्त $2n+1$ उपसूत्रे आहेत.
\end{prop}

\begin{prob}
\ref{fol:syn:sbf:prop:count-subfrms} सिद्ध करा.
\end{prob}












\section{रचनाक्रमिका}\label{fol:syn:fseq:sec}

सूत्रांची विगमनाधारित व्याख्या देणे आणि त्याला पूरक तंत्र म्हणून
सूत्रांचे गुणधर्म विगमनाने सिद्ध करणे हा सुबक व कार्यक्षम
दृष्टिकोन आहे. तथापि, सूत्रांच्या रचनेचा तळापासून वर जाणारा,
टप्प्याटप्प्याचा दृष्टिकोनही उपयुक्त ठरू शकतो; येथे आपण तो
\emph{रचनाक्रमिका} या संकल्पनेच्या साहाय्याने मांडतो.

पदे आणि सूत्रे त्यांच्या विगमनाधारित व्याख्यांचा निर्देश न करता
या रीतीने कशी मांडता येतात हे दाखवण्यासाठी, आपण प्रथम एखाद्या
भाषा~$\Lang L$ मधील चिन्हांपासून बनलेल्या स्वेच्छ चिन्हमालेची संकल्पना
मांडतो.

\begin{defn}[चिन्हमाला]
\label{fol:syn:fseq:defn:string}
$\Lang L$ ही प्रथम-क्रम भाषा आहे असे समजा. \emph{$\Lang L$-चिन्हमाला}
म्हणजे $\Lang L$ च्या चिन्हांची सांत क्रमिका. संदर्भावरून भाषा~$\Lang L$
स्पष्टपणे ठरलेली असेल, तर $\Lang L$-चिन्हमालेला आपण अनेकदा फक्त
\emph{चिन्हमाला} म्हणू.
\end{defn}

\begin{ex}
कोणत्याही प्रथम-क्रम भाषा $\Lang L$ साठी सर्व $\Lang L$-सूत्रे
या $\Lang L$-चिन्हमाला असतात; पण याचे उलट खरे नसते. उदाहरणार्थ,
\[)(\Obj v_0\lif\lexists\] ही $\Lang L$-चिन्हमाला आहे, पण
$\Lang L$-सूत्र नाही.
\end{ex}

\begin{defn}[पदांसाठी रचनाक्रमिका]
\label{fol:syn:fseq:defn:fseq-trm}
$\Lang L$-चिन्हमालांची सांत क्रमिका $\tuple{t_0,\dotsc,t_n}$ ही पद~$t$
साठी \emph{रचनाक्रमिका} असते, जर $t \ident t_n$ आणि प्रत्येक $i \leq n$
साठी एकतर $t_i$ हे चल किंवा स्थिरांक असेल, अथवा $\Lang
L$ मध्ये एखादे $k$-स्थानी फलन~$f$ असेल आणि असे
$m_1,\dotsc,m_k < i$ अस्तित्वात असतील की
$t_i \ident f(t_{m_1},\dotsc,t_{m_k})$. फरक स्पष्ट करणे आवश्यक असेल,
तर पदांसाठीच्या रचनाक्रमिकांना आपण \emph{पद-रचनाक्रमिका} म्हणू.

\end{defn}

\begin{ex}
क्रमिका
\[    \tuple{\Obj c_0, \Obj v_0, \Atom{\Obj f^2_0}{\Obj c_0, \Obj v_0}, \Atom{\Obj f^1_0}{\Atom{\Obj f^2_0}{\Obj c_0, \Obj v_0}}}
\]
ही पद $\Atom{\Obj f^1_0}{\Atom{\Obj
f^2_0}{\Obj c_0, \Obj v_0}}$ साठी रचनाक्रमिका आहे; त्याचप्रमाणे
\[    \tuple{\Obj v_0, \Obj c_0, \Atom{\Obj f^2_0}{\Obj c_0, \Obj v_0}, \Atom{\Obj f^1_0}{\Atom{\Obj f^2_0}{\Obj c_0, \Obj v_0}}}.
\]
हीदेखील त्याची रचनाक्रमिका आहे.
\end{ex}

\begin{defn}[सूत्रांसाठी रचनाक्रमिका]
\label{fol:syn:fseq:defn:fseq-frm}
$\Lang L$-चिन्हमालांची सांत क्रमिका $\tuple{\varphi_0,\dotsc,\varphi_n}$ ही
$\varphi$ साठी \emph{रचनाक्रमिका} असते, जर $\varphi \ident \varphi_n$ आणि प्रत्येक
$i \leq n$ साठी एकतर $\varphi_i$ हे आण्विक सूत्र असेल, अथवा पुढीलपैकी
एखादी अट पूर्ण करणारे $j,k < i$ आणि एखादे चल~$x$ अस्तित्वात
असतील:
\begin{enumerate}
    \item $\varphi_i \ident \lnot \varphi_j$.
    \item $\varphi_i \ident (\varphi_j \land \varphi_k)$.
    \item $\varphi_i \ident (\varphi_j \lor \varphi_k)$.
    \item $\varphi_i \ident (\varphi_j \lif \varphi_k)$.
    \item $\varphi_i \ident (\varphi_j \liff \varphi_k)$.
    \item $\varphi_i \ident \lforall[x][\varphi_j]$.
    \item $\varphi_i \ident \lexists[x][\varphi_j]$.
\end{enumerate}
फरक स्पष्ट करणे आवश्यक असेल, तर सूत्रांसाठीच्या रचनाक्रमिकांना आपण
\emph{सूत्र-रचनाक्रमिका} म्हणू.
\end{defn}

\begin{ex}
\[\tuple{
    \Atom{\Obj A^1_0}{\Obj v_0},
    \Atom{\Obj A^1_1}{\Obj c_1},
    (\Atom{\Obj A^1_1}{\Obj c_1} \land \Atom{\Obj A^1_0}{\Obj v_0}),
    \lexists[\Obj v_0][(\Atom{\Obj A^1_1}{\Obj c_1} \land \Atom{\Obj A^1_0}{\Obj v_0})]
}
\]
ही $\lexists[\Obj v_0][(\Atom{\Obj A^1_1}{\Obj c_1} \land \Atom{\Obj
A^1_0}{\Obj v_0})]$ ची एक रचनाक्रमिका आहे; त्याचप्रमाणे
\begin{multline*}
\tuple{
    \Atom{\Obj A^1_0}{\Obj v_0},
    \Atom{\Obj A^1_1}{\Obj c_1},
    (\Atom{\Obj A^1_1}{\Obj c_1} \land \Atom{\Obj A^1_0}{\Obj v_0}),
    \Atom{\Obj A^1_1}{\Obj c_1},\\
    \lforall[\Obj v_1][\Atom{\Obj A^1_0}{\Obj v_0}],
    \lexists[\Obj v_0][(\Atom{\Obj A^1_1}{\Obj c_1} \land \Atom{\Obj A^1_0}{\Obj v_0})]
}.
\end{multline*}

दुसऱ्या उदाहरणावरून दिसते की रचनाक्रमिकेत ``अडगळ'' असू शकते:
अनावश्यक असलेली किंवा रचनेत काहीही हातभार न लावणारी सूत्रे.
\end{ex}

\begin{prop}\label{fol:syn:fseq:prop:formed}
$\Frm[L]$ मधील प्रत्येक सूत्र~$\varphi$ ला एक रचनाक्रमिका असते.
\end{prop}

\begin{proof}
$\varphi$ आण्विक आहे असे समजा. मग क्रमिका~$\tuple{\varphi}$ ही $\varphi$ साठी
रचनाक्रमिका आहे.

आता $\psi$ आणि~$\chi$ यांना अनुक्रमे $\tuple{\psi_0,\dotsc,\psi_n}$ आणि
$\tuple{\chi_0,\dotsc,\chi_m}$ या रचनाक्रमिका आहेत असे समजा.

\begin{enumerate}
  \item जर $\varphi \ident \lnot \psi$, तर
      $\tuple{\psi_0,\dotsc,\psi_n,\lnot \psi_n}$ ही $\varphi$ साठी
      रचनाक्रमिका आहे.
  \item जर $\varphi \ident (\psi \land \chi)$, तर
      $\tuple{\psi_0,\dotsc,\psi_n,\chi_0,\dotsc,\chi_m,(\psi_n \land \chi_m)}$
      ही $\varphi$ साठी रचनाक्रमिका आहे.
  \item जर $\varphi \ident (\psi \lor \chi)$, तर
      $\tuple{\psi_0,\dotsc,\psi_n,\chi_0,\dotsc,\chi_m,(\psi_n \lor \chi_m)}$
      ही $\varphi$ साठी रचनाक्रमिका आहे.
  \item जर $\varphi \ident (\psi \lif \chi)$, तर
      $\tuple{\psi_0,\dotsc,\psi_n,\chi_0,\dotsc,\chi_m,(\psi_n \lif \chi_m)}$
      ही $\varphi$ साठी रचनाक्रमिका आहे.
  \item जर $\varphi \ident (\psi \liff \chi)$, तर
      $\tuple{\psi_0,\dotsc,\psi_n,\chi_0,\dotsc,\chi_m,(\psi_n \liff \chi_m)}$
      ही $\varphi$ साठी रचनाक्रमिका आहे.
  \item जर $\varphi \ident \lforall[x][\psi]$, तर
      $\tuple{\psi_0,\dotsc,\psi_n,\lforall[x][\psi_n]}$
      ही $\varphi$ साठी रचनाक्रमिका आहे.
  \item जर $\varphi \ident \lexists[x][\psi]$, तर
      $\tuple{\psi_0,\dotsc,\psi_n,\lexists[x][\psi_n]}$
      ही $\varphi$ साठी रचनाक्रमिका आहे.
\end{enumerate}
सूत्रे वरील विगमनाच्या तत्त्वानुसार प्रत्येक सूत्राला एक
रचनाक्रमिका असते.
\end{proof}

याचे उलट विधानही आपण सिद्ध करू शकतो. ते महत्त्वाचे आहे, कारण त्यामुळे
सूत्रे परिभाषित करण्याच्या आपल्या दोन्ही पद्धती सममूल्य आहेत, म्हणजे
त्या समान परिणाम देतात, हे दिसते. तसेच रचनाक्रमिकांच्या लांबीवर साधे
विगमन वापरून सूत्रांविषयीची प्रमेये सिद्ध करता येतात.

\begin{lem}
\label{fol:syn:fseq:lem:fseq-init}
$\tuple{\varphi_0,\dotsc,\varphi_n}$ ही $\varphi_n$ साठी रचनाक्रमिका आहे आणि
$k \leq n$ असे समजा. मग $\tuple{\varphi_0,\dotsc,\varphi_k}$ ही $\varphi_k$ साठी
रचनाक्रमिका आहे.
\end{lem}

\begin{proof}
अभ्यासप्रश्न.
\end{proof}

\begin{prob}
\ref{fol:syn:fseq:lem:fseq-init} सिद्ध करा.
\end{prob}

\begin{thm}
\label{fol:syn:fseq:thm:fseq-frm-equiv}
$\Frm[L]$ हा अशा सर्व $\Lang L$-चिन्हमाला~$\varphi$ चा संच आहे की $\varphi$
साठी एखादी सूत्र-रचनाक्रमिका अस्तित्वात आहे.
\end{thm}

\begin{proof}
भाषा~$\Lang L$ मधील ज्या सर्व चिन्हमालांना रचनाक्रमिका आहे, त्यांचा
संच~$F$ असू द्या. $\Frm[L] \subseteq F$ हे
\ref{fol:syn:fseq:prop:formed} मध्ये आपण पाहिले आहे; म्हणून आता उलट
समावेश सिद्ध करू.

$\varphi$ ला $\tuple{\varphi_0,\dotsc,\varphi_n}$ ही रचनाक्रमिका आहे असे समजा.
$n$ वरील प्रबल विगमनाने $\varphi \in \Frm[L]$ हे सिद्ध करू. लांबी
$m < n$ असलेली रचनाक्रमिका ज्या प्रत्येक चिन्हमालेला आहे, ती
$\Frm[L]$ मध्ये आहे, हे आपले विगमन गृहीतक आहे. रचनाक्रमिकेच्या
व्याख्येनुसार एकतर $\varphi \ident \varphi_n$ आण्विक आहे किंवा पुढीलपैकी एखादी
अट पूर्ण करणारे $j,k < n$ अस्तित्वात असले पाहिजेत:
\begin{enumerate}
    \item $\varphi \ident \lnot \varphi_j$.
    \item $\varphi \ident (\varphi_j \land \varphi_k)$.
    \item $\varphi \ident (\varphi_j \lor \varphi_k)$.
    \item $\varphi \ident (\varphi_j \lif \varphi_k)$.
    \item $\varphi \ident (\varphi_j \liff \varphi_k)$.
    \item $\varphi \ident \lforall[x][\varphi_j]$.
    \item $\varphi \ident \lexists[x][\varphi_j]$.
\end{enumerate}
आता प्रत्येक शक्यतेचा स्वतंत्र विचार करू. $\varphi$ आण्विक असेल, तर
$\varphi_n \in \Frm[L]$. त्याऐवजी $\varphi \ident (\varphi_j \land \varphi_k)$ असे समजा.


\ref{fol:syn:fseq:lem:fseq-init} नुसार
$\tuple{\varphi_0,\dotsc,\varphi_j}$ आणि $\tuple{\varphi_0,\dotsc,\varphi_k}$ या अनुक्रमे
$\varphi_j$ आणि~$\varphi_k$ साठी रचनाक्रमिका आहेत. या $\varphi$ च्या
रचनाक्रमिकेच्या उचित आरंभीच्या उपक्रमिका असल्यामुळे दोन्हींची लांबी
$n$ पेक्षा कमी आहे. म्हणून विगमन गृहीतकानुसार $\varphi_j$ आणि~$\varphi_k$ ही
$\Frm[L]$ मध्ये आहेत; त्यामुळे सूत्राच्या व्याख्येनुसार
$(\varphi_j \land \varphi_k)$ हेही त्यात आहे. इतर शक्यता समांतर युक्तिवादाने
सिद्ध होतात.
\end{proof}

पदांसाठीच्या रचनाक्रमिकांना सूत्रांच्या रचनाक्रमिकांसारखेच गुणधर्म
असतात.

\begin{prop}
\label{fol:syn:fseq:prop:fseq-trm-equiv}
$\Trm[L]$ हा अशा सर्व $\Lang L$-चिन्हमाला $t$ चा संच आहे की $t$
साठी एखादी पद-रचनाक्रमिका अस्तित्वात आहे.
\end{prop}

\begin{proof}
अभ्यासप्रश्न.
\end{proof}

\begin{prob}
\ref{fol:syn:fseq:prop:fseq-trm-equiv} सिद्ध करा.
सूचना: \ref{fol:syn:fseq:thm:fseq-frm-equiv} च्या सिद्धतेत वापरलेली
रणनीती वापरा.
\end{prob}

रचनाक्रमिकेत दोन प्रकारची ``अडगळ'' येऊ शकते: पुनरावृत्त घटक आणि सूत्र
किंवा पद यांच्या रचनेशी असंबद्ध घटक. न्यूनतम रचनाक्रमिका विचारात घेऊन
आपण दोन्ही प्रकार दूर करू शकतो.

\begin{defn}[न्यूनतम रचनाक्रमिका]
\label{fol:syn:fseq:defn:minimal-fseq}
सूत्र~$\varphi$ साठीची रचनाक्रमिका $\tuple{\varphi_0, \ldots, \varphi_n}$ ही
$\varphi$ साठी \emph{न्यूनतम रचनाक्रमिका} असते, जर $\varphi$ साठीची प्रत्येक
दुसरी रचनाक्रमिका~$s$ घेतली असता~$s$ ची लांबी $n+1$ पेक्षा अधिक किंवा
तिच्याइतकी असेल.

त्याचप्रमाणे, पद~$t$ साठीची रचनाक्रमिका $\tuple{t_0, \ldots, t_n}$ ही
$t$ साठी \emph{न्यूनतम रचनाक्रमिका} असते, जर $t$ साठीची प्रत्येक
दुसरी रचनाक्रमिका~$s$ घेतली असता~$s$ ची लांबी $n+1$ पेक्षा अधिक किंवा
तिच्याइतकी असेल.
\end{defn}

एका सूत्राला किंवा पदाला एकापेक्षा अधिक न्यूनतम रचनाक्रमिका असू शकतात;
पण त्यांत नेमक्या त्याच चिन्हमाला असतील, हे लक्षात घ्या.

\begin{prop}
\label{fol:syn:fseq:prop:subformula-equivs}
पुढील विधाने सममूल्य आहेत:
\begin{enumerate}
    \item $\psi$ हे $\varphi$ चे उप-सूत्र आहे.
    \item $\psi$ हे $\varphi$ च्या प्रत्येक रचनाक्रमिकेत येते.
    \item $\psi$ हे $\varphi$ च्या एखाद्या न्यूनतम रचनाक्रमिकेत येते.
\end{enumerate}
\end{prop}

\begin{proof}
अभ्यासप्रश्न.
\end{proof}

\begin{prob}
\ref{fol:syn:fseq:prop:subformula-equivs} सिद्ध करा.
\end{prob}

\begin{history}
रेमंड स्मलियन यांनी त्यांच्या \emph{First-Order Logic}
(Smullyan, 1968) या पाठ्यपुस्तकात रचनाक्रमिका प्रथम मांडल्या.
रचनाक्रमिकांचे आणखी गुणधर्म Zuckerman (1973) यांनी प्रस्थापित केले.
\end{history}












\section{मुक्त चल आणि वाक्य}\label{fol:syn:fvs:sec}

\begin{defn}[चल च्या मुक्त आस्थिती]
\label{fol:syn:fvs:defn:free-occ}
सूत्र मधील चल च्या \emph{मुक्त} आस्थिती पुढीलप्रमाणे
विगमनाने परिभाषित केल्या जातात:
\begin{enumerate}
\item \indcase*{\varphi}{$\varphi$ आण्विक आहे}{$\indfrm$ मधील सर्व
  चल-आस्थिती मुक्त असतात.}

\item \indcase{\varphi}{\lnot \psi}{$\indfrm$ मधील मुक्त
  चल-आस्थिती म्हणजे नेमक्या $\psi$ मधील मुक्त आस्थिती.}

\item \indcase{\varphi}{(\psi \ast \chi)}{$\indfrm$ मधील मुक्त
  चल-आस्थिती म्हणजे $\psi$ मधील आस्थिती आणि त्यांच्यासोबत
  $\chi$ मधील आस्थिती.}

\item \indcase{\varphi}{\lforall[x][\psi]}{$\indfrm$ मधील मुक्त
  चल-आस्थिती म्हणजे $\psi$ मधील $x$ च्या आस्थिती सोडून इतर
  सर्व मुक्त आस्थिती.}

\item \indcase{\varphi}{\lexists[x][\psi]}{$\indfrm$ मधील मुक्त
  चल-आस्थिती म्हणजे $\psi$ मधील $x$ च्या आस्थिती सोडून इतर
  सर्व मुक्त आस्थिती.}
\end{enumerate}
\end{defn}

\begin{defn}[बद्ध चरे]
सूत्र~$\varphi$ मधील चल ची आस्थिती मुक्त नसेल, तर ती
\emph{बद्ध} असते.
\end{defn}

\begin{prob}
\ref{fol:syn:fvs:defn:free-occ} च्या धर्तीवर बद्ध चर-आस्थितींची
विगमनाधारित व्याख्या द्या.
\end{prob}

\begin{defn}[व्याप्ती]
सूत्र~$\varphi$ मध्ये $\lforall[x][\psi]$ ची एखाद्या उपसूत्राची
  आस्थिती असेल, तर $\varphi$ मधील $\psi$ च्या संबंधित आस्थितीला
  $\lforall[x]$ च्या संबंधित आस्थितीची \emph{व्याप्ती} म्हणतात.
  $\lexists[x]$ साठीही त्याचप्रमाणे.

$\varphi$ मधील संख्यापकाच्या
$\lforall[x]$ किंवा
    $\lexists[x]$ या आस्थितीची व्याप्ती $\psi$ असेल,
तर $\psi$ मधील $x$ च्या मुक्त आस्थिती
$\lforall[x][\psi]$ आणि
$\lexists[x][\psi]$ मध्ये बद्ध असतात. या आस्थिती
त्या उल्लेख केलेल्या संख्यापक-आस्थितीने \emph{बद्ध केल्या आहेत} असे
आपण म्हणतो.
\end{defn}

\begin{ex}
पुढील सूत्राचा विचार करा:
\[\lexists[\Obj v_0][\underbrace{\Atom{\Obj A^2_0}{\Obj v_0,\Obj v_1}}_{\psi}]
\]
$\psi$ हे $\lexists[\Obj v_0]$ ची व्याप्ती दर्शवते. हा संख्यापक $\psi$
मधील $\Obj v_0$ ची आस्थिती बद्ध करतो; पण $\Obj v_1$ ची आस्थिती बद्ध
करत नाही. म्हणून या बाबतीत $\Obj v_1$ हे मुक्त चर आहे.


अधिक गुंतागुंतीच्या सूत्र~$\varphi$ मध्ये हे कसे कार्य करते ते आता
पाहू शकतो:
\[\lforall[\Obj v_0][\underbrace{(\Atom{\Obj A^1_0}{\Obj v_0} \lif
    \Atom{\Obj A^2_0}{\Obj v_0, \Obj v_1})}_{\psi}] \lif \lexists[\Obj
  v_1][\underbrace{(\Atom{\Obj A^2_1}{\Obj v_0, \Obj v_1} \lor \lforall[\Obj v_0][\overbrace{\lnot \Atom{\Obj A^1_1}{\Obj v_0}}^{\theta}])}_{\chi}]
\]
$\psi$ ही पहिल्या $\lforall[\Obj v_0]$ ची व्याप्ती, $\chi$ ही
$\lexists[\Obj v_1]$ ची व्याप्ती आणि $\theta$ ही दुसऱ्या
$\lforall[\Obj v_0]$ ची व्याप्ती आहे. पहिले $\lforall[\Obj v_0]$ हे
$\psi$ मधील $\Obj v_0$ च्या आस्थिती बद्ध करते, $\lexists[\Obj v_1]$ हे
$\chi$ मधील $\Obj v_1$ ची आस्थिती बद्ध करते आणि दुसरे
$\lforall[\Obj v_0]$ हे $\theta$ मधील $\Obj v_0$ ची आस्थिती बद्ध करते.
$\Obj v_1$ ची पहिली आस्थिती आणि $\Obj v_0$ ची चौथी आस्थिती $\varphi$ मध्ये
मुक्त आहेत. $\Obj v_0$ ची शेवटची आस्थिती $\theta$ मध्ये मुक्त, पण $\chi$
आणि~$\varphi$ मध्ये बद्ध आहे.
\end{ex}

\begin{defn}[वाक्य]
सूत्र~$\varphi$ मध्ये चलांच्या मुक्त आस्थिती नसतील तेव्हा
आणि केवळ तेव्हाच ते \emph{वाक्य} असते.
\end{defn}














\section{आदेशन}\label{fol:syn:sub:sec}

\begin{defn}[पदातील आदेशन]
$s$ मधील~$x$ च्या प्रत्येक आस्थितीच्या जागी $t$ \emph{आदेशित करून}
मिळणारा परिणाम $\Subst{s}{t}{x}$ आपण पुढीलप्रमाणे पुनरावर्ती रीतीने
परिभाषित करतो:
\begin{enumerate}
\item \indcase{s}{c}{$\Subst{\indfrm}{t}{x}$ म्हणजे फक्त $s$.}

\item \indcase{s}{y}{$y$ हे चर असून $y \not\ident x$ असेल, तर
  $\Subst{\indfrm}{t}{x}$ म्हणजेही फक्त~$s$.}

\item \indcase{s}{x}{$\Subst{\indfrm}{t}{x}$ म्हणजे~$t$.}

\item\indcase{s}{\Atom{f}{t_1, \dots, t_n}}{$\Subst{\indfrmp}{t}{x}$
  म्हणजे $\Atom{f}{\Subst{t_1}{t}{x}, \dots,
  \Subst{t_n}{t}{x}}$.}
\end{enumerate}
\end{defn}

\begin{defn}
$\varphi$ मधील~$x$ ची कोणतीही मुक्त आस्थिती $t$ मधील एखादे चर बद्ध करणाऱ्या
संख्यापकाच्या व्याप्तीत येत नसेल, तर पद~$t$ हे $\varphi$ मध्ये $x$ साठी
\emph{आदेशनासाठी मुक्त} असते.
\end{defn}

\begin{ex} ~
\begin{enumerate}
\item $\Obj v_8$ हे $\lexists[\Obj
  v_3]\Atom{\Obj A^2_4}{\Obj v_3,\Obj v_1}$ मध्ये $\Obj v_1$ साठी
  मुक्त आहे.

\item $\Obj f^2_1(\Obj v_1, \Obj v_2)$ हे
  $\lforall[\Obj v_2]\Atom{\Obj A^2_4}{\Obj v_0,\Obj v_2}$ मध्ये
  $\Obj v_0$ साठी \emph{मुक्त नाही}.
\end{enumerate}
\end{ex}

\begin{defn}[सूत्र मधील आदेशन]
$\varphi$ हे सूत्र, $x$ हे चल आणि $t$ हे $\varphi$ मध्ये~$x$
साठी आदेशनासाठी मुक्त असलेले पद असेल, तर $\varphi$ मधील~$x$ च्या सर्व मुक्त
आस्थितींच्या जागी $t$ आदेशित करून मिळणारा परिणाम म्हणजे
$\Subst{\varphi}{t}{x}$.
\begin{enumerate}
\item \indcase{\varphi}{\lfalse}{$\Subst{\indfrm}{t}{x}$
    म्हणजे $\lfalse$.}

\item \indcase{\varphi}{\ltrue}{$\Subst{\indfrm}{t}{x}$
    म्हणजे $\ltrue$.}

\item \indcase{\varphi}{\Atom{P}{t_1,\dots,
    t_n}}{$\Subst{\indfrm}{t}{x}$ म्हणजे
    $\Atom{P}{\Subst{t_1}{t}{x}, \dots, \Subst{t_n}{t}{x}}$.}

\item \indcase{\varphi}{\eq[t_1][t_2]}{$\Subst{\indfrmp}{t}{x}$ म्हणजे
  $\Subst{t_1}{t}{x} = \Subst{t_2}{t}{x}$.}

\item \indcase{\varphi}{\lnot \psi}{$\Subst{\indfrmp}{t}{x}$
    म्हणजे $\lnot \Subst{\psi}{t}{x}$.}

\item \indcase{\varphi}{(\psi \land
    \chi)}{$\Subst{\indfrmp}{t}{x}$ म्हणजे $(\Subst{\psi}{t}{x} \land
    \Subst{\chi}{t}{x})$.}

\item \indcase{\varphi}{(\psi \lor
    \chi)}{$\Subst{\indfrmp}{t}{x}$ म्हणजे $(\Subst{\psi}{t}{x} \lor
    \Subst{\chi}{t}{x})$.}

\item \indcase{\varphi}{(\psi \lif
    \chi)}{$\Subst{\indfrmp}{t}{x}$ म्हणजे $(\Subst{\psi}{t}{x} \lif
    \Subst{\chi}{t}{x})$.}

\item \indcase{\varphi}{(\psi \liff
    \chi)}{$\Subst{\indfrmp}{t}{x}$ म्हणजे $(\Subst{\psi}{t}{x} \liff
    \Subst{\chi}{t}{x})$.}

\item
\indcase{\varphi}{\lforall[y][\psi]}{$y$ हे $x$ पेक्षा वेगळे चर असेल, तर
    $\Subst{\indfrmp}{t}{x}$ म्हणजे
    $\lforall[y][\Subst{\psi}{t}{x}]$; अन्यथा
    $\Subst{\indfrmp}{t}{x}$ म्हणजे फक्त $\indfrm$.}

\item
\indcase{\varphi}{\lexists[y][\psi]}{$y$ हे $x$ पेक्षा वेगळे चर असेल, तर
    $\Subst{\indfrmp}{t}{x}$ म्हणजे
    $\lexists[y][\Subst{\psi}{t}{x}]$; अन्यथा
    $\Subst{\indfrmp}{t}{x}$ म्हणजे फक्त $\indfrm$.}
\end{enumerate}
\end{defn}

\begin{explain}
आदेशन निष्फळ असू शकते, हे लक्षात घ्या: $\varphi$ मध्ये $x$ मुळीच येत नसेल,
तर $\Subst{\varphi}{t}{x}$ म्हणजे फक्त~$\varphi$.

$\varphi$ मध्ये~$x$ साठी $t$ हे आदेशनासाठी मुक्त असले पाहिजे हा निर्बंध पुढील
प्रकारची प्रकरणे वगळण्यासाठी आवश्यक आहे. $\varphi \ident
\lexists[y][x < y]$ आणि $t \ident y$ असेल, तर
$\Subst{\varphi}{t}{x}$ हे $\lexists[y][y < y]$ होईल. या प्रकरणात आदेशन
करताना मुक्त चर~$y$ हे संख्यापक $\lexists[y]$ ने ``ग्रहित'' केले जाते;
हा परिणाम अनिष्ट आहे. उदाहरणार्थ, $\lforall[x][\psi]$ लागू असेल तेव्हा
$\Subst{\psi}{t}{x}$ हेही लागू असावे, अशी आपली अपेक्षा असते. पण
$\lforall[x][\lexists[y][x < y]]$ चा विचार करा (येथे $\psi$ म्हणजे
$\lexists[y][x < y]$). हे, उदाहरणार्थ, नैसर्गिक संख्यांविषयी सत्य
असलेले वाक्य आहे: प्रत्येक संख्या~$x$ पेक्षा मोठी एखादी संख्या~$y$
असते. $x$ च्या जागी $y$ चे आदेशन करण्याची परवानगी दिली, तर आपल्याला
$\Subst{\psi}{y}{x} \ident \lexists[y][y < y]$ मिळेल; हे असत्य आहे.
$t$ मधील कोणतेही मुक्त चर $\varphi$ मधील संख्यापकाने अखेरीस बद्ध होणार नाही,
अशी अट घालून आपण हे टाळतो.
\end{explain}

अवजड लेखन टाळण्यासाठी आपण पुढील संकेत अनेकदा वापरतो: सूत्र~$\varphi$
मध्ये चल~$x$ मुक्तपणे येऊ शकत असेल, तर हे दाखवण्यासाठी आपण
$\varphi(x)$ असेही लिहितो. कोणते $\varphi$ आणि~$x$ अभिप्रेत आहेत हे स्पष्ट असेल
आणि $t$ हे पद असेल ($\varphi(x)$ मध्ये $x$ साठी ते मुक्त आहे असे गृहीत
धरलेले असेल), तर $\Subst{\varphi}{t}{x}$ याचे संक्षिप्त रूप म्हणून आपण
$\varphi(t)$ लिहितो. म्हणून, उदाहरणार्थ, ``$\varphi(t)$ याला
$\lforall[x][\varphi(x)]$ चा दृष्टांत म्हणतो,'' असे आपण म्हणू शकतो. याचा
अर्थ असा: $\varphi$ हे कोणतेही सूत्र, $x$ हे चल आणि $t$ हे
$\varphi$ मध्ये~$x$ साठी मुक्त असलेले पद असेल, तर
$\Subst{\varphi}{t}{x}$ हा $\lforall[x][\varphi]$ चा दृष्टांत असतो.



\chapter{प्रथम-क्रम तर्कशास्त्राची चिन्हार्थमीमांसा}\label{fol:sem::chap}









\section{परिचय}\label{fol:syn:its:sec}

व्यंजकांना अर्थ देणे हे चिन्हार्थमीमांसेचे क्षेत्र आहे. चिन्हार्थमीमांसेतील
मध्यवर्ती संकल्पना म्हणजे संरचना मधील पूर्ति. संरचना
भाषेच्या मूलभूत घटकांना अर्थ देते: प्रांत म्हणजे वस्तूंचा अरिक्त
संच. संख्यापक या प्रांतातील घटकांवर फिरतात असा त्यांचा अर्थ लावला जातो,
स्थिरांक ना प्रांतातील घटक नेमून दिले जातात, फलने ना
प्रांत पासून त्याच प्रांताकडे जाणारी फलने नेमून दिली जातात आणि
विधेये ना प्रांत वरील संबंध नेमून दिले जातात. प्रांत
आणि मूलभूत शब्दसंग्रहातील चिन्हांना केलेल्या नेमणुका मिळून
संरचना बनते.
चले ही सूत्रे मध्ये येऊ शकतात; त्यांना चिन्हार्थ देण्यासाठी
आपल्याला प्रांत मधील घटक ही नेमून द्यावी लागतात---हेच
चर-मूल्यांकन होय. शेवटी, पूर्ति-संबंध या सर्व गोष्टी एकत्र आणतो.
सूत्राची संरचना~$\Struct{M}$ मध्ये चल च्या
मूल्यांकन~$s$ च्या सापेक्ष पूर्ति होऊ शकते; हे $\Sat{M}{\varphi}[s]$ असे
लिहितात. हा संबंधदेखील~$\varphi$ च्या संरचनेवरील विगमनाने परिभाषित केला जातो.
उदाहरणार्थ, $(\varphi \land \psi)$ ची पूर्ति परिभाषित करताना तार्किक संयोजकांचे
सत्यता-कोष्टक वापरून ती $\varphi$ आणि~$\psi$ यांची पूर्ति होते (किंवा होत नाही)
यांच्या आधारे सांगितली जाते. त्यानंतर असे निष्पन्न होते की सूत्र~$\varphi$
हे वाक्य, म्हणजे त्यात मुक्त चरे नसतील, तर चल चे
मूल्यांकन असंबद्ध ठरते; त्यामुळे संरचनांमध्ये वाक्यांची
सरळ पूर्ति होते (किंवा होत नाही) असे आपण म्हणू शकतो.

वाक्ये साठीच्या पूर्ति-संबंध~$\Sat{M}{\varphi}$ च्या आधारे आपण वैधता,
तार्किक निष्पन्नता आणि पूर्ततायोग्यता या मूलभूत चिन्हार्थविषयक संकल्पना
परिभाषित करू शकतो. वाक्याची प्रत्येक संरचनांमध्ये पूर्ति
होत असेल, तर ते वैध, $\Entails \varphi$, असते. वाक्यांच्या एखाद्या
संचातून ते तार्किकरीत्या निष्पन्न होते, $\Gamma \Entails \varphi$, जर~$\Gamma$ मधील
सर्व वाक्यांची पूर्ति करणारे प्रत्येक संरचना हे~$\varphi$ चीही
पूर्ति करत असेल. आणि एखाद्या वाक्यांच्या संचातील सर्व
वाक्यांची एकाच वेळी पूर्ति करणारे एखादे संरचना असेल, तर तो
संच पूर्ततायोग्य असतो. सूत्रे विगमनाधारित रीतीने परिभाषित केली आहेत
आणि पूर्ति हीही सूत्रांच्या संरचनेवरील विगमनाने परिभाषित केली जाते;
म्हणून आपल्या चिन्हार्थमीमांसेचे गुणधर्म सिद्ध करण्यासाठी आणि परिभाषित
चिन्हार्थविषयक संकल्पनांचा परस्परसंबंध लावण्यासाठी विगमन वापरता येते.












\section{प्रथम-क्रम भाषांसाठीच्या संरचना}\label{fol:syn:str:sec}

\begin{explain}
प्रथम-क्रम भाषांना स्वतःचा कोणताही \emph{अर्थ लावलेला नसतो:}
स्थिरांक, फलने आणि विधेये यांना कोणताही विशिष्ट
अर्थ जोडलेला नसतो. एखादी \emph{संरचना}
विनिर्दिष्ट करून अर्थ दिले जातात. ती \emph{प्रांत}, म्हणजे
स्थिरांक ज्या वस्तू निर्देशित करतात, फलने ज्यांवर क्रिया
करतात आणि संख्यापक ज्या वस्तूंवर फिरतात त्या वस्तू विनिर्दिष्ट करते.
याशिवाय, कोणते स्थिरांक कोणत्या वस्तू निर्देशित करतात,
फलन वस्तूंचे वस्तूंमध्ये प्रतिचित्रण कसे करते आणि विधेये
कोणत्या वस्तूंना लागू होतात हे ती विनिर्दिष्ट करते. संरचना या
तर्कशास्त्रातील \emph{चिन्हार्थविषयक} संकल्पनांचा---उदाहरणार्थ, तार्किक
निष्पन्नता, वैधता आणि पूर्ततायोग्यता या संकल्पनांचा---आधार आहेत. साहित्यात त्यांना
निरनिराळ्या ठिकाणी ``रचना,'' ``अर्थनिर्धारण'' किंवा ``प्रतिमान'' असे
म्हटले जाते.
\end{explain}

\begin{defn}[संरचना]
प्रथम-क्रम तर्कशास्त्राची भाषा~$\Lang{L}$ हिच्यासाठीची
\emph{संरचना}~$\Struct M$ पुढील घटकांनी बनते:
\begin{enumerate}
\item \emph{प्रांत:} अरिक्त संच~$\Domain M$.
\item \emph{स्थिरांकाचे अर्थनिर्धारण:} $\Lang{L}$ मधील प्रत्येक
  स्थिरांक~$c$ साठी घटक~$\Assign{c}{M} \in \Domain M$.
\item \emph{विधेयांचे अर्थनिर्धारण:} $\Lang{L}$ मधील प्रत्येक
  $n$-स्थानी विधेय~$R$ साठी ($\eq$ सोडून) $n$-स्थानी संबंध
  $\Assign{R}{M} \subseteq \Domain{M}^n$.
\item \emph{फलनांचे अर्थनिर्धारण:} $\Lang{L}$ मधील प्रत्येक
  $n$-स्थानी फलन~$f$ साठी $n$-स्थानी फलन
  $\Assign{f}{M} \colon \Domain{M}^n \to \Domain{M}$.
\end{enumerate}
\end{defn}

\begin{ex}
अंकगणिताच्या भाषेसाठीच्या संरचना~$\Struct M$ मध्ये एक संच,
स्थिरांक~$\Obj 0$ चे अर्थनिर्धारण म्हणून $\Domain M$ मधील
घटक~$\Assign{\Obj 0}{M}$, एकस्थानी फलन
$\Assign{\Obj \prime}{M} \colon \Domain{M} \to \Domain M$,
$\Domain M^2 \to \Domain M$ अशी दोन द्विस्थानी फलने
$\Assign{\Obj +}{M}$ आणि $\Assign{\Obj \times}{M}$, तसेच द्विस्थानी
संबंध~$\Assign{\Obj <}{M} \subseteq \Domain{M}^2$ असतात.

अशा रचनेचे एक स्वाभाविक उदाहरण पुढीलप्रमाणे आहे:
\begin{enumerate}
\item $\Domain N = \Nat$
\item $\Assign{\Obj 0}{N} = 0$
\item प्रत्येक $n \in \Nat$ साठी $\Assign{\Obj \prime}{N}(n) = n + 1$
\item प्रत्येक $n, m \in \Nat$ साठी $\Assign{\Obj +}{N}(n, m) = n + m$
\item प्रत्येक $n, m \in \Nat$ साठी $\Assign{\Obj \times}{N}(n, m) = n\cdot m$
\item $\Assign{\Obj <}{N} = \Setabs{\tuple{n, m}}{n \in \Nat, m \in
  \Nat, n < m}$
\end{enumerate}
अशा रीतीने परिभाषित केलेल्या $\Lang L_A$ साठीच्या रचनेला~$\Struct N$
\emph{अंकगणिताचे मानक प्रतिमान} म्हणतात, कारण ती~$\Lang L_A$ मधील
अतार्किक चिन्हांचा अर्थ आपण अपेक्षा करू त्याच रीतीने लावते.

तरीही, $\Lang L_A$ साठी इतर अनेक संरचना शक्य आहेत. उदाहरणार्थ,
प्रांत म्हणून~$\Nat$ ऐवजी पूर्णांकांचा संच~$\Int$ घेऊन $\Obj 0$,
$\Obj \prime$, $\Obj +$, $\Obj \times$, $\Obj <$ यांची अर्थनिर्धारणे
त्यानुसार परिभाषित करता येतील. पण संख्यांशी दुरान्वयानेही संबंध नसलेल्या
$\Lang L_A$ साठीच्या रचनाही आपण परिभाषित करू शकतो.
\end{ex}

\begin{ex}
संचसिद्धान्ताची भाषा~$\Lang L_Z$ हिच्यासाठीची रचना~$\Struct M$ हिला
फक्त एक संच आणि एकच द्विस्थानी संबंध आवश्यक असतात. म्हणून तांत्रिकदृष्ट्या,
उदाहरणार्थ, माणसांचा संच आणि ``$x$ हा~$y$ पेक्षा वयाने मोठा आहे'' हा संबंध
$\Lang L_Z$ साठी संरचना म्हणून वापरता येईल; त्याचप्रमाणे~$\Nat$
आणि $n, m \in \Nat$ साठीचा $n \ge m$ हाही वापरता येईल.


$\Lang L_Z$ साठीची एक विशेष रोचक संरचना, जिच्या प्रांतातील
घटक प्रत्यक्षात संच आहेत आणि जिच्यात $\Obj \in$ चे अर्थनिर्धारण
प्रत्यक्षात ``$x$ हा~$y$ चा घटक आहे'' हा संबंध आहे, ती म्हणजे
\emph{आनुवंशिकदृष्ट्या सांत संचांची} संरचना~$\Struct{{HF}}$:
\begin{enumerate}
\item $\Domain{{{HF}}} = \emptyset \cup \Pow{\emptyset} \cup
  \Pow{\Pow{\emptyset}} \cup \Pow{\Pow{\Pow{\emptyset}}} \cup \dots$;
\item $\Assign{\Obj \in}{{{HF}}} = \Setabs{\tuple{x, y}}{x, y \in
  \Domain{{{HF}}}, x \in y}$.
\end{enumerate}
\end{ex}

\begin{digress}
संरचना म्हणून काय गणले जाईल याविषयी आपण केलेल्या अटींचा आपल्या
तर्कशास्त्रावर परिणाम होतो. उदाहरणार्थ, प्रांत अरिक्त ठेवण्याची निवड
संख्यकीकृत वाक्यांच्या पूर्तिची (किंवा सत्यतेची) नेहमीची मांडणी गृहीत
धरल्यास $\lexists[x][(\varphi(x) \lor \lnot \varphi(x))]$ वैध---म्हणजे तार्किक
सत्य---असल्याची खात्री देते. तसेच, सर्व स्थिरांक नी प्रांतातील
वस्तूच निर्देशित केली पाहिजे ही अट अस्तित्ववाची सामान्यीकरण हा
निर्दोष अनुमान-रूपबंध असल्याची खात्री देते: $\varphi(a)$, म्हणून
$\lexists[x][\varphi(x)]$. नावांना प्रांताबाहेरील वस्तू निर्देशित करण्याची
किंवा काहीच निर्देशित न करण्याची मुभा दिली, तर आपण \emph{मुक्त
तर्कशास्त्राच्या} दिशेने जाऊ; त्यात अस्तित्ववाची सामान्यीकरणासाठी एक
अतिरिक्त आधारवाक्यालागते: $\varphi(a)$ आणि $\lexists[x][\eq[x][a]]$, म्हणून
$\lexists[x][\varphi(x)]$.
\end{digress}












\section{प्रथम-क्रम भाषांसाठी बंद पदांनी आच्छादित संरचना}\label{fol:syn:cov:sec}

\begin{explain}
एखाद्या पदात चले नसतील, तर ते \emph{बंद} असते, हे आठवा.
\end{explain}

\begin{defn}[बंद पदांचे मूल्य]
$t$ हे भाषा~$\Lang L$ हिचे बंद पद आणि $\Struct M$ ही~$\Lang L$ साठीची
संरचना असेल, तर \emph{मूल्य}~$\Value{t}{M}$ पुढीलप्रमाणे
परिभाषित केले जाते:
\begin{enumerate}
\item $t$ हे फक्त स्थिरांक~$c$ असेल, तर
  $\Value{c}{M} = \Assign{c}{M}$.
\item $t$ हे $\Atom{f}{t_1, \ldots, t_n}$ या रूपात असेल, तर
  \[  \Value{t}{M} = \Assign{f}{M}(\Value{t_1}{M}, \ldots,
  \Value{t_n}{M}).
  \]
\end{enumerate}
\end{defn}

\begin{defn}[आच्छादित संरचना]
प्रांतातील प्रत्येक घटक हा एखाद्या बंद पदाचे मूल्य असेल, तर ती
संरचना \emph{आच्छादित} असते.
\end{defn}

\begin{ex}
स्थिरांक $\Obj{zero}$, $\Obj{one}$, $\Obj{two}$, \dots, द्विस्थानी
विधेय~$<$ आणि द्विस्थानी फलने $+$ व~$\times$ असलेली
भाषा~$\Lang L$ घेऊ. मग~$\Lang L$ साठीची संरचना~$\Struct M$ अशी
आहे: तिचा प्रांत $\Domain M = \{0, 1, 2, \ldots \}$ असून
$\Assign{\Obj{zero}}{M} = 0$, $\Assign{\Obj{one}}{M} = 1$,
$\Assign{\Obj{two}}{M} = 2$ आणि असेच पुढे ही तिची अर्थनिर्धारणे आहेत.
द्विस्थानी संबंधचिन्ह~$<$ साठी $\Assign{<}{M}$ हा अशा सर्व जोड्यांचा
संच आहे की $\tuple{c_1, c_2} \in \Domain{M}^2$ आणि $c_1$ हा~$c_2$ पेक्षा
लहान आहे: उदाहरणार्थ, $\tuple{1, 3} \in \Assign{<}{M}$, पण
$\tuple{2, 2} \notin \Assign{<}{M}$. द्विस्थानी फलन~$+$ साठी
$\Assign{+}{M}$ नेहमीच्या रीतीने परिभाषित करू---उदाहरणार्थ,
$\Assign{+}{M}(2,3)$ हे~$5$ कडे पाठवते; आणि द्विस्थानी
फलन~$\times$ साठीही त्याचप्रमाणे. म्हणून $\Obj{four}$ चे
मूल्य फक्त~$4$ आहे आणि $\times(\Obj{two},
+(\Obj{three},\Obj{zero}))$ चे मूल्य (किंवा मध्यस्थ संकेतपद्धतीत,
$\Obj{two} \times (\Obj{three} + \Obj{zero})$) पुढीलप्रमाणे आहे:

\begin{multline*}
\Value{\times(\Obj{two}, +(\Obj{three},\Obj{zero}))}{M} \\
\begin{aligned}
& =\Assign{\times}{M}(\Value{\Obj{two}}{M}, \Value{+(\Obj{three}, \Obj{zero})}{M})\\
& = \Assign{\times}{M}(\Value{\Obj{two}}{M}, \Assign{+}{M}(\Value{\Obj{three}}{M},
\Value{\Obj{zero}}{M})) \\
& = \Assign{\times}{M}(\Assign{\Obj{two}}{M}, \Assign{+}{M}(\Assign{\Obj{three}}{M},
\Assign{\Obj{zero}}{M})) \\
& = \Assign{\times}{M}(2, \Assign{+}{M}(3, 0)) \\
& = \Assign{\times}{M}(2, 3) \\
& = 6
\end{aligned}
\end{multline*}
\end{ex}

\begin{prob}
अंकगणिताचे मानक प्रतिमान~$\Struct N$ आच्छादित आहे का? स्पष्टीकरण द्या.
\end{prob}












\section{संरचना मधील
  सूत्राची पूर्ति}\label{fol:syn:sat:sec}

\begin{explain}
एकीकडे पदे आणि सूत्रे यांसारख्या अभिव्यक्ती आणि दुसरीकडे
संरचना यांना जोडणाऱ्या मूलभूत संकल्पना म्हणजे पदाचे
\emph{मूल्य} आणि सूत्राची \emph{पूर्ति}. अनौपचारिकपणे,
पदाचे मूल्य म्हणजे संरचना मधील घटक असतो---पद
नुसते स्थिरचिन्ह असेल, तर त्या स्थिरचिन्हाला संरचना ने नेमून
दिलेली वस्तू हे त्याचे मूल्य असते; आणि ते फलने वापरून
बनवलेले असेल, तर पदातील स्थिरचिन्हांची मूल्ये आणि फलनचिन्हांना
नेमून दिलेली फलने यांपासून त्याचे मूल्य काढले जाते. विधेयांना
दिलेल्या अर्थनिर्धारणामुळे संरचनेच्या प्रांतात सूत्र सत्य
होत असेल, तर सूत्राची त्या संरचनांमध्ये \emph{पूर्ति
होते}. पूर्तिची ही संकल्पना विगमनाधारित रीतीने विनिर्दिष्ट केली जाते:
आण्विक सूत्रांची पूर्ति केव्हा होते हे संरचनेचे
विनिर्देशन थेट सांगते; आणि एखाद्या जटिल सूत्राची पूर्ति केव्हा
होते हे आपण तिचा मुख्य संयोजक किंवा संख्यापक आणि तिच्या तात्काळ
उपसूत्रांची पूर्ति होते की नाही यांवरून परिभाषित करतो.

इथे संख्यापकांची स्थिती थोडी अवघड आहे, कारण संख्यकीकृत सूत्र
च्या तात्काळ उपसूत्रामध्ये एक मुक्त चल असते आणि
संरचना ही चलांची मूल्ये विनिर्दिष्ट करत नाहीत.
ही अडचण हाताळण्यासाठी आपण \emph{चर-मूल्यांकने} ही आणतो आणि पूर्तिची
व्याख्या केवळ संरचनेच्या सापेक्ष न देता, संरचना आणि
चल चे मूल्यांकन या दोहोंच्या सापेक्ष देतो.
\end{explain}

\begin{defn}[चर-मूल्यांकन]
संरचना~$\Struct{M}$ हिच्यासाठीचे \emph{चर-मूल्यांकन}~$s$
म्हणजे प्रत्येक चल चे~$\Domain M$ मधील घटक कडे
प्रतिचित्रण करणारे फलन; म्हणजेच, $s\colon \Var \to \Domain M$.
\end{defn}

\begin{explain}
संरचना प्रत्येक स्थिरांक ला मूल्य नेमून देते, तर
चर-मूल्यांकन प्रत्येक चराला मूल्य नेमून देते. पण त्यांच्यापासून बनलेली
पदेही प्रांत मधील घटक ना नावे देण्यासाठी वापरायची आहेत.
त्यासाठी आपण पदांचे मूल्य विगमनाधारित रीतीने परिभाषित करतो.
स्थिरांक आणि चरांसाठी मूल्य अनुक्रमे संरचना किंवा
चर-मूल्यांकन विनिर्दिष्ट करते तेच असते; अधिक जटिल पदांसाठी ते
संरचना ने फलने ना नेमून दिलेली फलने वापरून पुनरावर्ती
रीतीने काढले जाते.
\end{explain}

\begin{defn}[पदांचे मूल्य]
$t$ हे भाषा~$\Lang L$ हिचे पद, $\Struct M$ ही~$\Lang L$ साठीची
संरचना आणि $s$ हे~$\Struct M$ साठीचे चल चे
मूल्यांकन असेल, तर \emph{मूल्य}~$\Value{t}{M}[s]$ पुढीलप्रमाणे
परिभाषित केले जाते:
\begin{enumerate}
\item \indcase{t}{c}{$\Value{\indfrm}{M}[s] = \Assign{\indcomplex}{M}$.}
\item \indcase{t}{x}{$\Value{\indfrm}{M}[s] = s(\indcomplex)$.}
\item \indcase{t}{\Atom{f}{t_1, \ldots, t_n}}{
\[\Value{\indfrm}{M}[s] = \Assign{f}{M}(\Value{t_1}{M}[s], \ldots,
\Value{t_n}{M}[s]).
\]}
\end{enumerate}
\end{defn}

\begin{defn}[$x$-पर्याय]
$s$ हे संरचना~$\Struct M$ साठीचे चल चे मूल्यांकन
असेल, तर~$s$ पासून फार तर~$x$ ला नेमून दिलेल्या मूल्यातच भिन्न असणारे
$\Struct M$ साठीचे कोणतेही चल चे मूल्यांकन~$s'$ याला~$s$ चा
\emph{$x$-पर्याय} म्हणतात. $s'$ हा~$s$ चा $x$-पर्याय असेल, तर आपण
$\varAssign{s'}{s}{x}$ असे लिहितो.
\end{defn}

\begin{explain}
मूल्यांकन~$s$ च्या $x$-पर्यायाने~$x$ ला काही वेगळेच नेमून दिले
\emph{पाहिजे} असे नाही, हे लक्षात घ्या. वस्तुतः, प्रत्येक मूल्यांकन
स्वतःचाच $x$-पर्याय मानला जातो.
\end{explain}

\begin{defn}
  $s$ हे संरचना~$\Struct M$ साठीचे चल चे मूल्यांकन
  आणि $m \in \Domain{M}$ असेल, तर मूल्यांकन~$\Subst{s}{m}{x}$ हे पुढील
  रीतीने परिभाषित केलेले चर-मूल्यांकन आहे:
  \[\Subst{s}{m}{x}(y) = \begin{cases}
    m & \text{जर } y \ident x\\
    s(y) & \text{अन्यथा}.
  \end{cases}\]
\end{defn}

दुसऱ्या शब्दांत, $\Subst{s}{m}{x}$ हा~$s$ चा असा विशिष्ट $x$-पर्याय
आहे, जो प्रांतातील घटक~$m$ हे~$x$ ला नेमून देतो आणि~$x$ शिवाय
इतर चले ना~$s$ ने नेमून दिलेल्याच गोष्टी नेमून देतो.

\begin{defn}[पूर्ति]
\label{fol:syn:sat:defn:satisfaction}
संरचना~$\Struct M$ मध्ये चल चे मूल्यांकन~$s$ याच्या
सापेक्ष सूत्र~$\varphi$ ची पूर्ति, संकेतांत $\Sat{M}{\varphi}[s]$, पुढीलप्रमाणे
पुनरावर्ती रीतीने परिभाषित केली जाते. (``$\Sat{M}{\varphi}[s]$ नाही'' या
अर्थाने आपण $\Sat/{M}{\varphi}[s]$ असे लिहितो.)
\begin{enumerate}
\item
  \indcase{\varphi}{\lfalse}{$\Sat/{M}{\indfrm}[s]$.}

\item
  \indcase{\varphi}{\ltrue}{$\Sat{M}{\indfrm}[s]$.}

\item \indcase{\varphi}{\Atom{R}{t_1, \dots, t_n}}{$\Sat{M}{\indfrm}[s]$
  तेव्हा आणि केवळ तेव्हाच, जेव्हा $\langle \Value{t_1}{M}[s], \dots,
  \Value{t_n}{M}[s] \rangle \in \Assign{R}{M}$.}

\item \indcase{\varphi}{\eq[t_1][t_2]}{$\Sat{M}{\indfrm}[s]$ तेव्हा आणि
  केवळ तेव्हाच, जेव्हा $\Value{t_1}{M}[s] = \Value{t_2}{M}[s]$.}

\item
  \indcase{\varphi}{\lnot \psi}{$\Sat{M}{\indfrm}[s]$ तेव्हा आणि केवळ तेव्हाच,
    जेव्हा $\Sat/{M}{\psi}[s]$.}

\item
  \indcase{\varphi}{(\psi \land \chi)}{$\Sat{M}{\indfrm}[s]$ तेव्हा आणि केवळ
    तेव्हाच, जेव्हा $\Sat{M}{\psi}[s]$ आणि $\Sat{M}{\chi}[s]$.}

\item
  \indcase{\varphi}{(\psi \lor \chi)}{$\Sat{M}{\indfrm}[s]$ तेव्हा आणि केवळ
    तेव्हाच, जेव्हा $\Sat{M}{\psi}[s]$ किंवा $\Sat{M}{\chi}[s]$ (किंवा
    दोन्ही).}

\item
  \indcase{\varphi}{(\psi \lif \chi)}{$\Sat{M}{\indfrm}[s]$ तेव्हा आणि केवळ
    तेव्हाच, जेव्हा $\Sat/{M}{\psi}[s]$ किंवा $\Sat{M}{\chi}[s]$ (किंवा
    दोन्ही).}

\item
  \indcase{\varphi}{(\psi \liff \chi)}{$\Sat{M}{\indfrm}[s]$ तेव्हा आणि केवळ
    तेव्हाच, जेव्हा एकतर $\Sat{M}{\psi}[s]$ आणि $\Sat{M}{\chi}[s]$ दोन्ही,
    किंवा $\Sat{M}{\psi}[s]$ व $\Sat{M}{\chi}[s]$ यांपैकी एकही नाही.}

\item
  \indcase{\varphi}{\lforall[x][\psi]}{$\Sat{M}{\indfrm}[s]$ तेव्हा आणि केवळ
    तेव्हाच, जेव्हा प्रत्येक घटक~$m \in \Domain M$ साठी
    $\Sat{M}{\psi}[\Subst{s}{m}{x}]$.}

\item
  \indcase{\varphi}{\lexists[x][\psi]}{$\Sat{M}{\indfrm}[s]$ तेव्हा आणि केवळ
  तेव्हाच, जेव्हा किमान एका घटक~$m \in \Domain M$ साठी
  $\Sat{M}{\psi}[\Subst{s}{m}{x}]$.}
\end{enumerate}
\end{defn}

\begin{explain}
शेवटच्या दोन उपवाक्यांत
चर-मूल्यांकने महत्त्वाची आहेत. ``प्रत्येक $m \in
\Domain{M}$ साठी $\Sat{M}{\psi(m)}$'' असे म्हणून आपण
$\lforall[x][\psi(x)]$ ची पूर्ति परिभाषित करू शकत नाही.
 ``किमान एका $m \in \Domain{M}$ साठी $\Sat{M}{\psi(m)}$''
असे म्हणून आपण $\lexists[x][\psi(x)]$ ची पूर्ति परिभाषित करू शकत नाही.
कारण $m \in \Domain M$ असेल, तर ते भाषेचे चिन्ह नसते आणि म्हणून
$\psi(m)$ हे सूत्र नसते (म्हणजे $\Subst{\psi}{m}{x}$ अपरिभाषित
असते).~$\Struct{M}$ मधील प्रत्येक घटक चे नाव देणारी
स्थिरांक किंवा पदे
आपल्याजवळ उपलब्ध आहेत असेही गृहीत धरता येत नाही, कारण संरचना
च्या व्याख्येत तशी कोणतीही अट नाही. प्रमाण भाषेतील स्थिरांकाचा
संच गणनीय अनंत आहे; म्हणून $\Domain{M}$ हा गणनीय नसेल,
तर प्रत्येक वस्तूला नाव देण्यासाठी पुरेशी स्थिरांक सुद्धा नसतात.

ही अडचण सोडवण्यासाठी आपण चल ची मूल्यांकने आणतो; त्यांमुळे
चरे प्रांतातील घटक शी थेट जोडता येतात. मग, उदाहरणार्थ,
$\lexists[x][\psi(x)]$ ची~$\Struct M$
मध्ये पूर्ति तेव्हा आणि केवळ तेव्हाच होते, जेव्हा
किमान एका $m \in \Domain{M}$ साठी $\Sat{M}{\psi(m)}$, असे न म्हणता आपण म्हणतो:
तिची~$\Struct M$ मध्ये~$s$ च्या \emph{सापेक्ष} पूर्ति तेव्हा आणि केवळ
तेव्हाच होते, जेव्हा $\psi(x)$ ची~$\Subst{s}{m}{x}$ च्या सापेक्ष
किमान एका $m \in \Domain M$ साठी पूर्ति होते.

\end{explain}

\begin{ex}
$\Lang{L} = \{a, b, f, R\}$ असू द्या, जिथे $a$ आणि $b$ ही
स्थिरांक, $f$ हे दोन-स्थानी फलन आणि $R$ हे दोन-स्थानी
विधेय आहे. पुढीलप्रमाणे परिभाषित केलेली संरचना~$\Struct{M}$
विचारात घ्या:
\begin{enumerate}
\item $\Domain M = \{1, 2, 3, 4\}$
\item $\Assign{a}{M} = 1$
\item $\Assign{b}{M} = 2$
\item $x+y \le 3$ असेल तर $\Assign{f}{M}(x, y) = x+y$, अन्यथा $= 3$.
\item $\Assign{R}{M} = \{\tuple{1, 1}, \tuple{1, 2}, \tuple{2, 3}, \tuple{2, 4}\}$
\end{enumerate}
प्रत्येक चल ला $1 \in \Domain{M}$ नेमून देणारे फलन
$s(x) = 1$ हे~$\Struct{M}$ साठीचे चर-मूल्यांकन आहे.

मग
\begin{align*}
\Value{f(a,b)}{M}[s] & = \Assign{f}{M}(\Value{a}{M}[s], \Value{b}{M}[s]).
\intertext{$a$ आणि $b$ ही स्थिरांक असल्यामुळे $\Value{a}{M}[s]
  = \Assign{a}{M} = 1$ आणि $\Value{b}{M}[s] = \Assign{b}{M} = 2$. म्हणून}
\Value{f(a,b)}{M}[s] & = \Assign{f}{M}(1, 2) = 1+2 = 3.
\intertext{$f(f(a,b),a)$ चे मूल्य काढण्यासाठी पुढील अभिव्यक्ती विचारात घ्यावी लागते:}
\Value{f(f(a,b),a)}{M}[s] & = \Assign{f}{M}(\Value{f(a, b)}{M}[s],
\Value{a}{M}[s]) = \Assign{f}{M}(3, 1) = 3,
\intertext{कारण $3+1 > 3$. $s(x) = 1$ आणि $\Value{x}{M}[s] =
  s(x)$ असल्यामुळे आपल्याकडे पुढील समीकरणही आहे:}
\Value{f(f(a,b),x)}{M}[s] & = \Assign{f}{M}(\Value{f(a, b)}{M}[s],
\Value{x}{M}[s]) = \Assign{f}{M}(3, 1) = 3,
\end{align*}

आण्विक सूत्र~$R(t_1, t_2)$ च्या फलसाधकांच्या मूल्यांची क्रमित
जोडी, म्हणजे $\tuple{\Value{t_1}{M}[s], \Value{t_2}{M}[s]}$, ही
$\Assign{R}{M}$ ची घटक असेल, तर त्या सूत्राची पूर्ति होते.
म्हणून, उदाहरणार्थ, $\tuple{\Value{b}{M}, \Value{f(a,b)}{M}} =
\tuple{2, 3} \in \Assign{R}{M}$ असल्यामुळे
$\Sat{M}{R(b,f(a,b))}[s]$; पण $\tuple{1, 3} \notin \Assign{R}{M}$
असल्यामुळे $\Sat/{M}{R(x, f(a,b))}[s]$.


अनाण्विक सूत्र~$\varphi$ ची पूर्ति होते का हे ठरवण्यासाठी त्याच्या मुख्य
संयोजकाला लागू होणारे विगमनाधारित व्याख्येतील उपवाक्य वापरावे. उदाहरणार्थ,
$R(a, a) \lif (R(b, x) \lor R(x, b))$ मधील मुख्य संयोजक~$\lif$ आहे आणि
\begin{align*}
  & \Sat{M}{R(a, a) \lif (R(b, x) \lor R(x, b))}[s] \text{ तेव्हा आणि केवळ तेव्हाच, जेव्हा }\\
  & \qquad
  \Sat/{M}{R(a,a)}[s] \text{ किंवा } \Sat{M}{R(b, x) \lor R(x, b)}[s]
  \intertext{$\Sat{M}{R(a,a)}[s]$ असल्यामुळे (कारण $\tuple{1,1} \in
    \Assign{R}{M}$), उत्तर अजून ठरवता येत नाही; प्रथम
    $\Sat{M}{R(b, x) \lor R(x, b)}[s]$ आहे का ते शोधावे लागते:}
  & \Sat{M}{R(b, x) \lor R(x, b)}[s] \text{ तेव्हा आणि केवळ तेव्हाच, जेव्हा }\\
  & \qquad \Sat{M}{R(b,x)}[s] \text{ किंवा } \Sat{M}{R(x, b)}[s]
  \intertext{आणि हे खरे आहे, कारण $\Sat{M}{R(x, b)}[s]$
    (कारण $\tuple{1,2} \in \Assign{R}{M}$).}
\end{align*}

आठवा, $s$ चा $x$-पर्याय म्हणजे~$s$ पासून फार तर~$x$ ला नेमून
दिलेल्या मूल्यातच भिन्न असणारे चर-मूल्यांकन. $\Domain{M}$ च्या प्रत्येक
घटक साठी~$s$ चा एक $x$-पर्याय आहे:
\begin{align*}
  s_1 & = \Subst{s}{1}{x}, &
  s_2 & = \Subst{s}{2}{x},\\
  s_3 & = \Subst{s}{3}{x}, &
  s_4 & = \Subst{s}{4}{x}.
\end{align*}
म्हणून, उदाहरणार्थ, $s_2(x) = 2$ आणि~$x$ शिवाय इतर सर्व चर~$y$ साठी
$s_2(y) = s(y) = 1$. $\Domain{M} = \{1, 2, 3, 4\}$ असल्यामुळे हेच
$\Struct{M}$ या रचनेसाठीचे~$s$ चे सर्व $x$-पर्याय आहेत. विशेषतः,
$s_1 = s$ हे लक्षात घ्या ($s$ हा नेहमीच स्वतःचा $x$-पर्याय असतो).

अस्तित्ववाची संख्यापक असलेल्या
  सूत्र~$\lexists[x][\varphi(x)]$ ची पूर्ति होते का हे ठरवण्यासाठी,
  किमान एका $m \in \Domain M$ साठी
  $\Sat{M}{\varphi(x)}[\Subst{s}{m}{x}]$ आहे का हे ठरवावे लागते. म्हणून,
  \[  \Sat{M}{\lexists[x][(R(b,x) \lor R(x,b))]}[s],
  \]
  कारण $\Sat{M}{R(b,x) \lor R(x, b)}[\Subst{s}{1}{x}]$
  ($\Subst{s}{3}{x}$ हेसुद्धा चालले असते). पण,
  \[  \Sat/{M}{\lexists[x][(R(b,x) \land R(x,b))]}[s]
  \]
  कारण कोणताही $m \in \Domain{M}$ निवडला, तरी
  $\Sat/{M}{R(b,x) \land R(x,b)}[\Subst{s}{m}{x}]$.

वैश्विक संख्यापक असलेल्या
  सूत्र~$\lforall[x][\varphi(x)]$ ची पूर्ति होते का हे ठरवण्यासाठी,
  प्रत्येक $m \in \Domain M$ साठी
  $\Sat{M}{\varphi(x)}[\Subst{s}{m}{x}]$ आहे का हे ठरवावे लागते. म्हणून,
  \[  \Sat{M}{\lforall[x][(R(x,a) \lif R(a,x))]}[s],
  \]
  कारण प्रत्येक $m \in \Domain M$ साठी
  $\Sat{M}{R(x,a) \lif R(a,x)}[\Subst{s}{m}{x}]$. $m = 1$ साठी
  $\Sat{M}{R(a,x)}[\Subst{s}{1}{x}]$ असल्यामुळे उत्तरवर्ती सत्य आहे;
  $m = 2$, $3$ आणि~$4$ साठी $\Sat/{M}{R(x,a)}[\Subst{s}{m}{x}]$
  असल्यामुळे पूर्ववर्ती असत्य आहे. पण,
  \[  \Sat/{M}{\lforall[x][(R(a,x) \lif R(x,a))]}[s]
  \]
  कारण $\Sat/{M}{R(a,x) \lif R(x,a)}[\Subst{s}{2}{x}]$ (कारण
  $\Sat{M}{R(a, x)}[\Subst{s}{2}{x}]$ आणि $\Sat/{M}{R(x,
  a)}[\Subst{s}{2}{x}]$).





अधिक जटिल उदाहरणासाठी पुढील सूत्र विचारात घ्या:
\[\lforall[x][(R(a,x) \lif \lexists[y][R(x,y)])].
\]
$\Sat/{M}{R(a,x)}[\Subst{s}{3}{x}]$ आणि
$\Sat/{M}{R(a,x)}[\Subst{s}{4}{x}]$ असल्यामुळे, सशर्त सूत्राच्या
उत्तरवर्तीची काळजी घ्यावी लागणारी रंजक प्रकरणे फक्त $m = 1$ आणि
$m = 2$ ही आहेत. $\Sat{M}{\lexists[y][R(x,y)]}[\Subst{s}{1}{x}]$ आहे
का? किमान एक $n \in \Domain M$ असा असेल की
$\Sat{M}{R(x,y)}[\Subst{\Subst{s}{1}{x}}{n}{y}]$, तर ते खरे आहे.
वस्तुतः, $n = 1$ घेतल्यास $\Subst{\Subst{s}{1}{x}}{n}{y} =
\Subst{s}{1}{y} = s$. $s(x) = 1$, $s(y) = 1$ आणि
$\tuple{1,1} \in \Assign{R}{M}$ असल्यामुळे उत्तर होकारार्थी आहे.


$\Sat{M}{\lexists[y][R(x,y)]}[\Subst{s}{2}{x}]$ आहे का हे ठरवण्यासाठी
आपल्याला चल ची मूल्यांकने
$\Subst{\Subst{s}{2}{x}}{n}{y}$ तपासावी लागतात. इथे $n = 1$ साठी हे
मूल्यांकन $s_2 = \Subst{s}{2}{x}$ आहे आणि ते $R(x,y)$ ची पूर्ति करत
नाही ($s_2(x) = 2$, $s_2(y) = 1$ आणि
$\tuple{2,1}\notin \Assign{R}{M}$). पण
$\Subst{\Subst{s}{2}{x}}{3}{y} = \Subst{s_2}{3}{y}$ विचारात घ्या.
$\tuple{2,3} \in \Assign{R}{M}$ असल्यामुळे
$\Sat{M}{R(x,y)}[\Subst{s_2}{3}{y}]$ आणि म्हणून
$\Sat{M}{\lexists[y][R(x,y)]}[s_2]$.

म्हणून, प्रत्येक $m \in \Domain M$ साठी एकतर
$\Sat/{M}{R(a,x)}[\Subst{s}{m}{x}]$ ($m = 3$, $4$ असेल तर) किंवा
$\Sat{M}{\lexists[y][R(x,y)]}[\Subst{s}{m}{x}]$ ($m = 1$, $2$ असेल तर),
आणि म्हणून
\[\Sat{M}{\lforall[x][(R(a,x) \lif \lexists[y][R(x,y)])]}[s].
\]
दुसरीकडे,
\[\Sat/{M}{\lexists[x][(R(a,x) \land \lforall[y][R(x,y)])]}[s].
\]
फक्त $m = 1$ आणि $m = 2$ यांच्यासाठी
$\Sat{M}{R(a,x)}[\Subst{s}{m}{x}]$. पण~$m$ च्या या दोन्ही मूल्यांपैकी
प्रत्येकासाठी एक $n \in \Domain M$ आहे---पहिल्या मूल्यासाठी $n = 4$
आणि दुसऱ्यासाठी $n = 1$---ज्यामुळे
$\Sat/{M}{R(x,y)}[\Subst{\Subst{s}{m}{x}}{n}{y}]$ आणि म्हणून $m = 1$
व $m = 2$ साठी
$\Sat/{M}{\lforall[y][R(x,y)]}[\Subst{s}{m}{x}]$. एकंदरीत, असा कोणताही
$m \in \Domain M$ नाही की
$\Sat{M}{R(a,x) \land \lforall[y][R(x,y)]}[\Subst{s}{m}{x}]$.


\end{ex}









\begin{prob}
एक स्थिरांक, एक एक-स्थानी फलन आणि एक दोन-स्थानी
विधेय असलेली $\Lang L = \{c, f, A\}$ असू द्या; आणि
संरचना~$\Struct{M}$ पुढीलप्रमाणे दिलेली असू द्या:
\begin{enumerate}
\item $\Domain M = \{1, 2, 3\}$
\item $\Assign{c}{M} = 3$
\item $\Assign{f}{M}(1) = 2, \Assign{f}{M}(2) = 3, \Assign{f}{M}(3) = 2$
\item $\Assign{A}{M} = \{\tuple{1, 2}, \tuple{2, 3}, \tuple{3, 3}\}$
\end{enumerate}
(a) सर्व चले~$v$ साठी $s(v) = 1$ असू द्या. पुढील विधान
खरे आहे की नाही हे शोधा:
\[\Sat{M}{\lexists[x][(A(f(z), c) \lif \lforall[y][(A(y, x) \lor A(f(y),
      x))])]}[s]
\]
तुमचे उत्तर का योग्य आहे ते स्पष्ट करा.

(b) ज्यात या सूत्राची पूर्ति होत नाही अशी वेगळी रचना आणि
चल चे मूल्यांकन द्या.
\end{prob}












\section{चर-मूल्यांकने}\label{fol:syn:ass:sec}

\begin{explain}
चल चे मूल्यांकन~$s$ हे \emph{प्रत्येक} चराला मूल्य देते---आणि
चरे अनंत असतात. अर्थात, हे आवश्यक नाही. आपण चल च्या
मूल्यांकनांनी सर्व चले ना मूल्ये नेमावीत अशी अट घालतो, कारण
त्यामुळे गोष्टी खूप सोप्या होतात. पद~$t$ चे मूल्य आणि सूत्र~$\varphi$
ची संरचनांमध्ये~$s$ च्या सापेक्ष पूर्ति होते की नाही, हे फक्त
$t$ मध्ये येणाऱ्या चले ना आणि~$\varphi$ च्या मुक्त चले ना
$s$ ने केलेल्या नेमणुकांवर अवलंबून असते. पुढील दोन विधानांचा आशय हाच
आहे. “अवलंबून असणे” ही कल्पना नेमकी करण्यासाठी आपण दोन गोष्टी दाखवतो:
$t$ मधील सर्व चरांवर एकमत असणारी कोणतीही दोन चर-मूल्यांकने तेच मूल्य
देतात; आणि दोन चर-मूल्यांकने~$\varphi$ च्या सर्व मुक्त चरांवर एकमत असतील, तर
$\varphi$ ची एकाच्या सापेक्ष पूर्ति होते तेव्हा आणि केवळ तेव्हाच तिची
दुसऱ्याच्या सापेक्ष पूर्ति होते.
\end{explain}

\begin{prop}\label{fol:syn:ass:prop:valindep}
पद~$t$ मधील चले ही $x_1$, \dots,~$x_n$ यांपैकी असतील आणि
$i = 1$, \dots,~$n$ साठी $s_1(x_i) = s_2(x_i)$ असेल, तर
$\Value{t}{M}[s_1] = \Value{t}{M}[s_2]$.
\end{prop}

\begin{proof}
$t$ च्या जटिलतेवरील विगमनाने सिद्ध करू. आधार बाबतीत~$t$ हे
स्थिरांक किंवा $x_1$, \dots,~$x_n$ यांपैकी एखादे चर असू शकते.
$t = c$ असेल, तर
$\Value{t}{M}[s_1] = \Assign{c}{M} = \Value{t}{M}[s_2]$. $t = x_i$
असेल, तर विधानाच्या गृहीतकानुसार $s_1(x_i) = s_2(x_i)$; म्हणून
$\Value{t}{M}[s_1] = s_1(x_i) = s_2(x_i) = \Value{t}{M}[s_2]$.

विगमन पायरीसाठी $t = \Atom{f}{t_1, \dots, t_k}$ असून दावा
$t_1$, \dots, $t_k$ साठी खरा आहे असे समजू. मग
\begin{align*}
  \Value{t}{M}[s_1] & = \Value{\Atom{f}{t_1, \dots, t_k}}{M}[s_1] \\
  & = \Assign{f}{M}(\Value{t_1}{M}[s_1], \dots, \Value{t_k}{M}[s_1]).
\intertext{$j = 1$, \dots,~$k$ साठी~$t_j$ मधील चले ही
    $x_1$, \dots,~$x_n$ यांपैकी आहेत. विगमन गृहीतकानुसार
    $\Value{t_j}{M}[s_1] = \Value{t_j}{M}[s_2]$. म्हणून,}
\Value{t}{M}[s_1] & = \Value{\Atom{f}{t_1, \dots, t_k}}{M}[s_1] \\
  & = \Assign{f}{M}(\Value{t_1}{M}[s_1], \dots, \Value{t_k}{M}[s_1]) \\
  & = \Assign{f}{M}(\Value{t_1}{M}[s_2], \dots, \Value{t_k}{M}[s_2]) \\
  & = \Value{\Atom{f}{t_1, \dots, t_k}}{M}[s_2] = \Value{t}{M}[s_2].
\end{align*}
\end{proof}

\begin{prop}\label{fol:syn:ass:prop:satindep}
$\varphi$ मधील मुक्त चले ही $x_1$, \dots,~$x_n$ यांपैकी असतील आणि
$i = 1$, \dots,~$n$ साठी $s_1(x_i) = s_2(x_i)$ असेल, तर
$\Sat{M}{\varphi}[s_1]$ तेव्हा आणि केवळ तेव्हाच, जेव्हा
$\Sat{M}{\varphi}[s_2]$.
\end{prop}

\begin{proof}
$\varphi$ च्या जटिलतेवर विगमन वापरू. आधार बाबतीत~$\varphi$ हे आण्विक असते आणि
$\varphi$ पुढीलपैकी असू शकते:
$\ltrue$,
$\lfalse$,
$k$-स्थानी विधेय~$R$ आणि पदे $t_1$, \dots,~$t_k$ यांसाठी
$\Atom{R}{t_1, \dots, t_k}$; किंवा पदे~$t_1$ व~$t_2$ यांसाठी
$\eq[t_1][t_2]$. शेवटच्या दोन बाबतींत आपण द्विशर्त ची फक्त
पुढची दिशा दाखवतो, कारण उलट दिशेची सिद्धता सममित आहे.

\begin{enumerate}
\item
  \indcase{\varphi}{\ltrue}{$\Sat{M}{\indfrm}[s_1]$ आणि
    $\Sat{M}{\indfrm}[s_2]$ दोन्ही.}

\item
  \indcase{\varphi}{\lfalse}{$\Sat/{M}{\varphi}[s_1]$ आणि
    $\Sat/{M}{\varphi}[s_2]$ दोन्ही.}

\item
  \indcase{\varphi}{\Atom{R}{t_1, \ldots, t_k}}{
    $\Sat{M}{\indfrm}[s_1]$ असे समजू. मग
    \[    \langle \Value{t_1}{M}[s_1], \ldots, \Value{t_k}{M}[s_1] \rangle
    \in \Assign{R}{M}.
    \]
    $i = 1$, \dots,~$k$ साठी \ref{fol:syn:ass:prop:valindep} नुसार
    $\Value{t_i}{M}[s_1] = \Value{t_i}{M}[s_2]$. म्हणून आपल्याकडे
    $\langle \Value{t_1}{M}[s_2], \ldots, \Value{t_k}{M}[s_2] \rangle
    \in \Assign{R}{M}$ हेसुद्धा आहे; आणि म्हणून
    $\Sat{M}{\indfrm}[s_2]$.
    }
\item
  \indcase{\varphi}{\eq[t_1][t_2]}{$\Sat{M}{\indfrm}[s_1]$ असे समजू.
    मग $\Value{t_1}{M}[s_1] = \Value{t_2}{M}[s_1]$. म्हणून,
    \begin{align*}
      \Value{t_1}{M}[s_2] & = \Value{t_1}{M}[s_1]
      & \text{(\ref{fol:syn:ass:prop:valindep} नुसार)} \\
      & = \Value{t_2}{M}[s_1]
      & \text{(कारण $\Sat{M}{\eq[t_1][t_2]}[s_1]$)}\\
      &= \Value{t_2}{M}[s_2]
      & \text{(\ref{fol:syn:ass:prop:valindep} नुसार),}
    \end{align*}
    म्हणून $\Sat{M}{\eq[t_1][t_2]}[s_2]$.}
\end{enumerate}

आता~$\varphi$ पेक्षा कमी जटिल असलेल्या सर्व सूत्रे~$\psi$ साठी
$\Sat{M}{\psi}[s_1]$ तेव्हा आणि केवळ तेव्हाच, जेव्हा
$\Sat{M}{\psi}[s_2]$, असे समजू. विगमन पायरी~$\varphi$ च्या मुख्य कारकाने
ठरणाऱ्या बाबींनुसार पुढे जाते. प्रत्येक बाबतीत आपण द्विशर्त ची
फक्त पुढची दिशा दाखवतो; उलट दिशेची सिद्धता सममित आहे. संख्यापकांच्या
बाबी सोडून इतर सर्व बाबतींत आपण~$\varphi$ च्या उप-सूत्रे~$\psi$ ना
विगमन गृहीतक लावतो. $\psi$ ची मुक्त चरे ही~$\varphi$ च्या मुक्त चरांपैकी
असतात. त्यामुळे $s_1$ व~$s_2$ ही~$\varphi$ च्या मुक्त चरांवर एकमत असतील,
तर ती~$\psi$ च्या मुक्त चरांवरही एकमत असतात आणि विगमन गृहीतक~$\psi$ ला
लागू होते.

\begin{enumerate}
\item

    \indcase{\varphi}{\lnot \psi}{$\Sat{M}{\indfrm}[s_1]$ असेल, तर
      $\Sat/{M}{\psi}[s_1]$; म्हणून विगमन गृहीतकानुसार
      $\Sat/{M}{\psi}[s_2]$ आणि त्यामुळे $\Sat{M}{\indfrm}[s_2]$.}

\item

    \indcase{\varphi}{\psi \land \chi}{$\Sat{M}{\indfrm}[s_1]$ असेल, तर
      $\Sat{M}{\psi}[s_1]$ आणि $\Sat{M}{\chi}[s_1]$; म्हणून विगमन
      गृहीतकानुसार $\Sat{M}{\psi}[s_2]$ आणि $\Sat{M}{\chi}[s_2]$.
      त्यामुळे $\Sat{M}{\indfrm}[s_2]$.}

\item

    \indcase{\varphi}{\psi \lor \chi}{$\Sat{M}{\indfrm}[s_1]$ असेल, तर
      $\Sat{M}{\psi}[s_1]$ किंवा $\Sat{M}{\chi}[s_1]$. विगमन गृहीतकानुसार
      $\Sat{M}{\psi}[s_2]$ किंवा $\Sat{M}{\chi}[s_2]$; म्हणून
      $\Sat{M}{\indfrm}[s_2]$.}

\item

    \indcase{\varphi}{\psi \lif \chi}{$\Sat{M}{\indfrm}[s_1]$ असेल, तर
    $\Sat/{M}{\psi}[s_1]$ किंवा $\Sat{M}{\chi}[s_1]$. विगमन गृहीतकानुसार
    $\Sat/{M}{\psi}[s_2]$ किंवा $\Sat{M}{\chi}[s_2]$; म्हणून
    $\Sat{M}{\indfrm}[s_2]$.}

\item

    \indcase{\varphi}{\psi \liff \chi}{$\Sat{M}{\indfrm}[s_1]$ असेल, तर एकतर
      $\Sat{M}{\psi}[s_1]$ आणि $\Sat{M}{\chi}[s_1]$, किंवा
      $\Sat/{M}{\psi}[s_1]$ आणि $\Sat/{M}{\chi}[s_1]$. विगमन गृहीतकानुसार
      एकतर $\Sat{M}{\psi}[s_2]$ आणि $\Sat{M}{\chi}[s_2]$, किंवा
      $\Sat/{M}{\psi}[s_2]$ आणि $\Sat/{M}{\chi}[s_2]$. दोन्हींपैकी
      कोणतीही बाब असली, तरी $\Sat{M}{\indfrm}[s_2]$.}

\item

    \indcase{\varphi}{\lexists[x][\psi]}{$\Sat{M}{\indfrm}[s_1]$ असेल, तर असा
      एक $m \in \Domain{M}$ असतो की $\Sat{M}{\psi}[\Subst{s_1}{m}{x}]$.
      $s_1' = \Subst{s_1}{m}{x}$ आणि $s_2' =
      \Subst{s_2}{m}{x}$ असू द्या. $\psi$ ची मुक्त चरे
      $x_1$, \dots, $x_n$ आणि~$x$ यांपैकी आहेत. $s_1'$ आणि~$s_2'$ हे
      अनुक्रमे $s_1$ व~$s_2$ यांचे $x$-पर्याय असल्यामुळे आणि गृहीतकानुसार
      $s_1(x_i) = s_2(x_i)$ असल्यामुळे $s_1'(x_i) = s_2'(x_i)$.
      $s_1'$ आणि~$s_2'$ आपण ज्या रीतीने परिभाषित केली आहे, तिच्यामुळे
      $s_1'(x) = s_2'(x) = m$. मग~$\psi$ आणि $s_1'$, $s_2'$ यांना विगमन
      गृहीतक लागू होते; म्हणून $\Sat{M}{\psi}[s_2']$. $s_2' =
      \Subst{s_2}{m}{x}$ असल्यामुळे असा एक $m \in \Domain{M}$ आहे की
      $\Sat{M}{\psi}[\Subst{s_2}{m}{x}]$; म्हणून
      $\Sat{M}{\indfrm}[s_2]$.}

\item

    \indcase{\varphi}{\lforall[x][\psi]}{$\Sat{M}{\indfrm}[s_1]$ असेल, तर
      प्रत्येक $m \in \Domain{M}$ साठी
      $\Sat{M}{\psi}[\Subst{s_1}{m}{x}]$. प्रत्येक
      $m \in \Domain{M}$ साठी $\Sat{M}{\psi}[\Subst{s_2}{m}{x}]$
      हेसुद्धा दाखवायचे आहे. म्हणून मनमाना $m \in \Domain{M}$ घेऊन
      $s_1' = \Subst{s_1}{m}{x}$ आणि $s_2' =
      \Subst{s_2}{m}{x}$ विचारात घ्या. आपल्याकडे
      $\Sat{M}{\psi}[s_1']$. $\psi$ ची मुक्त चरे $x_1$, \dots, $x_n$
      आणि~$x$ यांपैकी आहेत. $s_1'$ आणि~$s_2'$ हे अनुक्रमे $s_1$
      व~$s_2$ यांचे $x$-पर्याय असल्यामुळे आणि गृहीतकानुसार
      $s_1(x_i) = s_2(x_i)$ असल्यामुळे $s_1'(x_i) = s_2'(x_i)$.
      $s_1'$ आणि~$s_2'$ आपण ज्या रीतीने परिभाषित केली आहे, तिच्यामुळे
      $s_1'(x) = s_2'(x) = m$. मग~$\psi$ आणि~$s_1'$, $s_2'$ यांना विगमन
      गृहीतक लागू होते आणि आपल्याकडे $\Sat{M}{\psi}[s_2']$. हे प्रत्येक
      $m \in \Domain{M}$ ला लागू होते, म्हणजे सर्व~$m \in \Domain{M}$
      साठी $\Sat{M}{\psi}[\Subst{s_2}{m}{x}]$; म्हणून
      $\Sat{M}{\indfrm}[s_2]$.
      }
\end{enumerate}
विगमनाने आपल्याला पुढील निष्कर्ष मिळतो: $\varphi$ मधील मुक्त चले
ही $x_1$, \dots, $x_n$ यांपैकी असतील आणि $i = 1$, \dots,~$n$ साठी
$s_1(x_i)=s_2(x_i)$ असेल, तेव्हा $\Sat{M}{\varphi}[s_1]$ तेव्हा आणि केवळ
तेव्हाच, जेव्हा $\Sat{M}{\varphi}[s_2]$.
\end{proof}



\begin{explain}
वाक्ये मध्ये मुक्त चरे नसतात; म्हणून कोणतीही दोन चर-मूल्यांकने
कोणत्याही वाक्याच्या सर्व (शून्य) मुक्त चरांना त्याच गोष्टी नेमून देतात.
आत्ताच सिद्ध केलेल्या विधानाचा अर्थ मग असा होतो की वाक्याची
एखाद्या रचनेत चर-मूल्यांकनाच्या सापेक्ष पूर्ति होते की नाही, हे पूर्णपणे
त्या मूल्यांकनापासून स्वतंत्र असते. हे तथ्य आपण नोंदवू. पुढे येणारी
वाक्याची संरचना मधील पूर्तिची व्याख्या (चर-मूल्यांकनाचा
उल्लेख न करता) त्यामुळे समर्थित होते.
\end{explain}

\begin{cor}
\label{fol:syn:ass:cor:sat-sentence}
$\varphi$ हे वाक्य आणि $s$ हे चर-मूल्यांकन असेल, तर प्रत्येक
चर-मूल्यांकन~$s'$ साठी $\Sat{M}{\varphi}[s]$ तेव्हा आणि केवळ तेव्हाच,
जेव्हा $\Sat{M}{\varphi}[s']$.
\end{cor}

\begin{proof}
  $s'$ हे कोणतेही चर-मूल्यांकन असू द्या. $\varphi$ हे वाक्य
  असल्यामुळे त्यात मुक्त चरे नाहीत; म्हणून प्रत्येक चर-मूल्यांकन~$s'$
  हे~$\varphi$ च्या सर्व मुक्त चरांना~$s$ सारख्याच गोष्टी आपोआप नेमून देते.
  त्यामुळे \ref{fol:syn:ass:prop:satindep} ची अट पूर्ण होते आणि
  $\Sat{M}{\varphi}[s]$ तेव्हा आणि केवळ तेव्हाच, जेव्हा
  $\Sat{M}{\varphi}[s']$.
\end{proof}

\begin{defn}
\label{fol:syn:ass:defn:satisfaction}
$\varphi$ हे वाक्य असेल, तर संरचना~$\Struct M$ ही~$\varphi$ ची
\emph{पूर्ति करते}, $\Sat{M}{\varphi}$, असे आपण तेव्हा आणि केवळ तेव्हाच
म्हणतो, जेव्हा सर्व चर-मूल्यांकने~$s$ साठी $\Sat{M}{\varphi}[s]$.
\end{defn}

$\Sat{M}{\varphi}$ असेल, तर आपण सरळ \emph{$\varphi$ हे~$\Struct{M}$ मध्ये सत्य
आहे} असेही म्हणतो. पूर्तिची संकल्पना स्वतंत्र वाक्यांपासून
वाक्यांच्या संचांपर्यंत स्वाभाविकपणे विस्तारते.

\begin{defn}
\label{fol:syn:ass:defn:sat}
$\Gamma$ हा वाक्यांचा संच असेल, तर संरचना~$\Struct M$
ही~$\Gamma$ ची \emph{पूर्ति करते}, $\Sat{M}{\Gamma}$, असे आपण तेव्हा
आणि केवळ तेव्हाच म्हणतो, जेव्हा प्रत्येक $\varphi \in \Gamma$ साठी
$\Sat{M}{\varphi}$.

\end{defn}

\begin{prop}\label{fol:syn:ass:prop:sentence-sat-true}
  $\Struct{M}$ ही संरचना, $\varphi$ हे वाक्य आणि $s$ हे
  चर-मूल्यांकन असू द्या. $\Sat{M}{\varphi}$ तेव्हा आणि केवळ तेव्हाच,
  जेव्हा $\Sat{M}{\varphi}[s]$.
\end{prop}

\begin{proof}
सराव.
\end{proof}

\begin{prob}
\ref{fol:syn:ass:prop:sentence-sat-true} सिद्ध करा.
\end{prob}

\begin{prop}\label{fol:syn:ass:prop:sat-quant}
  $\varphi(x)$ मध्ये फक्त~$x$ मुक्त येतो आणि $\Struct M$ ही
  संरचना आहे असे समजा. मग:
  \begin{enumerate}
    \item $\Sat{M}{\lexists[x][\varphi(x)]}$ तेव्हा आणि केवळ
      तेव्हाच, जेव्हा किमान एका चर-मूल्यांकन~$s$ साठी
      $\Sat{M}{\varphi(x)}[s]$.
    \item $\Sat{M}{\lforall[x][\varphi(x)]}$ तेव्हा आणि केवळ
      तेव्हाच, जेव्हा सर्व चर-मूल्यांकने~$s$ साठी
      $\Sat{M}{\varphi(x)}[s]$.
\end{enumerate}
\end{prop}

\begin{proof}
सराव.
\end{proof}

\begin{prob}
\ref{fol:syn:ass:prop:sat-quant} सिद्ध करा.
\end{prob}

\begin{prob}
\DeclareRobustCommand{\VDash}{\mathrel{||}\joinrel\Relbar}
$\Lang L$ ही फलने नसलेली भाषा आहे असे समजा. संरचना~$\Struct M$,
स्थिरांक~$c$ आणि $a \in \Domain M$ दिले असता, $\Struct M[a/c]$
ही~$\Struct M$ सारखीच संरचना आहे, फक्त
$\Assign{c}{M[a/c]} = a$, अशी व्याख्या करा. वाक्ये~$\varphi$ साठी
$\Struct M \VDash \varphi$ ची पुढीलप्रमाणे व्याख्या करा:
\begin{enumerate}
\item
  \indcase{\varphi}{\lfalse}{$\Struct{M} \VDash \indfrm$ नाही.}

\item
  \indcase{\varphi}{\ltrue}{$\Struct{M} \VDash \indfrm$.}

\item \indcase{\varphi}{\Atom{R}{d_1, \dots, d_n}}{$\Struct M \VDash \indfrm$
  तेव्हा आणि केवळ तेव्हाच, जेव्हा
  $\langle \Assign{d_1}{M}, \dots, \Assign{d_n}{M} \rangle \in
  \Assign{R}{M}$.}

\item \indcase{\varphi}{\eq[d_1][d_2]}{$\Struct M \VDash \indfrm$ तेव्हा
  आणि केवळ तेव्हाच, जेव्हा
  $\Assign{d_1}{M} = \Assign{d_2}{M}$.}

\item
  \indcase{\varphi}{\lnot \psi}{$\Struct{M} \VDash \indfrm$ तेव्हा आणि केवळ
    तेव्हाच, जेव्हा $\Struct{M} \VDash {\psi}$ नाही.}

\item
  \indcase{\varphi}{(\psi \land \chi)}{$\Struct{M} \VDash
    \indfrm$ तेव्हा आणि केवळ तेव्हाच, जेव्हा
    $\Struct{M} \VDash\psi$ आणि $\Struct{M} \VDash \chi$.}

\item
  \indcase{\varphi}{(\psi \lor \chi)}{$\Struct{M} \VDash \indfrm$ तेव्हा आणि
    केवळ तेव्हाच, जेव्हा $\Struct{M} \VDash \psi$ किंवा
    $\Struct{M} \VDash \chi$ (किंवा दोन्ही).}

\item
  \indcase{\varphi}{(\psi \lif \chi)}{$\Struct{M} \VDash \indfrm$ तेव्हा आणि
    केवळ तेव्हाच, जेव्हा $\Struct {M} \VDash \psi$ नाही किंवा
    $\Struct M \VDash \chi$ (किंवा दोन्ही).}

\item
  \indcase{\varphi}{(\psi \liff \chi)}{$\Struct{M} \VDash {\indfrm}$ तेव्हा आणि
    केवळ तेव्हाच, जेव्हा एकतर $\Struct{M} \VDash {\psi}$ आणि
    $\Struct{M} \VDash {\chi}$ दोन्ही, किंवा $\Struct{M} \VDash {\psi}$ व
    $\Struct{M} \VDash {\chi}$ यांपैकी एकही नाही.}

\item
  \indcase{\varphi}{\lforall[x][\psi]}{$\Struct{M} \VDash {\indfrm}$ तेव्हा आणि
    केवळ तेव्हाच, जेव्हा प्रत्येक $a \in \Domain{M}$ साठी
    $\Struct{M[a/c]} \VDash \Subst{\psi}{c}{x}$, जर~$\psi$ मध्ये~$c$ येत
    नसेल.}

\item
  \indcase{\varphi}{\lexists[x][\psi]}{$\Struct{M} \VDash
  {\indfrm}$ तेव्हा आणि केवळ तेव्हाच, जेव्हा असा एक $a \in \Domain M$
  आहे की $\Struct{M[a/c]} \VDash \Subst{\psi}{c}{x}$, जर~$\psi$ मध्ये~$c$
  येत नसेल.}
\end{enumerate}
$x_1$, \dots, $x_n$ ही~$\varphi$ मधील सर्व मुक्त चले,
$c_1$, \dots, $c_n$ ही~$\varphi$ मध्ये नसलेली स्थिरचिन्हे,
$a_1$, \dots, $a_n \in \Domain M$ आणि $s(x_i) = a_i$ असू द्या.

$\Sat{M}{\varphi}[s]$ तेव्हा आणि केवळ तेव्हाच, जेव्हा
$\Struct M[a_1/c_1,\dots,a_n/c_n]
\VDash \Subst{\Subst{\varphi}{c_1}{x_1}\dots}{c_n}{x_n}$, हे दाखवा.

(या प्रश्नावरून असे दिसते की प्रथम-क्रम तर्कशास्त्राची चिन्हार्थमीमांसा
चर-मूल्यांकनांविना देणे शक्य आहे.)
\end{prob}

\begin{prob}
$f$ हे~$\varphi(x,y)$ मध्ये नसलेले फलनचिन्ह आहे असे समजा. अशी
संरचना~$\Struct{M}$ आहे की
$\Sat{M}{\lforall[x][\lexists[y][\varphi(x,y)]]}$, तेव्हा आणि केवळ तेव्हाच
अशी~$\Struct M'$ आहे की
$\Sat{M'}{\lforall[x][\varphi(x,f(x))]}$, हे दाखवा.

(हा प्रश्न स्कोलेमचे प्रमेय म्हणून परिचित असलेल्या प्रमेयाची एक विशेष
बाब आहे; $\lforall[x][\varphi(x,f(x))]$ याला
$\lforall[x][\lexists[y][\varphi(x,y)]]$ चे \emph{स्कोलेम सामान्य रूप}
म्हणतात.)
\end{prob}












\section{विस्तारात्मकता}\label{fol:syn:ext:sec}

\begin{explain}
विस्तारात्मकता, जिला कधीकधी संबद्धता म्हणतात, अनौपचारिकपणे पुढीलप्रमाणे
व्यक्त करता येते: संरचना~$\Struct M$ मध्ये चल च्या
मूल्यांकन~$s$ च्या सापेक्ष सूत्र~$\varphi$ ची पूर्ति होण्यावर परिणाम
करणारे घटक म्हणजे फक्त प्रांताचा आकार आणि~$\varphi$ मध्ये प्रत्यक्ष
येणाऱ्या भाषेच्या घटकांना~$\Struct M$ व~$s$ यांनी केलेल्या नेमणुका.

विस्तारात्मकतेचा एक तात्काळ परिणाम असा आहे: दोन संरचना
$\Struct M$ आणि~$\Struct M'$ यांचा प्रांत एकच असेल आणि वाक्य~$\varphi$ मध्ये
येणाऱ्या भाषेच्या सर्व घटकांबाबत त्या एकमत असतील, तर~$\varphi$ स्वतः सत्य
आहे की नाही याबाबतही $\Struct M$ आणि~$\Struct M'$ एकमत असले पाहिजेत.
\end{explain}

\begin{prop}[विस्तारात्मकता]
  \label{fol:syn:ext:prop:extensionality}
  $\varphi$ हे सूत्र, $\Struct M_1$ व~$\Struct M_2$ या
  $\Domain{M_1} = \Domain{M_2}$ असलेल्या संरचना, आणि~$s$ हे
  $\Domain{M_1} = \Domain{M_2}$ वरील चर-मूल्यांकन असू द्या. $\varphi$ मध्ये
  येणारे प्रत्येक स्थिरांक~$c$, संबंधचिन्ह~$R$ आणि फलन~$f$
  यांसाठी $\Assign{c}{M_1} = \Assign{c}{M_2}$,
  $\Assign{R}{M_1}=\Assign{R}{M_2}$ आणि
  $\Assign{f}{M_1} = \Assign{f}{M_2}$ असेल, तर
  $\Sat{M_1}{\varphi}[s]$ तेव्हा आणि केवळ तेव्हाच, जेव्हा
  $\Sat{M_2}{\varphi}[s]$.
\end{prop}

\begin{proof}
  प्रथम (पद~$t$ वरील विगमनाने) प्रत्येक पदासाठी
  $\Value{t}{M_1}[s] = \Value{t}{M_2}[s]$ हे सिद्ध करा. मग नुकताच
  सिद्ध केलेला दावा विगमनाच्या आधार बाबतीत ($\varphi$ आण्विक असताना) वापरून,
  $\varphi$ वरील विगमनाने विधान सिद्ध करा.
\end{proof}

\begin{prob}
\ref{fol:syn:ext:prop:extensionality} ची सिद्धता सविस्तर पूर्ण करा.
\end{prob}

\begin{cor}[वाक्ये साठी विस्तारात्मकता]
  \label{fol:syn:ext:cor:extensionality-sent}
  $\varphi$ हे वाक्य आणि $\Struct{M_1}$, $\Struct{M_2}$ या
  \ref{fol:syn:ext:prop:extensionality} मधीलप्रमाणे असू द्या. मग
  $\Sat{M_1}{\varphi}$ तेव्हा आणि केवळ तेव्हाच, जेव्हा $\Sat{M_2}{\varphi}$.
\end{cor}

\begin{proof}
\ref{fol:syn:ext:prop:extensionality} आणि \ref{fol:syn:ass:cor:sat-sentence} यांवरून
हे निष्पन्न होते.
\end{proof}

शिवाय, पदाचे मूल्य आणि संरचना ही सूत्राची पूर्ति करते की
नाही, हे फक्त तिच्या उपपदांच्या मूल्यांवर अवलंबून असते.

\begin{prop}\label{fol:syn:ext:prop:ext-terms}
$\Struct M$ ही संरचना, $t$ व~$t'$ ही पदे आणि $s$ हे
चर-मूल्यांकन असू द्या. मग $\Value{\Subst{t}{t'}{x}}{M}[s] =
\Value{t}{M}[\Subst{s}{\Value{t'}{M}[s]}{x}]$.
\end{prop}

\begin{proof}
$t$ वरील विगमनाने सिद्ध करू.
\begin{enumerate}
\item $t$ हे स्थिरचिन्ह, म्हणजे $t\ident c$, असेल, तर
  $\Subst{t}{t'}{x} = c$ आणि $\Value{c}{M}[s] = \Assign{c}{M} =
  \Value{c}{M}[\Subst{s}{\Value{t'}{M}[s]}{x}]$.

\item $t$ हे~$x$ शिवाय दुसरे चर, म्हणजे $t \ident y$, असेल, तर
  $\Subst{t}{t'}{x} = y$ आणि
  $\varAssign{s}{\Subst{s}{\Value{t'}{M}[s]}{x}}{x}$ असल्यामुळे
  $\Value{y}{M}[s] =
  \Value{y}{M}[\Subst{s}{\Value{t'}{M}[s]}{x}]$.

\item $t \ident x$ असेल, तर $\Subst{t}{t'}{x} = t'$. पण
  $\Subst{s}{\Value{t'}{M}[s]}{x}$ च्या व्याख्येनुसार
  $\Value{x}{M}[\Subst{s}{\Value{t'}{M}[s]}{x}] = \Value{t'}{M}[s]$.

\item $t \ident \Atom{f}{t_1,\dots,t_n}$ असेल, तर:
\begin{multline*}
  \Value{\Subst{t}{t'}{x}}{M}[s] \\
  \begin{aligned}[b]
& = \Value{\Atom{f}{\Subst{t_1}{t'}{x}, \dots, \Subst{t_n}{t'}{x}}}{M}[s]\\
& \qquad    \text{$\Subst{t}{t'}{x}$ च्या व्याख्येनुसार}\\
& = \Assign{f}{M}(\Value{\Subst{t_1}{t'}{x}}{M}[s], \dots,
    \Value{\Subst{t_n}{t'}{x}}{M}[s])\\
    & \qquad  \text{$\Value{\Atom{f}{\dots}}{M}[s]$ च्या व्याख्येनुसार}\\
& = \Assign{f}{M}(\Value{t_1}{M}[\Subst{s}{\Value{t'}{M}[s]}{x}], \dots,
   \Value{t_n}{M}[\Subst{s}{\Value{t'}{M}[s]}{x}])\\
& \qquad    \text{विगमन गृहीतकानुसार}\\
& = \Value{t}{M}[\Subst{s}{\Value{t'}{M}[s]}{x}]
    \text{$\Value{\Atom{f}{\dots}}{M}[\Subst{s}{\Value{t'}{M}[s]}{x}]$ च्या व्याख्येनुसार}
  \end{aligned}
\end{multline*}

\end{enumerate}
\end{proof}

\begin{prop}\label{fol:syn:ext:prop:ext-formulas} $\Struct M$ ही
संरचना, $\varphi$ हे सूत्र, $t'$ हे~A मध्ये~x साठी
आदेशनासाठी मुक्त असलेले पद आणि $s$ हे चर-मूल्यांकन असू द्या. मग
$\Sat{M}{\Subst{\varphi}{t'}{x}}[s]$ तेव्हा आणि केवळ तेव्हाच, जेव्हा
$\Sat{M}{\varphi}[\Subst{s}{\Value{t'}{M}[s]}{x}]$.

\end{prop}

\begin{proof}
सराव.
\end{proof}

\begin{prob}
\ref{fol:syn:ext:prop:ext-formulas} सिद्ध करा.
\end{prob}

\begin{explain}
  \ref{fol:syn:ext:prop:ext-terms} आणि \ref{fol:syn:ext:prop:ext-formulas} यांचा
  आशय पुढीलप्रमाणे आहे. आपल्याजवळ पद~$t$ किंवा सूत्र~$\varphi$ आणि
  दुसरे पद~$t'$ असून, $\Subst{t}{t'}{x}$ चे मूल्य किंवा
  $\Subst{\varphi}{t'}{x}$ ची संरचना~$\Struct M$ मध्ये
  चल च्या मूल्यांकन~$s$ च्या सापेक्ष पूर्ति होते की नाही हे
  जाणून घ्यायचे आहे असे समजा. मग आपण प्रथम आदेशन करून नंतर
  $\Struct{M}$ आणि~$s$ यांच्या सापेक्ष मूल्य किंवा पूर्ति विचारात घेऊ
  शकतो; किंवा प्रथम~$\Struct{M}$ मध्ये~$s$ च्या सापेक्ष~$t'$ चे
  मूल्य~$m = \Value{t'}{M}[s]$ ठरवून, चल चे मूल्यांकन
  बदलून~$\Subst{s}{m}{x}$ करू शकतो, आणि नंतर~$\Struct{M}$ व
  $\Subst{s}{m}{x}$ यांच्या सापेक्ष~$t$ चे मूल्य किंवा
  $\Sat{M}{\varphi}[\Subst{s}{m}{x}]$ आहे की नाही हे विचारात घेऊ शकतो.
  आपण यांपैकी कोणतीही पद्धत वापरली, तरी उत्तर एकच असेल, याची हमी
  \ref{fol:syn:ext:prop:ext-terms} आणि \ref{fol:syn:ext:prop:ext-formulas}
  देतात.
\end{explain}











\section{चिन्हार्थविषयक संकल्पना}\label{fol:syn:sem:sec}

\begin{explain}
प्रथम-क्रम भाषांसाठीच्या संरचनेची व्याख्या मिळाल्यावर वाक्यांचे
काही मूलभूत चिन्हार्थविषयक गुणधर्म आणि वाक्यांमधील संबंध परिभाषित करता
येतात. त्यांतील सर्वांत सोपी संकल्पना म्हणजे वाक्याची \emph{वैधता}.
प्रत्येक संरचनांमध्ये एखाद्या वाक्याची पूर्ति होत असेल, तर ते
वाक्य वैध असते. त्यातील अतार्किक चिन्हांचा अर्थ कसा लावला आहे याचा
विचार न करता ज्या वाक्यांची पूर्ति होते, तीच वैध वाक्ये होत. म्हणून वैध
वाक्यांना \emph{तार्किक सत्ये} असेही म्हणतात---ती कोणत्याही
संरचनांमध्ये सत्य, म्हणजे पूर्त, असतात. त्यामुळे त्यांची सत्यता
फक्त त्यांत येणारी तार्किक चिन्हे आणि त्यांची विन्यासात्मक संरचना
यांवर अवलंबून असते; अतार्किक चिन्हे किंवा त्यांचे अर्थनिर्धारण यांवर नाही.
\end{explain}

\begin{defn}[वैधता]
वाक्य~$\varphi$ हे \emph{वैध}, $\Entails \varphi$, तेव्हा आणि केवळ तेव्हाच असते,
जेव्हा प्रत्येक संरचना~$\Struct M$ साठी $\Sat{M}{\varphi}$.
\end{defn}

\begin{defn}[तार्किक निष्पन्नता]
वाक्यांच्या संच~$\Gamma$ मधून वाक्य~$\varphi$
\emph{तार्किकरीत्या निष्पन्न होते}, $\Gamma \Entails \varphi$, तेव्हा आणि
केवळ तेव्हाच, जेव्हा $\Sat{M}{\Gamma}$ असलेल्या प्रत्येक
संरचना~$\Struct M$ साठी $\Sat{M}{\varphi}$.
\end{defn}


\begin{defn}[पूर्ततायोग्यता]
एखाद्या संरचना~$\Struct M$ साठी $\Sat{M}{\Gamma}$ असेल, तर
वाक्यांचा संच~$\Gamma$ \emph{पूर्ततायोग्य} असतो. $\Gamma$
पूर्ततायोग्य नसेल, तर त्याला \emph{अपूर्ततायोग्य} म्हणतात.
\end{defn}

\begin{prop}
वाक्य~$\varphi$ वैध असते तेव्हा आणि केवळ तेव्हाच, जेव्हा
वाक्यांच्या प्रत्येक संच~$\Gamma$ साठी $\Gamma \Entails \varphi$.
\end{prop}

\begin{proof}
अग्र दिशेसाठी, $\varphi$ वैध आणि $\Gamma$ हा वाक्यांचा संच असू द्या.
$\Sat{M}{\Gamma}$ अशी संरचना~$\Struct M$ असू द्या. $\varphi$ वैध
असल्यामुळे $\Sat{M}{\varphi}$; म्हणून $\Gamma \Entails \varphi$.

उलट दिशेच्या व्यत्यासासाठी, $\varphi$ अवैध असू द्या; म्हणजे
$\Sat/{M}{\varphi}$ अशी संरचना~$\Struct M$ आहे. $\Gamma =
\{ \ltrue \}$ असताना $\ltrue$ वैध असल्यामुळे $\Sat{M}{\Gamma}$.
त्यामुळे अशी संरचना~$\Struct M$ आहे की $\Sat{M}{\Gamma}$ पण
$\Sat/{M}{\varphi}$; म्हणून $\Gamma$ मधून~$\varphi$ तार्किकरीत्या निष्पन्न होत
नाही.
\end{proof}

\begin{prop}
\label{fol:syn:sem:prop:entails-unsat}
$\Gamma \Entails \varphi$ तेव्हा आणि केवळ तेव्हाच, जेव्हा
$\Gamma \cup \{\lnot \varphi\}$ अपूर्ततायोग्य असतो.
\end{prop}

\begin{proof}
अग्र दिशेसाठी, $\Gamma \Entails \varphi$ असे समजू आणि, याउलट, अशी
संरचना~$\Struct M$ आहे की $\Sat{M}{\Gamma
  \cup \{ \lnot \varphi \}}$, असे समजू. $\Sat{M}{\Gamma}$ आणि
$\Gamma \Entails \varphi$ असल्यामुळे $\Sat{M}{\varphi}$. तसेच
$\Sat{M}{\Gamma\cup \{ \lnot \varphi \}}$ असल्यामुळे
$\Sat{M}{\lnot \varphi}$; म्हणून $\Sat{M}{\varphi}$ आणि $\Sat/{M}{\varphi}$ दोन्ही
मिळतात, हा विरोधाभास आहे. त्यामुळे अशी कोणतीही
संरचना~$\Struct M$ असू शकत नाही; म्हणून
$\Gamma \cup \{ \lnot \varphi \}$ अपूर्ततायोग्य आहे.

उलट दिशेसाठी, $\Gamma \cup \{ \lnot \varphi \}$ अपूर्ततायोग्य आहे असे
समजू. म्हणून प्रत्येक संरचना~$\Struct M$ साठी एकतर
$\Sat/{M}{\Gamma}$ किंवा $\Sat{M}{\varphi}$. त्यामुळे
$\Sat{M}{\Gamma}$ असलेल्या प्रत्येक संरचना~$\Struct M$ साठी
$\Sat{M}{\varphi}$; म्हणून $\Gamma \Entails \varphi$.
\end{proof}

\begin{prob}
\begin{enumerate}
\item $\Gamma \Entails \bot$ तेव्हा आणि केवळ तेव्हाच, जेव्हा $\Gamma$
  अपूर्ततायोग्य आहे, हे दाखवा.
\item $\Gamma \cup \{\varphi\} \Entails \bot$ तेव्हा आणि केवळ तेव्हाच,
  जेव्हा $\Gamma \Entails \lnot \varphi$, हे दाखवा.
\item $\varphi$ किंवा $\Gamma$ मध्ये~$c$ येत नाही असे समजा.
  $\Gamma \Entails \lforall[x][\varphi]$ तेव्हा आणि केवळ तेव्हाच, जेव्हा
  $\Gamma \Entails \Subst{\varphi}{c}{x}$, हे दाखवा.
\end{enumerate}
\end{prob}

\begin{prop}
$\Gamma \subseteq \Gamma'$ आणि $\Gamma \Entails \varphi$ असेल, तर
$\Gamma' \Entails \varphi$.
\end{prop}

\begin{proof}
$\Gamma \subseteq \Gamma'$ आणि $\Gamma \Entails \varphi$ असे समजा.
$\Sat{M}{\Gamma'}$ अशी रचना~$\Struct M$ असू द्या; मग
$\Sat{M}{\Gamma}$ आणि $\Gamma \Entails \varphi$ असल्यामुळे
$\Sat{M}{\varphi}$. म्हणून जेव्हा जेव्हा $\Sat{M}{\Gamma'}$, तेव्हा
$\Sat{M}{\varphi}$; त्यामुळे $\Gamma' \Entails \varphi$.
\end{proof}


\begin{thm}[चिन्हार्थविषयक निगमन प्रमेय]
\label{fol:syn:sem:thm:sem-deduction}
$\Gamma \cup \{\varphi\} \Entails \psi$ तेव्हा आणि केवळ तेव्हाच, जेव्हा
$\Gamma \Entails \varphi \lif \psi$.
\end{thm}

\begin{proof}
अग्र दिशेसाठी, $\Gamma \cup \{ \varphi \} \Entails \psi$ असू द्या आणि
$\Sat{M}{\Gamma}$ अशी संरचना~$\Struct M$ असू द्या.
$\Sat{M}{\varphi}$ असेल, तर $\Sat{M}{\Gamma \cup \{ \varphi \} }$; आणि
$\Gamma \cup \{ \varphi \}$ मधून~$\psi$ तार्किकरीत्या निष्पन्न होत असल्यामुळे
$\Sat{M}{\psi}$. म्हणून $\Sat{M}{\varphi \lif \psi}$; त्यामुळे
$\Gamma \Entails \varphi \lif \psi$.

उलट दिशेसाठी, $\Gamma \Entails \varphi \lif \psi$ असू द्या आणि
$\Sat{M}{\Gamma \cup \{ \varphi \}}$ अशी संरचना~$\Struct M$ असू
द्या. मग $\Sat{M}{\Gamma}$; म्हणून $\Sat{M}{\varphi \lif \psi}$ आणि
$\Sat{M}{\varphi}$ असल्यामुळे $\Sat{M}{\psi}$. त्यामुळे जेव्हा जेव्हा
$\Sat{M}{\Gamma \cup \{ \varphi \} }$, तेव्हा $\Sat{M}{\psi}$; म्हणून
$\Gamma \cup \{ \varphi \} \Entails \psi$.
\end{proof}

\begin{prop}\label{fol:syn:sem:prop:quant-terms}
  $\Struct{M}$ ही संरचना, $\varphi(x)$ हे एक मुक्त चर~$x$ असलेले
  सूत्र आणि $t$ हे बंद पद असू द्या. मग:
  \begin{enumerate}
    \item $\varphi(t) \Entails \lexists[x][\varphi(x)]$
    \item $\lforall[x][\varphi(x)] \Entails \varphi(t)$
  \end{enumerate}
\end{prop}

\begin{proof}
  \begin{enumerate}
    \item
      $\Sat{M}{\varphi(t)}$ असे समजा. $s(x) =
        \Value{t}{M}$ असे चर-मूल्यांकन~$s$ असू द्या. $\varphi(t)$ हे
        वाक्य असल्यामुळे $\Sat{M}{\varphi(t)}[s]$.
        \ref{fol:syn:ext:prop:ext-formulas} नुसार
        $\Sat{M}{\varphi(x)}[s]$. \ref{fol:syn:ass:prop:sat-quant} नुसार
        $\Sat{M}{\lexists[x][\varphi(x)]}$.
    \item
      $\Sat{M}{\lforall[x][\varphi(x)]}$ असे समजा.
        $s(x) = \Value{t}{M}$ असे चर-मूल्यांकन~$s$ असू द्या.
        \ref{fol:syn:ass:prop:sat-quant} नुसार $\Sat{M}{\varphi(x)}[s]$.
        \ref{fol:syn:ext:prop:ext-formulas} नुसार
        $\Sat{M}{\varphi(t)}[s]$. $\varphi(t)$ हे वाक्य असल्यामुळे
        \ref{fol:syn:ass:prop:sentence-sat-true} नुसार
        $\Sat{M}{\varphi(t)}$.
  \end{enumerate}
\end{proof}





\chapter{उपपत्ती आणि त्यांची प्रतिमाने}\label{fol:mat::chap}








\section{परिचय}\label{fol:mat:int:sec}

\begin{explain}
स्वयंसिद्धकीय पद्धतीचा विकास हा विज्ञानाच्या इतिहासातील एक महत्त्वपूर्ण
पराक्रम आहे आणि गणिताच्या इतिहासात त्याला विशेष महत्त्व आहे. एखाद्या
क्षेत्राचा स्वयंसिद्धकीय विकास करताना अनेक प्रश्न स्पष्ट करावे लागतात: हे
क्षेत्र कशाविषयी आहे? सर्वांत मूलभूत संकल्पना कोणत्या? त्यांचे परस्पर संबंध
कसे आहेत? क्षेत्रातील सर्व संकल्पना या मूलभूत संकल्पनांच्या साहाय्याने
परिभाषित करता येतात का? या संकल्पना कोणते नियम पाळतात आणि त्यांनी कोणते
नियम पाळलेच पाहिजेत?

स्वयंसिद्धकीय पद्धत आणि तर्कशास्त्र जणू एकमेकांसाठीच निर्माण झाले आहेत.
आकारिक तर्कशास्त्र स्वयंसिद्धकीय उपपत्ती मांडण्यासाठी, उपपत्तीच्या
स्वयंसिद्धकांपासून प्रमेये नेमक्या निर्दिष्ट रीतीने सिद्ध करण्यासाठी आणि
स्वयंसिद्धकांची पूर्ति करणाऱ्या सर्व प्रणालींच्या गुणधर्मांचा सुव्यवस्थित
अभ्यास करण्यासाठी साधने पुरवते.
\end{explain}

\begin{defn}
वाक्यांचा संच~$\Gamma$ \emph{बंद} असतो तेव्हा आणि केवळ तेव्हाच,
जेव्हा $\Gamma \Entails \varphi$ असल्यास $\varphi \in \Gamma$ असते.
वाक्यांच्या संच~$\Gamma$ चे \emph{संवरण} म्हणजे
$\Setabs{\varphi}{\Gamma \Entails \varphi}$ होय.

जर~$\Gamma$ हे~$\Delta$ चे संवरण असेल, तर~$\Gamma$ हा वाक्यांच्या
संच~$\Delta$ ने \emph{स्वयंसिद्धकांद्वारे निरूपित} केला आहे असे म्हणतो.
\end{defn}

\begin{explain}
एखादी स्वयंसिद्धकीय उपपत्ती म्हणजे तिच्या स्वयंसिद्धकांच्या संच~$\Delta$
ने निरूपित केलेल्या वाक्यांचा संच, असे आपण मानू शकतो. दुसऱ्या शब्दांत,
आपल्याला अभ्यासायच्या स्वयंसिद्धकीय रीतीने विकसित केलेल्या शास्त्रातील
मूलभूत संकल्पना दर्शवणारी अतार्किक चिन्हे असलेली प्रथम-क्रम भाषा आणि त्या
शास्त्राचे मूलभूत नियम व्यक्त करणारा वाक्यांचा संच दिलेला असेल,
तर स्वयंसिद्धकांपासून तार्किकरीत्या निष्पन्न होणारी या भाषेतील सर्व
वाक्ये उपपत्तीचे प्रतिनिधित्व करतात असे आपण मानू शकतो. याची व्याप्ती,
एकच आदिम संकल्पना आणि सोपी स्वयंसिद्धके असलेल्या अंशतः क्रमांच्या
उपपत्तीसारख्या साध्या उदाहरणांपासून न्यूटनीय यामिकीसारख्या गुंतागुंतीच्या
उपपत्तींपर्यंत आहे.

स्वयंसिद्धकीय पद्धतीकडे पाहण्याचा हा आकारिक दृष्टिकोन इतका महत्त्वाचा
ठरण्यामागील मुख्य तार्किक तथ्ये पुढीलप्रमाणे आहेत. $\Gamma$ ही एखाद्या
उपपत्तीची स्वयंसिद्धक-प्रणाली, म्हणजे वाक्यांचा संच, आहे असे समजा.
\begin{enumerate}
\item एखादी स्वयंसिद्धक-प्रणाली संरचनेचा अभिप्रेत वर्ग नेमका कधी
  पकडते, हे आपण काटेकोरपणे सांगू शकतो. म्हणजे, आपल्याला संरचनेच्या
  एखाद्या विशिष्ट वर्गात रस असेल, तर स्वयंसिद्धक-प्रणाली~$\Gamma$ तो वर्ग
  तेव्हा आणि केवळ तेव्हाच यशस्वीपणे पकडते, जेव्हा त्या संरचना म्हणजे
  नेमकी ती~$\Struct M$ असतात ज्यांच्यासाठी $\Sat{M}{\Gamma}$.
\item या बाबतीत आपल्याला अपयश येऊ शकते, कारण काही~$\Struct M$ साठी
  $\Sat{M}{\Gamma}$ असते, पण~$\Struct M$ ही आपण अभिप्रेत धरलेल्या
  संरचना पैकी नसते. त्यामुळे~$\Struct M$ मध्ये सत्य नसलेली
  स्वयंसिद्धके आपल्याला जोडावी लागू शकतात.
\item निदान~$\Gamma$ ही सर्व अभिप्रेत संरचनांमध्ये सत्य आहे एवढ्या
  बाबतीत आपण यशस्वी झालो, तर $\Gamma \Entails \varphi$ असताना वाक्य~$\varphi$ हे
  सर्व अभिप्रेत संरचनांमध्ये सत्य असते. म्हणून, वाक्ये सर्व
  अभिप्रेत संरचनांमध्ये सत्य आहेत हे दाखवण्यासाठी, ती
  स्वयंसिद्धकांपासून तार्किकरीत्या निष्पन्न होतात एवढे दाखवून आपण तार्किक
  साधने (उदाहरणार्थ, निष्पत्ती-पद्धती) वापरू शकतो.
\item कधीकधी आपल्या मनात अभिप्रेत संरचना नसतात; त्याऐवजी आपण
  स्वयंसिद्धकांपासूनच सुरुवात करतो: काही विशिष्ट नियमांची पूर्ति करावी अशा
  काही आदिम संकल्पना आपण घेतो आणि ते नियम स्वयंसिद्धक-प्रणालीत संकेतबद्ध
  करतो. स्वयंसिद्धके परस्परविरोधी नाहीत, ही गोष्ट आपल्याला लगेच पडताळून
  पाहायची असते. ती परस्परविरोधी असतील, तर त्या नियमांचे पालन करणाऱ्या
  संकल्पना असूच शकत नाहीत आणि आपण विसंगत उपपत्ती मांडण्याचा प्रयत्न केलेला
  असतो. $\Gamma$ चे प्रतिमान शोधून असे घडत नाही हे आपण पडताळू शकतो. आणि
  आपल्या उपपत्तीची प्रतिमाने असतील, तर त्यांचा तपास करण्यासाठी तसेच प्रतिमाने
  रचण्यासाठी आपण तार्किक पद्धती वापरू शकतो.
\item स्वयंसिद्धकांचे स्वातंत्र्य हाही एक महत्त्वाचा प्रश्न आहे. एखादे
  स्वयंसिद्धक प्रत्यक्षात इतरांचा परिणाम असल्यामुळे अनावश्यक असू शकते.
  $\Gamma$ मधील स्वयंसिद्धक~$\varphi$ अनावश्यक आहे हे आपण
  $\Gamma \setminus \{\varphi\} \Entails \varphi$ सिद्ध करून दाखवू शकतो. तसेच,
  $(\Gamma \setminus \{\varphi\}) \cup \{\lnot \varphi\}$ पूर्ततायोग्य आहे हे
  दाखवून ते स्वयंसिद्धक अनावश्यक नाही हेही सिद्ध करू शकतो. उदाहरणार्थ,
  भूमितीचे समांतर गृहीतक इतर स्वयंसिद्धकांपासून स्वतंत्र आहे हे अशाच प्रकारे
  दाखवले गेले.
\item उपपत्तीतील संकल्पनांची परिभाष्यता हाही आणखी एक महत्त्वाचा प्रश्न आहे:
  भाषेच्या निवडीवर उपपत्तीची प्रतिमाने कशाची बनलेली असतात हे ठरते. पण
  उपपत्तीचा प्रत्येक पैलू तिच्या प्रतिमानांत स्वतंत्रपणे दर्शवावाच लागतो असे
  नाही. उदाहरणार्थ, प्रत्येक क्रमसंबंध~$\le$ त्याच्याशी संबंधित काटेकोर
  क्रमसंबंध~$<$ ठरवतो---एक दिला असेल, तर दुसरा परिभाषित करता येतो. म्हणून
  असा क्रम असलेल्या उपपत्तीच्या प्रतिमानात त्याच्याशी संबंधित काटेकोर
  क्रमसंबंध \emph{वेगळ्याने} असणे आवश्यक नाही. सर्वसाधारणपणे, एक संबंध
  इतरांच्या आधारे नेमका कधी परिभाषित करता येतो? एखादा संबंध इतरांच्या
  आधारे परिभाषित करणे कधी अशक्य असते (आणि म्हणून तो भाषेच्या आदिम
  संकल्पनांत जोडावा लागतो)?
\end{enumerate}
\end{explain}












\section{संरचनेचे गुणधर्म व्यक्त करणे}\label{fol:mat:exs:sec}

\begin{explain}
फलने आणि संबंध यांच्यावरील अटी व्यक्त करणे, किंवा अधिक सर्वसाधारणपणे,
रचनेतील फलने आणि संबंध या अटींची पूर्ति करतात हे व्यक्त करणे अनेकदा
उपयुक्त आणि महत्त्वाचे असते. उदाहरणार्थ, एखाद्या भाषेसाठीच्या ज्या
संरचना आपल्याला विधेयांचा जो अर्थ ``पकडायचा'' आहे तो
पकडतात, त्या तसे न करणाऱ्या रचनांपासून वेगळ्या ओळखण्याचे मार्ग आपल्याला
हवे असतात. अर्थात, कोणती संरचना आपल्याला ``अभिप्रेत'' आहेत हे
ठरवण्याचे पूर्ण स्वातंत्र्य आपल्याला आहे. उदाहरणार्थ, आपण असे ठरवू शकतो
की विधेय~$\le$ चे अर्थनिर्धारण हा क्रमसंबंध असलाच पाहिजे; किंवा
आपल्याला फक्त~$\Lang L$ च्या अशा अर्थनिर्धारणांत रस आहे की ज्यांचा प्रांत
संचांचा बनलेला आहे आणि ज्यांत~$\Obj \in$ चे अर्थनिर्धारण ``चा
घटक आहे'' या संबंधाने केलेले आहे. पण भाषेतील वाक्ये वापरून
आपण हे करू शकतो का? दुसऱ्या शब्दांत, संरचना~$\Struct M$ वरील
कोणत्या अटी आपण~$\Struct M$ च्या भाषेतील वाक्य (किंवा कदाचित
वाक्यांचा संच) वापरून व्यक्त करू शकतो? काही अटी आपल्याला व्यक्त करता
येणार नाहीत. उदाहरणार्थ, फक्त $\Domain M = \Nat$ असलेल्या
संरचना~$\Struct M$ मध्येच सत्य असे~$\Lang L_A$ चे कोणतेही वाक्य
नाही. ``प्रांतात फक्त नैसर्गिक संख्या आहेत'' हे आपण व्यक्त करू शकत नाही.
पण संरचनेचे काही ``रचनात्मक गुणधर्म'' आपण कदाचित व्यक्त करू शकतो.
संरचनेचे कोणते गुणधर्म आपण वाक्ये वापरून व्यक्त करू शकतो?
किंवा, दुसऱ्या प्रकारे विचारल्यास, एखादे वाक्य (किंवा
वाक्यांचा संच) सत्य करणाऱ्या रचना म्हणून आपण
संरचनेच्या कोणत्या संग्रहांचे वर्णन करू शकतो?
\end{explain}

\begin{defn}[संचाचे प्रतिमान]
भाषा~$\Lang L$ मधील वाक्यांचा संच~$\Gamma$ असू द्या. प्रत्येक
$\varphi \in \Gamma$ साठी $\Sat{M}{\varphi}$ असेल, तर संरचना~$\Struct M$
\emph{हे~$\Gamma$ चे प्रतिमान आहे} असे म्हणतो.
\end{defn}

\begin{ex}
वाक्य $\lforall[x][x \le x]$ हे~$\Struct M$ मध्ये तेव्हा आणि केवळ तेव्हाच
सत्य असते, जेव्हा $\Assign{\le}{M}$ हा परावर्ती संबंध असतो. वाक्य
$\lforall[x][\lforall[y][((x \le y \land y \le x) \lif x = y)]]$ हे
~$\Struct M$ मध्ये तेव्हा आणि केवळ तेव्हाच सत्य असते, जेव्हा
$\Assign{\le}{M}$ प्रतिसममित असतो. वाक्य
$\lforall[x][\lforall[y][\lforall[z][((x \le y \land y \le z)
      \lif x \le z)]]]$ हे~$\Struct M$ मध्ये तेव्हा आणि केवळ तेव्हाच सत्य
असते, जेव्हा $\Assign{\le}{M}$ संक्रमक असतो. म्हणून पुढील संचाची प्रतिमाने
\begin{align*}
\{\quad &\lforall[x][x \le x], \\
   & \lforall[x][\lforall[y][((x \le y \land y \le
    x) \lif x = y)]], \\
   &\lforall[x][\lforall[y][\lforall[z][((x \le y
      \land y \le z) \lif x \le z)]]] \quad \}
\end{align*}
म्हणजे नेमकी ती रचना होत ज्यांत~$\Assign{\le}{M}$ हा परावर्ती,
प्रतिसममित आणि संक्रमक, म्हणजेच अंशतः क्रम, असतो. त्यामुळे ही वाक्ये
आपण \emph{अंशतः क्रमांच्या प्रथम-क्रम उपपत्तीची} स्वयंसिद्धके म्हणून घेऊ
शकतो.
\end{ex}












\section{प्रथम-क्रम उपपत्तींची उदाहरणे}\label{fol:mat:the:sec}

\begin{ex}
भाषा~$\Lang L_<$ मधील काटेकोर रेषीय क्रमांची उपपत्ती पुढील संचाने
स्वयंसिद्धकांद्वारे निरूपित होते:
\begin{align*}
\{\quad & \lforall[x][\lnot x < x], \\
& \lforall[x][\lforall[y][((x < y \lor y <
    x) \lor x = y)]], \\
& \lforall[x][\lforall[y][\lforall[z][((x < y
      \land y < z) \lif x < z)]]] \quad \}
\end{align*}
ही उपपत्ती अभिप्रेत संरचना पूर्णपणे पकडते: प्रत्येक काटेकोर
रेषीय क्रम या स्वयंसिद्धक-प्रणालीचे प्रतिमान असतो; आणि उलट, $R$ हा
संच~$X$ वरील रेषीय क्रम असेल, तर $\Domain M = X$ आणि
$\Assign{<}{M} = R$ असलेली रचना~$\Struct M$ या उपपत्तीचे प्रतिमान असते.
\end{ex}

\begin{ex}
$\Obj 1$ (स्थिरांक) आणि $\cdot$ (द्विस्थानी फलन) असलेल्या
भाषेतील गटांची उपपत्ती पुढील वाक्यांनी स्वयंसिद्धकांद्वारे निरूपित होते:
\begin{align*}
& \lforall[x][\eq[(x \cdot \Obj 1)][x]]\\
& \lforall[x][\lforall[y][\lforall[z][\eq[(x \cdot (y \cdot z))][((x
          \cdot y) \cdot z)]]]]\\
& \lforall[x][\lexists[y][\eq[(x \cdot y)][\Obj 1]]]
\end{align*}
\end{ex}

\begin{ex}
पेआनो अंकगणिताची उपपत्ती अंकगणिताची भाषा~$\Lang L_A$ मधील पुढील
वाक्यांनी स्वयंसिद्धकांद्वारे निरूपित होते.
\begin{align*}
& \lforall[x][\lforall[y][(\eq[x'][y'] \lif \eq[x][y])]]\\
& \lforall[x][\eq/[\Obj 0][x']]\\
& \lforall[x][\eq[(x + \Obj 0)][x]]\\
& \lforall[x][\lforall[y][\eq[(x + y')][(x + y)']]]\\
& \lforall[x][\eq[(x \times \Obj 0)][\Obj 0]]\\
& \lforall[x][\lforall[y][\eq[(x \times y')][((x \times y) + x)]]]\\
& \lforall[x][\lforall[y][(x < y \liff \lexists[z][\eq[(z' + x)][y])]]]\\
\intertext{आणि पुढील रूपाची सर्व वाक्ये}
& (\varphi(\Obj 0) \land \lforall[x][(\varphi(x) \lif \varphi(x'))]) \lif \lforall[x][\varphi(x)]
\end{align*}
शेवटच्या रूपाची अनंतपणे अनेक वाक्ये असल्यामुळे ही स्वयंसिद्धक-प्रणाली
अनंत आहे. शेवटच्या रूपाला \emph{विगमन-रूपबंध} म्हणतात. (प्रत्यक्षात,
विगमन-रूपबंधाचे स्वरूप आपण येथे दाखवले आहे त्यापेक्षा थोडे अधिक गुंतागुंतीचे
आहे.)

शेवटचे स्वयंसिद्धक हे~$<$ ची \emph{स्पष्ट व्याख्या} आहे.
\end{ex}

\begin{ex}
शुद्ध संचांची उपपत्ती गणिताच्या पायामध्ये (आणि गणिताच्या तत्त्वज्ञानातही)
महत्त्वाची भूमिका बजावते. एखाद्या संचाचे सर्व घटक हेही शुद्ध संच
असतील, तर तो संच शुद्ध असतो. म्हणून रिक्त संच शुद्ध मानला जातो; पण एखाद्या
संचात संच नसलेली कोणतीही वस्तू घटक म्हणून असेल, तर तो संच शुद्ध
नसतो. म्हणजे शुद्ध संच हे फक्त रिक्त संचापासून, कोणतेही ``मूलघटक'' न
वापरता बनवलेले संच होत; मूलघटक म्हणजे स्वतः संच नसलेल्या वस्तू.

पुढील वाक्ये शुद्ध संचांच्या उपपत्तीसाठी स्वयंसिद्धक-प्रणाली मानता येतील:
\begin{align*}
& \lexists[x][\lnot \lexists[y][y \in x]]\\
& \lforall[x][\lforall[y][(\lforall[z](z \in x \liff z \in y) \lif
\eq[x][y])]]\\
& \lforall[x][\lforall[y][\lexists[z][\lforall[u][(u \in z \liff
(\eq[u][x] \lor \eq[u][y]))]]]]\\
& \lforall[x][\lexists[y][\lforall[z][(z \in y \liff \lexists[u][(z \in
        u \land u \in x)])]]]\\
\intertext{आणि पुढील रूपाची सर्व वाक्ये} &
\lexists[x][\lforall[y][(y \in x \liff \varphi(y))]]
\end{align*}
पहिले स्वयंसिद्धक सांगते की कोणतेही घटक नसलेला एक संच आहे (म्हणजे
$\emptyset$ अस्तित्वात आहे); दुसरे सांगते की संच विस्तारात्मक आहेत; तिसरे
सांगते की कोणत्याही संच~$X$ आणि~$Y$ साठी संच~$\{X, Y\}$ अस्तित्वात आहे;
आणि चौथे सांगते की कोणत्याही संच~$X$ साठी संच~$\cup X$ अस्तित्वात आहे,
जिथे~$\cup X$ हा~$X$ च्या सर्व घटकांचा संयोग आहे.

सर्वांत शेवटी उल्लेखलेल्या वाक्यांना एकत्रितपणे
\emph{अनौपचारिक गुणधर्माधारित संचरचना-रूपबंध} म्हणतात. तो मूलतः असे सांगतो
की प्रत्येक~$\varphi(x)$ साठी संच~$\Setabs{x}{\varphi(x)}$ अस्तित्वात आहे---म्हणून
पहिल्या नजरेत ते सत्य, उपयुक्त आणि कदाचित आवश्यकही वाटणारे स्वयंसिद्धक आहे.
त्याला ``अनौपचारिक'' म्हणतात कारण, प्रत्यक्षात, त्यामुळे ही उपपत्ती
अपूर्ततायोग्य होते: $\varphi(y)$ म्हणून $\lnot y \in y$ घेतल्यास पुढील
वाक्य मिळते:
\[\lexists[x][\lforall[y][(y \in x \liff \lnot y \in y)]]
\]
आणि या वाक्याची कोणत्याही संरचनांमध्ये पूर्ति होत नाही.
\end{ex}

\begin{ex}
\emph{अवयवमीमांसेमध्ये} \emph{अवयवत्वाचा} संबंध हा एक मूलभूत संबंध आहे.
संचांच्या उपपत्तींप्रमाणेच अवयवत्वाच्याही अशा उपपत्ती आहेत, ज्या या
संबंधाविषयीच्या विविध (आणि कधीकधी परस्परविरोधी) संकल्पना
स्वयंसिद्धकांद्वारे निरूपित करतात.

अवयवमीमांसेच्या भाषेत एकच द्विस्थानी विधेयचिन्ह~$\Obj P$ असते आणि
$\Atom{\Obj P}{x, y}$ चा ``अर्थ'' $x$ हा~$y$ चा अवयव आहे असा असतो. हे
अर्थनिर्धारण मनात ठेवले असता, या भाषेतील संरचना ला
\emph{अवयवत्वाची रचना} म्हणतात. अर्थात, एकच द्विस्थानी विधेय असलेली प्रत्येक
रचना या नावास खरोखर पात्र असेलच असे नाही. ``अवयवत्व'' पकडण्याची शक्यता
असण्यासाठी $\Assign{\Obj P}{M}$ ने काही अटींची पूर्ति केली पाहिजे; त्या अटी
आपण अवयवत्वाच्या उपपत्तीची स्वयंसिद्धके म्हणून मांडू शकतो. उदाहरणार्थ,
वस्तूंवरील अवयवत्व हा अंशतः क्रम आहे: प्रत्येक वस्तू स्वतःचाच अवयव असते
(जरी तो \emph{अनुचित} अवयव असला तरी); कोणत्याही दोन भिन्न वस्तू एकमेकींचे
अवयव असू शकत नाहीत; आणि एखाद्या वस्तूच्या अवयवाचा अवयव हा स्वतः त्या
वस्तूचा अवयव असतो. लक्षात घ्या की या अर्थाने ``चा अवयव आहे'' हे ``चा
उपसंच आहे'' यासारखे आहे; पण ते ``चा घटक आहे'' यासारखे नाही, कारण घटक-संबंध
परावर्तीही नाही आणि संक्रमकही नाही.
\begin{align*}
& \lforall[x][\Atom{\Obj P}{x,x}] \\
& \lforall[x][\lforall[y][((\Part{x}{y} \land \Part{y}{x})
      \lif \eq[x][y])]] \\
& \lforall[x][\lforall[y][\lforall[z][((\Part{x}{y} \land
        \Part{y}{z}) \lif \Part{x}{z})]]]\\
\intertext{याखेरीज, कोणत्याही दोन वस्तूंचा अवयवी संयोग असतो (ज्याचे त्या दोन
  वस्तू अवयव असतात आणि जो या बाबतीत न्यूनतम असतो).}  &
\lforall[x][\lforall[y][\lexists[z][\lforall[u][(\Part{z}{u} \liff
        (\Part{x}{u} \land \Part{y}{u}))]]]]
\end{align*}
तत्त्वमीमांसकांनी विचारात घेतलेल्या अवयवत्वाच्या मूलभूत तत्त्वांपैकी ही
फक्त काही तत्त्वे आहेत. मात्र, आधी काही परिभाषित संबंध न मांडता पुढील
तत्त्वांची सूत्रे बनवणे किंवा ती लिहून काढणे लवकरच कठीण होते. उदाहरणार्थ,
अवयवमीमांसेत रस असलेले बहुतेक तत्त्वमीमांसक पुढील तत्त्वही वैध मानतात:
जेव्हा वस्तू~$x$ ला उचित अवयव~$y$ असतो, तेव्हा तिला असा अवयव~$z$ देखील
असतो की त्या अवयवाचा~$y$ शी कोणताही सामाईक अवयव नसतो आणि त्यामुळे~$y$ व
$z$ यांचा अवयवी संयोग~$x$ असतो.
\end{ex}












\section{संरचना मधील संबंध व्यक्त करणे}\label{fol:mat:exr:sec}

\begin{explain}
सूत्रांचा एक प्रमुख उपयोग म्हणजे संरचना~$\Struct M$ मधील
गुणधर्म आणि संबंध हे~$\Struct M$ ची भाषा~$\Lang L$ हिच्या आदिम संकल्पनांच्या
साहाय्याने व्यक्त करणे. याचा अर्थ पुढीलप्रमाणे आहे: $\Struct M$ चा
प्रांत हा वस्तूंचा संच असतो. स्थिरांक, फलने आणि
विधेये यांचे~$\Struct M$ मध्ये अनुक्रमे~$\Domain M$ मधील काही वस्तू,
$\Domain M$ वरील फलने आणि~$\Domain M$ वरील संबंध यांनी अर्थनिर्धारण केलेले
असते. उदाहरणार्थ, $\Obj A^2_0$ हे~$\Lang L$ मध्ये असेल, तर~$\Struct M$
त्याला संबंध~$R = \Assign{{\Obj A^2_0}}{M}$ नेमते. मग सूत्र
$\Atom{\Obj A^2_0}{\Obj v_1, \Obj v_2}$ पुढील अर्थाने नेमका तो संबंध
\emph{व्यक्त करते}: चर-मूल्यांकन~$s$ हे~$\Obj v_1$ ला
$a \in \Domain{M}$ आणि~$\Obj v_2$ ला $b \in \Domain M$ शी संगत करत असेल,
तर
\[Rab \text{\quad तेव्हा आणि केवळ तेव्हाच\quad} \Sat{M}{\Atom{\Obj A^2_0}{\Obj v_1, \Obj v_2}}[s].
\]
येथे चर-मूल्यांकनांचा समावेश करणे आवश्यक आहे हे लक्षात घ्या: आपण फक्त
``$Rab$ iff $\Sat{M}{\Atom{\Obj A^2_0}{a, b}}$'' असे म्हणू शकत नाही, कारण
$a$ आणि~$b$ ही आपल्या भाषेतील चिन्हे नाहीत; ती~$\Domain{M}$ ची
घटक आहेत.

आपल्याकडे केवळ अणुसूत्र सूत्रे नसून, तार्किक संयोजके आणि संख्यापक
वापरून ती एकत्र करता येतात; त्यामुळे अधिक गुंतागुंतीची सूत्रे अशी इतर
संबंध परिभाषित करू शकतात जी~$\Struct M$ मध्ये थेट अंतर्भूत नसतात. ते कसे
करायचे आणि विशेषतः संरचनांमध्ये कोणते संबंध परिभाषित करता येतात,
यात आपल्याला रस आहे.
\end{explain}

\begin{defn}
$\varphi(\Obj v_1,\dots, \Obj v_n)$ हे~$\Lang L$ चे असे सूत्र असू द्या
की ज्यात फक्त~$\Obj v_1$,\dots, $\Obj v_n$ हीच चरे मुक्त येतात; आणि
$\Struct M$ ही~$\Lang L$ साठीची संरचना असू द्या. कोणत्याही अशा
चर-मूल्यांकन~$s$ साठी की $s(\Obj v_i) = a_i$ ($i = 1, \dots, n$),
\[Ra_1\dots a_n \text{\quad तेव्हा आणि केवळ तेव्हाच\quad} \Sat{M}{\Atom{\varphi}{\Obj
    v_1,\dots, \Obj v_n}}[s]
\]
असे असेल, तर आणि केवळ तेव्हाच $\varphi(\Obj v_1,\dots, \Obj v_n)$ हे
$R \subseteq \Domain M^n$ हा \emph{संबंध व्यक्त करते}.
\end{defn}

\begin{ex}
अंकगणिताच्या मानक प्रतिमान~$\Struct N$ मध्ये सूत्र $\Obj
v_1 < \Obj v_2 \lor \eq[\Obj v_1][\Obj v_2]$ हे~$\Nat$ वरील~$\le$
संबंध व्यक्त करते. सूत्र $\eq[\Obj v_2][\Obj v_1']$ हे उत्तरवर्ती
संबंध व्यक्त करते; म्हणजे, $m$ हा~$n$ चा उत्तरवर्ती असेल तेव्हा
$Rnm$ असणारा संबंध~$R \subseteq \Nat^2$. सूत्र
$\eq[\Obj v_1][\Obj v_2']$ हे पूर्ववर्ती संबंध व्यक्त करते. सूत्रे
$\lexists[\Obj v_3][(\eq/[\Obj v_3][\Obj 0] \land
  \eq[\Obj v_2][(\Obj v_1 + \Obj v_3)])]$ आणि $\lexists[\Obj
  v_3][\eq[(\Obj v_1 + {\Obj v_3}')][\Obj v_2]]$ ही दोन्ही~$\Obj <$
संबंध व्यक्त करतात.

याचा अर्थ अंकगणिताच्या भाषेत विधेयचिन्ह~$<$ प्रत्यक्षात अनावश्यक आहे; ते
परिभाषित करता येते.
\end{ex}

\begin{explain}
ही कल्पना फक्त विशिष्ट संरचनांमध्येच रंजक नाही; एखादे अभिप्रेत
प्रतिमान किंवा प्रतिमाने वर्णन करण्यासाठी आपण भाषा वापरतो तेव्हा, म्हणजे
उपपत्तींचा विचार करताना, ती सर्वसाधारणपणेही रंजक ठरते. अशा उपपत्तींमध्ये
मूलभूत चिन्हे म्हणून अनेकदा फक्त काही विधेये असतात, पण ती ज्या
प्रांताचे वर्णन करतात त्यात इतर अनेक संबंध महत्त्वाची भूमिका बजावतात. या
इतर संबंधांना भाषेतील मूलभूत विधेयांचे अर्थनिर्धारण करणाऱ्या
संबंधांनी सुव्यवस्थित रीतीने व्यक्त करता येत असेल, तर आपण त्यांना त्या
भाषेत \emph{परिभाषित} करू शकतो असे म्हणतो.
\end{explain}

\begin{prob}
पुढील संबंध परिभाषित करणारी~$\Lang L_A$ मधील सूत्रे शोधा:
\begin{enumerate}
\item $n$ हे~$i$ आणि~$j$ यांच्या दरम्यान आहे;
\item $n$ ने~$m$ ला निःशेष भाग जातो (म्हणजे $m$ ही~$n$ ची पटीत संख्या आहे);
\item $n$ ही अविभाज्य संख्या आहे (म्हणजे $1$ आणि~$n$ यांखेरीज कोणत्याही
  संख्येने~$n$ ला निःशेष भाग जात नाही).
\end{enumerate}
\end{prob}

\begin{prob}
सूत्र $\varphi(\Obj v_1, \Obj v_2)$ हे संरचना~$\Struct M$ मध्ये संबंध
$R \subseteq \Domain M^2$ व्यक्त करते असे समजा. पुढील संबंध व्यक्त करणारी
सूत्रे शोधा:
\begin{enumerate}
\item $R$ चा व्युत्क्रम~$R^{-1}$;
\item सापेक्ष गुणाकार~$R \mid R$;
\end{enumerate}
$R$ चे संक्रमक संवरण~$R^+$ व्यक्त करण्याचा मार्ग तुम्हाला शोधता येतो का?
\end{prob}

\begin{prob}
फक्त एक द्विस्थानी विधेयचिन्ह~$<$ असलेली भाषा~$\Lang{L}$ असू द्या (इतर
कोणतीही स्थिरांक, फलने किंवा विधेये नाहीत---अर्थात
~$\eq$ सोडून). $\Domain{N} = \Nat$ आणि
$\Assign{<}{N} = \Setabs{\tuple{n,m}}{n < m}$ असणारी रचना~$\Struct{N}$
असू द्या. पुढील विधाने सिद्ध करा:
\begin{enumerate}
\item $\{ 0 \}$ हा~$\Struct{N}$ मध्ये परिभाष्य आहे;
\item $\{ 1 \}$ हा~$\Struct{N}$ मध्ये परिभाष्य आहे;
\item $\{ 2 \}$ हा~$\Struct{N}$ मध्ये परिभाष्य आहे;
\item प्रत्येक $n \in \Nat$ साठी संच~$\{ n \}$ हा~$\Struct{N}$ मध्ये
  परिभाष्य आहे;
\item $\Domain{N}$ चा प्रत्येक सांत उपसंच~$\Struct{N}$ मध्ये परिभाष्य आहे;
\item $\Domain{N}$ चा प्रत्येक सांत-पूरक उपसंच~$\Struct{N}$ मध्ये परिभाष्य
  आहे (जिथे $X \subseteq \Nat$ हा सांत-पूरक असतो तेव्हा आणि केवळ तेव्हाच,
  जेव्हा $\Nat \setminus X$ हा सांत असतो).
\end{enumerate}
\end{prob}












\section{संचांची उपपत्ती}\label{fol:mat:set:sec}

जवळजवळ संपूर्ण गणित संचांच्या उपपत्तीमध्ये विकसित करता येते. या
उपपत्तीमध्ये गणित विकसित करताना अनेक गोष्टी कराव्या लागतात. प्रथम,
संबंध~$\in$ साठी स्वयंसिद्धकांचा संच लागतो. वेगवेगळ्या अनेक
स्वयंसिद्धक-प्रणाली विकसित झाल्या असून, त्यांतील~$\in$ चे गुणधर्म कधीकधी
परस्परविरोधी असतात. निवडीच्या स्वयंसिद्धकासह झर्मेलो--फ्रेंकेल संच
उपपत्ती म्हणून ओळखली जाणारी स्वयंसिद्धक-प्रणाली~$\Log{ZFC}$ विशेष उठून
दिसते: ती आतापर्यंत सर्वाधिक वापरली आणि अभ्यासली गेलेली प्रणाली आहे, कारण
गणितज्ञांना ज्या जवळजवळ सर्व गोष्टी सिद्ध करता याव्यात अशी अपेक्षा असते,
त्या सिद्ध करण्यास तिची स्वयंसिद्धके पुरेशी ठरतात. पण हे स्थापित करण्यापूर्वी
गणितज्ञांना व्यक्त करायच्या सर्व गोष्टी आपण \emph{व्यक्त} तरी कशा करू शकतो,
हे प्रथम स्पष्ट करणे आवश्यक आहे. सुरुवातीलाच, भाषेत कोणतीही स्थिरांक
किंवा फलने नसतात; त्यामुळे विशिष्ट संचांविषयी (उदाहरणार्थ,
$\emptyset$ किंवा~$\Nat$), संचांवरील क्रियांविषयी (उदाहरणार्थ,
$X \cup Y$ आणि~$\Pow{X}$), इतकेच नव्हे तर संबंध आणि फलने यांसारख्या
संचांखेरीज इतर गोष्टींचा समावेश असलेल्या रचनांविषयी आपण बोलू शकतो हेही
पहिल्या नजरेत अस्पष्ट वाटते.

सुरुवातीला, ``चा घटक आहे'' हाच काही आपल्याला महत्त्वाचा असलेला एकमेव संबंध
नाही; ``चा उपसंच आहे'' हा जवळजवळ तितकाच महत्त्वाचा वाटतो. पण ``चा घटक
आहे'' याच्या आधारे ``चा उपसंच आहे'' हे आपण \emph{परिभाषित} करू शकतो. त्यासाठी
संच उपपत्तीच्या भाषेत असे सूत्र~$\varphi(x, y)$ शोधावे लागते की संचांची
जोडी~$\tuple{X, Y}$ त्याची पूर्ति तेव्हा आणि केवळ तेव्हाच करते, जेव्हा
$X \subseteq Y$. पण~$X$ हा~$Y$ चा उपसंच तेव्हा आणि केवळ तेव्हाच असतो,
जेव्हा~$X$ चे सर्व घटक हे~$Y$ चेही घटक असतात. म्हणून
पुढील सूत्राने~$\subseteq$ ची व्याख्या करता येते:
\[\lforall[z][(z \in x \lif z \in y)]
\]
आता, एखाद्या सूत्रात~$\subseteq$ संबंध वापरायचा असेल तेव्हा त्याऐवजी हे
सूत्र वापरता येईल ($x$ आणि~$y$ यांच्या जागी योग्य पदे घालून आणि आवश्यक
असल्यास बद्ध चर~$z$ चे नाव बदलून). उदाहरणार्थ, संचांची विस्तारात्मकता
म्हणजे कोणतेही संच~$x$ आणि~$y$ हे परस्परांचे उपसंच असतील, तर~$x$ आणि~$y$
एकच संच असले पाहिजेत. हे
$\lforall[x][\lforall[y][((x \subseteq y \land y
    \subseteq x) \lif x = y)]]$ ने व्यक्त करता येते; किंवा वरच्या व्याख्येने
$\subseteq$ बदलल्यास पुढील सूत्राने:
\[\lforall[x][\lforall[y][((\lforall[z][(z \in x \lif z \in y)] \land
    \lforall[z][(z \in y \lif z \in x)]) \lif x = y)]].
\]
हे खरे तर~$\Log{ZFC}$ च्या स्वयंसिद्धकांपैकी एक, ``विस्तारात्मकतेचे
स्वयंसिद्धक,'' आहे.

$\emptyset$ साठी कोणतेही स्थिरांक नाही, पण ``$x$ रिक्त आहे'' हे
$\lnot \lexists[y][y \in x]$ ने व्यक्त करता येते. मग ``$\emptyset$
अस्तित्वात आहे'' हे वाक्य~$\lexists[x][\lnot \lexists[y][y \in
    x]]$ बनते. हे~$\Log{ZFC}$ चे आणखी एक स्वयंसिद्धक आहे. (विस्तारात्मकतेच्या
स्वयंसिद्धकानुसार फक्त एकच रिक्त संच असतो, हे लक्षात घ्या.) संच उपपत्तीच्या
भाषेत~$\emptyset$ विषयी बोलायचे असेल तेव्हा आपण ते ``रिक्त असलेला एक संच
अस्तित्वात आहे आणि \dots'' असे लिहू. उदाहरणार्थ, $\emptyset$ हा प्रत्येक संचाचा
उपसंच आहे ही गोष्ट पुढीलप्रमाणे व्यक्त करता येईल:
\[\lexists[x][(\lnot\lexists[y][y \in x] \land \lforall[z][x \subseteq
    z])]
\]
येथे अर्थातच $x \subseteq z$ हेसुद्धा त्याच्या व्याख्येने बदलावे लागेल.

$X \cup Y$ आणि~$\Pow{X}$ यांसारख्या संचांवरील क्रियांविषयी बोलण्यासाठी
अशीच युक्ती वापरावी लागते. संच उपपत्तीच्या भाषेत फलनचिन्हे नसतात, पण
$X \cup Y = Z$ आणि~$\Pow{X} = Y$ हे फलनीय संबंध पुढील सूत्रांनी व्यक्त
करता येतात:
\begin{align*}
& \lforall[u][((u \in x \lor u \in y) \liff u \in z)]\\
& \lforall[u][(u \subseteq x \liff u \in y )]
\end{align*}
कारण $X \cup Y$ चे घटक म्हणजे नेमके~$X$ चे घटक किंवा
~$Y$ चे घटक असलेले संच, आणि~$\Pow{X}$ चे घटक म्हणजे नेमके~$X$
चे उपसंच होत. पण त्यामुळे आपण $x \cup y$ किंवा~$\Pow{x}$ ही पदे असल्यासारखी
वापरू शकत नाही; $X \cup Y = Z$ आणि~$\Pow{X} = Y$ हे संबंध परिभाषित करणारी
संपूर्ण सूत्रे एवढीच वापरता येतात. प्रत्यक्षात, या संबंधांची कधी तरी
पूर्ति होते, म्हणजे संयोग आणि घातसंच नेहमी अस्तित्वात असतात, हे आपल्याला
अजून माहीत नाही. उदाहरणार्थ, वाक्य
$\lforall[x][\lexists[y][\Pow{x} = y]]$ हे~$\Log{ZFC}$ चे आणखी एक
स्वयंसिद्धक (घातसंचाचे स्वयंसिद्धक) आहे.

आता क्रमित जोड्या किंवा फलने यांविषयी कसे बोलायचे? येथे क्रमित जोड्या आणि
फलने हे विशिष्ट प्रकारचे संच म्हणून कसे मानता येतात, हे स्पष्ट करावे लागते.
क्रमित जोडी~$\tuple{x, y}$ ची व्याख्या संच~$\{\{x\}, \{x, y\}\}$ म्हणून
करणे हा एक मार्ग आहे. पण आधीप्रमाणेच या संचाला नाव देणारे फलन
आपण समाविष्ट करू शकत नाही; फक्त $\tuple{x, y} = z$, म्हणजे
$\{\{x\}, \{x, y\}\} = z$, हा संबंध परिभाषित करता येतो:
\[\lforall[u][(u \in z \liff (\lforall[v][(v \in u \liff v = x)] \lor
  \lforall[v][(v \in u \liff (v = x \lor v = y))]))]
\]
याचा अर्थ~$z$ चे घटक~$u$ म्हणजे नेमके असे संच की ज्यांचा एकमेव
घटक हा~$x$ असतो किंवा ज्यांचे एकमेव घटक हे~$x$ आणि~$y$
असतात (दुसऱ्या शब्दांत, जे संच~$\{x\}$ शी किंवा~$\{x, y\}$ शी एकरूप
असतात). हे मिळाल्यानंतर आणखी गोष्टीही सांगता येतात; उदाहरणार्थ,
$X \times Y = Z$:
\[\lforall[z][(z \in Z \liff \lexists[x][\lexists[y][(x \in
        X \land y \in Y \land \tuple{x, y} = z)]])]
\]

फलन~$f \colon X \to Y$ याचा संबंध~$f(x) = y$, म्हणजे जोड्यांचा संच
$\Setabs{\tuple{x,y}}{f(x) = y}$, म्हणून विचार करता येतो. मग संच~$f$ हे
$X$ पासून~$Y$ पर्यंतचे फलन आहे असे आपण तेव्हा म्हणू शकतो, जेव्हा (a) तो
$\subseteq X \times Y$ असा संबंध आहे, (b) तो सर्वत्र परिभाषित आहे,
म्हणजे प्रत्येक $x \in X$ साठी असा काही $y \in Y$ आहे की
$\tuple{x, y} \in f$, आणि (c) तो एकमूल्यी आहे, म्हणजे
$\tuple{x, y}, \tuple{x, y'} \in f$ असताना $y = y'$ (कारण फलनांची मूल्ये
एकमेव असली पाहिजेत). म्हणून ``$f$ हे~$X$ पासून~$Y$ पर्यंतचे फलन आहे'' हे
पुढीलप्रमाणे लिहिता येते:
\begin{align*}
 \lforall[u][(u \in f \lif {}] & \lexists[x][\lexists[y][(x \in X \land y \in
      Y \land \tuple{x, y} = u)]]) \land {}\\
 \lforall[x][(x \in X \lif {}] &
      (\lexists[y][(y \in Y \land \mathrm{maps}(f, x, y))] \land {}\\
& (\lforall[y][\lforall[y'][((\mathrm{maps}(f, x, y) \land
    \mathrm{maps}(f, x, y')) \lif y = y')]]))
\end{align*}
येथे $\mathrm{maps}(f, x, y)$ हे $\lexists[v][(v \in f \land
  \tuple{x, y} = v)]$ चे संक्षिप्त रूप आहे (हे सूत्र ``$f(x) = y$''
व्यक्त करते).

आता, उदाहरणार्थ, $f\colon X \to Y$ हे एकास-एक आहे हे व्यक्त करणेही
कठीण नाही:
\begin{multline*}
 f \colon X \to Y \land \lforall[x][\lforall[x'][((x \in X \land x' \in
     X \land {}]] \\
 \lexists[y][(\mathrm{maps}(f, x, y) \land \mathrm{maps}(f,
        x', y))]) \lif x = x')
\end{multline*}
फलन~$f\colon X \to Y$ हे एकास-एक तेव्हा आणि केवळ तेव्हाच असते, जेव्हा
$f$ हे $x, x' \in X$ या दोन्हींना एकाच~$y$ शी संगत करत असल्यास $x = x'$.
हे सूत्र~$\mathrm{inj}(f, X, Y)$ असे संक्षिप्तपणे लिहिल्यास, आता आपण संच
उपपत्तीच्या भाषेत कँटरच्या प्रमेयासारखी साधी नसलेली गोष्टही मांडू शकतो:
$\Pow{X}$ पासून~$X$ मध्ये जाणारे कोणतेही एकास-एक फलन नाही:
\[\lforall[X][\lforall[Y][(\Pow{X} = Y \lif
    \lnot\lexists[f][\mathrm{inj}(f, Y, X)])]]
\]

प्रत्येक परिभाषक गुणधर्मासाठी संचाचे अस्तित्व हमीने देणारे आणखी एक
स्वयंसिद्धक संच उपपत्तीला आवश्यक आहे असे वाटू शकते. $\varphi(x)$ हे मुक्त
चर~$x$ असलेले संच उपपत्तीचे सूत्र असेल, तर पुढील वाक्य विचारात घेता
येते:
\[\lexists[y][\lforall[x][(x \in y \liff \varphi(x))]].
\]
हे वाक्य सांगते की असा संच~$y$ आहे ज्याचे घटक म्हणजे नेमके
ते सर्व~$x$ होत जे~$\varphi(x)$ ची पूर्ति करतात. या रूपबंधाला ``गुणधर्माधारित
संचरचनेचे तत्त्व'' म्हणतात. ते फार उपयुक्त वाटते; पण दुर्दैवाने ते विसंगत
आहे. $\varphi(x) \ident \lnot x \in x$ घ्या; मग गुणधर्माधारित संचरचनेचे तत्त्व
पुढील विधान सांगते:
\[\lexists[y][\lforall[x][(x \in y \liff x \notin x)]],
\]
म्हणजे स्वतःचे घटक नसलेल्या सर्व संचांच्या संचाचे अस्तित्व ते
सांगते. असा संच अस्तित्वात असू शकत नाही---ही रसेलची विरोधापत्ती आहे.
प्रत्यक्षात, $\Log{ZFC}$ मध्ये या तत्त्वाची मर्यादित---आणि सुसंगत---आवृत्ती,
म्हणजे विभक्तीकरण तत्त्व, समाविष्ट आहे:
\[\lforall[z][\lexists[y][\lforall[x][(x \in y \liff (x \in z \land
      \varphi(x))]]].
\]

\begin{prob}
पुढील निष्पन्नता दाखवणारी निष्पत्ती देऊन गुणधर्माधारित संचरचनेचे
तत्त्व विसंगत आहे, हे दाखवा:
\[\lexists[y][\lforall[x][(x \in y \liff x \notin x)]] \Proves \lfalse.
\]
प्रथम $(A \lif \lnot A) \land (\lnot A \lif A)
\Proves \lfalse$ दाखवल्यास मदत होईल.
\end{prob}











\section{संरचनेचे आकारमान व्यक्त करणे}\label{fol:mat:siz:sec}

\begin{explain}
भाषेतील अतार्किक चिन्हे न वापरताही रचनांचे काही गुणधर्म व्यक्त करता येतात.
उदाहरणार्थ, अशी वाक्ये आहेत की ती संरचनांमध्ये तेव्हा आणि केवळ
तेव्हाच सत्य असतात, जेव्हा त्या संरचनेच्या प्रांत मध्ये
किमान, जास्तीत जास्त किंवा नेमके दिलेल्या संख्येइतके~$n$ घटक असतात.
\end{explain}

\begin{prop}
वाक्य
\begin{multline*}
  \varphi_{\ge n} \ident \lexists[x_1][\lexists[x_2][\dots\lexists[x_n][{}]]]\\
  \begin{aligned}
  (\eq/[x_1][x_2] \land {}
  \eq/[x_1][x_3] \land \eq/[x_1][x_4] \land \dots \land \eq/[x_1][x_n] \land {}\\
\eq/[x_2][x_3] \land \eq/[x_2][x_4] \land \dots \land {} \eq/[x_2][x_n] \land {} \\
\vdots\\
\eq/[x_{n-1}][x_n])
\end{aligned}
\end{multline*}
हे संरचना~$\Struct M$ मध्ये तेव्हा आणि केवळ तेव्हाच सत्य असते,
जेव्हा $\Domain M$ मध्ये किमान $n$ घटक असतात. परिणामी,
$\Sat{M}{\lnot \varphi_{\ge n+1}}$ तेव्हा आणि केवळ तेव्हाच लागू होते, जेव्हा
$\Domain M$ मध्ये जास्तीत जास्त~$n$ घटक असतात.

\end{prop}

\begin{prop}
वाक्य
\begin{multline*}
\varphi_{= n} \ident \lexists[x_1][\lexists[x_2][\dots\lexists[x_n][{}]]] \\
\begin{aligned}
  (\eq/[x_1][x_2] \land {}
  \eq/[x_1][x_3] \land \eq/[x_1][x_4] \land \dots \land \eq/[x_1][x_n] \land {}\\
\eq/[x_2][x_3] \land \eq/[x_2][x_4] \land \dots \land {} \eq/[x_2][x_n] \land {} \\
\vdots\\
\eq/[x_{n-1}][x_n] \land {} \\
\lforall[y][(\eq[y][x_1] \lor \dots \lor \eq[y][x_n]]))
\end{aligned}
\end{multline*}
हे संरचना~$\Struct M$ मध्ये तेव्हा आणि केवळ तेव्हाच सत्य असते,
जेव्हा $\Domain M$ मध्ये नेमके $n$ घटक असतात.
\end{prop}

\begin{prop}
एखादी संरचना अनंत तेव्हा आणि केवळ तेव्हाच असते, जेव्हा ती पुढील
संचाचे प्रतिमान असते:
\[\{\varphi_{\ge 1}, \varphi_{\ge 2}, \varphi_{\ge 3}, \dots \}.
\]
\end{prop}

असे एकही पूर्णतः तार्किक वाक्य नाही की जे~$\Struct M$ मध्ये तेव्हा आणि केवळ
तेव्हाच सत्य असेल, जेव्हा $\Domain M$ अनंत असेल. पण अशी अतार्किक
विधेये असलेली वाक्ये देता येतात ज्यांची फक्त अनंत प्रतिमाने
असतात (जरी प्रत्येक अनंत संरचना त्यांचे प्रतिमान नसते). सांत रचना
असण्याचा गुणधर्म आणि अगणनीय रचना असण्याचा गुणधर्म
वाक्यांच्या अनंत संचानेही व्यक्त करता येत नाही. ही तथ्ये संहतता आणि
लोव्हेनहाइम--स्कोलेम प्रमेयांवरून निष्पन्न होतात.



\chapter{प्रथम-क्रम तर्कशास्त्रापलीकडे}\label{fol:byd::chap}
\begin{editorial}
  जेरेमी अविगॅड यांच्या तर्कशास्त्रावरील टिपणांवरून रूपांतरित केलेला हा
  अध्याय इतर कोणकोणत्या तार्किक प्रणाली अस्तित्वात आहेत यांची अत्यंत संक्षिप्त
  झलक देतो. ज्या अभ्यासक्रमात या प्रणाली शिकवल्या जात नाहीत, त्यात पुढील
  अध्ययनासाठी विषय सुचवणारा अध्याय म्हणून तो अभिप्रेत आहे. येथे उल्लेखिलेल्या
  प्रत्येक विषयाला पुढे---अशी आशा आहे---ओपन लॉजिक प्रोजेक्टमध्ये स्वतंत्र
  भागाच्या पातळीवरील विवेचन मिळेल.
\end{editorial}









\section{आढावा}\label{fol:byd:int:sec}

प्रथम-क्रम तर्कशास्त्र हीच काही महत्त्वाची एकमेव तार्किक प्रणाली नाही:
प्रथम-क्रम तर्कशास्त्राचे अनेक विस्तार आणि अनेक भेद आहेत. सामान्यतः एखाद्या
तर्कशास्त्रात भाषेचे आकारिक विनिर्देशन, बहुतेकदा---पण नेहमीच नाही---एक
निगामी प्रणाली आणि बहुतेकदा---पण नेहमीच नाही---एक अभिप्रेत चिन्हार्थमीमांसा
यांचा समावेश असतो. पण या संज्ञेच्या तांत्रिक वापरामुळे एक उघड प्रश्न निर्माण
होतो: सहजप्रज्ञ किंवा तात्त्विक अर्थाने वापरल्या जाणाऱ्या ``तर्कशास्त्र'' या
शब्दाशी प्रथम-क्रम तर्कशास्त्र नसलेल्या तर्कशास्त्रांचा संबंध काय? खाली
वर्णिलेल्या सर्व प्रणाली कोणत्या ना कोणत्या प्रकारच्या युक्तिविचाराचे प्रतिमान
देण्यासाठी रचल्या आहेत; पण त्या तार्किक कशामुळे ठरतात, हे आपण सांगू शकतो का?

याची सोपी उत्तरे मिळत नाहीत. ``तर्कशास्त्र'' हा शब्द वेगवेगळ्या अर्थांनी आणि
वेगवेगळ्या संदर्भांत वापरला जातो; आणि ``सत्य'' या संकल्पनेप्रमाणेच या
संकल्पनेचेही अनेक तात्त्विक भूमिकांतून विश्लेषण झाले आहे. उदाहरणार्थ, कोणती
विधाने अनिवार्यतः सत्य, अनुभवपूर्व सत्य, अतार्किक पदांच्या अर्थनिर्धारणापासून
स्वतंत्रपणे सत्य, त्यांच्या आकारामुळे सत्य किंवा भाषिक संकेतामुळे सत्य आहेत
हे ठरवणे हे तार्किक युक्तिविचाराचे ध्येय आहे, असे कोणी मानू शकेल; आणि यांपैकी
प्रत्येक संकल्पना बरीच स्पष्ट करावी लागते. केवळ गणितात वापरल्या जाणाऱ्या
तर्कशास्त्रापुरते लक्ष मर्यादित केले, तरी त्याची व्याप्ती किती याविषयी फारसे
एकमत नाही. उदाहरणार्थ, फ्रेगे यांनी सुरू केलेल्या {\em तर्कमीमांसावादी}
परंपरेत रसेल आणि व्हाइटहेड यांनी \textit{प्रिन्सिपिया मॅथेमॅटिका} मध्ये
तर्कशास्त्राच्या पायावर गणित विकसित करण्याचा प्रयत्न केला. त्यांची तार्किक
प्रणाली खाली वर्णिलेल्या प्रणालीसारख्या उच्च-प्रकार तर्कशास्त्राचे एक रूप
होती. शेवटी त्यांना बहुतेक निकषांनुसार पूर्णतः तार्किक न वाटणारी स्वयंसिद्धके
समाविष्ट करावी लागली (विशेषतः अनंतत्वाचे स्वयंसिद्धक आणि न्यूनीकरणक्षमतेचे
स्वयंसिद्धक); तरीही उच्च-क्रम युक्तिविचाराची काही रूपे तार्किक म्हणून
स्वीकारली पाहिजेत, असे कोणी मानू शकेल. याउलट, ज्यांच्या सत्ताशास्त्रात
``विधानांना'' संवादाच्या वैध वस्तू म्हणून स्थान नाही असे क्वाइन म्हणतात की,
द्वितीय-क्रम आणि उच्च-क्रम तर्कशास्त्र ही मेंढीचे कातडे पांघरलेल्या संच
उपपत्तीचीच रूपे आहेत; म्हणजेच, विधेयांवर संख्यकीकरण करणाऱ्या प्रणाली पूर्णतः
तार्किक नाहीत.

आत्तासाठी असे तात्त्विक प्रश्न पुढे कधीतरी विचारात घेणेच योग्य ठरेल; आणि
खालील प्रणाली तार्किक असोत वा नसोत, विविध प्रकारच्या युक्तिविचारांची आकारिक
आदर्शीकरणे म्हणून त्यांच्याकडे पाहूया.












\section{बहुवर्गीय तर्कशास्त्र}\label{fol:byd:msl:sec}

प्रथम-क्रम तर्कशास्त्रात चरे आणि संख्यापक एकाच प्रांत वर लागू होतात.
पण अनेक (परस्परविभक्त) प्रांत असणे बऱ्याचदा उपयुक्त ठरते: उदाहरणार्थ,
संख्यांचा प्रांत, भूमितीय वस्तूंचा प्रांत, संख्यांपासून संख्यांकडे
जाणाऱ्या फलनांचा प्रांत, आबेली गटांचा प्रांत, इत्यादी ठेवावेसे वाटू
शकतात.

बहुवर्गीय तर्कशास्त्र अशी चौकट पुरवते. सुरुवात ``वर्गांच्या'' यादीने होते---
वस्तूचा ``वर्ग'' ती कोणत्या ``प्रांत'' मध्ये असणे अपेक्षित आहे ते दर्शवतो.
मग प्रत्येक वर्गासाठी स्वतंत्र चले आणि संख्यापक, तसेच (सामान्यतः)
स्वतंत्र एकरूपता असते. फलने आणि संबंध कोणत्या वर्गांतील वस्तू फलसाधक
म्हणून घेऊ शकतात यानुसार तेही ``प्रकारबद्ध'' केलेले असतात. याखेरीज
प्रथम-क्रम तर्कशास्त्राचे नेहमीचे नियम तसेच ठेवले जातात; संख्यापकांचे नियम
प्रत्येक वर्गासाठी स्वतंत्रपणे पुनरावृत्त होतात.

उदाहरणार्थ, आंतरराष्ट्रीय संबंधांचा अभ्यास करण्यासाठी आपण फ्रेंच नागरिक आणि
जर्मन नागरिक असे वस्तूंचे दोन वर्ग असलेली भाषा निवडू शकतो. पहिल्या वर्गातील
वस्तूंसाठी ``द्राक्षारस पितो'' हा एकस्थानी संबंध; दुसऱ्या वर्गातील वस्तूंसाठी
``वुर्स्ट खातो'' हा दुसरा एकस्थानी संबंध; आणि दोन फलसाधक घेणारा ``बहुराष्ट्रीय
विवाहित जोडपे बनवतात'' हा द्विस्थानी संबंध असू शकतो, ज्याचा पहिला फलसाधक
पहिल्या वर्गातील आणि दुसरा फलसाधक दुसऱ्या वर्गातील असतो. फ्रेंच नागरिकांसाठी
$a$, $b$, $c$ ही चरे आणि जर्मन नागरिकांसाठी $x$, $y$, $z$ ही चरे वापरली, तर
\[\lforall[a][\lforall x][(\Atom{\Obj{MarriedTo}}{a,x} \lif
(\Atom{\Obj{DrinksWine}}{a} \lor \lnot \Atom{\Obj{EatsWurst}}{x})]]
\]
हे सूत्र असे सांगते की कोणत्याही फ्रेंच व्यक्तीचा विवाह एखाद्या जर्मन
व्यक्तीशी झाला असेल, तर एकतर ती फ्रेंच व्यक्ती द्राक्षारस पिते किंवा ती जर्मन
व्यक्ती वुर्स्ट खात नाही.

बहुवर्गीय तर्कशास्त्र प्रथम-क्रम तर्कशास्त्रात स्वाभाविक रीतीने अंतःस्थापित
करता येते: बहुवर्गीय प्रांत मधील सर्व वस्तू एकत्र करून प्रथम-क्रमाचा एक
प्रांत घ्यायचा, वर्गांची नोंद ठेवण्यासाठी एकस्थानी विधेये
वापरायची आणि संख्यापकांना योग्य वर्गापुरते सापेक्ष करायचे. उदाहरणार्थ, वरच्या
उदाहरणाशी संबंधित प्रथम-क्रम भाषेत वर्णिलेल्या इतर संबंधांबरोबर
``$\Obj{German}$'' आणि ``$\Obj{French}$'' ही एकस्थानी विधेये असतील;
मात्र वर्गाविषयीच्या अटी पुसलेल्या असतील. जर्मन वर्गातील चल~$x$
असलेला वर्गबद्ध संख्यापक $\lforall[x][\varphi]$ पुढीलप्रमाणे रूपांतरित होतो:
\[\lforall[x][(\Atom{\Obj{German}}{x} \lif \varphi)].
\]
वर्ग परस्परांपासून वेगळे आहेत याची खात्री देणारी स्वयंसिद्धकेही जोडावी
लागतात---उदाहरणार्थ,
$\lforall[x][\lnot (\Atom{\Obj{German}}{x} \land
  \Atom{\Obj{French}}{x})]$---तसेच ``द्राक्षारस पितो'' हा संबंध फक्त
$\Atom{\Obj{French}}{x}$ या विधेयाची पूर्ति करणाऱ्या वस्तूंनाच लागू होतो
याची हमी देणारी स्वयंसिद्धकेही जोडावी लागतात. या संकेतांमुळे आणि
स्वयंसिद्धकांमुळे बहुवर्गीय वाक्ये प्रथम-क्रम वाक्ये मध्ये आणि
बहुवर्गीय निष्पत्त्या प्रथम-क्रम निष्पत्त्यांमध्ये रूपांतरित होतात,
हे दाखवणे अवघड नाही. तसेच बहुवर्गीय संरचना संबंधित प्रथम-क्रम
संरचनांमध्ये आणि उलट दिशेनेही ``रूपांतरित'' होतात; म्हणून बहुवर्गीय
तर्कशास्त्रासाठीही आपल्याला संपूर्णता प्रमेय मिळते.












\section{द्वितीय-क्रम तर्कशास्त्र}\label{fol:byd:sol:sec}

द्वितीय-क्रम तर्कशास्त्राची भाषा व्यक्तींच्या प्रांत वरच नव्हे, तर त्या
प्रांत वरील संबंधांवरही संख्यापन करण्याची मुभा देते. प्रथम-क्रम
भाषा~$\Lang{L}$ दिली असता, प्रत्येक $k$ साठी $k$-स्थानी संबंधांवर मान घेणारी
$R$ ही चले जोडली जातात आणि त्या चरांवर संख्यापन करण्याची मुभा दिली
जाते. $R$ हे एखाद्या $k$-स्थानी संबंधाचे चल असेल आणि $t_1$,
\dots,~$t_k$ ही नेहमीची (प्रथम-क्रम) पदे असतील, तर
$\Atom{R}{t_1,\dots,t_k}$ हे अण्वीय सूत्र आहे. याखेरीज सूत्रांचा
संच प्रथम-क्रम तर्कशास्त्राप्रमाणेच, फक्त द्वितीय-क्रम संख्यापनासाठी अतिरिक्त
उपनियम घालून, परिभाषित केला जातो. लक्षात घ्या की आपल्याकडे एकरूपता फक्त
प्रथम-क्रम पदांसाठी आहे: $R$ आणि $S$ ही समान स्थानसंख्या~$k$ असलेल्या
संबंधांची चले असतील, तर $\eq[R][S]$ हे पुढील सूत्राचे संक्षिप्त रूप
म्हणून परिभाषित करता येते:
\[\lforall[x_1 \dots][\lforall[x_k][(\Atom{R}{x_1, \dots, x_k} \liff
  \Atom{S}{x_1, \dots, x_k})]].
\]

द्वितीय-क्रम तर्कशास्त्राचे नियम संख्यापकांचे नियम नव्या द्वितीय-क्रम चरांपर्यंत
विस्तारित करतात. मात्र ही चरे $\Lang{L}$ मधील विधेये शी आणि अधिक
सामान्यपणे $\Lang{L}$ मधील सूत्रे शी कशी परस्परक्रिया करतात, हे येथे
थोडे काळजीपूर्वक स्पष्ट करावे लागते. किमान चौकटीत संबंधचरे पदे म्हणून गणली
जातात; म्हणून पुढील प्रकारच्या अनुमानांना मुभा असते:
\[\varphi(R) \Proves \lexists[R][\varphi(R)]
\]
पण $\Lang{L}$ ही अंकगणिताची भाषा असून $<$ हे स्थिर संबंधचिन्ह असेल, तर पुढील
अनुमानही ग्राह्य असावे अशी अपेक्षा असते:
\[x < y \Proves \lexists[R][\Atom{R}{x,y}]
\]
किंवा दिलेल्या सूत्र~$\varphi$ साठी,
\[\varphi(x_1, \dots, x_k) \Proves \lexists[R][\Atom{R}{x_1,\dots,x_k}]
\]
आणखी सामान्यपणे, पुढील प्रकारच्या अनुमानांना मुभा द्यावीशी वाटू शकते:
\[\Subst{\varphi}{\lambd[\vec x][\psi(\vec x)]}{R} \Proves
\lexists[R][\varphi]
\]
येथे $\Subst{\varphi}{\lambd[\vec x][\psi(\vec x)]}{R}$ म्हणजे $\varphi$ मधील
$\Obj{R}{t_1,\dots,t_k}$ या रूपाचे प्रत्येक अण्वीय सूत्र
$\psi(t_1, \dots, t_k)$ ने बदलल्यावर मिळणारे सूत्र. हा शेवटचा नियम
{\em गुणधर्माधारित संबंधरचना-रूपबंध} असण्याशी समतुल्य आहे; म्हणजे
\[\lexists[R][\lforall[x_1, \dots, x_k][(\varphi(x_1, \dots, x_k) \liff
\Atom{R}{x_1, \dots, x_k})]],
\]
या रूपाचे, द्वितीय-क्रम भाषेतील प्रत्येक सूत्र~$\varphi$ साठी एक,
स्वयंसिद्धक असते; त्यात $R$ हे मुक्त चर नसते. (सराव: $R$ ला $\varphi$ मध्ये येऊ
दिल्यास हा रूपबंध विसंगत होतो, हे दाखवा!)

तर्कशास्त्रज्ञ “द्वितीय-क्रम तर्कशास्त्राची स्वयंसिद्धके” म्हणतात तेव्हा
सहसा प्रथम-क्रम तर्कशास्त्राचा द्वितीय-क्रम संख्यापक-नियमांनी केलेला किमान
विस्तार आणि गुणधर्माधारित संबंधरचना-रूपबंध त्यांना अभिप्रेत असतो. पण या
स्वयंसिद्धकांच्या आणि नियमांच्या दुर्बल उपप्रणालींचा अभ्यास करणेही अनेकदा
उपयुक्त असते. उदाहरणार्थ, गुणधर्माधारित संबंधरचनेचा स्वयंसिद्धक-रूपबंध पूर्ण
सामान्यतेत \emph{स्वसमावेशक} आहे: द्वितीय-क्रम संख्यापक असलेल्या
सूत्र ने “परिभाषित” होणारा $\Atom{R}{x_1, \dots, x_k}$ हा संबंध
अस्तित्वात आहे असे विधान करण्याची तो मुभा देतो; आणि हे संख्यापक अशा सर्व
संबंधांच्या संचावर मान घेतात---ज्या संचात $R$ स्वतःही येतो!{} विसाव्या
शतकाच्या सुमारास रसेलच्या विरोधापत्तीला दिलेली एक सामान्य प्रतिक्रिया म्हणजे
अशा परिभाषांनाच तिच्यासाठी दोषी धरणे आणि गणिताचा पाया विकसित करताना त्या
टाळणे. सूत्र~$\varphi$ मध्ये द्वितीय-क्रम संख्यापक वापरण्यास मनाई केली, तर
गुणधर्माधारित संबंधरचनेचे काहीसे दुर्बल असे {\em स्वसमावेशरहित} रूप मिळते.

चिन्हार्थाच्या दृष्टीने द्वितीय-क्रम संरचनांमध्ये त्या भाषेसाठीची एक
प्रथम-क्रम संरचना आणि ज्या संबंधांवर द्वितीय-क्रम संख्यापक मान घेतात
असा त्या प्रांत वरील संबंधांचा संच असतो (अधिक नेमके सांगायचे, तर प्रत्येक
$k$ साठी स्थानसंख्या~$k$ असलेल्या संबंधांचा एक संच असतो). अर्थात,
निष्पत्ती प्रणालीत गुणधर्माधारित संबंधरचना समाविष्ट असेल, तर
“द्वितीय-क्रम भागात” ती स्वयंसिद्धके पूर्ण करण्याइतके संबंध असावेत ही
अतिरिक्त अट येते---नाहीतर निष्पत्ती प्रणाली निर्दोष राहत नाही!{} पुरेसे
संबंध उपलब्ध असल्याची खात्री करण्याचा एक सोपा मार्ग म्हणजे द्वितीय-क्रम भागात
प्रथम-क्रम भागावरील \emph{सर्व} संबंध घेणे. अशा संरचना ला
\emph{पूर्ण} म्हणतात आणि एका अर्थाने तीच भाषेची “अभिप्रेत संरचना”
असते. आपण फक्त पूर्ण संरचनेचा विचार केला, तर त्याला
\emph{पूर्ण} द्वितीय-क्रम चिन्हार्थमीमांसा म्हणतात. अशा वेळी संरचना
निर्दिष्ट करणे म्हणजे फक्त तिचा प्रथम-क्रम भाग निर्दिष्ट करणे; कारण
द्वितीय-क्रम भागातील घटक त्यावरून अप्रत्यक्षपणे ठरतात.

सारांशतः, द्वितीय-क्रम तर्कशास्त्राविषयी बोलताना काही संदिग्धता असते.
निष्पत्ती प्रणालीच्या दृष्टीने पुढीलपैकी कोणतीही गोष्ट अभिप्रेत असू शकते:
\begin{enumerate}
\item काही गुणधर्माधारित संबंधरचना-स्वयंसिद्धकांसह “किमान” द्वितीय-क्रम
  निष्पत्ती प्रणाली.
\item पूर्ण गुणधर्माधारित संबंधरचना असलेली “मानक” द्वितीय-क्रम
  निष्पत्ती प्रणाली.
\end{enumerate}
चिन्हार्थाच्या दृष्टीने पुढीलपैकी कशातही स्वारस्य असू शकते:
\begin{enumerate}
\item “दुर्बल” चिन्हार्थमीमांसा, जिच्यात संरचनांमध्ये प्रथम-क्रम
  भाग आणि गुणधर्माधारित संबंधरचना-स्वयंसिद्धके पूर्ण करण्याइतका मोठा
  द्वितीय-क्रम भाग असतो.
\item “मानक” द्वितीय-क्रम चिन्हार्थमीमांसा, जिच्यात फक्त पूर्ण
  संरचनेचा विचार केला जातो.
\end{enumerate}
तर्कशास्त्रज्ञांनी कोणती निष्पत्ती प्रणाली किंवा कोणती चिन्हार्थमीमांसा
अभिप्रेत आहे हे सांगितले नाही, तर सहसा प्रत्येक यादीतील दुसरी बाब त्यांना
अभिप्रेत असते. पुढे दिसेल त्याप्रमाणे, ही चिन्हार्थमीमांसा वापरण्याचा फायदा असा
की तिच्यात अनेक स्वाभाविक गणितीय रचनांची समरूपतेपर्यंत एकमेव वर्णने मिळतात;
त्याच वेळी निष्पत्ती प्रणाली बरीच समर्थ असून या चिन्हार्थमीमांसेसाठी
निर्दोष असते. मात्र निष्पत्ती प्रणाली या चिन्हार्थमीमांसेसाठी
\emph{संपूर्ण नाही}; इतकेच नव्हे, तर प्रभावी रीतीने दिलेली \emph{कोणतीही}
निष्पत्ती प्रणाली पूर्ण द्वितीय-क्रम चिन्हार्थमीमांसेसाठी संपूर्ण नसते.
दुसरीकडे, दुर्बल चिन्हार्थमीमांसेसाठी निष्पत्ती प्रणाली \emph{संपूर्ण
आहे}, हे आपण पाहू; म्हणून एखादे वाक्य सिद्ध करता येत नसेल, तर ते असत्य असणारी
\emph{एखादी तरी} संरचना असते, जरी ती पूर्ण रचना असेलच असे नाही.

द्वितीय-क्रम तर्कशास्त्राची भाषा बरीच अभिव्यक्तिसमर्थ आहे. एकस्थानी संबंधांची
प्रांताच्या उपसंचांशी ओळख करता येते; म्हणून विशेषतः या संचांवर संख्यापनही
करता येते. उदाहरणार्थ, स्वाभाविक संख्यांसाठीचे आगमन एकाच स्वयंसिद्धकाने
व्यक्त करता येते:
\[\lforall[R][((\Atom{R}{\Obj{0}} \land \lforall[x][(\Atom{R}{x} \lif
    \Atom{R}{x'})]) \lif \lforall[x][\Atom{R}{x}])].
\]
अंकगणिताच्या भाषेत $\Obj 0, \Obj \prime, +, \times$ आणि $<$ ही चिन्हे घेतली,
तर त्यांचे वर्तन वर्णिण्यासाठी पुढील स्वयंसिद्धके जोडता येतात:
\begin{enumerate}
\item $\lforall[x][\lnot x' = \Obj 0]$
\item $\lforall[x][\lforall[y][(s(x) = s(y) \lif x = y)]]$
\item $\lforall[x][(x + \Obj 0) = x]$
\item $\lforall[x][\lforall[y][(x + y') = (x + y)']]$
\item $\lforall[x][(x \times \Obj 0) = \Obj 0]$
\item $\lforall[x][\lforall[y][(x \times y') = ((x \times y) + x)]]$
\item $\lforall[x][\lforall[y][(x < y \liff \lexists[z][y = (x + z')])]]$
\end{enumerate}
आपण पूर्ण द्वितीय-क्रम चिन्हार्थमीमांसा वापरत असू, तर ही स्वयंसिद्धके आणि
वरील आगमनाचे स्वयंसिद्धक मिळून अंकगणिताचे मानक प्रतिमान असलेल्या
संरचना~$\Struct{N}$ चे समरूपतेपर्यंत एकमेव वर्णन देतात, हे दाखवणे
अवघड नाही. ही स्वयंसिद्धके सत्य असलेली कोणतीही संरचना~$\Struct{M}$
दिली असता, $\Nat$ वरील नेहमीचे पुनरावर्तन वापरून $\Nat$ पासून
$\Struct{M}$ च्या प्रांत कडे जाणारे $f$ हे फलन असे परिभाषित करा की
$f(0) = \Assign{\Obj 0}{M}$ आणि $f(x+1) = \Assign{\prime}{M}(f(x))$.
$\Nat$ वरील नेहमीचे आगमन आणि $\Struct M$ मध्ये स्वयंसिद्धके (1) आणि~(2)
सत्य असण्याचा उपयोग केल्यावर $f$ हे एकास-एक आहे असे दिसते. $f$ हे
आच्छादक आहे हे पाहण्यासाठी, $f$ च्या व्याप्तीत असलेल्या
$\Domain{M}$ मधील घटकांचा संच $P$ घ्या. $\Struct M$ पूर्ण असल्यामुळे $P$
हा द्वितीय-क्रम प्रांत मध्ये असतो. $f$ च्या रचनेवरून
$\Assign{\Obj 0}{M}$ हा $P$ मध्ये असतो आणि $P$ हा
$\Assign{\prime}{M}$ खाली बंद असतो. $\Struct M$ मध्ये आगमनाचे स्वयंसिद्धक
सत्य असणे (विशेषतः $P$ साठी) $P$ हा $\Struct M$ च्या संपूर्ण प्रथम-क्रम
प्रांत इतकाच असल्याची हमी देते. यावरून $f$ हे एकास-एक आच्छादन आहे.
$\Nat$ वर नेहमीचे आगमन वारंवार वापरून $f$ हे रचनारक्षक प्रतिचित्रण आहे हे
दाखवणे यापेक्षा अवघड नाही.

संचसैद्धान्तिक भाषेत फलन हा संबंधाचाच एक विशेष प्रकार असतो; उदाहरणार्थ,
$\lforall[x][\lexists![y][R(x,y)]]$ पूर्ण करणाऱ्या $R$ या द्विस्थानी
संबंधाशी एकस्थानी फलन~$f$ ची ओळख करता येते. त्यामुळे फलनांवरही संख्यापन करता
येते. मग पूर्ण चिन्हार्थमीमांसा वापरून अनंत संरचनेचा वर्ग असा
परिभाषित करता येतो: त्या संरचना $\Struct M$ असतात ज्यांच्यासाठी
$\Struct M$ च्या प्रांत पासून त्याच्याच उचित उपसंचाकडे जाणारे एक
एकास-एक फलन असते:
\[\lexists[f][(\lforall[x][\lforall[y][(\eq[f(x)][f(y)] \lif \eq[x][y])]]
  \land \lexists[y][\lforall[x][\eq/[f(x)][y]]])].
\]
या वाक्याचा निषेध मग सांत संरचनेचा वर्ग परिभाषित करतो.

याखेरीज, रेषीय क्रमाच्या परिभाषेत पुढील अट जोडून सुक्रमांचा वर्ग परिभाषित
करता येतो:
\[\lforall[P][(\lexists[x][\Atom{P}{x}] \lif \lexists[x][(\Atom{P}{x}
    \land \lforall[y][(y < x \lif \lnot \Atom{P}{y})])])].
\]
“संच” आणि “एकस्थानी संबंध” यांची ओळख केल्यास, प्रत्येक रिक्तेतर संचाला
लघुतम घटक असतो असे हे सांगते. आणखी एका उदाहरणात, शिखरांची परस्पर न जोडलेल्या
भागांत कोणतीही अ-तुच्छ विभागणी नसते असे सांगून आलेखांची संयोजकता व्यक्त करता
येते:
\[\lnot \lexists[A][(\lexists[x][A(x)] \land \lexists[y][\lnot A(y)]
  \land \lforall[w][\lforall[z][((\Atom{A}{w} \land \lnot \Atom{A}{z})
      \lif \lnot \Atom{R}{w,z})]])].
\]
दुसऱ्या उदाहरणासाठी, ज्या सांत संरचनेच्या प्रांताचा आकार सम
आहे त्यांचा वर्ग परिभाषित करण्याचा सराव तुम्ही करू शकता. आणखी लक्षवेधक
म्हणजे, परिमेय संख्या अंतर्भूत करणारे संपूर्ण क्रमित क्षेत्र म्हणून वास्तव
संख्यांचे समरूपतेपर्यंत एकमेव वर्णन देता येते.

थोडक्यात, द्वितीय-क्रम तर्कशास्त्र प्रथम-क्रम तर्कशास्त्रापेक्षा खूपच अधिक
अभिव्यक्तिसमर्थ आहे. ही चांगली बाजू; आता अडचण पाहू. पूर्ण द्वितीय-क्रम
चिन्हार्थमीमांसेसाठी संपूर्ण असणारी प्रभावी निष्पत्ती प्रणाली नाही, हे
आपण आधीच नमूद केले आहे. प्रथम-क्रम तर्कशास्त्राचे अनेक गुणधर्म येथे उपलब्ध
नसतात; त्यांत संहतता आणि लोव्हेनहाइम--स्कोलेम प्रमेयांचा समावेश होतो.

दुसरीकडे, पूर्ण द्वितीय-क्रम चिन्हार्थमीमांसा सोडून दुर्बल चिन्हार्थमीमांसा
स्वीकारली, तर किमान द्वितीय-क्रम निष्पत्ती प्रणाली या
चिन्हार्थमीमांसेसाठी संपूर्ण असते. दुसऱ्या शब्दांत, $\Proves$ चा अर्थ
“किमान प्रणालीत सिद्ध होते” आणि $\Entails$ चा अर्थ “दुर्बल
चिन्हार्थमीमांसेत तार्किकरीत्या निष्पन्न होते” असा घेतल्यास, $\Gamma \Entails
\varphi$ असेल तेव्हा $\Gamma \Proves \varphi$ असते, हे दाखवता येते.
निष्पत्ती प्रणालीत ठरावीक गुणधर्माधारित संबंधरचना-स्वयंसिद्धके समाविष्ट
करायची असतील, तर ही स्वयंसिद्धके पूर्ण करणाऱ्या द्वितीय-क्रम रचनांपुरती
चिन्हार्थमीमांसा मर्यादित करावी लागते: उदाहरणार्थ, $\Delta$ हा
गुणधर्माधारित संबंधरचना-स्वयंसिद्धकांचा संच (शक्यतो त्यांपैकी सर्वांचा संच)
असेल, तर $\Gamma \cup \Delta \Entails \varphi$ असल्यास $\Gamma \cup \Delta
\Proves \varphi$ असते. विशेषतः, आपण विचारात घेत असलेली गुणधर्माधारित
संबंधरचना-स्वयंसिद्धके वापरून $\varphi$ सिद्ध करता येत नसेल, तर ती स्वयंसिद्धके
सत्य असूनही $\lnot \varphi$ सत्य असणारे एक प्रतिमान असते.

दुर्बल चिन्हार्थमीमांसेसाठी संपूर्णता प्रमेय का लागू होते हे पाहण्याचा सर्वांत
सोपा मार्ग म्हणजे द्वितीय-क्रम तर्कशास्त्राचा पुढीलप्रमाणे बहुवर्गीय
तर्कशास्त्र म्हणून विचार करणे. एका वर्गाचा अर्थ नेहमीचा “प्रथम-क्रम”
प्रांत असा लावला जातो; मग प्रत्येक $k$ साठी “स्थानसंख्या $k$ असलेल्या
संबंधांचा” प्रांत असतो. भाषेत
“$\Atom{\Obj{true}_k}{R,x_1,\dots,x_k}$” ही अंगभूत संबंधचिन्हे घेतो. येथे
$R$ हे “$k$-स्थानी संबंध” या वर्गाचे चर आणि $x_1$, \dots,~$x_k$ या
प्रथम-क्रम वर्गातील वस्तू असताना, $R$ हा संबंध $x_1$, \dots,~$x_k$ यांना
लागू होतो असे हे चिन्ह सांगते.

ही ओळख स्वीकारल्यावर दुर्बल द्वितीय-क्रम चिन्हार्थमीमांसा मूलतः बहुवर्गीय
तर्कशास्त्राच्या नेहमीच्या चिन्हार्थमीमांसेसारखीच असते; आणि बहुवर्गीय
तर्कशास्त्राचे प्रथम-क्रम तर्कशास्त्रात अंतःस्थापन करता येते, हे आपण आधीच
पाहिले आहे. त्यामुळे दोन्ही दिशांतील रूपांतरणांचा विचार करता, द्वितीय-क्रम
तर्कशास्त्राची दुर्बल संकल्पना हे प्रत्यक्षात प्रथम-क्रम तर्कशास्त्राचेच
छद्मरूप आहे; त्यात प्रांत मध्ये योग्य स्वयंसिद्धकांच्या अधीन असलेल्या
“वस्तू” आणि “संबंध” या दोन्हींचा समावेश असतो.












\section{उच्च-क्रम तर्कशास्त्र}\label{fol:byd:hol:sec}

प्रथम-क्रम तर्कशास्त्रातून द्वितीय-क्रम तर्कशास्त्राकडे गेल्यामुळे औपचारिक
भाषेच्या चौकटीत प्रथम-क्रम प्रांत मधील वस्तूंच्या संचांविषयी बोलणे शक्य
झाले. तेथेच का थांबावे? उदाहरणार्थ, तृतीय-क्रम तर्कशास्त्रामुळे वस्तूंच्या
संचांचे संच, किंवा वस्तू आणि वस्तूंचे संच हे दोन्ही अंतर्भूत करणारे संचही
हाताळता यावेत. चतुर्थ-क्रम तर्कशास्त्रामुळे अशा प्रकारच्या वस्तूंच्या
संचांविषयी बोलता येईल. तुमच्या लक्षात आले असेलच की ही कल्पना कितीही वेळा
पुनरावृत्त करता येते.

व्यवहारात उच्च-क्रम तर्कशास्त्र संबंधांऐवजी फलनांच्या मदतीने
सूत्रted केले जाते. (स्वाभाविक ओळखी गृहीत धरल्यास हा फरक महत्त्वाचा
नाही.) काही मूलभूत “वर्ग” $A$, $B$, $C$,~\dots (ज्यांना आता आपण “प्रकार”
म्हणू) दिले असता, पुढील अट घालून नवे प्रकार तयार करता येतात:
\begin{quote}
$\sigma$ आणि $\tau$ हे सांत प्रकार असतील, तर $\sigma \to \tau$ हाही सांत
प्रकार असतो.
\end{quote}
प्रकारांना विन्यासात्मक “लेबले” समजा; आपल्या प्रांत मध्ये हव्या असलेल्या
वस्तूंचे ती वर्गीकरण करतात. $\sigma \to \tau$ हा प्रकार, प्रकार~$\sigma$ च्या
वस्तू घेऊन प्रकार~$\tau$ च्या वस्तू देणारी फलने असलेल्या वस्तूंचे वर्णन करतो.
उदाहरणार्थ, “सत्य” आणि “असत्य” या सत्यतामूल्यांचा $\Omega$ हा प्रकार आणि
नैसर्गिक संख्यांचा $\Nat$ हा प्रकार घ्यावासा वाटू शकतो. अशा वेळी $\Nat \to
\Omega$ प्रकारच्या वस्तू एकस्थानी संबंध किंवा $\Nat$ चे उपसंच मानता येतात;
$\Nat \to \Nat$ प्रकारच्या वस्तू ही नैसर्गिक संख्यांपासून नैसर्गिक संख्यांकडे
जाणारी फलने असतात; आणि $(\Nat \to \Nat) \to \Nat$ प्रकारच्या वस्तू
“फलनके” असतात, म्हणजे फलने घेऊन संख्या देणारी उच्च-प्रकार फलने.

द्वितीय-क्रम तर्कशास्त्राप्रमाणेच, उच्च-क्रम तर्कशास्त्राकडे बहुवर्गीय
तर्कशास्त्राचा एक प्रकार म्हणून पाहता येते; त्यात विचारात घ्यायच्या प्रत्येक
वस्तुप्रकारासाठी एक वर्ग असतो. पण उच्च-प्रकार तर्कशास्त्राचा विन्यास
मुळापासूनच परिभाषित करणे सहसा अधिक स्पष्ट ठरते. उदाहरणार्थ, सांत प्रकारांचा
संच पुढीलप्रमाणे विगमनाने परिभाषित करता येतो:
\begin{enumerate}
\item $\Nat$ हा सांत प्रकार आहे.
\item $\sigma$ आणि $\tau$ हे सांत प्रकार असतील, तर $\sigma
  \to \tau$ हाही सांत प्रकार आहे.
\item $\sigma$ आणि $\tau$ हे सांत प्रकार असतील, तर $\sigma \times
  \tau$ हाही सांत प्रकार आहे.
\end{enumerate}
सहज अर्थाने, $\Nat$ नैसर्गिक संख्यांचा प्रकार दर्शवतो, $\sigma \to \tau$
हा $\sigma$ पासून $\tau$ कडे जाणाऱ्या फलनांचा प्रकार दर्शवतो आणि
$\sigma \times \tau$ हा वस्तूंच्या जोड्यांचा प्रकार दर्शवतो, ज्यातील एक वस्तू
$\sigma$ मधील आणि दुसरी $\tau$ मधील असते. मग पदांचा संच पुढीलप्रमाणे
विगमनाने परिभाषित करता येतो:
\begin{enumerate}
\item प्रत्येक प्रकार $\sigma$ साठी $x$, $y$, $z$, \dots ही प्रकार~$\sigma$
  ची अनेक चरे असतात.
\item $\Obj 0$ हे $\Nat$ प्रकारचे पद आहे.
\item $\Obj S$ (उत्तरवर्ती) हे $\Nat \to \Nat$ प्रकारचे पद आहे.
\item $s$ हे $\sigma$ प्रकारचे पद आणि $t$ हे
  $\Nat \to (\sigma \to \sigma)$ प्रकारचे पद असेल, तर $\Obj{R}_{st}$ हे
  $\Nat \to \sigma$ प्रकारचे पद आहे.
\item $s$ हे $\tau \to \sigma$ प्रकारचे पद आणि $t$ हे प्रकार~$\tau$ चे पद
  असेल, तर $s(t)$ हे $\sigma$ प्रकारचे पद आहे.
\item $s$ हे प्रकार~$\sigma$ चे पद आणि $x$ हे प्रकार~$\tau$ चे चर असेल,
  तर $\lambd[x][s]$ हे $\tau \to \sigma$ प्रकारचे पद आहे.
\item $s$ हे प्रकार~$\sigma$ चे पद आणि $t$ हे प्रकार~$\tau$ चे पद असेल,
  तर $\tuple{s, t}$ हे $\sigma \times \tau$ प्रकारचे पद आहे.
\item $s$ हे प्रकार~$\sigma \times \tau$ चे पद असेल, तर $p_1(s)$ हे
  प्रकार~$\sigma$ चे पद आणि $p_2(s)$ हे प्रकार~$\tau$ चे पद आहे.
\end{enumerate}
सहज अर्थाने, $\Obj{R}_{st}$ हे पुढीलप्रमाणे पुनरावर्ती रीतीने परिभाषित केलेले
फलन दर्शवते:
\begin{align*}
\Obj{R}_{st}(0) & = s \\
\Obj{R}_{st}(x+1) & = t(x, R_{st}(x)),
\end{align*}
$\tuple{s, t}$ ही ज्याचा पहिला घटक~$s$ आणि दुसरा घटक~$t$ आहे अशी जोडी
दर्शवते, तर $p_1(s)$ आणि~$p_2(s)$ हे $s$ चे पहिले व दुसरे घटक
(“प्रक्षेपणे”) दर्शवतात. शेवटी, $\lambd[x][s]$ हे
\[f(x) = s
\]
या अटीने प्रकार~$\tau$ च्या कोणत्याही~$x$ साठी परिभाषित झालेले फलन~$f$
दर्शवते.

म्हणून बाब (6) आपल्याला गुणधर्माधारित फलनरचनेचे एक रूप देते आणि
पदांच्या मदतीने फलने परिभाषित करणे शक्य करते. समान प्रकारच्या पदांमधील
एकरूपता विधाने $\eq[s][t]$, नेहमीचे विधानीय संयोजक आणि उच्च-प्रकार
संख्यापन यांपासून सूत्रे घडवली जातात. वर परिभाषित केलेल्या पदांचे
नियमन करणारी मूलभूत समीकरणे, तसेच संख्यापक आणि एकरूपता असलेले
तर्कशास्त्राचे नेहमीचे नियम, ही मग प्रणालीची स्वयंसिद्धके मानता येतात.

सांत प्रकारांच्या प्रणालीत सत्यतामूल्यांचा $\Omega$ हा प्रकार जोडला, तर
त्याचा वापर नियंत्रित करणारी स्वयंसिद्धकेही समाविष्ट करावी लागतात. खरे तर
योग्य युक्ती वापरून गुंतागुंतीची सूत्रे पूर्णपणे काढून टाकता येतात आणि
त्यांच्या जागी $\Omega$ प्रकारची पदे ठेवता येतात!{} त्यानुसार सिद्धता-पद्धतीतही
बदल करता येतो. त्यातून मूलतः १९३० च्या दशकात ॲलोन्झो चर्च यांनी मांडलेली
\emph{प्रकारांची साधी उपपत्ती} मिळते.

द्वितीय-क्रम तर्कशास्त्राप्रमाणेच, उच्च-प्रकार चिन्हार्थमीमांसेच्याही विविध
आवृत्त्या वापरता येतात. पूर्ण आवृत्तीत $\sigma \to \tau$ प्रकारची चरे,
प्रकार~$\sigma$ च्या वस्तूंपासून प्रकार~$\tau$ च्या वस्तूंकडे जाणाऱ्या
\emph{सर्व} फलनांच्या संचावर मान घेतात. अपेक्षेप्रमाणे ही चिन्हार्थमीमांसा
इतकी समर्थ आहे की तिच्यासाठी संपूर्ण आणि प्रभावी निष्पत्ती प्रणाली
असू शकत नाही. पण दुर्बल चिन्हार्थमीमांसेचाही विचार करता येतो; त्यात
संरचनांमध्ये प्रत्येक प्रकार $\tau$ साठी $T_\tau$ हे घटकांचे संच
आणि उपयोजन, प्रक्षेपण इत्यादींसाठी योग्य कृत्ये असतात. तपशील योग्य रीतीने
मांडल्यास वर वर्णिलेल्या प्रकारच्या निष्पत्ती प्रणालींसाठी संपूर्णता
प्रमेये मिळवता येतात.

उच्च-प्रकार तर्कशास्त्र आकर्षक आहे, कारण गणिताचा मोठा भाग स्वाभाविक रीतीने
अंतःस्थापित करता येईल अशी चौकट ते देते: $\Nat$ पासून सुरुवात करून वास्तव
संख्या, संतत फलने इत्यादी परिभाषित करता येतात. अंतःप्रज्ञावादी तर्कशास्त्राच्या
संदर्भातही ते विशेष आकर्षक आहे, कारण प्रकारांना स्पष्ट “रचनाशील” अर्थनिर्धारणे
आहेत. खरे तर, उच्च-प्रकार चिन्हार्थमीमांसेच्या रचनाशील आवृत्त्या
(अभिजात तर्कशास्त्राऐवजी अंतःप्रज्ञावादी तर्कशास्त्रावर आधारलेल्या) विकसित
करता येतात. त्या ही रचनाशील अर्थनिर्धारणे अतिशय स्पष्ट करतात आणि अनेक
दृष्टींनी त्यांच्या अभिजात समकक्षांपेक्षा अधिक रोचक असतात.












\section{अंतःप्रज्ञावादी तर्कशास्त्र}\label{fol:byd:il:sec}

द्वितीय-क्रम आणि उच्च-क्रम तर्कशास्त्राच्या उलट, अंतःप्रज्ञावादी प्रथम-क्रम
तर्कशास्त्र हे अभिजात आवृत्तीचे एक निर्बंधित रूप आहे; अधिक “रचनाशील” प्रकारच्या
तर्कविचाराचे प्रतिमान देणे हा त्याचा उद्देश आहे. त्यामागील काही प्रेरणा पुढील
उदाहरणांनी स्पष्ट होतील.

समजा, एखाद्या दिवशी कोणीतरी तुमच्याकडे येऊन सांगितले की त्यांनी अशी नैसर्गिक
संख्या~$x$ ठरवली आहे की $x$ अविभाज्य असेल तर रीमान गृहीतक सत्य असते आणि $x$
संयुक्त असेल तर रीमान गृहीतक असत्य असते. फारच चांगली बातमी!{} रीमान गृहीतक
सत्य आहे की नाही हा गणितातील मोठ्या अनिर्णित प्रश्नांपैकी एक आहे; आणि येथे तर
ही समस्या केवळ गणनेच्या समस्येत, म्हणजे एखादी ठरावीक संख्या अविभाज्य आहे की
नाही हे ठरवण्यात, रूपांतरित झाल्यासारखी दिसते.

$x$ चे ते अद्भुत मूल्य कोणते? ते त्याचे पुढीलप्रमाणे वर्णन करतात: रीमान गृहीतक
सत्य असेल तर $x$ ही $7$ इतकी आणि अन्यथा $9$ इतकी नैसर्गिक संख्या आहे.

तुम्ही चिडून आपले पैसे परत मागता. अभिजात दृष्टिकोनातून वरील वर्णन $x$ चे
एकमेव मूल्य खरोखर ठरवते; पण तुम्हाला प्रत्यक्षात हवे आहे ते
\emph{स्पष्टपणे} दिलेले $x$ चे मूल्य.

आणखी एक, कदाचित कमी कृत्रिम, उदाहरण पाहू. अपरिमेय संख्येचा परिमेय घात घेऊन
परिमेय उत्तर मिळू शकते, हे आपल्याला माहीत आहे. उदाहरणार्थ,
$\sqrt{2}^2 = 2$. अपरिमेय संख्येचा \emph{अपरिमेय} घात घेऊन परिमेय उत्तर
मिळू शकते की नाही, हे तितके स्पष्ट नाही. पुढील प्रमेय याचे होकारार्थी उत्तर
देते:

\begin{thm}
$a^b$ परिमेय असेल अशा $a$ आणि $b$ या अपरिमेय संख्या आहेत.
\end{thm}

\begin{proof}
$\sqrt{2}^{\sqrt{2}}$ चा विचार करा. ही संख्या परिमेय असेल, तर आपले काम झाले:
$a = b = \sqrt{2}$ घेता येईल. अन्यथा ती अपरिमेय आहे. मग
\[(\sqrt{2}^{\sqrt{2}})^{\sqrt{2}} = \sqrt{2}^{\sqrt{2} \cdot
  \sqrt{2}} = \sqrt{2}^2 = 2,
\]
आणि ही संख्या निश्चितच परिमेय आहे. म्हणून या प्रकरणात $a$ म्हणून
$\sqrt{2}^{\sqrt{2}}$ आणि $b$ म्हणून~$\sqrt 2$ घ्या.
\end{proof}

याला वैध सिद्धता म्हणता येईल का? बहुतेक गणितज्ञांच्या मते होय. पण यात पुन्हा
काहीसे असमाधानकारक असे काही आहे: विशिष्ट गुणधर्म असलेल्या वास्तव संख्यांच्या
जोडीचे अस्तित्व आपण सिद्ध केले, पण ती संख्यांची \emph{कोणती} जोडी आहे हे
सांगता आलेले नाही. हाच निष्कर्ष अशा रीतीने सिद्ध करता येतो की सिद्धतेत
$a$, $b$ ही जोडी \emph{प्रत्यक्ष दिलेली} असते: $a = \sqrt{3}$ आणि
$b = \log_3 4$ घ्या. मग
\[a^b = \sqrt{3}^{\log_3 4} = 3^{1/2 \cdot \log_3 4} = (3^{\log_3
  4})^{1/2} = 4^{1/2}= 2,
\]
कारण $3^{\log_3 x} = x$.

पहिल्या सिद्धतेतील पाऊल अमान्य असणाऱ्या तर्कविचाराचे प्रतिमान देण्यासाठी
अंतःप्रज्ञावादी तर्कशास्त्र रचले आहे. $\varphi(x)$ ची पूर्ति करणारा $x$ अस्तित्वात
आहे हे सिद्ध करणे म्हणजे दुसऱ्या सिद्धतेप्रमाणे एखादा ठरावीक~$x$ आणि तो
$\varphi$ ची पूर्ति करतो याची सिद्धता द्यावी लागते. $\varphi$ किंवा $\psi$ सत्य आहे हे
सिद्ध करण्यासाठी त्यांपैकी एकाची सिद्धता देता आली पाहिजे.

आकारिक भाषेत सांगायचे, तर प्रथम-क्रम तर्कशास्त्राची निष्पत्ती प्रणाली
ठरावीक रीतीने निर्बंधित केल्यावर अंतःप्रज्ञावादी प्रथम-क्रम तर्कशास्त्र मिळते.
त्याचप्रमाणे द्वितीय-क्रम किंवा उच्च-क्रम तर्कशास्त्राच्याही अंतःप्रज्ञावादी
आवृत्त्या असतात. गणितीय दृष्टिकोनातून या केवळ आकारिक निगमन-पद्धती आहेत; पण
आधी नमूद केल्याप्रमाणे, एका प्रकारच्या गणितीय तर्कविचाराचे प्रतिमान देणे हा
त्यांचा उद्देश आहे. हा तर्कविचार गणिताविषयीच्या एखाद्या तात्त्विक दृष्टिकोनाने
(उदाहरणार्थ, ब्रौवर यांच्या अंतःप्रज्ञावादाने) समर्थित मानता येतो; किंवा
बिशप यांच्या रचनाशीलतावादाच्या धर्तीवर अधिक “मूर्त” आणि समाधानकारक असा
गणितीय तर्कविचार मानता येतो; तसेच आकारिक वर्णन अनौपचारिक प्रेरणा खरोखर
पकडते की नाही यावर वादही घालता येतो. पण आपली तात्त्विक भूमिका कोणतीही असो,
अंतःप्रज्ञावादी तर्कशास्त्राचा आकारिक रीतीने मांडलेले तर्कशास्त्र म्हणून
अभ्यास करता येतो; आणि कारणे कोणतीही असोत, अनेक गणितीय तर्कशास्त्रज्ञांना तो
अभ्यास रंजक वाटतो.

अंतःप्रज्ञावादी संयोजकांचे एक अनौपचारिक रचनाशील अर्थनिर्धारण आहे. ते सहसा
BHK अर्थनिर्धारण म्हणून ओळखले जाते (ब्रौवर, हीटिंग आणि कोल्मोगोरोव यांच्या
नावांवरून). ते असे आहे: $\varphi \land \psi$ ची सिद्धता म्हणजे $\varphi$ ची सिद्धता
आणि $\psi$ ची सिद्धता यांची जोडी; $\varphi \lor \psi$ ची सिद्धता म्हणजे $\varphi$ ची
सिद्धता किंवा $\psi$ ची सिद्धता, तसेच यांपैकी नेमके कोणते प्रकरण आहे याची स्पष्ट
माहिती; $\varphi \lif \psi$ ची सिद्धता म्हणजे $\varphi$ च्या सिद्धतेचे $\psi$ च्या
सिद्धतेत रूपांतर करणारी प्रक्रिया; $\lforall[x][\varphi(x)]$ ची सिद्धता म्हणजे
$x$ च्या कोणत्याही मूल्यासाठी $\varphi(x)$ ची सिद्धता देणारी प्रक्रिया; आणि
$\lexists[x][\varphi(x)]$ ची सिद्धता म्हणजे $x$ चे एखादे मूल्य आणि ते मूल्य $\varphi$
ची पूर्ति करते याची सिद्धता. या अर्थनिर्धारणाचे संगणकीय वर्णन
“करी--हॉवर्ड समरूपण” किंवा “सूत्रे-हे-प्रकार प्रतिमान” या नावांनी करता
येते: सूत्र एखादा विशिष्ट दत्तांश-प्रकार ठरवते असे समजा आणि
सिद्धतांना त्या दत्तांश-प्रकारांच्या संगणकीय वस्तू समजा; या वस्तूंमुळे संबंधित
सूत्र सत्य आहे हे आपल्याला दिसते.

अंतःप्रज्ञावादी तर्कशास्त्राकडे अनेकदा “विमध्य सिद्धांत वजा केल्यावर उरणारे”
अभिजात तर्कशास्त्र म्हणून पाहिले जाते. पुढील प्रमेय हे अधिक नेमके करते.
\begin{thm}
अंतःप्रज्ञावादी तर्कशास्त्रात पुढील स्वयंसिद्धक-रूपबंध समतुल्य आहेत:
\begin{enumerate}
\item $(\lnot \varphi \lif \lfalse) \lif \varphi$.
\item $\varphi \lor \lnot \varphi$
\item $\lnot \lnot \varphi \lif \varphi$
\end{enumerate}
\end{thm}
कोणत्याही एका रूपबंधापासून इतर दोन्हींचे दृष्टांत मिळवणे हा अंतःप्रज्ञावादी
तर्कशास्त्रातील चांगला सराव आहे.

ब्रौवर यांच्या अंतःप्रज्ञावादाचे आकारीकरण म्हणून मांडलेल्या अंतःप्रज्ञावादी
विधानीय तर्कशास्त्राच्या पहिल्या निगमन-पद्धती कोल्मोगोरोव, ग्लिवेन्को आणि
हीटिंग यांनी स्वतंत्रपणे दिल्या. अंतःप्रज्ञावादी प्रथम-क्रम तर्कशास्त्राचे
(आणि अंतःप्रज्ञावादी गणिताच्या काही भागांचे) पहिले आकारीकरण हीटिंग यांनी
केले. अभिजात तर्कशास्त्रात वैध असलेले अनेक रूपबंध अंतःप्रज्ञावादी
तर्कशास्त्रात वैध नसले, तरी त्यांपैकी अनेक वैधही असतात.

\emph{द्विनकरण रूपांतरण} अभिजात आणि अंतःप्रज्ञावादी तर्कशास्त्रांतील एक
महत्त्वाचा संबंध वर्णिते. ते पुढीलप्रमाणे विगमनाने परिभाषित केले जाते
($\varphi^N$ हे अभिजात सूत्र~$\varphi$ चे “अंतःप्रज्ञावादी” रूपांतर समजा):
\begin{align*}
\varphi^N & \ident \lnot\lnot \varphi \quad \text{अण्वीय सूत्रे $\varphi$ साठी} \\
(\varphi \land \psi)^N & \ident (\varphi^N \land \psi^N) \\
(\varphi \lor \psi)^N & \ident  \lnot\lnot (\varphi^N \lor \psi^N) \\
(\varphi \lif \psi)^N & \ident (\varphi^N \lif \psi^N) \\
(\lforall[x][\varphi])^N & \ident \lforall[x][\varphi^N] \\
(\lexists[x][\varphi])^N & \ident \lnot\lnot\lexists[x][\varphi^N]
\end{align*}
कोल्मोगोरोव आणि ग्लिवेन्को यांनी विधानीय तर्कशास्त्रासाठी या रूपांतरणाच्या
आवृत्त्या दिल्या होत्या; विधेय तर्कशास्त्रासाठी ते गोडेल आणि गेंटझेन यांनी
स्वतंत्रपणे दिले. आपल्याला पुढील निष्कर्ष मिळतो.

\begin{thm}
\begin{enumerate}
\item $\varphi \liff \varphi^N$ हे अभिजात तर्कशास्त्रात सिद्धतायोग्य आहे.
\item $\varphi$ हे अभिजात तर्कशास्त्रात सिद्धतायोग्य असेल, तर $\varphi^N$ हे
  अंतःप्रज्ञावादी तर्कशास्त्रात सिद्धतायोग्य आहे.
\end{enumerate}
\end{thm}

आता पुढील संवादाची कल्पना करता येते. अभिजात गणितज्ञ: “मी $\varphi$ सिद्ध केले
आहे!” अंतःप्रज्ञावादी गणितज्ञ: “तुमची सिद्धता वैध नाही. तुम्ही प्रत्यक्षात
$\varphi^N$ सिद्ध केले आहे.” अभिजात गणितज्ञ: “मला तेही चालेल!” अभिजात
गणितज्ञाच्या दृष्टीने अंतःप्रज्ञावादी गणितज्ञ क्षुल्लक भेद करत आहे, कारण ती
दोन्ही समतुल्य आहेत. पण अंतःप्रज्ञावादी गणितज्ञाच्या मते त्यांत फरक आहे.

लक्षात घ्या की वरील रूपांतरण फक्त शुद्ध तर्कशास्त्राशी संबंधित आहे; अभिजात
आणि अंतःप्रज्ञावादी गणितासाठी योग्य \emph{अतार्किक} स्वयंसिद्धके कोणती किंवा
त्यांच्यात कोणता संबंध आहे, या प्रश्नाला ते हात घालत नाही. पण वरील प्रमेयाचा
पुढील किंचित विस्तार काही उपयुक्त माहिती देतो:

\begin{thm}
$\Gamma$ ने $\varphi$ अभिजात रीतीने सिद्ध होत असेल, तर $\Gamma^N$ ने $\varphi^N$
अंतःप्रज्ञावादी रीतीने सिद्ध होते.
\end{thm}

दुसऱ्या शब्दांत, काही गृहीतकांपासून $\varphi$ अभिजात रीतीने सिद्धतायोग्य असेल,
तर त्यांच्या द्विनकरण रूपांतरांपासून $\varphi^N$ अंतःप्रज्ञावादी रीतीने
सिद्धतायोग्य असते.

एखादे वाक्य किंवा विधानीय सूत्र अंतःप्रज्ञावादी रीतीने वैध आहे हे
दाखवण्यासाठी फक्त त्याची सिद्धता द्यावी लागते. पण ते वैध नाही हे कसे
दाखवणार? त्यासाठी आपल्याला निर्दोष, आणि शक्यतो संपूर्ण, चिन्हार्थमीमांसा हवी.
क्रिप्के यांनी दिलेली चिन्हार्थमीमांसा हे काम उत्तम रीतीने करते.

अभिजात तर्कशास्त्रासाठी केलेला खेळच येथेही खेळता येतो: चिन्हार्थमीमांसा
परिभाषित करायची आणि निर्दोषता व संपूर्णता सिद्ध करायची. मात्र पुढील भेद लक्षात
घेण्यासारखा आहे. अभिजात तर्कशास्त्रात चिन्हार्थमीमांसा एका अर्थाने
संयोजकांच्या अर्थातच अंतर्भूत असलेली “उघड” चिन्हार्थमीमांसा होती. क्रिप्के
चिन्हार्थमीमांसेसाठी काही अंतःप्रज्ञ स्पष्टीकरण देता येत असले, तरी तिच्यात
तशी अपरिहार्यता जाणवत नाही. शिवाय, अभिजात संरचना ही स्वाभाविक गणितीय
संकल्पना आहे; म्हणून संरचना ही संकल्पना अभिजात प्रथम-क्रम
तर्कशास्त्राच्या अभ्यासाचे साधन मानता येते, किंवा अभिजात प्रथम-क्रम
तर्कशास्त्र हे गणितीय संरचनेच्या अभ्यासाचे साधन मानता येते.
याउलट, क्रिप्के संरचना कडे फक्त तार्किक रचना म्हणूनच पाहता येते;
त्यांना स्वतंत्र गणितीय महत्त्व आहे असे दिसत नाही.

विधानीय भाषेसाठीची क्रिप्के संरचना~$\mModel M = \tuple{W, R, V}$
मध्ये $W$ हा संच, $W$ वरील लघुतम घटक असलेला $R$ हा अंशतः क्रम आणि
विधानीय चलांचे $W$ च्या घटक ना केलेले “एकस्वनिक” मूल्यांकन असते.
येथे $W$ चे घटक “जगे” किंवा “ज्ञानावस्था” दर्शवतात; $v \geq u$ हा
घटक $u$ ची “संभाव्य भावी अवस्था” दर्शवतो; आणि $u$ ला दिलेली विधानीय चरे ही
$u$ अवस्थेत सत्य असल्याचे माहीत असलेली विधाने असतात. बलन-संबंध
$\mSat{M}{\varphi}[w]$ मग हा संबंध भाषेतील कोणत्याही सूत्रे पर्यंत
विस्तारित करतो; $\mSat{M}{\varphi}[w]$ चा अर्थ “$w$ अवस्थेत $\varphi$ सत्य आहे” असा
वाचा. हा संबंध पुढीलप्रमाणे विगमनाने परिभाषित केला जातो:
\begin{enumerate}
\item $\mSat{M}{\Obj p_i}[w]$ तेव्हा आणि केवळ तेव्हाच, जेव्हा $\Obj p_i$
  हे $w$ ला दिलेल्या विधानीय चरांपैकी एक असते.
\item $\mSat/{M}{\lfalse}[w]$.
\item $\mSat{M}{(\varphi \land \psi)}[w]$ तेव्हा आणि केवळ तेव्हाच, जेव्हा
  $\mSat{M}{\varphi}[w]$ आणि $\mSat{M}{\psi}[w]$.
\item $\mSat{M}{(\varphi \lor \psi)}[w]$ तेव्हा आणि केवळ तेव्हाच, जेव्हा
  $\mSat{M}{\varphi}[w]$ किंवा $\mSat{M}{\psi}[w]$.
\item $\mSat{M}{(\varphi \lif \psi)}[w]$ तेव्हा आणि केवळ तेव्हाच, जेव्हा
  $w' \geq w$ आणि $\mSat{M}{\varphi}[w']$ असताना नेहमी $\mSat{M}{\psi}[w']$.
\end{enumerate}
प्रतिउदाहरण देणारी क्रिप्के संरचना रचून $\lnot (p \land q) \lif
(\lnot p \lor \lnot q)$ हे अंतःप्रज्ञावादी रीतीने वैध नाही हे दाखवण्याचा
प्रयत्न करणे हा चांगला सराव आहे.












\section{मोडल तर्कशास्त्रे}\label{fol:byd:mod:sec}

पुढील सशर्त वाक्याचे उदाहरण पाहा:
\begin{quote}
  जेरेमी त्या खोलीत एकटा असेल, तर त्याने मद्यपान केलेले आहे, तो विवस्त्र आहे
  आणि खुर्च्यांवर नाचत आहे.
\end{quote}
हे सशर्त विधान वास्तविक अभिव्यंजनाच्या अर्थाने सत्य असू शकते, पण तरीही
दिशाभूल करणारे असू शकते; कारण पूर्वांग असत्य आहे किंवा उत्तरांग सत्य आहे
एवढ्यापेक्षा पूर्वांग
आणि निष्कर्ष यांच्यात अधिक दृढ संबंध आहे असे ते सुचवताना दिसते. म्हणजे, हा
दावा फक्त या विशिष्ट जगातच सत्य नाही (जेरेमी खोलीत एकटा नसल्यामुळे येथे तो
क्षुल्लकपणे सत्य असू शकतो), तर पूर्वांग सत्य \emph{असते}, तर निष्कर्षही सत्य
\emph{असता}, असे त्याची शब्दरचना सुचवते. दुसऱ्या शब्दांत, हे विधान फक्त या
जगात नव्हे, तर कोणत्याही “संभाव्य” जगात सत्य आहे; किंवा ते या विशिष्ट जगात
केवळ सत्य असण्याऐवजी \emph{अनिवार्यपणे} सत्य आहे, असा त्याचा अर्थ घेता येतो.

अशा प्रकारच्या अनिवार्यतेचा अर्थ लावण्यासाठी मोडल तर्कशास्त्र रचले गेले.
नेहमीच्या विधानीय तर्कशास्त्रात चौकट-कारक जोडल्यावर मोडल विधानीय
तर्कशास्त्र मिळते; म्हणजे $\varphi$ हे सूत्र असेल, तर $\Box \varphi$ हेही
सूत्र असते. सहज अर्थाने, $\Box \varphi$ हे $\varphi$ \emph{अनिवार्यपणे} सत्य
आहे, किंवा कोणत्याही संभाव्य जगात सत्य आहे, असे सांगते. $\Diamond \varphi$ हे
सहसा $\lnot \Box \lnot \varphi$ चे संक्षिप्त रूप मानले जाते आणि $\varphi$
\emph{संभाव्यतः} सत्य आहे असे सांगणारे म्हणून वाचता येते. अर्थात, विधेय
तर्कशास्त्रातही मोडलता जोडता येते.

मोडल तर्कशास्त्राची चिन्हार्थमीमांसा देण्यासाठी क्रिप्के संरचना वापरता
येतात; खरे तर क्रिप्के यांनी ही चिन्हार्थमीमांसा प्रथम मोडल तर्कशास्त्र
डोळ्यांसमोर ठेवूनच रचली. अंशतः क्रमांपुरते मर्यादित न राहता, अधिक सामान्यपणे
“संभाव्य जगांचा” $P$ हा संच आणि जगांमधील $\Atom{R}{x,y}$ हा द्विपदी
“प्राप्यता” संबंध घेतला जातो. सहज अर्थाने, $\Atom{R}{p,q}$ हे जग~$q$ हे
जग~$p$ शी सुसंगत आहे असे सांगते; म्हणजे आपण जग~$p$ “मध्ये” असू, तर जग
$q$ सारखे असू शकले असते ही शक्यता विचारात घ्यावी लागते.

मोडल तर्कशास्त्राला कधी कधी “विस्तारात्मक” तर्कशास्त्राच्या विरुद्ध
“आशयात्मक” तर्कशास्त्र म्हणतात. अभिजात तर्कशास्त्रासारख्या विस्तारात्मक
तर्कशास्त्राची अभिप्रेत चिन्हार्थमीमांसा केवळ “प्रत्यक्ष” अशा एका जगाचा
निर्देश करते; पण “आशयात्मक” तर्कशास्त्राची चिन्हार्थमीमांसा अधिक विस्तृत
सत्ताशास्त्रावर अवलंबून असते. अनिवार्यतेची संरचना रचण्याखेरीज,
उपयोजनानुसार $\Box$ आणि $\Diamond$ यांचा नवा अर्थ लावून इतर भाषिक रचनांची
संरचना रचण्यासाठीही मोडलता वापरता येते. उदाहरणार्थ:
\begin{enumerate}
\item सिद्धतायोग्यता तर्कशास्त्रात $\Box \varphi$ चा अर्थ “$\varphi$ सिद्धतायोग्य
  आहे” आणि $\Diamond \varphi$ चा अर्थ “$\varphi$ सुसंगत आहे” असा घेतला जातो.
\item ज्ञानविषयक तर्कशास्त्रात $\Box \varphi$ चा अर्थ “मला $\varphi$ माहीत आहे” किंवा
  “माझा $\varphi$ वर विश्वास आहे” असा घेता येतो.
\item कालिक तर्कशास्त्रात $\Box \varphi$ चा अर्थ “$\varphi$ नेहमी सत्य आहे” आणि
  $\Diamond \varphi$ चा अर्थ “$\varphi$ कधीतरी सत्य असते” असा घेता येतो.
\end{enumerate}

मोडलता हाताळणारे नियम आणि स्वयंसिद्धके तर्कशास्त्रात जोडावीशी वाटतात.
उदाहरणार्थ, $\Log{S4}$ प्रणालीत विधानीय तर्कशास्त्राची नेहमीची स्वयंसिद्धके
आणि नियम, तसेच पुढील स्वयंसिद्धके असतात:
\begin{align*}
& \Box (\varphi \lif \psi) \lif (\Box \varphi \lif \Box \psi)\\
& \Box \varphi \lif \varphi \\
& \Box \varphi \lif \Box \Box \varphi
\intertext{आणि “$\varphi$ पासून $\Box \varphi$ हा निष्कर्ष काढा” हा नियमही असतो.
  $\Log{S5}$ मध्ये पुढील स्वयंसिद्धक जोडले जाते:}
& \Diamond \varphi \lif \Box \Diamond \varphi
\end{align*}
या स्वयंसिद्धकांच्या विविध आवृत्त्या वेगवेगळ्या उपयोजनांसाठी योग्य असू शकतात;
उदाहरणार्थ, S5 हे सहसा तार्किक अनिवार्यतेची संकल्पना लक्षणित करते असे मानतात.
चांगली गोष्ट अशी की क्रिप्के संरचना मधील प्राप्यता संबंध स्वाभाविक
रीतीने निर्बंधित करून, संबंधित निष्पत्ती प्रणाली निर्दोष आणि संपूर्ण
असेल अशी चिन्हार्थमीमांसा सहसा सापडते. उदाहरणार्थ, $\Log{S4}$ अशा क्रिप्के
संरचनेच्या वर्गाशी संबंधित आहे ज्यांत प्राप्यता संबंध परावर्ती आणि
संक्रमक असतो. $\Log{S5}$ अशा क्रिप्के संरचनेच्या वर्गाशी संबंधित
आहे ज्यांत प्राप्यता संबंध {\em वैश्विक} असतो, म्हणजे प्रत्येक जग इतर प्रत्येक
जगापासून प्राप्य असते; म्हणून $\Box \varphi$ तेव्हा आणि केवळ तेव्हाच लागू होते,
जेव्हा $\varphi$ प्रत्येक जगात लागू होते.












\section{इतर तर्कशास्त्रे}\label{fol:byd:oth:sec}

आतापर्यंत तुमच्या लक्षात आले असेलच की नवे तर्कशास्त्र रचणे अवघड नाही. तुम्हीही
स्वतःचा विन्यास तयार करू शकता, निगमन-पद्धती बनवू शकता आणि तिला साजेशी
चिन्हार्थमीमांसा घडवू शकता. निष्पत्ती प्रणाली चिन्हार्थमीमांसेसाठी
संपूर्ण हवी असेल, तर थोडे चातुर्य वापरावे लागू शकते; तसेच तुमचे तर्कशास्त्र
खरोखर रंजक आहे हे जगाला पटवून देण्यासाठी काही प्रयत्न करावे लागू शकतात. पण
त्याबदल्यात त्याच्या गणितीय आणि संगणकीय गुणधर्मांचा शोध घेताना अनेक तास
निर्भेळ आनंद मिळवता येईल.

अलीकडच्या दशकांत आकारिक तर्कशास्त्रांचा खरोखरच स्फोट झाला आहे. संदिग्ध
गुणधर्मांविषयीच्या तर्कविचाराचे प्रतिमान देण्यासाठी अस्पष्ट तर्कशास्त्र रचले
आहे. अनिश्चिततेविषयीच्या तर्कविचाराचे प्रतिमान देण्यासाठी संभाव्यताधारित
तर्कशास्त्र रचले आहे. खंडनीय प्रकारच्या तर्कविचाराचे प्रतिमान देण्यासाठी
पूर्वमान्य तर्कशास्त्रे आणि अ-एकस्वनिक तर्कशास्त्रे रचली आहेत; म्हणजे नवी
माहिती मिळाल्यावर नंतर उलटवता येतील असे “वाजवी” निष्कर्ष. ज्ञानाविषयीच्या
तर्कविचाराचे प्रतिमान देणारी ज्ञानविषयक तर्कशास्त्रे आहेत; कार्यकारण
संबंधांविषयीच्या तर्कविचाराचे प्रतिमान देणारी कार्यकारणविषयक तर्कशास्त्रे
आहेत; आणि नैतिक व नीतिविषयक कर्तव्यांविषयीच्या तर्कविचाराचे प्रतिमान देणारी
“कर्तव्यविषयक” तर्कशास्त्रेही आहेत. या प्रणाली मांडण्यामागील मुख्य प्रेरणा
तात्त्विक, गणितीय किंवा संगणकीय आहे यानुसार, अशा गोष्टींचा अभ्यास गणितीय
तर्कशास्त्र, तात्त्विक तर्कशास्त्र, कृत्रिम बुद्धिमत्ता, संज्ञानविज्ञान किंवा
इतर क्षेत्रांच्या नावाखाली झालेला आढळेल.

ही यादी सतत वाढत जाते आणि शक्यता अंतहीन दिसतात. सर्व मानवी तर्कबुद्धीचे गणनेत
रूपांतर करण्याचे लायप्निट्स यांचे स्वप्न आपण कदाचित कधीच साध्य करणार नाही---
पण त्यामुळे प्रयत्न करणे थांबवण्याचे कारण नाही.



\chapter{प्रतिमान उपपत्तीची मूलतत्त्वे}\label{mod:bas::chap}








\section{न्यूनीकृत आणि विस्तारित रचना}\label{mod:bas:red:sec}

ज्या भाषांत काही चिन्हे समान आहेत अशा भाषांची, तसेच त्या भाषांसाठीच्या
संरचनेची, तुलना करणे अनेकदा उपयुक्त किंवा आवश्यक असते. सर्वांत
नेहमीचे प्रकरण असे असते की भाषा~$\Lang{L}$ मधील सर्व चिन्हे
भाषा~$\Lang{L'}$ मध्येही असतात, म्हणजे $\Lang{L} \subseteq
\Lang{L'}$. मग अतिरिक्त चिन्हांची अर्थनिर्धारणे जोडून आणि समान चिन्हांची
अर्थनिर्धारणे तशीच ठेवून $\Lang{L}$-संरचना~$\Struct{M}$ चा विस्तार
नेहमीच $\Lang{L'}$-संरचनांमध्ये करता येतो. दुसरीकडे,
$\Lang{L}$ मध्ये नसलेल्या चिन्हांची अर्थनिर्धारणे केवळ ``विसरून''
$\Lang{L'}$-रचना~$\Struct{M'}$ पासून $\Lang{L}$-रचना मिळवता येते.

\begin{defn}
\label{mod:bas:red:defn:reduct}
समजा $\Lang L \subseteq \Lang L'$, $\Struct M$ ही
$\Lang L$-संरचना आणि $\Struct M'$ ही $\Lang L'$-संरचना आहे.
$\Struct M$ ही $\Struct M'$ ची $\Lang L$ पर्यंतची \emph{न्यूनीकृत रचना}
आणि $\Struct M'$ ही $\Struct M$ ची $\Lang L'$ पर्यंतची
\emph{विस्तारित रचना} असेल तेव्हा आणि केवळ तेव्हाच पुढील अटी पूर्ण होतात:
\begin{enumerate}
\item $\Domain{M} = \Domain{M'}$
\item प्रत्येक स्थिरांक~$c \in \Lang L$ साठी $\Assign{c}{M} =
  \Assign{c}{M'}$.
\item प्रत्येक फलन~$f \in \Lang L$ साठी $\Assign{f}{M} =
  \Assign{f}{M'}$.
\item प्रत्येक विधेय~$P \in \Lang L$ साठी $\Assign{P}{M} =
  \Assign{P}{M'}$.
\end{enumerate}
\end{defn}

\begin{prop}
\label{mod:bas:red:prop:reduct}
एखादी $\Lang{L}$-संरचना~$\Struct{M}$ ही एखाद्या
$\Lang{L'}$-संरचना~$\Struct{M'}$ ची न्यूनीकृत रचना असेल, तर सर्व
$\Lang{L}$-वाक्ये~$\varphi$ साठी,
\[\Sat{M}{\varphi} \text{ तेव्हा आणि केवळ तेव्हाच } \Sat{M'}{\varphi}.
\]
\end{prop}

\begin{proof}
  सराव.
\end{proof}

\begin{prob}
\ref{mod:bas:red:prop:reduct} सिद्ध करा.
\end{prob}

\begin{defn}
आपल्याकडे $\Lang{L}$-रचना $\Struct{M}$ असेल, एका $n$-स्थानी
विधेय~$P$ ची भर घालून मिळालेला $\Lang{L}$ चा विस्तार
$\Lang{L'} = \Lang{L} \cup \{P\}$ असेल, आणि $R \subseteq \Domain{M}^n$
हा $n$-स्थानी संबंध असेल, तर $\Assign{P}{M'} = R$ असलेली
$\Struct{M}$ ची विस्तारित रचना~$\Struct{M'}$ दर्शवण्यासाठी आपण
$\Expan{M}{R}$ असे लिहितो.
\end{defn}











\section{उपसंरचना}\label{mod:bas:sub:sec}

एखाद्या संरचना~$\Struct{M}$ चा प्रांत हा दुसऱ्या
$\Struct{M'}$ चा उपसंच असू शकतो. पण $\Struct{M}$ हा $\Struct{M'}$ चा
``भाग'' आहे असे आपण अर्थातच फक्त तेव्हाच मानावे, जेव्हा $\Domain{M}
\subseteq \Domain{M'}$ असण्याबरोबरच भाषेतील चिन्हांचे अर्थनिर्धारण करताना
$\Struct{M}$ आणि $\Struct{M'}$ किमान सामाईक भाग~$\Domain{M}$ वर
``सहमत'' असतात.

\begin{defn}
\label{mod:bas:sub:defn:substructure}
त्याच भाषेसाठी~$\Lang L$ संरचना $\Struct M$ आणि
$\Struct M'$ दिल्या असता, $\Struct M$ ही $\Struct M'$ ची
\emph{उपसंरचना} आणि $\Struct M'$ ही $\Struct M$ ची
\emph{वर्धित रचना} आहे, असे आपण म्हणतो; $\Struct M \substruct \Struct M'$
असे लिहितो; आणि तेव्हा आणि केवळ तेव्हाच पुढील अटी पूर्ण होतात:
\begin{enumerate}
\item $\Domain{M} \subseteq \Domain{M'}$,
\item प्रत्येक स्थिर $c \in \Lang L$ साठी $\Assign{c}{M} =
    \Assign{c}{M'}$;
\item प्रत्येक $n$-स्थानी फलन $f \in \Lang L$ साठी
  सर्व $a_1$, \dots, $a_n \in \Domain{M}$ यांकरिता
  $\Assign{f}{M}(a_1, \dots, a_n) = \Assign{f}{M'}(a_1, \dots, a_n)$.
\item प्रत्येक $n$-स्थानी विधेय $R \in \Lang L$ साठी, सर्व
  $a_1$, \dots, $a_n \in \Domain{M}$ यांकरिता $\langle
  a_1, \dots, a_n\rangle \in \Assign{R}{M}$ तेव्हा आणि केवळ तेव्हाच,
  जेव्हा $\langle a_1, \dots, a_n\rangle \in \Assign{R}{M'}$.
\end{enumerate}
\end{defn}

\begin{rem}
\label{mod:bas:sub:rem:substructure}
भाषेत एकही स्थिर किंवा फलने नसतील, तर कोणताही $N
\subseteq \Domain{M}$ हा $\Assign{R}{N} = \Assign{R}{M} \cap N^n$
असे ठेवून $\Struct M$ ची, प्रांत~$\Domain{N} = N$ असलेली,
उपसंरचना~$\Struct{N}$ निश्चित करतो.
\end{rem}













\section{सांत मर्यादेपलीकडील प्रसरण}\label{mod:bas:ove:sec}

\begin{thm}
\label{mod:bas:ove:overspill} वाक्यांचा एखादा संच $\Gamma$ याला हवे तितके मोठे
सांत प्रतिमान असतील, तर त्याला एक अनंत प्रतिमान असते.
\end{thm}

\begin{proof}
$c_0$, $c_1$, \dots ही गणनीय इतकी नवीन स्थिरे जोडून $\Gamma$ च्या भाषेचा
विस्तार करा आणि $\Gamma \cup \{c_i \neq c_j : i \neq j\}$ हा संच विचारात
घ्या. $\Gamma$ ला हवे तितके मोठे सांत प्रतिमान आहेत, याचा अर्थ प्रत्येक
$m >0$ साठी असा $n\ge m$ असतो की $\Gamma$ ला आकारमान~$n$ चे प्रतिमान
असते. यावरून $\Gamma \cup \{c_i \neq c_j : i \neq j\}$ सांततः
पूर्ततायोग्य आहे. संहततेनुसार $\Gamma \cup \{c_i \neq c_j : i \neq j\}$
याला प्रतिमान $\Struct M$ असते; त्याचा प्रांत अनंत असलाच पाहिजे, कारण ते
$c_i \neq c_j$ या सर्व असमानतांची पूर्ति करते.
\end{proof}

\begin{prop}
\label{mod:bas:ove:inf-not-fo}
कोणत्याही प्रथम-क्रम भाषेत असे कोणतेही वाक्य $\varphi$ नाही की ते
संरचना~$\Struct M$ मध्ये तेव्हा आणि केवळ तेव्हाच सत्य असते, जेव्हा
त्या संरचनेचा प्रांत $\Domain{M}$ अनंत असतो.
\end{prop}

\begin{proof}
असे एखादे $\varphi$ असते, तर त्याचे नकरण $\lnot \varphi$ सर्व आणि फक्त सांत
संरचनांमध्ये सत्य असते. म्हणून त्याला हवे तितके मोठे सांत प्रतिमान
असतात, पण अनंत प्रतिमान नसते; हे \ref{mod:bas:ove:overspill} शी विसंगत आहे.
\end{proof}











\section{समरूपी रचना}\label{mod:bas:iso:sec}

प्रथम-क्रमाच्या संरचना दोन प्रकारे एकसारख्या असू शकतात. त्यांच्या
सारखेपणाचा एक प्रकार असा की त्या तीच वाक्ये सत्य ठरवतात. अशा
संरचना ना आपण \emph{प्राथमिक सममूल्य} म्हणतो. परंतु रचना अतिशय
भिन्न असूनही तीच वाक्ये सत्य ठरवू शकतात---उदाहरणार्थ, एक रचना
गणनीय असू शकते आणि दुसरी नसू शकते. याचे कारण असे की
संरचनेची पुष्कळ वैशिष्ट्ये प्रथम-क्रम भाषांमध्ये व्यक्त करता येत
नाहीत: कधी भाषा पुरेशी समृद्ध नसते, तर कधी लोव्हेनहाइम--स्कोलेम
प्रमेयासारख्या प्रथम-क्रम तर्कशास्त्राच्या मूलभूत मर्यादा आड येतात. म्हणून
संरचनेच्या सारखेपणाचा दुसरा, अधिक काटेकोर प्रकार असा की त्या
मूलतः एकच असतात: त्यांना घडवणाऱ्या वस्तू वेगळ्या असतात, पण त्यांची
रचनात्मक वैशिष्ट्ये वेगळी नसतात. \emph{समरूपते} च्या संकल्पनेने हे नेमके
मांडता येते.

\begin{defn}
\label{mod:bas:iso:defn:elem-equiv}
एकाच भाषा~$\Lang{L}$ साठीच्या दोन संरचना $\Struct{M}$ आणि
$\Struct M'$ दिलेल्या असताना, $\Lang{L}$ मधील प्रत्येक वाक्य~$\varphi$
साठी $\Sat{M}{\varphi}$ तेव्हा आणि केवळ तेव्हाच जेव्हा $\Sat{M'}{\varphi}$ असेल,
तर $\Struct{M}$ ही $\Struct M'$ शी \emph{प्राथमिक सममूल्य} आहे असे म्हणतो
आणि $\Struct{M} \equiv \Struct M'$ असे लिहितो.
\end{defn}

\begin{defn}
\label{mod:bas:iso:defn:isomorphism} एकाच भाषा~$\Lang L$ साठीच्या दोन
संरचना $\Struct{M}$ आणि $\Struct M'$ दिलेल्या असताना, पुढील अटी
पूर्ण करणारे $h \colon \Domain{M} \to \Domain{M'}$ फलन असेल, तर
$\Struct{M}$ ही~$\Struct M'$ शी \emph{समरूपी} आहे असे म्हणतो आणि
$\Struct{M} \simeq \Struct M'$ असे लिहितो:
\begin{enumerate}
\item $h$ हे एकास-एक आहे: $h(x) =
  h(y)$ असेल, तर $x = y$;
\item $h$ हे आच्छादक आहे: प्रत्येक $y \in \Domain{M'}$ साठी
  $h(x) = y$ असा $x \in \Domain{M}$ असतो;
\item \label{mod:bas:iso:defn:iso-const}प्रत्येक स्थिरांक $c$ साठी:
  $h(\Assign{c}{M}) = \Assign{c}{M'}$;
\item \label{mod:bas:iso:defn:iso-pred}प्रत्येक $n$-स्थानी विधेय~$P$ साठी:
  \[  \tuple{a_1, \dots, a_n}\in \Assign{P}{M} \quad\text{तेव्हा आणि केवळ तेव्हाच जेव्हा}\quad
  \tuple{h(a_1), \dots, h(a_n)} \in \Assign{P}{M'};
  \]
\item \label{mod:bas:iso:defn:iso-func}प्रत्येक $n$-स्थानी फलन $f$ साठी:
  \[  h(\Assign{f}{M}(a_1, \dots, a_n)) =
  \Assign{f}{M'}(h(a_1), \dots, h(a_n)).
  \]
\end{enumerate}
\end{defn}

\begin{thm}
\label{mod:bas:iso:thm:isom}
$\Struct{M} \iso \Struct M'$ असेल, तर $\Struct{M} \elemequiv
\Struct{M'}$.
\end{thm}

\begin{proof}
$h$ हे $\Struct{M}$ पासून $\Struct M'$ वरचे समरूपण असू द्या. कोणत्याही
मूल्यांकन~$s$ साठी $h \circ s$ हे $h$ आणि $s$ यांचे संयोजन असते, म्हणजे
ते $\Struct{M'}$ मधील असे मूल्यांकन आहे की $(h \circ s)(x) = h(s(x))$.
$t$ आणि $\varphi$ यांवरील विगमनाने पुढील अधिक सबळ दावे सिद्ध करता येतात:
\begin{enumerate}
  \item[a.] $h(\Value{t}{M}[s]) = \Value{t}{M'}[h\circ s]$.
  \item[b.] $\Sat{M}{\varphi}[s]$ तेव्हा आणि केवळ तेव्हाच जेव्हा $\Sat{M'}{\varphi}[h \circ s]$.
\end{enumerate}
पहिला दावा~$t$ च्या गुंतागुंतीवरील विगमनाने सिद्ध होतो.
\begin{enumerate}
\item $t \ident c$ असेल, तर $\Value{c}{M}[s] = \Assign{c}{M}$ आणि
  $\Value{c}{M'}[h \circ s] = \Assign{c}{M'}$. म्हणून
  $h(\Value{t}{M}[s]) = h(\Assign{c}{M}) = \Assign{c}{M'}$
  (\ref{mod:bas:iso:defn:isomorphism} मधील \ref{mod:bas:iso:defn:iso-const} नुसार) $=
  \Value{t}{M'}[h \circ s]$.
\item $t \ident x$ असेल, तर $\Value{x}{M}[s] = s(x)$ आणि
  $\Value{x}{M'}[h \circ s] = h(s(x))$. म्हणून $h(\Value{x}{M}[s]) =
  h(s(x)) = \Value{x}{M'}[h \circ s]$.
\item $t \ident f(t_1, \dots, t_n)$ असेल, तर
  \begin{align*}
    \Value{t}{M}[s] & = \Assign{f}{M}(\Value{t_1}{M}[s], \dots, \Value{t_n}{M}[s]) \quad\text{आणि}\\
  \Value{t}{M'}[h \circ s] & = \Assign{f}{M'}(\Value{t_1}{M'}[h \circ
    s], \dots, \Value{t_n}{M'}[h \circ s]).
  \end{align*}

  प्रत्येक $i$ साठी $h(\Value{t_i}{M}[s])
  = \Value{t_i}{M'}[h\circ s]$ हे विगमन गृहीतक आहे. म्हणून
  \begin{align}
    h(\Value{t}{M}[s])
    & = h(\Assign{f}{M}(\Value{t_1}{M}[s], \dots, \Value{t_n}{M}[s]) \notag\\
    & = \Assign{f}{M'}(h(\Value{t_1}{M}[s]), \dots,
    h(\Value{t_n}{M}[s])) \label{mod:bas:iso:iso-1}\\
    & = \Assign{f}{M'}(\Value{t_1}{M'}[h \circ s], \dots,
    \Value{t_n}{M'}[h \circ s]) \label{mod:bas:iso:iso-2}\\
    & = \Value{t}{M'}[h\circ s] \notag
  \end{align}
  येथे \ref{mod:bas:iso:iso-1} हे \ref{mod:bas:iso:defn:isomorphism} मधील
  \ref{mod:bas:iso:defn:iso-func} वरून, तर \ref{mod:bas:iso:iso-2} हे विगमन गृहीतकावरून मिळते.
\end{enumerate}
(b) भागाचा पुरावा सरावासाठी सोडला आहे.

$\varphi$ हे वाक्य असेल, तर मूल्यांकने~$s$ आणि $h \circ s$ महत्त्वाची राहत
नाहीत; म्हणून $\Sat{M}{\varphi}$ तेव्हा आणि केवळ तेव्हाच जेव्हा
$\Sat{M'}{\varphi}$.
\end{proof}

\begin{prob}
\ref{mod:bas:iso:thm:isom} मधील (b) चा पुरावा सविस्तर पूर्ण करा.
\ref{mod:bas:iso:defn:isomorphism} मध्ये समरूपणाचे वैशिष्ट्य ठरवणाऱ्या
पाचही गुणधर्मांपैकी प्रत्येक गुणधर्म कुठे वापरला आहे, ते नोंदवण्याची
खात्री करा.
\end{prob}

\begin{defn}
रचना $\Struct{M}$ ची \emph{स्वयंरूपता} म्हणजे $\Struct{M}$ पासून
त्याच रचनेवर असलेले समरूपण होय.
\end{defn}

\begin{prob}
कोणतीही रचना $\Struct{M}$ असू द्या. जर $X$ हा $\Struct{M}$ चा परिभाष्य
उपसंच असेल आणि $h$ ही $\Struct{M}$ ची स्वयंरूपता असेल, तर $X
= \Setabs{h(x)}{x \in X}$ हे दाखवा (म्हणजे $h$ अंतर्गत $X$ अपरिवर्तित
राहतो).
\end{prob}













\section{संरचनेची उपपत्ती}\label{mod:bas:thm:sec}

प्रत्येक संरचना~$\Struct{M}$ काही वाक्ये सत्य ठरवते आणि
काही असत्य ठरवते. ती सत्य ठरवत असलेल्या सर्व वाक्यांच्या संचाला
तिची \emph{उपपत्ती} म्हणतात. तो संच खरोखरच एक उपपत्ती असतो, कारण
त्याच्यापासून तार्किकरीत्या निष्पन्न होणारे काहीही त्याच्या सर्व
प्रतिमानांमध्ये, आणि त्यांत~$\Struct{M}$ चाही समावेश होतो, सत्य असले
पाहिजे.

\begin{defn}
  संरचना~$\Struct M$ दिलेली असताना, $\Struct{M}$ ची
  \emph{उपपत्ती} म्हणजे $\Struct{M}$ मध्ये सत्य असलेल्या वाक्ये
  चा संच $\Theory{M}$ होय; म्हणजे, $\Theory{M} =
  \Setabs{\varphi}{\Sat{M}{\varphi}}$.
\end{defn}

अभिप्रेत अर्थनिर्धारण असलेल्या वाक्यांच्या संचांसाठीही आपण
``उपपत्ती'' हा शब्द अनौपचारिकपणे वापरतो, मग ते संच निगामीदृष्ट्या बंद
असोत किंवा नसोत.

\begin{prop}
कोणत्याही $\Struct{M}$ साठी $\Theory{M}$ संपूर्ण असते.
\end{prop}

\begin{proof}
कोणत्याही वाक्य~$\varphi$ साठी एकतर $\Sat{M}{\varphi}$ किंवा
$\Sat{M}{\lnot \varphi}$; म्हणून एकतर $\varphi \in \Theory{M}$ किंवा
$\lnot \varphi \in \Theory{M}$.
\end{proof}

\begin{prop}\label{mod:bas:thm:prop:equiv}
  प्रत्येक $\varphi \in \Theory{M}$ साठी $\Struct{N} \models \varphi$ असेल,
  तर $\Struct{M} \elemequiv \Struct{N}$.
\end{prop}

\begin{proof}
सर्व $\varphi \in \Theory{M}$ साठी $\Sat{N}{\varphi}$ असल्यामुळे
$\Theory{M} \subseteq \Theory{N}$. जर $\Sat{N}{\varphi}$, तर
$\Sat/{N}{\lnot \varphi}$; म्हणून $\lnot \varphi \notin \Theory{M}$.
$\Theory{M}$ संपूर्ण असल्यामुळे $\varphi \in \Theory{M}$. म्हणून
$\Theory{N} \subseteq \Theory{M}$ आणि आपल्याला
$\Struct{M} \elemequiv \Struct{N}$ मिळते.
\end{proof}

\begin{rem}\label{mod:bas:thm:remark:R}
  $\Struct{R} = \langle\Real, <\rangle$ विचारात घ्या. या
  संरचनेचा प्रांत म्हणजे वास्तव संख्यांचा संच~$\Real$ असून,
  ती फक्त एक 2-स्थानी विधेय असलेल्या भाषा साठीची रचना
  आहे; त्या विधेयाचे अर्थनिर्धारण वास्तव संख्यांवरील $<$ संबंध असे केले
  आहे. स्पष्टपणे $\Struct{R}$ ही अगणनीय आहे; पण
  $\Theory{R}$ उघडपणे सुसंगत असल्यामुळे लोव्हेनहाइम--स्कोलेम प्रमेयाने
  तिला एक गणनीय प्रतिमान असते, त्याला $\Struct{S}$ म्हणू.
  \ref{mod:bas:thm:prop:equiv} नुसार $\Struct{R} \equiv \Struct{S}$. शिवाय,
  $\Struct{R}$ आणि $\Struct{S}$ या समरूपी नसल्यामुळे यावरून
  \ref{mod:bas:iso:thm:isom} चा उलट दावा सर्वसाधारणपणे खोटा ठरतो.
\end{rem}











\section{आंशिक समरूपणे}\label{mod:bas:pis:sec}

\begin{defn}
  दोन संरचना $\Struct{M}$ आणि $\Struct{N}$ दिलेल्या असताना,
  $\Struct{M}$ पासून $\Struct{N}$ कडेचे \emph{आंशिक समरूपण} म्हणजे
  $\Domain M$ मधील घटक आदान म्हणून घेऊन $\Domain N$ मध्ये मूल्ये देणारे
  असे सांत अंशतः फलन $p$, जे आपल्या प्रांतावर
  \ref{mod:bas:iso:defn:isomorphism} मधील समरूपणाच्या अटी पूर्ण करते:
  \begin{enumerate}
  \item $p$ हे एकास-एक आहे;
  \item प्रत्येक स्थिरांक~$c$ साठी: $p(\Assign{c}{M})$ परिभाषित
    असेल, तर $p(\Assign{c}{M}) = \Assign{c}{N}$;
  \item प्रत्येक $n$-स्थानी विधेय $P$ साठी: $a_1$, \dots, $a_n$
    हे $p$ च्या प्रांतात असतील, तर $\langle a_1, \dots, a_n\rangle \in
    \Assign P M$ तेव्हा आणि केवळ तेव्हाच जेव्हा $\langle p(a_1), \dots, p(a_n) \rangle
    \in \Assign P N$;
  \item प्रत्येक $n$-स्थानी फलन $f$ साठी: $a_1$, \dots, $a_n$
    हे $p$ च्या प्रांतात असतील, तर
    \begin{multline*}
    p(\Assign f M (a_1, \dots,a_n))\\
    = \Assign f N (p(a_1), \dots, p(a_n)).
    \end{multline*}
  \end{enumerate}
  $p$ सांत असण्याचा अर्थ $\dom{p}$ सांत असतो.
\end{defn}

लक्षात घ्या की रिक्त फलन~$\emptyset$ हे कोणत्याही दोन संरचना
मधील आंशिक समरूपण नेहमीच असते.

\begin{defn}\label{mod:bas:pis:defn:partialisom}
  दोन संरचना $\Struct{M}$ आणि $\Struct{N}$ यांना
  \emph{आंशिकरीत्या समरूपी} म्हणतात आणि $\Struct{M} \iso[p]
  \Struct{N}$ असे लिहितात, तेव्हा आणि केवळ तेव्हाच जेव्हा
  $\Struct{M}$ आणि $\Struct{N}$ यांमधील आंशिक समरूपणांचा अरिक्त संच $I$
  पुढील \emph{पुढे-मागे} गुणधर्म पूर्ण करतो:
  \begin{enumerate}
  \item (\emph{पुढे}) प्रत्येक $p \in I$ आणि $a \in \Domain M$ साठी असा
    $q \in I$ असतो की $p \subseteq q$ आणि $a$ हा $q$ च्या प्रांतात
    असतो;
  \item (\emph{मागे}) प्रत्येक $p \in I$ आणि $b \in \Domain N$ साठी असा
    $q \in I$ असतो की $p \subseteq q$ आणि $b$ हा $q$ च्या व्याप्तीत
    असतो.
  \end{enumerate}
\end{defn}

\begin{thm}\label{mod:bas:pis:thm:p-isom1}
  $\Struct{M} \iso[p] \Struct{N}$ असून $\Struct{M}$ आणि
  $\Struct{N}$ या गणनीय असतील, तर $\Struct{M} \iso
  \Struct{N}$.
\end{thm}

\begin{proof}
  $\Struct{M}$ आणि $\Struct{N}$ या गणनीय असल्यामुळे
  $\Domain{M} = \{a_0, a_1, \ldots \}$ आणि $\Domain{N} = \{b_0, b_1,
  \ldots \}$ असे मानू. मनमाना $p_0 \in I$ घेऊन सुरुवात करून, आंशिक
  समरूपणांची वाढती क्रमिका $p_0 \subseteq p_1 \subseteq p_2
  \subseteq \cdots$ पुढीलप्रमाणे परिभाषित करतो:
  \begin{enumerate}
  \item $n+1$ विषम असेल, म्हणजे $n = 2r$ असेल, तर पुढे गुणधर्म वापरून
    असा $p_{n+1} \in I$ शोधा की $p_n \subseteq p_{n+1}$
    आणि $a_r$ हा $p_{n+1}$ च्या प्रांतात असेल;
  \item $n+1$ सम असेल, म्हणजे $n+1 =2r$ असेल, तर मागे गुणधर्म वापरून
    असा $p_{n+1} \in I$ शोधा की $p_n \subseteq p_{n+1}$
    आणि $b_r$ हा $p_{n+1}$ च्या व्याप्तीत असेल.
  \end{enumerate}
आता
\[p = \bigcup_{n\ge 0} p_n,
\]
असे ठेवले, तर $p$ हे $\Struct{M}$ आणि $\Struct{N}$ यांमधील समरूपण
असते.
\end{proof}

\begin{prob}
  \ref{mod:bas:pis:thm:p-isom1} मध्ये परिभाषित केलेले $p$ खरोखरच
  समरूपण आहे, हे सविस्तर दाखवा.
\end{prob}

\begin{thm}\label{mod:bas:pis:thm:p-isom2}
  $\Struct{M}$ आणि $\Struct{N}$ या निव्वळ संबंधात्मक भाषा
  साठीच्या संरचना आहेत असे समजा (म्हणजे फक्त विधेये
  असलेली आणि फलने किंवा स्थिरे नसलेली भाषा). मग
  $\Struct{M} \iso[p] \Struct{N}$ असेल, तर
  $\Struct{M} \elemequiv \Struct{N}$ सुद्धा असते.
\end{thm}

\begin{proof}
  सूत्रे वरील विगमनाने पुढील गोष्ट दाखवता येते. $a_1$, \dots,
  $a_n$ आणि $b_1$, \dots, $b_n$ यांसाठी प्रत्येक $a_i$ ला $b_i$ कडे
  नेणारे आंशिक समरूपण $p$ असेल, तसेच $s_1(x_i) =a_i$ आणि $s_2(x_i) =b_i$
  ($i =1$, \dots,~$n$ साठी) असेल, तर $\Sat{M}{\varphi}[s_1]$ तेव्हा आणि
  केवळ तेव्हाच जेव्हा $\Sat{N}{\varphi}[s_2]$. $n=0$ चे प्रकरण
  $\Struct{M} \elemequiv \Struct{N}$ देते.
\end{proof}

\begin{rem}
फलने उपस्थित असतील, तरी मागील निष्कर्ष सत्य असतो; पण $a_1$,
\dots, $a_n$ यांनी निर्मित $\Struct{M}$ च्या subसंरचना आणि
$b_1$, \dots, $b_n$ यांनी निर्मित $\Struct{N}$ च्या
subसंरचना यांमध्ये $p$ ने प्रेरित केलेले समरूपण विचारात घ्यावे
लागते.
\end{rem}

मागील निष्कर्षाचे टप्प्यांत ``विघटन'' करता येते: सूत्रामध्ये
एकमेकांत अंतर्भूत संख्यापकांची संख्या आणि संबंधित आंशिक समरूपणांचा
किती वेळा विस्तार करता येतो यांचा संबंध प्रस्थापित करावा लागतो.

\begin{defn}
  कोणतेही सूत्र~$\varphi$ दिले असता, $\varphi$ चा
  \emph{संख्यापक-प्रतांक}, $\QuantRank{\varphi} \in \Nat$, म्हणजे $\varphi$ मध्ये एकमेकांत अंतर्भूत
  संख्यापकांची कमाल संख्या; तो पुनरावर्ती रीतीने परिभाषित केला जातो.
  दोन संरचना $\Struct{M}$ आणि $\Struct{N}$ यांना
  \emph{$n$-सममूल्य} म्हणतात आणि $\Struct{M} \elemequiv[n] \Struct{N}$
  असे लिहितात, जर संख्यापक-प्रतांक जास्तीत जास्त~$n$ असलेल्या सर्व
  वाक्यांबाबत त्या एकमत असतील.
\end{defn}

\begin{prop}\label{mod:bas:pis:prop:qr-finite}
  $\Lang{L}$ ही सांत निव्वळ संबंधात्मक भाषा असू द्या, म्हणजे
  सांत इतकी विधेये आणि स्थिरांक असलेली, पण फलने
  नसलेली भाषा. मग प्रत्येक $n \in \Nat$ साठी, तार्किक
  सममूल्यतेपर्यंत, भाषा $\Lang{L}$ मध्ये संख्यापक-प्रतांक
  जास्तीत जास्त $n$ असलेली प्रथम-क्रम वाक्ये सांत इतकीच असतात.
\end{prop}

\begin{proof}
  $n$ वरील विगमनाने.
\end{proof}

\begin{defn}
  संरचना~$\Struct{M}$ दिलेली असताना, $\Domain M^{<\omega}$ हा
  $\Domain{M}$ वरील सर्व सांत क्रमिकांचा संच असू द्या. घटकांच्या सांत
  क्रमिकांवर $\mathbf{a}, \mathbf{b}, \mathbf{c}, \ldots$ ही चले वापरतो.
  जर $\mathbf{a} \in \Domain{M}^{<\omega}$ आणि $a \in \Domain{M}$,
  तर $\mathbf{a}a$ हे $\mathbf{a}$ आणि $a$ यांची \emph{जोडणी} दर्शवते.
\end{defn}

\begin{defn}
  संरचना $\Struct{M}$ आणि $\Struct{N}$ दिलेल्या असताना, समान
  लांबीच्या क्रमिकांमधील संबंध $I_n \subseteq \Domain M^{<\omega}
  \times \Domain N^{<\omega}$ हे $n$ वरील पुनरावर्तनाने पुढीलप्रमाणे
  परिभाषित करतो:
   \begin{enumerate}
   \item $I_0(\mathbf{a},\mathbf{b})$ तेव्हा आणि केवळ तेव्हाच जेव्हा
     $\mathbf{a}$ आणि $\mathbf{b}$ या अनुक्रमे $\Struct{M}$ आणि
     $\Struct{N}$ मध्ये तीच आण्विक सूत्रे पूर्ण करतात; म्हणजे,
     त्या क्रमिकांची समान लांबी k असेल, $s_1(x_i) = a_i$ आणि $s_2(x_i) =
     b_i$ असेल आणि $\varphi$ हे आण्विक असून त्यातील सर्व चले
     $x_1$, \dots,~$x_k$ यांपैकी असतील, तर $\Sat{M}{\varphi}[s_1]$ तेव्हा आणि
     केवळ तेव्हाच जेव्हा~$\Sat{N}{\varphi}[s_2]$.

   \item $I_{n+1} (\mathbf{a},\mathbf{b})$ तेव्हा आणि केवळ तेव्हाच
     जेव्हा प्रत्येक $a\in \Domain M$ साठी असा $b\in \Domain N$ असतो की
     $I_n (\mathbf{a}a,\mathbf{b}b)$, आणि उलटही.
   \end{enumerate}
\end{defn}


\begin{defn}
  $\Struct{M}$ आणि $\Struct{N}$ यांच्यासाठी
  $I_n(\emptyseq,\emptyseq)$ खरे असेल, तर $\Struct{M} \approx_n
  \Struct{N}$ असे लिहितो (येथे $\emptyseq$ ही रिक्त क्रमिका आहे).
\end{defn}

\begin{thm}\label{mod:bas:pis:thm:b-n-f}
  $\Lang{L}$ ही निव्वळ संबंधात्मक भाषा असू द्या. मग
  $I_n (\mathbf{a},\mathbf{b})$ असल्यास, $\QuantRank{\varphi} \le n$ असलेल्या
  प्रत्येक $\varphi$ साठी $\Sat{M}{\varphi}[\mathbf{a}]$ तेव्हा आणि केवळ तेव्हाच
  जेव्हा $\Sat{N}{\varphi}[\mathbf{b}]$ (येथे पुन्हा, $s(x_i) = a_i$ असणारे
  कोणतेही $s$ जर $\varphi$ पूर्ण करत असेल, तर $\mathbf{a}$ हे $\varphi$ पूर्ण
  करते). शिवाय, $\Lang{L}$ सांत असेल, तर उलट दावाही खरा असतो.
\end{thm}

\begin{proof}
  $I_n(\mathbf{a},\mathbf{b})$ यावरून $\mathbf{a}$ आणि $\mathbf{b}$
  संख्यापक-प्रतांक जास्तीत जास्त $n$ असलेली तीच सूत्रे पूर्ण करतात,
  हा पुरावा $\varphi$ वरील सोप्या विगमनाने मिळतो. उलट दाव्यासाठी आपण $n$
  वरील विगमन करतो आणि \ref{mod:bas:pis:prop:qr-finite} वापरतो; त्यावरून प्रत्येक
  $n$ साठी त्या संख्यापक-प्रतांकाची परस्पर तार्किकदृष्ट्या असममूल्य
  सूत्रे जास्तीत जास्त सांत इतकी असतात.

  $n=0$ साठी, $\mathbf{a}$ आणि $\mathbf{b}$ ही तीच संख्यापकमुक्त
  सूत्रे पूर्ण करतात या गृहीतकावरून ती तीच आण्विक सूत्रे पूर्ण करतात;
  म्हणून $I_0(\mathbf{a},\mathbf{b})$.

  $n+1$ च्या प्रकरणासाठी, $\mathbf{a}$ आणि $\mathbf{b}$ ही
  संख्यापक-प्रतांक जास्तीत जास्त $n+1$ असलेली तीच सूत्रे पूर्ण
  करतात असे समजा. $I_{n+1}(\mathbf{a},\mathbf{b})$ दाखवण्यासाठी प्रत्येक
  $a \in \Domain M$ साठी असा $b \in \Domain N$ असतो की
  $I_n(\mathbf{a}a,\mathbf{b}b)$, एवढे दाखवणे पुरेसे आहे; आणि विगमन
  गृहीतकाने प्रत्येक $a \in \Domain M$ साठी असा $b \in \Domain N$ असतो
  की $\mathbf{a}a$ आणि $\mathbf{b}b$ ही संख्यापक-प्रतांक जास्तीत जास्त
  $n$ असलेली तीच सूत्रे पूर्ण करतात, एवढे दाखवणे पुन्हा पुरेसे आहे.

  $a \in \Domain M$ दिलेला असताना, $\Struct{M}$ मध्ये $\mathbf{a}a$ ने
  पूर्ण केलेल्या, प्रतांक जास्तीत जास्त $n$ असलेल्या
  $\psi(x,\mathbf{y})$ या सूत्रांचा संच $\mathsf{T}^a_n$ असू द्या.
  $\mathsf{T}^a_n$ सांत आहे, म्हणून त्यातील सूत्रांचा संयोग असलेले एकच प्रथम-क्रम
  सूत्र म्हणून त्याचा विचार करू शकतो. यावरून $\mathbf{a}$ हे
  $\lexists[x][\mathsf{T}^a_n(x,\mathbf{y})]$ पूर्ण करते; या सूत्राचा
  संख्यापक-प्रतांक जास्तीत जास्त $n+1$ आहे. गृहीतकानुसार $\mathbf{b}$ हे
  $\Struct{N}$ मध्ये तेच सूत्र पूर्ण करते; म्हणून असा $b \in
  \Domain N$ असतो की $\mathbf{b}b$ हे $\mathsf{T}^a_n$ पूर्ण करते. विशेषतः,
  $\mathbf{b}b$ ही $\mathbf{a}a$ प्रमाणेच संख्यापक-प्रतांक जास्तीत जास्त
  $n$ असलेली सर्व सूत्रे पूर्ण करते. त्याचप्रमाणे, प्रत्येक $b \in
  \Domain N$ साठी असा $a\in \Domain M$ असतो की $\mathbf{a}a$ आणि
  $\mathbf{b}b$ ही संख्यापक-प्रतांक जास्तीत जास्त $n$ असलेली तीच
  सूत्रे पूर्ण करतात, हे दाखवता येते; त्यामुळे पुरावा पूर्ण होतो.
\end{proof}

\begin{cor}\label{mod:bas:pis:cor:b-n-f}
  $\Struct{M}$ आणि $\Struct{N}$ या सांत भाषा मधील निव्वळ
  संबंधात्मक संरचना असतील, तर $\Struct{M} \approx_n\Struct{N}$
  तेव्हा आणि केवळ तेव्हाच जेव्हा $\Struct{M} \elemequiv[n] \Struct{N}$.
  विशेषतः, $\Struct{M} \elemequiv \Struct{N}$ तेव्हा आणि केवळ तेव्हाच
  जेव्हा प्रत्येक $n$ साठी $\Struct{M} \approx_n \Struct{N}$.
\end{cor}











\section{सघन रेषीय क्रम}\label{mod:bas:dlo:sec}

\begin{defn}
  \emph{अंतबिंदुविरहित सघन रेषीय क्रम} म्हणजे एकच 2-स्थानी
  विधेय~$<$ असलेल्या भाषा साठीची आणि पुढील वाक्ये पूर्ण
  करणारी संरचना~$\Struct{M}$ होय:
  \begin{enumerate}
  \item $\lforall[x][\lnot x < x]$;
  \item $\lforall[x][\lforall[y][\lforall[z][(x < y \lif (y < z \lif x
    <z ))]]]$;
  \item $\lforall[x][\lforall[y][(x< y \lor \eq[x][y] \lor y < x)]]$;
  \item $\lforall[x][\lexists[y][x < y]]$;
  \item $\lforall[x][\lexists[y][y < x]]$;
  \item $\lforall[x][\lforall[y][(x < y \lif \lexists[z][(x < z \land
        z < y))]]]$.
 \end{enumerate}
\end{defn}

\begin{thm}\label{mod:bas:dlo:thm:cantorQ}
  कोणतेही दोन गणनीय अंतबिंदुविरहित सघन रेषीय क्रम समरूपी
  असतात.
\end{thm}

\begin{proof}
  $\Struct{M_1}$ आणि $\Struct{M_2}$ हे गणनीय अंतबिंदुविरहित
  सघन रेषीय क्रम असू द्या, जेथे ${<_1} = \Assign{<}{M_1}$ आणि ${<_2} =
  \Assign{<}{M_2}$; तसेच $\PIso{I}$ हा त्यांच्यामधील सर्व आंशिक
  समरूपणांचा संच असू द्या. किमान $\emptyset \in \PIso{I}$ असल्यामुळे
  $\PIso{I}$ रिक्त नाही. $\PIso{I}$ हा पुढे-मागे गुणधर्म पूर्ण करतो हे
  आपण दाखवतो. मग $\Struct{M_1} \iso[p] \Struct{M_2}$; आणि
  \ref{mod:bas:pis:thm:p-isom1} नुसार प्रमेय सिद्ध होते.

  $\PIso{I}$ हा पुढे गुणधर्म पूर्ण करतो हे दाखवण्यासाठी $p \in
  \PIso{I}$ असू द्या, $i = 1$, \dots,~$n$ साठी $p(a_i) = b_i$ असू द्या
  आणि सर्वसाधारणतेला बाधा न आणता $a_1 <_1 a_2 <_1 \cdots <_1 a_n$
  असे समजा. $a \in \Domain{M_1}$ दिलेला असताना, तो आधीच त्या नकाशाच्या
  प्रांतात असेल, तर तोच नकाशा अपेक्षित विस्तार आहे. नकाशा रिक्त असेल,
  तर दुसऱ्या रचनेच्या अरिक्त प्रांतातील कोणताही घटक प्रतिमा म्हणून
  निवडून त्याची दिलेल्या घटकाबरोबरची जोडी त्या नकाशात घालता येते.
  उरलेल्या प्रकरणांत
  पुढीलप्रमाणे $b \in \Domain{M_2}$ शोधा:

  \begin{enumerate}
  \item जर $a <_1 a_1$, तर $b <_2 b_1$ असेल असा $b \in
    \Domain{M_2}$ घ्या;
  \item जर $a_n <_1 a$, तर $b_n <_2 b$ असेल असा $b \in \Domain{M_2}$ घ्या;
 \item एखाद्या $i$ साठी $a_i <_1 a <_1 a_{i+1}$ असेल, तर $b_i <_2 b
   <_2 b_{i+1}$ असेल असा $b \in \Domain{M_2}$ घ्या.
  \end{enumerate}
  $\Struct{M_2}$ हा अंतबिंदुविरहित सघन रेषीय क्रम असल्यामुळे अपेक्षित
  गुणधर्म असलेला $b$ नेहमी शोधता येतो. $q = p \cup \{ \langle a, b
  \rangle \}$ असे परिभाषित करा; मग $q \in \PIso{I}$ हा $p$ चा अपेक्षित
  विस्तार आहे. यावरून पुढे गुणधर्म प्रस्थापित होतो. मागे गुणधर्मही
  याचप्रमाणे सिद्ध होतो. म्हणून $\Struct{M_1} \iso[p]
  \Struct{M_2}$; \ref{mod:bas:pis:thm:p-isom1} नुसार $\Struct{M_1} \iso
  \Struct{M_2}$.
\end{proof}

\begin{prob}
  $\PIso{I}$ हा मागे गुणधर्म पूर्ण करतो, हे तपासून
  \ref{mod:bas:dlo:thm:cantorQ} चा पुरावा पूर्ण करा.
\end{prob}

\begin{rem}
  $\Struct{S}$ हा कोणताही गणनीय अंतबिंदुविरहित सघन रेषीय क्रम
  असू द्या. मग (\ref{mod:bas:dlo:thm:cantorQ} नुसार) $\Struct{S} \iso
  \Struct{Q}$, जेथे $\Struct{Q} = (\Rat, <)$ हा परिमेय संख्यांचा संच
  $\Rat$ प्रांत म्हणून असलेला गणनीय सघन रेषीय क्रम आहे. आता
  \ref{mod:bas:thm:remark:R} मधील संरचना~$\Struct{R} = (\Real, <)$
  पुन्हा विचारात घ्या. असा गणनीय संरचना~$\Struct{S}$
  आहे की $\Struct{R} \elemequiv \Struct{S}$, हे आपण पाहिले. पण
  $\Struct{S}$ हा गणनीय अंतबिंदुविरहित सघन रेषीय क्रम आहे;
  म्हणून तो संरचना~$\Struct{Q}$ शी समरूपी (आणि म्हणून प्राथमिक
  सममूल्य) आहे. प्राथमिक सममूल्यतेच्या संक्रमणकतेने $\Struct{R}
  \elemequiv \Struct{Q}$. (हाच पुढे-मागे युक्तिवाद वापरून
  $\Struct{R} \iso[p] \Struct{Q}$ प्रस्थापित करून आपण हे थेटही दाखवू
  शकलो असतो.)
\end{rem}


\chapter{अंकगणिताची प्रतिमाने}\label{mod:mar::chap}








\section{परिचय}\label{mod:mar:int:sec}

अंकगणिताचे \emph{मानक प्रतिमान} म्हणजे अशी संरचना~$\Struct{N}$ की
$\Domain{N} = \Nat$ आणि तिच्यात $\Obj{0}$, $\prime$, $+$, $\times$ व
$<$ यांचे अपेक्षित नेहमीच्या अर्थाने अर्थनिर्धारण केलेले असते. म्हणजे,
$\Obj{0}$ हे $0$ असते, $\prime$ हे उत्तरवर्ती फलन असते, $+$ चा अर्थ
बेरीज आणि $\times$ चा अर्थ~$\Nat$ मधील संख्यांचा गुणाकार असा असतो.
अधिक स्पष्टपणे,
\begin{align*}
  \Assign{\Obj{0}}{N} & = 0\\
  \Assign{\prime}{N}(n) & = n + 1\\
  \Assign{+}{N}(n, m) & = n + m\\
  \Assign{\times}{N}(n, m) & = nm
\end{align*}
अर्थात, $\Lang{L_A}$ साठी अशाही रचना असतात की त्यांचे प्रांत~$\Nat$
नसतात. उदाहरणार्थ, $\Domain{M} = \{a\}^*$ (एकाच $a$ या चिन्हाच्या
सांत क्रमिका, म्हणजे $\emptyset$, $a$, $aa$, $aaa$, \dots) असा प्रांत
असलेली $\Struct{M}$ घेता येते आणि पुढील अर्थनिर्धारण करता येते:
\begin{align*}
  \Assign{\Obj{0}}{M} & = \emptyset\\
  \Assign{\prime}{M}(s) & = s \concat a\\
  \Assign{+}{M}(n, m) & = a^{n + m}\\
  \Assign{\times}{M}(n, m) & = a^{nm}
\end{align*}
या दोन रचना ``मूलतः एकच'' आहेत: फरक फक्त प्रांताचे
घटक कोणते आहेत यात आहे; अर्थनिर्धारण फलनांमुळे प्रांताचे
घटक परस्पर कसे संबंधित आहेत यात नाही. अशा दोन संरचना ना
\emph{समरूप} म्हणतात.

संहतता प्रमेयावरून सहज निष्पन्न होते की~$\Struct{N}$ मध्ये सत्य असलेल्या
कोणत्याही उपपत्तीला~$\Struct{N}$ शी समरूप नसलेली प्रतिमानेही असतात. अशा
रचनांना \emph{अमानक} म्हणतात. त्यांच्याबद्दलची रंजक गोष्ट अशी की मानक
प्रतिमानातील घटक (म्हणजे $\Struct{N}$ मधील, तसेच तिच्याशी समरूप
असलेल्या सर्व संरचना मधील घटक) मानक संख्यांक~$\num{n}$ यांच्या
मूल्यांनी पूर्णपणे व्यापलेले असतात, म्हणजे
\[\Domain{N} = \Setabs{\Value{\num{n}}{N}}{n \in \Nat}
\]
पण अमानक प्रतिमानांत तसे नसते: $\Struct{M}$ अमानक असेल, तर किमान एक
$x \in \Domain{M}$ असा असतो की प्रत्येक~$n$ साठी
$x \neq \Value{\num{n}}{M}$.

हे अमानक घटक फार रंजक आहेत: ते ``अनंत नैसर्गिक संख्या'' आहेत. पण त्यांच्या
अस्तित्वामुळे अपूर्णता-घटनांचे एका अर्थी स्पष्टीकरणही मिळते. उदाहरणार्थ,
पेआनो अंकगणितासाठीचे सुसंगतता-विधान $\OCon[\Th{PA}]$, म्हणजे $\lnot
\lexists[x][\OPrf[\Th{PA}](x, \gn{\lfalse})]$, विचारात घ्या.
$\Th{PA}$ कडून $\OCon[\Th{PA}]$ किंवा $\lnot \OCon[\Th{PA}]$ यांपैकी
काहीही सिद्ध होत नसल्यामुळे यांपैकी कोणतेही विधान~$\Th{PA}$ मध्ये सुसंगतपणे
जोडता येते. $\Th{PA}$ सुसंगत असल्यामुळे
$\Sat{N}{\OCon[\Th{PA}]}$, आणि म्हणून $\Sat/{N}{\lnot
  \OCon[\Th{PA}]}$. त्यामुळे $\Struct{N}$ हे $\Th{PA}
\cup \{\lnot \OCon[\Th{PA}]\}$ चे \emph{प्रतिमान नाही}, आणि त्या उपपत्तीची
सर्व प्रतिमाने अमानकच असली पाहिजेत. $\Th{PA} \cup \{\lnot
\OCon[\Th{PA}]\}$ च्या प्रतिमानात असा एखादा घटक असलाच पाहिजे की
तो $\lexists[x][\OPrf[\Th{PA}](x, \gn{\lfalse})]$ सत्य करणारा साक्षी
ठरेल, म्हणजे~$\Th{PA}$ पासून व्याघाताच्या निष्पत्तीची गोडेल संख्या.

असा घटक मानक असू शकत नाही---कारण प्रत्येक~$n$ साठी
$\Th{PA} \Proves \lnot \OPrf[\Th{PA}](\num{n}, \gn{\lfalse})$.











\section{अंकगणिताची मानक प्रतिमाने}\label{mod:mar:stm:sec}

अंकगणिताची भाषा~$\Lang{L_A}$ ही अर्थातच संख्यांविषयी, विशेषतः नैसर्गिक
संख्यांविषयी, अभिप्रेत आहे. म्हणून ``ते'' मानक प्रतिमान~$\Struct{N}$ विशेष
आहे: आपल्याला ज्याविषयी बोलायचे आहे तेच हे प्रतिमान. पण तर्कशास्त्रात
आपल्याला अनेकदा फक्त रचनात्मक गुणधर्मांत रस असतो, आणि समरूप असलेल्या कोणत्याही
दोन संरचनांमध्ये ते समान असतात. त्यामुळे व्याप्ती थोडी वाढवून
$\Struct{N}$ शी समरूप असलेली कोणतीही संरचना ``मानक'' मानता येते.

\begin{defn}
$\Lang{L_A}$ साठीची एखादी संरचना जर~$\Struct{N}$ शी समरूप असेल,
तर ती \emph{मानक} आहे.
\end{defn}

\begin{prop}
\label{mod:mar:stm:prop:standard-domain} एखादी संरचना~$\Struct{M}$ मानक असेल,
तर तिचा प्रांत म्हणजे मानक संख्यांकांच्या मूल्यांचा संच असतो, म्हणजे
\[\Domain{M} = \Setabs{\Value{\num{n}}{M}}{n \in \Nat}
\]
\end{prop}

\begin{proof}
प्रत्येक $\Value{\num{n}}{M} \in \Domain{M}$ हे स्पष्ट आहे. प्रत्येक
$x \in \Domain{M}$ हा कोणत्यातरी~$n$ साठी $\Value{\num{n}}{M}$ बरोबर आहे,
एवढेच दाखवायचे आहे. $\Struct{M}$ मानक असल्यामुळे ते~$\Struct{N}$ शी समरूप
आहे. $g\colon \Nat \to \Domain{M}$ हे समरूपण आहे असे समजा. मग
$g(n) = g(\Value{\num{n}}{N}) = \Value{\num{n}}{M}$. पण प्रत्येक
$x \in \Domain{M}$ साठी $g(n) = x$ असे एक~$n \in \Nat$ असते, कारण
$g$ हे आच्छादक आहे.
\end{proof}

\begin{explain}
$\Lang{L_A}$ साठीची एखादी रचना~$\Struct{M}$ मानक असेल, तर तिच्या
प्रांत मधील सर्व घटकांना $\num{0}$, $\num{1}$, $\num{2}$, \dots हे
मानक संख्यांक, म्हणजेच $\Obj{0}$, $\Obj{0}'$, $\Obj{0}''$ इत्यादी पदे,
नावे देऊ शकतात. अर्थात, याचा अर्थ $\Domain{M}$ मधील घटक
\emph{संख्याच आहेत} असा नाही; $\Domain{N}$ मधील संख्या जशा निवडून दाखवता
येतात, त्याच रीतीने हे घटक निवडून दाखवता येतात एवढाच त्याचा अर्थ आहे.
\end{explain}

\begin{prob}
\ref{mod:mar:stm:prop:standard-domain} चा व्यत्यास असत्य आहे हे दाखवा;
म्हणजे, $\Domain{M} = \Setabs{\Value{\num{n}}{M}}{n \in \Nat}$ असूनही
$\Struct{N}$ शी समरूप नसलेल्या एखाद्या संरचना~$\Struct{M}$ चे
उदाहरण द्या.
\end{prob}

\begin{prop}
\label{mod:mar:stm:prop:thq-standard}
जर $\Sat{M}{\Th{Q}}$ आणि $\Domain{M} =
\Setabs{\Value{\num{n}}{M}}{n \in \Nat}$, तर $\Struct{M}$ मानक आहे.
\end{prop}

\begin{proof}
$\Struct{M}$ हे~$\Struct{N}$ शी समरूप आहे हे दाखवायचे आहे.
$g(n) = \Value{\num{n}}{M}$ अशी व्याख्या असलेले
$g\colon \Nat \to \Domain{M}$ फलन विचारात घ्या. गृहीतकानुसार $g$ हे
आच्छादक आहे. ते एकास-एक सुद्धा आहे: $n \neq m$ असेल तेव्हा
$\Th{Q} \Proves \eq/[\num{n}][\num{m}]$. म्हणून
$\Sat{M}{\Th{Q}}$ असल्यामुळे, $n \neq m$ असेल तेव्हा
$\Sat{M}{\eq/[\num{n}][\num{m}]}$. त्यामुळे $n \neq m$ असल्यास
$\Value{\num{n}}{M} \neq \Value{\num{m}}{M}$, म्हणजे $g(n) \neq g(m)$.

$g$ हे समरूपण आहे हेही पडताळायचे आहे.
\begin{enumerate}
\item $\Assign{\Obj{0}}{N} = 0$ असल्यामुळे
  $g(\Assign{\Obj{0}}{N}) = g(0)$. $g$ च्या व्याख्येवरून
  $g(0) = \Value{\num{0}}{M}$. पण $\num{0}$ म्हणजे फक्त $\Obj{0}$;
  आणि एखादे पद स्थिरांक असेल, तर त्याचे मूल्य त्या संरचना ने
  त्या स्थिरांक ला नेमून दिलेले मूल्य असते, म्हणजे
  $\Value{\Obj{0}}{M} = \Assign{\Obj{0}}{M}$. म्हणून अपेक्षेप्रमाणे
  $g(\Assign{\Obj{0}}{N}) = \Assign{\Obj{0}}{M}$.
\item $\Struct{N}$ मध्ये $\prime$ हे~$\Nat$ वरील उत्तरवर्ती फलन असल्यामुळे
  $g(\Assign{\prime}{N}(n)) = g(n+1)$. पुढे, $g$ च्या व्याख्येवरून
  $g(n+1) = \Value{\num{n+1}}{M}$. पण $\num{n+1}$ आणि $\num{n}'$ हे
  एकच पद आहे, म्हणून $\Value{\num{n+1}}{M} = \Value{\num{n}'}{M}$.
  मूल्य-फलनाच्या व्याख्येवरून हे
  $= \Assign{\prime}{M}(\Value{\num{n}}{M})$. तसेच
  $\Value{\num{n}}{M} = g(n)$ असल्यामुळे
  $g(\Assign{\prime}{N}(n)) = \Assign{\prime}{M}(g(n))$ मिळते.
\item $\Struct{N}$ मध्ये $+$ हे~$\Nat$ वरील बेरीज-फलन असल्यामुळे
  $g(\Assign{+}{N}(n,m)) = g(n+m)$. पुढे, $g$ च्या व्याख्येवरून
  $g(n+m) = \Value{\num{n+m}}{M}$. पण $\Th{Q} \Proves
  \num{n+m} = (\num{n} + \num{m})$, म्हणून $\Value{\num{n+m}}{M} =
  \Value{\num{n}+\num{m}}{M}$. मूल्य-फलनाच्या व्याख्येवरून हे $=
  \Assign{+}{M}(\Value{\num{n}}{M},\Value{\num{m}}{M})$. तसेच
  $\Value{\num{n}}{M} = g(n)$ आणि $\Value{\num{m}}{M} = g(m)$ असल्यामुळे
  $g(\Assign{+}{N}(n, m)) = \Assign{+}{M}(g(n), g(m))$ मिळते.
\item $g(\Assign{\times}{N}(n, m)) = \Assign{\times}{M}(g(n), g(m))$:
  सराव.
\item $\tuple{n,m} \in \Assign{<}{N}$ हे तेव्हाच, जेव्हा $n < m$. जर
  $n < m$, तर $\Th{Q} \Proves \num{n} < \num{m}$ आणि
  $\Sat{M}{\num{n} < \num{m}}$ सुद्धा. त्यामुळे
  $\tuple{\Value{\num{n}}{M}, \Value{\num{m}}{M}} \in \Assign{<}{M}$,
  म्हणजे $\tuple{g(n), g(m)} \in \Assign{<}{M}$. जर $n \not< m$, तर
  $\Th{Q} \Proves \lnot \num{n} < \num{m}$ आणि परिणामी
  $\Sat/{M}{\num{n} < \num{m}}$. त्यामुळे पूर्वीप्रमाणे
  $\tuple{g(n), g(m)} \notin \Assign{<}{M}$. दोन्ही एकत्र केल्यास:
  $\tuple{n,m} \in \Assign{<}{N}$ हे तेव्हाच, जेव्हा
  $\tuple{g(n), g(m)} \in \Assign{<}{M}$.
\end{enumerate}
\end{proof}

\begin{explain}
$\Lang{L_A}$ साठीच्या इतर कोणत्याही संरचना~$\Struct{M}$ च्या
प्रांतात~$\Nat$ कडून संगती ठरवण्याचा सर्वांत साहजिक मार्ग म्हणजे फलन~$g$;
कारण अशा प्रत्येक $\Struct{M}$ मध्ये $\num{0}$, $\num{1}$, $\num{2}$
इत्यादींनी नाव दिलेले घटक असतात. त्यामुळे, $\Struct{M}$ ने
$\num{n}$ विषयीची किमान काही मूलभूत विधाने~$\Struct{N}$ प्रमाणेच सत्य
केली आणि $g$ हे एकास-एक व आच्छादकही असेल, तर $g$ हे समरूपण ठरेल यात नवल
नाही. किंबहुना, $\Domain{M}$ मध्ये $\num{n}$ यांनी नाव दिलेल्यांव्यतिरिक्त
दुसरे घटक नसतील, तर असे समरूपण एकमेव असते.
\end{explain}

\begin{prop}
\label{mod:mar:stm:prop:thq-unique-iso} जर $\Struct{M}$ मानक असेल, तर
\ref{mod:mar:stm:prop:thq-standard} च्या पुराव्यातील $g$ हे~$\Struct{N}$ कडून
$\Struct{M}$ कडे जाणारे एकमेव समरूपण आहे.
\end{prop}

\begin{proof}
$h\colon \Nat \to \Domain{M}$ हे $\Struct{N}$ आणि~$\Struct{M}$ यांमधील
समरूपण आहे असे समजा. $n$ वरील विगमनाने $g = h$ हे दाखवू. $n = 0$ असेल,
तर $g$ च्या व्याख्येवरून $g(0) = \Assign{\Obj{0}}{M}$. पण $h$ हे समरूपण
असल्यामुळे $h(0) = h(\Assign{\Obj{0}}{N}) =\Assign{\Obj{0}}{M}$;
म्हणून $g(0) = h(0)$.

आता $n+1$ चे प्रकरण विचारात घ्या. आपल्याला
\begin{align*}
  g(n+1) & = \Value{\num{n+1}}{M} \text{ हे~$g$ च्या व्याख्येवरून}\\
  & = \Value{\num{n}'}{M} \text{ कारण $\num{n+1}\ident \num{n}'$}\\
  & = \Assign{\prime}{M}(\Value{\num{n}}{M})
    \text{ हे $\Value{t'}{M}$ च्या व्याख्येवरून}\\
  & = \Assign{\prime}{M}(g(n))  \text{ हे~$g$ च्या व्याख्येवरून}\\
  & = \Assign{\prime}{M}(h(n)) \text{ हे विगमन गृहीतकावरून}\\
  & = h(\Assign{\prime}{N}(n)) \text{ कारण $h$ हे समरूपण आहे}\\
  & = h(n+1)
\end{align*}
मिळते.
\end{proof}

\begin{explain}
कोणत्याही गणनीय अनंत संच~$M$ साठी $\Nat$ आणि $M$ यांमध्ये
एकास-एक आच्छादन असते; म्हणून असा प्रत्येक संच~$M$ हा एखाद्या मानक
प्रतिमान~$\Struct{M}$ चा संभाव्य प्रांत आहे. किंबहुना, एकदा $z \in M$
ही वस्तू आणि योग्य फलन $s$ ही अनुक्रमे $\Assign{\Obj{0}}{M}$ आणि
$\Assign{\prime}{M}$ म्हणून निवडली, की $+$, $\times$ आणि $<$ यांचे
अर्थनिर्धारण आधीच निश्चित होते. $s\colon M \to M \setminus \{z\}$ अशी
एकास-एक आणि आच्छादक दोन्ही असलेली फलनेच मानक प्रतिमानात
$\Assign{\prime}{M}$ म्हणून योग्य असतात. $s$ च्या व्याप्तीत~$z$ असू शकत
नाही; अन्यथा $\lforall[x][\eq/[\Obj 0][x']]$ असत्य होईल. हे वाक्य
$\Struct{N}$ मध्ये सत्य आहे, म्हणून $\Struct{M}$ नेही ते सत्य केले पाहिजे.
फलन~$s$ हे एकास-एक असले पाहिजे, कारण $\Struct{N}$ मधील उत्तरवर्ती
फलन~$\Assign{\prime}{N}$ तसे आहे; आणि $\Assign{\prime}{N}$ हे
एकास-एक आहे हे~$\Struct{N}$ मध्ये सत्य असलेल्या वाक्य ने
व्यक्त होते. ते आच्छादक सुद्धा असले पाहिजे; अन्यथा असा एखादा
$x \in M \setminus \{z\}$ असेल की तो~$s$ च्या व्याप्तीत नसेल, म्हणजे
वाक्य $\lforall[x][(\eq[x][\Obj 0] \lor
\lexists[y][\eq[y'][x]])]$ हे~$\Struct{M}$ मध्ये असत्य होईल---पण ते
$\Struct{N}$ मध्ये सत्य आहे.

\end{explain}












\section{अमानक प्रतिमाने}\label{mod:mar:nst:sec}

\begin{explain}
$\Lang{L_A}$ साठीची एखादी संरचना~$\Struct{N}$ शी समरूप असेल,
तर तिला मानक म्हणतो. एखादी संरचना~$\Struct{N}$ शी समरूप नसेल,
तर तिला अमानक म्हणतो.
\end{explain}

\begin{defn}
$\Lang{L_A}$ साठीची संरचना~$\Struct{M}$ जर~$\Struct{N}$ शी समरूप
नसेल, तर ती \emph{अमानक} आहे. प्रांतातील जे घटक
$x \in \Domain{M}$ कोणत्यातरी $n \in \Nat$ साठी
$\Value{\num{n}}{M}$ बरोबर आहेत, त्यांना ($\Struct{M}$ मधील)
\emph{मानक संख्या} म्हणतात; आणि जे तसे नाहीत त्यांना \emph{अमानक संख्या}
म्हणतात.
\end{defn}

\begin{explain}
\ref{mod:mar:stm:prop:standard-domain} नुसार, $\Lang{L_A}$ साठीच्या कोणत्याही
मानक संरचनांमध्ये फक्त मानक घटक असतात. त्यामुळे अमानक
संरचनांमध्ये किमान एक अमानक घटक असलाच पाहिजे. किंबहुना, अमानक
घटक अस्तित्वात असणे हे ती संरचना अमानक असल्याची हमी देते.
\end{explain}

\begin{prop}
$\Lang{L_A}$ साठीच्या संरचना~$\Struct{M}$ मध्ये अमानक संख्या
असेल, तर $\Struct{M}$ अमानक आहे.
\end{prop}

\begin{proof}
तसे नाही असे समजा; म्हणजे $\Struct{M}$ मानक असूनही तिच्यात अमानक संख्या
$x$ आहे असे समजा. $g\colon \Nat \to \Domain{M}$ हे समरूपण असू द्या.
$n$ वरील विगमनाने $g(\Value{\num{n}}{N}) = \Value{\num{n}}{M}$ हे सहज
दिसते. दुसऱ्या शब्दांत, $g$ हे~$\Struct{N}$ मधील मानक संख्यांना
$\Struct{M}$ मधील मानक संख्यांवर पाठवते. $\Struct{M}$ मध्ये अमानक संख्या
असेल, तर गृहीतकाच्या विरुद्ध $g$ हे आच्छादक असू शकत नाही.
\end{proof}

\begin{prob}
$\Th{Q}$ मध्ये पुढील स्वयंसिद्धके आहेत हे आठवा:
\begin{align*}
& \lforall[x][\lforall[y][(\eq[x'][y'] \lif \eq[x][y])]] \tag{$\mathsf{Q}_1$}\\
& \lforall[x][\eq/[\Obj 0][x']] \tag{$\mathsf{Q}_2$}\\
& \lforall[x][(\eq[x][\Obj 0] \lor \lexists[y][\eq[x][y']])] \tag{$\mathsf{Q}_3$}
\end{align*}
पुढील अटी पूर्ण करणाऱ्या संरचना~$\Struct{M_1}$, $\Struct{M_2}$,
$\Struct{M_3}$ द्या:
\begin{enumerate}
\item $\Sat{M_1}{\mathsf{Q}_1}$, $\Sat{M_1}{\mathsf{Q}_2}$, $\Sat/{M_1}{\mathsf{Q}_3}$;
\item $\Sat{M_2}{\mathsf{Q}_1}$, $\Sat/{M_2}{\mathsf{Q}_2}$, $\Sat{M_2}{\mathsf{Q}_3}$; आणि
\item $\Sat/{M_3}{\mathsf{Q}_1}$, $\Sat{M_3}{\mathsf{Q}_2}$, $\Sat{M_3}{\mathsf{Q}_3}$.
\end{enumerate}
अर्थात, प्रत्येकासाठी फक्त~$\Assign{\Obj{0}}{M_i}$ आणि
$\Assign{\prime}{M_i}$ निश्चित करायची आहेत.
\end{prob}

\begin{explain}
$\Lang{L_A}$ साठी अमानक संरचना निश्चित करणे पुरेसे सोपे आहे.
उदाहरणार्थ, प्रांत~$\Int$ असलेली रचना घ्या आणि सर्व अतार्किक चिन्हांचे
नेहमीप्रमाणे अर्थनिर्धारण करा. ऋण संख्या कोणत्याही~$n$ साठी $\num{n}$ ची
मूल्ये नसल्यामुळे ही रचना अमानक आहे. अर्थात, $\Struct{N}$ प्रमाणेच ती
वाक्ये सत्य करते या अर्थाने ती अंकगणिताचे \emph{प्रतिमान} नसेल. उदाहरणार्थ,
$\lforall[x][\eq/[x'][\Obj{0}]]$ हे असत्य आहे. तरीसुद्धा, संहतता प्रमेय
वापरून अंकगणिताची अमानक प्रतिमाने अस्तित्वात आहेत हे सहज सिद्ध करता येते.
\end{explain}

\begin{prop}
$\Th{TA} = \Setabs{\varphi}{\Sat{N}{\varphi}}$ ही~$\Struct{N}$ ची उपपत्ती असू द्या.
$\Th{TA}$ ला गणनीय अमानक प्रतिमान आहे.
\end{prop}

\begin{proof}
$\Lang{L_A}$ मध्ये नवीन स्थिरांक~$c$ घालून तिचा विस्तार करा आणि
पुढील वाक्यांचा संच विचारात घ्या:
\[\Gamma = \Th{TA} \cup \{\eq/[c][\num{0}], \eq/[c][\num{1}],
\eq/[c][\num{2}], \dots\}
\]
$\Gamma$ च्या कोणत्याही प्रतिमान~$\Struct{M^c}$ मध्ये
$x = \Assign{c}{M^c}$ असा अमानक घटक असेल, कारण प्रत्येक
$n \in \Nat$ साठी $x \neq \Value{\num{n}}{M}$. तसेच
$\Th{TA} \subseteq \Gamma$ असल्यामुळे अर्थातच $\Sat{M^c}{\Th{TA}}$.
फक्त~$c$ विसरून $\Struct{M^c}$ पासून $\Lang{L_A}$ साठीची
संरचना~$\Struct{M}$ बनवली, तरी तिच्या प्रांतात तो अमानक~$x$ राहतो
आणि $\Sat{M}{\Th{TA}}$ सुद्धा. $\Th{TA}$ मध्ये $c$ येत नसल्यामुळे या
दुसऱ्या दाव्याची हमी मिळते. म्हणून $\Gamma$ ला प्रतिमान आहे एवढे दाखवणे
पुरेसे आहे.


$\Gamma$ ला प्रतिमान आहे हे दाखवण्यासाठी संहतता प्रमेय वापरू. $\Gamma$ चा
प्रत्येक सांत उपसंच पूर्ततायोग्य असेल, तर $\Gamma$ सुद्धा पूर्ततायोग्य आहे.
कोणताही सांत उपसंच $\Gamma_0 \subseteq \Gamma$ विचारात घ्या.
$\Gamma_0$ मध्ये~$\Th{TA}$ मधील काही वाक्ये आणि
$\eq/[c][\num{n}]$ या रूपातील सांत वाक्ये येतात. या दुसऱ्या रूपातील एकही
वाक्य नसेल, तर c ला शून्य मूल्य देणारा N चा विस्तार त्या उपसंचाला थेट
पूर्ण करतो; म्हणून पुढे किमान एक असे वाक्य आहे असे समजा.

$\eq/[c][\num{k}] \in \Gamma_0$ असा $k$ हा सर्वांत मोठा अंक समजा.
$\Assign{c}{N_k} = k+1$ हे अर्थनिर्धारण समाविष्ट करण्यासाठी~$\Struct{N}$
चा विस्तार करून $\Struct{N_k}$ ची व्याख्या करा. $\Sat{N_k}{\Gamma_0}$:
जर $\varphi \in \Th{TA}$, तर $\Struct{N_k}$ ही~$c$ वगळता प्रत्येक बाबतीत
$\Struct{N}$ सारखीच असल्यामुळे आणि~$\varphi$ मध्ये $c$ येत नसल्यामुळे
$\Sat{N_k}{\varphi}$. तसेच $n \le k$ आणि $\Value{c}{N_k} = k+1$ असल्यामुळे
$\Sat{N_k}{\eq/[c][\num{n}]}$. म्हणून $\Gamma$ चा प्रत्येक सांत उपसंच
पूर्ततायोग्य आहे आणि संहततेनुसार संपूर्ण संचाला प्रतिमान आहे. अधोगामी
लोव्हेनहाइम--स्कोलेम प्रमेयानुसार हे प्रतिमान गणनीय घेता येते.

\end{proof}











\section{$\Th{Q}$ ची प्रतिमाने}\label{mod:mar:mdq:sec}

\begin{explain}
$\Struct{N}$ जी वाक्ये सत्य करते तीच सत्य करणाऱ्या अमानक
संरचना आहेत, म्हणजे त्या~$\Th{TA}$ ची प्रतिमाने आहेत, हे आपल्याला
माहीत आहे. $\Sat{N}{\Th{Q}}$ असल्यामुळे $\Th{TA}$ चे कोणतेही प्रतिमान
$\Th{Q}$ चेही प्रतिमान असते. $\Th{Q}$ ही~$\Th{TA}$ पेक्षा फारच दुर्बल
आहे; उदाहरणार्थ, $\Th{Q} \Proves/ \lforall[x][\lforall[y][\eq[(x +
      y)][(y+x)]]]$. दुर्बलतर उपपत्ती पूर्ण करणे सोपे असते: त्यांना अधिक
प्रतिमाने असतात. उदाहरणार्थ, $\Th{Q}$ च्या काही प्रतिमानांत
$\lforall[x][\lforall[y][\eq[(x + y)][(y+x)]]]$ असत्य असते; पण ती
$\Th{TA}$ ची किंवा त्या दृष्टीने $\Th{PA}$ चीसुद्धा प्रतिमाने असू शकत
नाहीत. $\Th{Q}$ ची प्रतिमाने तुलनेने सोपीही असतात: ती स्पष्टपणे निश्चित
करता येतात.
\end{explain}

\begin{ex}
\label{mod:mar:mdq:ex:model-K-of-Q}
$\Domain{K} = \Nat \cup \{a\}$ असा प्रांत आणि पुढील अर्थनिर्धारणे असलेली
संरचना~$\Struct{K}$ विचारात घ्या:
\begin{align*}
  \Assign{\Obj{0}}{K} & = 0\\
  \Assign{\prime}{K}(x) & =
  \begin{cases}
    x+1 & \text{जर $x\in \Nat$}\\
    a & \text{जर $x = a$}
  \end{cases}\\
  \Assign{+}{K}(x, y) & =
  \begin{cases}
    x+y & \text{जर $x$, $y \in\Nat$}\\
    a & \text{इतर प्रसंगी}
  \end{cases}\\
  \Assign{\times}{K}(x, y) & =
  \begin{cases}
    xy & \text{जर $x$, $y \in\Nat$}\\
    0 & \text{जर $x = 0$ किंवा $y = 0$}\\
    a & \text{इतर प्रसंगी}\\
  \end{cases}\\
  \Assign{<}{K} & =
  \Setabs{\tuple{x,y}}{x, y \in \Nat \text{ आणि } x<y} \cup
  \Setabs{\tuple{x,a}}{x \in \Domain{K}}
\end{align*}
$\Sat{K}{\Th{Q}}$ हे दाखवण्यासाठी $\Th{Q}$ ची सर्व स्वयंसिद्धके
$\Struct{K}$ मध्ये सत्य आहेत हे पडताळावे लागेल. सोयीसाठी,
$\Assign{\prime}{K}(x)$ करिता $x^\nssucc$ ($\Struct{K}$ मधील $x$ चा
``उत्तरवर्ती''), $\Assign{+}{K}(x, y)$ करिता $x \nsplus y$
($\Struct{K}$ मधील $x$ आणि $y$ यांची ``बेरीज''),
$\Assign{\times}{K}(x, y)$ करिता $x \nstimes y$ ($\Struct{K}$ मधील
$x$ आणि~$y$ यांचा ``गुणाकार''), आणि $\tuple{x,y} \in \Assign{<}{K}$
करिता $x \nsless y$ लिहू. या संक्षेपांनी $\Struct{K}$ मधील क्रिया अधिक
स्पष्टपणे अशा देता येतात:
\[\begin{array}{c|c}
  x & x^\nssucc \\
  \hline
  n & n+1 \\
  a & a
\end{array}
\qquad
\begin{array}{c|ccc}
  x \nsplus y & 0 & m & a \\
  \hline
  0 & 0 & m & a \\
  n & n & n+m & a \\
  a & a & a & a \\
\end{array}
\qquad
\begin{array}{c|ccc}
  x \nstimes y & 0 & m & a \\
  \hline
  0 & 0 & 0 & 0 \\
  n & 0 & nm & a \\
  a & 0 & a & a \\
\end{array}
\]
$n$, $m \in \Nat$ यांच्यासाठी $n \nsless m$ हे तेव्हाच, जेव्हा $n<m$;
आणि प्रत्येक~$x \in \Domain{K}$ साठी $x \nsless a$.

$\nssucc$ हे एकास-एक असल्यामुळे
$\Sat{K}{\lforall[x][\lforall[y][(\eq[x'][y'] \lif \eq[x][y])]]}$.
$\Struct{K}$ मध्ये $0$ हा $\nssucc$-उत्तरवर्ती नसल्यामुळे
$\Sat{K}{\lforall[x][\eq/[\Obj 0][x']]}$. प्रत्येक $n>0$ साठी
$n = (n-1)^\nssucc$ आणि $a = a^\nssucc$ असल्यामुळे
$\Sat{K}{\lforall[x][(\eq[x][\Obj 0] \lor \lexists[y][\eq[x][y']])]}$.

$n \nsplus 0 = n+0 = n$ आणि~$\nsplus$ च्या व्याख्येवरून
$a\nsplus 0 = a$ असल्यामुळे $\Sat{K}{\lforall[x][\eq[(x + \Obj 0)][x]]}$.
$\Sat{K}{\lforall[x][\lforall[y][\eq[(x + y')][(x + y)']]]}$ हे थोडे
अधिक गुंतागुंतीचे आहे. $n$, $m$ या दोन्ही मानक असतील, तर:
\begin{align*}
(n \nsplus m^\nssucc) & = (n+(m+1)) = (n+m)+1 = (n \nsplus m)^\nssucc
\intertext{कारण मानक संख्यांवर $\nsplus$ आणि $^\nssucc$ ही अनुक्रमे $+$
  आणि $\prime$ यांच्याशी जुळतात. आता $x \in \Domain{K}$ असे समजा. मग}
(x \nsplus a^\nssucc) & = (x \nsplus a) = a = a^\nssucc = (x \nsplus a)^\nssucc
\intertext{उरलेले प्रकरण $y \in \Domain{K}$ पण $x =
  a$ असण्याचे आहे. येथे $y = n$ मानक आहे की $y = a$, यानुसारही प्रकरणे
  वेगळी करावी लागतात:}
(a \nsplus n^\nssucc) & = (a \nsplus (n+1)) = a = a^\nssucc = (a \nsplus n)^\nssucc\\
(a \nsplus a^\nssucc) & = (a \nsplus a) = a = a^\nssucc = (a \nsplus a)^\nssucc
\end{align*}

अर्थात, हा तपशील आवश्यकतेपेक्षा थोडा अधिक आहे. उदाहरणार्थ,
$z$ काहीही असले तरी $a \nsplus z = a$ असल्यामुळे
$a \nsplus a^\nssucc = a$ हे लगेच निष्पन्न होते. उरलेली स्वयंसिद्धकेही
याच रीतीने पडताळता येतात.

त्यामुळे $\Struct{K}$ हे~$\Th{Q}$ चे प्रतिमान आहे. तिची ``बेरीज''~$\nsplus$
ही क्रमनिरपेक्षसुद्धा आहे. पण $\Struct{N}$ मध्ये सत्य आणि $\Struct{K}$
मध्ये असत्य अशी इतर वाक्ये आहेत, तसेच याच्या उलटही आहे. उदाहरणार्थ,
$a \nsless a$, म्हणून $\Sat{K}{\lexists[x][x < x]}$ आणि
$\Sat/{K}{\lforall[x][\lnot x<x]}$. यावरून $\Th{Q} \Proves/
\lforall[x][\lnot x < x]$ हे दिसते.
\end{ex}

\begin{prob}
\ref{mod:mar:mdq:ex:model-K-of-Q} मधील $\Struct{K}$ ही~$\Th{Q}$ ची
उरलेली स्वयंसिद्धके पूर्ण करते हे सिद्ध करा:
\begin{align*}
  & \lforall[x][\eq[(x \times \Obj 0)][\Obj 0]] \tag{$\mathsf{Q}_6$}\\
  & \lforall[x][\lforall[y][\eq[(x \times y')][((x \times y) + x)]]] \tag{$\mathsf{Q}_7$}\\
  & \lforall[x][\lforall[y][(x < y \liff \lexists[z][\eq[(z' + x)][y]])]] \tag{$\mathsf{Q}_8$}
\end{align*}
फक्त~$\prime$ असलेले, $\Struct{N}$ मध्ये सत्य पण $\Struct{K}$ मध्ये असत्य
असे वाक्य शोधा.
\end{prob}

\begin{ex}
\label{mod:mar:mdq:ex:model-L-of-Q} $\Domain{L} = \Nat \cup \{a, b\}$ असा प्रांत
आणि पुढीलप्रमाणे दिलेली $\Assign{\prime}{L} = \nssucc$ व
$\Assign{+}{L} = \nsplus$ ही अर्थनिर्धारणे असलेली संरचना~$\Struct{L}$
विचारात घ्या:
\[\begin{array}{c|c}
  x & x^\nssucc \\
  \hline
  n & n+1 \\
  a & a\\
  b & b
\end{array}
\qquad
\begin{array}{c|ccc}
  x \nsplus y & m & a & b\\
  \hline
  n & n+m & b & a\\
  a & a & b & a\\
  b & b & b & a
\end{array}
\]
$\nssucc$ हे एकास-एक आहे, $0$ त्याच्या व्याप्तीत नाही, आणि
$0$ व्यतिरिक्त प्रत्येक~$x \in \Domain{L}$ त्याच्या व्याप्तीत आहे;
म्हणून $\mathsf{Q}_1$--$\mathsf{Q}_3$ ही स्वयंसिद्धके~$\Struct{L}$ मध्ये सत्य आहेत.
प्रत्येक $x$ साठी $x \nsplus 0 = x$, म्हणून $\mathsf{Q}_4$ सुद्धा सत्य आहे.
$\mathsf{Q}_5$ साठी $x \nsplus y^\nssucc$ आणि $(x \nsplus y)^\nssucc$ विचारात
घ्या. $x$ आणि $y$ दोन्ही मानक असतील, तर ते समान आहेत; कारण तेव्हा
$\nssucc$ आणि $\nsplus$ ही $\prime$ आणि $+$ यांच्याशी जुळतात. $x$ अमानक
आणि $y$ मानक असेल, तर $x \nsplus y^\nssucc = x = x^\nssucc =
(x \nsplus y)^\nssucc$. $x$ आणि $y$ दोन्ही अमानक असतील, तर चार प्रकरणे
आहेत:
\begin{align*}
&  a \nsplus a^\nssucc  = b = b^\nssucc = (a \nsplus a)^\nssucc\\
&  b \nsplus b^\nssucc  = a = a^\nssucc = (b \nsplus b)^\nssucc\\
&  b \nsplus a^\nssucc  = b = b^\nssucc = (b \nsplus a)^\nssucc\\
&  a \nsplus b^\nssucc  = a = a^\nssucc = (a \nsplus b)^\nssucc\\
\intertext{जर $x$ मानक पण $y$ अमानक असेल, तर}
&  n \nsplus a^\nssucc  = n \nsplus a = b = b^\nssucc = (n \nsplus a)^\nssucc\\
&  n \nsplus b^\nssucc  = n \nsplus b = a = a^\nssucc = (n \nsplus b)^\nssucc
\end{align*}

त्यामुळे $\Sat{L}{\mathsf{Q}_5}$. पण $a \nsplus 0 \neq 0 \nsplus a$, म्हणून
$\Sat/{L}{\lforall[x][\lforall[y][\eq[(x+y)][(y+x)]]]}$.
\end{ex}

\begin{prob}
\ref{mod:mar:mdq:ex:model-L-of-Q} मधील $\Struct{L}$ मध्ये
$\times$ आणि $<$ यांचे अर्थनिर्धारण करणारी $\nstimes$ आणि $\nsless$
समाविष्ट करून तिचा विस्तार करा. तुमची रचना~$\Th{Q}$ ची उरलेली
स्वयंसिद्धके पूर्ण करते हे दाखवा:
\begin{align*}
& \lforall[x][\eq[(x \times \Obj 0)][\Obj 0]] \tag{$\mathsf{Q}_6$}\\ &
  \lforall[x][\lforall[y][\eq[(x \times y')][((x \times y) + x)]]]
  \tag{$\mathsf{Q}_7$}\\ & \lforall[x][\lforall[y][(x < y \liff \lexists[z][\eq[(z'+x)][y]])]] \tag{$\mathsf{Q}_8$}
\end{align*}
\end{prob}

\begin{prob}
\ref{mod:mar:mdq:ex:model-L-of-Q} मधील $\Struct{L}$ मध्ये $a^\nssucc
= a$ आणि $b^\nssucc = b$. ज्या~$\Th{Q}$ च्या प्रतिमानात
$a^\nssucc = b$ आणि $b^\nssucc = a$ असे आहे, असे प्रतिमान अस्तित्वात आहे का?
\end{prob}

\begin{explain}
अमानक घटक हे मानक घटकांच्या ``पलीकडे'' असलेली~$\Th{Q}$ ची
प्रतिमाने आपण स्पष्टपणे रचली आहेत. किंबहुना, स्वयंसिद्धकांनुसार तितके असणे
आवश्यक आहे. अमानक घटक~$x$ हा ${} \nsless 0$ असू शकत नाही, कारण
$\Th{Q} \Proves \lforall[x][\lnot x<0]$
(\ref{inc:req:min:lem:less-zero} पाहा). तसेच, प्रत्येक $n$ साठी
$\Th{Q} \Proves \lforall[x][(x < \num{n}' \lif
(\eq[x][\num{0}] \lor \eq[x][\num{1}] \lor \dots \lor
\eq[x][\num{n}]))]$ (\ref{inc:req:min:lem:less-nsucc}); म्हणून
कोणत्याही~$n>0$ साठी $a \nsless n$ असू शकत नाही.
\end{explain}











\section{$\Th{PA}$ ची प्रतिमाने}\label{mod:mar:mpa:sec}

\begin{explain}
$\Th{TA}$ चे कोणतेही अमानक प्रतिमान हे~$\Th{PA}$ चेही अमानक प्रतिमान
असते. $\Th{TA}$ ची आणि म्हणून~$\Th{PA}$ ची अमानक प्रतिमाने अस्तित्वात
आहेत हे आपल्याला माहीत आहे. अशा अमानक प्रतिमानांत अमानक ``संख्या'', म्हणजे
सर्व मानक ``संख्यांच्या'' ``पलीकडे'' असलेले प्रांताचे घटक, असतात
हेही आपल्याला माहीत आहे. पण त्यांची मांडणी कशी असते? त्या किती असतात?
दुर्बलतर उपपत्ती~$\Th{Q}$ च्या प्रतिमानात फक्त एक अमानक संख्या असू शकते
हे आपण पाहिले. पण या साध्या संरचना~$\Th{PA}$ किंवा $\Th{TA}$ ची
प्रतिमाने नाहीत.

$\Th{PA}$ किंवा $\Th{TA}$ च्या प्रतिमानांची रचना समजून घेण्याची गुरुकिल्ली
म्हणजे या उपपत्तींमध्ये कोणती तथ्ये निष्पन्न करता येण्याजोगे आहेत हे पाहणे. उदाहरणार्थ,
$\Th{PA}$ कडून आधीच $\lforall[x][\eq/[x][x']]$ आणि
$\lforall[x][\lforall[y][\eq[(x+y)][(y+x)]]]$ सिद्ध होतात; म्हणून ही
वाक्ये असत्य असलेल्या साध्या रचना~$\Th{PA}$ ची प्रतिमाने असू शकत
नाहीत.

$\Struct{M}$ हे~$\Th{PA}$ चे प्रतिमान आहे असे समजा. मग
$\Th{PA} \Proves \varphi$ असल्यास $\Sat{M}{\varphi}$. पुन्हा
$\Assign{\Obj{0}}{M}$ करिता $\nszero$, $\Assign{\prime}{M}$ करिता
$\nssucc$, $\Assign{+}{M}$ करिता $\nsplus$,
$\Assign{\times}{M}$ करिता $\nstimes$ आणि $\Assign{<}{M}$ करिता
$\nsless$ वापरू. मग कोणतेही वाक्य~$\varphi$ हे $\nszero$, $\nssucc$,
$\nsplus$, $\nstimes$ आणि $\nsless$ यांविषयी एखादी अट सांगते; आणि
$\Sat{M}{\varphi}$ असेल, तर ती अट पूर्ण झाली पाहिजे. उदाहरणार्थ,
$\Sat{M}{\mathsf{Q}_1}$, म्हणजे
$\Sat{M}{\lforall[x][\lforall[y][(\eq[x'][y'] \lif \eq[x][y])]]}$,
असेल तर $\nssucc$ हे एकास-एक असले पाहिजे.
\end{explain}

\begin{prop}
$\Struct{M}$ मध्ये $\nsless$ हा काटेकोर रेषीय क्रम आहे, म्हणजे तो पुढील
अटी पूर्ण करतो:
\begin{enumerate}
\item कोणत्याही~$x \in \Domain{M}$ साठी $x \nsless x$ नाही.
\item $x \nsless y$ आणि $y \nsless z$ असतील, तर $x \nsless z$.
\item कोणत्याही $x \neq y$ साठी $x \nsless y$ किंवा $y \nsless x$.
\end{enumerate}
\end{prop}

\begin{proof}
$\Th{PA}$ कडून पुढील सूत्रे सिद्ध होतात:
\begin{enumerate}
\item $\lforall[x][\lnot x < x]$
\item $\lforall[x][\lforall[y][\lforall[z][((x < y \land y < z) \lif x < z)]]]$
\item $\lforall[x][\lforall[y][((x < y \lor y < x) \lor \eq[x][y]))]]$
\end{enumerate}
\end{proof}

\begin{prop}
\label{mod:mar:mpa:prop:M-discrete} $\nsless$-क्रमात $\nszero$ हा~$\Domain{M}$ मधील
लघुतम घटक आहे. प्रत्येक $x$ साठी $x \nsless x^\nssucc$, आणि हा
गुणधर्म असलेला $\nsless$-लघुतम घटक म्हणजे $x^\nssucc$. शून्याहून
वेगळ्या कोणत्याही~$x$ साठी $y^\nssucc = x$ असा एकमेव $y$ असतो. ($y$ ला
$\Struct{M}$ मधील~$x$ चा ``पूर्ववर्ती'' म्हणतो आणि तो~$^\nssucc x$ असा
दर्शवतो.)

\end{prop}

\begin{proof}
सराव.
\end{proof}

\begin{prob}
\ref{mod:mar:mpa:prop:M-discrete} मधील $\nszero$, $\nssucc$ आणि
$\nsless$ यांचे गुणधर्म सुनिश्चित करणारी, $\Th{PA}$ मध्ये निष्पन्न करता येण्याजोगे
(आणि म्हणून~$\Struct{N}$ मध्ये सत्य) असलेली~$\Lang{L_A}$ मधील
वाक्ये शोधा.
\end{prob}

\begin{prop}
$\Struct{M}$ चे सर्व मानक घटक हे सर्व अमानक घटक पेक्षा
($\nsless$ नुसार) लहान असतात.
\end{prop}

\begin{proof}
$\Struct{M}$ चा मानक घटक~$\Value{\num{n}}{M}$ याचा संक्षेप म्हणून
$n$ वापरू. कोणत्याही~$n \in \Nat$ साठी
$\lforall[x][(x < \num{n}' \lif (\eq[x][\num{0}] \lor
  \eq[x][\num{1}] \lor \dots \lor \eq[x][\num{n}]))]$ हे $\Th{Q}$
कडून आधीच सिद्ध होते. $\nsless \nszero$ असलेले घटक
नाहीत. म्हणून $n$ मानक आणि $x$ अमानक असेल, तर $x \nsless n$ असू शकत
नाही. व्याख्येनुसार, अमानक घटक म्हणजे कोणत्याही~$n \in \Nat$ साठी
$\Value{\num{n}}{M}$ नसलेला घटक; म्हणून $x \neq n$ सुद्धा. $\nsless$
हा रेषीय क्रम असल्यामुळे $n \nsless x$ असलेच पाहिजे.
\end{proof}

\begin{prop}
$\Domain{M}$ मधील प्रत्येक अमानक घटक~$x$ हा पुढील उपसंचाचा घटक
असतो:
\[\dots ^{\nssucc\nssucc\nssucc}x \nsless ^{\nssucc\nssucc}x \nsless
^{\nssucc}x \nsless x \nsless x^{\nssucc} \nsless x^{\nssucc\nssucc}
\nsless x^{\nssucc\nssucc\nssucc} \nsless \dots
\]
या उपसंचाला \emph{$x$ चा खंड} म्हणतो आणि तो $[x]$ असा लिहितो. त्याला
लघुतम किंवा महत्तम घटक नाही. एखाद्या मानक~$n$ साठी
$x \nsplus n = y$ किंवा $y \nsplus n = x$ असे असणाऱ्या
$y \in \Domain{M}$ चा संच, असे त्याचे लक्षण सांगता येते.
\end{prop}

\begin{proof}
असा संच~$[x]$ नेहमी अस्तित्वात असतो हे स्पष्ट आहे; कारण
$\Domain{M}$ मधील प्रत्येक अशून्य घटक~$y$ ला एकमेव
उत्तरवर्ती~$y^\nssucc$ आणि एकमेव पूर्ववर्ती~$^\nssucc y$ असतो. लागोपाठच्या
घटक $y$, $y^\nssucc$ साठी $y \nsless y^\nssucc$, आणि
$y$ ज्याच्यापेक्षा $\nsless$ आहे असा $\Domain{M}$ मधील $\nsless$-लघुतम
घटक म्हणजे $y^\nssucc$. नेहमी $^\nssucc y \nsless y$ आणि
$y \nsless y^\nssucc$ असल्यामुळे $[x]$ ला लघुतम किंवा महत्तम घटक
नाही. $y \in [x]$ असेल, तर $x \in [y]$; कारण मग एकतर
$y^{\nssucc\dots\nssucc} = x$ किंवा $x^{\nssucc\dots\nssucc} = y$.
$y^{\nssucc\dots\nssucc} = x$ ($n$ वेळा $\nssucc$) असेल, तर
$y \nsplus n = x$, आणि याचा व्यत्यासही खरा; कारण $n$ ही $\prime$ ची
संख्या असेल, तर $\Th{PA} \Proves
\lforall[x][\eq[x^{\prime\dots\prime}][(x + \num{n})]]$.
\end{proof}

\begin{prop}
$[x] \neq [y]$ आणि $x \nsless y$ असतील, तर कोणत्याही $u \in [x]$ आणि
कोणत्याही $v \in [y]$ साठी $u \nsless v$.
\end{prop}

\begin{proof}
$\Th{PA} \Proves \lforall[x][\lforall[y][(x < y \lif (x' < y
    \lor x' = y))]]$ हे लक्षात घ्या. म्हणून $u \nsless v$ आणि
$[u] \neq [v]$ असतील, तर कोणत्याही~$n$ साठी
$u \nsplus n^\nssucc \nsless v$ सुद्धा.

प्रत्येक $u \in [x]$ हा~$\nsless y$ आहे: गृहीतकानुसार
$x \nsless y$. $u \nsless x$ असेल, तर संक्रमणशीलतेने $u \nsless y$.
आणि $x \nsless u$ पण $u \in [x]$ असेल, तर कोणत्यातरी~$n$ साठी
$u = x\nsplus n^\nssucc$; म्हणून नुकत्याच सिद्ध केलेल्या तथ्यावरून
$u \nsless y$.

आता $v \in [y]$ हा~$\nsless y$ आहे, म्हणजे एखाद्या मानक~$m$
साठी $v \nsplus m^\nssucc = y$, असे समजा. यामुळे $v \nsless x$ असण्याची
शक्यता बाद होते; अन्यथा $y = v \nsplus m^\nssucc \nsless x$ झाले असते.
तसेच स्पष्टपणे $x \neq v$; अन्यथा
$x \nsplus m^\nssucc = v \nsplus m^\nssucc = y$ आणि $[x] = [y]$ झाले
असते. म्हणून $x \nsless v$. पण मग कोणत्याही~$n$ साठी
$x \nsplus n^\nssucc \nsless v$ सुद्धा. त्यामुळे $x \nsless u$ आणि
$u \in [x]$ असतील, तर $u \nsless v$. $u \nsless x$ असेल, तर
संक्रमणशीलतेने $u \nsless v$.

शेवटी, $y \nsless v$ असेल, तर आपण दाखवल्याप्रमाणे $u \nsless y$ आणि
$y \nsless v$ असल्यामुळे $u \nsless v$.
\end{proof}

\begin{cor}
$[x] \neq [y]$ असेल, तर $[x] \cap [y] = \emptyset$.
\end{cor}

\begin{proof}
$z \in [x]$ आणि $x \nsless y$ असे समजा. मग प्रत्येक $u \in [y]$ साठी
$z \nsless u$. जर $z \in [y]$ असेल, तर $z \nsless z$ मिळेल.
$y \nsless x$ असेल तेव्हाही याचप्रमाणे.
\end{proof}

\begin{explain}
याचा अर्थ खंडांनाच $\nsless$ चा आदर राखणारा क्रम देता येतो:
$[x] \nsless [y]$ हे तेव्हाच, जेव्हा $x \nsless y$; किंवा सममूल्य रीतीने,
कोणत्याही $u \in [x]$ आणि $v \in [y]$ साठी $u \nsless v$. मानक खंड
$[0]$ हा लघुतम खंड आहे हे स्पष्ट आहे. त्याचा कोणत्याही अमानक खंडाशी छेद
नाही, आणि कोणत्याही दोन अमानक खंडांचाही परस्परांशी छेद नाही. विशेषतः,
उत्तरवर्ती किंवा पूर्ववर्ती वारंवार घेऊन वेगळ्या खंडात ``पोहोचता'' येत नाही.
\end{explain}

\begin{prop}
$x$ आणि $y$ अमानक असतील, तर $x \nsless x \nsplus y$ आणि
$x \nsplus y \notin [x]$.
\end{prop}

\begin{proof}
$y$ अमानक असेल, तर $y \neq \nszero$. $\Th{PA} \Proves
\lforall[x][(y \neq \Obj{0} \lif x < (x+y))]$. आता
$x \nsplus y \in [x]$ असे समजा. $x \nsless x \nsplus y$ असल्यामुळे
$x \nsplus n^\nssucc = x \nsplus y$ असे झाले असते. पण $\Th{PA} \Proves
\lforall[x][\lforall[y][\lforall[z][(\eq[(x+y)][(x+z)] \lif y = z)]]]$
(बेरीजेसाठीचा निरसन नियम). याचा अर्थ एखाद्या मानक~$n$ साठी
$y = n^\nssucc$ असा झाला असता; पण $y$ अमानक आहे असे गृहीत आहे.
\end{proof}

\begin{prop}
लघुतम अमानक खंड नाही.
\end{prop}

\begin{proof}
$\Th{PA} \Proves \lforall[x][\lexists[y][(\eq[(y+y)][x] \lor
      \eq[(y+y)'][x])]]$; म्हणजे प्रत्येक $x$ ला~$2$ ने भाग जातो
(शक्यतो बाकी~$1$ ठेवून). $x$ अमानक असेल, तर $y$ सुद्धा अमानक असतो.
आधीच्या विधानावरून $y \nsless y \nsplus y$ आणि
$y \nsplus y \notin [y]$. मग $y \nsless (y \nsplus y)^\nssucc$ आणि
$(y \nsplus y)^\nssucc \notin [y]$ सुद्धा. पण $x = y \nsplus y$ किंवा
$x = (y \nsplus y)^\nssucc$, म्हणून $y \nsless x$ आणि $y \notin [x]$.
\end{proof}

\begin{prop}
महत्तम खंड नाही.
\end{prop}

\begin{proof}
सराव.
\end{proof}

\begin{prob}
$\Th{PA}$ च्या अमानक प्रतिमानात महत्तम खंड नसतो हे दाखवा.
\end{prob}

\begin{prop}
\label{mod:mar:mpa:prop:blocks-dense}
खंडांचा क्रम सघन आहे. म्हणजे, $x \nsless y$ आणि $[x] \neq [y]$ असतील,
तर त्या दोहोंपेक्षा वेगळा आणि त्यांच्यामध्ये असलेला एक खंड~$[z]$ असतो.
\end{prop}

\begin{proof}
$x \nsless y$ असे समजा. पूर्वीप्रमाणे, $x \nsplus y$ ला दोनने भाग जातो
(शक्यतो बाकी ठेवून): एखादा $z \in \Domain{M}$ असा असतो की एकतर
$x \nsplus y = z \nsplus z$ किंवा
$x \nsplus y = (z \nsplus z)^\nssucc$. घटक~$z$ हा $x$ आणि~$y$ यांची
``सरासरी'' आहे, आणि $x \nsless z$ व $z \nsless y$.

\end{proof}

\begin{prob}
\ref{mod:mar:mpa:prop:blocks-dense} चा तपशीलवार पुरावा लिहा.
$z$ चे अस्तित्व सुनिश्चित करण्यासाठी $\Th{PA}$ ने कोणते वाक्य
निष्पन्न केले पाहिजे? $x \nsless z$ आणि $z \nsless y$ का, तसेच
$[x] \neq [z]$ आणि $[z] \neq [y]$ का?
\end{prob}

\begin{explain}
प्रतिमान गणनीय असेल, तर अमानक खंडांचा क्रम परिमेय संख्यांच्या क्रमासारखा
असतो: ते अंतबिंदुविरहित गणनीय अनंत सघन रेषीय क्रम बनवतात. असे
कोणतेही दोन गणनीय अनंत क्रम समरूपी असतात हे दाखवता येते. त्यावरून
सत्य अंकगणिताच्या कोणत्याही दोन गणनीय अमानक प्रतिमान
$\Struct{M}_1$ आणि $\Struct{M_2}$ यांची फक्त $<$ आणि $=$ असलेल्या भाषेतील
न्यूनीकरणे समरूपी आहेत. खरेच, पुढीलप्रमाणे समरूपण~$h$ ची व्याख्या करता
येते: $\Struct{M_1}$ आणि $\Struct{M_2}$ यांचे मानक भाग मानक
प्रतिमान~$\Struct{N}$ शी आणि म्हणून परस्परांशी समरूपी आहेत. अमानक भाग
बनवणाऱ्या खंडांचा क्रम स्वतः परिमेय संख्यांसारखा असल्यामुळे ते समरूपी आहेत;
प्रत्येक खंडातील स्वेच्छ घटक परस्पर जुळवून आणि नंतर $\Struct{M_1}$ मधील
$x$ च्या उत्तरवर्तीची प्रतिमा ही $\Struct{M_2}$ मधील $x$ च्या प्रतिमेचा
उत्तरवर्ती घेऊन, खंडांचे समरूपण खंडांच्या \emph{आत} विस्तारित करता येते.
$\mathfrak{M}_1$ आणि $\mathfrak{M}_2$ या अंकगणिताच्या पूर्ण भाषेत समरूपी
आहेत असे यावरून \emph{निष्पन्न होत नाही} हे लक्षात घ्या (खरे तर समरूपण
नेहमी भाषा सापेक्ष असते), कारण $\Domain{M_1}$ आणि
$\Domain{M_2}$ वर बेरीज आणि गुणाकार परिभाषित करण्याचे परस्पर-असमरूपी मार्ग
आहेत. (पूर्ण उपपत्तीच्या गणनीय प्रतिमानांची संख्या~2 असू शकत नाही, या
वॉट यांच्या प्रसिद्ध प्रमेयावरूनही हे निष्पन्न होते.)

\end{explain}










\section{अंकगणिताची संगणनक्षम प्रतिमाने}\label{mod:mar:cmp:sec}

\begin{explain}
मानक प्रतिमान~$\Struct{N}$ ला दोन उपयुक्त वैशिष्ट्ये आहेत. त्याचा प्रांत
नैसर्गिक संख्यांचा संच~$\Nat$ आहे, म्हणजे त्याचे घटक हे नेमके तेच प्रकारचे
विषय आहेत ज्यांच्याविषयी अंकगणिताची भाषा वापरून आपल्याला बोलायचे असते, आणि
मानक संख्यांक~$\num{n}$ हा खरोखर~$n$ लाच निर्देशित करतो. दुसरे उपयुक्त
वैशिष्ट्य म्हणजे~$\Lang{L_A}$ मधील अतार्किक चिन्हांची सर्व अर्थनिर्धारणे
\emph{संगणनक्षम} आहेत. $\Assign{\prime}{N}$, $\Assign{+}{N}$ आणि
$\Assign{\times}{N}$ म्हणून काम करणारी उत्तरवर्ती, बेरीज आणि गुणाकार ही
फलने संख्यांवरील संगणनक्षम फलने आहेत. (म्हणजे ट्यूरिंग यंत्रांनी संगणन करता
येणारी किंवा, उदाहरणार्थ, आदिम पुनरावर्तनाने परिभाषित करता येणारी.) तसेच
$\Struct{N}$ वरील लघुतर संबंध, म्हणजे~$\Assign{<}{N}$, निर्णेय आहे.

$\Th{Q}$ आणि $\Th{PA}$ यांसारख्या अंकगणितीय उपपत्तींच्या अमानक प्रतिमानांत
अमानक घटक असलेच पाहिजेत. त्यामुळे त्यांच्या प्रांतांत सामान्यतः~$\Nat$
व्यतिरिक्तही घटक असतात. तरी कोणतीही गणनीय संरचना कोणत्याही
गणनीय अनंत संचावर, त्यात~$\Nat$ वरही, उभारता येते. म्हणून प्रांत
$\Nat$ असलेली अमानक प्रतिमानेही आहेत. अशा~$\Struct{M}$ प्रतिमानांत अर्थातच
किमान काही संख्या आपली नेहमीची भूमिका बजावू शकत नाहीत, कारण असा काही~$k$
असतो की प्रत्येक~$n \in \Nat$ साठी तो~$\Value{\num{n}}{M}$ पेक्षा
वेगळा असतो.
\end{explain}

\begin{defn}
$\Lang{L_A}$ साठीची संरचना~$\Struct{M}$ ही \emph{संगणनक्षम}
तेव्हाच, जेव्हा $\Domain{M} = \Nat$ असते, $\Assign{\prime}{M}$,
$\Assign{+}{M}$ आणि $\Assign{\times}{M}$ ही संगणनक्षम फलने असतात, आणि
$\Assign{<}{M}$ हा निर्णेय संबंध असतो.
\end{defn}

\begin{ex}\label{mod:mar:cmp:ex:comp-model-q}
\ref{mod:mar:mdq:ex:model-K-of-Q} मधील रचना~$\Struct{K}$ आठवा. तिचा प्रांत
$\Domain{K} = \Nat
\cup \{a\}$ होता आणि अर्थनिर्धारणे पुढीलप्रमाणे होती:
\begin{align*}
  \Assign{\Obj{0}}{K} & = 0\\
  \Assign{\prime}{K}(x) & =
  \begin{cases}
    x+1 & \text{जर $x\in \Nat$}\\
    a & \text{जर $x = a$}
  \end{cases}\\
  \Assign{+}{K}(x, y) & =
  \begin{cases}
    x+y & \text{जर $x$, $y \in\Nat$}\\
    a & \text{इतर प्रसंगी}
  \end{cases}\\
  \Assign{\times}{K}(x, y) & =
  \begin{cases}
    xy & \text{जर $x$, $y \in\Nat$}\\
    0 & \text{जर $x=0$ किंवा $y=0$}\\
    a & \text{इतर प्रसंगी}\\
  \end{cases}\\
  \Assign{<}{K} & =
  \Setabs{\tuple{x,y}}{x, y \in \Nat \text{ आणि } x<y} \cup
  \Setabs{\tuple{x,a}}{x \in \Domain{K}}
\end{align*}

पण $\Domain{K}$ हा गणनीय अनंत आहे आणि म्हणून तो~$\Nat$ शी तुल्यबल
आहे. उदाहरणार्थ, $g\colon \Nat \to \Domain{K}$ असे घेऊ, जिथे $g(0) =
a$ आणि $n>0$ साठी $g(n) = n-1$; हे एक एकास-एक आच्छादन आहे.

त्याचा वापर करून नव्या~$\Th{Q}$-प्रतिमान~$\Struct{K'}$ आणि
$\Struct{K}$ यांमध्ये समरूपता मिळवता येते. $\Struct{K'}$ मध्ये
$\Lang{L_A}$ च्या चिन्हांना वेगळी फलने आणि संबंध नेमावे लागतात, कारण
$\Nat$ मधील वेगवेगळे घटक मानक आणि अमानक संख्यांच्या भूमिका
बजावतात.

विशेषतः, आता~$0$ हा लघुतम मानक संख्या नसून~$a$ ची भूमिका बजावतो. आता~$1$
ही लघुतम मानक संख्या आहे. म्हणून $\Assign{\Obj{0}}{K'} = 1$ असे नेमू.
आता उत्तरवर्ती फलनही वेगळे आहे: एखादी मानक संख्या, म्हणजे $n > 0$, दिली
असता ते अजूनही $n+1$ परत करते. पण आता~$0$ हा~$a$ ची भूमिका बजावतो आणि
तो स्वतःचाच उत्तरवर्ती आहे. म्हणून $\Assign{\prime}{K'}(0) = 0$.
बेरीज आणि गुणाकारासाठी याचप्रमाणे
\begin{align*}
\Assign{+}{K'}(x, y) & =
  \begin{cases}
    x+y-1 & \text{जर $x$, $y >0$}\\
    0 & \text{इतर प्रसंगी}
  \end{cases}\\
  \Assign{\times}{K'}(x, y) & =
  \begin{cases}
    1 & \text{जर $x = 1$ किंवा $y = 1$}\\
    xy - x - y + 2 & \text{जर $x$, $y > 1$}\\
    0 & \text{इतर प्रसंगी}\\
  \end{cases}
\end{align*}
आणि $\tuple{x, y} \in \Assign{<}{K'}$ हे तेव्हाच, जेव्हा $x < y$, $x > 0$
आणि $y > 0$ असतात, किंवा $y = 0$ असतो.

ही सर्व फलने नैसर्गिक संख्यांची संगणनक्षम फलने आहेत आणि
$\Assign{<}{K'}$ हा~$\Nat$ वरील निर्णेय संबंध आहे---पण ती~$\Nat$ वरील
उत्तरवर्ती, बेरीज आणि गुणाकार हीच फलने नाहीत, तसेच $\Assign{<}{K'}$ हाही
$\Nat$ वरील~$<$ हाच संबंध नाही.
\end{ex}

\begin{prob}
\ref{mod:mar:mdq:ex:model-L-of-Q} मधील~$\Struct{L}$ शी समरूपी,
$\Domain{L'} = \Nat$ असलेली संरचना~$\Struct{L'}$ द्या.
\end{prob}

\begin{explain}
\ref{mod:mar:cmp:ex:comp-model-q} यावरून~$\Th{Q}$ ला प्रांत~$\Nat$ असलेली
संगणनक्षम अमानक प्रतिमाने आहेत हे दिसते. पण पुढील निष्कर्ष दाखवतो की
$\Th{PA}$ च्या प्रतिमानांसाठी (आणि म्हणून~$\Th{TA}$ च्या प्रतिमानांसाठीही)
हे खरे नाही.
\end{explain}

\begin{thm}[टेनेनबाउमचे प्रमेय]
$\Struct{N}$ हे~$\Th{PA}$ चे एकमेव संगणनक्षम प्रतिमान आहे.
\end{thm}



\chapter{अंतर्वेशन प्रमेय}\label{mod:int::chap}









\section{परिचय}\label{mod:int:int:sec}

अंतर्वेशन प्रमेय म्हणजे पुढील निष्कर्ष: समजा
$\Entails \varphi \lif \psi$. मग असे वाक्य~$\chi$ असते की
$\Entails \varphi \lif \chi$ आणि $\Entails \chi \lif \psi$. शिवाय,
$\chi$ मधील प्रत्येक स्थिरांक, फलन आणि विधेय
($\eq$ खेरीज) हे $\varphi$ आणि~$\psi$ या दोन्हींमध्ये आढळते. वाक्य~$\chi$
ला $\varphi$ आणि~$\psi$ यांचे \emph{अंतर्वेशक} म्हणतात.

अंतर्वेशन प्रमेय स्वतःही रंजक आहे; पण त्याचे मुख्य महत्त्व असे की,
उपपत्तीतील परिभाष्यतेविषयीचे निष्कर्ष आणि दोन सुसंगत उपपत्ती एकत्र केल्यावर
मिळणारी उपपत्ती कोणत्या अटींवर सुसंगत असते याविषयीचे निष्कर्ष सिद्ध करण्यासाठी
ते वापरता येते. पहिला निष्कर्ष बेथचे परिभाष्यता प्रमेय म्हणून, तर
दुसरा रॉबिन्सनचे संयुक्त सुसंगतता प्रमेय म्हणून ओळखला जातो.












\section{वाक्याचे विलगीकरण}\label{mod:int:sep:sec}

अंतर्वेशन प्रमेयाची सिद्धता सुरू करण्यापूर्वी थोडी पूर्वतयारी आवश्यक आहे.
$\varphi$ आणि~$\psi$ यांचे अंतर्वेशक म्हणजे असे वाक्य~$\chi$ की
$\varphi \Entails \chi$ आणि $\chi \Entails \psi$. प्रतिपरिवर्तनानुसार, यांपैकी
नंतरचे विधान तेव्हा आणि केवळ तेव्हाच खरे असते, जेव्हा $\lnot \psi
\Entails \lnot \chi$. हा गुणधर्म असलेले वाक्य~$\chi$ हे $\varphi$
आणि~$\lnot \psi$ यांना \emph{विलग करते} असे म्हणतात. म्हणून $\varphi$ आणि~$\psi$
यांचे अंतर्वेशक शोधणे म्हणजे $\varphi$ आणि~$\lnot \psi$ यांना विलग करणारे
वाक्य शोधणेच होय. नेहमीप्रमाणे, याचे व्यापक रूप विचारात घेणे उपयोगी
ठरेल: वाक्यांच्या दोन \emph{संचांना} विलग करणारे वाक्य.

\begin{defn}
वाक्य~$\chi$ हे वाक्यसंच~$\Gamma$ आणि~$\Delta$ यांना \emph{विलग करते}
तेव्हा आणि केवळ तेव्हाच, जेव्हा $\Gamma \Entails \chi$ आणि $\Delta
\Entails \lnot \chi$. असे कोणतेही वाक्य अस्तित्वात नसेल, तर
$\Gamma$ आणि~$\Delta$ यांना \emph{अविलगनीय} म्हणतात.
\end{defn}

$\Gamma$, $\Delta$ आणि~$\chi$ यांच्या प्रतिमान-वर्गांमधील समावेशन-संबंध
खाली दर्शवले आहेत:
\begin{figure}[h]
  \centering
  \begin{tikzpicture}[node distance=2cm, auto, thick]
    \draw [rounded corners] (0,0) -- (6,0) -- (6,3) -- (0,3) --  cycle;
    \draw (1.5,1.5) circle (0.9cm);
    \draw (4.5,1.5) circle (0.9cm);
    \path node at (1.5,1.5) {\Large $\Gamma$};
    \path node at (4.5,1.5) {\Large $\Delta$};
    \path node at (0.4,2.6) {\Large $\formula{C}$};
    \path node at (3.4,0.4) {\Large $\lnot \formula{C}$};
    \draw (2.5,0) .. controls (2.5,1.5) and (3.5,1.5) .. (3.5,3);
  \end{tikzpicture}
  \caption{$\chi$ हे $\Gamma$ आणि $\Delta$ यांना विलग करते}
  \label{mod:int:sep:fig:sep}
\end{figure}


\begin{lem}\label{mod:int:sep:lem:sep1}
समजा $\Lang{L}_0$ या भाषेत $\Gamma$ आणि~$\Delta$ या \emph{दोन्हींमध्ये}
आढळणारे प्रत्येक स्थिरांक, फलन आणि विधेय ($\doteq$
खेरीज) येते; आणि प्रत्येक $n \ge 0$ साठी अनंत नवीन स्थिरांक
$\Obj c_n$ जोडून $\Lang{L}'_0$ मिळवली आहे. मग $\Gamma$ आणि~$\Delta$ हे
$\Lang{L}_0$ मध्ये अविलगनीय असतील, तर ते $\Lang{L}'_0$ मध्येही अविलगनीय
असतात.
\end{lem}

\begin{proof}
विरोधाभास पद्धतीने सिद्ध करू: विरुद्धार्थी असे समजा की $\Gamma$ आणि
$\Delta$ हे $\Lang{L}'_0$ मध्ये विलगनीय आहेत. मग एखाद्या $\chi \in
\Lang{L}_0$ साठी $\Gamma \Entails \Subst{\chi}{c}{x}$ आणि $\Delta
\Entails \lnot \Subst{\chi}{c}{x}$ असते (येथे $c$ हे नवीन स्थिरांक
आहे---$\chi$ मध्ये अशी एकाहून अधिक नवीन स्थिरांक असलेले प्रकरणही यासारखेच
आहे). संहततेनुसार, $\Gamma$ चा \emph{सांत} उपसंच~$\Gamma_0$ आणि
$\Delta$ चा सांत उपसंच~$\Delta_0$ असे आहेत की $\Gamma_0 \Entails
\Subst{\chi}{c}{x}$ आणि $\Delta_0 \Entails \lnot \Subst{\chi}{c}{x}$.
$\Gamma_0$ मधील सर्व सूत्रांचा संयोग~$\varphi_{G}$ आणि $\Delta_0$ मधील सर्व
सूत्रांचा संयोग~$\varphi_{H}$ असू द्या. मग
\begin{align*}
  \varphi_{G} & \Entails \Subst{\chi}{c}{x}, & \varphi_{H}  \Entails \lnot \Subst{\chi}{c}{x}.
\end{align*}
पहिल्या विधानावर सामान्यीकरण लावल्यास $\varphi_{G} \Entails \lforall[x][\chi]$
मिळते; आणि दुसऱ्यावर प्रतिपरिवर्तन लावल्यास $\Subst{\chi}{c}{x}
\Entails \lnot \varphi_{H}$ मिळते, म्हणून $\lforall[x][\chi] \Entails \lnot \varphi_{H}$.

पुन्हा प्रतिपरिवर्तन केल्यास $\varphi_{H} \Entails \lnot \lforall[x][\chi]$
मिळते. एकस्वनिकतेनुसार,
\begin{align*}
  \Gamma &\Entails \lforall[x][\chi], &
  \Delta & \Entails \lnot \lforall[x][\chi],
\end{align*}
त्यामुळे $\lforall[x][\chi]$ हे $\Lang{L}_0$ मध्ये $\Gamma$ आणि~$\Delta$
यांना विलग करते.
\end{proof}

\begin{lem}\label{mod:int:sep:lem:sep2}
समजा $\Gamma \cup \{ \lexists[x][\varphi_{S}] \}$ आणि~$\Delta$ हे अविलगनीय
आहेत, आणि $c$ हे $\Gamma$, $\Delta$ किंवा~$\varphi_{S}$ यांपैकी कशातही नसलेले नवीन
स्थिरांक आहे. मग $\Gamma \cup \{ \lexists[x][\varphi_{S}],
\Subst{\varphi_{S}}{c}{x} \}$ आणि~$\Delta$ हेही अविलगनीय असतात.
\end{lem}

\begin{proof}
विरोधाभासासाठी असे समजा की $\chi$ हे $\Gamma \cup \{ \lexists[x][\varphi_{S}],
\Subst{\varphi_{S}}{c}{x}\}$ आणि~$\Delta$ यांना विलग करते, पण त्याच वेळी
$\Gamma \cup \{\lexists[x]{\varphi_{S}} \}$ आणि~$\Delta$ हे अविलगनीय आहेत. दोन
प्रकरणे वेगळी करू:
\begin{enumerate}
\item $\chi$ मध्ये $c$ आढळत नाही: या प्रकरणात $\Gamma \cup
  \{\lexists[x][\varphi_{S}], \lnot\chi \}$ हा संच पूर्ततायोग्य असतो (अन्यथा
  $\chi$ हे $\Gamma \cup \{\lexists[x][\varphi_{S}] \}$ आणि~$\Delta$ यांना विलग
  करेल). $\Subst{\varphi_{S}}{c}{x}$ ची भर घातल्यानंतरही तो पूर्ततायोग्य राहतो;
  म्हणून अखेरीस $\chi$ हे $\Gamma \cup \{ \lexists[x][\varphi_{S}],
  \Subst{\varphi_{S}}{c}{x} \}$ आणि~$\Delta$ यांना विलग करत नाही.
\item $\chi$ मध्ये $c$ आढळते, त्यामुळे $\chi$ चे रूप
  $\Subst{\chi}{c}{x}$ असते. मग
  \[  \Gamma \cup \{ \lexists[x][\varphi_{S}], \Subst{\varphi_{S}}{c}{x}\} \Entails \Subst{\chi}{c}{x},
  \]
  असे असते. यावरून निगमन प्रमेय आणि सामान्यीकरण यांच्या साहाय्याने
  $\Gamma, \lexists[x][\varphi_{S}] \Entails \lforall[x][(\varphi_{S} \lif \chi)]$ आणि
  अखेरीस $\Gamma \cup \{ \lexists[x][\varphi_{S}] \} \Entails
  \lexists[x][\chi]$ मिळते. दुसरीकडे, $\Delta \Entails \lnot
  \Subst{\chi}{c}{x}$; म्हणून सामान्यीकरणाने $\Delta \Entails \lnot
  \lexists[x][\chi]$. त्यामुळे $\Gamma \cup \{\lexists[x][\varphi_{S}] \}$
  आणि~$\Delta$ हे विलगनीय ठरतात, आणि हा विरोधाभास आहे.
\end{enumerate}
\end{proof}












\section{क्रेगचे अंतर्वेशन प्रमेय}\label{mod:int:prf:sec}

\begin{thm}[क्रेगचे अंतर्वेशन प्रमेय]
\label{mod:int:prf:thm:interpol}
$\Entails \varphi \lif \psi$ असेल, तर असे वाक्य~$\chi$ असते की
$\Entails \varphi \lif \chi$ आणि $\Entails \chi \lif \psi$; तसेच $\chi$ मधील
प्रत्येक स्थिरांक, फलन आणि विधेय ($\eq$ खेरीज) हे
$\varphi$ आणि~$\psi$ या दोन्हींमध्ये आढळते. वाक्य~$\chi$ ला $\varphi$ आणि~$\psi$
यांचे \emph{अंतर्वेशक} म्हणतात.
\end{thm}

\begin{proof}
समजा $\Lang{L}_1$ ही $\varphi$ ची भाषा आणि $\Lang{L}_2$ ही $\psi$ ची भाषा
आहे. $\Lang{L}_0 = \Lang{L}_1 \cap \Lang{L}_2$ असे ठेवू. प्रत्येक
$i \in \{0, 1, 2 \}$ साठी, $\Lang{L}_i$ मध्ये अनंत नवीन स्थिरांक
$\Obj c_0, \Obj c_1, \Obj c_2, \dots$ जोडून $\Lang{L}'_i$ मिळवू.

$\varphi$ अपूर्ततायोग्य असेल, तर $\lexists[x][\eq/[x][x]]$ हे अंतर्वेशक आहे.
$\lnot \psi$ अपूर्ततायोग्य असेल (आणि म्हणून $\psi$ वैध असेल), तर
$\lexists[x][\eq[x][x]]$ हे अंतर्वेशक आहे. म्हणून $\varphi$ आणि~$\lnot \psi$
ही दोन्ही पूर्ततायोग्य आहेत असेही आपण मानू शकतो.

अंतर्वेशन प्रमेयाचे प्रतिपरिवर्तित विधान सिद्ध करण्यासाठी, $\varphi$ आणि~$\psi$
यांचे कोणतेही अंतर्वेशक नाही असे गृहीत धरू. दुसऱ्या शब्दांत, $\{\varphi\}$
आणि~$\{\lnot \psi\}$ हे $\Lang{L}_0$ मध्ये अविलगनीय आहेत असे गृहीत धरू.

$(\{ \varphi \}, \{\lnot\psi\})$ या जोडीचा विस्तार करून महत्तम अविलगनीय जोडी
$(\Gamma^*, \Delta^*)$ मिळवणे हे आपले उद्दिष्ट आहे. $\varphi_0$, $\varphi_1$,
$\varphi_2$, \dots{} हे $\Lang{L}_1$ च्या वाक्यांचे प्रगणन आणि $\psi_0$,
$\psi_1$, $\psi_2$, \dots{} हे $\Lang{L}_2$ च्या वाक्यांचे प्रगणन असू
द्या. प्रत्येक $n \ge 0$ साठी वाक्यांच्या संचांचे दोन वर्धमान अनुक्रम
$(\Gamma_n, \Delta_n)$ पुढीलप्रमाणे परिभाषित करू. $\Gamma_0 = \{ \varphi\}$
आणि $\Delta_0 = \{\lnot \psi \}$ असे ठेवू. $(\Gamma_n, \Delta_n)$ आधीच
परिभाषित आहेत असे मानून $\Gamma_{n+1}$ आणि $\Delta_{n+1}$ पुढीलप्रमाणे
परिभाषित करू:
\begin{enumerate}
\item $\Gamma_n \cup \{\varphi_n \}$ आणि~$\Delta_n$ हे $\Lang{L}'_0$ मध्ये
  अविलगनीय असतील, तर $\varphi_n$ चा $\Gamma_{n+1}$ मध्ये समावेश करा. शिवाय,
  $\varphi_n$ हे $\lexists[x][\varphi_{S}]$ या रूपातील अस्तित्ववाची सूत्र असेल,
  तर $\Gamma_n$, $\Delta_n$, $\varphi_n$ किंवा~$\psi_n$ यांपैकी कशातही न आढळणारे
  नवीन स्थिरांक~$c$ निवडा आणि $\Subst{\varphi_{S}}{c}{x}$ चा $\Gamma_{n+1}$
  मध्ये समावेश करा.
\item $\Gamma_{n+1}$ आणि~$\Delta_n \cup \{\psi_n \}$ हे $\Lang{L}'_0$
  मध्ये अविलगनीय असतील, तर $\psi_n$ चा $\Delta_{n+1}$ मध्ये समावेश करा.
  शिवाय, $\psi_n$ हे $\lexists[x][\varphi_{S}]$ या रूपातील अस्तित्ववाची सूत्र
  असेल, तर $\Gamma_{n+1}$, $\Delta_n$, $\varphi_n$ किंवा~$\psi_n$ यांपैकी कशातही
  न आढळणारे नवीन स्थिरांक~$c$ निवडा आणि $\Subst{\varphi_{S}}{c}{x}$ चा
  $\Delta_{n+1}$ मध्ये समावेश करा.
\end{enumerate}
शेवटी, पुढील व्याख्या करू:
\begin{align*}
  \Gamma^* & = \bigcup_{n\ge 0} \Gamma_n, &
  \Delta^* & = \bigcup_{n\ge 0} \Delta_n.
\end{align*}
आता $n$ वरील युगपत् विगमनाने पुढील गोष्टी सिद्ध करता येतात:
\begin{enumerate}
\item\label{mod:int:prf:part-a} $\Gamma_n$ आणि~$\Delta_n$ हे $\Lang{L}'_0$ मध्ये
  अविलगनीय आहेत;
\item\label{mod:int:prf:part-b} $\Gamma_{n+1}$ आणि~$\Delta_n$ हे $\Lang{L}'_0$
  मध्ये अविलगनीय आहेत.
\end{enumerate}
\ref{mod:int:prf:part-a} चा आधार \ref{mod:int:sep:lem:sep1} मुळे मिळतो. \ref{mod:int:prf:part-b}
साठी तीन प्रकरणे वेगळी करावी लागतात:
\begin{enumerate}
\item $\Gamma_0 \cup \{\varphi_0 \}$ आणि~$\Delta_0$ हे विलगनीय असतील, तर
  $\Gamma_1 = \Gamma_0$; त्यामुळे \ref{mod:int:prf:part-b} म्हणजे फक्त
  \ref{mod:int:prf:part-a} होय.
\item $\Gamma_1 = \Gamma_0 \cup\{ \varphi_0\}$ असेल, तर रचनेनुसार
  $\Gamma_1$ आणि~$\Delta_0$ हे अविलगनीय आहेत.
\item आता फक्त $\varphi_0$ अस्तित्ववाची असलेले प्रकरण उरते; तेव्हा
  $\Gamma_1 = \Gamma_0 \cup \{ \lexists[x][\varphi_{S}], \Subst{\varphi_{S}}{c}{x}
  \}$. रचनेनुसार $\Gamma_0 \cup \{ \lexists[x][\varphi_{S}]\}$ आणि~$\Delta_0$
  हे अविलगनीय आहेत; म्हणून \ref{mod:int:sep:lem:sep2} नुसार $\Gamma_0 \cup
  \{ \lexists[x][\varphi_{S}], \Subst{\varphi_{S}}{c}{x} \}$ आणि~$\Delta_0$ हेही
  अविलगनीय आहेत.
\end{enumerate}
यावरून वरील \ref{mod:int:prf:part-a} आणि \ref{mod:int:prf:part-b} यांच्यासाठी विगमनाचा आधार
पूर्ण होतो. आता विगमन-पायरी पाहू. \ref{mod:int:prf:part-a} साठी, $\Delta_{n+1} =
\Delta_n \cup \{ \psi_n \}$ असेल, तर रचनेनुसार $\Gamma_{n+1}$ आणि
$\Delta_{n+1}$ हे अविलगनीय आहेत ($\psi_n$ अस्तित्ववाची असतानाही
\ref{mod:int:sep:lem:sep2} नुसार हेच मिळते); आणि $\Delta_{n+1} = \Delta_n$
असेल ($\Gamma_{n+1}$ आणि~$\Delta_n \cup \{\psi_n\}$ विलगनीय असल्यामुळे),
तर \ref{mod:int:prf:part-b} साठीचे विगमन-गृहीतक वापरतो. \ref{mod:int:prf:part-b} च्या
विगमन-पायरीसाठी, $\Gamma_{n+2} = \Gamma_{n+1} \cup \{\varphi_{n+1} \}$
असेल, तर रचनेनुसार $\Gamma_{n+2}$ आणि~$\Delta_{n+1}$ हे अविलगनीय आहेत
($\varphi_{n+1}$ अस्तित्ववाची असतानाही \ref{mod:int:sep:lem:sep2} नुसार हेच
मिळते); आणि $\Gamma_{n+2} = \Gamma_{n+1}$ असेल, तर नुकतेच सिद्ध केलेले
\ref{mod:int:prf:part-a} चे विगमन-प्रकरण वापरतो. अशा रीतीने \ref{mod:int:prf:part-a} आणि
\ref{mod:int:prf:part-b} यांवरील विगमन पूर्ण होते.

यावरून $\Gamma^*$ आणि~$\Delta^*$ हे अविलगनीय आहेत. तसे नसते, तर
संहततेनुसार असा $n \ge 0$ असता की $\Gamma_n$ आणि~$\Delta_n$ विलगनीय
असते, आणि हे \ref{mod:int:prf:part-a} च्या विरुद्ध आहे. विशेषतः, $\Gamma^*$ आणि
$\Delta^*$ हे सुसंगत आहेत: कारण यांपैकी अनुक्रमे पहिला किंवा दुसरा विसंगत
असता, तर $\lexists[x][\eq/[x][x]]$ किंवा
$\lforall[x][\eq[x][x]]$ यांनी ते विलग झाले असते.

आता $\Gamma^*$ हा $\Lang{L}'_1$ मध्ये महत्तम सुसंगत आहे आणि त्याचप्रमाणे
$\Delta^*$ हा $\Lang{L}'_2$ मध्ये महत्तम सुसंगत आहे, हे दाखवू. पहिल्यासाठी,
एखाद्या $n \ge 0$ करिता $\varphi_n \notin \Gamma^*$ आणि $\lnot \varphi_n
\notin \Gamma^*$ असे समजा. $\varphi_n \notin \Gamma^*$ असल्यास $\Gamma_n
\cup \{\varphi_n \}$ हा $\Delta_n$ पासून विलगनीय असतो; म्हणून असे
$\chi \in \Lang{L}'_0$ असते की पुढील दोन्ही विधाने खरी असतात:
\begin{align*}
  \Gamma^* & \Entails \varphi_n \lif \chi, &
  \Delta^* & \Entails \lnot \chi.
\end{align*}
त्याचप्रमाणे, $\lnot \varphi_n \notin \Gamma^*$ असेल, तर असे $\chi' \in
\Lang{L}'_0$ असते की पुढील दोन्ही विधाने खरी असतात:
\begin{align*}
  \Gamma^* & \Entails \lnot \varphi_n \lif \chi', &
  \Delta^* & \Entails \lnot \chi'.
\end{align*}
विधान तर्कशास्त्रानुसार $\Gamma^* \Entails \chi \lor \chi'$ आणि
$\Delta^* \Entails \lnot (\chi \lor \chi')$; म्हणून $\chi \lor \chi'$ हे
$\Gamma^*$ आणि~$\Delta^*$ यांना विलग करते. यासारख्याच युक्तिवादाने
$\Delta^*$ हा महत्तम असल्याचे सिद्ध होते.

शेवटी, $\Gamma^* \cap \Delta^*$ हा $\Lang{L}'_0$ मध्ये महत्तम सुसंगत
आहे हे दाखवू. तो सुसंगत संचांचा छेद असल्यामुळे उघडपणे सुसंगत आहे. महत्तमता
दाखवण्यासाठी $\varphi_{S} \in \Lang{L}'_0$ असू द्या. आता $\Gamma^*$ हा
$\Lang{L'_1} \supseteq \Lang{L'_0}$ मध्ये महत्तम आहे आणि त्याचप्रमाणे
$\Delta^*$ हा $\Lang{L'_2} \supseteq \Lang{L'_0}$ मध्ये महत्तम आहे.
त्यामुळे $\varphi_{S} \in \Gamma^*$ किंवा $\lnot \varphi_{S} \in \Gamma^*$, आणि
$\varphi_{S} \in \Delta^*$ किंवा $\lnot \varphi_{S} \in \Delta^*$. जर $\varphi_{S} \in
\Gamma^*$ आणि $\lnot \varphi_{S} \in \Delta^*$ असते, तर $\varphi_{S}$ ने $\Gamma^*$
आणि~$\Delta^*$ यांना विलग केले असते; आणि जर $\lnot \varphi_{S} \in \Gamma^*$
आणि $\varphi_{S} \in \Delta^*$ असते, तर $\lnot \varphi_{S}$ ने $\Gamma^*$ आणि~$\Delta^*$
यांना विलग केले असते. म्हणून एकतर $\varphi_{S} \in \Gamma^* \cap \Delta^*$
किंवा $\lnot \varphi_{S} \in \Gamma^* \cap \Delta^*$; त्यामुळे $\Gamma^* \cap
\Delta^*$ हा महत्तम आहे.

$\Gamma^*$ हा महत्तम सुसंगत असल्यामुळे त्याला असे प्रतिमान
$\Struct{M}'_1$ असते की त्याचा प्रांत~$\Domain{M'_1}$ हा
स्थिरांकाची अर्थनिर्धारणे असलेल्या सर्व आणि केवळ $\Assign{c}{M'_1}$
या घटकांचा बनलेला असतो---अगदी संपूर्णता प्रमेयाच्या सिद्धतेप्रमाणे
(\ref{fol:com:cth:thm:completeness}). त्याचप्रमाणे, $\Delta^*$ ला
असे प्रतिमान~$\Struct{M}'_2$ असते की त्याचा प्रांत~$\Domain{M'_2}$
हा स्थिरांकाची $\Assign{c}{M'_2}$ ही अर्थनिर्धारणे घेऊन बनलेला असतो.

$\Lang{L'_1} \setminus \Lang{L'_0}$ मधील स्थिरांक, फलने
आणि विधेये यांची अर्थनिर्धारणे वगळून $\Struct{M'_1}$ पासून
$\Struct{M_1}$ मिळवू; आणि त्याचप्रमाणे $\Struct{M_2}$ मिळवू. मग
$h(\Assign{c}{M'_1}) = \Assign{c}{M'_2}$ अशी व्याख्या असलेले प्रतिचित्रण
$h \colon M_1 \to M_2$ हे $\Lang{L}'_0$ मधील समरूपण आहे, कारण दाखवल्याप्रमाणे
$\Gamma^* \cap \Delta^*$ हा $\Lang{L}'_0$ मध्ये महत्तम सुसंगत आहे. हे
पुढील कारणाने मिळते: कोणतेही $\Lang{L}'_0$-वाक्य हे एकतर $\Gamma^*$
आणि~$\Delta^*$ या दोन्हींचा घटक असते, किंवा दोन्हींपैकी कशाचाही नसते;
म्हणून $\Assign{c}{M'_1} \in \Assign{P}{M'_1}$ तेव्हा आणि केवळ तेव्हाच,
जेव्हा $\Atom{P}{c} \in \Gamma^*$, तेव्हा आणि केवळ तेव्हाच, जेव्हा
$\Atom{P}{c} \in \Delta^*$, तेव्हा आणि केवळ तेव्हाच, जेव्हा
$\Assign{c}{M'_2} \in \Assign{P}{M'_2}$. समरूपणांनी पूर्ण केलेल्या इतर
अटीही याच रीतीने सिद्ध करता येतात.

आता भाषा~$\Lang{L_1} \cup \Lang{L_2}$ साठी प्रतिमान~$\Struct{M}$
पुढीलप्रमाणे परिभाषित करू:
\begin{enumerate}
\item प्रांत~$\Domain{M}$ म्हणजे $\Domain{M_2}$, म्हणजेच सर्व
  $\Assign{c}{M'_2}$ घटकांचा संच;
\item विधेय~$P$ हे $\Lang{L_2} \setminus \Lang{L_1}$ मध्ये
  असेल, तर $\Assign{P}{M} = \Assign{P}{M'_2}$;
\item विधेय~$P$ हे $\Lang{L}_1\setminus \Lang{L}_2$ मध्ये असेल, तर
  $\Assign{P}{M} = h(\Assign{P}{M'_1})$, म्हणजे
  $\tuple{\Assign{c_1}{M'_2}, \dots, \Assign{c_n}{M'_2}} \in
  \Assign{P}{M}$ तेव्हा आणि केवळ तेव्हाच, जेव्हा
  $\tuple{\Assign{c_1}{M'_1}, \dots, \Assign{c_n}{M'_1}} \in
  \Assign{P}{M'_1}$.

\item विधेय~$P$ हे $\Lang{L}_0$ मध्ये असेल, तर $\Assign{P}{M} =
  \Assign{P}{M'_2} = h(\Assign{P}{M'_1})$.
\item $\Lang{L}_1 \cup \Lang{L}_2$ ची फलने, त्यांत
  स्थिरांक धरून, याच रीतीने हाताळली जातात.
\end{enumerate}

शेवटी, सूत्रे वरील विगमनाने असे दाखवता येते की $\Struct{M}$ हे
$\Lang{L'_1}$ च्या सर्व सूत्रे बाबत $\Struct{M'_1}$ शी आणि
$\Lang{L'_2}$ च्या सर्व सूत्रे बाबत $\Struct{M'_2}$ शी सहमत असते.
विशेषतः, $\Struct{M} \Entails \Gamma^* \cup \Delta^*$; त्यामुळे
$\Struct{M} \Entails \varphi$ आणि $\Struct{M} \Entails \lnot\psi$, तसेच
$\not\Entails \varphi \lif \psi$. याने क्रेगच्या अंतर्वेशन प्रमेयाची सिद्धता
पूर्ण होते.
\end{proof}












\section{परिभाष्यता प्रमेय}\label{mod:int:def:sec}

अंतर्वेशन प्रमेयाचा एक महत्त्वाचा उपयोग म्हणजे बेथचे परिभाष्यता प्रमेय.
$n$-स्थानी संबंध~$R$ परिभाषित करण्यासाठी आपण $R$ चा समावेश नसलेले आणि
$n$ मुक्त चले असलेले सूत्र~$\chi$ देऊ शकतो. ही $\chi$ च्या
साहाय्याने केलेली $R$ ची \emph{स्पष्ट} परिभाषा असेल. मग असेही म्हणता येते
की $n$-स्थानी विधेय~$P$ असलेल्या भाषा मधील
उपपत्ती~$\Sigma(P)$ मध्ये औपचारिक स्पष्ट परिभाषा समाविष्ट असेल (किंवा निदान
ती परिभाषा त्या उपपत्तीतून निष्पन्न होत असेल), म्हणजे
\[\Sigma(P) \Entails \lforall[x_1][\dots
  \lforall[x_n][(\Atom{P}{x_1,\dots, x_n} \liff \chi(x_1, \dots,
    x_n))]].
\]
तर ती उपपत्ती $P$ ला स्पष्टपणे परिभाषित करते. परंतु एखादा संबंध
परिभाषित करण्याच्या---म्हणजे तो पूर्णपणे निर्धारित करण्याच्या---मार्गांपैकी
स्पष्ट परिभाषा हा केवळ एक मार्ग आहे. कोणत्याही प्रतिमानात भाषा च्या उर्वरित भागाचा अर्थबोध
दिला की $P$ चा अर्थबोध निश्चित होईल, अशीही एखादी उपपत्ती असू शकते.
परिभाष्यता प्रमेय असे सांगते की उपपत्ती जेव्हा या रीतीने $P$ चा अर्थबोध
निश्चित करते---जेव्हा ती $P$ ला \emph{अंतर्निहितपणे परिभाषित करते}---तेव्हा
ती त्याला स्पष्टपणेही परिभाषित करते.

\begin{defn}
समजा, विधेय~$P$ चा समावेश नसलेली $\Lang{L}$ ही भाषा आहे. $\Lang{L}
\cup \{P\}$ मधील वाक्यांचा संच $\Sigma(P)$ हा $P$ ला
\emph{स्पष्टपणे परिभाषित करतो}, जर आणि केवळ जर $\Lang{L}$ मधील असे
सूत्र~$\chi(x_1, \dots, x_n)$ असेल की
\[\Sigma(P) \Entails \lforall[x_1][\dots
  \lforall[x_n][(\Atom{P}{x_1,\dots, x_n} \liff \chi(x_1, \dots,
    x_n))]].
\]
\end{defn}

\begin{defn}
समजा, विधेये~$P$ आणि~$P'$ यांचा समावेश नसलेली
$\Lang{L}$ ही भाषा आहे. $\Lang{L} \cup \{P\}$ मधील
वाक्यांचा संच $\Sigma(P)$ हा $P$ ला \emph{अंतर्निहितपणे परिभाषित
करतो}, जर आणि केवळ जर
\[\Sigma(P) \cup \Sigma(P') \Entails \lforall[x_1][\dots
  \lforall[x_n][(\Atom{P}{x_1,\dots, x_n} \liff \Atom{P'}{x_1,\dots,
      x_n})]],
\]
जेथे $\Sigma(P')$ हा $\Sigma(P)$ मध्ये सर्वत्र $P$ च्या जागी $P'$ घालून
मिळणारा संच आहे.
\end{defn}

दुसऱ्या शब्दांत, कोणतेही प्रतिमान $\Struct{M}$ आणि $R, R' \subseteq
\Domain{M}^n$ यांच्यासाठी, $\Expan{M}{R} \Entails \Sigma(P)$ आणि
$\Expan{M}{R'} \Entails \Sigma(P')$ हे दोन्ही खरे असतील, तर $R=R'$;
येथे $\Expan{M}{R}$ ही $\Lang{L}$ च्या $\Lang{L} \cup \{P\}$ या
विस्तारासाठीची अशी संरचना~$\Struct{M'}$ आहे की
$\Assign{P}{M'} = R$, आणि $\Expan{M}{R'}$ साठीही तसेच.

\begin{thm}[बेथचे परिभाष्यता प्रमेय] $\Lang{L}
  \cup\{P\}$-सूत्रांचा संच $\Sigma(P)$ हा $P$ ला अंतर्निहितपणे
  परिभाषित करतो, जर आणि केवळ जर $\Sigma(P)$ हा $P$ ला स्पष्टपणे परिभाषित करतो.

\end{thm}

\begin{proof}
$\Sigma(P)$ हा $P$ ला स्पष्टपणे परिभाषित करत असेल, तर पुढील दोन्ही निष्पन्न होतात:
\begin{align*}
  \Sigma(P) & \Entails & \lforall[x_1][\dots \lforall[x_n]
    [(\Atom{P}{x_1,\dots, x_n} \liff \chi(x_1,\dots,x_n))]]\\
  \Sigma(P') & \Entails & \lforall[x_1][\dots \lforall[x_n]
    [(\Atom{P'}{x_1,\dots, x_n} \liff \chi(x_1,\dots,x_n))]]
\end{align*}
आणि निष्कर्ष मिळतो. उलट दिशेसाठी, समजा $\Sigma(P)$ हा $P$ ला
अंतर्निहितपणे परिभाषित करतो. प्रथम आपण $\Lang{L}$ मध्ये स्थिरांक
$c_1$, \dots,~$c_n$ समाविष्ट करतो. मग
\[\Sigma(P) \cup \Sigma(P') \Entails
\Atom{P}{c_1, \dots, c_n} \to  \Atom{P'}{c_1, \dots, c_n}.
\]
संहततेनुसार, असे सांत संच $\Delta_0 \subseteq \Sigma(P)$ आणि
$\Delta_1 \subseteq \Sigma(P')$ असतात की
\[\Delta_0 \cup \Delta_1 \Entails
\Atom{P}{c_1, \dots, c_n} \to \Atom{P'}{c_1, \dots, c_n}.
\]
$\theta(P)$ हा अशा सर्व वाक्ये $\varphi(P)$ चा संयोग समजू की एकतर
$\varphi(P) \in \Delta_0$ किंवा $\varphi(P') \in \Delta_1$; आणि $\theta(P')$ हा अशा
सर्व वाक्ये $\varphi(P')$ चा संयोग समजू की एकतर $\varphi(P) \in \Delta_0$
किंवा $\varphi(P') \in \Delta_1$. मग $\theta(P) \land
\theta(P') \Entails \Atom{P}{c_1, \dots, c_n} \to \Atom{P'}{c_1, \dots, c_n}$.

प्रत्येक विधेय $\Entails$ च्या एका बाजूला येईल अशा रीतीने याची
पुनर्मांडणी करता येते:
\[\theta(P) \land \Atom{P}{c_1, \dots, c_n} \Entails
\theta(P') \to \Atom{P'}{c_1, \dots, c_n}.
\]
क्रेगच्या अंतर्वेशन प्रमेयानुसार $P$ किंवा $P'$ यांचा समावेश नसलेले असे
वाक्य $\chi(c_1,\dots, c_n)$ असते की
\begin{align*}
  \theta(P) \land \Atom{P}{c_1, \dots, c_n} & \Entails \chi(c_1,\dots, c_n); \\
  \chi(c_1,\dots, c_n) & \Entails \theta(P') \to \Atom{P'}{c_1, \dots, c_n}.
\end{align*}
या दोन निष्पन्नतांपैकी पहिलीवरून $\theta(P) \Entails
\Atom{P}{c_1,\dots, c_n} \lif \chi(c_1,\dots, c_n)$ मिळते. आणि दुसरीवरून,
$\Lang{L} \cup \{P\}$-प्रतिमान $\Expan{M}{R} \Entails \varphi(P)$ असेल जर
आणि केवळ जर संबंधित $\Lang{L} \cup \{P'\}$-प्रतिमान
$\Expan{M}{R} \models \varphi(P')$ असेल, त्यामुळे
$\chi(c_1,\dots, c_n) \Entails \theta(P) \lif \Atom{P}{c_1,\dots, c_n}$ मिळते,
आणि त्यावरून
\[\theta(P) \Entails \chi(c_1,\dots,c_n) \to \Atom{P}{c_1,\dots, c_n}.
\]
दोन्ही एकत्र केल्यावर $\theta(P) \Entails \Atom{P}{c_1,\dots, c_n}
\liff \chi(c_1, \dots, c_n)$ मिळते; आणि एकस्वनिकता व सामान्यीकरण यांनुसार पुढीलही मिळते:
\[\Sigma(P) \Entails
\lforall[x_1][\dots\lforall[x_n][(\Atom{P}{x_1,\dots, x_n} \liff
    \chi(x_1,\dots, x_n))]].
\]
\end{proof}



\chapter{लिंडस्ट्रॉमचे प्रमेय}\label{mod:lin::chap}








\section{परिचय}\label{mod:lin:int:sec}

काही आणखी निर्बंध गृहीत धरल्यास संहतता आणि अधोगामी
लोव्हेनहाइम--स्कोलेम प्रमेये ज्या महत्तम तर्कशास्त्रासाठी खरी ठरतात, ते
प्रथम-क्रम तर्कशास्त्र आहे, असे सांगणारे लिंडस्ट्रॉमचे लक्षणचित्रण या
प्रकरणात सिद्ध करणे हे आपले उद्दिष्ट आहे
(\ref{fol:com:com:thm:compactness} आणि
\ref{fol:com:dls:thm:downward-ls}). प्रथम, हे प्रमेय ज्या
तर्कशास्त्रांच्या सामान्य वर्गाला लागू होते, त्याचे अधिक व्यापक वर्णन
आपल्याला आवश्यक आहे. आपण स्वतःला \emph{संबंधात्मक} भाषांपुरते मर्यादित
ठेवू, म्हणजे ज्यांत केवळ विधेये आणि व्यक्ती सूचित करणारी स्थिर
चिन्हे असतात, पण फलने नसतात अशा भाषा.











\section{अमूर्त तर्कशास्त्रे}\label{mod:lin:alg:sec}

\begin{defn}
\emph{अमूर्त तर्कशास्त्र} म्हणजे जोडी $
\tuple{L, \models_L}$; येथे $L$ हे प्रत्येक भाषा~$\Lang{L}$
ला $L(\Lang{L})$ हा वाक्यांचा संच नेमून देणारे फलन आहे, आणि
$\models_L$ हा भाषा~$\Lang{L}$ साठीच्या संरचना आणि
$L(\Lang{L})$ च्या घटक यांमधील संबंध आहे. विशेषतः,
$\tuple{F, \models}$ हे साधारण प्रथम-क्रम तर्कशास्त्र आहे; म्हणजे $F$ हे
भाषा~$\Lang{L}$ ला $\Lang{L}$ मधील स्थिर चिन्हांपासून बनवलेल्या प्रथम-क्रम
वाक्यांचा संच त्या भाषेला नेमून देणारे फलन आहे, आणि $\models$ हा
प्रथम-क्रम तर्कशास्त्राचा पूर्ति-संबंध आहे.
\end{defn}

लक्षात घ्या की दिलेल्या भाषा साठी आपण प्रथम-क्रम तर्कशास्त्रातील
संरचनेचीच संकल्पना अजूनही वापरत आहोत; परंतु $\Lang{L}$ मधील
मूलभूत चिन्हांपासून वाक्ये नेहमीच्या रीतीने बनवली जातात, किंवा
$\models_L$ हा संबंध प्रथम-क्रम तर्कशास्त्राप्रमाणेच प्रत्यावर्ती रीतीने
परिभाषित केला जातो, असे आपण पूर्वगृहीत धरत नाही. म्हणून, ही पूर्णतः
सामान्य व्याख्या अशा प्रकरणांचाही समावेश करण्यासाठी आहे ज्यांत
$\tuple{L,\models_L}$ मधील वाक्ये मध्ये अनंत लांबीची संधी किंवा
विकल्प, $\lexists$ आणि $\lforall$ यांखेरीज इतर संख्यापक (उदा.,
``असे अनंत अनेक~$x$ आहेत की \dots''), किंवा कदाचित अनंत लांबीचे
संख्यापक-उपसर्ग असतात. $L(\Lang{L})$ मधील ``वाक्ये'' ही
प्रथम-क्रम तर्कशास्त्राची साधारण वाक्ये असतीलच असे नाही, हे
ठळक करण्यासाठी या प्रकरणात त्यांवर व्यापणारी चले म्हणून $\varphi_{E}$,
$\varphi_{F}$,~\dots वापरतो, आणि साधारण प्रथम-क्रम सूत्रे साठी $\varphi$,
$\psi$,~\dots राखून ठेवतो.

\begin{defn}
$\Mod(L){\varphi_{E}}$ ही संज्ञा $\Setabs{\Struct{M}}{\Struct{M}
  \models_L \varphi_{E}}$ हा वर्ग दर्शवू दे. भाषा स्पष्ट सांगणे आवश्यक
असेल, तर आपण $\Mod[L](L){\varphi_{E}}$ असे लिहितो. $\Lang{L}$ साठीच्या दोन
संरचना $\Struct{M}$ आणि $\Struct{N}$ या $\tuple{L, \models_L}$
मध्ये \emph{प्राथमिक सममूल्य} आहेत, आणि हे $\Struct{M} \elemequiv[L]
\Struct{N}$ असे लिहितो, जर $L(\Lang{L})$ मधील तीच वाक्ये
दोन्हींमध्ये सत्य असतील.
\end{defn}

\begin{defn}
भाषा $\Lang{L}$ साठीचे अमूर्त तर्कशास्त्र
$\tuple{L,\models_L}$ पुढील गुणधर्म पूर्ण करत असेल, तर ते \emph{सामान्य}
आहे:
\begin{enumerate}
\item (\emph{$L$-एकस्वनिकता}) भाषा $\Lang{L}$ आणि
  $\Lang{L'}$ यांच्यासाठी, $\Lang{L} \subseteq \Lang{L'}$ असेल, तर
  $L(\Lang{L}) \subseteq L(\Lang{L'})$.
\item (\emph{विस्तार गुणधर्म}) प्रत्येक $\varphi_{E} \in L(\Lang{L})$
  साठी $\Lang{L}$ चा असा \emph{सांत} उपसंच $\Lang{L'}$ असतो की
  $\Struct{M} \models_L \varphi_{E}$ हा संबंध $\Struct{M}$ च्या
  $\Lang{L'}$-न्यूनीकृत रचनेवरच अवलंबून असतो; म्हणजे $\Struct{M}$ आणि
  $\Struct{N}$ यांची $\Lang{L'}$ वरील न्यूनीकृत रचना एकच असेल, तर
  $\Struct{M} \models_L \varphi_{E}$ तेव्हा आणि केवळ तेव्हाच, जेव्हा
  $\Struct{N} \models_L \varphi_{E}$.
\item (\emph{समरूपता गुणधर्म}) $\Struct{M} \models_L \varphi_{E}$ आणि
  $\Struct{M} \simeq \Struct{N}$ असेल, तर $\Struct{N} \models_L \varphi_{E}$
  हेही खरे असते.
\item (\emph{पुनर्नामकरण गुणधर्म}) पुनर्नामकरणाखाली $\models_L$ संबंध
  टिकून राहतो: भाषा $\Lang{L}'$ ही $\Lang{L}$ मधील प्रत्येक
  चिन्ह $P$ च्या जागी त्याच स्थानसंख्येचे चिन्ह $P'$ आणि प्रत्येक स्थिर
  $c$ च्या जागी वेगळे स्थिर $c'$ ठेवून मिळत असेल, तर प्रत्येक
  संरचना~$\Struct{M}$ आणि वाक्य~$\varphi_{E}$ साठी,
  $\Struct{M} \models_L \varphi_{E}$ तेव्हा आणि केवळ तेव्हाच, जेव्हा
  $\Struct{M}' \models_L \varphi_{E}'$; येथे $\Struct{M}'$ ही $\Lang{L}'$-संरचना
  असून ती $\Lang{L}$ मधील मूळ रचनेशी अनुरूप आहे, आणि $\varphi_{E}' \in
  L(\Lang{L}')$.

\item (\emph{बूलीय गुणधर्म}) अमूर्त तर्कशास्त्र $\tuple{L,
  \models_L}$ बूलीय संयोजकांखाली पुढील अर्थाने संवृत असते: प्रत्येक
  $\varphi_{E} \in L(\Lang{L})$ साठी असा $\varphi_{F} \in L(\Lang{L})$ असतो की
  $\Struct{M} \models_L \varphi_{F}$ तेव्हा आणि केवळ तेव्हाच, जेव्हा
  $\Struct{M} \not\models_L \varphi_{E}$; आणि प्रत्येक $\varphi_{E}$ आणि $\varphi_{F}$ साठी असा
  $\varphi_{G}$ असतो की $\Mod(L){\varphi_{G}} = \Mod(L){\varphi_{E}} \cap
  \Mod(L){\varphi_{F}}$. आण्विक सूत्रे आणि इतर संयोजकांसाठीही अशाच अटी
  असतात.
\item (\emph{संख्यापक गुणधर्म}) $\Lang{L}$ मधील प्रत्येक स्थिर $c$
  आणि $\varphi_{E} \in L(\Lang{L})$ साठी असा $\varphi_{F} \in L(\Lang{L})$ असतो की
  \[  \Mod[L'](L){\varphi_{F}} = \Setabs{\Struct{M}}{\Expan{M}{a} \in
  \Mod[L](L){\varphi_{E}} \text{ काही } a \in \Domain{M}},
  \]
  येथे $\Lang{L'} = \Lang{L} \setminus \{c\}$ आणि
  $\Expan{M}{a}$ ही $c$ ला $a$ नेमून देणारी $\Struct{M}$ ची
  $\Lang{L}$ वरील विस्तारित रचना आहे.
\item (\emph{सापेक्षीकरण गुणधर्म}) $\varphi_{E} \in L(\Lang{L})$ हे
  वाक्य आणि $R$, $c_1$, \dots, $c_n$ ही $\Lang{L}$ मध्ये
  नसलेली चिन्हे दिली असता, असे वाक्य
  $\varphi_{F} \in L(\Lang{L} \cup \{R,c_1,\ldots,c_n\})$ असते की त्याला $\varphi_{E}$ चे
  $\Atom{R}{x, c_1, \dots c_n}$ पर्यंतचे \emph{सापेक्षीकरण} म्हणतात,
  आणि प्रत्येक संरचना~$\Struct{M}$ साठी:
  \[  \Expan{M}{X, b_1, \ldots, b_n} \models_L \varphi_{F} \text{ तेव्हा आणि
    केवळ तेव्हाच, जेव्हा } \Struct{N} \models_L \varphi_{E},
  \]
  येथे $\Struct{N}$ ही $\Struct{M}$ ची अशी उपरचना आहे की तिचा
  प्रांत $\Domain{N} = \Setabs{a\in \Domain{M}}{\Assign{R}{M}(a,
  b_1, \dots, b_n)}$ (पहा \ref{mod:bas:sub:rem:substructure}), आणि
  $\Expan{M}{X, b_1, \ldots, b_n}$ ही $R$, $c_1$, \dots, $c_n$ यांना
  अनुक्रमे $X$, $b_1,$ \dots, $b_n$ अशी अर्थनिर्धारणे देणारी
  $\Struct{M}$ ची विस्तारित रचना आहे (येथे $X \subseteq M^{n+1}$).
\end{enumerate}
\end{defn}

\begin{defn}
दोन अमूर्त तर्कशास्त्रे $\tuple{L_1, \models_{L_1}}$ आणि
$\tuple{L_2, \models_{L_2}}$ दिली असता, प्रत्येक भाषा
$\Lang{L}$ आणि वाक्य $\varphi_{E} \in L_1(\Lang{L})$ साठी असे
वाक्य $\varphi_{F} \in L_2(\Lang{L})$ असेल की $\Mod[L](L_1){\varphi_{E}} =
\Mod[L](L_2){\varphi_{F}}$, तर दुसरे तर्कशास्त्र पहिल्याइतके \emph{किमान
  तितकेच अभिव्यक्तिक्षम} आहे, असे म्हणतो आणि $\tuple{L_1, \models_{L_1}}
\leq \tuple{L_2, \models_{L_2}}$ असे लिहितो. $\tuple{L_1,
  \models_{L_1}} \leq \tuple{L_2, \models_{L_2}}$ आणि $\tuple{L_2,
  \models_{L_2}} \leq \tuple{L_1, \models_{L_1}}$ ही दोन्ही विधाने
खरी असतील, तर $\tuple{L_1, \models_{L_1}}$ आणि $\tuple{L_2,
  \models_{L_2}}$ ही तर्कशास्त्रे \emph{सममूल्य} आहेत.
\end{defn}

\begin{rem}
  प्रथम-क्रम तर्कशास्त्र, म्हणजे अमूर्त तर्कशास्त्र $\tuple{F,
  \models}$, हे सामान्य आहे. वस्तुतः, वरील गुणधर्म प्रथम-क्रम
  तर्कशास्त्रासाठी बहुतांशी सरळ आहेत. विस्तार गुणधर्म हा
  विस्तारात्मकतेवर येऊन ठेपतो, आणि वाक्य $\varphi_{E}$ चे
  $\Atom{R}{x, c_1, \dots, c_n}$ पर्यंतचे सापेक्षीकरण प्रत्येक
  उपसूत्र $\lforall[x][\varphi_{F}]$ च्या जागी
  $\lforall[x][(\Atom{R}{x, c_1, \dots, c_n} \lif \ \varphi_{F})]$ ठेवून
  मिळते, एवढेच आपण नोंदवतो. शिवाय, $\tuple{L, \models_L}$ सामान्य
  असेल, तर $\tuple{F, \models} \leq \tuple{L, \models_{L}}$; हे
  प्रथम-क्रम सूत्रे वरील विगमनाने दाखवता येते. म्हणून कोणतीही
  सामान्यता न गमावता प्रत्येक प्रथम-क्रम वाक्य प्रत्येक सामान्य
  तर्कशास्त्रात आहे असे आपण मानू शकतो.
\end{rem}












\section{संहतता आणि लोव्हेनहाइम--स्कोलेम गुणधर्म}\label{mod:lin:lsp:sec}

आता आपण संहतता आणि लोव्हेनहाइम--स्कोलेम यांचे अमूर्त तर्कशास्त्रांसाठीचे
स्वाभाविक विस्तार देऊ.

\begin{defn}
अमूर्त तर्कशास्त्र $\tuple{L, \models_L}$ ला \emph{संहतता गुणधर्म}
असतो, जर $L(\Lang{L})$-वाक्यांचा प्रत्येक संच $\Gamma$ असा
असेल की प्रत्येक सांत $\Gamma_0 \subseteq \Gamma$ पूर्ततायोग्य असताना
तो संच पूर्ततायोग्य असतो.
\end{defn}

\begin{defn}
प्रत्येक पूर्ततायोग्य $\Gamma$ ला गणनीय प्रतिमान असेल, तर
$\tuple{L, \models_L}$ ला \emph{अधोगामी लोव्हेनहाइम--स्कोलेम गुणधर्म}
असतो.
\end{defn}


\ref{mod:bas:pis:defn:partialisom} मधील आंशिक समरूपणाची संकल्पना निव्वळ
``बैजिक'' आहे (म्हणजे भाषेतील वाक्यांचा संदर्भ न घेता, केवळ
संरचनेच्या भाषा~$\Lang{L}$ ने दिलेल्या स्थिरांच्या
संदर्भात ती मांडलेली आहे), म्हणून ती अमूर्त तर्कशास्त्रांनाही लागू होते.
प्रथम-क्रम तर्कशास्त्राच्या बाबतीत, दोन संरचना आंशिकरीत्या समरूपी
असतील, तर त्या प्राथमिक सममूल्य असतात, हे आपल्याला
\ref{mod:bas:pis:thm:p-isom2} वरून माहीत आहे. मनमाना $\varphi_{E} \in
L(\Lang{L})$ साठी सूत्रे वरील विगमन उपलब्ध असेलच असे नसल्यामुळे ती
सिद्धता अमूर्त तर्कशास्त्रांना लागू होत नाही; तरीही लोव्हेनहाइम--स्कोलेम
गुणधर्म असेल, तर प्रमेय खरे असते.

\begin{thm}
\label{mod:lin:lsp:thm:abstract-p-isom}
$\tuple{L, \models_L}$ हे लोव्हेनहाइम--स्कोलेम गुणधर्म असलेले सामान्य
तर्कशास्त्र आहे असे समजा. मग कोणत्याही दोन आंशिकरीत्या समरूपी
संरचना या $\tuple{L, \models_L}$ मध्ये प्राथमिक सममूल्य असतात.
\end{thm}

\begin{proof}
$\Struct{M} \simeq_p \Struct{N}$ असे समजा; पण एखाद्या $\varphi_{E}$ साठी
$\Struct{M} \models_L \varphi_{E}$ आणि $\Struct{N} \not\models_L \varphi_{E}$ असेही
समजा. समरूपता गुणधर्मामुळे $\Domain{M}$ आणि $\Domain{N}$ हे परस्पर
वियुक्त आहेत असे आपण मानू शकतो; आणि विस्तार गुणधर्मामुळे सांत
भाषा~$\Lang{L}$ साठी $\varphi_{E} \in L(\Lang{L})$ आहे असे मानू शकतो.
$\Struct{M}$ आणि $\Struct{N}$ यांमधील आंशिक समरूपणांचा संच
$\mathcal{I}$ असू द्या; तसेच कोणतीही सामान्यता न गमावता, $p \in
\mathcal{I}$ आणि $q \subseteq p$ असेल, तर $q \in \mathcal{I}$ हेही
मानू.

$\Domain{M}^{<\omega}$ हा $\Domain{M}$ च्या घटक च्या सांत
क्रमिकांचा संच आहे. $S$ हा $\Domain{M}^{<\omega}$ वरील जोडणी दर्शवणारा
त्रिस्थानी संबंध असू द्या; म्हणजे $\mathbf{a}, \mathbf{b},
\mathbf{c} \in \Domain{M}^{<\omega}$ असतील, तर $S(\mathbf{a},
\mathbf{b}, \mathbf{c})$ तेव्हा आणि केवळ तेव्हाच खरे असते, जेव्हा
$\mathbf{c}$ ही $\mathbf{a}$ आणि $\mathbf{b}$ यांची जोडणी असते. तसेच
$T$ हा असा त्रिस्थानी संबंध असू द्या की $b \in M$ आणि
$\mathbf{a}, \mathbf{c} \in \Domain{M}^{<\omega}$ यांच्यासाठी
$T(\mathbf{a}, b, \mathbf{c})$ तेव्हा आणि केवळ तेव्हाच खरे असते,
जेव्हा $\mathbf{a} = a_1, \dots a_n$ आणि $\mathbf{c} = a_1, \dots
a_n, b$. नवीन 3-स्थानी विधेये $P$ आणि $Q$ निवडा आणि विश्व
$\Domain{M} \cup \Domain{M}^{<\omega}$ असलेली संरचना
$\Struct{M}^*$ बनवा; तिच्यात $\Struct{M}$ ही उपरचना असते आणि $P$ व
$Q$ ची अर्थनिर्धारणे अनुक्रमे जोडणी-संबंध $S$ व $T$ असतात (म्हणून
$\Struct{M}^*$ ही भाषा $\Lang{L} \cup \{ P, Q\}$ मधील आहे).

$\Domain{N}^{<\omega}$, $S'$, $T'$, $P'$, $Q'$ आणि $\Struct{N}^*$
यांची व्याख्या याच रीतीने करा. गृहीतकानुसार $\Struct{M} \simeq_p
\Struct{N}$ असल्यामुळे $\Domain{M}^{<\omega}$ आणि
$\Domain{N}^{<\omega}$ यांमध्ये असा संबंध $I$ असतो की
$I(\mathbf{a}, \mathbf{b})$ तेव्हा आणि केवळ तेव्हाच खरे असते, जेव्हा
$\mathbf{a}$ आणि $\mathbf{b}$ समरूपी असतात आणि
\ref{mod:bas:pis:defn:partialisom} मधील पुढे-मागे अट पूर्ण करतात. आता
$\Struct{A}$ ही अशी संरचना असू द्या की तिचा प्रांत हा
$\Struct{M}^*$ आणि $\Struct{N}^*$ यांच्या प्रांताचा संयोग आहे,
$\Struct{M}^*$ आणि $\Struct{N}^*$ या तिच्या उपसंरचना आहेत,
आणि तिच्या भाषा मध्ये संबंध~$I$ ने अर्थनिर्धारित केलेले एक
अतिरिक्त द्विपदी विधेय~$R$ व प्रांत $\Domain{M}^*$ आणि
$\Domain{N}^*$ दर्शवणारी विधेये आहेत.


\begin{figure}[h]
  \centering
  \begin{tikzpicture}[node distance=2cm, auto, thick, >=stealth']
    \draw [rounded corners] (0,0) -- (8,0) -- (8,4) -- (0,4) --  cycle;
    \draw (2,2) circle (0.5cm);
    \draw (2,2) circle (1.25cm);
    \draw (6,2) circle (0.5cm);
    \draw (6,2) circle (1.25cm);
    \path node at (0.75,3.5) {\large $\Struct{A}$};
    \path node at (2,2) {\large $\Struct{M}$};
    \path node at (6,2) {\large $\Struct{N}$};
    \path node at (3.5,1) {\large $\Struct{M}^*$};
    \path node at (7.5,1) {\large $\Struct{N}^*$};
    \node (Idom) at (2.8,2) {};
    \node (Irng) at (5.2,2) {};
    \draw[<->, bend left] (Idom) to node {\large $I$} (Irng) ;
  \end{tikzpicture}
  \caption{अंतर्गत आंशिक समरूपण असलेली संरचना~$\Struct{A}$.}
\end{figure}

निर्णायक निरीक्षण असे आहे की संरचना~$\Struct{A}$ च्या
भाषा मध्ये $\Struct{A}$ मध्ये सत्य असलेले असे
\emph{अमूर्त-तर्कशास्त्रीय} वाक्य $\theta_1$ असते की ते
$\Struct{M} \models_L \varphi_{E}$ आणि $\Struct{N} \not\models_L \varphi_{E}$ असे
सांगते (यासाठी सापेक्षीकरण गुणधर्म आवश्यक आहे); तसेच $\Struct{A}$ मध्ये
सत्य असलेले असे \emph{प्रथम-क्रम} वाक्य $\theta_2$ असते की ते
$I$ या आंशिक समरूपणाद्वारे $\Struct{M} \simeq_p \Struct{N}$ असे
सांगते. लोव्हेनहाइम--स्कोलेम गुणधर्मामुळे $\theta_1$ आणि $\theta_2$ ही दोन्ही
एका गणनीय प्रतिमान $\Struct{A}_0$ मध्ये सत्य असतात; त्यात
आंशिकरीत्या समरूपी अशा उपरचना $\Struct{M}_0$ आणि $\Struct{N}_0$ असतात,
ज्यांच्यासाठी $\Struct{M}_0 \models_L \varphi_{E}$ आणि $\Struct{N}_0
\not\models_L \varphi_{E}$. पण गणनीय आणि आंशिकरीत्या समरूपी
संरचना या \ref{mod:bas:pis:thm:p-isom1} नुसार वस्तुतः समरूपी
असतात; हे सामान्य अमूर्त तर्कशास्त्रांच्या समरूपता गुणधर्माच्या विरुद्ध
आहे.


\end{proof}












\section{लिंडस्ट्रॉमचे प्रमेय}\label{mod:lin:prf:sec}

\begin{lem}
\label{mod:lin:prf:lem:lindstrom}
$\varphi_{E} \in L(\Lang{L})$ आणि $\Lang{L}$ सांत आहे असे समजा; तसेच असा
$n \in \Nat$ आहे असे समजा की कोणत्याही दोन संरचना
$\Struct{M}$ आणि~$\Struct{N}$ साठी, $\Struct{M} \equiv_n
\Struct{N}$ आणि $\Struct{M} \models_L \varphi_{E}$ असेल, तर $\Struct{N}
\models_L \varphi_{E}$ हेही खरे असते. मग $\varphi_{E}$ हे एका प्रथम-क्रम
वाक्य शी सममूल्य असते; म्हणजे असे प्रथम-क्रम $\theta$ असते की
$\Mod(L){\varphi_{E}} = \Mod(L){\theta}$.
\end{lem}

\begin{proof}
कोणत्याही दोन $n$-सममूल्य संरचना $\Struct{M}$ आणि
$\Struct{N}$ या $\varphi_{E}$ ला नेमलेल्या सत्यतामूल्याबाबत सहमत असतील, असा $n$
घ्या. \ref{mod:bas:pis:prop:qr-finite} आठवा: तार्किक सममूल्यतेपर्यंत,
सांत भाषा मध्ये संख्यापक-प्रतांक $n$ पेक्षा अधिक नसलेली प्रथम-क्रम
वाक्ये सांत इतकीच असतात. आता प्रत्येक स्थिर
संरचना~$\Struct{M}$ साठी, $\Struct{M}$ मध्ये सत्य आणि
$\QuantRank{\varphi_{E}} \le n$ असलेल्या सर्व प्रथम-क्रम वाक्ये~$\varphi_{E}$
यांची संधी $\theta_{\Struct{M}}$ असू द्या (ही संधी सांत आहे), म्हणजे
$\Struct{N} \models \theta_{\Struct{M}}$ तेव्हा आणि केवळ तेव्हाच, जेव्हा
$\Struct{N} \equiv_n \Struct{M}$. मग $\theta = \textstyle\bigvee
\Setabs{\theta_{\Struct{M}}}{\Struct{M} \models_L \varphi_{E}}$ असे ठेवा; हा विकल्पही
(तार्किक सममूल्यतेपर्यंत) सांत आहे.

यावरून $\Mod(L){\varphi_{E}} = \Mod(L){\theta}$ हा निष्कर्ष मिळतो. वस्तुतः,
$\Struct{N} \models_L \theta$ असेल, तर एखाद्या $\Struct{M} \models_L
\varphi_{E}$ साठी $\Struct{N} \models \theta_{\Struct{M}}$ असते; त्यामुळे
लेम्माच्या गृहीतकानुसार $\Struct{N} \models_L \varphi_{E}$ हेही खरे असते.
उलट, $\Struct{N} \models_L \varphi_{E}$ असेल, तर $\theta_\Struct{N}$ हा $\theta$ मधील
एक विकल्पांश आहे; आणि $\Struct{N} \models \theta_\Struct{N}$ असल्यामुळे
$\Struct{N} \models_L \theta$ हेही खरे असते.
\end{proof}

\begin{thm}[लिंडस्ट्रॉमचे प्रमेय]
  \label{mod:lin:prf:thm:lindstrom} सामान्य अमूर्त तर्कशास्त्र
  $\tuple{L, \models_L}$ ला संहतता आणि लोव्हेनहाइम--स्कोलेम गुणधर्म
  आहेत असे समजा. मग $\tuple{L, \models_L} \le \tuple{F, \models}$
  (म्हणून $\tuple{L, \models_L}$ हे प्रथम-क्रम तर्कशास्त्राशी सममूल्य
  आहे).

\end{thm}

\begin{proof}
\ref{mod:lin:prf:lem:lindstrom} नुसार, सांत $\Lang{L}$ साठी कोणतेही $\varphi_{E}
\in L(\Lang{L})$ दिले असता असा $n \in \Nat$ असतो की कोणत्याही दोन
संरचना $\Struct{M}$ आणि~$\Struct{N}$ साठी: $\Struct{M}
\equiv_n \Struct{N}$ असेल, तर $\Struct{M}$ आणि $\Struct{N}$ या $\varphi_{E}$
बाबत सहमत असतात, हे दाखवणे पुरेसे आहे. कारण मग $\varphi_{E}$ हे प्रथम-क्रम
वाक्य शी सममूल्य असते, आणि त्यावरून $\tuple{L, \models_L} \le
\tuple{F, \models}$ मिळते. आपण सांत आणि निव्वळ संबंधात्मक
भाषा मध्ये काम करत असल्यामुळे, \ref{mod:bas:pis:thm:b-n-f}
नुसार $\Struct{M} \equiv_n \Struct{N}$ या विधानाच्या जागी अनुरूप बैजिक
विधान $I_n(\emptyset,\emptyset)$ घेता येते.

$\varphi_{E}$ दिले असता, विरोधासाठी असे समजा की प्रत्येक $n$ साठी अशा
संरचना $\Struct{M}_n$ आणि $\Struct{N}_n$ असतात की
$I_n(\emptyset, \emptyset)$; पण (समजा) $\Struct{M}_n \models_L \varphi_{E}$
आणि $\Struct{N}_n \not\models_L \varphi_{E}$. भाषेत सांत इतकीच आण्विक
वाक्ये असल्यामुळे, गरज पडल्यास $\Struct{M}_n$ चा उपअनुक्रम
घेऊन त्या सर्व त्याच आण्विक वाक्यांची पूर्ति करतात असे मानू
शकतो. हा उपअनुक्रम अनंत निवडून नव्याने क्रमांकित केल्यास प्रत्येक मूळ
निर्देशांक नव्या निर्देशांकाएवढा किंवा त्याहून मोठा असतो; म्हणून आवश्यक
सममूल्यतेची खोली कायम राहते. आता स्थिरांतील समानतेचा आकृतिबंध सर्वत्र
सारखा आहे. त्यामुळे समरूपता गुणधर्म वापरून सर्व $\Struct{M}_n$ भाषेतील
स्थिरांना त्याच वस्तू नेमतात आणि त्या वस्तूंखेरीज त्यांचे प्रांत परस्पर
विभक्त आहेत असे मानू शकतो. सर्व $\Struct{M}_n$ चा संयोग,
म्हणजे प्रत्येक $\Struct{M}_n$ उपरचना असलेली एकमेव न्यूनतम
संरचना, $\Struct{M}$ असू द्या. \ref{mod:lin:lsp:thm:abstract-p-isom}
च्या सिद्धतेप्रमाणे, $\Struct{M}^*$ ही $\Struct{M}$ ची अशी वर्धित रचना
असू द्या की तिचा प्रांत $\Domain{M} \cup
\Domain{M}^{<\omega}$ आहे आणि तिच्या विस्तारित भाषा मध्ये
जोडणी-विधेये $P$ आणि~$Q$ आहेत.




$\Struct{N}_n$, $\Struct{N}$ आणि $\Struct{N}^*$ यांची व्याख्या याच
रीतीने करा. आता $\Struct{A}$ ही अशी संरचना असू द्या की तिचा
प्रांत हा $\Struct{M}^*$ आणि $\Struct{N}^*$ यांच्या प्रांत
बरोबर नैसर्गिक संख्या~$\Nat$ आणि त्यांचा नैसर्गिक क्रम~$\le$ यांचा
समावेश करतो. तिच्या भाषा मध्ये प्रांत $\Domain{M}$,
$\Domain{N}$, $\Domain{M}^{<\omega}$ आणि $\Domain{N}^{<\omega}$
दर्शवणारी अतिरिक्त विधेये असतात; तसेच $\Struct{M}_n$ आणि
$\Struct{N}_n$ यांचे प्रांत पुढील अर्थाने संकेतबद्ध करणारी विधेये असतात:
\begin{align*}
  \Domain{M_n} & = \Setabs{a \in \Domain{M}}{R(a, n)}; &
  \Domain{N_n} & = \Setabs{a \in \Domain{N}}{S(a,n)}; \\
  \Domain{M}^{<\omega}_n & = \Setabs{a \in \Domain{M}^{<\omega}}{R(a,n)}; &
  \Domain{N}^{<\omega}_n & = \Setabs{a \in \Domain{N}^{<\omega}}{S(a,n)}.
\end{align*}
संरचना~$\Struct{A}$ मध्ये असा त्रिस्थानी संबंध $J$ ही असतो की
$J(n, \mathbf{a}, \mathbf{b})$ तेव्हा आणि केवळ तेव्हाच खरे असते,
जेव्हा $I_n(\mathbf{a}, \mathbf{b})$.


आता $R$, $S$, $J$ इत्यादींनी वाढवलेल्या भाषा~$\Lang{L}$ मध्ये
असे वाक्य~$\theta$ असते की ते $\le$ हा पहिला घटक असलेला पण शेवटचा
घटक नसलेला विविक्त रेषीय क्रम आहे, $\Struct{M}_n \models_L \varphi_{E}$,
$\Struct{N}_n \not\models_L \varphi_{E}$, आणि क्रमातील प्रत्येक $n$ साठी
$J(n, \mathbf{a}, \mathbf{b})$ तेव्हा आणि केवळ तेव्हाच खरे असते,
जेव्हा $I_n(\mathbf{a}, \mathbf{b})$, असे सांगते.


संहतता गुणधर्म वापरून, $\theta$ बरोबर एका नवीन स्थिराला प्रत्येक मानक
घटकाच्या वर ठेवणारी वाक्ये घ्या. या संचाला असे प्रतिमान $\Struct{A}^*$
मिळते की त्यातील क्रमात अमानक घटक~$n^*$ असतो. विशेषतः,
$\Struct{A}^*$ मध्ये अशा उपसंरचना $\Struct{M_{n^*}}$ आणि
$\Struct{N_{n^*}}$ असतात की $\Struct{M_{n^*}} \models_L \varphi_{E}$ आणि
$\Struct{N_{n^*}} \not\models_L \varphi_{E}$. आता $\Domain{M_{n^*}}$ मधील
आणि $\Domain{N_{n^*}}$ मधील $k$-सदस्यी क्रमितकांच्या जोड्यांचा संच
$\mathcal{I}$ पुढीलप्रमाणे परिभाषित करता येतो: $\tuple{\mathbf{a},
\mathbf{b}} \in \mathcal{I}$ तेव्हा आणि केवळ तेव्हाच, जेव्हा
$J(n^*-k, \mathbf{a}, \mathbf{b})$; येथे $k$ ही $\mathbf{a}$ आणि
$\mathbf{b}$ यांची लांबी आहे. $n^*$ हा अमानक असल्यामुळे प्रत्येक मानक
$k$ साठी $n^* - k >0$ असते, आणि $\mathcal{I}$ हा संच
$\Struct{M_{n^*}} \simeq_p \Struct{N_{n^*}}$ या वस्तुस्थितीचा साक्षी
आहे. पण \ref{mod:lin:lsp:thm:abstract-p-isom} नुसार $\Struct{M_{n^*}}$ ही
$\Struct{N_{n^*}}$ शी $L$-सममूल्य आहे; हा विरोधाभास आहे.


\end{proof}



\chapter{पुनरावर्ती फलने}\label{cmp:rec::chap}
\begin{editorial}
  या जेरेमी अविगॅड यांच्या पुनरावर्ती फलनांवरील टिपा आहेत; रिचर्ड झॅक
  यांनी त्यांत सुधारणा करून त्यांचा विस्तार केला आहे. या प्रकरणात
  खरोखरच काही सराव आहेत आणि विन्यासमीमांसेच्या अंकगणितीकरणावरील
  चर्चेचा आधार देण्यासाठी ते स्वतंत्रपणे समाविष्ट करता येते.
\end{editorial}








\section{परिचय}\label{cmp:rec:int:sec}

संगणनक्षमतेची गणितीय उपपत्ती विकसित करण्यासाठी सर्वप्रथम
संगणनक्षमतेचे एक \emph{प्रतिमान} विकसित करावे लागते. आज आपण
संगणनक्षमता ही संगणक करतात त्या प्रकारची गोष्ट आहे असे मानतो, आणि संगणक
चिन्हांच्या साहाय्याने काम करतात. परंतु संगणनक्षमतेच्या उपपत्तींच्या
विकासाच्या आरंभी संगणनाचे आदर्श उदाहरण \emph{संख्यात्मक} संगणन होते.
गणितज्ञांना संख्या-सैद्धांतिक फलनांमध्ये, म्हणजे संगणित करता येणाऱ्या
$f\colon \Nat^n \to \Nat$ या फलनांमध्ये, नेहमीच रस होता. म्हणून
संगणनक्षमतेच्या उपपत्तीच्या आरंभी अशाच फलनांचा अभ्यास झाला, यात आश्चर्य
नाही. बेरीज, गुणाकार, (नैसर्गिक संख्यांची) घातांक क्रिया यांसारख्या
संगणनक्षम संख्यात्मक फलनांच्या सर्वाधिक परिचित उदाहरणांत एक रोचक समान
वैशिष्ट्य आहे: त्यांची \emph{पुनरावर्ती रीतीने} व्याख्या करता येते.
त्यामुळे पुनरावर्ती व्याख्यांच्या आधारे \emph{संगणनक्षम फलनाची}
सर्वसाधारण व्याख्या करण्याचा प्रयत्न अगदी स्वाभाविक आहे.
संख्या-सैद्धांतिक फलनांची पुनरावर्ती रीतीने व्याख्या करण्याच्या अनेक
संभाव्य मार्गांपैकी येथे एका विशेषतः सोप्या व्याख्या-आकृतिबंधाला
मध्यवर्ती स्थान मिळते: तथाकथित \emph{आदिम पुनरावर्तन}.

संगणनक्षम फलनांबरोबरच संगणनक्षम संच आणि संबंध यांतही आपल्याला रस असू
शकतो. दिलेली संख्या संचाचा घटक आहे की नाही याचे उत्तर आपण
संगणित करू शकत असू, तर तो संच संगणनक्षम असतो; आणि दिलेला क्रमित
घटकसमूह $\tuple{n_1, \dots, n_k}$ हा संबंधाचा घटक आहे
की नाही हे आपण संगणित करू शकत असू तेव्हा आणि केवळ तेव्हाच तो संबंध
संगणनक्षम असतो. संचाचे किंवा संबंधाचे \emph{जातिबोधक फल} विचारात
घेतल्यास संगणनक्षम संच व संबंध यांची चर्चा संगणनक्षम फलनांच्या चर्चेत
समाविष्ट करता येते. त्यामुळे आदिम पुनरावर्ती संबंधांचीही व्याख्या करता
येते; उदाहरणार्थ, ``$m$ ला $n$ ने निःशेष भाग जातो'' हा संबंध आदिम
पुनरावर्ती आहे.

तथापि, आदिम पुनरावर्ती फलने---केवळ आदिम पुनरावर्तन वापरून परिभाषित करता
येणारी फलने---हीच एकमेव संगणनक्षम संख्या-सैद्धांतिक फलने नाहीत. आदिम
पुनरावर्तनाची अनेक सामान्यीकरणे विचारात घेण्यात आली आहेत; परंतु फलने
संगणित करण्याची सर्वाधिक शक्तिशाली आणि व्यापक मान्यता असलेली अतिरिक्त
पद्धत म्हणजे अपरिबद्ध शोध. यातून \emph{आंशिक पुनरावर्ती फलनांची}
व्याख्या मिळते, आणि त्यासंबंधित व्याख्येतून \emph{सामान्य पुनरावर्ती
  फलनांची} व्याख्या मिळते. सामान्य पुनरावर्ती फलने संगणनक्षम आणि
सर्वत्र परिभाषित असतात; आणि ही व्याख्या आंशिक पुनरावर्ती फलनांपैकी
सर्वत्र परिभाषित असणारी नेमकी फलने लक्षणित करते. पुनरावर्ती फलने
संगणनाच्या प्रत्येक इतर प्रतिमानाचे (ट्यूरिंग यंत्रे, लॅम्डा कलन इत्यादींचे)
अनुकरण करू शकतात आणि म्हणून ती संगणनाच्या अनेक मान्य प्रतिमानांपैकी एक
प्रतिमान ठरतात.












\section{आदिम पुनरावर्तन}\label{cmp:rec:pre:sec}

नैसर्गिक संख्यांचे एक वैशिष्ट्य असे आहे की, उत्तरवर्ती क्रिया~$+1$
सांत वेळा लागू करून प्रत्येक नैसर्गिक संख्या~$0$ पासून गाठता येते---कोणतीही
नैसर्गिक संख्या एकतर~$0$ असते किंवा~$0$ च्या उत्तरवर्तीचा \dots{}
उत्तरवर्ती असते. या तथ्याचा उपयोग करणारे फलन~$h\colon\Nat \to \Nat$
विनिर्दिष्ट करण्याचा एक मार्ग असा: (a)~आदान~$0$ साठी $h$ चे मूल्य काय
आहे ते विनिर्दिष्ट करणे, आणि (b)~$h(x)$ चे मूल्य दिले असता $h(x+1)$ चे
मूल्य कसे संगणित करायचे तेही विनिर्दिष्ट करणे. कारण (a) आपल्याला
$h(0)$ काय आहे ते थेट सांगते, त्यामुळे~$0$ साठी $h$ परिभाषित होते.
आता (b) ने दिलेली सूचना $x=0$ साठी वापरून आपण~$h(0)$ वरून
$h(1) = h(0+1)$ संगणित करू शकतो. तीच सूचना $x=1$ साठी वापरून
आपण~$h(1)$ वरून $h(2) = h(1+1)$ संगणित करतो, आणि असेच पुढे.
प्रत्येक नैसर्गिक संख्या~$x$ साठी, आपण अखेरीस $h(x)$ वरून $h(x+1)$
परिभाषित करतो त्या पायरीपर्यंत पोहोचतो; आधारमूल्यापासून अशी पावले
घेतल्यामुळे सर्व $x \in \Nat$ साठी $h(x)$ परिभाषित होते.

उदाहरणार्थ, पुढील दोन समीकरणांनी आपण $h\colon \Nat \to \Nat$ विनिर्दिष्ट
करतो असे समजा:
\begin{align*}
h(0) & =  1\\
h(x+1) & =  2 \cdot h(x)
\end{align*}
गुणाकार कसा करायचा हे आपल्याला आधीच माहीत असेल, तर वरील (a) आणि~(b)
साठी आवश्यक माहिती ही समीकरणे देतात. दुसरे समीकरण क्रमाक्रमाने लागू
केल्यावर आपल्याला पुढील गोष्टी मिळतात:
\begin{align*}
  h(1) & = 2\cdot h(0) = 2,\\
  h(2) & = 2\cdot h(1) = 2\cdot 2,\\
  h(3) & = 2 \cdot h(2) = 2\cdot 2 \cdot 2,\\
  & \vdots
\end{align*}
आपण विनिर्दिष्ट केलेले फलन~$h$ हे $h(x) = 2^x$ आहे, हे दिसते.

नैसर्गिक संख्यांच्या वैशिष्ट्यपूर्ण गुणधर्मामुळे हे दोन निकष पूर्ण करणारे
एकमेव फलन~$h$ असते. अशा समीकरणांच्या जोडीला फलन~$h$ ची
\emph{आदिम पुनरावर्तनाने केलेली व्याख्या} म्हणतात. तिला तसे म्हणण्याचे
कारण म्हणजे आपण~$h$ ची व्याख्या ``पुनरावर्ती रीतीने'' करतो, म्हणजे
व्याख्येत, विशेषतः दुसऱ्या समीकरणाच्या उजव्या बाजूला, $h$ स्वतःच येते.
ती ``आदिम'' असते, कारण $h(x+1)$ परिभाषित करताना आपण केवळ~$h(x)$ चे
मूल्य, म्हणजे अगदी आधीचे मूल्य, वापरतो. फलनाची~$\Nat$ वर पुनरावर्ती
रीतीने व्याख्या करण्याचा हा सर्वांत सोपा मार्ग आहे.

बेरीज आणि गुणाकार यांसारखी आणखी मूलभूत फलनेही आपण आदिम पुनरावर्तनाने
परिभाषित करू शकतो. मात्र या प्रकरणांत संबंधित फलने $2$-स्थानी असतात.
आपण आदानस्थानांपैकी एक स्थिर ठेवतो आणि दुसरे पुनरावर्तनासाठी वापरतो.
उदाहरणार्थ, $\Add(x, y)$ परिभाषित करण्यासाठी आपण~$x$ स्थिर ठेवून आधी
$y=0$ साठी मूल्य परिभाषित करू शकतो आणि नंतर~$y$ च्या साहाय्याने $y+1$
साठीचे मूल्य परिभाषित करू शकतो. $x$ स्थिर असल्यामुळे ते परिभाषक
समीकरणांच्या डाव्या आणि उजव्या दोन्ही बाजूंना येईल.
\begin{align*}
\Add(x,0) & =  x\\
\Add(x,y+1) & =  \Add(x,y)+1
\end{align*}
ही समीकरणे सर्व $x$ \emph{आणि}~$y$ साठी $\Add$ चे मूल्य विनिर्दिष्ट
करतात. उदाहरणार्थ, $\Add(2,3)$ शोधण्यासाठी आपण $x = 2$ साठी परिभाषक
समीकरणे लागू करतो: पहिल्या समीकरणाने $\Add(2,0) = 2$ शोधतो, आणि नंतर
दुसरे समीकरण क्रमाक्रमाने वापरून $\Add(2,1) = 2 + 1 = 3$,
$\Add(2, 2) = 3 + 1 = 4$, $\Add(2, 3) = 4 + 1 = 5$ शोधतो.

$\Add$ च्या व्याख्येत आपण दुसऱ्या समीकरणाच्या उजव्या बाजूला $+$ वापरले,
पण केवळ~$1$ ची भर घालण्यासाठी. दुसऱ्या शब्दांत, आपण उत्तरवर्ती फलन
$\Succ(z) = z+1$ वापरले आणि $\Add(x,y+1)$ परिभाषित करण्यासाठी ते मागील
मूल्य~$\Add(x,y)$ ला लागू केले. म्हणून मागील मूल्याला लागू केलेल्या एका
फलनाच्या रूपात पुनरावर्ती व्याख्येचा विचार करता येतो. परंतु त्या फलनाला
केवळ मागील मूल्यावरच नव्हे, तर $x$ आणि~$y$ यांवरही अवलंबून राहण्याची
मुभा दिल्याने काही बिघडत नाही---आणि कधीकधी ते आवश्यक असते. पुढील
समीकरणे पाहा:
\begin{align*}
  \Mult(x,0) & =  0 \\
  \Mult(x,y+1) & =  \Add(\Mult(x,y),x)
\end{align*}
ही, आधीचे मूल्य $\Mult(x,y)$ आणि पहिले आदान~$x$ या दोहोंना फलन
$\Add$ लागू करून मिळणारी, फलन $\Mult$ ची आदिम पुनरावर्ती व्याख्या आहे.
ती सर्व आदाने $x$ आणि~$y$ यांच्यासाठी फलन~$\Mult(x,y)$ परिभाषितही करते.
उदाहरणार्थ, $\Mult(2,3)$ हे $\Mult(2,0)$, $\Mult(2,1)$,
$\Mult(2,2)$ आणि~$\Mult(2,3)$ क्रमाक्रमाने संगणित करून निर्धारित होते:
\begin{align*}
  \Mult(2,0) & = 0\\
  \Mult(2,1) & = \Mult(2,0+1) =
  \Add(\Mult(2,0), 2) = \Add(0, 2) = 2\\
  \Mult(2,2) & = \Mult(2,1+1) =
  \Add(\Mult(2,1), 2) = \Add(2, 2) = 4\\
  \Mult(2,3) & = \Mult(2,2+1) =
  \Add(\Mult(2,2), 2) = \Add(4, 2) = 6
\end{align*}

मग सर्वसाधारण आकृतिबंध असा आहे: फलन~$h(x_0, \dots, x_{k-1}, y)$ ची
आदिम पुनरावर्ती व्याख्या देण्यासाठी आपण दोन समीकरणे देतो. पहिले समीकरण
फलन~$h$ चा उल्लेख न करता $h(x_0, \dots, x_{k-1}, 0)$ चे मूल्य
परिभाषित करते. दुसरे समीकरण $h(x_0, \dots, x_{k-1}, y)$, इतर आदाने
$x_0$, \dots,~$x_{k-1}$ आणि~$y$ यांच्या साहाय्याने $h(x_0,
\dots, x_{k-1}, y+1)$ चे मूल्य परिभाषित करते. त्या दुसऱ्या समीकरणात
केवळ~$h$ चे अगदी आधीचे मूल्य वापरता येते. या दोन समीकरणांच्या उजव्या
बाजूंनी दिलेल्या क्रियाच स्वतः फलने $f$ आणि~$g$ आहेत असे मानल्यास,
आदिम पुनरावर्तनाने नवीन फलन~$h$ परिभाषित करण्याचा सर्वसाधारण आकृतिबंध
असा आहे:
\begin{align*}
  h(x_0, \dots, x_{k-1}, 0) & = f(x_0, \dots, x_{k-1})\\
  h(x_0, \dots, x_{k-1}, y+1) & =
  g(x_0, \dots, x_{k-1}, y, h(x_0, \dots, x_{k-1}, y))
\end{align*}
$\Add$ च्या बाबतीत $k=1$ आणि $f(x_0) = x_0$ (एकरूपता फलन) असते,
तसेच $g(x_0, y, z) = z + 1$ (आपल्या तिसऱ्या आदानाचा उत्तरवर्ती
परत करणारे $3$-स्थानी फलन) असते:
\begin{align*}
  \Add(x_0, 0) & = f(x_0) = x_0\\
  \Add(x_0, y+1) & = g(x_0, y, \Add(x_0, y)) =
  \Succ(\Add(x_0, y))
\end{align*}
$\Mult$ च्या बाबतीत $f(x_0) = 0$ (नेहमी~$0$ परत करणारे स्थिर फलन)
आणि $g(x_0, y, z) = \Add(z,x_0)$ (आपल्या शेवटच्या व पहिल्या आदानांची
बेरीज परत करणारे $3$-स्थानी फलन) असते:
\begin{align*}
  \Mult(x_0,0) & =  f(x_0) = 0 \\
  \Mult(x_0,y+1) & =  g(x_0, y, \Mult(x_0,y)) =
  \Add(\Mult(x_0,y), x_0)
\end{align*}











\section{संयोजन}\label{cmp:rec:com:sec}

$f$ आणि $g$ ही नैसर्गिक संख्यांची दोन एकस्थानी फलने असतील, तर त्यांचे
संयोजन करता येते: $h(x) = f(g(x))$. मग नवीन फलन~$h(x)$ हे $f$ आणि~$g$
या फलनांपासून \emph{संयोजनाने} परिभाषित होते. याचे एकाहून अधिक आदाने
असलेल्या फलनांसाठी सामान्यीकरण करायचे आहे.

ते करण्याचा एक मार्ग असा: $f$ हे $k$-स्थानी फलन आहे आणि $g_0$,
\dots, $g_{k-1}$ ही सर्व $n$-स्थानी असलेली $k$ फलने आहेत असे समजा.
मग पुढीलप्रमाणे नवीन $n$-स्थानी फलन $h$ परिभाषित करता येते:
\[h(x_0, \dots, x_{n-1}) =
f(g_0(x_0, \dots, x_{n-1}), \dots, g_{k-1}(x_0, \dots, x_{n-1}))
\]
$f$ आणि सर्व $g_i$ संगणनक्षम असतील, तर $h$ देखील संगणनक्षम असते:
$h(x_0, \dots, x_{n-1})$ संगणित करण्यासाठी प्रथम प्रत्येक $i = 0$,
\dots,~$k-1$ साठी $y_i = g_i(x_0, \dots, x_{n-1})$ ही मूल्ये संगणित
करा. नंतर ही मूल्ये $f$ ला आदाने म्हणून देऊन
$h(x_0, \dots, x_{n-1}) = f(y_0, \dots, y_{k-1})$ संगणित करा.

आधीपासून उपलब्ध असलेल्या काही फलनांचा उपयोग करून नवीन फलन संगणित करताना
जे घडते त्याचे हे लक्षणचित्रण अनावश्यकपणे बंधनकारक वाटू शकते. एक कारण
म्हणजे, कधीकधी आपण फलनाची सर्व आदाने वापरत नाही; उदाहरणार्थ,~$\Add$
च्या आदिम पुनरावर्ती व्याख्येत वापरण्यासाठी $g(x, y, z) = \Succ(z)$
परिभाषित केले तेव्हा तसेच झाले. पुढील फलने वापरण्याची आपल्याला मुभा आहे
असे समजा:
\[\Proj{n}{i}(x_0, \dots, x_{n-1}) = x_i
\]
फलनांना~$\Proj{k}{i}$ \emph{प्रक्षेपण} फलने म्हणतात:
$\Proj{n}{i}$ हे $n$-स्थानी फलन आहे. मग $g$ ची व्याख्या अशी देता येते:
\[g(x, y, z) = \Succ(\Proj{3}{2}(x, y, z)).
\]
येथे $f$ ची भूमिका $1$-स्थानी फलन $\Succ$ बजावते, म्हणून $k=1$.
तसेच $g_0$ ची भूमिका बजावणारे एक $3$-स्थानी फलन $\Proj{3}{2}$ आहे.
परिणाम म्हणजे तिसऱ्या आदानाचा उत्तरवर्ती परत करणारे $3$-स्थानी फलन.

प्रक्षेपण फलनांमुळे आदानांचा क्रम बदलून किंवा काही आदाने एकरूप करूनही
नवीन फलने परिभाषित करता येतात. उदाहरणार्थ, फलन $h(x) = \Add(x, x)$ ची
व्याख्या अशी देता येते:
\[h(x_0) = \Add(\Proj{1}{0}(x_0),\Proj{1}{0}(x_0)).
\]
येथे $k=2$, $n=1$, $f(y_0,y_1)$ ची भूमिका $\Add$ बजावते, आणि
$g_0(x_0)$ व $g_1(x_0)$ या दोहोंची भूमिका~$\Proj{1}{0}(x_0)$ हे
एकस्थानी प्रक्षेपण फलन (म्हणजे एकरूपता फलन) बजावते.

$f(y_0, y_1)$ हे फलन आपल्याकडे आधीच असेल, तर फलन
$h(x_0, x_1) = f(x_1, x_0)$ ची व्याख्या अशी देता येते:
\[h(x_0, x_1) = f(\Proj{2}{1}(x_0, x_1),\Proj{2}{0}(x_0, x_1)).
\]
येथे $k=2$, $n = 2$, आणि $g_0$ व $g_1$ यांच्या भूमिका अनुक्रमे
$\Proj{2}{1}$ आणि~$\Proj{2}{0}$ बजावतात.

$g_0$, \dots,~$g_{k-1}$ या सर्वांची स्थानसंख्या~$n$ समान असणे आवश्यक
आहे, याबद्दलही शंका वाटू शकते. (लक्षात ठेवा, फलनाची \emph{स्थानसंख्या}
म्हणजे त्याच्या आदानांची संख्या; $n$-स्थानी फलनाची स्थानसंख्या~$n$
असते.) पण प्रक्षेपण फलने समाविष्ट केल्याने अपेक्षित लवचीकता मिळते.
उदाहरणार्थ, $f$ आणि~$g$ ही $3$-स्थानी फलने आहेत आणि पुढीलप्रमाणे
परिभाषित केलेले $h$ हे $2$-स्थानी फलन आहे असे समजा:
\[h(x,y) = f(x,g(x,x,y),y).
\]
प्रक्षेपण फलने वापरून~$h$ ची व्याख्या पुढीलप्रमाणे पुन्हा लिहिता येते:
\[h(x,y) = f(\Proj{2}{0}(x,y), g(\Proj{2}{0}(x,y), \Proj{2}{0}(x,y),
\Proj{2}{1}(x,y)), \Proj{2}{1}(x,y)).
\]
मग $h$ हे $f$ चे $\Proj{2}{0}$, $l$ आणि $\Proj{2}{1}$ यांच्याशी
केलेले संयोजन आहे, येथे
\[l(x,y) = g(\Proj{2}{0}(x,y),\Proj{2}{0}(x,y),\Proj{2}{1}(x,y)),
\]
म्हणजे $l$ हे $g$ चे $\Proj{2}{0}$, $\Proj{2}{0}$ आणि~$\Proj{2}{1}$
यांच्याशी केलेले संयोजन आहे.











\section{आदिम पुनरावर्ती फलने}\label{cmp:rec:prf:sec}

आदिम पुनरावर्तन आणि संयोजन वापरून आधीच्या फलनांपासून नवीन फलने कशी
परिभाषित करता येतात, ते पुन्हा नोंदवू.

\begin{defn}
  \label{cmp:rec:prf:defn:primitive-recursion} $f$ हे $k$-स्थानी फलन ($k\ge 1$)
  आणि $g$ हे $(k+2)$-स्थानी फलन आहे असे समजा. \emph{$f$ आणि $g$ पासून
  आदिम पुनरावर्तनाने परिभाषित केलेले फलन} म्हणजे पुढील समीकरणांनी परिभाषित
  केलेले $(k+1)$-स्थानी फलन~$h$:
  \begin{align*}
    h(x_0,\dots,x_{k-1},0) & =  f(x_0,\dots,x_{k-1}) \\
    h(x_0,\dots,x_{k-1},y+1) & = g(x_0,\dots,x_{k-1}, y, h(x_0,\dots,x_{k-1}, y))
  \end{align*}
\end{defn}

\begin{defn}
  \label{cmp:rec:prf:defn:composition} $f$ हे $k$-स्थानी फलन आणि $g_0$, \dots,
  $g_{k-1}$ ही सर्व $n$-स्थानी असलेली $k$ फलने आहेत असे समजा.
  \emph{$f$ चे $g_0$, \dots,~$g_{k-1}$ यांच्याशी संयोजन करून परिभाषित
    केलेले फलन} म्हणजे पुढील समीकरणाने परिभाषित केलेले $n$-स्थानी
  फलन~$h$:
  \[  h(x_0, \dots, x_{n-1}) =
  f(g_0(x_0, \dots, x_{n-1}), \dots, g_{k-1}(x_0, \dots, x_{n-1})).
  \]
\end{defn}

$\Succ$ आणि, प्रत्येक नैसर्गिक संख्या $n$ व प्रत्येक $i < n$ साठी,
प्रक्षेपण फलने
\[\Proj{n}{i}(x_0,\dots,x_{n-1}) = x_i,
\]
यांच्याबरोबरच फलन~$\Zero(x) = 0$ चा आपण आदिम पुनरावर्ती फलनांत
समावेश करू.

\begin{defn}
  आदिम पुनरावर्ती फलनांचा संच म्हणजे $\Nat^n$ पासून $\Nat$ कडे जाणाऱ्या
  आणि पुढील कलमांनी विगमनाधारित रीतीने परिभाषित केलेल्या फलनांचा संच:
  \begin{enumerate}
  \item $\Zero$ हे आदिम पुनरावर्ती आहे.
  \item $\Succ$ हे आदिम पुनरावर्ती आहे.
  \item प्रत्येक प्रक्षेपण फलन $\Proj{n}{i}$ हे आदिम पुनरावर्ती आहे.
  \item $f$ हे $k$-स्थानी आदिम पुनरावर्ती फलन आणि $g_0$,
    \dots,~$g_{k-1}$ ही $n$-स्थानी आदिम पुनरावर्ती फलने असतील, तर
    $f$ चे $g_0$, \dots,~$g_{k-1}$ यांच्याशी केलेले संयोजन आदिम
    पुनरावर्ती असते.
  \item $f$ हे $k$-स्थानी आदिम पुनरावर्ती फलन आणि $g$ हे $k+2$-स्थानी
    आदिम पुनरावर्ती फलन असेल, तर $f$ आणि $g$ पासून आदिम पुनरावर्तनाने
    परिभाषित केलेले फलन आदिम पुनरावर्ती असते.
\end{enumerate}
\end{defn}

\begin{explain}
अधिक संक्षेपाने, आदिम पुनरावर्ती फलनांचा संच हा $\Zero$, $\Succ$ आणि
प्रक्षेपण फलने~$\Proj{n}{j}$ यांचा समावेश असलेला आणि संयोजन व आदिम
पुनरावर्तन यांखाली बंद असलेला लघुतम संच आहे.

आदिम पुनरावर्ती फलनांच्या संचाचे वर्णन करण्याचा आणखी एक मार्ग म्हणजे
त्याची ``टप्प्यांच्या'' रूपात व्याख्या करणे. आरंभीच्या फलनांचा---$\Zero$,
$\Succ$ आणि प्रक्षेपणांचा---संच $S_0$ ने दर्शवू. ही टप्पा~$0$ वरील
आदिम पुनरावर्ती फलने आहेत. टप्पा $S_i$ परिभाषित झाल्यावर, आधीच~$S_i$
मध्ये असलेल्या फलनांना संयोजन किंवा आदिम पुनरावर्तन यांपैकी एका क्रियेचा
एकदा उपयोग करून मिळणाऱ्या सर्व फलनांचा संच $S_{i+1}$ असू द्या. मग
\[S = \bigcup_{i \in \Nat} S_i
\]
हा सर्व आदिम पुनरावर्ती फलनांचा संच आहे.
\end{explain}

$\Add$ हे आदिम पुनरावर्ती फलन आहे, याची खात्री करू.

\begin{prop}
  बेरीज फलन $\Add(x,y) = x+y$ हे आदिम पुनरावर्ती आहे.
\end{prop}

\begin{proof}
$\Add$ ची $f$ आणि~$g$ या दोन फलनांच्या साहाय्याने केलेली व
\ref{cmp:rec:prf:defn:primitive-recursion} च्या आकृतिबंधाशी जुळणारी आदिम
पुनरावर्ती व्याख्या आपल्याकडे आधीच आहे:
\begin{align*}
  \Add(x_0, 0) & = f(x_0) = x_0\\
  \Add(x_0, y+1) & = g(x_0, y, \Add(x_0, y)) =
  \Succ(\Add(x_0, y))
\end{align*}
म्हणून $f$ आणि $g$ ही आदिम पुनरावर्ती असतील, तर $\Add$ देखील आदिम
पुनरावर्ती असते. $f(x_0) = x_0 = \Proj{1}{0}(x_0)$, आणि प्रक्षेपण
फलने आदिम पुनरावर्ती मानली जातात; म्हणून $f$ हे आदिम पुनरावर्ती आहे.
फलन $g$ हे पुढीलप्रमाणे परिभाषित केलेले तीन-स्थानी फलन $g(x_0, y, z)$ आहे:
\[g(x_0, y, z) = \Succ(z).
\]
यावरून $g$ आदिम पुनरावर्ती आहे असे अजून म्हणता येत नाही, कारण $g$ आणि
$\Succ$ ही अगदी समान फलने नाहीत: $\Succ$ हे एकस्थानी आहे, पण $g$ तीन-स्थानी
असणे आवश्यक आहे. मात्र संयोजनाने आपण $g$ ची ``औपचारिक'' व्याख्या अशी
देऊ शकतो:
\[g(x_0, y, z) = \Succ(\Proj{3}{2}(x_0, y, z))
\]
$\Succ$ आणि $\Proj{3}{2}$ ही आदिम पुनरावर्ती फलने मानली जात असल्यामुळे,
आदिम पुनरावर्ती फलनांपासून संयोजनाने परिभाषित करता येणारे $g$ देखील आदिम
पुनरावर्ती असते.
\end{proof}

\begin{prop}
  \label{cmp:rec:prf:prop:mult-pr}
  गुणाकार फलन $\Mult(x,y) = x \cdot y$ हे आदिम पुनरावर्ती आहे.
\end{prop}

\begin{proof}
  सराव.
\end{proof}

\begin{prob}
  $\Mult$ ची आदिम पुनरावर्ती व्याख्या
  \ref{cmp:rec:prf:defn:primitive-recursion} ने आवश्यक केलेल्या
  रूपात मांडता येते हे दाखवून आणि संबंधित $f$ व~$g$ ही फलने आदिम
  पुनरावर्ती आहेत हे दाखवून \ref{cmp:rec:prf:prop:mult-pr} सिद्ध करा.
\end{prob}

\begin{ex}
आदिम पुनरावर्ती व्याख्येचे आपले अगदी पहिले उदाहरण पुढीलप्रमाणे होते:
\begin{align*}
h(0) & =  1 \\
h(y+1) & =  2 \cdot h(y).
\end{align*}
हे फलन \ref{cmp:rec:prf:defn:primitive-recursion} ने आवश्यक केलेल्या रूपात बसू शकत
नाही, कारण $k=0$. व्याख्येत $1$ आणि $2$ हे स्थिरही येतात. पहिली अडचण
टाळण्यासाठी एक नाममात्र आदान आणू आणि फलन~$h'$ परिभाषित करू:
\begin{align*}
h'(x_0, 0) & =  f(x_0) = 1 \\
h'(x_0, y+1) & =  g(x_0, y, h'(x_0, y)) = 2 \cdot h'(x_0, y).
\end{align*}
फलन $f(x_0) = 1$ हे $\Succ$ आणि $\Zero$ यांपासून संयोजनाने परिभाषित
करता येते: $f(x_0) = \Succ(\Zero(x_0))$. फलन $g$ हे
$g'(z) = 2 \cdot z$ आणि प्रक्षेपणे यांपासून संयोजनाने परिभाषित करता येते:
\begin{align*}
g(x_0, y, z) & = g'(\Proj{3}{2}(x_0, y, z))
\intertext{आणि पुढे $g'$ ची स्वतःची व्याख्या संयोजनाने अशी देता येते:}
g'(z) & = \Mult(g''(z), \Proj{1}{0}(z))
\intertext{आणि}
g''(z) & = \Succ(f(z)),
\end{align*}
येथे $f$ वरीलप्रमाणेच आहे: $f(z) = \Succ(\Zero(z))$. आता आपल्याकडे~$h'$
आल्यावर, पुन्हा संयोजन वापरून $h(y) =
h'(\Proj{1}{0}(y),\Proj{1}{0}(y))$ असे ठेवता येते. यावरून मूलभूत
फलनांपासून संयोजने आणि आदिम पुनरावर्तने यांचा अनुक्रम वापरून $h$ ची
व्याख्या करता येते; म्हणून $h$ हे आदिम पुनरावर्ती आहे.
\end{ex}











\section{आदिम पुनरावर्ती फलनांची चिन्हांकने}\label{cmp:rec:not:sec}

आदिम पुनरावर्ती फलनांचे अचूक विगमनाधारित वर्णन उपलब्ध असल्याचा एक फायदा
म्हणजे त्यांचे वर्णन पद्धतशीरपणे करता येते. उदाहरणार्थ, पुढीलप्रमाणे
प्रत्येक अशा फलनाला एक ``चिन्हांकन'' देता येते. शून्य, उत्तरवर्ती आणि
प्रक्षेपणे यांसाठी अनुक्रमे $\Zero$, $\Succ$ आणि $\Proj{n}{i}$ ही चिन्हे
वापरा. आता $h$ हे $k$-स्थानी फलन~$f$ आणि $n$-स्थानी फलने $g_0$,
\dots,~$g_{k-1}$ यांच्यापासून संयोजनाने परिभाषित केलेले आहे, आणि या
फलनांना अनुक्रमे $F$, $G_0$, \dots,~$G_{k-1}$ ही चिन्हांकने दिली आहेत
असे समजा. मग नवीन चिन्ह $\fn{Comp}_{k,n}$ वापरून फलन $h$ हे
$\fn{Comp}_{k,n}[F,G_0,\dots,G_{k-1}]$ ने दर्शवता येते.

आदिम पुनरावर्तनाने परिभाषित केलेल्या फलनांसाठी समरूप चिन्हांकने वापरता
येतात. $(k+1)$-स्थानी फलन~$h$ हे $k$-स्थानी फलन~$f$ आणि $(k+2)$-स्थानी
फलन~$g$ यांच्यापासून आदिम पुनरावर्तनाने परिभाषित केलेले आहे, आणि $f$
व~$g$ यांना अनुक्रमे $F$ आणि~$G$ ही चिन्हांकने दिली आहेत असे समजा. मग
$h$ ला दिलेले चिन्हांकन $\fn{Rec}_k[F,G]$ असते.

बेरीज फलनाची आदिम पुनरावर्तनाने व्याख्या पुढीलप्रमाणे आहे, हे आठवा:
\begin{align*}
  \Add(x_0, 0) & = \Proj{1}{0}(x_0) = x_0\\
  \Add(x_0, y+1) & = \Succ(\Proj{3}{2}(x_0, y, \Add(x_0, y))) = \Add(x_0, y) +1
\end{align*}
येथे~$f$ ची भूमिका $\Proj{1}{0}$ बजावते आणि~$g$ ची भूमिका
$\Succ(\Proj{3}{2}(x_0, y, z))$ बजावते. हे $1$-स्थानी फलन~$\Succ$
आणि $3$-स्थानी फलन~$\Proj{3}{2}$ यांच्यापासून संयोजनाने परिभाषित
केलेल्या फलनाचा परिणाम असल्यामुळे त्याला
$\fn{Comp}_{1,3}[\Succ,\Proj{3}{2}]$ हे चिन्हांकन दिले जाते. या
व्यवस्थेत बेरीज फलन पुढील चिन्हांकनाने दर्शवता येते:
\[\fn{Rec}_1[\Proj{1}{0},\fn{Comp}_{1,3}[\Succ,\Proj{3}{2}]].
\]
अशी चिन्हांकने कधीकधी उपयोगी ठरतात; उदाहरणार्थ, आदिम पुनरावर्ती फलनांचे
प्रगणन करताना.

\begin{prob}
  $\Mult$ साठी संपूर्ण आदिम पुनरावर्ती चिन्हांकन द्या.
\end{prob}











\section{आदिम पुनरावर्ती फलने संगणनक्षम असतात}\label{cmp:rec:cmp:sec}

फलन $h$ हे आदिम पुनरावर्तनाने परिभाषित केले आहे असे समजा:
\begin{eqnarray*}
h(\vec x, 0)     & = & f(\vec x) \\
h(\vec x, y + 1) & = & g(\vec x, y, h(\vec x, y))
\end{eqnarray*}
आणि $f$ व $g$ ही फलने संगणनक्षम आहेत असे समजा. ($x_0$, \dots,
$x_{k-1}$ यांचे संक्षिप्त रूप म्हणून आपण $\vec x$ वापरतो.) मग
$h(\vec x, 0)$ संगणित करता येते हे उघड आहे, कारण ते फक्त $f(\vec x)$
आहे आणि ते संगणनक्षम आहे असे आपण गृहीत धरले आहे. त्यानंतर
$h(\vec x, 1)$ देखील संगणित करता येते, कारण $1 = 0 + 1$ आणि म्हणून
$h(\vec x, 1)$ चे मूल्य पुढीलप्रमाणे मिळते:
\begin{align*}
 h(\vec x, 1) & = g(\vec x, 0, h(\vec x, 0)) =  g(\vec x, 0, f(\vec x)).
\intertext{याच प्रकारे पुढे जाऊन आपण पुढील मूल्ये संगणित करू शकतो:}
h(\vec x, 2) & = g(\vec x, 1, h(\vec x, 1)) = g(\vec x, 1, g(\vec x, 0, f(\vec x)))\\
h(\vec x, 3) & = g(\vec x, 2, h(\vec x, 2)) = g(\vec x, 2, g(\vec x, 1, g(\vec x, 0, f(\vec x))))\\
h(\vec x, 4) & = g(\vec x, 3, h(\vec x, 3)) = g(\vec x, 3, g(\vec x, 2, g(\vec x, 1, g(\vec x, 0, f(\vec x)))))\\
& \vdots
\end{align*}
अशा रीतीने, सर्वसाधारणपणे $h(\vec x, y)$ संगणित करण्यासाठी
$h(\vec x, 0)$, $h(\vec x, 1)$, \dots, ही मूल्ये क्रमाक्रमाने
$h(\vec x, y)$ पर्यंत संगणित करा.

म्हणून $f$ आणि $g$ ही फलने संगणनक्षम असतील, तर आदिम पुनरावर्ती
व्याख्येपासून नवीन संगणनक्षम फलन मिळते. त्याचप्रमाणे $f$ आणि $g_i$ ही
फलने संगणनक्षम असतील, तर त्यांच्या संयोजनापासूनही संगणनक्षम फलन मिळते.

मूलभूत फलने $\Zero$, $\Succ$ आणि $\Proj{n}{i}$ ही संगणनक्षम आहेत, आणि
संयोजन व आदिम पुनरावर्तन यांमुळे संगणनक्षम फलनांपासून संगणनक्षम फलने
मिळतात. म्हणून प्रत्येक आदिम पुनरावर्ती फलन संगणनक्षम असते.











\section{आदिम पुनरावर्ती फलनांची उदाहरणे}\label{cmp:rec:exa:sec}


आदिम पुनरावर्ती फलनांची काही उदाहरणे आपल्याकडे आधीच आहेत: बेरीज आणि
गुणाकार ही फलने~$\Add$ आणि $\Mult$. एकरूपता फलन $\fn{id}(x) = x$ हे
आदिम पुनरावर्ती आहे, कारण ते फक्त~$\Proj{1}{0}$ आहे. स्थिर फलने
$\fn{const}_n(x) = n$ ही आदिम पुनरावर्ती आहेत, कारण $\Zero$ आणि $\Succ$
यांपासून क्रमिक संयोजनाने ती परिभाषित करता येतात. आदिम पुनरावर्ती
व्याख्यांत स्थिरांचा वापर करायचा असेल तेव्हा हे उपयोगी पडते. उदाहरणार्थ,
फलन $f(x) = 2 \cdot x$ परिभाषित करायचे असेल, तर $\fn{const}_n(x)$ आणि
गुणाकार यांपासून संयोजन करून ते $f(x) =
\Mult(\fn{const}_2(x), \Proj{1}{0}(x))$ असे मिळवता येते. यापुढे आपण
ही युक्ती वापरू.

\begin{prop}
  घातांक फलन $\fn{exp}(x, y) = x^y$ हे आदिम पुनरावर्ती आहे.
\end{prop}

\begin{proof}
  आपण $\fn{exp}$ ची आदिम पुनरावर्ती व्याख्या अशी देऊ शकतो:
  \begin{align*}
    \fn{exp}(x, 0) & = 1\\
    \fn{exp}(x, y+1) & = \Mult(x, \fn{exp}(x,y)).
    \intertext{काटेकोरपणे पाहता ही आदिम पुनरावर्ती फलनांपासून केलेली
      पुनरावर्ती व्याख्या नाही. मात्र औपचारिक रूपात आपल्याकडे पुढील
      व्याख्या आहे:}
    \fn{exp}(x, 0) & = f(x)\\
    \fn{exp}(x, y+1) & = g(x, y, \fn{exp}(x,y)).
    \intertext{येथे}
    f(x) & = \Succ(\Zero(x)) = 1\\
    g(x, y, z) & = \Mult(\Proj{3}{0}(x, y, z), \Proj{3}{2}(x, y, z)) = x \cdot z
  \end{align*}
  आणि म्हणून $f$ व $g$ ही आदिम पुनरावर्ती फलनांपासून संयोजनाने
  परिभाषित केलेली आहेत.
\end{proof}

\begin{prop}
  पुढीलप्रमाणे परिभाषित केलेले पूर्ववर्ती फलन $\fn{pred}(y)$,
  \[  \fn{pred}(y) = \begin{cases}
    0 & \text{जर $y=0$ असेल}\\
    y-1 & \text{अन्यथा}
  \end{cases}
  \]
  हे आदिम पुनरावर्ती आहे.
\end{prop}

\begin{proof}
 पुढील गोष्ट लक्षात घ्या:
 \begin{align*}
   \fn{pred}(0) & = 0 \text{ आणि}\\
   \fn{pred}(y+1) & = y.
 \end{align*}
 ही जवळजवळ आदिम पुनरावर्ती व्याख्या आहे. काटेकोरपणे पाहता ती आदिम
 पुनरावर्तनाने केलेल्या व्याख्येच्या आकृतिबंधात बसत नाही, कारण त्या
 आकृतिबंधासाठी किमान एक अतिरिक्त आदान~$x$ आवश्यक असते. तसेच
 $\fn{pred}(y+1)$ च्या व्याख्येत $\fn{pred}(y)$ प्रत्यक्षात वापरले जात
 नाही, हेही काहीसे वेगळे आहे. पण आपण प्रथम पुढीलप्रमाणे
 $\fn{pred}'(x, y)$ परिभाषित करू शकतो:
 \begin{align*}
   \fn{pred}'(x, 0) & = \Zero(x) = 0,\\
   \fn{pred}'(x, y+1) & = \Proj{3}{1}(x, y, \fn{pred'}(x, y)) = y.
 \end{align*}
आणि नंतर त्यापासून संयोजनाने $\fn{pred}$ परिभाषित करू शकतो; उदाहरणार्थ,
$\fn{pred}(x) = \fn{pred}'(\Zero(x), \Proj{1}{0}(x))$ असे.
\end{proof}

\begin{prop}
  क्रमगुणित फलन $\fn{fac}(x) = \fact{x} = 1 \cdot 2 \cdot 3
  \cdot \dots \cdot x$ हे आदिम पुनरावर्ती आहे.
\end{prop}

\begin{proof}
  त्याची स्पष्ट आदिम पुनरावर्ती व्याख्या अशी आहे:
  \begin{align*}
    \fn{fac}(0) &= 1\\
    \fn{fac}(y+1) & = \fn{fac}(y) \cdot (y+1).
    \intertext{औपचारिक रूपात आपल्याला प्रथम दोन-स्थानी फलन $h$ परिभाषित करावे लागते:}
    h(x, 0) & = \fn{const}_1(x)\\
    h(x, y+1) & = g(x, y, h(x, y))
    \intertext{येथे $g(x, y, z) = \Mult(\Proj{3}{2}(x, y, z),
      \Succ(\Proj{3}{1}(x, y, z)))$; आणि नंतर पुढीलप्रमाणे ठेवावे:}
    \fn{fac}(y) & = h(\Proj{1}{0}(y), \Proj{1}{0}(y)) = h(y,y).
  \end{align*}
  यापुढे आपण थोडी अधिक औपचारिक मोकळीक घेऊ आणि संयोजन व आदिम
  पुनरावर्तन यांच्या साहाय्याने केलेल्या पूर्ण औपचारिक व्याख्या प्रत्येक
  वेळी देणार नाही.
\end{proof}

\begin{prop}
  पुढीलप्रमाणे परिभाषित केलेली छाटलेली वजाबाकी, $x \tsub y$,
  \[  x \tsub y = \begin{cases}
    0 & \text{जर $x < y$ असेल}\\
    x-y & \text{अन्यथा}
  \end{cases}
  \]
  ही आदिम पुनरावर्ती आहे.
\end{prop}

\begin{proof}
  आपल्याकडे पुढील समीकरणे आहेत:
  \begin{align*}
    x \tsub 0 & = x\\
    x \tsub (y+1) & = \fn{pred}(x \tsub y)
  \end{align*}
\end{proof}

\begin{prop}
 $x$ आणि $y$ यांतील अंतर, $\left|x-y\right|$, हे आदिम पुनरावर्ती आहे.
\end{prop}

\begin{proof}
  आपल्याकडे $\left| x-y \right| = (x \tsub y) + (y \tsub x)$ आहे;
  म्हणून अंतर हे $+$ आणि $\tsub$ यांपासून संयोजनाने परिभाषित करता येते,
  आणि ही दोन्ही आदिम पुनरावर्ती आहेत.
\end{proof}

\begin{prop}
  $x$ आणि $y$ यांतील महत्तम, $\fn{max}(x,y)$, हे आदिम पुनरावर्ती आहे.
\end{prop}

\begin{proof}
  $+$ आणि $\tsub$ यांपासून संयोजनाने $\fn{max}(x,y)$ ची व्याख्या अशी
  देता येते:
  \[  \fn{max}(x,y) \defis x + (y \tsub x).
  \]
  $x$ हे महत्तम असेल, म्हणजे $x \ge y$ असेल, तर $y \tsub x = 0$;
  म्हणून $x + (y \tsub x) = x + 0 = x$. $y$ हे महत्तम असेल, तर
  $y \tsub x = y - x$, आणि म्हणून $x + (y \tsub x) = x + (y - x) = y$.
\end{proof}

\begin{prop}
  \label{cmp:rec:exa:prop:min-pr}
  $x$ आणि $y$ यांतील लघुतम, $\fn{min}(x,y)$, हे आदिम पुनरावर्ती आहे.
\end{prop}

\begin{proof}
  सराव.
\end{proof}

\begin{prob}
  \ref{cmp:rec:exa:prop:min-pr} सिद्ध करा.
\end{prob}

\begin{prob}
पुढील फलन आदिम पुनरावर्ती आहे हे दाखवा:
\[f(x, y) =
2^{(2^{\iddots^{2^{x}}})}\raisebox{1ex}{\bigg\rbrace}
\raisebox{1ex}{\text {$y$ वेळा $2$}}\]
\end{prob}

\begin{prob}
पूर्णांक भागाकार $d(x, y) = \lfloor x/y \rfloor$ (म्हणजे दशांश
बिंदूनंतर येणारे सर्व काही वगळलेला भागाकार) हा आदिम पुनरावर्ती आहे, हे
दाखवा. $y = 0$ असताना $d(x, y) = 0$ असे आपण ठरवतो. आदिम पुनरावर्तन
आणि संयोजन वापरून~$d$ ची स्पष्ट व्याख्या द्या.
\end{prob}

\begin{prop}
आदिम पुनरावर्ती फलनांचा संच पुढील दोन क्रियांखाली बंद आहे:
\begin{enumerate}
\item सांत बेरजा: $f(\vec x, z)$ हे आदिम पुनरावर्ती असेल, तर पुढील
फलनही आदिम पुनरावर्ती असते:
\[g(\vec x, y) \defis \sum_{z = 0}^y f(\vec x, z).
\]
\item सांत गुणाकार: $f(\vec x, z)$ हे आदिम पुनरावर्ती असेल, तर पुढील
फलनही आदिम पुनरावर्ती असते:
\[h(\vec x, y) \defis \prod_{z = 0}^y f(\vec x, z).
\]
\end{enumerate}
\end{prop}

\begin{proof}
उदाहरणार्थ, सांत बेरजा पुढील समीकरणांनी पुनरावर्ती रीतीने परिभाषित होतात:
\begin{align*}
  g(\vec x, 0) & = f(\vec x, 0)\\
  g(\vec x, y+1) & = g(\vec x, y) + f(\vec x, y+1).
\end{align*}
\end{proof}












\section{आदिम पुनरावर्ती संबंध}\label{cmp:rec:prr:sec}


\begin{defn}
संबंध $R(\vec x)$ चे जातिबोधक फल
\[\Char{R}(\vec x) = \left\{
  \begin{array}{ll}
  1 & \mbox{जर $R(\vec x)$ असेल} \\
  0 & \mbox{अन्यथा}
  \end{array}
\right.
\]
आदिम पुनरावर्ती असेल, तर त्या संबंधाला आदिम पुनरावर्ती म्हणतात.
\end{defn}

दुसऱ्या शब्दांत, आदिम पुनरावर्ती संबंध $R(\vec x)$ असा उल्लेख केला, तर
$\Char{R}(\vec x) = 1$ या रूपातील संबंध अभिप्रेत असतो; येथे $\Char{R}$
हे कोणत्याही आदानासाठी 1 किंवा 0 यांपैकी एक मूल्य परत करणारे आदिम
पुनरावर्ती फलन असते. उदाहरणार्थ, $x = 0$ असताना आणि केवळ तेव्हाच सत्य
असलेला संबंध $\fn{IsZero}(x)$ हा $\Char{\fn{IsZero}}$ या फलनाशी संबंधित
आहे; हे फलन आदिम पुनरावर्तन वापरून पुढीलप्रमाणे परिभाषित केले जाते:
\begin{align*}
\Char{\fn{IsZero}}(0) & = 1,\\
\Char{\fn{IsZero}}(x+1) & = 0.
\end{align*}

संबंधांचे इतर आदिम पुनरावर्ती फलनांशी संयोजन करता येते, हे स्पष्ट असावे.
म्हणून पुढील संबंधही आदिम पुनरावर्ती आहेत:
\begin{enumerate}
\item समानता संबंध, $x = y$, जो $\fn{IsZero}(\left|x -
  y\right|)$ ने परिभाषित होतो
\item लघुतर अथवा समान संबंध, $x \leq y$, जो $\fn{IsZero}(x
  \tsub y)$ ने परिभाषित होतो
\end{enumerate}


\begin{prop}
  आदिम पुनरावर्ती संबंधांचा संच बूलीय क्रियांखाली बंद आहे; म्हणजे
  $P(\vec x)$ आणि $Q(\vec x)$ आदिम पुनरावर्ती असतील, तर पुढील संबंधही
  आदिम पुनरावर्ती असतात:
  \begin{enumerate}
  \item $\lnot P(\vec x)$
  \item $P(\vec x) \land Q(\vec x)$
  \item $P(\vec x) \lor Q(\vec x)$
  \item $P(\vec x) \lif Q(\vec x)$
  \end{enumerate}
\end{prop}

\begin{proof}
  $P(\vec x)$ आणि $Q(\vec x)$ आदिम पुनरावर्ती आहेत, म्हणजे त्यांची
  जातिबोधक फले $\Char{P}$ आणि $\Char{Q}$ आदिम पुनरावर्ती आहेत, असे
  समजा. $\lnot P(\vec x)$ इत्यादींची जातिबोधक फलेही आदिम पुनरावर्ती
  आहेत हे दाखवायचे आहे.
  \[  \Char{\lnot P}(\vec x) = \begin{cases}
    0 & \text{जर $\Char{P}(\vec x) = 1$ असेल}\\
    1 & \text{अन्यथा}
  \end{cases}
  \]
  आपण $\Char{\lnot P}(\vec x)$ ची व्याख्या $1 \tsub \Char{P}(\vec x)$
  अशी देऊ शकतो.
  \[  \Char{P \land Q}(\vec x) = \begin{cases}
    1 & \text{जर $\Char{P}(\vec x) = \Char{Q}(\vec x) = 1$ असेल}\\
    0 & \text{अन्यथा}
  \end{cases}
  \]
  आपण $\Char{P \land Q}(\vec x)$ ची व्याख्या $\Char{P}(\vec x) \cdot
  \Char{Q}(\vec x)$ किंवा $\fn{min}(\Char{P}(\vec x), \Char{Q}(\vec
  x))$ अशी देऊ शकतो. त्याचप्रमाणे,
  \begin{align*}
    \Char{P \lor Q}(\vec x) & = \fn{max}(\Char{P}(\vec x), \Char{Q}(\vec x)) \text{ आणि}\\
    \Char{P \lif Q}(\vec x) & = \fn{max}(1 \tsub
  \Char{P}(\vec x), \Char{Q}(\vec x)).
  \end{align*}
\end{proof}

\begin{prop}
  आदिम पुनरावर्ती संबंधांचा संच परिबद्ध संख्यकीकरणाखाली बंद आहे; म्हणजे
  $R(\vec x, z)$ हा आदिम पुनरावर्ती संबंध असेल, तर पुढील संबंधही आदिम
  पुनरावर्ती असतात:
  \begin{align*}
    & \bforall{z < y}{R(\vec x, z)} \text{ आणि}\\
    & \bexists{z < y}{R(\vec x, z)}.
  \end{align*}
  $y$ पेक्षा लहान प्रत्येक~$z$ साठी $R(\vec x, z)$ सत्य असेल तेव्हा
  आणि केवळ तेव्हाच $\bforall{z < y}{R(\vec x, z)}$ हे $\vec x$ आणि
  $y$ यांसाठी सत्य असते; $\bexists{z < y}{R(\vec x, z)}$ साठीही समरूप
  अर्थ लागू होतो.
\end{prop}

\begin{proof}
  संकेताप्रमाणे आपण $\bforall{z < 0}{R(\vec x, z)}$ हे सत्य मानतो
  (कारण~$0$ पेक्षा लहान कोणतेही $z$ नसते) आणि
  $\bexists{z < 0}{R(\vec x, z)}$ हे असत्य मानतो. परिबद्ध सार्वत्रिक
  संख्यापक सांत गुणाकार किंवा पुनरावृत्त लघुतमाप्रमाणे कार्य करतो; म्हणजे
  $P(\vec x, y) \defiff \bforall{z < y}{R(\vec x, z)}$ असेल, तर
  $\Char{P}(\vec x, y)$ ची व्याख्या अशी देता येते:
  \begin{align*}
    \Char{P}(\vec x, 0) & = 1\\
    \Char{P}(\vec x, y+1) & =
    \fn{min}(\Char{P}(\vec x, y), \Char{R}(\vec x, y)).
  \end{align*}
  परिबद्ध अस्तित्ववाची संख्यकीकरणाची व्याख्या त्याचप्रमाणे $\fn{max}$
  वापरून देता येते. पर्यायाने, सममूल्यता
  $\bexists{z < y}{R(\vec x, z)}
  \liff \lnot \bforall{z < y}{\lnot R(\vec x, z)}$ वापरून ते परिबद्ध
  सार्वत्रिक संख्यकीकरणापासून परिभाषित करता येते. उदाहरणार्थ,
  $\bexists{x \leq y}{\dots x\dots}$ या रूपातील परिबद्ध संख्यापक
  $\bexists{x < y+1}{\dots x \dots}$ याच्याशी सममूल्य आहे, हे लक्षात घ्या.
\end{proof}

\begin{prob}
  तीन-स्थानी संबंध $x \equiv y \mod n$ (भाजक~$n$ समशेषता) हा आदिम
  पुनरावर्ती आहे, हे दाखवा.
\end{prob}

आणखी एक उपयोगी आदिम पुनरावर्ती फलन म्हणजे पुढीलप्रमाणे परिभाषित केलेले
सशर्त फलन $\fn{cond}(x,y,z)$:
\begin{align*}
  \fn{cond}(x,y,z) & = \begin{cases}
  y & \text{जर $x = 0$ असेल} \\
  z & \text{अन्यथा}.
\end{cases}
\intertext{याची पुनरावर्ती व्याख्या अशी आहे:}
\fn{cond}(0,y,z) & = y,\\
\fn{cond}(x+1,y,z) & = z.
\end{align*}
आदिम पुनरावर्ती संबंधांपासून आदिम पुनरावर्ती फलनांच्या प्रकरणांनुसार
केलेल्या व्याख्यांचे समर्थन करण्यासाठी हे फलन वापरता येते:

\begin{prop}
$g_0(\vec x)$, \dots,~$g_m(\vec x)$ ही आदिम पुनरावर्ती फलने आणि $R_0(\vec
x)$, \dots, $R_{m-1}(\vec x)$ हे आदिम पुनरावर्ती संबंध असतील, तर
पुढीलप्रमाणे परिभाषित केलेले फलन $f$ देखील आदिम पुनरावर्ती असते:
\[f(\vec x) = \begin{cases}
    g_0(\vec x) & \text{जर $R_0(\vec{x})$ असेल} \\
    g_1(\vec x) & \text{जर $R_1(\vec{x})$ असेल आणि $R_0(\vec{x})$ नसेल} \\
    \vdots & \\
    g_{m-1}(\vec x) & \text{जर $R_{m-1}(\vec{x})$ असेल आणि आधीच्या
      कोणत्याही अटी सत्य नसतील}
    \\
    g_m(\vec x) & \mbox{अन्यथा}
\end{cases}
\]
\end{prop}

\begin{proof}
  $m = 1$ असताना हे फक्त पुढीलप्रमाणे परिभाषित केलेले फलन आहे:
  \[  f(\vec x) = \fn{cond}(\Char{\lnot R_0}(\vec x),g_0(\vec x),g_1(\vec
  x)).
  \]
  $m$ हे $1$ पेक्षा मोठे असताना या रूपातील व्याख्यांचे केवळ संयोजन करता
  येते.
\end{proof}










\section{परिबद्ध लघुतमीकरण}\label{cmp:rec:bmi:sec}

\begin{explain}
एखादा गुणधर्म किंवा संबंध~$P$ पूर्ण करणारी लघुतम संख्या म्हणून फलनाची
व्याख्या करणे अनेकदा उपयोगी ठरते. $P$ निर्णेय असेल, तर $P$ पूर्ण करणारी
लघुतम संख्या सापडेपर्यंत $0$, $1$, $2$, \dots, या सर्व संभाव्य संख्यांची
चाचणी घेऊन आपण हे फलन सहज संगणित करू शकतो. या प्रकारच्या अपरिबद्ध
शोधामुळे आपण आदिम पुनरावर्ती फलनांच्या क्षेत्राबाहेर जातो. मात्र स्वतंत्रपणे
दिलेल्या एखाद्या परिबंधापेक्षा \emph{लहान असलेल्या} लघुतम संख्येतच रस
असेल, तर आपण आदिम पुनरावर्ती क्षेत्रात राहतो. दुसऱ्या शब्दांत, आणि थोडे
अधिक सर्वसाधारणपणे, आपल्याकडे आदिम पुनरावर्ती संबंध~$R(x,z)$ आहे असे
समजा. $x$ आणि~$y$ यांना $R(x, z)$ सत्य करणाऱ्या लघुतम $z < y$ कडे
पाठवणारे फलन विचारात घ्या. $R(x, 0)$, $R(x, 1)$, \dots,
$R(x, y-1)$ यांची चाचणी घेऊन तेही संगणित करता येते. पण ते आदिम
पुनरावर्ती का आहे?
\end{explain}

\begin{prop}
$R(\vec x, z)$ आदिम पुनरावर्ती असेल, तर $R(\vec x, z)$ सत्य करणारे
लघुतम~$z$, $y$ पेक्षा लहान असा साक्षीदार असल्यास, परत करणारे आणि अन्यथा
$y$ परत करणारे फलन $m_R(\vec{x}, y)$ देखील आदिम पुनरावर्ती असते.
फलन~$m_R$ आपण असे लिहू:
\[\bmin{z < y}{R(\vec{x}, z)},
\]
\end{prop}

\begin{proof}
$R(\vec{x}, z)$ सत्य करणारे~$z < 0$ असू शकत नाही, कारण मुळातच कोणतेही
$z < 0$ नसते. म्हणून $m_R(\vec x, 0) = 0$.

परिबंध $y + 1$ या रूपाचा असेल, तर तीन प्रकरणे असतात:
\begin{enumerate}
\item $R(\vec{x}, z)$ सत्य करणारे काही $z < y$ असते; अशा वेळी
$m_R(\vec{x}, y+1) = m_R(\vec{x}, y)$.
\item असे कोणतेही~$z<y$ नसते, पण $R(\vec{x}, y)$ सत्य असते; अशा वेळी
$m_R(\vec{x}, y+1) = y$.
\item $R(\vec{x}, z)$ सत्य करणारे कोणतेही $z < y+1$ नसते; अशा वेळी
$m_R(\vec{x}, y+1) = y+1$.
\end{enumerate}
म्हणून आदिम पुनरावर्तनाने $m_R(\vec x, 0)$ ची व्याख्या पुढीलप्रमाणे
देता येते:
\begin{align*}
m_R(\vec{x}, 0) & = 0\\
m_R(\vec{x}, y+1) & =
\begin{cases}
m_R(\vec{x}, y) & \text{जर $m_R(\vec x, y) \neq y$ असेल}\\
y & \text{जर $m_R(\vec x, y) = y$ आणि $R(\vec{x}, y)$ असेल}\\
y+1 & \text{अन्यथा.}
\end{cases}
\end{align*}
$R(\vec x, z)$ सत्य करणारे काही $z<y$ असते तेव्हा आणि केवळ तेव्हाच
$m_R(\vec x, y) \neq y$ असते, हे लक्षात घ्या.
\end{proof}

\begin{prob}
$R(\vec x, z)$ आदिम पुनरावर्ती आहे असे समजा. $R(\vec{x}, z)$ सत्य
करणारे लघुतम~$z$, $y$ पेक्षा लहान असा साक्षीदार असल्यास, परत करणारे आणि
अन्यथा $0$ परत करणारे फलन $m'_R(\vec{x}, y)$ हे~$\Char{R}$ पासून
आदिम पुनरावर्तनाने परिभाषित करा.
\end{prob}











\section{अविभाज्य संख्या}\label{cmp:rec:pri:sec}

परिबद्ध संख्यकीकरण आणि परिबद्ध लघुतमीकरण यांमुळे नेहमी आढळणारी अनेक
फलने व संबंध आदिम पुनरावर्ती आहेत हे दाखवण्यासाठी आपल्याला बरीच साधने
मिळतात. उदाहरणार्थ, ``$x$ ने $y$ ला निःशेष भाग जातो'' हा संबंध विचारात
घ्या; तो $x \mid y$ असा लिहितात. $y$ ला~$x$ ने बाकी न उरता भागता येत
असेल, म्हणजे $y$ ही~$x$ ची पूर्णांक पटीत संख्या असेल, तेव्हा $x \mid y$
हा संबंध सत्य असतो. (ते घडत नसेल, तर आपण $x \nmid y$ असे लिहितो;
भाजक धन असेल तेव्हा याचा अर्थ $y$ ला $x$ ने भागल्यावर बाकी $> 0$ उरते.)
दुसऱ्या शब्दांत, एखाद्या~$z$ साठी $x \cdot z = y$ असेल तेव्हा आणि केवळ
तेव्हाच $x \mid y$. असा~$z$ अस्तित्वात असेल, तर $\leq y$ असलेला एक
साक्षीदार निवडता येतो. म्हणून एखाद्या $z \le y$ साठी $x \cdot z = y$
असेल तेव्हा आणि केवळ तेव्हाच $x \mid y$. त्यामुळे $=$ आणि गुणाकार
यांच्यापासून परिबद्ध अस्तित्ववाची संख्यकीकरण वापरून $x \mid y$ या संबंधाची
व्याख्या पुढीलप्रमाणे करता येते:
\[x \mid y \defiff \bexists{z \leq y}{(x \cdot z) = y}.
\]
अशा रीतीने $x \mid y$ हा संबंध आदिम पुनरावर्ती आहे हे आपण दाखवले.

नैसर्गिक संख्या~$x$ ही $0$ किंवा $1$ नसल्यास आणि तिला फक्त $1$ ने व
स्वतःनेच भाग जात असल्यास तिला \emph{अविभाज्य} म्हणतात. दुसऱ्या शब्दांत,
अविभाज्य संख्यांचे वैशिष्ट्य असे की, जेव्हा $y \mid x$, तेव्हा एकतर
$y = 1$ किंवा~$y=x$. $x$ अविभाज्य आहे की नाही हे तपासण्यासाठी आपल्याला
फक्त सर्व $y \le x$ साठी $y \mid x$ आहे की नाही हे तपासावे लागते, कारण
$y > x$ असल्यास आपोआप~$y \nmid x$. म्हणून, $x$ अविभाज्य असेल तेव्हा आणि
केवळ तेव्हाच सत्य असणाऱ्या $\fn{Prime}(x)$ या संबंधाची व्याख्या
\[\fn{Prime}(x) \defiff x \geq 2 \land \bforall{y \leq x}{(y \mid x \lif y
  = 1 \lor y = x)}
\]
अशी करता येते आणि म्हणून हा संबंध आदिम पुनरावर्ती आहे.

अविभाज्य संख्या $2$, $3$, $5$, $7$, $11$, इत्यादी आहेत. या अनुक्रमात
स्थानांक~$x$ वरची अविभाज्य संख्या परत करणारे $p(x)$ हे फलन विचारात घ्या;
म्हणजे $p(0) = 2$, $p(1) = 3$, $p(2) = 5$, इत्यादी. (सोयीसाठी आपण
$p(x)$ हे अनेकदा $p_x$ असे लिहू; $p_0=2$, $p_1=3$, इत्यादी.)

$x$ पेक्षा मोठी असलेली पहिली अविभाज्य संख्या परत करणारे
$\fn{nextPrime(x)}$ हे फलन आपल्याकडे असेल, तर आदिम पुनरावर्तन वापरून
$p$ ची व्याख्या सहज करता येते:
\begin{align*}
  p(0) & = 2\\
  p(x+1) & = \fn{nextPrime}(p(x))
\end{align*}
$\fn{nextPrime}(x)$ ही $y > x$ आणि $y$ अविभाज्य अशी लघुतम संख्या~$y$ असल्यामुळे
तिचे संगणन अपरिबद्ध शोधाने सहज करता येते. पण यूक्लिडच्या एका निकालामुळे
तिची व्याख्या परिबद्ध लघुतमीकरणानेही करता येते: $x$ पेक्षा मोठी आणि
$\fact{x}+1$ पेक्षा मोठी नसलेली अविभाज्य संख्या नेहमीच अस्तित्वात असते.
\[  \fn{nextPrime(x)} =
  \bmin{y \leq \fact{x}+1}{(y > x \land \fn{Prime}(y))}.
\]
यावरून $\fn{nextPrime}(x)$ आणि म्हणून $p(x)$ ही फलने केवळ संगणनक्षमच नव्हे,
तर आदिम पुनरावर्तीही आहेत.

(उत्सुकता असेल तर यूक्लिडच्या निकालाचा संक्षिप्त पुरावा असा. दोनपेक्षा
लहान मूल्यांसाठी दोन हीच अपेक्षित अविभाज्य संख्या आहे. इतर मूल्यांसाठी
$p_n$ ही $\le x$ असलेली सर्वांत मोठी अविभाज्य संख्या आहे असे समजा आणि
सर्व $\le x$ अविभाज्य संख्यांचा $p =
p_0\cdot p_1 \cdot \dots \cdot p_n$ हा गुणाकार विचारात घ्या. एकतर $p+1$
अविभाज्य असते, किंवा $x$ आणि~$p+1$ यांच्या दरम्यान एखादी अविभाज्य संख्या
असते. ते का? $p+1$ अविभाज्य नाही असे समजा. मग $q \mid p+1$ अशी एखादी
अविभाज्य संख्या असते, जेथे $q < p+1$. $\le x$ असलेल्या कोणत्याही अविभाज्य
संख्येने $p+1$ ला भाग जात नाही. ($p$ च्या व्याख्येनुसार, प्रत्येक
$p_i \le x$ अविभाज्य संख्येने~$p$ ला भाग जातो, म्हणजे बाकी~$0$ उरते.
म्हणून प्रत्येक $p_i \le x$ अविभाज्य संख्येने $p+1$ ला भागल्यास बाकी~$1$
उरते आणि त्यामुळे $p_i \nmid p+1$.) म्हणून $q$ ही $> x$ आणि $< p+1$
असलेली अविभाज्य संख्या आहे. तसेच $p \le \fact{x}$, म्हणून $> x$ आणि $\le
\fact{x}+1$ असलेली अविभाज्य संख्या अस्तित्वात आहे.)

\begin{prob}
परिबद्ध लघुतमीकरण वापरून पूर्णांक भागाकार $d(x, y)$ ची व्याख्या करा.
\end{prob}











\section{क्रमिका}\label{cmp:rec:seq:sec}

आदिम पुनरावर्ती फलनांचा संच उल्लेखनीयरीत्या सुदृढ आहे. तरीही
\emph{क्रमिका} हाताळण्याची पुरेशी पद्धत विकसित केल्यावर आपल्याला आणखी बरेच
काही करता येईल. नैसर्गिक संख्यांच्या सांत क्रमिका आपण पुढील प्रकारे नैसर्गिक
संख्यांशी ओळखू: $\langle a_0, a_1, a_2, \dots, a_k \rangle$ ही क्रमिका
पुढील संख्येशी संगत असेल:
\[p_0^{a_0+1} \cdot p_1^{a_1+1} \cdot p_2^{a_2+1} \cdot \dots \cdot
p_k^{a_k+1}.
\]
उदाहरणार्थ, $\langle 2, 7, 3\rangle$ आणि $\langle 2, 7, 3, 0, 0 \rangle$
या क्रमिकांचे संख्यात्मक संकेतांक वेगळे राहतील याची खात्री करण्यासाठी आपण
घातांकांमध्ये एक मिळवतो. रिकाम्या क्रमिकेचा संकेतांक म्हणून $0$ आणि~$1$
दोन्ही घेता येतात; निश्चिततेसाठी $\emptyseq$ हे~$0$ दर्शवू दे.

क्रमिकांचे हे सांकेतीकरण चालते याचे कारण तथाकथित अंकगणिताचे मूलभूत प्रमेय
होय: प्रत्येक $n \ge 2$ नैसर्गिक संख्या पुढील रूपात एकमेव रीतीने लिहिता
येते:
\[n = p_0^{a_0} \cdot p_1^{a_1} \cdot \dots \cdot p_k^{a_k}
\]
येथे $a_k \ge 1$. त्यामुळे $\tuple{}(a_0,
\dots, a_k) = \tuple{a_0, \dots, a_k}$ हे चित्रण एकास-एक आहे: वेगवेगळ्या
क्रमिका वेगवेगळ्या संख्यांकडे पाठवल्या जातात आणि प्रत्येक संख्येशी जास्तीत
जास्त एकच क्रमिका संगत असते.

आता क्रमिकेची लांबी ठरवणे, तिचा $i$ वा घटक ठरवणे, क्रमिकेच्या शेवटी एक
घटक जोडणे आणि दोन क्रमिकांची जोडणी करणे या सर्व क्रिया आदिम पुनरावर्ती
आहेत हे आपण दाखवू.

\begin{prop}
  क्रमिका~$s$ ची लांबी परत करणारे $\len{s}$ हे फलन आदिम पुनरावर्ती आहे.
\end{prop}

\begin{proof}
  $R(i, s)$ हा संबंध पुढीलप्रमाणे परिभाषित करू:
  \[  R(i, s) \text{ तेव्हा आणि केवळ तेव्हाच }
  p_i \mid s \land p_{i+1} \nmid s.
  \]
  $R$ हा संबंध स्पष्टपणे आदिम पुनरावर्ती आहे. $s$ हा रिकाम्या नसलेल्या
  क्रमिकेचा संकेतांक असेल, म्हणजे
  \[  s = p_0^{a_0+1} \cdot \dots \cdot p_{k}^{a_{k}+1},
  \]
  तेव्हा $p_i \mid s$ असणारी $p_i$ ही सर्वांत मोठी अविभाज्य संख्या असल्यास,
  म्हणजे $i = k$ असल्यास, $R(i,s)$ सत्य असते. म्हणून $s$ ची लांबी $i+1$
  असते तेव्हा आणि केवळ तेव्हाच $p_i$ ही~$s$ ला भाग जाणारी सर्वांत मोठी
  अविभाज्य संख्या असते. त्यामुळे पुढील व्याख्या करता येते:
  \[  \len{s} =
  \begin{cases}
    0 & \text{जर $s = 0$ किंवा $s = 1$ असेल} \\
    1 + \bmin{i < s}{R(i, s)} & \text{अन्यथा.}
  \end{cases}
  \]
  येथे परिबद्ध लघुतमीकरण वापरता येते, कारण $s$ हा क्रमिकेचा संकेतांक असेल
  तेव्हा $R(i,s)$ पूर्ण करणारा फक्त एकच~$i$ असतो आणि असा $i$ अस्तित्वात
  असल्यास तो~$s$ पेक्षा लहान असतो.
\end{proof}

\begin{prop}
  क्रमिका~$s$ च्या शेवटी $a$ जोडल्यावर मिळणारी क्रमिका परत करणारे
  $\fn{append}(s,a)$ हे फलन आदिम पुनरावर्ती आहे.
\end{prop}

\begin{proof}
  $\fn{append}$ ची व्याख्या पुढीलप्रमाणे करता येते:
  \[  \fn{append}(s,a) =
  \begin{cases}
    2^{a+1} & \text{जर $s = 0$ किंवा $s = 1$ असेल} \\
    s \cdot p_{\len{s}}^{a+1} & \text{अन्यथा.}
  \end{cases}
  \]
\end{proof}

\begin{prop}
  $s$ चा $i$ वा घटक परत करणारे $\fn{element}(s,i)$ हे फलन आदिम
  पुनरावर्ती आहे; येथे पहिल्या घटकाला $0$ वा म्हणतात आणि $i$ हा $s$ च्या
  लांबीपेक्षा मोठा किंवा समान असल्यास फलन $0$ परत करते.
\end{prop}

\begin{proof}
$a$ हा~$s$ चा $i$ वा घटक असतो तेव्हा आणि केवळ तेव्हाच $p_i^{a+1}$ ही
$s$ ला भाग जाणारी~$p_i$ ची सर्वांत मोठी घातशक्ती असते, म्हणजे
$p_i^{a+1} \mid s$ पण $p_i^{a+2} \nmid s$. म्हणून:
  \[  \fn{element}(s,i) =
  \begin{cases}
    0 & \mbox{जर $i \geq \len{s}$ असेल} \\
    \bmin{a < s}{(p_i^{a+2} \nmid s)} & \text{अन्यथा.}
  \end{cases}
  \]
\end{proof}

वर परिभाषित केलेल्या फलनांची पूर्ण नावे वापरण्याऐवजी आपण अधिक संक्षिप्त
चिन्हांकन आणू. $\fn{element}(s,i)$ ऐवजी $(s)_i$ वापरू आणि
$\tuple{s_0, \dots, s_k}$ हे पुढील अभिव्यक्तीचे संक्षिप्त रूप मानू:
\[\fn{append}(\fn{append}(\dots \fn{append}(\emptyseq,s_0)
\dots),s_k).
\]
$s$ ची लांबी~$k$ असेल, तर $s$ चे घटक $(s)_0$, \dots,~$(s)_{k-1}$
असतात, हे लक्षात घ्या.

\begin{prop}
दोन क्रमिकांची जोडणी करणारे $\fn{concat}(s,t)$ हे फलन आदिम पुनरावर्ती
आहे.
\end{prop}

\begin{proof}
  आपल्याला पुढील गुणधर्म असलेले $\fn{concat}$ हे फलन हवे आहे:
  \[    \fn{concat}(\tuple{a_0, \dots, a_k}, \tuple{b_0, \dots, b_l}) =
    \tuple{a_0, \dots, a_k, b_0, \dots, b_l}.
  \]
  $t$ चे पहिले $n$ घटक~$s$ च्या शेवटी जोडणारे
  $\fn{hconcat}(s,t,n)$ हे ``साहाय्यक'' फलन वापरू. त्याची आदिम
  पुनरावर्तनाने पुढीलप्रमाणे व्याख्या करता येते:
  \begin{align*}
    \fn{hconcat}(s,t,0) & = s\\
    \fn{hconcat}(s,t,n+1) & = \fn{append}(\fn{hconcat}(s,t,n),(t)_n)
    \intertext{मग $\fn{concat}$ ची व्याख्या पुढीलप्रमाणे करता येते:}
    \fn{concat}(s,t) & = \fn{hconcat}(s,t,\len{t}).
  \end{align*}
\end{proof}

$\fn{concat}(s,t)$ ऐवजी आपण $s \concat t$ असे लिहू.

क्रमिकेच्या लांबीच्या आणि तिच्या सर्वांत मोठ्या घटकाच्या साहाय्याने तिच्या
संख्यात्मक संकेतांकावर परिबंध घालता येणे आपल्याला उपयोगी पडेल. $s$ ही
लांबी~$k$ असलेली क्रमिका आहे आणि तिचा प्रत्येक घटक एखाद्या~$x$ पेक्षा
लहान किंवा समान आहे असे समजा. मग $s$ चे जास्तीत जास्त $k$ अविभाज्य अवयव
असतात; त्यांपैकी प्रत्येक जास्तीत जास्त~$p_{k-1}$ असतो आणि~$s$ च्या
अविभाज्य अवयवीकरणात प्रत्येकाचा घातांक जास्तीत जास्त $x+1$ असतो. दुसऱ्या
शब्दांत, धन लांबीसाठी आपण
\[\fn{sequenceBound}(x,k) = p_{k-1}^{k \cdot (x+1)},
\]
अशी व्याख्या केली, तर वर वर्णन केलेल्या क्रमिका~$s$ चा संख्यात्मक संकेतांक
जास्तीत जास्त~$\fn{sequenceBound}(x,k)$ असतो. शून्य लांबीच्या क्रमिकेसाठी
क्रमिका-परिबंधाचे मूल्य एक मानू.

क्रमिकांवरील असा परिबंध आपल्याला परिबद्ध शोध वापरून नवीन फलने परिभाषित
करण्याची पद्धत देतो. उदाहरणार्थ, परिबद्ध शोध वापरून $\fn{concat}$ साठी
त्याच क्रमिकेचा संकेतांक देणारी पर्यायी व्याख्या करता येते. आपण शोधत असलेल्या वस्तूचे, म्हणजे जोडलेल्या क्रमिकेच्या
संख्यात्मक संकेतांकाचे, आदिम पुनरावर्ती \emph{विनिर्देश} आणि शोध किती दूरवर
करायचा त्यावरील परिबंध एवढेच लिहावे लागते. पुढील व्याख्या हे काम करते:
\begin{align*}
  \fn{concat}(s,t) = {} & \bmin{v < \fn{sequenceBound}(s+t,\len{s} +
    \len{t})}{} \\
  & \quad(\len{v} = \len{s} + \len{t} \land {}\\
    & \qquad (\bforall{i < \len{s}}{((v)_i = (s)_i)}) \land {} \\
      & \qquad (\bforall{j < \len{t}}{((v)_{\len{s}+j} = (t)_j)}))
\end{align*}
या पर्यायी व्याख्येत रिकाम्या क्रमिकेसाठी 0 हा संकेतांक मिळतो. आधीची
पुनरावर्ती व्याख्या दुसरी क्रमिका रिकामी असल्यास पहिल्या क्रमिकेचा
संकेतांक जसाच्या तसा ठेवते. म्हणून दोन्ही क्रमिका रिकाम्या आणि पहिल्या
क्रमिकेचा संकेतांक 1 असल्यास त्या व्याख्येत 1 मिळतो. 0 आणि 1 हे दोन्ही
रिकाम्या क्रमिकेचे संकेतांक असल्यामुळे क्रमिकेचा अर्थ तोच राहतो; मात्र
या एका प्रकरणात संख्यात्मक संकेतांक समान नसतात.


\begin{prob}
पुढील गुणधर्म असलेले आदिम पुनरावर्ती फलन~$\fn{sconcat}(s)$ अस्तित्वात आहे
हे दाखवा:
\[\fn{sconcat}(\tuple{s_0, \dots, s_k}) = s_0 \concat \dots \concat s_k.
\]
\end{prob}

\begin{prob}
पुढील गुणधर्म असलेले आदिम पुनरावर्ती फलन~$\fn{tail}(s)$ अस्तित्वात आहे
हे दाखवा:
\begin{align*}
  \fn{tail}(\emptyseq) & = 0 \text{ आणि}\\
  \fn{tail}(\tuple{s_0, \dots, s_{k}}) & = \tuple{s_1, \dots, s_{k}}.
\end{align*}
\end{prob}

\begin{prop}
  \label{cmp:rec:seq:prop:subseq}
  $s$ च्या $i$ व्या घटकापासून सुरू होणारी आणि लांबी~$n$ असलेली उपक्रमिका
  परत करणारे $\fn{subseq}(s, i, n)$ हे फलन आदिम पुनरावर्ती आहे.
\end{prop}

\begin{proof}
  सराव.
\end{proof}

\begin{prob}
  \ref{cmp:rec:seq:prop:subseq} सिद्ध करा.
\end{prob}











\section{वृक्ष}\label{cmp:rec:tre:sec}

ज्याप्रमाणे क्रमिका संख्यांनी आणि त्यांचे गुणधर्म व त्यांवरील क्रिया त्यांच्या
संकेतांकांवरील आदिम पुनरावर्ती संबंधांनी व फलनांनी दर्शवता येतात, त्याचप्रमाणे
वृक्ष नैसर्गिक संख्यांनी दर्शवणे कधीकधी उपयोगी ठरते. हे करण्यासाठी आपण
क्रमिका आणि त्यांचे संकेतांक वापरू. वृक्ष म्हणजे एकच गाठ (जिला नामांकन असू
शकते), किंवा अनेक उपवृक्षांशी जोडलेली गाठ (जिला नामांकन असू शकते). त्या
गाठीला वृक्षाचे \emph{मूळ} म्हणतात आणि तिला जोडलेल्या उपवृक्षांना त्याचे
\emph{तात्काळ उपवृक्ष} म्हणतात.

वृक्षाचा पुनरावर्ती संकेतांक $\tuple{k, d_1, \dots, d_k}$ या क्रमिकेने
देऊ; येथे $k$ ही तात्काळ उपवृक्षांची संख्या आणि $d_1$, \dots,~$d_k$ हे
त्यांच्या संकेतांकांचे मूल्य आहेत. गाठींना नामांकने असतील, तर ती तात्काळ
उपवृक्षांनंतर समाविष्ट करता येतात. म्हणून फक्त $l$ नामांकन असलेली एक गाठ
मिळून बनलेल्या वृक्षाचा संकेतांक $\tuple{0,l}$ असेल; आणि $l_1$ नामांकन
असलेले मूळ हे $l_2$, $l_3$ नामांकने असलेल्या दोन एक-गाठीच्या वृक्षांना
जोडलेले असेल, तर त्या वृक्षाचा संकेतांक $\tuple{2,
  \tuple{0, l_2}, \tuple{0, l_3}, l_1}$ असेल.

\begin{prop}
  \label{cmp:rec:tre:prop:subtreeseq}
  संकेतांक~$t$ असलेल्या वृक्षाच्या सर्व उपवृक्षांचे संकेतांक घटक म्हणून
  असलेल्या क्रमिकेचा संकेतांक परत करणारे $\fn{SubtreeSeq}(t)$ हे फलन
  आदिम पुनरावर्ती आहे.
\end{prop}

\begin{proof}
  प्रथम, $\fn{ISubtrees}(t) = \fn{subseq}(t, 1, (t)_0)$ हे आदिम पुनरावर्ती
  फलन आहे आणि वृक्ष~$t$ च्या तात्काळ उपवृक्षांचे संकेतांक परत करते, हे लक्षात
  घ्या. आता मूळापासून जास्तीत जास्त $n$ गाठी दूर असलेल्या सर्व उपवृक्षांची,
  संभाव्य पुनरुक्तींसह, क्रमिका काढणारे साहाय्यक फलन
  $\fn{hSubtreeSeq}(t,n)$ परिभाषित करू. मूळापासून जास्तीत जास्त $0$ गाठी
  दूर असलेल्या~$t$ च्या उपवृक्षांची क्रमिका---दुसऱ्या शब्दांत, $t$ च्या
  मुळापासून सुरू होणारी क्रमिका---म्हणजे फक्त~$t$ असलेली क्रमिका. $t$ च्या
  स्तर~$n+1$ पर्यंतच्या सर्व उपवृक्षांची क्रमिका मिळवण्यासाठी स्तर~$n$
  पर्यंतच्या उपवृक्षांची क्रमिका आणि स्तर~$n$ पर्यंतच्या त्या उपवृक्षांच्या
  तात्काळ उपवृक्षांची क्रमिका यांची जोडणी करतो. ही दुसरी क्रमिका मिळवण्यासाठी,
  $f(x)$ हे क्रमिकांचे संकेतांक परत करणारे आदिम पुनरावर्ती फलन असेल, तर
  $g_f(s, k) = f((s)_0) \concat
  \dots \concat f((s)_{k-1})$ हेही आदिम पुनरावर्ती आहे, हे लक्षात घ्या:
    \begin{align*}
      g(s, 0) & = \emptyseq\\
      g(s, k+1) & = g(s, k) \concat f((s)_k)
    \end{align*}
    उदाहरणार्थ, $s$ ही वृक्षांची क्रमिका असेल, तर
    $h(s) = g_{\fn{ISubtrees}}(s, \len{s})$ ही~$s$ च्या घटकांच्या सर्व
    तात्काळ उपवृक्षांची क्रमिका देते. तिच्या साहाय्याने
    $\fn{hSubtreeSeq}$ ची पुढीलप्रमाणे व्याख्या करता येते:
    \begin{align*}
      \fn{hSubtreeSeq}(t, 0) & = \tuple{t} \\
      \fn{hSubtreeSeq}(t, n+1) & = \fn{hSubtreeSeq}(t, n) \concat
      h(\fn{hSubtreeSeq}(t, n)).
    \end{align*}
    संकेतांक~$t$ असलेल्या वृक्षातील उपवृक्षांचा महत्तम स्तर, म्हणजे मूळ आणि
    एखादी पर्णगाठ यांमधील महत्तम अंतर, संकेतांक~$t$ ने परिबद्ध असतो. म्हणून
    संकेतांक~$t$ असलेल्या वृक्षाच्या सर्व उपवृक्षांच्या संकेतांकांची क्रमिका
    $\fn{hSubtreeSeq}(t, t)$ ने मिळते.
\end{proof}

\begin{prob}
  \ref{cmp:rec:tre:prop:subtreeseq} च्या सिद्धतेतील
  $\fn{hSubtreeSeq}$ च्या व्याख्येत सर्वसाधारणपणे पुनरुक्ती येतात. परिणामी
  क्रमिकेत प्रत्येक उपवृक्षाचा संकेतांक फक्त एकदाच येईल याची खात्री करणारी
  पर्यायी व्याख्या द्या.
\end{prob}











\section{पुनरावर्तनाचे इतर प्रकार}\label{cmp:rec:ore:sec}

जोडीकरण आणि क्रमिका वापरून आदिम पुनरावर्तनाच्या अधिक असामान्य (आणि
उपयोगी) रूपांचे समर्थन करता येते. उदाहरणार्थ, पुढील व्याख्येप्रमाणे दोन
फलने एकसामयिकपणे परिभाषित करणे अनेकदा उपयोगी ठरते:
\begin{align*}
h_0(\vec x, 0) & = f_0(\vec x) \\
h_1(\vec x, 0) & = f_1(\vec x) \\
h_0(\vec x, y+1) & = g_0(\vec x, y, h_0(\vec x, y), h_1(\vec x, y)) \\
h_1(\vec x, y+1) & = g_1(\vec x, y, h_0(\vec x, y), h_1(\vec x, y))
\end{align*}
हे \emph{एकसामयिक पुनरावर्तनाचे} एक उदाहरण आहे. फलने परिभाषित करण्याची
आणखी एक उपयोगी पद्धत म्हणजे $h(\vec x, y+1)$ चे मूल्य आधीच्या
$h(\vec x, 0)$, \dots,~$h(\vec x, y)$ या \emph{सर्व} मूल्यांच्या आधारे
देणे; उदाहरणार्थ:
\begin{align*}
h(\vec x, 0) & = f(\vec x) \\
h(\vec x, y+1) & = g(\vec x, y, \tuple{h(\vec x, 0), \dots, h(\vec x, y)}).
\end{align*}
पुढील योजना ही कल्पना अधिक संक्षिप्तपणे व्यक्त करते:
\[h(\vec x, y) = g(\vec x, y, \tuple{h(\vec x, 0), \dots, h(\vec x, y-1)})
\]
येथे $y$ हे $0$ असेल, तेव्हा $g$ चा शेवटचा आदान केवळ रिकामी क्रमिका
असतो, असे समजायचे. दोन्ही मांडण्यांमागील कल्पना अशी की, ``उत्तरवर्ती
टप्प्याचे'' संगणन करताना $h$ ला आतापर्यंत संगणित केलेल्या मूल्यांची संपूर्ण
क्रमिका वापरता येते. याला \emph{मूल्यक्रम पुनरावर्तन} म्हणतात. एका विशिष्ट
उदाहरणात, पुढील प्रकारच्या व्याख्येचे समर्थन करण्यासाठी ते वापरता येते:
\begin{align*}
h(\vec x, y) & = \begin{cases}
  g(\vec x, y, h(\vec x, k(\vec x, y))) & \text{जर $k(\vec x, y) < y$ असेल} \\
  f(\vec x) & \text{अन्यथा.}
\end{cases}
\end{align*}
दुसऱ्या शब्दांत, $y$ येथील $h$ चे मूल्य~$k$ ने दिलेल्या \emph{कोणत्याही}
आधीच्या मूल्यावरील $h$ च्या मूल्याच्या आधारे संगणित करता येते.


\begin{prob}
  शेष फल $r(x,y)$ ची मूल्यक्रम पुनरावर्तनाने व्याख्या करा. ($x$, $y$ या
  नैसर्गिक संख्या असतील आणि $y > 0$ असेल, तर $r(x,y)$ ही~$y$ पेक्षा लहान
  अशी संख्या असते की, एखाद्या~$z$ साठी $x = z\times y + r(x,y)$. निश्चिततेसाठी,
  $y=0$ असल्यास $r(x,0) = 0$ असे मानू.)
\end{prob}

ही फलने नेहमीच्या आदिम पुनरावर्तनाने कशी मिळवायची याचा विचार करा. आदिम
पुनरावर्तनाची आणखी एक शेवटची आवृत्ती अधिक लवचीक आहे: तिच्यात वाटेत
\emph{प्राचल} (बाजूची मूल्ये) बदलण्याची मुभा असते:
\begin{align*}
h(\vec x, 0) & = f(\vec x) \\
h(\vec x, y+1) & = g(\vec x, y, h(k(\vec x), y))
\end{align*}
याचाही नेहमीच्या आदिम पुनरावर्तनाने अनुकार करता येतो. (हे करणे अवघड आहे.
सूचनेसाठी, संगणन हाताने उलट क्रमाने उलगडून पाहा.)











\section{आदिम पुनरावर्ती नसलेली फलने}\label{cmp:rec:npr:sec}

आदिम पुनरावर्ती फलनांमध्ये सहज समजणाऱ्या अर्थाने संगणनक्षम असलेली सर्व
फलने समाविष्ट होत नाहीत. सर्व एकचलीय आदिम पुनरावर्ती फलनांची $f_0$,
$f_1$, $f_2$,~\dots अशी सूची बनवता येते आणि आदान~$y$ वरील $f_x$ चे मूल्य
प्रभावीपणे संगणित करता येते, हे सहज स्पष्ट असावे. दुसऱ्या शब्दांत,
\[g(x,y) = f_x(y)
\]
याने परिभाषित केलेले $g(x,y)$ हे फलन संगणनक्षम आहे. पण मग पुढील फलनही
संगणनक्षम आहे:
\begin{eqnarray*}
h(x) & = & g(x,x) + 1 \\
& = & f_x(x) +1.
\end{eqnarray*}
प्रत्येक आदिम पुनरावर्ती फलन $f_i$ साठी $h$ आणि $f_i$ यांची $i$ येथील
मूल्ये वेगळी असतात. म्हणून $h$ हे संगणनक्षम आहे, पण आदिम पुनरावर्ती नाही;
आणि $g$ विषयीही हेच म्हणता येते. ही कँटर यांच्या कर्णरेषीय युक्तिवादाची
``प्रभावी'' आवृत्ती आहे.

संगणनक्षम पण आदिम पुनरावर्ती नसलेल्या फलनांची अधिक स्पष्ट उदाहरणे देता
येतात. उदाहरणार्थ, $g^n(x)$ हे चिन्हांकन $g(g(\dots g(x)))$ दर्शवू दे;
त्यात एकूण $n$ वेळा $g$ येते. तसेच फलनांची $g_0,g_1,\dots$ ही क्रमिका
पुढीलप्रमाणे परिभाषित करू:
\begin{eqnarray*}
g_0(x) & = & x+1 \\
g_{n + 1}(x) & = & g_n^x(x)
\end{eqnarray*}
प्रत्येक $g_n$ हे फलन आदिम पुनरावर्ती आहे, हे पडताळून पाहता येईल. प्रत्येक
पुढचे फलन त्याच्या आधीच्या फलनापेक्षा खूप अधिक वेगाने वाढते; $g_1(x)$ हे
$2x$ इतके, $g_2(x)$ हे $2^x \cdot x$ इतके आणि $g_3(x)$ साधारणपणे $x$
इतक्या $2$ च्या घातांचा स्तंभ जितक्या वेगाने वाढतो तितक्या वेगाने वाढते.
ॲकरमन--पेटर फलन हे मूलतः $G(x) = g_x(x)$ हे फलन आहे आणि ते कोणत्याही
आदिम पुनरावर्ती फलनापेक्षा अधिक वेगाने वाढते, हे दाखवता येते.

आता आदिम पुनरावर्ती फलनांच्या प्रगणनाकडे परत येऊ. आपण प्रत्येक आदिम
पुनरावर्ती फलनाला सांकेतिक चिन्हांकन दिले आहे, हे लक्षात घ्या; म्हणून
चिन्हांकनांचे प्रगणन करणे पुरेसे आहे. प्रत्येक चिन्हांकन~$F$ ला नैसर्गिक
संख्या $\#(F)$ पुढीलप्रमाणे पुनरावर्ती रीतीने देता येते:
\begin{eqnarray*}
\#(0) & = & \langle 0 \rangle \\
\#(S) & = & \langle 1 \rangle \\
\#(\Proj{n}{i}) & = & \langle 2, n, i \rangle \\
\#(\fn{Comp}_{k,l}[H,G_0,\dots,G_{k-1}]) & = & \langle
3,k,l,\#(H),\#(G_0),\dots,\#(G_{k-1}) \rangle \\
\#(\fn{Rec}_l[G,H]) & = & \langle 4, l, \#(G), \#(H) \rangle
\end{eqnarray*}
येथे मागील विभागातील संकेतांक वापरून संख्यांची प्रत्येक क्रमिका नैसर्गिक
संख्या म्हणून पाहता येते, या तथ्याचा उपयोग केला आहे. परिणामी प्रत्येक
संकेतांकाला नैसर्गिक संख्या मिळते. काही क्रमिका (आणि म्हणून काही संख्या)
चिन्हांकनांशी संगत नसतात, हे उघड आहे. पण $i$ ने अशा चिन्हांकनाला सांकेतिक
केले असेल, तर $i$ ने सांकेतिक केलेल्या त्या चिन्हांकनाचे एकचलीय आदिम
पुनरावर्ती फलन $f_i$ असू दे; आणि
अन्यथा ते स्थिरांक $0$ फलन असू दे. अंतिम परिणाम असा की, एकचलीय आदिम
पुनरावर्ती फलनांचे प्रगणन करण्याची स्पष्ट पद्धत आपल्याला मिळते.

(प्रत्यक्षात, स्थिरांक शून्य फलनासारखी काही फलने सूचीमध्ये एकाहून अधिक वेळा
येतील. हा केवळ आपल्या सांकेतीकरणाचा कृत्रिम परिणाम नाही; स्थिरांक शून्य
फलनाला एकाहून अधिक चिन्हांकने असतात, याचाही तो परिणाम आहे. पुढे आपण पाहू की
या पुनरुक्ती संगणनक्षम रीतीने टाळता येत नाहीत. उदाहरणार्थ, दिलेले चिन्हांकन
स्थिरांक शून्य फलन दर्शवते की नाही हे ठरवणारे कोणतेही संगणनक्षम फलन नसते.)

आता आपण $g(x,y)$ हे फलन $f_x(y)$ ने दिलेले मानू शकतो; येथे $f_x$ म्हणजे
नुकत्याच वर्णन केलेल्या प्रगणनातील फलन. $g(x,y)$ संगणनक्षम आहे हे आपल्याला
कसे माहीत? सहजपणे हे स्पष्ट आहे: $g(x,y)$ संगणित करण्यासाठी प्रथम $x$
``उलगडून'' ते एकचलीय फलनाचे चिन्हांकन आहे की नाही हे पाहा. तसे असल्यास,
आदान~$y$ वरील त्या फलनाचे मूल्य संगणित करा.













\section{आंशिक पुनरावर्ती फलने}\label{cmp:rec:par:sec}


पुनरावर्ती फलनांच्या व्याख्येमागील प्रेरणा समजण्यासाठी पुढील गोष्ट लक्षात
घ्या: संगणनक्षम पण आदिम पुनरावर्ती नसलेली फलने अस्तित्वात आहेत या आपल्या
सिद्धतेतून प्रत्यक्षात आणखी बरेच काही सिद्ध होते. युक्तिवाद साधा होता:
$f_0,f_1,\dots$ या फलनांचे असे प्रगणन करता येते की, $x$ आणि $y$ यांचे फलन
म्हणून $f_x(y)$ संगणनक्षम असते, एवढेच तथ्य आपण वापरले. म्हणून हा युक्तिवाद
\emph{अशा प्रकारे प्रगणन करता येणाऱ्या फलनांच्या कोणत्याही वर्गाला} लागू
होतो. यामुळे आपण अडचणीत येतो: संगणनक्षम फलनांचे स्पष्ट वर्णन आपल्याला हवे
आहे; पण संगणनक्षम फलनांच्या संग्रहाचे अशा प्रकारचे कोणतेही प्रभावी स्पष्ट
वर्णन सर्वसमावेशक असू शकत नाही!

यातून बाहेर पडण्याचा मार्ग म्हणजे \emph{अंशतः} फलनांना स्थान देणे. अंशतः
संगणनक्षम फलनांचे प्रगणन करणे \emph{शक्य आहे}, हे आपण पाहू.
 पुढे आपण आपल्या
कर्णरेषीय युक्तिवादाकडे परत येऊ आणि अंशतः फलने समाविष्ट केल्यावर तो का
चालत नाही याचा शोध घेऊ.

आता प्रश्न असा आहे: सर्व आंशिक पुनरावर्ती फलने मिळवण्यासाठी आदिम पुनरावर्ती
फलनांमध्ये काय भर घालावी लागेल? आपल्याला दोन गोष्टी कराव्या लागतील:
\begin{enumerate}
\item आदिम पुनरावर्ती फलनांची आपली व्याख्या बदलून त्यात अंशतः फलनांनाही
  स्थान द्यावे.
\item काही नवीन अंशतः फलने समाविष्ट होतील अशी काही गोष्ट व्याख्येत
  \emph{जोडावी}.
\end{enumerate}

पहिली गोष्ट सोपी आहे. पूर्वीप्रमाणे आपण शून्य, उत्तरवर्ती आणि प्रक्षेपण
यांपासून सुरुवात करू आणि संयोजन व आदिम पुनरावर्तन यांखाली बंद करू. फरक
इतकाच की, व्याख्येतील काही पदे परिभाषित नसण्याची शक्यता सामावण्यासाठी
संयोजनाची आणि आदिम पुनरावर्तनाची व्याख्या बदलावी लागते. $f$ आणि $g$ ही
अंशतः फलने असतील, तर $x$ येथे $f$ परिभाषित आहे, म्हणजे $x$ हा $f$ च्या
परिभाषाक्षेत्रात आहे, हे दर्शवण्यासाठी $f(x) \fdefined$ असे लिहू; आणि
याच्या विरुद्ध अर्थासाठी, म्हणजे $x$ येथे $f$ परिभाषित नाही हे दर्शवण्यासाठी
$f(x) \fundefined$ असे लिहू. $f(x)$ आणि $g(x)$ ही दोन्ही अपरिभाषित आहेत,
किंवा दोन्ही परिभाषित असून समान आहेत, हे दर्शवण्यासाठी $f(x) \simeq g(x)$
असे लिहू. अधिक गुंतागुंतीच्या पदांसाठीही ही चिन्हांकने वापरू. $h$ आणि
$g_0$, \dots,~$g_k$ ही सर्व अंशतः फलने असतील, तर
\[h(g_0(\vec x),\dots,g_k(\vec x))
\]
हे पद तेव्हा आणि केवळ तेव्हाच परिभाषित असते, जेव्हा प्रत्येक $g_i$ हे
$\vec x$ येथे परिभाषित असते आणि $h$ हे $g_0(\vec x)$, \dots,~$g_k(\vec x)$
येथे परिभाषित असते, असा संकेत आपण स्वीकारू. हे समजल्यावर अंशतः फलनांसाठीच्या
संयोजनाच्या आणि आदिम पुनरावर्तनाच्या व्याख्या वरीलप्रमाणेच राहतात; फक्त
``$=$'' च्या जागी ``$\simeq$'' लिहावे लागते.

अंशतः फलने मिळवण्यासाठी आदिम पुनरावर्ती फलनांच्या व्याख्येत आपण
\emph{अपरिबद्ध शोध परिकर्मी} जोडू. $f(x,\vec z)$ हे नैसर्गिक संख्यांवरील
कोणतेही अंशतः फलन असेल, तर $\mu x \; f(x,\vec z)$ ची व्याख्या पुढीलप्रमाणे
करू:
\begin{quote}
  $f(0,\vec z), f(1,\vec z), \dots, f(x,\vec
  z)$ ही सर्व मूल्ये परिभाषित असतील आणि $f(x,\vec z) = 0$ असेल असा लघुतम
  $x$, असा $x$ अस्तित्वात असल्यास.
\end{quote}
अन्यथा $\mu x \; f(x,\vec z)$ अपरिभाषित आहे, असे समजायचे. यामुळे
$\mu x \; f(x,\vec z)$ ची एकमेव व्याख्या होते.

\begin{explain}
आपल्या व्याख्येत अल्गोरिदमचा
किंवा कोणत्याही विशिष्ट संगणन-प्रतिमानाचा संदर्भ नाही, हे लक्षात घ्या. पण
संयोजन आणि आदिम पुनरावर्तनाप्रमाणे अपरिबद्ध शोधामागेही एक क्रियात्मक,
संगणकीय सहजकल्पना आहे. अंशतः फलनाच्या संगणनक्षमतेच्या संदर्भात, ज्या
आदानांसाठी फलन अपरिभाषित असते ती आदाने संगणन न थांबणाऱ्या आदानांशी संगत
असतात. $\mu x \; f(x,\vec z)$ संगणित करण्याची प्रक्रिया अशी: शून्य मूल्य
परत येईपर्यंत $f(0,\vec z), f(1,\vec z), f(2,\vec z)$ संगणित करत राहा.
मात्र मधल्या संगणनांपैकी एखादे थांबले नाही, तर $\mu x \; f(x,\vec z)$ चे
संगणनही थांबत नाही.
\end{explain}

$R(x,\vec z)$ हा कोणताही संबंध असेल, तर $\mu x \; R(x,\vec z)$ ची व्याख्या
$\mu x \; (1 \tsub \Char{R}(x,\vec z))$ अशी करतात. दुसऱ्या शब्दांत,
$R(x,\vec z)$ सत्य असलेले $x$ चे लघुतम मूल्य $\mu x \; R(x,\vec z)$ परत
करते. म्हणून $f(x,\vec z)$ हे सर्वत्र परिभाषित फलन असेल, तर
$\mu x \; f(x,\vec z)$ हे $\mu x \; (f(x,\vec z) = 0)$ याच्याशी समान असते.
पण आपली मूळ व्याख्या अधिक सर्वसाधारण आहे, कारण ती $f(x,\vec z)$ सर्वत्र
परिभाषित नसण्याची शक्यता मान्य करते. (याउलट, संबंधाचे जातिबोधक फल नेहमीच
सर्वत्र परिभाषित असते.)

\begin{defn}
नैसर्गिक संख्यांकडून नैसर्गिक संख्यांकडे जाणाऱ्या (विविध स्थानसंख्या असलेल्या)
अंशतः फलनांचा असा लघुतम संच, ज्यात शून्य, उत्तरवर्ती आणि प्रक्षेपणे आहेत आणि
जो संयोजन, आदिम पुनरावर्तन व अपरिबद्ध शोध यांखाली बंद आहे, त्याला
\emph{आंशिक पुनरावर्ती फलनांचा} संच म्हणतात.
\end{defn}

अर्थात, काही आंशिक पुनरावर्ती फलने सर्वत्र परिभाषित असतील, म्हणजे प्रत्येक
आदानासाठी परिभाषित असतील.

\begin{defn}
\label{cmp:rec:par:defn:recursive-fn}
सर्वत्र परिभाषित असलेल्या आंशिक पुनरावर्ती फलनांच्या संचाला
\emph{पुनरावर्ती फलनांचा} संच म्हणतात.
\end{defn}

पुनरावर्ती फलन सर्वत्र परिभाषित आहे यावर भर देण्यासाठी त्याला कधीकधी
``सर्वत्र परिभाषित पुनरावर्ती'' असे म्हणतात.











\section{प्रमाण रूप प्रमेय}\label{cmp:rec:nft:sec}

\begin{thm}[क्लिनीचे प्रमाण रूप प्रमेय]
\label{cmp:rec:nft:thm:kleene-nf}
असा आदिम पुनरावर्ती संबंध $T(e, x, s)$ आणि असे आदिम पुनरावर्ती फलन
$U(s)$ असते की पुढील गुणधर्म पूर्ण होतो: $f$ हे कोणतेही आंशिक पुनरावर्ती
फलन असेल, तर एखाद्या~$e$ साठी सर्व $x$ बाबत
\[f(x) \simeq U(\umin{s}{T(e, x, s)})
\]
असते.
\end{thm}

\begin{explain}
प्रमाण रूप प्रमेयाची सिद्धता गुंतागुंतीची असली, तरी तिची मूलभूत कल्पना
सोपी आहे. प्रत्येक आंशिक पुनरावर्ती फलनाला एक \emph{निर्देशांक}~$e$
असतो; सहज अर्थाने, तो त्या फलनाचा कार्यक्रम किंवा व्याख्या संकेतबद्ध
करणारी संख्या असते. $f(x) \fdefined$ असेल, तर त्या संगणनाची पद्धतशीर
नोंद करून तिला एखाद्या संख्येने~$s$ संकेतबद्ध करता येते. केवळ $x$ आणि
व्याख्या~$e$ वापरून, $s$ ही संख्या आदान~$x$ वरील~$f$ चे संगणन
संकेतबद्ध करते का, हे आदिम पुनरावर्ती पद्धतीने तपासता येते. म्हणून
``निर्देशांक~$e$ असलेल्या फलनाचे आदान~$x$ साठी संगणन आहे आणि $s$
त्या संगणनाचा संकेतांक आहे'' हा संबंध~$T$ आदिम पुनरावर्ती असतो.
संगणनाची पूर्ण नोंद~$s$ दिली असता, $s$ चा परिणाम म्हणजे~$f(x)$ चे मूल्य;
तेही~$s$ पासून आदिम पुनरावर्ती पद्धतीने मिळवता येते.

प्रमाण रूप प्रमेय दाखवते की कोणत्याही आंशिक पुनरावर्ती फलनाच्या
व्याख्येसाठी एकच अपरिबद्ध शोध पुरेसा असतो. म्हणजे, निर्देशांक~$e$
असलेल्या फलनाचे आदान~$x$ वरील संगणन संकेतबद्ध करणारी संख्या सापडेपर्यंत
आपण सर्व संख्यांमध्ये शोध करू शकतो. आंशिक पुनरावर्ती फलनांची ``नावे''
म्हणून या संख्या~$e$ वापरता येतात आणि प्रमेयातील समीकरणाने परिभाषित
केलेल्या~$f$ फलनासाठी $\cfind{e}$ असे लिहिता येते. एका आंशिक
पुनरावर्ती फलनाला एकापेक्षा अधिक निर्देशांक असू शकतात, हे लक्षात घ्या;
खरे तर प्रत्येक आंशिक पुनरावर्ती फलनाला अनंत निर्देशांक असतात.
\end{explain}











\section{थांबण्याची समस्या}\label{cmp:rec:hlt:sec}

सर्वसाधारणपणे \emph{थांबण्याची समस्या} अशी आहे: संगणनक्षम फलनाचे
विनिर्देशन~$e$ (उदा., कार्यक्रम) आणि संख्या~$n$ दिली असता, आदान~$n$
वरील त्या फलनाचे संगणन थांबते का, म्हणजे काही परिणाम देते का, हे
ठरवणे. ही समस्या स्वतः संगणनक्षम फलनाने सोडवता येत नाही, हे ॲलन
ट्यूरिंग यांनी सिद्ध केले. म्हणजे, पुढील फलन संगणनक्षम नाही:
\[h(e, n) =
\begin{cases}
1 & \text{जर संगणन $e$ आदान $n$ वर थांबत असेल तर}\\
0 & \text{अन्यथा,}
\end{cases}
\]

आंशिक पुनरावर्ती फलनांच्या संदर्भात कार्यक्रमाच्या विनिर्देशनाची भूमिका
क्लिनीच्या प्रमाण रूप प्रमेयातील निर्देशांक~$e$ बजावू शकतो. $f$ हे
आंशिक पुनरावर्ती फलन असेल, तर प्रमाण रूप प्रमेयातील समीकरण ज्या
कोणत्याही $e$~साठी खरे असेल, तो~$f$ चा निर्देशांक असतो. संख्या~$e$
दिली असता, प्रमाण रूप प्रमेय सांगते की
\[\cfind{e}(x) \simeq U(\mu s \; T(e, x, s))
\]
हे आंशिक पुनरावर्ती असते; आणि प्रत्येक आंशिक पुनरावर्ती $f\colon \Nat
\to \Nat$ साठी असा $e \in \Nat$ असतो की सर्व~$x \in \Nat$ साठी
$\cfind{e}(x) \simeq f(x)$. खरे तर अशा प्रत्येक $f$ साठी एकच नव्हे,
तर अनंत निर्देशांक~$e$ असतात. \emph{थांबण्याचे फलन}~$h$ पुढीलप्रमाणे
परिभाषित करतात:
\[h(e, x) =
\begin{cases}
1 & \text{जर $\cfind{e}(x) \fdefined$ असेल तर}\\
0 & \text{अन्यथा.}
\end{cases}
\]
$\cfind{e}(x) \fundefined$ असेल तेव्हा $h(e, x) = 0$, हे लक्षात घ्या.
प्रमाण रूपातील सूत्राने प्रत्येक~$e$ साठी आंशिक पुनरावर्ती फलन मिळते;
म्हणून येथे निर्देशांक नसल्याची वेगळी अपवादस्थिती उद्भवत नाही.


\begin{thm}
\label{cmp:rec:hlt:thm:halting-problem}
थांबण्याचे फलन~$h$ आंशिक पुनरावर्ती नाही.
\end{thm}

\begin{proof}
$h$ आंशिक पुनरावर्ती आहे असे समजल्यास, आपण पुढील फलन परिभाषित
करू शकू:
\[d(y) =
\begin{cases}
1 & \text{जर $h(y, y) = 0$ असेल तर}\\
\umin{x}{x \neq x} & \text{अन्यथा.}
\end{cases}
\]
कोणतीही संख्या $x$ ही $x \neq x$ पूर्ण करत नसल्यामुळे $\umin{x}{x \neq
x}$ अस्तित्वात नाही. म्हणून $d(y) \fundefined$ तेव्हा आणि केवळ तेव्हाच,
जेव्हा $h(y,y) \neq 0$. या व्याख्येवरून पुढील गोष्टी मिळतात:
\begin{enumerate}
\item $d(y) \fdefined$ तेव्हा आणि केवळ तेव्हाच, जेव्हा $\cfind{y}(y)
  \fundefined$; प्रत्येक $y$ हा काही आंशिक पुनरावर्ती फलनाचा निर्देशांक
  असल्यामुळे वेगळी अपवादस्थिती नाही.
\item $d(y) \fundefined$ तेव्हा आणि केवळ तेव्हाच, जेव्हा $\cfind{y}(y)
  \fdefined$.
\end{enumerate}
$h$ आंशिक पुनरावर्ती असेल, तर $d$ ही आंशिक पुनरावर्ती असेल. त्यामुळे
क्लिनीच्या प्रमाण रूप प्रमेयानुसार तिला निर्देशांक~$e_d$ असेल.
$h(e_d, e_d)$ चे मूल्य पाहू. $0$ आणि~$1$ असे दोन पर्याय आहेत.
\begin{enumerate}
\item $h(e_d, e_d) = 1$ असेल, तर $\cfind{e_d}(e_d) \fdefined$.
  पण $\cfind{e_d} \simeq d$ आणि $d(e_d)$ हे $h(e_d, e_d) = 0$
  तेव्हा आणि केवळ तेव्हाच परिभाषित असते. म्हणून $h(e_d, e_d) \neq 1$.
\item $h(e_d, e_d) = 0$ असेल, तर $e_d$ हा त्या फलनाचा निर्देशांक असल्याने
  $\cfind{e_d}(e_d) \fundefined$. पण पुन्हा $\cfind{e_d} \simeq d$
  आणि $d(e_d)$ तेव्हा आणि केवळ तेव्हाच अपरिभाषित असते, जेव्हा
  $\cfind{e_d}(e_d) \fdefined$. त्यामुळे या प्रकरणातही विरोधाभास मिळतो.
\end{enumerate}
यावरून $e_d$ हा आंशिक पुनरावर्ती फलनाचा निर्देशांक असताना वरील
दोन्ही पर्याय अशक्य ठरतात. पण $h$ आंशिक पुनरावर्ती असेल, तर $d$ ही
आंशिक पुनरावर्ती असेल आणि तिचा निर्देशांक म्हणून $e_d$ निवडणे योग्य
असेल. म्हणून $h$ आंशिक पुनरावर्ती असू शकत नाही.
\end{proof}












\section{सामान्य पुनरावर्ती फलने}\label{cmp:rec:gen:sec}

सर्वत्र परिभाषित फलनांचा एक संच मिळवण्याचा आणखी एक मार्ग आहे.
सर्वत्र परिभाषित $f(x,\vec z)$ हे फलन \emph{नियमित} असे म्हणू, जर
नैसर्गिक संख्यांच्या प्रत्येक क्रमिका~$\vec z$ साठी असा $x$ असेल की
$f(x,\vec z) = 0$. दुसऱ्या शब्दांत, ज्या फलनांना अपरिबद्ध शोध लागू
केल्यावर सर्वत्र परिभाषित फलन मिळते, तीच नियमित फलने आहेत.
सावधगिरीने अपरिबद्ध शोध नियमित फलनांपुरताच मर्यादित करता येतो:

\begin{defn}
\label{cmp:rec:gen:defn:general-recursive}
\emph{सामान्य पुनरावर्ती फलनांचा} संच हा विविध चलसंख्यांची नैसर्गिक
संख्यांपासून नैसर्गिक संख्यांकडे जाणारी फलने असलेला सर्वांत लहान संच
आहे: त्यात शून्य, उत्तरवर्ती आणि प्रक्षेपणे आहेत, आणि तो संयोजन, आदिम
पुनरावर्तन आणि \emph{नियमित} फलनांना लागू केलेला अपरिबद्ध शोध यांखाली
बंद आहे.
\end{defn}

प्रत्येक सामान्य पुनरावर्ती फलन सर्वत्र परिभाषित आहे, हे स्पष्ट आहे.
\ref{cmp:rec:gen:defn:general-recursive} आणि \ref{cmp:rec:par:defn:recursive-fn}
यांतील फरक असा की दुसऱ्या व्याख्येत मधल्या टप्प्यांवर आंशिक पुनरावर्ती
फलने वापरण्यास परवानगी आहे; शेवटी मिळणारे फलन सर्वत्र परिभाषित असावे,
एवढीच अट आहे. म्हणून ``सामान्य'' हा ऐतिहासिक अवशेष असलेला शब्द
दिशाभूल करणारा आहे: वरवर पाहता \ref{cmp:rec:gen:defn:general-recursive} ही
व्याख्या \ref{cmp:rec:par:defn:recursive-fn} पेक्षा \emph{कमी} सामान्य
आहे. पण सुदैवाने हा फरक केवळ भासमान आहे. व्याख्या वेगळ्या असल्या,
तरी सामान्य पुनरावर्ती फलनांचा संच आणि पुनरावर्ती फलनांचा संच एकच आहे.




\chapter{संगणनक्षमतेचा सिद्धांत}\label{cmp:thy::chap}
\begin{editorial}
  या प्रकरणातील मजकुराचे पुनरावलोकन करून त्याचा विस्तार करणे आवश्यक आहे.
  विशेषतः, येथे अजून सराव दिलेले नाहीत.
\end{editorial}








\section{परिचय}\label{cmp:thy:int:sec}

\emph{संगणनक्षमता सिद्धांत} म्हणून ओळखली जाणारी तर्कशास्त्राची शाखा
फलने व संच यांची संगणनक्षमता किंवा सापेक्ष संगणनक्षमता यांच्याशी संबंधित
प्रश्नांचा अभ्यास करते. या विषयाला पूर्वी \emph{पुनरावर्तन सिद्धांत}
म्हणत असत, यावरून क्लिनींचा प्रभाव दिसतो; आज ही दोन्ही नावे
सर्वसाधारणपणे वापरली जातात.

एखाद्या संगणनाच्या प्रतिमानात फलन~$f\colon \Nat \pto \Nat$ चे संगणन
करता येत असेल, तर त्याला \emph{आंशिक संगणनक्षम} म्हणू. $f$ सर्वत्र
परिभाषित असेल, तर $f$ ला फक्त \emph{संगणनक्षम} म्हणू. ज्याचे
जातिबोधक फलन~$\Char{R}$ संगणनक्षम आहे, असा संबंध~$R$ देखील
संगणनक्षम म्हणतात. $f$ व~$g$ ही आंशिक फलने असतील, तर
$f(x) \fdefined$ याचा अर्थ $f$ हे~$x$ येथे परिभाषित आहे, म्हणजे
$x$ हे~$f$ च्या परिभाषाक्षेत्रात आहे; $f(x) \fundefined$ याचा
अर्थ उलट, म्हणजे $f$ हे~$x$ येथे परिभाषित नाही. $f(x) \simeq g(x)$
याचा अर्थ $f(x)$ आणि $g(x)$ ही दोन्ही अपरिभाषित आहेत, किंवा दोन्ही
परिभाषित असून समान आहेत.

संगणनाच्या एखाद्या विशिष्ट प्रतिमानाचा उल्लेख न करताही संगणनक्षमता
सिद्धांताचा अभ्यास करता येतो. त्यासाठी एक सार्वत्रिक आंशिक संगणनक्षम
फलन~$\fn{Un}(k, x)$ आहे, हे दाखवतात. त्याच्या आधारे आंशिक संगणनक्षम
फलनांचे प्रगणन करता येते. $k$-वे एकस्थानी आंशिक संगणनक्षम फलन दर्शवण्यासाठी
आपण~$\cfind{k}$ हे चिन्हांकन वापरू; त्याची व्याख्या
$\cfind{k}(x) \simeq \fn{Un}(k, x)$ अशी आहे. (क्लिनी यांनी यासाठी
$\{ k \}$ हे चिन्हांकन वापरले होते; पण अलीकडे ते तितकेसे वापरले
जात नाही.) यापेक्षा थोड्या अधिक सामान्य स्वरूपात, कोणत्याही
स्थानसंख्येच्या आंशिक संगणनक्षम फलनांचे एकसमान प्रगणन करता येते.
$k$-वे $n$-स्थानी आंशिक पुनरावर्ती फलन दर्शवण्यासाठी आपण
$\cfind{k}[n]$ वापरू.

$f(\vec x, y)$ हे सर्वत्र परिभाषित किंवा आंशिक फलन असेल, तर
$\umin{y}{f (\vec x, y)}$ हे~$\vec x$ चे असे फलन आहे जे
$f(\vec x, y) = 0$ होईल असा लघुतम~$y$ परत करते, मात्र
$f(\vec x, 0)$, \dots, $f(\vec x, y-1)$ ही सर्व मूल्ये परिभाषित
असली पाहिजेत. असा~$y$ नसल्यास $\umin{y}{f (\vec x, y)}$
अपरिभाषित असते. $R(\vec x, y)$ हा संबंध असेल, तर
$\umin{y}{R(\vec x, y)}$ ची व्याख्या $R(\vec x, y)$ सत्य
असेल असा लघुतम~$y$ अशी केली जाते; दुसऱ्या शब्दांत,
$1 \tsub \Char{R}(\vec x, y) = 0$ होईल असा लघुतम~$y$.

एखादे फलन संगणनक्षम आहे हे दाखवण्यासाठी पुढील दोन मार्ग आहेत:
\begin{enumerate}
\item काटेकोरपणे: ट्यूरिंग यंत्र किंवा आंशिक पुनरावर्ती फलन स्पष्टपणे
  वर्णन करणे आणि ते इच्छित फलनाचे संगणन करते हे दाखवणे;
\item अनौपचारिकपणे: त्याचे संगणन करणाऱ्या अल्गोरिदमचे वर्णन करणे आणि
  चर्च यांच्या प्रबंधाचा आधार घेणे.
\end{enumerate}
या दोन मार्गांमध्ये नेमकी सीमारेषा नाही. अल्गोरिदमचे तपशीलवार
वर्णन पुरेशी माहिती देणारे असावे, जेणेकरून तत्त्वतः योग्य ट्यूरिंग
यंत्र किंवा आंशिक पुनरावर्ती व्याख्यांची मालिका कशी तयार करायची हे
तुलनेने स्पष्ट होईल. पूर्णपणे काटेकोर व्याख्या फारशा माहितीपूर्ण
असण्याची शक्यता नाही. म्हणून आपण या दोन मार्गांमध्ये योग्य मध्यममार्ग
शोधू; थोडक्यात, अचूक व्याख्या मिळवता येतील याची खात्री देणारे
सहज समजणारे पण काटेकोर पुरावे शोधू.











\section{संगणनांचे सांकेतिकरण}\label{cmp:thy:cod:sec}

संगणनाच्या प्रत्येक प्रतिमानात पुढील गोष्टी करता येतात:
\begin{enumerate}
\item संगणनक्षम फलनांच्या \emph{व्याख्यांचे} पद्धतशीर वर्णन करता येते.
  उदाहरणार्थ, ट्यूरिंग यंत्रांचे विनिर्देश, पुनरावर्ती व्याख्या किंवा
  कार्यक्रमण भाषेतील कार्यक्रम अशा व्याख्या देतात असे मानता येते.
\item एखाद्या व्याख्येने दिलेल्या फलनाचे ठरावीक आदानासाठी झालेले
  संगणन पूर्णपणे नोंदवता येते. उदाहरणार्थ, संगणनाच्या प्रत्येक पायरीवरील
  विन्यासांची क्रमिका (यंत्राची अवस्था आणि पट्टीवरील मजकूर) ट्यूरिंग
  यंत्राचे संगणन वर्णन करते.
\item संगणनाची एखादी कथित नोंद खरोखरच, दिलेल्या व्याख्येच्या
  संगणनक्षम फलनाचे दिलेल्या आदानासाठी झालेले संगणन कसे घडले याची नोंद
  आहे का, हे तपासता येते (त्या आदानावर फलन परिभाषित आहे, म्हणजे
  संगणन थांबते).
\item संगणनाच्या अशा पूर्ण नोंदीच्या वर्णनातून दिलेल्या आदानासाठी
  फलनाचे मूल्य काढता येते. उदाहरणार्थ, थांबणाऱ्या ट्यूरिंग यंत्राच्या
  संगणनाच्या अगदी शेवटच्या पायरीत पट्टीवर असलेला मजकूर हेच मूल्य असते.
\end{enumerate}

सांकेतिकरण वापरून संगणनक्षम फलनाच्या प्रत्येक वर्णनाला संख्यात्मक
\emph{निर्देशांक} देता येतो, आणि त्या निर्देशांकावरून सूचना संगणनक्षम
रीतीने पुन्हा मिळवता येतात. त्याचप्रमाणे, संगणनाची पूर्ण नोंदही एका
संख्येने संकेतबद्ध करता येते. मग ``निर्देशांक~$e$ असलेल्या फलनाच्या
आदान~$x$ वरील संगणनाच्या नोंदीचा संकेतांक~$s$ आहे'' हा परिणामी
अंकगणितीय संबंध आणि ``संकेतांक~$s$ असलेल्या संगणनक्रमिकेचे निर्गत''
हे फलन संगणनक्षम असतात; किंबहुना ती आदिम पुनरावर्ती असतात.

हे मूलभूत तथ्य फार सामर्थ्यवान आहे. त्याच्या आधारे संगणनक्षमतेविषयी
अनेक महत्त्वाचे आणि आश्चर्यकारक परिणाम सिद्ध करता येतात, तेही निवडलेल्या
संगणन-प्रतिमानावर अवलंबून न राहता.











\section{प्रमाण रूप प्रमेय}\label{cmp:thy:nfm:sec}

संगणनक्षम फलनांच्या व्याख्यांचे वर्णन करता येते आणि एखाद्या आदानावर
त्या फलनाच्या संगणनाची पूर्ण नोंद असल्याचे सांगणारे कथित वर्णन खरे आहे
का हे तपासता येते, असे समजा. मग संगणनाच्या प्रतिमानावर अवलंबून न राहता,
कोणत्याही संगणनक्षम फलन~$f$ चे कोणत्याही आदान~$x$ वरील मूल्य पुढीलप्रमाणे
ठरवता येईल, असे वाटते:
\begin{enumerate}
  \item संगणनांच्या नोंदींच्या सर्व शक्य वर्णनांमध्ये शोध घेणे.
  \item एखादी नोंद~$f(x)$ च्या संगणनाची नोंद आहे का हे तपासणे.
  \item योग्य नोंद सापडल्यास तिच्यातून~$f(x)$ चे मूल्य काढणे.
\end{enumerate}
हे खरोखरच शक्य आहे, हाच क्लिनी यांच्या प्रमाण रूप प्रमेयाचा आशय आहे.

\begin{thm}[क्लिनीचे प्रमाण रूप प्रमेय]
\label{cmp:thy:nfm:thm:normal-form}
असा आदिम पुनरावर्ती संबंध~$T(e, x, s)$ आणि असे आदिम पुनरावर्ती
फलन~$U(s)$ असते की पुढील गुणधर्म पूर्ण होतो: $f$ हे कोणतेही
आंशिक संगणनक्षम फलन असेल, तर एखाद्या~$e$ साठी
\[f(x) \simeq U(\umin{s}{T(e, x, s)})
\]
हे सर्व~$x$ साठी असते.
\end{thm}

\begin{proof}[सिद्धतेचे रूपरेषात्मक वर्णन]
संगणनाच्या कोणत्याही प्रतिमानासाठी संगणनक्षम फलन~$f$ चे वर्णन
काटेकोरपणे परिभाषित करून ते वर्णन नैसर्गिक संख्या~$e$ वापरून
संकेतबद्ध करता येते. त्याचप्रमाणे, निर्देशांक~$e$ असलेल्या फलनाचे
आदान~$x$ वरील संगणन कसे घडले याची नोंद करणाऱ्या ``संगणनक्रमिके''ची
कल्पना काटेकोरपणे परिभाषित करता येते. अशा संगणनक्रमिकेला देखील
संख्या~$s$ ने संकेतबद्ध करता येते. हे अशा प्रकारे करता येते की
\begin{enumerate}
  \item संख्या~$s$ ही निर्देशांक~$e$ असलेल्या फलनाच्या आदान~$x$
  वरील संगणनक्रमिकेचा संकेतांक असेल तेव्हा आणि केवळ तेव्हाच खरा असणारा संबंध
  $T(e, x, s)$, आणि
  \item $s$ ने संकेतबद्ध केलेल्या संगणनक्रमिकेला त्या संगणनाचा अंतिम
  परिणाम देणारे फलन $U(s)$
\end{enumerate}
ही दोन्ही संगणनक्षम असतात. किंबहुना, संबंध~$T$ आणि फलन~$U$
आदिम पुनरावर्ती असतात.
\end{proof}

\begin{explain}
प्रमाण रूप प्रमेयाची काटेकोर सिद्धता द्यायची असेल, तर संगणनाचे एक
प्रतिमान ठरवून संगणनक्षम फलनांच्या वर्णनांचे आणि संगणनक्रमिकांचे
सांकेतिकरण तपशीलवार करावे लागेल; तसेच संबंध~$T$ आणि फलन~$U$
आदिम पुनरावर्ती आहेत हे पडताळावे लागेल. बहुतेक उपयोगांसाठी $T$
आणि~$U$ संगणनक्षम आहेत व $U$ सर्वत्र परिभाषित आहे, एवढे पुरेसे असते.

प्रमाण रूप प्रमेयाची सिद्धता एका संक्षिप्त सूत्रवाक्यात लक्षात ठेवणे
कदाचित उत्तम: $\umin{s}{T(e, x, s)}$ हे निर्देशांक~$e$ असलेल्या
फलनाच्या आदान~$x$ वरील संगणनक्रमिकेचा शोध घेते; अशी क्रमिका
सापडल्यास $U$ त्या संगणनाचे निर्गत परत करते.
\end{explain}

संगणनाचे प्रतिमान आंशिक पुनरावर्ती फलने हे असेल (क्लिनी यांनी
मूळतः हेच प्रतिमान वापरले होते), तर या प्रमेयावरून कोणतेही फलन
परिभाषित करण्यासाठी अपरिबद्ध शोधाचा, म्हणजे एकाच
$\umin{y}{f(\vec x, y)}$ प्रवर्तकाचा, एकदाच वापर पुरेसा असतो.
या अर्थाने प्रत्येक आंशिक पुनरावर्ती फलनाला प्रमाण रूप असते.

$T$ आणि $U$ वापरून $\cfind{0}$, $\cfind{1}$, $\cfind{2}$, \dots
हे प्रगणन परिभाषित करता येते. यापुढे $T$ आणि~$U$ यांची एक योग्य
निवड निश्चित केली आहे असे समजू, आणि
\[\cfind{e}(x) \simeq U(\umin{s}{T(e,x,s)})
\]
हे समीकरण $\cfind{e}$ ची \emph{व्याख्या} मानू.

आणखी एक उपयुक्त तथ्य पुढीलप्रमाणे आहे:

\begin{thm}
प्रत्येक आंशिक संगणनक्षम फलनाला अनंत निर्देशांक असतात.
\end{thm}

हेही सहज स्पष्ट होते. एखाद्या संगणनक्षम फलनाचे वर्णन दिले असता,
त्याच फलनाचे (समान आदान-निर्गत जोड्यांचे) संगणन करणारे वेगळे वर्णन
तयार करता येते. उदाहरणार्थ, संगणनावर काहीही परिणाम न करणारी क्रिया
आधी करता येते ($0 = 0$ आहे का ते तपासणे किंवा $5$ पर्यंत मोजणे इत्यादी).
बदललेल्या वर्णनाचा निर्देशांक मूळ निर्देशांकापेक्षा नेहमी वेगळा असेल.
दोन्ही त्याच फलनाचे निर्देशांक आहेत; फक्त संगणनाची पद्धत किंचित वेगळी आहे.











\section{$s$-$m$-$n$ प्रमेय}\label{cmp:thy:smn:sec}

\begin{explain}
पुढील प्रमेयाला ``$s$-$m$-$n$ प्रमेय'' म्हणतात; त्याचे कारण थोड्याच
वेळात स्पष्ट होईल. हे प्रमेय नेमके काय सांगते हे समजून घेणे हाच कठीण
भाग आहे; विधान समजल्यावर ते बरेचसे स्पष्ट वाटेल.
\end{explain}

\begin{thm}
\label{cmp:thy:smn:thm:s-m-n}
  नैसर्गिक संख्यांच्या प्रत्येक $n$ आणि~$m$ जोडीसाठी असे आदिम
  पुनरावर्ती फलन~$s^m_n$ असते की प्रत्येक क्रमिका
  $e$, $a_0$, \dots, $a_{m-1}$, $y_0$ ,\dots, $y_{n-1}$ साठी
  \[  \cfind{s^m_n(e, a_0, \dots, a_{m-1})}[n](y_0, \dots, y_{n-1}) \simeq
  \cfind{e}[m+n](a_0, \dots, a_{m-1}, y_0, \dots, y_{n-1}).
\]
\end{thm}

\begin{explain}
$s^m_n$ हे \emph{कार्यक्रमांवर} कार्य करते असा विचार करणे उपयुक्त आहे.
म्हणजे, $s^m_n$ हे $(m+n)$-स्थानी फलनाचा कार्यक्रम~$e$ आणि निश्चित
आदाने $a_0$, \dots, $a_{m-1}$ घेते; उरलेल्या आदानांच्या $n$-स्थानी
फलनासाठी ते कार्यक्रम
$s^m_n(e, a_0, \dots, a_{m-1})$ परत करते.


\end{explain}











\section{सार्वत्रिक आंशिक संगणनक्षम फलन}\label{cmp:thy:uni:sec}

\begin{thm}
\label{cmp:thy:uni:thm:univ-comp}
एक सार्वत्रिक आंशिक संगणनक्षम फलन~$\fn{Un}(e,x)$ असते. दुसऱ्या
शब्दांत, पुढील गुणधर्म असलेले फलन $\fn{Un}(e,x)$ असते:
\begin{enumerate}
\item $\fn{Un}(e,x)$ आंशिक संगणनक्षम आहे.
\item $f(x)$ हे कोणतेही आंशिक संगणनक्षम फलन असेल, तर अशी नैसर्गिक
संख्या $e$ असते की प्रत्येक~$x$ साठी $f(x) \simeq \fn{Un}(e,x)$.
\end{enumerate}
\end{thm}

\begin{proof}
क्लिनी यांच्या प्रमाण रूप प्रमेयातील (\ref{cmp:thy:nfm:thm:normal-form})
$U$ आणि $T$ घेऊन $\fn{Un}(e,x) \simeq U(\umin{s}{T(e,x,s)})$
अशी व्याख्या करा.
\end{proof}

\begin{explain}
आंशिक संगणनक्षम फलनांचे प्रभावी प्रगणन आहे, हे अचूकपणे सांगण्याचा
हा मार्ग आहे. कल्पना अशी की $f_e(x) = \fn{Un}(e,x)$ ने परिभाषित
केलेल्या फलनाला $f_e$ म्हटले, तर $f_0$, $f_1$, $f_2$, \dots या
क्रमिकेत सर्व आंशिक संगणनक्षम फलने येतात; शिवाय $f_e(x)$ चे
संगणन $e$ आणि~$x$ या दोहोंवर ``एकसमानपणे'' करता येते. सोपेपणासाठी
आपण एकस्थानी फलनांसाठी सार्वत्रिक असलेले द्विस्थानी फलन वापरत आहोत.
परंतु संख्यांच्या क्रमिकांचे सांकेतिकरण करून हे अधिक आदानांसाठी सहज
सामान्य करता येते. उदाहरणार्थ, $f(x,y,z)$ हे $3$-स्थानी आंशिक
पुनरावर्ती फलन असेल, तर $g(x) \simeq f((x)_0, (x)_1, (x)_2)$
हे एकस्थानी आंशिक पुनरावर्ती फलन असते.

\end{explain}











\section{सार्वत्रिक संगणनक्षम फलन नसते}\label{cmp:thy:nou:sec}

आंशिक संगणनक्षम फलनांसाठी सार्वत्रिक असे आंशिक संगणनक्षम फलन असले,
तरी सर्वत्र परिभाषित संगणनक्षम फलनांसाठी सार्वत्रिक असे सर्वत्र
परिभाषित संगणनक्षम फलन नसते.


\begin{thm}
\label{cmp:thy:nou:thm:no-univ}
सार्वत्रिक संगणनक्षम फलन नसते. दुसऱ्या शब्दांत, असे कोणतेही
फलन $\fn{Un}'(k, x)$ घ्या की $f(x)$ हे सर्वत्र परिभाषित संगणनक्षम
फलन असल्यास, अशी नैसर्गिक संख्या~$k$ असते की प्रत्येक~$x$ साठी
$f(x) = \fn{Un}'(k,x)$; तर असे फलन संगणनक्षम नसते.
\end{thm}

\begin{proof}
सिद्धता साध्या विकर्णीकरणाने मिळते: $\fn{Un}'(k,x)$ हे सर्वत्र
परिभाषित आणि संगणनक्षम असेल, तर
\[d(x) = \fn{Un}'(x, x) + 1
\]
हेही सर्वत्र परिभाषित आणि संगणनक्षम असेल. परंतु व्याख्येवरून
$d(k)$ हे $\fn{Un}'(k,k)$ च्या समान नाही. म्हणून प्रत्येक $k$
साठी $d(x)$ आणि~$\fn{Un}'(k, x)$ यांची मूल्ये किमान एका~$x$
साठी, म्हणजे $x = k$ साठी, वेगळी असतात.
\end{proof}

\begin{explain}
वरील \ref{cmp:thy:uni:thm:univ-comp} दाखवते की सार्वत्रिक फलन आंशिक
असण्याची परवानगी दिल्यास हे विकर्णीकरण टाळता येते. म्हणजे,
$\fn{Un}$ हे सर्वत्र परिभाषित संगणनक्षम फलनांसाठी सार्वत्रिक आहे;
पण ते स्वतः सर्वत्र परिभाषित नाही. आंशिक फलनांच्या बाबतीत
विकर्णीकरणाचा युक्तिवाद चालत नाही.
\end{explain}

\begin{prob}
  \ref{cmp:thy:nou:thm:no-univ} च्या सिद्धतेतील विकर्णीकरणाचा युक्तिवाद आंशिक
  फलनांच्या बाबतीत का चालत नाही हे समजण्यासाठी
  $f(x) \simeq \fn{Un}(x,x)+1$ हे फलन विचारात घ्या. ते आंशिक
  संगणनक्षम आहे का? असल्यास, त्याला निर्देशांक~$e$ आहे, म्हणजे
  $f(x) \simeq \fn{Un}(e,x)$. मग~$f(e)$ विषयी काय म्हणता येईल?
\end{prob}












\section{थांबण्याची समस्या}\label{cmp:thy:hlt:sec}

रचनेवरूनच, सार्वत्रिक आंशिक संगणनक्षम फलन~$\fn{Un}(e,x)$ तेव्हा आणि
केवळ तेव्हाच परिभाषित असते, जेव्हा~$e$ ने संकेतबद्ध केलेल्या फलनाच्या आदान~$x$
वरील संगणनातून मूल्य मिळते. असे होते की नाही हे आपण ठरवू शकतो का,
हा नैसर्गिक प्रश्न आहे. प्रत्यक्षात ते ठरवता येत नाही. ट्यूरिंग
यंत्रांच्या संगणन-प्रतिमानात याचा अर्थ, दिलेले ट्यूरिंग यंत्र दिलेल्या
आदानावर थांबते का हे संगणनाने ठरवता येत नाही. म्हणून पुढील प्रमेयाला
``थांबण्याच्या समस्येची अनिर्णेयता'' म्हणतात. खाली आपण दोन सिद्धता
देऊ. पहिली आपल्या आधीच्या चर्चेचा धागा पुढे नेते, तर दुसरी अधिक थेट आहे.

\begin{thm}
\label{cmp:thy:hlt:thm:halting-problem}
पुढीलप्रमाणे असू द्या:
\[h(e, x)  =
\begin{cases}
1 & \text{जर\/ $\fn{Un}(e, x)$ परिभाषित असेल तर} \\
0 & \text{अन्यथा.}
\end{cases}
\]
मग $h$ संगणनक्षम नाही.
\end{thm}

\begin{proof}
$h$ संगणनक्षम आहे असे समजा. मग सार्वत्रिक संगणनक्षम फलन परिभाषित
करता येईल, हे दाखवू. पुढीलप्रमाणे व्याख्या करा:
\[\fn{Un'}(e,x) =
\begin{cases}
\fn{Un}(e,x) & \text{जर $h(e,x) = 1$ असेल तर} \\
0 & \text{अन्यथा.}
\end{cases}
\]
आता $\fn{Un'}(e, x)$ हे सर्वत्र परिभाषित फलन आहे आणि $h$
संगणनक्षम असेल तर तेही संगणनक्षम आहे. उदाहरणार्थ, आदिम पुनरावर्तन
वापरून $g$ ची व्याख्या पुढीलप्रमाणे करता येते:
\begin{align*}
g(0, e, x) & \simeq 0 \\
g(y+1, e, x) & \simeq \fn{Un}(e,x);
\end{align*}
मग
\[\fn{Un'}(e,x) \simeq g(h(e,x),e,x).
\]
$\fn{Un'}(e,x)$ हे $\fn{Un}(e,x)$ परिभाषित असेल तेथे त्याच्याशी
सहमत असल्याने, $\fn{Un'}$ हे सर्वत्र परिभाषित असलेल्या आंशिक
संगणनक्षम फलनांसाठी सार्वत्रिक आहे. पण यामुळे
\ref{cmp:thy:nou:thm:no-univ} शी विरोध होतो.
\end{proof}

\begin{proof}
$h(e,x)$ संगणनक्षम आहे असे समजा. फलन $g$ ची पुढीलप्रमाणे व्याख्या करा:
\[g(x) =
\begin{cases}
  0                & \text{जर $h(x,x) = 0$ असेल तर} \\
  \fundefined & \text{अन्यथा.}
\end{cases}
\]
फलन $g$ आंशिक संगणनक्षम आहे. उदाहरणार्थ, त्याची व्याख्या
$\umin{y}{h(x,x) = 0}$ अशी करता येते. म्हणून एखाद्या~$e$ साठी
प्रत्येक~$x$ बाबत $g(x) \simeq \fn{Un}(e,
x)$ असते. $g$ हे $e$ येथे परिभाषित आहे का? असल्यास, $g$ च्या
व्याख्येवरून $h(e,e) = 0$ ($h$ परिभाषित असेल तेव्हाच ते~$0$ मूल्य
घेऊ शकते). $h$ च्या व्याख्येनुसार, याचा अर्थ $\fn{Un}(e, e)$
अपरिभाषित आहे. पण प्रत्येक~$x$ साठी $g(x) \simeq \fn{Un}(e, x)$
या गृहितकावरून $g(e)$ अपरिभाषित ठरते—विरोध. दुसरीकडे, $g(e)$
अपरिभाषित असेल, तर $h(e,e) \neq 0$, म्हणून $h(e,e) = 1$.
त्यामुळे $\fn{Un}(e, e)$ परिभाषित आहे. पण
$g(x) \simeq \fn{Un}(e, x)$ असल्यामुळे $g(e)$ देखील परिभाषित
असेल. पुन्हा विरोध मिळतो.
\end{proof}













\section{रसेलच्या विरोधापत्तीशी तुलना}\label{cmp:thy:rus:sec}

या विभागातील युक्तिवादांची रसेलच्या विरोधापत्तीशी तुलना केल्यास
त्यांतील साम्य आणि भेद स्पष्ट होतात:
\begin{enumerate}
\item रसेलची विरोधापत्ती: $S = \Setabs{x}{x \notin x}$ असे मानू.
  मग $S
  \in S$ तेव्हा आणि केवळ तेव्हाच, जेव्हा $S \notin S$; हा विरोध आहे.

  \emph{निष्कर्ष:} असा संच~$S$ अस्तित्वात नाही. ``सर्व संचांचा संच''
  अस्तित्वात आहे हे गृहितक संचसिद्धांताच्या इतर स्वयंसिद्धांशी विसंगत आहे.


\item रसेलच्या विरोधापत्तीचे एक रूपांतर: सर्व फलनांच्या संचापासून
  $\{ 0, 1 \}$ मध्ये जाणारे ``फलन'' $F$ पुढीलप्रमाणे परिभाषित करू:
  \[  F(f) =
  \begin{cases}
    1 & \text{जर $f$ हे $f$ च्या परिभाषाक्षेत्रात असेल आणि $f(f) = 0$ असेल तर} \\
    0 & \text{अन्यथा}
  \end{cases}
  \]
  तशाच युक्तिवादावरून $F(F) = 0$ तेव्हा आणि केवळ तेव्हाच, जेव्हा
  $F(F) = 1$; हा विरोध आहे.

  \emph{निष्कर्ष:} $F$ हे फलन नाही. ``सर्व फलनांचा संच'' इतका
  मोठा आहे की तो फलनाचे परिभाषाक्षेत्र असू शकत नाही.

\item विकर्णीकरणाचा युक्तिवाद: $f_0$, $f_1$, \dots हे आंशिक
  संगणनक्षम फलनांचे प्रगणन असू द्या आणि $G \colon \Nat \to
  \{ 0, 1 \}$ ची व्याख्या पुढीलप्रमाणे करू:
  \[  G(x) =
  \begin{cases}
    1 & \text{जर $f_x(x)\downarrow = 0$ असेल तर} \\
    0 & \text{अन्यथा}
  \end{cases}
  \]
  $G$ संगणनक्षम असेल, तर एखाद्या $k$ साठी ते फलन $f_k$
  असेल. पण मग $G(k) = 1$ तेव्हा आणि केवळ तेव्हाच, जेव्हा $G(k) = 0$;
  हा विरोध आहे.

  \emph{निष्कर्ष:} $G$ संगणनक्षम नाही. संचसिद्धांताच्या
  स्वयंसिद्धांनुसार $G$ हे तरीही फलन आहे, हे लक्षात घ्या;
  येथे विरोधापत्ती नाही, तर फरक स्पष्ट झाला आहे.
\end{enumerate}

आंशिक फलने, संगणनक्षम फलने, आंशिक संगणनक्षम फलने इत्यादींबद्दलची
चर्चा गोंधळात टाकू शकते. $\Nat$ पासून $\Nat$ मध्ये जाणाऱ्या सर्व
आंशिक फलनांचा संच हा वस्तूंचा मोठा संग्रह आहे. त्यांपैकी काही
सर्वत्र परिभाषित आहेत, काही आंशिक संगणनक्षम आहेत, काही सर्वत्र
परिभाषित आणि संगणनक्षम अशी दोन्ही आहेत, तर काही यांपैकी कोणतीच नाहीत.
आपण नुसते ``फलन'' म्हटले, तर सामान्यतः सर्वत्र परिभाषित फलन अभिप्रेत
असते, हे लक्षात ठेवा. म्हणून पुढील प्रकार मिळतात:
\begin{enumerate}
\item संगणनक्षम फलने
\item सर्वत्र परिभाषित नसलेली आंशिक संगणनक्षम फलने
\item संगणनक्षम नसलेली फलने
\item सर्वत्र परिभाषितही नसलेली आणि संगणनक्षमही नसलेली आंशिक फलने
\end{enumerate}
हा फरक समजण्यासाठी $\Nat$ पासून $\Nat$ मध्ये जाणारी सर्व आंशिक
फलने दर्शवणारा मोठा चौकोन काढा. त्यात सर्वत्र परिभाषित फलने
आणि आंशिक संगणनक्षम फलने दर्शवणारी, एकमेकांवर काही प्रमाणात
येणारी दोन क्षेत्रे दाखवा. अशा आकृतीत मिळणाऱ्या प्रत्येक
क्षेत्रातील एखाद्या वस्तूचे वर्णन करता येते का, हे पाहणे चांगला सराव आहे.











\section{संगणनक्षम संच}\label{cmp:thy:cps:sec}

संगणनक्षमतेची कल्पना संगणनक्षम फलनांपासून संगणनक्षम संचांपर्यंत
विस्तारता येते:

\begin{defn}
  $S$ हा नैसर्गिक संख्यांचा संच असू द्या. $S$ \emph{संगणनक्षम}
  असतो तेव्हा आणि केवळ तेव्हाच, जेव्हा त्याचे जातिबोधक फलन~$\Char{S}$
  संगणनक्षम असते. दुसऱ्या शब्दांत, $S$~संगणनक्षम असतो तेव्हा आणि केवळ
  तेव्हाच, जेव्हा पुढील फलन संगणनक्षम असते:
\[\Char{S}(x) =
\begin{cases}
1 & \text{जर $x \in S$ असेल तर} \\
0 & \text{अन्यथा}
\end{cases}
\]
त्याचप्रमाणे, संबंध $R(x_0, \dots, x_{k-1})$ संगणनक्षम असतो तेव्हा आणि
केवळ तेव्हाच, जेव्हा त्याचे जातिबोधक फलन संगणनक्षम असते.

संगणनक्षम संच आणि संबंध यांना \emph{निर्णेय} असेही म्हणतात.
\end{defn}

\begin{explain}
आता संगणनक्षमतेच्या अनेक कल्पना आहेत: आंशिक फलनांसाठी, फलनांसाठी
आणि संचांसाठी. त्यांची गल्लत करू नका!{}
\end{explain}











\section{संगणनक्षमपणे प्रगणनीय संच}\label{cmp:thy:ces:sec}

\begin{defn}
एखादा संच रिकामा असेल किंवा एखाद्या संगणनक्षम फलनाची व्याप्ती असेल,
तर तो \emph{संगणनक्षमपणे प्रगणनीय} असतो.
\end{defn}

\begin{history}
संगणनक्षमपणे प्रगणनीय संचांना \emph{पुनरावर्तीपणे प्रगणनीय}
असेही म्हणतात. ही मूळ संज्ञा आहे. आज ही दोन्ही नावे, तसेच
``c.e.'' आणि ``r.e.'' ही संक्षिप्त रूपे, सर्वसाधारणपणे वापरली जातात.
\end{history}

\begin{explain}
या व्याख्येचा अर्थ काय आणि ही संज्ञा योग्य का आहे, याचा विचार करा.
कल्पना अशी की $S$ हा संगणनक्षम फलन~$f$ ची व्याप्ती असेल, तर
\[S = \{ f(0), f(1), f(2), \dots \},
\]
म्हणून $f$ हे~$S$ च्या घटकांचे ``प्रगणन'' करते असे मानता येते.
व्याख्येनुसार $f$ वाढते फलन असणे आवश्यक नाही, म्हणजे प्रगणन
चढत्या क्रमाने असणे आवश्यक नाही. किंबहुना, $f$ एकास-एक असणेही
आवश्यक नाही; म्हणजे $f(0)$, $f(1)$, $f(2)$, \dots{} या~$S$
च्या प्रगणनात पुनरुक्ती चालते. उदाहरणार्थ, स्थिर फलन $f(x) = 0$
हे संच $\{ 0
\}$ चे प्रगणन करते.
\end{explain}

प्रत्येक संगणनक्षम संच संगणनक्षमपणे प्रगणनीय असतो. हे पाहण्यासाठी,
$S$~संगणनक्षम आहे असे समजा. $S$ रिकामा असेल, तर व्याख्येवरून
तो संगणनक्षमपणे प्रगणनीय आहे. अन्यथा $a$ हा $S$ मधील कोणताही
घटक घ्या. $f$ ची पुढीलप्रमाणे व्याख्या करा:
\[f(x) =
\begin{cases}
x & \text{जर $\Char{S}(x) = 1$ असेल तर} \\
a & \text{अन्यथा.}
\end{cases}
\]
मग $f$ हे संगणनक्षम फलन आहे आणि $S$ ही~$f$ ची व्याप्ती आहे.












\section{संगणनक्षमपणे
  प्रगणनीय संचांच्या समतुल्य व्याख्या}\label{cmp:thy:eqc:sec}


संगणनक्षमपणे प्रगणनीय असण्याचा अर्थ सांगणारी पुढील काही महत्त्वाची
समतुल्य विधाने आहेत.

\begin{thm}
\label{cmp:thy:eqc:thm:ce-equiv}
$S$ हा नैसर्गिक संख्यांचा संच असू द्या. मग पुढील विधाने समतुल्य आहेत:
\begin{enumerate}
\item\label{cmp:thy:eqc:case:ce} $S$ संगणनक्षमपणे प्रगणनीय आहे.
\item\label{cmp:thy:eqc:case:ran-pc} $S$ ही एखाद्या \emph{आंशिक}
  संगणनक्षम फलनाची व्याप्ती आहे.
\item\label{cmp:thy:eqc:case:ran-prim} $S$ रिकामा आहे किंवा एखाद्या आदिम
  पुनरावर्ती फलनाची व्याप्ती आहे.
\item\label{cmp:thy:eqc:case:ce-domain} $S$ हे एखाद्या आंशिक संगणनक्षम
  फलनाचे \emph{परिभाषाक्षेत्र} आहे.
\end{enumerate}
\end{thm}

\begin{explain}
पहिली तीन विधाने सांगतात की कोणत्याही रिकाम्या नसलेल्या
संगणनक्षमपणे प्रगणनीय संचाचे प्रगणन संगणनक्षम फलनाने, आंशिक
संगणनक्षम फलनाने किंवा आदिम पुनरावर्ती फलनाने करता येते.
चौथे विधान सांगते की $S$ संगणनक्षमपणे प्रगणनीय असेल, तर
एखाद्या निर्देशांक~$e$ साठी
\[S = \Setabs{x}{\cfind{e}(x) \fdefined}.
\]
दुसऱ्या शब्दांत, $S$ हा अशा आदानांचा संच आहे ज्यांवर
$\cfind{e}$ चे संगणन थांबते. म्हणून संगणनक्षमपणे प्रगणनीय
संचांना कधी कधी \emph{अर्धनिर्णेय} म्हणतात: संख्या त्या संचात
असेल, तर शेवटी ``होय'' असे उत्तर मिळते; पण ती संचात नसेल,
तर ``नाही'' असे उत्तर कधीच मिळत नाही!{}
\end{explain}

\begin{proof}
प्रत्येक आदिम पुनरावर्ती फलन संगणनक्षम असते आणि प्रत्येक
संगणनक्षम फलन आंशिक संगणनक्षमही असते. म्हणून
\ref{cmp:thy:eqc:case:ran-prim} वरून \ref{cmp:thy:eqc:case:ce} मिळते, आणि
\ref{cmp:thy:eqc:case:ce} वरून~\ref{cmp:thy:eqc:case:ran-pc} मिळते.
($S$ रिकामा असेल, तर $S$ कुठेही परिभाषित नसलेल्या आंशिक
संगणनक्षम फलनाची व्याप्ती असतो, हे लक्षात घ्या.)
\ref{cmp:thy:eqc:case:ran-pc} वरून \ref{cmp:thy:eqc:case:ran-prim} मिळते हे
दाखवले, तर पहिली तीन विधाने समतुल्य ठरतील.

म्हणून $S$ ही आंशिक संगणनक्षम फलन $\cfind{e}$ ची व्याप्ती आहे
असे समजा. $S$ रिकामा असेल, तर काम झाले. अन्यथा $a$ हा~$S$
मधील कोणताही घटक घ्या. क्लिनी यांच्या प्रमाण रूप प्रमेयावरून
पुढीलप्रमाणे लिहिता येते:
\[\cfind{e}(x) = U(\umin{s}{T(e, x, s)}).
\]
विशेषतः, $\cfind{e}(x) \fdefined$ असून त्याचे मूल्य $= y$ असते तेव्हा
आणि केवळ तेव्हाच, जेव्हा असा $s$ असतो की $T(e, x, s)$ आणि $U(s) = y$.
$f(z)$ ची व्याख्या पुढीलप्रमाणे करा:
\[f(z) = \begin{cases}
  U((z)_1) & \text{जर $T(e, (z)_0, (z)_1)$ असेल तर} \\
  a        & \text{अन्यथा.}
\end{cases}
\]
मग $f$ आदिम पुनरावर्ती आहे, कारण $T$ आणि $U$
तशी आहेत.

$S$ ही~$f$ ची व्याप्ती आहे हे दाखवायचे आहे; म्हणजे प्रत्येक
नैसर्गिक संख्या~$y$ साठी $y \in S$ तेव्हा आणि केवळ तेव्हाच, जेव्हा
ती~$f$ च्या व्याप्तीत असते. प्रथम $y \in S$ असे समजा. मग $y$ ही
$\cfind{e}$ च्या व्याप्तीत आहे. म्हणून काही $x$ आणि~$s$
साठी $T(e,x,s)$ लागू होतो आणि $U(s) = y$. त्यामुळे
$y = f(\tuple{x,s})$. उलट, $y$ ही~$f$ च्या व्याप्तीत
असेल, तर एकतर $y = a$ असेल, किंवा एखाद्या~$z$ साठी
$T(e,(z)_0,(z)_1)$ लागू होईल आणि $U((z)_1) = y$ असेल.
दुसऱ्या प्रकरणात $\cfind{e}((z)_0) \fdefined =
y$ असल्याने, दोन्ही प्रकारे $y$ ही~$S$ मध्येच असते.


($\cfind{e}(x) \fdefined = y$ या चिन्हांकनाचा अर्थ
``$\cfind{e}(x)$ परिभाषित आहे आणि $y$ च्या समान आहे'' असा आहे.
त्याऐवजी $\cfind{e}(x) =
y$ असेही लिहिता येईल; पण आपण आंशिक फलनाशी व्यवहार करत आहोत,
याची आठवण करून देण्यासाठी अतिरिक्त बाण कधी कधी उपयुक्त असतो.)

\ref{cmp:thy:eqc:thm:ce-equiv} ची सिद्धता पूर्ण करण्यासाठी
\ref{cmp:thy:eqc:case:ce} आणि~\ref{cmp:thy:eqc:case:ce-domain} समतुल्य आहेत
हे दाखवणे पुरेसे आहे. प्रथम \ref{cmp:thy:eqc:case:ce} वरून
\ref{cmp:thy:eqc:case:ce-domain} मिळते हे पाहू. संच रिकामा असेल,
तर कुठेही परिभाषित नसलेल्या आंशिक संगणनक्षम फलनाचे तो
परिभाषाक्षेत्र आहे. आता तो रिकामा नाही असे समजा. मग
$S$ ही संगणनक्षम फलन~$f$ ची व्याप्ती आहे, म्हणजे
\[S = \Setabs{y}{\text{अशा काही $x$ साठी, } f(x) = y}.
\]
पुढीलप्रमाणे घ्या:
\[g(y) = \umin{x}{(f(x) = y)}.
\]
मग $g$ आंशिक संगणनक्षम आहे आणि $g(y)$ तेव्हा आणि केवळ तेव्हाच
परिभाषित असते, जेव्हा काही~$x$ साठी $f(x) = y$ असते. दुसऱ्या शब्दांत, $g$
चे परिभाषाक्षेत्र ही~$f$ ची व्याप्ती आहे.



शेवटी \ref{cmp:thy:eqc:case:ce-domain} वरून~\ref{cmp:thy:eqc:case:ce} मिळते
हे दाखवू. $S$ हे आंशिक संगणनक्षम फलन~$\cfind{e}$ चे
परिभाषाक्षेत्र आहे असे समजा, म्हणजे
\[S = \Setabs{x}{\cfind{e}(x) \fdefined}.
\]
$S$ रिकामा असेल, तर काम झाले; अन्यथा $a$ हा~$S$
मधील कोणताही घटक घ्या. $f$ ची पुढीलप्रमाणे व्याख्या करा:
\[f(z) = \begin{cases}
(z)_0 & \text{जर $T(e,(z)_0,(z)_1)$ असेल तर} \\
a & \text{अन्यथा.}
\end{cases}
\]
मग वरच्याप्रमाणे, संख्या $x$ ही~$f$ च्या व्याप्तीत असते तेव्हा आणि
केवळ तेव्हाच, जेव्हा $\cfind{e}(x) \fdefined$; आणि ही अट तेव्हा आणि
केवळ तेव्हाच खरी असते, जेव्हा $x \in S$.

\end{proof}

\ref{cmp:thy:eqc:thm:ce-equiv} मधील विधान~\ref{cmp:thy:eqc:case:ce-domain}
संगणनक्षमपणे प्रगणनीय संचांचे प्रगणन करण्याचा सोयीचा मार्ग देते:
प्रत्येक~$e$ साठी $W_e$ हे $\cfind{e}$ चे परिभाषाक्षेत्र दर्शवू,
म्हणजे
\[W_e = \Setabs{x}{\cfind{e}(x) \fdefined}.
\]
मग $A$ हा कोणताही संगणनक्षमपणे प्रगणनीय संच असेल, तर
एखाद्या~$e$ साठी $A = W_e$.

पुढील प्रमेय संगणनक्षमपणे प्रगणनीय संचांचे आणखी एक वैशिष्ट्य
सांगते.

\begin{thm}
\label{cmp:thy:eqc:thm:exists-char}
$S$ हा संच संगणनक्षमपणे प्रगणनीय असतो तेव्हा आणि केवळ तेव्हाच, जेव्हा
असा संगणनक्षम संबंध $R(x,y)$ असतो की
\[S = \Setabs{ x }{ \lexists[y][R(x,y)] }.
\]
\end{thm}

\begin{proof}
प्रथम $S$ संगणनक्षमपणे प्रगणनीय आहे असे समजा. मग काही
$e$ साठी $S = W_e$. या~$e$ साठी $S$ असे लिहिता येते:
\[S = \Setabs{ x }{ \lexists[y][T(e, x, y)] }.
\]
उलट, $S = \Setabs{ x }{
  \lexists[y][R(x, y)] }$ असे समजा. $f$ ची पुढीलप्रमाणे व्याख्या करा:
\[f(x) \simeq \umin{y}{R(x, y)}.
\]
मग $f$ आंशिक संगणनक्षम आहे आणि $S$ हे~$f$ चे परिभाषाक्षेत्र आहे.
\end{proof}












\section{संगणनक्षम नसलेले संच अस्तित्वात आहेत}\label{cmp:thy:ncp:sec}


प्रत्येक संगणनक्षम संच संगणनक्षमपणे प्रगणनीय असतो, हे आपण वर पाहिले.
याचे उलटही खरे आहे का? पुढील प्रमेय दाखवते की सर्वसाधारणपणे तसे नाही.

\begin{thm}\label{cmp:thy:ncp:thm:K-0}
$K_0$ हा $\Setabs{\tuple{e, x}}{\cfind{e}(x) \fdefined}$ हा
संच असू द्या. मग $K_0$ संगणनक्षमपणे प्रगणनीय आहे, पण संगणनक्षम नाही.
\end{thm}

\begin{proof}
$K_0$ संगणनक्षमपणे प्रगणनीय आहे हे पाहण्यासाठी, तो पुढीलप्रमाणे
परिभाषित फलन~$f$ चे परिभाषाक्षेत्र आहे हे लक्षात घ्या:
\[f(z) = \umin{y}{(z = \tuple{(z)_0,(z)_1} \land T((z)_0, (z)_1, y))}.
\]
कारण $\cfind{e}(x)$ परिभाषित असेल, तर $f(\tuple{e, x})$
थांबणारी संगणनक्रमिका शोधते; $\cfind{e}(x)$ अपरिभाषित असेल,
तर $f(\tuple{e, x})$ देखील अपरिभाषित असते; आणि $z$ हा
जोडीचा संकेतांकसुद्धा नसेल, तर $f(z)$ देखील अपरिभाषित असते. केवळ लांबी
दोन येते एवढी कसोटी पुरेशी नाही: 42 या अवैध संकेतांकासाठी लांबीचे फलन
दोन देते, पण दोन्ही काढलेल्या घटकांना पुन्हा संकेतबद्ध केल्यास 6 मिळते.
म्हणून येथे जोडीचे अचूक पुनःसंकेतबद्धीकरण तपासले आहे.


$K_0$ संगणनक्षम नाही, ही बाब म्हणजे थांबण्याच्या समस्येचीच
अनिर्णेयता, \ref{cmp:thy:hlt:thm:halting-problem}.
\end{proof}

संच $K_0$ म्हणजे $\cfind{e}(x) \fdefined$ असा गुणधर्म असलेल्या
$\tuple{e,x}$ जोड्यांचा संच. म्हणजे $\tuple{e,x} \in K_0$
तेव्हा आणि केवळ तेव्हाच, जेव्हा $\cfind{e}$ हे आदान~$x$ वर परिभाषित असते
(संगणन थांबते). म्हणून त्याला ``थांबणारा संच'' असेही म्हणतात.
$K = \Setabs{e}{\cfind{e}(e) \fdefined}$ हा संच
``स्वतःच्या निर्देशांकावर थांबणारा संच'' आहे. तो अनिर्णेय संचाचे
मानक उदाहरण म्हणून वारंवार वापरला जातो.

\begin{thm}\label{cmp:thy:ncp:thm:K}
स्वतःच्या निर्देशांकावर थांबणारा संच $K = \Setabs{e}{\cfind{e}(e) \fdefined}$
हा संगणनक्षमपणे प्रगणनीय आहे, पण निर्णेय नाही.
\end{thm}

\begin{proof}
  $K$ निर्णेय आहे, म्हणजे त्याचे जातिबोधक फलन
  $\Char{K}$ संगणनक्षम आहे, असे समजा. पुढीलप्रमाणे घ्या:
  \[d(e) = \begin{cases}
  1 & \text{जर\/ $\Char{K}(e) = 0$ असेल तर}\\
  \fundefined & \text{अन्यथा.}
  \end{cases}
  \]
  $k$ हा~$d$ चा निर्देशांक असू द्या, म्हणजे $d \simeq \cfind{k}$.
  मग $d(k)
  \simeq \cfind{k}(k)$. पण $d(k)
  \fdefined$ तेव्हा आणि केवळ तेव्हाच, जेव्हा $\cfind{k}(k) \fundefined$, हे~$d$
  च्या व्याख्येवरून मिळते; त्यामुळे विरोध होतो.

  $K$ हे $f(x) = \umin{y}{T(x,x,y)}$ चे परिभाषाक्षेत्र आहे,
  म्हणून ते संगणनक्षमपणे प्रगणनीय आहे.
\end{proof}












\section{संगणनक्षमपणे
  प्रगणनीय संच संयोग आणि छेद यांखाली बंद असतात}\label{cmp:thy:clo:sec}

पुढील प्रमेय संगणनक्षमपणे प्रगणनीय संचांच्या काही संवरण गुणधर्मांबद्दल आहे.

\begin{thm}
$A$ आणि $B$ हे संगणनक्षमपणे प्रगणनीय असतील, तर
$A \cap B$ आणि $A \cup B$ देखील तसेच असतात.
\end{thm}

\begin{proof}
\ref{cmp:thy:eqc:thm:ce-equiv} मुळे संगणनक्षमपणे प्रगणनीय संचांची
वेगवेगळी वैशिष्ट्ये वापरता येतात. त्याचे उदाहरण म्हणून काही
वेगवेगळ्या सिद्धता देऊ.

पहिल्या दोन सिद्धतांत दोन्ही संच रिकामे नाहीत असे धरा. एखादा
रिकामा असेल, तर त्यांचा छेद रिकामा आणि संयोग दुसरा संच असतो;
निष्कर्ष त्या वेळी थेट मिळतो.

पहिल्या सिद्धतेसाठी $A$ चे प्रगणन संगणनक्षम फलन~$f$ करते
आणि $B$ चे प्रगणन संगणनक्षम फलन~$g$ करते असे समजा.
पुढीलप्रमाणे घ्या:
\begin{align*}
h(x) & = \umin{y}{(f(y) = x \lor g(y) = x)} \text{ आणि}\\
j(x) & = \umin{y}{(f((y)_0) = x \land g((y)_1) = x)}.
\end{align*}
मग $A \cup B$ हे $h$ चे परिभाषाक्षेत्र आहे आणि
$A \cap B$ हे~$j$ चे परिभाषाक्षेत्र आहे.

\begin{explain}
संगणनाच्या भाषेत येथे पुढील गोष्ट घडते: $A$ आणि $B$ यांचे
प्रगणन करणाऱ्या प्रक्रिया दिल्या असता, घटक $x$ हा $A \cup B$
मध्ये आहे का हे एका तरी प्रगणनात~$x$ शोधून अर्धनिर्णीत करता
येते. तसेच $x$ हा $A \cap B$ मध्ये आहे का हे दोन्ही
प्रगणनांत एकाच वेळी~$x$ शोधून अर्धनिर्णीत करता येते.
\end{explain}

दुसऱ्या सिद्धतेसाठी पुन्हा $A$ चे प्रगणन~$f$ करते आणि
$B$ चे प्रगणन~$g$ करते असे समजा. पुढीलप्रमाणे घ्या:
\[k(x) = \begin{cases}
f(x/2) & \text{जर $x$ सम असेल तर} \\
g((x-1)/2) & \text{जर $x$ विषम असेल तर.}
\end{cases}
\]
मग $k$ हे $A \cup B$ चे प्रगणन करते. कल्पना एवढीच की
$k$ हे $f$ आणि~$g$ यांच्या प्रगणनांचा आलटूनपालटून वापर करते.
$A \cap
B$ चे प्रगणन अधिक अवघड आहे. $A \cap B$ रिकामा असेल,
तर तो अर्थातच संगणनक्षमपणे प्रगणनीय आहे. अन्यथा $c$ हा
$A \cap B$ मधील कोणताही घटक घ्या आणि $l$ ची व्याख्या
पुढीलप्रमाणे करा:
\[l(x) = \begin{cases}
f((x)_0) & \text{जर $f((x)_0) = g((x)_1)$ असेल तर} \\
c & \text{अन्यथा.}
\end{cases}
\]
संगणनाच्या भाषेत, $l$ हे $f$ आणि $g$ यांच्या प्रगणनांतील
घटकांच्या जोड्यांमधून जाते आणि सापडलेली प्रत्येक जुळणी परत
करते; अन्यथा ते~$c$ परत करून प्रगणनात प्रगती करत नाही.

शेवटच्या सिद्धतेसाठी $A$ हे आंशिक फलन $m(x)$ चे
\emph{परिभाषाक्षेत्र} आणि $B$ हे आंशिक फलन $n(x)$ चे
परिभाषाक्षेत्र आहे असे समजा. मग $A \cap B$ हे आंशिक फलन
$m(x) +
n(x)$ चे परिभाषाक्षेत्र आहे.

\begin{explain}
संगणनाच्या भाषेत, $A$ हा ज्यांवर $m$ थांबते त्या
मूल्यांचा संच असेल आणि $B$ हा ज्यांवर $n$ थांबते त्या
मूल्यांचा संच असेल, तर $A \cap B$ हा दोन्ही प्रक्रिया
थांबतात अशा मूल्यांचा संच आहे.
\end{explain}

$A \cup B$ हा थांबणाऱ्या आदानांचा संच म्हणून व्यक्त करणे
अधिक कठीण आहे, कारण $m$ आणि $n$ यांचे समांतर अनुकरण
करावे लागते. $d$ हा $m$ चा निर्देशांक आणि $e$ हा $n$
चा निर्देशांक असू द्या; दुसऱ्या शब्दांत, $m =
\cfind{d}$ आणि $n = \cfind{e}$. मग $A \cup B$ हे
पुढील फलनाचे परिभाषाक्षेत्र आहे:
\[p(x) = \umin{y}{(T(d,x,y) \lor T(e,x,y))}.
\]
\begin{explain}
संगणनाच्या भाषेत, आदान $x$ वरील $p$ हे $m$ किंवा
$n$ यांपैकी एखाद्याचे थांबणारे संगणन शोधते आणि त्यांपैकी
एक सापडल्यास थांबते.
\end{explain}
\end{proof}











\section{पूरक संचाखाली संगणनक्षमपणे प्रगणनीय संच बंद नसतात}\label{cmp:thy:cmp:sec}

$A$ संगणनक्षमपणे प्रगणनीय आहे असे समजा. मग~$A$ चा पूरक संच,
$\Complement{A} = \Nat \setminus A$, हाही नेहमी संगणनक्षमपणे
प्रगणनीय असतो का? पुढील प्रमेय आणि अनुप्रमेय दाखवतात की उत्तर
``नाही'' असे आहे.

\begin{thm}
\label{cmp:thy:cmp:thm:ce-comp}
$A$ हा नैसर्गिक संख्यांचा कोणताही संच असू द्या. $A$ संगणनक्षम असतो
तेव्हा आणि केवळ तेव्हाच, जेव्हा $A$ आणि $\Complement{A}$ हे दोन्ही
संगणनक्षमपणे प्रगणनीय असतात.
\end{thm}

\begin{proof}
एका दिशेची सिद्धता सोपी आहे: $A$ संगणनक्षम असेल, तर
$\Complement{A}$ देखील संगणनक्षम आहे ($\Char{A} = 1 \tsub
\Char{\Complement{A}}$). म्हणून दोन्ही संच संगणनक्षमपणे प्रगणनीय आहेत.

उलट दिशेने, $A$ आणि~$\Complement{A}$ हे दोन्ही संगणनक्षमपणे
प्रगणनीय आहेत असे समजा. $A$ हे~$\cfind{d}$ चे परिभाषाक्षेत्र
आणि $\Complement{A}$ हे~$\cfind{e}$ चे परिभाषाक्षेत्र असू द्या.
$h$ ची व्याख्या अशी करा:
\[h(x) = \umin{s}{(T(d,x,s) \lor T(e,x,s))}.
\]
म्हणजे, आदान~$x$ वर $h$ हे~$\cfind{d}$ किंवा~$\cfind{e}$
यांचे थांबणारे संगणन शोधते. $x \in A$ असेल, तर पहिल्या
पर्यायात शोध यशस्वी होतो; $x \in
\Complement{A}$ असेल, तर दुसऱ्या पर्यायात. म्हणून $h$
हे सर्वत्र परिभाषित संगणनक्षम फलन आहे. आता प्रत्येक~$x$
साठी $x \in A$ तेव्हा आणि केवळ तेव्हाच, जेव्हा $T(d, x, h(x))$, म्हणजे
या आदानावर $\cfind{d}$ परिभाषित असते. $T(d, x, h(x))$
हा संगणनक्षम संबंध असल्यामुळे $A$ संगणनक्षम आहे.

\end{proof}

\begin{explain}
अनौपचारिक संगणकीय भाषेत ही गोष्ट अधिक सोपी आहे: $A$
निर्णीत करण्यासाठी आदान $x$ वर $\cfind{d}$ आणि $\cfind{e}$
यांच्या थांबणाऱ्या संगणनांचा शोध घ्या. त्यांपैकी एक नक्की थांबते;
$\cfind{d}$ थांबले, तर $x$ हा~$A$ मध्ये आहे; अन्यथा $x$
हा~$\Complement{A}$ मध्ये आहे.
\end{explain}

\begin{cor}
\label{cmp:thy:cmp:cor:comp-k}
$\Complement{K_0}$ संगणनक्षमपणे प्रगणनीय नाही.
\end{cor}

\begin{proof}
$K_0$ संगणनक्षमपणे प्रगणनीय आहे, पण संगणनक्षम नाही, हे
आपल्याला माहीत आहे. $\Complement{K_0}$ संगणनक्षमपणे प्रगणनीय
असता, तर \ref{cmp:thy:cmp:thm:ce-comp} मुळे $K_0$ संगणनक्षम ठरला असता;
हे \ref{cmp:thy:ncp:thm:K-0} ला विरोधी आहे.
\end{proof}











\section{न्यूनीकरणक्षमता}\label{cmp:thy:red:sec}

\begin{explain}
आता आपल्याला माहीत आहे की $K_0$ हा किमान एक संच संगणनक्षमपणे
प्रगणनीय आहे, पण संगणनक्षम नाही. असे इतरही संच असतील हे उघड
आहे. प्रत्येक वेळी मूलभूत व्याख्यांपासून पुन्हा सुरुवात न करता
इतर संचांचे असे गुणधर्म सिद्ध करण्याची प्रभावी पद्धत
न्यूनीकरणक्षमता देते.

सर्वसाधारणपणे, $A$ या संचाचे~$B$ कडे ``न्यूनीकरण'' म्हणजे
घटक हे~$B$ मध्ये आहेत की नाहीत या प्रश्नांची उत्तरे
रूपांतरित करून घटक हे~$A$ मध्ये आहेत की नाहीत या प्रश्नांची उत्तरे
मिळवण्याची पद्धत. येथे आपण ``अनेक-एक न्यूनीकरणक्षमता''
विचारात घेऊ; वेगवेगळे गुणधर्म असणाऱ्या न्यूनीकरणक्षमतेच्या इतरही
अनेक संकल्पना आहेत. संगणकीय गुंतागुंतीच्या अभ्यासातही अशा
संकल्पना महत्त्वाच्या असतात; तेथे कार्यक्षमतेचाही विचार करावा
लागतो. उदाहरणार्थ, एखादा संच NP मध्ये असेल आणि प्रत्येक
NP समस्या त्याच्याकडे न्यूनीत करता येत असेल, तर त्याला
``NP-पूर्ण'' म्हणतात. येथे पुढे येणाऱ्या न्यूनीकरणासारखेच
न्यूनीकरण तेथे वापरतात, पण त्यासाठी ते बहुपदी वेळेत संगणित
करता येण्याची अतिरिक्त अट असते.

न्यूनीकरणाची कल्पना आपण आधीच अप्रत्यक्षपणे वापरली आहे.
$K$ हा संच पुढीलप्रमाणे ठरवा:
\[K = \Setabs{x}{\cfind{x}(x) \fdefined},
\]
म्हणजे $K = \Setabs{x}{x \in W_x}$. थांबण्याची समस्या
अनिर्णेय असल्याचा आपला पुरावा (\ref{cmp:thy:hlt:thm:halting-problem})
सर्वांत थेटपणे $K$ संगणनक्षम नाही हे दाखवतो. $K_0$ हा
पुढील संच आहे, हे आठवा:
\[K_0 = \Setabs{\tuple{e, x}}{\cfind{e}(x) \fdefined },
\]
म्हणजे $K_0 = \Setabs{\tuple{e,x}}{x \in W_e}$.

$K$ च्या असंगणनक्षमतेच्या कोणत्याही पुराव्याचा विस्तार करून
$K_0$ ची असंगणनक्षमता दाखवणे सोपे आहे: $K_0$ संगणनक्षम असता,
तर घटक~$x$ हा $K$ मध्ये आहे की नाही हे फक्त
$\tuple{x, x}$ ही जोडी $K_0$ मध्ये आहे का ते विचारून
निर्णीत करता आले असते. $x$ ला $\tuple{x, x}$ कडे नेणारे
फलन~$f$ हे $K$ चे $K_0$ कडे \emph{न्यूनीकरण} करण्याचे
उदाहरण आहे.
\end{explain}

\begin{defn}
$A$ आणि $B$ हे नैसर्गिक संख्यांचे संच असू द्या.
संगणनक्षम फलन~$f\colon \Nat \to \Nat$ हे $A$ चे~$B$
कडे \emph{अनेक-एक न्यूनीकरण} असते जर आणि फक्त जर
प्रत्येक नैसर्गिक संख्या~$x$ साठी
\[x \in A \quad \text{तेव्हा आणि केवळ तेव्हाच} \quad f(x) \in B.
\]
असे न्यूनीकरण~$f$ अस्तित्वात असल्यास, $A$ हे~$B$ कडे
\emph{अनेक-एक न्यूनीत करता येते} असे म्हणतो; हे
$A \leq_m B$ असे लिहितो. $A$ चे~$B$ कडे आणि उलटही
अनेक-एक न्यूनीकरण शक्य असल्यास, $A$ आणि $B$ हे
\emph{अनेक-एक सममूल्य} आहेत असे म्हणतो; हे
$A \equiv_m B$ असे लिहितो.
\end{defn}

वरील व्याख्येतील फलन $f$ एकास-एक असेल, तर $A$ हे~$B$
कडे \emph{एकास-एक न्यूनीत करता येते} असे म्हणतात.
पुढे दिलेली बहुतेक न्यूनीकरणे ही अधिक मजबूत अट पूर्ण करतात;
पण आपण ती बाब वापरणार नाही.

\begin{digress}
एकास-एक न्यूनीकरणक्षमता ही अनेक-एक न्यूनीकरणक्षमतेपेक्षा खरोखर
अधिक मजबूत अट आहे; हे खरे आहे, पण अजिबात उघड नाही. म्हणजे,
असे अनंत संच $A$ आणि~$B$ आहेत की $A$ हे~$B$ कडे अनेक-एक
न्यूनीत करता येते, पण~$B$ कडे एकास-एक न्यूनीत करता येत नाही.
\end{digress}











\section{न्यूनीकरणक्षमतेचे गुणधर्म}\label{cmp:thy:ppr:sec}

$A$ चे~$B$ कडे न्यूनीकरण होत असेल, तर आपण $A \leq_m B$
असे लिहितो. या चिन्हांकनावरून असे सुचते की $A \leq_m B$
असल्यास $A$ हे~$B$ पेक्षा ``अधिक कठीण नाही'' आणि $B$
हे~$A$ इतके किंवा त्याहून ``अधिक कठीण'' आहे. तसेच नेहमीचे
संख्यांवरील $\le$ जसे क्रम लावते, तसे $\le_m$ संचांना
(किंवा निर्णय समस्यांना) त्यांच्या गुंतागुंतीनुसार क्रम लावते.
पुढील दोन विधाने ही अंतःप्रेरणा समर्थित करतात. पहिले विधान
$\le_m$ संक्रमक आहे असे सांगते.

\begin{prop}
\label{cmp:thy:ppr:prop:trans-red}
$A \leq_m B$ आणि $B \leq_m C$ असल्यास $A \leq_m C$.
\end{prop}

\begin{proof}
$A$ चे~$B$ कडे न्यूनीकरण आणि $B$ चे~$C$ कडे न्यूनीकरण
यांचे संयोजन केल्यास $A$ चे~$C$ कडे न्यूनीकरण मिळते.
\end{proof}

\begin{prob}
$f$ आणि~$g$ ही अनुक्रमे $A$ चे~$B$ कडे आणि $B$ चे~$C$
कडे अनेक-एक न्यूनीकरणे असतील, तर $\comp{f}{g}$ हे $A$
चे~$C$ कडे अनेक-एक न्यूनीकरण आहे, हे दाखवून
\ref{cmp:thy:ppr:prop:trans-red} सिद्ध करा.
\end{prob}

\begin{prop}
\label{cmp:thy:ppr:prop:reduce}
$A$ आणि $B$ हे कोणतेही संच असू द्या आणि $A \leq_m B$
असे समजा.
\begin{enumerate}
\item $B$ संगणनक्षमपणे प्रगणनीय असेल, तर $A$ देखील तसाच आहे.
\item $B$ संगणनक्षम असेल, तर $A$ देखील तसाच आहे.
\end{enumerate}
\end{prop}

\begin{proof}
$f$ हे $A$ चे~$B$ कडे अनेक-एक न्यूनीकरण असू द्या.
पहिल्या विधानासाठी, $B$ हे आंशिक फलन~$g$ चे परिभाषाक्षेत्र
असेल, तर $A$ हे~$\comp{f}{g}$ चे परिभाषाक्षेत्र आहे हे
तपासा:
\begin{align*}
x \in A & \text{ तेव्हा आणि केवळ तेव्हाच } f(x) \in B \\
& \text{ तेव्हा आणि केवळ तेव्हाच }  g(f(x)) \fdefined.
\end{align*}

दुसऱ्या विधानासाठी, $B$ संगणनक्षम असेल, तर $B$
आणि~$\Complement{B}$ हे दोन्ही संगणनक्षमपणे प्रगणनीय
असतात (\ref{cmp:thy:cmp:thm:ce-comp}), हे आठवा. $f$ हे
$\Complement{A}$ चे $\Complement{B}$ कडेही अनेक-एक
न्यूनीकरण आहे, हे तपासणे कठीण नाही. त्यामुळे या सिद्धतेच्या
पहिल्या भागानुसार $A$ आणि~$\Complement{A}$ हे
संगणनक्षमपणे प्रगणनीय आहेत. म्हणून \ref{cmp:thy:cmp:thm:ce-comp} मुळे
$A$ देखील संगणनक्षम आहे. (पर्यायाने,
$\Char{A} = \comp{f}{\Char{B}}$ हे तपासा; मग
$\Char{B}$ संगणनक्षम असेल, तर~$\Char{A}$ देखील आहे.)
\end{proof}

\begin{prob}
$f$ हे $A$ चे~$B$ कडे अनेक-एक न्यूनीकरण आहे असे
समजा. मग $f$ हे $\Complement{A}$ चे $\Complement{B}$
कडेही अनेक-एक न्यूनीकरण आहे हे दाखवा.
\end{prob}

\begin{prob}
$f\colon \Nat \to \Nat$ हे $A$ चे~$B$ कडे अनेक-एक न्यूनीकरण असेल, तर
$\Char{A}
= \comp{f}{\Char{B}}$ हे दाखवा.

\end{prob}

\begin{digress}
\emph{ट्यूरिंग न्यूनीकरणक्षमता} ही न्यूनीकरणक्षमतेची अधिक
सर्वसाधारण संकल्पना इतर संदर्भांत, विशेषतः अनिर्णेयतेचे निष्कर्ष
सिद्ध करण्यासाठी, उपयोगी आहे. \ref{cmp:thy:cmp:cor:comp-k} नुसार
$K_0$ चा पूरक संच~$K_0$ कडे न्यूनीत करता येत नाही, कारण
तो संगणनक्षमपणे प्रगणनीय नाही. पण अंतःप्रेरणेने, $K_0$
विषयीच्या प्रश्नांची उत्तरे माहीत असतील, तर त्याच्या पूरकाविषयीची
उत्तरेही माहीत असतील. संगणनक्षम प्रक्रिया~$B$ विषयी प्रश्न
विचारून~$A$ विषयीच्या प्रश्नांची उत्तरे ठरवू शकत असेल, तर
$A$ हे~$B$ कडे ट्यूरिंग न्यूनीत करता येते असे म्हणतात.
अनेक-एक न्यूनीकरणक्षमतेपेक्षा ही अधिक व्यापक संकल्पना आहे:
अनेक-एक न्यूनीकरणात (1)~$B$ विषयी एकच प्रश्न विचारता येतो,
आणि (2)~``होय'' या उत्तराचे $A$ विषयीच्या प्रश्नासाठी
``होय'' असेच, तसेच ``नाही'' चे ``नाही'' असेच रूपांतर व्हावे
लागते. $A$ हे~$B$ कडे ट्यूरिंग न्यूनीत करता येत असेल
आणि $B$ संगणनक्षम असेल, तर $A$ देखील संगणनक्षम असतो;
परंतु आपण पाहिल्याप्रमाणे, संगणनक्षमपणे प्रगणनीय असण्याबाबत
असेच विधान खरे ठरत नाही.

आपण चर्चिलेल्या न्यूनीकरणक्षमतेच्या वेगवेगळ्या संकल्पनांचा
आणि त्यांतील फरकांचा विचार करा. मात्र या प्रकरणात आपण फक्त
अनेक-एक न्यूनीकरणक्षमता वापरू. आधीच्या परिच्छेदातील दोन्ही
प्रकारांना संगणकीय गुंतागुंतीत समांतर संकल्पना आहेत; त्यांत
ट्यूरिंग यंत्रे बहुपदी वेळेत चालण्याची अतिरिक्त अट असते.
अनेक-एक न्यूनीकरणक्षमतेचे गुंतागुंत-संबंधी रूप
\emph{कार्प न्यूनीकरणक्षमता}, तर ट्यूरिंग न्यूनीकरणक्षमतेचे
तसे रूप \emph{कुक न्यूनीकरणक्षमता} म्हणून ओळखले जाते.
\end{digress}











\section{संपूर्ण संगणनक्षमपणे प्रगणनीय संच}\label{cmp:thy:cce:sec}

\begin{defn}
संच $A$ हा (अनेक-एक न्यूनीकरणक्षमतेखाली) \emph{संपूर्ण
संगणनक्षमपणे प्रगणनीय संच} असतो, जर
\begin{enumerate}
\item $A$ संगणनक्षमपणे प्रगणनीय असेल, आणि
\item इतर कोणत्याही संगणनक्षमपणे प्रगणनीय संच $B$ साठी $B \leq_m A$ असेल.
\end{enumerate}
\end{defn}

म्हणजे, संपूर्ण संगणनक्षमपणे प्रगणनीय संच हे अशा संचांमधील
शक्य तितके ``कठीण'' संच असतात. त्यांच्याद्वारे
\emph{कोणत्याही} संगणनक्षमपणे प्रगणनीय संचाविषयीच्या
प्रश्नांची उत्तरे मिळवता येतात.

\begin{thm}
$K$, $K_0$, आणि $K_1$ हे तिन्ही संपूर्ण संगणनक्षमपणे
प्रगणनीय संच आहेत.
\end{thm}

\begin{proof}
$K_0$ संपूर्ण आहे हे पाहण्यासाठी $B$ हा कोणताही
संगणनक्षमपणे प्रगणनीय संच असू द्या. मग काही निर्देशांक $e$
साठी
\[B = W_e = \Setabs{x}{\cfind{e}(x) \fdefined}.
\]
$f(x) = \tuple{e, x}$ हे फलन $f$ असू द्या. मग प्रत्येक
नैसर्गिक संख्या $x$ साठी $x \in B$ तेव्हा आणि केवळ तेव्हाच, जेव्हा
$f(x) \in K_0$. म्हणजे $f$ हे $B$ चे~$K_0$ कडे
न्यूनीकरण करते.

$K_1$ संपूर्ण आहे हे पाहण्यासाठी, \ref{cmp:thy:k1:prop:k1}
च्या सिद्धतेत आपण $K_0$ चे त्याच्याकडे न्यूनीकरण केले
आहे, हे लक्षात घ्या. म्हणून \ref{cmp:thy:ppr:prop:trans-red}
नुसार कोणताही संगणनक्षमपणे प्रगणनीय संच~$K_1$ कडेही
न्यूनीत करता येतो.

$K_0$ चे~$K$ कडेही अशाच प्रकारे न्यूनीकरण करता येते.

\end{proof}

\begin{prob}
$K$ चे $K_0$ कडे न्यूनीकरण द्या.
\end{prob}

\begin{digress}
आतापर्यंत आपण पाहिलेल्या संगणनक्षमपणे प्रगणनीय संचांची सर्व
उदाहरणे एकतर संगणनक्षम किंवा संपूर्ण आहेत. हे विचित्र वाटले
पाहिजे!{} संगणनक्षमही नाहीत आणि संपूर्णही नाहीत, अशी
संगणनक्षमपणे प्रगणनीय संचांची उदाहरणे आहेत का? उत्तर होय
आहे; मात्र 1950 च्या दशकाच्या मध्यापर्यंत हे सिद्ध झाले
नव्हते. फ्रीडबर्ग आणि मुछ्निक यांनी स्वतंत्रपणे ते सिद्ध केले.
\end{digress}











\section{न्यूनीकरणक्षमतेचे एक उदाहरण}\label{cmp:thy:k1:sec}

\ref{cmp:thy:ppr:prop:reduce} चा एक उपयोग पाहू.

\begin{prop}
\label{cmp:thy:k1:prop:k1}
पुढीलप्रमाणे घ्या:
\[K_1 = \Setabs{e}{\cfind{e}(0) \fdefined}.
\]
मग $K_1$ संगणनक्षमपणे प्रगणनीय आहे, पण संगणनक्षम नाही.
\end{prop}

\begin{proof}
$K_1 = \Setabs{e}{\lexists[s][T(e,0,s)]}$ असल्याने,
\ref{cmp:thy:eqc:thm:exists-char} नुसार $K_1$ संगणनक्षमपणे
प्रगणनीय आहे.

$K_1$ संगणनक्षम नाही हे दाखवण्यासाठी $K_0$ चे
त्याच्याकडे न्यूनीकरण होते हे दाखवू.

\begin{explain}
हे थोडे गुंतागुंतीचे आहे, कारण $K_1$ वापरून आपण फक्त
विशिष्ट आदान $0$ पासून सुरू होणाऱ्या संगणनांबद्दल प्रश्न
विचारू शकतो. अशा प्रश्नांची उत्तरे देणारा हुशार मित्र
तुमच्याकडे आहे असे समजा (अशा मित्रांना ``ओरॅकल'' म्हणतात).
आता कोणीतरी तुम्हाला $\tuple{e,
  x}$ ही जोडी $K_0$ मध्ये आहे का, म्हणजे यंत्र $e$ हे
आदान $x$ वर थांबते का, असे विचारते. तुम्ही दुसरे यंत्र,
$e_x$, तयार करू शकता. ते \emph{कोणतेही} आदान मिळाले तरी
त्याकडे दुर्लक्ष करते आणि त्याऐवजी~$e$ हे यंत्र आदान
$x$ वर चालवते. मग यंत्र $e$ आदान $x$ वर थांबते का हा
प्रश्न, यंत्र $e_x$ आदान $0$ वर (किंवा इतर कोणत्याही
आदानावर) थांबते का या प्रश्नाशी समतुल्य आहे. आता नवीन
यंत्र $e_x$ आदान $0$ वर थांबते का हे मित्राला विचारा;
बदललेल्या प्रश्नाचे त्याचे उत्तर मूळ प्रश्नाचे उत्तरही देते.
यातून $K_0$ चे~$K_1$ कडे हवे ते न्यूनीकरण मिळते.
\end{explain}

सार्वत्रिक आंशिक संगणनक्षम फलन वापरून $f$ हे पुढीलप्रमाणे
व्याख्यायित त्रिस्थानी फलन असू द्या:
\[f(x,y,z) \simeq \cfind{x}(y).
\]
$f$ आपल्या तिसऱ्या आदानाकडे पूर्णपणे दुर्लक्ष करते, हे
लक्षात घ्या. $f = \cfind{e}[3]$ होईल असा निर्देशांक $e$
निवडा; म्हणजे
\[\cfind{e}[3](x,y,z) \simeq \cfind{x}(y).
\]
$s$-$m$-$n$ प्रमेयानुसार असे फलन $s(e,x,y)$ अस्तित्वात
आहे की प्रत्येक $z$ साठी
\begin{align*}
\cfind{s(e,x,y)}(z) & \simeq \cfind{e}[3](x,y,z) \\
& \simeq \cfind{x}(y).
\end{align*}

\begin{explain}
वरील अनौपचारिक युक्तिवादाच्या दृष्टीने $s(e,x,y)$ हा
अशा यंत्राचा निर्देशांक आहे की ते कोणतेही आदान $z$
मिळाले तरी त्याकडे दुर्लक्ष करते आणि $\cfind{x}(y)$
संगणित करते.
\end{explain}

विशेषतः, आपल्याला
\[\cfind{s(e,x,y)}(0) \fdefined \quad \text{तेव्हा आणि केवळ तेव्हाच} \quad
\cfind{x}(y) \fdefined.
\]
म्हणजे $\tuple{x, y} \in K_0$ तेव्हा आणि केवळ तेव्हाच, जेव्हा $s(e,x,y) \in
K_1$. कुठेही परिभाषित नसलेल्या आंशिक संगणनक्षम फलनाचा निर्देशांक $u$
निश्चित करा. मग पुढीलप्रमाणे व्याख्यायित फलन $g$:
\[g(w) = \begin{cases}
s(e,(w)_0,(w)_1) & \text{जर $w = \tuple{(w)_0,(w)_1}$ असेल तर} \\
u & \text{अन्यथा.}
\end{cases}
\]
हे $K_0$ चे~$K_1$ कडे न्यूनीकरण आहे. जोडीचा वैध संकेतांक असल्याची कसोटी
आदिम पुनरावर्ती आहे. अवैध संकेतांक $K_0$ मध्ये नसतो आणि त्याला $K_1$
मध्ये नसलेला निर्देशांक $u$ मिळतो; वैध जोड्यांसाठी वरील सममूल्यत्व लागू होते.

\end{proof}











\section{सर्वत्र परिभाषितता अनिर्णेय आहे}\label{cmp:thy:tot:sec}

एखादी गोष्ट असंगणनक्षम आहे हे दाखवण्यासाठी $s$-$m$-$n$
प्रमेय वापरण्याचे आणखी एक उदाहरण पाहू. $\fn{Tot}$ हा
सर्वत्र परिभाषित संगणनक्षम फलनांच्या निर्देशांकांचा संच असू
द्या, म्हणजे
\[\fn{Tot} = \Setabs{x}{\text{प्रत्येक $y$ साठी, $\cfind{x}(y)\fdefined$}}.
\]

\begin{prop}
\label{cmp:thy:tot:prop:total}
$\fn{Tot}$ संगणनक्षम नाही.
\end{prop}

\begin{proof}
$\fn{Tot}$ संगणनक्षम नाही हे पाहण्यासाठी $K$ चे त्याच्याकडे
न्यूनीकरण होते हे दाखवणे पुरेसे आहे. $h(x,y)$ ची व्याख्या
पुढीलप्रमाणे करा:
\[h(x,y) \simeq
\begin{cases}
0 & \text{जर $x \in K$ असेल तर} \\
\fundefined & \text{अन्यथा}
\end{cases}
\]
$h(x,y)$ हे $y$ वर अजिबात अवलंबून नाही, हे लक्षात घ्या.
$h$ हे आंशिक संगणनक्षम आहे हे पाहणे कठीण नसावे: आदान
$x, y$ वर~$h$ संगणित करताना आधी~$\cfind{x}$ चे
आदान~$x$ वरील अनुकरण करतो; हे संगणन थांबले, तर
$h(x,y)$ हे $0$ परत करून थांबते. म्हणून $h(x,y)$ हे
$\Zero(\umin{s}{T(x,x,s)})$ इतकेच आहे; येथे $\Zero$
हे स्थिर शून्य फलन आहे.

$s$-$m$-$n$ प्रमेयानुसार असे आदिम पुनरावर्ती फलन~$k(x)$
अस्तित्वात आहे की प्रत्येक $x$ आणि~$y$ साठी
\[\cfind{k(x)}(y) =
\begin{cases}
0 & \text{जर $x \in K$ असेल तर} \\
\fundefined & \text{अन्यथा}
\end{cases}
\]
म्हणून $x \in K$ असल्यास $\cfind{k(x)}$ सर्वत्र परिभाषित
असते; अन्यथा ते अपरिभाषित असते. अशा प्रकारे $k$ हे $K$
चे~$\fn{Tot}$ कडे न्यूनीकरण आहे.
\end{proof}

\begin{digress}
खरे तर $\fn{Tot}$ संगणनक्षमपणे प्रगणनीयही नाही; त्याची
गुंतागुंत ``अंकगणितीय पदानुक्रमात'' अधिक वरच्या स्तरावर
येते. मात्र येथे आपण या अधिक मजबूत निष्कर्षाचा विचार करणार नाही.
\end{digress}











\section{राइसचे प्रमेय}\label{cmp:thy:rce:sec}

विचार केल्यास लक्षात येईल की \ref{cmp:thy:tot:prop:total}
च्या सिद्धतेत $\fn{Tot}$ चे विशिष्ट गुणधर्म वापरलेले नाहीत.
$h(x,y)$ हे स्थिर फलन $j(y) = 0$ प्रमाणे वागते तेव्हा आणि केवळ
तेव्हाच, जेव्हा $x$ हा $K$ मध्ये असतो, अशी रचना आपण केली; पण त्या
परिस्थितीत ते इतर कोणत्याही आंशिक संगणनक्षम फलनाप्रमाणे
वागेल अशीही रचना करता आली असती. या निरीक्षणामुळे अधिक
सर्वसाधारण प्रमेय मांडता येते. साधारणपणे ते सांगते की
संगणनक्षम फलनांचा कोणताही अतुच्छ गुणधर्म निर्णेय नसतो.

$\cfind{0}$, $\cfind{1}$, $\cfind{2}$,~\dots हे आंशिक
संगणनक्षम फलनांचे आपले प्रमाण प्रगणन आहे, हे लक्षात ठेवा.

\begin{thm}[राइसचे प्रमेय]
  $C$ हा आंशिक संगणनक्षम फलनांचा कोणताही संच असू द्या आणि $A =
  \Setabs{n}{\cfind{n} \in C}$ असू द्या. $A$ संगणनक्षम असेल,
  तर $C$ रिकामा असतो किंवा $C$ हा सर्व आंशिक संगणनक्षम
  फलनांचा संच असतो.
\end{thm}

{\em निर्देशांकसंच} म्हणजे असा संच $A$ की $n$ आणि $m$ हे
एकच फलन ``संगणित'' करणारे निर्देशांक असतील, तर $n$ आणि
$m$ हे दोन्ही $A$ मध्ये असतात किंवा दोन्ही नसतात. प्रमेयातील संच~$A$
हा गुणधर्म पूर्ण करतो हे पाहणे कठीण नाही. उलट, $A$
निर्देशांकसंच असेल आणि $C$ हा त्या निर्देशांकांनी संगणित
होणाऱ्या फलनांचा संच असेल, तर $A = \Setabs{n}{\cfind{n} \in C}$.

\begin{explain}
या संज्ञेत राइसचे प्रमेय म्हणजे कोणताही अतुच्छ निर्देशांकसंच
निर्णेय नसतो असे विधान. प्रमेयाचा अर्थ समजण्यासाठी
\emph{कार्यक्रम} (उदा., तुमच्या आवडत्या प्रोग्रामिंग भाषेतील)
आणि ते संगणित करत असलेली फलने यांतील फरक लक्षात घ्या.
कार्यक्रम हे विन्यासात्मक वस्तू आहेत, आणि त्यांच्याविषयीचे
काही प्रश्न नक्कीच संगणनक्षम आहेत: या कार्यक्रमात 150
पेक्षा अधिक चिन्हे आहेत का? त्यात 22 पेक्षा अधिक ओळी
आहेत का? त्यात ``while'' विधान आहे का? ``print'' विधानाचे
आदान म्हणून ``hello world'' ही चिन्हमाला कुठे येते का?
राइसचे प्रमेय सांगते की कार्यक्रमाच्या \emph{वर्तनाविषयीचा}
कोणताही अतुच्छ प्रश्न संगणनक्षम नसतो. उदाहरणार्थ: हा
कार्यक्रम आदान $0$ वर थांबतो का? तो कधी तरी थांबतो का?
तो कधी सम संख्या परत करतो का?
\end{explain}

\begin{proof}[राइसच्या प्रमेयाची सिद्धता]
$C$ रिकामा नाही आणि सर्व आंशिक संगणनक्षम फलनांचा संचही
नाही असे समजा. $A$ हा~$C$ मधील फलनांच्या निर्देशांकांचा
संच असू द्या. $A$ संगणनक्षम असता, तर थांबण्याची समस्या
सोडवता आली असती, हे आपण दाखवू; म्हणून $A$ संगणनक्षम नाही.

व्यापकता न गमावता, कुठेही परिभाषित नसलेले फलन $f$ हे
$C$ मध्ये नाही असे धरू शकतो (तसे नसेल, तर पुढील
युक्तिवादात $C$ आणि त्याचा पूरक संच यांची अदलाबदल करा).
$g$ हे~$C$ मधील कोणतेही फलन असू द्या. कल्पना अशी आहे:
$A$ निर्णीत करता आले, तर~$f$ संगणित करणारे निर्देशांक
आणि~$g$ संगणित करणारे निर्देशांक यांतील फरक ओळखता
येईल; आणि ती क्षमता वापरून थांबण्याची समस्या सोडवता येईल.

कसे ते पाहू. सार्वत्रिक आंशिक संगणनक्षम फलने वापरून
पुढील फलन व्याख्यायित करता येते:
\[h(x,y) \simeq
\begin{cases}
\text{अपरिभाषित} & \text{जर $\cfind{x}(x) \fundefined$ असेल तर} \\
g(y) & \text{अन्यथा.}
\end{cases}
\]
$h$ संगणित करण्यासाठी आधी $\cfind{x}(x)$ संगणित करण्याचा
प्रयत्न करतो; ते संगणन थांबले, तर पुढे~$g(y)$ संगणित करतो;
आणि {\em हे} संगणन थांबले, तर त्याचे मूल्य परत करतो.
अधिक औपचारिकपणे,
\[h(x,y) \simeq \Proj{2}{0}(g(y),\fn{Un}(x,x)).
\]
असे लिहिता येते. येथे $\Proj{2}{0}(z_0, z_1) = z_0$ हे
$0$ वे आदान परत करणारे, संगणनक्षम $2$-स्थानी प्रक्षेपण
फलन आहे.

मग $h$ हे आंशिक संगणनक्षम फलनांचे संयोजन आहे. उजवी बाजू परिभाषित
असून तिचे मूल्य~$g(y)$ च्या समान असते तेव्हा आणि केवळ तेव्हाच, जेव्हा
$\fn{Un}(x,x)$ आणि $g(y)$ ही दोन्ही परिभाषित असतात.

निश्चित~$x$ साठी, $\cfind{x}(x)$ अपरिभाषित असेल, तर
प्रत्येक~$y$ साठी $h(x,y)$ अपरिभाषित असते; आणि
$\cfind{x}(x)$ परिभाषित असेल, तर $h(x,y) \simeq g(y)$.
म्हणून निश्चित~$x$ साठी $h(x,y)$ हे एकतर $f$ प्रमाणे
किंवा $g$ प्रमाणे वागते. $\cfind{x}(x)$ परिभाषित आहे
की नाही हे ठरवणे म्हणजे या दोन स्थितींपैकी कोणती आहे
हे ठरवणे. पण हेच $h_x(y) \simeq h(x,y)$ हे~$C$
मध्ये आहे की नाही हे ठरवण्यासारखे आहे; $A$ संगणनक्षम
असता, तर तसे करता आले असते.

अधिक औपचारिकपणे, $h$ आंशिक संगणनक्षम असल्याने तो काही
निर्देशांक~$e$ साठी $\cfind{e}$ या फलनाइतकाच आहे.
$s$-$m$-$n$ प्रमेयानुसार असे आदिम पुनरावर्ती फलन~$s$
आहे की प्रत्येक~$x$ साठी $\cfind{s(e,x)}(y) = h_x(y)$.
आता प्रत्येक $x$ साठी $\cfind{x}(x) \fdefined$ असेल,
तर $\cfind{s(e,x)}$ हे~$g$ इतकेच फलन आहे; म्हणून
$s(e,x)$ हा $A$ मध्ये आहे. उलट, $\cfind{x}(x)
\uparrow$ असेल, तर $\cfind{s(e,x)}$ हे~$f$ इतकेच फलन
आहे; म्हणून $s(e,x)$ हा~$A$ मध्ये नाही. दुसऱ्या शब्दांत,
प्रत्येक~$x$ साठी $x
\in K$ तेव्हा आणि केवळ तेव्हाच, जेव्हा $s(e,x) \in A$. $A$ संगणनक्षम
असता, तर~$K$ देखील संगणनक्षम ठरला असता; हा विरोधाभास
आहे. म्हणून $A$ संगणनक्षम नाही.
\end{proof}

राइसचे प्रमेय अतिशय प्रभावी आहे. पुढील तत्काळ अनुप्रमेयात
त्याच्या उपयोगांची काही उदाहरणे आहेत.
\begin{cor}
पुढील संच अनिर्णेय आहेत.
\begin{enumerate}
\item $\Setabs{x}{\text{$17$ हे $\cfind{x}$ च्या व्याप्तीत आहे}}$
\item $\Setabs{x}{\text{$\cfind{x}$ स्थिर आहे}}$
\item $\Setabs{x}{\text{$\cfind{x}$ सर्वत्र परिभाषित आहे}}$
\item $\Setabs{x}{\text{जेव्हा $y < y'$ असेल तेव्हा $\cfind{x}(y) \fdefined$, आणि
    जर $\cfind{x}(y') \fdefined$ असेल, तर $\cfind{x}(y) < \cfind{x}(y')$}}$
\end{enumerate}
\end{cor}

\begin{proof}
  हे सर्व अतुच्छ निर्देशांकसंच आहेत.
\end{proof}











\section{स्थिरबिंदू प्रमेय}\label{cmp:thy:fix:sec}

थांबण्याच्या समस्येचा पुन्हा विचार करू. तात्पुरत्या चिन्हांकनासाठी,
$\tuple{x, y}$ ऐवजी $\gn{\cfind{x}(y)}$ असे लिहू;
याला $\cfind{x}(y)$ या मूल्याचे ``नाव'' दर्शवणारे चिन्ह समजा.
या चिन्हांकनाने थांबण्याची समस्या अनिर्णेय असल्याचा आपला एक
पुरावा पुन्हा मांडता येतो.

प्रश्न: पुढील गुणधर्म असलेले संगणनक्षम फलन $h$ आहे का?
प्रत्येक $x$ आणि $y$ साठी,
\[h(\gn{\cfind{x}(y)}) =
\begin{cases}
1 & \text{जर $\cfind{x}(y) \fdefined$ असेल तर} \\
0 & \text{अन्यथा.}
\end{cases}
\]

उत्तर: नाही. अन्यथा पुढील आंशिक फलन
\[g(x) \simeq
\begin{cases}
0 & \text{जर $h(\gn{\cfind{x}(x)}) = 0$ असेल तर} \\
\text{अपरिभाषित} & \text{अन्यथा}
\end{cases}
\]
आंशिक संगणनक्षम असते आणि त्याला काही निर्देशांक~$e$ असता.
पण मग
\[\cfind{e}(e) \simeq
\begin{cases}
0 & \text{जर $h(\gn{\cfind{e}(e)}) = 0$ असेल तर} \\
\text{अपरिभाषित} & \text{अन्यथा,}
\end{cases}
\]
असे मिळते. त्यानुसार $\cfind{e}(e)$ परिभाषित असेल
तेव्हा आणि केवळ तेव्हाच ते अपरिभाषित असेल; हा विरोधाभास आहे.

आता $\cfind{e}$ असलेल्या समीकरणाकडे पाहा. त्यात एका
अर्थाने स्वसंदर्भ आहे: $\cfind{e}(e)$ चे मूल्य विशिष्ट
पद्धतीने $\gn{\cfind{e}(e)}$ वर अवलंबून राहील अशी
रचना केली आहे. स्थिरबिंदू प्रमेय सांगते की हे सर्वसाधारणपणे
करता {\em येते}—फक्त विरोधाभास सिद्ध करण्यासाठीच नाही.

\ref{cmp:thy:fix:lem:fixed-equiv} मध्ये स्थिरबिंदू प्रमेय मांडण्याच्या
दोन समतुल्य पद्धती दिल्या आहेत. तार्किकदृष्ट्या दोन्ही
विधाने खरी असल्याने ती समतुल्य आहेत असे म्हणता येते;
पण आपला खरा अर्थ असा आहे की प्रत्येक विधानावरून दुसरे
सहज निष्पन्न होते. म्हणून ती एकाच प्रमेयाची पर्यायी विधाने
मानता येतात.

\begin{lem}
\label{cmp:thy:fix:lem:fixed-equiv}
पुढील विधाने समतुल्य आहेत:
\begin{enumerate}
\item प्रत्येक आंशिक संगणनक्षम फलन $g(x,y)$ साठी असा
  निर्देशांक~$e$ आहे की प्रत्येक~$y$ साठी
\[\cfind{e}(y) \simeq g(e,y).
\]
\item प्रत्येक संगणनक्षम फलन~$f(x)$ साठी असा निर्देशांक~$e$
  आहे की प्रत्येक~$y$ साठी
\[\cfind{e}(y) \simeq \cfind{f(e)}(y).
\]
\end{enumerate}
\end{lem}

\begin{proof}
$(1) \Rightarrow (2)$: $f$ दिल्यावर $g(x,y) \simeq
\fn{Un}(f(x),y)$ असे $g$ व्याख्यायित करा. (1) वापरून
प्रत्येक~$y$ साठी पुढील गुणधर्म असलेला निर्देशांक~$e$
मिळवा:
\begin{align*}
\cfind{e}(y) & = \fn{Un}(f(e),y) \\
& = \cfind{f(e)}(y).
\end{align*}

$(2) \Rightarrow (1)$: $g$ दिल्यावर $s$-$m$-$n$
प्रमेय वापरून प्रत्येक $x$ आणि~$y$ साठी
$\cfind{f(x)}(y) \simeq g(x,y)$ होईल असे $f$ मिळवा.
(2) वापरून असा निर्देशांक~$e$ मिळवा की
\begin{align*}
\cfind{e}(y) & = \cfind{f(e)}(y) \\
& = g(e,y).
\end{align*}
याने सिद्धता पूर्ण झाली.
\end{proof}

\begin{explain}
विधान (1) खरे आहे (आणि म्हणून (2) देखील) हे दाखवण्यापूर्वी,
ते किती विलक्षण आहे याचा विचार करा. $e$ हा संगणक
कार्यक्रम आहे असे समजा; विधान (1) सांगते की कोणतेही
आंशिक संगणनक्षम $g(x,y)$ दिले असता, $g_e(y) \simeq g(e,y)$
संगणित करणारा संगणक कार्यक्रम $e$ शोधता येतो. म्हणजे
स्वतःचाच संदर्भ असलेले फलन संगणित करणारा कार्यक्रम शोधता येतो.
\end{explain}

\begin{thm}
\ref{cmp:thy:fix:lem:fixed-equiv} मधील दोन्ही विधाने खरी आहेत.
विशेषतः, प्रत्येक आंशिक संगणनक्षम फलन $g(x,y)$ साठी
असा निर्देशांक~$e$ आहे की प्रत्येक~$y$ साठी
\[\cfind{e}(y) \simeq g(e,y).
\]
\end{thm}

\begin{proof}
आवश्यक घटक वरच्या थांबण्याच्या समस्येवरील चर्चेत अप्रत्यक्षपणे
आहेतच. $\fn{diag}(x)$ हे असे संगणनक्षम फलन असू द्या की
ते प्रत्येक $x$ साठी $f_x(y) \simeq \cfind{x}(x,y)$
या फलनाचा निर्देशांक परत करते, म्हणजे
\[\cfind{\fn{diag}(x)}(y) \simeq \cfind{x}(x,y).
\]
$\fn{diag}$ कडे द्विस्थानी फलनाचा कार्यक्रम एकस्थानी
फलनाच्या कार्यक्रमात रूपांतरित करणारे फलन म्हणून पाहा;
मूळ कार्यक्रमाचा पहिला युक्तिवाद त्याच कार्यक्रमावर निश्चित
करून हे फलन मिळते. $\fn{diag}$ ची औपचारिक व्याख्या अशी
करता येते: आधी $s$ ची व्याख्या
\[s(x,y) \simeq \fn{Un}^2(x,x,y),
\]
अशी करा. येथे $\fn{Un}^2$ हे आंशिक
संगणनक्षम द्विस्थानी फलनांसाठी सार्वत्रिक असलेले त्रिस्थानी
फलन आहे. मग $s$-$m$-$n$ प्रमेयानुसार पुढील अट पूर्ण करणारे
आदिम पुनरावर्ती फलन $\fn{diag}$ मिळते:
\[\cfind{\fn{diag}(x)}(y) \simeq s(x,y).
\]

आता $l$ फलनाची व्याख्या अशी करा:
\[l(x,y) \simeq g(\fn{diag}(x),y).
\]
आणि $\gn{l}$ हा $l$ चा निर्देशांक असू द्या. शेवटी
$e = \fn{diag}(\gn{l})$ असू द्या. मग प्रत्येक $y$ साठी
\begin{align*}
\cfind{e}(y) & \simeq \cfind{\fn{diag}(\gn{l})}(y) \\
& \simeq \cfind{\gn{l}}(\gn{l}, y) \\
& \simeq l(\gn{l}, y) \\
& \simeq g(\fn{diag}(\gn{l}),y) \\
& \simeq g(e, y),
\end{align*}
हवे तसे.
\end{proof}

\begin{explain}
येथे नेमके काय घडते? स्वतःचाच मजकूर छापणारा संगणक
कार्यक्रम लिहायचा आहे असे समजा. यासाठी तुमची प्रोग्रामिंग
भाषा चिन्हमालांवरील समृद्ध व काहीशी विचित्र फलनांचा संच
पुरवते असेही समजा. विशेषतः, तिच्यात पुढीलप्रमाणे काम करणारे
फलन $\fn{diag}$ आहे असे समजा: आदान चिन्हमाला~$s$
दिली असता $\fn{diag}$ $s$ मधील `x' या चिन्हाच्या प्रत्येक
घटनेचा शोध घेते आणि त्या जागी मूळ चिन्हमालेची अवतरणचिन्हांत
घेतलेली आवृत्ती बसवते. उदाहरणार्थ, पुढील चिन्हमाला
\begin{quote}
\begin{verbatim}
hello x world
\end{verbatim}
\end{quote}
आदान म्हणून दिल्यास हे फलन पुढील चिन्हमाला परत करते:
\begin{quote}
\begin{verbatim}
hello 'hello x world' world
\end{verbatim}
\end{quote}
मग अपेक्षित कार्यक्रम लिहिणे सोपे आहे; पुढील कार्यक्रम
हे काम करतो हे तपासा:
\begin{quote}
\begin{verbatim}
print(diag('print(diag(x))'))
\end{verbatim}
\end{quote}
C++ आणि Java सारख्या अधिक प्रचलित प्रोग्रामिंग भाषांतही
हीच कल्पना (अधिक गुंतागुंतीच्या अंमलबजावणीसह) काम करते.

आता आपण स्थिरबिंदू प्रमेयाच्या सिद्धतेपासून काही पावलेच
दूर आहोत. $\fn{print}(x,y)$ हे print फलनाचे एक रूप
चिन्हमाला $x$ आणि संख्यात्मक युक्तिवाद $y$ घेते व
$x$ ही चिन्हमाला $y$ वेळा छापते असे समजा. मग पुढील ``कार्यक्रम''
\begin{quote}
\begin{verbatim}
getinput(y);
print(diag('getinput(y); print(diag(x), y)'), y)
\end{verbatim}
\end{quote}
आदान $y$ मिळाल्यावर स्वतःचा मजकूर $y$ वेळा छापतो.
$\fn{getinput}$---$\fn{print}$---$\fn{diag}$ या सांगाड्याच्या
जागी कोणतेही फलन $g(x,y)$ घेतल्यास
\begin{quote}
\begin{verbatim}
g(diag('g(diag(x), y)'), y)
\end{verbatim}
\end{quote}
मिळते. हा कार्यक्रम आदान $y$ वर स्वतःच्या कार्यक्रमाला
आणि $y$ ला $g$ चे युक्तिवाद म्हणून देतो. ``अवतरणचिन्हांत
घेणे'' म्हणजे ``निर्देशांक वापरणे'' असे मानल्यास वरील
सिद्धता मिळते.

सध्या हा पुरावा औपचारिक युक्ती किंवा जादू वाटला तरी
हरकत नाही. पण वरील युक्तिवादाचे तपशील पुन्हा उभारता
आले पाहिजेत. अपूर्णतेची प्रमेये (आणि संबंधित ``स्थिरबिंदू
प्रमेय'') सिद्ध करताना ते का काम करते हे समजण्याच्या
इतर पद्धतींवर चर्चा करू.
\end{explain}













\section{स्थिरबिंदू प्रमेयाचा उपयोग}\label{cmp:thy:apf:sec}

स्थिरबिंदू प्रमेयामुळे आंशिक संगणनक्षम फलनांची व्याख्या
त्यांच्या निर्देशांकांच्या साहाय्याने करता येते. उदाहरणार्थ,
प्रत्येक $y$ साठी पुढील अट पूर्ण करणारा निर्देशांक $e$
शोधता येतो:
\[\cfind{e}(y) = e + y.
\]
आणखी एक उदाहरण म्हणजे स्थिरबिंदू प्रमेयाची सिद्धता
वापरून Java किंवा C++ मध्ये स्वतःचा मजकूर छापणारा
कार्यक्रम तयार करता येतो.

प्रत्येक $e$ साठी $W_e$ हे $\cfind{e}$ चे परिभाषाक्षेत्र
असू द्या. मग $W_0$, $W_1$, $W_2$,~\dots ही क्रमिका
संगणनक्षमपणे प्रगणनीय संचांचे प्रगणन करते, हे आठवा.
यांतील काही संच संगणनक्षम आहेत. आदान म्हणून $x$
घेऊन, $W_x$ संगणनक्षम असेल तेव्हा त्याच्या जातिबोधक
फलनाचा निर्देशांक परत करणारा अल्गोरिदम आहे का, असा
प्रश्न विचारता येतो. उत्तर ``नाही'' आहे; असा अल्गोरिदम नाही:

\begin{thm}
पुढील गुणधर्म असलेले कोणतेही आंशिक संगणनक्षम फलन $f$
नाही: $W_e$ संगणनक्षम असेल तेव्हा $f(e)$ परिभाषित
असते आणि $\cfind{f(e)}$ हे त्याचे जातिबोधक फलन असते.
\end{thm}

\begin{proof}
$f$ हे कोणतेही आंशिक संगणनक्षम फलन असू द्या. आपण असा
निर्देशांक $e$ तयार करू की $W_e$ संगणनक्षम असेल, पण
अपेक्षित निर्देशांक मिळतच नसेल किंवा मिळालेले $\cfind{f(e)}$
हे त्याचे जातिबोधक फलन नसेल.

स्थिरबिंदू प्रमेय वापरून असा निर्देशांक $e$ मिळतो की
\[\cfind{e}(y) \simeq
\begin{cases}
0 & \text{जर $y=0$ आणि $\cfind{f(e)}(0) \fdefined = 0$ असेल तर} \\
\text{अपरिभाषित} & \text{अन्यथा.}
\end{cases}
\]
म्हणजे पुढीलप्रमाणे व्याख्यायित फलनावर स्थिरबिंदू प्रमेय
लावून $e$ मिळतो:
\[g(x,y) \simeq
\begin{cases}
0 & \text{जर $y=0$ आणि $\cfind{f(x)}(0) \fdefined = 0$ असेल तर} \\
\text{अपरिभाषित} & \text{अन्यथा.}
\end{cases}
\]
$g$ आंशिक संगणनक्षम आहे हे अनौपचारिकपणे असे पाहता
येते: आदान $x$ आणि $y$ मिळाल्यावर अल्गोरिदम आधी
$y$ हे~$0$ शी समान आहे का ते तपासतो. समान असल्यास
तो $f(x)$ संगणित करतो आणि मग सार्वत्रिक यंत्र वापरून
$\cfind{f(x)}(0)$ संगणित करतो. हे शेवटचे संगणन थांबून
$0$ परत करत असेल, तर अल्गोरिदम~$0$ परत करतो;
अन्यथा अल्गोरिदम थांबत नाही.

आता $\cfind{f(e)}(0)$ परिभाषित असून $0$ च्या समान
असेल, तर $\cfind{e}(y)$ परिभाषित असते तेव्हा आणि केवळ
तेव्हाच $y$~$0$ च्या समान असते; म्हणून $W_e =
\{ 0 \}$. $\cfind{f(e)}(0)$ अपरिभाषित असेल किंवा
परिभाषित असूनही $0$ च्या समान नसेल, तर $W_e = \emptyset$.
दोन्ही स्थितींत अपेक्षित गुणधर्म फसतो: या आदानावर फलन
निर्देशांक परत करतच नसेल, तर अपेक्षित परिभाषितता मिळत
नाही; निर्देशांक परत करत असेल, तर $\cfind{f(e)}$ हे~$W_e$
चे जातिबोधक फलन नाही, कारण आदान $0$ वर ते चुकीचे
उत्तर देते.
\end{proof}











\section{स्वसंदर्भ वापरून फलनांची व्याख्या}\label{cmp:thy:slf:sec}

फलनांची व्याख्या त्यांच्याच साहाय्याने करता येणे साधारणपणे
उपयुक्त असते. उदाहरणार्थ, संगणनक्षम फलने $k$, $l$
आणि~$m$ दिली असता, स्थिरबिंदू उपप्रमेयानुसार प्रत्येक~$y$
साठी पुढील समीकरण पूर्ण करणारे आंशिक संगणनक्षम फलन~$f$
अस्तित्वात असते:
\[f(y) \simeq
\begin{cases}
k(y) & \text{जर $l(y) = 0$ असेल तर} \\
f(m(y)) & \text{अन्यथा.}
\end{cases}
\]
अधिक नेमके सांगायचे, तर
\[g(x,y) \simeq
\begin{cases}
k(y) & \text{जर $l(y) = 0$ असेल तर} \\
\cfind{x}(m(y)) & \text{अन्यथा}
\end{cases}
\]
असे घेऊन, नंतर $\cfind{e}(y) = g(e,y)$ होईल असा
निर्देशांक~$e$ स्थिरबिंदू उपप्रमेयाने मिळवल्यावर~$f$ मिळते.

ठोस उदाहरण म्हणून, ``महत्तम साधारण विभाजक'' फलन
$\fn{gcd}(u,v)$ ची व्याख्या अशी करता येते:
\[\fn{gcd}(u,v) \simeq
\begin{cases}
v & \text{जर $u = 0$ असेल तर} \\
\fn{gcd}(\fn{mod}(v, u), u) & \text{अन्यथा}
\end{cases}
\]
येथे $\fn{mod}(v, u)$ म्हणजे $v$ ला~$u$ ने भागल्यावर
उरणारी बाकी. स्थिरबिंदू उपप्रमेय वापरून $\fn{gcd}$
आंशिक संगणनक्षम आहे हे मिळते. (प्रत्यक्षात, $y$ मध्ये
$\tuple{u, v}$ या जोडीचा संकेतांक घेऊन हे वरच्या
रूपात मांडता येते.) मग~$u$ वरील गणितीय विगमनाने $\fn{gcd}$
हे प्रत्यक्षात सर्वत्र परिभाषित आहे हे दिसते.

अर्थात, नुकत्याच पाहिलेल्या उदाहरणांपेक्षा बरीच अधिक
गुंतागुंतीची स्वसंदर्भात्मक व्याख्या करता येते. बहुतेक
प्रोग्रामिंग भाषा फलनांची व्याख्या त्यांच्याच साहाय्याने
करण्याचे कोणते ना कोणते साधन देतात. पण एखादे फलन
संगणित करणाऱ्या अल्गोरिदमच्या \emph{निर्देशांका}च्या
साहाय्याने त्याची व्याख्या करता येण्यापेक्षा हे काहीसे
कमी विलक्षण आहे. स्थिरबिंदू प्रमेय पूर्ण सर्वसाधारणतेने
नेमके तेच करू देते.



\chapter{ट्यूरिंग यंत्रांवरील संगणने}\label{tur:mac::chap}








\section{परिचय}\label{tur:mac:int:sec}

$\Nat$ पासून $\Nat$ मध्ये जाणारे एखादे फलन \emph{संगणनक्षम} असणे म्हणजे
नेमके काय? या प्रश्नाला सर्वप्रथम दिलेल्या आणि सर्वाधिक परिचित
उत्तरांपैकी एक असे आहे: एखादे फलन ट्यूरिंग यंत्राने काढता येत असेल, तर ते
संगणनक्षम असते. ॲलन ट्यूरिंग यांनी 1936 मध्ये ही संकल्पना मांडली.
ट्यूरिंग यंत्र हे \emph{संगणनाचे प्रतिमान} आहे---म्हणजे ``संगणनाची
कार्यपद्धती'' ही कल्पना गणिती नेमकेपणाने परिभाषित करण्याचा एक मार्ग.
त्या कल्पनेचा नेमका अर्थ काय, यावर मतभेद असले तरी ट्यूरिंग यंत्रे ही
संगणनाच्या कार्यपद्धती निर्दिष्ट करण्याची एक पद्धत आहेत, यावर व्यापक
सहमती आहे. ``ट्यूरिंग यंत्र'' या नावामुळे हलणारे भाग असलेल्या भौतिक
यंत्राची प्रतिमा डोळ्यांसमोर येऊ शकते; पण काटेकोरपणे पाहता ट्यूरिंग
यंत्र ही पूर्णपणे गणिती रचना आहे. म्हणून ती संगणनाच्या कार्यपद्धतीची
आदर्शीकृत मांडणी करते. उदाहरणार्थ, ट्यूरिंग यंत्राला लागणाऱ्या वेळेवर
किंवा स्मृतीवर आपण कोणतेही बंधन घालत नाही: एखाद्या संगणनाला विश्वातील
अणूंपेक्षा जास्त साठवण-जागा किंवा पायऱ्या लागणार असल्या तरी ट्यूरिंग
यंत्र ते संगणन करू शकते.

\begin{explain}
ट्यूरिंग यंत्राचा विचार एका विशिष्ट प्रकारच्या काल्पनिक यंत्रणेसाठीचा
कार्यक्रम म्हणून करणे बहुधा सर्वांत सोयीचे ठरेल. या यंत्रणेत एक
\emph{फीत} आणि एक \emph{वाचन-लेखन शीर्ष} असते. आपल्या ट्यूरिंग
यंत्रांच्या आवृत्तीत फीत एका दिशेने (उजवीकडे) अनंत असते आणि ती
\emph{घरांमध्ये} विभागलेली असते. प्रत्येक घरात सांत \emph{वर्णमालेतील}
एखादे चिन्ह असू शकते. अशा वर्णमालेत कितीही भिन्न चिन्हे असू शकतात;
पण आपल्याला मुख्यतः तीनच पुरेशी पडतील: $\TMendtape$, $\TMblank$ आणि
$\TMstroke$. यंत्रणा सुरू करताना फीत रिकामी असते (म्हणजे प्रत्येक घरात
$\TMblank$ हे चिन्ह असते); अपवाद म्हणजे सर्वांत डावीकडील घर, ज्यात
$\TMendtape$ असते, आणि \emph{आदान} असलेली सांत संख्येने घरे. कोणत्याही
क्षणी यंत्रणा सांत संख्येतील एखाद्या \emph{अवस्थेत} असते. सुरुवातीला
शीर्ष आदानाच्या पहिल्या घरावर, म्हणजे फितीच्या डाव्या टोकाच्या
घराच्या लगेच उजवीकडील घरावर, वाचत असते आणि यंत्रणा निर्दिष्ट
\emph{प्रारंभिक अवस्थेत} असते. यंत्रणेच्या प्रत्येक पायरीवर शीर्ष
सध्या वाचत असलेल्या घरातील चिन्ह, यंत्रणेची अवस्था आणि ट्यूरिंग
यंत्राचा कार्यक्रम मिळून पुढे काय होईल ते ठरवतात. हा कार्यक्रम एका
आंशिक फलनाने दिलेला असतो: त्याला अवस्था~$q$ आणि चिन्ह~$\sigma$ ही
आदाने दिल्यावर ते $\tuple{q', \sigma',
  D}$ ही त्रयी देते. यंत्रणा $q$ अवस्थेत असताना $\sigma$ चिन्ह
वाचत असेल, तर ती सध्याच्या घरातील चिन्हाच्या जागी $\sigma'$ लिहिते;
$D$ हे $\TMleft$, $\TMright$ किंवा $\TMstay$ यांपैकी कोणते आहे,
त्यानुसार शीर्ष डावीकडे, उजवीकडे किंवा त्याच जागी राहते; आणि यंत्रणा
$q'$ अवस्थेत जाते.


उदाहरणार्थ, \ref{tur:mac:int:fig:tm} मधील स्थिती पाहा.
\begin{figure}
  \olasset{assets/diagrams/turing-machine.tikz}
  \caption{कार्यक्रम चालवत असलेले ट्यूरिंग यंत्र.}
  \label{tur:mac:int:fig:tm}
\end{figure}
ट्यूरिंग यंत्राच्या फितीच्या दिसणाऱ्या भागात सर्वांत डावीकडील घरात
फितीच्या टोकाचे $\TMendtape$ चिन्ह आहे; त्यानंतर तीन $1$, एक $0$ आणि
आणखी चार $1$ आहेत. शीर्ष डावीकडून तिसरे घर वाचत आहे; त्या घरात
$1$ आहे आणि यंत्रणा $q_1$ अवस्थेत आहे---म्हणजे ``यंत्र $q_1$
अवस्थेत $1$ वाचत आहे'' असे आपण म्हणतो. ट्यूरिंग यंत्राचा कार्यक्रम
$\tuple{q_1, 1}$ या आदानासाठी $\tuple{q_2,
0, \TMstay}$ ही त्रयी देत असेल, तर यंत्र तिसऱ्या घरातील $1$ च्या
जागी $0$ लिहील, वाचन-लेखन शीर्ष त्याच जागी ठेवील आणि $q_2$
अवस्थेत जाईल. त्यानंतर कार्यक्रम $\tuple{q_2, 0}$ या आदानासाठी
$\tuple{q_3, 0, \TMright}$ देत असेल, तर यंत्र त्या $0$ च्या
जागी पुन्हा $0$ लिहील (म्हणजे प्रत्यक्षात शीर्षाखालील फितीतील मजकूर
अपरिवर्तित राहील), एक घर उजवीकडे सरकेल आणि $q_3$ अवस्थेत जाईल.
अशाच रीतीने पुढे चालेल.

एखादी अवस्था $q_n$ आणि चिन्ह $\sigma$ अशी आढळली की
$\tuple{q_n, \sigma}$ साठी कोणतीही सूचना नाही---म्हणजे
$\tuple{q_n,\sigma}$ या आदानावर संक्रमण फलन अपरिभाषित आहे---तर यंत्र
\emph{थांबते} असे आपण म्हणतो.
दुसऱ्या शब्दांत, यंत्राला अमलात आणण्यासाठी कोणतीही सूचना उरत नाही
आणि त्या टप्प्यावर त्याचे कार्य थांबते. काही वेळा थांबणे एका विशिष्ट
थांबण्याच्या अवस्थेने~$h$ दर्शवतात. याचे अधिक तपशीलवार प्रदर्शन पुढे
केले जाईल.
\end{explain}

\begin{digress}
ट्यूरिंग यांच्या ``On computable numbers'' या शोधनिबंधाचे वैशिष्ट्य
असे की त्यात त्यांनी केवळ औपचारिक व्याख्याच दिली नाही, तर ती व्याख्या
संगणनक्षमतेची अंतर्ज्ञानी कल्पना योग्यरीत्या पकडते, असा युक्तिवादही
केला. व्याख्येवरून ट्यूरिंग यंत्राने काढता येणारे कोणतेही फलन अंतर्ज्ञानी
अर्थाने संगणनक्षम आहे, हे स्पष्ट व्हावे. ट्यूरिंग यांनी याचा उलट
दिशेचा दावा---म्हणजे आपण स्वाभाविकपणे संगणनक्षम मानू असे कोणतेही फलन
अशा यंत्राने काढता येते---यासाठी तीन प्रकारचे युक्तिवाद दिले. त्यांच्या
शब्दांत ते असे:
\begin{enumerate}
\item अंतःप्रज्ञेला थेट आवाहन.
\item दोन व्याख्या सममूल्य आहेत, याची सिद्धता (नवी व्याख्या
  अंतःप्रज्ञेला अधिक पटत असल्यास).
\item संगणनक्षम संख्यांच्या मोठ्या वर्गांची उदाहरणे देणे.
\end{enumerate}
आपला उद्देश संगणनक्षमता ``तत्त्वतः'' परिभाषित करण्याचा आहे: म्हणजे
वेळ आणि जागा यांच्या व्यावहारिक मर्यादा विचारात न घेता. अर्थात,
संगणनक्षमतेची ही सर्वांत व्यापक व्याख्या मिळाल्यानंतर मर्यादित
संसाधनांतील संगणनाचाही अभ्यास करता येतो. ``संगणकीय जटिलता'' म्हणून
ओळखल्या जाणाऱ्या विषयाचा हा मध्यवर्ती भाग आहे.
\end{digress}

\begin{history}
ॲलन ट्यूरिंग यांनी 1936 मध्ये ट्यूरिंग यंत्रांची कल्पना मांडली.
त्या काळी त्यांना प्रथम-क्रम तर्कशास्त्रातील निर्णयक्षमतेचा प्रश्न
विशेष महत्त्वाचा वाटत होता; तरी संगणकाच्या रचनेच्या पायावरील निर्णायक
शोधनिबंध म्हणून त्यांच्या त्या निबंधाचे वर्णन झाले आहे. त्या निबंधात
ट्यूरिंग यांनी संगणनक्षम वास्तव संख्यांवर---म्हणजे ज्यांचे दशांश
विस्तार संगणनक्षम आहेत अशा संख्यांवर---लक्ष केंद्रित केले. पण
नैसर्गिक संख्यांवरील संगणनक्षम फलनांसाठी आणि पुढील गोष्टींसाठीही या
संकल्पना रूपांतरित करणे कठीण नाही, असे त्यांनी नमूद केले. लक्षात घ्या:
पहिला कार्यरत सर्वसाधारण वापराचा संगणक 1941 मध्ये (जर्मन कॉनराड
झुसे यांनी बर्लिनमध्ये) बांधला जाण्याच्या पूर्ण पाच
वर्षे आधी हे झाले. 1936 नंतर सात वर्षांनी (1943) टॉमी फ्लावर्स आणि
डॉलिस हिलमधील त्यांच्या सहकाऱ्यांनी संकेतभेदनासाठी कोलॉसस बांधला;
तो ब्लेच्ली पार्कमध्ये 1944 मध्ये कार्यरत झाला.
अमेरिकन ENIAC
(1945) हा नऊ वर्षांनी आला; मँचेस्टरमधील प्रायोगिक
यंत्र---Manchester Small-Scale Experimental Machine---बारा
वर्षांनी (1948) बांधले गेले; आणि अमेरिकन BINAC (1949) ची पहिली
चाचणी तेरा वर्षांनी झाली. Manchester SSEM हा संग्रहित-कार्यक्रम
असलेला पहिला इलेक्ट्रॉनिक संगणक होता---त्यात कार्यक्रमाच्या सूचना
इलेक्ट्रॉनिक स्मृतीत ठेवता येत होत्या. त्याआधीच्या यंत्रांत सूचना
कागदी फितीवर दिल्या जात किंवा काही यंत्रांत तारा नव्याने जोडाव्या लागत.

\end{history}











\section{ट्यूरिंग यंत्रांचे निरूपण}\label{tur:mac:rep:sec}

\begin{explain}
ट्यूरिंग यंत्रे \emph{अवस्था-आलेखांद्वारे} चित्ररूपात दाखवता येतात.
या आलेखांत बाणांनी जोडलेली अवस्थांची घरे असतात. अर्थातच प्रत्येक
अवस्था-घर यंत्राची एक अवस्था दर्शवते. प्रत्येक बाण त्या अवस्थेतून
अमलात आणता येणारी एक सूचना दर्शवतो; सूचनेचे तपशील त्या बाणाच्या
वर किंवा खाली लिहिलेले असतात. पुढील यंत्र पाहा. त्याच्या $q_0$ आणि
$q_1$ अशा फक्त दोन अंतर्गत अवस्था आणि एकच सूचना आहे:
\[\begin{tikzpicture}[->,>=stealth',shorten >=1pt,auto,node distance=2.8cm,
                    semithick]
  \tikzstyle{every state}=[fill=none,draw=black,text=black]
  \node[initial,state]   (A)              {$q_0$};
  \node[state]   (B) [right of=A] {$q_1$};
  \path (A) edge  node {\TMtrans{\TMblank}{\TMstroke}{\TMright}} (B);
\end{tikzpicture}
\]
ट्यूरिंग यंत्राला वाचन-लेखन शीर्ष आणि त्यावर आदान लिहिलेली फीत
असते, हे आठवा. या सूचनेचा अर्थ असा: \emph{$q_0$ अवस्थेत
रिकामे चिन्ह~$\TMblank$ वाचत असाल, तर $\TMstroke$ लिहा,
उजवीकडे सरका आणि $q_1$ अवस्थेत जा}. हे संक्रमण फलनाने
$\tuple{q_0, \TMblank}$ याचा $\tuple{q_1, \TMstroke,
\TMright}$ शी केलेला संबंध याच्या सममूल्य आहे.
\end{explain}

\begin{ex}
\emph{सम यंत्र}: फितीवर $\TMstroke$ चिन्हांची संख्या सम असेल तेव्हा
आणि तेव्हाच पुढील ट्यूरिंग यंत्र थांबते (फितीवरील पहिल्या
$\TMblank$ आधी सर्व $\TMstroke$ चिन्हे येतात, हे गृहीत धरून).
\[\begin{tikzpicture}[->,>=stealth',shorten >=1pt,auto,node distance=2.8cm,
                    semithick]
  \tikzstyle{every state}=[fill=none,draw=black,text=black]
  \node[initial,state]         (A)              {$q_0$};
  \node[state]         (B) [right of=A] {$q_1$};
  \path (A) edge [bend left] node {\TMtrans{\TMstroke}{\TMstroke}{\TMright}} (B)
        (B) edge [loop above] node {\TMtrans{\TMblank}{\TMblank}{\TMright}} (B)
            edge [bend left] node {\TMtrans{\TMstroke}{\TMstroke}{\TMright}} (A);
\end{tikzpicture}
\]

हा अवस्था-आलेख पुढील संक्रमण फलनाशी संबंधित आहे:
\begin{align*}
\delta(q_0, \TMstroke) & = \tuple{q_1, \TMstroke, \TMright},\\
\delta(q_1, \TMstroke) & = \tuple{q_0, \TMstroke, \TMright},\\
\delta(q_1, \TMblank)  & = \tuple{q_1, \TMblank, \TMright}
\end{align*}
\end{ex}

\begin{explain}
वरील यंत्र आदानातील रेघांची संख्या सम असेल तेव्हाच थांबते.
अन्यथा यंत्र (सैद्धांतिकदृष्ट्या) अमर्याद काळ कार्यरत राहते.
कोणत्याही यंत्राच्या कोणत्याही आदानासाठी त्याच्या \emph{स्थितिवर्णनांचा}
क्रम पाहता येतो; यंत्र थांबल्यास या क्रमावरून त्याचे प्रदान निश्चित
करता येते. पण यंत्र थांबत नसेल तर स्थितिवर्णनांचा क्रम अनंत असतो;
त्यावरून प्रत्येक आदानासाठी सांत वेळेत प्रदान मिळते असा अर्थ नाही.

स्थितिवर्णनाची औपचारिक
व्याख्या पुढे देऊ. सध्या, यंत्र चालू असताना प्रत्येक क्षणी त्याची
स्थिती दाखवणाऱ्या आकृत्यांची मालिका असे स्थितिवर्णनांचे अंतर्ज्ञानी
आकलन करता येईल. स्थितिवर्णनात फितीवरील मजकूर, यंत्राची अवस्था
आणि वाचन-लेखन शीर्षाची जागा दिसते.

सम यंत्राला चार $\TMstroke$ चिन्हांचे आदान दिल्यास त्याची
स्थितिवर्णने पाहू. या वेळी यंत्र थांबेल, अशी अपेक्षा आहे. मग
तीन $\TMstroke$ चिन्हांचे आदान देऊ; त्या वेळी ते कायम चालू राहील.

यंत्र $q_0$ अवस्थेत सुरू होते आणि सर्वांत डावीकडील $\TMstroke$
वाचते. त्याची प्रारंभिक स्थिती अशी दर्शवता येते:
\[\TMendtape \TMstroke_0 \TMstroke \TMstroke \TMstroke \TMblank \ldots
\]
वरील स्थितिवर्णन सरळ आहे. दिसते त्याप्रमाणे यंत्र $q_0$ अवस्थेत
सुरू होते आणि सर्वांत डावीकडील $\TMstroke$ वाचते. पहिल्या

$\TMstroke$ च्या खाली लिहिलेला अवस्थेच्या नावाचा निर्देशांक हे
दर्शवतो. या टप्प्यावर लागू होणारी सूचना
$\delta(q_0, \TMstroke) = \tuple{q_1, \TMstroke, \TMright}$ आहे;
म्हणून यंत्र फितीवर उजवीकडे सरकते आणि $q_1$ अवस्थेत जाते.
\[\TMendtape \TMstroke \TMstroke_1 \TMstroke \TMstroke \TMblank \ldots
\]
आता यंत्र $q_1$ अवस्थेत $\TMstroke$ वाचत आहे, म्हणून
$\delta(q_1, \TMstroke) = \tuple{q_0,
  \TMstroke, \TMright}$ ही सूचना ``पाळावी'' लागते. त्यातून पुढील
स्थितिवर्णन मिळते:
\[\TMendtape \TMstroke \TMstroke \TMstroke_0 \TMstroke \TMblank \ldots
\]
यंत्र पुढे चालू राहिल्यावर हे नियम त्याच क्रमाने पुन्हा लागू होतात
आणि पुढील दोन स्थितिवर्णने मिळतात:
\[\TMendtape \TMstroke \TMstroke \TMstroke \TMstroke_1 \TMblank \ldots
\]
\[\TMendtape \TMstroke \TMstroke \TMstroke \TMstroke \TMblank_0 \ldots
\]
आता यंत्र $q_0$ अवस्थेत $\TMblank$ वाचत आहे. संक्रमण-आलेखावरून
येथे अमलात आणण्यासाठी कोणतीही सूचना नाही, हे सहज दिसते;
म्हणून यंत्र थांबले आहे. याचा अर्थ आदान स्वीकारले गेले.

आता यंत्राला तीन $\TMstroke$ चिन्हांचे आदान देऊन सुरू करू.
सुरुवातीची काही स्थितिवर्णने मिळतीजुळती असतील, कारण त्याच सूचना
अमलात येतात; फक्त फितीवरील आदान थोडे वेगळे आहे:
\[\TMendtape \TMstroke_0 \TMstroke \TMstroke \TMblank \ldots
\]
\[\TMendtape \TMstroke \TMstroke_1 \TMstroke \TMblank \ldots
\]
\[\TMendtape \TMstroke \TMstroke \TMstroke_0 \TMblank \ldots
\]
\[\TMendtape \TMstroke \TMstroke \TMstroke \TMblank_1 \ldots
\]
आता यंत्र सर्व $\TMstroke$ चिन्हांच्या पुढे गेले आहे आणि $q_1$
अवस्थेत $\TMblank$ वाचत आहे. आलेखात दाखवल्याप्रमाणे
$\delta(q_1, \TMblank) =\tuple{q_1,
\TMblank, \TMright}$ या रूपाची सूचना आहे. फीत उजवीकडे अनंतपर्यंत
$\TMblank$ चिन्हांनी भरलेली असल्याने यंत्र ही सूचना \emph{कायम}
अमलात आणत राहील; ते $q_1$ अवस्थेतच राहून सतत उजवीकडे सरकेल.
यंत्र कधीच थांबणार नाही आणि आदान स्वीकारणार नाही.
\end{explain}

\begin{explain}
सर्व यंत्रे थांबतातच असे नाही, हे लक्षात ठेवणे महत्त्वाचे आहे. थांबणे
म्हणजे अमलात आणण्यासाठी सूचना संपणे असेल, तर प्रत्येक अवस्थेत
प्रत्येक चिन्हासाठी बाहेर जाणारा बाण ठेवून कधीही न थांबणारे यंत्र
बनवता येते. $q_0$ अवस्थेत $\TMblank$ वाचण्यासाठी एक सूचना
जोडून सम यंत्र कायम चालू राहील असे बदलता येते.
\end{explain}

\begin{ex}
\[\begin{tikzpicture}[->,>=stealth',shorten >=1pt,auto,node distance=2.8cm,
                    semithick]
  \tikzstyle{every state}=[fill=none,draw=black,text=black]
  \node[initial,state] (A)              {$q_0$};
  \node[state]         (B) [right of=A] {$q_1$};
  \path
  (A) edge [bend left] node {\TMtrans{\TMstroke}{\TMstroke}{\TMright}} (B)
      edge [loop above] node {\TMtrans{\TMblank}{\TMblank}{\TMright}} (A)
  (B) edge [loop above] node {\TMtrans{\TMblank}{\TMblank}{\TMright}} (B)
      edge [bend left] node {\TMtrans{\TMstroke}{\TMstroke}{\TMright}} (A);
\end{tikzpicture}
\]
\end{ex}

\begin{explain}
यंत्र-सारणी हा ट्यूरिंग यंत्रांचे निरूपण करण्याचा दुसरा मार्ग आहे.
यंत्र-सारणीत फितीवरील वर्णमाला $x$-अक्षावर आणि यंत्राच्या अवस्था
$y$-अक्षावर दाखवलेल्या असतात. सारणीतील प्रत्येक अवस्था आणि चिन्ह
यांच्या छेदावर सूचनेचा उरलेला भाग---नवी अवस्था, नवे चिन्ह आणि
हालचालीची दिशा---लिहिलेला असतो. यंत्र कोणत्या अवस्थेत आणि कोणते
चिन्ह वाचताना थांबते हे यंत्र-सारणीवरून सहज ठरवता येते. सारणीत
जिथे नोंद नाही, तिथे यंत्र थांबू शकते. अवस्था-आलेख किंवा सूचनासंच
यांत यंत्र नेमके कुठे थांबेल हे लगेच स्पष्ट होत नाही; पण
यंत्र-सारणीत रिकामी घरे शोधून अशी ठिकाणे पटकन ओळखता येतात.
\end{explain}

\begin{ex}
सम यंत्राची यंत्र-सारणी अशी आहे:
\[\centering
\begin{tabular}{lllll}
\cline{2-4}
\multicolumn{1}{l|}{}      & \multicolumn{1}{c|}{$\TMblank$}
& \multicolumn{1}{c|}{$\TMstroke$}     & \multicolumn{1}{c|}{$\TMendtape$}   &  \\ \cline{1-4}
\multicolumn{1}{|l|}{$q_0$} & \multicolumn{1}{l|}{}
& \multicolumn{1}{l|}{\TMtrans{\TMstroke}{q_1}{\TMright}} &
\multicolumn{1}{l|}{\phantom{\TMtrans{\TMstroke}{q_1}{\TMright}}} &  \\ \cline{1-4}
\multicolumn{1}{|l|}{$q_1$} & \multicolumn{1}{l|}{\TMtrans{\TMblank}{q_1}{\TMright}}
& \multicolumn{1}{l|}{\TMtrans{\TMstroke}{q_0}{\TMright}} &
\multicolumn{1}{l|}{\phantom{\TMtrans{\TMstroke}{q_1}{\TMright}}} &  \\ \cline{1-4}
\end{tabular}
\]
दिसते त्याप्रमाणे यंत्र $q_0$ अवस्थेत $\TMblank$ वाचताना थांबते.
\end{ex}

\begin{explain}
आतापर्यंत आपण फक्त आदान वाचून ते स्वीकारणारी यंत्रे पाहिली.
परंतु ट्यूरिंग यंत्रे वाचूही शकतात आणि लिहूही शकतात. अशा यंत्राचे
एक उदाहरण (अशी अनेक उदाहरणे असली तरी) म्हणजे \emph{दुप्पट करणारे
यंत्र}. हे यंत्र फितीवरील $n$ $\TMstroke$ चिन्हांचा सलग गट
आदान म्हणून घेते आणि $2n$ $\TMstroke$ चिन्हांचा सलग गट प्रदान करते.
\end{explain}

\begin{ex}
  \label{tur:mac:rep:ex:doubler} दुप्पट करणारे यंत्र बनवण्यापूर्वी ही समस्या
  सोडवण्यासाठी \emph{धोरण} ठरवणे महत्त्वाचे आहे. यंत्राला (आपल्या
  व्याख्येनुसार) त्याने किती $\TMstroke$ चिन्हे वाचली हे लक्षात
  ठेवता येत नाही; म्हणून फितीवरील सर्व $\TMstroke$ चिन्हांचा मागोवा
  घेण्याचा मार्ग हवा. एक मार्ग म्हणजे आदान आणि प्रदान यांमध्ये
  $\TMblank$ ठेवून त्यांना वेगळे करणे. मग यंत्र आदानातील पहिले
  $\TMstroke$ पुसू शकते, आदानाचा उरलेला भाग पार करू शकते,
  $\TMblank$ सोडू शकते आणि दोन नवी $\TMstroke$ चिन्हे लिहू शकते.
  नंतर ते आदानातील दुसरे $\TMstroke$ शोधून त्याच्यासाठीही असेच
  दुपटीकरण करील. आदानातील प्रत्येक एका $\TMstroke$ साठी ते
  प्रदानात दोन $\TMstroke$ चिन्हे लिहील. पुढे जाताना आदान
  पुसल्यामुळे कोणतेही $\TMstroke$ चिन्ह सुटणार नाही किंवा त्याची
  दुप्पट मोजणी होणार नाही, याची खात्री होते. संपूर्ण आदान पुसल्यावर
  फितीवर $2n$ $\TMstroke$ चिन्हे उरतील. परिणामी ट्यूरिंग यंत्राचा
  अवस्था-आलेख \ref{tur:mac:rep:fig:doubler} मध्ये दाखवला आहे.
  \begin{figure}
\[\begin{tikzpicture}[->,>=stealth',shorten >=1pt,auto,node distance=2.8cm,
                    semithick]
  \tikzstyle{every state}=[fill=none,draw=black,text=black]
  \node[initial,state] (1)              {$q_0$};
  \node[state]         (2) [right of=1] {$q_1$};
  \node[state]         (3) [right of=2] {$q_2$};
  \node[state]         (4) [below of=3] {$q_3$};
  \node[state]         (5) [left of=4]  {$q_4$};
  \node[state]         (6) [left of=5]  {$q_5$};
  \path (1) edge node {\TMtrans{\TMstroke}{\TMblank}{\TMright}} (2)
    (2) edge [loop above] node {\TMtrans{\TMstroke}{\TMstroke}{\TMright}} (2)
      edge node {\TMtrans{\TMblank}{\TMblank}{\TMright}} (3)
    (3) edge [loop above] node {\TMtrans{\TMstroke}{\TMstroke}{\TMright}} (3)
        edge node {\TMtrans{\TMblank}{\TMstroke}{\TMright}} (4)
    (4) edge [loop below] node {\TMtrans{\TMblank}{\TMstroke}{\TMleft}} (4)
        edge node {\TMtrans{\TMstroke}{\TMstroke}{\TMleft}} (5)
    (5) edge [loop below]  node {\TMtrans{\TMstroke}{\TMstroke}{\TMleft}} (5)
        edge              node {\TMtrans{\TMblank}{\TMblank}{\TMleft}} (6)
    (6) edge [loop below] node {\TMtrans{\TMstroke}{\TMstroke}{\TMleft}} (6)
        edge              node {\TMtrans{\TMblank}{\TMblank}{\TMright}} (1);
\end{tikzpicture}
\]
\caption{दुप्पट करणारे यंत्र}
\label{tur:mac:rep:fig:doubler}
\end{figure}
\end{ex}

\begin{prob}
एखादे आदान निवडा आणि \ref{tur:mac:rep:ex:doubler} मधील दुप्पट
करणाऱ्या यंत्राची स्थितिवर्णने क्रमाने शोधा.
\end{prob}

\begin{prob}
$\{\TMendtape,\TMblank, A, B\}$ ही वर्णमाला असलेले ट्यूरिंग यंत्र
रचा. $A$ आणि $B$ यांच्या कोणत्याही चिन्हमालेत $A$ ची संख्या $B$
च्या संख्येइतकी \emph{आणि} सर्व $A$ सर्व $B$ च्या आधी असतील,
तर यंत्राने ती चिन्हमाला स्वीकारावी, म्हणजे तिच्यावर थांबावे.
$A$ आणि $B$ यांच्या संख्या असमान असतील किंवा सर्व $A$ सर्व $B$
च्या आधी नसतील, तर यंत्राने ती नाकारावी, म्हणजे तिच्यावर थांबू नये.
(उदाहरणार्थ, यंत्राने $AABB$ आणि $AAABBB$ स्वीकाराव्यात; पण
$AAB$ आणि $AABBAABB$ या दोन्ही नाकाराव्यात.)
\end{prob}

\begin{prob}
$\{\TMendtape,\TMblank, A, B\}$ ही वर्णमाला असलेले ट्यूरिंग यंत्र
रचा. त्याने $A$ आणि $B$ यांची कोणतीही चिन्हमाला $\alpha$ आदान
म्हणून घेऊन तिची प्रत तयार करावी आणि $\alpha\alpha$ या रूपाचे
प्रदान द्यावे. (उदाहरणार्थ, $ABBA$ या आदानापासून $ABBAABBA$
हे प्रदान मिळावे.)
\end{prob}

\begin{prob}
\emph{वर्णक्रमात आहे का?}: $\{\TMendtape,\TMblank, A, B\}$ ही
वर्णमाला असलेले ट्यूरिंग यंत्र रचा. $A$ आणि $B$ यांची सांत
क्रमिका आदान दिल्यावर सर्व $A$ सर्व $B$ यांच्या डावीकडे आहेत की
नाही, हे त्याने तपासावे. यंत्राने आदान-चिन्हमाला फितीवर ठेवावी;
ती ``वर्णक्रमात'' असेल तर थांबावे आणि नसेल तर कायम पुनरावर्तनात
चालू राहावे.
\end{prob}

\begin{prob}
\emph{वर्णक्रम लावणारे यंत्र}: $\{\TMendtape,\TMblank, A, B\}$
ही वर्णमाला असलेले ट्यूरिंग यंत्र रचा. $A$ आणि $B$ यांची सांत
क्रमिका आदान म्हणून घेऊन त्यांची पुनर्रचना अशा रीतीने करावी की
सर्व $A$ सर्व $B$ यांच्या डावीकडे येतील. (उदाहरणार्थ, $BABAA$
ही क्रमिका $AAABB$ व्हावी आणि $ABBABB$ ही क्रमिका $AABBBB$
व्हावी.)
\end{prob}











\section{ट्यूरिंग यंत्रे}\label{tur:mac:tur:sec}

\begin{explain}
ट्यूरिंग यंत्राची औपचारिक व्याख्या अमूर्त वाटते, पण प्रत्यक्षात
ती सोपी आहे: यंत्राचे कार्य कसे चालते हे निर्दिष्ट करण्यासाठी लागणारी
सर्व माहिती ती एका गणिती रचनेत एकत्र करते. यात (1) यंत्र कोणत्या
अवस्थांत असू शकते, (2) फितीवर कोणती चिन्हे असू शकतात, (3) यंत्राने
कोणत्या अवस्थेत सुरू व्हावे आणि (4) त्याचा सूचनासंच कोणता आहे,
या गोष्टी येतात.
\end{explain}

\begin{defn}[ट्यूरिंग यंत्र]
\emph{ट्यूरिंग यंत्र} $M$ म्हणजे पुढील घटक असलेली
$\langle Q, \Sigma, q_0,
\delta\rangle$ ही क्रमित चौकडी:
\begin{enumerate}
\item \emph{अवस्थांचा} सांत संच~$Q$,
\item $\TMendtape$ आणि $\TMblank$ यांचा समावेश असलेली सांत
  \emph{वर्णमाला} $\Sigma$,
\item \emph{प्रारंभिक अवस्था}~$q_0 \in Q$,
\item सांत \emph{सूचनासंच}~$\delta\colon Q \times \Sigma
  \pto Q \times \Sigma \times \{\TMleft, \TMright, \TMstay\}$.
\end{enumerate}
या आंशिक फलनाला~$\delta$ हे $M$ चे \emph{संक्रमण फलन} असेही म्हणतात.
\end{defn}

\begin{explain}
फीत फक्त एकाच दिशेने अनंत आहे, असे आपण मानतो. म्हणून फितीचे
डावे टोक दर्शवण्यासाठी $\TMendtape$ हे खास चिन्ह ठेवणे उपयुक्त ठरते.
त्यामुळे ट्यूरिंग यंत्राच्या कार्यक्रमाला फितीच्या बाहेर जाण्याचा
``धोका'' कधी आहे हे ओळखणे सोपे जाते. हे चिन्ह कधीही बदलून लिहिले
जात नाही, असे आपण गृहीत धरू शकलो असतो; म्हणजे
$\delta(q,\TMendtape)$ परिभाषित असेल, तर
$\delta(q,\TMendtape) = \tuple{q', \TMendtape, x}$ असेल.
काही पाठ्यपुस्तके असे करतात; आपण मात्र हे गृहीत धरत नाही. ट्यूरिंग
यंत्र बनवताना ते $\TMendtape$ कधीही बदलून लिहिणार नाही याची
सावधगिरी तुम्ही घेऊ शकता. शिवाय काही प्रसंगी असे बदलून लिहिण्यास
परवानगी दिल्यास सोयीची लवचीकता मिळते.
\end{explain}

\begin{ex}
\emph{सम यंत्र:} सम यंत्र औपचारिकरीत्या
$\tuple{Q, \Sigma, q_0, \delta}$ ही क्रमित चौकडी आहे. येथे
\begin{align*}
Q & = \{ q_0, q_1 \} \\
\Sigma & = \{ \TMendtape, \TMblank, \TMstroke \}, \\
\delta(q_0, \TMstroke) & = \tuple{q_1, \TMstroke, \TMright},\\
\delta(q_1, \TMstroke) & = \tuple{q_0, \TMstroke, \TMright},\\
\delta(q_1, \TMblank)  & = \tuple{q_1, \TMblank, \TMright}.
\end{align*}
\end{ex}











\section{स्थितिवर्णने आणि संगणने}\label{cmp:tur:con:sec}

\begin{explain}
मागे सम यंत्राच्या स्थितिवर्णनांचा क्रम कसा पाहिला, ते आठवा.
फीत, वाचन-लेखन शीर्ष आणि ट्यूरिंग यंत्राचा कार्यक्रम यांची बनलेली
काल्पनिक यंत्रणा ही ट्यूरिंग यंत्राचे संगणन कसे असते ते डोळ्यांसमोर
आणण्याचा केवळ अंतर्ज्ञानी मार्ग आहे. औपचारिकरीत्या दिलेल्या आदानावर
ट्यूरिंग यंत्राचे संगणन \emph{स्थितिवर्णनांची} क्रमिका म्हणून
परिभाषित करता येते. प्रत्येक स्थितिवर्णनात चिन्हांची एक क्रमिका
(संगणनाच्या त्या टप्प्यावर फितीतील मजकुराशी संबंधित), वाचन-लेखन
शीर्षाची जागा दर्शवणारी संख्या आणि एक अवस्था असते. त्यांच्या
आधारे ट्यूरिंग यंत्र~$M$ दिलेल्या आदानावर काय काढते हे परिभाषित
करता येते.
\end{explain}

\begin{defn}[स्थितिवर्णन]
ट्यूरिंग यंत्र $M = \tuple{Q, \Sigma, q_0,
\delta}$ चे \emph{स्थितिवर्णन} म्हणजे $\tuple{C, m, q}$ ही त्रयी;
येथे
\begin{enumerate}
\item $C \in \Sigma^*$ ही $\Sigma$ मधील चिन्हांची सांत क्रमिका आहे,
\item $m \in \Nat$ ही $< \len{C}$ असलेली संख्या आहे, आणि
\item $q \in Q$ आहे.
\end{enumerate}
अंतर्ज्ञानी अर्थाने क्रमिका~$C$ म्हणजे फितीतील मजकूर (सर्वांत
डावीकडील घरापासून शेवटच्या रिकाम्या नसलेल्या किंवा आधी भेट दिलेल्या
घरापर्यंतच्या सर्व घरांतील चिन्हे); $m$~म्हणजे वाचन-लेखन शीर्ष ज्या
घरावर आहे त्या घराचा क्रमांक (सर्वांत डावीकडील घराचा क्रमांक $0$
धरून); आणि $q$ म्हणजे यंत्राची सध्याची अवस्था.
\end{defn}

\begin{explain}
ट्यूरिंग यंत्राचे संभाव्य आदान ही चिन्हांची एक क्रमिका असते; बहुतेक
वेळा ती एखाद्या संख्येला काही रूपात सांकेतिक करणारी क्रमिका असते.
ट्यूरिंग यंत्राचे प्रारंभिक स्थितिवर्णन म्हणजे त्या आदानावर यंत्र
चालवायला सुरू करतानाचे स्थितिवर्णन: फितीत डाव्या टोकाचे चिन्ह असते
आणि त्याच्या लगेच उजवीकडील घरांत आदान लिहिलेले असते; वाचन-लेखन
शीर्ष आदानाच्या सर्वांत डावीकडील घरावर असते (म्हणजे डाव्या टोकाच्या
चिन्हाच्या उजवीकडील घरावर); आणि यंत्रणा निर्दिष्ट प्रारंभिक
अवस्थेत~$q_0$ असते.

\end{explain}

\begin{defn}[प्रारंभिक स्थितिवर्णन]
आदान रिकामे नसताना, $I \in \Sigma^*$ या आदानासाठी $M$ चे \emph{प्रारंभिक स्थितिवर्णन}
असे आहे:
\[\tuple{\TMendtape \frown I, 1, q_0}.
\]
रिकामे आदान घेतल्यास फितीवरील पहिल्या रिकाम्या घराचाही
स्थितिवर्णनात समावेश करावा; अन्यथा शीर्षाचा क्रमांक आणि क्रमिकेची
लांबी यांच्यावरील वरची अट पूर्ण होत नाही.

\end{defn}

\begin{explain}
$\frown$ हे चिन्ह \emph{जोडणीसाठी} आहे---आदान-चिन्हमाला डाव्या
टोकाच्या चिन्हाच्या लगेच उजवीकडे सुरू होते.

\end{explain}

\begin{defn}
स्थितिवर्णन $\tuple{C, m, q}$ पासून \emph{एका पायरीत
$\tuple{C', m', q'}$ हे स्थितिवर्णन मिळते} ($M$ नुसार), असे आपण
तेव्हा आणि केवळ तेव्हाच म्हणतो, जेव्हा
\begin{enumerate}
\item $C$ चे $m$-वे चिन्ह $\sigma$ असते,
\item $M$ च्या सूचनासंचात $\delta(q, \sigma) =
  \tuple{q', \sigma', D}$ निर्दिष्ट असते,
\item $C'$ चे $m$-वे चिन्ह $\sigma'$ असते, आणि
\item
\begin{enumerate}
\item $D = L$ असते; या स्थितीत $m>0$ असल्यास $m' = m - 1$ असते, अन्यथा $m'=0$ असते; किंवा
\item $D = R$ आणि $m' = m + 1$ असते, किंवा
\item $D = N$ आणि $m' = m$ असते,
\end{enumerate}
\item $m' = \len{C}$ असेल, तर $\len{C'} = \len{C} + 1$ असते आणि
  $C'$ चे $m'$-वे चिन्ह~$\TMblank$ असते. अन्यथा $\len{C'}=\len{C}$ असते.
\item $i < \len{C}$ आणि $i \neq m$ असलेल्या सर्व $i$ साठी $C'(i) = C(i)$ असते,
\end{enumerate}
\end{defn}

\begin{defn}\label{cmp:tur:con:defn:run-output}
\emph{आदान~$I$ वरील $M$ ची धाव} म्हणजे $M$ च्या स्थितिवर्णनांची
$C_i$ ही सांत किंवा अनंत क्रमिका; येथे $C_0$ हे आदान~$I$ साठी $M$ चे प्रारंभिक
स्थितिवर्णन असते आणि क्रमिकेतील अंतिम स्थितिवर्णन वगळता प्रत्येक $C_i$ पासून एका पायरीत $C_{i+1}$
मिळते.


$C_k = \tuple{C, m, q}$ असेल, $C$ चे $m$-वे चिन्ह~$\sigma$ असेल
आणि $\delta(q,
\sigma)$ अपरिभाषित असेल, तर $M$ \emph{आदान $I$ वर $k$
पायऱ्यांनंतर थांबते} असे म्हणतो. त्या वेळी आदान~$I$ साठी $M$ चे
\emph{प्रदान}~$O$ असते; येथे $O$ ही शेवटी~$\TMblank$ नसलेली
चिन्हमाला असून एखाद्या~$j \in \Nat$ साठी
$C = \TMendtape \concat O \concat \TMblank^j$ असते.
($\TMblank^j$ म्हणजे $j$ $\TMblank$ चिन्हांची क्रमिका.)
ही प्रदानाची अट डाव्या टोकाचे चिन्ह शेवटपर्यंत शाबूत राहिलेल्या
धावांसाठीच थेट लागू होते. ते चिन्ह बदलून लिहिले असल्यास प्रदान
ठरवण्यासाठी अतिरिक्त नियम लागेल.

\end{defn}

\begin{explain}
या व्याख्येनुसार $M$ चे प्रदान~$O$ नेहमी~$\TMblank$ व्यतिरिक्त
इतर चिन्हाने संपते; किंवा $k$~व्या वेळी संपूर्ण फीत (सर्वांत
डावीकडील~$\TMendtape$ सोडून)~$\TMblank$ चिन्हांनी भरलेली असेल,
तर $O$~ही रिकामी चिन्हमाला असते.
\end{explain}











\section{संख्यांचे एकचिन्ही निरूपण}\label{tur:mac:una:sec}

\begin{explain}
ट्यूरिंग यंत्रे त्यांच्या फितीवर लिहिलेल्या चिन्हांच्या क्रमिकांवर
कार्य करतात. यंत्र वापरत असलेल्या वर्णमालेनुसार या क्रमिका विविध
आदानांचे आणि प्रदानांचे निरूपण करू शकतात. विशेष रस आहे तो
\emph{अंकगणितीय} फलने, म्हणजे नैसर्गिक संख्यांची फलने, संगणित करणाऱ्या
ट्यूरिंग यंत्रांत. धन पूर्णांकांचे निरूपण करण्याचा एक सोपा मार्ग म्हणजे
एकाच चिन्हाच्या~$\TMstroke$ क्रमिकांनी त्यांचे सांकेतीकरण करणे.
$n \in \Nat$ असेल, तर $n = 0$ असताना $\TMstroke^n$ ही रिकामी
क्रमिका मानू; अन्यथा ती नेमक्या $n$ $\TMstroke$ चिन्हांची क्रमिका
असेल.
\end{explain}

\begin{defn}[संगणन]
ट्यूरिंग यंत्र~$M$ हे $f\colon \Nat^k \to \Nat$ फलन \emph{संगणित
करते}, असे तेव्हा आणि केवळ तेव्हाच म्हणतो, जेव्हा पुढील स्वरूपाचे
प्रत्येक आदान घेतल्यावर
\[\TMstroke^{n_1} \TMblank \TMstroke^{n_2} \TMblank \dots \TMblank \TMstroke^{n_k}
\]
$M$~थांबते आणि त्याचे प्रदान $\TMstroke^{f(n_1, \dots, n_k)}$ असते.
\end{defn}

\begin{prob}
ट्यूरिंग यंत्र~$M$ हे $f\colon \Nat^k \to \Nat^m$ फलन
कधी संगणित करते याची व्याख्या द्या.
\end{prob}

\begin{ex}\label{tur:mac:una:ex:adder}
\emph{बेरीज:}
$f(n,m) = n + m$ हे फलन संगणित करणारे यंत्र बनवू.
यासाठी सुरुवातीला फितीवर अनुक्रमे $n$ आणि~$m$ लांबीचे
$\TMstroke$ चिन्हांचे दोन खंड असले पाहिजेत आणि शेवटी
$n+m$~$\TMstroke$ चिन्हांचा एकच खंड राहिल्यावर यंत्र थांबले
पाहिजे. दोन $\TMstroke$ आदान-खंडांमध्ये~$\TMblank$ असते. म्हणून एक पद्धत अशी:
$\TMblank$ असलेल्या घरात एक रेघ लिहा आणि शेवटचे~$\TMstroke$
पुसून टाका.
\begin{figure}\[\begin{tikzpicture}[->,>=stealth',shorten >=1pt,auto,node distance=2.8cm,
                    semithick]
  \tikzstyle{every state}=[fill=none,draw=black,text=black]
  \node[initial,state]         (A)              {$q_0$};
  \node[state]         (B) [right of=A] {$q_1$};
  \node[state]         (C) [right of=B] {$q_2$};
  \path (A) edge node {\TMtrans{\TMblank}{\TMstroke}{\TMstay}} (B)
           edge [loop above] node {\TMtrans{\TMstroke}{\TMstroke}{\TMright}} (A)
        (B) edge [loop above] node {\TMtrans{\TMstroke}{\TMstroke}{\TMright}} (B)
            edge node {\TMtrans{\TMblank}{\TMblank}{\TMleft}} (C)
        (C) edge [loop above] node {\TMtrans{\TMstroke}{\TMblank}{\TMstay}} (C);
\end{tikzpicture}
\]\caption{$f(x,y) = x+y$ संगणित करणारे यंत्र}

\label{tur:mac:una:fig:adder}
\end{figure}
\end{ex}

\begin{prob}
$\tuple{3,2}$ हे आदान असताना \ref{tur:mac:una:ex:adder}
मधील यंत्राच्या स्थितिवर्णनांचा क्रम पाहा. यंत्राने $0+0$ संगणित
केल्यास काय घडते?
\end{prob}

\begin{explain}
\ref{tur:mac:rep:ex:doubler} मध्ये आपण एका ट्यूरिंग यंत्राचे उदाहरण दिले:
ते~$\TMstroke$ चिन्हांची क्रमिका आदान म्हणून घेते आणि फितीवर
दुप्पट~$\TMstroke$ चिन्हांची क्रमिका ठेवून थांबते. हे दुप्पट करणारे
यंत्र आहे. पण दुप्पट झालेल्या $\TMstroke$ खंडाच्या डावीकडेही
प्रदानात $\TMblank$ चिन्हे येतात. म्हणून, वाटू शकते तसे, ते
प्रत्यक्षात $f(x) = 2x$ हे फलन संगणित करत नाही. ही अडचण
दूर करण्याचे दोन मार्ग पाहू.
\end{explain}

\begin{ex}
\ref{tur:mac:una:fig:doubler-disc} मधील यंत्र $f(x) =
2x$ हे फलन संगणित करते. \ref{tur:mac:rep:ex:doubler} मधील यंत्र
प्रत्येक आदानातील~$\TMstroke$ पुसून फितीच्या अगदी उजवीकडे दोन
$\TMstroke$ चिन्हे लिहिते. त्याऐवजी हे यंत्र आदानातील प्रत्येक
$\TMstroke$ साठी उजवीकडे एकच~$\TMstroke$ वाढवते. आदान कुठे
संपते हे लक्षात ठेवण्यासाठी ते आदान आणि वाढवलेल्या रेघा यांच्यामध्ये
एक~$\TMblank$ ठेवते; अगदी शेवटी त्या जागी~$\TMstroke$ लिहिते.
तसेच आदानात सध्या कुठे आहोत हे ``लक्षात ठेवण्यासाठी'' आदान-खंडातील
एक~$\TMstroke$ तात्पुरते~$\TMblank$ ने बदलते.
  \begin{figure}
\[\begin{tikzpicture}[->,>=stealth',shorten >=1pt,auto,node distance=2.8cm,
                    semithick]
  \tikzstyle{every state}=[fill=none,draw=black,text=black]
  \node[initial,state] (0)              {$q_0$};
  \node[state]         (1) [above of=0] {$q_1$};
  \node[state]         (2) [above of=1] {$q_2$};
  \node[state]         (3) [right of=2] {$q_3$};
  \node[state]         (4) [below of=3] {$q_4$};
  \node[state]         (5) [below of=4] {$q_5$};
  \node[state]         (6) [above right of=4] {$q_6$};
  \node[state]         (7) [below of=6] {$q_7$};
  \node[state]         (8) [below of=7] {$q_8$};
  \path (0) edge node {\TMtrans{\TMstroke}{\TMblank}{\TMright}} (1)
    (1) edge [loop left] node {\TMtrans{\TMstroke}{\TMstroke}{\TMright}} (1)
      edge node {\TMtrans{\TMblank}{\TMblank}{\TMright}} (2)
    (2) edge [loop above] node {\TMtrans{\TMstroke}{\TMstroke}{\TMright}} (2)
        edge node {\TMtrans{\TMblank}{\TMstroke}{\TMleft}} (3)
    (3) edge [loop above] node {\TMtrans{\TMstroke}{\TMstroke}{\TMleft}} (3)
        edge node[left] {\TMtrans{\TMblank}{\TMblank}{\TMleft}} (4)
    (4) edge        node[left] {\TMtrans{\TMstroke}{\TMstroke}{\TMleft}} (5)
    (5) edge [loop below] node {\TMtrans{\TMstroke}{\TMstroke}{\TMleft}} (5)
        edge              node {\TMtrans{\TMblank}{\TMstroke}{\TMright}} (0)
    (4) edge      node[sloped] {\TMtrans{\TMblank}{\TMstroke}{\TMright}} (6)
    (6) edge              node {\TMtrans{\TMblank}{\TMstroke}{\TMright}} (7)
    (7) edge [loop right] node {\TMtrans{\TMstroke}{\TMstroke}{\TMright}} (7)
        edge              node {\TMtrans{\TMblank}{\TMblank}{\TMleft}} (8)
    (8) edge [loop right] node {\TMtrans{\TMstroke}{\TMblank}{\TMstay}} (8);
    \end{tikzpicture}
\]
\caption{$f(x) = 2x$ संगणित करणारे यंत्र}
\label{tur:mac:una:fig:doubler-disc}
\end{figure}
\end{ex}

\begin{ex}\label{tur:mac:una:ex:mover}
$f(x) = 2x$ संगणित करण्याचा दुसरा मार्ग म्हणजे मूळ दुप्पट
करणारे यंत्र तसेच ठेवून, शेवटी दुप्पट झालेला रेघांचा खंड फितीच्या
अगदी डावीकडे नेणाऱ्या अवस्था आणि सूचना जोडणे. \ref{tur:mac:una:fig:mover}
मधील यंत्र नेमका हा शेवटचा भाग करते: सुरुवातीला फितीवर
प्रथम~$\TMblank$ चिन्हांचा आणि नंतर~$\TMstroke$ चिन्हांचा खंड
असतो (आणि शीर्ष~$\TMblank$ खंडातील कोणत्याही घरावर असते);
मग ते $\TMstroke$ चिन्हे एकेक करून पुसते आणि फितीच्या सुरुवातीला
लिहिते. काम कधी संपले हे ओळखता यावे म्हणून ते प्रथम
$\TMstroke$ खंडाचा शेवट~$\TMendtape$ चिन्हाने दाखवते;
शेवटी हे चिन्ह पुसले जाते. अवस्था~$q_6$ पासून मोजायला
सुरुवात केली आहे, जेणेकरून त्या मूळ दुप्पट करणाऱ्या यंत्रात जोडता
येतील. त्यासाठी एक अतिरिक्त सूचना
$\delta(q_0,
\TMblank) = \tuple{q_6,\TMblank,\TMstay}$, म्हणजे~$q_0$ पासून
~$q_6$ कडे~$\TMtrans{\TMblank}{\TMblank}{\TMstay}$ हे लेबल असलेला
बाण, एवढीच हवी. (येथे एक सूक्ष्म अडचण आहे: परिणामी यंत्र
आदान~$x=0$ साठी चालत नाही. ही अडचण सरावासाठी ठेवू.)
\begin{figure}
  \[  \begin{tikzpicture}[->,>=stealth',shorten >=1pt,auto,node distance=2.8cm,
                      semithick]
    \tikzstyle{every state}=[fill=none,draw=black,text=black]
    \node[initial,state] (6)              {$q_6$};
    \node[state]         (7) [right of=6] {$q_7$};
    \node[state]         (8) [right of=7] {$q_8$};
    \node[state]         (9) [below of=8] {$q_9$};
    \node[state]         (10) [left of=9] {$q_{10}$};
    \node[state]         (11) [left of=10]  {$q_{11}$};
    \node[state]         (12) [below of=11]  {$q_{12}$};
    \node[state]         (13) [right of=12] {$q_{13}$};
    \node[state]         (14) [right of=13] {$q_{14}$};
    \path
    (6)  edge node {\TMtrans{\TMblank}{\TMblank}{\TMright}} (7)
    (7)  edge [loop above] node {\TMtrans{\TMblank}{\TMblank}{\TMright}} (7)
         edge node {\TMtrans{\TMstroke}{\TMstroke}{\TMright}} (8)
    (8)  edge [loop above] node {\TMtrans{\TMstroke}{\TMstroke}{\TMright}} (8)
         edge node {\TMtrans{\TMblank}{\TMendtape}{\TMleft}} (9)
    (9)  edge [loop right] node {\TMtrans{\TMstroke}{\TMstroke}{\TMleft}} (9)
         edge node {\TMtrans{\TMblank}{\TMblank}{\TMright}} (10)
    (10) edge node[above] {\TMtrans{\TMstroke}{\TMblank}{\TMleft}} (11)
    (11) edge [loop above] node {\TMtrans{\TMblank}{\TMblank}{\TMleft}} (11)
         edge node[left] {\begin{tabular}{@{}l@{}}
          \TMtrans{\TMendtape}{\TMendtape}{\TMright}\\
          \TMtrans{\TMstroke}{\TMstroke}{\TMright}
         \end{tabular}} (12)
    (12) edge node[] {\TMtrans{\TMblank}{\TMstroke}{\TMright}} (13)
    (13) edge [loop below] node {\TMtrans{\TMblank}{\TMblank}{\TMright}} (13)
         edge node[sloped] {\TMtrans{\TMstroke}{\TMblank}{\TMleft}} (11)
    (13) edge node {\TMtrans{\TMendtape}{\TMblank}{\TMstay}} (14);
  \end{tikzpicture}
  \]
  \caption{$\TMstroke$ चिन्हांचा खंड डावीकडे नेणे}
  \label{tur:mac:una:fig:mover}
\end{figure}
\end{ex}

\begin{prob}
\ref{tur:mac:una:ex:mover} मध्ये आपण
\ref{tur:mac:rep:ex:doubler} मधील मूळ दुप्पट करणारे यंत्र
आणि \ref{tur:mac:una:fig:mover} मधील खंड हलवणारे यंत्र
एकत्र करून बनणाऱ्या यंत्राचे वर्णन केले. हे संयुक्त यंत्र
आदान~$x=0$, म्हणजे रिकामी फीत, घेऊन सुरू केल्यास काय घडते?
या वेळी यंत्राने प्रदान~$2x=0$ देऊन थांबावे यासाठी त्यात
काय बदल कराल? (एक अवस्था आणि एक संक्रमण जोडून हे करता
आले पाहिजे.)

\end{prob}

\begin{prob}
\emph{वजाबाकी:} $n$ आणि~$m$ लांबीच्या रेघांच्या दोन अरिक्त
क्रमिका आदान म्हणून दिल्या असताना, जेथे $n >
m$ आहे, $f(n,m) = n - m$ हे फलन संगणित करणारे ट्यूरिंग यंत्र
रचा.
\end{prob}

\begin{prob}
\emph{समता:} पुढील फलन संगणित करणारे ट्यूरिंग यंत्र रचा:
\[\fn{equality}(n,m) =
\begin{cases}
  \text{1} & \text{$n = m$ असेल तर} \\
  \text{0} & \text{$n \neq m$ असेल तर}
\end{cases}
\]
येथे~$n$ आणि~$m \in \PosInt$ आहेत.
\end{prob}

\begin{prob}
$x$ आणि $y$ हे धन पूर्णांक फितीवर $\TMblank$ ने विभक्त केलेल्या
$\TMstroke$ चिन्हांच्या क्रमिकांनी निरूपित केले असताना
$\min(x,y)$ हे फलन संगणित करणारे ट्यूरिंग यंत्र रचा.
यंत्राच्या वर्णमालेत अतिरिक्त चिन्हे वापरू शकता.

$\min$ हे फलन आपल्या आदानांतील लहान मूल्य निवडते; म्हणून
$\min(3,5)=3$, $\min(20,16)=16$, $\min(4,4)=4$, इत्यादी.
\end{prob}

\begin{defn}
ट्यूरिंग यंत्र~$M$ हे $f\colon \Nat^k
  \pto \Nat$ हे आंशिक फलन तेव्हा आणि केवळ तेव्हाच संगणित
करते, जेव्हा
  \begin{enumerate}
    \item $f(n_1, \dots, n_k) = m$ असेल तर $M$ हे
    $\TMstroke^{n_1}\concat\TMblank\concat
    \dots \concat\TMblank\concat\TMstroke^{n_k}$ आदानावर
    थांबते आणि त्याचे प्रदान $\TMstroke^{m}$ असते.
    \item $f(n_1, \dots, n_k)$ अपरिभाषित असेल तर $M$
    मुळीच थांबत नाही, किंवा थांबले तरी त्याचे प्रदान
    $\TMstroke$ चिन्हांचा एकच खंड नसते.
  \end{enumerate}
\end{defn}











\section{थांबण्याच्या अवस्था}\label{tur:mac:hal:sec}

\begin{explain}
अमलात आणण्यासाठी कोणतीही सूचना उरली नसेल तेव्हाच आपली यंत्रे
थांबतात, अशी आपण व्याख्या केली आहे. पण ट्यूरिंग यंत्रांच्या
नेहमीच्या निरूपणांत थांबण्यासाठी खास \emph{थांबण्याची अवस्था}~$h$
असते; येथे $h \in Q$.

थांबण्याच्या अवस्थेमागील कल्पना सोपी आहे: यंत्राचे काम संपले की
(ते आदान स्वीकारण्यास तयार झाले किंवा त्याने प्रदान लिहून संपवले
की) ते~$h$ अवस्थेत जाते आणि तिथे थांबते. काही यंत्रांना अशा दोन
अवस्था असतात: एक आदान स्वीकारण्यासाठी आणि दुसरी ते नाकारण्यासाठी.
\end{explain}

\begin{ex}\emph{थांबण्याच्या अवस्था}.
ही कल्पना स्पष्ट करण्यासाठी सम यंत्रात बदल करून पाहू. आदानातील
रेघांची संख्या सम असेल तेव्हा यंत्र~$q_0$ अवस्थेत थांबण्याऐवजी,
त्याला थांबण्याच्या अवस्थेत पाठवणारी सूचना जोडू शकतो.
\[\begin{tikzpicture}[->,>=stealth',shorten >=1pt,auto,node distance=2.8cm,
                    semithick]
  \tikzstyle{every state}=[fill=none,draw=black,text=black]
  \node[initial, state]         (A)                     {$q_0$};
  \node[state]         (B) [right of=A] {$q_1$};
  \node[state]         (C) [below of=A] {$h$};
  \path (A) edge [bend left] node {\TMtrans{\TMstroke}{\TMstroke}{R}} (B)
  	    edge node {\TMtrans{\TMblank}{\TMblank}{N}} (C)
        (B) edge [loop above] node {\TMtrans{\TMblank}{\TMblank}{R}} (B)
            edge [bend left] node {\TMtrans{\TMstroke}{\TMstroke}{R}} (A);
\end{tikzpicture}
\]

उदाहरण आणखी वाढवू. यंत्राला आदानातील रेघांची संख्या विषम असल्याचे
आढळल्यास मूळ सम यंत्र कधीच थांबत नाही. चक्रात फिरत राहणाऱ्या
सूचनेच्या जागी यंत्राला \emph{नकार} अवस्था~$r$ मध्ये नेणारी सूचना
देऊन यंत्रात बदल करू शकतो.
\[\begin{tikzpicture}[->,>=stealth',shorten >=1pt,auto,node distance=2.8cm,
                    semithick]
  \tikzstyle{every state}=[fill=none,draw=black,text=black]
  \node[initial,state]         (A)                     {$q_0$};
  \node[state]         (B) [right of=A] {$q_1$};
  \node[state]         (C) [below of=A] {$h$};
  \node[state]         (D) [below of=B] {$r$};
  \path (A) edge [bend left] node {\TMtrans{\TMstroke}{\TMstroke}{R}} (B)
  	    edge node {\TMtrans{\TMblank}{\TMblank}{N}} (C)
        (B) edge node {\TMtrans{\TMblank}{\TMblank}{N}} (D)
            edge [bend left] node {\TMtrans{\TMstroke}{\TMstroke}{R}} (A);
\end{tikzpicture}
\]
\end{ex}

\begin{explain}
अशा प्रसंगी खास थांबण्याची अवस्था जोडल्याने यंत्र कोणती आदाने
स्वीकारते आणि कोणती नाकारते हे स्पष्ट होते. पण खास थांबण्याची
अवस्था जोडल्याने यंत्राची संगणनशक्ती वाढत नाही, हेही लक्षात
घ्यावे. त्याचप्रमाणे थांबण्याची कमी औपचारिक कल्पनाही उपयोगाची
आहे. या प्रकरणात आतापर्यंत वापरलेल्या थांबण्याच्या व्याख्येमुळे
\emph{थांबण्याची समस्या} याचा पुरावा अंतर्ज्ञानी रीतीने
सहज दाखवता येतो. म्हणून आपण आपली मूळ व्याख्या पुढेही वापरू.
\end{explain}











\section{शिस्तबद्ध यंत्रे}\label{tur:mac:dis:sec}

\begin{explain}
\ref{tur:mac:hal:sec} मध्ये आपण एकच निर्दिष्ट थांबण्याची अवस्था~$h$
असलेली ट्यूरिंग यंत्रे पाहिली: अशी यंत्रे थांबलीच तर~$h$
अवस्थेतच थांबतात. म्हणून एका थांबण्याच्या अवस्थेची यंत्रे, आपण
सर्वसाधारण ट्यूरिंग यंत्रांना जितकी मुभा देतो त्यापेक्षा अधिक
``शिस्तबद्ध'' आहेत. ट्यूरिंग यंत्रांच्या वर्तनावर आणखीही निर्बंध
घालता येतील. उदाहरणार्थ, घर~$0$ वरील फितीच्या डाव्या टोकाचे
चिन्ह पुसण्यास किंवा घर~$0$ वरून डावीकडे जाण्याचा प्रयत्न
करण्यास आपण त्यांना मनाई केलेली नाही. (आपल्या व्याख्येनुसार
अशा वेळी शीर्ष घर~$0$ वरच राहते; इतर व्याख्यांनुसार यंत्र
थांबते.) तसेच यंत्र थांबलेच तर घर~$1$ वर थांबेल, असे गृहीत
धरता आले तर काही वेळा सोयीचे ठरते.
\end{explain}

\begin{defn}\label{tur:mac:dis:defn:disciplined}
ट्यूरिंग यंत्र~$M$ \emph{शिस्तबद्ध} असते तेव्हा आणि केवळ तेव्हाच,
जेव्हा
\begin{enumerate}
    \item त्याला एकच निर्दिष्ट थांबण्याची अवस्था~$h$ असते,
    \item ते थांबलेच तर घर~$1$ वाचत असताना थांबते,
    \item ते घर~$0$ वरील $\TMendtape$ चिन्ह कधीही पुसत नाही, आणि
    \item ते घर~$0$ वरून डावीकडे जाण्याचा कधीही प्रयत्न करत नाही.
\end{enumerate}
\end{defn}

\begin{explain}
कोणतेही ट्यूरिंग यंत्र त्याच वर्तनाचे पण निर्दिष्ट थांबण्याची
अवस्था असलेले यंत्र बनवता येते, हे आपण आधी पाहिले. यासाठी
नवी अवस्था~$h$ जोडायची आणि मूळ~$\delta$ ज्या
$\tuple{q,\sigma}$ जोडीवर अपरिभाषित आहे, त्या प्रत्येक जोडीसाठी
$\delta(q, \sigma) = \tuple{h, \sigma, N}$ ही सूचना जोडायची.
पुरावा किचकट असला तरी कोणतेही ट्यूरिंग यंत्र~$M$ अशा
शिस्तबद्ध ट्यूरिंग यंत्रात~$M'$ बदलता येते की जे त्याच आदानांवर
थांबते आणि तेच प्रदान देते. उदाहरणार्थ, यंत्र थांबत असताना
घर~$1$ वर नसेल, तर शीर्ष फितीच्या डाव्या टोकाचे चिन्ह सापडेपर्यंत
डावीकडे नेणाऱ्या, त्यानंतर एक घर उजवीकडे नेणाऱ्या आणि मग यंत्र
थांबवणाऱ्या काही सूचना जोडता येतात. इतर अटी कशा पूर्ण करता
येतील याचा विचार तुम्ही करा.
मूळ यंत्राने डाव्या टोकाचे चिन्ह पुसले किंवा तेच चिन्ह दुसऱ्या घरावर
लिहिले असेल, तर एवढी साधी शोध-पद्धत पुरेशी नाही; सर्वसाधारण
रूपांतराच्या पुराव्यात हा प्रसंग वेगळा हाताळावा लागतो.

\end{explain}

\begin{ex}
    \ref{tur:mac:dis:fig:adder-disc} मध्ये आपण \ref{tur:mac:una:ex:adder}
    मधील बेरीज-यंत्राचे शिस्तबद्ध यंत्रात रूपांतर करतो.
    \begin{figure}\[\begin{tikzpicture}[->,>=stealth',shorten >=1pt,auto,node distance=2.8cm,
                    semithick]
  \tikzstyle{every state}=[fill=none,draw=black,text=black]
  \node[initial,state]         (A)              {$q_0$};
  \node[state]         (B) [right of=A] {$q_1$};
  \node[state]         (C) [below right of=B] {$q_2$};
  \node[state]         (D) [below left of=C] {$q_3$};
  \node[state]         (H) [left of=D] {$h$};
  \path (A) edge node {\TMtrans{\TMblank}{\TMstroke}{\TMstay}} (B)
           edge [loop below] node {\TMtrans{\TMstroke}{\TMstroke}{\TMright}} (A)
        (B) edge [loop below] node {\TMtrans{\TMstroke}{\TMstroke}{\TMright}} (B)
            edge node {\TMtrans{\TMblank}{\TMblank}{\TMleft}} (C)
        (C) edge node {\TMtrans{\TMstroke}{\TMblank}{\TMleft}} (D)
        (D) edge [loop above] node {\TMtrans{\TMstroke}{\TMstroke}{\TMleft}} (D)
        edge node {\TMtrans{\TMendtape}{\TMendtape}{\TMright}} (H);
\end{tikzpicture}
\]\caption{शिस्तबद्ध बेरीज-यंत्र}

\label{tur:mac:dis:fig:adder-disc}
\end{figure}
\end{ex}

\begin{prop}\label{tur:mac:dis:prop:disciplined} प्रत्येक ट्यूरिंग यंत्र~$M$
साठी असे शिस्तबद्ध ट्यूरिंग यंत्र~$M'$ असते की $M$~प्रदान~$O$
देऊन थांबत असेल तर ते देखील प्रदान~$O$ देऊन थांबते; आणि
$M$~थांबत नसेल तर ते देखील थांबत नाही. विशेषतः, ट्यूरिंग
यंत्राला संगणित करता येणारे कोणतेही $f\colon\Nat^n \to \Nat$
फलन शिस्तबद्ध ट्यूरिंग यंत्रालाही संगणित करता येते.
\end{prop}

\begin{prob}\label{tur:mac:dis:prob:disc-succ}
$f(x) = x+1$ संगणित करणारे शिस्तबद्ध यंत्र द्या.
\end{prob}


\begin{prob}\label{tur:mac:dis:prob:copier}
$\TMstroke^n$ आदानावर सुरू केल्यावर
$\TMstroke^n \concat \TMblank \concat \TMstroke^n$ हे प्रदान
देणारे शिस्तबद्ध यंत्र शोधा.
\end{prob}











\section{ट्यूरिंग यंत्रांचे संयोजन}\label{tur:mac:cmb:sec}

\begin{explain}
आतापर्यंत पाहिलेली ट्यूरिंग यंत्रांची उदाहरणे बरीच सोपी आहेत.
पण आधुनिक कार्यक्रमणभाषेत सोडवता येणारी कोणतीही समस्या ट्यूरिंग
यंत्रांनीही सोडवता येते. अधिक गुंतागुंतीची ट्यूरिंग यंत्रे बनवताना
ती एकत्र जोडता येतात हे समजून घेणे महत्त्वाचे आहे. मग प्रक्रिया
सोप्या भागांत विभागून अधिक गुंतागुंतीच्या समस्या सोडवणारी यंत्रे
बनवता येतात. गुंतागुंतीची समस्या घटक भागांत विभागण्याचा नैसर्गिक
मार्ग सापडला, तर अनेक सोपी ट्यूरिंग यंत्रे तयार करून त्यांचे
एकाच यंत्रात संयोजन करता येते. असे यंत्र अनेक टप्प्यांत समस्या
सोडवते. पुढील विभागातील थांबण्याची समस्या हाताळताना हे विशेष
महत्त्वाचे ठरेल.

$M = \tuple{Q, \Sigma, q_0, \delta}$ आणि
~$M' = \tuple{Q', \Sigma', q_0', \delta'}$ ही ट्यूरिंग यंत्रे
कशी जोडायची? $M$~थांबल्यानंतर त्याच फितीवर
यंत्र~$M'$ ची धाव सुरू करू; फितीवरील मजकूर आणि शीर्षाची जागा कायम राहील.
असे करणारे एकच ट्यूरिंग यंत्र $M \frown M'$ मिळवण्यासाठी पुढील
पायऱ्या घ्या:
\begin{enumerate}
    \item $M'$ च्या~$Q'$ मधील सर्व अवस्थांचे क्रमांक (किंवा नावे)
    बदला, म्हणजे $M$ आणि~$M'$ यांची कोणतीही अवस्था समान राहणार
    नाही ($Q \cap Q' = \emptyset$).
    \item $M \frown M'$ च्या अवस्था $Q \cup Q'$ असतील.
    \item फितीची वर्णमाला $\Sigma \cup \Sigma'$ असेल.
    \item प्रारंभिक अवस्था~$q_0$ असेल.
    \item संक्रमण फलन पुढील $\delta''$ असेल:
    \[\delta''(q,\sigma) =
    \begin{cases}
      \delta(q,\sigma) & \text{$q \in Q$ आणि $\delta(q,\sigma)$ परिभाषित असेल तर}\\
      \delta'(q,\sigma) & \text{$q \in Q'$ असेल तर}\\
      \tuple{q_0', \sigma, \TMstay} & \text{$q \in Q$ असेल आणि
      $\delta(q,\sigma)$ अपरिभाषित असेल तर}
    \end{cases}\]
\end{enumerate}

मिळालेले यंत्र $q \in Q$ अवस्थेत असताना~$M$ च्या सूचना,
तर $q \in Q'$ अवस्थेत असताना~$M'$ च्या सूचना वापरते.
$q \in Q$ अवस्थेत~$\sigma$ चिन्ह वाचताना $M$ मध्ये त्यासाठी
संक्रमण नसेल (म्हणजे $M$ थांबले असते), तर ते~$M'$ च्या
प्रारंभिक अवस्थेत जाते; फितीवरील मजकूर आणि शीर्षाची जागा
तशीच ठेवते.

यंत्र~$M$ शिस्तबद्ध नसेल तर $M$~थांबताना फितीचे शीर्ष कुठे
आहे हे आपल्याला माहीत नसते. म्हणून~$M$ च्या अंतिम
स्थितिवर्णनात शीर्ष घर~$1$ वाचत असेलच असे नाही. यंत्रे
जोडताना हे लक्षात ठेवणे महत्त्वाचे आहे.
\end{explain}

\begin{ex}
\emph{यंत्रांचे संयोजन:} सुरुवातीला $n$ आणि~$m$ लांबीचे
$\TMstroke$ चिन्हांचे दोन खंड आदान म्हणून घेऊन शेवटी
फितीवर $2(m+n)$ $\TMstroke$ चिन्हांचा एकच खंड ठेवून थांबणारे
यंत्र रचू. यासाठी परिचित अशी बेरीज-यंत्र आणि दुप्पट करणारे
यंत्र ही दोन यंत्रे जोडता येतील. आधी बेरीज-यंत्राचा अवस्था-आलेख
काढू.
\[\begin{tikzpicture}[->,>=stealth',shorten >=1pt,auto,node distance=2.8cm,
                    semithick]
  \tikzstyle{every state}=[fill=none,draw=black,text=black]
  \node[initial,state] (A)              {$q_0$};
  \node[state]         (B) [right of=A] {$q_1$};
  \node[state]         (C) [right of=B] {$q_2$};
  \path (A) edge node {\TMtrans{\TMblank}{\TMstroke}{\TMstay}} (B)
            edge [loop above] node {\TMtrans{\TMstroke}{\TMstroke}{\TMright}} (A)
        (B) edge [loop above] node {\TMtrans{\TMstroke}{\TMstroke}{\TMright}} (B)
            edge node {\TMtrans{\TMblank}{\TMblank}{\TMleft}} (C)
        (C) edge [loop above] node {\TMtrans{\TMstroke}{\TMblank}{\TMstay}} (C);
\end{tikzpicture}
\]
प्रदान दुप्पट करण्यासाठी अवस्था~$q_2$ मध्ये थांबण्याऐवजी
यंत्र पुढे चालवायचे आहे. दुप्पट करणारे यंत्र आदानातील पहिली
रेघ पुसते आणि स्वतंत्र प्रदानात दोन रेघा लिहिते, हे आठवा.
बेरीज-यंत्राच्या प्रदानातील पहिली रेघ शीर्ष वाचत असेल याची
खात्री करणारी सूचना जोडू.
\[\begin{tikzpicture}[->,>=stealth',shorten >=1pt,auto,node distance=2.8cm,
                    semithick]
  \tikzstyle{every state}=[fill=none,draw=black,text=black]
  \node[initial,state] (A)              {$q_0$};
  \node[state]         (B) [right of=A] {$q_1$};
  \node[state]         (C) [right of=B] {$q_2$};
  \node[state]         (D) [below left of=C] {$q_3$};
  \node[state]         (E) [below left of=D] {$q_4$};
  \path (A) edge node {\TMtrans{\TMblank}{\TMstroke}{\TMstay}} (B)
            edge [loop above] node {\TMtrans{\TMstroke}{\TMstroke}{\TMright}} (A)
        (B) edge [loop above] node {\TMtrans{\TMstroke}{\TMstroke}{\TMright}} (B)
            edge node {\TMtrans{\TMblank}{\TMblank}{\TMleft}} (C)
        (C) edge node {\TMtrans{\TMstroke}{\TMblank}{\TMleft}} (D)
        (D) edge [loop left] node {\TMtrans{\TMstroke}{\TMstroke}{\TMleft}} (D)
            edge node[left, xshift=-2mm] {\TMtrans{\TMendtape}{\TMendtape}{\TMright}} (E);
\end{tikzpicture}
\]
आता आदान दुप्पट करणे सोपे आहे: दुप्पट करणारे यंत्र
अवस्था~$q_4$ ला जोडायचे. त्यासाठी त्या यंत्राच्या अवस्थांना
पुन्हा नावे द्यावी लागतील, म्हणजे त्या~$q_0$ ऐवजी~$q_4$
पासून सुरू होतील. अशा रीतीने दोन प्रारंभिक अवस्था राहणार नाहीत.
अंतिम अवस्था-आलेख \ref{tur:mac:cmb:fig:combined} मध्ये दाखवल्याप्रमाणे
असेल.
\begin{figure}
\[\begin{tikzpicture}[->,>=stealth',shorten >=1pt,auto,node distance=2.8cm,
                    semithick]
  \tikzstyle{every state}=[fill=none,draw=black,text=black]
  \node[initial,state] (A)              {$q_0$};
  \node[state]         (B) [right of=A] {$q_1$};
  \node[state]         (C) [right of=B] {$q_2$};
  \node[state]         (D) [below left of=C] {$q_3$};
  \node[state]         (E) [below left of=D] {$q_4$};
  \node[state]         (2) [right of=E] {$q_5$};
  \node[state]         (3) [right of=2] {$q_6$};
  \node[state]         (4) [below of=3] {$q_7$};
  \node[state]         (5) [left of=4]  {$q_8$};
  \node[state]         (6) [left of=5]  {$q_9$};
  \path (A) edge node {\TMtrans{\TMblank}{\TMstroke}{\TMstay}} (B)
            edge [loop above] node {\TMtrans{\TMstroke}{\TMstroke}{\TMright}} (A)
        (B) edge [loop above] node {\TMtrans{\TMstroke}{\TMstroke}{\TMright}} (B)
            edge node {\TMtrans{\TMblank}{\TMblank}{\TMleft}} (C)
        (C) edge node {\TMtrans{\TMstroke}{\TMblank}{\TMleft}} (D)
        (D) edge [loop left] node {\TMtrans{\TMstroke}{\TMstroke}{\TMleft}} (D)
            edge node[left, xshift=-2mm] {\TMtrans{\TMendtape}{\TMendtape}{\TMright}} (E)
    (E) edge node {\TMtrans{\TMstroke}{\TMblank}{\TMright}} (2)
    (2) edge [loop above] node {\TMtrans{\TMstroke}{\TMstroke}{\TMright}} (2)
      edge node {\TMtrans{\TMblank}{\TMblank}{\TMright}} (3)
    (3) edge [loop above] node {\TMtrans{\TMstroke}{\TMstroke}{\TMright}} (3)
        edge node {\TMtrans{\TMblank}{\TMstroke}{\TMright}} (4)
    (4) edge [loop below] node {\TMtrans{\TMblank}{\TMstroke}{\TMleft}} (4)
        edge node {\TMtrans{\TMstroke}{\TMstroke}{\TMleft}} (5)
    (5) edge [loop below]  node {\TMtrans{\TMstroke}{\TMstroke}{\TMleft}} (5)
        edge              node {\TMtrans{\TMblank}{\TMblank}{\TMleft}} (6)
    (6) edge [loop below] node {\TMtrans{\TMstroke}{\TMstroke}{\TMleft}} (6)
        edge              node {\TMtrans{\TMblank}{\TMblank}{\TMright}} (E);
\end{tikzpicture}
\]
\caption{बेरीज-यंत्र आणि दुप्पट करणारे यंत्र यांचे संयोजन}

\label{tur:mac:cmb:fig:combined}
\end{figure}
\end{ex}

\begin{prop}
    $M$ आणि $M'$ ही शिस्तबद्ध यंत्रे अनुक्रमे $f\colon
    \Nat^k \to \Nat$ आणि $f'\colon \Nat \to \Nat$ ही फलने
    संगणित करत असतील, तर $M \frown M'$ हेही शिस्तबद्ध असून
    ~$\comp{f}{f'}$ संगणित करते.
\end{prop}

\begin{proof}
    $M$ शिस्तबद्ध असल्यामुळे ते
    प्रदान~$f(n_1,\dots,n_k) = m$ देऊन थांबते तेव्हा शीर्ष
    घर~$1$ वाचत असते. आता यंत्र~$M'$ च्या प्रारंभिक अवस्थेत
    गेल्यावर प्रदान $f'(m)$ देऊन यंत्र पुन्हा घर~$1$ वरच थांबते.
    \ref{tur:mac:dis:defn:disciplined} मधील इतर अटीही पूर्ण होतात.
\end{proof}

\begin{prob}
    \ref{tur:mac:dis:prob:disc-succ} मधील यंत्र~$M$ घेऊन
    $M \frown M$ रचा आणि $f(x) = x+2$ संगणित करणारे
    शिस्तबद्ध ट्यूरिंग यंत्र द्या.
\end{prob}










\section{ट्यूरिंग यंत्रांचे विविध प्रकार}\label{tur:mac:var:sec}

खरे तर ट्यूरिंग यंत्रांची व्याख्या अनेक प्रकारे करता येते;
आपली व्याख्या हा त्यांपैकी एकच प्रकार आहे. काही बाबतींत ती
इतर व्याख्यांपेक्षा अधिक मोकळी आहे. आपण कोणतीही सांत वर्णमाला
चालू देतो; अधिक मर्यादित व्याख्येत फितीवरील फक्त दोन चिन्हे,
$\TMstroke$ आणि~$\TMblank$, चालतील. आपण यंत्राला एकाच वेळी
फितीवर चिन्ह लिहू आणि शीर्ष हलवू देतो; इतर व्याख्यांत लिहिणे
किंवा शीर्ष हलवणे यांपैकी एकच क्रिया करता येते. शीर्ष न सरकवता
लिहिण्याची शक्यताही आपण ठेवतो; इतर व्याख्यांत $\TMstay$
ही ``सूचना'' नसते. दुसरीकडे, काही बाबतींत आपली व्याख्या अधिक
मर्यादित आहे. फीत फक्त एका दिशेला अनंत आहे असे आपण मानले;
इतर व्याख्यांत ती डावीकडे आणि उजवीकडे दोन्ही दिशांना अनंत
असते. एवढेच नव्हे, तर कितीही सांत संख्येने स्वतंत्र फिती किंवा घरांची
अनंत जाळीही चालू देता येईल. आपण ट्यूरिंग यंत्राचा सूचनासंच
संक्रमण फलनाने दर्शवतो; इतर व्याख्यांत संक्रमण संबंध वापरतात,
ज्यामुळे एखाद्या परिस्थितीत यंत्रासाठी एकाहून अधिक सूचना
शक्य असतात.

व्याख्येतील हा शेवटचा बदल विशेष लक्षवेधक आहे. आपल्या व्याख्येत
यंत्र~$q$ अवस्थेत~$\sigma$ चिन्ह वाचत असताना
$\delta(q, \sigma)$ पुढील चिन्ह, अवस्था आणि शीर्षाची जागा
निश्चित करते. पण सूचनासंच हा सध्याच्या अवस्था-चिन्ह जोड्या
$\tuple{q,
  \sigma}$ आणि पुढील अवस्था-चिन्ह-दिशा त्रयी $\tuple{q', \sigma',
  D}$ यांमधील संबंध म्हणून घेतला, तर यंत्राची कृती एकमेवपणे
निश्चित होणार नाही: सूचना-संबंधात $\tuple{q,
  \sigma, q', \sigma', D}$ आणि $\tuple{q, \sigma, q'', \sigma'', D'}$
या दोन्ही नोंदी असू शकतात. अशा वेळी आपल्याला
\emph{अनिर्धारक} ट्यूरिंग यंत्र मिळते. संगणकीय गुंतागुंत
सिद्धांतात अशा यंत्रांना महत्त्वाचे स्थान आहे.

ट्यूरिंग यंत्र कधी थांबते याबाबतही वेगवेगळे संकेत आहेत:
संक्रमण फलन अपरिभाषित असेल तेव्हा ते थांबते असे आपण मानतो;
इतर व्याख्यांत यंत्राने खास निर्दिष्ट थांबण्याच्या अवस्थेत
असणे आवश्यक असते. अशी निर्दिष्ट अवस्था मागितल्याने ट्यूरिंग
यंत्रांना काय संगणित करता येते यावर निर्बंध येत नाही, हे आपण
\ref{tur:mac:hal:sec} मध्ये स्पष्ट केले आहे. आपल्या यंत्रांची फीत
एकाच दिशेला अनंत असल्यामुळे कधी एखादी सूचना नीट अमलात
आणता येत नाही: यंत्र सर्वांत डावीकडील घर वाचत आहे आणि
डावीकडे सरकायची सूचना आहे. आपल्या व्याख्येनुसार शीर्ष
फितीबाहेर ``पडण्याऐवजी'' त्याच घरावर राहते; पण त्या वेळी
यंत्र थांबेल अशीही व्याख्या करता आली असती. ती व्याख्याही
सममूल्य आहे: घर~$0$ वरून डावीकडे जाण्याचा प्रयत्न केल्यावर
थांबणाऱ्या यंत्राचे वर्तन अनुकरण करण्यासाठी प्रत्येक
$\delta(q,\TMendtape) =
\tuple{q',\sigma,\TMleft}$ संक्रमण काढून टाकता येईल.
मग~$\TMendtape$ वरून डावीकडे जाण्याचा प्रयत्न करण्याऐवजी
यंत्र थांबेल.\footnote{ही पद्धत \emph{अगदी तशी} चालत नाही,
कारण घर~$0$ सोडून इतर घरांतही $\TMendtape$ लिहिण्यास किंवा
वाचण्यास मनाई नाही (\ref{tur:mac:una:ex:mover} पाहा). अशा वापरासाठी
दुसरे~$\TMendtape'$ चिन्ह जोडून ही अडचण दूर करता येते.
तसेच मूळ यंत्राने घर शून्यवरील चिन्ह बदलले असल्यास त्या सीमास्थानाची
ओळख जतन करण्यासाठी स्वतंत्र स्थिर सीमा-चिन्ह किंवा योग्य सांकेतिक
व्यवस्था आवश्यक आहे; फक्त इतर घरांसाठी दुसरे चिन्ह वापरणे पुरेसे नाही.}


संख्या निरूपित करण्याचेही वेगवेगळे मार्ग आहेत (आणि त्यामुळे
ट्यूरिंग यंत्र संगणित करत असलेल्या आदान-प्रदान फलनाचे रूपही
बदलते): आपण एकचिन्ही निरूपण वापरतो, पण द्विमान निरूपणही
वापरता येते. त्यासाठी $\TMblank$ आणि~$\TMendtape$ यांच्या
जोडीला आणखी दोन चिन्हे लागतात.

आता एक महत्त्वाची गोष्ट: यांपैकी कोणत्याही बदलामुळे नेमकी
कोणती फलने ट्यूरिंग-संगणनक्षम आहेत हे बदलत नाही.
\emph{एका व्याख्येनुसार एखादे फलन ट्यूरिंग-संगणनक्षम असेल,
तर त्या सर्व व्याख्यांनुसार ते ट्यूरिंग-संगणनक्षम असते.}

याची सविस्तर पडताळणी येथे करणार नाही. एक उदाहरण पाहू:
दोन दिशांना अनंत असलेली फीत किंवा अनेक फिती वापरू दिल्याने
संगणनशक्ती वाढत नाही. कारण साधारणपणे एकाच दिशेला अनंत असलेली
एक फीत वापरणारे ट्यूरिंग यंत्र अनेक फितींचे किंवा दोन्ही
दिशांना अनंत असलेल्या फितींचे अनुकरण करू शकते. उदाहरणार्थ,
अतिरिक्त अवस्था आणि सूचना वापरून अनेक फितींच्या किंवा
दोन दिशांना अनंत असलेल्या फितीच्या यंत्राचा कार्यक्रम, एका
दिशेला अनंत असलेली एकच फीत वापरणाऱ्या यंत्राच्या कार्यक्रमात
``रूपांतरित'' करता येतो. रूपांतरित यंत्र सम क्रमांकांची घरे
फीत~$1$ वरील घरांसाठी (किंवा द्विदिश अनंत फितीच्या ``धन''
बाजूच्या घरांसाठी), आणि विषम क्रमांकांची घरे फीत~$2$ वरील
घरांसाठी (किंवा त्या फितीच्या ``ऋण'' बाजूच्या घरांसाठी)
वापरू शकते.











\section{चर्च--ट्यूरिंग प्रबंध}\label{tur:mac:ctt:sec}

प्रभावी प्रक्रियेच्या संकल्पनेला अचूक पर्याय म्हणून ट्यूरिंग
यंत्रे मांडली जातात. ट्यूरिंग यांच्या मते, प्रभावी प्रक्रिया आणि
ट्यूरिंग यंत्र या दोन्ही संकल्पना समजलेल्या कोणालाही प्रभावी
प्रक्रियेने करता येणारी कोणतीही गोष्ट ट्यूरिंग यंत्रानेही करता
येईल असे अंतर्ज्ञान होईल. प्रभावी प्रक्रियेच्या संकल्पनेचे
इतर सर्व प्रस्तावित अचूक पर्याय व्याप्तीच्या दृष्टीने ट्यूरिंग
यंत्राच्या संकल्पनेशी सममूल्य ठरतात, ही गोष्ट या दाव्याला
आधार देते---म्हणजे ते नेमका तोच फलनांचा संच संगणित करू
शकतात. या दाव्याला \emph{चर्च--ट्यूरिंग प्रबंध} म्हणतात.

\begin{defn}[चर्च--ट्यूरिंग प्रबंध]
\emph{चर्च--ट्यूरिंग प्रबंध} असे सांगतो की प्रभावी प्रक्रियेने
संगणित करता येणारी कोणतीही गोष्ट ट्यूरिंग-संगणनक्षम असते.
\end{defn}

चर्च--ट्यूरिंग प्रबंधाचा दोन प्रकारे आधार घेतला जातो. पहिला
वापर काहीसा आळशीपणाला सबब देणारा आहे. समजा, एखादी गोष्ट
संगणित करण्याच्या प्रभावी प्रक्रियेचे वर्णन आपल्याकडे
``छद्म-संकेतलेखनात'' आहे. अशा वेळी प्रत्यक्षात ट्यूरिंग यंत्र
बनवले नसले तरी हेच फलन एखादे ट्यूरिंग यंत्र संगणित करते,
असा दावा चर्च--ट्यूरिंग प्रबंधाच्या आधारे समर्थित करता येतो.

चर्च--ट्यूरिंग प्रबंधाचा दुसरा वापर तत्त्वज्ञानाच्या दृष्टीने
अधिक रोचक आहे. काही फलने ट्यूरिंग यंत्रांनी संगणित करता
येत नाहीत, हे दाखवता येते. त्यावरून चर्च--ट्यूरिंग प्रबंध
वापरून कोणत्याही प्रक्रियेने ती प्रभावी रीतीने संगणित करता
येत नाहीत, असा निष्कर्ष काढता येतो. कारण तशी एखादी प्रक्रिया
असती तर चर्च--ट्यूरिंग प्रबंधानुसार त्यासाठी ट्यूरिंग यंत्रही
असले असते. म्हणून एखादे फलन संगणित करणारे ट्यूरिंग यंत्र
असू शकत नाही हे सिद्ध केले, तर त्यासाठी प्रभावी प्रक्रियाही
असू शकत नाही. विशेषतः, तथाकथित थांबण्याची समस्या केवळ
ट्यूरिंग यंत्रांनीच सोडवता येत नाही असे नव्हे; ती मुळीच
प्रभावी रीतीने सोडवता येत नाही, असा दावा करण्यासाठी
चर्च--ट्यूरिंग प्रबंध वापरला जातो.



\chapter{अनिर्णेयता}\label{tur:und::chap}








\section{प्रस्तावना}\label{tur:und:int:sec}

प्रत्येक फलन, अगदी प्रत्येक अंकगणितीय फलनही, संगणनक्षम
असेलच असे नाही, हे उघड वाटू शकते. अशी फलने खूप आहेत
आणि त्यांचे वर्तन फार गुंतागुंतीचे असू शकते. उदाहरणार्थ,
किरणोत्सर्गी कणांच्या क्षयावरून किंवा इतर अव्यवस्थित अथवा
यादृच्छिक वर्तनावरून परिभाषित केलेली फलने. समजा, एखाद्या
दिलेल्या क्षणापासून आपण एका सेकंदाच्या अंतरांची मोजणी सुरू
करतो आणि त्या आरंभीच्या क्षणानंतरच्या $n$-व्या एक-सेकंदाच्या
अंतरात विश्वातील जितक्या कणांचा क्षय होतो ती संख्या $f(n)$
या फलनाची किंमत म्हणून ठरवतो. हे फलन कधीही संगणित
करता येणार नाही असे वाटण्यासारखे आहे.

मात्र अशी फलने कशी संगणित करावीत याची कल्पना न सुचणे
ही एक गोष्ट आहे; ती संगणनक्षम नाहीत हे प्रत्यक्ष सिद्ध
करणे ही दुसरी गोष्ट आहे. खरे तर, अत्यंत गुंतागुंतीची
वाटणारी फलनेही अमूर्त अर्थाने संगणनक्षम असू शकतात.
उदाहरणार्थ, विश्वाचा कालावधी मर्यादित आहे असे समजा---काही
विश्वशास्त्रीय सिद्धांतांच्या भाकिताप्रमाणे फार दूरच्या भविष्यात
विश्व आकुंचन पावून एका बिंदूत सामावेल. मग त्या आरंभीच्या
क्षणापासून $f(n)$ परिभाषित असेल असे फक्त मर्यादित
(पण अविश्वसनीयरीत्या मोठ्या) संख्येने सेकंद असतील.
आणि केवळ मर्यादित संख्येच्या आदानांवर परिभाषित असलेले
कोणतेही फलन आंशिक संगणनक्षम असते: त्याची फलिते एका मोठ्या
सारणीमध्ये नोंदवता येतील किंवा ट्यूरिंग यंत्राच्या एका
अतिशय मोठ्या अवस्था-संक्रमण आलेखात सांकेतिक करता येतील.
यादीबाहेरच्या आदानांवर कार्यक्रम थांबणार नाही अशी व्यवस्था करता येते.


ज्या फलनांच्या किंमती होय/नाही या प्रश्नांची उत्तरे देतात,
अशा फलनांच्या विशेष प्रकारांत आपल्याला अनेकदा रस असतो.
उदाहरणार्थ, “$n$ ही अविभाज्य संख्या आहे का?” हा प्रश्न
पुढील फलनाशी संबंधित आहे:
\[\fn{isprime}(n) = \begin{cases}
  1 & \text{जर\/ $n$ अविभाज्य असेल तर}\\
  0 & \text{अन्यथा.}
  \end{cases}
\]
एखाद्या होय/नाही प्रश्नाशी संबंधित $1/0$-मूल्यी फलन
प्रभावी रीतीने संगणनक्षम असेल, तर त्या प्रश्नाचा
\emph{प्रभावी रीतीने निर्णय करता येतो} असे आपण म्हणतो.

प्रभावी रीतीने संगणित करता न येणारी फलने किंवा प्रभावी
रीतीने निर्णय करता न येणाऱ्या समस्या अस्तित्वात आहेत,
हे गणिती पद्धतीने सिद्ध करण्यासाठी संगणनाचे विशिष्ट
प्रतिमान निश्चित करणे आवश्यक आहे. त्यानंतर त्या
प्रतिमानाला काही फलने संगणित करता येत नाहीत किंवा
काही समस्यांचा निर्णय करता येत नाही, हे दाखवावे लागते.
उदाहरणार्थ, सर्वच फलने ट्यूरिंग यंत्रांनी संगणित करता
येत नाहीत आणि सर्वच समस्यांचा निर्णय ट्यूरिंग यंत्रांनी
करता येत नाही, हे आपण दाखवू शकतो. मग चर्च--ट्यूरिंग
प्रबंधाच्या आधारे केवळ ट्यूरिंग यंत्रे सर्व फलने संगणित
करण्याइतकी सामर्थ्यवान नाहीत असेच नव्हे, तर कोणतीही
प्रभावी प्रक्रियाही तसे करू शकत नाही, असा निष्कर्ष
काढता येतो.

असे निषेधात्मक निष्कर्ष सिद्ध करण्याची गुरुकिल्ली म्हणजे
ट्यूरिंग यंत्रांनाही संख्या देता येतात ही गोष्ट. याचा
सर्वांत सोपा मार्ग म्हणजे त्यांचे प्रगणन करणे: ट्यूरिंग
यंत्रे आणि त्यांचे कार्यक्रम लिहिण्याची एखादी निश्चित
पद्धत ठरवून त्यांची पद्धतशीर यादी करता येते. असे
करता येते हे लक्षात आल्यावर, केवळ संचसंख्येच्या
विचारांवरून ट्यूरिंग-असंगणनक्षम फलने अस्तित्वात आहेत
हे कळते: $\Nat$ पासून~$\Nat$ मध्ये जाणाऱ्या फलनांचा
संच (खरे तर, $\Nat$ पासून $\{0, 1\}$ मध्ये जाणाऱ्या
फलनांचा संचसुद्धा) अगणनीय म्हणजे अगणनीय
आहे; परंतु सर्व ट्यूरिंग यंत्रांचे प्रगणन करता येत
असल्यामुळे ट्यूरिंग-संगणनक्षम फलनांचा संच केवळ
गणनीय अनंत म्हणजे गणनीय अनंत आहे.

याशिवाय, संगणित करता येत नाहीत किंवा ज्यांचा निर्णय
करता येत नाही अशा अनुक्रमे \emph{विशिष्ट} फलनांची
आणि समस्यांचीही व्याख्या करता येते. त्यांपैकी
एक म्हणजे तथाकथित \emph{थांबण्याची समस्या.}
ट्यूरिंग यंत्रांच्या सूचना यादीत दिल्यास त्यांचे
मर्यादित वर्णन करता येते. ट्यूरिंग यंत्राचे असे
वर्णन, म्हणजेच ट्यूरिंग यंत्राचा कार्यक्रम, अर्थातच
दुसऱ्या ट्यूरिंग यंत्राला आदान म्हणून देता येते.
म्हणून इतर ट्यूरिंग यंत्रांबद्दलच्या प्रश्नांचा
निर्णय करणाऱ्या ट्यूरिंग यंत्रांचा विचार करता येतो.
विशेष रोचक प्रश्न असा आहे: “दिलेल्या ट्यूरिंग यंत्राला
आदान~$n$ दिल्यावर ते अखेरीस थांबते का?”
या प्रश्नाचा निर्णय करू शकणारे ट्यूरिंग यंत्र असते
तर बरे झाले असते: ट्यूरिंग यंत्रे अनंत चक्रात अडकत
नाहीत वगैरे तपासणारे गुणवत्ता-नियंत्रण यंत्र म्हणून
त्याचा विचार करा. ट्यूरिंग यांनी सिद्ध केलेली
रोचक गोष्ट म्हणजे असे यंत्र असूच शकत नाही.
ट्यूरिंग यंत्र~$M$ चे वर्णन आणि एखादी
संख्या~$n$ यांचे आदान दिल्यावर $M$ हे यंत्र
आदान $n$ वर थांबले असते की नाही यानुसार नेहमी
थांबून अनुक्रमे $1$ किंवा $0$ फलित देणारे एकही
ट्यूरिंग यंत्र असू शकत नाही.

अनिर्णेय समस्यांची विशिष्ट उदाहरणे मिळाल्यावर इतर
समस्याही अनिर्णेय आहेत हे दाखवण्यासाठी त्यांचा
उपयोग करता येतो. उदाहरणार्थ, “प्रथम-क्रम
सूत्र~$\varphi$ वैध आहे का?” हा एक प्रसिद्ध
अनिर्णेय प्रश्न आहे. आदान म्हणून प्रथम-क्रम
सूत्र~$\varphi$ दिल्यावर $\varphi$ वैध आहे की
नाही यानुसार $1$ किंवा $0$ फलित देत नेहमी
थांबेल असे कोणतेही ट्यूरिंग यंत्र नाही.
ऐतिहासिकदृष्ट्या, ही समस्या प्रभावी रीतीने
सोडवण्याची प्रक्रिया शोधण्याच्या प्रश्नाला
सरळ “निर्णय समस्या” म्हटले जात असे; म्हणून
निर्णय समस्या सोडवता येत नाही असे आपण म्हणतो.
ट्यूरिंग आणि चर्च यांनी हे फलित साधारणतः
एकाच काळात स्वतंत्रपणे सिद्ध केले; त्यामुळे
त्याला चर्च--ट्यूरिंग प्रमेय असेही म्हणतात.











\section{ट्यूरिंग यंत्रांचे प्रगणन}\label{tur:und:enu:sec}

\begin{explain}
सर्व ट्यूरिंग यंत्रांचा संच गणनीय म्हणजे गणनीय
आहे, असे आपण दाखवू शकतो. प्रत्येक ट्यूरिंग यंत्राचे
मर्यादित वर्णन करता येते, ही त्यामागची बाब आहे. अवस्थांचा
संच आणि फितीवरील चिन्हसंच हे दोन्ही सांत संच आहेत.
संक्रमण फलन हे $Q \times \Sigma$ पासून
$Q \times \Sigma \times \{\TMleft, \TMright, \TMstay\}$
मधील आंशिक फलन आहे. म्हणून ते परिभाषित असलेल्या
मर्यादित संख्येच्या आदान-जोड्यांवरील त्याची मूल्ये
यादीत दिल्यास त्याचेही वर्णन करता येते.

ही मांडणी एका मर्यादेपर्यंत खरी असली, तरी त्यात
एक सूक्ष्म फरक आहे. ट्यूरिंग यंत्राच्या व्याख्येत
अवस्थांच्या संचात आणि फितीच्या वर्णमालेत कोणते
घटक असावेत यावर बंधन घातलेले नाही.
म्हणून उदाहरणार्थ, प्रत्येक वास्तव संख्येला अवस्था
म्हणून वापरणारे ट्यूरिंग यंत्र तांत्रिकदृष्ट्या
असू शकते. मात्र ट्यूरिंग यंत्राचे \emph{वर्तन}
कोणत्या वस्तू अवस्था किंवा फितीवरील चिन्हे म्हणून
वापरल्या जातात यावर अवलंबून नसते.
\ref{tur:und:enu:fig:variants} मधील दोन ट्यूरिंग यंत्रे पाहा.
\begin{figure}
\begin{center}
\begin{tikzpicture}[->,>=stealth',shorten >=1pt,auto,node distance=2.8cm,
                    semithick]
  \tikzstyle{every state}=[fill=none,draw=black,text=black]

  \node[initial,state]         (A)              {$q_0$};
  \node[state]         (B) [right of=A] {$q_1$};

  \path (A) edge [bend left] node {\TMtrans{\TMstroke}{\TMstroke}{\TMright}} (B)
        (B) edge [loop above] node {\TMtrans{\TMblank}{\TMblank}{\TMright}} (B)
            edge [bend left] node {\TMtrans{\TMstroke}{\TMstroke}{\TMright}} (A);
\end{tikzpicture}\\
\begin{tikzpicture}[->,>=stealth',shorten >=1pt,auto,node distance=2.8cm,
  semithick]
\tikzstyle{every state}=[fill=none,draw=black,text=black]

\node[initial,state]         (A)              {$s$};
\node[state]         (B) [right of=A] {$h$};

\path (A) edge [bend left] node {\TMtrans{A}{A}{\TMright}} (B)
(B) edge [loop above] node {\TMtrans{\TMblank}{\TMblank}{\TMright}} (B)
edge [bend left] node {\TMtrans{A}{A}{\TMright}} (A);
\end{tikzpicture}
\end{center}
\caption{\emph{सम} यंत्राची रूपे}
\label{tur:und:enu:fig:variants}
\end{figure}
हे दोन आलेख दोन यंत्रे दर्शवतात: $M$ ची फितीची
वर्णमाला $\Sigma = \{\TMendtape,\TMblank,\TMstroke\}$
आणि अवस्थांचा संच $\{q_0,q_1\}$ आहे; तर $M'$ ची
वर्णमाला $\Sigma' = \{\TMendtape,\TMblank,A\}$ आणि
अवस्था $\{s,h\}$ आहेत. पण त्यांच्यातील सूचना
इतर बाबतीत सारख्याच आहेत: $n$ $\TMstroke$
चिन्हांची क्रमिका दिल्यावर $n$ सम असेल तेव्हा आणि केवळ तेव्हाच
$M$ थांबते; आणि $n$ $A$ चिन्हांची क्रमिका दिल्यावर
$n$ सम असेल तेव्हा आणि केवळ तेव्हाच $M'$ थांबते. आपण फक्त
$\TMstroke$ चे नाव~$A$, $q_0$ चे नाव~$s$
आणि $q_1$ चे नाव~$h$ असे बदलले आहे. हे
उदाहरण सर्वसाधारण करता येते: चिन्हे आणि अवस्था
अशा रीतीने पुनर्नामित केल्याने एक यंत्र दुसऱ्यात
रूपांतरित होत असेल, तर त्या ट्यूरिंग यंत्रांना
एकच मानता येते. खरे तर, ट्यूरिंग यंत्राची चिन्हे
आणि अवस्था या सरळ धन पूर्णांकच मानता येतात:
$\sigma_0$ ऐवजी~$1$, $\sigma_1$ ऐवजी~$2$
इत्यादी; $\TMendtape$ म्हणजे~$1$, $\TMblank$
म्हणजे~$2$ इत्यादी. अशा प्रकारे, \emph{सम}
यंत्र \ref{tur:und:enu:fig:standard-even} मध्ये दाखवलेल्या
यंत्रात रूपांतरित होते.
\begin{figure}
\[\begin{tikzpicture}[->,>=stealth',shorten >=1pt,auto,node distance=2.8cm,
  semithick]
\tikzstyle{every state}=[fill=none,draw=black,text=black]
\node[initial,state]         (A)              {$1$};
\node[state]         (B) [right of=A] {$2$};
\path (A) edge [bend left] node {\TMtrans{3}{3}{\TMright}} (B)
(B) edge [loop above] node {\TMtrans{2}{2}{\TMright}} (B)
edge [bend left] node {\TMtrans{3}{3}{\TMright}} (A);
\end{tikzpicture}
\]
\caption{एक मानक \emph{सम} यंत्र}
\label{tur:und:enu:fig:standard-even}
\end{figure}
ज्या ट्यूरिंग यंत्राच्या अवस्था आणि चिन्हे धन पूर्णांक
असतात त्याला \emph{मानक} यंत्र म्हणता येईल; यापुढे
केवळ मानक यंत्रांचाच विचार करू.\footnote{“मानक
यंत्र” ही संज्ञा स्वतः मानक संज्ञा नाही.}

ट्यूरिंग यंत्रांचा संच गणनीय म्हणजे गणनीय
आहे, हे दाखवण्याचे आपले उद्दिष्ट होते. वरील विचार
लक्षात घेतल्यास मानक ट्यूरिंग यंत्रांचा संच
गणनीय म्हणजे गणनीय
आहे हे दाखवणे पुरेसे आहे. समजा, एक मानक ट्यूरिंग
यंत्र $M = \tuple{Q, \Sigma, q_0, \delta}$ दिले
आहे. धन पूर्णांकांच्या एका सांत क्रमिकेत त्याचे
वर्णन कसे करावे? प्रथम अवस्थांची संख्या, मग
अवस्था, चिन्हांची संख्या, चिन्हे आणि आरंभीची
अवस्था अशी यादी करू. (लक्षात घ्या, $M$ मानक
यंत्र असल्यामुळे ही सर्व धन पूर्णांक आहेत.)
मग~$\delta$ चे काय? संभाव्य आदाने, म्हणजे
$\tuple{q,\sigma}$ या जोड्या, यांचा संच सांत आहे;
कारण $Q$ आणि~$\Sigma$ सांत आहेत. म्हणून~$\delta$
मधील माहिती म्हणजे ज्या सर्व बाबतीत
$\delta(q,\sigma) = \tuple{q', \sigma', D}$ आहे,
अशा $\tuple{q, \sigma, q', \sigma', d}$
या सर्व $5$-घटकी क्रमितांची सांत यादी होय.
येथे $d$ ही दिशा~$D$ सांकेतिक
करणारी संख्या आहे (उदाहरणार्थ,
~$\TMleft$ साठी $1$, ~$\TMright$ साठी $2$
आणि ~$\TMstay$ साठी $3$).

अशा प्रकारे प्रत्येक मानक ट्यूरिंग यंत्राचे
धन पूर्णांकांच्या सांत यादीने, म्हणजे
$s_M \in (\PosInt)^*$ या क्रमिकेने, वर्णन
करता येते. उदाहरणार्थ, मानक \emph{सम}
यंत्र पुढील क्रमिकेने सांकेतिक केले जाते:
\[2, \underbrace{1, 2}_Q, 3, \overbrace{1, 2, 3}^\Sigma, 1, \underbrace{1, 3, 2, 3, 2}_{\delta(1,3) = \tuple{2,3,R}},
\overbrace{2, 2, 2, 2, 2}^{\delta(2,2) = \tuple{2,2,R}},
\underbrace{2, 3, 1, 3, 2}_{\delta(2,3) = \tuple{1,3,R}}.
\]
\end{explain}

\begin{thm}
$\Nat$ पासून~$\Nat$ मध्ये जाणारी अशी फलने
आहेत जी ट्यूरिंग-संगणनक्षम नाहीत.
\end{thm}

\begin{proof}
धन पूर्णांकांच्या सांत क्रमिकांचा संच~$(\PosInt)^*$
गणनीय म्हणजे गणनीय आहे, हे आपल्याला माहीत आहे
(\ref{sfr:siz:zigzag:prob:posint-star}). मानक ट्यूरिंग
यंत्रांच्या वर्णनांचा संच हा~$(\PosInt)^*$ चा उपसंच
असल्याने तोही गणनीय आहे. $\Nat$ पासून~$\Nat$
मध्ये जाणारे प्रत्येक ट्यूरिंग-संगणनक्षम फलन कोणत्या
ना कोणत्या (खरे तर अनेक) ट्यूरिंग यंत्रांनी संगणित
केले जाते. त्याच्या अवस्था आणि चिन्हांना धन
पूर्णांकांची नावे देऊन (विशेषतः $\TMendtape$ ला~$1$,
$\TMblank$ ला~$2$ आणि $\TMstroke$ ला~$3$)
प्रत्येक ट्यूरिंग-संगणनक्षम फलन कोणत्या ना कोणत्या
मानक ट्यूरिंग यंत्रानेही संगणित केले जाते, हे
दिसते. म्हणून $\Nat$ पासून~$\Nat$ मध्ये जाणाऱ्या
सर्व ट्यूरिंग-संगणनक्षम फलनांचा संचही गणनीय आहे.

दुसरीकडे, $\Nat$ पासून~$\Nat$ मध्ये जाणाऱ्या
सर्व फलनांचा संच गणनीय म्हणजे गणनीय नाही
(\ref{sfr:siz:red:prob:nat-nat}). जर प्रत्येक
फलन एखाद्या ट्यूरिंग यंत्राने संगणित केले
जात असेल, तर ती फलने संगणित करणाऱ्या यंत्रांची
सर्व वर्णने यादीत दिल्याने सर्व फलनांचे प्रगणन
करता आले असते. त्यामुळे काही फलने
ट्यूरिंग-संगणनक्षम नाहीत.
\end{proof}

\begin{prob}
  अवस्था आणि वर्णमालेतील चिन्हांची स्वतंत्र
  स्पष्ट यादी न देता ट्यूरिंग यंत्रांचे वर्णन
  करण्याचा मार्ग तुम्हाला सुचतो का? “मानक”
  यंत्राची स्वतःची संकल्पना तुम्ही ठरवू शकता;
  पण तुमच्या नव्या अर्थाने प्रत्येक ट्यूरिंग
  यंत्राचे संगणन एखाद्या “मानक” यंत्राने
  कसे करता येईल, हेही स्पष्ट करा.
\end{prob}











\section{सार्वत्रिक ट्यूरिंग यंत्रे}\label{tur:und:uni:sec}

\ref{tur:und:enu:sec} मध्ये प्रत्येक ट्यूरिंग यंत्राचे
पूर्णांकांच्या सांत क्रमिकेने वर्णन कसे करता येते,
याची चर्चा केली. ही क्रमिका त्या यंत्राच्या अवस्था,
वर्णमाला, आरंभीची अवस्था आणि सूचना सांकेतिक करते.
अशा सर्व वर्णनांचा संच गणनीय म्हणजे गणनीय
आहे, हेही आपण दाखवले. अशा वर्णनांचा संच
गणनीय अनंत म्हणजे गणनीय अनंत असल्यामुळे
~$\Nat$ पासून या वर्णनांवर जाणारे
आच्छादक म्हणजे आच्छादक फलन आहे.
असे आच्छादक फलन उदाहरणार्थ कांतोरची
नागमोडी पद्धत वापरून मिळवता येते. त्यामुळे
ट्यूरिंग यंत्रांच्या (वर्णनांचे) प्रगणन करता येते.
असे एक प्रभावी प्रगणन शून्य निर्देशांकापासून निश्चित केल्यावर $0$-वे, $1$-वे,
$2$-रे, \dots, $e$-वे ट्यूरिंग यंत्र असे
बोलता येते. या संख्यांना \emph{निर्देशांक} म्हणतात.


\begin{defn}
$M$ हे (आपण निश्चित केलेल्या प्रगणनातील)
$e$-वे ट्यूरिंग यंत्र असेल, तर $e$ हा~$M$ चा
\emph{निर्देशांक} आहे असे म्हणतो. $e$-व्या
ट्यूरिंग यंत्रासाठी $M_e$ असे लिहितो.
\end{defn}

एका यंत्राला एकाहून अधिक निर्देशांक असू शकतात.
उदाहरणार्थ, $M$~च्या सूचनांची यादी कोणत्या
क्रमाने करतो यावरून त्याची दोन वर्णने वेगळी
असू शकतात; अशा वेगळ्या वर्णनांचे निर्देशांकही
वेगवेगळे असतात.

महत्त्वाची बाब अशी की ट्यूरिंग यंत्रांच्या वर्णनांचे
प्रगणन अशा रीतीने करता येते की निर्देशांकावरून
यंत्र~$M$ चे वर्णन प्रभावी रीतीने संगणित करता
येईल आणि यंत्र~$M$ च्या वर्णनावरून त्याचा
एक निर्देशांकही प्रभावी रीतीने संगणित करता
येईल. मग चर्च--ट्यूरिंग प्रबंधानुसार, निर्देशांक~$e$
असलेल्या ट्यूरिंग यंत्राचे वर्णन मिळवून ते
आपल्या फितीवर फलित म्हणून लिहिणारे ट्यूरिंग
यंत्र असू शकते. हे वर्णन म्हणजे
~$\TMstroke$ चिन्हांच्या खंडांची क्रमिका असेल
(हे खंड $M_e$ चे वर्णन करणाऱ्या क्रमिकेतील
धन पूर्णांक दर्शवतील).

आता असा प्रश्न स्वाभाविकपणे पडतो: ट्यूरिंग
यंत्रांच्या निर्देशांकांवरील कोणती फलने स्वतः
ट्यूरिंग यंत्रांनी संगणित करता येतात? ट्यूरिंग
यंत्रांच्या निर्देशांकांचे कोणते गुणधर्म ट्यूरिंग
यंत्रांनी निर्णीत करता येतात? उदाहरण पाहा:
निर्देशांक~$e$ वरून निर्देशांक~$e$ असलेल्या
ट्यूरिंग यंत्रातील अवस्थांची संख्या देणारे
फलन ट्यूरिंग यंत्राने संगणित करता येते.
असे यंत्र पुढीलप्रमाणे काम करेल: सुरुवातीला
त्याच्या फितीवर $e$~$\TMstroke$ चिन्हांचा
एकच खंड असेल. प्रथम ते $e$ वरून त्याचे
वर्णन उलगडेल. आता ते वर्णन फितीवरील
~$\TMstroke$ चिन्हांच्या खंडांच्या क्रमिकेने
दर्शवलेले असेल. या क्रमिकेतील पहिला
घटक म्हणजे अवस्थांची संख्या आहे.
म्हणून पहिला $\TMstroke$ खंड सोडून इतर
सर्व पुसून टाकायचे आणि मग थांबायचे, एवढेच
उरते.

पुढील फलित विशेष उल्लेखनीय आहे:

\begin{thm}\label{tur:und:uni:thm:universal-tm} असे एक \emph{सार्वत्रिक
  ट्यूरिंग यंत्र}~$U$ आहे की आदान $\tuple{e,n}$ वर
  सुरू केल्यास ते
  \begin{enumerate}
    \item $M_e$ हे आदान~$n$ वर थांबते तेव्हा आणि केवळ तेव्हाच थांबते, आणि
    \item $M_e$ हे फलित $m$ देऊन थांबले, तर~$U$ देखील
    त्याच फलितासह थांबते.
  \end{enumerate}
  म्हणून $U$ हे $f\colon \Nat \times \Nat \pto \Nat$
  फलन संगणित करते; $M_e$ हे आदान~$n$ वर सुरू
  होऊन फलित~$m$ देत थांबले तर $f(e,n) = m$,
  आणि अन्यथा ते अपरिभाषित असते.
\end{thm}

\begin{proof}
  $U$ चा संपूर्ण सूचनासंच येथे प्रत्यक्ष लिहून काढणे अव्यवहार्य आहे, कारण
  या रचनेतील यंत्र अत्यंत गुंतागुंतीचे आहे. पण ते कसे
  काम करते याची रूपरेषा देता येते आणि मग
  चर्च--ट्यूरिंग प्रबंधाचा आधार घेता येतो.
  सुरुवातीला $U$ च्या फितीवर $e$ $\TMstroke$
  चिन्हांचा खंड आणि त्यानंतर $n$~$\TMstroke$
  चिन्हांचा खंड असतो. ते प्रथम निर्देशांक~$e$
  आदान~$n$ च्या उजवीकडे “उलगडते.” यामुळे
  यंत्र~$M_e$ च्या सूचनांचे वर्णन करणाऱ्या
  संख्यांची यादी मिळते (म्हणजे~$\TMblank$
  चिन्हांनी वेगळे केलेले $\TMstroke$ चिन्हांचे
  खंड). त्यानंतर $U$ हे $M_e$ च्या वर्णनाच्या
  उजवीकडे $M_e$ च्या आरंभीच्या अवस्थेची
  संख्या आणि संख्या~$1$ लिहिते. (पुन्हा,
  या संख्या $\TMstroke$ चिन्हांच्या खंडांनी
  एकचिन्ही रूपात दर्शवलेल्या आहेत.) पुढे,
  आदानातील $n$~$\TMstroke$ चिन्हांचा खंड
  उजवीकडे नकलते---पण प्रत्येक $\TMstroke$
  चिन्हाच्या जागी तीन $\TMstroke$ चिन्हांचा
  खंड लिहिते ($\TMstroke$ या चिन्हाची संख्या~$3$
  आहे, तर~$\TMendtape$ ची संख्या $1$ आणि
  ~$\TMblank$ ची संख्या $2$ आहे, हे आठवा).
  अशा खंडांच्या क्रमिकेच्या डाव्या टोकाला
  ($U$ च्या फितीवर ते~$\TMblank$ चिन्हांनी
  वेगळे केलेले आहेत) ते एकच~$\TMstroke$
  चिन्ह लिहिते; हाच~$\TMendtape$ चा संकेतांक आहे.
  आदान रिकामे असेल, तर त्यानंतरच्या पहिल्या रिकाम्या घराचाही
  सांकेतिक खंड लिहिते, जेणेकरून प्रारंभिक शीर्षासाठी घर उपलब्ध असेल.

  आता $U$ च्या फितीवर निर्देशांक~$e$, संख्या~$n$,
  आरंभीच्या अवस्थेचा संकेतांक (“सध्याची अवस्था”),
  आरंभीच्या शीर्षस्थानाची संख्या~$1$
  (“सध्याचे शीर्षस्थान”) आणि अनुकरणातील
  “फिती”वरील आरंभीचा मजकूर नोंदलेला असतो
  (म्हणजे~$M_e$ च्या “फिती”वरील “चिन्हांचे”
  संकेतांक दर्शवणाऱ्या~$\TMstroke$ चिन्हांच्या
  खंडांची, ~$\TMblank$ चिन्हांनी वेगळे केलेली
  क्रमिका).

  आदान~$n$ वर सुरू केलेले $M_e$ जे काम
  करेल त्याचे अनुकरण ते पुढीलप्रमाणे करते:
  \begin{enumerate}
    \item “सध्याच्या शीर्षस्थानाची” संख्या~$k$ शोधणे
    (सुरुवातीला ती~$1$ असते),
    \item डाव्या टोकाचे चिन्ह दर्शवणारा खंड शून्य क्रमांकाचा धरून
    “फितीवरील” $k$-व्या खंडाकडे जाऊन
    तेथील “चिन्ह” कोणते आहे ते पाहणे,
    \item\label{tur:und:uni:find-inst}
    सध्याची “अवस्था” आणि “चिन्ह” यांना जुळणारी
    सूचना शोधणे,
    \item “फितीवरील” $k$-व्या खंडाकडे
    परत जाऊन तेथील “चिन्हा”ऐवजी $M_e$ लिहील
    त्या चिन्हाचा संकेतांक लिहिणे,
    \item सध्याची “अवस्था” नोंदवलेली आहे
    तेथे शीर्ष नेऊन त्या संख्येऐवजी नव्या
    अवस्थेची संख्या लिहिणे,
    \item “फितीवरील स्थान” नोंदवलेले आहे
    तेथे जाऊन एक~$\TMstroke$ पुसणे किंवा
    एक~$\TMstroke$ जोडणे (सूचना अनुक्रमे
    डावीकडे किंवा उजवीकडे जायला सांगत असेल तर). सध्याचे
    शीर्षस्थान शून्य असेल आणि डावीकडे जाण्याची सूचना असेल,
    तर ते स्थान बदलत नाही. उजवीकडे सरकल्यावर सांकेतिक फितीच्या
    सध्याच्या शेवटापलीकडे घर लागल्यास नवीन रिकाम्या घराचा खंड जोडते.
    \item पुन्हा हेच करणे.\footnote{येथे आपण काही सूक्ष्म
    अडचणींचा तपशील टाळत आहोत. उदाहरणार्थ,
    “सध्याचे शीर्षस्थान” मोजणाऱ्या संख्येत
    वाढ करताना $U$ ला अतिरिक्त जागा लागू शकते;
    अशा वेळी त्याला संपूर्ण “फीत” उजवीकडे
    हलवावी लागेल.}
  \end{enumerate}
  आदान~$n$ वर सुरू केलेले $M_e$ कधीच थांबत
  नसेल, तर $U$ देखील कधीच थांबत नाही;
  म्हणून त्याचे फलित अपरिभाषित असते.

  पायरी~\ref{tur:und:uni:find-inst} मध्ये $M_e$ च्या
  वर्णनात सध्याच्या “अवस्था”/“चिन्ह” जोडीला
  जुळणारी सूचना नसेल, तर $M_e$ थांबले असते.
  तसे झाल्यास $U$ प्रथम अनुकरणातील फितीवर प्रदानाची व्याख्या
  लागू होते का ते तपासते: सुरुवातीला मूळ डाव्या टोकाचे चिन्ह,
  त्यानंतर सलग रेघा आणि त्यानंतर फक्त रिकामी घरे असली पाहिजेत.
  ही अट पूर्ण असल्यास आपल्या फितीवरील “फिती”च्या
  डावीकडचा भाग पुसते. $M_e$ च्या फितीवरील
  एक~$\TMstroke$ दर्शवणाऱ्या प्रत्येक तीन
  ~$\TMstroke$ चिन्हांच्या खंडासाठी ते
  स्वतःच्या फितीच्या डाव्या टोकाच्या चिन्हाच्या लगेच उजवीकडे एक
  $\TMstroke$ लिहिते आणि त्याच वेळी क्रमाने
  “फीत” पुसत जाते. हे संपल्यावर $U$ च्या
  फितीवर लांबी~$m$ असलेला $\TMstroke$
  चिन्हांचा एकच खंड उरतो. रेघा नसल्यास प्रदान रिकामे असते.

  वरची पूर्ण प्रदान-अट मोडली असेल—उदाहरणार्थ, मूळ प्रदानातील
  रेघेसाठी तीन~$\TMstroke$ चिन्हांचा खंड नसणे, मध्ये रिकामे घर
  असणे किंवा आरंभीचे डाव्या टोकाचे चिन्ह नसणे—तर $U$
  एक रिकामे चिन्ह आणि त्यानंतर एक रेघ असा मुद्दाम अमान्य
  प्रदानाचा आकार लिहून थांबते. या बाबतीत $U$ च्या फितीवर
  $\TMstroke$ चिन्हांचा एकच खंड नसल्याने
  त्याचे फलित नैसर्गिक संख्या नसते; म्हणजे
  या बाबतीत $f(e,n)$ अपरिभाषित असते. आरंभीच्या डाव्या टोकाच्या
  चिन्हाचा सांकेतिक खंड आणि अखेरीच्या रिकाम्या घरांचे खंड
  हे वैध प्रदानात असू शकतात; ते रेघांचे खंड नाहीत म्हणून नाकारायचे नाहीत.

\end{proof}











\section{थांबण्याची समस्या}\label{tur:und:hal:sec}

\begin{explain}
ट्यूरिंग यंत्रांच्या वर्णनांचे एक प्रगणन आपण निश्चित
केले आहे असे समजा. त्यामुळे प्रत्येक ट्यूरिंग
यंत्राला एक \emph{निर्देशांक} मिळतो: ट्यूरिंग
यंत्रांच्या वर्णनांच्या $M_0$, $M_1$, $M_2$, $M_3$,
\dots{} या प्रगणनातील त्याचे स्थान.


ट्यूरिंग-संगणनक्षम नसलेली फलने असलीच पाहिजेत,
हे आपल्याला माहीत आहे: ट्यूरिंग यंत्रांच्या
वर्णनांचा---आणि म्हणून यंत्रांचा---संच
गणनीय म्हणजे गणनीय आहे; पण
$\Nat$ पासून $\Nat$ मध्ये जाणाऱ्या सर्व
फलनांचा संच तसा नाही. मात्र संगणनक्षम
नसलेल्या विशिष्ट फलनांची उदाहरणेही देता
येतात. त्यांपैकी एक म्हणजे थांबण्याचे फलन.
\end{explain}

\begin{defn}[थांबण्याचे फलन]
 \emph{थांबण्याचे फलन}~$h$ असे परिभाषित करतात:
\[h(e,n) =
\begin{cases}
  \text{0} & \text{जर यंत्र~$M_e$ आदान $n$ वर थांबत नसेल तर} \\
  \text{1} & \text{जर यंत्र~$M_e$ आदान $n$ वर थांबत असेल तर}
\end{cases}
\]
\end{defn}

\begin{defn}[थांबण्याची समस्या]
\emph{थांबण्याची समस्या} म्हणजे (कोणत्याही $e$, $n$
साठी) ट्यूरिंग यंत्र~$M_e$ हे $n$ रेघांच्या
आदानावर थांबते का हे ठरवण्याची समस्या.
\end{defn}

\begin{explain}
~$h$ हे ट्यूरिंग-संगणनक्षम नाही, हे आपण
त्याच्याशी संबंधित~$s$ हे फलन ट्यूरिंग-संगणनक्षम
नाही असे दाखवून सिद्ध करू. हा पुरावा कोणत्याही
ट्यूरिंग यंत्राचे संगणन शिस्तबद्ध ट्यूरिंग
यंत्राने करता येते (\ref{tur:mac:dis:sec}),
आणि दोन ट्यूरिंग यंत्रे जोडून एक यंत्र
बनवता येते (\ref{tur:mac:cmb:sec}),
या दोन गोष्टींवर आधारलेला आहे.
\end{explain}

\begin{defn} फलन~$s$ पुढीलप्रमाणे परिभाषित आहे:
\[s(e) =
\begin{cases}
  \text{0} & \text{जर यंत्र~$M_e$ आदान $e$ वर थांबत नसेल तर} \\
  \text{1} & \text{जर यंत्र~$M_e$ आदान $e$ वर थांबत असेल तर}
\end{cases}
\]
\end{defn}

\begin{lem}
फलन~$s$ ट्यूरिंग-संगणनक्षम नाही.
\end{lem}

\begin{proof}
विरोधासाठी फलन~$s$ ट्यूरिंग-संगणनक्षम आहे असे
समजा. मग~$s$ संगणित करणारे ट्यूरिंग यंत्र~$S$
असेल. व्यापकतेला बाधा न आणता, $S$ शिस्तबद्ध
आहे असे मानता येते. विशेषतः, ते थांबते
तेव्हा पहिल्या घराकडे पाहत असते.

या यंत्राला दुसरे यंत्र~$J$ जोडता येते.
त्याला आदान~$0$ दिल्यास ते थांबते
(म्हणजे आरंभीच्या अवस्थेत, फितीच्या टोकाच्या
चिन्हाच्या उजवीकडील घराकडे पाहताना त्याला
$\TMblank$ दिसल्यास); अन्यथा ते उजवीकडे
चालतच राहते आणि कधीच थांबत नाही.
$S$ ला~$J$ जोडून तयार केलेले
$S \concat J$ हे ट्यूरिंग यंत्र आहे;
म्हणून ते कोणत्यातरी~$e$ साठी $M_e$
आहे (म्हणजे प्रगणनात कुठेतरी येते).
$M_e$ हे $e$ $\TMstroke$ चिन्हांच्या
आदानावर सुरू करा. दोनच शक्यता आहेत:
$M_e$ थांबेल किंवा थांबणार नाही.
\begin{enumerate}
\item $M_e$ हे $e$ $\TMstroke$ चिन्हांच्या
  आदानावर थांबते असे समजा. मग $s(e)
  = 1$. म्हणून~$e$ वर सुरू केलेले $S$
  थांबते आणि फितीवर फलित म्हणून एकच
  $\TMstroke$ ठेवते. मग $J$ फितीवरील
  $\TMstroke$ चिन्हापासून सुरू होते. त्या
  बाबतीत $J$ थांबत नाही. पण $M_e$ हे
  $S \concat J$ यंत्रच आहे; म्हणून $S$
  नंतर~$J$ जे करेल तेच त्याने करायला
  हवे (म्हणजे या बाबतीत उजवीकडे चालत
  राहून कधीच न थांबणे). म्हणून $M_e$
  हे $e$ $\TMstroke$ चिन्हांच्या आदानावर
  थांबू शकत नाही.

\item आता $M_e$ हे $e$ $\TMstroke$
  चिन्हांच्या आदानावर थांबत नाही असे
  समजा. मग $s(e) = 0$ आणि आदान~$e$
  वर सुरू केलेले $S$ रिकामी फीत ठेवून
  थांबते. रिकाम्या फितीवर सुरू केलेले
  $J$ लगेच थांबते. पुन्हा, $M_e$ हे
  $S$ नंतर~$J$ जे करेल तेच करते;
  म्हणून $M_e$ हे $e$ $\TMstroke$
  चिन्हांच्या आदानावर थांबलेच पाहिजे.
\end{enumerate}
दोन्ही बाबतीत आपल्या गृहीतकाला विरोध
मिळतो. त्यामुळे असे ट्यूरिंग यंत्र~$S$
असू शकत नाही: $s$ ट्यूरिंग-संगणनक्षम नाही.
\end{proof}

\begin{thm}[थांबण्याच्या समस्येची असमाधेयता]
\label{tur:und:hal:thm:halting-problem} थांबण्याची समस्या
सोडवता येत नाही; म्हणजे फलन~$h$
ट्यूरिंग-संगणनक्षम नाही.
\end{thm}

\begin{proof}
$h$ हे ट्यूरिंग-संगणनक्षम आहे, आणि समजा
ट्यूरिंग यंत्र~$H$ ते संगणित करते.
मग~$H$ वापरून~$s$ संगणित करणारे ट्यूरिंग
यंत्र बनवता येईल: प्रथम आदानाची एक प्रत
बनवा (दोन्ही प्रतींमध्ये~$\TMblank$ चिन्ह
ठेवा). मग पुन्हा आरंभी जा आणि~$H$
चालवा. पहिले काम करणारे यंत्र स्पष्टपणे
बनवता येते (\ref{tur:mac:dis:prob:copier}
पाहा); आणि $H$ अस्तित्वात असेल, तर
ते अशा प्रतिकृती तयार करणाऱ्या यंत्राला
जोडून $M_e$ हे आदान~$e$ वर थांबते
का हे ठरवणारे, म्हणजे~$s$ संगणित
करणारे, नवे यंत्र मिळेल. पण असे यंत्र
असू शकत नाही, हे आपण आधीच दाखवले
आहे. म्हणून~$h$ देखील ट्यूरिंग-संगणनक्षम
नाही.
\end{proof}

\begin{prob}
तीन-रेघांची थांबण्याची (3-Halt) समस्या
म्हणजे अन्यथा रिकाम्या असलेल्या फितीवर
तीन $\TMstroke$ चिन्हांचे आदान दिल्यावर
स्वैरपणे निवडलेले ट्यूरिंग यंत्र थांबते
की नाही हे ठरवणारी निर्णय प्रक्रिया
देण्याची समस्या. 3-Halt समस्या
सोडवता येत नाही हे सिद्ध करा.
\end{prob}

\begin{prob}
ट्यूरिंग यंत्र आणि आदान यांच्या
$M_e$ आणि~$n$ या जोड्यांसाठी,
$e \neq n$ असताना थांबण्याची समस्या
सोडवता येत असेल, तर $e = n$
असलेल्या बाबतीतही ती सोडवता
येते, हे दाखवा.
\end{prob}

\begin{prob}
आदान~$e$ ही संख्या असेल आणि ती
सर्व ट्यूरिंग यंत्रांच्या प्रगणनाद्वारे
यंत्र~$M_e$ ओळखून देत असेल,
तर थांबण्याची समस्या सोडवता
येत नाही, हे आपण सिद्ध केले.
त्याऐवजी \ref{tur:und:enu:sec}
मधील ट्यूरिंग यंत्रांची वर्णनेच
थेट आदान म्हणून दिली तर काय?
निर्देशांकांऐवजी ट्यूरिंग यंत्रांची
वर्णने आदान म्हणून घेऊन थांबण्याच्या
समस्येचा निर्णय करणारे ट्यूरिंग
यंत्र असू शकते का? का किंवा
का नाही, ते स्पष्ट करा.
\end{prob}

\begin{prob} पुढीलप्रमाणे परिभाषित
  \emph{आंशिक} फलन~$s'$ हे
  ट्यूरिंग-संगणनक्षम \emph{आहे},
  हे दाखवा:
  \[  s'(e) =
  \begin{cases}
    \text{1} & \text{जर यंत्र~$M_e$ आदान $e$ वर थांबत असेल तर}\\
    \text{अपरिभाषित} & \text{जर यंत्र~$M_e$ आदान $e$ वर थांबत नसेल तर}
  \end{cases}
  \]
\end{prob}











\section{निर्णय समस्या}\label{tur:und:dec:sec}

दिलेल्या वाक्याची वैधता ठरवण्याची
प्रभावी पद्धत असेल तेव्हा आणि केवळ तेव्हाच
प्रथम-क्रम तर्कशास्त्र \emph{निर्णेय} आहे,
असे आपण म्हणतो. प्रत्यक्षात अशी पद्धत नाही:
प्रथम-क्रम वाक्यांच्या वैधतेचा निर्णय
करण्याची समस्या सोडवता येत नाही.

हा महत्त्वाचा निषेधात्मक निष्कर्ष स्थापित
करण्यासाठी, निर्णय समस्या ट्यूरिंग यंत्राने
सोडवता येत नाही हे आपण सिद्ध करतो.
म्हणजे, फितीवर प्रथम-क्रम वाक्य असताना सुरू केले की ते वाक्य वैध
आहे की नाही यानुसार अखेरीस थांबून
$1$ किंवा~$0$ फलित देईल, असे कोणतेही
ट्यूरिंग यंत्र नाही हे दाखवतो.
चर्च--ट्यूरिंग प्रबंधानुसार, संगणनक्षम
असलेले प्रत्येक फलन ट्यूरिंग-संगणनक्षम
असते. म्हणून हे “वैधतेचे फलन”
मुळात प्रभावी रीतीने संगणित करता
येत असेल, तर ते ट्यूरिंग-संगणनक्षमही
असेल. ते ट्यूरिंग-संगणनक्षम नसेल,
तर प्रभावी रीतीनेही संगणित करता
येणार नाही.

निर्णय समस्या सोडवता येत नाही हे
सिद्ध करण्यासाठी आपण थांबण्याच्या
समस्येचे तिच्याकडे न्यूनीकरण करू.
याचा अर्थ पुढीलप्रमाणे. ट्यूरिंग
यंत्राचे वर्णन~$e$ या संकेतांकाने होत असेल,
तर ते यंत्र आदान~$w$ वर थांबते
तेव्हा~$1$ आणि अन्यथा~$0$ देणारे
फलन~$h(e,w)$ ट्यूरिंग-संगणनक्षम
नाही: आधी सिद्ध केलेली एकचिन्ही नैसर्गिक-संख्या आदानांची
थांबण्याची समस्या हे याचे विशेष प्रकरण आहे.
प्रथम-क्रम वाक्यांची वैधता ठरवणारे
ट्यूरिंग यंत्र असेल, तर~$h$ संगणित
करणारे ट्यूरिंग यंत्रही असेल, हे
आपण दाखवू. $h$ ट्यूरिंग यंत्राने
संगणित करता येत नसल्यामुळे वैधतेचा
निर्णय करणारे ट्यूरिंग यंत्रही असू
शकत नाही.

या युक्तीतील पहिली पायरी अशी आहे:
प्रत्येक आदान~$w$ आणि ट्यूरिंग
यंत्र~$M$ साठी, $M$ चा सूचनासंच
आणि आदान~$w$ दर्शवणारे
वाक्य $\mathsf{T}(M, w)$,
तसेच “$M$ अखेरीस थांबते”
हे व्यक्त करणारे वाक्य~$\mathsf{E}(M, w)$ प्रभावी रीतीने
लिहिता येतात, असे दाखवायचे;
आणि ती अशी असतील की
\begin{quote}
  $\Entails \mathsf{T}(M, w) \lif \mathsf{E}(M,w)$ हे खरे असते तेव्हा आणि केवळ तेव्हाच, जेव्हा $M$ हे आदान~$w$ वर थांबते.
\end{quote}
आपल्या पुराव्याचा मुख्य भाग ही
वाक्ये $\mathsf{T}(M, w)$ आणि~$\mathsf{E}(M, w)$
यांचे वर्णन करण्यात आणि
$\mathsf{T}(M, w) \lif \mathsf{E}(M, w)$ हे
आदान~$w$ वर~$M$ थांबते तेव्हा आणि केवळ तेव्हाच
वैध असते हे पडताळण्यात जाईल.











\section{ट्यूरिंग यंत्रांचे निरूपण}\label{tur:und:rep:sec}

\begin{explain}
ट्यूरिंग यंत्रे आणि त्यांचे वर्तन प्रथम-क्रम तर्कशास्त्रातील
वाक्याने निरूपित करण्यासाठी योग्य भाषा
परिभाषित करावी लागते. त्या भाषेचे दोन भाग आहेत:
यंत्राच्या स्थितिवर्णनांचे वर्णन करणारी विधेये
म्हणजे विधेये, आणि संगणनाच्या पायऱ्या (“क्षण”)
तसेच फितीवरील स्थाने यांना संख्या देण्यासाठीचे
पदबंध.

दोन प्रकारची विधेये म्हणजे विधेये
आपण घेऊ; दोन्ही द्विस्थानी आहेत: प्रत्येक
अवस्था~$q$ साठी विधेय~$\Obj Q_q$,
आणि फितीवरील प्रत्येक चिन्ह~$\sigma$
साठी विधेय~$\Obj S_\sigma$.
पहिल्या प्रकाराने~$M$ ची अवस्था आणि
त्याच्या फितीवरील शीर्षाचे स्थान वर्णन
करता येते; दुसऱ्या प्रकाराने फितीवरील
मजकूर वर्णन करता येतो.

फितीवरील शीर्षाचे स्थान आणि केलेल्या
पायऱ्यांची संख्या व्यक्त करण्यासाठी
संख्या व्यक्त करण्याचा मार्ग हवा.
यासाठी स्थिरांक~$\Obj 0$
आणि उत्तरवर्ती फलन~$\prime$ हे
$1$-स्थानी फलन वापरतो. संकेताप्रमाणे
ते त्याच्या आदानाच्या \emph{नंतर}
लिहितो (आणि कंस गाळतो).

प्रत्येक संख्या $n$ साठी तिचे
~$\Lang L_M$ मधील निरूपण करणारे
मानक पद~$\num{n}$ असते; त्याला
~$n$ चा \emph{संख्यांक} म्हणतात.
$\num{0}$ म्हणजे $\Obj 0$,
$\num{1}$ म्हणजे $\Obj 0'$,
$\num{2}$ म्हणजे $\Obj 0''$,
इत्यादी. अधिक औपचारिकरीत्या:
\begin{align*}
\num{0} & = \Obj 0 \\
\num{n+1} &= \num{n}'
\end{align*}
$\num{0}$ हे पद, म्हणजे $\Obj 0$,
फितीवरील सर्वांत डावीकडील घराचे स्थान
आणि संगणनाची पहिली पायरी सुरू होण्यापूर्वीचा
क्षण (प्रारंभिक स्थितिवर्णन) या दोन्हींना
नाव देते. $\num{1}$ हे पद, म्हणजे
$\Obj 0'$, सर्वांत डावीकडील घराच्या
उजवीकडील घराला आणि संगणनाच्या
पहिल्या पायरीनंतरच्या क्षणाला नाव
देते; पुढेही याचप्रमाणे.

फितीवरील स्थानांचा क्रम (“याच्या डावीकडे”)
आणि संगणनाच्या पायऱ्यांचा क्रम (“याच्या
आधी”) हे दोन्ही व्यक्त करण्यासाठी
विधेय~$<$ सुद्धा घेऊ.

भाषा ठरल्यावर $\mathsf{T}(M,
w)$ ची “स्वयंसिद्धके” यादीत देऊ.
म्हणजे आदान~$w$ वर चालवलेल्या~$M$
चे वर्तन एकत्रितपणे वर्णन करणारी
वाक्ये. त्यांत
$\Obj 0$, $\prime$ आणि $<$ यांवर
अटी घालणारी वाक्ये,
आरंभीचे स्थितिवर्णन सांगणारी
वाक्ये, तसेच विशिष्ट सूचना
राबवल्यानंतर~$M$ चे स्थितिवर्णन
काय होते ते सांगणारी वाक्ये
असतील.
\end{explain}

\begin{defn}
  \label{tur:und:rep:defn:tm-descr}
$M = \tuple{Q, \Sigma, q_0, \delta}$
हे ट्यूरिंग यंत्र दिले असता
भाषा~$\Lang L_M$ मध्ये पुढील
गोष्टी असतात:
\begin{enumerate}
\item प्रत्येक अवस्था~$q \in Q$ साठी
  द्विस्थानी विधेय $\Obj Q_q(x, y)$.
  अंतर्ज्ञानाच्या दृष्टीने,
  $\Obj Q_q(\num{m}, \num{n})$
  याचा अर्थ “$n$ पायऱ्यांनंतर
  $M$ हे~$q$ अवस्थेत असून
  $m$-व्या घराकडे पाहत आहे.”
\item प्रत्येक चिन्ह~$\sigma\in \Sigma$
  साठी द्विस्थानी विधेय
  $\Obj S_\sigma(x, y)$.
  अंतर्ज्ञानाच्या दृष्टीने,
  $\Obj S_\sigma(\num{m},
  \num{n})$ याचा अर्थ “$n$
  पायऱ्यांनंतर $m$-व्या घरात
  चिन्ह~$\sigma$ आहे.”
\item एक स्थिरांक $\Obj 0$
\item एक एकस्थानी फलन $\prime$
\item एक द्विस्थानी विधेय $<$
\end{enumerate}
\end{defn}

आदान $w = \sigma_{i_1}\dots\sigma_{i_k}$
वर ट्यूरिंग यंत्र~$M$ कसे चालते याचे
वर्णन करणारी वाक्ये
पुढीलप्रमाणे आहेत:
\begin{enumerate}
\item संख्या आणि~$<$ यांचे वर्णन
करणारी स्वयंसिद्धके:
\begin{enumerate}
\item प्रत्येक संख्या तिच्या उत्तरवर्तीपेक्षा
लहान आहे असे सांगणारे वाक्य:
\[\lforall[x][x < x']
\]
\item $<$ संक्रमणशील आहे याची
खात्री करणारे वाक्य:
\[\lforall[x][\lforall[y][\lforall[z][
      ((x < y \land y < z) \lif x < z)]]]
\]
\end{enumerate}
\item आदानाच्या प्रारंभिक स्थितिवर्णनाचे
वर्णन करणारी स्वयंसिद्धके:
\begin{enumerate}
\item $0$~पायऱ्यांनंतर---म्हणजे
  यंत्र सुरू होण्यापूर्वी---$M$
  हे आरंभीच्या अवस्था~$q_0$
  मध्ये असून घर~$1$ कडे पाहते:
\[\Obj Q_{q_0}(\num{1}, \num{0})
\]
\item पहिल्या $k+1$ घरांत
  $\TMendtape$, $\sigma_{i_1}$,
  \dots, $\sigma_{i_k}$ ही चिन्हे
  असतात:
\[\Obj S_\TMendtape(\num{0}, \num{0}) \land
\Obj S_{\sigma_{i_1}}(\num{1}, \num{0}) \land
\dots \land
\Obj S_{\sigma_{i_k}}(\num{k}, \num{0})
\]
\item बाकी फीत रिकामी असते:
\[\lforall[x][(\num{k} < x \lif \Obj S_\TMblank(x, \num{0}))]
\]
\end{enumerate}
\item एका स्थितिवर्णनातून पुढच्या
स्थितिवर्णनाकडे जाणाऱ्या संक्रमणांचे
वर्णन करणारी स्वयंसिद्धके:

पुढे $\varphi(x, y)$ हे पुढील
रूपातील सर्व वाक्यांच्या
संयोगाने मिळते असे मानू:
\[\lforall[z][
  (((z < x \lor x < z) \land \Obj S_\sigma(z, y))
  \lif \Obj S_\sigma(z, y'))]
\]
येथे दोन चल मुक्त असल्याने हा स्वतः स्वतंत्र वाक्य
नसून सूत्रांचा रूपबंध आहे; पुढील संक्रमण-वाक्यांत
हे चल सार्वत्रिक संख्यापकांखाली येतात.

येथे $\sigma \in \Sigma$.
$\varphi(\num{m},\num{n})$ वापरून
“घर~$m$ सोडल्यास, $n+1$
पायऱ्यांनंतरची फीत $n$
पायऱ्यांनंतर होती तशीच आहे”
हे व्यक्त करतो.
\begin{enumerate}
\item \label{tur:und:rep:rep-right} प्रत्येक
  $\delta(q_i, \sigma) =
  \tuple{q_j, \sigma', \TMright}$
  सूचनेसाठी पुढील वाक्य:
\begin{align*}
& \lforall[x][\lforall[y][(
   (\Obj Q_{q_i}(x, y) \land \Obj S_{\sigma}(x, y)) \lif {}]] \\
&\qquad   (\Obj Q_{q_j}(x', y') \land \Obj S_{\sigma'}(x, y') \land
\varphi(x, y)))
\end{align*}
याचा अर्थ असा: $y$ पायऱ्यांनंतर
यंत्र~$q_i$ अवस्थेत असून
चिन्ह~$\sigma$ असलेल्या
घर~$x$ कडे पाहत असेल,
तर $y+1$ पायऱ्यांनंतर
ते घर~$x+1$ कडे पाहील,
~$q_j$ अवस्थेत असेल,
घर~$x$ मध्ये आता~$\sigma'$
असेल आणि घर~$x$ वगळता
इतर प्रत्येक घरात~$y$
पायऱ्यांनंतर होते तेच
चिन्ह असेल.

\item \label{tur:und:rep:rep-left} प्रत्येक
  $\delta(q_i, \sigma) =
  \tuple{q_j, \sigma', \TMleft}$
  सूचनेसाठी पुढील वाक्य:
\begin{align*}
& \lforall[x][\lforall[y][
    ((\Obj Q_{q_i}(x', y) \land \Obj S_{\sigma}(x', y)) \lif {}]]\\
& \qquad   (\Obj Q_{q_j}(x, y') \land \Obj S_{\sigma'}(x', y') \land
\varphi(x', y))) \land {}\\
& \lforall[y][((\Obj Q_{q_i}(\num{0}, y) \land \Obj S_{\sigma}(\num{0},
    y)) \lif {}]\\
& \qquad (\Obj Q_{q_j}(\num{0}, y') \land \Obj S_{\sigma'}(\num{0},
  y') \land \varphi(\num{0}, y)))
\end{align*}
ही रचना नीट समजून घ्या:
आता आपण “घर~$x$ कडे
पाहत असेल तर~\dots”
याऐवजी “घर $x+1$
कडे पाहत असेल तर~\dots”
असे म्हणतो. डावीकडे
जाण्याचा अर्थ पुढच्या
पायरीत यंत्र घर~$x$
कडे पाहते. पण लिहिले
जाते ते घर~$x+1$.
वजाबाकी किंवा पूर्ववर्ती
फलन आपल्याकडे नसल्यामुळे
ही रचना वापरतो.
पहिल्या अटीतील स्थिरता-नियम
लिहिले जाणारे घरच
वगळतो; डावीकडे गेलेले
नवे शीर्षस्थान नाही.


$x+1$ या रूपातील संख्या
$1$, $2$, \dots,
असतात हे लक्षात घ्या.
म्हणून घर~$0$ कडे पाहताना
डावीकडे जाण्याची सूचना
असल्याची बाब यात येत
नाही (आणि अर्थातच ते
डावीकडे जाऊ शकत नाही---त्याच
घरावर राहते). दुसऱ्या
संयुक्त अटीत ही विशेष
बाब येते: $y$ पायऱ्यांनंतर
यंत्र घर~$0$ कडे पाहत,
अवस्था $q_i$ मध्ये असेल
आणि घर~$0$ मध्ये
चिन्ह~$\sigma$ असेल, तर
$y+1$ पायऱ्यांनंतर ते
घर~$0$ कडेच पाहत असेल,
नव्या अवस्था~$q_j$ मध्ये
असेल, घर~$0$ वरील
चिन्ह $\sigma'$ असेल
आणि घर~$0$ वगळता
इतर घरांत $y$~पायऱ्यांनंतर
होती तीच चिन्हे असतील.
\item \label{tur:und:rep:rep-stay} प्रत्येक
  $\delta(q_i, \sigma) =
  \tuple{q_j, \sigma', \TMstay}$
  सूचनेसाठी पुढील वाक्य:
\begin{align*}
& \lforall[x][\lforall[y][(
   (\Obj Q_{q_i}(x, y) \land \Obj S_{\sigma}(x, y)) \lif {}]] \\
&\qquad   (\Obj Q_{q_j}(x, y') \land \Obj S_{\sigma'}(x, y') \land
\varphi(x, y)))
\end{align*}
\end{enumerate}
\end{enumerate}
ट्यूरिंग यंत्र~$M$ आणि आदान~$w$
यांच्यासाठी वरील सर्व वाक्यांच्या संयोगाला $\mathsf{T}(M, w)$ म्हणू.

~$M$ अखेरीस थांबते हे व्यक्त
करण्यासाठी “काही पायऱ्यांनंतर
संक्रमण फलन अपरिभाषित होईल”
असे सांगणारे वाक्य हवे. ज्या सर्व
$\tuple{q, \sigma}$ जोड्यांसाठी
~$\delta(q, \sigma)$ अपरिभाषित
आहे, त्यांचा संच~$X$
असू द्या. मग~$\mathsf{E}(M, w)$
हे पुढील वाक्य
असू द्या:
\[\lexists[x][\lexists[y][(\bigvee_{\tuple{q, \sigma} \in
      X}(\Obj Q_q(x, y) \land \Obj S_\sigma(x, y)))]]
\]

खास थांबण्याची अवस्था~$h$
असलेले ट्यूरिंग यंत्र
वापरले तर हे आणखी सोपे
होते: अशा वेळी पुढील
वाक्य~$\mathsf{E}(M, w)$
यंत्र अखेरीस थांबते हे
व्यक्त करते:
\[\lexists[x][\lexists[y][\Obj Q_h(x, y)]]
\]

\begin{prop}
\label{tur:und:rep:prop:mlessk}
$m < k$ असेल, तर
$\mathsf{T}(M, w) \Entails \num{m} < \num{k}$.
\end{prop}

\begin{proof}
सराव.
\end{proof}

\begin{prob}
\ref{tur:und:rep:prop:mlessk}
सिद्ध करा. (सूचना: $k-m$
वर गणितीय विगमन वापरा.)
\end{prob}











\section{निरूपणाची पडताळणी}\label{tur:und:ver:sec}

\begin{explain}
आपले निरूपण योग्य आहे हे तपासण्यासाठी दोन
गोष्टी सिद्ध कराव्या लागतील. प्रथम,
आदान~$w$ वर $M$ थांबत असेल, तर
$\mathsf{T}(M, w) \lif \mathsf{E}(M, w)$ वैध आहे
हे दाखवायचे. त्यानंतर उलट दिशाही
दाखवायची: $\mathsf{T}(M, w) \lif \mathsf{E}(M, w)$
वैध असेल, तर आदान~$w$ वर
चालवलेले $M$ खरोखर अखेरीस
थांबते.

या दोन सिद्धतांच्या युक्त्या खूप
वेगळ्या आहेत. पहिल्या फलितासाठी
प्रथम-क्रम तर्कशास्त्रातील
वाक्य, म्हणजे
$\mathsf{T}(M, w) \lif \mathsf{E}(M, w)$,
वैध आहे हे दाखवायचे. याचा
सर्वांत सोपा मार्ग म्हणजे
निष्पत्ती देणे.
मात्र आपला पुरावा सर्व $M$
आणि $w$ साठी लागू हवा;
म्हणून ज्या एका विशिष्ट
वाक्याची
निष्पत्ती
द्यायची आहे असे नाही,
तर अशी अनंत वाक्ये आहेत.
त्यामुळे $M$ आणि~$w$
कसेही असले आणि आदान~$w$
वर $M$ थांबायला कितीही
पायऱ्या लागत असल्या,
तरी $\mathsf{T}(M, w) \lif \mathsf{E}(M, w)$
याची निष्पत्ती
मिळेल, हे गणितीय विगमनाने
सिद्ध करणे हा योग्य मार्ग आहे.

अर्थातच, या विगमनाचा आधार
$M$ थांबणाऱ्या स्थितिवर्णनापर्यंत
पोहोचेपर्यंत केलेल्या पायऱ्यांची
संख्या असेल. विगमनात्मक पुराव्यात
आदान~$w$ वरील $M$ च्या
चालक्रमातील प्रत्येक पायरी~$n$
साठी $\mathsf{T}(M, w) \Entails \chi(M, w, n)$
हे सिद्ध करू. येथे $\chi(M, w, n)$
हे आदान~$w$ वर चालवलेल्या~$M$
चे $n$ पायऱ्यांनंतरचे स्थितिवर्णन
अचूक सांगते. आता $M$ आदान~$w$
वर, समजा, $n$ पायऱ्यांनंतर
थांबत असेल, तर $\chi(M, w, n)$
हे थांबणारे स्थितिवर्णन सांगेल.
असे स्थितिवर्णन सांगणारे
$\chi(M, w, n)$ असेल तेव्हा
$\chi(M, w, n) \Entails \mathsf{E}(M, w)$
हेही दाखवू. म्हणून $M$
आदान~$w$ वर थांबत असेल,
तर एखाद्या~$n$ साठी
$n$~पायऱ्यांनंतर $M$
थांबणाऱ्या स्थितिवर्णनात असेल.
त्या वेळी $\mathsf{T}(M, w)
\Entails \chi(M, w, n)$,
आणि $\chi(M, w, n)$
थांबणारे स्थितिवर्णन सांगत
असल्यामुळे $\chi(M, w, n) \Entails \mathsf{E}(M, w)$.
त्यामुळे $\mathsf{T}(M, w) \Entails \mathsf{E}(M, w)$,
म्हणजेच $\Entails \mathsf{T}(M, w)
\lif \mathsf{E}(M, w)$ मिळते.

उलट दिशेची युक्ती फार वेगळी
आहे. येथे $\Entails \mathsf{T}(M, w) \lif \mathsf{E}(M, w)$
असे गृहीत धरून $M$ आदान~$w$
वर थांबते हे सिद्ध करायचे.
गृहीतकावरून $\mathsf{T}(M, w) \Entails \mathsf{E}(M,
w)$ मिळते; म्हणजे $\mathsf{T}(M,
w)$ सत्य असलेल्या प्रत्येक
रचनेत
$\mathsf{E}(M, w)$ सत्य आहे.
म्हणून $\mathsf{T}(M, w)$ सत्य
असलेली एक रचना
~$\Struct{M}$ तयार करू:
तिचे परिभाषाक्षेत्र $\Nat$
असेल, आणि सर्व $\Obj Q_q$
व $\Obj S_\sigma$ यांचे
अर्थनिर्धारण आदान~$w$
वरील~$M$ च्या चालक्रमातील
स्थितिवर्णनांनी ठरेल.
उदाहरणार्थ,
$\Sat{M}{\Obj Q_q(\num{m}, \num{n})}$
तेव्हा आणि केवळ तेव्हाच, जेव्हा $M$ हे
आदान~$w$ वर $n$ पायऱ्या
चालवल्यानंतर अवस्था~$q$
मध्ये असून घर~$m$ कडे
पाहत असेल. आता गृहीतकानुसार
$\mathsf{T}(M, w) \Entails \mathsf{E}(M, w)$
आणि रचनेवरून
$\Sat{M}{\mathsf{T}(M,
  w)}$, म्हणून
$\Sat{M}{\mathsf{E}(M, w)}$.
पण $\Sat{M}{\mathsf{E}(M, w)}$
तेव्हा आणि केवळ तेव्हाच, जेव्हा एखाद्या
$n \in \Domain{M} = \Nat$
साठी आदान~$w$ वर
चालवलेले $M$ हे
~$n$ पायऱ्यांनंतर
थांबणाऱ्या स्थितिवर्णनात
असेल.

\end{explain}

\begin{defn}
$\chi(M, w, n)$ हे पुढील
वाक्य असू द्या:
\[\Obj Q_q(\num{m}, \num{n}) \land \Obj S_{\sigma_0}(\num{0}, \num{n})
\land \dots \land \Obj S_{\sigma_k}(\num{k}, \num{n}) \land
\lforall[x][(\num{k} < x \lif \Obj S_\TMblank(x, \num{n}))]
\]
येथे वेळ~$n$ ला~$q$
ही $M$ ची अवस्था आहे,
वेळ~$n$ ला $M$ घर~$m$
कडे पाहत आहे, वेळ~$n$
ला $0 \le i \le k$
साठी घर~$i$ मध्ये
चिन्ह~$\sigma_i$ आहे,
आणि $k$ हे वेळ~$0$
ला दिलेल्या संपूर्ण आदान-शब्दाच्या शेवटच्या
घराचा क्रमांक किंवा $n$ पायऱ्यांपर्यंत
शीर्षाने भेट दिलेल्या सर्वांत उजव्या
घराचा क्रमांक, यांपैकी अधिक मोठा क्रमांक आहे.
रिकाम्या आदानासाठी शब्दाच्या शेवटचा क्रमांक
शून्य मानतो; आरंभी शीर्ष पहिल्या घरावर असल्याने
ते घरही या मर्यादेत येते.

\end{defn}

\begin{lem}\label{tur:und:ver:lem:halt-config-implies-halt}
आदान~$w$ वर चालवलेले $M$
हे $n$ पायऱ्यांनंतर
थांबणाऱ्या स्थितिवर्णनात
असेल, तर $\chi(M, w, n) \Entails \mathsf{E}(M, w)$.
\end{lem}

\begin{proof}
आदान~$w$ वर $M$ हे $n$
पायऱ्यांनंतर थांबते असे
समजा. मग अशी अवस्था~$q$,
घर~$m$ आणि चिन्ह~$\sigma$
असतील की
\begin{enumerate}
\item $n$ पायऱ्यांनंतर $M$
  हे~$q$ अवस्थेत असून,
  ज्या घर~$m$ वर~$\sigma$
  आहे त्याकडे पाहत असेल.
\item संक्रमण फलन $\delta(q, \sigma)$
  अपरिभाषित असेल.
\end{enumerate}
$\chi(M, w, n)$ हे या
स्थितिवर्णनाचे वर्णन आहे;
त्यात $\Obj Q_{q}(\num{m}, \num{n})$
आणि $\Obj
S_{\sigma}(\num{m}, \num{n})$
ही संयुक्त विधाने असतील.
त्यांवरून एकत्रितपणे
$\mathsf{E}(M,
w)$ मिळते:
\[\lexists[x][\lexists[y][(\bigvee_{\tuple{q, \sigma} \in
      X}(\Obj Q_q(x, y) \land \Obj S_\sigma(x, y)))]]
\]
कारण $\tuple{q,\sigma} \in X$
असल्याने
$\Obj Q_{q}(\num{m}, \num{n}) \land \Obj S_{\sigma}(\num{m},
\num{n}) \Entails \bigvee_{\tuple{q, \sigma} \in X} (\Obj Q_q(\num{m},
\num{n}) \land \Obj S_{\sigma}(\num{m}, \num{n}))$.

\end{proof}

\begin{explain}
म्हणून $M$ आदान $w$
वर थांबत असेल, तर
एखाद्या~$n$ साठी
$\chi(M,
w, n) \Entails \mathsf{E}(M,w)$.
आता चालक्रम पोहोचतो अशा कोणत्याही वेळ~$n$
साठी $\mathsf{T}(M, w) \Entails \chi(M, w, n)$
हे दाखवू.
\end{explain}

\begin{lem}
\label{tur:und:ver:lem:config}
प्रत्येक $n$ साठी,
$M$ चा चालक्रम $n$ पायऱ्यांनंतरच्या
स्थितिवर्णनापर्यंत पोहोचत असेल, तर
$\mathsf{T}(M, w)
\Entails \chi(M, w, n)$.
\end{lem}


\begin{proof}
विगमनाची आरंभीची पायरी:
$n = 0$ असेल, तर
$\chi(M, w, 0)$ मधील अवस्था व फितीची विधाने
$\mathsf{T}(M, w)$ च्या आरंभीच्या अटींवरून मिळतात.
आदान रिकामे नसल्यास ही त्यातीलच संयुक्त विधाने आहेत.
आदान रिकामे असल्यास शीर्ष पहिल्या घरावर आहे;
तेथील रिकामे चिन्ह शून्यापेक्षा त्याचा क्रमांक मोठा
असल्याने मिळते. पहिल्या घरापलीकडील रिकामेपण
क्रमाच्या संक्रमणशीलतेने शून्यापलीकडील
रिकामेपणावरून मिळते.

विगमनाचे गृहीतक:
$n$-व्या पायरीपूर्वी $M$
थांबलेले नसेल, तर
$\mathsf{T}(M,w) \Entails \chi(M, w, n)$.
आता $\chi(M, w, n)$
थांबणारे स्थितिवर्णन सांगत
नसेल, तर $\mathsf{T}(M, w) \Entails
\chi(M, w, n+1)$ हे
दाखवायचे आहे.

$n \ge 0$ असे समजा
आणि आदान $w$ वर
सुरू केलेले $M$ हे
$n$ पायऱ्यांनंतर
अवस्था~$q$ मध्ये असून
घर~$m$ कडे पाहत
असेल. $n$ पायऱ्यांनंतर
$M$ थांबत नसल्यामुळे
$M$ च्या कार्यक्रमात
पुढील तीनपैकी एका
रूपाची सूचना असली पाहिजे:

\begin{enumerate}
\item \label{tur:und:ver:right} $\delta(q, \sigma) = \tuple{q', \sigma', \TMright}$

\item \label{tur:und:ver:left} $\delta(q, \sigma) = \tuple{q', \sigma', \TMleft}$

\item \label{tur:und:ver:stay} $\delta(q, \sigma) = \tuple{q', \sigma', \TMstay}$
\end{enumerate}

या तीनही बाबींचा आता
एकेक करून विचार करू.

\begin{enumerate}
\item \ref{tur:und:ver:right} या रूपाची सूचना
  आहे असे समजा. अशा वेळी
  \ref{tur:und:rep:defn:tm-descr}\ref{tur:und:rep:rep-right}
  यांनुसार
\begin{align*}
& \lforall[x][\lforall[y][((\Obj Q_{q}(x,
  y) \land \Obj S_\sigma(x, y)) \lif {}]]\\
& \qquad (\Obj Q_{q'}(x',
  y') \land \Obj S_{\sigma'}(x, y') \land \varphi(x, y)))
\intertext{हे $\mathsf{T}(M,w)$ चे एक संयुक्त विधान आहे.
  त्यावरून पुढील वाक्य
  मिळते (सार्वत्रिक विधानात~$x$ ऐवजी
  $\num{m}$ आणि~$y$ ऐवजी
  $\num{n}$ बसवून):}
& (\Obj Q_{q}(\num{m}, \num{n}) \land \Obj S_{\sigma}(\num{m},
\num{n})) \lif {}\\
& \qquad (\Obj Q_{q'}(\num{m}', \num{n}') \land
\Obj S_{\sigma'}(\num{m}, \num{n}') \land \varphi(\num{m}, \num{n})).
\intertext{विगमनाच्या गृहीतकानुसार
  $\mathsf{T}(M, w) \Entails \chi(M, w, n)$,
  म्हणजे}
& \Obj Q_q(\num{m}, \num{n}) \land \Obj S_{\sigma_0}(\num{0}, \num{n})
\land \dots \land \Obj S_{\sigma_k}(\num{k}, \num{n}) \land \\
& \qquad\lforall[x][(\num{k} < x \lif \Obj S_\TMblank(x, \num{n}))]\\
\intertext{$n$ पायऱ्यांनंतर फितीवरील घर~$m$ मध्ये
  चिन्ह~$\sigma$ असल्याने, त्याच्याशी संबंधित
  संयुक्त विधान~$\Obj S_\sigma(\num{m}, \num{n})$
  आहे; म्हणून यावरून पुढील मिळते:}
& \Obj Q_{q}(\num{m}, \num{n}) \land \Obj S_{\sigma}(\num{m},
   \num{n})
\intertext{आता आपल्याला पुढील मिळते:}
& \Obj Q_{q'}(\num{m}', \num{n}') \land \Obj S_{\sigma'}(\num{m},
  \num{n}') \land {}\\
& \qquad\Obj S_{\sigma_0^{+}}(\num{0}, \num{n}') \land \dots \land
  \Obj S_{\sigma_k^{+}}(\num{k}, \num{n}') \land {}\\
& \qquad \lforall[x][(\num{k} < x \lif \Obj S_\TMblank(x, \num{n}'))]
\end{align*}
येथे वरची अधिक खूण पुढील स्थितीतील चिन्ह
दर्शवते: लिहिलेल्या घरात नव्या सूचनेतील चिन्ह,
आणि इतर घरांत आधीची चिन्हे आहेत.
याचे कारण असे: पहिली ओळ आधीच्या
सोपाधिक विधानाच्या परिणामी भागावरून
निगमन नियमाने थेट मिळते.
मधल्या ओळीत लिहिलेल्या घरातील नवे चिन्ह
पहिल्या ओळीवरूनच मिळते; जुन्या चिन्हाचे
$S_{\sigma_m}(\num{m},\num{n}')$
हे विधान राखलेले नाही. इतर घरांतील प्रत्येक
विधान हे~$\chi(M, w, n)$ मधील संबंधित
संयुक्त विधान आणि $\varphi(\num{m},
\num{n})$ यांवरून मिळते.

$m < k$ असेल, तर
$\mathsf{T}(M,w) \Proves \num{m} < \num {k}$
(\ref{tur:und:rep:prop:mlessk}),
आणि~$<$ च्या संक्रमणशीलतेवरून
$\lforall[x][(\num{k} < x \lif \num{m} < x)]$
मिळते. $m = k$ असेल,
तर $\lforall[x][(\num{k} < x \lif \num{m} < x)]$
केवळ तर्कनियमांवरून मिळते.
म्हणून शेवटची ओळ
$\chi(M, w,
n)$ मधील संबंधित संयुक्त विधान,
$\lforall[x][(\num{k} < x \lif \num{m} < x)]$,
आणि $\varphi(\num{m},
\num{n})$ यांवरून मिळते.
$m<k$ असेल, तर हेच
$\chi(M, w, n+1)$ आहे.

आता $m=k$ असे समजा.
या वेळी $n+1$ पायऱ्यांनंतर
शीर्षाने घर~$k+1$ लाही
भेट दिलेली असेल; तेच
आता भेट दिलेले सर्वांत
उजवे घर आहे. म्हणून
$\chi(M, w, n+1)$ मध्ये
$\Obj
S_\TMblank(\num{k}',\num{n}')$
हे नवे संयुक्त विधान असेल,
आणि शेवटचे संयुक्त विधान
$\lforall[x][(\num{k}' < x \lif \Obj S_\TMblank(x, \num{n}'))]$
असेल. ही दोन्ही वाक्ये
देखील निष्पन्न होतात
हे तपासावे लागेल.

आपल्याकडे आधीच
$\lforall[x][(\num{k} < x \lif \Obj S_\TMblank(x,
  \num{n}'))]$
आहे. विशेषतः त्यावरून
$\num{k} < \num{k}' \lif
\Obj S_\TMblank(\num{k}', \num{n}')$
मिळते. स्वयंसिद्धक
$\lforall[x][x <
  x']$ यावरून
$\num{k} < \num{k}'$
मिळते. सोपाधिक
निगमनाने
$\Obj
S_\TMblank(\num{k}',\num{n}')$
मिळते.

शिवाय,
$\mathsf{T}(M,w) \Proves \num{k} < \num{k}'$
असल्यामुळे~$<$ च्या
संक्रमणशीलतेच्या
स्वयंसिद्धकावरून
$\lforall[x][(\num{k}' < x \lif \Obj
  S_\TMblank(x, \num{n}'))]$
मिळते. (याची पडताळणी
सरावासाठी सोडतो.)

\item \ref{tur:und:ver:left} या रूपाची
सूचना आहे असे समजा.
मग \ref{tur:und:rep:defn:tm-descr}\ref{tur:und:rep:rep-left}
यांनुसार
\begin{align*}
& \lforall[x][\lforall[y][((\Obj Q_{q}(x', y) \land \Obj
    S_{\sigma}(x', y)) \lif {}]]\\
& \qquad (\Obj Q_{q'}(x, y') \land \Obj
  S_{\sigma'}(x', y') \land \varphi(x', y))) \land {}\\
& \lforall[y][((\Obj Q_{q_i}(\num{0}, y) \land \Obj S_{\sigma}(\num{0},
    y)) \lif {}]\\
& \qquad (\Obj Q_{q_j}(\num{0}, y') \land \Obj S_{\sigma'}(\num{0}, y')
  \land \varphi(\num{0}, y)))
\intertext{हे $\mathsf{T}(M,w)$ चे एक
  संयुक्त विधान आहे. $m>0$
  असेल, तर $l = m - 1$
  असे घ्या (म्हणजे
  $m = l+1$). वरील
  वाक्याच्या
  पहिल्या संयुक्त विधानावरून
  पुढील मिळते:}
& (\Obj Q_{q}(\num{l}', \num{n}) \land \Obj S_{\sigma}(\num{l}', \num{n}))
\lif {} \\
& \qquad (\Obj Q_{q'}(\num{l}, \num{n}') \land \Obj S_{\sigma'}(\num{l}',
\num{n}') \land \varphi(\num{l}', \num{n}))
\intertext{अन्यथा $l = m = 0$
  असे घ्या आणि दुसऱ्या
  संयुक्त विधानावरून निष्पन्न
  होणारे पुढील वाक्य पाहा:}
& ((\Obj Q_{q_i}(\num{0}, \num{n}) \land \Obj S_{\sigma}(\num{0}, \num{n})) \lif {}\\
& \qquad   (\Obj Q_{q_j}(\num{0}, \num{n}') \land \Obj S_{\sigma'}(\num{0}, \num{n}') \land
\varphi(\num{0}, \num{n})))
\intertext{शून्यावरील शाखेतील पहिली अवस्था
  आधीची, तर दुसरी पुढची अवस्था आहे.
  दोन्हीपैकी कोणत्याही
  वाक्यावरून पुढील मिळते:}
&  \Obj Q_{q'}(\num{l}, \num{n}') \land \Obj S_{\sigma'}(\num{m},
  \num{n}') \land {}\\
&\qquad  \Obj S_{\sigma_0^{+}}(\num{0}, \num{n}')\land \dots \land
  \Obj S_{\sigma_k^{+}}(\num{k}, \num{n}')  \land {} \\
& \qquad  \lforall[x][(\num{k} < x
  \lif \Obj S_\TMblank(x, \num{n}'))]
\end{align*}
पहिल्या शाखेत स्थिरता-नियम लिहिले गेलेले
मूळ घर वगळतो; डावीकडे गेलेले नवे शीर्षस्थान नाही.
येथेही वरची अधिक खूण पुढील स्थितीतील
चिन्ह दर्शवते; लिहिलेल्या घरात नवे चिन्ह
आणि इतर घरांत आधीची चिन्हे आहेत.


आधीप्रमाणेच. (पहिल्या बाबतीत
$\num{l}' \ident \num{l+1}
\ident \num{m}$ आणि
दुसऱ्या बाबतीत
$\num{l} \ident \num{0}$
हे लक्षात घ्या.) पण
हेच $\chi(M, w, n+1)$
आहे.

\item \ref{tur:und:ver:stay} ही बाब
सरावासाठी सोडतो.
\end{enumerate}
म्हणून चालक्रम पोहोचतो अशा कोणत्याही~$n$
साठी $\mathsf{T}(M, w) \Entails \chi(M, w, n)$
हे दाखवले.
\end{proof}

\begin{prob}
\ref{tur:und:ver:lem:config}
च्या पुराव्यातील
~\ref{tur:und:ver:stay}
ही बाब पूर्ण करा.
\end{prob}

\begin{prob}
$\Obj S_{\sigma_i}(\num{i}, \num{n})$
आणि $\varphi(m, n)$ यांवरून
$\Obj S_{\sigma_i}(\num{i}, \num{n}')$
ची निष्पत्ती द्या
($i \neq
m$, म्हणजे $i < m$
किंवा $m < i$,
असे गृहीत धरून).
\end{prob}

\begin{prob}
$\lforall[x][(\num{k} < x \lif \Obj
  S_\TMblank(x, \num{n}'))]$,
$\lforall[x][x < x']$,
आणि
$\lforall[x][\lforall[y][\lforall[z][ ((x < y \land y < z) \lif x <
      z)]]]$
यांवरून
$\lforall[x][(\num{k}' < x \lif \Obj
  S_\TMblank(x, \num{n}'))]$
ची निष्पत्ती द्या.
\end{prob}


\begin{lem}
\label{tur:und:ver:lem:valid-if-halt}
$M$ आदान~$w$ वर
थांबत असेल, तर
$\mathsf{T}(M, w) \lif
\mathsf{E}(M, w)$ वैध आहे.
\end{lem}

\begin{proof}
\ref{tur:und:ver:lem:config} नुसार,
चालक्रम पोहोचतो अशा कोणत्याही वेळ~$n$
ला~$M$ चे वेळ~$n$ वरील स्थितिवर्णन
सांगणारे~$\chi(M, w, n)$
हे~$\mathsf{T}(M, w)$ वरून
निष्पन्न होते. समजा,
$M$ हे $k$ पायऱ्यांनंतर
थांबते. त्या वेळी
ते कोणत्यातरी~$m \in \Nat$
साठी घर~$m$ कडे
पाहत असेल. मग
$\chi(M, w, k)$ हे~$M$
चे थांबणारे स्थितिवर्णन
सांगते. म्हणजे त्यात
$\Obj Q_q(\num{m}, \num{k})$
आणि $\Obj
S_\sigma(\num{m}, \num{k})$
ही दोन्ही संयुक्त विधाने
असतील व $\delta(q,\sigma)$
अपरिभाषित असेल. म्हणून
\ref{tur:und:ver:lem:halt-config-implies-halt}
नुसार
$\chi(M, w, k) \Entails \mathsf{E}(M,
w)$. पण
$\mathsf{T}(M, w) \Entails \chi(M, w, k)$
असल्यामुळे
$\mathsf{T}(M, w)
\Entails \mathsf{E}(M, w)$;
म्हणून $\mathsf{T}(M, w) \lif \mathsf{E}(M, w)$
वैध आहे.
\end{proof}


\begin{explain}
आपला दावा पूर्णपणे
पडताळण्यासाठी उलटी
दिशाही स्थापित करावी
लागेल: $\mathsf{T}(M, w) \lif \mathsf{E}(M, w)$
वैध असेल, तर
आदान~$w$ वर सुरू
केलेले $M$ खरोखर
थांबते.
\end{explain}

\begin{lem}
\label{tur:und:ver:lem:halt-if-valid}
$\Entails \mathsf{T}(M, w) \lif \mathsf{E}(M, w)$
असेल, तर $M$
आदान~$w$ वर थांबते.
\end{lem}

\begin{proof}
परिभाषाक्षेत्र~$\Nat$ असलेली
$\Lang L_M$-रचना
~$\Struct M$ घ्या. ती
$\Obj 0$ चा अर्थ~$0$,
$\prime$ चा अर्थ उत्तरवर्ती
फलन आणि $<$ चा अर्थ
लहान-असणे संबंध असा
लावते. $\Obj Q_q$
आणि~$\Obj S_\sigma$
ही विधेये पुढीलप्रमाणे
अर्थनिर्धारित करतात:
\begin{align*}
  \Assign{\Obj Q_q}{M} & =
\Setabs{\tuple{m, n}}{\begin{array}{ll}\text{आदान~$w$ वर सुरू होऊन $n$ पायऱ्यांनंतर,}\\ \text{$M$ हे अवस्था $q$
  मध्ये असून घर~$m$ कडे पाहते}\end{array}} \\
\Assign{\Obj S_\sigma}{M} & = \Setabs{\tuple{m, n}}{\begin{array}{ll}
\text{आदान~$w$ वर सुरू होऊन $n$ पायऱ्यांनंतर,}\\ \text{$M$ च्या घर~$m$ मध्ये
  चिन्ह~$\sigma$ आहे}\end{array}}
\end{align*}
म्हणजे, आदान~$w$ वर
सुरू केलेले $M$ प्रत्यक्षात
प्रत्येक पायरीला काय करते
हे वर्णन करणारी रचना~$\Struct{M}$
आपण बांधतो. स्पष्टपणे,
$\Sat{M}{\mathsf{T}(M, w)}$.
$\Entails \mathsf{T}(M, w) \lif \mathsf{E}(M, w)$
असेल, तर
$\Sat{M}{\mathsf{E}(M, w)}$
देखील; म्हणजे
\[\Sat{M}{\lexists[x][\lexists[y][(\bigvee_{\tuple{q, \sigma} \in
      X}(\Obj Q_q(x, y) \land \Obj S_\sigma(x, y)))]]}.
\]
$\Domain{M} = \Nat$
असल्यामुळे अशी
$m$, $n \in \Nat$
असतील की काही~$q$
आणि~$\sigma$ साठी
$\delta(q, \sigma)$
अपरिभाषित असून
$\Sat{M}{\Obj Q_q(\num{m}, \num{n}) \land \Obj S_\sigma(\num{m},
\num{n})}$
असेल. ~$\Struct M$
च्या व्याख्येनुसार,
याचा अर्थ आदान~$w$
वर सुरू केलेले $M$
हे~$n$ पायऱ्यांनंतर
अवस्था~$q$ मध्ये असून
चिन्ह~$\sigma$ वाचत आहे,
आणि संक्रमण फलन
अपरिभाषित आहे;
म्हणजे~$M$ थांबले आहे.
\end{proof}











\section{निर्णय समस्या सोडवता येत नाही}\label{tur:und:uns:sec}

\begin{thm}
\label{tur:und:uns:thm:decision-prob}
निर्णय समस्या सोडवता येत नाही:
फितीवर प्रथम-क्रम तर्कशास्त्रातील
वाक्य~$\psi$ आदान म्हणून
असताना सुरू केले की $D$ अखेरीस थांबून,
$\psi$ वैध असेल तेव्हा आणि केवळ तेव्हाच~$1$ आणि
अन्यथा~$0$ फलित देणारे
कोणतेही ट्यूरिंग यंत्र~$D$ नाही.
\end{thm}

\begin{proof}
निर्णय समस्या सोडवता येते,
म्हणजे असे ट्यूरिंग यंत्र~$D$
आहे, असे समजा. मग पुढीलप्रमाणे
थांबण्याची समस्या सोडवता येईल.
ट्यूरिंग यंत्र~$M_e$ चा संकेतांक~$e$
आणि आदान~$w$ दिले असता
संबंधित वाक्य
$\mathsf{T}(M_e, w) \lif \mathsf{E}(M_e, w)$
संगणित करून अंतखूणेच्या लगेच उजवीकडील
पहिल्या आदानघराकडे पाहत थांबणारे
ट्यूरिंग यंत्र~$E$ बांधू.

मग आदान~$e$ आणि~$w$
दिले असता $E \concat D$
प्रथम $\mathsf{T}(M_e, w) \lif
\mathsf{E}(M_e, w)$ संगणित करेल,
आणि नंतर त्या आदानावर
निर्णय-समस्येचे यंत्र~$D$
चालवेल. $\mathsf{T}(M_e, w) \lif \mathsf{E}(M_e, w)$
वैध असेल तेव्हा आणि केवळ तेव्हाच $D$
$1$ देऊन थांबेल, आणि
अन्यथा~$0$ देईल.
\ref{tur:und:ver:lem:halt-if-valid}
आणि \ref{tur:und:ver:lem:valid-if-halt}
यांनुसार $\mathsf{T}(M_e, w) \lif \mathsf{E}(M_e, w)$
वैध असते तेव्हा आणि केवळ तेव्हाच~$M_e$
आदान~$w$ वर थांबते.
म्हणून आदान~$e$ आणि~$w$
दिले असता $E\concat D$
हे~$M_e$ आदान~$w$
वर थांबत असेल तेव्हा आणि केवळ तेव्हाच
$1$ देऊन थांबेल, आणि
अन्यथा~$0$ देऊन थांबेल.
दुसऱ्या शब्दांत, $E \concat D$
थांबण्याची समस्या सोडवेल.
पण \ref{tur:und:hal:thm:halting-problem}
नुसार असे कोणतेही
ट्यूरिंग यंत्र असू शकत नाही.
\end{proof}

\begin{cor}\label{tur:und:uns:cor:undecidable-sat}
प्रथम-क्रम तर्कशास्त्रातील
कोणतेही वाक्य
पूर्ततायोग्य आहे की नाही,
हे ठरवण्याची समस्या अनिर्णेय आहे.
\end{cor}

\begin{proof}
  पूर्ततायोग्यता ट्यूरिंग यंत्र~$S$
  ठरवू शकते असे समजा. मग
  निर्णय समस्या पुढीलप्रमाणे
  सोडवता येईल: आदान म्हणून
  वाक्य~$\psi$ दिले
  असता~$\psi$ एक घर उजवीकडे
  सरकवा. घर~$1$ कडे परत
  या आणि तेथे~$\lnot$
  हे चिन्ह लिहा.


  आता ट्यूरिंग यंत्र~$S$
  चालवा. फितीवर ते अखेरीस
  $1$ ($\lnot \psi$ पूर्ततायोग्य
  असेल तर) किंवा~$0$
  ($\lnot \psi$ अपूर्ततायोग्य
  असेल तर) फलित देऊन
  थांबेल. फलित मिळाल्यावर शीर्ष पहिल्या
  आदानघरावर परत आणा. घर~$1$ वर
  $\TMstroke$ असेल, तर
  तो पुसा; घर~$1$
  रिकामे असेल, तर तेथे
  $\TMstroke$ लिहा.
  त्यानंतर थांबा.

  हे ट्यूरिंग यंत्र नेहमी
  थांबते. $\lnot
  \psi$ अपूर्ततायोग्य असेल
  तेव्हा आणि केवळ तेव्हाच ते~$1$ देते;
  अन्यथा~$0$ देते.
  कारण $\psi$ वैध असते
  तेव्हा आणि केवळ तेव्हाच~$\lnot \psi$
  अपूर्ततायोग्य असते,
  यंत्र~$\psi$ वैध असेल
  तेव्हा आणि केवळ तेव्हाच~$1$ आणि
  अन्यथा~$0$ देईल.
  म्हणजे ते निर्णय
  समस्या सोडवेल.
\end{proof}

\begin{explain}
म्हणून “$\psi$ हे प्रथम-क्रम
तर्कशास्त्रातील वैध वाक्य आहे का?” या प्रश्नाला
नेहमी बरोबर “हो” किंवा
“नाही” उत्तर देणारे
ट्यूरिंग यंत्र नाही. मात्र
“हो” हे उत्तर बरोबर असेल
तेव्हा ते देणारे ट्यूरिंग
यंत्र \emph{आहे}; उत्तर
“नाही” असेल, तर ते
थांबतच नाही. हे प्रथम-क्रम
तर्कशास्त्राच्या निर्दोषता आणि
संपूर्णता प्रमेयांवरून, तसेच
निष्पत्त्या प्रभावी
रीतीने प्रगणित करता येतात
यावरून मिळते.
\end{explain}

\begin{thm}
  \label{tur:und:uns:thm:valid-ce}
  प्रथम-क्रम वाक्यांची वैधता अर्धनिर्णेय आहे:
  फितीवर प्रथम-क्रम
  तर्कशास्त्रातील वाक्य~$\psi$ आदान म्हणून
  असताना सुरू केले की
  $\psi$ वैध असेल तेव्हा आणि केवळ तेव्हाच
  अखेरीस थांबून~$1$
  देणारे, आणि अन्यथा
  न थांबणारे ट्यूरिंग
  यंत्र~$E$ आहे. हे $E$
  यंत्र वैधतेला होकार देते.
\end{thm}

\begin{proof}
  प्रथम-क्रम तर्कशास्त्रातील
  सर्व शक्य निष्पत्त्या
  एका प्रभावी अल्गोरिदमने
  एकामागून एक तयार करता
  येतात. यंत्र~$E$ तसे
  करते; $\Proves
  \psi$ हे दाखवणारी
  निष्पत्ती
  सापडली की ते~$1$
  देऊन थांबते. निर्दोषता
  प्रमेयामुळे $E$~हे $1$
  देऊन थांबले असेल,
  तर~$\Entails \psi$.
  संपूर्णता प्रमेयामुळे
  $\Entails \psi$ असेल,
  तर~$\Proves \psi$
  हे दाखवणारी निष्पत्ती असते. $E$
  सर्व शक्य निष्पत्त्या
  क्रमाने तयार करत असल्याने
  $\Proves \psi$ दाखवणारी
  अशी निष्पत्ती त्याला
  अखेरीस सापडेल; म्हणून
  ते~$1$ देऊन थांबेल.
\end{proof}











\section{ट्राख्टेनब्रोटचे प्रमेय}\label{tur:und:tra:sec}

\begin{explain}
  \ref{tur:und:rep:sec} मध्ये ट्यूरिंग
  यंत्र~$M$ आणि आदान
  चिन्हमाला~$w$ यांच्यासाठी
  $\mathsf{T}(M,w)$ आणि~$\mathsf{E}(M,w)$
  ही वाक्ये
  परिभाषित केली. त्यानंतर
  \ref{tur:und:ver:lem:valid-if-halt}
  आणि \ref{tur:und:ver:lem:halt-if-valid}
  मध्ये आदान~$w$ वर
  सुरू केलेले~$M$ अखेरीस
  थांबते तेव्हा आणि केवळ तेव्हाच
  $\mathsf{T}(M,w) \lif \mathsf{E}(M,w)$
  वैध असते, हे दाखवले.
  थांबण्याची समस्या अनिर्णेय
  असल्यामुळे प्रथम-क्रम
  तर्कशास्त्रातील वाक्यांची वैधता आणि
  पूर्ततायोग्यताही अनिर्णेय आहेत
  (\ref{tur:und:uns:thm:decision-prob} आणि \ref{tur:und:uns:cor:undecidable-sat}).

  पण वाक्यांची वैधता आणि
  पूर्ततायोग्यता कोणत्याही
  रचनांसाठी
  परिभाषित होतात, मग त्या
  सांत असोत वा अनंत.
  एखादे वाक्य
  सांत रचनेत
  पूर्ततायोग्य आहे का
  (किंवा सर्व सांत
  रचनांमध्ये
  वैध आहे का), हे ठरवणे
  सोपे असेल असे वाटू शकते.
  निर्णय समस्या सोडवता
  येत नाही याचा पुरावा
  आपण यासाठीही रूपांतरित
  करू शकतो; तसे नसते
  हे त्यातून दिसेल.

  प्रथम \ref{tur:und:ver:lem:halt-if-valid}
  च्या पुराव्याकडे परत पाहा.
  तेथे आपण~$\mathsf{T}(M,w)$ चे
  प्रतिरूप~$\Struct M$
  तयार केले; आदान~$w$
  वर सुरू केलेले यंत्र~$M$
  नेमके काय करते ते
  त्यात वर्णिले आहे.
  त्या प्रतिरूपाचे परिभाषाक्षेत्र
  $\Nat$, म्हणजे अनंत,
  होते. पण $M$ हे
  आदान~$w$ वर खरोखर
  थांबत असेल, तर त्याच
  पद्धतीने सांत प्रतिरूप
  ~$\Struct M'$ बांधता येते.
  समजा आदान~$w$ वर
  सुरू झालेले $M$ हे
  $k$ पायऱ्यांनंतर थांबते.
  $\Domain{M'}$ हे
  $\{0, \dots, n\}$ घ्या;
  येथे $n$ ही संख्या
  $k$ पेक्षा एकाने मोठी
  संख्या आणि~$w$ च्या
  लांबीपेक्षा एकाने मोठी
  संख्या यांपैकी मोठी
  आहे. तसेच
  \[    \Assign{\prime}{M'}(x) =
    \begin{cases}
      x + 1 &\text{$x < n$ असेल तर}\\
      n &\text{अन्यथा,}
    \end{cases}
  \]
  घ्या, आणि
  $\tuple{x,y} \in \Assign{<}{M'}$
  तेव्हा आणि केवळ तेव्हाच, जेव्हा $x < y$
  किंवा $x = y = n$.
  इतर बाबतीत $\Struct{M'}$
  ची व्याख्या~$\Struct{M}$
  प्रमाणेच आहे. $\Struct{M'}$
  च्या व्याख्येनुसार,
  \ref{tur:und:ver:lem:halt-if-valid}
  च्या पुराव्यातल्याप्रमाणे
  $\Sat{M'}{\mathsf{T}(M,w)}$.
  तसेच $M$ हे आदान~$w$
  वर थांबते असे गृहीत
  धरल्यामुळे
  $\Sat{M'}{\mathsf{E}(M,w)}$.
  म्हणून~$\Struct{M'}$
  हे $\mathsf{T}(M,w) \land \mathsf{E}(M,w)$
  चे सांत प्रतिरूप आहे
  (येथे $\lif$ ऐवजी
  $\land$ वापरले आहे).


  पुराव्याचा अर्धा भाग
  झाला आहे: $M$ हे
  आदान~$w$ वर थांबले,
  तर $\mathsf{T}(M,w) \land \mathsf{E}(M,w)$
  ला सांत प्रतिरूप
  असते, हे दाखवले.
  दुर्दैवाने याचा उलट
  निष्कर्ष खरा नाही:
  काही ट्यूरिंग यंत्रे
  आदान~$w$ वर थांबत
  नाहीत, तरी
  $\mathsf{T}(M,w)
  \land \mathsf{E}(M,w)$ ला
  सांत प्रतिरूप असू शकते.
  उदाहरणार्थ, एकच
  अवस्था $q_0$ आणि
  $\delta(q_0,\TMblank) = \tuple{q_0,\TMblank,\TMstay}$
  ही एकच सूचना
  असलेले यंत्र~$M$
  घ्या. रिकाम्या आदानावर
  $w = \emptyseq$
  सुरू केल्यावर ते
  कधीच थांबत नाही:
  ते अनंत फेरीत
  राहते, पण फीत
  बदलत नाही किंवा
  शीर्ष हलवत नाही.
  सर्व स्थितिवर्णने
  सारखीच असतात
  (तीच अवस्था,
  तेच शीर्षस्थान,
  तीच फीत).
  तरी $\mathsf{T}(M,\emptyseq) \land \mathsf{E}(M,\emptyseq)$
  सत्य करणारी सांत
  रचना
  ~$\Struct{M''}$
  परिभाषित करता येते
  (सराव). पण अशा
  रचना
  वगळण्यासाठी $\mathsf{T}(M,w)$
  योग्य रीतीने बदलता
  येते.

  \begin{prob}
    $q_0$ ही एकच
    अवस्था आणि
    $\delta(q_0,\TMblank) = \tuple{q_0,\TMblank,\TMstay}$
    ही एकच सूचना
    असलेले ट्यूरिंग
    यंत्र~$M$ घ्या.
    $\Domain{M''} = \{0, 1, 2\}$,
    $\Assign{\prime}{M''}(0) =
    \Assign{\prime}{M''}(1) = 1$
    आणि $\Assign{\prime}{M''}(2) = 2$,
    तसेच $\Assign{<}{M''} =
    \{\tuple{0,1}, \tuple{1,1}, \tuple{2,2}\}$
    असे घ्या.
    $\Assign{\Obj Q_{q_0}}{M''}$,
    $\Assign{\Obj
    S_{\TMblank}}{M''}$
    आणि $\Assign{\Obj S_{\TMendtape}}{M''}$
    अशी परिभाषित करा
    की $\mathsf{T}(M,\emptyseq)$
    आणि $\mathsf{E}(M,\emptyseq)$
    दोन्ही सत्य ठरतील;
    तसे का होते ते
    स्पष्ट करा.
    सूचना: $\delta(q_0, \TMendtape)$
    अपरिभाषित आहे हे
    लक्षात घ्या.
    पुढील दोन्ही विधाने
    ~$\Struct{M''}$
    मध्ये सत्य आहेत
    याची खात्री करा:
    \begin{align*}
    & \Obj Q_{q_0}(\num{1}, \num{n}) \land \Obj S_{\TMendtape}(\num{0}, \num{n})
      \land
      \lforall[x][(\num{0} < x \lif \Obj S_\TMblank(x, \num{n}))] \text{\quad सर्व $n \in \Nat$ साठी}\\
    & \lexists[y][(\Obj Q_{q_0}(\num{0}, y) \land \Obj S_{\TMendtape}(\num{0}, y))]
    \end{align*}
  \end{prob}

\end{explain}

ट्यूरिंग यंत्र~$M$ हे
आदान $w = \sigma_{i_1}\dots\sigma_{i_k}$
वर कसे चालते, याचे
वर्णन करणारी पुढील
वाक्ये पाहा:
\begin{enumerate}
\item संख्या आणि~$<$
यांचे वर्णन करणारी
स्वयंसिद्धके (\ref{tur:und:rep:sec}
मधील~$\mathsf{T}(M,w)$ च्या
व्याख्येप्रमाणेच).
\item आरंभीच्या स्थितिवर्णनाचे
वर्णन करणारी स्वयंसिद्धके:
तीही~$\mathsf{T}(M,w)$
च्या व्याख्येप्रमाणेच.
\item एका स्थितिवर्णनातून
पुढच्या स्थितिवर्णनात
होणाऱ्या संक्रमणाचे
वर्णन करणारी स्वयंसिद्धके:

पुढील बाबींसाठी $\varphi(x, y)$
आधीप्रमाणे असू द्या,
आणि
\[  \psi(y) \ident \lforall[x][(x < y \lif \eq/[x][y])].
\]
\begin{enumerate}
\item \label{tur:und:tra:rep-right}
$\delta(q_i, \sigma) =
  \tuple{q_j, \sigma', \TMright}$
रूपाच्या प्रत्येक सूचनेसाठी
पुढील वाक्य:
\begin{align*}
& \lforall[x][\lforall[y][(
   (\Obj Q_{q_i}(x, y) \land \Obj S_{\sigma}(x, y)) \lif {}]] \\
&\qquad   (\Obj Q_{q_j}(x', y') \land \Obj S_{\sigma'}(x, y') \land
\varphi(x, y) \land \psi(y')))
\end{align*}
\item \label{tur:und:tra:rep-left}
$\delta(q_i, \sigma) =
  \tuple{q_j, \sigma', \TMleft}$
रूपाच्या प्रत्येक सूचनेसाठी
पुढील वाक्य:
\begin{align*}
& \lforall[x][\lforall[y][
    ((\Obj Q_{q_i}(x', y) \land \Obj S_{\sigma}(x', y)) \lif {}]]\\
& \qquad   (\Obj Q_{q_j}(x, y') \land \Obj S_{\sigma'}(x', y') \land
\varphi(x', y) \land \psi(y'))) \land {}\\
& \lforall[y][((\Obj Q_{q_i}(\num{0}, y) \land \Obj S_{\sigma}(\num{0},
    y)) \lif {}]\\
& \qquad (\Obj Q_{q_j}(\num{0}, y') \land \Obj S_{\sigma'}(\num{0},
  y') \land \varphi(\num{0}, y) \land \psi(y')))
\end{align*}
\item \label{tur:und:tra:rep-stay}
$\delta(q_i, \sigma) =
  \tuple{q_j, \sigma', \TMstay}$
रूपाच्या प्रत्येक सूचनेसाठी
पुढील वाक्य:
\begin{align*}
& \lforall[x][\lforall[y][(
   (\Obj Q_{q_i}(x, y) \land \Obj S_{\sigma}(x, y)) \lif {}]] \\
&\qquad (\Obj Q_{q_j}(x, y') \land \Obj S_{\sigma'}(x, y') \land
\varphi(x, y) \land \psi(y')))
\end{align*}
\end{enumerate}
येथे दिसते त्याप्रमाणे~$M$
च्या संक्रमणांचे वर्णन
करणारी वाक्ये~$\mathsf{T}(M,w)$ मधील
संबंधित वाक्यांसारखीच आहेत;
फक्त शेवटी $\psi(y')$
जोडले आहे. $\psi(y')$
मुळे “पुढच्या” स्थितिवर्णनाचा
क्रमांक $y'$ हा आधीच्या
सर्व क्रमांकांहून वेगळा
असतो: $\Obj 0$,
$\Obj 0'$, \dots.
























\end{enumerate}
$\mathsf{T}'(M, w)$ हे ट्यूरिंग
यंत्र~$M$ आणि आदान~$w$
यांच्यासाठी वरील सर्व
वाक्यांचे
संयुक्त विधान असू द्या.

\begin{lem}\label{tur:und:tra:lem:halts-sat}
  आदान~$w$ वर सुरू
  केलेले $M$ थांबत असेल,
  तर $\mathsf{T}'(M,w) \land \mathsf{E}(M,w)$
  ला सांत प्रतिरूप असते.
\end{lem}

\begin{proof}
  \ref{tur:und:ver:lem:halt-if-valid}
  च्या पुराव्यातील अनंत प्रतिरूपाच्या बांधणीप्रमाणे
  सांत रचना $\Struct{M'}$ बांधा;
  मात्र पुढील बदल करा:
  \begin{align*}
    \Domain{M'} & = \{0, \dots, n\},\\
    \Assign{\prime}{M'}(x) & =
    \begin{cases}
      x + 1 &\text{$x < n$ असेल तर}\\
      n &\text{अन्यथा,}
    \end{cases} \\
    \tuple{x,y} \in \Assign{<}{M'} &\text{$x < y$ किंवा $x = y = n$ असेल तेव्हा आणि केवळ तेव्हाच,}
  \end{align*}
  येथे $n = \max(k+1,\len{w}+1)$
  आणि~$k$ ही अशी
  लघुतम संख्या आहे
  की आदान~$w$ वर
  सुरू झालेले $M$
  हे~$k$ पायऱ्यांनंतर
  थांबले आहे.
  $\Sat{M'}{\mathsf{T}'(M,w) \land \mathsf{E}(M,w)}$
  याची पडताळणी सरावासाठी
  सोडतो.

\end{proof}

\begin{prob}
  $\Sat{M'}{\mathsf{T}'(M,w) \land \mathsf{E}(M,w)}$
  हे सिद्ध करून
  \ref{tur:und:tra:lem:halts-sat}
  चा पुरावा पूर्ण करा.
\end{prob}

\begin{lem}\label{tur:und:tra:lem:sat-halts}
  $\mathsf{T}'(M,w) \land \mathsf{E}(M,w)$
  ला सांत प्रतिरूप
  असेल, तर आदान~$w$
  वर सुरू झालेले
  $M$ थांबते.
\end{lem}

\begin{proof}
  प्रतिपरिवर्तन सिद्ध करू.
  समजा आदान~$w$
  वर सुरू झालेले
  $M$ थांबत नाही.
  $\mathsf{T}'(M,w) \land \mathsf{E}(M,w)$
  ला प्रतिरूपच नसेल,
  तर पुरावा पूर्ण.
  म्हणून~$\Struct{M}$
  हे~$\mathsf{T}'(M,w) \land \mathsf{E}(M, w)$
  चे प्रतिरूप आहे
  असे समजा. ते
  सांत असू शकत
  नाही हे दाखवू.


  \ref{tur:und:ver:lem:config}
  प्रमाणे, आदान~$w$
  वर सुरू केलेले~$M$
  हे~$n$ पायऱ्यांनंतर
  थांबलेले नसेल आणि
  किमान एक पायरी
  झालेली असेल, तर
  $\mathsf{T}'(M,w)
  \Entails \chi(M, w, n) \land \psi(\num{n})$
  हे सिद्ध करता येते.
  आदान~$w$ वर
  सुरू झालेले~$M$
  थांबत नसल्यामुळे,
  प्रत्येक सकारात्मक
  $n \in \Nat$ साठी
  $\mathsf{T}'(M,w) \Entails \chi(M, w, n) \land
  \psi(\num{n})$.
  \ref{tur:und:rep:prop:mlessk}
  नुसार सर्व~$k < n$
  साठी $\mathsf{T}'(M,w) \Entails \num{k} < \num{n}$
  हेही लक्षात घ्या.
  तसेच $\psi(\num{n})
  \Entails \num{k} < \num{n} \lif
  \eq/[\num{k}][\num{n}]$.
  म्हणून सर्व~$k < n$
  साठी $\Sat{M}{\eq/[\num{k}][\num{n}]}$;
  म्हणजे अनंत संख्यांक
  पदे~$\num{k}$ यांना
  $\Struct{M}$ मध्ये
  परस्परभिन्न मूल्ये
  असलीच पाहिजेत.
  त्यासाठी $\Domain{M}$
  अनंत असावे लागते;
  म्हणून~$\Struct{M}$
  हे~$\mathsf{T}'(M,w) \land \mathsf{E}(M, w)$
  चे सांत प्रतिरूप
  असू शकत नाही.

\end{proof}

\begin{prob}
  आदान~$w$ वर
  सुरू केलेले~$M$
  हे~$n$ पायऱ्यांनंतर
  थांबलेले नसेल आणि
  किमान एक पायरी
  झालेली असेल,
  तर $\mathsf{T}'(M,w) \Entails \psi(\num{n})$
  हे सिद्ध करून
  \ref{tur:und:tra:lem:sat-halts}
  चा पुरावा पूर्ण करा.
\end{prob}

\begin{thm}[ट्राख्टेनब्रोटचे प्रमेय]
  \label{tur:und:tra:thm:trakhtenbrodt}
  प्रथम-क्रम तर्कशास्त्रातील
  कोणत्याही वाक्याला सांत प्रतिरूप
  आहे का (म्हणजे ते
  सांत प्रतिरूपात पूर्ततायोग्य आहे का),
  हे ठरवणे अनिर्णेय आहे.
\end{thm}

\begin{proof}
  सांत प्रतिरूपात पूर्ततायोग्यता ठरवणारे
  ट्यूरिंग यंत्र~$F$
  आहे असे समजा.
  मग कोणतेही ट्यूरिंग
  यंत्र~$M$ आणि
  आदान~$w$ दिले असता
  $\mathsf{T}'(M,w) \land \mathsf{E}(M,w)$
  हे वाक्य संगणित करून
  त्याला सांत प्रतिरूप
  आहे का हे~$F$
  ने ठरवता येईल.
  \ref{tur:und:tra:lem:halts-sat} आणि \ref{tur:und:tra:lem:sat-halts}
  नुसार त्याला सांत
  प्रतिरूप असते तेव्हा आणि केवळ तेव्हाच
  आदान~$w$ वर
  सुरू झालेले $M$
  थांबते. म्हणून~$F$
  वापरून थांबण्याची
  समस्या सोडवता येईल;
  पण ती सोडवता येत
  नाही हे आपल्याला
  माहीत आहे.
\end{proof}

\begin{cor}\label{tur:und:tra:cor:fproof-incomp}
  \emph{सांत} वैधतेसाठी
  निर्दोष आणि संपूर्ण, तसेच निष्पत्तींना
  प्रभावीरीत्या प्रगणित करता येईल अशी
  कोणतीही निष्पत्ती-पद्धती असू शकत नाही.
  म्हणजे $\Proves \psi$ तेव्हा आणि केवळ तेव्हाच,
  जेव्हा प्रत्येक सांत रचना~$\Struct{M}$ मध्ये $\Sat{M}{\psi}$,
  अशी प्रभावी निष्पत्ती-पद्धती असू शकत नाही.

\end{cor}

\begin{proof}
  सराव.
\end{proof}

\begin{prob}
  \ref{tur:und:tra:cor:fproof-incomp}
  सिद्ध करा. प्रत्येक
  सांत रचनेत $\psi$ सत्य
  असेल तेव्हा आणि केवळ तेव्हाच
  $\lnot \psi$ सांत प्रतिरूपात
  पूर्ततायोग्य नसेल,
  हे लक्षात घ्या.
  \ref{tur:und:uns:thm:valid-ce}
  मधील अर्थाने सांत
  प्रतिरूपात पूर्ततायोग्यता अर्धनिर्णेय
  का आहे, हे स्पष्ट
  करा. सांत वैधतेसाठी
  निष्पत्ती प्रभावीरीत्या प्रगणित करता येणारी
  निष्पत्ती-पद्धती
  असेल, तर सांत
  प्रतिरूपात पूर्ततायोग्यता निर्णेय
  होईल, असा युक्तिवाद
  करण्यासाठी हे वापरा.
\end{prob}




\part{अपूर्णता}
\chapter{अपूर्णतेचा परिचय}\label{inc:int::chap}
\begin{editorial}
 या भागातील साहित्य अपूर्णतेच्या प्रमेयांविषयी आहे. यासाठी
 प्रथम-क्रम तर्कशास्त्राच्या भागातील साहित्य (विशेषतः सिद्धता-पद्धती)
 आणि संगणनक्षमता भागातील पुनरावर्ती फलनांविषयीचे साहित्य आवश्यक आहे.
 हे Jeremy Avigad यांच्या टिपणांवर आधारित असून Richard Zach यांनी त्यांत
 सुधारणा केल्या आहेत.
\end{editorial}









\section{ऐतिहासिक पार्श्वभूमी}\label{inc:int:bgr:sec}

या विभागात अपूर्णतेच्या प्रमेयांची पार्श्वभूमी समजण्यासाठी उपयुक्त अशा
ऐतिहासिक घडामोडींची थोडक्यात चर्चा करू. विशेषतः, गणितीय तर्कशास्त्राच्या
इतिहासाचा अगदी स्थूल आढावा घेऊ आणि मग गणिताच्या पायाभरणीच्या इतिहासाविषयी
थोडे बोलू.

\begin{digress}
  ``गणितीय तर्कशास्त्र'' या संज्ञेला दोन अर्थ आहेत. ``गणितीय'' हा शब्द
  अभ्यासविषय दर्शवतो असे मानल्यास, ``गणिताचे तर्कशास्त्र'' म्हणजे गणितीय
  युक्तिवादाची तत्त्वे. तो अभ्यासाची पद्धत दर्शवतो असे मानल्यास,
  ``तर्कशास्त्राचे गणित'' म्हणजे युक्तिवादाच्या तत्त्वांचा गणितीय अभ्यास.
  पुढील विवेचनात गणितीय तर्कशास्त्र हे दोन्ही अर्थांनी, अनेकदा एकाच वेळी,
  येते.
\end{digress}

तर्कशास्त्राचा अभ्यास मूलतः ॲरिस्टॉटलपासून सुरू झाला. तो सुमारे
384--322 \textsc{bce} या काळात जगला. त्याच्या \emph{कॅटेगरीज}, \emph{प्रायर
  ॲनॅलिटिक्स} आणि \emph{पोस्टीरियर ॲनॅलिटिक्स} या ग्रंथांत वैज्ञानिक युक्तिवादाच्या
तत्त्वांचा पद्धतशीर अभ्यास आहे; त्यात संवाक्याचाही सखोल व
पद्धतशीर अभ्यास समाविष्ट आहे.

मध्ययुगभर विद्वत् तत्त्वज्ञानावर ॲरिस्टॉटलच्या तर्कशास्त्राचे वर्चस्व होते.
अठराव्या शतकातसुद्धा कांटने त्याचे तर्कशास्त्र परिपूर्ण असून त्यात सुधारणा
करण्याची गरज नाही, असे म्हटले. पण न्यायवाक्यरचनेचा सिद्धांत गणितीय
युक्तिवादाच्या अगदी वरवरच्या पैलूंपलीकडे काही दाखवण्यासाठी अत्यंत मर्यादित
आहे. त्याच्या एक शतक आधी, न्यूटनचा समकालीन लाइब्निझ याने तार्किक
युक्तिवादासाठी पूर्ण ``कलन'' असावे अशी कल्पना केली होती. त्याने अशा कलनाची
रचना करण्याच्या दिशेने काही प्राथमिक पावलेही टाकली; त्यातून मूलतः विधानीय
तर्कशास्त्राची एक आवृत्ती वर्णिली गेली.

एकोणिसावे शतक तर्कशास्त्रासाठी निर्णायक ठरले. 1854 मध्ये जॉर्ज बूलने
\emph{विचाराचे नियम} लिहिला; त्यातील विधानीय तर्कशास्त्राचा सखोल बीजगणितीय
अभ्यास आधुनिक मांडण्यांपासून फारसा दूर नाही. 1879 मध्ये गॉटलोब फ्रेगेने
\emph{Begriffsschrift} (संकल्पनालेखन) प्रसिद्ध केला. त्यात विधानीय
तर्कशास्त्राला संख्यापक आणि संबंध यांची भर घातली असल्याने प्रथम-क्रम
तर्कशास्त्र समाविष्ट होते. प्रत्यक्षात, फ्रेगेच्या तार्किक पद्धतींत
उच्च-क्रम तर्कशास्त्र आणि त्याहून अधिक गोष्टीही होत्या. \emph{अंकगणिताचे
  मूलभूत नियम} या ग्रंथात त्याने सर्व अंकगणित केवळ तार्किक गृहीतकांपासून
आपल्या Begriffsschrift मध्ये निष्पन्न करता येईल, हे दाखवण्याचा प्रयत्न
केला. दुर्दैवाने, रसेलने 1902 मध्ये दाखवून दिल्याप्रमाणे ही गृहीतके
विसंगत निघाली. तरी तो विसंगत स्वयंसिद्धक बाजूला ठेवल्यास, फ्रेगेने आधुनिक
तर्कशास्त्र जवळजवळ एकट्याने निर्माण केले; ही विलक्षण कामगिरी होती.
बूलनंतर पिअर्स आणि श्रॉडर यांच्यासह बीजगणितीय दृष्टिकोन असलेल्या इतर
विचारवंतांनीही संख्यापकाधारित तर्कशास्त्र स्वतंत्रपणे विकसित केले.

आता गणिताच्या पायाभरणीतील घडामोडींकडे वळू. तर्कशास्त्राची गणितात महत्त्वाची
भूमिका असल्यामुळे नुकत्याच वर्णिलेल्या घडामोडींशी त्यांचा मोठा संबंध आहे.
उदाहरणार्थ, सर्व अंकगणित केवळ आपल्या तार्किक चौकटीवर आधारता येईल हे
दाखवण्याच्या स्पष्ट उद्देशाने फ्रेगेने आपले तर्कशास्त्र विकसित केले.
विशेषतः, कांटने अंकगणितातील सत्ये अनुभवपूर्व \emph{संश्लेषक} मानली होती;
फ्रेगेला त्याऐवजी ती अनुभवपूर्व \emph{विश्लेषक} आहेत, हे दाखवायचे होते.
भूमितीतील सत्ये अनुभवपूर्व संश्लेषक आहेत हे मात्र फ्रेगेने मान्य केले.


अनेकांच्या मते गणिताचा खरा उदय ग्रीकांबरोबर झाला. सुमारे 300 इ.स.पू.
मध्ये लिहिलेला युक्लिडचा \emph{मूलतत्त्वे} हा ग्रंथ काटेकोरपणा आणि
अचूकता यांवर भर देणाऱ्या ग्रीक गणिताचे विकसित उदाहरण आहे. युक्लिडच्या
\emph{मूलतत्त्वे}मधील व्याख्या आणि सिद्धता आजही माध्यमिक शाळेच्या
भूमितीच्या पाठ्यपुस्तकांत जवळजवळ तशाच टिकून आहेत (जिथे अजून माध्यमिक
शाळेत भूमिती शिकवली जाते तिथे). गणितातच नव्हे तर तत्त्वज्ञानाच्या
शाखांतही काटेकोर युक्तिवादाचा आदर्श म्हणून गणितीय युक्तिवादाचे हे रूप
मानले गेले आहे. (स्पिनोझाने नैतिक आणि धार्मिक युक्तिवादसुद्धा
युक्लिडच्या शैलीत मांडले; ते पाहणे काहीसे विचित्र वाटते!)

न्यूटन आणि लाइब्निझ यांनी सतराव्या शतकात कलनाचा शोध लावला. (प्राधान्यावरून
शतकानुशतके तीव्र वाद चालला; पण आज बहुतेक अभ्यासकांच्या मते त्यांची
बहुतांश प्रगती स्वतंत्रपणे झाली.) कलनात, उदाहरणार्थ, अनंतसूक्ष्म
राशींच्या अनंत बेरजांबद्दल युक्तिवाद करावा लागतो. या वैशिष्ट्यांमुळे
बिशप बर्कलीने टीका केली: ईश्वरावरील विश्वास आपल्या काळातील गणितापेक्षा
कमी तर्कसंगत नाही, असे त्याचे म्हणणे होते. अठराव्या शतकात कलनाच्या
पद्धती मोठ्या प्रमाणावर वापरल्या गेल्या; उदाहरणार्थ, लिओनहार्ड
ऑयलरने अनंत बेरजांचा समावेश असलेली गणिते करून प्रभावी निष्कर्ष मिळवले.

एकोणिसाव्या शतकात गणितज्ञांनी कलनाला अधिक भक्कम पाया देऊन बर्कलीच्या
टीकेला उत्तर देण्याचा प्रयत्न केला. कॉशी, वायेरश्ट्रास, बोल्झानो आणि
इतरांच्या प्रयत्नांतून ``एप्सिलॉन आणि डेल्टा'' यांच्या साहाय्याने सीमा,
सातत्य, अवकलन आणि समाकलन यांच्या आजच्या व्याख्या निर्माण झाल्या; म्हणजेच
त्यात अनंतसूक्ष्म राशींचा संदर्भ उरला नाही. पुढे त्याच शतकात गणितज्ञांनी
वास्तव संख्यांसह कलनाचे सर्व पैलू नैसर्गिक संख्यांच्या आधारे स्पष्ट
करण्याचा अधिक प्रयत्न केला. (क्रोनेकर: ``पूर्ण संख्या ईश्वराने निर्माण
केल्या; उरलेले सर्व मानवाचे कार्य आहे.'') 1872 मध्ये डेडेकिंडने
``सातत्य आणि अपरिमेय संख्या'' लिहून वास्तव संख्या परिमेय संख्यांच्या
संचांच्या रूपात कशा ``रचता'' येतात हे दाखवले. (परिमेय संख्यांकडे, तुम्हाला
माहीत आहेच, नैसर्गिक संख्यांच्या जोड्यांच्या रूपात पाहता येते.) 1888 मध्ये
त्याने ``Was sind und was sollen die Zahlen'' (साधारणतः ``नैसर्गिक
संख्या म्हणजे काय आणि त्यांचे स्वरूप कसे असावे?'') लिहून नैसर्गिक संख्या
निखळ ``तार्किक'' संज्ञांत स्पष्ट करण्याचा प्रयत्न केला. 1887 मध्ये
क्रोनेकरने ``\"Uber den Zahlbegriff'' (``संख्येच्या संकल्पनेविषयी'')
लिहून सर्व गणितीय वस्तू पूर्ण संख्यांच्या आधारे दर्शवण्याविषयी सांगितले.
1889 मध्ये ज्यूझेप्पे पेआनोने नैसर्गिक संख्यांसाठी आकारिक, प्रतीकात्मक
स्वयंसिद्धके दिली.

एकोणिसाव्या शतकाच्या अखेरीस अनंताशी व्यवहार करण्याचे नवे धाडसही दिसले.
त्याआधी अनंत वस्तू आणि रचना (उदा., नैसर्गिक संख्यांचा संच) अतिशय
सावधपणे हाताळल्या जात. ``अनंत अनेक'' म्हणजे ``तुम्हाला हव्या तितक्या''
आणि ``सीमेत जवळ जाणे'' म्हणजे ``तुम्हाला हवे तितके जवळ जाणे'' असे समजले
जाई. पण गेऑर्ग कांटोरने अनंताचा प्रत्यक्ष अर्थाने विचार करणे शक्य आहे हे
दाखवले. कांटोर, डेडेकिंड आणि इतरांच्या कार्यामुळे गणिताची आज व्यापकपणे
स्वीकारली जाणारी सर्वसाधारण संचसिद्धान्तीय समज विकसित होऊ लागली.

यातून आपण तर्कशास्त्र आणि पायाभरणीतील विसाव्या शतकाच्या घडामोडींपर्यंत
पोहोचतो. 1902 मध्ये रसेलने फ्रेगेला पत्र लिहून त्याच्या तार्किक पद्धतीतील
विरोधाभास कळवला.

1904 मध्ये झर्मेलोने तथाकथित ``निवडीचे स्वयंसिद्धक'' वापरून कांटोरचे
सुव्यवस्थापन तत्त्व सिद्ध केले; या स्वयंसिद्धकाच्या स्वीकारार्हतेवर मोठा वाद
झाला. 1910 ते 1913 या काळात रसेल आणि व्हाइटहेड यांच्या \emph{Principia
  Mathematica}चे तीन खंड प्रसिद्ध झाले. गणिताला तार्किक आधार देण्याचा
फ्रेगेचा कार्यक्रम त्यांनी पुढे नेला. मात्र रसेल आणि व्हाइटहेड यांना
निखळ तार्किक म्हणून समर्थनीय वाटणे कठीण अशी दोन तत्त्वे स्वीकारावी
लागली: अनंततेचे स्वयंसिद्धक आणि ``न्यूनीकरणीयतेचे'' स्वयंसिद्धक. 1900
च्या दशकात पॉँकारेने गणितातील ``स्वसमावेशक व्याख्यां''च्या वापरावर
टीका केली. 1910 च्या दशकात ब्रौवरने विमध्य सिद्धांत
($\varphi \lor \lnot \varphi$) न वापरता संपूर्ण गणिताची ``अंतःप्रज्ञावादी''
पायावर पुनर्मांडणी करण्याचा प्रस्ताव मांडायला सुरुवात केली.

खरोखरच विलक्षण काळ!{} सर्व गणित तर्कशास्त्रात रूपांतरित करण्याच्या
कार्यक्रमाला आता ``तर्कमीमांसावाद'' म्हणतात. वर सांगितलेल्या अडचणींमुळे
तो अयशस्वी ठरला, असे सामान्यतः मानले जाते. मानसिक अंतःप्रज्ञावादी
रचनांवर गणित उभारण्याच्या कार्यक्रमाला ``अंतःप्रज्ञावाद'' म्हणतात;
दैनंदिन गणितावर तो अति कठोर निर्बंध घालतो, असे मानले जाते. शतकाच्या
वळणावर, सर्वकाळातील अत्यंत प्रभावशाली गणितज्ञांपैकी एक असलेला डेव्हिड
हिल्बर्ट हा कांटोर आणि डेडेकिंड यांनी आणलेल्या नव्या, अमूर्त पद्धतींचा
ठाम समर्थक होता: ``कांटोरने आपल्यासाठी निर्माण केलेल्या स्वर्गातून
आपल्याला कोणीही हाकलू शकणार नाही.'' त्याच वेळी, या नव्या पद्धतींवरील
पायाभूत टीकेचीही त्याला जाणीव होती (विचित्र म्हणजे, आता या पद्धतींना
``अभिजात'' म्हणतात). दोन्ही गोष्टी साधण्याचा त्याने असा मार्ग सुचवला:
\begin{enumerate}
\item अभिजात पद्धती आकारिक स्वयंसिद्धके आणि नियम यांच्या साहाय्याने
  दर्शवा; गणितीय प्रश्न स्वयंसिद्धकीय पद्धतीतील सूत्रे म्हणून दर्शवा.
\item या आकारिक निगमन-पद्धती सुसंगत आहेत हे सिद्ध करण्यासाठी सुरक्षित,
  ``सांतवादी'' पद्धती वापरा.
\end{enumerate}

पहिले उद्दिष्ट साधण्याच्या दिशेने हिल्बर्टच्या कार्याने मोठी मजल मारली.
1899 मध्ये त्याने आपल्या प्रसिद्ध \emph{भूमितीची पायाभरणी} या पुस्तकात
भूमितीसाठी हे काम केले होते. पुढील वर्षांत, हिल्बर्टने भूमितीसाठी जे
केले ते गणिताच्या इतर शाखांत करण्यासाठी तो, त्याचे अनेक विद्यार्थी आणि
सहकारी यांनी काम केले. हिल्बर्टने स्वतः अंकगणित आणि विश्लेषण यांसाठी
स्वयंसिद्धक-प्रणाली दिल्या. झर्मेलोने संचसिद्धान्ताची स्वयंसिद्धकीय मांडणी
दिली; फ्रेंकेल, स्कोलेम, फोन न्यूमन आणि इतरांनी ती पुढे विकसित केली.
1920 च्या दशकाच्या मध्यापर्यंत ``सर्व'' गणिताची स्वयंसिद्धकीय मांडणी
म्हणवणारे दोन दृष्टिकोन होते: रसेल आणि व्हाइटहेड यांचे \emph{Principia
  mathematica}, आणि पुढे झर्मेलो--फ्रेंकेल संच उपपत्ती म्हणून ओळखली
जाणारी मांडणी.

1921 मध्ये हिल्बर्टने या पद्धती सुसंगत असल्याचे सिद्ध करण्याच्या उद्देशाने
संशोधन कार्यक्रम सुरू केला. यात त्याला अनेक विद्यार्थ्यांची, विशेषतः
बेर्नाइस, आकरमान आणि पुढे गेन्ट्सेन यांची मदत झाली. या उद्दिष्टामागील
मूळ कल्पना अशी होती: गणितात विसंगतीची निष्पत्ती शक्य आहे का, हा
प्रश्न चिन्हांच्या संभाव्य क्रमांबद्दलच्या संयुक्तिक प्रश्नात रूपांतरित
करायचा. म्हणजे, अंकगणित, विश्लेषण किंवा संचसिद्धान्त यांच्या
स्वयंसिद्धक-प्रणालीतील स्वयंसिद्धकांपासून, उदाहरणार्थ, $\varphi \land
\lnot \varphi$ ची योग्य निष्पत्ती ठरणाऱ्या वाक्यांचे संभाव्य क्रम
तपासायचे. असा चिन्हक्रम अशक्य असल्याची सिद्धता ही स्वतः गणितीय सिद्धता
असल्यामुळे या स्वयंसिद्धकीय पद्धतींत आकारिकरीत्या मांडता आली असती.
दुसऱ्या शब्दांत, उदाहरणार्थ, अंकगणित सुसंगत आहे असे सांगणारे
$\OCon$ हे वाक्य असले असते. शिवाय, विशेषतः त्याची सिद्धता केवळ अतिशय
मर्यादित, ``सांतवादी'' साधनांनी होत असल्यास, ते वाक्य संबंधित पद्धतींत
सिद्ध करता आले पाहिजे होते.

विकसित स्वयंसिद्धक-प्रणाली प्रत्येक गणितीय प्रश्न सोडवतील, हे दुसरे
उद्दिष्ट दोन प्रकारे नेमके करता येते. पहिला प्रकार असा: गणिताच्या
स्वयंसिद्धक-प्रणालीच्या भाषेतील कोणतेही वाक्य~$\varphi$ घेतल्यास,
स्वयंसिद्धकांपासून $\varphi$ किंवा $\lnot \varphi$ सिद्ध करता येते. असे असल्यास,
स्वयंसिद्धकांच्या आधारे ज्यांची सिद्धता किंवा खंडन दोन्ही होत नाहीत अशी
वाक्ये, म्हणजे त्यांनी न सोडवलेले प्रश्न, उरणार नाहीत. हा गुणधर्म
असलेल्या स्वयंसिद्धक-प्रणालीला \emph{संपूर्ण} म्हणतात. अर्थात, दिलेल्या
एखाद्या वाक्यासाठी या दोन पर्यायांपैकी कोणता खरा आहे हे ठरवणे तरीही
कठीण असू शकते. पण तत्त्वतः ते ठरवण्याची पद्धत असली पाहिजे. प्रत्यक्षात,
हिल्बर्टने विचारात घेतलेल्या स्वयंसिद्धक आणि निष्पत्ती-पद्धतींसाठी
संपूर्णतेवरून अशी पद्धत अस्तित्वात आहे असे निष्पन्न होईल---जरी हिल्बर्टला
याची जाणीव नव्हती. दुसरा अर्थ अधिक प्रबळ मागणी करतो: दिलेले
वाक्य~$\varphi$ स्वयंसिद्धकांपासून निष्पन्न करता येते की नाही हे ठरवणारी
यांत्रिक, संगणकीय पद्धत असावी.

1931 मध्ये गोडेलने दोन ``अपूर्णतेची प्रमेये'' सिद्ध केली. त्यांवरून
हा कार्यक्रम मूळ स्वरूपात यशस्वी होऊ शकत नाही, हे दिसले. अंकगणितासाठी
पुरेशी प्रबळ, सुसंगत आणि प्रभावी स्वयंसिद्धकीय मांडणी असलेली प्रणाली
संपूर्ण असू शकत नाही. शिवाय, दुसऱ्या अपूर्णता प्रमेयाच्या अटी पूर्ण करणारी
अशी प्रणाली तिच्या प्रमाण सिद्धतायोग्यता-विधेयाने व्यक्त केलेली स्वतःची
सुसंगतता सिद्ध करू शकत नाही. त्या सुसंगतता-विधानाचे खंडनही होत नाही असे
म्हणण्यासाठी अधिक गृहीतक हवे; उदाहरणार्थ, प्रणालीने सिद्ध केलेली सर्व
अंकगणितीय वाक्ये नैसर्गिक संख्यांच्या प्रमाण प्रतिमानात सत्य असावीत.


यामुळे हिल्बर्टच्या मूळ कार्यक्रमाला घातक धक्का बसला. तरीही गणितात
अनेकदा घडते तसे, यातून संशोधनाच्या नव्या दिशाही उघडल्या. सर्व गणिताला
व्यापणारी एकच आकारिक पद्धत नसल्यास, अधिक मर्यादित पद्धती विकसित करून
त्यांत काय सिद्ध करता येते याचा अभ्यास करणे अर्थपूर्ण ठरते. या पद्धतींची
सुसंगतता सिद्ध करण्यासाठी कमी निर्बंध असलेल्या सिद्धता-पद्धती विकसित
करणे, आणि त्यांची सुसंगतता सिद्ध करणे किती कठीण आहे हे मोजण्याचे मार्ग
शोधणेही अर्थपूर्ण ठरते. अपूर्णतेच्या प्रमेयांची गृहीतके पूर्ण करणाऱ्या
आकारिक पद्धतींत असे प्रश्न असतात की जे त्या सोडवू शकत नाहीत. म्हणून
दिलेल्या आकारिक पद्धतीला न सोडवता येणारे ``महत्त्वाचे'' प्रश्न शोधणे,
आणि ते सोडवण्यासाठी पद्धत किती प्रबळ असावी लागते हे ठरवणे योग्य ठरते.
आजतागायत तर्कशास्त्रज्ञ या प्रश्नांचा सिद्धता-उपपत्ती या नव्या
गणितीय शाखेत अभ्यास करत आहेत.












\section{व्याख्या}\label{inc:int:def:sec}

गणिताला आकारिक रूप देण्याचा आणि असे रूप सुसंगत व संपूर्ण आहे हे
दाखवण्याचा हिल्बर्टचा कार्यक्रम राबवण्यासाठी प्रथम भाषा, तार्किक चौकट
आणि स्वयंसिद्धक-प्रणाली निवडावी लागेल. आपल्या उद्देशासाठी गणित प्रथम-क्रम
भाषेत आकारिकरीत्या मांडता येते असे समजू. म्हणजे स्थिरांक,
फलने आणि विधेये यांचा असा काही संच आहे की प्रथम-क्रम
तर्कशास्त्रातील संयोजके व संख्यापक यांच्यासह त्यातून गणितातील विधाने
व्यक्त करता येतात. अशी भाषा अस्तित्वात आहे याबद्दल बहुतेकांचे एकमत
आहे: ती संचसिद्धान्ताची भाषा आहे आणि त्यात $\in$ हे एकमेव अतार्किक
चिन्ह आहे. इतकी साधी भाषा इतकी अभिव्यक्तिक्षम असू शकते हा दावा
पहिल्या नजरेला नक्कीच अविश्वसनीय वाटतो; गणितातील जवळजवळ सर्व काही
या अत्यल्प शब्दसंग्रहात व्यक्त करता येते हे प्रस्थापित करण्यासाठी
मोठे काम करावे लागले. सध्या चर्चा सोपी ठेवण्यासाठी ती केवळ
अंकगणितापुरती, म्हणजे नैसर्गिक संख्या~$\Nat$ यांच्याशी संबंधित
गणिताच्या भागापुरती, मर्यादित करू. अंकगणितातील तथ्ये व्यक्त करण्यासाठी
योग्य भाषा~$\Lang L_A$ आहे. $\Lang L_A$ मध्ये एक द्विस्थानी
विधेय~$<$, एक स्थिरांक~$\Obj 0$, एक एकस्थानी
फलन~$\prime$, आणि $+$ व $\times$ ही दोन द्विस्थानी
फलने आहेत.

\begin{defn}
वाक्यांचा संच~$\Gamma$ हा तार्किक निष्पन्नतेखाली बंद असेल, म्हणजे
$\Gamma = \Setabs{\varphi}{\Gamma \Entails
\varphi}$ असेल, तर त्याला \emph{उपपत्ती} म्हणतात.
\end{defn}

उपपत्ती निर्दिष्ट करण्याचे दोन सोपे मार्ग आहेत. एक म्हणजे एखाद्या
संरचनांमध्ये सत्य असलेल्या वाक्यांचा संच म्हणून
तिला निर्दिष्ट करणे. उदाहरणार्थ, $\Lang L_A$ साठी अशी संरचना
घ्या की तिचे प्रांत~$\Nat$ आहे आणि सर्व अतार्किक चिन्हांचा
अपेक्षेप्रमाणे अर्थ लावला आहे.

\begin{defn}\label{inc:int:def:def:standard-model}
\emph{अंकगणिताचे प्रमाण प्रतिमान} ही पुढीलप्रमाणे परिभाषित
केलेली संरचना~$\Struct
N$ आहे:
\begin{enumerate}
\item $\Domain N = \Nat$
\item $\Assign{\Obj 0}{N} = 0$
\item $\Assign{\Obj \prime}{N}(n) = n + 1$ सर्व $n \in \Nat$ साठी
\item $\Assign{\Obj +}{N}(n, m) = n + m$ सर्व $n, m \in \Nat$ साठी
\item $\Assign{\Obj \times}{N}(n, m) = n\cdot m$ सर्व $n, m \in \Nat$ साठी
\item $\Assign{\Obj <}{N} = \Setabs{\tuple{n, m}}{n \in \Nat, m \in
  \Nat, n < m}$
\end{enumerate}
\end{defn}

`$\times$' आणि `$\cdot$' यांतील फरक लक्षात घ्या: $\times$ हे
अंकगणिताच्या भाषेतील चिन्ह आहे. गुणाकाराची आठवण व्हावी म्हणून ते
निवडले आहे; पण $\times$ ही गुणाकाराची क्रिया नसून द्विस्थानी फलनचिन्ह
आहे (औपचारिकरीत्या, $\Obj f^2_1$). याउलट, $\cdot$ हेच नेहमीचे
गुणाकाराचे फलन \emph{आहे}. $n \cdot m$ सारखे काही दिसल्यास त्याचा
अर्थ $n$ आणि~$m$ या संख्यांचा गुणाकार होतो; $x \times y$ सारखे
काही दिसल्यास आपण अंकगणिताच्या भाषेतील पदाबद्दल बोलत असतो. प्रमाण
प्रतिमानात फलनचिन्ह~$\times$ चा अर्थ नैसर्गिक संख्यांवरील फलन
$\cdot$ असा निश्चित केला जातो. बेरीजेसाठी मात्र अंकगणिताच्या
भाषेतील फलनचिन्ह आणि नैसर्गिक संख्यांवरील बेरीज हे दोन्ही दर्शवण्यासाठी
आपण $+$ वापरतो. त्याचा अर्थ संदर्भावरून ठरवावा लागतो.

\begin{defn}
\emph{सत्य अंकगणिताची} उपपत्ती म्हणजे अंकगणिताच्या प्रमाण प्रतिमानात
पूर्त होणाऱ्या वाक्यांचा संच; म्हणजे,
\[\Th{TA} = \Setabs{\varphi}{\Sat{N}{\varphi}}.
\]
\end{defn}

$\Th{TA}$ ही उपपत्ती आहे. कारण $\Th{TA} \Entails \varphi$ असेल तेव्हा
$\Th{TA}$ ची पूर्ती करणाऱ्या प्रत्येक संरचनांमध्ये $\varphi$ ची
पूर्ती होते. $\Sat{N}{\Th{TA}}$ असल्यामुळे $\Sat{N}{\varphi}$ मिळते;
म्हणून $\varphi \in \Th{TA}$.

उपपत्ती~$\Gamma$ निर्दिष्ट करण्याचा दुसरा मार्ग म्हणजे काही
वाक्यसंच~$\Gamma_0$ पासून तार्किकरीत्या निष्पन्न होणाऱ्या
वाक्यांचा संच म्हणून तिला सांगणे. तेव्हा $\Gamma$ हे
तार्किक निष्पन्नतेखाली $\Gamma_0$ चे ``संवरण'' असते. या मार्गाने उपपत्ती निर्दिष्ट करणे तेव्हाच
उपयुक्त ठरते जेव्हा $\Gamma_0$ स्पष्टपणे निर्दिष्ट केला असेल;
उदाहरणार्थ, $\Gamma_0$ मधील घटक ची यादी दिली असेल.
किमान एवढे तरी हवे की $\Gamma_0$ निर्णेय असावा: एखादे
वाक्य $\Gamma_0$ चे घटक आहे की नाही हे ठरवण्यासाठी
संगणनक्षम चाचणी असावी. $\Gamma_0$ मधील वाक्यांना
$\Gamma$ ची \emph{स्वयंसिद्धके}, आणि $\Gamma$ ला
$\Gamma_0$ या संचातील \emph{स्वयंसिद्धकांद्वारे निरूपित} म्हणतो.

\begin{defn}
उपपत्ती~$\Gamma$ ही $\Gamma_0$ या संचातील
\emph{स्वयंसिद्धकांद्वारे निरूपित} आहे तेव्हा आणि केवळ तेव्हाच
\[\Gamma = \Setabs{\varphi}{\Gamma_0 \Entails \varphi}
\]
असते.
\end{defn}

\begin{defn}
पुढील वाक्यांनी स्वयंसिद्धकांद्वारे निरूपित केलेली उपपत्ती
$\Th{Q}$ ``रॉबिन्सनची $\Th{Q}$'' म्हणून ओळखली जाते आणि ती
अंकगणिताची अतिशय साधी उपपत्ती आहे.
\begin{align*}
& \lforall[x][\lforall[y][(\eq[x'][y'] \lif \eq[x][y])]] \tag{$\mathsf{Q}_1$}\\
& \lforall[x][\eq/[\Obj 0][x']] \tag{$\mathsf{Q}_2$}\\
& \lforall[x][(\eq[x][\Obj 0] \lor \lexists[y][\eq[x][y']])] \tag{$\mathsf{Q}_3$}\\
& \lforall[x][\eq[(x + \Obj 0)][x]] \tag{$\mathsf{Q}_4$}\\
& \lforall[x][\lforall[y][\eq[(x + y')][(x + y)']]] \tag{$\mathsf{Q}_5$}\\
& \lforall[x][\eq[(x \times \Obj 0)][\Obj 0]] \tag{$\mathsf{Q}_6$}\\
& \lforall[x][\lforall[y][\eq[(x \times y')][((x \times y) + x)]]] \tag{$\mathsf{Q}_7$}\\
& \lforall[x][\lforall[y][(x < y \liff \lexists[z][\eq[(z' + x)][y]])]] \tag{$\mathsf{Q}_8$}
\end{align*}
$\{\mathsf{Q}_1, \dots, \mathsf{Q}_8\}$ हा वाक्यांचा संच $\Th{Q}$ ची
स्वयंसिद्धके आहेत. म्हणून त्यांच्यापासून तार्किकरीत्या निष्पन्न होणारी सर्व
वाक्ये $\Th{Q}$ मध्ये आहेत:
\[\Th{Q} = \Setabs{\varphi}{\{\mathsf{Q}_1, \dots, \mathsf{Q}_8\} \Entails \varphi}.
\]
\end{defn}

\begin{defn}
$\varphi(x)$ हे $\Lang L_A$ मधील असे सूत्र आहे असे समजा की
त्यात $x$ आणि $y_1$, \dots, $y_n$ ही मुक्त चरे आहेत. तर पुढील
रूपातील कोणतेही वाक्य हे \emph{विगमन-रूपबंधाचे} उदाहरण आहे:
\[\lforall[y_1][\dots\lforall[y_n][((\varphi(\Obj 0) \land \lforall[x][(\varphi(x)
\lif \varphi(x'))]) \lif \lforall[x][\varphi(x)])]]
\]

\emph{पेआनो अंकगणित}~$\Th{PA}$ ही $\Th{Q}$ च्या स्वयंसिद्धकांसह
विगमन-रूपबंधाच्या सर्व उदाहरणांनी स्वयंसिद्धकांद्वारे निरूपित
केलेली उपपत्ती आहे.
\end{defn}

\begin{explain}
विगमन-रूपबंधाचे प्रत्येक उदाहरण~$\Struct{N}$ मध्ये सत्य असते.
सूत्र~$\varphi$ मध्ये फक्त एकच मुक्त चल~$x$ असल्यास
हे सर्वात सहज दिसते. तेव्हा $\varphi(x)$ हे~$\Struct{N}$ मध्ये
$\Nat$ चा उपसंच~$X_{\varphi}$ परिभाषित करते. $s(x) = n$ असताना
$\Sat{N}{\varphi(x)}[s]$ होईल अशा सर्व~$n \in \Nat$ चा~$X_{\varphi}$ हा संच आहे.
विगमन-रूपबंधाचे संबंधित उदाहरण असे आहे:
\[((\varphi(\Obj 0) \land \lforall[x][(\varphi(x) \lif \varphi(x'))]) \lif
  \lforall[x][\varphi(x)]).
\]
याचा पूर्वपक्ष~$\Struct{N}$ मध्ये सत्य असल्यास $0 \in X_{\varphi}$;
आणि $n \in X_{\varphi}$ असेल तेव्हा $n+1$ देखील त्यात असते.
$0 \in X_{\varphi}$ असल्याने $1 \in X_{\varphi}$ मिळते. पुन्हा
$1 \in X_{\varphi}$ असल्याने $2 \in X_{\varphi}$ मिळते, आणि पुढे असेच.
म्हणून प्रत्येक
$n \in \Nat$ साठी $n \in X_{\varphi}$. याचा अर्थ
$\lforall[x][\varphi(x)]$ ची~$\Struct{N}$ मध्ये पूर्ती होते.
\end{explain}

$\Th{Q}$ आणि $\Th{PA}$ या दोन्ही स्वयंसिद्धकांद्वारे निरूपित
उपपत्ती आहेत. त्या किती प्रबळ आहेत, हा मुख्य प्रश्न आहे.
उदाहरणार्थ, $\Lang L_A$ मध्ये व्यक्त करता येणारी~$\Nat$ विषयीची
सर्व सत्ये $\Th{PA}$ सिद्ध करू शकते का? नेमके सांगायचे तर,
$\Lang L_A$ मध्ये मांडता येणारे सर्व प्रश्न $\Th{PA}$ ची
स्वयंसिद्धके सोडवतात का?

हा प्रश्न दुसऱ्या रीतीनेही विचारता येईल: $\Th{PA} = \Th{TA}$
आहे का? $\Th{TA}$ अर्थातच~$\Nat$ विषयीची सर्व सत्ये सिद्ध
करते (म्हणजे ती त्यात समाविष्ट असतात), आणि $\Lang L_A$ मधील
सर्व प्रश्न सोडवते. कारण $\varphi$ हे $\Lang L_A$ मधील वाक्य
असेल तर $\Sat{N}{\varphi}$ किंवा $\Sat{N}{\lnot \varphi}$ यांपैकी
एक सत्य असते; म्हणून $\Th{TA}
\Entails \varphi$ किंवा $\Th{TA} \Entails \lnot \varphi$. अशा उपपत्तीला
\emph{संपूर्ण} म्हणतात.

\begin{defn}
उपपत्ती $\Gamma$ ही \emph{संपूर्ण} आहे तेव्हा आणि केवळ तेव्हाच तिच्या
भाषेतील प्रत्येक वाक्य~$\varphi$ साठी $\Gamma \Entails \varphi$
किंवा $\Gamma \Entails \lnot \varphi$.
\end{defn}

\begin{explain}
संपूर्णता प्रमेयामुळे $\Gamma \Entails \varphi$ तेव्हा आणि केवळ तेव्हाच
$\Gamma \Proves
\varphi$. म्हणून $\Gamma$ संपूर्ण आहे तेव्हा आणि केवळ तेव्हाच तिच्या भाषेतील
प्रत्येक वाक्य~$\varphi$ साठी $\Gamma \Proves \varphi$ किंवा
$\Gamma \Proves \lnot \varphi$.
\end{explain}

आणखी एक प्रश्न असा आहे: एखादे वाक्य~$\Th{TA}$ मध्ये,
$\Th{PA}$ मध्ये, किंवा अगदी~$\Th{Q}$ मध्येही आहे की नाही हे
तपासण्यासाठी संगणकीय प्रक्रिया वापरता येईल का? संच (उदाहरणार्थ,
वाक्यांचा संच) निर्णेय कधी असतो हे परिभाषित करून हा
प्रश्न अधिक नेमका करता येईल.

\begin{defn}
संच $X$ हा \emph{निर्णेय} आहे तेव्हा आणि केवळ तेव्हाच अशी संगणकीय प्रक्रिया
असते की आदान~$x$ दिल्यावर $x \in X$ असल्यास ती~$1$ आणि
अन्यथा~$0$ परत करते.
\end{defn}

म्हणून प्रश्न असा होतो: $\Th{TA}$ ($\Th{PA}$, $\Th{Q}$) निर्णेय आहे का?

या सर्व प्रश्नांचे उत्तर नकारार्थी असेल. यांपैकी कोणतीही उपपत्ती
निर्णेय नाही. मात्र ही घटना केवळ या उपपत्तींनाच लागू नाही. प्रत्यक्षात,
विशिष्ट अटी पूर्ण करणाऱ्या \emph{कोणत्याही} उपपत्तीसाठी हेच निष्कर्ष
लागू होतात. यांपैकी एक अट, जी $\Th{Q}$ आणि $\Th{PA}$ पूर्ण करतात,
अशी आहे की उपपत्ती निर्णेय स्वयंसिद्धकसंचाने निरूपित झालेली असावी.

\begin{defn}
एखादी उपपत्ती निर्णेय स्वयंसिद्धकसंचाने निरूपित झाली असेल तर ती
\emph{निर्णेय स्वयंसिद्धकसंचाने निरूपणयोग्य} आहे.
\end{defn}

\begin{ex}
सांत वाक्यांच्या संचाने स्वयंसिद्धकांद्वारे निरूपित
झालेली कोणतीही उपपत्ती निर्णेय स्वयंसिद्धकसंचाने निरूपणयोग्य असते, कारण प्रत्येक
सांत संच निर्णेय असतो. म्हणून, उदाहरणार्थ, $\Th{Q}$ ही
निर्णेय स्वयंसिद्धकसंचाने निरूपणयोग्य आहे.

$\Th{PA}$ सारखी रूपबंधांद्वारे स्वयंसिद्धके दिलेली उपपत्तीदेखील
निर्णेय स्वयंसिद्धकसंचाने निरूपणयोग्य असते. कारण $\psi$ हे $\Th{PA}$ च्या स्वयंसिद्धकांपैकी
आहे का, म्हणजे $\psi$ हे $\Th{PA}$ चे स्वयंसिद्धक असल्यास
$\Char{X}(\psi) = 1$ आणि
अन्यथा $= 0$ असणारे फलन $\Char{X}$ कसे काढायचे, हे पुढीलप्रमाणे
तपासता येते. प्रथम, $\psi$ हे $\Th{Q}$ च्या स्वयंसिद्धकांपैकी आहे
का ते तपासा. असल्यास उत्तर ``होय'' आणि $\Char{X}(\psi) = 1$.
नसल्यास ते विगमन-रूपबंधाचे उदाहरण आहे का हे तपासा. हे पद्धतशीरपणे
करता येते. येथे पुढील पद्धत कदाचित सर्वात स्पष्ट आहे: विगमन-रूपबंधाचे
प्रत्येक उदाहरण सुरुवातीला काही वैश्विक संख्यापकांनी, त्यानंतर
सोपाधिक असलेल्या एका उप-सूत्र ने सुरू होते. त्या सोपाधिकाचा
उत्तरपक्ष $\lforall[x][\varphi(x, y_1, \dots, y_n)]$ असतो. येथे $x$
आणि $y_1$, \dots, $y_n$ ही~$\varphi$ ची सर्व मुक्त चरे असतात, आणि
$\psi$ च्या सुरुवातीचे संख्यापक $y_1$, \dots,~$y_n$ ही चरे बद्ध
करतात. हे~$\varphi$ काढून घेतल्यावर त्याची मुक्त चरे सुरुवातीच्या
वैश्विक संख्यापकांनी आणि~$\lforall[x]$ ने बद्ध केलेल्या चरांशी
जुळतात का ते तपासा. मग सोपाधिकाचा पूर्वपक्ष पुढील स्वरूपाशी जुळतो
का ते तपासा:
\[\varphi(\Obj 0, y_1, \dots, y_n) \land \lforall[x][(\varphi(x, y_1, \dots, y_n)
\lif \varphi(x', y_1, \dots, y_n))]
\]
तो जुळल्यास $\psi$ विगमन-रूपबंधाचे उदाहरण आहे; न जुळल्यास $\psi$ तसे नाही.
\end{ex}

या प्रश्नाचे---आणि कोणत्या उपपत्ती संपूर्ण किंवा निर्णेय असतात या
अधिक व्यापक प्रश्नाचे---उत्तर शोधताना पुढील व्याख्याही उपयुक्त ठरेल.
संच $X$ हा रिक्त असेल किंवा $f \colon \Nat
\to X$ असे आच्छादक फलन असेल तेव्हा आणि केवळ तेव्हाच तो गणनीय
असतो, हे आठवा. अशा फलनाला $X$ चे प्रगणन म्हणतात.

\begin{defn}
संच $X$ हा रिक्त असेल किंवा त्याचे संगणनक्षम प्रगणन असेल तेव्हा आणि केवळ तेव्हाच
त्याला \emph{संगणनक्षमपणे प्रगणनीय} (संक्षेपात संगणनक्षमपणे प्रगणनीय) म्हणतात.
\end{defn}

निर्णेय स्वयंसिद्धकसंचाने निरूपणयोग्यता व्यतिरिक्त, अपूर्णतेची प्रमेये लागू होण्यासाठी
उपपत्तीविषयी आणखी एक अट अशी असेल की ती संगणनक्षम फलने आणि
निर्णेय संबंध यांच्याबद्दलची मूलभूत तथ्ये सिद्ध करू शकेल इतकी
प्रबळ असावी. ``मूलभूत तथ्ये'' म्हणजे संगणनक्षम फलनांच्या प्रत्येक
आदानासाठी त्यांची मूल्ये काय आहेत हे व्यक्त करणारी वाक्ये.
आणि ``पुरेशी प्रबळ'' म्हणजे संबंधित उपपत्ती ही वाक्ये आपली प्रमेये
मानतात. उदाहरण म्हणून बेरीज हे सर्वसाधारण संगणनक्षम फलन घ्या.
$2$ आणि $3$ या आदानांसाठी $+$ चे मूल्य~$5$ आहे, म्हणजे
$2+3 = 5$. $2$ आणि $3$ या आदानांसाठी $+$ चे मूल्य~$5$ आहे
असे सांगणारे अंकगणिताच्या भाषेतील वाक्य आहे:
$(\num{2} + \num{3}) = \num{5}$. उदाहरणार्थ, $\Th{Q}$ हे वाक्य
सिद्ध करते. अधिक सर्वसाधारणपणे, प्रत्येक संगणनक्षम फलन
$f(x_1, x_2)$ साठी~$\Lang{L_A}$ मध्ये असे सूत्र~$\varphi_f(x_1,
x_2, y)$ हवे की $f(n_1, n_2) = m$ असेल तेव्हा $\Th{Q}
\Proves \varphi_f(\num{n_1}, \num{n_2}, \num{m})$. अशा रीतीने
$\Th{Q}$ हे $n_1$, $n_2$ या आदानांसाठी~$f$ चे मूल्य~$m$
आहे हे सिद्ध करते. प्रत्यक्षात, आपण याहून थोडे अधिक मागतो:
या $n_1$,~$n_2$ या आदानांसाठी~$f$ चे मूल्य दुसरी कोणतीही संख्या नाही हेदेखील
तिने सिद्ध करावे. निर्णेय संबंधांबाबतही तसेच आहे. पुढील
दोन व्याख्या हे नेमके करतात.

\begin{defn}
सूत्र~$\varphi(x_1, \dots, x_k, y)$ हे
$f\colon \Nat^k \to \Nat$ या फलनाचे~$\Gamma$ मध्ये
\emph{निरूपण} करते तेव्हा आणि केवळ तेव्हाच की $f(n_1,
\dots, n_k) = m$ असेल तेव्हा पुढील दोन्ही अटी पूर्ण होतात:
\begin{enumerate}
\item $\Gamma \Proves \varphi(\num{n_1}, \dots, \num{n_k}, \num{m})$, आणि
\item $\Gamma \Proves \lforall[y](\varphi(\num{n_1}, \dots, \num{n_k},
y) \lif y = \num{m})$.
\end{enumerate}
\end{defn}

\begin{defn}
सूत्र~$\varphi(x_1, \dots, x_k)$ हे
$R \subseteq \Nat^k$ या संबंधाचे \emph{निरूपण} करते तेव्हा आणि केवळ तेव्हाच की
\begin{enumerate}
\item $R(n_1, \dots, n_k)$ असेल तेव्हा $\Gamma \Proves \varphi(\num{n_1},
\dots, \num{n_k})$, आणि
\item $R(n_1, \dots, n_k)$ नसेल तेव्हा $\Gamma \Proves \lnot
\varphi(\num{n_1}, \dots, \num{n_k})$.
\end{enumerate}
\end{defn}

पहिले अपूर्णता प्रमेय लागू होण्याची सामर्थ्याची अट अशी आहे: उपपत्ती
सर्व संगणनक्षम फलनांचे आणि सर्व निर्णेय संबंधांचे निरूपण
करत असावी. $\Th{Q}$ आणि तिच्या विस्तारित उपपत्ती ही अट पूर्ण करतात.
पण हे सिद्ध करण्यास आपल्याला थोडा वेळ लागेल: $\Th{Q}$ काय
सिद्ध करू शकते याविषयीचे हे सोपे नसलेले तथ्य आहे, कारण $\Th{Q}$ कडे
थोडीच स्वयंसिद्धके आहेत आणि त्यांच्यापासून ही सर्व तथ्ये सिद्ध
करावी लागतील. तरी $\Th{Q}$ ही अतिशय दुर्बल उपपत्ती आहे.
म्हणून $\Th{Q}$ सर्व संगणनक्षम फलने निरूपित करते हे सिद्ध
करणे कठीण असले, तरी बहुतेक महत्त्वाच्या उपपत्ती~$\Th{Q}$ पेक्षा
प्रबळ असतात, म्हणजे $\Th{Q}$ पेक्षा अधिक गोष्टी सिद्ध करतात.
$\Th{Q}$ ने सिद्ध केलेली गोष्ट कोणतीही अधिक प्रबळ उपपत्तीही
सिद्ध करते; म्हणून $\Th{Q}$ सर्व संगणनक्षम फलने निरूपित करत
असल्यास प्रत्येक अधिक प्रबळ उपपत्तीदेखील तसे करते. त्यामुळे
पहिल्या अपूर्णता प्रमेयाची ही अट अनेक महत्त्वाच्या उपपत्ती पूर्ण करतात.
आपल्या कठीण परिश्रमांचे फळ म्हणजे ही प्रमेये उपपत्तींच्या विस्तृत
वर्गाला लागू होतात हे दिसेल. नक्कीच, ``सर्व गणित'' आकारिकरीत्या
मांडू पाहणाऱ्या कोणत्याही उपपत्तीने $\Th{Q}$ सिद्ध करते ते
सर्व सिद्ध केले पाहिजे; किमान प्राथमिक संगणनांचे निष्कर्ष तरी
तिला सामावून घेता आले पाहिजेत. म्हणून ``सर्व गणिताची'' उपपत्ती
म्हणवण्यास पात्र असलेली उपपत्ती सुसंगत व प्रभावी स्वयंसिद्धकीय
मांडणी असलेली असेल तर तिला पहिले अपूर्णता प्रमेय लागू होईल.
दुसऱ्या अपूर्णता प्रमेयासाठी पुढील आढाव्यात सांगितलेल्या अधिक
कडक अटीही हव्यात.














\section{अपूर्णतेच्या निष्कर्षांचा आढावा}\label{inc:int:ovr:sec}

गणिताला निर्णेय स्वयंसिद्धकसंचाने निरूपणयोग्य उपपत्तीत
आकारिक रूप देता येईल आणि ती संपूर्ण व निर्णेय आहे हे
सिद्ध करता येईल, अशी हिल्बर्टची अपेक्षा होती. शिवाय, अतिशय दुर्बल,
``सांतवादी'' साधनांनी त्या उपपत्तीची सुसंगतता सिद्ध करून
अंतःप्रज्ञावादाच्या आक्षेपांपासून अभिजात गणिताचे समर्थन करण्याचे
त्याचे ध्येय होते. गोडेलच्या अपूर्णता प्रमेयांनी ही उद्दिष्टे
साध्य होऊ शकत नाहीत हे दाखवले.

गोडेलच्या पहिल्या अपूर्णता प्रमेयाने रसेल आणि व्हाइटहेड यांच्या
\emph{Principia Mathematica} ची एक आवृत्ती संपूर्ण नाही हे
दाखवले. पण सिद्धता प्रत्यक्षात अतिशय व्यापक होती आणि अनेक प्रकारच्या
उपपत्तींना लागू होते. म्हणजे केवळ \emph{Principia Mathematica}
गणित पूर्णपणे पकडू शकला नाही असे नव्हे, तर कोणतीही \emph{स्वीकारार्ह}
उपपत्ती तसे करू शकत नाही. अपूर्णता प्रमेये लागू होण्यासाठी उपपत्तीत
कोणते गुणधर्म पुरेसे आहेत हे वेगळे ओळखण्यास आणि गोडेलची सिद्धता
केवळ त्यांवर अवलंबून राहील अशी व्यापक करण्यास काही काळ लागला.
आता मात्र अंकगणिताच्या $\Lang L_A$ भाषेतील उपपत्त्यांसाठी पहिले
अपूर्णता प्रमेय अतिशय व्यापक रूपात मांडता येते.

\begin{thm}
$\Gamma$ ही~$\Lang
L_A$ मधील सुसंगत आणि निर्णेय स्वयंसिद्धकसंचाने निरूपणयोग्य उपपत्ती असेल, आणि ती सर्व
संगणनक्षम फलनांचे व निर्णेय संबंधांचे निरूपण करत असेल,
तर $\Gamma$ संपूर्ण नाही.
\end{thm}

$\Gamma$ संपूर्ण नाही याचा अर्थ किमान एक असे
वाक्य~$\varphi$ आहे की $\Gamma \Proves/ \varphi$ आणि
$\Gamma \Proves/ \lnot
\varphi$. अशा वाक्याला ($\Gamma)$ पासून \emph{स्वतंत्र}
म्हणतात. अशी स्वतंत्र वाक्ये असलीच पाहिजेत हे प्रत्यक्षात तुलनेने
लवकर सिद्ध करता येते. परंतु गोडेलच्या सिद्धतेचे सामर्थ्य असे की
ती स्वतःच्या असिद्धतायोग्यतेचा दावा करणाऱ्या वाक्याचे
\emph{विशिष्ट उदाहरण} देते.
या रंजक रचनेत~$\Gamma$ साठीचे \emph{गोडेल-वाक्य} म्हणून ओळखले
जाणारे वाक्य~$\mathsf{G}_\Gamma$ तयार होते. ते सिद्ध करता येत
नाही, कारण $\Gamma$ मध्ये $\mathsf{G}_\Gamma$ हेच वाक्य~$\Gamma$ मध्ये
$\mathsf{G}_\Gamma$ सिद्ध करता येत नाही या दाव्याशी सममूल्य असते.
ही रचना \emph{रचनात्मकपणे} केली जाते: $\Gamma$ ची स्वयंसिद्धकीय
मांडणी आणि निष्पत्ती-पद्धतीचे वर्णन दिल्यास सिद्धतेतून
$\mathsf{G}_\Gamma$ प्रत्यक्ष लिहिण्याची पद्धत मिळते. केवळ सुसंगततेवरून हे मूळ गोडेल-वाक्य
असिद्धतायोग्य असते; त्याचे खंडनही अशक्य ठरवण्यासाठी अधिक प्रबळ
गृहीतक, उदाहरणार्थ ओमेगा-सुसंगतता, आवश्यक असते. मात्र रॉसरच्या
बदललेल्या रचनेत केवळ सुसंगततेवरून स्वतंत्र वाक्य मिळते; वरील
प्रमेय त्या अधिक व्यापक रूपात दिले आहे.


गोडेलच्या सिद्धतेतील रचनेसाठी $\Lang L_A$ मधील पदे व
सूत्रे यांचे गुणधर्म आणि त्यांवरील क्रिया त्याच $\Lang L_A$
मध्ये व्यक्त करण्याचा मार्ग हवा. यात ``$\varphi$ हे वाक्य आहे'',
``$\delta$ ही~$\varphi$ ची निष्पत्ती आहे'' यांसारखे गुणधर्म
आणि $\Subst{\varphi}{t}{x}$ सारख्या क्रिया येतात. या मार्गाने
(a) चिन्हे आणि त्यांचे अनुक्रम (म्हणजे पदे, सूत्रे,
निष्पत्त्या इत्यादी) नैसर्गिक संख्यांत सांकेतिक करून
त्या गुणधर्मांचे व संबंधांचे निरूपण करावे लागते, कारण $\Lang L_A$
नैसर्गिक संख्यांविषयी बोलू शकते. तसेच (b) $\Gamma$ ला संबंधित
तथ्ये सिद्ध करता येतील अशा रीतीने हे करावे लागते. म्हणून हे
गुणधर्म नैसर्गिक संख्यांच्या निर्णेय गुणधर्मांनी आणि क्रिया
त्यांवरील संगणनक्षम फलनांनी सांकेतिक होतात हे दाखवावे लागते.
याला ``विन्यासमीमांसेचे अंकगणितीकरण'' म्हणतात.

विन्यासमीमांसेचे अंकगणितीकरण कसे करावे हे पाहण्यापूर्वी, $\Gamma$
``पुरेशी प्रबळ'' आहे या अटीचा विचार करू; म्हणजे ती सर्व संगणनक्षम
फलनांचे आणि निर्णेय संबंधांचे निरूपण करते. यासाठी
``संगणनक्षम'' याची नेमकी व्याख्या द्यावी लागेल. ती अनेक मार्गांनी
देता येते; उदाहरणार्थ, ट्यूरिंग यंत्रांच्या प्रतिमानातून किंवा
सर्वसाधारण वापराच्या संगणक-भाषेतील कार्यक्रमांनी संगणित होणाऱ्या
फलनांच्या रूपात. मात्र आपले उद्दिष्ट या फलनांचे आणि संबंधांचे
भाषा~$\Lang L_A$ मधील उपपत्तीत निरूपण करणे असल्याने
केवळ संख्यात्मक फलनांच्या संगणनक्षमतेची सोपी व्याख्या निवडणे
उत्तम. ही \emph{पुनरावर्ती फलन} ही संकल्पना आहे. म्हणून प्रथम
पुनरावर्ती फलनांची चर्चा करू. मग $\Th{Q}$ आधीच सर्व पुनरावर्ती
फलनांचे आणि संबंधांचे निरूपण करते हे दाखवू. त्यामुळे $\Th{Q}$
आणि $\Th{PA}$ यांसारख्या विशिष्ट उपपत्त्यांना अपूर्णता प्रमेय
लागू करता येईल, कारण त्या ``पुरेशा प्रबळ'' उपपत्ती असल्याचे
आपण प्रस्थापित केलेले असेल.

विन्यासमीमांसेच्या अंकगणितीकरणातून अखेरीस सूत्र
$\OProv[\Gamma](x)$ मिळते. सूत्रे ना संख्यांनी सांकेतिक
करण्यामुळे हे सूत्र~$\Gamma$ च्या स्वयंसिद्धकांपासून सिद्धतायोग्यता
व्यक्त करते. विशेषतः, $\varphi$ चा संकेतांक~$n$ असेल आणि $\Gamma \Proves
\varphi$ असेल, तर $\Gamma \Proves \Prov[\Gamma](\num{n})$.
$\Gamma$ साठीचे हे ``सिद्धतायोग्यता विधेय'' विशिष्ट अर्थाने
$\Gamma$ ची सुसंगतताही~$\Lang L_A$ मधील वाक्य म्हणून
व्यक्त करू देते. $\Gamma$ साठीचे ``सुसंगतता-विधान'' हे
$\lnot \Prov[\Gamma](\num{n})$ हे वाक्य मानू, जिथे $n$ हा
एखाद्या विरोधाभासाचा, उदा. ~ $\lfalse$ चा, संकेतांक आहे.
दुसरे अपूर्णता प्रमेय सांगते की सुसंगत निर्णेय स्वयंसिद्धकसंचाने निरूपणयोग्य उपपत्ती
स्वतःची सुसंगतता-विधानेही सिद्ध करू शकत नाहीत. हे प्रमेय लागू
होण्यासाठी लागणाऱ्या अटी, उपपत्ती सर्व संगणनक्षम फलने आणि
निर्णेय संबंध निरूपित करते एवढ्यापेक्षा काहीशा अधिक
कडक आहेत; पण $\Th{PA}$ त्या पूर्ण करते हे आपण दाखवू.












\section{अनिर्णेयता आणि अपूर्णता}\label{inc:int:dec:sec}

गोडेलच्या अपूर्णता प्रमेयांच्या सिद्धतेसाठी विन्यासमीमांसेचे
अंकगणितीकरण आवश्यक आहे. पण ते न करताही, एखादी उपपत्ती सर्व
निर्णेय संबंधांचे निरूपण करते या गृहीतकावरून काही
चांगले निष्कर्ष मिळतात. ही सिद्धता थांबण्याच्या समस्येची अनिर्णेयता
सिद्ध करणाऱ्या विकर्ण युक्तिवादासारखी आहे.

\begin{thm}
$\Gamma$ ही प्रत्येक निर्णेय संबंधांचे निरूपण
करणारी सुसंगत उपपत्ती असेल, तर $\Gamma$ निर्णेय नाही.
\end{thm}

\begin{proof}
$\Gamma$ निर्णेय आहे असे समजू. $\Gamma$ प्रत्येक
निर्णेय संबंधांचे निरूपण करत असेल तर विसंगत असलीच
पाहिजे, हे दाखवू.

निर्णेय गुणधर्म (एकस्थानी संबंध) एक मुक्त चल असलेल्या
सूत्रांनी निरूपित होतात. येथे एकच मुक्त चल असलेल्या सर्व
सुविन्यस्त सूत्रांचे प्रगणन घ्यायचे आहे; फक्त निर्णेय गुणधर्मांचे
निरूपण करणाऱ्या सूत्रांचा उपवर्ग निवडायचा नाही. त्यांचे
$\varphi_0(x)$, $\varphi_1(x)$, \dots, हे संगणनक्षम प्रगणन असू द्या.
आता पुढील संच $D \subseteq \Nat$ पाहा:
\[D = \Setabs{n}{\Gamma \Proves \lnot \varphi_n(\num{n})}
\]
$D$ हा निर्णेय संच आहे. कारण प्रथम $\varphi_n(x)$ संगणित
करून, त्यावरून $\lnot \varphi_n(\num{n})$ मिळवून $n \in D$
आहे का हे तपासता येते. $\varphi_n(x)$ मधील $x$ च्या प्रत्येक मुक्त
आस्थितीऐवजी पद $\num{n}$ ठेवणे आणि $\varphi_n(\num{n})$ च्या
आधी $\lnot$ लावणे ही यांत्रिक प्रक्रिया आहे. गृहीतकानुसार
$\Gamma$ निर्णेय असल्यामुळे
$\lnot \varphi_n(\num{n}) \in \Gamma$ आहे का हे तपासता येते.
असल्यास $n \in D$, नसल्यास $n \notin D$. त्यामुळे $D$
देखील निर्णेय आहे.

$\Gamma$ सर्व निर्णेय गुणधर्मांचे निरूपण करत
असल्याने ती~$D$ चे देखील निरूपण करते. $\Gamma$ मध्ये
$D$ चे निरूपण करणारी सर्व सूत्रे ही $\varphi_0(x)$,
$\varphi_1(x)$, \dots, या यादीत आहेत. म्हणून $\varphi_d(x)$ हे
$\Gamma$ मध्ये $D$ चे निरूपण करेल असा क्रमांक $d$
घ्या. $d
\notin D$ असेल, तर $\varphi_d(x)$ हे~$D$ चे निरूपण
करत असल्यामुळे $\Gamma \Proves
\lnot \varphi_d(\num{d})$. पण मग $d$ हा~$D$ च्या व्याख्येतील
अट पूर्ण करतो, म्हणून $d \in D$. हे $d \notin D$ शी
विसंगत आहे. म्हणून अप्रत्यक्ष सिद्धतेने $d \in D$.

$d \in D$ असल्यामुळे~$D$ च्या व्याख्येनुसार $\Gamma \Proves \lnot
\varphi_d(\num{d})$. दुसरीकडे $\varphi_d(x)$ हे $\Gamma$ मध्ये~$D$ चे
निरूपण करत असल्यामुळे $\Gamma \Proves \varphi_d(\num{d})$.
म्हणून $\Gamma$ विसंगत आहे.
\end{proof}

\begin{explain}
वरील प्रमेय दाखवते की सर्व निर्णेय संबंधांचे निरूपण
करणारी कोणतीही सुसंगत उपपत्ती निर्णेय असू शकत नाही.
$\Th{Q}$ सर्व निर्णेय संबंधांचे निरूपण करते हे
आपण दाखवू. म्हणजे $\Th{Q}$ समाविष्ट असलेल्या $\Th{PA}$
आणि $\Th{TA}$ यांसारख्या सर्व उपपत्तीदेखील तसे करतात; त्यांपैकी
सुसंगत उपपत्ती निर्णेय नसतात. (येथे नाव घेतलेल्या उपपत्ती
प्रमाण प्रतिमानात सत्य असल्यामुळे सुसंगत आहेत.)

या निष्कर्षाचा वापर करून पहिल्या अपूर्णता प्रमेयाची दुर्बल आवृत्तीही
मिळते. निर्णेय स्वयंसिद्धकसंचाने निरूपणयोग्य आणि संपूर्ण असलेली कोणतीही
उपपत्ती निर्णेय असते. त्यामुळे सुसंगत, निर्णेय स्वयंसिद्धकसंचाने निरूपणयोग्य
आणि सर्व निर्णेय गुणधर्मांचे निरूपण करणारी उपपत्ती
संपूर्ण असू शकत नाही.
\end{explain}

\begin{thm}
$\Gamma$ निर्णेय स्वयंसिद्धकसंचाने निरूपणयोग्य आणि संपूर्ण असेल, तर ती
निर्णेय असते.
\end{thm}

\begin{proof}
कोणतीही विसंगत उपपत्ती निर्णेय असते, कारण विसंगत
उपपत्त्यांत सर्व वाक्ये असतात. म्हणून ``$\varphi \in
\Gamma$ आहे का?'' या प्रश्नाचे उत्तर नेहमी ``होय'' असते,
म्हणजे तो निर्णय करता येतो.

आता $\Gamma$ सुसंगत, निर्णेय स्वयंसिद्धकसंचाने निरूपणयोग्य आणि संपूर्ण आहे
असे समजू. $\Gamma$ निर्णेय स्वयंसिद्धकसंचाने निरूपणयोग्य असल्याने ती
संगणनक्षमपणे प्रगणनीय आहे. कारण~$\Gamma$ च्या
स्वयंसिद्धकांपासूनच्या सर्व योग्य निष्पत्त्यांचे संगणनक्षम
फलनाने प्रगणन करता येते. योग्य निष्पत्ती पासून ती ज्या
वाक्यानिष्पन्न करते ते संगणित करता येते. त्यामुळे
दोन्ही मिळून~$\Gamma$ च्या सर्व प्रमेयांचे प्रगणन करणारे
संगणनक्षम फलन मिळते. $\Gamma$ सुसंगत आणि संपूर्ण असल्यामुळे
वाक्य~$\varphi$ हे~$\Gamma$ चे प्रमेय असते तेव्हा आणि केवळ तेव्हाच
$\lnot \varphi$ प्रमेय नसते. म्हणून $\varphi \in
\Gamma$ आहे का हे पुढीलप्रमाणे ठरवता येते. $\Gamma$ च्या सर्व
प्रमेयांचे प्रगणन करा. या यादीत $\varphi$ आल्यास $\Gamma \Proves \varphi$
हे कळते. यादीत $\lnot
\varphi$ आल्यास $\Gamma \Proves/ \varphi$ हे कळते. $\Gamma$
संपूर्ण असल्यामुळे यांपैकी एक प्रकरण अखेरीस घडतेच;
म्हणून प्रक्रिया अखेरीस उत्तर देते.
\end{proof}

\begin{cor}
\label{inc:int:dec:cor:incompleteness}
$\Gamma$ सुसंगत, निर्णेय स्वयंसिद्धकसंचाने निरूपणयोग्य आणि प्रत्येक निर्णेय
गुणधर्मांचे निरूपण करणारी असेल, तर ती संपूर्ण नाही.
\end{cor}

\begin{proof}
$\Gamma$ संपूर्ण असेल, तर मागील प्रमेयामुळे ती
निर्णेय असेल (कारण ती निर्णेय स्वयंसिद्धकसंचाने निरूपणयोग्य आणि सुसंगत आहे).
पण $\Gamma$ प्रत्येक निर्णेय गुणधर्मांचे निरूपण
करत असल्यामुळे पहिल्या प्रमेयाने ती निर्णेय नाही.
\end{proof}

\begin{prob}
$\Th{TA} = \Setabs{\varphi}{\Sat{N}{\varphi}}$ ही निर्णेय स्वयंसिद्धकसंचाने निरूपणयोग्य
नाही, हे दाखवा. $\Th{TA}$ सर्व निर्णेय गुणधर्मांचे निरूपण करते
असे गृहीत धरू शकता.
\end{prob}

उदाहरणार्थ, $\Th{Q}$ सर्व निर्णेय गुणधर्मांचे निरूपण
करते हे आपण प्रस्थापित केल्यावर उपप्रमेयावरून $\Th{Q}$ अपूर्ण
असलीच पाहिजे असे कळेल. पण त्याची सिद्धता स्वतंत्र वाक्य
चे उदाहरण देत नाही; असे वाक्य अस्तित्वात असलेच पाहिजे
एवढेच दाखवते. त्यासाठी विन्यासमीमांसेचे अंकगणितीकरण करून
गोडेलच्या मूळ सिद्धतेची कल्पना वापरावी लागेल. आणि अर्थातच,
$\Th{Q}$ सर्व निर्णेय गुणधर्मांचे निरूपण करते हा
पहिला दावाही अजून सिद्ध करायचा आहे.

प्रत्येक \emph{रंजक} उपपत्ती अपूर्ण किंवा अनिर्णेय असेलच असे
नाही. अनेक उपपत्ती महत्त्वाची गणितीय तथ्ये व्यक्त करण्याइतक्या
प्रबळ असूनही गोडेलच्या निष्कर्षांच्या अटी पूर्ण करत नाहीत.
उदाहरणार्थ, $\Th{Pres} = \Setabs{\varphi \in \Lang{L_{A^+}}}{\Sat{N}{\varphi}}$
ही प्रमाण प्रतिमानात सत्य असलेल्या, पण~$\times$ नसलेल्या
अंकगणिताच्या भाषेतील वाक्यांची उपपत्ती संपूर्ण आणि
निर्णेय दोन्ही आहे. या उपपत्तीला प्रेस्बर्गर अंकगणित म्हणतात;
केवळ $\Obj{0}$, $\prime$ आणि~$+$ यांच्या साहाय्याने मांडता
येणारी नैसर्गिक संख्यांविषयीची सर्व सत्ये ती सिद्ध करते.



\chapter{विन्यासमीमांसेचे अंकगणितीकरण}\label{inc:art::chap}
\begin{editorial}
  चिन्हांकित टॅब्लोसाठी अंकगणितीकरण अद्याप उपलब्ध नाही, याची नोंद घ्यावी.
\end{editorial}








\section{परिचय}\label{inc:art:int:sec}

संगणनक्षमता आणि तर्कशास्त्र यांचा संबंध जोडण्यासाठी तर्कशास्त्रातील
वस्तू (चिन्हे, पदे, सूत्रे, निष्पत्त्या), त्यांवरील क्रिया,
आणि त्यांचे गुणधर्म व संबंध यांविषयी संगणकीय हाताळणीस अनुकूल रीतीने
बोलता आले पाहिजे. चिन्हे, चिन्हांचे अनुक्रम आणि त्यांपासून तयार
होणाऱ्या इतर वस्तूंवरील संगणनक्षम फलने व संबंध विचारात घेऊन हे थेट
करता येते. तार्किक विन्यासमीमांसेतील सर्व वस्तू सांत असतात आणि
चिन्हांच्या गणनीय संचांपासून बनतात; म्हणून काही संगणन-प्रतिमानांत
अशी थेट हाताळणी शक्य आहे. पण पुनरावर्ती फलनांसारखी इतर संगणन-प्रतिमाने
संख्या आणि त्यांवरील संबंध व फलने यांपुरती मर्यादित असतात. शिवाय,
अखेरीस काही उपपत्त्यांच्या आतही विन्यासमीमांसा हाताळायची आहे;
विशेषतः अंकगणिताच्या भाषेत मांडलेल्या उपपत्त्यांमध्ये. अशा वेळी
विन्यासमीमांसेचे \emph{अंकगणितीकरण} करावे लागते: विन्यासात्मक वस्तू,
त्यांवरील क्रिया आणि त्यांचे संबंध यांचे अनुक्रमे संख्या, अंकगणितीय
फलने आणि अंकगणितीय संबंध यांच्या रूपात निरूपण करावे लागते. लाइब्निझपर्यंत
मागे जाणारी यामागील कल्पना म्हणजे विन्यासात्मक वस्तूंना संख्या नेमून देणे.

चिन्हांना त्यांचे ``संकेतांक'' म्हणून संख्या देणे तुलनेने सोपे आहे.
काही चिन्हांमुळे थोडी अडचण येते; उदाहरणार्थ, चले अनंत असतात
आणि प्रत्येक स्थानसंख्या~$n$ साठी फलने देखील अनंत असतात.
तरीही चिन्हांना पद्धतशीरपणे संख्या देता येतात, जेणेकरून, उदाहरणार्थ,
$\Obj v_2$ आणि $\Obj v_3$ यांना वेगवेगळे संकेतांक मिळतील. चिन्हांचे
अनुक्रम (उदा. पदे आणि सूत्रे) सांकेतिक करणे अधिक अवघड आहे.
पण संख्यांचे अनुक्रम निव्वळ अंकगणितीय रीतीने हाताळता आले (उदा.
अविभाज्य संख्यांच्या घातांनी अनुक्रम सांकेतिक करून), तर एकेका चिन्हाचे
सांकेतीकरण चिन्हांच्या अनुक्रमांपर्यंत वाढवता येते; त्यानंतर
सूत्रांचे अनुक्रम किंवा इतर मांडण्या, जसे निष्पत्त्या,
यांपर्यंतही ते वाढवता येते. या विस्तारित सांकेतीकरणाला ``गोडेल
क्रमांकन'' म्हणतात. प्रत्येक पदाला, सूत्राला आणि
निष्पत्ती ला गोडेल संख्या नेमली जाते.

चिन्हांचे अनुक्रम त्यांच्या संकेतांकांच्या अनुक्रमांनी सांकेतिक करून,
आणि संगणनक्षम फलनांनी हाताळता येईल अशी अनुक्रम-सांकेतीकरणाची पद्धत
निवडून, गोडेल संख्याही संगणनक्षम फलनांनी हाताळता येतात. प्रत्यक्षात
येथे लागणारी सर्व फलने आदिम पुनरावर्ती असतील. उदाहरणार्थ, अनुक्रमाच्या
संकेतांकावरून त्याची लांबी आणि त्यातील $i$-वा घटक काढण्याची दोन्ही
फलने आदिम पुनरावर्ती आहेत. अनुक्रम सांकेतिक करणारी संख्या, उदाहरणार्थ,
सूत्र~$\varphi$ ची गोडेल संख्या असेल, तर सूत्राची लांबी
आणि सूत्र मधील $i$-व्या चिन्हाचा संकेतांक हेही~$\varphi$ च्या
गोडेल संख्येवरून काढता येतात, हे लगेच दिसते. एखादी संख्या योग्य
रीतीने रचलेल्या पदाची किंवा योग्य निष्पत्तीची गोडेल संख्या
आहे हा गुणधर्म आदिम पुनरावर्ती आहे, हे सिद्ध करणे थोडे अधिक अवघड आहे.
तरीही ते शक्य आहे, कारण आपल्याला अभिप्रेत अनुक्रमांची (पदे,
सूत्रे, निष्पत्त्या) व्याख्या विगमनाधारित रीतीने केलेली आहे.

उदाहरण म्हणून आदेशनाची क्रिया विचारात घेऊ. $\varphi$ हे सूत्र, $x$ हा
चर आणि $t$ हे पद असेल, तर $\varphi$ मधील~$x$ च्या प्रत्येक मुक्त
आस्थितीऐवजी~$t$ ठेवल्यावर $\Subst{\varphi}{t}{x}$ मिळते. आता $\varphi$,
$x$, $t$ यांना अनुक्रमे $k$, $l$, $m$ या गोडेल संख्या दिल्या आहेत
असे समजू. तीच पद्धत $\Subst{\varphi}{t}{x}$ लाही एखादी गोडेल संख्या,
समजा~$n$, देते. $k$, $l$, $m$ यांना~$n$ शी जोडणारे हे फलन म्हणजे
आदेशन-क्रियेचे अंकगणितीय समरूप आहे. आदेशन-क्रिया $\varphi$, $x$, $t$
यांपासून $\Subst{\varphi}{t}{x}$ देते, तर अंकगणितीकृत आदेशन-फलन
$k$, $l$, $m$ या गोडेल संख्यांपासून गोडेल संख्या~$n$ देते.
हे फलन आदिम पुनरावर्ती आहे, हे आपण पाहू.

विन्यासमीमांसेचे अंकगणितीकरण केवळ अमूर्त रुचीचे नाही. अंकगणिताची
भाषा यांसारख्या भाषांत भाषांविषयी ``बोलण्याची'' साधने आधीपासून
नसूनही व्यंजकांचे गुंतागुंतीचे गुणधर्म त्यांत आकारिक रीतीने व्यक्त
करता येतात, ही कल्पना सुरुवातीला सहज सुचण्यासारखी नव्हती. मग
अशा भाषेतील, उदाहरणार्थ पेआनो अंकगणितासारखी, उपपत्ती स्वतःच्या
भाषेविषयी काय \emph{सिद्ध} करू शकते, हा प्रश्न विचारणे स्वाभाविक
ठरते (उदा. वाक्ये सिद्ध करता येतात का किंवा ती सत्य आहेत
का). यातून गोडेलची (असिद्धतायोग्यतेविषयी) आणि टार्स्कीची
(सत्याच्या अव्याख्येयतेविषयी) मर्यादा दाखवणारी प्रसिद्ध प्रमेये
समोर येतात. पण संगणनक्षमता उपपत्तीतील काही महत्त्वाचे निष्कर्ष
सिद्ध करण्यासाठीही विन्यासमीमांसेचे अंकगणितीकरण महत्त्वाचे आहे;
उदा. उपपत्त्यांची संगणकीय शक्ती किंवा संगणनक्षमतेच्या विविध
प्रतिमानांतील संबंध यांविषयीचे निष्कर्ष. इतर वस्तू आणि गुणधर्मांचे
अंकगणितीकरण कसे करावे याचा नमुनाही विन्यासमीमांसेचे अंकगणितीकरण
देते. उदाहरणार्थ, (ट्यूरिंग यंत्रांच्या) स्थितिवर्णनांचे आणि
संगणनांचेही अशाच रीतीने अंकगणितीकरण करता येते. त्यामुळे एका
प्रतिमानातील संगणनांचे (उदा. ट्यूरिंग यंत्रांतील) दुसऱ्या प्रतिमानात
(उदा. पुनरावर्ती फलनांत) अनुकरण करता येते.











\section{चिन्हांचे सांकेतीकरण}\label{inc:art:cod:sec}

प्रथम-क्रम तर्कशास्त्राच्या मूलभूत भाषा~$\Lang L$ मध्ये ही चिन्हे वापरली जातात:
\[\lfalse \quad \lnot \quad \lor \quad \land \quad \lif \quad \lforall
\quad \lexists \quad \eq \quad ( \quad ) \quad ,
\]
यांबरोबर चरे व स्थिरांक यांचे गणनीय संच, तसेच कोणत्याही
स्थानसंख्येची फलने व विधेये यांचे गणनीय संचही
वापरले जातात. या प्रत्येक चिन्हाला \emph{संकेतांक} म्हणून एक संख्या
देता येते: प्रत्येक चिन्हाचा संकेतांक वेगळा असेल आणि दोन भिन्न
चिन्हांना एकच संख्या दिली जाणार नाही. सर्व चिन्हांचा संच गणनीय
असल्यामुळे त्याची नैसर्गिक संख्यांच्या संचाशी एकास-एक आच्छादन आहे;
म्हणून हे शक्य आहे. परंतु संकेतांकावरून चिन्ह आणि त्याविषयीची काही
माहिती (उदा. फलन ची स्थानसंख्या) संगणनक्षम रीतीने परत
मिळाली पाहिजे. असे करण्याचे अनेक मार्ग आहेत. पुढील मार्गात आदिम
पुनरावर्ती फलने वापरली आहेत. ($\tuple{n_0, \dots, n_k}$ ही
$n_0$, \dots, $n_k$ या संख्यांच्या अनुक्रमाला सांकेतिक करणारी संख्या
आहे, हे आठवा.)

\begin{defn}
$s$ हे~$\Lang L$ मधील चिन्ह असेल, तर त्याचा \emph{चिन्हसंकेतांक}~$\scode s$
पुढीलप्रमाणे परिभाषित करू:
\begin{enumerate}
\item $s$ हे तार्किक चिन्हांपैकी असेल, तर पुढील सारणीनुसार $\scode s$ ठरतो:
\[\begin{array}{cccccccccc}
  \lfalse & \lnot & \lor & \land & \lif & \lforall \\
\tuple{0, 0} & \tuple{0, 1} & \tuple{0, 2} & \tuple{0, 3} &
\tuple{0, 4} & \tuple{0, 5} \\
\lexists & \eq & ( & ) & ,\\
\tuple{0, 6} & \tuple{0, 7} &
\tuple{0, 8} & \tuple{0, 9} & \tuple{0, 10}
\end{array}
\]
\item $s$ हा $i$-वा चर $\Obj v_i$ असेल, तर $\scode s = \tuple{1, i}$.
\item $s$ हे $i$-वे स्थिरांक~$\Obj c_i$ असेल, तर
  $\scode s = \tuple{2, i}$.
\item $s$ हे $i$-वे $n$-स्थानी फलन~$\Obj f_i^n$ असेल, तर
  $\scode s = \tuple{3, n, i}$.
\item $s$ हे $i$-वे $n$-स्थानी विधेय~$\Obj P_i^n$ असेल, तर
  $\scode s = \tuple{4, n, i}$.
\end{enumerate}
\end{defn}

\begin{prop}
पुढील संबंध आदिम पुनरावर्ती आहेत:
\begin{enumerate}
\item $\fn{Fn}(x, n)$ तेव्हा आणि केवळ तेव्हाच जेव्हा $x$ हा एखाद्या~$i$ साठी
  $\Obj f^n_i$ चा संकेतांक असतो; म्हणजेच $x$ हा $n$-स्थानी
  फलन चा संकेतांक असतो.
\item $\fn{Pred}(x, n)$ तेव्हा आणि केवळ तेव्हाच जेव्हा $x$ हा एखाद्या~$i$ साठी
  $\Obj P^n_i$ चा संकेतांक असतो, किंवा $x$ हा $\eq$ चा संकेतांक
  असतो आणि $n = 2$ असते; म्हणजेच $x$ हा $n$-स्थानी विधेय
  चा संकेतांक असतो.
\end{enumerate}
\end{prop}

\begin{defn}
$s_0, \dots, s_{n-1}$ हा चिन्हांचा अनुक्रम असेल, तर त्याची
\emph{गोडेल संख्या} $\tuple{\scode{s_0}, \dots, \scode{s_{n-1}}}$ आहे.
\end{defn}

\begin{explain}
\emph{संकेतांक} आणि \emph{गोडेल संख्या} वेगवेगळ्या गोष्टी आहेत,
हे लक्षात घ्या. उदाहरणार्थ, $\Obj v_5$ या चराचा संकेतांक
$\scode{\Obj v_5} = \tuple{1, 5} = 2^2\cdot 3^6$ आहे.
पण पद म्हणून पाहिले असता~$\Obj v_5$ हाही लांबी~$1$ असलेला
चिन्हांचा अनुक्रम आहे. त्यामुळे $\Obj v_5$ या \emph{पदाची}
\emph{गोडेल संख्या}~$\Gn{\Obj v_5}$ ही
$\tuple{\scode{\Obj v_5}} = 2^{\scode{\Obj v_5} + 1} =
2^{2^2\cdot 3^6 + 1}$ आहे.
\end{explain}

\begin{ex}
$k_0$, \dots, $k_{n-1}$ हा संख्यांचा अनुक्रम असेल, तर
अविभाज्य संख्यांच्या घातांनी केलेल्या सांकेतीकरणात
$\tuple{k_0, \dots, k_{n-1}}$ हा त्याचा संकेतांक
\[2^{k_0+1}\cdot3^{k_1+1}\cdot \dots \cdot p_{n-1}^{k_{n-1}+1},
\]
आहे, हे आठवा. येथे $p_i$ ही $i$-वी अविभाज्य संख्या आहे
($p_0 = 2$ पासून सुरुवात होते). म्हणून, उदाहरणार्थ,
$\eq[\Obj v_0][\Obj 0]$ हे सूत्र, किंवा अधिक स्पष्टपणे
${\eq}(\Obj v_0,\Obj c_0)$, याची गोडेल संख्या
\[\tuple{\scode{\eq},\scode{(},\scode{\Obj v_0},\scode{,}, \scode{\Obj
    c_0},\scode{)}}.
\]
आहे. येथे $\scode{\eq}$ हा $\tuple{0,7} = 2^{0+1}\cdot
3^{7+1}$ आहे; $\scode{\Obj v_0}$ हा $\tuple{1,0} = 2^{1+1}\cdot3^{0+1}$
आहे; इत्यादी. म्हणून $\Gn{=(\Obj v_0,\Obj c_0)}$ ही संख्या
\begin{multline*}
2^{\scode{=} + 1}\cdot 3^{\scode{(}+1}\cdot 5^{\scode{\Obj v_0}+1}
\cdot 7^{\scode{,} + 1} \cdot 11^{\scode{\Obj c_0}+1} \cdot
13^{\scode{)}+1} = \\
2^{2^1\cdot 3^8 + 1}\cdot 3^{2^1\cdot 3^9+1}\cdot 5^{2^2\cdot 3^1+1}
\cdot 7^{2^1\cdot 3^{11} + 1} \cdot 11^{2^3\cdot3^1+1} \cdot
13^{2^1\cdot3^{10}+1} = \\
2^{13\,123}\cdot 3^{39\,367}\cdot 5^{13}\cdot 7^{354\,295}\cdot11^{25}\cdot13^{118\,099}.
\end{multline*}
\end{ex}











\section{पदांचे सांकेतीकरण}\label{inc:art:trm:sec}

\begin{explain}
पद हा चिन्हांच्या अनुक्रमाचा एक विशिष्ट प्रकार आहे: पदांसाठीच्या
रचनानियमांनुसार स्थिरचिन्हे आणि चरे यांपासून त्याची विगमनाधारित
रचना होते. चिन्हे सांकेतिक करण्याची पद्धत आणि संख्यांचे अनुक्रम
सांकेतिक करण्याची पद्धत वापरून चिन्हांचे अनुक्रम संख्यांनी सांकेतिक
करता येतात. त्यामुळे पदांना गोडेल संख्या देणे कठीण नाही. खरी
अडचण ही आहे की एखादी संख्या योग्य रीतीने रचलेल्या पदाची गोडेल
संख्या आहे हा गुणधर्म संगणनक्षम, किंबहुना आदिम पुनरावर्ती आहे, हे दाखवणे.
\end{explain}

चले आणि स्थिरांक ही सर्वांत सोपी पदे आहेत. $x$ ही
अशा पदाची गोडेल संख्या आहे का, हे तपासणे सोपे आहे:
$x$ ही एखाद्या~$i$ साठी $\Gn{\Obj v_i}$ असेल, तर आणि तरच
$\fn{Var}(x)$ खरे असते. दुसऱ्या शब्दांत, $x$ हा लांबी~$1$
असलेल्या अनुक्रमाचा संकेतांक आहे आणि त्यातील एकमेव घटक $(x)_0$
हा एखाद्या चल~$\Obj v_i$ चा संकेतांक आहे; म्हणजे एखाद्या~$i$
साठी $x$ हा $\tuple{\tuple{1, i}}$ आहे. तसेच, $x$ ही एखाद्या~$i$
साठी $\Gn{\Obj c_i}$ असेल, तर आणि तरच $\fn{Const}(x)$ खरे असते.
हे दोन्ही संबंध आदिम पुनरावर्ती आहेत; कारण असा~$i$ अस्तित्वात
असेल, तर तो $< x$ असलाच पाहिजे:
\begin{align*}
  \fn{Var}(x) & \defiff \bexists{i<x}{x = \tuple{\tuple{1, i}}}\\
  \fn{Const}(x) & \defiff \bexists{i<x}{x = \tuple{\tuple{2, i}}}
\end{align*}

\begin{prop}
\label{inc:art:trm:prop:term-primrec}
$x$ ही अनुक्रमे पदाची किंवा बंद पदाची गोडेल संख्या असेल, तर आणि तरच
खरे असणारे $\fn{Term}(x)$ आणि $\fn{ClTerm}(x)$ हे संबंध आदिम पुनरावर्ती आहेत.
\end{prop}

\begin{proof}
चिन्हांचा अनुक्रम~$s$ हा पद असेल, तर आणि तरच पदांची अशी
रचनाक्रमिका~$s_0$, \dots, $s_{k-1} = s$ अस्तित्वात असते की जिच्यात
पदांच्या रचनानियमांनुसार स्थिरांक आणि चले यांपासून
पद~$s$ कसे बनले याची नोंद असते. एखादी प्रस्तावित रचनाक्रमिका
रचनानियमांचे पालन करते ही अट व्यक्त करण्यासाठी, प्रत्येक $i < k$
बाबत पुढीलपैकी एक अट खरी असली पाहिजे:
\begin{enumerate}
\item $s_i$ हे चल~$\Obj v_j$ आहे; किंवा
\item $s_i$ हे स्थिरांक~$\Obj c_j$ आहे; किंवा
\item $i$-व्या स्थानापूर्वी आलेल्या $t_1$, \dots, $t_n$ या $n$
  पदांपासून $n$-स्थानी फलन~$\Obj f^n_j$ वापरून $s_i$ बनले आहे.
\end{enumerate}
गोडेल संख्यांवरील संबंधित संबंध आदिम पुनरावर्ती आहे हे दाखवण्यासाठी
ही अट आदिम पुनरावर्ती फलने, संबंध आणि परिबद्ध संख्यकीकरण वापरून
व्यक्त करावी लागेल.

$y$ ही $s_0$, \dots, $s_{k-1}$ या अनुक्रमाचा संकेतांक आहे असे समजू;
म्हणजे $y = \tuple{\Gn{s_0}, \dots, \Gn{s_{k-1}}}$. तो गोडेल
संख्या~$x$ असलेल्या पदाच्या रचनाक्रमिकेचा संकेतांक असेल, तर आणि तरच
प्रत्येक $i < k$ साठी पुढीलपैकी एक अट खरी असते:
\begin{enumerate}
\item $\fn{Var}((y)_i)$; किंवा
\item $\fn{Const}((y)_i)$; किंवा
\item एखादा~$n$ आणि $z = \tuple{z_1, \dots, z_n}$ ही संख्या
  अस्तित्वात असतात, अशी की प्रत्येक $z_l$ हा $i' < i$ असलेल्या
  एखाद्या निर्देशांकासाठी $(y)_{i'}$ बरोबर असतो आणि
\[(y)_i = \Gn{\Obj f^n_j(} \concat \fn{flatten}(z) \concat \Gn{)},
\]
\end{enumerate}
आणि याशिवाय $(y)_{k-1} = x$ असते. ($\fn{flatten}(z)$ हे फलन
$\tuple{\Gn{t_1}, \dots, \Gn{t_n}}$ या अनुक्रमापासून
$\Gn{t_1, \dots, t_n}$ तयार करते आणि ते आदिम पुनरावर्ती आहे.)

(3) मधील $j$, $n$ हे निर्देशांक, $t_l$ या पदांच्या गोडेल संख्या
$z_l$, आणि $\tuple{z_1, \dots, z_n}$ या अनुक्रमाचा संकेतांक~$z$
हे सर्व~$y$ पेक्षा लहान आहेत. वरच्या~$k$ ऐवजी $\len{y}$
वापरता येते. म्हणून ``$y$ ही गोडेल संख्या~$x$ असलेल्या पदाच्या
रचनाक्रमिकेचा संकेतांक आहे'' हे वाक्य त्या संबंधाचे आदिम पुनरावर्ती
असणे स्पष्ट करणाऱ्या रीतीने व्यक्त करता येते.

आता~$y$ साठी आदिम पुनरावर्ती मर्यादा आहे याचीच खात्री करून घ्यायची
आहे. $x$ ही एखाद्या पदाची गोडेल संख्या असेल, तर त्याला जास्तीत
जास्त $\len{x}$ पदांची रचनाक्रमिका असली पाहिजे: कारण~$s$ च्या
केवळ उपपदांची रचनाक्रमिका निवडता येते. तिच्यातील प्रत्येक पदाची
सुरुवात~$s$ मधील एखाद्या स्थानावर
होते, आणि दोन उपपदे एकाच स्थानावर सुरू होऊ शकत नाहीत. $s$ च्या
प्रत्येक उपपदाची गोडेल संख्या अर्थातच $\le x$ आहे. म्हणून
$k=\len{x}$ असेल, तर संकेतांक $\le p_{k-1}^{k(x+1)}$ असलेली
रचनाक्रमिका नेहमी अस्तित्वात असते.


$\fn{ClTerm}$ साठी चले संबंधी अट फक्त वगळावी.
\end{proof}

\begin{prob}
$\tuple{\Gn{t_1}, \dots, \Gn{t_n}}$ या अनुक्रमापासून
$\Gn{t_1, \dots, t_n}$ तयार करणारे $\fn{flatten}(z)$ हे फलन
आदिम पुनरावर्ती आहे, हे दाखवा.
\end{prob}

\begin{prop}\label{inc:art:trm:prop:num-primrec}
  $\fn{num}(n) = \Gn{\num{n}}$ हे फलन आदिम पुनरावर्ती आहे.
\end{prop}

\begin{proof}
  $\fn{num}(n)$ ची आदिम पुनरावर्तनाने व्याख्या करू:
  \begin{align*}
    \fn{num}(0) & = \Gn{\Obj 0}\\
    \fn{num}(n+1) & = \Gn{\prime(} \concat \fn{num}(n) \concat \Gn{)}.
  \end{align*}
\end{proof}











\section{सूत्र यांचे सांकेतीकरण}\label{inc:art:frm:sec}

$\fn{Term}(x)$ हा संबंध आदिम पुनरावर्ती रीतीने परिभाषित केल्यावर
सूत्रे साठीचा संबंधित संबंध $\fn{Frm}(x)$ हाही आदिम
पुनरावर्ती रीतीने परिभाषित करता येतो.

\begin{prop}
$x$ ही आण्विक सूत्राची गोडेल संख्या असेल, तर आणि तरच
खरे असणारा $\fn{Atom}(x)$ हा संबंध आदिम पुनरावर्ती आहे.
\end{prop}

\begin{proof}
$x$ ही आण्विक सूत्राची गोडेल संख्या असेल, तर आणि तरच
पुढीलपैकी एक अट खरी असते:
\begin{enumerate}
\item $n$, $j < x$ आणि $z \le B(x)$ अशी मूल्ये अस्तित्वात आहेत की
  $\len{z}=n$ आणि प्रत्येक $i < n$ साठी $\fn{Term}((z)_i)$ खरे आहे आणि $x = $
\[\Gn{\Obj P^n_j(} \concat \fn{flatten}(z) \concat \Gn{)}.
\]
\item $z_1, z_2 < x$ अशी मूल्ये अस्तित्वात आहेत की $\fn{Term}(z_1)$,
  $\fn{Term}(z_2)$ खरे आहेत आणि $x = $
\[\Gn{{\eq}(} \concat z_1 \concat \Gn{,} \concat z_2 \concat{\Gn{)}}.
\]
  \item $x = \Gn{\lfalse}$.
  \item $x = \Gn{\ltrue}$.
\end{enumerate}
येथे $k=\len{x}\ge1$ आणि $B(x)=p_{k-1}^{k(x+1)}$ घ्या.
शून्य लांबीचा अनुक्रम आण्विक सूत्र नाही. प्रत्येक आदानपदाचा संकेतांक
अंतिम सूत्राच्या संकेतांकापेक्षा लहान असतो; मात्र त्या संकेतांकांच्या
संपूर्ण अनुक्रमाचा संकेतांक तसाच लहान असेल असे नाही. आदानपदांची
संख्या $n\le k$ असल्याने $B(x)$ ही त्या अनुक्रमासाठी आदिम पुनरावर्ती
वरची मर्यादा आहे.

\end{proof}

\begin{prop}
\label{inc:art:frm:prop:frm-primrec}
$x$ ही सूत्राची गोडेल संख्या असेल, तर आणि तरच
खरे असणारा $\fn{Frm}(x)$ हा संबंध आदिम पुनरावर्ती आहे.
\end{prop}

\begin{proof}
चिन्हांचा अनुक्रम~$s$ हे सूत्र असेल, तर आणि तरच
सूत्राची अशी रचनाक्रमिका~$s_0$, \dots, $s_{k-1} = s$
अस्तित्वात असते की जिच्यात रचनानियमांनुसार आण्विक सूत्रे
पासून~$s$ कसे बनले याची नोंद असते. केवळ उपसूत्रांची रचनाक्रमिका
निवडल्यास प्रत्येक $s_i$ चा संकेतांक
$s$ च्या संकेतांक~$x$ पेक्षा मोठा नसतो; पण संपूर्ण रचनाक्रमिका
$\tuple{s_0, \dots, s_{k-1}}$ हिचा संकेतांक त्यापेक्षा लहान
असणे आवश्यक नाही: मूळ इंग्रजीतील उलट दावा एका आण्विक सूत्रासाठीही
खोटा ठरतो. त्याऐवजी मागील पदांच्या सिद्धतेप्रमाणे त्याला
त्या अंतिम सूत्राच्या संकेतांकावर अवलंबून आदिम पुनरावर्ती वरची
मर्यादा देता येते.
\end{proof}

\begin{prob}
\ref{inc:art:frm:prop:frm-primrec} ची,
\ref{inc:art:trm:prop:term-primrec} च्या पहिल्या सिद्धतेच्या
धर्तीवर, सविस्तर सिद्धता द्या.
\end{prob}

\begin{prop}
\label{inc:art:frm:prop:freeocc-primrec}
$x$ ही गोडेल संख्या असलेल्या सूत्रातील $i$-वे चिन्ह हे
गोडेल संख्या~$z$ असलेल्या चराची मुक्त आस्थिती असेल,
तर आणि तरच खरे असणारा $\fn{FreeOcc}(x, z, i)$ हा संबंध
आदिम पुनरावर्ती आहे.
\end{prop}

\begin{proof}
सराव.
\end{proof}

\begin{prob}
\ref{inc:art:frm:prop:freeocc-primrec} सिद्ध करा.
सूत्राच्या ज्या उपचिन्हमाला स्वतः सूत्र आहेत,
त्या त्याची उप-सूत्र आहेत, हे तथ्य वापरू शकता.
\end{prob}

\begin{prop}
$x$ ही वाक्याची गोडेल संख्या असेल, तर आणि तरच
खरे असणारा $\fn{Sent}(x)$ हा गुणधर्म आदिम पुनरावर्ती आहे.
\end{prop}

\begin{proof}
वाक्य म्हणजे चलांच्या मुक्त आस्थिती नसलेले
सूत्र. म्हणून $\fn{Sent}(x)$ तेव्हा आणि केवळ तेव्हाच खरे असते जेव्हा
\begin{multline*}
\fn{Frm}(x) \land \bforall{i<\len{x}}{\bforall{z<x}{}}\\
(\bexists{j<z}{z=\Gn{\Obj v_j}} \lif \lnot\fn{FreeOcc}(x,z,i)).
\end{multline*}

\end{proof}












\section{आदेशन}\label{inc:art:sub:sec}

आदेशन म्हणजे सूत्र~$\varphi$ मध्ये चल~$u$ च्या
सर्व मुक्त आस्थितींच्या जागी पद~$t$ ठेवण्याची क्रिया आहे;
ती $\Subst{\varphi}{t}{u}$ अशी लिहितात, हे आठवा. चले,
सूत्रे आणि पदे यांच्या गोडेल संख्यांवर ही क्रिया
केली, तर ती आदिम पुनरावर्ती असते.

\begin{prop}\label{inc:art:sub:prop:subst-primrec}
पुढील गुणधर्म असलेले आदिम पुनरावर्ती फलन $\fn{Subst}(x, y, z)$
अस्तित्वात आहे:
\[\fn{Subst}(\Gn{\varphi}, \Gn{t}, \Gn{u}) = \Gn{\Subst{\varphi}{t}{u}}.
\]
\end{prop}

\begin{proof}
आता $\fn{hSubst}$ हे फलन पुढीलप्रमाणे आदिम पुनरावर्तनाने
परिभाषित करता येते:
\begin{multline*}
\begin{aligned}
\fn{hSubst}(x, y, z, 0) & = \emptyseq \\
\fn{hSubst}(x, y, z, i+1) & =
\end{aligned}\\
\begin{cases}
\fn{hSubst}(x, y, z, i) \concat y & \text{जर $\fn{FreeOcc}(x, z, i)$ खरे असेल तर} \\
\fn{append}(\fn{hSubst}(x, y, z, i), (x)_{i}) & \text{अन्यथा.}
\end{cases}
\end{multline*}
आता $\fn{Subst}(x, y, z)$ ची व्याख्या
$\fn{hSubst}(x, y, z, \len{x})$ अशी करता येते.
\end{proof}

\begin{prop}
\label{inc:art:sub:prop:free-for}
गोडेल संख्या~$y$ असलेले पद, गोडेल संख्या~$x$ असलेल्या सूत्रात
गोडेल संख्या~$z$ असलेल्या चरासाठी आदेशनासाठी मुक्त असेल,
तर आणि तरच खरे असणारा $\fn{FreeFor}(x, y, z)$ हा संबंध
आदिम पुनरावर्ती आहे.
\end{prop}

\begin{proof} सराव. \end{proof}

\begin{prob}
\ref{inc:art:sub:prop:free-for} सिद्ध करा.
\end{prob}











\section{$\Log{LK}$ मधील निष्पत्ती}\label{inc:art:plk:sec}

\begin{explain}
निष्पत्त्यांचे अंकगणितीकरण करण्यासाठी निष्पत्त्या संख्यांनी दर्शवावे लागतात.
प्रत्येक निगमनाला नामांकन असलेले क्रमवर्तींचे वृक्ष म्हणजे निष्पत्त्या
असल्यामुळे, पुनरावर्ती निरूपण हा सहज मार्ग आहे. आपण निष्पत्ती
एका क्रमित समूहाने दर्शवतो: त्याचे घटक म्हणजे अंतिम क्रमवर्ती, नामांकन,
आणि अखेरच्या निगमनाच्या आधारविधानांपर्यंत येणाऱ्या उप-निष्पत्त्यांची निरूपणे.
\end{explain}

\begin{defn}
$\Gamma$ ही वाक्यांची सांत क्रमिका, $\Gamma =
\tuple{\varphi_1, \dots, \varphi_n}$, असेल, तर $\Gn{\Gamma} = \tuple{\Gn{\varphi_1},
  \dots, \Gn{\varphi_n}}$.

$\Gamma \Sequent \Delta$ ही क्रमवर्ती असेल, तर
$\Gamma \Sequent \Delta$ ची गोडेल संख्या
पुढीलप्रमाणे असते:
\[\Gn{\Gamma \Sequent \Delta} = \tuple{\Gn{\Gamma}, \Gn{\Delta}}
\]

$\pi$ ही $\Log{LK}$ मधील निष्पत्ती असेल, तर $\Gn{\pi}$ ची
व्याख्या पुढीलप्रमाणे करतो:
\begin{enumerate}
\item $\pi$ मध्ये फक्त $\Gamma \Sequent \Delta$ ही आरंभीची क्रमवर्ती
  असेल, तर $\Gn{\pi}$ म्हणजे
  \[    \tuple{0, \Gn{\Gamma \Sequent \Delta}}.
  \]
\item $\pi$ चा शेवट एक किंवा दोन आधारविधानांच्या निगमनाने होत असेल,
  त्या निगमनाचा निष्कर्ष $\Gamma \Sequent \Delta$ असेल, आणि $\pi_1$
  व $\pi_2$ या अखेरच्या निगमनाच्या आधारविधानांवर संपणाऱ्या तात्काळ
  उपनिष्पत्ती असतील, तर $\Gn{\pi}$ अनुक्रमे
  \begin{align*}
    & \tuple{1, \Gn{\pi_1}, \Gn{\Gamma \Sequent \Delta}, k} \text{ किंवा}\\
    & \tuple{2, \Gn{\pi_1}, \Gn{\pi_2}, \Gn{\Gamma \Sequent \Delta}, k},
  \end{align*}
  असते. अखेरच्या निगमनात कोणता नियम वापरला आहे त्यानुसार $k$ पुढील
  तक्त्यात दिले आहे:

  \begin{tabular}{lccccccc}
    \text{नियम:} & \LeftR{\Weakening} & \RightR{\Weakening} &
    \LeftR{\Contraction} & \RightR{\Contraction} &
    \LeftR{\Exchange} & \RightR{\Exchange} \\
    $k$: & 1 & 2 & 3 & 4 & 5 & 6 \\[2ex]
    \text{नियम:} & \LeftR{\lnot} & \RightR{\lnot} &
    \LeftR{\land} & \RightR{\land} &
    \LeftR{\lor} & \RightR{\lor} \\
    $k$: & 7 & 8 & 9 & 10 & 11 & 12 \\[2ex]
    \text{नियम:} & \LeftR{\lif} & \RightR{\lif} &
    \LeftR{\lforall} & \RightR{\lforall} &
    \LeftR{\lexists} & \RightR{\lexists} \\
    $k$: & 13 & 14 & 15 & 16 & 17 & 18 \\[2ex]
    \text{नियम:} & \Cut & = \\
    $k$: & 19 & 20
  \end{tabular}
\end{enumerate}
\end{defn}
\begin{ex}
  पुढील अगदी सोपी निष्पत्ती पाहा:
  \begin{prooftree}
    \Axiom$\varphi \fCenter \varphi$
    \RightLabel{\LeftR{\land}}
    \UnaryInf$\varphi \land \psi \fCenter \varphi$
    \RightLabel{\RightR{\lif}}
    \UnaryInf$\fCenter (\varphi \land \psi) \lif \varphi$
  \end{prooftree}
  आरंभीच्या क्रमवर्तीची गोडेल संख्या $p_0 = \tuple{0,
    \Gn{\varphi \Sequent \varphi}}$ असेल. $\LeftR{\land}$ च्या निष्कर्षावर
    संपणाऱ्या निष्पत्तीची गोडेल संख्या $p_1 =
    \tuple{1, p_0, \Gn{\varphi \land \psi \Sequent \varphi}, 9}$ असेल. येथे $1$
    हे $\LeftR{\land}$ ला एकच आधारविधान असल्यामुळे आले आहे;
    $\varphi \land \psi \Sequent \varphi$ या निष्कर्षाची गोडेल संख्या पुढे आहे;
    आणि $9$ हा $\LeftR{\land}$ चा संकेतांक आहे. मग संपूर्ण
    निष्पत्तीची गोडेल संख्या
    $\tuple{1, p_1, \Gn{\Sequent (\varphi \land \psi) \lif \varphi}, 14}$,
    म्हणजेच
  \[
  \tuple{1, \tuple{1, \tuple{0, \Gn{\varphi \Sequent \varphi}}, \Gn{\varphi \land \psi \Sequent \varphi}, 9},
    \Gn{\Sequent (\varphi \land \psi) \lif \varphi}, 14}.
  \]
\end{ex}

\begin{explain}
निष्पत्त्यांचे निरूपण ठरल्यानंतर, अशा निष्पत्त्यांवर आदिम पुनरावर्ती रीतीने
क्रिया करता येतात आणि त्यांचे आवश्यक गुणधर्म व संबंध तसेच व्यक्त करता
येतात, हेही दाखवावे लागेल. काही क्रिया सोप्या आहेत. उदाहरणार्थ,
निष्पत्तीची गोडेल संख्या~$p$ दिली असता,
$\fn{EndSequent}(p) = (p)_{(p)_0+1}$ हे तिच्या अंतिम क्रमवर्तीची
गोडेल संख्या देते आणि $\fn{LastRule}(p) = (p)_{(p)_0+2}$ हे तिच्या
अखेरच्या नियमाचा संकेतांक देते. पुढीलप्रमाणे परिभाषित केलेला
क्रमिकेच्या संकेतांकाची वैधता तपासण्यासाठी आदिम पुनरावर्ती अट घ्या:
\[\fn{Seq}(u) \defiff u=0 \lor u=\prod_{i<\len{u}}p_i^{(u)_i+1}.
\]
ती रिक्त क्रमिकेचे दोन्ही संकेतांक स्वीकारते आणि अतिरिक्त अविभाज्य
घटक असलेले अवैध संकेतांक वगळते.
$\fn{Sequent}(s)$ गुणधर्म
\[  \begin{gathered}
  s=\tuple{(s)_0,(s)_1} \land \fn{Seq}((s)_0) \land \fn{Seq}((s)_1)\\
  {}\land \bforall{i<\len{(s)_0} + \len{(s)_1}}{\fn{Sent}(((s)_0 \concat (s)_1)_i)}
  \end{gathered}
\]
$s$ साठी खरा असतो, तर आणि तरच $s$ ही वाक्ये असलेल्या
क्रमवर्तीची गोडेल संख्या असते. काही इतर क्रिया अधिक कठीण आहेत; त्या कशा कराव्यात याची
किमान रूपरेषा आपण पाहू. $\pi$ ही $\Gamma$ पासून $\varphi$ ची
निष्पत्ती आहे, हा संबंध $\pi$ आणि $\varphi$ यांच्या गोडेल संख्यांचा
आदिम पुनरावर्ती संबंध आहे, हे दाखवणे आपले उद्दिष्ट आहे.
\end{explain}

\begin{prop}\label{inc:art:plk:prop:followsby}
  गोडेल संख्या~$p$ असलेल्या निष्पत्ती~$\pi$ मधील अखेरचे निगमन
  योग्य असेल, तर आणि तरच खरा असणारा $\fn{Correct}(p)$ गुणधर्म
  आदिम पुनरावर्ती आहे.
\end{prop}

\begin{proof}
  $\Gamma \Sequent \Delta$ ही आरंभीची क्रमवर्ती दोनपैकी एका
  परिस्थितीत असते: एखादे वाक्य~$\varphi$ असे असते की
  $\Gamma \Sequent \Delta$ म्हणजे $\varphi \Sequent \varphi$ असते;
  किंवा एखादे बंद पद~$t$ असे असते की
  $\Gamma \Sequent \Delta$ म्हणजे $\emptyset \Sequent \eq[t][t]$ असते.
  गोडेल संख्यांच्या भाषेत, मूलभूत आरंभीच्या क्रमवर्तींसाठी
  $\fn{InitialSeq}(s)$ पुढील अट पूर्ण झाली तर आणि तरच खरे असते;
  चिन्हांकित आवृत्त्यांत खाली दिलेल्या अटीही जोडायच्या आहेत:
  \begin{align*}
  \bexists{x < s}{} (\fn{Sent}(x) & \land
  s = \tuple{\tuple{x},\tuple{x}})
  \lor {}\\
  \bexists{t<s}{} (\fn{ClTerm}(t) & \land
  s = \tuple{0, \tuple{\Gn{{\eq}(} \concat t \concat \Gn{,} \concat t \concat \Gn{)}}}).
  \end{align*}
  या चिन्हांकित आवृत्तीत $s=\tuple{0,\tuple{\Gn{\ltrue}}}$
  हीदेखील आरंभीची क्रमवर्ती आहे.
  या चिन्हांकित आवृत्तीत $s=\tuple{\tuple{\Gn{\lfalse}},0}$
  हीदेखील आरंभीची क्रमवर्ती आहे.

  प्रत्येक निगमन नियम~$R$ साठी $\fn{FollowsBy}_R(p)$ हा संबंधही
  आदिम पुनरावर्ती आहे, हे दाखवावे लागेल. $p$ ही निष्पत्ती~$\pi$
  ची गोडेल संख्या असेल आणि $\pi$ ची अंतिम क्रमवर्ती तिच्या तात्काळ
  $\pi$ च्या उप-निष्पत्त्यांच्या अंतिम क्रमवर्तींपासून $R$ च्या योग्य वापराने
  निष्पन्न होत असेल, तर आणि तरच $\fn{FollowsBy}_R(p)$ खरे असते.

  \RightR{\land} नियमाचे उदाहरण सोपे आहे. $\pi$ चा शेवट
  $\RightR{\land}$ च्या योग्य वापराने होत असेल, तर तिचे रूप असे असते:
  \begin{prooftree}
    \AxiomC{}
    \RightLabel{$\pi_1$}
    \Deduce$\Gamma \fCenter \Delta, \varphi$

    \AxiomC{}
    \RightLabel{$\pi_2$}
    \Deduce$\Gamma \fCenter \Delta, \psi$

    \RightLabel{\RightR\land}
    \BinaryInf$\Gamma \fCenter \Delta, \varphi \land \psi$
  \end{prooftree}
  त्यामुळे निष्पत्ती~$\pi$ मधील अखेरचे निगमन $\RightR{\land}$ चा योग्य वापर
  असेल, तर आणि तरच $\Gamma$ व $\Delta$ या वाक्यांच्या क्रमिका
  आणि $\varphi$, $\psi$ ही दोन वाक्ये अशी असतात की $\pi_1$ ची
  अंतिम क्रमवर्ती $\Gamma \Sequent \Delta, \varphi$, $\pi_2$ ची अंतिम
  क्रमवर्ती $\Gamma \Sequent \Delta, \psi$, आणि $\pi$ ची अंतिम
  क्रमवर्ती $\Gamma \Sequent \Delta, \varphi \land \psi$ असते.
  हेच आता गोडेल संख्यांत मांडायचे आहे. $s = \Gn{\Gamma \Sequent \Delta}$
  असेल, तर $(s)_0 = \Gn{\Gamma}$ आणि $(s)_1 = \Gn{\Delta}$.
  म्हणून $\fn{FollowsBy}_{\RightR{\land}}(p)$ पुढील अट पूर्ण झाली
  तर आणि तरच खरे असते:
\begin{align*}
& \bexists{g < p}{\bexists{d < p}{\bexists{a < p}{\bexists{b < p}{\quad}}}} \\
& \qquad \fn{Sequent}(\tuple{g,d}) \land \fn{Sent}(a) \land \fn{Sent}(b) \land {}\\
& \qquad \fn{EndSequent}(p) =
  \tuple{g, d \concat
    \tuple{\Gn{(} \concat a \concat \Gn{\land} \concat b \concat \Gn{)}}} \land {} \\
& \qquad \fn{EndSequent}((p)_1) = \tuple{g, d \concat \tuple{a}} \land {}\\
& \qquad \fn{EndSequent}((p)_2) = \tuple{g, d \concat \tuple{b}} \land {}\\
& \qquad (p)_0 = 2 \land \fn{LastRule}(p) = 10.
\end{align*}
या ओळी अनुक्रमे पुढील गोष्टी व्यक्त करतात: $\Gamma$ ची गोडेल संख्या
$g$, $\Delta$ ची गोडेल संख्या $d$, सूत्र~$\varphi$ ची गोडेल संख्या $a$, आणि
सूत्र~$\psi$ ची गोडेल संख्या $b$ आहे; $\pi$ ची अंतिम क्रमवर्ती
$\Gamma \Sequent \Delta, \varphi \land \psi$ आहे; $\pi_1$ ची अंतिम क्रमवर्ती
$\Gamma \Sequent \Delta, \varphi$ आहे; $\pi_2$ ची अंतिम क्रमवर्ती
$\Gamma \Sequent \Delta, \psi$ आहे; आणि $\pi$ ला दोन तात्काळ
उपनिष्पत्ती असून तिचा अखेरचा नियम क्रमांक~$10$ असलेला
$\RightR\land$ आहे.

$\pi$ मधील अखेरचे निगमन $\RightR\lexists$ चा योग्य वापर असेल,
तर आणि तरच $\Gamma$ व $\Delta$ या क्रमिका, $\varphi$ हे सूत्र,
$x$ हे चल आणि $t$ हे बंद पद अशी असतात की $\pi$ ची अंतिम क्रमवर्ती
$\Gamma \Sequent \Delta, \lexists[x][\varphi]$ आणि $\pi_1$ ची अंतिम
क्रमवर्ती $\Gamma \Sequent \Delta, \Subst{\varphi}{t}{x}$ असते.
म्हणून गोडेल संख्यांत $\fn{FollowsBy}_{\RightR\lexists}(p)$ पुढील
अट पूर्ण झाली तर आणि तरच खरे असते:
\begin{align*}
  & \bexists{g<p}{\bexists{d<p}{\bexists{a<p}{\bexists{x<p}{\bexists{t<p}{\quad}}}}}\\
  & \qquad \fn{Sequent}(\tuple{g,d}) \land \fn{Frm}(a) \land \fn{Var}(x) \land {}\\
  & \qquad \fn{ClTerm}(t) \land {}\\
  & \qquad \fn{EndSequent}(p) =
  \tuple{
    g,
    d \concat \tuple{\Gn{\lexists} \concat x \concat a}
  } \land {}\\
  & \qquad \fn{EndSequent}((p)_1) =
  \tuple{
    g,
    d \concat \tuple{\fn{Subst}(a, t, x)}
  } \land {}\\
  & \qquad (p)_0 = 1 \land \fn{LastRule}(p) = 18.
\end{align*}
प्रस्तावित निष्पत्तीगाठीचा संकेतांक नेमक्या पुढीलपैकी एका रूपात असावा:
\begin{align*}
\fn{ProofNode}(p) \defiff {}& p=\tuple{0,\fn{EndSequent}(p)} \lor {}\\
& p=\tuple{1,(p)_1,\fn{EndSequent}(p),\fn{LastRule}(p)} \lor {}\\
& p=\tuple{2,(p)_1,(p)_2,\fn{EndSequent}(p),\fn{LastRule}(p)}.
\end{align*}
ही अट अनावश्यक अतिरिक्त घटक किंवा चुकीची गाठीची स्थानसंख्या वगळते.
मग $\fn{Correct}(p)$ ची व्याख्या अशी करतो:
\begin{multline*}
  \fn{ProofNode}(p) \land \fn{Sequent}(\fn{EndSequent}(p)) \land {}\\
  [(\fn{LastRule}(p) = 1 \land
  \fn{FollowsBy}_{\LeftR\Weakening}(p)) \lor \dots \lor {}\\
  (\fn{LastRule}(p) = 20 \land \fn{FollowsBy}_{\eq}(p)) \lor {}\\
  (p)_0 = 0 \land \fn{InitialSeq}(\fn{EndSequent}(p))]
\end{multline*}
पहिली ओळ गाठीच्या संकेतांकाचे नेमके रूप तपासते आणि $p$ ची अंतिम क्रमवर्ती खरोखर वाक्यांचीच
क्रमवर्ती आहे, याची खात्री करते. शेवटची ओळ $p$ मध्ये फक्त आरंभीची
क्रमवर्ती असण्याचा प्रसंग हाताळते.

\end{proof}

\begin{prob}
  \ref{inc:art:plk:prop:followsby} प्रमाणे पुढील गुणधर्म
  परिभाषित करा:
  \begin{enumerate}
  \item $\fn{FollowsBy}_{\Cut}(p)$,
  \item $\fn{FollowsBy}_{\LeftR\lif}(p)$,
  \item $\fn{FollowsBy}_{\eq}(p)$,
  \item $\fn{FollowsBy}_{\RightR{\lforall}}(p)$.
  \end{enumerate}
  शेवटच्या गुणधर्मासाठी, गोडेल संख्या~$p$ असलेल्या
  निष्पत्तीमधील अखेरचे निगमन आयगेन चराची अट पूर्ण करते
  का, हे आदिम पुनरावर्ती रीतीने तपासता येते, हेही दाखवावे लागेल.
  म्हणजेच $\RightR{\lforall}$ चा आयगेन चर~$a$ अंतिम
  क्रमवर्तीत येत नाही.
\end{prob}

\begin{prop}
  \label{inc:art:plk:prop:deriv}
  $p$ ही योग्य निष्पत्ती~$\pi$ ची गोडेल संख्या असेल, तर खरा
  असणारा $\fn{Deriv}(p)$ संबंध आदिम पुनरावर्ती आहे.
\end{prop}

\begin{proof}
  निष्पत्ती~$\pi$ मधील प्रत्येक निगमन एखाद्या नियमाचा योग्य
  वापर असेल, म्हणजेच तिच्या प्रत्येक उप-निष्पत्त्यांचा शेवट
  योग्य निगमनाने होत असेल, तर ती योग्य असते. म्हणून
  $\fn{Deriv}(p)$ तर आणि तरच
  \[  \bforall{i<\len{\fn{SubtreeSeq}(p)}}{\fn{Correct}((\fn{SubtreeSeq}(p))_i)}.
  \]
\end{proof}

\begin{prop}
समजा $\Gamma$ हा वाक्यांचा आदिम पुनरावर्ती संच आहे.
मग $\Prf[\Gamma](x, y)$ हा संबंध आदिम पुनरावर्ती आहे. हा संबंध
असे व्यक्त करतो: $x$ हा $\Gamma_0 \Sequent \varphi$ च्या
निष्पत्ती~$\pi$ चा संकेतांक आहे, येथे काही सांत
$\Gamma_0 \subseteq \Gamma$ आहे, आणि $y$ ही $\varphi$ ची गोडेल संख्या आहे.
\end{prop}

\begin{proof}
समजा ``$y \in \Gamma$'' हे $R_\Gamma(y)$ या आदिम पुनरावर्ती
विधेयाने दिले आहे. $\Prf[\Gamma](x, y)$ हा संबंध आदिम पुनरावर्ती आहे,
हे दाखवायचे आहे. $y$ ही वाक्य~$\varphi$ ची गोडेल संख्या असेल
आणि $x$ हा अंतिम क्रमवर्ती $\Gamma_0 \Sequent \varphi$ असलेल्या
$\Log{LK}$-निष्पत्तीचा संकेतांक असेल, तर आणि तरच हा संबंध खरा असतो.

मागील प्रस्तावानुसार, $x$ हा $\Log{LK}$ मधील योग्य
निष्पत्ती~$\pi$ चा संकेतांक असेल, तर आणि तरच खरा असणारा
$\fn{Deriv}(x)$ गुणधर्म आदिम पुनरावर्ती आहे. $x$ असा संकेतांक असेल,
तर $\fn{EndSequent}(x)$ हा $\pi$ च्या अंतिम क्रमवर्तीचा संकेतांक आहे.
म्हणून $(\fn{EndSequent}(x))_0$ हा तिच्या डाव्या बाजूचा आणि
$(\fn{EndSequent}(x))_1$ हा उजव्या बाजूचा संकेतांक आहे. त्यामुळे
``$\pi$ च्या अंतिम क्रमवर्तीची उजवी बाजू $\varphi$ आहे'' हे
$\len{(\fn{EndSequent}(x))_1} = 1 \land ((\fn{EndSequent}(x))_1)_0 =
y$ असे व्यक्त करता येते. $\pi$ च्या अंतिम क्रमवर्तीची डावी बाजू
अर्थातच सांत असते; तिच्यातील प्रत्येक वाक्य $\Gamma$ मध्ये आहे,
हे व्यक्त करणे उरते. म्हणून $\Prf[\Gamma](x, y)$ ची व्याख्या अशी करता येते:
\begin{align*}
\Prf[\Gamma](x, y) \defiff {}&
\fn{Deriv}(x) \land {} \\
& \bforall{i <
  \len{(\fn{EndSequent}(x))_0}}{R_\Gamma(((\fn{EndSequent}(x))_0)_i)} \land {}\\
& \len{(\fn{EndSequent}(x))_1} = 1 \land ((\fn{EndSequent}(x))_1)_0 = y.
\end{align*}
\end{proof}











\section{नैसर्गिक निगमनातील निष्पत्ती}\label{inc:art:pnd:sec}

\begin{explain}
निष्पत्त्यांचे अंकगणितीकरण करण्यासाठी निष्पत्त्या संख्यांनी
दर्शवावे लागतात. प्रत्येक निगमनाला एक किंवा दोन लेबले असलेल्या
सूत्रांचे वृक्ष म्हणजे निष्पत्त्या असल्यामुळे, पुनरावर्ती
निरूपण हा सहज मार्ग आहे. आपण निष्पत्ती एका क्रमित समूहाने दर्शवतो.
त्याचे घटक म्हणजे अखेरच्या निगमनाच्या आधारविधानांकडे नेणाऱ्या तात्काळ
उप-निष्पत्त्यांची संख्या, त्या उप-निष्पत्त्यांची निरूपणे,
अंतिम सूत्र, अखेरच्या निगमनाचे मुक्तता-लेबल आणि त्या निगमनाचा
प्रकार दर्शवणारी संख्या.
\end{explain}

\begin{defn}
  $\delta$ ही नैसर्गिक निगमनातील निष्पत्ती असेल, तर
  $\Gn{\delta}$ ची विगमनाने व्याख्या पुढीलप्रमाणे करतो:
  \begin{enumerate}
  \item
    $\delta$ मध्ये फक्त गृहीतक~$\varphi$ असेल, तर $\Gn{\delta}$ म्हणजे
    $\tuple{0, \Gn{\varphi}, n}$. हे गृहीतक मुक्त न केलेले असेल,
    तर~$n$ ही संख्या~$0$ असते; अन्यथा ती त्याचे संख्यात्मक लेबल असते.
  \item
    $\delta$ चा शेवट शून्य, एक, दोन किंवा तीन आधारविधानांच्या
    निगमनाने होत असेल, तर $\Gn{\delta}$ अनुक्रमे
    \begin{align*}
      & \tuple{0, \Gn{\varphi}, n, k}, \\
      & \tuple{1, \Gn{\delta_1}, \Gn{\varphi}, n, k}, \\
      & \tuple{2, \Gn{\delta_1}, \Gn{\delta_2}, \Gn{\varphi}, n, k}, \text{ किंवा}\\
      & \tuple{3, \Gn{\delta_1}, \Gn{\delta_2}, \Gn{\delta_3}, \Gn{\varphi}, n,
        k},
    \end{align*}
    असते. येथे $\delta_1$, $\delta_2$, $\delta_3$ या $\delta$ मधील
    अखेरच्या निगमनाच्या आधारविधानांवर संपणाऱ्या उप-निष्पत्त्या आहेत;
    $\varphi$ हा $\delta$ मधील त्या निगमनाचा निष्कर्ष आहे; $n$ हे त्या निगमनाचे
    मुक्तता-लेबल आहे (निगमनाने एकही गृहीतक मुक्त होत नसेल, तर~$0$);
    आणि अखेरच्या निगमनात वापरलेल्या नियमानुसार $k$ पुढील तक्त्यात
    दिले आहे.

    \begin{tabular}{lccccccc}
      \text{नियम:} & \Intro{\land} & \Elim{\land} & \Intro{\lor} & \Elim{\lor} \\
      $k$: & 1 & 2 & 3 & 4  \\[1.5ex]
      \text{नियम:} &  \Intro{\lif} & \Elim{\lif} & \Intro{\lnot} & \Elim{\lnot} \\
      $k$: & 5 & 6 & 7 & 8 \\[1.5ex]
      \text{नियम:}  &  \FalseInt & \FalseCl & \Intro{\lforall} & \Elim{\lforall} \\
      $k$: & 9 & 10 & 11 & 12 \\[1.5ex]
      \text{नियम:} & \Intro{\lexists} & \Elim{\lexists} & \Intro{\eq} & \Elim{\eq} \\
      $k$: & 13 & 14 & 15 & 16
    \end{tabular}
  \end{enumerate}
\end{defn}
\begin{ex}
  पुढील अगदी सोपी निष्पत्ती पाहा:
  \begin{prooftree}
    \AxiomC{$\Discharge{\varphi \land \psi}{1}$}
    \RightLabel{\Elim{\land}}
    \UnaryInfC{$\varphi$}
    \DischargeRule{\Intro{\lif}}{1}
    \UnaryInfC{$(\varphi \land \psi) \lif \varphi$}
  \end{prooftree}
  गृहीतकाची गोडेल संख्या $d_0 = \tuple{0,
    \Gn{(\varphi \land \psi)}, 1}$ असेल. $\Elim{\land}$ च्या निष्कर्षावर
  संपणाऱ्या निष्पत्तीची गोडेल संख्या $d_1 = \tuple{1,
     d_0, \Gn{\varphi}, 0, 2}$ असेल. येथे $1$ हे $\Elim{\land}$ ला एकच
  आधारविधान असल्यामुळे आले आहे; निष्कर्ष~$\varphi$ ची गोडेल संख्या पुढे
  आहे; $0$ हे एकही गृहीतक मुक्त होत नसल्यामुळे आले आहे; आणि $2$
  हा $\Elim{\land}$ चा संकेतांक आहे. मग संपूर्ण निष्पत्तीची
  गोडेल संख्या $\tuple{1, d_1, \Gn{((\varphi \land \psi) \lif
      \varphi)}, 1, 5}$, म्हणजेच
  \[
  \tuple{1, \tuple{1, \tuple{0, \Gn{(\varphi \land \psi)}, 1}, \Gn{\varphi}, 0, 2},
    \Gn{((\varphi \land \psi) \lif \varphi)}, 1, 5}.
  \]
\end{ex}

\begin{explain}
निष्पत्त्यांचे निरूपण ठरल्यानंतर, अशा निष्पत्त्यांच्या
गोडेल संख्यांवर आदिम पुनरावर्ती रीतीने क्रिया करता येतात आणि त्यांचे
आवश्यक गुणधर्म व संबंध व्यक्त करता येतात, हेही दाखवावे लागेल.
काही क्रिया सोप्या आहेत. उदाहरणार्थ, निष्पत्तीची गोडेल
संख्या~$d$ दिली असता, $\fn{EndFmla}(d) = (d)_{(d)_0+1}$ हे
तिच्या अंतिम सूत्राची गोडेल संख्या देते;
$\fn{DischargeLabel}(d) = (d)_{(d)_0 + 2}$ हे मुक्तता-लेबल देते;
आणि $\fn{LastRule}(d) = (d)_{(d)_0 + 3}$ हे अखेरच्या निगमनाचा
प्रकार दर्शवणारी संख्या देते. काही इतर क्रिया अधिक कठीण आहेत;
त्या कशा कराव्यात याची किमान रूपरेषा आपण पाहू. $\delta$ ही
$\Gamma$ पासून $\varphi$ ची निष्पत्ती आहे, हा संबंध
$\delta$ आणि $\varphi$ यांच्या गोडेल संख्यांचा आदिम पुनरावर्ती संबंध
आहे, हे दाखवणे आपले उद्दिष्ट आहे.
\end{explain}

\begin{prop}
  पुढील संबंध आदिम पुनरावर्ती आहेत:
  \begin{enumerate}
  \item $\varphi$ हे $\delta$ मध्ये लेबल~$n$ असलेले गृहीतक म्हणून येते.
  \item $\delta$ मधील लेबल~$n$ असलेली सर्व गृहीतके $\varphi$ या रूपाची
    असतात (म्हणजे $\delta$ मध्ये लेबल~$n$ वापरून गृहीतक~$\varphi$
    मुक्त करता येते).
  \end{enumerate}
\end{prop}

\begin{proof}
  सूत्रांच्या गोडेल संख्या आणि निष्पत्त्यांच्या गोडेल
  संख्या यांच्यातील अनुरूप संबंध आदिम पुनरावर्ती आहेत, हे दाखवायचे आहे.
  \begin{enumerate}
  \item $\fn{Assum}(x, d, n)$ आदिम पुनरावर्ती आहे, हे दाखवू.
    गोडेल संख्या~$d$ असलेल्या निष्पत्तीमधील लेबल~$n$ च्या
    गृहीतकाची गोडेल संख्या~$x$ असेल, तर हा संबंध खरा असतो.
    गोडेल संख्या~$\tuple{0, x, n}$ असलेली निष्पत्ती ही
    $d$ ची उप-निष्पत्ती असेल, तर आणि तरच असे घडते. येथे
    निष्पत्त्यांचे आपले सांकेतीकरण हे
    \ref{cmp:rec:tre:sec} मध्ये दिलेल्या वृक्षांच्या
    सांकेतीकरणाचे विशेष प्रकरण आहे, हे लक्षात घ्या. म्हणून आदिम
    पुनरावर्ती फलन $\fn{SubtreeSeq}(d)$ हे $d$ च्या सर्व
    उप-निष्पत्त्यांच्या गोडेल संख्यांची क्रमिका देते
    ($d$ च्या या क्रमिकेत पुनरुक्तीही असू शकतात). म्हणून व्याख्या अशी करता येते:
    \[    \fn{Assum}(x, d, n) \defiff \bexists{i<\len{\fn{SubtreeSeq}(d)}}{(\fn{SubtreeSeq}(d))_i =
      \tuple{0, x, n}}.
    \]
  \item $\fn{Discharge}(x, d, n)$ आदिम पुनरावर्ती आहे, हे दाखवू.
    गोडेल संख्या~$d$ असलेल्या निष्पत्तीमधील लेबल~$n$
    असलेली सर्व गृहीतके गोडेल संख्या~$x$ असलेल्या सूत्र
    चीच असतील, तर हा संबंध खरा असतो. तो
    $\bforall{y<d}{(\fn{Assum}(y, d, n) \lif y
      = x)}$ असेल, तर आणि तरच खरा असतो.
  \end{enumerate}
\end{proof}

\begin{prop}
  \label{inc:art:pnd:prop:followsby}
  गोडेल संख्या~$d$ असलेल्या निष्पत्ती~$\delta$ मधील अखेरचे
  निगमन योग्य असेल, तर आणि तरच खरा असणारा $\fn{Correct}(d)$ गुणधर्म
  आदिम पुनरावर्ती आहे.
\end{prop}

\begin{proof}
  प्रत्येक निगमन नियम~$R$ साठी $\fn{FollowsBy}_R(d)$ हा संबंध
  आदिम पुनरावर्ती आहे, हे दाखवावे लागेल. $d$ ही निष्पत्ती~$\delta$
  ची गोडेल संख्या असेल आणि $\delta$ चे अंतिम सूत्र $\delta$ च्या
  तात्काळ उप-निष्पत्त्या पासून $R$ च्या योग्य वापराने
  निष्पन्न होत असेल, तर आणि तरच $\fn{FollowsBy}_R(d)$ खरे असते.

  \Intro{\land} नियमाचे उदाहरण सोपे आहे. $\delta$ चा शेवट
  $\Intro{\land}$ च्या योग्य वापराने होत असेल, तर तिचे रूप असे असते:
  \begin{prooftree}
    \AxiomC{}
    \RightLabel{$\delta_1$}
    \DeduceC{$\varphi$}

    \AxiomC{}
    \RightLabel{$\delta_2$}
    \DeduceC{$\psi$}

    \RightLabel{\Intro\land}
    \BinaryInfC{$\varphi \land \psi$}
  \end{prooftree}
  मग $\delta$ ची गोडेल संख्या~$d$ ही
  $\tuple{2, d_1, d_2, \Gn{(\varphi \land \psi)}, 0, k}$ असते. येथे
  $\fn{EndFmla}(d_1) = \Gn{\varphi}$,
  $\fn{EndFmla}(d_2) = \Gn{\psi}$, $n=0$ आणि $k=1$.
  म्हणून $\fn{FollowsBy}_{\Intro\land}(d)$ ची व्याख्या अशी करता येते:
  \begin{multline*}
    (d)_0 = 2 \land \fn{DischargeLabel}(d) = 0 \land \fn{LastRule}(d) = 1 \land {}\\
    \fn{EndFmla}(d) = {}\\
    \Gn{(} \concat \fn{EndFmla}((d)_1) \concat \Gn{\land}
    \concat \fn{EndFmla}((d)_2) \concat \Gn{)}.
  \end{multline*}

  आणखी एक सोपे उदाहरण $\Intro\eq$ नियमाचे आहे. त्याला
  आधारविधाने नसतात, म्हणून गृहीतकांप्रमाणेच $(d)_0 = 0$ असते.
  त्याचे मुक्तता-लेबलही नसते, म्हणजे $n=0$. मात्र $\varphi$ हे
  एखाद्या बंद पद~$t$ साठी $\eq[t][t]$ या रूपाचे असले पाहिजे.
  येथे आदिम पुनरावर्ती व्याख्या अशी आहे:
  \begin{multline*}
    (d)_0 = 0 \land \fn{DischargeLabel}(d) = 0 \land {}\\
    \bexists{t<d}{(\fn{ClTerm}(t) \land
      \fn{EndFmla}(d) = {}\\
      \Gn{{\eq}(} \concat t \concat \Gn{,} \concat t \concat \Gn{)})}.
  \end{multline*}

  अधिक गुंतागुंतीचे उदाहरण म्हणून $\fn{FollowsBy}_{\Intro{\lif}}(d)$
  पाहू. $\delta$ चे अंतिम सूत्र $(\varphi \lif \psi)$ या रूपाचे,
  $\delta_1$ चे अंतिम सूत्र~$\psi$ असेल आणि $\delta$ मधील
  लेबल~$n$ असलेले प्रत्येक गृहीतक $\varphi$ या रूपाचे असेल,
  तर आणि तरच हा संबंध खरा असतो. हे आदिम पुनरावर्ती रीतीने
  पुढीलप्रमाणे व्यक्त करता येते:
  \begin{multline*}
    (d)_0 = 1 \land {}\\
    \bexists{a<d}{(\fn{Discharge}(a, (d)_1, \fn{DischargeLabel}(d)) \land {}}\\
      \fn{EndFmla}(d) = (\Gn{(} \concat a \concat \Gn{\lif}
      \concat \fn{EndFmla}((d)_1) \concat \Gn{)}))
  \end{multline*}
  (येथे $a$ ही $\varphi$ ची गोडेल संख्या समजा.)

  आणखी एका उदाहरणासाठी \Intro{\lexists} पाहू. $\delta$ मधील
  अखेरचे निगमन योग्य असेल, तर आणि तरच सूत्र~$\varphi$,
  बंद पद~$t$ आणि चल~$x$ अशी असतात की
  $\Subst{\varphi}{t}{x}$ हे निष्पत्ती~$\delta_1$ चे अंतिम
  सूत्र असते आणि $\lexists[x][\varphi]$ हा अखेरच्या निगमनाचा
  निष्कर्ष असतो. म्हणून $\fn{FollowsBy}_{\Intro{\lexists}}(d)$
  पुढील अट पूर्ण झाली तर आणि तरच खरे असते:
  \begin{multline*}
    (d)_0 = 1 \land \fn{DischargeLabel}(d) = 0 \land {} \\
    \bexists{a < d}{\bexists{x<d}{\bexists{t<d}{
          (\fn{ClTerm}(t) \land \fn{Var}(x)  \land {}}}}\\
    \fn{Subst}(a,t,x) = \fn{EndFmla}((d)_1) \land
    \fn{EndFmla}(d) = (\Gn{\lexists} \concat x \concat a)).
  \end{multline*}

नेमका गाठीचा आकार तपासण्यासाठी मागील विभागातील वैध क्रमिकासंकेतांकाची
अट वापरा. गृहीतकाला तीन घटक असतात; नियमगाठीला आधारविधानांच्या
संख्येपेक्षा चार अधिक घटक असतात:
\[\fn{NDNode}(d) \defiff \fn{Seq}(d) \land
[((d)_0=0 \land \len{d}=3) \lor ((d)_0\le3 \land \len{d}=(d)_0+4)].
\]
  मग $\fn{Correct}(d)$ ची व्याख्या अशी करतो:
  \begin{multline*}
    \fn{NDNode}(d) \land \fn{Sent}(\fn{EndFmla}(d)) \land {}\\
    [(\fn{LastRule}(d) = 1 \land
    \fn{FollowsBy}_{\Intro\land}(d)) \lor \dots \lor {}\\
    (\fn{LastRule}(d) = 16 \land \fn{FollowsBy}_{\Elim\eq}(d)) \lor {}\\
    \bexists{n<d}{\bexists{x<d}{(d = \tuple{0, x, n})}}].
  \end{multline*}
  पहिली ओळ संकेतांकाचा नेमका गाठीचा आकार आणि $d$ चे अंतिम सूत्र वाक्य आहे, याची खात्री
  करते. शेवटची ओळ $d$ मध्ये फक्त गृहीतक असण्याचा प्रसंग हाताळते.
\end{proof}

\begin{prob}
  \ref{inc:art:pnd:prop:followsby} प्रमाणे पुढील गुणधर्म
  परिभाषित करा:
  \begin{enumerate}
  \item $\fn{FollowsBy}_{\Elim{\lif}}(d)$,
  \item $\fn{FollowsBy}_{\Elim{\eq}}(d)$,
  \item $\fn{FollowsBy}_{\Elim{\lor}}(d)$,
  \item $\fn{FollowsBy}_{\Intro{\lforall}}(d)$.
  \end{enumerate}
  शेवटच्या गुणधर्मासाठी, गोडेल संख्या~$d$ असलेल्या
  निष्पत्तीमधील अखेरचे निगमन आयगेन चराची अट पूर्ण करते
  का, हे आदिम पुनरावर्ती रीतीने तपासता येते, हेही दाखवा.
  म्हणजेच $\Intro{\lforall}$ निगमनाचा आयगेन चर~$a$ हा
  $d$ च्या अंतिम सूत्रामध्ये किंवा $d$ च्या मुक्त न केलेल्या
  गृहीतकात येत नाही. यासाठी
  \ref{inc:art:pnd:prop:openassum} मधील आदिम पुनरावर्ती
  विधेय $\fn{OpenAssum}$ वापरू शकता.
\end{prob}

\begin{prop}
  \label{inc:art:pnd:prop:deriv}
  $d$ ही योग्य निष्पत्ती~$\delta$ ची गोडेल संख्या असेल,
  तर खरा असणारा $\fn{Deriv}(d)$ संबंध आदिम पुनरावर्ती आहे.
\end{prop}

\begin{proof}
  निष्पत्ती~$\delta$ मधील प्रत्येक निगमन एखाद्या नियमाचा
  योग्य वापर असेल, म्हणजेच तिच्या प्रत्येक उप-निष्पत्त्यांचा
  शेवट योग्य निगमनाने होत असेल, तर ती योग्य असते. म्हणून
  $\fn{Deriv}(d)$ तर आणि तरच
  \[  \bforall{i<\len{\fn{SubtreeSeq}(d)}}{\fn{Correct}((\fn{SubtreeSeq}(d))_i)}
  \]
\end{proof}

\begin{prop}
  \label{inc:art:pnd:prop:openassum} $z$ ही गोडेल संख्या
  मुक्त न केलेले गृहीतक~$\varphi$ ची असेल, आणि हे गृहीतक
  गोडेल संख्या~$d$ असलेल्या निष्पत्ती~$\delta$ मध्ये असेल,
  तर खरा असणारा $\fn{OpenAssum}(z, d)$ संबंध आदिम पुनरावर्ती आहे.
\end{prop}

\begin{proof}
  गृहीतकाच्या एखाद्या आस्थितीला शून्येतर लेबल~$n$ असेल आणि ती
  $\delta$ च्या अशा उप-निष्पत्तीमध्ये येत असेल जिचा शेवट
  मुक्तता-लेबल~$n$ असलेल्या नियमाने होतो, तर ती आस्थिती
  मुक्त केलेले असते. म्हणून अशा गृहीतकाची किमान एक आस्थिती
  $\delta$ मध्ये मुक्त केलेले नसेल, तर $\varphi$ हे
  $\delta$ चे मुक्त न केलेले गृहीतक असते. येथे सावध राहिले
  पाहिजे: $\delta$ मध्ये $\varphi$ च्या मुक्त केलेले आणि
  मुक्त न केलेले अशा दोन्ही आस्थिती असू शकतात.

  $\delta_0$, \dots, $\delta_k$ ही क्रमिका पाहा. येथे $\delta_0 =
  \delta$, $\delta_k$ हे गृहीतक $\Discharge{\varphi}{n}$ आहे
  (एखाद्या~$n$ साठी), आणि $\delta_{i+1}$ ही $\delta_i$ ची तात्काळ
  उप-निष्पत्ती आहे. लेबल शून्य असेल, तर हे गृहीतक व्याख्येनेच
  मुक्त न केलेले असते. अन्यथा, अशा क्रमिकेत गृहीतकाच्या आधीच्या कोणत्याही
  $\delta_i$ चा शेवट मुक्तता-लेबल~$n$ असलेल्या निगमनाने होत नसेल,
  तर $\varphi$ हे $\delta$ चे मुक्त न केलेले गृहीतक आहे.

  आदिम पुनरावर्ती फलन $\fn{SubtreeSeq}(d)$ हे $\delta$ च्या
  सर्व उप-निष्पत्त्यांच्या गोडेल संख्यांची क्रमिका देते.
  $\delta$ च्या मूळापासून गृहीतकापर्यंतचा उप-निष्पत्त्यांचा मार्ग त्यातून
  निवडता येतो. अशी उपक्रमिका ओळखणे हा आदिम पुनरावर्ती संबंध आहे:
  $\fn{Subseq}(s, s')$ तर आणि तरच
  $\bforall{i<\len{s}}{\lexists[j<\len{s'}][(s)_i = (s')_j]}$.
  येथे दिलेली अट प्रत्येक घटकाचे सदस्यत्व तपासते; ती स्वतःहून
  क्रम-संरक्षण व्यक्त करत नाही. तात्काळ उप-निष्पत्ती असण्याचा
  संबंधही आदिम पुनरावर्ती आहे: $\fn{Subderiv}(d, d')$ तर आणि तरच
  $\bexists{j<(d')_0}{d = (d')_{j+1}}$. म्हणून
  $\fn{OpenAssum}(z, d)$ ची व्याख्या अशी करता येते:
  \begin{multline*}
    \bexists{s<1+\fn{sequenceBound}(d,d)}{(\fn{Seq}(s) \land \len{s}>0 \land {}\\
      \fn{Subseq}(s, \fn{SubtreeSeq}(d))
      \land (s)_0 = d \land {}} \\
    \bexists{n<d}{((s)_{\len{s} \tsub 1} = \tuple{0, z, n} \land {}}\\
      \bforall{i<(\len{s} \tsub 1)}{(\fn{Subderiv}((s)_{i+1}, (s)_i) \land {}}\\
      (n=0 \lor \fn{DischargeLabel}((s)_i) \neq n)))).
  \end{multline*}
मार्गात प्रत्येक पुढचा गाठीचा संकेतांक आधीच्यापेक्षा लहान असतो.
म्हणून वैध मार्गात जास्तीत जास्त $d$ गाठी आणि प्रत्येक संकेतांक
जास्तीत जास्त $d$ असतो; वरची आदिम पुनरावर्ती मर्यादा पुरेशी आहे.

\end{proof}

\begin{prop}
  \label{inc:art:pnd:prop:prf-prim-rec}
  समजा $\Gamma$ हा वाक्यांचा आदिम पुनरावर्ती संच आहे.
  मग $\Prf[\Gamma](x, y)$ हा संबंध आदिम पुनरावर्ती आहे. तो असे व्यक्त
  करतो: $x$ हा $\Gamma$ मधील मुक्त न केलेले गृहीतकांपासून
  $\varphi$ च्या निष्पत्ती~$\delta$ चा संकेतांक आहे, आणि
  $y$ ही $\varphi$ ची गोडेल संख्या आहे.
\end{prop}

\begin{proof}
  समजा ``$y \in \Gamma$'' हे $R_\Gamma(y)$ या आदिम पुनरावर्ती
  विधेयाने दिले आहे. $\Prf[\Gamma](x, y)$ आदिम पुनरावर्ती आहे,
  हे दाखवायचे आहे. $y$ ही वाक्य~$\varphi$ ची गोडेल संख्या असेल,
  आणि $x$ हा अंतिम सूत्र~$\varphi$ असलेल्या नैसर्गिक निगमनातील
  निष्पत्तीचा संकेतांक असेल, तसेच तिची सर्व
  मुक्त न केलेले गृहीतके $\Gamma$ मध्ये असतील, तर आणि तरच हा
  संबंध खरा असतो.

  \ref{inc:art:pnd:prop:deriv} नुसार, $x$ ही नैसर्गिक निगमनातील योग्य
  निष्पत्ती~$\delta$ ची गोडेल संख्या असेल, तर आणि तरच खरा
  असणारा $\fn{Deriv}(x)$ गुणधर्म आदिम पुनरावर्ती आहे. त्यामुळे
  $\Prf[\Gamma](x, y)$ ची व्याख्या अशी करता येते:
  \begin{align*}
    \Prf[\Gamma](x, y) \defiff {}
    & \fn{Deriv}(x) \land \fn{EndFmla}(x) = y \land {} \\
    & \bforall{z < x}{(\fn{OpenAssum}(z, x) \lif R_\Gamma(z))}.
  \end{align*}
\end{proof}











\section{स्वयंसिद्धकीय निष्पत्ती}\label{inc:art:pax:sec}

\begin{explain}
स्वयंसिद्धकीय निष्पत्त्यांचे अंकगणितीकरण करण्यासाठी
निष्पत्त्या संख्यांनी दर्शवावे लागतात. निष्पत्त्या या
सूत्रांच्या क्रमिका असल्यामुळे, प्रत्येक निष्पत्ती
तिच्यातील सूत्रांच्या संकेतांकांच्या क्रमिकेचा संकेतांक
म्हणून दर्शवणे हा सहज मार्ग आहे.
\end{explain}

\begin{defn}
$\delta$ ही $\varphi_1$, \dots,~$\varphi_n$ या सूत्रांची
स्वयंसिद्धकीय निष्पत्ती असेल, तर $\Gn{\delta}$ म्हणजे
\[\tuple{\Gn{\varphi_1}, \dots, \Gn{\varphi_n}}.
\]
\end{defn}

\begin{ex}
  पुढील अगदी सोपी निष्पत्ती पाहा:
  \begin{derivation}
  1. & $\psi \lif (\psi \lor \varphi)$ \\
  2. & $(\psi \lif (\psi \lor \varphi)) \lif (\varphi  \lif (\psi \lif (\psi \lor \varphi)))$\\
  3. & $\varphi  \lif (\psi \lif (\psi \lor \varphi))$
  \end{derivation}
  या निष्पत्तीची गोडेल संख्या
  \begin{align*}
  \openTuple\,
    & \Gn{\psi \lif (\psi \lor \varphi)}, \\
    &\Gn{(\psi \lif (\psi \lor \varphi)) \lif (\varphi  \lif (\psi \lif (\psi \lor \varphi)))},\\
    & \Gn{\varphi  \lif (\psi \lif (\psi \lor \varphi))} \,\closeTuple.
  \end{align*}
  अशी असेल.
\end{ex}

\begin{explain}
निष्पत्त्यांचे निरूपण ठरल्यानंतर, अशा निष्पत्त्यांवर
आदिम पुनरावर्ती रीतीने क्रिया करता येतात आणि त्यांचे आवश्यक
गुणधर्म व संबंध तसेच व्यक्त करता येतात, हेही दाखवावे लागेल.
काही क्रिया सोप्या आहेत. उदाहरणार्थ, निष्पत्तीची गोडेल
संख्या~$d$ दिली असता, $(d)_{\len{d}-1}$ हे तिच्या अंतिम
सूत्राची गोडेल संख्या देते. काही इतर क्रिया अधिक कठीण
आहेत; त्या कशा कराव्यात याची किमान रूपरेषा आपण पाहू.
``$\delta$ ही $\Gamma$ पासून $\varphi$ ची निष्पत्ती आहे''
हा संबंध $\delta$ आणि $\varphi$ यांच्या गोडेल संख्यांचा आदिम
पुनरावर्ती संबंध आहे, हे दाखवणे आपले उद्दिष्ट आहे.
\end{explain}

\begin{prop}
\label{inc:art:pax:prop:followsby}
पुढील संबंध आदिम पुनरावर्ती आहेत:
\begin{enumerate}
\item $\varphi$ हे स्वयंसिद्धक आहे.
\item $\delta$ मधील $i$ वी ओळ विध्यात्मक अनुमानाने समर्थित आहे.
\item $\delta$ मधील $i$ वी ओळ \QR{} ने समर्थित आहे.
\item $\delta$ ही योग्य निष्पत्ती आहे.
\end{enumerate}
\end{prop}

\begin{proof}
सूत्रांच्या गोडेल संख्या आणि निष्पत्त्यांच्या
गोडेल संख्या यांच्यातील अनुरूप संबंध आदिम पुनरावर्ती आहेत,
हे दाखवायचे आहे.
\begin{enumerate}
\item स्वयंसिद्धक-रूपबंधांची यादी आपल्याला दिली आहे; $\varphi$ हे
  स्वयंसिद्धक असेल, तर ते या रूपबंधांपैकी एखाद्याच्या रूपाचे असते.
  यादी सांत असल्यामुळे, प्रत्येक रूपबंधासाठी $\varphi$ त्या रूपाचे
  आहे का हे आदिम पुनरावर्ती रीतीने तपासता येते, एवढे दाखवणे पुरेसे आहे.
  उदाहरणार्थ पुढील स्वयंसिद्धक-रूपबंध पाहा:
  \[  \psi \lif (\chi \lif \psi).
  \]
  $\psi$ आणि $\chi$ ही सूत्रे अशी असतील की `$($', $\psi$,
  `$\lif$', `$($', $\chi$, `$\lif$', $\psi$ आणि~`$))$' या
  चिन्हमाला जोडून $\varphi$ मिळते, तर $\varphi$ हा या रूपबंधाचा नमुना आहे.
  $\varphi$ ची गोडेल संख्या $n$ असल्यास हा गुणधर्म तपासता येतो:
  क्रमिकांची जोडणी आदिम पुनरावर्ती आहे आणि हा संबंध खरा असेल,
  तर $\psi$ व $\chi$ ही $\varphi$ ची उप-सूत्रे असल्यामुळे $\psi$ व $\chi$
  यांच्या गोडेल संख्या $\varphi$ च्या गोडेल संख्येपेक्षा लहान असतात.
  म्हणून व्याख्या अशी करता येते:
  \begin{multline*}
  \fn{IsAx}_{\psi \lif (\chi \lif \psi)}(n) \defiff \bexists{b< n}{
    \bexists{c < n}{(\fn{Frm}(b) \land \fn{Frm}(c) \land {}}}\\ n =
  \Gn{(} \concat b \concat \Gn{\lif} \concat \Gn{(} \concat c \concat
  \Gn{\lif} \concat b \concat \Gn{))}).
  \end{multline*}
  प्रत्येक स्वयंसिद्धक-रूपबंधासाठी अशी व्याख्या केली, तर
  त्यांच्या विकल्पाने $\fn{IsAx}(n)$ हा गुणधर्म परिभाषित होतो:
  ``$n$ ही स्वयंसिद्धकाची गोडेल संख्या आहे.''
\item $\delta$ मधील $i$ वी ओळ विध्यात्मक अनुमानाने समर्थित असेल,
  तर आणि तरच $j$ आणि $k < i$ अशा आधीच्या ओळी असतात की
  $j$ व्या ओळीवरील सूत्र~$\varphi$ असते,
  $k$ व्या ओळीवरील सूत्र $\varphi \lif \psi$ असते, आणि
  $i$ व्या ओळीवरील सूत्र~$\psi$ असते.
  \begin{multline*}
    \fn{MP}(d, i) \defiff \bexists{j < i}{\bexists{k < i}{}}\\
    (d)_k = \Gn{(} \concat (d)_j
      \concat \Gn{\lif} \concat (d)_i \concat \Gn{)}
  \end{multline*}
  परिबद्ध संख्यकीकरण, जोडणी आणि $=$ हे आदिम पुनरावर्ती
  असल्यामुळे हा संबंधही आदिम पुनरावर्ती आहे.
\item $\delta$ मधील एखादी ओळ \QR{} ने समर्थित असेल,
  तर ती $\psi \lif \lforall[x][\varphi(x)]$ या रूपाची असते;
  आधीची एखादी ओळ, कोणत्यातरी स्थिरांक~$c$ साठी,
  $\psi \lif \varphi(c)$ असते; आणि $c$ हा $\psi$ मध्ये येत नाही.
  पुढील अटी पूर्ण झाल्या तर आणि तरच असे घडते:
  \begin{enumerate}
  \item एखादे वाक्य~$\psi$ असते (येथे हा विशेष प्रसंग;
    सर्वसाधारण नियमात सूत्रही चालते), आणि
  \item $x$ हाच मुक्त चल असू शकणारे सूत्र~$\varphi(x)$ असते, ज्यासाठी
  \item ओळ $i$ मध्ये $\psi \lif \lforall[x][\varphi(x)]$ असते,
  \item आधीच्या एखाद्या ओळीत, $j < i$, कोणत्यातरी अचर~$c$ साठी
    $\psi \lif \Subst{\varphi}{c}{x}$ असते,
  \item आणि हा अचर $\psi$ मध्ये येत नाही.
  \end{enumerate}
  या अटी आदिम पुनरावर्ती रीतीने तपासता येतात. कारण $\psi$,
  $\varphi(x)$ आणि $x$ यांच्या गोडेल संख्या $i$ व्या ओळीवरील
  सूत्राच्या गोडेल संख्येपेक्षा लहान आहेत; तसेच $c$ ची गोडेल
  संख्या $j$ व्या ओळीवरील सूत्राच्या गोडेल संख्येपेक्षा लहान आहे:
  \begin{multline*}
    \fn{QR}_1(d, i) \defiff \bexists{j<i}{\bexists{b < (d)_i}{\bexists{x <
        (d)_i}{\bexists{a < (d)_i}{\bexists{c < (d)_j}{(}}}}}
    \\ \fn{Var}(x) \land \fn{Const}(c) \land {} \\
    (d)_i = \Gn{(} \concat
    b \concat \Gn{\lif} \concat \Gn{\lforall} \concat x \concat a
    \concat \Gn{)} \land {}\\
    (d)_j = \Gn{(} \concat b \concat
    \Gn{\lif} \concat \fn{Subst}(a,c,x) \concat \Gn{)} \land {}\\ \fn{Frm}(b)
    \land \fn{Frm}(a) \land {} \bforall{k <
      \len{b}}{(b)_k \neq (c)_0})
  \end{multline*}
  येथे $c$ आणि $x$ यांना अचर आणि चल यांच्या पदरूपातील
  गोडेल संख्या मानले आहे (म्हणजे त्यांचे चिन्हसंकेतांक नाहीत).
  $x$ हा $\varphi(x)$ मधील आदेशनासाठीचा चल आहे;
  $\Subst{\varphi(x)}{c}{x}$ हे वाक्य आहे का ही मूळ
  वाक्यांपुरती तपासणी सर्वसाधारण संख्यक-नियमासाठी आवश्यक नाही.
  सूत्ररूपता तपासून, $\psi$ मधील प्रत्येक चिन्ह $c$ पेक्षा
  वेगळे आहे ही अट ठेवून $c$ हा $\psi$ मध्ये येत नाही,
  याची खात्री करतो.

  \QR{} ची दुसरी आवृत्ती सरावासाठी ठेवतो.
\item $d$ ही योग्य निष्पत्तीची गोडेल संख्या असेल, तर आणि
  तरच ती रिकामी नसते आणि तिच्यातील प्रत्येक ओळ स्वयंसिद्धक
  किंवा विध्यात्मक अनुमान अथवा \QR{} ने समर्थित असते. म्हणून:
  \[  \fn{Deriv}(d) \defiff \fn{Seq}(d) \land \len{d}>0 \land
  \bforall{i < \len{d}}{(\fn{IsAx}((d)_i) \lor
    \fn{MP}(d,i) \lor \fn{QR}(d, i))}
  \]
\end{enumerate}
\end{proof}

\begin{prob}
\ref{inc:art:pax:prop:followsby} प्रमाणे पुढील संबंध
परिभाषित करा:
\begin{enumerate}
\item $\fn{IsAx}_{\varphi \lif (\psi \lif (\varphi \land \psi))}(n)$,
\item $\fn{IsAx}_{\lforall[x][\varphi(x)] \lif \varphi(t)}(n)$,
\item $\fn{QR}_{2}(d, i)$ (\QR{} च्या दुसऱ्या आवृत्तीसाठी).
\end{enumerate}
\end{prob}

\begin{prop}
समजा $\Gamma$ हा वाक्यांचा आदिम पुनरावर्ती संच आहे.
मग $\Prf[\Gamma](x, y)$ हा संबंध आदिम पुनरावर्ती आहे.
तो असे व्यक्त करतो: $x$ हा $\Gamma$ पासून $\varphi$ च्या
निष्पत्ती~$\delta$ चा संकेतांक आहे आणि $y$ ही
$\varphi$ ची गोडेल संख्या आहे.
\end{prop}

\begin{proof}
समजा ``$y \in \Gamma$'' हे $R_\Gamma(y)$ या आदिम पुनरावर्ती
विधेयाने दिले आहे. $\Prf[\Gamma](x, y)$ हा संबंध आदिम पुनरावर्ती
आहे, हे दाखवायचे आहे. $y$ ही वाक्य~$\varphi$ ची गोडेल संख्या
असेल आणि $x$ हा $\Gamma$ पासून $\varphi$ च्या
निष्पत्तीचा संकेतांक असेल, तर आणि तरच
$\Prf[\Gamma](x, y)$ खरा असतो.

मागील प्रस्तावानुसार, $x$ हा योग्य निष्पत्ती~$\delta$
चा संकेतांक असेल, तर आणि तरच खरा असणारा $\fn{Deriv}(x)$
गुणधर्म आदिम पुनरावर्ती आहे. मात्र त्या व्याख्येत
निष्पत्तीमधील ओळी समर्थित करण्याचा अतिरिक्त मार्ग म्हणून
$\Gamma$ संच विचारात घेतलेला नाही. ओळ \QR{} ने समर्थित आहे
का हे तपासणाऱ्या आदिम पुनरावर्ती कसोटीतही अचर~$c$ हा
$\Gamma$ मध्ये येऊ नये ही अट घेतलेली नाही. $\Gamma$ चा
थेट विचार करून व्याख्या सुधारता येईल; पण $\fn{Deriv}$ आणि
निगमन प्रमेय वापरणे सोपे आहे. काही सांत वाक्यांची यादी
$\psi_1$, \dots, $\psi_n \in \Gamma$ अशी असेल की
$\{\psi_1, \dots, \psi_n\} \Proves \varphi$, तर आणि तरच
$\Gamma \Proves \varphi$. निगमन प्रमेयानुसार, हे
$\Proves (\psi_1 \lif (\psi_2 \lif
\cdots (\psi_n \lif \varphi)\cdots))$ असेल तेव्हा घडते.
गोडेल संख्या~$z$ असलेले वाक्य या रूपाचे आहे का,
हे आदिम पुनरावर्ती रीतीने तपासता येते. म्हणून $x$ ला
गोडेल संख्या~$y$ असलेल्या वाक्याची \emph{$\Gamma$ पासून}
निष्पत्ती असण्याचा थेट संकेतांक मानण्याऐवजी,
$x$ हा वरच्या रूपातील आवर्ती सोपाधिक वाक्याच्या
$\emptyset$ पासूनच्या निष्पत्तीचा संकेतांक मानतो.
यामुळे खालील कसोटी सिद्धतायोग्यतेसाठी समतुल्य असली,
तरी संकेतांकाचा अर्थ सुरुवातीच्या विधानातील थेट संकेतांकापेक्षा बदलतो.

प्रथम, वाक्यांची क्रमिका दिल्यास तिच्यातील सर्व वाक्ये
पूर्वांग म्हणून आणि दिलेले वाक्य उत्तरांग म्हणून असलेले
सोपाधिक वाक्य आदिम पुनरावर्ती रीतीने तयार करता येते:
\begin{align*}
  \fn{hCond}(s, y, 0) & = y \\
  \fn{hCond}(s, y, n+1) & = \Gn{(} \concat (s)_{n} \concat \Gn{\lif}
  \concat \fn{hCond}(s, y, n) \concat \Gn{)}\\
  \fn{Cond}(s, y) & = \fn{hCond}(s, y, \len{s})\\
  \intertext{म्हणून $\Prf[\Gamma](x, y)$ ची व्याख्या अशी करता येते:}
  \Prf[\Gamma](x, y) & \defiff \fn{Sent}(y) \land
  \bexists{s < \fn{sequenceBound}(x,x)}{(\fn{Seq}(s) \land {}} \\
  &\qquad (x)_{\len{x}-1} = \fn{Cond}(s,y) \land {} \\
  &\qquad\bforall{i<\len{s}}{(s)_i \in \Gamma} \land {} \\
  &\qquad\fn{Deriv}(x)).
\end{align*}

येथे पुनरावर्तनाने क्रमिकेतील पूर्वांगे आतून बाहेर या क्रमाने
जोडली जातात; आधी लिहिलेल्या बाहेरून आतच्या यादीसाठी ती उलट क्रमाने
निवडता येतात. $s$ वरील परिबंध असा मिळतो: प्रत्येक $(s)_i$ ही
निष्पत्तीच्या अखेरच्या ओळीतील उप-सूत्राची गोडेल संख्या
आहे; म्हणजेच ती $(x)_{\len{x}-1}$ पेक्षा लहान आहे. $\psi
\in \Gamma$ या पूर्वांगांची संख्या, म्हणजे $s$ ची लांबी,
$x$ च्या अखेरच्या ओळीच्या लांबीपेक्षा कमी आहे.
\end{proof}



\chapter{$\Th{Q}$ मधील निरूपणीयता}\label{inc:req::chap}








\section{परिचय}\label{inc:req:int:sec}

अपूर्णता प्रमेये अशा उपपत्तींना लागू होतात ज्यांत संगणनक्षम
फलनांविषयीची मूलभूत तथ्ये व्यक्त आणि सिद्ध करता येतात. अशा प्रकारची
अगदी किमान उपपत्ती आपण पाहणार आहोत; तिला ``$\Th{Q}$'' म्हणतात
(राफेल रॉबिन्सन यांच्या नावावरून तिला कधीकधी ``रॉबिन्सनची $Q$''
असेही म्हणतात). एखादे फलन $\Th{Q}$ मध्ये \emph{निरूपणीय} असणे म्हणजे
काय, हे आपण सांगू आणि मग पुढील विधान सिद्ध करू:
\begin{quote}
  फलन $\Th{Q}$ मध्ये निरूपणीय असते, तर आणि तरच ते संगणनक्षम असते.
\end{quote}
यामुळे आपल्याला संगणनक्षमतेचे आणखी एक प्रतिमान मिळते. शिवाय,
थांबण्याच्या समस्येचे त्याकडे न्यूनीकरण करून
$\Setabs{\varphi}{\Th{Q} \Proves \varphi}$ हा संच निर्णेय नाही, हे दाखवण्यासाठीही
त्याचा उपयोग करू. पुढे याहून अधिक बलवान निष्कर्ष आपण सिद्ध करू.

$\Th{Q}$ ची भाषा ही अंकगणिताची भाषा आहे; $\Th{Q}$ मध्ये पुढील
स्वयंसिद्धके आहेत. ती एकरूपता असलेल्या प्रथम-क्रम तर्कशास्त्रातील
इतर स्वयंसिद्धके आणि नियम यांच्यासोबत वापरायची आहेत:
\begin{align*}
& \lforall[x][\lforall[y][(\eq[x'][y'] \lif \eq[x][y])]] \tag{$\mathsf{Q}_1$}\\
& \lforall[x][\eq/[\Obj 0][x']] \tag{$\mathsf{Q}_2$}\\
& \lforall[x][(\eq[x][\Obj 0] \lor \lexists[y][\eq[x][y']])] \tag{$\mathsf{Q}_3$}\\
& \lforall[x][\eq[(x + \Obj 0)][x]] \tag{$\mathsf{Q}_4$}\\
& \lforall[x][\lforall[y][\eq[(x + y')][(x + y)']]] \tag{$\mathsf{Q}_5$}\\
& \lforall[x][\eq[(x \times \Obj 0)][\Obj 0]] \tag{$\mathsf{Q}_6$}\\
& \lforall[x][\lforall[y][\eq[(x \times y')][((x \times y) + x)]]] \tag{$\mathsf{Q}_7$}\\
& \lforall[x][\lforall[y][(x < y \liff \lexists[z][\eq[(z' + x)][y]])]] \tag{$\mathsf{Q}_8$}
\end{align*}
प्रत्येक नैसर्गिक संख्या $n$ साठी $\num{n}$ हा संख्यांक
$\Obj{0}^{\prime\prime\ldots\prime}$ हे पद म्हणून परिभाषित करा;
त्यात एकूण $n$ उत्तरवर्ती चिन्हे असतात. म्हणजे $\num{0}$ हा
स्थिरांक~$\Obj{0}$ एकटाच, $\num{1}$ हे $\Obj{0}'$,
$\num{2}$ हे $\Obj{0}''$, इत्यादी.

अंकगणिताची उपपत्ती म्हणून $\Th{Q}$ \emph{अत्यंत} दुर्बल आहे;
उदाहरणार्थ, $\lforall[x][\eq/[x][x']]$ किंवा
$\lforall[x][\lforall[y][(x + y) = (y + x)]]$ यांसारखी
अगदी साधी तथ्येसुद्धा त्यात सिद्ध करता येत नाहीत. पण $\Th{Q}$
इतकी रंजक असण्याचे मोठे कारण ती \emph{दुर्बल} आहे, हेच आहे.
अपूर्णता प्रमेय लागू होण्याइतकीच तिची शक्ती आहे. $\Th{Q}$
रंजक असण्याचे आणखी एक कारण म्हणजे तिचा स्वयंसिद्धकांचा संच
\emph{सांत} आहे.

$\Th{Q}$ पेक्षा अधिक बलवान उपपत्ती \emph{पेआनो अंकगणित}
$\Th{PA}$ हिला $\Th{Q}$ मध्ये विगमन-रूपबंध जोडून मिळवतात:
\[(\varphi(\Obj 0) \land \lforall[x][(\varphi(x) \lif \varphi(x'))]) \lif \lforall[x][\varphi(x)]
\]
येथे $\varphi(x)$ हे कोणतेही सूत्र असते. $\varphi(x)$ मध्ये $x$ व्यतिरिक्त
इतर चले मुक्त असतील, तर त्या सर्वांना बांधण्यासाठी
आरंभी सार्वत्रिक संख्यक जोडतो (म्हणजे विगमन-रूपबंधाचे संबंधित
उदाहरण वाक्य असेल). उदाहरणार्थ, $\varphi(x, y)$ मध्ये
चल~$y$ हेसुद्धा मुक्त असेल, तर संबंधित उदाहरण असे आहे:
\[\lforall[y][((\varphi(\Obj 0) \land \lforall[x][(\varphi(x) \lif \varphi(x'))]) \lif
  \lforall[x][\varphi(x)])]
\]
विगमन-रूपबंधाची उदाहरणे वापरून $\Th{Q}$ च्या स्वयंसिद्धकांपेक्षा
$\Th{PA}$ च्या स्वयंसिद्धकांपासून बरेच अधिक सिद्ध करता येते.
खरे तर, पेआनो अंकगणितात सिद्ध करता येत नाहीत अशी नैसर्गिक
संख्यांविषयीची ``स्वाभाविक'' विधाने शोधणे बरेच कठीण असते!{}

\begin{defn}
\label{inc:req:int:defn:representable-fn}
  नैसर्गिक संख्यांपासून नैसर्गिक संख्यांकडे जाणारे
  $f(x_0,\ldots,x_k)$ हे फलन $\Th{Q}$ मध्ये
  {\em निरूपणीय} आहे असे म्हणतात, जर असे सूत्र
  $\varphi_f(x_0,\dots,x_k,y)$ असेल की
  $f(n_0,\dots,n_k) = m$ असेल तेव्हा $\Th{Q}$ पुढील दोन्ही
  विधाने सिद्ध करते:
\begin{enumerate}
\item\label{inc:req:int:defn:rep:a} $\varphi_f(\num{n_0}, \dots, \num{n_k}, \num{m})$
\item\label{inc:req:int:defn:rep:b} $\lforall[y][(\varphi_f(\num{n_0}, \dots,
\num{n_k}, y) \lif \num{m} = y)]$.
\end{enumerate}
\end{defn}

ही व्याख्या इतर मार्गांनीही मांडता येते. उदाहरणार्थ, $\Th{Q}$
$\lforall[y][(\varphi_f(\num{n_0}, \dots,
\num{n_k}, y) \liff \eq[y][\num{m}])]$ सिद्ध करते, अशी
समतुल्य अटही आपण ठेवू शकलो असतो.

\begin{thm}
\label{inc:req:int:thm:representable-iff-comp}
फलन $\Th{Q}$ मध्ये निरूपणीय असते, तर आणि तरच ते संगणनक्षम असते.
\end{thm}

या प्रमेयाच्या पुराव्याला दोन दिशा आहेत. विन्यासमीमांसेचे अंकगणितीकरण
झाल्यावर डावीकडून उजवीकडे जाणारी दिशा तुलनेने सोपी आहे. दुसऱ्या
दिशेसाठी अधिक काम लागते. मूलभूत कल्पना अशी: ``संगणनक्षम'' याचा
अर्थ नेमका करण्यासाठी ``सामान्य पुनरावर्ती'' हा पर्याय निवडतो आणि
प्रत्येक सामान्य पुनरावर्ती फलन $\Th{Q}$ मध्ये निरूपणीय आहे,
हे दाखवतो. लक्षात घ्या, शून्य फलन~$\Zero$, उत्तरवर्ती फलन~$\Succ$
आणि प्रक्षेपण फलने~$\Proj{n}{i}$ यांच्यापासून संयोजन, आदिम पुनरावर्तन
आणि नियमित लघुतमीकरण वापरून परिभाषित करता येणारे फलन म्हणजे
सामान्य पुनरावर्ती फलन. म्हणून प्रत्येक सामान्य पुनरावर्ती फलन
$\Th{Q}$ मध्ये निरूपणीय आहे हे दाखवण्याचा एक मार्ग असा: मूलभूत
फलने निरूपणीय आहेत हे दाखवायचे आणि काही फलने निरूपणीय असताना
त्यांच्यापासून संयोजन, आदिम पुनरावर्तन किंवा नियमित लघुतमीकरणाने
बनणारी फलनेही निरूपणीय असतात, हे दाखवायचे. दुसऱ्या शब्दांत,
मूलभूत फलने निरूपणीय आहेत आणि निरूपणीय फलनांचा वर्ग संयोजन,
आदिम पुनरावर्तन व नियमित लघुतमीकरण यांखाली ``बंद'' आहे,
हे दाखवता येईल. त्यावरून प्रत्येक सामान्य पुनरावर्ती फलन
निरूपणीय असल्याची हमी मिळते.

पण निरूपणीय फलनांचा वर्ग आदिम पुनरावर्तनाखाली बंद आहे हे दाखवण्याची
पायरी कठीण ठरते. ती टाळण्यासाठी, आधी आपण दाखवतो की प्रत्यक्षात
आदिम पुनरावर्तनाशिवायही काम चालते. म्हणजे, प्रत्येक सामान्य
पुनरावर्ती फलन मूलभूत फलनांपासून केवळ संयोजन आणि नियमित
लघुतमीकरण वापरून परिभाषित करता येते, हे दाखवतो. त्यासाठी आदिम
पुनरावर्तन एका विशिष्ट नियमित लघुतमीकरणाने करता येते, हे दाखवतो.
मात्र हे साधण्यासाठी काही आणखी मूलभूत फलने घ्यावी लागतात:
बेरीज, गुणाकार आणि एकरूपता संबंधाचे जातिबोधक फल~$\Char{=}$. मग
या \emph{सर्व} मूलभूत फलनांची $\Th{Q}$ मधील निरूपणीयता आणि
निरूपणीय फलनांचा संयोजन व नियमित लघुतमीकरण यांखालील बंदपणा
दाखवून हे प्रमेय सिद्ध करता येते.











\section{$\Th{Q}$ मध्ये निरूपणीय फलने संगणनक्षम असतात}\label{inc:req:rpc:sec}

$\Th{Q}$ मध्ये निरूपणीय असलेले प्रत्येक फलन संगणनक्षम आहे,
हे आपण सिद्ध करू. त्याआधी $\Th{Q}$ मध्ये निरूपणीय फलनांविषयी
एक साहाय्यक विधान सिद्ध करावे लागेल.

\begin{lem}\label{inc:req:rpc:lem:rep-q}
  $f(x_0,
  \dots, x_k)$ हे $\Th{Q}$ मध्ये निरूपणीय असेल, तर
  असे सूत्र~$\varphi_f(x_0, \dots, x_k, y)$ असते की
  \[  \Th{Q} \Proves \varphi_f(\num{n_0}, \dots, \num{n_k}, \num{m})
  \quad\text{तेव्हा आणि केवळ तेव्हाच}\quad m = f(n_0, \dots, n_k).
  \]
\end{lem}

\begin{proof}
  ``तर'' हा भाग
\ref{inc:req:int:defn:representable-fn}\ref{inc:req:int:defn:rep:a} मधून मिळतो.
``तरच'' हा भाग पुढीलप्रमाणे दिसतो: समजा $\Th{Q} \Proves \varphi_f(\num{n_0},
\dots, \num{n_k}, \num{m})$, पण $m \neq f(n_0, \dots, n_k)$.
$l = f(n_0, \dots, n_k)$ असे धरू.
\ref{inc:req:int:defn:representable-fn}\ref{inc:req:int:defn:rep:a} नुसार $\Th{Q}
\Proves \varphi_f(\num{n_0}, \dots, \num{n_k}, \num{l})$.
\ref{inc:req:int:defn:representable-fn}\ref{inc:req:int:defn:rep:b} नुसार
$\Th{Q} \Proves \lforall[y][(\varphi_f(\num{n_0}, \dots, \num{n_k}, y) \lif \num{l} =
y)]$. तर्कशास्त्र वापरून आणि $\Th{Q} \Proves
\varphi_f(\num{n_0}, \dots, \num{n_k}, \num{m})$ या गृहीतकावरून
$\Th{Q} \Proves \eq[\num{l}][\num{m}]$ मिळते. दुसरीकडे,
\ref{inc:req:bre:lem:q-proves-neq} नुसार $\Th{Q} \Proves
\eq/[\num{l}][\num{m}]$. म्हणून $\Th{Q}$ विसंगत ठरेल.
पण ते अशक्य आहे: मानक प्रतिमान $\Th{Q}$ ची पूर्ती करते
(पाहा \ref{inc:int:def:def:standard-model}), म्हणजे
$\Sat{N}{\Th{Q}}$; आणि निर्दोषता प्रमेयानुसार प्रत्येक
पूर्ततायोग्य उपपत्ती सुसंगत असते
(\ref{fol:axd:sou:cor:consistency-soundness}, \ref{fol:seq:sou:cor:consistency-soundness}, \ref{fol:ntd:sou:cor:consistency-soundness}, \ref{fol:tab:sou:cor:consistency-soundness}).
\end{proof}

\begin{lem}
$\Th{Q}$ मध्ये निरूपणीय असलेले प्रत्येक फलन संगणनक्षम आहे.
\end{lem}

\begin{proof}
हे खरे असण्यामागील अंतर्ज्ञानी कल्पना आधी पाहू. $f$ चे मूल्य
काढण्यासाठी पुढील पद्धत वापरू. अंकगणिताच्या भाषेतील सर्व संभाव्य
निष्पत्त्या~$\delta$ यांची यादी करायची; हे यांत्रिकरीत्या
करता येते. प्रत्येकासाठी, ती $\varphi_f(\num{n_0}, \dots,
\num{n_k}, \num{m})$ या रूपाच्या सूत्राची
निष्पत्ती आहे का ते तपासायचे (हे $f$ चे $\Th{Q}$
मधील निरूपण करणारे सूत्र आहे; पाहा \ref{inc:req:rpc:lem:rep-q}).
असेल, तर \ref{inc:req:rpc:lem:rep-q} नुसार $m = f(n_0, \dots, n_k)$;
म्हणजे $f$ चे मूल्य सापडले. हा शोध संपतो, कारण
$\Th{Q} \Proves \varphi_f(\num{n_0}, \dots, \num{n_k},
\num{f(n_0, \dots, n_k)})$; त्यामुळे योग्य प्रकारची
$\delta$ अखेरीस सापडते.

ही मांडणी अजून पूर्णपणे नेमकी नाही: आपली पद्धत फक्त संख्यांवर
न चालता निष्पत्त्या आणि सूत्रे वर चालते, आणि
``सर्व संभाव्य निष्पत्त्यांची यादी'' यांत्रिकरीत्या कशी
करता येते हे आपण अजून स्पष्ट केलेले नाही. पण आपण पाहिलेच आहे
की पदे, सूत्रे आणि निष्पत्त्या यांना गोडेल संख्यांनी
संकेतबद्ध करता येते. संगणनाचे नेमके प्रतिमान म्हणून सामान्य
पुनरावर्ती फलनेही आपण मांडली आहेत. तसेच
$\Prf[\Th{Q}](d,y)$ हा संबंध (आदिम) पुनरावर्ती आहे:
नैसर्गिक निगमनातील थेट निष्पत्ती-संकेतांकाचा अर्थ घेतल्यास,
$d$ ही $\Th{Q}$ च्या स्वयंसिद्धकांपासून $y$ ही गोडेल संख्या
असलेल्या सूत्राच्या निष्पत्तीची गोडेल संख्या
असेल, तर आणि तरच हा संबंध खरा असतो. आपल्याला आणखी लागणारी आदिम
पुनरावर्ती फलने म्हणजे $\fn{num}$
(\ref{inc:art:trm:prop:num-primrec}) आणि $\fn{Subst}$
(\ref{inc:art:sub:prop:subst-primrec}). यांपासून लघुतमीकरणाने
$f$ परिभाषित करता येते; म्हणून $f$ पुनरावर्ती आहे.

प्रथम पुढील व्याख्या करा:
\begin{multline*}
  A(n_0, \dots, n_k, m) = \\
  \fn{Subst}(\fn{Subst}(\dots\fn{Subst}(\Gn{\varphi_f}, \fn{num}(n_0), \Gn{x_0}),\\
  \dots),
  \fn{num}(n_k),  \Gn{x_k}), \fn{num}(m), \Gn{y})
\end{multline*}
हे गुंतागुंतीचे दिसते, पण हे फक्त $A(n_0, \dots, n_k,
m) = \Gn{\varphi_f(\num{n_0}, \dots, \num{n_k}, \num{m})}$
हे फलन आहे.

आता $R(n_0, \dots, n_k, s)$ हा संबंध पाहा. $(s)_0$ ही
$\varphi_f(\num{n_0}, \dots, \num{n_k}, \num{(s)_1})$ च्या
$\Th{Q}$ पासूनच्या निष्पत्तीची गोडेल संख्या असेल
तेव्हा हा संबंध खरा असतो:
\[R(n_0, \dots, n_k, s) \quad\text{तेव्हा आणि केवळ तेव्हाच}\quad \Prf[\Th{Q}]((s)_0, A(n_0,
\dots, n_k, (s)_1))
\]
$R(n_0, \dots, n_k, s)$ खरा असणारा एखादा $s$ मिळाला,
तर $(s)_0$ आणि $(s)_1$ ही संख्यांची जोडी मिळाली आहे:
$(s)_0$ ही $\varphi_f(\num{n_0}, \dots,
\num{n_k}, \num{(s)_1})$ च्या निष्पत्तीची गोडेल संख्या
असते. म्हणून $s$ शोधणे म्हणजे अनौपचारिक पुराव्यातील $d$ आणि
$m$ ही जोडी शोधण्यासारखे आहे. असा $s$ ``शोधणारे'' संगणनक्षम
फलन नियमित लघुतमीकरणाने परिभाषित करता येते. लक्षात घ्या,
$R$ साठी हा शोध नियमित आहे: प्रत्येक $n_0$, \dots, $n_k$
साठी $\Th{Q} \Proves \varphi_f(\num{n_0}, \dots,
\num{n_k}, \num{f(n_0, \dots, n_k)})$ हे सिद्ध करणारी निष्पत्ती~$\delta$
असते. म्हणून $s = \tuple{\Gn{\delta}, f(n_0, \dots, n_k)}$
घेतल्यास $R(n_0, \dots, n_k, s)$ खरा असतो. त्यामुळे $f$
पुढीलप्रमाणे लिहिता येते:
\[f(n_0,\dots,n_{k}) = (\umin{s}{R(n_0, \dots, n_k, s)})_1.
\]
\end{proof}











\section{बीटा फलनावरील साहाय्यक विधान}\label{inc:req:bet:sec}


बेरीज, गुणाकार आणि $\Char{=}$ उपलब्ध असतील तर आदिम पुनरावर्तन
करता येते हे दाखवण्यासाठी क्रमिका हाताळणारी फलने विकसित करावी लागतील.
(घातांक फलनही उपलब्ध असते, तर हे काम सोपे झाले असते.) आदिम
पुनरावर्तन उपलब्ध असताना ``$n$ वी अविभाज्य संख्या'' यांसारख्या
गोष्टी परिभाषित करून तुलनेने सरळ सांकेतीकरण निवडता येत होते.
पण येथे आदिम पुनरावर्तन उपलब्ध नाही---उलट, लघुतमीकरण वापरून
आदिम पुनरावर्तन करता येते हेच दाखवायचे आहे---म्हणून अधिक
युक्तीने काम करावे लागेल.

\begin{lem}
\label{inc:req:bet:lem:beta}
असे फलन $\beta(d,i)$ आहे की प्रत्येक क्रमिका $a_0$,
\dots,~$a_n$ साठी अशी संख्या~$d$ असते की प्रत्येक $i \le n$
साठी $\beta(d,i) = a_i$. शिवाय, $\beta$ मूलभूत फलनांपासून
केवळ संयोजन आणि नियमित लघुतमीकरण वापरून परिभाषित करता येते.
\end{lem}

$d$ हा क्रमिका $\tuple{a_0, \dots, a_n}$ चा संकेतांक आहे
आणि $\beta(d,i)$ तिचा $i$ वा घटक देतो, असे समजा.
(लक्षात घ्या, या ``सांकेतीकरणात'' आपल्याला परिचित असलेले
अविभाज्य संख्यांच्या घातांवर आधारित सांकेतीकरण \emph{वापरलेले नाही}!)
हे साहाय्यक विधान किमान आवश्यक तेवढेच सांगते: क्रमिका जोडता
येतात, त्यांना घटक शेवटी जोडता येतो किंवा संयोजन व नियमित
लघुतमीकरणाने परिभाषित फलने वापरून $a_0$, \dots,~$a_n$
यांपासून $d$ \emph{काढता} येतो, असेही ते म्हणत नाही. ते फक्त
एवढेच सांगते की प्रत्येक क्रमिकेसाठी काही संकेतांक असतो आणि
त्यातून घटक मिळवणारे एक ``विसांकेतीकरण'' फलन आहे.

$\beta$ हे चिन्हांकन गोडेल यांचे आहे. पुन्हा सांगायचे तर,
या साहाय्यक विधानाच्या पुराव्यातील कठीण भाग म्हणजे वरवर
मर्यादित वाटणारी साधने---केवळ संयोजन आणि लघुतमीकरण---वापरून
योग्य $\beta$ परिभाषित करणे; अर्थात बेरीज, गुणाकार आणि
$\Char{=}$ वापरता येतात. या विधानाचे अनेक पुरावे आहेत; पण
त्यांपैकी एक सुबोध मार्ग आजही गोडेल यांची मूळ पद्धत आहे.
त्या पद्धतीत सुनझीचे प्रमेय (परंपरेने ``चिनी शेष प्रमेय'')
हे संख्या-सैद्धान्तिक तथ्य वापरले जाते.

\begin{defn}
दोन नैसर्गिक संख्या $a$ आणि $b$ या \emph{सापेक्षतः अविभाज्य}
(सहमूळ) असतात, तर आणि तरच त्यांचा महत्तम साधारण
विभाजक~$1$ असतो; म्हणजे त्यांचा दुसरा कोणताही सामाईक विभाजक नसतो.
\end{defn}

\begin{defn}
धन भाजक $c>0$ साठी नैसर्गिक संख्या $a$ आणि $b$ या \emph{भाजक~$c$ समशेष},
$a \equiv b \mod c$, असतात, तर आणि तरच
$c \mid (a-b)$; म्हणजे $a$ आणि $b$ यांना $c$ ने भागल्यावर
शेष समान मिळतो.

\end{defn}

सुनझीचे प्रमेय असे आहे:
\begin{thm}
समजा $x_0$, \dots,~$x_n$ या धन नैसर्गिक संख्या परस्पर जोड्यांनी सापेक्षतः
अविभाज्य आहेत. $y_0$, \dots,~$y_n$ या कोणत्याही संख्या असू द्या.
तर अशी संख्या $z$ आहे की
\begin{align*}
z & \equiv y_0 \mod x_0 \\
z & \equiv y_1 \mod x_1 \\
& \vdots  \\
z & \equiv y_n \mod x_n.
\end{align*}
\end{thm}

सुनझीचे प्रमेय आपण पुढीलप्रमाणे वापरू: $x_0$, \dots,~$x_n$
या अनुक्रमे $y_0$, \dots,~$y_n$ पेक्षा मोठ्या असतील, तर
$z$ हा क्रमिका $\tuple{y_0, \dots,y_n}$ चा संकेतांक घेता येतो.
$y_i$ परत मिळवण्यासाठी $z$ ला $x_i$ ने भागून शेष घ्यायचा.
या सांकेतीकरणासाठी $x_0$, \dots,~$x_n$ यांची योग्य मूल्ये
शोधावी लागतील.

यासाठी पुढील दोन निरीक्षणे उपयोगी पडतील. $y_0$, \dots,~$y_n$
दिली असता, पुढीलप्रमाणे घ्या:
\begin{align*}
j &= \max(n, y_0 + 1, \dots, y_n + 1), \\
m &= \lcm(1,\dots,j),
\end{align*}
दुसऱ्या ओळीत लघुतम साधारण विभाज्य घेतला आहे.
आणि पुढीलप्रमाणे घ्या:
\begin{align*}
x_0 & = 1 + m \\
x_1 & = 1 + 2 \cdot m \\
x_2 & = 1 + 3 \cdot m \\
& \vdots  \\
x_n & = 1 + (n+1) \cdot m
\end{align*}
मग पुढील दोन गोष्टी खऱ्या आहेत:
\begin{enumerate}
\item\label{inc:req:bet:rel-prime} $x_0,\dots,x_n$ या परस्पर जोड्यांनी
  सापेक्षतः अविभाज्य आहेत.
\item\label{inc:req:bet:less} प्रत्येक $i$ साठी $y_i < x_i$.
\end{enumerate}
\ref{inc:req:bet:rel-prime} हे खरे का, ते पाहू. दोन भिन्न निर्देशांक
निवडा. $p$ ही अविभाज्य संख्या
असेल आणि $p \mid x_i$ व $p \mid x_k$ असतील, तर
$p \mid 1 + (i+1) m$ आणि $p \mid 1 + (k+1) m$.
मग $p$ त्यांच्या फरकालाही विभाजित करते:
\[(1 + (i+1)m) - (1+ (k+1)m) = (i-k) m.
\]
$p$ ही $1 + (i+1)m$ ला विभाजित करते, म्हणून ती $m$ ला
विभाजित करू शकत नाही (अन्यथा पहिल्या भागाकारात शेष~$1$
उरेल). त्यामुळे $p$ ही $(i-k)m$ ला विभाजित करत असल्याने
$p$ ही $i-k$ लाही विभाजित करते. पण $\left|i-k\right|$ जास्तीत
जास्त~$n$ आहे, आणि $j \geq n$ असे घेतले आहे; म्हणून
पुन्हा $p \mid m$, हा विरोधाभास. म्हणून $x_i$ आणि
$x_k$ या दोन्हींना विभाजित करणारी अविभाज्य संख्या नाही.
\ref{inc:req:bet:less} ही अट सोपी आहे: $y_i < j \leq m < x_i$.

आता $\beta$ फलनावरील साहाय्यक विधान सिद्ध करू. लक्षात ठेवा,
$0$, उत्तरवर्ती फलन, बेरीज, गुणाकार, $\Char{=}$, प्रक्षेपणे,
आणि नियमित फलनांना लघुतमीकरण लागू करून व संयोजनाने
यांपासून परिभाषित केलेली कोणतीही फलने वापरता येतात.
एखाद्या संबंधाचे जातिबोधक फल असे परिभाषित करता येत असेल,
तर तो संबंधही वापरता येतो. आधीप्रमाणे, हे संबंध बूलीय
संयोग आणि परिबद्ध संख्यकीकरण यांखाली बंद आहेत हे दाखवता येते;
उदाहरणार्थ:
\begin{align*}
\fn{not}(x) & \defis \Char{=}(x,0)\\
\bmin{x \leq z}{R(x,y)} & \defis \umin{x}{(R(x,y) \lor x = z)}\\
\bexists{x \leq z}{R(x,y)} & \defiff R(\bmin{x \leq z}{R(x,y)}, y)
\end{align*}
मग पुढील सर्व गोष्टी आदिम पुनरावर्तनाशिवायही परिभाषित
करता येतात, हे दाखवता येते:
\begin{enumerate}
\item जोडीकरण फलन, $J(x,y) = \frac{1}{2}[(x+y)(x+y+1)] + x$;

\item प्रक्षेपण फलने
\begin{align*}
K(z) & = \bmin{x \leq z}{\bexists{y \leq z}{z = J(x,y)}},\\
L(z) & = \bmin{y \leq z}{\bexists{x \leq z}{z = J(x,y)}};
\end{align*}
\item $x < y$ हा लघुत्वाचा संबंध;
\item $x \mid y$ हा विभाज्यतेचा संबंध;






\item $y$ ला $x$ ने भागल्यावर मिळणारा शेष देणारे
  $\fn{rem}(x,y)$ हे फलन.
\end{enumerate}
शून्य भाजकासाठी या शेष फलाचे मूल्य येथे दिलेले नाही; आधीच्या
शेष फलाच्या परिपाठाप्रमाणे ते शून्य घेता येते. पुढील
वापरातील भाजक मात्र नेहमी धन असतो.
आता पुढील व्याख्या करा:
\begin{align*}
\beta^*(d_0,d_1,i) & = \fn{rem}(1+(i+1) d_1,d_0) \text{ आणि}\\
\beta(d,i) & = \beta^*(K(d),L(d),i).
\end{align*}
हेच आपल्याला हवे असलेले फलन आहे. वर दिलेल्या
$a_0,\dots,a_n$ साठी
\[j = \max(n,a_0+1,\dots,a_n+1),
\]
असे घ्या, आणि $d_1 = \lcm(1,\dots,j)$ असे घ्या.
वरील \ref{inc:req:bet:rel-prime} नुसार $1+d_1$, $1+2 d_1$,
\dots, $1+(n+1) d_1$ या सापेक्षतः अविभाज्य आहेत,
आणि \ref{inc:req:bet:less} नुसार त्या अनुक्रमे
$a_0,\dots,a_n$ पेक्षा मोठ्या आहेत. सुनझीच्या प्रमेयानुसार
प्रत्येक~$i$ साठी पुढील अट पूर्ण करणारे $d_0$ चे मूल्य आहे:
\[d_0 \equiv a_i \mod (1+(i+1)d_1)
\]
म्हणून ($d_1$ हे~$a_i$ पेक्षा मोठे असल्यामुळे)
\[a_i = \fn{rem}(1+(i+1)d_1,d_0).
\]
$d = J(d_0,d_1)$ असे घ्या. मग प्रत्येक $i \le n$ साठी
\begin{align*}
\beta(d,i) & =  \beta^*(d_0,d_1,i) \\
& =  \fn{rem}(1+(i+1) d_1,d_0) \\
& =  a_i
\end{align*}
हेच अपेक्षित होते. याने $\beta$-फलनावरील साहाय्यक विधानाचा
पुरावा पूर्ण होतो.

\begin{prob}
  $x < y$, $x \mid y$ हे संबंध आणि $\fn{rem}(x,y)$ हे
  फलन आदिम पुनरावर्तनाशिवाय परिभाषित करता येते, हे दाखवा.
  $0$, उत्तरवर्ती फलन, बेरीज, गुणाकार, $\Char{=}$,
  प्रक्षेपणे, तसेच परिबद्ध लघुतमीकरण आणि परिबद्ध संख्यकीकरण
  वापरू शकता.
\end{prob}











\section{आदिम पुनरावर्तनाचे अनुकरण}\label{inc:req:pri:sec}

बीटा फलन वापरून नियमित लघुतमीकरणाने आदिम पुनरावर्तनाने केलेल्या
व्याख्येचे ``अनुकरण'' करता येते, हे आता दाखवू. समजा आपल्याकडे
$f(\vec x)$ आणि $g(\vec x, y, z)$ ही फलने आहेत. मग $f$ आणि~$g$
यांपासून आदिम पुनरावर्तनाने परिभाषित केलेले फलन~$h(\vec x,y)$
असे आहे:
\begin{align*}
h(\vec x, 0) & =  f(\vec x) \\
h(\vec x, y+1) & =  g(\vec x, y, h(\vec x, y)).
\end{align*}
मूलभूत फलने आणि त्यांपासून संयोजन व नियमित लघुतमीकरणाने
परिभाषित केलेली फलने (उदाहरणार्थ~$\beta$) वापरून,
$f$ आणि~$g$ यांपासून केवळ संयोजन व नियमित लघुतमीकरणाने
$h$ परिभाषित करता येते, हे दाखवायचे आहे.

\begin{lem}
\label{inc:req:pri:lem:prim-rec}
$h$ हे $f$ आणि $g$ यांपासून आदिम पुनरावर्तनाने परिभाषित
करता येत असेल, तर $f$, $g$ आणि $\Zero$, $\Succ$,
$\Proj{n}{i}$, $\Add$, $\Mult$, $\Char{=}$ ही फलने
वापरून संयोजन व नियमित लघुतमीकरणानेही ते परिभाषित करता येते.
\end{lem}

\begin{proof}
प्रथम $\hat h(\vec x, y)$ हे साहाय्यक फलन परिभाषित करू.
ते लघुतम संख्या~$d$ परत करते; हा $d$ पुढील अटी
पूर्ण करणाऱ्या क्रमिकेचा संकेतांक असतो:
\begin{enumerate}
\item $(d)_0 = f(\vec x)$, आणि
\item प्रत्येक $i < y$ साठी $(d)_{i+1} = g(\vec x, i, (d)_i)$,
\end{enumerate}
येथे $(d)_i$ हा आता $\beta(d,i)$ चा संक्षेप आहे. दुसऱ्या
शब्दांत, $\hat h$ हे $\tuple{h(\vec x, 0), h(\vec x, 1),
\dots, h(\vec x, y)}$ या क्रमिकेचा संकेतांक परत करते.
$\hat h$ पुढीलप्रमाणे लिहिता येते:
\[\hat h(\vec x, y) = \umin{d}{(\beta(d,0) = f(\vec x) \land \bforall{i <
  y}{\beta(d,i+1) = g(\vec x, i,\beta(d,i)})}.
\]
लक्षात घ्या: येथे आदिम पुनरावर्तन लागत नाही; फक्त
लघुतमीकरण पुरेसे आहे. बीटा फलनावरील साहाय्यक विधान
\ref{inc:req:bet:lem:beta} नुसार या लघुतमीकरणातील शोध नियमित आहे.

आता आपल्याकडे
\[h(\vec x, y) = \beta(\hat h(\vec x, y), y),
\]
म्हणून मूळ दिलेली दोन फलने आणि मूलभूत फलने यांपासून केवळ संयोजन
व नियमित लघुतमीकरणाने $h$ परिभाषित करता येते.
\end{proof}











\section{मूलभूत फलने~$\Th{Q}$ मध्ये निरूपणीय आहेत}\label{inc:req:bre:sec}

सर्वप्रथम सर्व मूलभूत फलने~$\Th{Q}$ मध्ये निरूपणीय आहेत, हे
दाखवायचे आहे. अखेरीस, $k$ आदाने असलेल्या प्रत्येक मूलभूत
फलनासाठी $f(x_0,\dots,x_{k-1})$ त्याचे निरूपण करणारे
सूत्र $\varphi_f(x_0,\dots,x_{k-1},y)$ कसे द्यायचे,
हे दाखवावे लागेल.

शून्य, उत्तरवर्ती, बेरीज, गुणाकार, समतेचे जातिबोधक फल आणि
प्रक्षेपण फलने यांची निरूपणीयता दाखवता येईल. प्रत्येक बाबतीत
योग्य निरूपक सूत्र सहज सुचते; उदाहरणार्थ, शून्याचे निरूपण
$y = \Obj 0$ या सूत्राने, उत्तरवर्ती फलनाचे सूत्र
$x_0' = y$ या सूत्राने आणि बेरीजेचे सूत्र
$(x_0 + x_1) = y$ या सूत्राने होते. खरे काम म्हणजे संबंधित
वाक्ये $\Th{Q}$ मध्ये सिद्ध होतात, हे दाखवणे. उदाहरणार्थ,
वरील सूत्र बेरीजेचे निरूपण करते असे म्हणण्यासाठी
नैसर्गिक संख्यांच्या प्रत्येक जोडी $m$ आणि $n$ साठी $\Th{Q}$
पुढील विधाने सिद्ध करते, हे दाखवावे लागते:
\begin{align*}
& \eq[\num n + \num m][\num {n+m}] \text{ आणि}\\
& \lforall[y][(\eq[(\num n + \num m)][y] \lif \eq[y][\num{n+m}])].
\end{align*}

\begin{prop}
\label{inc:req:bre:prop:rep-zero}
शून्य फलन $\Zero(x) = 0$ चे~$\Th{Q}$ मध्ये
$\varphi_{\Zero}(x,y) \ident \eq[y][\Obj 0]$ या सूत्राने निरूपण होते.
\end{prop}

\begin{prop}
\label{inc:req:bre:prop:rep-succ}
उत्तरवर्ती फलन $\Succ(x) = x+1$ चे~$\Th{Q}$ मध्ये
$\varphi_{\Succ}(x,y) \ident \eq[y][x']$ या सूत्राने निरूपण होते.
\end{prop}

\begin{prop}
\label{inc:req:bre:prop:rep-proj}
प्रक्षेपण फलन $\Proj{n}{i}(x_0, \dots, x_{n-1}) = x_i$ चे
~$\Th{Q}$ मध्ये पुढील सूत्राने निरूपण होते:
\[\varphi_{\Proj{n}{i}}(x_0, \dots, x_{n-1}, y) \ident \eq[y][x_i].\]
\end{prop}

\begin{prob}
$\eq[y][\Obj 0]$, $\eq[y][x']$ आणि $\eq[y][x_i]$
ही सूत्रे अनुक्रमे $\Zero$, $\Succ$ आणि $\Proj{n}{i}$ यांचे
निरूपण करतात, हे सिद्ध करा.
\end{prob}

\begin{prop}
\label{inc:req:bre:prop:rep-id}
समता~$=$ हिचे जातिबोधक फल,
\[\Char{=}(x_0, x_1) =
\begin{cases}
  1 & \text{जर } x_0 =x_1\\
  0 & \text{अन्यथा}
\end{cases}
\]
याचे~$\Th{Q}$ मध्ये पुढील सूत्राने निरूपण होते:
\[  \varphi_{\Char{=}}(x_0, x_1, y) \ident (\eq[x_0][x_1] \land \eq[y][\num{1}]) \lor (\eq/[x_0][x_1] \land
\eq[y][\num{0}]).
\]
\end{prop}

या पुराव्यासाठी पुढील साहाय्यक विधान लागते.

\begin{lem}
\label{inc:req:bre:lem:q-proves-neq} $n$ आणि $m$ या नैसर्गिक संख्या असू द्या.
$n \neq m$ असेल, तर $\Th{Q} \Proves \eq/[\num n][\num m]$.
\end{lem}

\begin{proof}
$n$ वर विगमन वापरून, प्रत्येक $m$ साठी $n \neq m$ असेल तर
$Q \Proves \eq/[\num n][\num m]$, हे दाखवू.

आधार टप्प्यात $n = 0$. जर $m$ हे $0$ च्या बरोबर नसेल,
तर एखादी नैसर्गिक संख्या $k$ अशी असते की $m = k +
1$. $\lforall[x][\eq/[0][x']]$ असे सांगणारे स्वयंसिद्धक
आपल्याकडे आहे. संख्यकाच्या एका स्वयंसिद्धकानुसार $x$ च्या जागी
$\num k$ ठेवून $\eq/[0][\num k']$ हा निष्कर्ष मिळतो. पण
$\num k'$ म्हणजेच $\num m$.

विगमन टप्प्यात, $n$ साठी विधान खरे आहे असे धरून $n+1$ चा
विचार करू. $m$ ही कोणतीही नैसर्गिक संख्या असू द्या. दोन शक्यता
आहेत: $m = 0$, किंवा एखाद्या $k$ साठी $m = k+1$. पहिली
शक्यता वरप्रमाणे हाताळता येते. दुसऱ्या शक्यतेत $n+1 \neq
k+1$ असे समजा. मग $n \neq k$. $n$ साठीच्या विगमन गृहीतकाने
$\Th{Q} \Proves \eq/[\num n][\num k]$. आपल्याकडे
$\lforall[x][\lforall[y][\eq[x'][y'] \lif \eq[x][y]]]$
असे सांगणारे स्वयंसिद्धक आहे. संख्यकाच्या एका स्वयंसिद्धकाने
$\eq[\num n'][\num k'] \lif \eq[\num n][\num
  k]$ मिळते. विधानीय तर्कशास्त्र वापरून $\Th{Q}$ मध्ये
$\eq/[\num n][\num k] \lif \eq/[\num n'][\num k']$
हा निष्कर्ष मिळतो. आता अभिव्यंजन-निर्मूलन नियमाने
$\eq/[\num n'][\num k']$ मिळते; आणि $\num k'$ म्हणजेच
$\num m$ असल्याने आपल्याला हेच दाखवायचे होते.
\end{proof}

\begin{explain}
हे साहाय्यक विधान फार मोठी गोष्ट सांगत नाही: वेगवेगळे संख्यांक
वेगवेगळ्या वस्तू दर्शवतात, हे $\Th{Q}$ सिद्ध करू शकते,
असे त्याचे सार आहे. उदाहरणार्थ, $\Th{Q}$ हे
$0'' \neq 0'''$ सिद्ध करते. पण हे सर्वसाधारणपणे खरे आहे
हे दाखवताना काही काळजी घ्यावी लागते. येथे विगमन वापरले असले,
तरी ते $\Th{Q}$ च्या \emph{बाहेर} केलेले विगमन आहे, हेही लक्षात घ्या.
\end{explain}

\begin{proof}[\ref{inc:req:bre:prop:rep-id} चा पुरावा]
$n = m$ असेल, तर $\num{n}$ आणि $\num{m}$ ही एकच पदे आहेत
आणि $\Char{=}(n, m) = 1$. तसेच $\Th{Q} \Proves (\eq[\num{n}][\num{m}] \land
\eq[\num{1}][\num{1}])$, म्हणून ते $\varphi_=(\num{n}, \num{m},
\num{1})$ सिद्ध करते. $n \neq m$ असेल, तर $\Char=(n, m) = 0$.
\ref{inc:req:bre:lem:q-proves-neq} नुसार $\Th{Q} \Proves \eq/[\num{n}][\num{m}]$;
म्हणून $(\eq/[\num{n}][\num{m}] \land \Obj 0 = \Obj 0)$
हेही सिद्ध होते. त्यामुळे $\Th{Q}
\Proves \varphi_=(\num{n}, \num{m}, \num{0})$.

दुसऱ्या अटीसाठीही दोन शक्यता आहेत. $n = m$ असेल, तर
$\Th{Q} \Proves \lforall[y][(\varphi_=(\num{n}, \num{m}, y)
  \lif \eq[y][\num{1}])]$ हे दाखवायचे आहे. अनौपचारिकपणे
युक्तिवाद करू: $\varphi_=(\num{n}, \num{m}, y)$, म्हणजेच
\[(\eq[\num{n}][\num{n}] \land \eq[y][\num{1}]) \lor
(\eq/[\num{n}][\num{n}] \land \eq[y][\num{0}])
\]
असे समजा. डाव्या विकल्पातून तर्कशास्त्राने
$\eq[y][\num{1}]$ मिळते. उजवा विकल्प मात्र
तर्कशास्त्राने सिद्ध होणाऱ्या $\eq[\num{n}][\num{n}]$
शी विसंगत आहे.

दुसरीकडे, $n \neq m$ असेल, तर $\varphi_=(\num{n},
\num{m}, y)$ असे आहे:
\[(\eq[\num{n}][\num{m}] \land \eq[y][\num{1}]) \lor
(\eq/[\num{n}][\num{m}] \land \eq[y][\num{0}])
\]
येथे डावा विकल्प \ref{inc:req:bre:lem:q-proves-neq} नुसार
$\Th{Q}$ मध्ये सिद्ध होणाऱ्या $\eq/[\num{n}][\num{m}]$
शी विसंगत आहे; उजव्या विकल्पातून $\eq[y][\num{0}]$ मिळते.
\end{proof}

\begin{prop}
\label{inc:req:bre:prop:rep-add}
बेरीज फलन $\Add(x_0, x_1) = x_0+x_1$ चे~$\Th{Q}$ मध्ये
पुढील सूत्राने निरूपण होते:
\[  \varphi_{\Add}(x_0, x_1, y) \ident \eq[y][(x_0 + x_1)].
\]
\end{prop}

\begin{lem}
\label{inc:req:bre:lem:q-proves-add}
$\Th{Q} \Proves \eq[(\num{n} + \num{m})][\num{n+m}]$
\end{lem}

\begin{proof}
$m$ वर विगमन वापरून हे सिद्ध करू. $m = 0$ असेल, तर
$\Th{Q} \Proves \eq[(\num{n} + \Obj 0)][\num{n}]$
हे विधान आहे. ते स्वयंसिद्धक~$\mathsf{Q}_4$ वरून मिळते. आता
$m$ साठी विधान खरे आहे असे समजून $m+1$ साठी, म्हणजे
$\Th{Q} \Proves \eq[(\num{n} +
  \num{m+1})][\num{n+m+1}]$ हे सिद्ध करू. लक्षात घ्या,
$\num{m+1}$ म्हणजेच $\num{m}'$ आणि $\num{n+m+1}$ म्हणजेच
$\num{n+m}'$. स्वयंसिद्धक $\mathsf{Q}_5$ ने $\Th{Q}
\Proves \eq[(\num{n} + \num{m}')][(\num{n}+\num{m})']$.
विगमन गृहीतकाने $\Th{Q} \Proves \eq[(\num{n} + \num{m})][\num{n+m}]$.
म्हणून $\Th{Q} \Proves \eq[(\num{n} + \num{m}')][\num{n+m}']$.
\end{proof}

\begin{proof}[\ref{inc:req:bre:prop:rep-add} चा पुरावा]
$\Add$ चे निरूपण करणारे सूत्र~$\varphi_\Add(x_0, x_1, y)$
म्हणजे $\eq[y][(x_0 + x_1)]$. प्रथम, $\Add(n, m) = k$
असेल, तर $\Th{Q} \Proves \varphi_\Add(\num{n}, \num{m}, \num{k})$,
म्हणजे $\Th{Q}
\Proves \eq[\num{k}][(\num{n} + \num{m})]$, हे दाखवू.
पण $k = n + m$ असल्यामुळे $\num{k}$ म्हणजेच $\num{n+m}$;
आणि \ref{inc:req:bre:lem:q-proves-add} मध्ये आपण
$\Th{Q} \Proves \eq[(\num{n} +
  \num{m})][\num{n+m}]$ हे सिद्ध केले आहे. समतेच्या
सममितीने त्यावरून अपेक्षित उलट क्रमातील समता मिळते.

$\Add(n, m) = k$ असेल, तर पुढील विधानही दाखवायचे आहे:
\[\Th{Q} \Proves \lforall[y][(\varphi_\Add(\num{n}, \num{m}, y) \lif
  \eq[y][\num{k}])].
\]
निरूपक सूत्र गृहीत धरून, समतेच्या सममितीने
$\eq[(\num{n} + \num{m})][y]$ असे समजा. कारण
\[\Th{Q} \Proves \eq[(\num{n}+\num{m})][\num{n+m}],
\]
म्हणून डावीकडील पदाऐवजी $\num{n+m}$ ठेवून कोणत्याही~$y$
साठी $\eq[\num{n+m}][y]$ मिळते. समतेच्या सममितीने
अपेक्षित क्रमातील समता मिळते; उलट अभिव्यंजनही याच सिद्ध
समतेवरून आणि समतेच्या नियमांवरून मिळते.
\end{proof}

\begin{prop}
\label{inc:req:bre:prop:rep-mult}
गुणाकार फलन $\Mult(x_0, x_1) = x_0 \cdot x_1$ चे
~$\Th{Q}$ मध्ये पुढील सूत्राने निरूपण होते:
\[  \varphi_{\Mult}(x_0, x_1, y) \ident y = (x_0 \times x_1).
\]
\end{prop}

\begin{proof} सरावासाठी. \end{proof}

\begin{lem}
\label{inc:req:bre:lem:q-proves-mult}
$\Th{Q} \Proves \eq[(\num{n} \times \num{m})][\num{n \cdot m}]$
\end{lem}

\begin{proof} सरावासाठी. \end{proof}

\begin{prob}
\ref{inc:req:bre:lem:q-proves-mult} सिद्ध करा.
\end{prob}

\begin{prob}
\ref{inc:req:bre:lem:q-proves-mult} वापरून
\ref{inc:req:bre:prop:rep-mult} सिद्ध करा.
\end{prob}

\begin{explain}
  आठवून पाहा: अंकगणिताच्या भाषेतील फलनचिन्हासाठी आपण $\times$
  वापरतो, तर संख्यांवरील नेहमीच्या गुणाकारासाठी $\cdot$ वापरतो.
  म्हणून संख्यादर्शक व्यंजकांमध्ये (उदाहरणार्थ, $m \cdot n$)
  $\cdot$ येऊ शकते; पण अंकगणिताच्या भाषेतील पदांमध्येच
  (उदाहरणार्थ, $(\num{m} \times
  \num{n})$) $\times$ येते. आणखी गोंधळाची बाब म्हणजे $+$
  हे चिन्ह फलन आणि बेरीजक्रिया या दोन्हींसाठी वापरतो.
  ते पदांच्या मध्ये आल्यास---उदाहरणार्थ, $(\num{n} + \num{m})$
  मध्ये---अंकगणिताच्या भाषेतील $2$ आदानांचे फलन असते;
  आणि संख्यांच्या मध्ये आल्यास---उदाहरणार्थ, $n+m$ मध्ये---
  ती बेरीजक्रिया असते. $\num{n+m}$ हे उदाहरणही यांत येते:
  तो $n+m$ या संख्येला अनुरूप असलेला मानक संख्यांक आहे.
\end{explain}











\section{फलनांचे संयोजन~$\Th{Q}$ मध्ये निरूपणीय असते}\label{inc:req:cmp:sec}

समजा $h$ पुढीलप्रमाणे परिभाषित आहे:
\[h(x_0,\dots,x_{l-1}) = f(g_0(x_0,\dots,x_{l-1}), \dots,
g_{k-1}(x_0,\dots,x_{l-1})).
\]
येथे $f$ आणि अनुक्रमे $g_0$,
\dots,~$g_{k-1}$ या फलनांचे निरूपण करणारी
सूत्रे $\varphi_f, \varphi_{g_0}, \dots,
\varphi_{g_{k-1}}$ आपल्याला आधीच सापडली आहेत. आता $h$ चे
निरूपण करणारे सूत्र~$\varphi_h$ शोधायचे आहे.

प्रत्येक फलन $1$-स्थानी असलेली सोपी बाब आधी पाहू; म्हणजे
$h(x) = f(g(x))$ विचारात घेऊ. $\varphi_f(y, z)$ हे~$f$ चे,
आणि $\varphi_g(x, y)$ हे~$g$ चे निरूपण करत असेल, तर $h$ चे
निरूपण करणारे सूत्र~$\varphi_h(x, z)$ शोधायचे आहे.
$h(x) = z$, तर आणि तरच असा एखादा~$y$ असतो की
$z = f(y)$ आणि $y = g(x)$ ही दोन्ही विधाने खरी असतात.
($h(x) = z$ असेल तर $g(x)$ हा असा~$y$ असतो; आणि असा
$y$ असेल, तर $y = g(x)$ आणि $z = f(y)$ असल्याने
$z = f(g(x))$.) यावरून $\varphi_h(x, z)$ साठी
$\lexists[y][(\varphi_g(x, y) \land \varphi_f(y,
  z))]$ हे योग्य सूत्र असावे असे दिसते. आता फक्त
$\Th{Q}$ संबंधित सूत्रे सिद्ध करते, हे तपासायचे आहे.

\begin{prop}
\label{inc:req:cmp:prop:rep1}
$h(n) = m$ असेल, तर $\Th{Q} \Proves \varphi_h(\num{n}, \num{m})$.
\end{prop}

\begin{proof}
$h(n) = m$, म्हणजे $f(g(n)) = m$, असे समजा. $k = g(n)$ असू द्या.
मग
\begin{align*}
  \Th{Q} & \Proves \varphi_g(\num{n}, \num{k})
  \intertext{कारण $\varphi_g$ हे~$g$ चे निरूपण करते; तसेच}
  \Th{Q} & \Proves \varphi_f(\num{k}, \num{m})
  \intertext{कारण $\varphi_f$ हे~$f$ चे निरूपण करते. म्हणून}
  \Th{Q} & \Proves \varphi_g(\num{n}, \num{k}) \land \varphi_f(\num{k}, \num{m})
  \intertext{आणि त्यावरून}
  \Th{Q} & \Proves \lexists[y][(\varphi_g(\num{n}, y) \land \varphi_f(y, \num{m}))],
\end{align*}
म्हणजेच $\Th{Q} \Proves \varphi_h(\num{n}, \num{m})$.
\end{proof}

\begin{prop}
\label{inc:req:cmp:prop:rep2}
$h(n) = m$ असेल, तर $\Th{Q} \Proves \lforall[z][(\varphi_h(\num{n}, z) \lif
  z = \num{m})]$.
\end{prop}

\begin{proof}
$h(n) = m$, म्हणजे $f(g(n)) = m$, असे समजा. $k = g(n)$ असू द्या.
मग
\begin{align*}
  \Th{Q} & \Proves \lforall[y][(\varphi_g(\num{n}, y) \lif \eq[y][\num{k}])]
  \intertext{कारण $\varphi_g$ हे~$g$ चे निरूपण करते; तसेच}
  \Th{Q} & \Proves \lforall[z][(\varphi_f(\num{k}, z) \lif \eq[z][\num{m}])]
  \intertext{कारण $\varphi_f$ हे~$f$ चे निरूपण करते. थोडेसे
    तर्कशास्त्र वापरून पुढील विधानही दाखवता येते:}
  \Th{Q} & \Proves \lforall[z][(\lexists[y][(\varphi_g(\num{n}, y) \land
      \varphi_f(y, z))] \lif \eq[z][\num{m}])].
\end{align*}
म्हणजेच $\Th{Q} \Proves \lforall[y][(\varphi_h(\num n, y) \lif \eq[y][\num m])]$.
\end{proof}

$f$ आणि~$g_i$ यांची आदाने $1$ पेक्षा अधिक असलेल्या गुंतागुंतीच्या
बाबतीतही हीच कल्पना लागू पडते.

\begin{prop}
\label{inc:req:cmp:prop:rep-composition}
$\varphi_f(y_0, \dots, y_{k-1}, z)$ हे~$\Th{Q}$ मध्ये
$f(y_0, \dots, y_{k-1})$ चे, आणि
$\varphi_{g_i}(x_0, \dots, x_{l-1}, y)$ हे~$\Th{Q}$ मध्ये
$g_i(x_0, \dots, x_{l-1})$ चे निरूपण करत असेल, तर
\begin{multline*}
  \lexists[y_0\dots][\lexists[y_{k-1}][(\varphi_{g_0}(x_0,\dots,x_{l-1},y_0) \land
      \dots \land {}]]\\
  \varphi_{g_{k-1}}(x_0,\dots,x_{l-1},y_{k-1}) \land \varphi_f(y_0,\dots,y_{k-1},z))
\end{multline*}
हे सूत्र पुढील फलनाचे निरूपण करते:
\[h(x_0, \dots, x_{l-1}) = f(g_0(x_0, \dots, x_{l-1}), \dots, g_{k-1}(x_0,
\dots, x_{l-1})).
\]
\end{prop}

\begin{proof}
सरावासाठी.
\end{proof}

\begin{prob}
\ref{inc:req:cmp:prop:rep1} आणि
\ref{inc:req:cmp:prop:rep2} यांच्या पुराव्यांचा आधार घेऊन
\ref{inc:req:cmp:prop:rep-composition} चा पुरावा तपशीलवार द्या.
\end{prob}











\section{नियमित लघुतमीकरण~$\Th{Q}$ मध्ये निरूपणीय असते}\label{inc:req:min:sec}

अपरिबद्ध शोधाचा विचार करू. समजा $g(x, z)$ हे नियमित फलन
$\Th{Q}$ मध्ये, उदाहरणार्थ सूत्र~$\varphi_g(x, z, y)$ या
सूत्राने, निरूपणीय आहे. $f(z) = \umin{x}{[g(x, z) = 0]}$
असे $f$ परिभाषित करू. $f$ चे निरूपण करणारे
सूत्र~$\varphi_f(z, y)$ शोधायचे आहे. $f(z)$ चे मूल्य
म्हणजे अशी संख्या~$x$ की (अ) $g(x, z) = 0$ आणि (आ) तशी
ती लघुतम आहे, म्हणजे कोणत्याही $w < x$ साठी
$g(w, z) \neq 0$. म्हणून पुढील निवड स्वाभाविक आहे:
\[\varphi_f(z,y) \ident \varphi_g(y, z, \Obj 0) \land \lforall[w][(w < y \lif \lnot
  \varphi_g(w, z, \Obj 0))].
\]
अर्थात, सर्वसाधारण बाबतीत $z$ ऐवजी $z_0$, \dots, $z_k$
ही चलं घ्यावी लागतील.

या पुराव्यातही $\Th{Q}$ ची सिद्ध करण्याची शक्ती किती आहे,
याविषयी काही साहाय्यक विधाने लागतील.

\begin{lem}
\label{inc:req:min:lem:succ} प्रत्येक स्थिरांक~$a$ आणि प्रत्येक नैसर्गिक
संख्या~$n$ साठी,
\[\Th{Q} \Proves \eq[(a' + \num n)][(a + \num n)'].
\]
\end{lem}

\begin{proof}
नेहमीप्रमाणे $n$ वर विगमन वापरू. आधार टप्प्यात $n =
0$ असेल, तर $\Th{Q}$ हे $\eq[(a' + \Obj 0)][(a + \Obj
0)']$ सिद्ध करते, हे दाखवायचे आहे. पण पुढील विधाने मिळतात:
\begin{align}
  \Th{Q} & \Proves \eq[(a' + \Obj 0)][a'] \quad \text{स्वयंसिद्धक $Q_4$ ने}
  \label{inc:req:min:step1}\\
  \Th{Q} & \Proves  \eq[(a + \Obj 0)][a] \quad \text{स्वयंसिद्धक $Q_4$ ने}
  \label{inc:req:min:step2} \\
  \Th{Q} & \Proves \eq[(a + \Obj 0)'][a'] \quad \text{\ref{inc:req:min:step2} वरून}
  \label{inc:req:min:step3} \\
  \Th{Q} & \Proves \eq[(a' + \Obj 0)][(a + \Obj 0)'] \quad
  \text{\ref{inc:req:min:step1} आणि \ref{inc:req:min:step3} वरून}\notag
\end{align}
विगमन टप्प्यात, $\Th{Q}
\Proves \eq[(a' + \num n)][(a + \num n)']$ हे सिद्ध झाले आहे
असे समजू. $\num{n+1}$ म्हणजे $\num{n}'$ असल्याने, आता
$\Th{Q}$ हे $\eq[(a' + \num{n}')][(a + \num{n}')']$
सिद्ध करते, हे दाखवायचे आहे. आपल्याला मिळते:
\begin{align}
  \Th{Q} & \Proves \eq[(a' + \num n')][(a' + \num n)'] \quad
  \text{स्वयंसिद्धक $\mathsf{Q}_5$ ने} \label{inc:req:min:step5}\\
  \Th{Q} & \Proves \eq[(a' + \num n)'][(a + \num n')'] \quad
  \text{विगमन गृहीतक, बेरीजेचे स्वयंसिद्धक आणि समतेचे नियम वापरून} \label{inc:req:min:step6}\\
  \Th{Q} & \Proves \eq[(a' + \num n')][(a + \num n')'] \quad
  \text{\ref{inc:req:min:step5} आणि \ref{inc:req:min:step6} वरून.} \notag
\end{align}
\end{proof}

हे पुन्हा लक्षात घ्यावे की यावरून $\Th{Q}$ हे
$\lforall[x][\lforall[y][(x' + y) = (x + y)']]$ सिद्ध करते,
असे म्हणता येत नाही; वरचे विधान त्याहून दुर्बल आहे. हे
वाक्य~$\Struct{N}$ मध्ये खरे असले, तरी $\Th{Q}$
त्याला सिद्ध करू शकत नाही.

\begin{lem}
\label{inc:req:min:lem:less-zero}
$\Th{Q} \Proves \lforall[x][\lnot x < \Obj 0]$.
\end{lem}

\begin{proof}
पुरावा अनौपचारिक रीतीने देऊ (म्हणजे औपचारिक निष्पत्ती
कशी रचावी यासाठी फक्त संकेत देऊ).

कोणत्याही~$a$ साठी $\lnot a < \Obj 0$ हे सिद्ध करायचे आहे.
$<$ च्या व्याख्येनुसार $\Th{Q}$ मध्ये
$\lnot \lexists[y][\eq[(y'+a)][\Obj 0]]$ सिद्ध करायचे आहे.
$\lexists[y][\eq[(y'+a)][\Obj 0]]$ असे मानून व्याघात मिळवू.
$\eq[(b'+a)][\Obj 0]$ असे समजा. $\mathsf{Q}_3$ वापरल्यावर
$\eq[a][\Obj 0] \lor \lexists[y][\eq[a][y']]$ मिळते.
आता दोन बाबी पाहू.

पहिली बाब: $\eq[a][\Obj 0]$ खरे आहे. $\eq[(b'+a)][\Obj 0]$
वरून $\eq[(b' + \Obj 0)][\Obj 0]$ मिळते. $\Th{Q}$ च्या
स्वयंसिद्धक~$\mathsf{Q}_4$ ने $\eq[(b' + \Obj 0)][b']$ मिळते,
म्हणून $\eq[b'][\Obj 0]$ मिळते. पण स्वयंसिद्धक~$\mathsf{Q}_2$ ने
$\eq/[b'][\Obj 0]$ हेसुद्धा मिळते; हा व्याघात आहे.

दुसरी बाब: एखाद्या $c$ साठी $\eq[a][c']$. मग
$\eq[(b' +
c')][\Obj 0]$ मिळते. स्वयंसिद्धक~$\mathsf{Q}_5$ ने
$\eq[(b' + c)'][\Obj 0]$ मिळते; हेसुद्धा स्वयंसिद्धक~$Q_2$
शी विसंगत आहे.
\end{proof}

\begin{lem}
\label{inc:req:min:lem:less-nsucc}
प्रत्येक नैसर्गिक संख्या~$n$ साठी,
  \[  \Th{Q} \Proves
  \lforall[x][(x < \num {n+1} \lif (\eq[x][\Obj 0] \lor \dots \lor
    \eq[x][\num n]))].
  \]
\end{lem}

\begin{proof}
$n$ वर विगमन वापरू. प्रथम $n = 0$ हा आधार टप्पा पाहू.
त्या वेळी कोणत्याही~$a$ साठी $a < \num 1 \lif \eq[a][\Obj 0]$
हे दाखवायचे आहे. $a < \num 1$ असे समजा. मग $<$ च्या
परिभाषक स्वयंसिद्धकाने $\lexists[y][\eq[(y'+a)][\Obj 0']]$
मिळते (कारण $\num 1 \ident
\Obj 0'$).

$b$ कडे हा गुणधर्म आहे, म्हणजे $\eq[(b'+a)][\Obj 0']$,
असे समजा. $\eq[a][\Obj 0]$ दाखवायचे आहे. स्वयंसिद्धक~$\mathsf{Q}_3$
ने $\eq[a][\Obj 0]$ किंवा एखाद्या $c$ साठी
$\eq[a][c']$ मिळते. पहिल्या बाबतीत आणखी काही दाखवायचे नाही.
म्हणून $\eq[a][c']$ असे समजा. मग
$\eq[(b' + c')][\Obj 0']$ मिळते. $\Th{Q}$ च्या स्वयंसिद्धक~$\mathsf{Q}_5$
ने $\eq[(b'+c)'][\Obj 0']$ मिळते. स्वयंसिद्धक~$\mathsf{Q}_1$ ने
$\eq[(b' + c)][\Obj
0]$ मिळते. पण स्वयंसिद्धक~$\mathsf{Q}_8$ नुसार याचा अर्थ
$c < \Obj 0$; याला \ref{inc:req:min:lem:less-zero} ने विरोध होतो.

आता विगमन टप्पा पाहू. $n$ साठीचे विधान गृहीत धरून $n+1$
साठी ते सिद्ध करू. $a < \num {n+2}$ असे समजा. पुन्हा
$\mathsf{Q}_3$ वापरून दोन बाबी मिळतात: $\eq[a][\Obj 0]$, किंवा
एखाद्या $b$ साठी $\eq[a][b']$. पहिल्या बाबतीत
$\eq[a][\Obj 0] \lor \dots \lor
\eq[a][\num{n+1}]$ सहज मिळते. दुसऱ्या बाबतीत $b'
< \num {n+2}$, म्हणजे $b' < \num{n+1}'$. स्वयंसिद्धक~$\mathsf{Q}_8$
नुसार एखाद्या $c$ साठी
$\eq[(c'+b')][\num{n+1}']$. स्वयंसिद्धक $\mathsf{Q}_5$ ने
$\eq[(c'+b)'][\num{n+1}']$. स्वयंसिद्धक~$\mathsf{Q}_1$ ने
$\eq[(c'+b)][\num{n+1}]$; म्हणून स्वयंसिद्धक~$\mathsf{Q}_8$ ने
$b < \num{n+1}$. विगमन गृहीतकाने
$\eq[b][\Obj 0] \lor \dots \lor \eq[b][\num{n}]$.
यावरून तर्कशास्त्राने
$\eq[b'][\Obj 0'] \lor \dots \lor \eq[b'][\num{n}']$
मिळते; आणि $\eq[a][b']$ असल्याने
$\eq[a][\num{1}] \lor \dots \lor \eq[a][\num{n+1}]$
मिळते. उलट दिशेसाठी, उजवीकडील प्रत्येक निर्दिष्ट संख्यांकाशी
समता गृहीत धरल्यास, वरच्या मर्यादेपर्यंतचे उरलेले अंतर
दाखवणारी विशिष्ट संख्या साक्षी म्हणून घ्यावी. आधी सिद्ध
केलेल्या संख्यांकांच्या बेरीजविषयक विधानाने अपेक्षित बेरीज
सिद्ध होते आणि क्रमसंबंधाच्या परिभाषक स्वयंसिद्धकाने तो संख्यांक
वरच्या मर्यादेपेक्षा लहान ठरतो;
आधार टप्प्यातही याच रीतीने उलट दिशा मिळते.
\end{proof}

\begin{lem}
  \label{inc:req:min:lem:trichotomy} प्रत्येक नैसर्गिक संख्या~$m$ साठी,
  \[  \Th{Q} \Proves
  \lforall[y][((y < \num{m} \lor \num{m} < y) \lor \eq[y][\num{m}])].
  \]
\end{lem}

\begin{proof}
$m$ वर विगमन वापरू. प्रथम $m=0$ ही बाब पाहू. $\mathsf{Q}_3$ ने
$\Th{Q} \Proves
\lforall[y][(\eq[y][\Obj 0] \lor \lexists[z][\eq[y][z']])]$.
$a$ कोणताही असू द्या. मग $\eq[a][\Obj 0]$ किंवा एखाद्या~$b$
साठी $\eq[a][b']$. पहिल्या बाबतीत
$(a < \Obj 0 \lor \Obj
0 < a) \lor \eq[a][\Obj 0]$ हेसुद्धा मिळते. पण
$\eq[a][b']$ असेल, तर $\eq$ च्या तर्कशास्त्राने
$\eq[(b' +
\Obj 0)][(a + \Obj 0)]$ मिळते. $\mathsf{Q}_4$ ने
$\eq[(a +
\Obj 0)][a]$, म्हणून $\eq[(b' + \Obj 0)][a]$ आणि त्यामुळे
$\lexists[z][\eq[(z' + \Obj 0)][a]]$ मिळते. $\mathsf{Q}_8$
मधील $<$ च्या व्याख्येनुसार $\Obj 0 < a$. $\Obj 0 < a$
असेल, तर $(\Obj 0 < a \lor a
< \Obj 0) \lor \eq[a][\Obj 0]$ हेसुद्धा मिळते.

आता पुढील विधान गृहीत धरू:
\begin{align*}
  \Th{Q} & \Proves \lforall[y][((y < \num{m} \lor \num{m} < y) \lor
    \eq[y][\num{m}])]
  \intertext{आणि पुढील विधान सिद्ध करायचे आहे:}
  \Th{Q} & \Proves \lforall[y][((y < \num{m+1} \lor \num{m+1} < y) \lor
    \eq[y][\num{m+1}])]
\end{align*}
$a$ कोणताही असू द्या. $\mathsf{Q}_3$ ने $\eq[a][\Obj 0]$ किंवा
एखाद्या~$b$ साठी $\eq[a][b']$. पहिल्या बाबतीत $\mathsf{Q}_4$ ने
$\eq[\num{m}' +
a][\num{m+1}]$ मिळते; म्हणून $\mathsf{Q}_8$ ने
$a < \num{m+1}$.

आता दुसरी बाब, $\eq[a][b']$, पाहू. विगमन गृहीतकानुसार
$(b < \num{m} \lor \num{m} < b) \lor
    \eq[b][\num{m}]$.

पहिला विकल्प $b < \num{m}$ हा $\mathsf{Q}_8$ नुसार
$\lexists[z][\eq[(z' + b)][\num{m}]]$ याच्याशी समतुल्य आहे.
एखाद्या $c$ कडे हा गुणधर्म आहे असे समजा.
$\eq[(c' + b)][\num{m}]$ असेल, तर
$\eq[(c' + b)'][\num{m}']$ हेसुद्धा मिळते. $\mathsf{Q}_5$ ने
$\eq[(c' + b)'][(c' + b')]$ मिळते. म्हणून
$\eq[(c'+b')][\num{m}']$ मिळते. $\num{m}'
\ident \num{m+1}$ हे लक्षात ठेवून $c$ वर अस्तित्ववाचक
सामान्यीकरण केल्यावर
$\lexists[u][\eq[(u' + b')][\num{m+1}]]$ मिळते.
म्हणून $b < \num{m}$ असेल, तर $b' < \num{m+1}$, आणि
त्यामुळे $a < \num{m+1}$.

आता $\num{m} < b$ असे समजा, म्हणजे
$\lexists[z][\eq[(z' +
\num{m})][b]]$. $c$ हा असा~$z$ असेल, म्हणजे
$\eq[(c' +
\num{m})][b]$, असे समजा. तर्कशास्त्राने
$\eq[(c' + \num{m})'][b']$ मिळते. $\mathsf{Q}_5$ ने
$\eq[(c' + \num{m}')][b']$ मिळते. $\eq[a][b']$ आणि
$\num{m}' \ident
\num{m+1}$ असल्याने $\eq[(c' + \num{m+1})][a]$.
$\mathsf{Q}_8$ ने $\num{m+1} < a$.

शेवटी $\eq[b][\num{m}]$ असे समजू. मग तर्कशास्त्राने
$\eq[b'][\num{m}']$, आणि म्हणून
$\eq[a][\num{m+1}]$ मिळते.

अशा प्रकारे $m$ आणि~$b$ यांच्या बाबतीतील प्रत्येक विकल्पापासून
$m+1$ आणि~$a$ यांच्या बाबतीतील अनुरूप विकल्प मिळतो.
\end{proof}

\begin{prop}
\label{inc:req:min:prop:rep-minimization}
$\varphi_g(x, z, y)$ हे $g(x, z)$ चे~$\Th{Q}$ मध्ये निरूपण
करत असेल आणि ते फलन नियमित असेल, तर
\[\varphi_f(z,y) \ident \varphi_g(y, z, \Obj 0) \land \lforall[w][(w < y \lif \lnot
  \varphi_g(w, z, \Obj 0))]
\]
हे सूत्र $f(z) = \umin{x}{[g(x, z) = 0]}$ चे निरूपण करते.
\end{prop}

\begin{proof}
प्रथम $f(n) = m$ असेल, तर
$\Th{Q} \Proves \varphi_f(\num n, \num m)$, म्हणजे पुढील विधान,
हे दाखवू:
\begin{align}
  \Th{Q} & \Proves \varphi_g(\num{m}, \num{n}, \Obj 0) \land \lforall[w][(w <
    \num{m} \lif \lnot \varphi_g(w, \num{n}, \Obj 0))]. \notag
  \intertext{$\varphi_g(x, z, y)$ हे $g(x, z)$ चे निरूपण करते आणि
    $f(n) = m$ असेल तेव्हा $g(m, n) =
    0$; म्हणून}
\Th{Q} & \Proves \varphi_g(\num{m}, \num{n}, \Obj 0). \notag
\intertext{$f(n) = m$ असेल, तर प्रत्येक $k < m$ साठी
  $g(k, n) \neq 0$. निरूपणातील एकमेव-मूल्याची अट आणि
  भिन्न संख्यांकांची असमता वापरल्यावर}
\Th{Q} & \Proves \lnot \varphi_g(\num{k}, \num{n}, \Obj 0). \notag
\intertext{यावरून पुढील विधान मिळते:}
\Th{Q} & \Proves \lforall[w][(w < \num{m} \lif \lnot
  \varphi_g(w, \num{n}, \Obj 0))]. \label{inc:req:min:rep-less}
\end{align}
$m = 0$ असेल, तर हे \ref{inc:req:min:lem:less-zero} वरून;
अन्यथा \ref{inc:req:min:lem:less-nsucc} वरून मिळते.

आता $f(n) = m$ असेल, तर
$\Th{Q} \Proves
\lforall[y][(\varphi_f(\num{n}, y) \lif \eq[y][\num{m}])]$
हे दाखवू. औपचारिक सिद्धता वाचकावर सोडून युक्तिवाद पुन्हा
अनौपचारिकपणे सांगू.

$\varphi_f(\num{n}, b)$ असे समजा. त्यावरून (अ)
$\varphi_g(b, \num{n},
\Obj 0)$ आणि (आ) $\lforall[w][(w < b \lif \lnot \varphi_g(w, \num{n}, \Obj
  0))]$ मिळते. \ref{inc:req:min:lem:trichotomy} नुसार
$(b < \num{m} \lor \num{m} < b)
\lor \eq[b][\num{m}]$. $b < \num{m}$ किंवा $\num{m}
< b$ या दोन्ही शक्यतांपासून व्याघात मिळतो, हे दाखवू.

$\num{m} < b$ असेल, तर (आ) वरून
$\lnot \varphi_g(\num{m}, \num{n}, \Obj 0)$ मिळते. पण
$m = f(n)$ असल्याने $g(m, n) = 0$; आणि $\varphi_g$ हे~$g$
चे निरूपण करत असल्याने $\Th{Q} \Proves
\varphi_g(\num{m}, \num{n}, \Obj 0)$. हा व्याघात आहे.

आता $b < \num{m}$ असे समजा. \ref{inc:req:min:rep-less} नुसार
$\Th{Q} \Proves \lforall[w][(w <
  \num{m} \lif \lnot \varphi_g(w, \num{n}, \Obj 0))]$,
म्हणून $\lnot \varphi_g(b, \num{n}, \Obj 0)$ मिळते.
हे पुन्हा (अ) शी विसंगत आहे. उलट अभिव्यंजन आधी सिद्ध
केलेल्या निरूपक सूत्राच्या निर्दिष्ट मूल्याच्या उदाहरणावरून आणि
समतेच्या प्रतिस्थापन नियमावरून मिळते.
\end{proof}











\section{संगणनक्षम फलने~$\Th{Q}$ मध्ये निरूपणीय असतात}\label{inc:req:crq:sec}

\begin{thm}
प्रत्येक संगणनक्षम फलन~$\Th{Q}$ मध्ये निरूपणीय असते.
\end{thm}

\begin{proof}
संदिग्धता टाळण्यासाठी, चर्च--ट्यूरिंग प्रबंध वापरून फलन
संगणनक्षम असते तर आणि तरच ते सामान्य पुनरावर्ती असते,
असे म्हणू. सामान्य पुनरावर्ती फलने म्हणजे शून्य फलन~$\Zero$,
उत्तरवर्ती फलन~$\Succ$ आणि प्रक्षेपण फलन~$\Proj{n}{i}$
यांपासून संयोजन, आदिम पुनरावर्तन व नियमित लघुतमीकरण वापरून
परिभाषित करता येणारी फलने. \ref{inc:req:pri:lem:prim-rec} नुसार
$f$ आणि~$g$ यांपासून आदिम पुनरावर्तनाने परिभाषित होणारे
कोणतेही फलन~$h$ हे $f$, $g$, $\Zero$,
$\Succ$, $\Proj{n}{i}$, $\Add$, $\Mult$, $\Char{=}$
यांपासून संयोजन व नियमित लघुतमीकरण वापरूनही परिभाषित
करता येते. परिणामी, फलन सामान्य पुनरावर्ती असते तर आणि तरच ते $\Zero$, $\Succ$, $\Proj{n}{i}$, $\Add$,
$\Mult$, $\Char{=}$ यांपासून संयोजन व नियमित लघुतमीकरण
वापरून परिभाषित करता येते.

शिवाय, ही सर्व मूलभूत फलने~$\Th{Q}$ मध्ये निरूपणीय आहेत,
हे आपण सिद्ध केले आहे
(\ref{inc:req:bre:prop:rep-zero} आणि \ref{inc:req:bre:prop:rep-succ} आणि \ref{inc:req:bre:prop:rep-proj} आणि \ref{inc:req:bre:prop:rep-id} आणि \ref{inc:req:bre:prop:rep-add} आणि \ref{inc:req:bre:prop:rep-mult});
आणि निरूपणीय फलनांपासून संयोजन किंवा नियमित लघुतमीकरणाने
परिभाषित होणारे कोणतेही फलन
(\ref{inc:req:cmp:prop:rep-composition},
\ref{inc:req:min:prop:rep-minimization}) हेही निरूपणीय असते.
म्हणून प्रत्येक सामान्य पुनरावर्ती फलन~$\Th{Q}$ मध्ये
निरूपणीय असते.
\end{proof}

\begin{explain}
संगणनक्षम फलनांचा वर्ग म्हणजेच $\Th{Q}$ मध्ये निरूपणीय
फलनांचा वर्ग, असे लक्षणन आपल्याला मिळाले आहे. प्रत्यक्षात
हा पुरावा अधिक व्यापक आहे. निरूपणीयतेच्या व्याख्येवरून
$\Th{Q}$ चा विस्तार असलेल्या कोणत्याही उपपत्तीत (किंवा
ज्यात $\Th{Q}$ चे औपचारिक अर्थनिर्धारण करता येते अशा उपपत्तीत)
संगणनक्षम फलनांचे निरूपण करता येते, हे सहज दिसते. उलट,
ज्या सुसंगत निष्पत्ती व्यवस्थेत भिन्न संख्यांकांची असमता
सिद्ध होते आणि निष्पत्ती ओळखण्याची प्रक्रिया संगणनक्षम
असते, त्या व्यवस्थेत निरूपणीय असलेले प्रत्येक फलन संगणनक्षम
असते. विसंगत व्यवस्थेत प्रत्येक सूत्र सिद्ध होत असल्याने
ही उलटी दिशा तेथे लागू पडत नाही. या दिशेसाठी प्रभावी पुरावा-प्रगणन
पुरेसे आहे. म्हणून, उदाहरणार्थ, या उपपत्ती सुसंगत असतील, तर संगणनक्षम
फलनांचा वर्ग म्हणजे पेआनो अंकगणितात किंवा अगदी
झर्मेलो--फ्रेंकेल संच उपपत्तीत निरूपणीय फलनांचा वर्ग,
असेही लक्षणन करता येते. गोडेल यांनी नोंदवल्याप्रमाणे,
हे काहीसे आश्चर्यकारक आहे. सिद्धतायोग्यतेबाबत प्रश्नांचे
उत्तर कोणती उपपत्ती निवडली आहे यावर फार अवलंबून असते,
हे आपण पुढे पाहू; साधारणपणे स्वयंसिद्धके जितकी बलवान,
तितके अधिक सिद्ध करता येते. पण स्वयंसिद्धकीय उपपत्तींच्या
मोठ्या वर्गात निरूपणीय फलने नेमकी संगणनक्षम फलनेच असतात;
त्या उपपत्ती सुसंगत असतील, भिन्न संख्यांकांची असमता सिद्ध करत असतील
आणि त्यांतील पुराव्यांचे प्रभावी प्रगणन करता येत असेल,
तर अधिक बलवान उपपत्ती अधिक फलनांचे निरूपण करत नाहीत.
\end{explain}












\section{संबंधांचे निरूपण}\label{inc:req:rel:sec}

एखादा \emph{संबंध} निरूपणीय असणे म्हणजे काय, ते आता सांगू.

\begin{defn}
\label{inc:req:rel:defn:representing-relations} नैसर्गिक संख्यांवरील
$R(x_0,\dots,x_k)$ हा संबंध {\em $\Th{Q}$ मध्ये निरूपणीय}
असतो, जर असे सूत्र $\varphi_R(x_0,\dots,x_k)$ असेल की
$R(n_0,\dots,n_k)$ सत्य असेल तेव्हा $\Th{Q}$ हे
$\varphi_R(\num{n_0},\dots,\num{n_k})$ सिद्ध करते, आणि
$R(n_0,\dots,n_k)$ असत्य असेल तेव्हा $\Th{Q}$ हे
$\lnot \varphi_R(\num{n_0},
\dots, \num{n_k})$ सिद्ध करते.
\end{defn}

\begin{thm}
\label{inc:req:rel:thm:representing-rels} संबंध~$\Th{Q}$ मध्ये निरूपणीय
असतो तर आणि तरच तो संगणनक्षम असतो.
\end{thm}

\begin{proof}
प्रथम डावीकडून उजवीकडे जाणारी दिशा पाहू. समजा
$R(x_0,\dots,x_k)$ चे निरूपण $\varphi_R(x_0,\dots,x_k)$ या
सूत्राने होते. $R$ चे संगणन करण्यासाठी पुढील कार्यपद्धती
आहे: $n_0$, \dots,~$n_k$ ही आदाने मिळाल्यावर
$\varphi_R(\num{n_0}, \dots, \num{n_k})$ चा पुरावा आणि
$\lnot \varphi_R(\num{n_0}, \dots, \num{n_k})$ चा पुरावा
यांचा एकाच वेळी शोध घ्या. गृहीतकानुसार यांपैकी एक पुरावा
नक्की सापडेल. पहिला सापडल्यास ``होय'' आणि दुसरा सापडल्यास
``नाही'' असे उत्तर द्या. ही उपपत्ती सुसंगत असल्याने दोन्ही
विरुद्ध पुरावे मिळू शकत नाहीत.

उलट दिशेसाठी $R(x_0, \dots, x_k)$ संगणनक्षम आहे असे समजा.
व्याख्येनुसार याचा अर्थ $\Char{R}(x_0, \dots, x_k)$ हे
फलन संगणनक्षम आहे. \ref{inc:req:int:thm:representable-iff-comp} नुसार
$\Char{R}$ चे एखाद्या सूत्राने, समजा
$\varphi_{\Char{R}}(x_0, \dots, x_k,
y)$ ने, निरूपण होते. $\varphi_R(x_0, \dots, x_k)$ हे
$\varphi_{\Char{R}}(x_0,
\dots, x_k, \num{1})$ सूत्र असू द्या. मग कोणत्याही
$n_0$, \dots,~$n_k$ साठी $R(n_0,
\dots, n_k)$ सत्य असेल, तर $\Char{R}(n_0, \dots, n_k) = 1$.
त्या वेळी $\Th{Q}$ हे
$\varphi_{\Char{R}}(\num{n_0}, \dots, \num{n_k},
\num{1})$ सिद्ध करते; म्हणून $\Th{Q}$ हे
$\varphi_R(\num{n_0}, \dots,
\num{n_k})$ सिद्ध करते. दुसरीकडे $R(n_0, \dots, n_k)$
असत्य असेल, तर $\Char{R}(n_0, \dots, n_k) = 0$.
म्हणून $\Th{Q}$ पुढील विधान सिद्ध करते:
\[\lforall[y][(\varphi_{\Char{R}}(\num{n_0}, \dots, \num{n_k}, y) \lif y =
  \num{0})].
\]
$\Th{Q}$ हे $\eq/[\num{0}][\num{1}]$ सिद्ध करत असल्याने,
$\Th{Q}$ हे $\lnot \varphi_{\Char{R}}(\num{n_0}, \dots, \num{n_k}, \num{1})$
सिद्ध करते; म्हणून $\lnot \varphi_R(\num{n_0}, \dots, \num{n_k})$
हेही सिद्ध करते.
\end{proof}

\begin{prob}
$R$ हा संबंध~$\Th{Q}$ मध्ये निरूपणीय असेल, तर
$\Char{R}$ हेसुद्धा निरूपणीय असते, हे दाखवा.
\end{prob}











\section{अनिर्णेयता}\label{inc:req:und:sec}

एखाद्या उपपत्तीला $\Th{T}$ \emph{अनिर्णेय} म्हणतो, जर $\Th{T}$ च्या
भाषेतील कोणतेही वाक्य~$\varphi$ दिले असता ``$\Th{T}$ हे~$\varphi$
सिद्ध करते का?'' या प्रश्नाचे योग्य उत्तर सांत पायऱ्यांत आणि
नेहमी देणारी संगणकीय पद्धत अस्तित्वात नसेल. म्हणून अंकगणिताच्या
भाषेतील वाक्य~$\varphi$ दिल्यावर $\Th{Q} \Proves \varphi$ आहे की नाही,
हे ठरवणारी संगणकीय पद्धत असेल तर आणि तरच $\Th{Q}$
निर्णेय असेल. हे अधिक नेमके करण्यासाठी विचारू: $\Prov[\Th{Q}](y)$
हा संबंध पुनरावर्ती आहे का? येथे तो~$y$ साठी सत्य
असतो, तर आणि तरच $y$ ही~$\Th{Q}$ मध्ये सिद्ध होणाऱ्या वाक्याची
गोडेल संख्या असते. उत्तर नाही.

\begin{thm}
$\Th{Q}$ अनिर्णेय आहे; म्हणजे पुढील संबंध
\[\Prov[\Th{Q}](y) \defiff \fn{Sent}(y) \land
\lexists[x][\Prf[\Th{Q}](x, y)]
\]
पुनरावर्ती नाही.
\end{thm}

\begin{proof}
तो पुनरावर्ती आहे असे समजा. मग थांबण्याची समस्या अशी
सोडवता येईल: $e$ आणि $n$ दिले असता,
$\cfind{e}(n) \fdefined$ तर आणि तरच असा एखादा~$s$
असतो की $T(e, n, s)$, हे आपल्याला माहीत आहे. येथे $T$ हे
\ref{cmp:rec:nft:thm:kleene-nf} मधील क्लिनीचे विधेय
आहे. $T$ आदिम पुनरावर्ती असल्याने त्याचे~$\Th{Q}$ मध्ये
$\psi_T$ या सूत्राने निरूपण होते; म्हणजे
$\Th{Q}
\Proves \psi_T(\num{e}, \num{n}, \num{s})$ तर आणि तरच
$T(e, n, s)$. जर $\Th{Q}
\Proves \psi_T(\num{e}, \num{n}, \num{s})$ असेल, तर
$ \Th{Q} \Proves
\lexists[y][\psi_T(\num{e}, \num{n}, y)]$ हेसुद्धा मिळते.
असा एकही $s$ नसेल, तर प्रत्येक~$s$ साठी
$\Th{Q} \Proves \lnot \psi_T(\num{e}, \num{n}, \num{s})$.
पण $\Th{Q}$ ही $\omega$-सुसंगत आहे: प्रत्येक~$n \in \Nat$
साठी $\Th{Q}
\Proves \lnot \varphi(\num{n})$ असेल, तर
$\Th{Q}
\Proves/ \lexists[y][\varphi(y)]$. कारण $\Th{Q}$ ची स्वयंसिद्धके
मानक प्रतिमानात~$\Struct{N}$ खरी आहेत. म्हणून
$\Th{Q}
\Proves/ \lexists[y][\psi_T(\num{e}, \num{n}, y)]$.
दुसऱ्या शब्दांत, $\Th{Q}
\Proves \lexists[y][\psi_T(\num{e}, \num{n}, y)]$ तर आणि तरच असा एखादा $s$ असतो की $T(e, n, s)$; म्हणजे तर आणि तरच $\cfind{e}(n) \fdefined$. $e$ आणि~$n$ पासून
$\Gn{\lexists[y][\psi_T(\num{e},
    \num{n}, y)]}$ आपण संगणित करू शकतो; तसे करणारे
$g(e, n)$ हे आदिम पुनरावर्ती फलन असू द्या. मग
\[h(e, n) =
\begin{cases}
1 & \text{जर $\Prov[\Th{Q}](g(e, n))$ असेल}\\
0 & \text{अन्यथा}.
\end{cases}
\]
$\Prov[\Th{Q}]$ पुनरावर्ती असेल, तर यावरून $h$ सुद्धा
पुनरावर्ती ठरेल. पण~$h$ हे
\ref{cmp:rec:hlt:thm:halting-problem} नुसार पुनरावर्ती नाही.
म्हणून $\Prov[\Th{Q}]$ हाही पुनरावर्ती नाही.
\end{proof}

\begin{cor}
प्रथम-क्रम तर्कशास्त्र अनिर्णेय आहे.
\end{cor}

\begin{proof}
प्रथम-क्रम तर्कशास्त्र निर्णेय असते, तर~$\Th{Q}$ मधील
सिद्धतायोग्यताही निर्णेय असती: कारण $\Th{Q} \Proves \varphi$
तर आणि तरच $\Proves \mathsf{T} \lif \varphi$; येथे $\mathsf{T}$ हा
$\Th{Q}$ च्या सांत स्वयंसिद्धकांचा संयोग आहे.
\end{proof}











\section{\texorpdfstring{$\Sigma_1$}{सिग्मा-१} संपूर्णता}\label{inc:inp:s1c:sec}

$\Th{Q}$ आणि तिच्या सुसंगत, स्वयंसिद्धकीकरण करता येणाऱ्या
विस्तारांची अपूर्णता असूनही, संख्यांकांविषयी अनेक मूलभूत तथ्ये
$\Th{Q}$ सिद्ध करते, हे आपण पाहिले आहे. ही क्षमता प्रत्यक्षात
आणखी बरीच व्यापक आहे. $\Th{Q}$ मध्ये काय सिद्ध करता येते
याची व्याप्ती समजण्यासाठी $\Delta_0$, $\Sigma_1$ आणि
$\Pi_1$ सूत्रे यांच्या संकल्पना मांडू. स्थूलमानाने,
$\Sigma_1$ सूत्र हे $\lexists[x][\psi(x)]$ या रूपाचे
असते; येथे $\psi$ फक्त विधानीय संयोजके आणि परिबद्ध संख्यापक
वापरून तयार केलेले असते. $\varphi$ हे $\Struct{N}$ मध्ये सत्य
असलेले $\Sigma_1$ वाक्य असेल, तर
$\Th{Q} \Proves \varphi$, हे आपण दाखवू
(\ref{inc:inp:s1c:thm:sigma1-completeness}).

\begin{defn}
\label{inc:inp:s1c:defn:bd-quant}
\emph{परिबद्ध अस्तित्ववाची सूत्र} हे
$\lexists[x][(x < t \land \varphi(x))]$ या रूपाचे असते; येथे $t$
हे असे कोणतेही पद आहे की त्यात संबंधित संख्यापकाचे चल येत नाही. परंपरेने ते
$\bexists{x < t}{\varphi(x)}$ असे लिहितो.

\emph{परिबद्ध वैश्विक सूत्र} हे
$\lforall[x][(x < t \lif \varphi(x))]$ या रूपाचे असते; येथेही
$t$ या पदात संबंधित संख्यापकाचे चल नसते. परंपरेने ते
$\bforall{x < t}{\varphi(x)}$ असे लिहितो.
\end{defn}

\begin{defn}
\label{inc:inp:s1c:defn:delta0-sigma1-pi1-frm}
सूत्र $\psi$ हे $\Delta_0$ असते, जर ते आण्विक
सूत्रांपासून केवळ विधानीय संयोजके आणि परिबद्ध
संख्यकीकरण वापरून तयार झाले असेल.

सूत्र $\varphi$ हे $\Sigma_1$ असते, जर
$\varphi \ident \lexists[x][\psi(x)]$ आणि $\psi$ हे $\Delta_0$
असेल.

सूत्र $\varphi$ हे $\Pi_1$ असते, जर
$\varphi \ident \lforall[x][\psi(x)]$ आणि $\psi$ हे $\Delta_0$
असेल.
\end{defn}

\begin{lem}
\label{inc:inp:s1c:lem:q-proves-clterm-id} समजा $t$ हे बंद पद असून
$\Value{t}{N} = n$. तर $\Th{Q} \Proves \eq[t][\num n]$.
\end{lem}

\begin{proof}
$t$ च्या गुंतागुंतीवर विगमन करू. आधार प्रकरणात
$\Value{\Obj 0}{N} = 0$, आणि
$\num 0 \ident \Obj 0$ असल्यामुळे
$\Th{Q} \Proves \eq[\Obj 0][\num 0]$.

विगमनाच्या टप्प्यासाठी $t_1$ आणि $t_2$ ही पदे अशी घ्या की
$\Value{t_1}{N} = n_1$, $\Value{t_2}{N} = n_2$,
$\Th{Q} \Proves \eq[t_1][\num n_1]$, आणि
$\Th{Q} \Proves \eq[t_2][\num n_2]$.

मग $\Value{(t_1')}{N} = n_1 + 1$. विगमन गृहीतक आणि
$\eq[\num{n_1}'][\num{n_1}']$ या सूत्र वर
प्रथम-क्रम एकरूपतेचे नियम लावल्यास
$\Th{Q} \Proves \eq[t_1'][{\num n_1}']$ मिळते.
संख्यांकांच्या व्याख्येने
$\Th{Q} \Proves \eq[t_1'][\num{n_1 + 1}]$ मिळते.

बेरीजेसाठी,
\[      \Value{(t_1 + t_2)}{N}
    = \Value{t_1}{N} + \Value{t_2}{N}
    = n_1 + n_2.
\]
विगमन गृहीतक आणि एकरूपतेच्या नियमांनी प्रथम
$\Th{Q} \Proves \eq[t_1 + t_2][\num{n_1} + t_2]$,
आणि ते नियम पुन्हा लावल्यास
$\Th{Q} \Proves \eq[t_1 + t_2][\num{n_1} + \num{n_2}]$
मिळते. \ref{inc:req:bre:lem:q-proves-add} नुसार
$\Th{Q} \Proves \eq[\num{n_1} + \num{n_2}][\num{n_1 + n_2}]$;
म्हणून $\Th{Q} \Proves \eq[t_1 + t_2][\num{n_1 + n_2}]$.

$\times$ साठीही
\ref{inc:req:bre:lem:q-proves-mult} वापरून असाच
युक्तिवाद करता येतो.

अंकगणितातील बंद पदांचे हे सर्व प्रकार आहेत. म्हणून
$\Value{t}{N} = n$ असलेल्या प्रत्येक बंद पद~$t$ साठी
$\Th{Q} \Proves \eq[t][\num n]$.
\end{proof}

\begin{prob}
\ref{inc:inp:s1c:lem:q-proves-clterm-id} मधील $\times$ संबंधीचा
भाग सविस्तर सिद्ध करा.
\end{prob}

\begin{lem}
\label{inc:inp:s1c:lem:atomic-completeness}
समजा $t_1$ आणि $t_2$ ही बंद पदे आहेत. तर
\begin{enumerate}
\item $\Value{t_1}{N} = \Value{t_2}{N}$ असेल, तर
    $\Th{Q} \Proves \eq[t_1][t_2]$.
\item $\Value{t_1}{N} \neq \Value{t_2}{N}$ असेल, तर
    $\Th{Q} \Proves \eq/[t_1][t_2]$.
\item $\Value{t_1}{N} < \Value{t_2}{N}$ असेल, तर
    $\Th{Q} \Proves t_1 < t_2$.
\item $\Value{t_2}{N} \leq \Value{t_1}{N}$ असेल, तर
    $\Th{Q} \Proves \lnot(t_1 < t_2)$.
\end{enumerate}
\end{lem}

\begin{proof}
दिलेले $t_1$ आणि $t_2$ घेऊन
$n = \Value{t_1}{N}$ आणि $m = \Value{t_2}{N}$
निश्चित करू.

समजा $\varphi \ident t_1 = t_2$.
\ref{inc:inp:s1c:lem:q-proves-clterm-id} नुसार
$\Th{Q} \Proves \eq[t_1][\num n]$ आणि
$\Th{Q} \Proves \eq[t_2][\num m]$.
$n = m$ असेल, तर
$\Th{Q} \Proves \eq[\num n][\num m]$;
म्हणून एकरूपतेच्या संक्रमणशीलतेने
$\Th{Q} \Proves \eq[t_1][t_2]$.
$n \neq m$ असेल, तर
$\Th{Q} \Proves \eq/[\num n][\num m]$;
एकरूपतेचे नियम पुन्हा वापरल्यास
$\Th{Q} \Proves \eq/[t_1][t_2]$.

आता $\varphi \ident t_1 < t_2$ असे घेऊ. दोन्ही प्रकरणांत
$\mathsf{Q}_8$ या स्वयंसिद्धकाचा आधार घेऊ. त्यानुसार सर्व
$x,y$ साठी
$x < y \liff \lexists[z][\eq[z' + x][y]]$.

समजा $\Sat{N}{t_1 < t_2}$. मग असा $k \in \Nat$
असतो की $n + k + 1 = m$.
\ref{inc:inp:s1c:lem:q-proves-clterm-id} नुसार
$\Th{Q} \Proves \eq[t_1][\num n]$ आणि
$\Th{Q} \Proves \eq[t_2][\num m]$.
या प्रमेयाच्या पहिल्या भागाने निश्चित संख्यांकांची बेरीज
काढल्यास
$\Th{Q} \Proves \eq[{\num k}' + \num n][\num m]$.
एकरूपतेच्या संक्रमणशीलतेने मग
$\Th{Q} \Proves \eq[{\num k}' + t_1][t_2]$;
म्हणून
$\Th{Q} \Proves \lexists[z][\eq[z' + t_1][t_2]]$.
$\mathsf{Q}_8$ ची उजवीकडून डावीकडे दिशा वापरल्यास
$\Th{Q} \Proves t_1 < t_2$.

याउलट $\Sat/{N}{t_1 < t_2}$, म्हणजे $m \leq n$,
असे समजा.

$\Th{Q}$ मध्ये काम करून $t_1 < t_2$ असे गृहीत धरा.
$\mathsf{Q}_8$ ची डावीकडून उजवीकडे दिशा लावल्यास असा $z$
मिळतो की $\eq[z' + t_1][t_2]$.
$\Th{Q} \Proves \eq[t_1][\num n]$ आणि
$\Th{Q} \Proves \eq[t_2][\num m]$ असल्यामुळे
$\eq[z' + \num n][\num m]$.

$m$ वर बाह्य विगमन करू: प्रत्येक पायरीवर $\mathsf{Q}_5$
वापरून बेरीज उत्तरवर्तीच्या रूपात लिहू आणि उत्तरवर्तीचे
एकास-एक असणे वापरून दोन्ही बाजूंचे एक-एक उत्तरवर्ती
रद्द करू. अशा रीतीने
$\eq[z' + \num{n - m}][\Obj 0]$ मिळते.
$m = n$ असेल, तर शून्य जोडण्याच्या स्वयंसिद्धकाने
$\eq[z'][\Obj 0]$ मिळते; पण $\mathsf{Q}_2$ नुसार उत्तरवर्ती
शून्य नसतो, म्हणून विरोधाभास मिळतो.
$m < n$ असेल, तर $\mathsf{Q}_5$ पुन्हा वापरून
$\eq[(z' + \num{n - m - 1})'][\Obj 0]$ मिळते;
$\mathsf{Q}_2$ मुळे पुन्हा विरोधाभास मिळतो.
म्हणून $\Th{Q} \Proves \lnot(t_1 < t_2)$.
\end{proof}

\begin{lem}
\label{inc:inp:s1c:lem:bounded-quant-equiv}
समजा $\varphi$ हे सूत्र, $t$ हे बंद पद, आणि
$k=\Value{t}{N}$. तर
\begin{enumerate}
\item $\Th{Q} \Proves \bforall{x<t}{\varphi(x)}$ तर आणि तरच $\Th{Q} \Proves
    \varphi(\num 0) \land \dots \land \varphi(\num{k-1})$.
\item $\Th{Q} \Proves \bexists{x<t}{\varphi(x)}$ तर आणि तरच $\Th{Q} \Proves
    \varphi(\num 0) \lor \dots \lor \varphi(\num{k-1})$.
\end{enumerate}
\end{lem}

\begin{proof}
    परिबद्ध वैश्विक संख्यापकाचे प्रकरण सिद्ध करू.
    $\Value{t}{N} = 0$ असेल, तर
    \ref{inc:req:min:lem:less-zero} नुसार
    $x<\num 0$ अशी कोणतीही संख्या नसते; म्हणून
    समतुल्यतेची डावी बाजू $\Th{Q}$ मध्ये सिद्ध होते.
    उजवीकडील रिक्त संयोगाला $\ltrue$ मानतो;
    तेही $\Th{Q}$ मध्ये सिद्ध होते.

    म्हणून आता एखाद्या नैसर्गिक संख्या~$k$ साठी
    $\Value{t}{N} = k+1$ असे समजू.
    \ref{inc:inp:s1c:lem:q-proves-clterm-id} नुसार
    $\bforall{x<\num{k+1}}{\varphi(x)}$ या रूपाच्या
    सूत्र साठी काम करणे पुरेसे आहे.

    समजा
    $\Th{Q} \Proves \bforall{x<\num{k+1}}{\varphi(x)}$,
    आणि $n \leq k$ घेऊ.
    \ref{inc:inp:s1c:lem:atomic-completeness} नुसार
    $\Th{Q} \Proves \num n < \num{k+1}$;
    म्हणून तर्कनियमांनी
    $\Th{Q} \Proves \varphi(\num n)$.
    प्रत्येक $n \leq k$ साठी हे $k+1$ वेळा
    लावल्यास अपेक्षेप्रमाणे
    $\Th{Q} \Proves \varphi(\num 0)
    \land \dots \land \varphi(\num k)$ मिळते.

    उलट दिशेसाठी समजा
    $\Th{Q} \Proves
    \varphi(\num 0) \land \dots \land \varphi(\num k)$.
    $\Th{Q}$ मध्ये $x < \num{k+1}$ असे समजू.
    \ref{inc:req:min:lem:less-nsucc} नुसार
    $x = \num 0 \lor \dots \lor x = \num k$.
    मग तर्कनियमांनी $\varphi(x)$ मिळते आणि त्यावरून
    $\bforall{x<\num{k+1}}{\varphi(x)}$ हे वैश्विक विधान
    मिळते.

    परिबद्ध अस्तित्ववाची संख्यकीकरण असलेल्या
    सूत्रे साठी समतुल्यतेचा पुरावाही असाच आहे.
\end{proof}

\begin{prob}
\ref{inc:inp:s1c:lem:bounded-quant-equiv} मधील अस्तित्ववाची
प्रकरणाचा सविस्तर पुरावा द्या.
\end{prob}

\begin{lem}
\label{inc:inp:s1c:lem:delta0-completeness}
$\varphi$ हे $\Struct{N}$ मध्ये सत्य असलेले
$\Delta_0$ वाक्य असेल, तर
$\Th{Q} \Proves \varphi$.
\end{lem}

\begin{proof}
सूत्राच्या गुंतागुंतीवर विगमन करू.

आधार प्रकरण \ref{inc:inp:s1c:lem:atomic-completeness} ने
मिळते; आता विगमनाचा टप्पा पाहू. सोयीसाठी,
निषेधाच्या प्रकरणाची विभागणी निषिद्ध सूत्र
च्या रचनेनुसार करू. इतर विधानीय संयोजक पुढील
प्रकरणांमधून संक्षेपाने व्यक्त करता येतात.

\begin{enumerate}
\item $(\varphi \land \psi)$ हे $\Struct{N}$ मध्ये सत्य
असेल, तर $\varphi$ आणि $\psi$ दोन्ही $\Struct{N}$ मध्ये सत्य आहेत.
विगमन गृहीतकाने
$\Th{Q} \Proves \varphi$ आणि
$\Th{Q} \Proves \psi$; म्हणून तर्कनियमांनी
$\Th{Q} \Proves (\varphi \land \psi)$.

\item $\lnot (\varphi \land \psi)$ हे $\Struct{N}$
मध्ये सत्य असेल, तर $\lnot \varphi$ किंवा
$\lnot \psi$ हे $\Struct{N}$ मध्ये सत्य आहे. सामान्यत्वाला बाधा
न आणता पहिले सत्य आहे असे धरू.
विगमन गृहीतकाने $\Th{Q} \Proves \lnot \varphi$;
म्हणून तर्कनियमांनी
$\Th{Q} \Proves \lnot (\varphi \land \psi)$.

\item $(\varphi \lor \psi)$ हे $\Struct{N}$ मध्ये सत्य
असेल, तर $\varphi$ हे $\Struct{N}$ मध्ये सत्य आहे किंवा
$\psi$ हे $\Struct{N}$ मध्ये सत्य आहे.
सामान्यत्वाला बाधा न आणता पहिले सत्य आहे
असे धरू. विगमन गृहीतकाने
$\Th{Q} \Proves \varphi$;
म्हणून तर्कनियमांनी
$\Th{Q} \Proves (\varphi \lor \psi)$.

\item $\lnot(\varphi \lor \psi)$ हे $\Struct{N}$
मध्ये सत्य असेल, तर $\lnot \varphi$ आणि
$\lnot \psi$ दोन्ही $\Struct{N}$ मध्ये सत्य आहेत.
विगमन गृहीतकाने
$\Th{Q} \Proves \lnot \varphi$ आणि
$\Th{Q} \Proves \lnot \psi$.
त्यामुळे तर्कनियमांनी
$\Th{Q} \Proves \lnot(\varphi \lor \psi)$.

\item $\bforall{x<t}{\varphi(x)}$ हे
$\Struct{N}$ मध्ये सत्य आहे असे समजा;
येथे $t$ हे बंद पद आणि $k=\Value{t}{N}$.
सर्व $n < \Value{t}{N}$ साठी
$\varphi(\num n)$ हे $\Struct{N}$ मध्ये सत्य
असेल, तर विगमन गृहीतक आणि तर्कनियमांनी
$\Th{Q} \Proves
\varphi(\num 0) \land \dots \land \varphi(\num{k-1})$.
\ref{inc:inp:s1c:lem:bounded-quant-equiv} नुसार
$\Th{Q} \Proves \bforall{x<t}{\varphi(x)}$.

\item परिबद्ध अस्तित्ववाची संख्यापकाचे
प्रकरणही वैश्विक संख्यापकासारखेच आहे:
येथे $\bexists{x < t}{\varphi(x)}$ या रूपाचे
वाक्य असते.

\item $\lnot \bforall{x<t}{\varphi(x)}$ हे
$\Struct{N}$ मध्ये सत्य आहे असे समजा;
येथे $t$ हे बंद पद. हे वाक्य
$\bexists{x<t}{\lnot \varphi(x)}$ या वाक्य
शी समतुल्य आहे आणि ही समतुल्यता
$\Th{Q}$ मध्येही सिद्ध करता येते. म्हणून
परिबद्ध अस्तित्ववाची संख्यापकाचा युक्तिवाद
लागू करता येतो.

\item त्याचप्रमाणे, $t$ हे बंद पद असताना
$\lnot \bexists{x<t}{\varphi(x)}$ हे
$\Struct{N}$ मध्ये सत्य आहे असे समजा.
हे वाक्य $\Th{Q}$ मध्ये
$\bforall{x<t}{\lnot\varphi(x)}$ शी समतुल्य
आहे; म्हणून परिबद्ध वैश्विक संख्यापकाचा
युक्तिवाद लागू करता येतो.

\item अखेरीस $\lnot \varphi$ हे
$\Struct{N}$ मध्ये सत्य आहे असे समजा.
उरलेली प्रकरणे अशी: $\varphi$ आण्विक आहे,
किंवा एखाद्या $\Delta_0$ वाक्य $\psi$
साठी $\lnot \varphi \ident \lnot\lnot \psi$.
$\varphi$ आण्विक असेल, तर
\ref{inc:inp:s1c:lem:atomic-completeness} नुसार
$\Th{Q} \Proves \lnot \varphi$.
$\lnot \varphi \ident \lnot\lnot \psi$ असेल,
तर तर्कनियमांनी ते $\Th{Q}$ मध्ये
$\psi$ शी सिद्धतया समतुल्य आहे.
$\lnot \varphi$ हे $\Struct{N}$ मध्ये सत्य
असल्याने तेही $\Struct{N}$ मध्ये सत्य आहे.
म्हणून विगमन गृहीतकाने
$\Th{Q} \Proves \lnot \varphi$.
\end{enumerate}
\end{proof}

\begin{prob}
\ref{inc:inp:s1c:lem:delta0-completeness} मधील
अस्तित्ववाची प्रकरणाचा सविस्तर पुरावा द्या.
\end{prob}

\begin{thm}
\label{inc:inp:s1c:thm:sigma1-completeness}
$\varphi$ हे $\Struct{N}$ मध्ये सत्य असलेले
$\Sigma_1$ वाक्य असेल, तर
$\Th{Q} \Proves \varphi$.
\end{thm}

\begin{proof}
$\lexists[x][\varphi(x)]$ हे $\Struct{N}$ मध्ये
सत्य असलेले $\Sigma_1$ वाक्य असेल,
तर अशी नैसर्गिक संख्या~$n$ आणि चल-मूल्यन~$s$
असतात की $s(x) = n$ आणि
$\Sat{N}{\varphi(x)}[s]$.
पूर्ति-संबंधाच्या नेहमीच्या गुणधर्मांनुसार
$\Sat{N}{\varphi(\num n)}$.
परंतु $\varphi(\num n)$ हे बंद $\Delta_0$ सूत्र
आहे; म्हणून \ref{inc:inp:s1c:lem:delta0-completeness} ने
$\Th{Q} \Proves \varphi(\num n)$.
मग तर्कनियमांनी
$\Th{Q} \Proves \lexists[x][\varphi(x)]$.
\end{proof}



\chapter{उपपत्ती आणि संगणनक्षमता}\label{inc:tcp:chap}
\begin{editorial}
  हे प्रकरण संगणनक्षमता सिद्धांतावरील प्रकरणातील सामग्रीवर
  अवलंबून आहे; ते शिकवले नसेल, तर हे प्रकरण वगळता येईल.
  सध्या हा जेरेमी अविगॅड यांच्या टिपणांवरून केलेला प्राथमिक
  रूपांतरित मसुदा आहे. त्याचे अद्याप पुनरावलोकन झालेले नाही
  आणि सराव-प्रश्नही दिलेले नाहीत.
\end{editorial}









\section{परिचय}\label{inc:tcp:int:sec}

\begin{editorial}
हा विभाग पुन्हा लिहिणे आवश्यक आहे.
\end{editorial}

आपल्याकडे पुढील गोष्टी उपलब्ध आहेत:
\begin{enumerate}
\item एखादे फलन $\Th{Q}$ मध्ये निरूपणीय असणे म्हणजे काय,
  याची व्याख्या (\ref{inc:req:int:defn:representable-fn})
\item एखादा संबंध $\Th{Q}$ मध्ये निरूपणीय असणे म्हणजे काय,
  याची व्याख्या (\ref{inc:req:rel:defn:representing-relations})
\item $\Th{Q}$ मध्ये निरूपणीय फलने नेमकी संगणनक्षम
  फलनेच असतात, असे सांगणारे प्रमेय
  (\ref{inc:req:int:thm:representable-iff-comp})
\item $\Th{Q}$ मध्ये निरूपणीय संबंध नेमके संगणनक्षम
  संबंधच असतात, असे सांगणारे प्रमेय
  (\ref{inc:req:rel:thm:representing-rels})
\end{enumerate}

{ \em उपपत्ती} म्हणजे तार्किक निष्पन्नतेखाली बंद असलेला वाक्यसंच.
म्हणजे $T$ ने $\varphi$ सिद्ध केले, तर $\varphi$ हे $T$ मध्ये
असते. उपपत्तीचा विचार वाक्यांचा आणि त्यांच्यापासून
निष्पन्न होणाऱ्या सर्व वाक्यांचा संग्रह म्हणून करणे
उपयुक्त ठरेल. यापुढे $\Th{Q}$ हे चिन्ह
\ref{inc:req:int:sec} मधील आठ स्वयंसिद्धकांपासून
निष्पन्न होणाऱ्या सर्व वाक्यांच्या { \em उपपत्ती}साठी
वापरू. $\Th{Q}$ मधील सूत्रांना संख्यांनी
सांकेतिक करता येते, हे लक्षात ठेवा. $\varphi$ असे
सूत्र असेल, तर $\varphi$ ला सांकेतिक करणारी
संख्या $\Gn{\varphi}$ ने दर्शवू. या सांकेतीकरणाच्या
आधारे आता सूत्रांचे विविध संच संगणनक्षम आहेत की
नाहीत, हा प्रश्न विचारता येतो.












\section[Q हा संगणनक्षमपणे प्रगणनीय-संपूर्ण आहे]{$\Th{Q}$ हा संगणनक्षमपणे प्रगणनीय-संपूर्ण आहे}\label{inc:tcp:qce:sec}


\begin{thm}
$\Th{Q}$ हा संगणनक्षमपणे प्रगणनीय आहे, पण निर्णेय नाही.
खरे तर तो संपूर्ण संगणनक्षमपणे प्रगणनीय संच आहे.
\end{thm}

\begin{proof}
$\Th{Q}$ हा संगणनक्षमपणे प्रगणनीय आहे हे पाहणे अवघड नाही:
तो अशा वाक्यांच्या संकेतांकांचा संच आहे की
$y$ हा वाक्याचा संकेतांक असेल, तर $\Th{Q}$ मध्ये
$y$ ने सांकेतिक केलेल्या वाक्याची एखादी सिद्धता $x$ अस्तित्वात असते:
\[Q = \Setabs{y}{\lexists[x][\Prf[\Th{Q}](x,y)]}.
\]
पण $\Prf[\Th{Q}](x,y)$ संगणनक्षम आहे
(खरे तर आदिम पुनरावर्ती आहे), आणि वरच्या रूपात
लिहिता येणारा कोणताही संच संगणनक्षमपणे प्रगणनीय असतो.

तो संपूर्ण संगणनक्षमपणे प्रगणनीय संच आहे, असे म्हणणे
$K \leq_m Q$ असे म्हणण्यास समतुल्य आहे. येथे
$K = \Setabs{x}{\cfind{x}(x) \fdefined}$.
म्हणून $K$ चे~$\Th{Q}$ कडे न्यूनीकरण दाखवू.
क्लिनीचे विधेय $T(e,x,s)$ आदिम पुनरावर्ती असल्यामुळे
त्याचे $\Th{Q}$ मध्ये, समजा $\varphi_T$ या सूत्राने, निरूपण
होते. मग प्रत्येक $x$ साठी
\begin{align*}
x \in K & \lif \lexists[s][T(x,x,s)] \\
& \lif \lexists[s][(\Th{Q} \Proves \varphi_T(\num x,\num x, \num s))]
  \\
& \lif \Th{Q} \Proves \lexists[s][\varphi_T(\num x, \num x, s)].
\end{align*}
उलट $\Th{Q} \Proves \lexists[s][\varphi_T(\num x, \num x, s)]$
असेल, तर या सिद्धांताची स्वयंसिद्धके मानक नैसर्गिक-संख्या
प्रतिमानात सत्य असल्याने त्या प्रतिमानात हे अस्तित्ववाचक
वाक्य सत्य असते. म्हणून काही नैसर्गिक संख्या~$n$ साठी
$\varphi_T(\num x, \num x, \num n)$ सत्य असते.
$T(x,x,n)$ असत्य असेल, तर $\varphi_T$ हे $T$ चे
निरूपण करत असल्याने $\Th{Q}$ ही
$\lnot \varphi_T(\num x, \num x, \num n)$ सिद्ध करेल.
मात्र ते मानक प्रतिमानात असत्य सूत्र ठरेल आणि
$\Th{Q}$ ते सिद्ध करेल, हा विरोधाभास आहे.
म्हणून $T(x,x,n)$ सत्य असले पाहिजे; त्यावरून
$\cfind{x}(x) \fdefined$.

थोडक्यात, प्रत्येक $x$ साठी $x$ हा $K$ मध्ये
असतो तर आणि तरच $\Th{Q}$ हे
$\lexists[s][\varphi_T(\num x,\num x,s)]$ सिद्ध करते.
म्हणून $x$ पासून
$\lexists[s][\varphi_T(\num x,
  \num x, s)]$ या वाक्याचा संकेतांक देणारे
$f$ हे फलन $K$ चे~$\Th{Q}$ कडे न्यूनीकरण करते.
संख्यांकाचे प्रतिस्थापन आणि सूत्राचे सांकेतीकरण
संगणनक्षम असल्याने हे न्यूनीकरण संगणनक्षम आहे.
विकर्ण थांबण्याचा संच आधीच अनिर्णेय असल्याने
हे न्यूनीकरण या संचाचाही अनिर्णेयपणा दाखवते.
\end{proof}












\section[Q चा ओमेगा-सुसंगत विस्तार अनिर्णेय असतो]{$\Th{Q}$ चा $\omega$-सुसंगत विस्तार अनिर्णेय असतो}\label{inc:tcp:oqn:sec}

\begin{explain}
$\Th{Q}$ हा c.e.-संपूर्ण आहे या सिद्धतेत,
$\Th{Q}$ मध्ये सिद्ध होणारे प्रत्येक वाक्य नैसर्गिक संख्यांबाबत
``सत्य'' असते या तथ्याचा आधार घेतला होता. पुढील व्याख्या आणि
प्रमेय, वरील सिद्धतेत आवश्यक असलेले ``सत्य''चे नेमके पैलू
ओळखून, त्या प्रमेयाचे सबलीकरण करतात. मात्र या प्रमेयावर फार
रेंगाळू नका, कारण लवकरच आपण त्याहून अधिक सबल प्रमेय सिद्ध करू.
ते मुख्यतः ऐतिहासिक कारणासाठी समाविष्ट केले आहे: गोडेल यांच्या
मूळ लेखात $\omega$-सुसंगततेची संकल्पना वापरली होती; पण लगेचच
$\omega$-सुसंगततेच्या जागी साधी सुसंगतता घालून त्यांचा निकाल
अधिक सबल करण्यात आला.
\end{explain}

\begin{defn}
\label{inc:tcp:oqn:thm:oconsis-q}
उपपत्ती $\Th{T}$ ही $\omega$-सुसंगत असते, जर पुढील अट पूर्ण होत असेल:
$\lexists[x][\varphi(x)]$ हे कोणतेही वाक्य असो. $\Th{T}$ कडून
$\lnot \varphi(\num 0)$, $\lnot \varphi(\num 1)$,
$\lnot \varphi(\num 2)$, \dots ही सर्व वाक्ये सिद्ध होत असतील,
तर $\Th{T}$ कडून $\lexists[x][\varphi(x)]$ सिद्ध होत नाही.
\end{defn}

\begin{thm}
\label{inc:tcp:oqn:thm:oconsis-q-undec}
  $\Th{T}$ ही कोणतीही $\omega$-सुसंगत उपपत्ती असेल आणि तिच्यात
  $\Th{Q}$ समाविष्ट असेल, तर $\Th{T}$ अनिर्णेय असते.
\end{thm}

\begin{proof}
$\Th{T}$ मध्ये $\Th{Q}$ समाविष्ट असल्यास संगणनक्षम
फलनांचे आणि संबंधांचे $\Th{T}$ मध्ये निरूपण होते.
आधीची सिद्धता किंचित बदलली तरी पुरेसे आहे. पूर्वीप्रमाणे,
$x \in K$ असेल, तर $\Th{T}$ हे $\lexists[s][\varphi_T(\num
  x,\num x, s)]$ सिद्ध करते. उलट, $\Th{T}$ हे
$\lexists[s][\varphi_T(\num x, \num x, s)]$ सिद्ध करते असे धरू.
मग $x$ हा $K$ मध्ये असलाच पाहिजे: अन्यथा यंत्र $x$ चे
$x$ या आदानावरील थांबणारे संगणन अस्तित्वात नसते. $\varphi_T$ हे
क्लिनीच्या $T$ संबंधाचे निरूपण करत असल्याने $\Th{T}$ कडून
$\lnot \varphi_T(\num x, \num x, \num 0)$,
$\lnot \varphi_T(\num x, \num x, \num 1)$, \dots
सिद्ध होतील. यामुळे $\Th{T}$ ही $\omega$-विसंगत ठरेल.
\end{proof}












\section[Q चे सुसंगत विस्तार अनिर्णेय असतात]{$\Th{Q}$ चे सुसंगत विस्तार अनिर्णेय असतात}\label{inc:tcp:cqn:sec}

\begin{explain}
कोणत्याही सूत्र~$\varphi$ साठी $\varphi$ आणि $\lnot \varphi$ ही दोन्ही सूत्रे
सिद्ध होत नसतील, तर उपपत्ती \emph{सुसंगत} असते, हे लक्षात ठेवा.
व्याघातापासून कोणतेही वाक्य निष्पन्न होत असल्यामुळे विसंगत उपपत्ती
क्षुल्लक असते: तिच्यात प्रत्येक वाक्य सिद्ध होते.
उपपत्ती $\omega$-सुसंगत असेल, तर ती सुसंगतही असते, हे स्पष्ट आहे.
पण सुसंगतता ही अधिक दुर्बल अट आहे (म्हणजे सुसंगत असलेल्या, पण
$\omega$-सुसंगत नसलेल्या उपपत्ती आहेत). त्यामुळे
\ref{inc:tcp:oqn:thm:oconsis-q-undec} मधील गृहीतक साध्या सुसंगततेपर्यंत
शिथिल करून अधिक सबल प्रमेय मिळवता येते.
\end{explain}

\begin{lem}
``सार्वत्रिक संगणनक्षम संबंध'' अस्तित्वात नाही. म्हणजे पुढील गुणधर्म
असलेला कोणताही द्विस्थानी संगणनक्षम संबंध $R(x,y)$ नाही:
$S(y)$ हा कोणताही एकस्थानी संगणनक्षम संबंध असेल, तर असा काही $k$
असतो की प्रत्येक $y$ साठी $S(y)$ सत्य असते तर आणि तरच
$R(k,y)$ सत्य असते.
\end{lem}

\begin{proof}
$R(x,y)$ हा सार्वत्रिक संगणनक्षम संबंध आहे असे समजू.
$S(y)$ हा $\lnot R(y,y)$ हा संबंध घेऊ. $S(y)$ संगणनक्षम असल्याने
काही $k$ साठी $S(y)$ हे $R(k,y)$ शी सममूल्य असेल.
पण मग $S(k)$ हे $R(k,k)$ आणि $\lnot R(k,k)$ या
दोन्हींशी सममूल्य ठरेल; हा व्याघात आहे.
\end{proof}


\begin{thm}
$\Th{T}$ ही $\Th{Q}$ समाविष्ट असलेली कोणतीही सुसंगत उपपत्ती
असेल, तर $\Th{T}$ अनिर्णेय असते.
\end{thm}


\begin{proof}
$\Th{T}$ हा $\Th{Q}$ चा सुसंगत, निर्णेय विस्तार आहे असे समजू.
$\Th{T}$ वापरून सार्वत्रिक संगणनक्षम संबंध परिभाषित करू
आणि त्यामुळे व्याघात मिळवू.

$R(x,y)$ पुढील अटीवर, आणि फक्त त्याच अटीवर, सत्य असू द्या:
\begin{quote}
$x$ हा $\theta(u)$ या सूत्राचा संकेतांक आहे, आणि $\Th{T}$ हे
$\theta(\num y)$ सिद्ध करते.
\end{quote}
$\Th{T}$ निर्णेय आहे असे गृहीत धरले असल्यामुळे $R$ संगणनक्षम
आहे: सूत्र-संकेतांक ओळखणे व त्यात संख्यांकाचे प्रतिस्थापन करणेही
संगणनक्षम आहे. आता $R$ सार्वत्रिक असल्याचे दाखवू.
$S(y)$ हा कोणताही संगणनक्षम संबंध असेल, तर त्याचे $\Th{Q}$ मध्ये
(आणि म्हणून $\Th{T}$ मध्येही) $\theta_S(u)$ या सूत्राने निरूपण होते.
म्हणून प्रत्येक $n$ साठी
\begin{eqnarray*}
S(n) & \lif & T \vdash \theta_S(\num n) \\
& \lif & R(\Gn{\theta_S(u)}, n)
\end{eqnarray*}
आणि
\begin{eqnarray*}
\lnot S(n) & \lif & T \vdash \lnot \theta_S(\num n) \\
& \lif & T \not\vdash \theta_S(\num n) \quad \text{($\Th{T}$ सुसंगत असल्यामुळे)} \\
& \lif & \lnot R(\Gn{\theta_S(u)}, n).
\end{eqnarray*}
म्हणजे प्रत्येक $y$ साठी $S(y)$ सत्य असते तर आणि तरच
$R(\Gn{\theta_S(u)}, y)$ सत्य असते. त्यामुळे $R$ सार्वत्रिक
ठरतो आणि अपेक्षित व्याघात मिळतो.
\end{proof}

``सत्य अंकगणित'' ही $\Setabs{\varphi}{ \Sat{\Nat}{\varphi}}$ उपपत्ती
असू द्या; म्हणजे अंकगणिताच्या भाषेत प्रमाण अर्थनिर्धारणानुसार
सत्य असलेल्या वाक्यांचा हा संच आहे.
प्रमाण प्रतिमानात क्यूची स्वयंसिद्धके सत्य असल्याने
सत्य अंकगणित हा त्याचा सुसंगत विस्तार आहे.

\begin{cor}
सत्य अंकगणित अनिर्णेय आहे.
\end{cor}















\section{निर्णेय स्वयंसिद्धकसंचाने निरूपणयोग्य उपपत्ती}\label{inc:tcp:cax:sec}

उपपत्ती $\Th{T}$ साठी स्वयंसिद्धकांचा संगणनक्षम संच~$A$ असेल,
तर तिला \emph{निर्णेय स्वयंसिद्धकसंचाने निरूपणयोग्य} म्हणतात.
($A$ हा $\Th{T}$ चा स्वयंसिद्धकसंच आहे याचा अर्थ
$T = \Setabs{\varphi}{A \Proves \varphi}$.) नैसर्गिक संख्यांचे कोणतेही
``युक्तिसंगत'' स्वयंसिद्धकीकरण हा गुणधर्म पूर्ण करेल.
विशेषतः, सांत स्वयंसिद्धकसंच असलेली कोणतीही उपपत्ती
निर्णेय स्वयंसिद्धकसंचाने निरूपणयोग्य असते.

\begin{lem}
$\Th{T}$ ही निर्णेय स्वयंसिद्धकसंचाने निरूपणयोग्य असेल, तर $\Th{T}$ ही
संगणनक्षमपणे प्रगणनीय असते.
\end{lem}

\begin{proof}
$A$ हा $\Th{T}$ चा संगणनक्षम स्वयंसिद्धकसंच आहे असे समजू.
$\varphi \in T$ असल्याचे होकारार्थी प्रमाण शोधण्यासाठी
स्वयंसिद्धकांपासून $\varphi$ ची निष्पत्ती शोधत राहा;
ते वाक्य सदस्य नसेल, तर हा शोध थांबेलच असे नाही.

थोडे वेगळ्या शब्दांत, $\varphi$ हे $\Th{T}$ मध्ये असते तर आणि तरच $A$ मधील $\psi_1$, \dots, $\psi_k$ अशा स्वयंसिद्धकांची
सांत यादी आणि प्रथम-क्रम तर्कशास्त्रात
$(\psi_1 \land \dots \land \psi_k) \lif \varphi$ ची
निष्पत्ती असते. पण ``असा काही \dots आहे की \dots''
या रूपातील व्याख्या असलेला कोणताही संच संगणनक्षमपणे प्रगणनीय असतो,
जर दुसऱ्या ``\dots'' मधील अट संगणनक्षम असेल, हे आपण आधीच जाणतो.
\end{proof}












\section{निर्णेय स्वयंसिद्धकसंचाने निरूपणयोग्य संपूर्ण उपपत्ती निर्णेय असतात}\label{inc:tcp:cdc:sec}

प्रत्येक वाक्य~$\varphi$ साठी $\varphi$ किंवा $\lnot \varphi$ यांपैकी
किमान एक सिद्ध होत असेल, तर उपपत्तीला {\em संपूर्ण} म्हणतात.

\begin{lem}
उपपत्ती $\Th{T}$ संपूर्ण आणि निर्णेय स्वयंसिद्धकसंचाने निरूपणयोग्य आहे असे समजू.
मग $\Th{T}$ निर्णेय असते.
\end{lem}

\begin{proof}
$\Th{T}$ संपूर्ण आहे आणि $A$ हा तिचा संगणनक्षम स्वयंसिद्धकसंच
आहे असे समजू. $\Th{T}$ विसंगत असेल, तर ती स्पष्टपणे संगणनक्षम
आहे. (वाक्य-संकेतांकासाठी अल्गोरिदम: ``नेहमी होय म्हणा.'')
वाक्याचा संकेतांक नसलेले आदान प्रथम वगळावे. म्हणून $\Th{T}$
सुसंगतही आहे असे मानू शकतो.

वाक्य $\varphi$ हे $\Th{T}$ मध्ये आहे की नाही हे ठरवण्यासाठी,
$\Th{T}$ पासून $\varphi$ ची निष्पत्ती आणि
$\lnot \varphi$ ची निष्पत्ती यांचा एकाच वेळी शोध घ्या.
$\Th{T}$ संपूर्ण असल्यामुळे यांपैकी एक तरी निष्पत्ती मिळेलच;
आणि $\Th{T}$ सुसंगत असल्यामुळे $\lnot \varphi$ ची
निष्पत्ती मिळाल्यास $\varphi$ ची निष्पत्ती असणार नाही.

हेच दुसऱ्या प्रकारे पाहू: $\Th{T}$ ही संगणनक्षमपणे प्रगणनीय आहे, हे आपण
आधीच जाणतो. म्हणून आधी सिद्ध केलेल्या प्रमेयानुसार,
$\Th{T}$ चा पूरक संचदेखील संगणनक्षमपणे प्रगणनीय आहे हे दाखवणे पुरेसे आहे.
पण वाक्य~$\varphi$ हे $\overline{\Th{T}}$ मध्ये असते तर आणि तरच $\lnot \varphi$ हे $\Th{T}$
मध्ये असते. त्यामुळे $\overline{\Th{T}} \leq_m
\Th{T}$.
येथे उपपत्ती म्हणजे वाक्यांच्या संकेतांकांचा संच मानला आहे.
वाक्य-संकेतांकांवर हे न्यूनीकरण निषेध घालून मिळते. सर्व नैसर्गिक
संख्यांवर ते एकूण फलन करण्यासाठी, वाक्याचा संकेतांक नसलेले
प्रत्येक आदान शून्याची स्वतःशी एकरूपता सांगणाऱ्या निश्चित
तार्किक प्रमेयाच्या संकेतांकाकडे पाठवा. अशी आदाने पूरक संचात
असतात आणि त्यांची प्रतिमा उपपत्तीत असते. वाक्य-संकेतांक
ओळखणेही संगणनक्षम असल्याने हे एकूण न्यूनीकरण संगणनक्षम आहे.
\end{proof}












\section[Q चा संपूर्ण, सुसंगत, निर्णेय स्वयंसिद्धकसंचाने निरूपणयोग्य विस्तार नाही]{$\Th{Q}$ चा संपूर्ण, सुसंगत, निर्णेय स्वयंसिद्धकसंचाने निरूपणयोग्य
  विस्तार नाही}\label{inc:tcp:inc:sec}

\begin{thm}
\label{inc:tcp:inc:thm:first-incompleteness}
$\Th{Q}$ चा संपूर्ण, सुसंगत आणि निर्णेय स्वयंसिद्धकसंचाने निरूपणयोग्य विस्तार
अस्तित्वात नाही.
\end{thm}

\begin{proof}
$\Th{Q}$ चा कोणताही सुसंगत, निर्णेय विस्तार नसतो, हे
आपल्याला आधीच माहीत आहे. पण $\Th{T}$ संपूर्ण आणि
निर्णेय स्वयंसिद्धकसंचाने निरूपणयोग्य असेल, तर ती निर्णेय असते.
\end{proof}

\begin{explain}
हे प्रमेय गोडेल यांच्या १९३१ मधील पहिल्या अपूर्णता प्रमेयाच्या
मूळ मांडणीपासून फारसे दूर नाही. अधिक आधुनिक पारिभाषिक शब्द
सोडले, तर मुख्य फरक असे आहेत: गोडेल यांनी ``सुसंगत''ऐवजी
``$\omega$-सुसंगत'' ही अट वापरली; तसेच संगणनक्षमतेची आकारिक
संकल्पना तेव्हा प्रस्थापित झाली नव्हती, म्हणून ``निर्णेय स्वयंसिद्धकसंचाने निरूपणयोग्य''
हे सर्वसाधारणपणे म्हणता येत नव्हते. (पुढील दशकभरात गोडेल
यांच्यासह अनेकांनी संगणनक्षमतेची आकारिक प्रतिमाने विकसित केली;
गोडेल प्रकारची घटना कोणत्या उपपत्तींना लागू होते हे नेमके
सांगणे, हा त्यामागील एक मोठा उद्देश होता.)

संपूर्णता, सुसंगतता आणि निर्णेय स्वयंसिद्धकसंचाने निरूपणयोग्यता हे तिन्ही गुणधर्म
एकाच वेळी मिळू शकत नाहीत, असे प्रमेय सांगते. मात्र यांपैकी
कोणताही एक सोडल्यास उरलेले दोन मिळू शकतात: $\Th{Q}$ सुसंगत
आणि संगणनक्षम स्वयंसिद्धकांनी निरूपित आहे, पण संपूर्ण नाही;
विसंगत उपपत्ती संपूर्ण असून संगणनक्षम स्वयंसिद्धकांनी
(उदाहरणार्थ, $\{ \eq/[0][0] \}$ या संचाने) निरूपित होते,
पण सुसंगत नाही; आणि अंकगणितातील सर्व सत्य वाक्यांचा संच
संपूर्ण व सुसंगत आहे, पण संगणनक्षम स्वयंसिद्धकांनी
निरूपित होत नाही.
\end{explain}











\section[Q मध्ये सिद्ध होणारी आणि ज्यांचे निषेध सिद्ध होतात अशा वाक्यांचे संच संगणनक्षमपणे अविलगनीय आहेत]{$\Th{Q}$ मध्ये सिद्ध होणारी आणि ज्यांचे निषेध सिद्ध होतात
  अशा वाक्यांचे संच संगणनक्षमपणे अविलगनीय आहेत}\label{inc:tcp:ins:sec}


$\overline{\Th{Q}}$ हा ज्या वाक्यांचे {\em निषेध} $\Th{Q}$ मध्ये
सिद्ध होतात त्यांचा संच असू द्या, म्हणजे
$\overline{\Th{Q}} = \Setabs{\varphi}{\Th{Q} \Proves
  \lnot \varphi}$. विभक्त संच $A$ आणि $B$ यांना
संगणनक्षमपणे अविलगनीय म्हणतात, जर असा कोणताही
संगणनक्षम संच $C$ नसेल की $A \subseteq C$ आणि
$B \subseteq \Complement{C}$.

\begin{lem}
$\Th{Q}$ आणि $\overline{\Th{Q}}$ हे संगणनक्षमपणे अविलगनीय आहेत.
\end{lem}

\begin{proof}
$C$ हा असा संगणनक्षम संच आहे असे समजू की
$\Th{Q} \subseteq C$ आणि
$\overline{\Th{Q}} \subseteq \Complement{C}$.
$R(x,y)$ हा पुढील संबंध असू द्या:
\begin{quote}
$x$ हा $\theta(u)$ या सूत्राचा संकेतांक आहे आणि
$\theta(\num y)$ हे $C$ मध्ये आहे.
\end{quote}
सूत्र-संकेतांक ओळखणे, संख्यांकाचे प्रतिस्थापन करणे आणि
या संचातील सदस्यत्व तपासणे संगणनक्षम असल्यामुळे $R(x,y)$
संगणनक्षम आहे. आता तो सार्वत्रिक संगणनक्षम संबंध आहे असे दाखवू;
त्यातून व्याघात मिळेल.

$S(y)$ हा संगणनक्षम संबंध असून त्याचे $\Th{Q}$ मध्ये
$\theta_S(u)$ या सूत्राने निरूपण होते असे समजू. मग
\begin{eqnarray*}
S(n) & \lif & \Th{Q} \Proves \theta_S(\num n) \\
& \lif & \theta_S(\num n) \in C
\end{eqnarray*}
आणि
\begin{eqnarray*}
\lnot S(n) & \lif & \Th{Q} \Proves \lnot \theta_S(\num n) \\
& \lif & \theta_S(\num n) \in \overline{\Th{Q}} \\
& \lif & \theta_S(\num n) \not\in C
\end{eqnarray*}
असे मिळते. म्हणून $S(y)$ हे
$R(\#(\theta_S(u)),y)$ शी सममूल्य आहे.
\end{proof}












\section[Q शी सुसंगत उपपत्ती अनिर्णेय असतात]{$\Th{Q}$ शी सुसंगत उपपत्ती अनिर्णेय असतात}\label{inc:tcp:con:sec}

पुढील प्रमेय असे सांगते की केवळ $\Th{Q}$ अनिर्णेय आहे असे नाही,
तर $\Th{Q}$ शी विसंगती निर्माण न करणारी कोणतीही उपपत्ती अनिर्णेय असते.

\begin{thm}
$\Th{T}$ ही अंकगणिताच्या भाषेतील $\Th{Q}$ शी सुसंगत असलेली
कोणतीही उपपत्ती असू द्या (म्हणजे $\Th{T} \cup \Th{Q}$ सुसंगत आहे).
मग $\Th{T}$ अनिर्णेय असते.
\end{thm}

\begin{proof}
$\Th{Q}$ ची $\mathsf{Q}_1$, \dots, $\mathsf{Q}_8$ ही सांत स्वयंसिद्धके आहेत,
हे आठवा. त्यांच्या जागी $\mathsf{E} = \mathsf{Q}_1 \land
\dots \land \mathsf{Q}_8$ हे एकच स्वयंसिद्धकही घेता येते.

$\Th{T}$ ही $\Th{Q}$ शी सुसंगत असलेली निर्णेय उपपत्ती आहे
असे समजू. पुढील संच घेऊ:
\[C = \Setabs{\varphi}{\Th{T} \Proves \mathsf{E} \lif \varphi}.
\]
निश्चित सूत्राशी अभिव्यंजन तयार करणे संगणनक्षम आहे;
गृहीत उपपत्ती निर्णेय असल्याने या संचातील सदस्यत्वही निर्णेय आहे.
$C$ हा $\Th{Q}$ आणि $\overline{\Th{Q}}$ यांना संगणनक्षमपणे
विलग करणारा संच ठरेल, असे दाखवू; त्यामुळे व्याघात मिळेल.
प्रथम, $\varphi$ हे $\Th{Q}$ मध्ये असेल, तर $\Th{Q}$ च्या
स्वयंसिद्धकांपासून $\varphi$ सिद्ध होते. निगमन प्रमेयानुसार
प्रथम-क्रम तर्कशास्त्रात $\mathsf{E} \lif \varphi$ ची निष्पत्ती
असते. तार्किक प्रमेय असल्याने ते गृहीत निर्णेय उपपत्तीतही सिद्ध होते.
म्हणून $\varphi$ हे $C$ मध्ये आहे.

दुसरीकडे, $\varphi$ हे $\overline{\Th{Q}}$ मध्ये असेल, तर प्रथम-क्रम
तर्कशास्त्रात $\mathsf{E} \lif \lnot \varphi$ ची सिद्धता असते.
$\Th{T}$ हे $\mathsf{E} \lif \varphi$ सुद्धा सिद्ध करत असेल, तर
$\Th{T}$ हे $\lnot \mathsf{E}$ सिद्ध करेल; अशा वेळी
$\Th{T} \cup \Th{Q}$ विसंगत होईल. पण
$\Th{T} \cup \Th{Q}$ सुसंगत आहे असे आपण गृहीत धरले आहे.
म्हणून $\Th{T}$ हे $\mathsf{E} \lif \varphi$ सिद्ध करत नाही आणि $\varphi$
हे $C$ मध्ये नाही.

म्हणजे $\varphi$ हे $\Th{Q}$ मध्ये असेल तर $C$ मध्ये असते,
आणि $\varphi$ हे $\overline{\Th{Q}}$ मध्ये असेल तर $\Complement{C}$
मध्ये असते, हे दाखवले. त्यामुळे $C$ हा संगणनक्षम विलगक
ठरतो; हाच अपेक्षित व्याघात आहे.
\end{proof}

हे प्रमेय अतिशय प्रभावी आहे. उदाहरणार्थ, त्यावरून पुढील निष्कर्ष मिळतो:

\begin{cor}
  अंकगणिताच्या भाषेसाठी प्रथम-क्रम तर्कशास्त्र (म्हणजे
  $\Setabs{\varphi}{\text{$\varphi$ हे प्रथम-क्रम तर्कशास्त्रात सिद्ध होते}}$
  हा संच) अनिर्णेय आहे.
\end{cor}

\begin{proof}
प्रथम-क्रम तर्कशास्त्र हा $\emptyset$ पासून निष्पन्न होणाऱ्या
वाक्यांचा संच आहे; आणि तो रिक्त स्वयंसिद्धकसंच $\Th{Q}$ शी सुसंगत आहे.
\end{proof}












\section[Q च्या अर्थनिर्धारणाशी सुसंगत उपपत्ती अनिर्णेय असतात]{$\Th{Q}$ च्या अर्थनिर्धारणाशी सुसंगत उपपत्ती अनिर्णेय असतात}\label{inc:tcp:itp:sec}

हे निकाल आणखी सबल करता येतात. अनौपचारिकपणे सांगायचे तर,
भाषा $\Lang{L_1}$ चे दुसरी भाषा $\Lang{L_2}$ मध्ये
अर्थनिर्धारण करताना $\Lang{L_1}$ चे विश्व, संबंधचिन्हे आणि
फलनचिन्हे ही $\Lang{L_2}$ मधील सूत्रे वापरून
परिभाषित करावी लागतात. येथे त्यासाठी वेळ देणार नसलो,
तरी या संकल्पनेची नेमकी व्याख्या देता येते.

\begin{thm}
  $\Th{T}$ ही अशा भाषेतील उपपत्ती असू द्या की त्या भाषेत
  अंकगणिताच्या भाषेचे अर्थनिर्धारण करता येते, आणि त्या
  अर्थनिर्धारणातील $\Th{Q}$ शी $\Th{T}$ सुसंगत असते.
  मग $\Th{T}$ अनिर्णेय असते. $\Th{T}$ कडून $\Th{Q}$ च्या
  स्वयंसिद्धकांची अर्थनिर्धारित रूपे सिद्ध होत असतील, तर
  $\Th{T}$ चा कोणताही सुसंगत विस्तार निर्णेय नसतो.
\end{thm}

या प्रमेयाची सिद्धता मागील प्रमेयाच्या सिद्धतेत थोडा बदल करून
मिळते; प्रतिउदाहरणावरून $\Th{Q}$ आणि $\overline{\Th{Q}}$ यांचे
विलगीकरण करता येईल. निवडीच्या स्वयंसिद्धकासह
झर्मेलो--फ्रेंकेल संच उपपत्ती $\Th{ZFC}$ ही सामान्य
गणिताचा मोठा भाग विकसित करण्याइतकी समर्थ स्वयंसिद्धकीय
पायाभूत व्यवस्था मानता येते. $\Th{ZFC}$ मध्ये नैसर्गिक
संख्या परिभाषित करता येतात; आणि या अर्थनिर्धारणातून
$\Th{Q}$ ची स्वयंसिद्धके सत्य ठरतात. म्हणून पुढील निष्कर्ष मिळतो:

\begin{cor}
$\Th{ZFC}$ चा कोणताही सुसंगत, निर्णेय विस्तार नाही.
\end{cor}

\begin{cor}
$\Th{ZFC}$ चा संपूर्ण, सुसंगत आणि संगणनक्षमपणे
निर्णेय स्वयंसिद्धकसंचाने निरूपणयोग्य असा कोणताही विस्तार नाही.
\end{cor}

$\Th{ZFC}$ च्या भाषेत $\in$ हे एकमेव द्विस्थानी संबंधचिन्ह
आहे. (खरे तर समतेचीही गरज नाही.) म्हणून पुढील निष्कर्ष मिळतो:

\begin{cor}
द्विस्थानी संबंधचिन्ह असलेल्या कोणत्याही भाषेसाठी प्रथम-क्रम
तर्कशास्त्र अनिर्णेय आहे.
\end{cor}

दोन एकस्थानी फलनचिन्हे असलेल्या कोणत्याही भाषेलाही हा निकाल
लागू होतो, कारण त्यांचा वापर करून द्विस्थानी संबंधचिन्हाचे
अनुकरण करता येते. हे निकाल नेमक्या सीमेवरचे आहेत: केवळ
\emph{एकस्थानी} संबंधचिन्हे आणि जास्तीत जास्त एक
\emph{एकस्थानी} फलनचिन्ह असलेल्या भाषेसाठी प्रथम-क्रम
तर्कशास्त्र निर्णेय असते, असे सिद्ध होते.

आणखी एक माहितीची बाब. प्रमाण प्रतिमानात सत्य असलेल्या
$\Obj 0$, $'$, $+$, $\times$, $<$ या भाषेतील वाक्यांचा
संच अनिर्णेय आहे, हे आपल्याला माहीत आहे. प्रत्यक्षात,
उरलेल्या चिन्हांच्या साहाय्याने $<$ परिभाषित करता येते,
आणि $\times$ व $'$ यांच्या साहाय्याने $+$ परिभाषित
करता येते. म्हणून $\Obj 0$, $'$, $\times$ या भाषेतील
सत्य वाक्यांचा संच अनिर्णेय आहे. दुसरीकडे, प्रमाण
अंकगणितीय प्रतिमानात सत्य असलेल्या $\Obj 0$, $'$, $+$
या भाषेतील वाक्यांचा संच निर्णेय आहे, हे प्रेस्बुर्गर यांनी
दाखवले. मात्र त्याची निर्णयपद्धती संगणनदृष्ट्या अव्यवहार्य आहे.



\chapter{अपूर्णता आणि सिद्धतायोग्यता}\label{inc:inp:chap}









\section{परिचय}\label{inc:inp:int:sec}

हिल्बर्टच्या मते, नैसर्गिक संख्यांसारख्या गणितीय रचनेसाठीची
स्वयंसिद्धकपद्धती त्या रचनेविषयीची सर्व सत्य विधाने सिद्ध करू देत
नसेल, तर ती अपुरी आहे. निष्पत्तीच्या आकारिक पद्धतींविषयीच्या
त्याच्या नंतरच्या स्वारस्याचा विचार केला, तर नैसर्गिक संख्यांविषयी
युक्तिवाद करण्यासाठी वापरली जाणारी आकारिक पद्धत केवळ सुसंगतच
नव्हे, तर \emph{संपूर्ण}ही असावी, असे त्याला वाटत होते असे दिसते:
त्या भाषेतील प्रत्येक विधान किंवा त्याचा निषेध निष्पन्न करता येण्याजोगे
असावा. गोडेलचे पहिले अपूर्णता प्रमेय दाखवते की अशी
स्वयंसिद्धकपद्धती अस्तित्वात नाही: अंकगणितासाठी संपूर्ण, सुसंगत
आणि स्वयंसिद्धकीकरणीय अशी कोणतीही आकारिक पद्धत नाही. खरे तर,
“पुरेशी प्रबळ”, सुसंगत आणि स्वयंसिद्धकीकरणीय अशी कोणतीही
गणितीय उपपत्ती संपूर्ण नसते.

हिल्बर्टचा आणखी महत्त्वाचा उद्देश म्हणजे अभिजात गणिताला आधार
देण्याच्या कार्यक्रमाचा मुख्य भाग: अभिजात युक्तिवादाचे निरूपण
करणाऱ्या आकारिक पद्धतींच्या सुसंगततेच्या सांतवादी सिद्धता
शोधणे. म्हणून हिल्बर्टच्या कार्यक्रमासाठी गोडेलचे दुसरे
अपूर्णता प्रमेय हा अधिक मोठा धक्का होता. पहिल्या प्रमेयाप्रमाणे
दुसरेही प्रमेय काहीसे स्थूल शब्दांत सांगता येते. साधारणपणे ते
असे म्हणते की अंकगणिताची कोणतीही सुसंगत, प्रभावीपणे
स्वयंसिद्धकीकरण करता येणारी आणि पुरेशी प्रबळ उपपत्ती स्वतःची
सुसंगतता व्यक्त करणारे प्रमाण आकारिक वाक्य सिद्ध करू शकत नाही. येथे “पुरेशी प्रबळ” यामध्ये
केवळ~$\Th{Q}$ पेक्षा थोडे अधिक सामर्थ्य समाविष्ट करावे लागेल.

गोडेलच्या अपूर्णता प्रमेयाच्या मूळ सिद्धतेमागची कल्पना
एपिमेनिडीजच्या विरोधाभासात दिसते. क्रीटचा रहिवासी असलेल्या
एपिमेनिडीजने सर्व क्रीटवासी खोटारडे असतात असे म्हटले;
“हे वाक्य असत्य आहे” हे विधान त्या विरोधाभासाचे अधिक थेट
रूप आहे. सारतः, सत्याच्या जागी निष्पन्नता ठेवून गोडेलने
स्वतःच निष्पन्न करता येण्याजोगे नसल्याचे अप्रत्यक्षपणे सांगणारे
वाक्य आकारिक रीतीने तयार केले. ते वाक्य
निष्पन्न करता येण्याजोगे असेल, तर उपपत्ती विसंगत ठरेल. गोडेलने
विचारात घेतलेल्या स्वयंसिद्धकपद्धतीत त्या वाक्याचा
निषेधही निष्पन्न करता येण्याजोगे नसतो, असे त्याने दाखवले. (या दुसऱ्या
भागासाठी त्याला उपपत्ती~$\Th{T}$ ही तथाकथित
“$\omega$-सुसंगत” आहे असे गृहीत धरावे लागले.
$\omega$-सुसंगतता ही सुसंगततेशी संबंधित पण अधिक प्रबळ अट
आहे.\footnote{म्हणजे प्रत्येक $\omega$-सुसंगत उपपत्ती
सुसंगत असते; उलट विधान खरे नसते.} गोडेलनंतर काही वर्षांनी
रोसरने केवळ~$\Th{T}$ च्या साध्या सुसंगततेचे गृहीतक पुरेसे
आहे असे दाखवले.)

पहिली अडचण म्हणजे स्वतःचा निर्देश करणारे वाक्य
कसे तयार करायचे हे समजणे. $\Th{Q}$ च्या भाषेतील प्रत्येक
सूत्र~$\varphi$ साठी~$\gn{\varphi}$ हा~$\Gn{\varphi}$ या संख्येला
अनुरूप संख्यांक मानू. याचा अर्थ नीट पाहू: $\varphi$ हे
$\Th{Q}$ च्या भाषेतील सूत्र आहे, $\Gn{\varphi}$ ही
नैसर्गिक संख्या आहे, आणि $\gn{\varphi}$ हे $\Th{Q}$ च्या
भाषेतील \emph{पद} आहे. त्यामुळे $\Th{Q}$ च्या भाषेतील
प्रत्येक सूत्र~$\varphi$ ला~$\gn{\varphi}$ असे \emph{नाव}
मिळते; ते $\Th{Q}$ च्या भाषेतील पद आहे. यामुळे
$\Th{Q}$ च्या भाषेतील सूत्रे ना
इतर सूत्रे विषयी काहीतरी “सांगता” येते अशी
संकल्पनात्मक चौकट मिळते. पुढील उपपत्तीला स्थिरबिंदू
उपपत्ती म्हणतात.

\begin{lem}
$\Th{T}$ ही~$\Th{Q}$ चा विस्तार असलेली कोणतीही उपपत्ती
असू द्या, आणि $\psi(x)$ हे फक्त~$x$ हेच मुक्त चल असलेले
कोणतेही सूत्र असू द्या. मग असे
वाक्य~$\varphi$ असते की
$\Th{T} \Proves\varphi \liff \psi(\gn{\varphi})$.
\end{lem}

ही उपपत्ती असे सांगते की कोणत्याही गुणधर्मासाठी $\psi(x)$,
“$\psi(x)$ माझ्याबाबतीत खरा आहे” असे सांगणारे
वाक्य~$\varphi$ असते, आणि $\Th{T}$ ला हे
“माहीत” असते.

असे वाक्य कसे तयार करायचे? क्वाइनने मांडलेले
एपिमेनिडीजच्या विरोधाभासाचे पुढील रूप पाहू:
\begin{quote}
“उद्धृत रूप आधी जोडल्यास असत्य विधान मिळते” उद्धृत रूप
आधी जोडल्यास असत्य विधान मिळते.
\end{quote}
हे वाक्य थेट स्वतःचा निर्देश करत नाही. ते फक्त अवतरणचिन्हांतील
विन्यासात्मक वस्तूंविषयी विधान करते; त्या अर्थाने ते
पुढील वाक्यांसारखेच आहे:
\begin{enumerate}
\item “रॉबर्ट” हे छान नाव आहे.
\item “मी धावलो.” हे छोटे वाक्य आहे.
\item “तीन शब्द आहेत” यात तीन शब्द आहेत.
\end{enumerate}
पण “उद्धृत रूप आधी जोडल्यास असत्य विधान मिळते” या
वाक्यांशाच्या आधी त्याचेच उद्धृत रूप जोडले, तर काय होते?
मूळ वाक्यच मिळते!{} थोडक्यात, ते वाक्य स्वतःच असत्य
असल्याचे सांगते.












\section{स्थिरबिंदू उपपत्ती}\label{inc:inp:fix:sec}

\begin{explain}
स्थिरबिंदू उपपत्ती सांगते की कोणत्याही
सूत्र~$\psi(x)$ साठी असे वाक्य~$\varphi$ असते की
$\Th{T} \Proves \varphi \liff \psi(\gn{\varphi})$, मात्र
$\Th{T}$ हा~$\Th{Q}$ चा विस्तार असला पाहिजे.
खोटारड्याच्या वाक्यासाठी $\varphi$ हे~“$\gn{\varphi}$ असत्य आहे”
या विधानाशी (उपपत्ती~$\Th{T}$ मध्ये सिद्ध होईल अशा
रीतीने) सममूल्य हवे; म्हणजेच $\Gn{\varphi}$ ही एका असत्य
वाक्याची गोडेल संख्या आहे असे त्यात म्हटलेले हवे.
सिद्धतेची कल्पना समजून घेण्यासाठी $\varphi$ विषयी क्वाइनच्या
अनौपचारिक मांडणीशी तुलना करू: “उद्धृत रूप आधी जोडल्यास
असत्य विधान मिळते” उद्धृत रूप आधी जोडल्यास असत्य विधान
मिळते. एखादा शब्दसमूह घेऊन त्याच्या आधी त्याचेच उद्धृत
रूप जोडून वाक्य तयार करण्याच्या क्रियेला त्या शब्दसमूहाचे
\emph{विकर्णीकरण} म्हणू शकतो; तयार झालेला परिणाम हे
त्याचे विकर्णरूप. म्हणून “उद्धृत रूप आधी जोडल्यास असत्य
विधान मिळते” या शब्दसमूहाचे विकर्णरूप म्हणजे “उद्धृत रूप
आधी जोडल्यास असत्य विधान मिळते” उद्धृत रूप आधी
जोडल्यास असत्य विधान मिळते. आता लक्षात घ्या की क्वाइनचे
खोटारड्याचे वाक्य हे “असत्य विधान मिळते” याचे नव्हे,
तर “उद्धृत रूप आधी जोडल्यास असत्य विधान मिळते” याचे
विकर्णरूप आहे. म्हणजेच खोटारड्याचे वाक्य मिळवण्यासाठी
ज्याचे विकर्णीकरण करतो, त्या गुणधर्मातच विकर्णीकरण
समाविष्ट असते!{}

अंकगणिताच्या भाषेत, एक मुक्त चल असलेल्या
सूत्राची गोडेल संख्या काढून, त्या चलाच्या जागी
त्या संख्येचा प्रमाण संख्यांक बसवून त्याचे उद्धृत रूप
बनवतो. $\mathsf{E}(x)$ चे विकर्णरूप हे $\mathsf{E}(\num{n})$ आहे,
जिथे $n = \Gn{\mathsf{E}(x)}$. (यापुढे
$\num{\Gn{\mathsf{E}(x)}}$ याचे संक्षिप्त रूप
$\gn{\mathsf{E}(x)}$ असे लिहू.) म्हणून $\psi(x)$ हा
“असत्य आहे” हा गुणधर्म असेल, तर “स्वतःचे उद्धृत रूप
आधी जोडल्यास असत्य विधान मिळते” याचा अर्थ
“स्वतःच्या विकर्णरूपाच्या गोडेल संख्येला लागू केल्यावर
असत्य विधान मिळते” असा होईल. सूत्राची गोडेल
संख्या~$n$ दिल्यावर त्याच्या विकर्णरूपाची गोडेल संख्या
काढणाऱ्या $\fn{diag}(n)$ या फलनासाठी
$\Obj{diag}$ असे चिन्ह आपल्याकडे असेल, तर
$\mathsf{E}(x)$ हे $\psi(\Obj{diag}(x))$ असे लिहिता येईल.
क्वाइनच्या खोटारड्याच्या वाक्याचे रूप मग याचे
विकर्णीकरण, म्हणजे $\mathsf{E}(\gn{\mathsf{E}(x)})$ किंवा
$\psi(\Obj{diag}(\gn{\psi(\Obj{diag}(x))}))$, असे असेल.
अर्थात, आता $\psi(x)$ हा इतर कोणताही गुणधर्म असला,
तरी हीच रचना चालेल. अपूर्णता प्रमेयासाठी $\psi(x)$
हा “$x$ हे~$\Th{T}$ मध्ये निष्पन्न करता येण्याजोगे नाही” हा
गुणधर्म घेऊ. तेव्हा $\mathsf{E}(x)$ चे वर्णन असे होईल:
“स्वतःच्या विकर्णरूपाच्या गोडेल संख्येला लागू केल्यावर
$\Th{T}$ मध्ये निष्पन्न करता येण्याजोगे नसलेले
वाक्य मिळते.”

$\Th{T}$ मध्ये ही रचना आकारिक रीतीने मांडण्यासाठी
$\fn{diag}$ चे आकारिकीकरण करावे लागेल.
$\fn{diag}(n)$ हे संगणनक्षम, किंबहुना आदिम पुनरावर्ती
फलन आहे: $n$ ही~$\mathsf{E}(x)$ या सूत्राची गोडेल
संख्या असेल, तर $\fn{diag}(n)$ हे
$\mathsf{E}(\gn{\mathsf{E}(x)})$ ची गोडेल संख्या देते. अन्य आदानांसाठी शून्य मूल्य निवडून
हे फलन सर्व नैसर्गिक संख्यांवर परिभाषित करतो.
(आठवा, $\gn{\mathsf{E}(x)}$ हा~$\mathsf{E}(x)$ च्या गोडेल
संख्येचा प्रमाण संख्यांक, म्हणजे
$\num{\Gn{\mathsf{E}(x)}}$, आहे.) $\Obj{diag}$ हे
$\fn{diag}$ चे निरूपण करणारे $\Th{T}$ मधील
फलनचिन्ह असेल, तर $\varphi$ म्हणून
$\psi(\Obj{diag}(\gn{\psi(\Obj{diag}(x))}))$
हे सूत्र घेता येईल. लक्षात घ्या:
\begin{align*}
\fn{diag}(\Gn{\psi(\Obj{diag}(x))}) & =
\Gn{\psi(\Obj{diag}(\gn{\psi(\Obj{diag}(x))}))} \\
& = \Gn{\varphi}.
\end{align*}
$\Th{T}$ पुढील विधान निष्पन्न करू शकते असे धरल्यास,
\[\Obj{diag}(\gn{\psi(\Obj{diag}(x))}) = \gn{\varphi},
\]
त्याला $\psi(\Obj{diag}(\gn{\psi(\Obj{diag}(x))}))
\liff \psi(\gn{\varphi})$ हेही निष्पन्न करता येते.
पण डावी बाजू व्याख्येनुसार~$\varphi$ आहे.

अर्थात, सामान्यतः $\Obj{diag}$ हे
$\Th{T}$ मधील फलनचिन्ह नसेल; आणि
$\Th{Q}$ मधील तर नक्कीच नाही. पण
$\fn{diag}$ संगणनक्षम असल्यामुळे ते
$\Th{Q}$ मध्ये एखाद्या
$\theta_{\fn{diag}}(x,y)$ सूत्राने
\emph{निरूपित} करता येते. म्हणून
$\psi(\Obj{diag}(x))$ असे लिहिण्याऐवजी
$\lexists[y][(\theta_{\fn{diag}}(x,y) \land \psi(y))]$
असे लिहू शकतो. याखेरीज वर रेखाटलेली सिद्धता तशीच
चालते; किंबहुना ती आधीच~$\Th{Q}$ मध्ये चालते.
\end{explain}

\begin{lem}
\label{inc:inp:fix:lem:fixed-point} $\psi(x)$ हे~$x$ हे एकमेव मुक्त चल
असलेले कोणतेही सूत्र असू द्या. मग असे वाक्य~$\varphi$
असते की $\Th{Q} \Proves
\varphi \liff \psi(\gn{\varphi})$.
\end{lem}


\begin{proof}
$\psi(x)$ दिल्यावर $\mathsf{E}(x)$ हे
$\lexists[y][(\theta_{\fn{diag}}(x,y) \land \psi(y))]$
सूत्र मानू, आणि $\varphi$ हे त्याचे विकर्णरूप,
म्हणजे $\mathsf{E}(\gn{\mathsf{E}(x)})$, मानू.

$\theta_{\fn{diag}}$ हे $\fn{diag}$ चे निरूपण
करते आणि $\fn{diag}(\Gn{\mathsf{E}(x)}) = \Gn{\varphi}$
असल्यामुळे $\Th{Q}$ पुढील विधाने
निष्पन्न करू शकते:
\begin{align}
  & \theta_{\fn{diag}}(\gn{\mathsf{E}(x)}, \gn{\varphi}) \label{inc:inp:fix:repdiag1} \\
  & \lforall[y][(\theta_{\fn{diag}}(\gn{\mathsf{E}(x)},y) \lif
  \eq[y][\gn{\varphi}])]. \label{inc:inp:fix:repdiag2}
\end{align}
आता $\Th{Q} \Proves \varphi \liff \psi(\gn{\varphi})$
हे दाखवू. केवळ तर्कशास्त्र आणि $\Th{Q}$
मध्ये निष्पन्न करता येण्याजोगे असलेली तथ्ये वापरून
अनौपचारिक युक्तिवाद करू.

प्रथम~$\varphi$, म्हणजे $\mathsf{E}(\gn{\mathsf{E}(x)})$, खरे
मानू. $\mathsf{E}(x)$ च्या व्याख्येकडे परत पाहिल्यास,
$\mathsf{E}(\gn{\mathsf{E}(x)})$ हेच पुढील विधान आहे:
\[\lexists[y][(\theta_{\fn{diag}}(\gn{\mathsf{E}(x)},y) \land \psi(y))].
\]
अशा~$y$ चा विचार करा. आता
$\theta_{\fn{diag}}(\gn{\mathsf{E}(x)},y)$ असल्यामुळे
\ref{inc:inp:fix:repdiag2} वरून $y = \gn{\varphi}$ मिळते.
म्हणून $\psi(y)$ वरून~$\psi(\gn{\varphi})$ मिळते.

आता $\psi(\gn{\varphi})$ मानू. \ref{inc:inp:fix:repdiag1}
वरून पुढील मिळते:
\begin{align*}
& \theta_{\fn{diag}}(\gn{\mathsf{E}(x)}, \gn{\varphi}) \land \psi(\gn{\varphi}).
\intertext{त्यावरून पुढील निष्पन्न होते:}
& \lexists[y][(\theta_{\fn{diag}}(\gn{\mathsf{E}(x)},y) \land \psi(y))].
\end{align*}
पण हेच $\mathsf{E}(\gn{\mathsf{E}(x)})$, म्हणजेच~$\varphi$, आहे.
\end{proof}

\begin{digress}
या सिद्धतेची संगणनक्षमता सिद्धांतातील स्थिरबिंदू
उपपत्तीच्या सिद्धतेशी तुलना करा. येथे स्वतःच्याच
साहाय्याने \emph{विधान} परिभाषित करायचे आहे,
तर तिथे स्वतःच्याच साहाय्याने \emph{फलन}
परिभाषित करायचे होते. हा फरक सोडला, तर
मूळ कल्पना तीच आहे.
\end{digress}

\begin{prob}
  सूत्र~$\varphi(x)$ ला \emph{सत्याची व्याख्या}
  असे म्हणतात, जर सर्व वाक्ये~$\psi$
  साठी $\Th{Q}
  \Proves \psi \liff \varphi(\gn{\psi})$ असेल.
  स्थिरबिंदू उपपत्ती वापरून कोणतेही
  सूत्र सत्याची व्याख्या असू शकत नाही
  हे दाखवा.
\end{prob}












\section{पहिले अपूर्णता प्रमेय}\label{inc:inp:1in:sec}

आता गोडेलने पहिल्या अपूर्णता प्रमेयासाठी दिलेल्या मूळ
सिद्धतेचे वर्णन करता येईल. अंकगणिताच्या भाषेचा विस्तार
असलेल्या भाषेतील $\Th{T}$ ही कोणतीही संगणनक्षमपणे
स्वयंसिद्धकीकृत उपपत्ती असू द्या, आणि $\Th{T}$ मध्ये
$\Th{Q}$ ची स्वयंसिद्धके समाविष्ट असू द्या. यावरून
विशेषतः $\Th{T}$ संगणनक्षम फलने आणि संबंध
निरूपित करते.

सूत्रे आणि सिद्धता संख्या म्हणून योग्य रीतीने सांकेतिक
केल्यास $\Prf[\Th{T}](x,y)$ हा संबंध संगणनक्षम
असतो, असा आपण युक्तिवाद केला आहे. येथे
$\Prf[\Th{T}](x,y)$ खरा असतो, तर आणि तरच
$x$ ही~$\Th{T}$ मध्ये गोडेल संख्या~$y$
असलेल्या सूत्राच्या निष्पत्ती
ची गोडेल संख्या असते. गोडेलच्या मनात असलेल्या
विशिष्ट उपपत्तीसाठी हा संबंध आदिम पुनरावर्ती आहे
हेही त्याने आपल्या शोधपत्रातील 45 फलने आणि
संबंधांची यादी वापरून दाखवले. त्या यादीतील
45 वा संबंध, $x B y$, हा त्याने निवडलेल्या
$\Th{T}$ साठीचा $\Prf[\Th{T}](x,y)$
संबंधच आहे. गोडेल आपल्या शोधपत्रात जिथे
“पुनरावर्ती” हा शब्द वापरतो, तिथे आज आपण
“आदिम पुनरावर्ती” ही संज्ञा वापरू, हे लक्षात ठेवा.

$\Prf[\Th{T}](x,y)$ संगणनक्षम असल्यामुळे
तो $\Th{T}$ मध्ये निरूपणीय आहे. त्याचे
निरूपण करणाऱ्या सूत्रासाठी
$\OPrf[\Th{T}](x,y)$ हे चिन्ह वापरू.
$\OProv[\Th{T}](y)$ हे सूत्र
$\lexists[x][\OPrf[\Th{T}](x,y)]$
असे मानू. हे गोडेलच्या यादीतील 46 व्या
संबंधाचे, $ \fn{Bew}(y)$ चे, वर्णन करते.
गोडेलच्या नोंदीनुसार, हाच एक संबंध असा
आहे की तो “पुनरावर्ती आहे असे म्हणता येत
नाही.” त्याचा संभाव्य अर्थ असा: व्याख्येवरून
तो संगणनक्षम आहे हे स्पष्ट होत नाही; आणि
नंतरच्या प्रगतीतून, उपपत्ती क्यूचा सुसंगत विस्तार असेल,
तर हा संबंध संगणनक्षम नसतो हेही सिद्ध झाले.

$\Th{T}$ ही~$\Th{Q}$ समाविष्ट करणारी
स्वयंसिद्धकीकरणीय उपपत्ती असू द्या. मग
$\Prf[\Th{T}](x, y)$ निर्णेय असतो आणि
म्हणून $\Th{Q}$ मध्ये
सूत्र~$\OPrf[\Th{T}](x, y)$
याने निरूपणीय असतो. वर वर्णन केल्याप्रमाणे
$\OProv[\Th{T}](y)$ हे सूत्र मानू.
स्थिरबिंदू उपपत्तीवरून असे सूत्र
$\mathsf{G}_\Th{T}$ मिळते की $\Th{Q}$
(आणि म्हणून $\Th{T}$ सुद्धा) पुढील
निष्पन्न करते:
\begin{equation}
\label{inc:inp:1in:eqn:qpf}
\mathsf{G}_\Th{T} \liff \lnot \OProv[\Th{T}](\gn{\mathsf{G}_\Th{T}}).
\end{equation}
लक्षात घ्या की $\mathsf{G}_\Th{T}$ चे
आशयतः विधान असे आहे: “$\mathsf{G}_\Th{T}$
हे~$\Th{T}$ मध्ये निष्पन्न करता येण्याजोगे नाही.”

\begin{lem}\label{inc:inp:1in:lem:cons-G-unprov}
$\Th{T}$ ही~$\Th{Q}$ चा विस्तार असलेली
सुसंगत आणि स्वयंसिद्धकीकरणीय उपपत्ती
असेल, तर $\Th{T} \Proves/ \mathsf{G}_\Th{T}$.
\end{lem}

\begin{proof}
$\Th{T}$ हे $\mathsf{G}_\Th{T}$ ला
निष्पन्न करते असे मानू. मग
निष्पत्ती \emph{अस्तित्वात आहे}.
म्हणून एखाद्या संख्या $m$ साठी
$\Prf[\Th{T}](m,\Gn{\mathsf{G}_\Th{T}})$
हा संबंध खरा असतो. पण मग
$\Th{Q}$ हे
$\OPrf[\Th{T}](\num m, \gn{\mathsf{G}_\Th{T}})$
वाक्य निष्पन्न करते.
त्यामुळे $\Th{Q}$ हे
$\lexists[x][\OPrf[\Th{T}](x,\gn{\mathsf{G}_\Th{T}})]$
सुद्धा निष्पन्न करते; व्याख्येनुसार
हेच $\OProv[\Th{T}](\gn{\mathsf{G}_\Th{T}})$
आहे. \ref{inc:inp:1in:eqn:qpf} वरून
$\Th{Q}$ हे $\lnot \mathsf{G}_\Th{T}$
निष्पन्न करते. $\Th{T}$ हा
$\Th{Q}$ चा विस्तार असल्यामुळे
$\Th{T}$ सुद्धा तसेच करते.
म्हणजे $\Th{T}$ हे $\mathsf{G}_\Th{T}$
निष्पन्न करत असेल, तर
$\lnot \mathsf{G}_\Th{T}$ सुद्धा
निष्पन्न करते, असे आपण दाखवले.
त्यामुळे ते विसंगत ठरेल.
\end{proof}

\begin{defn}
\label{inc:inp:1in:thm:oconsis-q}
उपपत्ती $\Th{T}$ ला
\emph{$\omega$-सुसंगत} म्हणतात, जर
पुढील अट पूर्ण होत असेल: कोणतेही
वाक्य $\lexists[x][\varphi(x)]$ घेऊ.
$\Th{T}$ हे $\lnot \varphi(\num 0)$,
$\lnot \varphi(\num 1)$,
$\lnot \varphi(\num 2)$, \dots
निष्पन्न करत असेल, तर
$\Th{T}$ हे $\lexists[x][\varphi(x)]$ सिद्ध करत नाही.
\end{defn}

प्रत्येक $\omega$-सुसंगत उपपत्ती
सुसंगतसुद्धा असते, हे लक्षात घ्या.
कारण $\Th{T}$ विसंगत असेल, तर
प्रत्येक~$\varphi$ साठी $\Th{T} \Proves \varphi$
असते. विशेषतः $\Th{T}$ विसंगत
असेल, तर प्रत्येक~$n$ साठी
$\lnot \varphi(\num n)$ हे ते
निष्पन्न करते आणि
$\lexists[x][\varphi(x)]$ हेही
निष्पन्न करते.
म्हणून विसंगत $\Th{T}$ हे
$\omega$-विसंगत असते. याच्या
प्रतिवर्ती विधानावरून
$\omega$-सुसंगत $\Th{T}$
सुसंगत असलेच पाहिजे.

\begin{lem}\label{inc:inp:1in:lem:omega-cons-G-unref}
$\Th{T}$ ही~$\Th{Q}$ चा विस्तार
असलेली $\omega$-सुसंगत,
स्वयंसिद्धकीकरणीय उपपत्ती
असेल, तर
$\Th{T} \Proves/ \lnot \mathsf{G}_\Th{T}$.
\end{lem}

\begin{proof}
$\Th{T}$ हे $\lnot \mathsf{G}_\Th{T}$
निष्पन्न करत असेल, तर ते
$\omega$-विसंगत असते, हे दाखवू.
$\Th{T}$ हे $\lnot \mathsf{G}_\Th{T}$
निष्पन्न करते असे मानू.
$\Th{T}$ विसंगत असेल, तर ते
$\omega$-विसंगत आहेच.
नसेल, तर $\Th{T}$ सुसंगत आहे.
म्हणून \ref{inc:inp:1in:lem:cons-G-unprov}
वरून ते $\mathsf{G}_\Th{T}$ ला
निष्पन्न करत नाही. $\Th{T}$
मध्ये $\mathsf{G}_\Th{T}$ ची
निष्पत्ती नसल्यामुळे
$\Th{Q}$ पुढील विधाने
निष्पन्न करते:
\[\lnot \OPrf[\Th{T}](\num 0, \gn{\mathsf{G}_\Th{T}}), \lnot \OPrf[\Th{T}](\num 1,
\gn{\mathsf{G}_\Th{T}}), \lnot \OPrf[\Th{T}](\num 2, \gn{\mathsf{G}_\Th{T}}), \dots
\]
त्यामुळे $\Th{T}$ सुद्धा ती
सिद्ध करते. दुसरीकडे
\ref{inc:inp:1in:eqn:qpf} वरून
$\lnot \mathsf{G}_\Th{T}$ हे
$\lexists[x][\OPrf[\Th{T}](x,\gn{\mathsf{G}_\Th{T}})]$
याच्याशी सममूल्य असते.
म्हणून $\Th{T}$ हे
$\omega$-विसंगत आहे.
\end{proof}

\begin{prob}
  प्रत्येक $\omega$-सुसंगत उपपत्ती
  सुसंगत असते. याचे उलट खरे नसते
  हे दाखवा: सुसंगत पण
  $\omega$-विसंगत उपपत्ती असतात.
  यासाठी $\Th{Q} \cup
  \{\lnot \mathsf{G}_\Th{Q}\}$ हा संच
  सुसंगत पण $\omega$-विसंगत
  आहे असे दाखवा.
\end{prob}

\begin{thm}
\label{inc:inp:1in:thm:first-incompleteness}
$\Th{T}$ ही~$\Th{Q}$ चा
$\omega$-सुसंगत,
स्वयंसिद्धकीकरणीय विस्तार असलेली
कोणतीही उपपत्ती असू द्या.
मग $\Th{T}$ संपूर्ण नाही.
\end{thm}

\begin{proof}
  $\Th{T}$ हे $\omega$-सुसंगत
  असेल, तर ते सुसंगत असते.
  म्हणून \ref{inc:inp:1in:lem:cons-G-unprov}
  वरून $\Th{T}
  \Proves/ \mathsf{G}_\Th{T}$.
  \ref{inc:inp:1in:lem:omega-cons-G-unref}
  वरून $\Th{T} \Proves/
  \lnot \mathsf{G}_\Th{T}$. म्हणजे
  $\Th{T}$ अपूर्ण आहे: ते
  $\mathsf{G}_\Th{T}$ किंवा
  $\lnot \mathsf{G}_\Th{T}$ यांपैकी
  एकही विधान निष्पन्न करत नाही.
\end{proof}












\section{रोसरचे प्रमेय}\label{inc:inp:ros:sec}

गोडेलची सिद्धता बदलून अधिक सबळ निकाल मिळेल का,
म्हणजे “$\omega$-सुसंगत” ऐवजी फक्त
“सुसंगत” असे गृहीत धरता येईल का?
रोसरने शोधलेल्या युक्तीमुळे याचे उत्तर
“होय” असे आहे. रोसरची युक्ती म्हणजे
$\OProv[\Th{T}](y)$ च्या जागी
“सुधारित” निष्पन्नता विधेय
$\ORProv_T(y)$ वापरणे.

\begin{thm}
\label{inc:inp:ros:thm:rosser}
$\Th{T}$ ही~$\Th{Q}$ चा विस्तार
असलेली कोणतीही सुसंगत आणि
स्वयंसिद्धकीकरणीय उपपत्ती असू द्या.
मग $\Th{T}$ संपूर्ण नसते.
\end{thm}

\begin{proof}
आठवा, $\OProv[\Th{T}](y)$ ची व्याख्या
$\lexists[x][\OPrf[\Th{T}](x,
  y)]$ अशी आहे. येथे
$\OPrf[\Th{T}](x, y)$ हे अशा निर्णेय
संबंधाचे निरूपण करते की $x$ ही
गोडेल संख्या~$y$ असलेल्या
वाक्याच्या निष्पत्ती
ची गोडेल संख्या असते, तर आणि तरच तो
संबंध खरा असतो. $x$ आणि~$y$
यांच्यातील पुढील संबंधही निर्णेय
असतो: $x$ ही गोडेल
संख्या~$y$ असलेल्या वाक्याच्या
\emph{खंडनाची} गोडेल संख्या असते.
$\fn{not}(x)$ हे पुढील कार्य
करणारे आदिम पुनरावर्ती फलन
असू द्या: $x$ हा $\varphi$ या
सूत्राचा संकेतांक असेल, तर
$\fn{not}(x)$ हा $\lnot
\varphi$ चा संकेतांक असतो. सूत्रसंकेतांक नसलेल्या
आदानासाठी शून्य परत देऊन हे फलन सर्वत्र परिभाषित करतो. मग
$\Refut[\Th{T}](x, y)$
खरे असते, तर आणि तरच
$\Prf[\Th{T}](x, \fn{not}(y))$
खरे असते. त्याचे निरूपण
$\ORefut[\Th{T}](x, y)$
करू द्या. मग
$\Th{T} \Proves \lnot \varphi$
असेल आणि $\delta$ ही
संबंधित निष्पत्ती असेल,
तर $\Th{Q} \Proves
\ORefut[\Th{T}](\gn{\delta}, \gn{\varphi})$.
आता $\ORProv[\Th{T}](y)$ ची
व्याख्या अशी करू:
\[\lexists[x][(\OPrf[\Th{T}](x,y) \land \lforall[z][(z < x \lif \lnot
  \ORefut[\Th{T}](z,y))])].
\]
स्थूलपणे, $\ORProv[\Th{T}](y)$
असे सांगते की “$y$ ची
$\Th{T}$ मध्ये सिद्धता आहे,
आणि कमी संकेतांक असलेले
$y$ चे खंडन नाही.”
$\Th{T}$ सुसंगत असेल,
तर $\ORProv[\Th{T}](y)$ आणि
$\OProv[\Th{T}](y)$ ही दोन्ही
विधाने त्याच संख्यांसाठी
सत्य असतात. पण $\Th{T}$
मधील \emph{सिद्धतायोग्यतेच्या}
दृष्टीने (आता सत्य आणि
सिद्धतायोग्यता यांत फरक आहे
हे आपल्याला माहीत आहे)
त्यांचे गुणधर्म भिन्न असतात.
$\Th{T}$ \emph{वि}संगत
असेल, तर ती दोन्ही विधाने
त्याच संख्यांसाठी सत्य
\emph{नाहीत}!{}
($\ORProv[\Th{T}](y)$
याचा अर्थ अनेकदा “$y$ रोसर
अर्थाने सिद्धतायोग्य आहे”
असा घेतला जातो. पण आत्ताच
पाहिल्याप्रमाणे रोसर
सिद्धतायोग्यता हा
सिद्धतायोग्यतेचा काही
विशेष प्रकार नाही:
विसंगत उपपत्तीत काही
वाक्ये सिद्धतायोग्य
असतात, पण रोसर अर्थाने
सिद्धतायोग्य नसतात.
म्हणून हा शब्द गोंधळात
टाकू शकतो. तो टाळण्यासाठी
“$y$ शमोव्हेबल आहे”
असेही म्हणता येईल.)

स्थिरबिंदू उपपत्तीवरून
असे सूत्र $\mathsf{R}_\Th{T}$
असते की
\begin{equation}
  \Th{Q} \Proves \mathsf{R}_\Th{T} \liff \lnot \ORProv[\Th{T}](\gn{\mathsf{R}_\Th{T}}).
  \label{inc:inp:ros:RT}
\end{equation}
\ref{inc:inp:1in:thm:first-incompleteness}
च्या सिद्धतेपेक्षा इथे
आपला दावा अधिक सबळ आहे:
$\Th{T}$ सुसंगत असेल,
तर $\Th{T}$ हे
$\mathsf{R}_\Th{T}$ ला निष्पन्न
करत नाही, आणि $\Th{T}$
$\lnot \mathsf{R}_\Th{T}$ लाही
निष्पन्न करत नाही.
(म्हणजे $\omega$-सुसंगततेचे
गृहीतक लागत नाही.)

प्रथम $\Th{T} \Proves/
\mathsf{R}_{\Th{T}}$ हे दाखवू.
उलट ते सिद्ध होत असेल,
तर~$T$ पासून
$\mathsf{R}_{\Th{T}}$ ची
निष्पत्ती आहे;
तिची गोडेल संख्या $n$
मानू. $\OPrf[\Th{T}]$
हे $\Prf[\Th{T}]$ चे
$\Th{Q}$ मध्ये निरूपण
करत असल्यामुळे
$\Th{Q} \Proves
\OPrf[\Th{T}](\num{n}, \gn{\mathsf{R}_{\Th{T}}})$.
शिवाय प्रत्येक $k < n$
साठी, $\Th{T}$ सुसंगत
असल्यामुळे, $k$ ही
$\lnot \mathsf{R}_{\Th{T}}$
च्या निष्पत्तीची
गोडेल संख्या नसते.
म्हणून प्रत्येक $k < n$
साठी $\Th{Q} \Proves \lnot
\ORefut[\Th{T}](\num{k}, \gn{\mathsf{R}_{\Th{T}}})$.
\ref{inc:req:min:lem:less-nsucc}
वरून $\Th{Q} \Proves
\lforall[z][(z < \num{n} \lif \lnot \ORefut[\Th{T}](z,
  \gn{\mathsf{R}_{\Th{T}}}))]$.
त्यामुळे
\[\Th{Q} \Proves \lexists[x][(\OPrf[\Th{T}](x,\gn{\mathsf{R}_{\Th{T}}}) \land \lforall[z][(z
    < x \lif \lnot \ORefut[\Th{T}](z,\gn{\mathsf{R}_{\Th{T}}}))])],
\]
पण हेच
$\ORProv[\Th{T}](\gn{\mathsf{R}_{\Th{T}}})$
आहे. \ref{inc:inp:ros:RT} वरून
$\Th{Q}
\Proves \lnot \mathsf{R}_{\Th{T}}$.
$\Th{T}$ हा $\Th{Q}$ चा
विस्तार असल्यामुळे
$\Th{T}
\Proves \lnot \mathsf{R}_{\Th{T}}$
सुद्धा मिळते. पण आपण
$\Th{T} \Proves \mathsf{R}_{\Th{T}}$
असे गृहीत धरले होते.
त्यामुळे $\Th{T}$ विसंगत
ठरेल; हे प्रमेयाच्या
गृहीतकाच्या विरुद्ध आहे.

आता $\Th{T} \Proves/
\lnot \mathsf{R}_{\Th{T}}$ हे
दाखवू. उलट ते सिद्ध होत
असेल आणि
$\lnot \mathsf{R}_{\Th{T}}$
च्या निष्पत्तीची
गोडेल संख्या $n$ असेल
असे मानू. मग
$\Refut[\Th{T}](n, \Gn{\mathsf{R}_{\Th{T}}})$
खरे असते. $\ORefut[\Th{T}]$
हे $\Refut[\Th{T}]$ चे
$\Th{Q}$ मध्ये निरूपण
करत असल्यामुळे
$\Th{Q} \Proves
\ORefut[\Th{T}](\num{n}, \gn{\mathsf{R}_{\Th{T}}})$.
मग $\Th{T}$ हे
$\mathsf{R}_{\Th{T}}$ सुद्धा
निष्पन्न करेल आणि
त्यामुळे विसंगत ठरेल,
हे पुन्हा दाखवू. कारण
\begin{align*}
\Th{Q} & \Proves \mathsf{R}_{\Th{T}} \liff \lnot \ORProv[\Th{T}](\gn{\mathsf{R}_{\Th{T}}}),
\intertext{आणि $\Th{T}$ हा~$\Th{Q}$ चा विस्तार असल्यामुळे पुढील दाखवणे पुरेसे आहे:}
\Th{Q} & \Proves \lnot \ORProv[\Th{T}](\gn{\mathsf{R}_{\Th{T}}}).
\end{align*}
वाक्य $\lnot
\ORProv[\Th{T}](\gn{\mathsf{R}_{\Th{T}}})$,
म्हणजे
\begin{align*}
  \lnot & \lexists[x][(\OPrf[\Th{T}](x,\gn{\mathsf{R}_{\Th{T}}}) \land
    \lforall[z][(z < x \lif \lnot \ORefut[\Th{T}](z,\gn{\mathsf{R}_{\Th{T}}}))])],
  \intertext{हे तर्कशास्त्रीय दृष्ट्या पुढील विधानाशी सममूल्य आहे:}
  & \lforall[x][(\OPrf[\Th{T}](x,\gn{\mathsf{R}_{\Th{T}}}) \lif
    \lexists[z][(z < x \land \ORefut[\Th{T}](z,\gn{\mathsf{R}_{\Th{T}}}))])].
\end{align*}
$\Th{Q}$ काय निष्पन्न
करते याविषयीची तथ्ये वापरून
अनौपचारिक युक्तिवाद करू.
$x$ स्वेच्छेने निवडलेला
आणि $\OPrf[\Th{T}](x,
\gn{\mathsf{R}_{\Th{T}}})$ खरा
आहे असे मानू. आपल्याला
आधीच माहीत आहे की
$\Th{T} \Proves/ \mathsf{R}_{\Th{T}}$.
म्हणून प्रत्येक $k$
साठी $\Th{Q} \Proves \lnot
\OPrf[\Th{T}](\num{k}, \gn{\mathsf{R}_{\Th{T}}})$.
त्यामुळे प्रत्येक $k$
साठी $\eq/[x][\num{k}]$
मिळते. विशेषतः (a)
$\eq/[x][\num{n}]$
मिळते. तसेच
$\lnot(\eq[x][\num{0}]
\lor \eq[x][\num{1}] \lor \dots
\lor \eq[x][\num{n-1}])$
मिळते. म्हणून
\ref{inc:req:min:lem:less-nsucc}
वरून (b) $\lnot(x < \num{n})$.
\ref{inc:req:min:lem:trichotomy}
वरून $\num{n} < x$.
$\Th{Q} \Proves
\ORefut[\Th{T}](\num{n}, \gn{\mathsf{R}_{\Th{T}}})$
असल्यामुळे
$\num{n} < x \land
\ORefut[\Th{T}](\num{n}, \gn{\mathsf{R}_{\Th{T}}})$
मिळते, आणि त्यावरून
$\lexists[z][(z < x
\land \ORefut[\Th{T}](z,\gn{\mathsf{R}_{\Th{T}}}))]$
मिळते. $x$ स्वेच्छेने
निवडलेला असल्यामुळे
अपेक्षेप्रमाणे पुढील मिळते:
\[\lforall[x][(\OPrf[\Th{T}](x,\gn{\mathsf{R}_{\Th{T}}}) \lif
  \lexists[z][(z < x \land \ORefut[\Th{T}](z,\gn{\mathsf{R}_{\Th{T}}}))])].
\]
\end{proof}

\begin{prob}
नैसर्गिक संख्यांचे दोन संच
$A$ आणि $B$ यांना
\emph{संगणनक्षमपणे अविलगनीय}
म्हणतात, जर $A \subseteq X$
आणि $B \subseteq \Complement{X}$
असेल असा कोणताही
निर्णेय संच~$X$
नसेल ($\Complement{X}$
म्हणजे~$X$ चा पूरक
$\Nat \setminus X$).
$\Th{T}$ ही~$\Th{Q}$
ची सुसंगत,
स्वयंसिद्धकीकरणीय विस्तार-उपपत्ती
असू द्या. $\Th{T}$
मध्ये सिद्ध होणाऱ्या
वाक्यांच्या गोडेल
संख्यांचा संच $A$,
आणि $\Th{T}$ मध्ये
खंडित होणाऱ्या वाक्यांच्या
गोडेल संख्यांचा संच $B$
असेल असे मानू. $A$ आणि
$B$ हे संगणनक्षमपणे
अविलगनीय आहेत हे सिद्ध करा.
\end{prob}












\section{गोडेलच्या मूळ शोधपत्राशी तुलना}\label{inc:inp:gop:sec}

गोडेलच्या 1931 मधील शोधपत्राला थोडा वेळ
देणे उपयुक्त ठरेल. त्याच्या प्रस्तावनेत आपण
नुकत्याच पाहिलेल्या कल्पनांची रूपरेषा आहे.
शोधपत्र नुसते चाळले तरी प्रत्येक टप्प्यावर
काय चालले आहे हे सहज दिसते: प्रथम गोडेल
आकारिक पद्धत $P$ (विन्यास, स्वयंसिद्धके,
सिद्धतांचे नियम) वर्णन करतो; मग आदिम
पुनरावर्ती फलने आणि संबंध परिभाषित करतो;
त्यानंतर $x B y$ हा संबंध आदिम पुनरावर्ती
आहे हे दाखवतो आणि आदिम पुनरावर्ती फलने
व संबंध $\Th{P}$ मध्ये निरूपित होतात
असा युक्तिवाद करतो. मग वर दिल्याप्रमाणे
अपूर्णता प्रमेय सिद्ध करतो. तिसऱ्या विभागात,
सिद्ध करता न येणारे विधान अंकगणिताच्या
भाषेतील वाक्य म्हणून घेता येते हे दाखवतो.
यातूनच $\beta$-उपपत्ती उद्भवली; पुनरावर्ती
फलने $\Th{Q}$ मध्ये निरूपणीय आहेत हे
दाखवताना अनुक्रम हाताळण्यासाठी आपणही ती
वापरली. गोडेल एवढेच पुरेसे ठरणाऱ्या
किमान स्वयंसिद्धकसंचाची वेगळी ओळख करून
देत नाही, पण आता $\Th{Q}$ ते काम करते
हे आपल्याला माहीत आहे. अखेरीस चौथ्या
विभागात तो दुसऱ्या अपूर्णता प्रमेयाच्या
सिद्धतेची रूपरेषा देतो.












\section[PA साठीच्या निष्पन्नता अटी]{$\Th{PA}$ साठीच्या निष्पन्नता अटी}\label{inc:inp:prc:sec}

पेआनो अंकगणित, म्हणजे $\Th{PA}$, ही
सर्व सूत्रे साठी विगमनाची स्वयंसिद्धके
जोडून मिळणारी $\Th{Q}$ ची विस्तार-उपपत्ती
आहे. म्हणजे प्रत्येक सूत्र~$\varphi$
साठी $\Th{Q}$ मध्ये पुढील रूपाचे
स्वयंसिद्धक जोडतात:
\[(\varphi(0) \land \lforall[x][(\varphi(x) \lif \varphi(x'))]) \lif \lforall[x][\varphi(x)]
\]
या रूपबंधात इतर मुक्त चल असतील, तर आरंभी त्यांच्यावर
सार्विक संख्यापक जोडून वाक्यरूपाचे स्वयंसिद्धक घेतो.
हे खरे तर एक \emph{रूपबंध} आहे, म्हणजे
अनंत स्वयंसिद्धके आहेत, हे लक्षात घ्या
(आणि $\Th{PA}$ चे सांत
स्वयंसिद्धकसंचाने स्वयंसिद्धकीकरण करता येत
{\em नाही} असेही सिद्ध होते).
पण चिन्हमाला विगमनाच्या
स्वयंसिद्धकाचे उदाहरण आहे
की नाही हे प्रभावीपणे ठरवता
येत असल्यामुळे $\Th{PA}$
चा स्वयंसिद्धकसंच संगणनक्षम
आहे. $\Th{PA}$ ही~$\Th{Q}$
पेक्षा बरीच अधिक सबल
उपपत्ती आहे. उदाहरणार्थ,
नेहमीच्या पद्धतीने विगमन
वापरून बेरीज आणि गुणाकार
क्रमनिरपेक्ष आहेत हे
सहज सिद्ध करता येते.
खरे तर बहुतेक सांत
संख्यासिद्धांतिक आणि
संयोजनात्मक युक्तिवाद
$\Th{PA}$ मध्ये मांडता येतात.

$\Th{PA}$ ही संगणनक्षमपणे
स्वयंसिद्धकीकृत असल्यामुळे
$\Prf[\Th{PA}](x,y)$ हे
निष्पन्नता विधेय
संगणनक्षम आहे आणि म्हणून
$\Th{Q}$ मध्ये निरूपित
होते (म्हणूनच
$\Th{PA}$ मध्येही).
पूर्वीप्रमाणे, या संबंधाचे
निरूपण करणाऱ्या सूत्रासाठी
$\OPrf[\Th{PA}](x,y)$
वापरू. $\OProv[\Th{PA}](y)$
हे सूत्र
$\lexists[x][\OPrf[\Th{PA}](x,y)]$
असे मानू. त्याचा सहज
अर्थ असा: “$y$ हे
$\Th{PA}$ च्या स्वयंसिद्धकांपासून
निष्पन्न करता येण्याजोगे आहे.”
$\Th{Q}$ च्या स्वयंसिद्धकांपेक्षा
थोडे अधिक सामर्थ्य लागते,
कारण आपण वापरत असलेली
उपपत्ती या निष्पन्नता
विधेयाविषयीची काही मूलभूत
तथ्ये निष्पन्न करू शकली
पाहिजे. पुढील अटींसाठी प्रत्यक्ष आकारिक सिद्धतांच्या प्रमाण
संकेतांकीकरणावर आधारित विधेय निवडले आहे. केवळ निर्णेय
संबंधाचे निरूपण करणारे मनमानी सूत्र घेतल्याने या अटी
आपोआप मिळत नाहीत. नेमकी आवश्यक तथ्ये
पुढीलप्रमाणे:
\begin{enumerate}
\item[P1.] जर $\Th{PA} \Proves \varphi$, तर $\Th{PA} \Proves
  \OProv[\Th{PA}](\gn{\varphi})$.
\item[P2.] सर्व सूत्रे $\varphi$ आणि $\psi$ साठी,
  \[  \Th{PA} \Proves \OProv[\Th{PA}](\gn{\varphi \lif \psi}) \lif
  (\OProv[\Th{PA}](\gn{\varphi}) \lif \OProv[\Th{PA}](\gn{\psi})).
  \]
\item[P3.] प्रत्येक सूत्र~$\varphi$ साठी,
  \[  \Th{PA} \Proves \OProv[\Th{PA}](\gn{\varphi})
  \lif \OProv[\Th{PA}](\gn{\OProv[\Th{PA}](\gn{\varphi})}).
  \]
\end{enumerate}
हे तीन गुणधर्म खरे आहेत
याची खात्री करण्यासाठी
$\OProv[\Th{PA}](y)$
या सूत्राचे बारकाईने
वर्णन करून संबंधित आकारिक
निष्पत्त्यांचे वर्णन
$\Th{PA}$ च्या स्वयंसिद्धकांनी
करावे लागते. अटी (1) आणि~(2)
सोप्या आहेत; खरे काम
अट~(3) साठी लागते
(त्यासाठी नेमके काय
करावे लागेल याचा विचार करा
\dots). तपशील मांडणे
कंटाळवाणे आणि फारसे
रंजक नसल्यामुळे
$\Th{PA}$ मध्ये वरचे
तीन गुणधर्म आहेत हे
येथे आपण गृहीत धरू.
$\OProv[\Th{PA}](y)$
ची योग्य निवड पुढील
अटही पूर्ण करेल:
\begin{enumerate}
\item[P4.] जर $\Th{PA} \Proves \OProv[\Th{PA}](\gn{\varphi})$, तर
  $\Th{PA} \Proves \varphi$.
\end{enumerate}
पण हे तथ्य आपल्याला
लागणार नाही. ही शेवटची अट आधीच्या तीन आकारिक अटींचा
निष्कर्ष नाही; तिच्यासाठी पेआनो अंकगणिताचे मानक प्रतिमानातील
सत्य आणि प्रत्यक्ष सिद्धतासंकेतांकांचे प्रमाण विधेय वापरले आहे.

\begin{digress}
तपशील वगळण्याबाबत
गोडेलनेही आपल्यासारखाच
आळस केला. 1931 च्या
शोधपत्राच्या अखेरीस
त्याने दुसऱ्या अपूर्णता
प्रमेयाच्या सिद्धतेची
रूपरेषा दिली आणि
तपशील पुढील शोधपत्रात
देण्याचे आश्वासन दिले.
ते तपशील त्याने
नंतर प्रकाशित केले
नाहीत; युक्तिवाद समजणाऱ्या
सर्वांनाच तो पूर्ण करता
येईल असे वाटत असल्यामुळे
तपशील भरण्याची गरजही
त्याला पडली नाही.
\end{digress}












\section{दुसरे अपूर्णता प्रमेय}\label{inc:inp:2in:sec}

$\Th{PA}$ स्वतःची सुसंगतता सिद्ध करू
शकत नाही, हा दावा कसा व्यक्त करायचा?
$\Th{PA}$ विसंगत आहे असे म्हणणे
म्हणजे $\Th{PA} \Proves \eq[0][1]$
असे म्हणणे. सुसंगतता-विधानासाठी
असत्यता-स्थिराचा संकेतांक घेतला तरी
सुसंगततेचा आशय तोच राहतो; या
सिद्धतेत $\OCon[\Th{PA}]$ हे
वाक्य $\lnot
\OProv[\Th{PA}](\gn{\lfalse})$
मानू. मग पुढील प्रमेय अपेक्षित
निकाल देते:

\begin{thm}
\label{inc:inp:2in:thm:second-incompleteness}
$\Th{PA}$ सुसंगत असेल, तर
$\Th{PA}$ हे
$\OCon[\Th{PA}]$ निष्पन्न
करत नाही.
\end{thm}

हे प्रमेय $\OCon[\Th{PA}]$
च्या विशिष्ट निरूपणावर
(म्हणजे $\OProv[\Th{PA}](y)$
च्या विशिष्ट निरूपणावर)
अवलंबून आहे, हे लक्षात
घेणे महत्त्वाचे आहे.
आपण फक्त एवढेच वापरणार
आहोत की $\OProv[\Th{PA}](y)$
चे निरूपण तीन
निष्पन्नता अटी पूर्ण
करते. म्हणून असे गुणधर्म
असलेले निष्पन्नता
विधेय असलेल्या, क्यूचा विस्तार असलेल्या कोणत्याही
सुसंगत उपपत्तीला हे प्रमेय लागू
होते.

गोडेलच्या युक्तिवादाची
रूपरेषा वाचणे उपयुक्त आहे;
प्रमेय त्यातून नेमके
समारोपाप्रमाणे निघते.
युक्तिवाद असा आहे.
\ref{inc:inp:1in:thm:first-incompleteness}
च्या सिद्धतेत आपण तयार
केलेले गोडेल-वाक्य
$\mathsf{G}_\Th{PA}$ असू द्या.
“$\Th{PA}$ सुसंगत
असेल, तर $\Th{PA}$
हे $\mathsf{G}_\Th{PA}$ ला
निष्पन्न करत नाही”
हे आपण दाखवले आहे.
हे $\Th{PA}$
\emph{मध्येच} आकारिक
केल्यास पुढील विधानाची
सिद्धता मिळते:
\[\OCon[\Th{PA}] \lif \lnot \OProv[\Th{PA}](\gn{\mathsf{G}_\Th{PA}}).
\]
आता $\Th{PA}$ हे
$\OCon[\Th{PA}]$
निष्पन्न करते असे
मानू. मग ते $\lnot
\OProv[\Th{PA}](\gn{\mathsf{G}_\Th{PA}})$
सुद्धा निष्पन्न
करते. पण $\mathsf{G}_\Th{PA}$
हे गोडेल-वाक्य असल्यामुळे
हे $\mathsf{G}_\Th{PA}$ शी
सममूल्य आहे. म्हणून
$\Th{PA}$ हे
$\mathsf{G}_\Th{PA}$
निष्पन्न करते.

पण $\Th{PA}$ सुसंगत
असेल, तर ते
$\mathsf{G}_\Th{PA}$ ला
निष्पन्न करत
नाही, हे आपल्याला
माहीत आहे!{}
म्हणून $\Th{PA}$
सुसंगत असेल, तर
$\OCon[\Th{PA}]$
सुद्धा निष्पन्न
करू शकत नाही.

युक्तिवाद अधिक नेमका
करण्यासाठी $\mathsf{G}_\Th{PA}$
हे~$\Th{PA}$ चे
गोडेल-वाक्य मानू.
निष्पन्नता अटी
(P1)--(P3) वापरून
$\Th{PA}$ हे
$\OCon[\Th{PA}] \lif
\mathsf{G}_\Th{PA}$ निष्पन्न
करते हे दाखवू.
त्यावरून $\Th{PA}$
हे $\OCon[\Th{PA}]$
निष्पन्न करत नाही
हे मिळेल. पुढे
$\Th{PA}$ मधील
सिद्धतेची रूपरेषा आहे.
(सोयीसाठी
$\Th{PA}$ हे अधोलिखित
वगळू.)
\begin{align}
& \mathsf{G} \liff \lnot \OProv(\gn{\mathsf{G}}) \label{inc:inp:2in:G2-1}\\
& \qquad\text{$\mathsf{G}$ हे गोडेल-वाक्य आहे}\notag \\
& \mathsf{G} \lif \lnot \OProv(\gn{\mathsf{G}}) \label{inc:inp:2in:G2-2}\\
  & \qquad\text{\ref{inc:inp:2in:G2-1} वरून} \notag\\
& \mathsf{G} \lif
  (\OProv(\gn{\mathsf{G}}) \lif \lfalse) \label{inc:inp:2in:G2-3}\\
  & \qquad\text{\ref{inc:inp:2in:G2-2} व तर्कशास्त्रावरून}\notag\\
& \OProv(\gn{
    \mathsf{G} \lif
    (\OProv(\gn{\mathsf{G}}) \lif \lfalse)
  }) \label{inc:inp:2in:G2-4}\\
  & \qquad\text{\ref{inc:inp:2in:G2-3} आणि अट P1 वरून} \notag\\
& \OProv(\gn{\mathsf{G}}) \lif
  \OProv(\gn{
    (\OProv(\gn{\mathsf{G}}) \lif \lfalse)
    }) \label{inc:inp:2in:G2-5}\\
  & \qquad\text{\ref{inc:inp:2in:G2-4} आणि अट P2 वरून} \notag\\
& \OProv(\gn{\mathsf{G}}) \lif (\OProv(\gn{\OProv(\gn{\mathsf{G}})}) \lif \OProv(\gn{\lfalse})) \label{inc:inp:2in:G2-6}\\
  & \qquad\text{\ref{inc:inp:2in:G2-5}, अट P2 आणि तर्कशास्त्रावरून} \notag\\
& \OProv(\gn{\mathsf{G}}) \lif
  \OProv(\gn{\OProv(\gn{\mathsf{G}})}) \label{inc:inp:2in:G2-7}\\
   & \qquad\text{अट P3 वरून} \notag\\
& \OProv(\gn{\mathsf{G}}) \lif \OProv(\gn{\lfalse}) \label{inc:inp:2in:G2-8}\\
  & \qquad \text{\ref{inc:inp:2in:G2-6}, \ref{inc:inp:2in:G2-7} व तर्कशास्त्रावरून}\notag\\
& \OCon \lif \lnot \OProv(\gn{\mathsf{G}}) \label{inc:inp:2in:G2-9}\\
  & \qquad\text{\ref{inc:inp:2in:G2-8} च्या प्रतिवर्ती विधानावरून आणि $\OCon \ident \lnot \OProv(\gn{\lfalse})$ वरून}\notag \\
& \OCon \lif \mathsf{G} \notag\\
  & \qquad\text{\ref{inc:inp:2in:G2-1}, \ref{inc:inp:2in:G2-9} व तर्कशास्त्रावरून}\notag
\end{align}
वरील तर्कशास्त्राचा
वापर केवळ विधानीय
तर्कशास्त्रातील प्राथमिक
तथ्यांपुरता आहे.
उदाहरणार्थ,
\ref{inc:inp:2in:G2-3} मध्ये
$\Proves \lnot\varphi \liff (\varphi\lif
\lfalse)$ वापरले आहे,
आणि \ref{inc:inp:2in:G2-8} मध्ये
$\varphi \lif (\psi \lif \chi), \varphi \lif \psi
\Proves \varphi \lif \chi$ वापरले
आहे. \ref{inc:inp:2in:G2-5} आणि
\ref{inc:inp:2in:G2-6} मध्ये अट~P2
वापरताना तिची पुढील
उदाहरणे वापरली आहेत:
$\OProv(\gn{\varphi \lif \psi}) \lif
(\OProv(\gn{\varphi}) \lif \OProv(\gn{\psi}))$.
पहिल्यात $\varphi \ident
\mathsf{G}$ आणि $\psi \ident
\OProv(\gn{\mathsf{G}}) \lif \lfalse$;
दुसऱ्यात $\varphi
\ident \OProv(\gn{\mathsf{G}})$
आणि $\psi \ident \lfalse$.

दुसऱ्या अपूर्णता प्रमेयाचे
अधिक अमूर्त रूप असे:

\begin{thm}
\label{inc:inp:2in:thm:second-incompleteness-gen}
$\Th{T}$ ही~$\Th{Q}$ चा
विस्तार असलेली कोणतीही
सुसंगत, स्वयंसिद्धकीकृत
उपपत्ती असू द्या.
$\OProv[\Th{T}](y)$
हे सूत्र $\Th{T}$
साठीच्या
निष्पन्नता अटी
P1--P3 पूर्ण करत असेल,
तर $\Th{T}$ हे
$\OCon[\Th{T}]$
निष्पन्न करत नाही.
\end{thm}

\begin{prob}
$\Th{PA}$ हे
$\mathsf{G}_{\Th{PA}} \lif \OCon[\Th{PA}]$
निष्पन्न करते हे
दाखवा.
\end{prob}

\begin{digress}
या कथेतून असा बोध होतो:
गणितासाठीची कोणतीही
“स्वीकारार्ह” सुसंगत
उपपत्ती स्वतःचे
सुसंगतता-विधान
निष्पन्न करू शकत
नाही. $\Th{T}$ ही
गणिताची उपपत्ती असून
तिच्यात $\Th{Q}$ आणि
हिल्बर्टचे “सांतवादी”
युक्तिवाद (त्यांचा नेमका
अर्थ काहीही असो)
समाविष्ट आहेत असे
मानू. मागील प्रमेयातील सुसंगततेची आणि निवडलेल्या
सिद्धतायोग्यता विधेयाची सर्व अटीही पूर्ण होतात असे
गृहीत धरा. मग संपूर्ण
$\Th{T}$ हे
$\Th{T}$ चे
सुसंगतता-विधान
निष्पन्न करू
शकत नाही. त्यामुळे
तिचा सांतवादी भागही
$\Th{T}$ चे
सुसंगतता-विधान
निष्पन्न करू
शकत नाही, हे तर
अधिकच स्पष्ट आहे.
त्या अर्थाने
“सर्व गणितासाठी”
सुसंगततेची सांतवादी
सिद्धता असू शकत नाही.

“सांतवादी” या संज्ञेचा
अर्थ लावण्यात थोडी
मोकळीक आहे. गोडेलने
1931 च्या शोधपत्रात
हिल्बर्टला आकारिक
करायचे असलेल्या
गणिताबाहेरही आपण
“सांतवादी” म्हणू
असे काही असू शकेल,
ही शक्यता मान्य केली.
पण ही गोडेलची
उदार भूमिका होती:
आज सांतवादी म्हणता
येईल पण, उदाहरणार्थ,
निवडीच्या स्वयंसिद्धकासह
झर्मेलो--फ्रेंकेल संच
उपपत्ती $\Th{ZFC}$
मध्ये आकारिक करता
येणार नाही असे
काही सापडेल, हे
कल्पना करणे कठीण आहे.
\end{digress}












\section{लोबचे प्रमेय}\label{inc:inp:lob:sec}

उपपत्ती~$\Th{T}$ साठीचे गोडेल-वाक्य
हा $\lnot \OProv[\Th{T}](y)$ चा
स्थिरबिंदू असतो; म्हणजे असे
वाक्य~$\mathsf{G}$ असते की
\[\Th{T} \Proves \lnot \OProv[\Th{T}](\gn{\mathsf{G}}) \liff \mathsf{G}.
\]
सुसंगत उपपत्तीत ते निष्पन्न करता येण्याजोगे नसते. कारण
$\Th{T} \Proves \mathsf{G}$ असेल, तर
(a) निष्पन्नता अट~(1) वरून
$\Th{T} \Proves \OProv[\Th{T}](\gn{\mathsf{G}})$,
आणि (b) $\Th{T}
\Proves \mathsf{G}$ व $\Th{T} \Proves \lnot \OProv[\Th{T}](\gn{\mathsf{G}})
\liff \mathsf{G}$ यांवरून
$\Th{T} \Proves \lnot \OProv[\Th{T}](\gn{\mathsf{G}})$
मिळते. त्यामुळे $\Th{T}$ विसंगत
ठरेल. आता $\OProv[\Th{T}](y)$
च्या स्थिरबिंदूची स्थिती विचारणे
स्वाभाविक आहे. म्हणजे असे
वाक्य~$\mathsf{H}$ घेऊ की
\[\Th{T} \Proves \OProv[\Th{T}](\gn{\mathsf{H}}) \liff \mathsf{H}.
\]
ते निष्पन्न करता येण्याजोगे असेल, तर
अट~(1) वरून
$\Th{T} \Proves \OProv[\Th{T}](\gn{\mathsf{H}})$.
पण वरील सममूल्यतेला
निष्कर्षणाचा नियम लावूनही
हाच निष्कर्ष मिळतो.
म्हणून गोडेल-वाक्याच्या
बाबतीतला युक्तिवाद येथे
$\Th{T}$ विसंगत आहे
असे दाखवत नाही.
अर्थात यावरून
$\Th{T}$ हे
$\mathsf{H}$ ला \emph{खरोखर}
निष्पन्न करते असेही
सिद्ध होत नाही.

हा प्रश्न थोडा व्यापक
केल्यास पुढे जाता येते.
स्थिरबिंदूच्या
सममूल्यतेची
डावीकडून उजवीकडची
दिशा
$\OProv[\Th{T}](\gn{\mathsf{H}}) \lif \mathsf{H}$
ही \emph{प्रतिबिंबन तत्त्व}
म्हणवणाऱ्या सर्वसाधारण
रूपबंधाचे एक उदाहरण आहे:
$\OProv[\Th{T}](\gn{\varphi}) \lif \varphi$.
त्याला असे नाव दिले
आहे कारण ते एका
अर्थाने $\Th{T}$
स्वतः काय निष्पन्न
करू शकते यावर
“प्रतिबिंबन” करू शकते
असे व्यक्त करते.
मूलतः ते कोणत्याही~$\varphi$
साठी “$\Th{T}$
$\varphi$ ला निष्पन्न
करू शकत असेल, तर
$\varphi$ सत्य आहे”
असे म्हणते. हे
अर्थात फक्त सत्यनिष्ठ
उपपत्त्यांसाठीच खरे
असते. त्यामुळे
उपपत्त्या सर्वसाधारणपणे
त्याचे प्रत्येक उदाहरण
निष्पन्न करणार नाहीत
असे वाटते. मग
पुरेशी प्रबळ आणि
निष्पन्नता अटी
पूर्ण करणारी उपपत्ती
कोणती उदाहरणे
निष्पन्न करू शकते?
$\varphi$ स्वतः
निष्पन्न करता येण्याजोगे असेल
अशी सर्व उदाहरणे
नक्कीच. पुढील
निकाल दाखवतो की
एवढीच उदाहरणे शक्य
आहेत.

\begin{thm}\label{inc:inp:lob:thm:lob}
$\Th{T}$ ही~$\Th{Q}$
चा विस्तार असलेली
स्वयंसिद्धकीकरणीय
उपपत्ती असू द्या,
आणि $\OProv[\Th{T}](y)$
हे सूत्र \ref{inc:inp:prc:sec}
मधील P1--P3 अटी
पूर्ण करत असेल.
$\Th{T}$ हे
$\OProv[\Th{T}](\gn{\varphi}) \lif \varphi$
निष्पन्न करत
असेल, तर प्रत्यक्षात
$\Th{T}$ हे
$\varphi$ सुद्धा
निष्पन्न करते.
\end{thm}

दुसऱ्या शब्दांत,
$\Th{T} \Proves/ \varphi$
असेल, तर
$\Th{T} \Proves/
\OProv[\Th{T}](\gn{\varphi}) \lif \varphi$.
हा निकाल लोबचे
प्रमेय म्हणून ओळखला
जातो.

\begin{explain}
लोबच्या प्रमेयामागील
युक्ती समजण्यासाठी
सांताक्लॉज अस्तित्वात
आहे हे “सिद्ध”
करणारा चतुर युक्तिवाद
पाहा. (हा निष्कर्ष
आवडत नसेल, तर
तुम्हाला हवा असलेला
दुसरा निष्कर्ष
त्याच्या जागी
घेता येईल.)
युक्तिवाद असा:
\begin{enumerate}
\item $X$ हे “$X$
  सत्य असेल, तर
  सांताक्लॉज
  अस्तित्वात आहे”
  हे वाक्य मानू.
\item $X$ सत्य आहे
  असे धरू.
\item मग त्यात
  म्हटलेले खरे आहे:
  $X$ सत्य असेल,
  तर सांताक्लॉज
  अस्तित्वात आहे.
\item आपण $X$ सत्य
  आहे असे धरलेच
  आहे; म्हणून
  (2) आणि~(3)
  वर निष्कर्षणाचा
  नियम लावून
  सांताक्लॉज
  अस्तित्वात आहे
  असा निष्कर्ष
  काढू शकतो.
\item (2) मधील
  “$X$ सत्य आहे”
  या गृहीतकावरून
  (4) मधील
  “सांताक्लॉज
  अस्तित्वात आहे”
  हे आपण निष्पन्न
  केले. अटीची
  सिद्धता वापरून
  “$X$ सत्य असेल,
  तर सांताक्लॉज
  अस्तित्वात आहे”
  हे दाखवले.
\item पण हेच
  वाक्य~$X$ आहे.
  म्हणजे आपण
  $X$ सत्य आहे
  हे दाखवले.
\item मग वरील
  (2)--(4) च्या
  युक्तिवादाने
  सांताक्लॉज
  अस्तित्वात आहे.
\end{enumerate}
या कल्पनेचे
आकारिकीकरण करताना
“सत्य आहे” ऐवजी
“निष्पन्न करता येण्याजोगे आहे”
आणि “सांताक्लॉज
अस्तित्वात आहे”
ऐवजी~$\varphi$
घेतल्यास लोबच्या
प्रमेयाची सिद्धता
मिळते. युक्ती अशी:
स्थिरबिंदू उपपत्ती
सूत्र~
$\OProv[\Th{T}](y) \lif \varphi$
याला लावायची.
त्याचा स्थिरबिंदू
वरील रूपरेषेतील
वाक्य~$X$
याच्याशी जुळतो.
\end{explain}

\begin{proof}[याची सिद्धता: \ref{inc:inp:lob:thm:lob}]
$\varphi$ हे असे
वाक्य
असू द्या की
$\Th{T}$ हे
$\OProv[\Th{T}](\gn{\varphi}) \lif \varphi$
निष्पन्न करते.
$\psi(y)$ हे
सूत्र~
$\OProv[\Th{T}](y)
\lif \varphi$ मानू.
स्थिरबिंदू उपपत्ती
वापरून असे
वाक्य~$\theta$
शोधू की
$\Th{T}$ हे
$\theta \liff \psi(\gn{\theta})$
निष्पन्न करते.
मग पुढीलपैकी
प्रत्येक विधान
$\Th{T}$ मध्ये
निष्पन्न करता येण्याजोगे आहे:
\begin{align}
  & \theta \liff (\OProv[\Th{T}](\gn{\theta}) \lif \varphi) \label{inc:inp:lob:L-1}\\
  & \qquad \text{$\theta$ हा~$\psi(y)$ चा स्थिरबिंदू आहे}\notag \\
  & \theta \lif (\OProv[\Th{T}](\gn{\theta}) \lif \varphi) \label{inc:inp:lob:L-2}\\
  & \qquad\text{\ref{inc:inp:lob:L-1} वरून}\notag\\
  & \OProv[\Th{T}](\gn{\theta \lif (\OProv[\Th{T}](\gn{\theta}) \lif \varphi)}) \label{inc:inp:lob:L-3}\\
  & \qquad \text{\ref{inc:inp:lob:L-2} आणि अट P1 वरून}\notag \\
  & \OProv[\Th{T}](\gn{\theta}) \lif \OProv[\Th{T}](\gn{\OProv[\Th{T}](\gn{\theta}) \lif \varphi})
  \label{inc:inp:lob:L-4}\\
  &\qquad \text{\ref{inc:inp:lob:L-3} व अट P2 वरून}\notag \\
  & \OProv[\Th{T}](\gn{\theta}) \lif (\OProv[\Th{T}](\gn{\OProv[\Th{T}](\gn{\theta})}) \lif \OProv[\Th{T}](\gn{\varphi})) \label{inc:inp:lob:L-5}\\
  &\qquad \text{\ref{inc:inp:lob:L-4} व अट P2 पुन्हा वापरून} \notag\\
& \OProv[\Th{T}](\gn{\theta}) \lif \OProv[\Th{T}](\gn{\OProv[\Th{T}](\gn{\theta})}) \label{inc:inp:lob:L-6}\\
  & \qquad\text{निष्पन्नता अट P3 वरून} \notag\\
  & \OProv[\Th{T}](\gn{\theta}) \lif \OProv[\Th{T}](\gn{\varphi}) \label{inc:inp:lob:L-7} \\
  &\qquad\text{\ref{inc:inp:lob:L-5} आणि \ref{inc:inp:lob:L-6} वरून}\notag\\
  & \OProv[\Th{T}](\gn{\varphi}) \lif \varphi \label{inc:inp:lob:L-8}\\
  &\qquad\text{प्रमेयाच्या गृहीतकावरून} \notag\\
  & \OProv[\Th{T}](\gn{\theta}) \lif \varphi \label{inc:inp:lob:L-9}\\
  &\qquad\text{\ref{inc:inp:lob:L-7} आणि \ref{inc:inp:lob:L-8} वरून}\notag\\
  & (\OProv[\Th{T}](\gn{\theta}) \lif \varphi) \lif \theta \label{inc:inp:lob:L-10}\\
  & \qquad \text{\ref{inc:inp:lob:L-1} वरून}\notag \\
  & \theta \label{inc:inp:lob:L-11}\\
  & \qquad\text{\ref{inc:inp:lob:L-9} आणि \ref{inc:inp:lob:L-10} वरून}\notag \\
  & \OProv[\Th{T}](\gn{\theta}) \label{inc:inp:lob:L-12}\\
  & \qquad\text{\ref{inc:inp:lob:L-11} आणि अट~P1 वरून}\notag \\
  & \varphi \qquad\qquad\text{\ref{inc:inp:lob:L-9} आणि \ref{inc:inp:lob:L-12} वरून}\notag
\end{align}
\end{proof}

लोबचे प्रमेय मिळाल्यावर
दुसऱ्या अपूर्णता प्रमेयाची
लहान सिद्धता मिळते
(P1--P3 अटी पूर्ण
करणारे निष्पन्नता
विधेय असलेल्या उपपत्त्यांसाठी):
$\Th{T} \Proves
\OProv[\Th{T}](\gn{\lfalse}) \lif \lfalse$
असेल, तर
$\Th{T} \Proves \lfalse$.
$\Th{T}$ सुसंगत
असेल, तर
$\Th{T} \Proves/ \lfalse$.
म्हणून
$\Th{T}
\Proves/ \OProv[\Th{T}](\gn{\lfalse}) \lif \lfalse$,
म्हणजे
$\Th{T} \Proves/
\OCon[\Th{T}]$.
$\OProv[\Th{T}](x)$
चा स्थिरबिंदू~$\mathsf{H}$
हा निष्पन्न करता येण्याजोगे
आहे हेही यावरून
दाखवता येते.
कारण
\begin{align*}
  \Th{T} & \Proves \OProv[\Th{T}](\gn{\mathsf{H}}) \liff \mathsf{H}\\
  \intertext{यावरून विशेषतः}
    \Th{T} & \Proves \OProv[\Th{T}](\gn{\mathsf{H}}) \lif \mathsf{H}
\end{align*}
आणि म्हणून लोबच्या
प्रमेयावरून
$\Th{T} \Proves \mathsf{H}$.





\begin{prob}
$\Th{T}$ ही संगणनक्षमपणे
स्वयंसिद्धकीकृत उपपत्ती
असू द्या, आणि
$\OProv[\Th{T}]$ हे
$\Th{T}$ साठीचे
निष्पन्नता
विधेय असू द्या.
पुढील चार विधाने
विचारात घ्या:
\begin{enumerate}
\item जर $T \Proves \varphi$, तर $T \Proves \OProv[\Th{T}](\gn{\varphi})$.
\item $T \Proves \varphi \lif \OProv[\Th{T}](\gn{\varphi})$.
\item जर $T \Proves \OProv[\Th{T}](\gn{\varphi})$, तर $T \Proves \varphi$.
\item $T \Proves \OProv[\Th{T}](\gn{\varphi}) \lif \varphi$
\end{enumerate}
ही प्रत्येक विधाने
कोणत्या अटींवर
खरी असतात?
\end{prob}












\section{सत्याची अपरिभाष्यता}\label{inc:inp:tar:sec}

\emph{परिभाष्यता} ही संकल्पना
अंकगणिताच्या भाषेसाठी
आकारिक चिन्हार्थमीमांसा
उपलब्ध असण्यावर
अवलंबून असते.
अंकगणिताच्या भाषेतील
सूत्रे आणि वाक्ये
यांचा एक संच आपण
वर्णिला आहे.
ही वाक्ये नैसर्गिक
संख्यांबद्दल विधाने
करतात असे वाचणे
हे त्यांचे “अभिप्रेत
अर्थनिर्धारण” आहे;
असे विधान सत्य
किंवा असत्य असू
शकते.
$\Struct{N}$ ही
$\Nat$ आधारक्षेत्र
असलेली आणि
अंकगणिताच्या भाषेतील
चिन्हांचे मानक
अर्थनिर्धारण करणारी
संरचना असू द्या.
मग $\Sat{N}{\varphi}$
याचा अर्थ “$\varphi$
मानक अर्थनिर्धारणात
सत्य आहे” असा होतो.

\begin{defn}
नैसर्गिक संख्यांवरील
$R(x_1,\dots,x_k)$ हा
संबंध $\Struct{N}$
मध्ये \emph{परिभाष्य}
असतो, जर आणि फक्त जर
अंकगणिताच्या भाषेत
असे $\varphi(x_1,\dots,x_k)$
सूत्र असते की
प्रत्येक $n_1,\dots,n_k$
साठी $R(n_1,\dots,n_k)$
हे $\Sat{N}{\varphi(\num n_1,\dots,\num
  n_k)}$ असेल
तर आणि तरच खरे असते.
\end{defn}

दुसऱ्या शब्दांत,
संबंध $\Struct{N}$
मध्ये परिभाष्य
असतो, जर आणि फक्त
जर तो $\Th{TA}$
या उपपत्तीत निरूपणीय
असतो; येथे
$\Th{TA} =
\Setabs{\varphi}{\Sat{N}{\varphi}}$
हा अंकगणितातील
सत्य वाक्यांचा
संच आहे.
(हे लगेच स्पष्ट
झाले नसेल, तर
व्याख्या पुन्हा
तपासून याची
खात्री करून घ्या.)

\begin{lem}
प्रत्येक संगणनक्षम
संबंध $\Struct{N}$
मध्ये परिभाष्य
असतो.
\end{lem}

\begin{proof}
$\Th{Q}$ मध्ये
एखाद्या संबंधाचे
निरूपण करणारे
सूत्र $\Struct{N}$
मध्ये त्याच संबंधाची
परिभाषा देते, हे
सहज तपासता येते.
\end{proof}

आता उलट विधानही
खरे आहे का, असे
विचारता येते:
$\Struct{N}$ मध्ये
परिभाष्य असलेला
प्रत्येक संबंध
संगणनक्षम असतो
का? उत्तर नाही.
उदाहरणार्थ:

\begin{lem}
थांबण्याचा संबंध
$\Struct{N}$ मध्ये
परिभाष्य असतो.
\end{lem}

\begin{proof}
क्लिनीचे प्रमाण
रूप प्रमेय सांगते
की प्रत्येक आंशिक
संगणनक्षम फलन~$f$
साठी असा निर्देशांक~$e$
असतो की सर्व
$x \in \Nat$ साठी
$f(x) =
U(\umin{s}{T(e,x,s)})$;
येथे $U$ आणि $T$
आदिम पुनरावर्ती
आणि म्हणून सर्वत्र
परिभाषित आहेत.
त्यामुळे $f(x)$
परिभाषित असते
(म्हणजे संगणन
थांबते), जर आणि
फक्त जर असा~$s$
असतो की $T(e,x,s)$
खरे असते.

आता $H$ हा
थांबण्याचा संबंध
असू द्या, म्हणजे
\[H = \Setabs{\tuple{e,x}}{\lexists[s][T(e, x, s)]}.
\]
$\theta_T$ हे
$\Struct{N}$ मध्ये
$T$ ची परिभाषा
देऊ द्या. मग
\[H = \Setabs{\tuple{e,x}}{\Sat{N}{\lexists[s][\theta_T(\num e, \num x, s)]}},
\]
म्हणून $\lexists[s][\theta_T(z, x, s)]$
हे $\Struct{N}$
मध्ये~$H$ ची
परिभाषा देते.
\end{proof}

\begin{prob}
$Q(n) \defiff n \in \Setabs{\Gn{\varphi}}{\Th{Q} \Proves \varphi}$
हा संबंध अंकगणितात
परिभाष्य आहे
हे दाखवा.
\end{prob}

मग $\Th{TA}$
चे काय? ती
अंकगणितात परिभाष्य
आहे का? म्हणजे,
$\Setabs{\Gn{\varphi}}{\Sat{N}{\varphi}}$
हा संच अंकगणितात
परिभाष्य आहे
का? तार्स्कीचे
प्रमेय याचे
नकारार्थी उत्तर
देते.

\begin{thm}
\label{inc:inp:tar:thm:tarski}
अंकगणितातील सत्य
वाक्यांचा संच
अंकगणितात परिभाष्य
नाही.
\end{thm}

\begin{proof}
$\theta(x)$ त्याची
परिभाषा देते असे
समजू; म्हणजे
$\Sat{N}{\varphi}$ हे
$\Sat{N}{\theta(\gn{\varphi})}$
असेल तर आणि
तरच खरे असते.
स्थिरबिंदू उपपत्तीवरून
असे बंद सूत्र $\varphi$
असते की
$\Th{Q} \Proves \varphi \liff \lnot \theta(\gn{\varphi})$,
आणि म्हणून
$\Sat{N}{\varphi \liff \lnot \theta(\gn{\varphi})}$.
पण मग
$\Sat{N}{\varphi}$ हे
$\Sat{N}{\lnot \theta(\gn{\varphi})}$
असेल तर आणि
तरच खरे असते;
$\theta(y)$ अंकगणितातील
सत्य विधानांचा
संच परिभाषित करते
या गृहीतकाला यामुळे
विरोध होतो.
\end{proof}

तार्स्कीने हे
विश्लेषण सत्याच्या
अधिक व्यापक
तत्त्वज्ञानात्मक
संकल्पनेला लावले.
कोणतीही भाषा $L$
दिल्यास, तार्स्कीच्या
मते $L$ साठी
सत्याची पर्याप्त
संकल्पना तिच्या
प्रत्येक वाक्य
$X$ साठी पुढील
अट पूर्ण करायला
हवी:
\begin{quote}
‘$X$’ सत्य आहे
जर आणि फक्त जर
$X$.
\end{quote}
इंग्रजीसाठी
तार्स्कीने दिलेले
वारंवार उद्धृत
केले जाणारे
उदाहरण मराठीत
असे मांडता येते:
\begin{quote}
‘बर्फ पांढरा आहे’
हे सत्य आहे
जर आणि फक्त जर
बर्फ पांढरा आहे.
\end{quote}
परंतु विकर्ण
फलनाचे निरूपण
करण्याइतकी प्रबळ
असलेली कोणतीही
भाषा आणि कोणतेही
भाषिक विधेय $T(x)$
घेतल्यास, “$X$
जर आणि फक्त जर
$T(\text{‘$X$’})$
नाही” असे
वाक्य $X$
बांधता येते.
सत्यविषयक विधेयाने
एकाच वाक्याला
सत्य आणि असत्य
दोन्ही ठरवावे,
असे आपल्याला
नको आहे.
म्हणून भाषेच्या
मर्यादेबाहेर
काही ना काही
प्रकारे पाऊल
टाकल्याशिवाय
त्या भाषेतील
सर्व वाक्यांसाठी
सत्यविषयक विधेय
निर्दिष्ट करता
येत नाही,
असा तार्स्कीने
निष्कर्ष काढला.
दुसऱ्या शब्दांत,
अशा पुरेशा प्रबळ भाषेतील
सर्व वाक्यांचे सत्यविषयक विधेय
त्याच भाषेत
परिभाषित करता
येत नाही.



\part{द्वितीय-क्रम तर्कशास्त्र}
\chapter{विन्यासमीमांसा आणि चिन्हार्थमीमांसा}\label{sol:syn:chap}
\begin{editorial}
हा द्वितीय-क्रम
तर्कशास्त्रावरील
भागाचा आरंभ आहे.
\end{editorial}
\begin{editorial}
आतापर्यंत येथे
द्वितीय-क्रम
तर्कशास्त्राची
मूलभूत विन्यासमीमांसा
आणि चिन्हार्थमीमांसा
घेतली आहे.
स्वतंत्र प्रकरण
म्हणून हे फारच
लहान आहे.
द्वितीय-क्रम
चलांसाठीचे आदेशन
घेतल्याशिवाय
द्वितीय-क्रम
तर्कशास्त्राच्या
निष्पत्ती प्रणालींबद्दल
बोलता येणार नाही;
आणि तेथे काही
सूक्ष्म प्रश्नही
उद्भवतात.
\end{editorial}









\section{परिचय}\label{sol:syn:int:sec}

प्रथम-क्रम तर्कशास्त्रात,
दिलेल्या भाषेतील
अतार्किक चिन्हे,
म्हणजे तिचे
स्थिरांक,
फलने
आणि विधेये,
तार्किक चिन्हांबरोबर
जोडून आपण
प्रथम-क्रम
संरचना
विषयी विधाने
मांडतो.
हे पूर्ति
या संकल्पनेने
केले जाते.
ती संरचना~
$\Struct{M}$,
चर-मूल्यांकन~$s$
आणि सूत्र~$\varphi$
यांना जोडते:
$\Sat{M}{\varphi}[s]$
खरे असते, जर
आणि फक्त जर
त्यातील
स्थिरांक,
फलने
आणि विधेये
यांचा अर्थ
$\Struct{M}$ नुसार
आणि त्यातील मुक्त
चलांचा अर्थ~$s$
नुसार घेतल्यावर
$\varphi$ जे व्यक्त
करते ते सत्य
असते.
एकरूपता~$\eq$
चा अर्थ
$\Sat{M}{\varphi}[s]$
च्या व्याख्येतच
दिलेला असतो;
$\lforall$ आणि
$\lexists$ यांचाही
अर्थ तसाच
दिलेला असतो.
यांपैकी पहिले
चिन्ह रचनेच्या
प्रांत~$\Domain{M}$
वरील एकरूपता
संबंध म्हणूनच
नेहमी अर्थ लावले
जाते; संख्यकांचा
आवाका नेहमी
संपूर्ण प्रांत
असतो.
पण महत्त्वाचे
म्हणजे संख्यकीकरण
फक्त प्रांत
मधील घटकांवर
करता येते;
म्हणून संख्यकानंतर
फक्त वस्तू-चले
येऊ शकतात.

द्वितीय-क्रम
तर्कशास्त्रात
भाषा आणि पूर्तीची
व्याख्या या दोन्ही
विस्तारून त्यांत
मुक्त व बद्ध
फलन-चल, विधेय-चल
आणि त्यांवरील
संख्यकीकरण
समाविष्ट केले
जाते.
वस्तू-चलांचा
स्थिरांक शी
जसा संबंध असतो,
तसाच या चलांचा
फलने
आणि विधेये
शी असतो.
द्वितीय-क्रम
तर्कशास्त्रातील
पदे व सूत्रे
घडवताना त्यांची
तशीच भूमिका
असते आणि
त्यांवरील
संख्यकीकरणही
तत्सम रीतीने
हाताळले जाते.
\emph{मानक}
चिन्हार्थमीमांसेत
द्वितीय-क्रम
संख्यकांचा
आवाका योग्य
प्रकारच्या सर्व
संभाव्य वस्तूंवर
असतो (फलन-चलांसाठी
$\Domain{M}$ वरून~$\Domain{M}$
मध्ये जाणारी
$n$-स्थानी फलने,
आणि विधेय-चलांसाठी
$n$-स्थानी संबंध).
उदाहरणार्थ,
$\lforall[\Obj{v_0}][(\Obj{P^1_0}(\Obj{v_0}) \lor \lnot
  \Obj{P^1_0}(\Obj{v_0}))]$
हे प्रथम-क्रम
आणि द्वितीय-क्रम
या दोन्ही
तर्कशास्त्रांतील
सूत्र आहे;
पण दुसऱ्यामध्ये
$\lforall[\Obj{V^1_0}][\lforall[\Obj{v_0}][(\Obj{V^1_0}(\Obj{v_0}) \lor
    \lnot \Obj{V^1_0}(\Obj{v_0}))]]$
आणि
$\lexists[\Obj{V^1_0}][\lforall[\Obj{v_0}][(\Obj{V^1_0}(\Obj{v_0}) \lor
    \lnot \Obj{V^1_0}(\Obj{v_0}))]]$
यांचाही विचार
करता येतो.
त्यांत मुक्त
चले नसल्याने
ती द्वितीय-क्रम
तर्कशास्त्रातील
वाक्ये
आहेत.
येथे $\Obj{V^1_0}$
हे द्वितीय-क्रम
$1$-स्थानी
विधेय-चल आहे.
$\Obj{V^1_0}$
चे अनुमत
अर्थनिर्धारण
$1$-स्थानी
विधेय
असलेल्या
$\Obj{P^1_0}$
ला देता येणाऱ्या
अर्थनिर्धारणासारखेच
असते, म्हणजे
ते~$\Domain{M}$
चे उपसंच असतात.
त्यांवर संख्यकीकरण
केल्यास
$\lforall[\Obj{v_0}][(\Obj{V^1_0}(\Obj v_0) \lor \lnot
  \Obj{V^1_0}(v_0))]$
हे $\Obj{V^1_0}$
ला~$\Domain{M}$
चा उपसंच मूल्य
म्हणून देण्याच्या
प्रत्येक पद्धतीत
किंवा किमान
एका पद्धतीत
खरे आहे असे
म्हटले जाते.
प्रत्येक संचात
दिलेला घटक
असतो किंवा
नसतो; म्हणून
दोन्ही वाक्ये
कोणत्याही
संरचना
मध्ये सत्य असतात.











\section{पदे आणि सूत्र}\label{sol:syn:frm:sec}

प्रथम-क्रम
तर्कशास्त्राप्रमाणेच
द्वितीय-क्रम
तर्कशास्त्रातील
अभिव्यक्तींची
रचना मूलभूत
शब्दसंग्रहापासून
होते. त्यात
\emph{चले},
\emph{स्थिरांक},
\emph{विधेये}
आणि कधीकधी
\emph{फलने}
असतात.
या चिन्हांपासून,
तार्किक संयोजके,
संख्यक आणि कंस
व स्वल्पविराम
यांसारखी
विरामचिन्हे
यांच्या साहाय्याने,
\emph{पदे} आणि
\emph{सूत्रे}
तयार होतात.
फरक असा की
वस्तूंसाठीच्या
चलांव्यतिरिक्त
द्वितीय-क्रम
तर्कशास्त्रात
संबंधांसाठी
आणि फलनांसाठीही
चले असतात;
त्यांवर
संख्यकीकरण
करण्याचीही
मुभा असते.
म्हणून द्वितीय-क्रम
तर्कशास्त्राची
तार्किक चिन्हे
म्हणजे प्रथम-क्रम
तर्कशास्त्राची
चिन्हे आणि
यांखेरीज पुढील:

\begin{enumerate}
\item प्रत्येक
  स्थानसंख्या~$n$
  साठी द्वितीय-क्रम
  संबंध-चले
  चा गणनीय अनंत
  संच:
  $\Obj V_0^n$, $\Obj V_1^n$, $\Obj V_2^n$, \dots
\item द्वितीय-क्रम
  फलन-चले
  चा गणनीय अनंत
  संच:
  $\Obj u_0^n$, $\Obj u_1^n$, $\Obj u_2^n$, \dots
\end{enumerate}

प्रथम-क्रम
चल~$\Obj v_i$
यांसाठी आपण
$x$, $y$, $z$
ही अधिचले
वापरतो; त्याचप्रमाणे
$\Obj V_i^n$
यांसाठी
$X$, $Y$, $Z$
इत्यादी, आणि
$\Obj u_i^n$
यांसाठी
$u$, $v$
इत्यादी अधिचले
वापरू.

\begin{explain}
द्वितीय-क्रम
भाषेतील अतार्किक
चिन्हे प्रथम-क्रम
भाषेप्रमाणेच
निर्दिष्ट होतात:
तिची स्थिरांक,
फलने
आणि विधेये
यांची यादी देऊन.

प्रथम-क्रम
तर्कशास्त्रात
एकरूपता~$\eq$
सहसा समाविष्ट
असते.
प्रथम-क्रम
तर्कशास्त्रात
भाषा~$\Lang{L}$
मधील अतार्किक
चिन्हे काही
लक्षणीय गोष्ट
व्यक्त करण्यासाठी
अत्यावश्यक असतात.
अतार्किक चिन्हे
न वापरणारी
वाक्ये
अर्थात असतात;
पण फक्त~$\eq$
असल्यास काही
लक्षणीय गोष्ट
सांगणे कठीण
असते.
द्वितीय-क्रम
तर्कशास्त्रात
आपल्याकडे संबंध-
आणि फलन-चलांचा
अमर्याद पुरवठा
असल्यामुळे,
अतार्किक चिन्हांचा
विशेष पुरवठा
नसतानाही
प्रथम-क्रम
भाषेत जे काही
सांगता येते
ते सांगता येते.
\end{explain}

\begin{defn}[द्वितीय-क्रम पदे]
भाषा~$\Lang L$
मधील \emph{द्वितीय-क्रम
पदांचा} संच
$\TrmSOL[L]$
हा \ref{fol:syn:frm:defn:terms}
मध्ये पुढील
खंड वाढवून
परिभाषित होतो:
\begin{enumerate}
\item $u$ हे
  $n$-स्थानी
  फलन-चल
  आणि $t_1$,
  \dots, $t_n$
  ही पदे असतील,
  तर $\Atom{u}{t_1, \ldots, t_n}$
  हे पद असते.
\end{enumerate}
\end{defn}

\begin{explain}
म्हणून द्वितीय-क्रम
पद प्रथम-क्रम
पदासारखेच दिसते;
फक्त प्रथम-क्रम
पदात जिथे
फलन~$\Obj{f^n_i}$
असते, तिथे
द्वितीय-क्रम
पदात त्याच्या
जागी फलन-चल~$\Obj{u^n_i}$
असू शकते.
\end{explain}

\begin{defn}[द्वितीय-क्रम सूत्र]
भाषा~$\Lang L$
मधील \emph{द्वितीय-क्रम
सूत्रे} चा
संच~$\FrmSOL[L]$
हा \ref{fol:syn:frm:defn:formulas}
मध्ये पुढील
खंड वाढवून
परिभाषित होतो:
\begin{enumerate}
\item $X$ हे
  $n$-स्थानी
  विधेय-चल
  आणि $t_1$,
  \dots, $t_n$
  ही~$\Lang L$
मधील द्वितीय-क्रम
पदे असतील, तर
  $\Atom{X}{t_1,\ldots, t_n}$
  हे अणुरूप
  सूत्र
  असते.

\item $\varphi$ हे सूत्र आणि $u$ हे फलन-चल असेल,
  तर $\lforall[u][\varphi]$ हे सूत्र असते.

\item $\varphi$ हे सूत्र आणि $X$ हे विधेय-चल असेल,
  तर $\lforall[X][\varphi]$ हे सूत्र असते.

\item $\varphi$ हे सूत्र आणि $u$ हे फलन-चल असेल,
  तर $\lexists[u][\varphi]$ हे सूत्र असते.

\item $\varphi$ हे सूत्र आणि $X$ हे विधेय-चल असेल,
  तर $\lexists[X][\varphi]$ हे सूत्र असते.
\end{enumerate}
\end{defn}












\section{पूर्ति}\label{sol:syn:sat:sec}

\begin{explain}
द्वितीय-क्रम
सूत्रे साठी
पूर्ति-संबंध
$\Sat{M}{\varphi}[s]$
परिभाषित करताना
द्वितीय-क्रम
चले
समाविष्ट करण्यासाठी
व्याख्या वाढवाव्या
लागतात.
संरचना
ही संकल्पना
द्वितीय-क्रम
आणि प्रथम-क्रम
तर्कशास्त्रात
सारखीच असते.
फरक फक्त
चर-मूल्यांकन~$s$
बाबत असतो:
आता त्याने
प्रथम-क्रम
चले
बरोबरच
द्वितीय-क्रम
चले
यांनाही मूल्ये
द्यावी लागतात.
\end{explain}

\begin{defn}[चर-मूल्यांकन]
संरचना~$\Struct{M}$
साठीचे
\emph{चर-मूल्यांकन}~$s$
हे असे फलन
आहे जे पुढीलपैकी
प्रत्येकाला
दिल्याप्रमाणे
प्रतिचित्रित करते:
\begin{enumerate}
\item वस्तू-चल~$\Obj{v_i}$
  याला~$\Domain M$
  मधील घटकाकडे,
  म्हणजे
  $s(\Obj{v_i}) \in \Domain{M}$;
\item $n$-स्थानी
  संबंध-चल~$\Obj{V_i^n}$
  याला~$\Domain{M}$
  वरील $n$-स्थानी
  संबंधाकडे,
  म्हणजे
  $s(\Obj{V_i^n}) \subseteq \Domain{M}^n$;
\item $n$-स्थानी
  फलन-चल~$\Obj{u_i^n}$
  याला
  $\Domain{M}$
  वरून $\Domain{M}$
  मध्ये जाणाऱ्या
  $n$-स्थानी
  फलनाकडे,
  म्हणजे
  $s(\Obj{u_i^n})\colon \Domain{M}^n \to \Domain{M}$;
\end{enumerate}
\end{defn}

\begin{explain}
संरचना
प्रत्येक स्थिरांक
आणि फलन
ला मूल्य
देते, तर द्वितीय-क्रम
चर-मूल्यांकन
प्रत्येक वस्तू-चल
आणि फलन-चल
यांना अनुक्रमे
वस्तू आणि फलने
देते. या दोन्हींच्या
साहाय्याने
प्रत्येक पदाला
मूल्य देता येते.
\end{explain}

\begin{defn}[पदाचे मूल्य]
$t$ हे भाषा~$\Lang L$
मधील पद,
$\Struct M$ ही~$\Lang L$
साठी संरचना,
आणि $s$ हे~$\Struct M$
साठी चल
मूल्यांकन असेल,
तर \emph{मूल्य}~
$\Value{t}{M}[s]$
हे प्रथम-क्रम
पदांप्रमाणेच,
पुढील अतिरिक्त
खंडासह परिभाषित
होते:
\begin{quote}
\indcase{t}{\Atom{u}{t_1, \ldots, t_n}}{
\[\Value{\indfrm}{M}[s] = s(u)(\Value{t_1}{M}[s], \ldots,
\Value{t_n}{M}[s]).
\]}
\end{quote}
\end{defn}

\begin{defn}[$x$-पर्याय]
$s$ हे संरचना~$\Struct M$
साठी चल
मूल्यांकन असेल,
तर~$\Struct M$
साठीचे असे
कोणतेही चल
मूल्यांकन $s'$
जे $s$ पेक्षा
जास्तीत जास्त
$x$ ला दिलेल्या
मूल्यातच भिन्न
असेल, त्याला
$s$ चा
\emph{$x$-पर्याय}
म्हणतात.
$s'$ हा $s$
चा $x$-पर्याय
असेल, तर
$\varAssign{s'}{s}{x}$
असे लिहितो.
(द्वितीय-क्रम
चल $X$ किंवा
$u$ यांसाठीही
तसेच.)
\end{defn}

\begin{defn}

  $s$ हे संरचना~$\Struct M$
  साठी चल
  मूल्यांकन
  आणि $m \in \Domain{M}$
  असेल, तर
  मूल्यांकन~$\Subst{s}{m}{x}$
  हे पुढीलप्रमाणे
  परिभाषित केलेले
  चल
  मूल्यांकन आहे:
  \[\Subst{s}{m}{x}(y) = \begin{cases}
    m & \text{जर } y \ident x\\
    s(y) & \text{अन्यथा},
  \end{cases}\]
  $X$ हे $n$-स्थानी
  संबंध-चल
  आणि $M \subseteq
  \Domain{M}^n$
  असेल, तर
  $\Subst{s}{M}{X}$
  हे पुढीलप्रमाणे
  परिभाषित केलेले
  चल
  मूल्यांकन आहे:
  \[\Subst{s}{M}{X}(y) = \begin{cases}
    M & \text{जर } y \ident X\\
    s(y) & \text{अन्यथा}.
  \end{cases}\]
  $u$ हे $n$-स्थानी
  फलन-चल
  आणि $f\colon \Domain{M}^n
  \to \Domain{M}$
  असेल, तर
  $\Subst{s}{f}{u}$
  हे पुढीलप्रमाणे
  परिभाषित केलेले
  चल
  मूल्यांकन आहे:
  \[\Subst{s}{f}{u}(y) = \begin{cases}
    f & \text{जर } y \ident u\\
    s(y) & \text{अन्यथा}.
  \end{cases}\]
  प्रत्येक बाबतीत
  $y$ हे कोणतेही
  प्रथम-क्रम किंवा
  द्वितीय-क्रम
  चल
  असू शकते.
\end{defn}

\begin{defn}[पूर्ति]
द्वितीय-क्रम
सूत्रे~$\varphi$
साठी पूर्तीची
व्याख्या
\ref{fol:syn:sat:defn:satisfaction}
प्रमाणेच आहे;
त्यात पुढील
खंड वाढवले
आहेत:
\begin{enumerate}
\item \indcase{\varphi}{\Atom{X^n}{t_1, \dots, t_n}}{$\Sat{M}{\indfrm}[s]$
  हे $\langle \Value{t_1}{M}[s], \dots, \Value{t_n}{M}[s] \rangle \in
  s(X^n)$ असेल तर आणि तरच खरे.}
\item
  \indcase{\varphi}{\lforall[X][\psi]}{$\Sat{M}{\indfrm}[s]$ हे प्रत्येक $M
  \subseteq \Domain{M}^n$ साठी $\Sat{M}{\psi}[\Subst{s}{M}{X}]$ असेल
  तर आणि तरच खरे.}

\item
  \indcase{\varphi}{\lexists[X][\psi]}{$\Sat{M}{\indfrm}[s]$ हे किमान
  एका $M \subseteq \Domain{M}^n$ साठी
  $\Sat{M}{\psi}[\Subst{s}{M}{X}]$ असेल तर आणि तरच खरे.}

\item
  \indcase{\varphi}{\lforall[u][\psi]}{$\Sat{M}{\indfrm}[s]$ हे प्रत्येक
  $f\colon \Domain{M}^n \to \Domain{M}$ साठी
  $\Sat{M}{\psi}[\Subst{s}{f}{u}]$ असेल तर आणि तरच खरे.}

\item
  \indcase{\varphi}{\lexists[u][\psi]}{$\Sat{M}{\indfrm}[s]$ हे किमान
  एका $f\colon \Domain{M}^n \to \Domain{M}$ साठी
  $\Sat{M}{\psi}[\Subst{s}{f}{u}]$ असेल तर आणि तरच खरे.}
\end{enumerate}
\end{defn}

\begin{ex}
  सूत्र
  $\lforall[z][(\Atom{X}{z} \liff \lnot
    \Atom{Y}{z})]$
  विचारात घ्या.
  त्यात द्वितीय-क्रम
  संख्यक नाहीत,
  पण द्वितीय-क्रम
  चले
  $X$ आणि~$Y$
  आहेत (येथे ती
  एक-स्थानी मानली
  आहेत).
  याच्याशी संबंधित
  प्रथम-क्रम
  वाक्य
  $\lforall[z][(\Atom{P}{z} \liff \lnot \Atom{R}{z})]$
  असे म्हणते की
  आधारक्षेत्रातील प्रत्येक वस्तू
  $P$ आणि $R$ यांच्या
  अर्थनिर्धारणांपैकी नेमक्या एकात येते:
  ती पहिल्यात येते, तर आणि तरच
  ती दुसऱ्यात येत नाही.
  संरचना
  मध्ये
  विधेय~$P$
  चे अर्थनिर्धारण
  $\Assign{P}{M}$
  याने दिलेले
  असते. पण
  $X$ आणि~$Y$
  सारख्या द्वितीय-क्रम
  चले
  चे अर्थनिर्धारण
  संरचना
  स्वतः देत नाही;
  ते चल
  मूल्यांकन देते.
  हे द्वितीय-क्रम
  सूत्र
  वाक्य
  नाही (त्यात
  $X$ आणि~$Y$
  ही मुक्त
  चले
  आहेत);
  म्हणून त्याची
  पूर्ति फक्त
  संरचना~$\Struct{M}$
  आणि त्यासोबत
  चल
  मूल्यांकन~$s$
  यांच्या सापेक्ष
  ठरते.

  $\Sat{M}{\lforall[z][(\Atom{X}{z} \liff \lnot \Atom{Y}{z})]}[s]$
  हे $s(X)$ मधील
  घटक आणि $s(Y)$ मधील
  घटक एकमेकांपासून वेगळे असतील
  आणि दोन्ही संच मिळून संपूर्ण आधारक्षेत्र झाकतील,
  तर आणि तरच खरे असते; म्हणजे
  $s(Y) = \Domain{M} \setminus s(X)$
  असेल तर आणि
  तरच खरे असते.
  उदाहरणार्थ,
  $\Domain{M} = \{1, 2, 3\}$
  घ्या.
  येथे विधेये,
  फलने
  किंवा स्थिरांक
  येत नसल्यामुळे
  $\Struct{M}$
  चे आधारक्षेत्र
  एवढेच संबंधित
  आहे.
  आता
  $s_1(X) = \{1, 2\}$
  आणि
  $s_1(Y) = \{3\}$
  असताना
  $\Sat{M}{\lforall[z][(\Atom{X}{z}
      \liff \lnot \Atom{Y}{z})]}[s_1]$
  खरे आहे.

  याउलट,
  $s_2(X) = \{1, 2\}$
  आणि
  $s_2(Y) = \{2, 3\}$
  असल्यास
  $\Sat/{M}{\lforall[z][(\Atom{X}{z} \liff \lnot
      \Atom{Y}{z})]}[s_2]$.
  कारण
  $\Sat{M}{\Atom{X}{z}}[\Subst{s_2}{2}{z}]$
  (कारण
  $2 \in \Subst{s_2}{2}{z}(X)$),
  पण
  $\Sat/{M}{\lnot \Atom{Y}{z}}[\Subst{s_2}{2}{z}]$
  (कारण
  $2 \in
  \Subst{s_2}{2}{z}(Y)$
  सुद्धा).
\end{ex}

\begin{ex}
  $\Sat{M}{\lexists[Y][(\lexists[y][\Atom{Y}{y}] \land
  \lforall[z][(\Atom{X}{z} \liff \lnot \Atom{Y}{z})])]}[s]$
  हे तेव्हाच खरे
  असते जेव्हा
  असा $N \subseteq \Domain{M}$
  असतो की
  $\Sat{M}{(\lexists[y][\Atom{Y}{y}] \land \lforall[z][(\Atom{X}{z}
  \liff \lnot \Atom{Y}{z})])}[\Subst{s}{N}{Y}]$.
  आणि असे
  $N \neq \emptyset$
  असताना
  ($\Sat{M}{\lexists[y][\Atom{Y}{y}]}[\Subst{s}{N}{Y}]$
  मिळण्यासाठी)
  आणि आधीच्या
  उदाहरणाप्रमाणे
  $N = \Domain{M} \setminus s(X)$
  असताना घडते.
  दुसऱ्या शब्दांत,
  $\Sat{M}{\lexists[Y][(\lexists[y][\Atom{Y}{y}] \land
  \lforall[z][(\Atom{X}{z} \liff \lnot \Atom{Y}{z})])]}[s]$
  हे
  $\Domain{M} \setminus s(X)$
  रिकामे नसेल,
  म्हणजे
  $s(X) \neq
  \Domain{M}$
  असेल, तर आणि
  तरच खरे असते.
  म्हणून
  $\Domain{M}
  = \{1, 2, 3\}$
  आणि
  $s(X) = \{1, 2\}$
  असल्यास
  सूत्र
  ची पूर्ति होते;
  पण
  $s(X) = \{1, 2, 3\}
  = \Domain{M}$
  असल्यास होत नाही.

  $s(X) = \Domain{M}$
  असताना या
  सूत्र
  ची पूर्ति
  होत नसल्यामुळे,
  पुढील
  वाक्य
  \[  \lforall[X][\lexists[Y][(\lexists[y][\Atom{Y}{y}] \land
      \lforall[z][(\Atom{X}{z} \liff \lnot \Atom{Y}{z})])]]
  \]
  कधीच पूर्ण
  होत नाही:
  कोणतीही
  संरचना~$\Struct{M}$
  घेतली तरी
  $s(X) = \Domain{M}$
  हे मूल्यांकन
  त्या
  वाक्य
  ला असत्य
  ठरवेल.
  याउलट,
  पुढील वाक्य
  \[  \lexists[X][\lexists[Y][(\lexists[y][\Atom{Y}{y}] \land
      \lforall[z][(\Atom{X}{z} \liff \lnot \Atom{Y}{z})])]]
  \]
  कोणत्याही
  मूल्यांकन~$s$
  च्या सापेक्ष
  पूर्ण होते,
  कारण
  $M \subseteq \Domain{M}$
  पण
  $M \neq \Domain{M}$
  असा संच नेहमी
  निवडता येतो
  (उदा.,
  $M = \emptyset$).
\end{ex}

\begin{ex}
  द्वितीय-क्रम
  वाक्य~
  $\lforall[X][\lforall[y][X(y)]]$
  असे म्हणते की
  प्रत्येक
  $1$-स्थानी
  संबंध, म्हणजे
  प्रत्येक
  गुणधर्म, प्रत्येक
  वस्तूला लागू
  होतो.
  हे स्पष्टपणे
  कधीच खरे नसते:
  प्रत्येक~$\Struct{M}$
  मध्ये
  $s(X) = \emptyset$
  आणि
  $s(y) = a \in
  \Domain{M}$
  असे चर-मूल्यांकन~$s$
  घेतल्यास
  $\Sat/{M}{X(y)}[s]$
  मिळते.
  याचा अर्थ
  $\varphi \lif
  \lforall[X][\lforall[y][X(y)]]$
  हे द्वितीय-क्रम
  तर्कशास्त्रात
  $\lnot \varphi$
  शी सममूल्य
  आहे; म्हणजे
  $\Sat{M}{\varphi \lif
  \lforall[X][\lforall[y][X(y)]]}$
  हे
  $\Sat{M}{\lnot \varphi}$
  असेल तर आणि
  तरच खरे.
  दुसऱ्या शब्दांत,
  द्वितीय-क्रम
  तर्कशास्त्रात
  $\lforall$
  आणि~$\lif$
  वापरून
  $\lnot$
  परिभाषित
  करता येते.
\end{ex}

\begin{prob}
  द्वितीय-क्रम
  तर्कशास्त्रात
  $\lforall$
  आणि~$\lif$
  वापरून इतर
  संयोजके
  परिभाषित
  करता येतात
  हे दाखवा:
  \begin{enumerate}
    \item येथे X हे संयोजित करावयाच्या दोन्ही सूत्रांमध्ये मुक्त नसलेले
    एकस्थानी विधेय चल निवडा. द्वितीय-क्रम तर्कशास्त्रात $\varphi \land \psi$ हे
    $\lforall[X][(\varphi \lif (\psi \lif \lforall[x][X(x)])\lif
    \lforall[x][X(x)])]$ शी सममूल्य आहे हे सिद्ध करा.
    \item फक्त $\lforall$ आणि $\lif$ वापरून
    $\varphi \lor \psi$ शी सममूल्य द्वितीय-क्रम सूत्र शोधा.
  \end{enumerate}
\end{prob}











\section{चिन्हार्थविषयक संकल्पना}\label{sol:syn:sem:sec}

\begin{explain}
\emph{वैधता},
\emph{तार्किक निष्पन्नता}
आणि
\emph{पूर्ततायोग्यता}
या मध्यवर्ती
तार्किक संकल्पना
द्वितीय-क्रम
तर्कशास्त्रासाठी
प्रथम-क्रम
तर्कशास्त्राप्रमाणेच
परिभाषित होतात;
फरक इतकाच की
आधारभूत पूर्ति-संबंध
आता द्वितीय-क्रम
सूत्रे
साठीचा असतो.
अर्थात, द्वितीय-क्रम
वाक्य
हे असे
सूत्र
असते ज्यातील
विधेय-चल आणि
फलन-चल यांसह
सर्व चले
बद्ध असतात.
\end{explain}

\begin{defn}[वैधता]
वाक्य $\varphi$
\emph{वैध},
$\Entails \varphi$,
असते जर आणि
फक्त जर प्रत्येक
संरचना~$\Struct M$
साठी $\Sat{M}{\varphi}$.
\end{defn}

\begin{defn}[तार्किक निष्पन्नता]
वाक्यांच्या
संच~$\Gamma$
मधून वाक्य~$\varphi$
\emph{तार्किकरीत्या
निष्पन्न होते},
$\Gamma
\Entails \varphi$,
जर आणि फक्त जर
$\Sat{M}{\Gamma}$
असलेल्या प्रत्येक
संरचना~$\Struct M$
साठी
$\Sat{M}{\varphi}$.
\end{defn}

\begin{defn}[पूर्ततायोग्यता]
एखाद्या
संरचना~$\Struct M$
साठी
$\Sat{M}{\Gamma}$
असेल, तर वाक्यांचा
संच~$\Gamma$
\emph{पूर्ततायोग्य}
असतो.
$\Gamma$
पूर्ततायोग्य नसेल,
तर त्याला
\emph{अपूर्ततायोग्य}
म्हणतात.
\end{defn}











\section{अभिव्यक्तीचे सामर्थ्य}\label{sol:syn:exp:sec}

\begin{explain}
द्वितीय-क्रम
चलांवरील
संख्यकीकरणामुळे
भाषेचे अभिव्यक्तीचे
सामर्थ्य प्रथम-क्रम
तर्कशास्त्राच्या
तुलनेत प्रचंड
वाढते.
द्वितीय-क्रम
अस्तित्ववाचक
संख्यकीकरणाने
विशिष्ट गुणधर्म
असलेली फलने
किंवा संबंध
अस्तित्वात आहेत,
असे म्हणता येते.
प्रथम-क्रम
तर्कशास्त्रात
हे करण्याचा
एकमेव मार्ग
म्हणजे त्यासाठी
अतार्किक चिन्ह
(म्हणजे फलन
किंवा विधेय)
निर्दिष्ट करणे.
द्वितीय-क्रम
वैश्विक संख्यकीकरणाने
आधारक्षेत्राचे
सर्व उपसंच,
त्यावरील सर्व
संबंध किंवा
प्रांत वरून
प्रांत मध्ये
जाणारी सर्व
फलने यांना
एखादा गुणधर्म
आहे, असे
म्हणता येते.
प्रथम-क्रम
तर्कशास्त्रात
भाषेतील अतार्किक
चिन्हांपैकी
एखाद्याला
नेमून दिलेल्या
उपसंचांना,
संबंधांना किंवा
फलनांनाच
गुणधर्म आहे,
एवढेच म्हणता
येते.
एखादा गुणधर्म
असलेले उपसंच,
संबंध किंवा
फलने अस्तित्वात
आहेत किंवा
त्या सर्वांना
तो गुणधर्म
आहे असे म्हणताना,
त्या गुणधर्माचे
वर्णन करतानाही
द्वितीय-क्रम
संख्यकीकरण
वापरता येते.
यामुळे प्रथम-क्रम
तर्कशास्त्रात
परिभाषित न
होणारे संबंध
परिभाषित करता
येतात आणि
तेथे व्यक्त न
होणारे आधारक्षेत्राचे
गुणधर्म व्यक्त
करता येतात.
\end{explain}

\begin{defn}
$\Struct{M}$ ही
भाषा~$\Lang{L}$
साठी संरचना
असेल, तर
$R \subseteq \Domain{M}^2$
हा संबंध~$\Lang{L}$
मध्ये \emph{परिभाष्य}
असतो, जर फक्त
$x$ आणि~$y$
हीच मुक्त चले
असलेले असे
सूत्र~$\varphi_R(x, y)$
असते की
$s(x) = a$
आणि $s(y) = b$
असताना
$R(a, b)$
खरे असते
(म्हणजे
$\tuple{a,b} \in R$),
जर आणि फक्त
जर
$\Sat{M}{\varphi_R(x, y)}[s]$.
\end{defn}

\begin{ex}
प्रथम-क्रम
तर्कशास्त्रात
एकरूपता संबंध~
$\Id{\Domain{M}}$
(म्हणजे
$\Setabs{\tuple{a,a}}{a \in
  \Domain{M}}$)
सूत्र~$\eq[x][y]$
याने परिभाषित
करता येतो.
द्वितीय-क्रम
तर्कशास्त्रात
हा संबंध
\emph{$\eq$ शिवाय}
परिभाषित
करता येतो.
कारण $a$
आणि $b$
हे~$\Domain{M}$
चे तेच
घटक
असतील, तर
ते~$\Domain{M}$
च्या त्याच
उपसंचांचे
घटक
असतात (कारण
संच त्यांतील
घटक
नी ठरतात).
उलट $a$
आणि $b$
भिन्न असतील,
तर ते त्याच
उपसंचांचे
घटक
नसतात:
उदा.,
$a \in \{a\}$
पण $a \neq b$
असताना
$b
\notin \{a\}$.
म्हणून
“~$\Domain{M}$
च्या त्याच
उपसंचांचे
घटक
असणे” हा
संबंध $a$
आणि $b$
यांच्याबाबत
$a = b$
असेल तर आणि
तरच लागू
होतो.
हा संबंध
द्वितीय-क्रम
तर्कशास्त्रात
व्यक्त करता
येतो, कारण
~$\Domain{M}$
च्या सर्व
उपसंचांवर
संख्यकीकरण
करता येते.
म्हणून पुढील
सूत्र
$\Id{\Domain{M}}$
परिभाषित करते:
\[\lforall[X][(X(x) \liff X(y))]
\]
\end{ex}

\begin{prob}
$\lforall[X][(X(x) \lif X(y))]$
(लक्षात घ्या:
$\liff$ नव्हे,
$\lif$!)
हे सूत्र
$\Id{\Domain{M}}$
परिभाषित करते,
हे दाखवा.
\end{prob}

\begin{ex}
$R$ हे
द्विस्थानी
विधेय
असेल, तर
$\Assign{R}{M}$
हा~$\Domain{M}$
वरील द्विस्थानी
संबंध असतो.
थोडे गोंधळात
टाकणारे असले,
तरी~$R$ हेच
चिन्ह त्या
विधेय
साठी~$R$ आणि
संबंध~$\Assign{R}{M}$
साठीही वापरू.
$R$ चे
\emph{संक्रमक
संवरण}~$R^+$
हा $a$
आणि $b$
यांच्यात तेव्हाच
लागू होणारा
संबंध आहे
जेव्हा काही
$c_1$, \dots, $c_k$
साठी
$R(a,c_1)$,
$R(c_1, c_2)$,
\dots, $R(c_k,b)$
लागू होतात.
यात $k = 0$
हा प्रसंगही
येतो, म्हणजे
$R(a,b)$
लागू असेल,
तर $R^+(a,b)$
ही लागू असतो.
याचा अर्थ
$R \subseteq R^+$.
खरे तर
$R^+$ हा~$R$
समाविष्ट करणारा
सर्वांत लहान
संक्रमक संबंध
आहे.
द्वितीय-क्रम
तर्कशास्त्रात
$X$ हा~$R$
समाविष्ट करणारा
संक्रमक संबंध
आहे असे
म्हणता येते:
\begin{multline*}
  \psi_R(X) \ident \lforall[x][\lforall[y][(R(x,y) \lif X(x, y))]] \land {}\\
\lforall[x][\lforall[y][\lforall[z][((X(x,y) \land X(y,z)) \lif X(x,
      z))]]].
\end{multline*}
पहिला संयुक्तांश
$R \subseteq X$
सांगतो आणि
दुसरा $X$
संक्रमक आहे
असे सांगतो.

$X$ हा
असा सर्वांत
लहान संबंध
आहे असे
म्हणणे म्हणजे
$R$ समाविष्ट
करणाऱ्या
आणि संक्रमक
असलेल्या
प्रत्येक संबंधात
तो समाविष्ट
आहे असे
म्हणणे.
म्हणून
पुढील
सूत्र
याने~$R$
चे संक्रमक
संवरण
परिभाषित
करता येते:
\[R^+(X) \ident \psi_R(X) \land \lforall[Y][(\psi_R(Y) \lif
  \lforall[x][\lforall[y][(X(x, y) \lif Y(x,y))]])].
\]
$\Sat{M}{R^+(X)}[s]$
हे
$s(X) = R^+$
असेल तर आणि
तरच खरे.
$R$ चे
संक्रमक संवरण
प्रथम-क्रम
तर्कशास्त्रात
व्यक्त करता
येत नाही.
\end{ex}












\section{अनंत आणि गणनीय
  प्रांत यांचे वर्णन}\label{sol:syn:siz:sec}


येथे निवडीचे स्वयंसिद्धक मानणारा नेहमीचा
संचसैद्धान्तिक संदर्भ घेतला आहे. त्यात डेडेकिंड अनंतता
आणि सर्वसाधारण अनंतता सममूल्य असतात.
संच~$M$ हा
(डेडेकिंडच्या अर्थाने)
अनंत असतो,
जर आणि फक्त जर
$f\colon M
\to M$ असे
एकास-एक
फलन असते जे
आच्छादक
नसते; म्हणजे
$\ran{f} \neq M$.
प्रथम-क्रम
तर्कशास्त्रात
एक-स्थानी
फलन~$f$
घेऊन
संरचना~$\Struct{M}$
मध्ये त्याला
नेमलेले फलन
$\Assign{f}{M}$
हे एकास-एक
आहे आणि
$\ran{f} \neq
\Domain{M}$
आहे असे
म्हणता येते:
\[\lforall[x][\lforall[y][(\eq[f(x)][f(y)] \to \eq[x][y])]] \land
\lexists[y][\lforall[x][\eq/[y][f(x)]]].
\]
$\Struct{M}$
ही या
वाक्य
ची पूर्ति करत
असेल, तर
$\Assign{f}{M}:
\Domain{M} \to \Domain{M}$
हे एकास-एक
असते, म्हणून
$\Domain{M}$
अनंत असलेच
पाहिजे.
$\Domain{M}$
अनंत असेल
आणि म्हणून
असे फलन
अस्तित्वात असेल,
तर ते फलन
$\Assign{f}{M}$
म्हणून घेता
येते आणि
$\Struct{M}$
ही त्या
वाक्य
ची पूर्ति करेल.
मात्र यासाठी
आपल्या भाषेत
या हेतूने
वापरण्यात येणारे
अतार्किक चिन्ह~$f$
असणे आवश्यक
आहे.
द्वितीय-क्रम
तर्कशास्त्रात
असे फलन
\emph{अस्तित्वात आहे}
असे थेट
म्हणता येते.
यासाठी
आता~$f$
लागत नाही,
आणि शुद्ध
द्वितीय-क्रम
तर्कशास्त्रातील
पुढील
वाक्य
मिळते:
\[\fn{Inf} \ident \lexists[u][(\lforall[x][\lforall[y][(\eq[u(x)][u(y)]
      \lif \eq[x][y])]] \land \lexists[y][\lforall[x][\eq/[y][u(x)]]])].
\]
$\Sat{M}{\fn{Inf}}$
हे
$\Domain{M}$
अनंत असेल
तर आणि तरच
खरे असते.
मग
$\fn{Fin} \ident \lnot \fn{Inf}$
असे परिभाषित
करता येते;
$\Sat{M}{\fn{Fin}}$
हे
$\Domain{M}$
सांत असेल
तर आणि तरच
खरे असते.
शुद्ध प्रथम-क्रम
तर्कशास्त्रातील
कोणतेही एक
वाक्य
प्रांत
अनंत आहे
हे व्यक्त
करू शकत
नाही, मात्र
अशा वाक्यांचा
अनंत संच
ते करू शकतो.
शुद्ध प्रथम-क्रम
तर्कशास्त्रातील
वाक्ये
चा असा
कोणताही संच
नाही की
त्याची
संरचना
मध्ये पूर्ति
होते जर आणि
फक्त जर तिचे
आधारक्षेत्र
सांत आहे.

\begin{prop}
$\Sat{M}{\fn{Inf}}$
हे
$\Domain{M}$
अनंत असेल
तर आणि तरच
खरे असते.
\end{prop}

\begin{proof}
$\Sat{M}{\fn{Inf}}$
असेल तर आणि
तरच काही~$s$
साठी
  $\Sat{M}{\lforall[x][\lforall[y][(\eq[u(x)][u(y)] \lif \eq[x][y])]]
    \land \lexists[y][\lforall[x][\eq/[y][u(x)]]]}[s]$.
  असे असेल,
  तर $s(u)$
  हे एकास-एक
  फलन आहे
  आणि काही
  $y
  \in \Domain{M}$
  हे~$s(u)$
  च्या व्याप्तीत
  नाही.
  उलट
  $f\colon \Domain{M} \to \Domain{M}$
  असे
  एकास-एक
  फलन असेल
  की
  $\ran{f}
  \neq \Domain{M}$,
  तर
  $s(u) = f$
  असे चर-मूल्यांकन
  मिळते.
\end{proof}

संच $M$
गणनीय
असतो, जर
त्यातील
घटक
ची पुनरावृत्ती
नसलेली
(पण शक्यतो
सांत असलेली)
पुढील
प्रगणना असते:
\[m_0, m_1, m_2, \dots
\]
अशी प्रगणना
असते जर
आणि फक्त
जर
घटक
$z \in M$
आणि
$f\colon M \to M$
असे फलन
असते की
$z$, $f(z)$,
$f(f(z))$, \dots,
हे~$M$
चे सर्व
घटक
असतात.
कारण प्रगणना
असेल, तर
$z = m_0$
आणि
$f(m_k) = m_{k+1}$
(किंवा
$m_k$ हा
प्रगणनेतील
शेवटचा
घटक
असेल, तर
$f(m_k) = m_k$)
हे अपेक्षित
घटक
आणि फलन
ठरतात.
उलट असे
$z$ आणि
$f$ असतील,
तर
$z$, $f(z)$,
$f(f(z))$, \dots,
या यादीतील पुनरावृत्ती
काढल्यावर~$M$ ची प्रगणना
मिळते आणि
$M$ हा
गणनीय
असतो.
$z$ आणि~$f$
यांचे अस्तित्व
द्वितीय-क्रम
तर्कशास्त्रात
व्यक्त करून
असे
वाक्य
तयार करता
येते की
ते
संरचना
मध्ये खरे
असते जर
आणि फक्त
जर ती
संरचना
गणनीय
असते:
\[\fn{Count} \ident
\lexists[z][\lexists[u][\lforall[X][((X(z) \land
      \lforall[x][(X(x) \lif X(u(x)))]) \lif \lforall[x][X(x)])]]]
\]

\begin{prop}
$\Sat{M}{\fn{Count}}$
हे
$\Domain{M}$
गणनीय
असेल तर
आणि तरच
खरे असते.
\end{prop}

\begin{proof}
$\Domain{M}$
गणनीय
आहे असे
समजू आणि
$m_0$, $m_1$,
\dots, ही
प्रगणना घेऊ.
पुनरावृत्ती
काढून टाकल्यास
कोणतेही
$m_k$ दोनदा
येणार नाही
याची खात्री
करता येते.
$f(m_k) = m_{k+1}$
असे परिभाषित
करू (सांत
प्रगणनेच्या
शेवटच्या घटकावर
फलन तोच घटक
परत देते)
आणि
$s(z) = m_0$,
$s(u) = f$
घेऊ.
आता पुढील
गोष्ट दाखवू:
\[\Sat{M}{\lforall[X][((X(z) \land \lforall[x][(X(x) \lif X(u(x)))])
    \lif \lforall[x][X(x)])]}[s]
\]
$M \subseteq \Domain{M}$
हा कोणताही
संच असू
द्या.
तसेच
$\Sat{M}{(X(z) \land \lforall[x][(X(x) \lif
X(u(x)))])}[\Subst{s}{M}{X}]$
असे धरू.
मग
$\Subst{s}{M}{X}(z) \in M$
आणि
$x \in M$
असेल तेव्हा
$(\Subst{s}{M}{X}(u))(x) \in M$
ही.
दुसऱ्या
शब्दांत,
$\varAssign{\Subst{s}{M}{X}}{s}{X}$
असल्यामुळे
$m_0 \in M$
आणि
$x \in M$
असेल तर
$f(x) \in M$;
म्हणून
$m_0 \in M$,
$m_1 = f(m_0) \in M$,
$m_2 = f(f(m_0)) \in M$
इत्यादी.
त्यामुळे
$M = \Domain{M}$
आणि म्हणून
$\Sat{M}{\lforall[x][X(x)]}[\Subst{s}{M}{X}]$.
$M \subseteq
\Domain{M}$
कोणताही
असल्यामुळे
सिद्धता पूर्ण:
$\Sat{M}{\fn{Count}}$.

आता
$\Sat{M}{\fn{Count}}$,
म्हणजे
\[\Sat{M}{\lforall[X][((X(z) \land \lforall[x][(X(x) \lif X(u(x)))])
    \lif \lforall[x][X(x)])]}[s]
\]
काही~$s$
साठी आहे
असे समजू.
$m = s(z)$
आणि
$f = s(u)$
घेऊ आणि
$M = \{m,
f(m), f(f(m)), \dots\}$
हा संच
विचारात
घेऊ.
असा
परिभाषित
केलेला
$M$
स्पष्टपणे
गणनीय
आहे.
मग
\[\Sat{M}{(X(z) \land \lforall[x][(X(x) \lif X(u(x)))])
    \lif \lforall[x][X(x)]}[\Subst{s}{M}{X}]
\]
हे गृहीतकावरून
मिळते.
तसेच
$\Sat{M}{X(z)}[\Subst{s}{M}{X}]$,
कारण
$M \ni m =
\Subst{s}{M}{X}(z)$;
आणि
$\Sat{M}{\lforall[x][(X(x) \lif
X(u(x)))]}[\Subst{s}{M}{X}]$
सुद्धा,
कारण
$x \in M$
असेल तेव्हा
$f(x) \in
M$.
म्हणून
पूर्वपक्ष
आणि शर्त
दोन्हींची
पूर्ति
झाल्याने
उत्तरपक्षाचीही
पूर्ति झाली
पाहिजे:
$\Sat{M}{\lforall[x][X(x)]}[\Subst{s}{M}{X}]$.
पण याचा
अर्थ
$M = \Domain{M}$.
म्हणून
$\Domain{M}$
गणनीय
आहे, कारण
$M$ व्याख्येने
तसाच आहे.
\end{proof}

\begin{prob}
वाक्य
$\fn{Inf} \land \fn{Count}$
हे नेमक्या
गणनीय अनंत
आधारक्षेत्रांत
खरे असते.
$\fn{Count}$
ची व्याख्या
बदलून
आधारक्षेत्र
गणनीय अनंत
आहे हे
थेट व्यक्त
करणारे
वेगळे
वाक्य
मिळवा आणि
ते तसे
करते हे
सिद्ध करा.
\end{prob}



\chapter{द्वितीय-क्रम तर्कशास्त्राची अधिउपपत्ती}\label{sol:met:chap}









\section{परिचय}\label{sol:met:int:sec}


प्रथम-क्रम तर्कशास्त्राचे अनेक उपयुक्त गुणधर्म आहेत. ते निर्णेय नाही,
हे आपल्याला माहीत आहे; पण किमान ते स्वयंसिद्धकीकरणीय आहे. म्हणजे
प्रथम-क्रम तर्कशास्त्रासाठी निर्दोष आणि संपूर्ण अशा सिद्धता-प्रणाली
आहेत. त्या असा निष्पन्नता संबंध~$\Proves$ देतात की
वाक्यांचा कोणताही संच~$\Gamma$ आणि वाक्य~$\varphi$
यांसाठी $\Gamma \Entails \varphi$ जर आणि फक्त जर
$\Gamma \Proves \varphi$. विशेषतः, प्रथम-क्रम तर्कशास्त्रातील
वैध वाक्यांचा संच संगणनक्षमपणे प्रगणनीय आहे. या संगणनक्षम प्रगणनासाठी भाषा प्रभावी रीतीने
दिलेली असून सिद्धतांचेही प्रभावी प्रगणन करता येते, असे गृहीत आहे. एक संगणनक्षम फलन~$f\colon \Nat \to
\Sent[L]$ आहे, ज्यात~$f$ ची मूल्ये नेमकी~$\Lang{L}$ मधील सर्व
वैध वाक्ये आहेत. कारण निष्पत्त्यांचे प्रगणन
करता येते आणि त्यांपैकी एका~वाक्याला निष्पन्न
करणाऱ्या निष्पत्तींचे त्या वाक्य शी प्रतिचित्रण
करता येते. द्वितीय-क्रम तर्कशास्त्र प्रथम-क्रम तर्कशास्त्रापेक्षा
अधिक अभिव्यक्तिशील आहे; त्यामुळे त्याची वैधता ओळखणे
सर्वसाधारणपणे अधिक कठीण असते. प्रत्यक्षात, द्वितीय-क्रम
तर्कशास्त्र अनिर्णेय तर आहेच, पण त्याची वैधता
संगणनक्षमपणे प्रगणनीय सुद्धा नाही, हे आपण दाखवू. म्हणून
द्वितीय-क्रम तर्कशास्त्रासाठी निर्दोष, संपूर्ण आणि प्रभावी
सिद्धता-प्रणाली असू शकत नाही (मात्र निर्दोष पण अपूर्ण सिद्धता-प्रणाली उपलब्ध
आहेत आणि त्या संशोधनाचे महत्त्वाचे विषय आहेत).

प्रथम-क्रम तर्कशास्त्राचे आणखी दोन गुणधर्म आहेत: ते संहत
आहे (वाक्यांच्या संच~$\Gamma$ चा प्रत्येक सांत उपसंच
पूर्ततायोग्य असेल, तर $\Gamma$ स्वतःही पूर्ततायोग्य असतो)
आणि प्रगणनीय भाषेसाठी त्याला लोव्हेनहाइम--स्कोलेम प्रमेय लागू होते
($\Gamma$ चे अनंत प्रतिमान असेल, तर त्याचे
गणनीय अनंत प्रतिमानही असते). हे दोन्ही निष्कर्ष
द्वितीय-क्रम तर्कशास्त्रासाठी लागू होत नाहीत. कारण
द्वितीय-क्रम तर्कशास्त्राला प्रांताच्या आकाराविषयी
अशा गोष्टी व्यक्त करता येतात ज्या प्रथम-क्रम तर्कशास्त्राला
व्यक्त करता येत नाहीत.












\section{द्वितीय-क्रम अंकगणित}\label{sol:met:spa:sec}


पेआनो अंकगणिताची उपपत्ती~$\Th{PA}$ हिच्यात
$\Th{Q}$ ची आठ स्वयंसिद्धके समाविष्ट आहेत, हे आठवा:
\begin{align*}
& \lforall[x][\eq/[x'][\Obj 0]]\\
& \lforall[x][\lforall[y][(\eq[x'][y'] \lif \eq[x][y])]]\\
& \lforall[x][(\eq[x][\Obj 0] \lor \lexists[y][\eq[x][y']])]\\
& \lforall[x][\eq[(x + \Obj 0)][x]]\\
& \lforall[x][\lforall[y][\eq[(x + y')][(x + y)']]]\\
& \lforall[x][\eq[(x \times \Obj 0)][\Obj 0]]\\
& \lforall[x][\lforall[y][\eq[(x \times y')][((x \times y) + x)]]]\\
& \lforall[x][\lforall[y][(x < y \liff \lexists[z][\eq[(z' + x)][y]])]]\\
\intertext{आणि पुढील स्वरूपाची सर्व वाक्येही:}
& (\varphi(\Obj 0) \land \lforall[x][(\varphi(x) \lif \varphi(x'))]) \lif \lforall[x][\varphi(x)].
\end{align*}
शेवटचे हे एक ``रूपबंध'' आहे: अंकगणिताच्या भाषेतील
अमर्याद वाक्ये निर्माण करणारा नमुना, प्रत्येक
सूत्र~$\varphi(x)$ साठी एक वाक्य. इतर मुक्त चल असतील,
तर आरंभी त्यांच्यावर सार्विक संख्यापक जोडून वाक्यरूप घेतो. याला (प्रथम-क्रम)
\emph{विगमनाचा स्वयंसिद्धक-रूपबंध} म्हणतो.
\emph{द्वितीय-क्रम} पेआनो अंकगणित~$\Th{PA^2}$ मध्ये
विगमन एकाच वाक्यात मांडता येते. $\Th{PA^2}$ मध्ये
वरील पहिली आठ स्वयंसिद्धके आणि (द्वितीय-क्रम)
\emph{विगमनाचे स्वयंसिद्धक} असते:
\[\lforall[X][((X(\Obj 0) \land \lforall[x][(X(x) \lif X(x'))]) \lif \lforall[x][X(x)])].
\]
याचा अर्थ असा: प्रांताचा उपसंच~$X$ यामध्ये
$\Assign{\Obj{0}}{M}$ असेल आणि कोणत्याही~$x \in \Domain{M}$
बरोबर~$\Assign{\prime}{M}(x)$ ही असेल (म्हणजे तो
``उत्तराधिकारी क्रियेखाली बंद'' असेल), तर प्रांत मधील
सर्व घटक त्यात असतील (म्हणजे $X =
\Domain{M})$.

विगमनाचे स्वयंसिद्धक असा निर्वाळा देते की ते पूर्ण करणाऱ्या
कोणत्याही संरचनेच्या~$\Domain{M}$ मध्ये
स्वयंसिद्धकांना आवश्यक असलेलेच घटक असतात;
म्हणजे $n \in \Nat$ साठी~$\num{n}$ ची मूल्ये.
$\Th{PA^2}$ च्या प्रतिमानात अमानक संख्या नसतात.

\begin{thm}
\label{sol:met:spa:thm:sol-pa-standard}
जर $\Sat{M}{\Th{PA^2}}$ असेल, तर $\Domain{M} =
\Setabs{\Value{\num{n}}{M}}{n \in \Nat}$.
\end{thm}

\begin{proof}
$N = \Setabs{\Value{\num{n}}{M}}{n \in \Nat}$ असे मानू
आणि $\Sat{M}{\Th{PA^2}}$ असे गृहीत धरू. अर्थात,
कोणत्याही $n \in \Nat$ साठी
$\Value{\num{n}}{M} \in \Domain{M}$; म्हणून
$N \subseteq \Domain{M}$.

आता उलट दिशेचा अंतर्भाव दाखवू. $s(X) = N$ असे
चल-मूल्यांकन $s$ घेऊ. गृहीतकानुसार,
\begin{align*}
\Sat{M}{\lforall[X][((X(\Obj 0) \land \lforall[x][(X(x)
      \lif X(x'))]) \lif \lforall[x][X(x)])]}, & \text{ म्हणून}\\
\Sat{M}{(X(\Obj 0) \land \lforall[x][(X(x) \lif X(x'))]) \lif
  \lforall[x][X(x)]}[s]. &
\end{align*}
या सशर्त विधानाचे पूर्वांग पाहू. $\Value{\Obj{0}}{M} \in
N$, म्हणून $\Sat{M}{X(\Obj{0})}[s]$. दुसरे संयुक्तपद,
$\lforall[x][(X(x) \lif X(x'))]$, हेही पूर्ण होते.
कारण $x \in N$ असे समजू. $N$ च्या व्याख्येनुसार,
काही~$n$ साठी $x = \Value{\num{n}}{M}$. त्यावरून
$\Assign{\prime}{M}(x) = \Value{\num{n+1}}{M} \in N$.
म्हणून $\Assign{\prime}{M}(x) \in N$.

त्यामुळे $\Sat{M}{X(\Obj 0) \land \lforall[x][(X(x) \lif X(x'))]}[s]$.
म्हणून $\Sat{M}{\lforall[x][X(x)]}[s]$. याचा अर्थ,
प्रत्येक $x \in \Domain{M}$ साठी $x \in s(X) = N$.
म्हणून $\Domain{M} \subseteq N$.
\end{proof}

\begin{cor}
\label{sol:met:spa:cor:sol-pa-categorical}
$\Th{PA^2}$ ची कोणतीही दोन प्रतिमाने समरूप असतात.
\end{cor}

\begin{proof}
\ref{sol:met:spa:thm:sol-pa-standard} नुसार,
$\Th{PA^2}$ च्या कोणत्याही प्रतिमानाच्या प्रांतात
$\Value{\num{n}}{M}$ ही मूल्येच असतात. असे प्रत्येक
प्रतिमान~$\Th{Q}$ चेही प्रतिमान असते.
\ref{mod:mar:stm:prop:thq-standard} नुसार
ते मानक असते, म्हणजे~$\Struct{N}$ शी समरूप असते.
\end{proof}

वर $\Th{PA^2}$ ची व्याख्या पहिली आठ अंकगणितीय स्वयंसिद्धके
आणि द्वितीय-क्रम विगमनाचे स्वयंसिद्धक असलेली उपपत्ती म्हणून
केली. प्रत्यक्षात, द्वितीय-क्रम तर्कशास्त्राच्या
अभिव्यक्तीक्षमतेमुळे द्वितीय-क्रम पेआनो अंकगणितासाठी
अंकगणितीय स्वयंसिद्धकांपैकी फक्त \emph{पहिली दोन}
आणि विगमन पुरेसे असते.

\begin{prop}\label{sol:met:spa:prop:sol-pa-definable}
पहिली दोन अंकगणितीय स्वयंसिद्धके (उत्तराधिकारी स्वयंसिद्धके)
आणि द्वितीय-क्रम विगमनाचे स्वयंसिद्धक असलेली द्वितीय-क्रम
उपपत्ती $\Th{PA^{2\dagger}}$ असो. मग $\le$, $+$ आणि
$\times$ हे $\Th{PA^{2\dagger}}$ मध्ये परिभाष्य आहेत.
\end{prop}

\begin{proof}
$\le$ परिभाष्य आहे हे दाखवण्यासाठी असे
सूत्र~$\varphi_\le(x, y)$ शोधावे लागेल की
$\Sat{N}{\varphi_\le(\num n, \num m)}$ जर आणि फक्त जर
$n \le m$. पुढील सूत्र पाहू:
\[\psi(x, Y) \ident Y(x) \land \lforall[y][(Y(y) \lif Y(y'))]
\]
संच $Y \subseteq \Nat$ कडून $\psi(\num n, Y)$ ची
पूर्तता झाल्यास $\Setabs{m}{n \le m} \subseteq Y$ असते.
उलट, फक्त हा अंतर्भाव पुरेसा नाही; Y उत्तराधिकारी क्रियेखाली
बंदही असावा लागतो. n पासून पुढच्या सर्व नैसर्गिक संख्यांचा
संच स्वतः वरील सूत्राची पूर्तता करतो आणि अशा प्रत्येक Y मध्ये
तो समाविष्ट असतो. म्हणून
$\varphi_\le(x, y) \ident
\lforall[Y][(\psi(x, Y) \lif Y(y))]$ असे घेता येते.

बेरीज परिभाष्य आहे हे पाहण्यासाठी लक्षात घ्या की $k+l=m$
जर आणि फक्त जर असे फलन $u$ असेल की $u(0)=k$,
प्रत्येक $n$ साठी $u(n')=u(n)'$ आणि $m = u(l)$.
या समतुल्यतेने $\Th{PA^{2\dagger}}$ मध्ये पुढील
सूत्राद्वारे बेरीज परिभाषित करता येते:
\[\varphi_+(x,y,z) \ident \lexists[u][(u(\Obj 0)=x \land \lforall[w][u(w')=u
 (w)'] \land u(y)=z)]
\]
स्पष्टच, $\Sat{N}{\varphi_+(\num k, \num l, \num m)}$
जर आणि फक्त जर $k+l=m$.
\end{proof}

\begin{prob}
\ref{sol:met:spa:prop:sol-pa-definable} चा पुरावा पूर्ण करा.
\end{prob}












\section{द्वितीय-क्रम तर्कशास्त्र स्वयंसिद्धकीकरणीय नाही}\label{sol:met:nax:sec}


\begin{thm}
\label{sol:met:nax:thm:sol-undecidable}
द्वितीय-क्रम तर्कशास्त्र अनिर्णेय आहे.
\end{thm}

\begin{proof}
प्रथम-क्रम वाक्य प्रथम-क्रम तर्कशास्त्रात वैध असेल,
जर आणि फक्त जर ते द्वितीय-क्रम तर्कशास्त्रात वैध असेल.
प्रथम-क्रम तर्कशास्त्र अनिर्णेय आहे.
\end{proof}

\begin{thm}
\label{sol:met:nax:cor:sol-not-axiomatizable}
द्वितीय-क्रम तर्कशास्त्राच्या मानक चिन्हार्थमीमांसेसाठी
निर्दोष, संपूर्ण आणि प्रभावी निष्पत्ती प्रणाली अस्तित्वात नाही.
\end{thm}

\begin{proof}

अंकगणिताच्या प्रथम-क्रम भाषेतील वाक्य~$\varphi$ असो.
$\Sat{N}{\varphi}$ जर आणि फक्त जर $\Th{PA^2} \Entails \varphi$.
$\mathsf{P}$ हा~$\Th{PA^2}$ च्या नऊ स्वयंसिद्धकांचा संयोग असो.
$\Th{PA^2} \Entails \varphi$ जर आणि फक्त जर
$\Entails \mathsf{P} \lif \varphi$; म्हणजे प्रत्येक रचनेसाठी
$\Sat{M}{\mathsf{P} \lif \varphi}$. आता पुढील वाक्य पाहू:
$\lforall[z][\lforall[u][\lforall[u'][\lforall[u''][\lforall[L][(\mathsf{P}' \lif \varphi')]]]]]$.
हे $\Obj{0}$ च्या जागी $z$, $\prime$ च्या जागी
एकस्थानी फलनचल $u$, $+$ आणि $\times$ यांच्या जागी
अनुक्रमे द्विस्थानी फलनचले $u'$ आणि $u''$,
आणि $<$ च्या जागी द्विस्थानी संबंधचल~$L$ ठेवून
नंतर सार्वत्रिक संख्यापनाने मिळते. येथे नव्याने वापरलेले
सर्व चले आधीच्या सूत्रांत नसलेले निवडतो; आवश्यक असल्यास
आधी बद्ध चलांची नावे बदलतो, म्हणून capture होत नाही. हे शुद्ध द्वितीय-क्रम
तर्कशास्त्रातील वैध वाक्य असेल, जर आणि फक्त जर मूळ
वाक्य वैध असेल, जर आणि फक्त जर
$\Th{PA^2} \Entails \varphi$, जर आणि फक्त जर
$\Sat{N}{\varphi}$. म्हणून द्वितीय-क्रम तर्कशास्त्रासाठी
निर्दोष, संपूर्ण आणि प्रभावी सिद्धता-प्रणाली असती, तर
प्रथम-क्रम अंकगणितीय वाक्यांचे आणि त्यांच्या वरील निर्मित
वाक्यांसाठीच्या सिद्धतांचे समांतर प्रभावी प्रगणन वापरून
$\Struct{N}$ मध्ये सत्य असलेल्या वाक्यांचे
संगणनक्षम प्रगणन $f\colon \Nat \to \Sent[L_A]$
परिभाषित करता आले असते. हे फलन काही प्रथम-क्रम
सूत्र~$\psi_f(x, y)$ द्वारे~$\Th{Q}$ मध्ये निरूपित
करता आले असते. मग सूत्र~$\lexists[x][\psi_f(x, y)]$
हे~$\Struct{N}$ मधील सत्य प्रथम-क्रम वाक्यांचा संच
परिभाषित करील. हा टार्स्कीच्या प्रमेयाशी विरोध आहे.
\end{proof}














\section{द्वितीय-क्रम तर्कशास्त्र संहत नाही}\label{sol:met:com:sec}

\begin{explain}
वाक्यांचा संच~$\Gamma$ याचा प्रत्येक सांत उपसंच
पूर्ततायोग्य असेल, तर त्याला \emph{सांततः पूर्ततायोग्य}
म्हणतो. प्रथम-क्रम तर्कशास्त्राचा असा गुणधर्म आहे की
वाक्यांचा संच~$\Gamma$ सांततः पूर्ततायोग्य
असेल, तर तो पूर्ततायोग्य असतो. या गुणधर्माला
\emph{संहतता} म्हणतात. याचे परिणाम-संबंध वापरणारे
समतुल्य रूप असे: $\Gamma \Entails \varphi$ असेल,
तर काही सांत उपसंच~$\Gamma_0 \subseteq \Gamma$
यासाठी आधीच $\Gamma_0 \Entails
\varphi$ असते. या रूपात हा संपूर्णता प्रमेयाचा थेट
परिणाम आहे: $\Gamma \Entails \varphi$ असेल, तर
संपूर्णतेमुळे $\Gamma \Proves \varphi$. पण
निष्पत्तीमध्ये~$\Gamma$ मधील सांत
संख्येनेच वाक्ये वापरता येतात.

द्वितीय-क्रम तर्कशास्त्रासाठी संहतता लागू होत नाही.
द्वितीय-क्रम वाक्यांचे असे संच आहेत की ते
सांततः पूर्ततायोग्य असतात, पण पूर्ततायोग्य नसतात;
आणि ते एखाद्या~$\varphi$ चे परिणाम देतात, पण त्यांचा
कोणताही सांत उपसंच त्या~$\varphi$ चे परिणाम देत नाही.
\end{explain}


\begin{thm}
\label{sol:met:com:thm:sol-not-compact}
द्वितीय-क्रम तर्कशास्त्र संहत नाही.
\end{thm}

\begin{proof}
हे आठवा:
\[\fn{Inf} \ident \lexists[u][(\lforall[x][\lforall[y][(\eq[u(x)][u(y)]
      \lif \eq[x][y])]] \land \lexists[y][\lforall[x][\eq/[y][u(x)]]])]
\]
संरचनेचा प्रांत अनंत असेल, जर आणि फक्त जर
तिच्यात या सूत्राची पूर्ति होते. प्रांतात किमान
$n$ घटक आहेत, असे सांगणारे वाक्य
$\varphi^{\ge n}$ असो; उदाहरणार्थ,
\[\varphi^{\ge n} \ident \lexists[x_1][\dots\lexists[x_n][(\eq/[x_1][x_2]
    \land \eq/[x_1][x_3] \land \dots \land \eq/[x_{n-1}][x_n])]].
\]
वाक्यांचा पुढील संच पाहू:
\[\Gamma = \{\lnot \fn{Inf}, \varphi^{\ge 1}, \varphi^{\ge 2}, \varphi^{\ge 3}, \dots\}.
\]
तो सांततः पूर्ततायोग्य आहे: कोणताही सांत उपसंच~$\Gamma_0
\subseteq \Gamma$ घेतल्यास असा $k$ निवडता येतो की
$\varphi^{\ge k} \in \Gamma$ असते, पण $n>k$ साठी
कोणतेही $\varphi^{\ge n} \in \Gamma_0$ नसते.
$\Domain{M}$ मध्ये $k$ घटक असतील, तर
$\Sat{M}{\Gamma_0}$. मात्र $\Gamma$ पूर्ततायोग्य
नाही: $\Sat{M}{\lnot \fn{Inf}}$ असेल, तर
$\Domain{M}$ सांत असला पाहिजे; समजा त्याचा
आकार~$k$ आहे. मग $\Sat/{M}{\varphi^{\ge k+1}}$.
\end{proof}

\begin{prob}
संच~$\Gamma$ आणि वाक्य~$\varphi$ यांचे असे
उदाहरण द्या की $\Gamma \Entails \varphi$, पण प्रत्येक
सांत उपसंच~$\Gamma_0 \subseteq
\Gamma$ साठी $\Gamma_0 \Entails/ \varphi$.
\end{prob}













\section{द्वितीय-क्रम तर्कशास्त्रात लोव्हेनहाइम--स्कोलेम प्रमेय लागू होत नाही}\label{sol:met:lst:sec}

\begin{explain}
(अधोगामी) लोव्हेनहाइम--स्कोलेम प्रमेयानुसार, अनंत
प्रतिमान असलेल्या वाक्यांच्या प्रत्येक संचाला
गणनीय प्रतिमानही असते. हे प्रमेयदेखील
संपूर्णता प्रमेयाचा परिणाम आहे: संपूर्णतेचा पुरावा
वाक्यांच्या कोणत्याही सुसंगत संचासाठी प्रतिमान
तयार करतो आणि ते प्रतिमान गणनीय असते.
ऊर्ध्वगामी लोव्हेनहाइम--स्कोलेम प्रमेयही आहे.
त्यानुसार, वाक्यांच्या संचाला
गणनीय अनंत प्रतिमान असेल, तर त्याला
अगणनीय प्रतिमानही असते. ही दोन्ही
प्रमेये द्वितीय-क्रम तर्कशास्त्रात लागू होत नाहीत.
\end{explain}


\begin{thm}
\label{sol:met:lst:thm:sol-no-ls} लोव्हेनहाइम--स्कोलेम प्रमेय
द्वितीय-क्रम तर्कशास्त्रात लागू होत नाही:
काही वाक्यांना अनंत प्रतिमाने असतात,
पण गणनीय प्रतिमाने नसतात.
\end{thm}

\begin{proof}
हे आठवा:
\[\fn{Count} \ident \lexists[z][\lexists[u][\lforall[X][((X(z) \land
      \lforall[x][(X(x) \lif X(u(x)))]) \lif \lforall[x][X(x)])]]]
\]
हे संरचना~$\Struct{M}$ मध्ये सत्य असेल,
जर आणि फक्त जर $\Domain{M}$ गणनीय असेल.
म्हणून $\lnot\fn{Count}$ हे~$\Struct{M}$ मध्ये
सत्य असेल, जर आणि फक्त जर $\Domain{M}$
अगणनीय असेल. अशा संरचना
आहेत---कोणताही अगणनीय संच प्रांत
म्हणून घ्या; उदाहरणार्थ, $\Pow{\Nat}$ किंवा~$\Real$.
म्हणून $\lnot \fn{Count}$ ला अनंत प्रतिमाने
आहेत, पण गणनीय प्रतिमाने नाहीत.
\end{proof}

\begin{thm}
काही वाक्यांना गणनीय अनंत प्रतिमाने
असतात, पण अगणनीय प्रतिमाने नसतात.
\end{thm}

\begin{proof}
$\fn{Count} \land \fn{Inf}$ हे $\Nat$ मध्ये
सत्य असते; पण $\Domain{M}$ अगणनीय
असलेल्या कोणत्याही संरचना~$\Struct{M}$
मध्ये ते सत्य नसते.
\end{proof}




\chapter{द्वितीय-क्रम तर्कशास्त्र आणि संचसिद्धान्त}\label{sol:set:chap}
\begin{editorial}
या विभागात द्वितीय-क्रम तर्कशास्त्रात घातसंच आणि
सांतत्य यांचे संकेतीकरण केले आहे. निष्कर्ष मांडले
आहेत; पण त्यांचे पुरावे अजून भरायचे आहेत. अद्याप
सरावप्रश्न नाहीत---आणि व्याख्या व निष्कर्षांमध्येही
त्रुटी असू शकतात. सावधपणे वापरा; काही असत्य
किंवा अस्पष्ट आढळल्यास कळवा.
\end{editorial}









\section{परिचय}\label{sol:set:int:sec}

द्वितीय-क्रम तर्कशास्त्रात प्रांताच्या उपसंचांवर तसेच
फलनांवर संख्यापन करता येते; म्हणून त्यात संचसिद्धान्ताचा
काही भाग तरी मांडता येईल अशी अपेक्षा आहे. येथे
‘मांडणे’ याचा अर्थ संचसिद्धान्तीय गुणधर्म आणि विधाने
द्वितीय-क्रम तर्कशास्त्रात व्यक्त करता येणे, तेही
संचांसाठी कोणतीही खास अतार्किक चिन्हे (उदा.,
संचसिद्धान्तातील सदस्यत्वाचे विधेय) न वापरता.
उदाहरणार्थ, संचांचे संयोग आणि छेद तसेच उपसंच-संबंध
परिभाषित करता येतात; शिवाय संचांच्या आकारांची तुलना
करता येते आणि कँटरच्या प्रमेयासारखे निष्कर्षही मांडता येतात.












\section{संचांची तुलना}\label{sol:set:cmp:sec}

\begin{prop}
सूत्र $\lforall[x][(X(x) \lif Y(x))]$ उपसंच-संबंध
परिभाषित करते; म्हणजे
$\Sat{M}{\lforall[x][(X(x) \lif Y(x))]}[s]$
जर आणि फक्त जर $s(X) \subseteq s(Y)$.
\end{prop}

\begin{prop}
सूत्र $\lforall[x][(X(x) \liff Y(x))]$ संचांवरील
समानता-संबंध परिभाषित करते; म्हणजे
$\Sat{M}{\lforall[x][(X(x) \liff Y(x))]}[s]$
जर आणि फक्त जर $s(X) = s(Y)$.
\end{prop}

\begin{prop}
सूत्र $\lexists[x][X(x)]$ अरिक्त असण्याचा
गुणधर्म परिभाषित करते; म्हणजे
$\Sat{M}{\lexists[x][X(x)]}[s]$
जर आणि फक्त जर $s(X) \neq \emptyset$.
\end{prop}

संच~$X$ हा संच~$Y$ पेक्षा आकारमानाने मोठा नाही,
$\cardle{X}{Y}$, जर आणि फक्त जर
एकास-एक फलन $f\colon X \to Y$ असेल.
फलन एकास-एक आहे, तसेच~$X$ मधील निविष्टांसाठी
त्याची मूल्ये~$Y$ मध्ये आहेत, हे व्यक्त करता येते.
म्हणून प्रांताच्या उपसंचांवर आकारमानाने मोठा नाही
हा संबंधही परिभाषित करता येतो.

\begin{prop}

पुढील सूत्र
\[\lexists[u][(\lforall[x][(X(x) \lif Y(u(x)))] \land \lforall[x][\lforall[y][((X(x) \land X(y) \land \eq[u(x)][u(y)]) \lif
      \eq[x][y])]])]
\]
आकारमानाने मोठा नाही हा संबंध परिभाषित करते. येथे
एकास एक अट फक्त X मधील दोन घटकांसाठी मागितली आहे.
X वरील फलन पूर्ण प्रांतावर मनमानी वाढवता येते.
\end{prop}

दोन संचांचे आकारमान समान असते, म्हणजे ते
``तुल्यबल'' असतात, $\cardeq{X}{Y}$, जर आणि
फक्त जर एकास-एक व आच्छादक फलन~$f\colon X \to Y$ असेल.

\begin{prop}
पुढील सूत्र
\begin{multline*}
  \lexists[u][(\lforall[x][(X(x) \lif Y(u(x)))] \land {}\\
    \lforall[x][\lforall[y][((X(x) \land X(y) \land \eq[u(x)][u(y)]) \lif \eq[x][y])]]
  \land {} \\\lforall[y][(Y(y) \lif \lexists[x][(X(x)
      \land \eq[y][u(x)])])])]
\end{multline*}
तुल्यबल असण्याचा संबंध परिभाषित करते.
\end{prop}

या सूत्रे ना अनुक्रमे $X \subseteq Y$,
$X = Y$, $X \neq \emptyset$, $\cardle{X}{Y}$
आणि $\cardeq{X}{Y}$ अशी संक्षिप्त चिन्हे वापरू.
(संच $X$ आणि $Y$ यांच्याविषयी अनौपचारिकपणे
बोलतानाही हीच चिन्हे वापरल्यामुळे थोडा गोंधळ होऊ शकतो;
पण येथे ती एकस्थानी संबंधचले $X$ आणि~$Y$
असलेल्या द्वितीय-क्रम तर्कशास्त्रातील सूत्रे
साठीची संक्षिप्त चिन्हे आहेत.)

\begin{prop}
वाक्य $\lforall[X][\lforall[Y][((\cardle{X}{Y} \land
    \cardle{Y}{X}) \lif \cardeq{X}{Y})]]$ वैध आहे.
\end{prop}

\begin{proof}
कोणत्याही उपसंचांसाठी $X \subseteq \Domain{M}$
आणि $Y \subseteq \Domain{M}$, जर
$\cardle{X}{Y}$ आणि $\cardle{Y}{X}$, तर
$\cardeq{X}{Y}$ असेल, तर संरचना~$\Struct{M}$
मध्ये या वाक्याची पूर्ति होते. पण हे
\emph{कोणत्याही} संच $X$ आणि $Y$ साठी खरे आहे:
हेच श्रेडर--बर्नस्टाइन प्रमेय आहे.
\end{proof}













\section{संचांची आकारमाने}\label{sol:set:crd:sec}

\begin{explain}
प्रांत सांत किंवा अनंत, गणनीय किंवा
अगणनीय आहे हे जसे व्यक्त करता येते,
तसेच~$\Domain{M}$ चा उपसंच सांत किंवा अनंत,
गणनीय किंवा अगणनीय असण्याचा
गुणधर्मही परिभाषित करता येतो.
\end{explain}

\begin{prop}

सूत्र $\fn{Inf}(X) \ident$
\begin{multline*}
\lexists[u][(\lforall[x][(X(x) \lif X(u(x)))] \land {} \\
\lforall[x][\lforall[y][((X(x) \land X(y) \land \eq[u(x)][u(y)]) \lif
      \eq[x][y])]] \land {} \\
\lexists[y][(X(y) \land \lforall[x][(X(x) \lif \eq/[y][u(x)])])])]
\end{multline*}
याची चल-मूल्यांकन~$s$ नुसार पूर्ति होईल, जर आणि
फक्त जर $s(X)$ अनंत असेल. स्रोतदुरुस्ती: फलन
X ला स्वतःमध्येच पाठवते, X वर एकास एक असते आणि X मधील
किमान एक घटक त्याच्या X वरील व्याप्तीत नसतो, अशा सर्व अटी
घेतल्या आहेत. त्यामुळे मूळ सूत्राचा सांत उपसंचावरील counterexample
दूर होतो. अनंततेचा नेहमीचा निवडीसह संदर्भ येथेही लागू आहे.
\end{prop}

\begin{prop}
सूत्र $\fn{Count}(X) \ident $
\begin{multline*}
\lnot\lexists[z][X(z)] \lor {} \\
\lexists[z][\lexists[u][(X(z) \land
    \lforall[x][(X(x) \lif X(u(x)))] \land {} \\
\lforall[Y][((Y \subseteq X \land Y(z) \land
      \lforall[x][(Y(x) \lif Y(u(x)))]) \lif X = Y)])]]
\end{multline*}
याची चल-मूल्यांकन~$s$ नुसार पूर्ति होईल, जर आणि
फक्त जर $s(X)$ गणनीय असेल. स्रोतदुरुस्ती:
पहिला खंड रिकामा संच घेतो. अरिक्त X साठी आरंभीचा घटक
आणि X ला स्वतःमध्ये पाठवणारे फलन घेतले आहे. शेवटची अट
फक्त X मधील बंद उपसंचांबद्दल आहे; त्यामुळे X हा आरंभीच्या
घटकाचा पूर्ण orbit असतो. हे संपूर्ण प्रांतच X असण्याची
चुकीची सक्ती करत नाही.
\end{prop}

कँटरच्या प्रमेयावरून अगणनीय संच आहेत,
आणि अनंत आकारमानांचे निरनिराळे अमर्याद स्तरही
आहेत, हे आपल्याला माहीत आहे. संचसिद्धान्तात
संचांच्या आकारमानांचे संपूर्ण अंकगणित विकसित केले
जाते आणि संचांना अनंत संचांक दिले जातात. सांत
संचांची आकारमाने मोजणारे संचांक म्हणजे स्वाभाविक
संख्या. गणनीय अनंत संचांचा संचांक हा पहिला
अनंत संचांक असून त्याला~$\aleph_0$ (``अलेफ-शून्य'')
म्हणतात. पुढील अनंत आकारमान~$\aleph_1$ आहे.
गणनीय नसलेल्या संचाला असू शकणारे हे
सर्वांत लहान आकारमान आहे (गणनीय अनंत संचाचे
आकारमान~$\aleph_0$ असते). ``$X$ चे आकारमान
$\aleph_0$ आहे'' याची व्याख्या
$\fn{Aleph}_0(X) \liff \fn{Inf}(X) \land \fn{Count}(X)$
अशी करता येते. $X$ चे आकारमान $\aleph_1$
असेल, जर आणि फक्त जर त्याच्यापेक्षा काटेकोरपणे लहान
आकारमानाचे सर्व उपसंच सांत किंवा~$\aleph_0$ आकारमानाचे
असतील आणि तो स्वतः सांत किंवा~$\aleph_0$ आकारमानाचा
नसेल. म्हणून पुढील दुरुस्त सूत्र घेतो:
$\fn{Aleph_1}(X) \ident \lnot\fn{Count}(X) \land
\lforall[Y][((Y \subseteq X \land \lnot\cardeq{Y}{X}) \lif
  (\lnot\fn{Inf}(Y) \lor \fn{Aleph}_0(Y)))]$. $\aleph_2$ आकारमानाची
व्याख्याही अशाच प्रकारे करायची, असे त्यात म्हटले आहे.
स्रोतदुरुस्ती: X अगणनीय असावा लागतो. स्वतः X किंवा
त्याच्याशी तुल्यबल असलेला उपसंच हा लहान आकारमानाचा नाही;
त्याला finite किंवा aleph-zero असण्याची चुकीची अट लावलेली नाही.

एक विशेष महत्त्वाचे आकारमान म्हणजे तथाकथित
सांतत्याचा संचांक. ते $\Pow{\Nat}$ चे आकारमान,
किंवा समतुल्यपणे~$\Real$ चे आकारमान आहे.
एखादा संच सांतत्याच्या आकारमानाचा आहे, हेही
द्वितीय-क्रम तर्कशास्त्रात व्यक्त करता येते; मात्र
त्यासाठी आणखी थोडे काम करावे लागते.












\section{सांतत्याचे आकारमान}\label{sol:set:pow:sec}

\begin{explain}

द्वितीय-क्रम तर्कशास्त्रात प्रांताच्या उपसंचांवर
संख्यापन करता येते; पण प्रांताच्या उपसंचांच्या
संचांवर करता येत नाही. ते थेट करण्यासाठी
\emph{तृतीय-क्रम} तर्कशास्त्र लागेल. उदाहरणार्थ,
संचाच्या घातसंचापासून मूळ संचाकडे एकास-एक
फलन नसते हे कँटरचे प्रमेय मांडायचे असल्यास आपण
असे म्हणण्याचा प्रयत्न करू शकतो: ``प्रत्येक
संच~$X$ आणि प्रत्येक संच~$P$ यांसाठी, $P$
हा~$X$ चा घातसंच असेल, तर
$\cardle{P}{X}$ नसते.'' पण $P$ हा~$X$ चा
घातसंच आहे असे म्हणण्यासाठी $P$ चे घटक
नेमके~$X$ चे सर्व उपसंच आहेत, हे सूत्रबद्ध करावे
लागेल; उदाहरणार्थ,
$\lforall[Y][(P(Y) \liff Y \subseteq X)]$.
अडचण $P(Y)$ मध्ये आहे: ते द्वितीय-क्रम
तर्कशास्त्रातील सूत्र नाही, कारण
$P$ सारख्या एकस्थानी संबंधचलाचे निविष्ट फक्त
पदेच असू शकतात.

मात्र प्रांत पुरेसा मोठा असेल, तर संचांच्या संचांवरचे
संख्यापन \emph{अनुकरणाने} करता येते. कल्पना अशी
की द्विस्थानी संबंध~$R$ हा प्रांत मधील
घटक ना त्याच प्रांत मधील
घटक शी जोडतो. असा~$R$ दिल्यास,
एखाद्या~$x$ शी $R$-संबंधित सर्व घटक
गोळा करता येतात: $\Setabs{y \in
  \Domain{M}}{R(x,y)}$ हा~$x$ ने
``सांकेतिक केलेला'' संच आहे. उलट,
$Z \subseteq \Pow{\Domain{M}}$ हा~$\Domain{M}$
च्या उपसंचांचा संग्रह असेल आणि~$\Domain{M}$
मध्ये~$Z$ मधील संचांइतके तरी घटक
असतील, तर असा संबंध~$R \subseteq \Domain{M}^2$
असतो की प्रत्येक $Y \in Z$ हा
काही~$x$ ने~$R$ द्वारे सांकेतिक केला जातो.
\end{explain}

\begin{defn}
$R \subseteq \Domain{M}^2$ असेल, तर
$\Setabs{y \in \Domain{M}}{R(x, y)}$ या संचाचे
\emph{$x$ ने $R$-सांकेतीकरण} होते.
\end{defn}

घटक~$x \in \Domain{M}$ हा
$Z \subseteq \Domain{M}$ संचाचे $R$-सांकेतीकरण
करत असेल, तर संच $Y \subseteq \Domain{M}$ हा
संचांच्या संचाचे सांकेतीकरण करतो: ते म्हणजे~$Y$
मधील घटक नी सांकेतिक केलेले संच.
म्हणून संच $Y$ हा $\Pow{X}$ चे
$R$-सांकेतीकरण करू शकतो. असे होईल, जर
आणि फक्त जर प्रत्येक $Z \subseteq X$ साठी
काही $x \in Y$ हे~$Z$ चे $R$-सांकेतीकरण
करत असेल आणि प्रत्येक $x \in Y$ हा
काही~$Z \subseteq X$ चे $R$-सांकेतीकरण करत असेल.

\begin{prop}
सूत्र
  \[  \fn{Codes}(x, R, Z) \ident \lforall[y][(Z(y) \liff
    R(x, y))]
  \]
हे $s(x)$ हा $s(R)$ द्वारे $s(Z)$ चे
सांकेतीकरण करतो, असे व्यक्त करते. पुढील
सूत्र
\begin{multline*}
  \fn{Pow}(Y, R, X) \ident \\
  \lforall[Z][(Z \subseteq X \lif \lexists[x][(Y(x) \land
    \fn{Codes}(x, R, Z))])] \land {} \\ \lforall[x][(Y(x) \lif
  \lforall[Z][(\fn{Codes}(x, R, Z) \lif Z \subseteq X)])]
\end{multline*}
हे $s(Y)$ हा $s(R)$ द्वारे~$s(X)$ च्या
घातसंचाचे सांकेतीकरण करतो, असे व्यक्त करते;
म्हणजे $s(Y)$ चे घटक $s(R)$ द्वारे
नेमके $s(X)$ चे सर्व उपसंच सांकेतिक करतात.
\end{prop}

\begin{explain}
या युक्तीने उपसंचांवर थेट संख्यापन करण्याऐवजी
त्यांच्या संकेतांवर संख्यापन करून घातसंचाविषयी
विधाने व्यक्त करता येतात. उदाहरणार्थ,
आता कँटरचे प्रमेय असे मांडता येते की~$X$ चा
घातसंच सांकेतिक करणाऱ्या कोणत्याही संबंधाच्या
प्रांतापासून मूळ~$X$ कडे एकास-एक फलन नसते.
\end{explain}

\begin{prop}
पुढील वाक्य
\begin{multline*}
  \lforall[X][\lforall[Y][\lforall[R][(\fn{Pow}(Y, R, X) \lif \\
      \lnot \lexists[u][(\lforall[x][\lforall[y][((Y(x) \land Y(y) \land \eq[u(x)][u(y)]) \lif
            \eq[x][y])]] \land {}\\
        \lforall[x][(Y(x) \lif X(u(x)))])])]]]
\end{multline*}
वैध आहे.
\end{prop}

\begin{explain}
गणनीय अनंत संचाचा घातसंच अगणनीय
असतो; म्हणून त्याचा संचांक कोणत्याही
गणनीय अनंत संचाच्या संचांकापेक्षा
मोठा असतो (तो संचांक $\aleph_0$ असतो).
$\Pow{\Nat}$ चे आकारमान ``सांतत्याचे आकारमान''
म्हणतात, कारण ते वास्तव संख्यारेषेवरील
बिंदूंच्या संचाच्या, म्हणजे~$\Real$ च्या,
आकारमानाइतके असते. एखाद्या
गणनीय अनंत संचाचा घातसंच सांकेतिक
करण्याइतका प्रांत मोठा असेल, तर एखादा
संच सांतत्याच्या आकारमानाचा आहे हे असे
मांडता येते: तो~$\aleph_0$ आकारमानाच्या संच~$X$
चा घातसंच सांकेतिक करणाऱ्या कोणत्याही
संच~$Y$ शी तुल्यबल आहे. (प्रांत पुरेसा मोठा
नसेल, म्हणजे त्यात~$\Real$ शी तुल्यबल
उपसंचच नसेल, तर~$\Pow{X}$ चे सांकेतीकरण
करणारा संबंधही असू शकत नाही.)
\end{explain}

\begin{prop}
$\cardle{\Real}{\Domain{M}}$ असेल, तर पुढील
सूत्र
\begin{multline*}
\fn{Cont}(Y) \ident
\lexists[X][\lexists[R][((\fn{Aleph}_0(X) \land
      \fn{Pow}(Y, R, X)) \land \\
      \lforall[x][\lforall[y][((Y(x) \land {}
      Y(y) \land \lforall[z][(R(x,z) \liff R(y,z))]) \lif \eq[x][y])]])]]
\end{multline*}
हे $\cardeq{s(Y)}{\Real}$ असे व्यक्त करते.
\end{prop}

\begin{proof}
  $\fn{Pow}(Y, R, X)$ यानुसार $s(Y)$ हा $s(R)$
  द्वारे~$s(X)$ चा घातसंच सांकेतिक करतो;
  $\fn{Aleph}_0(X)$ नुसार हा संच गणनीय अनंत असतो.
  म्हणून $s(Y)$ चे आकारमान किमान सांतत्याच्या
  आकारमानाइतके असते; ते अधिकही असू शकते
  (जर $s(Y)$ चे अनेक घटक~$X$ चा तोच उपसंच
  सांकेतिक करत असतील). शेवटचे संयुक्तपद
  हे टाळते: $s(R)$ द्वारे $s(Y)$ चे घटक आणि
  $s(X)$ चे उपसंच यांच्यातील संगती
  एकास-एक असावी, अशी त्याची अट आहे.
  उलट, Y सांतत्याच्या आकारमानाचा असेल, तर प्रांतातील
  गणनीय अनंत उपसंच X च्या घातसंचाशी त्याची bijection निवडता
  येते. त्या bijection मधील सदस्यत्वाने R परिभाषित केल्यास
  Pow आणि unique-code या दोन्ही अटी मिळतात. येथे आधीच्या
  विभागातील दुरुस्त Count आणि Aleph-zero definitions वापरल्या आहेत.
\end{proof}

\begin{prop}
$\cardeq{\Domain{M}}{\Real}$ जर आणि फक्त जर
\begin{multline*}
  \Sat{M}{\lexists[Y][(\lforall[x][Y(x)] \land \fn{Cont}(Y))]}.
        \end{multline*}
  स्रोतदुरुस्ती: Y हाच पूर्ण प्रांत असणे आणि त्याचे
  आकारमान Cont ने ओळखले जाणे स्पष्ट मागितले आहे. फक्त
  काही घातसंच code करणे व जुन्या range अटी घेणे अपुरे होते:
  identity फलनाने त्या अटी अधिक मोठ्या प्रांतातही पूर्ण होतात.
  \end{prop}

\begin{explain}
सांतत्य परिकल्पनेनुसार सांतत्याचे आकारमान
हा पहिला अगणनीय संचांक असतो;
म्हणजे $\Pow{\Nat}$ चे आकारमान~$\aleph_1$ असते.
\end{explain}

\begin{prop}
सांतत्य परिकल्पना खरी असेल, जर आणि फक्त जर
\[\fn{CH} \ident
\lforall[X][(\fn{Aleph}_1(X) \liff \fn{Cont}(X))]\]
वैध असेल. येथे आधी दुरुस्त केलेल्या आकारमानांच्या
definitions वापरल्या आहेत. सांतत्य आणि अलेफ-एक समान असतील,
तर दोन्ही गुणधर्म प्रत्येक उपसंचासाठी समान ठरतात. उलट,
सांतत्याच्या आकारमानाचा प्रांत घेऊन त्यातील अलेफ-एक
आकारमानाच्या उपसंचाला हे विधान लावल्यास अपेक्षित समानता मिळते.
\end{prop}

हे लक्षात घ्या की $\lnot \fn{CH}$ वैध असेल,
जर आणि फक्त जर सांतत्य परिकल्पना खोटी असेल,
असे होत नाही. गणनीय प्रांतात
आकारमान~$\aleph_1$ असलेले उपसंच नसतात,
आणि सांतत्याच्या आकारमानाचे उपसंचही नसतात;
म्हणून अशा गणनीय प्रांतात
$\fn{CH}$ नेहमी खरे असते. मात्र सांतत्य
परिकल्पना खोटी असेल, जर आणि फक्त जर
वैध असेल असे पुढील वेगळे वाक्य देता येते:

\begin{prop}
सांतत्य परिकल्पना खोटी असेल, जर आणि
फक्त जर \[\fn{NCH} \ident
\lforall[X][(\fn{Cont}(X) \lif \lexists[Y][(Y \subseteq X \land \lnot
    \fn{Count}(Y) \land \lnot \cardeq{X}{Y})])]\]
वैध असेल. सांतत्य अलेफ-एकपेक्षा मोठे असेल, तर
प्रत्येक सांतत्य आकारमानाच्या संचात अलेफ-एक आकारमानाचा
अगणनीय, त्याच्याशी तुल्यबल नसलेला उपसंच निवडता येतो.
सांतत्य अलेफ-एक असेल, तर अशा लहान आकारमानाचा अगणनीय
उपसंच नसतो. लहान प्रांतांत Cont-संच नसल्याने शर्त रिक्ततः खरी असते.
\end{prop}



\part{लॅम्डा कलन}\label{lam:part}
\begin{editorial}
हा भाग लॅम्डा कलनाविषयी आहे. परिचयाचे प्रकरण जेरेमी अविगाड यांच्या
टिपणांवर आधारित आहे; त्यातील काही भाग आता पुनरुक्त झाला असून पुढील
प्रकरणांमध्ये समाविष्ट आहे. विन्यासमीमांसा, चर्च--रॉसर गुणधर्म आणि
लॅम्डा-परिभाष्यता यांवरील प्रकरणे झेसन चियान यांनी मिटॅक्सच्या उन्हाळी
प्रशिक्षणकाळात तयार केली. त्यांचे अजून परीक्षण आणि पुनरीक्षण करायचे आहे.
\end{editorial}
\chapter{परिचय}\label{lam:int:chap}
\begin{editorial}
या प्रकरणात जेरेमी यांनी लॅम्डा कलनावर लिहिलेली मूळ संक्षिप्त टिपणे आहेत.
त्यातील विभाग एकत्र करायचे आहेत; तसेच लॅम्डा-परिभाष्यतेविषयीचा मजकूर
त्या विषयावरील स्वतंत्र, अधिक तपशीलवार प्रकरणातील मजकुरात मिसळायचा आहे.
\end{editorial}








\section{आढावा}\label{lam:int:ovr:sec}

अॅलॉन्झो चर्च यांनी १९३०च्या दशकाच्या सुरुवातीला लॅम्डा कलनाची रचना
रचनाशील तर्कशास्त्राचा आधार म्हणून केली होती; संगणनक्षम फलनांचे
प्रतिमान म्हणून \emph{नव्हे}. पण ते ट्यूरिंग-संगणनक्षम फलने आणि आंशिक
पुनरावर्ती फलने यांसारख्या संगणनक्षमतेच्या इतर व्याख्यांशी समतुल्य
असल्याचे लवकरच दाखवले गेले. सुरुवातीला हे किंचित अनपेक्षित वाटले,
ही गोष्ट हे वैशिष्ट्यीकरण अधिकच रंजक बनवते.

फलनाला स्वतंत्र नाव देण्याऐवजी, त्याची व्याख्या करणाऱ्या सांकेतिक
व्यक्तीकरणाने थेट त्याचा निर्देश करण्याची लॅम्डा संकेतन ही सोयीची
पद्धत आहे. “$f$ हे $f(x) = x + 3$ ने परिभाषित फलन असू द्या,” असे
म्हणण्याऐवजी “$f$ हे $\lambd[x][(x + 3)]$ फलन असू द्या,” असे म्हणता
येते. दुसऱ्या शब्दांत, $\lambd[x][(x+3)]$ हे फलसाधकाच्या मूल्यात तीन मिळवणाऱ्या
फलनाचे केवळ \emph{नाव} आहे. या व्यक्तीकरणातील $x$ हे स्थानधारक आहे: त्याच फलनाचा निर्देश
$\lambd[y][(y + 3)]$ नेही करता येतो. इतर प्राचले असतानाही हे संकेतन
चालते. उदाहरणार्थ, $g(x, y)$ हे दोन चरांचे फलन आणि $k$ ही नैसर्गिक
संख्या असेल, तर $\lambd[x][g(x,k)]$ हे कोणत्याही~$x$ ला
~$g(x, k)$ शी जोडणारे फलन आहे.

सांकेतिक व्यक्तीकरणावरून फलन परिभाषित करण्याच्या या पद्धतीला
\emph{लॅम्डा अमूर्तीकरण} म्हणतात. लॅम्डा अमूर्तीकरणाची दुसरी बाजू
म्हणजे \emph{उपयोजन}: आपल्याकडे एखादे फलन $f$ असेल (समजा, नैसर्गिक
संख्यांवर परिभाषित), तर त्याचे 2 सारख्या कोणत्याही मूल्यावर
\emph{उपयोजन} करता येते. अर्थात, प्रचलित संकेतनात आपण मिळणाऱ्या
मूल्यासाठी~$f(2)$ लिहितो.

लॅम्डा अमूर्तीकरण आणि उपयोजन एकत्र केल्यास काय होते? मग अमूर्तीकृत
चराच्या जागी ज्यावर उपयोजन केले ते मूल्य “बसवून” परिणामी
व्यक्तीकरण सोपे करता येते. उदाहरणार्थ,
\[(\lambd[x][(x + 3)])(2)
\]
हे~$2 + 3$ असे सोपे करता येते.

आतापर्यंत आपण प्रचलित संकल्पनांसाठी नवे संकेतन मांडण्यापलीकडे काही
केलेले नाही. मात्र लॅम्डा कलन संचसैद्धान्तिक दृष्टिकोनापासून अधिक
मूलगामी फारकत घेते. या चौकटीत:
\begin{enumerate}
\item सर्वकाही फलन दर्शवते.
\item लॅम्डा अमूर्तीकरणाने फलने परिभाषित करता येतात.
\item कोणत्याही गोष्टीचे इतर कोणत्याही गोष्टीवर उपयोजन करता येते.
\end{enumerate}
उदाहरणार्थ, $F$ हे लॅम्डा कलनातील पद असेल, तर $F(F)$ अर्थपूर्ण
असल्याचे नेहमी गृहीत धरले जाते. या उदार चौकटीला
\emph{प्रकारविरहित} लॅम्डा कलन म्हणतात; येथे “प्रकारविरहित” याचा अर्थ
“कशाचे कशावर उपयोजन करता येईल यावर बंधन नाही” असा आहे.

\begin{digress}
\emph{प्रकारयुक्त} लॅम्डा कलनही आहे; तो प्रकारविरहित आवृत्तीचा
एक महत्त्वाचा पर्याय आहे. प्रकारयुक्त लॅम्डा कलन अनेक बाबतींत
प्रकारविरहित कलनासारखे असले, तरी त्याची अभिजात संचसैद्धान्तिक
चौकटीशी सांगड घालणे खूप सोपे आहे आणि त्याचे काही गुणधर्मही
खूप वेगळे आहेत.

लॅम्डा कलनावरील संशोधन सैद्धान्तिक संगणकशास्त्रात आणि
प्रोग्रामिंग भाषांच्या रचनेत मध्यवर्ती ठरले आहे. जॉन मॅकार्थी
यांनी १९५०च्या दशकात तयार केलेली LISP ही या कल्पनांचा
प्रभाव असलेल्या भाषेचे सुरुवातीचे उदाहरण आहे.
\end{digress}











\section{लॅम्डा कलनाची विन्यासमीमांसा}\label{lam:int:syn:sec}

सुरुवातीला $x$, $y$, $z$,~\dots अशी चरांची मालिका आणि $a$, $b$,
$c$,~\dots अशी काही स्थिरांक चिन्हे घेतली जातात. पदांचा संच पुढीलप्रमाणे
विगमनाधारित रीतीने परिभाषित केला जातो:
\begin{enumerate}
\item प्रत्येक चर हे पद आहे.
\item प्रत्येक स्थिरांक हे पद आहे.
\item $M$ आणि $N$ ही पदे असतील, तर $(MN)$ हेही पद आहे.
\item $M$ हे पद आणि $x$ हे चर असेल, तर $(\lambd[x][M])$ हे
  पद आहे.
\end{enumerate}
$(MN)$ या रूपातील पदांना \emph{उपयोजने}, तर $(\lambd[x][M])$
या रूपातील पदांना \emph{अमूर्तीकरणे} म्हणतात.

एकही स्थिरांक नसलेल्या प्रणालीला \emph{शुद्ध} लॅम्डा कलन म्हणतात.
आपण मुख्यतः शुद्ध $\lambd$-कलनात काम करणार असल्याने सर्व लहान इंग्रजी
अक्षरे चरे दर्शवतील. $\lambd$-कलनातील पदे दर्शवण्यासाठी आपण मोठी
इंग्रजी अक्षरे ($M$, $N$ इत्यादी) वापरू.

आपण संकेतनाच्या काही रूढी पाळू:

\begin{conv}
\begin{enumerate}
\item कंस वगळले असता उपयोजन डावीकडून उजवीकडे होते. उदाहरणार्थ,
  $M$, $N$, $P$ आणि $Q$ ही पदे असतील, तर $MNPQ$ हे
  $(((MN)P)Q)$ चे संक्षिप्त रूप आहे.
\item पुन्हा, कंस वगळले असता लॅम्डा अमूर्तीकरणाला शक्य तितकी
  व्यापक व्याप्ती द्यायची. उदाहरणार्थ, $\lambd[x][MNP]$ याचा
  अर्थ $(\lambd[x][((MN)P)])$ असा घ्यायचा.
\item एका लॅम्डाने अनेक चरे अमूर्त करता येतात. उदाहरणार्थ,
  $\lambd[xyz][M]$ हे
  $\lambd[x][\lambd[y][\lambd[z][M]]]$ चे संक्षिप्त रूप आहे.
\end{enumerate}
\end{conv}

उदाहरणार्थ,
\[\lambd[xy][xxyx \lambd[z][xz]]
\]
हे पुढीलचे संक्षिप्त रूप आहे:
\[\lambd[x][\lambd[y][((((xx)y)x)(\lambd[z][(xz)]))]].
\]
या रूढी पाठ करा. सुरुवातीला त्या तुमची डोकेदुखी वाढवतील;
पण हळूहळू त्यांची सवय होईल आणि मग कंसांच्या पसाऱ्यापेक्षा
त्या कमी त्रासदायक वाटतील.


इतर चरांचा बंधन-संबंध न बदलता ज्या दोन पदांत फक्त
बद्ध चरांची नावे बदलली आहेत, त्यांना
$\alpha$-सममूल्य म्हणतात; उदाहरणार्थ, $\lambd[x][x]$ आणि
$\lambd[y][y]$. त्यांना “तेच” पद मानणे सोयीचे ठरेल; म्हणजे
$M$ आणि $N$ ही पदे एकच आहेत असे म्हणताना “बद्ध चरांच्या
नामांतराचा फरक सोडून” ती एकच आहेत, असाही अर्थ अभिप्रेत असेल.
हे नामांतर करताना मुक्त चर पकडले जाणार नाहीत याची काळजी घेतली जाते.
एखाद्या चराचा वापर त्याच नावाच्या चराला अमूर्त करणाऱ्या
$\lambd$ च्या व्याप्तीत असेल, तर तो वापर “बद्ध” असतो;
असे बंधन नसलेला वापर “मुक्त” असतो. एकाच वापराला अनेक
अमूर्तीकरणे लागू असतील, तर सर्वांत आतले त्याच नावाचे अमूर्तीकरण
त्याला बांधते. आधीच्या उदाहरणात मुक्त चरे नाहीत; पण
\[(\lambd[z][yz])x
\]
यात $y$ आणि $x$ मुक्त आहेत, तर $z$ बद्ध आहे.











\section{लॅम्डा पदांचे न्यूनीकरण}\label{lam:int:red:sec}


लॅम्डा पदांचे काय करता येते? ती सोपी करता येतात. $M$ आणि $N$ ही
कोणतीही लॅम्डा पदे आणि $x$ हे कोणतेही चर असो. आदेशनानंतर $N$ मधील
मुक्त चरांशी अडथळा निर्माण करतील अशा $M$ मधील बद्ध चरांची आधी
नावे बदलून, मग $M$ मधील~$x$ च्या प्रत्येक मुक्त वापराच्या जागी $N$ आदेशल्यास जो परिणाम
मिळतो, तो आपण $\Subst{M}{N}{x}$ ने दर्शवू. उदाहरणार्थ,
\[\Subst{(\lambd[w][xxw])}{yyz}{x} = \lambd[w][(yyz)(yyz)w].
\]

\begin{digress}
आदेशनासाठी $[N/x]M$, $[x/N]M$ आणि $M[x/N]$ ही पर्यायी
संकेतनेही वापरली जातात. गोंधळ टाळा!
\end{digress}

सहजप्रज्ञेने पाहता, $(\lambd[x][M])N$ आणि $\Subst{M}{N}{x}$
यांचा अर्थ एकच आहे; पहिले पद बदलून दुसरे करण्याच्या क्रियेला
\emph{$\beta$-संकुचन} म्हणतात. $(\lambd[x][M])N$ ला
\emph{रेडेक्स}, आणि $\Subst{M}{N}{x}$ ला त्याचे
\emph{संकुचनफल} म्हणतात. सर्वसाधारणपणे, एखाद्या उपपदाचे
$\beta$-संकुचन करून $P$ हे पद~$P'$ असे बदलता येत असेल, तर
$P$ चे \emph{एका पायरीत $\beta$-न्यूनीकरण होऊन $P'$ मिळते}
असे म्हणतो आणि $P \redone P'$ लिहितो. $P$ पासून अशा
एक-पायरी न्यूनीकरणांच्या काही पायऱ्यांनी (एकही नसेल तरी)
~$P'$ मिळत असेल, तर $P$ चे~$P'$ पर्यंत
\emph{$\beta$-न्यूनीकरण होते}; संकेतांत, $P \red P'$.
ज्या पदाचे आणखी $\beta$-न्यूनीकरण करता येत नाही, त्याला
\emph{$\beta$-अन्यूनीकरणीय} किंवा \emph{$\beta$-सामान्य}
म्हणतात. संदर्भ स्पष्ट असेल तेव्हा “$\beta$-न्यूनीकरण होते”
इत्यादींऐवजी फक्त “न्यूनीकरण होते” असे म्हणू.

काही उदाहरणे पाहू.
\begin{enumerate}
\item आपल्याला मिळते:
\begin{align*}
(\lambd[x][xxy]) \lambd[z][z] & \redone (\lambd[z][z])(\lambd[z][z]) y \\
& \redone (\lambd[z][z]) y \\
& \redone y.
\end{align*}
\item पद “सोपे” केल्याने ते अधिक गुंतागुंतीचे होऊ शकते:
\begin{align*}
(\lambd[x][xxy])(\lambd[x][xxy]) & \redone (\lambd[x][xxy])(\lambd[x][xxy])y \\
& \redone (\lambd[x][xxy])(\lambd[x][xxy])yy \\
& \redone \dots
\end{align*}
\item पद तसेच राहूही शकते:
\[(\lambd[x][xx])(\lambd[x][xx]) \redone (\lambd[x][xx])(\lambd[x][xx]).
\]
\item शिवाय, काही पदांचे एकाहून अधिक प्रकारे न्यूनीकरण करता येते;
  उदाहरणार्थ,
\[(\lambd[x][(\lambd[y][yx]) z)] v \redone (\lambd[y][yv]) z
\]
हे सर्वांत बाहेरच्या उपयोजनाचे संकुचन केल्याने मिळते; आणि
\[(\lambd[x][(\lambd[y][yx]) z)] v \redone (\lambd[x][zx]) v
\]
हे सर्वांत आतल्या उपयोजनाचे संकुचन केल्याने मिळते. मात्र येथे
लक्षात घ्या की दोन्ही पदांचे पुढे न्यूनीकरण होऊन तेच पद, $zv$, मिळते.
\end{enumerate}

शेवटच्या उदाहरणात मिळणारा एकच अंतिम परिणाम हा योगायोग नाही;
तो लॅम्डा कलनाचा “चर्च--रॉसर गुणधर्म” म्हणून ओळखला जाणारा
सखोल आणि महत्त्वाचा गुणधर्म दाखवतो.











\section{चर्च--रॉसर गुणधर्म}\label{lam:int:cr:sec}

\begin{thm}
\label{lam:int:cr:thm:church-rosser}
$M$, $N_1$ आणि $N_2$ ही पदे असून $M \red N_1$ आणि $M \red
N_2$ असे समजा. मग असे $P$ पद आहे की $N_1 \red P$ आणि $N_2 \red P$.
\end{thm}

\begin{cor}
$M$ चे न्यूनीकरण सामान्य रूपापर्यंत करता येते असे समजा. मग हे
सामान्य रूप एकमेव आहे.
\end{cor}

\begin{proof}
$M \red N_1$ आणि $M \red N_2$ असेल, तर मागील प्रमेयामुळे
असे पद~$P$ आहे की $N_1$ आणि $N_2$ या दोन्हींचे न्यूनीकरण
होऊन~$P$ मिळते. $N_1$ आणि~$N_2$ ही दोन्ही सामान्य रूपात
असतील, तर हे केवळ $N_1 \ident P \ident N_2$ असेल तेव्हाच शक्य आहे.
\end{proof}

शेवटी, दोन पदे $M$ आणि~$N$ यांचे न्यूनीकरण होऊन तेच पद मिळत
असेल, तर त्यांना \emph{$\beta$-सममूल्य}, किंवा फक्त
\emph{सममूल्य}, म्हणू; म्हणजेच असे $P$ असेल की $M
\red P$ आणि $N \red P$. हे $M \equal[\beta] N$ असे लिहितात.
\ref{lam:int:cr:thm:church-rosser} वापरून $\equal[\beta]$ हा
समतुल्यता संबंध असल्याचे तपासता येते; शिवाय, कोणत्याही $M$
आणि~$N$ साठी $M \red N$ किंवा $N \red M$ असेल, तर
$M \equal[\beta] N$. (खरे तर, हा गुणधर्म असलेला
$\equal[\beta]$ हा \emph{लहानात लहान} समतुल्यता संबंध आहे,
हे दाखवता येते.)











\section{करीकरण}\label{lam:rep:cur:sec}

$\lambd[x][M]$ हे $\lambd$-अमूर्तीकरण एकच फलसाधक घेणारे
फलन दर्शवते; अनेक फलसाधक घेणारी फलने परिभाषित करायची असतील,
तर ही मोठी मर्यादा ठरते. एक मार्ग म्हणजे जोड्या, तिकड्या इत्यादी
घडवता येतील असा $\lambd$-कलनाचा विस्तार करणे. मग, उदाहरणार्थ,
त्रिस्थानी फलन $\lambd[x][M]$ ला त्याचा फलसाधक म्हणून
तिकडी अपेक्षित असेल. मात्र हे \emph{करीकरणाने} करणे अधिक सोयीचे आहे.

एक उदाहरण पाहू. क्षणभर $\lambd$-कलनात $+$ ही क्रिया आहे
असे मानू. बेरीज हे $2$-स्थानी फलन आहे, म्हणजे ते दोन फलसाधक
घेते. पण $\lambd$-अमूर्तीकरण आपल्याला एकच फलसाधक घेणारी फलने
देते: $\lambd[(x,y)][(x+y)]$ यांसारख्या व्यक्तीकरणांना विन्यासात
परवानगी नाही. तरी $\lambd[y][(x+y)]$ ने दिलेले एकस्थानी
फलन~$f_x(y)$ घेता येते; ते आपल्या एकमेव फलसाधक~$y$ मध्ये
$x$ मिळवते. खरे तर हे एक फलन नाही, तर प्रत्येक संख्या~$x$
साठी “$x$ मिळवा” अशा वेगवेगळ्या फलनांचे कुटुंब आहे. आता
$\lambd[x][f_x]$ असे आणखी एक एकस्थानी फलन~$g$ परिभाषित
करता येते. फलसाधक $x$ ला उपयोजल्यास $g(x)$ हे
फलन~$f_x$ परत करते---म्हणजे $g$ ची मूल्ये इतर फलने आहेत.
आता $g$ चे $x$ ला आणि मिळालेल्या फलनाचे~$y$ ला उपयोजन
केल्यास $(g(x))y = f_x(y) = x+y$ मिळते. अशा प्रकारे हे एकस्थानी फलन द्विस्थानी बेरीज फलनाचेच काम करू शकते.
“करीकरण” म्हणजे द्विस्थानी फलनाचे, ज्याची मूल्ये एकस्थानी
फलने आहेत अशा, एकस्थानी फलनात रूपांतर करण्याची ही युक्ती.


आता $\lambd$-कलनाच्या विन्यासात बसणारे उदाहरण पाहू.
या उदाहरणातील अमूर्तीकरणाचे दुसरे चर आदेशित पहिल्या पदात
मुक्त नसावे; आवश्यक असल्यास त्याचे आधी नामांतर करा.
$f(x,y) = x$ हे फलन कसे दर्शवायचे? दोन फलसाधक घेऊन
पहिला परत करणारे फलन परिभाषित करायचे असेल, तर
$\lambd[x][\lambd[y][x]]$ लिहिता येते. हे अक्षरशः
फलसाधक~$x$ घेऊन फलन~$\lambd[y][x]$ परत करणारे फलन आहे.
$\lambd[y][x]$ हे फलन आणखी एक फलसाधक~$y$ घेते,
पण तो बाजूला ठेवून नेहमी~$x$ परत करते. दोन फलसाधकांना
$\lambd[x][\lambd[y][x]]$ उपयोजल्यास काय होते ते पाहू:
\begin{align*}
  (\lambd[x][\lambd[y][x]])MN
  \bredone & (\lambd[y][M])N \\
  \bredone & M
\end{align*}

पुढील सर्वसाधारण मांडणीत अमूर्तीकरणातील चरे परस्पर भिन्न
आणि उपयोजल्या जाणाऱ्या सर्व पदांतील मुक्त चरांपासून वेगळी
निवडली आहेत; आवश्यक नामांतर आधी केले आहे.
सर्वसाधारणपणे, $x_1$, \dots,~$x_n$ ही प्राचले असलेले
आणि पद~$N$ ने परिभाषित होणारे फलन लिहायचे असेल,
तर $\lambd[x_1][\lambd[x_2][\ldots\lambd[x_n][N]]]$
लिहिता येते. त्याला $n$ फलसाधक उपयोजल्यास मिळते:
\begin{multline*}
  (\lambd[x_1][\lambd[x_2][\ldots\lambd[x_n][N]]]) M_1 \dots M_n \bredone\\
  \begin{aligned}
  \bredone {} & (\Subst{(\lambd[x_2][\ldots\lambd[x_n][N]])}{M_1}{x_1}) M_2
  \dots M_n\\
   \eqs {} & (\lambd[x_2][\ldots\lambd[x_n][\Subst{N}{M_1}{x_1}]]) M_2
            \dots M_n \\
  \vdots & \\
  \bredone {} &\Subst{\Subst{N}{M_1}{x_1}\ldots}{M_n}{x_n}
  \end{aligned}
\end{multline*}
शेवटच्या ओळीचा अक्षरशः अर्थ फलनाच्या व्याख्येच्या देहात
$x_i$ च्या जागी $M_i$ आदेशणे असा आहे; फलनाला अनेक
फलसाधक उपयोजल्यावर आपल्याला नेमके हेच अपेक्षित असते.










\section{लॅम्डा-परिभाष्य अंकगणितीय फलने}\label{lam:int:rep:sec}

लॅम्डा कलन संगणनाचे प्रतिमान कसे ठरू शकते? सुरुवातीला तर या
विधानाचा अर्थ काय, हेही स्पष्ट नाही. नैसर्गिक संख्यांवरील
संगणनक्षमतेची चर्चा करायची असेल, तर त्या संख्यांचे योग्य
निरूपण शोधावे लागेल. आश्चर्यकारकरीत्या चांगले काम करणारे
एक निरूपण पुढे दिले आहे.

\begin{defn}
प्रत्येक नैसर्गिक संख्या~$n$ साठी \emph{चर्च अंक}
$\num{n}$ हे लॅम्डा पद $\lambd[x][\lambd[y][(x(x(x(\dots
x(y)))))]]$ असे परिभाषित करा; अमूर्तीकरणाच्या देहात एकूण $n$ वेळा $x$ येते.
\end{defn}

$\num{n}$ ही पदे “पुनरावर्तक” आहेत: $f$ हा निविष्ट दिल्यास
$\num{n}$ हे $y$ ला $f^n(y)$ शी जोडणारे फलन परत करते.
प्रत्येक अंक सामान्य रूपात आहे, हे लक्षात घ्या. आता
एखादे लॅम्डा पद नैसर्गिक संख्यांवरील फलनाचे “संगणन करते”
याचा अर्थ सांगता येईल.

\begin{defn}
$f(x_0, \dots, x_{k-1})$ हे $\Nat$ पासून $\Nat$ मध्ये
जाणारे $k$-स्थानी आंशिक फलन असो. एखादे बंद $\lambd$-पद~$F$
हे~$f$ ला \emph{लॅम्डा-परिभाषित करते} असे म्हणतो, जर आणि फक्त जर
नैसर्गिक संख्यांच्या प्रत्येक $n_0$, \dots,~$n_{k-1}$
अनुक्रमासाठी पुढील दोन्ही अटी पूर्ण होत असतील: $f(n_0, \dots, n_{k-1})$ परिभाषित असेल तर
\[F\, \num{n_0}\, \num{n_1} \dots \num{n_{k-1}} \red \num{f(n_0, n_1, \dots,
  n_{k-1})}
\]
असे असते; आणि तसे नसेल तर $F\, \num{n_0}\, \num{n_1}
\dots \num{n_{k-1}}$ याला सामान्य रूप नसते.
\end{defn}

\begin{thm}
\label{lam:int:rep:thm:lambda-def}
फलन $f$ हे आंशिक संगणनक्षम फलन असते, जर आणि फक्त जर
एखाद्या लॅम्डा पदाने ते लॅम्डा-परिभाषित असते.
\end{thm}

\begin{explain}
हे प्रमेय काहीसे थक्क करणारे आहे. संगणनाचे प्रतिमान म्हणून
लॅम्डा कलन अगदी साधे कलन आहे; त्यात लॅम्डा अमूर्तीकरण आणि
उपयोजन एवढ्याच क्रिया आहेत!{} तरी या मोजक्या साधनांतून
कोणतीही संगणनप्रक्रिया साकार करता येते.
\end{explain}











\section{लॅम्डा-परिभाष्य फलने संगणनक्षम असतात}\label{lam:int:cmp:sec}

\begin{thm}
\label{lam:int:cmp:thm:lambda-computable}
आंशिक फलन $f$ हे एखाद्या लॅम्डा पदाने लॅम्डा-परिभाषित
असेल, तर ते संगणनक्षम आहे.
\end{thm}

\begin{proof}
फलन~$f$ हे लॅम्डा पद~$X$ ने लॅम्डा-परिभाषित आहे असे समजा.
~$f$ चे संगणन करण्याची एक अनौपचारिक पद्धत सांगू.
$m_0$, \dots,~$m_{n-1}$ ही निविष्ट मिळाल्यावर
$X \num m_0 \ldots \num m_{n-1}$ हे पद लिहा आणि आधी तेच चर्च अंक आहे का ते तपासा.
मग मूळ पदाची सर्व एका-पायरीतील न्यूनीकरणे लिहून वृक्ष तयार करा;
त्याखाली त्या प्रत्येक पदाची सर्व एका-पायरीतील न्यूनीकरणे
(म्हणजे मूळ पदाची दोन-पायरीतील न्यूनीकरणे) लिहा; आणि
प्रत्येक पातळी पूर्ण करून असेच पुढे चालू ठेवा. पदातील redex
जागा सांत असतात; प्रत्येक पायरीतील आवश्यक नामांतरासाठी
ठरलेली ताजी नावे घेतल्याने प्रत्येक पातळीही सांत राहते.
नामांतरापर्यंत कोणताही चर्च अंक मिळाला, तर संबंधित संख्या
उत्तर म्हणून परत करा. फलनाचे मूल्य अपरिभाषित असेल, तर
असा अंक मिळत नाही आणि शोध अखंड चालू राहतो.

चर्च यांच्या प्रबंधाचा आधार घेतल्यास हे फलन संगणनक्षम
असल्याचे कळते. प्रमेय सिद्ध करण्याचा अधिक चांगला मार्ग
म्हणजे या शोधपद्धतीचे पुनरावर्ती वर्णन देणे. उदाहरणार्थ,
“$\fn{IsASubterm}$,” “$\fn{Substitute}$,”
“$\fn{ReducesToInOneStep}$,” “$\fn{ReductionSequence}$,”
“$\fn{Numeral}$” इत्यादी आदिम पुनरावर्ती फलने आणि
संबंधांची मालिका परिभाषित करता येईल. मग
$f(m_0, \dots, m_{n-1})$ चे संगणन करणारी आंशिक
पुनरावर्ती पद्धत म्हणजे $X \num{m_0} \dots \num{m_{n-1}}$
पासून सुरू होऊन अंकापाशी संपणारा एका-पायरीतील
न्यूनीकरणांचा अनुक्रम शोधणे आणि त्या अंकाशी संबंधित
संख्या परत करणे. तपशील मोठे आणि कंटाळवाणे आहेत,
पण इतरथा नेहमीचे आहेत.
\end{proof}











\section{संगणनक्षम फलने लॅम्डा-परिभाष्य असतात}\label{lam:int:lrp:sec}

\begin{thm}
\label{lam:int:lrp:thm:computable-lambda}
प्रत्येक संगणनक्षम आंशिक फलन लॅम्डा-परिभाष्य असते.
\end{thm}

\begin{proof}

प्रत्येक आंशिक संगणनक्षम फलन~$f$ साठी बंद प्रतिनिधी~$F$
शोधायचा आहे. क्लिनीच्या प्रमाण रूप प्रमेयानुसार, total आदिम
पुनरावर्ती अंकफलने दर्शवणे, total आदिम पुनरावर्ती zero-test
चे अपरिबद्ध लघुतमीकरण करणे, आणि त्यानंतर योग्य output
काढणे पुरेसे आहे. Total आदिम पुनरावर्ती फलनांसाठी आरंभीची
फलने, total संयोजन आणि total आदिम पुनरावर्तन यांच्या पुढील
constructions वापरू. Search चे मूल्य अपरिभाषित असताना
output extractor ते टाकून देऊ नये म्हणून शेवटचे उपयोजन
ते search आधी force करेल. ही शेवटची construction
लघुतमीकरणाच्या विभागात देऊन या प्रमेयाची सिद्धता पूर्ण करू.
सर्व आंशिक फलनांसाठी सरळ संयोजन किंवा सरळ recursion
योग्य असल्याचे आधीच गृहीत धरण्याची गरज नाही.
\end{proof}

उर्वरित पुरावा अधिक वाचनीय व्हावा म्हणून आपण अधिक प्रचलित
संकेतन वापरू. उदाहरणार्थ, $Mxyz$ ऐवजी $M(x, y, z)$ लिहू.
यातून काही अर्थ सूचित होत असला, तरी प्रकारविरहित लॅम्डा
कलनातील पदांना निश्चित स्थानसंख्या नसते, हे लक्षात ठेवा;
म्हणून त्याच पदासाठी~$M$ आपण $M(x,y)$ आणि
$M(x,y,z,w)$ असे दोन्ही लिहू शकतो. पण या संकेतनाने
आपण $M$ ला तीन चरांचे फलन मानत आहोत, हे सूचित होते
आणि पुढील व्याख्यांमागील हेतू अधिक स्पष्ट होतो.
तसेच “$M$ ची व्याख्या $M =
\lambd[x][\lambd[y][\lambd[z][\dots]]]$ ने करा”
याऐवजी “$M$ ची व्याख्या $M(x,y,z) = \dots$ ने करा”
असे म्हणू.











\section{मूलभूत आदिम पुनरावर्ती फलने लॅम्डा-परिभाष्य असतात}\label{lam:int:bas:sec}

\begin{lem}
$\Zero$, $\Succ$ आणि $\Proj{n}{i}$ ही फलने
लॅम्डा-परिभाष्य आहेत.
\end{lem}

\begin{proof}

येथे $\Zero$ हे आधी परिभाषित केलेले एकस्थानी शून्य फलन आहे.
त्याला दर्शवणारे पद
$\lambd[w][\lambd[x][\lambd[y][y]]]$
असे घ्या. बाहेरचे अमूर्तीकरण निविष्ट बाजूला ठेवते आणि
चर्चचा शून्य अंक परत करते.

उत्तरवर्ती फलन~$\Succ$ ची
$\fn{Succ}(u) = \lambd[x][\lambd[y][x(uxy)]]$ अशी
व्याख्या करतात. हे कसे चालते याचा विचार करा: पुनरावर्तक
मानलेल्या प्रत्येक अंकासाठी $\num{n}$ आणि प्रत्येक
फलनासाठी~$f$, $\fn{Succ}(\num{n},f)$ हे असे फलन आहे
की निविष्ट~$y$ वर ते~$y$ पासून सुरू करून $f$ चे $n$
वेळा उपयोजन करते आणि त्यानंतर आणखी एकदा उपयोजन करते.

प्रक्षेपणांबद्दल अधिक काही सांगण्यासारखे नाही:
$\fn{Proj}^n_i(x_0, \dots,
x_{n-1}) = x_i$. दुसऱ्या शब्दांत, आपल्या रूढींनुसार,
$\fn{Proj}^n_i$ हे लॅम्डा पद
$\lambd[x_0][\dots \lambd[x_{n-1}][x_i]]$ आहे.
\end{proof}











\section{लॅम्डा-परिभाष्य फलने संयोजनाच्या अंतर्गत बंद असतात}\label{lam:int:com:sec}

\begin{lem}
लॅम्डा-परिभाष्य फलनांचा वर्ग संयोजनाच्या अंतर्गत बंद असतो.
\end{lem}

\begin{proof}

प्रथम या सरळ मांडणीत सर्व घटक फलने सर्वत्र परिभाषित आहेत
असे घ्या. $f$ हे $h$, $g_0,$ \dots,~$g_{k-1}$ यांच्या संयोजनाने
परिभाषित आहे असे समजा. $h$, $g_0$, \dots,~$g_{k-1}$
ही फलने अनुक्रमे $H$, $G_0$, \dots,~$G_{k-1}$
यांनी लॅम्डा-परिभाषित आहेत असे गृहीत धरल्यास,
~$f$ ला लॅम्डा-परिभाषित करते असे पद~$F$ शोधायचे आहे.
पण~$F$ ची व्याख्या सरळ अशी करता येते:
\[F(x_0, \dots, x_{l-1}) = H(G_0(x_0, \dots, x_{l-1}),
\dots, G_{k-1}(x_0, \dots, x_{l-1})).
\]
प्रत्येक घटक निविष्ट अंकांवर सामान्य रूपातील अंक देतो;
ते सर्व न्यूनीकरण करून मग बाह्य फलनाचे न्यूनीकरण केल्याने
संयोजनाचा योग्य अंक मिळतो. त्यामुळे total फलनांसाठी ही
मांडणी योग्य आहे.
आंशिक फलनांसाठी हा सरळ पुरावा पुरेसा नाही. बाह्य फलन
निविष्ट वापरत नसेल, तर ते अपरिभाषित अंतर्गत संगणनही
टाकून देऊ शकते. म्हणून सामान्य strict partial closure साठी
पुढे पूर्ण होणारे \ref{lam:int:lrp:thm:computable-lambda} वापरतो:
\ref{lam:int:cmp:thm:lambda-computable} ने सर्व घटक आंशिक संगणनक्षम
आहेत; त्यांचे strict संयोजन आंशिक संगणनक्षम आहे; आणि
पूर्ण झालेले पहिले प्रमेय त्याला लॅम्डा प्रतिनिधी देते.
त्या प्रमेयाची पुढील सिद्धता येथे फक्त total closure वापरते,
म्हणून यात चक्रीय अवलंबन नाही.
\end{proof}











\section{लॅम्डा-परिभाष्य फलने आदिम पुनरावर्तनाच्या अंतर्गत बंद असतात}\label{lam:int:pr:sec}


आदिम पुनरावर्तनासाठी आपण टप्प्याटप्प्याने construction करू.
या मुख्य सिद्धतेत $g$ आणि~$h$ ही सर्वत्र परिभाषित फलने असून
त्यांना अनुक्रमे लॅम्डा-परिभाषित करणारी बंद पदे $G'$ आणि $H'$
आधीच दिलेली आहेत. पुढील फलन~$f$ ला दर्शवणारे पद $F'$ हवे:
\begin{align*}
f(0, \vec z) & = g(\vec z) \\
f(x+1, \vec z) & = h(x, f(x,\vec z), \vec z).
\end{align*}
प्रत्येक नैसर्गिक संख्या~$n$ साठी पुढील अटी पुरेशा आहेत:
\begin{align*}
F'(\num{0}, \vec z) & \equiv G'(\vec z) \\
F'(\overline{n+1}, \vec z) & \equiv H'(\num{n}, F'(\num{n}, \vec z), \vec z).
\end{align*}
कारण निविष्ट प्राचलांच्या जागी अंक घातल्यावर विगमनाने
प्रत्येक पायरीवर योग्य सामान्य रूपातील अंक मिळतो.

प्राचले फलनमूल्यात सामावण्यासाठी पुढील अटी पूर्ण करणारे
पद~$F$ शोधणे पुरेसे आहे:
\begin{align*}
F(\num 0) & \equiv G \\
F(\overline {n+1}) & \equiv H(\num{n},F(\num{n}))
\intertext{प्रत्येक नैसर्गिक संख्या $n$ साठी; येथे}
G & = \lambd[\vec z][G'(\vec z)] \text{ आणि}\\
H(u,v) & = \lambd[\vec z][H'(u,v(\vec z),\vec z)].
\end{align*}
आता $F'(u,\vec z)=F(u)(\vec z)$ असे घेतल्यास वरची
प्राचलांसहित recurrence मिळते. येथे मागील value function
आधीच योग्य recursion index साठी घेतलेले आहे; त्याला
फक्त उरलेली प्राचले उपयोजायची आहेत.

पद $F$ परिभाषित करण्यापूर्वी क्रमित जोड्या हाताळण्याची
पद्धत हवी. पुढील उपपत्ती ती देते.

\begin{lem}
असे लॅम्डा पद $D$ आहे की लॅम्डा पदांच्या प्रत्येक जोडी
$M$ आणि~$N$ साठी $D(M,N)(\num{0}) \red M$ आणि
$D(M,N)(\num{1}) \red N$.
\end{lem}

\begin{proof}
आधी लॅम्डा पद $K$ ची पुढीलप्रमाणे व्याख्या करा:
\[K(y) = \lambd[x][y].
\]
म्हणजे $K$ हे पद $\lambd[y][\lambd[x][y]]$ आहे.
दुसऱ्या प्रकारे पाहिले, तर प्रत्येक $M$ साठी $K(M)$
हे कोणतीही निविष्ट घेतल्यावर~$M$ परत करणारे स्थिर फलन आहे.

आता $D(x,y,z)$ ची व्याख्या $D(x,y,z) = z (K(y))x$
अशी करा. मग आपल्याला मिळते:
\begin{align*}
D(M,N,\num 0) & \red \num 0 (K(N)) M \red M \text{ आणि}\\
D(M,N,\num 1) & \red \num 1 (K(N)) M \red K(N) M \red N,
\end{align*}
हेच अपेक्षित होते.
\end{proof}

येथे कल्पना अशी आहे की $D(M,N)$ ही $\tuple{M, N}$
जोडी दर्शवते. $P$ अशी जोडी दर्शवते असे गृहीत धरल्यास,
$P(\num 0)$ आणि $P(\num 1)$ हे अनुक्रमे डावे आणि उजवे
प्रक्षेपण, $(P)_0$ आणि $(P)_1$, दर्शवतात. पुढे आपण
हीच संकेतने वापरू.

\begin{lem}
लॅम्डा-परिभाष्य फलनांचा वर्ग आदिम पुनरावर्तनाच्या
अंतर्गत बंद असतो.
\end{lem}

\begin{proof}
कोणतीही पदे $G$ आणि $H$ दिली असता प्रत्येक नैसर्गिक
संख्या~$n$ साठी पुढील अटी पूर्ण करणारे पद $F$
शोधता येते, हे दाखवायचे आहे:
\begin{align*}
F(\num{0}) & \equiv G \\
F(\num{n+1}) & \equiv H(\num{n}, F(\num{n}))
\end{align*}
साधारण कल्पना म्हणजे अंकांचा पुनरावर्तक म्हणून वापर करून
\emph{जोड्यांचा} अनुक्रम संगणित करणे:
\[\tuple{\num 0, F(\num 0)}, \tuple{\num 1, F(\num 1)}, \dots,
\]
पहिली जोडी म्हणजे $\tuple{\num 0, G}$, हे लक्षात घ्या.
$\tuple{\num{n}, F(\num{n})}$ ही जोडी दिल्यास पुढची जोडी
$\tuple{\num{n+1}, F(\num{n+1})}$ ही
$\tuple{\num{n+1}, H(\num{n}, F(\num{n}))}$ शी
सममूल्य असायला हवी. हे एका पायरीतील संक्रमण करणारे
लॅम्डा पद~$T$ आपण तयार करू.

तपशील असे आहेत. येथे S हे आधी तयार केलेले उत्तरवर्ती
फलन दर्शवते. $T(u)$ ची व्याख्या पुढीलप्रमाणे करा:
\[T(u) = \tuple{S((u)_0), H((u)_0,(u)_1)}.
\]
मग कोणत्याही संख्या~$n$ साठी पुढील विधान सहज तपासता येते:
\[T(\tuple{\num{n}, M}) \red \tuple{\num{n+1}, H(\num{n}, M)}.
\]
वर सुचवल्याप्रमाणे, $G$ आणि $H$ दिली असता
~$F(u)$ ची व्याख्या पुढीलप्रमाणे करा:
\[F(u) = (u(T,\tuple{\num 0, G}))_1.
\]
म्हणजे निविष्ट~$\num{n}$ वर $F$ हे
$\tuple{\num 0, G}$ पासून सुरू करून $T$ चे $n$
वेळा पुनरावर्तन करते आणि मग दुसरा घटक परत करते.
सुरुवातीला आपल्याला मिळते:
\begin{enumerate}
\item $\num{0} (T, \tuple{\num 0, G}) \equiv \tuple{\num{0}, G}$
\item $F(\num{0}) \equiv G$
\end{enumerate}
~$n$ वरील विगमनाने प्रत्येक नैसर्गिक संख्येसाठी
पुढील गोष्टी दाखवता येतात:
\begin{enumerate}
\item $\num{n+1}(T, \tuple{\num{0}, G}) \equiv \tuple{\num{n+1}, F(\num{n+1})}$
\item $F(\num{n+1}) \equiv H(\num{n}, F(\num{n}))$
\end{enumerate}
दुसऱ्या विधानासाठी,
\begin{align*}
F(\num{n+1}) & \red (\num{n+1}(T, \tuple{\num 0, G}))_1 \\
& \equiv (T(\num{n} (T, \tuple{\num 0, G})))_1 \\
& \equiv (T(\tuple{\num{n}, F(\num{n})}))_1 \\
& \equiv (\tuple{\num{n+1}, H(\num{n}, F(\num{n}))})_1 \\
& \equiv H(\num{n}, F(\num{n})).
\end{align*}
येथे तिसऱ्या ओळीत विगमन गृहीतक वापरले आहे.
पहिल्या विधानासाठी,
\begin{align*}
\num{n+1} (T, \tuple{\num 0, G}) &
\equiv T(\num{n} (T, \tuple{\num 0, G})) \\
& \equiv T( \tuple{\num{n}, F(\num{n})}) \\
& \equiv \tuple{\num{n+1}, H(\num{n}, F(\num{n}))} \\
& \equiv \tuple{\num{n+1}, F(\num{n+1})}.
\end{align*}
येथे शेवटच्या ओळीत दुसरे विधान वापरले आहे.
अशा प्रकारे $F(\num 0) \equiv G$ आणि प्रत्येक $n$
साठी $F(\num {n+1}) \equiv H(\num
n, F(\num{n}))$ दाखवले; आपल्याला नेमके हेच हवे होते.
वरील beta-recursion equations कोणत्याही दिलेल्या पदांसाठी
योग्य आहेत. Numeric फलनाच्या प्रतिनिधित्वाचा येथे वापर
total मूळ व step फलनांसाठी केला आहे. सामान्य आंशिक
strict पुनरावर्तनाचा निष्कर्ष \ref{lam:int:cmp:thm:lambda-computable}
आणि पुढे स्वतंत्रपणे पूर्ण होणाऱ्या
\ref{lam:int:lrp:thm:computable-lambda} मधून मिळतो: आंशिक
संगणनक्षम फलने strict आदिम पुनरावर्तनाच्या अंतर्गत बंद आहेत.
Step फलन आधीचे अपरिभाषित value टाकून देऊ शकते,
म्हणून arbitrary पदांसाठीची equations स्वतःहून तो
आंशिक निष्कर्ष सिद्ध करतात असा दावा नाही.
\end{proof}











\section{स्थिरबिंदू संयोजक}\label{lam:int:fix:sec}

तुमच्याकडे लॅम्डा पद $g$ आहे आणि दुसरे पद $k$ हवे
आहे असे समजा; त्यात $k$ हे $gk$ शी
$\beta$-सममूल्य असावे.

आपल्या संकेतनरूढींनुसार पुढील पदे परिभाषित करा;
येथील अमूर्त केलेले चर दिलेल्या पदात मुक्त नसावे.
आवश्यक असल्यास त्याचे आधी नामांतर करा:
\[\fn{diag}(x) = xx
\]
आणि
\[l(x) = g(\fn{diag}(x))
\]
म्हणजे $l$ हे पद $\lambd[x][g(xx)]$ आहे.
$k$ हे पद $ll$ असू द्या. मग आपल्याला मिळते:
\begin{align*}
k & = (\lambd[x][g(xx)])(\lambd[x][g(xx)]) \\
& \red  g((\lambd[x][g(xx)])(\lambd[x][g(xx)])) \\
& = gk.
\end{align*}
आता
\[Y = \lambd[g][((\lambd[x][g(xx)])(\lambd[x][g(xx)]))]
\]
असे घेतल्यास $Yg$ आणि $g(Yg)$ यांचे न्यूनीकरण
होऊन एकच पद मिळते; म्हणून $Yg \equiv_\beta
g(Yg)$. याला “करीचा संयोजक” म्हणतात. त्याऐवजी
\[Y = (\lambd[xg][g(xxg)])(\lambd[xg][g(xxg)])
\]
असे घेतल्यास प्रत्यक्षात $Yg$ चे $g(Yg)$ पर्यंत
न्यूनीकरण होते; हे अधिक प्रबळ विधान आहे.
$Y$ च्या या दुसऱ्या रूपाला “ट्यूरिंगचा संयोजक”
म्हणतात.











\section{लॅम्डा-परिभाष्य फलने लघुतमीकरणाच्या अंतर्गत बंद असतात}\label{lam:int:min:sec}

\begin{lem}
$f(x,y)$ हे लॅम्डा-परिभाष्य आहे असे समजा.
$g$ ची पुढीलप्रमाणे व्याख्या करा:
\[g(x) \simeq \umin{y}{f(x,y)}.
\]
मग $g$ हे लॅम्डा-परिभाष्य आहे.
\end{lem}

\begin{proof}

प्रथम $f$ सर्वत्र परिभाषित आहे असे घ्या. मुख्य कल्पना साधारण अशी आहे. $x$ दिले असता,
~$n$ पासून सुरू करून~$y$ शोधणारे फलन $h_x(n)$
परिभाषित करण्यासाठी स्थिरबिंदू लॅम्डा पद $Y$ वापरू;
मग $g(x)$ म्हणजे $h_x(0)$. फलन~$h_x$ हे
स्थिरबिंदू समीकरणाचे उत्तर म्हणून मांडता येते:
\[h_x(n) \simeq
\begin{cases}
n & \text{जर $f(x,n) = 0$ असेल} \\
h_x(n+1) & \text{जर $f(x,n)>0$ असेल.}
\end{cases}
\]

या total प्रकरणाचे तपशील असे आहेत. $f$ हे लॅम्डा-परिभाष्य असल्याने
कोणते तरी बंद पद $F$ त्याला लॅम्डा-परिभाषित करते.
$D(M, N, \bar{0}) \red M$ आणि $D(M, N, \bar{1})
\red N$ असे असणारे लॅम्डा पद~$D$ आपल्याकडे आहे,
हे लक्षात ठेवा. क्षणभर $x$ स्थिर धरू. $h_x$ ला
लॅम्डा-परिभाषित करण्यासाठी~$x$ वर अवलंबून असणारे
आणि पुढील अट पूर्ण करणारे पद~$H$ शोधायचे आहे:
\[H(\num{n}) \equiv D(\num{n}, H(S(\num{n})), F(x, \num{n})).
\]
स्थिरबिंदू पद~$Y$ वापरून हे करता येते. प्रथम $U$
हे पुढील पद असू द्या:
\[\lambd[h][\lambd[z][D(z,(h(Sz)),F(x,z))]],
\]
आणि मग $H$ हे पद~$YU$ असू द्या. $H$ मधील
एकमेव मुक्त चर~$x$ आहे, हे लक्षात घ्या. वरचे
समीकरण $H$ पूर्ण करते, हे दाखवू.

$Y$ च्या व्याख्येनुसार,
\[H = YU \equiv U(YU) = U(H).
\]
विशेषतः, प्रत्येक नैसर्गिक संख्या~$n$ साठी,
\begin{align*}
H(\num{n}) & \equiv U(H, \num{n}) \\
& \red D(\num{n}, H(S(\num{n})), F(x, \num{n})),
\end{align*}
अपेक्षेप्रमाणे. शेवटच्या ओळीत~$x$ च्या जागी
$\num{m}$ हा अंक आदेशल्यास,
$F(\num{m}, \num{n})$ चे न्यूनीकरण होऊन
$\num{0}$ मिळते तेव्हा हे व्यक्तीकरण $\num{n}$
पर्यंत न्यूनीत होते; आणि $F(\num{m}, \num{n})$
चे न्यूनीकरण होऊन इतर कोणताही अंक मिळतो तेव्हा
ते $H(S(\num{n}))$ पर्यंत न्यूनीत होते.

पुरावा पूर्ण करण्यासाठी $G$ हे
$\lambd[x][H(\num 0)]$ असू द्या. मग $G$
हे~$g$ ला लॅम्डा-परिभाषित करते; म्हणजेच प्रत्येक~$m$
साठी $g(m)$ परिभाषित असेल तर $G(\num m)$ चे
न्यूनीकरण होऊन~$\overline {g(m)}$ मिळते;
Zero witness मिळाल्यास आधीच्या सांत positive tests नंतर
तोच अंक मिळतो. Witness नसल्यास प्रत्येक total test
positive असल्याने search पुढच्या candidate ला जात राहते.
सर्वांत डाव्या redex चे न्यूनीकरण या बाह्य search ला अखंड
उलगडते; बाह्य function position मधील search संपत नाही.
त्याला सामान्य रूप असते, तर normalizing leftmost strategy
ते सांत पायऱ्यांत शोधली असती. यासाठी
\href{https://www.cs.ox.ac.uk/andrew.ker/docs/lambdacalculus-lecture-notes-ht2009.pdf}{Andrew D. Ker, \emph{Lambda Calculus and Types}, प्रमेय 3.2.3,
मुद्रित पृष्ठ 35} मधील standardization परिणाम वापरला आहे.
या search ला बाहेरून अधिक पदे उपयोजली तरी ते बाह्य
function position मध्येच राहते; अनंत search टाकून दिले जात नाही.
आता आधीच्या \ref{lam:int:lrp:thm:computable-lambda} ची सिद्धता
पूर्ण करू. क्लिनी प्रमाण रूपानुसार
\[ f(\vec x)\simeq e(\umin{n}{r(\vec x,n)}),
\]
येथे $e$ आणि $r$ सर्वत्र परिभाषित आदिम पुनरावर्ती फलने
आहेत; zero test मान्य halting record ओळखते. त्यांची बंद
प्रतिनिधी पदे $E$ आणि $R$ आधीच्या total constructions ने
मिळतात. वरील search construction मध्ये अनेक निविष्ट प्राचले
तशीच राखून त्याच्या लघुतमीकरणाचा प्रतिनिधी $W$ मिळतो.
मग $I=\lambd[t][t]$ घेऊन
\[ F(\vec x)=(\lambd[v][v\,I\,(E v)])\,W(\vec x)
\]
असे घ्या; आवश्यक बद्ध चरे ताजी निवडली आहेत. Search
ने $s$ हा अंक दिल्यास
\[ (\lambd[v][v\,I\,(E v)])\,\num{s}
 \red \num{s}\,I\,(E\,\num{s})
 \red E\,\num{s}.
\]
कारण ओळख फलनाचे कितीही पुनरावर्तन केले तरी value बदलत
नाही. शेवटचे पद योग्य output अंक देते. Search अखंड चालत
असेल, तर ते बाह्य function position मध्ये force होत राहते;
ते टाकून देणाऱ्या extractor पर्यंत पोहोचत नाही आणि सामान्य
रूप नसते. याने \ref{lam:int:lrp:thm:computable-lambda} स्वतंत्रपणे
सिद्ध झाले.
शेवटी प्रारंभीच्या उपपत्तीत $f$ आंशिक असण्याची सामान्य
शक्यता घ्या. \ref{lam:int:cmp:thm:lambda-computable} ने ते आंशिक
संगणनक्षम आहे. Strict अपरिबद्ध लघुतमीकरणातील प्रत्येक
पूर्वीचे value परिभाषित व positive असावे लागते; हा शोध
आंशिक संगणनक्षम आहे. पूर्ण झालेल्या
\ref{lam:int:lrp:thm:computable-lambda} ने त्याला बंद लॅम्डा
प्रतिनिधी मिळतो. Arbitrary आंशिक प्रतिनिधीला थेट वरील
selector मध्ये बसवणे पुरेसे आहे असे गृहीत घेतलेले नाही.
\end{proof}



\chapter{विन्यासमीमांसा}\label{lam:syn:chap}








\section{पदे}\label{lam:syn:trm:sec}

लॅम्डा कलनातील पदे $\Obj{v_0}$, $\Obj{v_1}$, \dots या
अमर्याद चरांपासून, ~“$\lambd$” या चिन्हापासून आणि कंसांपासून
विगमनाधारित रीतीने घडवली जातात. चरे दर्शवण्यासाठी आपण
$x$, $y$, $z$, \dots{} आणि पदे दर्शवण्यासाठी
$M$, $N$, $P$, \dots{} वापरू.

\begin{defn}[पदे] \label{lam:syn:trm:defn:term}
लॅम्डा कलनातील \emph{पदांचा} संच पुढीलप्रमाणे
विगमनाधारित रीतीने परिभाषित करतात:
\begin{enumerate}
  \item \label{lam:syn:trm:defn:term-var} $x$ हे चर असेल, तर $x$ हे
    पद आहे.
  \item \label{lam:syn:trm:defn:term-abs} $x$ हे चर आणि $M$ हे
    पद असेल, तर $(\lambd[x][M])$ हे पद आहे.
  \item \label{lam:syn:trm:defn:term-app} $M$ आणि $N$ ही दोन्ही
    पदे असतील, तर $(MN)$ हे पद आहे.
\end{enumerate}
\end{defn}

पद $(\lambd[x][M])$ हे \ref{lam:syn:trm:defn:term-abs} नुसार
घडवले असेल, तर ते \emph{अमूर्तीकरणाचे} फलित आहे
असे म्हणतो; आणि $\lambd[x]$ मधील $x$ ला
\emph{प्राचल} म्हणतात. \ref{lam:syn:trm:defn:term-app}
नुसार घडवलेले पद $(MN)$ हे \emph{उपयोजनाचे}
फलित असते.

वर परिभाषित पदांमध्ये सर्व कंस स्पष्टपणे लिहिलेले
असतात. उदाहरणार्थ,
$(\lambd[x][((\lambd[x][x])(\lambd[x][(xx)]))])$
या पदावरून दिसते की हे खूपच अवजड होऊ शकते. कंस
टाळण्यासाठी आपण काही रूढी मांडू. मात्र अधिकृत
व्याख्येमुळे \ref{lam:syn:trm:defn:term} नुसार पद कसे
घडले हे ठरवणे सोपे जाते. उदाहरणार्थ,
$(\lambd[x][((\lambd[x][x])(\lambd[x][(xx)]))])$
हे पद घडवण्याची शेवटची पायरी अमूर्तीकरणच असली
पाहिजे; त्यात प्राचल हे~$x$ आहे.
ते $((\lambd[x][x])(\lambd[x][(xx)]))$ या पदावर
अमूर्तीकरण करून मिळते; हे पद दोन पदांच्या
उपयोजनातून घडले आहे. त्या दोन पदांपैकी प्रत्येक
अमूर्तीकरणातून घडले आहे; आणि पुढे असेच चालते.

\begin{prob}
$(\lambd[g][(\lambd[x][(g (x x))])
  (\lambd[x][(g (x x))])])$ हे पद कसे घडले ते वर्णन करा.
\end{prob}











\section{एकमेव वाचनीयता}\label{lam:syn:unq:sec}

प्रत्येक पद घडवण्याचा एकमेव मार्ग असतो का, असा प्रश्न
पडू शकतो; आणि तसेच आहे. प्रत्येक लॅम्डा पद घडवण्याची
आणि त्याचा अर्थ लावण्याची एकच पद्धत आहे. पुढील चर्चेत
\emph{रचनाप्रक्रिया} म्हणजे \ref{lam:syn:trm:defn:term}
मधील रचनानियम एकदा किंवा अनेकदा वापरून पद घडवण्याची
प्रक्रिया.

\begin{lem}\label{lam:syn:unq:lem:term-start}
पदाची सुरुवात एकतर चराने किंवा कंसाने होते.
\end{lem}

\begin{proof}
\ref{lam:syn:trm:defn:term} नुसार घडवले असेल तरच एखादी
गोष्ट पद ठरते. ते \ref{lam:syn:trm:defn:term-var} चे फलित
असेल, तर ते चर असलेच पाहिजे. ते
\ref{lam:syn:trm:defn:term-abs} किंवा
\ref{lam:syn:trm:defn:term-app} चे फलित असेल, तर
त्याची सुरुवात कंसाने होते.
\end{proof}

\begin{lem}\label{lam:syn:unq:lem:app-start}
उपयोजनाच्या फलिताची सुरुवात एकतर दोन कंसांनी किंवा
एक कंस आणि त्यानंतर चर याने होते.
\end{lem}

\begin{proof}
$M$ हे उपयोजनाचे फलित असेल, तर ते $(PQ)$ या रूपात
असते; म्हणून त्याची सुरुवात कंसाने होते. $P$ हे पद
असल्यामुळे \ref{lam:syn:unq:lem:term-start} नुसार त्याची
सुरुवात कंसाने किंवा चराने होते.
\end{proof}

\begin{lem}\label{lam:syn:unq:lem:initial}
पदाचा कोणताही उचित आरंभीचा भाग स्वतः पद नसतो.
\end{lem}

\begin{prob}
पदांच्या लांबीवरील विगमनाने
\ref{lam:syn:unq:lem:initial} सिद्ध करा.
\end{prob}

\begin{prop}[एकमेव वाचनीयता] \label{lam:syn:unq:prop:unq}
प्रत्येक पदाची रचनाप्रक्रिया एकमेव असते. दुसऱ्या
शब्दांत, एखाद्या रचनाप्रक्रियेने पद $M$ घडत असेल,
तर तेच पद घडवू शकणारी ती एकमेव रचनाप्रक्रिया असते.
\end{prop}

\begin{proof}
पदांच्या रचनाप्रक्रियेवरील विगमनाने हे सिद्ध करू.
  \begin{enumerate}
    \item $M$ हे $x$ या रूपात आहे, जिथे $x$ हे
      एखादे चर आहे. अमूर्तीकरण आणि उपयोजन यांची
      फलिते नेहमी कंसाने सुरू होतात, म्हणून $M$
      घडवताना त्यांचा वापर झालेला असू शकत नाही.
      त्यामुळे $M$ ची रचनाप्रक्रिया
      \ref{lam:syn:trm:defn:term}\ref{lam:syn:trm:defn:term-var}
      ची एकच पायरी असली पाहिजे.
    \item $M$ हे $(\lambd[x][N])$ या रूपात आहे,
      जिथे $x$ हे चर आणि $N$ हे पद आहे. ते
      \ref{lam:syn:trm:defn:term}\ref{lam:syn:trm:defn:term-var}
      नुसार घडलेले असू शकत नाही, कारण ते एकटे
      चर नाही. \ref{lam:syn:unq:lem:app-start} नुसार ते
      उपयोजनाचेही फलित नाही. त्यामुळे $M$ हे
      फक्त~$N$ वरील अमूर्तीकरणाचे फलित असू शकते.
      विगमन गृहीतकानुसार $N$ ची रचनाप्रक्रियाही
      एकमेव आहे.
    \item $M$ हे $(PQ)$ या रूपात आहे, जिथे $P$
      आणि $Q$ ही पदे आहेत. त्याची सुरुवात कंसाने
      होत असल्याने ते
      \ref{lam:syn:trm:defn:term}\ref{lam:syn:trm:defn:term-var}
      नुसारही घडलेले असू शकत नाही.
      \ref{lam:syn:unq:lem:term-start} नुसार $P$ ची सुरुवात
      $\lambd$ ने होऊ शकत नाही; म्हणून $(PQ)$ हे
      अमूर्तीकरणाचे फलितही असू शकत नाही. आता $M$
      उपयोजनाने घडवण्याचा आणखी एक मार्ग आहे असे
      समजा; उदाहरणार्थ, ते $(P'Q')$ या रूपातही
      आहे. सुरुवातीची पदे वेगळी असतील, तर $P$
      हा $P'$ चा उचित आरंभीचा भाग असेल किंवा
      उलट असेल; \ref{lam:syn:unq:lem:initial} नुसार हे
      अशक्य आहे. ती एकच असतील, तर उरलेले भागही
      एकच असतात. म्हणून $P$ आणि $Q$ एकमेव
      रीतीने ठरतात आणि विगमन गृहीतकानुसार $P$
      आणि $Q$ यांच्या रचनाप्रक्रियाही एकमेव असतात.
  \end{enumerate}
\end{proof}

वरील प्रतिज्ञेचा अधिक वाचनीय अर्थ असा आहे:
\begin{prop}
  पद $M$ पुढीलपैकी एकाच रूपात असू शकते:
  \begin{enumerate}
    \item $x$, जिथे $x$ हे $M$ ने एकमेव रीतीने
      ठरणारे चर आहे.
    \item $(\lambd[x][N])$, जिथे $x$ हे चर आणि
      $N$ हे दुसरे पद आहे; दोन्ही $M$ ने एकमेव
      रीतीने ठरतात.
    \item $(PQ)$, जिथे $P$ आणि $Q$ ही दोन पदे
      $M$ ने एकमेव रीतीने ठरतात.
  \end{enumerate}
\end{prop}











\section{संक्षिप्त विन्यास}\label{lam:syn:abb:sec}

\ref{lam:syn:trm:defn:term} मध्ये परिभाषित पदे लिहिताना
कधीकधी अवजड वाटतात; म्हणून अधिक संक्षिप्त विन्यास
मांडणे उपयुक्त आहे. अर्थात, या संक्षिप्त संकेतनातील
पदेही एकमेव रीतीने वाचता येतील याची काळजी
घ्यावी लागेल. \emph{संक्षिप्त पदे} म्हणून ओळखली
जाणारी एक प्रचलित पद्धत पुढीलप्रमाणे आहे.

\begin{enumerate}
\item कंस वगळले असता उपयोजन डावीकडून उजवीकडे
  होते. उदाहरणार्थ, $M$, $N$, $P$ आणि~$Q$
  ही पदे असतील, तर $MNPQ$ हे $(((MN)P)Q)$
  चे संक्षिप्त रूप आहे.
\item पुन्हा, कंस वगळले असता लॅम्डा अमूर्तीकरणाला
  शक्य तितकी व्यापक व्याप्ती द्यायची. उदाहरणार्थ,
  $\lambd[x][MNP]$ याचा अर्थ
  $(\lambd[x][((MN)P)])$ असा घ्यायचा.
\item एका लॅम्डाने अनेक चरे अमूर्त करता येतात.
  उदाहरणार्थ, $\lambd[xyz][M]$ हे
  $\lambd[x][\lambd[y][\lambd[z][M]]]$
  चे संक्षिप्त रूप आहे.
\end{enumerate}

उदाहरणार्थ,
\[\lambd[xy][xxyx \lambd[z][xz]]
\]
हे पुढीलचे संक्षिप्त रूप आहे:
\[(\lambd[x][(\lambd[y][((((xx)y)x)(\lambd[z][(xz)]))])]).
\]

\begin{prob}
संक्षिप्त पद
$\lambd[g][(\lambd[x][g (x x)])
  \lambd[x][g (x x)]]$
विस्तारून लिहा.
\end{prob}











\section{मुक्त चरे}\label{lam:syn:fv:sec}

लॅम्डा कलनात फलनांचा अभ्यास केला जातो, आणि लॅम्डा अमूर्तीकरणातून
फलने मिळतात. सहज समजुतीने, $\lambd[x][M]$ हे असे फलन आहे की
फलनाचा फलसाधक~$x$ ला दिल्यावर त्याची मूल्ये~$M$ ने
ठरतात. पण~$M$ मधील $x$ ची प्रत्येक आस्थिती यासाठी संबंधित नसते:
$M$ मध्ये आणखी एखादे अमूर्तीकरण $\lambd[x][N]$ असेल, तर~$N$
मधील $x$ च्या आस्थिती $\lambd[x][N]$ शी संबंधित असतात;
$\lambd[x][M]$ शी नाही. म्हणून $\lambd[x][M]$ मधील
लॅम्डा अमूर्तीकरण $\lambd[x]$ हे~$M$ मधील $x$ च्या अशा
आस्थिती \emph{बद्ध करते} ज्या दुसऱ्या लॅम्डा अमूर्तीकरणाने
आधीच बद्ध केलेल्या नाहीत---म्हणजे~$M$ मधील $x$ च्या
\emph{मुक्त} आस्थिती.

\begin{defn}[व्याप्ती]
पद~$N$ मध्ये $\lambd[x][M]$ आढळत असेल, तर त्यातील~$M$
ची संबंधित आस्थिती ही~$\lambd[x]$ ची
\emph{व्याप्ती} असते.
\end{defn}


\begin{defn}[मुक्त आणि बद्ध आस्थिती]
  पद~$M$ मधील चर $x$ ची एखादी आस्थिती $\lambd[x]$ च्या
  व्याप्तीत नसेल तर ती \emph{मुक्त} असते; अन्यथा ती
  \emph{बद्ध} असते. $\lambd[x][M]$ मधील चर~$x$ ची
  एखादी आस्थिती सुरुवातीच्या $\lambd[x]$ ने बद्ध केलेली
  असते, जर आणि फक्त जर~$M$ मधील $x$ ची ती आस्थिती मुक्त असते.
\end{defn}

\begin{ex}
  $\lambd[x][x y]$ मध्ये $x$ आणि $y$ हे दोन्ही $\lambd[x]$
  च्या व्याप्तीत आहेत; म्हणून $x$ हे $\lambd[x]$ ने बद्ध होते.
  $y$ कोणत्याही $\lambd[y]$ च्या व्याप्तीत नसल्याने ते मुक्त
  असते. $\lambd[x][x x]$ मध्ये $x$ च्या दोन्ही आस्थिती
  $\lambd[x]$ ने बद्ध होतात, कारण $xx$ मध्ये त्या दोन्ही मुक्त
  आहेत. $((\lambd[x][xx])x)$ मध्ये~$x$ ची शेवटची आस्थिती
  मुक्त आहे, कारण ती कोणत्याही $\lambd[x]$ च्या व्याप्तीत नाही.
  $\lambd[x][(\lambd[x][x])x]$ मध्ये पहिल्या $\lambd[x]$
  ची व्याप्ती $(\lambd[x][x])x$ आहे आणि दुसऱ्या $\lambd[x]$
  ची व्याप्ती शेवटून दुसरी असलेली~$x$ ची आस्थिती आहे.
  $(\lambd[x][x])x$ मध्ये $x$ ची शेवटची आस्थिती मुक्त,
  तर शेवटून दुसरी बद्ध आहे. म्हणून
  $\lambd[x][(\lambd[x][x])x]$ मधील $x$ ची शेवटून दुसरी
  आस्थिती दुसऱ्या $\lambd[x]$ ने, आणि शेवटची आस्थिती
  पहिल्या $\lambd[x]$ ने बद्ध होते.
\end{ex}

पद $P$ मधील चरांच्या सर्व आस्थिती तपासून मुक्त चरांचा संच
मिळवता येतो. हा संच $\FV{P}$ असा दर्शवतात; त्याची
स्वाभाविक व्याख्या पुढीलप्रमाणे आहे:

\begin{defn}[पदाचे मुक्त चरे] \label{lam:syn:fv:def:fv}
  पदाच्या \emph{मुक्त चरांचा} संच विगमनाने पुढीलप्रमाणे परिभाषित करतात:
  \begin{enumerate}
    \item $\FV{x} = \{x\}$ \label{lam:syn:fv:def:fv1}
    \item $\FV{\lambd[x][N]} = \FV{N} \setminus \{x\}$ \label{lam:syn:fv:def:fv2}
    \item $\FV{PQ} = \FV{P} \cup \FV{Q}$ \label{lam:syn:fv:def:fv3}
  \end{enumerate}
\end{defn}

\begin{prob}
  \begin{enumerate}
  \item या पदातील $\lambd[g]$ आणि दोन $\lambd[x]$ यांच्या
    व्याप्ती ओळखा: $\lambd[g][(\lambd[x][g (x x)]) \lambd[x][g (x x)]]$.
  \item $\lambd[g][(\lambd[x][g (x x)]) \lambd[x][g (x x)]]$
    मधील चरांच्या सर्व आस्थिती बद्ध आहेत का? असल्यास,
    त्यांना अनुक्रमे कोणती अमूर्तीकरणे बद्ध करतात?
    \item $\FV{\lambd[x][(\lambd[y][(\lambd[z][x y]) z]) y]}$
      काढा.
  \end{enumerate}
\end{prob}

\begin{explain}
मुक्त चर म्हणजे जणू बाह्य जगाचा (\emph{परिसराचा}) संदर्भ.
मुक्त चरे असलेले पद अंशतःच निश्चित झालेले पद मानता येते,
कारण त्याचे वर्तन आपण परिसर कसा ठरवतो यावर अवलंबून असते.
उदाहरणार्थ, $\lambd[x][f x]$ हे पद फलसाधक~$x$ स्वीकारून
त्यावर $f$ लावून मिळणारा परिणाम देते; येथे चर~$f$ मुक्त आहे.
या पदाचे मूल्य ते ज्या परिसरात आहे त्यावर, विशेषतः त्या
परिसरातील $f$ च्या मूल्यावर, अवलंबून असते.

या पदावर अमूर्तीकरण केल्यास $\lambd[f][\lambd[x][f
    x]]$ मिळते. आता हे पद परिसरातील चर $f$ वर अवलंबून
नसते, कारण ते दोन फलसाधक स्वीकारून पहिला दुसऱ्यावर
लावल्याचा परिणाम देणारे फलन दर्शवते. परिसरातील $f$
बदलला तरी या पदाच्या वर्तनावर परिणाम होत नाही: हे पद
फक्त फलसाधक म्हणून दिलेले मूल्य वापरते, परिसरातील
$f$ चे मूल्य वापरत नाही.
\end{explain}

\begin{defn}[बंद पद, संयोजक]
  मुक्त चरे नसलेल्या पदाला \emph{बंद पद} किंवा
  \emph{संयोजक} म्हणतात.
\end{defn}

\begin{lem}\label{lam:syn:fv:lem:fv}
  \begin{enumerate}
    \item \label{lam:syn:fv:lem:fv-abs} $y \neq x$ असेल, तर $y \in
      \FV{\lambd[x][N]}$ जर आणि फक्त जर $y \in \FV{N}$.
    \item \label{lam:syn:fv:lem:fv-app} $y \in \FV{PQ}$ जर आणि
      फक्त जर $y \in \FV{P}$ किंवा $y
      \in \FV{Q}$.
    \end{enumerate}
\end{lem}

\begin{proof}
  सरावासाठी.
\end{proof}

\begin{prob}
  \ref{lam:syn:fv:lem:fv} सिद्ध करा.
\end{prob}











\section{आदेशन}\label{lam:syn:sub:sec}

\begin{explain}
मुक्त चरे हे परिसरातील चरांचे संदर्भ असतात. म्हणून मुक्त
चराच्या जागी ठराविक मूल्य वापरणे अर्थपूर्ण ठरते. उदाहरणार्थ,
$\lambd[x][f x]$ मधील $f$ च्या जागी ओळख फलनासारखे
$\lambd[y][y]$ हे विशिष्ट पद ठेवायचे असू शकते. त्यातून
$\lambd[x][(\lambd[y][y]) x]$ मिळते. मुक्त चरांच्या जागी
लॅम्डा पदे ठेवण्याच्या प्रक्रियेला आदेशन म्हणतात.
\end{explain}

\begin{defn}[आदेशन] \label{lam:syn:sub:defn:substitution}
  पद $M$ मधील चर $x$ च्या जागी पद $N$ चे
  \emph{आदेशन}, $\Subst{M}{N}{x}$, विगमनाने पुढीलप्रमाणे
  परिभाषित करतात:
  \begin{enumerate}
    \item $\Subst{x}{N}{x} = N$. \label{lam:syn:sub:defn:substitution-1}
    \item $\Subst{y}{N}{x}  = y$, जर $x \neq y$. \label{lam:syn:sub:defn:substitution-2}
    \item $\Subst{PQ}{N}{x}  = (\Subst{P}{N}{x}) (\Subst{Q}{N}{x})$.
      \label{lam:syn:sub:def:substitution-3}
    \item $\Subst{(\lambd[y][P])}{N}{x} = \lambd[y][\Subst{P}{N}{x}]$,
      जर $x \neq y$ आणि $y \notin \FV{N}$; अन्यथा अपरिभाषित.
      \label{lam:syn:sub:defn:substitution-4}
  \end{enumerate}
\end{defn}

\begin{explain}
\ref{lam:syn:sub:defn:substitution}\ref{lam:syn:sub:defn:substitution-4} मध्ये
$x \neq y$ ही अट घालण्याचे कारण असे की~$M$ मधील
चर~$x$ च्या \emph{बद्ध} आस्थितींच्या जागी~$N$ ठेवायचे नाही.
उदाहरणार्थ, $\Subst{\lambd[x][x]}{y}{x}$ विचारात घेताना
$\lambd[x][y]$ मिळू नये; बद्ध आस्थिती बदलली नाही तर
$\lambd[x][x]$ हेच पद राहते. मात्र वरच्या औपचारिक
नियमांनुसार या प्रसंगी आदेशन अपरिभाषित आहे.

$\lambd[y][P]$ मध्ये~$x$ च्या जागी $N$ चे आदेशन करताना
$y \notin \FV{N}$ हीही अट घालतात. उदाहरणार्थ,
$\lambd[y][x]$ मध्ये $x$ च्या जागी $y$ चे आदेशन,
म्हणजे $\Subst{\lambd[y][x]}{y}{x}$, करता येत नाही:
त्यातून $\lambd[y][y]$ मिळेल. हे पद फलसाधक स्वीकारून
तोच परत करणारे फलन दर्शवते. पण $\lambd[y][x]$ हे पद
नेहमी~$x$ (किंवा $x$ ज्याचा संदर्भ देते ते) परत करणारे
फलन दर्शवते. म्हणून अपेक्षित परिणाम म्हणजे फलसाधक
स्वीकारून तो बाजूला ठेवणारे आणि परिसरातील चर~$y$
परत करणारे फलन. हे नीट साधण्यासाठी आधी बद्ध चर~$y$
चे नाव ``बदलावे'' लागेल.
\end{explain}

\begin{prob}
  पुढील आदेशनांचे परिणाम काय आहेत?
  \begin{enumerate}
  \item $\Subst{\lambd[y][x(\lambd[w][vwx])]}{(uv)}{x}$
  \item $\Subst{\lambd[y][x(\lambd[x][x])]}{(\lambd[y][xy])}{x}$
  \item $\Subst{y(\lambd[v][xv])}{(\lambd[y][vy])}{x}$
  \end{enumerate}
\end{prob}

\begin{thm} \label{lam:syn:sub:thm:notinfv}
  $x \notin \FV{M}$ असेल, तर डावी बाजू परिभाषित
  असल्यास $\FV{\Subst{M}{N}{x}} = \FV{M}$.
\end{thm}

\begin{proof}
  $M$ च्या रचनेवर विगमन करू.
  \begin{enumerate}
  \item $M$ हे चर आहे: सरावासाठी.
  \item $M$ हे $(PQ)$ या रूपाचे आहे: सरावासाठी.
  \item $M$ हे $\lambd[y][P]$ या रूपाचे आहे. तसेच
    $\Subst{\lambd[y][P]}{N}{x}$ परिभाषित असल्यामुळे
    ते $\lambd[y][\Subst{P}{N}{x}]$ असते. म्हणून
    $\Subst{P}{N}{x}$ परिभाषित असते; शिवाय $x \neq y$
    आणि $x \notin \FV{P}$. त्यामुळे:
    \begin{multline*}
      \FV{\Subst{\lambd[y][P]}{N}{x}} = \\
    \begin{aligned}
      & = \FV{\lambd[y][\Subst{P}{N}{x}]} &&
      \text{\ref{lam:syn:sub:defn:substitution-4} नुसार}\\
      & = \FV{\Subst{P}{N}{x}} \setminus \{y\} &&
      \text{\ref{lam:syn:fv:def:fv}\ref{lam:syn:fv:def:fv2} नुसार}\\
      & = \FV{P} \setminus \{y\} && \text{विगमन गृहीतकानुसार} \\
      & = \FV{\lambd[y][P]} && \text{\ref{lam:syn:fv:def:fv}\ref{lam:syn:fv:def:fv2} नुसार}
    \end{aligned}
    \end{multline*}
  \end{enumerate}
\end{proof}

\begin{prob}
  \ref{lam:syn:sub:thm:notinfv} ची सिद्धता पूर्ण करा.
\end{prob}

\begin{thm} \label{lam:syn:sub:thm:infv}
  $x \in \FV{M}$ असेल, तर डावी बाजू परिभाषित
  असल्यास $\FV{\Subst{M}{N}{x}} = (\FV{M} \setminus
  \{x\}) \cup \FV{N}$.
\end{thm}

\begin{proof}
  $M$ च्या रचनेवर विगमन करू.
  \begin{enumerate}
  \item $M$ हे चर आहे: सरावासाठी.
  \item $M$ हे $PQ$ या रूपाचे आहे: $\Subst{(PQ)}{N}{x}$
    परिभाषित असल्यामुळे ते
    $(\Subst{P}{N}{x})(\Subst{Q}{N}{x})$ असते आणि दोन्ही
    आदेशन परिभाषित असतात. शिवाय $x \in \FV{PQ}$
    असल्यामुळे $x \in \FV{P}$ किंवा $x \in \FV{Q}$,
    किंवा दोन्ही, असते. उरलेला भाग सरावासाठी आहे.
  \item $M$ हे $\lambd[y][P]$ या रूपाचे आहे.
    $\Subst{\lambd[y][P]}{N}{x}$ परिभाषित असल्यामुळे
    ते $\lambd[y][\Subst{P}{N}{x}]$ असते;
    $\Subst{P}{N}{x}$ परिभाषित असते, $x \neq y$ आणि
    $y \notin \FV{N}$. तसेच $x \in
    \FV{\lambd[y][P]}$ असल्यामुळे $x \in \FV{P}$.
    आता:
    \begin{multline*}
      \FV{\Subst{(\lambd[y][P])}{N}{x}} = \\
      \begin{aligned}
      & = \FV{\lambd[y][\Subst{P}{N}{x}]} \\
      & = \FV{\Subst{P}{N}{x}} \setminus \{y\} \\
      & = ((\FV{P} \setminus \{x\}) \cup \FV{N}) \setminus \{y\}
       && \text{विगमन गृहीतकानुसार} \\
      & = (\FV{P} \setminus \{x, y\}) \cup \FV{N}
       && y \notin \FV{N} \\
      & = (\FV{\lambd[y][P]} \setminus \{x\}) \cup \FV{N}
      \end{aligned}
      \end{multline*}
  \end{enumerate}
\end{proof}

\begin{prob}
  \ref{lam:syn:sub:thm:infv} ची सिद्धता पूर्ण करा.
\end{prob}

\begin{thm}\label{lam:syn:sub:thm:clr}
  $x \notin \FV{N}$ आणि आदेशन परिभाषित
  असेल, तर $x \notin \FV{\Subst{M}{N}{x}}$.
\end{thm}

\begin{proof}
  सरावासाठी.
\end{proof}

\begin{prob}
  \ref{lam:syn:sub:thm:clr} सिद्ध करा.
\end{prob}


\begin{thm}\label{lam:syn:sub:thm:inv}
  $\Subst{M}{y}{x}$ परिभाषित असेल आणि $y \notin \FV{M}$,
  तर $\Subst{\Subst{M}{y}{x}}{x}{y} = M$.
\end{thm}

\begin{proof}
  $M$ च्या रचनेवर विगमन करू.
  \begin{enumerate}
  \item $M$ हे चर $z$ आहे: सरावासाठी.
  \item $M$ हे $(PQ)$ या रूपाचे आहे. मग:
    \begin{align*}
      \Subst{\Subst{(PQ)}{y}{x}}{x}{y}
      &=\Subst{((\Subst{P}{y}{x})(\Subst{Q}{y}{x}))}{x}{y} \\
      &= (\Subst{\Subst{P}{y}{x}}{x}{y})(\Subst{\Subst{Q}{y}{x}}{x}{y}) \\
      &= (PQ) \text{ विगमन गृहीतकानुसार}
    \end{align*}
  \item $M$ हे $\lambd[z][N]$ या रूपाचे आहे.
    $\Subst{\lambd[z][N]}{y}{x}$ परिभाषित असल्यामुळे
    $z \neq y$ हे माहीत आहे. म्हणून:
    \begin{align*}
      \Subst{\Subst{(\lambd[z][N])}{y}{x}}{x}{y}\\
      & = \Subst{(\lambd[z][\Subst{N}{y}{x}])}{x}{y} \\
      & = \lambd[z][\Subst{\Subst{N}{y}{x}}{x}{y}] \\
      &= \lambd[z][N] \text{ विगमन गृहीतकानुसार}
    \end{align*}
  \end{enumerate}
\end{proof}

\begin{prob}
  \ref{lam:syn:sub:thm:inv} ची सिद्धता पूर्ण करा.
\end{prob}











\section{$\alpha$-परिवर्तन}\label{lam:syn:alp:sec}

$\lambd[x][x]$ आणि $\lambd[y][y]$ यांचा संबंध काय?
ही दोन्ही पदे ओळख फलन दर्शवतात. अर्थात, विन्यासाच्या
दृष्टीने ती भिन्न पदे आहेत. त्यांच्यात फक्त बद्ध चराच्या
नावाचा फरक आहे; एकातील बद्ध चराचे ``नामांतर'' केल्यावर
दुसरे पद मिळते. याला \emph{$\alpha$-परिवर्तन} म्हणतात.

\begin{defn}[बद्ध चराचे नामांतर, $\aconvone$]
  पद $M$ मध्ये $\lambd[x][N]$ ची आस्थिती असेल,
  $x \neq y$, $y \notin \FV{N}$ असेल आणि
  $\Subst{N}{y}{x}$ परिभाषित असेल, तर ती आस्थिती
  \begin{equation*}
    \lambd[y][\Subst{N}{y}{x}]
  \end{equation*}
  ने बदलून मिळणाऱ्या $M'$ ला \emph{बद्ध चराचे नामांतर}
  म्हणतात आणि $M \redone[\alpha] M'$ असे लिहितात.
\end{defn}

\begin{defn}[संबंधाची रचनांशी सुसंगती]
  पदांवरील संबंध $R$ पुढील अटी पूर्ण करत असेल, तर तो
  \emph{रचनांशी सुसंगत} म्हणतात:
  \begin{enumerate}
  \item $R N N'$ असेल, तर $R \lambd[x][N] \lambd[x][N']$.
  \item $R P P'$ असेल, तर $R (PQ) (P'Q)$.
  \item $R Q Q'$ असेल, तर $R (PQ) (PQ')$.
  \end{enumerate}
\end{defn}

आता ही व्याख्या दुसऱ्या रीतीने मांडू:
\begin{defn}[बद्ध चराचे नामांतर, $\aconvone$]
  \emph{बद्ध चराचे नामांतर} ($\redone[\alpha]$) हा पदांवरील
  सर्वांत लहान रचनांशी सुसंगत संबंध आहे जो पुढील अट पूर्ण करतो:
  \begin{align*}
    &\lambd[x][N] \redone[\alpha] \lambd[y][\Subst{N}{y}{x}] && \text{जर
      $x \neq y$, $y \notin \FV{N}$} \\
    & &&\text{आणि $\Subst{N}{y}{x}$ परिभाषित असेल}
  \end{align*}
\end{defn}

येथे ``सर्वांत लहान'' म्हणजे सुसंगती आणि अतिरिक्त अट
यांमुळे आवश्यक ठरणाऱ्या जोड्याच या संबंधात असतात;
इतर कोणत्याही नसतात. म्हणून हा संबंध पुढीलप्रमाणेही
परिभाषित करता येतो:

\begin{defn}[बद्ध चराचे नामांतर, $\aconvone$] \label{lam:syn:alp:defn:aconvone}
  \emph{बद्ध चराचे नामांतर} ($\aconvone$) विगमनाने
  पुढीलप्रमाणे परिभाषित करतात:
  \begin{enumerate}
  \item $N \aconvone N'$ असेल, तर $\lambd[x][N] \aconvone
    \lambd[x][N']$. \label{lam:syn:alp:defn:aconvone1}
  \item $P \aconvone P'$ असेल, तर $(PQ) \aconvone (P'Q)$. \label{lam:syn:alp:defn:aconvone2}
  \item $Q \aconvone Q'$ असेल, तर $(PQ) \aconvone (PQ')$. \label{lam:syn:alp:defn:aconvone3}
  \item $x \neq y$, $y \notin \FV{N}$ आणि $\Subst{N}{y}{x}$
    परिभाषित असेल, तर
    $\lambd[x][N] \redone[\alpha] \lambd[y][\Subst{N}{y}{x}]$.
    \label{lam:syn:alp:defn:aconvone4}
  \end{enumerate}
\end{defn}

या व्याख्या सममूल्य आहेत; त्यांची सिद्धता सरावासाठी ठेवतो.
यापुढे आपण विगमनाधारित व्याख्या वापरू.

\begin{defn}[$\alpha$-परिवर्तन, $\aconv$]
  \emph{$\alpha$-परिवर्तन} ($\aconv$) हा पदांवरील
  $\aconvone$ समाविष्ट करणारा सर्वांत लहान परावर्ती
  आणि संक्रमक संबंध आहे.
\end{defn}

वरच्याप्रमाणे, ``सर्वांत लहान'' म्हणजे परावर्तिता, संक्रमणशीलता आणि
$\aconvone$ यांमुळे आवश्यक ठरणाऱ्या जोड्याच या संबंधात
असतात. यावरून पुढील सममूल्य व्याख्या मिळते:

\begin{defn}[$\alpha$-परिवर्तन, $\aconv$] \label{lam:syn:alp:defn:aconv}
  \emph{$\alpha$-परिवर्तन} ($\aconv$) विगमनाने
  पुढीलप्रमाणे परिभाषित करतात:
  \begin{enumerate}
  \item $P \aconv Q$ आणि $Q \aconv R$ असतील, तर $P \aconv R$.
    \label{lam:syn:alp:defn:aconv1}
  \item $P \aconvone Q$ असेल, तर $P \aconv Q$. \label{lam:syn:alp:defn:aconv2}
  \item $P \aconv P$. \label{lam:syn:alp:defn:aconv3}
  \end{enumerate}
\end{defn}

\begin{ex}
  $\lambd[x][f x]$ चे $\alpha$-परिवर्तन करून
  $\lambd[y][f
  y]$ मिळते; उलट दिशेनेही परिवर्तन करता येते.
  अनौपचारिकपणे, ही दोन्ही पदे फलसाधक स्वीकारून
  त्यावर $f$ लावल्याचा परिणाम देणारी फलने आहेत;
  त्यांतील~$f$ हा परिसरातील चराचा संदर्भ आहे.

  $\lambd[x][f x]$ चे $\alpha$-परिवर्तन करून
  $\lambd[x][g
    x]$ मिळत नाही. ही पदे अनुक्रमे परिसरातील
  $f$ आणि~$g$ यांचा संदर्भ घेतात आणि म्हणून भिन्न
  फलने दर्शवतात: $f$ आणि $g$ भिन्न असलेल्या परिसरांत
  त्यांचे वर्तन भिन्न असते.
\end{ex}

\begin{prob}
  पुढील पदजोड्या $\alpha$-परिवर्तनीय आहेत का?
  \begin{enumerate}
  \item $\lambd[x][\lambd[y][x]]$ आणि $\lambd[y][\lambd[x][y]]$
  \item $\lambd[x][\lambd[y][x]]$ आणि $\lambd[c][\lambd[b][a]]$
  \item $\lambd[x][\lambd[y][x]]$ आणि $\lambd[c][\lambd[b][a]]$
  \end{enumerate}
\end{prob}

\begin{lem}\label{lam:syn:alp:lem:fv-one}
  $P \aconvone Q$ असेल, तर $\FV{P} = \FV{Q}$.
\end{lem}

\begin{proof}
  $P \aconvone Q$ च्या निष्पत्ती वर विगमन करू.
  \begin{enumerate}
  \item शेवटचा नियम \ref{lam:syn:alp:defn:aconvone4} असेल, तर $P$
    हे $\lambd[x][N]$ या रूपाचे आणि $Q$ हे
    $\lambd[y][\Subst{N}{y}{x}]$ या रूपाचे असते;
    येथे $x \neq y$, $y \notin \FV{N}$ आणि
    $\Subst{N}{y}{x}$ परिभाषित असते. $x \in \FV{N}$
    आहे की नाही यानुसार प्रकरणे वेगळी करू:
    \begin{enumerate}
    \item $x \in \FV{N}$ असेल, तर:
      \begin{align*}
        \FV{\lambd[y][\Subst{N}{y}{x}]} &=
        \FV{\Subst{N}{y}{x}} \setminus \{y\} \\
        & = ((\FV{N} \setminus \{x\}) \cup \{y\}) \setminus \{y\}
         && \text{\ref{lam:syn:sub:thm:infv} नुसार} \\
        & = \FV{N} \setminus \{x\} \\
        & = \FV{\lambd[x][N]}
      \end{align*}
    \item $x \notin \FV{N}$ असेल, तर:
      \begin{align*}
        \FV{\lambd[y][\Subst{N}{y}{x}]}
        & = \FV{\Subst{N}{y}{x}} \setminus \{y\} \\
        & = \FV{N} \setminus \{x\}
         && \text{\ref{lam:syn:sub:thm:notinfv} नुसार} \\
        & = \FV{\lambd[x][N]}.
      \end{align*}
    \end{enumerate}
  \item उरलेली तीन प्रकरणे सरावासाठी आहेत.
  \end{enumerate}
\end{proof}

\begin{prob}
  \ref{lam:syn:alp:lem:fv-one} ची सिद्धता पूर्ण करा.
\end{prob}

\begin{lem}\label{lam:syn:alp:lem:inv}
  $P \aconvone Q$ असेल, तर $Q \aconvone P$.
\end{lem}

\begin{proof}
  $P \aconvone Q$ च्या निष्पत्ती वर विगमन करू.
  \begin{enumerate}
  \item शेवटचा नियम \ref{lam:syn:alp:defn:aconvone4} असेल, तर $P$
    हे $\lambd[x][N]$ या रूपाचे आणि $Q$ हे
    $\lambd[y][\Subst{N}{y}{x}]$ या रूपाचे असते;
    येथे $x \neq y$, $y \notin \FV{N}$ आणि
    $\Subst{N}{y}{x}$ परिभाषित असते. प्रथम,
    \ref{lam:syn:sub:thm:clr} नुसार $x \notin
    \FV{\Subst{N}{y}{x}}$ आहे. \ref{lam:syn:sub:thm:inv}
    नुसार $\Subst{\Subst{N}{y}{x}}{x}{y}$ केवळ
    परिभाषितच नाही, तर~$N$ च्या बरोबरही आहे. म्हणून
    \ref{lam:syn:alp:defn:aconvone4} नुसार
    $\lambd[y][\Subst{N}{y}{x}]
    \aconvone \lambd[x][\Subst{\Subst{N}{y}{x}}{x}{y}] =
    \lambd[x][N]$.
  \end{enumerate}
\end{proof}

\begin{prob}
  \ref{lam:syn:alp:lem:inv} ची सिद्धता पूर्ण करा.
\end{prob}

\begin{thm}
  $\alpha$-परिवर्तन हा पदांवरील तुल्यता संबंध आहे;
  म्हणजे तो परावर्ती, सममित आणि संक्रमक आहे.
\end{thm}

\begin{proof}
  \begin{enumerate}
  \item प्रत्येक पद $M$ साठी, बद्ध चरांच्या \emph{शून्य}
    नामांतरांनी $M$ चे $M$ मध्ये परिवर्तन होते.
  \item $P$ चे बद्ध चरांच्या काही नामांतरांनी $Q$ मध्ये
    $\alpha$-परिवर्तन होत असेल, तर \ref{lam:syn:alp:lem:inv}
    वापरून ती नामांतरे उलट क्रमाने उलटवता येतात आणि
    $Q$ पासून~$P$ मिळते.
  \item $P$ चे बद्ध चरांच्या एका क्रमाने $Q$ मध्ये आणि
    $Q$ चे दुसऱ्या क्रमाने $R$ मध्ये $\alpha$-परिवर्तन
    होत असेल, तर दोन्ही क्रम सलग लावून $P$ पासून
    $R$ मिळते.
  \end{enumerate}
\end{proof}

यापुढे $M$ आणि $N$ ही पदे \emph{$\alpha$-सममूल्य},
$M \aeq N$, आहेत असे म्हणू, जर आणि फक्त जर $M$ चे $N$ मध्ये
$\alpha$-परिवर्तन होत असेल. आपण नुकतेच दाखवल्याप्रमाणे,
हे घडते, जर आणि फक्त जर $N$ चेही~$M$ मध्ये
$\alpha$-परिवर्तन होते.

\begin{thm}\label{lam:syn:alp:thm:fv}
  $M \aeq N$ असेल, तर $\FV{M} = \FV{N}$.
\end{thm}

\begin{proof}
  \ref{lam:syn:alp:lem:fv-one} वरून लगेच मिळते.
\end{proof}

\begin{lem}\label{lam:syn:alp:lem:sub:R}
  $R \aeq R'$ असेल आणि $\Subst{M}{R}{y}$ परिभाषित
  असेल, तर $\Subst{M}{R'}{y}$ परिभाषित असते आणि
  $\Subst{M}{R}{y}$ शी $\alpha$-सममूल्य असते.
\end{lem}

\begin{proof}
  सरावासाठी.
\end{proof}

\begin{prob}
  \ref{lam:syn:alp:lem:sub:R} सिद्ध करा.
\end{prob}

\ref{lam:syn:sub:sec} मध्ये पाहिल्याप्रमाणे, काही प्रसंगी
आदेशन अपरिभाषित असते. पण पदांवर $\alpha$-परिवर्तन
वापरून बद्ध चरांची नावे बदलल्यास आदेशन नेहमी परिभाषित
करता येते. परिणामी पदाची $\alpha$-सममूल्यता टिकते,
हे पुढील प्रमेय दाखवते:

\begin{thm}\label{lam:syn:alp:thm:sub}
  कोणत्याही $M$, $R$ आणि~$y$ साठी असे $M'$ असते की
  $M \aeq M'$ आणि $\Subst{M'}{R}{y}$ परिभाषित असते.
  आणखी एखादी जोडी $M'' \aeq M$ आणि $R''$ अशी असेल
  की $\Subst{M''}{R''}{y}$ परिभाषित आहे आणि $R'' \aeq R$,
  तर $\Subst{M'}{R}{y} \aeq \Subst{M''}{R''}{y}$.
\end{thm}

\begin{proof}

प्रथम एक सांत संच वर्ज्य नावे म्हणून दिला असता, कोणत्याही
पदातील सर्व बद्ध नावे त्या संचाबाहेर, परस्पर भिन्न आणि
पदाच्या सर्व मूळ नावांपासून वेगळी करता येतात, हे पाहू.
सर्वांत आतल्या अमूर्तीकरणांपासून बाहेरच्या दिशेने नामांतर
करा. प्रत्येक पायरीला आधी निवडलेल्या नावांपासूनही वेगळे
नवे नाव निवडा; अमर्याद चर असल्याने अशी निवड शक्य आहे.
या पायरीत आतल्या सर्व binder ची नावे आधीच ताजी आहेत,
म्हणून जुन्या बाह्य binder चे नाव त्यांच्याशी समान नाही.
नवे नावही त्यांच्याशी किंवा देहातील मुक्त नावांशी समान नाही.
त्यामुळे \ref{lam:syn:sub:defn:substitution} मधील आंशिक
नियमांनी ही renaming substitution परिभाषित असते आणि
\ref{lam:syn:alp:defn:aconvone4} लागू करता येतो. पदाच्या रचनेवरील
विगमनाने प्रत्येक पायरी आणि संपूर्ण सांत प्रक्रिया योग्य ठरते.
आता वर्ज्य नावांत $y$ आणि $\FV{R}$ मधील सर्व नावे घ्या
आणि ही प्रक्रिया $M$ वर करा; मिळालेल्या पदाला $M'$ म्हणा.
प्रत्येक पायरी alpha परिवर्तन असल्याने $M\aeq M'$.
आता कोणतेही बद्ध नाव $y$ नाही किंवा $R$ मध्ये मुक्त नाही.
म्हणून रचनेवरील सरळ विगमनाने $\Subst{M'}{R}{y}$
परिभाषित आहे. याने अस्तित्व सिद्ध होते.

एकमेव alpha-class मिळण्याचा उरलेला दावा पुढीलप्रमाणे
सिद्ध करू. \ref{lam:syn:unq:prop:unq} मधील एकमेव रचनावृक्षात
प्रत्येक मुक्त चराच्या पानावर त्याचे नाव ठेवा; प्रत्येक बद्ध
वापराला त्याला बांधणाऱ्या अमूर्तीकरणाच्या नोडचा निर्देश द्या.
अमूर्तीकरणांच्या नावांखेरीज हा वृक्ष आणि हे निर्देश प्रत्येक
alpha पायरीत जपले जातात. हे rename नियमाच्या ताजेपणाच्या
अटीवरून आणि तीन compatibility नियमांवरून तपासता येते.
उलट, दोन पदांचे असे नामनिरपेक्ष वृक्ष व निर्देश समान असतील,
तर संबंधित अमूर्तीकरणांच्या प्रत्येक स्थानाला दोन्ही पदांच्या
सर्व नावांबाहेरचे तेच ताजे नाव द्या. वरील आतून-बाहेर
नामांतरप्रक्रियेने दोन्ही पदे एकाच नामांकित पदात बदलतात.
म्हणून ती alpha-सममूल्य आहेत.
जेव्हा raw आदेशन परिभाषित असते, तेव्हा ते फक्त मुक्त
चराच्या संबंधित पानांच्या जागी आदेशित पदाचा असा वृक्ष बसवते.
आंशिक नियमातील अटींमुळे त्यातील मुक्त पानांना बाह्य binder
पकडत नाहीत; मूळ बद्ध पानांचे निर्देशही बदलत नाहीत.
त्यामुळे $M'\aeq M''$ आणि $R\aeq R''$ असताना,
दोन्ही आदेशन परिभाषित असतील तर त्यांचे नामनिरपेक्ष
परिणामी वृक्ष समान असतात. वरील उलट दिशेने
$\Subst{M'}{R}{y}\aeq\Subst{M''}{R''}{y}$ मिळते.
याने अस्तित्व आणि परिणामी alpha-class ची uniqueness
दोन्ही दावे सिद्ध होतात.
\end{proof}

\begin{prob}
  \ref{lam:syn:alp:thm:sub} ची सिद्धता पूर्ण करा.
\end{prob}

\begin{cor}\label{lam:syn:alp:cor:sub}
  कोणत्याही $M$, $R$ आणि~$y$ साठी $M'$ आणि $R'$ ची
  अशी जोडी असते की $M' \aeq M$, $R \aeq R'$
  आणि $\Subst{M'}{R'}{y}$ परिभाषित असते. आणखी
  $M'' \aeq M$ आणि $R'' \aeq R$ अशी जोडी असेल
  आणि $\Subst{M''}{R''}{y}$ परिभाषित असेल, तर
  $\Subst{M'}{R'}{y} \aeq
  \Subst{M''}{R''}{y}$.
\end{cor}

\begin{proof}
  \ref{lam:syn:alp:thm:sub} वरून लगेच मिळते.
\end{proof}












\section{डी ब्रुइन निर्देशांक}\label{lam:syn:deb:sec}

$\alpha$-सममूल्यता ही स्वाभाविक कल्पना आहे, कारण
$\alpha$-सममूल्य पदांचा ``अर्थ एकच'' असतो. प्रत्यक्षात,
एकमेकांत $\alpha$-परिवर्तन करता येणारी पदे वेगळी
न मानणारा लॅम्डा पदांचा विन्यास देता येतो. यांपैकी
सर्वाधिक परिचित पद्धतीत चरांच्या जागी त्यांचे
\emph{डी ब्रुइन निर्देशांक} वापरतात.

$\lambd[x][M]$ लिहिताना फलनाचा प्राचल $x$ आहे हे
स्पष्ट सांगतो, म्हणजे~$M$ मध्ये त्या प्राचलाचा संदर्भ
घेण्यासाठी $x$ वापरता येतो. मात्र डी ब्रुइन निर्देशांक
पद्धतीत प्राचलांना नावे नसतात. फलनाच्या मुख्य भागात
प्राचलाचा संदर्भ हा त्याच्या आस्थितीपासून त्याला बद्ध
करणाऱ्या अमूर्तीकरणापर्यंत मधे येणाऱ्या अमूर्तीकरणांची
पातळी मोजणाऱ्या संख्येने दर्शवतात. उदाहरणार्थ,
$\lambd[x][\lambd[y][y x]]$ पाहा: बाहेरील अमूर्तीकरण
चर~$x$ बद्ध करते; आतील अमूर्तीकरण चर~$y$ बद्ध करते.
उपपद $y x$ हे आतील अमूर्तीकरणाच्या व्याप्तीत आहे.
$y$ आणि त्याला बद्ध करणारे~$\lambd[y]$ यांच्यात
एकही अमूर्तीकरण नाही; पण $x$ आणि त्याला बद्ध
करणारे~$\lambd[x]$ यांच्यात एक अमूर्तीकरण आहे.
म्हणून $y x$ साठी $0\, 1$, आणि संपूर्ण पदासाठी
$\lambd[][\lambd[][01]]$ असे लिहितात.

\begin{defn}
  डी ब्रुइन पदे विगमनाने पुढीलप्रमाणे परिभाषित करतात:
  \begin{enumerate}
  \item $n$, जेथे $n$ ही कोणतीही नैसर्गिक संख्या आहे.
  \item $PQ$, जेथे $P$ आणि $Q$ ही दोन्ही डी ब्रुइन पदे आहेत.
  \item $\lambd[][N]$, जेथे $N$ हे डी ब्रुइन पद आहे.
  \end{enumerate}
\end{defn}


सामान्य लॅम्डा पदांचे डी ब्रुइन निर्देशांकित पदांत
औपचारिक रूपांतर पुढीलप्रमाणे करतात:
\begin{defn}
  \begin{align*}
    F_\Gamma(x) &= \Gamma(x) \\
    F_\Gamma(PQ) &= F_\Gamma(P)F_\Gamma(Q) \\
    F_\Gamma(\lambd[x][N]) &= \lambd[][F_{x,\Gamma}(N)]
  \end{align*}
  येथे $\Gamma$ ही शून्यापासून निर्देशांकित चरांची यादी
  आहे आणि $\Gamma(x)$ ही $\Gamma$ मधील चर $x$ ची
  डावीकडून पहिली जागा दर्शवते. त्यामुळे एकाच नावाची
  पुनरावृत्ती असली, तरी सर्वांत जवळच्या अमूर्तीकरणाचा
  निर्देशांक मिळतो. सुरुवातीच्या यादीत रूपांतर करायच्या
  पदाची सर्व मुक्त चरे असणे आवश्यक आहे. उदाहरणार्थ, $\Gamma$ ही $x,y,z$
  असेल, तर $\Gamma(x)$ हे $0$ आणि $\Gamma(z)$ हे $2$ आहे.

  $x,\Gamma$ म्हणजे $\Gamma$ च्या सुरुवातीला $x$
  ठेवून मिळालेली यादी. उदाहरणार्थ, मागील उदाहरणातील
  $\Gamma$ साठी $w,\Gamma$ ही $w,x,y,z$ आहे.
\end{defn}

डी ब्रुइन पदापासून नेहमीचे लॅम्डा पद पुढीलप्रमाणे
परत मिळवतात. हे रूपांतर तेव्हाच परिभाषित असते, जेव्हा
प्रत्येक संख्या त्या आस्थितीवर उपलब्ध असलेल्या यादीतील
एका जागेचा निर्देशांक असते. अमूर्तीकरणाच्या देहात
जाताना यादीच्या सुरुवातीला आणखी एक नाव येते.
उपलब्ध जागेबाहेरचा निर्देशांक असलेल्या पदासाठी
हे उलट रूपांतर परिभाषित नाही:

\begin{defn}
  \begin{align*}
    G_\Gamma(n) &= \Gamma[n] \\
    G_\Gamma(PQ) &= G_\Gamma(P) G_\Gamma(Q) \\
    G_\Gamma(\lambd[][N]) &= \lambd[x][G_{x,\Gamma}(N)]
  \end{align*}
  येथेही $\Gamma$ ही शून्यापासून निर्देशांकित चरांची
  यादी आहे, आणि $\Gamma[n]$ म्हणजे~$n$ व्या जागेवरील
  चर. उदाहरणार्थ, $\Gamma$ ही $x,y,z$ असेल, तर
  $\Gamma[1]$ हे $y$ आहे.

  शेवटच्या समीकरणातील चर $x$ हा $\Gamma$ मध्ये
  नसलेला कोणताही चर निवडतात.
\end{defn}

येथे काही निष्कर्ष त्यांच्या सिद्धता न देता मांडतो:

\begin{prop}
  $M \aconvone M'$ असेल, आणि $\Gamma$ ही
  $\FV{M}$ समाविष्ट करणारी कोणतीही यादी असेल,
  तर $F_\Gamma(M) \eqs F_\Gamma(M')$.
\end{prop}











\section{$\alpha$-सममूल्यतावर्ग म्हणून पदे}\label{lam:syn:tr:sec}

यापुढे आपण पदांचा विचार $\alpha$-सममूल्यतेपर्यंत करू.
म्हणजे एखादे पद लिहिले की ते ज्या $\alpha$-सममूल्यतावर्गात
आहे तो वर्ग अभिप्रेत असेल. उदाहरणार्थ,
$\lambd[a][\lambd[b][a c]]$ असे लिहिल्यावर त्याच्याशी
$\alpha$-सममूल्य असलेल्या सर्व पदांचा संच अभिप्रेत आहे;
त्यात $\lambd[a][\lambd[b][a
    c]]$, $\lambd[b][\lambd[a][b c]]$ इत्यादी येतात.

तसेच मागील विभागांत $N, Q$ अशी अक्षरे एखादे पद
दर्शवत होती; यापुढे ती एक वर्ग दर्शवतील. पुढील अभ्यासाचे
विषय स्वतंत्र पदांऐवजी हे वर्ग असतील. $x, y$ सारखी
अक्षरे मात्र चरच दर्शवतील.

वर्ग~$M$ मधील कोणताही घटक दर्शवण्यासाठी
$\rep{M}$ हे संकेतन वापरू. एकाहून अधिक घटक हवे
असतील, तर $\rep{M}[0], \rep{M}[1]$ इत्यादी असे लिहू.

मांडणी सोपी करण्यासाठी पदांसाठीची संकेते वर्गांसाठीही
वापरू. वर्गांवर पुढील व्याख्या करतो:
\begin{defn}
  \begin{enumerate}
  \item $\lambd[x][N]$ म्हणजे $\lambd[x][\rep{N}]$
    समाविष्ट करणारा वर्ग.
  \item $PQ$ म्हणजे $\rep{P}\rep{Q}$ समाविष्ट करणारा वर्ग.
  \end{enumerate}
\end{defn}

या व्याख्या नीट परिभाषित आहेत हे सहज दिसते,
कारण $\alpha$-परिवर्तन रचनांशी सुसंगत आहे.

\begin{defn} \label{lam:syn:tr:def:fv}
  $\alpha$-सममूल्यतावर्ग~$M$ ची \emph{मुक्त चरे},
  म्हणजे $FV(M)$, ही $FV(\rep{M})$ अशी परिभाषित करतो.
\end{defn}

ही व्याख्या नीट परिभाषित आहे, कारण
$FV(\rep{M}[0]) =
FV(\rep{M}[1])$ हे \ref{lam:syn:alp:thm:fv} मध्ये दाखवले आहे.

वर्गांमध्ये आदेशन करण्यासाठीही पूर्वीचे संकेतन वापरू:
\begin{defn} \label{lam:syn:tr:def:sub}
  $M$ मध्ये $y$ च्या जागी $R$ चे \emph{आदेशन}, म्हणजे
  $\Subst{M}{R}{y}$, हा $\Subst{\rep{M}}{\rep{R}}{y}$
  समाविष्ट करणारा वर्ग म्हणून परिभाषित करतो; येथे
  $\rep{M}$ आणि $\rep{R}$ असे निवडायचे की आदेशन
  परिभाषित असेल.
\end{defn}

\ref{lam:syn:alp:cor:sub} नुसार ही व्याख्याही नीट परिभाषित आहे.

या व्याख्येमुळे युक्तिवाद किती सोपा होतो ते पाहा.
उदाहरणार्थ, पुढील z हे आदेशित चर आणि त्याच्या जागी
बसवायचे चर या दोन्हींपासून वेगळे ताजे नाव निवडा:
\begin{align}
  \Subst{\lambd[x][x]}{y}{x} & =\label{lam:syn:tr:eq:1}\\
  &= \Subst{\lambd[z][z]}{y}{x} \label{lam:syn:tr:eq:2}\\
  &= \lambd[z][\Subst{z}{y}{x}] \\
  &= \lambd[z][z]
\end{align}

\ref{lam:syn:tr:eq:1} कडे अजूनही स्वतंत्र पदांवरील आदेशन म्हणून
पाहिले, तर ते अपरिभाषित आहे. पण आधी सांगितल्याप्रमाणे,
आता आपण वर्गांवरील आदेशन विचारात घेतो; म्हणून
\ref{lam:syn:tr:eq:2} करता येते. $\lambd[x][x]$ आणि
$\lambd[z][z]$ हे एकाच वर्गात असल्याने पहिल्या
पदाऐवजी दुसरे वापरता येते.

याच कारणाने यापुढे आदेशनासाठी आवश्यक अटी पूर्ण
करणारेच प्रतिनिधी आपण निवडले आहेत असे गृहीत धरू.
उदाहरणार्थ, $\Subst{\lambd[x][N]}{R}{y}$ दिसल्यास
त्यातील प्रतिनिधी $\lambd[x][N]$ असा निवडला आहे की
$x \neq y$ आणि $x \notin FV(R)$, असे मानू.
देहातील सर्व आतल्या बद्ध नावांनाही आदेशनाच्या चरापासून
आणि आदेशित पदाच्या मुक्त चरांपासून वेगळे निवडायचे आहे;
केवळ बाहेरच्या अमूर्तीकरणाच्या नावाच्या अटी पुरेशा नाहीत.

$\lambd[x][x]$ ला ``वर्ग'' म्हणणे थोडे विचित्र वाटते.
म्हणून यापुढे अशा वर्गांना $\Lambda$-पदे (लॅम्डा-वर्गपदे)
म्हणू; या भागात पुढे त्यांना साधे ``पदे'' असेही म्हणू.
यामुळे आधीची परिचित $\lambda$-पदे आणि हे वर्ग
यांतील फरक स्पष्ट राहील.

\begin{editorial}
  मात्र पूर्वीच्या पदांचा विचार अजून सोडता येत नाही:
  $\Lambda$-पदांची संपूर्ण व्याख्या $\lambda$-पदांवर
  आधारलेली आहे. तसेच $\Lambda$-पदांवरील फलने थेट
  परिभाषित करण्याची पद्धत आपण अजून दिलेली नाही.
  त्यामुळे अशी फलने प्रथम $\lambda$-पदांवर परिभाषित
  करून मग त्यांना $\Lambda$-पदांवर ``प्रक्षेपित'' करावे
  लागते; आदेशनासाठी आपण तेच केले. मात्र वाचकाला
  $\Lambda$-पदांवरील फलने कशी परिभाषित करता येतील
  हे सहज समजेल असे आपण गृहीत धरतो.
\end{editorial}










\section{$\beta$-न्यूनीकरण}\label{lam:int:bet:sec}

$(\lambd[m][(\lambd[y][y]) m])$ पाहिल्यावर त्याचा
$\lambd[m][m]$ शी काही संबंध असावा असे वाटते:
दुसरे पद पहिल्याचे ``सोपीकरण'' केल्यावर मिळायला हवे.
\emph{$\beta$-न्यूनीकरण} ही कल्पना औपचारिकपणे पकडते.

\begin{defn}[$\beta$-संकुचन, $\bredone$] \label{lam:int:bet:defn:betacontr}
  \emph{$\beta$-संकुचन} ($\bredone$) हा पदांवरील
  सर्वांत लहान रचनांशी सुसंगत संबंध आहे जो पुढील अट पूर्ण करतो:
  \[    (\lambd[x][N])Q \bredone \Subst{N}{Q}{x}
  \]
  $P \bredone Q$ असेल, तर $P$ चे $Q$ मध्ये
  \emph{$\beta$-संकुचन} झाले असे म्हणतो.
  $(\lambd[x][N])Q$ या रूपाच्या पदाला \emph{रेडेक्स} म्हणतात.
\end{defn}

\begin{prob} \label{lam:int:bet:prob:def}
  \ref{lam:syn:alp:defn:aconvone} मध्ये बद्ध चराच्या
  नामांतरासाठी जशा मांडल्या, तशाच $\beta$-संकुचनाच्या
  सममूल्य विगमनाधारित व्याख्या सविस्तर लिहा.
\end{prob}

\begin{defn}[$\beta$-न्यूनीकरण, $\bred$] \label{lam:int:bet:defn:betared}
  \emph{$\beta$-न्यूनीकरण} ($\bred$) हा $\bredone$
  समाविष्ट करणारा पदांवरील सर्वांत लहान परावर्ती आणि
  संक्रमक संबंध आहे. $P \bred Q$ असेल, तर $P$ चे
  $Q$ मध्ये $\beta$-न्यूनीकरण झाले असे म्हणतो.
\end{defn}

संदर्भ स्पष्ट असेल, तर $\bredone$ ऐवजी $\redone$
आणि $\bred$ ऐवजी $\red$ असे लिहू.

अनौपचारिकपणे, $M \bred N$ असते, जर आणि फक्त जर
$\beta$-संकुचनाच्या शून्य किंवा अधिक पायऱ्यांनी
$M$ चे $N$ मध्ये परिवर्तन करता येते.

\begin{defn}[$\beta$-सामान्य]
ज्या पदाचे आणखी $\beta$-संकुचन करता येत नाही,
त्याला \emph{$\beta$-सामान्य} म्हणतात.
\end{defn}

$M \bred N$ आणि $N$ हे $\beta$-सामान्य असेल, तर
$N$ ला~$M$ चे \emph{सामान्य रूप} म्हणतो. एखाद्या पदाचे
सामान्य रूप एकमेव असते का, असा प्रश्न पडू शकतो.
उत्तर होय आहे, हे आपण पुढे पाहू.

काही उदाहरणे पाहू.
\begin{enumerate}
\item आपल्याला मिळते:
\begin{align*}
(\lambd[x][xxy]) \lambd[z][z] & \redone (\lambd[z][z])(\lambd[z][z]) y \\
& \redone (\lambd[z][z]) y \\
& \redone y
\end{align*}
\item पदाचे ``सोपीकरण'' केल्यावर ते प्रत्यक्षात अधिक गुंतागुंतीचेही होऊ शकते:
\begin{align*}
(\lambd[x][xxy])(\lambd[x][xxy]) & \redone (\lambd[x][xxy])(\lambd[x][xxy])y \\
& \redone (\lambd[x][xxy])(\lambd[x][xxy])yy \\
& \redone \dots
\end{align*}
\item पद तसेचही राहू शकते:
\[(\lambd[x][xx])(\lambd[x][xx]) \redone (\lambd[x][xx])(\lambd[x][xx])
\]
\item तसेच काही पदांचे एकाहून अधिक मार्गांनी न्यूनीकरण
  करता येते. उदाहरणार्थ,
\[(\lambd[x][(\lambd[y][yx]) z]) v \redone (\lambd[y][yv]) z
\]
हे बाहेरील उपयोजनाचे संकुचन करून मिळते; आणि
\[(\lambd[x][(\lambd[y][yx]) z]) v \redone (\lambd[x][zx]) v
\]
हे आतील उपयोजनाचे संकुचन करून मिळते. तथापि,
या प्रसंगी दोन्ही पदांचे पुढील न्यूनीकरण एकाच
$zv$ या पदात होते.
\end{enumerate}

शेवटच्या उदाहरणातील समान परिणाम हा योगायोग नाही.
तो लॅम्डा कलनाचा चर्च–रॉसर गुणधर्म म्हणून ओळखला
जाणारा खोल आणि महत्त्वाचा गुणधर्म दाखवतो.

\begin{digress}
  सर्वसाधारणपणे, एखाद्या पदाचे $\beta$-न्यूनीकरण
  अनेक प्रकारे करता येते. म्हणून न्यूनीकरणाच्या अनेक
  पद्धती मांडल्या गेल्या आहेत. त्यांपैकी सर्वाधिक प्रचलित
  \emph{स्वाभाविक पद्धत} आहे. ती नेहमी \emph{डावीकडील
  सर्वप्रथम} रेडेक्सचे संकुचन करते; रेडेक्सची जागा
  पदात तो जिथे सुरू होतो त्यावरून ठरवली जाते.
  या पद्धतीचा उपयुक्त गुणधर्म असा: एखाद्या पदाचे
  कोणत्याही पद्धतीने सामान्य रूपात न्यूनीकरण करता येते
  तेव्हाच स्वाभाविक पद्धतीनेही ते करता येते. पुढे
  वेगळे स्पष्ट केलेले नसेल, तर आपण स्वाभाविक पद्धत वापरू.
\end{digress}

\begin{defn}[$\beta$-सममूल्यता, $\equal$]
  \emph{$\beta$-सममूल्यता} ($\equal$) हा संबंध विगमनाने
  पुढीलप्रमाणे परिभाषित करतात:
  \begin{enumerate}
  \item $M \equal M$.
  \item $M \equal N$ असेल, तर $N \equal M$.
  \item $M \equal N$, $N \equal O$ असतील, तर $M \equal O$.
  \item $M \equal N$ असेल, तर $PM \equal PN$.
  \item $M \equal N$ असेल, तर $MQ \equal NQ$.
  \item $M \equal N$ असेल, तर $\lambd[x][M] \equal \lambd[x][N]$.
  \item $(\lambd[x][N])Q \equal \Subst{N}{Q}{x}$.
  \end{enumerate}
\end{defn}

पहिले तीन नियम या संबंधाला तुल्यता संबंध बनवतात;
पुढील तीन त्याला रचनांशी सुसंगत बनवतात; आणि शेवटचा
नियम त्यात $\beta$-संकुचन समाविष्ट करतो.

अनौपचारिकपणे, दोन पदे $\beta$-सममूल्य असतात, जर आणि फक्त जर
एका पदाचे दुसऱ्यात शून्य किंवा अधिक पायऱ्यांनी
परिवर्तन करता येते, आणि प्रत्येक पायरी $\beta$-संकुचनाची
किंवा $\beta$-संकुचनाच्या ``उलट'' प्रकारची असते.
उलट $\beta$-संकुचन
असे परिभाषित करतात की $M$ चे उलट $\beta$-संकुचन
 $N$ मध्ये होते, जर आणि फक्त जर $N$ चे $\beta$-संकुचन $M$
मध्ये होते.

वरील नियमांखेरीज आणखी नियम जोडून हा संबंध विस्तारू;
विस्तारित सममूल्यता संबंध $\equal[X]$ असा दर्शवू,
जेथे $X$ हा जोडलेला नियम आहे.











\section{$\eta$-परिवर्तन}\label{lam:syn:eta:sec}

$\lambda$-पदांवर आणखी एक संबंध आहे. \ref{lam:syn:fv:sec}
मध्ये आपण $\lambd[x][(fx)]$ हे उदाहरण पाहिले: ते
फलसाधक स्वीकारून त्यावर $f$ लावते. दुसऱ्या शब्दांत,
ते~$f$ हेच फलन दर्शवते: $\lambd[x][(fx)]N$ आणि
$fN$ या दोन्हींचे न्यूनीकरण~$fN$ मध्ये होते. ही कल्पना
मांडण्यासाठी $\eta$-न्यूनीकरण (आणि $\eta$-विस्तार)
वापरतो.

\begin{defn}[$\eta$-संकुचन, $\eredone$]
  \label{lam:syn:eta:defn:beredone}
  \emph{$\eta$-संकुचन} ($\eredone$) हा पदांवरील
  सर्वांत लहान रचनांशी सुसंगत संबंध आहे जो पुढील अट पूर्ण करतो:
  \[    \lambd[x][M x] \eredone M \text{ जर } x \notin FV(M)
  \]
\end{defn}

\begin{defn}[$\beta\eta$-न्यूनीकरण, $\bered$] \label{lam:syn:eta:defn:bered}
  \emph{$\beta\eta$-न्यूनीकरण} ($\bered$) हा $\bredone$
  आणि $\eredone$ समाविष्ट करणारा पदांवरील सर्वांत लहान
  परावर्ती आणि संक्रमक संबंध आहे. म्हणजे परावर्तन आणि
  संक्रमणाचे नियम, तसेच पुढील दोन नियम, त्यात असतात:
  \begin{enumerate}
  \item $M \bredone N$ असेल, तर $M \bered N$. \label{lam:syn:eta:defn:bered3}
  \item $M \eredone N$ असेल, तर $M \bered N$. \label{lam:syn:eta:defn:bered4}
  \end{enumerate}

\end{defn}

\begin{defn}
  $\equal$ हा सममूल्यता संबंध $\eta$-परिवर्तनाच्या
  पुढील नियमाने विस्तारतो:
  \[  \lambd[x][M x] \equal M \text{ जर } x \notin FV(M)
  \]
  विस्तारित संबंध $\equal[\eta]$ असा दर्शवतो.
\end{defn}

$\eta$-सममूल्यता महत्त्वाची आहे, कारण तिचा लॅम्डा
पदांच्या विस्तारात्मकतेशी संबंध आहे:

\begin{defn}[विस्तारात्मकता]
  $\equal$ हा सममूल्यता संबंध (\ext) नियमाने विस्तारतो:
  \begin{center}
  $Mx \equal Nx$ असेल, तर $M \equal N$; येथे $x \notin FV(MN)$.
  \end{center}
  विस्तारित संबंध $\equal[\ext]$ असा दर्शवतो.
\end{defn}

ढोबळपणे, पदांकडे फलने म्हणून पाहिल्यास एकाच
ताज्या मुक्त चराला फलसाधक म्हणून दिल्यावर त्यांची
उपयोजने सममूल्य असतील, तर ती पदे सममूल्य मानावीत,
असे हा नियम सांगतो. येथे एखाद्या ठरावीक स्थिर आदानावर
मिळालेली समान मूल्ये पुरेशी नाहीत; नियमातील चर
दोन्ही पदांत मुक्त नसणे आवश्यक आहे.

आता $\eta$ नियमाने नेमकी विस्तारात्मकताच मिळते,
त्यापलीकडे काही नाही, हे सिद्ध करू.

\begin{thm}
  $M \equal[\ext] N$ असते, जर आणि फक्त जर $M \equal[\eta] N$.
\end{thm}

\begin{proof}
  प्रथम $\equal[\eta]$ हा संबंध विस्तारात्मकतेच्या
  नियमानुसार बंद आहे हे सिद्ध करू. म्हणजे \ext{} नियमाने
  $\equal[\eta]$ मध्ये नवीन काही जोडले जात नाही. त्यामुळे
  $\equal[\eta]$ मध्ये $\equal[\ext]$ समाविष्ट होतो;
  म्हणून $M \equal[\ext] N$ असेल, तर $M \equal[\eta] N$.

  $\equal[\eta]$ हा संबंध \ext{} नियमानुसार बंद आहे
  हे दाखवण्यासाठी लक्षात घ्या: त्या \ext{} नियमाने $M
  \equal N$ हा निष्कर्ष काढताना $Mx
  \equal[\eta] Nx$ हा पूर्वपक्ष असतो. मग $\equal$ च्या
  एका नियमानुसार $\lambd[x][Mx]
  \equal[\eta] \lambd[x][Nx]$ मिळते. दोन्ही बाजूंना
  $\eta$ नियम लावल्यावर $M \equal[\eta] N$ मिळते.

  त्याचप्रमाणे $\eta$ नियम $\equal[\ext]$ मध्ये
  समाविष्ट आहे हे सिद्ध करू. $x \notin
  FV(M)$ असलेली कोणतीही $\lambd[x][Mx]$ आणि $M$
  घेतल्यास $(\lambd[x][Mx])x \equal[\ext] Mx$ मिळते;
  म्हणून \ext{} नियमाने
  $\lambd[x][Mx] \equal[\ext] M$ मिळते.
\end{proof}



\chapter{चर्च–रॉसर गुणधर्म}\label{lam:cr:chap}









\section{व्याख्या आणि गुणधर्म}\label{lam:cr:dap:sec}

या प्रकरणात चर्च–रॉसर गुणधर्माची कल्पना आणि त्या
गुणधर्माचे काही सामान्य पैलू मांडू.

\begin{defn}[चर्च–रॉसर गुणधर्म, CR] पदांवरील संबंध $\xredone$
  पुढील अट पूर्ण करत असेल, तर त्याला \emph{चर्च–रॉसर गुणधर्म}
  आहे असे म्हणतात: $M \xredone P$ आणि $M \xredone Q$
  असतील, तर असे~$N$ अस्तित्वात असते की
  $P \xredone N$ आणि $Q \xredone N$.
\end{defn}

लॅम्डा कलनाकडे संगणनाचे एक प्रतिमान म्हणून पाहता येते.
त्यात सामान्य रूपातील पदे ``मूल्ये'' आहेत आणि पदांवरील
न्यूनीकरणसंबंध हे ``संगणननियम'' आहेत. चर्च–रॉसर गुणधर्म
असे सांगतो की संगणन पुढे नेण्याचे एकाहून अधिक मार्ग
असले, तरी निकाल मिळत असल्यास त्याचे मूल्य एकमेव असते.

प्राथमिक बीजगणितातील उदाहरण घ्या. $4 \times (1+2) + 3$
मोजण्याचे अनेक मार्ग आहेत. आधी $1+2$ चे~$3$ करून
$4
\times 3+3$ मिळवता येते. किंवा वितरणनियम वापरून
आधी $4 \times (1+2)$ चे $4 \times 1+4
\times 2+3$ असे रूप करता येते. या दोन्ही रूपांचे
पुढील न्यूनीकरण $12+3$ मध्ये होते.

$\xredone$ म्हणजे $\beta$-न्यूनीकरण घेतले, तर
चर्च–रॉसर गुणधर्माचा एक परिणाम लगेच दिसतो:
पदाचे सामान्य रूप असेल, तर ते एकमेव असते.
समजा $M$ चे $P$ आणि $Q$ मध्ये न्यूनीकरण होते,
आणि ही दोन्ही सामान्य रूपे आहेत. चर्च–रॉसर गुणधर्मामुळे
असे $N$ असते की $P$ आणि $Q$ या दोन्हींचे त्यात
न्यूनीकरण होते. पण गृहीतकानुसार $P$ आणि $Q$ ही
सामान्य रूपे असल्याने $P$ आणि $Q$ यांचे $N$ मध्ये
न्यूनीकरण फक्त तुच्छ न्यूनीकरणच असू शकते.
म्हणजे $P$, $Q$ आणि $N$ ही समान आहेत.
त्यामुळे पदाच्या सामान्य रूपाचा उल्लेख आपण
\emph{एकमेव} म्हणून करू शकतो.

लॅम्डा कलनाकडे संगणनाचे प्रतिमान म्हणून पाहिल्यास,
एखाद्या पदाचे सामान्य रूप हा त्या पदापासून सुरू झालेल्या
संगणनाचा ``अंतिम निकाल'' मानता येतो. वरील निष्कर्षानुसार
असा अंतिम निकाल अस्तित्वात असेल, तर तो एकच असतो.
उदाहरणार्थ, $4 \times (1+2)+3$ मोजल्याचा निकाल
एकच, म्हणजे~$15$, असतो.

\begin{thm} \label{lam:cr:dap:thm:str}
  संबंध $\xredone$ ला चर्च–रॉसर गुणधर्म असेल आणि
  $\xred$ हा $\xredone$ समाविष्ट करणारा सर्वांत लहान
  संक्रमक संबंध असेल, तर $\xred$ लाही चर्च–रॉसर
  गुणधर्म असतो.
\end{thm}

\begin{proof}
  समजा
  \begin{align*}
    M & \xredone P_1 \xredone \dots \xredone P_m \text{ आणि}\\
    M & \xredone Q_1 \xredone \dots \xredone Q_n.
  \end{align*}
  उंची $m + 1$ आणि रुंदी $n + 1$ असलेली पदांची
  $N$ जाळी रचून प्रमेय सिद्ध करू. $i$ व्या ओळीतील
  आणि $j$ व्या स्तंभातील पद $N_{i,j}$ असे दर्शवू.

  $N_{i,j} \xredone N_{i+1,j}$ आणि
  $N_{i,j} \xredone N_{i,j+1}$ होईल अशी $N$
  रचू. तिची व्याख्या पुढीलप्रमाणे:
  \begin{align*}
    N_{0,0} &= M \\
    N_{i,0} &= P_i && \text{जर } 1 \le i \le m \\
    N_{0,j} &= Q_j && \text{जर } 1 \le j \le n \\
  \intertext{आणि इतर प्रसंगी:}
    N_{i,j} &= R
  \end{align*}
  येथे $R$ हे असे पद आहे की $N_{i-1,j} \xredone R$
  आणि $N_{i,j-1}
  \xredone R$. $\xredone$ च्या चर्च–रॉसर गुणधर्मामुळे
  असे पद नेहमी अस्तित्वात असते.

  आता $N_{m,0} \xredone \dots \xredone N_{m,n}$
  आणि $N_{0,n}
  \xredone \dots \xredone N_{m,n}$ मिळते.
  येथे $N_{m,0}$ हे $P_m$ आणि $N_{0,n}$ हे $Q_n$ आहे.
  $\xred$ च्या व्याख्येनुसार प्रमेय सिद्ध होते.
\end{proof}












\section{समांतर $\beta$-न्यूनीकरण}\label{lam:cr:pb:sec}

या विभागात \emph{समांतर $\beta$-न्यूनीकरण} ही संकल्पना
मांडू आणि तिला चर्च–रॉसर गुणधर्म आहे हे सिद्ध करू.

\begin{defn}[समांतर $\beta$-न्यूनीकरण, $\bredpar$] \label{lam:cr:pb:defn:bredpar}
  पदांचे समांतर न्यूनीकरण ($\bredpar$) पुढीलप्रमाणे
  विगमनाने परिभाषित केले जाते:
  \begin{enumerate}
    \item \label{lam:cr:pb:defn:bredpar1} $x \bredpar x$.
    \item \label{lam:cr:pb:defn:bredpar2} $N \bredpar N'$ असेल, तर $\lambd[x][N] \bredpar
      \lambd[x][N']$.
    \item \label{lam:cr:pb:defn:bredpar3} $P \bredpar P'$ आणि $Q \bredpar Q'$ असतील, तर $PQ \bredpar
      P'Q'$.
    \item \label{lam:cr:pb:defn:bredpar4} $N \bredpar N'$ आणि $Q \bredpar Q'$ असतील, तर
      $(\lambd[x][N])Q \bredpar \Subst{N'}{Q'}{x}$.
  \end{enumerate}
\end{defn}

समांतर $\beta$-न्यूनीकरणात एका पायरीत पदातील कितीही
रेडेक्सचे संकुचन करता येते. ते $\beta$-न्यूनीकरणापेक्षा
वेगळे आहे: मूळ पदात असलेल्या रेडेक्सचेच संकुचन त्या
पायरीत करता येते; समांतर $\beta$-न्यूनीकरणामुळे नव्याने निर्माण
झालेल्या रेडेक्सचे नाही. उदाहरणार्थ,
$(\lambd[f][fx])(\lambd[y][y])$ या पदाचे समांतर
$\beta$-न्यूनीकरण केल्यास तेच पद किंवा
$(\lambd[y][y])x$ मिळू शकते; एका पायरीत पुढे~$x$
मिळू शकत नाही. मात्र नेहमीच्या $\beta$-न्यूनीकरणाने
ते~$x$ मध्ये न्यूनीत होते. कारण संबंधित रेडेक्स समांतर
$\beta$-न्यूनीकरणाची पहिली पायरी झाल्यावरच निर्माण
होतो. दुसरी समांतर $\beta$-न्यूनीकरणाची पायरी
घेतल्यास~$x$ मिळते.

\begin{thm}\label{lam:cr:pb:thm:refl}
  $M \bredpar M$.
\end{thm}

\begin{proof}
  सरावासाठी.
\end{proof}

\begin{prob}
  \ref{lam:cr:pb:thm:refl} सिद्ध करा.
\end{prob}

\begin{defn}[$\beta$-पूर्ण विकास]\label{lam:cr:pb:defn:bcd}
  $M$ चा \emph{$\beta$-पूर्ण विकास}~$\bcd{M}$ पुढीलप्रमाणे
  विगमनाने परिभाषित केला जातो:
  \begin{align}
    \bcd{x} &= x \label{lam:cr:pb:defn:bcd1} \\
    \bcd{(\lambd[x][N])} &= \lambd[x][\bcd{N}] \label{lam:cr:pb:defn:bcd2}\\
    \bcd{(PQ)} &= \bcd{P}\bcd{Q} && \text{जर $P$ हे $\lambd$-अमूर्तीकरण नसेल तर}
    \label{lam:cr:pb:defn:bcd3} \\
    \bcd{((\lambd[x][N])Q)} &= \Subst{\bcd{N}}{\bcd{Q}}{x} \label{lam:cr:pb:defn:bcd4}
  \end{align}
\end{defn}

पदाचा $\beta$-पूर्ण विकास म्हणजे नावाप्रमाणे
``संपूर्ण समांतर न्यूनीकरण''. समांतर $\beta$-न्यूनीकरणात
एखाद्या रेडेक्सचे संकुचन टाळता येते; पण पूर्ण विकासात
सर्व मूळ रेडेक्सचे संकुचन करावेच लागते. म्हणून
$(\lambd[f][fx])(\lambd[y][y])$ चा पूर्ण विकास तेच
पद नसून $(\lambd[y][y])x$ आहे.

\begin{editorial}
  या व्याख्येत अडचण अशी की ($\lambd$-)पदांवरील फलने
  पुनरावर्ती रीतीने कशी परिभाषित करायची हे अजून
  मांडलेले नाही. पुढे दुरुस्त करू.
\end{editorial}

\begin{lem}\label{lam:cr:pb:lem:comp}
  $M \bredpar M'$ आणि $R \bredpar R'$ असतील, तर $\Subst{M}{R}{y}
  \bredpar \Subst{M'}{R'}{y}$.
\end{lem}

\begin{proof}

अमूर्तीकरणांची नावे आवश्यक त्या सर्व मुक्त चरांपासून
आणि आदेशित चरांपासून वेगळी निवडू. पुढील आदेशनसूत्र वापरू:
\[\Subst{\Subst{N}{Q}{x}}{R}{y}
=\Subst{\Subst{N}{R}{y}}{\Subst{Q}{R}{y}}{x},
\qquad x\ne y,\quad x\notin\FV{R}.
\]
हे सूत्र पदाच्या रचनेवरील विगमनाने सिद्ध होते. चराचे नाव
$x$ असल्यास दोन्ही बाजू $\Subst{Q}{R}{y}$ होतात;
नाव $y$ असल्यास दोन्ही बाजू $R$ होतात, कारण $R$ मध्ये
$x$ मुक्त नाही. इतर चर तसेच राहतात. उपयोजनात दोन्ही
उपपदांवर विगमन गृहीतक लागू होते. अमूर्तीकरणात त्याचे
बद्ध नाव दोन्ही आदेशित पदांच्या मुक्त चरांपासून आणि
$x,y$ पासून वेगळे निवडल्याने तेच गृहीतक देहाला लागू
करता येते. याने वर्गपदांवरील सूत्र सिद्ध होते.
या सिद्धतेतील बाह्य अमूर्तीकरणाचे नाव $x$ हे $y$ पासून
आणि $R,R'$ यांच्या सर्व मुक्त चरांपासून वेगळे निवडू.
  $M \bredpar M'$ च्या निष्पत्ती वर विगमन करू.
  \begin{enumerate}
    \item अखेरची पायरी \ref{lam:cr:pb:defn:bredpar1} आहे: सरावासाठी.
    \item अखेरची पायरी \ref{lam:cr:pb:defn:bredpar2} आहे: मग $M$ हे
      $\lambd[x][N]$ आणि $M'$ हे $\lambd[x][N']$ आहे,
      जेथे $N \bredpar N'$. आपल्याला
      $\Subst{(\lambd[x][N])}{R}{y} \bredpar
      \Subst{(\lambd[x][N'])}{R'}{y}$, म्हणजे
      $\lambd[x][\Subst{N}{R}{y}] \bredpar
      \lambd[x][\Subst{N'}{R'}{y}]$, हे सिद्ध करायचे आहे.
      ते \ref{lam:cr:pb:defn:bredpar2} आणि विगमन गृहीतकावरून
      लगेच मिळते.
    \item अखेरची पायरी \ref{lam:cr:pb:defn:bredpar3} आहे: सरावासाठी.
    \item अखेरची पायरी \ref{lam:cr:pb:defn:bredpar4} आहे: $M$ हे
      $(\lambd[x][N])Q$ आणि $M'$ हे $\Subst{N'}{Q'}{x}$ आहे.
      आपल्याला $\Subst{((\lambd[x][N])Q)}{R}{y} \bredpar
      \Subst{\Subst{N'}{Q'}{x}}{R'}{y}$, म्हणजे
      $(\lambd[x][\Subst{N}{R}{y}])\Subst{Q}{R}{y} \bredpar
      \Subst{\Subst{N'}{R'}{y}}{\Subst{Q'}{R'}{y}}{x}$,
      हे सिद्ध करायचे आहे. वरील आदेशनसूत्राने या दोन
      उजव्या बाजू समान आहेत. मग \ref{lam:cr:pb:defn:bredpar4}
      आणि दोन्ही पूर्वपक्षांना लागू केलेल्या विगमन गृहीतकावरून
      आवश्यक समांतर न्यूनीकरण मिळते.
  \end{enumerate}
\end{proof}

\begin{prob}
  \ref{lam:cr:pb:lem:comp} ची सिद्धता पूर्ण करा.
\end{prob}

\begin{lem}\label{lam:cr:pb:lem:cont}
  $M \bredpar M'$ असेल, तर $M' \bredpar \bcd{M}$.
\end{lem}
\begin{proof}
  $M \bredpar M'$ च्या निष्पत्ती वर विगमन करू.
  \begin{enumerate}
    \item अखेरचा नियम \ref{lam:cr:pb:defn:bredpar1} आहे: सरावासाठी.
    \item अखेरचा नियम \ref{lam:cr:pb:defn:bredpar2} आहे: $M$ हे
      $\lambd[x][N]$ आणि $M'$ हे $\lambd[x][N']$ आहे,
      जेथे $N \bredpar N'$. आपल्याला $\lambd[x][N'] \bredpar
      \bcd{(\lambd[x][N])}$, म्हणजे \ref{lam:cr:pb:defn:bcd2} नुसार
      $\lambd[x][N'] \bredpar \lambd[x][\bcd{N}]$,
      दाखवायचे आहे. ते \ref{lam:cr:pb:defn:bredpar2} आणि
      विगमन गृहीतकावरून मिळते.
    \item अखेरचा नियम \ref{lam:cr:pb:defn:bredpar3} आहे: $M$ हे $PQ$
      आणि $M'$ हे $P'Q'$ आहे, जेथे $P$, $Q$, $P'$, $Q'$
      अशी पदे आहेत की $P \bredpar P'$ आणि $Q \bredpar Q'$.
      विगमन गृहीतकानुसार $P' \bredpar \bcd{P}$ आणि
      $Q' \bredpar \bcd{Q}$.
      \begin{enumerate}
        \item काही $x$ आणि $N$ साठी $P$ हे $\lambd[x][N]$
          असेल, तर काही $N'$ साठी $P'$ हे $\lambd[x][N']$
          असावे लागते आणि $N \bredpar N'$. विगमन गृहीतकानुसार
          $N' \bredpar \bcd{N}$ आणि $Q' \bredpar \bcd{Q}$.
          मग \ref{lam:cr:pb:defn:bredpar4} नुसार
          $(\lambd[x][N'])Q' \bredpar
          \Subst{\bcd{N}}{\bcd{Q}}{x}$.
        \item $P$ हे $\lambd$-अमूर्तीकरण नसेल, तर
          \ref{lam:cr:pb:defn:bredpar3} नुसार $P'Q' \bredpar
          \bcd{P}\bcd{Q}$; आणि \ref{lam:cr:pb:defn:bcd3} नुसार
          उजवीकडील पद $\bcd{PQ}$ आहे.
      \end{enumerate}
    \item अखेरचा नियम \ref{lam:cr:pb:defn:bredpar4} आहे: काही
      $x$, $N$, $Q$, $N'$ आणि~$Q'$ साठी $M$ हे
      $(\lambd[x][N])Q$ आणि $M'$ हे $\Subst{N'}{Q'}{x}$,
      जेथे $N \bredpar N'$ आणि $Q \bredpar Q'$.
      विगमन गृहीतकानुसार $N' \bredpar \bcd{N}$ आणि
      $Q' \bredpar \bcd{Q}$. \ref{lam:cr:pb:lem:comp} नुसार
      $\Subst{N'}{Q'}{x} \bredpar \Subst{\bcd{N}}{\bcd{Q}}{x}$;
      उजवीकडील पद नेमके $\bcd{((\lambd[x][N])Q)}$ आहे.
  \end{enumerate}
\end{proof}

\begin{prob}
  \ref{lam:cr:pb:lem:cont} ची सिद्धता पूर्ण करा.
\end{prob}

\begin{thm}\label{lam:cr:pb:thm:cr}
  $\bredpar$ ला चर्च–रॉसर गुणधर्म आहे.
\end{thm}

\begin{proof}
  \ref{lam:cr:pb:lem:cont} वरून लगेच मिळते.
\end{proof}












\section{$\beta$-न्यूनीकरण}\label{lam:cr:b:sec}

\begin{lem}\label{lam:cr:b:lem:one-par}
  $M \bredone M'$ असेल, तर $M \bredpar M'$.
\end{lem}
\begin{proof} $M \bredone M'$ च्या निष्पत्तीवर विगमन करू.
  मूळस्थानी संकुचनाच्या प्रकरणात काही $x$, $N$ आणि~$Q$ साठी
  $M$ हे $(\lambd[x][N])Q$ आणि $M'$ हे
  $\Subst{N}{Q}{x}$ आहे. \ref{lam:cr:pb:thm:refl} नुसार
  $N \bredpar N$ आणि $Q \bredpar Q$ असल्याने
  \ref{lam:cr:pb:defn:bredpar}\ref{lam:cr:pb:defn:bredpar4} नुसार
  $(\lambd[x][N])Q \bredpar \Subst{N}{Q}{x}$ लगेच मिळते.
  अमूर्तीकरण किंवा उपयोजनाच्या आत संकुचन झाल्यास,
  विगमन गृहीतक, अपरिवर्तित उपपदाची समांतर परावर्तिता
  आणि संबंधित रचना-नियमांवरून निष्कर्ष मिळतो.
\end{proof}

\begin{lem}\label{lam:cr:b:lem:par-red}
  $M \bredpar M'$ असेल, तर $M \bred M'$.
\end{lem}

\begin{proof} $M \bredpar M'$ च्या निष्पत्ती वर विगमन करू.
  \begin{enumerate}
  \item अखेरचा नियम \ref{lam:cr:pb:defn:bredpar1} आहे: मग $M$
    आणि $M'$ दोन्ही $x$ आहेत आणि $x \bred x$.
  \item अखेरचा नियम \ref{lam:cr:pb:defn:bredpar2} आहे: काही
    $x$, $N$, $N'$ साठी $M$ हे $\lambd[x][N]$
    आणि $M'$ हे $\lambd[x][N']$ आहे, जेथे
    $N \bredpar N'$. विगमन गृहीतकानुसार $N \bred N'$.
    मग $N \bred N'$ साठी लागणारी तीच $\bredone$
    संकुचने अमूर्तीकरणाच्या आत करून
    $\lambd[x][N] \bred \lambd[x][N']$ मिळते.
  \item अखेरचा नियम \ref{lam:cr:pb:defn:bredpar3} आहे: काही
    $P$, $Q$, $P'$, $Q'$ साठी $M$ हे $PQ$ आणि
    $M'$ हे $P'Q'$ आहे, जेथे $P \bredpar P'$
    आणि $Q \bredpar Q'$. विगमन गृहीतकानुसार
    $P \bred P'$ आणि $Q \bred Q'$. म्हणून आधी
    $P \bred P'$ आणि नंतर $Q \bred Q'$ असे
    न्यूनीकरण केल्यावर $PQ \bred P'Q'$ मिळते.
  \item अखेरचा नियम \ref{lam:cr:pb:defn:bredpar4} आहे: काही
    $x$, $N$, $N'$, $Q$, $Q'$ साठी $M$ हे
    $(\lambd[x][N])Q$ आणि $M'$ हे
    $\Subst{N'}{Q'}{x}$ आहे, जेथे $N \bredpar N'$
    आणि $Q \bredpar Q'$. विगमन गृहीतकानुसार
    $Q \bred Q'$ आणि $N \bred N'$. म्हणून आधी
    $N \bred N'$, मग $Q \bred Q'$, आणि शेवटी
    $(\lambd[x][N'])Q'$ चे $\Subst{N'}{Q'}{x}$
    मध्ये संकुचन करून
    $(\lambd[x][N])Q \bred \Subst{N'}{Q'}{x}$ मिळते.
  \end{enumerate}
\end{proof}

\begin{lem}\label{lam:cr:b:lem:str}
  $\bredpar$ समाविष्ट करणारा सर्वांत लहान संक्रमक
  संबंध $\bred$ आहे.
\end{lem}

\begin{proof}
  $\bredpar$ समाविष्ट करणारा सर्वांत लहान संक्रमक
  संबंध $\xred$ मानू.

  $\bred \subseteq \xred$: समजा $M \bred M'$,
  म्हणजे $M \ident M_1 \bredone \dots \bredone M_k \ident M'$.
  \ref{lam:cr:b:lem:one-par} नुसार $M \ident M_1 \bredpar
  \dots \bredpar M_k \ident M'$. $\xred$ मध्ये
  $\bredpar$ समाविष्ट आहे आणि तो संक्रमक आहे;
  म्हणून $M \xred M'$. शून्य संकुचनांच्या प्रकरणातही
  समांतर न्यूनीकरणाच्या परावर्ती गुणधर्मामुळे हाच निष्कर्ष मिळतो.

  $\xred \subseteq \bred$: समजा $M \xred M'$,
  म्हणजे $M \ident M_1 \bredpar \dots \bredpar M_k \ident M'$.
  \ref{lam:cr:b:lem:par-red} नुसार $M \ident M_1 \bred
  \dots \bred M_k \ident M'$. $\bred$ संक्रमक
  असल्याने $M \bred M'$.
\end{proof}

\begin{thm}\label{lam:cr:b:thm:cr}
  $\bred$ ला चर्च–रॉसर गुणधर्म आहे.
\end{thm}

\begin{proof}
  \ref{lam:cr:dap:thm:str}, \ref{lam:cr:pb:thm:cr} आणि
  \ref{lam:cr:b:lem:str} वरून लगेच मिळते.
\end{proof}












\section{समांतर $\beta\eta$-न्यूनीकरण}\label{lam:cr:pbe:sec}

या विभागात $\beta\eta$-न्यूनीकरणाशी निगडित समांतर
$\beta\eta$-न्यूनीकरणाला चर्च–रॉसर गुणधर्म आहे हे सिद्ध करू.

\begin{defn}[समांतर $\beta\eta$-न्यूनीकरण, $\beredpar$] \label{lam:cr:pbe:defn:beredpar}
  पदांवरील \emph{समांतर $\beta\eta$-न्यूनीकरण}
  ($\beredpar$) पुढीलप्रमाणे विगमनाने परिभाषित केले जाते:
  \begin{enumerate}
    \item \label{lam:cr:pbe:defn:beredpar1} $x \beredpar x$.
    \item \label{lam:cr:pbe:defn:beredpar2} $N \beredpar N'$ असेल, तर $\lambd[x][N] \beredpar
      \lambd[x][N']$.
    \item \label{lam:cr:pbe:defn:beredpar3} $P \beredpar P'$ आणि $Q \beredpar Q'$ असतील, तर $PQ \beredpar
      P'Q'$.
    \item \label{lam:cr:pbe:defn:beredpar4} $N \beredpar N'$ आणि $Q \beredpar Q'$ असतील, तर
      $(\lambd[x][N])Q \beredpar \Subst{N'}{Q'}{x}$.
    \item \label{lam:cr:pbe:defn:beredpar5} $N \beredpar N'$ असेल, तर $\lambd[x][Nx]
      \beredpar N'$, मात्र $x \notin FV(N)$ असले पाहिजे.
  \end{enumerate}
\end{defn}

\begin{thm}\label{lam:cr:pbe:thm:refl}
  $M \beredpar M$.
\end{thm}

\begin{proof}
  सरावासाठी.
\end{proof}

\begin{prob}
  \ref{lam:cr:pbe:thm:refl} सिद्ध करा.
\end{prob}

\begin{defn}[$\beta\eta$-पूर्ण विकास]\label{lam:cr:pbe:defn:becd}
  $M$ चा \emph{$\beta\eta$-पूर्ण विकास}~$\becd{M}$
  पुढीलप्रमाणे परिभाषित केला जातो. पाचवी अट लागू
  होत असल्यास दुसऱ्या अटीऐवजी ती वापरायची आहे:
  \begin{align}
    \becd{x} &= x \label{lam:cr:pbe:defn:becd1} \\
    \becd{(\lambd[x][N])} &= \lambd[x][\becd{N}] \label{lam:cr:pbe:defn:becd2}\\
    \becd{(PQ)} &= \becd{P}\becd{Q} && \text{जर $P$ हे $\lambd$-अमूर्तीकरण नसेल तर}
    \label{lam:cr:pbe:defn:becd3} \\
    \becd{((\lambd[x][N])Q)} &= \Subst{\becd{N}}{\becd{Q}}{x}
                             \label{lam:cr:pbe:defn:becd4} \\
    \becd{(\lambd[x][Nx])} &= \becd{N} \label{lam:cr:pbe:defn:becd5} & \text{जर $x
                                                      \notin FV(N)$}
  \end{align}
\end{defn}

\begin{lem}\label{lam:cr:pbe:lem:comp}
  $M \beredpar M'$ आणि $R \beredpar R'$ असतील, तर $\Subst{M}{R}{y}
  \beredpar \Subst{M'}{R'}{y}$.
\end{lem}

\begin{proof}
  सर्व अमूर्तीकरणांची नावे आदेशनाच्या चरापासून आणि
  दोन्ही आदेशित पदांच्या मुक्त चरांपासून वेगळी निवडू.
  $M \beredpar M'$ च्या निष्पत्ती वर विगमन करू.

  पहिली चार प्रकरणे \ref{lam:cr:pb:lem:comp} मधील
  प्रकरणांसारखीच आहेत. अखेरचा नियम
  \ref{lam:cr:pbe:defn:beredpar5} असेल, तर काही $x$ आणि~$N'$
  साठी $M$ हे $\lambd[x][Nx]$, $M'$ हे $N'$
  आहे, $x \notin FV(N)$ आणि $N \beredpar N'$.
  आपल्याला $\Subst{(\lambd[x][Nx])}{R}{y} \beredpar
  \Subst{N'}{R'}{y}$, म्हणजे
  $\lambd[x][\Subst{N}{R}{y} x] \beredpar \Subst{N'}{R'}{y}$,
  हे सिद्ध करायचे आहे. ते \ref{lam:cr:pbe:defn:beredpar}\ref{lam:cr:pbe:defn:beredpar5}
  आणि विगमन गृहीतकावरून मिळते. बाह्य बद्ध नाव मूळ
  देहात मुक्त नाही आणि आदेशित पदातही मुक्त नाही, म्हणून
  आदेशनानंतरच्या देहात ते मुक्त होत नाही. त्यामुळे वापरलेल्या
  इटा नियमाची ताजेपणाची अट पूर्ण होते.
\end{proof}

\begin{lem}\label{lam:cr:pbe:lem:cont}
  $M \beredpar M'$ असेल, तर $M' \beredpar \becd{M}$.
\end{lem}

\begin{proof}

प्रथम समांतर न्यूनीकरणाने नवीन मुक्त चरे निर्माण होत नाहीत,
हे त्याच्या पाच नियमांवरील विगमनाने मिळते. बीटा प्रकरणात
आदेशनाचे मुक्त-चर नियम वापरतात; इटा प्रकरणात बद्ध चर
उरलेल्या पदात मुक्त नसतो. पूर्ण विकास हाही एक समांतर
न्यूनीकरण आहे, हे त्याच्या व्याख्येवरील विगमनाने मिळते.
म्हणून पूर्ण विकासानेही मुक्त चरे वाढत नाहीत. खाली आवश्यक
तेथे प्रतिनिधींची बद्ध नावे सर्व संबंधित सांत मुक्त-चर
संचांपासून वेगळी निवडू.
पदाच्या रचनेवर पुढील दोन दावे एकत्र सिद्ध करू:
\begin{enumerate}
\item मुख्य दावा: $M\beredpar M'$ असल्यास
$M'\beredpar\becd{M}$.
\item सहाय्यक दावा: $A=\lambd[x][B]$, $A\beredpar A'$
आणि $S\beredpar T$ असल्यास
$A'S\beredpar\Subst{\becd{B}}{T}{x}$.
\end{enumerate}
सहाय्यक दावा अमूर्तीकरणांपुरताच आहे. तो प्रथम पाहू.
$A\beredpar A'$ ची शेवटची पायरी अमूर्तीकरणाचा नियम
असेल, तर $A'=\lambd[x][B']$ आणि $B\beredpar B'$.
लहान पदासाठीच्या मुख्य विगमन दाव्याने
$B'\beredpar\becd{B}$ मिळते. याला $S\beredpar T$
जोडून बीटा समांतर नियमाने सहाय्यक निष्कर्ष मिळतो.
शेवटची पायरी इटा नियमाची असेल, तर $B=Hx$,
$x\notin\FV{H}$, $A'=H'$ आणि $H\beredpar H'$.
$H$ अमूर्तीकरण नसेल, तर लहान पदासाठीच्या मुख्य दाव्याने
$H'\beredpar\becd{H}$ मिळते. उपयोजनाच्या नियमाने
$H'S\beredpar\becd{H}T$; ही उजवी बाजू
$\Subst{\becd{(Hx)}}{T}{x}$ आहे.
$H=\lambd[y][L]$ असेल, तर लहान अमूर्तीकरणासाठीचा
सहाय्यक विगमन दावा
$H'S\beredpar\Subst{\becd{L}}{T}{y}$ देतो.
येथे $y$ हे $x$ पासून आणि $T$ च्या मुक्त चरांपासून
वेगळे निवडले आहे. पूर्ण विकासाच्या चौथ्या नियमाने
$\becd{(Hx)}=\Subst{\becd{L}}{x}{y}$; तसेच $x$
हे $L$ मध्ये किंवा त्याच्या पूर्ण विकासात मुक्त नाही.
\ref{lam:cr:pb:lem:comp} च्या सिद्धतेत दिलेल्या आदेशनसूत्राने
\[\Subst{\Subst{\becd{L}}{x}{y}}{T}{x}
=\Subst{\becd{L}}{T}{y}.
\]
म्हणून या प्रकरणातही सहाय्यक दावा सिद्ध होतो.

आता मुख्य दावा सिद्ध करू. चरासाठी तो परावर्तितेने मिळतो.
$M=PQ$ अमूर्तीकरण नसलेल्या फलनस्थानील पदासह असेल,
तर समांतर उपयोजनाचाच नियम लागू होतो. दोन्ही उपपदांना
मुख्य विगमन दावा लावून $P'Q'\beredpar\becd{P}\becd{Q}$
मिळते; उजवी बाजू $\becd{M}$ आहे.
$M=(\lambd[x][N])Q$ असेल आणि शेवटची पायरी बीटा
नियमाची असेल, तर $M'=\Subst{N'}{Q'}{x}$,
$N\beredpar N'$ आणि $Q\beredpar Q'$.
दोन्ही मुख्य विगमन दावे आणि \ref{lam:cr:pbe:lem:comp} यांनी
$M'\beredpar\Subst{\becd{N}}{\becd{Q}}{x}=\becd{M}$
मिळते. शेवटची पायरी उपयोजनाची असेल, तर
$M'=P'Q'$, $\lambd[x][N]\beredpar P'$ आणि
$Q\beredpar Q'$. मुख्य विगमन दाव्याने
$Q'\beredpar\becd{Q}$. मग लहान फलनस्थानील
अमूर्तीकरणासाठीचा सहाय्यक दावा तोच आवश्यक निष्कर्ष देतो;
यात $P'$ अजून अमूर्तीकरण असणे गृहीत धरलेले नाही.
शेवटी $M=\lambd[x][N]$ पाहू. शेवटची पायरी इटा
नियमाची असेल, तर $N=Hx$, $x\notin\FV{H}$ आणि
$M'=H'$ जेथे $H\beredpar H'$. मुख्य विगमन दाव्याने
$H'\beredpar\becd{H}=\becd{M}$.
शेवटची पायरी अमूर्तीकरणाच्या नियमाची असेल, तर
$M'=\lambd[x][N']$ जेथे $N\beredpar N'$.
$M$ इटा रेडेक्स नसेल, तर मुख्य विगमन दावा आणि
अमूर्तीकरणाचा नियम थेट निष्कर्ष देतात.
इटा रेडेक्स असेल, तर $N=Hx$ आणि $x\notin\FV{H}$.
$Hx\beredpar N'$ ची शेवटची पायरी उपयोजनाची असेल,
तर $N'=H'x$, $H\beredpar H'$; मुख्य विगमन दाव्याने
$H'\beredpar\becd{H}$ आणि इटा नियमाने
$\lambd[x][H'x]\beredpar\becd{H}=\becd{M}$.
ही पायरी बीटाची असेल, तर $H=\lambd[y][L]$,
$L\beredpar L'$ आणि $N'=\Subst{L'}{x}{y}$.
नवीन मुक्त चरे निर्माण होत नसल्याने $x\notin\FV{L'}$,
आणि $x\ne y$ अशी नावांची निवड आहे. म्हणून
$\lambd[x][\Subst{L'}{x}{y}]$ हे $\lambd[y][L']$
याच वर्गपदाचे दुसरे रूप आहे. अमूर्तीकरणाच्या नियमाने
$H\beredpar\lambd[y][L']$, आणि लहान पद $H$ साठीचा
मुख्य विगमन दावा
$\lambd[y][L']\beredpar\becd{H}=\becd{M}$ देतो.
ही सर्व शक्य प्रकरणे आहेत. प्रत्येक वापरलेल्या विगमन
दाव्याचे मूलपद सध्याच्या पदाचे योग्य उपपद आहे,
म्हणून एकत्र विगमन पूर्ण झाले.
\end{proof}

\begin{thm}\label{lam:cr:pbe:thm:cr}
  $\beredpar$ ला चर्च–रॉसर गुणधर्म आहे.
\end{thm}

\begin{proof}
  \ref{lam:cr:pbe:lem:cont} वरून लगेच मिळते.
\end{proof}












\section{$\beta\eta$-न्यूनीकरण}\label{lam:cr:be:sec}

$\beta\eta$-न्यूनीकरणाला ($\bered$) चर्च–रॉसर
गुणधर्म आहे.

\begin{lem}\label{lam:cr:be:lem:one-par}
  $M \beredone M'$ असेल, तर $M \beredpar M'$.
\end{lem}

\begin{proof}
  $M \beredone M'$ च्या निष्पत्ती वर विगमन करू.
  $M \eredone M'$ हे $\eta$-परिवर्तनाचे प्रकरण
  असेल (म्हणजे \ref{lam:syn:eta:defn:beredone}),
  तर समांतर इटा-नियमासाठी \ref{lam:cr:pbe:thm:refl}
  वापरू. इतर प्रकरणे \ref{lam:cr:b:lem:one-par}
  प्रमाणेच आहेत.
\end{proof}

\begin{lem}\label{lam:cr:be:lem:par-red}
  $M \beredpar M'$ असेल, तर $M \bered M'$.
\end{lem}

\begin{proof} $M \beredpar M'$ च्या निष्पत्ती वर विगमन करू.

  पहिली चार प्रकरणे समांतर बीटा-न्यूनीकरणाच्या
  संबंधित प्रकरणांप्रमाणे हाताळता येतात. अखेरचा नियम
  \ref{lam:cr:pbe:defn:beredpar5} असेल, तर काही $x$, $N$,
  $N'$ साठी $M$ हे $\lambd[x][Nx]$ आणि $M'$ हे $N'$
  आहे, जेथे $x \notin FV(N)$ आणि $N \beredpar N'$.
  म्हणून आधी $\lambd[x][Nx]$ चे $N$ मध्ये
  $\eta$-परिवर्तनाने न्यूनीकरण करू आणि मग विगमन
  गृहीतकानुसार $N \bered N'$ दाखवणाऱ्या $\beredone$
  पायऱ्या घेऊ.
\end{proof}

\begin{lem}\label{lam:cr:be:lem:str}
  $\beredpar$ समाविष्ट करणारा सर्वांत लहान संक्रमक
  संबंध $\bered$ आहे.
\end{lem}

\begin{proof}
  \ref{lam:cr:b:lem:str} प्रमाणे.
\end{proof}

\begin{thm}\label{lam:cr:be:thm:cr}
  $\bered$ ला चर्च–रॉसर गुणधर्म आहे.
\end{thm}

\begin{proof}
  \ref{lam:cr:dap:thm:str}, \ref{lam:cr:pbe:thm:cr} आणि
  \ref{lam:cr:be:lem:str} वरून मिळते.
\end{proof}


\chapter{लॅम्डा-परिभाष्यता}\label{lam:rep:chap}
\begin{quote}
हे प्रकरण प्रायोगिक स्वरूपाचे आहे. अधिक स्पष्टीकरणाची
  गरज आहे; व्याख्या, सिद्धतांसह विधाने आणि अधिक
  उदाहरणे अशा रचनेत मजकूर अधिक चांगल्या प्रकारे
  मांडायला हवा.
\end{quote}








\section{परिचय}\label{lam:ldf:int:sec}

पहिल्या नजरेला लॅम्डा कलन हे फलने दर्शवणाऱ्या
अभिव्यक्ती आणि त्यांचे इतर अभिव्यक्तींना उपयोजन
यांचे अतिशय अमूर्त कलन वाटते. या कलनाच्या
विन्यासात ही फलने कोणत्या विशिष्ट प्रकारच्या वस्तूंवर
असतात याचा निर्देश नाही; विशेषतः नैसर्गिक संख्यांचा
उल्लेखही नाही. त्याच्या मूलभूत क्रिया—उपयोजन आणि
लॅम्डा अमूर्तीकरण—सर्व प्रकारच्या फलनांवर लागू
होतात, केवळ नैसर्गिक संख्यांवरील फलनांवर नाही.

तरीही थोड्या कल्पकतेने लॅम्डा कलनात अंकगणितीय
फलने, म्हणजे नैसर्गिक संख्यांवरील फलने, परिभाषित
करता येतात. त्यासाठी प्रत्येक नैसर्गिक संख्या~$n \in
\Nat$ करिता विशेष $\lambd$-पद~$\num n$ परिभाषित
करू. ते~$n$ साठीचे \emph{चर्च अंक} आहे. चर्च अंकांना
अलाँझो चर्च यांचे नाव दिले आहे.

\begin{defn}
  $n \in \Nat$ असेल, तर संबंधित \emph{चर्च अंक}
  $\num{n}$ हे~$n$ दर्शवते:
  \[    \num{n} \ident \lambd[fx][f^n(x)]
  \]
  येथे $f^n(x)$ म्हणजे $f$ हे~$x$ वर $n$ वेळा लागू
  केल्याचा परिणाम. उदाहरणार्थ, $\num{0}$ हे
  $\lambd[fx][x]$ आहे आणि $\num{3}$ हे
  $\lambd[fx][f(f(f\,x))]$ आहे.
\end{defn}

चर्च अंक $\num n$ हे दोन फलसाधक $f$ आणि $x$
स्वीकारून $f^n(x)$ परत करणारे फलन दर्शवणारे लॅम्डा
पद म्हणून संकेतित केलेले आहे. चर्च अंक स्पष्टपणे
सामान्य रूपात आहेत.

लॅम्डा कलनात नैसर्गिक संख्या दर्शवणे उपयोगी ठरते
तेव्हाच, जेव्हा त्यांच्यावर संगणनही करता येते.
चर्च अंकांवर लॅम्डा कलनात संगणन करणे म्हणजे
एखादे $\lambd$-पद~$F$ अशा चर्च अंकाला लागू करणे
आणि तयार झालेले पद~$F\, \num n$ सामान्य रूपात
न्यूनीत करणे. ते नेहमी सामान्य रूपात न्यूनीत झाले
आणि ते सामान्य रूप नेहमी एखादे चर्च अंक~$\num m$
असेल, तर संगणनाचा परिणाम~$m$ ही संख्या मानता येते.
मग~$F$ हे $f\colon \Nat \to \Nat$ फलन परिभाषित
करते असे मानता येते: $f(n) = m$ असते, जर आणि फक्त जर
$F\, \num n \red \num m$. चर्च–रॉसर गुणधर्मामुळे
सामान्य रूप अस्तित्वात असेल, तर ते एकमेव असते.
म्हणून $F\, \num n \red \num m$ असेल, तर
$F \, \num n$ चे दुसऱ्या कोणत्याही सामान्य रूपात,
विशेषतः दुसऱ्या चर्च अंकात, न्यूनीकरण होऊ शकत नाही.

उलट, $f\colon \Nat \to \Nat$ फलन दिले असता,
अशा रीतीने~$f$ परिभाषित करणारे पद $F$ आहे का,
असे विचारता येते. असे असेल, तर $F$ हे~$f$ ला
\emph{लॅम्डा-परिभाषित} असे म्हणतो आणि $f$ हे
लॅम्डा-परिभाष्य आहे असे म्हणतो. ही संकल्पना
अनेकस्थानी आणि आंशिक फलनांपर्यंत व्यापक करता येते.

\begin{defn}
$f\colon \Nat^k \to \Nat$ असे मानू; येथे आंशिक
फलनही मान्य आहे. या प्रकरणातील न्यूनीकरण आणि सामान्य
रूप हे बीटा-न्यूनीकरण आणि बीटा-सामान्य रूप या अर्थाने आहेत. सर्व
$n_0$, \dots,~$n_{k-1}$ साठी $f(n_0, \dots,
n_{k-1})$ परिभाषित असेल, तर
\[F \, \num{n_0} \, \num{n_1} \dots \num{n_{k-1}} \red \num{f(n_0, n_1, \dots,
  n_{k-1})}
\]
आणि अन्यथा $F \, \num{n_0} \,
\num{n_1} \dots \num{n_{k-1}}$ ला सामान्य रूप
नसेल, तेव्हा बंद लॅम्डा पद~$F$ हे $f$ ला
\emph{लॅम्डा-परिभाषित} असे म्हणतो.
\end{defn}

स्थिर फलने हे एक अत्यंत सोपे उदाहरण आहे.
$C_k \ident \lambd[x][\num{k}]$ हे पद
$c_k\colon \Nat \to \Nat$ फलनाला
लॅम्डा-परिभाषित करते; येथे $c_k(n) = k$.
कारण कोणत्याही~$n$ साठी $C_k \, \num n \ident
(\lambd[x][\num{k}])\num n \redone \num{k}$.
ओळख फलन $\lambd[x][x]$ ने लॅम्डा-परिभाषित आहे.
अधिक गुंतागुंतीची फलने परिभाषित करणे अर्थातच
कठीण असते आणि त्यासाठी बरीच कल्पकता लागते.
त्यामुळे प्रत्येक संगणनक्षम फलन लॅम्डा-परिभाष्य
आहे हे काहीसे आश्चर्यकारक वाटू शकते. याचा उलट
दिशेचा निष्कर्षही खरा आहे: फलन लॅम्डा-परिभाष्य
असेल, तर ते संगणनक्षम असते.











\section{लॅम्डा-परिभाष्य अंकगणितीय फलने}\label{lam:rep:arf:sec}

\begin{prop}
  \label{lam:rep:arf:prop:succ-ld}
  उत्तरवर्ती फलन~$\Succ$ हे लॅम्डा-परिभाष्य आहे.
\end{prop}

\begin{proof}
उत्तरवर्ती फलनाला लॅम्डा-परिभाषित करणारे पद असे आहे:
\begin{align*}
  \fn{Succ} & \ident \lambd[a][\lambd[fx][f ({a} f x)]].
  \intertext{आपल्या संकेतनरूढीनुसार याचा पूर्ण अर्थ असा:}
  \fn{Succ} & \ident \lambd[a][\lambd[f][\lambd[x][(f (({a} f) x))]]].
  \intertext{$\fn{Succ}$ हे~$a$ हा संख्यारूपी फलसाधक
    स्वीकारणारे फलन आहे आणि त्याचे मूल्य दुसरे फलन,
    $\lambd[fx][f ({a} f x)]$, आहे. ते फलन स्वतः चर्च अंक
    नाही. मात्र $a$ हा फलसाधक चर्च अंक असेल, तर त्याचे
    न्यूनीकरण चर्च अंकात होते. पाहा:}
   (\lambd[a][\lambd[fx][f ({a} f x)]])\,\num{n} & \redone
   \lambd[fx][f ({\num{n}} f x)].
   \intertext{आतील पद $\num{n}fx$ मध्ये रेडेक्स आहे, कारण
     $\num{n}$ हे $\lambd[fx][f^nx]$ आहे. म्हणून
     $\num{n}fx \red f^nx$. त्यामुळे संपूर्ण पदासाठी}
   \fn{Succ}\, \num n & \red \lambd[fx][f(f^n(x))],
\end{align*}
म्हणजे $\num{n+1}$.
\end{proof}

\begin{ex}
  $\fn{Succ}$ हे~$\num{0}$, म्हणजे
  $\lambd[fx][x]$, ला लागू केल्यावर काय होते ते पाहू.
  पदे पूर्णपणे लिहू:
  \begin{align*}
    \fn{Succ}\, \num{0} & \ident (\lambd[a][\lambd[f][\lambd[x][(f (({a} f) x))]]])(\lambd[f][\lambd[x][x]])\\
    & \redone \lambd[f][\lambd[x][(f (({(\lambd[f][\lambd[x][x]])} f) x))]] \\
    & \redone \lambd[f][\lambd[x][(f ({(\lambd[x][x])} x))]] \\
    & \redone \lambd[f][\lambd[x][(f x)]] \ident \num{1}\\
  \end{align*}
\end{ex}

\begin{prob}
  पुढील पद
  \begin{align*}
    \fn{Succ}' & \ident \lambd[{n}][\lambd[fx][{n} f (f x)]]
  \end{align*}
  उत्तरवर्ती फलनाला लॅम्डा-परिभाषित करते. का ते स्पष्ट करा.
\end{prob}

\begin{prop}
  \label{lam:rep:arf:prop:add-ld}
  बेरीज फलन~$\Add$ हे लॅम्डा-परिभाष्य आहे.
\end{prop}

\begin{proof}
बेरीज पुढील पदांनी लॅम्डा-परिभाषित होते:
\begin{align*}
  \fn{Add} & \ident \lambd[{a}{b}][\lambd[fx][{a} f ({b} f x)]]
\intertext{किंवा पर्यायाने,}
  \fn{Add}' & \ident \lambd[{a}{b}][{a}\, \fn{Succ} \, {b}].
\intertext{पहिले बेरीजपद असे कार्य करते: $\fn{Add}$ प्रथम
  ${a}$ आणि ${b}$ या दोन संख्या स्वीकारते. त्यातून
  $f$ आणि $x$ स्वीकारून $af(bfx)$ परत करणारे फलन
  मिळते. $a$ आणि $b$ हे चर्च अंक $\num{n}$ आणि
  $\num{m}$ असतील, तर त्याचे न्यूनीकरण $f^{n+m}(x)$
  मध्ये होते; ते $f^{n}(f^{m}(x))$ सारखेच आहे.
  टप्प्याटप्प्याने पाहू:}
  (\lambd[{a}{b}][\lambd[fx][{a} f ({b} f x)]])\num n\,\num m & \red
  \lambd[fx][\num{n}\, f (\num {m}\, f x)] \\
  & \red \lambd[fx][\num{n}\, f (f^m x)] \\
  & \red \lambd[fx][f^n (f^m x)] \ident \num {n+m}.
\intertext{बेरीजचे दुसरे निरूपण $\fn{Add'}$ वेगळे कार्य
  करते: ते चर्च अंक $\num{n}$ आणि $\num{m}$ यांना
  लागू केल्यास,}
\fn{Add}' \num n \,\num m
& \red \num{n}\, \fn{Succ}\, \num{m}.
\intertext{पण $\num{n} f x$ चे नेहमी $f^n(x)$ मध्ये
  न्यूनीकरण होते. म्हणून,}
  \num{n}\,  \fn{Succ}\, \num{m} & \red \fn{Succ}^n(\num{m}).
\end{align*}
$\fn{Succ}$ हे उत्तरवर्ती फलनाला लॅम्डा-परिभाषित
करते आणि~$m$ ला उत्तरवर्ती फलन $n$ वेळा लागू
केल्यावर $n+m$ मिळते. म्हणून याचे पुढे
$\num{n+m}$ मध्ये न्यूनीकरण होते.
\end{proof}

\begin{prop}
  \label{lam:rep:arf:prop:mult-ld}
  गुणाकार पुढील पदाने लॅम्डा-परिभाष्य आहे:
  \[  \fn{Mult} \ident \lambd[ab][\lambd[fx][a (b f) x]]
  \]
\end{prop}

\begin{proof}
  हे कसे कार्य करते ते पाहण्यासाठी $\fn{Mult}$ ला
  चर्च अंक $\num{n}$ आणि $\num{m}$ लागू करू.
  $\fn{Mult} \, \num{n} \, \num{m}$ चे न्यूनीकरण
  $\lambd[fx][\num{n}(\num{m}\, f)x]$ मध्ये होते.
  पद $\num{m} f$ हे आपल्या फलसाधकाला $f$ हे
  $m$ वेळा लागू करणारे फलन परिभाषित करते.
  परिणामी, $\num{n} (\num{m} f) x$ हे
  ``$f$ हे $m$ वेळा लागू करा'' या फलनाला~$x$ वर
  $n$ वेळा लागू करते. म्हणजे $f$ हे $x$ वर
  $n\cdot m$ वेळा लागू होते. मिळणारे सामान्य पद
  म्हणजे नेमका चर्च अंक $\num{nm}$.
\end{proof}

\begin{editorial}
$\eta$-न्यूनीकरणाने हे पद आणखी सोपे करता येते:
\[  \fn{Mult} \ident \lambd[ab][\lambd[f][a (b f)]].
\]

पण त्याआधी $\eta$-न्यूनीकरण स्पष्ट करावे लागेल.
\end{editorial}

\begin{prob}
पुढील पदाने गुणाकार लॅम्डा-परिभाषित होतो:
\[  \fn{Mult}' \ident \lambd[ab][b (\fn{Add}\, a) \num{0}].
\]
हे का कार्य करते ते स्पष्ट करा.
\end{prob}


घातांक फलनाची $\lambd$-पद म्हणून व्याख्या
आश्चर्यकारकपणे सोपी आहे:
\[  \fn{Exp} \ident \lambd[be][\lambd[fx][e b f x]].
\]
पहिला फलसाधक $b$ हा पाया आणि दुसरा~$e$ हा
घातांक आहे. आपल्या संख्या-संकेतनानुसार $e f$
हे सहजपणे $f^e$ आहे. मात्र चर्च अंकाचे पूर्ण
बीटा-सामान्य रूप मिळण्यासाठी दोन बाह्य प्राचल आवश्यक
आहेत. घातांक शून्य असल्यास आतले पद identity मध्ये
न्यूनीत होते आणि या दोन प्राचलांनी चर्च अंक एक मिळतो.
इतर घातांकांसाठी पद वारंवार फलनसंयोजन करून योग्य
चर्च अंक देते. येथे शून्याच्या शून्य घाताचे मूल्यही एक
मानले आहे. हे समजणे कठीण वाटल्यास,
वारंवार गुणाकारानेही घातांक फलन परिभाषित करता येते:
\[  \fn{Exp}' \ident \lambd[be][e (\fn{Mult}\, b) \num{1}].
\]

चर्च अंकांवरील पूर्ववर्ती फलन आणि वजाबाकी वाटतात
तितकी सोपी नाहीत: त्यांसाठी जोड्यांचे संकेतन आवश्यक आहे.











\section{जोड्या आणि पूर्ववर्ती फलन}\label{lam:ldf:pai:sec}

\begin{defn}
$M$ आणि $N$ यांची जोडी ($\tuple{M,N}$ असे लिहितात)
पुढीलप्रमाणे परिभाषित केली जाते. येथे अमूर्तीकरणाचे
नाव दोन्ही पदांच्या मुक्त चरांपासून वेगळे निवडायचे आहे:
\[\tuple{M,N} \ident \lambd[f][fMN].
\]
\end{defn}

सहज अर्थाने ही जोडी म्हणजे एक फलन: ते दुसरे फलन
स्वीकारते आणि ते जोडीतल्या दोन्ही घटकांना लागू करते.
या कल्पनेवरून दोन पदे घेऊन त्यांची जोडी परत करणारे
पुढील रचनाफलन मिळते:
\[  \fn{Pair} \ident \lambd[mn][\lambd[f][fmn]]
\]
जोडी दिल्यावर तिचे घटक पुन्हा मिळवायचे असतात.
त्यासाठी जोडी फलसाधक म्हणून स्वीकारून तिचा पहिला
किंवा दुसरा घटक परत करणारी दोन घटक-निवड फलने हवीत:
\begin{align*}
  \fn{Fst} & \ident \lambd[p][p(\lambd[mn][m])]\\
  \fn{Snd} & \ident \lambd[p][p(\lambd[mn][n])]
\end{align*}

\begin{prob}
  घटक-निवड फलने $\fn{Fst}$ आणि $\fn{Snd}$ का कार्य
  करतात ते स्पष्ट करा.
\end{prob}

आता जोड्यांच्या साहाय्याने पूर्ववर्ती फलन
लॅम्डा-परिभाषित करता येते:
\[  \fn{Pred} \ident \lambd[n][\fn{Fst}(n (\lambd[p][\tuple{\fn{Snd}\, {p}, \fn{Succ}(\fn{Snd}\, {p})}]) \tuple{\num 0, \num 0})]
\]
$\num n\, f x$ चे $f^{n}(x)$ मध्ये न्यूनीकरण
होते हे लक्षात ठेवा. येथे $f$ हे जोडी $p$ स्वीकारून
$p$ चा दुसरा घटक आणि त्या दुसऱ्या घटकाचा
उत्तरवर्ती यांची नवी जोडी परत करणारे फलन आहे;
$x$ ही $\tuple{\num{0},\num{0}}$ जोडी आहे. म्हणून $n=0$
असेल, तर परिणामी जोडी $\tuple{\num{0},\num{0}}$ असते;
अन्यथा ती $\tuple{\num{n-1}, \num n}$ असते.
मग $\fn{Pred}$ या परिणामी जोडीचा पहिला घटक
परत करते.

नैसर्गिक संख्यांवरील शून्यावर थांबणारी वजाबाकी
म्हणजे $a$ ला $\fn{Pred}$ हे $b$ वेळा लागू करणे:
\[\fn{Sub} \ident \lambd[ab][b \fn{Pred}\, a].
\]











\section{सत्यतामूल्ये आणि संबंध}\label{lam:ldf:tvr:sec}

शुद्ध लॅम्डा कलनात सत्यतामूल्यांचे पुढीलप्रमाणे
संकेतन करता येते:
\begin{align*}
  \fn{true} & \ident \lambd[x][\lambd[y][x]]\\
  \fn{false} & \ident \lambd[x][\lambd[y][y]]
\end{align*}

सत्यतामूल्ये \emph{निवडफलने} म्हणून दर्शवली जातात:
दोन फलसाधक स्वीकारून त्यांपैकी एक परत करणारी
फलने. सत्यतामूल्य $\fn{true}$ पहिला फलसाधक
निवडते आणि $\fn{false}$ दुसरा. उदाहरणार्थ,
$\fn{true}\, M N$ चे नेहमी $M$ मध्ये, तर
$\fn{false}\, M N$ चे नेहमी~$N$ मध्ये न्यूनीकरण होते.

\begin{defn}
संबंध $R \subseteq \Nat^k$ हा लॅम्डा-परिभाष्य
आहे असे म्हणतो, जर आणि फक्त जर असे बंद पद~$R$ असते की
प्रत्येक नैसर्गिक निविष्टांसाठी पुढील अटी पूर्ण होतात:
\begin{align*}
  R\, \num{n_1} \dots \num{n_k} & \bred \fn{true}
  \intertext{जेव्हा $R(n_1, \dots, n_k)$ सत्य असेल, आणि}
  R\, \num{n_1} \dots \num{n_k} & \bred \fn{false}
\end{align*}
अन्यथा.
\end{defn}

उदाहरणार्थ, $0$ साठी खरा आणि $0$ व्यतिरिक्त इतर
कोणत्याही संख्येसाठी खरा नसलेला संबंध
$\fn{IsZero} = \{0\}$ पुढील पदाने
लॅम्डा-परिभाष्य आहे:
\[  \fn{IsZero} \ident \lambd[n][n (\lambd[x][\fn{false}])\, \fn{true}].
\]
हे कसे कार्य करते? चर्च अंक हे पुनरावर्तक आहेत:
पहिला फलसाधक दुसऱ्यावर $n$~वेळा लागू करणारी
फलने. म्हणून सुरुवातीचे मूल्य $\fn{true}$ घेऊ;
पुनरावर्तनाच्या प्रत्येक पायरीत आधीच्या पायरीचा
निकाल काहीही असला, तरी $\fn{false}$ परत करू.
ही पायरी सुरुवातीच्या मूल्यावर $n$ वेळा लागू होते.
म्हणून ती एकदाही लागू झाली नाही, म्हणजे $n = 0$
असेल, तेव्हाच निकाल $\fn{true}$ असतो.

सत्यतामूल्यांच्या या संकेतनावरून आणखी सत्यता-फलने
परिभाषित करता येतात. नकरण आणि संधियोग यांची
दोन निरूपणे अशी:
\begin{align*}
  \fn{Not} & \ident \lambd[x][x\, \fn{false}\, \fn{true}]\\
  \fn{And} & \ident \lambd[x][\lambd[y][xy \,\fn{false}]]
\end{align*}
``$\fn{Not}$'' हे एक फलसाधक स्वीकारते:
तो $\fn{false}$ असेल, तर $\fn{true}$ परत करते;
आणि तो~$\fn{true}$ असेल, तर $\fn{false}$.
``$\fn{And}$'' हे दोन सत्यतामूल्ये स्वीकारते आणि
निकाल $\fn{true}$ असायला हवा, जर आणि फक्त जर
दोन्ही निविष्ट सत्यतामूल्ये~$\fn{true}$ असतील. वर सांगितल्याप्रमाणे सत्यतामूल्ये
निवडफलने आहेत. म्हणून सत्यतामूल्य $x$ ला दोन
फलसाधक लागू केल्यास, $x$ हे $\fn{true}$ असेल
तर पहिला फलसाधक आणि अन्यथा दुसरा मिळतो.
आता $\fn{And}$ आपले दोन फलसाधक $x$ आणि $y$
घेते व $x$ ला अनुक्रमे $y$ आणि $\fn{false}$
हे देते. $x$ सत्यतामूल्य आहे असे गृहीत धरल्यास,
$x$ हे $\fn{true}$ असेल तर निकाल~$y$,
आणि $x$ हे $\fn{false}$ असेल तर निकाल
$\fn{false}$ होतो; हाच अपेक्षित परिणाम आहे.

येथे $\fn{And}$ ला फलसाधक म्हणून फक्त
सत्यतामूल्येच दिली जातात असे आपण गृहीत धरतो.
इतर पदे दिल्यास निकाल—सामान्य रूप अस्तित्वात
असेल तर ते—सत्यतामूल्य असेलच असे नाही.

\begin{prob}
  सत्यतामूल्यांचे $\lambd$-पदांद्वारे केलेले संकेतन
  वापरून समावेशक आणि अपवर्जक विकल्पयोग दर्शवणारी
  $\fn{Or}$ आणि $\fn{Xor}$ ही फलने परिभाषित करा.
\end{prob}











\section{आदिम पुनरावर्ती फलने लॅम्डा-परिभाष्य असतात}\label{lam:ldf:prf:sec}

मूलभूत फलने $\Zero$, $\Succ$ आणि $\Proj{n}{i}$
यांपासून संयोजन व आदिम पुनरावर्तन वापरून परिभाषित
करता येणारी फलने म्हणजे आदिम पुनरावर्ती फलने, हे लक्षात घ्या.

\begin{lem}
  \label{lam:ldf:prf:lem:basic}
  मूलभूत आदिम पुनरावर्ती फलने $\Zero$, $\Succ$ आणि
  प्रक्षेपणे~$\Proj{n}{i}$ ही लॅम्डा-परिभाष्य आहेत.
\end{lem}

\begin{proof}
पुढील पदे त्यांना लॅम्डा-परिभाषित करतात:
\begin{align*}
  \fn{Zero} & \ident \lambd[a][\lambd[fx][x]]\\
  \fn{Succ} & \ident \lambd[a][\lambd[fx][f (a f x)]] \\
  \fn{Proj}^n_i & \ident \lambd[x_0\dots x_{n-1}][x_i]
\end{align*}
\end{proof}

\begin{lem}
  \label{lam:ldf:prf:lem:comp} सर्वत्र परिभाषित $k$-स्थानी फलन $f$ आणि $n$-स्थानी
  सर्वत्र परिभाषित फलने $g_0, \dots, g_{k-1}$ ही अनुक्रमे $F$, $G_0$,
  \dots, $G_{k-1}$ या पदांनी लॅम्डा-परिभाष्य आहेत,
  आणि $h$ हे त्यांच्यापासून संयोजनाने परिभाषित आहे असे समजा.
  मग $h$ हे लॅम्डा-परिभाष्य आहे.
\end{lem}

\begin{proof}
  पुढील पद $h$ ला लॅम्डा-परिभाषित करते:
  \[  H \ident \lambd[x_0 \dots x_{n-1}][F\, (G_0 x_0 \dots
    x_{n-1}) \dots (G_{k-1} x_0 \dots x_{n-1})]
  \]
  याची पडताळणी सरावासाठी सोडली आहे.
\end{proof}

\begin{prob}
  $H\num{n_0}\dots\num{n_{n-1}} \red \num{h(n_0, \dots,
    n_{n-1})}$ हे दाखवून \ref{lam:ldf:prf:lem:comp}
  चा पुरावा पूर्ण करा.
\end{prob}

\ref{lam:ldf:prf:lem:comp} मध्ये $f$ आणि $g_0$,
\dots,~$g_{k-1}$ ही आदिम पुनरावर्ती असणे आवश्यक
नव्हते, हे लक्षात घ्या; ती सर्वत्र परिभाषित आणि
लॅम्डा-परिभाष्य असणे पुरेसे आहे.

\begin{lem}\label{lam:ldf:prf:lem:prim}
  $f$ हे $n$-स्थानी आणि~$g$ हे $n+2$-स्थानी फलन
  असून ती सर्वत्र परिभाषित आणि अनुक्रमे $F$ आणि~$G$ या पदांनी
  लॅम्डा-परिभाष्य आहेत असे समजा. $h$ हे
  $f$ आणि $g$ यांच्यापासून आदिम पुनरावर्तनाने परिभाषित असेल,
  तर $h$ देखील लॅम्डा-परिभाष्य आहे.
\end{lem}

\begin{proof}
  $h$ ची व्याख्या अशी आहे, हे लक्षात घ्या:
  \begin{align*}
    h(x_1, \dots, x_n, 0) &= f(x_1, \dots, x_n)\\
    h(x_1, \dots, x_n, y+1) & = g(x_1, \dots, x_n, y, h(x_1, \dots, x_n, y)).
  \end{align*}
  सहज अर्थाने, आदिम पुनरावर्ती व्याख्येत $f(x_1,
  \dots, x_n)$ या सुरुवातीच्या मूल्यापासून $h$ चे
  मूल्य $y$ वेळा अद्ययावत होते. हे चर्च अंकांच्या
  व्याख्येसारखे आहे: तेही पुनरावर्तक आहेत.

  सोपेपणासाठी आपण एका अतिरिक्त फलसाधक~$x$ साठी
  व्याख्या व पुरावा देऊ. $h$ ला लॅम्डा-परिभाषित
  करणारे पद असे आहे:
  \begin{align*}
    H \ident & \lambd[x][\lambd[y][\fn{Snd} (y
        D \tuple{\num{0}, F x})]]
    \intertext{जेथे }
    D \ident & \lambd[p][\tuple{\fn{Succ} (\fn{Fst}\, p), (G x
          (\fn{Fst}\, p) (\fn{Snd}\, p))}]
  \end{align*}
  पुनरावर्तनात आपण ठेवलेली अवस्था ही जोडी आहे:
  पहिला घटक सध्याचे $y$ आणि दुसरा त्या स्थितीसाठीचे
  $h$ चे मूल्य. प्रत्येक पायरीत $y$ आणि $h$
  यांच्या नव्या मूल्यांची जोडी तयार होते. पुनरावर्तन
  संपल्यावर तिचा दुसरा घटक परत करतो; तेच $h$ चे
  अंतिम मूल्य असते. एकाहून अधिक अतिरिक्त फलसाधक
  असतील, तर ती सर्व निविष्टे सुरुवातीच्या फलनाला आणि
  प्रत्येक पायरीच्या फलनाला त्याच क्रमाने दिली जातात;
  प्रत्येक बाह्य प्राचलासाठी स्वतंत्र अमूर्तीकरण घ्यायचे.
  पुनरावर्तनाच्या जोडीतील मोजणी आणि पूर्वमूल्याचे घटक
  तेच राहतात, म्हणून खालील विगमन तसाच लागू होतो.
  हे निरूपण आदिम पुनरावर्तनाचेच
  आहे, हे आता सिद्ध करू.

  कोणत्याही $n$ आणि $m$ साठी $H\,\num{n}\,\num{m} \red
  \num{h(n,m)}$ दाखवायचे आहे. त्यासाठी आधी $D_n \ident
  \Subst{D}{\num{n}}{x}$ असेल, तर $D_n^m \tuple{\num{0}, F\, \num{n}} \red
  \tuple{\num{m}, \num{h(n, m)}}$ हे दाखवू. $m$ वरील
  विगमनाने पुढे जाऊ.

  $m=0$ असेल, तर $D_n^0 \tuple{\num{0}, F\, \num{n}} \red
  \tuple{\num{0}, \num{h(n, 0)}}$ दाखवायचे आहे.
  पण $D_n^0 \tuple{\num{0}, F\,
    \num{n}}$ म्हणजेच $\tuple{\num{0}, F\, \num{n}}$.
  $F$ हे $f$ ला लॅम्डा-परिभाषित करते; त्यामुळे याचे
  $\tuple{\num{0}, \num{f(n)}}$ मध्ये न्यूनीकरण होते;
  आणि $f(n) = h(n, 0)$ असल्याने हेच
  $\tuple{\num{0}, \num{h(n,0)}}$ आहे.

  आता $D_n^m \tuple{\num{0}, F\, \num{n}} \red
  \tuple{\num{m}, \num{h(n, m)}}$ असे समजा.
  $D_n^{m+1} \tuple{\num{0}, F\, \num{n}} \red
  \tuple{\num{m+1}, \num{h(n, m+1)}}$ हे दाखवू:
  \begin{align*}
    D_n^{m+1} \tuple{\num{0}, F\, \num{n}}
    & \ident D_n(D_n^{m} \tuple{\num{0}, F\, \num{n}})\\
    & \red D_n\,\tuple{\num{m}, \num{h(n, m)}} \text{\qquad (विगमन गृहीतकाने)}\\
    & \ident (\lambd[p][
      \tuple{
        \fn{Succ} (\fn{Fst}\, p), (G \, \num{n} (\fn{Fst}\, p) (\fn{Snd}\, p))
    }]) \tuple{\num{m}, \num{h(n,m)}}\\
    & \redone
    \tuple{\fn{Succ} (\fn{Fst}\, \tuple{\num{m}, \num{h(n,m)}}),\\
      & \qquad (G \, \num{n} (\fn{Fst}\, \tuple{\num{m}, \num{h(n,m)}}) (\fn{Snd}\, \tuple{\num{m}, \num{h(n,m)}}))}\\
    & \red
    \tuple{\fn{Succ}\, \num{m}, (G \, \num{n} \,\num{m} \, \num{h(n,m)})}\\
    & \red \tuple{\num{m+1}, \num{g(n, m, h(n, m))}}
  \end{align*}
  $g(n, m, h(n, m)) = h(n, m+1)$ असल्याने हे सिद्ध झाले.

  शेवटी पुढील न्यूनीकरण पाहा:
  \begin{align*}
    H\,\num n\,\num m &\ident \lambd[x][\lambd[y][\fn{Snd} (y
        (\lambd[p].\tuple{\fn{Succ} (\fn{Fst}\, p), (G\, x\,
          (\fn{Fst}\, p)\, (\fn{Snd}\, p))}) \tuple{\num{0}, F x})]]\\
    & \qquad\,\num n\, \num m\\
    & \red \fn{Snd} (\num m \,
    \underbrace{(\lambd[p].\tuple{\fn{Succ} (\fn{Fst}\, p), (G \,\num n\,
      (\fn{Fst}\, p) (\fn{Snd}\, p))})}_{D_n} \tuple{\num{0}, F \num n})\\
    & \ident \fn{Snd} (\num m \, D_n\, \tuple{\num{0}, F \num n})\\
    & \red \fn{Snd} \,(D_n^m  \tuple{\num{0}, F \num n})\\
    & \red \fn{Snd} \,\tuple{\num{m}, \num{h(n, m)}} \\
    & \red \num{h(n,m)}.
  \end{align*}
\end{proof}

\begin{prop}
  प्रत्येक आदिम पुनरावर्ती फलन लॅम्डा-परिभाष्य आहे.
\end{prop}

\begin{proof}
  \ref{lam:ldf:prf:lem:basic} नुसार सर्व मूलभूत फलने
  लॅम्डा-परिभाष्य आहेत. तसेच \ref{lam:ldf:prf:lem:comp}
  आणि \ref{lam:ldf:prf:lem:prim} नुसार सर्वत्र परिभाषित लॅम्डा-परिभाष्य
  फलनांचा वर्ग संयोजन व आदिम पुनरावर्तन यांच्या
  अंतर्गत बंद आहे.
\end{proof}












\section{स्थिरबिंदू}\label{lam:ldf:fp:sec}

क्रमगुणित फलनाचे पुनरावर्तनाने असे पद $\fn{Fac}$
म्हणून परिभाषण करायचे आहे असे समजा:
\[\fn{Fac} \ident \lambd[n][\fn{IsZero}\, n\, \num 1 (\fn{Mult}\, n (\fn{Fac}(\fn{Pred} \, n)))]
\]
म्हणजे $n = 0$ असेल, तर $n$ चे क्रमगुणित $1$;
अन्यथा ते $n$ गुणिले $n-1$ चे क्रमगुणित.
पण $\fn{Fac}$ ची अशी व्याख्या करता येत नाही,
कारण उजव्या बाजूस $\fn{Fac}$ हेच पद येते.
अशा स्वसंदर्भाचा समावेश असलेल्या पुनरावर्ती
व्याख्या लॅम्डा कलनाचा भाग नाहीत. उदाहरणार्थ,
पुढील पद परिभाषित करताना
\[\fn{Mult} \ident \lambd[ab][b (\fn{Add}\, a) \num{0}]
\]
उजव्या बाजूस $\fn{Add}$ यांसारखी पूर्वी परिभाषित
पदेच वापरली जातात. $\fn{Add}$ च्या जागी त्याचे
परिभाषक पद ठेवून ते काढून टाकता येते. मग
$\fn{Mult}$ शुद्ध लॅम्डा पद म्हणून मिळेल;
$\fn{Add}$ मध्येही इतर परिभाषित पदे असतील
(उदा., $\fn{Add}'$ मध्ये आहेत), तर अशी जागापालट
सुरू ठेवून शेवटी शुद्ध लॅम्डा पद मिळवता येते.

वरच्या $\fn{Fac}$ सारख्या पुनरावर्ती व्याख्येच्या
बाबतीत मात्र तसे होत नाही. उजवीकडील $\fn{Fac}$
च्या जागी $\fn{Fac}$ चीच व्याख्या ठेवल्यास पुढील पद मिळते:
\begin{align*}
  B &\ident \lambd[m][\fn{IsZero}\,m\,\num{1}
       (\fn{Mult}\,m\,(\fn{Fac}(\fn{Pred}\,m)))]\\
  \fn{Fac} &\ident \lambd[n][\fn{IsZero}\,n\,\num{1}
       (\fn{Mult}\,n\,(B(\fn{Pred}\,n)))]
\end{align*}
येथे B हे फक्त आतले पद थोडक्यात लिहिण्यासाठीचे नाव आहे.
त्याच्या व्याख्येतील $\fn{Fac}$ अजूनही गेलेले नाही.
बाह्य अमूर्तीकरणाची व्याप्ती दोन्ही conditional शाखांवर आहे.
ही प्रक्रिया पुन्हा केल्यास व्याख्या वाढतच राहते
आणि कधीही शुद्ध लॅम्डा पद मिळत नाही. म्हणून
क्रमगुणित, किंवा सर्वसाधारणपणे पुनरावर्ती फलने,
अशा रीतीने परिभाषित करता येत नाहीत.

तरी या पुनरावर्ती व्याख्येतून एक गोष्ट समजते:
$f$ हे क्रमगुणित फलन दर्शवणारे पद असेल, तर
\[\fn{Fac}' \ident \lambd[g][\lambd[n][\fn{IsZero} \, n \, \num 1\, (\fn{Mult} \, n\, (g (\fn{Pred} n)))]]
\]
या पदाला $f$ लागू केल्यावर मिळणारे
$\fn{Fac}'\,f$ हेही क्रमगुणित फलन दर्शवते.
म्हणजे $\fn{Fac}'$ याकडे फलन स्वीकारून फलन
परत करणारे फलन म्हणून पाहिल्यास, $\fn{Fac'}\, f$
आणि~$f$ यांचे सर्व चर्च अंकांवरील परिणाम समान असतात.
यावरून ही दोन पदे स्वतः बीटा-सममूल्य आहेत असे मात्र
ठरत नाही; संख्यात्मक फलनाचे निरूपण आणि पदांची
बीटा-सममूल्यता या वेगळ्या अटी आहेत. $\fn{Fac}' \, f \equal[\beta] f$
हा गुणधर्म असलेल्या फलन~$f$ ला $\fn{Fac}'$
चा \emph{स्थिरबिंदू} म्हणतात. पुढील संयोजकाने
$\fn{Fac}'$ चा असा स्थिरबिंदू मिळवून त्याचे सर्व
अंकांवरील परिणाम क्रमगुणिताचे आहेत हे सिद्ध करू.

लॅम्डा कलनात दिलेल्या पदाचा स्थिरबिंदू काढणारी
पदे आहेत. त्यांच्या साहाय्याने $\fn{Fac}'$
सारख्या पदापासून क्रमगुणिताची व्याख्या मिळवता येते.

\begin{defn}
  \label{lam:ldf:fp:defn:Turing-Y}
  \emph{Y-संयोजक} हे पुढील पद आहे:
  \[  Y \ident (\lambd[ux][x(uux)])(\lambd[ux][x(uux)]).
  \]
\end{defn}

\begin{thm}
  $Y$ साठी कोणत्याही पद $g$ बाबत $Yg \red g(Yg)$; म्हणून
  $Yg$ हा नेहमी~$g$ चा स्थिरबिंदू असतो.
\end{thm}

\begin{proof}
  $(\lambd[ux][x(uux)])$ ला~$U$ असे संक्षेपाने
  लिहू. मग $Y \ident UU$. म्हणून
  \begin{align*}
    Y g &\ident (\lambd[ux][x(uux)])U\,g \\
    &\red (\lambd[x][x(UUx)])g\\
    & \red g(UUg) \ident g(Yg).
  \end{align*}
  $g(Yg)$ आणि $Yg$ या दोघांचेही $g(Yg)$ मध्ये
  न्यूनीकरण होते; म्हणून $g(Yg) \equal[\beta]
  Yg$ आणि $Yg$ हा~$g$ चा स्थिरबिंदू आहे.
\end{proof}

$Yg$ मध्ये रेडेक्स असल्याने न्यूनीकरण अनंतकाळ
पुढे चालू राहू शकते:
\begin{align*}
  Y g &\red g (Y g) \\
    &\red g (g (Y g)) \\
    &\red g(g (g (Y g)))\\
    &\ldots
\end{align*}
हा सांत पदांच्या अधिकाधिक उकलांचा अंतहीन अनुक्रम आहे;
तो स्वतः एक अनंत लॅम्डा पद नाही. म्हणून $Yg$ म्हणजे $g$
वारंवार लावलेल्या उकलांनी व्यक्त केलेला स्थिरबिंदू असे
सहज अर्थाने समजू शकतो. त्याला $g$ आणखी एकदा
लावल्याने, बोलायचे तर, काहीच जास्त होत नाही:
$Yg$ वर $g$ अनंत वेळा लावल्यावर पुन्हा $g$
लावले, तरी ते $Yg$ वर $g$ अनंत वेळा लावण्यासारखेच आहे.

वरील $Yg$ पासून सुरू होणारा $\beta$-न्यूनीकरणाचा
अनुक्रम अनंत आहे, हे लक्षात घ्या. त्यामुळे $Yg$
एखाद्या पदाला लागू केल्यावरचे $(Yg)N$ हे पदही
$(g(Yg))N$, $(g(g(Yg)))N$, \dots अशा अनंत
भिन्न पदांमध्ये न्यूनीत होऊ शकते. तरी न्यूनीकरणाचा
एखादा \emph{दुसरा} अनुक्रम सामान्य रूपात संपणे
शक्य आहे.

क्रमगुणिताचे उदाहरण घ्या. $\fn{Fac}$ ची व्याख्या
$Y\, \fn{Fac}'$ अशी करा (म्हणजे $\fn{Fac'}$
चा स्थिरबिंदू). मग:
\begin{align*}
  \fn{Fac}\,\num 3
  &\red Y\, \fn{Fac}' \, \num 3 \\
  &\red \fn{Fac}' (Y \,\fn{Fac}') \, \num 3 \\
  &\ident (\lambd[x][\lambd[n][\fn{IsZero} \, n\, \num 1\,
      (\fn{Mult}\, n\, (x (\fn{Pred}\, n)))]]) \, \fn{Fac} \, \num 3\\
  & \red \fn{IsZero} \, \num 3\, \num 1\,
      (\fn{Mult}\, \num 3\, (\fn{Fac} (\fn{Pred}\, \num 3))) \\
  &\red  \fn{Mult}\, \num 3 \, (\fn{Fac} \, \num 2).
  \intertext{त्याचप्रमाणे,}
  \fn{Fac}\,\num 2
  &\red  \fn{Mult}\, \num 2 \, (\fn{Fac} \, \num 1) \\
  \fn{Fac}\,\num 1
  &\red  \fn{Mult}\, \num 1 \, (\fn{Fac} \, \num 0)
  \intertext{पण}
  \fn{Fac}\,\num 0
  &\red \fn{Fac}' (Y \,\fn{Fac}') \, \num 0 \\
  &\ident (\lambd[x][\lambd[n][\fn{IsZero} \, n\, \num 1\,
      (\fn{Mult}\, n\, (x (\fn{Pred}\, n)))]]) \, \fn{Fac} \, \num 0\\
  & \red \fn{IsZero} \, \num 0\, \num 1\,
      (\fn{Mult}\, \num 0\, (\fn{Fac} (\fn{Pred}\, \num 0))).\\
  &\red \num 1.
  \intertext{म्हणून एकत्रितपणे}
  \fn{Fac}\,\num 3
  &\red  \fn{Mult}\, \num 3 \,
  (\fn{Mult} \, \num 2\, (\fn{Mult} \, \num 1\, \num 1)).
\end{align*}

$\fn{Fac'}$ साठी जे खरे आहे ते कोणत्याही
पुनरावर्ती व्याख्येसाठी खरे आहे. पुढील पुनरावर्ती
समीकरण दिले आहे असे समजा. पुनरावर्तनाचे नाव आणि
सर्व निविष्ट प्राचलांची नावे परस्पर भिन्न निवडायची आहेत:
\begin{align*}
g\,x_1\dots x_n & \equal[\beta] N
\intertext{येथे $N$ मध्ये~$g$ आणि $x_1$, \dots,~$x_n$
असू शकतात. मग असे पद~$G \ident (Y \lambd[g][\lambd[x_1\dots x_n][N]])$
नेहमी असते की}
G\,x_1 \dots x_n & \equal[\beta] \Subst{N}{G}{g}.
\intertext{कारण स्थिरबिंदू प्रमेयानुसार,}
G \ident (Y \lambd[g][\lambd[x_1\dots x_n][N]]) & \red \lambd[g][\lambd[x_1\dots x_n][N]](Y \lambd[g][\lambd[x_1\dots x_n][N]])\\
& \ident (\lambd[g][\lambd[x_1\dots x_n][N]])\,G
\intertext{आणि परिणामी,}
G\,x_1 \dots x_n
& \red (\lambd[g][\lambd[x_1\dots x_n][N]])\,G\,x_1\dots x_n\\
& \red (\lambd[x_1\dots x_n][\Subst{N}{G}{g}])\,x_1\dots x_n\\
  & \red \Subst{N}{G}{g}.
\end{align*}

\ref{lam:ldf:fp:defn:Turing-Y} मधील $Y$-संयोजक
ॲलन ट्यूरिंगने दिला होता. ॲलॉन्झो चर्चने
दुसरे रूप सुचवले होते; त्याला~$Y_C$ म्हणू:
\[Y_C \ident \lambd[g][(\lambd[x][g(xx)])(\lambd[x][g(xx)])].
\]
चर्चचा संयोजक ट्यूरिंगच्या संयोजकापेक्षा काहीसा
कमकुवत आहे. औपचारिक चर g वापरल्यास $Y_Cg \equal[\beta] g(Y_Cg)$, पण
$Y_Cg \bred g(Y_Cg)$ नाही. $V$ हे
$\lambd[x][g(xx)]$ पद असू द्या; येथे त्याचे बद्ध नाव
आदेशित पदाच्या मुक्त चरांपासून वेगळे निवडायचे आहे. मग
$Y_C \ident \lambd[g][VV]$. त्यामुळे
\begin{align*}
  VV & \ident (\lambd[x][g(xx)])V \red g(VV)
  \text{ आणि त्यामुळे}\\
  Y_C g & \ident (\lambd[g][VV])g \red VV \red g(VV), \text{ पण तसेच}\\
  g(Y_C g) & \ident g((\lambd[g][VV])g) \red g(VV).
\end{align*}
दुसऱ्या शब्दांत, $Y_Cg$ आणि $g(Y_Cg)$ यांचे
न्यूनीकरण होऊन एकच पद $g(VV)$ मिळते;
म्हणून $Y_Cg \equal[\beta] g(Y_Cg)$. अनेक
उपयोजनांसाठी एवढे पुरेसे असते.











\section{लघुतमीकरण}\label{lam:ldf:min:sec}

मूलभूत फलने $\Zero$, $\Succ$, $\Proj{n}{i}$
यांपासून संयोजन, आदिम पुनरावर्तन आणि नियमित
लघुतमीकरण वापरून मिळणारी फलने म्हणजे सामान्य
पुनरावर्ती फलने. सर्व सामान्य पुनरावर्ती फलने
लॅम्डा-परिभाष्य आहेत हे दाखवण्यासाठी,
लॅम्डा-परिभाष्य फलनापासून नियमित लघुतमीकरणाने
परिभाषित केलेले कोणतेही फलनही लॅम्डा-परिभाष्य
आहे हे दाखवावे लागेल.

\begin{lem}
  \label{lam:ldf:min:lem:min} $f(x_1, \dots, x_k, y)$ हे सर्वत्र परिभाषित आणि नियमित
  आणि लॅम्डा-परिभाष्य असेल, तर
  \[  h(x_1, \dots, x_k) = \umin{y}{f(x_1,\dots,x_k, y) = 0}
  \]
  याने परिभाषित केलेले $h$~देखील
  लॅम्डा-परिभाष्य आहे.
\end{lem}

\begin{proof}
  लॅम्डा पद~$F$ हे नियमित फलन $f(\vec x, y)$ ला
  $\lambda$-परिभाषित करते असे समजा. $h$ ला
  लॅम्डा-परिभाषित करण्यासाठी शोधफलन आणि
  स्थिरबिंदू संयोजक वापरू:
  \begin{align*}
    \fn{Search} & \ident \lambd[g][\lambd[f\,\vec{x}\,y][
        \fn{IsZero} (f\, \vec{x}\, y)\, y\, (g\, f\, \vec{x} (\fn{Succ}\, y))]]\\
    H & \ident \lambd[\vec x][(Y \, \fn{Search}) F\, \vec{x}\, \num{0}],
  \end{align*}
  येथे $Y$ हा \ref{lam:ldf:fp:defn:Turing-Y} मधील
  ट्यूरिंग स्थिरबिंदू संयोजक आहे. Recursive call मध्ये
  predicate प्रतिनिधी पुन्हा पहिला फलसाधक म्हणून द्यायचा
  आहे; तो वगळल्यास पुढील calls चे सर्व फलसाधक सरकतात.
  सहज अर्थाने, $\fn{Search}$ हे स्वतःला पुन्हा
  बोलावणारे फलन आहे: $y$ पासून सुरू करून
  $f\, \vec x\, y$ शून्य आहे का हे तपासते;
  शून्य असेल तर~$y$ परत करते, अन्यथा
  $\fn{Succ}\, y$ घेऊन स्वतःला पुन्हा बोलावते.
  म्हणून $(Y \, \fn{Search}) F
  \num{n_1}\dots\num{n_k}\,\num{0}$ हे
  $f(n_1, \dots, n_k, m) = 0$ पूर्ण करणारा
  लघुतम~$m$ परत करते.

  ट्यूरिंग संयोजकाच्या थेट उकलीने आणि परिभाषित
  प्रत्येक predicate चाचणीच्या न्यूनीकरणाने पुढील गोष्ट मिळते:
  \begin{align*}
    (Y \, \fn{Search}) F \num{n_1}
    \dots\num{n_k}\, \num{m} & \red \num{m}
    \intertext{जर $f(n_1, \dots,
      n_k, m) = 0$ असेल; अन्यथा}
    & \red (Y \, \fn{Search}) F\, \num{n_1} \dots\num{n_k}\, \num{m+1}
    \intertext{मिळते. $f$ नियमित असल्याने काही $y$ साठी
      $f(n_1, \dots, n_k, y) = 0$ असते; म्हणून}
    (Y \, \fn{Search}) F \num{n_1} \dots\num{n_k}\,\num{0}
    & \red \num{h(n_1, \dots, n_k)}.
    \end{align*}
\end{proof}

\begin{prop}
  प्रत्येक सामान्य पुनरावर्ती फलन लॅम्डा-परिभाष्य आहे.
\end{prop}

\begin{proof}
 \ref{lam:ldf:prf:lem:basic} नुसार सर्व मूलभूत फलने
 लॅम्डा-परिभाष्य आहेत. तसेच \ref{lam:ldf:prf:lem:comp},
 \ref{lam:ldf:prf:lem:prim} आणि \ref{lam:ldf:min:lem:min} नुसार
 सर्वत्र परिभाषित लॅम्डा-परिभाष्य फलनांचा वर्ग संयोजन,
 आदिम पुनरावर्तन आणि नियमित लघुतमीकरण यांच्या
 अंतर्गत बंद असतो.
\end{proof}











\section{आंशिक पुनरावर्ती फलने लॅम्डा-परिभाष्य असतात}\label{lam:ldf:par:sec}

मूलभूत फलनांपासून संयोजन, आदिम पुनरावर्तन आणि
अपरिबद्ध लघुतमीकरण वापरून मिळणारी फलने म्हणजे
आंशिक पुनरावर्ती फलने. सामान्य पुनरावर्ती फलनांपेक्षा
त्यांचा फरक असा की अपरिबद्ध शोधात वापरलेली फलने
नियमित असणे आवश्यक नाही. नियमिततेची अट नसल्याने
लघुतमीकरणाने परिभाषित केलेली फलने काही वेळा
अपरिभाषित राहू शकतात.

प्रथमदर्शनी, (सर्वत्र परिभाषित) सामान्य पुनरावर्ती
फलने लॅम्डा-परिभाष्य आहेत हे दाखवण्याच्या
पद्धतीनेच सर्व आंशिक पुनरावर्ती फलनेही
लॅम्डा-परिभाष्य आहेत हे सिद्ध करता येईल असे
वाटू शकते. उदाहरणार्थ, $f$ आणि $g$ ही अनुक्रमे
$F$ आणि~$G$ या पदांनी लॅम्डा-परिभाषित
असतील, तर $f$ चे $g$ शी संयोजन
$\lambd[x][F (G x)]$ ने लॅम्डा-परिभाषित होते,
जर दोन्ही फलने सर्वत्र परिभाषित असतील.
पण फलने आंशिक असतील, तेव्हा अडचण येते.
$g(x)$ अपरिभाषित असल्यास $G x$ ला सामान्य रूप
नसते. बहुतेक वेळा मग $F (G x)$ लाही सामान्य रूप
नसते, आणि आपल्याला तेच हवे आहे. पण
$F$ हे $\lambd[x][\num{0}]$ असेल, तर
$F (G x)$ ला $\num{0}$ हे सामान्य रूप मिळते.

ही अडचण दूर करता येते. सर्व आंशिक पुनरावर्ती
फलने अशी लॅम्डा-परिभाषित करण्याचे मार्ग आहेत
की अपरिभाषित मूल्ये सामान्य रूप नसलेल्या पदांनी
दर्शवली जातात. मात्र हे मार्ग सामान्य पुनरावर्ती
फलनांसाठी वापरलेल्या पद्धतीपेक्षा काहीसे अधिक
गुंतागुंतीचे आणि कमी सहज आहेत. येथे पुरावा न देता
प्रमेय नोंदवतो. या आवृत्तीच्या आधीच्या परिचयात
\ref{lam:int:lrp:thm:computable-lambda} ची दुरुस्त
सिद्धता आणि लघुतमीकरणातील forced-search extraction
यांनी ही construction आधीच पूर्ण केली आहे:

\begin{thm}
  सर्व आंशिक पुनरावर्ती फलने लॅम्डा-परिभाष्य आहेत.
\end{thm}











\section{लॅम्डा-परिभाष्य फलने पुनरावर्ती असतात}\label{lam:dfl:ldr:sec}

सर्व आंशिक पुनरावर्ती फलने लॅम्डा-परिभाष्य
आहेत एवढेच नव्हे; याचे व्यस्त विधानही खरे आहे.
म्हणजे सर्व लॅम्डा-परिभाष्य फलने आंशिक
पुनरावर्ती असतात.

\begin{thm}
  \label{lam:dfl:ldr:thm:lambda-computable} आंशिक फलन $f$
  हे लॅम्डा-परिभाष्य असेल, तर ते आंशिक
  पुनरावर्ती असते.
\end{thm}

\begin{proof}
  पुराव्याची फक्त रूपरेषा देऊ. आधी $\lambd$-पदांचे
  अंकगणितीकरण करू: अनुक्रमांच्या नेहमीच्या अभाज्य
  संख्यांच्या घातांवरील संकेतनाचा वापर करून
  $\lambd$-पदांना पद्धतशीरपणे गोडेल संख्या देऊ.
  मग आंशिक पुनरावर्ती फलन $\fn{normalize}(t)$
  परिभाषित करू. ते लॅम्डा पदाची गोडेल संख्या~$t$
  फलसाधक म्हणून घेते आणि पदाला सामान्य रूप
  असेल, तर त्या रूपाची गोडेल संख्या परत करते;
  अन्यथा ते अपरिभाषित राहते. हे फलन प्रभावी आहे:
  प्रत्येक पायरीत डावीकडील सर्वप्रथम बीटा रेडेक्स
  निवडून, capture टाळण्यासाठी ठरावीक क्रमाने ताजी
  बद्ध नावे घेऊन, त्याचे संकुचन करायचे. सामान्य रूप
  असल्यास ही पद्धत तिथे पोहोचते; अन्यथा शोध संपत नाही.
  अवैध input codes वरही फलन अपरिभाषित ठेवू. पुढे
  $\fn{toChurch}$ आणि $\fn{fromChurch}$ ही दोन
  आंशिक पुनरावर्ती फलने परिभाषित करू. ती नैसर्गिक
  संख्या आणि त्या संख्यांच्या अनुरूप चर्च अंकांच्या
  गोडेल संख्या यांमध्ये पुढे व मागे रूपांतर करतात.
  पहिले रूपांतर सर्व नैसर्गिक संख्यांवर परिभाषित आहे.
  दुसरे रूपांतर फक्त चर्च अंकाच्या सामान्य रूपाचा वैध
  code दिल्यास परिभाषित आहे; alpha-सममूल्य बद्ध
  नावांची निवड बदलली, तरी अंक तोच ओळखायचा.

  या पुनरावर्ती फलनांचा वापर करून फलन~$f$
  आंशिक पुनरावर्ती म्हणून परिभाषित करता येते.
  $f$ ला लॅम्डा-परिभाषित करणारे $\lambd$-पद~$F$
  आहे. $f(n_1, \dots, n_k)$ संगणित करण्यासाठी
  आधी $\fn{toChurch}(n_i)$ वापरून अनुरूप चर्च
  अंकांच्या गोडेल संख्या मिळवा. $\Gn{F}$ आणि या
  codes पासून संकेतनातील प्रभावी उपयोजन रचनाकार
  वापरून $F \num{n_1}\dots\num{n_k}$ या पदाची
  गोडेल संख्या मिळवा; येथे साधी numeric बेरीज किंवा
  codes चे अंध जोडणे अभिप्रेत नाही. आता या गोडेल संख्येवर
  $\fn{normalize}$ वापरा. $f(n_1, \dots, n_k)$
  परिभाषित असेल, तर
  $F \num{n_1}\dots\num{n_k}$ ला सामान्य रूप
  असते (ते चर्च अंकच असले पाहिजे); अन्यथा सामान्य
  रूप नसते (आणि म्हणून
  \[\fn{normalize}(\Gn{F\num{n_1}\dots\num{n_k}})\]
  हे अपरिभाषित असते). शेवटी सामान्यीकृत पदाच्या
  गोडेल संख्येवर $\fn{fromChurch}$ वापरा.
\end{proof}



\part{बहुमूल्य तर्कशास्त्र}\label{mvl:part}
\begin{editorial}
  या भागात विधानीय बहुमूल्य तर्कशास्त्रांवरील मसुदा साहित्य आहे.
\end{editorial}
\chapter{विन्यासमीमांसा आणि चिन्हार्थमीमांसा}\label{mvl:syn:chap}









\section{परिचय}\label{mvl:syn:int:sec}

अभिजात तर्कशास्त्रात $\lnot$, $\land$, $\lor$, $\lif$ आणि
$\liff$ ही विधानीय संयोजके वापरून विधानीय चले पासून
घडवलेल्या सूत्रांचा विचार करतो. अभिजात तर्कशास्त्राची
चिन्हार्थमीमांसा परिभाषित करताना आपण $\True$ आणि $\False$ ही दोन
सत्यतामूल्ये वापरतो. सत्य-मूल्यांकन~$\pAssign{v}$ मध्ये
विधानीय चले ना अर्थ देतो; ते विधानीय चले ना $\True$ आणि
$\False$ ही सत्यतामूल्ये नेमून देते. मग कोणतेही सत्य-मूल्यांकन
कोणत्याही सूत्र~$\varphi$ चे $\pValue{v}(\varphi)$ हे सत्यतामूल्य
ठरवते; आणि सूत्राची सत्य-मूल्यांकन~$\pAssign{v}$ मध्ये
पूर्ति होते, $\pSat{v}{\varphi}$, जर आणि फक्त जर
$\pValue{v}(\varphi) = \True$.

बहुमूल्य तर्कशास्त्रे $\True$ आणि $\False$ यांपलीकडील
सत्यतामूल्यांना परवानगी देऊन अभिजात द्विमूल्य तर्कशास्त्राचे
सामान्यीकरण करतात. म्हणून बहुमूल्य तर्कशास्त्रात
सत्य-मूल्यांकन~$\pAssign{v}$ हे प्रत्येक विधानीय चल~$p$
ला संभाव्य सत्यतामूल्यांपैकी एक मूल्य नेमून देणारे फलन असते.
अनुमत सत्यतामूल्यांच्या संचाला आपण साधारणपणे~$V$ म्हणू.
$V = \{\True, \False\}$ असलेले अभिजात तर्कशास्त्र हेही
बहुमूल्य तर्कशास्त्र आहे; त्यात $\pValue{v}(\varphi)$ हे सत्यतामूल्य
संयोजकांच्या परिचित वैशिष्ट्यपूर्ण सत्यताकोष्टकांनुसार ठरते.

अधिक सत्यतामूल्ये समाविष्ट केल्यावर ``आणि,'' ``किंवा,'' ``नाही,''
आणि ``जर---तर'' असे वाचत असलेल्या संयोजकांसाठी
$\pValue{v}(\varphi)$ कसे काढायचे, याचे एकाहून अधिक स्वाभाविक
पर्याय असतात. म्हणून बहुमूल्य तर्कशास्त्र केवळ सत्यतामूल्यांच्या
संचानेच नव्हे, तर प्रत्येक संयोजकासाठी निवडलेल्या
\emph{सत्यता-फलनांनी}ही निश्चित होते. तथापि, तर्कशास्त्र~$\Log L$
साठी ही निवड झाल्यावर, अभिजात तर्कशास्त्राप्रमाणेच
$\pValue{v}(\varphi)[\Log L]$ हे सत्यतामूल्य त्या सत्य-मूल्यांकनाने
एकमेव रीतीने ठरते. अभिजात तर्कशास्त्राप्रमाणे बहुमूल्य
तर्कशास्त्रेही \emph{सत्यता-फलनात्मक} आहेत.

हे चिन्हार्थविषयक पायाभूत घटक मिळाल्यावर सर्वतः सत्यता,
तार्किक निष्पन्नता आणि पूर्ततायोग्यता यांच्या समकक्ष
चिन्हार्थविषयक संकल्पना परिभाषित करता येतात. अभिजात
तर्कशास्त्रात सूत्र सर्वतः सत्य असते, जर
कोणत्याही~$\pAssign{v}$ साठी त्याचे सत्यतामूल्य
$\pValue{v}(\varphi) = \True$ असेल. बहुमूल्य तर्कशास्त्रात ही
संकल्पना थोडी व्यापक करावी लागते. प्रथम, सत्यतामूल्यांच्या~$V$
संचात $\True$ असणे आवश्यक नाही. उदाहरणार्थ,
काही बहुमूल्य तर्कशास्त्रे $0$ आणि~$1$ यांदरम्यानच्या सर्व
परिमेय संख्या सत्यतामूल्ये म्हणून वापरतात. अशा वेळी $1$
सहसा $\True$ ची भूमिका बजावते. इतर तर्कशास्त्रांत
एकच नव्हे, तर अनेक सत्यतामूल्ये ती भूमिका बजावतात.
म्हणून प्रत्येक बहुमूल्य तर्कशास्त्रात \emph{निर्दिष्ट
सत्यतामूल्यां}चा संच~$V^+$ असावा अशी अट घालतो.
मग सूत्राची सत्य-मूल्यांकन~$\pAssign{v}$ मध्ये
पूर्ति होते, $\pSat{v}{\varphi}[\Log L]$, जर आणि फक्त जर
$\pValue{v}(\varphi)[\Log L] \in V^+$. सूत्र~$\varphi$
हे त्या तर्कशास्त्रातील सर्वतः सत्य सूत्र असते,
$\Entails[\Log L] \varphi$, जर आणि फक्त जर
कोणत्याही~$\pAssign{v}$ साठी $\pValue{v}(\varphi) \in V^+$.
शेवटी, सूत्रांच्या संचापासून $\varphi$ निष्पन्न होते,
$\Gamma \Entails[\Log L] \varphi$, जर~$\Gamma$ मधील सर्व
सूत्रांची पूर्ति करणारे प्रत्येक सत्य-मूल्यांकन~$\varphi$
चीही पूर्ति करत असेल.












\section{भाषा आणि संयोजके}\label{mvl:syn:con:sec}

अभिजात विधानीय तर्कशास्त्रात आणि इतर अनेक तर्कशास्त्रांत
\emph{विधानीय स्थिर} व \emph{संयोजक} यांचा संच वापरला जातो.
उदाहरणार्थ, पुढील चिन्हे आपण आदिम मानतो:
\begin{enumerate}
\item असत्यता दर्शवणारा विधानीय स्थिर~$\lfalse$.
\item सत्यता दर्शवणारा विधानीय स्थिर~$\ltrue$.
\item तार्किक संयोजके:
  \startycommalist
  \ycomma $\lnot$ (नकारण)
  \ycomma $\land$ (संधियोग)
  \ycomma $\lor$ (विकल्पयोग)
  \ycomma $\lif$ (शर्त)
  \ycomma $\liff$ (द्विशर्त)
\end{enumerate}


बहुमूल्य तर्कशास्त्रांतही हीच संयोजके वापरली जातात. तथापि,
उदाहरणार्थ संधियोगाच्या वेगवेगळ्या आवृत्त्या एकाच तर्कशास्त्रात
असणे अनेकदा उपयुक्त ठरते; त्या आवृत्त्या वेगळ्या दाखवण्यासाठी
भिन्न चिन्हे लागतात. काही बहुमूल्य तर्कशास्त्रांत अभिजात
तर्कशास्त्रात समतुल्य संयोजक नसलेली संयोजकेही असतात.
म्हणून आपण नेहमीपेक्षा थोडी अधिक सामान्य चौकट घेऊ.

\begin{defn}
  \emph{विधानीय भाषा} ही \emph{संयोजकां}च्या $\Lang L$
  या संचाने बनते. प्रत्येक संयोजक~$\star$ ची \emph{स्थानसंख्या}
  असते; स्थानसंख्या~$n$ असलेला संयोजक \emph{$n$-स्थानी}
  म्हणतात. स्थानसंख्या~$0$ असलेल्या संयोजकांना
  \emph{स्थिर} असेही म्हणतात; स्थानसंख्या~$1$ असलेल्या
  संयोजकांना \emph{एकस्थानी}, आणि स्थानसंख्या~$2$ असलेल्या
  संयोजकांना \emph{द्विस्थानी} म्हणतात.
\end{defn}

\begin{ex}
  विधानीय तर्कशास्त्राच्या मानक भाषेत~$\Lang L_0$ पुढील
  संयोजके (त्यांच्या स्थानसंख्यांसह) असतात:
  $\lfalse$~($0$),
  $\lnot$~($1$),
  $\land$~($2$),
  $\lor$~($2$),
  $\lif$~($2$). आपण विचारात घेणारी बहुतेक तर्कशास्त्रे
  ही भाषा वापरतात. काही तर्कशास्त्रे प्रस्थापित संकेतांनुसार
  काही संयोजकांसाठी वेगळी चिन्हे वापरतात. उदाहरणार्थ,
  गुणाकार तर्कशास्त्रात संधियोगाचे चिन्ह $\land$ ऐवजी
  अनेकदा~$\odot$ असते. कधी नवीन कारक जोडणेही सोयीचे
  ठरते; उदाहरणार्थ, $\triangle$ हा निश्चितता कारक
  ($1$-स्थानी).
\end{ex}












\section{सूत्र}\label{mvl:syn:fml:sec}

\begin{defn}[सूत्र]
\label{mvl:syn:fml:defn:formulas}
विधानीय भाषा~$\Lang L$ मधील \emph{सूत्रे} चा संच~$\Frm[L]$
पुढीलप्रमाणे विगमनाने परिभाषित करतो:
\begin{enumerate}
\item प्रत्येक विधानीय चल~$\Obj p_i$ हे आण्विक
  सूत्र असते.
\item $\Lang L$ मधील प्रत्येक $0$-स्थानी संयोजक (विधानीय स्थिर)
  हे आण्विक सूत्र असते.
\item $\star$ हा $\Lang L$ मधील सकारात्मक स्थानसंख्येचा
  $n$-स्थानी संयोजक असेल आणि
  $\varphi_1$, \dots, $\varphi_n$ ही सूत्रे असतील, तर
  $\star(\varphi_1, \dots, \varphi_n)$ हे सूत्र असते.
\item यांखेरीज काहीही सूत्र नाही.
\end{enumerate}
$\star$ $1$-स्थानी असेल, तर $\star(\varphi_1)$ हे अनेकदा फक्त
$\star \varphi_1$ असे लिहितात. $\star$ $2$-स्थानी असेल, तर
$\star(\varphi_1,\varphi_2)$ हे अनेकदा $(\varphi_1 \star \varphi_2)$ असे लिहितात.
\end{defn}

नेहमीप्रमाणे, सर्वांत बाहेरचे कंस आपण अनेकदा न सांगता वगळू.

\begin{ex}
  मानक भाषेत~$\Lang{L_0}$,
  $\Obj p_1 \lif (\Obj p_1 \land \lnot \Obj p_2)$ हे सूत्र आहे.
  गुणाकार तर्कशास्त्राच्या भाषेत ते याऐवजी
  $\Obj p_1 \lif (\Obj p_1 \odot \lnot \Obj p_2)$ असे लिहितात.
  भाषेत $1$-स्थानी $\triangle$ जोडल्यास,
  $\triangle (\Obj p_1 \land \Obj p_2) \lif (\triangle \Obj p_1 \land \triangle \Obj p_2)$
  यांसारखी सूत्रेही मिळतात.
\end{ex}












\section{मॅट्रिक्स}\label{mvl:syn:mat:sec}

बहुमूल्य तर्कशास्त्र त्याची भाषा, सत्यतामूल्यांचा~$V$ संच,
निर्दिष्ट सत्यतामूल्यांचा उपसंच आणि त्याच्या संयोजकांसाठीची
सत्यता-फलने यांवरून परिभाषित होते. या घटकांना एकत्रितपणे
\emph{मॅट्रिक्स} म्हणतात.

\begin{defn}[मॅट्रिक्स]
\label{mvl:syn:mat:defn:matrix}
तर्कशास्त्र~$\Log L$ साठीच्या \emph{मॅट्रिक्स} मध्ये पुढील घटक असतात:
\begin{enumerate}
\item भाषा~$\Lang L$ बनवणारा संयोजकांचा संच;
\item सत्यतामूल्यांचा $V \neq \emptyset$ संच;
\item निर्दिष्ट सत्यतामूल्यांचा $V^+ \subseteq V$ संच;
\item $\Lang L$ मधील प्रत्येक $n$-स्थानी संयोजक~$\star$ साठी
  सत्यता-फलन~$\tf{\star} : V^n \to V$. $n = 0$ असल्यास
  $\tf{\star}$ हे फक्त~$V$ मधील एक मूल्य असते.
\end{enumerate}
\end{defn}

\begin{ex}
अभिजात तर्कशास्त्र~\LogCL{} साठीच्या मॅट्रिक्समध्ये पुढील घटक असतात:
\begin{enumerate}
  \item $\lfalse$, $\lnot$, $\land$, $\lor$, $\lif$ असलेली
  मानक विधानीय भाषा~$\Lang L_0$.
  \item सत्यतामूल्यांचा संच $V = \{\True, \False\}$.
  \item $\True$ हेच एकमेव निर्दिष्ट मूल्य, म्हणजे $V^+ = \{\True\}$.
  \item $\lfalse$ साठी $\tf{\lfalse} = \False$. इतर सत्यता-फलने
  नेहमीच्या सत्यताकोष्टकांनी दिली आहेत (पहा
  \ref{mvl:syn:mat:fig:tf-CL}).
\end{enumerate}
\begin{figure}
  \begin{center}
      \begin{tabular}{c|c}
        $\tf{\lnot}$ & \\
        \hline
        $\True$ & $\False$ \\
        $\False$ & $\True$
      \end{tabular}
      \quad
      \begin{tabular}{c|cc}
        $\tf{\land}$ & $\True$ & $\False$ \\
        \hline
        $\True$ & $\True$ & $\False$ \\
        $\False$ & $\False$ & $\False$
      \end{tabular}
      \quad
      \begin{tabular}{c|cc}
        $\tf{\lor}$ & $\True$ & $\False$ \\
        \hline
        $\True$ & $\True$ & $\True$ \\
        $\False$ & $\True$ & $\False$
      \end{tabular}
      \quad
      \begin{tabular}{c|cc}
        $\tf{\lif}$ & $\True$ & $\False$ \\
        \hline
        $\True$ & $\True$ & $\False$ \\
        $\False$ & $\True$ & $\True$
      \end{tabular}
    \end{center}
    \caption{अभिजात तर्कशास्त्र~$\LogCL$ साठीची सत्यता-फलने.}
    \label{mvl:syn:mat:fig:tf-CL}
  \end{figure}
\end{ex}












\section{सत्य-मूल्यांकन आणि पूर्ति}\label{mvl:syn:val:sec}

\begin{defn}[सत्य-मूल्यांकने]
$V$ हा सत्यतामूल्यांचा संच असू द्या. $\Lang{L}$ साठी~$V$
मध्ये जाणारे \emph{सत्य-मूल्यांकन} म्हणजे भाषेतील
विधानीय चले ना~$V$ मधील घटक नेमून
देणारे फलन~$\pAssign{v}$ होय; म्हणजेच
$\pAssign{v} \colon \PVar \to V$.
\end{defn}

\begin{defn}\label{mvl:syn:val:defn:pValue}
  बहुमूल्य तर्कशास्त्र~$\Log L$ च्या सत्यतामूल्यांच्या~$V$
  संचात जाणारे सत्य-मूल्यांकन~$\pAssign{v}$ दिले असता,
  मूल्यन फलन $\pValue{v} \colon \Frm[L] \to V$
  पुढीलप्रमाणे विगमनाने परिभाषित करतो:
  \begin{enumerate}
    \item $\pValue{v}(\Obj p_n) = \pAssign{v}(\Obj p_n)$;
    \item $\star$ हा $0$-स्थानी संयोजक असेल, तर
    $\pValue{v}(\star)  = \tf{\star}[\Log L]$;
    \item $\star$ हा सकारात्मक स्थानसंख्येचा
    $n$-स्थानी संयोजक असेल, तर
    \[      \pValue{v}(\star(\varphi_1, \dots, \varphi_n)) = \tf{\star}[\Log L]
      (\pValue{v}(\varphi_1), \dots, \pValue{v}(\varphi_n)).
    \]
  \end{enumerate}
\end{defn}

\begin{defn}[पूर्ति]
\label{mvl:syn:val:defn:satisfaction} सूत्र~$\varphi$ ची
  सत्य-मूल्यांकन~$\pAssign{v}$ ने \emph{पूर्ति होते},
  $\pSat{v}{\varphi}[\Log L]$, जर आणि फक्त जर
  $\pValue{v}(\varphi)[\Log L] \in V^+$; येथे $V^+$ हा
  तर्कशास्त्र~$\Log L$ मधील निर्दिष्ट सत्यतामूल्यांचा संच आहे.

  ``$\pSat{v}{\varphi}[\Log L]$ नाही'' हे दर्शवण्यासाठी आपण
  $\pSat/{v}{\varphi}[\Log L]$ असे लिहितो. $\Gamma$ हा
  सूत्रांचा संच असेल, तर $\pSat{v}{\Gamma}[\Log L]$
  जर आणि फक्त जर प्रत्येक~$\varphi \in \Gamma$ साठी
  $\pSat{v}{\varphi}[\Log L]$.
\end{defn}












\section{चिन्हार्थविषयक संकल्पना}\label{mvl:syn:sem:sec}

बहुमूल्य तर्कशास्त्र~$\Log L$ एखाद्या मॅट्रिक्सने दिलेले
आहे असे समजू. मग~$\Log L$ साठी नेहमीच्या चिन्हार्थविषयक
संकल्पना परिभाषित करता येतात.

\begin{defn}
\begin{enumerate}
\item सूत्र~$\varphi$ \emph{पूर्ततायोग्य} असते, जर
  एखाद्या~$\pAssign{v}$ साठी $\pSat{v}{\varphi}$;
  कोणत्याही $\pAssign{v}$ साठी $\pSat{v}{\varphi}$ नसेल, तर ते
  \emph{अपूर्ततायोग्य} असते;
\item सूत्र~$\varphi$ हे \emph{सर्वतः सत्य सूत्र}
  असते, जर सर्व सत्य-मूल्यांकने~$v$ साठी $\pSat{v}{\varphi}$;
\item $\Gamma$ हा सूत्रांचा संच असेल, तर
  $\Gamma \Entails \varphi$ (``$\Gamma$ पासून $\varphi$
  निष्पन्न होते'') जर आणि फक्त जर
  $\pSat{v}{\Gamma}$ असलेल्या प्रत्येक
  सत्य-मूल्यांकन~$\pAssign{v}$ साठी $\pSat{v}{\varphi}$;
\item $\Gamma$ हा सूत्रांचा संच असेल, तर
  $\pSat{v}{\Gamma}$ असलेले सत्य-मूल्यांकन~$\pAssign{v}$
  अस्तित्वात असल्यास $\Gamma$ \emph{पूर्ततायोग्य} असतो;
  अन्यथा $\Gamma$ \emph{अपूर्ततायोग्य} असतो.
\end{enumerate}
\end{defn}

अभिजात तर्कशास्त्रात जशी काही तथ्ये या संकल्पनांसाठी
लागू होतात, तशीच ती येथेही लागू होतात:

\begin{prop}
\label{mvl:syn:sem:prop:semanticalfacts}
\begin{enumerate}
\item $\varphi$ हे सर्वतः सत्य सूत्र असते, जर आणि फक्त जर
  $\emptyset \Entails \varphi$;
\item $\Gamma$ पूर्ततायोग्य असेल, तर $\Gamma$ चा प्रत्येक सांत
  उपसंचही पूर्ततायोग्य असतो;
\item\label{mvl:syn:sem:def:monotonicity}
एकस्वनिकता: $\Gamma \subseteq \Delta$ आणि
  $\Gamma \Entails \varphi$ असल्यास $\Delta \Entails \varphi$ देखील;
\item\label{mvl:syn:sem:def:Cut}
संक्रमकता: $\Gamma \Entails \varphi$ आणि
  $\Delta \cup \{ \varphi\} \Entails \psi$ असल्यास
  $\Gamma \cup \Delta \Entails \psi$;
\end{enumerate}
\end{prop}

\begin{proof}
सराव.
\end{proof}

\begin{prob}
\ref{mvl:syn:sem:prop:semanticalfacts} सिद्ध करा.
\end{prob}

अभिजात तर्कशास्त्रात निष्पन्नता आणि सशर्त संयोजक यांचा संबंध
जोडता येतो. उदाहरणार्थ, \emph{विध्यात्मक अनुमान} वैध असते:
जर $\Gamma \Entails \varphi$ आणि $\Gamma \Entails \varphi \lif \psi$
असेल, तर $\Gamma \Entails \psi$. अभिजात तर्कशास्त्रात
$\Entails$ आणि $\lif$ यांच्यातील दुसरा महत्त्वाचा संबंध म्हणजे
चिन्हार्थविषयक निगमन प्रमेय: $\Gamma \Entails \varphi \lif \psi$
जर आणि फक्त जर $\Gamma \cup \{\varphi\} \Entails \psi$.
हे निष्कर्ष बहुमूल्य तर्कशास्त्रांत \emph{नेहमीच लागू होतात असे नाही}.
ते लागू होणे $\tf{\lif}$ या सत्यता-फलनावर अवलंबून असते.












\section[अभिजात तर्कशास्त्राची उपतर्कशास्त्रे]{अभिजात तर्कशास्त्र~$\LogCL$ ची उपतर्कशास्त्रे म्हणून बहुमूल्य तर्कशास्त्रे}\label{mvl:syn:sub:sec}

नेहमीची बहुमूल्य तर्कशास्त्रे अशा मॅट्रिक्सांनी परिभाषित होतात की
$\{\True, \False\}$ मधील आदानांवर त्यांच्या सत्यता-फलनांची मूल्ये
अभिजात सत्यता-फलनांशी जुळतात. नेमके म्हणजे, या तर्कशास्त्रांत
$x \in \{\True, \False\}$ असल्यास $\tf{\lnot}[\Log L](x) =
\tf{\lnot}[\LogCL](x)$; आणि $\star$ हे $\land$, $\lor$,
$\lif$ यांपैकी कोणतेही असेल व $x, y \in \{\True, \False\}$
असतील, तर $\tf{\star}[\Log L](x,y) =
\tf{\star}[\LogCL](x,y)$. दुसऱ्या शब्दांत,
$\lnot$, $\land$, $\lor$, $\lif$ यांची सत्यता-फलने
$\{\True,\False\}$ वर मर्यादित केली असता ती नेमकी
अभिजात सत्यता-फलनेच असतात.

\begin{prop}\label{mvl:syn:sub:prop:mvl-cl}
  बहुमूल्य तर्कशास्त्र~$\Log L$ च्या भाषेत
  $\lnot$, $\land$, $\lor$, $\lif$ ही संयोजके आहेत,
  $\True, \False \in V$ आहेत आणि त्याची सत्यता-फलने
  पुढील अटी पूर्ण करतात असे समजा:
  \begin{enumerate}
    \item\label{mvl:syn:sub:prop:not} $x = \True$ किंवा $x = \False$
    असल्यास $\tf{\lnot}[\Log L](x) = \tf{\lnot}[\LogCL](x)$;
    \item\label{mvl:syn:sub:prop:land} $\tf{\land}[\Log L](x,y) = \tf{\land}[\LogCL](x,y)$,
    \item\label{mvl:syn:sub:prop:lor} $\tf{\lor}[\Log L](x,y) = \tf{\lor}[\LogCL](x,y)$,
    \item\label{mvl:syn:sub:prop:lif} $\tf{\lif}[\Log L](x,y) = \tf{\lif}[\LogCL](x,y)$,
    जेव्हा $x, y \in \{\True, \False\}$.
  \end{enumerate}
  मग $V$ मध्ये जाणाऱ्या कोणत्याही सत्य-मूल्यांकन~$\pAssign v$
  साठी, प्रत्येक विधानीय चलाबाबत $\pAssign
  v(p) \in \{\True,\False\}$ असल्यास,
  या चार संयोजकांनी आणि विधानीय चलांनीच बनलेल्या प्रत्येक
  सूत्रासाठी $\pValue v[\Log L](\varphi) = \pValue
  v[\LogCL](\varphi)$.
\end{prop}

\begin{proof}
  $\varphi$ वरील विगमनाने.
  \begin{enumerate}
  \item $\varphi \ident p$ हे आण्विक असेल, तर
  $\pValue v[\Log L](\varphi) = \pAssign v(p) = \pValue
  v[\LogCL](\varphi)$.
  \item $\varphi \ident \lnot B$ असेल, तर
  \begin{align*}
    \pValue v[\Log L](\varphi) & = \tf{\lnot}[\Log L](\pValue v[\Log L](\psi)) & &\text{\ref{mvl:syn:val:defn:pValue} नुसार}\\
    & = \tf{\lnot}[\Log L](\pValue v[\LogCL](\psi)) && \text{विगमन-गृहीतकानुसार}\\
    & = \tf{\lnot}[\LogCL](\pValue v[\LogCL](\psi))
&&\text{\ref{mvl:syn:sub:prop:not} या गृहीतकानुसार,}\\
&&&\text{कारण $\pValue v[\LogCL](\psi) \in \{\True, \False\}$,}\\
& = \pValue v[\LogCL](\varphi) &&\text{\ref{mvl:syn:val:defn:pValue} नुसार}.
\end{align*}
\item $\varphi \ident (\psi \land \chi)$ असेल, तर
\begin{align*}
  \pValue v[\Log L](\varphi) & = \tf{\land}[\Log L](\pValue v[\Log L](\psi), \pValue v[\Log L](\chi)) & &\text{\ref{mvl:syn:val:defn:pValue} नुसार}\\
  & = \tf{\land}[\Log L](\pValue v[\LogCL](\psi),\pValue v[\LogCL](\chi)) && \text{विगमन-गृहीतकानुसार}\\
  & = \tf{\land}[\LogCL](\pValue v[\LogCL](\psi),\pValue v[\LogCL](\chi))
&&\text{\ref{mvl:syn:sub:prop:land} या गृहीतकानुसार,}\\
&&&\text{कारण $\pValue v[\LogCL](\psi),\pValue v[\LogCL](\chi) \in \{\True, \False\}$,}\\
& = \pValue v[\LogCL](\varphi) &&\text{\ref{mvl:syn:val:defn:pValue} नुसार}.
\end{align*}
\end{enumerate}
$\varphi \ident (\psi \lor \chi)$ आणि $\varphi \ident (\psi \lif \chi)$
ही प्रकरणेही यांसारखीच आहेत.
\end{proof}

\begin{cor}
  या समान चार-संयोजक खंडातील सूत्रेच विचारात घेतली असता,
  बहुमूल्य तर्कशास्त्राने \ref{mvl:syn:sub:prop:mvl-cl} च्या अटी पूर्ण केल्या,
  $\True \in V^+$ आणि $\False \notin V^+$ असेल, तर
  ${\Entails[\Log L]} \subseteq {\Entails[\LogCL]}$;
  म्हणजे $\Gamma \Entails[\Log L] \psi$ असल्यास
  $\Gamma \Entails[\LogCL] \psi$. विशेषतः $\Log L$ मधील
  या खंडातील प्रत्येक सर्वतः सत्य सूत्र अभिजात तर्कशास्त्रातही
  सर्वतः सत्य असते.
\end{cor}

\begin{proof}
  व्यत्यय-विधान सिद्ध करू. $\Gamma \Entails/[\LogCL] \psi$
  असे समजा. मग असे सत्य-मूल्यांकन~$\pAssign v\colon \PVar \to
  \{\True, \False\}$ अस्तित्वात असते की सर्व
  $\varphi \in \Gamma$ साठी $\pValue v[\LogCL](\varphi) = \True$
  आणि $\pValue v[\LogCL](\psi) = \False$.
  $\True, \False \in V$ असल्याने सत्य-मूल्यांकन~$\pAssign v$
  हे~$\Log L$ साठीही सत्य-मूल्यांकन आहे.
  \ref{mvl:syn:sub:prop:mvl-cl} नुसार सर्व $\varphi \in \Gamma$ साठी
  $\pValue v[\Log L](\varphi) = \True$ आणि
  $\pValue v[\Log L](\psi) = \False$.
  $\True \in V^+$ आणि $\False \notin V^+$ असल्याने
  $\pSat{v}{\Gamma}[\Log L]$ आणि
  $\pSat/{v}{\psi}[\Log L]$; म्हणजे
  $\Gamma \Entails/[\Log L] \psi$.
\end{proof}


\chapter{त्रिमूल्य तर्कशास्त्रे}\label{mvl:thr:chap}









\section{परिचय}\label{mvl:thr:int:sec}

$\True$ आणि $\False$ या मूल्यांच्या जोडीला आणखी एक मूल्य~$\Undef$
जोडले, तर त्रिमूल्य तर्कशास्त्र मिळते. हे एकच अतिरिक्त
सत्यतामूल्य असले, तरी $\lnot$, $\land$, $\lor$ आणि $\lif$
यांसाठी सत्यता-फलने परिभाषित करण्याचे अनेक पर्याय असतात.
मग एखादे तर्कशास्त्र या सत्यता-फलनांचा कोणताही संयोग वापरू
शकते; तसेच फक्त $\True$ ला निर्दिष्ट मानायचे की $\True$
आणि~$\Undef$ या दोन्हींना, हाही पर्याय असतो.

येथे आपण सर्वाधिक परिचित त्रिमूल्य तर्कशास्त्रांपैकी काही,
त्यांमागील प्रेरणा आणि त्यांचे काही गुणधर्म मांडतो.













\section{\L ukasiewicz यांचे तर्कशास्त्र}\label{mvl:thr:luk:sec}

बहुमूल्य तर्कशास्त्राचे प्रकाशित झालेले आणि सविस्तर मांडलेले
सर्वांत आरंभीच्या प्रस्तावांपैकी एक पोलिश तत्त्वज्ञ Jan \L ukasiewicz
यांनी 1920 मध्ये मांडला.\footnote{Jan \L ukasiewicz,
``O logice trójwartościowej, \emph{Ruch Filozoficzny} 5
(1920), पृ. 170--171; 19 जून 1920 च्या व्याख्यानाचा वृत्तांत.
मूळ अंकाचा स्कॅन: \url{https://kpbc.umk.pl/Content/247342/PDF/EOD_OPEN_006_12_nr_09.pdf}.} अरिस्टॉटलच्या सागरी युद्धाच्या समस्येतून
\L ukasiewicz यांना प्रेरणा मिळाली: \emph{आज} ``उद्या सागरी युद्ध
होईल'' हे विधान सत्यही नाही आणि असत्यही नाही असे दिसते; त्याचे
सत्यतामूल्य अजून निश्चित झालेले नाही. अशा ``भविष्यातील
परिस्थितीवश'' विधानांसाठी \L ukasiewicz यांनी तिसरे सत्यतामूल्य
आणण्याचा प्रस्ताव मांडला.
\begin{quote}
पुढील वर्षी एखाद्या ठराविक क्षणी, उदाहरणार्थ 21 डिसेंबरला दुपारी बारा वाजता,
मी वॉर्सामध्ये असेन, हे सध्या होकारार्थी किंवा नकारार्थी अशा
कोणत्याही प्रकारे निश्चित झालेले नाही, असे मी विरोधाभासाशिवाय
गृहीत धरू शकतो. म्हणून त्या वेळी मी वॉर्सामध्ये असेन हे संभाव्य
आहे, पण अनिवार्य नाही. या गृहितकावर ``पुढील वर्षी 21 डिसेंबरला
दुपारी बारा वाजता मी वॉर्सामध्ये असेन'' हे विधान सध्या सत्यही नाही आणि
असत्यही नाही. कारण ते आताच सत्य असेल, तर भविष्यात मी
वॉर्सामध्ये असणे अनिवार्य ठरेल; हे गृहितकाशी विसंगत आहे.
दुसरीकडे ते आताच असत्य असेल, तर भविष्यात मी वॉर्सामध्ये असणे
अशक्य ठरेल; हेही गृहितकाशी विसंगत आहे. म्हणून विचाराधीन
विधान या क्षणी सत्यही नाही आणि असत्यही नाही; त्याला ``0''
म्हणजे असत्यता आणि ``1'' म्हणजे सत्यता यांपेक्षा वेगळे तिसरे
मूल्य असले पाहिजे. या मूल्याला आपण ``$\frac{1}{2}$.'' असे
दर्शवू शकतो. ते ``संभाव्य'' याचे प्रतिनिधित्व करते आणि
``सत्य'' व ``असत्य'' यांच्याबरोबर तिसरे मूल्य म्हणून सामील होते.
\end{quote}
\L ukasiewicz यांच्या तिसऱ्या सत्यतामूल्यासाठी आपण $\Undef$
वापरू.\footnote{\L ukasiewicz येथे ``संभाव्य'' हा शब्द आज रूढ
नसलेल्या अर्थाने, म्हणजे संभाव्य पण अनिवार्य नाही, या अर्थाने
वापरतात.}

या अर्थानुसार $\lnot$, $\land$ आणि $\lor$ या संयोजकांची
सत्यता-फलने ठरवणे सोपे आहे: भविष्यातील परिस्थितीवश विधानाचे
निषेधनही भविष्यातील परिस्थितीवश विधान असते; म्हणून
$\tf{\lnot}(\Undef) = \Undef$. संधीमधील एक घटक अनिर्धारित
असेल आणि दुसरा सत्य असेल, तर ती संधीही अनिर्धारित असते.
कारण परिस्थितीवश घटक पुढे कसा ठरतो त्यानुसार संधी सत्यही
ठरू शकते आणि असत्यही ठरू शकते. म्हणून \[    \tf{\land}(\True, \Undef) =
\tf{\land}(\Undef, \True) =
\Undef.
\]
पण दुसरा घटक असत्य असेल, तर संधी सत्य ठरू शकत नाही;
म्हणून \[\tf{\land}(\False, \Undef) =
\tf{\land}(\Undef, \False) = \False.\]
दोन्ही निविष्टे $\True$ किंवा $\False$ ही निश्चित सत्यतामूल्ये
असतील, तर उरलेली फलनमूल्ये अभिजात तर्कशास्त्राप्रमाणेच असतात.

सशर्त विधानाच्या बाबतीत परिस्थिती थोडी अधिक गुंतागुंतीची आहे.
$q$ हे भविष्यातील परिस्थितीवश विधान आहे असे समजा. $p$ असत्य
असेल, तर $q$ पुढे कसेही ठरले तरी $p \lif q$ सत्य असेल; म्हणून
$\tf{\lif}(\False, \Undef) = \True$ असे ठेवावे. तसेच $p$
सत्य असेल, तर $q$ पुढे कसेही ठरले तरी $q \lif p$ सत्य असेल;
म्हणून $\tf{\lif}(\Undef, \True) = \True$. $p$ सत्य असेल,
तर $p \lif q$ सत्य किंवा असत्य ठरू शकते; म्हणून
$\tf{\lif}(\True, \Undef) = \Undef$. त्याचप्रमाणे $p$
असत्य असेल, तर $q \lif p$ सत्य किंवा असत्य ठरू शकते;
म्हणून $\tf{\lif}(\Undef, \False) = \Undef$. आता $p$
आणि $q$ दोन्ही भविष्यातील परिस्थितीवश विधाने असण्याचे
प्रकरण उरते. मूळ प्रेरणेनुसार या प्रकरणात खरे तर $\Undef$
हेच मूल्य द्यायला हवे. मात्र तसे केल्यास $\varphi \lif \varphi$
हे \emph{सर्वतः सत्य सूत्र राहणार नाही}. $\varphi \lor \lnot \varphi$
आणि $\lnot(\varphi \land \lnot \varphi)$ यांचा त्याग करायला
\L ukasiewicz यांना अडचण वाटली नाही; पण $\varphi \lif \varphi$
चा त्याग करणे त्यांना मान्य नव्हते. म्हणून त्यांनी
$\tf{\lif}(\Undef, \Undef) = \True$ असे ठरवले.

\begin{defn}\label{mvl:thr:luk:def:lukasiewicz}
त्रिमूल्य \L ukasiewicz तर्कशास्त्र पुढील मॅट्रिक्सने परिभाषित होते:
\begin{enumerate}
  \item मानक विधानीय भाषा $\Lang L_0$ चा केवळ
  $\lnot$, $\land$, $\lor$, $\lif$ हे संयोजक असलेला खंड.
  \item सत्यतामूल्यांचा संच $V = \{\True, \Undef, \False\}$.
  \item $\True$ हेच एकमेव निर्दिष्ट मूल्य आहे, म्हणजे $V^+ = \{\True\}$.
  \item सत्यता-फलने पुढील कोष्टकांनी दिली आहेत:
  \begin{center}
    \begin{tabular}{c|c}
      $\tf{\lnot}$ & \\
      \hline
      $\True$ & $\False$ \\
      $\Undef$ & $\Undef$ \\
      $\False$ & $\True$
    \end{tabular}
    \quad
    \begin{tabular}{c|ccc}
      $\tf{\land}[\LogLuk[3]]$ & $\True$ & $\Undef$ & $\False$ \\
      \hline
      $\True$ & $\True$ & $\Undef$ & $\False$ \\
      $\Undef$ & $\Undef$ & $\Undef$ & $\False$\\
      $\False$ & $\False$ & $\False$ & $\False$
    \end{tabular}
    \\[2ex]
    \begin{tabular}{c|ccc}
      $\tf{\lor}[\LogLuk[3]]$ & $\True$ & $\Undef$ & $\False$ \\
      \hline
      $\True$ & $\True$ & $\True$ & $\True$ \\
      $\Undef$ & $\True$ & $\Undef$ & $\Undef$ \\
      $\False$ & $\True$ & $\Undef$ & $\False$
    \end{tabular}
    \quad
    \begin{tabular}{c|ccc}
      $\tf{\lif}[\LogLuk[3]]$ & $\True$ & $\Undef$ & $\False$ \\
      \hline
      $\True$ & $\True$ & $\Undef$ & $\False$ \\
      $\Undef$ & $\True$ & $\True$ & $\Undef$  \\
      $\False$ & $\True$ & $\True$ & $\True$
    \end{tabular}
  \end{center}
\end{enumerate}
\end{defn}
सहज लक्षात येते की केवळ $\lnot$, $\land$ आणि $\lor$
असलेल्या कोणत्याही सूत्र~$\varphi$ मधील सर्व
विधानीय चले ना $\Undef$ नेमले, तर त्या
सूत्राचे सत्यतामूल्य~$\Undef$ असेल. म्हणून, उदाहरणार्थ,
अभिजात तर्कशास्त्रातील $p \lor \lnot p$ आणि $\lnot(p
\land \lnot p)$ ही सर्वतः सत्य सूत्रे $\LogLuk[3]$
मध्ये सर्वतः सत्य नाहीत; कारण $\pAssign v(p) = \Undef$
असताना $\pValue{v}(\varphi) = \Undef$.
संबंधित प्रत्येक विधानीय चलासाठी $\pAssign v(p) = \True$ किंवा
$\False$ असलेल्या सत्य-मूल्यांकने मध्ये $\pValue v(\varphi)$
हे त्या सूत्राच्या अभिजात सत्यतामूल्याशी जुळते.
\begin{prop}
  $\varphi$ मधील प्रत्येक $p$ साठी $\pAssign v(p) \in \{\True, \False\}$ असेल, तर
  $\pValue v(\varphi)[\LogLuk[3]] = \pValue v(\varphi)[\LogCL]$.
\end{prop}
\begin{prob}\label{mvl:thr:luk:prob:luk-iff} असे समजा की आपण
  $\pValue v(\varphi \liff \psi) = \pValue v((\varphi \lif \psi) \land (\psi \lif
  \varphi))$ अशी $\LogLuk[3]$ मध्ये व्याख्या करतो. मग $\liff$
  चे सत्यताकोष्टक कसे असेल?
\end{prob}
अभिजात तर्कशास्त्रातील अनेक सर्वतः सत्य सूत्रे \LogLuk[3]
मध्येही \emph{सर्वतः सत्य आहेत}; उदाहरणार्थ,
$\lnot p \lif (p \lif q)$. अभिजात तर्कशास्त्राप्रमाणेच
हे सत्यताकोष्टकाने पडताळता येते:
\begin{center}
\begin{tabular}{cc|cccccc}
  $p$ & $q$ & $\lnot$ & $p$ & $\lif$ & $\smash{(}p$ & $\lif$ & $q\smash{)}$ \\ \hline
  \True & \True & \False & \True & \True & \True & \True & \True \\
  \True & \Undef & \False & \True & \True & \True & \Undef & \Undef \\
  \True & \False & \False & \True & \True & \True & \False & \False \\
  \Undef & \True & \Undef & \Undef & \True & \Undef & \True & \True \\
  \Undef & \Undef & \Undef & \Undef & \True & \Undef & \True & \Undef \\
  \Undef & \False & \Undef & \Undef & \True & \Undef & \Undef & \False \\
  \False & \True & \True & \False & \True & \False & \True & \True \\
  \False & \Undef & \True & \False & \True & \False & \True & \Undef \\
  \False & \False & \True & \False & \True & \False & \True & \False \\
\end{tabular}
\end{center}
\begin{prob}
  पुढील सूत्रे \LogLuk[3] मध्ये सर्वतः सत्य आहेत हे दाखवा:
  \begin{enumerate}
    \item $p \lif (q \lif p)$
    \item\label{mvl:thr:luk:prob:luk-taut-2} $\lnot(p \land q) \liff (\lnot p \lor \lnot q)$
    \item\label{mvl:thr:luk:prob:luk-taut-3} $\lnot(p \lor q) \liff (\lnot p \land \lnot q)$
  \end{enumerate}
  (\ref{mvl:thr:luk:prob:luk-taut-2} आणि
  \ref{mvl:thr:luk:prob:luk-taut-3} मध्ये $\varphi \liff \psi$
  हा $(\varphi \lif \psi) \land (\psi \lif \varphi)$ चा संक्षेप माना किंवा
  \ref{mvl:thr:luk:prob:luk-iff} साठीचे तुमचे उत्तर वापरा.)
\end{prob}
\begin{prob}
  पुढील अभिजात सर्वतः सत्य सूत्रे \LogLuk[3] मध्ये सर्वतः सत्य
  नाहीत हे दाखवा:
  \begin{enumerate}
    \item $(\lnot p \land p) \lif q$
    \item $((p \lif q) \lif p) \lif p$
    \item $(p \lif (p \lif q)) \lif (p \lif q)$
  \end{enumerate}
\end{prob}
त्यामुळे अभिजात तर्कशास्त्रातील सर्व सर्वतः सत्य सूत्रे
$\LogLuk[3]$ मध्ये सर्वतः सत्य नसली, तरी प्रत्येक
सत्य-मूल्यांकन मध्ये त्यांना निदान $\True$ किंवा $\Undef$
यांपैकी एक मूल्य मिळेल, असे वाटू शकते. पण तसे नाही.
पुढील सूत्र हे प्रतिउदाहरण आहे:
\[
  \lnot(p \lif \lnot p) \lor \lnot(\lnot p \lif p)
\]
$p$ चे मूल्य~$\Undef$ असेल, तर या सूत्राचे मूल्य $\False$ येते.

\begin{prob}
  पुढीलपैकी कोणते संबंध \L ukasiewicz तर्कशास्त्रात खरे ठरतात?
  प्रत्येकासाठी सत्यताकोष्टक द्या.
  \begin{enumerate}
    \item $p, p \lif q \Entails q$
    \item $\lnot\lnot p \Entails p$
    \item $p \land q \Entails p$
    \item $p \Entails p \land p$
    \item $p \Entails p \lor q$
  \end{enumerate}
\end{prob}

\L ukasiewicz यांना आपल्या त्रिमूल्य पद्धतीच्या आधारे
संभाव्यतेचे तर्कशास्त्र उभारायचे होते. त्यासाठी त्यांनी ``$\varphi$
संभाव्य आहे'' याकरिता एकस्थानी संयोजक $\Diamond \varphi$ आणि
``$\varphi$ अनिवार्य आहे'' याकरिता त्याला अनुरूप $\Box \varphi$ घेतले:
\begin{center}
  \begin{tabular}{c|c}
    $\tf{\Diamond}$ & \\
    \hline
    $\True$ & $\True$ \\
    $\Undef$ & $\True$ \\
    $\False$ & $\False$
  \end{tabular}
  \quad
  \begin{tabular}{c|c}
    $\tf{\Box}$ & \\
    \hline
    $\True$ & $\True$ \\
    $\Undef$ & $\False$ \\
    $\False$ & $\False$
  \end{tabular}
\end{center}
म्हणजे $p$ संभाव्य आहे, जर आणि फक्त जर ते असत्य असल्याचे
आधीच निश्चित झालेले नसेल; आणि $p$ अनिवार्य आहे, जर आणि
फक्त जर ते सत्य असल्याचे आधीच निश्चित झालेले असेल.

\begin{prob}
  $\Box$ आणि~$\Diamond$ यांची सत्यताकोष्टके जोडून विस्तारलेल्या
  \LogLuk[3] मध्ये $\Box p \liff \lnot\Diamond \lnot p$ आणि
  $\Diamond p \liff \lnot \Box \lnot p$ ही सर्वतः सत्य सूत्रे
  आहेत हे दाखवा.
\end{prob}

पण या प्रस्तावित मोडल तर्कशास्त्राच्या मर्यादा लवकरच स्पष्ट
झाल्या: पुढे काहीही घडले तरी $p \land \lnot p$ कधीच सत्य
ठरू शकत नाही. म्हणून सध्या त्याचे मूल्य निश्चित नसले
(आणि त्यामुळे ते अनिर्धारित असले), तरी ते अशक्य मानायला
हवे; म्हणजे $\lnot \Diamond(p \land \lnot p)$ हे सर्वतः सत्य
सूत्र असायला हवे. तथापि, $\pAssign v(p) = \Undef$ असेल,
तर $\pValue v(\lnot \Diamond(p \land \lnot p)) = \False$.
कारण या वेळी $p \land \lnot p$ चे मूल्य $\Undef$,
त्याच्या $\Diamond$ चे मूल्य $\True$, आणि शेवटच्या निषेधनाचे
मूल्य $\False$ असते.
संभाव्यता आणि अनिवार्यता यांसारखे मोडल भेद व्यक्त करण्यासाठी
दोन सत्यतामूल्ये पुरेशी नाहीत, हा \L ukasiewicz यांचा निष्कर्ष
बरोबर होता; पण तिसरे सत्यतामूल्य जोडणेही पुरेसे नाही.













\section{क्लिनी तर्कशास्त्रे}\label{mvl:thr:skl:sec}

Stephen Kleene यांनी संगणन प्रक्रियांच्या परिणामांच्या दृष्टीने
सत्यतामूल्यांचा विचार करणाऱ्या तर्कशास्त्रातून प्रेरित होऊन दोन
त्रिमूल्य तर्कशास्त्रे मांडली: एखादी प्रक्रिया $\True$ किंवा
$\False$ हा परिणाम देऊ शकते; पण ती समाप्तही होणार नाही असे
घडू शकते. त्या वेळी संबंधित सत्यतामूल्य अपरिभाषित असते;
ते सत्यतामूल्य~$\Undef$ ने दर्शवले जाते.

विधान~$\varphi$ चे निषेधन संगणित करण्यासाठी आधी~$\varphi$ चे मूल्य
काढायचे आणि मग त्याच्या विरुद्ध मूल्याचा परिणाम द्यायचा.
$\varphi$ चे संगणन समाप्त होत नसेल, तर ही संपूर्ण प्रक्रियाही
समाप्त होत नाही; म्हणून $\Undef$ चे निषेधन $\Undef$ आहे.

संधी $\varphi \land \psi$ संगणित करण्याचे दोन मार्ग आहेत: आधी~$\varphi$
आणि नंतर $\psi$ संगणित करता येईल. दोन्हींचा परिणाम~$\True$
असेल तर संधीचा परिणाम $\True$, अन्यथा $\False$ असेल.
दोन्हींपैकी कोणतेही संगणन थांबले नाही, तर संपूर्ण प्रक्रियाही
थांबत नाही. त्यामुळे या पद्धतीत संधीचा एक घटक अपरिभाषित
असेल, तर संधीही अपरिभाषित असते. विसंधीबाबतही हेच लागू होते.

पण $\varphi$ आणि $\psi$ यांचे मूल्यांकन समांतरपणे करता आले, तर
अधिक चांगले करता येते. दोन प्रक्रियांपैकी एक थांबून $\False$
हा परिणाम देताच आपण थांबू शकतो, कारण संधी असत्य असलीच
पाहिजे. म्हणून या पद्धतीत एक घटक असत्य असेल, तर दुसरा
अपरिभाषित असला तरी संधी असत्य असते. त्याचप्रमाणे विसंधीचे
समांतर संगणन करताना दोन घटकांपैकी एक सत्य परिणाम देताच
थांबता येते; विसंधी सत्य असलीच पाहिजे. त्यामुळे संयुक्त
विधानाचा एक घटक अपरिभाषित असला तरी त्याचा परिणाम कळू
शकतो. या अर्थानुसार $\Undef$ ला ``अपरिभाषित'' याऐवजी
``अज्ञात'' असे वाचता येईल.

या दोन अर्थांमधून क्लिनी यांची प्रबळ आणि दुर्बल तर्कशास्त्रे
मिळतात. सशर्त विधानाची व्याख्या $\lnot \varphi \lor \psi$ शी
समतुल्य अशी केली जाते.

\begin{defn}
\emph{प्रबळ क्लिनी तर्कशास्त्र}~$\LogKs$ पुढील मॅट्रिक्सने
परिभाषित होते:
\begin{enumerate}
  \item मानक विधानीय भाषा $\Lang L_0$ चा केवळ
  $\lnot$, $\land$, $\lor$, $\lif$ हे संयोजक असलेला खंड.
  \item सत्यतामूल्यांचा संच $V = \{\True, \Undef, \False\}$.
  \item $\True$ हेच एकमेव निर्दिष्ट मूल्य आहे, म्हणजे $V^+ = \{\True\}$.
  \item सत्यता-फलने पुढील कोष्टकांनी दिली आहेत:
  \begin{center}
    \begin{tabular}{c|c}
      $\tf{\lnot}$ & \\
      \hline
      $\True$ & $\False$ \\
      $\Undef$ & $\Undef$ \\
      $\False$ & $\True$
    \end{tabular}
    \quad
    \begin{tabular}{c|ccc}
      $\tf{\land}[\LogKs]$ & $\True$ & $\Undef$ & $\False$ \\
      \hline
      $\True$ & $\True$ & $\Undef$ & $\False$ \\
      $\Undef$ & $\Undef$ & $\Undef$ & $\False$\\
      $\False$ & $\False$ & $\False$ & $\False$
    \end{tabular}
    \\[2ex]
    \begin{tabular}{c|ccc}
      $\tf{\lor}[\LogKs]$ & $\True$ & $\Undef$ & $\False$ \\
      \hline
      $\True$ & $\True$ & $\True$ & $\True$ \\
      $\Undef$ & $\True$ & $\Undef$ & $\Undef$ \\
      $\False$ & $\True$ & $\Undef$ & $\False$
    \end{tabular}
    \quad
    \begin{tabular}{c|ccc}
      $\tf{\lif}[\LogKs]$ & $\True$ & $\Undef$ & $\False$ \\
      \hline
      $\True$ & $\True$ & $\Undef$ & $\False$ \\
      $\Undef$ & $\True$ & $\Undef$ & $\Undef$  \\
      $\False$ & $\True$ & $\True$ & $\True$
    \end{tabular}
  \end{center}
\end{enumerate}
\end{defn}
\begin{defn}
\emph{दुर्बल क्लिनी तर्कशास्त्र}~$\LogKw$ पुढील मॅट्रिक्सने
परिभाषित होते:
\begin{enumerate}
  \item मानक विधानीय भाषा $\Lang L_0$ चा केवळ
  $\lnot$, $\land$, $\lor$, $\lif$ हे संयोजक असलेला खंड.
  \item सत्यतामूल्यांचा संच $V = \{\True, \Undef, \False\}$.
  \item $\True$ हेच एकमेव निर्दिष्ट मूल्य आहे, म्हणजे $V^+ = \{\True\}$.
  \item सत्यता-फलने पुढील कोष्टकांनी दिली आहेत:
  \begin{center}
    \begin{tabular}{c|c}
      $\tf{\lnot}$ & \\
      \hline
      $\True$ & $\False$ \\
      $\Undef$ & $\Undef$ \\
      $\False$ & $\True$
    \end{tabular}
    \quad
    \begin{tabular}{c|ccc}
      $\tf{\land}[\LogKw]$ & $\True$ & $\Undef$ & $\False$ \\
      \hline
      $\True$ & $\True$ & $\Undef$ & $\False$ \\
      $\Undef$ & $\Undef$ & $\Undef$ & $\Undef$\\
      $\False$ & $\False$ & $\Undef$ & $\False$
    \end{tabular}
    \\[2ex]
    \begin{tabular}{c|ccc}
      $\tf{\lor}[\LogKw]$ & $\True$ & $\Undef$ & $\False$ \\
      \hline
      $\True$ & $\True$ & $\Undef$ & $\True$ \\
      $\Undef$ & $\Undef$ & $\Undef$ & $\Undef$ \\
      $\False$ & $\True$ & $\Undef$ & $\False$
    \end{tabular}
    \quad
    \begin{tabular}{c|ccc}
      $\tf{\lif}[\LogKw]$ & $\True$ & $\Undef$ & $\False$ \\
      \hline
      $\True$ & $\True$ & $\Undef$ & $\False$ \\
      $\Undef$ & $\Undef$ & $\Undef$ & $\Undef$  \\
      $\False$ & $\True$ & $\Undef$ & $\True$
    \end{tabular}
  \end{center}
\end{enumerate}
\end{defn}
\begin{prop}
  $\LogKs$ आणि $\LogKw$ यांत कोणतेही सर्वतः सत्य सूत्र नाही.
\end{prop}
\begin{proof}
  सर्व विधानीय चले~$p$ साठी
  $\pAssign v(p) = \Undef$ असेल, तर कोणत्याही
  सूत्र~$\varphi$ चे सत्यतामूल्य~$\pValue v(\varphi) =
  \Undef$ असेल; कारण दोन्ही तर्कशास्त्रांत
  \[
    \tf{\lnot}(\Undef) = \tf{\lor}(\Undef, \Undef) = \tf{\land}(\Undef,
  \Undef) = \tf{\lif}(\Undef, \Undef) = \Undef
  \]
  असे आहे. $\LogKs$ आणि $\LogKw$ या दोन्हींमध्ये
  $\Undef \notin V^+$ असल्यामुळे या सत्य-मूल्यांकन मध्ये
  $\varphi$ चे सत्यतामूल्य निर्दिष्ट नसेल.
\end{proof}

दुर्बल आणि प्रबळ क्लिनी तर्कशास्त्रांत सर्वतः सत्य सूत्रे
नसली, तरी त्यांचे तार्किक निष्पन्नता संबंध अतुच्छ आहेत.

\begin{prob}
  पुढीलपैकी कोणते संबंध (a) प्रबळ आणि (b) दुर्बल क्लिनी
  तर्कशास्त्रात खरे ठरतात? प्रत्येकासाठी सत्यताकोष्टक द्या.
  \begin{enumerate}
    \item $p, p \lif q \Entails q$
    \item $p \lor q, \lnot p \Entails q$
    \item $p \land q \Entails p$
    \item $p \Entails p \land p$
    \item $p \Entails p \lor q$
  \end{enumerate}
\end{prob}

Dmitry Bochvar यांनी $\Undef$ चा अर्थ ``अर्थहीन'' असा घेतला
आणि विरोधाभास निर्माण करणाऱ्या विधानांना~$\Undef$ हे मूल्य
देऊन खोटारड्याच्या विरोधाभासासारखे विरोधाभास सोडवण्याचा
प्रयत्न केला. त्यांनी मूलतः दुर्बल क्लिनी तर्कशास्त्राला
अतिरिक्त संयोजके जोडून तयार होणारे तर्कशास्त्र मांडले.
त्यांतील दोन संयोजके ``बाह्य निषेधन'' आणि ``अपरिभाषित
आहे'' हा परिकर्मी अशी आहेत:
\begin{center}
  \begin{tabular}{c|c}
    $\tf{\sim}$ & \\
    \hline
    $\True$ & $\False$ \\
    $\Undef$ & $\True$ \\
    $\False$ & $\True$
  \end{tabular}
\quad
  \begin{tabular}{c|c}
    $\tf{+}$ & \\
    \hline
    $\True$ & $\False$ \\
    $\Undef$ & $\True$ \\
    $\False$ & $\False$
  \end{tabular}
\end{center}

\begin{prob}
  Bochvar यांच्या तर्कशास्त्रात $\lnot$ आणि~$+$ यांच्या
  साहाय्याने $\sim$ ची व्याख्या करता येते का? म्हणजे, फक्त
  विधानीय चल~$p$ असलेले, $\sim$ नसलेले
  आणि नेहमी~$\mathord{\sim}p$ इतकेच सत्यतामूल्य घेणारे
  सूत्र शोधा. तुमचे उत्तर योग्य आहे हे
  सत्यताकोष्टकाने दाखवा.
\end{prob}













\section{G\"odel तर्कशास्त्रे}\label{mvl:thr:god:sec}

Kurt G\"odel यांनी $n$-मूल्य तर्कशास्त्रांची एक क्रमिका
मांडली. त्यांतील प्रत्येकात अंतःप्रज्ञावादी तर्कशास्त्रात
वैध असलेली सर्व सूत्रे येतात आणि प्रत्येक अभिजात
तर्कशास्त्राचे उपतर्कशास्त्र आहे. त्यांपैकी पहिले
लक्षवेधी तर्कशास्त्र पुढीलप्रमाणे आहे:

\begin{defn}\label{mvl:thr:god:defn:goedel}
\emph{$3$-मूल्य G\"odel तर्कशास्त्र}~$\LogGod$ पुढील
मॅट्रिक्सने परिभाषित होते:
\begin{enumerate}
  \item $\lfalse$, $\lnot$, $\land$, $\lor$, $\lif$ असलेली मानक विधानीय भाषा $\Lang L_0$.
  \item सत्यतामूल्यांचा संच $V = \{\True, \Undef, \False\}$.
  \item $\True$ हेच एकमेव निर्दिष्ट मूल्य आहे, म्हणजे $V^+ = \{\True\}$.
  \item $\lfalse$ साठी $\tf{\lfalse} = \False$. उरलेल्या
  संयोजकांची सत्यता-फलने पुढील कोष्टकांनी दिली आहेत:
  \begin{center}
    \begin{tabular}{c|c}
      $\tf{\lnot}[\LogGod]$ & \\
      \hline
      $\True$ & $\False$ \\
      $\Undef$ & $\False$ \\
      $\False$ & $\True$
    \end{tabular}
    \quad
    \begin{tabular}{c|ccc}
      $\tf{\land}[\LogGod]$ & $\True$ & $\Undef$ & $\False$ \\
      \hline
      $\True$ & $\True$ & $\Undef$ & $\False$ \\
      $\Undef$ & $\Undef$ & $\Undef$ & $\False$\\
      $\False$ & $\False$ & $\False$ & $\False$
    \end{tabular}
    \\[2ex]
    \begin{tabular}{c|ccc}
      $\tf{\lor}[\LogGod]$ & $\True$ & $\Undef$ & $\False$ \\
      \hline
      $\True$ & $\True$ & $\True$ & $\True$ \\
      $\Undef$ & $\True$ & $\Undef$ & $\Undef$ \\
      $\False$ & $\True$ & $\Undef$ & $\False$
    \end{tabular}
    \quad
    \begin{tabular}{c|ccc}
      $\tf{\lif}[\LogGod]$ & $\True$ & $\Undef$ & $\False$ \\
      \hline
      $\True$ & $\True$ & $\Undef$ & $\False$ \\
      $\Undef$ & $\True$ & $\True$ & $\False$  \\
      $\False$ & $\True$ & $\True$ & $\True$
    \end{tabular}
  \end{center}
\end{enumerate}
\end{defn}
$\land$ आणि~$\lor$ यांची सत्यताकोष्टके \L ukasiewicz
आणि प्रबळ क्लिनी तर्कशास्त्रांप्रमाणेच आहेत, हे लक्षात येईल;
पण $\lnot$ आणि~$\lif$ यांच्या कोष्टकांत फरक आहे.
G\"odel तर्कशास्त्रात $\tf{\lnot}(\Undef) = \False$.
\L ukasiewicz आणि क्लिनी तर्कशास्त्रांच्या विपरीत,
$\tf{\lif}(\Undef, \False) = \False$; तर क्लिनी
तर्कशास्त्राच्या विपरीत (पण \L ukasiewicz तर्कशास्त्राप्रमाणे)
$\tf{\lif}(\Undef, \Undef) = \True$.
वर सूचित केलेल्या अंतःप्रज्ञावादी तर्कशास्त्राशी असलेल्या
संबंधावरून अपेक्षित आहे त्याप्रमाणे $\LogGod[3]$ हे
अंतःप्रज्ञावादी तर्कशास्त्राच्या जवळचे आहे. त्या तर्कशास्त्रात
वैध असलेली सर्व सूत्रे $\LogGod[3]$ मध्ये सर्वतः सत्य
असतात; आणि अंतःप्रज्ञावादी तर्कशास्त्रात वैध नसलेली अनेक
अभिजात सर्वतः सत्य सूत्रे $\LogGod[3]$ मध्येही सर्वतः
सत्य नसतात. उदाहरणार्थ, पुढील सूत्रे सर्वतः सत्य नाहीत:
\begin{align*}
  & p \lor \lnot p && (p \lif q) \lif (\lnot p \lor q) \\
  & \lnot\lnot p \lif p && \lnot(\lnot p \land \lnot q) \lif (p \lor q) \\
  & ((p \lif q) \lif p) \lif p && \lnot(p \lif q) \lif (p \land \lnot q)
\end{align*}
मात्र $\LogGod[3]$ मधील प्रत्येक सर्वतः सत्य सूत्र
अंतःप्रज्ञावादी तर्कशास्त्रातही वैध असेलच असे नाही;
उदाहरणार्थ, $\lnot\lnot p \lor \lnot p$ किंवा $(p
\lif q) \lor (q \lif p)$.
\begin{prob}
  पुढील सूत्रे $\LogGod[3]$ मध्ये सर्वतः सत्य आहेत हे
  सत्यताकोष्टकांनी दाखवा:
  \begin{align*}
    & \lnot\lnot p \lor \lnot p\\
    & (p \lif q) \lor (q \lif p) \\
    & \lnot(p \land q) \lif (\lnot p \lor \lnot q) \\
    & (p \lif q) \lor (q \lif r) \lor (r \lif s)
  \end{align*}
\end{prob}
\begin{prob}
  पुढील सूत्रे $\LogGod[3]$ मध्ये सर्वतः सत्य नाहीत हे
  सत्यताकोष्टकांनी दाखवा:
  \begin{align*}
    & (p \lif q) \lif (\lnot p \lor q) \\
    & \lnot(\lnot p \land \lnot q) \lif (p \lor q) \\
    & ((p \lif q) \lif p) \lif p \\
    & \lnot(p \lif q) \lif (p \land \lnot q)
  \end{align*}
\end{prob}
\begin{prob}
  पुढीलपैकी कोणते संबंध G\"odel तर्कशास्त्रात खरे ठरतात?
  प्रत्येकासाठी सत्यताकोष्टक द्या.
  \begin{enumerate}
    \item $p, p \lif q \Entails q$
    \item $p \lor q, \lnot p \Entails q$
    \item $p \land q \Entails p$
    \item $p \Entails p \land p$
    \item $p \Entails p \lor q$
  \end{enumerate}
\end{prob}
\section[इतर सत्यतामूल्येही निर्दिष्ट करणे]{$\True$ व्यतिरिक्तही मूल्य निर्दिष्ट करणे}\label{mvl:thr:mul:sec}
आतापर्यंत पाहिलेल्या सर्व तर्कशास्त्रांत निर्दिष्ट
सत्यतामूल्यांचा संच $V^+ = \{\True\}$ होता; म्हणजे एखाद्या
विधान निर्दिष्ट मानले जाई, जर आणि फक्त जर त्याचे
सत्यतामूल्य~$\True$ असेल.
पण एखादे विधान फक्त~$\False$ नसले, तरी ते सत्य मानता येईल.
अशा वेळी, उदाहरणार्थ, मॅट्रिक्समध्ये
$V^+ = \{\True, \Undef\}$ असे ठरवून तर्कशास्त्र मिळेल.
\begin{defn}
\emph{विरोधापत्तीचे तर्कशास्त्र}~$\LogLP$ पुढील मॅट्रिक्सने
परिभाषित होते:
\begin{enumerate}
  \item मानक विधानीय भाषा $\Lang L_0$ चा केवळ
  $\lnot$, $\land$, $\lor$, $\lif$ हे संयोजक असलेला खंड.
  \item सत्यतामूल्यांचा संच $V = \{\True, \Undef, \False\}$.
  \item $\True$ आणि $\Undef$ निर्दिष्ट आहेत, म्हणजे $V^+ = \{\True, \Undef\}$.
  \item सत्यता-फलने प्रबळ क्लिनी तर्कशास्त्रातील फलनांसारखीच आहेत.
\end{enumerate}
\end{defn}
\begin{defn}
Halld\'en यांचे \emph{निरर्थकतेचे तर्कशास्त्र}~$\LogHal$
पुढील मॅट्रिक्सने परिभाषित होते:
\begin{enumerate}
  \item मानक विधानीय भाषा $\Lang L_0$ च्या
  $\lnot$, $\land$, $\lor$, $\lif$ या संयोजकांच्या खंडाचा
  $1$-स्थानी संयोजक~$+$ जोडून केलेला विस्तार.
  \item सत्यतामूल्यांचा संच $V = \{\True, \Undef, \False\}$.
  \item $\True$ आणि $\Undef$ निर्दिष्ट आहेत, म्हणजे $V^+ = \{\True, \Undef\}$.
  \item सत्यता-फलने दुर्बल क्लिनी तर्कशास्त्रातील फलनांसारखीच
  आहेत; त्याखेरीज ``अर्थहीन आहे'' हा परिकर्मी आहे:
  \begin{center}
    \begin{tabular}{c|c}
    $\tf{+}$ & \\
    \hline
    $\True$ & $\False$ \\
    $\Undef$ & $\True$ \\
    $\False$ & $\False$
    \end{tabular}
  \end{center}
\end{enumerate}
\end{defn}
यांची सत्यताकोष्टके क्लिनी तर्कशास्त्रांशी जुळतात; पण
त्यांच्या विपरीत यांत सर्वतः सत्य सूत्रे \emph{आहेत}.
\begin{prop}\label{mvl:thr:mul:prop:LP-taut-CL} $\LogLP$ मधील सर्वतः सत्य
  सूत्रे आणि अभिजात विधानीय तर्कशास्त्रातील सर्वतः सत्य सूत्रे
  एकच आहेत.
\end{prop}
\begin{proof}
  \ref{mvl:syn:sub:prop:mvl-cl} नुसार $\Entails[\LogLP] \varphi$
  असेल, तर $\Entails[\LogCL] \varphi$. उलट दिशा दाखवण्यासाठी
  आपण सिद्ध करू की $\pValue v(\varphi)[\LogKs] = \False$ असेल असे
  सत्य-मूल्यांकन~$\pAssign v\colon \PVar \to \{\False, \True,
  \Undef\}$ असल्यास $\pValue {v'}(\varphi)[\LogCL] = \False$
  असे सत्य-मूल्यांकन~$\pAssign {v'}\colon \PVar \to \{\False, \True\}$
  असते. यावरून $\LogLP$ साठी अपेक्षित निष्कर्ष मिळतो. कारण
  $\LogKs$ आणि $\LogLP$ यांची वैशिष्ट्यपूर्ण सत्यता-फलने
  सारखीच आहेत आणि $\LogLP$ मध्ये फक्त $\False$ हेच मूल्य
  निर्दिष्ट नाही ($\LogLP$ आणि~$\LogKs$ यांतील तोच एक फरक
  आहे). म्हणून $\Entails/[\LogLP] \varphi$ असेल, तर कोणत्यातरी
  सत्य-मूल्यांकन~$\pAssign v$ साठी $\pValue v(\varphi)[\LogLP] = \pValue
  v[\LogKs](\varphi) = \False$. सिद्ध करावयाच्या विधानानुसार
  $\pValue{v'}[\LogCL](\varphi) = \False$, म्हणजे $\Entails/[\LogCL] \varphi$.
  हे सिद्ध करण्यासाठी आधी $\pAssign {v'}$ ची व्याख्या अशी करू:
  \[
    \pAssign {v'}(p) =
    \begin{cases}
    \True & \text{जर } \pAssign {v}(p) \in \{\True, \Undef\}\\
    \False & \text{इतर वेळी}
  \end{cases}
  \]
  आता $\varphi$ वरील विगमनाने दाखवू की (a)~$\pValue v(\varphi)[\LogKs]
  = \False$ असेल, तर $\pValue {v'}(\varphi)[\LogCL] = \False$;
  आणि (b)~$\pValue v(\varphi)[\LogKs] = \True$ असेल, तर
  $\pValue {v'}(\varphi)[\LogCL] = \True$.
  \begin{enumerate}
    \item विगमनाचा आधार: $\varphi \ident p$. (a) किंवा (b) चा
    पूर्वपक्ष लागू असेल, तेव्हा मूळ मूल्य अनुक्रमे असत्य किंवा
    सत्य असते; त्या प्रकरणांतच
    \ref{mvl:syn:val:defn:pValue} नुसार $\pValue v(\varphi)[\LogKs]
    = \pAssign v(p) = \pValue {v'}(\varphi)[\LogCL]$ ही समानता
    मिळते, आणि त्यावरून (a) व~(b) दोन्ही सिद्ध होतात.

    विगमनाच्या पायरीसाठी पुढील प्रकरणे पाहू:
    \item $\varphi \ident \lnot \psi$.
    \begin{enumerate}
      \item $\pValue v(\lnot\psi)[\LogKs] = \False$ असे समजा.
      $\tf{\lnot}[\LogKs]$ च्या व्याख्येनुसार
      $\pValue v(\psi)[\LogKs]  = \True$. विगमन गृहितकातील
      प्रकरण~(b) नुसार $\pValue {v'}(\psi)[\LogCL]  = \True$;
      म्हणून $\pValue {v'}(\lnot \psi)[\LogCL] = \False$.
      \item $\pValue v(\lnot\psi)[\LogKs] = \True$ असे समजा.
      $\tf{\lnot}[\LogKs]$ च्या व्याख्येनुसार
      $\pValue v(\psi)[\LogKs]  = \False$. विगमन गृहितकातील
      प्रकरण~(a) नुसार $\pValue {v'}(\psi)[\LogCL]  = \False$;
      म्हणून $\pValue {v'}(\lnot \psi)[\LogCL] = \True$.
    \end{enumerate}

    \item $\varphi \ident (\psi \land \chi)$.
    \begin{enumerate}
      \item $\pValue v(\psi \land \chi)[\LogKs] = \False$ असे समजा.
      $\tf{\land}[\LogKs]$ च्या व्याख्येनुसार
      $\pValue v(\psi)[\LogKs]  = \False$ किंवा
      $\pValue v(\chi)[\LogKs]  = \False$. विगमन गृहितकातील
      प्रकरण~(a) नुसार $\pValue {v'}(\psi)[\LogCL]  = \False$
      किंवा $\pValue {v'}(\chi)[\LogCL]  = \False$;
      म्हणून $\pValue {v'}(\psi \land \chi)[\LogCL] = \False$.
      \item $\pValue v(\psi \land \chi)[\LogKs] = \True$ असे समजा.
      $\tf{\land}[\LogKs]$ च्या व्याख्येनुसार
      $\pValue v(\psi)[\LogKs]  = \True$ आणि
      $\pValue v(\chi)[\LogKs]  = \True$. विगमन गृहितकातील
      प्रकरण~(b) नुसार $\pValue {v'}(\psi)[\LogCL]  = \True$
      आणि $\pValue {v'}(\chi)[\LogCL]  = \True$;
      म्हणून $\pValue {v'}(\psi \land \chi)[\LogCL] = \True$.
    \end{enumerate}
  \end{enumerate}

  उरलेली दोन प्रकरणे याच प्रकारची आहेत आणि सरावासाठी
  सोडली आहेत. दुसरा मार्ग असा: वरील पुराव्याने फक्त $\lnot$
  आणि~$\land$ असलेल्या सर्व सूत्रे साठी निष्कर्ष
  सिद्ध केला आहे. आता $\LogKs$ आणि $\LogCL$ या दोन्हींमध्ये
  कोणत्याही $\pAssign v$ साठी $\pValue v(\psi \lor
  \chi) = \pValue v(\lnot(\lnot\psi \land \lnot \chi))$ आणि $\pValue v(\psi \lif
  \chi) = \pValue v(\lnot(\psi \land \lnot \chi))$ या समानता
  वापरता येतात.
\end{proof}

\begin{prob}
  \ref{mvl:thr:mul:prop:LP-taut-CL} चा पुरावा पूर्ण
  करा; म्हणजे $\varphi \ident (\psi \lor \chi)$ आणि
  $\varphi \ident (\psi \lif \chi)$ या प्रकरणांसाठी (a) व~(b)
  सिद्ध करा.
\end{prob}

\begin{prob}
प्रत्येक अभिजात सर्वतः सत्य सूत्र $\LogHal$ मध्येही
सर्वतः सत्य आहे हे सिद्ध करा.
\end{prob}

दोन्ही तर्कशास्त्रांच्या समान चार-संयोजकीय भाषाखंडात त्यांची
सर्वतः सत्य सूत्रे अभिजात तर्कशास्त्रासारखीच असली, तरी त्यांचे
तार्किक निष्पन्नता संबंध वेगळे आहेत. Halld\'en यांच्या अतिरिक्त
संयोजकाचा समावेश असलेल्या सूत्रांसाठी हे समानतेचे विधान नाही. उदाहरणार्थ,
$\LogLP$ हे \emph{व्याघातसहिष्णु} आहे: त्यात
$\lnot p, p \Entails/ q$. म्हणून
$\lnot \varphi, \varphi \Entails \psi$ हे स्फोटाचे तत्त्व सर्वसाधारणपणे
लागू होत नाही. ($\varphi$ आणि $\psi$ यांच्या काही प्रकरणांत
ते लागू होते; उदाहरणार्थ, $\psi$~हे सर्वतः सत्य सूत्र असेल तर.)

\begin{prob}
  पुढीलपैकी कोणते संबंध (a)~$\LogLP$ आणि (b)~$\LogHal$
  मध्ये खरे ठरतात? प्रत्येकासाठी सत्यताकोष्टक द्या.
  \begin{enumerate}
    \item $p, p \lif q \Entails q$
    \item $\lnot q, p \lif q \Entails \lnot p$
    \item $p \lor q, \lnot p \Entails q$
    \item $\lnot p, p \Entails q$
    \item $p \Entails p \lor q$
    \item $p \lif q, q\lif r \Entails p \lif r$
  \end{enumerate}
\end{prob}

$\LogLuk[3]$ मध्ये $\Undef$ हे मूल्यही निर्दिष्ट मानल्यास
काय होईल?
केवळ निर्दिष्ट मूल्यांचा हा बदल केल्याने पुढील R-Mingle
मॅट्रिक्स मिळत नाही; त्यासाठी सशर्त विधानाचे सत्यता-फलनही बदलते.

\begin{defn}
  \emph{त्रिमूल्य R-Mingle}~$\LogRM[3]$ चा असत्यता-स्थिरासह
  विस्तार पुढील मॅट्रिक्सने परिभाषित होतो:
  \begin{enumerate}
    \item $\lfalse$, $\lnot$, $\land$, $\lor$, $\lif$ असलेली मानक विधानीय भाषा $\Lang L_0$.
    \item सत्यतामूल्यांचा संच $V = \{\True, \Undef, \False\}$.
    \item $\True$ आणि $\Undef$ निर्दिष्ट आहेत, म्हणजे $V^+ = \{\True, \Undef\}$.
    \item $\tf{\lfalse}=\False$. $\lnot$, $\land$ आणि $\lor$
    यांची सत्यता-फलने \L ukasiewicz तर्कशास्त्र~$\LogLuk[3]$
    मधील फलनांसारखीच आहेत. सशर्त विधानाचे सत्यताकोष्टक असे आहे:
    \begin{center}
      \begin{tabular}{c|ccc}
        $\tf{\lif}[\LogRM[3]]$ & $\True$ & $\Undef$ & $\False$ \\
        \hline
        $\True$ & $\True$ & $\False$ & $\False$ \\
        $\Undef$ & $\True$ & $\Undef$ & $\False$ \\
        $\False$ & $\True$ & $\True$ & $\True$
      \end{tabular}
    \end{center}
  \end{enumerate}
  येथे $\Undef$ हे मधले मूल्य सत्य आणि असत्य दोन्ही असल्याचे
  दर्शवते. सशर्त विधानाचे कोष्टक \L ukasiewicz यांच्या कोष्टकापेक्षा
  वेगळे आहे.\footnote{Barteld Kooi आणि Allard Tamminga,
  ``Completeness via Correspondence for Extensions of the Logic of
  Paradox,'' (2012), प्रथम ऑनलाइन आवृत्ती पृ. 10, विभाग 4;
  \url{https://doi.org/10.1017/S1755020312000196}.}
\end{defn}

\begin{prob}
  पुढीलपैकी कोणते संबंध $\LogRM[3]$ मध्ये खरे ठरतात?
  \begin{enumerate}
    \item $p, p \lif q \Entails q$
    \item $p \lor q, \lnot p \Entails q$
    \item $\lnot p, p \Entails q$
    \item $p \Entails p \lor q$
  \end{enumerate}
\end{prob}

निर्दिष्ट मूल्ये बदलल्यामुळे वेगवेगळी सत्यताकोष्टके कधी कधी
एकच तर्कशास्त्र, म्हणजे एकच तार्किक निष्पन्नता संबंध, देऊ
शकतात. उदाहरणार्थ, G\"odel तर्कशास्त्रात~$\{\True\}$ ऐवजी
$V^+ = \{\True, \Undef\}$ घेतल्यास असे घडते.

\begin{prop}\label{mvl:thr:mul:prop:gl-udes}
  $V = \{\False, \Undef,\True\}$,
  $V^+=\{\True, \Undef\}$ आणि $3$-मूल्य G\"odel
  तर्कशास्त्राची सत्यता-फलने असलेला मॅट्रिक्स अभिजात
  तर्कशास्त्र परिभाषित करतो.
\end{prop}

\begin{proof}
  सराव.
\end{proof}

\begin{prob}
  G\"odel तर्कशास्त्राप्रमाणेच पण
  $V^+=\{\True,\Undef\}$ घेऊन परिभाषित केलेल्या
  तर्कशास्त्र~$\Log L$ साठी $\Gamma \Entails/[\Log L] \psi$
  असेल, तर $\Gamma \Entails/[\LogCL] \psi$ हे दाखवून
  \ref{mvl:thr:mul:prop:gl-udes} सिद्ध करा.
  \ref{mvl:thr:mul:prop:LP-taut-CL} मधील कल्पना
  वापरा; मात्र गुणधर्म (a) आणि~(b) सिद्ध करण्याऐवजी
  $\pValue v(\varphi)[\LogGod] = \False$ जर आणि फक्त जर
  $\pValue {v'}(\varphi)[\LogCL] = \False$ हे दाखवा (आणि म्हणून
  $\pValue v(\varphi)[\LogGod] \in \{\True,\Undef\}$ जर आणि
  फक्त जर $\pValue {v'}(\varphi)[\LogCL] = \True$).
  यावरून विधान का सिद्ध होते ते स्पष्ट करा.
\end{prob}



\chapter{अनंतमूल्य तर्कशास्त्रे}\label{mvl:inf:chap}









\section{परिचय}\label{mvl:inf:int:sec}

मॅट्रिक्समधील सत्यतामूल्यांची संख्या ससीम असण्याची गरज नाही.
अनंत सत्यतामूल्यांसाठी एक नैसर्गिक पर्याय म्हणजे $0$ आणि~$1$
यांदरम्यानच्या परिमेय संख्यांचा संच,
$V_\infty = [0,1] \cap \Rat$; म्हणजे
\begin{align*}
    V_\infty & = \Setabs{\frac{n}{m}}{n,m \in \Nat \text{ आणि } 0<m \text{ आणि } n\le m}.
\intertext{हा अनंत सत्यतामूल्यांचा संच विचारात घेताना दोन किंवा
अधिक मूल्ये असलेले पुढील उपसंचही पाहणे अनेकदा उपयोगाचे असते:}
V_m & = \Setabs{\frac{n}{m-1}}{n \in \Nat \text{ आणि } n<m}
\intertext{उदाहरणार्थ, $V_5$ मध्ये समान अंतरावर असलेली $5$ सत्यतामूल्ये आहेत:}
V_5 & = \{0, \frac{1}{4}, \frac{1}{2}, \frac{3}{4}, 1\}.
\end{align*}
या सत्यतामूल्यांच्या संचांवर आधारलेल्या तर्कशास्त्रांत
सहसा फक्त $1$ निर्दिष्ट असते, म्हणजे $V^+ = \{1\}$.
दुसऱ्या शब्दांत, $1$ ला निरपेक्ष सत्याची आणि $0$ ला
निरपेक्ष असत्यतेची भूमिका देतो; पण सूत्रे ना~$V$
मधील कोणतेही मधले मूल्य मिळू शकते.

याशिवाय $0$ आणि~$1$ यांदरम्यानच्या सर्व \emph{वास्तव}
संख्यांचा संच $V_{[0,1]} = [0,1]$, किंवा $[0,1]$ चे इतर
अनंत उपसंचही विचारात घेता येतात. हा सत्यतामूल्यांचा संच
वापरणाऱ्या तर्कशास्त्रांना अनेकदा \emph{अस्पष्ट} म्हणतात.














\section{\L ukasiewicz यांचे तर्कशास्त्र}\label{mvl:inf:luk:sec}

\begin{editorial}
  अनंतमूल्य \L ukasiewicz तर्कशास्त्रावरील हा विभाग
  संक्षिप्त ``प्रारूप'' आहे.
\end{editorial}

\begin{defn}\label{mvl:inf:luk:def:lukasiewicz} अनंतमूल्य \L ukasiewicz
तर्कशास्त्र~$\LogLuk[\infty]$ पुढील मॅट्रिक्सने परिभाषित होते:
\begin{enumerate}
  \item मानक विधानीय भाषा $\Lang L_0$ चा केवळ
  $\lnot$, $\land$, $\lor$, $\lif$ हे संयोजक असलेला खंड.
  \item सत्यतामूल्यांचा संच $V_\infty$.
  \item $1$ हेच एकमेव निर्दिष्ट मूल्य आहे, म्हणजे $V^+ = \{1\}$.
  \item सत्यता-फलने पुढील सूत्रांनी दिली आहेत:
  \begin{align*}
    \tf{\lnot}[\LogLuk](x) & = 1 - x\\
    \tf{\land}[\LogLuk](x,y) & = \min(x,y)\\
    \tf{\lor}[\LogLuk](x,y) & = \max(x,y)\\
    \tf{\lif}[\LogLuk](x,y) & = \min(1,1-(x-y)) = \begin{cases}
      1 & \text{जर } x \le y\\
      1-(x-y) & \text{इतर वेळी.}
    \end{cases}
    \end{align*}
\end{enumerate}
$m$-मूल्य \L ukasiewicz तर्कशास्त्राची व्याख्या याचप्रमाणे,
फक्त $V = V_m$ घेऊन केली जाते.
\end{defn}

\begin{prop}
  \ref{mvl:thr:luk:def:lukasiewicz} मध्ये परिभाषित
  $\LogLuk[3]$ आणि \ref{mvl:inf:luk:def:lukasiewicz} मध्ये
  परिभाषित $\LogLuk[3]$ हे एकच तर्कशास्त्र आहे.
\end{prop}

\begin{proof}
  \ref{mvl:thr:luk:def:lukasiewicz} मधील संयोजकांची
  सत्यताकोष्टके आणि \ref{mvl:inf:luk:def:lukasiewicz} मधील
  समीकरणांनी निश्चित होणारी सत्यताकोष्टके यांची तुलना
  केल्यास हे दिसते:
  \begin{center}
    \begin{tabular}{c|c}
      $\tf{\lnot}$ & \\
      \hline
      $1$ & $0$ \\
      $1/2$ & $1/2$ \\
      $0$ & $1$
    \end{tabular}
    \quad
    \begin{tabular}{c|ccc}
      $\tf{\land}[\LogLuk[3]]$ & $1$ & $1/2$ & $0$ \\
      \hline
      $1$ & $1$ & $1/2$ & $0$ \\
      $1/2$ & $1/2$ & $1/2$ & $0$\\
      $0$ & $0$ & $0$ & $0$
    \end{tabular}
    \\[2ex]
    \begin{tabular}{c|ccc}
      $\tf{\lor}[\LogLuk[3]]$ & $1$ & $1/2$ & $0$ \\
      \hline
      $1$ & $1$ & $1$ & $1$ \\
      $1/2$ & $1$ & $1/2$ & $1/2$ \\
      $0$ & $1$ & $1/2$ & $0$
    \end{tabular}
    \quad
    \begin{tabular}{c|ccc}
      $\tf{\lif}[\LogLuk[3]]$ & $1$ & $1/2$ & $0$ \\
      \hline
      $1$ & $1$ & $1/2$ & $0$ \\
      $1/2$ & $1$ & $1$ & $1/2$  \\
      $0$ & $1$ & $1$ & $1$
    \end{tabular}
  \end{center}
\end{proof}
\begin{prop}\label{mvl:inf:luk:prop:luk-infty-m}
  $\Gamma \Entails[\LogLuk[\infty]] \psi$ असेल, तर
  सर्व~$m \ge 2$ साठी $\Gamma \Entails[\LogLuk[m]] \psi$.
\end{prop}
\begin{proof}
  सराव.
\end{proof}
\begin{prob}
  \ref{mvl:inf:luk:prop:luk-infty-m} सिद्ध करा.
\end{prob}
खरे तर आधारसूत्रांचा संच ससीम असेल, तर याचे प्रतिलोम
विधानही खरे आहे.
अनंतमूल्य \L ukasiewicz तर्कशास्त्र हे सर्वाधिक प्रचलित
अस्पष्ट तर्कशास्त्र आहे. अस्पष्ट तर्कशास्त्रावरील साहित्यात
सशर्त विधानाची व्याख्या अनेकदा $\lnot \varphi \lor \psi$ अशी
केली जाते. त्यातून अनंतमूल्य प्रबळ क्लिनी तर्कशास्त्र मिळेल.
\begin{prob}
  $(p \lif q) \lor (q \lif p)$ हे $\LogLuk[\infty]$
  मध्ये सर्वतः सत्य सूत्र आहे हे दाखवा.
\end{prob}
\section{G\"odel तर्कशास्त्रे}\label{mvl:inf:god:sec}
\begin{editorial}
  अनंतमूल्य G\"odel तर्कशास्त्रावरील हा विभाग
  संक्षिप्त ``प्रारूप'' आहे.
\end{editorial}
\begin{defn}\label{mvl:inf:god:def:goedel} अनंतमूल्य G\"odel
तर्कशास्त्र~$\LogGod[\infty]$ पुढील मॅट्रिक्सने परिभाषित होते:
\begin{enumerate}
  \item $\lfalse$, $\lnot$, $\land$, $\lor$, $\lif$ असलेली मानक विधानीय भाषा $\Lang L_0$.
  \item सत्यतामूल्यांचा संच $V_\infty$.
  \item $1$ हेच एकमेव निर्दिष्ट मूल्य आहे, म्हणजे $V^+ = \{1\}$.
  \item सत्यता-फलने पुढील सूत्रांनी दिली आहेत:
  \begin{align*}
    \tf{\lfalse} & = 0\\
    \tf{\lnot}[\LogGod](x) & = \begin{cases}
      1 & \text{जर } x =0\\
      0 & \text{इतर वेळी}
    \end{cases}\\
    \tf{\land}[\LogGod](x,y) & = \min(x,y)\\
    \tf{\lor}[\LogGod](x,y) & = \max(x,y)\\
    \tf{\lif}[\LogGod](x,y) & = \begin{cases}
      1 & \text{जर } x \le y\\
      y & \text{इतर वेळी.}
    \end{cases}
    \end{align*}
\end{enumerate}
$m$-मूल्य G\"odel तर्कशास्त्राची व्याख्या याचप्रमाणे,
फक्त $V = V_m$ घेऊन केली जाते.
\end{defn}
\begin{prop}
  \ref{mvl:thr:god:defn:goedel} मध्ये परिभाषित
  $\LogGod[3]$ आणि \ref{mvl:inf:god:def:goedel} मध्ये परिभाषित
  $\LogGod[3]$ हे एकच तर्कशास्त्र आहे.
\end{prop}
\begin{proof}
  \ref{mvl:thr:god:defn:goedel} मधील संयोजकांची
  सत्यताकोष्टके आणि \ref{mvl:inf:god:def:goedel} मधील
  समीकरणांनी निश्चित होणारी सत्यताकोष्टके यांची तुलना
  केल्यास हे दिसते:
  \begin{center}
    \begin{tabular}{c|c}
      $\tf{\lnot}[\LogGod[3]]$ & \\
      \hline
      $1$ & $0$ \\
      $1/2$ & $0$ \\
      $0$ & $1$
    \end{tabular}
    \quad
    \begin{tabular}{c|ccc}
      $\tf{\land}[\LogGod]$ & $1$ & $1/2$ & $0$ \\
      \hline
      $1$ & $1$ & $1/2$ & $0$ \\
      $1/2$ & $1/2$ & $1/2$ & $0$\\
      $0$ & $0$ & $0$ & $0$
    \end{tabular}
    \\[2ex]
    \begin{tabular}{c|ccc}
      $\tf{\lor}[\LogGod]$ & $1$ & $1/2$ & $0$ \\
      \hline
      $1$ & $1$ & $1$ & $1$ \\
      $1/2$ & $1$ & $1/2$ & $1/2$ \\
      $0$ & $1$ & $1/2$ & $0$
    \end{tabular}
    \quad
    \begin{tabular}{c|ccc}
      $\tf{\lif}[\LogGod]$ & $1$ & $1/2$ & $0$ \\
      \hline
      $1$ & $1$ & $1/2$ & $0$ \\
      $1/2$ & $1$ & $1$ & $0$  \\
      $0$ & $1$ & $1$ & $1$
    \end{tabular}
  \end{center}
\end{proof}
\begin{prop}\label{mvl:inf:god:prop:god-infty-m}
  $\Gamma \Entails[\LogGod[\infty]] \psi$ असेल, तर
  सर्व~$m \ge 2$ साठी $\Gamma \Entails[\LogGod[m]] \psi$.
\end{prop}
\begin{proof}
  सराव.
\end{proof}
\begin{prob}
  \ref{mvl:inf:god:prop:god-infty-m} सिद्ध करा.
\end{prob}
खरे तर आधारसूत्रांचा संच ससीम असेल, तर याचे प्रतिलोम
विधानही खरे आहे.
$\LogGod[3]$ प्रमाणेच $\LogGod[\infty]$ मध्येही
अंतःप्रज्ञावादी तर्कशास्त्रात वैध असलेली सर्व सूत्रे
सर्वतः सत्य आहेत; आणि आधीची सर्वतः सत्य नसलेल्या सूत्रांची
उदाहरणे $\LogGod[\infty]$ मध्येही सर्वतः सत्य नाहीत:
\begin{align*}
  & p \lor \lnot p && (p \lif q) \lif (\lnot p \lor q) \\
  & \lnot\lnot p \lif p && \lnot(\lnot p \land \lnot q) \lif (p \lor q) \\
  & ((p \lif q) \lif p) \lif p && \lnot(p \lif q) \lif (p \land \lnot q)
\end{align*}
अंतःप्रज्ञावादी तर्कशास्त्रात वैध नसलेले, पण
$\LogGod[3]$ मध्ये सर्वतः सत्य असलेले
सूत्र~$(p \lif q) \lor (q \lif p)$ हे
$\LogGod[\infty]$ मध्येही सर्वतः सत्य आहे. खरे तर
$(\varphi \lif \psi) \lor (\psi \lif \varphi)$ ही योजना
अंतःप्रज्ञावादी तर्कशास्त्रात जोडून $\LogGod[\infty]$
चे वैशिष्ट्यीकरण करता येते. हे Michael Dummett यांनी
दाखवले; म्हणून $\LogGod[\infty]$ ला अनेकदा
G\"odel--Dummett तर्कशास्त्र~$\Log{LC}$ म्हणतात.
\begin{prob}
  $(p \lif q) \lor (q \lif p)$ हे $\LogGod[\infty]$
  मध्ये सर्वतः सत्य सूत्र आहे हे दाखवा.
\end{prob}
\begin{prob}
  $(p \lif q) \lor (q \lif r) \lor (r \lif s)$ हे
  $\LogGod[3]$ मध्ये सर्वतः सत्य असूनही
  $\LogGod[\infty]$ मध्ये सर्वतः सत्य नाही हे दाखवा.
\end{prob}
\chapter{क्रमवर्ती कलन}\label{mvl:seq:chap}
\section{परिचय}\label{mvl:seq:int:sec}
अभिजात तर्कशास्त्रासाठीचे क्रमवर्ती कलन ही कार्यक्षम आणि सोपी
निष्पत्ती-पद्धत आहे. बहुमूल्य तर्कशास्त्राची व्याख्या ससीम
सत्यतामूल्ये असलेल्या मॅट्रिक्सने केली असेल, म्हणजे $V$ ससीम असेल,
तर त्यासाठी क्रमवर्ती कलन देता येते. हे कसे करावे याची कल्पना
अभिजात प्रसंगातील क्रमवर्तींचे अर्थ आणि निगमन नियमांचे रूप पाहिल्यावर
मिळते.
आता पुढील क्रमवर्ती लक्षात घ्या:
\begin{align*}
    \varphi_1, \dots, \varphi_m & \Sequent \psi_1, \dots, \psi_n
\intertext{हिचा अर्थ पुढील सूत्र म्हणून लावता येतो:}
    (\varphi_1 \land \cdots \land \varphi_m) & \lif (\psi_1 \lor \cdots \lor
\psi_n)
\end{align*}
दुसऱ्या शब्दांत, क्रमवर्ती $\Gamma \Sequent \Delta$ ची
सत्य-मूल्यांकन~$\pAssign{v}$ ने तेव्हा आणि केवळ तेव्हाच
\emph{पूर्ति होते},
जेव्हा काही $\varphi \in \Gamma$ साठी $\pValue{v}(\varphi) = \False$ असते
किंवा काही $\varphi \in \Delta$ साठी $\pValue{v}(\varphi) = \True$ असते.
या अर्थलक्षणानुसार, आरंभीच्या क्रमवर्ती $\varphi \Sequent \varphi$ ची
नेहमी पूर्ति होते; कारण एकतर $\pValue v(\varphi) = \True$ असते
किंवा $\pValue v(\varphi) = \False$ असते.
बाजूची सूत्रे $\Gamma$, $\Delta$ वगळल्यास, $\Log{LK}$ मधील
सशर्त संयोजकाचे निगमन नियम असे आहेत:
\begin{defish}
    \Axiom$ \fCenter \varphi$
    \Axiom$ \psi \fCenter $
    \RightLabel{\LeftR{\lif}}
    \BinaryInf$ \varphi \lif \psi \fCenter $
    \DisplayProof
    \hfill
    \Axiom$ \varphi \fCenter \psi$
    \RightLabel{\RightR{\lif}}
    \UnaryInf$ \fCenter \varphi \lif \psi $
    \DisplayProof
\end{defish}
क्रमवर्तीच्या वरील चिन्हार्थलक्षी अर्थलक्षणाचा उपयोग केल्यास,
$\LeftR{\lif}$ नियम असे सांगतो की $\pValue v(\varphi) = \True$
आणि $\pValue v(\psi) = \False$ असतील, तर $\pValue v(\varphi \lif \psi) = \False$
असते. त्याचप्रमाणे $\RightR{\lif}$ नियम असे सांगतो की
$\pValue v(\varphi) = \False$ किंवा $\pValue v(\psi) = \True$
असेल, तर $\pValue v(\varphi \lif \psi) = \True$ असते. प्रत्यक्षात,
येथील ही दोन्ही सशर्त विधाने द्विशर्त विधाने आहेत.
$\LeftR{\land}$ आणि $\RightR{\lor}$ नियमांच्या प्रमाणित रूपात
त्यांना अनुरूप सशर्त विधाने द्विशर्त नसतील. पण या नियमांची
अशी पर्यायी रूपे आहेत की त्यात ती द्विशर्त ठरतात:
\begin{defish}
    \Axiom$\varphi, \psi, \Gamma \fCenter \Delta$
    \RightLabel{\LeftR{\land}}
    \UnaryInf$\varphi \land \psi, \Gamma \fCenter \Delta$
    \DisplayProof
    \hfill
    \Axiom$ \Gamma \fCenter \Delta, \varphi, \psi$
    \RightLabel{\RightR{\lor}}
    \UnaryInf$ \Gamma \fCenter \Delta, \varphi \lor \psi$
    \DisplayProof
\end{defish}
ही मूलभूत कल्पना $n$-मूल्य तर्कशास्त्राला लावल्यास, दोन जागांऐवजी
$n$ जागा असलेले क्रमवर्ती कलन मिळते; प्रत्येक सत्यतामूल्यासाठी
एक जागा असते. $V = \{\False, \Undef, \True\}$ असलेल्या त्रिमूल्य
तर्कशास्त्रात क्रमवर्ती ही $\Gamma \mid \Pi \mid \Delta$ या रूपाची
अभिव्यक्ती असते. तिची सत्य-मूल्यांकन~$\pAssign v$ ने तेव्हा आणि
केवळ तेव्हाच पूर्ति होते, जेव्हा काही $\varphi \in \Gamma$ साठी
$\pValue{v}(\varphi) = \False$ असते, किंवा काही $\varphi \in \Delta$ साठी
$\pValue{v}(\varphi) = \True$ असते, किंवा काही $\varphi \in \Pi$ साठी
$\pValue{v}(\varphi) = \Undef$ असते. त्यामुळे आरंभीच्या क्रमवर्ती
$\varphi \mid \varphi \mid \varphi$ ची नेहमी पूर्ति होते.
\section{नियम आणि निष्पत्ती}\label{mvl:seq:rul:sec}
पुढे $\Gamma, \Delta, \Pi, \Lambda$ या वाक्यांच्या
सांत क्रमिका दर्शवतात असे समजू.
\begin{defn}[क्रमवर्ती]
\emph{$n$-बाजूंची क्रमवर्ती} म्हणजे पुढील रूपाची अभिव्यक्ती:
\[
\Gamma_1 \nSequent \dots \nSequent \Gamma_n
\]
येथे प्रत्येक $\Gamma_i$ ही भाषा~$\Lang L$ मधील
वाक्यांची सांत, शक्यतो रिकामी, क्रमिका असते.
\end{defn}

\begin{defn}[आरंभीची क्रमवर्ती]
\emph{$n$-बाजूंची आरंभीची क्रमवर्ती} म्हणजे भाषेतील कोणत्याही
वाक्य~$\varphi$ साठी $\varphi \nSequent \dots \nSequent \varphi$
या रूपाची $n$-बाजूंची क्रमवर्ती.

भाषेत $0$-स्थानी संयोजक~$\star$, म्हणजे विधानीय स्थिरांक,
असेल तर $\dots \nSequent \star \nSequent \dots$ ही क्रमवर्तीही
आरंभीची मानतो. यात $\tf{\star} \in V$ शी संबंधित
सत्यतामूल्याच्या जागी $\star$ असते; इतर जागा रिकाम्या असतात.
\end{defn}

$n$-मूल्य तर्कशास्त्र~$\Log{L}$ मधील प्रत्येक संयोजकासाठी,
त्या संयोजकाला $\Log{L}$ मध्ये मिळू शकणाऱ्या प्रत्येक
सत्यतामूल्यासाठी एक तार्किक नियम असतो. $\Log{L}$ च्या
$n$-बाजूंच्या क्रमवर्ती कलनातील निष्पत्त्या ही क्रमवर्तींची
सांत वृक्षरूप मांडणी असते: सर्वांत वरच्या क्रमवर्ती आरंभीच्या असतात;
आणि एखादी क्रमवर्ती एक किंवा अधिक क्रमवर्तींच्या खाली असेल,
तर ती या कलनातील निगमन नियमाने त्यांच्यापासून योग्यरीत्या
निष्पन्न झालेली असावी लागते. यात संयोजकांचे तार्किक नियम
आणि पुढील विभागात दिलेले संरचनात्मक नियम दोन्ही येतात.

\begin{defn}[प्रमेये]
वाक्य~$\varphi$ हे $n$-मूल्य तर्कशास्त्र~$\Log{L}$ चे
\emph{प्रमेय} तेव्हा असते, जेव्हा $\Log{L}$ मधील प्रत्येक निर्दिष्ट
सत्यतामूल्याशी संबंधित जागेत $\varphi$ असलेल्या $n$-बाजूंच्या
क्रमवर्तीची निष्पत्ती असते; या क्रमवर्तीतील
इतर सर्व जागा रिकाम्या असाव्यात. $\varphi$ प्रमेय असेल तर
$\Proves[\Log{L}] \varphi$ आणि नसेल तर $\Proves/[\Log{L}] \varphi$ लिहितो.
\end{defn}

\begin{defn}[निष्पन्नता]
$n$-मूल्य तर्कशास्त्र~$\Log{L}$ मध्ये वाक्य~$\varphi$
हे वाक्यांच्या संच~$\Gamma$ पासून \emph{निष्पन्न करता येण्याजोगे}
असते, $\Gamma \Proves[\Log{L}] \varphi$, तेव्हा आणि केवळ तेव्हाच,
जेव्हा एखादा ससीम उपसंच~$\Gamma_0 \subseteq \Gamma$
आणि $\Gamma_0$ मधील वाक्यांची एखादी क्रमिका~$\Gamma_0'$
अशी असते की पुढील क्रमवर्तीची निष्पत्ती असते:
\[ \Lambda_1 \nSequent \dots \nSequent \Lambda_n \]
येथे जागा~$i$ निर्दिष्ट सत्यतामूल्याशी संबंधित असेल तर
$\Lambda_i$ हे $\varphi$ असते; अन्यथा ते $\Gamma_0'$ असते.
$\varphi$ हे $\Gamma$ पासून निष्पन्न करता येण्याजोगे नसेल तर
$\Gamma \Proves/ \varphi$ असे लिहितो.
\end{defn}

उदाहरणार्थ, $3$-मूल्य \L ukasiewicz तर्कशास्त्रासाठी
$3$-बाजूंचे क्रमवर्ती कलन आहे. $3$-बाजूंच्या क्रमवर्ती
$\Gamma \nSequent \Pi \nSequent \Delta$ मध्ये $\Gamma$
हे~$\False$ शी, $\Delta$ हे~$\True$ शी आणि $\Pi$
हे~$\Undef$ शी संबंधित असतात. स्वयंसिद्धके
$\varphi \nSequent \varphi \nSequent \varphi$ ही आहेत. फक्त $\True$
निर्दिष्ट असल्यामुळे, $\Gamma$ ही सांत क्रमिका असल्यास
$\Gamma \Proves[\LogLuk[3]] \varphi$
तेव्हा आणि केवळ तेव्हाच, जेव्हा $\Gamma \nSequent \Gamma \nSequent \varphi$
या क्रमवर्तीची निष्पत्ती असते. ($\Undef$ देखील निर्दिष्ट
असते, तर $\Gamma \nSequent \varphi \nSequent \varphi$ या क्रमवर्तीची
निष्पत्ती आवश्यक झाले असते.)
अनंत आधारसूत्रसंचासाठी, आधीच्या व्याख्येनुसार त्यातील सांत
उपसंचाच्या क्रमिकेचा वापर करायचा; अनंत क्रमवर्तीचा नाही.












\section{संरचनात्मक नियम}\label{mvl:seq:str:sec}

$n$-बाजूंच्या क्रमवर्ती कलनातील संरचनात्मक नियम
अभिजात प्रसंगाप्रमाणेच कार्य करतात; मात्र ते प्रत्येक जागा~$i$
साठी स्वतंत्रपणे लागू होतात.

\begin{defish}
\begin{center}
\AxiomC{$ \Gamma_1 \nSequent \dots \nSequent \phantom{\varphi,}\Gamma_i \nSequent \dots \nSequent \Gamma_n $}
\RightLabel{$\iR{\Weakening}{i}$}
\UnaryInfC{$\Gamma_1 \nSequent \dots \nSequent \varphi, \Gamma_i \nSequent \dots \nSequent \Gamma_n$}
\DisplayProof
\\[2ex]
\AxiomC{$ \Gamma_1 \nSequent \dots \nSequent \varphi, \varphi, \Gamma_i \nSequent \dots \nSequent \Gamma_n $}
\RightLabel{$\iR{\Contraction}{i}$}
\UnaryInfC{$\Gamma_1 \nSequent \dots \nSequent \phantom{\varphi,}\varphi, \Gamma_i \nSequent
\dots \nSequent \Gamma_n$}
\DisplayProof
\\[2ex]
\AxiomC{$ \Gamma_1 \nSequent \dots \nSequent \Gamma_i, \varphi, \psi, \Gamma_i' \nSequent \dots \nSequent \Gamma_n $}
\RightLabel{$\iR{\Exchange}{i}$}
\UnaryInfC{$\Gamma_1 \nSequent \dots \nSequent \Gamma_i, \psi, \varphi, \Gamma_i' \nSequent \dots
\nSequent \Gamma_n$}
\DisplayProof
\end{center}
\end{defish}
दुर्बलीकरण, आकुंचन आणि अदलाबदल यांच्या सलग निगमनांची मालिका
अनेकदा दुहेरी निगमन रेषांनी दर्शवली जाते.
\Cut{} नियमाची अनेक रूपे असतात: क्रमवर्तीतील भिन्न जागांच्या
प्रत्येक जोडी~$i \neq j$ साठी एक रूप:
\begin{defish}
\[
\AxiomC{$ \Gamma_1 \nSequent \dots \nSequent \varphi, \Gamma_i \nSequent \dots \nSequent \Gamma_n $}
\AxiomC{$ \Delta_1 \nSequent \dots \nSequent \varphi, \Delta_j \nSequent \dots \nSequent \Delta_n $}
\RightLabel{$\iR{\Cut}{i,j}$}
\BinaryInfC{$\Gamma_1,\Delta_1 \nSequent \dots \nSequent \Gamma_n, \Delta_n$}
\DisplayProof
\]
\end{defish}












\section{निवडक तर्कशास्त्रांसाठी विधानीय नियम}\label{mvl:seq:prl:sec}

एखाद्या संयोजकासाठीचे $n$-बाजूंच्या क्रमवर्ती कलनातील निगमन नियम
फक्त त्या संयोजकाच्या वैशिष्ट्यपूर्ण सत्यता-फलनावर अवलंबून असतात.
म्हणून भिन्न तर्कशास्त्रांत एखाद्या संयोजकाची व्याख्या एकाच
सत्यता-फलनाने केली असेल, तर त्या संयोजकाचे $n$-बाजूंचे
क्रमवर्ती नियम त्या तर्कशास्त्रांत सारखेच असतात.

\subsection{$\lnot$ साठीचे नियम}

$\lnot$ साठीचे पुढील नियम \L ukasiewicz आणि क्लिनी यांच्या
तर्कशास्त्रांना तसेच त्यांच्या प्रकारांना लागू होतात.

\begin{defish}
  \begin{center}
\AxiomC{$ \Gamma \nSequent \Pi \nSequent \Delta, \varphi $}
\RightLabel{\iR{\lnot}{\False}}
\UnaryInfC{$ \lnot \varphi, \Gamma \nSequent \Pi \nSequent \Delta$}
\DisplayProof
\\[2ex]
\AxiomC{$ \Gamma \nSequent \varphi, \Pi \nSequent \Delta $}
\RightLabel{\iR{\lnot}{\Undef}}
\UnaryInfC{$\Gamma \nSequent \lnot \varphi, \Pi \nSequent \Delta$}
\DisplayProof
\\[2ex]
\AxiomC{$\varphi, \Gamma \nSequent \Pi \nSequent \Delta$}
\RightLabel{\iR{\lnot}{\True}}
\UnaryInfC{$\Gamma \nSequent \Pi \nSequent \Delta,  \lnot \varphi$}
\DisplayProof
  \end{center}
\end{defish}
$\lnot$ साठीचे पुढील नियम G\"odel तर्कशास्त्राला लागू होतात.
\begin{defish}
\AxiomC{$ \Gamma \nSequent \varphi, \Pi \nSequent \Delta, \varphi $}
\RightLabel{\iR{\lnot}{\False}[\LogGod]}
\UnaryInfC{$ \lnot \varphi, \Gamma \nSequent \Pi \nSequent \Delta$}
\DisplayProof
\hfill
\AxiomC{$\varphi, \Gamma \nSequent \Pi \nSequent \Delta$}
\RightLabel{\iR{\lnot}{\True}[\LogGod]}
\UnaryInfC{$\Gamma \nSequent \Pi \nSequent \Delta,  \lnot \varphi$}
\DisplayProof
\hfill
\end{defish}
(G\"odel तर्कशास्त्रात $\lnot \varphi$ ला~$\Undef$ हे मूल्य
कधीही मिळत नाही; म्हणून मधल्या जागेसाठी नियम नाही.)
\subsection{$\land$ साठीचे नियम}
\L ukasiewicz, प्रबळ क्लिनी आणि G\"odel यांच्या
तर्कशास्त्रांतील $\land$ साठीचे नियम पुढीलप्रमाणे आहेत.
\begin{defish}
\begin{center}
  \AxiomC{$\varphi, \psi, \Gamma \nSequent \Pi \nSequent \Delta$}
  \RightLabel{$\iR{\land}{\False}$}
  \UnaryInfC{$\varphi \land \psi, \Gamma \nSequent \Pi \nSequent \Delta$}
  \DisplayProof
\\[2ex]
  \AxiomC{$\Gamma \nSequent \varphi, \Pi \nSequent \varphi, \Delta$}
  \AxiomC{$\Gamma \nSequent \psi, \Pi \nSequent \psi, \Delta$}
  \AxiomC{$\Gamma \nSequent \varphi, \psi, \Pi \nSequent \Delta$}
  \RightLabel{$\iR\land\Undef$}
  \TrinaryInfC{$\Gamma \nSequent \varphi \land \psi, \Pi \nSequent \Delta$}
  \DisplayProof
  \\[2ex]
  \AxiomC{$\Gamma \nSequent \Pi \nSequent \Delta, \varphi$}
  \AxiomC{$\Gamma \nSequent \Pi \nSequent \Delta, \psi$}
  \RightLabel{$\iR\land\True$}
  \BinaryInfC{$\Gamma \nSequent \Pi \nSequent \Delta, \varphi \land \psi$}
  \DisplayProof
\end{center}
\end{defish}
\subsection{$\lor$ साठीचे नियम}
\L ukasiewicz, प्रबळ क्लिनी आणि G\"odel यांच्या
तर्कशास्त्रांतील $\lor$ साठीचे नियम पुढीलप्रमाणे आहेत.
\begin{defish}
\begin{center}
 \AxiomC{$\varphi, \Gamma \nSequent \Pi \nSequent \Delta$}
  \AxiomC{$\psi, \Gamma \nSequent \Pi \nSequent \Delta$}
  \RightLabel{$\iR\lor\False$}
  \BinaryInfC{$\varphi \lor \psi, \Gamma \nSequent \Pi \nSequent \Delta$}
  \DisplayProof
\\[2ex]
  \AxiomC{$\varphi, \Gamma \nSequent \varphi, \Pi \nSequent \Delta$}
  \AxiomC{$\psi, \Gamma \nSequent \psi, \Pi \nSequent \Delta$}
  \AxiomC{$\Gamma \nSequent \varphi, \psi, \Pi \nSequent \Delta$}
  \RightLabel{$\iR\lor\Undef$}
  \TrinaryInfC{$\Gamma \nSequent \varphi \lor \psi, \Pi \nSequent \Delta$}
  \DisplayProof
  \\[2ex]
  \AxiomC{$\Gamma \nSequent \Pi \nSequent \Delta, \varphi, \psi$}
  \RightLabel{$\iR\lor\True$}
  \UnaryInfC{$\Gamma \nSequent \Pi \nSequent \Delta, \varphi \lor \psi$}
  \DisplayProof
\end{center}
\end{defish}
\subsection{$\lif$ साठीचे नियम}
\L ukasiewicz तर्कशास्त्रातील $\lif$ साठीचे नियम पुढीलप्रमाणे आहेत.
\begin{defish}
\begin{center}
  \AxiomC{$\Gamma \nSequent \Pi \nSequent \Delta, \varphi$}
  \AxiomC{$\psi, \Gamma \nSequent \Pi \nSequent \Delta$}
  \RightLabel{$\iR\lif\False[\LogLuk[3]]$}
  \BinaryInfC{$\varphi \lif \psi, \Gamma \nSequent \Pi \nSequent \Delta$}
  \DisplayProof
  \\[2ex]
  \AxiomC{$\Gamma \nSequent \varphi, \psi, \Pi \nSequent \Delta$}
  \AxiomC{$\psi, \Gamma \nSequent \Pi \nSequent \Delta, \varphi$}
  \RightLabel{$\iR\lif\Undef[\LogLuk[3]]$}
  \BinaryInfC{$\Gamma \nSequent \varphi \lif \psi, \Pi \nSequent \Delta$}
  \DisplayProof
\\[2ex]
  \AxiomC{$\varphi, \Gamma \nSequent \psi, \Pi \nSequent \Delta, \psi$}
  \AxiomC{$\varphi, \Gamma \nSequent \varphi, \Pi \nSequent \Delta, \psi$}
  \RightLabel{$\iR{\lif}{\True}[\LogLuk[3]]$}
  \BinaryInfC{$\Gamma \nSequent \Pi \nSequent \Delta, \varphi \lif \psi$}
  \DisplayProof
\end{center}
\end{defish}
प्रबळ क्लिनी तर्कशास्त्रातील $\lif$ साठीचे नियम पुढीलप्रमाणे आहेत.
\begin{defish}
\begin{center}
  \AxiomC{$\Gamma \nSequent \Pi \nSequent \Delta, \varphi$}
  \AxiomC{$\psi, \Gamma \nSequent \Pi \nSequent \Delta$}
  \RightLabel{$\iR\lif\False[\LogKs]$}
  \BinaryInfC{$\varphi \lif \psi, \Gamma \nSequent \Pi \nSequent \Delta$}
  \DisplayProof
  \\[2ex]
  \AxiomC{$\psi, \Gamma \nSequent \psi, \Pi \nSequent \Delta$}
  \AxiomC{$\Gamma \nSequent \varphi, \psi, \Pi \nSequent \Delta$}
  \AxiomC{$\Gamma \nSequent \varphi, \Pi \nSequent \Delta, \varphi$}
  \RightLabel{$\iR\lif\Undef[\LogKs]$}
  \TrinaryInfC{$\Gamma \nSequent \varphi \lif \psi, \Pi \nSequent \Delta$}
  \DisplayProof
\\[2ex]
  \AxiomC{$\varphi, \Gamma \nSequent \Pi \nSequent \Delta, \psi$}
  \RightLabel{$\iR{\lif}{\True}[\LogKs]$}
  \UnaryInfC{$\Gamma \nSequent \Pi \nSequent \Delta, \varphi \lif \psi$}
  \DisplayProof
\end{center}
\end{defish}
G\"odel तर्कशास्त्रातील $\lif$ साठीचे नियम पुढीलप्रमाणे आहेत.
\begin{defish}
\begin{center}
  \AxiomC{$\Gamma \nSequent \varphi, \Pi \nSequent \Delta, \varphi$}
  \AxiomC{$\psi, \Gamma \nSequent \Pi \nSequent \Delta$}
  \RightLabel{$\iR\lif\False[\LogGod[3]]$}
  \BinaryInfC{$\varphi \lif \psi, \Gamma \nSequent \Pi \nSequent \Delta$}
  \DisplayProof
  \\[2ex]
  \AxiomC{$\Gamma \nSequent \psi, \Pi \nSequent \Delta$}
  \AxiomC{$\Gamma \nSequent \Pi \nSequent \Delta, \varphi$}
  \RightLabel{$\iR\lif\Undef[\LogGod[3]]$}
  \BinaryInfC{$\Gamma \nSequent \varphi \lif \psi, \Pi \nSequent \Delta$}
  \DisplayProof
\\[2ex]
  \AxiomC{$\varphi, \Gamma \nSequent \psi, \Pi \nSequent \Delta, \psi$}
  \AxiomC{$\varphi, \Gamma \nSequent \varphi, \Pi \nSequent \Delta, \psi$}
  \RightLabel{$\iR{\lif}{\True}[\LogGod[3]]$}
  \BinaryInfC{$\Gamma \nSequent \Pi \nSequent \Delta, \varphi \lif \psi$}
  \DisplayProof
\end{center}
\end{defish}
\begin{sidewaysfigure}
  \begin{center}
  \AxiomC{$A \nSequent A \nSequent A$}
  \RightLabel{\iR \Weakening \True}
  \UnaryInfC{$A \nSequent A \nSequent B, A$}
  \RightLabel{\iR \Weakening \Undef}
  \UnaryInfC{$A \nSequent B, A \nSequent B, A$}
  \RightLabel{\iR \Weakening \Undef}
  \UnaryInfC{$A \nSequent A, B, A \nSequent B, A$}
  \AxiomC{$A \nSequent A \nSequent A$}
  \RightLabel{\iR \Weakening \True}
  \UnaryInfC{$A \nSequent A \nSequent A, A$}
  \RightLabel{\iR \Weakening \True}
  \UnaryInfC{$A \nSequent A \nSequent B, A, A$}
  \RightLabel{\iR \Weakening \False}
  \UnaryInfC{$B, A \nSequent A \nSequent B, A, A$}
  \RightLabel{$\iR\lif\Undef$}
  \BinaryInfC{$A \nSequent A \lif B, A \nSequent B, A$}
  \AxiomC{$B \nSequent B \nSequent B$}
  \RightLabel{\iR \Weakening \Undef}
  \UnaryInfC{$B \nSequent A, B \nSequent B$}
  \RightLabel{\iR \Exchange \Undef}
  \UnaryInfC{$B \nSequent B, A \nSequent B$}
  \RightLabel{\iR \Weakening \Undef}
  \UnaryInfC{$B \nSequent A, B, A \nSequent B$}
  \RightLabel{\iR \Weakening \False}
  \UnaryInfC{$A, B \nSequent A, B, A \nSequent B$}
  \RightLabel{\iR \Exchange \False}
  \UnaryInfC{$B, A \nSequent A, B, A \nSequent B$}
  \AxiomC{$A \nSequent A \nSequent A$}
  \RightLabel{\iR \Weakening \True}
  \UnaryInfC{$A \nSequent A \nSequent B, A$}
  \RightLabel{\iR \Weakening \False}
  \UnaryInfC{$B, A \nSequent A \nSequent B, A$}
  \RightLabel{\iR \Weakening \False}
  \UnaryInfC{$B, B, A \nSequent A \nSequent B, A$}
  \RightLabel{$\iR\lif\Undef$}
  \BinaryInfC{$B, A \nSequent A \lif B, A \nSequent B$}
  \RightLabel{$\iR\lif\False$}
  \BinaryInfC{$A \lif B, A \nSequent A \lif B, A \nSequent B$}
  \DisplayProof
  \end{center}
  \caption{$\LogLuk[3]$ मधील निष्पत्तीचे उदाहरण}
\end{sidewaysfigure}
\part{सामान्य मोडल तर्कशास्त्रे}\label{nml:part}
\begin{editorial}
  या भागात सामान्य मोडल तर्कशास्त्रांच्या अधिसिद्धांताचा विचार केला आहे.
  सध्या यात मूलभूत मोडल तर्कशास्त्राच्या अभिजात अनुरूपता
  सिद्धांतावरील Aldo Antonelli यांची टिपणे आहेत.
\end{editorial}
\chapter{विन्यासमीमांसा आणि चिन्हार्थमीमांसा}\label{nml:syn:chap}
\section{परिचय}\label{nml:syn:int:sec}
मोडल तर्कशास्त्र \emph{मोडल विधाने} आणि त्यांच्यातील तार्किक
निष्पन्नता-संबंध यांचा अभ्यास करते. मोडल विधानांची पुढील
उदाहरणे पाहा:
\begin{enumerate}
\item $2+2=4$ हे अनिवार्य आहे.
\item उद्या पाऊस पडणे अनिवार्यपणे संभाव्य आहे.
\item जर~$\varphi$ अनिवार्यपणे संभाव्य असेल, तर~$\varphi$ संभाव्य आहे.
\end{enumerate}
संभाव्यता आणि अनिवार्यता या एकमेव मोडलता नाहीत: इतर एकस्थानी
संयोजकही मोडलता मानले जातात. उदाहरणार्थ, ``$\varphi$ असे असायला हवे,''
``पुढे~$\varphi$ घडेल,'' ``Dana ला~$\varphi$ माहीत आहे,'' किंवा
``Dana चा~$\varphi$ वर विश्वास आहे.''
मोडल तर्कशास्त्र प्रथम अरिस्टॉटल यांच्या \emph{De
  Interpretatione} मध्ये दिसते. अनिवार्य असल्यास संभाव्य असते,
पण उलट नाही; संभाव्यता आणि अनिवार्यता एकमेकींच्या साहाय्याने
परिभाषित करता येतात; $\varphi \land \psi$ संभाव्यतः सत्य असेल तर
$\varphi$ संभाव्यतः सत्य असते आणि $\psi$ संभाव्यतः सत्य असते,
पण उलट नाही; तसेच $\varphi \to \psi$ अनिवार्य असेल, तर $\varphi$
अनिवार्य असताना $\psi$ ही अनिवार्य असते, या गोष्टी प्रथम
त्यांच्या लक्षात आल्या.
मोडल तर्कशास्त्राकडे पाहण्याचा पहिला आधुनिक दृष्टिकोन
C.~I. Lewis यांच्या कार्यातून पुढे आला. Lewis आणि Langford
यांच्या \emph{Symbolic Logic} (1932) या ग्रंथात या कार्याची
परिणती झाली. वास्तविक अभिव्यंजन वापरून अभिव्यंजनाचे प्रतिनिधित्व
करण्याबाबत Lewis आणि Langford समाधानी नव्हते: $\varphi \lif \psi$
हे ``$\varphi$ पासून $\psi$ निष्पन्न होते'' याचे अपुरे प्रतिनिधित्व आहे.
त्याऐवजी ``अनिवार्यतः, जर $\varphi$ तर $\psi$'' असे अभिव्यंजनाचे
लक्षण मांडून त्यांनी ते $\varphi \strictif \psi$ असे दर्शवले.
या विविध गुणधर्मांची मांडणी करताना Lewis यांनी पाच भिन्न
मोडल प्रणाली ओळखल्या: \Log{S1}, \ldots, \Log{S4}, \Log{S5};
त्यांतील शेवटच्या दोन अजूनही वापरात आहेत.
Lewis आणि Langford यांचा दृष्टिकोन पूर्णतः \emph{विन्यासमीमांसक}
होता: त्यांनी योग्य वाटणारी स्वयंसिद्धके व नियम निवडले आणि
त्यांच्या साहाय्याने काय सिद्ध करता येते याचा अभ्यास केला.
चिन्हार्थमीमांसक दृष्टिकोन बराच काळ उपलब्ध नव्हता.
नंतर Rudolf Carnap यांनी \emph{Meaning and Necessity} (1947)
मध्ये \emph{अवस्था-वर्णन} ही संकल्पना वापरून पहिला प्रयत्न केला.
अवस्था-वर्णन म्हणजे त्या वर्णनात ``सत्य'' असलेल्या आण्विक
वाक्यांचा संग्रह. कोणत्याही वाक्य~$\varphi$ साठी सत्यतेची व्याख्या
विस्तारित केल्यानंतर, $\varphi$ सर्व अवस्था-वर्णनांत सत्य असेल
तर ते \emph{अनिवार्यतः सत्य} आहे, अशी Carnap यांची व्याख्या आहे.
मात्र त्यांचा दृष्टिकोन \emph{पुनरावृत्त} मोडलता हाताळू शकत
नव्हता: ``संभाव्यतः अनिवार्यतः \ldots संभाव्यतः $\varphi$''
या रूपाची वाक्ये नेहमी सर्वांत आतील मोडलतेपुरती सीमित होत.
मोडल चिन्हार्थमीमांसेतील मोठी प्रगती Saul Kripke यांच्या
``A Completeness Theorem in Modal Logic'' (JSL 1959) या
लेखामुळे झाली. एखादे विधान ``सर्व संभाव्य जगांत'' सत्य
असेल तर ते अनिवार्यतः सत्य असते, या लाइब्निझ यांच्या कल्पनेवर
Kripke यांनी आपले कार्य आधारले. परंतु त्या साध्या अर्थलक्षणातही
Carnap यांच्या दृष्टिकोनासारखीच अडचण आहे: जग~$w$ येथे एखादे
विधान अनिवार्यतः सत्य आहे की नाही, हे $w$ च्या निवडीवर
(किंवा अवस्था-वर्णन~$s$ वर) मुळीच अवलंबून नसते.
म्हणून Kripke यांनी जगे \emph{प्राप्यता संबंध}~$R$ ने
संबंधित आहेत असे मानले; आणि ``अनिवार्यतः $\varphi$'' या रूपाचे
विधान जग~$w$ येथे तेव्हा आणि केवळ तेव्हाच सत्य असते,
जेव्हा $w$ पासून \emph{प्राप्य असलेल्या} प्रत्येक जग~$w'$
येथे $\varphi$ सत्य असते. या दृष्टिकोनाचे कोणतेही रूप वापरणाऱ्या
चिन्हार्थमीमांसेला क्रिप्के चिन्हार्थमीमांसा म्हणतात;
तिच्यामुळे मोडल तर्कशास्त्रांचा, अनेकवचनात, झपाट्याने विकास
होऊ शकला.
क्रिप्के चिन्हार्थमीमांसेनुसार अर्थ लावल्यास, मोडल तर्कशास्त्र
\emph{संबंधी रचना} ``आतून'' कशा दिसतात हे दाखवते. संबंधी
रचना म्हणजे द्विस्थानी संबंध असलेला संच; उदाहरणार्थ,
वर्गातील विद्यार्थ्यांचा सामाजिक सुरक्षा क्रमांकानुसार
लावलेला संच ही संबंधी रचना आहे. पण संबंधी रचना अनेक
क्षेत्रांत आढळतात. जगाच्या अवस्थांमधील सापेक्ष संभाव्यतेखेरीज,
एखाद्या व्यक्तीच्या ज्ञानाच्या अवस्था ज्ञानविषयक संभाव्यतेने
संबंधित असू शकतात किंवा गतिशील प्रणालीच्या अवस्थांमध्ये
स्थित्यंतरे असू शकतात. या सर्वांचे प्रतिरूपण मोडल तर्कशास्त्राने
करता येते: पहिल्या प्रकारातून नेहमीचे सत्यविषयक मोडल तर्कशास्त्र,
तर इतर प्रकारांतून ज्ञानविषयक तर्कशास्त्र, गतिशील तर्कशास्त्र
इत्यादी मिळतात.
मोडल तर्कशास्त्रज्ञ ``अनुरूपता सिद्धांत'' म्हणतात त्या एका
विशिष्ट दृष्टीकोनावर आपण लक्ष केंद्रित करतो. Kripke यांच्या
आरंभीच्या महत्त्वाच्या निष्कर्षांपैकी एक असा की प्राप्यता
संबंध~$R$ चे अनेक गुणधर्म, उदाहरणार्थ संक्रमणशीलता किंवा
सममिती, योग्य ``मोडल रूपबंधां''च्या साहाय्याने
\emph{मोडल भाषेतच} व्यक्त करता येतात. उदाहरणार्थ,
$R$ ची परावर्तिता ``जर अनिवार्यतः~$\varphi$, तर~$\varphi$'' या
रूपबंधाशी ``अनुरूप'' आहे, असे मोडल तर्कशास्त्रज्ञ म्हणतात.
\Ax{D}, \Ax{T}, \Ax{B}, \Ax{4} आणि~\Ax{5} हे रूपबंध
एकत्र करून मिळणाऱ्या अभिजात मोडल तर्कशास्त्रांपैकी
काहींच्या, उदाहरणार्थ \Log{S4} आणि \Log{S5}, अनुरूपता
सिद्धांताचा आपण मुख्यतः अभ्यास करू.
\section{मूलभूत मोडल तर्कशास्त्राची भाषा}\label{nml:syn:lan:sec}
\begin{defn}
मोडल तर्कशास्त्राच्या मूलभूत भाषेत पुढील गोष्टी असतात:
\begin{enumerate}
  \item असत्यता साठीचा विधानीय स्थिरांक~$\lfalse$.
  \item सत्यता साठीचा विधानीय स्थिरांक~$\ltrue$.
  \item विधानीय चलांचा गणनीय अनंत संच: $\Obj
    p_0$, $\Obj p_1$, $\Obj p_2$, \dots
  \item विधानीय संयोजक: \startycommalist
  \ycomma $\lnot$ (निषेधन)
  \ycomma $\land$ (संधी)
  \ycomma $\lor$ (विकल्प)
  \ycomma $\lif$ (शर्त)
  \ycomma $\liff$ (द्विशर्त).
  \item मोडल परिकर्मी $\Box$.
  \item मोडल परिकर्मी $\Diamond$.
\end{enumerate}
\end{defn}
\begin{defn}
मूलभूत मोडल भाषेतील \emph{सूत्रे} पुढीलप्रमाणे
  विगमनाने परिभाषित होतात:
\begin{enumerate}
\item $\lfalse$ हे आण्विक सूत्र आहे.
\item $\ltrue$ हे आण्विक सूत्र आहे.
\item प्रत्येक विधानीय चल $\Obj p_i$ हे (आण्विक) सूत्र आहे.
\item $\varphi$ हे सूत्र असेल, तर $\lnot \varphi$
  हे सूत्र आहे.
\item $\varphi$ आणि $\psi$ ही सूत्रे असतील, तर $(\varphi \land
  \psi)$ हे सूत्र आहे.
\item $\varphi$ आणि $\psi$ ही सूत्रे असतील, तर $(\varphi \lor \psi)$
  हे सूत्र आहे.
\item $\varphi$ आणि $\psi$ ही सूत्रे असतील, तर $(\varphi \lif \psi)$
  हे सूत्र आहे.
\item $\varphi$ आणि $\psi$ ही सूत्रे असतील, तर $(\varphi \liff \psi)$
  हे सूत्र आहे.
\item $\varphi$ हे सूत्र असेल, तर $\Box \varphi$
  हे सूत्र आहे.
\item $\varphi$ हे सूत्र असेल, तर $\Diamond \varphi$
  हे सूत्र आहे.
\item यांशिवाय दुसरे काहीही सूत्र नाही.
\end{enumerate}
\end{defn}
सूत्र~$\varphi$ मध्ये $\Box$ किंवा $\Diamond$ नसेल,
तर त्याला \emph{मोडलतारहित} म्हणतो.
\section{एकाच वेळी केलेले आदेशन}\label{nml:syn:sub:sec}
सूत्र~$\varphi$ मधील विधानीय चल च्या सर्व
आस्थितींच्या जागी दुसरे सूत्र घालून मिळणारे सूत्र हा
$\varphi$ चा \emph{दृष्टांत} असतो. वैधता आणि निष्पन्नता
यांवर चर्चा करताना सूत्रांच्या दृष्टांतांचा वारंवार
उल्लेख करावा लागेल. म्हणून ही संकल्पना नेमकी परिभाषित
करणे उपयुक्त आहे.
\begin{defn}\label{nml:syn:sub:def:subst-inst}
  $\varphi$ हे असे मोडल सूत्र असू द्या की त्यातील सर्व
  विधानीय चले ही $p_1$, \dots, $p_n$ यांपैकीच
  आहेत; तसेच $\theta_1$, \dots, $\theta_n$ हीदेखील मोडल सूत्रे
  असू द्या. मग
  $\SSubst{\varphi}{\subst{\theta_1}{p_1}, \dots, \subst{\theta_n}{p_n}}$
  म्हणजे $\varphi$ मध्ये प्रत्येक $p_i$ च्या जागी $\theta_i$ चे
  एकाच वेळी आदेशन करून मिळणारे सूत्र, अशी व्याख्या करतो.
  औपचारिकरीत्या ही $\varphi$ वरील विगमनाने केलेली व्याख्या आहे:
  \begin{enumerate}
    \item \indcase{\varphi}{\lfalse}{
        $\SSubst{\indfrm}{\subst{\theta_1}{p_1}, \dots,
          \subst{\theta_n}{p_n}}$ हे $\lfalse$ असते.}
    \item \indcase{\varphi}{\ltrue}{
        $\SSubst{\indfrm}{\subst{\theta_1}{p_1}, \dots,
          \subst{\theta_n}{p_n}}$ हे $\ltrue$ असते.}
    \item \indcase{\varphi}{q}{$\SSubst{\indfrm}{\subst{\theta_1}{p_1}, \dots,
        \subst{\theta_n}{p_n}}$ हे $q$ असते, जर $i
      = 1$, \dots,~$n$ साठी $q \not\ident p_i$ असेल.}
    \item \indcase{\varphi}{p_i}{$\SSubst{\indfrm}{\subst{\theta_1}{p_1},
        \dots, \subst{\theta_n}{p_n}}$ हे $\theta_i$ असते.}
    \item \indcase{\varphi}{\lnot \psi}{
        $\SSubst{\indfrm}{\subst{\theta_1}{p_1}, \dots, \subst{\theta_n}{p_n}}$ हे
        $\lnot \SSubst{\psi}{\subst{\theta_1}{p_1}, \dots, \subst{\theta_n}{p_n}}$ असते.}
    \item \indcase{\varphi}{(\psi \land
        \chi)}{$\SSubst{\indfrm}{\subst{\theta_1}{p_1}, \dots,
          \subst{\theta_n}{p_n}}$ हे पुढील सूत्र असते:
        \[(\SSubst{\psi}{\subst{\theta_1}{p_1}, \dots, \subst{\theta_n}{p_n}} \land
          \SSubst{\chi}{\subst{\theta_1}{p_1}, \dots, \subst{\theta_n}{p_n}}).\]}
    \item \indcase{\varphi}{(\psi \lor
          \chi)}{$\SSubst{\indfrm}{\subst{\theta_1}{p_1}, \dots,
          \subst{\theta_n}{p_n}}$ हे पुढील सूत्र असते:
        \[(\SSubst{\psi}{\subst{\theta_1}{p_1}, \dots, \subst{\theta_n}{p_n}} \lor
          \SSubst{\chi}{\subst{\theta_1}{p_1}, \dots, \subst{\theta_n}{p_n}}).\]}
    \item \indcase{\varphi}{(\psi \lif
          \chi)}{$\SSubst{\indfrm}{\subst{\theta_1}{p_1}, \dots,
          \subst{\theta_n}{p_n}}$ हे पुढील सूत्र असते:
        \[(\SSubst{\psi}{\subst{\theta_1}{p_1}, \dots, \subst{\theta_n}{p_n}} \lif
          \SSubst{\chi}{\subst{\theta_1}{p_1}, \dots, \subst{\theta_n}{p_n}}).\]}
    \item \indcase{\varphi}{(\psi \liff
          \chi)}{$\SSubst{\indfrm}{\subst{\theta_1}{p_1}, \dots,
          \subst{\theta_n}{p_n}}$ हे पुढील सूत्र असते:
        \[(\SSubst{\psi}{\subst{\theta_1}{p_1}, \dots, \subst{\theta_n}{p_n}} \liff
          \SSubst{\chi}{\subst{\theta_1}{p_1}, \dots, \subst{\theta_n}{p_n}}).\]}
    \item \indcase{\varphi}{\Box \psi}{
        $\SSubst{\indfrm}{\subst{\theta_1}{p_1}, \dots, \subst{\theta_n}{p_n}}$ म्हणजे
        $\Box
        \SSubst{\psi}{\subst{\theta_1}{p_1}, \dots, \subst{\theta_n}{p_n}}.$}
    \item \indcase{\varphi}{\Diamond \psi}{
        $\SSubst{\indfrm}{\subst{\theta_1}{p_1}, \dots, \subst{\theta_n}{p_n}}$ म्हणजे
        $\Diamond
        \SSubst{\psi}{\subst{\theta_1}{p_1}, \dots, \subst{\theta_n}{p_n}}.$}
  \end{enumerate}
  सूत्र $\SSubst{\varphi}{\subst{\theta_1}{p_1}, \dots,
    \subst{\theta_n}{p_n}}$ याला $\varphi$ चा \emph{आदेशन-दृष्टांत}
  म्हणतात.
\end{defn}

\begin{ex}
  समजा $\varphi$ हे $p_1 \lif \Box(p_1 \land p_2)$,
  $\theta_1$ हे $\Diamond(p_2 \lif p_3)$ आणि $\theta_2$ हे
  $\lnot\Box p_1$ आहे. मग
  $\SSubst{\varphi}{\subst{\theta_1}{p_1}, \subst{\theta_2}{p_2}}$ म्हणजे
  \begin{align*}
    \Diamond(p_2 \lif p_3) &
    \lif \Box(\Diamond(p_2 \lif p_3) \land \lnot\Box p_1)
    \intertext{तर $\SSubst{\varphi}{\subst{\theta_2}{p_1}, \subst{\theta_1}{p_2}}$ म्हणजे}
    \lnot\Box p_1 &
    \lif \Box(\lnot\Box p_1 \land \Diamond(p_2 \lif p_3))
    \intertext{लक्षात घ्या: एकाच वेळी केलेले आदेशन
      सर्वसाधारणपणे क्रमाने केलेल्या आदेशनासारखे नसते.
      उदाहरणार्थ, वरील
      $\SSubst{\varphi}{\subst{\theta_1}{p_1}, \subst{\theta_2}{p_2}}$ ची
      $\Subst{(\Subst{\varphi}{\theta_1}{p_1})}{\theta_2}{p_2}$ शी तुलना करा; नंतरचे असे आहे:}
    & \Subst{(\Diamond(p_2 \lif p_3) \lif
      \Box(\Diamond(p_2 \lif p_3) \land p_2))}{\lnot\Box p_1}{p_2},
      \text{ म्हणजे}\\
    \Diamond(\lnot\Box p_1 \lif p_3) &
    \lif \Box(\Diamond(\lnot\Box p_1
    \lif p_3) \land \lnot\Box p_1)\\
    \intertext{आणि $\Subst{(\Subst{\varphi}{\theta_2}{p_2})}{\theta_1}{p_1}$ शीही तुलना करा:}
    & \Subst{(p_1 \lif \Box(p_1 \land \lnot\Box p_1))}
      {\Diamond(p_2 \lif p_3)}{p_1}, \text{ म्हणजे}\\
    \Diamond(p_2 \lif p_3) &
    \lif \Box(\Diamond(p_2
    \lif p_3) \land \lnot\Box\Diamond(p_2 \lif p_3)).
  \end{align*}
\end{ex}












\section{संबंधी प्रतिमाने}\label{nml:syn:rel:sec}

सामान्य मोडल तर्कशास्त्रांच्या चिन्हार्थमीमांसेची मूलभूत संकल्पना
म्हणजे \emph{संबंधी प्रतिमान}. त्यात द्विस्थानी ``प्राप्यता संबंधाने''
संबंधित जगांचा संच आणि कोणते विधानीय चले कोणत्या
जगांत ``सत्य'' मानायची हे ठरवणारे मूल्यांकन असते.

\begin{defn}
  मूलभूत मोडल भाषेचे \emph{प्रतिमान} म्हणजे त्रिकूट $\mModel{M}
  = \tuple{W, R, V}$; येथे
  \begin{enumerate}
  \item $W$ हा ``जगांचा'' रिकामा नसलेला संच आहे,
  \item $R$ हा~$W$ वरील द्विस्थानी प्राप्यता संबंध आहे, आणि
  \item $V$ हे प्रत्येक विधानीय चल~$p$ ला संभाव्य जगांचा संच~$V(p)$ नेमून देणारे फलन आहे.
  \end{enumerate}
  $Rww'$ असेल तेव्हा $w'$ हे~$w$ पासून \emph{प्राप्य}
  आहे असे म्हणतो. $w \in V(p)$ असेल तेव्हा $p$ हे~$w$ येथे
  \emph{सत्य} आहे असे म्हणतो.
\end{defn}

संबंधी चिन्हार्थमीमांसेचा मोठा फायदा असा की प्रतिमाने
\ref{nml:syn:rel:fig:simple} मधील साध्या आकृतीसारख्या आकृत्यांनी दाखवता
येतात. जगांसाठी गाठी वापरतात; $w$ पासून~$w'$ कडे बाण
असेल तेव्हा आणि केवळ तेव्हाच जग~$w'$ हे~$w$ पासून प्राप्य
असते. शिवाय $w \in V(p)$ असेल तर गाठीला (जगाला)~$\mTrue{p}$,
अन्यथा~\mFalse{p} अशी खूण करतो. \ref{nml:syn:rel:fig:simple} मध्ये
$W = \{w_1, w_2, w_3\}$, $R = \{\tuple{w_1, w_2},\tuple{w_1,
  w_3}\}$, $V(p) = \{w_1, w_2\}$ आणि $V(q) = \{w_2\}$
असलेले प्रतिमान दाखवले आहे.

\begin{figure}
  \begin{center}
    \begin{tikzpicture}[modal]
      \node[world] (w1) [label={[align=right]right:\mTrue{p}\\ \mFalse{q}}]{$w_1$};
      \node[world] (w2) [label={[align=right]right:\mTrue{p}\\ \mTrue{q}},
        above right=of w1]{$w_2$};
      \node[world] (w3) [label={[align=right]right:\mFalse{p}\\ \mFalse{q}},
        below right=of w1] {$w_3$};
      \draw[->] (w1) to (w2);
      \draw[->] (w1) to (w3);
    \end{tikzpicture}
  \end{center}
  \caption{साधे प्रतिमान.}
  \label{nml:syn:rel:fig:simple}
\end{figure}












\section{एखाद्या जगात सत्य असणे}\label{nml:syn:trw:sec}

प्रत्येक मोडल प्रतिमान त्यातील कोणती मोडल सूत्रे कोणत्या
जगांत सत्य मानायची हे ठरवते. ``प्रतिमान~$\mModel{M}$ मध्ये
सूत्र~$\varphi$ जग~$w$ येथे सत्य आहे'' हा संबंधी
चिन्हार्थमीमांसेतील मूलभूत संबंध आहे. हा संबंध विगमनाने
परिभाषित केला जातो; मोडल नसलेल्या परिकर्मींसाठी तो
सत्यताकोष्टकांनी दिलेल्या नेहमीच्या लक्षणाशी जुळतो.

\begin{defn}\label{nml:syn:trw:defn:mmodels}
  प्रतिमान~$\mModel M$ मध्ये \emph{सूत्र~$\varphi$ हे~$w$ येथे
  सत्य असणे}, म्हणजे $\mSat{M}{\varphi}[w]$, पुढीलप्रमाणे
  विगमनाने परिभाषित करतो:
  \begin{enumerate}
  \item \indcase{\varphi}{\lfalse}{$\mSat{M}{\lfalse}[w]$ कधीच नसते}.
  \item \indcase{\varphi}{\ltrue}{$\mSat{M}{\ltrue}[w]$ नेहमी असते}.
  \item $\mSat{M}{p}[w]$ तेव्हा आणि केवळ तेव्हाच, जेव्हा $w \in V(p)$.
  \item \indcase{\varphi}{\lnot \psi}{$\mSat{M}{\indfrm}[w]$
    तेव्हा आणि केवळ तेव्हाच, जेव्हा $\mSat/{M}{\psi}[w]$}.
  \item \indcase{\varphi}{(\psi \land \chi)}{$\mSat{M}{\indfrm}[w]$
    तेव्हा आणि केवळ तेव्हाच, जेव्हा $\mSat{M}{\psi}[w]$ आणि
    $\mSat{M}{\chi}[w]$}.
  \item \indcase{\varphi}{(\psi \lor \chi)}{$\mSat{M}{\indfrm}[w]$
    तेव्हा आणि केवळ तेव्हाच, जेव्हा $\mSat{M}{\psi}[w]$ किंवा
    $\mSat{M}{\chi}[w]$} (किंवा दोन्ही).
  \item \indcase{\varphi}{(\psi \lif \chi)}{$\mSat{M}{\indfrm}[w]$
    तेव्हा आणि केवळ तेव्हाच, जेव्हा $\mSat/{M}{\psi}[w]$ किंवा
    $\mSat{M}{\chi}[w]$}.
  \item \indcase{\varphi}{(\psi \liff \chi)}{$\mSat{M}{\indfrm}[w]$
    तेव्हा आणि केवळ तेव्हाच, जेव्हा $\mSat{M}{\psi}[w]$ आणि
    $\mSat{M}{\chi}[w]$ ही दोन्ही असतात, किंवा
    $\mSat{M}{\psi}[w]$ आणि $\mSat{M}{\chi}[w]$ यांपैकी कोणतेही नसते}.
  \item\label{nml:syn:trw:defn:sub:mmodels-box}
    \indcase{\varphi}{\Box \psi}{$\mSat{M}{\indfrm}[w]$ तेव्हा आणि
    केवळ तेव्हाच, जेव्हा $Rww'$ असलेल्या प्रत्येक $w' \in W$
    साठी $\mSat{M}{\psi}[w']$.}
  \item\label{nml:syn:trw:defn:sub:mmodels-diamond}
    \indcase{\varphi}{\Diamond \psi}{$\mSat{M}{\indfrm}[w]$ तेव्हा आणि
    केवळ तेव्हाच, जेव्हा $Rww'$ असलेल्या किमान एका $w' \in W$
    साठी $\mSat{M}{\psi}[w']$.}
  \end{enumerate}
\end{defn}

\ref{nml:syn:trw:defn:sub:mmodels-box} या खंडानुसार, $Rww'$ असलेले
एकही~$w'$ नसेल तेव्हा सूत्र~$\Box \psi$ हे~$w$ येथे
सत्य असते, हे लक्षात घ्या. अशा वेळी $\Box \psi$ हे~$w$ येथे
\emph{रिक्तपणे} सत्य असते. शिवाय $\psi$ सत्य नसतानाही
$\Box \psi$ ची~$w$ येथे पूर्ति होऊ शकते. $\psi$ हे~$w$ येथे
सत्य असल्यामुळेच $\Diamond \psi$ हेही~$w$ येथे सत्य असेल
असे नाही. मात्र $Rww$ असेल, उदाहरणार्थ $R$ परावर्ती असेल,
तर ते खरे ठरते. $Rww'$ असलेले $w'$ मुळीच नसेल, तर कोणत्याही
$\varphi$ साठी $\mSat/{M}{\Diamond \varphi}[w]$ असते.

\begin{prob}
  \ref{nml:syn:rel:fig:simple} मधील प्रतिमान पाहा. पुढीलपैकी
  कोणती विधाने खरी आहेत?
    \begin{enumerate}
    \item $\mSat{M}{q}[w_1]$;
    \item $\mSat{M}{\lnot q}[w_3]$;
    \item $\mSat{M}{p \lor q}[w_1]$;
    \item $\mSat{M}{\Box (p \lor q)}[w_1]$;
    \item $\mSat{M}{\Box q}[w_3]$;
    \item $\mSat{M}{\Box \bot}[w_3]$;
    \item $\mSat{M}{\Diamond q}[w_1]$;
    \item $\mSat{M}{\Box q}[w_1]$;
    \item $\mSat{M}{\lnot \Box\Box \lnot q}[w_1]$.
    \end{enumerate}
\end{prob}



\begin{prop}\label{nml:syn:trw:prop:dual}
  \begin{enumerate}
  \item $\mSat{M}{\Box \varphi}[w]$ तेव्हा आणि केवळ तेव्हाच, जेव्हा $\mSat{M}{\lnot\Diamond\lnot \varphi}[w]$.
  \item $\mSat{M}{\Diamond \varphi}[w]$ तेव्हा आणि केवळ तेव्हाच, जेव्हा $\mSat{M}{\lnot\Box\lnot \varphi}[w]$.
  \end{enumerate}
\end{prop}

\begin{proof}
  \begin{enumerate}
    \item $\mSat{M}{\lnot\Diamond\lnot \varphi}[w]$ तेव्हा आणि केवळ
      तेव्हाच, जेव्हा $\mSat/{M}{\Diamond\lnot \varphi}[w]$;
      हे $\mSat{M}{}[w]$ च्या व्याख्येवरून मिळते.
      $\mSat{M}{\Diamond\lnot \varphi}[w]$ तेव्हा आणि केवळ तेव्हाच,
      जेव्हा $Rww'$ असलेल्या एखाद्या $w'$ साठी
      $\mSat{M}{\lnot \varphi}[w']$. म्हणून
      $\mSat/{M}{\Diamond\lnot \varphi}[w]$ तेव्हा आणि केवळ तेव्हाच,
      जेव्हा $Rww'$ असलेल्या प्रत्येक $w'$ साठी
      $\mSat/{M}{\lnot \varphi}[w']$. तसेच
      $\mSat/{M}{\lnot \varphi}[w']$ तेव्हा आणि केवळ तेव्हाच,
      जेव्हा $\mSat{M}{\varphi}[w']$. मिळून,
      $\mSat{M}{\lnot\Diamond\lnot \varphi}[w]$ तेव्हा आणि केवळ
      तेव्हाच, जेव्हा $Rww'$ असलेल्या प्रत्येक $w'$ साठी
      $\mSat{M}{\varphi}[w']$. $\mSat{M}{}[w]$ च्या व्याख्येवरून
      ही अट तेव्हा आणि केवळ तेव्हाच असते, जेव्हा
      $\mSat{M}{\Box \varphi}[w]$.
    \item $\mSat{M}{\lnot\Box\lnot \varphi}[w]$
      तेव्हा आणि केवळ तेव्हाच, जेव्हा
      $\mSat/{M}{\Box\lnot \varphi}[w]$.
      $\mSat{M}{\Box\lnot \varphi}[w]$ तेव्हा आणि केवळ तेव्हाच,
      जेव्हा $Rww'$ असलेल्या प्रत्येक $w'$ साठी
      $\mSat{M}{\lnot \varphi}[w']$. म्हणून
      $\mSat/{M}{\Box\lnot \varphi}[w]$ तेव्हा आणि केवळ तेव्हाच,
      जेव्हा $Rww'$ असलेल्या एखाद्या $w'$ साठी
      $\mSat/{M}{\lnot \varphi}[w']$. तसेच
      $\mSat/{M}{\lnot \varphi}[w']$ तेव्हा आणि केवळ तेव्हाच,
      जेव्हा $\mSat{M}{\varphi}[w']$. मिळून,
      $\mSat{M}{\lnot\Box\lnot \varphi}[w]$ तेव्हा आणि केवळ
      तेव्हाच, जेव्हा $Rww'$ असलेल्या एखाद्या~$w'$ साठी
      $\mSat{M}{\varphi}[w']$. पुन्हा $\mSat{M}{}[w]$ च्या
      व्याख्येवरून ही अट तेव्हा आणि केवळ तेव्हाच असते,
      जेव्हा $\mSat{M}{\Diamond \varphi}[w]$.
  \end{enumerate}
\end{proof}




\begin{prob}
  $\mModel{M} = \tuple{W, R, V}$ हे प्रतिमान असू द्या आणि
  $w_1, w_2 \in W$ साठी पुढील अटी पूर्ण होतात असे समजा:
  \begin{enumerate}
  \item प्रत्येक विधानीय चल~$p$ साठी $w_1 \in V(p)$
    तेव्हा आणि केवळ तेव्हाच, जेव्हा $w_2 \in V(p)$; आणि
  \item प्रत्येक $w \in W$ साठी: $Rw_1w$ तेव्हा आणि केवळ
    तेव्हाच, जेव्हा $Rw_2w$.
  \end{enumerate}
  सूत्रे वर विगमन करून दाखवा की सर्व सूत्रे~$\varphi$
  साठी $\mSat{M}{\varphi}[w_1]$ तेव्हा आणि केवळ तेव्हाच,
  जेव्हा $\mSat{M}{\varphi}[w_2]$.
\end{prob}

\begin{prob}
  $\mModel{M} = \tuple{W, R, V}$ असू द्या. दाखवा की
  $\mSat{M}{\lnot\Diamond \varphi}[w]$ तेव्हा आणि केवळ तेव्हाच,
  जेव्हा $\mSat{M}{\Box\lnot\varphi}[w]$.
\end{prob}










\section{प्रतिमानात सत्य असणे}\label{nml:syn:tru:sec}

कधीकधी दिलेल्या प्रतिमानातील प्रत्येक जगात कोणती सूत्रे सत्य
आहेत, हे जाणून घ्यायचे असते. यासाठी एक संकेतन देऊ.

\begin{defn}
  सूत्र~$\varphi$ हे $M = \tuple{W, R,
    V}$ या \emph{प्रतिमानात सत्य आहे}, म्हणजे $\mSat{M}{\varphi}$,
  तेव्हा आणि केवळ तेव्हाच, जेव्हा प्रत्येक $w \in W$ साठी
  $\mSat{M}{\varphi}[w]$.
\end{defn}

\begin{prop}\label{nml:syn:tru:prop:truthfacts}
  \begin{enumerate}
  \item $\mSat{M}{\varphi}$ असेल, तर $\mSat/{M}{\lnot \varphi}$;
    पण याचा \emph{व्यत्यास नाही}.
  \item $\mSat{M}{\varphi \lif \psi}$ असेल, तर $\mSat{M}{\varphi}$
    फक्त $\mSat{M}{\psi}$ असेल तरच; पण याचा \emph{व्यत्यास नाही}.
  \end{enumerate}
\end{prop}

\begin{proof}
  \begin{enumerate}
  \item $\mSat{M}{\varphi}$ असेल, तर $\varphi$ हे $W$ मधील सर्व
    जगांत सत्य आहे. $W \neq \emptyset$ असल्यामुळे
    $\mSat{M}{\lnot\varphi}$ असू शकत नाही; तसे असल्यास
    एखाद्या जगात $\varphi$ सत्यही आणि असत्यही असावे लागेल.

    दुसरीकडे, $\mSat/{M}{\lnot \varphi}$ असेल, तर एखाद्या
    $w \in W$ येथे $\varphi$ सत्य आहे. यावरून \emph{प्रत्येक}
    $w \in W$ साठी $\mSat{M}{\varphi}[w]$ मिळत नाही. उदाहरणार्थ,
    \ref{nml:syn:rel:fig:simple} मधील प्रतिमानात $\mSat/{M}{\lnot p}$
    आणि $\mSat/{M}{p}$ ही दोन्ही आहेत.
  \item $\mSat{M}{\varphi \lif \psi}$ आणि $\mSat{M}{\varphi}$ असे समजा.
    $\mSat{M}{\psi}$ दाखवण्यासाठी कोणतेही $w \in W$ घ्या.
    मग $\mSat{M}{\varphi \lif \psi}[w]$ आणि $\mSat{M}{\varphi}[w]$;
    म्हणून $\mSat{M}{\psi}[w]$. $w$ कोणतेही असल्यामुळे
    $\mSat{M}{\psi}$.

    व्यत्यास असत्य ठरवण्यासाठी असे प्रतिमान $\mModel{M}$
    शोधावे लागेल की $\mSat{M}{\varphi}$ हे फक्त $\mSat{M}{\psi}$
    असेल तरच, पण $\mSat/{M}{\varphi \lif \psi}$. पुन्हा
    \ref{nml:syn:rel:fig:simple} मधील प्रतिमान पाहा: $\mSat/{M}{p}$,
    आणि म्हणून $\mSat{M}{p}$ हे फक्त $\mSat{M}{q}$ असेल
    तरच, हे रिक्तपणे खरे आहे. मात्र $w_1$ येथे $p$ सत्य
    आणि $q$ असत्य असल्यामुळे $\mSat/{M}{p \lif
      q}$.
  \end{enumerate}
\end{proof}

\begin{prob}
  $p_1$, $p_2$, $p_3$ हीच विधानीय चले असलेल्या
  भाषेसाठी पुढील प्रतिमान $\mModel{M}$ पाहा:
  \begin{center}
    \begin{tikzpicture}[modal]
      \node[world] (w1) [label={[align=right]left:\mTrue{p_1}\\\mFalse{p_2}\\\mFalse{p_3}}]
            {$w_1$} ;
      \node[world] (w2) [label={[align=right]right:\mTrue{p_1}\\\mTrue{p_2}\\\mFalse{p_3}},
        below right=of w1]  {$w_2$};
      \node[world] (w3) [label={[align=right]right:\mTrue{p_1}\\\mTrue{p_2}\\\mTrue{p_3}},
        above right=of w2] {$w_3$};
      \draw[reflexive above] (w3) to (w3);
      \draw[->] (w1) to (w2);
      \draw[->] (w2) to (w3);
      \draw[->] (w1) to (w3);
    \end{tikzpicture}
  \end{center}
  पुढील सूत्रे आणि रूपबंध $\mModel{M}$ या प्रतिमानात,
  म्हणजे $\mModel{M}$ मधील प्रत्येक जगात, सत्य आहेत का? स्पष्ट करा.
  \begin{enumerate}
  \item $p\lif \Diamond p$ ($p$ अणुरूप असेल तेव्हा);
  \item $\varphi\lif \Diamond \varphi$ ($\varphi$ कोणतेही असेल तेव्हा);
  \item $\Box p \lif p$ ($p$ अणुरूप असेल तेव्हा);
  \item $\lnot p \lif \Diamond \Box p$ ($p$ अणुरूप असेल तेव्हा);
  \item $\Diamond \Box \varphi$ ($\varphi$ कोणतेही असेल तेव्हा);
  \item $\Box \Diamond p$ ($p$ अणुरूप असेल तेव्हा).
  \end{enumerate}
\end{prob}












\section{वैधता}\label{nml:syn:val:sec}

\begin{explain}
  सर्व प्रतिमानांत, म्हणजे प्रत्येक प्रतिमानातील प्रत्येक जगात, सत्य
  असणारी सूत्रे विशेष महत्त्वाची आहेत. संभाव्य जगांवरील
  एखाद्या प्राप्यता संबंधापासून अर्थलक्षण मिळते, या अर्थाने
  $\Box$ आणि $\Diamond$ यांचे अर्थलक्षण ``सामान्य'' असेल,
  तर ते कोणतेही असो, सत्य असणारी मोडल विधाने ही सूत्रे व्यक्त
  करतात. अशा सूत्रे ना आपण \emph{वैध} म्हणतो. उदाहरणार्थ,
  $\Box(p \land q) \lif \Box p$ हे वैध आहे. मात्र $\Box$
  च्या अनिवार्यतेशी संबंधित अर्थावरून वैध वाटणारी काही
  सूत्रे, उदाहरणार्थ $\Box p \lif p$, वैध नसतात.
  संबंधी प्रतिमानांचे एक महत्त्व असे की, वेगवेगळ्या प्रकारच्या
  प्राप्यता संबंधांनी $\Box$ आणि $\Diamond$ यांचे वेगवेगळे
  अर्थलक्षण व्यक्त करता येते. त्यामुळे वैधता केवळ \emph{सर्व}
  प्रतिमानांच्या संदर्भात नव्हे, तर \emph{विशिष्ट प्रकारच्या}
  सर्व प्रतिमानांच्या संदर्भातही परिभाषित करावी. उदाहरणार्थ,
  ज्या प्रतिमानांत प्रत्येक जग स्वतःलाच प्राप्य आहे, म्हणजे
  $R$ परावर्ती आहे, त्या सर्व प्रतिमानांत $\Box p \lif p$
  सत्य आहे, हे पुढे दिसेल. प्रतिमानांच्या वर्गांच्या संदर्भात
  वैधता परिभाषित केल्यास हे थोडक्यात सांगता येते:
  $\Box p \lif p$ हे परावर्ती प्रतिमानांच्या वर्गात वैध आहे.
\end{explain}

\begin{defn}
  सूत्र~$\varphi$ हे प्रतिमानांच्या $\mClass{C}$ या वर्गात
  \emph{वैध} आहे, जर ते $\mClass{C}$ मधील प्रत्येक प्रतिमानात
  सत्य असेल (म्हणजे $\mClass{C}$ मधील प्रत्येक प्रतिमानातील
  प्रत्येक जगात सत्य असेल). $\varphi$ हे $\mClass{C}$ मध्ये वैध
  असेल, तर $\mClass{C} \Entails \varphi$ असे लिहितो; आणि $\varphi$
  हे \emph{सर्व} प्रतिमानांच्या वर्गात वैध असेल, तर
  $\Entails \varphi$ असे लिहितो.
\end{defn}

\begin{prop}\label{nml:syn:val:prop:subset-class}
  $\varphi$ हे $\mClass{C}$ मध्ये वैध असेल, तर प्रत्येक
  $\mClass{C}' \subseteq \mClass{C}$ या वर्गातही ते वैध असते.
\end{prop}

\begin{prop}\label{nml:syn:val:prop:Nec-rule}
  $\varphi$ वैध असेल, तर $\Box\varphi$ देखील वैध असते.
\end{prop}

\begin{proof}
  $\Entails \varphi$ असे समजा. $\Entails \Box\varphi$ दाखवण्यासाठी
  कोणतेही प्रतिमान $\mModel{M} = \tuple{W, R, V}$ आणि
  $w \in W$ घ्या. $Rww'$ असेल, तर $\varphi$ वैध असल्यामुळे
  $\mSat{M}{\varphi}[w']$; म्हणून $\mSat{M}{\Box\varphi}[w]$ देखील.
  $\mModel{M}$ आणि $w$ कोणतेही असल्यामुळे
  $\Entails \Box\varphi$.
\end{proof}

\begin{prob}
  पुढील सूत्रे वैध आहेत हे दाखवा:
  \begin{enumerate}
  \item $\Box p \lif \Box (q \lif p)$;
  \item $\Box \lnot \lfalse$;
  \item $\Box p \lif (\Box q \lif \Box p)$.
  \end{enumerate}
\end{prob}

\begin{prob}
  $W = \{w\}$ असलेल्या $\mModel{M} = \tuple{W, R, V}$
  प्रतिमानांच्या $\mClass{C}$ या वर्गात $\varphi \lif \Box\varphi$
  वैध आहे, हे दाखवा. तसेच $R = \emptyset$ असलेल्या
  $\mModel{M} = \tuple{W, R, V}$ प्रतिमानांच्या वर्गात
  $\psi \lif \Box \varphi$ आणि $\Diamond \varphi \lif \psi$ वैध
  आहेत, हे दाखवा.
\end{prob}












\section{सर्वतः सत्य सूत्रांचे आदेशन-दृष्टांत}\label{nml:syn:tau:sec}

\begin{explain}
मोडल परिकर्मी नसलेले सूत्र प्रत्येक सत्यतामूल्य-नियुक्तीखाली
सत्य असेल, तर ते सर्वतः सत्य सूत्र असते. असे प्रत्येक सूत्र
प्रत्येक प्रतिमानातील प्रत्येक जगात सत्य असते, हे स्पष्ट आहे.
पण $\Box$ आणि~$\Diamond$ असलेल्या सूत्रे साठी
सर्वतः सत्य असण्याची संकल्पना परिभाषित केलेली नाही.
उदाहरणार्थ, विमध्य सिद्धांताचे दृष्टांत सूत्र असलेले
$\Box p \lor \lnot \Box p$ वैध आहे का? येथे
\emph{सर्वतः सत्य सूत्राचा आदेशन-दृष्टांत} ही
संकल्पना उपयोगी पडते: ते मोडल नसलेल्या सर्वतः सत्य
सूत्रात आदेशन करून मिळणारे सूत्र असते. असे
प्रत्येक आदेशन-दृष्टांत सूत्र वैध असणे अपेक्षित आहे;
तरीही ते सिद्ध करावे लागते.
\end{explain}

\begin{defn}
  मोडल सूत्र~$\psi$ हे \emph{सर्वतः सत्य सूत्राचा
  आदेशन-दृष्टांत} तेव्हा आणि केवळ तेव्हाच असते,
  जेव्हा $p_1$, \dots,~$p_n$ ही विधानीय चले असलेले मोडल नसलेले सर्वतः सत्य सूत्र~$\varphi$
  आणि सूत्रे $\theta_1$, \dots,~$\theta_n$ अशी अस्तित्वात
  असतात की $\psi \ident \SSubst{\varphi}{\subst{\theta_1}{p_1},
    \dots, \subst{\theta_n}{p_n}}$.
\end{defn}

\begin{lem}\label{nml:syn:tau:lem:valid-taut}
  $\varphi$ हे मोडल नसलेले सूत्र असून त्यातील
  विधानीय चले $p_1$, \dots, $p_n$ आहेत,
  आणि $\theta_1$, \dots, $\theta_n$ ही मोडल सूत्रे आहेत असे
  समजा. मग कोणतीही नियुक्ती $\pAssign{v}$, कोणतेही प्रतिमान
  $\mModel{M} = \tuple{W, R, V}$, आणि असे कोणतेही $w \in W$
  घ्या की $\pAssign{v}(p_i) = \True$ तेव्हा आणि केवळ तेव्हाच,
  जेव्हा $\mSat{M}{\theta_i}[w]$. अशा वेळी $\pSat{v}{\varphi}$
  तेव्हा आणि केवळ तेव्हाच, जेव्हा
  $\mSat{M}{\SSubst{\varphi}{\subst{\theta_1}{p_1}, \dots,
      \subst{\theta_n}{p_n}}}[w]$.
\end{lem}

\begin{proof}
  $\varphi$ वर विगमन करू.
  \begin{enumerate}
  \item \indcase{\varphi}{\lfalse}{$\pSat/{v}{\lfalse}$
      आणि $\mSat/{M}{\lfalse}[w]$ ही दोन्ही असतात}.
  \item \indcase{\varphi}{\ltrue}{$\pSat{v}{\ltrue}$
      आणि $\mSat{M}{\ltrue}[w]$ ही दोन्ही असतात}.
  \item \indcase{\varphi}{p_i}{
    \begin{align*}
      \pSat{v}{p_i} \Leftrightarrow {} & \pAssign{v}(p_i) = \True \\
      & \qquad \text{$\pSat{v}{p_i}$ च्या व्याख्येवरून}\\
      \Leftrightarrow {} & \mSat{M}{\theta_i}[w] \\
      &\qquad \text{गृहीतकावरून}\\
      \Leftrightarrow {} & \mSat{M}{\SSubst{p_i}{\subst{\theta_1}{p_1}, \dots,
          \subst{\theta_n}{p_n}}}[w]\\
      &\qquad \text{$\SSubst{p_i}{\subst{\theta_1}{p_1}, \dots,
          \subst{\theta_n}{p_n}} \ident \theta_i$ असल्यामुळे}.
      \end{align*}}
  \item \indcase{\varphi}{\lnot \psi}{
      \begin{align*}
        \pSat{v}{\lnot \psi} \Leftrightarrow {} & \pSat/{v}{\psi}\\
        &\qquad \text{$\pSat{v}{}$ च्या व्याख्येवरून};\\
        \Leftrightarrow {} & \mSat/{M}{\SSubst{\psi}{\subst{\theta_1}{p_1}, \dots,
          \subst{\theta_n}{p_n}}}[w]\\
        &\qquad \text{विगमन गृहीतकावरून}\\
        \Leftrightarrow {} &
        \mSat{M}{\SSubst{\lnot \psi}{\subst{\theta_1}{p_1}, \dots,
          \subst{\theta_n}{p_n}}}[w]\\
        &\qquad \text{मोडल पूर्तिच्या $\mSat{M}{}[w]$ व्याख्येवरून}.
      \end{align*}}
  \item \indcase{\varphi}{(\psi \land \chi)}{
      \begin{align*}
        \pSat{v}{\psi \land \chi} \Leftrightarrow {} &
        \pSat{v}{\psi} \text{ आणि } \pSat{v}{\chi}\\
        &\qquad \text{$\pSat{v}{}$ च्या व्याख्येवरून}\\
        \Leftrightarrow {} &
        \mSat{M}{\SSubst{\psi}{\subst{\theta_1}{p_1}, \dots,
            \subst{\theta_n}{p_n}}}[w] \text{ आणि } \\
        & \mSat{M}{\SSubst{\chi}{\subst{\theta_1}{p_1}, \dots,
          \subst{\theta_n}{p_n}}}[w]\\
        &\qquad \text{विगमन गृहीतकावरून}\\
        \Leftrightarrow{} &
        \mSat{M}{\SSubst{(\psi \land \chi)}{\subst{\theta_1}{p_1}, \dots,
          \subst{\theta_n}{p_n}}}[w]\\
        &\qquad \text{$\mSat{M}{}[w]$ च्या व्याख्येवरून}.
      \end{align*}}
  \item \indcase{\varphi}{(\psi \lor \chi)}{
      \begin{align*}
        \pSat{v}{\psi \lor \chi} \Leftrightarrow{} &
        \pSat{v}{\psi} \text{ किंवा } \pSat{v}{\chi}\\
        &\qquad \text{$\pSat{v}{}$ च्या व्याख्येवरून};\\
        \Leftrightarrow{} & \mSat{M}{\SSubst{\psi}{\subst{\theta_1}{p_1}, \dots,
            \subst{\theta_n}{p_n}}}[w] \text{ किंवा }\\
        & \mSat{M}{\SSubst{\chi}{\subst{\theta_1}{p_1}, \dots,
          \subst{\theta_n}{p_n}}}[w]\\
        &\qquad \text{विगमन गृहीतकावरून}\\
        \Leftrightarrow{} &
        \mSat{M}{\SSubst{(\psi \lor \chi)}{\subst{\theta_1}{p_1}, \dots,
          \subst{\theta_n}{p_n}}}[w]\\
        &\qquad \text{$\mSat{M}{}[w]$ च्या व्याख्येवरून}.
      \end{align*}}
  \item \indcase{\varphi}{(\psi \lif \chi)}{
      \begin{align*}
        \pSat{v}{\psi \lif \chi} \Leftrightarrow{} &
        \pSat/{v}{\psi} \text{ किंवा } \pSat{v}{\chi}\\
        &\qquad \text{$\pSat{v}{}$ च्या व्याख्येवरून}\\
        \Leftrightarrow{} &
        \mSat/{M}{\SSubst{\psi}{\subst{\theta_1}{p_1}, \dots,
            \subst{\theta_n}{p_n}}}[w] \text{ किंवा }\\
        & \mSat{M}{\SSubst{\chi}{\subst{\theta_1}{p_1}, \dots,
          \subst{\theta_n}{p_n}}}[w]\\
        &\qquad \text{विगमन गृहीतकावरून}\\
        \Leftrightarrow{} &
        \mSat{M}{\SSubst{(\psi \lif \chi)}{\subst{\theta_1}{p_1}, \dots,
          \subst{\theta_n}{p_n}}}[w]\\
        &\qquad \text{$\mSat{M}{}[w]$ च्या व्याख्येवरून}.
      \end{align*}}
  \item \indcase{\varphi}{(\psi \liff \chi)}{
      \begin{align*}
        \pSat{v}{\psi \liff \chi} \Leftrightarrow {} &
        \text{एकतर } \pSat{v}{\psi} \text{ आणि } \pSat{v}{\chi}\\
        & \text{किंवा } \pSat/{v}{\psi} \text{ आणि } \pSat/{v}{\chi}\\
        &\qquad \text{$\pSat{v}{}$ च्या व्याख्येवरून}\\
        \Leftrightarrow {} &
        \text{एकतर } \mSat{M}{\SSubst{\psi}{\subst{\theta_1}{p_1}, \dots,
            \subst{\theta_n}{p_n}}}[w] \text{ आणि }\\
        & \qquad \mSat{M}{\SSubst{\chi}{\subst{\theta_1}{p_1}, \dots,
          \subst{\theta_n}{p_n}}}[w]\\
        &  \text{ किंवा } \mSat/{M}{\SSubst{\psi}{\subst{\theta_1}{p_1}, \dots,
            \subst{\theta_n}{p_n}}}[w] \text{ आणि }\\
        & \qquad \mSat/{M}{\SSubst{\chi}{\subst{\theta_1}{p_1}, \dots,
          \subst{\theta_n}{p_n}}}[w]\\
        &\qquad\text{विगमन गृहीतकावरून}\\
        \Leftrightarrow {} &
        \mSat{M}{\SSubst{(\psi \liff \chi)}{\subst{\theta_1}{p_1}, \dots,
          \subst{\theta_n}{p_n}}}[w]\\
        &\qquad \text{$\mSat{M}{}[w]$ च्या व्याख्येवरून}.
      \end{align*}}
  \end{enumerate}
  \end{proof}

\begin{prop}\label{nml:syn:tau:prop:valid-taut}
  सर्वतः सत्य सूत्रांचे सर्व आदेशन-दृष्टांत वैध असतात.
\end{prop}

\begin{proof}
  प्रतिपरिवर्तनाने सिद्ध करू. काही प्रतिमान $\mModel{M}$ आणि जग~$w$
  यांसाठी पुढील अट असलेले $\varphi$ आहे असे समजा:
  $\mSat/{M}{\SSubst{\varphi}{\subst{\theta_1}{p_1}, \dots,
  \subst{\theta_n}{p_n}}}[w]$.
  $\pAssign{v}(p_i) = \True$ तेव्हा आणि केवळ तेव्हाच,
  जेव्हा $\mSat{M}{\theta_i}[w]$, अशी नियुक्ती $\pAssign{v}$
  परिभाषित करा ($q \notin \{p_1, \dots, p_n \}$ या
  चलांना $\pAssign{v}$ कोणतीही मूल्ये देते). मग
  \ref{nml:syn:tau:lem:valid-taut} वरून $\pSat/{v}{\varphi}$; म्हणून
  $\varphi$ सर्वतः सत्य सूत्र नाही.
\end{proof}













\section{रूपबंध आणि वैधता}\label{nml:syn:sch:sec}

\begin{defn}
  \emph{रूपबंध} म्हणजे एखाद्या मोडल सूत्र~$\chi$ चे
  नेमके सर्व आदेशन-दृष्टांत मिळून बनलेला सूत्रांचा संच; म्हणजे
  \[  \Setabs{\psi}{\lexists[\theta_1], \dots, \lexists[\theta_n] \left(\psi =
  \SSubst{\chi}{\subst{\theta_1}{p_1}, \dots, \subst{\theta_n}{p_n}} \right) }.
  \]
  सूत्र~$\chi$ याला रूपबंधाचे \emph{वैशिष्ट्यदर्शक} सूत्र
  म्हणतात. विधानीय चलांची नावे बदलण्याचा फरक सोडल्यास ते एकमेव असते.
  सूत्र~$\varphi$ हे त्या संचाचा घटक असेल, तर तो
  रूपबंधाचा \emph{दृष्टांत} असतो.
\end{defn}

रूपबंध दर्शवण्यासाठी $\chi$ मधील विधानीय चलांच्या जागी
`$\varphi$', `$\psi$',~\dots ही चिन्हे बसवून मिळणारी
अधिभाषेतील अभिव्यक्ती वापरणे सोयीचे असते. तार्किक स्थिरांक बदलत नाहीत. उदाहरणार्थ,
`$\varphi$', `$\varphi \lif \Box\varphi$', `$\varphi \lif (\psi \lif \varphi)$'
ही रूपबंध दर्शवतात. त्यांची वैशिष्ट्यदर्शक सूत्रे
अनुक्रमे $p$, $p \lif \Box p$, $p \lif (q
\lif p)$ आहेत. `$\varphi$' हा रूपबंध \emph{सर्व}
सूत्रांचा संच दर्शवतो.

\begin{defn}
  रूपबंधाचे सर्व दृष्टांत एखाद्या प्रतिमानात सत्य असतील तेव्हा
  आणि केवळ तेव्हाच तो रूपबंध त्या प्रतिमानात \emph{सत्य}
  असतो; आणि प्रत्येक प्रतिमानात सत्य असेल तेव्हा आणि केवळ
  तेव्हाच तो \emph{वैध} असतो.
\end{defn}

\begin{prop}\label{nml:syn:sch:prop:Kvalid}
  पुढील \Ax{K} हा रूपबंध वैध आहे:
  \begin{equation*}
  \Box(\varphi \lif \psi) \lif (\Box \varphi \lif \Box \psi). \tag{\Ax{K}}
  \end{equation*}
\end{prop}

\begin{proof}
  या रूपबंधाचा प्रत्येक दृष्टांत प्रत्येक प्रतिमानातील प्रत्येक
  जगात सत्य आहे, हे दाखवायचे आहे. म्हणून कोणतेही
  $\mModel{M} = \tuple{W,R,V}$ आणि $w \in W$ घ्या.
  एखादे अभिव्यंजन एखाद्या जगात सत्य असल्याचे दाखवताना,
  त्याचे पूर्वांग सत्य मानून उत्तरांगही सत्य आहे हे दाखवतो.
  येथे $\mSat{M}{\Box(\varphi \lif \psi)}[w]$ आणि
  $\mSat{M}{\Box \varphi}[w]$ असू द्या. आता
  $\mSat{M}{\Box \psi}[w]$ दाखवायचे आहे. म्हणून $Rww'$
  असलेले कोणतेही $w'$ घ्या. पहिल्या गृहीतकावरून
  $\mSat{M}{\varphi \lif \psi}[w']$ आणि दुसऱ्यावरून
  $\mSat{M}{\varphi}[w']$. त्यामुळे $\mSat{M}{\psi}[w']$.
  $w'$ कोणतेही असल्यामुळे $\mSat{M}{\Box \psi}[w]$.
\end{proof}

\begin{prop}\label{nml:syn:sch:prop:Dual-valid}
  पुढील \Dual{} हा रूपबंध वैध आहे:
  \begin{equation*}
  \Diamond \varphi \liff \lnot\Box\lnot \varphi. \tag{\Dual}
  \end{equation*}
\end{prop}

\begin{proof}
  सराव.
\end{proof}

\begin{prob}
  \ref{nml:syn:sch:prop:Dual-valid} सिद्ध करा.
\end{prob}

\begin{prop}\label{nml:syn:sch:prop:soundMP}
  प्रतिमानातील एखाद्या जगात $\varphi$ आणि $\varphi \lif \psi$ सत्य
  असतील, तर $\psi$ देखील सत्य असते. म्हणून वैध सूत्रे
  विध्यात्मक अनुमानाखाली बंद असतात.
\end{prop}

\begin{prop}\label{nml:syn:sch:prop:valid-instances}
  सूत्र~$\varphi$ तेव्हा आणि केवळ तेव्हाच वैध असते,
  जेव्हा त्याचे सर्व आदेशन-दृष्टांत वैध असतात. दुसऱ्या शब्दांत,
  रूपबंध तेव्हा आणि केवळ तेव्हाच वैध असतो, जेव्हा त्याचे
  वैशिष्ट्यदर्शक सूत्र वैध असते.
\end{prop}

\begin{proof}
  ``जर'' ही दिशा स्पष्ट आहे, कारण $\varphi$ स्वतःचाच
  आदेशन-दृष्टांत आहे.

  ``फक्त तरच'' ही दिशा सिद्ध करण्यासाठी पुढील दावा दाखवू.
  $\mModel{M} = \tuple{W, R, V}$ हे मोडल प्रतिमान असू द्या,
  आणि $\psi \ident
  \SSubst{\varphi}{\subst{\theta_1}{p_1},\dots,\subst{\theta_n}{p_n}}$ हा
  $\varphi$ चा आदेशन-दृष्टांत असू द्या. $1 \leq i \leq n$ साठी
  $V'(p_i) = \Setabs{w \in W}{\mSat{M}{\theta_i}[w]}$ ठेवा;
  इतर प्रत्येक विधानीय चल~$q \notin \{p_1,\dots,p_n\}$
  साठी $V'(q)=V(q)$ ठेवा. या सर्व चलांवरील नियुक्तीने
  $\mModel{M'} = \tuple{W, R,V'}$ परिभाषित करा.
  मग कोणत्याही~$w \in W$ साठी
  $\mSat{M}{\psi}[w]$ तेव्हा आणि केवळ तेव्हाच, जेव्हा
  $\mSat{M'}{\varphi}[w]$. (याची सिद्धता सरावासाठी ठेवली आहे.)
  आता $\varphi$ वैध आहे, पण $\varphi$ चा एखादा आदेशन-दृष्टांत
  $\psi$ वैध नाही, असे समजा. मग काही
  $\mModel{M} = \tuple{W, R, V}$ आणि काही $w \in W$
  साठी $\mSat/{M}{\psi}[w]$. पण दाव्यावरून
  $\mSat/{M'}{\varphi}[w]$; म्हणजे $\varphi$ वैध नाही. हा विरोधाभास.
\end{proof}

\begin{prob}
  \ref{nml:syn:sch:prop:valid-instances} च्या सिद्धतेतील
  ``फक्त तरच'' भागातला दावा सिद्ध करा. (सूचना:
  $\varphi$ वर विगमन करा.)
\end{prob}

मात्र एखादा रूपबंध प्रतिमानात सत्य असणे आणि त्याचे
वैशिष्ट्यदर्शक सूत्र त्या प्रतिमानात सत्य असणे या सममूल्य
अटी नाहीत, हे लक्षात घ्या. ``फक्त तरच'' ही दिशा खरी आहे:
$\varphi$ चा प्रत्येक दृष्टांत~$\mModel{M}$ मध्ये सत्य असेल,
तर $\varphi$~स्वतःही~$\mModel{M}$ मध्ये सत्य असते. पण
$\varphi$ हे~$\mModel{M}$ मध्ये सत्य असूनही $\varphi$ चा एखादा
दृष्टांत~$\mModel{M}$ मधील काही जगात असत्य असू शकतो.
अगदी साध्या प्रतिदृष्टांतासाठी~$p$ घेऊन एकच जग~$w$ असलेले आणि
$V(p) = \{w\}$ असलेले प्रतिमान घ्या. तेथे~$p$ हे~$w$
येथे सत्य आहे. पण $\lfalse$ हा~$p$ चा दृष्टांत असून
तो~$w$ येथे सत्य नाही.

\begin{prob}
  पुढीलपैकी कोणतेही सूत्रे वैध नाही, हे दाखवा:
  \begin{enumerate}
    \item[\Ax{D}:] \quad $\Box p \lif \Diamond p$;
    \item[\Ax{T}:] \quad $\Box p \lif p$;
    \item[\Ax{B}:] \quad $p \lif \Box\Diamond p$;
    \item[\Ax{4}:] \quad $\Box p \lif \Box \Box p$;
    \item[\Ax{5}:] \quad $\Diamond p \lif \Box \Diamond
      p$.
  \end{enumerate}
\end{prob}

\begin{table}[t]
    \centering
    \begin{tabular}{| l || l |}
      \hline
      {\emph{वैध रूपबंध}} & {\emph{अवैध रूपबंध}} \\
      \hline\hline
      $\Box(\varphi \lif \psi) \lif (\Diamond \varphi \lif \Diamond \psi)$
      & $\Box (\varphi \lor \psi) \lif (\Box \varphi \lor \Box \psi)$ \\
      $\Diamond (\varphi \lif \psi) \lif (\Box \varphi \lif \Diamond
      \psi)$
      & $(\Diamond \varphi \land \Diamond \psi) \lif \Diamond (\varphi
      \land \psi)$\\
      $\Box (\varphi \land \psi) \liff (\Box \varphi \land \Box \psi)$
      & $\varphi \lif \Box \varphi$ \\
      $\Box \varphi \lif \Box (\psi \lif \varphi)$
      & $\Box \Diamond \varphi \lif \psi$ \\
      $\lnot \Diamond \varphi \lif \Box (\varphi \lif \psi)$
      & $\Box \Box \varphi \lif \Box \varphi$ \\
      $\Diamond (\varphi \lor \psi) \liff (\Diamond \varphi \lor
      \Diamond \psi)$
      & $\Box \Diamond \varphi \lif \Diamond \Box \varphi$. \\
      \hline
    \end{tabular}
    \caption{वैध आणि अवैध रूपबंध.}
    \label{nml:syn:sch:tab:valid-invalidSchemas}
\end{table}

\begin{prob}
  \ref{nml:syn:sch:tab:valid-invalidSchemas} च्या
  पहिल्या स्तंभातील रूपबंध वैध आहेत, आणि दुसऱ्या
  स्तंभातील रूपबंध वैध नाहीत, हे सिद्ध करा.
\end{prob}

\begin{prob}
  पुढील रूपबंध वैध आहेत की अवैध, ते ठरवा:
  \begin{enumerate}
  \item $(\Diamond \varphi \lif \Box \psi) \lif (\Box \varphi \lif \Box
    \psi)$;
  \item $\Diamond(\varphi \lif \psi) \lor \Box(\psi \lif \varphi)$.
  \end{enumerate}
\end{prob}

\begin{prob}
  पुढील प्रत्येक रूपबंधासाठी असे प्रतिमान $\mModel{M}$ शोधा की
  त्या सूत्राचा प्रत्येक दृष्टांत $\mModel{M}$ मध्ये सत्य आहे:
  \begin{enumerate}
  \item $p \lif \Diamond\Diamond p$;
  \item $\Diamond p \lif \Box p$.
  \end{enumerate}
\end{prob}












\section{तार्किक निष्पन्नता}\label{nml:syn:ent:sec}

\begin{explain}
  एखाद्या जगात सत्य असण्याच्या व्याख्येने सूत्रे मधील
  तार्किक निष्पन्नतेचा संबंध परिभाषित करता येतो.
  सूत्र~$\psi$ पासून~$\varphi$ तार्किकरीत्या निष्पन्न
  होते तेव्हा आणि केवळ तेव्हाच, जेव्हा $\psi$ सत्य असेल
  तेव्हा $\varphi$ देखील सत्य असते. येथे ``तेव्हा'' म्हणताना
  कोणतेही प्रतिमान आणि त्यातील कोणतेही जग अभिप्रेत आहे.
\end{explain}

\begin{defn}
  $\Gamma$ हा सूत्रांचा संच आणि $\varphi$ हे सूत्र
  असेल, तर $\Gamma$ पासून~$\varphi$ \emph{तार्किकरीत्या
  निष्पन्न होते}, म्हणजे $\Gamma \Entails \varphi$, तेव्हा आणि
  केवळ तेव्हाच, जेव्हा प्रत्येक प्रतिमान
  $\mModel{M} = \tuple{W, R, V}$ आणि प्रत्येक $w \in W$
  साठी, सर्व $\psi \in \Gamma$ बाबत $\mSat{M}{\psi}[w]$
  असेल तर $\mSat{M}{\varphi}[w]$ असते. $\Gamma$ मध्ये
  एकच सूत्र~$\psi$ असेल, तर $\psi \Entails \varphi$ लिहितो.
\end{defn}

\begin{ex}
  एका सूत्र पासून दुसरे तार्किकरीत्या निष्पन्न होते हे
  दाखवण्यासाठी $\mSat{M}{}[w]$ ची व्याख्या वापरून सर्व
  प्रतिमानांबद्दल विचार करावा लागतो. उदाहरणार्थ,
  $p \lif \Diamond p \Entails \Box\lnot p \lif \lnot p$
  दाखवण्यासाठी असा युक्तिवाद करता येईल: प्रतिमान
  $\mModel{M} = \tuple{W, R,
    V}$ आणि $w \in W$ घ्या, आणि $\mSat{M}{p \lif \Diamond
    p}[w]$ असे समजा. $\mSat{M}{\Box\lnot p \lif \lnot
    p}[w]$ दाखवायचे आहे. तसे नाही असे समजा. मग
  $\mSat{M}{\Box\lnot p}[w]$ आणि $\mSat/{M}{\lnot p}[w]$.
  $\mSat/{M}{\lnot p}[w]$ असल्यामुळे $\mSat{M}{
    p}[w]$. गृहीतकावरून $\mSat{M}{p \lif \Diamond p}[w]$;
  म्हणून $\mSat{M}{\Diamond p}[w]$.
  $\mSat{M}{\Diamond
    p}[w]$ च्या व्याख्येवरून $Rww'$ असलेले असे $w'$
  आहे की $\mSat{M}{p}[w']$. पण $\mSat{M}{\Box \lnot p}[w]$
  देखील असल्यामुळे $\mSat{M}{\lnot p}[w']$; हा विरोधाभास.

  सूत्र~$\psi$ पासून दुसरे~$\varphi$ तार्किकरीत्या
  निष्पन्न होत नाही हे दाखवण्यासाठी प्रतिदृष्टांत द्यावा
  लागतो: असे प्रतिमान $\mModel{M} = \tuple{W, R,
    V}$ की त्यातील एखाद्या जगात~$w \in W$ येथे
  $\mSat{M}{\psi}[w]$, पण $\mSat/{M}{\varphi}[w]$.
  $p \lif \Diamond p \Entails/
  \Box p \lif p$ हे दाखवू. \ref{nml:syn:ent:fig:counterex} मधील
  प्रतिमान पाहा. येथे $\mSat{M}{\Diamond p}[w_1]$
  आणि म्हणून $\mSat{M}{p \lif
    \Diamond p}[w_1]$. मात्र $\mSat{M}{\Box p}[w_1]$
  असूनही $\mSat/{M}{p}[w_1]$ असल्यामुळे
  $\mSat/{M}{\Box p \lif p}[w_1]$.
  \begin{figure}
  \begin{center}
    \begin{tikzpicture}[modal]
      \node[world] (w1) [label=right:\mFalse{p}]{$w_1$};
      \node[world] (w2) [label=right:\mTrue{p}, above left=of w1]{$w_2$};
      \node[world] (w3) [label=right:\mTrue{p}, above right=of w1] {$w_3$};
      \draw[->] (w1) to (w2);
      \draw[->] (w1) to (w3);
    \end{tikzpicture}
  \end{center}
  \caption{$p \lif \Diamond p
  \Entails \Box p \lif p$ याचा प्रतिदृष्टांत.}
  \label{nml:syn:ent:fig:counterex}
  \end{figure}

  अनेकदा अगदी साधे प्रतिदृष्टांत पुरेसे असतात.
  $W' = \{w\}$, $R' = \emptyset$, आणि $V'(p) =
  \emptyset$ असलेले प्रतिमान $\mModel{M'} =
  \tuple{W', R', V'}$ हाही प्रतिदृष्टांत आहे:
  $\mSat/{M'}{p}[w]$ असल्यामुळे
  $\mSat{M'}{p \lif \Diamond p}[w]$. $w$ पासून एकही
  जग प्राप्य नसल्यामुळे $\mSat{M'}{\Box p}[w]$; म्हणून
  $\mSat/{M'}{\Box p
    \lif p}[w]$.
\end{ex}

\begin{prob}
  $\Box (\varphi \land \psi) \Entails \Box \varphi$ हे दाखवा.
\end{prob}

\begin{prob}
  $\Box (p \lif q) \Entails/ p \lif \Box q$ आणि $p \lif \Box
  q \Entails/\Box (p \lif q)$ हे दाखवा.
\end{prob}



\chapter{मोडल चौकटींची परिभाष्यता}\label{nml:frd:chap}









\section{प्रस्तावना}\label{nml:frd:int:sec}

मोडल तर्कशास्त्रात प्राप्यता संबंध आणि तो संबंध असलेल्या
प्रतिमानांतील काही सूत्रांची सत्यता यांमधील संबंध
महत्त्वाचा असतो. उदाहरणार्थ, प्राप्यता संबंध परावर्ती
आहे, म्हणजे प्रत्येक $w
\in W$ साठी $Rww$ आहे, असे समजा. दुसऱ्या शब्दांत,
प्रत्येक जग स्वतःपासून प्राप्य आहे. मग एखाद्या जगात~$w$
येथे $\Box \varphi$ सत्य असेल, तर ज्या प्राप्य जगांत~$\varphi$
सत्य असणे आवश्यक आहे त्यांत $w$ स्वतःही येते. म्हणून
$\mModel{M}$ चा प्राप्यता संबंध~$R$ परावर्ती असेल, तर
कोणतेही जग $w$ आणि कोणतेही सूत्र $\varphi$ घेतले, तरी
$\Box \varphi
\lif \varphi$ तेथे सत्य असेल. दुसऱ्या शब्दांत, $\Box p \lif
p$ हा रूपबंध आणि त्याचे सर्व आदेशन-दृष्टांत~$\mModel{M}$
मध्ये सत्य आहेत.

मात्र व्यत्यास असत्य आहे. उदाहरणार्थ, $\Box p \lif p$
हे~$\mModel{M}$ मध्ये सत्य असल्यावरून $R$ परावर्ती आहे
असे ठरत नाही. कारण $\Box p \lif
p$ सर्व जगांत सत्य असलेले अपरावर्ती प्रतिमान~$\mModel{M}$
सहज मिळते: स्वतःपासून प्राप्य नसलेले एकच जग~$w$ असलेले,
पण $w \in V(p)$ असलेले प्रतिमान घ्या. $p$ चे सत्यतामूल्य
योग्य निवडल्यास अपरावर्ती प्रतिमानात एखाद्या विशिष्ट
सूत्रासाठी $\Box \varphi \lif \varphi$ सत्य करता येते.

यासाठी मूल्यांकनाचा घटक~$V$ बाजूला ठेवू. $V(p)$ मध्ये
कोणती जगं आहेत याची पर्वा न करता, $\Box p \lif p$ हे
$\mModel{M}$ मधील सर्व जगांत सत्य असावे अशी अट घातल्यास
$R$ परावर्ती असणे आवश्यक ठरते. कारण अपरावर्ती
प्रतिमानात $Rww$ नसलेले किमान एक जग~$w$ असते.
$V(p) =
W \setminus \{w\}$ असे ठेवा. मग~$w$ सोडून इतर
सर्व जगांत $p$ सत्य असते; म्हणून~$w$ पासून प्राप्य
असलेल्या सर्व जगांतही ते सत्य असते ($w$ हे~$w$ पासून
प्राप्य नसल्याची खात्री आहे, आणि $p$ असत्य असलेले
$w$ हेच एकमेव जग आहे). दुसरीकडे, $w$ येथे $p$ असत्य
आहे; म्हणून $\Box
p \lif p$ हे~$w$ येथे असत्य आहे.

यावरून मूल्यांकनाशिवाय प्रतिमानाची रचना दाखवण्यासाठी
स्वतंत्र संकेतन हवे, असे सुचते; अशा रचनांना \emph{चौकटी}
म्हणतो. चौकट $\mModel{F}$ म्हणजे जगांचा संच आणि
प्राप्यता संबंध यांची जोडी $\tuple{W, R}$.
प्रत्येक प्रतिमान $\tuple{W, R, V}$ हे $\tuple{W, R}$
या चौकटीवर \emph{आधारित} असते. उलट, एखादी चौकट
तिच्यावर आधारित प्रतिमानांचा वर्ग ठरवते; आणि चौकटींचा
वर्ग त्यांपैकी कोणत्यातरी चौकटीवर आधारित प्रतिमानांचा
वर्ग ठरवतो. सूत्राचे चौकटीत \emph{वैध} असणे,
म्हणजे $\mModel{F} \Entails \varphi$, अशी व्याख्या करता येते:
चौकट~$\mModel{F}$ वर आधारित प्रत्येक $\mModel{M}$
साठी $\mSat{M}{\varphi}$.

या संकेतनाने सूत्रे आणि चौकटींचे वर्ग यांतील
अनुरूपता सांगता येते: उदाहरणार्थ,
$\mModel{F} \Entails \Box p \lif p$ तेव्हा आणि केवळ
तेव्हाच, जेव्हा $\mModel{F}$ परावर्ती आहे.













\section{प्राप्यता संबंधांचे गुणधर्म}\label{nml:frd:acc:sec}

अनेक मोडल सूत्रे प्राप्यता संबंधाच्या साध्या, कधीकधी
परिचित, गुणधर्मांचे वैशिष्ट्य दर्शवतात. एका दिशेने याचा
अर्थ असा की, एखादा गुणधर्म असलेल्या प्रत्येक प्रतिमानात
संबंधित सूत्र आणि त्याचे सर्व आदेशन-दृष्टांत सत्य
असतात. प्राप्यता संबंधांचे पाच अभिजात प्रकार आणि त्यांनी
सत्य ठरवलेली सूत्रे प्रथम पाहू.

\begin{thm}\label{nml:frd:acc:thm:soundschemas}
  $\mModel{M} = \tuple{W, R, V}$ हे प्रतिमान असू द्या.
  $R$ ला \ref{nml:frd:acc:tab:five} च्या डाव्या बाजूचा गुणधर्म
  असेल, तर उजव्या बाजूच्या सूत्राचा प्रत्येक
  दृष्टांत~$\mModel{M}$ मध्ये सत्य असतो.
\end{thm}

\begin{table}[t]
    \begin{tabular}{| l || l |}
      \hline
      {\emph{$R$ पुढील प्रकारचा असेल \dots}} & {\emph{तर \dots सूत्र~$\mModel{M}$ मध्ये सत्य:}} \\
      \hline \hline
      \emph{सर्वत्र उत्तरवर्ती}: $\forall u \exists v Ruv$ & \hbox to.43\textwidth
           {$\Box p \lif \Diamond p$ \hfill (\Ax{D})} \\
      \hline
      \emph{परावर्ती}: $\forall w Rww$
      & \hbox to.43\textwidth
          {$\Box p \lif p$ \hfill (\Ax{T})} \\
      \hline
      \emph{सममित}: &  \hbox to .43\textwidth
           {$p \lif \Box\Diamond p$ \hfill (\Ax{B})} \\
           $\forall u\forall v(Ruv \lif Rvu)$ & \\
      \hline
      \emph{संक्रमक}: & \hbox to.43\textwidth
           {$\Box p \lif \Box \Box p$ \hfill (\Ax{4})} \\
      $\forall u \forall v \forall w ((Ruv \land Rvw) \lif Ruw)$ & \\
      \hline
      \emph{युक्लिडी}: & \hbox to .43\textwidth
        {$\Diamond p \lif \Box\Diamond p$ \hfill (\Ax{5})} \\
      $\forall w \forall u \forall v ((Rwu \land Rwv) \lif Ruv)$ &\\
      \hline
    \end{tabular}
    \caption{अनुरूपतेविषयी पाच तथ्ये.}
    \label{nml:frd:acc:tab:five}
  \end{table}


\begin{proof}
  \Ax{B} चे प्रकरण पाहू. रूपबंध प्रतिमानात सत्य आहे हे
  दाखवण्यासाठी त्याचे सर्व दृष्टांत त्या प्रतिमानातील
  प्रत्येक जगात सत्य आहेत हे दाखवावे लागते. म्हणून
  $\varphi \lif \Box\Diamond \varphi$ हा \Ax{B} चा कोणताही
  दृष्टांत, आणि $w \in W$ हे कोणतेही जग असू द्या.
  $\Box \Diamond
  \varphi$ हे $w$ येथे सत्य असल्याचे दाखवण्यासाठी पूर्वांग
  $\varphi$ हे $w$ येथे सत्य आहे असे समजा. मग $w$ पासून
  प्राप्य असलेल्या प्रत्येक $w'$ येथे $\Diamond \varphi$
  सत्य आहे हे दाखवायचे आहे. आता $Rww'$ असलेल्या
  कोणत्याही $w'$ साठी, सममितीच्या गृहीतकावरून $Rw'w$
  देखील असते (\ref{nml:frd:acc:fig:Bsymm} पाहा).
  $\mSat{M}{\varphi}[w]$ असल्यामुळे
  $\mSat{M}{\Diamond \varphi}[w']$. $Rww'$ असलेले $w'$
  कोणतेही असल्यामुळे $\mSat{M}{\Box\Diamond\varphi}[w]$.

  इतर प्रकरणे सरावासाठी ठेवली आहेत.
\end{proof}

\begin{prob}
  \ref{nml:frd:acc:thm:soundschemas} ची सिद्धता पूर्ण करा.
\end{prob}

\begin{figure}
  \begin{center}
    \begin{tikzpicture}[modal]
      \node[world] (w1) [label={below:$\mSat{{}}{\formula{A}}$\\
          $\mSat{{}}{\Box\Diamond \formula{A}}$}] {$w$};
      \node[world] (w2) [label=below:{$\mSat{{}}{\Diamond \formula{A}}$},
        right=of w1] {$w'$};
      \draw[->, bend left] (w1) to (w2);
      \draw[->, bend left] (w2) to (w1);
    \end{tikzpicture}
  \end{center}
\caption{सममितीवर आधारित युक्तिवाद.}
\label{nml:frd:acc:fig:Bsymm}
\end{figure}

\ref{nml:frd:acc:thm:soundschemas} मधील व्यत्यास खरे नाहीत, हे लक्षात
घ्या: एखादा रूपबंध प्रतिमानात सत्य आहे म्हणून त्या प्रतिमानाच्या
प्राप्यता संबंधाला संबंधित गुणधर्म असतोच असे नाही.
\Ax{T} आणि परावर्ती प्रतिमाने या प्रकरणात, \Ax{T} स्वतःच
असत्य असलेले प्रतिमान सहज देता येते: $W = \{w\}$ आणि
$V(p) = \emptyset$ घ्या; प्राप्यता संबंध रिकामा असू द्या.
मग $R$ परावर्ती नाही, पण $\mSat{M}{\Box
  p}[w]$ आणि $\mSat/{M}{p}[w]$. तरी येथे~$\mModel{M}$
मध्ये \Ax{T} चा असत्य असलेला एक दृष्टांत आपण दाखवला आहे; काही इतर दृष्टांत,
उदाहरणार्थ $\Box \lnot p \lif \lnot p$, सत्य आहेत.
\Ax{T} चे \emph{प्रत्येक आदेशन-दृष्टांत}~$\mModel{M}$
मध्ये सत्य असूनही $\mModel{M}$ परावर्ती नसते, असे
प्रतिमान शोधणे अधिक कठीण आहे. पण अशी प्रतिमानेही आहेत:

\begin{prop}\label{nml:frd:acc:prop:reflexive}
  $\mModel{M} = \tuple{W, R, V}$ हे प्रतिमान असू द्या,
  जिथे $W = \{u, v
  \}$, जगं $u$ आणि $v$ यांचा $R$ संबंध दोन्ही दिशांनी
  आहे, म्हणजे $Ruv$ आणि $Rvu$, आणि कोणतेही जग स्वतःशी
  संबंधित नाही. प्रत्येक $p$ साठी $u \in V(p) \Leftrightarrow v
  \in V(p)$ असे समजा. मग:
  \begin{enumerate}
  \item प्रत्येक $\varphi$ साठी: $\mSat{M}{\varphi}[u]$ तेव्हा आणि
    केवळ तेव्हाच, जेव्हा $\mSat{M}{\varphi}[v]$
    ($\varphi$ वर विगमन वापरा).
  \item \Ax{T} चा प्रत्येक दृष्टांत $\mModel{M}$ मध्ये सत्य आहे.
  \end{enumerate}
  $\mModel{M}$ परावर्ती नाही (खरेतर ते
  \emph{अपरावर्ती} आहे), म्हणून \Ax{T} च्या बाबतीत
  \ref{nml:frd:acc:thm:soundschemas} चा व्यत्यास असत्य आहे.
  \ref{nml:frd:acc:thm:soundschemas} मधील काही, पण सर्वच नव्हे,
  इतर रूपबंधांसाठीही असेच युक्तिवाद देता येतात.
\end{prop}

\begin{prob}
  \ref{nml:frd:acc:prop:reflexive} मधील दावे सिद्ध करा.
\end{prob}

आपण \Ax{D}, \Ax{T}, \Ax{B}, \Ax{4}, आणि \Ax{5}
या पाच अभिजात सूत्रे वर भर देणार असलो, तरी
\ref{nml:frd:acc:tab:anotherfive} मध्ये प्राप्यता संबंधांचे आणखी
काही गुणधर्म नोंदवले आहेत. प्रत्येक जगापासून जास्तीत
जास्त एक जग प्राप्य असेल, तर प्राप्यता संबंध~$R$
आंशिक फलनीय असतो. प्रत्येक जगापासून नेमके
एकच जग प्राप्य असेल, तर त्याला फलनीय म्हणतो.
(म्हणून फलनीय संबंध नेमके तेच आहेत जे सर्वत्र उत्तरवर्ती
आणि आंशिक फलनीय दोन्ही आहेत.) त्यांना ``फलनीय''
म्हणतात कारण प्राप्यता संबंध एखाद्या (आंशिक) फलनासारखा
वागतो. $Ruv$ असेल तेव्हा $u$ आणि~$v$ यांच्या
``मध्ये'' एखादे~$w$ असेल, तर संबंध दुर्बल सघन
असतो. एका अर्थाने दुर्बल सघन संबंध संक्रमक संबंधांच्या
उलट आहेत: संक्रमक संबंधात~$u$ पासून~$w$ मार्गे
$v$ पर्यंत जाता येत असेल, तर $u$ पासून $v$ पर्यंत
थेट जाता येते; दुर्बल सघन संबंधात~$u$ पासून $v$
पर्यंत थेट जाता येत असेल, तर एखाद्या~$w$ मार्गेही
जाता येते. एखाद्या जग~$w$ पासून जगं $u$ आणि~$v$
प्राप्य असतील तेव्हा $u$ आणि~$v$ या दोन्हींपासून
एकच जग~$t$ प्राप्य असेल, तर संबंध दुर्बल दिशित असतो;
यालाच कधीकधी ``हिऱ्याचा गुणधर्म'' किंवा ``संगमन''
म्हणतात.

\begin{table}[t]
    \setlength{\tabcolsep}{4pt}
\begin{tabular}{| p{.50\textwidth} || p{.44\textwidth} |}
      \hline
      {\emph{$R$ पुढील प्रकारचा असेल \dots}} & {\emph{तर \dots सूत्र~$\mModel{M}$ मध्ये सत्य:}} \\
      \hline \hline
      \emph{आंशिक फलनीय}: &
      \multirow{2}{*}{$\Diamond p \lif \Box p$} \\
      $ \forall w \forall u \forall v ((Rwu \land Rwv) \lif u =v )$ & \\
      \hline
      \multirow{2}{*}{\emph{फलनीय}: $\forall w \exists v \forall u(Rwu \liff \eq[u][v]) $}
      & \multirow{2}{*}{$\Diamond p \liff \Box p$} \\
      & \\
      \hline
      \emph{दुर्बल सघन}: &  \multirow{2}{*}{$\Box \Box p \lif
        \Box p$} \\
      $\forall u\forall v(Ruv \lif \exists w(Ruw \land Rwv))$ & \\
      \hline
      \emph{दुर्बल संयोजित}: &
      \multirow{3}{*}{\hbox to.43\textwidth{$
        \begin{array}{@{}l@{}}
          \Box ((p \land \Box p) \lif q) \lor {} \\
          \qquad \Box ((q \land \Box q) \lif p)
        \end{array}$ \hfill (\Ax{L})}
      } \\
      $\forall w \forall u \forall v ((Rwu \land Rwv) \lif {}$ &\\
      \qquad $(Ruv \lor u=v \lor Rvu))$ & \\
      \hline
      \emph{दुर्बल दिशित}: & \multirow{3}{*}{\hbox to .43\textwidth
        {$\Diamond\Box p \lif \Box\Diamond p$\hfill (\Ax{G})}} \\
      $\forall w \forall u \forall v ((Rwu \land Rwv) \lif {}$ & \\
      \qquad $\exists t (Rut \land Rvt))$ & \\
      \hline
    \end{tabular}
    \caption{अनुरूपतेविषयी आणखी पाच तथ्ये.}
    \label{nml:frd:acc:tab:anotherfive}
  \end{table}

\begin{prob}
  $\mModel{M} = \tuple{W, R, V}$ हे प्रतिमान असू द्या.
  $R$ ला \ref{nml:frd:acc:tab:anotherfive} मधील
  डावीकडील गुणधर्म असेल, तर उजवीकडील संबंधित सूत्राचा
  प्रत्येक दृष्टांत~$\mModel{M}$ मध्ये सत्य असतो, हे दाखवा.
\end{prob}














\section{चौकटी}\label{nml:frd:fra:sec}

\begin{defn}
  \emph{चौकट} म्हणजे $\mModel{F} = \tuple{W,R}$ ही जोडी, जिथे
  $W$ हा जगांचा रिक्त नसलेला संच आणि $R$ हा~$W$ वरील
  द्विस्थानी संबंध आहे. प्रतिमान~$\mModel{M}$ हे
  $\mModel{F} = \tuple{W,R}$ या चौकटीवर \emph{आधारित}
  आहे तेव्हा आणि केवळ तेव्हाच, जेव्हा काही मूल्यांकन
  $V$ साठी $\mModel{M} = \tuple{W, R, V}$.
\end{defn}

\begin{defn}
  $\mModel{F}$ ही चौकट असेल, तर~$\mModel{F}$ वर आधारित
  प्रत्येक प्रतिमान~$\mModel{M}$ साठी $\mSat{M}{\varphi}$
  असल्यास $\varphi$ हे \emph{$\mModel{F}$ मध्ये वैध,}
  म्हणजे $\mModel{F} \Entails \varphi$, असे म्हणतो.

  $\mClass{F}$ हा चौकटींचा वर्ग असेल, तर त्या वर्गातील
  प्रत्येक चौकट~$\mModel{F} \in \mClass{F}$ साठी
  $\mModel{F} \Entails \varphi$ असल्यास, आणि तेव्हाच,
  $\varphi$ हे \emph{$\mClass{F}$ मध्ये वैध,}
  म्हणजे $\mClass{F} \Entails \varphi$, असे म्हणतो.
\end{defn}

चौकटी महत्त्वाच्या आहेत, कारण रूपबंध आणि प्राप्यता
संबंध~$R$ याचे गुणधर्म यांमधील अनुरूपता \emph{प्रतिमानांच्या
नव्हे, तर चौकटींच्या} पातळीवर असते. उदाहरणार्थ,
\Ax{T} सर्व परावर्ती प्रतिमानांत सत्य असला,
तरी \Ax{T} सत्य असलेल्या प्रत्येक प्रतिमानाचा संबंध
परावर्ती असतोच असे नाही. परंतु असे \emph{खरे} आहे की,
\Ax{T} सर्व परावर्ती \emph{चौकटींवर} \emph{वैध} आहे,
आणि \Ax{T} वैध
असलेली प्रत्येक चौकट परावर्ती आहे, ही दोन्ही विधाने
खरी आहेत.

\begin{rem}
चौकटींच्या वर्गातील वैधता ही प्रतिमानांच्या वर्गातील
वैधतेच्या संकल्पनेचे विशेष प्रकरण आहे: $\mClass{F} \Entails \varphi$
तेव्हा आणि केवळ तेव्हाच, जेव्हा $\mClass{C} \Entails \varphi$,
जिथे $\mClass{C}$ हा~$\mClass{F}$ मधील एखाद्या
चौकटीवर आधारित सर्व प्रतिमानांचा वर्ग आहे.

स्पष्टच आहे की, सूत्र किंवा रूपबंध वैध असेल,
म्हणजे \emph{सर्व} प्रतिमानांच्या वर्गाच्या संदर्भात वैध
असेल, तर तो चौकटींच्या कोणत्याही वर्ग~$\mClass{F}$
च्या संदर्भातही वैध असतो.
\end{rem}













\section{चौकटींची परिभाष्यता}\label{nml:frd:def:sec}

\ref{nml:frd:acc:thm:soundschemas} मधील व्यत्यास प्रतिमानांसाठी
असत्य असले, तरी ``प्रतिमान'' याऐवजी ``चौकट'' घेतल्यास
ते खरे ठरतात: \ref{nml:frd:acc:thm:soundschemas} मध्ये पाहिलेल्या
गुणधर्मांसाठी, सूत्र एखाद्या \emph{चौकटीत} वैध
असेल, तर त्या चौकटीच्या प्राप्यता संबंधाला संबंधित गुणधर्म
\emph{असतोच}. म्हणून संबंधित गुणधर्म असलेल्या चौकटींचे
वर्ग हे सूत्रे \emph{परिभाषित करतात}.

\begin{defn}
  $\mClass{F}$ हा चौकटींचा वर्ग असेल, तर नेमक्या
  त्यातील चौकटी~$\mModel{F} \in \mClass{F}$
  आणि त्यांच्याच बाबतीत $\mModel{F} \Entails \varphi$ असेल,
  तर $\varphi$ हा $\mClass{F}$ वर्ग \emph{परिभाषित करतो},
  असे म्हणतो.
\end{defn}

आता चौकटींसाठी परिभाष्यतेचे पूर्ण निष्कर्ष सिद्ध करू.

\begin{thm}\label{nml:frd:def:thm:fullCorrespondence}
\ref{nml:frd:acc:tab:five} च्या उजव्या बाजूचे सूत्र
चौकट~$\mModel{F}$ मध्ये वैध असेल, तर $\mModel{F}$ ला
डाव्या बाजूचा गुणधर्म असतो.
\end{thm}

\begin{proof}
  \begin{enumerate}
  \item \Ax{D} हे $\mModel{F} = \tuple{W, R}$ मध्ये वैध आहे,
    म्हणजे $\mModel{F} \Entails \Box p \lif \Diamond p$,
    असे समजा. $\mModel{F}$ वर आधारित
    $\mModel{M} = \tuple{W, R, V}$ हे प्रतिमान आणि
    $w \in W$ घ्या. $Rwv$ असलेले एखादे $v$ आहे हे
    दाखवायचे आहे. तसे नाही असे समजा. मग~$p$ सहित
    कोणत्याही~$\varphi$ साठी $\mSat{M}{\Box \varphi}[w]$ आणि
    $\mSat/{M}{\Diamond \varphi}[w]$ ही दोन्ही विधाने खरी
    असतात. पण मग $\mSat/{M}{\Box p \lif
      \Diamond p}[w]$, आणि हे $\mModel{F}
    \Entails \Box p \lif \Diamond p$ या गृहीतकाला
    विरोधी आहे.
  \item \Ax{T} हे $\mModel{F}$ मध्ये वैध आहे,
    म्हणजे $\mModel{F} \Entails \Box
    p \lif p$, असे समजा. कोणतेही जग $w \in W$ घ्या;
    $Rww$ हे दाखवायचे आहे. $u \in V(p)$ तेव्हा आणि
    केवळ तेव्हाच, जेव्हा $Rwu$, असे मूल्यांकन
    ठरवा (आणि $p$ पेक्षा वेगळ्या $q$ साठी $V(q)$
    मनमाना ठेवा, उदाहरणार्थ $V(q) = \emptyset$).
    $\mModel{M} = \tuple{W, R, V}$ घ्या. या रचनेनुसार
    $Rwu$ असलेल्या सर्व $u$ साठी $\mSat{M}{p}[u]$;
    म्हणून $\mSat{M}{\Box p}[w]$. पण गृहीतकावरून
    $\Box p \lif p$ हे~$w$ येथे सत्य आहे; म्हणून
    $\mSat{M}{p}[w]$. $V$ च्या व्याख्येनुसार हे
    $Rww$ असेल तेव्हाच शक्य आहे.
  \item प्रतिपरिवर्तनाने सिद्ध करू: $\mModel{F}$ सममित
    नाही असे समजा. मग \Ax{B}, म्हणजे
    $p \lif \Box\Diamond p$, हे
    $\mModel{F}= \tuple{W, R}$ मध्ये वैध नाही हे
    दाखवू. $\mModel{F}$ सममित नसल्यामुळे
    $u$, $v \in W$ अशी जगं आहेत की $Ruv$ आहे,
    पण $Rvu$ नाही. $w \in V(p)$ तेव्हा आणि केवळ
    तेव्हाच, जेव्हा $Rvw$ नाही, असे $V$ ठरवा (इतरत्र
    $V$ मनमाना आहे). $\mModel{M} =\tuple{W, R,
      V}$ घ्या. $V$ च्या व्याख्येनुसार, $Rvw$
    नसलेल्या सर्व $w$ साठी $\mSat{M}{p}[w]$;
    विशेषतः $Rvu$ नसल्यामुळे $\mSat{M}{p}[u]$.
    तसेच $Rvw$ तेव्हा आणि केवळ तेव्हाच, जेव्हा
    $w \notin V(p)$; म्हणून $Rvw$ आणि
    $\mSat{M}{p}[w]$ या दोन्ही अटी पूर्ण करणारे
    कोणतेही~$w$ नाही. त्यामुळे
    $\mSat/{M}{\Diamond p}[v]$. $Ruv$ असल्याने
    $\mSat/{M}{\Box\Diamond p}[u]$ देखील. म्हणून
    $\mSat/{M}{p \lif
      \Box\Diamond p}[u]$, आणि त्यामुळे
    $\Ax{B}$ हे~$\mModel{F}$ मध्ये वैध नाही.
  \item \Ax{4} हे $\mModel{F} = \tuple{W,R}$ मध्ये
    वैध आहे, म्हणजे $\mModel{F} \Entails \Box p
    \lif \Box\Box p$, असे समजा. $u$, $v$, $w \in W$
    ही कोणतीही जगं $Ruv$ आणि $Rvw$ अशा अटी पूर्ण
    करत असू द्या; $Ruw$ सिद्ध करायचे आहे.
    $z \in V(p)$ तेव्हा आणि केवळ तेव्हाच, जेव्हा
    $Ruz$, असे $V$ ठरवा (इतरत्र $V$ मनमाना आहे).
    $\mModel{M} =\tuple{W, R, V}$ घ्या. $V$ च्या
    व्याख्येनुसार $Ruz$ असलेल्या सर्व $z$ साठी
    $\mSat{M}{p}[z]$; म्हणून $\mSat{M}{\Box p}[u]$.
    गृहीतकावरून \Ax{4}, म्हणजे $\Box p
    \lif \Box \Box p$, हे $u$ येथे सत्य आहे;
    म्हणून $\mSat{M}{\Box \Box
      p}[u]$. $Ruv$ आणि $Rvw$ असल्यामुळे
    $\mSat{M}{p}[w]$; पण $V$ च्या व्याख्येनुसार
    हे $Ruw$ असेल तेव्हाच शक्य आहे.
  \item पुन्हा प्रतिपरिवर्तन वापरू. चौकट
    $\mModel{F} = \tuple{W, R}$ युक्लिडी नाही असे
    मानून ती \Ax{5} असत्य ठरवते, म्हणजे
    $\mModel{F} \Entails/ \Diamond p \lif
      \Box\Diamond p$, हे दाखवू. $u$, $v$, $w \in W$
    अशी जगं असू द्या की $Rwu$ आणि $Rwv$ आहेत,
    पण $Ruv$ नाही. सर्व जगं $z$ साठी $z \in V(p)$
    तेव्हा आणि केवळ तेव्हाच, जेव्हा $Ruz$ असे
    \emph{नाही}, असे $V$ ठरवा. $\mModel{M}
    =\tuple{W, R, V}$ घ्या. मग गृहीतकावरून
    $\mSat{M}{p}[v]$, आणि $Rwv$ असल्यामुळे
    $\mSat{M}{\Diamond p}[w]$ देखील. परंतु
    $Ruy$ आणि $\mSat{M}{p}[y]$ असलेले कोणतेही
    जग~$y$ नाही; म्हणून $\mSat/{M}{\Diamond
      p}[u]$. $Rwu$ असल्यामुळे
    $\mSat/{M}{\Box\Diamond
      p}[w]$; म्हणून \Ax{5}, म्हणजे
    $\Diamond p \lif \Box\Diamond p$, हे~$w$
    येथे असत्य आहे.
  \end{enumerate}
\end{proof}

\Ax{D} ची सिद्धता आणि इतर प्रकरणे यांत एक फरक दिसेल:
त्यात मूल्यांकन~$V$ ची कोणतीही विशेष निवड करावी लागली नाही. प्रत्यक्षात,
$\mSat{M}{\Ax{D}}$ असल्यास $\mModel{M}$ सर्वत्र
उत्तरवर्ती आहे, हे आपण सिद्ध केले. म्हणून \Ax{D}
हा रूपबंध केवळ चौकटींचा नव्हे, तर सर्वत्र उत्तरवर्ती
\emph{प्रतिमानांचा} वर्ग परिभाषित करतो.

\begin{cor}\label{nml:frd:def:prop:D-serial}
  \Ax{D} सत्य असलेले कोणतेही प्रतिमान सर्वत्र उत्तरवर्ती असते.
\end{cor}

\begin{cor}
\ref{nml:frd:acc:tab:five} च्या उजव्या बाजूचे प्रत्येक
सूत्र डावीकडील गुणधर्म असलेल्या चौकटींचा
वर्ग परिभाषित करते.
\end{cor}

\begin{proof}
  \ref{nml:frd:acc:thm:soundschemas} मध्ये आपण सिद्ध केले की,
  प्रतिमानाला डावीकडील गुणधर्म असेल, तर उजवीकडचे
  सूत्र त्यात सत्य असते. म्हणून चौकट~$\mModel{F}$
  ला तो गुणधर्म असेल, तर उजवीकडचे सूत्र
  $\mModel{F}$ मध्ये वैध असते.
  \ref{nml:frd:def:thm:fullCorrespondence} मध्ये आपण व्यत्यास
  सिद्ध केले: उजवीकडचे सूत्र~$\mModel{F}$
  मध्ये वैध असेल, तर $\mModel{F}$ ला डावीकडचा
  गुणधर्म असतो.
\end{proof}

\begin{prob}
\ref{nml:frd:acc:tab:anotherfive} च्या उजव्या
बाजूचे सूत्र चौकट~$\mModel{F}$ मध्ये वैध असेल,
तर $\mModel{F}$ ला डावीकडील गुणधर्म असतो, हे दाखवा.
यासाठी डावीकडील गुणधर्म \emph{नसलेली} चौकट घ्या,
आणि उजवीकडचे सूत्र एखाद्या जगात असत्य ठरेल
असे योग्य~$V$ ठरवा.
\end{prob}

\ref{nml:frd:def:thm:fullCorrespondence} वरून गुणधर्म एकत्रही
करता येतात हे दिसते: उदाहरणार्थ \Ax{B} आणि
\Ax{4} हे दोन्ही~$\mModel{F}$ मध्ये वैध असतील,
तर चौकट सममित आणि संक्रमक दोन्ही असते.
अनेक महत्त्वाची मोडल तर्कशास्त्रे काही चौकट-गुणधर्म
एकत्र असलेल्या सर्व चौकटींमध्ये वैध असलेल्या
सूत्रांचे संच म्हणून निश्चित होतात. म्हणून
त्यांना संबंधित परिभाषक सूत्रे वैध असलेल्या
सर्व चौकटींमध्ये वैध असलेल्या सूत्रांचे संच
म्हणूनही सांगता येते. उदाहरणार्थ, अभिजात व्यवस्था
\Log{S4} म्हणजे सर्व परावर्ती आणि संक्रमक चौकटींमध्ये
वैध असलेल्या सूत्रांचा संच; म्हणजे \Ax{T}
आणि~\Ax{4} दोन्ही वैध असलेल्या सर्व चौकटींमध्ये
वैध असलेल्या सूत्रांचा संच.
\Log{S5} म्हणजे सर्व परावर्ती, सममित आणि युक्लिडी
चौकटींमध्ये वैध असलेल्या सूत्रांचा संच; म्हणजे
\Ax{T}, \Ax{B}, आणि \Ax{5} ही तिन्ही वैध
असलेल्या सर्व चौकटींमध्ये वैध असलेल्या सूत्रांचा संच.

$R$ च्या गुणधर्मांमधील तार्किक संबंध सामान्यतः
संबंधित परिभाषक सूत्रे मधील संबंधांशी जुळतात.
उदाहरणार्थ, प्रत्येक परावर्ती संबंध सर्वत्र उत्तरवर्ती
असतो; म्हणून चौकटीत \Ax{T} वैध असेल, तर
\Ax{D} देखील वैध असतो. (हा संबंध \emph{तार्किक
निष्पन्नतेचा नाही}, हे लक्षात घ्या. $\mSat{M}{\Ax{T}}[w]$
असल्यावर नेहमी $\mSat{M}{\Ax{D}}[w]$ असते, असे
नाही.) असे काही संबंध पुढे नोंदवतो.

\begin{prop}\label{nml:frd:def:prop:relation-facts}
  $R$ हा $W$ या संचावरील द्विस्थानी संबंध असू द्या. मग:
  \begin{enumerate}
  \item $R$ परावर्ती असेल, तर तो सर्वत्र उत्तरवर्ती असतो.
  \item $R$ सममित असेल, तर तो संक्रमक असतो तेव्हा आणि
    केवळ तेव्हाच, जेव्हा तो युक्लिडी असतो.
  \item $R$ सममित किंवा युक्लिडी असेल, तर तो दुर्बल
    दिशित असतो (त्याला ``हिऱ्याचा गुणधर्म'' असतो).
  \item $R$ युक्लिडी असेल, तर तो दुर्बल संयोजित असतो.
  \item $R$ फलनीय असेल, तर तो सर्वत्र उत्तरवर्ती असतो.
  \end{enumerate}
\end{prop}

\begin{prob}
  \ref{nml:frd:def:prop:relation-facts} सिद्ध करा.
\end{prob}












\section{प्रथम-क्रमातील परिभाष्यता}\label{nml:frd:fol:sec}

चौकटींच्या प्राप्यता संबंधांचे अनेक गुणधर्म मोडल
सूत्रे नी परिभाषित करता येतात, हे आपण पाहिले.
उदाहरणार्थ, चौकटींची सममिती \Ax{B}, म्हणजे
$p \lif \Box\Diamond
p$, या सूत्र ने परिभाषित होते. आतापर्यंत
पाहिलेल्या सर्व अटी एकच द्विस्थानी विधेय
असलेल्या भाषा तील प्रथम-क्रम सूत्रे
नी व्यक्त करता येतात. उदाहरणार्थ, सममिती
$\lforall[x][\lforall[y][(\Atom{Q}{x,y} \lif \Atom{Q}{y, x})]]$
याने परिभाषित होते: $\Domain{M} =
W$ आणि $\Assign{Q}{M} = R$ असलेली प्रथम-क्रम
संरचना~$\Struct{M}$ वरील सूत्राची पूर्ति करते
तेव्हा आणि केवळ तेव्हाच, जेव्हा $R$ सममित असतो.
यातून पुढील व्याख्या सुचते:

\begin{defn}
  चौकटींचा वर्ग~$\mClass{F}$ \emph{प्रथम-क्रमात
  परिभाष्य} असतो, जर एकच द्विस्थानी विधेय~$Q$
  असलेल्या प्रथम-क्रम भाषेतील असे वाक्य~$\varphi$
  असेल की, $\Domain{M} = W$ आणि $\Assign{Q}{M}
  = R$ असलेल्या प्रथम-क्रम संरचना~$\Struct{M}$
  मध्ये $\Sat{M}{\varphi}$ तेव्हा आणि केवळ तेव्हाच,
  जेव्हा $\mModel{F} =
  \tuple{W, R} \in \mClass{F}$.
\end{defn}

आतापर्यंत पाहिलेले गुणधर्म आणि त्यांना परिभाषित करणारी
मोडल सूत्रे ही अपवादात्मक उदाहरणे ठरतात.
प्रत्येक सूत्र प्रथम-क्रमात परिभाष्य चौकटींचाच
वर्ग परिभाषित करते असे नाही; तसेच प्रथम-क्रमात
परिभाष्य चौकटींचा प्रत्येक वर्ग मोडल सूत्राने
परिभाषित होतो असेही नाही.

पहिल्या दाव्याचे प्रतिउदाहरण L\"ob सूत्र देते:
\begin{equation}
\Box(\Box p \lif p) \lif \Box p. \tag{\Ax{W}}
\end{equation}
\Ax{W} संक्रमक आणि प्रतिलोम-सुस्थापित चौकटींचा वर्ग
परिभाषित करते. $Rw_2w_1$, $Rw_3w_2$, \dots{}
अशी $w_1$, $w_2$, \dots{} ही अनंत श्रेणी नसल्यास
संबंध सुस्थापित असतो. उदाहरणार्थ, $\Nat$ वरील
$<$ हा संबंध सुस्थापित आहे; पण $\Int$ वरील
$<$ सुस्थापित नाही. एखाद्या संबंधाचा प्रतिलोम
सुस्थापित असेल, तर आणि तेव्हाच, तो संबंध प्रतिलोम-सुस्थापित
असतो. म्हणजे प्रतिलोम-सुस्थापित संबंधात $Rw_1w_2$,
$Rw_2w_3$, \dots{} अशी $w_1$, $w_2$, \dots{}
अनंत श्रेणी नसते.

मात्र संक्रमक प्रतिलोम-सुस्थापित संबंध परिभाषित करणारे
प्रथम-क्रम सूत्र नाही. असे काही सूत्र असल्याचे
समजून, $R =
\Assign{Q}{M}$ संक्रमक आणि प्रतिलोम-सुस्थापित असेल
तेव्हा आणि केवळ तेव्हाच $\Sat{M}{\phi}$, असे माना.
नवे स्थिरांक असलेल्या विस्तारित भाषेत $\varphi_n$ हे
पुढील सूत्र असू द्या:
\[(\Atom{Q}{a_1,a_2} \land \dots \land
    \Atom{Q}{a_{n},a_{n+1}})
\]
आता पुढील सूत्रांचा संच पाहू:
\[\Gamma = \{\phi, \varphi_1, \varphi_2, \dots\}.
\]
$\Gamma$ चा प्रत्येक सांत उपसंच पूर्ततायोग्य आहे.
उपसंचात $\varphi_k$ असेल, तर $k$ हा त्यातील सर्वात मोठा
निर्देशांक घ्या; असे सूत्र नसेल, तर त्याचे मूल्य एक
घ्या. $\Domain{M_k} = \{1, \dots, k+1\}$,
आवश्यक स्थिरांकांसाठी $\Assign{a_i}{M_k} = i$,
आणि $\Assign{Q}{M_k} = <$ असे ठेवा; उर्वरित
स्थिरांकांना प्रांतातील कोणतेही मूल्य द्या.
$\{1,
\dots, k+1\}$ वरील $<$ संक्रमक आणि प्रतिलोम-सुस्थापित
असल्यामुळे $\Sat{M_k}{\phi}$. रचनेनुसार प्रत्येक
$i \le k$ साठी $\Sat{M_k}{\varphi_i}$.
प्रथम-क्रम तर्कशास्त्राच्या संहतता प्रमेयावरून
$\Gamma$ ची कोणत्यातरी संरचना~$\Struct{M}$
मध्ये पूर्ति होते. गृहीतकावरून $\Sat{M}{\phi}$
असल्यामुळे $\Assign{Q}{M}$ हा संबंध प्रतिलोम-सुस्थापित
आहे. पण $\Assign{a_1}{M}$, $\Assign{a_2}{M}$,
\dots{} यांपासून प्रतिलोम-सुस्थापितपणाने निषिद्ध
ठरवलेली अनंत श्रेणी मिळते.

दुसऱ्या दाव्याचे प्रतिउदाहरण म्हणजे वैश्विकता:
प्रत्येक $u$ आणि $v$ साठी $Ruv$. वैश्विक चौकटी
$\lforall[x][\lforall[y][\Atom{Q}{x,y}]]$ या
सूत्राने प्रथम-क्रमात परिभाषित होतात. मात्र एखादे
मोडल सूत्र नेमक्या वैश्विक चौकटींमध्येच वैध
असते असे नाही. हा निष्कर्ष स्वतंत्रपणे महत्त्वाच्या
एका तथ्यावरून येतो: वैश्विक चौकटींमध्ये वैध
असलेली सूत्रे आणि परावर्ती, सममित व संक्रमक
चौकटींमध्ये वैध असलेली सूत्रांनीमकी तीच आहेत.
परावर्ती, सममित आणि संक्रमक असूनही वैश्विक
नसलेल्या चौकटी आहेत. म्हणून सर्व वैश्विक चौकटींमध्ये
वैध असलेले प्रत्येक सूत्र काही अवैश्विक
चौकटींमध्येही वैध असते.












\section{तुल्यता संबंध आणि \Log{S5}}\label{nml:frd:es5:sec}

सर्व वैश्विक चौकटींमध्ये वैध असलेल्या सूत्रे
चा संच म्हणून मोडल तर्कशास्त्र~\Log{S5} निश्चित होते;
अशा चौकटीत प्रत्येक जग स्वतःसह प्रत्येक जगापासून
प्राप्य असते. येथे $\Box$ आवश्यकतेशी आणि
$\Diamond$ शक्यतेशी जुळतो: $\varphi$ \emph{प्रत्येक}
जगात सत्य असेल तर $\Box \varphi$ सत्य असते, आणि
$\varphi$ \emph{किमान एका} जगात सत्य असेल तर
$\Diamond \varphi$ सत्य असते. \Log{S5} हे सर्व परावर्ती,
सममित आणि संक्रमक चौकटींमध्ये, म्हणजे सर्व
\emph{तुल्यता संबंधांवर}, वैध असलेल्या सूत्रे
चा संच म्हणूनही सांगता येते.

\begin{defn}
  $W$ वरील द्विस्थानी संबंध $R$ हा \emph{तुल्यता संबंध}
  असतो तेव्हा आणि केवळ तेव्हाच, जेव्हा तो परावर्ती,
  सममित आणि संक्रमक असतो. $W$ वरील संबंध $R$
  \emph{वैश्विक} असतो तेव्हा आणि केवळ तेव्हाच, जेव्हा
  सर्व $u,v \in
  W$ साठी $Ruv$ असते.
\end{defn}

\Ax{T}, \Ax{B}, आणि \Ax{4} हे अनुक्रमे परावर्ती,
सममित आणि संक्रमक चौकटींचे वैशिष्ट्य दर्शवतात.
म्हणून प्राप्यता संबंध तुल्यता संबंध असलेल्या चौकटी
नेमक्या त्या आहेत ज्यांत ही तिन्ही सूत्रे वैध
आहेत. $R$ चे इतर परिचित गुणधर्म एकत्र केल्यास
तुल्यता संबंध परिभाषित करणाऱ्या अटींच्या समतुल्य
अटी मिळतात; म्हणून सूत्रांच्या इतर संयोगांनीही
तुल्यता संबंधांचे वैशिष्ट्य सांगता येते.

\begin{prop}\label{nml:frd:es5:prop:equivalences}
  पुढील अटी समतुल्य आहेत:
  \begin{enumerate}
  \item $R$ हा तुल्यता संबंध आहे;
  \item $R$ परावर्ती आणि युक्लिडी आहे;
  \item $R$ सर्वत्र उत्तरवर्ती, सममित आणि युक्लिडी आहे;
  \item $R$ सर्वत्र उत्तरवर्ती, सममित आणि संक्रमक आहे.
  \end{enumerate}
\end{prop}

\begin{proof}
  सरावासाठी.
\end{proof}

\begin{prob}
  पुढील विधाने दाखवून \ref{nml:frd:es5:prop:equivalences}
  सिद्ध करा:
  \begin{enumerate}
  \item $R$ सममित आणि संक्रमक असेल, तर तो युक्लिडी आहे.
  \item $R$ परावर्ती असेल, तर तो सर्वत्र उत्तरवर्ती आहे.
  \item $R$ परावर्ती आणि युक्लिडी असेल, तर तो सममित आहे.
  \item $R$ सममित आणि युक्लिडी असेल, तर तो संक्रमक आहे.
  \item $R$ सर्वत्र उत्तरवर्ती, सममित आणि संक्रमक असेल,
    तर तो परावर्ती आहे.
  \end{enumerate}
  या विधानांवरून सर्व अटी समतुल्य आहेत, हे कसे
  सिद्ध होते ते स्पष्ट करा.
\end{prob}

\ref{nml:frd:es5:prop:equivalences} हे \ref{nml:prf:prs:prop:S5}
चे चिन्हार्थविषयक प्रतिरूप आहे: $R$ तुल्यता संबंध
असलेल्या चौकटींच्या मोडल तर्कशास्त्राचे, परंपरेने
\Log{S5} म्हणतात त्याचे, ते समतुल्य वैशिष्ट्यदर्शन देते.

वैश्विक संबंध आणि तुल्यता संबंध यांचे नाते काय?
प्रत्येक वैश्विक संबंध तुल्यता संबंध असला, तरी
प्रत्येक तुल्यता संबंध वैश्विक असेलच असे नाही.
तरी सर्व वैश्विक संबंधांवर वैध असलेली सूत्रे
आणि सर्व तुल्यता संबंधांवर वैध असलेली सूत्रे
नेमकी तीच आहेत.

\begin{prop}
  $R$ हा तुल्यता संबंध असू द्या. प्रत्येक $w \in W$
  साठी $w$ चा \emph{तुल्यतावर्ग} म्हणजे
  $[w] = \{w'\in W :
  Rww'\}$ हा संच ठरवा. मग:
  \begin{enumerate}
  \item $w \in [w]$;
  \item प्रत्येक तुल्यतावर्ग~$[w]$ वर $R$ वैश्विक आहे;
  \item तुल्यतावर्ग $W$ ला परस्पर विसंयुक्त आणि
    एकत्रितपणे संपूर्ण W व्यापणाऱ्या उपसंचांत विभागतात.
  \end{enumerate}
\end{prop}

\begin{prop}\label{nml:frd:es5:prop:S5=univ}
  सूत्र~$\varphi$ हे $R$ तुल्यता संबंध असलेल्या
  सर्व चौकटींमध्ये $\mModel{F} =\tuple{W,R}$ वैध
  असते तेव्हा आणि केवळ तेव्हाच, जेव्हा ते $R$
  वैश्विक असलेल्या सर्व चौकटींमध्ये
  $\mModel{F} =\tuple{W,R}$ वैध असते. म्हणून
  वैश्विक चौकटींचे तर्कशास्त्र म्हणजे \Log{S5}.
\end{prop}

\begin{proof}
  $W$ वरील वैश्विक संबंध~$R$ हा तुल्यता संबंध असतो,
  हे लगेच तपासता येते. म्हणून $\varphi$ हे $R$ तुल्यता
  संबंध असलेल्या सर्व चौकटींमध्ये वैध असेल, तर ते
  सर्व वैश्विक चौकटींमध्येही वैध असते. उलट दिशेसाठी
  प्रतिपरिवर्तनाने युक्तिवाद करू. $W$ वरील $R$
  तुल्यता संबंध असलेल्या चौकट~$\tuple{W,R}$ वर
  आधारित $\mModel{M} = \tuple{W, R, V}$ या
  प्रतिमानात $\psi$ हे सूत्र एखाद्या जगात~$w$
  असत्य आहे असे समजा. म्हणजे $\mSat/{M}{\psi}[w]$.
  पुढीलप्रमाणे $\mModel{M}' = \tuple{W',
    R', V'}$ हे प्रतिमान ठरवा:
  \begin{enumerate}
  \item $W' = [w]$;
  \item $R'$ हा $W'$ वर वैश्विक आहे;
  \item $V'(p) = V(p) \cap W'$.
  \end{enumerate}
  (\ref{nml:frd:es5:fig:partition} मधील छायांकित भाग
  $\mModel{M}'$ चा जगांचा संच~$W'$ दाखवतो.)
  $W'$ वर $R$ आणि $R'$ हे संबंध जुळतात, हे सहज
  दिसते. मग सूत्रे वर विगमन करून, प्रत्येक
  $\varphi$ आणि सर्व $w' \in W'$ साठी
  $\mSat{M'}{\varphi}[w']$ तेव्हा आणि केवळ तेव्हाच,
  जेव्हा $\mSat{M}{\varphi}[w']$, हे दाखवता येते
  (कारण $W'
  \subseteq W$). विशेषतः $\mSat/{M'}{\psi}[w]$;
  म्हणून $\psi$ हे वैश्विक चौकटीवर आधारित
  प्रतिमानात असत्य आहे.
\end{proof}

\begin{figure}[t]
  \centering
  \begin{tikzpicture}[node distance=2cm, auto, thick]
    \clip [rounded corners] (0,0) -- (6,0) -- (6,4) -- (0,4) --  cycle;
    \begin{scope}
      \clip (0,0) -- (2,0)  .. controls (1.5,1.5) and (4.5,1.5)
      .. (4,4) -- (0,4) -- (0,0) -- cycle;
      \filldraw[fill=gray!40] (0,4) -- (0,2) .. controls (1.5,1.5)
      and (4.5,1.5) .. (4.5,0) -- (6,0) -- (6,4) -- (0,4) --  cycle;
    \end{scope}
    \draw [name path=line1] (2,0) .. controls (1.5,1.5) and (4.5,1.5) .. (4,4);
    \draw [name path=line2] (0,2)  .. controls (1.5,1.5) and (4.5,1.5) .. (4.5,0);
    \path [name intersections={of=line1 and line2}];
    \draw [rounded corners, use as bounding box, very thick] (0,0)
    -- (6,0) -- (6,4) -- (0,4) --  cycle;
    \node at (2,3) {$[w]$} ;
    \node at (1,1) {$[u]$} ;
    \node at (3,0.75) {$[v]$} ;
    \node at (4.75,2) {$[z]$} ;
  \end{tikzpicture}
  \caption{$W$ चा तुल्यतावर्गांमध्ये विभाग.}
  \label{nml:frd:es5:fig:partition}
\end{figure}













\section{द्वितीय-क्रमातील परिभाष्यता}\label{nml:frd:st:sec}

मोडल सूत्रांनी परिभाषित होणारा चौकटींचा प्रत्येक गुणधर्म
प्रथम-क्रमात परिभाष्य असेलच असे नाही. पण एकस्थानी
विधेयांवर संख्यापन करण्याची परवानगी दिल्यास (म्हणजे
एकस्थानी द्वितीय-क्रम संख्यापन), मोडल सूत्रांनी परिभाष्य
असलेले सर्व चौकट-गुणधर्म परिभाषित करता येतात.
एखादे मोडल सूत्र जगात सत्य असण्याच्या अटी
प्रथम-क्रम सूत्रे शी ज्या पद्धतशीर रीतीने
संबंधित असतात, तिचा उपयोग करणे ही युक्ती आहे.
या पद्धतीला मोडल सूत्रांचे प्रथम-क्रम सूत्रांत
मानक भाषांतर म्हणतात. त्यासाठीच्या भाषेत
प्राप्यता संबंधासाठी द्विस्थानी विधेय~$Q$
तर असतेच, शिवाय~$\varphi$ मध्ये येणाऱ्या
विधानीय चले~$\Obj p_i$ साठी
एकस्थानी विधेय~$P_i$ देखील असते.

\begin{defn}
  \emph{मानक भाषांतर}~$\ST_x(\varphi)$ ची विगमनाने
  पुढीलप्रमाणे व्याख्या करतो. मोडल उपवाक्यांत $x$ पेक्षा वेगळा
  नवीन व्यक्तिचल $y$ निवडा. आतील भाषांतरातील संख्यापकांसाठीही
  मुक्त चल बद्ध होणार नाहीत अशी नवीन बद्ध चलांची नावे निवडावीत:
  \begin{enumerate}
  \item \indcase{\varphi}{\lfalse}{$\ST_x(\indfrm) = \lfalse$.}
  \item \indcase{\varphi}{\ltrue}{$\ST_x(\indfrm) = \ltrue$.}
  \item \indcase{\varphi}{\Obj p_i}{$\ST_x(\indfrm) = \Atom{P_i}{x}$.}
  \item \indcase{\varphi}{\lnot \psi}{$\ST_x(\indfrm) = \lnot \ST_x(\psi)$.}
  \item \indcase{\varphi}{(\psi \land \chi)}{$\ST_x(\indfrm) = (\ST_x(\psi)
      \land \ST_x(\chi))$.}
  \item \indcase{\varphi}{(\psi \lor \chi)}{$\ST_x(\indfrm) = (\ST_x(\psi)
      \lor \ST_x(\chi))$.}
  \item \indcase{\varphi}{(\psi \lif \chi)}{$\ST_x(\indfrm) = (\ST_x(\psi)
      \lif \ST_x(\chi))$.}
  \item \indcase{\varphi}{(\psi \liff \chi)}{$\ST_x(\indfrm) = (\ST_x(\psi)
      \liff \ST_x(\chi))$.}
  \item \indcase{\varphi}{\Box \psi}{$\ST_x(\indfrm) =
      \lforall[y][(\Atom{Q}{x,y} \lif \ST_y(\psi))]$.}
  \item \indcase{\varphi}{\Diamond \psi}{$\ST_x(\indfrm) =
      \lexists[y][(\Atom{Q}{x,y} \land \ST_y(\psi))]$.}
  \end{enumerate}
\end{defn}

उदाहरणार्थ, $\ST_x(\Box p \lif p)$ हे
$\lforall[y][(\Atom{Q}{x,y}
  \lif \Atom{P}{y})] \lif \Atom{P}{x}$ असे आहे.
~$\ST_x(\varphi)$ च्या भाषेतील कोणत्याही संरचना
साठी प्रांत, $Q$ ला नेमून दिलेला द्विस्थानी संबंध,
आणि एकस्थानी विधेये~$P_i$ यांना नेमून
दिलेले प्रांताचे उपसंच लागतात. दुसऱ्या शब्दांत,
अशा संरचनेचे घटक हे~$\varphi$ च्या प्रतिमानाच्या
घटकांशी तंतोतंत जुळतात: प्रांत म्हणजे जगांचा संच;
$Q$ ला नेमलेला द्विस्थानी संबंध म्हणजे प्राप्यता
संबंध; आणि $P_i$ यांना नेमलेले उपसंच म्हणजे
नेमणुका~$V(\Obj p_i)$. त्यामुळे मोडल प्रतिमानातील
$\varphi$ ची पूर्ति आणि अनुरूप संरचना मधील
$\ST_x(\varphi)$ ची पूर्ति जुळते, हे अपेक्षितच आहे:

\begin{prop}\label{nml:frd:st:prop:st}
  $\mModel{M} = \tuple{W, R, V}$ असू द्या.
  $\Domain{M'} = W$, $\Assign{Q}{M'} = R$, आणि
  $\Assign{P_i}{M'} = V(\Obj p_i)$ असलेली
  $\Struct{M'}$ ही प्रथम-क्रम संरचना असू द्या,
  तसेच $s(x) = w$ असू द्या. मग
  \[  \mSat{M}{\varphi}[w] \text{ तेव्हा आणि केवळ तेव्हाच } \Sat{M'}{\ST_x(\varphi)}[s]
  \]
\end{prop}

\begin{proof}
  $\varphi$ वर विगमनाने.
\end{proof}

\begin{prop}
  $\varphi$ हे मोडल सूत्र आणि
  $\mModel{F} = \tuple{W, R}$ ही चौकट असू द्या.
  $\Domain{F'} = W$ आणि $\Assign{Q}{F'} = R$
  असलेली $\Struct{F'}$ ही प्रथम-क्रम संरचना
  असू द्या. पुढील द्वितीय-क्रम सूत्राला $\varphi'$
  म्हणा:
  \[  \lforall[X_1][\dots\lforall[X_n][\lforall[x][
        \SSubst{\ST_x(\varphi)}{\subst{X_1}{P_1}, \dots,
          \subst{X_n}{P_n}}]]],
  \]
  जिथे $P_1$, \dots, $P_n$ ही~$\ST_x(\varphi)$
  मधील सर्व एकस्थानी विधेये आहेत. मग
  \[  \mModel{F} \Entails \varphi \text{ तेव्हा आणि केवळ तेव्हाच } \Sat{F'}{\varphi'}
  \]
\end{prop}

\begin{proof}
  $\Sat{F'}{\varphi'}$ तेव्हा आणि केवळ तेव्हाच, जेव्हा
  $i = 1$, \dots,~$n$ साठी
  $\Assign{P_i}{M'} \subseteq W$ असलेल्या
  मूळ रचनेच्या प्रत्येक विस्ताररूप
  संरचना~$\mModel{M'}$ मध्ये, आणि
  $s(x) \in W$ असलेल्या प्रत्येक $s$ साठी,
  $\Sat{M'}{\ST_x(\varphi)}[s]$ असते.
  \ref{nml:frd:st:prop:st} नुसार, हे तेव्हा आणि केवळ
  तेव्हाच, जेव्हा $\mModel{F}$ वर आधारित सर्व
  प्रतिमाने~$\mModel{M}$ आणि प्रत्येक जग~$w \in W$
  यांसाठी $\mSat{M}{\varphi}[w]$ असते; म्हणजे
  $\mModel{F} \Entails \varphi$.
\end{proof}

\begin{defn}
  चौकटींचा वर्ग~$\mClass{F}$ \emph{द्वितीय-क्रमात
  परिभाष्य} असतो, जर एकच द्विस्थानी
  विधेय~$P$ असलेल्या आणि द्वितीय-क्रमाचे संख्यापक
  केवळ एकस्थानी संचचलांवर असलेल्या भाषेत असे
  वाक्य~$\varphi$ असेल की,
  $\Domain{M} = W$ आणि $\Assign{P}{M} = R$
  असलेल्या संरचना~$\Struct{M}$ मध्ये
  $\Sat{M}{\varphi}$ तेव्हा आणि केवळ तेव्हाच,
  जेव्हा $\mModel{F} = \tuple{W, R} \in \mClass{F}$.
  व्यक्तिचलांवरील प्रथम-क्रम संख्यापकही अनुमत आहेत.
\end{defn}

\begin{cor}
  चौकटींचा वर्ग एखाद्या सूत्र~$\varphi$ ने
  परिभाष्य असेल, तर त्याच्या अनुरूप प्राप्यता
  संबंधांचा वर्ग एकस्थानी द्वितीय-क्रम वाक्याने
  परिभाष्य असतो.
\end{cor}

\begin{proof}
  आधीच्या सिद्धतेतील एकस्थानी द्वितीय-क्रम
  वाक्य~$\varphi'$ ला आवश्यक गुणधर्म आहे.
\end{proof}

उदाहरण म्हणून पुन्हा सूत्र~$\Box p \lif p$
पाहा. ते परावर्तितेची व्याख्या करते. अर्थात,
परावर्तिता वाक्य~$\lforall[x][\Atom{Q}{x,x}]$
याने प्रथम-क्रमात परिभाष्य आहे. पण ती पुढील
एकस्थानी द्वितीय-क्रम वाक्य नेही
परिभाष्य आहे:
\[\lforall[X][\lforall[x][(\lforall[y][(\Atom{Q}{x,y} \lif \Atom{X}{y})]
    \lif \Atom{X}{x})]].
\]
म्हणजे ही दोन वाक्ये समतुल्य आहेत. आधी द्वितीय-क्रम वाक्य
संरचना~$\Struct{M}$ मध्ये सत्य आहे असे समजा.
कोणतेही $w \in W$ घ्या. $R=\Assign{Q}{M}$ असू द्या आणि
$S=\Setabs{z \in W}{Rwz}$ ठेवा. $s(x)=w$ आणि $s(X)=S$
असलेली नेमणूक $s$ निवडा. वैश्विक संख्यापनाच्या व्याख्येवरून
\[\Sat{M}{\lforall[y][(\Atom{Q}{x,y} \lif \Atom{X}{y})]
    \lif \Atom{X}{x}}[s].
\]
$S$ च्या निवडीवरून या अभिव्यंजनाचे पूर्वांग $s$ नुसार सत्य आहे:
$w$ पासून प्राप्य असलेले प्रत्येक जग $S$ मध्ये आहे.
म्हणून $\Sat{M}{\Atom{X}{x}}[s]$; म्हणजे $w \in S$.
$S$ च्या व्याख्येनुसार $Rww$. $w$ कोणतेही असल्यामुळे
$\Sat{M}{\lforall[x][\Atom{Q}{x,x}]}$ मिळते.

आता $\Sat{M}{\lforall[x][\Atom{Q}{x,x}]}$
असे समजा आणि
$\Sat{M}{\lforall[X][\lforall[x][(\lforall[y][(\Atom{Q}{x,y} \lif
        \Atom{X}{y})] \lif \Atom{X}{x})]]}$
हे दाखवू. कोणतीही नेमणूक~$s$ घ्या आणि
$\Sat{M}{\lforall[y][(\Atom{Q}{x,y} \lif
    \Atom{X}{y})]}[s]$ असे माना. $s'$ ही~$s$
ची $y$-भिन्नरूप नेमणूक असू द्या, जिथे $s'(y)
= s(x)$; मग
$\Sat{M}{\Atom{Q}{x,y} \lif \Atom{X}{y}}[s']$,
म्हणजे $\Sat{M}{\Atom{Q}{x,x} \lif \Atom{X}{x}}[s]$.
$\Sat{M}{\lforall[x][\Atom{Q}{x,x}]}$ असल्यामुळे
पूर्वांग सत्य आहे; म्हणून $\Sat{M}{\Atom{X}{x}}[s]$,
आणि हेच दाखवायचे होते.

चौकटींचे काही परिभाष्य वर्ग प्रथम-क्रमात परिभाष्य
नसल्यामुळे, $\varphi'$ या रूपातील प्रत्येक एकस्थानी
द्वितीय-क्रम वाक्य हे प्रथम-क्रम
वाक्य शी समतुल्य असेलच असे नाही.
कोणती अशी आहेत हे ठरवण्याची प्रभावी पद्धत नाही.



\chapter{स्वयंसिद्धकीय निष्पत्ती}\label{nml:prf:chap}










\section{प्रस्तावना}\label{nml:axs:int:sec}

मूलभूत मोडल भाषेसाठी मोडल प्रतिमानांद्वारे चिन्हार्थमीमांसा
आपल्याकडे आहे. सूत्र वैध असणे, म्हणजे सर्व
प्रतिमानांतील सर्व जगांत सत्य असणे, किंवा प्रतिमानांच्या
किंवा चौकटींच्या एखाद्या वर्गाच्या संदर्भात वैध असणे,
म्हणजे त्या वर्गातील अथवा त्या चौकटीवर आधारित सर्व
प्रतिमानांच्या सर्व जगांत सत्य असणे, या संकल्पनाही
आपल्याकडे आहेत. तर्कशास्त्रात अशा चिन्हार्थविषयक
वैधतेला सहसा निष्पन्नता या सिद्धताविषयक
संकल्पनेशी जोडले जाते. एखाद्या पद्धतीत
निष्पन्नता अशी परिभाषित करायची की
सूत्र वैध असेल तेव्हा आणि केवळ तेव्हाच
निष्पन्न करता येण्याजोगे असेल, हे आपले उद्दिष्ट आहे.

सर्वात सोप्या आणि ऐतिहासिकदृष्ट्या सर्वात जुन्या
निष्पत्ती-पद्धती म्हणजे हिल्बर्ट-प्रकारच्या किंवा
स्वयंसिद्धकीय निष्पत्ती-पद्धती. अनेक मोडल
तर्कशास्त्रांसाठी हिल्बर्ट-प्रकारच्या निष्पत्ती-पद्धती तुलनेने सहज रचता येतात: अधितात्त्विक अभ्यासाच्या
दृष्टीने (उदा., निर्दोषता व संपूर्णता सिद्ध करताना)
त्या सोप्या असतात. पण त्यांमध्ये सूत्रे सिद्ध
करणे, उदाहरणार्थ नैसर्गिक निगमन पद्धतींपेक्षा, बरेच
कठीण असते.

हिल्बर्ट-प्रकारच्या निष्पत्ती-पद्धतीत
सूत्राची निष्पत्ती म्हणजे काही स्वयंसिद्धकांपासून
मोजक्या अनुमान नियमांनी मिळणारी, संबंधित
सूत्र पर्यंत पोचणारी सूत्रांची मालिका.
निष्पत्ती-पद्धत चिन्हार्थमीमांसेशी जुळावी अशी
आपली अपेक्षा असल्याने, निष्पन्न करता येण्याजोगे सूत्रे
सर्व प्रतिमानांत (किंवा स्वयंसिद्धकांचे सर्व आदेशन-दृष्टांत
सर्व जगांत सत्य असलेल्या प्रतिमानांत) सत्य आहेत याची खात्री करावी लागेल.
यासाठी आवश्यक असलेले मोडल तर्कशास्त्रांचे काही
गुणधर्म आधी वेगळे पाहू: ``सामान्य'' मोडल तर्कशास्त्रे.
सामान्य मोडल तर्कशास्त्रांसाठी फक्त दोन अनुमान नियम
गृहीत धरावे लागतात: विध्यात्मक अनुमान आणि
अनिवार्यीकरण. स्वयंसिद्धके म्हणून सर्वतः-सत्य
सूत्रांचे सर्व आदेशन-दृष्टांत आणि संबंधित मोडल
तर्कशास्त्रानुसार काही मोडल स्वयंसिद्धके घेतो.
सर्व प्रतिमानांच्या वर्गातच आपल्याला रस असला,
तरी~$\Ax{K}$ आणि~$\Ax{Dual}$
यांचे सर्व आदेशन-दृष्टांत स्वयंसिद्धके म्हणून घ्यावे
लागतात. एवढ्यानेच न्यूनतम सामान्य मोडल
तर्कशास्त्र~$\Log K$ मिळते.

\begin{defn}
\emph{विध्यात्मक अनुमान} हा नियम पुढील अनुमान-रूपबंध आहे:
\begin{prooftree}
\AxiomC{$\varphi$}
\AxiomC{$\varphi \lif \psi$}
\RightLabel{\MP}
\BinaryInfC{$\psi$}
\end{prooftree}
सूत्र~$\psi$ हे सूत्रे $\varphi$, $\chi$
यांपासून विध्यात्मक अनुमानाने \emph{निष्पन्न होते}
तेव्हा आणि केवळ तेव्हाच, जेव्हा
$\chi \ident \varphi \lif \psi$.
\end{defn}

\begin{defn}
\emph{अनिवार्यीकरण} हा नियम पुढील अनुमान-रूपबंध आहे:
\begin{prooftree}
\AxiomC{$\varphi$}
\RightLabel{\Nec}
\UnaryInfC{$\Box \varphi$}
\end{prooftree}
सूत्र~$\psi$ हे सूत्रे $\varphi$
पासून अनिवार्यीकरणाने निष्पन्न होते तेव्हा आणि
केवळ तेव्हाच, जेव्हा $\psi \ident \Box \varphi$.
\end{defn}

\begin{defn}
स्वयंसिद्धकांच्या संच~$\Sigma$ पासूनची
\emph{निष्पत्ती} म्हणजे सूत्रांची
$\psi_1$, $\psi_2$, \dots, $\psi_n$ अशी मालिका,
ज्यात प्रत्येक $\psi_i$ पुढीलपैकी एक असते:
\begin{enumerate}
\item सर्वतः-सत्य सूत्राचा आदेशन-दृष्टांत; किंवा
\item $\Sigma$ मधील सूत्राचा आदेशन-दृष्टांत; किंवा
\item $j$, $k < i$ असलेल्या दोन सूत्रे
  $\psi_j$, $\psi_k$ यांपासून विध्यात्मक अनुमानाने
  निष्पन्न होणारे; किंवा
\item $j < i$ असलेल्या सूत्र~$\psi_j$
  पासून अनिवार्यीकरणाने निष्पन्न होणारे.
\end{enumerate}
$\psi_n \ident \varphi$ असलेली अशी निष्पत्ती
असल्यास $\varphi$ हे $\Sigma$ पासून
\emph{निष्पन्न करता येण्याजोगे आहे}, असे म्हणतो;
संकेतात $\Sigma \Proves \varphi$.
\end{defn}

वरील परिच्छेदाप्रमाणे $\Sigma$ मध्ये
$\Ax{K}$ आणि~$\Ax{Dual}$ घेतल्यावर,
या व्याख्येनुसार निष्पन्न करता येण्याजोगे सूत्रांचा संच सामान्य मोडल
तर्कशास्त्र बनतो, हे पुढे दिसेल. तसेच, $\Sigma$ मधील
स्वयंसिद्धकांचे सर्व आदेशन-दृष्टांत सर्व जगांत सत्य असलेल्या
प्रत्येक प्रतिमानात प्रत्येक निष्पन्न करता येण्याजोगे सूत्र सत्य असते.
निष्पत्त्यांच्या या गुणधर्माला \emph{निर्दोषता} म्हणतात.
त्याचा व्यत्यास, \emph{संपूर्णता}, सिद्ध करणे अधिक कठीण आहे.












\section{सामान्य मोडल तर्कशास्त्रे}\label{nml:prf:nor:sec}

मोडल सूत्रांचा प्रत्येक संच स्वयंसिद्धकांच्या
एखाद्या संचापासून सिद्ध करता येणाऱ्या
सूत्रांचा संच म्हणून सहज वर्णन करता येत
नाही. मोडल तर्कशास्त्रे सुस्थित असावीत अशी
आपली अपेक्षा आहे. सर्वप्रथम, अभिजात विधानीय
तर्कशास्त्रात आपण निष्पन्न करू शकतो ते सर्व
अजूनही निष्पन्न करता येण्याजोगे असावे; आता सूत्रे
मध्ये $\Box$
    आणि~$\Diamond$ देखील
येऊ शकतात, हे अर्थातच लक्षात घेतले पाहिजे.
यासाठी मोडल तर्कशास्त्रात सर्व सर्वतः-सत्य दृष्टांत
असावेत आणि ते विध्यात्मक अनुमानाखाली बंद असावे,
अशी अट घालतो.

\begin{defn}
  \emph{मोडल तर्कशास्त्र} म्हणजे मोडल सूत्रे
  चा असा संच~$\Sigma$ की, तो
  \begin{enumerate}
  \item सर्व सर्वतः-सत्य सूत्रे सामावतो; आणि
  \item आदेशनाखाली बंद असतो; म्हणजे $\varphi \in \Sigma$
    आणि $\theta_1$, \dots, $\theta_n$ ही सूत्रे
    असतील, तर
    \[    \SSubst{\varphi}{\subst{\theta_1}{p_1}, \dots, \subst{\theta_n}{p_n}} \in \Sigma,
    \]
    \item \emph{विध्यात्मक अनुमानाखाली} बंद असतो;
    म्हणजे $\varphi$ आणि $\varphi
      \lif \psi \in \Sigma$ असतील, तर $\psi \in \Sigma$.
  \end{enumerate}
\end{defn}

मोडल तर्कशास्त्रांसाठी संबंधी चिन्हार्थमीमांसा
वापरायची असेल, तर सर्व मोडल प्रतिमानांत वैध
असलेली सर्व सूत्रे त्यांत समाविष्ट असावीत,
अशीही अट हवी. \Ax{K} आणि~\Dual{}
चे सर्व दृष्टांत निष्पन्न करता येण्याजोगे असतील, आणि
सूत्र~$\varphi$ निष्पन्न करता येण्याजोगे असेल तेव्हा
$\Box \varphi$ देखील तसेच असेल, तर ही अट पूर्ण होते.
या अटी पाळणाऱ्या मोडल तर्कशास्त्राला
\emph{सामान्य} म्हणतात. (असामान्य मोडल तर्कशास्त्रेही
आहेत; पण नेहमीची संबंधी प्रतिमाने त्यांच्यासाठी
पुरेशी नाहीत.)

\begin{defn}
  मोडल तर्कशास्त्र $\Sigma$ पुढील सूत्रे सामावत असेल,
  \begin{align*}
    \tag{\Ax{K}} & \Box(p \lif q) \lif (\Box p \lif \Box q),
    \\
      \tag{\Dual} & \Diamond p \liff \lnot\Box\lnot p
  \end{align*}
  आणि \emph{अनिवार्यीकरणाखाली} बंद असेल, म्हणजे
  $\varphi \in
  \Sigma$ असेल, तर $\Box \varphi \in \Sigma$,
  तर त्याला \emph{सामान्य} म्हणतात.
\end{defn}

सर्वतः-सत्य अभिव्यंजन \emph{जगातील सत्यता}
टिकवण्यासाठी पुरेसे सूक्ष्म असते; पण \Nec{} हा
नियम फक्त \emph{प्रतिमानातील सत्यता} टिकवतो
(आणि म्हणून चौकटीतील किंवा चौकटींच्या वर्गातील
वैधताही टिकवतो), हे लक्षात घ्या.

\begin{prop}\label{nml:prf:nor:prop:rk}
  प्रत्येक सामान्य मोडल तर्कशास्त्र \RK
  या नियमाखाली बंद असते:
  \begin{prooftree}
    \AxiomC{$\varphi_1 \lif (\varphi_2 \lif \cdots (\varphi_{n-1} \lif \varphi_n)\cdots)$}
    \RightLabel{\RK}
    \UnaryInfC{$\Box\varphi_1 \lif (\Box\varphi_2 \lif \cdots (\Box\varphi_{n-1}
      \lif \Box\varphi_n)\cdots).$}
  \end{prooftree}
\end{prop}

\begin{proof}
  $n$ वर विगमनाने सिद्ध करू. $n = 1$ असेल,
  तर हा नियम नेमका \Nec असतो, आणि प्रत्येक
  सामान्य मोडल तर्कशास्त्र \Nec खाली बंद असते.

  आता निष्कर्ष $n-1$ साठी खरा आहे असे समजा;
  तो~$n$ साठी सिद्ध करू.

  पुढील गृहीत धरा:
  \begin{align*}
  & \varphi_1 \lif (\varphi_2 \lif \cdots (\varphi_{n-1} \lif \varphi_n)\cdots) \in \Sigma
  \intertext{विगमन गृहीतकावरून मिळते:}
  & \Box\varphi_1 \lif (\Box\varphi_2 \lif \cdots \Box(\varphi_{n-1} \lif \varphi_n)\cdots)
  \in \Sigma
  \intertext{$\Sigma$ सामान्य मोडल तर्कशास्त्र असल्यामुळे
    त्यात~$\Ax{K}$ चे सर्व दृष्टांत आहेत; विशेषतः:}
  & \Box(\varphi_{n-1} \lif \varphi_n) \lif (\Box\varphi_{n-1} \lif \Box\varphi_n) \in \Sigma
  \intertext{विध्यात्मक अनुमान आणि योग्य सर्वतः-सत्य दृष्टांत
    वापरून मिळते:}
  & \Box\varphi_1 \lif (\Box\varphi_2 \lif \cdots (\Box\varphi_{n-1}
  \lif \Box\varphi_n)\cdots) \in \Sigma.
  \end{align*}
\end{proof}

\begin{prop}\label{nml:prf:nor:prop:notDiamondBot}
  प्रत्येक सामान्य मोडल तर्कशास्त्र $\Sigma$
  हे~$\lnot\Diamond\lfalse$ सामावते.
\end{prop}

\begin{prob}
  \ref{nml:prf:nor:prop:notDiamondBot} सिद्ध करा.
\end{prob}

\begin{prop}
  $\varphi_1$, \dots, $\varphi_n$ ही सूत्रे असू द्या.
  मग $\varphi_1$, \dots, $\varphi_n$ यांचे सर्व दृष्टांत
  सामावणारे लघुतम सामान्य मोडल तर्कशास्त्र $\Sigma$
  अस्तित्वात असते.
\end{prop}

\begin{proof}
  $\varphi_1$, \dots, $\varphi_n$ दिल्यावर
  $\varphi_1$, \dots, $\varphi_n$ यांचे सर्व दृष्टांत सामावणाऱ्या
  सर्व सामान्य मोडल तर्कशास्त्रांचा
  छेद~$\Sigma$ म्हणून ठरवा. अशा तर्कशास्त्रांचा
  वर्ग रिकामा नाही: सर्व सूत्रांचा संच
  $\Frm[L]$ हेही असे सामान्य मोडल तर्कशास्त्र आहे.
\end{proof}

\begin{defn}
$\varphi_1$, \dots, $\varphi_n$ सामावणाऱ्या लघुतम
सामान्य मोडल तर्कशास्त्राला \emph{मोडल प्रणाली}
म्हणतात; ते $\Log{K} \varphi_1 \dots
\varphi_n$ असे दर्शवतात. लघुतम सामान्य मोडल
तर्कशास्त्र~\Log{K} असे दर्शवतात.
\end{defn}












\section{निष्पत्ती आणि मोडल प्रणाली}\label{nml:prf:prf:sec}

सामान्य मोडल तर्कशास्त्रांसाठी निष्पत्ती
म्हणजे काय ते आधी ठरवू. साधारणपणे,
निष्पत्ती म्हणजे सूत्रांची अशी मालिका
की तिच्यातील प्रत्येक घटक काही
\emph{स्वयंसिद्धकांपैकी} एक (किंवा त्याचा आदेशन-दृष्टांत)
असतो, किंवा आधीच्या घटक पासून मोजक्या
अनुमान नियमांपैकी एखाद्याने मिळतो. सामान्य मोडल
तर्कशास्त्रात सर्व सर्वतः-सत्य सूत्रांचे दृष्टांत
, \Ax{K}, आणि~\Dual{}
स्वयंसिद्धके मानली जातात. त्यातून लघुतम सामान्य
मोडल तर्कशास्त्र, म्हणजे मोडल प्रणाली~$\Log{K}$,
मिळते. इतर प्रणाली मिळवण्यासाठी आणखी स्वयंसिद्धके
जोडता येतात. नियम मात्र नेहमी विध्यात्मक अनुमान~\MP{}
आणि अनिवार्यीकरण~\Nec हेच असतात.

\begin{defn}
  मोडल प्रणाली $\Log{K} \varphi_1 \dots \varphi_n$ आणि
  सूत्र~$\psi$ दिले असता, $\chi_1$, \dots,
  $\chi_k$ अशी सूत्रे असतील, जिथे $\chi_k = \psi$
  आणि प्रत्येक $\chi_i$ हा सर्वतः-सत्य दृष्टांत,
  किंवा $\Ax{K}$, $\Dual$,
  $\varphi_1$, \dots, $\varphi_n$ यांपैकी एखाद्याचा दृष्टांत,
  किंवा आधीच्या सूत्रांपासून \MP{} अथवा~\Nec
  नियमाने मिळणारे सूत्र असेल, तर आणि तेव्हाच
  $\psi$ हे $\Log{K} \varphi_1 \dots
  \varphi_n$ मध्ये \emph{निष्पन्न करता येण्याजोगे} आहे असे म्हणतो;
  संकेतात $\Log{K} \varphi_1 \dots \varphi_n \Proves \psi$.
\end{defn}

पुढील प्रतिज्ञेनुसार $\psi$ ची $\Sigma$-निष्पत्ती
दाखवून $\psi \in \Sigma$ हे सिद्ध करता येते.

\begin{prop}
  $\Log{K} \varphi_1 \dots \varphi_n = \Setabs{\psi}{\Log{K} \varphi_1 \dots \varphi_n \Proves \psi}$.
\end{prop}

\begin{proof}
  निष्पत्त्यांच्या लांबीवर विगमन करून
  $\Setabs{\psi}{\Log{K} \varphi_1 \dots \varphi_n \Proves \psi} \subseteq
  \Log{K} \varphi_1 \dots \varphi_n$ हे दाखवू.

  $\psi$ च्या निष्पत्तीची लांबी~$1$ असेल,
  तर त्यात एकच सूत्र असते. ते सूत्र आधीच्या
  सूत्रांपासून \MP{} किंवा \Nec ने मिळालेले
  असू शकत नाही. म्हणून ते सर्वतः-सत्य दृष्टांत,
  \Ax{K} चा, $\Dual$ चा,
  किंवा $\varphi_1$, \dots,~$\varphi_n$ यांपैकी एखाद्याचा
  दृष्टांत असला पाहिजे. $\Log{K}\varphi_1\dots\varphi_n$
  हे सर्व सामावते; म्हणून $\psi \in \Log{K}\varphi_1
  \dots \varphi_n$.

  $\psi$ च्या निष्पत्तीची लांबी~$> 1$
  असेल, तर आधीच्या प्रकरणांखेरीज निष्पत्तीमधील अंतिम ओळ नसलेल्या सूत्रांपासून
  \MP{} किंवा \Nec{} नेही
  $\psi$ मिळू शकते. $\psi$ हे $\chi$ आणि
  $\chi \lif \psi$ यांपासून (\MP ने) मिळत असेल,
  तर विगमन गृहीतकावरून $\chi$ आणि
  $\chi \lif \psi \in
  \Log{K}\varphi_1 \dots \varphi_n$. प्रत्येक मोडल
  तर्कशास्त्र विध्यात्मक अनुमानाखाली बंद असते;
  म्हणून $\psi \in \Log{K}\varphi_1 \dots
  \varphi_n$. $\psi \equiv \Box \chi$ हे $\chi$
  पासून \Nec ने मिळत असेल, तर विगमन
  गृहीतकावरून $\chi
  \in \Log{K}\varphi_1 \dots \varphi_n$. प्रत्येक
  सामान्य मोडल तर्कशास्त्र $\Nec$ खाली बंद असते;
  म्हणून $\psi \in
  \Log{K}\varphi_1\dots\varphi_n$.

  उलट समावेशासाठी $\Sigma = \Setabs{\psi}{\Log{K} \varphi_1 \dots \varphi_n \Proves \psi }$
  हे $\varphi_1$, \dots, $\varphi_n$ चे सर्व दृष्टांत
  सामावणारे सामान्य मोडल तर्कशास्त्र आहे, हे
  दाखवू. मग $\Log{K} \varphi_1 \dots \varphi_n$
  हे व्याख्येने असे लघुतम तर्कशास्त्र आहे,
  या तथ्याचा उपयोग करू.
  \begin{enumerate}
    \item प्रत्येक सर्वतः-सत्य सूत्र~$\psi$ हाही
      सर्वतः-सत्य दृष्टांत आहे. म्हणून
      $\Log{K}\varphi_1\dots\varphi_n \Proves \psi$;
      त्यामुळे $\Sigma$ सर्व सर्वतः-सत्य सूत्रे सामावते.
    \item $\Log{K}\varphi_1\dots\varphi_n \Proves \chi$
      आणि $\Log{K}\varphi_1\dots\varphi_n \Proves \chi \lif \psi$
      असल्यास $\Log{K}\varphi_1\dots\varphi_n \Proves \psi$:
      $\chi$ ची निष्पत्ती आणि $\chi \lif \psi$
      ची निष्पत्ती जोडून अंतिम ओळ~$\psi$
      घाला. या ओळीला \MP{} चे समर्थन मिळते.
      म्हणून $\Sigma$ विध्यात्मक अनुमानाखाली बंद आहे.
    \item $\psi$ ची निष्पत्ती असेल, तर
      $\psi$ च्या प्रत्येक आदेशन-दृष्टांताचीही
      निष्पत्ती असते: संपूर्ण निष्पत्तीमधील प्रत्येक सूत्र वर तेच आदेशन लावा.
      (सराव: निष्पत्त्यांच्या लांबीवर विगमन
      करून परिणामी मालिकाही योग्य निष्पत्ती
      आहे हे सिद्ध करा.) म्हणून $\Sigma$ एकरूप
      आदेशनाखाली बंद आहे. ($\Sigma$ मोडल तर्कशास्त्राच्या
      सर्व अटी पूर्ण करते, हे आता सिद्ध झाले.)
    \item $\Log{K}\varphi_1\dots\varphi_n \Proves \Ax{K}$;
      म्हणून $\Ax{K} \in \Sigma$.
    \item $\Log{K}\varphi_1\dots\varphi_n \Proves \Dual$
      असल्यामुळे $\Dual \in \Sigma$.
    \item $\Log{K}\varphi_1\dots\varphi_n \Proves \chi$
      असेल, तर अतिरिक्त ओळ~$\Box \chi$ ला~\Nec
      चे समर्थन मिळते. त्यामुळे $\Sigma$ हा
      \Nec खाली बंद आहे; म्हणजे $\Sigma$ सामान्य आहे.
  \end{enumerate}
\end{proof}













\section{\Log{K} मधील सिद्धता}\label{nml:prf:prk:sec}

लघुतम मोडल प्रणालीतील सिद्धतेचा सराव म्हणून
\ref{nml:syn:sch:tab:valid-invalidSchemas} च्या
डाव्या बाजूची वैध सूत्रे सर्व
\Log{K} मध्ये सिद्ध करता येतात, हे दाखवू.

\begin{prop}
  $\Log{K} \Proves \Box\varphi \lif \Box (\psi\lif \varphi)$
\end{prop}

\begin{proof}
  \begin{derivation}
    1. & $\varphi \lif (\psi \lif \varphi)$ & \Taut \\
    2. & $\Box(\varphi \lif (\psi \lif \varphi))$ & \Nec, 1 \\
    3. & $\Box(\varphi \lif (\psi \lif \varphi)) \lif
    (\Box\varphi \lif \Box(\psi \lif \varphi))$ & \Ax{K}\\
    4. & $\Box\varphi \lif \Box (\psi\lif \varphi)$ & \MP, 2, 3
  \end{derivation}
\end{proof}

\begin{prop}
  $\Log{K} \Proves \Box(\varphi \land \psi) \lif (\Box \varphi \land \Box\psi)$
\end{prop}

\begin{proof}
\begin{derivation}
  1. & $(\varphi \land \psi) \lif \varphi$ & \Taut \\
  2. & $\Box((\varphi \land \psi) \lif \varphi)$ & \Nec \\
  3. & $\Box((\varphi \land \psi) \lif \varphi) \lif (\Box(\varphi \land \psi) \lif \Box\varphi)$
  & \Ax{K} \\
  4. & $\Box(\varphi \land \psi) \lif \Box\varphi$ & \MP, 2, 3 \\
  5. & $(\varphi \land \psi) \lif \psi$ & \Taut \\
  6. & $\Box((\varphi \land \psi) \lif \psi)$ & \Nec \\
  7. & $\Box((\varphi \land \psi) \lif \psi) \lif (\Box(\varphi \land \psi) \lif \Box\psi)$
  & \Ax{K} \\
  8. & $\Box(\varphi \land \psi) \lif \Box\psi$ & \MP, 6, 7 \\
  9. & $(\Box(\varphi \land \psi) \lif \Box\varphi) \lif{}$ \\
  & \qquad $((\Box(\varphi \land \psi) \lif \Box\psi) \lif{}$ \\
  & \qquad $(\Box(\varphi \land \psi) \lif (\Box \varphi \land \Box\psi)))$ & \Taut\\
  10. &  $(\Box(\varphi \land \psi) \lif \Box\psi) \lif{}$ \\
  & \qquad $(\Box(\varphi \land \psi) \lif (\Box \varphi \land \Box\psi))$ & \MP, 4, 9\\
  11. & $\Box(\varphi \land \psi) \lif (\Box \varphi \land \Box\psi)$ & \MP, 8, 10.
\end{derivation}
ओळ~$9$ वरील सूत्र पुढील सर्वतः-सत्य सूत्राचा दृष्टांत आहे:
\[(p \lif q) \lif ((p \lif r) \lif (p \lif (q \land r))).
\]
\end{proof}

\begin{prop}
  $\Log{K}\Proves (\Box\varphi \land \Box\psi) \lif \Box (\varphi \land \psi)$
\end{prop}

\begin{proof}
  \begin{derivation}
    1. & $\varphi \lif (\psi \lif (\varphi \land \psi))$ & \Taut \\
    2. & $\Box(\varphi \lif (\psi \lif (\varphi \land \psi)))$ & \Nec, 1 \\
    3. & $\Box(\varphi \lif (\psi \lif (\varphi \land \psi))) \lif (\Box\varphi \lif \Box(\psi \lif (\varphi \land \psi)))$ & \Ax{K}\\
    4. & $\Box\varphi \lif \Box(\psi \lif (\varphi \land \psi))$ & \MP, 2, 3\\
    5. & $\Box(\psi \lif (\varphi \land \psi)) \lif (\Box\psi \lif \Box(\varphi \land \psi))$ & \Ax{K} \\
    6. & $(\Box\varphi \lif \Box(\psi \lif (\varphi \land \psi))) \lif {}$ \\
    & \qquad $((\Box(\psi \lif (\varphi \land \psi)) \lif (\Box\psi \lif \Box(\varphi \land \psi))) \lif {}$\\
    & \qquad $(\Box\varphi \lif (\Box \psi \lif \Box(\varphi \land \psi))))$ & \Taut\\
    7. & $(\Box(\psi \lif (\varphi \land \psi)) \lif (\Box\psi \lif \Box(\varphi \land \psi))) \lif {}$\\
    & \qquad $(\Box\varphi \lif (\Box \psi \lif \Box(\varphi \land \psi)))$ & \MP, 4, 6\\
    8. & $\Box\varphi \lif (\Box \psi \lif \Box(\varphi \land \psi))$ & \MP, 5, 7\\
    9. & $(\Box\varphi \lif (\Box \psi \lif \Box(\varphi \land \psi))) \lif {}$\\
    & \qquad $((\Box\varphi \land \Box\psi) \lif \Box(\varphi \land \psi))$ & \Taut\\
    10. & $(\Box\varphi \land \Box\psi) \lif \Box(\varphi \land \psi)$ & \MP, 8, 9
  \end{derivation}
  ओळी $6$ आणि $9$ वरील सूत्रे अनुक्रमे पुढील
  सर्वतः-सत्य सूत्रांचे दृष्टांत आहेत:
  \begin{align*}
    (p \lif q) & \lif ((q \lif r) \lif (p \lif r)) \\
    (p \lif (q \lif r)) & \lif ((p \land q) \lif r)
  \end{align*}
\end{proof}

\begin{prop}
  $\Log{K} \Proves \lnot\Box p \lif \Diamond \lnot
    p$
\end{prop}

\begin{proof}

  \begin{derivation}
    1. & $\Diamond \lnot p \liff \lnot\Box\lnot\lnot p$ & \Dual\\
    2. & $(\Diamond \lnot p \liff \lnot\Box\lnot\lnot p) \lif {}$ \\
    & \qquad $(\lnot \Box \lnot\lnot p \lif \Diamond \lnot p)$& \Taut\\
    3. & $\lnot \Box \lnot\lnot p \lif \Diamond \lnot p$ & \MP, 1, 2\\
    4. & $\lnot\lnot p \lif p$ & \Taut\\
    5. & $\Box(\lnot\lnot p \lif p)$ & \Nec, 4\\
    6. & $\Box(\lnot\lnot p \lif p) \lif (\Box \lnot\lnot p \lif \Box p)$ & \Ax{K}\\
    7. & $(\Box \lnot\lnot p \lif \Box p)$ & \MP, 5, 6\\
    8. & $(\Box \lnot\lnot p \lif \Box p) \lif (\lnot \Box p \lif \lnot\Box\lnot\lnot p)$ & \Taut\\
    9. & $\lnot \Box p \lif \lnot\Box\lnot\lnot p$ & \MP, 7, 8\\
    10. & $(\lnot \Box p \lif \lnot\Box\lnot\lnot p) \lif {}$ \\
    & \qquad $((\lnot \Box \lnot\lnot p \lif \Diamond \lnot p) \lif (\lnot \Box p \lif \Diamond\lnot p))$ & \Taut\\
    11. & $(\lnot \Box \lnot\lnot p \lif \Diamond \lnot p) \lif (\lnot \Box p \lif \Diamond\lnot p)$ & \MP, 9, 10\\
    12. & $\lnot\Box p \lif \Diamond \lnot p$ & \MP, 3, 11\\
  \end{derivation}
  ओळी $8$ आणि $10$ वरील सूत्रे अनुक्रमे पुढील
  सर्वतः-सत्य सूत्रांचे दृष्टांत आहेत:
  \begin{align*}
    & (p \lif q) \lif (\lnot q \lif \lnot p) \\
    & (p \lif q) \lif ((q \lif r) \lif (p \lif r)).
  \end{align*}

\end{proof}

\begin{prob}
  पुढील सूत्रे साठी $\Log{K}$ मध्ये
  निष्पत्त्या शोधा:
  \begin{enumerate}
    \item $\Box \lnot p \lif \Box(p \lif q)$
    \item $(\Box p \lor \Box q) \lif \Box(p \lor q)$
    \item $\Diamond p \lif \Diamond(p \lor q)$
  \end{enumerate}
\end{prob}













\section{व्युत्पन्न नियम}\label{nml:prf:der:sec}

निष्पत्त्या शोधणे आणि लिहिणे कठीण, किचकट
आणि पुनरावृत्तीचे काम आहे. उदाहरणार्थ, $\varphi \lif \psi$
पासून $\Box \varphi \lif \Box \psi$ कडे जायचे असेल,
म्हणजे नियम~\RK लावायचा असेल, तर आधी \Nec लावावा लागतो,
मग~\Ax{K} चा योग्य दृष्टांत लिहावा लागतो आणि नंतर~\MP
लावावा लागतो. तसेच $\varphi \lif \psi$ आणि $\psi \lif \chi$
पासून $\varphi \lif \chi$ कडे जाण्यासाठी पुढील (लांब)
सर्वतः-सत्य दृष्टांत लिहावा लागतो:
\[(\varphi \lif \psi) \lif ((\psi \lif \chi) \lif (\varphi \lif \chi))
\]
आणि \MP{} दोनदा लावावा लागतो. अनेकदा एखादे उप-सूत्र
त्याच्याशी समतुल्य असलेल्या सूत्राने बदलायचे असते;
उदाहरणार्थ, $\Diamond
  \varphi$ च्या जागी $\lnot\Box\lnot \varphi$, किंवा
$\lnot\lnot \varphi$ च्या जागी $\varphi$. म्हणून संपूर्ण
निष्पत्ती लिहिण्याऐवजी मध्यवर्ती पायऱ्या
निष्पन्न करता येण्याजोगे का आहेत, एवढेच नोंदवणे सोयीचे आहे.
त्यासाठी निष्पन्नता विषयी काही तथ्ये गोळा करू.

\begin{prop}
  $\Log{K} \Proves \varphi_1$, \dots, $\Log{K} \Proves \varphi_n$
  असेल आणि विधानीय तर्कशास्त्रानुसार $\varphi_1$, \dots,
  $\varphi_n$ पासून $\psi$ निष्पन्न होत असेल, तर
  $\Log{K} \Proves \psi$.
\end{prop}

\begin{proof}
  विधानीय तर्कशास्त्रानुसार $\varphi_1$, \dots,~$\varphi_n$
  पासून $\psi$ निष्पन्न होत असेल, तर
  \[  \varphi_1 \lif (\varphi_2 \lif \cdots (\varphi_n \lif \psi)\dots)
  \]
  हा सर्वतः-सत्य दृष्टांत असतो. \MP{} $n$ वेळा
  लावल्याने $\psi$ ची निष्पत्ती मिळते.
\end{proof}

या प्रतिज्ञेचा वापर~\PL ने दर्शवू.

\begin{prop}
  $\Log{K} \Proves \varphi_1 \lif (\varphi_2 \lif \cdots (\varphi_{n-1} \lif
  \varphi_n)\dots)$ असेल, तर $\Log{K} \Proves \Box \varphi_1 \lif (\Box \varphi_2 \lif
  \cdots (\Box \varphi_{n-1} \lif \Box \varphi_n)\dots)$.
\end{prop}

\begin{proof}
  \ref{nml:prf:nor:prop:rk} च्या सिद्धतेप्रमाणेच
  $n$ वर विगमन करा.
\end{proof}

या प्रतिज्ञेचा वापर~\RK ने दर्शवू. या परिणामांमुळे
निष्पन्नता विषयीची विधाने किती सोप्या
रीतीने सिद्ध करता येतात, ते पाहू.

\begin{prop}
  $\Log{K} \Proves (\Box\varphi \land \Box\psi) \lif \Box (\varphi \land \psi)$
\end{prop}

\begin{proof}
  \begin{derivation}
    1. & $\Log{K} \Proves \varphi \lif (\psi \lif (\varphi \land \psi))$ & \Taut \\
    2. & $\Log{K} \Proves \Box\varphi \lif (\Box \psi \lif \Box(\varphi \land \psi))$ & \RK, 1\\
    3. & $\Log{K} \Proves (\Box\varphi \land \Box\psi) \lif \Box(\varphi \land \psi)$ & \PL, 2
  \end{derivation}
\end{proof}

\begin{prop}\label{nml:prf:der:prop:rewriting}
  $\Log{K} \Proves \varphi \liff \psi$ आणि $\Log{K} \Proves
  \Subst{\chi}{\varphi}{q}$ असेल, तर $\Log{K} \Proves \Subst{\chi}{\psi}{q}$.
\end{prop}

\begin{proof}
  सराव.
\end{proof}

\begin{prob}
  $\chi$ च्या जटिलतेवर विगमन करून
  \ref{nml:prf:der:prop:rewriting} सिद्ध करा:
  $\Log{K} \Proves \varphi \liff \psi$ असेल, तर
  $\Log{K} \Proves \Subst{\chi}{\varphi}{q} \liff \Subst{\chi}{\psi}{q}$
  हे दाखवा.
\end{prob}

$\Diamond$ चे $\Box$ मध्ये (किंवा उलट) रूपांतर करायचे
असेल, किंवा सूत्र मधील दुहेरी निषेध काढायचा
असेल, तेव्हा ही प्रतिज्ञा विशेष उपयोगी ठरते.
पुढे $\chi(\varphi)$ च्या जागी सूत्र~$\chi(\psi)$ लिहिताना
\ref{nml:prf:der:prop:rewriting} चा वापर ``$\psi$ for $\varphi$''
असा दाखवू. दुसऱ्या शब्दांत, ``$\psi$ for $\varphi$''
याचा संक्षिप्त अर्थ असा:
  \begin{derivation}
    & $\Proves \chi(\varphi)$ \\
    & $\Proves \varphi \liff \psi$\\
    & $\Proves \chi(\psi)$ & \ref{nml:prf:der:prop:rewriting} वरून
  \end{derivation}
उदाहरणार्थ:

\begin{prop}
  $\Log{K} \Proves \lnot\Box p \lif \Diamond \lnot p$
\end{prop}

\begin{proof}

  \begin{derivation}
    1. & $\Log{K} \Proves \Diamond \lnot p \liff \lnot\Box\lnot\lnot p$ &
    \Dual\\
    2. & $\Log{K} \Proves \lnot \Box \lnot\lnot p \lif \Diamond \lnot p$ & \PL, 1\\
    3. & $\Log{K} \Proves \lnot\Box p \lif \Diamond \lnot p$ & $p$ for $\lnot\lnot p$\\
  \end{derivation}

\end{proof}

वरील निष्पत्तीमधील अंतिम पायरी
``$p$ for $\lnot\lnot p$'' ही पुढील पायऱ्यांचे
संक्षिप्त रूप आहे:
  \begin{derivation}
    & $\Log{K} \Proves \lnot \Box \lnot\lnot p \lif \Diamond \lnot
    p$\\
    & $\Log{K} \Proves \lnot\lnot p \liff p$ & \Taut\\
    & $\Log{K} \Proves \lnot\Box p \lif \Diamond \lnot p$ & \ref{nml:prf:der:prop:rewriting} वरून
  \end{derivation}
\ref{nml:prf:der:prop:rewriting} मधील $\chi(q)$, $\varphi$ आणि $\psi$
यांच्या जागी येथे अनुक्रमे
$\lnot \Box q \lif \Diamond \lnot p$,
$\lnot\lnot p$ आणि~$p$ आहेत.

सूत्रामध्ये $\lnot\Diamond \varphi$ हे उप-सूत्र
असेल, तर \ref{nml:prf:der:prop:rewriting} वापरून त्याच्या जागी
$\Box\lnot \varphi$ लिहिता येते, कारण
$\Log{K} \Proves \lnot\Diamond \varphi \liff \Box\lnot \varphi$.
हा आणि यांसारखे बदल थोडक्यात
``$\Box\lnot$ for $\lnot\Diamond$'' असे दर्शवू.

पुढील प्रतिज्ञेमुळे निष्पन्नता विषयीचे परिणाम
रूपबंधात्मक पद्धतीने सिद्ध करता येतात. उदाहरणार्थ,
आधीची प्रतिज्ञा केवळ $\Log{K} \Proves \lnot\Box p \lif
\Diamond \lnot p$ एवढेच दाखवत नाही; कोणत्याही~$\varphi$
साठी $\Log{K} \Proves \lnot\Box \varphi \lif \Diamond
\lnot \varphi$ हेही दाखवते.

\begin{prop}
  $\varphi$ हा $\psi$ चा आदेशन-दृष्टांत असेल आणि
  $\Log{K} \Proves \psi$ असेल, तर $\Log{K} \Proves \varphi$.
\end{prop}

\begin{proof}
  संपूर्ण निष्पत्ती वर आदेशन लावल्यानेही योग्य
  निष्पत्ती मिळते, हे तपासणे किचकट पण नेहमीचे
  काम आहे ($\psi$ च्या निष्पत्तीच्या लांबीवर
  विगमन करा). विशेषतः, सर्वतः-सत्य दृष्टांतांचे
  आदेशन-दृष्टांतही सर्वतः-सत्य दृष्टांत असतात;
  \Dual{} आणि~\Ax{K}
  यांच्या दृष्टांतांचे आदेशन-दृष्टांतही
  \Dual{} आणि~\Ax{K}
  यांचेच दृष्टांत असतात; आणि आधारसूत्रांमध्ये
  व निष्कर्षात विधानीय चलांच्या
  जागी सूत्रे ठेवले तरी \MP{} व~\Nec{}
  यांचा वापर योग्यच राहतो.
\end{proof}













\section{\Log{K} मधील आणखी सिद्धता}\label{nml:prf:mpr:sec}

\ref{nml:prf:der:sec} मध्ये दिलेली सुलभ पद्धत वापरून
\Log{K} मधील निष्पन्नता ची आणखी काही उदाहरणे पाहू.

\begin{prop}
  $\Log{K} \Proves \Box (\varphi \lif \psi) \lif (\Diamond \varphi \lif \Diamond
  \psi)$
\end{prop}

\begin{proof}
\begin{derivation}
1. & $\Log{K} \Proves (\varphi \lif \psi) \lif (\lnot \psi \lif \lnot \varphi)$ & \PL \\
2. &  $\Log{K} \Proves \Box(\varphi \lif \psi) \lif (\Box\lnot\psi \lif \Box\lnot\varphi)$ & \RK, 1 \\
3. & $\Log{K} \Proves (\Box\lnot\psi \lif \Box\lnot\varphi) \lif (\lnot
\Box\lnot\varphi \lif \lnot\Box\lnot\psi)$ & \Taut\\
4. & $\Log{K} \Proves \Box(\varphi \lif \psi) \lif (\lnot
\Box\lnot\varphi \lif \lnot\Box\lnot\psi)$ & \PL, 2, 3 \\
5. &  $\Log{K} \Proves \Box(\varphi \lif \psi) \lif (\Diamond\varphi \lif \Diamond\psi)$ & $\Diamond$ for $\lnot\Box\lnot$.
\end{derivation}
\end{proof}

\begin{prop}
$\Log{K}\Proves \Box\varphi \lif (\Diamond(\varphi \lif \psi) \lif
  \Diamond \psi)$
\end{prop}

\begin{proof}
  \begin{derivation}
  1. & $\Log{K} \Proves \varphi \lif (\lnot\psi \lif \lnot (\varphi \lif \psi))$ & \Taut \\
  2. & $\Log{K} \Proves \Box\varphi \lif  (\Box\lnot\psi \lif \Box\lnot (\varphi\lif \psi))$ &
  \RK, 1 \\
  3. &  $\Log{K} \Proves \Box\varphi \lif (\lnot \Box\lnot (\varphi\lif \psi) \lif \lnot
  \Box\lnot\psi)$ & \PL, 2 \\
  4. & $\Log{K} \Proves \Box\varphi \lif (\Diamond(\varphi \lif \psi) \lif
  \Diamond \psi)$ & $\Diamond$ for $\lnot\Box\lnot$.
  \end{derivation}
\end{proof}

\begin{prop}
  $\Log{K} \Proves (\Diamond\varphi \lor \Diamond\psi) \lif \Diamond(\varphi \lor \psi)$
\end{prop}

\begin{proof}
  \begin{derivation}
    1. & $\Log{K} \Proves \lnot(\varphi \lor \psi) \lif \lnot\varphi$ & \Taut \\
    2. & $\Log{K} \Proves \Box\lnot(\varphi \lor \psi) \lif \Box\lnot\varphi$ & \RK, 1 \\
    3. & $\Log{K} \Proves \lnot\Box\lnot \varphi \lif \lnot\Box\lnot(\varphi \lor\psi)$ &
    \PL, 2\\
    4. & $\Log{K} \Proves \Diamond\varphi \lif \Diamond(\varphi \lor \psi)$ &
    $\Diamond$ for $\lnot\Box\lnot$\\
    5. & $\Log{K} \Proves \Diamond\psi \lif \Diamond(\varphi \lor \psi)$ & तसेच\\
    6. & $\Log{K} \Proves (\Diamond\varphi \lor\Diamond\psi) \lif \Diamond(\varphi \lor \psi)$
    & \PL, 4, 5.
  \end{derivation}
\end{proof}

\begin{prop}
  $\Log{K} \Proves \Diamond(\varphi \lor\psi) \lif (\Diamond\varphi \lor \Diamond\psi)$
\end{prop}

\begin{proof}
  \begin{derivation}
    1. & $\Log{K} \Proves \lnot \varphi \lif (\lnot \psi \lif \lnot (\varphi \lor \psi))$ & \Taut \\
    2. & $\Log{K} \Proves \Box\lnot \varphi \lif
    (\Box\lnot\psi \lif \Box \lnot (\varphi \lor \psi))$ & \RK\\
    3. & $\Log{K} \Proves \Box\lnot \varphi \lif (\lnot \Box \lnot (\varphi \lor\psi)
    \lif \lnot\Box\lnot \psi)$ & \PL, 2\\
    4. & $\Log{K} \Proves \lnot \Box \lnot(\varphi \lor \psi) \lif (\Box \lnot \varphi \lif
    \lnot\Box\lnot \psi)$ & \PL, 3\\
    5. & $\Log{K} \Proves \lnot \Box \lnot(\varphi\lor\psi) \lif (\lnot
    \lnot\Box\lnot\psi \lif \lnot\Box\lnot\varphi)$ & \PL, 4\\
    6. & $\Log{K} \Proves \Diamond(\varphi \lor \psi) \lif (\lnot
    \Diamond\psi \lif \Diamond\varphi)$ & $\Diamond$ for $\lnot\Box\lnot$\\
    7. & $\Log{K} \Proves \Diamond(\varphi\lor\psi) \lif (\Diamond\psi \lor \Diamond\varphi)$ & \PL, 6. \\
  \end{derivation}
\end{proof}

\begin{prob}
  पुढील निष्पन्नता विषयीची विधाने सिद्ध करा:
  \begin{enumerate}
  \item $\Log{K} \Proves \Diamond \lnot \lfalse \lif (\Box \varphi \lif
    \Diamond \varphi)$;
  \item $\Log{K} \Proves \Box(\varphi \lor \psi) \lif (\Diamond \varphi \lor \Box
    \psi)$;
  \item $\Log{K} \Proves (\Diamond \varphi \lif \Box \psi) \lif \Box(\varphi \lif
    \psi)$.
  \end{enumerate}
\end{prob}












\section{द्वयी सूत्र}\label{nml:prf:dua:sec}

\begin{defn}\label{nml:prf:dua:def:duals}
  \Ax{T}, \Ax{B}, \Ax{4} आणि \Ax{5} ही
  सूत्रे आहेत; त्यांपैकी प्रत्येकाला
  \emph{द्वयी} सूत्र आहे.
  ते तळसूचक हिऱ्याने पुढीलप्रमाणे दर्शवले जाते:
  \begin{align}
    \tag{\Ax{T_\Diamond}} p & \lif \Diamond p\\
    \tag{\Ax{B_\Diamond}} \Diamond\Box p & \lif p\\
    \tag{\Ax{4_\Diamond}} \Diamond\Diamond p & \lif \Diamond p\\
    \tag{\Ax{5_\Diamond}} \Diamond\Box p & \lif \Box p
    \end{align}
\end{defn}

वरील द्वयी सूत्रांपैकी प्रत्येक सूत्र संबंधित
सूत्र मधील $p$ च्या जागी $\lnot p$ ठेवून,
प्रतिपरिवर्तन करून,
$\lnot\Box\lnot$ च्या जागी $\Diamond$ आणि
$\lnot\Diamond\lnot$ च्या जागी~$\Box$ ठेवून मिळते.
\Ax{D}, म्हणजे $\Box\varphi \lif \Diamond\varphi$, हे या अर्थाने
स्वतःचेच द्वयी सूत्र आहे.

\begin{prop}\label{nml:prf:dua:prop:dualsys}
  \ref{nml:prf:dua:def:duals} मधील प्रत्येक
  सूत्र~$\varphi$ साठी:
  $\Log{K}\varphi = \Log{K}\varphi_{\Diamond}$.
\end{prop}
\begin{proof}
  सराव.
\end{proof}
\begin{prob}
  \ref{nml:prf:dua:prop:dualsys} सिद्ध करा.
\end{prob}












\section{मोडल प्रणालींमधील सिद्धता}\label{nml:prf:prs:sec}

आता~\Log{K} व्यतिरिक्त इतर मोडल तर्कशास्त्रीय
प्रणालींमधील सिद्धता पाहू.

\begin{prop}\label{nml:prf:prs:prop:S5facts}
पुढील सिद्धताविषयक परिणाम मिळतात:
  \begin{enumerate}
  \item $\Log{KT5} \Proves \Ax{B}$;
  \item $\Log{KT5} \Proves \Ax{4}$;
  \item $\Log{KDB4} \Proves \Ax{T}$;
  \item $\Log{KB4} \Proves \Ax{5}$;
  \item $\Log{KB5} \Proves \Ax{4}$;
  \item \label{nml:prf:prs:prop:S5facts-KT-D}$\Log{KT} \Proves \Ax{D}$.
  \end{enumerate}
\end{prop}

\begin{proof}
  प्रत्येक परिणामाची सिद्धता देऊ.
  \begin{enumerate}
    \item $\Log{KT5} \Proves \Ax{B}$:
      \begin{derivation}
        1. & $\Log{KT5} \Proves \Diamond\varphi \lif \Box\Diamond\varphi$ & \Ax{5}\\
        2. & $\Log{KT5} \Proves \varphi \lif \Diamond\varphi$ & $\Ax{T_\Diamond}$\\
        3. & $\Log{KT5} \Proves \varphi \lif \Box\Diamond\varphi$ & \PL, 2, 1.
      \end{derivation}
    \item $\Log{KT5} \Proves \Ax{4}$:
      \begin{derivation}
        1. &$\Log{KT5} \Proves \Diamond\Box\varphi \lif \Box\Diamond\Box\varphi$ & \Ax{5}
        मध्ये $p$ च्या जागी $\Box\varphi$\\
        2. & $\Log{KT5} \Proves \Box\varphi \lif \Diamond\Box\varphi$ & $\Ax{T_\Diamond}$
        मध्ये $p$ च्या जागी $\Box\varphi$\\
        3. & $\Log{KT5} \Proves \Box\varphi \lif \Box\Diamond\Box\varphi$ & \PL, 2, 1\\
        4. & $\Log{KT5} \Proves \Diamond\Box\varphi \lif \Box\varphi$ & $\Ax{5_\Diamond}$\\
        5. & $\Log{KT5} \Proves \Box\Diamond\Box\varphi \lif \Box\Box\varphi$ & \RK{}, 4 \\
        6. & $\Log{KT5} \Proves \Box\varphi \lif \Box\Box\varphi$ & \PL, 3, 5. \\
      \end{derivation}
    \item $\Log{KDB4} \Proves \Ax{T}$:
      \begin{derivation}
        1. & $\Log{KDB4} \Proves \Diamond\Box\varphi \lif \varphi $ & $\Ax{B_\Diamond}$ \\
        2. & $\Log{KDB4} \Proves \Box\Box\varphi \lif \Diamond\Box\varphi$ & $\Ax{D}$
        मध्ये $p$ च्या जागी $\Box\varphi$\\
        3. & $\Log{KDB4} \Proves \Box\Box\varphi \lif \varphi$ & \PL, 1, 2\\
        4. & $\Log{KDB4} \Proves \Box\varphi \lif \Box\Box\varphi$ & \Ax{4} \\
        5. & $\Log{KDB4} \Proves \Box\varphi \lif \varphi$ & \PL, 4, 3. \\
      \end{derivation}
    \item $\Log{KB4} \Proves \Ax{5}$:
      \begin{derivation}
        1. & $\Log{KB4} \Proves \Diamond\varphi \lif \Box \Diamond\Diamond\varphi$ & \Ax{B}
        मध्ये $p$ च्या जागी $\Diamond\varphi$\\
        2. & $\Log{KB4} \Proves \Diamond\Diamond \varphi \lif \Diamond \varphi$ &
        $\Ax{4_\Diamond}$ \\
        3. & $\Log{KB4} \Proves \Box\Diamond\Diamond \varphi \lif \Box \Diamond\varphi$ &
        \RK{}, 2 \\
        4. & $\Log{KB4} \Proves \Diamond\varphi \lif \Box\Diamond\varphi$ & \PL, 1, 3.
      \end{derivation}
    \item $\Log{KB5} \Proves \Ax{4}$:
      \begin{derivation}
        1. & $\Log{KB5} \Proves \Box\varphi \lif \Box\Diamond\Box \varphi$ & \Ax{B} मध्ये
        $p$ च्या जागी $\Box\varphi$ \\
        2. & $\Log{KB5} \Proves  \Diamond\Box\varphi \lif \Box \varphi$ &
        $\Ax{5_\Diamond}$\\
        3. & $\Log{KB5} \Proves  \Box\Diamond\Box\varphi \lif \Box\Box \varphi$ & \RK{}, 2 \\
        4. & $\Log{KB5} \Proves  \Box\varphi \lif \Box\Box\varphi$ & \PL, 1, 3.
      \end{derivation}
    \item $\Log{KT} \Proves \Ax{D}$:
      \begin{derivation}
        1. & $\Log{KT} \Proves \Box \varphi \lif \varphi$ & \Ax{T} \\
        2. & $\Log{KT} \Proves \varphi \lif \Diamond\varphi$ & $\Ax{T_\Diamond}$ \\
        3. & $\Log{KT} \Proves \Box \varphi \lif \Diamond \varphi$ &  \PL, 1, 2
      \end{derivation}
  \end{enumerate}
\end{proof}

\begin{defn}
  प्रचलित संकेतानुसार \Log{S4} ही प्रणाली
  \Log{KT4} आणि \Log{S5} ही प्रणाली
  \Log{KTB4} अशी परिभाषित करतो.
\end{defn}

पुढील प्रतिज्ञा दाखवते की अभिजात प्रणाली~\Log{S5}
अनेक समतुल्य स्वयंसिद्धक-प्रणालींनी मिळते. हे
आश्चर्यकारक नाही; स्वयंसिद्धकांचे हे वेगवेगळे संयोग
समतुल्यता संबंधांनाच लक्षणित करतात
(\ref{nml:frd:es5:prop:equivalences} पाहा).

\begin{prop}\label{nml:prf:prs:prop:S5}
  $\Log{KTB4} = \Log{KT5} = \Log{KDB4} = \Log{KDB5}$.
\end{prop}

\begin{proof}
  सराव.
\end{proof}

\begin{prob}
  \ref{nml:prf:prs:prop:S5} सिद्ध करा.
\end{prob}














\section{निर्दोषता}\label{nml:prf:snd:sec}

ज्या निष्पत्ती प्रणालीमध्ये निष्पन्न
करता येणारे प्रत्येक सूत्र वैध असते, तिला निर्दोष
म्हणतात. मोडल प्रणाली विचारात घेताना, म्हणजे
\Ax{K} बरोबरच काही सूत्रे $\varphi_1$,
\dots,~$\varphi_n$ यांचे दृष्टांत वापरता येणाऱ्या
निष्पत्त्या विचारात घेताना, $\varphi_1$, \dots,
$\varphi_n$ यांचे सर्व आदेशन-दृष्टांत सर्व जगांत सत्य असलेल्या
प्रत्येक प्रतिमानात प्रत्येक निष्पन्न करता येण्याजोगे सूत्र सत्य असावे,
अशी अपेक्षा आहे.

\begin{thm}[निर्दोषता प्रमेय]\label{nml:prf:snd:thm:soundness}
  $\varphi_1$, \dots, $\varphi_n$ यांचा प्रत्येक दृष्टांत
  अनुक्रमे प्रतिमानांच्या $\mClass{C}_1$, \dots,
  $\mClass{C}_n$ या वर्गांमध्ये वैध असेल,
  तर $\Log{K}\varphi_1\dots \varphi_n
  \Proves \psi$ यावरून $\psi$ हे प्रतिमानांच्या
  $\mClass{C}_1 \cap \dots \cap \mClass{C}_n$
  या वर्गात वैध ठरते.
\end{thm}

\begin{proof}
  सिद्धतांच्या लांबीवर विगमन करू. संक्षेपासाठी
  $\mClass{C} =
  \mClass{C}_1 \cap \dots \cap \mClass{C}_n$ असे ठेवा.
  \begin{enumerate}
  \item विगमनाचा पाया: $\psi$ ची सिद्धता लांबी~$1$
    ची असेल, तर ते सर्वतः-सत्य दृष्टांत,
    \Ax{K} चा, किंवा~\Dual{} चा,
    अथवा $\varphi_1$, \dots,~$\varphi_n$ यांपैकी
    एखाद्याचा दृष्टांत असतो. पहिल्या प्रकरणात
    \ref{nml:syn:tau:prop:valid-taut} वरून $\psi$
    हा $\mClass{C}$ मध्ये वैध आहे, कारण सर्वतः-सत्य
    दृष्टांत प्रतिमानांच्या \emph{कोणत्याही} वर्गात
    वैध असतात. दुसऱ्या प्रकरणात
    \ref{nml:syn:sch:prop:Kvalid} आणि
      \ref{nml:syn:sch:prop:Dual-valid} वरून तेच
    मिळते. शेवटी तिसऱ्या प्रकरणात $\psi$ हा
    $\mClass{C}_i$ मध्ये वैध आहे आणि
    $\mClass{C} \subseteq \mClass{C}_i$ आहे;
    म्हणून \ref{nml:syn:val:prop:subset-class}
    वरून $\psi$ हा $\mClass{C}$ मध्येही वैध आहे.
  \item विगमन पायरी: $\psi$ ची सिद्धता लांबी $k>1$
    ची आहे असे समजा. $\psi$ हा सर्वतः-सत्य दृष्टांत,
    मागील पायरीत नमूद केलेल्या मोडल स्वयंसिद्धकांपैकी
    एखाद्याचा दृष्टांत, किंवा $\varphi_1$, \dots,
    $\varphi_n$ यांपैकी एखाद्याचा दृष्टांत असेल,
    तर मागील पायरीप्रमाणे युक्तिवाद लागू होतो.
    म्हणून आधीच्या
    $\chi \lif \psi$ आणि $\chi$ या सूत्रे वर
    \MP{} लावून $\psi$ मिळाला असे समजा. मग
    $\chi \lif \psi$ आणि $\chi$ यांच्या सिद्धतांची
    लांबी $<k$ आहे; विगमन गृहीतकावरून ती
    $\mClass{C}$ मध्ये वैध आहेत.
    \ref{nml:syn:sch:prop:soundMP} वरून $\psi$
    हाही $\mClass{C}$ मध्ये वैध आहे. शेवटी
    $\chi$ वर \Nec{} लावून $\psi$ मिळाला असे समजा
    (म्हणजे $\psi = \Box\chi$). विगमन गृहीतकावरून
    $\chi$ हा $\mClass{C}$ मध्ये वैध आहे आणि
    \ref{nml:syn:val:prop:Nec-rule} वरून $\psi$ ही
    वैध आहे.
  \end{enumerate}
\end{proof}












\section{प्रणाली भिन्न आहेत हे दाखवणे}\label{nml:prf:dis:sec}

\ref{nml:prf:prs:sec} मध्ये दोन मोडल तर्कशास्त्रीय
प्रणाली प्रत्यक्षात एकच आहेत, हे कसे सिद्ध करायचे
ते पाहिले. आता~\ref{nml:prf:snd:thm:soundness}
वापरून दोन मोडल प्रणाली $\Sigma$ आणि $\Sigma'$
भिन्न आहेत, हे दाखवता येते: असे सूत्र~$\varphi$
शोधायचे की $\Sigma' \Proves \varphi$, पण ते~$\Sigma$
च्या प्रतिमानात असत्य ठरते.

\begin{prop}
  $\Log{KD} \subsetneq \Log{KT}$
\end{prop}

\begin{proof} प्रत्येक परावर्ती संबंध सर्वत्र उत्तरवर्ती असतो,
  या अर्थविषयक तथ्याचा हा वाक्यरचनात्मक समकक्ष आहे.
  $\Log{KD} \subseteq \Log{KT}$ दाखवण्यासाठी
  $\Log{KD} \Proves \psi$ यावरून
  $\Log{KT} \Proves \psi$ मिळते हे पाहावे लागेल;
  \ref{nml:prf:prs:prop:S5facts}\ref{nml:prf:prs:prop:S5facts-KT-D}
  मध्ये दाखवल्याप्रमाणे $\Log{KT} \Proves
  \Ax{D}$ यावरून हे मिळते. हा समावेश कडक आहे हे
  दाखवण्यासाठी निर्दोषतेनुसार
  (\ref{nml:prf:snd:thm:soundness}) \Log{KD} चे असे
  प्रतिमान दाखवणे पुरेसे आहे की त्यात
  \Ax{T}, म्हणजे $\Box p \lif p$, असत्य ठरते
  (हे सोपे काम सरावासाठी ठेवले आहे). मग निर्दोषतेवरून
  $\Log{KD} \Proves/ \Box p \lif p$.
\end{proof}

\begin{prop}
  $\Log{KB} \neq \Log{K4}$.
\end{prop}

\begin{proof}
  \Ax{4} चा एखादा दृष्टांत असत्य ठरणारे सममित
  प्रतिमान रचू. तो दृष्टांत \Log{K4} मध्ये
  निष्पन्न करता येण्याजोगे आहे, पण \Log{KB} मध्ये नाही;
  म्हणून $\Log{K4} \nsubseteq \Log{KB}$ मिळेल.
  \ref{nml:prf:dis:fig:Bnot4} मधील सममित प्रतिमान
  $\mModel{M}$ पाहा. प्रतिमान सममित असल्याने
  \Ax{K} आणि~\Ax{B} ही $\mModel{M}$ मध्ये सत्य
  आहेत (अनुक्रमे \ref{nml:syn:sch:prop:Kvalid}
  आणि \ref{nml:frd:acc:thm:soundschemas} वरून).
  मात्र $\mSat/{M}{\Box p \lif \Box\Box p}[w_1]$.
\end{proof}

\begin{figure}[htpb]
  \centering
  \begin{tikzpicture}[modal]
    \node[world] (w1) [label=above:\mFalse{p},
    label={below:$\mSat{{}}{\Box p}$\\ $\mSat/{{}}{\Box\Box p}$}] {$w_1$};
    \node[world] (w2) [label=above:\mTrue{p},
      label=below:$\mSat/{{}}{\Box p}$, right=of w1]{$w_2$};
    \draw[->, bend left] (w1) to (w2);
    \draw[->, bend left] (w2) to (w1);
  \end{tikzpicture}
  \caption{\Ax{4} चा दृष्टांत असत्य ठरणारे सममित प्रतिमान.}
  \label{nml:prf:dis:fig:Bnot4}
\end{figure}

\begin{thm}\label{nml:prf:dis:thm:KTBnot45}
  $\Log{KTB} \Proves/ \Ax{4}$ आणि $\Log{KTB} \Proves/ \Ax{5}$.
\end{thm}

\begin{proof}
  \ref{nml:frd:acc:thm:soundschemas} वरून
  \Ax{T} आणि \Ax{B} यांचे सर्व दृष्टांत
  अनुक्रमे प्रत्येक परावर्ती आणि सममित प्रतिमानात
  सत्य असतात. म्हणून निर्दोषतेनुसार असे
  परावर्ती सममित प्रतिमान शोधणे पुरेसे आहे की
  त्यातील एका जगात~\Ax{4} चा एखादा दृष्टांत
  असत्य ठरेल; \Ax{5} साठीही तसेच.
  दोन्ही दाव्यांसाठी एकच प्रतिमान वापरू.
  \ref{nml:prf:dis:fig:KTBnot45} मधील सममित, परावर्ती
  प्रतिमान पाहा. येथे $\mSat/{M}{\Box p \to \Box \Box
    p}[w_1]$; म्हणून \Ax{4} हा $w_1$ येथे
  असत्य ठरतो. तसेच $\mSat/{M}{\Diamond
    \lnot p \lif \Box \Diamond \lnot p}[w_2]$;
  म्हणून $\varphi = \lnot p$ असलेला~\Ax{5} चा
  दृष्टांत $w_2$ येथे असत्य ठरतो.
\end{proof}

\begin{figure}[htpb]
  \centering
  \begin{tikzpicture}[modal]
    \node[world] (w1) [label={right:\mTrue{p}},
      label={below: $\mSat{{}}{\Box p}$\\
        $\mSat/{{}}{\Box\Box p}$\\
        $\mSat/{{}}{\Diamond\lnot p}$}] {$w_1$};
    \draw[reflexive above] (w1) to (w1);
    \node[world] (w2) [label={right:\mTrue{p}}, label={below:
        $\mSat{{}}{\Diamond\lnot p}$\\
        $\mSat/{{}}{\Box\Diamond\lnot p}$}, right=of w1] {$w_2$} ;
    \draw[reflexive above] (w2) to (w2);
    \node[world] (w3) [label={right:\mFalse{p}},right=of w2] {$w_3$};
    \draw[reflexive above] (w3) to (w3);
    \draw[->, bend left] (w1) to (w2);
    \draw[->, bend left] (w2) to (w3);
    \draw[->, bend left] (w3) to (w2);
    \draw[->, bend left] (w2) to (w1);
  \end{tikzpicture}
  \caption{\ref{nml:prf:dis:thm:KTBnot45} साठीचे प्रतिमान.}
  \label{nml:prf:dis:fig:KTBnot45}
\end{figure}

\begin{thm}\label{nml:prf:dis:thm:KD5not4}
  $\Log{KD5} \neq \Log{KT4} = \Log{S4}$.
\end{thm}

\begin{proof}
  \ref{nml:frd:acc:thm:soundschemas} वरून
  \Ax{D} आणि~\Ax{5} यांचे सर्व दृष्टांत
  प्रत्येक सर्वत्र उत्तरवर्ती युक्लिडी प्रतिमानात
  सत्य असतात. म्हणून असे प्रतिमान शोधणे पुरेसे
  आहे की त्यातील एका जगात~\Ax{4} चा एखादा
  दृष्टांत असत्य ठरेल. \ref{nml:prf:dis:fig:KD5not4}
  मधील प्रतिमान पाहा; येथे
  $\mSat/{M}{\Box p
  \lif \Box\Box p}[w_1]$.
\end{proof}

\begin{figure}[t]
  \centering
  \begin{tikzpicture}[modal]
    \node[world] (w2) [label=north west:\mTrue{p}]{$w_2$} ;
    \draw[reflexive left] (w2) to (w2);
    \node[world] (w1) [label=right:\mFalse{p},
      label=below:{$\mSat{{}}{\Box p}, \mSat/{{}}{\Box\Box p}$},
      below right=of w2]{$w_1$};
    \node[world] (w3) [label=north east:\mTrue{p}, above right=of w1]{$w_3$} ;
    \draw[reflexive right] (w3) to (w3);
    \node[world] (w4) [label=right:\mFalse{p},above right=of w2]{$w_4$};
    \draw[reflexive above] (w4) to (w4);
    \draw[->] (w1) to (w2);
    \draw[->] (w1) to (w3);
    \draw[->, bend left=15] (w2) to (w3);
    \draw[->, bend left=15] (w3) to (w2);
    \draw[->, bend left=15] (w2) to (w4);
    \draw[->, bend left=15] (w4) to (w2);
    \draw[->, bend left=15] (w3) to (w4);
    \draw[->, bend left=15] (w4) to (w3);
  \end{tikzpicture}
  \caption{\ref{nml:prf:dis:thm:KD5not4} साठीचे प्रतिमान.}
  \label{nml:prf:dis:fig:KD5not4}
\end{figure}

\begin{prob}
  $3$ जगांचे प्रतिमान वापरून
  \ref{nml:prf:dis:thm:KD5not4} ची पर्यायी
  सिद्धता द्या.
\end{prob}

\begin{prob}
  $\Log{KT4} \Proves/ \Ax{B}$ आणि
  $\Log{KT4} \Proves/ \Ax{5}$ हे दोन्ही दाखवणारे
  एकच परावर्ती संक्रमक प्रतिमान द्या.
\end{prob}













\section{संचातील सूत्र गृहीत धरून निष्पन्नता}\label{nml:prf:prg:sec}

\ref{nml:prf:prf:sec} मध्ये प्रणाली~$\Sigma$ मधील
सूत्राच्या सिद्धतेची संकल्पना परिभाषित केली.
आता संच~$\Gamma$ मधील सूत्रांपासून
$\Sigma$ मध्ये सिद्धतेसाठी ती विस्तारित करू.

\begin{defn}\label{nml:prf:prg:defn:Gammaproves}
  संच~$\Gamma$ मधील सूत्रांपासून
  प्रणाली~$\Sigma$ मध्ये सूत्र~$\varphi$
  हे निष्पन्न करता येण्याजोगे आहे, असे तेव्हा आणि तेव्हाच
  म्हणतो व $\Gamma \Proves[\Sigma] \varphi$ असे लिहितो,
  जेव्हा $n \geq 0$ आणि $\psi_1$, \dots, $\psi_n \in \Gamma$ अशी
  सूत्रे असतात की
  $\Sigma \Proves \psi_1 \lif (\psi_2 \lif \cdots (\psi_n \lif \varphi)
  \cdots)$.
येथे $n=0$ असताना ही implication मालिका फक्त $\varphi$ असते;
म्हणून कोणतीही गृहीतके न वापरता $\Sigma \Proves \varphi$ मिळते.
\end{defn}












\section{निष्पन्नतेचे गुणधर्म}\label{nml:prf:prp:sec}

\begin{prop}\label{nml:prf:prp:prop:derivabilityfacts}
  $\Sigma$ ही मोडल प्रणाली आणि $\Gamma$ हा मोडल
  सूत्रांचा संच असू द्या. पुढील गुणधर्म लागू होतात:
  \begin{enumerate}
  \item\label{nml:prf:prp:prop:derivabilityfacts-monotonicity} \emph{एकस्वनिकता}:
    $\Gamma \Proves[\Sigma] \varphi$ आणि
    $\Gamma \subseteq \Delta$ असेल, तर
    $\Delta \Proves[\Sigma] \varphi$;
  \item\label{nml:prf:prp:prop:derivabilityfacts-reflexivity}
    \emph{परावर्तिता}: $\varphi \in \Gamma$ असेल, तर $\Gamma
    \Proves[\Sigma] \varphi$;
  \item\label{nml:prf:prp:prop:derivabilityfacts-cut} \emph{कट}: $\Gamma
    \Proves[\Sigma] \varphi$ आणि $\Delta \cup \{\varphi\} \Proves[\Sigma] \psi$
    असेल, तर $\Gamma \cup \Delta \Proves[\Sigma] \psi$;
  \item\label{nml:prf:prp:prop:derivabilityfacts-deduction} \emph{निगमन प्रमेय}: $\Gamma \cup \{\psi\}
    \Proves[\Sigma] \varphi$ असेल तेव्हा आणि तेव्हाच $\Gamma \Proves[\Sigma] \psi \lif
      \varphi$;
  \item \label{nml:prf:prp:prop:derivabilityfacts-ruleT} \emph{नियम T}: जर
    $\Gamma \Proves[\Sigma] \varphi_1$ आणि \dots आणि $\Gamma
    \Proves[\Sigma] \varphi_n$ असेल आणि $\varphi_1 \to (\varphi_2 \lif \cdots (\varphi_n \lif
    \psi)\cdots)$ हा सर्वतः-सत्य दृष्टांत असेल, तर $\Gamma
    \Proves[\Sigma] \psi$.
  \end{enumerate}
\end{prop}

सिद्धता हा सोपा सराव आहे.
\ref{nml:prf:prp:prop:derivabilityfacts} मधील
\ref{nml:prf:prp:prop:derivabilityfacts-ruleT} या भागावरून,
उदाहरणार्थ, $\Gamma \Proves[\Sigma] \varphi \lor \psi$
आणि $\Gamma \Proves[\Sigma] \lnot \varphi$ असेल, तर
$\Gamma \Proves[\Sigma] \psi$ मिळते.
तसेच पुढे $\Gamma \cup \{ \varphi\} \Proves[\Sigma] \psi$
ऐवजी $\Gamma, \varphi \Proves[\Sigma] \psi$ असे लिहू.

\begin{defn}
  $\Gamma \Proves[\Sigma] \varphi$ असेल तेव्हा नेहमी
  $\varphi \in \Gamma$ असेल, तर आणि तेव्हाच संच
  $\Gamma$ हा प्रणाली~$\Sigma$ च्या सापेक्ष
  \emph{निगामीदृष्ट्या बंद} आहे, असे म्हणतो.
\end{defn}













\section{सुसंगतता}\label{nml:prf:con:sec}

सूत्रांच्या संचांचा सुसंगतता हा महत्त्वाचा
गुणधर्म आहे. सूत्रांच्या संचापासून
$\lfalse$ सारखा व्याघात निष्पन्न करता येण्याजोगे असेल,
तर तो संच विसंगत असतो; अन्यथा सुसंगत असतो.
संबंधित प्रणालीच्या सर्व स्वयंसिद्धकांचे सर्व आदेशन-दृष्टांत
प्रतिमानाच्या सर्व जगांत सत्य असतील, तर त्या प्रणालीच्या सापेक्ष
विसंगत संचातील सर्व सूत्रे प्रतिमानाच्या एका जगात
एकत्र सत्य असू शकत नाहीत. संपूर्णता प्रमेयासाठी आपण याचा
व्यत्यास सिद्ध करू: प्रत्येक सुसंगत संच
एका प्रतिमानातील एका जगात सत्य असतो; ते
प्रतिमान म्हणजे ``कॅनॉनिकल प्रतिमान.''

\begin{defn}
  संच~$\Gamma$ हा प्रणाली~$\Sigma$ च्या सापेक्ष
  \emph{सुसंगत}, किंवा आपण म्हणू त्याप्रमाणे
  $\Sigma$-सुसंगत, असेल तेव्हा आणि तेव्हाच
  $\Gamma \Proves/[\Sigma] \lfalse$.
\end{defn}

उदाहरणार्थ, $\{ \Box(p \lif q), \Box p, \lnot\Box q \}$
हा संच विधानीय तर्कशास्त्राच्या सापेक्ष सुसंगत आहे,
पण \Log{K}-सुसंगत नाही. तसेच
$\{ \Diamond p, \Box\Diamond
p \lif q, \lnot q \}$ हा संच \Log{K5}-सुसंगत नाही.

\begin{prop}\label{nml:prf:con:prop:consistencyfacts}
  $\Gamma$ हा सूत्रांचा संच असू द्या. मग:
  \begin{enumerate}
  \item $\Gamma$ हा $\Sigma$-सुसंगत असेल तेव्हा आणि
    तेव्हाच असे एखादे सूत्र~$\varphi$ असते की
    $\Gamma \Proves/[\Sigma] \varphi$.
  \item \label{nml:prf:con:prop:consistencyfacts-b}
    $\Gamma \Proves[\Sigma] \varphi$ असेल तेव्हा आणि
    तेव्हाच $\Gamma \cup \{
    \lnot\varphi \}$ हा $\Sigma$-सुसंगत नाही.
  \item \label{nml:prf:con:prop:consistencyfacts-c}
    $\Gamma$ हा $\Sigma$-सुसंगत असेल, तर
    कोणत्याही सूत्र~$\varphi$ साठी
    $\Gamma \cup \{ \varphi \}$ हा
    $\Sigma$-सुसंगत असतो किंवा
    $\Gamma \cup \{ \lnot\varphi \}$ हा
    $\Sigma$-सुसंगत असतो.
  \end{enumerate}
\end{prop}

\begin{proof}
  ही तथ्ये अभिजात विधानीय तर्कशास्त्रावरून सहज
  मिळतात. \ref{nml:prf:con:prop:consistencyfacts-c} साठी
  युक्तिवाद देऊ. प्रतिपरिवर्तनाने सिद्ध करू:
  $\Gamma \cup \{ \varphi \}$ आणि
  $\Gamma \cup \{ \lnot\varphi \}$ यांपैकी कोणताही
  संच $\Sigma$-सुसंगत नाही, असे समजा. मग
  सुसंगततेच्या व्याख्येवरून
  $\Gamma, \varphi \Proves[\Sigma]
  \lfalse$ आणि $\Gamma, \lnot \varphi \Proves[\Sigma] \lfalse$
  हे दोन्ही मिळतात. निगमन प्रमेयावरून
  $\Gamma \Proves[\Sigma] \varphi \to \lfalse$ आणि
  $\Gamma \Proves[\Sigma] \lnot\varphi \lif \lfalse$.
  पण $(\varphi \lif
  \lfalse) \lif ((\lnot\varphi \lif \lfalse) \lif \lfalse)$
  हा सर्वतः-सत्य दृष्टांत आहे; म्हणून
  \ref{nml:prf:prp:prop:derivabilityfacts}\ref{nml:prf:prp:prop:derivabilityfacts-ruleT}
  वरून $\Gamma \Proves[\Sigma] \lfalse$.
\end{proof}



\chapter{संपूर्णता आणि कॅनॉनिकल प्रतिमाने}\label{nml:com:chap}










\section{प्रस्तावना}\label{nml:com:int:sec}

$\Sigma$ ही मोडल प्रणाली असेल, तर निर्दोषता
प्रमेय असे सांगते: $\Sigma \Proves \varphi$ असल्यास,
ज्या प्रतिमान-वर्ग~$\mClass{C}$ मध्ये~$\Sigma$
मधील सर्व सूत्रांचे सर्व दृष्टांत वैध
आहेत, त्या वर्गात $\varphi$ वैध असते.
विशेषतः, $\Log{K}
\Proves \varphi$ असेल, तर $\varphi$ सर्व प्रतिमानांत
सत्य असते; $\Log{KT} \Proves \varphi$ असेल,
तर $\varphi$ सर्व परावर्ती प्रतिमानांत सत्य असते;
$\Log{KD} \Proves \varphi$ असेल, तर $\varphi$
सर्व सर्वत्र उत्तरवर्ती प्रतिमानांत सत्य असते; इत्यादी.

संपूर्णता हा निर्दोषतेचा व्यत्यास आहे:
उदाहरणार्थ, \Log{K} संपूर्ण आहे याचा अर्थ
सूत्र~$\varphi$ वैध असेल, तर $\Proves \varphi$.
संपूर्णता सिद्ध करणे निर्दोषता सिद्ध करण्यापेक्षा
खूप कठीण आहे. सुरुवातीला प्रतिपरिवर्तन विचारात
घेणे उपयुक्त ठरते: $\Proves/ \varphi$ असेल तेव्हा
नेहमी प्रतिप्रतिमान, म्हणजे
$\mSat/{M}{\varphi}$ असे~$\mModel{M}$ प्रतिमान,
अस्तित्वात असेल, तर आणि तेव्हाच \Log{K}
संपूर्ण असते. समतुल्यपणे ($\varphi$ चा निषेध
घेतल्यास), $\Proves/ \lnot \varphi$ असेल तेव्हा
$\varphi$ चे प्रतिमान असते, हे सिद्ध करता येईल.
असे प्रतिमान रचताना $\varphi$ मधील माहिती
वापरता येते. विशिष्ट सूत्रे साठी
प्रतिमाने शोधतानाही आपण अनेकदा असेच करतो.
उदाहरणार्थ, $p \lif \Box q$ चे प्रतिप्रतिमान
शोधायचे असेल, तर त्यात $p$ सत्य आणि
$\Box q$ असत्य असलेले जग हवे, हे माहीत असते.
ज्या जगात $\Box q$ असत्य आहे, त्या जगापासून
प्राप्य असे एखादे जग असले पाहिजे की त्यात $q$
असत्य आहे. आपल्याला एवढेच ठरवायचे असते:
कोणत्या जगांत विधानीय चले
सत्य आहेत आणि कोणती जगे कोणत्या जगांपासून
प्राप्य आहेत.

पण संपूर्णता सिद्ध करताना प्रतिमान रचण्यासाठी
आपल्यापाशी एखादे ठरावीक सूत्र~$\varphi$
नसते. $\Proves/[\Sigma] \lnot \varphi$ असलेल्या
प्रत्येक $\varphi$ साठी प्रतिमान अस्तित्वात आहे,
हे दाखवायचे असते. ही किमान आवश्यक अट आहे:
\emph{जर} $\Proves[\Sigma] \lnot \varphi$ असेल,
तर निर्दोषतेनुसार~$\varphi$ चे
($\Sigma$ सत्य असलेले) प्रतिमान असू शकत नाही.
आता लक्षात घ्या की $\Proves/[\Sigma]
\lnot \varphi$ असेल तेव्हा आणि तेव्हाच $\varphi$ हे
$\Sigma$-सुसंगत असते. ($\Sigma
\Proves/[\Sigma] \lnot \varphi$ आणि
$\varphi \Proves/[\Sigma] \lfalse$ समतुल्य
आहेत, हे आठवा.) म्हणून प्रत्येक
$\Sigma$-सुसंगत सूत्र साठी प्रतिमान
रचणे हे आपले काम आहे.

आपण पुढील युक्ती वापरू: $\varphi$ सामावणारा
$\Sigma$-सुसंगत सूत्रांचा संच शोधू;
त्यात $\varphi$ सत्य करणारे जग कसे असावे
हे सांगणारी इतर सूत्रेही असतील. असे संच
\emph{संपूर्ण} $\Sigma$-सुसंगत संच असतात.
प्रत्येक सुसंगत $\varphi$ साठी काम करणाऱ्या सर्वसाधारण रचनेत
फक्त एकच जग ठेवणे नेहमी पुरेसे नसते. अनेक जगे आणि
त्यांच्यातील प्राप्यता संबंध असण्यास रचनेने परवानगी दिली पाहिजे. $\varphi$ सामावणाऱ्या
संपूर्ण $\Sigma$-सुसंगत संचात
$\Box \psi$ आणि~$\Diamond \chi$ या रूपांतील
इतर सूत्रे ही असतील. सर्व प्राप्य
जगांत $\psi$ सत्य असले पाहिजे; किमान एका
प्राप्य जगात $\chi$ सत्य असले पाहिजे. यासाठी
शक्य असलेले \emph{सर्व} संपूर्ण
$\Sigma$-सुसंगत संच जगांच्या संचाचा
आधार म्हणून घेऊ. आपल्या प्रतिमानात एक
संपूर्ण $\Sigma$-सुसंगत संच दुसऱ्यापासून
केव्हा प्राप्य मानायचा, हे ठरवणे अवघड
ठरेल.

असे परिभाषित केलेल्या प्रतिमानात एखाद्या
जगात---जे स्वतः संपूर्ण $\Sigma$-सुसंगत
संच आहे---$\varphi$ सत्य असेल तेव्हा आणि
तेव्हाच $\varphi$ त्या संचाचा घटक
असेल, हे दाखवू. $\varphi$ हे
$\Sigma$-सुसंगत असेल, तर किमान एका
संपूर्ण $\Sigma$-सुसंगत संचाचा ते
घटक असेल (हे आपण सिद्ध करू);
म्हणून $\varphi$~सत्य असलेले जग मिळेल. अशा
रीतीने प्रत्येक $\Sigma$-सुसंगत
सूत्र~$\varphi$ काहीतरी जगात सत्य असलेले
एकच प्रतिमान आपल्याला मिळेल.
हेच~$\Sigma$ साठीचे \emph{कॅनॉनिकल}
प्रतिमान आहे.













\section{संपूर्ण $\Sigma$-सुसंगत संच}\label{nml:com:ccs:sec}

$\Sigma$ हा मोडल सूत्रांचा संच आहे असे
समजा---सामान्य मोडल तर्कशास्त्राची स्वयंसिद्धके
किंवा परिभाषक तत्त्वे म्हणून त्यांच्याकडे पाहा.
संच~$\Gamma$ हा $\Sigma$-सुसंगत असेल तेव्हा
आणि तेव्हाच $\Gamma \Proves/[\Sigma] \lfalse$;
म्हणजे प्रत्येक $\varphi_i \in \Gamma$ असताना
$\Sigma$ पासून~$\varphi_1 \lif (\varphi_2 \lif \cdots (\varphi_n \lif
\lfalse)\dots)$ ची निष्पत्ती नसते.
आपण ``कॅनॉनिकल'' प्रतिमान रचू; त्यात प्रत्येक
जग विशिष्ट प्रकारचा $\Sigma$-सुसंगत संच
असेल. तो केवळ $\Sigma$-सुसंगत नसून
महत्तम अर्थाने तसा असेल: प्रत्येक मोडल
सूत्राचे सत्यत्व तो ठरवेल.
प्रत्येक $\varphi$ साठी $\varphi
\in \Gamma$ किंवा $\lnot \varphi \in \Gamma$:

\begin{defn}
  संच~$\Gamma$ हा \emph{संपूर्ण $\Sigma$-सुसंगत}
  असेल तेव्हा आणि तेव्हाच तो $\Sigma$-सुसंगत
  असेल आणि प्रत्येक~$\varphi$ साठी
  $\varphi \in \Gamma$ किंवा $\lnot \varphi \in \Gamma$
  यांपैकी एक विधान खरे असेल.
\end{defn}

संपूर्ण $\Sigma$-सुसंगत संच~$\Gamma$ चे अनेक
उपयुक्त गुणधर्म आहेत. प्रथम, ते निगामीदृष्ट्या
बंद असतात: $\Gamma
\Proves[\Sigma] \varphi$ असेल, तर $\varphi \in \Gamma$.
विशेषतः, सूत्र~$\varphi \in \Sigma$ चा प्रत्येक
दृष्टांतही $\in
\Gamma$ असतो. शिवाय $\Gamma$ मधील सदस्यत्व
विधानीय संयोगकांच्या सत्यत्व-अटींशी जुळते.
आपण ``कॅनॉनिकल प्रतिमान'' परिभाषित करू तेव्हा
हे महत्त्वाचे ठरेल.

\begin{prop}\label{nml:com:ccs:prop:ccs-properties}
  $\Gamma$ हा संपूर्ण $\Sigma$-सुसंगत आहे असे
  समजा. मग:
  \begin{enumerate}
  \item \label{nml:com:ccs:prop:ccs-closed}
    $\Gamma$ हा~$\Sigma$ च्या सापेक्ष
    निगामीदृष्ट्या बंद आहे.
  \item \label{nml:com:ccs:prop:ccs-sigma}
    $\Sigma \subseteq \Gamma$.
  \item \label{nml:com:ccs:prop:ccs-lfalse}
    $\lfalse \notin \Gamma$
  \item \label{nml:com:ccs:prop:ccs-ltrue}
    $\ltrue \in \Gamma$
  \item \label{nml:com:ccs:prop:ccs-lnot}
    $\lnot\varphi \in \Gamma$ असेल तेव्हा आणि
    तेव्हाच $\varphi \notin
    \Gamma$.
  \item \label{nml:com:ccs:prop:ccs-land}
    $\varphi \land \psi \in \Gamma$ असेल तेव्हा आणि तेव्हाच
    $\varphi \in \Gamma$ आणि $\psi \in \Gamma$ असतील.
  \item \label{nml:com:ccs:prop:ccs-lor}
    $\varphi \lor \psi \in \Gamma$ असेल तेव्हा आणि तेव्हाच
    $\varphi \in \Gamma$ किंवा $\psi \in \Gamma$ असेल.
  \item \label{nml:com:ccs:prop:ccs-lif}
    $\varphi \lif \psi \in \Gamma$ असेल तेव्हा आणि तेव्हाच
    $\varphi \notin \Gamma$ किंवा $\psi \in \Gamma$ असेल.
  \item \label{nml:com:ccs:prop:ccs-liff}
    $\varphi \liff \psi \in \Gamma$ असेल तेव्हा आणि तेव्हाच
    $\varphi \in \Gamma$ आणि $\psi \in
    \Gamma$ दोन्ही असतील, किंवा
    $\varphi \notin \Gamma$ आणि $\psi \notin \Gamma$
    दोन्ही असतील.
  \end{enumerate}
\end{prop}

\begin{proof}
  \begin{enumerate}
  \item $\Gamma \Proves[\Sigma] \varphi$ पण $\varphi \notin
    \Gamma$ असे समजा. $\Gamma$ संपूर्ण
    $\Sigma$-सुसंगत असल्याने
    $\lnot\varphi \in \Gamma$. मग $\Gamma$ विसंगत
    ठरेल, कारण $\varphi, \lnot \varphi \Proves[\Sigma] \lfalse$.

  \item $\varphi \in \Sigma$ असेल, तर
    $\Gamma \Proves[\Sigma] \varphi$; आणि
    निगामी बंदतेवरून, म्हणजे
    \ref{nml:com:ccs:prop:ccs-closed} वरून, $\varphi
    \in \Gamma$.

  \item $\lfalse \in \Gamma$ असेल,
    तर $\Gamma \Proves[\Sigma] \lfalse$;
    त्यामुळे $\Gamma$ हा $\Sigma$-विसंगत
    ठरेल.

  \item $\Gamma \Proves[\Sigma] \ltrue$;
    म्हणून निगामी बंदतेवरून, म्हणजे
    \ref{nml:com:ccs:prop:ccs-closed} वरून,
    $\ltrue \in \Gamma$.

\item $\lnot\varphi \in \Gamma$ असेल, तर
    सुसंगततेमुळे $\varphi \notin \Gamma$.
    उलट $\varphi \notin \Gamma$ असेल, तर
    $\Gamma$ संपूर्ण $\Sigma$-सुसंगत
    असल्याने $\lnot\varphi \in \Gamma$.

  \item $\varphi \land \psi \in
      \Gamma$ असे समजा. $(\varphi \land \psi) \lif \varphi$
      हा सर्वतः-सत्य दृष्टांत असल्याने,
      निगामी बंदतेवरून, म्हणजे
      \ref{nml:com:ccs:prop:ccs-closed} वरून, $\varphi \in \Gamma$.
      $\psi \in \Gamma$ साठीही तसेच. उलट
      $\varphi \in \Gamma$ आणि $\psi \in
      \Gamma$ दोन्ही असतील, तर निगामी
      बंदतेवरून $(\varphi \land \psi) \in
      \Gamma$; कारण $\varphi \lif (\psi \lif (\varphi \land \psi))$
      हा सर्वतः-सत्य दृष्टांत आहे.

  \item $\varphi \lor \psi \in
      \Gamma$, पण $\varphi \notin \Gamma$ आणि
      $\psi \notin \Gamma$ असे समजा. $\Gamma$
      संपूर्ण $\Sigma$-सुसंगत असल्याने
      $\lnot\varphi \in \Gamma$ आणि
      $\lnot \psi \in \Gamma$. मग
      $\lnot(\varphi \lor \psi) \in \Gamma$;
      कारण $\lnot \varphi \lif (\lnot \psi \lif \lnot (\varphi \lor \psi))$
      हा सर्वतः-सत्य दृष्टांत आहे. त्यामुळे
      $\Gamma$ हा $\Sigma$-विसंगत ठरेल,
      हा व्याघात आहे. उलट $\varphi \in \Gamma$ असेल,
      तर $\varphi \lif (\varphi \lor \psi)$ हा सर्वतः-सत्य दृष्टांत
      आणि निगामी बंदता यांवरून $\varphi \lor \psi \in \Gamma$.
      तसेच $\psi \in \Gamma$ असेल, तर
      $\psi \lif (\varphi \lor \psi)$ वापरून तोच निष्कर्ष मिळतो.

  \item $\varphi \lif \psi \in
      \Gamma$ आणि $\varphi \in \Gamma$ असे समजा.
      मग $\Gamma \Proves[\Sigma] \psi$;
      त्यामुळे निगामी बंदतेवरून $\psi \in \Gamma$.
      उलट $\varphi
      \lif \psi \notin \Gamma$ असेल, तर
      $\Gamma$ संपूर्ण $\Sigma$-सुसंगत
      असल्याने $\lnot (\varphi \lif \psi) \in \Gamma$.
      $\lnot(\varphi \lif \psi) \lif \varphi$ हा
      सर्वतः-सत्य दृष्टांत असल्याने निगामी
      बंदतेवरून $\varphi \in
      \Gamma$. तसेच $\lnot(\varphi \lif \psi) \lif
      \lnot \psi$ हा सर्वतः-सत्य दृष्टांत असल्याने
      $\lnot \psi \in
      \Gamma$. $\Gamma$ हा $\Sigma$-सुसंगत
      असल्यामुळे $\psi \notin \Gamma$.

  \item $\varphi \liff \psi \in
      \Gamma$ असे समजा. $\varphi \in \Gamma$
      असेल, तर $\psi \in \Gamma$; कारण $(\varphi
      \liff \psi) \lif (\varphi \lif \psi)$ हा
      सर्वतः-सत्य दृष्टांत आहे. तसेच
      $\psi \in \Gamma$ असेल, तर $\varphi \in \Gamma$.
      म्हणून $\varphi \in \Gamma$ आणि $\psi \in \Gamma$
      दोन्ही असतील, किंवा $\varphi \in \Gamma$
      आणि $\psi \in \Gamma$ यांपैकी कोणतेही नसेल.

      उलट $\varphi \liff \psi \notin \Gamma$
      असे समजा. $\Gamma$ संपूर्ण
      $\Sigma$-सुसंगत असल्याने
      $\lnot (\varphi \liff \psi) \in
      \Gamma$. $\lnot(\varphi \liff \psi) \lif (\varphi \lif \lnot \psi)$
      हा सर्वतः-सत्य दृष्टांत असल्याने
      $\varphi \in \Gamma$ असेल, तर
      $\lnot \psi \in
      \Gamma$; आणि $\Gamma$ हा
      $\Sigma$-सुसंगत असल्याने $\psi \notin
      \Gamma$. तसेच $\psi \in \Gamma$
      असेल, तर $\varphi \notin \Gamma$.
      म्हणून $\varphi \in \Gamma$ आणि $\psi \in \Gamma$
      ही दोन्ही विधाने एकत्र खरी असण्याची शक्यता नाही;
      तसेच $\varphi \notin \Gamma$ आणि $\psi \notin \Gamma$
      ही दोन्ही विधाने एकत्र खरी असण्याचीही शक्यता नाही.
  \end{enumerate}
\end{proof}















\section{लिंडेनबाउमचे पूर्वप्रमेय}\label{nml:com:lin:sec}

लिंडेनबाउमचे पूर्वप्रमेय असे सांगते की सूत्रांचा प्रत्येक
$\Sigma$-सुसंगत संच एखाद्या \emph{संपूर्ण} $\Sigma$-सुसंगत
संचात समाविष्ट असतो. कॅनॉनिकल प्रतिमानाची आपली रचना दाखवेल की
प्रत्येक संपूर्ण $\Sigma$-सुसंगत संच~$\Delta$ साठी कॅनॉनिकल
प्रतिमानात असे एक जग असते की नेमके~$\Delta$ मधील
सूत्रे त्या जगात सत्य असतात. म्हणून लिंडेनबाउमचे पूर्वप्रमेय
प्रत्येक $\Sigma$-सुसंगत संच कॅनॉनिकल प्रतिमानातील कोणत्या तरी
जगात सत्य असल्याची हमी देते.

\begin{thm}[लिंडेनबाउमचे पूर्वप्रमेय]\label{nml:com:lin:thm:lindenbaum}
  $\Gamma$ हा $\Sigma$-सुसंगत असेल, तर~$\Gamma$ चा विस्तार असलेला
  संपूर्ण $\Sigma$-सुसंगत संच~$\Delta$ अस्तित्वात असतो.
\end{thm}

\begin{proof}
भाषेतील सर्व सूत्रांची $\varphi_0$, $\varphi_1$, \dots{} अशी यादी असू द्या;
त्यात एखादे सूत्र पुन्हा येऊ शकते. उदाहरणार्थ, आधी
$\Obj p_0$ नोंदवा आणि नंतर प्रत्येक $n \ge 1$ टप्प्यावर
$\Obj p_0$, \dots,~$\Obj p_n$ यांमधीलच चल वापरून
बनणारी व लांबी~$n$ पर्यंतची सर्व सूत्रे नोंदवा;
अशा सूत्रांची संख्या सांत असते. आपण~$n$ वरील विगमनाने
सूत्रांचे संच~$\Delta_n$ परिभाषित करू आणि मग
$\Delta = \bigcup_n \Delta_n$ ठेवू. प्रथम $\Delta_0 = \Gamma$
ठेवू. $\Delta_n$ परिभाषित झाला आहे असे धरून
$\Delta_{n+1}$ ची व्याख्या अशी करू:
\[\Delta_{n+1} =
\begin{cases}
  \Delta_n \cup \{\varphi_n\}, & \text{जर $\Delta_n \cup \{ \varphi_n\}$
    हा $\Sigma$-सुसंगत असेल;} \\
  \Delta_n \cup \{ \lnot \varphi_n\}, & \text{अन्यथा.}
\end{cases}
\]
आता $\Delta = \bigcup_{n=0}^\infty \Delta_n$ ठेवू.

या व्याख्येतून आवश्यक गुणधर्म असलेला संच~$\Delta$
खरोखर मिळतो हे दाखवायचे आहे: $\Gamma \subseteq \Delta$
आणि~$\Delta$ हा संपूर्ण $\Sigma$-सुसंगत आहे.

$\Gamma \subseteq \Delta$ हे उघड आहे, कारण रचनेनुसार
$\Delta_0 \subseteq \Delta$ आणि $\Delta_0 = \Gamma$.
खरे तर प्रत्येक~$n$ साठी $\Delta_n \subseteq \Delta$,
कारण~$\Delta$ हा सर्व $\Delta_n$ चा संघ आहे. (रचनेच्या
प्रत्येक पायरीवर आधीच्या संचात सूत्र जोडतो; म्हणून
$\Delta_n \subseteq \Delta_{n+1}$. तसेच
$\subseteq$ संक्रमणशील असल्याने $n \le m$ असेल तेव्हा
$\Delta_n \subseteq \Delta_{m}$.) रचनेच्या प्रत्येक
टप्प्यावर आपण $\varphi_n$ किंवा $\lnot \varphi_n$ जोडतो आणि
प्रत्येक सूत्र हे सर्व~$\varphi_n$ च्या यादीत किमान
एकदा येते. म्हणून प्रत्येक $\varphi$ साठी $\varphi
\in \Delta$ किंवा $\lnot \varphi \in \Delta$;
त्यामुळे व्याख्येनुसार~$\Delta$ संपूर्ण आहे.

शेवटी~$\Delta$ हा $\Sigma$-सुसंगत आहे हे दाखवायचे आहे.
त्यासाठी (a)~$\Delta$ हा $\Sigma$-विसंगत असता, तर
एखादा~$\Delta_n$ $\Sigma$-विसंगत असता, आणि
(b)~सर्व $\Delta_n$ हे $\Sigma$-सुसंगत आहेत,
असे आपण दाखवू.

समजा~$\Delta$ हा $\Sigma$-विसंगत आहे. मग
$\Delta \Proves[\Sigma] \lfalse$; म्हणजे
$\varphi_1$, \dots,~$\varphi_k \in \Delta$ अशी सूत्रे असतात की
$\Sigma \Proves \varphi_1 \lif (\varphi_2 \lif \cdots (\varphi_k
\lif \lfalse)\dots)$. $\Delta = \bigcup_{n=0}^\infty \Delta_n$
असल्याने प्रत्येक $\varphi_i \in \Delta_{n_i}$ असे कोणत्यातरी~$n_i$
साठी असते. $n=\max(\{0\}\cup\{n_1,\dots,n_k\})$ घ्या;
यामुळे शून्य गृहीतकांच्या $k=0$ प्रकरणातही $n=0$ ठरतो. $n_i \le n$ असल्यामुळे
$\Delta_{n_i} \subseteq \Delta_n$. म्हणून ही सर्व
$\varphi_i$ सूत्रे एखाद्या~$\Delta_n$ मध्ये असतात.
याचा अर्थ $\Delta_n \Proves[\Sigma] \lfalse$,
म्हणजेच~$\Delta_n$ हा $\Sigma$-विसंगत आहे.

प्रत्येक $\Delta_n$ $\Sigma$-सुसंगत आहे हे दाखवण्यासाठी
आपण~$n$ वरील सोपे विगमन वापरू. $\Delta_0 = \Gamma$
आणि~$\Gamma$ $\Sigma$-सुसंगत असल्याचे गृहीत आहे;
म्हणून $n = 0$ साठी दावा खरा आहे. आता तो~$n$
साठी खरा आहे, म्हणजे $\Delta_n$ $\Sigma$-सुसंगत
आहे, असे समजा. $\Delta_n \cup \{\varphi_n\}$ हा
$\Sigma$-सुसंगत असेल, तर $\Delta_{n+1}$ हाच
संच असतो; अन्यथा तो $\Delta_n \cup \{\lnot\varphi_n\}$
असतो. पहिल्या प्रकरणात $\Delta_{n+1}$
$\Sigma$-सुसंगत असणे स्पष्ट आहे. परंतु
\ref{nml:prf:con:prop:consistencyfacts}\ref{nml:prf:con:prop:consistencyfacts-c}
नुसार $\Delta_n \cup \{\varphi_n\}$ किंवा
$\Delta_n \cup \{\lnot\varphi_n\}$ यांपैकी एक
सुसंगत असतो; म्हणून दुसऱ्या प्रकरणातही
$\Delta_{n+1}$ सुसंगत असतो.
\end{proof}

\begin{cor}\label{nml:com:lin:cor:provability-characterization}
  $\Gamma \Proves[\Sigma] \varphi$ असेल तेव्हा आणि तेव्हाच
  $\Gamma$ चा विस्तार असलेल्या प्रत्येक संपूर्ण
  $\Sigma$-सुसंगत संच~$\Delta$ मध्ये $\varphi \in \Delta$
  असते. ($\Gamma = \emptyset$ असतानाही हे लागू होते;
  त्या बाबतीत मोडल प्रणाली~$\Sigma$ चे आणखी एक
  लक्षणन मिळते.)
\end{cor}

\begin{proof}
  $\Gamma \Proves[\Sigma] \varphi$ असे समजा आणि
  $\Gamma$ चा विस्तार असलेला कोणताही संपूर्ण
  $\Sigma$-सुसंगत संच~$\Delta$ घ्या. जर $\varphi
  \notin \Delta$ असेल, तर महत्तमतेमुळे
  $\lnot\varphi \in \Delta$. त्यामुळे एकस्वनिकतेने
  $\Delta \Proves[\Sigma] \varphi$ आणि स्वतःच्या गृहीतकावरून
  $\Delta \Proves[\Sigma] \lnot\varphi$; म्हणून~$\Delta$
  विसंगत ठरेल. उलट $\Gamma \Proves/[\Sigma] \varphi$
  असेल, तर $\Gamma \cup \{ \lnot\varphi\}$ हा
  $\Sigma$-सुसंगत संच असतो. लिंडेनबाउमच्या
  पूर्वप्रमेयानुसार $\Gamma \cup \{ \lnot\varphi \}$
  चा विस्तार असलेला संपूर्ण सुसंगत संच~$\Delta$
  अस्तित्वात असतो. सुसंगततेमुळे $\varphi
  \notin \Delta$.
\end{proof}













\section{मोडल कारक आणि संपूर्ण सुसंगत संच}\label{nml:com:mod:sec}

\begin{explain}
एखाद्या सामान्य मोडल तर्कशास्त्र~$\Sigma$ मधील संपूर्ण
$\Sigma$-सुसंगत संच~$\Delta$ हे ज्याचे जग आहेत,
असे प्रतिमान~$\mModel{M^\Sigma}$ रचताना अशा “जगां”मधील
प्राप्यता संबंध~$R^\Sigma$ देखील परिभाषित करावा लागेल.
प्राप्यता संबंध (आणि मूल्यांकन~$V^\Sigma$) यांची व्याख्या
अशी हवी की $\mSat{M^\Sigma}{\varphi}[\Delta]$ असेल तेव्हा
आणि तेव्हाच $\varphi \in \Delta$. हे कसे करावे?

प्राप्यता संबंध परिभाषित केल्यावर, जगातील सत्यतेच्या
व्याख्येनुसार $\mSat{M^\Sigma}{\Box\varphi}[\Delta]$
  असेल तेव्हा आणि तेव्हाच $R^\Sigma\Delta\Delta'$ असलेल्या
  प्रत्येक $\Delta'$ साठी $\mSat{M^\Sigma}{\varphi}[\Delta']$.
$\mSat{M^\Sigma}{\varphi}[\Delta]$ असेल तेव्हा आणि तेव्हाच
$\varphi \in \Delta$ हे सिद्ध करण्यासाठी हे विशेषतः
मोडल कारकाने सुरू होणाऱ्या सूत्रे साठी खरे असणे
आवश्यक आहे; म्हणजे
$\mSat{M^\Sigma}{\Box\varphi}[\Delta]$ असेल
  तेव्हा आणि तेव्हाच $\Box \varphi \in \Delta$.
ही आवश्यकता आणि
$\Box \varphi$
साठी जगातील सत्यतेची व्याख्या एकत्र केल्यास:
\[  \Box\varphi \in \Delta \text{ तेव्हा आणि तेव्हाच } \varphi \in \Delta' \text{
    प्रत्येक $\Delta'$ साठी, जेथे } R^\Sigma\Delta\Delta'
\]
डावीकडून उजवीकडे जाणारी दिशा पाहा:
  $\Box \varphi \in \Delta$ असेल, तर कोणत्याही $\varphi$
  आणि $R^\Sigma\Delta\Delta'$ असलेल्या कोणत्याही
  $\Delta'$ साठी $\varphi \in \Delta'$ असे ती सांगते.
  प्रत्येक $\Box\varphi \in \Delta$ साठी $\varphi \in \Delta'$
  असेल तेव्हा आणि तेव्हाच $R^\Sigma\Delta\Delta'$
  अशी अट ठेवल्यास हे खरे ठरते. या द्विसशर्ताच्या
  उजव्या बाजूची अट संक्षिप्तपणे
  $\Setabs{\varphi}{\Box\varphi \in \Delta} \subseteq \Delta'$
  अशी लिहिता येते.

मग प्रश्न असा आहे: $R^\Sigma$ ची ही व्याख्या
$\Box \varphi \in \Delta$ असेल तेव्हा आणि
  तेव्हाच $\mSat{M^\Sigma}{\Box \varphi}[\Delta]$ याची
  खरोखर हमी देते का? ती
    $\Diamond \varphi \in \Delta$
  असेल तेव्हा आणि तेव्हाच
  $\mSat{M^\Sigma}{\Diamond \varphi}[\Delta]$ याचीही
  हमी देते का? पुढील काही निष्कर्ष हे सिद्ध करतील.
\end{explain}

\begin{defn}
  $\Gamma$ हा सूत्रांचा संच असेल, तर
  \begin{align*}
    \Box\Gamma & = \Setabs{\Box \psi}{\psi \in \Gamma}\\
    \Diamond\Gamma & = \Setabs{\Diamond \psi}{\psi \in \Gamma}\\
    \intertext{आणि}
    \Box^{-1}\Gamma & = \Setabs{\psi}{\Box \psi \in \Gamma}\\
    \Diamond^{-1}\Gamma & = \Setabs{\psi}{\Diamond \psi \in \Gamma}\\
  \end{align*}
\end{defn}

दुसऱ्या शब्दांत, $\Box\Gamma$ म्हणजे~$\Gamma$ मधील
प्रत्येक सूत्राच्या पुढे $\Box$ लावलेला संच;
ही सर्व सूत्रे~$\Gamma$ मधूनच घेतली जातात.
$\Box^{-1}\Gamma$ मध्ये~$\Gamma$ मधील $\Box$ ने
सुरू होणाऱ्या सर्व सूत्रे मधील सुरुवातीचा
$\Box$ काढून टाकला जातो. ही व्याख्या स्वतंत्रपणे
फारशी महत्त्वाची नाही, पण संकेतलेखन बरेच सोपे करते.

लक्षात घ्या की $\Box\Box^{-1}\Gamma \subseteq \Gamma$:
\[\Box\Box^{-1}\Gamma = \Setabs{\Box\psi}{\Box\psi \in \Gamma}
\]
म्हणजे हा~$\Gamma$ मधील $\Box$ ने सुरू होणाऱ्या
सर्व सूत्रांचाच संच आहे.

\begin{lem}\label{nml:com:mod:lem:box1}
  $\Gamma \Proves[\Sigma] \varphi$ असेल, तर
  $\Box\Gamma \Proves[\Sigma] \Box \varphi$.
\end{lem}

\begin{proof}
  $\Gamma \Proves[\Sigma] \varphi$ असेल, तर
  $\psi_1$, \dots, $\psi_k \in \Gamma$ अशी सूत्रे
  असतात. $k=0$ असेल, तर $\Sigma \Proves \varphi$;
  \Nec{} ने $\Sigma \Proves \Box\varphi$ मिळते आणि
  शून्य गृहीतकांच्या प्रकरणानुसार $\Box\Gamma \Proves[\Sigma] \Box\varphi$.
  आता $k\geq1$ असे समजा. मग $\Sigma \Proves \psi_1 \lif (\psi_2 \lif \cdots
  (\psi_k \lif \varphi)\cdots)$. $\Sigma$ सामान्य असल्याने
  \RK{} नियमानुसार $\Sigma \Proves \Box\psi_1 \lif
  (\Box\psi_2 \lif \cdots (\Box\psi_k \lif \Box\varphi)\cdots)$;
  येथे अर्थातच $\Box\psi_1$, \dots, $\Box\psi_k \in
  \Box\Gamma$. म्हणून व्याख्येनुसार
  $\Box\Gamma \Proves[\Sigma] \Box \varphi$.
\end{proof}

\begin{lem}\label{nml:com:mod:lem:box2}
  $\Box^{-1} \Gamma \Proves[\Sigma] \varphi$ असेल,
  तर $\Gamma \Proves[\Sigma] \Box\varphi$.
\end{lem}

\begin{proof}
  $\Box^{-1}\Gamma \Proves[\Sigma] \varphi$ असे
  समजा. मग \ref{nml:com:mod:lem:box1} नुसार
  $\Box\Box^{-1}\Gamma \Proves[\Sigma] \Box\varphi$.
  पण $\Box\Box^{-1}\Gamma \subseteq \Gamma$
  असल्याने एकस्वनिकतेनुसार
  $\Gamma \Proves[\Sigma] \Box\varphi$.
\end{proof}




\begin{prop}\label{nml:com:mod:prop:box}
  $\Gamma$ हा संपूर्ण $\Sigma$-सुसंगत असेल,
  तर $\Box \varphi \in \Gamma$ असेल तेव्हा आणि
  तेव्हाच $\Box^{-1} \Gamma \subseteq \Delta$
  असलेल्या प्रत्येक संपूर्ण $\Sigma$-सुसंगत
  $\Delta$ साठी $\varphi \in \Delta$ असेल.
\end{prop}

\begin{proof}
  $\Gamma$ हा संपूर्ण $\Sigma$-सुसंगत आहे
  असे समजा. डावीकडून उजवीकडे जाणारी दिशा सोपी आहे:
  $\Box \varphi \in \Gamma$ आणि
  $\Box^{-1}\Gamma \subseteq \Delta$ असे
  समजा. $\Box\varphi \in \Gamma$ असल्याने
  $\varphi \in \Box^{-1}\Gamma \subseteq \Delta$;
  म्हणून $\varphi \in \Delta$.

  उजवीकडून डावीकडे जाणाऱ्या दिशेसाठी आपण प्रतिपरिवर्तन सिद्ध करू:
  $\Box\varphi \notin \Gamma$ असे समजा.
  $\Gamma$ हा संपूर्ण $\Sigma$-सुसंगत
  असल्याने तो निगामीदृष्ट्या बंद आहे;
  म्हणून $\Gamma \Proves/[\Sigma] \Box \varphi$.
  \ref{nml:com:mod:lem:box2} नुसार
  $\Box^{-1}\Gamma \Proves/[\Sigma] \varphi$.
  \ref{nml:prf:con:prop:consistencyfacts}\ref{nml:prf:con:prop:consistencyfacts-b}
  नुसार $\Box^{-1}\Gamma \cup \{ \lnot\varphi \}$
  हा $\Sigma$-सुसंगत आहे. लिंडेनबाउमच्या
  पूर्वप्रमेयानुसार $\Box^{-1}\Gamma \cup
  \{ \lnot\varphi \} \subseteq \Delta$ असा
  संपूर्ण $\Sigma$-सुसंगत संच~$\Delta$
  असतो. सुसंगततेमुळे $\varphi \notin \Delta$.
\end{proof}


\begin{lem}\label{nml:com:mod:lem:box-iff-diamond}
  $\Gamma$ आणि $\Delta$ हे संपूर्ण
  $\Sigma$-सुसंगत असोत. मग
  $\Box^{-1}\Gamma \subseteq \Delta$
  असेल तेव्हा आणि तेव्हाच
  $\Diamond\Delta \subseteq \Gamma$.
\end{lem}

\begin{proof}
  डावीकडून उजवीकडे जाणारी दिशा: $\Box^{-1}\Gamma
  \subseteq \Delta$ असे समजा आणि
  $\Diamond\varphi \in \Diamond\Delta$
  (म्हणजे $\varphi \in \Delta$) असे धरा.
  $\Diamond\varphi \in \Gamma$ दाखवण्यासाठी
  $\Box\lnot\varphi \notin \Gamma$ दाखवणे
  पुरेसे आहे; कारण मग महत्तमतेमुळे
  $\lnot\Box\lnot \varphi \in \Gamma$.
  आता $\Box\lnot\varphi \in \Gamma$
  असेल, तर गृहीतकानुसार
  $\lnot\varphi \in \Delta$; पण
  $\varphi \in \Delta$ असल्याने हे
  $\Delta$ च्या सुसंगततेच्या
  विरुद्ध आहे. म्हणून आवश्यकतेनुसार
  $\Box\lnot\varphi \notin \Gamma$.

  उजवीकडून डावीकडे जाणारी दिशा: $\Diamond\Delta
  \subseteq \Gamma$ असे समजा.
  प्रतिपरिवर्तन वापरू: $\Box\varphi
  \notin \Gamma$ दाखवण्यासाठी
  $\varphi \notin \Delta$ असे समजा.
  $\varphi \notin \Delta$ असेल,
  तर महत्तमतेमुळे
  $\lnot\varphi \in \Delta$;
  म्हणून गृहीतकानुसार
  $\Diamond\lnot\varphi \in \Gamma$.
  पण सामान्य मोडल तर्कशास्त्रात
  $\Diamond\lnot\varphi$ हे
  $\lnot\Box \varphi$ शी समतुल्य आहे;
  आणि दुसरे सूत्र~$\Gamma$ मध्ये
  असेल, तर सुसंगततेमुळे
  $\Box\varphi \notin\Gamma$,
  हेच दाखवायचे होते.
\end{proof}






\begin{prop}\label{nml:com:mod:prop:diamond}
  $\Gamma$ हा संपूर्ण $\Sigma$-सुसंगत
  असेल, तर $\Diamond \varphi \in \Gamma$
  असेल तेव्हा आणि तेव्हाच
  $\Diamond\Delta \subseteq
  \Gamma$ असा एखादा संपूर्ण
  $\Sigma$-सुसंगत $\Delta$
  असेल की $\varphi \in \Delta$.
\end{prop}

\begin{proof}
$\Gamma$ हा संपूर्ण $\Sigma$-सुसंगत
आहे असे समजा. \Dual{} आणि निगामी बंदतेनुसार
$\Diamond\varphi \in \Gamma$ असेल तेव्हा
आणि तेव्हाच $\lnot\Box\lnot \varphi \in
\Gamma$. $\Gamma$ संपूर्ण
$\Sigma$-सुसंगत असल्याने
\ref{nml:com:ccs:prop:ccs-properties}\ref{nml:com:ccs:prop:ccs-lnot}
नुसार $\lnot\Box\lnot \varphi \in \Gamma$
असेल तेव्हा आणि तेव्हाच
$\Box\lnot \varphi \notin \Gamma$.
\ref{nml:com:mod:prop:box} नुसार
$\Box\lnot \varphi \notin \Gamma$
असेल तेव्हा आणि तेव्हाच
$\Box^{-1}\Gamma \subseteq \Delta$
आणि $\lnot\varphi \notin \Delta$
असा एखादा संपूर्ण $\Sigma$-सुसंगत
$\Delta$ असतो. असा कोणताही~$\Delta$
घ्या. \ref{nml:com:mod:lem:box-iff-diamond}
नुसार $\Box^{-1}\Gamma \subseteq \Delta$
असेल तेव्हा आणि तेव्हाच
$\Diamond\Delta \subseteq \Gamma$.
तसेच
\ref{nml:com:ccs:prop:ccs-properties}\ref{nml:com:ccs:prop:ccs-lnot}
नुसार $\lnot \varphi \notin \Delta$
असेल तेव्हा आणि तेव्हाच
$\varphi \in \Delta$. म्हणून
$\Diamond\varphi \in \Gamma$
असेल तेव्हा आणि तेव्हाच
$\Diamond\Delta \subseteq \Gamma$
आणि $\varphi \in \Delta$ असा एखादा
संपूर्ण $\Sigma$-सुसंगत $\Delta$
असतो.
\end{proof}



\begin{prob}
  $\Gamma$ हा संपूर्ण $\Sigma$-सुसंगत
  असेल, तर पुढील विधान दाखवा:
  $\Diamond \varphi \in \Gamma$
    असेल तेव्हा आणि तेव्हाच
    $\Box^{-1}\Gamma \subseteq \Delta$
    आणि $\varphi \in \Delta$ असा संपूर्ण
    $\Sigma$-सुसंगत $\Delta$ असतो.
  \emph{हे
    \ref{nml:com:mod:lem:box-iff-diamond}
    न वापरता करा}.
\end{prob}













\section{कॅनॉनिकल प्रतिमाने}\label{nml:com:cmd:sec}

सुसंगत मोडल प्रणाली~$\Sigma$ साठीचे \emph{कॅनॉनिकल प्रतिमान}
हे विशिष्ट प्रतिमान~$\mModel{M^\Sigma}$ आहे; त्याची
जगे म्हणजे सर्व संपूर्ण $\Sigma$-सुसंगत संच.
त्याचा प्राप्यता संबंध~$R^\Sigma$ आणि
मूल्यांकन~$V^\Sigma$ यांची व्याख्या अशी केली
जाते की जग~$\Delta$ येथे सत्य असलेली
सूत्रे म्हणजे नेमकी~$\Delta$ मध्ये
असलेली सूत्रे असतात.

\begin{defn}
  $\Sigma$ हे सुसंगत सामान्य मोडल तर्कशास्त्र असू द्या,
  म्हणजे $\Sigma \Proves/ \lfalse$.
  $\Sigma$ साठीचे \emph{कॅनॉनिकल प्रतिमान}
  म्हणजे $\mModel{M}^\Sigma = \tuple{W^\Sigma, R^\Sigma,
    V^\Sigma}$, जेथे:
  \begin{enumerate}
  \item $W^\Sigma = \Setabs{ \Delta }{\Delta \text{
      हा संपूर्ण $\Sigma$-सुसंगत आहे} }$.
  \item $R^\Sigma \Delta\Delta'$ असेल तेव्हा आणि
    तेव्हाच $\Box^{-1}\Delta \subseteq
      \Delta'$.
  \item $V^\Sigma(p) = \Setabs{\Delta \in W^\Sigma}{p \in \Delta}$.
  \end{enumerate}
  सुसंगततेमुळे रिकामा संच $\Sigma$-सुसंगत आहे.
  \ref{nml:com:lin:thm:lindenbaum} नुसार त्याचा संपूर्ण
  $\Sigma$-सुसंगत विस्तार मिळतो; म्हणून $W^\Sigma$ रिकामा नाही.
\end{defn}














\section{सत्यता पूर्वप्रमेय}\label{nml:com:tru:sec}

कॅनॉनिकल प्रतिमान~$\mModel{M^\Sigma}$ ची व्याख्या
अशी केली आहे की $\mSat{M^\Sigma}{\varphi}[\Delta]$
असेल तेव्हा आणि तेव्हाच $\varphi \in \Delta$.
विधानीय चले साठी
$V^\Sigma$ ची व्याख्याच हे थेट दाखवते.
मात्र ही समतुल्यता सर्व सूत्रे साठी
खरी आहे याची पडताळणी करावी लागेल.
हे आपण विगमनाने करू. विगमन-पायरीत
विधानीय कारक असलेल्या सूत्रे साठी
समतुल्यता सिद्ध करावी लागते (येथे
\ref{nml:com:ccs:prop:ccs-properties} वापरतो),
तसेच मोडल कारक असलेल्या सूत्रांसाठीही
(येथे \ref{nml:com:mod:sec} मधील निष्कर्ष वापरतो).

\begin{prop}[सत्यता पूर्वप्रमेय]
  \label{nml:com:tru:prop:truthlemma}
  $\Sigma$ हे सुसंगत सामान्य मोडल तर्कशास्त्र असू द्या.
  प्रत्येक $\Delta \in W^\Sigma$ आणि प्रत्येक सूत्र~$\varphi$ साठी
  $\mSat{M^\Sigma}{\varphi}[\Delta]$ असेल तेव्हा
  आणि तेव्हाच $\varphi \in \Delta$.
\end{prop}

\begin{proof}
  $\varphi$ च्या रचनेवर विगमन करू.
  \begin{enumerate}
    \item \indcase{\varphi}{\lfalse}{$\mSat/{M^\Sigma}{\lfalse}[\Delta]$
        हे \ref{nml:syn:trw:defn:mmodels} नुसार, आणि
        $\lfalse \notin \Delta$ हे
        \ref{nml:com:ccs:prop:ccs-properties}\ref{nml:com:ccs:prop:ccs-lfalse}
        नुसार.}

    \item \indcase{\varphi}{\ltrue}{$\mSat{M^\Sigma}{\ltrue}[\Delta]$
        हे \ref{nml:syn:trw:defn:mmodels} नुसार, आणि
        $\ltrue \in \Delta$ हे
        \ref{nml:com:ccs:prop:ccs-properties}\ref{nml:com:ccs:prop:ccs-ltrue}
        नुसार.}

    \item \indcase{\varphi}{p}{$\mSat{M^\Sigma}{p}[\Delta]$
      असेल तेव्हा आणि तेव्हाच $\Delta \in
      V^\Sigma(p)$, हे \ref{nml:syn:trw:defn:mmodels}
      नुसार. तसेच $\Delta
      \in V^\Sigma(p)$ असेल तेव्हा आणि तेव्हाच
      $p \in \Delta$, हे~$V^\Sigma$ च्या
      व्याख्येनुसार.}

    \item \indcase{\varphi}{\lnot \psi}{
        $\mSat{M^\Sigma}{\lnot \psi}[\Delta]$
          असेल तेव्हा आणि तेव्हाच
          $\mSat/{M^\Sigma}{\psi}[\Delta]$
          (\ref{nml:syn:trw:defn:mmodels} नुसार);
          असेल तेव्हा आणि तेव्हाच $\psi \notin \Delta$
          (विगमन गृहीतकानुसार);
          असेल तेव्हा आणि तेव्हाच
          $\lnot \psi \in \Delta$
          (\ref{nml:com:ccs:prop:ccs-properties}\ref{nml:com:ccs:prop:ccs-lnot}
          नुसार).}

    \item \indcase{\varphi}{\psi \land \chi}{
        $\mSat{M^\Sigma}{\psi \land
            \chi}[\Delta]$ असेल तेव्हा आणि तेव्हाच
          $\mSat{M^\Sigma}{\psi}[\Delta]$ आणि
          $\mSat{M^\Sigma}{\chi}[\Delta]$
          (\ref{nml:syn:trw:defn:mmodels} नुसार);
          असेल तेव्हा आणि तेव्हाच $\psi \in \Delta$
          आणि $\chi \in \Delta$
          (विगमन गृहीतकानुसार);
          असेल तेव्हा आणि तेव्हाच $\psi \land \chi \in
          \Delta$
          (\ref{nml:com:ccs:prop:ccs-properties}\ref{nml:com:ccs:prop:ccs-land}
          नुसार).}

    \item \indcase{\varphi}{\psi \lor \chi}{
        $\mSat{M^\Sigma}{\psi \lor
            \chi}[\Delta]$ असेल तेव्हा आणि तेव्हाच
          $\mSat{M^\Sigma}{\psi}[\Delta]$ किंवा
          $\mSat{M^\Sigma}{\chi}[\Delta]$
          (\ref{nml:syn:trw:defn:mmodels} नुसार);
          असेल तेव्हा आणि तेव्हाच $\psi \in \Delta$
          किंवा $\chi \in \Delta$
          (विगमन गृहीतकानुसार);
          असेल तेव्हा आणि तेव्हाच $\psi \lor \chi \in
          \Delta$
          (\ref{nml:com:ccs:prop:ccs-properties}\ref{nml:com:ccs:prop:ccs-lor}
          नुसार).}

    \item \indcase{\varphi}{\psi \lif \chi}{
        $\mSat{M^\Sigma}{\psi \lif
            \chi}[\Delta]$ असेल तेव्हा आणि तेव्हाच
          $\mSat/{M^\Sigma}{\psi}[\Delta]$ किंवा
          $\mSat{M^\Sigma}{\chi}[\Delta]$
          (\ref{nml:syn:trw:defn:mmodels} नुसार);
          असेल तेव्हा आणि तेव्हाच $\psi \notin \Delta$
          किंवा $\chi \in \Delta$
          (विगमन गृहीतकानुसार);
          असेल तेव्हा आणि तेव्हाच $\psi \lif \chi
          \in \Delta$
          (\ref{nml:com:ccs:prop:ccs-properties}\ref{nml:com:ccs:prop:ccs-lif}
          नुसार).}

    \item \indcase{\varphi}{\psi \liff \chi}{
        $\mSat{M^\Sigma}{\psi \liff
            \chi}[\Delta]$ असेल तेव्हा आणि तेव्हाच
          $\mSat{M^\Sigma}{\psi}[\Delta]$ आणि
          $\mSat{M^\Sigma}{\chi}[\Delta]$
          ही दोन्ही, किंवा
          $\mSat/{M^\Sigma}{\psi}[\Delta]$ आणि
          $\mSat/{M^\Sigma}{\chi}[\Delta]$
          ही दोन्ही
          (\ref{nml:syn:trw:defn:mmodels} नुसार);
          असेल तेव्हा आणि तेव्हाच $\psi \in \Delta$
          आणि $\chi \in \Delta$ ही दोन्ही, किंवा
          $\psi \notin \Delta$ आणि $\chi \notin
          \Delta$ ही दोन्ही
          (विगमन गृहीतकानुसार);
          असेल तेव्हा आणि तेव्हाच $\psi \liff \chi \in
          \Delta$
          (\ref{nml:com:ccs:prop:ccs-properties}\ref{nml:com:ccs:prop:ccs-liff}
          नुसार).}

    \item \indcase{\varphi}{\Box \psi}{
        प्रथम
          $\mSat{M^\Sigma}{\Box \psi}[\Delta]$
          असे समजा.
          \ref{nml:syn:trw:defn:mmodels} नुसार
          $R^\Sigma \Delta\Delta'$ असलेल्या
          प्रत्येक $\Delta'$ साठी
          $\mSat{M^\Sigma}{\psi}[\Delta']$.
          विगमन गृहीतकानुसार
          $R^\Sigma \Delta\Delta'$ असलेल्या
          प्रत्येक $\Delta'$ साठी $\psi \in
          \Delta'$. $R^\Sigma$ च्या
          व्याख्येनुसार $\Box^{-1} \Delta
          \subseteq \Delta'$ असलेल्या प्रत्येक
          $\Delta'$ साठी $\psi \in \Delta'$.
          \ref{nml:com:mod:prop:box} नुसार
          $\Box \psi \in \Delta$.

          आता $\Box \psi \in \Delta$ असे समजा.
          $R^\Sigma \Delta\Delta'$ म्हणजेच
          $\Box^{-1}\Delta \subseteq \Delta'$
          असलेला कोणताही $\Delta' \in W^\Sigma$
          घ्या. $\Box \psi \in \Delta$
          असल्याने $\psi \in \Box^{-1} \Delta$.
          परिणामी $\psi \in \Delta'$.
          विगमन गृहीतकानुसार
          $\mSat{M^\Sigma}{\psi}[\Delta']$.
          $R^\Sigma \Delta\Delta'$ असलेला
          $\Delta'$ स्वेच्छ असल्याने
          $R^\Sigma \Delta\Delta'$ असलेल्या
          सर्व $\Delta' \in W^\Sigma$ साठी
          $\mSat{M^\Sigma}{\psi}[\Delta']$.
          \ref{nml:syn:trw:defn:mmodels} नुसार
          $\mSat{M^\Sigma}{\Box
            \psi}[\Delta]$.}

    \item \indcase{\varphi}{\Diamond \psi}{
        प्रथम
          $\mSat{M^\Sigma}{\Diamond \psi}[\Delta]$
          असे समजा.
          \ref{nml:syn:trw:defn:mmodels} नुसार
          $R^\Sigma \Delta\Delta'$ असलेल्या
          एखाद्या $\Delta'$ साठी
          $\mSat{M^\Sigma}{\psi}[\Delta']$.
          विगमन गृहीतकानुसार
          $R^\Sigma \Delta\Delta'$ असलेल्या
          एखाद्या $\Delta'$ साठी
          $\psi \in \Delta'$.
          $R^\Sigma$ च्या व्याख्येनुसार,
          $\Box^{-1} \Delta
            \subseteq \Delta'$ असलेल्या
            एखाद्या $\Delta'$ साठी
            $\psi \in \Delta'$.
            \ref{nml:com:mod:lem:box-iff-diamond} नुसार,
          $\Diamond \Delta' \subseteq \Delta$
          असलेल्या एखाद्या $\Delta'$ साठी
          $\psi \in \Delta'$.
          $\psi \in \Delta'$ असल्याने
          $\Diamond \psi \in \Diamond
          \Delta'$; म्हणून
          $\Diamond \psi \in \Delta$.

          आता $\Diamond \psi \in \Delta$
          असे समजा. \ref{nml:com:mod:prop:diamond}
          नुसार $\Diamond \Delta' \subseteq
          \Delta$ आणि $\psi \in \Delta'$
          असलेला संपूर्ण $\Sigma$-सुसंगत
          $\Delta' \in W^\Sigma$ असतो.
          \ref{nml:com:mod:lem:box-iff-diamond}
            नुसार $\Box^{-1}
            \Delta \subseteq \Delta'$
            आणि $\psi \in \Delta'$ असलेला
            $\Delta' \in W^\Sigma$ असतो.
          $R^\Sigma$ च्या व्याख्येनुसार
          $R^\Sigma \Delta\Delta'$;
          म्हणून $R^\Sigma
          \Delta\Delta'$ आणि
          $\psi \in \Delta'$ असलेला
          $\Delta' \in W^\Sigma$ असतो.
          विगमन गृहीतकानुसार $\mSat{M^\Sigma}{\psi}[\Delta']$.
          \ref{nml:syn:trw:defn:mmodels}
          नुसार $\mSat{M^\Sigma}{\Diamond
            \psi}[\Delta]$.}
  \end{enumerate}
\end{proof}















\section{\Log{K} चे निर्धारण आणि संपूर्णता}\label{nml:com:cmk:sec}

आता संपूर्णता सिद्ध करण्यासाठी कॅनॉनिकल प्रतिमान
वापरता येईल. $\Sigma$ साठीच्या कॅनॉनिकल
प्रतिमानात सत्य असलेली सूत्रे म्हणजे
नेमकी $\Sigma$-निष्पन्न करता येण्याजोगे सूत्रे असतात;
या तथ्यावरून संपूर्णता मिळते. हा गुणधर्म
असलेले प्रतिमान~$\Sigma$ चे
\emph{निर्धारण करते} असे म्हणतात.

\begin{defn}
  सर्व सूत्रांपैकी प्रत्येक~$\varphi$ साठी
  $\mSat{M}{\varphi}$ असेल तेव्हा आणि तेव्हाच
  $\Sigma \Proves \varphi$ असेल, तर
  प्रतिमान~$\mModel{M}$ हे सामान्य मोडल
  तर्कशास्त्र~$\Sigma$ चे
  \emph{निर्धारण करते}.
\end{defn}

\begin{thm}[निर्धारण]\label{nml:com:cmk:thm:determination}
   $\Sigma$ हे सुसंगत सामान्य मोडल तर्कशास्त्र असू द्या. मग
   $\mSat{M^\Sigma}{\varphi}$ असेल तेव्हा
   आणि तेव्हाच $\Sigma \Proves \varphi$.
\end{thm}

\begin{proof}
  $\mSat{M^\Sigma}{\varphi}$ असेल, तर प्रत्येक
  संपूर्ण $\Sigma$-सुसंगत $\Delta$ साठी
  $\mSat{M^\Sigma}{\varphi}[\Delta]$.
  म्हणून सत्यता पूर्वप्रमेयानुसार प्रत्येक
  संपूर्ण $\Sigma$-सुसंगत $\Delta$ साठी
  $\varphi \in \Delta$. त्यामुळे
  \ref{nml:com:lin:cor:provability-characterization}
  वरून ($\Gamma = \emptyset$ घेऊन)
  $\Sigma \Proves \varphi$.

  उलट $\Sigma \Proves \varphi$ असेल, तर
  \ref{nml:com:ccs:prop:ccs-properties}\ref{nml:com:ccs:prop:ccs-closed}
  नुसार प्रत्येक संपूर्ण $\Sigma$-सुसंगत
  $\Delta$ मध्ये~$\varphi$ असते. म्हणून
  सत्यता पूर्वप्रमेयानुसार प्रत्येक~$\Delta \in
  W^\Sigma$ साठी
  $\mSat{M^\Sigma}{\varphi}[\Delta]$;
  म्हणजे $\mSat{M^\Sigma}{\varphi}$.
\end{proof}

निर्दोषतेनुसार \Log{K} सुसंगत आहे, कारण असत्य सूत्र वैध नाही.
\Log{K} साठीचे कॅनॉनिकल प्रतिमान~\Log{K}
चे निर्धारण करत असल्याने
\Log{K} ची संपूर्णता लगेच निष्पन्नतेने मिळते:

\begin{cor}\label{nml:com:cmk:cor:Kcomplete}
  मूलभूत मोडल तर्कशास्त्र~\Log{K}
सर्व प्रतिमानांच्या वर्गाच्या सापेक्ष
संपूर्ण आहे; म्हणजे
$\Entails \varphi$ असेल, तर $\Log{K}
\Proves \varphi$.
\end{cor}

\begin{proof}
  प्रतिपरिवर्तनाने: $\Log{K} \Proves/ \varphi$
  असेल, तर निर्धारणावरून
  $\mSat/{M^{\Log{K}}}{\varphi}$;
  म्हणून~$\varphi$ वैध नाही.
\end{proof}

सर्वसाधारणपणे एखादी प्रणाली~$\Sigma$
प्रतिमानांच्या एखाद्या वर्गाच्या सापेक्ष
संपूर्ण आहे हे सिद्ध करण्यासाठी
केवळ निर्धारण पुरेसे नसते. उदाहरणार्थ,
\Log{KTB4} ची परावर्ती, सममित आणि
संक्रमक प्रतिमानांच्या वर्गाच्या सापेक्ष
संपूर्णता विचारात घ्या. प्रणाली~$\Sigma$
साठीच्या कॅनॉनिकल प्रतिमानातून ही संपूर्णता सिद्ध करायची असेल,
तर ते प्रतिमान त्या वर्गातील सदस्य आहे हेही दाखवावे लागते.
कॅनॉनिकल प्रतिमानाच्या रचनेवरून हे
स्पष्टपणे मिळत नाही---काही वेळा
तर ते खरेही नसते!













\section{चौकट-संपूर्णता}\label{nml:com:fra:sec}

दिलेल्या तर्कशास्त्राच्या कॅनॉनिकल प्रतिमानाला
संबंधित चौकट-गुणधर्म आहे हे दाखवल्यावर
\Log{K} साठीचे संपूर्णता प्रमेय इतर
मोडल प्रणालींनाही लागू करता येते.

\begin{thm}\label{nml:com:fra:thm:completeframeprops}
  सुसंगत सामान्य मोडल तर्कशास्त्र~$\Sigma$ मध्ये
  \ref{nml:com:fra:tab:correspondencetable} च्या
  डाव्या बाजूच्या सूत्रांपैकी एखादे असेल,
  तर~$\Sigma$ साठीच्या कॅनॉनिकल प्रतिमानाला
  उजव्या बाजूचा संबंधित गुणधर्म असतो.
\end{thm}

\begin{table}[htp]
  \centering
    \begin{tabular}{| l || l |}
      \hline
      \emph{$\Sigma$ मध्ये \dots असेल, तर} & \emph{$\Sigma$ साठीचे
        कॅनॉनिकल प्रतिमान \dots असते:} \\
      \hline \hline
      \Ax{D}:\;\; $\Box\varphi \lif \Diamond \varphi$ & \quad सर्वत्र उत्तरवर्ती; \\
      \hline
      \Ax{T}:\;\; $\Box\varphi \lif \varphi$ &\quad  परावर्ती;\\
      \hline
      \Ax{B}: \;\;$\varphi \lif \Box\Diamond\varphi$ &\quad  सममित; \\
      \hline
      \Ax{4}: \;\; $\Box\varphi \lif \Box\Box\varphi$ & \quad संक्रमक; \\
      \hline
      \Ax{5}: \;\; $\Diamond \varphi \lif \Box\Diamond\varphi$& \quad युक्लिडी.\\
      \hline
    \end{tabular}
    \caption{मूलभूत अनुरूपता तथ्ये.}
    \label{nml:com:fra:tab:correspondencetable}
\end{table}

\begin{proof}
  यांपैकी प्रत्येक प्रकरण स्वतंत्रपणे पाहू.

  $\Sigma$ मध्ये \Ax{D} आहे असे समजा आणि
  $\Delta \in W^\Sigma$ घ्या.
  $R^\Sigma \Delta\Delta'$ असा
  $\Delta'$ अस्तित्वात आहे हे दाखवायचे आहे.
  त्यासाठी $\Box^{-1}\Delta$ हा
  $\Sigma$-सुसंगत आहे हे दाखवणे पुरेसे;
  कारण मग लिंडेनबाउमच्या पूर्वप्रमेयानुसार
  त्याचा विस्तार असलेला संपूर्ण
  $\Sigma$-सुसंगत संच
  $\Delta' \supseteq \Box^{-1}\Delta$
  असतो आणि $R^\Sigma$ च्या व्याख्येनुसार

  $R^\Sigma \Delta\Delta'$.
  म्हणून व्याघातासाठी $\Box^{-1}\Delta$
  हा $\Sigma$-सुसंगत \emph{नाही}
  असे समजा; म्हणजे
  $\Box^{-1}\Delta \Proves[\Sigma] \lfalse$.
  \ref{nml:com:mod:lem:box2} नुसार
  $\Delta \Proves[\Sigma] \Box\lfalse$;
  आणि $\Sigma$ मध्ये \Ax{D} असल्याने
  $\Delta \Proves[\Sigma] \Diamond\lfalse$
  हेही मिळते. पण $\Sigma$ सामान्य असल्याने
  $\Sigma \Proves \lnot\Diamond \lfalse$
  (\ref{nml:prf:nor:prop:notDiamondBot});
  म्हणून $\Delta
  \Proves[\Sigma] \lnot\Diamond \lfalse$
  हेही मिळते. हा~$\Delta$ च्या
  सुसंगततेशी व्याघात आहे.

  आता $\Sigma$ मध्ये \Ax{T} आहे
  असे समजा आणि $\Delta \in W^\Sigma$
  घ्या. $R^\Sigma \Delta\Delta$
  दाखवायचे आहे; म्हणजेच

  $\Box^{-1}\Delta \subseteq \Delta$.
  पण $\Box\varphi \in \Delta$ असेल,
  तर \Ax{T} नुसार $\varphi \in \Delta$;
  हेच अपेक्षित होते.

  आता $\Sigma$ मध्ये \Ax{B} आहे
  असे समजा आणि $\Delta$, $\Delta' \in
  W^\Sigma$ साठी $R^\Sigma
  \Delta\Delta'$ असे धरा.
  $R^\Sigma \Delta'\Delta$
  दाखवायचे आहे; म्हणजे
  $\Box^{-1}\Delta'
    \subseteq \Delta$.
    \ref{nml:com:mod:lem:box-iff-diamond} नुसार
    ही अट पुढील अटीशी समतुल्य आहे:
  $\Diamond\Delta \subseteq \Delta'$.
  म्हणून $\varphi \in \Delta$ असे
  समजा. \Ax{B} नुसार
  $\Box\Diamond\varphi \in \Delta$.
  $R^\Sigma \Delta\Delta'$ या
  गृहीतकावरून
  $\Box^{-1}\Delta \subseteq \Delta'$;
  म्हणून आवश्यकतेनुसार
  $\Diamond\varphi \in \Delta'$.

  आता $\Sigma$ मध्ये \Ax{4} आहे असे
  समजा आणि $R^\Sigma \Delta_1\Delta_2$
  व $R^\Sigma \Delta_2\Delta_3$
  असे धरा. $R^\Sigma
  \Delta_1\Delta_3$ दाखवायचे आहे.
  गृहीतकावरून
  $\Box^{-1}\Delta_1 \subseteq \Delta_2$
  आणि $\Box^{-1}\Delta_2 \subseteq \Delta_3$
  या दोन्ही अटी मिळतात.
  $R^\Sigma \Delta_1\Delta_3$
  साठी $\Box^{-1}\Delta_1
  \subseteq \Delta_3$ दाखवणे पुरेसे.
  म्हणून $\psi \in \Box^{-1}\Delta_1$
  घ्या; म्हणजे $\Box\psi \in \Delta_1$.
  \Ax{4} नुसार $\Box\Box\psi \in
  \Delta_1$ हेही मिळते. गृहीतकांवरून
  आधी $\Box\psi \in \Delta_2$
  आणि नंतर $\psi \in \Delta_3$,
  हेच अपेक्षित होते.

  आता $\Sigma$ मध्ये \Ax{5} आहे
  असे समजा आणि $R^\Sigma
  \Delta_1\Delta_2$ व
  $R^\Sigma \Delta_1\Delta_3$
  असे धरा. $R^\Sigma
  \Delta_2\Delta_3$ दाखवायचे आहे.
  पहिल्या गृहीतकावरून
  $\Box^{-1}\Delta_1 \subseteq \Delta_2$

  मिळते; दुसरे गृहीतक
  $\Diamond\Delta_3 \subseteq \Delta_1$
  याच्याशी समतुल्य आहे
  \ref{nml:com:mod:lem:box-iff-diamond} नुसार.
  $R^\Sigma \Delta_2\Delta_3$
  दाखवण्यासाठी
    \ref{nml:com:mod:lem:box-iff-diamond} नुसार,
    पुढील अट दाखवणे पुरेसे:
  $\Diamond\Delta_3 \subseteq \Delta_2$.
  म्हणून $\Diamond\varphi \in
  \Diamond\Delta_3$ घ्या;
  म्हणजे $\varphi \in \Delta_3$.
  दुसऱ्या गृहीतकानुसार
  $\Diamond\varphi \in \Delta_1$;
  \Ax{5} नुसार
  $\Box\Diamond\varphi \in \Delta_1$
  हेही मिळते. आता पहिल्या
  गृहीतकावरून $\Diamond\varphi \in
  \Delta_2$, हेच अपेक्षित होते.
  \end{proof}

या निष्कर्षांवरून अनेक प्रणालींसाठी
संपूर्णतेची निष्पन्नता मिळते.
उदाहरणार्थ, $\Log{S5} = \Log{KT5} =
\Log{KTB4}$ हे सर्व परावर्ती युक्लिडी
प्रतिमानांच्या वर्गाच्या सापेक्ष संपूर्ण
आहे, हे आपल्याला माहीत आहे.
हा वर्ग सर्व परावर्ती, सममित आणि
संक्रमक प्रतिमानांच्या वर्गाइतकाच आहे.

\begin{thm}\label{nml:com:fra:thm:generaldet}
  $\mClass{C}_\Ax{D}$,
  $\mClass{C}_\Ax{T}$,
  $\mClass{C}_\Ax{B}$,
  $\mClass{C}_\Ax{4}$ आणि
  $\mClass{C}_\Ax{5}$
  हे अनुक्रमे सर्व सर्वत्र उत्तरवर्ती,
  परावर्ती, सममित, संक्रमक आणि
  युक्लिडी प्रतिमानांचे वर्ग असू द्या.
  मग \Ax{D}, \Ax{T}, \Ax{B},
  \Ax{4} आणि \Ax{5} यांमधील
  कोणत्याही योजना $\varphi_1$, \dots, $\varphi_n$
  साठी प्रणाली
  $\Log{K}\varphi_1 \dots \varphi_n$
  हिचे निर्धारण प्रतिमानांचा वर्ग
  $\mClass{C} = \mClass{C}_{\varphi_1} \cap \dots
  \cap \mClass{C}_{\varphi_n}$ करतो.
\end{thm}

\begin{prop}
  $\Sigma$ हे सुसंगत सामान्य मोडल तर्कशास्त्र
  असू द्या. मग:
  \begin{enumerate}
  \item\label{nml:com:fra:prop:anotherfive-a}
    $\Sigma$ मध्ये $\Diamond\varphi \lif \Box
    \varphi$ ही योजना असेल, तर~$\Sigma$
    साठीचे कॅनॉनिकल प्रतिमान
    आंशिक फलनीय असते.
  \item $\Sigma$ मध्ये
    $\Diamond\varphi \liff \Box
    \varphi$ ही योजना असेल, तर~$\Sigma$
    साठीचे कॅनॉनिकल प्रतिमान
    फलनीय असते.
  \item $\Sigma$ मध्ये
    $\Box\Box\varphi \lif \Box
    \varphi$ ही योजना असेल, तर~$\Sigma$
    साठीचे कॅनॉनिकल प्रतिमान
    दुर्बल सघन असते.
  \end{enumerate}
  (या चौकट-गुणधर्मांच्या व्याख्यांसाठी
  \ref{nml:frd:acc:tab:anotherfive} पाहा.)
\end{prop}

\begin{proof}
  \begin{enumerate}
  \item $\Sigma$ मध्ये
    $\Diamond \varphi \lif \Box \varphi$
    ही योजना आहे असे समजा.
    $R^\Sigma$ आंशिक फलनीय आहे हे
    दाखवण्यासाठी कोणत्याही
    $\Delta_1$, $\Delta_2$,
    $\Delta_3 \in W^\Sigma$
    साठी $R^\Sigma \Delta_1\Delta_2$
    आणि $R^\Sigma \Delta_1\Delta_3$
    असतील, तर $\Delta_2=\Delta_3$
    हे सिद्ध करायचे आहे.
    $R^\Sigma \Delta_1\Delta_2$
    असल्याने $\Box^{-1}\Delta_1
    \subseteq \Delta_2$;
    आणि $R^\Sigma \Delta_1\Delta_3$
    असल्याने $\Box^{-1}\Delta_1
    \subseteq \Delta_3$
    .
    $\Delta_2=\Delta_3$ साठी
    $\Delta_2 \subseteq \Delta_3$
    आणि $\Delta_3 \subseteq \Delta_2$
    हे दोन्ही समावेश सिद्ध करणे पुरेसे.
    पहिल्या समावेशासाठी
    $\varphi \in \Delta_2$ घ्या;
    मग $\Diamond\varphi \in \Delta_1$.
    या योजनेवरून आणि~$\Delta_1$ च्या
    निगामी बंदतेवरून
    $\Box\varphi \in \Delta_1$ हेही मिळते.
    $R^\Sigma \Delta_1\Delta_3$
    या गृहीतकावरून $\varphi \in \Delta_3$.
    दुसरा समावेशही तसाच सिद्ध होतो.

  \item $\Diamond\varphi \liff \Box\varphi$ या योजनेतून
    $\Box\varphi \lif \Diamond\varphi$, म्हणजे \Ax{D}, मिळते.
    म्हणून हे \ref{nml:com:fra:prop:anotherfive-a}
    आणि \ref{nml:com:fra:thm:completeframeprops}
    मधील सर्वत्र उत्तरवर्तित्वाच्या
    सिद्धतेवरून लगेच मिळते.

  \item $\Sigma$ मध्ये
    $\Box\Box\varphi \lif \Box\varphi$
    ही योजना आहे असे समजा.
    $R^\Sigma$ दुर्बल सघन आहे हे
    दाखवण्यासाठी $R^\Sigma
    \Delta_1\Delta_2$ घ्या.
    $R^\Sigma \Delta_1\Delta_3$
    आणि $R^\Sigma \Delta_3\Delta_2$
    असा संपूर्ण $\Sigma$-सुसंगत
    संच~$\Delta_3$ अस्तित्वात आहे
    हे दाखवायचे आहे. असे ठेवू:
    \[    \Gamma = \Box^{-1}\Delta_1 \cup \Diamond\Delta_2.
    \]
    $\Gamma$ हा $\Sigma$-सुसंगत
    आहे हे दाखवणे पुरेसे. मग
    लिंडेनबाउमच्या पूर्वप्रमेयानुसार
    त्याचा विस्तार असा संपूर्ण
    $\Sigma$-सुसंगत संच~$\Delta_3$
    मिळतो की $\Box^{-1}\Delta_1
    \subseteq \Delta_3$ आणि
    $\Diamond\Delta_2 \subseteq \Delta_3$.
    म्हणजे $R^\Sigma \Delta_1\Delta_3$

    आणि $R^\Sigma \Delta_3\Delta_2$
    (
      \ref{nml:com:mod:lem:box-iff-diamond} नुसार).

    व्याघातासाठी $\Gamma$ सुसंगत
    नाही असे समजा. $\Box^{-1}\Delta_1\subseteq\Delta_2$
    आणि $\Delta_2$ सुसंगत असल्याने फक्त पहिल्या संचातील
    गृहीतके वापरून व्याघात मिळू शकत नाही. म्हणून
    पुढील finite witness मध्ये $n\geq0$ आणि $m\geq1$ निवडता येतात.
    मग $\Box\varphi_1$, \dots,
    $\Box\varphi_n \in \Delta_1$
    आणि $\psi_1$, \dots,~$\psi_m \in \Delta_2$
    अशी सूत्रे असतात की
    \[\varphi_1, \dots,
    \varphi_n, \Diamond\psi_1, \dots, \Diamond\psi_m \Proves[\Sigma]
    \lfalse.\]
    $\Diamond (\psi_1 \land \dots \land \psi_m) \to
    (\Diamond\psi_1 \land \dots \land \Diamond\psi_m)$
    हे प्रत्येक सामान्य मोडल तर्कशास्त्रात
    निष्पन्न करता येण्याजोगे असल्याने
    पुढीलप्रमाणे युक्तिवाद करून
    आपण~$\Delta_2$ च्या
    सुसंगततेशी व्याघात मिळवतो:
    \begin{align*}
      \varphi_1, \dots, \varphi_n, & \Diamond\psi_1, \dots,
      \Diamond\psi_m \Proves[\Sigma] \lfalse \\
      \varphi_1, \dots,\varphi_n & \Proves[\Sigma] (\Diamond\psi_1 \land
      \dots \land \Diamond\psi_m) \lif \lfalse\\
      & \qquad\text{निगमन प्रमेयानुसार}\\
      & \qquad\text{\ref{nml:prf:prp:prop:derivabilityfacts}\ref{nml:prf:prp:prop:derivabilityfacts-deduction}, आणि \Taut} \\
      \varphi_1, \dots,\varphi_n & \Proves[\Sigma]
      \Diamond(\psi_1 \land
      \dots \land \psi_m) \lif \lfalse\\
      & \qquad \text{$\Sigma$ सामान्य असल्याने} \\
      \varphi_1, \dots,\varphi_n & \Proves[\Sigma]
      \lnot\Diamond (\psi_1 \land \dots \land \psi_m)\\
      & \qquad\text{\PL नुसार} \\
      \varphi_1, \dots,\varphi_n & \Proves[\Sigma]
      \Box\lnot (\psi_1 \land \dots \land \psi_m)\\
      & \qquad\text{$\lnot\Diamond$ साठी $\Box\lnot$} \\
      \Box\varphi_1, \dots,\Box\varphi_n
      & \Proves[\Sigma] \Box\Box \lnot (\psi_1 \land \dots \land \psi_m)\\
      & \qquad\text{\ref{nml:com:mod:lem:box1} नुसार} \\
      \Box\varphi_1, \dots,\Box\varphi_n
      & \Proves[\Sigma]
      \Box\lnot (\psi_1 \land \dots \land \psi_m)\\
      &\qquad\text{$\Box\Box\varphi \lif \Box\varphi$ योजनेनुसार} \\
      \Delta_1 & \Proves[\Sigma]
      \Box\lnot (\psi_1 \land \dots \land \psi_m)\\
      &\qquad\text{एकस्वनिकतेनुसार, \ref{nml:prf:prp:prop:derivabilityfacts}\ref{nml:prf:prp:prop:derivabilityfacts-monotonicity}} \\
      & \Box\lnot (\psi_1 \land \dots \land \psi_m) \in \Delta_1\\
      &\qquad\text{निगामी बंदतेनुसार}; \\
      & \lnot (\psi_1 \land \dots \land \psi_m) \in \Delta_2\\
      &\qquad \text{कारण }
      R^\Sigma \Delta_1\Delta_2.
    \end{align*}
  \end{enumerate}
\end{proof}

या उदाहरणांवरून प्रत्येक मोडल
प्रणाली~$\Sigma$ \emph{संपूर्ण}
असेल असे वाटू शकते: म्हणजे~$\Sigma$
ची सर्व प्रमेये ज्या प्रत्येक चौकटीत
वैध आहेत, त्या प्रत्येक चौकटीत
वैध असलेले प्रत्येक सूत्रही
ती प्रणाली सिद्ध करत असेल.
पण दुर्दैवाने या अर्थाने संपूर्ण
नसलेल्या अनेक प्रणाली आहेत.



\chapter{निस्यंदने आणि निर्णेयता}\label{nml:fil:chap}










\section{प्रस्तावना}\label{nml:fil:int:sec}

तर्कशास्त्राविषयी एक महत्त्वाचा प्रश्न नेहमीच
असा असतो: ते निर्णेय आहे का? म्हणजे,
“हे सूत्र वैध आहे का?” या प्रश्नाचे
उत्तर देणारी प्रभावी पद्धत आहे का?
विधानीय तर्कशास्त्र निर्णेय आहे:
सत्यता-कोष्टक रचून एखादे
सूत्र सर्वतः-सत्य आहे का हे
आपण प्रभावीपणे तपासू शकतो;
दिलेल्या सूत्र साठी असे
सत्यता-कोष्टक सांत असते.
पण मोडल सूत्र सर्व प्रतिमानांत सत्य
आहे का हे अशाच रीतीने तपासता
येईल असे स्पष्ट नाही, कारण
प्रतिमाने अनंत संख्येने आहेत.
दिलेल्या सूत्राशी संबंधित सर्व
सांत प्रतिमानांची यादी करता येते;
कारण प्रत्यक्ष त्या सूत्रामध्ये
येणाऱ्या विधानीय चले
ना जगांचे कोणते उपसंच नेमले
आहेत, एवढेच महत्त्वाचे असते.
प्राप्यता संबंध निश्चित असेल,
तर $V(p)$ ची संभाव्य वेगवेगळी
नेमणूक म्हणजे~$W$ चे सर्व
उपसंच; $\card{W} = n$ असेल,
तर असे $2^n$ उपसंच असतात.
आपल्या सूत्र~$\varphi$ मध्ये
$m$ वेगवेगळी विधानीय चले असतील, तर निश्चित जगांचा संच आणि निश्चित
प्राप्यता संबंध यांवर या चलांची $2^{nm}$ वेगवेगळी
मूल्यांकने मिळतात; इतर चलांची मूल्यांकने या सूत्रासाठी अप्रस्तुत असतात. प्रत्येक
प्रतिमानात~$\varphi$ सर्व जगांत
सत्य आहे का हे प्रत्येक जगात
$\varphi$ चे सत्यतामूल्य काढून
तपासता येते. अर्थात सर्व
संभाव्य प्राप्यता संबंधही
तपासावे लागतात; पण~$n$
जगांवर असे संबंधही सांत
संख्येनेच असतात (नेमके
$W \times W$ चे उपसंच,
म्हणजे $2^{n^2}$ संबंध).

आपल्याला \Log{K} ऐवजी
प्रतिमानांच्या एखाद्या वर्गाने
परिभाषित तर्कशास्त्रात रस
असेल (उदाहरणार्थ,
परावर्ती संक्रमक प्रतिमानांत),
तर प्राप्यता संबंध योग्य प्रकारचा
आहे का हेही तपासता यायला हवे.
आपल्याला हव्या असलेल्या चौकटी
दिलेल्या सांत मोडल सूत्रांच्या यादीने परिभाष्य
असतील तेव्हा हे करता येते
(उदाहरणार्थ, चौकटीत \Ax{T}
आणि \Ax{4} वैध आहेत का
हे तपासून). म्हणून सर्व सांत
चौकटी क्रमाने घ्याव्यात,
प्रत्येक चौकट इच्छित वर्गात
आहे का हे तपासावे,
त्या चौकटीवरील सर्व संभाव्य
प्रतिमाने नोंदवावीत आणि
प्रत्येकात $\varphi$ सत्य आहे का
हे तपासावे. नसेल, तर थांबावे:
$\varphi$ हा संबंधित प्रतिमानांच्या
वर्गात वैध नाही.

या कल्पनेत अडचण आहे:
पाहणे केव्हा थांबवता येईल,
किंवा कधी थांबवता येईलच का,
हे आपल्याला माहीत नाही.
त्या सूत्राला सांत
प्रतिप्रतिमान असेल, तर आपली
पद्धत ते शोधेल. पण सांत
प्रतिप्रतिमान नसेल, तर
उत्तर मिळणार नाही.
ते सूत्र वैध असू शकते
(म्हणजे कोणतेच प्रतिप्रतिमान
नसते), किंवा त्याला केवळ
अनंत प्रतिप्रतिमान असू शकते,
ज्यापर्यंत आपली तपासणी
कधीच पोहोचणार नाही.
प्रतिप्रतिमान असलेल्या प्रत्येक
सूत्राला सांत प्रतिप्रतिमानही
असते हे दाखवता आले, तर
ही अडचण दूर होऊ शकते.
असे असेल, तर त्या तर्कशास्त्राला
\emph{सांत प्रतिमान गुणधर्म}
आहे असे म्हणतात.

पण तर्कशास्त्राला सांत प्रतिमान
गुणधर्म आहे हे कसे दाखवायचे?
एक मार्ग म्हणजे~$\varphi$ चे
अनंत (प्रति)प्रतिमान सांत
प्रतिमानात रूपांतरित करण्याची
पद्धत शोधणे. तसे करता आले,
तर जेव्हा एखाद्या प्रतिमानात
$\varphi$ सत्य नसते, तेव्हा
त्यातून मिळालेल्या सांत
प्रतिमानातही $\varphi$ सत्य
नसेल. हे सांत प्रतिमान
सर्व सांत प्रतिमानांच्या
यादीत येईल; आणि वैध
नसलेल्या प्रत्येक सूत्राची
अवैधता आपण अखेरीस ठरवू.
मात्र सूत्र वैध असेल, तर
आपली पद्धत थांबणार नाही.
याशिवाय या पद्धतीतून
मिळणाऱ्या सांत प्रतिमानाच्या
आकाराला कमाल मर्यादा आहे
आणि ती मर्यादा केवळ
सूत्र~$\varphi$ पासून प्रभावीपणे काढता येते
हे दाखवता आले, तर
प्रतिप्रतिमान शोधताना
किती आकारापर्यंत सांत
प्रतिमाने तपासायची हे कळेल.
तेवढ्यापर्यंत प्रतिप्रतिमान
सापडले नाही, तर कोणतेच
नाही. अशा वेळी आपली
पद्धत कोणत्याही
सूत्र~$\varphi$ साठी
“$\varphi$ वैध आहे का?” हा
प्रश्न खरोखर निर्णीत करेल.

अनंत रचनांचे सांत रचनांत
रूपांतर करण्याची अनेकदा
यशस्वी होणारी रणनीती म्हणजे
संबंधित बाबतीत सारखे
वागणारे रचनेचे घटक
“एकरूप” मानणे. प्रतिमान~$\mModel{M}$
मध्ये संबंधित बाबतीत सारखे
वागणारी अनंत जगे असतील,
तर ती \emph{सर्व} जगे अशा
जगांच्या सांत संख्येतील
(कदाचित अनंत सदस्य असलेल्या)
“वर्गां”त विभागता येतील.
दुसऱ्या शब्दांत, जगांच्या
संचाचे योग्य प्रकारे विभाजन
करायचे: प्रत्येक वर्गात
कदाचित अनंत जगे असतील,
पण वर्गांची संख्या सांत असेल.
मग नवे प्रतिमान~$\mModel{M^*}$
परिभाषित करू; त्याची जगे
म्हणजे हे वर्ग. जुन्या
प्रतिमानातील सांत संख्येतील
वर्गांमुळे नव्या प्रतिमानात
सांत संख्येने जगे असतील,
म्हणजे सांत प्रतिमान मिळेल.
जग~$w$ ज्या वर्गात आहे
त्याला $[w]$ म्हणू.
जुन्या प्रतिमानात~$\varphi$ चे
प्रतिप्रतिमान असेल, तर नवे
प्रतिमानही प्रतिप्रतिमान
हवे असल्याने
$\mSat{M}{\varphi}[w]$ असेल
तेव्हा आणि तेव्हाच
$\mSat{M^*}{\varphi}[{[w]}]$
असावे. यासाठी विभाजन
आणि~$\mModel{M^*}$ चा
प्राप्यता संबंध योग्य
प्रकारे परिभाषित करावा लागतो.

हे कसे घडेल ते पाहण्यासाठी
प्रथम प्रत्येक जग प्रत्येक
जगाला प्राप्य आहे असे
समजू. मग
$\mSat{M}{\Box \psi}[w]$
असेल तेव्हा आणि तेव्हाच
एखाद्या $v \in W$ साठी
$\mSat{M}{\Box \psi}[v]$;
$\mModel{M^*}$ साठीही
$[w]$ आणि $[v]$ वापरून
तसेच होईल. पहिली कल्पना
म्हणून $P$ हा $\varphi$ मध्ये येणाऱ्या विधानचलांचा सांत संच घेऊ.
दोन जगे~$u$ आणि~$v$
एकाच वर्गात, म्हणजे
समतुल्य, असतील तेव्हा आणि
तेव्हाच प्रतिमान~$\mModel{M}$
मधील $P$ मधील सर्व विधानीय चले
वर त्यांचे सत्यत्व एकसारखे
असेल, अशी अट घेऊ.
म्हणजे प्रत्येक $p\in P$ साठी $\mSat{M}{p}[u]$
असेल तेव्हा आणि तेव्हाच
$\mSat{M}{p}[v]$.
सर्व विधानचलांसाठी $V^*(p) = \Setabs{[w]}{w\in W
  \text{ आणि }\mSat{M}{p}[w]}$ ठेवू. हा वर्गांचा प्रतिमा-संच
सर्व चलांसाठी परिभाषित आहे; मात्र खालील सत्यतेच्या जतनाचा दावा
केवळ $P$ मधील चलांनी बनलेल्या सूत्रांसाठी आहे. आपले ध्येय असे
दाखवणे आहे की
$\mSat{M}{\varphi}[w]$
असेल तेव्हा आणि तेव्हाच
$\mSat{M^*}{\varphi}[{[w]}]$.
हे आपण विगमनानेच सिद्ध
करू. आधार प्रकरण
$\varphi \equiv p$ असेल.
प्रथम $\mSat{M}{p}[w]$
असे समजा. मग व्याख्येनुसार
$[w] \in V^*(p)$;
म्हणून $\mSat{M^*}{p}[{[w]}]$.
आता $\mSat{M^*}{p}[{[w]}]$
असे समजा. याचा अर्थ
$[w] \in V^*(p)$;
म्हणजे~$w$ शी समतुल्य
अशा एखाद्या~$v$ साठी
$\mSat{M}{p}[v]$.
पण “$w$ हे~$v$ शी
समतुल्य आहे” याचा अर्थ
“$w$ आणि~$v$ येथे
नेमकी तीच
$P$ मधील विधानीय चले
सत्य आहेत”; म्हणून
$\mSat{M}{p}[w]$.
आता विगमन-पायरी पाहू;
उदाहरणार्थ $\varphi
\ident \lnot \psi$.
मग $\mSat{M}{\lnot \psi}[w]$
असेल तेव्हा आणि तेव्हाच
$\mSat/{M}{\psi}[w]$;
असेल तेव्हा आणि तेव्हाच
$\mSat/{M^*}{\psi}[{[w]}]$
(विगमन गृहीतकानुसार);
असेल तेव्हा आणि तेव्हाच
$\mSat{M^*}{\lnot \psi}[{[w]}]$.
इतर गैर-मोडल कारकांसाठीही
तसेच. $\Box$ साठीही हे
चालते: $\mSat{M^*}{\Box
  \psi}[{[w]}]$ असे समजा.
याचा अर्थ प्रत्येक $[u]$
साठी $\mSat{M^*}{\psi}[{[u]}]$.
विगमन गृहीतकानुसार प्रत्येक
$u$ साठी $\mSat{M}{\psi}[u]$.
म्हणून $\mSat{M}{\Box \psi}[w]$.
उलट, $\mSat{M}{\Box\psi}[w]$ असेल, तर प्रत्येक $u\in W$ साठी
$\mSat{M}{\psi}[u]$. विगमन गृहीतकानुसार प्रत्येक वर्ग $[u]$ येथे
$\mSat{M^*}{\psi}[{[u]}]$; म्हणून $\mSat{M^*}{\Box\psi}[{[w]}]$.
दोन्ही प्रतिमानांतील प्राप्यता संबंध येथे वैश्विक आहे.

सर्वसाधारण प्रकरणात
$\mModel{M^*}$ साठी
प्राप्यता संबंधही परिभाषित
करावा लागतो; तेव्हा
गोष्टी अधिक गुंतागुंतीच्या
होतात. प्रतिमान~$\mModel{M^*}$
च्या प्राप्यता संबंध~$R^*$
ने वरील विगमन-सिद्धता
चालण्यासाठी आवश्यक अटी
पूर्ण केल्या, तर त्या
प्रतिमानाला
\emph{निस्यंदन} म्हणू.
मग कोणतेही निस्यंदन~$\mModel{M^*}$
येथे $[w]$ वर~$\varphi$ सत्य
असेल तेव्हा आणि तेव्हाच
$\mModel{M}$ येथे~$w$
वर~$\varphi$ सत्य असेल.
पण अशी निस्यंदने
\emph{अस्तित्वात आहेत}
हेही दाखवावे लागते;
म्हणजे आवश्यक अटी
पूर्ण करणारा~$R^*$
परिभाषित करता आला पाहिजे.
हे घडण्यासाठी~$u$ आणि~$v$
ही जगे~$\varphi$ च्या सर्व उप-सूत्रे वर
एकमत असतील तेव्हा आणि तेव्हाच समतुल्य मानू.
यात $\varphi$ मध्ये येणाऱ्या विधानीय चले वरील
एकमतही समाविष्ट आहे; भाषेतील इतर सर्व चलांवरील एकमत मागितलेले नाही.
$\varphi$ ची उप-सूत्रे
सांत संख्येनेच असल्याने
निस्यंदन सांत राहील.
निस्यंदन परिभाषित
करण्याचा एकच मार्ग नाही;
निस्यंदनाचा प्राप्यता संबंध
आवश्यक गुणधर्म
(उदाहरणार्थ परावर्तित्व,
संक्रमकता इत्यादी)
राखेल यासाठी~$R^*$
ची व्याख्या कल्पकतेने
करावी लागते.













\section{पूर्वतयारी}\label{nml:fil:pre:sec}

निस्यंदनांमुळे आपल्या मोडल तर्कशास्त्रांच्या
प्रणाली निर्णेय आहेत हे सिद्ध करता येते:
त्यांना \emph{सांत प्रतिमान गुणधर्म}
आहे हे त्याने दाखवता येते. म्हणजे,
एखाद्या प्रतिमानात सत्य (असत्य)
असलेले कोणतेही सूत्र
एखाद्या \emph{सांत} प्रतिमानातही
सत्य (असत्य) असते. निस्यंदनांची
व्याख्या उपसूत्रांखाली बंद असलेल्या
सूत्रांच्या संचांच्या
सापेक्ष केली जाते.

\begin{defn}\label{nml:fil:pre:defn:modallyclosed}
  सूत्रांचा संच~$\Gamma$
  हा \emph{उपसूत्रांखाली बंद}
  असेल, जर त्यात~$\Gamma$
  मधील प्रत्येक सूत्र
  चे प्रत्येक उपसूत्र असेल.
  शिवाय, $\Gamma$ हा
  \emph{मोडलदृष्ट्या बंद}
  असेल, जर तो उपसूत्रांखाली
  बंद असेल आणि
  $\varphi \in \Gamma$ असल्यास
  $\Box\varphi,
  \Diamond\varphi \in \Gamma$
  हेही खरे असेल.
\end{defn}

उदाहरणार्थ, सूत्र~$\varphi$
दिल्यावर त्याच्या सर्व
उप-सूत्रांचा संच
उप-सूत्रे खाली बंद
असतो. $\varphi$ च्या उप-सूत्रे
च्या संचातून एखाद्या प्रतिमानाचे
निस्यंदन परिभाषित करताना
हवा असलेला गुणधर्म त्याला
असेल: मूळ प्रतिमानात~$\varphi$
सत्य (असत्य) असेल तेव्हा
आणि तेव्हाच निस्यंदनातही
ते सत्य (असत्य) असेल.

प्रतिमान~$\mModel{M}$ चे
$\Gamma$ मधून मिळणाऱ्या
निस्यंदनातील जगांचा संच
पुढील तुल्यता संबंधाच्या
सर्व तुल्यतावर्गांचा संच
म्हणून परिभाषित होतो.

\begin{defn}
$\mModel{M} =\tuple{W, R, V}$
असू द्या आणि $\Gamma$ हा
उप-सूत्रे खाली बंद
आहे असे समजा.
$\Gamma$ मधील तीच
सूत्रे सत्य ठरवणाऱ्या
कोणत्याही दोन जगांमध्ये
संबंध~$\equiv$ लागू होईल
अशी~$W$ वर त्याची व्याख्या
करू; म्हणजे:
\[u \equiv v \quad \text{तेव्हा आणि तेव्हाच} \quad
\forall \varphi \in \Gamma : \mSat{M}{\varphi}[u] \Leftrightarrow \mSat{M}{\varphi}[v].
\]
जग~$w$ चा तुल्यतावर्ग
$[w]_\equiv$, किंवा संक्षेपाने
$[w]$, म्हणजे~$w$ शी
$\equiv$-समतुल्य असलेल्या
सर्व जगांचा संच:
\[[w] = \Setabs{v}{v \equiv w}.
\]
\end{defn}

\begin{prop}
  $\mModel{M}$ आणि $\Gamma$
  दिल्यावर वरीलप्रमाणे
  परिभाषित $\equiv$ हा
  तुल्यता संबंध आहे;
  म्हणजे तो परावर्ती,
  सममित आणि संक्रमक आहे.
\end{prop}

\begin{proof}
  $\equiv$ हा परावर्ती आहे,
  कारण~$w$ येथे सत्य असलेली
  $\Gamma$ मधील सूत्रे
स्वतः येथील सूत्रांशी
नेमकी जुळतात.
  तो सममित आहे: $u$ आणि~$v$
  येथे $\Gamma$ मधील तीच
  सूत्रे सत्य असतील,
  तर~$v$ आणि~$u$
  यांच्याबाबतीतही तसेच.
  तो संक्रमकही आहे:
  $u$ आणि~$v$ येथे
  $\Gamma$ मधील तीच
  सूत्रे सत्य असतील,
  आणि~$v$ व~$w$ येथेही
  तीच असतील, तर~$u$
  व~$w$ येथेही
  $\Gamma$ मधील तीच
  सूत्रे सत्य असतील.
\end{proof}

इतर प्रत्येक तुल्यता
संबंधाप्रमाणे $\equiv$
हा~$W$ चे \emph{विभाजन}
करतो; म्हणजे~$W$ चे
परस्पर असंयुक्त उपसंच
मिळतात, आणि त्यांचा
संयोग संपूर्ण~$W$ असतो.
प्रत्येक $w \in W$
हा त्यांपैकी एका
विभागाचा, म्हणजे~$[w]$
चा, घटक असतो,
कारण $w \equiv w$.
म्हणून $[w]$ हे विभाग
सर्व~$W$ व्यापतात.
ते परस्पर असंयुक्तही
आहेत: $u \in [w]$
आणि $u \in [v]$
असतील, तर $u \equiv w$
आणि $u \equiv v$.
सममिती व संक्रमकतेवरून
$w \equiv v$; म्हणून
$[w] = [v]$.












\section{निस्यंदने}\label{nml:fil:fil:sec}

$\Gamma$ मधून मिळणारे $\mModel{M}$ चे “एकमेव” निस्यंदन
परिभाषित करण्याऐवजी कोणते प्रतिमान~$\mModel{M^*}$ हे
$\mModel{M}$ चे निस्यंदन ठरते ते परिभाषित करू.
सर्व निस्यंदनांमध्ये जगांचा संच~$W^*$ आणि मूल्यांकन~$V^*$
समान असतात. मात्र वेगवेगळ्या निस्यंदनांचे प्राप्यता
संबंध~$R^*$ वेगळे असू शकतात. निस्यंदन ठरण्यासाठी
$R^*$ ने काही अटी पूर्ण कराव्या लागतात. आपला मुख्य
निष्कर्ष, म्हणजे $\varphi \in \Gamma$ असल्यास
$\mSat{M}{\varphi}[w]$ तेव्हा आणि तेव्हाच
$\mSat{M^*}{\varphi}[{[w]}]$, सिद्ध करण्यासाठी
आपण नेमक्या याच अटी वापरू.

\begin{defn}\label{nml:fil:fil:defn:filtration}
  $\Gamma$ उपसूत्रांखाली बंद आहे आणि
  $\mModel{M} = \tuple{W, R, V}$ असे समजा.
  $\mModel{M}$ चे $\Gamma$ मधून मिळणारे
  \emph{निस्यंदन} म्हणजे पुढील अटी पूर्ण करणारे
  कोणतेही प्रतिमान $\mModel{M^*} = \tuple{W^*,R^*,V^*}$:
  \begin{enumerate}
  \item $W^* = \Setabs{[w]}{w \in W}$;
  \item \label{nml:fil:fil:defn:filtration-R}
    कोणत्याही $u,v \in W$ साठी:
    \begin{enumerate}
    \item \label{nml:fil:fil:defn:filtration-R1}
      $Ruv$ असेल, तर $R^*[u][v]$;
    \item \label{nml:fil:fil:defn:filtration-R2}
      $R^*[u][v]$ असेल, तर कोणत्याही $\Box \varphi \in \Gamma$ साठी,
      $\mSat{M}{\Box\varphi}[u]$ असल्यास $\mSat{M}{\varphi}[v]$;
    \item \label{nml:fil:fil:defn:filtration-R3}
      $R^*[u][v]$ असेल, तर कोणत्याही $\Diamond \varphi \in \Gamma$ साठी,
      $\mSat{M}{\varphi}[v]$ असल्यास $\mSat{M}{\Diamond\varphi}[u]$.
    \end{enumerate}
  \item $V^*(p) = \Setabs{[u]}{u \in V(p)}$.
  \end{enumerate}
\end{defn}

$V^*(p)$ म्हणजे काय ते नीट पाहू: $\mModel{M}$ मध्ये
ज्या प्रत्येक जग~$w$ वर~$p$ सत्य आहे, त्याच्या
तुल्यतावर्ग~$[w]$ यांचा तो संच आहे. एका बाजूने,
$w \in V(p)$ असेल, तर
या व्याख्येवरून $[w] \in V^*(p)$. पण उलट,
$[w] \in V^*(p)$ असल्यावर $w \in V(p)$ असेलच असे नाही.
$[w] \in V^*(p)$ असेल, तर \emph{एखाद्या}~$u \in V(p)$
साठी $[w] = [u]$ एवढीच खात्री मिळते. अर्थात
$[w] = [u]$ याचा अर्थ $w \equiv u$. म्हणून
$[w] \in V^*(p)$ असल्यावर एखाद्या $u \in V(p)$
साठी $w \equiv u$ आहे, एवढाच निष्कर्ष काढता येतो.

\begin{thm}\label{nml:fil:fil:thm:filtrations}
  $\mModel{M^*}$ हे $\mModel{M}$ चे $\Gamma$ मधून
  मिळणारे निस्यंदन असेल, तर प्रत्येक $\varphi \in \Gamma$
  आणि $w \in W$ साठी $\mSat{M}{\varphi}[w]$ असेल तेव्हा
  आणि तेव्हाच $\mSat{M^*}{\varphi}[{[w]}]$ असेल.
\end{thm}

\begin{proof}
  $\Gamma$ उपसूत्रांखाली बंद आहे हे वापरून~$\varphi$ वर
  विगमन करू. $\varphi \in \Gamma$ आणि $\Gamma$ उप-सूत्रे
  खाली बंद असल्याने~$\varphi$ ची सर्व उप-सूत्रे
  देखील $\in \Gamma$. त्यामुळे प्रत्येक विगमन-पायरीत
  $\varphi$ च्या उप-सूत्रे ना विगमन गृहीतक लागू होते.
  \begin{enumerate}
    \item \indcase{\varphi}{\lfalse}{$\mSat{M}{\indfrm}[w]$
        किंवा $\mSat{M^*}{\indfrm}[{[w]}]$ यांपैकी काहीच खरे नाही.}
    \item \indcase{\varphi}{\ltrue}{$\mSat{M}{\indfrm}[w]$
        आणि $\mSat{M^*}{\indfrm}[{[w]}]$ दोन्ही खरे आहेत.}
    \item \indcase{\varphi}{p}{डावीकडून उजवीकडे जाणारा निष्कर्ष
      तात्काळ मिळतो: $\mSat{M}{\indfrm}[w]$ असेल, तर
      $w \in V(p)$; त्यामुळे $[w] \in V^*(p)$,
      म्हणजेच $\mSat{M^*}{\indfrm}[{[w]}]$.
      उलट, $\mSat{M^*}{\indfrm}[{[w]}]$ असे समजा;
      म्हणजे $[w] \in V^*(p)$. मग एखाद्या
      $v \in V(p)$ साठी $w \equiv v$. म्हणून
      $\mSat{M}{p}[v]$ सुद्धा. $w \equiv v$ असल्याने
      $w$ व $v$ येथे $\Gamma$ मधील नेमकी तीच
      सूत्रे सत्य आहेत. गृहीतकानुसार
      $p \in \Gamma$ आणि $\mSat{M}{p}[v]$;
      म्हणून $\mSat{M}{\indfrm}[w]$.}

    \item \indcase{\varphi}{\lnot
          \psi}{$\mSat{M}{\indfrm}[w]$ तेव्हा आणि तेव्हाच
          $\mSat/{M}{\psi}[w]$. विगमन गृहीतकानुसार,
          $\mSat/{M}{\psi}[w]$ तेव्हा आणि तेव्हाच
          $\mSat/{M^*}{\psi}[{[w]}]$. अखेरीस,
          $\mSat/{M^*}{\psi}[{[w]}]$ तेव्हा आणि तेव्हाच
          $\mSat{M^*}{\indfrm}[{[w]}]$.}

    \item \indcase{\varphi}{(\psi \land
          \chi)}{$\mSat{M}{\indfrm}[w]$ तेव्हा आणि तेव्हाच
          $\mSat{M}{\psi}[w]$ आणि $\mSat{M}{\chi}[w]$.
          विगमन गृहीतकानुसार, $\mSat{M}{\psi}[w]$ तेव्हा आणि तेव्हाच
          $\mSat{M^*}{\psi}[{[w]}]$, तसेच $\mSat{M}{\chi}[w]$
          तेव्हा आणि तेव्हाच $\mSat{M^*}{\chi}[{[w]}]$.
          आणि $\mSat{M^*}{\indfrm}[{[w]}]$ तेव्हा आणि तेव्हाच
          $\mSat{M^*}{\psi}[{[w]}]$ आणि
          $\mSat{M^*}{\chi}[{[w]}]$.}

    \item \indcase{\varphi}{(\psi \lor
          \chi)}{$\mSat{M}{\indfrm}[w]$ तेव्हा आणि तेव्हाच
          $\mSat{M}{\psi}[w]$ किंवा $\mSat{M}{\chi}[w]$.
          विगमन गृहीतकानुसार, $\mSat{M}{\psi}[w]$ तेव्हा आणि तेव्हाच
          $\mSat{M^*}{\psi}[{[w]}]$, तसेच $\mSat{M}{\chi}[w]$
          तेव्हा आणि तेव्हाच $\mSat{M^*}{\chi}[{[w]}]$.
          आणि $\mSat{M^*}{\indfrm}[{[w]}]$ तेव्हा आणि तेव्हाच
          $\mSat{M^*}{\psi}[{[w]}]$ किंवा
          $\mSat{M^*}{\chi}[{[w]}]$.}

    \item \indcase{\varphi}{(\psi \lif
          \chi)}{$\mSat{M}{\indfrm}[w]$ तेव्हा आणि तेव्हाच
          $\mSat/{M}{\psi}[w]$ किंवा $\mSat{M}{\chi}[w]$.
          विगमन गृहीतकानुसार, $\mSat{M}{\psi}[w]$ तेव्हा आणि तेव्हाच
          $\mSat{M^*}{\psi}[{[w]}]$, तसेच $\mSat{M}{\chi}[w]$
          तेव्हा आणि तेव्हाच $\mSat{M^*}{\chi}[{[w]}]$.
          आणि $\mSat{M^*}{\indfrm}[{[w]}]$ तेव्हा आणि तेव्हाच
          $\mSat/{M^*}{\psi}[{[w]}]$ किंवा
          $\mSat{M^*}{\chi}[{[w]}]$.}

    \item \indcase{\varphi}{(\psi \liff
          \chi)}{$\mSat{M}{\indfrm}[w]$ तेव्हा आणि तेव्हाच
          $\mSat{M}{\psi}[w]$ आणि $\mSat{M}{\chi}[w]$,
          किंवा $\mSat/{M}{\psi}[w]$ आणि $\mSat/{M}{\chi}[w]$.
          विगमन गृहीतकानुसार, $\mSat{M}{\psi}[w]$ तेव्हा आणि तेव्हाच
          $\mSat{M^*}{\psi}[{[w]}]$, तसेच $\mSat{M}{\chi}[w]$
          तेव्हा आणि तेव्हाच $\mSat{M^*}{\chi}[{[w]}]$.
          आणि $\mSat{M^*}{\indfrm}[{[w]}]$ तेव्हा आणि तेव्हाच
          $\mSat{M^*}{\psi}[{[w]}]$ आणि
          $\mSat{M^*}{\chi}[{[w]}]$, किंवा
          $\mSat/{M^*}{\psi}[{[w]}]$ आणि
          $\mSat/{M^*}{\chi}[{[w]}]$.}

    \item \indcase{\varphi}{\Box
          \psi}{$\mSat{M}{\indfrm}[w]$ असे समजा.
          $\mSat{M^*}{\indfrm}[{[w]}]$ दाखवण्यासाठी
          $R^*[w][v]$ असेल असे कोणतेही~$v$ घ्या.
          \ref{nml:fil:fil:defn:filtration}\ref{nml:fil:fil:defn:filtration-R2}
          वरून $\mSat{M}{\psi}[v]$; आणि विगमन गृहीतकानुसार
          $\mSat{M^*}{\psi}[{[v]}]$. $v$ कोणतेही असल्याने
          $\mSat{M^*}{\indfrm}[{[w]}]$.

          उलट, $\mSat{M^*}{\indfrm}[{[w]}]$ असे समजा
          आणि $Rwv$ असेल असे कोणतेही~$v$ घ्या.
          \ref{nml:fil:fil:defn:filtration}\ref{nml:fil:fil:defn:filtration-R1}
          वरून $R^*[w][v]$; म्हणून
          $\mSat{M^*}{\psi}[{[v]}]$. विगमन गृहीतकानुसार
          $\mSat{M}{\psi}[v]$. $v$ कोणतेही असल्याने
          $\mSat{M}{\indfrm}[w]$.}

    \item \indcase{\varphi}{\Diamond
          \psi}{$\mSat{M}{\indfrm}[w]$ असे समजा. मग एखाद्या
          $v \in W$ साठी $Rwv$ आणि $\mSat{M}{\psi}[v]$.
          विगमन गृहीतकानुसार $\mSat{M^*}{\psi}[{[v]}]$;
          तसेच \ref{nml:fil:fil:defn:filtration}\ref{nml:fil:fil:defn:filtration-R1}
          वरून $R^*[w][v]$. म्हणून
          $\mSat{M^*}{\indfrm}[{[w]}]$.

          आता $\mSat{M^*}{\indfrm}[{[w]}]$ असे समजा.
          मग $R^*[w][v]$ आणि $\mSat{M^*}{\psi}[{[v]}]$
          असेल असा एखादा $[v] \in W^*$ आहे.
          विगमन गृहीतकानुसार $\mSat{M}{\psi}[v]$.
          \ref{nml:fil:fil:defn:filtration}\ref{nml:fil:fil:defn:filtration-R3}
          वरून $\mSat{M}{\indfrm}[w]$.}
  \end{enumerate}
\end{proof}

\begin{prob}
  \ref{nml:fil:fil:thm:filtrations} ची सिद्धता पूर्ण करा.
\end{prob}

प्रतिमानातील प्रत्येक जगावरील सत्यतेविषयीचा निष्कर्ष
प्रतिमानातील सत्यतेसाठी आणि प्रतिमानांच्या वर्गातील
वैधतेसाठीही लागू होतो.

\begin{cor}
  $\Gamma$ उपसूत्रांखाली बंद आहे असे समजा. मग:
  \begin{enumerate}
  \item $\mModel{M^*}$ हे $\mModel{M}$ चे $\Gamma$
    मधून मिळणारे निस्यंदन असेल, तर कोणत्याही
    $\varphi \in \Gamma$ साठी $\mSat{M}{\varphi}$ असेल तेव्हा
    आणि तेव्हाच $\mSat{M^*}{\varphi}$ असेल.
  \item $\mClass{C}$ हा प्रतिमानांचा वर्ग असेल आणि
    $\Gamma(\mClass{C})$ हा $\mClass{C}$ मधील प्रतिमानांच्या
    $\Gamma$-निस्यंदनांचा वर्ग असेल, तर कोणतेही
    सूत्र~$\varphi \in \Gamma$ हे $\mClass{C}$ मध्ये
    वैध असेल तेव्हा आणि तेव्हाच ते
    $\Gamma(\mClass{C})$ मध्ये वैध असेल.
  \end{enumerate}
\end{cor}














\section{निस्यंदनांची उदाहरणे}\label{nml:fil:exf:sec}

निस्यंदने अस्तित्वात आहेत हे आपण अजून दाखवलेले नाही.
पण प्रत्यक्षात कोणतेही प्रतिमान~$\mModel{M}$ घेतले,
तरी $\mModel{M}$ ची $\Gamma$ मधून मिळणारी अनेक निस्यंदने असतात.
त्यांपैकी दोन विशेष निस्यंदने पाहू: सर्वांत सूक्ष्म
आणि सर्वांत स्थूल. त्याच प्रतिमानाच्या निस्यंदनांचे
प्राप्यता संबंध वेगळे असू शकतात
(\ref{nml:fil:fil:defn:filtration} मध्ये $W^*$ आणि~$V^*$
थेट निश्चित केले आहेत). सर्वांत सूक्ष्म निस्यंदनात
संबंधित जगांच्या जोड्या शक्य तितक्या कमी असतात,
तर सर्वांत स्थूल निस्यंदनात त्या शक्य तितक्या अधिक असतात.

\begin{defn}
  $\Gamma$ उपसूत्रांखाली बंद असेल, तर प्रतिमान~$\mModel{M}$
  चे \emph{सर्वांत सूक्ष्म} निस्यंदन~$\mModel{M^*}$ पुढीलप्रमाणे
  परिभाषित होते:
  \[  R^*[u][v] \quad \text{तेव्हा आणि तेव्हाच} \quad \exists u'\in [u] \;
  \exists v' \in [v] : Ru'v'.
  \]
\end{defn}

\begin{prop}\label{nml:fil:exf:prop:finest}
  सर्वांत सूक्ष्म निस्यंदन~$\mModel{M^*}$ हे खरोखर निस्यंदन आहे.
\end{prop}

\begin{proof}
  अशा प्रकारे परिभाषित $R^*$ हा
  \ref{nml:fil:fil:defn:filtration}\ref{nml:fil:fil:defn:filtration-R}
  मधील अटी पूर्ण करतो हे तपासायचे आहे. तिन्ही अटी क्रमाने पाहू.

  $Ruv$ असेल, तर $u \in [u]$ आणि $v \in [v]$ असल्याने
  $R^*[u][v]$ सुद्धा. म्हणून
  \ref{nml:fil:fil:defn:filtration-R1} ही अट पूर्ण होते.


        \ref{nml:fil:fil:defn:filtration-R2} साठी $\Box\varphi \in \Gamma$,
        $R^*[u][v]$ आणि $\mSat{M}{\Box\varphi}[u]$ असे समजा.
        $R^*$ च्या व्याख्येवरून $u' \equiv u$ आणि $v' \equiv v$
        अशी जगे आहेत की $Ru'v'$. $u$ आणि $u'$ हे $\Gamma$
        वरील सत्यतेबाबत एकमत असल्याने
        $\mSat{M}{\Box\varphi}[u']$; म्हणून $\mSat{M}{\varphi}[v']$.
        $\Gamma$ उप-सूत्रे खाली बंद असल्याने
        $v$ आणि $v'$ यांच्यात~$\varphi$ च्या सत्यतेबाबतही
        एकमत आहे. म्हणून अपेक्षेप्रमाणे $\mSat{M}{\varphi}[v]$.


        \ref{nml:fil:fil:defn:filtration-R3} तपासण्यासाठी
        $\Diamond\varphi \in \Gamma$, $R^*[u][v]$ आणि
        $\mSat{M}{\varphi}[v]$ असे समजा. $R^*$ च्या व्याख्येवरून
        $u' \equiv u$ आणि $v' \equiv v$ अशी जगे आहेत की
        $Ru'v'$. $v$ आणि $v'$ यांच्यात~$\Gamma$ वरील
        सत्यतेबाबत एकमत आहे आणि $\Gamma$ उप-सूत्रे
        खाली बंद आहे; म्हणून $\mSat{M}{\varphi}[v']$ सुद्धा.
        त्यामुळे $\mSat{M}{\Diamond \varphi}[u']$.
        $u$ आणि $u'$ यांच्यातही $\Gamma$ वरील
        सत्यतेबाबत एकमत असल्याने
        $\mSat{M}{\Diamond \varphi}[u]$.
\end{proof}



\begin{defn}
  $\Gamma$ उपसूत्रांखाली बंद असेल, तर प्रतिमान~$\mModel{M}$
  चे \emph{सर्वांत स्थूल} निस्यंदन~$\mModel{M^*}$
  पुढीलप्रमाणे परिभाषित करतो:
  $R^*[u][v]$ तेव्हा आणि तेव्हाच

    पुढील \emph{दोन्ही} अटी पूर्ण होतात:
  \begin{enumerate}
  \item \label{nml:fil:exf:defn:coarsest-Box}$\Box\varphi \in \Gamma$
    आणि $\mSat{M}{\Box\varphi}[u]$ असेल, तर
    $\mSat{M}{\varphi}[v]$;
  \item \label{nml:fil:exf:defn:coarsest-Diamond}$\Diamond\varphi
    \in \Gamma$ आणि $\mSat{M}{\varphi}[v]$ असेल, तर
    $\mSat{M}{\Diamond\varphi}[u]$.
  \end{enumerate}
\end{defn}

\begin{prop}
  सर्वांत स्थूल निस्यंदन~$\mModel{M^*}$ हे खरोखर निस्यंदन आहे.
\end{prop}

\begin{proof}
  $R^*$ च्या व्याख्येनुसार आता फक्त $Ruv$ वरून
  $R^*[u][v]$ मिळते का हे तपासायचे आहे. म्हणून
  $Ruv$ असे समजा. $\Box\varphi \in \Gamma$
    आणि $\mSat{M}{\Box\varphi}[u]$ असे समजा; मग स्पष्टपणे
    $\mSat{M}{\varphi}[v]$, आणि
      \ref{nml:fil:exf:defn:coarsest-Box} ही अट पूर्ण होते.
    $\Diamond\varphi \in \Gamma$ आणि $\mSat{M}{\varphi}[v]$
    असे समजा. $Ruv$ असल्याने $\mSat{M}{\Diamond\varphi}[u]$
    , आणि \ref{nml:fil:exf:defn:coarsest-Diamond}
      ही अट पूर्ण होते.
\end{proof}

\begin{ex}
  $W = \PosInt$, $Rnm$ तेव्हा आणि तेव्हाच $m = n + 1$,
  आणि $V(p) = \Setabs{2n}{n \in
    \PosInt}$ असे समजा. प्रतिमान~$\mModel{M} = \tuple{W, R, V}$
  हे \ref{nml:fil:exf:fig:ex-filtration} मध्ये दाखवले आहे.
  त्याची जगे $1$, $2$ इत्यादी आहेत. प्रत्येक जगाला
  नेमके एकच जग, म्हणजे त्याचे पुढचे जग, प्राप्य आहे;
  आणि $p$ तेव्हा आणि तेव्हाच सत्य आहे जेव्हा संख्या सम असते.

  \begin{figure}
    \centering
    \begin{tikzpicture}[modal]
      \node[world] (1) [label=below:\mFalse{p}] {$1$};
      \node[world] (2) [label=below:\mTrue{p}, right=of 1] {$2$};
      \node[world] (3) [label=below:\mFalse{p}, right=of 2] {$3$};
      \node[world] (4) [label=below:\mTrue{p}, right=of 3] {$4$};
      \node[phantom] (5) [right=of 4] {};
      \draw[->] (1) to (2);
      \draw[->] (2) to (3);
      \draw[->] (3) to (4);
      \draw[dotted] (4) to (5);
    \end{tikzpicture}

    \begin{tikzpicture}[modal]
      \node[world] (1) [label=below:\mFalse{p}] {$[1]$};
      \node[world] (2) [label=below:\mTrue{p}, right=of 1] {$[2]$};
      \draw[->,bend left] (1) to (2);
      \draw[->,bend left] (2) to (1);

      \node[world] (3) [label=below:\mFalse{p},right=of 2] {$[1]$};
      \node[world] (4) [label=below:\mTrue{p}, right=of 3] {$[2]$};
      \draw[->,bend left] (3) to (4);
      \draw[->,bend left] (4) to (3);
      \draw[->,reflexive above] (4) to (4);
    \end{tikzpicture}
    \caption{एक अनंत प्रतिमान आणि त्याची निस्यंदने.}
    \label{nml:fil:exf:fig:ex-filtration}
  \end{figure}

  आता $\Gamma$ हा~$\Box p \lif p$ च्या सर्व
  उप-सूत्रांचा संच, म्हणजे
  $\{p, \Box p, \Box p \lif p\}$, असू द्या.
  $p$ सम संख्यांवरच सत्य आहे; $\Box p$ विषम
  संख्यांवरच सत्य आहे. म्हणून $\Box p \lif p$
  हे सम संख्यांवरच सत्य आहे. दुसऱ्या शब्दांत,
  प्रत्येक विषम संख्येवर $\Box p$ सत्य, पण $p$
  आणि $\Box p \lif p$ असत्य आहेत; प्रत्येक सम
  संख्येवर $p$ आणि $\Box p \lif p$ सत्य, पण
  $\Box p$ असत्य आहे. म्हणून
  $W^* = \{ [1], [2] \}$, जेथे
  $[1] = \{1, 3, 5, \dots\}$ आणि
  $[2] = \{2, 4, 6, \dots\}$.
  $2 \in V(p)$ असल्याने $[2] \in V^*(p)$;
  $1 \notin V(p)$ असल्याने $[1] \notin V^*(p)$.
  म्हणून $V^*(p) = \{[2]\}$.

  $W^*$ वर आधारलेल्या कोणत्याही निस्यंदनाच्या
  प्राप्यता संबंधात $\tuple{[1], [2]}, \tuple{[2],[1]}$
  या जोड्या असायलाच हव्यात: $R12$ असल्याने
  \ref{nml:fil:fil:defn:filtration}\ref{nml:fil:fil:defn:filtration-R1}
  वरून $R^*[1][2]$; तसेच $R23$ असल्याने
  $R^*[2][3]$, आणि $[3]=[1]$.
  त्यात $\tuple{[1],[1]}$ मात्र असू शकत नाही:
  ती जोडी असेल, तर $R^*[1][1]$ होईल;
  पण $\mSat{M}{\Box p}[1]$ आणि
  $\mSat/{M}{p}[1]$ असल्याने
  \ref{nml:fil:fil:defn:filtration-R2} शी विरोध येईल.
  $R^*[2][2]$ असावे किंवा नसावे, अशी कोणतीही
  अट नाही. म्हणून~$\mModel{M}$ ची दोन संभाव्य
  निस्यंदने आहेत; त्यांचे प्राप्यता संबंध अनुक्रमे
  \[  \{\tuple{[1],[2]}, \tuple{[2],[1]}\} \text{आणि}
  \{\tuple{[1],[2]}, \tuple{[2],[1]}, \tuple{[2],[2]}\}.
  \]
  असे आहेत. दोन्ही प्रकरणांत~$[1]$ येथे $p$ आणि
  $\Box p \lif p$ असत्य, तर $\Box p$ सत्य आहे;
  $[2]$ येथे $p$ आणि $\Box p \lif p$ सत्य,
  तर $\Box p$ असत्य आहे.
\end{ex}


\begin{prob}
  पुढील प्रतिमान~$\mModel{M} = \tuple{W, R, V}$ पाहा.
  येथे $W = \Setabs{0\sigma}{\sigma \in \Bin^*}$,
  म्हणजे~$0$ ने सुरू होणाऱ्या $0$ व~$1$ यांच्या
  सांत क्रमांचा संच; $R\sigma\sigma'$ तेव्हा आणि तेव्हाच
  $\sigma' = \sigma 0$ किंवा $\sigma' = \sigma 1$;
  तसेच $V(p) = \{0\} \cup \Setabs{\sigma
    0}{\sigma \in W}$ आणि $V(q) = \Setabs{\sigma 1}{\sigma \in
    W}$. त्याचे चित्र असे आहे:
  \begin{center}
    \begin{tikzpicture}[modal,
        node distance=2cm
      ]

      \node[world] (0) [label={below:\mTrue{p}\\\mFalse{q}},font=\tiny] {$0$};
      \node[world] (00) [label={below:\mTrue{p}\\\mFalse{q}},
        above right=of 0,font=\tiny] {$00$};
      \node[world] (000) [label={below:\mTrue{p}\\\mFalse{q}},
        right=of 00, yshift=.8cm, font=\tiny] {$000$};
      \node[phantom] (0000) [right=of 000, yshift=.5cm] {};
      \node[phantom] (0001) [right=of 000, yshift=-.5cm] {};
      \node[world] (001) [label={below:\mFalse{p}\\\mTrue{q}},
        right=of 00, yshift=-.8cm, font=\tiny] {$001$};
      \node[phantom](0010) [right=of 001, yshift=.5cm] {};
      \node[phantom](0011) [right=of 001, yshift=-.5cm] {};
      \node[world] (01) [label={below:\mFalse{p}\\\mTrue{q}},
        below right=of 0,font=\tiny] {$01$};
      \node[world] (010) [label={below:\mTrue{p}\\\mFalse{q}},
        right=of 01, yshift=.8cm,font=\tiny] {$010$};
      \node[phantom] (0100) [right=of 010, yshift=.5cm] {};
      \node[phantom] (0101) [right=of 010, yshift=-.5cm] {};
      \node[world] (011) [label={below:\mFalse{p}\\\mTrue{q}},
        right=of 01, yshift=-.8cm,font=\tiny] {$011$};
      \node[phantom] (0110) [right=of 011, yshift=.5cm] {};
      \node[phantom] (0111) [right=of 011, yshift=-.5cm] {};

      \draw[->] (0) to (00);
      \draw[->] (0) to (01);
      \draw[->] (00) to (000);
      \draw[dotted] (000) to (0000);
      \draw[dotted] (000) to (0001);
      \draw[->] (00) to (001);
      \draw[dotted] (001) to (0010);
      \draw[dotted] (001) to (0011);
      \draw[->] (01) to (010);
      \draw[dotted] (010) to (0100);
      \draw[dotted] (010) to (0101);
      \draw[->] (01) to (011);
      \draw[dotted] (011) to (0110);
      \draw[dotted] (011) to (0111);
    \end{tikzpicture}
  \end{center}
  प्रत्येक~$w$ साठी
  $\mSat/{M}{\Box(p \lor q) \lif (\Box p \lor \Box q)}[w]$
  असे आहे.

  $\Gamma$ हा~$\Box(p \lor q) \lif
  (\Box p \lor \Box q)$ च्या उप-सूत्रांचा संच
  असू द्या. $W^*$ आणि~$V^*$ काय असतील?
  $\mModel{M}$ च्या सर्वांत सूक्ष्म निस्यंदनाचा
  प्राप्यता संबंध कोणता असेल? सर्वांत स्थूल निस्यंदनाचा कोणता?
\end{prob}












\section{निस्यंदने सांत असतात}\label{nml:fil:fin:sec}

उप-सूत्रे खाली बंद असलेल्या कोणत्याही संच~$\Gamma$
साठी आपण निस्यंदनांची व्याख्या केली. केवळ त्या
व्याख्येवरून निस्यंदने सांत असतील याची खात्री मिळत नाही.
खरे तर~$\Gamma$ अनंत असेल (उदाहरणार्थ, सर्व
सूत्रांचा संच असेल), तर निस्यंदनही अनंत असू शकते.
मात्र~$\Gamma$ सांत असेल (उदाहरणार्थ, दिलेल्या
सूत्र~$\varphi$ च्या उप-सूत्रांचा संच असेल),
तर~$\Gamma$ मधून मिळणारे प्रत्येक निस्यंदनही सांत असते.

\begin{prop}\label{nml:fil:fin:prop:filt-are-finite}
  $\Gamma$ सांत असेल, तर प्रतिमान~$\mModel{M}$ चे
  $\Gamma$ मधून मिळणारे कोणतेही निस्यंदन~$\mModel{M^*}$
  सुद्धा सांत असते.
\end{prop}

\begin{proof}
  $W^*$ चा आकार म्हणजे तुल्यता संबंध~$\equiv$ अंतर्गत
  वेगवेगळ्या वर्ग~$[w]$ ची संख्या. अशा एकाच वर्गातील
  कोणतीही दोन जगे $u$, $v$, म्हणजे $u$ आणि $v$
  अशी जगे की $u \equiv v$, ही $\Gamma$ मधील सर्व सूत्रे~$\varphi$
  यांच्या सत्यतेबाबत एकमत असतात: $\varphi \in \Gamma$ असेल,
  तर~$\varphi$ हे $u$ आणि~$v$ या दोन्ही जगांवर सत्य असते
  किंवा दोन्हीवर असत्य असते. म्हणून प्रत्येक वर्ग~$[w]$
  हा~$\Gamma$ च्या एका उपसंचाशी जुळतो:~$[w]$ मधील
  जगांवर~$\varphi$ सत्य असेल अशा सर्व $\varphi \in \Gamma$ च्या संचाशी.
  दोन वेगवेगळे वर्ग $[u]$ आणि $[v]$ हे $\Gamma$ च्या
  एकाच उपसंचाशी जुळू शकत नाहीत. कारण~$u$ येथे सत्य
  असलेल्या सूत्रांचा संच आणि~$v$ येथे सत्य असलेल्या
  सूत्रांचा संच एकच असेल, तर~$u$ आणि~$v$ यांच्यात
  $\Gamma$ मधील सर्व सूत्रांच्या सत्यतेबाबत एकमत असते;
  म्हणजे $u \equiv v$. पण मग $[u] = [v]$.
  म्हणून $W^*$ पासून~$\Pow{\Gamma}$ मध्ये एकास-एक
  फल आहे आणि त्यामुळे $\card{W^*} \le \card{\Pow{\Gamma}}$.
  म्हणून~$\Gamma$ मध्ये $n$ वाक्ये असतील, तर
  $W^*$ ची गणनसंख्या~$2^n$ पेक्षा अधिक नाही.
\end{proof}













\section{\Log{K} आणि \Log{S5} यांना सांत प्रतिमान गुणधर्म आहे}\label{nml:fil:fmp:sec}

\begin{defn}
  मोडल तर्कशास्त्राच्या प्रणाली~$\Sigma$ ला
  \emph{सांत प्रतिमान गुणधर्म} आहे असे म्हणतात, जर
  सूत्र~$\varphi$ हे~$\Sigma$ च्या एखाद्या
  प्रतिमानातील एखाद्या जगावर सत्य असेल, तर~$\varphi$
  हे~$\Sigma$ च्या एखाद्या \emph{सांत} प्रतिमानातील
  एखाद्या जगावरही सत्य असेल. येथे प्रणालीचे प्रतिमान म्हणजे
  त्या प्रणालीतील प्रत्येक सूत्र प्रत्येक जगावर सत्य असलेले प्रतिमान.
\end{defn}

\begin{prop}\label{nml:fil:fmp:prop:K-fmp}
  \Log{K} ला सांत प्रतिमान गुणधर्म आहे.
\end{prop}

\begin{proof}
  \Log{K} हा वैध सूत्रांचा संच आहे;
  म्हणजे कोणतेही प्रतिमान हे~\Log{K} चे प्रतिमान आहे.
  \ref{nml:fil:fil:thm:filtrations} नुसार,
  $\mSat{M}{\varphi}[w]$ असेल, तर~$\varphi$ च्या
  उप-सूत्रांच्या संच~$\Gamma$ मधून मिळणाऱ्या
  $\mModel{M}$ च्या कोणत्याही निस्यंदनासाठी
  $\mSat{M^*}{\varphi}[{[w]}]$. कोणत्याही सूत्र
  ची उप-सूत्रे सांत संख्येनेच असतात; म्हणून
  $\Gamma$ सांत आहे. \ref{nml:fil:fin:prop:filt-are-finite}
  नुसार $\card{W^*} \le 2^n$, जेथे $n$ ही~$\Gamma$
  मधील सूत्रांची संख्या आहे. शिवाय \Log{K}
  प्रतिमानांवर कोणतेही निर्बंध घालत नसल्याने
  $\mModel{M^*}$ हे \Log{K}-प्रतिमान आहे.
\end{proof}

तर्कशास्त्र~\Log{L} याला निस्यंदनांद्वारे
सांत प्रतिमान गुणधर्म आहे हे दाखवण्यासाठी,
\Log{L}-प्रतिमानाचे निस्यंदन हे स्वतः
\Log{L}-प्रतिमान असणे आवश्यक आहे. यासाठी
अनेकदा बरेच काम करावे लागते आणि कोणतेही
निस्यंदन नेहमीच \Log{L}-प्रतिमान देत नाही.
मात्र वैश्विक प्रतिमानांच्या बाबतीत हे खरे आहे.

\begin{prop}\label{nml:fil:fmp:prop:univ-fin}
  $\mClass{U}$ हा वैश्विक प्रतिमानांचा वर्ग
  (\ref{nml:frd:es5:prop:S5=univ} पाहा) आणि
  $\mClass{U}_\mathrm{Fin}$ हा सर्व सांत वैश्विक
  प्रतिमानांचा वर्ग असू द्या. मग कोणतेही
  सूत्र~$\varphi$ हे $\mClass{U}$ मध्ये वैध असेल
  तेव्हा आणि तेव्हाच ते $\mClass{U}_\mathrm{Fin}$
  मध्ये वैध असेल.
\end{prop}

\begin{proof}
  सांत वैश्विक प्रतिमाने ही वैश्विक प्रतिमानेच
  असल्याने डावीकडून उजवीकडे जाणारा निष्कर्ष
  तात्काळ मिळतो. उलट दिशेसाठी, वैश्विक
  प्रतिमान~$\mModel{M}$ मधील एखाद्या जगावर~$w$
  $\varphi$ असत्य आहे असे समजा. $\Gamma$ मध्ये $\varphi$
  आणि त्याची सर्व उपसूत्रे घ्या; स्पष्टपणे $\Gamma$
  सांत आहे. $\mModel{M}$ चे एखादे निस्यंदन~$\mModel{M^*}$
  घ्या. \ref{nml:fil:fin:prop:filt-are-finite} नुसार
  $\mModel{M^*}$ सांत आहे; आणि
  \ref{nml:fil:fil:thm:filtrations} नुसार~$\mModel{M^*}$
  मधील $[w]$ येथे $\varphi$ असत्य आहे. $\mModel{M^*}$
  हेही वैश्विक आहे, एवढे आता पाहायचे आहे:
  कोणतीही $u$ आणि $v$ दिली, तर गृहीतकानुसार
  $Ruv$; आणि
  \ref{nml:fil:fil:defn:filtration}\ref{nml:fil:fil:defn:filtration-R}
  नुसार $R^*[u][v]$ सुद्धा.
\end{proof}

\begin{cor}\label{nml:fil:fmp:cor:S5fmp}
  \Log{S5} ला सांत प्रतिमान गुणधर्म आहे.
\end{cor}

\begin{proof}
  प्रथम, $\varphi$ हे \Log{S5} च्या कोणत्याही प्रतिमानातील एखाद्या
  जगावर सत्य आहे असे समजा. मग $\lnot\varphi$ हे \Log{S5} मध्ये
  निष्पन्न करता येण्याजोगे नसते; अन्यथा ते त्या प्रतिमानाच्या प्रत्येक जगावर
  सत्य असते. \ref{nml:com:fra:thm:generaldet} नुसार $\lnot\varphi$
  हे परावर्ती युक्लिडी प्रतिमानांच्या वर्गात वैध नसते. म्हणून $\varphi$
  हे त्या वर्गातील एखाद्या प्रतिमानाच्या एखाद्या जगावर सत्य असते.
  \ref{nml:frd:es5:prop:S5=univ} नुसार,
  $\varphi$ हे एखाद्या परावर्ती व युक्लिडी प्रतिमानातील
  एखाद्या जगावर सत्य असेल, तर ते एखाद्या वैश्विक
  प्रतिमानातील जगावर सत्य असते.
  \ref{nml:fil:fmp:prop:univ-fin} च्या सिद्धतेतील निस्यंदन वापरल्यावर
  ते एखाद्या सांत
  वैश्विक प्रतिमानातील जगावर सत्य असते
  (म्हणजे त्या प्रतिमानाचे~$\varphi$ च्या उप-सूत्रे
  च्या संचातून मिळणारे निस्यंदन).
  प्रत्येक वैश्विक प्रतिमान परावर्ती व युक्लिडीही असते;
  म्हणून $\varphi$ हे सांत परावर्ती युक्लिडी प्रतिमानातील
  एखाद्या जगावर सत्य असते.
\end{proof}

\begin{prob}
  सर्वत्र उत्तरवर्ती किंवा परावर्ती प्रतिमानाचे
  कोणतेही निस्यंदन अनुक्रमे सर्वत्र उत्तरवर्ती किंवा
  परावर्ती असते, हे दाखवा.
\end{prob}

\begin{prob}
  सममित (संक्रमक, युक्लिडी) प्रतिमानाचे
  सममित नसलेले (संक्रमक नसलेले, युक्लिडी नसलेले)
  निस्यंदन शोधा.
\end{prob}













\section{\Log{S5} निर्णेय आहे}\label{nml:fil:dec:sec}
दिलेल्या सांत रूपबंध-यादीने प्रभावीपणे परिभाषित मोडल तर्कशास्त्राच्या
प्रणाली \emph{निर्णेय} आहेत हे दाखवण्याचा
सांत प्रतिमान गुणधर्म हा सोपा मार्ग देतो.
म्हणजे, सूत्र त्या प्रणालीत
निष्पन्न करता येण्याजोगे आहे की नाही हे ठरवणारी
संगणनक्षम पद्धत मिळते.

\begin{thm}
  \Log{S5} निर्णेय आहे.
\end{thm}

\begin{proof}
  $\varphi$ दिले आहे असे समजा आणि~$\varphi$ मध्ये येणारी
  विधानीय चले ही $p_1$, \dots, $p_k$
  यांपैकीच आहेत असे समजा. प्रत्येक~$n\geq1$ साठी
  जगांचा निश्चित संच $W=\{1,\ldots,n\}$ आणि वैश्विक संबंध
  $R=W\times W$ घेऊ. $p_1$, \dots, $p_k$ यांच्या $2^{nk}$
  मूल्यांकनांची यादी सांत आहे; इतर सर्व विधानचलांना रिक्त संच नेमू.
  त्यामुळे प्रत्येक आकारासाठी तपासायच्या वैश्विक प्रतिमानांची सांत
  प्रभावी यादी मिळते.
  म्हणून दोन प्रक्रिया \emph{समांतर} चालवू शकतो:
  एका बाजूने सिद्धता पद्धतशीरपणे तयार करून
  \Log{S5} ची सर्व प्रमेये क्रमाने नोंदवायची;
  दुसऱ्या बाजूने $1$, $2$, \dots जगे असलेली सर्व
  वैश्विक प्रतिमाने क्रमाने नोंदवायची आणि त्यांपैकी
  एखाद्या प्रतिमानातील एखाद्या जगावर~$\varphi$ असत्य
  आहे का हे तपासायचे. अखेरीस या दोनपैकी एक
  प्रक्रिया उत्तर देईल. कारण
  \ref{nml:com:fra:thm:generaldet} आणि
  \ref{nml:fil:fmp:cor:S5fmp} नुसार, एकतर~$\varphi$
  निष्पन्न करता येण्याजोगे आहे किंवा ते एखाद्या सांत
  वैश्विक प्रतिमानात असत्य आहे.
\end{proof}

वरील सिद्धता \Log{S5} साठी चालते, कारण
वैश्विक प्रतिमानांची निस्यंदने आपोआप वैश्विक
असतात. परावर्तित्व आणि सर्वत्र उत्तरवर्तित्व
यांबाबतही तसेच आहे; इतर गुणधर्मांसाठी अधिक
काम करावे लागते.













\section{निस्यंदने आणि प्राप्यतेचे गुणधर्म}\label{nml:fil:acc:sec}

पाहिल्याप्रमाणे, वैश्विक, सर्वत्र उत्तरवर्ती आणि परावर्ती
प्रतिमानांची निस्यंदने अनुक्रमे वैश्विक, सर्वत्र उत्तरवर्ती
आणि परावर्ती असतात. पण सममित किंवा संक्रमक प्रतिमानाचे
प्रत्येक निस्यंदन अनुक्रमे सममित किंवा संक्रमक असेलच असे नाही.
तरी काही बाबतीत हे गुणधर्म टिकतील अशी निस्यंदने
परिभाषित करता येतात. त्यासाठी सर्वांत स्थूल निस्यंदनाच्या
व्याख्येसारखीच पद्धत वापरू, पण~$R^*$ च्या व्याख्येत
अतिरिक्त अटी घालू. $\Gamma$ उप-सूत्रे खाली बंद
आहे असे समजा. प्रतिमान~$\mModel{M} =\tuple{W, R, V}$
मधील जगांमध्ये $u$, $v$ लागू होणारे
\ref{nml:fil:acc:tab:Cn-filtrations} मधील संबंध~$C_i(u,v)$ पाहा.
या अटींच्या संयोगांवरून $R^*[u][v]$ परिभाषित करता येतो.
उदाहरणार्थ, $R^*[u][v]$ तेव्हा आणि तेव्हाच $C_1(u,v)$
ही अट पूर्ण होते असे ठरवल्यास नेमके सर्वांत स्थूल
निस्यंदन मिळते. $R^*[u][v]$ तेव्हा आणि तेव्हाच
$C_1(u,v)$ आणि $C_2(u, v)$ या दोन्ही अटी पूर्ण होतात
असे ठरवल्यास अधिक सूक्ष्म किंवा तेवढेच सूक्ष्म निस्यंदन मिळते. केवळ $C_1(u,v)$
पेक्षा $C_1(u,v)$ आणि $C_2(u,v)$ या दोन्ही अटी पूर्ण
करणाऱ्या जगांच्या जोड्या कमी किंवा तितक्याच असल्याने हे निस्यंदन
सर्वांत स्थूल निस्यंदनापेक्षा “किमान तितके सूक्ष्म” असते.

\begin{table}[ht]
  \centering
  \begin{tabular}{|ll|}
    \hline
    \multirow{2}{*}{$C_1(u,v)$:}
    &
      $\Box\varphi \in \Gamma$ आणि $\mSat{M}{\Box\varphi}[u]$ असेल,
      तर $\mSat{M}{\varphi}[v]$; आणि\\
    &
      $\Diamond\varphi \in \Gamma$ आणि $\mSat{M}{\varphi}[v]$ असेल,
      तर $\mSat{M}{\Diamond\varphi}[u]$; \\
    \hline
    \multirow{2}{*}{$C_2(u,v)$:}
    &
      $\Box\varphi \in \Gamma$ आणि $\mSat{M}{\Box\varphi}[v]$ असेल,
      तर $\mSat{M}{\varphi}[u]$; आणि\\
    &
      $\Diamond\varphi \in \Gamma$ आणि $\mSat{M}{\varphi}[u]$ असेल,
      तर $\mSat{M}{\Diamond\varphi}[v]$; \\
    \hline
    \multirow{2}{*}{$C_3(u,v)$:}
    &
      $\Box\varphi \in \Gamma$ आणि $\mSat{M}{\Box\varphi}[u]$ असेल,
      तर $\mSat{M}{\Box\varphi}[v]$; आणि\\
    &
      $\Diamond\varphi \in \Gamma$ आणि $\mSat{M}{\Diamond\varphi}[v]$ असेल,
      तर $\mSat{M}{\Diamond\varphi}[u]$; \\
    \hline
    \multirow{2}{*}{$C_4(u,v)$:}
    &
      $\Box\varphi \in \Gamma$ आणि $\mSat{M}{\Box\varphi}[v]$ असेल,
      तर $\mSat{M}{\Box\varphi}[u]$; आणि\\
    &
      $\Diamond\varphi \in \Gamma$ आणि $\mSat{M}{\Diamond\varphi}[u]$ असेल,
      तर $\mSat{M}{\Diamond\varphi}[v]$; \\
    \hline
    \end{tabular}
  \caption{निस्यंदने परिभाषित करण्यासाठी संभाव्य जगांवरील अटी.}
  \label{nml:fil:acc:tab:Cn-filtrations}
\end{table}

\begin{thm}\label{nml:fil:acc:thm:more-filtrations}
  $\mModel{M} =\tuple{W,R,V}$ हे प्रतिमान आणि
  $\Gamma$ हा उप-सूत्रे खाली बंद संच असू द्या.
  $W^*$ आणि $V^*$ यांची व्याख्या
  \ref{nml:fil:fil:defn:filtration} प्रमाणे करा. मग:
  \begin{enumerate}
  \item $R^*[u][v]$ तेव्हा आणि तेव्हाच $C_1(u, v) \land
    C_2(u,v)$ असे समजा. मग $R^*$ सममित आहे; आणि
    $\mModel{M}$ सममित असेल, तर $\mModel{M^*} =
    \tuple{W^*,R^*,V^*}$ हे निस्यंदन आहे.
  \item $R^*[u][v]$ तेव्हा आणि तेव्हाच $C_1(u, v)
    \land C_3(u,v)$ असे समजा. मग $R^*$ संक्रमक आहे;
    आणि $\mModel{M}$ संक्रमक असेल, तर
    $\mModel{M^*}=\tuple{W^*,R^*,V^*}$ हे निस्यंदन आहे.
  \item $R^*[u][v]$ तेव्हा आणि तेव्हाच $C_1(u, v) \land C_2(u,v)
    \land C_3(u,v) \land C_4(u,v)$ असे समजा. मग $R^*$
    सममित आणि संक्रमक आहे; आणि $\mModel{M}$ सममित
    व संक्रमक असेल, तर $\mModel{M^*}=\tuple{W^*,R^*,V^*}$
    हे निस्यंदन आहे.
  \item $R^*$ ची व्याख्या $R^*[u][v]$ तेव्हा आणि तेव्हाच $C_1(u,
    v) \land C_3(u,v) \land C_4(u,v)$ अशी करा. मग $R^*$
    संक्रमक आणि युक्लिडी आहे; आणि $\mModel{M}$ संक्रमक
    व युक्लिडी असेल, तर $\mModel{M^*}=\tuple{W^*,R^*,V^*}$
    हे निस्यंदन आहे.
  \end{enumerate}
\end{thm}

\begin{proof}
  \begin{enumerate}
    \item $R^*$ सममित आहे हे लगेच दिसते, कारण $C_1(u,v)
      \Leftrightarrow C_2(v,u)$ आणि $C_2(u,v) \Leftrightarrow
      C_1(v,u)$. त्यामुळे $\mModel{M}$ सममित असेल,
      तर $\mModel{M^*}$ हे $\Gamma$ मधून मिळणारे
      निस्यंदन आहे, एवढेच दाखवायचे आहे.
      $C_1(u,v)$ या अटीमुळे


          \ref{nml:fil:fil:defn:filtration-R2} आणि
          \ref{nml:fil:fil:defn:filtration-R3} या
          \ref{nml:fil:fil:defn:filtration} मधील अटी
      पूर्ण होते. म्हणून आता फक्त
      \ref{nml:fil:fil:defn:filtration}\ref{nml:fil:fil:defn:filtration-R1},
      म्हणजे $Ruv$ वरून $R^*[u][v]$ मिळते हे तपासायचे आहे.

      $Ruv$ असे समजा. $R^*[u][v]$ दाखवण्यासाठी
      $C_1(u,v)$ आणि $C_2(u,v)$ सिद्ध कराव्या लागतील.
      $C_1$ साठी: $\Box\varphi
        \in\Gamma$ आणि $\mSat{M}{\Box\varphi}[u]$ असेल,
        तर $Ruv$ असल्याने $\mSat{M}{\varphi}[v]$ सुद्धा.
         त्याचप्रमाणे,
        $\Diamond\varphi \in \Gamma$
        आणि $\mSat{M}{\varphi}[v]$ असेल, तर $Ruv$
        असल्याने $\mSat{M}{\Diamond\varphi}[u]$.
      $C_2$ साठी: $\Box\varphi \in \Gamma$
        आणि $\mSat{M}{\Box\varphi}[v]$ असेल, तर सममितीमुळे
        $Ruv$ वरून $Rvu$; म्हणून $\mSat{M}{\varphi}[u]$.
         त्याचप्रमाणे,
        $\Diamond\varphi \in\Gamma$
        आणि $\mSat{M}{\varphi}[u]$ असेल, तर
        $\mSat{M}{\Diamond\varphi}[v]$ (कारण सममितीमुळे $Rvu$).
    \item सराव.
    \item सराव.
    \item सराव.
  \end{enumerate}
\end{proof}

\begin{prob}
  \ref{nml:fil:acc:thm:more-filtrations} ची सिद्धता पूर्ण करा.
\end{prob}













\section{युक्लिडी प्रतिमानांची निस्यंदने}\label{nml:fil:euc:sec}

\ref{nml:fil:acc:sec} मधील पद्धत युक्लिडी, किंवा
सर्वत्र उत्तरवर्ती आणि युक्लिडी, प्रतिमानांसाठी
चालत नाही. \ref{nml:fil:euc:fig:ser-eucl} च्या वरच्या भागातील
प्रतिमान पाहा; ते युक्लिडी आणि सर्वत्र उत्तरवर्ती आहे.
$\Gamma = \{p, \Box p \}$ असू द्या. $\Gamma$ मधून
निस्यंदन घेतल्यावर $[w_1] = [w_3]$, कारण $w_1$
आणि $w_3$ हीच दोन जगे~$\Gamma$ वरील सत्यतेबाबत
एकमत आहेत. कोणत्याही निस्यंदनात~$\mModel{M}$
मधून वारशाने आलेले बाणही असतील; ते
\ref{nml:fil:euc:fig:ser-eucl2} मध्ये दाखवले आहेत. ते प्रतिमान
युक्लिडी नाही. शिवाय ते युक्लिडी करण्यासाठी
त्यात आणखी बाण घालता येत नाहीत. $[w_2]$
आणि $[w_4]$ यांच्यामध्ये दोन्ही दिशांनी बाण
घालावे लागतील; मग $[w_2]$ आणि~$[w_5]$
यांच्यामध्येही. पण~$w_2$ येथे $\Box p$ सत्य
असायला हवे, तर~$w_5$ येथे $p$ असत्य आहे.

\begin{figure}[htpb]
  \centering
  \begin{tikzpicture}[modal]
    \node[world] (w1)
          [label={left: \mFalse{p}},
            label={below: $\mSat{{}}{\Box p}$}] {$w_1$};
    \node[world] (w2)
          [label={right: \mTrue{p}},
            label = {below: $\mSat{{}}{\Box p}$},
            right=of w1] {$w_2$};
    \draw[reflexive above] (w2) to (w2);
    \draw[->] (w1) to (w2);
    \node[world] (w3)
          [label={left: \mFalse{p}},
            label={below: $\mSat{{}}{\Box p}$},
            below=of w1] {$w_3$};
    \node[world] (w4)
          [label={right:\mTrue{p}},
            label={below: $\mSat/{{}}{\Box p}$},
            right=of w3] {$w_4$};
    \draw[->] (w3) to (w4);
    \draw[reflexive above] (w4) to (w4);
    \node[world] (w5)
          [label={right: \mFalse{p}},
            label=below:$\mSat/{{}}{\Box p}$,
            right=of w4] {$w_5$};
    \draw[->, bend left] (w4) to (w5);
    \draw[->, bend left] (w5) to (w4);
    \draw[reflexive above] (w5) to (w5);
  \end{tikzpicture}
  \caption{सर्वत्र उत्तरवर्ती आणि युक्लिडी प्रतिमान.}
  \label{nml:fil:euc:fig:ser-eucl}
\end{figure}

\begin{figure}[ht]
  \centering
  \begin{tikzpicture}[modal,every text node part/.style={align=center}]
    \node[world] (w1)
          [label={left: \mFalse{p}},
            label={right:$[w_1]=[w_3]$},
            label={below: $\mSat{{}}{\Box p}$}] {$[w_1]$};
    \node[world] (w2)
         [label={right: \mTrue{p}},
           label = {below: $\mSat{{}}{\Box p}$},
           right=of w1, yshift=1.5cm] {$[w_2]$};
    \draw[reflexive above] (w2) to (w2);
    \draw[->] (w1) to (w2);
    \node[world] (w4)
         [label={right:\mTrue{p}},
           label={below: $\mSat/{{}}{\Box p}$},
           right=of w1,yshift=-1.5cm] {$[w_4]$};
    \draw[->] (w1) to (w4);
    \draw[reflexive above] (w4) to (w4);
    \node[world] (w5)
         [label={right: \mFalse{p}},
             label=below:$\mSat/{{}}{\Box p}$,
             right=of w4] {$[w_5]$};
    \draw[->, bend left] (w4) to (w5);
    \draw[->, bend left] (w5) to (w4);
    \draw[reflexive above] (w5) to (w5);
  \end{tikzpicture}
  \caption{\ref{nml:fil:euc:fig:ser-eucl} मधील प्रतिमानाचे निस्यंदन.}
  \label{nml:fil:euc:fig:ser-eucl2}
\end{figure}

विशेषतः युक्लिडी निस्यंदन मिळवण्यासाठी
उप-सूत्रे खाली बंद असलेल्या मनमानी $\Gamma$
मधून निस्यंदन घेणे पुरेसे नाही. त्याऐवजी
\emph{मोडलदृष्ट्या बंद} असलेले संच~$\Gamma$
घ्यावे लागतात (\ref{nml:fil:pre:defn:modallyclosed} पाहा).
असे रिक्त नसलेले वाक्यांचे संच अनंत असतात; म्हणून त्यांच्यापासून
संबंधित प्रणालीचा सांत प्रतिमान गुणधर्म किंवा
निर्णेयता तात्काळ मिळत नाही.

\begin{thm}\label{nml:fil:euc:thm:modal-closed-filt}
  $\Gamma$ मोडलदृष्ट्या बंद असू द्या,
  $\mModel{M}=\tuple{W,R,V}$ हे प्रतिमान असू द्या,
  आणि $\mModel{M^*} = \tuple{W^*,R^*,V^*}$ हे
  $\mModel{M}$ चे सर्वांत स्थूल निस्यंदन असू द्या.
  \begin{enumerate}
  \item $\mModel{M}$ सममित असेल, तर $\mModel{M^*}$ सुद्धा सममित आहे.
  \item $\mModel{M}$ संक्रमक असेल, तर $\mModel{M^*}$ सुद्धा संक्रमक आहे.
  \item $\mModel{M}$ युक्लिडी असेल, तर $\mModel{M^*}$ सुद्धा युक्लिडी आहे.
  \end{enumerate}
\end{thm}

\begin{proof}
\begin{enumerate}
  \item सराव. \Ax{B} आणि $\Ax{B_\Diamond}$ ही दोन्ही
    सर्व सममित प्रतिमानांत वैध आहेत, हे वापरा.
  \item $\mModel{M^*}$ हे सर्वांत स्थूल निस्यंदन असेल,
    तर व्याख्येनुसार $R^*[u][v]$ तेव्हा आणि तेव्हाच
    $C_1(u,v)$. संक्रमकतेसाठी $C_1(u,v)$ आणि
    $C_1(v,w)$ असे समजा; $C_1(u,w)$ दाखवायचे आहे.
    $\Box\varphi\in\Gamma$ आणि
      $\mSat{M}{\Box \varphi}[u]$ असे समजा.
      सर्व संक्रमक प्रतिमानांत \Ax{4} वैध असल्याने
      $\mSat{M}{\Box\Box\varphi}[u]$. मोडल बंदतेवरून
      $\Box\Box\varphi \in \Gamma$; म्हणून $C_1(u,v)$ वरून
      $\mSat{M}{\Box\varphi}[v]$, आणि $C_1(v,w)$ वरून
      $\mSat{M}{\varphi}[w]$.
      $\Diamond\varphi\in\Gamma$ आणि $\mSat{M}{\varphi}[w]$ असे समजा.
      उपसूत्रांखालील बंदतेवरून $\varphi\in\Gamma$; मोडल बंदतेवरून
      $\Diamond \varphi \in \Gamma$; म्हणून $C_1(v,w)$ वरून
      $\mSat{M}{\Diamond \varphi}[v]$. तसेच मोडल बंदतेवरून
      $\Diamond\Diamond\varphi \in \Gamma$; म्हणून $C_1(u,v)$ वरून
      $\mSat{M}{\Diamond\Diamond \varphi}[u]$.
      सर्व संक्रमक प्रतिमानांत $\Ax{4}_\Diamond$
      वैध असल्याने $\mSat{M}{\Diamond\varphi}[u]$.
  \item सराव. \Ax{5} आणि $\Ax{5_\Diamond}$ ही दोन्ही
    सर्व युक्लिडी प्रतिमानांत वैध आहेत, हे वापरा.
\end{enumerate}
\end{proof}

\begin{prob}
  \ref{nml:fil:euc:thm:modal-closed-filt} ची सिद्धता पूर्ण करा.
\end{prob}



\chapter{मोडल टॅब्लो}\label{nml:tab:chap}









\section{प्रस्तावना}\label{nml:tab:int:sec}

टॅब्लो म्हणजे चिन्हांकित सूत्रांचे (खाली शाखा फुटणारे) काही वृक्ष.
चिन्हांकित सूत्र म्हणजे सत्यतामूल्याचे चिन्ह
($\True$ किंवा $\False$) आणि वाक्य
यांची जोडी:
\[\sFmla{\True}{\varphi} \text{किंवा} \sFmla{\False}{\varphi}.
\]
टॅब्लोची सुरुवात काही \emph{गृहीतकां}पासून
होते. पुढील प्रत्येक चिन्हांकित सूत्र हे अनुमानाचा
एखादा नियम लावून मिळते. काही नियम वृक्षाच्या
टोकावर एक किंवा अधिक चिन्हांकित सूत्रे जोडतात;
इतर नियम दोन नवी टोके जोडून दोन शाखा तयार करतात.
नियम लावल्यावर मिळणाऱ्या चिन्हांकित सूत्रे मधील
सूत्र, नियम ज्या चिन्हांकित सूत्राला लावला
त्यातील सूत्रापेक्षा कमी गुंतागुंतीचे असते.
एखाद्या शाखेत $\sFmla{\True}{\varphi}$ आणि
$\sFmla{\False}{\varphi}$ ही दोन्ही असतील, तर ती शाखा
\emph{बंद} आहे असे म्हणतो. टॅब्लो मधील प्रत्येक
शाखा बंद असेल, तर संपूर्ण टॅब्लो बंद असतो.
बंद टॅब्लो ही अशी निष्पत्ती असते की
तिच्यावरून त्या टॅब्लोची सुरुवात ज्या
चिन्हांकित सूत्रांनी केली तो संच अपूर्ततायोग्य
आहे हे दिसते. यावरून $\Proves$ संबंधाची व्याख्या
करता येते: $\Gamma \Proves \varphi$ तेव्हा आणि तेव्हाच
$\Gamma_0 = \{\psi_1,
\dots, \psi_n\} \subseteq \Gamma$ असा एखादा सांत
संच असतो की पुढील गृहीतकांसाठी बंद टॅब्लो
मिळतो:
\[\{\sFmla{\False}{\varphi}, \sFmla{\True}{\psi_1}, \dots, \sFmla{\True}{\psi_n}\}.
\]

मोडल तर्कशास्त्रांसाठी चिन्हांकित सूत्र ही संकल्पना विस्तारित करावी लागते,
तसेच~$\Box$ आणि
    $\Diamond$ यांसाठी नियम जोडावे लागतात.
चिन्हाबरोबर ($\True$ किंवा $\False$) मोडल
टॅब्लो मधील सूत्रे ना
\emph{पूर्वचिन्हे}~$\sigma$ ही असतात.
ही पूर्वचिन्हे धन पूर्णांकांच्या रिकाम्या नसलेल्या
क्रमिका असतात; म्हणजे $\sigma \in (\PosInt)^* \setminus
\{\emptyseq\}$. अशी पूर्वचिन्हे लिहिताना भोवतालचे
$\tuple{\ }$ वगळतो आणि वेगवेगळ्या घटक
मध्ये $,$ ऐवजी $.$ वापरतो.
$\sigma$ हे पूर्वचिन्ह असेल, तर $\sigma.n$ म्हणजे
$\sigma \concat \tuple{n}$; उदाहरणार्थ
$\sigma = 1.2.1$ असेल, तर $\sigma.3$ म्हणजे
$1.2.1.3$. उदाहरणार्थ,
\[\sFmla{\True}{\Box \varphi \lif \varphi}[1.2]
\]
हे \emph{पूर्वचिन्हयुक्त चिन्हांकित सूत्र}
(किंवा संक्षेपाने \emph{पूर्वचिन्हयुक्त सूत्र}) आहे.

सहज अर्थाने, टॅब्लोच्या एखाद्या शाखेवरील
सूत्रे पूर्ण करू शकेल अशा प्रतिमानातील
जगाचे नाव पूर्वचिन्ह देते. आणि $\sigma$ हे
एखाद्या जगाचे नाव असेल, तर $\sigma.n$ हे
म्हणजे~$\sigma$ ने दर्शवलेल्या जगापासून प्राप्य
असलेल्या जगाचे नाव देते.













\section{\Ax{K} साठीचे नियम}\label{nml:tab:rul:sec}

नेहमीच्या विधानीय कारकांचे नियम नेहमीच्या विधानीय
चिन्हांकित टॅब्लो मधील नियमांसारखेच आहेत;
फक्त त्यांना पूर्वचिन्हे जोडली जातात. प्रत्येक वेळी,
$\sFmla{S}{\varphi}[\sigma]$ या चिन्हांकित सूत्र
ला नियम लावल्याने मिळणाऱ्या नव्या सूत्रे
वरही~$\sigma$ हेच पूर्वचिन्ह असते. हे सहज लक्षात
येते: उदाहरणार्थ, $\varphi \land \psi$ हे~$\sigma$ ने
नाव दिलेल्या जगात सत्य असेल, तर $\varphi$ आणि $\psi$
ही दोन्ही~$\sigma$ येथे सत्य असतात (इतर कोणत्याही
जगात असतीलच असे नाही). विधानीय नियम
\ref{nml:tab:rul:tab:prop-rules} मध्ये दिले आहेत.

\begin{table}
  \[\def\arraystretch{3}\begin{array}{|c|c|}
    \hline
    \AxiomC{\sFmla{\True}{\lnot \varphi}[\sigma]}
    \RightLabel{\TRule{\True}{\lnot}}
    \UnaryInfC{\sFmla{\False}{\varphi}[\sigma]}
    \DisplayProof
    &
    \AxiomC{\sFmla{\False}{\lnot \varphi}[\sigma]}
    \RightLabel{\TRule{\False}{\lnot}}
    \UnaryInfC{\sFmla{\True}{\varphi}[\sigma]}
    \DisplayProof
    \\[1ex]
    \hline
    \AxiomC{\sFmla{\True}{\varphi \land \psi}[\sigma]}
    \RightLabel{\TRule{\True}{\land}}
    \UnaryInfC{\sFmla{\True}{\varphi}[\sigma]}
    \noLine
    \UnaryInfC{\sFmla{\True}{\psi}[\sigma]}
    \DisplayProof
    &
    \AxiomC{\sFmla{\False}{\varphi \land \psi}[\sigma]}
    \RightLabel{\TRule{\False}{\land}}
    \UnaryInfC{$\sFmla{\False}{\varphi}[\sigma] \quad \mid \quad
      \sFmla{\False}{\psi}[\sigma]$}
    \DisplayProof
    \\[2ex]
    \hline
    \AxiomC{\sFmla{\True}{\varphi \lor \psi}[\sigma]}
    \RightLabel{\TRule{\True}{\lor}}
    \UnaryInfC{$\sFmla{\True}{\varphi}[\sigma] \quad \mid \quad
      \sFmla{\True}{\psi}[\sigma]$}
    \DisplayProof
    &
    \AxiomC{\sFmla{\False}{\varphi \lor \psi}[\sigma]}
    \RightLabel{\TRule{\False}{\lor}}
    \UnaryInfC{\sFmla{\False}{\varphi}[\sigma]}
    \noLine
    \UnaryInfC{\sFmla{\False}{\psi}[\sigma]}
    \DisplayProof
    \\[2ex]
    \hline
    \AxiomC{\sFmla{\True}{\varphi \lif \psi}[\sigma]}
    \RightLabel{\TRule{\True}{\lif}}
    \UnaryInfC{$\sFmla{\False}{\varphi}[\sigma] \quad \mid
      \quad \sFmla{\True}{\psi}[\sigma]$}
    \DisplayProof
    &
    \AxiomC{\sFmla{\False}{\varphi \lif \psi}[\sigma]}
    \RightLabel{\TRule{\False}{\lif}}
    \UnaryInfC{\sFmla{\True}{\varphi}[\sigma]}
    \noLine
    \UnaryInfC{\sFmla{\False}{\psi}[\sigma]}
    \DisplayProof
    \\[2ex]
    \hline
  \end{array}\]
  \caption{विधानीय कारकांसाठी पूर्वचिन्हयुक्त
    टॅब्लो नियम}
  \label{nml:tab:rul:tab:prop-rules}
\end{table}

शाखा बंद होण्याची अट नेहमीच्या टॅब्लो
प्रमाणेच आहे; पण येथे सूत्रे बरोबर त्यांची
पूर्वचिन्हेही जुळली पाहिजेत. म्हणजे एखाद्या शाखेत
पुढील दोन्ही असतील, तर ती बंद असते:
\[\sFmla{\True}{\varphi}[\sigma] \quad\text{आणि}\quad \sFmla{\False}{\varphi}[\sigma]
\]
येथे $\sigma$ हे एखादे पूर्वचिन्ह आणि $\varphi$ हे
सूत्र आहे.

गृहीतके मांडण्याचे नियमही नेहमीच्या टॅब्लो
प्रमाणेच आहेत; मात्र गृहीतकांसाठी आपण नेहमी
$1$ हे पूर्वचिन्ह वापरतो. (सर्व गृहीतकांसाठी एकच
पूर्वचिन्ह वापरले, तर नेमके कोणते पूर्वचिन्ह
वापरले याने फरक पडत नाही.) म्हणून, उदाहरणार्थ,
\[\psi_1, \dots, \psi_n \Proves \varphi
\]
असे तेव्हा आणि तेव्हाच म्हणतो जेव्हा पुढील
गृहीतकांसाठी बंद टॅब्लो असतो:
\[\sFmla{\True}{\psi_1}[1], \dots, \sFmla{\True}{\psi_n}[1],
\sFmla{\False}{\varphi}[1].
\]

मोडल कारक~$\Box$
    आणि~$\Diamond$
यांसाठी,~$\sigma$ हे पूर्वचिन्ह असलेल्या सूत्र
ला नियम लावल्यावर निष्कर्षाचे पूर्वचिन्ह
$\sigma.n$ असते. पण कोणता $n$ चालेल हे
चिन्ह~$\True$ आहे की~$\False$ यावर अवलंबून असते.

$\TRule{\True}{\Box}$ हा नियम
  $\sFmla{\True}{\Box \varphi}[\sigma]$ असलेल्या शाखेत
  $\sFmla{\True}{\varphi}[\sigma.n]$ जोडतो.
    त्याचप्रमाणे, $\TRule{\False}{\Diamond}$
  हा नियम $\sFmla{\False}{\Diamond \varphi}[\sigma]$
  असलेल्या शाखेत $\sFmla{\False}{\varphi}[\sigma.n]$
  जोडतो.
हे नियम
लावताना $\sigma.n$ हे पूर्वचिन्ह संबंधित शाखेत
\emph{आधीपासूनच} आलेले असले पाहिजे. अशा
पूर्वचिन्हाला त्या शाखेवर ‘वापरलेले’ म्हणू.

$\TRule{\False}{\Box}$ हा नियम
  $\sFmla{\False}{\Box \varphi}[\sigma]$ असलेल्या शाखेत
  $\sFmla{\False}{\varphi}[\sigma.n]$ जोडतो.
    त्याचप्रमाणे, $\TRule{\True}{\Diamond}$
  हा नियम $\sFmla{\True}{\Diamond \varphi}[\sigma]$
  असलेल्या शाखेत $\sFmla{\True}{\varphi}[\sigma.n]$
  जोडतो.
परंतु हे नियम
लावताना $\sigma.n$ हे पूर्वचिन्ह संबंधित शाखेत
\emph{आधीपासून आलेले नसले} पाहिजे. अशा
पूर्वचिन्हांना त्या शाखेसाठी ‘नवी’ म्हणतो.

हे नियम \ref{nml:tab:rul:tab:rules-K} मध्ये दिले आहेत.

\begin{table}
  \begin{center}
    \def\arraystretch{3}\def\fCenter{}
    \begin{tabular}{|c|c|}
    \hline

      \AxiomC{\sFmla{\True}{\Box \varphi}[\sigma]}
      \RightLabel{\TRule{\True}{\Box}}
      \UnaryInfC{\sFmla{\True}{\varphi}[\sigma.n]}
      \DisplayProof
      &
      \AxiomC{\sFmla{\False}{\Box \varphi}[\sigma]}
      \RightLabel{\TRule{\False}{\Box}}
      \UnaryInfC{\sFmla{\False}{\varphi}[\sigma.n]}
      \DisplayProof\\
      $\sigma.n$ वापरलेले आहे & $\sigma.n$ नवे आहे
      \\[1ex]
      \hline
      \AxiomC{\sFmla{\True}{\Diamond \varphi}[\sigma]}
      \RightLabel{\TRule{\True}{\Diamond}}
      \UnaryInfC{\sFmla{\True}{\varphi}[\sigma.n]}
      \DisplayProof
      &
      \AxiomC{\sFmla{\False}{\Diamond \varphi}[\sigma]}
      \RightLabel{\TRule{\False}{\Diamond}}
      \UnaryInfC{\sFmla{\False}{\varphi}[\sigma.n]}
      \DisplayProof\\
      $\sigma.n$ नवे आहे & $\sigma.n$ वापरलेले आहे
      \\[1ex]
      \hline
    \end{tabular}
  \end{center}
  \caption{\Ax{K} साठीचे मोडल नियम.}
  \label{nml:tab:rul:tab:rules-K}
\end{table}
\TRule{\True}{\Box}
या नियमासाठी पूर्वचिन्ह वापरलेले असणे आवश्यक आहे.
ही अट नसेल, तर पुढील टॅब्लो बंद असल्याचे
आपण चुकीने मानू:
  \begin{oltableau}
    [\pFmla{\True}{\Box \formula{A}}{1}, just = \TAss
      [\pFmla{\False}{\Diamond \formula{A}}{1}, just = \TAss
        [\pFmla{\True}{\formula{A}}{1.1}, just = {\TRule{\True}{\Box}[1]}
          [\pFmla{\False}{\formula{A}}{1.1},
            just = {\TRule{\False}{\Diamond}[2]}, close]
        ]
      ]
    ]
\end{oltableau}
पण $\Box \formula{A} \Entails/ \Diamond \formula{A}$;
म्हणून असे झाल्यास आपली प्रमाणपद्धती अप्रमाण असेल.
त्याचप्रमाणे, $\Diamond \formula{A} \Entails/
\Box \formula{A}$; परंतु
\TRule{\False}{\Box}
या नियमासाठी पूर्वचिन्ह नवे असण्याची अट नसेल,
तर पुढील टॅब्लो बंद होईल:
  \begin{oltableau}
    [\pFmla{\True}{\Diamond \formula{A}}{1}, just = \TAss,name=one
      [\pFmla{\False}{\Box \formula{A}}{1}, just = \TAss, name=two
        [\pFmla{\True}{\formula{A}}{1.1}, just = {\TRule{\True}{\Diamond}[1]}
          [\pFmla{\False}{\formula{A}}{1.1},
            just = {\TRule{\False}{\Box}[2]}, close]
        ]
      ]
    ]
  \end{oltableau}
\section{\Log{K} साठी टॅब्लो}\label{nml:tab:prk:sec}
\begin{ex}
  $\Proves (\Box\varphi \land \Box\psi)
  \lif \Box (\varphi \land \psi)$ हे दाखवणारा बंद टॅब्लो
  पुढीलप्रमाणे आहे.
  \begin{oltableau}
    [\pFmla{\False}{(\Box\formula{A} \land \Box\formula{B}) \lif
        \Box (\formula{A} \land \formula{B})}{1},
      just =\TAss
      [\pFmla{\True}{\Box\formula{A} \land \Box\formula{B}}{1},
        just = {\TRule{\False}{\lif}[1]}
        [\pFmla{\False}{\Box(\formula{A} \land \formula{B})}{1},
          just = {\TRule{\False}{\lif}[1]}
          [\pFmla{\True}{\Box\formula{A}}{1},
            just = {\TRule{\True}{\land}[2]}
            [\pFmla{\True}{\Box\formula{B}}{1},
              just = {\TRule{\True}{\land}[2]}
              [\pFmla{\False}{\formula{A} \land \formula{B}}{1.1},
                just = {\TRule{\False}{\Box}[3]}
                [\pFmla{\False}{\formula{A}}{1.1},
                  just = {\TRule{\False}{\land}[6]}
                  [\pFmla{\True}{\formula{A}}{1.1},
                    just= {\TRule{\True}{\Box}[4]}, close]]
                [\pFmla{\False}{\formula{B}}{1.1},
                  just = {\TRule{\False}{\land}[6]}
                  [\pFmla{\True}{\formula{B}}{1.1},
                    just= {\TRule{\True}{\Box}[5]}, close]]
              ]
            ]
          ]
        ]
      ]
    ]
  \end{oltableau}
\end{ex}
\begin{ex}
  $\Proves \Diamond(\varphi \lor \psi)
  \lif (\Diamond \varphi \lor \Diamond \psi)$ हे दाखवणारा
  बंद टॅब्लो पुढीलप्रमाणे आहे:
  \begin{oltableau}
    [\pFmla{\False}{\Diamond(\formula{A} \lor \formula{B}) \lif
        (\Diamond \formula{A} \lor \Diamond \formula{B})}{1},
      just =\TAss
      [\pFmla{\True}{\Diamond(\formula{A} \lor \formula{B})}{1},
        just = {\TRule{\False}{\lif}[1]}
        [\pFmla{\False}{\Diamond\formula{A} \lor \Diamond\formula{B}}{1},
          just = {\TRule{\False}{\lif}[1]}
          [\pFmla{\False}{\Diamond\formula{A}}{1},
            just = {\TRule{\False}{\lor}[3]}
            [\pFmla{\False}{\Diamond\formula{B}}{1},
              just = {\TRule{\False}{\lor}[3]}
              [\pFmla{\True}{\formula{A} \lor \formula{B}}{1.1},
                just = {\TRule{\True}{\Diamond}[2]},
                [\pFmla{\True}{\formula{A}}{1.1},
                  just = {\TRule{\True}{\lor}[6]}
                  [\pFmla{\False}{\formula{A}}{1.1},
                    just= {\TRule{\False}{\Diamond}[4]}, close]]
                [\pFmla{\True}{\formula{B}}{1.1},
                  just = {\TRule{\True}{\lor}[6]}
                  [\pFmla{\False}{\formula{B}}{1.1},
                    just= {\TRule{\False}{\Diamond}[5]}, close]]
              ]
            ]
          ]
        ]
      ]
    ]
  \end{oltableau}
\end{ex}
\begin{prob}
  पुढील सूत्रे साठी~$\Log{K}$ मध्ये बंद
  टॅब्लो शोधा:
  \begin{enumerate}
    \item $\Box \lnot p \lif \Box(p \lif q)$
    \item $(\Box p \lor \Box q) \lif \Box(p \lor q)$
    \item $\Diamond p \lif \Diamond(p \lor q)$
    \item $\Box(p \land q) \lif \Box p$
  \end{enumerate}
\end{prob}
\section{\Log{K} साठी निर्दोषता}\label{nml:tab:sou:sec}
\begin{editorial}
  निर्दोषतेचा हा पुरावा अभिजात विधानीय तर्कशास्त्राचा
  निर्दोषतेचा पुरावा पुन्हा मांडतो; म्हणजे सर्व काही
  सुरुवातीपासून सिद्ध करतो. स्वयंपूर्ण पुरावा हवा
  असेल, तर हे योग्य आहे. नेहमीच्या टॅब्लोची निर्दोषता
  तुम्ही आधी पाहिली असेल, तर येथे पुनरुक्ती वाटेल.
  पुढे स्वयंपूर्ण आवृत्ती आणि अमोडल प्रकरणावर आधारलेली
  आवृत्ती यांपैकी एक निवडता यावी, अशी योजना आहे.
\end{editorial}
\begin{explain}
  पूर्वचिन्हयुक्त टॅब्लो निर्दोष आहेत हे
  दाखवण्यासाठी, पुढील गृहीतकांचा
  \[
  \sFmla{\True}{\psi_1}[1], \dots, \sFmla{\True}{\psi_n}[1], \sFmla{\False}{\varphi}[1]
  \]
  बंद टॅब्लो असेल, तर $\psi_1, \dots, \psi_n
  \Entails \varphi$ हे सिद्ध करावे लागेल. व्यत्यास सिद्ध
  करणे सोपे आहे: एखाद्या $\mModel{M}$ मध्ये
  $w$ या जगात सर्व $i=1$, \dots,~$n$ साठी
  $\mSat{M}{\psi_i}[w]$, पण $\mSat/{M}{\varphi}[w]$
  असेल, तर कोणताही टॅब्लो बंद होऊ शकत नाही.
  अशा प्रतिदृष्टांत प्रतिमानावरून टॅब्लोची सुरुवातीची
  गृहीतके पूर्ततायोग्य आहेत हे दिसते. पुराव्याची युक्ती
  अशी: टॅब्लोच्या शाखेवरील सर्व
  पूर्वचिन्हयुक्त सूत्रे पूर्ततायोग्य असतील, तर
  कोणताही नियम लावल्यावर तयार होणाऱ्या विस्तारित
  शाखांपैकी किमान एक शाखा पूर्ततायोग्य असते. बंद शाखा
  अपूर्ततायोग्य असतात. म्हणून पूर्वचिन्हयुक्त सूत्रे
  च्या पूर्ततायोग्य संचासाठीच्या कोणत्याही टॅब्लो
  मध्ये किमान एक खुली शाखा असली पाहिजे.

  मोडल प्रकरणात ही युक्ती वापरण्यासाठी ‘पूर्ततायोग्य’
  ही व्याख्या मोडल प्रतिमाने आणि पूर्वचिन्हे यांच्यापर्यंत
  विस्तारित करावी लागते. त्यानंतर पुरावा सरळ आहे.
\end{explain}

\begin{defn}
  $P$ हा पूर्वचिन्हांचा एखादा संच असू द्या; म्हणजे
  $P \subseteq (\PosInt)^* \setminus
  \{\emptyseq\}$. तसेच $\mModel{M}$ हे प्रतिमान असू द्या.
  $\sigma$ आणि $\sigma.n$ ही दोन्ही~$P$ मध्ये असताना
  नेहमी $Rf(\sigma)f(\sigma.n)$ असेल, तर
  $f\colon P \to W$ या फलनाला $\mModel{M}$
  मधील~$P$ चे \emph{अर्थनिर्धारण} म्हणतात.

  पूर्वचिन्हांच्या~$P$ या संचाच्या अर्थनिर्धारणाच्या
  संदर्भात पुढील व्याख्या करता येतात:
  \begin{enumerate}
  \item $\mModel{M}$ हे $\sFmla{\True}{\varphi}[\sigma]$
    पूर्ण करते तेव्हा आणि तेव्हाच
    $\mSat{M}{\varphi}[f(\sigma)]$.
  \item $\mModel{M}$ हे $\sFmla{\False}{\varphi}[\sigma]$
    पूर्ण करते तेव्हा आणि तेव्हाच
    $\mSat/{M}{\varphi}[f(\sigma)]$.
  \end{enumerate}
\end{defn}

\begin{defn}
  $\Gamma$ हा पूर्वचिन्हयुक्त सूत्रांचा
  संच असू द्या आणि त्यात येणाऱ्या पूर्वचिन्हांचा
  संच $P(\Gamma)$ असू द्या. $f$ हे $\mModel{M}$
  मधील~$P(\Gamma)$ चे अर्थनिर्धारण असेल, तर
  $\mModel{M}$ ने~$f$ च्या संदर्भात~$\Gamma$
  पूर्ण करणे, म्हणजे $\mSat{M}{\Gamma}[f]$,
  याचा अर्थ $\mModel{M}$ ने~$f$ च्या संदर्भात
  $\Gamma$ मधील प्रत्येक पूर्वचिन्हयुक्त सूत्र
  पूर्ण करणे असा आहे. $\mSat{M}{\Gamma}[f]$
  होईल असे एखादे प्रतिमान~$\mModel{M}$ आणि
  $P(\Gamma)$ चे अर्थनिर्धारण~$f$ अस्तित्वात असेल,
  तर आणि तरच $\Gamma$ हा संच \emph{पूर्ततायोग्य} असतो.
\end{defn}

\begin{prop}
  एखाद्या सूत्र~$\varphi$ आणि पूर्वचिन्ह~$\sigma$
  यांसाठी $\Gamma$ मध्ये
  $\sFmla{\True}{\varphi}[\sigma]$ आणि
  $\sFmla{\False}{\varphi}[\sigma]$ ही दोन्ही असतील,
  तर $\Gamma$ अपूर्ततायोग्य आहे.
\end{prop}

\begin{proof}
  $\mSat{M}{\varphi}[f(\sigma)]$ आणि
  $\mSat/{M}{\varphi}[f(\sigma)]$ ही दोन्ही सत्य होतील
  असे प्रतिमान~$\mModel{M}$ आणि $P(\Gamma)$ चे
  अर्थनिर्धारण~$f$ असू शकत नाही.
\end{proof}

\begin{thm}[निर्दोषता]
  \label{nml:tab:sou:thm:tableau-soundness}
  $\Gamma$ साठी बंद टॅब्लो असेल, तर
  $\Gamma$ अपूर्ततायोग्य आहे.
\end{thm}

\begin{proof}
टॅब्लोवरील एखादी शाखा पूर्ततायोग्य असणे म्हणजे
तिच्यावरील चिन्हांकित सूत्रांचा संच पूर्ततायोग्य असणे.
आणि टॅब्लोमध्ये किमान एक पूर्ततायोग्य शाखा
असेल, तर त्याला पूर्ततायोग्य म्हणू.
येथे नियमानुसार तयार केलेल्या टॅब्लोंचाच विचार करतो. सुरुवातीच्या
गृहीतकांचे पूर्वचिन्ह $1$ असते आणि प्रत्येक नवे पूर्वचिन्ह आधी
वापरलेल्या पूर्वचिन्हाच्या शेवटी एक संख्या जोडून तयार होते.
म्हणून शाखेवरील पूर्वचिन्हांचा संच प्रत्येक रिक्त नसलेल्या आरंभीच्या
खंडाखाली बंद असतो.

आपण पुढील विधान दाखवू: पूर्ततायोग्य टॅब्लो ला
अनुमानाचा कोणताही नियम लावून विस्तारित केले, तरी
मिळणारा टॅब्लो पूर्ततायोग्य असतो. यावरून प्रमेय
सिद्ध होईल. कारण बंद टॅब्लो हा
सुरुवातीला फक्त~$\Gamma$ मधील गृहीतके असलेल्या
टॅब्लो ला अनुमानाचे नियम लावत तयार होतो.
म्हणून $\Gamma$ पूर्ततायोग्य असेल, तर त्यासाठीचा
प्रत्येक टॅब्लो पूर्ततायोग्य राहील. पण बंद
टॅब्लो पूर्ततायोग्य असू शकत नाही: त्याच्या सर्व
शाखा बंद असतात, आणि बंद शाखा अपूर्ततायोग्य असतात.

आपल्याकडे पूर्ततायोग्य टॅब्लो आहे असे समजू;
म्हणजे किमान एक पूर्ततायोग्य शाखा असलेला टॅब्लो
आहे. अनुमानाचा
नियम लावल्याने शाखेवर चिन्हांकित सूत्रे
जोडली जातात किंवा शाखेच्या दोन शाखा फुटतात.
नियमामुळे न बदलणारी एखादी पूर्ततायोग्य शाखा
टॅब्लोमध्ये असेल, तर विस्तारित
टॅब्लोमध्येही ती पूर्ततायोग्य राहते. म्हणून
नियम पूर्ततायोग्य शाखेला लावला जातो त्या
प्रकरणाचा विचार पुरेसा आहे.

त्या शाखेवरील चिन्हांकित सूत्रांचा संच
$\Gamma$ असू द्या आणि नियम
$\sFmla{S}{\varphi}[\sigma] \in \Gamma$ या
चिन्हांकित सूत्राला लावला जातो असे समजू.
नियमाने शाखा फुटत नसेल, तर $\Gamma$ मध्ये
नियमाचे निष्कर्ष जोडून मिळालेली विस्तारित
शाखा अजूनही पूर्ततायोग्य आहे हे दाखवावे लागते.
शाखा फुटत असेल, तर मिळणाऱ्या दोन शाखांपैकी
किमान एक पूर्ततायोग्य आहे हे दाखवावे लागते.

प्रथम शाखा न फुटणाऱ्या नियमांचा विचार करू.
\begin{enumerate}
\item $\sFmla{\True}{\lnot \psi}[\sigma] \in \Gamma$
  याला
  $\TRule{\True}{\lnot}$ लावून शाखा विस्तारित
  केली. आता शाखेवरील चिन्हांकित सूत्रांचा
  संच $\Gamma \cup
  \{\sFmla{\False}{\psi}[\sigma]\}$ आहे.
  $\mSat{M}{\Gamma}[f]$ असे समजू. विशेषतः
  $\mSat{M}{\lnot \psi}[f(\sigma)]$; म्हणून
  $\mSat/{M}{\psi}[f(\sigma)]$. म्हणजे
  $\mModel{M}$ हे~$f$ च्या संदर्भात
  $\sFmla{\False}{\psi}[\sigma]$ पूर्ण करते.
\item $\sFmla{\False}{\lnot \psi}[\sigma] \in \Gamma$
  याला
  $\TRule{\False}{\lnot}$ लावून शाखा विस्तारित
  केली: सरावासाठी.
\item $\sFmla{\True}{\psi \land \chi}[\sigma] \in \Gamma$
  याला
  $\TRule{\True}{\land}$ लावून शाखा विस्तारित
  केली. यात शाखेवर
  $\sFmla{\True}{\psi}[\sigma]$ आणि
  $\sFmla{\True}{\chi}[\sigma]$ ही दोन नवी
  चिन्हांकित सूत्रे मिळतात.
  $\mSat{M}{\Gamma}[f]$ असे समजू; विशेषतः
  $\mSat{M}{\psi \land \chi}[f(\sigma)]$.
  म्हणून $\mSat{M}{\psi}[f(\sigma)]$ आणि
  $\mSat{M}{\chi}[f(\sigma)]$. याचा अर्थ
  $\mModel{M}$ हे~$f$ च्या संदर्भात
  $\sFmla{\True}{\psi}[\sigma]$ आणि
  $\sFmla{\True}{\chi}[\sigma]$ ही दोन्ही
  पूर्ण करते.
\item $\sFmla{\False}{\psi \lor \chi}[\sigma] \in \Gamma$
  याला
  $\TRule{\False}{\lor}$ लावून शाखा विस्तारित
  केली: सरावासाठी.
\item $\sFmla{\False}{\psi \lif \chi}[\sigma] \in \Gamma$
  याला
  $\TRule{\False}{\lif}$ लावून शाखा विस्तारित
  केली. यात शाखेवर
  $\sFmla{\True}{\psi}[\sigma]$ आणि
  $\sFmla{\False}{\chi}[\sigma]$ ही दोन नवी
  चिन्हांकित सूत्रे मिळतात.
  $\mSat{M}{\Gamma}[f]$ असे समजू; विशेषतः
  $\mSat/{M}{\psi \lif \chi}[f(\sigma)]$.
  म्हणून $\mSat{M}{\psi}[f(\sigma)]$ आणि
  $\mSat/{M}{\chi}[f(\sigma)]$. याचा अर्थ
  $\mModel{M}, f$ हे
  $\sFmla{\True}{\psi}[\sigma]$ आणि
  $\sFmla{\False}{\chi}[\sigma]$ ही दोन्ही
  पूर्ण करते.


\item $\sFmla{\True}{\Box \psi}[\sigma] \in \Gamma$
  याला
  $\TRule{\True}{\Box}$ लावून शाखा विस्तारित केली:
  शाखेवर
    $\sFmla{\True}{\psi}[\sigma.n]$ हे नवे
    चिन्हांकित सूत्र मिळते. येथे
    $\sigma.n \in P(\Gamma)$, कारण
    $\sigma.n$ वापरलेले असले पाहिजे.
    $\mSat{M}{\Gamma}[f]$ असे समजू; विशेषतः
    $\mSat{M}{\Box \psi}[f(\sigma)]$.
    $f$ हे पूर्वचिन्हांचे अर्थनिर्धारण आहे आणि
    $\sigma$, $\sigma.n \in P(\Gamma)$,
    म्हणून $Rf(\sigma)f(\sigma.n)$.
    त्यामुळे $\mSat{M}{\psi}[f(\sigma.n)]$;
    म्हणजे $\mModel{M}, f$ हे
    $\sFmla{\True}{\psi}[\sigma.n]$ पूर्ण करते.

\item $\sFmla{\False}{\Box \psi}[\sigma] \in \Gamma$
  याला
  $\TRule{\False}{\Box}$ लावून शाखा विस्तारित केली:
  शाखेवर
    $\sFmla{\False}{\psi}[\sigma.n]$ हे नवे
    चिन्हांकित सूत्र मिळते. येथे
    $\sigma.n$ हे शाखेवरील नवे पूर्वचिन्ह,
    म्हणजे $\sigma.n \notin P(\Gamma)$.
    $\Gamma$ पूर्ततायोग्य असल्याने,
    $\mSat{M}{\Gamma}[f]$ होईल असे
    $\mModel{M}$ आणि $P(\Gamma)$ चे
    अर्थनिर्धारण~$f$ असते; विशेषतः
    $\mSat/{M}{\Box \psi}[f(\sigma)]$.
    $\Gamma \cup
    \{\sFmla{\False}{\psi}[\sigma.n]\}$
    पूर्ततायोग्य आहे हे दाखवायचे आहे.
    यासाठी $P(\Gamma) \cup \{\sigma.n\}$
    चे अर्थनिर्धारण पुढीलप्रमाणे ठरवू:

  $\mSat/{M}{\Box \psi}[f(\sigma)]$ असल्याने,
  $Rf(\sigma)w$ आणि $\mSat/{M}{\psi}[w]$
  होईल असे $w \in W$ असते. $f'$ हे~$f$
  प्रमाणेच असू द्या, फक्त
  $f'(\sigma.n) = w$ असे ठरवा.
  $f'(\sigma) = f(\sigma)$ आणि
  $Rf(\sigma)w$ असल्याने
  $Rf'(\sigma)f'(\sigma.n)$.
  पूर्वचिन्हांच्या आरंभीच्या खंडांखालील बंदतेमुळे नव्या
  $\sigma.n$ पासून वाढवलेले कोणतेही पूर्वचिन्ह आधी वापरलेले
  नसते. त्यामुळे अर्थनिर्धारणासाठी तपासायची नवी जोडी फक्त
  $(\sigma,\sigma.n)$ हीच आहे; म्हणून $f'$
  हे $P(\Gamma) \cup \{\sigma.n\}$ चे
  अर्थनिर्धारण आहे. स्पष्टपणे
  $\mSat/{M}{\psi}[f'(\sigma.n)]$.
  प्रत्येक $\sigma' \in P(\Gamma)$ साठी
  $f(\sigma') = f'(\sigma')$ असल्याने
  $\mSat{M}{\Gamma}[f']$. म्हणून
  $\mModel{M}, f'$ हे $\Gamma \cup
  \{\sFmla{\False}{\psi}[\sigma.n]\}$ पूर्ण करते.



\item $\sFmla{\True}{\Diamond \psi}[\sigma] \in \Gamma$
  याला
  $\TRule{\True}{\Diamond}$ लावून शाखा विस्तारित
  केली: शाखेवर
    $\sFmla{\True}{\psi}[\sigma.n]$ हे नवे
    चिन्हांकित सूत्र मिळते. येथे
    $\sigma.n$ हे शाखेवरील नवे पूर्वचिन्ह,
    म्हणजे $\sigma.n \notin P(\Gamma)$.
    $\Gamma$ पूर्ततायोग्य असल्याने,
    $\mSat{M}{\Gamma}[f]$ होईल असे
    $\mModel{M}$ आणि $P(\Gamma)$ चे
    अर्थनिर्धारण~$f$ असते; विशेषतः
    $\mSat{M}{\Diamond \psi}[f(\sigma)]$.
    $\Gamma \cup
    \{\sFmla{\True}{\psi}[\sigma.n]\}$
    पूर्ततायोग्य आहे हे दाखवायचे आहे.
    यासाठी $P(\Gamma) \cup \{\sigma.n\}$
    चे अर्थनिर्धारण पुढीलप्रमाणे ठरवू:

  $\mSat{M}{\Diamond \psi}[f(\sigma)]$ असल्याने,
  $Rf(\sigma)w$ आणि $\mSat{M}{\psi}[w]$
  होईल असे $w \in W$ असते. $f'$ हे~$f$
  प्रमाणेच असू द्या, फक्त
  $f'(\sigma.n) = w$ असे ठरवा.
  $f'(\sigma) = f(\sigma)$ आणि
  $Rf(\sigma)w$ असल्याने
  $Rf'(\sigma)f'(\sigma.n)$.
  पूर्वचिन्हांच्या आरंभीच्या खंडांखालील बंदतेमुळे नव्या
  $\sigma.n$ पासून वाढवलेले कोणतेही पूर्वचिन्ह आधी वापरलेले
  नसते. त्यामुळे अर्थनिर्धारणासाठी तपासायची नवी जोडी फक्त
  $(\sigma,\sigma.n)$ हीच आहे; म्हणून $f'$
  हे $P(\Gamma) \cup \{\sigma.n\}$ चे
  अर्थनिर्धारण आहे. स्पष्टपणे
  $\mSat{M}{\psi}[f'(\sigma.n)]$.
  प्रत्येक $\sigma' \in P(\Gamma)$ साठी
  $f(\sigma') = f'(\sigma')$ असल्याने
  $\mSat{M}{\Gamma}[f']$. म्हणून
  $\mModel{M}, f'$ हे $\Gamma \cup
  \{\sFmla{\True}{\psi}[\sigma.n]\}$ पूर्ण करते.
\item $\sFmla{\False}{\Diamond \psi}[\sigma] \in \Gamma$
  याला
  $\TRule{\False}{\Diamond}$ लावून शाखा विस्तारित
  केली: शाखेवर
    $\sFmla{\False}{\psi}[\sigma.n]$ हे नवे
    चिन्हांकित सूत्र मिळते. येथे
    $\sigma.n \in P(\Gamma)$, कारण
    $\sigma.n$ वापरलेले असले पाहिजे.
    $\mSat{M}{\Gamma}[f]$ असे समजू; विशेषतः
    $\mSat/{M}{\Diamond \psi}[f(\sigma)]$.
    $f$ हे पूर्वचिन्हांचे अर्थनिर्धारण आहे आणि
    $\sigma$, $\sigma.n \in P(\Gamma)$,
    म्हणून $Rf(\sigma)f(\sigma.n)$.
    त्यामुळे $\mSat/{M}{\psi}[f(\sigma.n)]$;
    म्हणजे $\mModel{M}, f$ हे
    $\sFmla{\False}{\psi}[\sigma.n]$ पूर्ण करते.

\end{enumerate}
आता शाखा फुटणाऱ्या नियमांचा विचार करू.
\begin{enumerate}
\item $\sFmla{\False}{\psi \land \chi}[\sigma] \in \Gamma$
  याला
  $\TRule{\False}{\land}$ लावून शाखा विस्तारित
  केली. यात दोन शाखा मिळतात: डावीकडची
  $\sFmla{\False}{\psi}[\sigma]$ मधून, आणि
  उजवीकडची $\sFmla{\False}{\chi}[\sigma]$
  मधून पुढे जाते. $\mSat{M}{\Gamma}[f]$
  असे समजू; विशेषतः
  $\mSat/{M}{\psi \land \chi}[f(\sigma)]$.
  मग $\mSat/{M}{\psi}[f(\sigma)]$ किंवा
  $\mSat/{M}{\chi}[f(\sigma)]$.
  पहिल्या प्रकरणात $\mModel{M}, f$ हे
  $\sFmla{\False}{\psi}[\sigma]$ पूर्ण करते,
  म्हणून डावीकडची शाखा पूर्ततायोग्य असते.
  दुसऱ्या प्रकरणात $\mModel{M}, f$ हे
  $\sFmla{\False}{\chi}[\sigma]$ पूर्ण करते,
  म्हणून उजवीकडची शाखा पूर्ततायोग्य असते.
\item $\sFmla{\True}{\psi \lor \chi}[\sigma] \in \Gamma$
  याला
  $\TRule{\True}{\lor}$ लावून शाखा विस्तारित
  केली: सरावासाठी.
\item $\sFmla{\True}{\psi \lif \chi}[\sigma] \in \Gamma$
  याला
  $\TRule{\True}{\lif}$ लावून शाखा विस्तारित
  केली: सरावासाठी.
\end{enumerate}
\end{proof}

\begin{prob}
\ref{nml:tab:sou:thm:tableau-soundness} चा
पुरावा पूर्ण करा.
\end{prob}

\begin{cor}
\label{nml:tab:sou:cor:entailment-soundness}
$\Gamma \Proves \varphi$ असेल, तर
$\Gamma \Entails \varphi$.
\end{cor}

\begin{proof}
  $\Gamma \Proves \varphi$ असेल, तर
  $\psi_1$, \dots, $\psi_n \in \Gamma$ अशी काही सूत्रे
  असतात की $\Delta =
  \{\sFmla{\False}{\varphi}[1], \sFmla{\True}{\psi_1}[1],
  \dots, \sFmla{\True}{\psi_n}[1]\}$ साठी
  बंद टॅब्लो असतो.
  $\Gamma \Entails \varphi$ हे दाखवायचे आहे.
  तसे नाही असे समजू. मग एखाद्या
  $\mModel{M}$ आणि~$w$ साठी सर्व
  $i=1$, \dots,~$n$ बाबत
  $\mSat{M}{\psi_i}[w]$, पण
  $\mSat/{M}{\varphi}[w]$. $f(1) = w$
  असे ठेवा. मग $f$ हे~$\mModel{M}$ मधील
  $P(\Delta)$ चे अर्थनिर्धारण आहे आणि
  $\mModel{M}$ हे~$f$ च्या संदर्भात
  $\Delta$ पूर्ण करते. पण
  \ref{nml:tab:sou:thm:tableau-soundness} नुसार
  $\Delta$ साठी बंद टॅब्लो असल्यामुळे
  तो अपूर्ततायोग्य आहे. हा विरोधाभास आहे.
  म्हणून $\Gamma \Entails \varphi$.
\end{proof}

\begin{cor}
\label{nml:tab:sou:cor:weak-soundness}
$\Proves \varphi$ असेल, तर $\varphi$ सर्व
प्रतिमानांमध्ये सत्य आहे.
\end{cor}













\section{इतर प्राप्यता संबंधांसाठीचे नियम}\label{nml:tab:mru:sec}

विशिष्ट प्राप्यता संबंधांनी निर्धारित होणाऱ्या
तर्कप्रणालींसाठी \ref{nml:tab:mru:tab:more-rules} मधील
अतिरिक्त नियमांचा विचार करू.

\begin{table}
  \begin{center}
    \def\arraystretch{3}\def\fCenter{}
    \begin{tabular}{|c|c|}
    \hline

      \AxiomC{\sFmla{\True}{\Box \varphi}[\sigma]}
      \RightLabel{T$\Box$}
      \UnaryInfC{\sFmla{\True}{\varphi}[\sigma]}
      \DisplayProof
    &

      \AxiomC{\sFmla{\False}{\Diamond \varphi}[\sigma]}
      \RightLabel{T$\Diamond$}
      \UnaryInfC{\sFmla{\False}{\varphi}[\sigma]}
      \DisplayProof
    \\[1ex]
    \hline
      \AxiomC{\sFmla{\True}{\Box \varphi}[\sigma]}
      \RightLabel{D$\Box$}
      \UnaryInfC{\sFmla{\True}{\Diamond\varphi}[\sigma]}
      \DisplayProof
    &
      \AxiomC{\sFmla{\False}{\Diamond \varphi}[\sigma]}
      \RightLabel{D$\Diamond$}
      \UnaryInfC{\sFmla{\False}{\Box\varphi}[\sigma]}
      \DisplayProof
    \\[1ex]
    \hline
      \AxiomC{\sFmla{\True}{\Box \varphi}[\sigma.n]}
      \RightLabel{B$\Box$}
      \UnaryInfC{\sFmla{\True}{\varphi}[\sigma]}
      \DisplayProof
    &
      \AxiomC{\sFmla{\False}{\Diamond \varphi}[\sigma.n]}
      \RightLabel{B$\Diamond$}
      \UnaryInfC{\sFmla{\False}{\varphi}[\sigma]}
      \DisplayProof
    \\[1ex]
    \hline
      \AxiomC{\sFmla{\True}{\Box \varphi}[\sigma]}
      \RightLabel{4$\Box$}
      \UnaryInfC{\sFmla{\True}{\Box\varphi}[\sigma.n]}
      \DisplayProof
    &
      \AxiomC{\sFmla{\False}{\Diamond \varphi}[\sigma]}
      \RightLabel{4$\Diamond$}
      \UnaryInfC{\sFmla{\False}{\Diamond\varphi}[\sigma.n]}
      \DisplayProof
    \\
    $\sigma.n$ वापरलेले आहे & $\sigma.n$ वापरलेले आहे
    \\[1ex]
    \hline
      \AxiomC{\sFmla{\True}{\Box \varphi}[\sigma.n]}
      \RightLabel{4r$\Box$}
      \UnaryInfC{\sFmla{\True}{\Box\varphi}[\sigma]}
      \DisplayProof
    &
      \AxiomC{\sFmla{\False}{\Diamond \varphi}[\sigma.n]}
      \RightLabel{4r$\Diamond$}
      \UnaryInfC{\sFmla{\False}{\Diamond\varphi}[\sigma]}
      \DisplayProof
    \\[1ex]
    \hline
    \end{tabular}
  \end{center}
  \caption{अधिक मोडल नियम.}
  \label{nml:tab:mru:tab:more-rules}
\end{table}
हे नियम जोडल्यावर \ref{nml:tab:mru:tab:logics-rules} मध्ये
दिलेल्या तर्कप्रणालींसाठी निर्दोष आणि पूर्ण
अशा पद्धती मिळतात.
\begin{table}
  \begin{center}
    \begin{tabular}{lll}
      \hline
      तर्कप्रणाली & $R$ हा \dots & नियम\\
      \hline
      $\Log{T} = \Log{KT}$ & परावर्ती &
      T$\Box$
      ,
      T$\Diamond$
      \\ \hline
      $\Log{D} = \Log{KD}$ & सर्वत्र उत्तरवर्ती &
      D$\Box$
      ,
      D$\Diamond$
      \\ \hline
      $\Log{K4}$ & संक्रमक &
      4$\Box$
      ,
      4$\Diamond$
      \\ \hline
      $\Log{B} = \Log{KTB}$ & परावर्ती, &
      T$\Box$
      ,
      T$\Diamond$\\
      & सममित &
      B$\Box$
      ,
      B$\Diamond$
      \\ \hline
      $\Log{S4} = \Log{KT4}$ & परावर्ती, &
      T$\Box$
      ,
      T$\Diamond$,\\
      & संक्रमक &
      4$\Box$
      ,
      4$\Diamond$
      \\ \hline
      $\Log{S5} = \Log{KT4B}$ & परावर्ती, &
      T$\Box$
      ,
      T$\Diamond$,\\
      & संक्रमक, &
      4$\Box$
      ,
      4$\Diamond$,\\
      &  युक्लिडी &
      4r$\Box$
      ,
      4r$\Diamond$
      \\ \hline
    \end{tabular}
  \end{center}
  \caption{विविध मोडल तर्कप्रणालींसाठीचे टॅब्लो नियम.}
  \label{nml:tab:mru:tab:logics-rules}
\end{table}
\begin{ex}
  $\Log{S5} \Proves \Ax{5}$, म्हणजे
  $\Diamond\varphi \lif \Box\Diamond\varphi$, हे दाखवणारा
  बंद टॅब्लो पुढीलप्रमाणे आहे.
  \begin{oltableau}
    [\pFmla{\False}{\Diamond\formula{A} \lif \Box\Diamond \formula{A}}{1},
      just = \TAss
      [\pFmla{\True}{\Diamond \formula{A}}{1}, just = {\TRule{\False}{\lif}[1]}
        [\pFmla{\False}{\Box\Diamond \formula{A}}{1},
          just = {\TRule{\False}{\lif}[1]}
          [\pFmla{\False}{\Diamond \formula{A}}{1.1},
            just = {\TRule{\False}{\Box}[3]}
            [\pFmla{\False}{\Diamond \formula{A}}{1},
              just = {4r$\Diamond$ 4}
              [\pFmla{\True}{\formula{A}}{1.2},
                just = {\TRule{\True}{\Diamond}[2]}
                [\pFmla{\False}{\formula{A}}{1.2},
                  just = {\TRule{\False}{\Diamond}[5]}, close]
              ]
            ]
          ]
        ]
      ]
    ]
  \end{oltableau}
\end{ex}
\begin{prob}
पुढील विधाने दाखवणारे बंद टॅब्लो द्या:
  \begin{enumerate}
  \item $\Log{KT5} \Proves \Ax{B}$;
  \item $\Log{KT5} \Proves \Ax{4}$;
  \item $\Log{KDB4} \Proves \Ax{T}$;
  \item $\Log{KB4} \Proves \Ax{5}$;
  \item $\Log{KB5} \Proves \Ax{4}$;
  \item $\Log{KT} \Proves \Ax{D}$.
  \end{enumerate}
\end{prob}
\section{अतिरिक्त नियमांची निर्दोषता}\label{nml:tab:msn:sec}
प्रतिमानांच्या एखाद्या वर्गातील एका प्रतिमानात
टॅब्लोची शाखा पूर्ततायोग्य असेल, तेव्हा
नियम लावल्यावर मिळणारी शाखाही त्या वर्गातील
प्रतिमानात पूर्ततायोग्य असेल, तर तो नियम त्या
वर्गासाठी निर्दोष आहे असे म्हणतो.
\begin{prop}
  \label{nml:tab:msn:prop:soundness-T}
  \Ax{T}$\Box$
   आणि
  \Ax{T}$\Diamond$
  परावर्ती प्रतिमानांसाठी
  निर्दोष आहेत.
\end{prop}
\begin{proof}
\begin{enumerate}
\item $\sFmla{\True}{\Box \psi}[\sigma] \in \Gamma$
  याला T$\Box$ लावून शाखा विस्तारित केली:
  शाखेवर
    $\sFmla{\True}{\psi}[\sigma]$ हे नवे चिन्हांकित सूत्र मिळते. $\mSat{M}{\Gamma}[f]$
    असे समजू; विशेषतः $\mSat{M}{\Box
      \psi}[f(\sigma)]$. $R$ परावर्ती असल्याने
    $Rf(\sigma)f(\sigma)$. म्हणून
    $\mSat{M}{\psi}[f(\sigma)]$; म्हणजे
    $\mModel{M}, f$ हे
    $\sFmla{\True}{\psi}[\sigma]$ पूर्ण करते.
\item $\sFmla{\False}{\Diamond \psi}[\sigma] \in \Gamma$
  याला T$\Diamond$ लावून शाखा विस्तारित केली:
  शाखेवर
    $\sFmla{\False}{\psi}[\sigma]$ हे नवे चिन्हांकित सूत्र मिळते. $\mSat{M}{\Gamma}[f]$
    असे समजू; विशेषतः $\mSat/{M}{\Diamond
      \psi}[f(\sigma)]$. $R$ परावर्ती असल्याने
    $Rf(\sigma)f(\sigma)$. म्हणून
    $\mSat/{M}{\psi}[f(\sigma)]$; म्हणजे
    $\mModel{M}, f$ हे
    $\sFmla{\False}{\psi}[\sigma]$ पूर्ण करते.
\end{enumerate}
\end{proof}
\begin{prop}
  \label{nml:tab:msn:prop:soundness-D}
  \Ax{D}$\Box$
   आणि
  \Ax{D}$\Diamond$
  सर्वत्र उत्तरवर्ती प्रतिमानांसाठी
  निर्दोष आहेत.
\end{prop}
\begin{proof}
\begin{enumerate}
\item $\sFmla{\True}{\Box \psi}[\sigma] \in \Gamma$
  याला D$\Box$ लावून शाखा विस्तारित केली:
  शाखेवर
    $\sFmla{\True}{\Diamond \psi}[\sigma]$ हे नवे
    चिन्हांकित सूत्र मिळते. $\mSat{M}{\Gamma}[f]$
    असे समजू; विशेषतः
    $\mSat{M}{\Box \psi}[f(\sigma)]$.
    $R$ सर्वत्र उत्तरवर्ती असल्याने
    $Rf(\sigma)w$ होईल असे $w
    \in W$ असते. मग $\mSat{M}{\psi}[w]$,
    आणि म्हणून $\mSat{M}{\Diamond \psi}[f(\sigma)]$.
    त्यामुळे $\mModel{M}, f$ हे
    $\sFmla{\True}{\Diamond\psi}[\sigma]$ पूर्ण करते.
\item $\sFmla{\False}{\Diamond \psi}[\sigma] \in \Gamma$
  याला D$\Diamond$ लावून शाखा विस्तारित केली:
  शाखेवर
    $\sFmla{\False}{\Box \psi}[\sigma]$ हे नवे
    चिन्हांकित सूत्र मिळते. $\mSat{M}{\Gamma}[f]$
    असे समजू; विशेषतः
    $\mSat/{M}{\Diamond \psi}[f(\sigma)]$.
    $R$ सर्वत्र उत्तरवर्ती असल्याने
    $Rf(\sigma)w$ होईल असे $w
    \in W$ असते. मग $\mSat/{M}{\psi}[w]$,
    आणि म्हणून $\mSat/{M}{\Box \psi}[f(\sigma)]$.
    त्यामुळे $\mModel{M}, f$ हे
    $\sFmla{\False}{\Box\psi}[\sigma]$ पूर्ण करते.
\end{enumerate}
\end{proof}
\begin{prop}
  \label{nml:tab:msn:prop:soundness-B}
  \Ax{B}$\Box$
   आणि
  \Ax{B}$\Diamond$
  सममित प्रतिमानांसाठी
  निर्दोष आहेत.
\end{prop}
\begin{proof}
\begin{enumerate}
\item $\sFmla{\True}{\Box \psi}[\sigma.n] \in \Gamma$
  याला B$\Box$ लावून शाखा विस्तारित केली:
  शाखेवर
    $\sFmla{\True}{\psi}[\sigma]$ हे नवे चिन्हांकित सूत्र मिळते. $\mSat{M}{\Gamma}[f]$
    असे समजू; विशेषतः $\mSat{M}{\Box
      \psi}[f(\sigma.n)]$. $f$ हे शाखेवरील
    पूर्वचिन्हांचे~$\mModel{M}$ मधील
    अर्थनिर्धारण असल्याने
    $Rf(\sigma)f(\sigma.n)$.
    $R$ सममित असल्याने
    $Rf(\sigma.n)f(\sigma)$.
    $\mSat{M}{\Box \psi}[f(\sigma.n)]$
    असल्यामुळे $\mSat{M}{\psi}[f(\sigma)]$.
    म्हणून $\mModel{M}, f$ हे
    $\sFmla{\True}{\psi}[\sigma]$ पूर्ण करते.
\item $\sFmla{\False}{\Diamond \psi}[\sigma.n] \in \Gamma$
  याला B$\Diamond$ लावून शाखा विस्तारित केली:
  शाखेवर
    $\sFmla{\False}{\psi}[\sigma]$ हे नवे चिन्हांकित सूत्र मिळते. $\mSat{M}{\Gamma}[f]$
    असे समजू; विशेषतः $\mSat/{M}{\Diamond
      \psi}[f(\sigma.n)]$. $f$ हे शाखेवरील
    पूर्वचिन्हांचे~$\mModel{M}$ मधील
    अर्थनिर्धारण असल्याने
    $Rf(\sigma)f(\sigma.n)$.
    $R$ सममित असल्याने
    $Rf(\sigma.n)f(\sigma)$.
    $\mSat/{M}{\Diamond \psi}[f(\sigma.n)]$
    असल्यामुळे $\mSat/{M}{\psi}[f(\sigma)]$.
    म्हणून $\mModel{M}, f$ हे
    $\sFmla{\False}{\psi}[\sigma]$ पूर्ण करते.
\end{enumerate}
\end{proof}
\begin{prop}
  \label{nml:tab:msn:prop:soundness-4}
  \Ax{4}$\Box$
   आणि
  \Ax{4}$\Diamond$
  संक्रमक प्रतिमानांसाठी
  निर्दोष आहेत.
\end{prop}
\begin{proof}
\begin{enumerate}
\item $\sFmla{\True}{\Box \psi}[\sigma] \in \Gamma$
  याला 4$\Box$ लावून शाखा विस्तारित केली:
  शाखेवर
    $\sFmla{\True}{\Box\psi}[\sigma.n]$ हे नवे
    चिन्हांकित सूत्र मिळते. $\mSat{M}{\Gamma}[f]$
    असे समजू; विशेषतः
    $\mSat{M}{\Box \psi}[f(\sigma)]$.
    $f$ हे शाखेवरील पूर्वचिन्हांचे
    $\mModel{M}$ मधील अर्थनिर्धारण आहे
    आणि $\sigma.n$ वापरलेले असले पाहिजे;
    म्हणून $Rf(\sigma)f(\sigma.n)$.
    आता $Rf(\sigma.n)w$ होईल असे कोणतेही
    जग~$w$ घ्या. $R$ संक्रमक असल्याने
    $Rf(\sigma)w$. तसेच
    $\mSat{M}{\Box \psi}[f(\sigma)]$
    असल्यामुळे $\mSat{M}{\psi}[w]$.
    म्हणून $\mSat{M}{\Box \psi}[f(\sigma.n)]$
    आणि $\mModel{M}, f$ हे
    $\sFmla{\True}{\Box \psi}[\sigma.n]$ पूर्ण करते.
\item $\sFmla{\False}{\Diamond \psi}[\sigma] \in \Gamma$
  याला 4$\Diamond$ लावून शाखा विस्तारित केली:
  शाखेवर
    $\sFmla{\False}{\Diamond\psi}[\sigma.n]$ हे नवे
    चिन्हांकित सूत्र मिळते. $\mSat{M}{\Gamma}[f]$
    असे समजू; विशेषतः
    $\mSat/{M}{\Diamond \psi}[f(\sigma)]$.
    $f$ हे शाखेवरील पूर्वचिन्हांचे
    $\mModel{M}$ मधील अर्थनिर्धारण आहे
    आणि $\sigma.n$ वापरलेले असले पाहिजे;
    म्हणून $Rf(\sigma)f(\sigma.n)$.
    आता $Rf(\sigma.n)w$ होईल असे कोणतेही
    जग~$w$ घ्या. $R$ संक्रमक असल्याने
    $Rf(\sigma)w$. तसेच
    $\mSat/{M}{\Diamond \psi}[f(\sigma)]$
    असल्यामुळे $\mSat/{M}{\psi}[w]$.
    म्हणून $\mSat/{M}{\Diamond \psi}[f(\sigma.n)]$
    आणि $\mModel{M}, f$ हे
    $\sFmla{\False}{\Diamond \psi}[\sigma.n]$ पूर्ण करते.
\end{enumerate}
\end{proof}
\begin{prop}
  \label{nml:tab:msn:prop:soundness-4r}
  \Ax{4r}$\Box$
   आणि
  \Ax{4r}$\Diamond$
  युक्लिडी प्रतिमानांसाठी
  निर्दोष आहेत.
\end{prop}
\begin{proof}
\begin{enumerate}
\item $\sFmla{\True}{\Box \psi}[\sigma.n] \in \Gamma$
  याला 4r$\Box$ लावून शाखा विस्तारित केली:
  शाखेवर
    $\sFmla{\True}{\Box\psi}[\sigma]$ हे नवे
    चिन्हांकित सूत्र मिळते. $\mSat{M}{\Gamma}[f]$
    असे समजू; विशेषतः $\mSat{M}{\Box
      \psi}[f(\sigma.n)]$.
    $f$ हे शाखेवरील पूर्वचिन्हांचे
    $\mModel{M}$ मधील अर्थनिर्धारण असल्याने
    $Rf(\sigma)f(\sigma.n)$.
    आता $Rf(\sigma)w$ होईल असे कोणतेही
    जग~$w$ घ्या. $R$ युक्लिडी असल्याने
    $Rf(\sigma.n)w$.
    $\mSat{M}{\Box \psi}[f(\sigma.n)]$
    असल्याने $\mSat{M}{\psi}[w]$.
    म्हणून $\mSat{M}{\Box \psi}[f(\sigma)]$;
    आणि $\mModel{M}, f$ हे
    $\sFmla{\True}{\Box \psi}[\sigma]$ पूर्ण करते.
\item $\sFmla{\False}{\Diamond \psi}[\sigma.n] \in \Gamma$
  याला 4r$\Diamond$ लावून शाखा विस्तारित केली:
  शाखेवर
    $\sFmla{\False}{\Diamond\psi}[\sigma]$ हे नवे
    चिन्हांकित सूत्र मिळते. $\mSat{M}{\Gamma}[f]$
    असे समजू; विशेषतः $\mSat/{M}{\Diamond
      \psi}[f(\sigma.n)]$.
    $f$ हे शाखेवरील पूर्वचिन्हांचे
    $\mModel{M}$ मधील अर्थनिर्धारण असल्याने
    $Rf(\sigma)f(\sigma.n)$.
    आता $Rf(\sigma)w$ होईल असे कोणतेही
    जग~$w$ घ्या. $R$ युक्लिडी असल्याने
    $Rf(\sigma.n)w$.
    $\mSat/{M}{\Diamond \psi}[f(\sigma.n)]$
    असल्याने $\mSat/{M}{\psi}[w]$.
    म्हणून $\mSat/{M}{\Diamond \psi}[f(\sigma)]$;
    आणि $\mModel{M}, f$ हे
    $\sFmla{\False}{\Diamond \psi}[\sigma]$ पूर्ण करते.
\end{enumerate}
\end{proof}
\begin{cor}
\label{nml:tab:msn:cor:soundness-logics}
\ref{nml:tab:mru:tab:logics-rules} मध्ये दिलेल्या
टॅब्लो पद्धती त्यांच्या संबंधित
प्रतिमानवर्गांसाठी निर्दोष आहेत.
\end{cor}
\section{\Log{S5} साठी सोपे टॅब्लो}\label{nml:tab:s5:sec}
\Log{S5} ही प्रणाली वैश्विक प्रतिमानांच्या
वर्गासाठी निर्दोष आणि पूर्ण आहे; अशा प्रतिमानांत
प्रत्येक जग प्रत्येक जगापासून प्राप्य असते.
वैश्विक प्रतिमानांत प्राप्यता संबंध स्वतंत्रपणे
तपासण्याची गरज नसते: ‘$\mSat{M}{\varphi}[w]$
होईल असे जग~$w$ आहे’ हे विधान तेव्हा आणि
तेव्हाच खरे असते जेव्हा असे~$w$ हे~$u$
पासून प्राप्य असते. म्हणून \Log{S5} मध्ये
प्रतिमानाची व्याख्या केवळ जगांचा संच आणि
सत्य-मूल्यांकन~$V$ यांद्वारे करता येते. यावरून
टॅब्लोचे नियमही सोपे करता येतील असे
सुचते. सर्वसाधारण प्रकरणात पूर्वचिन्हे म्हणून
धन पूर्णांकांच्या क्रमिका वापरतो, म्हणजे कोणत्या
पूर्वचिन्हाने नाव दिलेले जग दुसऱ्या जगापासून
प्राप्य आहे याची नोंद ठेवता येते:
$\sigma.n$ हे~$\sigma$ पासून प्राप्य
जगाचे नाव देते. पण \Log{S5} मध्ये कोणतेही
जग कोणत्याही जगापासून प्राप्य असते, त्यामुळे
अशी नोंद ठेवण्याची गरज नाही. त्याऐवजी
धन पूर्णांक पूर्वचिन्हे म्हणून वापरता येतात.
सोपे केलेले नियम \ref{nml:tab:s5:tab:rules-S5}
मध्ये दिले आहेत.
\begin{table}
  \begin{center}
    \def\arraystretch{3}\def\fCenter{}
    \begin{tabular}{|c|c|}
    \hline
      \AxiomC{\sFmla{\True}{\Box \varphi}[n]}
      \RightLabel{\TRule{\True}{\Box}}
      \UnaryInfC{\sFmla{\True}{\varphi}[m]}
      \DisplayProof
      &
      \AxiomC{\sFmla{\False}{\Box \varphi}[n]}
      \RightLabel{\TRule{\False}{\Box}}
      \UnaryInfC{\sFmla{\False}{\varphi}[m]}
      \DisplayProof\\
      $m$ वापरलेले आहे & $m$ नवे आहे
      \\[1ex]
      \hline
      \AxiomC{\sFmla{\True}{\Diamond \varphi}[n]}
      \RightLabel{\TRule{\True}{\Diamond}}
      \UnaryInfC{\sFmla{\True}{\varphi}[m]}
      \DisplayProof
      &
      \AxiomC{\sFmla{\False}{\Diamond \varphi}[n]}
      \RightLabel{\TRule{\False}{\Diamond}}
      \UnaryInfC{\sFmla{\False}{\varphi}[m]}
      \DisplayProof\\
      $m$ नवे आहे & $m$ वापरलेले आहे
      \\[1ex]
      \hline
    \end{tabular}
  \end{center}
  \caption{\Log{S5} साठी सोपे केलेले नियम.}
  \label{nml:tab:s5:tab:rules-S5}
\end{table}
\begin{ex}
  $\Log{S5} \Proves \Ax{5}$, म्हणजे
  $\Diamond\varphi \lif \Box\Diamond\varphi$, हे दाखवणारा
  सोपा केलेला बंद टॅब्लो पुढीलप्रमाणे आहे.
  \begin{oltableau}
    [\pFmla{\False}{\Diamond\formula{A} \lif \Box\Diamond \formula{A}}{1},
      just = \TAss
      [\pFmla{\True}{\Diamond \formula{A}}{1}, just = {\TRule{\False}{\lif}[1]}
        [\pFmla{\False}{\Box\Diamond \formula{A}}{1},
          just = {\TRule{\False}{\lif}[1]}
          [\pFmla{\False}{\Diamond \formula{A}}{2},
            just = {\TRule{\False}{\Box}[3]}
            [\pFmla{\True}{\formula{A}}{3},
              just = {\TRule{\True}{\Diamond}[2]}
                [\pFmla{\False}{\formula{A}}{3},
                  just = {\TRule{\False}{\Diamond}[4]}, close]
            ]
          ]
        ]
      ]
    ]
  \end{oltableau}
\end{ex}
\section{\Log{K} साठी संपूर्णता}\label{nml:tab:cpl:sec}
\begin{explain}
  टॅब्लोची पद्धत संपूर्ण आहे हे
  दाखवण्यासाठी पुढील विधान सिद्ध करावे लागते:
  $\Gamma \Proves \varphi$ दाखवणारा बंद टॅब्लो
  नसेल, तर $\Gamma \Entails/ \varphi$; म्हणजे
  प्रतिदृष्टांत प्रतिमान अस्तित्वात असते.
  पण ‘बंद टॅब्लो नाही’ याचा अर्थ
  तो बांधण्याचा प्रत्येक प्रयत्न बंद होण्यात
  अयशस्वी होतो. युक्ती अशी की प्रत्येक प्रयत्न
  अयशस्वी होत असेल, तर एक विशिष्ट
  \emph{पद्धतशीर आणि सर्वसमावेशक} प्रयत्नही
  अयशस्वी होतो. बंद टॅब्लो अस्तित्वात
  असेल, तर या पद्धतशीर प्रयत्नाने तो बंद
  झाला असता. अशा एकाच टॅब्लोमधील खुल्या
  शाखांत प्रतिदृष्टांत प्रतिमानाची व्याख्या
  करण्यासाठी लागणारी सर्व माहिती असते.
  त्यातील एका खुल्या शाखेवरून मिळणाऱ्या
  प्रतिदृष्टांत प्रतिमानाची जगे म्हणजे त्या
  शाखेवर वापरलेली सर्व पूर्वचिन्हे; आणि
  विधानीय चल~$p$ हे~$\sigma$ येथे
  तेव्हा आणि तेव्हाच सत्य असते जेव्हा
  $\sFmla{\True}{p}[\sigma]$ हे त्या
  शाखेवर आलेले असते.
\end{explain}
\begin{defn}
  टॅब्लो मधील शाखेला संपूर्ण
  म्हणतात, जर नियम लावता येईल असे
  $\sFmla{S}{\varphi}[\sigma]$ हे पूर्वचिन्हयुक्त
  सूत्र शाखेवर असताना पुढील गोष्टीही
  तिच्यावर असतील:
  \begin{enumerate}
    \item विधानीय, एकाच शाखेवर सूत्रे जोडणाऱ्या
      नियमाच्या बाबतीत त्याच्या निष्कर्षांतील
      पूर्वचिन्हयुक्त सूत्रे;
    \item विधानीय, शाखा फोडणाऱ्या नियमाच्या
      बाबतीत संबंधित निष्कर्षरूप
      सूत्रांपैकी एक;
    \item नवे पूर्वचिन्ह लागणाऱ्या मोडल नियमाच्या
      बाबतीत किमान एक शक्य निष्कर्ष;
    \item वापरलेले पूर्वचिन्ह लागणाऱ्या मोडल
      नियमाच्या बाबतीत शाखेवर असलेल्या
      प्रत्येक योग्य पूर्वचिन्हासाठी संबंधित
      निष्कर्ष.
  \end{enumerate}
\end{defn}
\begin{explain}
उदाहरणार्थ, संपूर्ण शाखेवर
$\sFmla{\True}{\psi \land \chi}[\sigma]$
असेल, तर $\sFmla{\True}{\psi}[\sigma]$
आणि $\sFmla{\True}{\chi}[\sigma]$
ही दोन्ही असतात. तिच्यावर
$\sFmla{\True}{\psi \lor \chi}[\sigma]$
असेल, तर $\sFmla{\True}{\psi}[\sigma]$
आणि $\sFmla{\True}{\chi}[\sigma]$
यांपैकी किमान एक असते.
शाखेवर $\sFmla{\False}{\Box \psi}[\sigma]$
असेल, तर काही~$n$ साठी अनुक्रमे
$\sFmla{\False}{\psi}[\sigma.n]$
हेही असते. आणि शाखेवर
$\sFmla{\True}{\Box \psi}[\sigma]$
असेल, तर शाखेवर $\sigma.n$ वापरलेले
असेल अशा प्रत्येक~$n$ साठी अनुक्रमे
$\sFmla{\True}{\psi}[\sigma.n]$
हेही असते.
\end{explain}
\begin{prop}\label{nml:tab:cpl:prop:complete-tableau}
  $\Gamma$ हा $1$ या एकाच पूर्वचिन्हाने चिन्हांकित सूत्रांचा
  सांत संच असू द्या. त्यासाठी असा टॅब्लो असतो की
  त्याची प्रत्येक शाखा बंद किंवा संपूर्ण असते.
\end{prop}
\begin{proof}
  $\Gamma$ साठीच्या टॅब्लो मधील
  एखादी खुली शाखा घ्या. त्या शाखेवर नियम
  लावता येतील अशी सांत पूर्वचिन्हयुक्त
  सूत्रे आहेत. ठरलेल्या एखाद्या क्रमाने
  (उदाहरणार्थ, वरून खाली), ज्या
  पूर्वचिन्हयुक्त सूत्रे साठी
  (1)--(4) या अटी आधीच पूर्ण होत नाहीत,
  त्यांपैकी प्रत्येकाला लागू होणारे नियम
  लावून शाखा विस्तारित करा. काही वेळा शाखा
  फुटतील; मग उरलेल्या पूर्वचिन्हयुक्त
  सूत्रे साठी प्रत्येक नव्या शाखेच्या
  टोकावर नियम लावा. पूर्वचिन्हयुक्त
  सूत्रे आणि शाखेवर वापरलेली
  पूर्वचिन्हे यांची संख्या सांत असल्याने
  या प्रक्रियेतून मूळ शाखेपासून पुढे जाणाऱ्या
  (कदाचित अनेक) शाखा मिळतात. त्यांपैकी
  प्रत्येकावर ही प्रक्रिया पुन्हा लावा.
  बंद झालेल्या शाखांवर पुढील नियम लावू नयेत.
  ही प्रक्रिया सांत पायऱ्यांत थांबते: नव्या पूर्वचिन्हावर येणारी
  सूत्रे त्याच्या पालकावरील मोडल सूत्रांची उपसूत्रे असतात.
  त्यामुळे प्रत्येक पूर्वचिन्ह-विस्ताराने उपलब्ध मोडल खोली कमी
  होते; प्रारंभीच्या सांत सूत्रसंचाची कमाल मोडल खोली ही शाखेवरील
  पूर्वचिन्हांच्या उंचीची मर्यादा ठरते. प्रत्येक पूर्वचिन्हावर फक्त
  प्रारंभीच्या सूत्रांची सांत चिन्हांकित उपसूत्रे येऊ शकतात.
  प्रत्येक नवेपणा मागणाऱ्या चिन्हांकित मोडल सूत्रासाठी एकच
  साक्षीदार पूर्वचिन्ह पुरते, म्हणून प्रत्येक पूर्वचिन्हाची संततीही
  सांत असते. मर्यादित उंची आणि प्रत्येक ठिकाणी सांत संतती
  यांमुळे शाखेवरील सर्व पूर्वचिन्हे सांत संख्येची असतात.
  उर्वरित प्रत्येक नियमाचे आवश्यक अनुप्रयोगही सांत असतात.
  रचनेमुळे उरलेली प्रत्येक खुली शाखा संपूर्ण असते.
\end{proof}
\begin{thm}[संपूर्णता]
  \label{nml:tab:cpl:thm:tableau-completeness}
  $\Gamma$ हा $1$ या एकाच पूर्वचिन्हाने चिन्हांकित सूत्रांचा संच असू द्या.
  $\Gamma$ साठी बंद टॅब्लो नसेल,
  तर $\Gamma$ पूर्ततायोग्य आहे.
\end{thm}
\begin{proof}
$\Gamma$ रिक्त असेल, तर एक जग असलेले कोणतेही प्रतिमान
ते पूर्ण करते. प्रथम $\Gamma$ रिक्त नसलेला सांत संच आहे
असे समजू. वरील विधानानुसार, $\Gamma$ साठी असा
टॅब्लो असतो की त्याची प्रत्येक शाखा
बंद किंवा संपूर्ण असते. बंद टॅब्लो नसल्याने,
त्या टॅब्लोमध्ये किमान एक शाखा
खुली आणि संपूर्ण असते. त्या शाखेवरील
पूर्वचिन्हयुक्त सूत्रांचा संच
$\Delta$ आणि त्यात आलेल्या पूर्वचिन्हांचा
संच $P(\Delta)$ असू द्या.
आता $\mModel{M(\Delta)} = \tuple{P(\Delta), R, V}$
हे प्रतिमान ठरवू. त्यातील जगे म्हणजे~$\Delta$
मधील पूर्वचिन्हे, आणि प्राप्यता संबंध असा आहे:
\[
R\sigma\sigma' \quad \text{तेव्हा आणि तेव्हाच} \quad
\sigma'=\sigma.n \quad \text{काही~$n$ साठी}
\]
आणि
\[V(p) = \Setabs{\sigma}{\sFmla{\True}{p}[\sigma] \in \Delta}.
\]
आता~$\varphi$ च्या रचनेवरून आगमनाने पुढील
दाखवू: $\sFmla{\True}{\varphi}[\sigma] \in
\Delta$ असेल, तर
$\mSat{M(\Delta)}{\varphi}[\sigma]$; आणि
$\sFmla{\False}{\varphi}[\sigma] \in \Delta$
असेल, तर $\mSat/{M(\Delta)}{\varphi}[\sigma]$.

\begin{enumerate}
  \item \indcase{\varphi}{p}{$\sFmla{\True}{\indfrm}[\sigma] \in \Delta$
    असेल, तर~$V$ च्या व्याख्येनुसार
    $\sigma \in V(p)$; म्हणून
    $\mSat{M(\Delta)}{\indfrm}[\sigma]$.

    $\sFmla{\False}{\indfrm}[\sigma] \in \Delta$ असेल, तर
    $\sFmla{\True}{\indfrm}[\sigma] \notin \Delta$,
    कारण अन्यथा शाखा बंद झाली असती.
    म्हणून $\sigma \notin V(p)$, आणि
    $\mSat/{M(\Delta)}{\indfrm}[\sigma]$.}
  \item \indcase{\varphi}{\lnot \psi}{जर
      $\sFmla{\True}{\indfrm}[\sigma] \in \Delta$,
      तर शाखा संपूर्ण असल्याने
      $\sFmla{\False}{\psi}[\sigma] \in \Delta$.
      आगमनाच्या गृहीतकानुसार
      $\mSat/{M(\Delta)}{\psi}[\sigma]$;
      म्हणून $\mSat{M(\Delta)}{\indfrm}[\sigma]$.

      $\sFmla{\False}{\indfrm}[\sigma] \in \Delta$
      असेल, तर शाखा संपूर्ण असल्याने
      $\sFmla{\True}{\psi}[\sigma] \in \Delta$.
      आगमनाच्या गृहीतकानुसार
      $\mSat{M(\Delta)}{\psi}[\sigma]$;
      म्हणून $\mSat/{M(\Delta)}{\indfrm}[\sigma]$.}
  \item \indcase{\varphi}{\psi \land \chi}{जर
      $\sFmla{\True}{\indfrm}[\sigma] \in \Delta$,
      तर शाखा संपूर्ण असल्याने
      $\sFmla{\True}{\psi}[\sigma] \in \Delta$
      आणि $\sFmla{\True}{\chi}[\sigma] \in \Delta$
      ही दोन्ही. आगमनाच्या गृहीतकानुसार
      $\mSat{M(\Delta)}{\psi}[\sigma]$ आणि
      $\mSat{M(\Delta)}{\chi}[\sigma]$.
      म्हणून $\mSat{M(\Delta)}{\indfrm}[\sigma]$.

      $\sFmla{\False}{\indfrm}[\sigma] \in \Delta$
      असेल, तर शाखा संपूर्ण असल्याने
      $\sFmla{\False}{\psi}[\sigma] \in \Delta$
      किंवा $\sFmla{\False}{\chi}[\sigma] \in \Delta$.
      आगमनाच्या गृहीतकानुसार
      $\mSat/{M(\Delta)}{\psi}[\sigma]$ किंवा
      $\mSat/{M(\Delta)}{\chi}[\sigma]$.
      म्हणून $\mSat/{M(\Delta)}{\indfrm}[\sigma]$.}
  \item \indcase{\varphi}{\psi \lor \chi}{जर
      $\sFmla{\True}{\indfrm}[\sigma] \in \Delta$,
      तर शाखा संपूर्ण असल्याने
      $\sFmla{\True}{\psi}[\sigma] \in \Delta$
      किंवा $\sFmla{\True}{\chi}[\sigma] \in \Delta$.
      आगमनाच्या गृहीतकानुसार
      $\mSat{M(\Delta)}{\psi}[\sigma]$ किंवा
      $\mSat{M(\Delta)}{\chi}[\sigma]$.
      म्हणून $\mSat{M(\Delta)}{\indfrm}[\sigma]$.

      $\sFmla{\False}{\indfrm}[\sigma] \in \Delta$
      असेल, तर शाखा संपूर्ण असल्याने
      $\sFmla{\False}{\psi}[\sigma] \in \Delta$
      आणि $\sFmla{\False}{\chi}[\sigma] \in \Delta$
      ही दोन्ही. आगमनाच्या गृहीतकानुसार
      $\mSat/{M(\Delta)}{\psi}[\sigma]$ आणि
      $\mSat/{M(\Delta)}{\chi}[\sigma]$.
      म्हणून $\mSat/{M(\Delta)}{\indfrm}[\sigma]$.}
  \item \indcase{\varphi}{\psi \lif \chi}{जर
      $\sFmla{\True}{\indfrm}[\sigma] \in \Delta$,
      तर शाखा संपूर्ण असल्याने
      $\sFmla{\False}{\psi}[\sigma] \in \Delta$
      किंवा $\sFmla{\True}{\chi}[\sigma] \in \Delta$.
      आगमनाच्या गृहीतकानुसार
      $\mSat/{M(\Delta)}{\psi}[\sigma]$ किंवा
      $\mSat{M(\Delta)}{\chi}[\sigma]$.
      म्हणून $\mSat{M(\Delta)}{\indfrm}[\sigma]$.

      $\sFmla{\False}{\indfrm}[\sigma] \in \Delta$
      असेल, तर शाखा संपूर्ण असल्याने
      $\sFmla{\True}{\psi}[\sigma] \in \Delta$
      आणि $\sFmla{\False}{\chi}[\sigma] \in \Delta$
      ही दोन्ही. आगमनाच्या गृहीतकानुसार
      $\mSat{M(\Delta)}{\psi}[\sigma]$ आणि
      $\mSat/{M(\Delta)}{\chi}[\sigma]$.
      म्हणून $\mSat/{M(\Delta)}{\indfrm}[\sigma]$.}
  \item \indcase{\varphi}{\Box \psi}{जर
      $\sFmla{\True}{\indfrm}[\sigma] \in \Delta$,
      तर शाखा संपूर्ण असल्याने तिच्यावर
      वापरलेल्या प्रत्येक $\sigma.n$ साठी
      $\sFmla{\True}{\psi}[\sigma.n] \in \Delta$.
      म्हणजे $R\sigma\sigma'$ होईल अशा
      प्रत्येक $\sigma' \in P(\Delta)$ साठी
      हे खरे आहे. आगमनाच्या गृहीतकानुसार,
      $R\sigma\sigma'$ असलेल्या प्रत्येक
      $\sigma'$ साठी
      $\mSat{M(\Delta)}{\psi}[\sigma']$.
      म्हणून $\mSat{M(\Delta)}{\indfrm}[\sigma]$.

      $\sFmla{\False}{\indfrm}[\sigma] \in \Delta$
      असेल, तर शाखा संपूर्ण असल्याने
      काही $\sigma.n$ साठी
      $\sFmla{\False}{\psi}[\sigma.n] \in \Delta$.
      आगमनाच्या गृहीतकानुसार
      $\mSat/{M(\Delta)}{\psi}[\sigma.n]$.
      $R\sigma(\sigma.n)$ असल्याने
      $\mSat/{M(\Delta)}{\psi}[\sigma']$
      होईल असे~$\sigma'$ असते. म्हणून
      $\mSat/{M(\Delta)}{\indfrm}[\sigma]$.}
  \item \indcase{\varphi}{\Diamond \psi}{जर
      $\sFmla{\True}{\indfrm}[\sigma] \in \Delta$,
      तर शाखा संपूर्ण असल्याने
      काही $\sigma.n$ साठी
      $\sFmla{\True}{\psi}[\sigma.n] \in \Delta$.
      आगमनाच्या गृहीतकानुसार
      $\mSat{M(\Delta)}{\psi}[\sigma.n]$.
      $R\sigma(\sigma.n)$ असल्याने
      $\mSat{M(\Delta)}{\psi}[\sigma']$
      होईल असे~$\sigma'$ असते. म्हणून
      $\mSat{M(\Delta)}{\indfrm}[\sigma]$.

      $\sFmla{\False}{\indfrm}[\sigma] \in \Delta$
      असेल, तर शाखा संपूर्ण असल्याने तिच्यावर
      वापरलेल्या प्रत्येक $\sigma.n$ साठी
      $\sFmla{\False}{\psi}[\sigma.n] \in \Delta$.
      म्हणजे $R\sigma\sigma'$ होईल अशा
      प्रत्येक $\sigma' \in P(\Delta)$ साठी
      हे खरे आहे. आगमनाच्या गृहीतकानुसार,
      $R\sigma\sigma'$ असलेल्या प्रत्येक
      $\sigma'$ साठी
      $\mSat/{M(\Delta)}{\psi}[\sigma']$.
      म्हणून $\mSat/{M(\Delta)}{\indfrm}[\sigma]$.}
\end{enumerate}
$\Gamma \subseteq \Delta$ असल्याने, प्रत्येक पूर्वचिन्हाला तेच जग
नेमणारे अर्थनिर्धारण घेऊन $\mSat{M(\Delta)}{\Gamma}[\mathrm{id}]$.
याने सांत प्रकरण सिद्ध झाले.
असांत $\Gamma$ साठी त्याच्या प्रत्येक सांत उपसंचाला बंद टॅब्लो
नसतो: तसा टॅब्लो असता, तर त्याची सांत गृहीतके $\Gamma$
मध्येही असल्याने तो $\Gamma$ साठी बंद टॅब्लो ठरला असता.
म्हणून सांत प्रकरणावरून प्रत्येक सांत उपसंच पूर्ततायोग्य आहे.
चिन्हांकित गृहीतकांचा साध्या मोडल सूत्रांचा संच असा ठरवू:
\[\Upsilon = \Setabs{\psi}{\sFmla{\True}{\psi}[1]\in\Gamma}
\cup \Setabs{\lnot\psi}{\sFmla{\False}{\psi}[1]\in\Gamma}.
\]
याच्या प्रत्येक सांत उपसंचाला मूळ $\Gamma$ चा सांत उपसंच
पूर्ण करणाऱ्या प्रतिमानात एकाच जगावर सत्य करता येते.
त्यामुळे $\Upsilon$ हा \Log{K}-सुसंगत संच आहे: त्यावरून
असत्य सूत्र निष्पन्न होत असते, तर त्या सांत सिद्धतेतील सांत
गृहीतकांच्या पूर्ततायोग्यतेला \Log{K} ची निर्दोषता विरोधी ठरली
असती. \ref{nml:com:lin:thm:lindenbaum} नुसार
$\Upsilon$ ला संपूर्ण \Log{K}-सुसंगत संच $\Lambda$ पर्यंत
वाढवता येते. \ref{nml:com:tru:prop:truthlemma} नुसार
कॅनॉनिकल प्रतिमानात $\Lambda$ या जगावर $\Upsilon$ मधील
प्रत्येक सूत्र सत्य असते. $f(1)=\Lambda$ हे अर्थनिर्धारण मूळ
चिन्हांकित संच $\Gamma$ पूर्ण करते. याने असांत प्रकरणही सिद्ध झाले.
\end{proof}

\begin{prob}
\ref{nml:tab:cpl:thm:tableau-completeness} चा
पुरावा पूर्ण करा.
\end{prob}

\begin{cor}
\label{nml:tab:cpl:cor:entailment-completeness}
$\Gamma \Entails \varphi$ असेल, तर
$\Gamma \Proves \varphi$.
\end{cor}

\begin{cor}
\label{nml:tab:cpl:cor:weak-completeness}
$\varphi$ सर्व प्रतिमानांमध्ये सत्य असेल,
तर $\Proves \varphi$.
\end{cor}













\section{टॅब्लो मधून प्रतिदृष्टांत प्रतिमाने}\label{nml:tab:cou:sec}

\begin{explain}
  संपूर्णता प्रमेयाचा पुरावा
  $\Entails \varphi$ असेल तर $\Proves \varphi$
  हे दाखवतोच; शिवाय $\Entails/ \varphi$
  असेल, तर~$\varphi$ साठी प्रतिदृष्टांत
  प्रतिमान बांधण्याची पद्धतही देतो.
  $\Log{K}$ साठी ही पद्धत
  \emph{निर्णय प्रक्रिया} आहे. कारण
  $\Entails/ \varphi$ असे समजू. मग
  \ref{nml:tab:cpl:prop:complete-tableau} चा
  पुरावा संपूर्ण टॅब्लो बांधण्याची
  पद्धत देतो. ती नेहमी थांबतेही.
  $\Log{K}$ चे विधानीय नियम कमी गुंतागुंतीची
  पूर्वचिन्हयुक्त सूत्रे जोडतात;
  म्हणजे प्रत्येक
  $\sFmla{S}{\varphi}[\sigma]$ या चिन्हांकित
  सूत्रासाठी प्रत्येक विधानीय नियम
  शाखेवर एकदाच लावावा लागतो.
  नवी पूर्वचिन्हे फक्त
  $\TRule{\False}{\Box}$
    आणि $\TRule{\True}{\Diamond}$
  या नियमांनी
  तयार होतात. शाखेवरील प्रत्येक पूर्वचिन्हयुक्त
  चिन्हांकित मोडल सूत्रासाठी संबंधित नियम
  एकदाच लावावा लागतो; त्या अनुप्रयोगात
  एकच नवे पूर्वचिन्ह मिळते.
  $\TRule{\True}{\Box}$
    आणि $\TRule{\False}{\Diamond}$
  हे नियम
  अनेकदा लावावे लागू शकतात; पण प्रत्येक
  पूर्वचिन्हयुक्त चिन्हांकित मोडल सूत्रासाठी
  त्याच्या पूर्वचिन्हाचा शाखेवर वापरलेला प्रत्येक
  तात्काळ विस्तार घेऊन एकदाच लावणे पुरते.
  \ref{nml:tab:cpl:prop:complete-tableau} च्या पुराव्यात
  दाखवल्याप्रमाणे, सांत प्रारंभिक सूत्रसंचाची मोडल
  खोली पूर्वचिन्हांची उंची मर्यादित करते; प्रत्येक
  पूर्वचिन्हाची संततीही सांत असते. म्हणून नवी
  पूर्वचिन्हे सांत संख्येनेच निर्माण होतात. म्हणून
  सांत पायऱ्यांनंतर शाखा बंद होते किंवा
  संपूर्ण होते.

  खुली संपूर्ण शाखा असलेला टॅब्लो
  बांधल्यावर \ref{nml:tab:cpl:thm:tableau-completeness}
  चा पुरावा मूळ पूर्वचिन्हयुक्त सूत्रे
  चा संच पूर्ण करणारे स्पष्ट प्रतिमान देतो.
  म्हणून $\Gamma \Entails \varphi$ असेल, तर
  बंद टॅब्लो अस्तित्वात असतो आणि
  $\Gamma \Proves \varphi$. याशिवाय, $\Gamma$ सांत
  असेल, तर योग्य पद्धतीने बंद टॅब्लो शोधताना
  त्याऐवजी ‘संपूर्ण’ टॅब्लो मिळाला,
  तर $\Gamma \Entails/ \varphi$ हे कळते
  आणि प्रतिदृष्टांत प्रतिमान प्रत्यक्ष
  बांधताही येते.
\end{explain}


\begin{ex}
  $\Proves/ \Box(p \lor q) \lif (\Box p \lor \Box
  q)$ हे आपल्याला माहीत आहे. टॅब्लो बांधण्याची
  सुरुवात पुढीलप्रमाणे होते:
  \begin{oltableau}
    [\pFmla{\False}{\Box(p \lor q) \lif (\Box p \lor \Box q)}{1},
      just = \TAss, checked
      [\pFmla{\True}{\Box(p \lor q)}{1},
        just = {\TRule{\False}{\lif}[1]},
        [\pFmla{\False}{\Box p \lor \Box q}{1},
          just = {\TRule{\False}{\lif}[1]}, checked
          [\pFmla{\False}{\Box p}{1},
            just = {\TRule{\False}{\lor}[3]}, checked
            [\pFmla{\False}{\Box q}{1},
              just = {\TRule{\False}{\lor}[3]}, checked
              [\pFmla{\False}{p}{1.1},
                just = {\TRule{\False}{\Box}[4]}, checked
                [\pFmla{\False}{q}{1.2},
                  just = {\TRule{\False}{\Box}[5]}, checked
                ]
              ]
            ]
          ]
        ]
      ]
    ]
  \end{oltableau}
  अर्थात हा टॅब्लो अजून पूर्ण झालेला नाही.
  पुढे खूण नसलेल्या एकमेव ओळीचा विचार करू:
  ओळ~$2$ वरील पूर्वचिन्हयुक्त सूत्र
  $\sFmla{\True}{\Box(p \lor q)}[1]$.
  पद्धतशीर टॅब्लो बांधताना शाखेवर वापरलेल्या
  $1.n$ या आकाराच्या प्रत्येक पूर्वचिन्हासाठी
  $\TRule{\True}{\Box}$ लावायचा असतो;
  म्हणजे $1.1$ आणि $1.2$ या दोन्हींसाठी:
  \begin{oltableau}
    [\pFmla{\False}{\Box(p \lor q) \lif (\Box p \lor \Box q)}{1},
      just = \TAss, checked
      [\pFmla{\True}{\Box(p \lor q)}{1},
        just = {\TRule{\False}{\lif}[1]},
        [\pFmla{\False}{\Box p \lor \Box q}{1},
          just = {\TRule{\False}{\lif}[1]}, checked
          [\pFmla{\False}{\Box p}{1},
            just = {\TRule{\False}{\lor}[3]}, checked
            [\pFmla{\False}{\Box q}{1},
              just = {\TRule{\False}{\lor}[3]}, checked
              [\pFmla{\False}{p}{1.1},
                just = {\TRule{\False}{\Box}[4]}, checked
                [\pFmla{\False}{q}{1.2},
                  just = {\TRule{\False}{\Box}[5]}, checked
                  [\pFmla{\True}{p \lor q}{1.1},
                    just = {\TRule{\True}{\Box}[2]}
                    [\pFmla{\True}{p \lor q}{1.2},
                      just = {\TRule{\True}{\Box}[2]}
                    ]
                  ]
                ]
              ]
            ]
          ]
        ]
      ]
    ]
  \end{oltableau}
  आता ओळी 2, 8 आणि 9 यांवर खूण नाही.
  पण नवे पूर्वचिन्ह जोडलेले नाही. म्हणून
  ओळी~8 आणि~9 वर
  $\TRule{\True}{\lor}$ लावू; तयार होणाऱ्या
  सर्व शाखांवर, त्या बंद होत नसतील तोवर,
  हाच नियम लागू करू:
  \begin{oltableau}
    [\pFmla{\False}{\Box(p \lor q) \lif (\Box p \lor \Box q)}{1},
      just = \TAss, checked
      [\pFmla{\True}{\Box(p \lor q)}{1},
        just = {\TRule{\False}{\lif}[1]}, checked
        [\pFmla{\False}{\Box p \lor \Box q}{1},
          just = {\TRule{\False}{\lif}[1]}, checked
          [\pFmla{\False}{\Box p}{1},
            just = {\TRule{\False}{\lor}[3]}, checked
            [\pFmla{\False}{\Box q}{1},
              just = {\TRule{\False}{\lor}[3]}, checked
              [\pFmla{\False}{p}{1.1},
                just = {\TRule{\False}{\Box}[4]}, checked
                [\pFmla{\False}{q}{1.2},
                  just = {\TRule{\False}{\Box}[5]}, checked
                  [\pFmla{\True}{p \lor q}{1.1},
                    just = {\TRule{\True}{\Box}[2]}, checked
                    [\pFmla{\True}{p \lor q}{1.2},
                      just = {\TRule{\True}{\Box}[2]}, checked
                      [\pFmla{\True}{p}{1.1},
                        just = {\TRule{\True}{\lor}[8]}, checked, close
                      ]
                      [\pFmla{\True}{q}{1.1},
                        just = {\TRule{\True}{\lor}[8]}, checked
                        [\pFmla{\True}{p}{1.2},
                          just = {\TRule{\True}{\lor}[9]}, checked]
                        [\pFmla{\True}{q}{1.2},
                          just = {\TRule{\True}{\lor}[9]}, checked, close]
                      ]
                    ]
                  ]
                ]
              ]
            ]
          ]
        ]
      ]
    ]
  \end{oltableau}
  आता एकच खुली शाखा उरते; ती संपूर्ण आहे.
  तिच्यावरून प्रतिमान असे ठरवू: जगे
  $W = \{1, 1.1, 1.2\}$ (खुल्या शाखेवर
  येणारी हीच पूर्वचिन्हे), प्राप्यता संबंध
  $R = \{\tuple{1, 1.1}, \tuple{1, 1.2}\}$,
  आणि सत्य-मूल्यांकन $V(p) = \{1.2\}$
  (कारण ओळ~11 वर
  $\sFmla{\True}{p}[1.2]$ आहे) व
  $V(q) = \{1.1\}$ (कारण ओळ~10 वर
  $\sFmla{\True}{q}[1.1]$ आहे).
  \ref{nml:tab:cou:fig:counter-Box} मध्ये हे प्रतिमान
  दाखवले आहे. ते $\Box(p \lor q) \lif
  (\Box p \lor \Box q)$ चे प्रतिदृष्टांत
  प्रतिमान आहे हे तपासता येईल.
  \begin{figure}
  \begin{center}
    \begin{tikzpicture}[modal]
      \node[world] (w1) [label={[align=right]right:\mFalse{p}\\ \mFalse{q}}]{$1$};
      \node[world] (w2) [label={[align=right]right:\mFalse{p}\\ \mTrue{q}},
        above left=of w1]{$1.1$};
      \node[world] (w3) [label={[align=right]right:\mTrue{p}\\ \mFalse{q}},
        above right=of w1] {$1.2$};
      \draw[->] (w1) to (w2);
      \draw[->] (w1) to (w3);
    \end{tikzpicture}
  \end{center}
  \caption{$\Box(p \lor q) \lif (\Box p \lor \Box
    q)$ चे प्रतिदृष्टांत प्रतिमान.}
  \label{nml:tab:cou:fig:counter-Box}
\end{figure}
\end{ex}




\chapter{मोडल क्रमवर्ती कलन}\label{nml:seq:chap}










\section{प्रस्तावना}\label{nml:seq:int:sec}

विधानीय तर्कशास्त्राच्या क्रमवर्ती कलनात
~$\Box$ आणि~
$\Diamond$ यांच्याशी संबंधित
अतिरिक्त नियम घालून त्याचा विस्तार करता येतो.
उदाहरणार्थ,~$\Log{K}$ साठी \Log{LK} बरोबर पुढील
नियम मिळतात:
\[    \Axiom$\Gamma \fCenter \Delta, \varphi$
    \RightLabel{$\Box$}
    \UnaryInf$\Box\Gamma \fCenter \Diamond\Delta, \Box\varphi$
    \DisplayProof
    \qquad
    \Axiom$\varphi, \Gamma \fCenter \Delta$
    \RightLabel{$\Diamond$}
    \UnaryInf$\Diamond\varphi, \Box\Gamma \fCenter \Diamond\Delta$
    \DisplayProof
\]
~$\Log{K}$ च्या विस्तारांसाठीही आणखी नियम
जोडावे लागतात.

प्रत्येक मोडल तर्कशास्त्रासाठी अशा प्रकारचे
क्रमवर्ती कलन उपलब्ध असेलच असे नाही. अर्थदृष्ट्या
साध्या असलेल्या~$\Log{S5}$ साठीसुद्धा (त्याची
व्याख्या प्राप्यता-संबंध मुळीच न वापरता करता येते),
मूळ मसुद्यात~$\Log{LK}$ पासून मिळणारे आणि~\Cut
नियमाशिवाय संपूर्ण असणारे कलन ज्ञात नसल्याचे म्हटले आहे.
हे सध्याच्या सर्व क्रमवर्ती कलनांविषयीचे विधान समजू नये.
उदाहरणार्थ, \href{https://arxiv.org/abs/1711.04634v3}{Aghaei
आणि Mohammadi यांच्या 2018 मधील लेखात} $\Log{G3c}$ चा
विस्तार असलेले $\Log{G3S5}$ दिले आहे. कलम~3 मधील नियम
आणि प्रमेये~5.1 व~5.2 त्यातील कटची अनुमतता आणि संपूर्णता
देतात; सिद्धतेसाठी त्याच कलनाचे अर्धविराम वापरणारे रूप घेतले आहे.
या प्रकरणात पुढे दिलेल्या विशिष्ट $\Log{LK}$-विस्तारासाठी
कटची गरज असल्याचा निष्कर्ष वेगळा आहे.
$\Log{S5}$ साठी कट-मुक्त संपूर्ण \emph{हायपरक्रमवर्ती}
कलनही आहे.












\section{\Log{K} साठी नियम}\label{nml:seq:rul:sec}

नेहमीच्या विधानीय संयोजकांसाठीचे नियम हे
नेहमीच्या क्रमवर्ती कलनातील~$\Log{LK}$ नियमांसारखेच
आहेत. स्वयंसिद्धकेही तीच आहेत: $\varphi \Sequent \varphi$
या रूपातील कोणतीही क्रमवर्ती स्वयंसिद्धक असते.

मोडल संकारक~$\Box$
आणि~$\Diamond$ यांच्यासाठी
पुढील अतिरिक्त नियम
आहेत:
\[    \Axiom$\Gamma \fCenter \Delta, \varphi$
    \RightLabel{$\Box$}
    \UnaryInf$\Box\Gamma \fCenter \Diamond\Delta, \Box\varphi$
    \DisplayProof
    \qquad
    \Axiom$\varphi, \Gamma \fCenter \Delta$
    \RightLabel{$\Diamond$}
    \UnaryInf$\Diamond\varphi, \Box\Gamma \fCenter \Diamond\Delta$
    \DisplayProof
\]
येथे $\Box\Gamma$ म्हणजे~$\Gamma$
मधील प्रत्येक सूत्राच्या आधी~$\Box$
लावून~$\Gamma$ पासून मिळणारा सूत्रांचा अनुक्रम;
आणि $\Diamond\Delta$ म्हणजे~$\Delta$
मधील प्रत्येक सूत्राच्या आधी~$\Diamond$
लावून~$\Delta$ पासून मिळणारा सूत्रांचा अनुक्रम.


$\Gamma$~आणि~$\Delta$
रिकामे असू शकतात; अशा वेळी निष्कर्ष-क्रमवर्तीतील
त्यांचे अनुक्रमे
$\Box\Gamma$~आणि~$\Diamond\Delta$
हे भागही रिकामे असतात.

एका सूत्र~$\varphi$ पुरतेच
उजव्या बाजूला~$\Box$ जोडणे आणि
डाव्या बाजूला~$\Diamond$ जोडणे
मर्यादित ठेवणे आवश्यक आहे. जर
उजवीकडे कोणत्याही संख्येने सूत्रे ना~$\Box$
जोडण्याची किंवा डावीकडे
कोणत्याही संख्येने सूत्रे ना~$\Diamond$ जोडण्याची
मुभा असती, तर पुढील निष्पन्न करता आले असते:
  \[      \Axiom$\varphi \fCenter \varphi$
      \RightLabel{\RightR{\lnot}}
      \UnaryInf$\fCenter \varphi, \lnot \varphi$
      \RightLabel{$\Box*$}
      \UnaryInf$\fCenter \Box\varphi, \Box\lnot \varphi$
      \doubleLine
      \RightLabel{\RightR{\lor}}
      \UnaryInf$\fCenter \Box\varphi \lor \Box \lnot \varphi$
      \DisplayProof
    \qquad
      \Axiom$\varphi \fCenter \varphi$
      \RightLabel{\LeftR{\lnot}}
      \UnaryInf$\lnot \varphi, \varphi \fCenter$
      \RightLabel{$\Diamond*$}
      \UnaryInf$\Diamond\lnot \varphi,\Diamond\varphi \fCenter $
      \RightLabel{\RightR{\lnot}}
      \UnaryInf$\Diamond\varphi \fCenter \lnot\Diamond\lnot \varphi$
      \RightLabel{\RightR{\lif}}
      \UnaryInf$ \fCenter \Diamond\varphi \lif \lnot\Diamond\lnot \varphi$
      \DisplayProof
  \]
मात्र $\Box\varphi \lor \Box \lnot \varphi$ आणि
$\Diamond\varphi \lif \lnot\Diamond\lnot \varphi$ यांपैकी कोणतेही सूत्र~$\Log{K}$ मध्ये
वैध नाही.

जर आधारविधानात~$\varphi$ व्यतिरिक्त बाजूची सूत्रेही
असू दिली, आणि $\Box$ नियमाने उजवीकडे
फक्त~$\varphi$ ला~$\Box$ जोडू दिला किंवा
$\Diamond$ नियमाने डावीकडे
फक्त~$\varphi$ ला~$\Diamond$ जोडू दिला (पण बाजूच्या
सूत्रांवर कोणतीही क्रिया केली नाही), तर पुढील
निष्पन्न करता आले असते:
\[    \Axiom$\varphi \fCenter \varphi$
    \RightLabel{\RightR{\lnot}}
    \UnaryInf$\fCenter \varphi, \lnot \varphi$
    \RightLabel{\RightR{\Exchange}}
    \UnaryInf$\fCenter \lnot \varphi, \varphi$
    \RightLabel{$\Box*$}
    \UnaryInf$\fCenter \lnot\varphi, \Box \varphi$
    \doubleLine
    \RightLabel{\RightR{\lor}}
    \UnaryInf$\fCenter \lnot\varphi \lor \Box \varphi$
    \DisplayProof
  \qquad
    \Axiom$\varphi \fCenter \varphi$
    \RightLabel{\LeftR{\lnot}}
    \UnaryInf$\lnot \varphi, \varphi \fCenter$
    \RightLabel{$\Diamond*$}
    \UnaryInf$\Diamond\lnot \varphi, \varphi \fCenter $
    \RightLabel{\RightR{\lnot}}
    \UnaryInf$\varphi \fCenter \lnot\Diamond\lnot \varphi$
    \RightLabel{\RightR{\lif}}
    \UnaryInf$ \fCenter \varphi \lif \lnot\Diamond\lnot \varphi$
    \DisplayProof
\]
मात्र $\lnot\varphi \lor \Box \varphi$ (जे $\varphi
\lif \Box\varphi$ शी सममूल्य आहे) आणि $\varphi \lif \lnot\Diamond\lnot \varphi$
 यांपैकी कोणतेही
सूत्र~$\Log{K}$ मध्ये वैध नाही.













\section{\Log{K} साठी क्रमवर्ती निष्पत्ती}\label{nml:seq:prk:sec}



\begin{ex}
  $\Proves (\Box\varphi \land \Box\psi)
  \lif \Box (\varphi \land \psi)$ हे दाखवणारी क्रमवर्ती कलनातील
  निष्पत्ती पुढे दिली आहे.
  \begin{prooftree}
    \Axiom$\varphi \fCenter \varphi$
    \doubleLine
    \UnaryInf$\psi, \varphi \fCenter \varphi$
    \Axiom$\psi \fCenter \psi$
    \doubleLine
    \UnaryInf$\psi, \varphi \fCenter \psi$
    \RightLabel{\RightR{\land}}
    \BinaryInf$\psi, \varphi \fCenter \varphi \land \psi$
    \RightLabel{$\Box$}
    \UnaryInf$\Box\psi, \Box\varphi \fCenter \Box
    (\varphi \land \psi)$
    \RightLabel{\LeftR{\land}}
    \UnaryInf$\Box\varphi \land \Box\psi, \Box\varphi \fCenter \Box
    (\varphi \land \psi)$
    \RightLabel{\LeftR{\Exchange}}
    \UnaryInf$\Box\varphi, \Box\varphi \land \Box\psi \fCenter \Box
    (\varphi \land \psi)$
    \RightLabel{\LeftR{\land}}
    \UnaryInf$\Box\varphi \land \Box\psi, \Box\varphi \land \Box\psi \fCenter \Box (\varphi \land \psi)$
    \RightLabel{\LeftR{\Contraction}}
    \UnaryInf$\Box\varphi \land \Box\psi \fCenter \Box (\varphi \land \psi)$
    \RightLabel{\RightR{\lif}}
    \UnaryInf$\fCenter (\Box\varphi \land \Box\psi)
    \lif \Box (\varphi \land \psi)$
  \end{prooftree}
\end{ex}



  \begin{ex}
    $\Proves \Diamond(\varphi
    \lor \psi) \lif (\Diamond \varphi \lor \Diamond\psi)$ हे दाखवणारी
    क्रमवर्ती कलनातील निष्पत्ती पुढे दिली आहे.
    \begin{prooftree}
      \Axiom$\varphi \fCenter \varphi$
      \doubleLine
      \UnaryInf$\varphi \fCenter \varphi, \psi$
      \Axiom$\psi \fCenter \psi$
      \doubleLine
      \UnaryInf$\psi \fCenter \varphi, \psi$
      \RightLabel{\LeftR{\lor}}
      \BinaryInf$\varphi \lor \psi \fCenter \varphi, \psi$
      \RightLabel{$\Diamond$}
      \UnaryInf$\Diamond(\varphi \lor \psi) \fCenter
      \Diamond \varphi, \Diamond\psi$
      \RightLabel{\RightR{\lor}}
      \UnaryInf$\Diamond(\varphi \lor \psi) \fCenter
      \Diamond \varphi, \Diamond \varphi \lor \Diamond\psi$
      \RightLabel{\RightR{\Exchange}}
      \UnaryInf$\Diamond(\varphi \lor \psi) \fCenter \Diamond \varphi \lor \Diamond\psi, \Diamond \varphi$
      \RightLabel{\RightR{\lor}}
      \UnaryInf$\Diamond(\varphi \lor \psi) \fCenter
      \Diamond \varphi \lor \Diamond\psi, \Diamond \varphi \lor \Diamond\psi$
      \RightLabel{\RightR{\Contraction}}
      \UnaryInf$\Diamond(\varphi
      \lor \psi) \fCenter \Diamond \varphi \lor \Diamond\psi$
      \RightLabel{\RightR{\lif}}
      \UnaryInf$\fCenter \Diamond(\varphi
      \lor \psi) \lif (\Diamond \varphi \lor \Diamond\psi)$
    \end{prooftree}
  \end{ex}



  येथे~\Dual ची निष्पत्ती दिली आहे.
  या सिद्धतेत $\psi \liff \chi$ हे
  $(\psi \lif \chi) \land (\chi \lif \psi)$ चे संक्षिप्त रूप घेतले आहे;
  म्हणून शेवटची पायरी दोन सशर्त सूत्रांचा संयोग बनवते.
  \begin{prooftree}
    \Axiom$\varphi \fCenter \varphi$
    \RightLabel{\LeftR{\lnot}}
    \UnaryInf$\lnot \varphi, \varphi \fCenter $
    \RightLabel{$\Diamond$}
    \UnaryInf$\Diamond\lnot \varphi, \Box \varphi \fCenter $
    \RightLabel{\RightR{\lnot}}
    \UnaryInf$\Box \varphi \fCenter \lnot\Diamond\lnot \varphi$
    \RightLabel{\RightR{\lif}}
    \UnaryInf$\fCenter \Box \varphi \lif \lnot\Diamond\lnot \varphi$
    \Axiom$\varphi \fCenter \varphi$
    \RightLabel{\RightR{\lnot}}
    \UnaryInf$\fCenter \varphi, \lnot \varphi $
    \RightLabel{\RightR{\Exchange}}
    \UnaryInf$\fCenter \lnot \varphi, \varphi $
    \RightLabel{$\Box$}
    \UnaryInf$\fCenter \Diamond\lnot \varphi, \Box \varphi$
    \RightLabel{\RightR{\Exchange}}
    \UnaryInf$\fCenter \Box \varphi, \Diamond\lnot \varphi$
    \RightLabel{\RightR{\lnot}}
    \UnaryInf$\lnot\Diamond\lnot \varphi \fCenter \Box \varphi$
    \RightLabel{\RightR{\lif}}
    \UnaryInf$\fCenter \lnot \Diamond \lnot \varphi \lif \Box \varphi$
    \RightLabel{\RightR{\land}}
    \BinaryInf$\fCenter \Box \varphi \liff \lnot\Diamond\lnot \varphi$
  \end{prooftree}








\begin{prob}
  पुढील सूत्रे साठी~$\Log{K}$ मध्ये
  क्रमवर्ती कलनातील सिद्धता शोधा:
  \begin{enumerate}
    \item $\Box \lnot p \lif \Box(p \lif q)$
    \item $(\Box p \lor \Box q) \lif \Box(p \lor q)$
    \item $\Diamond p \lif \Diamond(p \lor q)$
    \item $\Box(p \land q) \lif \Box p$
  \end{enumerate}
\end{prob}













\section{इतर प्राप्यता-संबंधांसाठी नियम}\label{nml:seq:mru:sec}

विशिष्ट प्राप्यता-संबंधांनी निर्धारित होणाऱ्या
तर्कशास्त्रांचा विचार करण्यासाठी आपण
\ref{nml:seq:mru:tab:more-rules} मधील अतिरिक्त नियम पाहू.


\begin{table}
  \begin{center}
    \def\arraystretch{3}
    \begin{tabular}{cc}
    \hline
      \Axiom$\varphi, \Gamma \fCenter \Delta$
      \RightLabel{T$\Box$}
      \UnaryInf$\Box\varphi, \Gamma \fCenter \Delta$
      \DisplayProof
    &
      \Axiom$\Gamma \fCenter \Delta, \varphi$
      \RightLabel{T$\Diamond$}
      \UnaryInf$\Gamma \fCenter \Delta, \Diamond\varphi$
      \DisplayProof
    \\[1ex]
    \multicolumn{2}{c}{
      \Axiom$\Gamma \fCenter \Delta$
      \RightLabel{D}
      \UnaryInf$\Box\Gamma \fCenter \Diamond\Delta$
      \DisplayProof}
    \\[1ex]
      \Axiom$\Gamma, \Diamond\Pi \fCenter \Box\Delta, \Lambda, \varphi$
      \RightLabel{B$\Box$}
      \UnaryInf$\Box\Gamma, \Pi \fCenter \Delta, \Diamond\Lambda, \Box\varphi$
      \DisplayProof
      &
      \Axiom$\varphi, \Diamond\Gamma, \Pi \fCenter \Box\Lambda, \Delta$
      \RightLabel{B$\Diamond$}
      \UnaryInf$\Diamond\varphi, \Gamma, \Box\Pi \fCenter \Lambda, \Diamond\Delta$
      \DisplayProof
    \\[1ex]
      \Axiom$\Box\Gamma \fCenter \Diamond\Delta, \varphi$
      \RightLabel{4$\Box$}
      \UnaryInf$\Box\Gamma \fCenter \Diamond\Delta, \Box\varphi$
      \DisplayProof
      &
      \Axiom$\varphi, \Box\Gamma \fCenter \Diamond\Delta$
      \RightLabel{4$\Diamond$}
      \UnaryInf$\Diamond\varphi, \Box\Gamma \fCenter \Diamond\Delta$
      \DisplayProof
    \\[1ex]
      \Axiom$\Box\Gamma, \Diamond\Pi \fCenter \Box\Delta,
      \Diamond\Lambda, \varphi$
      \RightLabel{5$\Box$}
      \UnaryInf$\Box\Gamma, \Diamond\Pi \fCenter \Box\Delta,
      \Diamond\Lambda, \Box\varphi$
      \DisplayProof
&
      \Axiom$\varphi, \Diamond\Gamma, \Box\Pi \fCenter \Diamond\Delta, \Box\Lambda$
      \RightLabel{5$\Diamond$}
      \UnaryInf$\Diamond\varphi,\Diamond\Gamma,\Box\Pi \fCenter \Diamond\Delta,\Box\Lambda$
      \DisplayProof
    \\[1ex]
    \hline
    \end{tabular}
  \end{center}
  \caption{अधिक मोडल नियम.}
  \label{nml:seq:mru:tab:more-rules}
\end{table}
हे नियम जोडल्यावर \ref{nml:seq:mru:tab:logics-rules} मध्ये
दिलेल्या तर्कशास्त्रांसाठी निर्दोष आणि संपूर्ण
पद्धती मिळतात.
\begin{table}
  \begin{center}
    \begin{tabular}{lll}
      \hline
      तर्कशास्त्र & $R$ हा \dots आहे & नियम\\
      \hline
      $\Log{T} = \Log{KT}$ & परावर्ती & $\Box$,
      $\Diamond$,
      T$\Box$
      ,
      T$\Diamond$
      \\ \hline
      $\Log{D} = \Log{KD}$ & सर्वत्र उत्तरवर्ती & $\Box$,
      $\Diamond$,
        D
      \\ \hline
      $\Log{K4}$ & संक्रमक & $\Box$,
      $\Diamond$,
      4$\Box$
      ,
      4$\Diamond$
      \\ \hline
      $\Log{B} = \Log{KTB}$ & परावर्ती, & $\Box$,
      $\Diamond$,
      T$\Box$
      ,
      T$\Diamond$\\
      & सममित &
      B$\Box$
      ,
      B$\Diamond$
      \\ \hline
      $\Log{S4} = \Log{KT4}$ & परावर्ती, & $\Box$,
      $\Diamond$,
      T$\Box$
      ,
      T$\Diamond$\\
      & संक्रमक &
      4$\Box$
      ,
      4$\Diamond$
      \\ \hline
      $\Log{S5} = \Log{KT5}$ & परावर्ती, & $\Box$,
      $\Diamond$,
      T$\Box$
      ,
      T$\Diamond$\\
      & संक्रमक, &
      5$\Box$
      ,
      5$\Diamond$\\
      &  युक्लिडी &
      \\ \hline
    \end{tabular}
  \end{center}
  \caption{विविध मोडल तर्कशास्त्रांसाठी क्रमवर्ती नियम.}
  \label{nml:seq:mru:tab:logics-rules}
\end{table}
\begin{ex}
  $\Log{K4} \Proves \Ax{4}$, म्हणजे
  $\Box\varphi \lif \Box\Box\varphi$, हे दाखवणारी क्रमवर्ती
  निष्पत्ती पुढे दिली आहे.
  \begin{prooftree}
    \Axiom$\Box\varphi \fCenter \Box\varphi$
    \RightLabel{4$\Box$}
    \UnaryInf$\Box\varphi \fCenter \Box\Box\varphi$
    \RightLabel{\RightR{\lif}}
    \UnaryInf$\fCenter \Box\varphi \lif \Box\Box\varphi$
  \end{prooftree}
\end{ex}
\begin{ex}
  $\Log{S5} \Proves \Ax{5}$, म्हणजे
  $\Diamond\varphi \lif \Box\Diamond\varphi$, हे दाखवणारी
  क्रमवर्ती निष्पत्ती पुढे दिली आहे.
  \begin{prooftree}
    \Axiom$\Diamond\varphi \fCenter \Diamond\varphi$
    \RightLabel{5$\Box$}
    \UnaryInf$\Diamond\varphi \fCenter \Box\Diamond\varphi$
    \RightLabel{\RightR{\lif}}
    \UnaryInf$\fCenter \Diamond\varphi \lif \Box\Diamond\varphi$
  \end{prooftree}
\end{ex}
\begin{ex}
या विभागात दिलेले \Log{S5} साठीचे क्रमवर्ती कलन \Cut{} नियमाशिवाय
संपूर्ण नाही; उदा., $\Diamond\Box \varphi \lif \varphi$ हे~$\Log{S5}$ मध्ये वैध असलेले
सूत्र~\Cut शिवाय सिद्ध करता येत नाही. पुढे~\Cut
वापरणारी निष्पत्ती दिली आहे:
\begin{prooftree}
\Axiom$\Box \varphi \fCenter \Box\varphi$
\RightLabel{5$\Diamond$}
\UnaryInf$\Diamond\Box\varphi \fCenter \Box\varphi$
\Axiom$\varphi \fCenter \varphi$
\RightLabel{T$\Box$}
\UnaryInf$\Box\varphi \fCenter \varphi$
\RightLabel{\Cut}
\BinaryInf$\Diamond\Box\varphi \fCenter \varphi$
\RightLabel{\RightR{\lif}}
\UnaryInf$\fCenter \Diamond\Box\varphi \lif \varphi$
\end{prooftree}
\end{ex}
\begin{prob}
पुढील निष्कर्ष दाखवणाऱ्या क्रमवर्ती निष्पत्त्या द्या:
  \begin{enumerate}
  \item $\Log{KT5} \Proves \Ax{B}$;
  \item $\Log{KT5} \Proves \Ax{4}$;
  \item $\Log{KDB4} \Proves \Ax{T}$;
  \item $\Log{KB4} \Proves \Ax{5}$;
  \item $\Log{KB5} \Proves \Ax{4}$;
  \item $\Log{KT} \Proves \Ax{D}$.
  \end{enumerate}
\end{prob}
\part{उपयोजित मोडल तर्कशास्त्र}\label{aml:part}
\begin{editorial}
या भागात कालिक आणि ज्ञानविषयक तर्कशास्त्रांसारख्या
मोडल तर्कशास्त्राच्या काही उपयोजनांवरील प्रायोगिक
मसुदा साहित्य आहे.
\end{editorial}
\chapter{कालिक तर्कशास्त्रे}\label{aml:tl:chap}
\section{प्रस्तावना}\label{aml:tl:int:sec}
कालिक तर्कशास्त्रे यापुढे काय घडेल किंवा यापूर्वी
काय घडले आहे यांविषयीच्या विधानांचा विचार करतात.
आर्थर प्रायर यांना कालिक तर्कशास्त्राचे जनक मानले
जाते; त्यांनी त्याला \emph{काळवाचक तर्कशास्त्र}
असे नाव दिले. विधानीय तर्कशास्त्राची मूलभूत चौकट
विधाने साधारणतः काळवाचक भेद नसलेली मानते. त्या
चौकटीत कालिक संकारक घालण्याच्या प्रायर यांच्या
मूळ मोडल पद्धतीचे येथे आपण मुख्यतः अनुसरण करू.
उदाहरणार्थ, विधानीय तर्कशास्त्रात मी बीझी नावाच्या
कुत्र्याबद्दल बोलू शकतो. कुत्र्यांच्या स्वभावाप्रमाणे
ते कधी बसते आणि कधी बसत नाही. बीझी बसले आहे
आणि बीझी बसलेले नाही असे एकत्र म्हणणे अभिजात
तर्कशास्त्रात विरोधाभासी ठरेल. पण ही दोन्ही विधाने
खरी असू शकतात, फक्त एकाच वेळी नाहीत. भाषेत
कालिक संकारक घातल्यास ही गोष्ट तुलनेने सहज
व्यक्त करता येते. कालिक संकारकांमुळे ``बीझीला
खाऊ किंवा चेंडू मिळेल'' यावरून ``बीझीला खाऊ
मिळेल किंवा बीझीला चेंडू मिळेल'' यांसारख्या
अनुमानांची वैधताही स्पष्ट करता येते.
मात्र कालिक तर्कशास्त्रामुळे अनेक तत्त्वज्ञानविषयक
प्रश्न निर्माण होतात आणि त्यामुळे त्याची एक चौकट
निवडण्याऐवजी दुसरी निवडावी लागू शकते. उदाहरणार्थ,
भविष्यविषयक परिस्थितीवश विधान म्हणजे भविष्याबद्दलचे
असे विधान की जे अनिवार्यही नाही आणि अशक्यही नाही.
``रिचर्ड उद्या किराणा दुकानात जाईल'' असे म्हटल्यास,
आपण अजून न घडलेल्या घटनेविषयी विधान करतो; त्या
विधानाचे सत्यतामूल्य वादग्रस्त असते. किंबहुना,
त्याला मुळात सत्यतामूल्य \emph{देता} तरी येते का,
हाही वादाचा विषय आहे. आपण कठोर निर्धारणवादी
असू, तर संबंधित घटना घडण्यापूर्वीच हे विधान
खरे किंवा खोटे असते, एवढेच की त्याचे सत्यतामूल्य
आपल्याला अजून माहीत नसते, असे मानणे आपल्याला
मान्य होऊ शकते. याउलट, भविष्य खऱ्या अर्थाने
खुले आहे आणि भविष्यविषयक परिस्थितीवश विधानांची
सत्यतामूल्ये अजून अनिर्धारित आहेत, असेही आपण
मानू शकतो.
काळाची रचना आणि स्वरूप यांविषयीच्या अशा अनेक
भूमिका कालिक तर्कशास्त्रासाठी आपण निवडलेल्या
प्रतिमानांत आणि सैद्धांतिक चौकटींत अंतर्भूत असतात.
उदाहरणार्थ, काळ रेषीय, शाखायुक्त किंवा अगदी
वर्तुळाकार आहे अशी प्रतिमाने तयार करावीत का,
हा प्रश्न आपण विचारू शकतो. कालिक प्रतिमानांना
आरंभबिंदू आणि अंतबिंदू असावेत का, तसेच काळाचे
निरूपण विविक्त क्षणांनी करावे की सांतत्य म्हणून
करावे, हेही ठरवावे लागेल.
\section{कालिक तर्कशास्त्राचे अर्थनिर्धारण}\label{aml:tl:sem:sec}
\begin{defn}
कालिक तर्कशास्त्राच्या मूलभूत भाषेत पुढील गोष्टी असतात:
\begin{enumerate}
  \item असत्यता दर्शवणारा विधानीय स्थिरांक~$\lfalse$.
  \item सत्यता दर्शवणारा विधानीय स्थिरांक~$\ltrue$.
  \item विधानीय चलांचा गणनीय अनंत संच: $\Obj
    p_0$, $\Obj p_1$, $\Obj p_2$, \dots
  \item विधानीय संयोजके: \startycommalist
  \ycomma $\lnot$ (नकार)
  \ycomma $\land$ (संयोग)
  \ycomma $\lor$ (विकल्पयोग)
  \ycomma $\lif$ (शर्त)
  \ycomma $\liff$ (द्विशर्त).
  \item भूतकाळी संकारक $\Ptemp$ आणि $\Htemp$.
  \item भविष्यकाळी संकारक $\Ftemp$ आणि $\Gtemp$.
\end{enumerate}
\end{defn}
पुढे इतर प्रकारचे मोडल संकारक जोडण्याची शक्यताही पाहू.
\begin{defn}
कालिक भाषेतील \emph{सूत्रे} ची पुढीलप्रमाणे
  आगमनाने व्याख्या करतो:
\begin{enumerate}
\item $\lfalse$ हे अण्विक सूत्र आहे.
\item $\ltrue$ हे अण्विक सूत्र आहे.
\item प्रत्येक विधानीय चल~$\Obj p_i$ हे (अण्विक) सूत्र आहे.
\item $\varphi$ हे सूत्र असेल, तर $\lnot \varphi$
  हेही सूत्र आहे.
\item $\varphi$ आणि $\psi$ ही सूत्रे असतील, तर $(\varphi \land
  \psi)$ हे सूत्र आहे.
\item $\varphi$ आणि $\psi$ ही सूत्रे असतील, तर $(\varphi \lor \psi)$
  हे सूत्र आहे.
\item $\varphi$ आणि $\psi$ ही सूत्रे असतील, तर $(\varphi \lif \psi)$
  हे सूत्र आहे.
\item $\varphi$ आणि $\psi$ ही सूत्रे असतील, तर $(\varphi \liff \psi)$
  हे सूत्र आहे.
\item $\varphi$ हे सूत्र असेल, तर $\Ptemp  \varphi$, $\Htemp  \varphi$, $\Ftemp \varphi$, $\Gtemp  \varphi$
  ही सर्व सूत्रे आहेत.
\item याखेरीज दुसरे काहीही सूत्र नाही.
\end{enumerate}
\end{defn}
इतर आशयात्मक तर्कशास्त्रांप्रमाणे कालिक तर्कशास्त्रांचे
अर्थनिर्धारणही संबंधी प्रतिमानांच्या आधारे केले जाते.
\begin{defn}
  कालिक भाषेसाठीचे \emph{प्रतिमान} हे
  $\mModel{M} = \tuple{T, \prec, V}$ असे त्रिक असते. येथे
  \begin{enumerate}
  \item $T$ हा रिक्तेतर संच असून त्यातील घटकांचा अर्थ
    काळातील बिंदू असा लावला जातो.
  \item $\prec$ हा $T$ वरील द्विपदी संबंध आहे.
  \item $V$ हे प्रत्येक विधानीय चल~$p$ ला कालबिंदूंचा संच $V(p)$ देणारे फलन आहे.
  \end{enumerate}
  $t \prec t'$ असेल, तर $t$ हा~$t'$ च्या \emph{आधी येतो}
  असे म्हणतो. $t \in V(p)$ असेल, तर $p$ हे~$t$ येथे
  \emph{सत्य आहे} असे म्हणतो.
\end{defn}
सध्या आपण~$\prec$ या आधी येण्याच्या संबंधावर कोणतीही अट
घातलेली नाही, हे लक्षात येईल. त्यामुळे कालिक
प्रतिमानांच्या रचनेवर अजून कोणतेही निर्बंध नाहीत:
काळ रेषीय, शाखायुक्त, वर्तुळाकार किंवा इतर
कोणत्याही रचनेचा असलेली प्रतिमाने आपण घेऊ शकतो.
सामान्य मोडल तर्कशास्त्राप्रमाणे, प्रत्येक कालिक
प्रतिमानात कोणती सूत्रे कोणत्या बिंदूंवर
सत्य ठरतात हे निश्चित होते. संबंधी अर्थनिर्धारणाच्या
मूलभूत कल्पनेसाठी आपण तीच संज्ञा वापरतो:
``प्रतिमान~$\mModel{M}$ मध्ये सूत्र~$\varphi$
बिंदू~$t$ येथे सत्य ठरते.'' या संबंधाची आगमनाने
व्याख्या केली जाते; मोडल नसलेल्या सर्व संकारकांसाठी
ती सामान्य मोडल तर्कशास्त्रातील व्याख्येसारखीच आहे.
\begin{defn}\label{aml:tl:sem:defn:tmodels}
  प्रतिमान~$\mModel M$ मध्ये~$t$ येथे सूत्र~$\varphi$
  ची \emph{सत्यता}, म्हणजेच $\mSat{M}{\varphi}[t]$,
  पुढीलप्रमाणे आगमनाने परिभाषित होते:
  \begin{enumerate}
  \item \indcase{\varphi}{\lfalse}{$\mSat{M}{\lfalse}[t]$ कधीही नाही}.
  \item \indcase{\varphi}{\ltrue}{$\mSat{M}{\ltrue}[t]$ नेहमी असते}.
  \item $\mSat{M}{p}[t]$ जर आणि फक्त जर $t \in V(p)$
  \item \indcase{\varphi}{\lnot \psi}{$\mSat{M}{\indfrm}[t]$ जर आणि फक्त जर
    $\mSat/{M}{\psi}[t]$}.
  \item \indcase{\varphi}{(\psi \land \chi)}{$\mSat{M}{\indfrm}[t]$ जर आणि फक्त जर
    $\mSat{M}{\psi}[t]$ आणि $\mSat{M}{\chi}[t]$}.
  \item \indcase{\varphi}{(\psi \lor \chi)}{$\mSat{M}{\indfrm}[t]$ जर आणि फक्त जर
    $\mSat{M}{\psi}[t]$ किंवा $\mSat{M}{\chi}[t]$} (किंवा दोन्ही).
  \item \indcase{\varphi}{(\psi \lif \chi)}{$\mSat{M}{\indfrm}[t]$ जर आणि फक्त जर
    $\mSat/{M}{\psi}[t]$ किंवा $\mSat{M}{\chi}[t]$}.
  \item \indcase{\varphi}{(\psi \liff \chi)}{$\mSat{M}{\indfrm}[t]$ जर आणि फक्त जर
    एकतर $\mSat{M}{\psi}[t]$ आणि $\mSat{M}{\chi}[t]$ दोन्ही, किंवा
    $\mSat{M}{\psi}[t]$ आणि $\mSat{M}{\chi}[t]$ दोन्ही नाहीत}.
  \item\label{aml:tl:sem:defn:sub:mmodels-p}
    \indcase{\varphi}{\Ptemp  \psi}{$\mSat{M}{\indfrm}[t]$ जर आणि फक्त जर
    असा एखादा $t' \in T$ असेल की $t' \prec t$ आणि $\mSat{M}{\psi}[t']$}
\item\label{aml:tl:sem:defn:sub:mmodels-h}
    \indcase{\varphi}{\Htemp  \psi}{$\mSat{M}{\indfrm}[t]$ जर आणि फक्त जर
    $t' \prec t$ असलेल्या प्रत्येक $t' \in T$ साठी $\mSat{M}{\psi}[t']$}
   \item\label{aml:tl:sem:defn:sub:mmodels-f}
    \indcase{\varphi}{\Ftemp  \psi}{$\mSat{M}{\indfrm}[t]$ जर आणि फक्त जर
    असा एखादा $t' \in T$ असेल की $t \prec t'$ आणि $\mSat{M}{\psi}[t']$}
\item\label{aml:tl:sem:defn:sub:mmodels-g}
    \indcase{\varphi}{\Gtemp  \psi}{$\mSat{M}{\indfrm}[t]$ जर आणि फक्त जर
    $t \prec t'$ असलेल्या प्रत्येक $t' \in T$ साठी $\mSat{M}{\psi}[t']$}
  \end{enumerate}
\end{defn}
अर्थनिर्धारणावरून $\Ptemp$ आणि~$\Htemp$ हे
परस्पर द्वैती संकारक आहेत, तसेच $\Ftemp$ आणि
$\Gtemp$ हेही. म्हणून $\Htemp  \varphi$ ची व्याख्या
$\lnot \Ptemp \lnot \varphi$ अशी करता येईल; $\Gtemp$
आणि~$\Ftemp$ यांनाही हेच लागू होते.
\section{कालिक चौकटींचे गुणधर्म}\label{aml:tl:acc:sec}
आपली कालिक प्रतिमाने~$\prec$ या संबंधावर कोणतीही
अट घालत नसल्यामुळे, आपल्या परिचित स्वयंसिद्धकांपैकी
सर्व प्रतिमानांत लागू असणारे एकमेव स्वयंसिद्धक~$K$
आहे; किंवा त्याची $K_{\Gtemp}$ आणि~$K_{\Htemp}$
ही अनुरूप रूपे आहेत:
\begin{align*}
\tag{$K_{\Gtemp}$}  \Gtemp (p \to q) & \to (\Gtemp p \to \Gtemp q)\\
\tag{$K_{\Htemp}$}  \Htemp (p \to q) & \to (\Htemp p \to \Htemp q)
\end{align*}
मात्र प्रतिमानांनी काळाच्या निरूपणावर अधिक कडक
अटी घालाव्यात असे आपल्याला वाटत असेल—उदा.,
$\prec$ हा रेषीय क्रम असावा—तर त्या प्रतिमानांत
इतर सूत्रेही वैध ठरतील.
\begin{table}[t]
  \begin{tabular}{| p{.465\textwidth} || p{.465\textwidth} |}
    \hline
    {\emph{$\prec$ हा \dots असेल}} & {\emph{तर~$\mModel{M}$ मध्ये \dots सत्य आहे:}} \\
    \hline \hline
    \emph{संक्रमक}: \newline
    $\forall u \forall v \forall w ((u \prec v \land v \prec w) \lif u \prec w)$ &
    $\Ftemp \Ftemp p \lif \Ftemp p$  \\
    \hline
    \emph{रेषीय}: \newline
    $\forall w \forall v (w \prec v \lor w = v \lor v \prec w)$ &
    $(\Ftemp \Ptemp  p \lor \Ptemp \Ftemp p) \lif (\Ptemp p \lor p \lor \Ftemp p)$\\
    \hline
    \emph{सघन}: \newline
    $\forall w \forall v (w \prec v \to \exists u(w \prec u \land u \prec v))$ &
    $\Ftemp p \lif \Ftemp \Ftemp p$ \\
    \hline
    \emph{अपरिबद्ध (भूतकाळाकडे)}: \newline
    $\forall w \exists v( v \prec w)$ &
    $\Htemp p \to \Ptemp p$ \\
    \hline
    \emph{अपरिबद्ध (भविष्याकडे)}: \newline
    $\forall w \exists v( w \prec v)$ &
    $\Gtemp p \to \Ftemp p$ \\
    \hline
  \end{tabular}
  \caption{कालिक चौकटींचे काही अनुरूपता-गुणधर्म.}
  \label{aml:tl:acc:tab:correspondence}
\end{table}
\ref{aml:tl:acc:tab:correspondence} मधील अनेक गुणधर्म काळाचे
निरूपण करण्यासाठी बनवलेल्या प्रतिमानांत अपेक्षित
वाटू शकतात. तरीही, $\prec$ संबंधावर आपल्याला
हव्या त्या अटी घालता येत असल्या, तरी सर्वच अटी
कालिक तर्कशास्त्राच्या भाषेत व्यक्त करता येणाऱ्या
सूत्रे द्वारे दर्शवता येतात असे नाही. उदाहरणार्थ,
अपरावर्तीपणा, म्हणजे $\forall w \lnot (w \prec w)$
ही अट, कालिक तर्कशास्त्रातील कोणत्याही सूत्राशी
अनुरूप नाही.
\section{कालिक तर्कशास्त्रासाठी अतिरिक्त संकारक}\label{aml:tl:ext:sec}
भूतकाळ व भविष्यकाळासाठीच्या एकस्थानी संकारकांखेरीज,
कालिक तर्कशास्त्रात कधीकधी ``पासून'' आणि ``पर्यंत''
दर्शवणारे द्विस्थानी संकारक $\Since $ आणि~$\Until$
यांचाही समावेश असतो. म्हणजे कालिक तर्कशास्त्राच्या
भाषेत $\Since $ आणि~$\Until$ जोडून, कालिक
सूत्राच्या व्याख्येत पुढील नियम जोडायचा:
\begin{itemize}
\item[] $\varphi$ आणि $\psi$ ही सूत्रे असतील, तर
  $(\Since  \varphi \psi)$ आणि $(\Until \varphi \psi)$ ही दोन्ही
  सूत्रे आहेत.
\end{itemize}
या संकारकांचे अर्थनिर्धारण पुढीलप्रमाणे केले जाते:
\begin{defn}\label{aml:tl:ext:defn:since-until}
  प्रतिमान~$\mModel M$ मध्ये~$t$ येथे सूत्र~$\varphi$
  ची \emph{सत्यता}:
  \begin{enumerate}
  \item\label{aml:tl:ext:defn:sub:mmodels-since} \indcase{\varphi}{\Since  \psi
    \chi}{$\mSat{M}{\indfrm}[t]$ जर आणि फक्त जर असा एखादा $t' \in
    T$ असेल की $t' \prec t$ आणि $\mSat{M}{\psi}[t']$; तसेच
    $t' \prec s \prec t$ असलेल्या प्रत्येक $s$ साठी
    $\mSat{M}{\chi}[s]$}
   \item\label{aml:tl:ext:defn:sub:mmodels-until} \indcase{\varphi}{\Until \psi
    \chi}{$\mSat{M}{\indfrm}[t]$ जर आणि फक्त जर असा एखादा $t' \in
    T$ असेल की $t \prec t'$ आणि $\mSat{M}{\psi}[t']$; तसेच
    $t \prec s \prec t'$ असलेल्या प्रत्येक $s$ साठी
    $\mSat{M}{\chi}[s]$}
  \end{enumerate}
\end{defn}
$\Since  \psi \chi$ चा सहज अर्थ असा: ``ज्या
वेळेपासून~$\psi$ सत्य होते, त्या वेळेपासून~$\chi$
सत्य राहिले आहे.'' आणि $\Until \psi \chi$ चा सहज
अर्थ: ``ज्या वेळेपर्यंत~$\psi$ सत्य होईल, त्या
वेळेपर्यंत~$\chi$ सत्य राहील.''
\section{संभाव्य इतिहास}\label{aml:tl:poss:sec}
आतापर्यंत वापरलेली कालिक तर्कशास्त्राची संबंधी
प्रतिमाने अत्यंत लवचिक आहेत, कारण प्राप्यता संबंधावर
कोणतेही निर्बंध घालावे लागत नाहीत. त्यामुळे कालिक
प्रतिमाने भूतकाळात आणि भविष्यकाळात दोन्हीकडे
शाखायुक्त असू शकतात. पण शाखा फुटण्याची अधिक
``मोडल'' कल्पना आपण विचारात घेऊ शकतो: घटनांचे
अनुक्रम हे संभाव्य इतिहास मानायचे. यासाठी भाषा
बदलणे आवश्यकच नाही; तरीही आपण ``नेहमीचे''
मोडल संकारक $\Box$ आणि~$\Diamond$ जोडू शकतो.
कालांतराने कर्त्यांच्या ज्ञानात होणारे बदल दर्शवण्यासाठी
ज्ञानविषयक प्राप्यता संबंध जोडण्याचाही विचार करता येतो.
\begin{defn}
  कालिक भाषेसाठीचे \emph{संभाव्य इतिहासांचे प्रतिमान}
  हे $\mModel{M} = \tuple{T, C, V}$ असे त्रिक असते. येथे
  \begin{enumerate}
    \item $T$ हा रिक्तेतर संच असून त्यातील घटकांचा अर्थ
      काळातील अवस्था असा लावला जातो.
    \item $C$ हा प्रणालीच्या संगणकीय मार्गांचा, म्हणजे
      \emph{संभाव्य इतिहासांचा}, संच आहे. दुसऱ्या शब्दांत,
      $C$ हा $s_1$, $s_2$,~$s_3$, \dots अशा अवस्थांच्या
      $\sigma$ या अनुक्रमांचा संच आहे; येथे
      प्रत्येक $s_i \in T$. येथे $s_i$ या काळातील स्वतंत्र
      अवस्था-घटना मानतो: एका इतिहासात एकाच $T$-घटकाची
      पुनरावृत्ती होत नाही. प्रणालीची तीच स्थिती पुन्हा आली,
      तरी तिच्या वेगवेगळ्या काळातील घटनांना वेगळे घटक मानायचे.
      त्यामुळे $t,\sigma$ याने त्या इतिहासातील एकच स्थान
      निश्चित होते आणि $\prec_\sigma$ चा पुढील अर्थ संदिग्ध राहत नाही.
    \item $V$ हे प्रत्येक विधानीय चल~$p$ ला कालबिंदूंचा संच~$V(p) \subseteq T$ देणारे फलन आहे.
  \end{enumerate}
  गोष्टी सोप्या करण्यासाठी, एखादा इतिहास~$C$ मध्ये
  असेल, तर त्याचे सुरुवातीचे घटक वगळून मिळणारे
  सर्व उर्वरित अनुक्रमही त्यात आहेत, असे आपण साधारणतः
  मानू. उदाहरणार्थ, $s_1$, $s_2$,~$s_3$ हा अनुक्रम~$C$
  मध्ये असेल, तर $s_2$, $s_3$ आणि~$s_3$ हेही त्यात
  असतात. तसेच, अवस्था $s_i$ आणि~$s_j$ या
  अनुक्रम~$\sigma$ मध्ये असतील, तर $i < j$ असताना
  $s_i \prec_\sigma s_j$ असे म्हणतो.
  $t \in V(p)$ असेल, तर $p$ हे~$t$ येथे
  \emph{सत्य आहे} असे म्हणतो.
\end{defn}
येथे संबंधित बदल असा आहे: प्रतिमान~$\mModel M$
मध्ये कालबिंदू~$t$ येथे सूत्राची सत्यता
तपासताना आपण ती $t$ ही अवस्था-घटना असलेल्या
इतिहास~$\sigma \in C$ च्या सापेक्ष तपासतो.
विधानीय चलांची सत्यता फक्त $t$ वर अवलंबून असते:
$\mSat{M}{p}[t,\sigma]$ जर आणि फक्त जर $t\in V(p)$.
सत्यता-फलनात्मक संयोजकांच्या अटी
\ref{aml:tl:sem:defn:tmodels} मधीलप्रमाणेच राहतात;
मात्र त्यांतील प्रत्येक उपसूत्राची सत्यता त्याच
$t,\sigma$ येथे तपासायची. उदाहरणार्थ,
$\mSat{M}{\lnot \psi}[t,\sigma]$ जर आणि फक्त जर
$\mSat/{M}{\psi}[t,\sigma]$; आणि
$\mSat{M}{\psi\land \chi}[t,\sigma]$ जर आणि फक्त जर
$\mSat{M}{\psi}[t,\sigma]$ आणि $\mSat{M}{\chi}[t,\sigma]$.
संयोजक स्वतः इतिहास बदलत नसले, तरी त्यांची कालिक किंवा
मोडल उपसूत्रे त्या इतिहासावर अवलंबून असू शकतात.
आता या इतिहासांच्या सापेक्ष आपण
भविष्यकाळी संकारक~$\Ftemp$ ची पुनर्व्याख्या करतो
आणि~$\Diamond$ संकारक जोडतो.
\begin{defn}\label{aml:tl:poss:defn:phmodels}
  प्रतिमान~$\mModel M$ मध्ये~$t, \sigma$ येथे
  सूत्र~$\varphi$ ची \emph{सत्यता}, म्हणजेच
  $\mSat{M}{\varphi}[t, \sigma]$:
  \begin{enumerate}
  \item\label{aml:tl:poss:defn:sub:phmodels-f} \indcase{\varphi}{\Ftemp
    \psi}{$\mSat{M}{\indfrm}[t, \sigma]$ जर आणि फक्त जर
    असा एखादा $t' \in T$ असेल की $t \prec_\sigma t'$
    आणि $\mSat{M}{\psi}[t', \sigma]$.}
  \item\label{aml:tl:poss:defn:sub:phmodels-diamond} \indcase{\varphi}{\Diamond
    \psi}{$\mSat{M}{\indfrm}[t, \sigma]$ जर आणि फक्त जर
    असा एखादा $\sigma' \in C$ असेल की त्यात $t$ येतो
    आणि $\mSat{M}{\psi}[t, \sigma']$.}
  \end{enumerate}
\end{defn}
इतर कालिक आणि मोडल संकारकांचीही अशीच व्याख्या
करता येते. आता आपण काळवाचकता आणि मोडलता
एकत्र करणारी विधाने दर्शवू शकतो. उदाहरणार्थ,
``$p$ पुढे घडणार नाही, पण घडू शकला असता'' हे
विधान $\lnot \Ftemp p \land \Diamond \Ftemp p$
या सूत्राने दर्शवता येईल. ज्या बिंदू $t$ आणि इतिहास $\sigma$ येथे
त्या इतिहासातील कोणत्याही पुढील अवस्थेत $p$ सत्य नसते,
पण $t$ असलेल्या एखाद्या पर्यायी इतिहासात $p$ पुढील
एखाद्या अवस्थेत सत्य होते, तेथे हे सूत्र सत्य ठरेल.
\chapter{ज्ञानविषयक तर्कशास्त्रे}\label{aml:el:chap}
\begin{editorial}
  या प्रकरणात ज्ञानविषयक तर्कशास्त्रांचा अधिसिद्धांत
  मांडला आहे. याची रचना अल्डो अँटोनेली यांच्या
  अभिजात मूलभूत मोडल तर्कशास्त्रावरील नोंदींसारखी
  आहे; मात्र द्विअनुकरण आणि गतिशील ज्ञानविषयक
  तर्कशास्त्रांवरील साहित्य जोडण्यासाठी ऑड्री यॅप
  यांनी त्याचे पुनर्लेखन केले आहे.
\end{editorial}
\section{प्रस्तावना}\label{aml:el:int:sec}
मोडल तर्कशास्त्र जसे \emph{मोडल विधानांचा} आणि
त्यांमधील तार्किक निष्पन्नतेच्या संबंधांचा विचार
करते, तसे ज्ञानविषयक तर्कशास्त्र \emph{ज्ञानविषयक
विधानांचा} आणि त्यांमधील निष्पन्नता संबंधांचा विचार
करते. मोडल संकारकांचा अर्थ संभाव्यता आणि
अनिवार्यता असा लावण्याऐवजी, एकस्थानी संयोजकांना
ज्ञानविषयक किंवा विश्वासविषयक अर्थ दिला जातो;
त्यायोगे ज्ञान आणि विश्वासाचे प्रतिमानीकरण करता
येते. उदाहरणार्थ, पुढील विधाने व्यक्त करायची
असू शकतात:
\begin{enumerate}
\item कॅलगरी हे अल्बर्टात आहे, हे रिचर्डला माहीत आहे.
\item सोफ्यावर कुत्रा असू शकतो, असे ऑड्रीला वाटते.
\item मंगळवारी तिचा वर्ग असतो, हे ऑड्रीला माहीत
  आहे, हे रिचर्डला माहीत आहे.
\item एका वर्षात १२ महिने असतात, हे सर्वांना माहीत आहे.
\end{enumerate}
आधुनिक ज्ञानविषयक तर्कशास्त्राचा उगम अनेकदा
याक्को हिंटिक्का यांच्या १९६२ मधील \emph{Knowledge
and Belief} या ग्रंथाशी जोडला जातो. तो ग्रंथ
लिहिला गेला तेव्हा संभाव्य जगांचे अर्थनिर्धारण
तर्कशास्त्रात वाढत्या प्रमाणात वापरले जात होते.
वास्तविक, ज्ञानविषयक तर्कशास्त्रे इतर मोडल
तर्कशास्त्रांसारखीच बहुतेक अर्थनिर्धारणाची साधने
वापरतात; मात्र त्यांचा अर्थ वेगळा लावतात. मुख्य
बदल म्हणजे \emph{प्राप्यता संबंध} कशाचे प्रतिनिधित्व
करतो, हे. ज्ञानविषयक तर्कशास्त्रांत तो काही प्रकारची
\emph{ज्ञानविषयक संभाव्यता} दर्शवतो. आपण कोणत्या
ज्ञानविषयक संकल्पनेचे प्रतिमानीकरण करतो त्यावर
प्राप्यता संबंधावर कोणत्या अटी घालायच्या हे अवलंबून
असेल, हे पुढे पाहू. अनुरूपता सिद्धांताला ज्ञानविषयक
अर्थ दिल्यावर काय घडते तेही पाहू. वरील उदाहरणांत
रिचर्ड आणि ऑड्री हे दोन कर्ते आहेत, आणि प्रत्येकीला
काय माहीत आहे यांतील संबंध येतो. आपण विचारात
घेणार असलेल्या ज्ञानविषयक तर्कशास्त्रांत अनेक कर्ते
असतील, म्हणून अशा गोष्टी व्यक्त करता येतील.
याउलट, एका कर्त्याचे ज्ञानविषयक तर्कशास्त्र फक्त
एकाच व्यक्तीला काय माहीत आहे किंवा काय वाटते
याबद्दल बोलेल.
\section{ज्ञानविषयक तर्कशास्त्राची भाषा}\label{aml:el:lan:sec}
\begin{defn}
$G$ हा कर्तृचिन्हांचा संच असू द्या. अनेक कर्त्यांच्या
ज्ञानविषयक तर्कशास्त्राच्या मूलभूत भाषेत पुढील गोष्टी
असतात:
\begin{enumerate}
  \item असत्यता दर्शवणारा विधानीय स्थिरांक~$\lfalse$.
  \item सत्यता दर्शवणारा विधानीय स्थिरांक~$\ltrue$.
  \item विधानीय चलांचा गणनीय अनंत संच: $\Obj
    p_0$, $\Obj p_1$, $\Obj p_2$, \dots
  \item विधानीय संयोजके: \startycommalist
  \ycomma $\lnot$ (नकार)
  \ycomma $\land$ (संयोग)
  \ycomma $\lor$ (विकल्पयोग)
  \ycomma $\lif$ (शर्त)
  \ycomma $\liff$ (द्विशर्त).
  \item $a \in G$ असताना ज्ञान-संकारक $\Knows_a$.
\end{enumerate}
\end{defn}
आपल्या व्यवस्थेत आपल्याला एकाच कर्त्याच्या ज्ञानाचा
विचार करायचा असेल, तर संच~$G$ आणि स्वतंत्र कर्ते
यांचा निर्देश वगळता येतो. त्या वेळी आपल्याकडे फक्त
मूलभूत संकारक~$\Knows$ असतो.
\begin{defn}
ज्ञानविषयक भाषेतील \emph{सूत्रे} ची पुढीलप्रमाणे
  आगमनाने व्याख्या करतो:
\begin{enumerate}
\item $\lfalse$ हे अण्विक सूत्र आहे.
\item $\ltrue$ हे अण्विक सूत्र आहे.
\item प्रत्येक विधानीय चल~$\Obj p_i$ हे (अण्विक) सूत्र आहे.
\item $\varphi$ हे सूत्र असेल, तर $\lnot \varphi$
  हेही सूत्र आहे.
\item $\varphi$ आणि $\psi$ ही सूत्रे असतील, तर $(\varphi \land
  \psi)$ हे सूत्र आहे.
\item $\varphi$ आणि $\psi$ ही सूत्रे असतील, तर $(\varphi \lor \psi)$
  हे सूत्र आहे.
\item $\varphi$ आणि $\psi$ ही सूत्रे असतील, तर $(\varphi \lif \psi)$
  हे सूत्र आहे.
\item $\varphi$ आणि $\psi$ ही सूत्रे असतील, तर $(\varphi \liff \psi)$
  हे सूत्र आहे.
\item $\varphi$ हे सूत्र असेल आणि $a \in G$ असेल,
  तर $\Knows_a \varphi$ हे सूत्र आहे.
\item याखेरीज दुसरे काहीही सूत्र नाही.
\end{enumerate}
\end{defn}
सूत्र~$\varphi$ मध्ये कोणत्याही $a\in G$ साठी
$\Knows_a$ नसेल, तर त्याला \emph{मोडल-विरहित} म्हणतो.
\begin{defn}
  $\Knows$ हा संकारक वैयक्तिक ज्ञान दर्शवण्यासाठी
  आहे; $\EKnows$ हा ``सर्वांना माहीत आहे'' असा
  वाचला जाणारा संकारक समूहज्ञान दर्शवतो. रिकामेतर
  सीमित $G' \subseteq G$ साठी $\EKnows_{G'} \varphi$ हे
  पुढील सीमित संयोगाचे संक्षिप्त रूप म्हणून परिभाषित करतो:
  \[\bigwedge_{b \in G'} \Knows_b \varphi.\]
  रिकाम्या समूहासाठी हे सत्य असलेला रिकामा संयोग मानतो.
  $G'$ अनंत असेल, तर हा संयोग मूलभूत सीमित सूत्रांच्या भाषेतील
  संक्षेप नाही. अशा समूहासाठी स्वतंत्र $\EKnows_{G'}$
  संकारक जोडल्यास त्याची सत्यतेची अट अशी घ्यायची:
  $\mSat{M}{\EKnows_{G'} \varphi}[w]$ जर आणि फक्त जर
  प्रत्येक $b\in G'$ साठी $\mSat{M}{\Knows_b \varphi}[w]$.
\end{defn}

कर्त्यांच्या~$G$ समूहातील ज्ञानाची अधिक प्रबळ
संकल्पनाही विचारात घेता येते: \emph{सामायिक ज्ञान}.
यात त्या समूहातील कर्त्यांच्या प्रत्येक रिकामेतर सीमित
क्रमासाठी, एकाला माहीत आहे की दुसऱ्याला माहीत आहे की
आणखी एकाला माहीत आहे, अशा प्रकारचे ज्ञान आवश्यक असते;
अशी अट प्रत्येक सीमित लांबीसाठी घेतली जाते.
सामायिक ज्ञानातून समूहज्ञान मिळते, पण उलट दिशा सर्व
प्रतिमानांत खरी नसते. तरीही प्रत्येक समूह व चौकटीत
हा भेद काटेकोर असेलच असे नाही; उदाहरणार्थ, एकाच
कर्त्याचा परावर्ती संक्रमक संबंध असेल, तर त्याचे
वैयक्तिक आणि सामायिक ज्ञान जुळते.
``$\varphi$ हे~$G$ समूहात सामायिक ज्ञान आहे'' हे
दर्शवण्यासाठी आपण $\CKnows_G \varphi$ वापरू. सर्व सीमित
लांबींच्या अटींचा संयोग स्वतः मूलभूत भाषेतील सीमित
सूत्र नाही; म्हणून $\CKnows_G$ हा विस्तारित भाषेतील
स्वतंत्र संकारक मानतो. त्याचे संबंधी अर्थनिर्धारण
\ref{aml:el:trw:sec} मध्ये दिले आहे.













\section{संबंधी प्रतिमाने}\label{aml:el:rel:sec}

ज्ञानविषयक तर्कशास्त्रांची मूलभूत चिन्हार्थमीमांसक
संकल्पना सामान्य मोडल तर्कशास्त्रांतील संकल्पनेसारखीच
आहे. संबंधी प्रतिमानांत अजूनही जगांचा संच आणि कोणती
विधानीय चले कोणत्या जगांत ``सत्य''
मानायची हे ठरवणारे मूल्यांकन असते. आपण फक्त एकाच
कर्त्याचा विचार करत असू, तर नेहमीप्रमाणे एकच
प्राप्यता संबंध असतो. मात्र अनेक कर्त्यांच्या ज्ञानविषयक
तर्कशास्त्रात एकाच प्राप्यता संबंधाऐवजी प्राप्यता
संबंधांचा संच असतो: कर्तृचिन्हांच्या~$G$ संचातील
प्रत्येक~$a$ साठी एक संबंध.

\emph{संबंधी प्रतिमानात} जगांचा संच असतो; प्रत्येक
कर्त्यासाठी असलेल्या द्विस्थानी प्राप्यता संबंधांनी
ही जगे संबंधित असतात. त्यात कोणती विधानीय चले कोणत्या जगांत सत्य आहेत हे ठरवणारे
मूल्यांकनही असते.

\begin{defn}
  अनेक कर्त्यांच्या ज्ञानविषयक भाषेचे \emph{प्रतिमान}
  म्हणजे $\mModel{M} = \tuple{W, R, V}$ असे त्रिकूट;
  येथे
  \begin{enumerate}
  \item $W$ हा ``जगांचा'' रिकामा नसलेला संच आहे,
  \item प्रत्येक $a \in G$ साठी ${R}_a$ हा~$W$ वरील
  द्विस्थानी प्राप्यता संबंध आहे; $R=(R_a)_{a\in G}$
  हा या संबंधांचा कर्त्यांनुसार निर्देशित परिवार आहे, आणि
  \item $V$ हे प्रत्येक विधानीय चल~$p$ ला संभाव्य जगांचा संच~$V(p)\subseteq W$ नेमून
    देणारे फलन आहे.
  \end{enumerate}
  $R_a ww'$ असेल तेव्हा $w'$ हे कर्ता~$a$ साठी~$w$
  पासून \emph{प्राप्य} आहे असे म्हणतो. $w \in V(p)$
  असेल तेव्हा $p$ हे~$w$ येथे \emph{सत्य} आहे असे म्हणतो.
\end{defn}

याची कार्यपद्धती सामान्य मोडल तर्कशास्त्रासारखीच
आहे; फक्त आणखी प्राप्यता संबंध जोडलेले असतात.
दिलेल्या कर्त्याचा प्राप्यता संबंध सामान्यतः त्या
कर्त्याच्या माहितीच्या स्थितीबद्दल काहीतरी दर्शवतो,
असा अर्थ आपण लावू. उदाहरणार्थ, $R_a ww'$ हे
अनेकदा $w'$ हे जग~$w$ येथे~$a$ कडे असलेल्या
माहितीशी सुसंगत आहे, असे व्यक्त करते असे मानतो.
दुसऱ्या शब्दांत, $w$ येथे त्या कर्त्याला $w$ आणि
जग~$w'$ यांमधील फरक ओळखता येत नाही.













\section{एखाद्या जगात सत्य असणे}\label{aml:el:trw:sec}

सामान्य मोडल तर्कशास्त्राप्रमाणेच प्रत्येक ज्ञानविषयक
प्रतिमान त्यात कोणती सूत्रे कोणत्या जगांत सत्य
मानायची हे ठरवते. संबंधी चिन्हार्थमीमांसेच्या मूलभूत
संकल्पनेसाठी ``प्रतिमान~$\mModel{M}$ मध्ये
सूत्र~$\varphi$ हे जग~$w$ येथे सत्य आहे'' हीच संज्ञा
वापरतो. या संबंधाची आगमनाने व्याख्या केली जाते;
मोडल नसलेल्या सर्व संकारकांसाठी ती सामान्य मोडल
प्रकरणातील व्याख्येसारखीच असते.

\begin{defn}\label{aml:el:trw:defn:mmodels}
  प्रतिमान~$\mModel M$ मध्ये \emph{सूत्र~$\varphi$ हे~$w$ येथे
  सत्य असणे}, म्हणजे $\mSat{M}{\varphi}[w]$, पुढीलप्रमाणे
  आगमनाने परिभाषित करतो:
  \begin{enumerate}
  \item \indcase{\varphi}{\lfalse}{$\mSat{M}{\lfalse}[w]$ कधीच नसते}.
  \item \indcase{\varphi}{\ltrue}{$\mSat{M}{\ltrue}[w]$ नेहमी असते}.
  \item $\mSat{M}{p}[w]$ तेव्हा आणि केवळ तेव्हाच, जेव्हा $w \in V(p)$
  \item \indcase{\varphi}{\lnot \psi}{$\mSat{M}{\indfrm}[w]$
    तेव्हा आणि केवळ तेव्हाच, जेव्हा $\mSat/{M}{\psi}[w]$}.
  \item \indcase{\varphi}{(\psi \land \chi)}{$\mSat{M}{\indfrm}[w]$
    तेव्हा आणि केवळ तेव्हाच, जेव्हा $\mSat{M}{\psi}[w]$ आणि
    $\mSat{M}{\chi}[w]$}.
  \item \indcase{\varphi}{(\psi \lor \chi)}{$\mSat{M}{\indfrm}[w]$
    तेव्हा आणि केवळ तेव्हाच, जेव्हा $\mSat{M}{\psi}[w]$ किंवा
    $\mSat{M}{\chi}[w]$} (किंवा दोन्ही).
  \item \indcase{\varphi}{(\psi \lif \chi)}{$\mSat{M}{\indfrm}[w]$
    तेव्हा आणि केवळ तेव्हाच, जेव्हा $\mSat/{M}{\psi}[w]$ किंवा
    $\mSat{M}{\chi}[w]$}.
  \item \indcase{\varphi}{(\psi \liff \chi)}{$\mSat{M}{\indfrm}[w]$
    तेव्हा आणि केवळ तेव्हाच, जेव्हा $\mSat{M}{\psi}[w]$ आणि
    $\mSat{M}{\chi}[w]$ ही दोन्ही असतात, किंवा
    $\mSat{M}{\psi}[w]$ आणि $\mSat{M}{\chi}[w]$ यांपैकी
    कोणतेही नसते}.
  \item\label{aml:el:trw:defn:sub:mmodels-box}
    \indcase{\varphi}{\Knows_a \psi}{$\mSat{M}{\indfrm}[w]$ तेव्हा आणि
    केवळ तेव्हाच, जेव्हा $R_a ww'$ असलेल्या प्रत्येक $w' \in W$
    साठी $\mSat{M}{\psi}[w']$}
  \end{enumerate}
\end{defn}

इथे मात्र प्राप्यता संबंधांवर घालायच्या अटींचा विचार
करावा लागतो. खंड~\ref{aml:el:trw:defn:sub:mmodels-box} नुसार,
$R_a ww'$ असलेले एकही~$w'$ नसेल, तेव्हा
सूत्र~$\Knows_a \psi$ हे~$w$ येथे सत्य असते.
हा खंड सामान्य मोडल तर्कशास्त्रातील खंडासारखाच
आहे: एखाद्या जगाला उत्तरवर्ती जगे नसतील, तर सर्व
$\Box$-सूत्रे तेथे रिक्तपणे सत्य असतात. $\Knows$ हा
\emph{ज्ञान} दर्शवतो असे मानल्यास हे अत्यंत
प्रतिसहजभासी वाटते. कारण एखाद्या कर्त्याला एकाच
वेळी $\varphi$ आणि $\lnot \varphi$ ही दोन्ही माहीत असू शकतील,
अशी \emph{कोणतीही} परिस्थिती नसते, असे आपण
सहसा मानतो.

एक उपाय म्हणजे ज्ञानविषयक तर्कशास्त्रातील प्राप्यता
संबंध नेहमी \emph{परावर्ती} असेल याची खात्री करणे.
यातून वास्तविक जग कर्त्याच्या माहितीशी सुसंगत आहे,
ही कल्पना ढोबळमानाने व्यक्त होते. वस्तुतः ज्ञानविषयक
तर्कशास्त्रांत सहसा S5 वापरतात; पण $\Knows_a$ संबंधाने
नेमके काय दर्शवायचे आहे यानुसार काही जण त्याहून
दुर्बल व्यवस्था वापरू शकतात.

\begin{figure}
  \begin{center}
    \begin{tikzpicture}[modal]
      \node[world] (w1) [label={[align=right]right:\mTrue{p}\\ \mFalse{q}}]{$w_1$};
      \node[world] (w2) [label={[align=right]right:\mTrue{p}\\ \mTrue{q}},
        above right=of w1]{$w_2$};
      \node[world] (w3) [label={[align=right]right:\mFalse{p}\\ \mFalse{q}},
        below right=of w1] {$w_3$};
      \draw[<->] (w1) to node {$a$} (w2);
      \draw[->,loop left] (w1) to node {$a,b$} (w1);
      \draw[<->] (w1) to [swap] node {$b$} (w3);
      \draw[->,loop above] (w2) to node {$a,b$} (w2) ;
      \draw[->,loop below] (w3) to node {$a,b$} (w3) ;
    \end{tikzpicture}
  \end{center}
  \caption{साधे ज्ञानविषयक प्रतिमान.}
  \label{aml:el:trw:fig:simple}
\end{figure}

\begin{prob}
  \ref{aml:el:trw:fig:simple} मधील प्रतिमानात
  पुढीलपैकी कोणती विधाने सत्य आहेत ते ठरवा:
  \begin{enumerate}
    \item $\mSat{M}{\lnot q}[w_1]$;
    \item $\mSat{M}{\Knows_a \lnot q}[w_1]$;
    \item $\mSat{M}{\Knows_b \lnot q}[w_1]$;
    \item $\mSat{M}{\Knows_b q \lor \Knows_b \lnot q}[w_2]$;
    \item $\mSat{M}{\Knows_a( \Knows_b q \lor \Knows_b \lnot q)}[w_2]$;
    \item $\mSat{M}{\EKnows_{\{a,b\}} \lnot q}[w_3]$;
    \end{enumerate}
\end{prob}

जगात सत्य असण्याची मूलभूत व्याख्या दिल्यानंतर,
मोडल वैधता आणि तार्किक निष्पन्नता यांसारख्या सामान्य
मोडल तर्कशास्त्रातील इतर चिन्हार्थमीमांसक संकल्पनाही
इथे लागू होतात; फक्त मोडल संकारकांच्या अर्थलक्षणाकडे
पाहण्याचा हा नवा दृष्टिकोन वापरावा लागतो.

आता सामायिक-ज्ञान संकारक~$\CKnows_G$ साठी सत्यतेच्या
अटीही देता येतात. \ref{sfr:rel:ops:sec} मधील
व्याख्येनुसार, संबंध~$R$ चे \emph{संक्रमक संवरण}~$R^+$
पुढीलप्रमाणे असते:
\begin{align*}
  R^+ &= \bigcup_{0 < n \in \mathbb{N}} R^n,
\intertext{येथे}
  R^1 & = R \text{ आणि}\\
  R^{n+1} & = \Setabs{\tuple{x, z}}{\exists y (R^n xy \land Ryz)}\quad(n\geq1).
\end{align*}
मग $G$ हा कर्त्यांचा समूह असेल, तर सर्व कर्त्यांच्या
प्राप्यता संबंधांच्या संघाचे संक्रमक संवरण म्हणून
$R_G = ( \bigcup_{b
\in G} R_b )^+$ परिभाषित करतो.
यात एक किंवा अधिक प्राप्यता-पायऱ्या असलेले मार्ग घेतले
आहेत; शून्य पायऱ्यांचा मार्ग आपोआप जोडलेला नाही.

\begin{defn}
$G' \subseteq G$ असेल, तर $\mSat{M}{\CKnows_{G'} \varphi}[ w]$
तेव्हा आणि केवळ तेव्हाच, जेव्हा $R_{G'} w w'$ असलेल्या
प्रत्येक $w'$ साठी $\mSat{M}{\varphi}[w']$.
\end{defn}
ही अट सर्व रिकामेतर सीमित कर्तृक्रमांसाठी ज्ञानाची
अट घेण्यास समतुल्य आहे. येथे सामान्य $R_a$ संबंधांसाठी
संक्रमक संवरण घेतले आहे; परावर्ती संक्रमक संवरण नाही.
समूह रिकामेतर आणि त्यातील संबंध परावर्ती असतील, तर
$R_{G'}ww$ असल्याने सामायिक ज्ञान असलेले सूत्र~$w$
येथेही सत्य असते. रिकाम्या समूहासाठी संबंध रिकामा
आणि ही सार्विक अट रिक्तपणे सत्य असते.













\section{प्राप्यता संबंध आणि ज्ञानविषयक तत्त्वे}\label{aml:el:acc:sec}

सामान्य मोडल तर्कशास्त्रांतील चौकट-अनुरूपतेबद्दल आपल्याला
आधीच जे माहीत आहे, त्या आधारावर ज्ञानविषयक अर्थलक्षण
दिल्यास वैशिष्ट्यदर्शक सूत्रे कशी दिसतात हे
पाहता येईल. ज्ञानविषयक तर्कशास्त्रांचा अर्थ सहसा S5
मध्ये लावतात, हे आपण सांगितले आहे. आता विविध
ज्ञानविषयक तत्त्वे कशी व्यक्त होतात आणि त्यांचा
विविध चौकट-अटींशी कसा अनुरूप संबंध आहे ते पाहू.

सामान्य मोडल तर्कशास्त्रात विविध मोडल सूत्रे
प्राप्यता संबंधांचे निरनिराळे गुणधर्म वैशिष्ट्यित
करतात, हे आठवा. विशिष्ट ज्ञानविषयक तत्त्वांशी
संबंधित अशी काही सूत्रे या तक्त्यात दिली आहेत.

\begin{table}[t]
    \begin{tabular}{| p{.465\textwidth} || p{.465\textwidth} |}
      \hline
      {\emph{$R$ असा असेल \dots}} & {\emph{तर~$\mModel{M}$ मध्ये \dots सत्य असते:}} \\
      \hline \hline

      & $\Knows (p \lif q) \lif (\Knows p \lif \Knows q)$
      \hfill \newline{(संवरण)} \\
      \hline
      \emph{परावर्ती}: $\forall w Rww$
      & $\Knows p \lif p$ \hfill \newline{(ज्ञानाची सत्यता)} \\
      \hline
      \emph{संक्रमक}: \newline
      $\forall u \forall v \forall w ((Ruv \land Rvw) \lif Ruw)$ &
      $\Knows p \lif \Knows \Knows p$ \hfill
           \newline{(सकारात्मक आत्मनिरीक्षण)} \\
      \hline
      \emph{युक्लिडी}: \newline
      $\forall w \forall u \forall v ((Rwu \land Rwv) \lif Ruv)$ &
      $\lnot \Knows p \lif \Knows \neg \Knows p$ \hfill
        \newline{(नकारात्मक आत्मनिरीक्षण)} \\
      \hline
    \end{tabular}
    \caption{चार ज्ञानविषयक तत्त्वे.}
    \label{aml:el:acc:tab:four}
  \end{table}

$T$ स्वयंसिद्धाशी संबंधित ज्ञानाची सत्यता हे या
तत्त्वांपैकी सर्वात कमी विवाद्य मानतात, कारण
सूत्र माहीत असेल, तर ते सत्य असलेच पाहिजे,
असा तिचा आशय आहे. संवरण तसेच सकारात्मक आणि
नकारात्मक आत्मनिरीक्षण ही तत्त्वे त्यापेक्षा बरीच
अधिक विवाद्य आहेत.

$K$ स्वयंसिद्धाशी संबंधित संवरण म्हणजे कर्त्याचे
ज्ञान निहितार्थाखाली बंद असते, ही कल्पना. काही
प्रसंगी ती रास्त वाटू शकते. उदाहरणार्थ, मी
व्हिक्टोरियात असेन, तर व्हॅन्कुव्हर बेटावर आहे,
हे मला माहीत असू शकते. काही विलक्षण संशयवादी
परिस्थिती बाजूला ठेवल्या, तर मी व्हिक्टोरियात
आहे हे मला माहीत आहे; म्हणून मी व्हॅन्कुव्हर
बेटावर आहे हेही मला माहीत असावे. या बाबतीत
माझ्या ज्ञानाचे तार्किक संवरण बऱ्यापैकी सहज
पटण्यासारखे वाटेल. दुसरीकडे, आपल्या ज्ञानाचे
परिणाम आपण नेहमी विचारात घेत नाही; म्हणून इतर
प्रसंगी हा निष्कर्ष तितका सहज पटणार नाही.

सकारात्मक आत्मनिरीक्षणाला कधीकधी KK-तत्त्व म्हणतात:
मला एखादी गोष्ट माहीत असेल, तर ती मला माहीत आहे
हेही मला माहीत असते, असे ते सांगते. ते 4 स्वयंसिद्धाचे
ज्ञानविषयक समरूप आहे. त्याचप्रमाणे नकारात्मक
आत्मनिरीक्षण असे सांगते की एखादी गोष्ट मला
\emph{माहीत नसेल}, तर ती मला माहीत नाही हे मला
माहीत असते; ते 5 स्वयंसिद्धाचे समरूप आहे. या
तत्त्वांतील फरक प्रतिउदाहरणातही जपावा. सकारात्मक
आत्मनिरीक्षण मोडण्यासाठी एखादी गोष्ट मला माहीत असूनही
ती मला माहीत आहे हे माहीत नसणे आवश्यक आहे.
नकारात्मक आत्मनिरीक्षण मोडण्यासाठी ती गोष्ट मला
माहीत नसूनही, ती मला माहीत नाही हे माहीत नसणे
आवश्यक आहे. प्रत्यक्षात माहीत असलेल्या गोष्टीबद्दल
साशंक असणे हे पहिल्या तत्त्वासाठी संबंधित ठरते;
दुसऱ्याच्या पूर्वांगात मात्र माहीत नसणे येते.














\section{द्विअनुकरणे}\label{aml:el:bsd:sec}

आपल्या ज्ञानविषयक भाषेच्या अभिव्यक्तिक्षमतेविषयी
उरलेला एक प्रश्न प्रतिमाने आणि त्यांत सत्य असलेली
सूत्रे यांमधील संबंधाशी निगडित आहे. चौकट-अनुरूपतेच्या
निष्कर्षांवरून आपण पाहिले की काही सूत्रे एखाद्या
चौकटीत वैध असतील, तर त्या चौकटीला विशिष्ट गुणधर्म
असतात. पण उदाहरणार्थ, एखाद्या जगापासून त्याच
जगाकडे जाणारा परावर्ती बाण आणि प्रत्येक जग पुढच्या
जगाकडे जाणारी अनंत साखळी यांमध्ये भेद करता येईल
अशी आपली मोडल भाषा आहे का? म्हणजे, या दोन
जगांपैकी फक्त एकाच जगात सत्य असेल असे एखादे
सूत्र~$A$ आहे का?

संबंधी प्रतिमानांतील संबंधित जगांत मूलभूत
ज्ञानविषयक सूत्रांनी भेद करता येत नाही, ही कल्पना
द्विअनुकरण संबंध पकडतो. यासाठी दोन प्रतिमानांतील
जगांची संख्या किंवा त्यांचे ग्राफ समरूप असण्याची
अट नाही; पुढील अटी योग्य अनुरूपता सांगतात.

\begin{defn}[द्विअनुकरण]
$M_1 = \tuple{W_1, R_1, V_1}$ आणि $M_2 = \tuple{W_2, R_2, V_2}$
ही एकाच कर्तृचिन्ह-संच~$G$ व एकाच विधानीय
भाषेची दोन संबंधी प्रतिमाने असू द्या. तसेच
$\mathcal{R} \subseteq W_1 \times W_2$ हा द्विस्थानी
संबंध असू द्या. प्रत्येक $\tuple{w_1, w_2} \in
\mathcal{R}$ साठी पुढील अटी पूर्ण होत असतील, तर
$\mathcal{R}$ हे \emph{द्विअनुकरण} आहे असे म्हणतो:

\begin{enumerate}
  \item सर्व विधानीय चले~$p$ साठी
  $w_1 \in V_1(p)$ तेव्हा आणि केवळ तेव्हाच, जेव्हा
  $w_2 \in V_2(p)$.
  \item सर्व कर्ते $a \in G$ आणि सर्व जगे $v_1 \in W_1$
  यांसाठी, $R_{1_a} w_1 v_1$ असेल, तर $R_{2_a} w_2 v_2$
  आणि $\tuple{v_1, v_2} \in \mathcal{R}$ असलेले काही
  $v_2 \in W_2$ असते.
  \item सर्व कर्ते $a \in G$ आणि सर्व जगे $v_2 \in W_2$
  यांसाठी, $R_{2_a} w_2 v_2$ असेल, तर $R_{1_a} w_1 v_1$
  आणि $\tuple{v_1, v_2} \in \mathcal{R}$ असलेले काही
  $v_1 \in W_1$ असते.
\end{enumerate}

$M_1$ आणि $M_2$ यांमधील एखादे द्विअनुकरण जगे
$w_1$ आणि $w_2$ यांना जोडत असेल, तर आपण
$\tuple{M_1, w_1} \leftrightarroweq
\tuple{M_2, w_2}$ असेही लिहू शकतो. तेव्हा
$\tuple{M_1, w_1}$ आणि $\tuple{M_2, w_2}$ यांना
\emph{द्विअनुकरणसदृश} म्हणतो.
\end{defn}

द्विअनुकरण संबंधातील निरनिराळ्या अटी वेगवेगळ्या
गोष्टी सुनिश्चित करतात. पहिली अट सर्व विधानीय चले
बाबतीत एकमत असल्याचे सुनिश्चित करते;
म्हणून द्विअनुकरणसदृश जगांत तीच मोडल-विरहित सूत्रे
सत्य असतात. उरलेल्या दोन अटींना अनुक्रमे ``पुढे''
आणि ``मागे'' अशा अटी म्हणतात. संबंधित जगांच्या प्रत्येक
प्राप्य उत्तरवर्ती जगाला दुसऱ्या प्रतिमानात संबंधित
प्राप्य उत्तरवर्ती जग दोन्ही दिशांनी मिळते; उत्तरवर्ती
जगांची संख्या सारखी किंवा बाणांच्या रचना समरूप असण्याची अट नाही.

\begin{thm}
$\tuple{M_1, w_1} \leftrightarroweq \tuple{M_2, w_2}$
असेल, तर प्रत्येक सूत्र~$\varphi$ साठी
$\mSat{M_1}{\varphi}[w_1]$ तेव्हा आणि केवळ तेव्हाच,
जेव्हा $\mSat{M_2}{\varphi}[w_2]$.
\end{thm}

\begin{figure}
  \begin{center}
    \begin{tikzpicture}[modal]
      \node[world] (w1) {$w_1$};
      \node[world] (w2) [above left=of w1]{$w_2$};
      \node[world] (w3) [above right=of w1] {$w_3$};
      \draw[<->] (w1) to node {$a$} (w2);
      \draw[->,loop below] (w1) to node {$a$} (w1);
      \draw[<->] (w1) to [swap] node {$a$} (w3);
      \draw[->,loop above] (w2) to node {$a$} (w2) ;
      \draw[->,loop above] (w3) to node {$a$} (w3) ;

      \node[world](v1)[right of=w1, xshift=1.5in]{$v_1$};
      \node[world](v2)[above of=v1]{$v_2$};
      \draw[<->] (v1) to node {$a$} (v2);
      \draw[->,loop below] (v1) to node {$a$} (v1);
      \draw[->,loop above] (v2) to node {$a$} (v2) ;

      \draw[-,dotted] (w1) to (v1) ;
      \draw[-,dotted,bend left=45] (w2) to (v2) ;
      \draw[-,dotted,bend left=30] (w3) to (v2) ;
    \end{tikzpicture}
  \end{center}
  \caption{दोन द्विअनुकरणसदृश प्रतिमाने.}
  \label{aml:el:bsd:fig:bisimilar}
\end{figure}

\ref{aml:el:bsd:fig:bisimilar} मध्ये बाणांची रचना दाखवली आहे;
द्विअनुकरणासाठी मूल्यांकनांची अनुरूपताही आवश्यक आहे.
येथे प्रत्येक विधानीय चल~$p$ साठी
$w_1$ आणि $v_1$ येथे त्याचे सत्यतामूल्य सारखे,
आणि $w_2,w_3,v_2$ या तिन्ही जगांत त्याचे सत्यतामूल्य
सारखे घेऊ. इतर कर्ते असतील, तर त्यांचे संबंध या
दोन प्रतिमानांत रिकामे घेऊ. मग
$\{\tuple{w_1,v_1},\tuple{w_2,v_2},\tuple{w_3,v_2}\}$
हा संबंध द्विअनुकरण आहे. म्हणून संबंधित जगांत मूलभूत
मोडल भाषेतील सूत्रांनी भेद करता येत नाही. मात्र
$w_2$ आणि~$w_3$ यांत निरनिराळी विधानीय चले
सत्य असतील, तर ते दोन्ही~$v_2$ शी जोडणारा हाच
संबंध पहिली अट पूर्ण करत नाही.












\section{सार्वजनिक घोषणा तर्कशास्त्र}\label{aml:el:pal:sec}

गतिशील ज्ञानविषयक तर्कशास्त्रांमुळे कर्त्यांचे ज्ञान
काळानुसार किंवा नवी माहिती मिळाल्यावर कसे बदलते हे
दर्शवता येते. यांपैकी अनेक तर्कशास्त्रे ज्ञानातील बदल
माहितीपर \emph{घटना} किंवा \emph{अद्यतने} यांद्वारे
दर्शवतात. अद्यतनाचा सर्वांत मूलभूत प्रकार म्हणजे
सार्वजनिक घोषणा: एखादे सूत्र सत्य असल्याची घोषणा
होते आणि सर्व कर्ते ती एकत्रितपणे पाहतात. यासाठी
भाषेचा पुढीलप्रमाणे विस्तार करतो.

\begin{defn}
$G$ हा कर्तृचिन्हांचा संच असू द्या. सार्वजनिक घोषणा
असलेल्या अनेक कर्त्यांच्या ज्ञानविषयक तर्कशास्त्राच्या
मूलभूत भाषेत पुढील गोष्टी असतात:
\begin{enumerate}
  \item असत्यता दर्शवणारा विधानीय स्थिरांक~$\lfalse$.
  \item सत्यता दर्शवणारा विधानीय स्थिरांक~$\ltrue$.
  \item विधानीय चलांचा गणनीय अनंत संच: $\Obj
    p_0$, $\Obj p_1$, $\Obj p_2$, \dots
  \item विधानीय संयोजके: \startycommalist
  \ycomma $\lnot$ (नकार)
  \ycomma $\land$ (संयोग)
  \ycomma $\lor$ (विकल्पयोग)
  \ycomma $\lif$ (शर्त)
  \ycomma $\liff$ (द्विशर्त).
  \item $a \in G$ असताना ज्ञान-संकारक $\Knows_a$.
  \item $\psi$ हे सूत्र असताना सार्वजनिक घोषणा-संकारक $[\psi]$.
\end{enumerate}
\end{defn}

सार्वजनिक घोषणा-संकारक बॉक्स-संकारकासारखे कार्य
करतो. त्यानुसार भाषेची आगमनाधारित व्याख्या पुढीलप्रमाणे:

\begin{defn}
  ज्ञानविषयक भाषेतील \emph{सूत्रे} ची पुढीलप्रमाणे
  आगमनाने व्याख्या करतो:
  \begin{enumerate}
  \item $\lfalse$ हे अण्विक सूत्र आहे.

  \item $\ltrue$ हे अण्विक सूत्र आहे.

  \item प्रत्येक विधानीय चल~$\Obj p_i$ हे (अण्विक)
  सूत्र आहे.

  \item $\varphi$ हे सूत्र असेल, तर $\lnot \varphi$
  हेही सूत्र आहे.

  \item $\varphi$ आणि $\psi$ ही सूत्रे असतील, तर $(\varphi \land
  \psi)$ हे सूत्र आहे.

  \item $\varphi$ आणि $\psi$ ही सूत्रे असतील, तर $(\varphi \lor \psi)$
  हे सूत्र आहे.

  \item $\varphi$ आणि $\psi$ ही सूत्रे असतील, तर $(\varphi \lif \psi)$
  हे सूत्र आहे.

  \item $\varphi$ आणि $\psi$ ही सूत्रे असतील, तर $(\varphi \liff \psi)$
  हे सूत्र आहे.

  \item $\varphi$ हे सूत्र असेल आणि $a \in G$ असेल,
  तर $\Knows_a \varphi$ हे सूत्र आहे.

  \item $\varphi$ आणि $\psi$ ही सूत्रे असतील,
  तर $[\varphi] \psi$ हे सूत्र आहे.

  \item याखेरीज दुसरे काहीही सूत्र नाही.
\end{enumerate}
\end{defn}

सूत्र $[\varphi] \psi$ चा अभिप्रेत अर्थ ``$\varphi$ सत्य
असल्याची सार्वजनिक घोषणा झाल्यानंतर $\psi$~सत्य असते''
असा आहे. सार्वजनिक घोषणांच्या संदर्भात सामायिक
ज्ञानाविषयी बोलणेही कधी उपयोगी ठरते; म्हणून भाषेत
सामायिक-ज्ञान संकारक~$\CKnows_G
\varphi$ देखील समाविष्ट करता येतो.













\section{सार्वजनिक घोषणा तर्कशास्त्राचे अर्थनिर्धारण}\label{aml:el:psm:sec}

सार्वजनिक घोषणा तर्कशास्त्रांची संबंधी प्रतिमाने
ज्ञानविषयक तर्कशास्त्रांतील प्रतिमानांसारखीच असतात.
मात्र सार्वजनिक घोषणा-संकारकाचे अर्थनिर्धारण नवे आहे.

\begin{defn}\label{aml:el:psm:defn:mmodels} प्रतिमान~$\mModel M = \tuple{W, R, V}$
  मध्ये \emph{सूत्र~$\varphi$ हे~$w$ येथे सत्य असणे},
  म्हणजे $\mSat{M}{\varphi}[w]$, पुढीलप्रमाणे आगमनाने
  परिभाषित करतो:
  \begin{enumerate}
  \item \indcase{\varphi}{\lfalse}{$\mSat{M}{\lfalse}[w]$ कधीच नसते}.
  \item \indcase{\varphi}{\ltrue}{$\mSat{M}{\ltrue}[w]$ नेहमी असते}.
  \item $\mSat{M}{p}[w]$ तेव्हा आणि केवळ तेव्हाच, जेव्हा $w \in V(p)$
  \item \indcase{\varphi}{\lnot \psi}{$\mSat{M}{\indfrm}[w]$
    तेव्हा आणि केवळ तेव्हाच, जेव्हा $\mSat/{M}{\psi}[w]$}.
  \item \indcase{\varphi}{(\psi \land \chi)}{$\mSat{M}{\indfrm}[w]$
    तेव्हा आणि केवळ तेव्हाच, जेव्हा $\mSat{M}{\psi}[w]$ आणि
    $\mSat{M}{\chi}[w]$}.
  \item \indcase{\varphi}{(\psi \lor \chi)}{$\mSat{M}{\indfrm}[w]$
    तेव्हा आणि केवळ तेव्हाच, जेव्हा $\mSat{M}{\psi}[w]$ किंवा
    $\mSat{M}{\chi}[w]$} (किंवा दोन्ही).
  \item \indcase{\varphi}{(\psi \lif \chi)}{$\mSat{M}{\indfrm}[w]$
    तेव्हा आणि केवळ तेव्हाच, जेव्हा $\mSat/{M}{\psi}[w]$ किंवा
    $\mSat{M}{\chi}[w]$}.
  \item \indcase{\varphi}{(\psi \liff \chi)}{$\mSat{M}{\indfrm}[w]$
    तेव्हा आणि केवळ तेव्हाच, जेव्हा $\mSat{M}{\psi}[w]$ आणि
    $\mSat{M}{\chi}[w]$ ही दोन्ही असतात, किंवा
    $\mSat{M}{\psi}[w]$ आणि $\mSat{M}{\chi}[w]$ यांपैकी
    कोणतेही नसते}.
  \item\label{aml:el:psm:defn:sub:mmodels-box}
    \indcase{\varphi}{\Knows_a \psi}{$\mSat{M}{\indfrm}[w]$ तेव्हा आणि
    केवळ तेव्हाच, जेव्हा $R_a ww'$ असलेल्या प्रत्येक $w' \in W$
    साठी $\mSat{M}{\psi}[w']$}
  \item\label{aml:el:psm:defn:sub:mmodels-pal}
    \indcase{\varphi}{[\psi] \chi}{$\mSat{M}{\indfrm}[w]$ तेव्हा आणि
    केवळ तेव्हाच, जेव्हा एकतर $\mSat/{M}{\psi}[w]$,
    किंवा $\mSat{M}{\psi}[w]$ आणि $\mSat{M \mid \psi}{\chi}[w]$}


  येथे $\mModel M \mid \psi = \tuple{W', R', V'}$ ची
  व्याख्या पुढीलप्रमाणे आहे:
  \begin{enumerate}
    \item $W' = \Setabs{ u \in W}{\mSat{M}{\psi}[u]}$.
      म्हणजे $\mModel M \mid \psi$ मधील जगे ही~$\mModel M$
      मध्ये $\psi$ सत्य असलेली जगेच आहेत.
    \item $R'_a = R_a \cap (W' \times W')$. प्रत्येक कर्त्याचा
      प्राप्यता संबंध केवळ~$W'$ मध्ये उरलेल्या जगांपुरता
      मर्यादित केला जातो.
    \item $V'(p) = \Setabs{ u \in W'}{u \in V(p) }$.
      त्याचप्रमाणे उरलेल्या जगांतील विधानीय मूल्यांकने
      तीच राहतात. माहितीपर घटनांमुळे विधानीय चलांचे
      सत्यतामूल्य बदलत नाही, ही कल्पना त्यातून दिसते.
  \end{enumerate}
  $W'$ रिकामा असू शकतो. पण $\psi$ हे~$w$ येथे असत्य
  असेल, तर घोषणा-सूत्र पहिल्या पर्यायाने सत्य आहे आणि
  त्या प्रतिबंधित रचनेत~$w$ येथे $\chi$ तपासायचे नाही.
  $\psi$ हे~$w$ येथे सत्य असेल, तर $w\in W'$ असल्याने
  $W'$ रिक्तेतर आहे आणि प्रतिबंधित प्रतिमानातील
  पुढील मूल्यांकन नीट परिभाषित आहे.

  \end{enumerate}
\end{defn}

म्हणून सार्वजनिक घोषणा तर्कशास्त्रांचे वैशिष्ट्य असे
की, प्रतिमान~$\mModel M$ येथे सूत्र सत्य आहे
की नाही हे कधीकधी~$\mModel M$ व्यतिरिक्त दुसरे
प्रतिमान पाहिल्याशिवाय ठरवता येत नाही.

हेही लक्षात घ्या की आपल्या अर्थनिर्धारणात घोषणा-संकारक
हा $\Box$~संकारक आहे. म्हणून एखाद्या जगात
सूत्र~$\varphi$ सत्य असल्याची घोषणा करता येत
नसेल, तर $[\varphi] \psi$ तेथे आपोआप सत्य असते; जसे
अंतिम जगांत सर्व $\Box$ सूत्रे सत्य असतात.

\begin{figure}
  \begin{center}
    \begin{tikzpicture}[modal]
      \node[world] (w1) [label={below left:$p, \lnot q$}]{$w_1$};
      \node[world] (w2) [left=of w1, label={left:$\lnot p, \lnot q$}]{$w_2$};
      \node[world] (w3) [above=of w1, label={right:$p, q$}] {$w_3$};
      \draw[<->] (w1) to [swap] node {$b$} (w2);
      \draw[->,loop below] (w1) to node {$a, b$} (w1);
      \draw[<->] (w1) to [swap] node {$a$} (w3);
      \draw[->,loop above] (w2) to node {$a, b$} (w2) ;
      \draw[->,loop above] (w3) to node {$a, b$} (w3) ;

      \node[world](v1)[right of=w1, xshift=1.25in, label={right:$p, \lnot q$}]{$w'_1$};
      \node[world](v3)[above of=v1,label={right:$p,q$}]{$w'_3$};
      \draw[<->] (v1) to node {$a$} (v3);
      \draw[->,loop below] (v1) to node {$a,b$} (v1);
      \draw[->,loop above] (v3) to node {$a,b$} (v3) ;

      \node(m)[below=of w1]{$\mModel M$};
      \node(m')[below=of v1]{$\mModel M \mid p$};

      \draw[-,dotted] (w1) to node {\footnotesize $p$ ची घोषणा}(v1) ;

    \end{tikzpicture}
  \end{center}
  \caption{$p$ ची सार्वजनिक घोषणा होण्यापूर्वी आणि झाल्यानंतर.}
  \label{aml:el:psm:fig:announcement-example}
\end{figure}

सूत्राची सार्वजनिक घोषणा म्हणजे प्रतिमान
संकुचित करणे, किंवा ज्या जगांत ते सूत्र सत्य
होते तितक्याच जगांपुरते ते मर्यादित करणे, असे पाहता
येते. \ref{aml:el:psm:fig:announcement-example} मध्ये~$p$ ची
सार्वजनिक घोषणा केल्याचा परिणाम दाखवला आहे.
या प्रतिमानाचे एक वैशिष्ट्य असे की घोषणेमुळे
कर्ता~$b$ ला~$p$ माहीत होते, पण कर्ता~$a$ ला
तसे नव्याने कळत नाही (कारण~$p$ सत्य आहे हे
$a$ ला आधीच माहीत होते).

अधिक औपचारिकपणे, $\mSat{M}{\lnot \Knows_b p}[w_1]$
असते; पण $\mSat{M
\mid p}{\Knows_b p}[w'_1]$ असते. यावरून
$\mSat{M}{[p] \Knows_b
p}[w_1]$ मिळते. पण घोषणेच्या परिणामाबद्दल
याहूनही अधिक प्रबळ विधान करता येते. प्रत्यक्षात
$\mSat{M}{[p]\CKnows_{\{a,b\}} p}[w_1]$ असते.
म्हणजे $p$ ची घोषणा झाल्यावर ते \emph{सामायिक ज्ञान}
बनते.

पण हे सर्वसाधारणपणे खरे आहे का? $\varphi$ सत्य असल्याची
घोषणा केल्यावर $\varphi$ हे \emph{नेहमीच} सामायिक ज्ञान बनेल
का? आश्चर्य वाटेल, पण उत्तर नाही असे आहे. प्रत्यक्षात,
काही सूत्रे सत्य असताना त्यांची घोषणा करणे
शक्य आहे, आणि घोषणेनंतर ती सत्य राहत नाहीत.
उदाहरणार्थ, \ref{aml:el:psm:fig:announcement-example} मधील
$w_1$ येथे $p \land \lnot \Knows_b p$ ची घोषणा
केल्याचा परिणाम पाहा. येथे $p$ हे $w_1,w_3$ या
दोन्ही जगांत सत्य आहे, म्हणून $\mModel M\mid p$ मध्ये
ही दोन्ही जगे उरतात. पण $w_3$ येथे $b$ चा एकमेव
प्राप्य पर्याय $w_3$ असल्याने $\Knows_b p$ सत्य आहे;
त्यामुळे $p\land\lnot\Knows_b p$ तेथे असत्य आहे.
हे संपूर्ण सूत्र फक्त $w_1$ येथे सत्य असते. म्हणून
$\mModel M\mid(p\land\lnot\Knows_b p)$ मध्ये केवळ
$w_1$ उरते आणि दोन्ही कर्त्यांसाठी फक्त त्याचे स्वबाण
उरतात; ते $\mModel M\mid p$ सारखे प्रतिमान नाही.
या एक-जगी प्रतिमानात $p$ आणि $\Knows_b p$ दोन्ही
सत्य आहेत. म्हणून
$\mSat{M \mid (p \land \lnot \Knows_b p)}{\lnot
(p \land \lnot \Knows_b p)}[w_1]$ असते. म्हणजे
घोषणा केल्यावर असत्य होणारे हे सूत्र आहे.



\part{अंतःप्रज्ञावादी तर्कशास्त्र}\label{int:part}
\chapter{प्रस्तावना}\label{int:int:chap}










\section{रचनाशील तर्कविचार}\label{int:int:cr:sec}

मोडल संकारक किंवा द्वितीय-क्रम संख्यापक जोडून अभिजात
तर्कशास्त्राचा विस्तार करता येतो; त्याउलट
अंतःप्रज्ञावादी तर्कशास्त्र अभिजात तर्कशास्त्रावर
मर्यादा घालते, म्हणून ते ``अभिजातेतर'' आहे.
अभिजात तर्कशास्त्र अनेक अर्थांनी \emph{अरचनाशील}
आहे. एका प्रकारच्या रचनाशील गणिताला वैशिष्ट्यपूर्ण
असणारा अधिक ``रचनाशील'' तर्कविचार अंतःप्रज्ञावादी
तर्कशास्त्रात पकडायचा आहे. यामागील काही प्रेरणा
पुढील उदाहरणांतून स्पष्ट होतील.

समजा, एखादी व्यक्ती असा दावा करते की तिने एक नैसर्गिक
संख्या~$n$ शोधली आहे: $n$ सम असेल, तर रीमान
गृहीतक सत्य आहे; आणि $n$ विषम असेल, तर ते असत्य
आहे. छान बातमी!{} रीमान गृहीतक सत्य आहे की नाही
हा गणितातील मोठ्या अनिर्णीत प्रश्नांपैकी एक आहे.
या व्यक्तीने तो प्रश्न गणनेचा, म्हणजे विशिष्ट संख्या
सम आहे की नाही हे ठरवण्याचा, प्रश्न केला असे दिसते.

मग~$n$ चे ते जादुई मूल्य कोणते? ती व्यक्ती असे
सांगते: रीमान गृहीतक सत्य असेल, तर~$n$ ही नैसर्गिक
संख्या $2$ आहे; अन्यथा ती $3$ आहे.

तुम्ही रागाने तुमचे पैसे परत मागता. अभिजात दृष्टिकोनातून
वरील वर्णन खरोखरच~$n$ चे एकमेव मूल्य ठरवते;
पण तुम्हाला हवे आहे ते~$n$ चे \emph{स्पष्टपणे}
दिलेले मूल्य.

आता आणखी एक, कदाचित कमी कृत्रिम, उदाहरण पाहू.
अपरिमेय संख्येचा परिमेय घात घेऊन परिमेय उत्तर मिळू
शकते, हे माहीत आहे. उदाहरणार्थ, $\sqrt{2}^2 = 2$.
पण अपरिमेय संख्येचा \emph{अपरिमेय} घात घेतल्यावरही
परिमेय उत्तर मिळू शकते का, हे तितके स्पष्ट नाही.
पुढील प्रमेय या प्रश्नाचे होकारार्थी उत्तर देते:

\begin{thm}
$a$ आणि $b$ अशा अपरिमेय संख्या आहेत की $a^b$ परिमेय आहे.
\end{thm}

\begin{proof}
$\sqrt{2}^{\sqrt{2}}$ पाहा. ही संख्या परिमेय असेल,
तर काम झाले: $a = b = \sqrt{2}$ घेता येईल. अन्यथा
ती अपरिमेय आहे. मग
\[(\sqrt{2}^{\sqrt{2}})^{\sqrt{2}} = \sqrt{2}^{\sqrt{2} \cdot
  \sqrt{2}} = \sqrt{2}^2 = 2,
\]
आणि हे परिमेय आहे. म्हणून या वेळी $a$ म्हणून
$\sqrt{2}^{\sqrt{2}}$ आणि $b$ म्हणून~$\sqrt 2$ घ्या.
\end{proof}

ही सिद्धता वैध आहे का? बहुतेक गणितज्ञांना ती वैध
वाटते. तरीही इथे थोडे असमाधानकारक काहीतरी आहे:
विशिष्ट गुणधर्म असलेल्या वास्तव संख्यांच्या जोडीचे
अस्तित्व सिद्ध केले, पण ती जोडी \emph{कोणती} हे
सांगता आले नाही. हाच निष्कर्ष अशा रीतीनेही सिद्ध
करता येतो की $a$, $b$ ही जोडी सिद्धतेत \emph{दिली}
जाते: $a = \sqrt{3}$ आणि $b = \log_3 4$ घ्या. मग
\[a^b = \sqrt{3}^{\log_3 4} = 3^{1/2 \cdot \log_3 4} = (3^{\log_3
  4})^{1/2} = 4^{1/2}= 2,
\]
कारण $3^{\log_3 x} = x$.

पहिल्या सिद्धतेतील पावलांसारखी पावले ग्राह्य नसतील
असा तर्कविचार अंतःप्रज्ञावादी तर्कशास्त्रात पकडायचा
आहे. $\varphi(x)$ ची पूर्तता करणारा $x$ अस्तित्वात आहे
हे सिद्ध करण्यासाठी दुसऱ्या सिद्धतेप्रमाणे एखादा
विशिष्ट~$x$ आणि तो $\varphi$ ची पूर्तता करतो याची
सिद्धता द्यावी लागते. $\varphi$ किंवा $\psi$ सत्य आहे
हे सिद्ध करायचे असेल, तर या दोहोंपैकी एकाची
सिद्धता देता आली पाहिजे.

औपचारिक भाषेत सांगायचे तर, अभिजात तर्कशास्त्राची
निष्पत्ती व्यवस्था विशिष्ट प्रकारे मर्यादित
केल्यावर अंतःप्रज्ञावादी तर्कशास्त्र मिळते. गणिती
दृष्टिकोनातून या केवळ औपचारिक निगमन व्यवस्था
आहेत; पण आधी सांगितल्याप्रमाणे त्यांत एका प्रकारचा
गणिती तर्कविचार पकडायचा आहे. तो तर्कविचार
गणिताविषयीच्या विशिष्ट तात्त्विक दृष्टिकोनातून
(उदाहरणार्थ, ब्रौवर यांच्या अंतःप्रज्ञावादातून)
समर्थनीय मानता येईल; किंवा बिशप यांच्या
रचनाशीलतावादाच्या धर्तीवर अधिक ``मूर्त'' आणि
समाधानकारक गणिती तर्कविचार मानता येईल. औपचारिक
वर्णन अनौपचारिक प्रेरणा खरोखर पकडते का, यावरही
वाद करता येईल. पण आपली तात्त्विक भूमिका काहीही
असो, अंतःप्रज्ञावादी तर्कशास्त्राचा औपचारिकरीत्या
मांडलेल्या तर्कशास्त्राप्रमाणे अभ्यास करता येतो;
आणि अनेक गणिती तर्कशास्त्रज्ञांना तो अभ्यास
रसपूर्ण वाटतो.












\section{अंतःप्रज्ञावादी तर्कशास्त्राची विन्यासमीमांसा}\label{int:int:syn:sec}

अंतःप्रज्ञावादी तर्कशास्त्राचा विन्यास विधानीय
तर्कशास्त्राच्या विन्यासासारखाच आहे. अभिजात विधानीय
तर्कशास्त्रात काही संयोजकांची व्याख्या इतर
संयोजकांनी करता येते. उदाहरणार्थ, $\varphi \lif \psi$ ची
$\lnot \varphi \lor \psi$ ने, किंवा $\varphi \lor \psi$ ची
$\lnot(\lnot \varphi \land \lnot \psi)$ ने व्याख्या करता
येते. म्हणून अभिजात तर्कशास्त्राच्या मांडणीत काही
संयोजक अशा व्याख्यांची संक्षिप्त रूपे म्हणून येतात.
अंतःप्रज्ञावादी तर्कशास्त्रात मात्र दोन अपवाद सोडता
तसे होत नाही: $\lnot \varphi$ हे $\varphi \lif \lfalse$ चे
संक्षिप्त रूप म्हणून परिभाषित करता येते---आणि
सहसा तसे करतात. अर्थात मग $\lfalse$ ची स्वतःची
व्याख्या इतर संयोजकांच्या साहाय्याने केलेली नसली पाहिजे!{} तसेच
$\varphi \liff \psi$ ची अभिजात तर्कशास्त्राप्रमाणे
$(\varphi \lif \psi) \land (\psi \lif \varphi)$ अशी व्याख्या
करता येते.

विधानीय अंतःप्रज्ञावादी तर्कशास्त्रातील सूत्रे
ही \emph{विधानीय चले} आणि विधानीय
स्थिरांक~$\lfalse$ यांपासून \emph{तार्किक संयोजके}
वापरून तयार होतात. आपल्याकडे पुढील गोष्टी आहेत:

\begin{enumerate}
\item विधानीय चलांचा गणनीय अनंत संच~$\PVar$: $\Obj p_0$,
  $\Obj p_1$, \dots
  \item असत्यता दर्शवणारा विधानीय स्थिरांक~$\lfalse$.
\item तार्किक संयोजके: $\land$
  (संयोग), $\lor$ (विकल्पयोग), $\lif$ (सशर्त)
\item विरामचिन्हे: (, ), आणि स्वल्पविराम.
\end{enumerate}

\begin{defn}[सूत्र]
\label{int:int:syn:defn:formulas} विधानीय अंतःप्रज्ञावादी
तर्कशास्त्रातील \emph{सूत्रे} चा संच~$\Frm[L_0]$
पुढीलप्रमाणे विगमनाने परिभाषित करतो:
\begin{enumerate}
\item $\lfalse$ हे अण्विक सूत्र आहे.

\item प्रत्येक विधानीय चल~$\Obj p_i$ हे अण्विक
  सूत्र आहे.

\item $\varphi$ आणि $\psi$ ही सूत्रे असतील, तर $(\varphi \land
  \psi)$ हे सूत्र आहे.

\item $\varphi$ आणि $\psi$ ही सूत्रे असतील, तर $(\varphi \lor \psi)$
  हे सूत्र आहे.

\item $\varphi$ आणि $\psi$ ही सूत्रे असतील, तर $(\varphi \lif \psi)$
  हे सूत्र आहे.

\item याखेरीज दुसरे काहीही सूत्र नाही.
\end{enumerate}
\end{defn}

वर दिलेल्या मूलभूत संयोजकांखेरीज पुढील
\emph{परिभाषित} चिन्हेही वापरतो: $\lnot$ (नकार)
आणि $\liff$ (द्विशर्त). परिभाषित
संकारकांनी तयार झालेल्या सूत्रांचा अर्थ पुढीलप्रमाणे:
\begin{enumerate}
\item $\lnot \varphi$ हे $\varphi \lif \lfalse$ चे संक्षिप्त रूप आहे.

\item $\varphi \liff \psi$ हे $(\varphi \lif \psi) \land (\psi
  \lif \varphi)$ चे संक्षिप्त रूप आहे.
\end{enumerate}

$\lnot$ ला औपचारिकरीत्या संक्षिप्त रूप मानले असले,
तरी कधीकधी ते मूलभूत असल्यासारखे~$\lnot$ चे स्पष्ट नियम
आणि व्याख्यांतील खंड देऊ. याचा मुख्य उद्देश सरावाचे
प्रश्न मांडता यावेत हा आहे.












\section{ब्रौवर--हीटिंग--कोल्मोगोरोव अर्थनिर्धारण}\label{int:int:bhk:sec}

\begin{editorial}
  BHK अर्थनिर्धारण वापरून अंतःप्रज्ञावादी विधानांची
  वैधता सिद्ध करणाऱ्या सिद्धता गोंधळात टाकतात;
  त्यांचे अधिक चांगले स्पष्टीकरण द्यायला हवे.
\end{editorial}

अंतःप्रज्ञावादी संयोजकांचे एक अनौपचारिक रचनाशील
अर्थनिर्धारण आहे; त्याला बहुधा ब्रौवर--हीटिंग--कोल्मोगोरोव
अर्थनिर्धारण म्हणतात. त्यात ``रचना'' ही संकल्पना
वापरली जाते; तिला रचनाशील सिद्धता समजू शकता.
(औपचारिक निष्पत्ती व्यवस्थेतील
निष्पत्ती शी गल्लत होऊ नये म्हणून BHK
अर्थनिर्धारणात ``सिद्धता'' हा शब्द वापरत नाही.)
या सहज समजणाऱ्या संकल्पनेच्या आधारे BHK अर्थनिर्धारण
अंतःप्रज्ञावादी संयोजकांचे अर्थ स्पष्ट करते.

\begin{enumerate}
\item अण्विक विधानाची रचना कशाला म्हणायचे हे आपल्याला
  माहीत आहे असे गृहीत धरू.
\item $\varphi_1 \land \varphi_2$ ची रचना म्हणजे $\tuple{M_1, M_2}$
  ही जोडी; येथे $M_1$ ही~$\varphi_1$ ची आणि $M_2$ ही~$\varphi_2$
  ची रचना आहे.
\item $\varphi_1 \lor \varphi_2$ ची रचना म्हणजे $\tuple{s, M}$
  ही जोडी; यात $s$ हे~$1$ आणि $M$ ही~$\varphi_1$ ची रचना
  असते, किंवा $s$ हे~$2$ आणि $M$ ही~$\varphi_2$ ची रचना
  असते.
\item $\varphi \lif \psi$ ची रचना म्हणजे~$\varphi$ ची रचना
  घेऊन तिला~$\psi$ च्या रचनेत रूपांतरित करणारे फलन.
\item $\lfalse$ (असंगती) ची कोणतीही रचना नसते.
\item $\lnot \varphi$ हे $\varphi \lif \lfalse$ चे समानार्थी रूप
  म्हणून परिभाषित केले आहे. म्हणजे $\lnot \varphi$ ची रचना
  ही~$\varphi$ ची रचना घेऊन तिला~$\lfalse$ च्या रचनेत
  रूपांतरित करणारे फलन असते.
\end{enumerate}

\begin{ex}
उदाहरणार्थ $\lnot \lfalse$ घ्या. त्याची रचना हे
असे फलन असते जे~$\lfalse$ ची कोणतीही रचना इनपुट
म्हणून घेऊन~$\lfalse$ ची रचना आउटपुट म्हणून देते.
ओळख फलन~$\Id{}$ ही अशीच एक रचना आहे:
$\lfalse$ ची रचना~$M$ दिली, तर $\Id{}(M) = M$
ही~$\lfalse$ ची रचना देते.
\end{ex}

साधारणपणे $\lnot \varphi$ चा अर्थ ``$\varphi$ ची रचना अशक्य आहे'' असा होतो.

\begin{ex}
कोणत्याही विधान~$\varphi$ साठी $\varphi \lif \lnot \lnot \varphi$
सिद्ध करू. हेच सूत्र $\varphi \lif ((\varphi \lif \lfalse) \lif
\lfalse)$ आहे. याची रचना असे फलन~$f$ असली पाहिजे
की $\varphi$ ची रचना~$M$ दिल्यावर ते $(\varphi \lif \lfalse)
\lif \lfalse$ ची रचना~$f(M)$ देते. हे फलन~$f$
$(\varphi \lif \lfalse) \lif \lfalse$ ची रचना कशी
तयार करते ते पाहू: $\varphi \lif \lfalse$
ची रचना~$h$ इनपुट म्हणून घेऊन~$\lfalse$ ची रचना
आउटपुट म्हणून देणारे फलन $g$ परिभाषित करायचे आहे.
त्यासाठी $g$ ची व्याख्या अशी करा: इनपुट~$h$ हे
आधी मिळालेल्या~$\varphi$ च्या
रचनेवर~$M$ लावा. $h$ चे आउटपुट~$h(M)$ ही
$\lfalse$ ची रचना असल्याने, $M$ ही~$\varphi$ ची रचना
असेल तेव्हा $f(M)(h) = h(M)$ ही~$\lfalse$ ची रचना
असते.
\end{ex}

\begin{ex}
आता $\lnot(\varphi \land \lnot \varphi)$, म्हणजेच $(\varphi
\land (\varphi \lif \lfalse)) \lif \lfalse$, याची रचना
देऊ. ती असे फलन~$f$ आहे की $\varphi \land (\varphi \lif
\lfalse)$ ची रचना~$M$ इनपुट म्हणून दिल्यास ते
$\lfalse$ ची रचना देते. संयोग $\psi_1 \land \psi_2$
ची रचना म्हणजे $\tuple{N_1, N_2}$ ही जोडी; येथे
$N_1$ ही~$\psi_1$ ची आणि $N_2$ ही~$\psi_2$ ची
रचना आहे. $\psi_1 \land \psi_2$ च्या रचनेतून
अनुक्रमे $\psi_1$ आणि $\psi_2$ यांच्या रचना मिळवणारी
फलने $p_1$ आणि $p_2$ अशी परिभाषित करू:
\begin{align*}
  p_1(\tuple{N_1, N_2}) &= N_1\\
  p_2(\tuple{N_1, N_2}) &= N_2
\end{align*}
फलन~$f$ पुढीलप्रमाणे काम करते. प्रथम ते आपल्या
इनपुट~$M$ ला $p_1$ लावते; त्यामुळे~$\varphi$ ची रचना
मिळते. नंतर~$M$ ला $p_2$ लावते; त्यामुळे
$\varphi \lif \lfalse$ ची रचना मिळते. ही रचना म्हणजे
फलन~$p_2(M)$; त्याला~$\varphi$ ची रचना इनपुट म्हणून
दिल्यास ते~$\lfalse$ ची रचना देते. दुसऱ्या शब्दांत,
$p_2(M)$ हे $p_1(M)$ ला लावल्यास~$\lfalse$ ची
रचना मिळते. म्हणून $f(M) = p_2(M)(p_1(M))$ अशी
व्याख्या करता येते.
\end{ex}

\begin{ex}
आता $((\varphi \land \psi) \lif \chi) \lif (\varphi \lif
(\psi \lif \chi))$ ची रचना देऊ. म्हणजे, $(\varphi \land
\psi) \lif \chi$ ची रचना~$g$ घेऊन $(\varphi \lif (\psi \lif
\chi))$ ची रचना देणारे फलन $f$ हवे. रचना~$g$
ही स्वतः फलन आहे: $\varphi \land \psi$ च्या रचनांपासून
$\chi$ च्या रचनांकडे जाणारे. आणि आउटपुट~$f(g)$
हे फलन~$h_g$ आहे: $\varphi$ च्या रचनांपासून, $\psi$
च्या रचनांना~$\chi$ च्या रचनांत नेणाऱ्या फलनांकडे जाणारे.

ठीक आहे, हे गोंधळात टाकणारे आहे. इनपुट~$g$ साठी
$f$ चे आउटपुट ठरणारे फलन $h_g$ तयार करायचे आहे.
$h_g$ चा इनपुट ही~$\varphi$ ची रचना~$M$ आहे.
$h_g(M)$ चे आउटपुट हे~$\psi$ च्या रचनांपैकी~$N$
पासून~$\chi$ च्या रचनांकडे जाणारे फलन~$k_M$ असले
पाहिजे. $k_{g,M}(N) = g(\tuple{M, N})$ असे ठेवू.
$\tuple{M, N}$ ही~$\varphi \land \psi$ ची रचना आहे हे
आठवा. म्हणून $k_{g,M}$ ही $\psi \lif \chi$ ची रचना
आहे: ती~$\psi$ च्या रचना~$N$ यांना~$\chi$ च्या
रचनांत नेते. आता $h_g(M) = k_{g,M}$ असे ठेवू.
हे फलन~$\varphi$ च्या रचना~$M$ यांना $\psi
\lif \chi$ च्या रचना~$k_{g,M}$ यांत नेते. पुढे
$f(g) = h_g$ असे ठेवू. हे फलन $(\varphi \land \psi)
\lif \chi$ च्या रचना~$g$ यांना $\varphi \lif (\psi \lif \chi)$
च्या रचनांत नेते. हुश्श!{}
\end{ex}

$\varphi \lor \lnot \varphi$ या विधानाला विमध्य सिद्धांत
म्हणतात. काही विशिष्ट~$\varphi$ साठी (उदाहरणार्थ,
$\lfalse \lor \lnot \lfalse$) तो सिद्ध करता येतो;
पण सर्वसाधारणपणे नाही. कारण अंतःप्रज्ञावादी
विकल्पयोगासाठी दोन विकल्पांपैकी एकाची रचना लागते;
पण काही विधाने अशी आहेत की सध्या ती सिद्धही
करता येत नाहीत किंवा खोडताही येत नाहीत
(उदाहरणार्थ, गोल्डबाखचे अनुमान). मात्र विमध्य
सिद्धांत खोडताही येत नाही: म्हणजे $\lnot \lnot (\varphi
\lor \lnot \varphi)$ सत्य असते.

\begin{ex}
$\lnot \lnot (\varphi \lor \lnot \varphi)$ सिद्ध करण्यासाठी
$\lnot(\varphi \lor \lnot \varphi)$, म्हणजे $(\varphi
\lor (\varphi \lif \lfalse)) \lif \lfalse$, याची रचना
घेऊन तिला~$\lfalse$ च्या रचनेत रूपांतरित करणारे
फलन~$f$ हवे. दुसऱ्या शब्दांत, $g$ ही $\lnot(\varphi
\lor \lnot \varphi)$ ची रचना असेल, तर $f(g)$ ही
$\lfalse$ ची रचना ठरेल, असे फलन $f$ हवे.

$g$ ही $\lnot(\varphi \lor \lnot \varphi)$ ची रचना आहे असे
समजा; म्हणजे ते $\varphi \lor \lnot \varphi$ ची रचना घेऊन
तिला~$\lfalse$ च्या रचनेत नेणारे फलन आहे.
$\varphi \lor \lnot \varphi$ ची रचना म्हणजे $\tuple{s, M}$
ही जोडी: यात एकतर $s = 1$ आणि $M$ ही~$\varphi$ ची
रचना असते, किंवा $s = 2$ आणि $M$ ही $\lnot \varphi$
ची रचना असते. $h_1$ हे~$\varphi$ ची रचना~$M_1$ घेऊन
$\varphi \lor \lnot \varphi$ ची रचना देणारे फलन असू द्या:
ते $M_1$ ला $\tuple{1, M_1}$ कडे नेते. तसेच $h_2$
हे $\lnot \varphi$ ची रचना~$M_2$ घेऊन
$\varphi \lor \lnot \varphi$ ची रचना देणारे फलन असू द्या:
ते $M_2$ ला $\tuple{2, M_2}$ कडे नेते.

$k$ हे $\comp{h_1}{g}$ असू द्या. ते~$\varphi$ ची रचना
दिल्यावर~$\lfalse$ ची रचना देणारे फलन आहे; म्हणजे
ते~$\varphi \lif \lfalse$ किंवा $\lnot \varphi$ ची रचना
आहे. आता $l$ हे $\comp{h_2}{g}$ असू द्या. हे
फलन $\lnot \varphi$ ची रचना दिल्यावर~$\lfalse$ ची
रचना देते. $k$ ही $\lnot \varphi$ ची रचना असल्याने,
$l(k)$ ही~$\lfalse$ ची रचना आहे.

अशा प्रकारे $\lnot(\varphi \lor \lnot \varphi)$ ची
रचना~$g$ घेऊन तिला~$\lfalse$ च्या रचनेत कसे
रूपांतरित करायचे ते सांगितले. म्हणजे
$\lnot(\varphi \lor \lnot \varphi)$ ची रचना~$g$
घेऊन~$\lfalse$ ची रचना~$l(k)$ देणारे फलन $f$
ही $\lnot\lnot(\varphi \lor \lnot \varphi)$ ची रचना आहे.
\end{ex}

पाहिल्याप्रमाणे, सूत्रांची अंतःप्रज्ञावादी
वैधता दाखवण्यासाठी BHK अर्थनिर्धारण वापरणे लवकरच
अवजड आणि गोंधळात टाकणारे होते. सुदैवाने,
अंतःप्रज्ञावादी तर्कशास्त्रासाठी अधिक चांगल्या
निष्पत्ती व्यवस्था आणि अधिक नेमकी
चिन्हार्थमीमांसक अर्थलक्षणे उपलब्ध आहेत.











\section{नैसर्गिक निगमन}\label{int:int:ntd:sec}

$\FalseCl$ चे नियम वगळलेले नैसर्गिक निगमन ही
अंतःप्रज्ञावादी तर्कशास्त्रासाठीची प्रमाणित
निष्पत्ती व्यवस्था आहे. येथे ते नियम पुन्हा
देऊन BHK अर्थनिर्धारणाच्या साहाय्याने त्यामागची
प्रेरणा स्पष्ट करू. प्रत्येक नियम असा समजू शकतो
की आधारविधानांच्या रचना असतील, तर निष्कर्षाचीही
रचना आहे असे त्यावरून ठरवता येते.

नैसर्गिक निगमनातील निष्पत्त्यांमध्ये मुक्त न
केलेली गृहीतके असतात. म्हणून~$\Gamma$ या
मुक्त न केलेले गृहीतकांपासून~$\varphi$ ची
निष्पत्ती असेल, तर तिला प्रत्येक $\psi \in
\Gamma$ च्या रचना घेऊन~$\varphi$ ची रचना देणारे
फलन मानावे. कोणत्याही मुक्त न केलेले गृहीतकांशिवाय
$\varphi$ ची निष्पत्ती असेल, तर BHK अर्थनिर्धारणाच्या
अर्थाने~$\varphi$ ची रचना असते. मात्र पुढील चर्चेत गरज
नसेल तेव्हा~$\Gamma$ चा उल्लेख वगळू.

गृहीतक~$\varphi$ स्वतःच त्या मुक्त न केलेले
गृहीतक~$\varphi$ पासून~$\varphi$ ची निष्पत्ती आहे.
हे BHK अर्थनिर्धारणाशी जुळते: रचनांवरील ओळख फलन
$\varphi$ ची कोणतीही रचना घेऊन~$\varphi$ चीच रचना देते.

\subsection{संयोग}

\begin{defish}
\AxiomC{$\varphi$}
\AxiomC{$\psi$}
\RightLabel{\Intro{\land}}
\BinaryInfC{$\varphi \land \psi$}
\DisplayProof
\hfill
\begin{tabular}{l}
\AxiomC{$\varphi \land \psi$}
\RightLabel{\Elim{\land}}
\UnaryInfC{$\varphi$}
\DisplayProof
\\[3ex]
\AxiomC{$\varphi \land \psi$}
\RightLabel{\Elim{\land}}
\UnaryInfC{$\psi$}
\DisplayProof
\end{tabular}
\end{defish}
आपल्याकडे अनुक्रमे $\varphi_1$ आणि $\varphi_2$ यांच्या
रचना $N_1$, $N_2$ आहेत असे समजा. मग
$\varphi_1 \land \varphi_2$ चीही रचना आहे: जोडी
$\tuple{N_1, N_2}$.
BHK अर्थनिर्धारणानुसार $\varphi_1 \land \varphi_2$ ची
रचना म्हणजे $\tuple{N_1, N_2}$ ही जोडी. अशी जोडी
आपल्याकडे आहे असे समजा. मग प्रत्येक संयुक्त घटकाची
रचना आहे: $N_1$ ही~$\varphi_1$ ची आणि $N_2$ ही
$\varphi_2$ ची रचना आहे.
\subsection{सशर्त}
\begin{defish}
  \AxiomC{$\Discharge{\varphi}{u}$}
  \DeduceC{$\psi$}
  \DischargeRule{$\Intro{\lif}$}{u}
  \UnaryInfC{$\varphi \lif \psi$}
  \DisplayProof
\hfill
  \AxiomC{$\varphi \lif \psi$}
  \AxiomC{$\varphi$}
  \RightLabel{$\Elim{\lif}$}
  \BinaryInfC{$\psi$}
  \DisplayProof
\end{defish}
मुक्त न केलेले गृहीतक~$\varphi$ पासून~$\psi$ ची
निष्पत्ती असेल, तर~$\varphi$ च्या रचनांना~$\psi$
च्या रचनांत नेणारे फलन~$f$ असते. तेच फलन
$\varphi \lif \psi$ ची रचना आहे. म्हणून $\Intro\lif$
च्या आधारविधानाची रचना~$\varphi$ च्या रचनेवर अवलंबून
असेल, तर निष्कर्ष $\varphi
\lif \psi$ ची रचना असते.
दुसरीकडे, $\varphi$ ची रचना~$N$ आणि $\varphi \lif \psi$
ची रचना~$f$ आहेत असे समजा. $\varphi \lif \psi$ ची
रचना हे~$\varphi$ च्या रचनांना~$\psi$ च्या रचनांत नेणारे
फलन असते. म्हणून $f(N)$ ही~$\psi$ ची रचना आहे;
म्हणजे $\Elim\lif$ च्या निष्कर्षाची रचना आहे.
\subsection{विकल्पयोग}
\begin{defish}
\begin{tabular}{l}
\AxiomC{$\varphi$}
\RightLabel{\Intro{\lor}}
\UnaryInfC{$\varphi \lor \psi$}
\DisplayProof
\\[3ex]
\AxiomC{$\psi$}
\RightLabel{\Intro{\lor}}
\UnaryInfC{$\varphi \lor \psi$}
\DisplayProof
\end{tabular}
\hfill
\AxiomC{$\varphi \lor \psi$}
\AxiomC{$\Discharge{\varphi}{n}$}
\DeduceC{$\chi$}
\AxiomC{$\Discharge{\psi}{n}$}
\DeduceC{$\chi$}
\DischargeRule{\Elim{\lor}}{n}
\TrinaryInfC{$\chi$}
\DisplayProof
\end{defish}
$\varphi_i$ ची रचना~$N_i$ असेल, तर तिच्यापासून
$\varphi_1 \lor \varphi_2$ ची रचना~$\tuple{i, N_i}$ करता
येते. उलट, $\varphi_1 \lor \varphi_2$ ची रचना म्हणजे
$\tuple{i, N_i}$ ही जोडी आपल्याकडे असेल, जिथे
$N_i$ ही~$\varphi_i$ ची रचना आहे; तसेच अनुक्रमे
$\varphi_1$, $\varphi_2$ यांच्या रचनांना~$\chi$ च्या रचनांत
नेणारी फलने $f_1$, $f_2$ असतील, तर $f_i(N_i)$
ही~$\chi$ ची रचना मिळते. हाच~$\Elim\lor$ चा निष्कर्ष.
\subsection{असंगती}
\begin{defish}
  \AxiomC{$\lfalse$}
  \RightLabel{\FalseInt}
  \UnaryInfC{$\varphi$}
  \DisplayProof
\end{defish}
मुक्त न केलेले गृहीतके~$\psi_1$, \dots, $\psi_n$
यांपासून~$\lfalse$ ची निष्पत्ती असेल, तर
$\psi_1$, \dots, $\psi_n$ यांच्या रचनांना~$\lfalse$
च्या रचनेत नेणारे फलन~$f(M_1,
\dots, M_n)$ असते. $\lfalse$ ची रचना नसल्याने
$\psi_1$, \dots, $\psi_n$ या सर्वांच्या रचनाही एकत्र
असू शकत नाहीत. म्हणून $f$ बाबत पुढीलही खरे असते:
$M_1$, \dots, $M_n$ या अनुक्रमे $\psi_1$, \dots,
$\psi_n$ यांच्या रचना \emph{असतील तर}, $f(M_1, \dots,
M_n)$ ही~$\varphi$ ची रचना \emph{असते}.
\subsection{$\lnot$ साठीचे नियम}
$\lnot \varphi$ ची व्याख्या $\varphi \lif \lfalse$ अशी
असल्याने काटेकोरपणे पाहता~$\lnot$ साठी वेगळे
नियम लागत नाहीत. तरी ते द्यायचे झाले, तर असे दिसतील:
\begin{defish}
\AxiomC{$\Discharge{\varphi}{n}$}
\noLine
\DeduceC{$\lfalse$}
\DischargeRule{\Intro{\lnot}}{n}
\UnaryInfC{$\lnot \varphi$}
\DisplayProof
\hfill
\AxiomC{$\lnot \varphi$}
\AxiomC{$\varphi$}
\RightLabel{\Elim{\lnot}}
\BinaryInfC{$\lfalse$}
\DisplayProof
\end{defish}
\subsection{निष्पत्तीची उदाहरणे}
\begin{enumerate}
\item $\Proves \varphi \lif (\lnot \varphi \lif \lfalse)$, i.e., $\Proves \varphi
  \lif ((\varphi \lif \lfalse) \lif \lfalse)$
  \begin{prooftree}
    \AxiomC{$\Discharge{\varphi}{2}$}
    \AxiomC{$\Discharge{\varphi \lif \lfalse}{1}$}
    \RightLabel{\Elim\lif}
    \BinaryInfC{$\lfalse$}
    \DischargeRule{\Intro\lif}{1}
    \UnaryInfC{$(\varphi \lif \lfalse)\lif \lfalse$}
    \DischargeRule{\Intro\lif}{2}
    \UnaryInfC{$\varphi \lif (\varphi \lif \lfalse) \lif \lfalse$}
  \end{prooftree}
\item $\Proves ((\varphi \land \psi) \lif \chi) \lif (\varphi \lif (\psi \lif \chi))$
  \begin{prooftree}
    \AxiomC{$\Discharge{(\varphi \land \psi) \lif \chi}{3}$}
    \AxiomC{$\Discharge{\varphi}{2}$}
    \AxiomC{$\Discharge{\psi}{1}$}
    \RightLabel{\Intro\land}
    \BinaryInfC{$\varphi \land \psi$}
    \RightLabel{\Elim\lif}
    \BinaryInfC{$\chi$}
    \DischargeRule{\Intro\lif}{1}
    \UnaryInfC{$\psi \lif \chi$}
    \DischargeRule{\Intro\lif}{2}
    \UnaryInfC{$\varphi \lif (\psi \lif \chi)$}
    \DischargeRule{\Intro\lif}{3}
    \UnaryInfC{$((\varphi \land \psi) \lif \chi) \lif (\varphi \lif (\psi \lif \chi))$}
  \end{prooftree}
\item $\Proves \lnot(\varphi \land \lnot \varphi)$, i.e., $\Proves (\varphi \land (\varphi
  \lif \lfalse))\lif \lfalse$
  \begin{prooftree}
    \AxiomC{$\Discharge{\varphi \land (\varphi \lif \lfalse)}{1}$}
    \RightLabel{\Elim\land}
    \UnaryInfC{$\varphi \lif \lfalse$}
    \AxiomC{$\Discharge{\varphi \land (\varphi \lif \lfalse)}{1}$}
    \RightLabel{\Elim\land}
    \UnaryInfC{$\varphi$}
    \RightLabel{\Elim\lif}
    \BinaryInfC{$\lfalse$}
    \DischargeRule{\Intro\lif}{1}
    \UnaryInfC{$(\varphi \land (\varphi \lif \lfalse))\lif \lfalse$}
  \end{prooftree}
\item $\Proves \lnot\lnot(\varphi \lor \lnot \varphi)$, i.e., $\Proves ((\varphi \lor
  (\varphi \lif \lfalse))\lif \lfalse)\lif \lfalse$
  \begin{prooftree}\hskip-3em
    \AxiomC{$\Discharge{(\varphi \lor (\varphi \lif \lfalse))\lif \lfalse}{2}$}
    \AxiomC{$\Discharge{(\varphi \lor (\varphi \lif \lfalse))\lif \lfalse}{2}$}
    \AxiomC{$\Discharge{\varphi}{1}$}
    \RightLabel{\Intro\lor}
    \UnaryInfC{$\varphi \lor (\varphi \lif \lfalse)$}
    \RightLabel{\Elim\lif}
    \BinaryInfC{$\lfalse$}
    \DischargeRule{\Intro\lif}{1}
    \UnaryInfC{$\varphi \lif \lfalse$}
    \RightLabel{\Intro\lor}
    \UnaryInfC{$\varphi \lor (\varphi \lif \lfalse)$}
    \RightLabel{\Elim\lif}
    \insertBetweenHyps{\hskip -4em}
    \BinaryInfC{$\lfalse$}
    \DischargeRule{\Intro\lif}{2}
    \UnaryInfC{$((\varphi \lor (\varphi \lif \lfalse))\lif \lfalse)\lif \lfalse$}
  \end{prooftree}
\end{enumerate}
\begin{prop}
  अंतःप्रज्ञावादी तर्कशास्त्रातील नैसर्गिक निगमनात
  $\Gamma \Proves \varphi$ असेल, तर अभिजात तर्कशास्त्रातही
  $\Gamma \Proves \varphi$ असते. विशेषतः, $\varphi$ हे
  अंतःप्रज्ञावादी प्रमेय असेल, तर ते अभिजात प्रमेयही असते.
\end{prop}
\begin{proof}
  नैसर्गिक निगमनाचा प्रत्येक नियम अभिजात नैसर्गिक
  निगमनाचाही नियम आहे. म्हणून अंतःप्रज्ञावादी
  तर्कशास्त्रातील प्रत्येक निष्पत्ती ही अभिजात
  तर्कशास्त्रातीलही निष्पत्ती आहे.
\end{proof}
\begin{prob}
  पुढील सूत्रांची अंतःप्रज्ञावादी तर्कशास्त्रात
  निष्पत्त्या द्या:
  \begin{enumerate}
    \item $(\lnot \varphi \lor \psi) \lif (\varphi \lif \psi)$
    \item $\lnot\lnot\lnot \varphi \lif \lnot \varphi$
    \item $\lnot\lnot (\varphi \land \psi) \liff (\lnot\lnot \varphi\land \lnot \lnot \psi)$
    \item $\lnot (\varphi \lor \psi) \liff (\lnot \varphi \land \lnot \psi)$
    \item $(\lnot \varphi \lor \lnot \psi) \lif \lnot (\varphi \land \psi)$
    \item $\lnot\lnot ( \varphi \land \psi) \lif  (\lnot\lnot \varphi \lor
    \lnot\lnot \psi)$
  \end{enumerate}
\end{prob}
\section{स्वयंसिद्धकीय निष्पत्ती}\label{int:int:axd:sec}
अंतःप्रज्ञावादी विधानीय तर्कशास्त्रातील स्वयंसिद्धकीय
निष्पत्त्या या संकल्पनात्मकदृष्ट्या सर्वांत सोप्या
आणि ऐतिहासिकदृष्ट्या पहिल्या निष्पत्ती-पद्धती आहेत. त्यांचे कार्य अभिजात विधानीय
तर्कशास्त्रातल्याप्रमाणेच चालते.
\begin{defn}[निष्पन्नता]
जर $\Gamma$ हा $\Lang L$ मधील सूत्रांचा संच
असेल, तर $\Gamma$ पासूनची \emph{निष्पत्ती} म्हणजे
सूत्रांचा सांत अनुक्रम $\varphi_1$, \dots,~$\varphi_n$
असा की प्रत्येक $i \le n$ साठी पुढीलपैकी एक अट
पूर्ण होते:
\begin{enumerate}
\item $\varphi_i \in \Gamma$; किंवा
\item $\varphi_i$ हे स्वयंसिद्धक आहे; किंवा
\item एखाद्या $\varphi_j$ आणि $\varphi_k$ पासून, जिथे $j < i$
  आणि $k < i$, विध्यात्मक अनुमानाने $\varphi_i$ निष्पन्न होते;
  म्हणजे $\varphi_k \ident \varphi_j \lif \varphi_i$.
\end{enumerate}
\end{defn}
\begin{defn}[स्वयंसिद्धके]
अंतःप्रज्ञावादी विधानीय तर्कशास्त्रातील
\emph{स्वयंसिद्धकांचा} संच $\PAx$ हा पुढील रूपांच्या
सर्व सूत्रांचा बनलेला आहे:
\begin{align}
  & (\varphi \land \psi) \lif \varphi \label{int:int:axd:ax:land1}\\
  & (\varphi \land \psi) \lif \psi \label{int:int:axd:ax:land2}\\
  & \varphi \lif (\psi \lif (\varphi \land \psi)) \label{int:int:axd:ax:land3}\\
  & \varphi \lif (\varphi \lor \psi) \label{int:int:axd:ax:lor1}\\
  & \varphi \lif (\psi \lor \varphi) \label{int:int:axd:ax:lor2}\\
  & (\varphi \lif \chi) \lif ((\psi \lif \chi) \lif ((\varphi \lor \psi) \lif \chi)) \label{int:int:axd:ax:lor3}\\
  & \varphi \lif (\psi \lif \varphi) \label{int:int:axd:ax:lif1}\\
  & (\varphi \lif (\psi \lif \chi)) \lif ((\varphi \lif \psi) \lif (\varphi \lif \chi)) \label{int:int:axd:ax:lif2}\\
  & \lfalse \lif \varphi \label{int:int:axd:ax:lfalse1}
\end{align}
\end{defn}
\begin{defn}[निष्पन्नता]
सूत्र~$\varphi$ हे $\Gamma$ पासून \emph{निष्पन्न करता येण्याजोगे}
असते, म्हणजे $\Gamma \Proves \varphi$, जर $\Gamma$
पासून सुरू होऊन $\varphi$ वर संपणारी निष्पत्ती असेल.
\end{defn}
\begin{defn}[प्रमेये]
रिक्त संचापासून $\varphi$ ची निष्पत्ती असेल,
तर सूत्र~$\varphi$ हे \emph{प्रमेय} असते.
$\varphi$ हे प्रमेय असेल तर आपण $\Proves \varphi$ लिहितो
आणि ते प्रमेय नसेल तर $\Proves/ \varphi$ लिहितो.
\end{defn}
\begin{prop}
  अंतःप्रज्ञावादी तर्कशास्त्राच्या स्वयंसिद्धकीय
  निष्पत्ती-पद्धतीत $\Gamma \Proves \varphi$
  असेल, तर अभिजात तर्कशास्त्रातही
  $\Gamma \Proves \varphi$ असते. विशेषतः, $\varphi$ हे
  अंतःप्रज्ञावादी प्रमेय असेल, तर ते अभिजात प्रमेयही असते.
\end{prop}
\begin{proof}
  प्रत्येक अंतःप्रज्ञावादी स्वयंसिद्धक हे अभिजात
  स्वयंसिद्धकही आहे. म्हणून अंतःप्रज्ञावादी
  तर्कशास्त्रातील प्रत्येक निष्पत्ती ही अभिजात
  तर्कशास्त्रातीलही निष्पत्ती आहे.
\end{proof}
\chapter{चिन्हार्थमीमांसा}\label{int:sem:chap}
\begin{editorial}
  या प्रकरणात अंतःप्रज्ञावादी तर्कशास्त्राच्या
  चिन्हार्थमीमांसेच्या व्याख्या एकत्र दिल्या
  आहेत. सध्या केवळ क्रिप्के चिन्हार्थमीमांसा आणि
  सांस्थितिक चिन्हार्थमीमांसा यांचा समावेश आहे.
  प्रतिमानांमध्ये सूत्रे सत्य कशी ठरतात किंवा
  सूत्रांची वैधता कशी सिद्ध करतात, यांची उदाहरणे
  अद्याप दिलेली नाहीत.
\end{editorial}
\section{प्रस्तावना}\label{int:sem:int:sec}
चिन्हार्थमीमांसेशिवाय कोणत्याही तर्कशास्त्राचे
समाधानकारक वर्णन होत नाही; अंतःप्रज्ञावादी
तर्कशास्त्रही त्याला अपवाद नाही. अभिजात
तर्कशास्त्रात सत्य-मूल्यांकने वर आधारलेली
चिन्हार्थमीमांसा प्रमाणभूत मानली जाते. पण
अंतःप्रज्ञावादी तर्कशास्त्रासाठी अनेक स्पर्धक
चिन्हार्थमीमांसा आहेत. संयोजकांच्या अर्थाचे
अंतःप्रज्ञावादी दृष्टिकोनातून मान्य होईल असे
पूर्ण स्पष्टीकरण देण्याच्या बाबतीत त्यांपैकी
एकही पूर्णपणे समाधानकारक नाही.
मोडल तर्कशास्त्रातील चिन्हार्थमीमांसेसारखी,
संबंधी प्रतिमाने वर आधारलेली चिन्हार्थमीमांसा
बहुधा सर्वाधिक प्रचलित आहे. यात विधानीय चले
जगांशी संबद्ध केली जातात आणि जगांना प्राप्यता
संबंध जोडतो. हा संबंध नेहमीच अंशतः क्रम असतो;
म्हणजे तो परावर्ती, प्रतिसममित आणि संक्रमक असतो.
सहज कल्पना करण्यासाठी या जगांना ज्ञानावस्था किंवा
“पुराव्याची परिस्थिती” समजू शकता. $w'$ ही अवस्था
$w$ पासून प्राप्य असेल तेव्हा आणि केवळ तेव्हाच,
आपल्याला सध्या माहीत असलेल्या गोष्टींनुसार $w'$
ही ज्ञानाची संभाव्य (भावी) अवस्था असते; म्हणजे
$w$ मध्ये जे माहीत आहे त्याच्याशी ती सुसंगत असते.
एकदा एखादे विधान माहीत झाले की ते नंतर अज्ञात
होऊ शकत नाही: $\varphi$ हे~$w$ मध्ये माहीत असेल
आणि~$Rww'$ असेल, तर $\varphi$ हे~$w'$ मध्येही माहीत
असते. त्यामुळे प्राप्यता संबंधाच्या दृष्टीने
“ज्ञान” एकस्वनिक आहे.
येथे “माहीत आहे” याचा अर्थ रचनाशील पडताळणीच्या
अर्थाने घ्यायचा आहे: सूत्र त्या ज्ञानावस्थेत सिद्ध झालेले
किंवा पडताळलेले आहे. ही पुढील विभागातील अंतःप्रज्ञावादी
सत्यतेची अट आहे; मागील ज्ञानविषयक तर्कशास्त्रातील
साधारण $\Knows_a$ संकारकाचा तोच अर्थ नाही.
$\varphi\land \psi$ हे~$w$ मध्ये पडताळलेले असते तेव्हा
आणि केवळ तेव्हाच $\varphi$ आणि $\psi$ दोन्ही तेथे पडताळलेले
असतात. $\varphi\lor \psi$ हे~$w$ मध्ये पडताळलेले असते
तेव्हा आणि केवळ तेव्हाच त्यांपैकी किमान एक तेथे
पडताळलेले असते. म्हणून $\land$ आणि $\lor$ यांच्या
स्थानिक सत्यतेच्या अटी अभिजात तर्कशास्त्रातील अटींशी जुळतात.
परावर्तिता आणि एकस्वनिकता यांमुळे आधी पडताळलेली सूत्रे
भावी अवस्थांतही पडताळलेली राहतात; मात्र “सर्व ज्ञानविषयक
पर्यायांत सत्य” या साधारण ज्ञान-संकारकाच्या अटीवरून
विकल्पासाठी ही स्थानिक अट निष्पन्न होत नाही.
उदाहरणार्थ, \ref{aml:el:trw:fig:simple} मधील $w_1$
येथे $\Knows_a(q\lor\lnot q)$ सत्य आहे, पण
$\Knows_a q\lor\Knows_a\lnot q$ असत्य आहे:
$a$ ला प्राप्य असलेल्या $w_1$ आणि $w_2$ या अवस्थांत
$q$ चे सत्यतामूल्य भिन्न आहे.
सशर्त विधानाच्या सत्यतेच्या अटी मात्र अभिजात
तर्कशास्त्राहून वेगळ्या आहेत. $\varphi \lif \psi$
हे~$w$ मध्ये माहीत असते तेव्हा आणि केवळ तेव्हाच
$Rww'$ असलेली अशी कोणतीही~$w'$ अवस्था नसते
जिथे $\varphi$ माहीत आहे पण $\psi$ माहीत नाही.
ही अट, $w$ मध्ये $\varphi$ माहीत नाही किंवा
$\psi$ माहीत आहे, या अटीसारखी नाही. कारण
$w$ मध्ये $\varphi$ किंवा $\psi$ यांपैकी काहीही माहीत
नसेल, तरी $Rww'$ असलेली अशी भावी ज्ञानावस्था
$w'$ असू शकते; त्या~$w'$ मध्ये $\varphi$ माहीत होते
पण $\psi$ माहीत होत नाही.
ज्या कोणत्याही संभाव्य भावी ज्ञानावस्थेत~$\varphi$
माहीत होते अशी अवस्था नसेल, तरच आपल्याला
$\lnot \varphi$ माहीत असते. यामागील कल्पना अशी:
$\varphi$ जाणून घेणे शक्य असेल, तर एखाद्या
संभाव्य भावी ज्ञानावस्थेत~$\varphi$ माहीत होते.
पण~$\lfalse$ माहीत असणे शक्य नाही; म्हणून
त्या भावी ज्ञानावस्थेत आपल्याला $\varphi$ माहीत
असेल, पण~$\lfalse$ माहीत नसेल.
या अर्थलक्षणात वर्जित मध्याचे तत्त्व वैध ठरत नाही.
कारण अशी काही $\varphi$ सूत्रे आहेत की जी अद्याप
माहीत नाहीत, पण पुढे माहीत होऊ शकतात.
अशा सूत्र~$\varphi$ साठी $\varphi$ आणि~$\lnot \varphi$
दोन्ही माहीत नसतात; म्हणून $\varphi \lor \lnot \varphi$
माहीत नसते. पण उदाहरणार्थ, $\lnot(\varphi
\land \lnot \varphi)$ हे आपल्याला माहीत असते.
कारण $\varphi$ आणि~$\lnot \varphi$ दोन्ही माहीत असलेली
कोणतीही भावी अवस्था शक्य नसते; आणि हे
$\varphi$ किंवा~$\lnot \varphi$ माहीत आहे की नाही
यापासून स्वतंत्रपणे आपल्याला माहीत असते.
अंतःप्रज्ञावादी तर्कशास्त्रासाठी संबंधी प्रतिमाने
हीच एकमेव उपलब्ध चिन्हार्थमीमांसा नाही.
सांस्थितिक चिन्हार्थमीमांसा हा आणखी एक पर्याय आहे:
त्यात विधानांचे अर्थलक्षण सांस्थितिक अवकाशातील
विवृत संच म्हणून केले जाते आणि संयोजकांचे
अर्थलक्षण त्या संचांवरील क्रिया म्हणून केले जाते
(उदा., $\land$ ला छेद अनुरूप असतो).
\section{संबंधी प्रतिमान}\label{int:sem:rel:sec}
अंतःप्रज्ञावादी विधानीय तर्कशास्त्राची अचूक
चिन्हार्थमीमांसा द्यायची असेल, तर ज्याच्या
संदर्भात सूत्रांचे मूल्यांकन करता येईल
असे प्रतिमान कोणते, याची व्याख्या द्यावी लागते.
तिच्या आधारे वैधता आणि निष्पन्नता यांसारख्या
चिन्हार्थविषयक संकल्पनांचीही व्याख्या करता येते.
अशी एक चिन्हार्थमीमांसा संबंधी प्रतिमाने
देतात.
\begin{defn}
  अंतःप्रज्ञावादी विधानीय तर्कशास्त्रासाठी
  संबंधी प्रतिमान म्हणजे त्रिक
  $\mModel{M} = \tuple{W, R, V}$, जिथे
  \begin{enumerate}
  \item $W$ हा अरिक्त संच आहे,
  \item $R$ हा~$W$ वरील अंशतः क्रम आहे (म्हणजे
    परावर्ती, प्रतिसममित आणि संक्रमक द्विस्थानी संबंध), आणि
  \item $V$ हे प्रत्येक विधानीय चल~$p$ ला~$W$ चा उपसंच नेमणारे फलन आहे;
    तसेच
  \item $V$ हे $R$ च्या दृष्टीने एकस्वनिक आहे:
    $w \in V(p)$ आणि $Rww'$ असेल, तर
    $w' \in V(p)$.
  \end{enumerate}
\end{defn}
\begin{defn}\label{int:sem:rel:defn:true-at-w}
  $\varphi$ हे \emph{$\mModel{M}$ मध्ये $w$ येथे सत्य असणे},
  म्हणजे $\mSat{M}{\varphi}[w]$, ही संकल्पना आपण पुढीलप्रमाणे
  विगमनाने परिभाषित करतो:
  \begin{enumerate}
  \item \indcase{\varphi}{p}{$\mSat{M}{\indfrm}[w]$ तेव्हा
    आणि केवळ तेव्हाच $w \in V(p)$.}
  \item \indcase{\varphi}{\lfalse}{$\mSat{M}{\indfrm}[w]$ नाही}.
  \item \indcase{\varphi}{\lnot \psi}{$\mSat{M}{\indfrm}[w]$
    तेव्हा आणि केवळ तेव्हाच $Rww'$ असलेले आणि
    $\mSat{M}{\psi}[w']$ पूर्ण करणारे कोणतेही
    $w'$ नाही}.
  \item \indcase{\varphi}{\psi \land \chi}{$\mSat{M}{\indfrm}[w]$
    तेव्हा आणि केवळ तेव्हाच $\mSat{M}{\psi}[w]$
    आणि $\mSat{M}{\chi}[w]$}.
  \item \indcase{\varphi}{\psi \lor \chi}{$\mSat{M}{\indfrm}[w]$
    तेव्हा आणि केवळ तेव्हाच $\mSat{M}{\psi}[w]$
    किंवा $\mSat{M}{\chi}[w]$ (किंवा दोन्ही)}.
  \item \indcase{\varphi}{\psi \lif \chi}{$\mSat{M}{\indfrm}[w]$
    तेव्हा आणि केवळ तेव्हाच $Rww'$ असलेल्या प्रत्येक
    $w'$ साठी $\mSat{M}{\psi}[w']$ नाही किंवा
    $\mSat{M}{\chi}[w']$} (किंवा दोन्ही).
  \end{enumerate}
  $\mSat{M}{\varphi}[w]$ नसल्यास आपण
  $\mSat/{M}{\varphi}[w]$ लिहितो. $\Gamma$ हा
  सूत्रांचा संच असेल, तर $\mSat{M}{\Gamma}[w]$
  याचा अर्थ प्रत्येक $\psi \in \Gamma$ साठी
  $\mSat{M}{\psi}[w]$ असा आहे.
\end{defn}
\begin{prob}
  \ref{int:sem:rel:defn:true-at-w} नुसार
  $\mSat{M}{\lnot \varphi}[w]$ तेव्हा आणि केवळ तेव्हाच
  $\mSat{M}{\varphi \lif \lfalse}[w]$, हे दाखवा.
\end{prob}
\begin{prop}\label{int:sem:rel:prop:true-monotonic}
  जगातील सत्यता $R$ च्या दृष्टीने एकस्वनिक आहे:
  $\mSat{M}{\varphi}[w]$ आणि $Rww'$ असेल, तर
  $\mSat{M}{\varphi}[w']$.
\end{prop}
\begin{proof}
  सराव.
\end{proof}
\begin{prob}
  \ref{int:sem:rel:prop:true-monotonic} सिद्ध करा.
\end{prob}
\section{चिन्हार्थविषयक संकल्पना}\label{int:sem:sem:sec}
\begin{defn}
  प्रत्येक $w \in W$ साठी $\mSat{M}{\varphi}[w]$
  असेल तेव्हा आणि केवळ तेव्हाच $\varphi$ हे
  \emph{$\mModel{M} = \tuple{W,R,V}$ या प्रतिमानात
  सत्य आहे}, म्हणजे $\mSat{M}{\varphi}$, असे म्हणतो.
  $\varphi$ हे सर्व प्रतिमानांत सत्य असेल तेव्हा आणि
  केवळ तेव्हाच ते \emph{वैध}, म्हणजे
  $\Entails \varphi$, असते. सूत्रांच्या
  संच~$\Gamma$ पासून~$\varphi$ \emph{तार्किकरीत्या
  निष्पन्न होते}, म्हणजे $\Gamma \Entails \varphi$,
  असे तेव्हा आणि केवळ तेव्हाच म्हणतो जेव्हा
  प्रत्येक प्रतिमान~$\mModel{M}$ आणि
  $\mSat{M}{\Gamma}[w]$ असलेल्या प्रत्येक~$w$
  साठी $\mSat{M}{\varphi}[w]$ असते.
\end{defn}
\begin{prop}\label{int:sem:sem:prop:sat-entails}
  \begin{enumerate}
  \item\label{int:sem:sem:prop:sat-entails1}
    $\mSat{M}{\Gamma}[w]$ आणि
    $\Gamma \Entails \varphi$ असेल, तर
    $\mSat{M}{\varphi}[w]$.
  \item\label{int:sem:sem:prop:sat-entails2}
    $\mSat{M}{\Gamma}$ आणि $\Gamma
    \Entails \varphi$ असेल, तर $\mSat{M}{\varphi}$.
  \end{enumerate}
\end{prop}
\begin{proof}
  \begin{enumerate}
  \item $\mSat{M}{\Gamma}[w]$ आणि
    $\Gamma \Entails \varphi$ असे समजा.
    निष्पन्नतेच्या व्याख्येवरून लगेच
    $\mSat{M}{\varphi}[w]$ मिळते.
  \item $\mSat{M}{\Gamma}$ असे समजा.
    मग प्रत्येक $u \in W$ साठी
    $\mSat{M}{\Gamma}[u]$; विशेषतः
    $\mSat{M}{\Gamma}[w]$. त्यामुळे
    \ref{int:sem:sem:prop:sat-entails1} नुसार
    $\mSat{M}{\varphi}[w]$. ते जग स्वेच्छ असल्याने
    निष्कर्ष प्रतिमानात सर्वत्र सत्य आहे.
  \end{enumerate}
\end{proof}
\begin{defn}\label{int:sem:sem:defn:restrict}
  $\mModel{M}$ हे संबंधी प्रतिमान आणि $w \in W$
  असे समजा. $w$ पासून प्राप्य जगांपुरते
  $\mModel{M}$ चे \emph{मर्यादन}
  $\mModel{M}_w=\tuple{W_w, R_w, V_w}$
  पुढीलप्रमाणे मिळते:
  \begin{align*}
    W_w & = \Setabs{u \in W}{Rwu},\\
    R_w  & = R \cap (W_w)^2, \text{ आणि}\\
    V_w(p) & = V(p) \cap W_w.
  \end{align*}
\end{defn}
\begin{prop}\label{int:sem:sem:prop:restrict}
  $\mSat{M}{\varphi}[w]$ तेव्हा आणि केवळ तेव्हाच
  $\mSat{M_w}{\varphi}$.
\end{prop}
\begin{prob}
  \ref{int:sem:sem:prop:restrict} सिद्ध करा.
\end{prob}
\begin{prop}
  $\mSat{M}{\Gamma}$ असलेल्या प्रत्येक
  प्रतिमान~$\mModel{M}$ मध्ये
  $\mSat{M}{\varphi}$ असेल, तर
  $\Gamma \Entails \varphi$.
\end{prop}
\begin{proof}
  $\mSat{M}{\Gamma}[w]$ असे समजा.
  प्रत्येक~$\psi \in \Gamma$ ला
  \ref{int:sem:sem:prop:restrict} लावल्यास
  $\mSat{M_w}{\Gamma}$ मिळते. गृहीतकावरून
  $\mSat{M_w}{\varphi}$ मिळते.
  पुन्हा \ref{int:sem:sem:prop:restrict} लावल्यास
  $\mSat{M}{\varphi}[w]$ मिळते.
\end{proof}
\section{सांस्थितिक चिन्हार्थमीमांसा}\label{int:sem:top:sec}
अंतःप्रज्ञावादी तर्कशास्त्राला चिन्हार्थमीमांसा
देण्याचा आणखी एक मार्ग म्हणजे संस्थिति ही
गणिती संकल्पना वापरणे.
\begin{defn}
  $X$ हा संच असू द्या. $X$ वरील
  \emph{संस्थिति} म्हणजे खालील गुणधर्म पूर्ण
  करणारा संच $\Top{O}
  \subseteq \Pow{X}$. $\Top{O}$ च्या
  घटक ना त्या संस्थितीचे
  \emph{विवृत संच} म्हणतात. $\Top{O}$
  सहित~$X$ या संचाला \emph{सांस्थितिक अवकाश}
  म्हणतात.
  \begin{enumerate}
  \item रिकामा संच आणि संपूर्ण अवकाश विवृत आहेत:
    $\emptyset$, $X \in \Top{O}$.
  \item विवृत संच सांत छेदांखाली संवृत असतात:
    $U$, $V \in \Top{O}$ असेल, तर
    $U \cap V \in \Top{O}$.
  \item विवृत संच स्वेच्छ संयोगांखाली संवृत असतात:
    प्रत्येक $i \in I$ साठी $U_i \in
    \Top{O}$ असेल, तर $\bigcup \Setabs{U_i}{i \in
      I} \in \Top{O}$.
  \end{enumerate}
\end{defn}
विवृत संचांचा संग्रह संदर्भावरून कळत असेल,
तर सांस्थितिक अवकाश दर्शवण्यासाठी फक्त $X$
लिहू शकतो. मात्र $X$ ला विवृत संचांचा संग्रह
दिल्यानंतरच त्याला सांस्थितिक अवकाश म्हणता येते.
\begin{defn}
  अंतःप्रज्ञावादी विधानीय तर्कशास्त्राचे
  \emph{सांस्थितिक प्रतिमान} म्हणजे
  $\mModel{X} = \tuple{X, \Top{O}, V}$ हे
  त्रिक, जिथे $\Top{O}$ ही~$X$ वरील
  संस्थिति आहे आणि $V$ हे प्रत्येक विधानीय
  चलाला~$\Top{O}$ मधील विवृत संच नेमणारे
  फलन आहे.
  सांस्थितिक प्रतिमान~$\mModel{X}$ दिले असता
  $\Prop{X}{\varphi}$ ची विगमनाने पुढीलप्रमाणे
  व्याख्या करता येते:
  \begin{enumerate}
  \item $\Prop{X}{\lfalse} = \emptyset$
  \item $\Prop{X}{p} = V(p)$
  \item $\Prop{X}{\varphi \land \psi} = \Prop{X}{\varphi} \cap \Prop{X}{\psi}$
  \item $\Prop{X}{\varphi \lor \psi} = \Prop{X}{\varphi} \cup \Prop{X}{\psi}$
  \item $\Prop{X}{\varphi \lif \psi} = \Interior{(X \setminus \Prop{X}{\varphi}) \cup
    \Prop{X}{\psi}}$
  \end{enumerate}
  येथे $\Interior{V}$ हे $V \subseteq X$
  या संचाला त्याच्या \emph{अंतर्भागाकडे} नेणारे
  फलन आहे; म्हणजे त्या संचात सामावलेल्या सर्व
  विवृत संचांचा संयोग. दुसऱ्या शब्दांत,
  \[
  \Interior{V} = \bigcup \Setabs{U}{U \subseteq V \text{ आणि } U \in \Top{O}}.
  \]
\end{defn}

कोणत्याही संचाचा अंतर्भाग हा विवृत संच असतो,
कारण तो विवृत संचांचा संयोग असतो. म्हणून
$\Prop{X}{\varphi}$ हा नेहमी विवृत संच असतो.

सांस्थितिक चिन्हार्थमीमांसा अत्यंत अमूर्त
असली, तरी तिच्यामागची प्रेरणा समजून घेण्याचे
मार्ग आहेत. $X$ चे घटक, म्हणजे
“बिंदू”, ही अशी स्थळे आहेत की जिथे विधानांचे
मूल्यांकन करता येते, असे समजा. ज्या सर्व
बिंदूंवर $\varphi$ सत्य आहे त्यांचा संच म्हणजे~$\varphi$
ने व्यक्त केलेले विधान. बिंदूंचा प्रत्येक संच
संभाव्य विधान असतोच असे नाही; फक्त $\Top{O}$ चे
घटक तसे असतात. $\varphi \Entails \psi$
तेव्हा आणि केवळ तेव्हाच $\varphi$ सत्य असलेल्या
प्रत्येक बिंदूवर $\psi$ सत्य असते; म्हणजे
प्रत्येक~$X$ आणि त्यावरील प्रत्येक अनुमत मूल्यांकनासाठी $\Prop{X}{\varphi}
\subseteq \Prop{X}{\psi}$. असंगत विधान~$\lfalse$
कधीच सत्य नसते; म्हणून
$\Prop{X}{\lfalse} = \emptyset$.

अंतःप्रज्ञावादी तर्कशास्त्रात वैध असणारे
नियम टिकावेत, म्हणजे सर्व सांस्थितिक प्रतिमानांसाठी
$\varphi \Proves \psi$ तेव्हा आणि केवळ तेव्हाच
$\Prop{X}{\varphi} \subseteq \Prop{X}{\psi}$
असावे, यासाठी $\psi \land \chi$,
$\psi \lor \chi$ आणि $\psi \lif \chi$ यांनी
व्यक्त केलेली विधाने $\psi$ आणि~$\chi$
यांनी व्यक्त केलेल्या विधानांशी कशी संबंधित
असावीत? कोणत्याही $\varphi$ साठी
$\lfalse \Proves \varphi$ अपेक्षित आहे; कारण
प्रत्येक $U$ साठी $\emptyset
\subseteq U$ असल्याने ते पूर्ण होते.
$\psi \land \chi \Proves \psi$ असल्याने
$\Prop{X}{\psi \land \chi} \subseteq \Prop{X}{\psi}$
ही अट हवी; तसेच
$\Prop{X}{\psi \land \chi} \subseteq \Prop{X}{\chi}$.
$W \subseteq U$ आणि $W \subseteq V$
या अटी पूर्ण करणारा सर्वांत मोठा संच
$U \cap V$ आहे. उलट,
$\psi \Proves \psi \lor \chi$ आणि
$\chi \Proves \psi \lor \chi$, म्हणून
$\Prop{X}{\psi} \subseteq \Prop{X}{\psi \lor \chi}$
आणि $\Prop{X}{\chi} \subseteq \Prop{X}{\psi \lor \chi}$
या अटी हव्यात. $U \subseteq W$ आणि
$V \subseteq W$ असलेला सर्वांत लहान
संच~$W$ म्हणजे $U \cup V$.

$\lif$ साठीची व्याख्या अधिक सूक्ष्म आहे:
$\varphi \lif \psi$ हे, $\varphi$ शी संयोग केल्यावर
$\psi$ निष्पन्न होईल, असे सर्वांत दुर्बल
विधान व्यक्त करते. $\varphi \lif \psi$ चा
$\varphi$ शी संयोग केल्यावर~$\psi$ निष्पन्न होते,
हे $(\varphi \lif \psi)
\land \varphi \Proves \psi $ वरून स्पष्ट आहे.
म्हणून $\Prop{X}{\varphi \lif \psi}$ हा
$\Prop{X}{\varphi \lif \psi} \cap \Prop{X}{\varphi}
\subseteq \Prop{X}{\psi}$ ही अट पूर्ण करणारा
सर्वांत मोठा विवृत संच असला पाहिजे;
यावरून आपली व्याख्या मिळते.



\chapter{निर्दोषता आणि संपूर्णता}\label{int:sc:chap}
\begin{editorial}
  या प्रकरणात अंतःप्रज्ञावादी विधानीय तर्कशास्त्रासाठी
  निर्दोषता आणि संपूर्णता यांविषयीचे निष्कर्ष
  एकत्र दिले आहेत. याला प्रस्तावना हवी आहे.
  संपूर्णतेची सिद्धता सिद्धतायोग्यतेबद्दलच्या
  अशा तथ्यांचा आधार घेते की जी कुठेतरी
  स्पष्टपणे मांडून सिद्ध करायला हवीत.
\end{editorial}









\section{स्वयंसिद्धकीय निष्पत्तीची निर्दोषता}\label{int:sc:sax:sec}

\begin{editorial}
  निर्दोषतेची सिद्धता सर्व स्वयंसिद्धके
  अंतःप्रज्ञावादी अर्थाने वैध आहेत, या
  तथ्यावर अवलंबून आहे. हे अजून सिद्ध
  करायचे आहे; उदाहरणार्थ, चिन्हार्थमीमांसेच्या
  प्रकरणात ते करता येईल.
\end{editorial}

\begin{thm}[निर्दोषता]
  \label{int:sc:sax:thm:soundness}
  $\Gamma \Proves \varphi$ असेल, तर
  $ \Gamma \Entails \varphi$.
\end{thm}

\begin{proof}
  $\Gamma \Proves \varphi$ असेल, तर
  $\Gamma \Entails \varphi$, हे सिद्ध करू.
  $\Gamma$ पासून~$\varphi$ च्या
  निष्पत्तीमधील सूत्रांची
  संख्या~$n$ यावर विगमनाने सिद्धता करतो.
  $\varphi_1$, \dots, $\varphi_n = \varphi$ ही
  $\Gamma$ पासूनची निष्पत्ती असेल,
  तर $\Gamma \Entails \varphi_n$ हे दाखवू.
  $\varphi_1$, \dots, $\varphi_n$ ही
  निष्पत्ती असेल, तर कोणत्याही $k<n$
  साठी $\varphi_1$, \dots, $\varphi_k$ हीही
  निष्पत्ती आहे, हे लक्षात घ्या.

  लांबी~$0$ च्या निष्पत्त्या नसतात; म्हणून
  $n=0$ साठी दावा रिक्तपणे खरा आहे.
  त्यामुळे लांबी~$<n$ असलेल्या सर्व
  निष्पत्त्या साठी दावा खरा आहे, असे
  विगमन-गृहीतक घेऊ. आता~$\varphi_n$ चे
  समर्थन कशावर आहे त्यानुसार प्रकार पाहू.

  \begin{enumerate}
  \item $\varphi_n$ हे स्वयंसिद्धक आहे.
    सर्व स्वयंसिद्धके वैध असल्याने कोणत्याही~$\Gamma$
    साठी $\Gamma \Entails
    \varphi_n$.
  \item $\varphi_n \in \Gamma$. मग कोणतेही
    $\mModel{M}$ आणि $w$ घ्या.
    $\mSat{M}{\Gamma}[w]$ असेल, तर उघडच
    $\mSat{M}{\varphi_n}[w]$; म्हणजे
    $\Gamma \Entails \varphi$.
  \item $\varphi_n$ हे $\varphi_i$ आणि
    $\varphi_j \ident \varphi_i \lif
    \varphi_n$ यांपासून \MP{} ने निष्पन्न होते.
    $\varphi_1$, \dots, $\varphi_i$ आणि
    $\varphi_1$, \dots, $\varphi_j$ या
    $\Gamma$ पासूनच्या निष्पत्त्या आहेत.
    म्हणून विगमन-गृहीतकानुसार
    $\Gamma \Entails \varphi_i$ आणि
    $\Gamma \Entails \varphi_i \lif \varphi_n$.

    $\mSat{M}{\Gamma}[w]$ असे समजा.
    $\mSat{M}{\Gamma}[w]$ आणि
    $\Gamma \Entails \varphi_i \lif \varphi_n$
    असल्यामुळे $\mSat{M}{\varphi_i \lif
      \varphi_n}[w]$. व्याख्येनुसार याचा अर्थ
    असा की $Rww'$ असलेल्या प्रत्येक $w'$
    साठी $\mSat{M}{\varphi_i}[w']$ असेल, तर
    $\mSat{M}{\varphi_n}[w']$. $R$ परावर्ती
    असल्याने $w$ हेही $Rww'$ पूर्ण करणाऱ्या
    $w'$ पैकी एक आहे. म्हणजे
    $\mSat{M}{\varphi_i}[w]$ असेल, तर
    $\mSat{M}{\varphi_n}[w]$.
    $\Gamma \Entails \varphi_i$ असल्याने
    $\mSat{M}{\varphi_i}[w]$. म्हणून
    अपेक्षेप्रमाणे $\mSat{M}{\varphi_n}[w]$.
  \end{enumerate}
\end{proof}












\section{नैसर्गिक निगमनाची निर्दोषता}\label{int:sc:snd:sec}

आता संबंधी चिन्हार्थमीमांसेच्या दृष्टीने नैसर्गिक निगमनाची
निर्दोषता सिद्ध करू. म्हणजेच, गृहीतकांच्या संचापासून
एखादे सूत्र निष्पन्न करता येण्याजोगे असेल, तर त्या
गृहीतकांच्या संचापासून ते सूत्र
तार्किकरीत्या निष्पन्न होते, हे दाखवू.

\begin{thm}[निर्दोषता]
  \label{int:sc:snd:thm:soundness}
  $\Gamma \Proves \varphi$ असेल, तर $ \Gamma \Entails \varphi$.
\end{thm}

\begin{proof}
  $\Gamma \Proves \varphi$ असेल, तर $\Gamma \Entails \varphi$,
  हे सिद्ध करू. $\Gamma$ पासून~$\varphi$ च्या
  निष्पत्ती वर विगमनाने सिद्धता करतो.

  \begin{enumerate}
  \item निष्पत्तीमध्ये केवळ गृहीतक~$\varphi$ असेल,
    तर $\varphi \Proves \varphi$. आता $\varphi \Entails \varphi$ हे
    दाखवायचे आहे. $ \mSat{M}{\varphi}[w]$ असे समजा.
    मग उघडच $\mSat{M}{\varphi}[w]$.

  \item निष्पत्तीचा शेवट $\Intro{\land}$ ने होतो:
    आधारविधान~$\psi$ ची
      मुक्त न केलेले गृहीतके~$\Gamma$ यांपासून आणि
      आधारविधान~$\chi$ ची मुक्त न केलेले
      गृहीतके~$\Delta$ यांपासून असलेल्या
      निष्पत्त्यांवरून $\Gamma \Proves
      \psi$ आणि $\Delta \Proves \chi$ हे मिळते.
      विगमन-गृहीतकानुसार $\Gamma \Entails \psi$
      आणि $\Delta \Entails \chi$. संपूर्ण
      निष्पत्तीची मुक्त न केलेले गृहीतके
      $\Gamma$ व $\Delta$ मिळून असल्यामुळे
      $\Gamma \cup \Delta \Entails \psi \land \chi$
      हे दाखवायचे आहे. म्हणून
      $\mSat{M}{\Gamma \cup \Delta}[w]$ असे समजा.
      मग $\mSat{M}{\Gamma}[w]$ सुद्धा.
      $\Gamma \Entails \psi$ असल्याने
      $\mSat{M}{\psi}[w]$. त्याचप्रमाणे
      $\mSat{M}{\chi}[w]$. म्हणून
      $\mSat{M}{\psi \land \chi}[w]$.

  \item निष्पत्तीचा शेवट $\Elim{\land}$ ने होतो:
    आधारविधान $\psi
        \land \chi$ ची मुक्त न केलेले गृहीतके
        $\Gamma$ यांपासून असलेली निष्पत्ती
        $\Gamma \Proves \psi \land \chi$ हे दाखवते.
        विगमन-गृहीतकानुसार
        $\Gamma \Entails \psi \land \chi$.
        $\Gamma \Entails \psi$ हे दाखवायचे आहे.
        म्हणून $\mSat{M}{\Gamma}[w]$ असे समजा.
        $\Gamma \Entails \psi \land \chi$
        असल्याने $\mSat{M}{\psi \land \chi}[w]$.
        मग $\mSat{M}{\psi}[w]$ सुद्धा.
        त्याचप्रमाणे $\Elim\land$ चा शेवट~$\chi$
        ने होत असेल, तर $\Gamma \Entails \chi$.

  \item निष्पत्तीचा शेवट $\Intro{\lor}$ ने होतो:
    आधारविधान $\psi$ आहे,
      आणि~$\psi$ येथे संपणाऱ्या निष्पत्तीची
      मुक्त न केलेले गृहीतके $\Gamma$ आहेत,
      असे समजा. मग $\Gamma \Proves \psi$;
      विगमन-गृहीतकानुसार $\Gamma \Entails \psi$.
      $\Gamma \Entails \psi \lor \chi$ हे दाखवायचे आहे.
      $\mSat{M}{\Gamma}[w]$ असे समजा.
      $\Gamma \Entails \psi$ असल्याने
      $\mSat{M}{\psi}[w]$. पण मग
      $\mSat{M}{\psi \lor \chi}[w]$ सुद्धा.
      त्याचप्रमाणे आधारविधान~$\chi$ असेल, तर
      $\Gamma \Entails \chi$.

  \item निष्पत्तीचा शेवट~$\Elim{\lor}$ ने होतो:
    आधारविधानांवर संपणाऱ्या
      निष्पत्त्या अनुक्रमे $\psi \lor \chi$ ची
      मुक्त न केलेले गृहीतके~$\Gamma$,
      $\theta$ ची मुक्त न केलेले गृहीतके
      $\Delta_1 \cup \{\psi\}$, आणि $\theta$ ची
      मुक्त न केलेले गृहीतके
      $\Delta_2 \cup \{\chi\}$ यांपासून आहेत.
      म्हणून $\Gamma \Proves \psi \lor \chi$,
      $\Delta_1 \cup \{\psi\} \Proves \theta$,
      आणि $\Delta_2 \cup \{\chi\} \Proves \theta$.
      विगमन-गृहीतकानुसार
      $\Gamma \Entails \psi \lor \chi$,
      $\Delta_1 \cup \{\psi\} \Entails \theta$,
      आणि $\Delta_2 \cup \{\chi\} \Entails \theta$.
      $\Gamma \cup \Delta_1 \cup \Delta_2 \Entails \theta$
      हे सिद्ध करायचे आहे.

    $\mSat{M}{\Gamma \cup \Delta_1 \cup \Delta_2}[w]$
    असे समजा. मग $\mSat{M}{\Gamma}[w]$;
    $\Gamma \Entails \psi \lor \chi$ असल्याने
    $\mSat{M}{\psi \lor \chi}[w]$.
    $\mSat{M}{}$ च्या व्याख्येनुसार
    $\mSat{M}{\psi}[w]$ किंवा
    $\mSat{M}{\chi}[w]$. म्हणून दोन प्रकार पाहू:
    (a)~$\mSat{M}{\psi}[w]$. मग
    $\mSat{M}{\Delta_1 \cup \{\psi\}}[w]$.
    $\Delta_1 \cup \{\psi\} \Entails \theta$
    असल्याने $\mSat{M}{\theta}[w]$.
    (b)~$\mSat{M}{\chi}[w]$. मग
    $\mSat{M}{\Delta_2 \cup \{\chi\}}[w]$.
    $\Delta_2 \cup \{\chi\} \Entails \theta$
    असल्याने $\mSat{M}{\theta}[w]$.
    म्हणून दोन्ही प्रकारांत
    $\mSat{M}{\theta}[w]$, हेच दाखवायचे होते.

  \item निष्पत्तीचा शेवट $\Intro{\lif}$ ने होतो
    आणि निष्कर्ष~$\psi\lif \chi$ आहे. मग
    आधारविधान~$\chi$ आहे, आणि त्या आधारविधानावर
    संपणाऱ्या निष्पत्तीची
    मुक्त न केलेले गृहीतके~$\Gamma \cup
    \{\psi\}$ आहेत. म्हणून
    $\Gamma \cup \{\psi\} \Proves \chi$;
    विगमन-गृहीतकानुसार
    $\Gamma \cup \{\psi\} \Entails \chi$.
    $\Gamma \Entails \psi \lif \chi$ हे दाखवायचे आहे.

    $\mSat{M}{\Gamma}[w]$ असे समजा.
    $Rww'$ असे असलेल्या सर्व $w'$ साठी
    $\mSat{M}{\psi}[w']$ असेल, तर
    $\mSat{M}{\chi}[w']$, हे दाखवायचे आहे.
    म्हणून $Rww'$ आणि $\mSat{M}{\psi}[w']$
    असे गृहीत धरा.
    \ref{int:sem:rel:prop:true-monotonic} नुसार
    $\mSat{M}{\Gamma}[w']$.
    $\Gamma \cup \{\psi\} \Entails \chi$
    असल्याने $\mSat{M}{\chi}[w']$,
    हेच दाखवायचे होते.

  \item निष्पत्तीचा शेवट $\Elim{\lif}$ ने होतो
    आणि निष्कर्ष~$\chi$ आहे. आधारविधाने
    $\psi \lif \chi$ आणि $\psi$ आहेत;
    त्यांची निष्पत्त्या अनुक्रमे
    मुक्त न केलेले गृहीतके $\Gamma$ आणि
    $\Delta$ यांपासून आहेत. म्हणून
    $\Gamma \Proves \psi \lif \chi$ आणि
    $\Delta \Proves \psi$.
    विगमन-गृहीतकानुसार
    $\Gamma \Entails \psi \lif \chi$
    आणि $\Delta \Entails \psi$.
    $\Gamma \cup \Delta \Entails \chi$
    हे दाखवायचे आहे.

    $\mSat{M}{\Gamma \cup \Delta}[w]$ असे समजा.
    $\mSat{M}{\Gamma}[w]$ आणि
    $\Gamma \Entails \psi \lif \chi$ असल्याने
    $\mSat{M}{\psi \lif \chi}[w]$.
    व्याख्येनुसार, $Rww'$ असे असलेल्या
    सर्व $w'$ साठी $\mSat{M}{\psi}[w']$
    असेल, तर $\mSat{M}{\chi}[w']$.
    $R$ परावर्ती असल्याने $w$ हे
    $Rww'$ असलेल्या $w'$ पैकी एक आहे.
    म्हणजे $\mSat{M}{\psi}[w]$ असेल,
    तर $\mSat{M}{\chi}[w]$.
    $\mSat{M}{\Delta}[w]$ आणि
    $\Delta \Entails \psi$ असल्याने
    $\mSat{M}{\psi}[w]$. म्हणून
    अपेक्षेप्रमाणे $\mSat{M}{\chi}[w]$.

  \item निष्पत्तीचा शेवट $\FalseInt$ ने होतो
    आणि निष्कर्ष~$\varphi$ आहे. आधारविधान
    $\lfalse$ आहे; त्या आधारविधानाची
    निष्पत्ती ज्या
    मुक्त न केलेले गृहीतकांपासून आहे,
    ती~$\Gamma$ आहेत. मग
    $\Gamma \Proves \lfalse$.
    विगमन-गृहीतकानुसार
    $\Gamma \Entails \lfalse$.
    $\Gamma \Entails \varphi$ हे दाखवायचे आहे.

    अप्रत्यक्ष सिद्धता करू.
    $\Gamma \Entails/ \varphi$ असेल, तर
    असे प्रतिमान~$\mModel{M}$ आणि जग~$w$
    आहेत की $\mSat{M}{\Gamma}[w]$
    आणि $\mSat/{M}{\varphi}[w]$.
    $\Gamma \Entails \lfalse$
    असल्याने $\mSat{M}{\lfalse}[w]$.
    पण हे अशक्य आहे; कारण व्याख्येनुसार
    $\mSat/{M}{\lfalse}[w]$.
    म्हणून $\Gamma \Entails \varphi$.
  \item निष्पत्तीचा शेवट $\Intro\lnot$ ने होतो: सराव.
  \item निष्पत्तीचा शेवट $\Elim\lnot$ ने होतो: सराव.
  \end{enumerate}
\end{proof}

\begin{prob}
  \ref{int:sc:snd:thm:soundness} ची सिद्धता
  पूर्ण करा. $\Intro\lnot$ आणि $\Elim\lnot$
  या प्रकारांसाठी
  \ref{int:sem:rel:defn:true-at-w} मधील
  $\mSat{M}{\lnot \varphi}[w]$ ची व्याख्या वापरा;
  म्हणजे $\lnot \varphi$ हे
  $\varphi \lif \lfalse$ ने व्याख्यित आहे
  असे मानू नका.
\end{prob}

\begin{prob}
  खालील सूत्रे अंतःप्रज्ञावादी तर्कशास्त्रात
  निष्पन्न करता येण्याजोगे नाहीत, हे दाखवा:
  \begin{enumerate}
    \item $(\varphi \lif \psi) \lor (\psi \lif \varphi)$
    \item $(\lnot\lnot \varphi \lif \varphi) \lif (\varphi \lor \lnot \varphi)$
    \item $(\varphi \lif \psi \lor \chi) \lif \bigl((\varphi \lif \psi)\lor(\varphi \lif \chi)\bigr)$
  \end{enumerate}
\end{prob}






\begin{quote}\small\textbf{स्रोतनोंद.} SOL6-A200: शून्य ते एक ही विगमन-पायरी स्पष्ट केली. सांत proof support रिकामा असल्यास शून्य टप्पा घेतला; नवीन सूत्र न मिळवता स्थिर राहणारी शाखा जपली. आधीचे fairness argument बदलले नाही.\end{quote}






\section{लिंडेनबाउमचे पूर्वप्रमेय}\label{int:sc:lin:sec}

अंतःप्रज्ञावादी तर्कशास्त्राचे संपूर्णता प्रमेय सिद्ध करताना
$\Gamma \Proves/ \varphi$ असे गृहीत धरून
$\mSat{M}{\Gamma}$ आणि~$\mSat/{M}{\varphi}$
असे प्रतिमान रचले जाते.

अभिजात तर्कशास्त्रात निष्पन्नता चा संबंध
सुसंगततेच्या संकल्पनेत आणता येतो. कारण सूत्रांच्या
एखाद्या संचापासून सूत्र~$\varphi$
निष्पन्न करता येण्याजोगे असते तेव्हा आणि केवळ तेव्हाच
त्या संचात~$\varphi$ चे नकरण घातल्यावर मिळणारा
संच विसंगत असतो. अंतःप्रज्ञावादी तर्कशास्त्रात
हे शक्य नाही. अंतःप्रज्ञावादी तर्कशास्त्रात
$\lnot\varphi$ विसंगत असल्यास केवळ
$\Proves \lnot\lnot \varphi$ हे मिळते.
$\lnot\lnot\varphi \lif \varphi$ हे सर्वसाधारणपणे
अंतःप्रज्ञावादी रीतीने खरे नसल्याने
$\Proves \varphi$ असा निष्कर्ष काढता येत नाही.

म्हणून प्रतिमान~$\mModel{M}$ रचताना
सूत्र~$\varphi$ ची निष्पन्नता नसणे
याची नोंद ठेवावी लागेल. त्यामुळे अभिजात
पद्धतीप्रमाणे प्रतिमान~$\mModel{M}$
रचण्यासाठी $\Gamma^* \supseteq \Gamma$
असा संपूर्ण संच वापरता येणार नाही;
कारण प्रत्येक संपूर्ण संच~$\Gamma^*$ साठी
$\Gamma^* \Proves \varphi \lor \lnot \varphi$.

संपूर्ण संच~$\Gamma^*$ वापरण्याऐवजी
सूत्रांच्या अभाज्य संचाची संकल्पना वापरू:

\begin{defn}\label{int:sc:lin:defn:prime}
  सूत्रांचा संच~$\Gamma$
  \emph{अभाज्य} असतो तेव्हा आणि केवळ तेव्हाच
  पुढील अटी पूर्ण होतात:
  \begin{enumerate}
  \item\label{int:sc:lin:defn:prime1} $\Gamma$ सुसंगत असतो,
    म्हणजे $\Gamma \Proves/ \lfalse$;
  \item\label{int:sc:lin:defn:prime2} $\Gamma \Proves \varphi$
    असेल, तर $\varphi \in \Gamma$; आणि
  \item\label{int:sc:lin:defn:prime3} $\varphi \lor \psi \in \Gamma$
    असेल, तर $\varphi \in \Gamma$ किंवा
    $\psi \in \Gamma$.
  \end{enumerate}
\end{defn}

\begin{lem}[लिंडेनबाउमचे पूर्वप्रमेय]
  \label{int:sc:lin:lem:lindenbaum}
  $\Gamma \Proves/ \varphi$ असेल, तर
  $\Gamma^* \supseteq \Gamma$ असा
  अभाज्य~$\Gamma^*$ अस्तित्वात असतो
  आणि $\Gamma^* \Proves/ \varphi$.
\end{lem}

\begin{proof}
  $\psi \lor \chi$ या रूपातील सर्व
  सूत्रांचे एक प्रगणन
  $\psi_1 \lor \chi_1$, $\psi_2 \lor \chi_2$, \dots,
  असू द्या. सूत्रांच्या संचांची
  वाढती श्रेणी~$\Gamma_n$ परिभाषित करू;
  प्रत्येक $\Gamma_{n+1}$ हा $\Gamma_n$
  मध्ये जास्तीतजास्त एक नवे सूत्र मिळवून
  परिभाषित केला जाईल.
  $\Gamma^*$ हा सर्व~$\Gamma_n$ चा
  संयोग असेल. नवी सूत्रे अशी निवडू
  की $\Gamma^*$ अभाज्य राहील आणि
  $\Gamma^* \Proves/ \varphi$ सुद्धा राहील.
  यासाठी प्रत्येक पायरीवर पुढील अटी
  पूर्ण करणारा पहिला विकल्पयोग~
  $\psi_i \lor \chi_i$ शोधला पाहिजे:
  \begin{enumerate}
  \item\label{int:sc:lin:gamma-1}
    $\Gamma_n \Proves \psi_i \lor \chi_i$
  \item\label{int:sc:lin:gamma-2}
    $\psi_i \notin \Gamma_n$ आणि
    $\chi_i \notin \Gamma_n$
  \end{enumerate}
  $\Gamma_n \cup \{\psi_i\} \Proves/ \varphi$
  असेल, तर $\Gamma_n$ मध्ये $\psi_i$
  मिळवू; अन्यथा $\chi_i$ मिळवू.
  हे यशस्वी होते ते दाखवायचे आहे.
  सध्या \ref{int:sc:lin:gamma-1} आणि~
  \ref{int:sc:lin:gamma-2} पूर्ण करणारा सर्वात
  लहान $i$ म्हणजे $i(n)$ असे ठरवू.

  $\Gamma_0 = \Gamma$ असे ठेवून
  \[  \Gamma_{n+1} = \begin{cases}
    \Gamma_n \cup \{\psi_{i(n)}\} &
    \text{जर $\Gamma_n \cup \{\psi_{i(n)}\} \Proves/ \varphi$ असेल} \\
    \Gamma_n \cup \{\chi_{i(n)}\} & \text{अन्यथा}
    \end{cases}
  \]
  अशी व्याख्या करू. $i(n)$ अपरिभाषित
  असेल, म्हणजे $\Gamma_n \Proves \psi \lor \chi$
  असताना नेहमी $\psi \in \Gamma_n$
  किंवा $\chi \in \Gamma_n$ असेल, तर
  $\Gamma_{n+1} = \Gamma_n$ ठेवू.
  आता $\Gamma^* = \bigcup_{n=0}^\infty \Gamma_n$
  असे ठेवू.

  प्रथम प्रत्येक $n$ साठी
  $\Gamma_n \Proves/ \varphi$ हे दाखवू.
  $n$ वर विगमन करू. $n = 0$ साठी
  हा दावा प्रमेयाच्या गृहीतकामुळे,
  म्हणजे $\Gamma \Proves/ \varphi$ मुळे,
  खरा आहे. कोणत्याही $n\geq0$ साठी
  $\Gamma_n \Proves/ \varphi$ असल्यास
  $\Gamma_{n+1} \Proves/ \varphi$
  हे दाखवायचे आहे.
  $i(n)$ अपरिभाषित असेल, तर
  $\Gamma_{n+1} = \Gamma_n$;
  काहीच सिद्ध करावे लागत नाही.
  म्हणून $i(n)$ परिभाषित आहे असे
  समजा. सोयीसाठी $i = i(n)$ असे लिहू.

  दाव्याचे व्यत्यस्त रूप सिद्ध करू.
  $\Gamma_{n+1} \Proves \varphi$ असे समजा.
  रचनेनुसार $\Gamma_n \cup
  \{\psi_i\} \Proves/ \varphi$ असेल, तर
  $\Gamma_{n+1} = \Gamma_n \cup
  \{\psi_i\}$; अन्यथा
  $\Gamma_{n+1} = \Gamma_n \cup \{\chi_i\}$.
  पहिला प्रकार असूच शकत नाही;
  कारण तेव्हा $\Gamma_{n+1} \Proves/ \varphi$.
  म्हणून $\Gamma_n \cup
  \{\psi_i\} \Proves \varphi$ आणि
  $\Gamma_{n+1} = \Gamma_n \cup \{\chi_i\}$.
  $i(n)$ च्या व्याख्येनुसार
  $\Gamma _n \Proves \psi_i \lor \chi_i$.
  $\Gamma_n \cup \{\psi_i\} \Proves \varphi$
  आहे. तसेच
  $\Gamma_{n+1} = \Gamma_n \cup
  \{\chi_i\} \Proves \varphi$ आहे.
  त्यामुळे $\Gamma_n \Proves \varphi$,
  हेच दाखवायचे होते.

  $\Gamma^* \Proves \varphi$ असेल, तर
  $\Gamma' \subseteq \Gamma^*$
  असा काही परिमित उपसंच असेल की
  $\Gamma' \Proves \varphi$.
  प्रत्येक $\theta \in \Gamma'$ कोणत्या
  तरी~$i$ साठी $\Gamma_i$ मध्ये असले
  पाहिजे. $\Gamma'$ रिकामा असेल, तर $n=0$ घ्या;
  अन्यथा या निर्देशांकांपैकी सर्वात मोठा $n$ घ्या. $i \le n$ असेल, तर
  $\Gamma_i \subseteq \Gamma_n$;
  म्हणून $\Gamma' \subseteq \Gamma_n$.
  पण मग $\Gamma_n \Proves \varphi$,
  आणि वर सिद्ध केलेल्या
  $\Gamma_n \Proves/ \varphi$ शी हा विरोधाभास.

  शेवटी $\Gamma^*$ अभाज्य आहे,
  म्हणजे \ref{int:sc:lin:defn:prime} मधील
  \ref{int:sc:lin:defn:prime1},
  \ref{int:sc:lin:defn:prime2} आणि~
  \ref{int:sc:lin:defn:prime3} या अटी
  पूर्ण करतो, हे दाखवू.

  प्रथम $\Gamma^* \Proves/ \varphi$;
  म्हणून $\Gamma^*$ सुसंगत आहे,
  आणि \ref{int:sc:lin:defn:prime1} पूर्ण होते.

  आता $\Gamma^* \Proves \psi \lor \chi$
  असेल, तर $\psi \in \Gamma^*$
  किंवा $\chi \in \Gamma^*$,
  हे दाखवू. यावरून
  \ref{int:sc:lin:defn:prime3} सिद्ध होते;
  कारण $\psi \lor \chi \in \Gamma^*$
  असेल, तर $\Gamma^* \Proves \psi
  \lor \chi$ सुद्धा. म्हणून
  $\Gamma^* \Proves \psi \lor \chi$
  पण $\psi \notin \Gamma^*$
  आणि $\chi \notin \Gamma^*$
  असे समजा.
  $\Gamma^* \Proves \psi \lor \chi$
  असल्याने कोणत्या तरी~$n$ साठी
  $\Gamma_n \Proves \psi \lor \chi$.
  $\psi \lor \chi$ हा विकल्पयोग
  प्रगणनात, समजा $\psi_j \lor \chi_j$
  म्हणून, येतो.
  $\psi_j \lor \chi_j$ हा
  $i(n)$ च्या व्याख्येतील गुणधर्म
  पूर्ण करतो: $\Gamma_n \Proves \psi_j
  \lor \chi_j$, तर
  $\psi_j \notin \Gamma_n$ आणि
  $\chi_j \notin \Gamma_n$.
  त्या निर्देशांकापेक्षा आधीचे निर्देशांक
  परिमित संख्येने आहेत. हा विकल्पयोग
  निवडला जात नाही तोवर अटी पूर्ण करत
  राहतो; तोवर निवडलेला विकल्पयोग
  $\psi_i \lor \chi_i$ त्याच्या आधीचाच असतो.
  त्यातील $\psi_i$ किंवा $\chi_i$
  मिळवल्यावर तो पुन्हा अटी पूर्ण करत नाही.
  म्हणून काही टप्प्यावर~$m$
  असे मिळेल की $j = i(m)$.
  पण तेव्हा $\psi \in \Gamma_{m+1}$
  किंवा $\chi \in \Gamma_{m+1}$;
  आणि हे $\psi \notin \Gamma^*$
  व $\chi \notin \Gamma^*$ या
  गृहीतकाच्या विरुद्ध आहे.

  आता $\Gamma^* \Proves \psi$
  असे समजा. मग
  $\Gamma^* \Proves \psi \lor \psi$.
  पण $\Gamma^* \Proves \psi \lor \psi$
  असेल, तर $\psi \in \Gamma^*$,
  हे आत्ताच सिद्ध केले आहे.
  त्यामुळे $\Gamma^*$ हा
  \ref{int:sc:lin:defn:prime} मधील
  \ref{int:sc:lin:defn:prime2} पूर्ण करतो.
\end{proof}

\begin{prob}
$\Gamma \Proves/ \bot$ असेल, तर
अभिजात तर्कशास्त्रात $\Gamma$ सुसंगत
असतो, म्हणजे~$\Gamma$ मधील सर्व सूत्रांना
सत्य करणारे मूल्यांकन अस्तित्वात असते,
हे दाखवा.
\end{prob}






\begin{quote}\small\textbf{स्रोतनोंद.} SOL6-A201: मूळ one-based pair list आणि शून्यासहित natural-number worlds यांतील विसंगती दुरुस्त केली. सर्व नैसर्गिक निर्देशांकांसाठी formula-pair enumeration शून्यापासून सुरू केली; शाखा आणि truth-lemma साठी लागणाऱ्या गुणधर्मांना बदल नाही.\end{quote}






\section{कॅनॉनिकल प्रतिमान}\label{int:sc:mod:sec}

आपल्या प्रतिमानातील जगे नैसर्गिक संख्यांच्या
सांत क्रमिका~$\sigma$ असतील, म्हणजे
$\sigma \in \Nat^*$. $\Nat^*$ ची
विगमनाने पुढीलप्रमाणे व्याख्या होते:
\begin{enumerate}
\item $\emptyseq \in \Nat^*$.
\item $\sigma \in \Nat^*$ आणि $n \in \Nat$
  असतील, तर $\sigma.n \in
  \Nat^*$ (येथे $\sigma.n$ म्हणजे
  $\sigma \concat \tuple{n}$, आणि
  $\sigma \concat \sigma'$ म्हणजे
  $\sigma$ व $\sigma'$ यांची जोडणी).
\item इतर काहीही $ \Nat^*$ मध्ये नाही.
\end{enumerate}
म्हणून $\Nat^*$ वापरून विगमनात्मक
व्याख्या देता येतात.

सूत्रांच्या सर्व जोड्यांचे
$\tuple{\psi_0, \chi_0}$,
$\tuple{\psi_1, \chi_1}$,
$\tuple{\psi_2, \chi_2}$, \dots,
असे $n\in\Nat=\{0,1,2,\dots\}$ ने निर्देशांकित
प्रगणन असू द्या. सूत्रे
चा संच~$\Delta$ दिला असता
$\Delta(\sigma)$ ची विगमनाने
पुढीलप्रमाणे व्याख्या करा:
\begin{enumerate}
\item $\Delta(\emptyseq) = \Delta$
\item $\Delta(\sigma.n) = {}$
  \[  \begin{cases}
    (\Delta(\sigma) \cup \{\psi_n\})^* &
    \text{जर $\Delta(\sigma) \cup \{\psi_n\} \Proves/ \chi_n$ असेल} \\
    \Delta(\sigma) & \text{अन्यथा}
  \end{cases}
  \]
\end{enumerate}
येथे $(\Delta(\sigma) \cup \{\psi_n\})^*$
म्हणजे $\Delta(\sigma) \cup \{\psi_n\}$
या संचाला आणि सूत्र~$\chi_n$
ला \ref{int:sc:lin:lem:lindenbaum} लावल्यावर
मिळणारा सूत्रांचा अभाज्य संच.
या व्याख्येनुसार
$\Delta(\sigma) \cup \{\psi_n\} \Proves/ \chi_n$
असेल, तर $\Delta(\sigma.n)
\Proves \psi_n$ आणि
$\Delta(\sigma.n) \Proves/ \chi_n$,
हे लक्षात घ्या. तसेच कोणत्याही~$n$
साठी $\Delta(\sigma) \subseteq \Delta(\sigma.n)$.
$\Delta$ अभाज्य असेल, तर सर्व~$\sigma$
साठी $\Delta(\sigma)$ अभाज्य असतो.

\begin{defn}\label{int:sc:mod:defn:canonical-model}
  $\Delta$ अभाज्य आहे असे समजा.
  मग $\Delta$ साठीचे
  \emph{कॅनॉनिकल प्रतिमान}
  $\mModel{M(\Delta)}$ पुढीलप्रमाणे परिभाषित होते:
  \begin{enumerate}
  \item $W = \Nat^*$, म्हणजे नैसर्गिक
    संख्यांच्या सांत क्रमिकांचा संच.
  \item $R$ हा असा अंशतः क्रम आहे की
    $R\sigma\sigma'$ तेव्हा आणि केवळ
    तेव्हाच जेव्हा $\sigma$ हा~$\sigma'$
    चा आरंभीचा खंड असतो (म्हणजे
    काही क्रमिका~$\sigma''$ साठी
    $\sigma' = \sigma \concat \sigma''$).
  \item $V(p) = \Setabs{\sigma}{p \in \Delta(\sigma)}$.
  \end{enumerate}
\end{defn}

$R$ हा खरोखर अंशतः क्रम आहे,
हे सहज तपासता येते. तसेच~$V$ वरील
एकस्वनिकतेची अट पूर्ण होते.
$\Delta(\sigma) \subseteq \Delta(\sigma.n)$
असल्याने $R\sigma\sigma'$ असेल तेव्हा
$\Delta(\sigma) \subseteq \Delta(\sigma')$;
हे~$\sigma$ पासूनच्या विस्ताराच्या
लांबीवर विगमनाने मिळते.












\section{सत्यता पूर्वप्रमेय}\label{int:sc:tru:sec}

पुढील पूर्वप्रमेय कॅनॉनिकल प्रतिमानातील
पूर्तीचा संबंध अभाज्य संच
$\Delta(\sigma)$ मध्ये कोणती
सूत्रे घटक आहेत
याच्याशी जोडते.

\begin{lem}\label{int:sc:tru:lem:truth}
  $\Delta$ अभाज्य असेल, तर
  $\mSat{M(\Delta)}{\varphi}[\sigma]$
  तेव्हा आणि केवळ तेव्हाच
  $\Delta(\sigma) \Proves \varphi$.
\end{lem}

\begin{proof}
  $\varphi$ वर विगमन करू.
  \begin{enumerate}
    \item \indcase{\varphi}{\lfalse}{$\Delta(\sigma)$
      अभाज्य असल्याने सुसंगत आहे;
      म्हणून $\Delta(\sigma) \Proves/ \indfrm$.
      व्याख्येनुसार
      $\mSat/{M(\Delta)}{\indfrm}[\sigma]$.}
    \item \indcase{\varphi}{p}{$\mSat{{}}{}$ च्या
        व्याख्येनुसार
        $\mSat{M(\Delta)}{\indfrm}[\sigma]$
        तेव्हा आणि केवळ तेव्हाच
        $\sigma \in V(p)$;
        म्हणजे $\Delta(\sigma) \Proves \indfrm$;
        कारण हा अभाज्य संच निष्पत्तीखाली संवृत आहे.}
    \item \indcase!{\varphi}{\lnot \psi}{}
    \item \indcase{\varphi}{\psi \land
        \chi}{$\mSat{M(\Delta)}{\indfrm}[\sigma]$
        तेव्हा आणि केवळ तेव्हाच
        $\mSat{M(\Delta)}{\psi}[\sigma]$
        आणि $\mSat{M(\Delta)}{\chi}[\sigma]$.
        विगमन-गृहीतकानुसार
        $\mSat{M(\Delta)}{\psi}[\sigma]$
        तेव्हा आणि केवळ तेव्हाच
        $\Delta(\sigma) \Proves
        \psi$; आणि~$\chi$ साठीही तसेच.
        पण $\Delta(\sigma) \Proves \psi$
        आणि $\Delta(\sigma) \Proves \chi$
        हे तेव्हा आणि केवळ तेव्हाच,
        जेव्हा $\Delta(\sigma)
        \Proves \indfrm$.}
    \item \indcase{\varphi}{\psi \lor
        \chi}{$\mSat{M(\Delta)}{\indfrm}[\sigma]$
        तेव्हा आणि केवळ तेव्हाच
        $\mSat{M(\Delta)}{\psi}[\sigma]$
        किंवा $\mSat{M(\Delta)}{\chi}[\sigma]$.
        विगमन-गृहीतकानुसार हे
        तेव्हा आणि केवळ तेव्हाच,
        जेव्हा $\Delta(\sigma) \Proves \psi$
        किंवा $\Delta(\sigma) \Proves \chi$.
        आता हे पुन्हा
        $\Delta(\sigma) \Proves \indfrm$
        च्या तुल्य आहे, हे दाखवायचे आहे.
        डावीकडून उजवीकडील दिशा उघड आहे.
        उजवीकडून डावीकडील दिशा
        $\Delta(\sigma)$ अभाज्य
        असल्याने मिळते.}
    \item \indcase{\varphi}{\psi \lif \chi}{आधी
        डावीकडून उजवीकडील दिशेचे
        व्यत्यस्त रूप पाहू:
        $\Delta(\sigma) \Proves/ \psi
        \lif \chi$ असे समजा.
        मग $\Delta(\sigma) \cup \{\psi\} \Proves/
        \chi$ सुद्धा.
        कोणत्या तरी~$n$ साठी
        $\tuple{\psi, \chi}$ ही
        $\tuple{\psi_n, \chi_n}$ असल्याने
        $\Delta(\sigma.n) = (\Delta(\sigma) \cup
        \{\psi\})^*$; आणि
        $\Delta(\sigma.n) \Proves \psi$
        पण $\Delta(\sigma.n) \Proves/
        \chi$. विगमन-गृहीतकानुसार
        $\mSat{M(\Delta)}{\psi}[\sigma.n]$
        आणि $\mSat/{M(\Delta)}{\chi}[\sigma.n]$.
        $R\sigma(\sigma.n)$ असल्यामुळे
        याचा अर्थ
        $\mSat/{M(\Delta)}{\indfrm}[\sigma]$.

        आता $\Delta(\sigma) \Proves \psi \lif \chi$
        असे समजा आणि $R\sigma\sigma'$ असू द्या.
        $\Delta(\sigma) \subseteq
        \Delta(\sigma')$ असल्याने
        $\Delta(\sigma') \Proves
        \psi$ असेल, तर
        $\Delta(\sigma') \Proves \chi$.
        दुसऱ्या शब्दांत, $R\sigma\sigma'$
        असलेल्या प्रत्येक $\sigma'$ साठी
        $\Delta(\sigma') \Proves/ \psi$
        किंवा $\Delta(\sigma') \Proves
        \chi$. विगमन-गृहीतकानुसार
        $R\sigma\sigma'$ असेल तेव्हा
        $\mSat/{M(\Delta)}{\psi}[\sigma']$
        किंवा $\mSat{M(\Delta)}{\chi}[\sigma']$;
        म्हणजे
        $\mSat{M(\Delta)}{\indfrm}[\sigma]$.}
  \end{enumerate}
\end{proof}












\section{संपूर्णता प्रमेय}\label{int:sc:cpl:sec}

\begin{thm}\label{int:sc:cpl:thm:completeness}
  $\Gamma \Entails \varphi$ असेल, तर
  $\Gamma \Proves \varphi$.
\end{thm}

\begin{proof}
  व्यत्यस्त रूप सिद्ध करू:
  $\Gamma \Proves/ \varphi$ असे समजा.
  मग \ref{int:sc:lin:lem:lindenbaum} नुसार
  $\Gamma^* \supseteq \Gamma$ असा
  अभाज्य संच आहे की
  $\Gamma^* \Proves/ \varphi$.
  \ref{int:sc:mod:defn:canonical-model}
  मध्ये व्याख्यित केलेले $\Gamma^*$ साठीचे
  कॅनॉनिकल प्रतिमान~$\mModel{M(\Gamma^*)}$
  घ्या. प्रत्येक $\psi \in \Gamma$
  साठी $\Gamma^* \Proves \psi$.
  $\Gamma^*(\emptyseq) = \Gamma^*$,
  हे लक्षात घ्या. सत्यता पूर्वप्रमेयानुसार
  (\ref{int:sc:tru:lem:truth})
  सर्व $\psi \in \Gamma$ साठी
  $\mSat{M(\Gamma^*)}{\psi}[\emptyseq]$
  आणि
  $\mSat/{M(\Gamma^*)}{\varphi}[\emptyseq]$.
  यावरून $\Gamma \Entails/ \varphi$.
\end{proof}

\begin{prob}
  $\varphi$ मध्ये फक्त विधानीय चले,
  $\lor$ आणि~$\land$ असतील, तर
  $\Entails/ \varphi$ हे दाखवा.
  यावरून अंतःप्रज्ञावादी तर्कशास्त्रात
  $\lor$ आणि~$\land$ वापरून
  $\lif$ ची व्याख्या करता येत नाही,
  असा निष्कर्ष काढा.
\end{prob}

\begin{prob}
  संपूर्णता प्रमेय वापरून
  $\Proves \varphi \lor \psi$ असेल,
  तर $\Proves \varphi$ किंवा
  $\Proves \psi$, हे सिद्ध करा.
  (सूचना: $\mSat/{M_1}{\varphi}$
  आणि $\mSat/{M_2}{\psi}$
  असे गृहीत धरून
  $\mSat/{M}{\varphi \lor \psi}$
  असेल असे नवे प्रतिमान
  $\mModel{M}$ रचा.)
\end{prob}

\begin{prob}
  $\mModel{M}$ हे रेषीय क्रम
  वापरणारे संबंधी प्रतिमान असेल,
  तर $\mSat{M}{(\varphi \lif \psi)\lor(\psi \lif \varphi)}$,
  हे दाखवा.
\end{prob}






\begin{quote}\small\textbf{स्रोतनोंद.} SOL6-A203: मूळ atom-only quotient चुकीचा असल्याची OLINC-226 सूचना योग्य होती; तरी त्याच अवैध truth-equivalence ची सिद्धता exercise मध्ये मागितली होती. आता सर्व उपसूत्रांवरील योग्य गाळणी, implication आणि primitive negation सहित truth-preservation proof, finite bound आणि bounded decision procedure दिली. Exercise योग्य गाळणीच्या तपशिलांसाठी केला.\end{quote}






\section{निर्णेयता}\label{int:sc:dec:sec}
संपूर्णता प्रमेयाची सिद्धता
$\Gamma \Proves/ \varphi$ असलेल्या
प्रत्येक वेळी $\Gamma \Entails/ \varphi$
याचा साक्षीदार म्हणून अनंत जगे
असलेले प्रतिमान देते, हे लक्षात घ्या.
पुढील विधानानुसार
$\Entails \varphi$ सिद्ध करण्यासाठी
सर्व सांत प्रतिमानांमध्ये (म्हणजे
जगांचा सांत संच असलेल्या प्रतिमानांमध्ये)
$\mSat{M}{\varphi}$ हे सिद्ध करणे पुरेसे आहे.

\begin{thm}\label{int:sc:dec:thm:decidability}
  $\Entails/ \varphi$ असेल, तर
  $\mSat/{M'}{\varphi}$ असे सांत प्रतिमान
  अस्तित्वात असते.
\end{thm}

\begin{proof}
  $\mModel{M}=\tuple{W,R,V}$ हे $\mSat/{M}{\varphi}$
  असे प्रतिमान असू द्या. $S$ हा $\varphi$ च्या सर्व उपसूत्रांचा
  सांत संच असू द्या; विशेषतः $\varphi\in S$ आणि $S$
  उपसूत्रे घेण्याखाली संवृत आहे. $P$ हा $\varphi$ मध्ये
  येणाऱ्या विधानीय चलांचा संच आहे.
  प्रत्येक $w\in W$ साठी
  \[    [w]=\Setabs{\psi\in S}{\mSat{M}{\psi}[w]}
  \]
  असे ठरवून $\mModel{M'}=\tuple{W',R',V'}$ परिभाषित करू:
  \[    W'=\Setabs{[w]}{w\in W},\qquad
    R'[w][u]\ \text{तेव्हा आणि केवळ तेव्हाच}\ [w]\subseteq[u],
  \]
  आणि $p\in P$ साठी $V'(p)=\Setabs{[w]\in W'}{p\in[w]}$;
  इतर विधानचरांसाठी $V'(p)=\emptyset$.
  $W'$ अरिक्त असून $|W'|\leq 2^{|S|}$.
  उपसंच-संबंध हा अंशतः क्रम आहे, आणि $V'$ ची
  एकस्वनिकता व्याख्येवरून मिळते. म्हणून $\mModel{M'}$
  हे सांत संबंधी प्रतिमान आहे.

\begin{quote}\small\textbf{स्रोतनोंद.} OLINC-226: मूळ atom-only quotient उपसूत्रांची सत्यता जपत नव्हता. खाली त्याच्या जागी सर्व उपसूत्रांवरील योग्य गाळणीची सिद्धता आहे; समतुल्यता S मधील सूत्रांसाठीच सांगितली आहे.\end{quote}
  प्रत्येक $\psi\in S$ साठी रचनेवरील विगमनाने
  \[    \mSat{M}{\psi}[w]\ \text{तेव्हा आणि केवळ तेव्हाच}\
    \mSat{M'}{\psi}[{[w]}]
  \]
  हे सिद्ध करू. विधानचरांसाठी दावा $V'$ च्या व्याख्येवरून
  मिळतो; $\lfalse$ दोन्ही प्रतिमानांत कोठेही सत्य नसतो.
  $\land$ आणि $\lor$ यांचे प्रकार त्यांच्या स्थानिक
  सत्यता-अटी आणि विगमन-गृहीतकावरून मिळतात.
  $\psi\ident \chi\lif \theta$ असू द्या. मूळ प्रतिमानात
  $\mSat{M}{\psi}[w]$ असेल आणि $[w]\subseteq[u]$,
  तर $\psi\in[w]\subseteq[u]$; म्हणून $\mSat{M}{\psi}[u]$.
  $\mSat{M'}{\chi}[{[u]}]$ असेल, तर विगमन-गृहीतकामुळे
  $\mSat{M}{\chi}[u]$; मूळ संबंधाच्या परावर्तितेने
  $\mSat{M}{\theta}[u]$, आणि पुन्हा विगमन-गृहीतकामुळे
  $\mSat{M'}{\theta}[{[u]}]$. म्हणून $\mSat{M'}{\psi}[{[w]}]$.
  उलट, $\mSat/{M}{\psi}[w]$ असेल, तर काही $u$ साठी
  $Rwu$, $\mSat{M}{\chi}[u]$ आणि $\mSat/{M}{\theta}[u]$.
  मूळ प्रतिमानातील सत्यतेच्या एकस्वनिकतेमुळे
  $[w]\subseteq[u]$; विगमन-गृहीतकाने $[u]$ ही नव्या
  प्रतिमानातही त्या सशर्त सूत्राची प्रतिउदाहरण-अवस्था आहे.
  $\psi\ident\lnot \chi$ असेल, तर त्याच पद्धतीने:
  $\mSat{M}{\psi}[w]$ आणि $[w]\subseteq[u]$ वरून
  $\mSat{M}{\psi}[u]$, म्हणून $\mSat/{M}{\chi}[u]$ आणि
  $\mSat/{M'}{\chi}[{[u]}]$. त्यामुळे नव्या प्रतिमानातही
  $\lnot \chi$ सत्य आहे. मूळ प्रतिमानात $\lnot \chi$ असत्य
  असेल, तर $Rwu$ आणि $\mSat{M}{\chi}[u]$ असलेली अवस्था
  $u$ मिळते; तिच्यासाठी $[w]\subseteq[u]$ आणि
  $\mSat{M'}{\chi}[{[u]}]$, म्हणून नव्या प्रतिमानातही नकरण
  असत्य आहे. प्रत्येक प्रकारात आवश्यक घटक-उपसूत्रे $S$
  मध्ये असल्याने विगमन-गृहीतक लागू आहे.
  अखेरीस $\mSat/{M}{\varphi}$ असल्याने काही $w\in W$
  साठी $\mSat/{M}{\varphi}[w]$; समतुल्यतेवरून
  $\mSat/{M'}{\varphi}[{[w]}]$. म्हणून $\mSat/{M'}{\varphi}$.
\end{proof}

\begin{prob}
\ref{int:sc:dec:thm:decidability} मधील योग्य
उपसूत्र-गाळणीचे तपशील तपासा: $V'$ नेमके परिभाषित आहे,
अंशतः क्रमाच्या दृष्टीने एकस्वनिक आहे, आणि $|W'|\leq2^{|S|}$
हे दाखवा. $S$ उपसूत्रे घेण्याखाली संवृत असण्याची अट
विगमन-सिद्धतेत कुठे वापरली आहे ते स्पष्ट करा.
\end{prob}

\ref{int:sc:dec:thm:decidability} आणि त्याच्या सिद्धतेतील
मर्यादेवरून $\Entails \varphi$ ठरवणारा अल्गोरिदम मिळतो.
$S$ आणि $P$ प्रभावीपणे मोजा. $1\leq n\leq2^{|S|}$
असलेल्या प्रत्येक $n$ साठी $n$ जगांवरील सर्व अंशतः क्रम
आणि $P$ वरील सर्व एकस्वनिक मूल्यांकने यांचा सांत शोध करा.
इतर विधानचरांना रिकामा संच दिल्याने $\varphi$ ची सत्यता बदलत नाही.
प्रत्येक अशा प्रतिमानातील प्रत्येक जगात $\varphi$ च्या सत्यतेचे
recursive clauses वापरून मूल्यांकन करा. प्रतिउदाहरण मिळाल्यास
$\Entails/ \varphi$; कोणतेही मिळाले नाही, तर प्रमेयानुसार
$\Entails \varphi$. हा शोध नेहमी संपतो.




\chapter{अंतःप्रज्ञावादी टॅब्लो}\label{int:tab:chap}
\begin{editorial}
  अंतःप्रज्ञावादी तर्कशास्त्रातील पूर्वचिन्हयुक्त
  टॅब्लोंविषयी हे मसुदा प्रकरण आहे.
  अधिक उदाहरणे, संपूर्णतेच्या सिद्धता,
  आणि बंद टॅब्लो शोधण्याचा प्रयत्न
  अयशस्वी झाल्यावर प्रतिप्रतिमाने कशी
  मिळवता येतात याची चर्चा हवी आहे.
\end{editorial}









\section{प्रस्तावना}\label{int:tab:int:sec}

अभिजात तर्कशास्त्रातील टॅब्लो म्हणजे चिन्हांकित सूत्रांचे (खाली शाखा फुटणारे) काही वृक्ष.
चिन्हांकित सूत्र म्हणजे सत्यतामूल्याचे चिन्ह
($\True$ किंवा $\False$) आणि वाक्य
यांची जोडी:
\[
\sFmla{\True}{\varphi} \text{किंवा} \sFmla{\False}{\varphi}.
\]
टॅब्लोची सुरुवात काही
\emph{गृहीतकां}पासून होते. पुढील प्रत्येक
चिन्हांकित सूत्र हे अनुमानाचा एखादा
नियम लावून मिळते. काही अनुमान-नियम
वृक्षाच्या टोकावर एक किंवा अधिक
चिन्हांकित सूत्रे जोडतात; इतर
नियम दोन नवी टोके जोडून दोन शाखा
तयार करतात. नियम लावल्यावर मिळणाऱ्या
चिन्हांकित सूत्रे मधील
सूत्र, नियम ज्या
चिन्हांकित सूत्राला लावला
त्यातील सूत्रापेक्षा कमी गुंतागुंतीचे असते.
एखाद्या शाखेत
$\sFmla{\True}{\varphi}$ आणि
$\sFmla{\False}{\varphi}$ ही दोन्ही
असतील, तर ती शाखा \emph{बंद}
आहे असे म्हणतो. टॅब्लो
मधील प्रत्येक शाखा बंद असेल,
तर संपूर्ण टॅब्लो बंद असतो.
बंद टॅब्लो ही अशी
निष्पत्ती असते की तिच्यावरून
त्या टॅब्लोची सुरुवात ज्या
चिन्हांकित सूत्रांनी केली तो
संच अपूर्ततायोग्य आहे हे दिसते.
यावरून $\Proves$ संबंधाची व्याख्या
करता येते: $\Gamma \Proves \varphi$
तेव्हा आणि केवळ तेव्हाच
$\Gamma_0 = \{\psi_1,
\dots, \psi_n\} \subseteq \Gamma$
असा एखादा सांत संच असतो की
पुढील गृहीतकांसाठी बंद
टॅब्लो मिळतो:
\[
\{\sFmla{\False}{\varphi}, \sFmla{\True}{\psi_1}, \dots, \sFmla{\True}{\psi_n}\}.
\]

अंतःप्रज्ञावादी तर्कशास्त्रासाठी
चिन्हांकित सूत्र ही संकल्पना
विस्तारित करावी लागते आणि
तार्किक कारकांचे नियम बदलावे लागतात.
चिन्हाबरोबर ($\True$ किंवा
$\False$) अंतःप्रज्ञावादी
टॅब्लो मधील सूत्रे ना
\emph{पूर्वचिन्हे}~$\sigma$ ही असतात.
ही पूर्वचिन्हे धन पूर्णांकांच्या
रिकाम्या नसलेल्या क्रमिका असतात;
म्हणजे $\sigma \in
(\PosInt)^* \setminus \{\emptyseq\}$.
अशी पूर्वचिन्हे लिहिताना
भोवतालचे $\tuple{\ }$ वगळतो
आणि वेगवेगळ्या घटक मध्ये
$,$ ऐवजी $.$ वापरतो.
$\sigma$ हे पूर्वचिन्ह असेल,
तर $\sigma.n$ म्हणजे
$\sigma \concat \tuple{n}$;
उदाहरणार्थ $\sigma = 1.2.1$
असेल, तर $\sigma.3$ म्हणजे
$1.2.1.3$. म्हणून उदाहरणार्थ,
\[
\sFmla{\True}{\varphi \lif (\psi \lif \chi)}[1.2]
\]
हे \emph{पूर्वचिन्हयुक्त
चिन्हांकित सूत्र} (किंवा संक्षेपाने
\emph{पूर्वचिन्हयुक्त सूत्र})
आहे.

सहज अर्थाने, टॅब्लोच्या
एखाद्या शाखेवरील सूत्रे
पूर्ण करू शकेल अशा प्रतिमानातील
जगाचे नाव पूर्वचिन्ह देते.
आणि $\sigma$ हे एखाद्या जगाचे
नाव असेल, तर $\sigma.n$ हे
म्हणजे~$\sigma$ ने दर्शवलेल्या
जगापासून प्राप्य असलेल्या
जगाचे नाव देते.

अंतःप्रज्ञावादी प्रतिमानांमध्ये
प्राप्यता संबंध परावर्ती आणि
संक्रमक असतो. पूर्वचिन्हांच्या
भाषेत याचा अर्थ असा की
$\sigma$ हा स्वतः~$\sigma$ पासून प्राप्य
असतो; आणि~$\sigma$ चा कोणताही
विस्तार, म्हणजे
$\sigma.n_1.\cdots.n_k$
या रूपातील पूर्वचिन्हसुद्धा,
त्यापासून प्राप्य असते.
$\sigma$ आणि त्याचा कोणताही
विस्तार दर्शवण्यासाठी
$\sigma.*$ ही संज्ञा वापरू.
म्हणजे $\sigma.*$ ही पूर्वचिन्हे
नेमकी~$\sigma$ पासून प्राप्य
असलेली सर्व पूर्वचिन्हे आहेत.






\begin{quote}\small\textbf{स्रोतनोंद.} SOL6-A206: उलट चिन्हांची जोडी आणि संपूर्ण शाखेची बंदता यांचे scopes वेगळे केले. नवे साक्षी-पूर्वचिन्ह किंवा त्याचा विस्तार शाखेवर आधी नसावा अशी अट स्पष्ट केली, जेणेकरून आधीच्या जगांचे अर्थनिर्धारण जपता येते.\end{quote}






\section{अंतःप्रज्ञावादी तर्कशास्त्रासाठीचे नियम}\label{int:tab:rul:sec}

$\land$ आणि $\lor$ या कारकांचे नियम नेहमीच्या विधानीय
चिन्हांकित टॅब्लो मधील नियमांसारखेच आहेत;
फक्त त्यांना पूर्वचिन्हे जोडली जातात. प्रत्येक वेळी,
$\sFmla{S}{\varphi}[\sigma]$ या चिन्हांकित सूत्र
ला नियम लावल्याने मिळणाऱ्या नव्या सूत्रे
वरही~$\sigma$ हेच पूर्वचिन्ह असते. हे सहज लक्षात
येते: उदाहरणार्थ, $\varphi \land \psi$ हे~$\sigma$ ने
नाव दिलेल्या जगात सत्य असेल, तर $\varphi$ आणि $\psi$
ही दोन्ही~$\sigma$ येथे सत्य असतात (आणि त्या
जगापासून प्राप्य असलेल्या प्रत्येक जगातही सत्य
असतात). $\land$ आणि $\lor$ यांचे नियम
\ref{int:tab:rul:tab:prop-rules} मध्ये दिले आहेत.

\begin{table}
  \[\def\arraystretch{3}\begin{array}{|c|c|}
    \hline
    \AxiomC{\sFmla{\True}{\varphi \land \psi}[\sigma]}
    \RightLabel{\TRule{\True}{\land}}
    \UnaryInfC{\sFmla{\True}{\varphi}[\sigma]}
    \noLine
    \UnaryInfC{\sFmla{\True}{\psi}[\sigma]}
    \DisplayProof
    &
    \AxiomC{\sFmla{\False}{\varphi \land \psi}[\sigma]}
    \RightLabel{\TRule{\False}{\land}}
    \UnaryInfC{$\sFmla{\False}{\varphi}[\sigma] \quad \mid \quad
      \sFmla{\False}{\psi}[\sigma]$}
    \DisplayProof
    \\[2ex]
    \hline
    \AxiomC{\sFmla{\True}{\varphi \lor \psi}[\sigma]}
    \RightLabel{\TRule{\True}{\lor}}
    \UnaryInfC{$\sFmla{\True}{\varphi}[\sigma] \quad \mid \quad
      \sFmla{\True}{\psi}[\sigma]$}
    \DisplayProof
    &
    \AxiomC{\sFmla{\False}{\varphi \lor \psi}[\sigma]}
    \RightLabel{\TRule{\False}{\lor}}
    \UnaryInfC{\sFmla{\False}{\varphi}[\sigma]}
    \noLine
    \UnaryInfC{\sFmla{\False}{\psi}[\sigma]}
    \DisplayProof
    \\[2ex]
    \hline
  \end{array}\]
  \caption{$\land$ आणि $\lor$ साठीचे पूर्वचिन्हयुक्त
    टॅब्लो नियम}
  \label{int:tab:rul:tab:prop-rules}
\end{table}

शाखा बंद होण्याची अट नेहमीच्या टॅब्लो
प्रमाणेच आहे; पण येथे सूत्रे बरोबर त्यांची
पूर्वचिन्हेही जुळली पाहिजेत. खरे तर ही अट
थोडी शिथिल करता येते: अंतःप्रज्ञावादी प्रतिमानात
सूत्रे एकदा सत्य झाली की पुढेही सत्य राहतात.
म्हणून $\varphi$ हे~$\sigma$ येथे सत्य असताना
ते~$\sigma.{*}$ या प्राप्य पूर्वचिन्हावर असत्य
असू शकत नाही. म्हणून एखाद्या शाखेत पुढील
दोन्ही असतील, तर ती बंद असते:
\[
\sFmla{\True}{\varphi}[\sigma] \quad\text{आणि}\quad \sFmla{\False}{\varphi}[\sigma.{*}]
\]
येथे $\sigma$ हे एखादे पूर्वचिन्ह आणि $\varphi$ हे
सूत्र आहे. उलट चिन्हे असतील, म्हणजे
शाखेत पुढील दोन्ही असतील,
\[
\sFmla{\False}{\varphi}[\sigma] \quad\text{आणि}\quad \sFmla{\True}{\varphi}[\sigma.{*}]
\]
तर सदर सूत्र-जुळणीची अट या उलट चिन्हांच्या जोडीला
तेव्हा आणि केवळ तेव्हाच लागू होते, जेव्हा $*$ ही
रिकामी क्रमिका असते. शाखा इतर कोणत्या बंदतेच्या
अटीमुळे बंद होऊ शकते, हे याने नाकारलेले नाही.

तसेच, एखाद्या शाखेत~$\sFmla{\True}{\bot}[\sigma]$
असेल, तरी ती बंद असते.

गृहीतके मांडण्याचे नियमही नेहमीच्या टॅब्लो
प्रमाणेच आहेत; मात्र गृहीतकांसाठी आपण नेहमी
$1$ हे पूर्वचिन्ह वापरतो. (सर्व गृहीतकांसाठी
एकच पूर्वचिन्ह वापरले, तर नेमके कोणते
पूर्वचिन्ह वापरले याने फरक पडत नाही.)
म्हणून, उदाहरणार्थ,
\[
\psi_1, \dots, \psi_n \Proves \varphi
\]
असे तेव्हा आणि केवळ तेव्हाच म्हणतो जेव्हा पुढील
गृहीतकांसाठी बंद टॅब्लो असतो:
\[
\sFmla{\True}{\psi_1}[1], \dots, \sFmla{\True}{\psi_n}[1],
\sFmla{\False}{\varphi}[1].
\]

शर्त~$\lif$ साठीचे नियम अभिजात आणि मोडल
प्रकरणांपेक्षा वेगळे आहेत. $\TRule{\True}{\lif}$
हा नियम $\sFmla{\True}{\varphi \lif \psi}[\sigma]$
असलेल्या शाखेतून दोन वेगळ्या शाखा काढतो:
एकीवर $\sFmla{\False}{\varphi}[\sigma.{*}]$ आणि
दुसरीवर $\sFmla{\True}{\psi}[\sigma.{*}]$ जोडतो.
हा नियम लावताना $\sigma.{*}$ हे पूर्वचिन्ह
संबंधित शाखेत \emph{आधीपासूनच} आलेले
असले पाहिजे. अशा पूर्वचिन्हाला त्या शाखेवर
‘वापरलेले’ म्हणू. ($\sigma.{*}$ मध्ये
$\sigma$ स्वतःही येत असल्याने, वेगवेगळ्या
शाखांवर $\sFmla{\False}{\varphi}[\sigma]$ आणि
$\sFmla{\True}{\psi}[\sigma]$ जोडून हा नियम
नेहमी लावता येतो.)

$\TRule{\False}{\lif}$ हा नियम
$\sFmla{\False}{\varphi \lif \psi}[\sigma]$
असलेल्या शाखेवर
$\sFmla{\True}{\varphi}[\sigma.n]$ आणि
$\sFmla{\False}{\psi}[\sigma.n]$ ही दोन्ही
सूत्रे जोडतो; येथे $\sigma.n$ हे त्या
शाखेसाठी नवे पूर्वचिन्ह असते: ते किंवा त्याचा कोणताही
विस्तार त्या शाखेवर आधी आलेला नसतो. पुढील नकरणाच्या
नियमांतही ‘नवे’ याच अर्थाने वापरतो.

$\lnot$ साठीचे नियम याच धर्तीवर व्याख्यायित
केले आहेत ($\lnot \varphi$ म्हणजे
$\varphi \lif \lfalse$ ही व्याख्या वापरून).

हे नियम \ref{int:tab:rul:tab:rules-lif-lnot} मध्ये दिले आहेत.

\begin{table}
  \[\def\arraystretch{3}\begin{array}{|c|c|}
    \hline
    \AxiomC{\sFmla{\True}{\lnot \varphi}[\sigma]}
    \RightLabel{\TRule{\True}{\lnot}}
    \UnaryInfC{$\sFmla{\False}{\varphi}[\sigma.{*}]$}
    \DisplayProof
    &
    \AxiomC{\sFmla{\False}{\lnot \varphi}[\sigma]}
    \RightLabel{\TRule{\False}{\lnot}}
    \UnaryInfC{\sFmla{\True}{\varphi}[\sigma.n]}
    \DisplayProof
    \\[1ex]
    \text{$\sigma.{*}$ वापरलेले आहे} & \text{$\sigma.n$ नवे आहे}\\
    \hline
    \AxiomC{\sFmla{\True}{\varphi \lif \psi}[\sigma]}
    \RightLabel{\TRule{\True}{\lif}}
    \UnaryInfC{$\sFmla{\False}{\varphi}[\sigma.{*}] \quad \mid \quad
      \sFmla{\True}{\psi}[\sigma.{*}]$}
    \DisplayProof
    &
    \AxiomC{\sFmla{\False}{\varphi \lif \psi}[\sigma]}
    \RightLabel{\TRule{\False}{\lif}}
    \UnaryInfC{\sFmla{\True}{\varphi}[\sigma.n]}
    \noLine
    \UnaryInfC{\sFmla{\False}{\psi}[\sigma.n]}
    \DisplayProof
    \\[1ex]
    \text{$\sigma.{*}$ वापरलेले आहे} & \text{$\sigma.n$ नवे आहे}\\
    \hline
  \end{array}\]
  \caption{$\lnot$ आणि $\lif$ साठीचे
    पूर्वचिन्हयुक्त टॅब्लो नियम}
  \label{int:tab:rul:tab:rules-lif-lnot}
\end{table}












\section{अंतःप्रज्ञावादी तर्कशास्त्रासाठीचे टॅब्लो}\label{int:tab:prf:sec}

\begin{ex}
  $(\varphi \land \psi) \lif \chi \Proves
  \varphi \lif (\psi \lif \chi)$ हे दाखवणारा
  बंद टॅब्लो पुढे दिला आहे.
  \begin{oltableau}
    [\pFmla{\True}{(\formula{A} \land \formula{B}) \lif \formula{C}}{1},
      just =\TAss
      [\pFmla{\False}{\formula{A} \lif (\formula{B} \lif \formula{C})}{1},
        just =\TAss
        [\pFmla{\True}{\formula{A}}{1.1},
          just = {\TRule{\False}{\lif}[2]}
          [\pFmla{\False}{\formula{B} \lif \formula{C}}{1.1},
            just = {\TRule{\False}{\lif}[2]}
            [\pFmla{\True}{\formula{B}}{1.1.1},
              just = {\TRule{\False}{\lif}[4]}
              [\pFmla{\False}{\formula{C}}{1.1.1},
                just = {\TRule{\False}{\lif}[4]}
                [\pFmla{\False}{\formula{A} \land \formula{B}}{1.1.1},
                  just = {\TRule{\True}{\lif}[1]}
                  [\pFmla{\False}{\formula{A}}{1.1.1},
                    just= {\TRule{\False}{\land}[7]}, close]
                  [\pFmla{\False}{\formula{B}}{1.1.1},
                    just= {\TRule{\False}{\land}[7]}, close]]
                [\pFmla{\True}{\formula{C}}{1.1.1},
                  just = {\TRule{\True}{\lif}[1]}, close]]
            ]
          ]
        ]
      ]
    ]
  \end{oltableau}
\end{ex}



\begin{prob}
  पुढील निष्कर्ष दाखवण्यासाठी बंद अंतःप्रज्ञावादी
  टॅब्लो शोधा:
  \begin{enumerate}
    \item $\Proves \varphi \lif (\psi \lif \varphi)$
    \item $\Proves \lnot(\varphi \land \lnot \varphi)$
    \item $\varphi \lif (\psi \lif \chi) \Proves (\varphi \land \psi) \lif \chi$
    \item $\lnot \varphi \lor \lnot \psi \Proves \lnot(\varphi \land \psi)$
  \end{enumerate}
\end{prob}






\begin{quote}\small\textbf{स्रोतनोंद.} SOL6-A207: मूळ adjacent-only अर्थनिर्धारणातील gap दुरुस्त करून सर्व वापरलेल्या पूर्वचिन्ह-विस्तारांसाठी प्राप्यता आवश्यक केली. नवा witness prefix घेताना सर्व जुने ancestors आणि जगांचे अर्थनिर्धारण जपले जाते हे सिद्ध केले; OLINC-233 ही ऐतिहासिक स्रोतदोषाची नोंद जपली.\end{quote}






\section{अंतःप्रज्ञावादी टॅब्लोची निर्दोषता}\label{int:tab:sou:sec}

\begin{explain}
  अंतःप्रज्ञावादी टॅब्लो निर्दोष आहेत हे
  दाखवण्यासाठी पुढील गृहीतकांना
  \[
  \sFmla{\True}{\psi_1}[1], \dots, \sFmla{\True}{\psi_n}[1], \sFmla{\False}{\varphi}[1]
  \]
  बंद टॅब्लो असेल, तर $\psi_1, \dots, \psi_n \Entails \varphi$
  हे दाखवावे लागेल. व्यत्यस्त रूप सिद्ध करणे सोपे आहे:
  एखाद्या $\mModel{M}$ प्रतिमानात आणि जगात~$w$,
  सर्व $i=1$, \dots,~$n$ साठी
  $\mSat{M}{\psi_i}[w]$ असेल, पण
  $\mSat/{M}{\varphi}[w]$ असेल, तर कोणताही
  टॅब्लो बंद होऊ शकत नाही. असे
  प्रतिप्रतिमान टॅब्लोची सुरुवातीची
  गृहीतके पूर्ततायोग्य असल्याचे दाखवते.
  पुराव्याची युक्ती अशी: टॅब्लो
  च्या शाखेवरील सर्व पूर्वचिन्हयुक्त
  सूत्रे पूर्ततायोग्य असतील, तर
  कोणताही नियम लावल्यानंतर मिळणाऱ्या
  विस्तारित शाखांपैकी किमान एक शाखा
  पूर्ततायोग्य असते. बंद शाखा अपूर्ततायोग्य
  असल्याने, पूर्ततायोग्य पूर्वचिन्हयुक्त
  सूत्रांच्या संचावरील कोणत्याही
  टॅब्लोमध्ये किमान एक खुली शाखा
  असलीच पाहिजे.

  ही युक्ती अंतःप्रज्ञावादी प्रकरणात
  वापरण्यासाठी ‘पूर्ततायोग्य’ या संकल्पनेची
  व्याख्या प्राप्यता-संबंध आणि पूर्वचिन्हे
  यांना सामावून घेईल अशी वाढवावी लागते.
  त्यानंतर पुरावा सरळ आहे.
\end{explain}

\begin{defn}
  $P$ हा पूर्वचिन्हांचा एखादा संच असू द्या,
  म्हणजे $P \subseteq (\PosInt)^*
  \setminus \{\emptyseq\}$; तसेच
  $\mModel{M}$ हे प्रतिमान असू द्या.
  $f\colon P\to W$ हे फलन $\mModel{M}$ मधील
  $P$ चे \emph{अर्थनिर्धारण} असते तेव्हा आणि केवळ तेव्हाच,
  जेव्हा प्रत्येक $\sigma,\tau\in P$ साठी $\sigma$ हा
  $\tau$ चा आरंभीचा खंड असेल, तर $Rf(\sigma)f(\tau)$.
  ही अट सर्व विस्तारांसाठी आहे; मधली पूर्वचिन्हे
  $P$ मध्ये असणे आवश्यक नाही.

  पूर्वचिन्हे~$P$ यांच्या अर्थनिर्धारणाच्या
  संदर्भात पुढील व्याख्या करता येतात:
  \begin{enumerate}
  \item $\mModel{M}$ हे $\sFmla{\True}{\varphi}[\sigma]$
    पूर्ण करते तेव्हा आणि केवळ तेव्हाच
    $\mSat{M}{\varphi}[f(\sigma)]$.
  \item $\mModel{M}$ हे $\sFmla{\False}{\varphi}[\sigma]$
    पूर्ण करते तेव्हा आणि केवळ तेव्हाच
    $\mSat/{M}{\varphi}[f(\sigma)]$.
  \end{enumerate}
\end{defn}

\begin{quote}\small\textbf{स्रोतनोंद.} OLINC-233: मूळ adjacent-prefix व्याख्येमुळे मधली पूर्वचिन्हे नसताना दीर्घ विस्तारापर्यंत प्राप्यता मिळत नव्हती. सध्याच्या मराठी व्याख्येत सर्व पूर्वचिन्ह-विस्तारांसाठी प्राप्यता थेट आवश्यक केली आहे; खालील बंदता-सिद्धतेला आता ती अट उपलब्ध आहे.\end{quote}
या व्याख्येनुसार $\sigma$ आणि $\sigma.{*}$ ही दोन्ही
पूर्वचिन्हे वापरलेली असतील, तर $Rf(\sigma)f(\sigma.{*})$.
रिकाम्या विस्तारासाठी हे $R$ च्या परावर्तितेनेही मिळते.

\begin{defn}
  $\Gamma$ हा पूर्वचिन्हयुक्त सूत्रे
  चा संच असू द्या आणि त्यात आलेल्या
  पूर्वचिन्हांचा संच $P(\Gamma)$ असू द्या.
  $f$ ही $P(\Gamma)$ ची
  $\mModel{M}$ मधील अर्थनिर्धारण
  असेल, तर $\mModel{M}$ हे
  $f$ च्या संदर्भात $\Gamma$
  पूर्ण करते, म्हणजे
  $\mSat{M}{\Gamma}[f]$, असे
  तेव्हा म्हणतो जेव्हा
  $\mModel{M}$ हे $\Gamma$ मधील
  प्रत्येक पूर्वचिन्हयुक्त
  सूत्र हे~$f$ च्या संदर्भात
  पूर्ण करते. $\mModel{M}$ असे
  प्रतिमान आणि $P(\Gamma)$ ची
  $f$ असे अर्थनिर्धारण अस्तित्वात
  असतील की $\mSat{M}{\Gamma}[f]$,
  तेव्हा आणि केवळ तेव्हाच $\Gamma$ \emph{पूर्ततायोग्य}
  आहे.
\end{defn}

\begin{prop}
  $\Gamma$ मध्ये एखाद्या
  सूत्र~$\varphi$ आणि पूर्वचिन्ह~$\sigma$
  साठी $\sFmla{\True}{\varphi}[\sigma]$ आणि
  $\sFmla{\False}{\varphi}[\sigma.{*}]$
  ही दोन्ही असतील, किंवा
  $\sFmla{\True}{\lfalse}[\sigma]$
  असेल, तर $\Gamma$ अपूर्ततायोग्य आहे.
\end{prop}

\begin{proof}
  नेहमीच $\mSat/{M}{\lfalse}[f(\sigma)]$
  असल्याने, $\sFmla{\True}{\lfalse}$
  असलेला $\Gamma$ अपूर्ततायोग्य असतो.

  आता दोन्ही चिन्हांकित सूत्रे पूर्ण
  करणारे $\mModel{M}$ प्रतिमान आणि
  $P(\Gamma)$ चे $f$ अर्थनिर्धारण
  असू शकत नाहीत. कारण
  $\mSat{M}{\varphi}[f(\sigma)]$
  असेल, तर
  \ref{int:sem:rel:prop:true-monotonic}
  आणि $Rf(\sigma)f(\sigma.{*})$
  यांवरून
  $\mSat{M}{\varphi}[f(\sigma.{*})]$
  मिळते. म्हणून
  $\mSat{M}{\varphi}[f(\sigma)]$ आणि
  $\mSat/{M}{\varphi}[f(\sigma.{*})]$
  ही दोन्ही एकत्र असू शकत नाहीत.
\end{proof}

\begin{thm}[निर्दोषता]
  \label{int:tab:sou:thm:tableau-soundness}
  $\Gamma$ ला बंद टॅब्लो
  असेल, तर $\Gamma$ अपूर्ततायोग्य
  आहे.
\end{thm}

\begin{proof}
टॅब्लोच्या एखाद्या शाखेवरील
चिन्हांकित सूत्रांचा संच
पूर्ततायोग्य असेल, तर ती शाखा
पूर्ततायोग्य म्हणू. किमान एक
पूर्ततायोग्य शाखा असलेला
टॅब्लो पूर्ततायोग्य म्हणू.

आपण पुढील विधान दाखवू:
पूर्ततायोग्य टॅब्लो ला अनुमानाचा
कोणताही एक नियम लावल्यावरही
पूर्ततायोग्य टॅब्लो मिळतो.
यावरून प्रमेय सिद्ध होते: कोणताही
बंद टॅब्लो हा~$\Gamma$
मधील गृहीतके एवढीच असलेल्या
टॅब्लो ला नियम लावून मिळतो.
म्हणून $\Gamma$ पूर्ततायोग्य
असेल, तर त्यावरील प्रत्येक
टॅब्लो पूर्ततायोग्य असेल.
पण बंद टॅब्लो पूर्ततायोग्य
नसतो, कारण त्याच्या सर्व शाखा
बंद असतात आणि बंद शाखा
अपूर्ततायोग्य असतात.

आता पूर्ततायोग्य टॅब्लो,
म्हणजे किमान एक पूर्ततायोग्य शाखा
असलेला टॅब्लो, घेऊ.
नियम लावताना एखाद्या शाखेला
चिन्हांकित सूत्रे जोडली जातात
किंवा तिचे दोन शाखांत विभाजन
होते. नियम न लावलेली एखादी
पूर्ततायोग्य शाखा आधीपासून
असेल, तर ती विस्तारित
टॅब्लोमध्येही तशीच
राहते आणि विस्तारित टॅब्लो
पूर्ततायोग्य असतो. म्हणून
पूर्ततायोग्य शाखेलाच नियम
लावल्याचे प्रकरण पाहणे पुरेसे आहे.

त्या शाखेवरील चिन्हांकित सूत्रे
चा संच $\Gamma$ असू द्या आणि
नियम ज्या चिन्हांकित सूत्र
ला लावतो ते
$\sFmla{S}{\varphi}[\sigma] \in \Gamma$
असू द्या. नियमामुळे शाखा फुटत
नसेल, तर नियमाच्या निष्कर्षांसह
$\Gamma$ असलेली विस्तारित शाखा
अजूनही पूर्ततायोग्य आहे हे
दाखवावे लागते. शाखा फुटत
असेल, तर मिळणाऱ्या दोन शाखांपैकी
किमान एक पूर्ततायोग्य आहे हे
दाखवावे लागते.
\begin{enumerate}
\item $\sFmla{\True}{\lnot \psi}[\sigma] \in \Gamma$
  याला $\TRule{\True}{\lnot}$ लावून
  शाखा वाढवू. मग विस्तारित
  शाखेत $\Gamma \cup
  \{\sFmla{\False}{\psi}[\sigma.{*}]\}$
  ही चिन्हांकित सूत्रे असतात.
  $\mSat{M}{\Gamma}[f]$ असे समजा.
  विशेषतः
  $\mSat{M}{\lnot \psi}[f(\sigma)]$.
  म्हणून $Rf(\sigma)w$ असलेल्या
  प्रत्येक $w$ साठी
  $\mSat/{M}{\psi}[w]$; यात
  $f(\sigma.{*})$ सुद्धा येते.
  त्यामुळे $\mModel{M}$ हे
  $\sFmla{\False}{\psi}[\sigma.{*}]$
  $f$ च्या संदर्भात पूर्ण करते.
\item $\sFmla{\False}{\lnot \psi}[\sigma] \in \Gamma$
  याला $\TRule{\False}{\lnot}$
  लावून शाखा वाढवण्याचे प्रकरण:
  सरावासाठी.
\item $\sFmla{\True}{\psi \land \chi}[\sigma] \in \Gamma$
  याला $\TRule{\True}{\land}$
  लावल्यावर त्याच शाखेत
  $\sFmla{\True}{\psi}[\sigma]$
  आणि $\sFmla{\True}{\chi}[\sigma]$
  ही दोन नवी चिन्हांकित सूत्रे
  मिळतात. $\mSat{M}{\Gamma}[f]$
  असे समजा; विशेषतः
  $\mSat{M}{\psi \land
    \chi}[f(\sigma)]$. मग
  $\mSat{M}{\psi}[f(\sigma)]$
  आणि $\mSat{M}{\chi}[f(\sigma)]$.
  म्हणजे $\mModel{M}$ हे
  $\sFmla{\True}{\psi}[\sigma]$
  आणि $\sFmla{\True}{\chi}[\sigma]$
  दोन्ही~$f$ च्या संदर्भात
  पूर्ण करते.
\item $\sFmla{\False}{\psi \lor \chi}[\sigma] \in \Gamma$
  याला $\TRule{\False}{\lor}$
  लावण्याचे प्रकरण: सरावासाठी.
\item $\sFmla{\False}{\psi \lif \chi}[\sigma] \in \Gamma$
  याला $\TRule{\False}{\lif}$
  लावल्यावर शाखेत
  $\sFmla{\True}{\psi}[\sigma.n]$
  आणि $\sFmla{\False}{\chi}[\sigma.n]$
  ही दोन नवी चिन्हांकित सूत्रे
  मिळतात. येथे $\sigma.n$
  हे शाखेसाठी नवे पूर्वचिन्ह आहे: $\sigma.n\notin P(\Gamma)$,
  आणि त्याचा कोणताही विस्तारही $P(\Gamma)$ मध्ये नाही.

  $\Gamma$ पूर्ततायोग्य असल्याने,
  $\mSat{M}{\Gamma}[f]$ असेल
  असे $\mModel{M}$ प्रतिमान
  आणि $P(\Gamma)$ ची~$f$
  अर्थनिर्धारण आहेत. विशेषतः
  $\mSat/{M}{\psi \lif \chi}[f(\sigma)]$.
  पुढील संच पूर्ततायोग्य आहे हे
  दाखवावे लागेल:
  $\Gamma \cup
  \{\sFmla{\True}{\psi}[\sigma.n],
  \sFmla{\False}{\chi}[\sigma.n]\}$.
  त्यासाठी $P(\Gamma)
  \cup \{\sigma.n\}$
  याचे अर्थनिर्धारण पुढीलप्रमाणे
  ठरवू:

  $\mSat/{M}{\psi \lif \chi}[f(\sigma)]$
  असल्याने $Rf(\sigma)w$,
  $\mSat{M}{\psi}[w]$ आणि
  $\mSat/{M}{\chi}[w]$
  असा $w \in W$ असतो.
  $f'$ हे $f$ सारखेच, फक्त
  $f'(\sigma.n) = w$
  असे ठरवू. $f'(\sigma) = f(\sigma)$
  आणि $Rf(\sigma)w$ असल्याने $Rf'(\sigma)f'(\sigma.n)$.
  जुन्या संचातील $\tau$ हा $\sigma.n$ चा आरंभीचा खंड
  असेल, तर तो $\sigma$ चाही आरंभीचा खंड आहे.
  म्हणून जुन्या अर्थनिर्धारणाने $Rf(\tau)f(\sigma)$,
  आणि संक्रमणामुळे $Rf'(\tau)f'(\sigma.n)$.
  नव्या पूर्वचिन्हाचा कोणताही विस्तार जुन्या संचात नाही;
  त्यामुळे त्यापासून जुन्या विस्ताराकडे अतिरिक्त अट येत नाही.
  नव्या पूर्वचिन्हाची स्वतःशी अट परावर्तितेने मिळते.
  म्हणून $f'$ हे $P(\Gamma)\cup\{\sigma.n\}$ चे
  सर्व पूर्वचिन्ह-विस्तारांसाठी अर्थनिर्धारण आहे.
  तसेच $\mSat{M}{\psi}[f'(\sigma.n)]$
  आणि $\mSat/{M}{\chi}[f'(\sigma.n)]$.
  सर्व पूर्वचिन्हे
  $\sigma' \in P(\Gamma)$ साठी
  $f(\sigma') = f'(\sigma')$
  असल्याने
  $\mSat{M}{\Gamma}[f']$.
  म्हणून $\mModel{M}, f'$ हे
  $\Gamma \cup
  \{\sFmla{\True}{\psi}[\sigma.n],
  \sFmla{\False}{\chi}[\sigma.n]\}$
  पूर्ण करते.
\end{enumerate}
आता दोन शाखा निर्माण करणारे
नियम पाहू.
\begin{enumerate}
\item $\sFmla{\False}{\psi \land \chi}[\sigma] \in \Gamma$
  याला $\TRule{\False}{\land}$
  लावल्यावर दोन शाखा मिळतात:
  डावीकडील शाखेत
  $\sFmla{\False}{\psi}[\sigma]$
  आणि उजवीकडील शाखेत
  $\sFmla{\False}{\chi}[\sigma]$.
  $\mSat{M}{\Gamma}[f]$
  असे समजा; विशेषतः
  $\mSat/{M}{\psi \land \chi}[f(\sigma)]$.
  मग $\mSat/{M}{\psi}[f(\sigma)]$
  किंवा $\mSat/{M}{\chi}[f(\sigma)]$.
  पहिल्या प्रकरणात
  $\mModel{M}, f$ हे
  $\sFmla{\False}{\psi}[\sigma]$
  पूर्ण करते; म्हणजे डावी
  शाखा पूर्ततायोग्य आहे.
  दुसऱ्या प्रकरणात
  $\mModel{M}, f$ हे
  $\sFmla{\False}{\chi}[\sigma]$
  पूर्ण करते; म्हणजे उजवी
  शाखा पूर्ततायोग्य आहे.
\item $\sFmla{\True}{\psi \lor \chi}[\sigma] \in \Gamma$
  याला $\TRule{\True}{\lor}$
  लावण्याचे प्रकरण: सरावासाठी.
\item $\sFmla{\True}{\psi \lif \chi}[\sigma] \in \Gamma$
  याला $\TRule{\True}{\lif}$
  लावण्याचे प्रकरण: सरावासाठी.
\end{enumerate}
\end{proof}

\begin{prob}
\ref{int:tab:sou:thm:tableau-soundness}
या प्रमेयाचा पुरावा पूर्ण करा.
\end{prob}

\begin{cor}
\label{int:tab:sou:cor:entailment-soundness}
$\Gamma \Proves \varphi$ असेल,
तर $\Gamma \Entails \varphi$.
\end{cor}

\begin{proof}
  $\Gamma \Proves \varphi$ असेल,
  तर काही
  $\psi_1$, \dots, $\psi_n \in
  \Gamma$ साठी
  $\Delta = \{\sFmla{\False}{\varphi}[1], \sFmla{\True}{\psi_1}[1],
  \dots, \sFmla{\True}{\psi_n}[1]\}$
  याला बंद टॅब्लो असतो.
  $\Gamma \Entails \varphi$
  हे दाखवायचे आहे. तसे नाही
  असे समजा. मग एखाद्या
  $\mModel{M}$ प्रतिमानात
  आणि जगात~$w$,
  $i=1$, \dots,~$n$ साठी
  $\mSat{M}{\psi_i}[w]$,
  पण $\mSat/{M}{\varphi}[w]$.
  $f(1) = w$ असे ठरवू;
  मग $f$ ही
  $P(\Delta)$ ची
  $\mModel{M}$ मधील
  अर्थनिर्धारण आहे आणि
  $\mModel{M}$ हे
  $\Delta$ ला~$f$ च्या
  संदर्भात पूर्ण करते.
  पण \ref{int:tab:sou:thm:tableau-soundness}
  नुसार, बंद टॅब्लो
  असल्याने $\Delta$
  अपूर्ततायोग्य आहे; ही
  विसंगती आहे. म्हणून
  $\Gamma \Entails \varphi$
  हेच मिळते.
\end{proof}

\begin{cor}
\label{int:tab:sou:cor:weak-soundness}
$\Proves \varphi$ असेल, तर
$\varphi$ हे सर्व प्रतिमानांत
सत्य असते.
\end{cor}



\part{प्रतिवास्तविक विधाने}\label{cnt:part}
\chapter{प्रस्तावना}\label{cnt:int:chap}









\section{वास्तविक अभिव्यंजन}\label{cnt:int:mat:sec}

इंग्रजीत सशर्त विधानाचे सर्वात सोपे रूप
‘‘जर \dots तर \dots’’ असे असते; येथे
\dots{} यांच्या जागी स्वतः विधाने येतात,
उदाहरणार्थ, ‘‘जर हे काम सेवकाने केले
असेल, तर माळी निर्दोष आहे.’’ प्राथमिक
तर्कशास्त्रात अशी सशर्त विधाने
$\lif$ या कारकाने दर्शवायला शिकतो:
\dots{} ने दाखवलेले भाग, उदाहरणार्थ,
सूत्रे $\varphi$ आणि~$\psi$ यांनी
दर्शवायचे; मग संपूर्ण सशर्त विधान
$\varphi \lif \psi$ असे दर्शवतात.

$\lif$ हा कारक \emph{सत्यता-फलनात्मक}
आहे: $\varphi \lif \psi$ चे सत्यतामूल्य---$\True$
किंवा $\False$---हे $\varphi$ आणि~$\psi$
यांच्या सत्यतामूल्यांवरून ठरते.
$\varphi$ असत्य असेल किंवा $\psi$ सत्य
असेल तेव्हा आणि केवळ तेव्हाच $\varphi \lif \psi$
सत्य असते; अन्यथा असत्य असते.
$\pAssign v$ या सत्यतामूल्य-नियुक्तीच्या
संदर्भात $\pSat{v}{\varphi \lif \psi}$
तेव्हा आणि केवळ तेव्हाच, जेव्हा
$\pSat/{v}{\varphi}$ किंवा
$\pSat{v}{\psi}$, अशी व्याख्या
करतो. ही चिन्हार्थमीमांसा
असलेल्या $\lif$ कारकाला
\emph{वास्तविक अभिव्यंजन}
म्हणतात.

या व्याख्येवरून काही प्राथमिक तार्किक
तथ्ये मिळतात. सर्वप्रथम,
वास्तविक अभिव्यंजनासाठी निगमन प्रमेय लागू होते:
\begin{align}
  \text{जर } \Gamma, \varphi \Entails \psi \text{ तर } \Gamma \Entails \varphi \lif \psi
\end{align}
ते सत्यता-फलनात्मक आहे:
$\varphi \lif \psi$ आणि
$\lnot \varphi \lor \psi$ सममूल्य आहेत:
\begin{align}
  \varphi \lif \psi & \Entails \lnot \varphi \lor \psi\\
  \lnot \varphi \lor \psi & \Entails \varphi \lif \psi
  \intertext{वास्तविक अभिव्यंजन त्याच्या
    उत्तरांगावरून आणि पूर्वांगाच्या
    निषेधावरून तार्किकरीत्या निष्पन्न होते:}
  \psi & \Entails \varphi \lif \psi\\
  \lnot \varphi & \Entails \varphi \lif \psi
  \intertext{असत्य वास्तविक अभिव्यंजन
    हे त्याच्या पूर्वांगाचा आणि उत्तरांगाच्या
    निषेधाचा संधियोग याच्याशी सममूल्य असते:
    $\varphi \lif \psi$ असत्य असेल, तर
    $\varphi \land \lnot \psi$ सत्य असते,
    आणि उलटही:}
  \lnot(\varphi \lif \psi) & \Entails \varphi \land \lnot \psi\\
  \varphi \land \lnot \psi & \Entails \lnot(\varphi \lif \psi)
  \intertext{वास्तविक अभिव्यंजनामुळे
    विध्यात्मक अनुमान लागू होते:}
  \varphi, \varphi \lif \psi & \Entails \psi
  \intertext{वास्तविक अभिव्यंजन
    निष्कर्षांचे संयुग्मीकरण करते:}
  \varphi \lif \psi,  \varphi \lif \chi & \Entails \varphi \lif (\psi \land \chi)
  \intertext{पूर्वांग अधिक बलवान
    नेहमी करता येते; म्हणजे
    हे अभिव्यंजन \emph{एकस्वनिक} आहे:}
  \varphi \lif \psi & \Entails (\varphi \land \chi) \lif \psi
  \intertext{वास्तविक अभिव्यंजन
    संक्रमक आहे; म्हणजे साखळी-नियम वैध आहे:}
  \varphi \lif \psi, \psi \lif \chi & \Entails \varphi \lif \chi
  \intertext{वास्तविक अभिव्यंजन
    त्याच्या प्रतिपरिवर्तित विधानाशी सममूल्य आहे:}
  \varphi \lif \psi & \Entails \lnot \psi \lif \lnot \varphi\\
  \lnot \psi \lif \lnot \varphi & \Entails \varphi \lif \psi
\end{align}

हे सर्व निष्कर्ष गणिती युक्तिवादात
उपयुक्त आणि निर्विवाद आहेत.
मात्र गणिताबाहेरील संदर्भांतील
सममूल्य निष्कर्षांना कथित
प्रतिदृष्टांत तत्त्वज्ञान आणि भाषाविज्ञान
यांच्या साहित्यात विपुल आहेत.
त्यांवरून असे सुचते की
इंग्रजीतील ‘‘जर \dots तर \dots’’
विधाने दर्शवताना
वास्तविक अभिव्यंजन~$\lif$
हा कारक योग्य नसतो---किंवा
किमान नेहमीच योग्य नसतो.












\section{वास्तविक अभिव्यंजनाचे विरोधाभास}\label{cnt:int:par:sec}

इंग्रजीतील ‘‘जर \dots तर \dots’’
विधानांचे प्रतीकीकरण $\varphi \lif \psi$
या रूपात करण्यावर टीका करणाऱ्यांपैकी
C.~I. Lewis हे सुरुवातीचे एक होते.
व्हाइटहेड आणि रसेल यांच्या
\emph{Principia Mathematica} मध्ये
वास्तविक अभिव्यंजनाचा वापर करताना
$\lif$ चा उच्चार ‘‘यावरून
तार्किकरीत्या निष्पन्न होते’’ असा
केला जात असे; त्यावर लुईस यांची
टीका होती. $\lif$ चा अर्थ
‘‘तार्किकरीत्या निष्पन्न होते’’
असा घेतला, तर कोणतेही असत्य
विधान~$p$ यावरून $p$ यावरून~$q$
निष्पन्न होते असे ठरेल, कारण
$p$ असत्य असताना
$p \lif (p \lif q)$ सत्य असते.
त्याचप्रमाणे कोणतेही सत्य
विधान~$q$ यावरून $p$ यावरून~$q$
निष्पन्न होते असे ठरेल, कारण
$q$ सत्य असताना
$q \lif (p \lif q)$ सत्य असते,
ही लुईस यांची रास्त तक्रार होती.

तर्कशास्त्रज्ञांना अर्थातच माहीत
आहे की \emph{तार्किक निष्पन्नता}
हा कारक नाही; तो सूत्रे
किंवा विधानांमधील संबंध आहे.
म्हणून गोंधळ टाळण्यासाठी
$\lif$ ला ‘‘तार्किकरीत्या निष्पन्न
होते’’ असे वाचूच नये.\footnote{
  गणितज्ञ आणि संगणकशास्त्रज्ञ
  ‘‘$\lif$’’ चा उच्चार ‘‘यावरून
  निष्पन्न होते’’ असा अजूनही
  वारंवार करतात; पण लुईस यांनी
  दाखवलेला गोंधळ टाळण्यासाठी
  तत्त्वज्ञ ‘‘तरच’’ असा
  उच्चार करतात.} तसे
वाचत नसू, तर लुईस यांची
ती विशिष्ट चिंता उरत नाही:
$p$ हे एखादे असत्य इंग्रजी
विधान दर्शवत असले, तरी
$p$ यावरून $q$ तार्किकरीत्या
‘‘निष्पन्न’’ होत नाही.
$p \Entails q$ आहे का हे
ठरवण्यासाठी \emph{सर्व}
सत्य-मूल्यांकने विचारात घ्याव्या
लागतात; योगायोगाने असत्य
ठरणारे विधान दर्शवण्यासाठी
$p$ वापरले, तरी
$p \Entails/ q$.

तरीही लुईस यांच्या पुढीलसारख्या
‘‘जर \dots तर\dots’’ विधानात
काहीतरी विचित्र वाटते:
\begin{quote}
चंद्र हिरव्या चीजचा बनलेला असेल,
तर $2+2=4$.
\end{quote}
आणि पुढील अनुमानांतही:
\begin{quote}
  चंद्र हिरव्या चीजचा बनलेला नाही.
  म्हणून चंद्र हिरव्या चीजचा
  बनलेला असेल, तर $2+2=4$.

  $2+2 = 4$. म्हणून चंद्र
  हिरव्या चीजचा बनलेला
  असेल, तर $2+2=4$.
\end{quote}
पण ‘‘जर \dots तर \dots’’
याचा अर्थ केवळ $\lif$
असा असेल, तर ते विधान
निर्विवाद सत्य आणि ही
अनुमाने निर्विवाद वैध ठरतील.

दुसरा प्रश्न $(\varphi \lif \psi) \lor (\psi \lif
\varphi)$ या सर्वतः सत्य सूत्राबद्दल
आहे. त्यावरून असे सुचेल की
वृत्तपत्रातून $S$ आणि $T$
ही कोणतीही दोन निश्चयार्थी
वाक्ये निवडली, तर ‘‘जर
$S$ तर $T$, किंवा जर
$T$ तर~$S$’’ हे विधान
सत्य असले पाहिजे.












\section{कठोर अभिव्यंजन}\label{cnt:int:str:sec}

लुईस यांनी \emph{कठोर अभिव्यंजन}~$\strictif$
हा कारक मांडला आणि वास्तविक अभिव्यंजनाऐवजी
तो तार्किक निष्पन्नतेशी जुळतो, असा युक्तिवाद
केला. अवश्यतावाचक मोडल तर्कशास्त्रात
$\varphi \strictif \psi$ ची व्याख्या
$\Box(\varphi \lif \psi)$ अशी करता येते.
म्हणून एखाद्या जगात कठोर अभिव्यंजन
सत्य असते तेव्हा आणि केवळ तेव्हाच संबंधित
वास्तविक अभिव्यंजन अवश्य सत्य असते.

वास्तविक अभिव्यंजनाच्या विरोधाभासांच्या
बाबतीत कठोर अभिव्यंजन कसे वागते?
पूर्वांग असत्य असलेले कठोर अभिव्यंजन,
तसेच उत्तरांग सत्य असलेले कठोर अभिव्यंजन,
सत्यही असू शकते किंवा असत्यही.
शिवाय $(\varphi \strictif \psi) \lor (\psi \strictif \varphi)$
वैध नाही. $\varphi \strictif \psi$
हे कठोर अभिव्यंजन
$\lnot \varphi \lor \psi$ शीही
सममूल्य नाही; त्यामुळे ते
सत्यता-फलनात्मक नाही.

S5 च्या संदर्भात पुढील संबंध मिळतात:
\begin{align}
  \varphi \strictif \psi & \Entails \lnot \varphi \lor \psi
  \text{ पण:}\\
  \lnot \varphi \lor \psi & \Entails/ \varphi \strictif \psi\\
  \psi & \Entails/ \varphi \strictif \psi\\
  \lnot \varphi & \Entails/ \varphi \strictif \psi\\
  \lnot(\varphi \strictif \psi) & \Entails/ \varphi \land \lnot \psi
  \text{ पण:}\\
  \varphi \land \lnot \psi & \Entails \lnot(\varphi \strictif \psi)
  \intertext{तरीही कठोर अभिव्यंजनामुळे
    विध्यात्मक अनुमान लागू होते:}
  \varphi, \varphi \strictif \psi & \Entails \psi
  \intertext{कठोर अभिव्यंजन
    निष्कर्षांचे संयुग्मीकरण करते:}
  \varphi \strictif \psi,  \varphi \strictif \chi & \Entails \varphi \strictif (\psi \land \chi)
  \intertext{कठोर अभिव्यंजनासाठी
    पूर्वांग अधिक बलवान करता येते:}
  \varphi \strictif \psi & \Entails (\varphi \land \chi) \strictif \psi
  \intertext{कठोर अभिव्यंजन
    संक्रमकही आहे:}
  \varphi \strictif \psi, \psi \strictif \chi & \Entails \varphi \strictif \chi
  \intertext{शेवटी, कठोर अभिव्यंजन
    त्याच्या प्रतिपरिवर्तित विधानाशी
    सममूल्य आहे:}
  \varphi \strictif \psi & \Entails \lnot \psi \strictif \lnot \varphi\\
  \lnot \psi \strictif \lnot \varphi & \Entails \varphi \strictif \psi
\end{align}

\begin{prob}
  कठोर अभिव्यंजनासाठी लागू न होणाऱ्या
  तार्किक निष्पन्नता-संबंधांना
  \Log{S5} मधील प्रतिदृष्टांत द्या;
  म्हणजे पुढील प्रत्येकासाठी:
  \begin{enumerate}
  \item $\lnot p \Entails/ \Box(p \lif q)$
  \item $q \Entails/ \Box(p \lif q)$
  \item $\lnot\Box(p \lif q) \Entails/ p \land \lnot q$
  \item $\Entails/ \Box(p \lif q) \lor \Box(q \lif p)$
  \end{enumerate}
\end{prob}

\begin{prob}
  वैध तार्किक निष्पन्नता-संबंध
  कठोर अभिव्यंजनासाठी लागू होतात
  हे दाखवण्यासाठी पुढील
  विधानांचे \Log{S5} मध्ये
  पुरावे द्या:
  \begin{enumerate}
  \item  $\Box(\varphi \lif \psi) \Entails \lnot \varphi \lor \psi$
  \item $\varphi \land \lnot \psi \Entails \lnot\Box(\varphi \lif \psi)$
  \item $\varphi, \Box(\varphi \lif \psi) \Entails \psi$
  \item $\Box(\varphi \lif \psi), \Box(\varphi \lif \chi) \Entails \Box(\varphi \lif (\psi
    \land \chi))$
  \item $\Box(\varphi \lif \psi) \Entails \Box((\varphi \land \chi) \lif \psi)$
  \item $\Box(\varphi \lif \psi), \Box(\psi \lif \chi) \Entails \Box(\varphi
    \lif \chi)$
  \item $\Box(\varphi \lif \psi) \Entails \Box(\lnot \psi \lif \lnot \varphi)$
  \item $\Box(\lnot \psi \lif \lnot \varphi) \Entails \Box(\varphi \lif \psi)$
  \end{enumerate}
  \end{prob}

तरी कठोर अभिव्यंजनाचे स्वतःचेही
‘‘विरोधाभास’’ आहेत. पूर्वांग असत्य
असेल किंवा उत्तरांग सत्य असेल,
तर वास्तविक अभिव्यंजन सत्य
असते. त्याचप्रमाणे, पूर्वांग
\emph{अवश्य} असत्य असेल
किंवा उत्तरांग अवश्य सत्य असेल,
तर कठोर अभिव्यंजन सत्य असते.
S5 मध्ये कोणतेही सत्य कठोर
अभिव्यंजन अवश्य सत्य असते
आणि कोणतेही असत्य कठोर
अभिव्यंजन अवश्य असत्य असते.
दुसऱ्या शब्दांत,
\begin{align}
  \Box \lnot \varphi & \Entails \varphi \strictif \psi\\
  \Box \psi & \Entails \varphi \strictif \psi\\
  \varphi \strictif \psi& \Entails \Box(\varphi \strictif \psi)\\
  \lnot(\varphi \strictif \psi) & \Entails \Box\lnot(\varphi \strictif \psi)
\end{align}
$\strictif$ चा अर्थ ‘‘तार्किकरीत्या
निष्पन्न होते’’ असा घेतल्यास
यात अडचण नाही. शेवटी,
तार्किक निष्पन्नतेचे संबंध ही
गणिती तथ्ये आहेत; म्हणून
ती नैमित्तिक असू शकत नाहीत.
पण ‘‘जर \dots तर \dots’’
याचा अर्थ पकडणारा तार्किक
कारक म्हणून $\strictif$
वापरायचा असेल, तर या
गोष्टी प्रश्न निर्माण करतात,
विशेषतः शेवटच्या दोन.
कारण काही ‘‘जर \dots तर \dots’’
विधाने नैमित्तिकरीत्या
सत्य किंवा असत्य असतात;
प्रत्यक्षात ती सहसा
ना अवश्य सत्य असतात,
ना अशक्य.

\begin{prob}
  पुढील विधानांचे \Log{S5}
  मध्ये पुरावे द्या:
  \begin{enumerate}
    \item  $\Box \lnot \varphi \Entails \varphi \strictif \psi$
    \item $\varphi \strictif \psi \Entails \Box(\varphi \strictif \psi)$
    \item $\lnot(\varphi \strictif \psi)  \Entails \Box\lnot(\varphi \strictif \psi)$
  \end{enumerate}
  त्यासाठी $\strictif$ ची
  व्याख्या वापरा.
\end{prob}












\section{प्रतिवास्तविक विधाने}\label{cnt:int:cnt:sec}

इंग्रजीतील ‘‘जर \dots तर \dots’’
या रचनेचा एक अतिशय प्रचलित आणि
महत्त्वाचा प्रकार \emph{to be}
या क्रियापदाच्या भूतकालीन
संकेतार्थी रूपाने बनतो:
‘‘जर असे घडले असते की \dots,
तर असे घडले असते की \dots’’.
अशा सशर्त विधानाचे पूर्वांग
सहसा असत्य, म्हणजे वस्तुस्थितीच्या
विरुद्ध असल्याने, त्यांना
\emph{प्रतिवास्तविक सशर्त विधाने}
म्हणतात. \emph{to be} चे
संकेतार्थी रूप त्यांत असल्याने
त्यांना \emph{संकेतार्थी सशर्त विधाने}
असेही म्हणतात. याउलट
\emph{निश्चयार्थी} सशर्त विधाने
‘‘जर असे आहे की \dots,
तर असे आहे की \dots’’
या रूपात असतात. प्रतिवास्तविक
आणि निश्चयार्थी सशर्त विधानांच्या
सत्यतेच्या अटी वेगळ्या असतात.
अॅडम्स यांचे प्रसिद्ध उदाहरण पाहा:
\begin{quote}
  जर ऑस्वाल्डने केनेडींची हत्या केली नाही,
  तर दुसऱ्या कोणीतरी केली.

  जर ऑस्वाल्डने केनेडींची हत्या केली
  नसती, तर दुसऱ्या कोणीतरी केली असती.
\end{quote}
पहिले विधान निश्चयार्थी, तर दुसरे
प्रतिवास्तविक आहे. पहिले विधान
स्पष्टपणे सत्य आहे: अध्यक्ष जॉन
एफ. केनेडी यांची हत्या
\emph{कोणीतरी} केली हे आपल्याला
माहीत आहे. ती व्यक्ती,
वॉरेन अहवालात म्हटल्याच्या विरुद्ध,
ली हार्वे ऑस्वाल्ड नसेल,
तर दुसऱ्या कोणीतरी हत्या
केली. दुसरे विधान मात्र
वेगळे काही सांगते. त्याचा
दावा असा की ऑस्वाल्डने
केनेडींची हत्या केली नसती
म्हणजे डॅलस येथील गोळीबार
टळला असता किंवा अयशस्वी
ठरला असता, तरी पुढचा
इतिहास अशा प्रकारे घडला
असता की दुसरा हत्येचा
प्रयत्न यशस्वी झाला असता.
हे प्रतिवास्तविक विधान सत्य असण्यासाठी, ऑस्वाल्डने
हत्या केली नसती तरी दुसऱ्याने हत्या केली असती,
अशी त्या पर्यायी परिस्थितीविषयीची अट आवश्यक आहे.
मृत्यू निश्चित व्हावा म्हणून काही प्रबळ शक्तींनी
कट रचलेला असणे हे तिचे एक संभाव्य उदाहरण आहे
(अनेक कटसिद्धांतवादी जसे मानतात तसे); ते एकमेव
संभाव्य कारण असल्याचा तार्किक निष्कर्ष मात्र निघत नाही.

\emph{निश्चयार्थी} सशर्त विधान
वास्तविक अभिव्यंजनाने अचूक
दर्शवले जाते का, हा अजून
वादाचा विषय आहे. विशेषतः,
इंग्रजी निश्चयार्थी सशर्त
विधानांच्या सत्यतेच्या अटी
वास्तविक अभिव्यंजन देते
या मताशी सुसंगत राहील
अशा प्रकारे त्याचे
विरोधाभास ‘‘समजावता’’
येतात का, हा प्रश्न
आहे. याउलट प्रतिवास्तविक
सशर्त विधाने वास्तविक
अभिव्यंजनाने अचूक
दर्शवता येत नाहीत,
यावर वाद नाही.
कारण प्रतिवास्तविक
विधानांची पूर्वांगे
सहसा असत्य असली,
तरी असत्य पूर्वांग
असलेले प्रत्येक
प्रतिवास्तविक विधान
सत्य नसते. उदाहरणार्थ,
ऑस्वाल्डने हत्या केली नसती तर इतर कोणाकडूनही
हत्या झाली नसती, अशी पर्यायी परिस्थिती गृहीत धरल्यास,
अॅडम्स यांचे प्रतिवास्तविक विधान असत्य ठरते.
वॉरेन अहवाल खरा मानणे आणि कट नसणे या दोन
गृहीतकांवरूनच ही अतिरिक्त प्रतिवास्तविक अट
तार्किकरीत्या निष्पन्न होते, असा दावा येथे करत नाही.

प्रतिवास्तविक सशर्त विधाने
कार्यकारणाविषयीच्या
युक्तिवादात महत्त्वाची
असतात: प्रतिवास्तविक
विधानांनी कारणसंबंध
दर्शवणे हे त्याचे
प्रमुख उदाहरण आहे.
उदाहरणार्थ, आगकाडी
ओढल्याने ती पेटते;
हे ‘‘ही आगकाडी
ओढली असती, तर ती
पेटली असती’’ असे
म्हणून दर्शवता येते.
वास्तविक अभिव्यंजनाने,
किंवा सर्वसाधारणपणे
निश्चयार्थी सशर्त विधानांनी,
हे दर्शवता येत नाही:
‘‘आगकाडी ओढली जाते
$\lif$ आगकाडी पेटते’’
हे विधान आगकाडी
कधीच ओढली गेली
नाही तर सत्य असते,
ती ओढली असती
तर काय झाले असते
याचा काहीही संबंध
नसतो. त्याहूनही
वाईट म्हणजे,
‘‘आगकाडी ओढली जाते
$\lif$ आगकाडीचे
फुलांच्या गुच्छात
रूपांतर होते’’ हेही
आगकाडी कधीच
ओढली गेली नाही
तर सत्य असते;
पण ती ओढली
असती तर नक्कीच
तिचे फुलांच्या
गुच्छात रूपांतर
झाले नसते.

\begin{quote}\small\textbf{स्रोतनोंद.} OLINC-238; SOL6-A212: एका सर्वाधिक जवळच्या जगाची मूळ सामायिक सत्यता-अट दुरुस्त केली आहे. स्टालनाकरच्या निवडलेल्या जगाची अट आणि लुईसची सर्व पूर्वांग-जगांना लागू होणारी गोलक-अट वेगळी दिली आहे; सर्वात जवळचे जग अस्तित्वातच नसेल तरी लुईसची अट लागू होते. मागील हत्या-उदाहरणातील कट हा एकमेव आवश्यक कारण असल्याचा अतिरिक्त दावा काढला आहे.\end{quote}
प्रतिवास्तविक विधानांचे योग्य तर्कशास्त्र नेमके काय,
यावर अजूनही वाद आहे. स्टालनाकर आणि लुईस यांनी
प्रभावी, पण वेगवेगळी विश्लेषणे दिली.
‘‘जर~$S$ असे घडले असते, तर~$T$ असे घडले असते’’
या विधानासाठी, स्टालनाकरच्या मांडणीत~$S$ सत्य
असलेले प्राप्य जग उपलब्ध असल्यास एक जवळचे जग
निवडले जाते; त्या निवडलेल्या जगात~$T$ सत्य असेल
तेव्हा आणि केवळ तेव्हाच प्रतिवास्तविक विधान सत्य असते.
अशक्य पूर्वांगासाठी त्याच्या औपचारिक मांडणीत
विशेष असंभव जग वापरले जाते.
लुईसची अट एका अद्वितीय जवळच्या जगावर अवलंबून नाही:
एकतर संबंधित गोलकांत~$S$ सत्य असलेले जगच नसते,
किंवा~$S$ सत्य असलेले किमान एक जग समाविष्ट करणारा
असा गोलक असतो की त्यातील प्रत्येक $S$-जगात~$T$
सत्य असते. सर्वाधिक जवळची $S$-जगे अस्तित्वात
असतील, तर ही अट त्या सर्व जगांत~$T$ सत्य असण्याशी
सममूल्य आहे; त्यांत अनेक जगे असू शकतात.
सर्वाधिक जवळचे $S$-जग नसतानाही गोलक-अट लागू होते.
हे भेद स्टालनाकर यांच्या मूळ लेखाच्या विभाग~II मध्ये
आणि लुईस यांच्या मूळ ग्रंथातील विभाग~1.3 मध्ये दिले आहेत
(\url{https://philosophy.hku.hk/courses/dm/phil2520/ATheoryOfConditionals-Stalnaker.pdf};
\url{https://fitelson.org/piksi/piksi_18/lewis_counterfactuals.pdf}, पृ.~16, 20).
संभाव्य जगांच्या
सत्ताशास्त्राचा
संदर्भ घेत असल्याने
याला ‘‘सत्ताविषयक’’
विश्लेषण म्हणतात.
इतर विश्लेषणे सशर्त
संभाव्यता किंवा
विश्वास-पुनरीक्षणाचे
सिद्धांत वापरतात.
प्रतिवास्तविक विधानांसाठी
अनेक वेगवेगळी
तर्कशास्त्रे सुचवली
गेली आहेत. लुईस--स्टालनाकर
यांचे एकच तर्कशास्त्रही
नाही: सर्वाधिक जवळच्या
संभाव्य जगांच्या आधारे
त्यांनी साधर्म्य असलेली
मांडणी केली, तरी
त्यांच्या भिन्न
गृहीतकांमुळे भिन्न
वैध अनुमाने मिळतात.



\chapter{किमान बदलाची चिन्हार्थमीमांसा}\label{cnt:min:chap}










\section{प्रस्तावना}\label{cnt:min:int:sec}

‘‘ही आगकाडी ओढली असती, तर ती पेटली असती’’
अशा प्रतिवास्तविक सशर्त विधानांचे
स्टालनाकर आणि लुईस यांनी विश्लेषण सुचवले.
अशा वाक्यांच्या सत्यतेच्या अटी
योग्य रीतीने कशा समजाव्यात,
याविषयी त्यांची मांडणी होती.
दोन्ही मांडण्यांमागील कल्पना अशी आहे:
प्रतिवास्तविक सशर्त विधान सत्य आहे का
हे ठरवताना त्याचे पूर्वांग सत्य होण्यासाठी
वास्तविक जगापासून किमान फरक असलेली
संभाव्य जगे विचारात घ्यावी लागतात.
या संभाव्य जगांत उत्तरांग सत्य असेल,
तर प्रतिवास्तविक विधान सत्य असते.
उदाहरणार्थ, माझ्या हातात आगकाडी
आणि आगपेटी आहे असे समजा.
वास्तविक जगात मी फक्त त्यांच्याकडे
पाहतो आणि आगकाडी ओढली तर काय
होईल, याचा विचार करतो.
मी आगकाडी ओढतो अशा पर्यायी जगातला
किमान बदल म्हणजे मी कृती करण्याचे
ठरवतो आणि आगकाडी ओढतो.
हा बदल किमान आहे, कारण आगकाडी ओढणे आणि त्याचे
आवश्यक परिणाम यांखेरीज दुसरे अनावश्यक बदल होत नाहीत: मी त्याच वेळी
हवेत उडी मारत नाही; आगकाडी
ओढल्याने माझ्या केसांनाही आग
लागत नाही; माझ्या बोटांतील सारी
ताकद अचानक नाहीशी होत नाही;
त्याच वेळी कुणी सुपर-सोकरने
माझ्यावर पाणी मारत नाही;
इत्यादी. त्या पर्यायी शक्यतेत
आगकाडी पेटते. म्हणून मी
आगकाडी ओढली असती,
तर ती पेटली असती,
हे सत्य आहे.

या सहज समजणाऱ्या मांडणीला
प्रतिवास्तविक तर्कशास्त्रांची औपचारिक
चिन्हार्थमीमांसा जोडता येते.
प्रतिवास्तविक विधानासाठी लुईस यांनी
‘‘$\boxright$’’ हे चिन्ह, तर
स्टालनाकर यांनी ‘‘$>$’’ हे चिन्ह
वापरले. आपण~$\cif$ वापरू आणि
ते विधानीय तर्कशास्त्रात द्विपदी
संयोजक म्हणून जोडू.
म्हणून $\varphi \lif \psi$ या रूपाच्या
सूत्रे बरोबरच
$\varphi \cif \psi$ या रूपाची
सूत्रे ही असतील.
मोडल तर्कशास्त्राच्या संबंधात्मक
चिन्हार्थमीमांसेप्रमाणे, औपचारिक
चिन्हार्थमीमांसा अशा प्रतिमानांवर
आधारलेली असते की ज्यांत
सूत्रे यांचे मूल्यांकन जगांच्या
संदर्भात केले जाते. त्यात
$\mSat{M}{\varphi \cif \psi}[w]$
याची पूर्ति-अट काही (इतर) जगे~$w'$
यांच्या संदर्भातील
$\mSat{M}{\varphi}[w']$ आणि
$\mSat{M}{\psi}[w']$
यांच्या साहाय्याने दिली जाते.
कोणती $w'$ जगे?
सहज कल्पनेनुसार,
$\mSat{M}{\varphi}[w']$ सत्य असलेली
आणि~$w$ च्या सर्वाधिक जवळ
असलेली जगे. त्यासाठी
‘‘जवळीक’’ दर्शवणारा
संबंधही प्रतिमानात
समाविष्ट करावा लागतो.

प्रतिवास्तविक परिस्थिती चित्ररूपात
दाखवण्याची एक बोधक पद्धत लुईस
यांनी आणली. प्रत्येक संभाव्य जग
हे इतर जगे समाविष्ट करणाऱ्या
एकमेकांच्या आत असलेल्या गोलकांच्या
प्रणालीच्या केंद्रस्थानी असते;
आपण हे गोलक समकेंद्री
वर्तुळांनी दाखवतो. या आकृतीत शेजारच्या दोन गोलकांच्या मधला थर
समान जवळीक दाखवतो. अधिक सूक्ष्म गोलकांची सामान्य
प्रणाली असताना मनमानी दोन गोलकांमधील सर्व जगे
सारख्याच जवळ असण्याची अट नसते.
आतल्या गोलकात असलेली जगे
अधिक जवळ, तर बाहेरच्या
गोलकात असलेली जगे अधिक
दूर असतात.
\begin{center}
\begin{tikzpicture}[scale=.7]\small
  \spheresystem{5}
  \draw (0,0) node {$w$};
  \propositionintersect{0}{4}{40}{3.5}
  {\tikzset{proposition/.append={smooth,tension=1.4}}
  \proposition[shift={(1.6,-1)}]{90}{1}{30}{4}}
  \path (1.5,3) node[below] {$\formula{B}$};
  \path (3,0) node {$\formula{A}$};
\end{tikzpicture}
\end{center}
या आकृतीतील सर्वाधिक जवळची $\varphi$-जगे म्हणजे
$\varphi$ सत्य असलेले किमान एक जग धारण करणाऱ्या
सर्वांत लहान गोलकातील $\varphi$-जगे~$w'$ (राखाडी भाग).
तो प्रणालीतील सर्वांत लहान गोलकच असावा असे नाही.
असा $\varphi$-जग धारण करणारा सर्वांत लहान गोलक
सर्व प्रतिमानांत अस्तित्वात असतोच असे नाही;
पुढील विभागातील गोलक-अट अशा प्रकरणांनाही लागू होते.
सहज कल्पनेनुसार,
सर्वाधिक जवळच्या सर्व
$\varphi$-जगांत $\psi$ सत्य
असेल, तर~$w$ येथे
$\varphi \cif \psi$ ची
पूर्ति होते.













\section{गोलक प्रतिमाने}\label{con:min:sph:sec}

प्रतिवास्तविक विधानांना औपचारिक
चिन्हार्थमीमांसा देण्याचा एक मार्ग
म्हणजे लुईस यांच्या अनौपचारिक
मांडणीचे गणिती रचनेत रूपांतर करणे.
मग जग~$w$ भोवतीचे गोलक
हे जगांचे संच असतात.
गोलक एकमेकांच्या आत
असल्यामुळे, जग~$w$ भोवतीच्या
या जगांच्या संचांना उपसंच-संबंधाने
रेषीय क्रमात असावे लागते.

\begin{defn}
  \emph{गोलक प्रतिमान} म्हणजे
  $\mModel{M} = \tuple{W, O, V}$
  हे त्रिक: येथे $W$ हा जगांचा
  रिक्त नसलेला संच आहे,
  $V\colon \PVar \to \Pow{W}$
  हे सत्य-मूल्यांकन आहे, आणि
  $O\colon W \to \Pow{\Pow{W}}$
  प्रत्येक जग~$w$ ला
  \emph{गोलकांची प्रणाली}~$O_w$
  नेमून देते. प्रत्येक~$w$ साठी
  $O_w$ हा जगांच्या संचांचा संच
  असून त्याने पुढील अटी पूर्ण
  केल्या पाहिजेत:
  \begin{enumerate}
  \item $O_w$ चे \emph{केंद्र}~$w$
    असते: $\{w\} \in O_w$.
  \item $O_w$ मधील गोलक
    \emph{एकमेकांच्या आत} असतात:
    $S_1$, $S_2 \in O_w$ असतील,
    तर $S_1 \subseteq S_2$ किंवा
    $S_2 \subseteq S_1$; म्हणजे
    $O_w$ ला~$\subseteq$ ने
    रेषीय क्रम असतो.
  \item प्रत्येक रिक्त नसलेल्या उपकुल $\mathcal S\subseteq O_w$
    साठी $\bigcup\mathcal S\in O_w$.
  \item प्रत्येक रिक्त नसलेल्या उपकुल $\mathcal S\subseteq O_w$
    साठी $\bigcap\mathcal S\in O_w$.
    या अटी अनंत उपकुलांनाही लागू होतात.
  \end{enumerate}
\end{defn}

$O_w$ मागील कल्पना अशी:
$w$ भोवतीची जगे~$w$ पासून
किती दूर आहेत यानुसार
थरांत मांडली जातात.
रिकामे गोलक वगळल्यास सर्वांत आतला गोलक म्हणजे
एकटे~$w$, अर्थात~$\{w\}$: इतर कोणत्याही जगापेक्षा
$w$ हे स्वतःच्या अधिक जवळ असते. व्याख्या रिकामा
गोलकही परवानगी देते; त्याने पुढील सत्यता-अट बदलत नाही.
$S \subsetneq S'$ असेल,
तर $S' \setminus S$
मधील जगे~$S$ मधील
जगांपेक्षा~$w$ पासून
अधिक दूर असतात:
$S'\setminus S$ हा $S'$ च्या आत पण $S$ च्या बाहेर
असलेल्या जगांचा ‘‘थर’’ आहे.
विशेषतः, गोलकांत
त्यांच्या बाह्य पृष्ठभागाच्या
आतील सर्व जगे असतात
असे समजावे; गोलक
म्हणजे फक्त स्वतंत्र
थर नव्हेत.

\begin{figure}
\begin{center}
\begin{tikzpicture}[layerwidth=1.5,scale=.6]\tiny
  \spheresystem{3}
  \propositionintersect{0}{3}{60}{6}
  \path (0,0) node[world] {$w$};
  \spherepos{-30}{2}{node[world] {$w_2$}}
  \spherepos{220}{2}{node[world] {$w_3$}}
  \spherepos{90}{2}{node[world] {$w_1$}}
  \spherepos[shift={(.1,0)}]{10}{3}{node[world] {$w_5$}}
  \spherepos[shift={(.1,0)}]{-10}{3}{node[world] {$w_6$}}
  \spherepos{225}{3}{node[world] {$w_4$}}
  \spherepos{20}{4}{node[world] {$w_7$}}
  \path (5,0) node[left] {$p$};
\end{tikzpicture}
\caption{गोलक प्रतिमानाची आकृती}
\label{con:min:sph:fig:sphere-model}
\end{center}
\end{figure}

\ref{con:min:sph:fig:sphere-model} मधील
आकृती पुढील गोलक प्रतिमान
दाखवते: $W = \{w, w_1, \dots, w_7\}$,
$V(p) = \{w_5, w_6,
w_7\}$. सर्वांत आतला गोलक
$S_1 = \{w\}$ आहे.
$w$ च्या सर्वाधिक जवळची
जगे $w_1, w_2, w_3$
आहेत; म्हणून पुढचा मोठा
गोलक $S_2 = \{w, w_1,
w_2, w_3\}$ आहे.
त्याहून बाहेरची जगे
$w_4$, $w_5$, $w_6$
आहेत; म्हणून सर्वांत
बाहेरचा गोलक
$S_3 = \{w, w_1, \dots, w_6\}$
आहे. $w$ भोवतीची
गोलक-प्रणाली
$O_w = \{S_1, S_2, S_3\}$
आहे. जग~$w_7$ हे~$w$
भोवतीच्या कोणत्याही गोलकात
नाही. ज्या जगांत~$p$
सत्य आहे त्यांपैकी
सर्वाधिक जवळची जगे
$w_5$ आणि~$w_6$
आहेत. त्यामुळे~$p$
सत्य असलेले जग
समाविष्ट करणारा सर्वांत
लहान गोलक~$S_3$ आहे.

गोलक प्रतिमान~$\mModel M$
मधील जग~$w$ येथे
सूत्र~$\varphi$ ची पूर्ति,
$\mSat{M}{\varphi}[w]$,
याची व्याख्या करताना
मोडल सूत्रे
यांच्या व्याख्येत
$\psi \cif \chi$
साठी पुढील अट जोडतो:

\begin{defn}
  $\mSat{M}{\psi \cif \chi}[w]$
  तेव्हा आणि केवळ तेव्हाच,
  जेव्हा पुढीलपैकी एक घडते:
  \begin{enumerate}
  \item\label{con:min:sph:sphere-vac}
    सर्व $u \in \bigcup O_w$
    साठी $\mSat/{M}{\psi}[u]$; किंवा
  \item\label{con:min:sph:sphere-nonvac}
    एखाद्या $S \in O_w$ साठी,
    \begin{enumerate}
      \item एखाद्या $u \in
        S$ साठी
        $\mSat{M}{\psi}[u]$; आणि
      \item सर्व $v \in S$
        साठी, एकतर
        $\mSat/{M}{\psi}[v]$
        किंवा
        $\mSat{M}{\chi}[v]$.
    \end{enumerate}
  \end{enumerate}
\end{defn}

या व्याख्येनुसार
$\mSat{M}{\psi \cif \chi}[w]$
तेव्हा आणि केवळ तेव्हाच, जेव्हा
एकतर पूर्वांग~$\psi$
हे~$w$ भोवतीच्या सर्व
गोलकांत सर्वत्र असत्य असते,
किंवा एखाद्या गोलक~$S$
मध्ये~$\psi$ सत्य असते आणि
त्या ‘‘$\psi$ समाविष्ट
करणाऱ्या’’ गोलकातील
सर्व जगांत
$\psi \lif \chi$
हे वास्तविक अभिव्यंजन
सत्य असते.
सहज स्पष्टीकरणावरून
अपेक्षा वाटू शकते,
तरी व्याख्येत~$S$
हा~$\psi$ समाविष्ट
करणारा \emph{सर्वांत
आतला} गोलक असावा
अशी अट घातलेली नाही.
पण \ref{con:min:sph:sphere-nonvac}
मधील अट एखाद्या
गोलक~$S$ साठी
पूर्ण होत असेल,
तर~$S$ मध्ये
समाविष्ट असलेल्या आणि
पूर्वांग सत्य असलेले जग
धारण करणाऱ्या प्रत्येक
लहान गोलकासाठीही
ती पूर्ण होते.
म्हणून असा सर्वांत
आतला गोलक अस्तित्वात
असेल, तर त्यासाठीही
अट पूर्ण होते.

गोलक प्रतिमानाच्या व्याख्येनुसार $\psi$ सत्य असलेले
जग धारण करणारा सर्वांत आतला गोलक असतोच असे नाही.
केवळ अनंत उतरती क्रमिका दिल्याने तो नसल्याचे सिद्ध होत नाही;
प्रणालीत तिच्यापेक्षा लहान पूर्वांग-गोलकही असू शकतो.
एक पूर्ण उदाहरण असे: $W=\{w\}\cup\{u_n:n\geq1\}$,
$S_n=\{w\}\cup\{u_j:j\geq n\}$ आणि
$O_w=\{\{w\}\}\cup\{S_n:n\geq1\}$ घ्या.
इतर जगांसाठी $O_{u_n}=\{\{u_n\}\}$ घ्या आणि
पूर्वांग म्हणून $\psi=p$, $V(p)=\{u_n:n\geq1\}$ ठेवा.
ही प्रणाली केंद्रित, एकमेकांच्या आत असलेली आणि प्रत्येक
रिक्त नसलेल्या उपकुलाच्या संघटनाखाली व छेदाखाली बंद आहे.
गोलक $S_n$ हे सर्व $p$-जग धारण करतात,
$S_{n+1}\subsetneq S_n$ आणि $\bigcap_{n\geq1}S_n=\{w\}$,
पण $p$ हे $w$ येथे असत्य आहे. म्हणून सर्वांत लहान
$p$-जग धारण करणारा गोलक नाही.
या प्रतिमानात $\mSat{M}{p\cif \chi}[w]$ तेव्हा आणि केवळ तेव्हाच,
जेव्हा एखाद्या $i\geq1$ साठी सर्व $j\geq i$ येथे
$\mSat{M}{\chi}[u_j]$. सममूल्यपणे, अशा $i$ पासून
$S_i,S_{i+1},\dots$ या सर्व गोलकांत $p\lif \chi$ सत्य असते.
म्हणून सर्वांत लहान पूर्वांग-गोलक नसतानाही मूळ सत्यता-अट
निर्णायक आहे.













\section{प्रतिवास्तविक विधानांचे सत्य आणि असत्य}\label{con:min:tf:sec}

सर्वाधिक जवळची $\varphi$-जगे अस्तित्वात असतील आणि
ती सर्व $\psi$-जगे असतील, तर प्रतिवास्तविक
विधान $\varphi \cif \psi$
(रिक्त नसलेल्या अर्थाने) सत्य असते.
\ref{con:min:tf:fig:true} मध्ये हे दाखवले आहे.
\begin{figure}
\begin{center}
\begin{tikzpicture}[scale=.7]\small
  \spheresystem{5}
  \draw (0,0) node {$w$};
  \propositionintersect{0}{4}{40}{3.5}
  {\tikzset{proposition/.append={smooth,tension=1.4}}
  \proposition[shift={(1.6,-1)}]{90}{1}{30}{4}}
  \path (1.5,3) node[below] {$\formula{B}$};
  \path (3,0) node {$\formula{A}$};
\end{tikzpicture}
\caption{रिक्त नसलेल्या अर्थाने सत्य प्रतिवास्तविक विधान}
\label{con:min:tf:fig:true}
\end{center}
\end{figure}
जग~$w$ भोवतीच्या गोलकांच्या
प्रणालीत $\varphi$ समाविष्ट
करणारा एकही गोलक नसेल,
तरी~$w$ येथे प्रतिवास्तविक
विधान सत्य असते. अशा वेळी
ते \emph{रिक्तपणे} सत्य
असते (\ref{con:min:tf:fig:vacuous} पाहा).
\begin{figure}
\begin{center}
\begin{tikzpicture}[scale=.7]\small
  \spheresystem{5}
  \draw (0,0) node {$w$};
  \proposition{0}{6.5}{40}{3.5}
  {\tikzset{proposition/.append={smooth,tension=1.4}}
  \proposition[shift={(1.2,-1)}]{90}{1}{30}{4}}
  \path (1.5,3) node[below] {$\formula{B}$};
  \spherepos{0}{7}{node {$\formula{A}$}}
\end{tikzpicture}
\caption{रिक्तपणे सत्य प्रतिवास्तविक विधान}
\label{con:min:tf:fig:vacuous}
\end{center}
\end{figure}

सर्वाधिक जवळची $\varphi$-जगे अस्तित्वात असलेल्या या
आकृत्यांत प्रतिवास्तविक विधानाच्या असत्यतेचे दोन
प्रकार पाहू. ही मांडणी सर्वाधिक जवळचे जग नसलेल्या
प्रतिमानांची पूर्ण विभागणी असल्याचा दावा नाही.
पहिल्या प्रकारात सर्वाधिक
जवळची $\varphi$-जगे
सर्वच $\psi$-जगे नसतात,
पण त्यांपैकी काही
तशी असतात.
या वेळी $\varphi \cif
\lnot \psi$ हे विधानही
असत्य असते
(\ref{con:min:tf:fig:false} पाहा).
\begin{figure}
\begin{center}
\begin{tikzpicture}[scale=.7]\small
  \spheresystem{5}
  \draw (0,0) node {$w$};
  \propositionintersect{0}{4}{40}{3.5}
  {\tikzset{proposition/.append={smooth,tension=1.4}}
  \proposition[shift={(1.6,0)}]{90}{1}{30}{3}}
  \path (1.5,3) node[below] {$\formula{B}$};
  \path (3,0) node {$\formula{A}$};
\end{tikzpicture}
\caption{असत्य प्रतिवास्तविक विधान; विरुद्ध उत्तरांगाचे विधानही असत्य}
\label{con:min:tf:fig:false}
\end{center}
\end{figure}
सर्वाधिक जवळच्या
$\varphi$-जगांचा
$\psi$-जगांशी
अजिबात छेद नसेल,
तर $\varphi \cif \psi$
असत्य असते. पण
या वेळी सर्वाधिक
जवळची सर्व
$\varphi$-जगे
$\lnot \psi$-जगे
असतात; म्हणून
$\varphi \cif \lnot \psi$
सत्य असते
(\ref{con:min:tf:fig:false-opposite} पाहा).
\begin{figure}
\begin{center}
\begin{tikzpicture}[scale=.7]\small
  \spheresystem{5}
  \draw (0,0) node {$w$};
  \propositionintersect{0}{4}{40}{3.5}
  {\tikzset{proposition/.append={smooth,tension=1.4}}
  \proposition[shift={(-1,0)}]{90}{1}{30}{3}}
  \path (-1,3) node[below] {$\formula{B}$};
  \path (3,0) node {$\formula{A}$};
\end{tikzpicture}
\caption{असत्य प्रतिवास्तविक विधान; विरुद्ध उत्तरांगाचे विधान सत्य}
\label{con:min:tf:fig:false-opposite}
\end{center}
\end{figure}

S5 मधील कठोर अभिव्यंजनाच्या सत्यतेच्या अवश्यतेशी
तुलना करता, प्रतिवास्तविक
विधाने परिस्थितीवश
असू शकतात.
\ref{con:min:tf:fig:contingent} मधील
गोलक प्रतिमान पाहा.
$u$ च्या सर्वाधिक
जवळची $\varphi$-जगे
ही सर्व $\psi$-जगे
आहेत; म्हणून
$\mSat{M}{\varphi \cif
  \psi}[u]$. पण~$v$
च्या सर्वाधिक जवळच्या
काही $\varphi$-जगांत
$\psi$ सत्य नाही;
म्हणून
$\mSat/{M}{\varphi \cif \psi}[v]$.
\begin{figure}
\begin{center}
\begin{tikzpicture}[scale=.6]\small
  \spheresystem{7}
  \spheresystem[shift={(0:3.1)},dashed]{4}
  \propositionintersect{45}{4}{40}{4.5}
  \begin{scope}
    \clip \propositionplot{45}{4}{40}{4.5} ;
    \spherelayer[shift={(0:3.1)}]{3}
  \end{scope}
  \proposition[shift={(1.4,-.5)}]{90}{1}{35}{5}
  \path (0,0) node {$u$};
  \path (0:3.1) node {$v$};
  \spherepos{35}{9}{node {$\formula{A}$}}
  \spherepos{80}{9}{node {$\formula{B}$}}
\end{tikzpicture}
\end{center}
\caption{परिस्थितीवश प्रतिवास्तविक विधान}
\label{con:min:tf:fig:contingent}
\end{figure}











\section{पूर्वांग बळकट करणे}\label{cnt:min:agg:sec}

‘‘पूर्वांग बळकट करणे’’ म्हणजे
$\varphi \lif \chi \Entails
(\varphi \land \psi) \lif \chi$
हे अनुमान. वास्तविक
अभिव्यंजनासाठी ते वैध,
पण प्रतिवास्तविक
विधानांसाठी अवैध आहे.
मी ही आगकाडी ओढली असती,
तर ती पेटली असती,
हे सत्य आहे असे समजा.
(म्हणजे आगकाडीत किंवा
आगपेटीच्या घासण्याच्या
पृष्ठभागात काही दोष नाही,
मी आगकाडी मोडणार नाही,
इत्यादी.) पण ही
आगकाडी मी बाह्य
अवकाशात ओढली असती,
तर ती पेटली असती,
हे सत्य नाही. म्हणून
पुढील अनुमान अवैध आहे:
\begin{quote}
  आगकाडी ओढली असती,
  तर ती पेटली असती.

  म्हणून आगकाडी बाह्य
  अवकाशात ओढली असती,
  तर ती पेटली असती.
\end{quote}

लुईस--स्टालनाकर यांच्या
सशर्त विधानांच्या
विश्लेषणातून हे
समजते: मी आगकाडी
बाह्य अवकाशात
ओढतो असे सर्वाधिक
जवळचे जग, मी
आगकाडी ओढतो
अशा सर्वाधिक
जवळच्या जगापेक्षा
वास्तविक जगापासून
खूप दूर आहे.
म्हणून दुसऱ्या
जगात आगकाडी
पेटते हे सत्य
असले, तरी
पहिल्या जगात
ती पेटत नाही.
असेच असायला हवे.

\begin{ex}
  गोलक चिन्हार्थमीमांसेत हे अनुमान
  अवैध ठरते; म्हणजे
  $p \cif r \Entails/
  (p \land q) \cif r$.
  $\mModel{M} = \tuple{W, O, V}$
  हे प्रतिमान पाहा: येथे
  $W = \{w, w_1, w_2\}$,
  $O_w = \{\{w\}, \{w,
  w_1\}, \{w, w_1, w_2\}\}$,
  $V(p) = \{w_1, w_2\}$,
  $V(q) = \{w_2\}$, आणि
  $V(r) = \{w_1\}$.
  इतर जगांसाठी $O_{w_1}=\{\{w_1\}\}$ आणि
  $O_{w_2}=\{\{w_2\}\}$ घ्या; त्यामुळे $O$ हे
  सर्व $W$ वर परिभाषित आहे.
  $p$ समाविष्ट करणारा
  $S = \{w, w_1\}$
  हा गोलक आहे आणि
  त्यातील सर्व जगांत
  $p \lif r$ सत्य आहे;
  म्हणून
  $\mSat{M}{p \cif r}[w]$.
  तसेच $(p \land q)$
  समाविष्ट करणारा
  $S' = \{w, w_1, w_2\}$
  हा गोलक आहे; पण
  $\mSat/{M}{(p \land q) \lif r}[w_2]$.
  $S'$ हा या $O_w$ मधील एकमेव $(p\land q)$-जग
  धारण करणारा गोलक आहे आणि त्यात अभिव्यंजन असत्य आहे.
  म्हणून कोणताही साक्षी-गोलक नाही; रिक्तपणे सत्य होण्याची
  अटही लागू होत नाही. त्यामुळे
  $\mSat/{M}{(p \land q) \cif r}[w]$
  (\ref{cnt:min:agg:fig:antecedent-strengthening} पाहा).

\begin{figure}
\begin{center}
\begin{tikzpicture}[layerwidth=1.5,scale=.6]\tiny
  \spheresystem{3}
  \propositionintersect{-10}{3}{30}{5}
  \begin{scope}
    \clip plot[smooth,tension=1] coordinates
          {(0,4.5) (30:.95) (4.5,-3.5)} -- (4.5,4.5) ;
    \spherefill{2}
  \end{scope}
  \draw[proposition] plot[dashed,smooth,tension=1] coordinates
          {(0,4.5) (30:.95) (4.5,-3.5)} ;
  \proposition{30}{2}{30}{5}
  \path (0,0) node[world] {$w$};
  \spherepos{30}{2}{node[world] {$w_1$}}
  \spherepos{-10}{3}{node[world] {$w_2$}}
  \spherepos{-10}{4}{node {$q$}}
  \spherepos{30}{4}{node {$r$}}
  \draw (1,4.5) node[below] {$p$};
\end{tikzpicture}
\caption{पूर्वांग बळकट करणे या नियमाचे प्रतिउदाहरण}
\label{cnt:min:agg:fig:antecedent-strengthening}
\end{center}
\end{figure}
\end{ex}












\section{संक्रमकता}\label{cnt:min:tra:sec}

वास्तविक अभिव्यंजनासाठी
साखळी-नियम लागू होतो:
$\varphi \lif \psi, \psi
\lif \chi \Entails \varphi \lif \chi$.
दुसऱ्या शब्दांत,
वास्तविक अभिव्यंजन
संक्रमक आहे.
प्रतिवास्तविक विधानांसाठीही
हे सत्य आहे का?
स्टालनाकर यांनी दिलेले
पुढील उदाहरण पाहा.
\begin{quote}
  जे.~एडगर हूव्हर रशियात
  जन्मले असते, तर ते
  कम्युनिस्ट झाले असते.

  जे.~एडगर हूव्हर कम्युनिस्ट
  असते, तर ते देशद्रोही
  झाले असते.

  म्हणून जे.~एडगर हूव्हर
  रशियात जन्मले असते,
  तर ते देशद्रोही
  झाले असते.
\end{quote}
हूव्हर यांचा जन्म
(प्रत्यक्षात ज्या वेळी झाला
त्याच वेळी) अमेरिकेत
न होता रशियात झाला
असता, तर ते
सोव्हिएत संघात वाढले
असते आणि कम्युनिस्ट
झाले असते, असे
समजू या. म्हणून
पहिले गृहीतक सत्य
आहे. त्याचप्रमाणे,
दुसरे गृहीतक
स्वतंत्रपणे पाहिल्यास
सत्य आहे. पण
निष्कर्ष असत्य आहे:
हूव्हर यांचा जन्म
सोव्हिएत संघात
झाला असता, तर
ते बहुधा निष्ठावान
कम्युनिस्ट झाले
असते; आपल्या
देशाशी विश्वासघात
केला नसता.
सत्य-मूल्यांची ही
सहज वाटणारी
नेमणूक स्टालनाकर--लुईस
यांच्या विश्लेषणातूनही
मिळते. हूव्हर यांच्या
जन्मस्थळातील बदल आणि त्याचे आवश्यक परिणाम
यांखेरीज अनावश्यक बदल नसलेले आपल्या जगाच्या
सर्वाधिक जवळचे
संभाव्य जग म्हणजे
हूव्हर सोव्हिएत
संघाचे चांगले नागरिक
म्हणून वाढतात
असे जग. पहिल्या
गृहीतकाचे आणि
निष्कर्षाचे पूर्वांग
सत्य असलेले हे
सर्वाधिक जवळचे
संभाव्य जग आहे.
तेथे हूव्हर
कम्युनिस्ट पक्षाचे
निष्ठावान सदस्य
असतात, म्हणून
देशद्रोही नसतात.
पण दुसरे गृहीतक
तपासताना वेगळे
जग पाहावे लागते:
हूव्हर कम्युनिस्ट
असलेले सर्वाधिक
जवळचे जग असे
की ते अमेरिकेत
जन्मले, नंतर
कम्युनिस्ट झाले,
आणि त्यामुळे
देशद्रोही ठरले.\footnote{
  अर्थात या उदाहरणाचा
  मुद्दा समजण्यासाठी
  काही सत्ताविषयक
  आणि राजकीय
  गृहीतके मानावी
  लागतात. उदाहरणार्थ,
  हूव्हर यांचा जन्म
  रशियन पालकांपासून
  होणे शक्य होते,
  किंवा १९५० च्या
  दशकातील अमेरिकेतील
  कम्युनिस्ट हे आपल्या
  देशाचे देशद्रोही होते,
  अशी गृहीतके.}

\begin{prob}
  प्रतिवास्तविक विधानांची
  संक्रमकता अपयशी
  ठरते याचे पटण्याजोगे,
  सहज उदाहरण शोधा.
\end{prob}

\begin{ex}\label{cnt:min:tra:ex:trans-counterex}
  गोलक चिन्हार्थमीमांसेत हे
  अनुमान अवैध ठरते;
  म्हणजे
  $p \cif q, q \cif r
  \Entails/ p \cif r$.
  $\mModel{M} = \tuple{W, O, V}$
  हे प्रतिमान पाहा: येथे
  $W = \{w, w_1, w_2\}$,
  $O_w = \{\{w\}, \{w,
  w_1\}, \{w, w_1, w_2\}\}$,
  $V(p) = \{w_2\}$,
  $V(q) = \{w_1, w_2\}$,
  आणि $V(r) = \{w_1\}$.
  तसेच $O_{w_1}=\{\{w_1\}\}$ आणि $O_{w_2}=\{\{w_2\}\}$
  घ्या, म्हणजे सर्व जगांचे गोलक-फलन पूर्ण होते.
  $p$ समाविष्ट करणारा
  $S = \{w, w_1, w_2\}$
  हा गोलक आहे आणि
  त्यातील सर्व जगांत
  $p \lif q$ सत्य आहे;
  म्हणून
  $\mSat{M}{p \cif q}[w]$.
  तसेच $q$ समाविष्ट
  करणारा
  $S' = \{w, w_1\}$
  हा गोलक आहे आणि
  त्यातील सर्व जगांत
  $\mSat{M}{q \lif r}$
  सत्य आहे; म्हणून
  $\mSat{M}{q \cif r}[w]$.
  पण $p$ समाविष्ट
  करणाऱ्या
  $\{w, w_1, w_2\}$
  गोलकात~$w_2$
  हे जग आहे,
  जिथे
  $\mSat/{M}{p \lif
    r}[w_2]$. हा $O_w$ मधील एकमेव $p$-जग धारण करणारा
  गोलक आहे; म्हणून $\mSat/{M}{p\cif r}[w]$.
  दोन्ही गृहीतके सत्य असून निष्कर्ष असत्य आहे.
\end{ex}

\begin{prob}
  \ref{cnt:min:tra:ex:trans-counterex}
  मधील प्रतिउदाहरणाला
  अनुरूप गोलकांची
  आकृती काढा.
\end{prob}

\begin{prob}
  \ref{cnt:min:tra:ex:trans-counterex}
  मध्ये जग~$w_2$
  येथे हूव्हर यांचा
  जन्म रशियात झाला
  आहे, ते कम्युनिस्ट
  आहेत आणि देशद्रोही
  नाहीत. जग~$w_1$
  येथे त्यांचा जन्म
  अमेरिकेत झाला
  आहे, ते कम्युनिस्ट
  आहेत आणि देशद्रोही
  आहेत. या प्रतिमानात
  $w_1$ हे~$w_2$
  पेक्षा~$w$ च्या
  अधिक जवळ आहे.
  हे आवश्यक आहे
  का? हूव्हर यांचा
  रशियात जन्म होणे
  हे त्यांच्या
  कम्युनिस्ट असण्यापेक्षा
  अधिक दूरची
  शक्यता आहे असे
  न मानणारे
  प्रतिउदाहरण
  देता येईल का?
\end{prob}












\section{प्रतिपरिवर्तन}\label{cnt:min:cpo:sec}

वास्तविक आणि कठोर
अभिव्यंजने त्यांच्या
प्रतिपरिवर्तनांशी सममूल्य
असतात. प्रतिवास्तविक
विधानांच्या बाबतीत
तसे नाही. क्रात्सर
यांनी दिलेले उदाहरण
पाहा:
\begin{quote}
  गटे यांचा मृत्यू १८३२ मध्ये
  झाला नसता, तरी आता
  ते मृतच असते.

  गटे आता मृत नसते,
  तर त्यांचा मृत्यू
  १८३२ मध्ये झाला असता.
\end{quote}
पहिले वाक्य सत्य आहे:
माणसे शेकडो वर्षे
जगत नाहीत. दुसरे
स्पष्टपणे असत्य आहे:
गटे आता मृत नसते,
तर ते अजून जिवंत
असते; त्यामुळे
त्यांचा मृत्यू
१८३२ मध्ये
झालेला नसता.

\begin{ex}\label{cnt:min:cpo:ex:contraposition-counterex}
  गोलक चिन्हार्थमीमांसेत
  प्रतिपरिवर्तन अवैध ठरते;
  म्हणजे
  $p \cif q \Entails/
  \lnot q \cif \lnot p$.
  $p$ चा अर्थ ‘‘गटे
  यांचा मृत्यू १८३२
  मध्ये झाला नाही’’
  आणि~$q$ चा अर्थ
  ‘‘गटे आता मृत
  आहेत’’ असा घ्या.
  हे पुढील प्रतिमानात
  दाखवता येते:
  $\mModel{M_1} = \tuple{W, O, V}$,
  येथे $W =
  \{w, w_1, w_2\}$,
  $O_w = \{\{w\}, \{w, w_1\},
  \{w, w_1, w_2\}\}$,
  $V(p) = \{w_1, w_2\}$
  आणि $V(q) = \{w, w_1\}$.
  $O_{w_1}=\{\{w_1\}\}$ आणि $O_{w_2}=\{\{w_2\}\}$ घ्या;
  $O$ आता संपूर्ण $W$ वर परिभाषित आहे.
  म्हणून~$w$ हे वास्तविक
  जग आहे, ज्यात गटे
  यांचा मृत्यू १८३२
  मध्ये झाला आणि
  ते अजूनही मृत आहेत;
  $w_1$ हे (जवळचे)
  जग आहे, ज्यात
  त्यांचा मृत्यू,
  समजा, १८३३ मध्ये
  झाला आणि ते
  अजूनही मृत आहेत;
  तर~$w_2$ हे
  (दूरचे) जग आहे,
  ज्यात गटे
  अजून जिवंत आहेत.
\begin{figure}
\begin{center}
\begin{tikzpicture}[scale=.6,layerwidth=1.5]\tiny
  \spheresystem{3}
  \begin{scope}
    \clip \propositionplot{0}{.8}{50}{4.5} -- (4.5,-4.5) -- (-4.5,-4.5) -- (-4.5,4.5) -- (4.5,4.5);
    \spherelayer{3}
  \end{scope}
  \propositionintersect{0}{2}{110}{5.5}
  \proposition{0}{.8}{50}{4.5}
  \path (0,0) node[world] {$w$};
  \spherepos{0}{2}{node[world] {$w_1$}}
  \spherepos{-50}{3}{node[world] {$w_2$}}
  \spherepos{28}{3}{node {$q$}}
  \spherepos{42}{3}{node {$\lnot q$}}
  \spherepos{-55}{4}{node {$p$}}
  \spherepos{-67}{4}{node {$\lnot p$}}
\end{tikzpicture}
\caption{प्रतिपरिवर्तनाचे प्रतिउदाहरण}
\label{cnt:min:cpo:fig:contraposition}
\end{center}
\end{figure}
$p$ समाविष्ट करणारा
$S = \{w, w_1\}$
हा गोलक आहे आणि
त्यातील सर्व जगांत
$p \lif q$ सत्य आहे;
म्हणून
$\mSat{M_1}{p \cif q}[w]$.
पण $\lnot q$
समाविष्ट करणाऱ्या
$\{w, w_1, w_2\}$
गोलकात~$w_2$
हे जग आहे; तेथे
$q$ असत्य आणि
$p$ सत्य आहे.
म्हणून
$\mSat/{M_1}{\lnot q
\lif \lnot p}[w_2]$. हा $O_w$ मधील एकमेव $\lnot q$-जग
धारण करणारा गोलक असल्याने
$\mSat/{M_1}{\lnot q\cif\lnot p}[w]$.
मूळ प्रतिवास्तविक विधान सत्य असून त्याचे प्रतिपरिवर्तन असत्य आहे.
\end{ex}



\part{संच उपपत्ती}\label{sth:::part}\label{sth:part}
\chapter{संचांची क्रमिक संकल्पना}\label{sth:story:chap}




\section{विस्तारात्मकता}\label{sth:story:extensionality:sec}

सर्वांत आधी सांगायचे म्हणजे
संचांची ओळख त्यांच्या
घटक वरून ठरते.
अधिक नेमकेपणाने:

\begin{axiom}[विस्तारात्मकता]
$A$ आणि~$B$ या संचांत
तेच घटक
असतील, तर~$A$
आणि~$B$ हा
एकच संच आहे.
\[
  \lforall[A][\lforall[B][(\lforall[x][(x \in A \liff x \in B)] \lif
  \eq[A][B])]]
\]
\end{axiom}

\ref{sfr:::part}
या संपूर्ण भागात
आपण हे गृहीत धरले
होते. तरी ते
पुन्हा सांगण्यासारखे
आहे. विस्तारात्मकतेचे
स्वयंसिद्धक संचाची
ओळख त्याच्या
घटक वरून
ठरते ही मूलभूत
कल्पना व्यक्त
करते. (म्हणून
संचांची तुलना
\emph{संकल्पनांशी}
करता येते:
अगदी त्याच
वस्तू अनेक
भिन्न संकल्पनांच्या
अंतर्गत येऊ
शकतात.)

हे तत्त्व का
स्वीकारावे? कारण
विस्तारात्मकता
नाकारणे म्हणजे
\emph{संच उपपत्ती}
म्हणता येईल
अशा गोष्टीचाच
त्याग करणे,
असे म्हणणे
विश्वसनीय वाटते.
संच उपपत्ती
म्हणजे विस्तारात्मक
संग्रहांचीच
उपपत्ती.

पण \ref{sth:::part}
मध्ये खरे आव्हान
\emph{कोणते संच
अस्तित्वात आहेत}
हे सांगणारी
तत्त्वे ठरवणे
आहे. या प्रश्नाचे
एकमेव खरोखर
‘‘उघड’’ वाटणारे
उत्तर सिद्धपणे
चुकीचे ठरते.







\section{रसेलची विरोधापत्ती (पुन्हा)}\label{sth:story:rus:sec}

\ref{sfr:::part}
मध्ये आपण अनौपचारिक
संचसिद्धान्त वापरला.
पण एका \emph{अतिशय}
भोळ्या कल्पनेनुसार
संच म्हणजे विधेयांचे
विस्तारच असतात.
या कल्पनेतून पुढील
तत्त्व स्वीकारावे लागेल:

\begin{defish}
  \emph{अनौपचारिक
  गुणधर्माधारित संचरचना.}
  कोणत्याही सूत्र~$\phi$
  साठी
  $\Setabs{x}{\phi(x)}$
  अस्तित्वात असतो.
\end{defish}

हे तत्त्व मोहक
असले, तरी ते
सिद्धपणे विसंगत
आहे. \ref{sfr:set:rus:sec}
मध्ये आपण हे
पाहिले आहे.
पण निष्कर्ष
इतका महत्त्वाचा
आणि इतका
सरळ आहे
की तो
शब्दशः पुन्हा
मांडण्यासारखा आहे.

\begin{thm}[रसेलची विरोधापत्ती]
$R = \Setabs{x}{x \notin x}$
असा कोणताही संच नाही.
\end{thm}

\begin{proof}
$R = \Setabs{x}{x \notin x}$
अस्तित्वात असेल, तर
$R \in R$ तेव्हा
आणि केवळ तेव्हाच
$R \notin R$;
हा व्याघात आहे.
\end{proof}

रसेल यांनी हा
निष्कर्ष जून १९०१
मध्ये शोधला. (मात्र
त्यांनी विरोधापत्ती
आपण आत्ताच
मांडलेल्या रूपात
दिलेली नव्हती,
कारण ते
\emph{Grundgesetze}
मध्ये मांडलेल्या
फ्रेग यांच्या
संच उपपत्तीचा
विचार करत होते.
\ref{sth:story:blv:sec}
मध्ये आपण
याकडे परत
येऊ.) फ्रेग
यांच्या प्रणालीतील
विसंगती स्पष्ट
करण्यासाठी रसेल
यांनी त्यांना
१६ जून १९०२
रोजी पत्र
लिहिले. त्या
पत्रव्यवहारासाठी
आणि थोड्या
पार्श्वभूमीसाठी
van Heijenoort (1967, पृ.~124--8)
पाहा.

या दोन ओळींची
सिद्धता \emph{शुद्ध
तर्कशास्त्रातून}
मिळते, यावर
भर द्यायला हवा.
मान्य आहे की
$\Setabs{x}{x \notin x}$
या लेखनात
आपण विस्तारात्मकतेचे
एक (तार्किकेतर?)
स्वयंसिद्धक अप्रत्यक्षपणे
वापरले आहे:
$\Setabs{x}{\phi(x)}$
म्हणजे, असा संच
असेल तर,
$\phi$ पूर्ण
करणाऱ्या वस्तूंचा
\emph{एकमेव}
(विस्तारात्मकतेमुळे)
संच. पण निष्कर्ष
पुढीलप्रमाणे
मांडल्यास
विस्तारात्मकतेचा
आभासही टाळता
येतो: \emph{ज्या
संचांचे स्वतःत
सदस्यत्व नाही,
नेमके तेच
ज्याचे सदस्य
आहेत असा
संच नाही}.
याचा संचांशीही
फारसा विशेष
संबंध नाही.
रसेल यांनी
स्वतः नमूद
केल्याप्रमाणे,
अगदी तसाच
युक्तिवाद
दाखवतो:
\emph{जे पुरुष
स्वतःची दाढी
करत नाहीत,
नेमक्या त्याच
पुरुषांची दाढी
करणारा पुरुष
नाही}. किंवा:
\emph{जी पग
जातीची कुत्री
स्वतःचा वास
घेत नाहीत,
नेमक्या त्याच
पग कुत्र्यांचा
वास घेणारा
पग कुत्रा
नाही}. इत्यादी.
सांकेतिक रूपात
हा निष्कर्ष
इतकाच:
\[
\lnot \exists x \forall z(Rzx \liff \lnot R zz).
\]
हे प्रथम-क्रम
तर्कशास्त्राचे
प्रमेय(-रूपबंध)
आहे. त्यामुळे
संच उपपत्तीत
फक्त फेरबदल
करून रसेलची
विरोधापत्ती
टाळता येत
नाही: संच
उपपत्तीपर्यंत
\emph{पोहोचण्याआधीच}
ती उद्भवते.
आपण (अभिजात)
प्रथम-क्रम
तर्कशास्त्र
वापरणार असू,
तर
$R = \Setabs{x}{x\notin x}$
असा संच नाही
हे \emph{स्वीकारावेच}
लागते.

सार असा:
विसंगती
\emph{टाळूनही}
अनौपचारिक
गुणधर्माधारित
संचरचना
स्वीकारायची
असेल, तर
केवळ \emph{संच
उपपत्तीत}
किरकोळ फेरबदल
पुरेसे नाहीत.
त्याऐवजी
\emph{तर्कशास्त्र}
मुळापासून
बदलावे लागेल.

अर्थात,
अभिजातेतर
तर्कशास्त्रे
वापरणाऱ्या
संच उपपत्ती
मांडल्या गेल्या
आहेत. पण
त्या, सौम्य
शब्दांत सांगायचे
तर, मानक
नाहीत.
रसेलच्या
विरोधापत्तीला
सामोरे जाण्याची
मानक पद्धत
म्हणजे तिला
संच अस्तित्वात
नसल्याची
सरळ सिद्धता
मानणे, आणि
मग या निष्कर्षासह
काम कसे करावे
हे शिकणे.
आपण हीच
पद्धत वापरू.









\section{स्वसमावेशरहित आणि स्वसमावेशक व्याख्या}\label{sth:story:predicative:sec}

रसेलचा संच~$R$ हा
$\Setabs{x}{x \notin
x}$ वापरून परिभाषित
केला होता. अधिक
स्पष्टपणे सांगायचे
तर~$R$ म्हणजे
स्वतःचे सदस्य
नसलेले, आणि
फक्त असेच,
सर्व संच सदस्य
म्हणून असलेला
संच ठरला असता.
म्हणून~$R$ ची
व्याख्या करताना
आपण~$R$
(अस्तित्वात असता)
ज्या प्रांतात
आला असता
त्या प्रांतावर
परिमाणन करतो.

ही \emph{स्वसमावेशक}
व्याख्या आहे.
सर्वसाधारणपणे,
परिभाषित होणारी
वस्तू ज्या
प्रांतात समाविष्ट
आहे त्या
प्रांतावर
परिमाणन करणारी
व्याख्या तेव्हा
आणि केवळ तेव्हाच
स्वसमावेशक
असते असे
म्हणता येईल.

या विरोधापत्तींच्या
पार्श्वभूमीवर
व्हाइटहेड, रसेल,
प्वाँकारे आणि
वाइल यांनी
अशा स्वसमावेशक
व्याख्यांना
‘‘दुष्टचक्र निर्माण
करणाऱ्या’’ म्हणून
नाकारले:
\begin{quote}
  टाळायच्या
  विरोधापत्तींचे
  विश्लेषण
  दाखवते की
  त्या सर्व
  एका प्रकारच्या
  दुष्टचक्रातून
  निर्माण होतात.
  एखाद्या
  वस्तुसंग्रहात
  असे सदस्य
  असू शकतात
  जे फक्त
  त्या संपूर्ण
  संग्रहाच्या
  साहाय्यानेच
  परिभाषित
  करता येतात,
  असे मानल्याने
  संबंधित दुष्टचक्रे
  तयार होतात
  [\ldots. \textparagraph]

  अनुचित समग्रता
  टाळू देणारे
  तत्त्व असे
  मांडता येते:
  ‘संग्रहातील
  \emph{सर्वांचा}
  उल्लेख ज्यात
  येतो, ती
  वस्तू त्या
  संग्रहातील
  एक वस्तू
  असू नये’;
  किंवा उलट:
  ‘एखाद्या
  संग्रहाला
  समग्रता असती
  तर त्यातील
  काही सदस्य
  फक्त त्या
  समग्रतेच्या
  साहाय्यानेच
  परिभाषित
  करता आले
  असते, तर
  त्या संग्रहाची
  समग्रता
  अस्तित्वात
  नाही.’ याला
  आपण
  ‘दुष्टचक्र
  तत्त्व’
  म्हणू,
  कारण
  अनुचित
  समग्रता
  गृहीत
  धरल्याने
  उद्भवणारी
  दुष्टचक्रे
  ते टाळू
  देते.
  (Whitehead आणि Russell, 1910, पृ.~37)
\end{quote}
त्यांच्याप्रमाणे
\emph{दुष्टचक्र तत्त्व}
स्वीकारले, तर
\ref{sth:story:rus:sec}
मधील विनाशकारी
अनौपचारिक
गुणधर्माधारित
संचरचना-रूपबंधाऐवजी
पुढील प्रकारचे
तत्त्व वापरण्याचा
प्रयत्न करता
येईल:

\begin{defish}
\emph{स्वसमावेशरहित
गुणधर्माधारित
संचरचना.}
फक्त आधीच्या स्तरातील संचांवर परिमाणन करणाऱ्या
प्रत्येक सूत्र~$\phi$ साठी, त्याच स्तरातील $x$ घेऊन बनणारा
संच$^\prime$ $\Setabs{x}{\phi(x)}$ पुढील स्तरात अस्तित्वात असतो.
त्याचे सदस्य आधीच्या स्तरातील संचच असतात.
ही स्तरांची अनौपचारिक रूपरेषा आहे;
पूर्ण औपचारिक भाषा आणि स्वयंसिद्धे येथे दिलेली नाहीत.
\end{defish}

संच$^{\prime}$ यांची निर्मिती आधीच्या स्तरातील संचांपासून
केल्यास नव्याने परिभाषित होणारी वस्तू त्या सदस्य-प्रांतात येत नाही.
त्यामुळे वरचा रसेलचा स्वसदस्यत्वाचा युक्तिवाद तसाच लागू होत नाही.
एवढ्यावरून कोणत्याही संपूर्ण औपचारिक उपपत्तीची
सुसंगतता सिद्ध होत नाही; त्यासाठी तिची अचूक भाषा,
स्वयंसिद्धे आणि स्वतंत्र सुसंगतता-युक्तिवाद आवश्यक आहेत.

दुर्दैवाने,
स्वसमावेशरहित
गुणधर्माधारित
संचरचना फारशी
\emph{व्यापक}
नाही. कारण
तिच्यामुळे
संच$^\prime$
या नव्या
वस्तू मिळतात.
म्हणून संच$^\prime$
यांच्यावर
परिमाणन करणारी
सूत्रेही विचारात
घ्यावी लागतील.
त्या सूत्रांसाठी सदस्य-प्रांत आणि नवीन वस्तूंचा स्तर
पुन्हा एकच केला आणि त्या स्तरात अमर्याद गुणधर्माधारित
संचरचना मानली, तर स्वतःचे
सदस्य नसलेल्या
सर्व संच$^\prime$
यांचा संच$^\prime$
पाहताच रसेलची
विरोधापत्ती
पुन्हा उद्भवेल.
त्यामुळे याच
विचारानुसार
संच$^\prime$
यांच्यावर
परिमाणन करणारे
सूत्र अनुरूप
संच$^{\prime\prime}$
देते असे
म्हणावे लागेल.
पुढे
संच$^{\prime\prime\prime}$,
संच$^{\prime\prime\prime\prime}$
इत्यादी लागतील.
अशी कितीतरी
प्राइम चिन्हे
टाळण्यासाठी
यांना संच$_0$,
संच$_1$,
संच$_2$,
संच$_3$,
संच$_4$,\ldots
असे मानणे
सोपे ठरेल.
यातून (सरळ)
प्रकार-उपपत्तीचा
मार्ग मिळतो.

अशा उपपत्तीवर
काही उघड
आक्षेप आहेत
(मात्र ते
\emph{निर्णायक}
आहेतच असे
स्पष्ट नाही).
थोडक्यात:
तिचा वापर
जड जातो;
ती अनेक
भिन्न प्रकारच्या
वस्तू गृहीत
धरते; आणि
शेवटी
स्वसमावेशक
व्याख्या
\emph{खरोखर
इतक्या वाईट}
आहेत का,
हेही स्पष्ट
नाही.

शेवटचा मुद्दा
स्पष्ट करण्यासाठी
रॅम्झी
यांची पुढील
टिप्पणी पाहा:
\begin{quote}
  एखाद्या
  गटातील
  सर्वांत
  उंच माणूस
  म्हणून आपण
  एखाद्याचा
  उल्लेख करू
  शकतो. असे
  केल्याने तो
  स्वतः ज्या
  समग्रतेचा
  सदस्य आहे
  तिच्याच
  साहाय्याने
  त्याची ओळख
  पटते; तरीही
  त्यात दुष्टचक्र
  नसते.
  (Ramsey, 1925)
\end{quote}
रॅम्झी यांचा
मुद्दा असा
की ‘‘गटातील
सर्वांत उंच
माणूस’’ ही
\emph{स्वसमावेशक}
व्याख्या आहे;
पण ती
स्पष्टपणे
निर्दोष आहे.

त्यावर असे
उत्तर देता
येईल की
या उदाहरणात
सर्वांत उंच
व्यक्तीला
\emph{स्वसमावेशरहित}
मार्गानेही
ओळखता येईल.
उदाहरणार्थ,
आपण त्या
माणसाकडे
सरळ बोट
दाखवू
शकतो. मग
स्वसमावेशक
व्याख्यांवरील
आक्षेप फक्त
\emph{स्वसमावेशक
मार्गानेच}
ओळखता येणाऱ्या
वस्तूंपुरता
मर्यादित करावा
लागेल. पण
अशा
‘‘मूलभूत
स्वसमावेशकतेत’’
नेमके वाईट
काय आहे,
हेही पुढे
स्पष्ट करावे
लागेल.\footnote{
  अधिक माहितीसाठी
  Linnebo (2010)
  पाहा.}

वस्तू
\emph{निर्माण}
करण्यासाठी
आपल्या
व्याख्यांनी
कृतीसूचना
द्यावी असे
अपेक्षित
असेल, तर
स्वसमावेशक
व्याख्या
निश्चितच
मोठी अडचण
ठरतात.
कारण अशा
व्याख्येवरून
वस्तू बनवताना
खरेच
दुष्टचक्रात
अडकावे
लागेल:
स्वसमावेशकपणे
निर्दिष्ट
वस्तू निर्माण
करण्यासाठी
\emph{आधी} सर्व
वस्तू
(त्यात ती
निर्दिष्ट
वस्तूही येते)
निर्माण कराव्या
लागतील,
कारण तिची
व्याख्या
सर्व वस्तूंच्या
संदर्भात आहे.
म्हणजे त्या
वस्तूची
निर्मिती
होण्यापूर्वीच
तिची निर्मिती
करावी लागेल.
पण पुन्हा,
‘‘मूलभूत
स्वसमावेशक’’
रीतीने
निर्दिष्ट
संचांवरील
हा गंभीर
आक्षेप
तेव्हाच
ठरतो,
जेव्हा
संच म्हणजे
आपण
\emph{निर्माण}
करत असलेल्या
वस्तू आहेत
असे मानतो.
आपण
(बहुधा)
तसे मानत
नाही.

त्यामुळे,
चांगल्या किंवा
वाईट कारणांनी,
रूढ झालेल्या
पद्धतीत
(अ)स्वसमावेशकतेबाबत
कठोर भूमिका
घेतलेली नाही.
त्याऐवजी
आता संचयी-क्रमिक
पद्धत मानली
जाते. शेवटी
तिच्यामुळे
संचांचे
‘‘टप्प्यां’’मध्ये
स्तरीकरण करता
येते---स्वसमावेशरहित
पद्धतीत वस्तू
संच$_0$,
संच$_1$,
संच$_2$,
\ldots मध्ये
स्तरीकृत
होतात तसेच
\emph{थोड्याफार}
प्रमाणात.
मात्र या
टप्प्यांतील
संचांच्या
मूलभूत प्रकारांत
कोणताही
फरक गृहीत
धरला जात
नाही.








\section{संचयी-क्रमिक पद्धत}\label{sth:story:approach:sec}

संचांचे टप्प्यांमध्ये
स्तरीकरण कसे
करायचे, याचे
थोडे अधिक
तपशीलवार
वर्णन असे:
\begin{quote}
  संच \emph{टप्प्यांमध्ये}
  तयार होतात.
  प्रत्येक टप्पा~$S$
  यासाठी, $S$ च्या \emph{आधी}
  कोणते टप्पे येतात, हे निश्चित असते.
  $S$ टप्प्यावर,
  $S$ पूर्वीच्या
  टप्प्यांवर
  तयार झालेले
  संच घेऊन
  बनलेल्या
  प्रत्येक
  संग्रहाचा
  संच तयार
  होतो.
  टप्प्यांवर
  तयार झालेले
  संच सोडून
  इतर संच
  नाहीत.
  (Shoenfield, 1977, पृ.~323)
\end{quote}
ही \emph{संचांची
संचयी-क्रमिक
संकल्पना}
याची रूपरेषा
आहे. \ref{sth:::part}
मध्ये मांडल्या
जाणाऱ्या
औपचारिक
संच उपपत्तीला
ती आधार
देईल.

थोडा अधिक
तपशील पाहू.
शोनफिल्ड
यांच्या
वर्णनानुसार,
प्रत्येक
टप्प्यावर
मागील
टप्प्यांपासून
आपल्याला
उपलब्ध
असलेल्या
संचांपासून
नवे संच
तयार होतात.
म्हणून
आरंभीच्या
टप्प्यावर,
म्हणजे
टप्पा~$0$
येथे,
\emph{आधीचा}
टप्पाच
नसतो;
\emph{मग तर}
पूर्वीच्या
टप्प्यांतून
आपल्याला
संचही
उपलब्ध
नसतात.\footnote{
  पहिला
  टप्पा
  \emph{असतो}
  असे का
  गृहीत
  धरावे?
  \ref{sth:z:story:sec}
  मधील
  \stagesord{}
  जवळची
  तळटीप
  पाहा.}
त्यामुळे
आपण
एकच संच
तयार करतो:
एकही
घटक
नसलेला
रिक्त संच
$\emptyset$.
टप्पा~$1$
येथे
मागील
टप्प्यांपासून
बरोबर
एक संच
उपलब्ध
असतो,
म्हणून
फक्त एक
नवा संच
$\{\emptyset\}$
तयार होतो.
टप्पा~$2$
येथे मागील
टप्प्यांपासून
दोन संच
उपलब्ध
असतात;
आपण दोन
नवे संच
$\{\{\emptyset\}\}$
आणि
$\{\emptyset,
\{\emptyset\}\}$
तयार करतो.
टप्पा~$3$
येथे मागील
टप्प्यांपासून
चार संच
उपलब्ध
असतात;
म्हणून
बारा
नवे संच
तयार
होतात\ldots.
त्यामुळे
संचांचे
संचयी-क्रमिक
चित्र
साधारण
असे
दिसेल
(अंक टप्पे
दर्शवतात):
\begin{center}
  \begin{tikzpicture}[scale=0.6]
  \tikzset{cut_here/.style={densely dotted}}
  \draw (-4,6) -- (-1, 0)-- (2,6);
  \draw[cut_here] (-5,8)--(-4,6);
  \draw[cut_here] (2,6)--(3,8);
  \draw(-1.5, 1)--(-0.5, 1);
  \draw(-2, 2)--(0, 2);
  \draw(-2.5, 3)--(0.5,3);
  \draw(-3, 4)--(1, 4);
  \draw(-3.5, 5)--(1.5, 5);
  \draw(-4, 6)--(2, 6);
  \node[label] at (0, 0) {\small 0};
  \node[label] at (0.5, 1) {\small 1};
  \node[label] at (1, 2) {\small 2};
  \node[label] at (1.5, 3) {\small 3};
  \node[label] at (2, 4) {\small 4};
  \node[label] at (2.5, 5) {\small 5};
  \node[label] at (3, 6) {\small 6};
  \end{tikzpicture}
\end{center}
मग ही मांडणी
का स्वीकारावी?

एक कारण:
ती चांगली
आणि हाताळायला
सुलभ मांडणी
आहे. सर्वांत
उघड वाटणारी
मांडणी,
म्हणजे
अनौपचारिक
गुणधर्माधारित
संचरचना,
अपयशी ठरल्यावर
आपल्याला
एखादी चांगली
मांडणी हवीच.

पण ही मांडणी \emph{फक्त} सुलभ नाही.
ती रसेलच्या विरोधापत्तीला कारणीभूत असलेली
अमर्याद संचरचना का मानू नये हे स्पष्ट करते
आणि संच उपपत्तीची स्वयंसिद्धे निवडण्यास आधार देते.
खालील युक्तिवाद हा रसेलचा संच का बनत नाही यासाठी आहे;
या मांडणीवर आधारित कोणतीही उपपत्ती \emph{सुसंगत} असते,
असे तो सिद्ध करत नाही. पूर्ण सुसंगततेसाठी
नेमकी स्वयंसिद्धे आणि स्वतंत्र युक्तिवाद आवश्यक असतो.

संचांची
संचयी-क्रमिक
संकल्पना
मानल्यास
संच
टप्प्याटप्प्याने
बनतात;
त्यांचे
घटक
म्हणजे
\emph{आधीच}
उपलब्ध
असलेल्या
वस्तू असतात.
म्हणून
कोणत्याही
टप्प्या~$S$
साठी
पुढील संच
तयार करता
येतो:
\[
  R_S = \Setabs{x}{x \notin x \text{ आणि $x$ हा $S$ पूर्वी उपलब्ध होता}}
\]
रसेलची
विरोधापत्ती
सिद्ध करताना
वापरलेला
युक्तिवाद
दाखवतो
की~$R_S$
हा संच
स्वतः
टप्पा~$S$
पूर्वी
उपलब्ध
नसतो.
आणि
त्यात
व्याघात
नाही.
शिवाय
संचांची
संचयी-क्रमिक
संकल्पना
स्वीकारल्यास
रसेलचा
मूळ संच
तयार होईल
अशी
\emph{अपेक्षाच}
ठेवायला
नको होती.
कारण तो
‘‘कधीही
उपलब्ध
होतील’’
अशा
स्वतःचे
सदस्य
नसलेल्या
सर्व
संचांचा
संच
असता.
थोडक्यात,
संचयी-क्रमिक
मांडणी
मानल्यावर
रसेलचा
संच तयार
करता
येत
नाही
ही
(सिद्ध
होणारी)
गोष्ट
\emph{आश्चर्यकारक}
नाही;
याचाच
\emph{अंदाज}
आपण
केला
असता.













\section{मूलघटक असावेत का?}\label{sth:story:urelements:sec}

पुढील काही प्रकरणांत आपण संचांच्या संचयी-क्रमिक संकल्पनेतून स्वयंसिद्धके मिळवण्याचा प्रयत्न करू. पण पुढे जाण्याआधी \emph{मूलघटकांबद्दल} आणखी काही सांगणे गरजेचे आहे.

\ref{sth:story:approach:sec} मधील चित्रात आपण जुन्या \emph{संचांपासून} नवे संचच तयार करत होतो. पण काही \emph{संच नसलेल्या} वस्तू---गायी, डुकरे, वाळूचे कण किंवा इतर काही---संचांचे घटक असू शकतात, असे मानायचे असेल. तसे असल्यास सुरुवातीला काही मूलभूत वस्तू, म्हणजे \emph{मूलघटक}, उपलब्ध मानू. मग प्रत्येक टप्प्या~$S$ वर त्या मूलघटकांपासून किंवा $S$ पूर्वीच्या टप्प्यांवर तयार झालेल्या संचांपासून बनणारे ``सर्व शक्य'' संच तयार होतात, असे म्हणू; यात दोन्ही प्रकारच्या वस्तू एकत्रही असू शकतात. अशा वेळी चित्र साधारण असे दिसेल:
\begin{center}
	\begin{tikzpicture}[scale=0.6]
	\tikzset{cut_here/.style={densely dotted}}
	\draw (-4,6) -- (-1, 0) -- (1,0) -- (4,6);
	\draw[cut_here] (-5,8)--(-4,6);
	\draw[cut_here] (5,8)--(4,6);
	\draw(-1.5, 1)--(1.5, 1);
	\draw(-2, 2)--(2, 2);
	\draw(-2.5, 3)--(2.5,3);
	\draw(-3, 4)--(3, 4);
	\draw(-3.5, 5)--(3.5, 5);
	\draw(-4, 6)--(4, 6);
	\node[label] at (2, 0) {\small 0};
	\node[label] at (2.5, 1) {\small 1};
	\node[label] at (3, 2) {\small 2};
	\node[label] at (3.5, 3) {\small 3};
	\node[label] at (4, 4) {\small 4};
	\node[label] at (4.5, 5) {\small 5};
	\node[label] at (5, 6) {\small 6};
	\end{tikzpicture}
\end{center}
आता आपल्याला एक निर्णय घ्यायचा आहे: \emph{मूलघटकांना परवानगी द्यावी का?}

तात्त्विक दृष्टीने आपल्या सिद्धान्तात मूलघटक समाविष्ट करणे सयुक्तिक आहे. त्यामागचे मुख्य कारण म्हणजे संच उपपत्ती \emph{लागू करता येण्याजोगी} बनवणे. हे स्पष्ट करण्यासाठी \ref{sfr:siz::chap} आठवा: दोन संच $A$ आणि~$B$ यांचा आकार समान, म्हणजे $\cardeq{A}{B}$, तेव्हा आणि केवळ तेव्हाच त्यांच्यात एकास-एक व आच्छादक फलन असते. आता शेतातल्या गायी आणि गोठ्यातली डुकरे हे दोन्ही संच बनवत असतील, तर ``गायी जितक्या आहेत तितकीच डुकरे आहेत'' या विधानाचा संच उपपत्तीच्या आधारे विचार करता येतो. पण मूलघटकांना मनाई केली आणि गायी व डुकरे यांचे संच \emph{बनतच नाहीत} असे मानले, तर हा विचार करता येणार नाही. खरे तर संचांशिवाय इतर कशालाही संच उपपत्ती लागू करण्याचा सरळ मार्ग आपल्याकडे उरणार नाही. (मूलघटक समाविष्ट करण्याची आणखी कारणे Potter (2004, पृ.~vi, 24, 50--1) येथे पाहा.)

मात्र गणितात मूलघटकांना परवानगी देणे फारसे प्रचलित नाही. त्याचे एक कारण म्हणजे मूलघटकांशिवाय संच उपपत्ती मांडणे \emph{अगदी किंचित} सोपे असते. पण कधीकधी त्यांना वगळण्याचे अधिक लक्षवेधी समर्थनही आढळते:
\begin{quote}
  संच उपपत्ती हा गणिताचा पाया आहे या मतानुसार, केवळ संचांविषयी बोलून आपल्याला संपूर्ण गणित व्यक्त करता आले पाहिजे. म्हणून आपल्या चलांची व्याप्ती गायी आणि डुकरे यांसारख्या वस्तूंवर असता कामा नये.

  (Kunen, 1980, पृ.~8)
\end{quote}
म्हणून उपयोज्यतेला महत्त्व दिल्यास मूलघटक \emph{समाविष्ट} करावेसे वाटतील; पण गणित शुद्ध संच उपपत्तीत आणण्याच्या पायाभूत ध्येयाला महत्त्व दिल्यास त्यांना \emph{वगळावेसे} वाटेल. थोडा आळसही मूलघटक वगळण्याकडे झुकवतो.

आपण सर्वांत सोपा मार्ग निवडू. पण त्यामागे काही प्रमाणात शैक्षणिक कारणही आहे. संच उपपत्तीच्या तत्त्वज्ञानाचा अभ्यास सुरू करण्यासाठी तुम्हाला तिच्या ज्या मूलभूत गोष्टी लागतील, त्यांची ओळख करून देणे हे आपले उद्दिष्ट आहे. या विषयावरील बहुतेक तात्त्विक साहित्य मूलघटक \emph{नसलेल्या} संच उपपत्तीची चर्चा करते. म्हणून त्या साहित्याशी साधारण जुळून राहिल्यास हे पुस्तक अधिक उपयोगी ठरेल.









\section{परिशिष्ट: फ्रेगे यांचा मूलनियम पाच}\label{sth:story:blv:sec}

\ref{sth:story:rus:sec} मध्ये आपण पाहिले की रसेल यांनी आपली विरोधापत्ती फ्रेगे यांनी \emph{Grundgesetze} मध्ये मांडलेल्या प्रणालीसमोरील अडचण म्हणून सांगितली. फ्रेगे यांच्या प्रणालीत अनौपचारिक गुणधर्माधारित संचरचनेची थेट मांडणी नव्हती. म्हणून या परिशिष्टात त्या प्रणालीत \emph{नेमके काय} होते, त्याचा अनौपचारिक गुणधर्माधारित संचरचनेशी कसा संबंध आहे आणि रसेल यांच्या विरोधापत्तीशी कसा संबंध आहे, हे अगदी थोडक्यात पाहू.

फ्रेगे यांची प्रणाली \emph{द्वितीय-क्रमाची} आहे. तिचा हेतू \emph{संकल्पनेचा विस्तार} ही संकल्पना औपचारिकरीत्या मांडणे हा होता.\footnote{काटेकोरपणे सांगायचे तर फ्रेगे यांनी अधिक व्यापक संकल्पना, म्हणजे फलनाची ``मूल्यव्याप्ती'', औपचारिकरीत्या मांडण्याचा प्रयत्न केला. संकल्पनांचे विस्तार हे त्या व्यापक संकल्पनेचे विशेष उदाहरण आहे. तपशीलांसाठी Heck (2012, पृ.~8--9) पाहा.} फ्रेगे यांच्या संकेतपद्धतीच्या धर्तीवर $\fregeext{x}{F(x)}$ हे \emph{संकल्पना~$F$ चा विस्तार} दर्शवेल. हे साधन ``$F$'' या \emph{विधेयाचे} (प्रथम-क्रमातील) ``$\fregeext{x}{F(x)}$'' या \emph{पदात} रूपांतर करते. या साधनाच्या आधारे फ्रेगे यांनी सदस्यत्वाची पुढील \emph{व्याख्या} दिली:
\[
	a \in b =_\text{df} \exists G(b = \fregeext{x}{G(x)} \land Ga)
\]
स्थूलमानाने, $a \in b$ तेव्हा आणि केवळ तेव्हाच जेव्हा $a$ अशी एखादी संकल्पना पूर्ण करते की तिचा विस्तार $b$ आहे. (येथे ``$\exists G$'' हा संख्यापक द्वितीय-क्रमाचा आहे, हे लक्षात घ्या.) फ्रेगे यांनी \emph{मूलनियम पाच} म्हणून ओळखले जाणारे पुढील तत्त्वही मानले:
$$\fregeext{x}{F(x)} = \fregeext{x}{G(x)} \liff \forall x (Fx \liff Gx)$$
स्थूलमानाने, दोन संकल्पनांचे विस्तार एकच असतात तेव्हा आणि केवळ तेव्हाच त्या नेमक्या त्याच वस्तूंना लागू पडतात. (पुन्हा, ``$F$'' आणि ``$G$'' ही दोन्ही विधेयस्थानी आहेत.) आता एक साधे तत्त्व सदस्यत्वाचा संबंध गुणधर्मपूर्तीशी जोडते:

\begin{lem}[\emph{Grundgesetze} मध्ये]\label{sth:story:blv:lem:Fregeextensions}
$\forall F \forall a(a \in \fregeext{x}{F(x)} \liff Fa)$
\end{lem}

\begin{proof}
$F$ आणि $a$ निश्चित करू. सदस्यत्वाच्या व्याख्येनुसार $a \in \fregeext{x}{F(x)}$ तेव्हा आणि केवळ तेव्हाच जेव्हा $\exists G(\fregeext{x}{F(x)}
= \fregeext{x}{G(x)} \land Ga)$; मूलनियम पाचनुसार तेव्हा आणि केवळ तेव्हाच जेव्हा $\exists G(\forall x(Fx \liff Gx) \land Ga)$; आणि प्राथमिक द्वितीय-क्रम तर्कशास्त्रानुसार तेव्हा आणि केवळ तेव्हाच जेव्हा $Fa$.
\end{proof}

यातून अनौपचारिक गुणधर्माधारित संचरचना जवळजवळ लगेच मिळते:

\begin{lem}[\emph{Grundgesetze} मध्ये.]
$\forall F \exists s \forall a (a \in s \liff Fa)$
\end{lem}

\begin{proof}
$F$ निश्चित करू. मग \ref{sth:story:blv:lem:Fregeextensions} वरून $\forall a (a \in
\fregeext{x}{F(x)} \liff Fa)$ मिळते. म्हणून अस्तित्ववाचक सामान्यीकरणाने $\exists s\forall a(a \in s \liff Fa)$ मिळते. $F$ स्वैर असल्यामुळे अपेक्षित निष्कर्ष मिळतो.
\end{proof}

येथे पूर्ण द्वितीय-क्रम गुणधर्माधारित संकल्पनारचना
वापरून $\exists F\forall x(Fx\liff x\notin x)$ हे गृहीतक मिळते.
सदस्यत्वाच्या वरील व्याख्येत द्वितीय-क्रम संख्यापक आहे;
म्हणून हा संकल्पनारचनेचा उपयोग फक्त प्रथम-क्रम सूत्रांपुरता मर्यादित नाही.
अशी $F$ घेऊन $R=\fregeext{x}{F(x)}$ ठेवले, तर
वरील उपप्रमेयानुसार $R\in R\liff FR$, तर $F$ च्या व्याख्येनुसार
$FR\liff R\notin R$. त्यामुळे रसेलची विरोधापत्ती मिळते.
या युक्तिवादाला मूलनियम पाचासोबत किमान
वरील विशिष्ट संकल्पनारचनेचे उदाहरण आवश्यक आहे;
पूर्ण रूपबंधामुळे ते उपलब्ध होते.
संकल्पनारचना मर्यादित करून सामान्य द्वितीय-क्रम प्रतिमाने वापरल्यास
मूलनियम पाचाच्या सुसंगततेचा वेगळा प्रश्न उभा राहतो;
प्रत्येक मर्यादित रूपबंध असुसंगत असल्याचा हा युक्तिवाद नाही.



\chapter{$\Z$ कडे वाटचाल}\label{sth:z:chap}




\section{मांडणीचे आणखी तपशील}\label{sth:z:story:sec}

\ref{sth:story:approach:sec} मध्ये आपण संच तयार होण्याच्या प्रक्रियेबद्दल शोनफिल्ड यांचे वर्णन उद्धृत केले. आता ही मांडणी थोडी अधिक नेमकी करण्यासाठी आणखी काही तत्त्वे लिहू. ती अशी:
\begin{enumerate}
\item[] \stageshier. प्रत्येक संच कोणत्या ना कोणत्या टप्प्यावर तयार होतो.
\item[] \stagesord. टप्प्यांचा क्रम असतो: काही टप्पे इतरांच्या \emph{आधी} येतात.\footnote{प्रत्यक्षात टप्पे \emph{सुक्रमित} आहेत, असे आपण अप्रकटपणे गृहीत धरू. याचा अर्थ \ref{sth:ordinals:chap} मध्ये स्पष्ट केला आहे. हे महत्त्वाचे गृहीतक आहे. खरे तर Scott (1974) यांच्या अतिशय कल्पक पद्धतीने हे गृहीतक प्रथम \emph{टाळता} येते आणि मग \emph{सिद्ध} करता येते. (आरंभीचा टप्पा असावा असे का मानावे, हेही त्यातून स्पष्ट होईल.) येथे त्यात जाता येणार नाही; अधिक माहितीसाठी Button (2021) पाहा.}
\item[] \stagesacc. कोणताही टप्पा $S$ आणि $S$ च्या \emph{आधी} तयार झालेल्या संचांचा कोणताही संग्रह घेतल्यास, टप्पा~$S$ वर नेमके तेच संच सदस्य असलेला संच तयार होतो. टप्पा~$S$ वर दुसरे काही तयार होत नाही.
\end{enumerate}
ही अनौपचारिक तत्त्वे आहेत; पण त्यांच्या आधारे झर्मेलो यांच्या संच उपपत्तीतील अनेक स्वयंसिद्धकांचे समर्थन करता येईल.

(एक सावधानतेची नोंद. पुढे आपण सर्वमान्य नावे असलेली प्रमाणित स्वयंसिद्धके मांडणार आहोत; पण नुकतीच तिरपी अक्षरे वापरून मांडलेली तत्त्वे साहित्यात कोणत्याही विशिष्ट नावांनी ओळखली जात नाहीत. आपण त्यांना केवळ उपयोगी वाटतील अशी नावे दिली आहेत.)








\section{विभक्तीकरण}\label{sth:z:sep:sec}

अनौपचारिक गुणधर्माधारित संचरचनेच्या जागी पुढील तत्त्व घेऊ:

\begin{axiom}[विभक्तीकरण-रूपबंध] प्रत्येक सूत्र $\phi(x)$ साठी पुढील एक स्वयंसिद्धक आहे: कोणताही $A$ दिल्यास $\Setabs{x \in A}{\phi(x)}$ हा संच अस्तित्वात असतो.
\end{axiom}

हे एकच स्वयंसिद्धक नाही, हे लक्षात घ्या. हा स्वयंसिद्धकांचा \emph{रूपबंध} आहे. विभक्तीकरणाची स्वयंसिद्धके \emph{अनंत} आहेत; प्रत्येक सूत्र $\phi(x)$ साठी एक. हा रूपबंध पुढीलप्रमाणेही लिहिता येतो (आणि सहसा तसाच लिहिला जातो):

\begin{defish}
``$S$'' नसलेल्या कोणत्याही सूत्र $\phi(x)$ साठी पुढील एक स्वयंसिद्धक आहे:
\[
	\forall A \exists S \forall x(x \in S \liff (\phi(x) \land x \in A)).
\]
\end{defish}

\ref{sth:::part} च्या सुरुवातीला सांगितलेल्या संकेतानुसार, विभक्तीकरणाच्या स्वयंसिद्धकांतील सूत्रांना~$\phi$ प्राचले असू शकतात.\footnote{याचा अर्थ काय, यासाठी \ref{sfr:infinite:induction:natinductionschema} नंतरची चर्चा पाहा.}

आपण सांगत आलेली संचांची संचयी-क्रमिक संकल्पना विभक्तीकरणाचे लगेच समर्थन करते. ते पाहण्यासाठी $A$ हा संच मानू. \stageshier{} नुसार $A$ कोणत्या तरी टप्प्या~$S$ वर तयार होतो. $A$ टप्पा~$S$ वर तयार झाला असल्याने $A$ चे सर्व सदस्य टप्पा~$S$ च्या आधी तयार झालेले असतात (\stagesacc{} नुसार). आता $A$ चे सदस्य असलेले आणि $\phi$ पूर्ण करणारे सर्व संच घ्या; हेसुद्धा टप्पा~$S$ च्या आधीच तयार झालेले असतात. म्हणून \stagesacc{} नुसार टप्पा~$S$ वर त्यांचा $\Setabs{x \in A}{\phi(x)}$ हा संच तयार होतो.

अनौपचारिक गुणधर्माधारित संचरचनेच्या विपरीत, विभक्तीकरणामुळे रसेलची विरोधापत्ती टळते. कारण $\Setabs{x}{x \notin
x}$ हा संच आहे, असे थेट म्हणता येत नाही. पण एखादा संच~$A$ \emph{दिला} असेल, तर $R_A = \Setabs{x \in A}{x \notin x}$ हा संच आहे, असे म्हणता येते. यातून फक्त $R_A \notin R_A$ आणि $R_A \notin A$ हे निष्पन्न होते; त्यांत चिंतेचे काही नाही.

मात्र विभक्तीकरणाचा एक तात्काळ आणि महत्त्वाचा परिणाम आहे:

\begin{thm}\label{sth:z:sep:thm:NoUniversalSet}
\emph{सर्वव्यापी} संच अस्तित्वात नाही; म्हणजे $\Setabs{x}{x = x}$ हा संच अस्तित्वात नाही.
\end{thm}

\begin{proof}
व्यतिरेकासाठी $V$ हा सर्वव्यापी संच आहे असे मानू. मग विभक्तीकरणानुसार $R =
\Setabs{x \in V}{x \notin x} = \Setabs{x}{x \notin x}$ हा संच अस्तित्वात असतो; पण याने रसेलची विरोधापत्ती निर्माण होते.
\end{proof}

सर्वव्यापी संच नसणे---खरे तर संचांच्या पदानुक्रमाला शेवट नसणे---ही संचांच्या संचयी-क्रमिक संकल्पनेमागील एक मूलभूत कल्पना आहे. त्यामुळे सहजज्ञानाने याच निष्कर्षापर्यंत दुसऱ्या मार्गाने कसे पोहोचता येते, हे पाहणे उपयुक्त ठरेल. सर्वव्यापी संच स्वतःचा घटक असलाच पाहिजे. पण टप्पे सुक्रमित असल्याच्या गृहीतकानुसार प्रत्येक संच प्रथम तयार होण्याचा एक टप्पा असतो. त्या टप्प्यावर तयार होणाऱ्या संचाचे सर्व घटक आधीच्या टप्प्यांवर तयार झालेले असतात. म्हणून \emph{कोणताही} संच स्वतःचा सदस्य नसतो: तो स्वतःचा सदस्य असता, तर प्रथम तयार होण्याच्या टप्प्याच्या आधीच तो तयार झाला असता; हा व्याघात आहे. या युक्तिवादाला सर्व सदस्यांचा एखादा शेवटचा टप्पा किंवा त्याचा लगतचा उत्तरवर्ती टप्पा आवश्यक नाही. (संचाच्या ``क्रमस्थानाची'' संकल्पना वापरून हा युक्तिवाद अधिक काटेकोर कसा करायचा, ते \ref{sth:spine:rank:defnsetrank} मध्ये पाहू. मात्र त्यासाठी आधी आणखी काही स्वयंसिद्धके लागतील.)

विभक्तीकरण आणि विस्तारात्मकता यांचे आणखी काही परिणाम पाहू.

\begin{prop}\label{sth:z:sep:prop:emptyexists}
कोणताही संच अस्तित्वात असेल, तर $\emptyset$ अस्तित्वात असतो.
\end{prop}

\begin{proof}
$A$ हा संच असेल, तर विभक्तीकरणानुसार $\emptyset = \Setabs{x \in A}{x \neq x}$ अस्तित्वात असतो.
\end{proof}

\begin{prop}
कोणतेही संच $A$ आणि $B$ यांच्यासाठी $A \setminus B$ अस्तित्वात असतो.
\end{prop}

\begin{proof}
विभक्तीकरणानुसार $A \setminus B = \Setabs{x \in A}{x \notin B}$ अस्तित्वात असतो.
\end{proof}

याशिवाय (जवळजवळ) मनमानी छेदही अस्तित्वात असतात:

\begin{prop}\label{sth:z:sep:prop:intersectionsexist}
$A \neq \emptyset$ असेल, तर $\bigcap A = \Setabs{x}{(\forall y \in A)x \in y}$ अस्तित्वात असतो.
\end{prop}

\begin{proof}
$A \neq \emptyset$ मानू; मग एखादा $c \in A$ असतो. म्हणून $\bigcap A =
\Setabs{x}{(\forall y \in A)x \in y} = \Setabs{x \in c}{(\forall y \in
A)x \in y}$; हा संच विभक्तीकरणानुसार अस्तित्वात असतो.
\end{proof}

पण $A \neq \emptyset$ ही अट लक्षात घ्या. कारण $\bigcap
\emptyset$ हा रिक्ततेमुळे सर्वव्यापी संच ठरेल; ते \ref{sth:z:sep:thm:NoUniversalSet} शी विसंगत आहे.







\section{संयोग}\label{sth:z:union:sec}

\ref{sth:z:sep:prop:intersectionsexist} मधून आपल्याला छेद मिळाले. पण मनमानी संयोगही अस्तित्वात असावेत, तर आणखी एक स्वयंसिद्धक मांडावे लागेल:

\begin{axiom}[संयोग] कोणत्याही संच $A$ साठी $\bigcup A =
\Setabs{x}{(\exists b \in A) x \in b}$ हा संच अस्तित्वात असतो.
\[
	\forall A \exists U \forall x(x \in U \liff (\exists b \in A)x \in b)
\]
\end{axiom}

संचांची संचयी-क्रमिक संकल्पनाही या स्वयंसिद्धकाचे समर्थन करते. $A$ हा संच मानू; \stageshier{} नुसार $A$ कोणत्या तरी टप्प्या $S$ वर तयार होतो. \stagesacc{} नुसार $A$ चा प्रत्येक सदस्य $S$ च्या \emph{आधी} तयार झालेला असतो; त्याचप्रमाणे विचार केल्यास $A$ च्या प्रत्येक सदस्याचा प्रत्येक सदस्यही $S$ च्या आधी तयार झालेला असतो. त्यामुळे \emph{हे सर्व} संच $S$ च्या आधी उपलब्ध असतात आणि $S$ वर त्यांचा एक संच तयार होतो. तोच $\bigcup A$ आहे.







\section{जोड्या}\label{sth:z:pairs:sec}

आता पुढील स्वयंसिद्धक पाहू:

\begin{axiom}[जोड्या]
कोणतेही संच $a, b$ दिल्यास $\{a, b\}$ हा संच अस्तित्वात असतो.
\[
	\forall a \forall b \exists P \forall x (x \in P \liff (x = a \lor x = b))
\]
\end{axiom}

क्रमिक संकल्पनेच्या आधारे या स्वयंसिद्धकाचे समर्थन असे करता येईल. $a$ टप्पा $S$ वर आणि $b$ टप्पा $T$ वर उपलब्ध आहेत असे मानू. $S$ आणि $T$ यांपैकी नंतर येणारा टप्पा $M$ घ्या; दोन्ही टप्पे समान असतील, तर $M=S=T$ घ्या. मग $a$ आणि $b$ दोन्ही टप्पा $M$ वर उपलब्ध असल्याने, $\{a,b\}$ हा $M$ नंतरच्या कोणत्याही टप्प्यावर उपलब्ध वस्तूंपासून बनू शकणारा संग्रह आहे.

पण थांबा!{} $M$ नंतर टप्पे \emph{आहेतच} असे का मानावे? तसे नसतील तर आपले समर्थन अपुरे पडेल. म्हणून जोड्यांचे स्वयंसिद्धक समर्थित करण्यासाठी \ref{sth:z:story:sec} मधील मांडणीत आणखी एक तत्त्व घालावे लागेल:
\begin{enumerate}
\item[] \stagessucc. शेवटचा टप्पा नसतो.
\end{enumerate}
या तत्त्वाचे समर्थन करता येते का? शोनफिल्ड यांच्या मांडणीत शेवटचा टप्पा नाही असे \emph{स्पष्टपणे} म्हटलेले नाही. तरी काटेकोरपणे पाहता हा अतिरिक्त भाग असला, तरी टप्प्याटप्प्याने संच तयार होतात या मूळ कल्पनेशी तो सुसंगत आहे. पुढील चर्चेत आपण हे तत्त्व स्वीकारू. त्याबरोबर जोड्यांचे स्वयंसिद्धकही स्वीकारू.

या नव्या स्वयंसिद्धकाच्या साहाय्याने आणखी अनेक संचांचे अस्तित्व सिद्ध करता येते. उदाहरणार्थ:

\begin{prop}\label{sth:z:pairs:prop:pairsconsequences}
कोणतेही संच $a$ आणि $b$ दिल्यास पुढील संच अस्तित्वात असतात:
\begin{enumerate}
\item\label{sth:z:pairs:singleton} $\{a\}$
\item\label{sth:z:pairs:binunion} $a \cup b$
\item\label{sth:z:pairs:tuples} $\tuple{a, b}$
\end{enumerate}
\end{prop}

\begin{proof}
\ref{sth:z:pairs:singleton}. जोड्यांच्या स्वयंसिद्धकानुसार $\{a, a\}$ अस्तित्वात असतो; विस्तारात्मकतेनुसार तो $\{a\}$ च असतो.

\ref{sth:z:pairs:binunion}. जोड्यांच्या स्वयंसिद्धकानुसार $\{a, b\}$ अस्तित्वात असतो. आता संयोगाच्या स्वयंसिद्धकानुसार $a \cup b = \bigcup
\{a, b\}$ अस्तित्वात असतो.

\ref{sth:z:pairs:tuples}. \ref{sth:z:pairs:singleton} नुसार $\{a\}$ अस्तित्वात असतो. जोड्यांच्या स्वयंसिद्धकानुसार $\{a,
b\}$ अस्तित्वात असतो. मग पुन्हा जोड्यांच्या स्वयंसिद्धकानुसार $\{\{a\}, \{a, b\}\} = \tuple{a, b}$ अस्तित्वात असतो.
\end{proof}

\begin{prob}
कोणतेही संच $a, b, c$ दिल्यास $\{a, b, c\}$ अस्तित्वात असतो, हे दाखवा.
\end{prob}

\begin{prob}
कोणतेही संच $a_1, \ldots, a_n$ दिल्यास $\{a_1, \ldots,
a_n\}$ अस्तित्वात असतो, हे दाखवा.
\end{prob}












\section{घातसंच}\label{sth:z:power:sec}

आता आणखी एक स्वयंसिद्धक पाहू:

\begin{axiom}[घातसंच]
कोणत्याही संच $A$ साठी $\Pow{A} = \Setabs{x}{x \subseteq A}$ हा संच अस्तित्वात असतो.
\[
	\forall A \exists P \forall x(x \in P \liff (\forall z \in x)z \in A)
\]
\end{axiom}

याचे समर्थन अगदी सरळ आहे. $A$ हा संच टप्पा $S$ वर तयार झाला असे मानू. मग \stagesacc{} नुसार $A$ चे सर्व सदस्य $S$ च्या आधी उपलब्ध होते. म्हणून विभक्तीकरणाच्या समर्थनाप्रमाणे विचार केल्यास $A$ चा प्रत्येक उपसंच टप्पा $S$ पर्यंत तयार झालेला असतो. ते सर्व उपलब्ध असल्याने $S$ नंतरच्या कोणत्याही टप्प्यावर त्यांचा एक संच तयार करता येतो. \stagessucc{} नुसार $S$ हा शेवटचा टप्पा नाही, म्हणून असा पुढील टप्पा आहेच. त्यामुळे $\Pow{A}$ अस्तित्वात असतो.

घातसंचाच्या स्वयंसिद्धकाचा एक चांगला परिणाम असा:

\begin{prop}\label{thm:Products}
कोणतेही संच $A, B$ दिल्यास त्यांचा कार्तीय गुणाकार $A \times B$ अस्तित्वात असतो.
\end{prop}

\begin{proof}
घातसंचाचे स्वयंसिद्धक आणि \ref{sth:z:pairs:prop:pairsconsequences} यांनुसार $\Pow{\Pow{A \cup B}}$ अस्तित्वात असतो. म्हणून विभक्तीकरणानुसार पुढील संच अस्तित्वात असतो:
\[
	C = \Setabs{z \in \Pow{\Pow{A \cup B}}}{(\exists x \in A)(\exists y \in B) z = \tuple{x, y}}.
\]
आता कोणतेही $x \in A$ आणि $y \in B$ दिल्यास \ref{sth:z:pairs:prop:pairsconsequences} नुसार $\tuple{x, y}$ अस्तित्वात असतो. शिवाय $x, y
\in A \cup B$ असल्याने $\{x\}, \{x, y\} \in \Pow{A \cup B}$ आणि $\tuple{x,y} \in \Pow{\Pow{A \cup B}}$ मिळते. म्हणून $A \times B = C$.
\end{proof}

या सिद्धतेत घातसंच आणि विभक्तीकरण एकत्र काम करतात. त्यात आश्चर्य नाही. विभक्तीकरण नसते, तर घातसंचाचे तत्त्व फारसे \emph{सामर्थ्यवान} ठरले नसते. कारण कोणते उपसंच अस्तित्वात असतात हे विभक्तीकरण सांगते; त्यामुळे प्रत्येक घातसंच किती ``भरलेला'' असेल, हेही त्यावर अवलंबून असते.

\begin{prob}
कोणतेही संच $A, B$ दिल्यास दाखवा: (i) प्रांत $A$ आणि व्याप्ती $B$ असलेल्या सर्व संबंधांचा संच अस्तित्वात असतो; आणि (ii) $A$ पासून $B$ कडे जाणाऱ्या सर्व फलनांचा संच अस्तित्वात असतो.
\end{prob}

\begin{prob}
$A$ हा संच आणि $\sim$ हा $A$ वरील तुल्यता संबंध मानू. $\sim$ नुसार $A$ वरील तुल्यतावर्गांचा संच, म्हणजे $\equivclass{A}{\sim}$, अस्तित्वात असतो हे सिद्ध करा.
\end{prob}








\section{अनंतता}\label{sth:z:infinity-again:sec}

(एखादा संच अस्तित्वात असेल, तर) अनंत संख्येने संच अस्तित्वात आहेत, हे दाखवण्यासाठी आपल्याकडे आधीच पुरेशी स्वयंसिद्धके आहेत. समजा एखादा संच अस्तित्वात आहे. मग \ref{sth:z:sep:prop:emptyexists} नुसार $\emptyset$ अस्तित्वात असतो. आता कोणताही संच~$x$ दिल्यास \ref{sth:z:pairs:prop:pairsconsequences} नुसार $x \cup \{x\}$ अस्तित्वात असतो. ही प्रक्रिया काही वेळा केल्यास पुढील संच मिळतात:
\begin{enumerate}
	\item[0.] $\emptyset$
	\item[1.] $ \{\emptyset\}$
	\item[2.] $\{\emptyset, \{\emptyset\}\}$
	\item[3.] $\{\emptyset, \{\emptyset\}, \{\emptyset, \{\emptyset\}\}\}$
	\item[4.] $\{\emptyset, \{\emptyset\}, \{\emptyset, \{\emptyset\}\}, \{\emptyset, \{\emptyset\}, \{\emptyset, \{\emptyset\}\}\}\}$
\end{enumerate}
हे प्रत्येक संच इतरांहून वेगळे आहेत, हे तपासता येते.

क्रमांकन $0$ पासून सुरू करण्यामागे काही कारणे आहेत. त्यांपैकी एक असे: ``$n$'' असे नाव दिलेल्या संचात नेमके $n$ सदस्य आहेत आणि सहजज्ञानानुसार तो $n$ व्या टप्प्यावर तयार होतो, हे तपासणे अवघड नाही.

पण यातून आपल्याला \emph{अनंत संख्येने} संच मिळतात; \emph{अनंत संच}, म्हणजे अनंत सदस्य असलेला एक संच, अस्तित्वात आहे याची खात्री मिळत नाही. हा फरक महत्त्वाचा आहे. कारण डेडेकिंड-अनंत संच मिळाल्याशिवाय डेडेकिंड बीजसंरचना बनवता येणार नाही. आणि नैसर्गिक संख्यांचा संच म्हणून वापरता यावी, यासाठी अशी बीजसंरचना आपल्याला हवी आहे. (\ref{sfr:infinite:dedekindsproof:sec} शी तुलना करा.)

आतापर्यंत मांडलेली स्वयंसिद्धके कोणत्याही अनंत संचाचे अस्तित्व \emph{सिद्ध करत नाहीत}. म्हणून नवे स्वयंसिद्धक मांडावे लागेल:

\begin{axiom}[अनंतता]
असा एक संच $I$ अस्तित्वात असतो की $\emptyset \in I$ आणि $x \in I$ असेल तेव्हा $x \cup \{x\} \in I$.
\begin{align*}
	\exists I( & (\exists o \in I)\forall x\ x \notin o \land {}\\
	& (\forall x \in I)(\exists s \in I)\forall z(z \in s \liff (z \in x \lor z = x)))
\end{align*}
\end{axiom}

अनंततेच्या स्वयंसिद्धकातील $I$ साठी $s(x)=x\cup\{x\}$ हे
$I$ पासून $I$ कडे जाणारे फलन आहे. मात्र सर्व $I$ वर त्याचा
एकास-एकपणा थेट गृहीत धरण्याऐवजी प्रथम संवरण-संच घेऊ:
\begin{defn}\label{sth:z:infinity-again:defnomega}
अनंततेच्या स्वयंसिद्धकातून मिळालेला कोणताही संच $I$ घ्या.
$s$ हे $s(x)=x\cup\{x\}$ असे फलन आणि
$\omega=\closureofunder{s}{\emptyset}$ हा $I$ मधील सर्वांत लहान
$s$-बंद, $\emptyset$ धारण करणारा उपसंच असू द्या.
$\omega$ चे सदस्य \emph{नैसर्गिक संख्या} म्हणू.
$n$ ही $\emptyset$ वर $s$ च्या $n$ वेळा केलेल्या उपयोगाचे फलित आहे,
असे म्हणू.
\end{defn}
हे संवरण विभक्तीकरणाने अस्तित्वात येते: $I$ मधील,
$\emptyset$ धारण करणाऱ्या प्रत्येक $s$-बंद उपसंचात असणारे सदस्य घ्या.
असा किमान एक उपसंच $I$ स्वतः आहे. मिळणारा संच
$\emptyset$ धारण करतो, $s$-बंद असतो आणि या सर्व उपसंचांत समाविष्ट असतो.
याच लघुत्वामुळे $\omega$ वर विगमन मिळते: एखादा गुणधर्म
$\emptyset$ साठी सत्य आणि $s$ च्या उपयोगाने जतन होत असेल,
तर त्या गुणधर्माचे $\omega$ मधील सदस्य विभक्तीकरणाने
असा बंद उपसंच बनवतात; म्हणून तो संपूर्ण $\omega$ असतो.
यासाठी आधी $s$ एकास-एक आहे असे मानावे लागत नाही.
या विगमनाने प्रत्येक $n\in\omega$ हा संक्रमक संच,
म्हणजे $(\forall u\in n)u\subseteq n$, असून $n\notin n$ हे सिद्ध होते.
रिक्त संचासाठी दोन्ही अटी स्पष्ट आहेत.
संक्रमक $n$ पासून $s(n)$ संक्रमक होतो: त्याचा सदस्य
$n$ असतो किंवा $n$ चा सदस्य असतो; दोन्ही प्रकरणांत
त्या सदस्याचे सर्व सदस्य $s(n)$ मध्ये असतात.
शिवाय $s(n)\in s(n)$ मानल्यास एकतर $s(n)=n$ किंवा $s(n)\in n$.
पहिल्या प्रकरणात $n\in s(n)$ वरून $n\in n$;
दुसऱ्यात $n$ च्या संक्रमकतेमुळे, पुन्हा $n\in s(n)$ वरून $n\in n$.
दोन्ही प्रकरणांत व्याघात आहे.
आता $s(n)=s(m)$ आणि $n\neq m$ असेल, तर
$n\in m$ आणि $m\in n$; $n$ च्या संक्रमकतेने $n\in n$,
हा व्याघात. त्यामुळे $s$ चे $\omega$ वरील बंधन एकास-एक आहे.
तसेच $\emptyset$ कोणत्याही $s(n)$ चे मूल्य नाही,
कारण $n\in s(n)$. म्हणून $\omega$ डेडेकिंड-अनंत आहे,
आणि $(\omega,s,\emptyset)$ ही डेडेकिंड बीजसंरचना आहे;
\ref{sfr:infinite:dedekind:thm:DedekindInfiniteAlgebra}
शी तुलना करा. $I$ स्वतःही डेडेकिंड-अनंत आहे:
$\omega$ वर $s$ आणि $I\setminus\omega$ वरील प्रत्येक घटकाला तोच घटक देणारे फलन वापरल्यास
$I$ वर एकास-एक फलन मिळते, ज्याच्या व्याप्तीत $\emptyset$ नसतो. या दोन भागांची मूल्ये अनुक्रमे $\omega$ आणि
$I\setminus\omega$ मध्ये असल्याने ती एकमेकांत जुळत नाहीत.
या फलनाचा आलेख $I\times I$ मधून विभक्तीकरणाने मिळतो.

आधीच्या काही परिच्छेदांत ``$n$'' अशी खूण केलेला संच आता $n$ ही संख्याच \emph{मानला} जाईल, हे मागे पाहून तपासता येते.

या निर्णयाचे महत्त्व \ref{sth:z:nat:sec} मध्ये पाहू. सध्या त्याच्या आधारे एक सहज वाटणारा निष्कर्ष सिद्ध करता येतो:

\begin{prop}\label{sth:z:infinity-again:naturalnumbersarentinfinite}
कोणतीही नैसर्गिक संख्या डेडेकिंड-अनंत नाही.
\end{prop}

\begin{proof}
सिद्धता विगमनाने करायची आहे; \ref{sfr:infinite:induction:thm:dedinfiniteinduction} पाहा. स्पष्टच, $0
= \emptyset$ डेडेकिंड-अनंत नाही. विगमनाच्या पायरीसाठी आपण व्यत्यस्त रूप सिद्ध करू: $s(n)$ डेडेकिंड-अनंत असल्याचे (अशक्य असूनही) मानल्यास $n$ डेडेकिंड-अनंत ठरेल.

म्हणून $s(n)$ डेडेकिंड-अनंत आहे असे मानू. म्हणजे असा एखादा एकास-एक फलन $f$ आहे की $\ran{f}\subsetneq \dom{f} = s(n) = n \cup
\{n\}$. आता दोन प्रकरणे पाहावी लागतील.

\emph{प्रकरण १: $n \notin \ran{f}$.} मग $\ran{f} \subseteq n$ आणि $f(n) \in n$. $g = \funrestrictionto{f}{n}$ मानू; मग $\ran{g} =
\ran{f} \setminus \{f(n)\} \subsetneq n = \dom{g}$. म्हणून $n$ डेडेकिंड-अनंत आहे.

\emph{प्रकरण २: $n \in \ran{f}$.} एखादा $m \in \dom{f} \setminus \ran{f}$ निश्चित करू आणि प्रांत $s(n) = n \cup \{n\}$ असलेले फलन $h$ पुढीलप्रमाणे ठरवू:
\[
h(x) =
	\begin{cases}
		f(x) & \text{जर }f(x) \neq n\\
		m & \text{जर }f(x)=n
	\end{cases}
\]
म्हणून $h$ आणि $f$ सर्वत्र जुळतात; फक्त $h(f^{-1}(n)) = m
\neq n = f(f^{-1}(n))$ या ठिकाणी वेगळे आहेत. $f$ हे एकास-एक फलन असल्याने $n \notin
\ran{h}$; तसेच $\ran{h} \subsetneq \dom{h} = s(n)$.
हे $h$ एकास-एक आहे: $f$ चे फक्त एक मूल्य बदलले आहे,
आणि नवे मूल्य $m$ आधीच्या व्याप्तीत नव्हते.
आता प्रकरण १ मधील युक्तिवादाने $n$ डेडेकिंड-अनंत ठरतो.
\end{proof}

पण अनंततेच्या स्वयंसिद्धकाचे \emph{समर्थन} कसे करायचे, हा प्रश्न उरतो. थोडक्यात उत्तर असे की आपण सांगत आलेल्या मांडणीत आणखी एक तत्त्व जोडावे लागेल. ते असे:
\begin{enumerate}
\item[] \stagesinf. अनंत टप्पा असतो. म्हणजे असा टप्पा असतो की (a) तो पहिला नसतो, (b) त्याच्या आधी काही टप्पे असतात आणि (c) त्याचा लगतचा पूर्ववर्ती टप्पा नसतो.
\end{enumerate}
या तत्त्वावरून अनंततेचे स्वयंसिद्धक सरळ मिळते. नैसर्गिक संख्या $n$ ही टप्पा $n$ वर तयार होते, हे आपल्याला ठाऊक आहे. म्हणून $\omega$ हा संच पहिल्या अनंत टप्प्यावर तयार होतो. $\omega$ स्वतःच अनंततेच्या स्वयंसिद्धकासाठी अपेक्षित उदाहरण ठरतो.

मात्र यामुळे \stagesinf{} चे समर्थन कसे करायचे हा प्रश्नच पुढे ढकलला जातो. \stagessucc{} प्रमाणे हे तत्त्वही संचांच्या संचयी-क्रमिक पदानुक्रमाबद्दलच्या मूळ मांडणीत स्पष्टपणे नव्हते. त्याहून महत्त्वाचे म्हणजे संच टप्प्याटप्प्याने तयार होतात या कल्पनेतूनच प्रक्रियेत \emph{अनंत} टप्पा असलाच पाहिजे, असे निष्पन्न होत नाही. टप्प्यांचा क्रम नैसर्गिक संख्यांसारखा आहे, असे मानणे पूर्णपणे सुसंगत वाटते.

मग एक उघड अडचण उभी राहते. \ref{sfr:infinite:dedekindsproof:sec} मध्ये विचारांचा डेडेकिंड-अनंत संच आहे या डेडेकिंड यांच्या ``सिद्धतेचा'' आपण विचार केला. ती तुम्हाला फार समाधानकारक वाटली नसेल. पण संचांच्या क्रमिक संकल्पनेतून (किंवा ``विचारांच्या नियमां''तून) \stagesinf{} आपोआप स्वीकारणे भाग पडत नसेल, तर डेडेकिंड-अनंत संच अस्तित्वात आहे या दाव्याचे अंतर्गत समर्थन अजूनही मिळत नाही.

या विषयावर आणखी बरेच काही सांगता येईल. पण एखादे स्वयंसिद्धक किंवा तत्त्व समर्थित करण्यासाठी \emph{काय लागते}, याचा विचार सुरू करता येईल अशा टप्प्यावर तुम्ही आता पोहोचला असाल, अशी आशा आहे. पुढे आपण \stagesinf{} गृहीत धरू.







\section{$\Z^-$: एक महत्त्वाचा टप्पा}\label{sth:z:milestone:sec}

पुढील विभागात \stagesinf{} कडे पुन्हा वळू. पण अनंततेचे स्वयंसिद्धक मिळाल्यावर आपण एक महत्त्वाचा टप्पा गाठला आहे. आता $\Zminus$ या उपपत्तीसाठी लागणारी सर्व स्वयंसिद्धके आपल्याकडे आहेत. नेमके सांगायचे तर:

\begin{defn}
$\Zminus$ या उपपत्तीची स्वयंसिद्धके अशी आहेत: विस्तारात्मकता, संयोग, जोड्या, घातसंच, अनंतता आणि विभक्तीकरण-रूपबंधाची सर्व उदाहरणे.
\end{defn}

या नावातील \emph{झर्मेलो} म्हणजे झर्मेलो यांची संच उपपत्ती; \emph{वजा} केलेली गोष्ट पुढे पाहू. झर्मेलो यांना हे श्रेय उचित आहे, कारण त्यांनी 1908 मधील कामात ही उपपत्ती मूलत: मांडली होती.\footnote{इतिहास आणि तांत्रिक तपशीलांवरील रंजक टिप्पण्यांसाठी पॉटर (2004, परिशिष्ट A) पाहा.}

या उपपत्तीत गणिताचा फार मोठा भाग मांडता येतो. विशेषत: \ref{sfr:::part} पुन्हा पाहून, तेथे अनौपचारिकरीत्या केलेल्या प्रत्येक गोष्टीची अधिक औपचारिक मांडणी $\Zminus$ मध्ये करता येते, याची तुम्ही \emph{खात्री करून घ्यायला हवी}. (थोडे असे केल्यावर पुढील \ref{sth:z:arbintersections:sec} कडे थेट वळावेसे वाटू शकते.) म्हणून पुढे वेगळी नोंद न करता आपण किमान $\Zminus$ मध्ये काम करत आहोत, असे मानू.







\section{आपल्या नैसर्गिक संख्यांची निवड}\label{sth:z:nat:sec}





\ref{sth:z:infinity-again:defnomega} मध्ये आपण ``नैसर्गिक संख्या'' या शब्दप्रयोगाची स्पष्ट व्याख्या केली. हा निर्णय कसा समजायचा? तो आधिभौतिक दावा नाही; काही संचांना नैसर्गिक संख्या म्हणून \emph{वागवण्याचा} तो निर्णय आहे. असे मानण्याची कारणे आपण \ref{sfr:rel:ref:sec}, \ref{sfr:arith:ref:sec} आणि \ref{sfr:infinite:dedekindsproof:sec} मध्ये पाहिली. पण ती आणखी नेमकी सांगता येतील.

आपले अनंततेचे स्वयंसिद्धक फोन न्यूमन (1925) यांच्या पद्धतीनुसार आहे. मात्र त्याऐवजी पुढील स्वयंसिद्धकही स्वीकारता आले असते:

\begin{defish}
\emph{झर्मेलो यांचे 1908 मधील अनंततेचे स्वयंसिद्धक.} असा एक संच $A$ आहे की $\emptyset \in A$ आणि $(\forall x \in A)\{x\} \in A$.
\end{defish}

आपण फोन न्यूमन यांच्या धर्तीवरचे अनंततेचे स्वयंसिद्धक न घेता झर्मेलो यांचे घेतले असते, तरी डेडेकिंड-अनंत संच आणि म्हणून डेडेकिंड बीजसंरचना मिळाली असती. झर्मेलो यांच्या पद्धतीत आपल्या बीजसंरचनेचा निवडक घटक पुन्हा $\emptyset$ हाच असता ($0$ चे आपले प्रतिरूप); पण एकास-एक फलन $x \mapsto x \cup \{x\}$ ऐवजी $x \mapsto \{x\}$ असे असते. याचा सर्वांत सरळ परिणाम असा की झर्मेलो यांच्या पद्धतीत $2$ म्हणजे $\{\{\emptyset\}\}$, तर आपण फोन न्यूमन यांच्यासह $2$ म्हणजे $\{\emptyset, \{\emptyset\}\}$ मानतो.

अनंततेच्या या दोन स्वयंसिद्धकांपैकी एकच का निवडावे? मुख्य व्यावहारिक कारण असे की फोन न्यूमन यांची पद्धत अनंतपार संख्यांसाठी चांगली विस्तारते. \ref{sth:ordinals:chap} पासून हे पाहू. मात्र केवळ \emph{अंकगणित करण्याच्या} दृष्टीने दोन्ही पद्धती तितक्याच उपयोगी आहेत. म्हणून कोणी नैसर्गिक संख्या \emph{संचच आहेत} असे सांगितले, तर लगेच प्रश्न पडतो: \emph{त्या नेमक्या कोणत्या संच आहेत?}

हा नेमका प्रश्न Benacerraf (1965) यांनी प्रसिद्ध केला. पण आपण यापूर्वी अनेकदा पाहिलेल्या व्यापक बाबीचे ते केवळ सर्वांत प्रसिद्ध उदाहरण आहे. मूलभूत मुद्दा असा: संच उपपत्तीच्या साहाय्याने अनेक ``सहजगत्या'' मानल्या जाणाऱ्या प्रकारच्या वस्तूंचे \emph{अनुकरण} करता येते: वास्तव संख्या, परिमेय संख्या, पूर्णांक आणि नैसर्गिक संख्या; तसेच क्रमित जोड्या, फलने आणि संबंध. पण संच उपपत्ती अशा अनुकरणाची \emph{एकमेव} निवड देत नाही. त्याच कामाला सरळपणे तितकेच उपयोगी पडले असते असे पर्याय \emph{नेहमीच} असतात.







\section{परिशिष्ट: संवरण, गुणधर्माधारित संचरचना आणि छेद}\label{sth:z:arbintersections:sec}

\ref{sth:z:milestone:sec} मध्ये आपण \ref{sfr:::part} मधील अनौपचारिक काम मागे पाहून ते~$\Zminus$ मध्ये करता येते का हे तपासायला सांगितले. तसे केले असल्यास \emph{एक} मुद्दा अडला असेल: डेडेकिंड यांच्या \emph{संवरणांच्या} मांडणीत वापरलेला \emph{छेद}.

\ref{sfr:infinite:dedekind:Closure} आठवा:
\begin{align*}
  \closureofunder{f}{o} & = \bigcap\Setabs{X}{o \in X \text{ आणि $X$ हा
$f$-बंद आहे}}.
\intertext{याचे सर्वसाधारण रूप पुढीलप्रमाणे असते:}
  C & = \bigcap\Setabs{X}{\phi(X)}.
\end{align*}
पण यामुळे सावध व्हायला हवे. अनौपचारिक गुणधर्माधारित संचरचना अपयशी ठरते, म्हणून $\Setabs{X}{\phi(X)}$ अस्तित्वात आहे याची हमी नाही. मग अशा व्याख्यांत काही \emph{अवैध शॉर्टकट} घेतला जातो आहे असे वाटते.

सुदैवाने तसे नाही; किंवा सध्याच्या रूपात त्या व्याख्या \emph{तशा असल्या} तरी थोडे काटेकोर काम करून त्या वैध करता येतात. \ref{sth:z:sep:prop:intersectionsexist} मध्ये $A \neq \emptyset$ असताना $\bigcap A$ का अस्तित्वात असतो हे सांगताना या कामाची पूर्वसूचना दिली होती. आता ते स्पष्टपणे पाहू.

विस्तारात्मकता मानल्यास $C$ ची $\bigcap\Setabs{X}{\phi(X)}$ अशी व्याख्या करताना आपल्याला खरे तर पुढील अट पूर्ण करणारी वस्तू $C$ हवी असते:
\begin{equation}\label{sth:z:arbintersections:bicondelimarbintersection}
	\forall x(x \in C \liff \forall X(\phi(X) \lif x \in X))\tag{*}
\end{equation}
आता असा \emph{किमान एक} संच $S$ आहे की $\phi(S)$, असे मानू. मग \ref{sth:z:arbintersections:bicondelimarbintersection} मिळवण्यासाठी \emph{विभक्तीकरण} वापरून $C$ ची पुढीलप्रमाणे व्याख्या करता येते:
\[
	C = \Setabs{x \in S}{\forall X(\phi(X) \lif x \in X)}.
\]
या व्याख्येवरून अपेक्षेप्रमाणे \ref{sth:z:arbintersections:bicondelimarbintersection} मिळते, हे तपासण्याचा सराव करा.






याच सर्वसाधारण पद्धतीने छेदाच्या व्याख्यांत अनौपचारिक गुणधर्माधारित संचरचनेचा दिसणारा वापर टाळता येईल. ज्या उदाहरणामुळे ही चर्चा सुरू झाली, त्या $\closureofunder{f}{o}$ बाबतीत हे कसे करायचे ते पाहू. \ref{sfr:infinite:dedekind:closureproperties} च्या सिद्धतेच्या सुरुवातीला $o \in \ran{f}\cup\{o\}$ आणि $\ran{f} \cup \{o\}$ हा $f$-बंद संच आहे, असे आपण पाहिले. म्हणून अपेक्षित संचाची व्याख्या अशी करता येते:
\[
	\closureofunder{f}{o} = \Setabs{x \in \ran{f} \cup \{o\}}{(\forall X \ni o)(X \text{ हा $f$-बंद आहे} \lif x \in X)}.
\]



\chapter{क्रमसूचक संख्या}\label{sth:ordinals:chap}









\section{प्रस्तावना}\label{sth:ordinals:intro:sec}

\ref{sth:z:chap} मध्ये \stagesinf{} या रूपात पदानुक्रमात अनंत टप्पा आहे असे आपण गृहीत धरले (अनंततेचे स्वयंसिद्धकही पाहा). पण \stagessucc{} मानल्यावर त्या टप्प्यावर थांबता येत नाही; पुढे जावेच लागते. म्हणून पहिल्या अनंत टप्प्यानंतरच्या टप्प्यावर त्या पहिल्या अनंत टप्प्यावर उपलब्ध असलेल्या संचांपासून बनणारे सर्व शक्य संग्रह तयार होतात; मग पुन्हा तेच; पुन्हा; पुन्हा; \dots

यात अप्रकटपणे आपण संख्येची एक ``सहज'' संकल्पना वापरू लागलो आहोत: सर्व नैसर्गिक संख्यांच्या \emph{नंतरही} संख्या असू शकतात. विशेषतः येथे अभिप्रेत आहे ती \emph{अनंतपार क्रमसूचक संख्या}. या कल्पनेला अधिक काटेकोर रूप देणे हे या प्रकरणाचे उद्दिष्ट आहे. आपण क्रमसूचक संख्यांची सर्वसाधारण संकल्पना पाहू आणि मग काही विशिष्ट संचांची स्पष्ट व्याख्या आपल्या क्रमसूचक संख्या म्हणून करू.







\section{क्रमसूचक संख्येची सर्वसाधारण कल्पना}\label{sth:ordinals:idea:sec}

नैसर्गिक संख्यांचा नेहमीचा क्रम पाहा:
\begin{center}
	\begin{tikzpicture}
	\foreach \x/\xtext in {0, 1, 2, 3, 4, 5}
	{
		\node (\x a) at (\x, 1) {\small{$\x$}};
		\node (\x a) at (\x.5, 1) {\small{$<$}};
	}
	\node (ldots) at (6, 1) {\small{$\ldots$}};
	\end{tikzpicture}
\end{center}
याला तांत्रिक भाषेत $\omega$-अनुक्रम म्हणतात. संचांच्या पदानुक्रमातील सुरुवातीचे टप्पे बांधताना हीच सर्वसाधारण क्रमरचना दिसते. आता या अनुक्रमातील $0$ बाकी सर्व संख्यांच्या शेवटी नेऊ:
\begin{center}
	\begin{tikzpicture}
	\foreach \x/\xtext in {1, 2, 3, 4, 5}
	{
		\node (\x a) at (\x, 1) {\small{$\x$}};
		\node (\x a) at (\x.5, 1) {\small{$<$}};
	}
	\node (ldots) at (6, 1) {\small{$\ldots$}};
	\node (bea) at (6.5, 1) {\small{$<$}};
	\node (noa) at (7, 1) {\small{$0$}};
	\end{tikzpicture}
\end{center}
येथे घटक तेच आहेत; पण त्यांचा क्रम मुळात वेगळा आहे. पहिल्या क्रमाला शेवटचा घटक नव्हता; नव्या क्रमाला आहे. या नव्या क्रमात घटकांचा $\omega$-अनुक्रम ($1, 2, 3, 4, 5, \ldots$) आहे आणि त्यानंतर आणखी एक घटक येतो. म्हणून तो $\omega+1$-अनुक्रम आहे.

हेच घटक वापरून आणखी वेगवेगळ्या प्रकारचे क्रम बनवता येतात. उदाहरणार्थ, सर्व सम संख्या त्यांच्या नेहमीच्या क्रमात आणि त्यानंतर सर्व विषम संख्या त्यांच्या नेहमीच्या क्रमात घ्या:
\begin{center}
	\begin{tikzpicture}
	\node(a1) at (1,1) {\small{$0$}};
	\node(a2) at (1.5,1) {\small{$<$}};
	\node(a3) at (2,1) {\small{$2$}};
	\node(a4) at (2.5,1) {\small{$<$}};
	\node(a5) at (3,1) {\small{$4$}};
	\node(a6) at (3.5,1) {\small{$<$}};
	\node(dots) at (4,1) {\small{$\ldots$}};
	\node(aaoe) at (4.5,1) {\small{$<$}};
	\node(b1) at (5,1) {\small{$1$}};
	\node(b2) at (5.5,1) {\small{$<$}};
	\node(b3) at (6,1) {\small{$3$}};
	\node(b4) at (6.5,1) {\small{$<$}};
	\node(b5) at (7,1) {\small{$\ldots$}};
	\end{tikzpicture}
\end{center}
हा एका $\omega$-अनुक्रमानंतर दुसरा $\omega$-अनुक्रम आहे; म्हणजे $\omega+\omega$-अनुक्रम.

असे पुढेही करता येईल. पण \emph{क्रमरचनांविषयी} ही चर्चा समजून घेण्यासाठी आपल्याला एक सर्वसाधारण पद्धत हवी आहे.








\section{सुक्रमण}\label{sth:ordinals:wo:sec}

मूलभूत संकल्पना पुढीलप्रमाणे आहे:

\begin{defn}
$<\subseteq A\times A$ हा संबंध $A$ ला \emph{सुक्रमित करतो}
तेव्हा आणि केवळ तेव्हाच, जेव्हा पुढील दोन्ही अटी पूर्ण होतात:
\begin{enumerate}
\item $<$ हा जोडलेला असतो; म्हणजे सर्व $a, b \in A$ साठी $a < b$ किंवा $a = b$ किंवा $b < a$;
\item $A$ च्या प्रत्येक रिकाम्या नसलेल्या उपसंचात असा घटक असतो की त्या उपसंचातील कोणताही सदस्य त्याच्या $<$-आधी येत नाही; म्हणजे $\emptyset \neq X \subseteq A$ असल्यास $(\exists m \in
X)(\forall z \in X)z \nless m$.
\end{enumerate}
\end{defn}

आपण नुकत्याच पाहिलेल्या तिन्ही उदाहरणांतील क्रमसंबंध सुक्रमणच करतात, हे सहज दिसते.

\begin{prob}
\ref{sth:ordinals:idea:sec} मध्ये नैसर्गिक संख्यांवरील तीन क्रमांची उदाहरणे दिली. प्रत्येक क्रम सुक्रमण आहे, हे तपासा.
\end{prob}

सुक्रमणाबद्दल काही प्राथमिक पण अत्यंत महत्त्वाची निरीक्षणे पाहू.

\begin{prop}\label{sth:ordinals:wo:wo:strictorder}
$<$ हा $A$ ला सुक्रमित करत असेल, तर $A$ च्या प्रत्येक रिकाम्या नसलेल्या उपसंचात एकमेव $<$-लघुतम सदस्य असतो; शिवाय $<$ हा अपरावर्ती, असममित आणि संक्रमक असतो.
\end{prop}

\begin{proof}
$X$ हा $A$ चा रिकामा नसलेला उपसंच असेल, तर त्यात असा घटक $m$ असतो की $X$ मधील कोणताही सदस्य त्याच्या $<$-आधी येत नाही; म्हणजे $(\forall z \in X)z \nless m$. $<$ जोडलेला असल्याने $(\forall z \in X)m \leq z$. म्हणून $m$ हा $X$ मधील $<$-लघुतम घटक आहे.

अपरावर्तितेसाठी एखादा $a \in A$ निश्चित करू. $\{a\}$ चा $<$-लघुतम घटक $a$ असल्याने $a \nless a$. संक्रमकतेसाठी $a < b < c$ असेल, तर $\{a, b, c\}$ मध्ये $<$-लघुतम घटक असल्यामुळे $a < c$. अपरावर्तिता आणि संक्रमकता यांवरून असममिती मिळते.
\end{proof}

\begin{prop}\label{sth:ordinals:wo:propwoinduction}
$<$ हा $A$ ला सुक्रमित करत असेल, तर कोणत्याही सूत्र $\phi(x)$ साठी:
\[
\text{जर }(\forall a \in A)((\forall b < a)\phi(b) \lif
\phi(a))\text{, तर }(\forall a \in A)\phi(a).
\]
\end{prop}

\begin{proof}
व्यत्यस्त रूप सिद्ध करू. $\lnot(\forall a \in
A)\phi(a)$, म्हणजे $X = \Setabs{x \in A}{\lnot\phi(x)} \neq
\emptyset$, असे मानू. मग $X$ मध्ये पूर्ववर्ती नसलेला घटक $a$ असतो. म्हणून $(\forall
b < a)\phi(b)$; पण $\lnot \phi(a)$.
\end{proof}\noindent हा शेवटचा गुणधर्म नैसर्गिक संख्यांवरील प्रबल विगमनाच्या तत्त्वाची आठवण करून देतो: $(\forall n \in \omega)((\forall m < n)\phi(m)
\lif \phi(n))$ असेल, तर $(\forall n \in \omega)\phi(n)$. त्यामुळे सुक्रमण ही फार \emph{भक्कम} संकल्पना ठरते.\footnote{आठवण: स्पष्टपणे अन्यथा सांगितले नसेल, तर सर्व सूत्रांत प्राचले असू शकतात.}








\section{क्रम-समरूपणे}\label{sth:ordinals:iso:sec}

सुक्रमण किती \emph{भक्कम} आहे, हे स्पष्ट करण्यासाठी आपण सुक्रमणांची तुलना करण्याची एक पद्धत आधी मांडू.

\begin{defn}
$<$ हा $A$ ला सुक्रमित करत असेल, तर $\tuple{A, <}$ या जोडीला \emph{सुक्रमण} म्हणतात. $\tuple{A, <}$ आणि $\tuple{B,
\lessdot}$ ही सुक्रमणे \emph{क्रम-समरूपी} असतात तेव्हा आणि केवळ तेव्हाच, जेव्हा $f \colon A \to B$ असे एकास-एक आच्छादन असते की: $x < y$ तेव्हा आणि केवळ तेव्हाच $f(x) \lessdot  f(y)$. अशा वेळी $\ordeq{\tuple{A, <}}{\tuple{B, \lessdot}}$ असे लिहितो आणि $f$ ला \emph{क्रम-समरूपण} म्हणतो.
\end{defn}
\noindent
पुढे संक्षेपासाठी आपण ``क्रम-समरूपणे'' यांऐवजी ``समरूपणे'' असे म्हणू. सहज अर्थाने, समरूपणे म्हणजे रचना जपणारी एकास-एक आच्छादने असतात. त्यांच्याबद्दल काही साधी तथ्ये पुढे दिली आहेत.

\begin{lem}\label{sth:ordinals:iso:isoscompose}
समरूपणांची संयुक्तीही समरूपण असते. म्हणजे, $f \colon A
\to B$ आणि $g \colon B \to C$ ही समरूपणे असतील, तर $(g \circ f)
\colon A \to C$ हेही समरूपण असते.
\end{lem}
\begin{prob}
	\ref{sth:ordinals:iso:isoscompose} सिद्ध करा.
\end{prob}
\begin{proof}
हा पुरावा सरावासाठी सोडला आहे.
\end{proof}
\begin{cor}\label{sth:ordinals:iso:ordisoisequiv}
	$\ordeq{X}{Y}$ हा समतुल्यता संबंध आहे.
\end{cor}
\begin{prop}\label{sth:ordinals:iso:ordisounique}
$\tuple{A, <}$ आणि $\tuple{B, \lessdot}$ ही समरूपी सुक्रमणे असतील, तर त्यांच्यामधील समरूपण एकमेव असते.
\end{prop}

\begin{proof}
$f$ आणि $g$ ही $A \to B$ समरूपणे मानू. \ref{sth:ordinals:wo:propwoinduction} वापरून विगमनाने हा निष्कर्ष सिद्ध करू.
$a\in A$ निश्चित करू आणि विगमनाच्या गृहीतकानुसार $(\forall b < a)f(b) = g(b)$ मानू. $x \in B$ निश्चित करू.

$x \lessdot f(a)$ असेल, तर $f^{-1}(x) < a$. $f$ आणि $g$ ही समरूपणे असल्यामुळे $g(f^{-1}(x)) \lessdot
g(a)$. पण $f^{-1}(x) < a$ असल्याने आपल्या गृहीतकावरून $x =f(f^{-1}(x)) = g(f^{-1}(x))$ मिळते. म्हणून $x \lessdot g(a)$. त्याचप्रमाणे, $x \lessdot g(a)$ असल्यास $x \lessdot
f(a)$.

म्हणून $(\forall x \in B)(x \lessdot f(a) \liff x \lessdot
g(a))$. \ref{sfr:rel:ord:prop:extensionality-strictlinearorders} वरून $f(a) = g(a)$ मिळते. त्यामुळे \ref{sth:ordinals:wo:propwoinduction} वरून $(\forall
a \in A)f(a) = g(a)$.
\end{proof}\noindent
यावरून सुक्रमणाचे भक्कम स्वरूप काहीसे स्पष्ट होते. ते अधिक समजण्यासाठी आणखी काही संज्ञा निश्चित करू.

\begin{defn}
$\tuple{A, <}$ हे सुक्रमण आणि $a \in A$ असेल, तर $A_a = \Setabs{x \in A}{x
< a}$ असे मानू. $A_a$ हा $A$ चा यथार्थ \emph{आरंभीचा खंड} आहे असे म्हणू; $A$ स्वतःही $A$ चा अयथार्थ आरंभीचा खंड असू शकतो. $<_a$ हा आरंभीच्या खंडापुरता $<$ चा संकोच, म्हणजे $\funrestrictionto{\mathord{<}}{A_a^2}$, मानू.
\end{defn}
\noindent
या संज्ञांद्वारे कोणतेही सुक्रमण स्वतःच्या यथार्थ आरंभीच्या खंडाशी समरूपी नसते, हे मांडता आणि सिद्ध करता येते.

\begin{lem}\label{sth:ordinals:iso:wellordnotinitial}
$\tuple{A, <}$ हे सुक्रमण आणि $a \in A$ असेल, तर
$\ordneq{\tuple{A, <}}{\tuple{A_a, <_a}}$.
\end{lem}

\begin{proof}
विसंगतीसाठी $f \colon A \to A_a$ हे समरूपण आहे असे मानू. $f$ हे एकास-एक आच्छादन असून $A_a \subsetneq A$ असल्याने, \ref{sth:ordinals:wo:wo:strictorder} वापरून $A$ मधील $b \neq f(b)$ पूर्ण करणारा $<$-लघुतम घटक $b \in A$ घेऊ. $(\forall x \in A)(x<b \liff x < f(b))$ हे दाखवू. मग \ref{sfr:rel:ord:prop:extensionality-strictlinearorders} वरून $b = f(b)$ मिळेल आणि विसंगती सिद्ध होईल.

$x < b$ मानू. $b$ च्या निवडीवरून $x = f(x)$. तसेच $f$ समरूपण असल्याने $f(x) < f(b)$. म्हणून $x < f(b)$.

$x < f(b)$ मानू. $f$ समरूपण असल्याने $f^{-1}(x) < b$; म्हणून $b$ च्या निवडीवरून $f^{-1}(x) = x$. त्यामुळे $x < b$.
\end{proof}

पुढील निष्कर्ष साधारणपणे असे सांगतो की समरूपणाचा ``आरंभीचा खंड'' हाही समरूपण असतो:

\begin{lem}\label{sth:ordinals:iso:wellordinitialsegment}
$\tuple{A, <}$ आणि $\tuple{B, \lessdot}$ ही सुक्रमणे मानू. $f
\colon A \to B$ हे समरूपण आणि $a \in A$ असेल, तर
$\funrestrictionto{f}{A_{a}} : A_a \to B_{f(a)}$ हे समरूपण असते.
\end{lem}

\begin{proof}
$f$ हे समरूपण असल्यामुळे:

\begin{align*}
	\funimage{f}{A_a} &= \funimage{f}{\Setabs{x \in A}{x < a}}\\
	&= \funimage{f}{\Setabs{f^{-1}(y) \in A}{f^{-1}(y) < a}} \\
	&= \Setabs{y \in B}{y \lessdot f(a)} \\
	&=B_{f(a)}
\end{align*}
मिळते. $f$ क्रम जपत असल्यामुळे $\funrestrictionto{f}{A_a}$ सुद्धा क्रम जपते.
\end{proof}

पुढील दोन निष्कर्षांवरून कोणतीही दोन सुक्रमणे तुलना करता येतात, हे सिद्ध होते:

\begin{lem}\label{sth:ordinals:iso:lemordsegments}
$\tuple{A, <}$ आणि $\tuple{B, \lessdot}$ ही सुक्रमणे मानू. जर
$\ordeq{\tuple{A_{a_1}, <_{a_1}}}{\tuple{B_{b_1}, \lessdot_{b_1}}}$
आणि $\ordeq{\tuple{A_{{a_2}}, <_{a_2}}}{\tuple{B_{{b_2}},
\lessdot_{b_2}}}$ असेल, तर ${a_1}  < {a_2} \text{ तेव्हा आणि केवळ तेव्हाच }{b_1} \lessdot
{b_2}$.
\end{lem}

\begin{proof}
\emph{डावीकडून उजवीकडे} हे सिद्ध करू; उलट दिशा तशीच आहे.
$\ordeq{\tuple{A_{a_1}, <_{a_1}}}{\tuple{B_{b_1},
\lessdot_{b_1}}}$ आणि $\ordeq{\tuple{A_{{a_2}},
<_{a_2}}}{\tuple{B_{{b_2}}, \lessdot_{b_2}}}$ दोन्ही मानू आणि $f \colon
A_{{a_2}} \to B_{{b_2}}$ हे समरूपण असू द्या. ${a_1} < {a_2}$ असेल, तर
\ref{sth:ordinals:iso:wellordinitialsegment} वरून
$\ordeq{\tuple{A_{a_1}, <_{a_1}}}{\tuple{B_{f({a_1})},
\lessdot_{f({a_1})}}}$ मिळते. म्हणून
$\ordeq{\tuple{B_{b_1}, \lessdot_{b_1}}}{\tuple{B_{f({a_1})},
\lessdot_{f({a_1})}}}$; आणि \ref{sth:ordinals:iso:wellordnotinitial} वरून ${b_1} = f({a_1})$. आता $f$ चे सहप्रांत $B_{{b_2}}$ असल्याने ${b_1} \lessdot {b_2}$.
\end{proof}

\begin{thm}\label{sth:ordinals:iso:thm:woalwayscomparable}
कोणतीही दोन सुक्रमणे दिल्यास, त्यांपैकी एक दुसऱ्याच्या आरंभीच्या खंडाशी समरूपी असते; तो खंड यथार्थ असलाच पाहिजे असे नाही.
\end{thm}

\begin{proof}
$\tuple{A, <}$ आणि $\tuple{B, \lessdot}$ ही सुक्रमणे मानू. विभक्तीकरण वापरून
\[
	f = \Setabs{\tuple{a, b} \in A \times B}{
		\ordeq{\tuple{A_a, <_a}}{\tuple{B_b, \lessdot_b}}}.
\]
असे ठरवू. \ref{sth:ordinals:iso:lemordsegments} वरून $\tuple{a_1, b_1}, \tuple{a_2, b_2} \in f$ अशा सर्व जोड्यांसाठी $a_1 < a_2$ तेव्हा आणि केवळ तेव्हाच $b_1 \lessdot b_2$. म्हणून $f \colon \dom{f} \to
\ran{f}$ हे समरूपण आहे.

$a_2 \in \dom{f}$ आणि $a_1 < a_2$ असेल, तर \ref{sth:ordinals:iso:wellordinitialsegment} वरून $a_1 \in \dom{f}$; म्हणून $\dom{f}$ हा $A$ चा आरंभीचा खंड आहे. त्याचप्रमाणे $\ran{f}$ हा $B$ चा आरंभीचा खंड आहे. विसंगतीसाठी हे दोन्ही \emph{यथार्थ} आरंभीचे खंड आहेत असे मानू. मग $A \setminus \dom{f}$ मधील $<$-लघुतम घटक $a$ घेऊ; त्यामुळे $\dom{f} = A_a$. तसेच $B \setminus
\ran{f}$ मधील $\lessdot$-लघुतम घटक $b$ घेऊ; त्यामुळे $\ran{f} = B_b$. म्हणून $f \colon A_a \to B_b$ हे समरूपण आहे आणि त्यामुळे $\tuple{a, b} \in f$; ही विसंगती आहे.
\end{proof}







\section{क्रमसूचक संख्यांची फोन न्यूमनची रचना}\label{sth:ordinals:vn:sec}

\ref{sth:ordinals:iso:thm:woalwayscomparable} वरून एक कल्पना सुचते. सुक्रमणांचे प्रतिनिधी म्हणून \emph{क्रमप्रकार} नावाच्या काही वस्तू आपण मांडू शकतो. $\tuple{A, <}$ या सुक्रमणाचा क्रमप्रकार $\ordtype{A, <}$ असा लिहू; मग पुढील दोन तत्त्वे मिळतील अशी अपेक्षा असेल:
\begin{align*}
	\ordtype{A, <} = \ordtype{B, \lessdot} &
	\text{ तेव्हा आणि केवळ तेव्हाच } \ordeq{\tuple{A, <}}{\tuple{B, \lessdot}}\\
	\ordtype{A, <} < \ordtype{B, \lessdot}&
	\text{ तेव्हा आणि केवळ तेव्हाच }\ordeq{\tuple{A, <}}{\tuple{B_b, \lessdot_b}}\text{ असे काही }b \in B
\end{align*}
शिवाय, नैसर्गिक संख्या जशा विशिष्ट संचांच्या रूपात मांडता येतात, तसे क्रमप्रकारही \emph{विशिष्ट संचांच्या रूपात} मांडता येतील अशी अपेक्षा असू शकते.

ते करण्याची सर्वाधिक प्रचलित पद्धत---आणि आपण अनुसरणार असलेली पद्धत---म्हणजे विशिष्ट \emph{प्रमाणभूत} सुक्रमित संचांच्या आधारे हे क्रमप्रकार परिभाषित करणे. असे प्रमाणभूत संच प्रथम फोन न्यूमन यांनी मांडले:

\begin{defn}
$(\forall x \in A)x \subseteq A$ असेल तेव्हा आणि केवळ तेव्हाच $A$ हा संच \emph{संक्रमक} असतो. $A$ हा संक्रमक असून त्याच्या सदस्यांपुरत्या $\in$ च्या बंधनाने सुक्रमित असेल तेव्हा आणि केवळ तेव्हाच $A$ ही \emph{क्रमसूचक संख्या} असते.
\end{defn}
\noindent
यापुढे क्रमसूचक संख्यांसाठी ग्रीक अक्षरे वापरू. व्याख्येवरून लगेच दिसते की $\alpha$ ही क्रमसूचक संख्या असल्यास $\tuple{\alpha, \in_\alpha}$ हे सुक्रमण आहे; येथे $\in_\alpha =
\Setabs{\tuple{x, y} \in \alpha^2}{x \in y}$. म्हणून संकेतलेखनात किंचित मोकळीक घेऊन $\alpha$ \emph{स्वतःच} एक सुक्रमण आहे असेही म्हणू शकतो.

पहिल्या काही क्रमसूचक संख्या पुढीलप्रमाणे आहेत:
\[
	\emptyset, \{\emptyset\},
	\{\emptyset, \{\emptyset\}\},
	\{\emptyset, \{\emptyset\}, \{\emptyset, \{\emptyset\}\}\}, \ldots
\]
अनंततेच्या स्वयंसिद्धकात, म्हणजे फोन न्यूमन यांच्या $\omega$ च्या व्याख्येत, आपल्याला भेटलेल्या पहिल्या काही क्रमसूचक संख्या याच आहेत (\ref{sth:z:infinity-again:sec} पाहा). हा योगायोग नाही. फोन न्यूमन यांच्या व्याख्येत नैसर्गिक संख्या क्रमसूचक संख्या असतात; पण अनंतपार क्रमसूचक संख्यांनाही जागा असते.

नेहमीप्रमाणे आता विचारता येईल: \emph{याच} क्रमसूचक संख्या आहेत का? की फोन न्यूमन यांनी आपल्याला असे काही संच दिले आहेत ज्यांना आपण क्रमसूचक संख्या \emph{मानू शकतो}? या प्रश्नावरील चर्चा \ref{sfr:rel:ref:sec}, \ref{sfr:arith:ref:sec}, \ref{sfr:infinite:dedekindsproof:sec} आणि \ref{sth:z:nat:sec} मधील चर्चांप्रमाणेच आहे; म्हणून येथे ती पुन्हा विस्ताराने करणार नाही. पुढे ``क्रमसूचक संख्या'' म्हटल्यावर आपण ``फोन न्यूमन यांच्या क्रमसूचक संख्या'' अभिप्रेत धरू.







\section{क्रमसूचक संख्यांचे मूलभूत गुणधर्म}\label{sth:ordinals:basic:sec}

पहिल्या काही क्रमसूचक संख्या नैसर्गिक संख्याच आहेत, हे आपण पाहिले. क्रमसूचक संख्यांची उपपत्ती विकसित करण्यामागचा मुख्य उद्देश म्हणजे नैसर्गिक संख्यांसाठी लागू असलेले विगमनाचे तत्त्व पुढे नेणे. काही प्राथमिक निष्कर्षांच्या साहाय्याने आपण ते साध्य करू.

\begin{lem}\label{sth:ordinals:basic:ordmemberord}
क्रमसूचक संख्येचा प्रत्येक घटक ही क्रमसूचक संख्या असतो.
\end{lem}

\begin{proof}
$\alpha$ ही क्रमसूचक संख्या आणि $b \in \alpha$ असे मानू. $\alpha$ संक्रमक असल्याने $b \subseteq \alpha$. त्यामुळे $\in$ जसा $\alpha$ ला सुक्रमित करतो, तसाच $\in$ हा $b$ लाही सुक्रमित करतो.

$b$ संक्रमक आहे हे पाहण्यासाठी $x \in c \in b$ मानू. $b
\subseteq \alpha$ असल्याने $c \in \alpha$. पुन्हा $\alpha$ संक्रमक असल्याने $c \subseteq
\alpha$ आणि म्हणून $x \in \alpha$. अशा प्रकारे $x, c, b \in \alpha$. पण $\in$ हा $\alpha$ ला सुक्रमित करतो; म्हणून \ref{sth:ordinals:wo:wo:strictorder} वरून $\in$ हा $\alpha$ वरील संक्रमक संबंध आहे. त्यामुळे $x
\in c \in b$ वरून $x \in b$. याचे सामान्यीकरण केल्यास $c \subseteq b$.
\end{proof}

\begin{cor}\label{sth:ordinals:basic:ordissetofsmallerord}
कोणत्याही क्रमसूचक संख्या~$\alpha$ साठी $\alpha = \Setabs{\beta \in \alpha}{\beta \text{ ही क्रमसूचक संख्या आहे}}$.
\end{cor}

\begin{proof}
\ref{sth:ordinals:basic:ordmemberord} वरून थेट मिळते.
\end{proof}

पुढील दोन मुख्य निष्कर्षांचा, \ref{sth:ordinals:basic:ordinductionschema} आणि \ref{sth:ordinals:basic:ordtrichotomy} यांचा, साधारण आशय असा की क्रमसूचक संख्याच सदस्यत्वाने सुक्रमित होतात:

\begin{thm}[अनंतपार विगमन]\label{sth:ordinals:basic:ordinductionschema}
कोणत्याही सूत्र $\phi(x)$ साठी:
\[
	\text{जर }\exists \alpha \phi(\alpha)\text{, तर }\exists \alpha(\phi(\alpha)
	\land  (\forall \beta \in \alpha) \lnot \phi(\beta))
\]
येथे दाखवलेली परिमाणके अव्यक्तपणे क्रमसूचक संख्यांपुरती मर्यादित आहेत.
\end{thm}

\begin{proof}

एखाद्या क्रमसूचक संख्या $\alpha$ साठी $\phi(\alpha)$ मानू. $ (\forall \beta
\in \alpha) \lnot \phi(\beta)$ असेल, तर सिद्धता झाली. अन्यथा, $\alpha$ ही क्रमसूचक संख्या असल्याने तिच्यात $\phi$ पूर्ण करणारा $\in$-लघुतम घटक मिळतो; \ref{sth:ordinals:basic:ordmemberord} वरून तोही क्रमसूचक संख्या आहे.






\end{proof}
\noindent
\ref{sth:ordinals:basic:ordinductionschema} हेच तत्त्व पुढील रूपबंधातही मांडता येते:
\[
\text{जर }\forall \alpha((\forall \beta \in \alpha)\phi(\beta) \lif
\phi(\alpha))\text{, तर }\forall \alpha\phi(\alpha)
\]
यासाठी \ref{sth:ordinals:basic:ordinductionschema} मध्ये $\lnot\phi(\alpha)$ घेऊन प्राथमिक तार्किक रूपांतरे करावी लागतात.

\begin{thm}[त्रिपर्याय]\label{sth:ordinals:basic:ordtrichotomy}
कोणत्याही क्रमसूचक संख्या $\alpha$ आणि $\beta$ साठी $\alpha \in \beta \lor \alpha = \beta \lor \beta \in \alpha$.
\end{thm}

\begin{proof}
पुरावा द्विगुण विगमनाने, म्हणजे \ref{sth:ordinals:basic:ordinductionschema} दोनदा वापरून, करता येतो. $x \in y \lor x = y \lor y \in x$ असेल तेव्हा आणि केवळ तेव्हाच $x$ आणि $y$ \emph{तुलनीय} आहेत असे म्हणू.

विगमनासाठी, $\alpha$ मधील प्रत्येक क्रमसूचक संख्या \emph{प्रत्येक} क्रमसूचक संख्येशी तुलनीय आहे असे मानू. आणखी एका विगमनासाठी, $\alpha$ ही $\beta$ मधील प्रत्येक क्रमसूचक संख्येशी तुलनीय आहे असे मानू. आता $\alpha$ ही $\beta$ शी तुलनीय आहे हे दाखवू. मग $\beta$ वरील विगमनाने $\alpha$ प्रत्येक क्रमसूचक संख्येशी तुलनीय ठरेल; आणि $\alpha$ वरील विगमनाने \emph{प्रत्येक} क्रमसूचक संख्या \emph{प्रत्येक} क्रमसूचक संख्येशी तुलनीय ठरेल, हेच हवे आहे. $\alpha \notin \beta$ आणि $\beta \notin
\alpha$ असे मानून $\alpha = \beta$ दाखवणे पुरेसे आहे.

$\alpha \subseteq \beta$ दाखवण्यासाठी $\gamma \in \alpha$ निश्चित करू; \ref{sth:ordinals:basic:ordmemberord} वरून ती क्रमसूचक संख्या आहे. म्हणून पहिल्या विगमन गृहीतकानुसार $\gamma$ ही $\beta$ शी तुलनीय आहे. पण $\gamma
= \beta$ किंवा $\beta \in \gamma$ यांपैकी काहीही असेल, तर $\beta \in \alpha$ मिळते (गरज पडल्यास $\alpha$ च्या संक्रमकतेचा वापर करून); हे आपल्या गृहीतकाविरुद्ध आहे. म्हणून $\gamma \in \beta$. याचे सामान्यीकरण केल्यास $\alpha \subseteq
\beta$.

दुसरे विगमन गृहीतक वापरून अगदी त्याच प्रकारच्या तर्काने $\beta \subseteq \alpha$ मिळते. म्हणून $\alpha = \beta$.
\end{proof}\noindent त्यामुळे कधी कधी $\alpha \in \beta$ ऐवजी $\alpha <\beta$ असे लिहू; कारण $\in$ येथे क्रमसंबंधाप्रमाणे वागतो. यामागे परिचय आणि $\alpha \in
\beta \lor \alpha = \beta$ याऐवजी $\alpha \leq \beta$ लिहिण्याची सोय यांपलीकडे विशेष कारण नाही.\footnote{$\alpha
\mathrel{\underline{\in}} \beta$ असेही लिहिता येईल; पण हा संकेत सर्वमान्य नाही.}

आपल्या मागील निष्कर्षांवरून पुढील दोन परिणाम लगेच मिळतात; पहिल्यात आपला नवा संकेत वापरला आहे:

\begin{cor}\label{sth:ordinals:basic:ordordered}
$\exists \alpha\phi(\alpha)$ असेल, तर $\exists \alpha(\phi(\alpha)
\land \forall \beta(\phi(\beta) \lif \alpha \leq \beta))$.
शिवाय, कोणत्याही क्रमसूचक संख्या $\alpha, \beta, \gamma$ साठी $\alpha
\notin \alpha$ आणि $\alpha \in \beta \in \gamma \lif \alpha \in
\gamma$.
\end{cor}

\begin{proof}
\ref{sth:ordinals:wo:wo:strictorder} प्रमाणेच.
\end{proof}

\begin{prob}
\ref{sth:ordinals:basic:ordtrichotomy} च्या पुराव्यातील ``अगदी त्याच प्रकारचा तर्क'' पूर्ण लिहा.
\end{prob}

\begin{cor}\label{sth:ordinals:basic:corordtransitiveord}
$A$ ही क्रमसूचक संख्या असते तेव्हा आणि केवळ तेव्हाच $A$ हा क्रमसूचक संख्यांचा संक्रमक संच असतो.
\end{cor}

\begin{proof}
\emph{डावीकडून उजवीकडे.} \ref{sth:ordinals:basic:ordmemberord} वरून. \emph{उजवीकडून डावीकडे.} $A$ हा क्रमसूचक संख्यांचा संक्रमक संच असेल, तर \ref{sth:ordinals:basic:ordinductionschema} आणि \ref{sth:ordinals:basic:ordtrichotomy} वरून $\in$ हा $A$ ला सुक्रमित करतो.
\end{proof}

\ref{sth:ordinals:basic:ordinductionschema} आणि \ref{sth:ordinals:basic:ordtrichotomy} यांचा आशय आपण क्रमसूचक संख्या $\in$ ने सुक्रमित होतात असा सांगितला. पण पुढील निष्कर्षामुळे असे विधान करताना \emph{फार सावध} राहिले पाहिजे:

\begin{thm}[बुराली-फोर्ती विरोधापत्ती]\label{sth:ordinals:basic:buraliforti}
सर्व क्रमसूचक संख्यांचा संच अस्तित्वात नाही.
\end{thm}

\begin{proof}
विसंगतीसाठी $O$ हा सर्व क्रमसूचक संख्यांचा संच आहे असे मानू. $\alpha \in
\beta \in O$ असेल, तर \ref{sth:ordinals:basic:ordmemberord} वरून $\alpha$ ही क्रमसूचक संख्या आहे आणि म्हणून $\alpha \in O$. त्यामुळे $O$ संक्रमक आहे; म्हणून \ref{sth:ordinals:basic:corordtransitiveord} वरून $O$ ही क्रमसूचक संख्या ठरते. मग $O \in O$, हे \ref{sth:ordinals:basic:ordordered} ला विरोधी आहे.
\end{proof}
\noindent
हा निष्कर्ष बुराली-फोर्ती यांच्या नावाने ओळखला जातो. परंतु सर्व क्रमसूचक संख्यांचा संच आहे असे मानल्यास \emph{विसंगती} निर्माण होते, हे १८९९ मध्ये डेडेकिंड यांना लिहिलेल्या पत्रात प्रथम स्पष्टपणे कांटोर यांच्या लक्षात आले. व्हान हायेनूर्ट यांनी सांगितल्याप्रमाणे:
\begin{quote}
  बुराली-फोर्ती यांनी स्वतः या विसंगतीतून \emph{विसंगतीद्वारे सिद्धतेने} असा निष्कर्ष काढला की क्रमसूचक संख्यांचा नैसर्गिक क्रम हा केवळ आंशिक क्रम आहे.
  (व्हान हायेनूर्ट 1967, पृ.~105)
\end{quote}
बुराली-फोर्ती यांची चूक बाजूला ठेवून आतापर्यंतचे निष्कर्ष असे सांगता येतील: क्रमसूचक संख्या या प्रत्येकी सदस्यत्वाने सुक्रमित संच असतात आणि एकत्रितपणेही सदस्यत्वाने सुक्रमित असतात; मात्र त्या एकत्रितपणे संच बनवत नाहीत.

हा भाग पूर्ण करण्यासाठी \ref{sth:ordinals:basic:ordinductionschema} आणि \ref{sth:ordinals:basic:ordtrichotomy} वरून मिळणारे क्रमसूचक संख्यांचे आणखी काही मूलभूत गुणधर्म पाहू.

\begin{prop}
क्रमसूचक संख्यांचा कोणताही काटेकोरपणे उतरणारा अनुक्रम परिमित असतो.
\end{prop}

\begin{proof}
क्रमसूचक संख्यांचा अनंत काटेकोरपणे उतरणारा अनुक्रम $\alpha_0 > \alpha_1 > \alpha_2 > \ldots$ असता, तर त्याच्या सदस्यांत $<$-लघुतम सदस्य नसता; हे \ref{sth:ordinals:basic:ordinductionschema} शी विसंगत आहे.
\end{proof}

\begin{prop}\label{sth:ordinals:basic:ordinalsaresubsets}
कोणत्याही क्रमसूचक संख्या $\alpha, \beta$ साठी $\alpha \subseteq \beta \lor \beta \subseteq \alpha$.
\end{prop}

\begin{proof}
$\alpha \in \beta$ असेल, तर $\beta$ संक्रमक असल्यामुळे $\alpha \subseteq \beta$. त्याचप्रमाणे $\beta \in \alpha$ असेल, तर $\beta \subseteq
\alpha$. आणि $\alpha = \beta$ असेल, तर $\alpha \subseteq \beta$ आणि $\beta \subseteq \alpha$. म्हणून \ref{sth:ordinals:basic:ordtrichotomy} वरून निष्कर्ष मिळतो.
\end{proof}

\begin{prop}\label{sth:ordinals:basic:ordisoidentity}
कोणत्याही क्रमसूचक संख्या $\alpha, \beta$ साठी $\alpha = \beta$ तेव्हा आणि केवळ तेव्हाच $\ordeq{\alpha}{\beta}$.
\end{prop}

\begin{proof}
क्रमसूचक संख्या या सुक्रमणे आहेत; म्हणून त्रिपर्याय (\ref{sth:ordinals:basic:ordtrichotomy}) आणि \ref{sth:ordinals:iso:wellordnotinitial} वरून हे लगेच मिळते.
\end{proof}

\begin{prob}\label{sth:ordinals:basic:probunionordinalsordinal}
$X$ चा प्रत्येक सदस्य क्रमसूचक संख्या असेल, तर $\bigcup X$ ही क्रमसूचक संख्या असते, हे सिद्ध करा.
\end{prob}








\section{प्रतिस्थापन}\label{sth:ordinals:replacement:sec}

\ref{sth:ordinals:vn:sec} मध्ये क्रमसूचक संख्यांना क्रमप्रकार म्हणून, म्हणजे सुक्रमणांचे प्रमाणभूत प्रतिनिधी म्हणून, घेता येईल असे सुचवून त्यांची ओळख करून दिली. हे साधण्यासाठी \emph{प्रत्येक सुक्रमण कोणत्या तरी क्रमसूचक संख्येशी समरूपी आहे} हे सिद्ध करावे लागेल. मग $\tuple{A, <} \isomorphic \alpha$ अशी क्रमसूचक संख्या $\alpha$ घेऊन $\ordtype{A, <}$ परिभाषित करता येईल.

दुर्दैवाने, आतापर्यंत आपण दिलेल्या स्वयंसिद्धकांवरून हा अपेक्षित निष्कर्ष \emph{सिद्ध करता येत नाही}. का ते \ref{sth:replacement:strength:sec} मध्ये पाहू; सध्या इतके लक्षात ठेवू की तो सिद्ध होत नाही. म्हणून आपल्याला एक नवे तत्त्व हवे आहे:

\begin{axiom}[प्रतिस्थापन-रूपबंध]
कोणत्याही सूत्र $\phi(x, y)$ साठी पुढील विधान हे एक स्वयंसिद्धक आहे:
\begin{quote}
	कोणत्याही $A$ साठी, $(\forall x \in A)\lexists![y][\phi(x,y)]$ असेल, तर $\Setabs{y}{(\exists x \in A)\phi(x,y)}$ अस्तित्वात असतो.
\end{quote}
\end{axiom}
\noindent
विभक्तीकरणाप्रमाणे हाही एक रूपबंध आहे: अनंत संख्येने असलेल्या निरनिराळ्या $\phi$ साठी त्यातून अनंत संख्येने स्वयंसिद्धके मिळतात. हेच तत्त्व पुढीलप्रमाणेही लिहिता येते, आणि सामान्यतः तसेच लिहिले जाते:

\begin{defish}
``$B$'' नसलेल्या कोणत्याही सूत्र $\phi(x,y)$ साठी पुढील विधान हे एक स्वयंसिद्धक आहे:
\[
\forall A[(\forall x \in A)\lexists![y][\phi(x,y)] \lif \exists B\forall y (y \in B \liff (\exists x \in A)\phi(x,y))]
\]
\end{defish}

पहिल्यांदा पाहताना हे सूत्र गुंतागुंतीचे वाटते. प्रतिस्थापनाचा पुढील सोपा परिणाम मात्र त्यामागील कल्पना \emph{अधिक स्पष्टपणे} व्यक्त करतो:\footnote{पद म्हणजे अशी अभिव्यक्ती की जिची मुक्त चरे कोणतीही भरली तरी ती नेमकी एकच वस्तू दर्शवते. उदाहरणार्थ, ``$\emptyset$'' हे रिकामा संच दर्शवणारे पद आहे; ``$\{x\}$'' हे $x$ चा एकसदस्यीय संच दर्शवणारे पद आहे ($x$ काहीही असो); ``$x \cup y$'' हे $x$ आणि $y$ यांचा संयोग दर्शवणारे पद आहे (ते काहीही असोत).}

\begin{cor}
कोणत्याही पद $\tau(x)$ आणि कोणत्याही संच $A$ साठी पुढील संच अस्तित्वात असतो:
\[
	\Setabs{\tau(x)}{x \in A} = \Setabs{y}{(\exists x \in A)y = \tau(x)}.
\]
\end{cor}

\begin{proof}
$\tau$ हे \emph{पद} असल्याने $\forall x \lexists![y][\tau(x) = y]$. विशेषतः $(\forall x \in A)\lexists![y][\tau(x) = y]$. म्हणून प्रतिस्थापनावरून $\Setabs{y}{(\exists x \in A)\tau(x) = y}$ अस्तित्वात असतो.
\end{proof}
\noindent
यावरून या स्वयंसिद्धकाला ``प्रतिस्थापन'' हे नाव योग्य वाटते: संच $A$ दिला असता, $A$ च्या प्रत्येक सदस्याच्या जागी $\tau$ अंतर्गत त्याची प्रतिमा ठेवून $\Setabs{\tau(x)}{x \in A}$ हा नवा संच तयार करता येतो. फलनाखालील संचाच्या प्रतिमेचा संकेत अनुसरून $\Setabs{\tau(x)}{x \in A}$ साठी $\funimage{\tau}{A}$ असेही लिहिता येईल.

पण येथे महत्त्वाची बाब म्हणजे $\tau$ हे \emph{पद} आहे. \ref{sfr:fun:rel:sec} मधील अर्थाने, म्हणजे क्रमित जोड्यांचा विशिष्ट संच या अर्थाने, ते \emph{फलनाचे} नाव असणे आवश्यक नाही. कारण $f$ हे त्या अर्थाने फलन असेल आणि $A\subseteq\dom{f}$,
तर $\funimage{f}{A}=\Setabs{f(x)}{x\in A}$ हा संच $\ran{f}$ चा
केवळ विशिष्ट उपसंच आहे. मनमाना $A$ घेतल्यास $f(x)$ हा संकेत
फक्त $A\cap\dom{f}$ मधील $x$ साठी वापरता येतो; प्रतिमा
क्रमित जोड्यांच्या आलेखावरून त्याच बंधनाने घेतली जाते. त्याचे अस्तित्व तर फक्त $\Zminus$ च्या स्वयंसिद्धकांवरून आधीच मिळते.\footnote{$\Setabs{y\in\bigcup\bigcup f}{(\exists x\in A)\tuple{x,y}\in f}$ पहा.} याउलट, प्रतिस्थापन हे आपल्या स्वयंसिद्धकांमध्ये \emph{सामर्थ्यवान} भर घालते; हे \ref{sth:replacement:chap} मध्ये दिसेल.






\section{$\ZFminus$: एक महत्त्वाचा टप्पा}\label{sth:ordinals:zfm:sec}

प्रतिस्थापनाला समर्थन कसे द्यावे, किंवा देता येईलच का, हा प्रश्न सरळ नाही. म्हणून तो \ref{sth:replacement:chap} साठी राखून ठेवू. पण प्रतिस्थापन जोडल्यामुळे आपण आणखी एक महत्त्वाचा टप्पा गाठला आहे. आता $\ZFminus$ या उपपत्तीसाठी लागणारी सर्व स्वयंसिद्धके आपल्याकडे आहेत. ती अशी:

\begin{defn}
$\ZFminus$ उपपत्तीत पुढील स्वयंसिद्धके आहेत: विस्तारात्मकता, संयोग, जोड्या, घातसंच, अनंतता, आणि विभक्तीकरण व प्रतिस्थापन या रूपबंधांतील सर्व स्वयंसिद्धके. दुसऱ्या शब्दांत, $\Zminus$ मध्ये प्रतिस्थापन जोडले की $\ZFminus$ मिळते.
\end{defn}
\noindent
यातील नाव \emph{झर्मेलो--फ्रेंकेल} संच उपपत्ती दर्शवते (त्यातून \emph{वगळलेली} गोष्ट आपण पुढे पाहू). प्रतिस्थापनाची मांडणी 1922 मध्ये फ्रेंकेल यांनी केली असे मानले जाते; म्हणून त्यांच्या नावाला येथे स्थान आहे. मात्र त्याचे पहिले काटेकोर स्वरूप स्कोलेम (1922) यांनी दिले.







\section{क्रमप्रकार म्हणून क्रमसूचक संख्या}\label{sth:ordinals:ordtype:sec}

आता प्रतिस्थापन उपलब्ध आहे आणि आपण $\ZFminus$ मध्ये काम करत आहोत. त्यामुळे ज्या निष्कर्षाकडे आपण वाटचाल केली, तो अखेर सिद्ध करता येईल:

\begin{thm}\label{sth:ordinals:ordtype:thmOrdinalRepresentation}
प्रत्येक सुक्रमण एकमेव क्रमसूचक संख्येशी समरूपी असते.
\end{thm}

\begin{proof}
$\tuple{A, <}$ हे सुक्रमण मानू. \ref{sth:ordinals:basic:ordisoidentity} वरून ते जास्तीत जास्त एका क्रमसूचक संख्येशी समरूपी असू शकते. म्हणून विसंगतीसाठी $\tuple{A, <}$ हे \emph{कोणत्याही} क्रमसूचक संख्येशी समरूपी नाही असे मानू. आधी $\tuple{A, <}$ ला ``शक्य तितके लहान'' करू. तपशील असा: $\tuple{A_a,
<_a}$ हा एखादा यथार्थ आरंभीचा खंड कोणत्याही क्रमसूचक संख्येशी समरूपी नसेल, तर हा गुणधर्म असणारा $a \in A$ मधील लघुतम सदस्य निवडू; मग $B = A_a$ आणि $\mathord{\lessdot} =
\mathord{<_a}$ ठेवू. तसे नसल्यास $B = A$ आणि $\mathord{\lessdot} =
\mathord{<}$ ठेवू.

या निवडीमुळे $B$ चा प्रत्येक यथार्थ आरंभीचा खंड कोणत्या तरी क्रमसूचक संख्येशी समरूपी असतो; वर पाहिल्याप्रमाणे ती एकमेव असते. म्हणून प्रतिस्थापनावरून पुढील संच अस्तित्वात असतो आणि तो फलन असतो:
\[
	f = \Setabs{\tuple{\beta, b}}{b \in B\text{ आणि }\ordeq{\beta}{\tuple{B_b, \lessdot_b}}}
\]
विसंगतीचा पुरावा पूर्ण करण्यासाठी कोणत्या तरी क्रमसूचक संख्या $\alpha$ साठी $f$ हे $\alpha \to B$ समरूपण आहे, हे दाखवू.

$\ran{f}
= B$ हे उघड आहे. \ref{sth:ordinals:iso:lemordsegments} वरून $f$ क्रम जपते; म्हणजे $\gamma \in \beta$ तेव्हा आणि केवळ तेव्हाच $f(\gamma) \lessdot f(\beta)$. $\dom{f}$ ही क्रमसूचक संख्या आहे हे दाखवण्यासाठी \ref{sth:ordinals:basic:corordtransitiveord} वरून $\dom{f}$ संक्रमक आहे हे दाखवणे पुरेसे आहे. म्हणून $\beta \in \dom{f}$ निश्चित करू; म्हणजे एखाद्या $b$ साठी $\ordeq{\beta}{\tuple{B_b, \lessdot_b}}$. $\gamma \in
\beta$ असेल, तर \ref{sth:ordinals:iso:wellordinitialsegment} वरून $\gamma \in \dom{f}$. म्हणून $\beta \subseteq
\dom{f}$.
\end{proof}

या निष्कर्षामुळे \ref{sth:ordinals:vn:sec} पासून अपेक्षित असलेली पुढील व्याख्या आता देता येते:

\begin{defn}
$\tuple{A, <}$ हे सुक्रमण असेल, तर त्याचा क्रमप्रकार $\ordtype{A, <}$ ही $\ordeq{\tuple{A, < }}{\alpha}$ पूर्ण करणारी एकमेव क्रमसूचक संख्या $\alpha$ असते.
\end{defn}

या व्याख्येवरून पुढील दोन सुंदर तत्त्वेही मिळतात:

\begin{cor}\label{sth:ordinals:ordtype:ordtypesworklikeyouwant}
$\tuple{A, <}$ आणि $\tuple{B, \lessdot}$ ही सुक्रमणे असतील, तर:
\begin{align*}
	\ordtype{A, <} = \ordtype{B, \lessdot}&\text{ तेव्हा आणि केवळ तेव्हाच }\ordeq{\tuple{A, <}}{\tuple{B, \lessdot}}\\
	\ordtype{A, <} \in \ordtype{B, \lessdot}&\text{ तेव्हा आणि केवळ तेव्हाच }\ordeq{\tuple{A, <}}{\tuple{B_b, \lessdot_b}}\text{ अशा काही }b \in B
\end{align*}
\end{cor}

\begin{proof}
पहिली समता \ref{sth:ordinals:basic:ordisoidentity} वरून मिळते. दुसरा दावा सिद्ध करण्यासाठी $\ordtype{A, <} = \alpha$ आणि $\ordtype{B,
\lessdot} = \beta$ मानू; तसेच $f \colon \beta \to \tuple {B, \lessdot}$ हे समरूपण घेऊ. मग:
\begin{align*}
	\alpha \in \beta&\text{ तेव्हा आणि केवळ तेव्हाच }\funrestrictionto{f}{\alpha} \colon \alpha \to B_{f(\alpha)}\text{ हे समरूपण आहे}\\
	&\text{ तेव्हा आणि केवळ तेव्हाच }\ordeq{\tuple{A, <}}{\tuple{B_{f(\alpha)}, \lessdot_{f(\alpha)}}}\\
	&\text{ तेव्हा आणि केवळ तेव्हाच }\ordeq{\tuple{A, <}}{\tuple{B_b, \lessdot_b}}\text{ अशा काही $b \in B$ साठी}
\end{align*}
हे \ref{sth:ordinals:iso:ordisounique}, \ref{sth:ordinals:iso:wellordinitialsegment} आणि \ref{sth:ordinals:basic:ordissetofsmallerord} वरून मिळते.
\end{proof}







\section{उत्तरवर्ती आणि सीमा क्रमसूचक संख्या}\label{sth:ordinals:opps:sec}

पुढील काही प्रकरणांत आपण क्रमसूचक संख्यांचा भरपूर वापर करू. म्हणून आधी काही सोप्या संकल्पना मांडू.

\begin{defn}
कोणत्याही क्रमसूचक संख्या $\alpha$ ची \emph{उत्तरवर्ती} संख्या $\ordsucc{\alpha}
=\alpha \cup \{\alpha\}$ असते. एखाद्या क्रमसूचक संख्या~$\beta$ साठी $\ordsucc{\beta} = \alpha$ असेल, तर $\alpha$ ही \emph{उत्तरवर्ती} क्रमसूचक संख्या आहे असे म्हणू. $\alpha$ ही \emph{सीमा} क्रमसूचक संख्या असते तेव्हा आणि केवळ तेव्हाच $\alpha$ रिकामीही नसते आणि उत्तरवर्ती क्रमसूचक संख्याही नसते.
\end{defn}
\noindent
पुढील निष्कर्षावरून \emph{उत्तरवर्ती} ही योग्य संकल्पना आहे हे दिसते:

\begin{prop}
कोणत्याही क्रमसूचक संख्या $\alpha$ साठी:
\begin{enumerate}
	\item $\alpha \in \ordsucc{\alpha}$;
	\item $\ordsucc{\alpha}$ ही क्रमसूचक संख्या आहे;
	\item $\alpha \in \beta \in
	\ordsucc{\alpha}$ अशी कोणतीही क्रमसूचक संख्या $\beta$ नाही.
	\end{enumerate}
\end{prop}

\begin{proof}
$\alpha \in \alpha \cup \{\alpha\} = \ordsucc{\alpha}$ हे स्पष्ट आहे. तसेच $\ordsucc{\alpha}$ हा क्रमसूचक संख्यांचा संक्रमक संच असल्याने \ref{sth:ordinals:basic:corordtransitiveord} वरून क्रमसूचक संख्या आहे. $\alpha \in \beta \in \ordsucc{\alpha}$ अशक्य आहे: तसे असल्यास $\beta \in \alpha$ किंवा $\beta = \alpha$ असेल, आणि त्यावरून \ref{sth:ordinals:basic:ordordered} शी विसंगती येईल.
\end{proof}
\noindent
यामुळे अनंतपार विगमनाने पुराव्याचा आणखी एक प्रकारही मिळतो:

\begin{thm}[सुलभ अनंतपार विगमन]\label{sth:ordinals:opps:simpletransrecursion}
$\phi(x)$ हे पुढील अटी पूर्ण करणारे सूत्र असू द्या:
\begin{enumerate}
	\item $\phi(\emptyset)$; आणि
	\item कोणत्याही क्रमसूचक संख्या $\alpha$ साठी, $\phi(\alpha)$ असेल, तर $\phi(\ordsucc{\alpha})$; आणि
	\item $\alpha$ ही सीमा क्रमसूचक संख्या असून $(\forall \beta \in
	\alpha)\phi(\beta)$ असेल, तर $\phi(\alpha)$.
\end{enumerate}
मग $\forall \alpha \phi(\alpha)$.
\end{thm}

\begin{proof}
व्यत्यस्त रूप सिद्ध करू. एखाद्या क्रमसूचक संख्येसाठी $\lnot\phi$ असेल असे मानू; अशी लघुतम क्रमसूचक संख्या $\gamma$ घेऊ. मग $\gamma = \emptyset$ असेल; किंवा $\phi(\alpha)$ पूर्ण करणाऱ्या एखाद्या $\alpha$ साठी $\gamma = \ordsucc{\alpha}$ असेल; किंवा $\gamma$ ही सीमा क्रमसूचक संख्या असून $(\forall
\beta \in \gamma)\phi(\beta)$ असेल.
\end{proof}
\noindent
शेवटी, पुढे उपयोगी पडणारा आणखी एक संकेत पाहू:

\begin{defn}\label{sth:ordinals:opps:defsupstrict}
$X$ हा क्रमसूचक संख्यांचा संच असेल, तर $\supstrict(X) = \bigcup_{\alpha
\in X} \ordsucc{\alpha}$.
\end{defn}
\noindent
येथे ``lsub'' म्हणजे ``लघुतम काटेकोर ऊर्ध्वबंध''.\footnote{काही पुस्तकांत यासाठी ``$\text{sup}(X)$'' लिहितात. पण इतर पुस्तकांत ``$\text{sup}(X)$'' हा संकेत लघुतम \emph{अकाटेकोर} ऊर्ध्वबंधासाठी, म्हणजे केवळ $\bigcup X$ साठी, वापरतात. $X$ मध्ये महत्तम सदस्य $\alpha$ असेल, तर या दोन संकल्पना वेगळ्या ठरतात: लघुतम \emph{काटेकोर} ऊर्ध्वबंध $\ordsucc{\alpha}$ असतो, तर लघुतम \emph{अकाटेकोर} ऊर्ध्वबंध फक्त $\alpha$ असतो.} पुढील निष्कर्ष हे स्पष्ट करतो:

\begin{prop}
$X$ हा क्रमसूचक संख्यांचा संच असेल, तर $\supstrict(X)$ ही $X$ मधील प्रत्येक क्रमसूचक संख्येपेक्षा मोठी असलेली लघुतम क्रमसूचक संख्या आहे.
\end{prop}

\begin{proof}
$Y = \Setabs{\ordsucc{\alpha}}{\alpha \in X}$ ठेवू; मग $\supstrict(X) = \bigcup Y$. क्रमसूचक संख्या संक्रमक असतात आणि क्रमसूचक संख्येचा प्रत्येक सदस्य क्रमसूचक संख्या असतो. म्हणून $\supstrict(X)$ हा क्रमसूचक संख्यांचा संक्रमक संच आहे आणि \ref{sth:ordinals:basic:corordtransitiveord} वरून क्रमसूचक संख्या आहे.

$\alpha \in X$ असेल, तर $\ordsucc{\alpha} \in Y$. म्हणून $\ordsucc{\alpha}
\subseteq \bigcup Y = \supstrict(X)$ आणि त्यामुळे $\alpha \in
\supstrict(X)$. म्हणजे $\supstrict(X)$ ही $X$ मधील प्रत्येक क्रमसूचक संख्येपेक्षा काटेकोरपणे मोठी आहे.

उलट, $\alpha \in \supstrict(X)$ असेल, तर एखाद्या $\beta \in X$ साठी $\alpha \in
\ordsucc{\beta} \in Y$; म्हणून $\alpha \leq
\beta \in X$. त्यामुळे $\supstrict(X)$ हा $X$ चा \emph{लघुतम} काटेकोर ऊर्ध्वबंध आहे.
\end{proof}



\chapter{टप्पे आणि प्रतांक}\label{sth:spine:chap}





\section{टप्प्यांची $V_\alpha$ या रूपात व्याख्या}\label{sth:spine:valpha:sec}

\ref{sth:ordinals:chap} मध्ये आपण सुक्रमणे आणि (फोन न्यूमन यांच्या) क्रमसूचक संख्या परिभाषित केल्या. या प्रकरणात त्यांचा वापर करून संचांच्या पदानुक्रमाचे \emph{स्वतःचे} वर्णन करू. यासाठी \ref{sth:ordinals:opps:sec} मधील उत्तरवर्ती आणि सीमा क्रमसूचक संख्यांच्या संकल्पना आठवू. पुढील व्याख्येत त्या वापरल्या आहेत:

\begin{defn}\label{sth:spine:valpha:defValphas}
	\begin{align*}
	V_\emptyset &\defis \emptyset\\
	V_{\ordsucc{\alpha}} &\defis \Pow{V_\alpha} & &
	\text{प्रत्येक क्रमसूचक संख्या }\alpha\\
	V_{\alpha} &\defis \bigcup_{\gamma < \alpha} V_\gamma & &
	\text{जेव्हा }\alpha\text{ ही सीमा क्रमसूचक संख्या असते}
\end{align*}
\end{defn}
\noindent
ही क्रमसूचक संख्यांवरील \emph{अनंतपार पुनरावर्तनाने} केलेली व्याख्या ठरेल. तिची तुलना नैसर्गिक संख्यांवरील फलनांच्या पुनरावर्ती व्याख्यांशी करावी.\footnote{\ref{sfr:infinite:dedekind:sec} मधील बेरीज, गुणाकार आणि घातांकनाच्या व्याख्या पाहा.} नैसर्गिक संख्यांच्या बाबतीत आधारप्रकरण आणि उत्तरवर्ती प्रकरणे परिभाषित करतो; पण क्रमसूचक संख्यांच्या बाबतीत \emph{सीमा} प्रकरणांतील वर्तनही सांगावे लागते.

$V_\alpha$ ची ही व्याख्या एक महत्त्वाचा टप्पा ठरेल. संचांची रचना \emph{टप्प्याटप्प्याने} होते, या कल्पनेने आपण संचांच्या पदानुक्रमाचे अनौपचारिक समर्थन केले. $V_\alpha$ हे प्रत्यक्षात तेच टप्पे आहेत. पण येथे महत्त्वाची बाब म्हणजे टप्प्यांचे हे \emph{आंतरिक} वर्णन आहे. उपपत्तीचे एखादे संभाव्य \emph{प्रतिमान} सुचवण्याऐवजी आपण टप्पे आपल्या संच उपपत्तीच्या \emph{आतच} परिभाषित करू.







\section{अनंतपार पुनरावर्तनाची प्रमेये}\label{sth:spine:recursion:sec}

मात्र आधी \ref{sth:spine:valpha:defValphas} ही यशस्वी व्याख्या आहे, हे निश्चित करावे लागेल. अधिक सर्वसाधारणपणे, अनंतपार पुनरावर्तनाने व्याख्या करण्याचा कोणताही प्रयत्न यशस्वी होतो, हे सिद्ध करणे गरजेचे आहे. हाच या विभागाचा उद्देश आहे.

\emph{सूचना: येथील विषय अवघड आहे}. तरी मुख्य मुद्दा सोपा आहे: अनंतपार विगमन आणि प्रतिस्थापन मिळून अनंतपार पुनरावर्तनाच्या (अनेक रूपांच्या) वैधतेची हमी देतात.\footnote{आठवण: स्पष्टपणे अन्यथा सांगितले नसेल, तर सर्व सूत्रे आणि पदे यांत प्राचले असू शकतात.}

\begin{defn}
	$\tau(x)$ हे पद, $f$ हे फलन आणि $\alpha$ ही क्रमसूचक संख्या मानू. $\dom{f} = \alpha$ आणि $(\forall \beta \in \alpha)f(\beta) = \tau(\funrestrictionto{f}{\beta})$ या दोन्ही अटी पूर्ण होत असतील तेव्हा आणि केवळ तेव्हाच $f$ हे $\tau$ साठी \emph{$\alpha$-निकटन} आहे असे म्हणू.
\end{defn}

\begin{lem}[परिबद्ध पुनरावर्तन]\label{sth:spine:recursion:transrecursionfun}
कोणत्याही पद $\tau(x)$ आणि कोणत्याही क्रमसूचक संख्या $\alpha$ साठी $\tau$ चे एकमेव $\alpha$-निकटन असते.
\end{lem}
\begin{proof}
कोणत्याही $\gamma \leq \alpha$ साठी एकमेव $\gamma$-निकटन असते, हे दाखवू.

आधी एकमेवता सिद्ध करू. अनुक्रमे $g$ आणि $h$ ही $\gamma$- आणि $\delta$-निकटने असू द्या. त्यांच्या आदानांवरील अनंतपार विगमनाने $\beta \in \dom{g} \cap \dom{h} =
\gamma \cap \delta = \min(\gamma, \delta)$ असेल, तर $g(\beta) = h(\beta)$ मिळते. म्हणून निकटने अस्तित्वात असतील, तर ती एकमेव असतात आणि सर्व सामायिक मूल्यांवर जुळतात.

अस्तित्व सिद्ध करण्यासाठी आता $\delta \leq \alpha$ या क्रमसूचक संख्यांवर सुलभ अनंतपार विगमन (\ref{sth:ordinals:opps:simpletransrecursion}) वापरू.

रिकामे फलन स्पष्टपणे $\emptyset$-निकटन असते.

$g$ हे $\gamma$-निकटन असेल, तर $g \cup
\{\tuple{\gamma, \tau(g)}\}$ हे $\ordsucc{\gamma}$-निकटन असते.

$\gamma$ ही सीमा क्रमसूचक संख्या असेल आणि प्रत्येक $\delta < \gamma$ साठी $g_\delta$ हे $\delta$-निकटन असेल, तर $g = \bigcup_{\delta \in \gamma} g_\delta$ ठेवू. विविध $g_\delta$ सर्व मूल्यांवर जुळत असल्याने हे फलन आहे. तसेच $\delta \in \gamma$ असेल, तर $g(\delta) = g_{\ordsucc{\delta}}(\delta) =
\tau(\funrestrictionto{g_{\ordsucc{\delta}}}{\delta}) =
\tau(\funrestrictionto{g}{\delta})$.

याने अनंतपार विगमनाचा पुरावा पूर्ण होतो.
\end{proof}

फलनाऐवजी \emph{पद} परिभाषित करण्याची मुभा घेतली, तर मागील निष्कर्षातील $\alpha$ ही मर्यादा काढून टाकता येते. पुढील निष्कर्षाचे विधान आणि पुरावा यांत $\sigma$ हे पद असताना $\funrestrictionto{\sigma}{\alpha} = \Setabs{\tuple{\beta,
\sigma(\beta)}}{\beta \in \alpha}$ असे मानू.

\begin{thm}[सामान्य पुनरावर्तन]\label{sth:spine:recursion:transrecursionschema}
कोणत्याही पद $\tau(x)$ साठी $\sigma(x)$ हे पद स्पष्टपणे परिभाषित करता येते, असे की कोणत्याही क्रमसूचक संख्या~$\alpha$ साठी $\sigma(\alpha) = \tau(\funrestrictionto{\sigma}{\alpha})$.
\end{thm}

\begin{proof}
प्रत्येक $\alpha$ साठी \ref{sth:spine:recursion:transrecursionfun} वरून $\tau$ चे एकमेव $\alpha$-निकटन $f_\alpha$ असते. $\sigma(\alpha)$ ची व्याख्या $f_{\ordsucc{\alpha}}(\alpha)$ अशी करू. आता:
	\begin{align*}
		\sigma(\alpha) &=
		f_{\ordsucc{\alpha}}(\alpha) \\&=
		\tau(\funrestrictionto{f_{\ordsucc{\alpha}}}{\alpha}) \\&=
		\tau(\Setabs{\tuple{\beta, f_{\ordsucc{\alpha}}(\beta)}}{\beta \in \alpha}) \\&=
		\tau(\Setabs{\tuple{\beta, f_{\ordsucc{\beta}}(\beta)}}{\beta \in \alpha})\\&=
		\tau(\funrestrictionto{\sigma}{\alpha})
	\end{align*}
येथे \ref{sth:spine:recursion:transrecursionfun} प्रमाणे सर्व $\beta < \alpha$ साठी $f_{\ordsucc{\beta}}(\beta) = f_{\ordsucc{\alpha}}(\beta)$, हे वापरले आहे.
\end{proof}
\noindent
\ref{sth:spine:recursion:transrecursionschema} हा एक \emph{रूपबंध} आहे, हे लक्षात घ्या. $\sigma$ ने फलन, म्हणजे विशिष्ट प्रकारचा \emph{संच}, परिभाषित व्हावा अशी अपेक्षा करता येत नाही. तसे असेल, तर $\dom{\sigma}$ हा सर्व क्रमसूचक संख्यांचा संच ठरेल; हे बुराली-फोर्ती विरोधापत्तीशी (\ref{sth:ordinals:basic:buraliforti}) विसंगत आहे.

तरी \ref{sth:spine:recursion:transrecursionschema} आपली $V_\alpha$ ची व्याख्या योग्य ठरवतो, हे अजून दाखवायचे आहे. हे लगेच स्पष्ट होणार नाही; पण अनंतपार पुनरावर्तनाचे शेवटचे, सोपे रूप पाहिल्यावर ते उघड होईल.

\begin{thm}[सुलभ पुनरावर्तन]\label{sth:spine:recursion:simplerecursionschema}
कोणतीही पदे $\tau(x)$ आणि $\theta(x)$ तसेच कोणताही संच $A$ दिल्यास, पुढील अटी पूर्ण करणारे पद $\sigma(x)$ स्पष्टपणे परिभाषित करता येते:
\begin{align*}
	\sigma(\emptyset) &= A\\
	\sigma(\ordsucc{\alpha}) &= \tau(\sigma(\alpha)) &&
		\text{प्रत्येक क्रमसूचक संख्या }\alpha\\
	\sigma(\alpha) &= \theta(\ran{\funrestrictionto{\sigma}{\alpha}})&&
	\text{जेव्हा }\alpha\text{ ही सीमा क्रमसूचक संख्या असते}
\end{align*}
\end{thm}

\begin{proof}
आधी $\xi(x)$ हे पद पुढीलप्रमाणे परिभाषित करू:


\[
	\xi(x) =
	\begin{cases}
			A &\text{जर $x$ क्रमसूचक प्रांताचे फलन नसेल, किंवा}\\
			&\hspace{1em}\text{त्याचा प्रांत रिकामा असेल; अन्यथा:}\\
			\tau(x(\alpha)) & \text{जर $\dom{x} = \ordsucc{\alpha}$}\\
			\theta(\ran{x}) & \text{जर $\dom{x}$ सीमा क्रमसूचक संख्या असेल}
	\end{cases}
\]
\ref{sth:spine:recursion:transrecursionschema} वरून प्रत्येक क्रमसूचक संख्या $\alpha$ साठी $\sigma(\alpha) = \xi(\funrestrictionto{\sigma}{\alpha})$ असे पद $\sigma(x)$ मिळते. शिवाय $\funrestrictionto{\sigma}{\alpha}$ हे $\alpha$ प्रांताचे फलन आहे. सुलभ अनंतपार विगमनाने (\ref{sth:ordinals:opps:simpletransrecursion}) $\sigma$ मध्ये अपेक्षित गुणधर्म आहेत हे दाखवू.

प्रथम, $\sigma(\emptyset) = \xi(\emptyset) = A$.

पुढे, $\sigma(\ordsucc{\alpha}) = \xi(\funrestrictionto{\sigma}{\ordsucc{\alpha}}) = \tau(\funrestrictionto{\sigma}{\ordsucc{\alpha}}(\alpha)) = \tau(\sigma(\alpha))$.

शेवटी, $\alpha$ ही सीमा क्रमसूचक संख्या असल्यास $\sigma(\alpha) = \xi(\funrestrictionto{\sigma}{\alpha}) = \theta(\ran{\funrestrictionto{\sigma}{\alpha}})$.
\end{proof}
\noindent
आता \ref{sth:spine:valpha:defValphas} योग्य ठरवण्यासाठी $A
= \emptyset$, $\tau(x) = \Pow{x}$ आणि $\theta(x) = \bigcup x$ एवढेच घ्यावे लागते. अखेर यामुळे $V_\alpha$ ची व्याख्या सिद्ध होते!{}







\section{टप्प्यांचे मूलभूत गुणधर्म}\label{sth:spine:Valphabasic:sec}

$V_\alpha$ च्या व्याख्येचे पायाभूत महत्त्व स्पष्ट करण्यासाठी त्यांच्याविषयी काही मूलभूत निष्कर्ष मांडू. आधी एक व्याख्या:\footnote{``उपसंच-समावेशक'' यासाठी प्रस्थापित संज्ञा नाही; मूळ इंग्रजीत Button (2021) यांनी वापरलेले ``potent'' हे नाव आहे.}
\begin{defn}
	$\forall x((\exists y \in A)x \subseteq y \lif x \in A)$ असेल तेव्हा आणि केवळ तेव्हाच संच $A$ \emph{उपसंच-समावेशक} आहे असे म्हणू.
\end{defn}
\begin{lem}\label{sth:spine:Valphabasic:Valphabasicprops}
प्रत्येक क्रमसूचक संख्या $\alpha$ साठी:
\begin{enumerate}
	\item\label{sth:spine:Valphabasic:Valphatrans} प्रत्येक $V_\alpha$ संक्रमक असतो.
	\item\label{sth:spine:Valphabasic:Valphapotent} प्रत्येक $V_\alpha$ उपसंच-समावेशक असतो.
	\item\label{sth:spine:Valphabasic:Valphacum} $\gamma \in \alpha$ असेल, तर $V_\gamma
	\in V_\alpha$ (आणि म्हणून \ref{sth:spine:Valphabasic:Valphatrans} वरून $V_\gamma \subseteq V_\alpha$ सुद्धा).
\end{enumerate}
\end{lem}

\begin{proof}
हे (एकसामयिक) अनंतपार विगमनाने सिद्ध करू. विगमनासाठी प्रत्येक क्रमसूचक संख्या $\beta < \alpha$ साठी \ref{sth:spine:Valphabasic:Valphatrans}--\ref{sth:spine:Valphabasic:Valphacum} लागू आहेत असे मानू.

$\alpha = \emptyset$ हे प्रकरण उघड आहे.

$\alpha = \ordsucc{\beta}$ मानू. \ref{sth:spine:Valphabasic:Valphacum} साठी, $\gamma \in \alpha$ असेल, तर गृहीतकावरून $V_\gamma \subseteq V_\beta$; म्हणून $V_\gamma \in \Pow{V_\beta} = V_\alpha$. \ref{sth:spine:Valphabasic:Valphapotent} साठी $A \subseteq B \in V_\alpha$, म्हणजे $A
\subseteq B \subseteq V_\beta$, असे मानू. मग $A \subseteq V_\beta$ आणि म्हणून $A \in V_\alpha$. \ref{sth:spine:Valphabasic:Valphatrans} साठी, $x \in A \in
V_\alpha$ असेल, तर $A \subseteq V_\beta$ आणि म्हणून $x \in V_\beta$. गृहीतकानुसार $V_\beta$ संक्रमक असल्याने $x \subseteq V_\beta$; म्हणून $x \in V_\alpha$.

$\alpha$ ही सीमा क्रमसूचक संख्या मानू. \ref{sth:spine:Valphabasic:Valphacum} साठी, $\gamma \in \alpha$ असेल, तर $\gamma \in \ordsucc{\gamma} \in \alpha$; म्हणून गृहीतकावरून $V_\gamma \in V_{\ordsucc{\gamma}}$ आणि त्यामुळे $V_\gamma \in \bigcup_{\beta \in \alpha} V_\beta = V_\alpha$. \ref{sth:spine:Valphabasic:Valphatrans} आणि \ref{sth:spine:Valphabasic:Valphapotent} साठी एवढेच पाहणे पुरेसे आहे की संक्रमक (अनुक्रमे उपसंच-समावेशक) संचांचा संयोग संक्रमक (अनुक्रमे उपसंच-समावेशक) असतो.
\end{proof}

\begin{lem}\label{sth:spine:Valphabasic:Valphanotref}
प्रत्येक क्रमसूचक संख्या $\alpha$ साठी $V_\alpha \notin V_\alpha$.
\end{lem}

\begin{proof}
अनंतपार विगमन वापरू. $V_\emptyset \notin V_\emptyset$ हे उघड आहे.

$V_{\ordsucc{\alpha}} \in V_{\ordsucc{\alpha}} = \Pow{V_\alpha}$ असेल, तर $V_{\ordsucc{\alpha}} \subseteq V_\alpha$. तसेच \ref{sth:spine:Valphabasic:Valphabasicprops} वरून $V_\alpha
\in V_{\ordsucc{\alpha}}$; म्हणून $V_\alpha \in V_\alpha$. व्यत्यस्त रूपाने, $V_\alpha \notin V_\alpha$ असेल, तर $V_{\ordsucc{\alpha}} \notin V_{\ordsucc{\alpha}}$.

$\alpha$ ही सीमा क्रमसूचक संख्या असून $V_\alpha \in V_\alpha = \bigcup_{\beta \in
\alpha}V_\beta$ असेल, तर एखाद्या $\beta \in \alpha$ साठी $V_\alpha \in V_\beta$. पण मग $V_\beta \in V_\alpha$ सुद्धा; म्हणून \ref{sth:spine:Valphabasic:Valphabasicprops} दोनदा वापरून $V_\beta \in V_\beta$ मिळते. व्यत्यस्त रूपाने, सर्व $\beta \in \alpha$ साठी $V_\beta \notin V_\beta$ असेल, तर $V_\alpha \notin
V_\alpha$.
\end{proof}

\begin{cor}
कोणत्याही क्रमसूचक संख्या $\alpha, \beta$ साठी: $\alpha \in \beta$ तेव्हा आणि केवळ तेव्हाच $V_\alpha \in V_\beta$.
\end{cor}

\begin{proof}
\ref{sth:spine:Valphabasic:Valphabasicprops} वरून एक दिशा मिळते. उलट, $V_\alpha \in V_\beta$ मानू. \ref{sth:spine:Valphabasic:Valphanotref} वरून $\alpha \neq \beta$. तसेच $\beta \notin \alpha$; कारण अन्यथा $V_\beta \in V_\alpha$ मिळेल आणि मग \ref{sth:spine:Valphabasic:Valphabasicprops} दोनदा वापरून $V_\beta \in V_\beta$ मिळेल, हे \ref{sth:spine:Valphabasic:Valphanotref} शी विसंगत आहे. म्हणून त्रिपर्यायावरून $\alpha \in \beta$.
\end{proof}
\noindent
या सर्वांवरून प्रत्येक $V_\alpha$ हा पदानुक्रमातील $\alpha$ वा टप्पा आहे असे मानता येते. का ते पाहू.

आपले $V_\alpha$ \emph{क्रमिक} प्रक्रियेत तयार होतात असे नक्कीच मानता येते; कारण क्रमसूचक संख्यांचा वापर पुनरावृत्तीच्या संकल्पनेशी जुळतो. शिवाय, एक टप्पा दुसऱ्याआधी तयार होतो, म्हणजे $V_\beta \in V_\alpha$, म्हणजे $\beta \in \alpha$, असेल, तर $V_\beta \subseteq V_\alpha$ असल्यामुळे रचना \emph{संचयी} असते. शेवटी, आधीच्या कोणत्याही टप्प्यावर उपलब्ध असलेल्या संचांचे \emph{सर्व} संभाव्य संग्रह आपण प्रत्यक्ष तयार करतो; कारण प्रत्येक उत्तरवर्ती टप्पा $V_{\ordsucc{\alpha}}$ हा त्याच्या पूर्ववर्ती $V_\alpha$ चा घातसंच असतो.

म्हणजे $\ZFminus$ मिळाल्यावर संचांच्या संचयी-क्रमिक पदानुक्रमाची आपली संकल्पना मांडण्याचे काम \emph{जवळजवळ} पूर्ण झाले आहे. (अर्थात प्रतिस्थापनाचे समर्थन अजून करायचे आहे.)








\section{पायाभूतता}\label{sth:spine:foundation:sec}

आपले काम फक्त \emph{जवळजवळ} पूर्ण झाले आहे, \emph{पूर्णपणे} नाही. कारण $\ZFminus$ मधील कोणत्याही तत्त्वावरून \emph{प्रत्येक} संच कोणत्या तरी $V_\alpha$ मध्ये आहे, म्हणजे प्रत्येक संच कोणत्या तरी टप्प्यावर तयार होतो, याची खात्री मिळत नाही.

पदानुक्रमाबाहेर काही संच आहेत की नाहीत, याची आपल्याला गणिती दृष्टीने \emph{पर्वा नाही} असा एक सरळ अर्थ आहे: तसे संच असतील, तर त्यांच्याकडे दुर्लक्ष करता येईल. पण \emph{संचाची संकल्पना} मांडताना प्रत्येक संच कोणत्या तरी टप्प्यावर तयार होतो, हीच कल्पना आपण वापरली आहे (\ref{sth:z:story:sec} मधील \stageshier{} पाहा). म्हणून पदानुक्रमाबाहेर संच असण्याची शक्यता आपल्याला वगळायची आहे. त्यासाठी प्रत्येक संच पदानुक्रमात कुठे तरी येतो याची हमी देणारे नवे स्वयंसिद्धक घालावे लागेल.

$V_\alpha$ हेच आपले टप्पे असल्यामुळे पुढील विधान थेट स्वयंसिद्धक म्हणून घेता येईल:

\begin{defish}
\emph{नियमितता.} $\forall A \exists \alpha\, A \subseteq V_\alpha$
\end{defish}

ही पद्धत पूर्णपणे योग्य ठरेल. मात्र पुढील विभागात स्पष्ट होणाऱ्या कारणांमुळे आपण याऐवजी दुसरे स्वयंसिद्धक स्वीकारू:

\begin{axiom}[पायाभूतता]
$(\forall A \neq \emptyset)(\exists B \in A)A \cap B = \emptyset$.
\end{axiom}

काही प्रयत्नांनंतर $\ZFminus$ मध्ये पायाभूततेवरून नियमितता मिळते, हे दाखवता येते:
\begin{defn}
प्रत्येक संच $A$ साठी पुढीलप्रमाणे ठेवू:
\begin{align*}
	\text{cl}_0(A) &= A,\\
	\text{cl}_{n+1}(A) &= \bigcup \text{cl}_n(A),\\
	\text{trcl}(A) &= \bigcup_{n < \omega} \text{cl}_{n}(A).
\end{align*}
$\text{trcl}(A)$ ला $A$ चे \emph{संक्रमक आवरण} म्हणू.
\end{defn}
\noindent
``संक्रमक आवरण'' हे नाव योग्य आहे:
\begin{prop}\label{sth:spine:foundation:subsetoftrcl}
$A \subseteq \trcl{A}$ आणि $\trcl{A}$ हा संक्रमक संच आहे.
\end{prop}

\begin{proof}
$A = \text{cl}_0(A) \subseteq \text{trcl}(A)$ हे उघड आहे. आणि $x
\in b \in \trcl{A}$ असेल, तर एखाद्या $n$ साठी $b \in \text{cl}_n(A)$; म्हणून $x
\in \text{cl}_{n+1}(A) \subseteq \text{trcl}(A)$.
\end{proof}

\begin{lem}\label{sth:spine:foundation:lem:TransitiveWellFounded}
$A$ हा संक्रमक संच असेल, तर $A \subseteq V_\alpha$ असेल अशी काही $\alpha$ अस्तित्वात असते.
\end{lem}

\begin{proof}
\ref{sth:ordinals:opps:defsupstrict} मधील ``$\supstrict(X)$'' ची व्याख्या आठवून पुढील दोन संच परिभाषित करू:
\begin{align*}
	D  &= \Setabs{x \in A}{\forall \delta\ x \nsubseteq V_\delta}\\
	\alpha &= \supstrict\Setabs{\delta}{(\exists x \in A)
	(x \subseteq V_\delta \land (\forall \gamma \in \delta)x \nsubseteq V_\gamma)}
\end{align*}
$D = \emptyset$ मानू. मग $x \in A$ असेल, तर $x \subseteq V_\delta$ असेल अशी काही $\delta$ असते. क्रमसूचक संख्या सुक्रमित असल्याने $(\forall \gamma \in \delta)x \nsubseteq V_\gamma$ ही अटही पूर्ण करणारी सर्वांत लहान अशी संख्या निवडता येते. म्हणून $\delta \in \alpha$; त्यामुळे \ref{sth:spine:Valphabasic:Valphabasicprops} वरून $x \in V_\alpha$. यावरून अपेक्षेप्रमाणे $A \subseteq
V_{\alpha}$.

म्हणून $D = \emptyset$ हे दाखवणे पुरेसे आहे. विसंगतीसाठी तसे नाही असे मानू. पायाभूततेवरून $B \in D \subseteq A$ असा सदस्य मिळतो की $D \cap B
= \emptyset$. $x \in B$ असेल, तर $A$ संक्रमक असल्याने $x \in A$; आणि $x \notin D$ असल्याने $\exists \delta\ x \subseteq
V_\delta$. आता
\[
	\beta = \supstrict\Setabs{\delta}{(\exists x \in B)(x \subseteq V_\delta \land (\forall \gamma < \delta)x \nsubseteq V_\gamma)}.
\]
असे ठेवू. पूर्वीप्रमाणे $B \subseteq V_\beta$ मिळते; पण ते $B \in
D$ या विधानाशी विसंगत आहे.
\end{proof}

\begin{thm}\label{sth:spine:foundation:zfentailsregularity}
नियमितता लागू आहे.
\end{thm}

\begin{proof}
$A$ निश्चित करू. \ref{sth:spine:foundation:subsetoftrcl} वरून $A \subseteq \trcl{A}$ आणि हा आवरण-संच संक्रमक आहे. म्हणून \ref{sth:spine:foundation:lem:TransitiveWellFounded} वरून $A \subseteq \trcl{A} \subseteq V_\alpha$ असेल अशी काही $\alpha$ असते.
\end{proof}

या निष्कर्षांवरून $\ZFminus$ मध्ये $\emph{पायाभूतता}\Rightarrow\emph{नियमितता}$ हे सशर्त विधान सिद्ध होते. \ref{sth:spine:rank:zfminusregularityfoundation} मध्ये $\ZFminus$ मध्ये $\emph{नियमितता}\Rightarrow\emph{पायाभूतता}$ हेही सिद्ध होते असे दाखवू. म्हणून $\ZFminus$ च्या संदर्भात पायाभूतता आणि नियमितता \emph{समतुल्य} आहेत. याचा अर्थ, $\ZFminus$ दिल्यावर नियमिततेशी असलेल्या समतुल्यतेवरून पायाभूततेचे समर्थन करता येते. आणि \stageshier{} च्या आधारे नियमिततेचे समर्थन लगेच करता येते.







\section{$\Z$ आणि $\ZF$: एक महत्त्वाचा टप्पा}\label{sth:spine:zf:sec}

पायाभूततेमुळे आपण आणखी एक महत्त्वाचा टप्पा गाठतो. आपण $\Zminus$ आणि $\ZFminus$ या उपपत्ती पाहिल्या; त्यांतून एक विशिष्ट गोष्ट ``वगळलेली'' आहे असे म्हटले होते. ती गोष्ट म्हणजे पायाभूतता. म्हणून:

\begin{defn}
$\Zminus$ मध्ये पायाभूतता जोडली की $\Z$ ही उपपत्ती मिळते. म्हणून तिची स्वयंसिद्धके विस्तारात्मकता, संयोग, जोड्या, घातसंच, अनंतता, पायाभूतता आणि विभक्तीकरण-रूपबंधातील सर्व स्वयंसिद्धके ही आहेत.

$\ZFminus$ मध्ये पायाभूतता जोडली की $\ZF$ ही उपपत्ती मिळते. दुसऱ्या शब्दांत, $\Z$ मध्ये प्रतिस्थापनाची सर्व स्वयंसिद्धके जोडली की $\ZF$ मिळते.
\end{defn}

तरी एक प्रश्न पडू शकतो: $\ZFminus$ मध्ये नियमितता पायाभूततेशी समतुल्य आहे आणि नियमिततेचे समर्थन स्पष्ट आहे, तर थेट नियमिततेला मूलभूत स्वयंसिद्धक म्हणून न घेता पायाभूततेचा आडमार्ग का घ्यावा?

ऐतिहासिक कारणे (कोणी कोणते तत्त्व कधी मांडले, ही) बाजूला ठेवली, तर मुख्य कारण असे: $V_\alpha$ च्या व्याख्येचा वापर न करता पायाभूतता मांडता येते. पण त्या व्याख्येसाठी \ref{sth:spine:recursion:sec} मधील सर्व काम आवश्यक होते: ती व्याख्या योग्य आहे हे दाखवण्यासाठी अनंतपार पुनरावर्तन सिद्ध करावे लागले. आणि त्या पुराव्यात \emph{प्रतिस्थापन} वापरले. म्हणून $\ZFminus$ च्या संदर्भात पायाभूतता आणि नियमितता समतुल्य असल्या, तरी $\Zminus$ च्या संदर्भात त्या समतुल्य नाहीत.

खरे तर प्रश्न या साध्या निरीक्षणापेक्षा अधिक गंभीर आहे. या पुस्तकाच्या व्याप्तीपलीकडील निष्कर्ष असा की $\Zminus$ आणि $\Z$ या दोन्ही उपपत्ती $V_\alpha$ परिभाषित करण्यास अपुऱ्या आहेत. म्हणून फक्त $\Z$ मध्ये काम करत असाल, तर आपण मांडलेल्या रूपातील नियमिततेला \emph{अर्थच} नाही. म्हणून नियमिततेऐवजी पायाभूतता हेच आपले अधिकृत स्वयंसिद्धक आहे.

यापुढे अन्यथा सांगितले नसेल, तर आपण अधिक काही न सांगता $\ZF$ मध्ये काम करू.







\section{प्रतांक}\label{sth:spine:rank:sec}

आता टप्पे $V_\alpha$ यांच्या रूपात परिभाषित केले आहेत आणि प्रत्येक संच कोणत्या तरी टप्प्याचा उपसंच असतो हे आपल्याला माहीत आहे. त्यामुळे संचाचा \emph{प्रतांक} परिभाषित करता येतो. सहज अर्थाने, $A$ चा प्रतांक म्हणजे $A$ ज्या पहिल्या क्षणी तयार होतो तो क्षण. अधिक नेमकेपणाने:

\begin{defn}\label{sth:spine:rank:defnsetrank}
प्रत्येक संच $A$ साठी $\setrank{A}$ ही $A
\subseteq V_\alpha$ पूर्ण करणारी लघुतम क्रमसूचक संख्या $\alpha$ असते.
\end{defn}
\begin{prop}\label{sth:spine:rank:ranksexist}
	कोणत्याही $A$ साठी $\setrank{A}$ अस्तित्वात असतो.
\end{prop}
\begin{proof}
	हा पुरावा सरावासाठी सोडला आहे.
\end{proof}
\begin{prob}
	\ref{sth:spine:rank:ranksexist} सिद्ध करा.
\end{prob}
\noindent
प्रतांकांचे सुक्रमण आपल्याला काही महत्त्वाचे निष्कर्ष सिद्ध करू देते:
\begin{prop}\label{sth:spine:rank:valphalowerrank}
कोणत्याही क्रमसूचक संख्या $\alpha$ साठी $V_\alpha = \Setabs{x}{\setrank{x} \in \alpha}$.
\end{prop}

\begin{proof}
$\setrank{x} \in \alpha$ असेल, तर $x \subseteq V_{\setrank{x}} \in
V_\alpha$; म्हणून $V_\alpha$ उपसंच-समावेशक असल्याने $x \in V_\alpha$ (येथे \ref{sth:spine:Valphabasic:Valphabasicprops} अनेकदा वापरले आहे). उलट, $x \in V_\alpha$ असेल, तर $x \subseteq V_\alpha$ आणि त्यामुळे $\setrank{x} \leq \alpha$. आता सुलभ अनंतपार विगमनाने समतेची शक्यता नाकारता येते: समता असेल, तर $x \notin V_\alpha$.
\end{proof}
\begin{prob}
	\ref{sth:spine:rank:valphalowerrank} मध्ये उल्लेखलेले सुलभ अनंतपार विगमन पूर्ण करा.
\end{prob}

\begin{prop}\label{sth:spine:rank:rankmemberslower}
$B \in A$ असेल, तर $\setrank{B} \in \setrank{A}$.
\end{prop}

\begin{proof}
\ref{sth:spine:rank:valphalowerrank} वरून $A \subseteq V_{\setrank{A}} = \Setabs{x}{\setrank{x} \in
\setrank{A}}$.
\end{proof}
\noindent
याचा वापर करून \emph{सर्व संचांविषयी} विगमनाच्या एका प्रकाराने निष्कर्ष सिद्ध करता येतो:

\begin{thm}[$\in$-विगमन रूपबंध]
कोणत्याही सूत्र $\phi$ साठी:
\[
	\forall A((\forall x \in A)\phi(x) \lif \phi(A)) \lif \forall A \phi(A).
\]
\end{thm}

\begin{proof}
व्यत्यस्त रूप सिद्ध करू. $\lnot \forall A
\phi(A)$ असे मानू. अनंतपार विगमनाने (\ref{sth:ordinals:basic:ordinductionschema}) $\phi$ पूर्ण न करणारा लघुतम संभाव्य प्रतांकाचा संच मिळतो; म्हणजे $\lnot \phi(A)$ आणि $\forall x(\setrank{x} \in \setrank{A}
\lif \phi(x))$ पूर्ण करणारा एखादा $A$. आता $x \in A$ असेल, तर \ref{sth:spine:rank:rankmemberslower} वरून $\setrank{x} \in
\setrank{A}$; म्हणून $\phi(x)$. म्हणजे $(\forall x \in
A)\phi(x) \land \lnot \phi(A)$.
\end{proof}\noindent या सामर्थ्यवान निष्कर्षाचा अनौपचारिक अर्थ असा सांगता येईल: $\phi$ हा गुणधर्म \emph{आनुवंशिक} असतो तेव्हा आणि केवळ तेव्हाच, जेव्हा कोणत्याही संचाचा प्रत्येक घटक $\phi$ पूर्ण करत असेल, तर तो संचही $\phi$ पूर्ण करतो. मग $\in$-विगमन सांगते: $\phi$ आनुवंशिक असेल, तर प्रत्येक संच $\phi$ पूर्ण करतो.

प्रतांकांवरील चर्चा (सध्या तरी) पूर्ण करण्यासाठी याआधी सूचित केलेले काही दावे सिद्ध करू.

\begin{prop}\label{sth:spine:rank:ranksupstrict}
$\setrank{A} = \supstrict_{x \in A}\setrank{x}$.
\end{prop}

\begin{proof}
$\alpha = \supstrict_{x \in A}\setrank{x}$ ठेवू. \ref{sth:spine:rank:rankmemberslower} वरून $\alpha \leq \setrank{A}$. पण $x \in A$ असेल, तर $\setrank{x} \in \alpha$; म्हणून \ref{sth:spine:rank:valphalowerrank} वरून $x \in V_\alpha$. त्यामुळे $A
\subseteq V_\alpha$, म्हणजे $\setrank{A} \leq \alpha$. म्हणून $\setrank{A} = \alpha$.
\end{proof}

\begin{cor}\label{sth:spine:rank:ordsetrankalpha}
कोणत्याही क्रमसूचक संख्या $\alpha$ साठी $\setrank{\alpha} = \alpha$.
\end{cor}

\begin{proof}
अनंतपार विगमनासाठी सर्व $\beta \in \alpha$ साठी $\setrank{\beta} = \beta$ असे मानू. आता \ref{sth:spine:rank:ranksupstrict} वरून $\setrank{\alpha} = \supstrict_{\beta \in
\alpha}\setrank{\beta} = \supstrict_{\beta \in \alpha}\beta = \alpha$.











\end{proof}

शेवटी, \ref{sth:spine:foundation:sec} च्या अखेरीस सांगितलेला निष्कर्ष झटपट सिद्ध करू: $\ZFminus$ मध्ये $\emph{नियमितता} \Rightarrow \emph{पायाभूतता}$ हे सशर्त विधान सिद्ध होते. (या पुराव्यात ``प्रतांक'' ही संकल्पना आणि \ref{sth:spine:rank:rankmemberslower} वापरता येतात. कारण विभागाच्या सुरुवातीला म्हटल्याप्रमाणे, ते $\ZFminus + \text{नियमितता}$ मध्ये मांडता येतात.)

\begin{prop}[$\ZFminus + \text{नियमितता}$ मध्ये काम करताना]\label{sth:spine:rank:zfminusregularityfoundation} पायाभूतता लागू आहे.
\end{prop}

\begin{proof}
$A \neq \emptyset$ निश्चित करू आणि त्या संचातील लघुतम संभाव्य प्रतांकाचा सदस्य $B \in A$ घेऊ. $c \in B$ असेल, तर \ref{sth:spine:rank:rankmemberslower} वरून $\setrank{c} \in \setrank{B}$; म्हणून $B$ च्या निवडीवरून $c \notin A$.
\end{proof}



\chapter{प्रतिस्थापन}\label{sth:replacement:chap}




\section{परिचय}\label{sth:replacement:intro:sec}

प्रतिस्थापन हा स्वयंसिद्धक-रूपबंधच $\ZF$ आणि~$\Z$ यांतील फरक ठरवतो. \ref{sth:ordinals:chap}--\ref{sth:spine:chap} या प्रकरणांमध्ये आपण त्याचा वापर केला. आता अखेर या प्रकरणात आपण विचारू: प्रतिस्थापनाचे समर्थन करता येते का?

हा प्रश्न नेमका करण्यासाठी प्रतिस्थापन खरेच बरेच \emph{सामर्थ्यवान} आहे, हे लक्षात घ्यावे लागेल. ते किती सामर्थ्यवान आहे याची कल्पना या प्रकरणात (आणि पुन्हा \ref{sth:card-arithmetic:fix:sec} मध्ये) येईल. त्यावरून त्याच्या समर्थनाची खरोखर गरज आहे, हेही दिसेल.

समर्थनाचे दोन प्रकार पाहू. साधारणपणे, स्वयंसिद्धकाच्या फलितांवरून त्याचे समर्थन करण्याचा प्रयत्न म्हणजे \emph{बाह्य} समर्थन; तर संबंधित गणिती संकल्पनांतूनच स्वयंसिद्धकाला पाठबळ मिळते असे दाखवण्याचा प्रयत्न म्हणजे \emph{आंतरिक} समर्थन. याचा अर्थ या प्रकरणात अधिक स्पष्ट होईल; तरी ही चर्चा फार मोठ्या विषयाची केवळ सुरुवात आहे. अधिक माहितीसाठी विशेषतः मॅडी (1988a आणि 1988b) पाहा.







\section{प्रतिस्थापनाचे सामर्थ्य}\label{sth:replacement:strength:sec}

प्रतिस्थापनाच्या सामर्थ्याविषयी एक सोपे निरीक्षण आधी पाहू: $\Z$ च्या पुढे गेल्याशिवाय $\omega + \omega$ इतकी किंवा त्याहून मोठी कोणतीही फोन न्यूमन क्रमसूचक संख्या अस्तित्वात आहे, हे सिद्ध करता येत नाही.

का, याचा आराखडा असा. $\ZF$ मध्ये काम करताना $V_{\omega+\omega}$ हा संच पाहू. हा संच $\Z$ च्या एका \emph{प्रतिमानाचा} प्रांत म्हणून काम करतो. हे समजण्यासाठी सूत्राच्या \emph{सापेक्षीकरणाचा} संकेत ठरवू:
\begin{defn}\label{sth:replacement:strength:formularelativization}
	कोणताही संच $M$ आणि कोणतेही सूत्र $\phi$ दिल्यास, $\phi$ मधील सर्व संख्यापक $M$ पुरते मर्यादित केल्यावर मिळणारे सूत्र $\phi^M$ असे मानू. म्हणजे ``$\lexists[x]$'' ऐवजी ``$(\lexists[x \in M])$'' आणि ``$\lforall[x]$'' ऐवजी ``$(\lforall[x \in M])$'' लिहू.
\end{defn}\noindent
$\Z$ च्या प्रत्येक स्वयंसिद्धक $\phi$ साठी $\ZF \vdash
\phi^{V_{\omega+\omega}}$ हे दाखवता येते. पण \ref{sth:spine:rank:ordsetrankalpha} वरून $\omega+\omega$ ही संख्या $V_{\omega+\omega}$ मध्ये \emph{नसते}. म्हणून $\omega+\omega$ अस्तित्वात नसण्याशी $\Z$ सुसंगत आहे.

म्हणूनच \ref{sth:ordinals:replacement:sec} मध्ये आपण म्हटले की \ref{sth:ordinals:ordtype:thmOrdinalRepresentation} प्रतिस्थापनाशिवाय सिद्ध करता येत नाही. कारण $\Z$ मध्येच असा स्पष्ट सुक्रमण-संबंध परिभाषित करणे सोपे आहे की ज्याचा क्रमप्रकार सहजपणे $\omega+\omega$ \emph{असायला हवा}. नैसर्गिक संख्यांवरील पुढील क्रम दाखवताना \ref{sth:ordinals:idea:sec} मध्ये आपण याचे अनौपचारिक उदाहरण दिले होते:
\begin{align*}
	n \lessdot m \text{ तेव्हा आणि केवळ तेव्हाच }&\text{एकतर }n < m\text{ आणि }m-n\text{ सम आहे,}\\
	& \text{किंवा $n$ सम आणि $m$ विषम आहेत.}
\end{align*}
पण $\omega+\omega$ अस्तित्वात नसेल, तर हे सुक्रमण कोणत्याही क्रमसूचक संख्येशी समरूपी नसते. म्हणून $\Z$ वरून \ref{sth:ordinals:ordtype:thmOrdinalRepresentation} \emph{सिद्ध होत नाही}.

उलट पाहिल्यास: प्रतिस्थापनामुळे $\omega+\omega$ चे अस्तित्व, आणि म्हणून $V_{\omega+\omega}$ चे अस्तित्व, सिद्ध करता येते. एवढेच नव्हे. आपण परिभाषित करू शकणाऱ्या \emph{कोणत्याही} सुक्रमणासाठी \ref{sth:ordinals:ordtype:thmOrdinalRepresentation} सांगते की त्याच्याशी समरूपी अशी काही $\alpha$ असते; म्हणून $V_\alpha$ अस्तित्वात असते. अशा थेट अर्थाने प्रतिस्थापनामुळे संचांचा पदानुक्रम \emph{फार उंच} असण्याची हमी मिळते.

पुढील काही विभागांत आणि नंतर \ref{sth:card-arithmetic:fix:sec} मध्ये प्रतिस्थापनामुळे पदानुक्रम नेमका \emph{किती} उंच जातो याची अधिक चांगली कल्पना येईल. सध्या एवढाच साधा मुद्दा: प्रतिस्थापनाचे समर्थन \emph{खरोखर आवश्यक} आहे!







\section{प्रतिस्थापनाविषयी बाह्य विचार}\label{sth:replacement:extrinsic:sec}

आधी प्रतिस्थापनाचे \emph{बाह्य} समर्थन करण्याचा एक प्रयत्न पाहू. बूलोस यांनी असा युक्तिवाद सुचवला आहे:
\begin{quote}
  [\ldots] प्रतिस्थापनाची स्वयंसिद्धके स्वीकारण्याचे कारण अगदी सोपे आहे: त्यांचे अनेक इष्ट परिणाम आहेत आणि (दिसते तसे) अनिष्ट परिणाम नाहीत. क्रमिक संकल्पनेविषयीच्या प्रमेयांखेरीज त्या परिणामांत अनंत संख्यांची आदर्श नसली तरी समाधानकारक उपपत्ती आणि सुस्थापित संबंधांवरील विगमनाधारित व्याख्यांना समर्थन देणारा अत्यंत उपयुक्त निष्कर्ष येतो.
  (बूलोस 1971, पृ.~229)
\end{quote}
बूलोस यांच्या कल्पनेचा सारांश असा की प्रतिस्थापनाचे समर्थन त्याच्या फलितांवरून करावे. त्यांनी सांगितलेली विशिष्ट फलिते आपण मागील काही प्रकरणांत पाहिली आहेत. सुक्रमणाच्या क्रमप्रकारासाठी फोन न्यूमन क्रमसूचक संख्या उत्तम प्रतिनिधी ठरतात, हे प्रतिस्थापनामुळे सिद्ध करता आले (हीच आपली ``अनंत संख्यांची आदर्श नसली तरी समाधानकारक उपपत्ती''). प्रतिस्थापनामुळे $V_\alpha$ परिभाषित करता आले, प्रतांकाची संकल्पना मांडता आली आणि $\in$-विगमन सिद्ध करता आले (हीच ``क्रमिक संकल्पनेविषयीची प्रमेये''). अखेरीस, प्रतिस्थापनामुळे अनंतपार पुनरावर्तनाचे प्रमेय सिद्ध करता येते (हीच ``सुस्थापित संबंधांवरील विगमनाधारित व्याख्या'').

हे परिणाम खरोखर इष्ट आहेत. पण इतके इष्ट परिणाम प्रतिस्थापनाला \emph{समर्थन} देण्यासाठी पुरेसे आहेत का? \emph{नाही}. निदान सरळपणे तरी नाही.

एक सोपी अडचण अशी आहे: प्रतिस्थापनाचे काही इष्ट परिणाम आपण मांडले, पण त्यांपैकी अनेक परिणाम दुसऱ्या मार्गांनीही मिळू शकले असते. ही बाब जितकी माहीत असायला हवी तितकी नाही; म्हणून परिस्थिती स्पष्ट करू.

संचांची एक सोपी उपपत्ती आहे: स्तर उपपत्ती, संक्षेपाने $\LT$.\footnote{$\LT$ ची पहिली रूपे मॉन्टेग्यू (1965) आणि स्कॉट (1974) यांनी दिली. पॉटर (2004) यांनी ती सोपी करून पुस्तकरूपात सविस्तर मांडली; अलीकडे बटन (2021) यांनी $\LT$ आणखी सोपी केली.} $\LT$ ची स्वयंसिद्धके फक्त विस्तारात्मकता, विभक्तीकरण आणि प्रत्येक संच कोणत्या तरी \emph{स्तराचा} उपसंच असतो हे विधान इतकीच आहेत. येथे ``स्तर'' अशी कल्पक व्याख्या आहे की हे स्तर आपल्या परिचयाच्या $V_\alpha$ प्रमाणे वागतात. म्हणून $\ZF$ वरून $\LT$ सिद्ध होते; पण $\LT$ ही $\ZF$ पेक्षा \emph{खूप} दुर्बल आहे. खरे तर $\LT$ वरून जोड्या, घातसंच, अनंतता किंवा प्रतिस्थापन यांपैकी काहीही मिळत नाही. $\LT$ मध्ये अनंतता आणि घातसंच जोडून मिळणाऱ्या उपपत्तीला $\Zr$ म्हणू. यावरून जोड्याही मिळतात; म्हणून $\Zr$ निदान $\Z$ इतकी सामर्थ्यवान आहे. पण प्रत्यक्षात $\Zr$ ही $\Z$ पेक्षा काटेकोरपणे अधिक सामर्थ्यवान आहे; कारण प्रत्येक संचाला प्रतांक असतो हे विधान ती जोडते (म्हणूनच तिला $\Zr$ म्हणावे असे मी सुचवतो). $\Zr$ वरून क्रमसूचक संख्यांची पूर्णपणे समाधानकारक उपपत्ती, पदानुक्रमाचे सुक्रमित टप्प्यांत विभाजन करणारे निष्कर्ष, $\in$-विगमनाचा पुरावा आणि अनंतपार पुनरावर्तनाचे एक \emph{रूप} मिळते.

थोडक्यात, बूलोस यांना ही बाब माहीत नसली, तरी त्यांनी सांगितलेले सर्व इष्ट परिणाम प्रतिस्थापनाशिवाय \emph{मिळू शकले असते}; त्यासाठी $\Z$ ऐवजी $\Zr$ वापरणे पुरेसे होते.

(मग आपण प्रचलित पद्धत वापरून $\LT$ आणि $\Zr$ ऐवजी $\ZF$ का शिकवली? दोन कारणे आहेत. पहिले: निव्वळ ऐतिहासिक कारणांमुळे $\LT$ पासून सुरुवात करणे अपारंपरिक आहे; संच उपपत्तीवरील अधिक प्रमाणित चर्चा वाचता याव्यात अशी आपली इच्छा होती. दुसरे: $\LT$ आणि $\Zr$ समजून घेण्याची तयारी झाल्यावर पॉटर (2004) आणि बटन (2021) थेट वाचता येतील.)

अर्थात $\Zr$ ही $\ZF$ पेक्षा काटेकोरपणे दुर्बल असल्याने, $\ZF$ सिद्ध करते पण $\Zr$ अनिर्णित ठेवते असे निष्कर्ष आहेत. त्यांपैकी एखाद्या निष्कर्षाच्या आधारे प्रतिस्थापनाचे बाह्य समर्थन करण्याचा प्रयत्न करता येईल. पण $\Zr$ कशी वापरावी हे समजल्यावर (अ) प्रतिस्थापनामुळेच, दुसऱ्या मार्गाने नाही, ठरणाऱ्या आणि (ब) सहज अर्थाने खऱ्या वाटणाऱ्या विधानांची अनेक उदाहरणे शोधणे बरेच कठीण असते. (अधिक माहितीसाठी पॉटर (2004, \S13.2) पाहा.)

मुख्य निष्कर्ष असा: प्रतिस्थापनाचे पटणारे बाह्य समर्थन द्यायचे असेल, तर प्रतिस्थापनाशिवाय \emph{मिळूच शकणार नाही} असा निष्कर्ष शोधावा लागेल. हे सोपे काम नाही.

प्रतिस्थापनाच्या निव्वळ बाह्य समर्थनाच्या कोणत्याही प्रयत्नासमोर आणखी एक अडचण उभी राहते; कदाचित ती पहिल्यापेक्षा अधिक मूलभूत आहे. बूलोस फक्त प्रतिस्थापनाचे अनेक इष्ट परिणाम सांगत नाहीत; त्याचे ``(दिसते तसे) अनिष्ट'' परिणाम नाहीत, असेही म्हणतात. पण कंसातील ``दिसते तसे'' ही सावधगिरीची नोंद अतिशय महत्त्वाची आहे.

आपण येथे कसे पोहोचलो ते आठवू. अनौपचारिक गुणधर्माधारित संचरचनेत विसंगती निर्माण झाली; त्याला प्रतिसाद म्हणून आपण संचांची संचयी-क्रमिक संकल्पना स्वीकारली. या संकल्पनेबरोबर एक मांडणी येते आणि ती आपल्याला तिच्या सुसंगततेची खात्री देईल, अशी आशा आहे. पण त्या मांडणीच्या आतून प्रतिस्थापनाचे समर्थन करता आले नाही, तर $\ZF$ सुसंगत आहे असे मानण्याचे (अजून) कोणतेही कारण आपल्याकडे नाही. किंवा अधिक नेमके: आत्तापर्यंत कोणालाही \emph{विसंगती सापडलेली नाही} या (कदाचित निव्वळ योगायोगाच्या) वस्तुस्थितीखेरीज $\ZF$ च्या सुसंगततेवर विश्वास ठेवण्याचे कारण नाही. नेमक्या याच अर्थाने बूलोस यांचे विधान ``(दिसते तसे) $\ZF$ सुसंगत आहे'' इतकेच ठरते. सुसंगततेबाबत याहून अधिक खात्री आपण मागितली पाहिजे.

$\ZF$ च्या (ज्ञात) परिणामांपुरतेच मर्यादित असलेले, म्हणजे प्रतिस्थापनाचे \emph{निव्वळ} बाह्य समर्थन करण्याचे, कोणतेही प्रयत्न या अडचणीला सामोरे जातील. म्हणून आपण प्रतिस्थापनाचे \emph{आंतरिक} समर्थन शोधू: संचांविषयी आपण सांगितलेली मांडणी आपल्याला आधीपासूनच प्रतिस्थापन स्वीकारण्यास बांधील करते, असे सुचवणारे समर्थन.








\section{आकारमानाच्या मर्यादेचे तत्त्व}\label{sth:replacement:limofsize:sec}

प्रतिस्थापनाचे ``आंतरिक'' समर्थन देण्याचा बहुधा सर्वाधिक प्रचलित प्रयत्न पुढील कल्पनेवर आधारलेला असतो:
\begin{enumerate}
  \item[] \limofsize. कोणत्याही वस्तूंचा संच तयार होतो, जर त्या फार जास्त नसतील तर.
\end{enumerate}
आकारमानाच्या पुढील तुलनेसाठी निवडीचे स्वयंसिद्धक गृहीत धरून, हे तत्त्व प्रतिस्थापनाला लगेच समर्थन देईल. कारण प्रतिस्थापनाने तयार झालेला कोणताही संच ज्या संचापासून तयार झाला त्याच्यापेक्षा मोठा असू शकत नाही. नेमके सांगायचे तर: समजा प्रतिस्थापन वापरून तुम्ही
$\funimage{\tau}{A} = \Setabs{\tau(x)}{x \in A}$ हा संच तयार केला; मग $\cardle{\funimage{\tau}{A}}{A}$; म्हणून $A$ चे सदस्य संच तयार होण्याइतके कमी असतील, तर त्यांच्या प्रतिमाही $\funimage{\tau}{A}$ तयार होण्याइतक्या कमी असतील.

प्रतिस्थापनाचे समर्थन करण्यासाठी \limofsize{} वापरण्यातील उघड अडचण अशी की \limofsize{} सारखे कोणतेही तत्त्व आपण अजून मांडलेलेच \emph{नाही}. शिवाय, \ref{sth:story:chap}--\ref{sth:z:chap} मध्ये संचांच्या संचयी-क्रमिक संकल्पनेची मांडणी करताना \limofsize{} कडे निर्देश करणारा \emph{संकेतही} त्यात नव्हता. म्हणूनच बूलोस यांनी एकदा असे लिहिले: ``संच उपपत्तीच्या `मागे' निदान दोन विचार आहेत, असा निष्कर्ष कदाचित काढता येईल'' (1989, पृ.~19). एका बाजूला, संचांच्या संचयी-क्रमिक संकल्पनेशी संबंधित कल्पना $\Z$ ला समर्थन देण्यासाठी आहेत. दुसऱ्या बाजूला, \limofsize{} हे प्रतिस्थापनाला समर्थन देण्यासाठी आहे.

पण अडचण एवढीच नाही की \limofsize{} बाबत आपण आतापर्यंत \emph{मौन} बाळगले आहे. उलट, नुकतेच मांडलेले \limofsize{} तत्त्व संचांच्या संचयी-क्रमिक संकल्पनेशी जुळणेही अवघड वाटते. कारण संच \emph{टप्प्यांत} तयार होतात, या कल्पनेचा त्यात अजिबात उल्लेख नाही.

या ``दोन विचारांचा'' (म्हणजे संचांची संचयी-क्रमिक संकल्पना आणि \limofsize) इतिहास पाहता हे फारसे आश्चर्यकारक नाही. संच उपपत्तीतील विरोधाभास रोखण्याचे हे मुळात दोन भिन्न प्रकल्प आहेत. संचांची संचयी-क्रमिक संकल्पना रसेलचा विरोधाभास साधारणपणे असे म्हणून रोखते: \emph{रसेलचा संच अस्तित्वात असेल अशी अपेक्षाच आपण करायला नको होती, कारण तो कोणत्याही टप्प्यावर ``तयार'' होणार नाही}. याउलट, \limofsize{} रसेलचा संच साधारणपणे असे म्हणून नाकारते: \emph{रसेलचा संच अस्तित्वात असेल अशी अपेक्षाच आपण करायला नको होती, कारण तो फार मोठा असता}.

असे मांडल्यावर स्पष्टच सांगू: विरोधाभासांना उत्तर म्हणून \limofsize{} ला \emph{खूपच} अधिक समर्थन आवश्यक आहे. उदाहरणार्थ, रसेलच्या विरोधाभासाचे हे रूप पाहा: \emph{स्वतःचाच वास न घेणाऱ्या पग कुत्र्यांचाच नेमका वास घेणारा कोणताही पग कुत्रा नाही} (\ref{sth:story:rus:sec} पाहा). ``असा पग कुत्रा का नाही?'' असे विचारल्यावर त्याला फारच जास्त पग कुत्र्यांचा वास घ्यावा लागला असता, हे चांगले उत्तर ठरणार नाही. मग रसेलचा संच अस्तित्वात नाही याचे सहज पटणारे स्पष्टीकरण तो अस्तित्वात असण्यासाठी ``फार मोठा'' झाला असता, हे कसे ठरेल?

थोडक्यात, \limofsize{} मागील ``सहजस्फूर्त'' प्रेरणेविषयी तुम्हाला थोडे गूढ वाटले, तर ते समजण्यासारखे आहे.







\section{प्रतिस्थापन आणि ``निरपेक्ष अनंतता''}\label{sth:replacement:absinf:sec}

आता \limofsize{} बाजूला ठेवून प्रतिस्थापनाचे समर्थन करण्याचे वेगळ्या प्रकारचे (आंतरिक) प्रयत्न पाहू. हे प्रयत्न संच टप्प्याटप्प्याने तयार होतात या कल्पनेला गांभीर्याने घेतात.

क्रमिक प्रक्रिया प्रथम सांगताना आपण प्रत्येक टप्प्यावर काय घडते हे स्पष्ट करणारी काही तत्त्वे मांडली होती: \stageshier, \stagesord आणि \stagesacc. नंतर टप्प्यांच्या संख्येविषयी काही सांगणारी तत्त्वे जोडली: \stagessucc{} नुसार संच तयार होण्याची प्रक्रिया कधीही संपत नाही, आणि \stagesinf{} नुसार ती अनंताव्या टप्प्यापर्यंत पोहोचते.

ही नंतरची दोन तत्त्वे आणखी व्यापक तत्त्वातून निष्पन्न होतात असे सुचवणे रास्त आहे; उदाहरणार्थ:
\begin{enumerate}
  \item[] \stagesinex. टप्प्यांची संख्या निरपेक्ष अर्थाने अनंत आहे; पदानुक्रम जितका उंच असू शकतो तितका उंच आहे.
\end{enumerate}
अर्थात हे एक अनौपचारिक तत्त्व आहे. पण संचांच्या संचयी-क्रमिक संकल्पनेतून ते लगेच \emph{निष्पन्न} होत नसले, तरी त्या संकल्पनेशी निश्चितच \emph{सुसंगत} वाटते. निदान \limofsize{} च्या विपरीत, संच टप्प्याटप्प्याने तयार होतात ही कल्पना त्यात टिकून राहते.

आता \stagesinex{} वापरून प्रतिस्थापनाचे समर्थन करता येईल अशी आशा आहे. ते कसे शक्य होईल ते पाहू.

\ref{sth:ordinals:idea:sec} मध्ये आपण पाहिले की (अनौपचारिक अर्थाने) $\omega+\omega$ शी समरूप असायला हवे असे सुक्रमण सहज तयार करता येते. दुसऱ्या शब्दांत, (अनौपचारिक अर्थाने) $\omega+\omega$ हा क्रमप्रकार असलेली टप्प्याटप्प्यांची क्रमिक प्रक्रिया आपण सहज कल्पू शकतो. म्हणून \stagesinex{} स्वीकारले असेल, तर पदानुक्रमात किमान $\omega+\omega$-वा टप्पा, म्हणजे $V_{\omega+\omega}$, आहे हेही स्वीकारले पाहिजे; कारण पदानुक्रम इतक्यापर्यंत \emph{नक्कीच जाऊ शकतो}.

हा विचार सर्वसाधारणपणे असा मांडता येतो: कोणतेही सुक्रमण घ्या; क्रमिक पदानुक्रम बांधण्याची प्रक्रिया किमान त्या सुक्रमणाइतकी पुढे गेली पाहिजे. \ref{sth:ordinals:ordtype:thmOrdinalRepresentation} याला \emph{स्वयंसिद्धक} मानल्यास याची खात्री मिळेल. मग प्रत्येक सुक्रमण एखाद्या फोन न्यूमन क्रमसूचक संख्येशी समरूप आहे असे ठरेल. प्रत्येक फोन न्यूमन क्रमसूचक संख्येचा प्रतांक ती स्वतःच असल्याने, \ref{sth:ordinals:ordtype:thmOrdinalRepresentation} यावरून असे ठरेल की संच उपपत्तीत आपण कोणतेही सुक्रमण वर्णन करू शकलो, तर संच बांधण्याची क्रमिक प्रक्रिया त्या सुक्रमणाच्याही पुढे गेली पाहिजे.

हा विचार \stagesinex{} चा परिणाम वाटतो. दुर्दैवाने, त्यावरून प्रतिस्थापन मिळवायचे असेल, तर एक सोपा तांत्रिक अडथळा आहे: प्रतिस्थापन हे \ref{sth:ordinals:ordtype:thmOrdinalRepresentation} पेक्षा काटेकोरपणे अधिक सामर्थ्यवान आहे. (ही बाब पॉटर (2004, \S13.2) यांनी नोंदवली आहे; तिचा पुरावा आपण \ref{sth:replacement:finiteaxiomatizability:sec} मध्ये देऊ.)

म्हणून \stagesinex{} चा अर्थ प्रतिस्थापन मिळेल अशा रीतीने घ्यायचा असेल, तर पदानुक्रम प्रत्येक सुक्रमणाच्या पुढे जातो \emph{एवढेच} त्याने सांगून चालणार नाही. त्याहून बलवान दावा आवश्यक आहे. यासाठी शोनफेल्ड (1977) यांनी या कल्पनेचा एक अत्यंत स्वाभाविक विस्तार सुचवला: पदानुक्रम कोणत्याही संचाशी \emph{सहअंतिम} नाही.\footnote{G\"odel यांनीही असाच विचार सुचवला असावा; पॉटर (2004, पृ.~223) पाहा. G\"odel आणि शोनफेल्ड यांच्या चर्चेसाठी इनकुर्वाती (2020, पृ.~90--95) पाहा.} अधिक स्पष्टपणे: $\tau$ हे संचांना पदानुक्रमातील टप्प्यांकडे पाठवणारे प्रतिचित्रण असेल, तर कोणत्याही संच $A$ ची $\tau$ खालील प्रतिमा संपूर्ण पदानुक्रम संपवू शकत नाही. दुसऱ्या शब्दांत (रूपबंधाच्या स्वरूपात):
\begin{enumerate}
  \item[] \stagescofin. $A$ हा संच असेल आणि प्रत्येक $x \in A$ साठी $\tau(x)$ हा टप्पा असेल, तर प्रत्येक $x \in A$ च्या $\tau(x)$ नंतर येणारा एक टप्पा असतो.
\end{enumerate}
$\ZF$ मध्ये \stagescofin{} चे योग्य औपचारिक रूप सिद्ध होते, हे स्पष्ट आहे. उलट, \stagescofin{} प्रतिस्थापनाचे समर्थन करते, असा अनौपचारिक युक्तिवाद करता येतो.\footnote{\stagescofin{} च्या एखाद्या औपचारिक रूपावरून प्रतिस्थापन सिद्ध करणे अधिक कठीण ठरेल; कारण केवळ $\Z$ मध्ये टप्प्यांची व्याख्या करण्याइतके सामर्थ्य नाही. त्यामुळे \stagescofin{} चे औपचारिकीकरण कसे करावे हे स्पष्ट नाही. मात्र \ref{sth:replacement:extrinsic:sec} मध्ये चर्चिल्याप्रमाणे $\LT$ च्या एखाद्या विस्तारात काम करणे हा एक पर्याय आहे.} समजा $(\forall x \in A)\lexists![y][\phi(x,y)]$. मग प्रत्येक $x \in A$ साठी $\phi(x,y)$ पूर्ण करणारा $y$ म्हणजे $\sigma(x)$ मानू, आणि $\sigma(x)$ पहिल्यांदा ज्या टप्प्यावर तयार होतो तो $\tau(x)$ मानू. \stagescofin{} नुसार असा टप्पा $V$ आहे की $(\forall x \in A)\tau(x)\in V$. आता प्रत्येक $\tau(x) \in V$ आणि $\sigma(x) \subseteq \tau(x)$ असल्याने, विभक्तीकरण वापरून $\Setabs{y
\in V}{(\exists x \in A)\sigma(x) = y} = \Setabs{y}{(\exists x \in
A)\phi(x,y)}$ हा संच मिळवता येतो.

\begin{prob}
  $\ZF$ मध्ये \stagescofin{} चे औपचारिकीकरण करा.
\end{prob}

म्हणून \stagescofin{} प्रतिस्थापनाला समर्थन देते. आणि \stagesinex{} हे \stagescofin{} ला समर्थन देते, हे निदान संभवनीय वाटते. समजा \stagescofin{} खोटे ठरते. म्हणजे पदानुक्रम एखाद्या संच~$A$ शी सहअंतिम आहे: असे $\tau$ प्रतिचित्रण आहे की प्रत्येक टप्पा~$S$ साठी असा काही $x \in A$ आहे की $S \in \tau(x)$. अशा परिस्थितीत पदानुक्रमाची कथित ``निरपेक्ष अनंतता'' पकडण्याचा मार्ग आपल्याकडे आहे: $A$ वर लागू केलेल्या $\tau$ च्या परासाने ती \emph{संपूर्णपणे व्यापली} जाते. त्यामुळे पदानुक्रम ``निरपेक्षपणे अनंत'' आहे या विचाराला बाधा येते. व्यत्यासाने: \stagesinex{} वरून \stagescofin{} मिळते; आणि ते प्रतिस्थापनाचे समर्थन करते.

प्रतिस्थापनाचे आंतरिक समर्थन देण्याचा हा खरोखर आशादायक प्रयत्न आहे. मात्र तो अखेरीस सफल होतो की नाही, हे ठरवण्याचे काम आम्ही तुम्हालाच सोपवतो.







\section{प्रतिस्थापन आणि प्रतिबिंबन}\label{sth:replacement:ref:sec}

\stagesinex{} च्या साहाय्याने प्रतिस्थापनाचे समर्थन करण्याचा आपला शेवटचा प्रयत्न एका सखोल आणि सुंदर निष्कर्षापासून सुरू होतो:\footnote{आठवण: स्पष्टपणे अन्यथा म्हटले नसेल, तर सर्व सूत्रांमध्ये प्राचले असू शकतात.}


\begin{thm}[प्रतिबिंबन रूपबंध]\label{sth:replacement:ref:reflectionschema}
कोणत्याही सूत्र $\phi$ साठी:
\[
\forall \alpha \exists \beta > \alpha (\forall x_1 \ldots, x_n \in
V_\beta)(\phi(x_1, \ldots, x_n) \liff \phi^{V_\beta}(x_1, \ldots, x_n))
\]
\end{thm}
\noindent
\ref{sth:replacement:strength:formularelativization} प्रमाणे, $\phi$ मधील प्रत्येक संख्यक $V_\beta$ या संचापुरता मर्यादित केल्यावर $\phi^{V_\beta}$ मिळते. म्हणून सहज अर्थाने प्रतिबिंबन असे सांगते: $\phi$ संपूर्ण पदानुक्रमात सत्य असेल, तर पदानुक्रमाच्या हवे तितक्या अनेक \emph{आरंभीच्या खंडांमध्ये}ही $\phi$ सत्य असते.

मॉन्टेग्यू (1961) आणि लेव्ही (1960) यांनी दाखवले की (योग्यरीत्या मांडलेली) प्रतिस्थापन आणि प्रतिबिंबन ही $\Z$ च्या साहाय्याने समतुल्य आहेत; म्हणजे त्यांपैकी कोणतेही एक जोडल्यास $\ZF$ मिळते. (हे निष्कर्ष आपण \ref{sth:replacement:refproofs:sec} मध्ये सिद्ध करू.) ही समतुल्यता लक्षात घेऊन, \stagesinex{} वरून प्रतिबिंबन आणि प्रतिस्थापनाचे समर्थन करण्याचा पुढील प्रयत्न करता येईल: \stagesinex{} असेल, तर पदानुक्रम खूपच उंच असायला हवा; इतका उंच की त्याविषयी आपण म्हणू शकतो अशा कोणत्याही गोष्टीने त्याच्या उंचीला बंध घालता कामा नये. हा विचार आपण असा समजू शकतो: कोणतेही वाक्य~$\phi$ संपूर्ण पदानुक्रमात सत्य असेल, तर ते पदानुक्रमाच्या हवे तितक्या अनेक आरंभीच्या खंडांतही सत्य असते. हेच प्रतिबिंबन आहे.

पुन्हा, प्रतिस्थापनाचे आंतरिक समर्थन देण्याचा हा खरोखर आशादायक प्रयत्न वाटतो. पण याविषयी सांगण्यासारखे येथे मावेल त्याहून खूप अधिक आहे. तो सफल होतो की नाही, हे आता तुम्ही स्वतः ठरवावे.\footnote{पुढे वाचायचे असल्यास इनकुर्वाती (2020, पृ.~95--100) उपयुक्त ठरेल.}







\section{परिशिष्ट: प्रतिस्थापनाशी संबंधित निष्कर्ष}\label{sth:replacement:refproofs:sec}

या विभागात आपण $\ZF$ मध्ये प्रतिबिंबन सिद्ध करू. प्रतिबिंबन आणि प्रतिस्थापन कोणत्या अर्थाने समतुल्य आहेत हेही सिद्ध करू. तसेच प्रतिबिंबन/प्रतिस्थापन यांच्या सामर्थ्याविषयी या सर्वांचा एक रोचक परिणाम सिद्ध करू. \emph{सूचना: या पाठ्यपुस्तकातील गणिताचा हा सहजपणे सर्वांत प्रगत भाग आहे.}

सुरुवातीला एक साहाय्यक सिद्धांत पाहू. चलांच्या किंवा वस्तूंच्या अनुक्रमांसाठी संक्षेप म्हणून तो \emph{वरची रेघ} ही लेखनपद्धत वापरतो. म्हणजे ``$\overline{a}_k$'' हे ``$a_{k_1}$, \dots,
$a_{k_n}$'' यांचे संक्षिप्त लेखन आहे; येथे $n$ संदर्भावरून ठरतो.

\begin{lem}\label{sth:replacement:refproofs:lemreflection}
प्रत्येक $1 \leq i \leq k$ साठी $\phi_i(\overline{v}_i, x)$ हे सूत्र असू द्या. मग प्रत्येक $\alpha$ साठी काही $\beta > \alpha$ असा असतो की, कोणतेही $\overline{a}_1,
\ldots, \overline{a}_k \in V_\beta$ आणि प्रत्येक $1 \leq i \leq k$ यांच्यासाठी:
\[
  \exists x\phi_i(\overline{a}_i, x) \rightarrow (\exists x \in V_\beta) \phi_i(\overline{a}_i, x)
\]
\end{lem}

\begin{proof}
पुढीलप्रमाणे $\mu$ हे पद परिभाषित करू: $\mu(\overline{a}_1, \ldots,
\overline{a}_k)$ हा पुढील सर्व अटी, प्रत्येक $1 \leq i \leq k$ साठी, पूर्ण करणारा लघुतम टप्पा $V$ आहे:
\[
\exists x\phi_i(\overline{a}_i, x) \rightarrow (\exists x \in V) \phi_i(\overline{a}_i, x)
\]
सर्व $\overline{a}_1, \ldots, \overline{a}_k$ साठी $\mu(\overline{a}_1, \ldots, \overline{a}_k)$ अस्तित्वात आहे, हे सहज तपासता येते. आता प्रतिस्थापन आणि आपले पुनरावर्तन प्रमेय वापरून परिभाषित करू:
\begin{align*}
  S_0 & = V_{\alpha+1}\\
  S_{n+1} & = S_n \cup \bigcup
  \Setabs{\mu(\overline{a}_1, \ldots, \overline{a}_k)}
  {\overline{a}_1, \ldots, \overline{a}_k \in S_n} \\
  S &= \bigcup_{m < \omega} S_m.
\end{align*}
प्रत्येक $S_n$, आणि त्यामुळे $S$ सुद्धा, $V_\alpha$ नंतरचा टप्पा आहे. आता $\overline{a}_1$, \dots,~$\overline{a}_k \in S$ निश्चित करा; म्हणून काही $n <
\omega$ साठी $\overline{a}_1$, \dots, $\overline{a}_k \in S_n$. काही $1 \leq i \leq k$ निश्चित करा आणि $\exists x
\phi_i(\overline{a}_i,x)$ मानू. मग रचनेनुसार $(\exists x \in \mu(\overline{a}_1,
\ldots, \overline{a}_k))\phi_i(\overline{a}_i, x)$; म्हणून $(\exists x \in S_{n+1})\phi_i(\overline{a}_i, x)$ आणि परिणामी $(\exists x \in S)\phi_i(\overline{a}_i, x)$. म्हणून हाच $S$ आपला $V_\beta$ आहे.
\end{proof}
\noindent
आता \ref{sth:replacement:ref:reflectionschema} सहज सिद्ध करता येतो:

\begin{proof}[पुरावा]
$\alpha$ निश्चित करा. व्यापकता न गमावता, $\phi$ मध्ये फक्त $\exists$, $\lnot$ आणि $\land$ हेच तार्किक संयोगक आहेत असे मानता येते (कारण अभिव्यक्तीसाठी ते पुरेसे आहेत). $\psi_1, \ldots, \psi_k$ या $\phi$ च्या सर्व उपसूत्रांची जटिलतेनुसार यादी असू द्या; त्यामुळे $\psi_k = \phi$. \ref{sth:replacement:refproofs:lemreflection} नुसार असा $\beta > \alpha$ आहे की, कोणतेही $\overline{a}_i \in V_\beta$ आणि प्रत्येक $1 \leq i \leq k$ साठी:
\begin{align}\label{reflectionnicelybehaved}
  \exists x\psi_i(\overline{a}_i, x) \rightarrow
  (\exists x \in V_\beta) \psi_i(\overline{a}_i, x)\tag{*}
\end{align}
$\psi_i$ च्या जटिलतेवर विगमन करून, कोणत्याही $\overline{a}_i \in V_\beta$ साठी $\psi_i(\overline{a}_i) \leftrightarrow
\psi_i^{V_\beta}(\overline{a}_i)$ हे दाखवू. $\psi_i$ अणुसूत्र असेल, तर हे उघड आहे. ही द्विशर्त असेही दाखवते की, हा गुणधर्म असलेल्या उपसूत्रांचे निषेध किंवा संधीकरण $\psi_i$ असेल, तर $\psi_i$ लाही हा गुणधर्म असतो. म्हणून संख्यकाचे प्रकरणच लक्ष देण्यासारखे आहे. $\overline{a}_i \in V_\beta$ निश्चित करा; मग:
\begin{align*}
  (\exists x \psi_i(\overline{a}_i, x))^{V_\beta}
  &\text{ तेव्हा आणि केवळ तेव्हाच }
  (\exists x \in V_\beta)\psi_i^{V_\beta}(\overline{a}_i, x)
  &&\text{व्याख्येनुसार}\\
  &\text{ तेव्हा आणि केवळ तेव्हाच }
  (\exists x \in V_\beta)\psi_i(\overline{a}_i,  x)
  &&\text{विगमनगृहीतकानुसार}\\
  &\text{ तेव्हा आणि केवळ तेव्हाच }
  \exists x \psi_i(\overline{a}_i, x)
  &&\text{\eqref{reflectionnicelybehaved} नुसार}
\end{align*}
याने विगमन पूर्ण होते; $\psi_k = \phi$ असल्याने अपेक्षित निष्कर्ष मिळतो.
\end{proof}

आपण $\ZF$ मध्ये प्रतिबिंबन सिद्ध केले आहे. आपला पुरावा मूलतः मॉन्टेग्यू (1961) यांच्या पुराव्यानुसार आहे. आता $\Z$ मध्ये प्रतिबिंबनावरून प्रतिस्थापन मिळते हे सिद्ध करायचे आहे. हा पुरावा लेव्ही (1960) यांच्या पुराव्यावर आधारित आहे, पण थोडा सोपा केला आहे.

आपण $\Z$ मध्ये काम करत असल्याने वर दिलेल्या अचूक रूपात प्रतिबिंबन मांडता येणार नाही. कारण प्रतिबिंबन मांडताना आपण ``$V_\alpha$'' हे लेखन वापरले, आणि त्याची व्याख्या $\Z$ मध्ये करता येत नाही (\ref{sth:spine:zf:sec} पाहा). म्हणून त्याऐवजी प्रतिबिंबनाचे वरकरणी दुर्बल रूप पुढीलप्रमाणे मांडू:

\begin{defish}
\emph{दुर्बल प्रतिबिंबन.} कोणतेही सूत्र $\phi$ घ्या. मग असा संक्रमक संच $S$ असतो की $0$, $1$ आणि $\phi$ मधील कोणतीही प्राचले हे $S$ चे सदस्य असतात, आणि $(\forall \overline{x} \in S)(\phi \liff
\phi^S)$.
\end{defish}

यावरून प्रतिस्थापन सिद्ध करण्यासाठी आधी लेव्ही (1960, प्रमेय 2 चा पहिला भाग) यांचे अनुसरण करून, एकाच वेळी दोन सूत्रांचे ``प्रतिबिंबन'' करता येते हे दाखवू:

\begin{lem}[$\Z + \text{दुर्बल प्रतिबिंबन}$ मध्ये.]\label{sth:replacement:refproofs:lem:reflect}
कोणतीही सूत्रे $\psi, \chi$ घ्या. मग असा संक्रमक संच $S$ असतो की $0$ आणि $1$ (आणि सूत्रांमधील कोणतीही प्राचले) हे $S$ चे सदस्य असतात, आणि $(\forall \overline{x} \in S)((\psi \liff \psi^S) \land (\chi
\liff \chi^S))$.
\end{lem}

\begin{proof}
$\phi$ हे $(z = 0 \land \psi) \lor (z = 1 \land \chi)$ सूत्र असू द्या.

येथे आपण संक्षिप्त लेखन वापरले आहे: ``$z = 0$'' चे ``$\forall t\, t \notin z$'' असे, आणि ``$z =1$'' चे ``$\forall s(s \in z
\liff \forall t\, t \notin s)$'' असे पूर्ण लेखन करायला हवे. पण $0, 1 \in S$ आणि $S$ संक्रमक असल्याने, ही सूत्रे $S$ साठी \emph{निरपेक्ष} आहेत; म्हणजे त्यांचे संख्यक $S$ पुरते मर्यादित केले तरी ती त्याच वस्तूला लागू होतात.\footnote{अधिक औपचारिकपणे, $\xi$ हे या दोन सूत्रांपैकी कोणतेही एक असेल, तर $\xi(z) \liff \xi^S(z)$.}

दुर्बल प्रतिबिंबनानुसार योग्य असा काही $S$ आपल्याला मिळतो की:
\begin{align*}
  (\forall z, \overline{x} \in S)(&\phi \liff \phi^S)\\
  \text{म्हणजे }(\forall z, \overline{x} \in S)(&((z = 0 \land \psi) \lor (z = 1 \land \chi)) \liff {}\\
  &\phantom{(}((z = 0 \land \psi) \lor (z = 1 \land \chi))^S)\\
  \text{म्हणजे }(\forall z, \overline{x} \in S)(&((z = 0 \land \psi) \lor (z = 1 \land \chi))\liff {}\\
  &\phantom{(}((z = 0 \land \psi^S) \lor (z = 1 \land \chi^S)))\\
  \text{म्हणजे }(\forall \overline{x} \in S)(&(\psi \liff \psi^S) \land (\chi \liff \chi^S))
\end{align*}
दुसऱ्या विधानावरून तिसरे मिळते; कारण ``$z = 0$'' आणि ``$z=1$'' ही $S$ साठी निरपेक्ष आहेत. $0 \neq 1$ असल्याने चौथे विधान मिळते.
\end{proof}\noindent आता लेव्ही (1960, प्रमेय 6) यांचे अनुसरण करून, आणि त्यांचा पुरावा थोडा सोपा करून, प्रतिस्थापन मिळवू:

\begin{thm}[$\Z$ + दुर्बल प्रतिबिंबन मध्ये]\label{thm:replacement}
कोणतेही सूत्र $\phi(v,w)$ आणि कोणताही $A$ घ्या. $(\forall x \in A)\lexists![y][\phi(x,y)]$ असेल, तर $\Setabs{y}{(\exists x \in A)\phi(x,y)}$ अस्तित्वात असतो.
\end{thm}

\begin{proof}
$(\forall x \in A)\lexists![y][\phi(x,y)]$ पूर्ण करणारा $A$ निश्चित करा आणि सूत्रे परिभाषित करा:
\begin{align*}
  \psi &\text{ म्हणजे } (\phi(x, z) \land A = A)\\
  \chi &\text{ म्हणजे } \lexists[y][\phi(x, y)]
\end{align*}
$A$ हे $\psi$ चे प्राचल असल्याने, \ref{sth:replacement:refproofs:lem:reflect} वापरून असा संक्रमक~$S$ मिळतो की $0, 1, A \in S$ (तसेच इतर कोणतीही प्राचले त्या संचात असतात) आणि:
\[
  (\forall x,z \in S)((\psi \liff \psi^S) \land (\chi \liff \chi^S))
\]
म्हणून विशेषतः:
\begin{align*}
  (\forall  x, z \in S)(&\phi(x, z) \liff \phi^S(x, z))\\
  (\forall x \in S)(&\exists y\phi(x, y) \liff (\exists y \in S)\phi^S(x, y))
\end{align*}
हे एकत्र करून, आणि $A \in S$ व $S$ संक्रमक असल्यामुळे $A \subseteq S$ हे लक्षात घेऊन:
\begin{align*}
  (\forall x \in A)(&\exists y\phi(x, y) \liff (\exists y \in S)\phi(x, y))
\end{align*}
आता $(\forall x \in A)(\lexists![y \in S])\phi(x, y)$; कारण $(\forall x \in A)\lexists![y][\phi(x, y)]$. म्हणून विभक्तीकरणावरून $\Setabs{y \in S}{(\exists x \in A) \phi(x, y)} = \Setabs{y}{(\exists
x \in A) \phi(x, y)}$ मिळतो.
\end{proof}








\section{परिशिष्ट: सांत स्वयंसिद्धकीकरणीयता}\label{sth:replacement:finiteaxiomatizability:sec}
प्रतिस्थापनापासून काही निष्कर्ष काढून हे प्रकरण संपवू. पहिला निष्कर्ष मॉन्टेग्यू (1961) यांचा आहे. लक्षात घ्या: हा $\ZF$ च्या \emph{अंतर्गत} पुरावा नाही, तर $\ZF$ \emph{विषयीचा} पुरावा आहे:
\begin{thm}\label{sth:replacement:finiteaxiomatizability:zfnotfinitely}
  $\ZF$ सुसंगत असेल, तर ती सांत स्वयंसिद्धकांनी मांडता येत नाही. अधिक सर्वसाधारणपणे: $\Th{T}$ सांत असेल आणि $\Th{T} \Proves \ZF$ असेल, तर $\Th{T}$ विसंगत असते.

  (येथे आपण गृहीत धरतो की विचाराधीन प्रथम-क्रम वाक्यांमध्ये $\in$ हेच एकमेव अतार्किक मूलचिन्ह आहे; तसेच प्रत्येक $\phi \in \ZF$ साठी $\Th{T} \Proves \phi$ आहे हे दाखवण्यासाठी $\Th{T} \Proves \ZF$ असे लिहितो.)
\end{thm}
\begin{proof}
  $\Th{T} \Proves \ZF$ पूर्ण करणारी सांत $\Th{T}$ निश्चित करा. म्हणून $\Th{T}$ प्रतिबिंबन, म्हणजे \ref{sth:replacement:ref:reflectionschema}, सिद्ध करते. $\Th{T}$ सांत असल्याने तिचे एका संधीकरणाच्या रूपात, $\theta$ असे, पुनर्लेखन करता येते. $\Th{T}$ अनंतता सिद्ध करत असल्याने, अनंततेचे स्वयंसिद्धकही $\theta$ मध्ये संधीघटक म्हणून घेऊ; त्यामुळे $\theta^X$ पूर्ण करणारा $X$ रिक्त नसतो. या सूत्रावर प्रतिबिंबन वापरल्यास $\Th{T} \Proves \exists \beta(\theta \liff \theta^{V_\beta})$. स्पष्टपणे $\Th{T} \Proves \theta$ असल्याने $\Th{T} \Proves \exists \beta\ \theta^{V_\beta}$ मिळते.

  आता $\psi(X)$ हे पुढीलचे संक्षिप्त लेखन असू द्या:
  \[
    \theta^X \land X\text{ संक्रमक आहे} \land (\forall Y \in X)(Y\text{ संक्रमक आहे}\lif \lnot \theta^{Y})
  \]
साधारणपणे याचा अर्थ: $X$ हे $\theta$ चे संक्रमक प्रतिमान आहे आणि या बाबतीत $\in$-लघुतम आहे. आता $\Th{T} \Proves \exists \beta\ \theta^{V_\beta}$ हे आठवा. $\ZF$ मधील, आणि त्यामुळे $\Th{T}$ मधील, प्रतांकांविषयीच्या मूलभूत तथ्यांवरून:
  \begin{equation}
    \Th{T} \Proves \exists M \psi(M). \tag{*}\label{Mpsi}
  \end{equation}
$\psi(X)$ चे पहिले संधीघटक वापरल्यास, $\Th{T} \Proves \sigma$ असेल तेव्हा $\Th{T} \Proves \forall X(\psi(X) \lif \sigma^X)$ मिळते. म्हणून \eqref{Mpsi} वरून:
  \begin{align*}
    \Th{T} &\Proves \forall X(\psi(X) \lif (\exists N \psi(N))^X)\\
  \intertext{हे आणि पुन्हा \eqref{Mpsi} वापरल्यास:}
    \Th{T} &\Proves \exists M(\psi(M) \land (\exists N \psi(N))^M)
  \intertext{म्हणून विशेषतः:}
    \Th{T} &\Proves \exists M(\psi(M) \land (\exists N \in M)((N\text{ संक्रमक आहे})^M \land (\theta^N)^M))
  \intertext{आता संक्रमणाविषयीच्या प्राथमिक युक्तिवादावरून:}
    \Th{T} &\Proves \exists M(\psi(M) \land (\exists N \in M)(N\text{ संक्रमक आहे} \land \theta^N))
  \end{align*}
म्हणून $\Th{T}$ विसंगत आहे.\footnote{या ``प्राथमिक युक्तिवादा''त संक्रमक संचांसाठी काही ``निरपेक्षतेविषयीची तथ्ये'' सिद्ध करावी लागतात.}
\end{proof}

पॉटर (2004, पृ.~223) यांनी नोंदवलेला पुढील निष्कर्ष यासारखाच आहे:

\begin{prop}\label{sth:replacement:finiteaxiomatizability:finiteextensionofZ}
  $\Th{T}$ ही $\Z$ मध्ये सांत नवीन स्वयंसिद्धके जोडून मिळालेली उपपत्ती असू द्या. $\Th{T} \Proves \ZF$ असेल, तर $\Th{T}$ विसंगत आहे. (येथे \ref{sth:replacement:finiteaxiomatizability:zfnotfinitely} प्रमाणेच गृहीत निर्बंध आहेत.)
\end{prop}
\begin{proof}
  विभक्तीकरणाची (अनंत अनेक) उदाहरणे \emph{वगळून} $\Th{T}$ च्या सर्व स्वयंसिद्धकांचे संधीकरण $\theta$ ने दर्शवा. या पुराव्यासाठी $\psi_*(X)$ हे पुढीलचे संक्षिप्त लेखन असू द्या:
\[
  \theta^X \land X\text{ संक्रमक व उपसंच-समावेशक आहे}
  \land (\forall Y \in X)(Y\text{ संक्रमक व उपसंच-समावेशक आहे}\lif\lnot\theta^Y).
\]
उपसंच-समावेशकत्वाची व्याख्या \ref{sth:spine:Valphabasic:sec} मध्ये दिली आहे. प्रतिबिंबनातून मिळणारे टप्पे या दोन्ही अटी पूर्ण करतात. म्हणून \ref{sth:replacement:finiteaxiomatizability:zfnotfinitely} च्या प्रतांक-युक्तिवादात या अटी वापरून $\Th{T}\Proves\exists M\psi_*(M)$ मिळते.

  \ref{sth:replacement:finiteaxiomatizability:zfnotfinitely} प्रमाणे, $\Th{T} \Proves \sigma$ असेल तेव्हा $\Th{T} \Proves \forall X(\psi_*(X) \lif \sigma^X)$ मिळते हा रूपबंध सिद्ध करता येतो. मग \ref{sth:replacement:finiteaxiomatizability:zfnotfinitely} प्रमाणेच पुरावा पूर्ण होतो; येथे संक्रमक व उपसंच-समावेशक $M$ मध्ये आतील $N$ साठी या दोन्ही अटींची निरपेक्षता वापरतो.

  तथापि, हा रूपबंध सिद्ध करण्यासाठी \ref{sth:replacement:finiteaxiomatizability:zfnotfinitely} पेक्षा थोडे अधिक काम लागते. कारण विभक्तीकरणाची उदाहरणे $\Th{T}$ मध्ये असली, तरी ती $\theta$ ची संधीघटके नाहीत. पण प्रत्येक विभक्तीकरण-उदाहरण $\sigma$ साठी $\Th{T} \Proves \forall X(X\text{ संक्रमक व उपसंच-समावेशक आहे} \lif \sigma^X)$ हे सिद्ध करून हा अडथळा दूर करता येतो. हे वाचकांसाठी सोडतो.
\end{proof}
\begin{prob}
  प्रत्येक विभक्तीकरण-उदाहरण $\sigma$ साठी $\Z \Proves \forall X(X\text{ संक्रमक व उपसंच-समावेशक आहे} \lif \sigma^X)$ आहे हे दाखवा. (हा रूपबंध आपण \ref{sth:replacement:finiteaxiomatizability:finiteextensionofZ} मध्ये वापरला.)
\end{prob}
\begin{prob}
  प्रत्येक $\phi \in \Z$ साठी $\ZF \Proves \phi^{V_{\omega+\omega}}$ आहे हे दाखवा.
\end{prob}
\begin{prob}
  \ref{sth:replacement:finiteaxiomatizability:zfnotfinitely} आणि \ref{sth:replacement:finiteaxiomatizability:finiteextensionofZ} यांच्या पुराव्यांत वापरलेले उर्वरित रूपबंधात्मक निष्कर्ष तपासा.
\end{prob}

\ref{sth:replacement:absinf:sec} मध्ये म्हटल्याप्रमाणे, यावरून प्रतिस्थापन हे \ref{sth:ordinals:ordtype:thmOrdinalRepresentation} पेक्षा काटेकोरपणे अधिक सामर्थ्यवान आहे असे दिसते. किंवा थोडे अधिक नेमके सांगायचे तर: $\Z$ + ``प्रत्येक सुक्रमण एका अद्वितीय क्रमसूचक संख्येशी समरूप आहे'' ही उपपत्ती सुसंगत असेल, तर प्रतिस्थापनाचे काही उदाहरण तिला सिद्ध करता येत नाही.














\chapter{क्रमसूचक अंकगणित}\label{sth:ord-arithmetic:chap}




\section{परिचय}\label{sth:ord-arithmetic:intro:sec}

\ref{sth:ordinals:chap} मध्ये आपण क्रमसूचक संख्यांची उपपत्ती विकसित केली. \ref{sth:spine:chap} मध्ये आपण पाहिले की पदानुक्रमाचा उर्वरित भाग ज्याभोवती बांधला जातो असा कणा म्हणून क्रमसूचक संख्यांकडे पाहता येते. पण क्रमसूचक संख्यांची तेवढीच भूमिका नाही. क्रमसूचक अंकगणित करण्याचे कामही आहे.

\ref{sth:ordinals:idea:sec} मध्ये $\omega$, $\omega+1$ आणि $\omega+\omega$ यांचा उल्लेख करताना आपण याची आधीच थोडी कल्पना दिली होती. तेव्हा आपण अनौपचारिकपणे बोललो; आता ते नीट स्पष्ट करण्याची वेळ आली आहे. मात्र या प्रकरणात तत्त्वज्ञान फारसे नाही, हे सांगायला हवे. येथे मुख्यतः तांत्रिक विकास आहे; त्याबरोबर एकच गोष्ट दोन वेगवेगळ्या मार्गांनी करता येते हे दाखवणारे (किंचित) रोचक निरीक्षण आहे.







\section{क्रमसूचक बेरीज}\label{sth:ord-arithmetic:add:sec}

समजा आपल्याला $\alpha$ आणि $\beta$ यांची बेरीज करायची आहे. $\alpha$ च्या एका {प्रतीनंतर} लगेच $\beta$ ची एक प्रत ठेवता येईल. (येथे \emph{प्रती} घ्याव्या लागतात; कारण \ref{sth:ordinals:basic:ordinalsaresubsets} वरून $\alpha \subseteq \beta$ किंवा $\beta \subseteq \alpha$ यांपैकी एक खरे असते.) याचा सहज अर्थ असा की आधी $\alpha$-लांबीचा पायऱ्यांचा अनुक्रम पार करायचा आणि \emph{त्यानंतर} $\beta$-लांबीचा अनुक्रम पार करायचा. मिळालेला अनुक्रम सुक्रमित असेल; म्हणून \ref{sth:ordinals:ordtype:thmOrdinalRepresentation} नुसार तो एका (अद्वितीय) क्रमसूचक संख्येशी समरूप असेल. त्या क्रमसूचक संख्येला $\alpha$ आणि $\beta$ यांची \emph{बेरीज} मानता येईल.

क्रमसूचक बेरीजेमागील सहज कल्पना ही आहे. तिची काटेकोर व्याख्या करण्यासाठी संचांच्या \emph{प्रती} घेण्याच्या कल्पनेपासून सुरुवात करू. कोणती वस्तू कुठून आली हे ओळखण्यासाठी $0$ आणि $1$ या मनमानी खूणा वापरायच्या:

\begin{defn}\label{sth:ord-arithmetic:add:defdissum}
$A$ आणि $B$ यांची \emph{विभक्त बेरीज} म्हणजे $A \disjointsum B = (A\times
\{0\}) \cup (B \times \{1\})$.
\end{defn}

यानंतर क्रमसूचक संख्यांच्या जोड्यांवर क्रम परिभाषित करू:

\begin{defn}
कोणत्याही क्रमसूचक संख्या $\alpha_1, \alpha_2, \beta_1, \beta_2$ यांच्यासाठी पुढीलप्रमाणे म्हणू:
\begin{align*}
  \tuple{\alpha_1, \alpha_2} \rlexless \tuple{\beta_1, \beta_2}\text{ तेव्हा आणि केवळ तेव्हाच }&
  \text{एकतर $\alpha_2 \in \beta_2$}\\
  & \text{किंवा $\alpha_2 = \beta_2$ आणि $\alpha_1 \in \beta_1$ या दोन्ही अटी}
\end{align*}
\end{defn}

हा \emph{उलट शब्दकोशीय} क्रम आहे; कारण आधी दुसऱ्या घटकानुसार आणि मग पहिल्या घटकानुसार क्रम लावला जातो. आता $\alpha$ ची प्रत आणि तिच्यानंतर $\beta$ ची प्रत घेतल्यावर मिळणारा क्रमप्रकार म्हणजे $\alpha \ordplus \beta$ अशी व्याख्या आपल्याला करायची होती, हे आठवा. त्यासाठी पुढील व्याख्या करू:

\begin{defn}\label{sth:ord-arithmetic:add:defordplus}
कोणत्याही क्रमसूचक संख्या $\alpha$, $\beta$ यांची बेरीज म्हणजे $\alpha \ordplus
\beta = \ordtype{\alpha \disjointsum \beta, \rlexless}$.
\end{defn}
\noindent
येथे आपण लेखनाचा किंचित सैल वापर केला आहे. काटेकोरपणे “$\rlexless$” ऐवजी “$\Setabs{\tuple{x,y}\in (\alpha\disjointsum\beta)\times(\alpha\disjointsum\beta)}{x \rlexless y}$” असे लिहायला हवे. तरी संक्षेपासाठी पुढेही हा सैल वापर करू.

पुढील निष्कर्ष आणि \ref{sth:ordinals:ordtype:thmOrdinalRepresentation} मिळून आपली व्याख्या नीट ठरते याची खात्री देतात:

\begin{lem}\label{sth:ord-arithmetic:add:ordsumlessiswo}
कोणत्याही क्रमसूचक संख्या $\alpha$ आणि $\beta$ यांच्यासाठी $\tuple{\alpha \disjointsum \beta, \rlexless}$ हे सुक्रमण आहे.
\end{lem}

\begin{proof}
$\alpha \disjointsum \beta$ वर $\rlexless$ हा संबंध संयोजित आहे, हे स्पष्ट आहे. तो सुस्थापित आहे हे दाखवण्यासाठी रिकामा नसलेला $X \subseteq \alpha
\disjointsum \beta$ निश्चित करा. $X$ मधील ज्या घटकांचा दुसरा निर्देशांक शक्य तितका लहान आहे त्यांचा उपसंच $Y$ असू द्या; म्हणजे $Y = \Setabs{\tuple{\gamma, i} \in X}{(\forall \tuple{\delta, j} \in X)i \leq j}$. आता $Y$ मधील सर्वांत लहान पहिला निर्देशांक असलेला घटक निवडा.
\end{proof}
\noindent
अशा रीतीने क्रमसूचक बेरीजची एक नीट आणि स्पष्ट व्याख्या मिळाली. आता एक अपेक्षित निष्कर्ष पाहू (\ref{sth:z:infinity-again:defnomega} नुसार $1 = \{0\}$ हे आठवा):
\begin{prop}
कोणत्याही क्रमसूचक संख्या $\alpha$ साठी $\alpha \ordplus  1 = \ordsucc{\alpha}$.
\end{prop}

\begin{proof}
$\ordsucc{\alpha} = \alpha \cup
\{\alpha\}$ पासून $\alpha\disjointsum1 = (\alpha \times \{0\})
\cup (\{0\} \times \{1\})$ कडे जाणारे समरूपण $f$ पाहा: $\gamma \in \alpha$ साठी $f(\gamma) =
\tuple{\gamma, 0}$, आणि $f(\alpha) = \tuple{0,
1}$.
\end{proof}
\noindent
शिवाय, बेरीज पुढील पुनरावर्ती अटी पूर्ण करते हे सहज दाखवता येते:

\begin{lem}\label{sth:ord-arithmetic:add:ordadditionrecursion}
कोणत्याही क्रमसूचक संख्या $\alpha, \beta$ यांच्यासाठी:
\begin{align*}
  \alpha\ordplus 0 &= \alpha\\
  \alpha \ordplus  (\beta\ordplus 1) &= (\alpha \ordplus  \beta) \ordplus  1\\
  \alpha  \ordplus  \beta &= \supstrict_{\delta < \beta}(\alpha \ordplus  \delta) && \text{जर $\beta $ सीमा क्रमसूचक संख्या असेल तर}
\end{align*}
\end{lem}

\begin{proof}
प्रत्येक प्रकरण तपासू. प्रथम:
\begin{align*}
  \alpha \ordplus  0
  & = \ordtype{(\alpha \times \{0\}) \cup (0 \times \{1\}), \rlexless} \\
  &= \ordtype{(\alpha \times \{0\}) \cup 0, \rlexless}\\
  &= \alpha\\
  \alpha \ordplus (\beta \ordplus  1)
  &= \ordtype{(\alpha\times \{0\}) \cup (\ordsucc{\beta}\times \{1\}), \rlexless} \\
  &= \ordtype{(\alpha\times \{0\}) \cup (\beta \times \{1\}), \rlexless} \ordplus 1\\
  &= (\alpha \ordplus  \beta) \ordplus  1
\end{align*}
आता $\beta \neq \emptyset$ ही सीमा क्रमसूचक संख्या असू द्या. $\delta < \beta$ असेल, तर $\delta\ordplus 1 < \beta$ सुद्धा; म्हणून $\alpha \ordplus  \delta$ हा $\alpha \ordplus  \beta$ चा उचित आरंभीचा खंड आहे. त्यामुळे $\alpha
\ordplus  \beta$ हा $X = \Setabs{\alpha
\ordplus  \delta}{\delta < \beta}$ चा काटेकोर ऊर्ध्वबंध आहे. शिवाय, $\alpha \leq \gamma <
\alpha \ordplus  \beta$ असेल, तर काही $\delta < \beta$ साठी $\gamma = \alpha \ordplus
\delta$ असते, हे स्पष्ट आहे. म्हणून $\alpha \ordplus  \beta =
\supstrict_{\delta< \beta}(\alpha\ordplus \delta)$.
\end{proof}

पण आता एक लक्षवेधी बाब आहे. क्रमसूचक बेरीज परिभाषित करण्यासाठी आपण \emph{त्याऐवजी} अनंतपार पुनरावर्तनाचे प्रमेय थेट वापरू शकलो असतो आणि \ref{sth:ord-arithmetic:add:ordadditionrecursion} मध्ये दिलेलीच पुनरावर्तन समीकरणे मांडू शकलो असतो (फक्त “$\beta \ordplus 1$” ऐवजी “$\ordsucc{\beta}$” वापरून).

म्हणून क्रमसूचक संख्यांवरील क्रिया परिभाषित करण्याचे दोन मार्ग आहेत. सुक्रमित संच स्पष्टपणे रचून त्याचा क्रमप्रकार घेता येतो; ही \emph{संश्लेषणात्मक} पद्धत. किंवा केवळ पुनरावर्तन समीकरणे देऊन \emph{पुनरावर्ती} पद्धतीने व्याख्या करता येते. योग्य रीतीने केल्यास दोन्हींचा निष्कर्ष मात्र एकच असतो. कारण \ref{sth:ordinals:ordtype:thmOrdinalRepresentation} नुसार या पुनरावर्तन समीकरणांनी \emph{अद्वितीय} क्रमसूचक संख्या ठरतात.

अनेक बाबतींत क्रमसूचक अंकगणित नैसर्गिक संख्यांच्या बेरीजेसारखेच वागते. उदाहरणार्थ, पुढील निष्कर्ष सिद्ध करता येतो:

\begin{lem}\label{sth:ord-arithmetic:add:ordinaladditionisnice}
$\alpha, \beta, \gamma$ या क्रमसूचक संख्या असतील, तर:
\begin{enumerate}
  \item\label{sth:ord-arithmetic:add:ordaddition1} $\beta < \gamma$ असेल, तर $\alpha
  \ordplus  \beta < \alpha \ordplus  \gamma$
  \item\label{sth:ord-arithmetic:add:ordaddition2} $\alpha \ordplus  \beta =
  \alpha\ordplus \gamma$ असेल, तर $\beta = \gamma$
  \item\label{sth:ord-arithmetic:add:ordaddition3} $\alpha \ordplus  (\beta \ordplus
  \gamma) = (\alpha \ordplus  \beta) \ordplus  \gamma$; म्हणजे बेरीज सहचारी आहे
  \item\label{sth:ord-arithmetic:add:ordaddition4} $\alpha \leq \beta$ असेल, तर $\alpha
  \ordplus  \gamma \leq \beta \ordplus \gamma$
\end{enumerate}
\end{lem}

\begin{proof}
इतर विधाने सरावासाठी ठेवून \ref{sth:ord-arithmetic:add:ordaddition3} सिद्ध करू. \ref{sth:ord-arithmetic:add:ordadditionrecursion} वापरून $\gamma$ वर साध्या अनंतपार विगमनाने पुरावा देऊ. $\gamma = 0$ असेल, तर:
\[
(\alpha \ordplus  \beta) \ordplus  0 = \alpha \ordplus  \beta  = \alpha \ordplus  (\beta \ordplus  0)
\]
$\gamma = \delta\ordplus 1$ असेल, तर विगमनगृहीतक म्हणून $(\alpha
\ordplus  \beta) \ordplus  \delta = \alpha \ordplus  (\beta \ordplus
\delta)$ मानू. आता \ref{sth:ord-arithmetic:add:ordadditionrecursion} तीनदा वापरल्यास:
\begin{align*}
  (\alpha \ordplus  \beta) \ordplus  (\delta \ordplus  1) & = ((\alpha \ordplus  \beta) \ordplus  \delta)\ordplus 1\\
  & = (\alpha \ordplus  (\beta \ordplus  \delta)) \ordplus  1\\
  & = \alpha \ordplus  ((\beta \ordplus  \delta)\ordplus 1)\\
  & = \alpha \ordplus  (\beta \ordplus  (\delta\ordplus 1))
\end{align*}
$\gamma$ ही सीमा क्रमसूचक संख्या असेल, तर प्रत्येक $\delta \in \gamma$ साठी $(\alpha \ordplus  \beta) \ordplus  \delta =
\alpha \ordplus  (\beta \ordplus  \delta)$ असे विगमनगृहीतक मानू. मग:
\begin{align*}
  (\alpha \ordplus  \beta) \ordplus  \gamma & = \supstrict_{\delta < \gamma}((\alpha \ordplus  \beta) \ordplus  \delta) \\
  &= \supstrict_{\delta < \gamma}(\alpha \ordplus  (\beta \ordplus  \delta))\\
  &= \alpha \ordplus  \supstrict_{\delta < \gamma}(\beta \ordplus  \delta)\\
  & = \alpha \ordplus  (\beta \ordplus  \gamma)
\end{align*}
\end{proof}

\begin{prob}
\ref{sth:ord-arithmetic:add:ordinaladditionisnice} मधील उर्वरित विधाने सिद्ध करा.
\end{prob}

या बाबतींत क्रमसूचक बेरीज परिचित वाटावी. पण एका अत्यंत महत्त्वाच्या बाबतीत ती नैसर्गिक संख्यांच्या बेरीजेसारखी \emph{नाही}.

\begin{prop}\label{sth:ord-arithmetic:add:ordsumnotcommute}
क्रमसूचक बेरीज {क्रमनिरपेक्ष नाही}; $1 \ordplus  \omega = \omega <
\omega \ordplus  1$.
\end{prop}

\begin{proof}
$1 \ordplus  \omega = \supstrict_{n < \omega} (1 \ordplus n)
= \omega \in \omega \cup \{\omega\} = \ordsucc{\omega} = \omega
\ordplus  1$ हे लक्षात घ्या.
\end{proof}
\noindent
सुरुवातीला आश्चर्य वाटू शकते; पण \emph{वाटायला नको}. एका बाजूला, $1 \ordplus  \omega$ पाहताना आपण सर्व नैसर्गिक संख्यांच्या \emph{आधी} एक जादा घटक ठेवून मिळणाऱ्या क्रमप्रकाराचा विचार करतो. \ref{sfr:infinite:hilbert:sec} मधील हिल्बर्टच्या हॉटेलप्रमाणे विचार केल्यास, सहज अर्थाने या जादा पहिल्या घटकामुळे एकूण क्रमप्रकारात फरक पडू नये. दुसऱ्या बाजूला, $\omega \ordplus  1$ पाहताना आपण सर्व नैसर्गिक संख्यांच्या \emph{नंतर} एक जादा घटक ठेवून मिळणाऱ्या क्रमप्रकाराचा विचार करतो. आणि तो अगदीच वेगळा प्रकार आहे!







\section{क्रमसूचक बेरीजचा उपयोग}\label{sth:ord-arithmetic:using-addition:sec}

क्रमसूचक संख्यांची बेरीज वापरून \ref{sth:spine:rank:defnsetrank} च्या अर्थाने विविध संचांचे प्रतांक स्पष्टपणे मोजता येतात:

\begin{lem}\label{sth:ord-arithmetic:using-addition:rankcomputation}
$\setrank{A} = \alpha$ आणि $\setrank{B} = \beta$ असतील, तर:
\begin{enumerate}
  \item\label{sth:ord-arithmetic:using-addition:exrankpow} $\setrank{\Pow{A}} = \alpha\ordplus 1$
  \item\label{sth:ord-arithmetic:using-addition:exrankpair} $\setrank{\{A, B\}} = \max(\alpha,
  \beta) \ordplus 1$
  \item\label{sth:ord-arithmetic:using-addition:exrankcup} $\setrank{A \cup B} = \max(\alpha,
  \beta)$
  \item\label{sth:ord-arithmetic:using-addition:exranktuple} $\setrank{\tuple{A,B}} = \max(\alpha,
  \beta) \ordplus  2$
  \item\label{sth:ord-arithmetic:using-addition:exranktimes} $\setrank{A \times B} \leq \max(\alpha,
  \beta) \ordplus  2$
  \item\label{sth:ord-arithmetic:using-addition:exrankunion} $\alpha$ रिकामी किंवा सीमा क्रमसूचक संख्या असेल, तर $\setrank{\bigcup A} = \alpha$; $\alpha = \gamma\ordplus 1$ असेल, तर $\setrank{\bigcup A} = \gamma$
\end{enumerate}
\end{lem}

\begin{proof}
संपूर्ण पुराव्यात आपण \ref{sth:spine:rank:ranksupstrict} वारंवार वापरू.

\emph{\ref{sth:ord-arithmetic:using-addition:exrankpow}.} $x \subseteq A$ असेल, तर $\setrank{x} \leq
\setrank{A}$. म्हणून $\setrank{\Pow{A}} \leq \alpha \ordplus  1$. विशेषतः $A
\in \Pow{A}$ असल्याने $\setrank{\Pow{A}} = \alpha \ordplus  1$.

\emph{\ref{sth:ord-arithmetic:using-addition:exrankpair}.} \ref{sth:spine:rank:ranksupstrict} नुसार.

\emph{\ref{sth:ord-arithmetic:using-addition:exrankcup}.} \ref{sth:spine:rank:ranksupstrict} नुसार.

\emph{\ref{sth:ord-arithmetic:using-addition:exranktuple}.} \ref{sth:ord-arithmetic:using-addition:exrankpair} दोनदा वापरल्यास.

\emph{\ref{sth:ord-arithmetic:using-addition:exranktimes}.} $A \times B \subseteq
\Pow{\Pow{A \cup B}}$ हे लक्षात घेऊन \ref{sth:ord-arithmetic:using-addition:exranktuple} वापरा.

\emph{\ref{sth:ord-arithmetic:using-addition:exrankunion}.} $\alpha = \gamma\ordplus 1$ असेल, तर काही $c \in A$ साठी $\setrank{c} = \gamma$ असते आणि $A$ च्या कोणत्याही घटक चा प्रतांक याहून मोठा नसतो; म्हणून $\setrank{\bigcup A} = \gamma$. $\alpha$ ही सीमा क्रमसूचक संख्या असेल, तर $A$ मध्ये $\alpha$ च्या मनमान्या जवळ (पण काटेकोरपणे खाली) प्रतांक असलेले घटक असतात. त्यामुळे $\bigcup A$ मध्येही $\alpha$ च्या मनमान्या जवळ (पण काटेकोरपणे खाली) प्रतांक असलेले घटक असतात, आणि म्हणून $\setrank{\bigcup A} = \alpha$.
\end{proof}
\noindent
\ref{sth:ord-arithmetic:using-addition:exranktimes} मध्ये \emph{अ}समानता का येते, हे दाखवण्याचे काम सरावासाठी सोडतो.

\begin{prob}
$\setrank{A \times B}=
\max(\setrank{A}, \setrank{B})$ असेल असे संच $A$ आणि $B$ द्या. तसेच $\setrank{A \times B}= \max(\setrank{A}, \setrank{B}) \ordplus  2$ असेल असे संच $A$ आणि $B$ द्या. इतर कोणते प्रतांक शक्य आहेत का?
\end{prob}

क्रमसूचक संख्या “सांत” किंवा “अनंत” असण्याचा अर्थ लावण्याचे अनेक रास्त मार्ग समतुल्य आहेत, हे दाखवण्याच्या स्थितीतही आता आपण आहोत:

\begin{lem}\label{sth:ord-arithmetic:using-addition:ordinfinitycharacter}
कोणत्याही क्रमसूचक संख्या $\alpha$ साठी पुढील विधाने समतुल्य आहेत:
\begin{enumerate}
  \item\label{sth:ord-arithmetic:using-addition:ord:notinomega} $\alpha\notin \omega$; म्हणजे $\alpha$ नैसर्गिक संख्या नाही
  \item\label{sth:ord-arithmetic:using-addition:ord:omegaplus} $\omega \leq \alpha$
  \item\label{sth:ord-arithmetic:using-addition:ord:oneplus} $1 \ordplus  \alpha = \alpha$

  \item\label{sth:ord-arithmetic:using-addition:ord:plusone} $\alpha \approx \alpha\ordplus 1$; म्हणजे $\alpha$ आणि $\alpha\ordplus 1$ यांचे संख्याबल समान आहे
  \item\label{sth:ord-arithmetic:using-addition:ord:infinite} $\alpha$ डेडेकिंड-अनंत आहे
  \end{enumerate}
\end{lem}
\noindent
म्हणून क्रमसूचक संख्या सांत किंवा अनंत असण्याचा अर्थ समजण्याचे पाच सिद्ध करता येण्याजोगे समतुल्य मार्ग आपल्याकडे आहेत.

\begin{proof}
\emph{\ref{sth:ord-arithmetic:using-addition:ord:notinomega} $\Rightarrow$ \ref{sth:ord-arithmetic:using-addition:ord:omegaplus}.} त्रिपर्यायानुसार.

\emph{\ref{sth:ord-arithmetic:using-addition:ord:omegaplus} $\Rightarrow$ \ref{sth:ord-arithmetic:using-addition:ord:oneplus}.} $\alpha \geq \omega$ निश्चित करा. अनंतपार विगमनानुसार अशी लघुतम क्रमसूचक संख्या $\gamma$ असते (ती $0$ असू शकते) की काही सीमा क्रमसूचक संख्या $\beta$ साठी $\alpha = \beta \ordplus \gamma$. मग:
\[
  1 \ordplus \alpha =
  1 \ordplus (\beta \ordplus \gamma) =
  (1 \ordplus \beta) \ordplus \gamma =
  \supstrict_{\delta < \beta} (1 \ordplus  \delta) \ordplus  \gamma =
  \beta \ordplus  \gamma =
  \alpha.
\]
\emph{\ref{sth:ord-arithmetic:using-addition:ord:oneplus} $\Rightarrow$ \ref{sth:ord-arithmetic:using-addition:ord:plusone}.} $f \colon (\alpha \disjointsum 1) \to (1
\disjointsum \alpha)$ हे एकास-एक आच्छादन सहज मिळते. $1 \ordplus \alpha = \alpha$ असेल, तर $g \colon (1 \disjointsum \alpha) \to \alpha$ हे समरूपण असते. आता $\comp{f}{g}$ पाहा.

\emph{\ref{sth:ord-arithmetic:using-addition:ord:plusone} $\Rightarrow$ \ref{sth:ord-arithmetic:using-addition:ord:infinite}.} $\alpha \approx \alpha \ordplus 1$ असेल, तर $f \colon
(\alpha \disjointsum 1) \to \alpha$ हे एकास-एक आच्छादन असते. प्रत्येक $\gamma < \alpha$ साठी $g(\gamma) = f(\gamma, 0)$ अशी व्याख्या करा. हे एकास-एक फलन $\alpha$ डेडेकिंड-अनंत आहे याची साक्ष देते; कारण $f(0,1) \in \alpha \setminus \ran{g}$.

\emph{\ref{sth:ord-arithmetic:using-addition:ord:infinite} $\Rightarrow$ \ref{sth:ord-arithmetic:using-addition:ord:notinomega}.} हा \ref{sth:z:infinity-again:naturalnumbersarentinfinite} आहे.
\end{proof}








\section{क्रमसूचक गुणाकार}\label{sth:ord-arithmetic:mult:sec}

आता क्रमसूचक गुणाकाराकडे वळू; त्यासाठीची पद्धत क्रमसूचक बेरीजेसारखीच आहे. समजा $\alpha$ ला~$\beta$ ने गुणायचे आहे. रुंदी $\alpha$ आणि उंची $\beta$ असलेली आयताकृती जाळी कल्पा. प्रत्येक ओळीतून पुढे जावे, मग पुढील ओळीतून\ldots आणि अखेरीस संपूर्ण जाळीतून जावे; त्यातून $\alpha$ आणि $\beta$ यांचा गुणाकार मिळतो. दुसऱ्या शब्दांत, $\beta$ मधील \emph{प्रत्येक} घटकाऐवजी $\alpha$ ची एक प्रत ठेवून $\alpha$ आणि $\beta$ यांचा गुणाकार तयार होतो.

हे औपचारिक करण्यासाठी $\alpha$ आणि $\beta$ यांच्या कार्तीय गुणाकारावर उलट शब्दकोशीय क्रम वापरू:

\begin{defn}
कोणत्याही क्रमसूचक संख्या $\alpha, \beta$ यांचा गुणाकार म्हणजे $\alpha \ordtimes \beta = \ordtype{\alpha \times \beta, \rlexless}$.
\end{defn}
\noindent
ही व्याख्या नीट ठरते याची पुन्हा खात्री करावी लागेल:

\begin{lem}\label{sth:ord-arithmetic:mult:ordtimeslessiswo}
कोणत्याही क्रमसूचक संख्या $\alpha$ आणि $\beta$ यांच्यासाठी $\tuple{\alpha \times \beta, \rlexless}$ हे सुक्रमण आहे.
\end{lem}

\begin{proof}
\ref{sth:ord-arithmetic:add:ordsumlessiswo} प्रमाणेच.
\end{proof}
\noindent
गुणाकार पुढीलप्रमाणे वागतो, हे सिद्ध करणेही कठीण नाही:

\begin{lem}\label{sth:ord-arithmetic:mult:ordtimesrecursion}
कोणत्याही क्रमसूचक संख्या $\alpha, \beta$ यांच्यासाठी:
\begin{align*}
  \alpha \ordtimes 0 &= 0\\
  \alpha \ordtimes (\beta \ordplus 1) &=
    (\alpha \ordtimes \beta) \ordplus \alpha\\
  \alpha  \ordtimes \beta &=
    \bigcup_{\delta < \beta}(\alpha \ordtimes \delta) &&
    \text{$\beta$ ही सीमा क्रमसूचक संख्या असेल तेव्हा}.
\end{align*}
\end{lem}

\begin{proof}
सरावासाठी सोडले आहे.
\end{proof}

खरे तर बेरीजप्रमाणेच, प्रत्यक्ष व्याख्येऐवजी या पुनरावर्तन समीकरणांवरून क्रमसूचक गुणाकार परिभाषित करता आला असता. त्याचप्रमाणे, काही गुणधर्म परिचितही वाटतात:

\begin{lem}\label{sth:ord-arithmetic:mult:ordinalmultiplicationisnice}
$\alpha, \beta, \gamma$ या क्रमसूचक संख्या असतील, तर:
\begin{enumerate}
  \item\label{sth:ord-arithmetic:mult:ordtimes1} $\alpha \neq 0$ आणि $\beta < \gamma$ असतील,
  तर $\alpha \ordtimes \beta < \alpha \ordtimes \gamma$;
  \item\label{sth:ord-arithmetic:mult:ordtimes2} $\alpha \neq 0$ आणि $\alpha \ordtimes
  \beta = \alpha\ordtimes\gamma$ असतील, तर $\beta = \gamma$;
  \item\label{sth:ord-arithmetic:mult:ordtimes3} $\alpha \ordtimes (\beta \ordtimes
  \gamma) = (\alpha \ordtimes \beta) \ordtimes \gamma$;
  \item\label{sth:ord-arithmetic:mult:ordtimes4} $\alpha \leq \beta$ असेल, तर $\alpha
  \ordtimes \gamma \leq \beta \ordtimes \gamma$;
  \item\label{sth:ord-arithmetic:mult:ordtimes5} $\alpha \ordtimes (\beta \ordplus
  \gamma) = (\alpha \ordtimes \beta )\ordplus (\alpha\ordtimes
  \gamma)$.
\end{enumerate}
\end{lem}

\begin{proof}
सरावासाठी सोडले आहे.
\end{proof}

मनसोक्त इतर निष्कर्ष सिद्ध करा (किंवा शोधून वाचा). पण \ref{sth:ord-arithmetic:add:ordsumnotcommute} लक्षात घेतल्यास पुढील गोष्ट आश्चर्यकारक वाटू नये:

\begin{prop}
क्रमसूचक गुणाकार क्रमनिरपेक्ष नाही: $2 \ordtimes \omega  =
\omega < \omega \ordtimes 2$
\end{prop}

\begin{proof}
$2 \ordtimes \omega = \supstrict_{n < \omega}(2\ordtimes  n) = \omega \in \supstrict_{n < \omega}(\omega \ordplus n) = \omega \ordplus \omega = \omega \ordtimes 2$.
\end{proof}
\noindent
यामागील सहज कारणही सरळ आहे. $2
\ordtimes \omega$ मोजताना प्रत्येक नैसर्गिक संख्येऐवजी दोन घटक ठेवता. सर्व “अर्ध्या” संख्या नैसर्गिक संख्यांमध्ये घातल्या, म्हणजे $\Setabs{\nicefrac{n}{2}}{n \in \omega}$ वरील नेहमीच्या क्रमाचा विचार केला, तरी तोच क्रमप्रकार मिळेल. या रूपात पाहिल्यास तो क्रमप्रकार स्वतः $\omega$ च्याच क्रमप्रकारासारखा आहे, हे स्पष्ट होते. पण $\omega \ordtimes 2$ मोजताना $\omega$ च्या दोन प्रती एकीमागोमाग ठेवता.

\begin{prob}
\ref{sth:ord-arithmetic:mult:ordtimeslessiswo},
\ref{sth:ord-arithmetic:mult:ordtimesrecursion},
आणि
\ref{sth:ord-arithmetic:mult:ordinalmultiplicationisnice}
सिद्ध करा.
\end{prob}








\section{क्रमसूचक घातांकन}\label{sth:ord-arithmetic:expo:sec}

आता क्रमसूचक घातांकनाकडे वळू. दुर्दैवाने, क्रमसूचक घातांकनाची \emph{सोपी} संश्लेषणात्मक व्याख्या नाही.

अर्थात स्पष्ट संश्लेषणात्मक व्याख्या \emph{आहेत}. त्यांपैकी एक पाहू. $\text{finfun}(\alpha,\beta)$ हा अशा सर्व फलन $f
\colon \alpha \to \beta$ यांचा संच असू द्या की $\Setabs{\gamma \in
\alpha}{f(\gamma) \neq 0}$ हा संच एखाद्या नैसर्गिक संख्येइतका संख्याबल असलेला आहे. $\text{finfun}(\alpha,\beta)$ वर पुढीलप्रमाणे सुक्रमण परिभाषित करा: $f
\sqsubset g$ तेव्हा आणि केवळ तेव्हाच $f \neq g$ आणि $f(\gamma_0) < g(\gamma_0)$; येथे $\gamma_0 = \text{max}\Setabs{\gamma \in \alpha}{f(\gamma) \neq
g(\gamma)}$. शून्येतर आधार $\alpha$ साठी $\ordexpo{\alpha}{\beta}$ ची व्याख्या $\ordtype{\text{finfun}(\beta, \alpha), \sqsubset}$ अशी करता येते. पॉटर यांनी ही स्पष्ट व्याख्या वापरली आहे आणि लगेच असे स्पष्ट केले:
\begin{quote}
  हा क्रम निवडण्यामागे एकमेव कारण म्हणजे योग्य पुनरावर्ती अट पूर्ण करणारी क्रमसूचक घातांकनाची व्याख्या मिळावी ही आपली इच्छा\ldots; आणि क्रमित बेरीज किंवा क्रमित गुणाकार यांपेक्षा या क्रमाची कल्पना करणे खूप कठीण आहे.
  (पॉटर 2004, पृ.~199)
\end{quote}
अगदी तसेच. बेरीज समजावताना आपण “$\alpha$ ची प्रत आणि तिच्यानंतर $\beta$ ची प्रत” असे म्हटले; गुणाकार समजावताना “$\alpha$ च्या प्रतींचा $\beta$-लांबीचा अनुक्रम” असे म्हटले. पण $\text{finfun}(\alpha, \gamma)$ विषयी इतके थोडक्यात सांगण्यासारखे आपल्याकडे काही नाही. म्हणून त्याऐवजी क्रमसूचक घातांकनाची व्याख्या थेट अनंतपार पुनरावर्तनाने \emph{देऊ}; म्हणजे:

\begin{defn}\label{sth:ord-arithmetic:expo:ordexporecursion}
शून्येतर क्रमसूचक आधार $\alpha$ साठी:
\begin{align*}
  \ordexpo{\alpha}{0} &= 1\\
  \ordexpo{\alpha}{\beta\ordplus 1} &=\ordexpo{\alpha}{\beta} \ordtimes \alpha\\
  \ordexpo{\alpha}{\beta} &= \bigcup_{\delta < \beta}\ordexpo{\alpha}{\delta}& & \text{$\beta$ ही सीमा क्रमसूचक संख्या असेल तेव्हा}
\end{align*}
\end{defn}

आपण संच उपपत्तीचे \emph{संशोधक म्हणून} काम करत असतो, तर क्रमसूचक घातांकनाचे आणखी काही गुणधर्म तपासण्याची इच्छा झाली असती. पण त्याविषयी येथे अधिक काही सांगायचे नाही; फक्त क्रमसूचक घातांकन क्रमनिरपेक्ष नाही, ही अपेक्षित बाब नोंदवू. कारण $\ordexpo{2}{\omega} =
\bigcup_{\delta < \omega}\ordexpo{2}{\delta} = \omega$, तर $\ordexpo{\omega}{2} = \omega \ordtimes \omega$. पण घातांकन क्रमनिरपेक्ष असावे अशी \emph{अपेक्षाही} ठेवू नये; नैसर्गिक संख्यांच्या बाबतीतही तसे नाही: $\ordexpo{2}{3} = 8 < 9 = \ordexpo{3}{2}$.

\begin{prob}
अनंतपार विगमन वापरून सिद्ध करा की, आधार शून्येतर असेल आणि $\ordexpo{\alpha}{\beta} = \ordtype{\text{finfun}(\beta, \alpha),
\sqsubset}$ अशी व्याख्या केली, तर \ref{sth:ord-arithmetic:expo:ordexporecursion} मधील पुनरावर्तन समीकरणे मिळतात.
\end{prob}



\chapter{संचांक}\label{sth:cardinals:chap}




\section{कँटरचे तत्त्व}\label{sth:cardinals:cp:sec}

मागे \ref{sth:ordinals:vn:sec} कडे मन वळवा. आपण सुक्रमित संचांविषयी चर्चा करत होतो आणि सुक्रमणांचे प्रतिनिधित्व करणाऱ्या वस्तू असाव्यात, असे सुचवले होते. या उद्देशाने आपण क्रमसूचक संख्या मांडल्या; नंतर \ref{sth:ordinals:ordtype:ordtypesworklikeyouwant} मध्ये दाखवले की त्यांचे गुणधर्म अपेक्षेप्रमाणे आहेत, म्हणजे:
\[
  \ordtype{A, <} = \ordtype{B, \lessdot}
  \text{ तेव्हा आणि केवळ तेव्हाच } \tuple{A, <} \isomorphic \tuple{B, \lessdot}.
\]
आता आणखी मागे, \ref{sfr:siz:equ:sec} कडे वळा. तेथे संचाच्या “आकारमाना”ची कल्पना आपण सहजरीत्या मांडली होती. विशेषतः, दोन संच तुल्यबल आहेत, म्हणजे $\cardeq{A}{B}$, असे आपण तेव्हा आणि केवळ तेव्हाच म्हणालो जेव्हा $f \colon A \to
B$ हे एकास-एक आच्छादन असते. सुक्रमणाच्या कल्पनेपेक्षा ही स्वभावतः सोपी कल्पना आहे: आपल्याला फक्त एकास-एक आच्छादन महत्त्वाची आहेत; क्रम-समरूपणांच्या बाबतीत जसे पाहतो तसे, ती एकास-एक आच्छादन कोणतीही रचना जतन करतात का, हे येथे पाहायचे नाही.

यावरून एक स्वाभाविक विचार सुचतो. सुक्रमणांचे मोजमाप करण्यासाठी जशा काही वस्तू, म्हणजे \emph{क्रमसूचक संख्या}, आपण मांडल्या, तशाच आकारमानाचे मोजमाप करण्यासाठी \emph{संचांक} मांडू शकतो. हेच या प्रकरणाचे उद्दिष्ट आहे.

हे संचांक नेमके काय असतील ते सांगण्याआधी, त्यांनी पूर्ण करायला हवे असे एक तत्त्व मांडू. संच $X$ चा संचांक $\card{X}$ असा लिहू; मग पुढील अट अपेक्षित आहे:
\[
  \card{A} = \card{B} \text{ तेव्हा आणि केवळ तेव्हाच } \cardeq{A}{B}.
\]
याला आपण \emph{कँटरचे} तत्त्व म्हणू, कारण ही कल्पना कँटर यांच्या मनात बहुधा सर्वप्रथम अगदी स्पष्टपणे होती. (तिचा \emph{ह्यूमच्या} तत्त्वाशी संबंध काय आहे, याविषयी \ref{sth:cardinals:hp:sec} मध्ये आणखी सांगू.) म्हणून प्रत्येक $X$ साठी $\card{X}$ ची अशी व्याख्या करणे हे आपले उद्दिष्ट आहे की कँटरचे तत्त्व सिद्ध होईल.







\section{क्रमसूचक संख्या म्हणून संचांक}\label{sth:cardinals:cardsasords:sec}

प्रत्यक्षात, संचांकांची आपली उपपत्ती क्रमसूचक संख्यांच्या उपपत्तीचाच (बिनदिक्कत) आधार घेईल. म्हणजे काही विशिष्ट क्रमसूचक संख्यांनाच आपण संचांक म्हणून परिभाषित करू. नेमकी व्याख्या अशी:

\begin{defn}\label{sth:cardinals:cardsasords:defcardinalasordinal}
$A$ ला सुक्रमित करता येत असेल, तर $\cardeq{A}{\gamma}$ असेल अशी लघुतम क्रमसूचक संख्या $\gamma$ म्हणजे $\card{A}$. कोणत्याही क्रमसूचक संख्या $\gamma$ साठी, $\gamma$ ला $\gamma = \card{\gamma}$ असेल तेव्हा आणि केवळ तेव्हाच \emph{संचांक} म्हणू.
\end{defn}

आत्ताच आपण “$A$ ला सुक्रमित करता येते” असा शब्दप्रयोग केला. गणितात बहुतेक वेळा घडते तसे, येथे शक्यतादर्शक भाषा हा अस्तित्वाच्या दाव्याचा अनौपचारिक पर्याय आहे: “$A$ ला सुक्रमित करता येते” याचा अर्थ फक्त “$A$ ला सुक्रमित करणारा संबंध अस्तित्वात आहे” असा आहे.

पण \ref{sth:cardinals:cardsasords:defcardinalasordinal} मध्ये एक अडचण आहे. \emph{प्रत्येक} संचाला आकारमान असावे, म्हणजे प्रत्येक~$A$ साठी $\card{A}$ अस्तित्वात असावा, अशी आपली इच्छा आहे. मात्र नुकतीच दिलेली व्याख्या अटीने सुरू होते: “$A$ ला सुक्रमित करता येत \emph{असेल}, तर\ldots”. एखादा संच $A$ सुक्रमित करता येत नसेल, तर ही व्याख्या $\card{A}$ ही वस्तू निश्चितच करणार नाही.

म्हणून \ref{sth:cardinals:cardsasords:defcardinalasordinal} वापरायची असेल, तर प्रत्येक संच सुक्रमित करता येतो याची खात्री हवी. दुर्दैवाने, ही खात्री~$\ZF$ मध्ये उपलब्ध नाही. त्यामुळे \ref{sth:cardinals:cardsasords:defcardinalasordinal} वापरायची असल्यास पुढीलसारखे नवे स्वयंसिद्धक जोडावे लागेल:
\begin{axiom}[सुक्रमण]
प्रत्येक संच सुक्रमित करता येतो.
\end{axiom}
सुक्रमणाचे स्वयंसिद्धक स्वीकारण्याजोगे आहे का, याची चर्चा \ref{sth:choice:chap} मध्ये करू. पण आतापासून आपण ते गृहीत धरू. ते वापरल्यावर \ref{sth:cardinals:cardsasords:defcardinalasordinal} प्रमाणे परिभाषित केलेले संचांक अस्तित्वात असतात आणि अपेक्षित गुणधर्म दाखवतात, हे सहज सिद्ध होते:

\begin{lem}\label{sth:cardinals:cardsasords:lem:CardinalsExist}
प्रत्येक संच $A$ साठी:
\begin{enumerate}
  \item\label{sth:cardinals:cardsasords:cardaexists} $\card{A}$ अस्तित्वात असतो आणि एकमेव असतो;
  \item\label{sth:cardinals:cardsasords:cardaapprox} $\cardeq{\card{A}}{A}$;
  \item\label{sth:cardinals:cardsasords:cardaidem} $\card{A}$ हा संचांक असतो, म्हणजे
  $\card{A} = \card{\card{A}}$;
\end{enumerate}
\end{lem}

\begin{proof}
$A$ निश्चित करा. सुक्रमणाच्या स्वयंसिद्धकानुसार $\tuple{A, R}$ हे सुक्रमण आहे. \ref{sth:ordinals:ordtype:thmOrdinalRepresentation} नुसार $\tuple{A,
R}$ हे एकमेव क्रमसूचक संख्या $\beta$ शी क्रम-समरूपी आहे. म्हणून $\cardeq{A}{\beta}$. अनंतपार विगमनानुसार $\cardeq{A}{\gamma}$ असेल अशी एकमेव लघुतम क्रमसूचक संख्या $\gamma$ आहे. म्हणून $\card{A}
= \gamma$; यावरून \ref{sth:cardinals:cardsasords:cardaexists} आणि \ref{sth:cardinals:cardsasords:cardaapprox} सिद्ध होतात.
\ref{sth:cardinals:cardsasords:cardaidem} साठी लक्षात घ्या की $\delta \in \gamma$ असेल, तर $\gamma$ लघुतम निवडल्यामुळे
$\cardless{\delta}{A}$; आणि तुल्यबलता हा समतुल्य संबंध असल्याने (\ref{sfr:siz:equ:equinumerosityisequi}) $\cardless{\delta}{\gamma}$ सुद्धा. म्हणून $\gamma =
\card{\gamma}$.

\end{proof}

पुढील निष्कर्ष कँटरच्या तत्त्वाची, आणि त्याहून अधिक गोष्टींची, खात्री देतो. (संचांकांना त्यांचा क्रम क्रमसूचक संख्यांकडूनच मिळतो; म्हणजे $\cardfont{a} < \cardfont{b}$ तेव्हा आणि केवळ तेव्हाच जेव्हा $\cardfont{a} \in \cardfont{b}$. संचांक असल्याचे माहीत असलेल्या वस्तूंसाठी आपण फ्राक्तूर अक्षरे वापरू. ही पद्धत बरीच रूढ आहे. दुसरी रूढ पद्धत म्हणजे संचांक क्रमसूचक संख्याच असल्यामुळे त्यांसाठी ग्रीक अक्षरे वापरणे, पण वर्णमालेच्या मधल्या भागातील अक्षरे निवडणे; उदा. $\kappa, \lambda$.):
\begin{lem}\label{sth:cardinals:cardsasords:lem:CardinalsBehaveRight}
कोणत्याही संच $A$ आणि $B$ साठी:
\begin{align*}
  \cardeq{A}{B} &\text{ तेव्हा आणि केवळ तेव्हाच } \card{A} = \card{B}\\
  \cardle{A}{B} &\text{ तेव्हा आणि केवळ तेव्हाच } \card{A} \leq \card{B}\\
  \cardless{A}{B}&\text{ तेव्हा आणि केवळ तेव्हाच } \card{A} < \card{B}
\end{align*}
\end{lem}

\begin{proof}
दुसऱ्या विधानातील डावीकडून उजवीकडे जाणारी दिशा सिद्ध करू (इतर प्रकरणे सारखीच आहेत आणि सरावासाठी सोडली आहेत). पुढील आकृती पाहा:
\begin{center}
  \begin{tikzpicture}
  \node (nodea) {$A$};
  \node[right = 6em of nodea] (nodeb) {$B$};
  \node[below = 2em of nodea] (nodecarda) {$\card{A}$};
  \node[below = 2em of nodeb] (nodecardb) {$\card{B}$};
  \draw[->] (nodea)--(nodeb);
  \draw[<->] (nodea)--(nodecarda);
  \draw[<->] (nodeb)--(nodecardb);
  \draw[->, dashed] (nodecarda)--(nodecardb);
\end{tikzpicture}
\end{center}
दुतोंडी बाण एकास-एक आच्छादन दाखवतात; त्यांचे अस्तित्व \ref{sth:cardinals:cardsasords:lem:CardinalsExist} मुळे मिळते. $\cardle{A}{B}$ गृहीत धरल्यावर $A\to B$ हे एकास-एक फलन असते. आता $\card{A}$ पासून $A$, तेथून $B$, आणि तेथून $\card{B}$ असा बाणांचा पाठपुरावा केल्यास $\card{A} \to \card{B}$ हे एकास-एक फलन मिळते (आकृतीतील तुटक बाण).
\end{proof}\noindent \ref{sth:cardinals:cardsasords:lem:CardinalsBehaveRight} वापरून श्रेडर--बर्नस्टाइन प्रमेय पुन्हा सिद्ध करता येते. त्याचे विधान असे: $\cardle{A}{B}$ आणि $\cardle{B}{A}$ असतील, तर $\cardeq{A}{B}$. हे विधान आपण \ref{sfr:siz:sb:thm:schroder-bernstein} म्हणून मांडले; पण आधी \ref{sfr:infinite:card-sb:sec} मध्ये ते सिद्ध करण्यासाठी बरीच मेहनत घेतली. आता पाहा:

\begin{proof}[श्रेडर--बर्नस्टाइनची नवी सिद्धता]
$\cardle{A}{B}$ आणि $\cardle{B}{A}$ असतील, तर \ref{sth:cardinals:cardsasords:lem:CardinalsBehaveRight} नुसार $\card{A} \leq \card{B}$ आणि $\card{B} \leq \card{A}$. म्हणून त्रिपर्याय आणि \ref{sth:cardinals:cardsasords:lem:CardinalsBehaveRight} नुसार $\card{A} = \card{B}$ आणि $\cardeq{A}{B}$.
\end{proof}
\noindent
ही सिद्धता फार सोपी असली, तरी तिच्यामागे प्रतिस्थापन (\ref{sth:ordinals:ordtype:thmOrdinalRepresentation} मिळवण्यासाठी) आणि सुक्रमण (\ref{sth:cardinals:cardsasords:lem:CardinalsBehaveRight} निश्चित करण्यासाठी) ही दोन्ही गृहीतके अप्रत्यक्षपणे आहेत. याउलट, \ref{sfr:infinite:card-sb:sec} मधील सिद्धता अधिक स्वयंपूर्ण होती (खरे तर ती~$\Zminus$ मध्येही करता येते).







\section{$\ZFC$: एक महत्त्वाचा टप्पा}\label{sth:cardinals:zfc:sec}

सुक्रमण जोडल्यावर आपण उपपत्तीचा अंतिम महत्त्वाचा टप्पा गाठला आहे. आता~$\ZFC$ साठी आवश्यक असलेली सर्व स्वयंसिद्धके आपल्याकडे आहेत. सविस्तर यादी अशी:

\begin{defn}
$\ZFC$ उपपत्तीत पुढील स्वयंसिद्धके आहेत: विस्तारात्मकता, संयोग, जोड्या, घातसंच, अनंतता, पायाभूतता, सुक्रमण आणि विभक्तीकरण व प्रतिस्थापन या रूपबंधांतील सर्व स्वयंसिद्धके. दुसऱ्या शब्दांत, $\ZF$ मध्ये सुक्रमण जोडले की $\ZFC$ मिळते.
\end{defn}

$\ZFC$ म्हणजे \emph{झर्मेलो--फ्रेंकेल} संच उपपत्ती, \emph{निवडीच्या स्वयंसिद्धकासह}. आपण जोडलेल्या स्वयंसिद्धकाला “निवड” नव्हे तर “सुक्रमण” म्हटले, म्हणून हे थोडे विचित्र वाटू शकते. पण पुढे आपण {निवडीचे स्वयंसिद्धक} मांडल्यावर, सुक्रमण आणि निवड ही $\ZF$ गृहीत धरल्यास समतुल्य आहेत हे दिसेल (\ref{sth:choice:woproblem:thmwochoice} पाहा). त्यामुळे यांपैकी कोणते “मूलभूत” स्वयंसिद्धक मानायचे, याने फरक पडत नाही. आणि “$\ZFC$” हे नाव साहित्यात सर्वमान्य आहे.







\section{सांत, गणनीय, अगणनीय}\label{sth:cardinals:classing:sec}

आता संचांकांची ओळख झाली आहे; म्हणून त्यांच्या विविध प्रकारांविषयी थोडी चर्चा करू. विशेषतः सांत, गणनीय आणि अगणनीय संचांक पाहू.

पहिल्या दोन निष्कर्षांवरून सांत संचांक हे नेमके सांत क्रमसूचक संख्याच असतील; त्यांची व्याख्या आपण मागे \ref{sth:z:infinity-again:defnomega} मध्ये आपल्या \emph{नैसर्गिक संख्या} म्हणून केली होती:

\begin{prop}\label{sth:cardinals:classing:finitecardisoequal}
$n, m \in \omega$ असू द्या. मग $n = m$ तेव्हा आणि केवळ तेव्हाच जेव्हा $\cardeq{n}{m}$.
\end{prop}

\begin{proof}
\emph{डावीकडून उजवीकडे} हे उघड आहे. \emph{उजवीकडून डावीकडे} सिद्ध करण्यासाठी, $n \neq m$ असूनही $\cardeq{n}{m}$ असे गृहीत धरा. त्रिपर्यायानुसार $n \in m$ किंवा $m \in n$; व्यापकता न गमावता $n \in m$ असे समजू. मग $n \subsetneq m$ आणि $f \colon m \to n$ हे एकास-एक आच्छादन असते; म्हणून $m$ डेडेकिंड-अनंत ठरतो. हे \ref{sth:z:infinity-again:naturalnumbersarentinfinite} शी विसंगत आहे.
\end{proof}

\begin{cor}\label{sth:cardinals:classing:naturalsarecardinals}
$n \in \omega$ असेल, तर $n$ हा संचांक आहे.
\end{cor}

\begin{proof}
हे लगेच मिळते.
\end{proof}
\noindent
यावरून संचांक “सांत” किंवा “अनंत” असण्याच्या अनेक स्वाभाविक कल्पना एकाच ठरतात:
\begin{thm}\label{sth:cardinals:classing:generalinfinitycharacter}कोणत्याही संच $A$ साठी पुढील विधाने समतुल्य आहेत:
\begin{enumerate}
  \item\label{sth:cardinals:classing:card:notinomega} $\card{A} \notin \omega$; म्हणजे
  $A$ सांत नाही;
  \item\label{sth:cardinals:classing:card:omegaplus} $\omega \leq \card{A}$;
  \item\label{sth:cardinals:classing:card:infinite} $A$ डेडेकिंड-अनंत आहे.
\end{enumerate}
\end{thm}

\begin{proof}
\ref{sth:ord-arithmetic:using-addition:ordinfinitycharacter},
\ref{sth:cardinals:cardsasords:lem:CardinalsBehaveRight} आणि
\ref{sth:cardinals:classing:naturalsarecardinals} यांवरून.
\end{proof}

म्हणून \ref{sfr:::part} मध्ये काहीशा अनौपचारिकपणे वापरलेल्या कल्पनांची पुढील \emph{व्याख्या} करता येते:

\begin{defn}\label{sth:cardinals:classing:defnfinite}
$\card{A}$ ही नैसर्गिक संख्या, म्हणजे $\card{A} \in \omega$, असेल तेव्हा आणि केवळ तेव्हाच $A$ ला \emph{सांत} म्हणू. अन्यथा $A$ ला \emph{अनंत} म्हणू.
\end{defn}
\noindent
पण ही व्याख्या $\ZFC$ च्या पार्श्वभूमीवर दिली आहे, हे लक्षात ठेवा. शेवटी, प्रत्येक संचाला संचांक असतो याची खात्री मिळवण्यासाठी सुक्रमणाचे स्वयंसिद्धक लागले होते. खरे तर सुक्रमणाशिवाय सांतही नाही आणि डेडेकिंड-अनंतही नाही असा संच असू शकतो. या प्रकारच्या प्रश्नाकडे \ref{sth:choice:chap} मध्ये परत येऊ. सध्या सुक्रमण गृहीत धरून पुढे जाऊ.

आता सांत संचांकांवरून अनंत संचांकांकडे वळू. दोन प्राथमिक निरीक्षणे अशी:

\begin{cor}\label{sth:cardinals:classing:omegaisacardinal}
$\omega$ हा लघुतम अनंत संचांक आहे.
\end{cor}

\begin{proof}
$\omega$ हा संचांक आहे: कारण $\omega$ डेडेकिंड-अनंत आहे; आणि कोणत्याही $n \in \omega$ साठी $\cardeq{\omega}{n}$ असेल, तर $n$ डेडेकिंड-अनंत ठरेल, जे \ref{sth:z:infinity-again:naturalnumbersarentinfinite} शी विसंगत आहे. आता व्याख्येनुसार $\omega$ हा लघुतम अनंत संचांक आहे.
\end{proof}

\begin{cor}
प्रत्येक अनंत संचांक ही सीमा क्रमसूचक संख्या असते.
\end{cor}

\begin{proof}
$\alpha$ ही अनंत उत्तरवर्ती क्रमसूचक संख्या असू द्या; म्हणजे काही $\beta$ साठी $\alpha = \beta
\ordplus 1$. पूर्ववर्ती सांत असती तर दिलेली उत्तरवर्ती सुद्धा सांत असती; म्हणून $\beta$ अनंत आहे (\ref{sth:cardinals:classing:finitecardisoequal} मधील सांत प्रकरण लक्षात घ्या). त्यामुळे \ref{sth:ord-arithmetic:using-addition:ordinfinitycharacter} नुसार $\cardeq{\beta}{\beta \ordplus 1}$. आता
\ref{sth:cardinals:cardsasords:lem:CardinalsBehaveRight} नुसार
$\card{\beta} = \card{\beta\ordplus 1} = \card{\alpha}$; म्हणून
$\alpha \neq \card{\alpha}$.
\end{proof}

मागे \ref{sfr:siz:enm-alt:defn:enumerable} मध्येच आपण गणनीय आणि अगणनीय अनंत संच वेगळे ओळखता येतात, हे दाखवले होते. त्या व्याख्येवरून पुढील निष्कर्ष स्वाभाविकपणे मिळतो:

\begin{prop}
$A$ हा गणनीय असेल तेव्हा आणि केवळ तेव्हाच जेव्हा $\card{A} \leq \omega$; तसेच $A$ हा अगणनीय असेल तेव्हा आणि केवळ तेव्हाच जेव्हा $\omega < \card{A}$.
\end{prop}

\begin{proof}
त्रिपर्यायानुसार ही दोन विधाने समतुल्य आहेत; म्हणून $A$ हा गणनीय तेव्हा आणि केवळ तेव्हाच जेव्हा $\card{A} \leq \omega$, एवढेच सिद्ध करणे पुरेसे आहे. \emph{उजवीकडून डावीकडे}: $\card{A} \leq \omega$ असेल, तर
\ref{sth:cardinals:cardsasords:lem:CardinalsBehaveRight} आणि
\ref{sth:cardinals:classing:omegaisacardinal} नुसार $\cardle{A}{\omega}$. \emph{डावीकडून उजवीकडे}: $A$ हा गणनीय असेल, तर \ref{sfr:siz:enm-alt:defn:enumerable} नुसार पुढील तीन प्रकरणांपैकी एक घडते:
\begin{enumerate}
  \item $A = \emptyset$ असेल, तर \ref{sth:cardinals:classing:naturalsarecardinals} आणि
  \ref{sth:cardinals:cardsasords:lem:CardinalsBehaveRight} नुसार $\card{A} = 0 \in \omega$.
  \item $\cardeq{n}{A}$ असेल, तर \ref{sth:cardinals:classing:naturalsarecardinals} आणि
  \ref{sth:cardinals:cardsasords:lem:CardinalsBehaveRight} नुसार $\card{A} = n \in \omega$.
  \item $\cardeq{\omega}{A}$ असेल, तर \ref{sth:cardinals:classing:omegaisacardinal} नुसार $\card{A} = \omega$.
\end{enumerate}
म्हणून तिन्ही प्रकरणांत $\card{A} \leq \omega$.
\end{proof}
\noindent
खरे तर $\omega$ चे स्थान विशेष आहे. गणनीय क्रमसूचक संख्या अनेक असल्या तरी:

\begin{cor}
$\omega$ हा एकमेव गणनीय अनंत संचांक आहे.
\end{cor}

\begin{proof}
$\cardfont{a}$ हा गणनीय अनंत संचांक असू द्या. $\cardfont{a}$ अनंत असल्यामुळे $\omega \leq \cardfont{a}$. तसेच $\cardfont{a}$ हा गणनीय संचांक असल्यामुळे $\cardfont{a} =
\card{\cardfont{a}} \leq \omega$. म्हणून त्रिपर्यायानुसार $\cardfont{a} = \omega$. \end{proof}

अर्थात, संचांक अनंत संख्येने आहेत. म्हणून प्रश्न पडेल: \emph{एकूण संचांक किती आहेत?} पुढील निष्कर्षांवरून या प्रश्नाचा पुनर्विचार करावा लागेल.

\begin{prop}\label{sth:cardinals:classing:unioncardinalscardinal}
$X$ चा प्रत्येक सदस्य संचांक असेल, तर $\bigcup X$ हाही संचांक आहे.
\end{prop}

\begin{proof}
$\bigcup X$ ही क्रमसूचक संख्या आहे, हे तपासणे सोपे आहे. $\alpha \in
\bigcup X$ ही क्रमसूचक संख्या असू द्या; मग काही संचांक $\cardfont{b}$ साठी $\alpha \in \cardfont{b} \in X$. $\cardfont{b}$ संचांक असल्यामुळे
$\cardless{\alpha}{\cardfont{b}}$. तसेच $\cardfont{b} \subseteq
\bigcup X$ असल्यामुळे $\cardle{\cardfont{b}}{\bigcup X}$; म्हणून
$\cardneq{\alpha}{\bigcup X}$. हा निष्कर्ष सर्वत्र लागू केल्यावर $\bigcup X$ हा संचांक ठरतो.
\end{proof}

\begin{thm}\label{sth:cardinals:classing:lem:NoLargestCardinal}
सर्वांत मोठा संचांक नाही.
\end{thm}

\begin{proof}
कोणत्याही संचांक $\cardfont{a}$ साठी कँटरचे प्रमेय
(\ref{sfr:siz:car:thm:cantor}) आणि
\ref{sth:cardinals:cardsasords:lem:CardinalsExist} यांवरून
$\cardfont{a} < \card{\Pow{\cardfont{a}}}$ मिळते.
\end{proof}

\begin{thm}
सर्व संचांकांचा संच अस्तित्वात नाही.
\end{thm}

\begin{proof}
विसंगतीने सिद्ध करण्यासाठी $C = \Setabs{\cardfont{a}}{\cardfont{a} \text{
संचांक आहे}}$ असे समजा. आता \ref{sth:cardinals:classing:unioncardinalscardinal} नुसार $\bigcup C$ हा संचांक आहे; म्हणून \ref{sth:cardinals:classing:lem:NoLargestCardinal} नुसार $\cardfont{b} > \bigcup C$ असा संचांक आहे. व्याख्येनुसार
$\cardfont{b} \in C$; म्हणून $\cardfont{b} \subseteq \bigcup{C}$ आणि त्यामुळे
$\cardfont{b} \leq \bigcup C$; ही विसंगती आहे.
\end{proof}

याची रसेलचा विरोधाभास आणि बुराली-फॉर्ती निष्कर्ष या दोन्हींशी तुलना करा.







\section{परिशिष्ट: ह्यूमचे तत्त्व}\label{sth:cardinals:hp:sec}

\ref{sth:cardinals:cp:sec} मध्ये आपण कँटरचे तत्त्व मांडले. ते असे:
\begin{align*}
  \card{A} = \card{B} & \text{ तेव्हा आणि केवळ तेव्हाच } A \approx B.
\intertext{आज ज्याला \emph{ह्यूमचे तत्त्व} म्हणतात त्याच्याशी हे खूप मिळतेजुळते आहे. ह्यूमचे तत्त्व असे:}
  \fregenum{x} {F(x)} = \fregenum{x}{G(x)} & \text{ तेव्हा आणि केवळ तेव्हाच } F \sim G
\end{align*}
येथे ‘$F \sim G$’ याचा अर्थ $F$ पूर्ण करणाऱ्या आणि $G$ पूर्ण करणाऱ्या वस्तूंची संख्या नेमकी समान आहे; म्हणजे $F$ पूर्ण करणाऱ्या व $G$ पूर्ण करणाऱ्या वस्तूंमध्ये द्विएक संबंध लावता येतो; म्हणजेच:
\begin{align*}
  \exists R(&\forall v\forall y(Rvy \lif (Fv \land Gy)) \land {}\\
    &\forall v(Fv \lif \lexists![y][Rvy]) \land {}\\
    &\forall y(Gy \lif \lexists![v][Rvy]))
\end{align*}
पण ह्यूमचे तत्त्व आणि कँटरचे तत्त्व यांत प्रकारभेद आहे. कँटरच्या तत्त्वातील चल “$A$” आणि “$B$” ही \emph{संच} दर्शवणारी प्रथम-क्रमाची पदे आहेत. ह्यूमच्या तत्त्वातील “$F$”, “$G$” आणि “$R$” ही प्रथम-क्रमाची पदे \emph{नाहीत}; ती \emph{विधेयस्थानी} येतात. (कदाचित ती \emph{गुणधर्म} दर्शवत असतील.) म्हणून ह्यूमचे तत्त्व साध्या भाषेत असे सांगता येईल: $F$ पूर्ण करणाऱ्या आणि $G$ पूर्ण करणाऱ्या वस्तूंची संख्या समान असते तेव्हा आणि केवळ तेव्हाच जेव्हा $F$ पूर्ण करणाऱ्या व $G$ पूर्ण करणाऱ्या वस्तूंमध्ये द्विएक संबंध असतो. याला \emph{ह्यूमचे तत्त्व} म्हणतात, कारण ह्यूम यांनी एकदा असे लिहिले:
\begin{quote}
  दोन संख्या अशा रीतीने जोडलेल्या असतील की एकीतील प्रत्येक एककाला दुसरीतील नेहमी एक एकक जुळते, तर आपण त्या समान आहेत असे म्हणतो.
  (ह्यूम 1740, भाग III, पुस्तक 1, \S1)
\end{quote}
फ्रेगे यांनी फ्रेगे (1884, \S63) मध्ये ह्यूम यांचा हा उतारा उद्धृत केला आणि नंतर (वस्तुतः) आपण ज्याला ह्यूमचे तत्त्व म्हटले आहे त्याची रूपरेषा मांडली; त्यामुळे आधुनिक गणितीय तर्कशास्त्रात या तत्त्वाला महत्त्व प्राप्त झाले.

ह्यूमचे तत्त्व आणि मूलनियम पाच यांतील रचनात्मक साम्य लक्षात घ्या. \ref{sth:story:blv:sec} मध्ये आपण मूलनियम पाच असे मांडले होते:
\[
  \fregeext{x}{F(x)} = \fregeext{x}{G(x)} \text{तेव्हा आणि केवळ तेव्हाच } \lforall[x][(F(x) \liff G(x))].
\]
या टप्प्यावर थोडे स्पष्टीकरण आणि तुलना उपयोगी ठरेल.

ह्यूमचे तत्त्व किंवा मूलनियम पाच यांसारख्या तत्त्वांचा अर्थ दोन प्रकारे घेता येतो: \emph{स्वसमावेशरहित} किंवा \emph{स्वसमावेशक} (\ref{sth:story:predicative:sec} आठवा). मूलनियम पाचच्या स्वसमावेशक अर्थानुसार, प्रत्येक~$F$ साठी $\fregeext{x}{F(x)}$ ही वस्तू, मूलनियम पाच मांडताना आपण ज्यावर संख्यापन केले त्याच प्रांतात येते. त्याचप्रमाणे ह्यूमच्या तत्त्वाच्या स्वसमावेशक अर्थानुसार, प्रत्येक~$F$ साठी $\fregenum{x}{F(x)}$ ही वस्तू, ह्यूमचे तत्त्व मांडताना ज्यावर संख्यापन केले त्याच प्रांतात येते. याउलट, \emph{स्वसमावेशरहित} अर्थानुसार $\fregeext{x}{F(x)}$ आणि~$\fregenum{x}{F(x)}$ या वस्तू एखाद्या \emph{वेगळ्या} प्रांतातील असतील.

मूलनियम पाचचा स्वसमावेशक अर्थ घेतल्यास सहज गुणधर्माधारित संचरचनेमुळे विसंगती निर्माण होते (तपशीलांसाठी \ref{sth:story:blv:sec} पाहा). सहज गुणधर्माधारित संचरचनेप्रमाणेच, त्याचा \emph{स्वसमावेशरहित} अर्थ घेतल्यास तो सुसंगत करता येतो. पण मग त्याच्याकडून अपेक्षित सर्व काही बहुधा साध्य होणार नाही.

ह्यूमचे तत्त्व मात्र सुसंगतपणे स्वसमावेशक अर्थाने घेता \emph{येते}. असा अर्थ घेतल्यावर ते बरेच सामर्थ्यवान ठरते.

उदाहरणार्थ, “$x \neq x$” हे विधेय पाहा; स्पष्टच कोणतीही वस्तू ते पूर्ण करत नाही. ह्यूमच्या तत्त्वावरून आता $\# x( x\neq
x)$ ही वस्तू मिळते. तिला आपण~$0$ ही संख्या मानू शकतो. आता \emph{स्वसमावेशक} अर्थानुसार—पण \emph{फक्त} त्याच अर्थानुसार—ही वस्तू $0$ आपल्या मूळ संख्यापन-प्रांतात येते. म्हणून “$x = 0$” या विधेयाला ह्यूमचे तत्त्व लावून $\#x (x =
0)$ ही वस्तू मिळवता येते. तिला आपण~$1$ ही संख्या मानू शकतो. शिवाय ह्यूमच्या तत्त्वावरून $0 \neq 1$ मिळते, कारण स्वतःशी असमान वस्तू आणि $0$ शी समान वस्तू यांच्यात द्विएक संबंध असू शकत नाही (पहिल्या प्रकारची एकही नाही, दुसऱ्या प्रकारची एक आहे). आता पुन्हा स्वसमावेशक अर्थाने $1$ आपल्या मूळ संख्यापन-प्रांतात येतो. म्हणून “$(x = 0 \lor x = 1)$” या विधेयाला ह्यूमचे तत्त्व लावून $\#x(x = 0 \lor
x = 1)$ ही वस्तू मिळते. तिला आपण~$2$ ही संख्या मानू शकतो; मग $0\neq 2$ आणि $1 \neq 2$ वगैरे दाखवता येते.

थोडक्यात, ह्यूमचे तत्त्व स्वसमावेशक अर्थाने घेतल्यास \emph{अनंत संख्येने वस्तू आहेत} हे त्यावरून मिळते. त्यामुळे \emph{नव-फ्रेगेवादी तर्कमीमांसावाद्यांना} अंकगणिताचा पाया म्हणून ह्यूमचे तत्त्व घ्यावेसे वाटले आहे.

पण फ्रेगे यांनी \emph{स्वतः} ह्यूमचे तत्त्व अंकगणिताचा पाया म्हणून घेतले नाही. उलट, एका स्पष्ट व्याख्येवरून त्यांनी ह्यूमचे तत्त्व सिद्ध केले: $\fregenum{x}{F(x)}$ ची व्याख्या $F \sim \Phi$ या संकल्पनेचा विस्तार अशी केली जाते. आधुनिक संज्ञेत हे $\fregenum{x}{F(x)} = \Setabs{G}{F \sim G}$ असे मांडण्याचा प्रयत्न करता येईल; पण मग सहज गुणधर्माधारित संचरचनेतील अडचणी पुन्हा उभ्या राहतील.



\chapter{संचांकांचे अंकगणित}\label{sth:card-arithmetic:chap}




\section{मूलभूत क्रियांची व्याख्या}\label{sth:card-arithmetic:opps:sec}

येथे क्रम लक्षात ठेवण्याची गरज नाही; म्हणून संचांकांचे अंकगणित परिभाषित करणे क्रमसूचक अंकगणितापेक्षा सोपे आहे. बेरीज, गुणाकार आणि घातांकन या तिन्ही क्रियांची व्याख्या एकाच वेळी करू.

\begin{defn}
$\cardfont{a}$ आणि $\cardfont{b}$ हे संचांक असतील, तर:
\begin{align*}
  \cardfont{a} \cardplus \cardfont{b} &\defis
  \card{\cardfont{a} \disjointsum \cardfont{b}}\\
  \cardfont{a} \cardtimes \cardfont{b} &\defis
  \card{\cardfont{a} \times \cardfont{b}}\\
  \cardexpo{\cardfont{a}}{\cardfont{b}} &\defis
  \card{\funfromto{\cardfont{b}}{\cardfont{a}}}
\end{align*}
येथे $\funfromto{X}{Y} = \Setabs{f}{f \text{ हे फलन आहे } X \to Y}$. (कोणत्याही संच $X$ आणि $Y$ साठी $\funfromto{X}{Y}$ अस्तित्वात असतो, हे दाखवणे सोपे आहे; तो सरावासाठी सोडतो.)
\end{defn}

\begin{prob}
कोणत्याही संच $X$ आणि~$Y$ साठी $\funfromto{X}{Y}$ अस्तित्वात असतो, हे $\Zminus$ मध्ये सिद्ध करा. $\ZF$ मध्ये काम करून, \ref{sth:ord-arithmetic:using-addition:rankcomputation} च्या धर्तीवर $\setrank{X}$ आणि $\setrank{Y}$ यांवरून $\setrank{\funfromto{X}{Y}}$ मोजा.
\end{prob}

या व्याख्येचे स्पष्टीकरण उपयोगी ठरेल. बेरीजेसाठी आपण \ref{sth:ord-arithmetic:add:defdissum} मध्ये परिभाषित केलेली विभक्त बेरीज $\disjointsum$ वापरतो; सांत प्रकरणांत या व्याख्येतून योग्य उत्तर मिळते, हे सहज दिसते. गुणाकारासाठी, \ref{sfr:set:pai:cardnmprod} सांगते की $A$ मध्ये $n$ आणि $B$ मध्ये $m$ घटक असतील, तर $A \times B$ मध्ये $n \cdot m$ घटक असतात; म्हणून आपली व्याख्या ही कल्पना अनंतपार गुणाकारापर्यंत विस्तारते. घातांकनाची बाबही तशीच: सांत प्रकरणातील कल्पना अनंतपार प्रकरणात नेतो. खरे तर काही बाबतींत अनंतपार संचांकांचे अंकगणित हे अनंतपार क्रमसूचक अंकगणितापेक्षा “सामान्य” अंकगणिताशी अधिक जुळते:

\begin{prop}\label{sth:card-arithmetic:opps:cardplustimescommute}
$\cardplus$ आणि $\cardtimes$ या क्रिया क्रमनिरपेक्ष व सहयोगी आहेत.
\end{prop}

\begin{proof}
क्रमनिरपेक्षतेसाठी \ref{sth:cardinals:cardsasords:lem:CardinalsBehaveRight} नुसार
$\cardeq{(\cardfont{a} \disjointsum
\cardfont{b})}{(\cardfont{b} \disjointsum \cardfont{a})}$ आणि
$\cardeq{(\cardfont{a} \times \cardfont{b})}{(\cardfont{b} \times
\cardfont{a})}$ एवढे पाहिले की पुरे. सहयोगित्व सरावासाठी सोडतो.
\end{proof}

\begin{prob}
$\cardplus$ आणि $\cardtimes$ या क्रिया सहयोगी आहेत, हे सिद्ध करा.
\end{prob}

\begin{prop}
$A$ अनंत असेल तेव्हा आणि केवळ तेव्हाच जेव्हा $\card{A} \cardplus 1 = 1 \cardplus \card{A} = \card{A}$.
\end{prop}

\begin{proof}
\ref{sth:ord-arithmetic:using-addition:ordinfinitycharacter} आणि
\ref{sth:cardinals:cardsasords:lem:CardinalsBehaveRight} यांवरून,
\ref{sth:cardinals:classing:generalinfinitycharacter} प्रमाणे.
\end{proof}

यावरून क्रमसूचक आणि संचांक बेरीज/गुणाकारासाठी वेगवेगळी चिन्हे का वापरावी लागतात हे स्पष्ट होते: या क्रिया खरोखर \emph{भिन्न} आहेत. पुढील दोन निष्कर्ष क्रमसूचक आणि संचांक घातांकनही भिन्न क्रिया आहेत हे दाखवतात. (\ref{sth:z:infinity-again:defnomega} नुसार $2 = \{0,
1\}$ हे आठवा):

\begin{lem}\label{sth:card-arithmetic:opps:lem:SizePowerset2Exp}
कोणत्याही $A$ साठी $\card{\Pow{A}} = \cardexpo{2}{\card{A}}$.
\end{lem}

\begin{proof}
प्रत्येक उपसंच $B \subseteq A$ साठी $\chi_B \in \funfromto{A}{2}$ असे पुढीलप्रमाणे ठरवा:
\begin{align*}
  \chi_{B}(x) &\defis
  \begin{cases}
    1 & \text{जर }x\in B\\
    0 & \text{अन्यथा.}
  \end{cases}
\end{align*}
आता $f(B) = \chi_B$ असे ठेवा; यातून $f \colon \Pow{A}
\to \funfromto{A}{2}$ हे एकास-एक आच्छादन मिळते. म्हणून $\cardeq{\Pow{A}}{\funfromto{A}{2}}$. त्यामुळे
$\cardeq{\Pow{A}}{\funfromto{\card{A}}{2}}$; आणि म्हणून
$\card{\Pow{A}} = \card{\funfromto{\card{A}}{2}} =
\cardexpo{2}{\card{A}}$.
\end{proof}

या संक्षिप्त पुराव्यात \ref{sfr:siz:red-alt:sec} मधील चर्चा बहुतांशी सामावली आहे. तेथे आपण $\Pow{\omega}$ ची अगणनीयता अनंत द्विमान चिन्हमालांच्या संचाच्या, म्हणजे $\Bin^\omega$ च्या, अगणनीयतेकडे कशी “न्यून” करता येते ते दाखवले. प्रत्यक्षात $\Bin^{\omega}$ म्हणजेच $\funfromto{\omega}{2}$; आणि आधीच्या पुराव्यावरून \ref{sfr:siz:red-alt:sec} मधील युक्तिवाद~$\omega$ ऐवजी कोणताही संच~$A$ वापरून करता येतो. यावरून संचांक घातांकनाविषयीही एक झटपट निष्कर्ष मिळतो:

\begin{cor}\label{sth:card-arithmetic:opps:cantorcor}
कोणत्याही संचांक~$\cardfont{a}$ साठी $\cardfont{a} < \cardexpo{2}{\cardfont{a}}$.
\end{cor}

\begin{proof}
कँटरचे प्रमेय (\ref{sfr:siz:car:thm:cantor}) आणि
\ref{sth:card-arithmetic:opps:lem:SizePowerset2Exp} यांवरून.
\end{proof}
\noindent
म्हणून $\omega < \cardexpo{2}{\omega}$. पण लक्षात घ्या: हा निष्कर्ष \emph{संचांक} घातांकनाबद्दल आहे. त्याची \emph{क्रमसूचक} घातांकनाशी तुलना केली पाहिजे; त्या बाबतीत $\omega =
\ordexpo{2}{\omega}$ असते (\ref{sth:ord-arithmetic:expo:sec} पाहा).

संचांक घातांकनाच्या चर्चेत असतानाच $\Real$ हा नेमक्या कोणत्या “अर्थाने” अगणनीय आहे, हेही अधिक स्पष्ट सांगता येईल.

\begin{thm}\label{sth:card-arithmetic:opps:continuumis2aleph0}
$\card{\Real} = \cardexpo{2}{\omega}$
\end{thm}

\begin{proof}[पुराव्याची रूपरेषा]
हे सिद्ध करण्याचे अनेक मार्ग आहेत. सर्वांत सरळ म्हणजे $\cardle{\Pow{\omega}}{\Real}$ आणि
$\cardle{\Real}{\Pow{\omega}}$ दाखवणे; मग श्रेडर--बर्नस्टाइन वापरून
$\cardeq{\Real}{\Pow{\omega}}$ मिळते, आणि
\ref{sth:card-arithmetic:opps:lem:SizePowerset2Exp} वरून
$\card{\Real} = \cardexpo{2}{\omega}$ मिळते. $f \colon
\Pow{\omega} \to \Real$ आणि $g \colon \Real \to \Pow{\omega}$ ही एकास-एक फलने परिभाषित करण्याचा (बोधप्रद) सराव आपण सोडतो.
\end{proof}

\begin{prob}
$\cardle{\Pow{\omega}}{\Real}$ आणि
$\cardle{\Real}{\Pow{\omega}}$ दाखवून \ref{sth:card-arithmetic:opps:continuumis2aleph0} चा पुरावा पूर्ण करा.
\end{prob}







\section{बेरीज आणि गुणाकार सुलभ करणे}\label{sth:card-arithmetic:simp:sec}

अनंतपार संचांक बेरीज आणि गुणाकार \emph{फारच} सोपे ठरतात. कारण संचांक या (काही विशिष्ट) क्रमसूचक संख्या आहेत; म्हणून त्या सुक्रमित आहेत आणि त्यांच्यावर एका विशिष्ट पद्धतीने काम करता येते. पण हे दाखवणे \emph{तितके} सोपे नाही. सुरुवातीला एक कल्पक व्याख्या हवी:

\begin{defn}
क्रमसूचक संख्यांच्या जोड्यांवरील \emph{कॅनॉनिकल क्रम} $\canonord$ पुढीलप्रमाणे परिभाषित करू: $\tuple{\alpha_1, \alpha_2} \canonord
\tuple{\beta_1, \beta_2}$ तेव्हा आणि केवळ तेव्हाच जेव्हा पुढीलपैकी एक अट पूर्ण होते:
\begin{enumerate}
  \item $\max(\alpha_1, \alpha_2) < \max(\beta_1, \beta_2)$; किंवा
  \item $\max(\alpha_1, \alpha_2) = \max(\beta_1, \beta_2)$ आणि
  $\alpha_1 < \beta_1$; किंवा
  \item $\max(\alpha_1, \alpha_2) = \max(\beta_1, \beta_2)$ आणि
  $\alpha_1 = \beta_1$ आणि $\alpha_2 < \beta_2$.
\end{enumerate}
\end{defn}

\begin{lem}
कोणत्याही क्रमसूचक संख्या $\alpha$ साठी $\tuple{\alpha \times \alpha, \canonord}$ हे सुक्रमण आहे.
\end{lem}

\begin{proof}
$\alpha \times \alpha$ वर $\canonord$ हा क्रम तुलनीय आहे, हे स्पष्ट आहे. समजा $\tuple{\alpha_1, \alpha_2}$ आणि $\tuple{\beta_1,
\beta_2}$ यांपैकी कोणतीही जोडी दुसरीपेक्षा $\canonord$-लहान नाही. मग $\max(\alpha_1,
\alpha_2) = \max(\beta_1, \beta_2)$ आणि $\alpha_1 = \beta_1$ आणि
$\alpha_2 = \beta_2$; म्हणून $\tuple{\alpha_1, \alpha_2} =
\tuple{\beta_1,  \beta_2}$.

सुक्रमण सिद्ध करण्यासाठी रिकामा नसलेला $X \subseteq \alpha\times\alpha$ घ्या. $\alpha$ क्रमसूचक संख्या असल्यामुळे
$\Setabs{\max(\gamma_1, \gamma_2)}{\tuple{\gamma_1,
\gamma_2} \in X}$ या संचाचा लघुतम सदस्य एखादा $\delta$ असतो. आता
$\Setabs{\tuple{\gamma_1,\gamma_2} \in X}{\max(\gamma_1, \gamma_2) =
\delta}$ मधून पहिला निर्देशांक लघुतम नसलेल्या सर्व जोड्या काढून टाका; उरलेल्यांपैकी दुसरा निर्देशांक लघुतम असलेली जोडी हा $X$ चा $\canonord$-लघुतम सदस्य आहे.
\end{proof}
\noindent
आता एक अगदी छोटे आणि सोपे निरीक्षण:

\begin{prop}\label{sth:card-arithmetic:simp:simplecardproduct}
$\cardeq{\alpha}{\beta}$ असेल, तर $\cardeq{\alpha \times \alpha}{\beta
\times \beta}$.
\end{prop}

\begin{proof}
$f \colon \alpha \to \beta$ वरून $\tuple{\gamma_1,
\gamma_2} \mapsto \tuple{f(\gamma_1), f(\gamma_2)}$ हे फलन मिळते.
\end{proof}
\noindent
आता एका महत्त्वाच्या सहायक प्रमेयासाठी ही सगळी निरीक्षणे वापरू:
\begin{lem}\label{sth:card-arithmetic:simp:alphatimesalpha}
कोणत्याही अनंत क्रमसूचक संख्या $\alpha$ साठी $\cardeq{\alpha}{\alpha \times \alpha}$.
\end{lem}

\begin{proof}
विसंगतीने सिद्ध करण्यासाठी, हे विधान खोटे ठरते अशी लघुतम अनंत क्रमसूचक संख्या $\alpha$ घ्या. \ref{sfr:siz:zigzag:natsquaredenumerable} नुसार
$\cardeq{\omega}{\omega\times\omega}$, म्हणून $\omega \in \alpha$.
याशिवाय $\alpha$ हा संचांक आहे: नसल्याचे गृहीत धरू; मग
$\card{\alpha} \in \alpha$ आणि विगमन-गृहीतकानुसार
$\cardeq{\card{\alpha}}{\card{\alpha} \times \card{\alpha}}$.
व्याख्येनुसार $\cardeq{\card{\alpha}}{\alpha}$; त्यामुळे
\ref{sth:card-arithmetic:simp:simplecardproduct} नुसार $\cardeq{\alpha}{\alpha\times\alpha}$, ही विसंगती आहे.

आता प्रत्येक $\tuple{\gamma_1, \gamma_2} \in \alpha \times \alpha$ साठी पुढील आरंभीचा खंड पाहा:
\begin{align*}
  \text{Seg}(\gamma_1, \gamma_2) &= \Setabs{\tuple{\delta_1, \delta_2} \in \alpha \times \alpha}{\tuple{\delta_1, \delta_2} \canonord \tuple{\gamma_1, \gamma_2}}
\end{align*}
$\gamma = \max(\gamma_1, \gamma_2)$ असे ठेवा; मग $\tuple{\gamma_1, \gamma_2} \canonord \tuple{\gamma+1, \gamma + 1}$. म्हणून $\gamma$ अनंत असेल तेव्हा:
\begin{align*}
  \text{Seg}(\gamma_1, \gamma_2) &
  \precsim ((\gamma \ordplus 1)\ordtimes (\gamma \ordplus 1))\\
  &\approx (\gamma \ordtimes \gamma)
  \text{, \ref{sth:ord-arithmetic:using-addition:ordinfinitycharacter} आणि
  \ref{sth:card-arithmetic:simp:simplecardproduct} नुसार}\\
  &\approx \gamma \text{, विगमन-गृहीतकानुसार}\\
  & \prec \alpha\text{, कारण $\alpha$ संचांक आहे}
\end{align*}
निर्देशांकातील मोठी संख्या सांत असेल तर हा खंडही सांत असतो; त्यामुळे तो दिलेल्या अनंत संचांकापेक्षा लहान आहे. म्हणून $\ordtype{\alpha\times \alpha, \canonord} \leq \alpha$, आणि त्यामुळे
$\cardle{\alpha \times \alpha}{\alpha}$. अर्थातच
$\cardle{\alpha}{\alpha \times \alpha}$ असल्याने श्रेडर--बर्नस्टाइननुसार अपेक्षित निष्कर्ष मिळतो.
\end{proof}

अखेरीस, सुलभीकरणाचा आपला निष्कर्ष मिळतो:

\begin{thm}\label{sth:card-arithmetic:simp:cardplustimesmax}
$\cardfont{a}, \cardfont{b}$ हे अनंत संचांक असतील, तर:
\[
  \cardfont{a}
\cardtimes \cardfont{b} = \cardfont{a} \cardplus \cardfont{b} =
\text{max}(\cardfont{a}, \cardfont{b}).
\]
\end{thm}

\begin{proof}
व्यापकता न गमावता $\cardfont{a} = \max(\cardfont{a},
\cardfont{b})$ असे समजा. मग \ref{sth:card-arithmetic:simp:alphatimesalpha} वापरून
$\cardfont{a}\cardtimes\cardfont{a} = \cardfont{a} \leq \cardfont{a}
\cardplus \cardfont{b} \leq \cardfont{a} \cardplus \cardfont{a} \leq
\cardfont{a} \cardtimes \cardfont{a}$ मिळते. \end{proof}\noindent त्याचप्रमाणे,
$\cardfont{a}$ अनंत असेल, तर $\leq\cardfont{a}$ आकारमानाच्या $\cardfont{a}$ इतक्या संचांच्या संयोगाचे आकारमानही $\leq\cardfont{a}$ असते:

\begin{prop}\label{sth:card-arithmetic:simp:kappaunionkappasize}
$\cardfont{a}$ हा अनंत संचांक असू द्या. प्रत्येक क्रमसूचक संख्या $\beta
\in \cardfont{a}$ साठी $X_\beta$ हा $\card{X_\beta} \leq
\cardfont{a}$ असलेला संच असू द्या. मग $\card{\bigcup_{\beta \in \cardfont{a}} X_\beta}
\leq \cardfont{a}$.
\end{prop}

\begin{proof}
प्रत्येक $\beta \in \cardfont{a}$ साठी $f_\beta \colon
X_\beta \to \cardfont{a}$ हे एकास-एक फलन निश्चित करा.\footnote{ही फलने “निश्चित” कशी केली? \ref{sth:choice:countablechoice:sec} पाहा.} आता $g \colon
\bigcup_{\beta \in \cardfont{a}} X_\beta \to \cardfont{a} \times
\cardfont{a}$ हे एकास-एक फलन पुढीलप्रमाणे ठरवा: $g(v) = \tuple{\beta, f_\beta(v)}$, येथे $v \in
X_\beta$ आणि कोणत्याही $\gamma \in \beta$ साठी $v \notin X_\gamma$. मग
\ref{sth:card-arithmetic:simp:cardplustimesmax} नुसार
$\bigcup_{\beta \in \cardfont{a}} X_\beta \preceq \cardfont{a} \times
\cardfont{a} \approx \cardfont{a}$.
\end{proof}







\section{संचांक घातांकनाची काही सुलभीकरणे}\label{sth:card-arithmetic:expotough:sec}

$\canonord$ ची व्याख्या थोडी गुंतागुंतीची होती; पण त्यातून संचांक बेरीज आणि गुणाकाराविषयी \ref{sth:card-arithmetic:simp:cardplustimesmax} हा उपयोगी निष्कर्ष मिळाला. अनंतपार घातांकन मात्र इतक्या सरळ रीतीने सुलभ करता येत नाही. असे का, हे समजावण्यासाठी सांत प्रकरणातील परिचित नमुन्याचा विस्तार करणाऱ्या निष्कर्षापासून सुरुवात करू (त्याचा पुरावा मात्र बऱ्याच अमूर्त पातळीवर आहे):

\begin{prop}\label{sth:card-arithmetic:expotough:simplecardexpo}
कोणत्याही संचांक $\cardfont{a}, \cardfont{b}, \cardfont{c}$ साठी
$\cardexpo{\cardfont{a}}{\cardfont{b} \cardplus \cardfont{c}} =
\cardexpo{\cardfont{a}}{\cardfont{b}} \cardtimes
\cardexpo{\cardfont{a}}{\cardfont{c}}$ आणि
$\cardexpo{(\cardexpo{\cardfont{a}}{\cardfont{b}})}{\cardfont{c}} =
\cardexpo{\cardfont{a}}{\cardfont{b} \cardtimes \cardfont{c}}$.
\end{prop}

\begin{proof}
पहिल्या विधानासाठी $f \colon
(\cardfont{b}\disjointsum\cardfont{c}) \to \cardfont{a}$ हे फलन पाहा. प्रत्येक
$\beta \in \cardfont{b}$ साठी $f_\cardfont{b}(\beta) = f(\beta, 0)$ आणि प्रत्येक
$\gamma \in \cardfont{c}$ साठी $f_\cardfont{c}(\gamma) = f(\gamma, 1)$ ठरवून त्याचे दोन भाग करा. मग $f \mapsto
\tuple{f_{\cardfont{b}}, f_{\cardfont{c}}}$ हे
$\funfromto{\cardfont{b} \disjointsum \cardfont{c}}{\cardfont{a}} \to
(\funfromto{\cardfont{b}}{\cardfont{a}} \times
\funfromto{\cardfont{c}}{\cardfont{a}})$ यांतील एकास-एक आच्छादन आहे.

दुसऱ्या विधानासाठी $f \colon \cardfont{c} \to
(\funfromto{\cardfont{b}}{\cardfont{a}})$ हे फलन पाहा; म्हणजे प्रत्येक $\gamma \in
\cardfont{c}$ साठी $f(\gamma) \colon \cardfont{b} \to
\cardfont{a}$ हे फलन मिळते. आता प्रत्येक $\tuple{\beta, \gamma} \in \cardfont{b} \times \cardfont{c}$ साठी $f^*(\beta, \gamma) = (f(\gamma))(\beta)$ ठरवा.
मग $f \mapsto f^*$ हे
$\funfromto{\cardfont{c}}{(\funfromto{\cardfont{b}}{\cardfont{a}})}
\to \funfromto{\cardfont{b} \times \cardfont{c}}{\cardfont{a}}$ यांतील एकास-एक आच्छादन आहे. येथे प्रांत प्रत्यक्ष कार्तीय गुणाकार आहे; त्या प्रांताच्या संचांकानेच संचांक गुणाकार परिभाषित होतो.
\end{proof}

अनंत संचांकांच्या बाबतीत $\cardexpo{\cardfont{a}}{\cardfont{b}}$ सहज मोजता यावा, असे आपल्याला \emph{हवे} आहे. या दिशेने पुढील पाऊल उपयोगी आहे:

\begin{prop}\label{sth:card-arithmetic:expotough:cardexpo2reduct}
$2 \leq \cardfont{a} \leq \cardfont{b}$ आणि $\cardfont{b}$ अनंत असेल, तर $\cardexpo{\cardfont{a}}{\cardfont{b}} =
\cardexpo{2}{\cardfont{b}}$.
\end{prop}

\begin{proof}
\begin{align*}
  \cardexpo{2}{\cardfont{b}} &\leq
  \cardexpo{\cardfont{a}}{\cardfont{b}}\text{, कारण $2 \leq \cardfont{a}$}\\
  &\leq \cardexpo{(\cardexpo{2}{\cardfont{a}})}{\cardfont{b}}
  \text{, \ref{sth:card-arithmetic:opps:lem:SizePowerset2Exp} नुसार}\\
  &= \cardexpo{2}{\cardfont{a} \cardtimes \cardfont{b}}
  \text{, \ref{sth:card-arithmetic:expotough:simplecardexpo} नुसार} \\
  &= \cardexpo{2}{\cardfont{b}}
  \text{, \ref{sth:card-arithmetic:simp:cardplustimesmax} नुसार}
\end{align*}
\end{proof}

यापुढे आणखी सुलभीकरण मिळेल अशी फारशी अपेक्षा ठेवू नये: \ref{sth:card-arithmetic:opps:lem:SizePowerset2Exp} नुसार $\cardfont{b} < \cardexpo{2}{\cardfont{b}}$. पण $\cardfont{b} <
\cardfont{a}$ असेल, तर $\cardexpo{\cardfont{a}}{\cardfont{b}}$ बद्दल काय सांगावे, हे यावरून ठरत नाही. अर्थात $\cardfont{b}$ \emph{सांत} असेल, तर काय करायचे ते माहीत आहे.

\begin{prop}
$\cardfont{a}$ अनंत असेल आणि $n \in \omega$ शून्येतर असेल, तर
$\cardexpo{\cardfont{a}}{n} = \cardfont{a}$.
\end{prop}

\begin{proof}
$\cardexpo{\cardfont{a}}{n} = \cardfont{a} \cardtimes \cardfont{a}
\cardtimes \ldots \cardtimes \cardfont{a} = \cardfont{a}$, हे \ref{sth:card-arithmetic:simp:cardplustimesmax} नुसार.
\end{proof}
\noindent
याशिवाय, आणखी काही प्रकरणांत $\cardexpo{\cardfont{a}}{\cardfont{b}}$ चे आकारमान ठरवता येते:

\begin{prop}
$2 \leq \cardfont{b} < \cardfont{a} \leq
\cardexpo{2}{\cardfont{b}}$ आणि $\cardfont{b}$ अनंत असेल, तर
$\cardexpo{\cardfont{a}}{\cardfont{b}} = \cardexpo{2}{\cardfont{b}}$.
\end{prop}

\begin{proof}
$\cardexpo{2}{\cardfont{b}}\leq \cardexpo{\cardfont{a}}{\cardfont{b}}
\leq \cardexpo{(\cardexpo{2}{\cardfont{b}})}{\cardfont{b}} =
\cardexpo{2}{\cardfont{b}\cardtimes\cardfont{b}} =
\cardexpo{2}{\cardfont{b}}$; युक्तिवाद \ref{sth:card-arithmetic:expotough:cardexpo2reduct} प्रमाणे.
\end{proof}
\noindent
पण यापुढे गोष्टी अधिक सूक्ष्म होतात.








\section{सांतत्य गृहीतक}\label{sth:card-arithmetic:ch:sec}

मागील निष्कर्ष योग्यच सूचित करतो की अनंत संचांकांचा $2$ च्या घातांशी, म्हणजे (\ref{sth:card-arithmetic:opps:lem:SizePowerset2Exp} नुसार) घातसंच घेण्याशी, संबंध सरळ असेल, तर संचांक घातांकन अगदी \emph{सोपे} होईल. पण अनंत संचांकांचा घातसंचांशी असा सरळ संबंध \emph{असतोच}, असे आपण गृहीत धरू शकत नाही.

हे उलगडण्यास सुरुवात करण्यासाठी काही सोयीची संकेतरचना मांडू.

\begin{defn}
$\cardsucc{\cardfont{a}}$ हा $\cardfont{a}$ पेक्षा काटेकोरपणे मोठा असलेला लघुतम संचांक, म्हणजे त्याचा उत्तरवर्ती संचांक, असू द्या. पुढील दोन अनंत अनुक्रम परिभाषित करू:
\begin{align*}
	\aleph_{0} &\defis \omega &
	\beth_{0} &\defis \omega\\
	\aleph_{\alpha \ordplus 1} &\defis \cardsucc{(\aleph_{\alpha})} &
	\beth_{\alpha+1} &\defis \cardexpo{2}{\beth_{\alpha}}\\
	\aleph_{\alpha} &\defis \bigcup_{\beta< \alpha} \aleph_{\beta} &
	\beth_{\alpha} &\defis \bigcup_{\beta < \alpha}\beth_{\beta} & \text{जेव्हा $\alpha$ सीमा क्रमसूचक संख्या असते}.
	\end{align*}
\end{defn}

$\cardsucc{\cardfont{a}}$ ची व्याख्या योग्य आहे. कारण \ref{sth:cardinals:classing:lem:NoLargestCardinal} नुसार प्रत्येक संचांक $\cardfont{a}$ साठी $\cardfont{a}$ पेक्षा मोठा संचांक असतो, आणि अनंतपार विगमनामुळे $\cardfont{a}$ पेक्षा मोठा असा \emph{लघुतम} संचांक असल्याची खात्री मिळते. $\cardfont{a}$ यांसारखे संचांक देणाऱ्या या दोन अनुक्रमांची उर्वरित व्याख्या अनंतपार पुनरावर्तनाने मिळते.

कँटरने “$\aleph$” ही संकेतरचना आणली. हे \emph{अलेफ} आहे: हिब्रू वर्णमालेतील पहिले अक्षर आणि हिब्रू भाषेतील “अनंत” या शब्दाचेही पहिले अक्षर. पर्सने “$\beth$” ही संकेतरचना आणली. हे \emph{बेथ} आहे: हिब्रू वर्णमालेतील दुसरे अक्षर.\footnote{डिसेंबर 1900 मध्ये कँटरला लिहिलेल्या पत्रात पर्सने ही संकेतरचना वापरली. दुर्दैवाने त्याच पत्रात $\alpha \geq \omega$ साठी $\beth_\alpha$ अस्तित्वात नसल्याचा चुकीचा युक्तिवादही त्याने दिला.} या संकेतरचना आपल्याला अनंत संचांक देतात.

\begin{prop}
प्रत्येक क्रमसूचक संख्या $\alpha$ साठी $\aleph_\alpha$ आणि $\beth_\alpha$ हे संचांक आहेत.
\end{prop}

\begin{proof}
दोन्ही निष्कर्ष साध्या अनंतपार विगमनाने मिळतात. \ref{sth:cardinals:classing:omegaisacardinal} नुसार $\aleph_0 =
\beth_0 = \omega$ हा संचांक आहे. $\aleph_\alpha$ आणि $\beth_\alpha$ हे दोन्ही संचांक आहेत असे गृहीत धरल्यास, $\aleph_{\alpha+1}$ आणि $\beth_{\alpha+1}$ हे व्याख्येनुसारच संचांक आहेत. शिवाय \ref{sth:cardinals:classing:unioncardinalscardinal} नुसार संचांकांच्या संचाचे युतीकरण हा संचांक असतो.
\end{proof}
\noindent
तसेच प्रत्येक अनंत संचांक हा $\aleph$ च्या स्वरूपाचा असतो:

\begin{prop}
$\cardfont{a}$ हा अनंत संचांक असेल, तर एकमेव $\gamma$ साठी $\cardfont{a} =
\aleph_\gamma$.
\end{prop}

\begin{proof}
अनंत संचांकांवर अनंतपार विगमन वापरू. अलेफ-शून्याच्या व्याख्येमुळे लघुतम अनंत संचांकाचे आधारप्रकरण मिळते. पुढील विगमनात $\cardfont{b} < \cardfont{a}$ असलेल्या प्रत्येक अनंत संचांकासाठी $\cardfont{b} =
\aleph_{\gamma_\cardfont{b}}$ असे गृहीत धरा; खालील युतीकरणांतही अशाच अनंत संचांकांचा विचार केला आहे. एखाद्या $\cardfont{b}$ साठी $\cardfont{a} =
\cardsucc{\cardfont{b}}$ असल्यास $\cardfont{a} =
\cardsucc{(\aleph_{\gamma_\cardfont{b}})}=
\aleph_{\gamma_\cardfont{b}+1}$. जर $\cardfont{a}$ हा कोणत्याही संचांकाचा उत्तरवर्ती नसेल, तर संचांक या क्रमसूचक संख्या असल्यामुळे $\cardfont{a} = \bigcup_{\cardfont{b} < \cardfont{a}} \cardfont{b} =
\bigcup_{\cardfont{b} < \cardfont{a}}{\aleph_{\gamma_\cardfont{b}}}$.
म्हणून $\gamma =
\bigcup_{\cardfont{b} < \cardfont{a}}\gamma_\cardfont{b}$ घेतल्यास $\cardfont{a} = \aleph_\gamma$.
\end{proof}

प्रत्येक अनंत संचांक हा $\aleph$ च्या स्वरूपाचा असल्यामुळे पुढील प्रश्न पडतो: प्रत्येक अनंत संचांक हा $\beth$ च्याही स्वरूपाचा असतो का? असे \emph{असते}, तर अनंत संचांकांचा घातसंच घेण्याच्या क्रियेशी सरळ संबंध असता. खरे तर पुढील विधान मिळाले असते:

\begin{defish}
\emph{सामान्यीकृत सांतत्य गृहीतक} (GCH). प्रत्येक $\alpha$ साठी $\aleph_\alpha = \beth_\alpha$.
\end{defish}

तसेच GCH खरे असेल, तर संचांक घातांकनात बरीच सुलभीकरणे करता येतील. विशेषतः अनंत आधार आणि शून्येतर घातांकांच्या बाबतीत, $\cardfont{b} < \cardfont{a}$ असताना $\cardexpo{\cardfont{a}}{\cardfont{b}}$ चे मूल्य
$\cardfont{a}\leq \cardexpo{\cardfont{a}}{\cardfont{b}} \leq
\cardsucc{\cardfont{a}}$ या मर्यादांत असल्याचे दाखवता येईल. त्यानंतर या दोन शक्यतांपैकी कोणती घडते (म्हणजे $\cardfont{a} = \cardexpo{\cardfont{a}}{\cardfont{b}}$ की $\cardexpo{\cardfont{a}}{\cardfont{b}} =
\cardsucc{\cardfont{a}}$), हे ठरवणाऱ्या नेमक्या अटीही देता येतील.\footnote{ही अट \emph{सहअंतिमतेवर} ठरते.}

पण GCH हे \emph{गृहीतक} आहे, \emph{प्रमेय} नाही. खरे तर गोडेल (1938) यांनी असे सिद्ध केले की $\ZFC$ सुसंगत असेल, तर $\ZFC + \text{GCH}$ सुद्धा सुसंगत आहे. पण नंतर असेही आढळले की $\ZFC$ मध्ये $\lnot$GCH तितक्याच रीतीने जोडता येते. यासाठी GCH चे सर्वांत सोपे, क्षुल्लक नसलेले \emph{विशेषप्रकरण} पाहा:

\begin{defish}
\emph{सांतत्य गृहीतक} (CH). $\aleph_1 = \beth_1$.
\end{defish}

कोहेन (1963) यांनी असे सिद्ध केले की $\ZFC$ सुसंगत असेल, तर $\ZFC
+ \lnot\text{CH}$ सुद्धा सुसंगत आहे. म्हणून सांतत्य गृहीतक हे $\ZFC$ पासून स्वतंत्र आहे.

सांतत्य गृहीतकाला हे नाव पडण्याचे कारण असे की “सांतत्य” हे वास्तव संख्यारेषा $\Real$ हिचे दुसरे नाव आहे. \ref{sth:card-arithmetic:opps:continuumis2aleph0} नुसार $\card{\Real} =
\beth_1$. म्हणून सांतत्य गृहीतक असे सांगते की नैसर्गिक संख्यांचा संचांक $\aleph_0 = \beth_0$ आणि सांतत्याचा संचांक $\beth_1$ यांच्या दरम्यान कोणताही संचांक नाही.

(G)CH हे $\ZFC$ पासून \emph{स्वतंत्र} असल्यामुळे त्यांच्या \emph{सत्यतेविषयी} काय म्हणावे? याविषयी \emph{बरेच} काही सांगता येते. संच उपपत्तीतील अलीकडील अनेक फलदायी संशोधनांचा रोख या प्रश्नांच्या तपासाकडे आहे. पण येथे दोन मुद्दे थोडक्यात अधोरेखित करू.

पहिला मुद्दा: या औपचारिक स्वातंत्र्यनिष्कर्षांवरून GCH किंवा CH यांचे सत्यमूल्य \emph{अनिर्धारित} आहे, असा निष्कर्ष \emph{लगेच} मिळत नाही. अखेर, आपल्याला स्वाभाविक वाटणारी आणखी स्वयंसिद्धके जोडल्यावर हा प्रश्न एका किंवा दुसऱ्या बाजूने सुटू शकेल. खुद्द गोडेल यांनीही हा योग्य प्रतिसाद असल्याचे सुचवले होते.

दुसरा मुद्दा: CH हे $\ZFC$ पासून स्वतंत्र असणे नक्कीच \emph{लक्षणीय} आहे; पण ते शब्दशः \emph{अविश्वसनीय} नाही. मुद्दा एवढाच की $\ZFC$ आपल्याला जे सांगते त्यानुसार संचांकापासून त्याच्या उत्तरवर्ती संचांकाकडे जाण्यासाठी, केवळ घातसंच घेण्यापेक्षा अधिक सूक्ष्म साधन लागू शकते.




हे दोन मुद्दे मांडल्यानंतर, आणखी जाणून घ्यायचे असल्यास या चर्चेत आपापली भूमिका मांडणाऱ्या विविध तत्त्वज्ञ आणि गणितज्ञांच्या लेखनाचा आधार घ्यावा लागेल.\footnote{तरी सुरुवात पॉटर (2004, \S15.6) यांच्या विवेचनापासून करणे उपयुक्त ठरेल.}







\section{$\aleph$-स्थिरबिंदू}\label{sth:card-arithmetic:fix:sec}

\ref{sth:spine:chap} मध्ये आपण असे सुचवले होते की प्रतिस्थापनाला समर्थनाची गरज आहे, कारण त्याच्यामुळे श्रेणीची उंची बरीच मोठी असावी लागते. संचांकांचे काही अंकगणित केल्यानंतर या उंचीचे एक छोटे उदाहरण देता येईल.

स्पष्टच $0 < \aleph_0$, $1 < \aleph_1$, $2 < \aleph_2$\ldots आणि प्रत्येक पावलागणिक आकारमानातील फरक आणखी \emph{वाढतो}. म्हणून प्रत्येक क्रमसूचक संख्या $\kappa$ साठी $\kappa< \aleph_\kappa$ असा तर्क करण्याचा मोह होतो.

पण $\ZFC$ स्वीकारल्यास हा तर्क \emph{खोटा} ठरतो. खरे तर \emph{$\aleph$-स्थिरबिंदू} अस्तित्वात असल्याचे सिद्ध करता येते; म्हणजे $\kappa=\aleph_\kappa$ असलेले संचांक $\kappa$ आहेत.

\begin{prop}\label{sth:card-arithmetic:fix:alephfixed}
एक $\aleph$-स्थिरबिंदू अस्तित्वात आहे.
\end{prop}

\begin{proof}
पुनरावर्तनाने पुढील व्याख्या करा:
\begin{align*}
	\kappa_0 &= 0\\
	\kappa_{n+1} &= \aleph_{\kappa_n}\\
	\kappa&= \bigcup_{n < \omega}\kappa_n
\end{align*}
\ref{sth:cardinals:classing:unioncardinalscardinal} नुसार $\kappa$ हा संचांक आहे. आता:
\[
	\kappa= \bigcup_{n < \omega} \kappa_{n+1} =
	\bigcup_{n < \omega}\aleph_{\kappa_n} =
	\bigcup_{\alpha < \kappa}\aleph_\alpha = \aleph_\kappa
\]






\end{proof}

आत्ताच रचलेल्या नेमक्या $\aleph$-स्थिरबिंदूविषयी बूलोस यांनी एक लेख लिहिला होता. लेखाच्या सुरुवातीला $\kappa$ चे अस्तित्व नोंदवल्यावर ते म्हणाले:
\begin{quote}
	संच उपपत्तीशी यापूर्वी परिचय नसलेल्यांच्या दृष्टीने [$\kappa$ ही] \emph{बरीच मोठी} संख्या आहे. ती इतकी मोठी आहे की, निदान $\kappa$ वस्तू आहेत असे ज्यातील एखादे विधान सांगते, अशा कोणत्याही उपपत्तीच्या सत्यतेविषयी प्रश्न उपस्थित करावा असे मला वाटते.
	(बूलोस 2000, पृ.~257)
\end{quote}
आणि शेवटी त्यांनी आपला लेख पुढील प्रश्नाने संपवला:
\begin{quote}
	कथेच्या सुरुवातीला काहीही असले, तरी इथवर येईपर्यंत चाके नुसतीच फिरत आहेत आणि प्रत्यक्षात जे काही आहे त्याचे वर्णन आपण आता ऐकतच नाही, अशी शंका आपल्याला येते का?
	(बूलोस 2000, पृ.~268)
\end{quote}
प्रत्यक्षात “जे काही आहे” त्याच्या पलीकडे आपण गेलो असू, तर दोषाचे बोट थेट प्रतिस्थापनाकडे दाखवावे लागेल. कारण आपला पुरावा याच स्वयंसिद्धकामुळे चालतो. असे असल्यास, काही दशकांपूर्वी केलेल्या प्रतिस्थापनाचे “कोणतेही अनिष्ट” परिणाम नाहीत या दाव्याचा बूलोस यांना पुनर्विचार करावा लागेल (पाहा \ref{sth:replacement:extrinsic:sec}).

पण $\kappa$ चे अस्तित्व इतके वाईट आहे का? येथे रसेलचा \emph{ट्रिस्ट्रम शॅंडी विरोधाभास} पाहणे उपयुक्त ठरेल. ट्रिस्ट्रम शॅंडी आपल्या जीवनाची नोंद दैनंदिनीत करतो, पण एका दिवसाची नोंद करायला त्याला एक वर्ष लागते. प्रत्येक वर्षागणिक तो अधिकाधिक मागे पडतो: एका वर्षानंतर त्याने केवळ एका दिवसाची नोंद केली आहे आणि नोंद न झालेले 364 दिवस जगला आहे; दोन वर्षांनंतर केवळ दोन दिवसांची नोंद केली आहे आणि नोंद न झालेले 728 दिवस जगला आहे; तीन वर्षांनंतर केवळ तीन दिवसांची नोंद केली आहे आणि नोंद न झालेले 1092 दिवस जगला आहे \dots\footnote{येथे लीप वर्षांचा विचार केलेला नाही.} तरी ट्रिस्ट्रम \emph{अमर} असेल, तर तो प्रत्येक दिवसाची नोंद करू शकेल, कारण आपल्या आयुष्याच्या $n$ व्या वर्षी तो $n$ व्या दिवसाची नोंद करेल. म्हणून “काळाच्या अखेरीस” त्याची दैनंदिनी पूर्ण असेल.

आता, $\kappa$ येईपर्यंत $\alpha$ हा $\aleph_\alpha$ पेक्षा लहान असतो---आणि वाढत्या प्रमाणात, अगदी हतबल करणाऱ्या प्रमाणात लहान असतो---पण त्या बिंदूवर आपण गाठतो आणि $\kappa
= \aleph_\kappa$ मिळते, या विचारापेक्षा वरील गोष्ट इतकी वेगळी का आहे?

हा प्रश्न बाजूला ठेवून, आणि $\ZFC$ स्वीकारतो असे गृहीत धरून, स्थिरबिंदू रचनांविषयीच्या आणखी काही रंजक गोष्टींनी हे प्रकरण संपवू. पुढील तीन निष्कर्षांचा सहज अर्थ असा की श्रेणीची उंची आणि रुंदी समान असलेला एक क्षुल्लक नसलेला टप्पा आहे:

\begin{prop}\label{sth:card-arithmetic:fix:bethfixed}
एक $\beth$-स्थिरबिंदू अस्तित्वात आहे; म्हणजे $\kappa=
\beth_\kappa$ असलेला $\kappa$ आहे.
\end{prop}

\begin{proof}
\ref{sth:card-arithmetic:fix:alephfixed} प्रमाणे, “$\aleph$” ऐवजी “$\beth$” वापरा.
\end{proof}

\begin{prop}\label{sth:card-arithmetic:fix:stagesize}
$\card{V_{\omega+\alpha}} = \beth_{\alpha}$. जर $\omega \ordtimes
\omega \leq \alpha$ असेल, तर $\card{V_\alpha} = \beth_\alpha$.
\end{prop}

\begin{proof}
पहिला दावा साध्या अनंतपार विगमनाने मिळतो. दुसरा दावा यावरून मिळतो की $\omega \ordtimes \omega \leq \alpha$ असेल, तर $\omega + \alpha = \alpha$. हे सिद्ध करण्यासाठी \ref{sth:ord-arithmetic:chap} मधील क्रमसूचक अंकगणिताची तथ्ये वापरू. आधी पाहा की $\omega \ordtimes \omega = \omega \ordtimes (1 \ordplus \omega) =
(\omega \ordtimes 1) \ordplus (\omega\ordtimes\omega) = \omega
\ordplus (\omega \ordtimes \omega)$. आता $\omega \ordtimes \omega
\leq \alpha$ असल्यास, म्हणजे एखाद्या $\beta$ साठी $\alpha = (\omega\ordtimes\omega) \ordplus \beta$ असल्यास, $\omega \ordplus \alpha = \omega \ordplus
((\omega \ordtimes \omega) \ordplus \beta) = (\omega \ordplus (\omega
\ordtimes \omega)) \ordplus \beta = (\omega \ordtimes \omega) \ordplus
\beta = \alpha$.
\end{proof}

\begin{cor}
$\card{V_\kappa} = \kappa$ असलेला $\kappa$ अस्तित्वात आहे.
\end{cor}

\begin{proof}
\ref{sth:card-arithmetic:fix:bethfixed} नुसार $\kappa$ हा एक $\beth$-स्थिरबिंदू असू द्या. स्पष्टच $\omega \ordtimes \omega < \kappa$. म्हणून \ref{sth:card-arithmetic:fix:stagesize} नुसार $\card{V_\kappa} =
\beth_\kappa= \kappa$.
\end{proof}

$V_\kappa$ च्या खाली जितके टप्पे आहेत, तितकेच घटक $V_\kappa$ मध्ये आहेत. त्यामुळे सहज अर्थाने $V_\kappa$ ची रुंदी त्याच्या उंचीइतकी आहे. हे ट्रिस्ट्रम शॅंडीच्या गोष्टीसारखेच आहे: \emph{घातसंच} घेऊन आपण एका टप्प्यापासून पुढच्या टप्प्यावर जातो, आणि प्रत्येक पावलावर श्रेणी \emph{खूप} मोठी होते. पण “अखेरीस”, म्हणजे $\kappa$ टप्प्यावर, श्रेणीची रुंदी तिच्या उंचीला गाठते.

प्रश्न पडू शकतो: \emph{श्रेणीची रुंदी तिच्या उंचीशी किती वेळा जुळते?} उत्तर असे: \emph{जितक्या क्रमसूचक संख्या आहेत, तितक्या वेळा.} पण याचे थोडे स्पष्टीकरण हवे.

पुढीलप्रमाणे $\tau$ हे पद परिभाषित करू. कोणत्याही $A$ साठी:
\begin{align*}
	\tau_0(A) & \defis \cardsucc{\card{A}}\\
	\tau_{n+1}(A) & \defis \beth_{\tau_n(A)}\\
	\tau(A) & \defis \bigcup_{n < \omega}\tau_n(A)
\intertext{\ref{sth:card-arithmetic:fix:bethfixed} प्रमाणे, कोणत्याही $A$ साठी $\tau(A)$ हा
$\beth$-स्थिरबिंदू आहे. सुरुवात उत्तरवर्ती संचांकापासून केल्यामुळे $\card{A} < \tau(A)$ मिळते. म्हणून पुढील पुनरावर्ती व्याख्या पाहा:}
	W_0 &\defis 0\\
	W_{\alpha + 1} & \defis \tau(W_\alpha)\\
	W_\alpha & \defis \bigcup_{\beta < \alpha} W_\beta
	\text{, जेव्हा $\alpha$ सीमा क्रमसूचक संख्या असते}
\end{align*}
ही रचना सर्व क्रमसूचक संख्यांसाठी परिभाषित आहे. सहज अर्थाने $W$ हे क्रमसूचक संख्यांपासून $\beth$-स्थिरबिंदूंकडे जाणारे “एकास-एक फलन” आहे; यातील शून्यव्या मूल्याचा अपवाद स्पष्ट ठेवावा. आणि आधीप्रमाणेच प्रत्येक $\alpha$ साठी $V_{W_\alpha}$ ची रुंदी त्याच्या उंचीइतकी आहे. शून्यव्या टप्प्यावर दोन्ही आकार शून्य आहेत; पुढील टप्प्यांवर हे बेथ-स्थिरबिंदूंच्या गुणधर्मामुळे मिळते.



\chapter{निवड}\label{sth:choice:chap}





\section{प्रस्तावना}\label{sth:choice:intro:sec}

\ref{sth:cardinals:chap}--\ref{sth:card-arithmetic:chap} मध्ये संचांकांना क्रमसूचक संख्या मानून संचांकांची उपपत्ती विकसित केली. ही पद्धत सुक्रमणाच्या स्वयंसिद्धकावर अवलंबून आहे. सुक्रमण हे दुसऱ्या एका तत्त्वाशी---निवडीच्या स्वयंसिद्धकाशी---सममूल्य ठरते; आणि त्या तत्त्वाच्या स्वीकारार्हतेविषयी गंभीर तात्त्विक चर्चा झाली आहे. या प्रकरणातील आपले प्रश्न असे आहेत: हे स्वयंसिद्धक कसे वापरले जाते आणि त्याचे समर्थन करता येते का?







\section{टार्स्की--स्कॉटची युक्ती}\label{sth:choice:tarskiscott:sec}

\ref{sth:cardinals:cardsasords:defcardinalasordinal} मध्ये संचांकांना क्रमसूचक संख्या म्हणून परिभाषित केले. यासाठी सुक्रमणाचे स्वयंसिद्धक गृहीत धरले होते. हीच प्रचलित प्रमाणपद्धत आहे, एवढ्याच कारणासाठी आपण तसे केले.

म्हणून सुक्रमणाशी संबंधित तात्त्विक प्रश्नांची चर्चा करण्याआधी, प्रचलित प्रमाणपद्धत सोडून सुक्रमण गृहीत \emph{न धरता} संचांकांची उपपत्ती विकसित करणे आपल्याला \emph{शक्य} आहे, हे स्पष्ट असणे महत्त्वाचे आहे. \ref{sfr:siz::chap} मध्ये आलेल्या $\cardeq{A}{B}$, $\cardle{A}{B}$ आणि $\cardless{A}{B}$ या व्याख्या तरीही वापरता येतील. फक्त \emph{संचांकाची} नवी संकल्पना लागेल.

$A$ चा संचांक पुढीलप्रमाणे परिभाषित करावा, अशी साधी कल्पना सुचू शकते:
\[
	\Setabs{x}{\cardeq{A}{x}}.
\]
\ref{sth:cardinals:hp:sec} च्या अगदी शेवटी मांडलेल्या फ्रेगेच्या $\# x Fx$ या व्याख्येशी याची तुलना करता येईल. तिथे सूचित केलेल्या कारणांमुळे ही सर्वसाधारण व्याख्या अपयशी ठरते. कोणताही एकघटकी संच $\{\emptyset\}$ शी तुल्यबल असतो. पण श्रेणीच्या प्रत्येक उत्तरवर्ती टप्प्यावर नवे एकघटकी संच तयार होतात (आधीच्या टप्प्याचा एकघटकी संचच पाहा). त्यामुळे अरिक्त आधारसंचाच्या बाबतीत $\Setabs{x}{\cardeq{A}{x}}$ अस्तित्वात नसतो, कारण त्याला प्रतांक असू शकत नाही. रिक्त आधारसंचाचे प्रकरण या आक्षेपाच्या कक्षेत येत नाही.

ही अडचण टाळण्यासाठी टार्स्की आणि स्कॉट यांची युक्ती वापरू:\footnote{आठवण: स्पष्टपणे वेगळे सांगितले नसेल, तर सर्व सूत्रांत प्राचल असू शकतात.}

\begin{defn}[टार्स्की--स्कॉट]
कोणत्याही सूत्र $\phi(x)$ साठी, $\phi(x)$ पूर्ण करणाऱ्या आणि शक्य तितक्या लघुतम प्रतांकाच्या सर्व $x$ चा संच $[ x : \phi(x)] $ असू द्या ($\phi$ पूर्ण करणारी कोणतीही वस्तू नसल्यास $\emptyset$ घ्या).
\end{defn}

ही व्याख्या योग्य आहे हे तपासले पाहिजे. $\ZF$ मध्ये काम करताना \ref{sth:spine:foundation:zfentailsregularity} नुसार प्रत्येक $x$ साठी $\setrank{x}$ अस्तित्वात असतो. आता $\phi$ पूर्ण करणाऱ्या काही वस्तू असतील, तर $(\exists x\subseteq V_\alpha)\phi(x)$ पूर्ण करणारा लघुतम प्रतांक $\alpha$ घेता येतो; त्याच्या लघुतमतेचा अर्थ $(\forall \beta
\in \alpha)(\forall x \subseteq V_\beta)\lnot \phi(x)$ असा आहे. त्यानंतर विभक्तीकरणाने $[x : \phi(x)]$ हा संच $\Setabs{x \in
V_{\alpha+1}}{\phi(x)}$ म्हणून परिभाषित करता येतो.

टार्स्की--स्कॉटच्या युक्तीचे समर्थन केले आहे; आता ती वापरून संचांकाची संकल्पना परिभाषित करता येईल:

\begin{defn}
$A$ चा \textsc{ts}-संचांक म्हणजे $\text{tsc}(A) = [x :
\cardeq{A}{x}]$.
\end{defn}

\textsc{ts}-संचांकाच्या व्याख्येत सुक्रमण वापरलेले नाही. पण ते स्वयंसिद्धक नसतानाही \emph{\textsc{ts}-संचांकांचे} वर्तन \ref{sth:cardinals:cardsasords:defcardinalasordinal} मध्ये परिभाषित केलेल्या \emph{संचांकांसारखेच} असल्याचे दाखवता येते. उदाहरणार्थ, \ref{sth:cardinals:cardsasords:lem:CardinalsBehaveRight} आणि \ref{sth:card-arithmetic:opps:lem:SizePowerset2Exp} हे निष्कर्ष \textsc{ts}-संचांकांच्या भाषेत पुन्हा मांडल्यास, सुक्रमण गृहीत न धरता त्यांचे पुरावे $\ZF$ मध्ये व्यवस्थित चालतात.

याच संदर्भात, टार्स्की--स्कॉटची युक्ती वापरून क्रमसूचक संख्यांची उपपत्तीही विकसित करता येते. $\tuple{A, <}$ हे सुक्रमण असताना $\text{tso}(A, <) = [\tuple{X,
R} : \ordeq{\tuple{A, <}}{\tuple{X, R}}]$ घ्या. संचांक आणि क्रमसूचक संख्या यांच्या या मांडणीविषयी अधिक माहितीसाठी पॉटर (2004, प्रकरणे~9--12) पाहा.







\section{तुलनाक्षमता आणि हार्टोग्सचे पूर्वप्रमेय}\label{sth:choice:hartogs:sec}

ही झाली अनुकूल बाजू. आता प्रतिकूल बाजू पाहू. निवडीचे स्वयंसिद्धक नसल्यास गोष्टी \emph{गुंतागुंतीच्या} होतात. असे का, हे पाहण्यासाठी (हार्टोग्स 1915) यांचा हा सुंदर निष्कर्ष पाहा:

\begin{lem}[\emph{$\ZF$ मध्ये}]\label{sth:choice:hartogs:HartogsLemma}
कोणत्याही संच $A$ साठी $\cardnless{\alpha}{A}$ असलेली क्रमसूचक संख्या $\alpha$ अस्तित्वात आहे.
\end{lem}

\begin{proof}
$B \subseteq A$ आणि $R \subseteq B^2$ असल्यास \ref{sth:ord-arithmetic:using-addition:rankcomputation} नुसार $\tuple{B, R} \subseteq
V_{\setrank{A}+4}$. म्हणून विभक्तीकरण वापरून पुढील संच पाहू:
\[
	C = \Setabs{\tuple{B, R} \in V_{\setrank{A}+5}}{B\subseteq A
	\text{ आणि $\tuple{B, R}$ हे सुक्रमण आहे}}
\]
प्रतिस्थापन आणि \ref{sth:ordinals:ordtype:thmOrdinalRepresentation} वापरून पुढील संच तयार करा:
\[
	\alpha = \Setabs{\ordtype{B, R}}{\tuple{B, R} \in C}.
\]
\ref{sth:ordinals:basic:corordtransitiveord} नुसार $\alpha$ ही क्रमसूचक संख्या आहे, कारण तो क्रमसूचक संख्यांचा संक्रमक संच आहे. कारण $\gamma \in \beta \in \alpha$ असेल, तर एखाद्या संचासाठी $\beta = \ordtype{B, R}$, जेथे $B \subseteq A$, आणि मग \ref{sth:ordinals:iso:wellordinitialsegment} नुसार एखाद्या $b
\in B$ साठी $\gamma = \ordtype{B_b, R_b}$. त्यामुळे $\gamma \in \alpha$.


विरोधसिद्धीसाठी एकास-एक फलन $f \colon \alpha \to A$ अस्तित्वात आहे असे समजा. मग पुढील व्याख्या करा:
\begin{align*}
	B &= \ran{f}\\
	R &= \Setabs{\tuple{f(\xi), f(\eta)} \in A \times A}{\xi \in \eta \land \xi \in \alpha \land \eta \in \alpha}.
\end{align*}
येथे अंतर्गत निर्देशांक फलनाच्या प्रांतात घेतले आहेत. स्पष्टच $\alpha = \ordtype{B, R}$ आणि $\tuple{B, R} \in C$. म्हणून $\alpha
\in \alpha$, ही विसंगती मिळते.
\end{proof}

यावरून एक सखोल निष्कर्ष मिळतो:

\begin{thm}[\emph{$\ZF$ मध्ये}]
पुढील विधाने सममूल्य आहेत:
\begin{enumerate}
	\item\label{sth:choice:hartogs:equivwo} सुक्रमणाचे स्वयंसिद्धक
	\item\label{sth:choice:hartogs:equivcompare} कोणत्याही संच $A$ आणि $B$ साठी $\cardle{A}{B}$ किंवा
	$\cardle{B}{A}$
\end{enumerate}
\end{thm}

\begin{proof}
\emph{\ref{sth:choice:hartogs:equivwo} $\Rightarrow$ \ref{sth:choice:hartogs:equivcompare}.} $A$ आणि $B$ निश्चित करा. \ref{sth:choice:hartogs:equivwo} वापरल्यास $\tuple{A, R}$ आणि $\tuple{B, S}$ ही सुक्रमणे मिळतात. \ref{sth:ordinals:ordtype:thmOrdinalRepresentation} वापरून $f \colon
\alpha \to \tuple{A, R}$ आणि $g \colon \beta \to \tuple{B, S}$ ही समरूपणे असू द्या. \ref{sth:ordinals:basic:ordinalsaresubsets} नुसार $\alpha \subseteq \beta$ किंवा $\beta \subseteq \alpha$. $\alpha \subseteq \beta$ असेल, तर $\comp{f^{-1}}{g} \colon A \to B$ हे एकास-एक फलन आहे आणि म्हणून
$\cardle{A}{B}$; त्याचप्रमाणे $\beta \subseteq \alpha$ असेल, तर $\cardle{B}{A}$.

\emph{\ref{sth:choice:hartogs:equivcompare} $\Rightarrow$ \ref{sth:choice:hartogs:equivwo}.} $A$ निश्चित करा. \ref{sth:choice:hartogs:HartogsLemma} नुसार $\cardnless{\beta}{A}$ असलेली क्रमसूचक संख्या $\beta$ आहे. \ref{sth:choice:hartogs:equivcompare} वापरल्यास $\cardle{A}{\beta}$ मिळते. म्हणून एकास-एक फलन $f \colon A \to
\beta$ अस्तित्वात आहे, आणि $\Setabs{\tuple{a, b} \in A \times A}{f(a)
\in f(b)}$ हा क्रम परिभाषित करून त्या फलनाने $A$ च्या घटकांचे सुक्रमण करता येते.
\end{proof}
\noindent
याचे तात्काळ फलित असे: सुक्रमण अपयशी ठरल्यास काही संच त्यांच्या आकारमानाच्या दृष्टीने \emph{अक्षरशः अतुलनीय} असतात. म्हणून सुक्रमण अपयशी ठरल्यास अनंतपार संचांकांचे अंकगणित गुंतागुंतीचे होते. उदाहरणार्थ, $A$ आणि $B$ अनंत असल्यास $\cardeq{\cardeq{A \disjointsum B}{A \times
B}}{M}$, जेथे $M$ हा $A$ आणि $B$ यांपैकी मोठा संच आहे, ही कल्पना सोडावी लागेल (पाहा \ref{sth:card-arithmetic:simp:cardplustimesmax}). अडचण साधी आहे: $A$ आणि $B$ यांच्या आकारमानांची \emph{तुलनाच} करता येत नसेल, तर कोणता मोठा आहे असा प्रश्न विचारण्याला अर्थ नाही.







\section{सुक्रमणाचा प्रश्न}\label{sth:choice:woproblem:sec}

सुक्रमण स्वीकारतो की नाही, यावर बऱ्याच गोष्टी अवलंबून आहेत. पण या तत्त्वाची चर्चा मुख्यतः त्याच्याशी सममूल्य असलेल्या निवडीच्या स्वयंसिद्धकावर केंद्रित झाली आहे. म्हणून आता त्याकडे वळू (आणि ही सममूल्यता सिद्ध करू).

1883 मध्ये कँटरने सुक्रमणाच्या स्वयंसिद्धकाला पाठिंबा दर्शवला. त्याच्या मते तो “विचाराचा असा नियम आहे की जो मूलभूत, फलितांनी समृद्ध, आणि विशेषतः त्याच्या सार्वत्रिक युक्ततेमुळे उल्लेखनीय वाटतो” (पॉटर 2004, पृ.~243 मध्ये उद्धृत). पण शेवटी या “स्वयंसिद्धकाला” पुराव्याची गरज आहे, असे कँटरचे मत झाले. गणितज्ञ समुदायाचेही तसेच मत झाले.

1904 मध्ये झर्मेलोने या प्रश्नाचे “समाधान” केले. त्याचे समाधान स्पष्ट करण्यासाठी काही व्याख्या लागतील.

\begin{defn}
प्रत्येक $x \in \dom{f}$ साठी $f(x) \in x$ असेल तेव्हा आणि केवळ तेव्हाच फलन $f$ हे \emph{निवडफलन} असते. $f$ हे निवडफलन असून $\dom{f} = A \setminus \{\emptyset\}$ असेल तेव्हा आणि केवळ तेव्हाच $f$ हे \emph{$A$ साठीचे निवडफलन} आहे, असे म्हणतो.
\end{defn}

सहज अर्थाने, प्रत्येक अरिक्त संच $x \in A$ साठी, $A$ चे निवडफलन त्या $x$ मधून विशिष्ट घटक $f(x)$ \emph{निवडते}. मग निवडीचे स्वयंसिद्धक असे आहे:

\begin{axiom}[निवड]
	प्रत्येक संचाला निवडफलन असते.
\end{axiom}

निवडीवरून सुक्रमण मिळते आणि उलटही मिळते, हे झर्मेलोने दाखवले:

\begin{thm}[$\ZF$ मध्ये]\label{sth:choice:woproblem:thmwochoice}
सुक्रमण आणि निवड सममूल्य आहेत.
\end{thm}

\begin{proof}
\emph{डावीकडून उजवीकडे.} $A$ हा संचांचा संच असू द्या. युतीकरणाच्या स्वयंसिद्धकामुळे $\bigcup A$ अस्तित्वात असतो; म्हणून सुक्रमणामुळे $\bigcup A$ ला सुक्रमित करणारा $<$ आहे. आता अरिक्त सदस्यांसाठी $f(x) = \text{$x$ चा $<$-लघुतम घटक}$ असे घ्या. हे $A$ साठीचे निवडफलन आहे.

\emph{उजवीकडून डावीकडे.} $A$ निश्चित करा. रिक्त संचाचे सुक्रमण रिक्त क्रमाने मिळते; म्हणून पुढे आधारसंच अरिक्त मानू. निवडीमुळे $\Pow{A} \setminus \{\emptyset\}$ साठी निवडफलन $f$ आहे. अनंतपार पुनरावर्तन वापरून पुढील फलनाची रचना करू:
\begin{align*}
	g(0) &= f(A)\\
	g(\alpha) &=
		\begin{cases}
			\text{थांबा!{}} &\text{जर }A = \funimage{g}{\alpha}\\
			f(A \setminus \funimage{g}{\alpha}) & \text{अन्यथा}\\
		\end{cases}
\end{align*}
“थांबा!” ही सूचना अधिक लांबलचक व्याख्येचे संक्षिप्त रूप आहे. म्हणजे प्रथम $A =
\funimage{g}{\alpha}$ मिळाल्यावर प्रत्येक $\delta \geq \alpha$ साठी $g(\delta) = A$ घ्या. ही स्थिर पुढील मूल्ये आधीच्या निवडी बदलत नाहीत; पुराव्यातील एकास-एकपणाचा दावा थांबण्यापूर्वीच्या भागापुरता आहे. सुरुवातीला, याने केवळ एक \emph{पद} परिभाषित होते एवढीच खात्री आहे (\ref{sth:spine:recursion:transrecursionschema} नंतरच्या नोंदी पाहा); पण शेवटी थांबण्यापूर्वीचा भाग एका संच-प्रांतावर मर्यादित होतो आणि त्यामुळे आपल्याला फलन मिळते, हे दाखवू.

$f$ हे निवडफलन असल्यामुळे प्रत्येक $\alpha$ साठी (थांबण्यापूर्वी) $g(\alpha) =
f(A \setminus \funimage{g}{\alpha}) \in A \setminus
\funimage{g}{\alpha}$; म्हणजे $g(\alpha) \notin \funimage{g}{\alpha}$.
म्हणून $g(\alpha) = g(\beta)$ असल्यास $g(\beta) \notin
\funimage{g}{\alpha}$, म्हणजे $\beta \notin \alpha$; आणि त्याचप्रमाणे $\alpha \notin \beta$. म्हणून त्रिपर्यायाने $\alpha = \beta$. त्यामुळे या भागावर $g$ हे एकास-एक आहे.

आता पाहा की आपण थांबतो!{}; म्हणजे $A = g[\alpha]$ असलेली एक लघुतम क्रमसूचक संख्या $\alpha$ आहे. तसे नसल्यास $g$ हे एकास-एक असल्याने प्रत्येक क्रमसूचक संख्या $\alpha$ साठी $\cardle{\alpha}{A}$ मिळेल, जे \ref{sth:choice:hartogs:HartogsLemma} शी विसंगत आहे. त्यामुळे थांबण्यापूर्वीच्या भागाला मर्यादित केल्यावर $\ran{g} = A$ सुद्धा मिळते.

ही तथ्ये एकत्र केल्यावर, एखाद्या क्रमसूचक संख्येपासून $A$ पर्यंत $g$ हे एकास-एक आच्छादन आहे; येथे फक्त थांबण्यापूर्वीचा भाग घेतला आहे. आता $g$ वापरून $A$ चे सुक्रमण करता येते.
\end{proof}

म्हणून सुक्रमण आणि निवड हे एकत्रच स्वीकारले किंवा नाकारले जातात. पण प्रश्न उरतो: त्यांचा स्वीकार करावा की नकार?








\section{गणनीय निवड}\label{sth:choice:countablechoice:sec}

निवड किंवा सुक्रमण अजिबात न वापरता पुढील निष्कर्ष सहज सिद्ध करता येतो:

\begin{lem}[$\Zminus$ मध्ये]
प्रत्येक सांत संचाला निवडफलन असते.
\end{lem}

\begin{proof}
$a = \{b_1, \ldots, b_n\}$ असू द्या; सदस्यांची पुनरावृत्ती करू नका. सुलभतेसाठी प्रत्येक $b_i
\neq \emptyset$ आहे असे समजा. मग $c_1 \in b_1, \ldots, c_n \in b_n$ असलेल्या वस्तू $c_1, \ldots, c_n$ आहेत. आता \ref{sth:z:pairs:prop:pairsconsequences} नुसार $\{\langle b_1,
c_1\rangle , \ldots, \langle b_n, c_n\rangle\}$ हा संच अस्तित्वात आहे; आणि हे $a$ साठीचे निवडफलन आहे.
\end{proof}

पण अनंत संचांचा विचार केल्याबरोबर गोष्टी अधिक अस्पष्ट होतात. उदाहरणार्थ, वरील विधानाचा हा “किमान” विस्तार पाहा:

\begin{defish}
\emph{गणनीय निवड.} प्रत्येक \emph{गणनीय} संचाला निवडफलन असते.
\end{defish}

हे निवडीचे एक विशेषप्रकरण आहे. हे तत्त्व वापरत असल्याची स्पष्ट जाणीव नसतानाही ते बऱ्याचदा वापरले गेले होते, असे दिसते. पुढील दोन सुंदर उदाहरणे पाहा.\footnote{ही उदाहरणे पॉटर (2004, \S9.4) आणि लूका इन्कुर्वाती यांच्याकडून घेतली आहेत.}

\begin{ex}
एक स्वाभाविक विचार असा: कोणत्याही संच $A$ साठी $\cardle{\omega}{A}$ किंवा एखाद्या $n \in \omega$ साठी $\cardeq{A}{n}$. प्रत्येक संच सांत किंवा अनंत असतो, ही सहज कल्पना मांडण्याचा हा एक मार्ग आहे. कँटर आणि इतर अनेक गणितज्ञांनी पुराव्याशिवाय हा दावा केला. आपण सावध असल्यामुळे \ref{sth:cardinals:classing:generalinfinitycharacter} मध्ये तो सिद्ध केला. पण त्या पुराव्यात कोणत्याही संच $A$ चे सुक्रमण करता येते आणि म्हणून $\card{A}$ अस्तित्वात असल्याची खात्री आहे, असे मानून आपण $\ZFC$ मध्ये काम केले. म्हणजे आपण निवड स्पष्टपणे गृहीत धरली होती.

खरे तर डेडेकिंड (1888) यांनी या दाव्याचा स्वतःचा पुरावा पुढीलप्रमाणे दिला होता:

\begin{thm}[$\Zminus + \text{गणनीय निवड}$ मध्ये]
कोणत्याही $A$ साठी $\cardle{\omega}{A}$ किंवा एखाद्या $n \in \omega$ साठी $\cardeq{A}{n}$.
\end{thm}

\begin{proof}
प्रत्येक $n \in \omega$ साठी $\cardneq{A}{n}$ असे समजा. मग विशेषतः प्रत्येक $n < \omega$ साठी नेमके $\cardexpo{2}{n}$ घटक असलेला उपसंच $A_n \subseteq A$ आहे. $A_0, A_1, A_2,
\ldots$ हा अनुक्रम वापरून प्रत्येक $n$ साठी पुढील व्याख्या करू:
\[
	B_n = A_n \setminus \bigcup_{i < n} A_i.
\]
शून्यव्या निर्देशांकासाठी आधीचे युतीकरण रिक्त आहे आणि पहिला उपसंच अरिक्त आहे. धन निर्देशांकांसाठी पुढील मोजणी पाहा:
\begin{align*}
	\card{\bigcup_{i < n}A_i}
	&\leq \card{A_0} + \card{A_1} + \ldots + \card{A_{n-1}}\\
	&=1 + 2 + \ldots + 2^{n-1}\\
	& = 2^n - 1\\
	& < 2^n = \card{A_n}
\end{align*}
म्हणून प्रत्येक $B_n$ मध्ये किमान एक घटक $c_n$ आहे. शिवाय $B_n$ हे संच जोडीजोडीने वियुक्त आहेत; त्यामुळे $c_n = c_m$ असल्यास $n = m$. पण प्रत्येक $c_n
\in A$. म्हणून $f(n) = c_n$ हे फलन $\omega \to A$ असे एकास-एक फलन आहे.
\end{proof}
\noindent
गणनीय निवड वापरली असल्याची डेडेकिंडने नोंद केली नाही. पण \emph{तुम्हाला} तिचा वापर दिसला का? पुन्हा पाहा. (खरेच: \emph{पुन्हा पाहा}.)

या पुराव्यात गणनीय निवड दोनदा वापरली. $A_0$, $A_1$, $A_2$, \dots\@ या संचांचा अनुक्रम मिळवण्यासाठी ती एकदा वापरली. त्यानंतर प्रत्येक $B_n$ मधून $c_n$ हे घटक निवडण्यासाठी ती पुन्हा वापरली. शिवाय येथे कोणत्याही प्रकारची निवड न वापरता हा निष्कर्ष सिद्ध करता येत नाही. कोहेन (1966, पृ.~138) यांनी निवडीचे कोणतेही रूप नसल्यास हा निष्कर्ष अपयशी ठरतो, असे सिद्ध केले. म्हणजे $\omega$ शी \emph{अतुलनीय} असलेले संच आहेत, ही बाब $\ZF$ शी सुसंगत आहे.
\end{ex}

\begin{ex}
1878 मध्ये कँटरने म्हटले होते की गणनीय संचांचे गणनीय युतीकरण गणनीय असते. त्याने पुरावा दिला नाही; कदाचित तो पुरावा स्वयंस्पष्ट आहे, असे त्याला वाटले असावे. आपण सावध असल्यामुळे \ref{sth:card-arithmetic:simp:kappaunionkappasize} मध्ये या निष्कर्षाचे अधिक सामान्य रूप सिद्ध केले. पण आपल्या पुराव्यात निवड स्पष्टपणे गृहीत धरली होती. आणि कमी सामान्य निष्कर्षाचा खालील पुरावाही गणनीय निवड वापरतो.

\begin{thm}[$\Zminus + \text{गणनीय निवड}$ मध्ये]
प्रत्येक $n \in \omega$ साठी $A_n$ गणनीय असेल, तर $\bigcup_{n <
\omega} A_n$ गणनीय असतो.
\end{thm}

\begin{proof}
सामान्यतेला बाधा न आणता प्रत्येक $A_n \neq \emptyset$ आहे असे मानता येते. मग प्रत्येक $n \in \omega$ साठी आच्छादन $f_n \colon \omega
\to A_n$ आहे. $f(m, n) = f_n(m)$ घेऊन $f \colon \omega \times \omega \to \bigcup_{n <
\omega} A_n$ परिभाषित करा. $\omega
\times \omega$ गणनीय आहे (\ref{sfr:siz:zigzag:natsquaredenumerable}) आणि $f$ हे आच्छादन आहे, म्हणून निष्कर्ष मिळतो.
\end{proof}
\noindent
गणनीय निवडीचा वापर दिसला का? $f_0$, $f_1$, $f_2$, \dots\footnote{\ref{sth:card-arithmetic:simp:kappaunionkappasize} मध्ये “प्रत्येक $\beta \in \cardfont{a}$ साठी एकास-एक फलन $f_\beta$ निश्चित करा” अशी सूचना देताना निवडीचा असाच वापर झाला होता.} हा फलनांचा अनुक्रम निवडण्यासाठी ती वापरली आहे. पुन्हा, कोणतेही निवडतत्त्व नसल्यास हा निष्कर्ष अपयशी ठरू शकतो. विशेषतः फेफरमन आणि लेव्ही (1963) यांनी असे सिद्ध केले की $\ZF$ सुसंगत असेल, तर वास्तव संख्यांचा संच $\Real$ हा गणनीय संचांचे गणनीय युतीकरण आहे, हे विधानही $\ZF$ शी सुसंगत आहे. पण याच मुद्द्याचे रसेल यांनी केलेले अधिक मजेशीर वर्णन पाहा:
\begin{quote}
  एक कोट्यधीश जेव्हा बुटांची जोडी विकत घेत असे, तेव्हाच मोज्यांचीही जोडी घेत असे, इतर वेळी कधीच नाही. दोन्ही गोष्टी विकत घेण्याची त्याला इतकी हौस होती की शेवटी त्याच्याकडे बुटांच्या $\aleph_0$ जोड्या आणि मोज्यांच्या $\aleph_0$ जोड्या झाल्या\dots\@ बुटांत उजवा आणि डावा असा फरक करता येतो. त्यामुळे प्रत्येक जोडीतून एक बूट निवडता येतो: सर्व उजवे किंवा सर्व डावे बूट निवडू शकतो. पण मोज्यांच्या बाबतीत असे कोणतेही निवडतत्त्व सुचत नाही; आणि गुणक स्वयंसिद्धक [म्हणजे प्रत्यक्षात निवड] गृहीत धरल्याशिवाय, प्रत्येक जोडीतून एक मोजा असलेला वर्ग आहे याची खात्री देता येत नाही.
  (रसेल 1919, पृ.~126)
\end{quote}
म्हणजे मोज्यांच्या गणनीय संख्येने जोड्या असतील, तर मोजेही फक्त गणनीय संख्येने असतात, हे सिद्ध करण्यासाठी काही प्रकारची निवड लागते. गणनीय निवड अशा युतीकरणाची गणनीयता सिद्ध करण्यास पुरेशी आहे; पण निवडतत्त्वांचा अभाव असताना गणनीय संचांचे गणनीय युतीकरण गणनीय असण्यात अपयश येऊ शकते. येथे गणनीय निवडीशी उलटी सममूल्यता सिद्ध केलेली नाही.
\end{ex}

यातून धडा असा मिळतो की वापरणाऱ्यांना फारशी जाणीव नसतानाही गणनीय निवड वारंवार वापरली गेली. तात्त्विक प्रश्न असा आहे: गणनीय निवडीचे \emph{समर्थन} कसे करावे?

सहज समर्थनाचा एखादा प्रयत्न सांत वेळेत अनंत क्रियांचा क्रम पूर्ण करण्याच्या कल्पनेचा आधार घेऊ शकतो. पहिली निवड $\nicefrac{1}{2}$ मिनिटात, दुसरी $\nicefrac{1}{4}$ मिनिटात, \dots, आणि $n$ वी निवड $\nicefrac{1}{2^n}$ मिनिटात, \dots\@ केली असे समजा. मग $1$ मिनिटात निवडींचा $\omega$-अनुक्रम पूर्ण करून निवडफलन परिभाषित केले असेल.

पण असा विचारप्रयोग खरोखर काय सांगू शकतो? प्रथम, “निवड करणे” ही कल्पना अगदी शब्दशः घेतल्यावर तो अवलंबून आहे. दुसरे म्हणजे त्यात गणित सत्तामीमांसक शक्यतेशी बांधले जात असल्यासारखे दिसते.

अधिक महत्त्वाचे असे की यामुळे \emph{केवळ} गणनीय निवडीऐवजी निवडीच्या \emph{पूर्ण रूपाचे} समर्थन मिळणार नाही. कारण \emph{प्रत्येक} संचाला निवडफलन हवे असेल, तर “कोणत्याही क्रमसूचक लांबीचा अनंत क्रियाक्रम” पूर्ण करता आला पाहिजे. स्पष्ट सांगायचे तर ही कल्पना हास्यास्पद आहे.







\section{निवडीविषयीचे आंतरिक विचार}\label{sth:choice:justifications:sec}

मग व्यापक प्रश्न असा आहे की सुक्रमणाचे, निवडीचे किंवा सर्व संचांच्या आकारमानांच्या तुलनाक्षमतेचे---यांपैकी कोणतेही म्हणा---समर्थन करता येते का?

पुढे \emph{आंतरिक} समर्थनाचा एक प्रयत्न पाहू. \ref{sth:z:story:sec} मध्ये श्रेणीविषयी काही तत्त्वे मांडली होती. त्यांपैकी एक पुन्हा सांगण्याजोगे आहे:
\begin{enumerate}
	\item[] \stagesacc. कोणत्याही टप्पा $S$ साठी, आणि $S$ च्या \emph{आधी} तयार झालेल्या कोणत्याही संचांसाठी: नेमके ते संच घटक असलेला संच $S$ या टप्प्यावर तयार होतो. $S$ या टप्प्यावर दुसरे काहीही तयार होत नाही.
\end{enumerate}
खरे तर या तत्त्वाद्वारे किंवा यासारख्या तत्त्वाद्वारे निवडीच्या स्वयंसिद्धकाचे समर्थन करता येते, असे अनेक लेखकांनी सुचवले आहे. या पद्धतीचे थोडक्यात स्पष्टीकरण देऊ.

निवडीशी सममूल्य असलेले \emph{आणखी एक} तत्त्व देणाऱ्या साध्या निष्कर्षापासून सुरुवात करू:

\begin{thm}[$\ZF$ मध्ये]\label{sth:choice:justifications:choiceset}
निवड पुढील तत्त्वाशी सममूल्य आहे. $A$ चे घटक वियुक्त आणि अरिक्त असतील, तर प्रत्येक $x \in A$ साठी $C
\cap x$ एकघटकी असलेला $C$ अस्तित्वात आहे. (अशा $C$ ला $A$ साठीचा {निवडसंच} म्हणतो.)
\end{thm}

या निष्कर्षाचा पुरावा सरळ आहे आणि वाचकासाठी सराव म्हणून सोडला आहे.

\begin{prob}
\ref{sth:choice:justifications:choiceset} सिद्ध करा. अडचण आल्यास (पॉटर 2004, पृ.~242--3) मध्ये पुरावा पाहता येईल.
\end{prob}

मुख्य मुद्दा असा की $A$ साठीच्या निवडसंचातून आधारकुटुंबाच्या युतीकरणाबाहेरील घटक वगळल्यावर तो $A$ साठीच्या निवडफलनाची प्रतिमा ठरतो. म्हणून निवडीचे समर्थन करण्यासाठी निवडसंचांच्या अस्तित्वाच्या भाषेतील तिच्या सममूल्य मांडणीचे समर्थन करण्याचा प्रयत्न करता येईल. आता नेमके तेच करू.

$A$ चे घटक वियुक्त आणि अरिक्त असू द्या. \stageshier{} नुसार (पाहा \ref{sth:z:story:sec}) $A$ कोणत्या तरी टप्पा $S$ वर तयार होतो. $\bigcup A$ चे सर्व घटक $S$ टप्प्याच्या आधी उपलब्ध असतात. आता \stagesacc{} नुसार, $S$ च्या आधी तयार झालेल्या \emph{कोणत्याही} संचांचे नेमके तेच घटक असलेला संच तयार होतो. दुसऱ्या शब्दांत: आधी उपलब्ध झालेल्या संचांचे प्रत्येक \emph{शक्य} संकलन $S$ वर अस्तित्वात असते. पण $A$ साठीच्या निवडसंचात एकत्र करता येतील अशा वस्तू निवडणे नक्कीच \emph{शक्य} आहे, असा या समर्थनाचा दावा आहे; ते $\bigcup A$ च्या अगदी विशिष्ट उपसंचातच येतील. म्हणून आवश्यक असा एखादा निवडसंच अस्तित्वात आहे, असे या युक्तिवादातून म्हटले जाते.

आंतरिक आधारावर निवडीचे समर्थन करण्याचा हा \emph{अतिशय} थोडक्यात मांडलेला प्रयत्न आहे. ही कल्पना पुढे नेण्यासाठी पॉटर यांनी (2004, \S14.8) केलेली सुबक मांडणी वाचावी.







\section{बानाख--टार्स्की विरोधाभास}\label{sth:choice:banach:sec}

बूलोस यांनी प्रतिस्थापनाचे समर्थन करण्याचा प्रयत्न केला तसा, \emph{बाह्य} विचारांचा आधार घेऊन निवडीचे समर्थन करण्याचाही प्रयत्न करता येईल (पाहा \ref{sth:replacement:extrinsic:sec}). अखेर निवड स्वीकारल्याने अनेक इष्ट परिणाम मिळतात: प्रत्येक संचांकाची तुलना करता येते, प्रत्येक संचाचे सुक्रमण करता येते, संचांकांना विशिष्ट प्रकारच्या क्रमसूचक संख्या मानता येते, इत्यादी.

पण कधी कधी निवडीचे \emph{अनिष्ट} परिणाम आहेत, असा दावा केला जातो. मुख्यतः (बानाख आणि टार्स्की 1924) यांच्या एका निष्कर्षामुळे हा दावा केला जातो.

\begin{thm}[बानाख--टार्स्की विरोधाभास ($\ZFC$ मध्ये)]
त्रिमितीय अवकाशातील धन त्रिज्येच्या कोणत्याही घनगोलाचे सांत संख्येने तुकडे करता येतात; घूर्णन आणि स्थानांतरणाने त्यांची पुनर्जुळणी करून त्या घनगोलाच्या दोन प्रती तयार करता येतात.
\end{thm}
\noindent
पहिल्या दृष्टीने हे थोडे थक्क करणारे आहे. या दोन घनगोलांचे घनफळ मूळ घनगोलाच्या \emph{दुप्पट} आहे. पण दृढ गती---घूर्णन आणि स्थानांतरण---घनफळ बदलत नाहीत. म्हणून बानाख--टार्स्कीमुळे जादूने नवे द्रव्य अस्तित्वात आणता येते, असे दिसते.

यापुढे आणखी आश्चर्य आहे.\footnote{पाहा वॅगन (2016, प्रमेय 3.12).} यासारखा युक्तिवाद दाखवतो की वाटाण्याचे सांत संख्येने तुकडे करून घूर्णन आणि स्थानांतरणाने त्यांची पुनर्जुळणी केल्यास बिग बेनच्या आकाराची आणि आकारमानाची वस्तू तयार करता येते.

पण केवळ $\ZF$ मध्ये यापैकी काहीही सिद्ध होत नाही.\footnote{तरी निवडीपेक्षा काटेकोरपणे दुर्बल तत्त्वांनी बानाख--टार्स्की सिद्ध करता येतो; पाहा वॅगन (2016, पृ.~303).} म्हणून निवड नाकारावी की या “विरोधाभासासह” राहायला शिकावे, हा निर्णय घ्यावा लागतो.

या “विरोधाभासासह” राहायला शिकावे, असे आपण सुचवू. खरे तर हा फारसा विरोधाभास आहे, असे आम्हाला वाटत नाही. विशेषतः पुढील निष्कर्षांपेक्षा तो कमी किंवा अधिक विरोधाभासी का आहे, हे दिसत नाही:\footnote{पॉटर (2004, पृ.~276--7), वेस्टन (2003, पृ.~16) आणि वॅगन (2016, 31, 308--9) इतर उदाहरणांनी असेच मुद्दे मांडतात.}
\begin{enumerate}
	\item $(0,1)$ या अंतरालात जितके बिंदू आहेत, तितकेच $\Real$ मध्ये आहेत.
	\\\emph{पुरावा}: $\tan(\pi(r-\nicefrac{1}{2}))$ पाहा.
	\item रेषेत जितके बिंदू आहेत, तितकेच चौरसात आहेत.
	\\पाहा \ref{his:set:pathology:sec} आणि \ref{his:set:cantorplane:sec}.
	\item अवकाशभरण करणाऱ्या वक्ररेषा आहेत.
	\\पाहा \ref{his:set:pathology:sec} आणि \ref{his:set:hilbertcurve:sec}.
\end{enumerate}
या तिन्ही निष्कर्षांना निवडीची गरज नाही. आज आपण त्यांना गणिताचे आश्चर्यकारक, सुंदर भाग मानतो. बानाख--टार्स्की विरोधाभासाविषयीही तीच भूमिका घ्यावी, असे वाटू शकते.

येथे एक तांत्रिक निरीक्षण आवश्यक आहे; पण त्यासाठी शांतपणे विचार करणे पुरेसे आहे. दृढ गती घनफळ जतन करतात. त्यामुळे घनगोलाचे जे पाच\footnote{विरोधाभासाचे विधान “सांत संख्येने तुकडे” असे केले होते. खरे तर रॉबिन्सन (1947) यांनी \emph{पाच} तुकड्यांत (त्याहून कमी तुकड्यांत नाही) हे विघटन करता येते, असे सिद्ध केले. पुराव्यासाठी वॅगन (2016, पृ.~66--7) पाहा.} तुकडे केले जातात, ते सर्व \emph{मापनीय} असू शकत नाहीत. साधारणपणे, या स्वतंत्र तुकड्यांना घनफळ देण्यात अर्थ नाही. त्यांना चित्रात दाखवता न येणारे बिंदूंचे “अनंत विखुरणे” समजावे. अशा “अनंत विखुरलेल्या” संचांची कल्पना करणे विचित्र वाटू शकते. पण आधी उपलब्ध झालेल्या संचांची \emph{सर्व शक्य} संकलने तयार करावीत, या \stagesacc{} मध्ये व्यक्त झालेल्या निर्देशावरून त्यांचे अस्तित्व मिळत असल्यासारखे दिसते.

हे पटत नसेल, तर बानाख--टार्स्की विरोधाभास स्वीकारण्याच्या बाजूने शेवटचा एक बाह्य युक्तिवाद पाहा. त्यावरून गणितातील उत्कृष्ट विनोद लगेच मिळतो:
\begin{enumerate}
	\item[] \emph{प्रश्न}. “बानाख--टार्स्की”च्या अक्षरांची पुनर्रचना केल्यावर काय मिळते?
	\item[] \emph{उत्तर}. “बानाख--टार्स्की बानाख--टार्स्की”.
\end{enumerate}







\section{परिशिष्ट: वितालीचा विरोधाभास}\label{sth:choice:vitali:sec}

बानाख--टार्स्कीची रचना स्वीकारार्ह आहे की नाही, हे नीट समजून घेण्यासाठी तिचा \emph{पुरावा} तपासला पाहिजे. दुर्दैवाने त्यासाठी येथे मांडता येईल त्याहून बरेच अधिक बीजगणित लागेल. तरी पुढील निष्कर्षावर लक्ष केंद्रित करून बानाख--टार्स्कीच्या पुराव्यावर प्रकाश टाकतील अशा काही थोडक्यात नोंदी देता येतील.\footnote{अधिक पूर्ण विवेचनासाठी (वेस्टन 2003) किंवा (वॅगन 2016) पाहा.}

\begin{thm}[वितालीचा विरोधाभास ($\ZFC$ मध्ये)]\label{sth:choice:vitali:vitaliparadox}
धन त्रिज्येच्या कोणत्याही वर्तुळाचे गणनीय संख्येने तुकडे करता येतात; घूर्णन आणि स्थानांतरणाने त्यांची पुनर्जुळणी करून त्या वर्तुळाच्या दोन प्रती तयार करता येतात.
\end{thm}

वितालीचा विरोधाभास बानाख--टार्स्की विरोधाभासापेक्षा सिद्ध करणे खूप सोपे आहे. वितालीने 1905 मध्ये केलेल्या अमापनीय संचाच्या रचनेवरून तो मिळत असल्यामुळे त्याला “वितालीचा विरोधाभास” म्हटले आहे. पण विताली आणि बानाख--टार्स्की यांच्या पुराव्यांचे संचसैद्धान्तिक पैलू अगदी सारखे आहेत. मुख्य फरक एवढाच की बानाख--टार्स्कीमध्ये \emph{सांत} विघटन आहे, तर वितालीमध्ये \emph{गणनीय अनंत} विघटन आहे. वेस्टन (2003) यांच्या शब्दांत, वितालीचा विरोधाभास “बानाख--टार्स्की विरोधाभासाइतका नक्कीच थक्क करणारा नाही; पण अगदी ‘साध्या’ परिस्थितीतही भूमितीय विरोधाभास येऊ शकतात, हे तो दाखवतो.”

वितालीचा विरोधाभास वर्तुळ या द्विमितीय आकृतीविषयी आहे. म्हणून $\Real^2$ या प्रतलावर काम करू. मूलबिंदूभोवती पूर्ण फेरीच्या \emph{परिमेय} भागांनुसार, $[0,2\pi)$ या मर्यादेत रेडियनमूल्य असलेली दक्षिणावर्त घूर्णने $\rotationsgroup$ मध्ये घ्या. येथे पूर्ण फेरीचा परिमेय भाग अभिप्रेत आहे; परिमेय संख्येइतके रेडियन नव्हेत. $\rotationsgroup$ विषयी पुढील बीजगणितीय तथ्ये पाहा (विधान समजले नाही, तर पुराव्यात त्याचा अर्थ स्पष्ट होईल):

\begin{lem}\label{sth:choice:vitali:rotationsgroupabelian}
फलनांच्या संयोजनाखाली $\rotationsgroup$ हा अबेली, म्हणजे क्रमनिरपेक्ष, {गट} होतो.
\end{lem}

\begin{proof}
$0$ रेडियनचे घूर्णन $0_{\rotationsgroup}$ असे लिहिल्यास हा $\rotationsgroup$ चा अविकारी घटक आहे. कारण प्रत्येक $\rho \in \rotationsgroup$ साठी $\comp{0_{\rotationsgroup}}{\rho} = \comp{\rho}{0_{\rotationsgroup}} =
\rho$.

प्रत्येक घटकाला व्यस्त असतो. शून्य घूर्णनाचा व्यस्त तोच असतो. इतर प्रकरणांत, $\rho \in \rotationsgroup$ हे $r$ रेडियनचे घूर्णन असेल, तर $\rho^{-1} \in \rotationsgroup$ हे $2\pi - r$ रेडियनचे घूर्णन असते. त्यामुळे $\rho \circ \rho^{-1} = 0_\rotationsgroup$.

संयोजनात सहचारिता आहे: कोणत्याही $\rho, \sigma, \tau \in
\rotationsgroup$ साठी $\comp{\rho}{(\comp{\sigma}{\tau})} =
\comp{(\comp{\rho}{\sigma})}{\tau}$.

संयोजन क्रमनिरपेक्ष आहे: कोणत्याही $\rho, \sigma \in \rotationsgroup$ साठी $\comp{\rho}{\sigma} =
\comp{\sigma}{\rho}$.
\end{proof}

खरे तर $\rotationsgroup$ या गटाचे दोन भाग करून प्रत्येक भागावरून पूर्ण गट पुन्हा मिळवता येतो:

\begin{lem}\label{sth:choice:vitali:disjointgroup}
$\rotationsgroup$ चे $\rotationsgroup_{1}$ आणि $\rotationsgroup_{2}$ अशा दोन वियुक्त संचांत विभाजन करता येते; दोन्ही $\rotationsgroup$ चे जनक संच आहेत.
\end{lem}

\begin{proof}
पूर्ण फेरीच्या परिमेय भागांशी जुळणारी आणि रेडियनमूल्य $[0, \pi)$ मध्ये असलेली घूर्णने $\rotationsgroup_{1}$ मध्ये घ्या; आणि $\rotationsgroup_{2} = \rotationsgroup
\setminus \rotationsgroup_1$ घ्या. प्राथमिक बीजगणिताने
$\Setabs{\comp{\rho}{\rho}}{\rho \in \rotationsgroup_1} =
\rotationsgroup$. $\rotationsgroup_2$ साठीही तसाच निष्कर्ष मिळतो.
\end{proof}

गटांविषयीचे हे तथ्य \ref{sth:choice:vitali:vitaliparadox} सिद्ध करण्यासाठी वापरू. $\onesphere$ हे एकक वर्तुळ, म्हणजे प्रतलाच्या मूलबिंदूपासून नेमके $1$ एकक अंतरावरील बिंदूंचा संच, म्हणजे $\Setabs{\tuple{r,s} \in \Real^2}{\sqrt{r^2+s^2}=1}$ असू द्या. $\onesphere$ वरील पुढील संबंध पाहून $\onesphere$ चे भाग करू:
\[
	r \sim s \emph{ असेल तेव्हा आणि केवळ तेव्हाच }(\exists \rho \in \rotationsgroup)\rho(r) = s.
\]
म्हणजे मूलबिंदूभोवती पूर्ण फेरीच्या परिमेय भागाचे घूर्णन करून एका बिंदूपासून दुसऱ्याकडे जाता येत असेल तेव्हा आणि केवळ तेव्हाच $\onesphere$ चे ते बिंदू या संबंधाने जोडलेले आहेत. अपेक्षेप्रमाणे:

\begin{lem}
$\sim$ हा तुल्यता संबंध आहे.
\end{lem}

\begin{proof}
\ref{sth:choice:vitali:rotationsgroupabelian} वापरून हे तात्काळ मिळते.
\end{proof}

आता निवड वापरून $\sim$ नुसार $\onesphere$ च्या प्रत्येक तुल्यतावर्गातून नेमका एक घटक असलेला संच $C$ मिळवू. म्हणजे तुल्यतावर्गांच्या पुढील संचावर निवडफलन $f$ घेऊ:\footnote{$\rotationsgroup$ हे गणनीय असल्यामुळे $E$ चा प्रत्येक घटक गणनीय आहे. $\onesphere$ हे अगणनीय असल्यामुळे \ref{sth:choice:vitali:vitalicover} आणि \ref{sth:card-arithmetic:simp:kappaunionkappasize} नुसार $E$ हे अगणनीय आहे. म्हणून येथे \emph{अ}गणनीय निवड वापरली आहे.}
\[
	E = \Setabs{\equivrep{r}{\sim}}{r \in \onesphere},
\]
आणि $C = \ran{f}$ घ्या. प्रत्येक घूर्णन $\rho \in \rotationsgroup$ साठी, $C$ च्या प्रत्येक बिंदूवर $\rho$ लागू करून मिळणाऱ्या बिंदूंचा संच $\funimage{\rho}{C}$ आहे. पुढील दोन निष्कर्ष दाखवतात की हे संच वर्तुळ पूर्णपणे व्यापतात आणि एकमेकांवर आच्छादित होत नाहीत:

\begin{lem}\label{sth:choice:vitali:vitalicover}
$\onesphere = \bigcup_{\rho \in \rotationsgroup} \funimage{\rho}{C}$.
\end{lem}

\begin{proof}
$s \in \onesphere$ निश्चित करा. $r \in \equivrep{s}{\sim}$ असलेला $r \in C$ आहे; म्हणजे $r \sim s$, म्हणजे एखाद्या $\rho \in \rotationsgroup$ साठी $\rho(r) = s$.
\end{proof}

\begin{lem}\label{sth:choice:vitali:vitalinooverlap}
$\rho_1 \neq \rho_2$ असल्यास $\funimage{\rho_1}{C} \cap \funimage{\rho_2}{C} = \emptyset$.
\end{lem}

\begin{proof}
$s \in \funimage{\rho_{1}}{C} \cap \funimage{\rho_{2}}{C}$ असे समजा. मग एखाद्या $r_{1}, r_{2} \in C$ साठी $s = \rho_{1}(r_{1}) = \rho_{2}(r_{2})$.
त्यामुळे $\rho^{-1}_2(\rho_1(r_1)) = r_2$, आणि $\comp{\rho_1}{\rho^{-1}_2} \in \rotationsgroup$; म्हणून $r_{1} \sim
r_{2}$. $\sim$ नुसार प्रत्येक तुल्यतावर्गातून $C$ नेमका एक घटक निवडतो, म्हणून $r_1 = r_2$. त्यामुळे $s = \rho_1(r_1) = \rho_2(r_1)$ आणि म्हणून $\rho_1 = \rho_2$.
\end{proof}

आता आधीची बीजगणितीय तथ्ये वर्तुळाला लागू करू:

\begin{lem}\label{sth:choice:vitali:pseudobanachtarski}
$\onesphere$ चे $D_{1}$ आणि $D_{2}$ अशा दोन वियुक्त संचांत विभाजन करता येते, ज्यामध्ये $D_{1}$ चे गणनीय संख्येने भाग करून ते घूर्णित केल्यास $\onesphere$ ची प्रत तयार होते (आणि $D_{2}$ साठीही तसेच).
\end{lem}

\begin{proof}
\ref{sth:choice:vitali:disjointgroup} मधील $\rotationsgroup_{1}$ आणि $\rotationsgroup_{2}$ वापरून पुढील व्याख्या करा:
\begin{align*}
	D_{1} &= \bigcup_{\rho \in \rotationsgroup_1} \funimage{\rho}{C} &
	D_{2} &= \bigcup_{\rho \in \rotationsgroup_2} \funimage{\rho}{C}
\end{align*}
\ref{sth:choice:vitali:vitalicover} नुसार हे $\onesphere$ चे विभाजन आहे, आणि \ref{sth:choice:vitali:vitalinooverlap} नुसार $D_1$ आणि $D_2$ वियुक्त आहेत. रचनेनुसार $D_1$ चे गणनीय संख्येने भाग करता येतात: प्रत्येक $\rho \in \rotationsgroup_1$ साठी $\funimage{\rho}{C}$ हा एक भाग. \ref{sth:choice:vitali:disjointgroup} आणि \ref{sth:choice:vitali:vitalicover} नुसार $\onesphere = \bigcup_{\rho \in
\rotationsgroup}\funimage{\rho}{C} = \bigcup_{\rho \in
\rotationsgroup_1}\funimage{(\comp{\rho}{\rho})}{C}$ असल्यामुळे हे भाग घूर्णित करून $\onesphere$ ची प्रत तयार करता येते. $D_2$ साठीही तोच युक्तिवाद लागू होतो. \end{proof}\noindent यावरून वितालीचा विरोधाभास लगेच मिळतो. वर्तुळाचे गणनीय संख्येने तुकडे (विविध $\funimage{\rho}{C}$) करून दृढ गतीने ते हलवल्यास $\onesphere$ पासून $\onesphere$ च्या \emph{दोन} प्रती तयार करता येतात. फक्त $D_1$ चा प्रत्येक भाग घूर्णित करा; आणि $D_2$ च्या प्रत्येक भागाचे आधी दुसरीकडे स्थानांतरण करून नव्या केंद्राभोवती घूर्णन करा.

पुराव्याची पद्धत पुन्हा पाहू. प्रतलावरील घूर्णन गटाच्या काही बीजगणितीय तथ्यांपासून सुरुवात केली. हा गट वापरून $\onesphere$ चे तुल्यतावर्गांत विभाजन केले. त्यानंतर प्रत्येक वर्गातून घटक निवडण्यासाठी निवड वापरून “विरोधाभास” मिळवला.

बानाख--टार्स्की सिद्ध करण्यासाठी नेमकी हीच पद्धत वापरतो. मुख्य फरक असा की त्यासाठीची बीजगणितीय तथ्ये वितालीच्या पुराव्यातील तथ्यांपेक्षा बरीच गुंतागुंतीची आहेत. पण त्या बीजगणितीय तथ्यांचा निवडीशी संबंध नाही. त्यांचा थोडक्यात सारांश देऊ.

बानाख--टार्स्की सिद्ध करण्यासाठी आधी \ref{sth:choice:vitali:disjointgroup} सारखा निष्कर्ष मिळवतो: दोन जनकांवरील \emph{मुक्त गटाचे} चार भाग करता येतात, आणि सहज अर्थाने ते “हलवून” पूर्ण गटाच्या दोन प्रती मिळवता येतात.\footnote{\emph{चार} भाग वापरता येतात, हा निष्कर्ष (रॉबिन्सन 1947) यांचा आहे. अलीकडील पुराव्यासाठी वॅगन (2016, प्रमेय 5.2) पाहा. गटाचे भाग “हलवणे” असे म्हणताना वेस्टन (2003, पृ.~3) यांच्या वर्णनाचा आधार घेतला आहे.} त्यानंतर $\Real^3$ च्या मूलबिंदूभोवतीची दोन विशिष्ट घूर्णने वापरून घूर्णनांचा मुक्त गट $F$ तयार करता येतो, हे दाखवतो.\footnote{पाहा वॅगन (2016, प्रमेय 2.1).} (अजून निवड वापरलेली नाही.) आता $F$ मधील घूर्णनाने एका बिंदूपासून दुसरा मिळत असेल तेव्हा आणि केवळ तेव्हाच गोलकाच्या पृष्ठभागावरील ते बिंदू “सदृश” मानतो. मग “सदृश” बिंदूंच्या प्रत्येक तुल्यतावर्गातून नेमका एक बिंदू निवडण्यासाठी \emph{निवड वापरतो}. \ref{sth:choice:vitali:pseudobanachtarski} प्रमाणे $F$ चे विभाजन पृष्ठभागाला लागू करण्याची कल्पना अशी की चार भाग “हलवून” पृष्ठभागाच्या दोन प्रती मिळवाव्यात. मात्र या त्रिमितीय रचनेत अक्षांवरील काही बिंदू घूर्णनांनी स्थिर राहतात; अशा बिंदूंचा स्वतंत्र विचार आणि आवश्यक अतिरिक्त युक्तिवाद या संक्षिप्त रूपरेषेत दिलेला नाही. म्हणून ही पूर्ण पुराव्याची जागा घेऊ शकत नाही. संबंधित निष्कर्ष (हाउसडॉर्फ 1914) असा मांडला जातो:

\begin{thm}[हाउसडॉर्फचा विरोधाभास ($\ZFC$ मध्ये)]
धन त्रिज्येच्या कोणत्याही गोलकाच्या पृष्ठभागाचे सांत संख्येने तुकडे करता येतात; घूर्णन आणि स्थानांतरणाने त्यांची पुनर्जुळणी करून त्या पृष्ठभागाच्या दोन वियुक्त प्रती तयार करता येतात.
\end{thm}

पूर्ण बानाख--टार्स्की प्रमेय मिळवण्यासाठी (जे फक्त गोलकाच्या पृष्ठभागाविषयी नाही, तर अंतर्भागाविषयीही आहे) आणखी काही बीजगणितीय युक्त्या लागतात. तरी लेखकाच्या मते मुख्य बीजगणितीय काम आधीच्या टप्प्यावर आहे. म्हणून वेस्टन लिहितात:
\begin{quote}
  [\ldots] मुक्त गटांविषयीचा निष्कर्ष बानाख--टार्स्की विरोधाभासाच्या पुराव्यातील \emph{मुख्य टप्पा} आहे. या दृष्टीने बानाख--टार्स्की विरोधाभास हा $\Real^3$ विषयीच्या विधानापेक्षा [$\Real^3$ मधील स्थानांतरणे आणि घूर्णने यांच्या] गटाच्या गुंतागुंतीविषयीचे विधान आहे. (वेस्टन 2003, पृ.~16)
\end{quote}
म्हणजे बानाख--टार्स्कीप्रमाणे \emph{सांत} विघटन मिळते की वितालीप्रमाणे \emph{गणनीय अनंत} विघटन मिळते, हे दोन किंवा तीन मितींत काम करताना गटांविषयी लागू होणाऱ्या काही तथ्यांवर ठरते.

हे शेवटचे निरीक्षण \ref{sth:choice:banach:sec} च्या अखेरचा विनोद थोडा बिघडवते. लेखन द्विमितीय असल्यामुळे “बानाख--टार्स्की बानाख--टार्स्की” अशी पुनर्जुळणी करायची असेल, तर “बानाख--टार्स्की”चे गणनीय \emph{अनंत} तुकडे करावे लागतील. विनोद दुरुस्त करण्यासाठी त्रिमितीय लेखन करावे. हा सराव वाचकासाठी सोडला आहे.

शेवटची एक नोंद. \ref{sth:choice:banach:sec} मध्ये गोलकाचे मिळणारे “तुकडे” सर्व \emph{मापनीय} असू शकत नाहीत; त्यात चित्रात दाखवता न येणारे “अनंत विखुरणे” असते, असे म्हटले होते. \ref{sth:choice:vitali:pseudobanachtarski} मिळवण्यासाठी केलेल्या निवडीच्या वापराविषयीही हेच लागू होते. हे स्पष्ट करणे उपयुक्त आहे.

पुन्हा थोडी पार्श्वभूमी सांगावी लागेल (ही \emph{केवळ} रूपरेषा आहे; \emph{मापाविषयी} पाठ्यपुस्तकातील विवेचन पाहावे). येथे वापरल्या जाणाऱ्या सांत मापात, संच $X$ वरील एखाद्या “$\sigma$-बीजगणितातील” प्रत्येक $E$ ला ऋणेतर मूल्य $\mu(E) \in
\Real$ दिले जाते; हे $X$ चे माप ठरवते. तपशील येथे आवश्यक नाहीत; पण फलन $\mu$ ने गणनीय योगशीलतेचे तत्त्व पाळले पाहिजे: वियुक्त संचांच्या गणनीय युतीकरणाचे माप त्यांच्या स्वतंत्र मापांची बेरीज असते. म्हणजे $X_n$ हे संच वियुक्त असताना $\mu(\bigcup_{n < \omega} X_n) = \sum_{n < \omega}\mu(X_n)$.
संच “अमापनीय” आहे याचा अर्थ तो या निश्चित मापाच्या मापनीय क्षेत्रात येत नाही; कोणतेही माप कधीच देता येत नाही, असा सार्वत्रिक अर्थ नाही. आता आधीचा $\rotationsgroup$ वापरू:

\begin{cor}[विताली]
$\mu$ हे असे माप असू द्या की $\mu(\onesphere) = 1$, आणि $X$ व $Y$ एकरूप असतील, तर $\mu(X) = \mu(Y)$. मग प्रत्येक $\rho \in
\rotationsgroup$ साठी $\funimage{\rho}{C}$ अमापनीय आहे.
\end{cor}

\begin{proof}
विरोधसिद्धीसाठी तसे नाही असे समजा. मग एखाद्या $\sigma \in \rotationsgroup$ आणि $r \in \Real$ साठी $\mu(\funimage{\sigma}{C}) =
r$ असू द्या. कोणत्याही $\rho \in \rotationsgroup$ साठी, $\funimage{\rho}{C}$ आणि $\funimage{\sigma}{C}$ एकरूप आहेत; म्हणून प्रत्येक $\rho \in \rotationsgroup$ साठी $\mu(\funimage{\rho}{C}) = r$.
\ref{sth:choice:vitali:vitalicover} आणि \ref{sth:choice:vitali:vitalinooverlap} नुसार $\onesphere =
\bigcup_{\rho \in \rotationsgroup}\funimage{\rho}{C}$ हे जोडीजोडीने वियुक्त संचांचे गणनीय युतीकरण आहे. म्हणून गणनीय योगशीलतेनुसार $\mu(\onesphere) = 1$ ही प्रत्येक $\funimage{\rho}{C}$ च्या मापांची बेरीज आहे; म्हणजे
\[
	1 = \mu(\onesphere) = \sum_{\rho \in \rotationsgroup}\mu(\funimage{\rho}{C}) = \sum_{\rho \in \rotationsgroup}r
\]
पण $r = 0$ असल्यास $\sum_{\rho \in \rotationsgroup}r = 0$, आणि $r
> 0$ असल्यास $\sum_{\rho \in \rotationsgroup}r = \infty$.
\end{proof}



\part{पद्धती}\label{mth:part}
\begin{editorial}
  या भागात सर्वसाधारण आणि पद्धतिविषयक साहित्य आहे. विशेषतः गणिताचा विद्यार्थी नसलेल्या वाचकाला अपरिचित असू शकणाऱ्या विविध सिद्धता-पद्धतींची स्पष्टीकरणे आहेत. या स्रोतभागात सध्या सिद्धता कशा लिहाव्यात यावरील एक प्रकरण आणि विगमनावरील एक प्रकरण आहे. त्यांत आणखी विभाग व सराव जोडणे, तसेच गणितीय संज्ञांविषयी एक प्रकरण तयार करणे नियोजित आहे.
\end{editorial}
\chapter{सिद्धता}\label{mth:prf:chap}









\section{प्रस्तावना}\label{mth:prf:int:sec}

प्रास्ताविक तर्कशास्त्राच्या अभ्यासामुळे तुम्हाला एखादी निष्पत्ती-पद्धती परिचित असेल---बहुधा नैसर्गिक निगमनाची किंवा फिच शैलीतील निष्पत्ती-पद्धती, किंवा कदाचित सिद्धता-वृक्षांची पद्धती. या पद्धतींत सूत्र सिद्ध करणे किंवा दिलेला युक्तिवाद युक्त आहे हे दाखवणे, अशा सिद्धता तुम्ही केल्याचे आठवत असेल. त्यासाठी अपेक्षित अंतिम निष्कर्ष मिळेपर्यंत तुम्ही त्या पद्धतीचे नियम वापरले होते. तर्कशास्त्रा\emph{विषयी} विचार करतानाही आपण गोष्टी सिद्ध करतो; पण बहुतेक वेळा निष्पत्ती-पद्धती वापरत नाही. आपण पाहणार असलेल्या बहुतेक सिद्धता संपूर्णपणे प्रथम-क्रम तर्कशास्त्राच्या भाषेत लिहिलेल्या नसून नैसर्गिक भाषेत असतात; या आवृत्तीत त्या मराठीत आहेत आणि त्यांत काही सांकेतिक भाषाही येते. अशा सिद्धता रचताना सुरुवातीला तुम्हाला प्रश्न पडू शकतात: निष्पत्ती-पद्धतीशिवाय एखादी गोष्ट कशी सिद्ध करायची? सुरुवात कुठून करायची? माझी सिद्धता बरोबर आहे हे कसे ठरवायचे?

सिद्धता करण्याचा प्रयत्न सुरू करण्यापूर्वी सिद्धता म्हणजे काय आणि ती कशी रचायची, हे समजणे महत्त्वाचे आहे. नावावरूनच सूचित होते तसे, एखादी गोष्ट सत्य आहे हे दाखवणे हा \emph{सिद्धते}चा उद्देश असतो. संवादाच्या रूपात याचा विचार करता येईल: कोणी तुम्हाला विचारते की एखादी गोष्ट खरी आहे का---उदाहरणार्थ, दोन वगळता प्रत्येक अविभाज्य संख्या विषम आहे का. केवळ “होय” असे उत्तर देणे पुरेसे नाही; त्या व्यक्तीला \emph{का}, हेही जाणून घ्यायचे असेल. अशा वेळी तुम्ही सिद्धता द्याल.

दैनंदिन संभाषणात उत्तराची ढोबळ दिशा सांगणे किंवा अपूर्ण उत्तर देणे पुरेसे ठरू शकते. पण तर्कशास्त्रात आणि गणितात आपल्याला काटेकोर सिद्धता हवी असते: एखादी गोष्ट खरी असल्याविषयी \emph{कोणत्याही} शंकेला जागा राहू नये, असे दाखवायचे असते. याचा अर्थ सिद्धतेतील प्रत्येक पायरीचे समर्थन केले पाहिजे आणि ते समर्थन योग्य असले पाहिजे. उदाहरणार्थ, तुम्ही वापरत असलेले गृहीतक हे सिद्ध करावयाच्या प्रमेयाच्या विधानात खरोखर गृहीत धरलेले असले पाहिजे; व्याख्यांचा वापर अचूक असला पाहिजे; समर्थनासाठी वापरलेली अनुमाने योग्य असली पाहिजेत; इत्यादी.

सहसा आपण एखादे विधान सिद्ध करत असतो. सिद्ध करावयाच्या विधानांना आपण विविध नावे देतो: विधान, प्रमेय, पूर्वप्रमेय किंवा अनुप्रमेय. विधान म्हणजे सिद्ध करावे असे मूलभूत म्हणणे: नोंदवण्याइतके महत्त्वाचे, पण कदाचित फार सखोल नसलेले किंवा वारंवार न वापरले जाणारे. प्रमेय म्हणजे लक्षणीय, महत्त्वाचे विधान. त्याची सिद्धता अनेकदा काही पायऱ्यांत विभागलेली असते. कधी त्याला ते प्रथम सिद्ध करणाऱ्या व्यक्तीचे नाव दिले जाते---उदाहरणार्थ, कँटरचे प्रमेय किंवा लोव्हेनहाइम--स्कोलेम प्रमेय (L\"owenheim--Skolem)---तर कधी त्याच्या विषयावरून नाव दिले जाते---उदाहरणार्थ, संपूर्णता प्रमेय. पूर्वप्रमेय म्हणजे अधिक महत्त्वाचा निष्कर्ष सिद्ध करण्यासाठी वापरले जाणारे विधान किंवा प्रमेय. नावामुळे गोंधळ होऊ शकतो: कधी पूर्वप्रमेये स्वतःही महत्त्वाचे निष्कर्ष असतात आणि ती मांडणाऱ्या व्यक्तीच्या नावाने ओळखली जातात---उदाहरणार्थ, झोर्नचे पूर्वप्रमेय. अनुप्रमेय म्हणजे दुसऱ्या निष्कर्षावरून सहज मिळणारा निष्कर्ष.

सिद्ध करावयाच्या विधानात अनेकदा काही गृहीतके असतात. आपण कोणत्या प्रकारच्या गोष्टींविषयी सिद्धता करतो, हे त्यांमुळे स्पष्ट होते. विधानाची सुरुवात “$\varphi$ हे $\psi \lif \chi$ या रूपाचे सूत्र असू द्या” किंवा “$\Gamma \Proves \varphi$ असे समजा”, किंवा अशाच शब्दांनी होऊ शकते. ही त्या विधानाची, प्रमेयाची किंवा पूर्वप्रमेयाची \emph{गृहीतके} आहेत; सिद्धतेत ती सत्य आहेत असे धरता येते. काय सिद्ध करायचे आहे याची व्याप्ती ती निश्चित करतात आणि आपण ज्या वस्तूंविषयी बोलतो त्यांच्यासाठी काही नावेही देतात. उदाहरणार्थ, तुमच्या विधानाची सुरुवात “$\varphi$ हे $\psi \lif \chi$ या रूपाचे सूत्र असू द्या” अशी असेल, तर तुम्ही केवळ एका विशिष्ट प्रकारच्या सर्व सूत्रांविषयी---म्हणजे सशर्त सूत्रांविषयी---काहीतरी सिद्ध करत आहात. तुमची सिद्धता ज्या $\psi \lif \chi$ विषयी बोलते ते त्या प्रकारचे स्वेच्छ सशर्त सूत्र आहे, असे समजले जाते.












\section{सिद्धतेची सुरुवात}\label{mth:prf:str:sec}

पण सुरुवात करायची तरी कुठून?

तुम्हाला काहीतरी सिद्ध करायला दिले आहे. म्हणून सिद्धतेत शेवटी त्याचाच उल्लेख आला पाहिजे. अर्थात, सुरुवातीलाच आपण ते सिद्ध करणार आहोत असे \emph{जाहीर} करता येते; पण ते गृहीतक म्हणून वापरायचे नाही. कोऱ्या कागदाच्या तळाशी काय सिद्ध करायचे आहे ते लिहा. त्यामुळे तुमचे ध्येय नजरेआड होणार नाही.

यानंतर, वापरता येण्याजोगी काही गृहीतके असतील. तुम्ही करत असलेल्या सिद्धतेच्या \emph{प्रकारा}विषयी पुढील विभागांत बोलताना हे अधिक स्पष्ट होईल. ही गृहीतके पानाच्या वरच्या बाजूला लिहा आणि ती गृहीतके आहेत हे स्पष्टपणे दर्शवा. उदाहरणार्थ, $p$ गृहीत धरत असाल, तर “$p$ असे गृहीत धरा” किंवा “$p$ असे समजा” असे लिहा. शेवटी, प्रश्नात काही व्याख्या असतील ज्या तुम्हाला समजणे आवश्यक आहे. एखादी विशिष्ट व्याख्या वापरायला सांगितलेली असेल; किंवा गृहीतकांत आणि ज्या निष्कर्षापर्यंत पोहोचायचे आहे त्यात विविध व्याख्या आलेल्या असतील. \emph{या व्याख्या लिहून काढा आणि त्यांचा अर्थ तुम्हाला समजतो आहे याची खात्री करा.}

सिद्धतेची मांडणी कशी करायची हे प्रश्नाच्या रूपावरही अवलंबून असते. वाक्याच्या प्रकारानुसार सिद्धतेची मांडणी कशी करायची याचा तपशील पुढील विभागांत दिला आहे.












\section{व्याख्यांचा वापर}\label{mth:prf:def:sec}

सिद्धतेत वापरल्या जाऊ शकणाऱ्या सर्व व्याख्या तुम्हाला परिचित असल्या पाहिजेत आणि त्यांचा योग्य वापर करता आला पाहिजे, असे आपण म्हटले होते. हा मुद्दा फार महत्त्वाचा आहे; त्याचा थोडा अधिक तपशील पाहू. गुणधर्म आणि संबंधांविषयी थोडक्यात बोलता यावे म्हणून व्याख्यांद्वारे त्यांना संक्षिप्त नावे दिली जातात. असे नाव म्हणजे \emph{व्याख्येय}; ते ज्याचा संक्षेप करते तो मजकूर म्हणजे \emph{व्याख्यापक}. सिद्धतेत अनेकदा व्याख्येय कसे मांडले होते याकडे परत जावे लागते. कारण सिद्धता पुढे नेण्यासाठी व्याख्यापकाची तार्किक रचना वापरावी लागते: व्याख्येय संज्ञा ज्याचा संक्षेप आहे तोच हा विस्तृत मजकूर असतो. व्याख्या उलगडल्यामुळे प्रत्यक्ष तार्किक कार्य जिथे होते तिथपर्यंत आपण पोहोचतो.

एका उदाहरणाने सुरुवात करू. पुढील विधान सिद्ध करायचे आहे असे समजा:

\begin{prop}
कोणत्याही $A$ आणि $B$ संचांसाठी, $A \cup B = B \cup A$.
\end{prop}

सिद्धतेची सुरुवात करण्यासाठीच दोन संच एकरूप असणे म्हणजे काय हे समजणे आवश्यक आहे. म्हणजे त्या समीकरणातील “$=$” चा संचांसाठी काय अर्थ आहे हे समजले पाहिजे. दोन संचांचे घटक तेच असतील तेव्हा ते संच एकरूप असतात, अशी व्याख्या आहे. म्हणून उलगडायची व्याख्या ही आहे:

\begin{defn}
$A$ आणि $B$ संच \emph{एकरूप} आहेत, $A = B$, तेव्हाच आणि फक्त तेव्हा जेव्हा~$A$ मधील प्रत्येक घटक हा~$B$ मधील घटक असतो आणि उलटही तसेच असते.
\end{defn}

या व्याख्येत $A$ आणि~$B$ ही स्वेच्छ संचांसाठी वापरलेली प्रतिनिधी नावे आहेत. ही व्याख्या ज्या गोष्टीची व्याख्या करते---म्हणजे \emph{व्याख्येय}---ती “$A = B$” ही अभिव्यक्ती आहे. $A = B$ कोणत्या अटीवर सत्य असते हे व्याख्या सांगते. ती अट---“~$A$ मधील प्रत्येक घटक हा~$B$ मधील घटक असतो आणि उलटही तसेच असते”---म्हणजे \emph{व्याख्यापक}.\footnote{या विशिष्ट प्रकरणात---आणि येथे फार गोंधळ होऊ शकतो!---$A = B$ असते तेव्हा $A$ आणि $B$ हे प्रत्यक्षात एकच संच असतात, जरी त्यासाठी डावीकडे आणि उजवीकडे वेगवेगळी अक्षरे वापरली असली तरी. पण त्या संचाची ओळख करून देण्याचे मार्ग वेगळे असू शकतात; त्यामुळे व्याख्या निरर्थक पुनरुक्ती ठरत नाही.} ही अट पूर्ण होते तेव्हाच आणि फक्त तेव्हा (इंग्रजी संक्षेप “iff”) $A = B$ सत्य असते, असे ही व्याख्या सांगते.

व्याख्या वापरताना तिच्यातील $A$ आणि $B$ यांचा संबंध प्रत्यक्ष हाताळत असलेल्या प्रकरणाशी योग्य रीतीने जुळवला पाहिजे. आपल्या प्रकरणात $A \cup B = B \cup A$ सत्य असण्यासाठी प्रत्येक $z \in A \cup B$ हा $B \cup A$ मध्येही असला पाहिजे आणि उलटही तसेच असले पाहिजे. विधानातील $A \cup B$ ही अभिव्यक्ती व्याख्येतील~$A$ ची भूमिका बजावते आणि $B \cup A$ ही व्याख्येतील~$B$ ची. व्याख्येत आणि आपण सिद्ध करत असलेल्या विधानात $A$ आणि $B$ हीच नावे वेगवेगळ्या भूमिकांत येतात; म्हणून या दोन वापरांचा गोंधळ होणार नाही याची काळजी घ्यावी. उदाहरणार्थ,~$A$ मधील प्रत्येक घटक हा~$B$ मधील घटक आहे आणि उलटही तसेच आहे, हे दाखवून विधान सिद्ध करता येईल असे समजणे चुकीचे होईल. त्यामुळे $A = B$ सिद्ध होईल; $A \cup B = B \cup A$ सिद्ध होणार नाही. शिवाय, $A$ आणि $B$ हे कोणतेही दोन संच असू शकतात. $A$ आणि~$B$ यांच्याविषयी काहीच गृहीत धरलेले नसेल, तर ते वेगळे संचही असू शकतात; म्हणून त्या चुकीच्या मार्गाने फार पुढे जाता येणार नाही.

सिद्धतेत संयोगासारख्या संचसैद्धान्तिक संकल्पना येतात. त्यामुळे सिद्धता कशी पुढे न्यायची हे समजण्यासाठी $\cup$ या चिन्हाचाही अर्थ समजला पाहिजे. कधी एखादी व्याख्या उलगडल्यावर आणखी व्याख्या उलगडाव्या लागतात. उदाहरणार्थ, $A \cup B$ ची व्याख्या $\Setabs{z}{z \in A \text{ किंवा } z \in B}$ अशी आहे. म्हणून $x \in A \cup B$ सिद्ध करायचे असेल, तर $\cup$ ची व्याख्या उलगडल्यावर $x \in \Setabs{z}{z \in A \text{ किंवा } z \in B}$ सिद्ध करावे लागेल असे दिसते. आता $x \in \Setabs{z}{\dots z\dots}$ हे $\dots x\dots$ असते तेव्हाच आणि फक्त तेव्हाच खरे असते, हेही आठवले पाहिजे. म्हणून $\Setabs{z}{\dots z \dots}$ या लेखनाची व्याख्या आणखी उलगडल्यावर असे दिसते की $x \in A$ किंवा $x \in B$ हे दाखवायचे आहे. त्यामुळे “$A \cup B$ मधील प्रत्येक घटक हा $B \cup A$ मधील घटक देखील असतो” याचा प्रत्यक्ष अर्थ असा आहे: “प्रत्येक $x$ साठी, $x \in A$ किंवा $x \in B$ असेल, तर $x \in B$ किंवा $x \in A$ असते.” विधानातील व्याख्या पूर्ण उलगडल्यावर आपण हे दाखवायचे आहे असे दिसते:

\begin{prop}
कोणत्याही $A$ आणि $B$ संचांसाठी: (a) प्रत्येक $x$ साठी, $x \in A$ किंवा $x \in B$ असेल, तर $x \in B$ किंवा $x \in A$ असते; आणि (b) प्रत्येक $x$ साठी, $x \in B$ किंवा $x \in A$ असेल, तर $x \in A$ किंवा $x \in B$ असते.
\end{prop}

व्याख्या उलगडणे हा सिद्धता रचण्याचा आवश्यक भाग आहे, हा महत्त्वाचा मुद्दा आहे. हे योग्य रीतीने करणे कधी अवघड असते: व्याख्येतील चल आणि सिद्ध करावयाच्या दाव्यातील पदे यांतील फरक ओळखून त्यांची योग्य जुळवणी केली पाहिजे. यशस्वी होण्यासाठी प्रश्न नेमके काय विचारतो आणि त्यातील सर्व संज्ञांचा अर्थ काय आहे हे समजले पाहिजे. अनेकदा एकाहून अधिक व्याख्या उलगडाव्या लागतील. पुढे येणाऱ्या साध्या सिद्धतांत उत्तर जवळजवळ थेट व्याख्यांवरूनच मिळते. अर्थात, प्रत्येक वेळी हे इतके सोपे असेलच असे नाही.

\begin{prob}
$A \cap B \neq \emptyset$ सिद्ध करायला सांगितले आहे असे समजा. येथे येणाऱ्या सर्व व्याख्या उलगडा; म्हणजे “$\cap$”, “=” किंवा “$\emptyset$” यांचा उल्लेख न करता हे विधान पुन्हा मांडा.
\end{prob}












\section{अनुमानाचे नमुने}\label{mth:prf:pat:sec}

सिद्धता स्वतंत्र अनुमानांनी बनतात. अनुमान करताना आपण सहसा “म्हणून”, “त्यामुळे” किंवा “यावरून” असे शब्द वापरून ते दर्शवतो. त्या अनुमानाचा आधार अनेकदा सिद्धतेत आधीच उपलब्ध असलेल्या एक-दोन गोष्टी असतात: एखादे गृहीतक किंवा आधीच्या अनुमानाने मिळालेला निष्कर्ष. स्पष्टतेसाठी या गोष्टींना नावे किंवा क्रमांक देता येतात आणि अनुमानात कोणती इतर विधाने वापरली आहेत हे दर्शवता येते. नव्या निष्कर्षाची सत्यता आधीच्या गोष्टींवरून \emph{का} मिळते याचे स्पष्टीकरणही अनुमानात अनेकदा असते. सिद्धतांत वारंवार वापरले जाणारे अनुमानाचे काही सामान्य नमुने आहेत; त्यांपैकी काही पुढे पाहू. विगमनाने सिद्धता करण्यासारखे काही नमुने अधिक गुंतागुंतीचे आहेत; त्यांची चर्चा नंतर होईल.

अनुमानाचा एक नमुना आपण आधीच पाहिला आहे: व्याख्या उलगडणे किंवा तिचा वापर करणे. व्याख्या उलगडताना व्याख्येय असलेले म्हणणे आपण व्याख्यापकाच्या रूपात पुन्हा मांडतो. उदाहरणार्थ, सिद्धतेत आधीच $D = E$ हे (a) प्रस्थापित केले आहे असे समजा. मग संचांसाठीच्या $=$ च्या व्याख्येवरून असे अनुमान करता येते: “म्हणून (a) आणि व्याख्येवरून,~$D$ मधील प्रत्येक घटक हा~$E$ मधील घटक असतो आणि उलटही तसेच असते.”

थोडा गोंधळ होऊ शकतो: अनेकदा अनुमानाचे समर्थन प्रत्यक्ष अनुमान करताना लिहिण्याऐवजी त्याआधीच लिहिले जाते. $D = E$ अद्याप सिद्ध केलेले नाही, पण ते सिद्ध करायचे आहे असे समजा. अपेक्षित निष्कर्ष $D = E$ असेल, तर व्याख्या वापरून हे ध्येयही पुन्हा मांडता येते: $D = E$ सिद्ध करण्यासाठी~$D$ मधील प्रत्येक घटक हा~$E$ मधील घटक आहे आणि उलटही तसेच आहे, हे सिद्ध करावे लागते. म्हणून सिद्धतेचे रूप असे असेल: (a)~$D$ मधील प्रत्येक घटक हा~$E$ मधील घटक आहे हे सिद्ध करा; (b)~$E$ मधील प्रत्येक घटक हा~$D$ मधील घटक आहे हे सिद्ध करा; (c) म्हणून (a), (b) आणि $=$ च्या व्याख्येवरून $D = E$. पण सहसा आपण ते अशा रीतीने लिहीत नाही. त्याऐवजी असे लिहू शकतो:
\begin{quote}
$D = E$ दाखवायचे आहे. ~$=$ च्या व्याख्येनुसार त्यासाठी~$D$ मधील प्रत्येक घटक हा~$E$ मधील घटक आहे आणि उलटही तसेच आहे, हे दाखवणे पुरेसे आहे.

(a) \dots (~$D$ मधील प्रत्येक घटक हा~$E$ मधील घटक आहे याची सिद्धता) \dots

(b) \dots (~$E$ मधील प्रत्येक घटक हा~$D$ मधील घटक आहे याची सिद्धता) \dots
\end{quote}

\subsection{संधियोगाचा वापर}

संधियोगातील एखादा संधिघटक निष्कर्ष म्हणून काढणे हा कदाचित अनुमानाचा सर्वात साधा नमुना आहे. म्हणजे $p$ आणि~$q$ असे गृहीत धरले असेल किंवा आधीच सिद्ध केले असेल, तर~$p$ चे (आणि~$q$ चेही) अनुमान करता येते. हे इतके मूलभूत अनुमान आहे की अनेकदा त्याचा स्वतंत्र उल्लेख होत नाही. उदाहरणार्थ, $D = E$ ची व्याख्या उलगडल्यावर~$D$ मधील प्रत्येक घटक हा~$E$ मधील घटक आहे आणि उलटही तसेच आहे, हे आपल्याला मिळते. यावरून~$E$ मधील प्रत्येक घटक हा~$D$ मधील घटक आहे असा निष्कर्ष काढता येतो; तोच “उलटही तसेच” हा भाग आहे.

\subsection{संधियोग सिद्ध करणे}

कधी सिद्ध करावयाचे विधान संधियोगाच्या रूपात असेल: “$p$ आणि $q$ सिद्ध करा” असे सांगितले जाईल. अशा वेळी दोन गोष्टी कराव्या लागतात: $p$ सिद्ध करा आणि नंतर $q$ सिद्ध करा. सिद्धता दोन भागांत विभागून स्पष्टतेसाठी त्या भागांना नावे किंवा क्रमांक देता येतात. सुरुवातीच्या टिपणांत पानाच्या वरच्या बाजूला “(1) $p$ सिद्ध करा” आणि मध्यभागी “(2) $q$ सिद्ध करा” असे लिहिता येईल. अर्थात, संधियोग सिद्ध करायला स्पष्टपणे सांगितलेले नसले तरी सिद्धतेत संधियोग सिद्ध करणे आवश्यक ठरू शकते. उदाहरणार्थ, $D = E$ सिद्ध करायला सांगितले असेल, तर~$=$ ची व्याख्या उलगडल्यावर~$D$ मधील प्रत्येक घटक हा~$E$ मधील घटक आहे \emph{आणि}~$E$ मधील प्रत्येक घटक हा~$D$ मधील घटक आहे, हे सिद्ध करावे लागते असे दिसेल.

\subsection{विकल्पयोग सिद्ध करणे}

सिद्ध करावयाचे विधान विकल्पयोगाच्या रूपात असेल---म्हणजे “$p$ किंवा $q$” या रूपाचे असेल---तर एका विकल्पघटकाची सत्यता दाखवणे पुरेसे असते. पण प्रमेयाच्या गृहीतकांवरून कोणताही एक विकल्पघटक आपोआप मिळतो, असे सहसा घडत नाही. उलट, अनेकदा प्रमेयाची गृहीतकेच विकल्पयोगात्मक असतात; किंवा एका प्रकारच्या सर्व वस्तूंत दोनपैकी एक गुणधर्म आहे हे दाखवायचे असते, पण काही वस्तूंत पहिला आणि इतरांत दुसरा गुणधर्म असतो. अशा वेळी प्रकरणांनुसार सिद्धता उपयुक्त ठरते; ती पुढे पाहू.

\subsection{सोपाधिक सिद्धता}

अनेक प्रमेये सशर्त रूपात असतात: $p$ असेल, तर $q$ देखील सत्य आहे, हे दाखवायचे असते. अशा सिद्धतेची मांडणी सहज करता येते: सशर्त विधानाचे पूर्वांग---येथे $p$---गृहीत धरा आणि त्यावरून निष्कर्ष~$q$ सिद्ध करा. म्हणून प्रमेय “$p$ असेल, तर $q$” असे असेल, तर सिद्धतेची सुरुवात “$p$ असे गृहीत धरा” अशी करा आणि शेवटी~$q$ सिद्ध झालेले असले पाहिजे.

सशर्त विधाने वेगवेगळ्या रीतीने मांडता येतात. “$p$ असेल, तर $q$” याऐवजी प्रमेय “$q$ असेल तरच $p$”, “$p$ असेल तर $q$” किंवा “$p$ या अटीवर $q$” असे म्हणू शकते. या सर्वांचा एकच अर्थ आहे: $p$ गृहीत धरून त्यावरून~$q$ सिद्ध करावे लागते. द्विसशर्त विधान---“$p$ सत्य असते तेव्हाच आणि फक्त तेव्हा जेव्हा $q$ सत्य असते” (इंग्रजी संक्षेप “iff”)---हे दोन सशर्त विधानांचे एकत्रीकरण असते, हे आठवा: $p$ असेल तर $q$ आणि $q$ असेल तर~$p$. म्हणून सोपाधिक सिद्धता दोनदा करावी लागते: एकदा पहिल्या सशर्त विधानासाठी आणि एकदा दुसऱ्यासाठी. मात्र कधी इतर “तेव्हाच आणि फक्त तेव्हा” विधानांची साखळी जोडूनही असे विधान सिद्ध करता येते: सुरुवात “$p$” पासून होते आणि शेवट “$q$” येथे. अशा वेळी प्रत्येक पायरी खरोखर “तेव्हाच आणि फक्त तेव्हा” अशीच आहे याची खात्री केली पाहिजे.

\subsection{वैश्विक विधाने}

वैश्विक विधानाचा वापर सोपा आहे: एखादी गोष्ट प्रत्येक वस्तूसाठी खरी असेल, तर ती प्रत्येक विशिष्ट वस्तूसाठी खरी असते. उदाहरणार्थ, सिद्धतेचे गृहीतक $A \subseteq B$ असेल, तर~$\subseteq$ ची व्याख्या उलगडल्यावर त्याचा अर्थ असा होतो: प्रत्येक $x \in A$ साठी $x \in B$ असते. त्यामुळे $z \in A$ हे आधीच माहीत असेल, तर $z \in B$ असा निष्कर्ष काढता येतो.

वैश्विक विधान सिद्ध करणे थोडे अवघड वाटू शकते. सहसा ही विधाने “$x$ मध्ये~$P$ हा गुणधर्म असेल, तर त्यात~$Q$ हाही गुणधर्म असतो” किंवा “सर्व $P$ असलेल्या वस्तू $Q$ असलेल्या असतात” या रूपात असतात. अर्थात, प्रत्येक विधान नेमके या रूपात असेलच असे नाही; नेमके काय सिद्ध करायला सांगितले आहे हे ओळखण्यासाठी थोडा सराव लागतो. पण अनेकदा एखादा गुणधर्म असलेल्या सर्व वस्तूंमध्ये आणखी एक विशिष्ट गुणधर्म आहे, हे सिद्ध करावे लागते.

वैश्विक विधान सिद्ध करण्यासाठी पहिला गुणधर्म असलेल्या वस्तूंना नावे किंवा चल द्यायचे आणि त्यांत दुसराही गुणधर्म आहे हे दाखवायचे. हे असेही म्हणता येईल: \emph{सर्व}~$P$ असलेल्या वस्तूंसाठी काहीतरी सिद्ध करायचे असेल, तर एका \emph{स्वेच्छ}~$P$ असलेल्या वस्तूसाठी ते सिद्ध करावे लागते. दिलेले नाव हे अशा स्वेच्छ~$P$ वस्तूचे नाव असते. स्वेच्छ वस्तूंसाठी सहसा एकेक अक्षर वापरतो आणि त्या अक्षरांसाठी काही रूढी असतात: नैसर्गिक संख्यांसाठी $n$, सूत्रे साठी $\varphi$, संचांसाठी $A$, फलनांसाठी $f$, इत्यादी.

$P$ या गुणधर्मापलीकडे विशेष अटी न घालता संपूर्ण सिद्धतेत सामान्यता टिकवणे हा मुख्य मुद्दा आहे. एका स्वेच्छ वस्तूत (“$x$”)~$P$ हा गुणधर्म आहे असे गृहीत धरून सुरुवात करायची आणि केवळ व्याख्या किंवा अनुमत गृहीतकांच्या आधारे $x$ मध्ये~$Q$ हा गुणधर्म आहे हे दाखवायचे. $x$ नेमके काय आहे याविषयी आणखी कोणतीही विशेष अट घातलेली नसल्यामुळे, $P$ असलेल्या प्रत्येक वस्तूत~$Q$ आहे असे म्हणता येते. म्हणजे $x$ हे~$P$ गुणधर्म असलेल्या \emph{सर्व} वस्तूंचे प्रतिनिधी नाव आहे.

\begin{prop}
सर्व $A$ आणि $B$ संचांसाठी, $A \subseteq A \cup B$.
\end{prop}

\begin{proof}
$A$ आणि $B$ हे स्वेच्छ संच असू द्या. $A \subseteq A \cup B$ दाखवायचे आहे. $\subseteq$ च्या व्याख्येनुसार याचा अर्थ असा: प्रत्येक $x$ साठी, $x \in A$ असेल तर $x \in A \cup B$ असते. म्हणून $x \in A$ हा~$A$ मधील स्वेच्छ घटक असू द्या. $x \in A \cup B$ दाखवायचे आहे. $x \in A$ असल्यामुळे $x \in A$ किंवा $x \in B$ असते. म्हणून $x \in \Setabs{x}{x \in A \lor x \in B}$. पण $\cup$ च्या व्याख्येनुसार याचाच अर्थ $x \in A \cup B$ असा होतो.
\end{proof}

\subsection{प्रकरणांनुसार सिद्धता}

गृहीतक म्हणून किंवा आधीच मिळालेला निष्कर्ष म्हणून विकल्पयोग उपलब्ध आहे असे समजा: $p$ किंवा $q$ सत्य आहे असे गृहीत धरलेले आहे किंवा सिद्ध झालेले आहे. तुम्हाला $r$ सिद्ध करायचे आहे. हे दोन पायऱ्यांत करता येते: प्रथम $p$ सत्य आहे असे गृहीत धरून~$r$ सिद्ध करा; नंतर $q$ सत्य आहे असे गृहीत धरून पुन्हा~$r$ सिद्ध करा. दोन पर्यायांपैकी एक खरा आहे असे गृहीत धरलेले किंवा माहीत असल्यामुळे ही पद्धत चालते. या दोन पायऱ्यांमुळे कोणताही एक पर्याय~$r$ च्या सत्यतेसाठी पुरेसा आहे हे मिळते. दोन्ही पर्याय खरे असतील, तर $r$~खरे असण्याची एकाऐवजी दोन कारणे असतील. $p$ आणि~$q$ दोन्ही गृहीत धरून $r$~सत्य आहे हे वेगळे सिद्ध करण्याची गरज नाही. काय करत आहोत हे दर्शवण्यासाठी “प्रकरणे वेगळी पाहू” असे जाहीर करतो. उदाहरणार्थ, $x \in B \cup C$ माहीत आहे असे समजा. $B \cup C$ ची व्याख्या $\Setabs{x}{x \in B \text{ किंवा } x \in C}$ अशी आहे. म्हणजे व्याख्येनुसार $x \in B$ किंवा $x \in C$. यावरून $x \in A$ सिद्ध करण्यासाठी प्रथम $x \in B$ गृहीत धरून त्यावरून $x \in A$ सिद्ध करू; नंतर $x \in C$ गृहीत धरून पुन्हा $x \in A$ सिद्ध करू. गृहीतकाखाली “प्रकरणे वेगळी पाहू” लिहा; त्याखाली “प्रकरण (1): $x \in B$” लिहा आणि पानाच्या मध्यावर “प्रकरण (2): $x \in C$” लिहा. मग वरच्या आणि खालच्या भागांतील सिद्धता पूर्ण करा.

सिद्ध करावयाचे विधान स्वतः विकल्पयोगात्मक असेल, तर प्रकरणांनुसार सिद्धता विशेष उपयुक्त ठरते. हे एक साधे उदाहरण:

\begin{prop}
$B \subseteq D$ आणि $C \subseteq E$ असे समजा. मग $B \cup C \subseteq D \cup E$.
\end{prop}

\begin{proof}
(a) $B \subseteq D$ आणि (b) $C \subseteq E$ असे गृहीत धरा. व्याख्येनुसार, कोणताही $x \in B$ हा $\in D$ देखील असतो (c) आणि कोणताही $x \in C$ हा $\in E$ देखील असतो (d). $B \cup C \subseteq D \cup E$ दाखवण्यासाठी $x \in B \cup C$ असेल, तर $x \in D \cup E$ असते हे दाखवावे लागते; हे $\subseteq$ च्या व्याख्येवरून मिळते. ~$ \cup$ च्या व्याख्येनुसार $x \in B \cup C$ हे $x \in B$ किंवा $x \in C$ असते तेव्हाच आणि फक्त तेव्हाच सत्य असते. त्याचप्रमाणे $x \in D \cup E$ हे $x \in D$ किंवा $x \in E$ असते तेव्हाच आणि फक्त तेव्हाच सत्य असते. म्हणून हे दाखवावे लागते: कोणत्याही $x$ साठी, $x \in B$ किंवा $x \in C$ असेल, तर $x \in D$ किंवा $x \in E$ असते.

\begin{quote}
आत्तापर्यंत फक्त व्याख्या उलगडल्या आहेत!{} आपल्या विधानाची $\subseteq$ आणि $\cup$ न वापरता पुनर्मांडणी केली आहे; आता एक वैश्विक सशर्त विधान सिद्ध करायचे उरले आहे. आधीच्या चर्चेनुसार, $x$ ही अशी वस्तू आहे की तिच्याविषयी “जर” भाग सत्य आहे असे गृहीत धरायचे आणि मग “तर” भागही सत्य आहे हे दाखवायचे. म्हणजे $x \in B$ किंवा $x \in C$ असे गृहीत धरून $x \in D$ किंवा $x \in E$ दाखवू.\footnote{हा परिच्छेद आपण काय करत आहोत याचे स्पष्टीकरण देतो; तो सिद्धतेचा भाग नाही. स्वतःच्या सिद्धता लिहिताना हा सगळा तपशील देण्याची गरज नाही.}
\end{quote}

$x \in B$ किंवा $x \in C$ असे समजा. $x \in D$ किंवा $x \in E$ दाखवायचे आहे. प्रकरणे वेगळी पाहू.

प्रकरण 1: $x \in B$. (c) वरून $x \in D$. म्हणून $x \in D$ किंवा $x \in E$. येथे एका विकल्पघटकावरून विकल्पयोग मिळवण्याचे आधी चर्चिलेले अनुमान केले आहे.

प्रकरण 2: $x \in C$. (d) वरून $x \in E$. म्हणून $x \in D$ किंवा $x \in E$.
\end{proof}

\subsection{अस्तित्ववाची विधान सिद्ध करणे}

अस्तित्ववाची विधान सिद्ध करायला सांगितले असेल, तर प्रश्नाचे रूप सहसा “असा~$x$ आहे की $\dots x \dots$, हे सिद्ध करा” असे असते. म्हणजे “$\dots x \dots$” ने वर्णन केलेला गुणधर्म असलेली वस्तू अस्तित्वात आहे, हे दाखवायचे असते. अशा वेळी योग्य वस्तू ओळखून तिच्यात आवश्यक गुणधर्म आहे हे दाखवावे लागते. हे सरळ वाटते; पण अशा सिद्धता अवघडही असू शकतात. सहसा त्यात एखादी वस्तू \emph{रचावी} किंवा \emph{परिभाषित} करावी लागते आणि अशा रीतीने परिभाषित केलेल्या वस्तूत आवश्यक गुणधर्म आहे हे सिद्ध करावे लागते. योग्य वस्तू शोधणे अवघड असू शकते; तिच्यात आवश्यक गुणधर्म आहे हे सिद्ध करणे अवघड असू शकते; आणि कधी आपण मुळात एखादी वस्तू परिभाषित करण्यात यशस्वी झालो आहोत हेच दाखवणे अवघड असते!{}

सहसा वस्तू स्पष्टपणे सांगून अशी सिद्धता लिहितो. उदाहरणार्थ, “$x$ म्हणजे \dots{} असू द्या” असे लिहितो; येथे \dots{} ने कोणती वस्तू अभिप्रेत आहे ते सांगितले जाते. आवश्यक असल्यास $\dots$ खरोखर अस्तित्वात असलेल्या वस्तूचे वर्णन करते हे सिद्ध करतो; मग $x$ मध्ये~$Q$ हा गुणधर्म आहे हे दाखवतो. हे एक साधे उदाहरण:

\begin{prop}
$x \in B$ असे समजा. मग असा~$A$ आहे की $A \subseteq B$ आणि $A \neq \emptyset$.
\end{prop}

\begin{proof}
$x \in B$ असे गृहीत धरा. $A = \{x\}$ असू द्या.
\begin{quote}
येथे~$A$ संच त्याच्या घटक ची यादी देऊन परिभाषित केला आहे. $x$ ही वस्तू आहे असे गृहीत धरले आहे; आणि सांत यादीत घटक देऊन संच तयार करता येतो. म्हणून येथे~$A$ संच परिभाषित करण्यात यश आले आहे हे वेगळे दाखवण्याची गरज नाही. मात्र $A$ मध्ये विधानाने अपेक्षित केलेले गुणधर्म आहेत हे अद्याप दाखवावे लागते. त्याशिवाय सिद्धता पूर्ण होत नाही!{}
\end{quote}
$x \in A$ असल्यामुळे $A \neq \emptyset$.
\begin{quote}
याचा आधार म्हणजे $A$ ची $\{x\}$ अशी व्याख्या आणि $x \in \{x\}$ तसेच $x \notin \emptyset$ या स्पष्ट गोष्टी.
\end{quote}
$x$ हा~$\{x\}$ मधील एकमेव घटक आहे आणि $x \in B$ आहे. म्हणून~$A$ मधील प्रत्येक घटक हा~$B$ मधील घटक देखील असतो. ~$ \subseteq$ च्या व्याख्येनुसार $A \subseteq B$.
\end{proof}

\subsection{अस्तित्ववाची विधानांचा वापर}

एखादे अस्तित्ववाची विधान सत्य आहे हे माहीत आहे असे समजा: ते सिद्ध केलेले असेल किंवा वापरता येण्याजोगे गृहीतक असेल. उदाहरणार्थ, “काही~$x$ साठी $x \in A$” किंवा “असा $x \in A$ आहे”. सिद्धतेत हे वापरायचे असेल, तर गृहीतकानुसार अस्तित्वात असलेल्या वस्तूंपैकी एका वस्तूला नाव दिले आहे असे धरता येते. $A$ मध्ये किमान एक वस्तू असल्यामुळे त्या नावाने दर्शवता येण्याजोग्या वस्तू आहेत. एखादी वस्तू नेमकी निवडून दाखवता येईलच असे नाही किंवा ती $\in A$ आहे यापलीकडे तिचे वर्णन करता येईलच असे नाही. पण सिद्धतेपुरते एका अशा वस्तूला नाव देता येते. आधी वापरलेले किंवा गृहीतकांत आलेले नाव पुन्हा वापरू नये; तसे केल्यास चुका होऊ शकतात. सिद्धतेत “काही $x$ साठी $x \in A$” यावरून “$a \in A$ असू द्या” असे लिहून हे दर्शवतो. आता~$a$ विषयी विचार करता येतो आणि इतर अनुमत गृहीतके वापरून $p$ हा निष्कर्ष मिळवता येतो. $p$ मध्ये~$a$ चा उल्लेख उरलेला नसेल आणि या नव्या नावाशी संबंधित अतिरिक्त, मुक्त न केलेले गृहीतक वापरलेले नसेल, तर $p$ मिळवण्यासाठीचे तात्पुरते $a \in A$ हे गृहीतक मुक्त करता येते. त्यामुळे “काही $x$ साठी $x \in A$” आणि वापरलेली इतर अनुमत गृहीतके यांवरून निष्कर्ष मिळतो.

\begin{prop}
$A \neq \emptyset$ असेल, तर $A \cup B \neq \emptyset$.
\end{prop}

\begin{proof}
$A \neq \emptyset$ असे समजा. म्हणून काही $x$ साठी $x \in A$.
\begin{quote}
येथे प्रथम विधानाचे गृहीतक फक्त पुन्हा मांडले आहे. या गृहीतकात, म्हणजे $A \neq \emptyset$ यात, एक अस्तित्ववाची विधान दडलेले आहे. काही व्याख्या उलगडल्यावर ते दिसते. $=$ च्या व्याख्येनुसार $A = \emptyset$ हे तेव्हाच आणि फक्त तेव्हा सत्य असते जेव्हा प्रत्येक $x \in A$ हा $\in \emptyset$ देखील असतो आणि प्रत्येक $x \in \emptyset$ हा $\in A$ देखील असतो. दोन्ही बाजूंचे नकरण केल्यावर हे मिळते: $A \neq \emptyset$ तेव्हाच आणि फक्त तेव्हा जेव्हा काही $x \in A$ हा $\notin \emptyset$ असतो किंवा काही $x \in \emptyset$ हा $\notin A$ असतो. कोणतीही वस्तू $\in \emptyset$ नसल्यामुळे दुसरा विकल्पघटक कधीही सत्य असू शकत नाही. तसेच “$x \in A$ आणि $x \notin \emptyset$” याचा अर्थ फक्त $x \in A$ इतकाच उरतो. म्हणून $A \neq \emptyset$ हे काही $x$ साठी $x \in A$ असते तेव्हाच आणि फक्त तेव्हाच सत्य असते. हे अस्तित्ववाची विधान आहे. आता~$A$ मधील घटक पैकी एकाला नाव देऊन या अस्तित्ववाची विधानाचा वापर करू:
\end{quote}
$a \in A$ असू द्या.
\begin{quote}
आता~$\in A$ असलेल्या एका वस्तूला नाव दिले आहे. पुढे~$a$ विषयी युक्तिवाद करू; पण केवळ $a \in A$ आहे एवढेच गृहीत धरू, आणखी काही नाही, याची काळजी घेऊ:
\end{quote}
$a \in A$ असल्यामुळे~$\cup$ च्या व्याख्येवरून $a \in A \cup B$. म्हणून काही $x$ साठी $x \in A \cup B$; म्हणजे $A \cup B \neq \emptyset$.
\begin{quote}
शेवटच्या पायरीत “$a \in A \cup B$” यावरून “काही $x$ साठी $x \in A \cup B$” असा निष्कर्ष काढला. त्यात आता $a$ चा उल्लेख नाही. म्हणून “काही $x$ साठी $x \in A \cup B$” हे केवळ “काही $x$ साठी $x \in A$” यावरून मिळते हे समजते. पण याचाच अर्थ $A \cup B \neq \emptyset$ असा होतो.
\end{quote}
\end{proof}

आपण केले तसे “$x$” सारखे बद्ध चल आणि $a$ सारखी तात्पुरती नावे वेगळी ठेवणे हा चांगला सराव आहे. प्रत्यक्षात मात्र अनेकदा असे करत नाही आणि $x$ हेच वापरतो:
\begin{quote}
$A \neq \emptyset$ असे समजा; म्हणजे असा $x \in A$ आहे. $\cup$ च्या व्याख्येनुसार $x \in A \cup B$. म्हणून $A \cup B \neq \emptyset$.
\end{quote}
मात्र असे करताना वेगवेगळ्या अस्तित्ववाची विधानांसाठी वेगवेगळी $x$ आणि $y$ नावे वापरतो आहोत याची विशेष काळजी घ्यावी. उदाहरणार्थ, पुढील युक्तिवाद ही “$A \neq \emptyset$ आणि $B \neq \emptyset$ असतील, तर $A \cap B \neq \emptyset$” याची योग्य सिद्धता \emph{नाही}; ते विधानच खरे नाही.
\begin{quote}
$A \neq \emptyset$ आणि $B \neq \emptyset$ असे समजा. म्हणून काही $x$ साठी $x \in A$ आणि काही $x$ साठी $x \in B$ देखील असते. $x \in A$ आणि $x \in B$ असल्यामुळे~$\cap$ च्या व्याख्येनुसार $x \in A \cap B$. म्हणून $A \cap B \neq \emptyset$.
\end{quote}
चुकीची पायरी कुठे येते ते ओळखता येईल का? हा निष्कर्ष का खरा नाही हे स्पष्ट करता येईल का?












\section{एक उदाहरण}\label{mth:prf:ex1:sec}

पहिले उदाहरण म्हणजे संचांच्या संयोग आणि छेदाविषयीचे पुढील साधे तथ्य. त्यातून व्याख्या उलगडणे, संधियोग आणि वैश्विक विधाने सिद्ध करणे, तसेच प्रकरणांनुसार सिद्धता करणे या पद्धती स्पष्ट होतील.

\begin{prop}
कोणत्याही $A$, $B$ आणि $C$ संचांसाठी, $A \cup (B \cap C) = (A \cup B) \cap (A \cup C)$.
\end{prop}

हे सिद्ध करू!{}

\begin{proof}
कोणत्याही $A$, $B$ आणि $C$ संचांसाठी $A \cup (B \cap C) = (A \cup B) \cap (A \cup C)$ दाखवायचे आहे.
\begin{quote}
प्रथम विधानातील “$=$” ची व्याख्या उलगडू. संच एकरूप आहेत हे सिद्ध करणे म्हणजे त्यांचे घटक तेच आहेत हे दाखवणे, हे आठवा. म्हणजे $A \cup (B \cap C)$ मधील सर्व घटक हे $(A \cup B) \cap (A \cup C)$ मधील घटक देखील असले पाहिजेत आणि उलटही तसेच असले पाहिजे. “उलटही तसेच” म्हणजे $(A \cup B) \cap (A \cup C)$ मधील प्रत्येक घटक हा $A \cup (B \cap C)$ मधील घटक देखील असला पाहिजे. म्हणून व्याख्या उलगडल्यावर संधियोग सिद्ध करावा लागतो असे दिसते. हे नोंदवू:
\end{quote}
व्याख्येनुसार, $A \cup (B \cap C) = (A \cup B) \cap (A \cup C)$ हे तेव्हाच आणि फक्त तेव्हा सत्य असते जेव्हा $A \cup (B \cap C)$ मधील प्रत्येक घटक हा $(A \cup B) \cap (A \cup C)$ मधील घटक देखील असतो आणि $(A \cup B) \cap (A \cup C)$ मधील प्रत्येक घटक हा $A \cup (B \cap C)$ मधील घटक असतो.
\begin{quote}
हा संधियोग असल्यामुळे प्रत्येक संधिघटक वेगळा सिद्ध करावा लागेल. पहिल्यापासून सुरुवात करू: $A \cup (B \cap C)$ मधील प्रत्येक घटक हा $(A \cup B) \cap (A \cup C)$ मधील घटक देखील असतो हे सिद्ध करू.

हे वैश्विक विधान आहे. म्हणून $A \cup (B \cap C)$ मधील एक स्वेच्छ घटक घेऊन तो $(A \cup B) \cap (A \cup C)$ मधील घटक देखील असला पाहिजे हे दाखवू. या स्वेच्छ घटक साठी एखादे चल निवडू; उदाहरणार्थ,~$z$. सिद्धता पुढे अशी चालते:
\end{quote}
प्रथम $A \cup (B \cap C)$ मधील प्रत्येक घटक हा $(A \cup B) \cap (A \cup C)$ मधील घटक देखील असतो हे सिद्ध करू. $z \in A \cup (B \cap C)$ असू द्या. $z \in (A \cup B) \cap (A \cup C)$ दाखवायचे आहे.
\begin{quote}
आता $\cup$ आणि~$\cap$ यांच्या व्याख्या उलगडायच्या आहेत. उदाहरणार्थ, $\cup$ ची व्याख्या अशी आहे: $A \cup B = \Setabs{z}{z \in A \text{ किंवा } z \in B}$. “$A \cup (B \cap C)$” ला ही व्याख्या लावताना व्याख्येतील “$B$” ची भूमिका आता “$B \cap C$” बजावते. म्हणून $A \cup (B \cap C) = \Setabs{z}{z \in A \text{ किंवा } z \in B \cap C}$. त्यामुळे $z \in A \cup (B \cap C)$ या गृहीतकाचा अर्थ $z \in \Setabs{z}{z \in A \text{ किंवा } z \in B \cap C}$ असा होतो. तसेच $z \in \Setabs{z}{\dots z\dots}$ हे \dots $z$ \dots{} असते तेव्हाच आणि फक्त तेव्हाच सत्य असते. म्हणजे येथे $z \in A$ किंवा $z \in B \cap C$ असते.
\end{quote}
$\cup$ च्या व्याख्येनुसार $z \in A$ किंवा $z \in B \cap C$ असते.
\begin{quote}
हा विकल्पयोग असल्यामुळे प्रकरणांनुसार सिद्धता उपयुक्त ठरेल. दोन्ही प्रकरणे घेऊन प्रत्येक प्रकरणात अपेक्षित निष्कर्ष---म्हणजे “$z \in (A \cup B) \cap (A \cup C)$”---मिळतो हे दाखवू.
\end{quote}
प्रकरण 1: $z \in A$ असे समजा.
\begin{quote}
या गृहीतकांवरून आणखी फार काही मिळत नाही. म्हणून अपेक्षित निष्कर्षात कोणती रचना आहे ते पाहू. $z \in (A \cup B) \cap (A \cup C)$ दाखवायचे आहे. $\cap$ च्या व्याख्येनुसार $z \in (A \cup B) \cap (A \cup C)$ दाखवण्यासाठी तो $(A \cup B)$ आणि $(A \cup C)$ या दोन्हींत आहे हे दाखवावे लागते. पण $z \in A \cup B$ हे $z \in A$ किंवा $z \in B$ असते तेव्हाच आणि फक्त तेव्हाच सत्य असते. तसेच प्रकरण~1 चे गृहीतक $z \in A$ आधीच उपलब्ध आहे. हाच युक्तिवाद $B$ ऐवजी $C$ वापरून केल्यास $z \in A \cup C$ मिळते. हा विचार अपेक्षित निष्कर्षापासून मागे जात केलेला होता; आता सिद्धतेत आवश्यक असलेल्या पुढच्या दिशेने तो नोंदवू.
\end{quote}
$z \in A$ असल्यामुळे $z \in A$ किंवा $z \in B$. म्हणून~$\cup$ च्या व्याख्येनुसार $z \in A \cup B$. त्याचप्रमाणे $z \in A \cup C$. पण~$\cap$ च्या व्याख्येनुसार याचाच अर्थ $z \in (A \cup B) \cap (A \cup C)$ असा होतो.
\begin{quote}
यामुळे प्रकरणांनुसार सिद्धतेचे पहिले प्रकरण पूर्ण झाले. आता $z \in B \cap C$ असलेल्या दुसऱ्या प्रकरणात निष्कर्ष मिळवायचा आहे.
\end{quote}
प्रकरण 2: $z \in B \cap C$ असे समजा.
\begin{quote}
पुन्हा दोन संचांचा छेद हाताळत आहोत. ~$ \cap$ ची व्याख्या वापरू:
\end{quote}
$z \in B \cap C$ असल्यामुळे~$\cap$ च्या व्याख्येनुसार $z$ हा $B$ आणि $C$ या दोन्हींचा घटक असला पाहिजे.
\begin{quote}
पुन्हा अपेक्षित निष्कर्ष पाहू. $z$ हा $(A \cup B)$ आणि $(A \cup C)$ या दोन्हींत आहे हे दाखवायचे आहे. पुन्हा उत्तर लगेच मिळते.
\end{quote}
$z \in B$ असल्यामुळे $z \in (A \cup B)$. $z \in C$ असल्यामुळे $z \in (A \cup C)$ देखील. म्हणून $z \in (A \cup B) \cap (A \cup C)$.
\begin{quote}
येथे पुन्हा $\cup$ आणि $\cap$ यांच्या व्याख्या वापरल्या. पण त्या व्याख्या आधीच आठवून घेतल्या आहेत आणि $z$ हा दोन संचांपैकी एका संचात असेल तर त्यांच्या संयोगात असतो, हेही आधीच दाखवले आहे. म्हणून आता प्रत्येक वापराचा तितका तपशील लिहिण्याची गरज नाही.

प्रकरणांनुसार सिद्धतेचे दुसरे प्रकरण पूर्ण झाले. म्हणून पहिला अपेक्षित निष्कर्ष आता सांगता येतो.
\end{quote}
म्हणून $z \in A \cup (B \cap C)$ असेल, तर $z \in (A \cup B) \cap (A \cup C)$.
\begin{quote}
आता उलटी दिशा दाखवायची आहे: $(A \cup B) \cap (A \cup C)$ मधील प्रत्येक घटक हा $A \cup (B \cap C)$ मधील घटक आहे. आधीप्रमाणे, पहिल्या संचातील एक स्वेच्छ घटक घेतला आहे असे गृहीत धरून तो दुसऱ्या संचात असला पाहिजे हे दाखवून हे वैश्विक विधान सिद्ध करू. काय करणार आहोत ते मांडू.
\end{quote}
आता $z \in (A \cup B) \cap (A \cup C)$ असे गृहीत धरा. $z \in A \cup (B \cap C)$ दाखवायचे आहे.
\begin{quote}
आता $z \in (A \cup B) \cap (A \cup C)$ या गृहीतकावरून विचार करत आहोत. सिद्धतेच्या पहिल्या भागात वापरलेले~$z$ हेच येथे वापरल्यामुळे फार गोंधळ होत नाही अशी आशा आहे. तो भाग पूर्ण झाल्यावर त्या भागातील तात्पुरती गृहीतके लागू राहिली नाहीत; म्हणून आता $z$ विषयी नवी गृहीतके घेता येतात. हे गोंधळात टाकत असेल, तर पुढील भागात $z$ ऐवजी वेगळे चल वापरा.

~$\cap$ च्या व्याख्येनुसार $z$ हा $A \cup B$ आणि $A \cup C$ या दोन्हींत आहे. पुढे $\cup$ ची व्याख्या उलगडल्यावर असे मिळते: $z \in A$ किंवा $z \in B$, आणि $z \in A$ किंवा $z \in C$ देखील. पुन्हा प्रकरणांनुसार सिद्धता दिसते; पण “आणि” मुळे गोंधळ होऊ शकतो. दोन्ही अटींऐवजी एकच अट पुरेशी आहे---$z$ हा $A$, $B$ किंवा $C$ पैकी एखाद्यात आहे---असे वाटू शकते. पण ती अदलाबदल चुकीची ठरेल. म्हणून काळजीपूर्वक एकेक विकल्पयोग पाहू.
\end{quote}
$\cap$ च्या व्याख्येनुसार $z \in A \cup B$ आणि $z \in A \cup C$. $\cup$ च्या व्याख्येनुसार $z \in A$ किंवा $z \in B$. प्रकरणे वेगळी पाहू.
\begin{quote}
पहिल्या विकल्पयोगावर लक्ष केंद्रित केल्यामुळे $A \cup C$ उलगडून मिळणारा दुसरा विकल्पयोग अजून नोंदवलेला नाही. आत्ता त्याची गरजही नाही. पहिले प्रकरण $z \in A$ हे आहे. संचाचा घटक हा त्या संचाचा दुसऱ्या कोणत्याही संचाशी संयोग केल्यावर मिळणाऱ्या संचाचा घटक देखील असतो. म्हणून प्रकरण~1 सोपे आहे:
\end{quote}
प्रकरण 1: $z \in A$ असे समजा. यावरून $z \in A \cup (B \cap C)$ मिळते.
\begin{quote}
आता दुसरे प्रकरण, $z \in B$, पाहू. येथे दुसऱ्या $\cup$ ची व्याख्या उलगडून आणखी एक प्रकरणांनुसार सिद्धता करू:
\end{quote}
प्रकरण 2: $z \in B$ असे समजा. $z \in A \cup C$ असल्यामुळे $z \in A$ किंवा $z \in C$. आणखी प्रकरणे वेगळी पाहू:

प्रकरण 2a: $z \in A$. मग पुन्हा $z \in A \cup (B \cap C)$.
\begin{quote}
हे थोडे विचित्र वाटेल: या प्रकरणात~$z \in B$ हे गृहीतक वापरण्याची गरजच पडली नाही. पण त्यात काही अडचण नाही.
\end{quote}
प्रकरण 2b: $z \in C$. मग $z \in B$ आणि $z \in C$. म्हणून $z \in B \cap C$ आणि परिणामी $z \in A \cup (B \cap C)$.
\begin{quote}
दोन्ही प्रकरणांनुसार सिद्धता पूर्ण झाल्या; त्यामुळे दुसरा अर्धा भागही पूर्ण झाला.
\end{quote}
म्हणून $z \in (A \cup B) \cap (A \cup C)$ असेल, तर $z \in A \cup (B \cap C)$.
\end{proof}












\section{आणखी एक उदाहरण}\label{mth:prf:ex2:sec}

\begin{prop}
$A \subseteq C$ असेल, तर $A \cup (C \setminus A) = C$.
\end{prop}

\begin{proof}
$A \subseteq C$ असे समजा. $A \cup (C \setminus A) = C$ दाखवायचे आहे.
\begin{quote}
सुरुवातीला हे सशर्त विधान आहे हे लक्षात घेऊ. त्याचे वैश्विक संख्यापन अप्रकट आहे: हे विधान सर्व $A$ आणि $C$ संचांसाठी खरे आहे. म्हणून $A$ आणि $C$ हे स्वेच्छ संचांचे चल आहेत. असे विधान सिद्ध करण्यासाठी पूर्वांग गृहीत धरून उत्तरांग सिद्ध करतो.

आता $A \subseteq C$ हे गृहीतक वापरू. ~$\subseteq$ ची व्याख्या उलगडू: गृहीतकाचा अर्थ असा की~$A$ मधील सर्व घटक हे~$C$ मधील घटक देखील आहेत. हे नोंदवू; संपूर्ण सिद्धतेत वापरणार असलेले हे महत्त्वाचे तथ्य आहे.
\end{quote}
~$\subseteq$ च्या व्याख्येनुसार, $A \subseteq C$ असल्यामुळे प्रत्येक $z$ साठी $z \in A$ असेल, तर $z \in C$ असते.
\begin{quote}
गृहीतकातील सर्व व्याख्या उलगडल्या आहेत. आता निष्कर्षाकडे वळू. $A \cup (C \setminus A) = C$ दाखवायचे आहे. म्हणून मागच्या उदाहरणाप्रमाणे सिद्धतेची मांडणी करू: $A \cup (C \setminus A)$ मधील प्रत्येक घटक हा~$C$ मधील घटक आहे आणि उलट $C$ मधील प्रत्येक घटक हा $A \cup (C \setminus A)$ मधील घटक आहे, हे दाखवू. हे थोडक्यात असे लिहिता येईल: $A \cup (C \setminus A) \subseteq C$ आणि $C \subseteq A \cup (C \setminus A)$. येथे व्याख्या उलगडण्याच्या उलट क्रिया करत आहोत; पण त्यामुळे सिद्धता वाचणे थोडे सोपे होते. हा संधियोग असल्यामुळे दोन्ही भाग सिद्ध करावे लागतात. पहिला भाग---म्हणजे $A \cup (C \setminus A)$ मधील प्रत्येक घटक हा~$C$ मधील घटक आहे---दाखवण्यासाठी, स्वेच्छ~$z$ साठी $z \in A \cup (C \setminus A)$ गृहीत धरून $z \in C$ दाखवू. $\cup$ च्या व्याख्येनुसार $z \in A \cup (C \setminus A)$ यावरून $z \in A$ किंवा $z \in C \setminus A$ असा निष्कर्ष काढता येतो. आता तुम्हाला या पद्धतीचा अंदाज येत असेल.
\end{quote}
$A \cup (C \setminus A) = C$ हे तेव्हाच आणि फक्त तेव्हा सत्य असते जेव्हा $A \cup (C \setminus A) \subseteq C$ आणि $C \subseteq (A \cup (C \setminus A))$ असतात. प्रथम $A \cup (C \setminus A) \subseteq C$ सिद्ध करू. $z \in A \cup (C \setminus A)$ असू द्या. मग $z \in A$ किंवा $z \in (C \setminus A)$.
\begin{quote}
एक विकल्पयोग मिळाला आहे; त्यावरून $z \in C$ सिद्ध करायचे आहे. यासाठी प्रकरणांनुसार सिद्धता करू.
\end{quote}
प्रकरण 1: $z \in A$. प्रत्येक $z$ साठी $z \in A$ असेल तर $z \in C$ असते, हे माहीत असल्यामुळे $z \in C$ मिळते.
\begin{quote}
येथे $A \subseteq C$ या गृहीतकावरून मिळालेले आणि आधी नोंदवलेले तथ्य वापरले. पहिले प्रकरण पूर्ण झाले; आता दुसऱ्या प्रकरणाकडे, $z \in (C \setminus A)$, वळू. $C \setminus A$ हे दोन संचांचा \emph{फरक} दर्शवते, हे आठवा:~$C$ मधील जे घटक हे~$A$ मधील घटक नाहीत त्यांचा हा संच आहे. पण $C$ मधील जो घटक हा~$A$ मध्ये नाही तो विशेषतः~$C$ मधील घटक असतोच.
\end{quote}
प्रकरण 2: $z \in (C \setminus A)$. याचा अर्थ $z \in C$ आणि $z \notin A$. म्हणून विशेषतः $z \in C$.
\begin{quote}
पहिली दिशा सिद्ध झाली. आता दुसरी दिशा पाहू. येथे $C \subseteq A \cup (C \setminus A)$ सिद्ध करायचे आहे. म्हणून $z \in C$ गृहीत धरून $z \in A \cup (C \setminus A)$ सिद्ध करू.
\end{quote}
आता $z \in C$ असू द्या. $z \in A$ किंवा $z \in C \setminus A$ दाखवायचे आहे.
\begin{quote}
$A$ मधील सर्व घटक हे $C$ मधील घटक देखील आहेत; आणि $C \setminus A$ हा $C$ मधील, पण $A$ मधील नसलेल्या, घटक चा संच आहे. म्हणून $z$ हा $A$ मध्ये किंवा $C \setminus A$ मध्ये असतो असे दिसते. निष्कर्ष का खरा आहे हे आधीच माहीत नसेल, तर हा युक्तिवाद थोडा अस्पष्ट वाटेल. तो पायरीपायरीने सिद्ध करणे अधिक चांगले. सिद्धता न देता मांडता येणारे एक साधे तथ्य वापरल्यास मदत होईल: $z \in A$ किंवा $z \notin A$. याला “विमध्य सिद्धांत” म्हणतात: कोणत्याही~$p$ विधानासाठी $p$ सत्य असते किंवा त्याचे नकरण सत्य असते. येथे $p$ हे $z \in A$ हे विधान आहे. हा विकल्पयोग असल्यामुळे पुन्हा प्रकरणांनुसार सिद्धता वापरता येते.
\end{quote}
$z \in A$ किंवा $z \notin A$. पहिल्या प्रकरणात $z \in A \cup (C \setminus A)$. दुसऱ्या प्रकरणात $z \in C$ आणि $z \notin A$. म्हणून $z \in C \setminus A$. पण मग $z \in A \cup (C \setminus A)$.
\begin{quote}
सिद्धता पूर्ण झाली: $A \cup (C \setminus A) = C$ दाखवले आहे.
\end{quote}
\end{proof}












\section{व्याघाताने सिद्धता}\label{mth:prf:con:sec}

व्याघाताने सिद्धता हा प्रथमतः नकारात्मक विधाने सिद्ध करण्यासाठी वापरला जाणारा अनुमानाचा नमुना आहे. एखादे~$p$ विधान \emph{असत्य} आहे हे दाखवायचे आहे असे समजा; म्हणजे~$\lnot p$ दाखवायचे आहे. यासाठी सर्वात उपयोगी मार्ग असा: (a) $p$~सत्य आहे असे समजा आणि (b) या गृहीतकावरून असत्य असल्याचे माहीत असलेली एखादी गोष्ट मिळते हे दाखवा. “असत्य असल्याचे माहीत असलेली गोष्ट” ही $p$ शीच व्याघात करणारा निष्कर्ष असेल किंवा संपूर्ण विधानाच्या दुसऱ्या एखाद्या गृहीतकाशी व्याघात करणारा निष्कर्ष असेल. उदाहरणार्थ, “$q$ असेल, तर $\lnot p$” सिद्ध करताना $q$~सत्य आहे असे गृहीत धरून त्यावरून~$\lnot p$ सिद्ध करतो. $\lnot p$ व्याघाताने सिद्ध करायचे असेल, तर~$q$ बरोबर $p$ हेदेखील गृहीत धरायचे. $p$ वरून $\lnot q$ सिद्ध करता आले, तर~$p$ या गृहीतकावरून दुसऱ्या~$q$ गृहीतकाशी व्याघात करणारा निष्कर्ष मिळतो हे दाखवले आहे. कारण $q$~आणि $\lnot q$ दोन्ही सत्य असू शकत नाहीत. अर्थात, व्याघात सिद्ध करताना इतर अनुमानाचे नमुने वापरावे लागतील आणि व्याख्याही उलगडाव्या लागतील. एक उदाहरण पाहू.

\begin{prop}
$A \subseteq B$ आणि $B = \emptyset$ असतील, तर $A$ मध्ये कोणतेही घटक नसतात.
\end{prop}

\begin{proof}
$A \subseteq B$ आणि $B = \emptyset$ असे समजा. $A$ मध्ये कोणतेही घटक नाहीत हे दाखवायचे आहे.
\begin{quote}
हे सशर्त विधान असल्यामुळे पूर्वांग गृहीत धरून उत्तरांग सिद्ध करायचे आहे. उत्तरांग असे आहे: $A$ मध्ये कोणतेही घटक नाहीत. हे आणखी स्पष्टपणे असे मांडता येते: असा~$x \in A$ आहे हे खरे नाही.
\end{quote}
$A$ मध्ये कोणतेही घटक नसतात तेव्हाच आणि फक्त तेव्हा जेव्हा असा~$x$ आहे की $x \in A$, हे खरे नसते.
\begin{quote}
म्हणून प्रत्यक्षात~$\lnot p$ हे नकारात्मक विधान सिद्ध करायचे आहे असे ठरले: असा $x \in A$ आहे हे खरे नाही. व्याघाताने सिद्धता करण्यासाठी त्याच्याशी संबंधित~$p$ हे होकारात्मक विधान---म्हणजे असा $x \in A$ आहे---गृहीत धरून त्यावरून व्याघात सिद्ध करावा लागेल. “व्याघात मिळवण्यासाठी असे गृहीत धरा” किंवा नुसते “उलट समजा” असे लिहून~$p$ हे गृहीतक मांडले, की आपण व्याघाताने सिद्धता करत आहोत हे दर्शवता येते.
\end{quote}
उलट समजा: असा $x \in A$ आहे.
\begin{quote}
व्याघात मिळवण्यासाठी हे नवे गृहीतक वापरू. आणखी दोन गृहीतके आहेत: $A \subseteq B$ आणि $B = \emptyset$. पहिल्यावरून $x \in B$ मिळते:
\end{quote}
$A \subseteq B$ असल्यामुळे $x \in B$.
\begin{quote}
पण $B = \emptyset$ असल्यामुळे $B$ मधील प्रत्येक घटक---उदाहरणार्थ $x$---हा~$\emptyset$ मधील घटक देखील असला पाहिजे.
\end{quote}
$B = \emptyset$ असल्यामुळे $x \in \emptyset$. हा व्याघात आहे; कारण व्याख्येनुसार $\emptyset$ मध्ये कोणतेही घटक नसतात.
\begin{quote}
यानेच सिद्धता पूर्ण होते. घेतलेल्या गृहीतकांवरून आवश्यक व्याघात मिळाला आहे; म्हणजे ती सर्व गृहीतके एकाच वेळी खरी असू शकत नाहीत. पहिली दोन गृहीतके ($A \subseteq B$ आणि $B = \emptyset$) कायम ठेवलेली असल्यामुळे शेवटी आणलेले गृहीतक---असा $x \in A$ आहे---असत्य असले पाहिजे. मात्र पूर्ण तपशील द्यायचा असल्यास हे स्पष्ट लिहिता येते.
\end{quote}
म्हणून असा $x \in A$ आहे हे गृहीतक असत्य असले पाहिजे. त्यामुळे व्याघाताने सिद्धतेवरून $A$ मध्ये कोणतेही घटक नाहीत.
\end{proof}

शास्त्रीय तर्कशास्त्रात प्रत्येक होकारात्मक विधान एका नकारात्मक विधानाशी सममूल्य असते: $p$ हे $\lnot\lnot p$ असते तेव्हाच आणि फक्त तेव्हाच सत्य असते. म्हणून व्याघाताने सिद्धता वापरून होकारात्मक विधानेही “अप्रत्यक्षपणे” सिद्ध करता येतात. ~$p$ सिद्ध करण्यासाठी त्याला $\lnot\lnot p$ हे नकारात्मक विधान समजा. $\lnot p$ वरून व्याघात सिद्ध करता आला, तर व्याघाताने सिद्धतेद्वारे $\lnot\lnot p$ प्रस्थापित होते आणि म्हणून~$p$ मिळते.

मागील उदाहरणात नकारात्मक विधान सिद्ध करायचे होते: $A$ मध्ये कोणतेही घटक नाहीत. म्हणून व्याघाताने सिद्धतेसाठी घेतलेले गृहीतक---म्हणजे असा $x \in A$ आहे---हे होकारात्मक होते. त्यामुळे विचार करण्यासाठी तात्पुरता $x \in A$ मिळाला; त्याच्याविषयी विचार पुढे नेत $x \in \emptyset$ पर्यंत पोहोचलो.

होकारात्मक विधान अप्रत्यक्षपणे सिद्ध करताना व्याघात मिळवण्यासाठी नकारात्मक गृहीतक घ्यावे लागते. पण अनेकदा होकारात्मक विधान नकारात्मक रूपात आणि नकारात्मक विधान होकारात्मक रूपात सहज पुन्हा मांडता येते. मागच्या सिद्धतेत “$A$ मध्ये कोणतेही घटक नाहीत” या नकारात्मक उत्तरांगाऐवजी “$A = \emptyset$” सिद्ध केले असते, तरी सिद्धता मूलतः तशीच असती. $=$ च्या व्याख्येनुसार “$A = \emptyset$” हे सामान्य विधान आहे; कारण ते उलगडल्यावर “~$A$ मधील प्रत्येक घटक हा~$\emptyset$ मधील घटक आहे आणि उलटही तसेच आहे” असे मिळते. पण ते “असा $x \in A$ आहे हे खरे नाही” या नकारात्मक विधानाशी सममूल्य आहे हे सहज दिसते.

म्हणून कधी~$p$ प्रत्यक्ष सिद्ध करण्यापेक्षा $\lnot p$ हे गृहीतक घेऊन काम करणे सोपे असते. पुढील उदाहरणांप्रमाणे प्रत्यक्ष सिद्धता तितकीच सोपी किंवा अधिक सोपी असली, तरी काही जण अप्रत्यक्ष मार्ग पसंत करतात. द्विनकरणामुळे गोंधळ होत असेल, तर व्याघाताने सिद्धता म्हणजे \emph{विरुद्ध} विधान गृहीत धरून व्याघात सिद्ध करणे, असा विचार करा. म्हणून $\lnot p$ व्याघाताने सिद्ध करणे म्हणजे~$p$ गृहीत धरून व्याघात सिद्ध करणे; आणि~$p$ व्याघाताने सिद्ध करणे म्हणजे~$\lnot p$ वरून व्याघात सिद्ध करणे.

\begin{prop}
$A \subseteq A \cup B$.
\end{prop}

\begin{proof}
$A \subseteq A \cup B$ दाखवायचे आहे.
\begin{quote}
प्रथमदर्शनी हे होकारात्मक विधान आहे: प्रत्येक $x \in A$ हा $A \cup B$ मध्येही असतो. त्याचे नकरण असे: काही $x \in A$ हा $\notin A \cup B$ असतो. म्हणून हे नकरण गृहीत धरून त्यावरून व्याघात मिळतो असे दाखवल्यास विधान अप्रत्यक्षपणे सिद्ध होते.
\end{quote}
उलट समजा; म्हणजे $A \nsubseteq A \cup B$.
\begin{quote}
$A \subseteq A \cup B$ ची व्याख्या अशी आहे: प्रत्येक $x \in A$ हा $\in A \cup B$ देखील असतो. $A \nsubseteq A \cup B$ चा अर्थ समजण्यासाठी उलगडलेल्या व्याख्येवर काही मूलभूत तार्किक क्रिया कराव्या लागतात: प्रत्येक $x \in A$ हा $\in A \cup B$ देखील असतो हे असत्य असते तेव्हाच आणि फक्त तेव्हा जेव्हा असा \emph{काही}~$x \in A$ असतो की तो $\notin A \cup B$ असतो. येथे फार काळजी घ्यावी: “सर्व $A$ हे $B$ आहेत” अशा विधानाचे नकरण चुकीचे उलगडल्यामुळे अनेक विद्यार्थ्यांच्या व्याघाताने सिद्धता करण्याच्या प्रयत्नांत चूक होते. म्हणजे $A \nsubseteq A \cup B$ हे तेव्हाच आणि फक्त तेव्हा सत्य असते जेव्हा असा $x$ आहे की $x \in A$ आणि $x \notin A \cup B$. यापुढे सिद्धता सोपी आहे.
\end{quote}
म्हणून असा $x \in A$ आहे की $x \notin A \cup B$. $\cup$ च्या व्याख्येनुसार $x \in A \cup B$ हे $x \in A$ किंवा $x \in B$ असते तेव्हाच आणि फक्त तेव्हाच सत्य असते. $x \in A$ असल्यामुळे $x \in A \cup B$ मिळते. हा $x \notin A \cup B$ या गृहीतकाशी व्याघात आहे.
\end{proof}

\begin{prob}
$A \cap B \subseteq A$ हे \emph{अप्रत्यक्षपणे} सिद्ध करा.
\end{prob}

\begin{prop}
$A \subseteq B$ आणि $B \subseteq C$ असतील, तर $A \subseteq C$.
\end{prop}

\begin{proof}
$A \subseteq B$ आणि $B \subseteq C$ असे समजा. $A \subseteq C$ दाखवायचे आहे.
\begin{quote}
अप्रत्यक्ष मार्ग वापरू: जे प्रस्थापित करायचे आहे त्याचे नकरण गृहीत धरू.
\end{quote}
उलट समजा; म्हणजे $A \nsubseteq C$.
\begin{quote}
आधीप्रमाणे, $A \nsubseteq C$ हे तेव्हाच आणि फक्त तेव्हा सत्य असते जेव्हा प्रत्येक $x \in A$ हा $\in C$ देखील असतो हे खरे नसते; म्हणजे काही $x \in A$ हा $\notin C$ असतो. सराव झाल्यावर अशी नकरणे उलगडण्यासाठी फार विचार करावा लागणार नाही.
\end{quote}
म्हणजे असा~$x$ आहे की $x \in A$ आणि $x \notin C$.
\begin{quote}
आता यावरून व्याघात मिळवता येईल. अर्थात, त्यासाठी इतर दोन गृहीतके वापरावी लागतील.
\end{quote}
$A \subseteq B$ असल्यामुळे $x \in B$. $B \subseteq C$ असल्यामुळे $x \in C$. पण हा $x \notin C$ शी व्याघात आहे.
\end{proof}

\begin{prop}
$A \cup B = A \cap B$ असेल, तर $A = B$.
\end{prop}

\begin{proof}
$A \cup B = A \cap B$ असे समजा. $A = B$ दाखवायचे आहे.
\begin{quote}
आता सुरुवातीची पद्धत परिचित आहे:
\end{quote}
व्याघात मिळवण्यासाठी $A \neq B$ असे गृहीत धरा.
\begin{quote}
व्याघाताने सिद्धतेसाठीचे गृहीतक $A \neq B$ आहे. $A = B$ हे $A \subseteq B$ आणि $B \subseteq A$ असतात तेव्हाच आणि फक्त तेव्हाच सत्य असल्यामुळे, $A \neq B$ हे $A \nsubseteq B$ \emph{किंवा} $B \nsubseteq A$ असते तेव्हाच आणि फक्त तेव्हाच सत्य असते. नकरणांवरील क्रिया काळजीपूर्वक करणे किती महत्त्वाचे आहे ते पाहा!{} या विकल्पयोगावरून व्याघात सिद्ध करण्यासाठी प्रकरणांनुसार सिद्धता वापरून प्रत्येक प्रकरणात व्याघात मिळतो हे दाखवू.
\end{quote}
$A \neq B$ हे $A \nsubseteq B$ किंवा $B \nsubseteq A$ असते तेव्हाच आणि फक्त तेव्हाच सत्य असते. प्रकरणे वेगळी पाहू.
\begin{quote}
पहिल्या प्रकरणात $A \nsubseteq B$ गृहीत धरतो; म्हणजे काही $x$ साठी $x \in A$, पण $\notin B$. $A \cap B$ ची व्याख्या $A$ आणि $B$ यांच्या समान घटक चा संच अशी आहे. म्हणून एखादी वस्तू त्यांपैकी एका संचात नसेल, तर ती छेदात नाही. $A \cup B$ हा $A$ बरोबर $B$ घेऊन मिळालेला संच आहे. म्हणून कोणत्याही एका संचातील वस्तू संयोगातही असते. यावरून $x \in A \cup B$, पण $x \notin A \cap B$, असे मिळते; म्हणून $A \cap B \neq A \cup B$.
\end{quote}

प्रकरण 1: $A \nsubseteq B$. मग काही $x$ साठी $x \in A$, पण $x \notin B$. $x \notin B$ असल्यामुळे $x \notin A \cap B$. $x \in A$ असल्यामुळे $x \in A \cup B$. म्हणून $A \cap B \neq A \cup B$; हा $A \cap B = A \cup B$ या गृहीतकाशी व्याघात आहे.

प्रकरण 2: $B \nsubseteq A$. मग काही $y$ साठी $y \in B$, पण $y \notin A$. आधीप्रमाणे $y \in A \cup B$, पण $y \notin A \cap B$. म्हणून $A \cap B \neq A \cup B$; हा पुन्हा $A \cap B = A \cup B$ शी व्याघात आहे.
\end{proof}











\section{सिद्धता वाचणे}\label{mth:prf:rea:sec}

पाठ्यपुस्तकांतील आणि लेखांतील सिद्धतांत आपण आत्तापर्यंत उदाहरणांत दिलेला सर्व तपशील क्वचितच असतो. प्रकरणे वेगळी पाहत आहोत, अप्रत्यक्ष सिद्धता देत आहोत किंवा व्याख्या वापरत आहोत, याकडे लेखक अनेकदा स्वतंत्र लक्ष वेधत नाहीत. म्हणून पाठ्यपुस्तकातील सिद्धता वाचताना ती समजण्यासाठी हा तपशील अनेकदा स्वतः भरून काढावा लागतो. सिद्धतेत कराव्या लागणाऱ्या विविध क्रियांचा सराव होण्यासाठीही हे उपयुक्त आहे. एक उदाहरण पाहू.

\begin{prop}[शोषणाचा नियम]
सर्व $A$, $B$ संचांसाठी,
\[
A \cap (A \cup B) = A
\]
\end{prop}

\begin{proof}
$z \in A \cap (A \cup B)$ असेल, तर $z \in A$. म्हणून $A \cap (A \cup B) \subseteq A$. आता $z \in A$ असे समजा. मग $z \in A \cup B$ देखील; म्हणून $z \in A \cap (A \cup B)$ देखील.
\end{proof}

शोषणाच्या नियमाची वरील सिद्धता फार संक्षिप्त आहे. वापरलेल्या व्याख्यांचा उल्लेख नाही; सिद्ध करण्याआधी “हे सिद्ध करायचे आहे” असे म्हटलेले नाही; इत्यादी. ती उलगडू. सिद्ध केलेले विधान हे कोणत्याही $A$ आणि $B$ संचांविषयीचे सामान्य विधान आहे. सिद्धतेत $A$ किंवा $B$ येतात तेव्हा ते स्वेच्छ संचांचे चल आहेत. सिद्धतेत मिळवलेली सामान्य विधाने ही संचांची एकरूपता सिद्ध करण्यासाठी आवश्यक असलेली विधाने आहेत: समीकरणाच्या डाव्या बाजूतील प्रत्येक घटक हा उजव्या बाजूतील घटक असतो आणि उलटही तसेच असते.

\begin{quote}
“$z \in A \cap (A \cup B)$ असेल, तर $z \in A$. म्हणून $A \cap (A \cup B) \subseteq A$.”
\end{quote}

हा एकरूपतेच्या सिद्धतेचा पहिला अर्धा भाग आहे. त्यात असे दाखवले आहे: स्वेच्छ~$z$ हा डाव्या बाजूतील घटक असेल, तर तो उजव्या बाजूतील घटक देखील असतो; म्हणजे $A \cap (A \cup B) \subseteq A$. $z \in A \cap (A \cup B)$ असे गृहीत धरा. $z$ हा दोन संचांच्या छेदाचा घटक असतो तेव्हाच आणि फक्त तेव्हा जेव्हा तो दोन्ही संचांचा घटक असतो. म्हणून $z \in A$ आणि $z \in A \cup B$ देखील, असा निष्कर्ष काढता येतो. विशेषतः $z \in A$ मिळते; हेच दाखवायचे होते. पहिल्या अर्ध्या भागासाठी इतकेच आवश्यक असल्यामुळे उरलेली सिद्धता दुसऱ्या अर्ध्या भागाची असली पाहिजे हे समजते; म्हणजे $A \subseteq A \cap (A \cup B)$ ची सिद्धता.

\begin{quote}
“आता $z \in A$ असे समजा. मग $z \in A \cup B$ देखील; म्हणून $z \in A \cap (A \cup B)$ देखील.”
\end{quote}

सुरुवातीला $z \in A$ असे गृहीत धरतो; कारण कोणत्याही~$z$ साठी $z \in A$ असेल, तर $z \in A \cap (A \cup B)$ असते हे दाखवत आहोत. $z \in A \cap (A \cup B)$ दाखवण्यासाठी “$\cap$” च्या व्याख्येनुसार (i) $z \in A$ आणि (ii) $z \in A \cup B$ दोन्ही दाखवावे लागतात. येथे (i) हे आपले गृहीतकच आहे; म्हणून त्यासाठी आणखी काही सिद्ध करायचे नाही. म्हणूनच सिद्धतेत त्याचा पुन्हा उल्लेख नाही. (ii) साठी हे आठवा: $z$ हा संचांच्या संयोगाचा घटक असतो तेव्हाच आणि फक्त तेव्हा जेव्हा तो त्यांपैकी किमान एका संचाचा घटक असतो. येथे $z \in A$ आहे आणि $A \cup B$ हा $A$ आणि $B$ यांचा संयोग आहे; म्हणून ही अट पूर्ण होते. त्यामुळे $z \in A \cup B$. (i) $z \in A$ आणि (ii) $z \in A \cup B$ दोन्ही दाखवले आहेत. म्हणून “$\cap$” च्या व्याख्येनुसार $z \in A \cap (A \cup B)$. सिद्धतेत या व्याख्यांचा उल्लेख नाही; त्या वाचकाला पूर्णपणे परिचित आहेत असे मानलेले आहे. तसे नसल्यास मागे जाऊन त्या काय आहेत हे आठवून घ्यावे लागेल. मग $z \in A$ वरून $z \in A \cup B$ का मिळते आणि $z \in A$ तसेच $z \in A \cup B$ यांवरून $z \in A \cap (A \cup B)$ का मिळते, हेही ओळखावे लागेल.

वरील सिद्धतेचे सर्व तपशील स्पष्ट केलेले आणखी एक रूप असे:
\begin{proof}{}
[संचांसाठीच्या $=$ च्या व्याख्येनुसार, $A \cap (A \cup B) = A$ दाखवण्यासाठी (a) $A \cap (A \cup B) \subseteq A$ आणि (b) $A \subseteq A \cap (A \cup B)$ दाखवावे लागते. (a): $\subseteq$ च्या व्याख्येनुसार $z \in A \cap (A \cup B)$ असेल, तर $z \in A$ असते हे दाखवावे लागते.] $z \in A \cap (A \cup B)$ असेल, तर $z \in A$ [$\cap$ च्या व्याख्येनुसार $z \in A \cap (A \cup B)$ हे $z \in A$ आणि $z \in A \cup B$ असते तेव्हाच आणि फक्त तेव्हाच सत्य असते]. म्हणून $A \cap (A \cup B) \subseteq A$. [(b): $\subseteq$ च्या व्याख्येनुसार $z \in A$ असेल, तर $z \in A \cap (A \cup B)$ असते हे दाखवावे लागते.] आता [(1)] $z \in A$ असे समजा. मग [(2)] $z \in A \cup B$ देखील [(1) वरून $z \in A$ किंवा $z \in B$; आणि $\cup$ च्या व्याख्येनुसार याचा अर्थ $z \in A \cup B$]. म्हणून $z \in A \cap (A \cup B)$ देखील [$\cap$ च्या व्याख्येनुसार $z \in A$, म्हणजे (1), आणि $z \in A \cup B$, म्हणजे (2), दोन्ही आवश्यक आहेत].
\end{proof}

\begin{prob}
$A \cup (A \cap B) = A$ ची पुढील सिद्धता विस्तृत करा. वापरलेल्या सर्व अनुमानाच्या नमुन्यांचा उल्लेख करा; प्रत्येक पायरी गृहीतकांवरून किंवा आधीच्या निष्कर्षांवरून का मिळते आणि कुठे कोणत्या व्याख्यांचा आधार घ्यावा लागतो, हे स्पष्ट करा.
\begin{proof}
$z \in A \cup (A \cap B)$ असेल, तर $z \in A$ किंवा $z \in A \cap B$. $z \in A \cap B$ असेल, तर $z \in A$. कोणताही $z \in A$ हा $\in A \cup (A \cap B)$ देखील असतो.
\end{proof}
\end{prob}












\section{मला जमत नाही!{}}\label{mth:prf:cnt:sec}

कधी ना कधी आपल्या सर्वांना हार मानावीशी वाटते. पण तुम्हाला हे \emph{जमू शकते}. शिक्षक, अध्यापन-सहायक आणि सहाध्यायी तुम्हाला मदत करू शकतात. त्यांच्याकडे मदत मागा!{} अडचणींचा मोठा पेच होऊ नये म्हणून आणि हार मानावीशी वाटल्यास काय करावे यासाठी काही सूचना पुढे दिल्या आहेत.

प्रश्न यशस्वीपणे सोडवता यावेत यासाठी पुढील गोष्टी करा:
\begin{enumerate}
\item \emph{शक्य तितक्या आधी सुरुवात करा.} सत्रभर आपण व्यस्त असतो आणि आपल्यापैकी अनेकांना काम पुढे ढकलण्याची सवय अडवते. गृहपाठाची सुरुवात लवकर करणे ही करता येण्याजोग्या सर्वात उपयुक्त गोष्टींपैकी एक आहे. त्यामुळे अडकलात तरी रडत बसण्यापलीकडे एखादा उपाय शोधण्यासाठी वेळ मिळतो.
\item \emph{सहाध्यायांशी बोला}. तुम्ही एकटे नाही. वर्गातील इतरांनाही अडचणी येत असतील; पण त्यांच्या अडचणी वेगळ्या असू शकतात. सहाध्यायांशी चर्चा केल्यामुळे प्रश्नाकडे पाहण्याची वेगळी दृष्टी मिळू शकते आणि कदाचित मार्ग सापडेल. अर्थात, त्यांचे उत्तर नुसते उतरवू नका: एखादी खूण मागा किंवा कुठे अडकलात हे समजावून पुढची पायरी विचारा. तुम्हाला समजले की त्यांनाही मदत करा. दुसऱ्याला पुढे जाण्यास मदत केल्याने आणि समजावून सांगितल्याने तुमची स्वतःची समजही सुधारते.
\item \emph{मदत मागा.} तुमच्यासाठी अनेक साधने उपलब्ध आहेत. शिक्षक आणि अध्यापन-सहायक तुमच्या मदतीसाठी आहेत आणि तुम्ही यशस्वी \emph{व्हावे अशी त्यांची इच्छा आहे}. प्रश्न सोडवण्यात मदत करणे आणि प्रक्रियेत नेमकी कुठे अडचण येते हे ओळखणे त्यांना शक्य असले पाहिजे.
\item \emph{थोडी विश्रांती घ्या.} अडकलात, तर त्या प्रश्नाकडे खूप वेळ पाहत बसल्यामुळे तसे झाले असण्याची \emph{शक्यता} आहे. थोडी विश्रांती घ्या, चहा प्या किंवा काही वेळ दुसरा प्रश्न सोडवा; मग ताज्या मनाने पहिल्या प्रश्नाकडे परत या. एखादी रात्र झोप घेऊन दुसऱ्या दिवशी पुन्हा पाहा.
\end{enumerate}

या उपायांसाठी सिद्धतेवर काम बराच वेळ आधी सुरू केलेले असणे आवश्यक आहे, हे लक्षात आले का? गृहपाठ देण्याच्या आदल्या दिवशी पहाटे दोन वाजता सिद्धता सुरू केली असेल, तर हे उपाय फार उपयोगी ठरणार नाहीत.

हे निराशाजनक वाटेल; पण सिद्धता करणे हे शेवटी समाधान देणारे आव्हान आहे. काही लोक हेच व्यवसाय म्हणून करतात; म्हणजे त्यात आनंद देणारे काहीतरी असले पाहिजे. जवळजवळ इतर सर्व गोष्टींप्रमाणे प्रश्न सोडवणे आणि सिद्धता करणे यांनाही सराव लागतो. काही सहाध्यायांना हे सोपे वाटत असेल; बहुधा त्यांनी आधीच बराच सराव केला असेल. फार लवकर हार मानू नका.

एखाद्या प्रश्नासाठी वेळ किंवा धीर संपला, तरी ठीक आहे. त्याचा अर्थ तुम्ही मूर्ख आहात किंवा तुम्हाला कधीच जमणार नाही असा नाही. शिक्षक किंवा दुसरा विद्यार्थी यांच्याकडून तो कसा सोडवतात ते समजून घ्या. कुठे चूक झाली किंवा कुठे अडकलात ते ओळखा, म्हणजे पुढे तशी अडचण आली तर ती टाळता येईल. मग उत्तर न पाहता पुन्हा सोडवण्याचा प्रयत्न करा. पुढच्या वेळी आधी सुरुवात करा आणि मदतही लवकर मागा.












\section{इतर साधने}\label{mth:prf:res:sec}

गणितात सिद्धता कशा कराव्यात याविषयी अनेक उपयुक्त पुस्तके आहेत. मूळ स्रोताने विशेषतः \emph{How to Read and do Proofs: An Introduction to Mathematical Thought Processes} (सोलो 2013) आणि \emph{How to Prove It: A Structured Approach} (व्हेलमन 2019) पाहण्याची सूचना केली आहे. \href{http://www.people.vcu.edu/~rhammack/BookOfProof/BookOfProof.pdf}{\emph{Book of Proof}} (हॅमॅक 2013) आणि \href{https://scholarworks.gvsu.edu/books/7/}{\emph{Mathematical Reasoning}} (सँडस्ट्रम 2019) ही सिद्धतांवरील पुस्तके मूळ स्रोतामध्ये इंटरनेटवर मुक्तपणे उपलब्ध म्हणून नमूद केलेली आहेत. तत्त्वज्ञानाच्या विद्यार्थ्यांसाठी \emph{More Precisely: The Math you need to do Philosophy} (स्टाइनहार्ट 2018) हे गणितीय विचाराचे उपयुक्त प्रास्ताविक पुस्तक ठरू शकते.

इंटरनेटवर सिद्धतांविषयी काही छोटी मार्गदर्शकेही आहेत. उदाहरणार्थ, \href{https://math.berkeley.edu/~hutching/teach/proofs.pdf}{“Introduction to Mathematical Arguments”} (हचिंग्ज 2003) आणि \href{https://eugeniacheng.com/wp-content/uploads/2017/02/cheng-proofguide.pdf}{“How to write proofs”} (चेंग 2004).

\subsection{प्रेरणादायी चित्रफिती}

गृहपाठ करण्याची प्रेरणाच उरली नाही असे वाटते का? मन खिन्न आहे का? या चित्रफिती मदत करू शकतील!

\begin{itemize}
\item \url{https://www.youtube.com/watch?v=ZXsQAXx_ao0}
\item \url{https://www.youtube.com/watch?v=BQ4yd2W50No}
\item \url{https://www.youtube.com/watch?v=StTqXEQ2l-Y}
\end{itemize}



\chapter{विगमन}\label{mth:ind:chap}









\section{प्रस्तावना}\label{mth:ind:int:sec}

विगमन ही सिद्धतेची महत्त्वाची पद्धती आहे. तर्कशास्त्र, सैद्धान्तिक संगणकशास्त्र आणि गणित यांच्या जवळजवळ सर्व शाखांत ती वेगवेगळ्या रूपांत वापरली जाते. तर्कशास्त्रातील अनेक निष्कर्ष सिद्ध करण्यासाठी ती आवश्यक आहे.

विगमनाचा अनेकदा निगमनाशी विरोध दाखवला जातो आणि विशिष्ट प्रकरणांकडून सामान्य विधानाकडे जाणारे अनुमान, असे त्याचे वर्णन केले जाते. उदाहरणार्थ, अनेक हिरवी पाचूची रत्ने आपण पाहिली आणि पाचू म्हणता येईल असे कोणतेही रत्न हिरवे नसलेले आढळले नाही, तर सर्व पाचू हिरवे असतात असा निष्कर्ष काढू शकतो. हे विगमनाधारित अनुमान आहे: ते अनेक विशिष्ट प्रकरणांकडून---हा पाचू हिरवा आहे, तो पाचू हिरवा आहे, इत्यादी---एका सामान्य विधानाकडे---सर्व पाचू हिरवे असतात---जाते. \emph{गणितीय} विगमनातही सामान्य विधान निष्कर्ष म्हणून मिळते; पण ते निरीक्षणांवरून केलेल्या या “साध्या विगमना”पेक्षा फार वेगळे असते.

स्थूलपणे सांगायचे तर, गणितातील विगमनाधारित सिद्धतेने एखाद्या प्रकारच्या सर्व गणितीय वस्तूंमध्ये विशिष्ट गुणधर्म आहे असा निष्कर्ष मिळतो. सर्वात साध्या प्रकरणात या सिद्धतेतील गणितीय वस्तू नैसर्गिक संख्या असतात. अशा वेळी सर्व नैसर्गिक संख्यांत एखादा गुणधर्म आहे हे सिद्ध करण्यासाठी विगमन वापरतो; त्यात असे दाखवतो:
\begin{enumerate}
\item $0$ मध्ये तो गुणधर्म आहे; आणि
\item एखाद्या~$k$ संख्येत तो गुणधर्म असेल, तर~$k+1$ मध्येही तो असतो.
\end{enumerate}
ज्या गणितीय वस्तूंना संख्या देता येतात त्यांच्याविषयीची सामान्य विधाने सिद्ध करण्यासाठीही नैसर्गिक संख्यांवरील विगमन अनेकदा वापरता येते. उदाहरणार्थ, प्रत्येक सांत संचात सांत संख्येने, समजा~$n$, घटक असतात. विगमन वापरून प्रत्येक~$n$ संख्येसाठी “~$n$ आकारमानाचे सर्व सांत संच \dots{} असतात” हा गुणधर्म आहे असे दाखवता आले, तर सर्व सांत संचांविषयी काहीतरी दाखवले असेल.

\emph{विगमनाधारित व्याख्येने दिलेल्या} गणितीय वस्तूंपर्यंतही विगमनाचा विस्तार करता येतो. उदाहरणार्थ, प्रथम-क्रम तर्कशास्त्रासारख्या औपचारिक भाषेतील अभिव्यक्ती विगमनाधारित व्याख्येने दिलेल्या असतात. \emph{संरचनात्मक विगमन} ही अशा सर्व अभिव्यक्तींविषयी निष्कर्ष सिद्ध करण्याची पद्धती आहे. विशेषतः तर्कशास्त्रात संरचनात्मक विगमन फार उपयुक्त आहे आणि मोठ्या प्रमाणात वापरले जाते.












\section{~$\Nat$ वरील विगमन}\label{mth:ind:inN:sec}

सर्वात साध्या रूपात विगमन ही सर्व नैसर्गिक संख्यांसाठी निष्कर्ष सिद्ध करण्याची पद्धती आहे. $0$ पासून सुरुवात करून वारंवार~$1$ वाढवत गेल्यास प्रत्येक नैसर्गिक संख्येपर्यंत अखेर पोहोचता येते, या तथ्याचा ती वापर करते. म्हणून एखादी गोष्ट प्रत्येक संख्येसाठी खरी आहे हे सिद्ध करण्यासाठी (1) ती $0$ साठी खरी आहे हे प्रस्थापित करता येते आणि (2) ती एखाद्या~$n$ संख्येसाठी खरी असेल तेव्हा पुढच्या~$n+1$ संख्येसाठीही खरी असते हे दाखवता येते. “~$n$ संख्येत $P$ हा गुणधर्म आहे” याचा संक्षेप $P(n)$ असा करू; तसेच “~$k$ संख्येत $P$ हा गुणधर्म आहे” याचा $P(k)$, इत्यादी. मग सर्व $n \in \Nat$ साठी $P(n)$ सिद्ध करणाऱ्या विगमनाधारित सिद्धतेत पुढील गोष्टी असतात:
\begin{enumerate}
\item $P(0)$ ची सिद्धता; आणि
\item कोणत्याही~$k$ साठी, $P(k)$ असेल तर $P(k+1)$ असते याची सिद्धता.
\end{enumerate}
हे पूर्ण स्पष्ट करण्यासाठी (1) आणि~(2) दोन्ही आहेत असे समजा. मग (1) वरून $P(0)$ सत्य आहे असे समजते. ~(2) देखील असल्यामुळे विशेषतः $P(0)$ असेल तर $P(0+1)$, म्हणजे $P(1)$, असते हे समजते. “कोणत्याही~$k$ साठी, $P(k)$ असेल तर $P(k+1)$” या सामान्य विधानात~$k$ च्या जागी~$0$ ठेवल्यावर हे मिळते. म्हणून विध्यात्मक अनुमानाने~$P(1)$ मिळते. पुन्हा (2) मध्ये, आता~$k$ च्या जागी $1$ घेतल्यावर, $P(1)$ असेल तर~$P(2)$ मिळते. ~$P(1)$ नुकतेच प्रस्थापित केलेले असल्यामुळे विध्यात्मक अनुमानाने~$P(2)$ मिळते. असेच पुढे चालू राहते. कोणत्याही~$n$ संख्येसाठी ही क्रिया $n$~वेळा केल्यावर अखेर~$P(n)$ मिळते. म्हणून (1) आणि~(2) मिळून कोणत्याही $n \in \Nat$ साठी~$P(n)$ प्रस्थापित करतात.

एक उदाहरण पाहू. $n$ फासे टाकल्यावर किती वेगवेगळ्या बेरजा मिळू शकतात हे शोधायचे आहे असे समजा. थोडे विचित्र वाटले, तरी $0$ फाशांपासून सुरुवात करू. फासेच नसतील, तर “टाकून” मिळणारी एकच बेरीज आहे: कोणतेही ठिपके नाहीत; म्हणजे बेरीज~$0$. म्हणून बेरजेच्या वेगवेगळ्या शक्यतांची संख्या~$1$ आहे. एकच फासा असेल, म्हणजे $n=1$, तर $1$ ते~$6$ अशी सहा मूल्ये शक्य आहेत. दोन फाशांनी $2$ ते $12$ पर्यंत कोणतीही बेरीज मिळू शकते; म्हणजे $11$ शक्यता. तीन फाशांनी $3$ ते $18$ पर्यंत कोणतीही संख्या मिळू शकते; म्हणजे $16$ वेगवेगळ्या शक्यता. $1$, $6$, $11$, $16$: एक नमुना दिसतो. कदाचित उत्तर $5n+1$ असेल? अर्थात, $5n+1$ ही कमाल शक्य संख्या आहे. कारण $n$ फाशांनी मिळणाऱ्या सर्वात लहान $n$ मूल्यापासून---सर्व फाशांवर $1$---सर्वात मोठ्या $6n$ मूल्यापर्यंत---सर्व फाशांवर $6$---फक्त $5n+1$ संख्या आहेत.

\begin{thm}
$n$ फाशांनी $n$ ते $6n$ यांमधील सर्व $5n+1$ शक्य मूल्ये मिळवता येतात.
\end{thm}

\begin{proof}
$P(n)$ हे विधान असू द्या: “$n$~फासे वापरून $n$ ते~$6n$ यांमधील कोणतीही संख्या मिळवता येते.” शून्य फाशांचे प्रकरण वर पाहिले आहे: रिक्त बेरीज शून्य असून एकच शक्यता असते. धन संख्यांसाठी विगमन वापरण्यास पुढील गोष्टी सिद्ध करू:
\begin{enumerate}
\item \emph{विगमनाचा आधार} $P(1)$: एकच फासा वापरून $1$ ते $6$ यांमधील कोणतीही संख्या मिळवता येते.
\item \emph{विगमनाची पायरी}: प्रत्येक धन नैसर्गिक संख्या $k$ साठी, $P(k)$ असेल तर~$P(k+1)$.
\end{enumerate}

(1) हे $6$ पृष्ठे असलेला फासा पाहून सिद्ध होते. त्याला सहाही पृष्ठे आहेत आणि $1$ ते~$6$ मधील प्रत्येक संख्या एका पृष्ठावर दिसते. म्हणून एक फासा वापरून $1$ ते~$6$ मधील कोणतीही संख्या मिळवता येते.

(2) सिद्ध करण्यासाठी सशर्त विधानाचे पूर्वांग, म्हणजे $P(k)$, गृहीत धरतो. या गृहीतकाला \emph{विगमन गृहीतक} म्हणतात. त्याचा वापर करून~$P(k+1)$ सिद्ध करतो. अवघड भाग असा: $k+1$ फाशांनी मिळणाऱ्या शक्य बेरजांचा विचार $k$~फाशांनी मिळणाऱ्या बेरजा आणि अतिरिक्त $k+1$ व्या फाशाचे मूल्य यांच्या आधारे करण्याचा मार्ग शोधायचा. विगमन गृहीतक वापरायचे असेल, तर असेच करावे लागते.

विगमन गृहीतकानुसार $k$~फाशांनी $k$ ते~$6k$ मधील कोणतीही संख्या मिळवता येते. ~$(k+1)$ व्या फाशावर~$1$ आला, तर एकूण बेरजेत $1$ वाढतो. म्हणून $k$~फासे टाकून नंतर $(k+1)$ व्या फाशावर~$1$ आणल्यास $k+1$ ते $6k+1$ मधील कोणतेही मूल्य मिळवता येते. काय उरले? $6k+2$ ते $6k+6$ ही मूल्ये. $k$ फाशांवर $6$ आणून आणि $(k+1)$ व्या फाशावर $2$ ते $6$ मधील संख्या आणून ही मूल्ये मिळवता येतात. या दोन्ही गोष्टी मिळून असे दाखवतात की $k+1$ फाशांनी $k+1$ ते $6(k+1)$ मधील कोणतीही संख्या मिळवता येते. म्हणजे~$P(k)$ हे विगमन गृहीतक वापरून~$P(k+1)$ सिद्ध केले आहे.
\end{proof}

अनेकदा वस्तूंच्या अनुक्रमाविषयी---संख्या, संच इत्यादी---काही सिद्ध करायचे असते आणि तो अनुक्रमच “विगमनाधारित” व्याख्येने दिलेला असतो. म्हणजे $(n+1)$ वी वस्तू $n$ व्या वस्तूच्या आधारे परिभाषित केलेली असते. अशा वेळी विगमन वापरतो. उदाहरणार्थ,~$n$ पर्यंतच्या नैसर्गिक संख्यांची बेरीज~$s_n$ अशी परिभाषित करता येते:
\begin{align*}
  s_0 & = 0\\
  s_{n+1} & = s_n + (n+1)
\end{align*}
या व्याख्येवरून असे मिळते:
\begin{align*}
  s_0 & = 0,\\
  s_1 & = s_0 + 1 && = 1,\\
  s_2 & = s_1 + 2 && = 1 + 2 = 3\\
  s_3 & = s_2 + 3 && = 1 + 2 + 3 = 6, \text{ इत्यादी.}
\end{align*}
आता विगमनाने $s_n = n(n+1)/2$ सिद्ध करता येते.

\begin{prop}
  $s_n = n(n+1)/2$.
\end{prop}

\begin{proof}
(1) $s_0 = 0\cdot(0 + 1)/2$ आणि (2) $s_k = k(k+1)/2$ असेल तर $s_{k+1} = (k+1)(k+2)/2$, हे सिद्ध करायचे आहे. (1) स्पष्ट आहे. (2) सिद्ध करण्यासाठी $s_k = k(k+1)/2$ हे विगमन गृहीतक घेऊ. त्याचा वापर करून $s_{k+1} = (k+1)(k+2)/2$ दाखवायचे आहे.

$s_{k+1}$ काय आहे? व्याख्येनुसार $s_{k+1} = s_k + (k+1)$. विगमन गृहीतकानुसार $s_k = k(k+1)/2$. आधीच्या समीकरणात ही किंमत ठेवता येते; मग अपूर्णांकांचे थोडे अंकगणित पुरते:
\begin{align*}
    s_{k+1} & = \frac{k(k+1)}{2} + (k+1) = {}\\
    & = \frac{k(k+1)}{2} + \frac{2(k+1)}{2} = {}\\
    & = \frac{k(k+1) + 2(k+1)}{2} = {}\\
    & = \frac{(k+2)(k+1)}{2}.
\end{align*}
\end{proof}

महत्त्वाचा धडा असा: विगमनाधारित व्याख्येने दिलेल्या $a_n$ अनुक्रमाविषयी काही सिद्ध करत असाल, तर विगमन हा स्वाभाविक मार्ग आहे. तसे नसले तरी---फाशांच्या शक्य बेरजांच्या उदाहरणाप्रमाणे---~$k+1$ चे प्रकरण~$k$ च्या प्रकरणाशी काही रीतीने जोडता आले, तर विगमन वापरता येते.












\section{प्रबल विगमन}\label{mth:ind:str:sec}

आधी चर्चिलेल्या विगमन तत्त्वात $P(0)$ सिद्ध करतो आणि $P(k)$ असेल तर $P(k+1)$ असते हेही सिद्ध करतो. दुसऱ्या भागात $P(k)$ सत्य आहे असे गृहीत धरून त्या गृहीतकावरून $P(k+1)$ सिद्ध करतो. धन निर्देशांकासाठी याऐवजी $P(k-1)$ गृहीत धरून $P(k)$ सिद्ध करता येते. महत्त्वाचे असे की कोणत्याही संख्येपासून तिच्या उत्तरवर्ती संख्येपर्यंतचे हे अनुमान करता आले पाहिजे: एखाद्या संख्येच्या पूर्ववर्ती संख्येसाठी विधान सत्य आहे असे गृहीत धरून तिच्यासाठी ते सिद्ध करता आले पाहिजे.

विगमन तत्त्वाचा आणखी एक प्रकार आहे. त्यात केवळ $k$ च्या पूर्ववर्ती $k-1$ साठी विधान खरे आहे असे न मानता,~$k$ पेक्षा लहान असलेल्या सर्व नैसर्गिक संख्यांसाठी ते खरे आहे असे मानतो आणि त्यावरून~$k$ साठी ते प्रस्थापित करतो. यानेही सर्व~$n \in \Nat$ साठी $P(n)$ मिळते. कारण $P(0)$ प्रस्थापित केल्यावर~$1$ पेक्षा लहान सर्व नैसर्गिक संख्यांसाठी $P$ खरे आहे असे प्रस्थापित झाले आहे. प्रत्येक $l<k$ साठी $P(l)$ असेल तर $P(k)$ असते, हे माहीत असल्यास विशेषतः $k=1$ साठीही ते माहीत असते. म्हणून $P(1)$ असा निष्कर्ष काढता येतो. यामुळे $P(0)$ आणि $P(1)$ सिद्ध झाले; म्हणजे सर्व $l<2$ साठी $P(l)$. तसेच प्रत्येक $l<2$ साठी $P(l)$ असेल तर $P(2)$ असते, हे सशर्त विधानही उपलब्ध असल्यामुळे~$P(2)$ निष्कर्ष म्हणून मिळते; आणि असेच पुढे.

खरे तर “प्रत्येक $k$ साठी, प्रत्येक $l<k$ साठी $P(l)$ असेल तर $P(k)$ असते” हे सामान्य सशर्त विधान प्रस्थापित करता आले, तर $P(0)$ वेगळे प्रस्थापित करण्याची गरज नसते; ते यावरूनच मिळते. “प्रत्येक $l<k$ साठी $P(l)$” सारखे सामान्य विधान कोणतेही $l<k$ नसतील तेव्हा सत्य असते, हे आठवा. हे रिक्त संख्यापनाचे प्रकरण आहे: कोणतेही $A$ नसतील, तर “सर्व $A$ हे $B$ आहेत” हे सत्य असते; कोणताही $x$ हा~$\varphi(x)$ पूर्ण करत नसेल, तर $\lforall[x][(\varphi(x) \lif \psi(x))]$ सत्य असते. या प्रकरणात औपचारिक रूप “$\lforall[l][(l < k \lif P(l))]$” असे असेल; आणि कोणतेही $l < k$ नसतील, तर ते सत्य असते. $k=0$ असेल तेव्हा नेमके हेच घडते: कोणतीही नैसर्गिक संख्या $l<0$ नाही. म्हणून $P$ काहीही असले, तरी “प्रत्येक $l<0$ साठी $P(l)$” सत्य असते. त्यामुळे “प्रत्येक $l<k$ साठी $P(l)$ असेल तर $P(k)$ असते” याची सिद्धता आपोआप~$P(0)$ प्रस्थापित करते.

~$k$ साठी विधान प्रस्थापित करताना केवळ $k-1$ साठीच्या विधानावर अवलंबून राहता येत नसेल, पण एक किंवा अधिक $l<k$ साठी ते सत्य असल्याचे गृहीतक आवश्यक असेल, तर हा प्रकार उपयुक्त ठरतो.












\section{विगमनाधारित व्याख्या}\label{mth:ind:idf:sec}

तर्कशास्त्रात अनेकदा वस्तूंचे प्रकार \emph{विगमनाधारित रीतीने} परिभाषित करतो. म्हणजे त्या प्रकारची वस्तू कशाला म्हणायचे याचे नियम देतो; त्यांत त्या प्रकारच्या जुन्या वस्तूंपासून नव्या वस्तू कशा मिळवायच्या हे सांगितलेले असते. उदाहरणार्थ, भाषेतील पदे आणि सूत्रे यांसारखे चिन्हांचे विशिष्ट अनुक्रम अनेकदा विगमनाने परिभाषित करतो. एक साधे उदाहरण पाहू: $\mathrm{a}$, $\mathrm{b}$, $\mathrm{c}$, $\mathrm{d}$ ही अक्षरे,~$\circ$ हे चिन्ह आणि $[$, $]$ हे कंस यांपासून बनलेल्या चिन्हमाला घ्या. उदाहरणार्थ, “$[[\mathrm{c} \circ \mathrm{d}][$”, “$[\mathrm{a}[]\circ]$”, “$\mathrm{a}$” किंवा “$[[\mathrm{a} \circ \mathrm{b}]\circ \mathrm{d}]$”. पहिल्या दोन चिन्हमालांत काहीतरी “चुकीचे” आहे असे वाटत असेल: पहिल्या चिन्हमालेत कंस मुळीच जुळत नाहीत; आणि “$\circ$” ने स्वतः अर्थपूर्ण असलेल्या अभिव्यक्ती “जोडल्या” पाहिजेत असे वाटेल. तिसरी आणि चौथी चिन्हमाला अधिक नीट दिसतात: प्रत्येक “$[$” साठी बंद करणारा “$]$” आहे---कंस असतील तर---आणि प्रत्येक $\circ$ च्या दोन्ही बाजूंना स्वतः नीट असलेल्या अभिव्यक्ती आहेत, ज्यांभोवती कंसांची जोडी आहे.

“नीटस पद” नेमके कशाला म्हणायचे हे निश्चित करायचे आहे. प्रथम, प्रत्येक स्वतंत्र अक्षर नीटस आहे. एकटे अक्षर नसलेले प्रत्येक नीटस पद “$[t \circ s]$” या रूपाचे असले पाहिजे; येथे $s$ आणि $t$ ही स्वतः नीटस असतात. उलट, $t$ आणि $s$ नीटस असतील, तर त्यांच्यामध्ये $\circ$ ठेवून आणि भोवती कंसांची जोडी घालून नवे नीटस पद बनवता येते. या क्रिया वापरून नीटस पदांचा संच \emph{परिभाषित} करता येईल. ही \emph{विगमनाधारित व्याख्या} आहे.

\begin{defn}[नीटस पदे]
\emph{नीटस पदांचा} संच पुढील विगमनाधारित व्याख्येने दिलेला आहे:
\begin{enumerate}
\item $\mathrm{a}$, $\mathrm{b}$, $\mathrm{c}$, $\mathrm{d}$ यांपैकी कोणतेही अक्षर नीटस पद आहे.
\item $s_1$ आणि $s_2$ नीटस पदे असतील, तर $[s_1 \circ s_2]$ देखील नीटस पद आहे.
\item इतर काहीही नीटस पद नाही.
\end{enumerate}
\end{defn}

या व्याख्येनुसार एखादी गोष्ट नीटस पद असते तेव्हाच आणि फक्त तेव्हा जेव्हा (1) आणि~(2) या दोन अटींनुसार सांत संख्येने पायऱ्यांत ती रचता येते. पहिल्या पायरीत केवळ स्वतंत्र अक्षरे असलेली सर्व नीटस पदे रचतो:
\[
\mathrm{a}, \mathrm{b}, \mathrm{c}, \mathrm{d}
\]
दुसऱ्या पायरीत रचलेल्या पदांना (2) लागू करतो. दोन अक्षरांच्या सर्व जोड्यांसाठी पुढील पदे मिळतील:
\[
[\mathrm{a} \circ \mathrm{a}], [\mathrm{a} \circ \mathrm{b}],
[\mathrm{b} \circ \mathrm{a}], \dots, [\mathrm{d} \circ \mathrm{d}]
\]
तिसऱ्या पायरीत आत्तापर्यंत रचलेल्या कोणत्याही दोन नीटस पदांना पुन्हा (2) लागू करतो. उदाहरणार्थ, $[\mathrm{a} \circ [\mathrm{a} \circ \mathrm{a}]]$ असे नवे नीटस पद मिळते: येथे $t$ हे पायरी~1 मधील $\mathrm{a}$ आहे आणि $s$ हे पायरी~2 मधील $[\mathrm{a} \circ \mathrm{a}]$ आहे. तसेच पायरी~2 मधील $[\mathrm{b} \circ \mathrm{c}]$ आणि $[\mathrm{d} \circ \mathrm{b}]$ या दोन पदांपासून $[[\mathrm{b} \circ \mathrm{c}] \circ [\mathrm{d} \circ \mathrm{b}]]$ रचता येते. असेच पुढे. या रीतीने न रचलेली कोणतीही गोष्ट नीटस पदांच्या संचात येऊ नये, हे खंड (3) निश्चित करतो.

“नीट वाटणारा” प्रत्येक चिन्हानुक्रम या व्याख्येनुसार नीटस आहे, हे अद्याप सिद्ध केलेले नाही, हे लक्षात घ्या. मात्र आपण रचू शकणारी प्रत्येक गोष्ट खरोखर “नीट वाटते”, हे स्पष्ट असले पाहिजे: कंस जुळतात आणि $\circ$ ने स्वतः नीटस असलेले भाग जोडलेले असतात.

विगमनाधारित व्याख्यांचे महत्त्वाचे वैशिष्ट्य असे की सर्व नीटस पदांविषयी काही सिद्ध करायचे असेल, तर कोणती प्रकरणे पाहावी लागतील हे व्याख्या सांगते. उदाहरणार्थ, $t$ हे नीटस पद आहे असे दिले असेल, तर $t$ कसे असू शकते हे व्याख्या सांगते: $t$ हे अक्षर असेल किंवा ते $[s_1 \circ s_2]$ या रूपाचे असेल, जिथे $s_1$ आणि~$s_2$ ही नीटस पदे आहेत. खंड (3) मुळे या दोनच शक्यता आहेत.

विगमनाधारित व्याख्येने दिलेल्या संचातील सर्व वस्तूंविषयी विधाने सिद्ध करताना प्रबल विगमन विशेष महत्त्वाचे ठरते. उदाहरणार्थ,~$n$ लांबीच्या प्रत्येक नीटस पदात $[$ ची संख्या~$< n/2$ आहे हे सिद्ध करायचे आहे असे समजा. याचा सर्व~$n$ विषयीचे विधान म्हणून विचार करता येतो: प्रत्येक $n$ साठी,~$n$ लांबीच्या कोणत्याही नीटस पदात $[$ ची संख्या $< n/2$ आहे.

\begin{prop}
कोणत्याही $n$ साठी,~$n$ लांबीच्या नीटस पदात $[$ ची संख्या $< n/2$ आहे.
\end{prop}

\begin{proof}
प्रबल विगमनाने हा निष्कर्ष सिद्ध करण्यासाठी पुढील सशर्त विधान सत्य आहे हे दाखवावे लागेल:
\begin{quote}
प्रत्येक $l < k$ साठी,~$l$ लांबीच्या कोणत्याही नीटस पदात $[$ ची संख्या $< l/2$ असेल, तर~$k$ लांबीच्या कोणत्याही नीटस पदात $[$ ची संख्या $< k/2$ असते.
\end{quote}
हे सशर्त विधान दाखवण्यासाठी त्याचे पूर्वांग सत्य आहे असे गृहीत धरा. म्हणजे कोणत्याही $l<k$ साठी,~$l$ लांबीच्या नीटस पदात $[$ ची संख्या $< l/2$ आहे असे गृहीत धरा. याला विगमन गृहीतक म्हणतो. ~$k$ लांबीच्या नीटस पदांसाठीही तेच सत्य आहे हे दाखवायचे आहे.

म्हणून $t$ हे~$k$ लांबीचे नीटस पद आहे असे समजा. नीटस पदांची व्याख्या विगमनाधारित असल्यामुळे दोन प्रकरणे आहेत: (1)~$t$ हे स्वतंत्र अक्षर आहे; किंवा (2)~$t$ हे $[s_1 \circ s_2]$ या रूपाचे आहे, जिथे $s_1$ आणि~$s_2$ ही नीटस पदे आहेत.
\begin{enumerate}
\item $t$ हे अक्षर आहे. मग $k = 1$ आणि $t$ मध्ये $[$ ची संख्या~$0$ आहे. $0 < 1/2$ असल्यामुळे विधान खरे आहे.
\item $t$ हे $[s_1 \circ s_2]$ या रूपाचे आहे, जिथे $s_1$ आणि~$s_2$ ही नीटस पदे आहेत. ~$s_1$ ची लांबी $l_1$ आणि~$s_2$ ची लांबी $l_2$ असू द्या. मग $t$ ची लांबी~$k$ ही $l_1+l_2+3$ आहे: $s_1$ आणि $s_2$ यांच्या लांबी अधिक $[$, $\circ$, $]$ ही तीन चिन्हे. $l_1+l_2+3$ हे नेहमी $l_1$ पेक्षा मोठे असल्यामुळे $l_1 < k$. त्याचप्रमाणे $l_2 < k$. म्हणजे $s_1$ आणि~$s_2$ या पदांना विगमन गृहीतक लागू होते: $s_1$ मध्ये $[$ ची संख्या~$m_1$ ही $< l_1/2$ आहे; आणि~$s_2$ मध्ये $[$ ची संख्या~$m_2$ ही $< l_2/2$ आहे.

$t$ मध्ये $[$ ची संख्या म्हणजे~$s_1$ मधील $[$ ची संख्या, अधिक~$s_2$ मधील $[$ ची संख्या, अधिक~$1$; म्हणजे $m_1 + m_2 + 1$. $m_1 < l_1/2$ आणि $m_2 < l_2/2$ असल्यामुळे:
\[
  m_1 + m_2 + 1 < \frac{l_1}{2} + \frac{l_2}{2} + 1 = \frac{l_1+l_2+2}{2} < \frac{l_1+l_2+3}{2} = k/2.
\]
\end{enumerate}
प्रत्येक प्रकरणात विगमन गृहीतकाच्या आधारे $t$ मध्ये $[$ ची संख्या $< k/2$ आहे हे दाखवले. प्रबल विगमनाने विधान सिद्ध होते.
\end{proof}

\begin{prob}
अतिनीटस पदांचा संच पुढील रीतीने परिभाषित करा:
\begin{enumerate}
\item $\mathrm{a}$, $\mathrm{b}$, $\mathrm{c}$, $\mathrm{d}$ यांपैकी कोणतेही अक्षर अतिनीटस पद आहे.
\item $s$ हे अतिनीटस पद असेल, तर $[s]$ देखील अतिनीटस पद आहे.
\item $s_1$ आणि $s_2$ अतिनीटस पदे असतील, तर $[s_1 \circ s_2]$ देखील अतिनीटस पद आहे.
\item इतर काहीही अतिनीटस पद नाही.
\end{enumerate}
~$n$ ही लांबी असलेल्या अतिनीटस पद~$t$ मध्ये $[$ ची संख्या $\le n/2 +1$ आहे हे दाखवा.
\end{prob}












\section{संरचनात्मक विगमन}\label{mth:ind:sti:sec}

आत्तापर्यंत सर्व नैसर्गिक संख्यांविषयी निष्कर्ष प्रस्थापित करण्यासाठी विगमन वापरले. पण त्याच्याशी संबंधित तत्त्व थेट वापरून विगमनाधारित व्याख्येने दिलेल्या संचातील सर्व घटक विषयी निष्कर्ष सिद्ध करता येतात. याला अनेकदा \emph{संरचनात्मक} विगमन म्हणतात; कारण ते त्या व्याख्येने दिलेल्या वस्तूंच्या संरचनेवर अवलंबून असते.

सहसा विगमनाधारित व्याख्येत (a) संचातील “आरंभीच्या” घटक ची यादी आणि (b) जुन्यांपासून नवे घटक निर्माण करणाऱ्या क्रियांची यादी दिलेली असते. उदाहरणार्थ, नीटस पदांच्या प्रकरणात आरंभीच्या वस्तू म्हणजे अक्षरे. एकच क्रिया आहे:
\begin{align*}
  o(s_1, s_2) = & [s_1 \circ s_2]
\end{align*}
नैसर्गिक संख्यांचा संच~$\Nat$ देखील विगमनाधारित व्याख्येने दिलेला आहे असा विचार करता येतो: आरंभीची वस्तू~$0$ आणि क्रिया म्हणजे~$x + 1$ हे उत्तरवर्ती फलन.

विगमनाधारित व्याख्येने दिलेल्या संचातील सर्व घटकांविषयी काही सिद्ध करायचे असेल---म्हणजे संचातील प्रत्येक घटक मध्ये~$P$ हा गुणधर्म आहे हे सिद्ध करायचे असेल---तर पुढील गोष्टी कराव्या लागतात:
\begin{enumerate}
\item आरंभीच्या वस्तूंमध्ये~$P$ आहे हे सिद्ध करा.
\item प्रत्येक~$o$ क्रियेसाठी, तिच्या सर्व आदानांत~$P$ असेल, तर फलितातही तो असतो हे सिद्ध करा.
\end{enumerate}
उदाहरणार्थ, सर्व नीटस पदांविषयी काही सिद्ध करण्यासाठी ते सर्व अक्षरांविषयी सत्य आहे हे सिद्ध करू; तसेच $s_1$ आणि $s_2$ यांविषयी ते स्वतंत्रपणे सत्य असेल, तर $[s_1 \circ s_2]$ विषयीही सत्य असते हे सिद्ध करू.

\begin{prop}
कोणत्याही नीटस पद~$t$ मध्ये $[$ ची संख्या $]$ च्या संख्येइतकी असते.
\end{prop}

\begin{proof}
संरचनात्मक विगमन वापरू. नीटस पदांची व्याख्या विगमनाधारित आहे: अक्षरे या आरंभीच्या वस्तू आणि जुन्या पदांपासून नवी नीटस पदे रचण्यासाठी $o$ ही क्रिया.
\begin{enumerate}
\item प्रत्येक अक्षरासाठी विधान सत्य आहे; कारण स्वतंत्र अक्षरात $[$ ची संख्या~$0$ आणि $]$ ची संख्याही~$0$ आहे.
\item $s_1$ मध्ये $[$ ची संख्या $]$ च्या संख्येइतकी आहे आणि $s_2$ साठीही तेच सत्य आहे असे समजा. $o(s_1, s_2)$ मध्ये---म्हणजे $[s_1 \circ s_2]$ मध्ये---$[$ ची संख्या ही $s_1$ आणि $s_2$ यांमधील $[$ च्या संख्यांची बेरीज अधिक एक आहे. $o(s_1, s_2)$ मध्ये $]$ ची संख्या ही $s_1$ आणि $s_2$ यांमधील $]$ च्या संख्यांची बेरीज अधिक एक आहे. म्हणून $o(s_1, s_2)$ मधील $[$ ची संख्या $o(s_1,s_2)$ मधील $]$ च्या संख्येइतकी आहे.
\end{enumerate}
\end{proof}

\begin{prob}
कोणत्याही नीटस पदाची सुरुवात~$]$ ने होत नाही, हे संरचनात्मक विगमनाने सिद्ध करा.
\end{prob}

संरचनात्मक विगमनाने आणखी एक सिद्धता देऊ. चिन्हांच्या~$t$ चिन्हमालेचा अरिक्त उचित पूर्वखंड म्हणजे अशी अरिक्त~$s$ चिन्हमाला की डावीकडून वाचताना प्रत्येक चिन्हानुसार ती $t$ शी जुळते, पण $t$ अधिक लांब असते. उदाहरणार्थ, $[a \circ {}$ हा $[a \circ b]$ चा अरिक्त उचित पूर्वखंड आहे. मात्र $[b \circ {}$ तसा नाही---दुसऱ्या चिन्हावर फरक आहे---आणि $[a \circ b]$ देखील तसा नाही---दोन्हींची लांबी समान आहे. येथे अरिक्त पूर्वखंडच घेतले आहेत: रिकाम्या पूर्वखंडात उघडणारे आणि बंद करणारे दोन्ही कंस शून्य असल्यामुळे पुढील काटेकोर असमानता त्यासाठी खरी नाही.

\begin{prop}\label{mth:ind:sti:prop:initial}
नीटस पद~$t$ च्या प्रत्येक अरिक्त उचित पूर्वखंडात $[$ ची संख्या $]$ पेक्षा जास्त असते.
\end{prop}

\begin{proof}
~$t$ वर विगमन करू:
\begin{enumerate}
\item $t$ हे स्वतंत्र अक्षर आहे: मग $t$ ला कोणतेही अरिक्त उचित पूर्वखंड नाहीत.
\item $t = [s_1 \circ s_2]$, जिथे $s_1$ आणि~$s_2$ ही नीटस पदे आहेत. $r$ हा $t$ चा अरिक्त उचित पूर्वखंड असेल, तर काही शक्यता आहेत:
\begin{enumerate}
\item $r$ हे केवळ $[$ आहे: मग $r$ मध्ये $[$ ची संख्या~$]$ पेक्षा एकाने जास्त आहे.
\item $r$ हे $[r_1$ आहे, जिथे $r_1$ हा~$s_1$ चा अरिक्त उचित पूर्वखंड आहे. $s_1$ नीटस पद असल्यामुळे विगमन गृहीतकानुसार $r_1$ मध्ये $[$ हे $]$ पेक्षा जास्त आहेत; आणि $[r_1$ साठीही तेच खरे आहे.
\item $r$ हे $[s_1$ किंवा $[s_1 \circ {}$ आहे. आधीच्या निष्कर्षानुसार~$s_1$ मध्ये $[$ आणि $]$ यांच्या संख्या समान आहेत. म्हणून $[s_1$ किंवा $[s_1 \circ {}$ मध्ये $[$ ची संख्या~$]$ पेक्षा एकाने जास्त आहे.
\item $r$ हे $[s_1 \circ r_2$ आहे, जिथे $r_2$ हा~$s_2$ चा अरिक्त उचित पूर्वखंड आहे. विगमन गृहीतकानुसार $r_2$ मध्ये $[$ हे $]$ पेक्षा जास्त आहेत. आधीच्या निष्कर्षानुसार~$s_1$ मध्ये $[$ आणि $]$ यांच्या संख्या समान आहेत. म्हणून $[s_1 \circ r_2$ मध्ये $[$ ची संख्या~$]$ पेक्षा मोठी आहे.
\item $r$ हे $[s_1 \circ s_2$ आहे. आधीच्या निष्कर्षानुसार $s_1$ मध्ये $[$ आणि $]$ यांच्या संख्या समान आहेत आणि~$s_2$ मध्येही त्या समान आहेत. म्हणून $[s_1 \circ s_2$ मध्ये $[$ ची संख्या~$]$ पेक्षा एकाने जास्त आहे.
\end{enumerate}
\end{enumerate}
\end{proof}












\section{संबंध आणि फलने}\label{mth:ind:rel:sec}

नैसर्गिक संख्या किंवा नीटस पदे यांसारख्या वस्तूंचा संच विगमनाधारित रीतीने परिभाषित केल्यावर, त्या वस्तूंवरील \emph{संबंध}ही विगमनाने परिभाषित करता येतात. उदाहरणार्थ, एका नीटस पद~$t_2$ मध्ये दुसरे नीटस पद~$t_1$ भाग म्हणून येत असेल, तर पहिले पद दुसऱ्याचे उपपद आहे, अशी कल्पना करा. त्यासाठी $t_1 \sqsubseteq t_2$ हे चिन्ह वापरू. अर्थात, प्रत्येक नीटस पद स्वतःचे उपपद असते: $t \sqsubseteq t$. या संबंधाची विगमनाधारित व्याख्या अशी देता येते:

\begin{defn}
नीटस पद~$t_1$ हे~$t_2$ चे उपपद आहे हा संबंध, $t_1 \sqsubseteq t_2$, पुढीलप्रमाणे~$t_2$ वर विगमनाने परिभाषित केला आहे:
\begin{enumerate}
\item $t_2$ हे अक्षर असेल, तर $t_1 \sqsubseteq t_2$ हे $t_1 = t_2$ असेल तेव्हाच आणि फक्त तेव्हाच सत्य असते.
\item $t_2$ हे $[s_1 \circ s_2]$ असेल, तर $t_1 \sqsubseteq t_2$ हे $t_1 = t_2$, $t_1 \sqsubseteq s_1$ किंवा $t_1 \sqsubseteq s_2$ यांपैकी एक असेल तेव्हाच आणि फक्त तेव्हाच सत्य असते.
\end{enumerate}
\end{defn}

उदाहरणार्थ, या व्याख्येनुसार $\mathrm{a} \sqsubseteq [\mathrm{b} \circ \mathrm{a}]$ मिळते. कारण (2) नुसार $\mathrm{a} \sqsubseteq [\mathrm{b} \circ \mathrm{a}]$ हे $\mathrm{a} = [\mathrm{b} \circ \mathrm{a}]$, किंवा $\mathrm{a} \sqsubseteq b$, किंवा $\mathrm{a} \sqsubseteq \mathrm{a}$ यांपैकी एक असेल तेव्हाच आणि फक्त तेव्हाच सत्य असते. पहिले दोन पर्याय असत्य आहेत: $\mathrm{a}$ हे स्पष्टपणे $[\mathrm{b} \circ \mathrm{a}]$ शी एकरूप नाही; आणि (1) नुसार $\mathrm{a} \sqsubseteq \mathrm{b}$ हे $\mathrm{a} = \mathrm{b}$ असेल तेव्हाच आणि फक्त तेव्हाच सत्य असते; तेही असत्य आहे. पण पुन्हा (1) नुसार $\mathrm{a} \sqsubseteq \mathrm{a}$ हे $\mathrm{a} = \mathrm{a}$ असेल तेव्हाच आणि फक्त तेव्हाच सत्य असते; ते सत्य आहे.

ही व्याख्या यशस्वी होण्यासाठी एक अजून न सिद्ध केलेले तथ्य आवश्यक आहे: प्रत्येक नीटस पद~$t$ हे स्वतंत्र अक्षर असते, किंवा अशी \emph{एकमेव रीतीने निश्चित} होणारी नीटस पदे $s_1$ आणि $s_2$ असतात की $t = [s_1 \circ s_2]$. येथे “एकमेव रीतीने निश्चित” याचा अर्थ असा: $t = [s_1 \circ s_2]$ असेल, तर $s_1 \neq r_1$ किंवा $s_2 \neq r_2$ असताना ते \emph{त्याचबरोबर} $= [r_1 \circ r_2]$ असू शकत नाही. असे घडले, तर खंड~(2) चाच स्वतःशी संघर्ष होऊ शकतो: $t_2$ ला $[s_1 \circ s_2]$ असे वाचल्यास $t_1 \sqsubseteq t_2$ मिळेल; पण $t_2$ ला $[r_1 \circ r_2]$ असे वाचल्यास $t_1 \sqsubseteq t_2$ असत्य मिळेल. हे घडू शकत नाही हे सिद्ध करण्याआधी, जिथे ते \emph{घडू शकते} असे एक उदाहरण पाहू.

\begin{defn}
\emph{कंसविरहित पदे} विगमनाधारित रीतीने अशी परिभाषित करा:
\begin{enumerate}
\item प्रत्येक अक्षर कंसविरहित पद आहे.
\item $s_1$ आणि $s_2$ कंसविरहित पदे असतील, तर $s_1 \circ s_2$ हे कंसविरहित पद आहे.
\item इतर काहीही कंसविरहित पद नाही.
\end{enumerate}
\end{defn}

कंसविरहित पदांची उदाहरणे म्हणजे $\mathrm{a}$, $\mathrm{b} \circ \mathrm{d}$ आणि $\mathrm{b} \circ \mathrm{a} \circ \mathrm{b}$. यांच्यासाठी “उपपद” ची व्याख्या वरच्या पद्धतीने केली, तर दुसरा खंड असा असेल:
\begin{center}
$t_2 = s_1 \circ s_2$ असेल, तर $t_1 \sqsubseteq t_2$ हे $t_1 = t_2$, $t_1 \sqsubseteq s_1$ किंवा $t_1 \sqsubseteq s_2$ असेल तेव्हाच आणि फक्त तेव्हाच सत्य असते.
\end{center}

आता $\mathrm{b} \circ \mathrm{a} \circ \mathrm{b}$ हे $s_1 \circ s_2$ या रूपाचे आहे, जिथे
\begin{align*}
s_1 & = \mathrm{b} \text{ आणि} & s_2 & = \mathrm{a} \circ \mathrm{b}.
\intertext{तेच पद $r_1 \circ r_2$ या रूपाचेही आहे, जिथे}
r_1 & = \mathrm{b} \circ \mathrm{a} \text{ आणि} & r_2 &= \mathrm{b}.
\end{align*}
आता $\mathrm{a} \circ \mathrm{b}$ हे $\mathrm{b} \circ \mathrm{a} \circ \mathrm{b}$ चे उपपद आहे का? पहिल्या वाचनानुसार उत्तर होय; दुसऱ्यानुसार नाही.

एखादे नीटस पद इतर नीटस पदांपासून कसे रचले आहे हे एकमेव रीतीने निश्चित होते, या गुणधर्माला \emph{एकमेव रीतीने वाचनीयता} म्हणतात. अशा विगमनाधारित वस्तूंवरील संबंधांच्या विगमनाधारित व्याख्या महत्त्वाच्या असल्यामुळे हा गुणधर्म सिद्ध केला पाहिजे.

\begin{prop}
$t$ हे नीटस पद आहे असे समजा. मग $t$ हे स्वतंत्र अक्षर आहे; किंवा अशी एकमेव रीतीने निश्चित होणारी नीटस पदे $s_1$, $s_2$ आहेत की~$t = [s_1 \circ s_2]$.
\end{prop}

\begin{proof}
$t$ स्वतंत्र अक्षर असेल, तर अट पूर्ण होते. म्हणून $t$ स्वतंत्र अक्षर नाही असे समजा. विगमनाधारित व्याख्येवरून मग $t$ हे काही नीटस पदे $s_1$ आणि~$s_2$ यांसाठी $[s_1 \circ s_2]$ या रूपाचे असले पाहिजे. ही पदे एकमेव रीतीने निश्चित होतात हे दाखवायचे आहे: $t = [r_1 \circ r_2]$ असेल, तर $s_1 = r_1$ आणि $s_2 = r_2$.

म्हणून $t = [s_1 \circ s_2]$ आणि $t = [r_1 \circ r_2]$ असेही समजा; येथे $s_1$, $s_2$, $r_1$, $r_2$ ही नीटस पदे आहेत. $s_1 = r_1$ आणि $s_2 = r_2$ दाखवायचे आहे. प्रथम $s_1$ आणि $r_1$ एकरूप असले पाहिजेत; तसे नसतील, तर त्यांपैकी एक दुसऱ्याचा अरिक्त उचित पूर्वखंड असेल. पण \ref{mth:ind:sti:prop:initial} नुसार $s_1$ आणि~$r_1$ दोन्ही नीटस पदे असतील, तर हे अशक्य आहे. $s_1 = r_1$ असल्यावर $s_2 = r_2$ देखील स्पष्टपणे मिळते.
\end{proof}

फलनेही विगमनाधारित रीतीने परिभाषित करता येतात. उदाहरणार्थ, कोणत्याही नीटस पदातील आतल्या आत असलेल्या $[\dots]$ जोड्यांची कमाल खोली देणारे~$f$ फलन पुढीलप्रमाणे परिभाषित करता येते:

\begin{defn}
\label{mth:ind:rel:defn:depth} नीटस पदाची \emph{खोली}, $f(t)$, पुढीलप्रमाणे विगमनाधारित व्याख्येने दिली आहे:
\[
f(t) = \begin{cases}
0 & \text{ जर $t$ अक्षर असेल}\\
\max(f(s_1), f(s_2)) + 1 & \text{ जर $t = [s_1 \circ s_2]$ असेल.}
\end{cases}
\]
\end{defn}

उदाहरणार्थ,
\begin{align*}
f([\mathrm{a} \circ \mathrm{b}]) & =
\max(f(\mathrm{a}),f(\mathrm{b})) + 1 = \\
&= \max(0, 0) + 1 = 1, \text{ आणि}\\
f([[\mathrm{a} \circ \mathrm{b}] \circ \mathrm{c}]) & =
\max(f([\mathrm{a} \circ \mathrm{b}]), f(\mathrm{c})) + 1 = \\
& = \max(1,0) + 1 = 2.
\end{align*}

अर्थात, येथे $s_1$ आणि $s_2$ ही नीटस पदे आहेत असे मानतो. प्रत्येक नीटस पद स्वतंत्र अक्षर किंवा $[s_1 \circ s_2]$ या रूपाचे असते, हे तथ्य वापरतो. ते या रूपात एकाच रीतीने असू शकते, हेही पुन्हा महत्त्वाचे आहे. का ते पाहण्यासाठी आधी परिभाषित केलेल्या कंसविरहित पदांचा पुन्हा विचार करा. त्यांच्यासाठी यासारखी “व्याख्या” अशी असेल:
\[
g(t) =
\begin{cases}
0 & \text{ जर $t$ अक्षर असेल}\\
\max(g(s_1), g(s_2)) + 1 & \text{ जर $t = s_1 \circ s_2$ असेल.}
\end{cases}
\]
आता $\mathrm{a} \circ \mathrm{b} \circ \mathrm{c} \circ \mathrm{d}$ हे कंसविरहित पद पाहा. ते अनेक रीतींनी वाचता येते. उदाहरणार्थ, $s_1 \circ s_2$ या रूपात, जिथे
\begin{align*}
s_1 & = \mathrm{a} \text{ आणि} &
s_2 & = \mathrm{b} \circ \mathrm{c} \circ \mathrm{d},
\intertext{किंवा $r_1 \circ r_2$ या रूपात, जिथे}
r_1 & = \mathrm{a} \circ b \text{ आणि} &
r_2 &= \mathrm{c} \circ \mathrm{d}.
\end{align*}
पहिल्या वाचनानुसार $g$ मोजल्यास असे मिळेल:
\begin{align*}
g(s_1 \circ s_2) & =
\max(g(\mathrm{a}), g(\mathrm{b} \circ \mathrm{c} \circ \mathrm{d})) + 1 =\\ & =
\max(0,2) + 1 = 3
\intertext{तर दुसऱ्या वाचनानुसार असे मिळते}
g(r_1 \circ r_2) & =
\max(g(\mathrm{a} \circ \mathrm{b}), g(\mathrm{c} \circ \mathrm{d})) + 1=\\ &
= \max(1,1) + 1 = 2
\end{align*}
पण फलनाने नेहमी एकमेव मूल्य दिले पाहिजे. म्हणून~$g$ ची ही “व्याख्या” मुळात फलनच परिभाषित करत नाही.

\begin{prob}
$l$ फलनाची विगमनाधारित व्याख्या द्या; येथे $l(t)$ म्हणजे नीटस पद~$t$ मधील चिन्हांची संख्या.
\end{prob}

\begin{prob}
नीटस पद~$t$ वर संरचनात्मक विगमन वापरून $f(t) < l(t)$ सिद्ध करा. येथे $l(t)$ म्हणजे~$t$ मधील चिन्हांची संख्या आणि $f(t)$ म्हणजे \ref{mth:ind:rel:defn:depth} मध्ये परिभाषित केलेली~$t$ ची खोली.
\end{prob}



\part{सिद्धता उपपत्ती}\label{pt:part}
\begin{editorial}
  सिद्धता उपपत्तीवरील साहित्य अपूर्ण आणि प्रायोगिक आहे;
  ते अद्याप मुख्य PDF संकलनात समाविष्ट केलेले नाही.
\end{editorial}
\chapter{कट-निर्मूलन}\label{pt:cut:chap}
\begin{editorial}हे पूर्ण अनुवादित पूरक प्रकरण स्थानिक निदानासाठी मांडले आहे. मूळ माएहारा-सिद्धतेतील गाळलेला चर-पदाचा प्रसंग SOL6-A296 मध्ये भाषा-सीमेच्या आधारे पूर्ण केला आहे; OLINC-344 ही ऐतिहासिक स्रोत-नोंद आहे.\end{editorial}










\section{परिचय}\label{pt:cut:int:sec}

\Cut{} नियम असा आहे:
\[
\Axiom$\Gamma \fCenter \Delta, A$
\Axiom$A, \Pi \fCenter \Lambda$
\RightLabel{\Cut}
\BinaryInf$\Gamma, \Pi \fCenter \Delta, \Lambda$
\DisplayProof
\]
(सूत्र~$A$ याला
\emph{कट-सूत्र} म्हणतात.)

\Cut{} हा नियम गेंटझेन यांच्या
मूळ क्रमवर्ती कलनात~\Log{LK} होता.
गेंटझेन यांचे \emph{Hauptsatz}
(“मुख्य प्रमेय”) असे सांगते की
~\Log{LK} मध्ये सिद्धतायोग्य असलेली
प्रत्येक क्रमवर्ती \Cut{} न वापरताही
सिद्धतायोग्य आहे; म्हणजे
~\Log{LK} मधील सिद्धतांतून
\Cut{} “काढून टाकता” येतो.
आपल्या \Log{G1c} आणि \Log{G3c}
पद्धतींत \Cut{} नाही; तरी त्यांच्याबाबतही
प्रश्न पडतो: \Cut{} आणि
~\Log{G1c} चे (किंवा \Log{G3c} चे)
इतर नियम वापरणाऱ्या सिद्धता मधून
\Cut{} काढून टाकता येईल का?
आपल्या आधी निश्चित केलेल्या
परिभाषेत हाच प्रश्न असा आहे:
\Cut{} हा \Log{G1c} चा
(किंवा \Log{G3c} चा)
ग्राह्य नियम आहे का?

कट-निर्मूलन प्रमेय हा सिद्धता
उपपत्तीतील कदाचित सर्वाधिक महत्त्वाचा
आणि प्रसिद्ध निष्कर्ष आहे.
सुरुवातीला आवश्यक वाटू शकणारा
नियम \Log{G1c} आणि \Log{G3c}
यांना प्रत्यक्षात लागत नाही, हे
तो दाखवतो. कारण प्रत्यक्ष सिद्धता
औपचारिक करताना उपप्रमेयांचा
वापर कसा दाखवायचा हा प्रश्न
पडतो. गणितज्ञ प्रथम एक निष्कर्ष
सिद्ध करतात आणि मग दुसऱ्या
निष्कर्षाच्या सिद्धतेत त्याचा
आधार घेतात. म्हणून उपप्रमेय
$L$ ची सिद्धता प्रथम
$\Sequent L$ या क्रमवर्तीच्या
सिद्धता म्हणून औपचारिक करता येईल.
मग $L$ वापरणाऱ्या दुसऱ्या
निष्कर्ष~$R$ ची सिद्धता
$L \Sequent
R$ या क्रमवर्तीच्या
सिद्धता म्हणून औपचारिक करू.
केवळ $R$ ची, म्हणजे
$\Sequent R$ ची
सिद्धता आहे, हे कसे ठरवायचे?
त्या दोन सिद्धता जोडण्याचा
मार्ग हवा; \Cut{} नियमाने
ते करता येईल.

\Log{G1c} आणि \Log{G3c}
संपूर्ण आहेत, असे आपण सांगितले:
म्हणून सिद्ध होण्याची अपेक्षा
असलेल्या सर्व वैध क्रमवर्ती
ते सिद्ध करतात. हे थेटही
सिद्ध करता येते. पण गेंटझेन
लिहीत होते तेव्हा केवळ स्वयंसिद्धकीय
निष्पत्ती-पद्धती संपूर्ण आहेत,
हे ज्ञात होते. \Log{LK} ची
संपूर्णता दाखवण्यासाठी गेंटझेन
यांनी स्वयंसिद्धकीय पद्धतीतील
निष्पत्तीचे
~\Log{LK} मधील सिद्धता मध्ये
भाषांतर दिले; त्यात
\Cut{} अत्यावश्यक होता.
कट-निर्मूलन केवळ अभिजात
तर्कशास्त्रापुरते महत्त्वाचे नाही:
इतर अनेक तर्कशास्त्रांनाही
क्रमवर्ती कलने आहेत. काहींसाठी
त्या कलनांची संपूर्णता थेट
सिद्ध करणे कठीण किंवा अशक्य
असते; त्यामुळे इतर पद्धतींमधून
\Cut{} वापरून केलेले भाषांतर
हाच संपूर्णता दाखवण्याचा मार्ग असतो.
मग \Cut{} अनावश्यक आहे आणि
\Cut{} नसलेली पद्धतही संपूर्ण
आहे, हे स्थापित करण्यासाठी
सिद्धता-उपपत्तीय कट-निर्मूलन
सिद्ध करणे आवश्यक ठरते.

कट-निर्मूलनाचा उपयोग इतर
स्वतंत्रपणे महत्त्वाचे निष्कर्ष
मिळवण्यासाठीही होतो. त्यांपैकी
क्रेगचे अंतर्वेशन उपप्रमेय आणि
एरब्राँचे प्रमेय ही दोन उदाहरणे
आपण पाहू.

कट-निर्मूलनाच्या सिद्धतेत
बहुधा \Cut{} वापरणाऱ्या
सिद्धता चे त्याविना
सिद्धतेत रूपांतर कसे करावे,
हे दाखवले जाते. ही रूपांतरे
सहसा “स्थानिक” असतात:
सिद्धतेचा एक भाग
(साधारणपणे \Cut{} अनुमानाने
संपणारी उप-सिद्धता)
घेऊन त्याच्या जागी
त्याच क्रमवर्तीची दुसरी
उप-सिद्धता ठेवतो;
तिच्यातील \Cut{} अनुमान
काढलेले किंवा इतर
अनुमानांनी बदललेले असते.
अशा युक्तिवादांनी \Cut{}
असलेल्या सिद्धता चे
त्याविना सिद्धतांत रूपांतर
करणारी पायरी-पायरीची
अल्गोरिदम मिळते. उपयोगांत
कट-निर्मूलन प्रक्रिया वापरताना
यातून अधिक माहिती मिळते.
उदाहरणार्थ, अंतर्वेशन उपप्रमेयाची
अशी प्रतिमान-उपपत्तीय सिद्धता
आहे की “$A \lif B$”
वैध असेल, तर अंतर्वेशक अस्तित्वात
असतो (आत्ता अंतर्वेशक म्हणजे
काय याकडे दुर्लक्ष करा).
कट-निर्मूलन वापरणारा
सिद्धता-उपपत्तीय युक्तिवाद
$A
\lif B$ ची सिद्धता
घेऊन अंतर्वेशक देणारी
अल्गोरिदम पुरवतो. प्रतिमान
उपपत्तीच्या युक्तिवादाच्या
तुलनेत ही सिद्धता-उपपत्तीय
पद्धत \emph{रचनात्मक} आहे.

कट-निर्मूलन प्रमेय सिद्ध
करण्याच्या दोन प्रचलित पद्धती
आहेत. पहिली, गेंटझेनपासून
आलेली, सिद्धतेतील सर्वोच्च
स्थानी असलेले कट काढता
येतात हे दाखवते आणि मग
असे कट संपेपर्यंत काढत जाते.
दुसरी, मूळची श्यूटे आणि
टेट यांची, सिद्धता मधील
सर्वाधिक गुंतागुंतीच्या कटांवर
लक्ष केंद्रित करते; इतर
कोणत्याही कटच्या सूत्र
ची खोली अधिक नाही, या
अर्थाने महत्तम असलेले
कट ती एकामागोमाग काढते.
दोन्ही पद्धतींचे फायदे
आणि मर्यादा आहेत. काही
पद्धतींमध्ये एक मार्ग
चालतो, तर दुसरा अवघड
किंवा अशक्य असतो.
म्हणून आपण दोन्ही मार्ग
वर्णन करू. ~\Log{G3c}
साठी कट-निर्मूलन सिद्ध करू.
खरे तर नेहमीचा \Cut{}
नव्हे, तर \emph{संदर्भ-सामायिक कट}~\CutCS:
\[
\Axiom$\Gamma \fCenter \Delta, A$
\Axiom$A, \Gamma \fCenter \Delta$
\RightLabel{\CutCS}
\BinaryInf$\Gamma \fCenter \Delta$
\DisplayProof
\]
एखाद्या क्रमवर्ती कलनात
दुर्बलीकरण आणि आकुंचन
हे प्रत्यक्ष नियम म्हणून
(जसे \Log{G1c} मध्ये)
किंवा ग्राह्य नियम म्हणून
(जसे \Log{G3c} मध्ये)
असतील, तर \Cut{} द्वारे
\CutCS चे अनुकरण करता
येते, आणि उलटही.
आपण \Cut ऐवजी \CutCS{}
निवडतो, कारण पहिली
कट-निर्मूलन पद्धत \CutCS साठी
वर्णन करणे सोपे आहे,
आणि दुसरी पद्धत केवळ
\CutCS साठीच चालते.













\section{सर्वोच्च स्थानी असलेले कट काढणे}\label{pt:cut:top:sec}

\begin{defn}
एका \CutCS{} नियमाने संपणारी सिद्धता~$\pi$
विचारात घ्या:
\[
\AxiomC{}
\RightLabel{$\pi_1$}
\Deduce$\Gamma \fCenter \Delta, A$
\AxiomC{}
\RightLabel{$\pi_2$}
\Deduce$A, \Gamma \fCenter \Delta$
\RightLabel{\CutCS}
\BinaryInf$\Gamma \fCenter \Delta$
\DisplayProof
\]
तिच्यात अन्य \CutCS{} अनुमाने नसावीत.
$\pi$ ची \emph{कट-उंची} म्हणजे
$\cheight{\pi} = \pheight{\pi_1} + \pheight{\pi_2}$.
$\pi$ ची \emph{कट-कोटी} म्हणजे
$\cutrank{\pi} = \depth{A}$.
\end{defn}

\begin{lem}\label{pt:cut:top:lem:cut-adm-G3c}
\CutCS{} नियम~\Log{G3c} मध्ये ग्राह्य आहे.
दुसऱ्या शब्दांत, सिद्धता~$\pi$ ही एका
\CutCS{} ने संपत असेल आणि इतरत्र
फक्त~$\Log{G3c}$ चे नियम वापरत असेल,
तर त्याच अंतिम क्रमवर्तीची
\Log{G3c} मधील सिद्धता~$\pi'$ असते.
\end{lem}

\begin{proof}
कटची उंची आणि कोटी यांवरील
दुहेरी विगमनाने सिद्धता करू.
सिद्धता~$\pi$ चा शेवट असा आहे:
\[
\AxiomC{}
\RightLabel{$\pi_1$}
\Deduce$\Gamma \fCenter \Delta, A$
\AxiomC{}
\RightLabel{$\pi_2$}
\Deduce$A, \Gamma \fCenter \Delta$
\RightLabel{\CutCS}
\BinaryInf$\Gamma \fCenter \Delta$
\DisplayProof
\]

विगमनाचा आधार: $\cheight{\pi} = 0$
आणि $\cutrank{\pi} = 0$.
$\cheight{\pi} = \pheight{\pi_1}
+ \pheight{\pi_2} = 0$ असल्याने
$\Gamma \Sequent \Delta, A$ आणि
$A, \Gamma \Sequent \Delta$
ही दोन्ही स्वयंसिद्धके आहेत.
दोन प्रसंग आहेत: (a)~$A$ हा
त्यांपैकी एखाद्यात मुख्य नाही,
किंवा (b)~$A$ हा दोन्हींत मुख्य आहे.
खाली दोन्ही प्रसंगांत
$\Gamma \Sequent \Delta$ हे स्वयंसिद्धक
असलेच पाहिजे हे दाखवू
(त्यासाठी विगमन गृहीतक लागत नाही).

आता $\pi_1$ आणि $\pi_2$ स्वयंसिद्धके
आहेत की अनुमानाने संपतात, तसेच त्या
अनुमानात~$A$ मुख्य आहे की नाही,
यावरून प्रसंग वेगळे करू. प्रसंग A व~B
विगमनाचा आधार समाविष्ट करतात.
विगमन पायरीत $\cheight{\pi} > 0$
किंवा $\cutrank{\pi} > 0$
(किंवा दोन्ही) असे मानतो. विगमन
गृहीतकानुसार, एका \CutCS{} ने
संपणाऱ्या सिद्धता मधील अंतिम
\CutCS{} काढता येतो, जर तिची
कट-कोटी $= \cutrank{\pi}$ आणि
कटची उंची $< \cheight{\pi}$
असेल, किंवा तिची कट-कोटी
$< \cutrank{\pi}$ असेल.
$\cheight{\pi} = 0$ चा प्रसंग
(दोन्ही आधार स्वयंसिद्धके) A व~B
मध्ये येतो. $\cheight{\pi} > 0$
चे प्रसंग C--D मध्ये येतात.

\paragraph{प्रसंग A:}
काही $B$ हे $\Gamma$ आणि~$\Delta$
या दोन्हींत येते. $B$ अणुसूत्र
असेल, तर $\Gamma \Sequent \Delta$
हे स्वयंसिद्धक असते; अन्यथा
\ref{pt:seq:ex:prop:G3c-axioms} नुसार
त्याची \Log{G3c} मधील \CutCS{}
विरहित सिद्धता~$\pi'$ असते.
यात विगमन आधाराचा प्रसंग~(a) येतो:
$\Gamma \Sequent \Delta, A$
किंवा $A, \Gamma \Sequent \Delta$
हे स्वयंसिद्धक असून त्यात $A$
मुख्य नसेल, तर अन्य अणुसूत्र~$B$
हे $\Gamma$ आणि~$\Delta$ या
दोन्हींत असते. विगमन पायरीतही
एखादा आधार स्वयंसिद्धक असून त्यात
$A$ मुख्य नसेल, तर हा प्रसंग
लागू होतो; दुसरा आधार नियमाचा
निष्कर्ष असला तरी चालते.

\paragraph{प्रसंग B:}
दोन्ही आधार स्वयंसिद्धके आहेत
आणि $A$ मुख्य, म्हणून अणुसूत्र,
आहे. डावा आधार
$\Gamma \Sequent \Delta, A$
पाहा. या $\Gamma \Sequent \Delta, A$
मध्ये $A$ मुख्य असल्यामुळे ते
$\Gamma$ मध्ये असलेच पाहिजे
(नसते तर ही क्रमवर्ती स्वयंसिद्धक
ठरली नसती). त्याचप्रमाणे,
$A$ हे $\Delta$ मध्येही असलेच
पाहिजे; अन्यथा
$A, \Gamma \Sequent \Delta$
स्वयंसिद्धक ठरणार नाही.
त्यामुळे $A$ अणुसूत्र असून
$\Gamma$ आणि $\Delta$ या
दोन्हींत येते; म्हणजे
$\Gamma \Sequent \Delta$
हे स्वयंसिद्धक आहे. हा
विगमन आधाराचा प्रसंग~(b) आहे.

\begin{center}
प्रसंग A आणि B चे पर्यायी रूप
\end{center}

\paragraph{प्रसंग A:}
कटच्या आधारांपैकी एक
$C, \Gamma' \Sequent \Delta', C$
या रूपाचे स्वयंसिद्धक आहे,
जेथे $C$ अणुसूत्र आहे.
$C$ हे कटचे सूत्र~$A$
नसेल, तर $C$ हे $\Gamma$
आणि~$\Delta$ या दोन्हींत
असलेच पाहिजे. मग
$\Gamma \Sequent \Delta$
हे स्वयंसिद्धक असून तेच~$\pi'$
घेता येते. $C \ident A$
कटचे सूत्र असेल, तर
डावा आधार स्वयंसिद्धक असल्यास
$A \in \Gamma$ आणि
$\Gamma = \Gamma', A$;
उजवा आधार स्वयंसिद्धक असल्यास
$A \in \Delta$ आणि
$\Delta = \Delta', A$.
पहिल्या प्रसंगात $\pi_2$
ही $A, A, \Gamma' \Sequent \Delta$
ची सिद्धता असते;
\LeftR{\Contraction} ची ग्राह्यता
(\ref{pt:seq:inv:prop:G3c-cont-adm})
वापरून $\Gamma \Sequent \Delta$
ची सिद्धता मिळते.
दुसऱ्या प्रसंगात $\pi_1$
ही $\Gamma \Sequent \Delta', A, A$
ची सिद्धता असते;
\RightR{\Contraction} ची ग्राह्यता
वापरून $\Gamma \Sequent \Delta$
ची सिद्धता मिळते.


\paragraph{प्रसंग B:}
एखादा आधार
$\lfalse, \Gamma' \Sequent \Delta'$
या रूपाचा स्वयंसिद्धक आहे.
डावा आधार असा असेल, तर
$\lfalse \in \Gamma$ आणि
$\Gamma \Sequent \Delta$
हेही स्वयंसिद्धक आहे. उजवा
आधार असा असेल, तर
$\lfalse \in \Gamma$
(मग पुन्हा $\Gamma \Sequent \Delta$
स्वयंसिद्धक असते) किंवा
$A \ident \lfalse$.
दुसऱ्या उपप्रसंगात $\pi_1$
ही $\Gamma \Sequent \Delta, \lfalse$
ची सिद्धता आहे. $\pi_1$
मधील प्रत्येक क्रमवर्तीच्या
उत्तरांगातून $\lfalse$ ची
एक प्रत काढून
$\Gamma \Sequent \Delta$
ची सिद्धता मिळते.
(सराव: $\pheight{\pi_1}$
वरील विगमनाने हे तपासा.)

आता प्रसंग~A किंवा B लागू
होत नाही असे समजा. उर्वरित
प्रसंगांत $\cheight{\pi} > 0$;
म्हणजे किमान एका आधाराचा
शेवट नियमाच्या अनुमानाने होतो.

\paragraph{प्रसंग C आणि D: कटचे स्थान बदलणे.}
प्रसंग~C आणि~D मधील सर्व शक्यतांत
एक सामाईक नमुना आहे. सुरुवातीला
\CutCS{} ने संपणारी सिद्धता असते;
तिच्या एका आधाराचा निष्कर्ष नियम~$R$
ने मिळतो, पण कटचे सूत्र~$A$
त्या नियमात मुख्य नसते (काही अन्य
सूत्र~$D$ मुख्य असते).
या सिद्धता मध्ये \CutCS{}
हा नियम~$R$ नंतर येतो; तिचे
योजनात्मक रूप असे:
\[
\AxiomC{}
\RightLabel{$\pi_1$}
\Deduce$\Gamma \fCenter \Delta', D, A$
\AxiomC{}
\RightLabel{$\pi_3$}
\Deduce$A, \Gamma, \Pi \fCenter \Delta', \Lambda$
\RightLabel{$R$}
\UnaryInf$A, \Gamma \fCenter \Delta', D$
\RightSubproofLabel{$\pi_2$}
\RightLabel{\CutCS}
\BinaryInf$\Gamma \fCenter \Delta', D$
\DisplayProof
\]
हे रूप फक्त त्या प्रसंगाचे आहे
जेथे $R$ हा एक-आधाराचा उजवा
नियम आहे, $A$ हे $R$ चे
मुख्य सूत्र नाही, आणि
\emph{मुख्य} असलेले सूत्र~$D$
हे \CutCS{} च्या उजव्या
आधाराच्या उत्तरांगात आहे.
$D$ पूर्वांगात असेल (म्हणजे
$R$ डावा नियम असेल) किंवा
\CutCS{} च्या डाव्या आधारात
असेल, तर रूप अर्थातच बदलेल.
$\Pi$ आणि~$\Lambda$ मधील
सूत्रे या नियम~$R$
च्या सक्रिय सूत्रे दर्शवतात.

आता हा क्रम उलटून
$R$ हा~\CutCS{} नंतर येणारी
सिद्धता विचारात घ्या:
\[
\AxiomC{}
\RightLabel{$\pi_1'$}
\Deduce$\Gamma, \Pi \fCenter \Delta', D, \Lambda, A$
\AxiomC{}
\RightLabel{$\pi_3'$}
\Deduce$A, \Gamma, \Pi \fCenter \Delta', D, \Lambda$
\RightLabel{\CutCS}
\BinaryInf$\Gamma, \Pi \fCenter \Delta', D, \Lambda$
\RightLabel{$R$}
\UnaryInf$\Gamma, \fCenter \Delta', D, D$
\DisplayProof
\]
$\pi_1$ आणि $\pi_3$ यांवर
उंची न वाढवणारे दुर्बलीकरण
लावून अनुक्रमे $\pi_1'$ आणि
$\pi_3'$ मिळवतो.

\CutCS{} आणि~$R$ यांचा क्रम
बदलत असल्यामुळे याला
“कटला~$R$ च्या वर नेणे” म्हणतात.
याने \CutCS{} च्या उजव्या
आधारावर संपणाऱ्या सिद्धता
ची उंची कमी होते; म्हणून कटची
उंची कमी होते. म्हणजे
नव्या \CutCS{} वरील
उप-सिद्धता ची उंची
$\pi_1$ आणि $\pi_3$ यांच्या
उंचीपेक्षा मोठी नसते, म्हणून
कटची उंची कमी होते
(उजवीकडील एक अनुमान काढले आहे).
आता विगमन गृहीतक लावून \CutCS{}
अनुमान काढता येते. शेवटी
$\RightR{\Contraction}$ ची ग्राह्यता
वापरून~$D$ ची अतिरिक्त
प्रत काढतो.

\paragraph{प्रसंग C:}
$\pi_1$ आणि $\pi_2$ यांपैकी
किमान एकाचा शेवट एक-आधाराच्या
नियम~$R$ ने होतो, आणि त्यात
$A$ मुख्य नसते. एकूण 18 प्रसंग
आहेत: \CutCS{} च्या डाव्या की
उजव्या आधारात $A$ मुख्य नाही
आणि नऊ एक-आधाराच्या नियमांपैकी
(\LeftR{\lnot}, \RightR{\lnot},
\RightR{\lor}, \LeftR{\land},
\RightR{\lif}, आणि चार संख्यक नियम)
$R$ कोणता आहे, यावर ते अवलंबून
असतात. आपण फक्त दोन उदाहरणे
पाहू: $\pi_2$ च्या अंतिम
अनुमानात $A$ मुख्य नाही आणि ते
अनुमान $\RightR{\lif}$ किंवा
\LeftR{\lexists} आहे.
इतर प्रसंग यांसारखेच असून
सरावासाठी ठेवतो.

समजा $\pi$ चा शेवट असा आहे:
\[
\AxiomC{}
\RightLabel{$\pi_1$}
\Deduce$\Gamma \fCenter \Delta', B \lif C, A$
\AxiomC{}
\RightLabel{$\pi_3$}
\Deduce$A, B, \Gamma \fCenter \Delta', C$
\RightLabel{$\RightR{\lif}$}
\UnaryInf$A, \Gamma \fCenter \Delta', B \lif C$
\RightSubproofLabel{$\pi_2$}
\RightLabel{\CutCS}
\BinaryInf$\Gamma \fCenter \Delta', B \lif C$
\DisplayProof
\]
येथे $\Delta = \Delta', B \lif C$.
$\pi_1$ ला उंची न वाढवणारे
दुर्बलीकरण
(\ref{pt:seq:adm:prop:weak-G3c-adm})
लावून $B, \Gamma \fCenter
\Delta', B \lif C, C, A$
ची सिद्धता $\pi_1'$ मिळवतो
(डावीकडे $B$ आणि उजवीकडे
$C$ घालतो). त्याचप्रमाणे
\ref{pt:seq:adm:prop:weak-G3c-adm}
हे~$\pi_3$ ला लावून
$A, B, \Gamma \fCenter
\Delta', B \lif C, C$
ची सिद्धता~$\pi_3'$ मिळते
(उजवीकडे $B \lif C$ घालतो).
विगमन गृहीतक पुढील सिद्धतेला
लागू होते:
\[
\AxiomC{}
\RightLabel{$\pi_1'$}
\Deduce$B, \Gamma \fCenter \Delta', B \lif C, C, A$
\AxiomC{}
\RightLabel{$\pi_3'$}
\Deduce$A, B, \Gamma \fCenter \Delta', B \lif C, C$
\RightLabel{\CutCS}
\BinaryInf$B, \Gamma \fCenter \Delta', B \lif C, C$
\DisplayProof
\]
यातील कटचे सूत्र~$A$
आहे: या सिद्धता ची कट-कोटी
$\pi$ इतकीच, पण कटची उंची
कमी आहे. त्यामुळे \CutCS{}
विरहित \Log{G3c} मधील
$B, \Gamma \fCenter
\Delta', B \lif C, C$
ची सिद्धता~$\pi_4$ मिळते.
$\RightR{\lif}$ लावून
$\Gamma \fCenter \Delta',
B \lif C, B \lif C$
ची सिद्धता मिळते:
\[
\AxiomC{}
\RightLabel{$\pi_4$}
\Deduce$B, \Gamma \fCenter \Delta', B \lif C, C$
\RightLabel{$\RightR{\lif}$}
\UnaryInf$\Gamma \fCenter \Delta', B \lif C, B \lif C$
\DisplayProof
\]
$\Gamma \Sequent \Delta$,
म्हणजे $\Gamma \fCenter
\Delta', B \lif C$,
याची सिद्धता~$\pi'$
मिळवण्यासाठी
\ref{pt:seq:inv:prop:G3c-cont-adm},
म्हणजे आकुंचनाची ग्राह्यता,
वापरा.

आता $\pi$ चा शेवट
पुढीलप्रमाणे आहे असे समजा:
\[
\AxiomC{}
\RightLabel{$\pi_1$}
\Deduce$\lexists[x][B(x)], \Gamma' \fCenter \Delta, A$
\AxiomC{}
\RightLabel{$\pi_3$}
\Deduce$A, B(c), \Gamma' \fCenter \Delta$
\RightLabel{$\LeftR{\lexists}$}
\UnaryInf$A, \lexists[x][B(x)], \Gamma' \fCenter \Delta$
\RightSubproofLabel{$\pi_2$}
\RightLabel{\CutCS}
\BinaryInf$\lexists[x][B(x)], \Gamma' \fCenter \Delta$
\DisplayProof
\]
पुन्हा पुढील सिद्धतेला
विगमन गृहीतक लावू:
\[
\AxiomC{}
\RightLabel{$\pi_1[B(c) \Sequent {}]$}
\Deduce$\lexists[x][B(x)], B(c), \Gamma' \fCenter \Delta, A$
\AxiomC{}
\LeftLabel{$\pi_3[\lexists[x][B(x)] \Sequent {}]$}
\Deduce$A, \lexists[x][B(x)], B(c), \Gamma' \fCenter \Delta$
\RightLabel{\CutCS}
\BinaryInf$\lexists[x][B(x)], B(c), \Gamma' \fCenter \Delta$
\RightLabel{\LeftR{\lexists}}
\UnaryInf$\lexists[x][B(x)], \lexists[x][B(x)], \Gamma' \fCenter \Delta$
\DisplayProof
\]
येथे आयगेन चराची अट पूर्ण होते.
कारण $\pi_2$ मधील
\LeftR{\lexists} वरील
आयगेन चराच्या अटीनुसार
$c$ हे $\lexists[x][B(x)]$,
$\Gamma'$ किंवा~$\Delta$
यांत येत नाही. शेवटी
\ref{pt:seq:inv:prop:G3c-cont-adm}
लावा.

\paragraph{प्रसंग D:}
$\pi_1$ आणि $\pi_2$ यांपैकी
किमान एकाचा शेवट दोन-आधारांच्या
नियम~$R$ ने होतो आणि त्यात
$A$ मुख्य नसते. असे सहा
प्रसंग आहेत. $\pi_2$ चा
शेवट~$\RightR{\land}$ ने होत
असताना त्याच्या अंतिम अनुमानात
$A$ मुख्य नसतो हा प्रसंग
पाहू; इतर तसेच आहेत.

समजा $\pi$ चा शेवट असा आहे:
\[
\AxiomC{}
\RightLabel{$\pi_1$}
\Deduce$\Gamma \fCenter \Delta', B \land C, A$
\AxiomC{}
\RightLabel{$\pi_3$}
\Deduce$A, \Gamma \fCenter \Delta', B$
\AxiomC{}
\RightLabel{$\pi_4$}
\Deduce$A, \Gamma \fCenter \Delta', C$
\RightLabel{$\RightR{\land}$}
\BinaryInf$A, \Gamma \fCenter \Delta', B \land C$
\RightSubproofLabel{$\pi_2$}
\RightLabel{\CutCS}
\BinaryInf$\Gamma \fCenter \Delta$
\DisplayProof
\]
विगमन गृहीतक पुढील
सिद्धता ना लागू होते:
\[
\AxiomC{}
\RightLabel{$\pi_1'$}
\Deduce$\Gamma \fCenter \Delta', B \land C, B, A$
\AxiomC{}
\RightLabel{$\pi_3'$}
\Deduce$A, \Gamma \fCenter \Delta', B \land C, B$
\RightLabel{\CutCS}
\BinaryInf$\Gamma \fCenter \Delta', B \land C, B$
\DisplayProof
\]
आणि
\[
\AxiomC{}
\RightLabel{$\pi_1''$}
\Deduce$\Gamma \fCenter \Delta', B \land C, C, A$
\AxiomC{}
\RightLabel{$\pi_4'$}
\Deduce$A, \Gamma \fCenter \Delta', B \land C, C$
\RightLabel{\CutCS}
\BinaryInf$\Gamma \fCenter \Delta', B \land C, C$
\DisplayProof
\]
येथे \Log{G3c} मधील
सिद्धता $\pi_1'$ आणि
$\pi_1''$ या $\pi_1$
पासून उंची न वाढवणाऱ्या
दुर्बलीकरणाने
(\ref{pt:seq:adm:prop:weak-G3c-adm})
मिळतात: एकदा उजवीकडे $B$
आणि एकदा $C$ घालून.
सिद्धता $\pi_3'$ आणि
$\pi_4'$ या अनुक्रमे $\pi_3$
आणि $\pi_4$ पासून
उंची न वाढवणाऱ्या
दुर्बलीकरणाने मिळतात:
यात $\pi_3$ व $\pi_4$ यांची
उंची टिकवून उजवीकडे $B \land C$
घालून.


विगमन गृहीतकानुसार
\Log{G3c} मधील
$\Gamma \fCenter \Delta', B \land C, B$
आणि
$\Gamma \fCenter \Delta', B \land C, C$
यांच्या सिद्धता~$\pi_5$ आणि~$\pi_6$
मिळतात. $\RightR{\land}$
नियमाने त्यांना जोडून
$\Gamma \Sequent \Delta',
B \land C, B \land C$
ची सिद्धता मिळते:
\[
\AxiomC{}
\RightLabel{$\pi_5$}
\Deduce$\Gamma \fCenter \Delta', B \land C, B$
\AxiomC{}
\RightLabel{$\pi_6$}
\Deduce$\Gamma \fCenter \Delta', B \land C, C$
\RightLabel{\RightR{\land}}
\BinaryInf$\Gamma \fCenter \Delta', B \land C, B \land C$
\DisplayProof
\]
$\Gamma \Sequent \Delta$,
म्हणजे $\Gamma \fCenter
\Delta', B \land C$,
याची सिद्धता मिळवण्यासाठी
\ref{pt:seq:inv:prop:G3c-cont-adm},
म्हणजे आकुंचनाची ग्राह्यता,
लावा.


\paragraph{प्रसंग E: कटची कोटी कमी करणे.}
शेवटचा प्रसंग असा की $\pi_1$
आणि $\pi_2$ दोन्हींचा शेवट
अशा अनुमानांनी होतो ज्यांत
$A$ मुख्य असते. $A$ च्या
संभाव्य मुख्य संकारक नुसार
सहा प्रसंग होतात.

प्रत्येक प्रसंगात कटचे
सूत्र~$A$ असलेले
\CutCS{} अनुमान आपण $A$
च्या उप-सूत्रे वरील
एक किंवा अधिक कटांत बदलतो;
म्हणजे ज्या अनुमानांत $A$
हे मुख्य सूत्र असते,
त्यांच्या सक्रिय सूत्रे
वर कट करतो. म्हणून या
प्रसंगांना कट-अनुमानाचे
“न्यूनीकरण” किंवा “परिवर्तन”
म्हणतात.

\begin{enumerate}
\item $A \ident \lnot B$
असेल, तर $\pi$ चा शेवट असा:
\[
\AxiomC{}
\RightLabel{$\pi_1'$}
\Deduce$B, \Gamma \fCenter \Delta$
\RightLabel{\RightR{\lnot}}
\UnaryInf$\Gamma \fCenter \Delta, \lnot B$
\LeftSubproofLabel{$\pi_1$}
\AxiomC{}
\RightLabel{$\pi_2'$}
\Deduce$\Gamma \fCenter \Delta, B$
\RightLabel{\LeftR{\lnot}}
\UnaryInf$\lnot B, \Gamma \fCenter \Delta$
\RightSubproofLabel{$\pi_2$}
\RightLabel{\CutCS}
\BinaryInf$\Gamma \fCenter \Delta$
\DisplayProof
\]
विगमन गृहीतकानुसार पुढील
सिद्धतेतून \CutCS{} काढता येतो:
\[
\AxiomC{}
\RightLabel{$\pi_2'$}
\Deduce$\Gamma \fCenter \Delta, B$
\AxiomC{}
\RightLabel{$\pi_1'$}
\Deduce$B, \Gamma \fCenter \Delta$
\RightLabel{\CutCS}
\BinaryInf$\Gamma \fCenter \Delta$
\DisplayProof
\]

\item
$A \ident (B \lor C)$ असेल,
तर $\pi$ चा शेवट असा आहे:
\[
\AxiomC{}
\RightLabel{$\pi_1'$}
\Deduce$\Gamma \fCenter \Delta, B, C$
\RightLabel{\RightR{\lor}}
\UnaryInf$\Gamma \fCenter \Delta, B \lor C$
\LeftSubproofLabel{$\pi_1$}
\AxiomC{}
\RightLabel{$\pi_2'$}
\Deduce$B, \Gamma \fCenter \Delta$
\AxiomC{}
\RightLabel{$\pi_2''$}
\Deduce$C, \Gamma \fCenter \Delta$
\RightLabel{\LeftR{\lor}}
\BinaryInf$B \lor C, \Gamma \fCenter \Delta$
\RightSubproofLabel{$\pi_2$}
\RightLabel{\CutCS}
\BinaryInf$\Gamma \fCenter \Delta$
\DisplayProof
\]
\CutCS{} ने संपणारी पुढील
सिद्धता विचारात घ्या:
\[
\AxiomC{}
\RightLabel{$\pi_1'$}
\Deduce$\Gamma \fCenter \Delta, B, C$
\AxiomC{}
\RightLabel{$\pi_2'''$}
\Deduce$C, \Gamma \fCenter \Delta, B$
\RightLabel{\CutCS}
\BinaryInf$\Gamma \fCenter \Delta, B$
\DisplayProof
\]
येथे $\pi_2'''$ ही $\pi_2''$
पासून उंची न वाढवणारे
दुर्बलीकरण लावून मिळते.
$\depth{C} < \depth{B \lor C}$
असल्यामुळे तिची कट-कोटी
$< \cutr{\pi}$ आहे.
विगमन गृहीतकानुसार
$\Gamma \fCenter \Delta, B$
ची \Log{G3c} मधील
सिद्धता~$\pi_1''$ मिळते.
आता \CutCS{} ने संपणारी
पुढील सिद्धता पाहा:
\[
\AxiomC{}
\RightLabel{$\pi_1''$}
\Deduce$\Gamma \fCenter \Delta, B$
\AxiomC{}
\RightLabel{$\pi_2'$}
\Deduce$B, \Gamma \fCenter \Delta$
\RightLabel{\CutCS}
\BinaryInf$\Gamma \fCenter \Delta$
\DisplayProof
\]
हिची कट-कोटी
$= \depth{B} < \depth{B \lor C}$
असल्यामुळे विगमन गृहीतक
पुन्हा लागू होते. त्यातून
$\Gamma \fCenter \Delta$
ची सिद्धता $\pi'$ मिळते.

\item $A \ident (B \land C)$.
हा प्रसंग सरावासाठी आहे.

\item $A \ident (B \lif C)$
असेल, तर $\pi$ चा शेवट असा:
\[
\AxiomC{}
\RightLabel{$\pi_1'$}
\Deduce$B, \Gamma \fCenter \Delta, C$
\RightLabel{\RightR{\lif}}
\UnaryInf$\Gamma \fCenter \Delta, B \lif C$
\LeftSubproofLabel{$\pi_1$}
\AxiomC{}
\RightLabel{$\pi_2'$}
\Deduce$\Gamma \fCenter \Delta, B$
\AxiomC{}
\RightLabel{$\pi_2''$}
\Deduce$C, \Gamma \fCenter \Delta$
\RightLabel{\LeftR{\lif}}
\BinaryInf$B \lif C, \Gamma \fCenter \Delta$
\RightSubproofLabel{$\pi_2$}
\RightLabel{\CutCS}
\BinaryInf$\Gamma \fCenter \Delta$
\DisplayProof
\]
$\pi_2''$ च्या पूर्वांगात $B$ घालून मिळालेल्या
सिद्धतेला $\pi_2'''$ लिहू. आता दोन्ही आधारांचा सामायिक
संदर्भ $B,\Gamma$ आहे. \CutCS{} ने संपणाऱ्या
पुढील सिद्धता ची कट-कोटी $\pi$ पेक्षा कमी आहे:
\[
\AxiomC{}
\RightLabel{$\pi_1'$}
\Deduce$B, \Gamma \fCenter \Delta, C$
\AxiomC{}
\RightLabel{$\pi_2'''$}
\Deduce$C, B, \Gamma \fCenter \Delta$
\RightLabel{\CutCS}
\BinaryInf$B, \Gamma \fCenter \Delta$
\DisplayProof
\]
विगमन गृहीतकानुसार
$B, \Gamma \fCenter \Delta$
ची सिद्धता~$\pi_1''$ मिळते.
आता \CutCS{} ने संपणारी
पुढील सिद्धता पाहा:
\[
\AxiomC{}
\RightLabel{$\pi_2'$}
\Deduce$\Gamma \fCenter \Delta, B$
\AxiomC{}
\RightLabel{$\pi_1''$}
\Deduce$B, \Gamma \fCenter \Delta$
\RightLabel{\CutCS}
\BinaryInf$\Gamma \fCenter \Delta$
\DisplayProof
\]
तिची कट-कोटीही $\pi$
च्या कट-कोटीपेक्षा कमी आहे.
म्हणून विगमन गृहीतकानुसार
$\Gamma \fCenter \Delta$
ची \Log{G3c}-सिद्धता~$\pi'$
मिळते.
\item $A \ident \lforall[x][B(x)]$
असेल, तर $\pi$ चा शेवट असा:
\[
\AxiomC{}
\RightLabel{$\pi_1'$}
\Deduce$\Gamma \fCenter \Delta, B(c)$
\RightLabel{\RightR{\lforall}}
\UnaryInf$\Gamma \fCenter \Delta, \lforall[x][B(x)]$
\LeftSubproofLabel{$\pi_1$}
\AxiomC{}
\RightLabel{$\pi_2'$}
\Deduce$B(t), \lforall[x][B(x)], \Gamma \fCenter \Delta$
\RightLabel{\LeftR{\lforall}}
\UnaryInf$\lforall[x][B(x)], \Gamma \fCenter \Delta$
\RightSubproofLabel{$\pi_2$}
\RightLabel{\CutCS}
\BinaryInf$\Gamma \fCenter \Delta$
\DisplayProof
\]
विगमन गृहीतकानुसार
पुढील सिद्धतेतून \CutCS{}
काढता येतो:
\[
\AxiomC{}
\RightLabel{$\pi_1''$}
\Deduce$B(t), \Gamma \fCenter \Delta, \lforall[x][B(x)]$
\AxiomC{}
\RightLabel{$\pi_2'$}
\Deduce$B(t), \lforall[x][B(x)], \Gamma \fCenter \Delta$
\RightLabel{\CutCS}
\BinaryInf$B(t), \Gamma \fCenter \Delta$
\DisplayProof
\]
येथे $\pi_1''$ ही $\pi_1$
पासून (फक्त~$\pi_1'$ पासून
नव्हे!) उंची न वाढवणारे
दुर्बलीकरण लावून मिळते.
(या सिद्धता ची कट-कोटी
तेवढीच, पण कटची उंची
कमी आहे.) त्यामुळे
$B(t), \Gamma \fCenter \Delta$
ची सिद्धता~$\pi_2''$ मिळते.

सिद्धता~$\pi_1'$ ची
अंतिम क्रमवर्ती
$\Gamma \Sequent \Delta, B(c)$
आहे, जेथे $c$ हा
$\RightR{\lforall}$ अनुमानाचा
आयगेन चर आहे. म्हणून $c$ हे
$B(x)$, $\Gamma$ किंवा~$\Delta$
यांत येत नाही. सिद्धता~$\pi$
नियमित आहे असे मानू; त्यामुळे
$c$ हा फक्त पुढील
$\RightR{\lforall}$ अनुमानाचा
आयगेन चर आहे आणि फक्त त्याच्या
वर येतो. म्हणजे तो फक्त~$\pi_1'$
मध्ये येतो आणि $\pi_1'$ मधील
कोणत्याही अनुमानाचा आयगेन चर
नसतो. $c$ हे $\pi_2$ मध्ये
येत नसल्यामुळे $t$ मध्येही
येत नाही.
\ref{pt:seq:qua:lem:subst-regular}
नुसार सिद्धता~$\Subst{\pi_1'}{t}{c}$
ही $\Gamma \Sequent \Delta, B(t)$
ची सिद्धता आहे. आता पुढील
सिद्धता ला विगमन गृहीतक लावू:
\[
\AxiomC{}
\RightLabel{$\Subst{\pi_1'}{t}{c}$}
\Deduce$\Gamma \fCenter \Delta, B(t)$
\AxiomC{}
\RightLabel{$\pi_2''$}
\Deduce$B(t), \Gamma \fCenter \Delta$
\RightLabel{\CutCS}
\BinaryInf$\Gamma \fCenter \Delta$
\DisplayProof
\]
(कटचे सूत्र~$B(t)$
असल्यामुळे विगमन गृहीतक
लागू होते; या सिद्धता ची
कट-कोटी $\pi$ पेक्षा कमी आहे.)
\item $A \ident \lexists[x][B(x)]$
हा प्रसंग सरावासाठी ठेवतो.
\end{enumerate}
\end{proof}

\begin{prob}
\ref{pt:cut:top:lem:cut-adm-G3c} च्या
सिद्धतेतील प्रसंग~D चा
पुढील उपप्रसंग तपासा:
$\pi_1$ चा शेवट
$\LeftR{\lif}$ ने होतो,
पण त्याच्या अंतिम अनुमानात
$A$ मुख्य नसते.
(टीप: येथे नेमके काय सिद्ध
करायचे, म्हणजे या प्रसंगात
$\pi$ चे रूप काय, हे स्पष्ट
मांडणेही सरावाचा भाग आहे.)
\end{prob}

\begin{prob}
\ref{pt:cut:top:lem:cut-adm-G3c} च्या
सिद्धतेतील प्रसंग~D चा
पुढील उपप्रसंग तपासा:
$\pi_1$ चा शेवट
$\LeftR{\lforall}$ ने होतो,
पण त्याच्या अंतिम अनुमानात
$A$ मुख्य नसते.
\end{prob}

\begin{prob}
\ref{pt:cut:top:lem:cut-adm-G3c} च्या
सिद्धतेतील प्रसंग~E चा
पुढील उपप्रसंग तपासा:
$\pi_1$ आणि $\pi_2$
या दोन्हींच्या अंतिम अनुमानात
$A$ मुख्य आहे आणि
$A \ident (B \land C)$.
\end{prob}

\begin{prob}
\ref{pt:cut:top:lem:cut-adm-G3c} च्या
सिद्धतेतील प्रसंग~E चा
पुढील उपप्रसंग तपासा:
$\pi_1$ आणि $\pi_2$
या दोन्हींच्या अंतिम अनुमानात
$A$ मुख्य आहे आणि
$A \ident \lexists[x][B(x)]$.
\end{prob}

प्रसंग~E मध्ये विगमन
गृहीतक ज्या \CutCS{} ने
संपणाऱ्या सिद्धता वर
लावतो, तिची कटची उंची
मूळ सिद्धतेच्या कटच्या
उंची पेक्षा मोठी
असू शकते. उदाहरणार्थ,
$A \ident (B \lif C)$
असताना $\pi$ चे दोन
\CutCS{} अनुमाने असणाऱ्या
सिद्धता मध्ये रूपांतर
केले:
\[
\AxiomC{}
\RightLabel{$\pi_2'$}
\Deduce$\Gamma \fCenter \Delta, B$
\AxiomC{}
\RightLabel{$\pi_1'$}
\Deduce$B, \Gamma \fCenter \Delta, C$
\AxiomC{}
\RightLabel{$\pi_2'''$}
\Deduce$C, B, \Gamma \fCenter \Delta$
\RightLabel{\CutCS}
\BinaryInf$B, \Gamma \fCenter \Delta$
\RightSubproofLabel{$\pi_1''$}
\RightLabel{\CutCS}
\BinaryInf$\Gamma \fCenter \Delta$
\DisplayProof
\]
वरचा \CutCS{}, जो
$\pi_1'$ आणि $\pi_2''$
यांच्यामध्ये आहे, त्याची
कट-कोटी आणि कटची उंची
दोन्ही $\pi$ पेक्षा कमी आहेत.
पण विगमन गृहीतक लावल्याने
मिळालेली सिद्धता~$\pi_1''$
हिची कटची उंची आता
$\pi$ पेक्षा कमी असण्याची
गरज नाही. तरी खालच्या
\CutCS{} चे कटचे सूत्र
$B$ असल्यामुळे त्याची
\emph{कोटी} $\pi$ च्या
कट-कोटीपेक्षा कमी आहे.
येथे $\pi_2'''$ हे आधीप्रमाणे $\pi_2''$ चे
दुर्बलीकरण आहे; त्यामुळे वरच्या कटाच्या दोन्ही आधारांत
$B,\Gamma$ हा सामायिक संदर्भ आहे. खालचा कट $B$
काढतो; त्याचा निष्कर्ष $\Gamma\Sequent\Delta$ आहे.

\ref{pt:cut:top:lem:cut-adm-G3c} नुसार
सिद्धता मधील एक सर्वोच्च
\CutCS{} अनुमान काढता येते.
हे उपप्रमेय सर्वोच्च~\CutCS{}
ने संपणाऱ्या उप-सिद्धता
वर एकामागोमाग लावून
सिद्धता मधील कितीही
\CutCS{} अनुमाने काढता
येतात.

\begin{thm}\label{pt:cut:top:thm:cut-elim-G3c-topmost}
$\Log{G3c} + \CutCS$ मधील
कोणतीही सिद्धता $\pi$
ही~\Log{G3c} मधील
सिद्धता मध्ये रूपांतरित
करता येते.
\end{thm}

\begin{proof}
  $\pi$ मधील \CutCS{}
  अनुमानांची संख्या~$n$
  हिच्यावर विगमन करू.
  $n=0$ असेल, तर काही
  करायचे नाही: $\pi$
  ही आधीच~\Log{G3c}
  मधील सिद्धता आहे.
  अन्यथा, सर्वोच्च
  \CutCS{} अनुमानावर
  संपणारी उप-सिद्धता~$\pi'$
  विचारात घ्या. तिचे
  एकमेव \CutCS{} अंतिम
  अनुमान असते.
  \ref{pt:cut:top:lem:cut-adm-G3c}
  लावून तिचे \CutCS{}
  विरहित सिद्धता~$\pi''$
  मध्ये रूपांतर करा आणि
  मूळ सिद्धतेत
  उप-सिद्धता~$\pi'$
  च्या जागी~$\pi''$ ठेवा.
  मिळालेल्या सिद्धता
  मध्ये $n-1$ \CutCS{}
  अनुमाने उरतात; म्हणून
  विगमन गृहीतक लागू होते.
\end{proof}













\section{सर्वोच्च कोटीचे कट काढून टाकणे}\label{pt:cut:inv:sec}

\ref{pt:seq:inv:lem:G3c-invert} च्या सिद्धतेत दिसले
की \Log{G3c} ने एखाद्या तार्किक नियमाचा निष्कर्ष
(तो नियम \RightR{\lexists} किंवा \LeftR{\lforall}
नसेल तर) सिद्ध केला, तर त्या नियमाची प्रत्येक
आधार-क्रमवर्तीसुद्धा सिद्धतेची उंची न वाढवता
सिद्ध करता येते. याहून अधिक दाखवता येते:
हेच $\Log{G3c} + \CutCS$ साठीही लागू होते.

आधीप्रमाणे, \CutCS{} अनुमानाची \emph{कट-कोटी}
म्हणजे त्यातील कटच्या सूत्राची खोली.

\begin{lem}\label{pt:cut:inv:lem:inv-G3c-cut}
$\Log{G3c} + \CutCS$ मध्ये \RightR{\lexists}
किंवा \LeftR{\lforall} सोडून अन्य तार्किक नियम~$R$
चा निष्कर्ष सिद्धता~$\pi$ ने सिद्ध होत असेल,
तर त्या नियमाची प्रत्येक आधार-क्रमवर्ती
सिद्धता~$\pi'$ ने (किंवा नियमाला दोन आधार
असल्यास सिद्धता~$\pi'$, $\pi''$ यांनी)
सिद्ध करता येते. शिवाय,
(a)~$\pheight{\pi'}\le\pheight{\pi}$ (आणि
$\pheight{\pi''}\le\pheight{\pi}$);
तसेच (b)~प्रत्येक कोटीसाठी $\pi'$ मधील
(आणि स्वतंत्रपणे $\pi''$ मधील) त्या कोटीच्या
\CutCS{} अनुमानांची संख्या $\pi$ मधील संख्येपेक्षा
मोठी नसते.
\end{lem}

\begin{proof}
पूर्वप्रमेय~\ref{pt:seq:inv:lem:G3c-invert} आणि~\ref{pt:seq:inv:lem:invert-quant}
यांच्या सिद्धता तपासल्यास हा दावा मिळतो. ती
सिद्धता~$\pheight{\pi}$ वरील विगमनाने होती.
आधार प्रसंगात एक स्वयंसिद्धक दुसऱ्या
स्वयंसिद्धकाने बदलले होते; म्हणून (a) आणि (b)
लगेच मिळतात. $\pi$ चा शेवटचा नियम~$R$
असेल, तर त्याच्या तत्काळ उप-सिद्धता या
आपल्याला हव्या असलेल्या सिद्धता~$\pi_i$
आहेत; त्यांची उंची लहान असते आणि प्रत्येक कोटीच्या
कटांची संख्या संपूर्ण $\pi$ मधील संख्येपेक्षा मोठी नसते.
दोन आधारांच्या नियमात एक आधार निवडल्यास दुसऱ्या
शाखेतील कट राहतीलच असे नाही; (b) साठी समानता नको,
ही वरची सीमाच आवश्यक आहे.
इतर प्रसंगांत $\pi$ च्या तत्काळ उप-सिद्धता वर,
म्हणजे शेवटच्या अनुमान~$R$ च्या आधारांकडे
नेणाऱ्या सिद्धतांवर, विगमन गृहीतक लावून
परिणामांवर तोच नियम $R$ लावला होता.
नव्या सिद्धतेच्या तत्काळ उप-सिद्धता साठी
(a) व (b) खऱ्या असल्याने संपूर्ण
सिद्धतेसाठीही त्या खऱ्या आहेत. शेवटचे
अनुमान~\CutCS{} असेल तो प्रसंग तपासणे
उरते. पुन्हा फक्त एक उदाहरण पाहू:
$\pi$ ही \LeftR{\land} च्या
$B \land C, \Gamma \Sequent \Delta$
या निष्कर्षाची सिद्धता असून
\CutCS{} ने संपते असे समजा:
\[
\AxiomC{}
\RightLabel{$\pi_1$}
\Deduce$B \land C, \Gamma \fCenter \Delta, A$
\AxiomC{}
\RightLabel{$\pi_2$}
\Deduce$A, B \land C, \Gamma \fCenter \Delta$
\RightLabel{\CutCS}
\BinaryInf$B \land C, \Gamma \fCenter \Delta$
\DisplayProof
\]
विगमन गृहीतक $\pi_1$ आणि~$\pi_2$
यांना लागू होते; त्या दोन्ही
$\LeftR{\land}$ च्या निष्कर्षाच्या
उदाहरणांच्या सिद्धता आहेत.
त्यातून अनुक्रमे $B, C, \Gamma \fCenter
\Delta, A$ आणि $A, B, C, \Gamma \fCenter \Delta$
यांच्या सिद्धता $\pi_1'$ आणि~$\pi_2'$
मिळतात. \CutCS{} वापरून त्यांना जोडल्यास
पुढील~$\pi'$ मिळते:
\[
\AxiomC{}
\RightLabel{$\pi_1'$}
\Deduce$B, C, \Gamma \fCenter \Delta, A$
\AxiomC{}
\RightLabel{$\pi_2'$}
\Deduce$A, B, C, \Gamma \fCenter \Delta$
\RightLabel{\CutCS}
\BinaryInf$B, C, \Gamma \fCenter \Delta$
\DisplayProof
\]
येथे (a) आणि (b) स्पष्टपणे पूर्ण होतात.
\end{proof}

\CutCS{} ऐवजी \Cut{} वापरला असता ही
पद्धत चालली नसती: शेवटच्या सिद्धता
च्या निष्कर्षातील पूर्वांगात
$B, C, \Gamma$ च्या दोन प्रती आणि
उत्तरांगात $\Delta$ च्या दोन प्रती आल्या
असत्या. मग आकुंचनाची ग्राह्यता वापरावी
लागली असती; पण
\ref{pt:seq:inv:prop:G3c-cont-adm}
ची सिद्धता या उलटवण्याच्या उपप्रमेयावरच
आधारित आहे. त्यामुळे युक्तिवाद चक्रीय ठरला असता.
या उलटवण्याच्या उपप्रमेयानंतर दुर्बलीकरण आणि आकुंचनाची
ग्राह्यताही $\Log{G3c}+\CutCS$ पर्यंत वाढवता येते;
त्यामुळे पुढील सिद्धतेत कमी कोटीचे कट असतानाही आकुंचन
वापरता येते. आधीच्या उंचीवरील विगमनात एवढाच नवा प्रसंग
आहे: अंतिम नियम $\CutCS$ असेल, तर दोन्ही आधारांतील
सामायिक संदर्भावर विगमन गृहीतक लावून तोच कट पुन्हा
जोडायचा. कट-सूत्र आणि त्याची कोटी बदलत नाहीत. इतर
नियमांसाठी आधीची सिद्धता वापरायची; आकुंचनाच्या मुख्य
प्रसंगांत आता वरील विस्तारित उलटवणे वापरायचे. त्यातून
उंची आणि कटांच्या सर्वोच्च कोटीत वाढ होत नाही.
संख्यक नियमांच्या आयगेन चरांची नावे आवश्यकतेनुसार आधी
बदलून नव्या संदर्भाच्या चरांपासून वेगळी ठेवायची.

पूर्वप्रमेय~\ref{pt:cut:inv:lem:inv-G3c-cut} वापरून
कट-निर्मूलनाची पर्यायी सिद्धता देता येते.
या सिद्धतेत आपण सर्वोच्च स्थानी असलेली
\CutCS{} अनुमाने एकामागोमाग काढत नाही;
त्याऐवजी सर्वोच्च कोटीची \CutCS{}
अनुमाने काढतो. म्हणजे
$\Log{G3c} + \CutCS$ मधील सिद्धता~$\pi$
घेऊन तिच्यातील कोणताही कट काढतो
(कदाचित त्याजागी कमी कोटीचे कट येतील)
आणि एकही कट उरेपर्यंत असेच करत राहतो.
अट एवढीच की निवडलेल्या कटच्या वर
त्याच किंवा त्याहून मोठ्या कोटीचे
कट नसावेत.

\begin{lem}\label{pt:cut:inv:lem:max-cut-red-G3c}
$\Log{G3c}+\CutCS$ मधील सिद्धता~$\pi$
ही $n$ कोटीच्या \CutCS{} अनुमानाने
संपत असेल आणि इतरत्र फक्त
$\Log{G3c}$ चे नियम व $<n$ कोटीची
\CutCS{} अनुमाने वापरत असेल, तर
$\Log{G3c} + \CutCS$ मध्ये
त्याच अंतिम क्रमवर्तीची अशी सिद्धता~$\pi'$
असते की तिच्यातील प्रत्येक \CutCS{}
अनुमानाची कोटी~$<n$ असते.
\end{lem}

\begin{proof}
सिद्धता~$\pi$ चा शेवट पुढीलप्रमाणे असतो:
\[
\AxiomC{}
\RightLabel{$\pi_1$}
\Deduce$\Gamma \fCenter \Delta, A$
\AxiomC{}
\RightLabel{$\pi_2$}
\Deduce$A, \Gamma \fCenter \Delta$
\RightLabel{\CutCS}
\BinaryInf$\Gamma \fCenter \Delta$
\DisplayProof
\]
आता~$A$ च्या रूपानुसार प्रसंग वेगळे करू.

$A$ हे अणुसूत्र आहे असे समजा. $\pi_1$
आणि $\pi_2$ मधील सर्व \CutCS{} अनुमानांची कोटी
$< \depth{A} = 0$ असली पाहिजे; म्हणून त्यांपैकी
कोणत्याही सिद्धतेत \CutCS{} अनुमान असू शकत नाही.
म्हणून $\pi_1$ आणि $\pi_2$ दोन्ही कटविरहित आहेत.
$\pheight{\pi_2}$ वरील विगमनाने
$\Gamma \Sequent \Delta$ ची कटविरहित सिद्धता
आहे हे दाखवू. $\pheight{\pi_2} = 0$ असेल,
तर $A, \Gamma \Sequent \Delta$ हे स्वयंसिद्धक आहे.
$\lfalse$ हे $\Gamma$ मध्ये असेल, तर निष्कर्ष स्वतःच
असत्य-स्वयंसिद्धक आहे. अन्यथा $A$ मुख्य नसेल, तर
काही अणुसूत्र $C$ हे $\Gamma$ आणि $\Delta$ या दोन्हींत
घटक असते; म्हणून
$\Gamma \Sequent \Delta$ हे स्वतःच स्वयंसिद्धक
आहे. $A\ident\lfalse$ असल्याचा प्रसंग
\ref{pt:cut:top:lem:cut-adm-G3c} च्या सिद्धतेतील प्रसंग B प्रमाणे
हाताळता येतो: कटविरहित $\pi_1$ च्या उत्तरांगातील कटासाठीची
असत्याची प्रत काढल्यास हवी ती सिद्धता मिळते.
उरलेल्या मुख्य अणुसूत्राच्या प्रसंगात
$\Delta = \Delta', A$. या प्रसंगात $\pi_1$
चा शेवट $\Gamma \Sequent \Delta', A, A$
असा असतो. \ref{pt:seq:inv:prop:G3c-cont-adm},
म्हणजे आकुंचनाची ग्राह्यता, लावल्यास
$\Gamma \Sequent \Delta', A$ ची
कटविरहित सिद्धता मिळते.
$\pheight{\pi_2} > 0$ असेल, तर
सिद्धता नियम~$R$ ने संपते. त्या नियमाची
एक आधार-क्रमवर्ती घ्या: तिचे रूप
$A, \Gamma', \Pi \Sequent \Delta', \Lambda$
असते, जेथे $\Pi$ व $\Lambda$ ही
$R$ ची सक्रिय सूत्रे आहेत.
\ref{pt:seq:adm:prop:weak-G3c-adm}
हे $\pi_1$ आणि निवडलेल्या आधारावर संपणाऱ्या
$\pi_2$ च्या तत्काळ उपसिद्धतेला लावल्यास
$\Gamma, \Gamma', \Pi \Sequent
\Delta, \Delta', \Lambda, A$ आणि
$A, \Gamma, \Gamma', \Pi \Sequent
\Delta, \Delta', \Lambda$ यांच्या
सिद्धता मिळतात. दुसरी सिद्धता $\pi_2$ पेक्षा
कमी उंचीची राहते; म्हणून त्यांना विगमन गृहीतक लागू होते. त्यामुळे~$R$ च्या प्रत्येक
आधारासाठी $\Gamma, \Gamma', \Pi \Sequent
\Delta, \Delta', \Lambda$ ची सिद्धता
मिळते. $R$ लावल्यास
$\Gamma, \Gamma \Sequent \Delta, \Delta$
मिळते; आणि \ref{pt:seq:inv:prop:G3c-cont-adm}
मुळे $\Gamma \Sequent \Delta$ ची
कटविरहित सिद्धता मिळते.


$A$ अणुसूत्र नसेल, तर त्याच्या मुख्य संकारक नुसार प्रसंग पाहू. येथे $A$ हाच
कट-सूत्र आहे. प्रत्येक प्रसंगात
\ref{pt:cut:inv:lem:inv-G3c-cut} वापरून $A$ मुख्य सूत्र
असलेल्या डाव्या व उजव्या नियमांच्या आधारांच्या
सिद्धता मिळवू. त्यानंतर
\ref{pt:cut:top:lem:cut-adm-G3c} च्या सिद्धतेतील
प्रसंग~E प्रमाणे पुढे जाऊ.
\begin{enumerate}
  \item $A \ident \lnot B$. सिद्धता $\pi_1$
  आणि $\pi_2$ यांचे शेवट अनुक्रमे
  $\Gamma \Sequent \Delta, \lnot B$ आणि
  $\lnot B, \Gamma \Sequent \Delta$ असे आहेत.
  \ref{pt:cut:inv:lem:inv-G3c-cut} नुसार
  $B, \Gamma \Sequent \Delta$ आणि
  $\Gamma \Sequent \Delta, B$ यांच्या
  सिद्धता $\pi_1'$ आणि $\pi_2'$ मिळतात.
  सिद्धता~$\pi'$ पुढीलप्रमाणे आहे:
  \[
  \AxiomC{}
  \RightLabel{$\pi_2'$}
  \Deduce$\Gamma \fCenter \Delta, B$
  \AxiomC{}
  \RightLabel{$\pi_1'$}
  \Deduce$B, \Gamma \fCenter \Delta$
  \RightLabel{\CutCS}
  \BinaryInf$\Gamma \fCenter \Delta$
  \DisplayProof
  \]
  \item $A \ident (B \lor C)$.
  सिद्धता $\pi_1$ आणि $\pi_2$
  यांचे शेवट
  $\Gamma \Sequent \Delta, B \lor C$ आणि
  $B \lor C, \Gamma \Sequent \Delta$ असे आहेत.
  $\pi_1$ ला \RightR{\lor} चे उलटवणे
  लावून \ref{pt:cut:inv:lem:inv-G3c-cut} नुसार
  $\Gamma \Sequent \Delta, B, C$
  ची सिद्धता $\pi_1'$ मिळते.
  $\pi_2$ ला \LeftR{\lor} चे उलटवणे
  लावून \ref{pt:cut:inv:lem:inv-G3c-cut} नुसार
  $B, \Gamma \Sequent \Delta$ ची
  सिद्धता $\pi_2'$ आणि
  $C, \Gamma \Sequent \Delta$ ची
  सिद्धता $\pi_2''$ मिळते.
  \ref{pt:seq:adm:prop:weak-G3c-adm}
  हे $\pi_2''$ ला लावल्यास
  $C, \Gamma \Sequent \Delta, B$
  ची सिद्धता~$\pi_2'''$ मिळते.
  सिद्धता~$\pi'$ पुढीलप्रमाणे आहे:
  \[
  \AxiomC{}
  \RightLabel{$\pi_1'$}
  \Deduce$\Gamma \fCenter \Delta, B, C$
  \AxiomC{}
  \RightLabel{$\pi_2'''$}
  \Deduce$C, \Gamma \fCenter \Delta, B$
  \RightLabel{\CutCS}
  \BinaryInf$\Gamma \fCenter \Delta, B$
  \AxiomC{}
  \RightLabel{$\pi_2'$}
  \Deduce$B, \Gamma \fCenter \Delta$
  \RightLabel{\CutCS}
  \BinaryInf$\Gamma \fCenter \Delta$
  \DisplayProof
  \]
  \item $A \ident (B \land C)$.
  हा प्रसंग सरावासाठी आहे.
  \item $A \ident (B \lif C)$.
  सिद्धता $\pi_1$ आणि $\pi_2$
  यांचे शेवट
  $\Gamma \Sequent \Delta, B \lif C$ आणि
  $B \lif C, \Gamma \Sequent \Delta$ असे आहेत.
  $\pi_1$ ला \RightR{\lif} चे उलटवणे
  लावून \ref{pt:cut:inv:lem:inv-G3c-cut} नुसार
  $B, \Gamma \Sequent \Delta, C$
  ची सिद्धता $\pi_1'$ मिळते.
  $\pi_2$ ला \LeftR{\lif} चे उलटवणे
  लावून \ref{pt:cut:inv:lem:inv-G3c-cut} नुसार
  $\Gamma \Sequent \Delta, B$ ची
  सिद्धता $\pi_2'$ आणि
  $C, \Gamma \Sequent \Delta$ ची
  सिद्धता $\pi_2''$ मिळते.
  \ref{pt:seq:adm:prop:weak-G3c-adm}
  हे $\pi_2''$ ला लावल्यास
  $C, B, \Gamma \Sequent \Delta$
  ची सिद्धता~$\pi_2'''$ मिळते.
  सिद्धता~$\pi'$ पुढीलप्रमाणे आहे:
  \[
  \AxiomC{}
  \RightLabel{$\pi_2'$}
  \Deduce$\Gamma \fCenter \Delta, B$
  \AxiomC{}
  \RightLabel{$\pi_1'$}
  \Deduce$B, \Gamma \fCenter \Delta, C$
  \AxiomC{}
  \RightLabel{$\pi_2'''$}
  \Deduce$C, B, \Gamma \fCenter \Delta$
  \RightLabel{\CutCS}
  \BinaryInf$B, \Gamma \fCenter \Delta$
  \RightLabel{\CutCS}
  \BinaryInf$\Gamma \fCenter \Delta$
  \DisplayProof
  \]
  \item $A \ident \lforall[x][B(x)]$.
  सिद्धता $\pi_1$ आणि $\pi_2$
  यांचे शेवट
  $\Gamma \Sequent \Delta, \lforall[x][B(x)]$
  आणि $\lforall[x][B(x)], \Gamma \Sequent \Delta$
  असे आहेत. \ref{pt:cut:inv:lem:inv-G3c-cut} नुसार
  $\Gamma \Sequent \Delta, B(c)$ ची
  सिद्धता~$\pi_1'$ मिळते.

आता सिद्धता~$\pi_2$ पाहू. तिचा शेवट
$\lforall[x][B(x)], \Gamma \Sequent \Delta$
असा आहे. $\pi_2$ च्या अंतिम क्रमवर्तीपासून
वर जाताना, डाव्या बाजूला आलेली
$\lforall[x][B(x)]$ ची अशी प्रत्येक आस्थिती
चिन्हांकित करा की ती अंतिम क्रमवर्तीतील
$\lforall[x][B(x)]$ च्या आस्थितीपर्यंत
“पोहोचते”. म्हणजे:
(a)~$\pi_2$ च्या अंतिम क्रमवर्तीतील
$\lforall[x][B(x)]$ चिन्हांकित करा.
(b)~एखाद्या नियमाच्या निष्कर्षाच्या
संदर्भात $\lforall[x][B(x)]$ चिन्हांकित
असेल, तर त्या नियमाच्या आधारात किंवा
आधारांत तिच्याशी अनुरूप आस्थिती चिन्हांकित
करा. (c)~\LeftR{\lforall} नियमाच्या
निष्कर्षात $\lforall[x][B(x)]$ हे मुख्य
सूत्र असून चिन्हांकित असेल, तर त्याच्या
आधारातील $\lforall[x][B(x)]$ आणि
$B(t)$ यांच्या अनुरूप सक्रिय आस्थिती
चिन्हांकित करा. हा मागोवा $\pi_2$
च्या सिद्धतावृक्षावरच घेतला आहे.

आता $\lforall[x][B(x)]$ च्या
या सर्व चिन्हांकित आस्थिती पाहू. त्यांपैकी
कोणतीही \LeftR{\lforall} व्यतिरिक्त
अन्य नियमात सक्रिय नसते
(फक्त \LeftR{\lforall} ची सक्रिय
सूत्रे चिन्हांकित होतात).
$\pi_2$ मधील $\lforall[x][B(x)]$
च्या सर्व चिन्हांकित आस्थिती काढून टाका.
यातून योग्य सिद्धता मिळाली, तर तिचा शेवट
$\Gamma \Sequent \Delta$ असा आहे आणि
ती $\pi$ च्या जागी ठेवता येते.
याने सर्वोच्च कोटीचा एक कट काढला.

अन्यथा मिळालेली सिद्धता योग्य नाही:
काही \LeftR{\lforall} अनुमानांत सक्रिय
असलेली $\lforall[x][B(x)]$ ही
सूत्रे आपण काढून टाकली आहेत.
त्यामुळे पुढील रूपाची “अनुमाने” मिळाली:
\[
  \AxiomC{}
  \RightLabel{$\pi_t$}
  \Deduce$B(t), \Pi \fCenter \Lambda$
  \RightLabel{\LeftR{\lforall}}
  \UnaryInf$\Pi \fCenter \Lambda$
  \DisplayProof
  \]
अशा प्रत्येक सर्वोच्च “अनुमाना”च्या जागी
पुढील सिद्धता ठेवा:
\[
  \AxiomC{}
  \RightLabel{$\Subst{\pi_1'}{t}{c}[\Pi \Sequent \Lambda]$}
  \Deduce$\Gamma, \Pi \fCenter \Delta, \Lambda, B(t)$
  \AxiomC{}
  \RightLabel{$\pi_t[\Gamma \Sequent \Delta]$}
  \Deduce$B(t), \Gamma, \Pi \fCenter \Delta, \Lambda$
  \RightLabel{\CutCS}
  \BinaryInf$\Gamma, \Pi \fCenter \Delta, \Lambda$
  \DisplayProof
  \]
तिच्या खालील प्रत्येक क्रमवर्तीच्या
डावीकडे $\Gamma$ आणि उजवीकडे
$\Delta$ घाला. आधारात चिन्हांकित
$B(t)$ असलेली सर्व
\LeftR{\lforall} “अनुमाने”
काही~$B(t)$ वरील \CutCS{}
अनुमानाने बदलेपर्यंत हे पुन्हा करा.
मिळालेल्या सिद्धता ची अंतिम क्रमवर्ती
(आता ती खरोखर योग्य सिद्धता आहे)
$\Gamma, \dots, \Gamma \Sequent
\Delta, \dots, \Delta$ या रूपाची असते.
आकुंचनाची ग्राह्यता लावून तिचे
$\Gamma \Sequent \Delta$ वरील
सिद्धता मध्ये रूपांतर करा आणि
ती $\pi$ च्या जागी ठेवा. आपण
$\lforall[x][B(x)]$ वरील एक कट
$B(t)$ वरील (कदाचित अनेक) कटांनी
बदलला आहे; पण या सर्व कटांची कोटी
कमी आहे. $\pi_1$ किंवा $\pi_2$ मधील उरलेल्या किंवा प्रत करून
वापरलेल्या कटांची कोटीही कमीच आहे. पुढील आकुंचनासाठी
वर सिद्ध केलेली, सर्वोच्च कट-कोटी न वाढवणारी विस्तारित
ग्राह्यता वापरली आहे.
\end{enumerate}
\end{proof}

\begin{prob}
\ref{pt:cut:inv:lem:max-cut-red-G3c} च्या सिद्धतेतील
$A \ident (B \land C)$ हा प्रसंग तपासा.
म्हणजे, दोन उपसिद्धतांसह पुढीलप्रमाणे संपणारी
सिद्धता~$\pi$ असेल:
\[
\AxiomC{}
\RightLabel{$\pi_1$}
\Deduce$\Gamma \fCenter \Delta, B \land C$
\AxiomC{}
\RightLabel{$\pi_2$}
\Deduce$B \land C, \Gamma \fCenter \Delta$
\RightLabel{\CutCS}
\BinaryInf$\Gamma \fCenter \Delta$
\DisplayProof
\]
आणि $\pi_1$ व $\pi_2$ यांतील
सर्व कटांची कोटी $<\depth{B
\land C}$ असेल, तर
$\Gamma \Sequent \Delta$ ची
फक्त $<\depth{B \land C}$
कोटीचे कट असलेली सिद्धता~$\pi'$
असते हे दाखवा.
\end{prob}






\ref{pt:cut:inv:lem:max-cut-red-G3c} नुसार
सिद्धता मधील सर्वोच्च कोटीचे
\CutCS{} अनुमान कमी कोटीच्या
\CutCS{} अनुमानांनी बदलता येते.
हे उपप्रमेय पुन्हा वापरून,
सर्वोच्च~\CutCS{} ने संपणाऱ्या
सर्वोच्च उप-सिद्धता वर ते एकामागोमाग
लावल्यास कोणत्याही सिद्धता मधून कितीही
\CutCS{} अनुमाने काढता येतात.


\begin{thm}\label{pt:cut:inv:thm:cut-elim-G3c-largest}
$\Log{G3c} + \CutCS$ मधील कोणतीही
सिद्धता $\pi$ ही~\Log{G3c} मधील
सिद्धतेत रूपांतरित करता येते.
\end{thm}

\begin{proof}
  प्रथम $\pi$ मधील सर्वोच्च कोटीचे
  सर्व कट काढता येतात हे सिद्ध करू.
  $d$ ही $\pi$ मधील कटांची सर्वोच्च
  कोटी आणि $n$ ही $d$ कोटीच्या कटांची
  संख्या असू द्या. $n$ वरील विगमनाने
  सिद्धता करू. $n=0$ असेल, तर
  सिद्ध करायचे काही नाही. $n>0$
  असेल, तर $d$ कोटीचा सर्वोच्च कट
  निवडा आणि त्यावर संपणारी उप-सिद्धता
  $\pi_1$ असू द्या.
  \ref{pt:cut:inv:lem:max-cut-red-G3c} नुसार
  त्याच अंतिम क्रमवर्तीची, फक्त
  $<d$ कोटीचे कट असलेली,
  सिद्धता~$\pi_1'$ आहे.
  $\pi$ मध्ये $\pi_1$ च्या जागी
  $\pi_1'$ ठेवा. मिळालेल्या
  सिद्धता मध्ये $d$ कोटीचे
  $n-1$ कट असतात; म्हणून
  विगमन गृहीतक लागू होते.

आता $d$ वरील विगमनाने निष्कर्ष
मिळतो. $d=0$ असेल, तर सर्व कट
अणुसूत्रांवरील आहेत आणि मागील
निष्कर्षाने ते सर्व काढता येतात.
$d>0$ असेल, तर मागील निष्कर्षाने
$d$ कोटीचे सर्व कट $<d$ कोटीच्या
कटांनी बदललेली सिद्धता मिळते.
$\pi$ मधील कटांची सर्वोच्च कोटी
$d$ असल्यामुळे नव्या सिद्धतेतील
कटांची सर्वोच्च कोटी $<d$ होते;
म्हणून विगमन गृहीतक लागू होते.
\end{proof}













\section{मध्य-क्रमवर्ती प्रमेय आणि एरब्राँचे प्रमेय}\label{pt:cut:mid:sec}

कट-निर्मूलन प्रमेयाचा एक महत्त्वाचा
परिणाम म्हणजे मध्य-क्रमवर्ती प्रमेय.
त्यात असे म्हटले आहे: सर्व सूत्र
अग्र-संख्यापक रूपातील असलेल्या
$\Gamma \Sequent \Delta$ ची
सिद्धता~$\pi$ असेल, तर अशी
सिद्धता~$\pi'$ असते की तिच्यात
$S \ident \Pi \Sequent
\Lambda$ ही क्रमवर्ती येते
(तिला \emph{मध्य-क्रमवर्ती} म्हणतात)
आणि~$\pi'$ मध्ये $S$ च्या
वरची सर्व अनुमाने विधानीय,
तर $S$ च्या खालची सर्व अनुमाने
संख्यापकांची असतात.
(सूत्र~$A$ हे
\emph{अग्र-संख्यापक} (किंवा
\emph{अग्र-संख्यापक रूपातील}) असते,
जर ते पुढील रूपात असेल:
\[\mathsf{Q}_1x_1\dots\mathsf{Q}_nx_n\,B(x_1,\dots,x_n)\]
येथे प्रत्येक~$\mathsf{Q}_i$
हा $\lforall$ किंवा~$\lexists$
असतो आणि $B$ संख्यापकविरहित असते.) मध्य-क्रमवर्ती प्रमेयाचा
एक महत्त्वाचा परिणाम म्हणजे
एरब्राँचे प्रमेय. त्याचे व्यापक
रूप थोडे कठीण आहे; साधारणतः
ते असे सांगते: प्रथम-क्रम
तर्कशास्त्रातील प्रत्येक
सिद्धतायोग्य वाक्य~$A$
साठी असे विधानीय
सर्वतः सत्य सूत्र~$A'$
असते की त्यावरून $A$
सिद्ध करता येते.
$B$ संख्यापकविरहित
असताना $\lexists[x][B(x)]$
या रूपातील वाक्ये साठी
पुढील विधान आहे:

\begin{thm}[एरब्राँचे प्रमेय]
$B(x)$ संख्यापकविरहित असू द्या
आणि $\Proves
\lexists[x][B(x)]$ असे समजा.
मग अशी पदे~$t_1$, \dots, $t_n$
असतात की
$\Proves B(t_1) \lor \dots \lor B(t_n)$.
\end{thm}

$B(t_1) \lor \dots \lor B(t_n)$
या सूत्राला
$\lexists[x][B(x)]$ चा
\emph{एरब्राँचा विकल्पयोग} म्हणतात.
$B$ संख्यापकविरहित असल्याने
हा विकल्पयोगही संख्यापकविरहित आहे.
म्हणून तो सिद्धतायोग्य असेल,
तर कटविरहित सिद्धतेने तो
सिद्धतायोग्य आहे; परिणामी,
केवळ विधानीय अनुमाने
वापरणारी सिद्धता आहे.
म्हणून तो सर्वतः सत्य आहे.

एरब्राँच्या प्रमेयातील कठीण
भाग म्हणजे प्रत्येक
सिद्धतायोग्य
सूत्र~$\lexists[x][B(x)]$
साठी एरब्राँचा विकल्पयोग
अस्तित्वात आहे हे दाखवणे.
उलट दिशा सोपी आहे:
$\lexists[x][B(x)]$ चा
एरब्राँचा विकल्पयोग असेल,
तर हे सूत्र सिद्धतायोग्य आहे.
कारण एरब्राँचा विकल्पयोग
$\lexists[x][B(x)]$
ला निष्पन्न करतो. उदाहरणार्थ,
$n=2$ असताना~$\Log{G3c}$
मधील ही सिद्धता पाहा:
\[
\Axiom$B(t_1) \fCenter \lexists[x][B(x)], B(t_1), B(t_2)$
\Axiom$B(t_2) \fCenter \lexists[x][B(x)], B(t_1), B(t_2)$
\RightLabel{\LeftR{\lor}}
\BinaryInf$B(t_1) \lor B(t_2) \fCenter \lexists[x][B(x)], B(t_1), B(t_2)$
\RightLabel{\RightR{\lexists}}
\UnaryInf$B(t_1) \lor B(t_2) \fCenter \lexists[x][B(x)], B(t_2)$
\RightLabel{\RightR{\lexists}}
\UnaryInf$B(t_1) \lor B(t_2) \fCenter \lexists[x][B(x)]$
\DisplayProof
\]

$\Sequent
\lexists[x][B(x)]$ ची
सिद्धता~$\pi$ दिली असता
त्यातून एरब्राँचा विकल्पयोग
मिळवण्यासाठी, कटविरहित
सिद्धता~$\pi$ मधील सर्व
$\RightR{\lexists}$
अनुमाने शेवटी सलग येतात
हा प्रसंग पाहा:
\[
\AxiomC{}
\RightLabel{$\pi_1$}
\Deduce$ \fCenter \lexists[x][B(x)], B(t_1), \dots, B(t_n)$
\RightLabel{\RightR{\lexists}}
\UnaryInf$ \fCenter \lexists[x][B(x)], B(t_2), \dots, B(t_n)$
\RightLabel{$(n-1) \times \RightR{\lexists}$}
\Deduce$ \fCenter \lexists[x][B(x)]$
\DisplayProof
\]
~$\pi$ मध्ये
$\lexists[x][B(x)]$
सक्रिय असलेली हीच सर्व
$\RightR{\lexists}$ अनुमाने
असतील, तर~$\pi_1$
मध्ये इतर संख्यापक अनुमाने
नसतात. कारण $B(x)$
मध्ये अन्य संख्यापक नाहीत
आणि अंतिम क्रमवर्तीच्या
डावीकडे सूत्रे नाहीत.
म्हणजे $ \fCenter \lexists[x][B(x)], B(t_1),
\dots, B(t_n)$ ही~$\pi$
ची मध्य-क्रमवर्ती आहे.
शिवाय मध्य-क्रमवर्तीच्या
वर संख्यापक अनुमाने
नसल्यामुळे तिच्या वरच्या
सर्व क्रमवर्तींत
$\lexists[x][B(x)]$
असते, पण ते सक्रिय किंवा
मुख्य नसते. म्हणून
~$\pi_1$ या सिद्धता
मधून ते काढून टाकता येते
आणि $\fCenter B(t_1), \dots, B(t_n)$
याची सिद्धता मिळते.
मग $\RightR{\lor}$ चा
$n-1$ वेळा उपयोग करून
एरब्राँच्या विकल्पयोगाची
सिद्धता,
$\Sequent B(t_1) \lor \dots \lor B(t_n)$,
मिळते.

अडचण अशी की
$\Sequent \lexists[x][B(x)]$
ची कटविरहित सिद्धता
असली तरी तिची सर्व
\RightR{\lexists} अनुमाने
शेवटी एका सलग गटात
येतात असे गृहीत धरता येत
नाही. त्या \RightR{\lexists}
अनुमानांमध्ये एखादे
$B(t_i)$ मुख्य असलेली
अनुमाने येऊ शकतात.
म्हणून मध्य-क्रमवर्ती
प्रमेयाची गरज आहे.

\begin{thm}[मध्य-क्रमवर्ती प्रमेय]\label{pt:cut:mid:thm:midsequent}
$\Gamma \Sequent \Delta$ ची
~\Log{G3c} मधील
सिद्धता~$\pi$ असेल आणि
$\Gamma$ व $\Delta$ मधील
सर्व सूत्रे
अग्र-संख्यापक असतील,
तर अशी सिद्धता~$\pi'$
असते की तिच्यात
$\Pi \Sequent \Lambda$
ही क्रमवर्ती येते;
तिच्या खालची सर्व
अनुमाने संख्यापकांची,
आणि वरची सर्व अनुमाने
विधानीय असतात.
\end{thm}

\begin{proof}
  ~$\pi$ मधील प्रत्येक संख्यापक
  अनुमानाचे \emph{क्रममान} म्हणजे
  त्याच्या खालच्या विधानीय
  अनुमानांची संख्या असे
  परिभाषित करू; आणि
  $\pi$ चे क्रममान~$o(\pi)$
  म्हणजे त्यातील सर्व
  संख्यापक अनुमानांच्या
  क्रममानांची बेरीज.
  $o(\pi) = 0$ असेल, तर
  कोणत्याही संख्यापक
  अनुमानाच्या खाली विधानीय
  अनुमान नसते आणि काम संपते:
  संख्यापक अनुमानच नसल्यास
  अंतिम क्रमवर्ती मध्य-क्रमवर्ती
  घेता येते. अन्यथा सर्व
  संख्यापक अनुमानांना एकच
  आधारक्रमवर्ती असल्याने
  एक सर्वोच्च संख्यापक
  अनुमान असते; त्याची
  आधारक्रमवर्ती मध्य-क्रमवर्ती असते.

  $o(\pi) > 0$ असेल, तर
  क्रममान $>0$ असलेले सर्वात
  खालचे संख्यापक अनुमान
  निवडा. त्याचे क्रममान $>0$
  असल्याने त्याच्या खाली
  विधानीय अनुमान असलेच पाहिजे.
  ते अशा अनुमानांपैकी
  सर्वात खालचे असल्याने
  त्याचा निष्कर्ष दुसऱ्या
  संख्यापक अनुमानाची
  आधारक्रमवर्ती असू शकत
  नाही: तसे असेल, तर
  त्या दुसऱ्या, अधिक
  खालच्या संख्यापक अनुमानाच्या
  खालीही विधानीय अनुमान
  आले असते. कोणते
  संख्यापक अनुमान आणि
  त्याखाली कोणते विधानीय
  अनुमान आहे, यानुसार
  (अनेक!) प्रसंग ठरतात.
  अंतिम क्रमवर्तीमध्ये
  केवळ अग्र-संख्यापक
  सूत्रे असल्यामुळे,
  संख्यापक अनुमानाचे
  मुख्य सूत्र पुढील
  विधानीय अनुमानात
  सक्रिय असू शकत नाही.
  अन्यथा त्या अनुमानाच्या
  मुख्य सूत्रात विधानीय
  संयोजकाच्या व्याप्तीत
  संख्यापक आला असता.
  उप-सूत्र गुणधर्मामुळे
  ते अंतिम क्रमवर्तीतील
  एखाद्या सूत्राचे
  उप-सूत्र असावे लागले
  असते. म्हणून संख्यापक
  अनुमानाचे मुख्य
  सूत्र हे खालच्या
  विधानीय अनुमानाच्या
  संदर्भातील घटक आहे.

  पुढील दोन उदाहरणे पाहू:

प्रथम $\RightR{\lforall}$
हे $\LeftR{\lif}$ च्या
उजव्या आधारक्रमवर्तीमध्ये
संपते असे समजा. मग
~$\pi$ ची संबंधित
उप-सिद्धता
ही उप-सिद्धता~$\pi_1$
मध्ये संपते:
\[
\AxiomC{}
\RightLabel{$\pi_2$}
\Deduce$\Gamma' \fCenter \Delta', \lforall[x][A(x)], B$
\AxiomC{}
\RightLabel{$\pi_3$}
\Deduce$C, \Gamma' \fCenter \Delta', A(c)$
\RightLabel{\RightR{\lforall}}
\UnaryInf$C, \Gamma' \fCenter \Delta', \lforall[x][A(x)]$
\RightLabel{\LeftR{\lif}}
\BinaryInf$B \lif C, \Gamma' \fCenter \Delta', \lforall[x][A(x)]$
\DisplayProof
\]
~$\pi$ नियमित आहे असे
मानता येते; म्हणून
आयगेन चर~$c$ हे~$\pi_3$
च्या बाहेर येत नाही.
विशेषतः ते $B
\lif C$ मध्ये किंवा
~$\pi_2$ मध्ये येत नाही.
आता सिद्धता~$\pi_1'$ पाहा:
\[
\AxiomC{}
\RightLabel{$\pi_2'$}
\Deduce$\Gamma' \fCenter \Delta', \lforall[x][A(x)], A(c), B$
\AxiomC{}
\RightLabel{$\pi_3$}
\Deduce$C, \Gamma' \fCenter \Delta', \lforall[x][A(x)], A(c)$
\RightLabel{\LeftR{\lif}}
\BinaryInf$B \lif C, \Gamma' \fCenter \Delta', \lforall[x][A(x)], A(c)$
\RightLabel{\RightR{\lforall}}
\UnaryInf$B \lif C, \Gamma' \fCenter \Delta', \lforall[x][A(x)], \lforall[x][A(x)]$
\DisplayProof
\]
येथे~$\pi_2'$ ही
~$\pi_2$ च्या उजवीकडे
~$A(c)$ जोडून दुर्बलीकरणाने
मिळते. (त्यामुळे नवीन
अनुमाने, विशेषतः क्रममानावर
परिणाम करू शकणारी
संख्यापक अनुमाने, जोडली जात
नाहीत. ~$\pi_2$ मध्ये
आयगेन चरे म्हणून
वापरलेले स्थिरांक~$A(c)$
मध्ये नसल्याने~$\pi_2$
मधील आयगेन चरांच्या
अटीही मोडत नाहीत.)
दुसऱ्या शाखेला आवश्यक
संदर्भ जोडण्यासाठीही
दुर्बलीकरण वापरता येते.
\RightR{\lforall} ची
आयगेन चराची अट अजूनही
पूर्ण आहे, कारण $c$
हे~$B \lif C$
मध्ये येत नाही.

आकुंचनाची ग्राह्यता वापरून
$B \lif C, \Gamma' \fCenter \Delta',
\lforall[x][A(x)]$ याची
सिद्धता~$\pi_1''$
मिळवतो. $\lforall[x][A(x)]$
साठी $\RightR{\Contraction}$
ची ग्राह्यता सिद्ध करताना
आणि त्या सिद्धतेत वापरलेल्या
$\RightR{\lforall}$ च्या
व्युत्क्रमणीयतेच्या सिद्धतेत
अनुमानांची क्रममाने वाढत नाहीत:
अनुमाने काढली जाऊ शकतात, म्हणून संबंधित
उरलेल्या अनुमानांच्या खालील विधानीय
अनुमानांची संख्या कमीही होऊ शकते.
म्हणून~$\pi_1''$ मध्ये
~$\pi_1$ पेक्षा अधिक
संख्यापक अनुमाने नाहीत,
आणि~$\pi_1''$ मधील
प्रत्येक संख्यापक अनुमानाच्या
खाली त्याच्या~$\pi_1$
मधील संबंधित अनुमानापेक्षा अधिक
विधानीय अनुमाने नाहीत.
अपवाद फक्त अंतिम
\RightR{\lforall} चा:
त्याच्या खाली~$\pi$
मध्ये \LeftR{\lif} होते,
ते आपण~$\pi_1''$ मध्ये
त्याच्या वर हलवले आहे.
~$\pi$ मधील $\pi_1$
च्या जागी $\pi_1''$ ठेवा.
मिळालेल्या सिद्धता
चे क्रममान कमी आहे, कारण
किमान $\RightR{\lforall}$
अनुमानाचे क्रममान (किमान)~$1$
ने कमी झाले. आता
विगमन गृहीतक लावा.

संख्यापक अनुमान
\RightR{\lexists} असल्याचे
प्रसंग आणखी सोपे आहेत:
आकुंचन लागत नाही आणि
आयगेन चराच्या अटीची
काळजीही करावी लागत नाही.
उदाहरणार्थ,
$\RightR{\lexists}$
नंतर $\RightR{\land}$
येतो असे समजा
(यावेळी डाव्या आधारक्रमवर्तीमध्ये).
संबंधित उप-सिद्धता
ही $\pi_1$ आहे:
\[
\AxiomC{}
\RightLabel{$\pi_2$}
\Deduce$\Gamma' \fCenter \Delta', \lexists[x][A(x)], A(t), B$
\RightLabel{\RightR{\lexists}}
\UnaryInf$\Gamma' \fCenter \Delta', \lexists[x][A(x)], B$
\AxiomC{}
\RightLabel{$\pi_3$}
\Deduce$\Gamma' \fCenter \Delta', \lexists[x][A(x)], C$
\RightLabel{\RightR{\land}}
\BinaryInf$\Gamma' \fCenter \Delta', \lexists[x][A(x)], B \land C$
\DisplayProof
\]
~$\pi$ मधील $\pi_1$
च्या जागी $\pi_1'$ ठेवा:
\[
\AxiomC{}
\RightLabel{$\pi_2$}
\Deduce$\Gamma' \fCenter \Delta', \lexists[x][A(x)], A(t), B$
\AxiomC{}
\RightLabel{$\pi_3'$}
\Deduce$\Gamma' \fCenter \Delta', \lexists[x][A(x)], A(t), C$
\RightLabel{\RightR{\land}}
\BinaryInf$\Gamma' \fCenter \Delta', \lexists[x][A(x)], A(t), B \land C$
\RightLabel{\RightR{\lexists}}
\UnaryInf$\Gamma' \fCenter \Delta', \lexists[x][A(x)], B \land C$
\DisplayProof
\]
येथे $\pi_3'$ ही
$\pi_3$ च्या उजवीकडे
~$A(t)$ जोडून
दुर्बलीकरणाने मिळते.
या दुर्बलीकरणात
~$\pi_3$ मधील
प्रत्येक क्रमवर्तीच्या
उजवीकडे केवळ $A(t)$
जोडले जाते; म्हणून
कोणत्याही संख्यापक
अनुमानाचे क्रममान बदलत
नाही. ~$\pi'$ आणि
~$\pi$ मधील
संख्यापक अनुमानांच्या
क्रममाने सारखीच आहेत;
फक्त~\RightR{\land}
च्या वर हलवलेल्या
\RightR{\lexists}
अनुमानाचे क्रममान~$1$
ने कमी होते. म्हणून
विगमन गृहीतक लागू होते.
\end{proof}

\begin{prob}
\ref{pt:cut:mid:thm:midsequent} च्या
सिद्धतेतील पुढील प्रसंग
वर्णन करा: \LeftR{\lexists}
हे~\RightR{\lif} च्या
आधारक्रमवर्तीमध्ये संपते;
आणि \LeftR{\lforall}
हे~\LeftR{\lor} च्या
डाव्या आधारक्रमवर्तीमध्ये
संपते.
\end{prob}

\begin{proof}[एरब्राँच्या प्रमेयाची सिद्धता]
$\pi$ चा शेवट
$\Sequent \lexists[x][A(x)]$
मध्ये होतो असे समजा.
\ref{pt:cut:mid:thm:midsequent} नुसार
$\pi$ चे मध्य-क्रमवर्ती~$S$
असलेल्या सिद्धता~$\pi'$
मध्ये रूपांतर करता येते;
तिचे रूप पुढीलप्रमाणे असले
पाहिजे: \[\Sequent
\lexists[x][A(x)], A(t_1), \dots, A(t_n).\]
$S$ च्या वरची सर्व
अनुमाने विधानीय असल्याने,
$\lexists[x][A(x)]$
हे $S$ च्या वरच्या
अनुमानात मुख्य असू शकत नाही.
ते अणुरूप नसल्याने
स्वयंसिद्धकातही मुख्य
असू शकत नाही. शिवाय,
$A(x)$ संख्यापकविरहित
आहे हे लक्षात घेतल्यास
$\lexists[x][A(x)]$
हे $A(t_1)$ चे
योग्य उपसूत्र असू शकत नाही;
म्हणून $S$ च्या वर
ते सक्रियही नसते.
त्यामुळे $S$ वर संपणाऱ्या
उप-सिद्धता मधून
$\lexists[x][A(x)]$
काढल्यावर
$\Sequent A(t_1), \dots, A(t_n)$
याची योग्य सिद्धता
मिळते. त्यातून
$\RightR{\lor}$ ची
$n-1$ अनुमाने लावून
$\Sequent A(t_1) \lor \dots \lor A(t_n)$
याची सिद्धता मिळते.
\end{proof}






















\section{माएहाराचे उपप्रमेय आणि क्रेगचे अंतर्वेशन प्रमेय}\label{pt:cut:itp:sec}

विल्यम क्रेग यांचे अंतर्वेशन प्रमेय असे सांगते:
$A \Entails B$ असेल, तर असे सूत्र~$C$
असते (त्याला \emph{अंतर्वेशक} म्हणतात) की
$A \Entails C$ आणि $C \Entails B$.
विशेष म्हणजे, $C$ मधील सर्व स्थिरांक आणि
विधेये हे $A$ व $B$ या दोन्हींमध्ये येतील
असे $C$ निवडता येते. (येथे कोणतीही फलने
नाहीत असे गृहीत धरतो.) ही अट नसेल, तर
$A$ आणि $B$ हेच सहजपणे अंतर्वेशक ठरतील.

अंतर्वेशन प्रमेय अभिजात प्रथम-क्रम तर्कशास्त्राला लागू होते.
त्याचा शोध प्रथम-क्रम तर्कशास्त्राचा “शेवटचा महत्त्वाचा गुणधर्म”
शोधला गेल्याप्रमाणे वर्णिला गेला आहे.
प्रतिमान उपपत्ती आणि स्वयंचलित तर्कपद्धती यांत त्याचे
महत्त्वाचे उपयोग आहेत. त्यामागील मूळ प्रेरणा आणि
उपयोग विज्ञानाच्या तत्त्वज्ञानात होता: एखाद्या उपपत्तीचे
स्वयंसिद्धकीकरण कोणत्या मूलभूत संज्ञांनी करता येते
किंवा येत नाही, हे तपासण्यासाठी ते वापरता येते.
यावरून बेथचे परिभाष्यता प्रमेयही निष्पन्न होते:
एखादी उपपत्ती संकल्पना अंतर्निहितपणे परिभाषित करू
शकत असेल, तर ती स्पष्टपणेही परिभाषित करू शकते.

गणिती तर्कशास्त्रातील अनेक निष्कर्षांप्रमाणे,
अंतर्वेशन प्रमेयाची सिद्धता प्रतिमान उपपत्तीच्या
तसेच सिद्धता उपपत्तीच्या पद्धतींनी करता येते.
सिद्धता-उपपत्तीय पद्धतीचा लाभ असा की ती
केवळ अंतर्वेशकाचे अस्तित्व दाखवत नाही, तर
$A \Entails B$ याच्या सिद्धता पासून
अंतर्वेशक मोजून देणारी अल्गोरिदमही देते;
उदाहरणार्थ, $A \Sequent B$ याची~\Log{G3c}
मधील सिद्धता वापरता येते.

$A$ मधील स्थिरांक आणि विधेये
यांचा संच $\Lang L(A)$ असू द्या आणि
$\Lang L(\Gamma) = \bigcup \Setabs{\Lang L(A)}{A \in
\Gamma}$ असू द्या. या संपूर्ण विभागात
फलने नसलेली सूत्रे आपण विचारात घेतो.

\begin{thm}[क्रेगचे अंतर्वेशन प्रमेय]\label{pt:cut:itp:thm:interpolation}
$A$ ची भाषा $\Lang L(A)$ आणि
$B$ ची $\Lang L(B)$ असू द्या.
$\Log{G3c} \Proves A \Sequent B$ असेल, तर
$\Lang L(C) \subseteq \Lang L(A) \cap \Lang L(B)$
या भाषा मधील असे सूत्र~$C$ असते की
$\Log{G3c} \Proves A \Sequent C$
आणि $\Log{G3c} \Proves C \Sequent B$.
\end{thm}

~\Log{G3c} मधील सिद्धता च्या रचनेचा
उपयोग करण्यासाठी आपण अंतर्वेशन प्रमेयाचे
अधिक व्यापक विधान सिद्ध करू.
त्यासाठी क्रमवर्तींचे पूर्वांग व उत्तरांग
प्रत्येकी दोन भागांत विभागलेले मानू.
$\maeh{\Gamma_1}{\Gamma_2} \Sequent
\maeh{\Delta_1}{\Delta_2}$ या लेखनात
$\Gamma_1$ आणि $\Delta_1$ पहिल्या भागात,
तर $\Gamma_2$ आणि $\Delta_2$ दुसऱ्या
भागात आहेत. पुढील उपप्रमेयावरून
अंतर्वेशन प्रमेय लगेच मिळते.

\begin{lem}[माएहाराचे उपप्रमेय]\label{pt:cut:itp:lem:maehara}
$\Log{G3c}
\Proves \maeh{\Gamma_1}{\Gamma_2} \Sequent \maeh{\Delta_1}{\Delta_2}$
असेल, तर $\Lang L(C)
\subseteq \Lang L(\Gamma_1,
\Delta_1) \cap \Lang L(\Gamma_2, \Delta_2)$
या भाषा मधील असे सूत्र~$C$ असते की
$\Log{G3c}
\Proves \Gamma_1 \Sequent \Delta_1, C$ आणि
$\Log{G3c} \Proves C,
\Gamma_2 \Sequent \Delta_2$.
\end{lem}
चारही संदर्भांतील सूत्रे वाक्ये असतील, तर अंतर्वेशकातील
मुक्त चरे सार्वत्रिकपणे बांधून वाक्यरूप अंतर्वेशक घेता येतो.
एका मुक्त चरासाठी पहिल्या सिद्धतेत तो चर आधी नव्या
स्थिरांकाने बदलून $\RightR{\lforall}$ लावायचा; बंद संदर्भांमुळे
आयगेन चराची अट पूर्ण होते. दुसऱ्या सिद्धतेत त्या चराचे
पद वापरून $\LeftR{\lforall}$ आणि आवश्यक दुर्बलीकरण लावायचे.
प्रत्येक मुक्त चरासाठी हे केल्याने दोन्ही क्रमवर्ती सिद्धतायोग्य
राहतात आणि स्थिरांक-विधेयांची भाषा-सीमा बदलत नाही.

\begin{proof}
  $\pi \Proves \maeh{\Gamma_1}{\Gamma_2} \Sequent
  \maeh{\Delta_1}{\Delta_2}$ असे गृहीत धरा.
  $n=\pheight{\pi}$ यावर विगमन करून विधान सिद्ध करू.

  $n=0$ असेल, तर
  $\maeh{\Gamma_1}{\Gamma_2} \Sequent
  \maeh{\Delta_1}{\Delta_2}$ हे स्वयंसिद्धक आहे
  आणि एखादे अणुरूप सूत्र~$A$
  हे $\Gamma_1, \Gamma_2$ आणि
  $\Delta_1, \Delta_2$ या दोन्ही बाजूंना येते.
  डावीकडे आणि उजवीकडे $A$ पहिल्या की
  दुसऱ्या भागात आहे यावर चार प्रसंग ठरतात.
  \begin{enumerate}
    \item $A \in \Gamma_1$ आणि $A \in \Delta_1$.
    $\Gamma_1 \Sequent \Delta_1, C$ आणि
    $C, \Gamma_2 \Sequent \Delta_2$ या दोन्ही
    सिद्धतायोग्य असतील असे~$C$ हवे.
    $A$ दोन्ही बाजूंना असल्याने, $C$ काहीही
    निवडले तरी पहिली क्रमवर्ती आधीच स्वयंसिद्धक आहे.
    $C, \Gamma_2 \Sequent \Delta_2$ ही दुसरी
    क्रमवर्ती सिद्धतायोग्य असण्याची खात्री
    करण्यासाठी $C \ident \lfalse$ घ्यावे लागते.
    \item $A \in \Gamma_1$ आणि $A \in \Delta_2$.
    मग $\Gamma_1 \Sequent \Delta_1, A$ आणि
    $A, \Gamma_2 \Sequent \Delta_2$ ही दोन्ही
    स्वयंसिद्धके आहेत; म्हणून $C \ident A$ घेता येते.
    \item $A \in \Gamma_2$ आणि $A \in \Delta_1$.
    मग $A, \Gamma_1 \Sequent \Delta_1$ आणि
    $\Gamma_2 \Sequent \Delta_2, A$ ही स्वयंसिद्धके
    असतील; पण त्यांत~$A$ चुकीच्या बाजूंना आहे.
    म्हणून $A$ अंतर्वेशक नाही.
    तथापि $\RightR{\lnot}$ आणि $\LeftR{\lnot}$
    वापरून $\Gamma_1
    \Sequent \Delta_1, \lnot A$ आणि
    $\lnot A, \Gamma_2 \Sequent
    \Delta_2$ या क्रमवर्ती सिद्ध करता येतात.
    \item $A \in \Gamma_2$ आणि $A \in \Delta_2$.
    हा सरावासाठीचा प्रसंग आहे.
  \end{enumerate}
  $\lfalse \in \Gamma_1$ किंवा
  $\lfalse \in \Gamma_2$ यांमुळे
  $\maeh{\Gamma_1}{\Gamma_2} \Sequent
  \maeh{\Delta_1}{\Delta_2}$ स्वयंसिद्धक असण्याचे
  प्रसंगही सरावासाठी सोडले आहेत.

  $n>0$ असेल, तर $\pi$ च्या शेवटी
  तार्किक अनुमान आहे. कोणता नियम वापरला
  आणि मुख्य सूत्र कोणत्या भागात आहे, यानुसार
  प्रसंग वेगळे करावे लागतील. प्रत्येक प्रसंगात
  $\pheight{\pi_1} < n$ असलेल्या सर्व
  सिद्धता~$\pi_1$ साठी, त्यांची अंतिम
  क्रमवर्ती कोणत्याही प्रकारे दोन भागांत
  विभागली असली तरी, उपप्रमेय खरे आहे हे
  विगमन गृहीतक वापरता येते. काही विशिष्ट
  प्रसंग पाहू.
  \begin{enumerate}
    \item $\pi$ च्या शेवटी $\RightR{\lnot}$
    असून मुख्य सूत्र~$\lnot
    A$ आहे. $\lnot A$ पहिल्या की
    दुसऱ्या भागात आहे, यानुसार $\pi$ चा
    शेवट पुढीलपैकी एक असतो:
    \[
    \AxiomC{}
    \RightLabel{$\pi_1$}
    \Deduce$\maeh{A, \Gamma_1}{\Gamma_2} \fCenter
      \maeh{\Delta_1}{\Delta_2}$
    \RightLabel{\RightR{\lnot}}
    \UnaryInf$\maeh{\Gamma_1}{\Gamma_2} \fCenter
      \maeh{\Delta_1, \lnot A}{\Delta_2}$
    \DisplayProof
    \]
    किंवा
    \[
    \AxiomC{}
    \RightLabel{$\pi_1$}
    \Deduce$\maeh{\Gamma_1}{A, \Gamma_2} \fCenter
      \maeh{\Delta_1}{\Delta_2}$
    \RightLabel{\RightR{\lnot}}
    \UnaryInf$\maeh{\Gamma_1}{\Gamma_2} \fCenter
      \maeh{\Delta_1}{\Delta_2, \lnot A}$
    \DisplayProof
    \]
    मुख्य सूत्र~$\lnot A$ पहिल्या
    भागात असेल, तर आधारक्रमवर्तीतील सक्रिय
    सूत्र~$A$ देखील आपण पहिल्या
    भागातच ठेवले आहे. ही निवड आपल्याला
    करता येते. सक्रिय सूत्र दुसऱ्या भागात
    ठेवले तरी $\pi_1$ च्या
    अंतिम क्रमवर्तीला विगमन गृहीतक लागू होईल;
    मात्र मग अंतर्वेशक शोधता येणार नाही
    (हे करून पाहा).

    आधी $\lnot A$ पहिल्या भागात
    असलेला प्रसंग पाहू. विगमन गृहीतकानुसार
    $\pi_1$ च्या अंतिम क्रमवर्तीला
    अंतर्वेशक असतो, म्हणजे असे
    सूत्र~$C_1$ असते की पुढील दोन्ही
    क्रमवर्ती
    \begin{align*}
    A, \Gamma_1 & \Sequent \Delta_1, C_1 &&\text{आणि} &
    C_1, \Gamma_2 & \Sequent \Delta_2
    \intertext{सिद्ध करता येतात. आता आपल्याला असे सूत्र~$C$ हवे की
    पुढील दोन्ही सिद्धतायोग्य असतील:}
    \Gamma_1 & \Sequent \Delta_1, \lnot A, C &&\text{आणि} &
    C, \Gamma_2 & \Sequent \Delta_2
    \end{align*}
    स्पष्टपणे $C \ident C_1$ घेता येते:
    उजवीकडील क्रमवर्ती सिद्धतायोग्य आहे
    हे विगमन गृहीतकाने आधीच मिळते, आणि
    डावीकडील क्रमवर्ती विगमन गृहीतकातून
    मिळालेल्या क्रमवर्तीवर~$\RightR{\lnot}$
    लावून मिळते. विगमन गृहीतकाने
    $\Lang L(C_1) \subseteq \Lang L(A, \Gamma_1,
    \Delta_1) \cap \Lang L(\Gamma_2, \Delta_2)$
    हेदेखील मिळते. $\Lang L(\lnot
    A) = \Lang L(A)$ आणि
    $\Lang L(C) = \Lang L(C_1)$ असल्याने
    $\Lang L(C_1) \subseteq \Lang L(\Gamma_1, \Delta_1, \lnot A) \cap \Lang
    L(\Gamma_2, \Delta_2)$.

    दुसऱ्या प्रसंगात $\lnot A$
    दुसऱ्या भागात असल्याने स्थिती थोडी
    वेगळी आहे. आधारक्रमवर्तीतील सक्रिय
    सूत्र देखील दुसऱ्या भागात
    ठेवून विगमन गृहीतक लावतो; त्यातून
    \begin{align*}
    \Gamma_1 & \Sequent \Delta_1, C_1 &&\text{आणि} &
    C_1, A, \Gamma_2 & \Sequent \Delta_2
    \intertext{सिद्ध करता येतात. यांतील दुसऱ्या क्रमवर्तीवर पुन्हा \RightR{\lnot} लावल्यास पुढील सिद्धता मिळतात:}
    \Gamma_1 & \Sequent \Delta_1, C_1 &&\text{आणि} &
    C_1, \Gamma_2 & \Sequent \Delta_2, \lnot A
    \end{align*}
    $C_1$ अंतर्वेशक असेल तर नेमक्या
    याच क्रमवर्ती सिद्धतायोग्य हव्या होत्या;
    म्हणून पुन्हा $C \ident C_1$ घेऊ.
    पूर्वीप्रमाणे $\Lang
    L(C_1) \subseteq \Lang L(\Gamma_1, \Delta_1) \cap \Lang L(\Gamma_2,
    \Delta_2, \lnot A)$.
    \item $\pi$ च्या शेवटी $\RightR{\lif}$
    असून मुख्य सूत्र~$A
    \lif B$ आहे. $A \lif
    B$ पहिल्या की दुसऱ्या भागात आहे
    यावर पुन्हा दोन प्रसंग मिळतात.
    सिद्धता~$\pi$ चा शेवट असा असतो:
    \[
    \AxiomC{}
    \RightLabel{$\pi_1$}
    \Deduce$\maeh{A, \Gamma_1}{\Gamma_2} \fCenter
      \maeh{\Delta_1, B}{\Delta_2}$
    \RightLabel{\RightR{\lif}}
    \UnaryInf$\maeh{\Gamma_1}{\Gamma_2} \fCenter
      \maeh{\Delta_1, A \lif B}{\Delta_2}$
    \DisplayProof
    \]
    किंवा
    \[
    \AxiomC{}
    \RightLabel{$\pi_1$}
    \Deduce$\maeh{\Gamma_1}{A, \Gamma_2} \fCenter
      \maeh{\Delta_1}{\Delta_2, B}$
    \RightLabel{\RightR{\lif}}
    \UnaryInf$\maeh{\Gamma_1}{\Gamma_2} \fCenter
      \maeh{\Delta_1}{\Delta_2, A \lif B}$
    \DisplayProof
    \]
    येथेही आधारक्रमवर्तीतील सक्रिय
    सूत्रे निष्कर्षातील मुख्य
    सूत्राच्याच भागात ठेवली आहेत.
    पहिल्या प्रसंगात विगमन गृहीतकाने
    पुढील सिद्धता मिळतात:
    \begin{align*}
    A, \Gamma_1 & \Sequent \Delta_1, B, C_1 &&\text{आणि} & C_1, \Gamma_2 & \Sequent \Delta_2
    \intertext{डावीकडील क्रमवर्ती \RightR{\lif} साठी योग्य आधारक्रमवर्ती आहे. म्हणून पुढील क्रमवर्ती सिद्धतायोग्य आहेत:}
    \Gamma_1 & \Sequent \Delta_1, A \lif B, C_1 &&\text{आणि} & C_1, \Gamma_2 & \Sequent \Delta_2
    \end{align*}
    म्हणून $C_1$ हा $\pi$ च्या
    अंतिम क्रमवर्तीचाही अंतर्वेशक आहे;
    पुन्हा $C \ident C_1$ घेता येते.

    दुसरा प्रसंगही याच प्रकारे हाताळता येतो.
    \item $\pi$ च्या शेवटी $\LeftR{\lif}$
    असून मुख्य सूत्र~$A
    \lif B$ आहे. $A \lif B$
    पहिल्या भागात असेल, तर
    सिद्धता~$\pi$ चा शेवट असा असतो:
    \[
    \AxiomC{}
    \RightLabel{$\pi_1$}
    \Deduce$\maeh{\Gamma_1}{\Gamma_2} \fCenter
      \maeh{\Delta_1, A}{\Delta_2}$
    \AxiomC{}
    \RightLabel{$\pi_2$}
    \Deduce$\maeh{B, \Gamma_1}{\Gamma_2} \fCenter
      \maeh{\Delta_1}{\Delta_2}$
    \BinaryInf$\maeh{A \lif B, \Gamma_1}{\Gamma_2} \fCenter
      \maeh{\Delta_1}{\Delta_2}$
    \DisplayProof
    \]
    विगमन गृहीतकाने आता चार
    सिद्धतायोग्य क्रमवर्ती मिळतात:
    डाव्या आधारक्रमवर्तीच्या
    अंतर्वेशक~$C_1$ साठी दोन आणि
    दुसऱ्या आधारक्रमवर्तीच्या
    अंतर्वेशक~$C_2$ साठी दोन:
    \begin{align*}
    \Gamma_1 & \Sequent \Delta_1, A, C_1 &&\text{(1)} &
    C_1, \Gamma_2 & \Sequent \Delta_2 && \text{(2)}\\
    B, \Gamma_1 & \Sequent \Delta_1, C_2 &&\text{(3)} &
    C_2, \Gamma_2 & \Sequent \Delta_2 && \text{(4)}
    \end{align*}
    दुर्बलीकरण लावून (1) च्या
    उजवीकडे~$C_2$ आणि (3) च्या
    उजवीकडे~$C_1$ जोडल्यास
    (1) आणि (3) या क्रमवर्ती
    \LeftR{\lif} च्या आधारक्रमवर्ती
    म्हणून वापरता येतात:
    \[
    \AxiomC{}
    \Deduce$\Gamma_1 \fCenter \Delta_1, A, C_1, C_2$
    \AxiomC{}
    \Deduce$B, \Gamma_1 \fCenter \Delta_1, C_1, C_2$
    \RightLabel{\LeftR{\lif}}
    \BinaryInf$A \lif B, \Gamma_1 \fCenter \Delta_1, C_1, C_2$
    \DisplayProof
    \]
    $\RightR{\lor}$ लावल्याने पुढील
    क्रमवर्ती मिळते:
    \[
    A \lif B, \Gamma_1 \fCenter \Delta_1, C_1 \lor C_2
    \]
    त्यामुळे या प्रसंगात
    $C_1 \lor C_2$ अंतर्वेशक असू
    शकतो असे सुचते. उरलेल्या
    (2) आणि~(4) क्रमवर्ती वापरून
    आवश्यक दुसरी क्रमवर्तीही
    सिद्ध करता येते:
    \[
    \AxiomC{}
    \Deduce$C_1, \Gamma_2 \fCenter \Delta_2$
    \AxiomC{}
    \Deduce$C_2, \Gamma_2 \fCenter \Delta_2$
    \RightLabel{\LeftR{\lor}}
    \BinaryInf$C_1 \lor C_2, \Gamma_2 \fCenter \Delta_2$
    \DisplayProof
    \]
    $C_1 \lor C_2$ योग्य
    भाषेत आहे हे तपासणे उरते.
    विगमन गृहीतकानुसार
    \begin{align*}
    \Lang L(C_1) & \subseteq
    \Lang L(\Gamma_1, \Delta_1, A) \cap \Lang L(\Gamma_2, \Delta_2) &\text{आणि}\\
    \Lang L(C_2) & \subseteq
    \Lang L(B, \Gamma_1, \Delta_1) \cap \Lang L(\Gamma_2, \Delta_2)
    \intertext{कारण}
    \Lang L(A) & \subseteq \Lang L(A \lif B) &\text{आणि}\\
    \Lang L(B) & \subseteq \Lang L(A \lif B)
    \intertext{म्हणून या दोन्ही अटी मिळतात}
    \Lang L(C_1) & \subseteq
    \Lang L(\Gamma_1, \Delta_1, A \lif B) \cap \Lang L(\Gamma_2, \Delta_2) &\text{आणि}\\
    \Lang L(C_2) & \subseteq
    \Lang L(\Gamma_1, \Delta_1, A \lif B) \cap \Lang L(\Gamma_2, \Delta_2)
    \intertext{आणि परिणामी, $\Lang L (C_1 \lor C_2) = {}$}
    \Lang L(C_1) \cup \Lang L(C_2) & \subseteq
    \Lang L(\Gamma_1, \Delta_1, A \lif B) \cap \Lang L(\Gamma_2, \Delta_2),
    \end{align*}
    आणि हेच अपेक्षित आहे.

    $A \lif B$ दुसऱ्या
    भागात असण्याचा प्रसंग
    वेगळा आहे, पण तशाच
    पद्धतीने हाताळता येतो.
    विगमन गृहीतकाने पुन्हा
    चार सिद्धतायोग्य क्रमवर्ती मिळतात:
    \begin{align*}
    \Gamma_1 & \Sequent \Delta_1, C_1 &&\text{(1)} &
    C_1, \Gamma_2 & \Sequent \Delta_2, A && \text{(2)}\\
    \Gamma_1 & \Sequent \Delta_1, C_2 &&\text{(3)} &
    C_2, B, \Gamma_2 & \Sequent \Delta_2 && \text{(4)}
    \end{align*}
    आता (2) आणि (4) या क्रमवर्ती,
    काही सूत्रे दुर्बल करून
    जोडल्यावर,~\LeftR{\lif}
    च्या आधारक्रमवर्ती ठरतात:
    \[
    \AxiomC{}
    \Deduce$C_1, C_2, \Gamma_2 \fCenter \Delta_2, A$
    \AxiomC{}
    \Deduce$B, C_1, C_2, \Gamma_2 \fCenter \Delta_2$
    \RightLabel{\LeftR{\lif}}
    \BinaryInf$A \lif B, C_1, C_2, \Gamma_2 \fCenter \Delta_2$
    \DisplayProof
    \]
    यांतून $C_1$ आणि $C_2$
    वापरून बनलेले एकच सूत्र
    असणारी क्रमवर्ती हवी आहे;
    पण यावेळी ते पूर्वांगात हवे
    (कारण आता दुसऱ्या भागाशी
    संबंध आहे). $\LeftR{\land}$
    लावल्याने ते साधते; म्हणून
    $C \ident
    (C_1 \land C_2)$ अंतर्वेशक घेऊ.
    उरलेल्या (1) आणि (3)
    क्रमवर्ती वापरून पहिल्या
    भागाची संबंधित क्रमवर्तीही
    सिद्ध करता येते:
    \[
    \AxiomC{}
    \Deduce$\Gamma_1 \fCenter \Delta_1, C_1$
    \AxiomC{}
    \Deduce$\Gamma_1 \fCenter \Delta_1, C_2$
    \RightLabel{\RightR{\land}}
    \BinaryInf$\Gamma_1 \fCenter \Delta_1, C_1 \land C_2$
    \DisplayProof
    \]
    पूर्वीप्रमाणे
    $\Lang L(C_1 \land C_2) \subseteq \Lang
    L(\Gamma_1, \Delta_1) \cap \Lang L(\Gamma_2, \Delta_2, A \lif
    B)$ मिळते.

    \item $\pi$ च्या शेवटी $\RightR{\lforall}$
    असून मुख्य सूत्र~$\lforall[x][A(x)]$ आहे:
    \[
    \AxiomC{}
    \RightLabel{$\pi_1$}
    \Deduce$\maeh{\Gamma_1}{\Gamma_2} \fCenter
    \maeh{\Delta_1, A(c)}{\Delta_2}$
    \RightLabel{\RightR{\lforall}}
    \UnaryInf$\maeh{\Gamma_1}{\Gamma_2} \fCenter
      \maeh{\Delta_1, \lforall[x][A(x)]}{\Delta_2}$
    \DisplayProof
    \]
    किंवा
    \[
    \AxiomC{}
    \RightLabel{$\pi_1$}
    \Deduce$\maeh{\Gamma_1}{\Gamma_2} \fCenter
      \maeh{\Delta_1}{\Delta_2, A(c)}$
    \RightLabel{\RightR{\lforall}}
    \UnaryInf$\maeh{\Gamma_1}{\Gamma_2} \fCenter
      \maeh{\Delta_1}{\Delta_2, \lforall[x][A(x)]}$
    \DisplayProof
    \]
    पहिल्या प्रसंगात विगमन गृहीतकाने
    $\Gamma_1 \Sequent \Delta_1, A(c), C_1$
    आणि $C_1, \Gamma_2
    \Sequent \Delta_2$ यांच्या सिद्धता मिळतात.
    पहिल्या क्रमवर्तीवर $\RightR{\lforall}$
    लावून $\Gamma_1 \Sequent \Delta_1, \lforall[x][A(x)], C_1$
    मिळते. ($\pi$ च्या अंतिम
    \RightR{\lforall} अनुमानात आयगेन चराची
    अट आधीच पूर्ण असल्याने येथेही ती पूर्ण आहे.)
    $\Lang
    L(C_1) \subseteq \Lang L(\Gamma_1, \Delta_1, \lforall[x][A(x)])
    \cap \Lang L(\Gamma_2, \Delta_2)$
    ही अट पूर्ण असेल, तर $C_1$ अंतर्वेशक ठरेल.
    विगमन गृहीतकाने
    $\Lang L(C_1) \subseteq \Lang L(\Gamma_1,
    \Delta_1, A(c)) \cap \Lang L(\Gamma_2, \Delta_2)$
    मिळते. म्हणून~$c$ ची काळजी घ्यावी लागते:
    $C_1$ मध्ये $c$ येऊ शकतो का? येऊ शकत नाही.
    $C_1$ मध्ये $c$ आल्यास
    $c \in \Lang L(\Gamma_1, \Delta_1,
    A(c))$ (जे खरे आहे) आणि
    $c \in \Lang L(\Gamma_2,
    \Delta_2)$ हेही दोन्ही खरे ठरतील.
    पण मग $\pi$ मधील
    $\RightR{\lforall}$ अनुमानाच्या
    निष्कर्षात $c$ येईल आणि
    आयगेन चराची अट मोडेल.

    दुसरा प्रसंगही यासारखाच आहे.

    \item $\pi$ च्या शेवटी
    $\RightR{\lexists}$ असून मुख्य
    सूत्र~$\lexists[x][A(x)]$ आहे.
    येथेही दोन प्रसंग आहेत.
    $\lexists[x][A(x)]$ पहिल्या
    भागात असलेला प्रसंग पाहू;
    दुसरा सरावासाठी सोडला आहे.

    पहिल्या प्रसंगात $\pi$ चा
    शेवट असा असतो:
    \[
    \AxiomC{}
    \RightLabel{$\pi_1$}
    \Deduce$\maeh{\Gamma_1}{\Gamma_2} \fCenter
      \maeh{\Delta_1, \lexists[x][A(x)], A(t)}{\Delta_2}$
    \RightLabel{\RightR{\lexists}}
    \UnaryInf$\maeh{\Gamma_1}{\Gamma_2} \fCenter
      \maeh{\Delta_1, \lexists[x][A(x)]}{\Delta_2}$
    \DisplayProof
    \]
    \begingroup\sloppy
    विगमन गृहीतकाने
    $\Gamma_1 \Sequent
    \Delta_1, \lexists[x][A(x)], A(t), C_1$
    आणि $C_1, \Gamma_2
    \Sequent \Delta_2$ यांच्या सिद्धता मिळतात.
    पहिल्या क्रमवर्तीवर $\RightR{\lexists}$
    लावल्याने $\Gamma_1 \Sequent \Delta_1, \lexists[x][A(x)], C_1$
    मिळते. म्हणून
    $\Lang L(C_1) \subseteq \Lang L(\Gamma_1, \Delta_1,
    \lexists[x][A(x)]) \cap \Lang L(\Gamma_2, \Delta_2)$
    ही अट पूर्ण असल्यास $C_1$ पुन्हा
    अंतर्वेशक ठरेल. येथे विगमन गृहीतकाने
    $\Lang L(C_1)
    \subseteq \Lang L(\Gamma_1, \Delta_1, \lexists[x][A(x)], A(t))
    \cap \Lang L(\Gamma_2, \Delta_2)$
    मिळते. आपण फलने नाहीत असे
    गृहीत धरले आहे, हे लक्षात घ्या.
    म्हणून पद~$t$ हे चर किंवा स्थिरांक असते.
    $t$ चर असेल, तर त्याच्या प्रतिस्थापनाने नवीन
    स्थिरांक किंवा विधेय-चिन्ह येत नाही:
    $\Lang L(A(t))\subseteq\Lang L(A(x))$.
    त्यामुळे वरील विगमन गृहीतकातील भाषा-सीमा
    अंतिम क्रमवर्तीच्या दोन भागांच्या भाषांच्या
    छेदातच राहते आणि $C\ident C_1$ घेता येते.
    $\RightR{\lexists}$ साठी आयगेन चराची अट नाही.
    आता $t$ हे स्थिरांक असल्याचा प्रसंग,
    म्हणजे $t\ident b$, पाहू.
    $b \notin \Lang L(C_1)$ असेल,
    किंवा $b \in \Lang L(\Gamma_1, \Delta_1, \lexists[x][A(x)]) \cap
    \Lang L(\Gamma_2, \Delta_2)$ असेल,
    तर अडचण नाही: $C_1$ अंतर्वेशक घेता येते.
    पण $b \in \Lang
    L(C_1)$ आणि $b \notin \Lang L(\Gamma_1, \Delta_1,
    \lexists[x][A(x)]) \cap \Lang L(\Gamma_2, \Delta_2)$
    असेल तर? प्रथम~$C_1$ मधील
    $b$ ची प्रत्येक उपस्थिती
    $C_1$ मध्ये न येणारा नवा चल~$y$
    निवडून त्याने बदलल्यावर
    $C_1 \ident \Subst{C_1'}{b}{y}$
    असे लिहिता येते; येथे $C_1'$
    मध्ये~$b$ येत नाही. शिवाय,
    $b \notin
    \Lang L(\Gamma_1, \Delta_1, \lexists[x][A(x)])$
    किंवा $b \notin
    \Lang L(\Gamma_2, \Delta_2)$.
    त्यानुसार $C \ident
    \lforall[y][C_1'(y)]$ किंवा
    $C \ident \lexists[y][C_1'(y)]$
    घेता येते. पहिल्या प्रसंगात:\par\endgroup
    \[
    \AxiomC{}
    \Deduce$\Gamma_1 \fCenter
      \Delta_1, \lexists[x][A(x)], A(b), C_1'(b)$
    \RightLabel{\RightR{\lexists}}
    \UnaryInf$\Gamma_1 \fCenter \Delta_1, \lexists[x][A(x)], C_1'(b)$
    \RightLabel{\RightR{\lforall}}
    \UnaryInf$\Gamma_1 \fCenter
      \Delta_1, \lexists[x][A(x)], \lforall[y][C_1'(y)]$
    \DisplayProof
    \]
    आणि
    \[
    \AxiomC{}
    \Deduce$C_1'(b), \lforall[y][C_1'(y)], \Gamma_2 \fCenter \Delta_2$
    \RightLabel{\LeftR{\lforall}}
    \UnaryInf$\lforall[y][C_1'(y)], \Gamma_2 \fCenter \Delta_2$
    \DisplayProof
    \]
    सर्वात वरच्या क्रमवर्तींच्या
    सिद्धता विगमन गृहीतकातून मिळतात
    (दुसऱ्या सिद्धतेत पूर्वांगात
    $\lforall[y][C_1'(y)]$ जोडण्यासाठी
    दुर्बलीकरणाची ग्राह्यता वापरतो).
    पहिल्या सिद्धता मधील
    \RightR{\lforall} साठी आयगेन चराची
    अट पूर्ण आहे, कारण
    $b \notin \Lang L(\Gamma_1, \Delta_1,
    \lexists[x][A(x)])$ आणि
    $C_1'$ मध्ये~$b$ येत नाही,
    अशीच त्याची व्याख्या केली आहे.

    दुसऱ्या प्रसंगात, विगमन गृहीतक
    आणि $\lexists[y][C_1'(y)]$
    जोडणाऱ्या दुर्बलीकरणाच्या
    ग्राह्यतेने पुढील मिळते:
    \[
    \AxiomC{}
    \Deduce$\Gamma_1 \fCenter
      \Delta_1, \lexists[x][A(x)], A(b), \lexists[y][C_1'(y)], C_1'(b)$
    \RightLabel{\RightR{\lexists}}
    \UnaryInf$\Gamma_1 \fCenter
      \Delta_1, \lexists[x][A(x)], \lexists[y][C_1'(y)], C_1'(b)$
    \RightLabel{\RightR{\lexists}}
    \UnaryInf$\Gamma_1 \fCenter
      \Delta_1, \lexists[x][A(x)], \lexists[y][C_1'(y)]$
    \DisplayProof
    \]
    आणि
    \[
    \AxiomC{}
    \Deduce$C_1'(b), \Gamma_2 \fCenter \Delta_2$
    \RightLabel{\LeftR{\lexists}}
    \UnaryInf$\lexists[y][C_1'(y)], \Gamma_2 \fCenter \Delta_2$
    \DisplayProof
    \]
    दुसऱ्या सिद्धता मधील
    \LeftR{\lexists} साठी आयगेन चराची
    अट पूर्ण आहे, कारण
    $b \notin \Lang L(\Gamma_2,
    \Delta_2)$ आणि $C_1'$ मध्ये~$b$
    येत नाही, अशीच त्याची व्याख्या केली आहे.
  \end{enumerate}
  उरलेले प्रसंगही यांसारखे आहेत;
  ते सरावासाठी सोडले आहेत.
\end{proof}

\begin{prob}
$\lnand$ हा संयोजक पुढील नियमांसह
आपल्याकडे आहे असे समजा:
\[
\Axiom$\Gamma \fCenter \Delta, A$
\Axiom$\Gamma \fCenter \Delta, B$
\RightLabel{\LeftR{\lnand}}
\BinaryInf$A \lnand B, \Gamma \fCenter \Delta$
\DisplayProof
\quad
\Axiom$A, B, \Gamma \fCenter \Delta$
\RightLabel{\RightR{\lnand}}
\UnaryInf$\Gamma \fCenter \Delta, A \lnand B$
\DisplayProof
\]
ज्या प्रसंगात $\pi$ चा शेवट
यांपैकी एखाद्या नियमाने होतो,
ते \ref{pt:cut:itp:lem:maehara} च्या सिद्धतेत
पूर्ण करून माएहाराचे उपप्रमेय
अजूनही खरे आहे हे दाखवा.
\end{prob}

माएहाराचे उपप्रमेय सिद्ध झाल्यावर
त्यावरून क्रेगचे अंतर्वेशन प्रमेय
सिद्ध करता येते.

\begin{proof}[\ref{pt:cut:itp:thm:interpolation} ची सिद्धता]
  $A \Sequent B$ सिद्धतायोग्य आहे असे
  समजा. $\maeh{A}{\ } \Sequent \maeh{\ }{B}$
  याला \ref{pt:cut:itp:lem:maehara} लावून
  अंतर्वेशक~$C$ मिळवा. मग
  $\Proves A \Sequent C$ आणि
  $\Proves C
  \Sequent B$.
\end{proof}

$\Gamma$ ही वाक्यांचा संच असलेली
उपपत्ती, भाषा~$\Lang L$ मध्ये आहे
आणि त्यात $n$-स्थानी {विधेय}~$R$ आहे
असे समजा. पुढील अट पूर्ण असेल, तेव्हा
आणि केवळ तेव्हाच, $\Gamma$ ही
$R$ ला \emph{अंतर्निहितपणे परिभाषित करते}:
\begin{align*}
\Gamma, \Gamma'
& \Entails \lforall[x_1][\dots\lforall[x_n][(R(x_1,\dots,x_n) \liff
R'(x_1, \dots, x_n))]],
\intertext{येथे $\Gamma'$ म्हणजे $\Gamma$ मध्ये $R$ ची प्रत्येक उपस्थिती
$R' \notin \Lang L$ ने बदलल्यावर मिळणारी उपपत्ती. $\Gamma$ ही $R$ ला
\emph{स्पष्टपणे परिभाषित करते}, तेव्हा आणि केवळ तेव्हाच, $R$ नसलेले असे
$A(x_1, \dots, x_n)$ असते की}
\Gamma &\Entails
\lforall[x_1][\dots\lforall[x_n][(R(x_1,\dots,x_n) \liff A(x_1,
\dots, x_n))]].
\end{align*}
स्पष्टपणे, $\Gamma$ ही $R$ ला
स्पष्टपणे परिभाषित करत असेल, तर
ती $R$ ला अंतर्निहितपणेही परिभाषित करते.
उलट विधान लगेच स्पष्ट होत नसले,
तरी ते खरे आहे:

\begin{lem}[बेथचे परिभाष्यता प्रमेय]
  $\Gamma$ ही $R$ ला अंतर्निहितपणे
  परिभाषित करत असेल, तर ती त्याला
  स्पष्टपणेही परिभाषित करते.
\end{lem}

\begin{proof}
  $\Gamma$ ही $R$ ला अंतर्निहितपणे
  परिभाषित करते असे समजा. वरीलप्रमाणे
  $R'$ आणि $\Gamma'$ घ्या;
  $\Lang L' = \Lang L(\Gamma') = \Lang L \setminus \{R\}
  \cup \{R'\}$ असू द्या. गृहीतकानुसार
  \begin{align*}
    \Gamma, \Gamma' &
    \Sequent \lforall[x_1][\dots\lforall[x_n][(R(x_1,\dots,x_n)
    \liff R'(x_1, \dots, x_n))]]
  \intertext{$c_1$, \dots,~$c_n$ हे $\Gamma$ मध्ये न येणारे स्थिरांक असू द्या. व्युत्क्रमणीयतेच्या उपप्रमेयाने मिळते:}
    \Gamma, \Gamma', R(c_1,\dots,c_n) & \Sequent  R'(c_1, \dots, c_n)
  \intertext{माएहाराचे उपप्रमेय पुढील क्रमवर्तीला लावा:}
    \maeh{\Gamma, R(c_1,\dots,c_n)}{\Gamma'} & \Sequent
      \maeh{\ }{R'(c_1, \dots, c_n)}
  \intertext{या बंद संदर्भांसाठी आवश्यक मुक्त चरे सार्वत्रिकपणे बांधून वाक्यरूप अंतर्वेशक निवडा. मग नव्या स्थिरांकांच्या जागी संबंधित चरे ठेवून असे~$A \ident A(x_1, \dots, x_n)$ मिळते की}
    \Gamma, R(c_1,\dots,c_n) & \Sequent A(c_1, \dots, c_n) \\
    A(c_1, \dots, c_n), \Gamma' & \Sequent R'(c_1, \dots, c_n)
  \intertext{यांच्या सिद्धता असतात. $\Lang L(A) \subseteq
  (\Lang L(\Gamma) \setminus \{R\}) \cup \{c_1,\dots,c_n\}$ असल्याने $A$ मध्ये~$R$ नसतो.
  दुसऱ्या क्रमवर्तीच्या सिद्धता मध्ये सर्वत्र $R'$ च्या जागी~$R$ ठेवून
  पुढील सिद्धता मिळते:}
    A(c_1, \dots, c_n), \Gamma & \Sequent R(c_1, \dots, c_n)
  \end{align*}
  येथील नव्या स्थिरांकांच्या जागी
  संबंधित चरे ठेवता येतात.
  \RightR{\lif} आणि
  \RightR{\lforall} लावून पुढील
  सिद्धता मिळते:
  \[
    \Gamma \Sequent \lforall[x_1][\dots\lforall[x_n][(R(x_1,\dots,x_n) \liff A(x_1, \dots, x_n))]].
  \]
  म्हणून $A(x_1, \dots, x_n)$ हे
  $\Gamma$ मध्ये $R$ ला स्पष्टपणे
  परिभाषित करते.
\end{proof}

अंतर्वेशनाशी समतुल्य असलेला
आणि त्याप्रमाणेच माएहाराचे उपप्रमेय वापरून
सिद्ध करता येणारा तिसरा निष्कर्ष असा:
दोन सुसंगत उपपत्ती~$\Gamma_1$ आणि
$\Gamma_2$ यांत $\Lang L(\Gamma) = \Lang L(\Gamma_1)
\cap \Lang L(\Gamma_2)$ या भाषेतील
पूर्ण उपपत्ती~$\Gamma$ समाविष्ट असेल,
तर $\Gamma_1 \cup \Gamma_2$ सुसंगत असते.

पूर्ण उपपत्ती म्हणजे तिच्या भाषेतील
प्रत्येक वाक्य~$A$ साठी
$\Gamma \Proves A$ किंवा
$\Gamma \Proves \lnot A$ सिद्ध होणारी उपपत्ती,
हे आठवा. $\Gamma_1$ आणि~$\Gamma_2$
या $\Gamma$ च्या सुसंगत विस्तारे
असल्याने $\Gamma$ सुसंगतच असली पाहिजे.
म्हणून $A$ हे
$\Lang L(\Gamma_1) \cap \Lang
L(\Gamma_2)$ या भाषेतील
वाक्य असेल, तर
$\Gamma_1 \Entails A$ आणि
$\Gamma_2
\Entails \lnot A$ दोन्ही शक्य नाहीत.
तसे एखादे~$A$ आहे असे समजा.
$\Gamma_1 \Entails A$ असेल, तर
$\Gamma \Entails A$.
अन्यथा, $\Gamma$ पूर्ण असल्याने
$\Gamma \Entails \lnot A$;
आणि $\Gamma \subseteq \Gamma_1$
असल्याने $\Gamma_1 \Entails \lnot A$,
म्हणजे $\Gamma_1$ विसंगत ठरेल.
म्हणून $\Gamma \Entails A$;
आता $\Gamma
\subseteq \Gamma_2$ असल्याने
$\Gamma_2 \Entails A$.
त्यामुळे $\Gamma \Entails/
\lnot A$; अन्यथा $\Gamma_2$
विसंगत ठरेल.

संयुक्त भाषा~
$\Lang L(\Gamma_1) \cap \Lang L(\Gamma_2)$
मधील असे वाक्य~$A$
नसणे की $\Gamma_1
\Entails A$ आणि
$\Gamma_2 \Entails \lnot A$,
हेच सिद्धतेसाठी पुरेसे आहे.
आता माएहाराचे उपप्रमेय वापरून
ही सिद्धता सिद्धता-उपपत्तीय
पद्धतीने देऊ.

\begin{thm}[रॉबिन्सनचे संयुक्त सुसंगतता प्रमेय]
$\Gamma_1$ आणि $\Gamma_2$
या सुसंगत उपपत्ती असू द्या.
$\Lang L(\Gamma_1) \cap \Lang L(\Gamma_2)$
या भाषा मधील कोणत्याही~$A$
साठी $\Gamma_1 \Entails A$
आणि $\Gamma_2 \Entails \lnot A$
दोन्ही होत नसतील, तर
$\Gamma_1 \cup \Gamma_2$ सुसंगत असते.
\end{thm}

\begin{proof}
  उलट, $\Gamma_1 \cup \Gamma_2$
  विसंगत आहे असे समजा.
  मग $\Gamma_1, \Gamma_2 \Sequent \ $
  या क्रमवर्तीची सिद्धता आहे.
  $\maeh{\Gamma_1}{\Gamma_2}
  \Sequent \maeh{\ }{\ }$ याला
  माएहाराचे उपप्रमेय लावून, आवश्यक असल्यास वर
  सांगितल्याप्रमाणे अंतर्वेशकातील मुक्त चरे बांधल्याने,
  $\Lang L(\Gamma_1) \cap \Lang L(\Gamma_2)$
  या भाषा मधील असे
  वाक्य~$A$ मिळते की
  $\Proves \Gamma_1 \Sequent A$
  आणि $\Proves A, \Gamma_2 \Sequent
  \quad$. दुसऱ्या क्रमवर्तीवरून
  $\Proves \Gamma_2 \Sequent \lnot A$
  मिळते. पण असे वाक्य
  अस्तित्वात नाही, हेच आपण
  गृहीत धरले होते.
\end{proof}

वर आपण उपपत्ती $\Gamma$,
$\Gamma_1$ आणि $\Gamma_2$
सांत आहेत असे नकळत गृहीत धरले.
संहततेमुळे हे निष्कर्ष
अनंत उपपत्तींनाही लागू होतात.



\chapter{नैसर्गिक निगमन}\label{pt:nat:chap}









\section{परिचय}\label{pt:nat:int:sec}

नैसर्गिक निगमन हा सिद्धता
पद्धतींचा एक परिवार आहे.
त्यात प्रत्येक तार्किक संयोजकासाठी
किंवा संख्यापकासाठी दोन
अनुमाने असतात: एक संकारकाचे
“प्रवेशन” करणारे
(परिचय नियम)
आणि दुसरे त्याचे “विलोपन”
करणारे (निरसन नियम).
उदाहरणार्थ, \Intro{\lif}
च्या निष्कर्षात
$A \lif B$ या रूपाची
सूत्रे असतात,
तर \Elim{\lif} च्या
आधारविधानांत
$A \lif B$ या रूपाची
सूत्रे असतात.

नैसर्गिक निगमनाच्या
पहिल्या दोन आवृत्त्या
१९३० च्या दशकात गरहार्ड
गेंटझेन आणि स्टॅनिस्लॉ
यास्कोव्हस्की यांनी मांडल्या.
सिद्धता-उपपत्तीचे अभ्यासक
बहुधा गेंट्सेनच्या पद्धतीची
चर्चा करतात; अध्यापनात
सर्वाधिक वापरल्या जाणाऱ्या
पद्धतींचा उगम याश्कोव्स्कीच्या
पद्धतीत आहे. दोन्हीत
\emph{गृहीतकांपासून} सुरुवात
करता येते: अशी सूत्रे
की ज्यांना सिद्ध करावे
लागत नाही, पण इतर
सूत्रे सिद्ध करणे
करताना वापरता येते.
संपूर्ण सिद्धता अशा
गृहीतकांवर अवलंबून असू शकते.
ही गृहीतके असिद्ध असल्यामुळे
ती सिद्धता तिचा निष्कर्ष
(तिचे \emph{अंतिम-सूत्र})
सिद्ध करत नाही; निष्कर्ष
गृहीतकांपासून निष्पन्न होतो,
एवढेच ती दाखवते.

काही अनुमान नियमांचा
निष्कर्ष काही गृहीतकांवर
आता अवलंबून राहत नाही:
असे नियम गृहीतके “मुक्त”
करतात. उदाहरणार्थ,
\Intro{\lif} नियम
गृहीतक~$A$ पासून
निष्कर्ष~$B$ ची
सिद्धता घेऊन
$A \lif B$ हा
निष्कर्ष देतो.
$A \lif B$ मध्ये
संपणारी सिद्धता
आता~$A$ वर अवलंबून
नसते: ते गृहीतक
मुक्त झालेले असते.
गेंट्सेनच्या पद्धतीत
सिद्धता हे
सूत्रांचे वृक्ष
असतात; गृहीतके सर्वात
वरची सूत्रे असतात,
आणि \Intro{\lif}
सारख्या नियमांवरील खुणा
कोणती गृहीतके मुक्त होतात
ते दाखवतात.
याश्कोव्स्कीच्या पद्धतीत
एखाद्या गृहीतकावर अवलंबून
असलेले सिद्धता चे भाग
वेगळे दाखवतात: उदाहरणार्थ,
उभी रेषा लावून आत सरकवलेले
किंवा चौकटीत बंद केलेले.
त्या चौकटीची पहिली ओळ
गृहीतक~$A$ असते;
अंतिम ओळ सोपाधिक विधानाचे
उत्तरांग $B$ असते;
$A$ मुक्त केले जाते;
आणि \Intro{\lif} चा
निष्कर्ष~$A \lif B$
चौकटीबाहेर असतो.

या पद्धती “नैसर्गिक”
आहेत, कारण त्यांचे अनुमान
नियम आपण (निदान गणितज्ञ)
सहजपणे जसे युक्तिवाद करतो
त्याशी जुळतात. एखादे
सोपाधिक विधान सिद्ध
करण्यासाठी त्याचे पूर्वांग
गृहीत धरून त्यावरून उत्तरांग
सिद्ध करतो: तो \Intro{\lif}.
सोपाधिक विधान $A \lif B$
आणि त्याचे पूर्वांग~$A$
माहीत असेल, तर~$B$
निष्पन्न करतो: तो
\Elim{\lif}.
विकल्पयोग~$A \lor B$
गृहीत धरून काही निष्कर्ष
दाखवण्यासाठी प्रकरणांनुसार
सिद्धता वापरतो: तो
\Elim{\lor}.
सार्वत्रिक विधान
$\lforall[x][A(x)]$
सिद्ध करण्यासाठी
कोणताही~$c$ घेऊन
$A(c)$ सिद्ध करतो:
तो \Intro{\lforall}.
अशाच प्रकारे पुढे जाता येते.

उदाहरणार्थ, गेंट्सेन-पद्धतीच्या
नैसर्गिक निगमनात
$A \lif (A \lor B)$
ची ही सोपी सिद्धता पाहा:
\[
\AxiomC{$\Discharge{A}{x}$}
\RightLabel{\Intro{\lor}}
\UnaryInfC{$A \lor B$}
\DischargeRule{\Intro{\lif}}{x}
\UnaryInfC{$A \lif (A \lor B)$}
\DisplayProof
\]
ती गृहीतक~$A$ पासून
सुरू होते; \Intro{\lor}
वापरून~$A \lor B$
निष्पन्न करते; मग
\Intro{\lif} वापरून
$A \lif (A \lor B)$
निष्पन्न करते.
\Intro{\lif} अनुमान
गृहीतक~$A$ मुक्त करते;
खूण~$x$ ते दर्शवते.

सिद्धता-उपपत्तीचे अभ्यासक
सूत्रांच्या वृक्षांवर
काम करणाऱ्या गेंट्सेनच्या
मूळ पद्धती पसंत करतात;
तरी एक रूपांतरही
सिद्धता-उपपत्तीत महत्त्वाचे आहे.
गृहीतके~$\Gamma$ यांपासून
$A$ ची सिद्धता
म्हणजे $A$ हे~$\Gamma$
पासून निष्पन्न होते,
अर्थात $\Gamma \Entails A$,
हे दाखवते.
नैसर्गिक निगमनाचे अनुमान
नियम अशा निष्पन्नतेबद्दल
विधान करतात असेही
स्वाभाविकपणे वाचता येते.
उदाहरणार्थ, \Intro{\lif}
म्हणतो की जर $\Gamma + A
\Entails B$, तर
$\Gamma \Entails A \lif B$.
म्हणून नियम थेट गृहीतके
आणि त्यांपासून निष्पन्न
होणारा निष्कर्ष या दोहोंवर
लागू होतील असे त्यांचे
रूप देता येते: म्हणजे
ते $\Gamma \Sequent A$
या क्रमवर्तींवर काम करतात.
वरील सोपी सिद्धता अशा
पद्धतीत पुढीलप्रमाणे दिसेल:
\[
\Axiom$x:A \fCenter A$
\RightLabel{\Intro{\lor}}
\UnaryInf$x:A \fCenter A \lor B$
\DischargeRule{\Intro{\lif}}{x}
\UnaryInf$\fCenter A \lif (A \lor B)$
\DisplayProof
\]
येथे नियम तेच आहेत:
ते~$\Sequent$ च्या
उजवीकडील सूत्र
वर काम करतात;
डावीकडील सूत्रे
प्रत्येक टप्प्यावर कोणत्या
गृहीतकांवर अवलंबित्व आहे
ते फक्त दाखवतात.

परिचय/निरसन
या नियम-जोड्यांचा पूर्ण संच
स्वतःहून अभिजात तर्कशास्त्रासाठी
पूर्ण नसतो. त्यात~$\lfalse$
शी संबंधित आणखी दोन
नियम जोडावे लागतात.
पहिला “अंतःप्रज्ञावादी
असत्यता नियम”~\FalseInt,
म्हणजे “स्फोट”:
$\lfalse$ पासून काहीही
निष्पन्न करता येते.
दुसरा अप्रत्यक्ष सिद्धतेचा
अभिजात सिद्धांत~\FalseCl{}
(किंवा “अभिजात असत्यता नियम”):
$\lnot A$ पासून
$\lfalse$ सिद्ध
करता येत असेल, तर
$A$ सिद्ध झाले.
\FalseCl वगळल्यास
अंतःप्रज्ञावादी तर्कशास्त्राला
योग्य अशी पद्धत मिळते.
दोन्ही नियम वगळल्यास
“किमान तर्कशास्त्र”
मिळते.

अभिजात आणि
अंतःप्रज्ञावादी तर्कशास्त्रासाठी
गेंट्सेन-पद्धतीच्या
नैसर्गिक निगमनाची आपण
चर्चा करू. या
\Log{N1c} आणि \Log{N1i}
पद्धती आहेत.
\Log{N2c} आणि \Log{N2i}
या संबंधित पद्धती
केवळ सूत्रे~$A$
यांवर काम करण्याऐवजी
$\Gamma \Sequent A$
या क्रमवर्तींवर काम करतात.












\section{नियम आणि सिद्धता}\label{pt:nat:rp:sec}

प्रथम काही व्याख्या करू
आणि संज्ञा निश्चित करू.

~\Log{N1c} आणि~\Log{N1i}
मधील $A \lif B$
उल्लेखणारे दोन नियम पाहू:
\[
\bottomAlignProof
\AxiomC{$\Discharge{A}{x}$}
\shortDeduce
\DeduceC{$B$}
\DischargeRule{\Intro{\lif}}{x}
\UnaryInfC{$A \lif B$}
\DisplayProof
\qquad
\bottomAlignProof
\AxiomC{$A \lif B$}
\AxiomC{$A$}
\RightLabel{\Elim{\lif}}
\BinaryInfC{$B$}
\DisplayProof
\]
\Elim{\lif} नियम सोपा आहे.
रेषेच्या वरची सूत्रे
ही \emph{आधारविधाने} होत.
डावीकडील सूत्र
हे \emph{मुख्य}
सूत्र~$A \lif B$
आहे. सर्व निरसन
नियमांत असे एक आधारविधान
असते आणि ते नेहमी
सर्वात डावीकडे असते.
त्याला अनुमानाचे
\emph{प्रमुख} आधारविधान
म्हणतात. या उदाहरणाप्रमाणे
इतर आधारविधाने असतील,
तर त्यांना \emph{गौण}
आधारविधाने म्हणतात.
परिचय नियमांची
आधारविधानेही प्रमुख
आधारविधाने म्हणतात.
रेषेखालील सूत्र
हा \emph{निष्कर्ष} आहे.

\Intro{\lif} नियमाला
एकच आधारविधान असते;
येथे निष्कर्ष हे
मुख्य सूत्र~$A \lif B$
असते. सर्व परिचय
नियम असेच असतात:
निष्कर्ष हे मुख्य सूत्र
असते, आणि त्याचा
मुख्य संकारकच “प्रवेशित”
केला जातो.

~$\Discharge{A}{x}$
या लेखनाचा अर्थ पुढीलप्रमाणे.
एखाद्या सिद्धता मध्ये
\Intro{\lif} अनुमान
आधारविधान म्हणून
सूत्र~$B$
वर लावतात.
~$B$ मध्ये संपणाऱ्या
उप-सिद्धता च्या
सुरुवातीला काही
सूत्रे असतात;
त्यांना \emph{गृहीतके}
म्हणतात. ~$A$
रूपाची एक किंवा अधिक
गृहीतके धरून $B$
ची सिद्धता म्हणजे
$A$ वरून~$B$
निष्पन्न होते हे दाखवणारी
सिद्धता होय.
\Intro{\lif} लावल्यावर
$A \lif B$
निष्पन्न होते; आता
निष्कर्ष गृहीतक~$A$
वर अवलंबून राहत नाही.
हे दाखवण्यासाठी,
\Intro{\lif} असलेल्या
सिद्धता मधील
$A$ रूपाची संबंधित
गृहीतके \emph{मुक्त}
किंवा \emph{बंद} केली
जातात. कोणती गृहीतके
मुक्त करता येतात, हे
नियम सांगतो.
$A \lif B$ निष्कर्ष
असलेल्या \Intro{\lif}
मध्ये ती $A$ रूपाची,
म्हणजे सोपाधिक विधानाच्या
पूर्वांगाशी जुळणारी
गृहीतके असतात.
कोणत्या नियमाने कोणती
गृहीतके मुक्त होतात हे
दाखवणे कधी महत्त्वाचे
असते; त्यासाठी मुक्तीची
खूण~$x$ वापरतात.
मात्र गृहीतके मुक्त
करणाऱ्या नियमाने योग्य
रूपाची \emph{सर्वच}
गृहीतके, किंवा एकही
गृहीतक, मुक्त करणे
\emph{बंधनकारक नाही}.
उदाहरणार्थ, पहिल्या
\Intro{\lif} ने एकही
गृहीतक मुक्त केलेले
नसले तरी पुढील सिद्धता
योग्य आहे:
\[
\AxiomC{$\Discharge{A}{y}$}
\DischargeRule{\Intro{\lif}}{x}
\UnaryInfC{$B \lif A$}
\DischargeRule{\Intro{\lif}}{y}
\UnaryInfC{$A \lif (B \lif A)$}
\DisplayProof
\]
नियमाने एकही गृहीतक
मुक्त केले नाही, तर
मुक्तीची खूण वगळतो
(येथेही तसे केले आहे).

पुढील सिद्धता पाहा:
\[
\AxiomC{$\Discharge{A}{x}$}
\AxiomC{$\Discharge{A}{y}$}
\RightLabel{\Intro{\land}}
\BinaryInfC{$A \land A$}
\RightLabel{\Elim{\land}}
\UnaryInfC{$A$}
\DischargeRule{\Intro{\lif}}{x}
\UnaryInfC{$A \lif A$}
\DischargeRule{\Intro{\lif}}{y}
\UnaryInfC{$A \lif (A \lif A)$}
\DisplayProof
\]
पहिले \Intro{\lif}
अनुमान डावीकडील
गृहीतक~$A^x$ मुक्त करते;
दुसरे उजवीकडील
गृहीतक~$A^y$ मुक्त करते.
म्हणजे खूण~$x$ असलेली
सर्व गृहीतके खूण~$x$
असलेल्या \Intro{\lif}
ने, आणि खूण~$y$
असलेली सर्व गृहीतके
खूण~$y$ असलेल्या
त्या अनुमानाने
मुक्त होतात.
यासाठी एकाच खूणेची
सर्व गृहीतके तेच
सूत्र असणे
महत्त्वाचे आहे.

संख्यापकांचे नियम
किंचित अधिक गुंतागुंतीचे आहेत.
उदाहरणार्थ,~\Log{N1i}
मध्ये $\lexists$
साठीचे नियम असे:
\[
\bottomAlignProof
\AxiomC{$A(t)$}
\RightLabel{\Intro{\lexists}}
\UnaryInfC{$\lexists[x][A(x)]$}
\DisplayProof
\qquad
\bottomAlignProof
\AxiomC{$\lexists[x][A(x)]$}
\AxiomC{$\Discharge{A(c)}{x}$}
\shortDeduce
\DeduceC{$B$}
\DischargeRule{\Elim{\lexists}}{x}
\BinaryInfC{$B$} \DisplayProof
\bottomAlignProof
\]
\Intro{\lexists} मध्ये
आधारविधान~$A(t)$ असते;
येथे $t$ हे बंद पद,
म्हणजे चल नसलेले पद आहे.
असे अनुमान योग्य,
म्हणजे \Intro{\lexists}
या नियम-रूपबंधाचे उदाहरण,
तेव्हाच असते जेव्हा
एखादे सूत्र~$A$,
चल~$x$ आणि बंद पद~$t$
असते ज्यामुळे निष्कर्ष
$\lexists[x][A]$
आणि आधारविधान
~$\Subst{A}{t}{x}$
असते. \Elim{\lexists}
नियमाने $A(c)$
रूपाची गृहीतके मुक्त
करता येतात: प्रत्येक
अनुमानासाठी असा एखादा
स्थिरांक~$c$ असतो
की $\Subst{A}{c}{x}$
रूपाची गृहीतके मुक्त
होतात. एका अनुमानाने
मुक्त होणाऱ्या सर्व
गृहीतकांत तोच~$c$
वापरला पाहिजे.
या स्थिरांकाला
\Elim{\lexists}
अनुमानाचे
\emph{आयगेन चर} म्हणतात.
~$c$ हे मुक्त होणाऱ्या
~$A(c)$ व्यतिरिक्त
इतर उघड्या गृहीतकात,
~$B$ मध्ये किंवा
~$A$ मध्ये येत
नसेल तेव्हाच अनुमान
योग्य असते. याला
\emph{आयगेन चराची अट}
म्हणतात.

~\Log{N1c} आणि~\Log{N1i}
मधील सिद्धता म्हणजे
नियम वापरून स्वयंसिद्धकांपासून
विगमनाने तयार होणारे
सूत्रांचे सांत वृक्ष.
प्रत्येक पान हे मुक्तीच्या
खुणेसह किंवा तिच्याविना
असलेले गृहीतक असते;
इतर प्रत्येक सूत्र
त्याच्या लगत वरच्या
सूत्रांपासून
पद्धतीच्या अनुमान नियमाने
निष्पन्न होते.
अनुमानाच्या वरच्या
वृक्षातील गृहीतक म्हणून
आलेले एखादे सूत्र
त्या अनुमानापर्यंत मुक्त
झालेले नसेल, तर त्याला
(त्या अनुमानाशी संबंधित)
\emph{उघडे} म्हणतात.

\begin{defn}\label{pt:nat:rp:defn:proof}
~\Log{N1c} आणि~\Log{N1i}
मधील \emph{सिद्धता}
हे पुढीलप्रमाणे विगमनाने
परिभाषित सूत्रे
चे सांत वृक्ष आहेत:
\begin{enumerate}
  \item\label{pt:nat:rp:defn:prf:assumption}
  मुक्तीची $x$ किंवा $y$
  सारखी खूण असलेले
  प्रत्येक सूत्र
  स्वतःच सिद्धता असते.
  एकाच सूत्र
  ची सिद्धता असेल,
  तर ते सूत्र
  उघडे गृहीतक मानतात.
  \item\label{pt:nat:rp:defn:prf:one-prem}
  $R$ हा एक, दोन किंवा
  तीन आधारविधाने
  $A_1$, $A_2$, $A_3$
  आणि निष्कर्ष~$A$
  असलेला नियम असू द्या.
  $\delta_1$, $\delta_2$,
  $\delta_3$ या अनुक्रमे
  $A_1$, $A_2$,
  $A_3$ मध्ये संपणाऱ्या
  सिद्धता असतील, तर
  पुढीलपैकी योग्य रचना
  अनुक्रमे सिद्धता असते:
  \[
  \AxiomC{}
  \RightLabel{$\delta_1$}
  \DeduceC{$A_1$}
  \RightLabel{$R$}
  \UnaryInfC{$A$}
  \DisplayProof
\qquad
  \AxiomC{}
  \RightLabel{$\delta_1$}
  \DeduceC{$A_1$}
  \AxiomC{}
  \RightLabel{$\delta_2$}
  \DeduceC{$A_2$}
  \RightLabel{$R$}
  \BinaryInfC{$A$}
  \DisplayProof
\qquad
  \AxiomC{}
  \RightLabel{$\delta_1$}
  \DeduceC{$A_1$}
  \AxiomC{}
  \RightLabel{$\delta_2$}
  \DeduceC{$A_2$}
  \AxiomC{}
  \RightLabel{$\delta_3$}
  \DeduceC{$A_3$}
  \RightLabel{$R$}
  \TrinaryInfC{$A$}
  \DisplayProof
  \]
  मात्र त्या सिद्धता
  मधील गृहीतक-खुणा
  सुसंगत असल्या पाहिजेत
  (दोन भिन्न
  सूत्रे ना
  तीच मुक्तीची खूण
  नसावी), आणि ही
  अनुमाने~$R$ चे
  योग्य उपयोग असले पाहिजेत.
  \item नव्या सिद्धता
  मधील सर्वात वरची
  सूत्र उदाहरणे
  ही त्या सिद्धता
  ची \emph{उघडी गृहीतके}
  मानली जातात, जर ती
  $\delta_1$, $\delta_2$
  किंवा~$\delta_3$
  मधील उघडी गृहीतके
  असतील आणि नियमाने
  मुक्त केलेली नसतील.
  नियमावर खूण~$x$
  (किंवा $x\, y$)
  असेल, तर \Intro{\lif}
  आणि~\Intro{\lnot}
  साठी $\delta_1$
  मधील खूण~$x$
  असलेली सर्व उघडी
  गृहीतके मुक्त होतात;
  \Elim{\lexists}
  साठी $\delta_2$
  मधील खूण~$x$
  असलेली सर्व उघडी
  गृहीतके मुक्त होतात;
  \Elim{\lor} साठी
  $\delta_2$ मधील
  खूण~$x$ असलेली
  आणि~$\delta_3$ मधील
  खूण~$y$ असलेली
  सर्व उघडी गृहीतके
  मुक्त होतात.
\end{enumerate}
\end{defn}

सिद्धता ची
विगमनात्मक रचना करताना
वापरलेल्या सिद्धता ना
तिच्या \emph{उप-सिद्धता}
म्हणतात. अनुमानाच्या
आधारविधानांत संपणाऱ्या
उप-सिद्धता ना त्या
अनुमानाच्या \emph{निकटच्या
उप-सिद्धता} म्हणतात.
\Log{N1c} आणि \Log{N1i}
यांचे नियम
\ref{pt:nat:rp:tab:N1} मध्ये आहेत.








\begin{table}
\begin{center}
\begin{tabular}[t]{@{}cc@{}}
\bottomAlignProof
\AxiomC{$\Discharge{A}{x}$}
\shortDeduce
\DeduceC{$B$}
\DischargeRule{\Intro{\lif}}{x}
\UnaryInfC{$A \lif B$}
\DisplayProof
&
\bottomAlignProof
\AxiomC{$A \lif B$}
\AxiomC{$A$}
\RightLabel{\Elim{\lif}}
\BinaryInfC{$B$}
\DisplayProof
\\[4ex]
\bottomAlignProof
\AxiomC{$A$}
\AxiomC{$B$}
\RightLabel{\Intro{\land}}
\BinaryInfC{$A \land B$}
\DisplayProof
&
\begin{tabular}[b]{@{}l@{}}
\AxiomC{$A \land B$}
\RightLabel{\Elim{\land}[1]}
\UnaryInfC{$A$}
\DisplayProof
\\[3ex]
\bottomAlignProof
\AxiomC{$A \land B$}
\RightLabel{\Elim{\land}[2]}
\UnaryInfC{$B$}
\DisplayProof
\end{tabular}
\\[4ex]
\begin{tabular}{@{}l@{}}
\AxiomC{$A$}
\RightLabel{\Intro{\lor}[1]}
\UnaryInfC{$A \lor B$}
\DisplayProof
\\[3ex]
\AxiomC{$B$}
\RightLabel{\Intro{\lor}[2]}
\UnaryInfC{$A \lor B$}
\DisplayProof
\end{tabular}
&
\AxiomC{$A \lor B$}
\AxiomC{$\Discharge{A}{x}$}
\shortDeduce
\DeduceC{$C$}
\AxiomC{$\Discharge{B}{y}$}
\shortDeduce
\DeduceC{$C$}
\DischargeRule{\Elim{\lor}}{x\,y}
\TrinaryInfC{$C$}
\DisplayProof
\bottomAlignProof
\\[6ex]
\bottomAlignProof
\AxiomC{$\Discharge{A}{x}$}
\shortDeduce
\DeduceC{$\lfalse$}
\DischargeRule{\Intro{\lnot}}{x}
\UnaryInfC{$\lnot A$}
\DisplayProof
&
\bottomAlignProof
\AxiomC{$\lnot A$}
\AxiomC{$A$}
\RightLabel{\Elim{\lnot}}
\BinaryInfC{$\lfalse$}
\DisplayProof
\\[4ex]
\bottomAlignProof
\AxiomC{$A(a)$}
\RightLabel{\Intro{\lforall}}
\UnaryInfC{$\lforall[x][A(x)]$}
\DisplayProof
&
\bottomAlignProof
\AxiomC{$\lforall[x][A(x)]$}
\RightLabel{\Elim{\lforall}}
\UnaryInfC{$A(t)$}
\DisplayProof
\\[4ex]
\bottomAlignProof
\AxiomC{$A(t)$}
\RightLabel{\Intro{\lexists}}
\UnaryInfC{$\lexists[x][A(x)]$}
\DisplayProof
&
\bottomAlignProof
\AxiomC{$\lexists[x][A(x)]$}
\AxiomC{$\Discharge{A(a)}{x}$}
\shortDeduce
\DeduceC{$C$}
\DischargeRule{\Elim{\lexists}}{x}
\BinaryInfC{$C$}
\DisplayProof
\\[4ex]
\bottomAlignProof
\AxiomC{$\lfalse$}
\RightLabel{\FalseInt}
\UnaryInfC{$A$}
\DisplayProof
&
\bottomAlignProof
\AxiomC{$\Discharge{\lnot A}{x}$}
\shortDeduce
\DeduceC{$\lfalse$}
\DischargeRule{\FalseCl}{x}
\UnaryInfC{$A$}
\DisplayProof
\end{tabular}
\end{center}
\caption{\Log{N1i} आणि \Log{N1c}
यांचे नियम. $\Intro{\lforall}$
मध्ये स्थिरांक~$a$
निष्कर्षात किंवा उघड्या
गृहीतकात येता कामा नये.
$\Elim{\lexists}$ मध्ये
$a$ दोन्ही आधारविधानांच्या
निष्कर्षांत किंवा अनुमानानंतर
उरलेल्या उघड्या गृहीतकांत
येता कामा नये.
\Log{N1i} मध्ये \FalseCl
नसतो.}
\label{pt:nat:rp:tab:N1}
\end{table}




\begin{defn}
सिद्धता~$\delta$ ची
\emph{उंची}~
$\pheight{\delta}$
पुढीलप्रमाणे विगमनाने
परिभाषित करतात.
\begin{enumerate}
  \item $\delta$ केवळ
  गृहीतकाची असेल, तर
  $\pheight{\delta} = 0$.
  \item $\delta$ चा शेवट
  एका आधारविधानाच्या
  अनुमानात होत असेल
  आणि निकटची
  उप-सिद्धता~$\delta_1$
  असेल, तर
  $\pheight{\delta} =
  \pheight{\delta_1} + 1$.
  \item $\delta$ चा शेवट
  दोन आधारविधानांच्या
  अनुमानात होत असेल
  आणि निकटच्या
  उप-सिद्धता~$\delta_1$
  आणि~$\delta_2$ असतील,
  तर $\pheight{\delta} =
  \max(\pheight{\delta_1},
  \pheight{\delta_2}) + 1$.
  \item $\delta$ चा शेवट
  \Elim{\lor} मध्ये होत
  असेल आणि निकटच्या
  उप-सिद्धता~$\delta_1$,
  $\delta_2$ आणि~$\delta_3$
  असतील, तर
  $\pheight{\delta} =
  \max(\pheight{\delta_1},
  \pheight{\delta_2},
  \pheight{\delta_3}) + 1$.
\end{enumerate}
क्रमवर्ती~$S$ ची उंची
$\le n$ असलेली
सिद्धता असेल,
तर $ \Proves[n] S$
असे लिहितो.
\end{defn}

\begin{ex}
~\Log{N1i} मध्ये कोणतेही
सूत्र हे स्वतःपासून
उघडे गृहीतक घेऊन
सिद्धता असते.
म्हणून
\[
A \land B^x
\]
ही सिद्धता~$\delta_1$
आहे; तिची उंची~$0$.
~$\delta_1$ वर
~\Elim{\land} च्या
दोनपैकी एक आवृत्ती
लावून सिद्धता~$\delta_2$
आणि~$\delta_3$ मिळतात:
\[
\AxiomC{$A \land B^x$}
\RightLabel{\Elim{\land}[2]}
\UnaryInfC{$B$}
\DisplayProof
\qquad
\AxiomC{$A \land B^x$}
\RightLabel{\Elim{\land}[1]}
\UnaryInfC{$A$}
\DisplayProof
\]
~$\delta_2$ आणि
~$\delta_3$ मधील
अंतिम अनुमानांच्या निकटच्या
उप-सिद्धता सारख्याच:
$A \land B$ या
गृहीतकापासून तयार झालेल्या.
दोन्ही सिद्धता मध्ये
$A\land B$ हे
उघडे गृहीतक म्हणून येते.
आपल्याला
$\pheight{\delta_2} =
\pheight{\delta_3} = 1$
मिळते.

~\Intro{\land} नियमात
$B$ आणि $A$
या रूपाची दोन
आधारविधाने असतील,
तर $B \land A$
या रूपाचा निष्कर्ष
निघतो. त्यामुळे
पुढील सिद्धता~$\delta_4$
आहे:
\[
\AxiomC{$A \land B^x$}
\RightLabel{\Elim{\land}[2]}
\UnaryInfC{$B$}
\LeftSubproofLabel{$\delta_2$}
\AxiomC{$A \land B^x$}
\RightLabel{\Elim{\land}[1]}
\UnaryInfC{$A$}
\RightSubproofLabel{$\delta_3$}
\RightLabel{\Intro{\land}}
\BinaryInfC{$B \land A$}
\DisplayProof
\]
तिची उंची
$\pheight{\delta_4} =
\max(\pheight{\delta_2},
\pheight{\delta_3})+1
= \max(1,1) + 1 = 2$.
तिच्यात $A \land B$
या गृहीतकाची दोन
उदाहरणे आहेत; दोन्ही
उघडी आहेत.

शेवटी~\Intro{\lif}
लावून $\delta_5$
मिळवू:
\[
\AxiomC{$\Discharge{A \land B}{x}$}
\RightLabel{\Elim{\land}[2]}
\UnaryInfC{$B$}
\LeftSubproofLabel{$\delta_2$}
\AxiomC{$\Discharge{A \land B}{x}$}
\RightLabel{\Elim{\land}[1]}
\UnaryInfC{$A$}
\RightSubproofLabel{$\delta_3$}
\RightLabel{\Intro{\land}}
\BinaryInfC{$B \land A$}
\DischargeRule{\Intro{\lif}}{x}
\UnaryInfC{$(A \land B) \lif (B \land A)$}
\DisplayProof
\]
तिची उंची
$\pheight{\delta_4} + 1 = 3$.
\Intro{\lif} अनुमान
$A \land B$
या रूपाची, खूण~$x$
असलेली सर्व गृहीतके
मुक्त करते; म्हणजे
निकटच्या उपसिद्धतेतील
आधी उघडी असलेली
दोन्ही गृहीतके.
~$\delta_5$ ची
तीच दोन गृहीतके असल्याने
तिच्यात उघडे गृहीतक
उरत नाही; म्हणून ती
तिच्या अंतिम-सूत्र
$(A \land B)
\lif (B \land A)$
\emph{ची} सिद्धता
आहे.
\end{ex}












\section{क्रमवर्तींसह नैसर्गिक निगमन: \Log{N2c} आणि~\Log{N2i}}\label{pt:nat:seq:sec}

~\Log{N1c} किंवा
\Log{N1i} मधील
सिद्धता~$\delta$
तिचे अंतिम-सूत्र~$A$
हे तिच्या उघड्या गृहीतकांपासून
निष्पन्न होते हे दाखवते.
~$\delta$ मध्ये $B^x$
हे उघडे गृहीतक असलेल्या
सूत्रे~$B$
यांचा संच~$\Gamma$
असेल, तर हे
$\delta \Proves \Gamma
\Sequent A$ असे
लिहिता येते.
यावरून क्रमवर्तींसाठीही
नैसर्गिक निगमन पद्धत
बांधता येईल असे दिसते.
अशा पद्धतीत प्रत्येक
क्रमवर्ती तिच्या उजवीकडील
सूत्र कोणत्या
उघड्या गृहीतकांवर अवलंबून
आहे, म्हणजे त्या
क्रमवर्तीत संपणाऱ्या
उप-सिद्धता मध्ये
कोणती गृहीतके उघडी
आहेत, हे सांगते.
यालाच क्रमवर्तींसह,
किंवा “क्रमवर्ती-पद्धतीचे”
नैसर्गिक निगमन म्हणतो.
मिळणाऱ्या पद्धती
\Log{N2c} आणि~\Log{N2i}
अशा आहेत.

~\Log{N1} आणि
\Log{N2} यांतील
संबंध अधिक निकट करण्यासाठी
डावीकडील संच~$\Gamma$
हे खूण लावलेल्या
सूत्रे~$x:B$
यांचे धरू. ~\Log{N1}
चे नियम केवळ
आधारविधानांतील
सूत्रे वर,
म्हणजे निकटच्या
उप-सिद्धता च्या
अंतिम-सूत्रे वर
लागतात. त्यामुळे
\Log{N2} चे नियम
क्रमवर्तीच्या उजव्या
भागावरच लागतात.
क्रमवर्तींच्या डावीकडील
खूण लावलेल्या सूत्रांना
म्हणून फक्त संदर्भ
म्हणता येते.

~\Log{N2} सिद्धतांत
गृहीतकांच्या जागी
सर्वात वर अशा रूपाची
स्वयंसिद्धक-क्रमवर्ती असते:
\[x:A \Sequent A.\]
नियम~\Log{N1}
च्या नियमांशी तंतोतंत
जुळतात; ते
\ref{pt:nat:seq:tab:N2} मध्ये आहेत.








\begin{table}
\begin{tabular}{@{}c@{}c@{}}
\Axiom$x: A, \Gamma \fCenter B$
\RightLabel{\Intro{\lif}}
\UnaryInf$\Gamma \fCenter A \lif B $
\DisplayProof
&
\Axiom$\Gamma_1 \fCenter A \lif B$
\Axiom$\Gamma_2 \fCenter A$
\RightLabel{\Elim{\lif}}
\BinaryInf$\Gamma_1, \Gamma_2 \fCenter B$
\DisplayProof
\\[4ex]
\Axiom$\Gamma_1 \fCenter A$
\Axiom$\Gamma_2 \fCenter B$
\RightLabel{\Intro{\land}}
\BinaryInf$\Gamma_1, \Gamma_2 \fCenter A \land B $
\DisplayProof
&
\begin{tabular}[b]{@{}l@{}}
\Axiom$\Gamma \fCenter A \land B$
\RightLabel{\Elim{\land}[1]}
\UnaryInf$\Gamma \fCenter A$
\DisplayProof
\\[3ex]
\Axiom$\Gamma \fCenter A \land B$
\RightLabel{\Elim{\land}[2]}
\UnaryInf$\Gamma \fCenter B$
\DisplayProof\\[3ex]
\end{tabular}
\\[4ex]
\Axiom$\Gamma \fCenter A$
\RightLabel{\Intro{\lor}[1]}
\UnaryInf$\Gamma \fCenter A \lor B$
\DisplayProof
&
\Axiom$ \Gamma \fCenter B$
\RightLabel{\Intro{\lor}[2]}
\UnaryInf$\Gamma \fCenter A \lor B$
\DisplayProof\\[3ex]
\multicolumn{2}{c}{
\begin{tabular}[b]{@{}l@{}}
\Axiom$\Gamma_1 \fCenter A \lor B$
\Axiom$x:A, \Gamma_2 \fCenter C$
\Axiom$y:B, \Gamma_3 \fCenter C$
\RightLabel{\Elim{\lor}}
\TrinaryInf$\Gamma_1, \Gamma_2, \Gamma_3 \fCenter C$
\DisplayProof
\end{tabular}
}
\\[4ex]
\Axiom$x:A, \Gamma \fCenter \lfalse$
\RightLabel{\Intro{\lnot}}
\UnaryInf$ \Gamma \fCenter \lnot A $
\DisplayProof
&
\Axiom$\Gamma_1 \fCenter \lnot A $
\Axiom$\Gamma_2 \fCenter A $
\RightLabel{\Elim{\lnot}}
\BinaryInf$\Gamma_1, \Gamma_2 \fCenter \lfalse$
\DisplayProof
\\[4ex]
\Axiom$\Gamma \fCenter A(a) $
\RightLabel{\Intro{\lforall}}
\UnaryInf$\Gamma \fCenter \lforall[x][A(x)]$
\DisplayProof
&
\Axiom$\Gamma \fCenter \lforall[x][A(x)]$
\RightLabel{\Elim{\lforall}}
\UnaryInf$\Gamma \fCenter A(t)$
\DisplayProof
\\[4ex]
\Axiom$\Gamma \fCenter A(t) $
\RightLabel{\Intro{\lexists}}
\UnaryInf$\Gamma \fCenter \lexists[x][A(x)]$
\DisplayProof
&
\Axiom$\Gamma_1 \fCenter  \lexists[x][A(x)]$
\Axiom$x: A(a), \Gamma_2 \fCenter C$
\RightLabel{\Elim{\lexists}}
\BinaryInf$\Gamma_1, \Gamma_2 \fCenter C$
\DisplayProof
\\[4ex]
\Axiom$\Gamma \fCenter \lfalse$
\RightLabel{\FalseInt}
\UnaryInf$\Gamma \fCenter A$
\DisplayProof
&
\Axiom$x: \lnot A, \Gamma \fCenter \lfalse$
\RightLabel{\FalseCl}
\UnaryInf$\Gamma \fCenter A$
\DisplayProof
\end{tabular}
\caption{\Log{N2c} आणि \Log{N2i}
यांचे नियम. $\Intro{\forall}$
आणि $\Elim{\exists}$ मध्ये
स्थिरांक~$a$
निष्कर्ष-क्रमवर्तीमध्ये
येता कामा नये.
$\Gamma_i$ हे खूण
लावलेल्या सूत्रे~$x :
A$ चे संच आहेत.
\Intro{\lforall} आणि
\Elim{\lexists} नियमांसाठी
आयगेन चराची अट पूर्ण
झाली पाहिजे:
$a$ निष्कर्ष-क्रमवर्तीमध्ये
येता कामा नये.
\Log{N2i} मध्ये \FalseCl
नसतो.}
\label{pt:nat:seq:tab:N2}
\end{table}




~\Log{N1c} आणि
\Log{N1i} मध्ये
\Intro{\lif},
\Intro{\lnot},
\Elim{\lor} आणि
\Elim{\lexists} हे
नियम गृहीतके मुक्त करू
शकतात, पण तसे करणे
बंधनकारक नाही.
म्हणून~\Log{N2c}
आणि~\Log{N2i}
मध्येही संबंधित आधारविधानाच्या
संदर्भात अपेक्षित
सूत्रे $A$,
$B$ आणि $A(a)$
नसली तरी त्या
नियमांचा योग्य उपयोग
मान्य करतो.

\begin{prop}
$\delta$ ही~\Log{N1c}
किंवा \Log{N1i}
मध्ये~$A$ ची
सिद्धता असेल,
तर~\Log{N2c}
(\Log{N2i}) मध्ये
$\Gamma \Sequent A$
ची सिद्धता~$\delta'$
असते. येथे~$\Gamma$
हा असा सर्व~$x:B$
यांचा संच आहे की
$B^x$ हे~$\delta$
चे उघडे गृहीतक असते.
\end{prop}

\begin{proof}
  $n=\pheight{\delta}$
  वर विगमन करू.
  $n=0$ असेल, तर
  $\delta$ ही एकच
  उघडे गृहीतक~$A^x$.
  मग $\Gamma = \{x:A\}$
  आणि $\Gamma
  \Sequent A$ हे
  \Log{N2c} किंवा
  ~\Log{N2i} चे
  स्वयंसिद्धक आहे.

  $n>0$ असेल, तर
  ~$\delta$ मधील अंतिम
  अनुमानानुसार प्रसंग
  पाडू. येथे एकच प्रसंग
  पाहू; उरलेले सरावासाठी.

  अंतिम अनुमान
  $\Intro{\lif}$ असल्याचे
  समजा. मग~$\delta$
  चा शेवट असा होतो:
  \[
  \AxiomC{$\Discharge{B}{x}$}
  \RightLabel{$\delta_1$}
  \DeduceC{$C$}
  \DischargeRule{\Intro{\lif}}{x}
  \UnaryInfC{$B \lif C$}
  \DisplayProof
  \]
  विगमन गृहीतकानुसार
  ~\Log{N2c} (\Log{N2i})
  मध्ये
  $\Gamma' \Sequent C$
  ची सिद्धता~$\delta_1'$
  आहे; $\Gamma'$ म्हणजे
  ~$\delta_1$ मधील
  खूण लावलेली उघडी
  गृहीतके. $B^x$
  हे $\delta_1$
  चे उघडे गृहीतक असेल,
  तर \Intro{\lif}
  अनुमानात ते मुक्त होते,
  आणि~$\delta$ ची
  उघडी गृहीतके
  $\Gamma =
  \Gamma' \setminus \{x:B\}$
  होतात. म्हणजे
  ~$\delta_1'$ ही
  $x:B, \Gamma \Sequent C$
  ची सिद्धता
  असते; मग
  $\delta'$ पुढीलप्रमाणे
  घेता येते:
  \[
  \AxiomC{}
  \RightLabel{$\delta_1'$}
  \Deduce$x:B, \Gamma \fCenter C$
  \RightLabel{\Intro{\lif}}
  \UnaryInf$\Gamma \fCenter B \lif C$
  \DisplayProof
  \]
  $B^x$ हे~$\delta_1$
  चे उघडे गृहीतक
  नसेल, तर
  \Intro{\lif}
  कोणतेही गृहीतक
  मुक्त करत नाही;
  ~$\delta$ आणि
  $\delta_1$ यांची
  उघडी गृहीतके तीच
  राहतात. ~\Log{N2c}
  आणि~\Log{N2i}
  मध्ये \Intro{\lif}
  च्या आधारविधानाच्या
  संदर्भात $x:B$
  असणे आवश्यक नसल्याने
  $\delta'$ पुढीलप्रमाणे
  घेता येते:
  \[
  \AxiomC{}
  \RightLabel{$\delta_1'$}
  \Deduce$\Gamma \fCenter C$
  \RightLabel{\Intro{\lif}}
  \UnaryInf$\Gamma \fCenter B \lif C$
  \DisplayProof
  \]

  ~$\delta$ चा शेवट
  इतर नियमांत होणारे
  प्रसंगही असेच आहेत.
\end{proof}

\begin{prop}
$\delta$ ही \Log{N2c}
किंवा \Log{N2i} मध्ये
$\Gamma \Sequent A$
ची सिद्धता असेल,
तर~\Log{N1c}
(\Log{N1i}) मध्ये
अंतिम-सूत्र~$A$
आणि प्रत्येक
$x:B \in \Gamma$
साठी उघडे गृहीतक~$B^x$
असलेली सिद्धता~$\delta'$
मिळते.
\end{prop}

\begin{proof}
  पुन्हा~$\delta$
  ची उंची~$n$
  यावर विगमन करू.
  $n=0$ असेल, तर
  $\delta$ हे
  $x:A \Sequent A$
  स्वयंसिद्धक असते.
  सिद्धता~$\delta'$
  म्हणजे गृहीतक~$A^x$.

  $n>0$ असेल,
  तर~$\delta$ मधील
  अंतिम अनुमानानुसार
  प्रसंग पाडू.
  उदाहरणार्थ, ते
  ~$\Intro{\lif}$
  असेल, तर~$\delta$
  चा शेवट असा:
  \[
  \bottomAlignProof
  \AxiomC{}
  \RightLabel{$\delta_1$}
  \Deduce$x:B, \Gamma \fCenter C$
  \RightLabel{\Intro{\lif}}
  \UnaryInf$\Gamma \fCenter B \lif C$
  \DisplayProof
  \quad\text{ किंवा }\quad
  \bottomAlignProof
  \AxiomC{}
  \RightLabel{$\delta_1$}
  \Deduce$\Gamma \fCenter C$
  \RightLabel{\Intro{\lif}}
  \UnaryInf$\Gamma \fCenter B \lif C$
  \DisplayProof
  \]
  विगमन गृहीतकानुसार
  ~\Log{N1c} (\Log{N1i})
  मध्ये~$C$ ची
  सिद्धता मिळते.
  पहिल्या प्रसंगी
  $B^x$ हे~$\delta_1'$
  चे उघडे गृहीतक असते;
  दुसऱ्या प्रसंगी नसते.
  मग सिद्धता~$\delta'$
  अनुक्रमे पुढीलप्रमाणे
  घेता येते:
  \[
  \bottomAlignProof
  \AxiomC{$\Discharge{B}{x}$}
  \RightLabel{$\delta_1'$}
  \DeduceC{$C$}
  \DischargeRule{\Intro{\lif}}{x}
  \UnaryInfC{$B \lif C$}
  \DisplayProof
  \quad\text{ किंवा }\quad
  \bottomAlignProof
  \AxiomC{}
  \RightLabel{$\delta_1'$}
  \DeduceC{$C$}
  \RightLabel{\Intro{\lif}}
  \UnaryInfC{$B \lif C$}
  \DisplayProof
  \]
\end{proof}













\section{नियमित सिद्धता आणि आदेशन}\label{pt:nat:qua:sec}

संख्यापकांच्या नियमांमध्ये
~\Elim{\lexists} आणि
~\Intro{\lforall} यांना
लागणाऱ्या “आयगेन चराच्या
अटींमुळे” काही गुंतागुंत येते.
या अनुमानांतील
स्थिरांक~$c$ पुन्हा
वापरला जाणार नाही याची
खात्री केल्यास बहुतांश
गुंतागुंत हाताळता येते.
पुढील व्याख्या करू.

\begin{defn}
नैसर्गिक निगमनातील
सिद्धता~$\delta$
\emph{नियमित} असते, जर
\begin{enumerate}
  \item कोणतेही आयगेन
  चर~$a$ एकाहून अधिक
  \Elim{\lexists} किंवा
  ~\Intro{\lforall} अनुमानांचे
  आयगेन चर म्हणून वापरलेले
  नसेल; आणि
  \item $a$ हे
  \Elim{\lexists} किंवा
  ~\Intro{\lforall} अनुमानाचे
  आयगेन चर असेल, तर ते
  अनुक्रमे गौण आधारविधानाकडे
  किंवा आधारविधानाकडे
  जाणाऱ्या निकटच्या
  उप-सिद्धता मध्येच
  येत असेल.
\end{enumerate}
\end{defn}

\begin{lem}\label{pt:nat:qua:lem:subst-regular}
$\delta$ ही $C$ मध्ये
संपणारी नियमित सिद्धता असू द्या.
$c$ हा~$\delta$ मध्ये
आयगेन चर म्हणून न वापरलेला
स्थिरांक असू द्या;
आणि~$t$ हे $\delta$
चे कोणतेही आयगेन चर नसलेले
बंद पद असू द्या. या प्रकरणातील संख्यापक-नियमांत
वापरलेल्या बंद-पदाच्या अटीनुसार ही मर्यादा ठेवली आहे. मग~$\delta$
मधील $c$ च्या प्रत्येक
उपस्थितीच्या जागी $t$
बसवून मिळणारी
$\Subst{\delta}{t}{c}$
ही नियमित सिद्धता आहे.
$\Subst{\delta}{t}{c}$
चे अंतिम-सूत्र
$\Subst{C}{t}{c}$ असते.
$A^x$ हे~$\delta$
चे उघडे गृहीतक असेल,
तर $\Subst{A}{t}{c}^x$
हे $\Subst{\delta}{t}{c}$
चे उघडे गृहीतक असते.
\end{lem}

\begin{proof}
  $\pheight{\delta}$ वर
  विगमन करू.
  $\pheight{\delta} = 0$
  असेल, तर सिद्धतेत केवळ
  गृहीतक~$A^x$ असते.
  म्हणून $\Subst{\delta}{t}{c}$
  हीच~$\Subst{A}{t}{c}^x$
  असते.

  $\pheight{\delta} > 0$
  असेल, तर~$\delta$
  मधील अंतिम अनुमानानुसार
  प्रसंग पाडू.
  विधानीय संयोजकांच्या
  नियमांचे प्रसंग सोपे
  असून सरावासाठी ठेवले आहेत.
  $\delta$ चा शेवट
  संख्यापक अनुमानात होतो
  ते प्रसंग पाहू.
  \begin{enumerate}
    \item अंतिम अनुमान
    $\Intro{\lforall}$
    असेल, तर~$\delta$
    चे रूप असे आहे:
    \[
    \AxiomC{}
    \RightLabel{$\delta'$}
    \DeduceC{$A(a,c)$}
    \RightLabel{\Intro{\lforall}}
    \UnaryInfC{$\lforall[x][A(x,c)]$}
    \DisplayProof
    \]
    विगमन गृहीतकानुसार
    $\Subst{\delta'}{t}{c}$
    ही $A(a,t)$ ची
    नियमित सिद्धता आहे.
    $a$ हे~$\delta$
    चे आयगेन चर असल्याने
    $a$ हे~$t$ मध्ये येत नाही.
    म्हणून
    $\Subst{\delta}{t}{c}$
    मधील अंतिम अनुमान
    \[
    \AxiomC{}
    \RightLabel{$\Subst{\delta'}{t}{c}$}
    \DeduceC{$A(a,t)$}
    \RightLabel{\Intro{\lforall}}
    \UnaryInfC{$\lforall[x][A(x,t)]$}
    \DisplayProof
    \]
    योग्य आहे: $t$
    मध्ये~$a$ नसल्याने
    निष्कर्षात किंवा कोणत्याही
    उघड्या गृहीतकात~$a$
    येत नाही.
    \item अंतिम अनुमान
    $\Elim{\lforall}$
    असेल, तर~$\delta$
    चे रूप असे आहे:
    \[
    \AxiomC{}
    \RightLabel{$\delta'$}
    \DeduceC{$\lforall[x][A(x,c)]$}
    \RightLabel{\Elim{\lforall}}
    \UnaryInfC{$A(s(c),c)$}
    \DisplayProof
    \]
    विगमन गृहीतकानुसार
    $\Subst{\delta'}{t}{c}$
    ही $\lforall[x][A(x,t)]$
    ची नियमित सिद्धता
    आहे. (निष्कर्षातील
    पद~$s(c)$ मध्ये~$c$
    असूही शकतो किंवा नसूही
    शकतो.) आता
    $\Subst{\delta}{t}{c}$
    मधील अंतिम अनुमान
    \[
    \AxiomC{}
    \RightLabel{$\Subst{\delta'}{t}{c}$}
    \DeduceC{$\lforall[x][A(x,t)]$}
    \RightLabel{\Elim{\lforall}}
    \UnaryInfC{$A(s(t),t)$}
    \DisplayProof
    \]
    योग्य आहे, आणि
    $A(s(t),t) \ident
    \Subst{A(s(c),c)}{t}{c}$.
    \Intro{\lexists} चा
    प्रसंगही असाच आहे.
  \item अंतिम अनुमान
    $\Elim{\lexists}$
    असेल, तर~$\delta$
    चे रूप असे आहे:
    \[
    \AxiomC{}
    \RightLabel{$\delta_1'$}
    \DeduceC{$\lexists[x][A(x,c)]$}
    \AxiomC{$\Discharge{A(a,c)}{x}$}
    \RightLabel{$\delta_2'$}
    \DeduceC{$C$}
    \DischargeRule{\Elim{\lexists}}{x}
    \BinaryInfC{$C$}
    \DisplayProof
    \]
    विगमन गृहीतकानुसार
    $\Subst{\delta_1'}{t}{c}$
    आणि
    $\Subst{\delta_2'}{t}{c}$
    या अनुक्रमे
    $\lexists[x][A(x,t)]$
    ची आणि उघड्या
    गृहीतक~$A(a,t)$ पासून
    $\Subst{C}{t}{c}$
    ची नियमित सिद्धता
    आहेत. $a$ हे~$\delta$
    चे आयगेन चर असल्याने
    $a$ हे~$t$ मध्ये येत नाही.
    म्हणून
    $\Subst{\delta}{t}{c}$
    मधील अंतिम अनुमान
    \[
    \AxiomC{}
    \RightLabel{$\Subst{\delta_1'}{t}{c}$}
    \DeduceC{$\lexists[x][A(x,t)]$}
    \AxiomC{$\Discharge{A(a,t)}{x}$}
    \RightLabel{$\Subst{\delta_2'}{t}{c}$}
    \DeduceC{$\Subst{C}{t}{c}$}
    \DischargeRule{\Elim{\lexists}}{x}
    \BinaryInfC{$\Subst{C}{t}{c}$}
    \DisplayProof
    \]
    योग्य आहे:
    $\Subst{A(x,t)}{a}{x}
    \ident \Subst{A(a,c)}{t}{c}$;
    निष्कर्ष~$\Subst{C}{t}{c}$
    किंवा कोणत्याही उघड्या
    गृहीतकात~$a$ येत नाही.
  \end{enumerate}
\end{proof}

\begin{prop}\label{pt:nat:qua:prop:regular}
$\delta$ ही \Log{N1c}
किंवा \Log{N1i} मधील
सिद्धता असेल,
तर त्याच रचनेची
(विशेषतः त्याच
उंची ची)
नियमित सिद्धता~$\delta'$
असते; ती~$\delta$
पेक्षा आयगेन चरांची आणि आवश्यक असल्यास
बद्ध गृहीतक-खुणांची नावे बदलल्यामुळे भिन्न असते;
उघडी गृहीतके आणि त्यांची नावे जपली जातात.
\end{prop}

\begin{proof}
या सिद्धतेपुरते,
~$\delta$ मधील एखादे
आयगेन-चर अनुमान
\emph{स्वच्छ} म्हणू,
जर त्याचे आयगेन चर
इतर कोणत्याही अनुमानाचे
आयगेन चर नसेल आणि
त्या अनुमानाच्या निकटच्या
उप-सिद्धता व्यतिरिक्त
~$\delta$ मध्ये
कुठेही येत नसेल;
अन्यथा ते \emph{अस्वच्छ}
म्हणू. ~$n$ ही~$\delta$
मधील अस्वच्छ आयगेन-चर
अनुमानांची संख्या असू द्या.
~$n$ वर विगमन करून
~$\delta$ चे अपेक्षित
नियमित सिद्धता~$\delta'$
मध्ये रूपांतर करता येते,
हे दाखवू. $n=0$ असेल,
तर~$\delta$ मधील सर्व
आयगेन-चर अनुमाने स्वच्छ
आहेत; म्हणून~$\delta$
नियमित आहे आणि आणखी
काही सिद्ध करायचे नाही.
अन्यथा सर्वात वरचे
अस्वच्छ आयगेन-चर अनुमान
निवडा; उदाहरणार्थ
\Intro{\lforall}.
त्यात संपणारी~$\delta$
ची उप-सिद्धता~$\delta_1$
या रूपाची आहे:
    \[
    \AxiomC{}
    \RightLabel{$\delta_2$}
    \DeduceC{$A(c)$}
    \RightLabel{\Intro{\lforall}}
    \UnaryInfC{$\lforall[x][A(x)]$}
    \DisplayProof
    \]
~$c$ हे आयगेन चर असल्याने
ते~$\lforall[x][A(x)]$
मध्ये किंवा उघड्या गृहीतकात
येत नाही. ~$\delta_2$
नियमित आहे: निवडलेल्या
अनुमानाच्या वर असलेली
सर्व आयगेन-चर अनुमाने
स्वच्छ आहेत. ~$a$ हा
~$\delta$ मध्ये कुठेही
न येणारा स्थिरांक
असू द्या. आदेशनापूर्वी $\delta_2$ मधील
बद्ध गृहीतक-खुणा व त्यांना बांधलेली गृहीतके एकत्र
नवी नावे देऊन उर्वरित $\delta$ च्या खुणांपासून
वेगळी करा; उघड्या गृहीतकांच्या खुणा बदलू नका.
त्यामुळे आदेशनाने बदललेले एखादे अंतर्गत गृहीतक
बाहेरच्या वेगळ्या सूत्राबरोबर तीच खूण वापरणार नाही.
\ref{pt:nat:qua:lem:subst-regular}
नुसार
$\Subst{\delta_2}{a}{c}$
ही $A(a)$ ची नियमित
सिद्धता आहे.
आयगेन चराच्या अटीने
~$\delta_2$ च्या कोणत्याही
उघड्या गृहीतकात~$c$
नसतो; म्हणून
~$\Subst{\delta_2}{a}{c}$
च्या उघड्या गृहीतकात
~$a$ नसतो. त्यामुळे
~$\Intro{\lforall}$
लावता येतो.
~$\delta$ मधील
$\delta_1$ च्या जागी
पुढील सिद्धता~$\delta_1'$
ठेवा:
    \[
    \AxiomC{}
    \RightLabel{$\Subst{\delta_2}{a}{c}$}
    \DeduceC{$A(a)$}
    \RightLabel{\Intro{\lforall}}
    \UnaryInfC{$\lforall[x][A(x)]$}
    \DisplayProof
    \]
यात एक अस्वच्छ आयगेन-चर
अनुमान स्वच्छ केले आहे,
आयगेन चर~$c$ च्या जागी~$a$ आले आहे,
आणि आवश्यक बद्ध गृहीतक-खुणा नवी नावे देऊन जपल्या आहेत.
परिणामी सिद्धतेत
जास्तीत जास्त $n-1$ अस्वच्छ आयगेन-चर
अनुमाने राहतात; जुने चर इतरत्र न उरल्यास
आणखी अनुमानेही स्वच्छ होऊ शकतात. म्हणून
विगमन गृहीतकाने अपेक्षित
निष्कर्ष मिळतो.
\end{proof}

\begin{prob}
\ref{pt:nat:qua:prop:regular} च्या
सिद्धतेत सर्वात वरचे
आयगेन-चर अनुमान
~\Elim{\lexists} असते
तो प्रसंग वर्णन करा.
\end{prob}

नुकतेच सिद्ध केलेले
विधान आपण पुढे स्पष्टपणे
वापरणार नाही; पण चर्चेतील
सर्व सिद्धता नियमित
आहेत असे मानू.
तसे नसल्यास आयगेन
चरांची नावे बदलून त्या
नियमित करता येतात.

\ref{pt:nat:qua:lem:subst-regular}
आणि \ref{pt:nat:qua:prop:regular}
ही विधाने \Log{N2c}
आणि~\Log{N2i}
यांनाही लागू होतात.
त्यांच्या सिद्धता
सरावासाठी ठेवतो.













\section{सिद्धता यांचे कलम-संयोजन}\label{pt:nat:gra:sec}

~\Log{N1c} आणि~\Log{N1i}
मधील सिद्धता म्हणजे
गृहीतके~$B$ यांपासून
सूत्रे~$A$ यांच्या
सिद्धता होत. आता $B$
चीही सिद्धता आहे असे समजा.
$B$ ची सिद्धता आणि
$B$ पासून $A$ ची सिद्धता
एकत्र करून, $B$ हे गृहीतक
न वापरणारी $A$ ची
सिद्धता कशी मिळवता येईल?

एक मार्ग म्हणजे
~\Intro{\lif} लावून $A$
च्या सिद्धता मधील
$B$ हे गृहीतक काढून टाकणे.
मग मिळालेली $B \lif A$
ची सिद्धता आणि $B$
ची सिद्धता एकत्र
केल्यास पुढील सिद्धता मिळते:
\[
\AxiomC{$\Discharge{B}{x}$}
\DeduceC{$A$}
\DischargeRule{\Intro{\lif}}{x}
\UnaryInfC{$B \lif A$}
\AxiomC{}
\DeduceC{$B$}
\RightLabel{\Elim{\lif}}
\BinaryInfC{$A$}
\DisplayProof
\]
हे सरळ आहे, पण काहीसे
अनावश्यक वळण घेते:
$\lif$ आणून लगेच त्याचे
निराकरण का करावे?
$B \lif A$ मधून जाणारा
असा “वळसा” टाळता यायला हवा.
(असे वळसे नेहमी टाळता
येतात, हाच प्रत्यक्षात
सामान्यरूपीकरण प्रमेयाचा आशय आहे.)

~\Log{N1c} आणि~\Log{N1i}
मध्ये या प्रसंगी तो वळसा
सहज टाळता येतो: $B^x$
या गृहीतकांच्या जागी $B$
ची संपूर्ण सिद्धता बसवता येते.
इतर सिद्धतांच्या गृहीतकांवर
सिद्धता चे असे
“कलम-संयोजन” उपयुक्त आहे;
म्हणून त्याची व्याख्या नोंदवू.

\begin{defn}
समजा $\delta_1$ ही उघडे
गृहीतक~$B^x$ यापासून
$A$ ची \Log{N1c}-सिद्धता
(किंवा \Log{N1i}-सिद्धता) आहे,
आणि $\delta_2$ ही $B$
ची \Log{N1c}-सिद्धता
(\Log{N1i}-सिद्धता) आहे.
मग $\delta_1$ मधील
$B^x$ च्या प्रत्येक
निरसन न झालेल्या उदाहरणाच्या
जागी $\delta_2$ बसवून
मिळणारा परिणाम म्हणजे
$\Subst{\delta_1}{\delta_2}{B^x}$.
\end{defn}

~$\delta_1$ आणि $\delta_2$
मध्ये एकाच खूणेची, पण वेगवेगळी
गृहीतक-सूत्रे नसतील,
आणि $\delta_2$ च्या कोणत्याही
उघड्या गृहीतकात $\delta_1$
चे आयगेन चर येत नसेल,
आणि $\delta_2$ च्या उघड्या गृहीतकाची कोणतीही खूण
$\delta_1$ मधील मुक्ती-नियमाची खूण नसेल,
तर परिणामी
$\Subst{\delta_1}{\delta_2}{B^x}$
ही योग्य सिद्धता आहे.
तिची उघडी गृहीतके म्हणजे
$\delta_1$ मधील बसवलेल्या
उदाहरणांव्यतिरिक्त
उरलेली उघडी गृहीतके,
आणि $\delta_2$ मधील
उघडी गृहीतके.
सिद्धता चे कलम-संयोजन
करताना या अटी सहसा पूर्ण होतात.
~$\delta_1$ आणि $\delta_2$
या मोठ्या सिद्धता
च्या उपसिद्धता असतील,
तर पहिली अट पूर्ण होते.
दुसरी अट पूर्ण नसेल,
तर $\delta_1$ ची आयगेन चरे
पुन्हा नाव देऊन ती पूर्ण
करता येते.
ही खुणांची नवी अटही आवश्यक आहे: दोन्हीकडे तेच सूत्र
असले तरी आत बसवलेले उघडे गृहीतक बाहेरच्या मुक्ती-नियमाने
अनवधानाने मुक्त होता कामा नये. कलम-संयोजनापूर्वी
$\delta_1$ मधील मुक्ती-खुणा आणि त्यांना बांधलेली गृहीतके
एकत्र नवी नावे देऊन, $\delta_2$ च्या खुणांपासून वेगळी
करता येतात; मूळ उघड्या गृहीतकांची नावे जपायची.

तशीच व्याख्या
\Log{N2c} आणि~\Log{N2i}
यांनाही लागू होते.

\begin{defn}
समजा $\delta_1$ ही
$x:B, \Gamma_1 \Sequent A$
ची \Log{N2c}-सिद्धता
(किंवा \Log{N2i}-सिद्धता) आहे;
आणि $\delta_2$ ही
$\Gamma_2 \Sequent B$
ची \Log{N2c}-सिद्धता
(\Log{N2i}-सिद्धता) आहे.
मग $\delta_1$ मधील
अंतिम क्रमवर्तीपर्यंत मुक्त न झालेल्या
$x:B$ शी संबंधित प्रत्येक $x:B \Sequent B$
स्वयंसिद्धकाच्या जागी $\delta_2$ बसवून,
त्याखालील प्रत्येक क्रमवर्तीच्या संदर्भातील
संबंधित $x:B$ काढून त्याऐवजी $\Gamma_2$ जोडून
मिळणारा परिणाम म्हणजे
$\Subst{\delta_1}{\delta_2}{x:B}$.
\end{defn}

\begin{prop}
जर $\delta_1 \Proves
x:B, \Gamma_1 \Sequent
A$ आणि $\delta_2 \Proves
\Gamma_2 \Sequent B$,
तर $\Subst{\delta_1}{\delta_2}{x:B}
\Proves \Gamma_1, \Gamma_2 \Sequent
A$, मात्र दोन्ही सिद्धतांतील गृहीतक-खुणा सुसंगत असाव्यात,
$\Gamma_2$ मधील कोणतीही खूण $\delta_1$ मधील
मुक्ती-नियमाची खूण नसावी आणि $\delta_1$ चे
कोणतेही आयगेन चर $\Gamma_2$ मध्ये येता कामा नये.
\end{prop}

\begin{proof}
  $\pheight{\delta_1}$ वर
  विगमन करून.
\end{proof}












\section{\Log{G2i} पासून \Log{N2i} मध्ये रूपांतर}\label{pt:nat:tgi:sec}

~\Log{G2i} मधील क्रमवर्तींचे पूर्वपक्ष हे
सूत्रे यांचे बहुसंच~$\Gamma$ असतात; तर
\Log{N2i} मधील क्रमवर्तींचे पूर्वपक्ष हे
खूण लावलेल्या सूत्रांचे संच~$\Gamma'$ असतात.
त्यांतील प्रत्येक खूण फक्त एकदाच येते.
खूण लावलेल्या सूत्रांचा संच~$\Gamma'$
हा बहुसंच~$\Gamma$ शी \emph{अनुरूप आहे} असे
तेव्हाच म्हणू, जेव्हा प्रत्येक सूत्र~$A$
साठी $x : A$ या रूपातील~$\Gamma'$ चे
घटक जितके आहेत तितकी संख्या
~$\Gamma$ मध्ये~$A$ जितक्या वेळा येते
(म्हणजे~$A$ च्या प्रती जितक्या आहेत) त्या
संख्येपेक्षा अधिक नसते.

\begin{prop}\label{pt:nat:tgi:prop:G2i-to-N2i}
$\Log{G2i} + \Cut \Proves \Gamma \Sequent \Delta$
असेल, तर
$\Log{N2i} \Proves \Gamma' \Sequent C$.
येथे~$\Delta$ रिकामा असेल तर
$C \ident \lfalse$; अन्यथा
$\Delta = \{A\}$ आणि $C \ident A$.
तसेच~$\Gamma'$ हा~$\Gamma$ शी अनुरूप आहे.
\end{prop}

\begin{proof}
  ~$\pi$ ही~\Log{G2i} + \Cut मध्ये
  $\Gamma \Sequent \Delta$ ची
  सिद्धता असेल, तर~$\Gamma$
  शी अनुरूप असा~$\Gamma'$ आणि
  ~\Log{N2i} मध्ये
  ~$\Gamma' \Sequent C$ ची
  सिद्धता~$\delta$ मिळते,
  हे दाखवू. त्यासाठी
  $n = \pheight{\pi}$
  यावर विगमन करू.

  $n = 0$ असेल, तर~$\pi$
  हे स्वयंसिद्धक आहे.
  ते $\lfalse \Sequent \quad$
  असे असेल, तर~$\delta$
  हे $x:\lfalse \Sequent \lfalse$
  स्वयंसिद्धक; म्हणजे
  $\Gamma' = \{x:\lfalse\}$
  आणि $C \ident \lfalse$.
  ~$\pi$ हे $A \Sequent A$
  या रूपाचे स्वयंसिद्धक असेल,
  तर~$\delta$ हे
  $x: A \Sequent A$.

  $n>0$ असेल, तर~$\pi$
  चा शेवट~$R$ या नियमात होतो.
  विगमन गृहीतकानुसार त्या नियमाच्या
  आधारविधानांशी अनुरूप क्रमवर्तींच्या
  सिद्धता ~\Log{N2i} मध्ये मिळतात.
  या सिद्धतांतील सर्व मुक्तीच्या खुणा
  वेगवेगळ्या आहेत, आणि एखादी खूण~$x$
  मुक्त होत असेल तर तिला मुक्त करणाऱ्या
  अनुमानाच्या वरच येते, असे गृहीत धरू
  शकतो; गरज पडल्यास खुणा बदलू शकतो.
  या सिद्धता एकत्र करून योग्य
  $\Gamma' \Sequent C$ ची सिद्धता
  कशी मिळते ते आता~$R$ नुसार पाहू.

  \begin{enumerate}
    \item $R$ हा \RightR{\Weakening} असेल,
    तर~$\pi$ चा शेवट असा होतो:
    \[
    \AxiomC{}
    \RightLabel{$\pi_1$}
    \Deduce$\Gamma \fCenter$
    \RightLabel{\RightR{\Weakening}}
    \UnaryInf$\Gamma \fCenter A$
    \DisplayProof
    \]
    ~$\pi_1$ वर विगमन गृहीतक लावल्यास
    $\Gamma' \Sequent \lfalse$ ची
    सिद्धता मिळते. तिला~$\FalseInt$
    लावून $\Gamma' \Sequent A$ ची
    सिद्धता मिळते.
    \item $R$ हा \LeftR{\Weakening} असेल,
    तर~$\pi$ चा शेवट असा होतो:
    \[
    \AxiomC{}
    \RightLabel{$\pi_1$}
    \Deduce$\Gamma \fCenter \Delta$
    \RightLabel{\LeftR{\Weakening}}
    \UnaryInf$A, \Gamma \fCenter \Delta$
    \DisplayProof
    \]
    ~$\pi_1$ वर विगमन गृहीतक लावल्यास
    $\Gamma' \Sequent C$ ची
    सिद्धता मिळते; तीच~$\delta$ घेऊ.
    ~$\Gamma'$ हा~$\Gamma$ शी अनुरूप असेल,
    तर तो $A, \Gamma$ शीही अनुरूप असतो.
    कारण~$\Gamma'$ मधील $x : A$ ची संख्या
    ही~$\Gamma$ मधील~$A$ च्या प्रतींपेक्षा
    अधिक नाही; त्यामुळे ती
    $A, \Gamma$ मधील~$A$
    च्या प्रतींपेक्षाही अधिक नाही.
    \item $R$ हा \LeftR{\Contraction} असेल,
    तर~$\pi$ चा शेवट असा होतो:
    \[
    \AxiomC{}
    \RightLabel{$\pi_1$}
    \Deduce$A, A, \Gamma \fCenter \Delta$
    \RightLabel{\LeftR{\Contraction}}
    \UnaryInf$A, \Gamma \fCenter \Delta$
    \DisplayProof
    \]
    ~$\pi_1$ वर विगमन गृहीतक लावल्यास
    $\Pi, \Gamma' \Sequent C$ ची
    सिद्धता~$\delta_1$ मिळते;
    येथे $\Pi \subseteq \{x : A,
    y:A\}$. $\Pi = \{x : A, y:A\}$
    असेल, तर~$\delta_1$ मधील
    सर्व खूण लावलेली सूत्रे~$y:A$
    ही~$x:A$ ने बदला.
    ~$\delta_1$ च्या अंतिम क्रमवर्तीच्या
    संदर्भात~$x$ असल्यामुळे
    ~$\delta_1$ मधील कोणतेही अनुमान
    $x:B$ या रूपाचे सूत्र
    मुक्त करत नाही, असे धरता येते;
    म्हणजे~$\delta_1$ मधील
    कोणत्याही क्रमवर्तीच्या डाव्या बाजूवरील
    $x:B$ सक्रिय नसते.
    म्हणून बदललेली सिद्धता अजूनही
    ~\Log{N2i} मधील सिद्धता असते.
    \item $R$ हा $\RightR{\lif}$ असेल,
    तर~$\pi$ चा शेवट असा होतो:
    \[
    \AxiomC{}
    \RightLabel{$\pi_1$}
    \Deduce$A, \Gamma \fCenter B$
    \RightLabel{\RightR{\lif}}
    \UnaryInf$\Gamma \fCenter A \lif B$
    \DisplayProof
    \]
    विगमन गृहीतकानुसार
    $x: A, \Gamma' \Sequent B$
    किंवा $\Gamma' \Sequent B$
    ची सिद्धता~$\delta_1$ मिळते;
    येथे~$\Gamma'$ हा~$\Gamma$ शी
    अनुरूप आहे. ~$\delta_1$ ला
    \Intro{\lif} चे एक अनुमान जोडून
    मिळणारी सिद्धता~$\delta$ घेऊ.
    \item $R$ हा \LeftR{\lif} असेल,
    तर~$\pi$ चा शेवट असा होतो:
    \[
    \AxiomC{}
    \RightLabel{$\pi_1$}
    \Deduce$\Gamma_1 \fCenter A$
    \AxiomC{}
    \RightLabel{$\pi_2$}
    \Deduce$B, \Gamma_2 \fCenter \Delta$
    \RightLabel{\LeftR{\lif}}
    \BinaryInf$A \lif B, \Gamma_1, \Gamma_2 \fCenter \Delta$
    \DisplayProof
    \]
    विगमन गृहीतकानुसार सिद्धता मिळतात:
    $\delta_1$ ही $\Gamma_1' \Sequent A$ ची;
    आणि~$\delta_2$ ही
    $x: B, \Gamma_2' \Sequent C$
    किंवा $\Gamma_2' \Sequent C$ ची.
    ~$\delta_2$ दुसऱ्या रूपात असेल,
    तर~$\delta_2$ हीच~$\delta$ घेऊ.
    अन्यथा~$\delta_1$ मधील
    सर्व सूत्रांच्या आणि मुक्तीच्या खुणा
    बदलून~$\delta_1$ व~$\delta_2$
    यांत कोणतीही खूण समान नसल्याची
    खात्री करू शकतो. मग
    $\Gamma_1' \cup \Gamma_2'$ हा
    $\Gamma_1 \cup \Gamma_2$ शी अनुरूप आहे.
    ~$\delta_2$ च्या अंतिम क्रमवर्तीत
    $x:B$ असल्यामुळे
    ~$\delta_2$ मधील~$x$ या
    मुक्ती-खुणेच्या कोणत्याही अनुमानात
    $x:B$ सक्रिय नसते.
    ~$\delta_2$ मधील प्रत्येक
    $x:B \Sequent B$ स्वयंसिद्धकाच्या
    जागी पुढील सिद्धता~$\delta_3$ ठेवा:
    \[
    \Axiom$x: A \lif B \fCenter A \lif B$
    \AxiomC{}
    \RightLabel{$\delta_1$}
    \Deduce$\Gamma_1' \fCenter A$
    \RightLabel{\Elim{\lif}}
    \BinaryInf$x: A \lif B, \Gamma_1' \fCenter B$
    \DisplayProof
    \]
    आणि~$\delta_2$ मधील प्रत्येक
    $x:B$ च्या जागी $x:A \lif B$
    ठेवा. परिणामी
    $\Subst{\delta_2}{\delta_3}{x:B}$
    ही $x:A \lif B, \Gamma_1', \Gamma_2' \Sequent C$
    ची सिद्धता आहे.
    \item $R$ हा \RightR{\land} असेल,
    तर~$\pi$ चा शेवट असा होतो:
    \[
    \AxiomC{}
    \RightLabel{$\pi_1$}
    \Deduce$\Gamma_1 \fCenter A$
    \AxiomC{}
    \RightLabel{$\pi_2$}
    \Deduce$\Gamma_2 \fCenter B$
    \RightLabel{\RightR{\land}}
    \BinaryInf$\Gamma_1, \Gamma_2 \fCenter A \land B$
    \DisplayProof
    \]
    विगमन गृहीतकानुसार
    $\Gamma_1' \Sequent A$ ची
    सिद्धता~$\delta_1$ आणि
    $\Gamma_2' \Sequent B$ ची
    सिद्धता~$\delta_2$ मिळते.
    ~$\delta_2$ मधील सर्व सूत्रांच्या
    आणि मुक्तीच्या खुणा बदलून
    ~$\delta_1$ व~$\delta_2$ यांत
    कोणतीही खूण समान नसल्याची
    खात्री करता येते. मग
    $\Gamma_1' \cup \Gamma_2'$ हा
    $\Gamma_1 \cup \Gamma_2$ शी अनुरूप आहे.
    ~$\delta_1$ आणि~$\delta_2$
    यांच्या अंतिम क्रमवर्तींना
    $\Intro{\land}$ लावून
    सिद्धता~$\delta$ मिळते.
    \item $R$ हा $\LeftR{\land}$ असेल,
    तर, उदाहरणार्थ, ~$\pi$ चा
    शेवट असा होतो:
    \[
    \AxiomC{}
    \RightLabel{$\pi_1$}
    \Deduce$A, \Gamma \fCenter \Delta$
    \RightLabel{\LeftR{\land}}
    \UnaryInf$A \land B, \Gamma \fCenter \Delta$
    \DisplayProof
    \]
    विगमन गृहीतकानुसार
    $x: A, \Gamma' \Sequent C$
    किंवा $\Gamma' \Sequent C$
    ची सिद्धता~$\delta_1$ मिळते;
    येथे~$\Gamma'$ हा~$\Gamma$ शी
    अनुरूप आहे. दुसऱ्या प्रसंगी
    ~$\delta_1$ हीच~$\delta$ घेऊ.
    पहिल्या प्रसंगी~$x:A$
    अंतिम क्रमवर्तीत असल्यामुळे
    ~$\delta_1$ मधील~$x$ या
    मुक्ती-खुणेच्या कोणत्याही अनुमानात
    $x:A$ सक्रिय नसते.
    प्रत्येक $x:A \Sequent A$
    स्वयंसिद्धकाच्या जागी
    पुढील सिद्धता~$\delta_2$ ठेवा:
    \[
    \Axiom$x: A \land B \fCenter A \land B$
    \RightLabel{\Elim{\land}}
    \UnaryInf$x: A \land B \fCenter A$
    \DisplayProof
    \]
    आणि प्रत्येक~$x:A$ च्या जागी
    $x:A \land B$ ठेवा; म्हणजे
    $\Subst{\delta_1}{\delta_2}{x:A}$
    घ्या. ही $x:A \land B, \Gamma' \Sequent C$
    ची सिद्धता आहे.
    \item $R$ हा \Cut असेल,
    तर~$\pi$ चा शेवट असा होतो:
    \[
    \AxiomC{}
    \RightLabel{$\pi_1$}
    \Deduce$\Gamma_1 \fCenter A$
    \AxiomC{}
    \RightLabel{$\pi_2$}
    \Deduce$A, \Gamma_2 \fCenter \Delta$
    \RightLabel{\Cut}
    \BinaryInf$\Gamma_1, \Gamma_2 \fCenter \Delta$
    \DisplayProof
    \]
    विगमन गृहीतकानुसार
    $\Gamma_1' \Sequent A$ ची
    सिद्धता~$\delta_1$ आणि
    $x:A, \Gamma_2' \Sequent C$
    किंवा $\Gamma_2' \Sequent C$
    ची सिद्धता~$\delta_2$ मिळते.
    दुसऱ्या प्रसंगी~$\delta$
    म्हणून~$\delta_2$ घ्या.
    अन्यथा, गरज असल्यास
    ~$\delta_1$ मधील खुणा बदलून
    दोन सिद्धतांच्या खुणा वेगळ्या करा.
    ~$\delta_2$ मधील प्रत्येक
    $x:A \Sequent A$ स्वयंसिद्धकाच्या
    जागी~$\delta_1$ ठेवून
    $\Subst{\delta_2}{\delta_1}{x:A}$
    ही $\Gamma_1', \Gamma_2' \Sequent C$
    ची सिद्धता~$\delta$ मिळते.
  \end{enumerate}
\end{proof}













\section{\Log{N2} पासून \Log{G2} मध्ये रूपांतर}\label{pt:nat:tni:sec}

\begin{prop}\label{pt:nat:tni:prop:N2-to-G2}
$\Log{N2c}\ (\Log{N2i}) \Proves
\Gamma \Sequent A$ असेल, तर
$\Log{G2c}\ (\Log{G2i}) + \Cut \Proves \Gamma'
\Sequent \Delta$. येथे~$\Gamma$ मधील
खुणा काढून मिळणारा सूत्रांचा
बहुसंच~$\Gamma'$ आहे; तसेच~$A$
हा~$\lfalse$ नसेल तर
$\Delta = \{A\}$, आणि असेल तर
$\Delta = \emptyset$.
\end{prop}

\begin{proof}
पुन्हा~$\Gamma \Sequent A$ च्या
सिद्धता~$\delta$ च्या उंचीवर
विगमन करू. ~$\delta$ हे
$x:A \Sequent A$ स्वयंसिद्धक
असेल, तर~$\pi$ हे अनुरूप
$A \Sequent A$ स्वयंसिद्धक
आहे; किंवा~$A \ident \lfalse$
असेल तर
$\lfalse \Sequent \ $ स्वयंसिद्धक आहे.

~$\delta$ चा शेवट
अनुमाननियमात होत असेल, तर
पुढील प्रसंग वेगळे पाहू:
\begin{enumerate}
  \item $R$ हा
  मुख्य सूत्र~$B \lif C$
  असलेला~\Intro{\lif} नियम असेल,
  आणि~$x:B$ हे आधारविधानाच्या
  संदर्भात असेल पण निष्कर्षाच्या
  संदर्भात नसेल, तर~$\delta$
  चा शेवट असा होतो:
  \[
  \AxiomC{}
  \RightLabel{$\delta_1$}
  \Deduce$x:B, \Gamma \fCenter C$
  \RightLabel{\Intro{\lif}}
  \UnaryInf$\Gamma \fCenter B \lif C$
  \DisplayProof
  \]
  विगमन गृहीतकानुसार
  $B, \Gamma' \Sequent C$
  ची सिद्धता~$\pi_1$
  मिळते; $C \ident \lfalse$
  असेल तर तिचे रूप
  $B, \Gamma' \Sequent \ $
  असते. त्यावर~\RightR{\lif}
  लावून~$\pi$ मिळते;
  $C \ident \lfalse$ असताना
  त्याआधी~\RightR{\Weakening}
  लावावे लागते.
  \item $R$ हा~\Intro{\lif}
  असेल, पण~$x:B$
  आधारविधानाच्या संदर्भात नसेल,
  तर~$\delta$ चा शेवट असा होतो:
  \[
  \AxiomC{}
  \RightLabel{$\delta_1$}
  \Deduce$\Gamma \fCenter C$
  \RightLabel{\Intro{\lif}}
  \UnaryInf$\Gamma \fCenter B \lif C$
  \DisplayProof
  \]
  विगमन गृहीतकानुसार
  निष्पन्न सूत्र असत्य
  नसेल तेव्हा~$\Gamma' \Sequent C$
  ची सिद्धता~$\pi_1$ मिळते;
  असत्याच्या प्रसंगी तिचा
  उत्तरपक्ष रिकामा असतो.
  गरज असल्यास आधी
  उजवीकडील दुर्बलीकरणाने
  असत्य उजवीकडे आणा.
  मग~$B$ साठी
  \LeftR{\Weakening}
  आणि त्यानंतर
  \RightR{\lif} लावून~$\pi$
  मिळते.
  \item $R$ हा~\Elim{\lif}
  असेल, तर~$\delta$ चा
  शेवट असा होतो:
  \[
  \AxiomC{}
  \RightLabel{$\delta_1$}
  \Deduce$\Gamma_1 \fCenter B \lif A$
  \AxiomC{}
  \RightLabel{$\delta_2$}
  \Deduce$\Gamma_2 \fCenter B$
  \RightLabel{\Elim{\lif}}
  \BinaryInf$\Gamma_1, \Gamma_2 \fCenter A$
  \DisplayProof
  \]
  येथे~$\Gamma = \Gamma_1 \cup \Gamma_2$.
  विगमन गृहीतकानुसार
  सिद्धता मिळतात:
  $\pi_1$ ही
  $\Gamma_1' \Sequent B \lif A$
  ची आणि~$\pi_2$ ही
  $\Gamma_2' \Sequent B$ ची.
  दुसऱ्या निष्कर्षसूत्राचा असत्य
  हा प्रसंग असेल, तर दुसरी
  सिद्धता मिळवण्यासाठी
  विगमनातून आलेल्या रिकाम्या
  उत्तरपक्षाच्या सिद्धतेवर
  उजवीकडील दुर्बलीकरण लावा.
  पुढीलप्रमाणे सिद्धता~$\pi$
  मिळते:
  \[
  \AxiomC{}
  \RightLabel{$\pi_1$}
  \Deduce$\Gamma_1' \fCenter B \lif A$
  \AxiomC{}
  \RightLabel{$\pi_2$}
  \Deduce$\Gamma_2' \fCenter B$
  \Axiom$A \fCenter A$
  \RightLabel{\LeftR{\lif}}
  \BinaryInf$B \lif A, \Gamma_2' \fCenter A$
  \RightLabel{\Cut}
  \BinaryInf$\Gamma_1', \Gamma_2' \fCenter A$
  \DisplayProof
  \]
  बहुसंच~$\Gamma_1' \cup \Gamma_2'$
  हा $\Gamma' =
  (\Gamma_1 \cup \Gamma_2)'$
  याच्याशी समान असेलच असे नाही.
  उदाहरणार्थ,~$\Gamma_1$
  आणि~$\Gamma_2$ दोन्हींमध्ये
  $x:C$ असेल, तर
  $\Gamma_1' \cup \Gamma_2'$
  मध्ये~$C$ च्या दोन प्रती
  असतात; पण
  $(\Gamma_1 \cup \Gamma_2)'$
  मध्ये एकच प्रत असते.
  योग्य~\LeftR{\Contraction}
  अनुमाने जोडून हे दुरुस्त करता येते.

  $A \ident \lfalse$
  असेल, तर~$\delta_2$
  वरून विगमनाने मिळणारी
  सिद्धता एकटी रिकामा
  उत्तरपक्ष देत नाही.
  वरच्या रचनेत
  $A$ च्या ओळख-स्वयंसिद्धकाऐवजी
  रिकामा उत्तरपक्ष असलेले
  असत्य-स्वयंसिद्धक वापरा.
  डाव्या अभिव्यंजन-नियमापुढे
  कट लावल्यावर
  $\Gamma_1', \Gamma_2' \Sequent \ $
  ची सिद्धता~$\pi$ मिळते.
  त्यामुळे यासाठी
  \LeftR{\Weakening} किंवा
  \RightR{\Weakening}
  लावण्याची गरज नाही.
\item $R$ हा~\FalseCl असेल,
तर~$\delta$ चा शेवट असा:
  \[
  \AxiomC{}
  \RightLabel{$\delta_1$}
  \Deduce$x: \lnot A, \Gamma \fCenter \lfalse$
  \RightLabel{\FalseCl}
  \UnaryInf$\Gamma \fCenter A$
  \DisplayProof
  \]
  विगमन गृहीतकानुसार
  $\lnot A, \Gamma' \Sequent \ $
  ची सिद्धता~$\pi_1$
  आहे. पुढील सिद्धता~$\pi$ घेऊ:
  \[
  \Axiom$A \fCenter A$
  \RightLabel{\RightR{\lnot}}
  \UnaryInf$\fCenter A, \lnot A$
  \AxiomC{}
  \RightLabel{$\pi_1$}
  \Deduce$\lnot A, \Gamma' \fCenter $
  \RightLabel{\Cut}
  \BinaryInf$\Gamma' \fCenter A$
  \DisplayProof
  \]
  $A\ident\lfalse$ असेल, तर या शेवटच्या क्रमवर्तीचा
  उत्तरपक्षही रिकामा हवा. त्यासाठी वरील सिद्धतेला
  $\lfalse\Sequent\ $ या स्वयंसिद्धकासह आणखी एक
  $\Cut$ लावून $\Gamma'\Sequent\ $ मिळवायचे.
  $\FalseCl$ मध्ये कोणतेही गृहीतक मुक्त होत नसेल,
  तर आधाराची सिद्धता विगमनाने थेट
  $\Gamma'\Sequent\ $ देते; $A\not\ident\lfalse$
  असल्यासच उजवीकडील दुर्बलीकरणाने $A$ जोडायचे.
\end{enumerate}
फक्त~\FalseCl{} च्या
रूपांतरातून मिळणाऱ्या
सिद्धता~$\pi$ मध्ये
एखाद्या क्रमवर्तीच्या उजव्या बाजूस
एकापेक्षा अधिक सूत्र
येऊ शकते. म्हणजे
~\Log{N2i} मधील सिद्धता
रूपांतरित केल्यास
\Log{G2i} मधील सिद्धता मिळते.
\end{proof}



\chapter{सामान्यरूपीकरण}\label{pt:nor:chap}










\section{परिचय}\label{pt:nor:int:sec}

~$A \lif B$ ही शर्त तुम्ही
सिद्ध केली असेल, तर
\Elim{\lif} वापरून~$B$ ची
सिद्धता करताना तिचा उपयोग
होऊ शकतो. अर्थात त्यासाठी~$A$
ची सिद्धता~$\delta_1$ ही हवीच.
~$A \lif B$ ची तुमची
सिद्धता~$\delta_3$ बहुधा
\Intro{\lif} ने संपते;
तो नियम उघड्या गृहीतक~$A$
पासून~$B$ ची
सिद्धता~$\delta_2$
घेऊन~$A \lif B$
ची सिद्धता बनवतो.
तसे असल्यास तुमची अंतिम
सिद्धता पुढील रूपाची असते:
\[
\AxiomC{$\Discharge{A}{x}$}
\RightLabel{$\delta_2$}
\DeduceC{$B$}
\DischargeRule{\Intro{\lif}}{x}
\UnaryInfC{$A \lif B$}
\AxiomC{}
\RightLabel{$\delta_1$}
\DeduceC{$A$}
\RightLabel{\Elim{\lif}}
\BinaryInfC{$B$}
\DisplayProof
\]
आधीच मिळालेली
$A \lif B$ ची सिद्धता
पुन्हा वापरली म्हणून
तुमची पद्धत कार्यक्षम असली,
तरी सिद्धता स्वतः
कार्यक्षम नाही.
$A \lif B$ आणून लगेच
निरसित करण्याचा हा
वळसा वाटतो. ~$B$
सिद्ध करण्याचा अधिक
थेट मार्ग असायला हवा.

तसा मार्ग नेहमीच असतो.
अशा “वळशांना”—परिचय
नियमाच्या लगेच नंतर
निरसन नियम लावल्यास
तयार होणाऱ्या रचनांना—
नैसर्गिक निगमनाच्या
सिद्धता मधून नेहमी
काढून टाकता येते.
याला \emph{सामान्यरूपीकरण}
म्हणतात; परिणामी मिळणारी
सिद्धता ही \emph{सामान्यरूप}
असते (किंवा ती सिद्धता
\emph{सामान्य रूपात} आहे
असे म्हणतात).

या निष्कर्षाची सिद्धता
क्रमवर्ती कलनातील
कट-निर्मूलन प्रमेयाप्रमाणेच
काहीशी गुंतागुंतीची आहे;
तिच्यात अनेक प्रसंग
पाहावे लागतात.
अंतःप्रज्ञावादी~\Log{N1i}
पेक्षा अभिजात~\Log{N1c}
साठी ती सिद्ध करणे
अधिक कठीण असते.
क्रमवर्ती कलनात~\Cut{}
नियम वापरल्यामुळे विना-कट
सिद्ध करता येत नाही असे
अधिक काही सिद्ध करता येत
नाही, हे कट-निर्मूलन दाखवते.
त्याचप्रमाणे वळसे टाळून
मिळू शकत नाहीत असे
अधिक निष्कर्ष वळसे वापरूनही
मिळत नाहीत, हे
सामान्यरूपीकरण दाखवते.
इतर साम्येही आहेत:
एखाद्या निष्कर्षाची
वळसांसह सर्वात लहान
सिद्धता जितकी असेल,
त्याच निष्कर्षाची
सामान्यरूप सिद्धता
त्याहून खूप मोठी असू शकते.
तसेच क्रमवर्ती कलनात
विना-कट सर्वात लहान
सिद्धता पेक्षा
कट असलेली सिद्धता
खूप लहान असू शकतात.
उप-सूत्र गुणधर्मही
यातून मिळतो: नैसर्गिक
निगमनाच्या सिद्धता मध्ये
निष्कर्ष आणि उघडी गृहीतके
यांच्या उप-सूत्रे, संख्यापकांमुळे लागणारी
त्यांची पद-आदेशनाची उदाहरणे आणि आवश्यक असत्य
वापरून निष्कर्ष नेहमी
सिद्ध करता येतो.
क्रमवर्ती कलनात हा गुणधर्म
कट-निर्मूलनातून मिळतो,
तसाच येथे तो
सामान्यरूपीकरणातून मिळतो.












\section{खंड आणि कट}\label{pt:nor:seg:sec}

परिचय नियमाच्या
नंतर निरसन नियम
येणे हा सिद्धता मधला
वळसा असतो; सिद्धता
सामान्यरूप करणे म्हणजे
असे वळसे काढणे.
पण~\Elim{\lor} आणि
\Elim{\lexists} मुळे
“काढायचा वळसा” याची
व्याख्याही थोडी
गुंतागुंतीची होते.
अडचण अशी: वळशात
निरसन नियम
परिचय नियमाच्या
\emph{लगेच} नंतर येईलच
असे नाही. त्यांच्यामध्ये
\Elim{\lor} किंवा
\Elim{\lexists} असू
शकतो; उदाहरणार्थ:
\[
\AxiomC{}
\DeduceC{$\lexists[x][C(x)]$}
\AxiomC{$\Discharge{A}{x}$}
\DeduceC{$B$}
\DischargeRule{\Intro{\lif}}{x}
\UnaryInfC{$A \lif B$}
\DischargeRule{\Elim{\lexists}}{y}
\BinaryInfC{$A \lif B$}
\AxiomC{}
\DeduceC{$A$}
\RightLabel{\Elim{\lif}}
\BinaryInfC{$B$}
\DisplayProof
\]
इतकेच नव्हे,~\Intro{\lif}
आणि~\Elim{\lif} यांच्या
मध्ये कितीही
\Elim{\lor} किंवा
\Elim{\lexists} अनुमाने
असू शकतात; उदाहरणार्थ:
\[
\AxiomC{}
\DeduceC{$\lexists[x][C(x)]$}
\AxiomC{}
\DeduceC{$D \lor E$}
\AxiomC{$\Discharge{A}{x}$}
\DeduceC{$B$}
\DischargeRule{\Intro{\lif}}{x}
\UnaryInfC{$A \lif B$}
\AxiomC{$\Discharge{A}{x}$}
\DeduceC{$B$}
\DischargeRule{\Intro{\lif}}{x}
\UnaryInfC{$A \lif B$}
\RightLabel{\Elim{\lor}}
\TrinaryInfC{$A \lif B$}
\DischargeRule{\Elim{\lexists}}{y}
\BinaryInfC{$A \lif B$}
\AxiomC{}
\DeduceC{$A$}
\RightLabel{\Elim{\lif}}
\BinaryInfC{$B$}
\DisplayProof
\]
या उदाहरणात तोच
\Elim{\lif} एकाहून
अधिक वळशांत भाग
घेतो: त्याच्या आधी
एकाहून अधिक~\Intro{\lif}
आहेत. अशा सर्व
शक्यता सामावण्यासाठी
~\Log{N1i}-सिद्धतेतील
सूत्र आस्थितींच्या
विशिष्ट अनुक्रमांनी
वळसे परिभाषित करतो;
त्यांना \emph{खंड}
म्हणतात. येथे त्या
~$A \lif B$
च्या आस्थितींच्या
अनुक्रमिका आहेत:
$\Elim{\lor}$ किंवा
$\Elim{\lexists}$
च्या गौण आधारविधानानंतर
त्याच अनुमानाचा
निष्कर्ष येतो.
या उदाहरणात असे
दोन खंड आहेत;
प्रत्येकात~$A \lif B$
च्या तीन आस्थिती
आहेत. आता औपचारिक
व्याख्या पाहू.

\begin{defn}
~\Log{N1i} मधील
सिद्धता~$\delta$
मध्ये एकाच सूत्र~$A$
च्या $A_1$, \dots,
$A_n$ या~$\delta$
मधील आस्थितींचा
अनुक्रम पुढील अटी
पूर्ण करत असेल,
तर तो \emph{खंड} आहे:
\begin{enumerate}
  \item $A_1$ हा
  \Elim{\lor} किंवा
  \Elim{\lexists}
  चा निष्कर्ष नसतो;
  \item $A_n$ हे
  \Elim{\lor} किंवा
  \Elim{\lexists}
  चे गौण आधारविधान नसते;
  \item $i = 1$, \dots,~$n-1$
  साठी~$A_i$ हे
  \Elim{\lor} किंवा
  \Elim{\lexists}
  चे गौण आधारविधान
  असते, आणि~$A_{i+1}$
  हा त्याच अनुमानाचा
  निष्कर्ष असतो.
\end{enumerate}
त्या खंडाची
\emph{लांबी}~$n$ आहे.
\end{defn}

प्रत्येक खंड काढून
टाकण्याजोगा वळसा
असतोच असे नाही.
अशा वळशांना
\emph{कट} म्हणतात.

\begin{defn}
एखाद्या खंडातील
$A_n$ हे
निरसन अनुमानाचे
प्रमुख आधारविधान असेल,
आणि पुढीलपैकी एक
अट पूर्ण होत असेल,
तर तो \emph{कट} आहे:
$n=1$ असून
$A_1 = A_n$ हे
परिचय अनुमानाचा
किंवा~\FalseInt चा
निष्कर्ष असेल;
किंवा~$n>1$ असेल.
\end{defn}

कट-खंडाची सुरुवात
परिचय नियमाच्या
निष्कर्षाने व्हावीच
लागते असे नाही.
तो निरसन
नियमाच्या प्रमुख
आधारविधानात संपणे
हीच आवश्यक अट आहे.
मात्र लांबी~$1$
असलेला खंड कट
म्हणून मोजायचा असेल,
तर तो परिचय
नियमाचा किंवा~\FalseInt
चा निष्कर्ष असणेही
आवश्यक आहे.

\begin{defn}
कटाची \emph{कोटी}
ही~$\depth{A}$
असते. सिद्धता~$\delta$
ची \emph{कट-कोटी}
$\cutrank{\delta}$
म्हणजे~$\delta$
मधील कटांच्या
कोटींपैकी महत्तम कोटी.
कटाची कोटी
$\cutrank{\delta}$
इतकी असेल, म्हणजे
ती महत्तम असेल,
तर तो~$\delta$
मधील \emph{महत्तम} कट.
~$\delta$ ची
\emph{कट-लांबी}
म्हणजे तिच्यातील
सर्व महत्तम कटांच्या
लांबींची बेरीज.
\end{defn}













\section{क्रमपरिवर्तन रूपांतरणे}\label{pt:nor:per:sec}

~\Log{N1i} मधील कट
हे निरसन नियमाच्या
प्रमुख आधारविधानात संपणारे
खंड असतात. सामान्यरूपीकरण
सिद्धतेतील एक प्रसंग म्हणजे
महत्तम कटांची लांबी आणि
त्याबरोबर सिद्धता ची
कट-लांबी कमी करणे.
$>1$ लांबीचा कट
अनेक प्रकारे संपू शकतो:
कटमधील~$A$ ची शेवटची
सूत्र आस्थिती
(\Elim{\lor} किंवा
\Elim{\lexists}) चा
निष्कर्ष आहे का, आणि
ती कोणत्या निरसन
अनुमानाचे प्रमुख आधारविधान
आहे, यांवर ते अवलंबून असते.

उदाहरणार्थ, संबंधित
अनुमाने~\Elim{\lor} आणि
\Elim{\land} असतील,
$A \ident B \land C$
असेल, आणि~$\delta_1$
ही उप-सिद्धता पुढील
रूपात संपत असेल:
\[
\AxiomC{}
\RightLabel{$\delta_2$}
\DeduceC{$D \lor E$}
\AxiomC{$\Discharge{D}{x}$}
\RightLabel{$\delta_3$}
\DeduceC{$B \land C$}
\AxiomC{$\Discharge{E}{y}$}
\RightLabel{$\delta_4$}
\DeduceC{$B \land C$}
\DischargeRule{\Elim{\lor}}{x\,y}
\TrinaryInfC{$B \land C$}
\RightLabel{\Elim{\land}[1]}
\UnaryInfC{$B$}
\DisplayProof
\]
तर या कटमधील शेवटची
सूत्र आस्थिती~$A_n$
ही~\Elim{\lor} चा
निष्कर्ष, म्हणजे खालचे
$B \land C$, आहे.
त्याआधीची सूत्र
आस्थिती~$A_{n-1}$
ही~\Elim{\lor}
च्या गौण आधारविधानांपैकी
एक आहे.

या कटाची लांबी
कमी करण्यासाठी
\Elim{\lor} आणि
\Elim{\land} अनुमानांचा
क्रम बदलू (“क्रमपरिवर्तन”
करू) शकतो. म्हणजे~$\delta$
मधील उपसिद्धता~$\delta_1$
पुढील रचनेने बदलतो:
\[
\AxiomC{}
\RightLabel{$\delta_2$}
\DeduceC{$D \lor E$}
\AxiomC{$\Discharge{D}{x}$}
\RightLabel{$\delta_3$}
\DeduceC{$B \land C$}
\RightLabel{\Elim{\land}[1]}
\UnaryInfC{$B$}
\AxiomC{$\Discharge{E}{y}$}
\RightLabel{$\delta_4$}
\DeduceC{$B \land C$}
\RightLabel{\Elim{\land}[1]}
\UnaryInfC{$B$}
\DischargeRule{\Elim{\lor}}{x\,y}
\TrinaryInfC{$B$}
\DisplayProof
\]
आधी~\Elim{\lor} च्या
खाली~\Elim{\land}
च्या आधारविधानात संपणारे
कट आता~\Elim{\lor}
च्या गौण आधारविधानांच्या
वर असलेल्या~\Elim{\land}
अनुमानांपैकी एका
आधारविधानात संपतात.
म्हणजे प्रत्येकाची लांबी
~$1$ ने कमी झाली.

याखेरीज~$\delta_2$,
$\delta_3$ किंवा~$\delta_4$
यांपैकी एखाद्यात पूर्णपणे
असलेला कट, किंवा~$\delta$
मधील~$\delta_1$
बाहेर पूर्णपणे असलेला कट,
बदलत नाही. त्यात तीच
सूत्रे आणि तीच
अनुमाने असतात; म्हणून
विशेषतः त्याची कोटी व
लांबी~$\delta$ मधील
अनुरूप कटाइतकीच असते.
त्यामुळे नव्या सिद्धता
ची कट-कोटी वाढत नाही,
आणि नव्या सिद्धता
ची कट-लांबी
जुन्या सिद्धता च्या
कट-लांबीपेक्षा कमी होते.

\begin{prob}
क्रमपरिवर्तन रूपांतरणानंतर
किमान एका कटाची लांबी~$1$
ने कमी होते.
\Elim{\land} ला
एखाद्या~\Elim{\lor}
च्या पलीकडे हलवले,
तर प्रत्यक्षात दोन
कट-खंडांची लांबी कमी
होते; त्यामुळे या
प्रसंगी सिद्धता ची
कट-लांबी किमान~$2$
ने कमी होते.
अशा एकाच क्रमपरिवर्तनाने
सिद्धता ची कट-लांबी
$2$ पेक्षा अधिक कमी
होऊ शकते का?
कटाची \emph{कोटी}
कमी होऊ शकते का?
\end{prob}

निरसन नियम
\Elim{\lor} किंवा
\Elim{\lexists} नियमाच्या
वर हलवण्याच्या प्रक्रियेला
\emph{क्रमपरिवर्तन रूपांतरण}
म्हणतात. काही नियम-जोड्यांचे
नमुने~\ref{pt:nor:per:tab:permutation-conversions}
मध्ये दिले आहेत.
(उरलेले सरावासाठी सोडले आहेत.)
हे नमुने वापरण्यापूर्वी मुक्ती-नियमांच्या बद्ध खुणा आणि
त्यांना बांधलेली गृहीतके एकत्र नवी नावे द्यायची,
जेणेकरून हलवलेल्या आधारसिद्धतेची उघडी गृहीतके
त्या नियमांनी अनवधानाने मुक्त होणार नाहीत.
उपसिद्धतेच्या अनेक प्रती लागल्यास प्रत्येक प्रतीतील
आयगेन स्थिरांक आणि त्यांच्याशी संबंधित बद्ध खुणा
स्वतंत्र नवी नावे देऊन पुन्हा नियमित सिद्धता मिळवायची;
उघडी गृहीतके आणि अंतिम सूत्रे जपायची.
या नावबदलांमुळे कटांची कोटी किंवा लांबी बदलत नाही.





























































































































{\small
\begin{longtable}{@{}l@{}}
\AxiomC{$D \lor E$}
\AxiomC{$\Discharge{D}{x}$}
\shortDeduce
\DeduceC{$B \land C$}
\AxiomC{$\Discharge{E}{y}$}
\shortDeduce
\DeduceC{$B \land C$}
\DischargeRule{\Elim{\lor}}{x\,y}
\TrinaryInfC{$B \land C$}
\RightLabel{\Elim{\land}[1]}
\UnaryInfC{$B$}
\DisplayProof
$\longrightarrow$\\[6ex]
\quad\AxiomC{$D \lor E$}
\AxiomC{$\Discharge{D}{x}$}
\shortDeduce
\DeduceC{$B \land C$}
\RightLabel{\Elim{\land}[1]}
\UnaryInfC{$B$}
\AxiomC{$\Discharge{E}{y}$}
\shortDeduce
\DeduceC{$B \land C$}
\RightLabel{\Elim{\land}[1]}
\UnaryInfC{$B$}
\DischargeRule{\Elim{\lor}}{x\,y}
\TrinaryInfC{$B$}
\DisplayProof
\\[10ex]
\AxiomC{$D \lor E$}
\AxiomC{$\Discharge{D}{x}$}
\shortDeduce
\DeduceC{$B \lif C$}
\AxiomC{$\Discharge{E}{y}$}
\shortDeduce
\DeduceC{$B \lif C$}
\DischargeRule{\Elim{\lor}}{x\,y}
\TrinaryInfC{$B \lif C$}
\AxiomC{$B$}
\RightLabel{\Elim{\lif}}
\BinaryInfC{$C$}
\DisplayProof
$\longrightarrow$\\[6ex]
\AxiomC{$D \lor E$}
\AxiomC{$\Discharge{D}{x}$}
\shortDeduce
\DeduceC{$B \lif C$}
\AxiomC{$B$}
\RightLabel{\Elim{\lif}}
\BinaryInfC{$C$}
\AxiomC{$\Discharge{E}{y}$}
\shortDeduce
\DeduceC{$B \lif C$}
\AxiomC{$B$}
\RightLabel{\Elim{\lif}}
\BinaryInfC{$C$}
\DischargeRule{\Elim{\lor}}{x\,y}
\TrinaryInfC{$C$}
\DisplayProof
\\[10ex]
\AxiomC{$D \lor E$}
\AxiomC{$\Discharge{D}{x}$}
\shortDeduce
\DeduceC{$B \lif C$}
\AxiomC{$\Discharge{E}{y}$}
\shortDeduce
\DeduceC{$B \lif C$}
\DischargeRule{\Elim{\lor}}{x\,y}
\TrinaryInfC{$B \lif C$}
\AxiomC{$B$}
\RightLabel{\Elim{\lif}}
\BinaryInfC{$C$}
\DisplayProof
$\longrightarrow$\\[6ex]
\AxiomC{$D \lor E$}
\AxiomC{$\Discharge{D}{x}$}
\shortDeduce
\DeduceC{$B \lif C$}
\AxiomC{$B$}
\RightLabel{\Elim{\lif}}
\BinaryInfC{$C$}
\AxiomC{$\Discharge{E}{y}$}
\shortDeduce
\DeduceC{$B \lif C$}
\AxiomC{$B$}
\RightLabel{\Elim{\lif}}
\BinaryInfC{$C$}
\DischargeRule{\Elim{\lor}}{x\,y}
\TrinaryInfC{$C$}
\DisplayProof
\\[10ex]
\AxiomC{$D \lor E$}
\AxiomC{$\Discharge{D}{x}$}
\shortDeduce
\DeduceC{$\lforall[x][B(x)]$}
\AxiomC{$\Discharge{E}{y}$}
\shortDeduce
\DeduceC{$\lforall[x][B(x)]$}
\DischargeRule{\Elim{\lor}}{x\,y}
\TrinaryInfC{$\lforall[x][B(x)]$}
\RightLabel{\Elim{\lforall}}
\UnaryInfC{$B(t)$}
\DisplayProof
\quad$\longrightarrow$\\[8ex]
\quad\AxiomC{$D \lor E$}
\AxiomC{$\Discharge{D}{x}$}
\shortDeduce
\DeduceC{$\lforall[x][B(x)]$}
\RightLabel{\Elim{\lforall}}
\UnaryInfC{$B(t)$}
\AxiomC{$\Discharge{E}{y}$}
\shortDeduce
\DeduceC{$\lforall[x][B(x)]$}
\RightLabel{\Elim{\lforall}}
\UnaryInfC{$B(t)$}
\DischargeRule{\Elim{\lor}}{x\,y}
\TrinaryInfC{$B(t)$}
\DisplayProof
\\
\caption{\Log{N1i} साठी काही क्रमपरिवर्तन रूपांतरणे.}
\label{pt:nor:per:tab:permutation-conversions}
\end{longtable}
}

\begin{prob}
\Elim{\lor} नंतर
\Elim{\lor} किंवा
\Elim{\lnot} येत असेल,
तर क्रमपरिवर्तन
रूपांतरणे द्या.
\end{prob}

\begin{prob}
\Elim{\lexists} नंतर
\Elim{\lexists} येत असेल,
तर क्रमपरिवर्तन
रूपांतरणे द्या.
दुसऱ्या प्रसंगी
आयगेन चराची कोणतीही
अट का मोडत नाही,
हे स्पष्ट करा.
\end{prob}

क्रमपरिवर्तन रूपांतरण
लावताना आयगेन चराची
अट मोडत नाही ना,
हे तपासले पाहिजे.
पुढील उदाहरण घ्या:
\[
\AxiomC{}
\RightLabel{$\delta_2$}
\DeduceC{$D \lor E$}
\AxiomC{$\Discharge{D}{x}$}
\RightLabel{$\delta_3$}
\DeduceC{$\lexists[x][B(x)]$}
\AxiomC{$\Discharge{E}{y}$}
\RightLabel{$\delta_4$}
\DeduceC{$\lexists[x][B(x)]$}
\DischargeRule{\Elim{\lor}}{x\,y}
\TrinaryInfC{$\lexists[x][B(x)]$}
\AxiomC{$\Discharge{B(c)}{z}$}
\RightLabel{$\delta_5$}
\DeduceC{$C$}
\DischargeRule{\Elim{\lexists}}{z}
\BinaryInfC{$C$}
\DisplayProof
\]
याच्या जागी पुढील
रचना बसवली, तर
\[
\AxiomC{}
\RightLabel{$\delta_2$}
\DeduceC{$D \lor E$}
\AxiomC{$\Discharge{D}{x}$}
\RightLabel{$\delta_3$}
\DeduceC{$\lexists[x][B(x)]$}
\AxiomC{$\Discharge{B(c)}{z}$}
\RightLabel{$\delta_5$}
\DeduceC{$C$}
\DischargeRule{\Elim{\lexists}}{z}
\BinaryInfC{$C$}
\AxiomC{$\Discharge{E}{y}$}
\RightLabel{$\delta_4$}
\DeduceC{$\lexists[x][B(x)]$}
\AxiomC{$\Discharge{B(c)}{z}$}
\RightLabel{$\delta_5$}
\DeduceC{$C$}
\DischargeRule{\Elim{\lexists}}{z}
\BinaryInfC{$C$}
\DischargeRule{\Elim{\lor}}{x\,y}
\TrinaryInfC{$C$}
\DisplayProof
\]
आयगेन चर~$c$ हा,
उदाहरणार्थ,~$D$
मध्ये येईल का,
अशी शंका येऊ शकते.
मूळ सिद्धता मध्ये
गृहीतक~$D$ हे
मूळ~$\Elim{\lexists}$
च्या वर मुक्त होते.
म्हणून ते
\Elim{\lexists}
अनुमानाच्या प्रमुख
आधारविधानाच्या वर
उघड्या गृहीतकात
येत नाही; त्यामुळे
तेथे आयगेन चराची अट
पूर्ण होते. पण नव्या
सिद्धता मध्ये~$D$
हे~\Elim{\lexists}
अनुमानांनंतरच मुक्त
होते, म्हणून ती अट
मोडेल असे वाटते.
तथापि, आपल्या
सिद्धता नियमित
आहेत हे लक्षात घ्या:
प्रत्येक आयगेन चर
एकाच अनुमानाचा
आयगेन चर असतो आणि
तो सिद्धता मध्ये
फक्त~\Elim{\lexists}
च्या गौण आधारविधानाच्या
वरच येतो. विशेषतः
$c$ हा~$\delta_3$
किंवा~$\delta_4$
मध्ये येऊ शकत नाही.
दाखवलेल्या दोन प्रतींतील $c$ हे एकच आयगेन स्थिरांक
दोन अनुमानांत वापरल्याचे अंतिम रूप मानायचे नाही:
वरच्या नियमानुसार प्रत्येक प्रतीला स्वतंत्र नवा स्थिरांक
आणि त्याच्या बद्ध गृहीतकांना सुसंगत नव्या खुणा द्यायच्या.

या उदाहरणातून आणखी
एक गोष्ट दिसते.
नव्या सिद्धता
मध्ये~$\delta_5$
च्या दोन प्रती आहेत.
म्हणून~$\delta_5$
मध्ये महत्तम कट असेल,
तर आपण कट-लांबी
\emph{कमी केलेली नाही}!
त्यामुळे~$\delta_5$
मध्ये महत्तम कट नाही
हे माहीत असेल तेव्हाच
क्रमपरिवर्तन रूपांतरण
लावण्याची काळजी
घ्यावी लागते.
नेहमी \emph{सर्वात वरचा}
महत्तम कट निवडून
हे साधतो: म्हणजे असा
महत्तम कट की त्याच्या
उप-सिद्धता मध्ये
(येथे~$\delta_1$)
उप-सिद्धता च्या
शेवटच्या अनुमानाच्या
वर संपणारा दुसरा
महत्तम कट नसतो.
यामुळे~$\delta_5$
मध्ये, किंबहुना
$\delta_2$, $\delta_3$
किंवा~$\delta_4$
मध्येही, इतर महत्तम
कट असण्याची शक्यता
नाहीशी होते.














\section{न्यूनीकरण रूपांतरणे}\label{pt:nor:red:sec}

~\Log{N1i} मधील
सिद्धता मध्ये एखाद्या
कटाची लांबी~$1$ असेल,
तर तो एका सूत्र
आस्थितीचा खंड असतो.
तीच सूत्र आस्थिती
परिचय अनुमानाचा किंवा~\FalseInt
चा निष्कर्ष आणि
निरसन अनुमानाचे
प्रमुख आधारविधान असते.
सामान्यरूपीकरणाच्या
सिद्धतेत अशा कटांचे
“न्यूनीकरण” करतात:
अशा परिचय/निरसन
जोडीत संपणारी उप-सिद्धता
त्या दोन्ही अनुमानांशिवाय
असलेल्या नव्या सिद्धता
ने बदलतात. याने नवे
कट निर्माण होऊ शकतात.
पण त्या नव्या कटांची
कोटी न्यूनीकृत कटाच्या
कोटीपेक्षा कमी असते.
त्यामुळे नव्या सिद्धता
ची कट-लांबी कमी होते
(कारण लांबी~$1$
असलेला महत्तम कट
निघून गेला), किंवा
सिद्धता ची कट-कोटी
कमी होते (संबंधित
महत्तम कट एकटाच
महत्तम कट असल्यास).

उदाहरणार्थ, संबंधित
अनुमाने~\Intro{\lif}
आणि~\Elim{\lif} असतील,
$A \ident B \lif C$
असेल, आणि~$\delta_1$
ही उप-सिद्धता
पुढीलप्रमाणे संपत असेल:
\[
\AxiomC{$\Discharge{B}{x}$}
\RightLabel{$\delta_2$}
\DeduceC{$C$}
\DischargeRule{\Intro{\lif}}{x}
\UnaryInfC{$B \lif C$}
\AxiomC{}
\RightLabel{$\delta_3$}
\DeduceC{$B$}
\RightLabel{\Elim{\lif}}
\BinaryInfC{$C$}
\DisplayProof
\]
हा कट काढण्यासाठी
~$\delta$ मधील
उपसिद्धता~$\delta_1$
पुढील रचनेने बदलतो:
\[
\AxiomC{}
\RightLabel{$\delta_3$}
\DeduceC{$B$}
\RightLabel{$\Subst{\delta_2}{\delta_3}{B^x}$}
\DeduceC{$C$}
\DisplayProof
\]
~$\Subst{\delta_2}{\delta_3}{B^x}$
म्हणजे~$\delta_2$
मधील~$x$ खूण असलेल्या
प्रत्येक गृहीतक~$B$
च्या जागी संपूर्ण
सिद्धता~$\delta_3$
ठेवून मिळणारी सिद्धता.
कलम-संयोजनापूर्वी $\delta_2$ च्या बद्ध मुक्ती-खुणा
आणि त्यांना बांधलेली गृहीतके एकत्र नवी नावे देऊन
$\delta_3$ च्या उघड्या खुणांपासून वेगळी ठेवायची.
त्यानंतर नव्या सिद्धतेची उघडी गृहीतके म्हणजे
$\delta_2$ मधील बसवलेल्या $B^x$ आस्थितींव्यतिरिक्त
उरलेली उघडी गृहीतके, आणि प्रत्यक्ष बसवलेल्या
$\delta_3$ च्या प्रतींची उघडी गृहीतके.
$\delta_3$ मध्ये स्वतः $B^x$ उघडे असू शकते;
म्हणून नव्या सिद्धतेत खूण $x$ अजिबात उरत नाही
असे म्हणता येत नाही. तसेच मुक्ती निष्फळ असेल किंवा
एखादी शाखा काढली गेली असेल, तर मूळची काही
उघडी गृहीतके गळू शकतात. आवश्यक निष्कर्ष असा की
नवी उघडी गृहीतके मूळ $\delta_1$ च्या उघड्या
गृहीतकांच्या पलीकडे जात नाहीत.
$\delta_2$ मधील मुक्ती-नियम आता आत बसवलेल्या
उघड्या गृहीतकांना अनवधानाने मुक्त करत नाहीत.
आणि~$\delta_3$
मधील गृहीतके मुक्त
करणारी अनुमाने अनेक
प्रतींमध्ये येऊ शकतात;
त्या प्रती
~$\Subst{\delta_2}{\delta_3}{B^x}$
मधील संबंधित गृहीतके
मुक्त करतात.
सर्व न्यूनीकरण नमुन्यांत हीच सुरक्षित कलम-संयोजनाची
अट वापरायची. अनेक प्रतींतील आयगेन स्थिरांक आणि
त्यांच्याशी संबंधित बद्ध खुणा स्वतंत्र नवी नावे देऊन
नियमितता जपायची. संख्यापकांच्या नमुन्यांतील $t$ हे
घोषित बंद पदे वापरण्याच्या संकेतानुसार बंद पद आहे.

~$\delta_2$ किंवा
$\delta_3$ यांच्या
आतमध्ये पूर्णपणे असलेला
कट, किंवा~$\delta$
मध्ये~$\delta_1$
बाहेर पूर्णपणे असलेला
कट, बदलत नाही.
त्यात तीच सूत्रे
आणि तीच अनुमाने
असतात; त्यामुळे त्याची
कोटी आणि लांबी
~$\delta$ मधील
अनुरूप कटासारखीच
राहते. आपण लांबी~$1$
असलेला महत्तम कट
काढला आहे; म्हणून
कट-लांबी कमी झाली,
किंवा~$\delta$
ची कट-लांबी~$1$
असून हा एकटाच महत्तम
कट होता, तर कट-कोटी
कमी झाली.

तथापि, दोन गोष्टी
घडू शकतात:
\begin{enumerate}
  \item सर्व~$B^x$
  गृहीतकांच्या जागी
  $\delta_3$ ठेवल्याने
  ~$\delta_3$ मधील
  सर्व कटांच्या प्रती
  वाढतात. त्यांपैकी
  काही महत्तम कट
  असतील, तर नव्या
  सिद्धता ची कट-लांबी
  जुन्या सिद्धता
  पेक्षा अधिक होऊ
  शकते. हे होणार
  नाही याची खात्री
  हवी. म्हणून न्यूनीकरण
  रूपांतरणे फक्त
  \emph{सर्वात उजव्या}
  महत्तम कटांनाच
  लावतो: न्यूनीकृत
  होत असलेल्या कटाच्या
  उजवीकडे~$\delta$
  मधील कोणत्याही
  शाखेवर दुसरा महत्तम
  कट नसतो. ~$\delta_3$
  मध्ये महत्तम कट असता, तर
  तो येथे न्यूनीकृत
  होणाऱ्या कटाच्या
  उजवीकडे असता;
  म्हणून हा सर्वात
  उजवा कट नसता.
  सर्वात उजवे कट
  निवडल्याने सिद्धता
  ची कट-लांबी किंवा
  अगदी कट-कोटीही
  कमी होते.
  \item ~$\delta_2$
  मधील एखादे गृहीतक~$B^x$
  निरसन अनुमानाचे
  प्रमुख आधारविधान
  असेल, आणि~$\delta_3$
  चे अंतिम अनुमान
  \Elim{\lor}, \Elim{\lexists}
  किंवा परिचय
  अनुमान असेल, तर
  त्या गृहीतकावर
  ~$\delta_3$ चे
  कलम-संयोजन केल्यामुळे
  नवा कट तयार होतो.
  ~$\delta_2$ मध्ये
  जे केवळ गृहीतक
  होते, ते आता लांबी~$1$
  असलेला कट
  (परिचय
  अनुमानाचा निष्कर्ष आणि
  निरसन अनुमानाचे
  प्रमुख आधारविधान)
  किंवा कट-खंडाची
  शेवटची सूत्र
  आस्थिती
  (निरसन
  अनुमानाचे प्रमुख
  आधारविधान आणि
  \Elim{\lor}
  किंवा~\Elim{\lexists}
  चा निष्कर्ष)
  होते. $B^x$ हे
  \Elim{\lor}
  किंवा~\Elim{\lexists}
  चे गौण आधारविधान
  असेल, तर ते आधीच
  एखाद्या खंडाची,
  कदाचित कट-खंडाचीही,
  पहिली सूत्र
  आस्थिती होती.
  नव्या सिद्धता
  मध्ये तो खंड
  आता अधिक लांब
  होऊ शकतो. पण
  या नव्या कटाचे
  किंवा विद्यमान
  कट-खंडाचे सूत्र
  हे~$B$ आहे, आणि
  $\depth{B} < \depth{B \lif C}$.
  म्हणून नवे कट
  आले किंवा काही
  कटांची लांबी
  वाढली, तरी त्यांपैकी
  कोणताही महत्तम कट
  नसतो. अशा रीतीने
  कट-कोटी कमी होते,
  किंवा समान कोटीचे
  कट उरले असतील
  तर कट-लांबी कमी
  होते.
\end{enumerate}

लांबी~$1$ असलेला
कट कमी करण्याच्या
प्रक्रियेला
\emph{न्यूनीकरण रूपांतरण}
(किंवा
\emph{वळसा रूपांतरण})
म्हणतात. काही
नियम-जोड्यांचे नमुने
~\ref{pt:nor:red:tab:reduction-conversions}
मध्ये दिले आहेत.
(उरलेले सरावासाठी
सोडले आहेत.)

\begin{table}
\begin{multline*}
\bottomAlignProof
\AxiomC{$\Discharge{B}{x}$}
\RightLabel{$\delta_2$}
\shortDeduce
\DeduceC{$C$}
\DischargeRule{\Intro{\lif}}{x}
\UnaryInfC{$B \lif C$}
\AxiomC{}
\RightLabel{$\delta_3$}
\shortDeduce
\DeduceC{$B$}
\RightLabel{\Elim{\lif}}
\BinaryInfC{$C$}
\DisplayProof
\quad \longrightarrow\quad
\shoveright{\bottomAlignProof
\AxiomC{}
\RightLabel{$\delta_3$}
\shortDeduce
\DeduceC{$B$}
\RightLabel{$\Subst{\delta_2}{\delta_3}{B^x}$}
\shortDeduce
\DeduceC{$C$}
\DisplayProof}
\\[4ex]
\shoveleft{\bottomAlignProof
\AxiomC{}
\RightLabel{$\delta_2$}
\shortDeduce
\DeduceC{$B$}
\AxiomC{}
\RightLabel{$\delta_3$}
\shortDeduce
\DeduceC{$C$}
\RightLabel{\Intro{\land}}
\BinaryInfC{$B \land C$}
\RightLabel{\Elim{\land}[1]}
\UnaryInfC{$B$}
\DisplayProof
\quad \longrightarrow\quad
\bottomAlignProof
\AxiomC{}
\RightLabel{$\delta_2$}
\shortDeduce
\DeduceC{$B$}
\DisplayProof}\qquad
\shoveright{\bottomAlignProof
\AxiomC{}
\RightLabel{$\delta_2$}
\shortDeduce
\DeduceC{$B$}
\AxiomC{}
\RightLabel{$\delta_3$}
\shortDeduce
\DeduceC{$C$}
\RightLabel{\Intro{\land}}
\BinaryInfC{$B \land C$}
\RightLabel{\Elim{\land}[2]}
\UnaryInfC{$C$}
\DisplayProof
\quad \longrightarrow\quad
\bottomAlignProof
\AxiomC{}
\RightLabel{$\delta_3$}
\shortDeduce
\DeduceC{$C$}
\DisplayProof}
\\[4ex]
\bottomAlignProof
\AxiomC{}
\RightLabel{$\delta_2$}
\shortDeduce
\DeduceC{$B$}
\RightLabel{\Intro{\lor}[1]}
\UnaryInfC{$B \lor C$}
\AxiomC{$\Discharge{B}{x}$}
\RightLabel{$\delta_3$}
\shortDeduce
\DeduceC{$D$}
\AxiomC{$\Discharge{C}{y}$}
\RightLabel{$\delta_4$}
\shortDeduce
\DeduceC{$D$}
\DischargeRule{\Elim{\lor}}{x\,y}
\TrinaryInfC{$D$}
\DisplayProof
\quad \longrightarrow\quad
\shoveright{\bottomAlignProof
\AxiomC{}
\RightLabel{$\delta_2$}
\shortDeduce
\DeduceC{$B$}
\RightLabel{$\Subst{\delta_3}{\delta_2}{B^x}$}
\shortDeduce
\DeduceC{$D$}
\DisplayProof}
\\[2ex]
\shoveleft{\bottomAlignProof
\AxiomC{}
\RightLabel{$\delta_2$}
\shortDeduce
\DeduceC{$B(c)$}
\RightLabel{\Intro{\lforall}}
\UnaryInfC{$\lforall[x][B(x)]$}
\RightLabel{\Elim{\lforall}}
\UnaryInfC{$B(t)$}
\DisplayProof
\quad \longrightarrow\quad
\bottomAlignProof
\AxiomC{}
\RightLabel{$\Subst{\delta_2}{t}{c}$}
\shortDeduce
\DeduceC{$B(t)$}
\DisplayProof}
\\[4ex]
\shoveright{\bottomAlignProof
\AxiomC{}
\RightLabel{$\delta_2$}
\shortDeduce
\DeduceC{$B(t)$}
\RightLabel{\Intro{\lexists}}
\UnaryInfC{$\lexists[x][B(x)]$}
\AxiomC{$\Discharge{B(c)}{x}$}
\RightLabel{$\delta_3$}
\shortDeduce
\DeduceC{$D$}
\DischargeRule{\Elim{\lexists}}{x}
\BinaryInfC{$D$}
\DisplayProof
\quad \longrightarrow\quad
\bottomAlignProof
\AxiomC{}
\RightLabel{$\delta_2$}
\shortDeduce
\DeduceC{$B(t)$}
\RightLabel{$\Subst{\Subst{\delta_3}{t}{c}}{\delta_2}{B(t)^x}$}
\shortDeduce
\DeduceC{$D$}
\DisplayProof}
\end{multline*}
\caption{\Log{N1i} साठी काही न्यूनीकरण रूपांतरणे.}
\label{pt:nor:red:tab:reduction-conversions}
\end{table}

\begin{prob}
\Intro{\lnot} नंतर
\Elim{\lnot} येत
असेल, तर न्यूनीकरण
रूपांतरण द्या.
\end{prob}

आणखी एक प्रसंग
पाहावा लागतो. कटाच्या
व्याख्येत असा प्रसंग
येऊ शकतो की कटाची
लांबी~$1$ आहे,
आणि~$A$ हे
\FalseInt{} चा निष्कर्ष
तसेच निरसन
अनुमानाचे प्रमुख
आधारविधान आहे.
$A$ हा
परिचय
नियमाचा निष्कर्ष
असताना करतो तसेच
येथेही करू. उदाहरणार्थ,
पुढील बदल करा:
\[
\bottomAlignProof
\AxiomC{}
\RightLabel{$\delta_2$}
\DeduceC{$\lfalse$}
\UnaryInfC{$B \lif C$}
\AxiomC{}
\RightLabel{$\delta_3$}
\DeduceC{$B$}
\RightLabel{\Elim{\lif}}
\BinaryInfC{$C$}
\DisplayProof
\quad\text{ऐवजी}\quad
\bottomAlignProof
\AxiomC{}
\RightLabel{$\delta_2$}
\DeduceC{$\lfalse$}
\UnaryInfC{$C$}
\DisplayProof
\]
~$A \ident (B \lor
C)$ किंवा
$A \ident \lexists[x][B(x)]$
असेल, तर थोडी काळजी
हवी. उदाहरणार्थ,
पुढील रचनेच्या जागी
\[
\bottomAlignProof
\AxiomC{}
\RightLabel{$\delta_2$}
\DeduceC{$\lfalse$}
\RightLabel{\FalseInt}
\UnaryInfC{$B \lor C$}
\AxiomC{$\Discharge{B}{x}$}
\RightLabel{$\delta_3$}
\DeduceC{$D$}
\AxiomC{$\Discharge{C}{y}$}
\RightLabel{$\delta_4$}
\DeduceC{$D$}
\DischargeRule{\Elim{\lor}}{x\,y}
\TrinaryInfC{$D$}
\DisplayProof
\quad\text{ऐवजी}\quad
\bottomAlignProof
\AxiomC{}
\RightLabel{$\delta_2$}
\DeduceC{$\lfalse$}
\RightLabel{\FalseInt}
\UnaryInfC{$D$}
\DisplayProof
\]
असा बदल केला,
तर कट-सूत्र~$D$
असलेला नवा कट
निर्माण होऊ शकतो.
पण
$\depth{D} < \depth{B \lor C}$
असेलच याची हमी नाही!{}
म्हणून~\Intro{\lor}/\Elim{\lor}
न्यूनीकरणात जसे
केले, तसेच येथेही
करावे आणि पुढील
रचनेने बदलावे:
\[
\bottomAlignProof
\AxiomC{}
\RightLabel{$\delta_2$}
\DeduceC{$\lfalse$}
\RightLabel{\FalseInt}
\UnaryInfC{$B$}
\RightLabel{$\delta_3$}
\DeduceC{$D$}
\DisplayProof
\quad\text{किंवा}\quad
\bottomAlignProof
\AxiomC{}
\RightLabel{$\delta_2$}
\DeduceC{$\lfalse$}
\RightLabel{\FalseInt}
\UnaryInfC{$C$}
\RightLabel{$\delta_4$}
\DeduceC{$D$}
\DisplayProof
\]












\section{सामान्यरूपीकरण प्रमेय}\label{pt:nor:thm:sec}

एखाद्या सिद्धता मध्ये
कट-खंड नसतील, तर ती
\emph{सामान्य} रूपात असते.
म्हणजे तिच्यावर क्रमपरिवर्तन
रूपांतरण किंवा न्यूनीकरण
रूपांतरण यांपैकी कोणतेही
लावता येत नाही; आणि
हेच या स्थितीचे लक्षण आहे.
एखादी सिद्धता
\emph{सामान्यरूप करणे}
म्हणजे कट शिल्लक
राहणार नाहीत तोपर्यंत
क्रमपरिवर्तन आणि न्यूनीकरण
रूपांतरणे लावून तिला
सामान्यरूप सिद्धतेत
बदलणे. हे नेहमी शक्य
असते, हाच
\emph{सामान्यरूपीकरण प्रमेय}.

\begin{thm}\label{pt:nor:thm:thm:normalization}
कोणतीही~\Log{N1i}-सिद्धता
$\delta$ सामान्यरूप करता
येते; म्हणजे तिचे कट
नसलेल्या सिद्धता~$\delta'$
मध्ये रूपांतर करता येते.
\end{thm}

\begin{proof}
  ~$\delta$ ची कट-कोटी
  आणि कट-लांबी यांवर
  विगमन करू. विगमन
  गृहीतक असे: कोणतीही
  सिद्धता जिची
  कट-कोटी कमी असेल,
  किंवा कट-कोटी समान
  पण कट-लांबी कमी असेल,
  सामान्यरूप करता येते.

  कट-लांबी~$0$
  असेल, तर कट नसतात
  आणि~$\delta$ आधीच
  सामान्यरूप असते.
  म्हणून~$\delta$ ची
  कट-कोटी आणि कट-लांबी
  दोन्ही~$>0$ आहेत असे
  समजू. सर्वात वरचा
  आणि सर्वात उजवा
  महत्तम कट-खंड निवडा.
  त्या कट-खंडाची लांबी
  $>1$ असेल, तर
  क्रमपरिवर्तन रूपांतरण
  लावा. लांबी~$1$
  असेल, तर संबंधित
  न्यूनीकरण रूपांतरण लावा.
  याने मिळणाऱ्या नव्या
  सिद्धता ची कट-कोटी
  तीच राहून कट-लांबी
  कमी होते, किंवा
  कट-कोटीच कमी होते;
  म्हणून विगमन गृहीतक
  लागू होते.
\end{proof}

वरील सिद्धतेतील
पद्धत—नेहमी सर्वात
वरचा आणि सर्वात
उजवा महत्तम कट
हाताळणे—सामान्यरूप
सिद्धता मिळवण्यासाठी
रूपांतरणे एका ठरावीक
क्रमाने लावावी लागतात
असे सुचवते. पण आश्चर्य
म्हणजे क्रम महत्त्वाचा
नसतो. क्रमपरिवर्तन आणि
न्यूनीकरण रूपांतरणे
कोणत्याही क्रमाने
लावली तरी शेवटी
सामान्यरूप सिद्धता
मिळते.

खंड आणि कट यांच्या
व्याख्या~\Log{N2i}
साठीही सहज देता येतात;
आणि~\Log{N1i} मधील
क्रमपरिवर्तन व न्यूनीकरण
रूपांतरणे थेट~\Log{N2i}
मध्ये रूपांतरित होतात.
त्यामुळे सामान्यरूपीकरण
प्रमेय~\Log{N2i} साठीही
खरे आहे.

\begin{thm}\label{pt:nor:thm:thm:normalization-N2i}
कोणतीही~\Log{N2i}-सिद्धता
$\delta$ सामान्यरूप करता
येते; म्हणजे तिचे कट
नसलेल्या सिद्धता~$\delta'$
मध्ये रूपांतर करता येते.
\end{thm}















\section{सामान्यरूप \Log{N2i}-सिद्धता आणि \Log{G2i} यांतील रूपांतरे}\label{pt:nor:tr:sec}

\Log{N2i} मधील सिद्धता चे \Log{G2i} मधील सिद्धता मध्ये
आणि उलट दिशेने रूपांतर करता येते. प्रथम,
कट-विरहित \Log{G2i} सिद्धता चे रूपांतर केल्यावर
सामान्यरूप \Log{N2i}-सिद्धता मिळतात.

हा निष्कर्ष मांडण्यासाठी~\Log{N2i} आणि~\Log{G2i}
यांतील क्रमवर्तींचे पूर्वपक्ष जोडणारी काही परिभाषा हवी.
खूण लावलेल्या सूत्रांचा संच~$\Gamma'$ हा खूण नसलेल्या
सूत्रांच्या बहुसंच~$\Gamma$ शी \emph{अनुरूप} असतो,
जर प्रत्येक सूत्र~$A$ साठी~$\Gamma'$ मधील
$x : A$ या रूपातील घटक ची संख्या ही~$\Gamma$
मधील~$A$ च्या प्रतींच्या संख्येहून जास्त नसेल.

\begin{cor}
$\Log{G2i} \Proves \Gamma \Sequent \Delta$ असेल, तर
$\Gamma' \Sequent C$ ची सामान्यरूप
\Log{N2i}-सिद्धता~$\delta$ असते; येथे खूण लावलेल्या
सूत्रांचा संच~$\Gamma'$ हा बहुसंच~$\Gamma$ शी अनुरूप असतो.
$\Delta=\emptyset$ असल्यास~$C\ident\lfalse$,
आणि अन्यथा~$\Delta=\{C\}$ असते.
म्हणजे, प्रत्येक सूत्र~$A$ साठी~$\Gamma'$ मधील
$x : A$ या रूपातील घटक ची संख्या ही~$\Gamma$
मधील~$A$ च्या प्रतींच्या संख्येहून (म्हणजे~$A$ च्या
पुनरावृत्तींच्या संख्येहून) जास्त नसते.
\end{cor}

\begin{proof}
\ref{pt:nat:tgi:prop:G2i-to-N2i} च्या सिद्धतेतील
रचनेकडे पाहिल्यास दिसते: परिचय नियम
सिद्धता च्या शेवटी, तर निरसन नियम
फक्त सिद्धता च्या सुरुवातीला जोडतो.
म्हणून निरसन नियमाच्या वर परिचय
नियम कधी येत नाही.
\end{proof}

मात्र~\ref{pt:nat:tgi:prop:G2i-to-N2i} मधील
\Cut{} नियमाचे रूपांतर~\Log{N2i}-सिद्धता~$\delta$ मध्ये
कट निर्माण करू शकते. एखादी~\Log{G2i} सिद्धता
शेवटी~\Cut{} ने संपत असेल, आणि तिच्या तत्काळ उप-सिद्धता
$\pi_1$ व~$\pi_2$ अनुक्रमे~$\Gamma_1 \Sequent A$
आणि~$A, \Gamma_2 \Sequent \Delta_2$ यांच्या असतील,
तर~$\pi_1$ चे रूपांतर~$\delta_1$ हे~$\pi_2$ च्या
रूपांतर~$\delta_2$ मधील $x:A \Sequent A$
स्वयंसिद्धकांवर कलम-संयोजित करतो. जर~$\delta_1$
मुख्य सूत्र~$A$ असलेल्या परिचय
नियमाने संपत असेल आणि~$\delta_2$ मधील
$x:A \Sequent A$ स्वयंसिद्धक हे निरसन
नियमाचे प्रमुख आधारविधान असेल, तर परिणामी~$\delta$
मध्ये कट तयार होतो.

\Log{N2i} मधील सिद्धता चे~\Log{G2i} मधील
सिद्धता मध्येही रूपांतर करता येते.
\ref{pt:nat:tni:prop:N2-to-G2} च्या सिद्धतेत
\Elim{\land}, \Elim{\lif}, \Elim{\lnot},
\Elim{\lforall} यांचे रूपांतर परिणामी~\Log{G2i}
सिद्धता मध्ये~\Cut{} अनुमाने निर्माण करते.
पण सुरुवातीची~\Log{N2i}-सिद्धता सामान्यरूप असेल,
तर असे कट टाळता येतात.

~\Log{N2i}-सिद्धता~$\delta$ मधील
क्रमवर्तींच्या आस्थितींच्या $S_1$, \dots,~$S_n$
या यादीला \emph{शाखा} म्हणू, जर~$S_1$ स्वयंसिद्धक
असेल आणि~$S_i$ जेव्हा एखाद्या अनुमानाचे प्रमुख
आधारविधान असेल तेव्हा~$S_{i+1}$ हा त्याचा निष्कर्ष
असेल. दुसऱ्या शब्दांत, शाखा स्वयंसिद्धकापासून
~$\delta$ मध्ये खाली जात क्रमवर्तींचा मागोवा घेते;
पण निरसन नियमांच्या गौण आधारविधानांवर थांबते.
जिचा~$S_n$ अंतिम क्रमवर्ती असतो ती~$\delta$ ची
\emph{मुख्य शाखा}. प्रत्येक सिद्धता~$\delta$ ला
किमान एक मुख्य शाखा असते, कारण प्रत्येक नियमाला
किमान एक प्रमुख आधारविधान असते (फक्त~\Intro{\land}
ला दोन असतात).

\begin{prop}\label{pt:nor:tr:prop:normal-N2i-to-G2i}
~$\delta$ ही~$\Gamma \Sequent A$ ची सामान्यरूप
$\Log{N2c}$ किंवा~$\Log{N2i}$-सिद्धता असेल,
तर~$\Gamma' \Sequent \Delta$ ची कट-विरहित
$\Log{G2c}$-सिद्धता (अनुक्रमे~$\Log{G2i}$-सिद्धता)~$\pi$
असते. येथे~$\Gamma$ मधील खुणा काढून मिळणारा
सूत्रांचा बहुसंच~$\Gamma'$ आहे;
$A$ हा~$\lfalse$ नसेल तर~$\Delta = \{A\}$,
आणि असेल तर~$\Delta = \emptyset$.
\end{prop}

\begin{proof}
यावेळी~$\Gamma \Sequent A$ च्या सिद्धता~$\delta$
मधील अनुमानांच्या संख्येवर विगमन करतो.
~$\delta$ मध्ये अनुमान नसेल, तर ती
$x:A \Sequent A$ हा स्वयंसिद्धक आहे;
तेव्हा~$\pi$ हा अनुरूप~$A \Sequent A$
स्वयंसिद्धक आहे, किंवा~$A \ident \lfalse$
असेल तर~$\lfalse \Sequent \ $ स्वयंसिद्धक आहे.

~$\delta$ चा शेवट अनुमाननियमात होत असेल,
तर पुढील प्रसंग वेगळे पाहू:
\begin{enumerate}
  \item $R$ हा परिचय नियम किंवा~\FalseInt
  असेल, तर रूपांतर~\ref{pt:nat:tni:prop:N2-to-G2}
  च्या सिद्धतेप्रमाणेच होते आणि~\Cut लागत नाही.

  \item ~$\delta$ चा शेवटचा नियम~$R$ हा निरसन नियम
  असेल, तर पुढील बाबी लक्षात घ्या:
  \begin{enumerate}
    \item ~$\delta$ ची मुख्य शाखा कोणत्याही
    परिचय नियमातून जात नाही; अन्यथा त्या
    शाखेवर परिचय नियम किंवा~\FalseInt{}
    यानंतर निरसन नियम येईल आणि~$\delta$
    सामान्यरूप राहणार नाही.
    \item ~$\delta$ च्या मुख्य शाखांवर
    परिचय नियम नसल्याने एकच मुख्य शाखा
    असते. (फक्त~\Intro{\land} ला दोन प्रमुख
    आधारविधाने आहेत; म्हणून अनेक मुख्य शाखा असतील,
    तर त्या~\Intro{\land} च्या आधारविधानांतून जाव्या लागतील.)
    \item मुख्य शाखेत~\Elim{\lor} किंवा~\Elim{\lexists}
    चे प्रमुख आधारविधान असेल, तर असा एकच नियम
    असतो आणि तो शेवटचा नियम~$R$ असतो. (नसल्यास
    शाखेत~\Elim{\lor} किंवा~\Elim{\lexists}
    चा निष्कर्ष पुढील निरसन नियमाचे प्रमुख
    आधारविधान होईल; क्रमपरिवर्तन रूपांतरण शक्य असल्यामुळे
    सिद्धता सामान्यरूप राहणार नाही.)
  \item \Elim{\lor} आणि~\Elim{\lexists} यांखेरीज
  कोणताही निरसन नियम गृहीतके निरस्त करत नाही,
  म्हणजे क्रमवर्तींचा डावा संदर्भ बदलत नाही.
  \Elim{\lor} आणि~\Elim{\lexists} हे दोन्ही नियम
  फक्त गौण आधारविधानाच्या वरची
  गृहीतके निरस्त करतात; म्हणजे फक्त गौण
  आधारविधानांचे डावे संदर्भ बदलतात.
  त्यामुळे~$\delta$ च्या मुख्य शाखेवरील
  क्रमवर्ती~$S_i$ चा संदर्भ~$\Gamma_1$ असेल,
  तर~$S_i$ हे प्रमुख आधारविधान असलेल्या अनुमानाच्या
  निष्कर्ष~$S_{i+1}$ चा संदर्भ~$\Gamma_1' \supseteq \Gamma_1$
  असतो.
  \item म्हणून मुख्य शाखेतील सर्वांत वरचा
  क्रमवर्ती~$S_1$ हा~$x: C \Sequent C$
  असेल, तर~$x:C \in \Gamma$.
  \end{enumerate}
  आता~$C$ च्या रूपानुसार प्रसंग वेगळे पाहू.
  मुख्य शाखेवरील प्रत्येक~$S_i$ हे निरसन
  नियमाचे प्रमुख आधारविधान असल्याने
  $x:C \Sequent C$ बद्दलही तेच खरे आहे.

  \item $C \ident D \land E$ असेल, तर~$\delta$
  चे रूप असे आहे:
  \[
  \Axiom$x:D \land E \fCenter D \land E$
  \RightLabel{\Elim{\land}[1]}
  \UnaryInf$x:D \land E \fCenter D$
  \Deduce$x:D \land E, \Gamma_1 \fCenter A$
  \DisplayProof
  \]
  मुख्य शाखेत~$x:D \land E \Sequent D$ आहे.
  त्या शाखेत~\Intro{\lif}, \Intro{\lnot} यांचे
  प्रमुख आधारविधान किंवा~\Elim{\lor}, \Elim{\lexists}
  यांचे गौण आधारविधान येत नाही; म्हणून
  $x:D \land E$ मधील संदर्भ-सूत्रे
  मुख्य शाखेवर बदलत नाहीत. त्यामुळे अंतिम
  क्रमवर्तीच्या संदर्भात~$x:D \land E$ असलेच पाहिजे.
  शिवाय,~\Elim{\land} ने संपणाऱ्या उप-सिद्धता
  ऐवजी~$w:D \Sequent D$ ठेवून, आणि मुख्य
  शाखेवरील प्रत्येक क्रमवर्तीच्या संदर्भात
  $x:D \land E$ ऐवजी~$w:D$ ठेवून,
  सिद्धता~$\delta$ चे सिद्धता~$\delta_1$
  मध्ये रूपांतर करता येते; म्हणजे:
  \[
  \bottomAlignProof
  \Axiom$x:D \land E \fCenter D \land E$
  \RightLabel{\Elim{\land}[1]}
  \UnaryInf$x:D \land E \fCenter D$
  \Deduce$x:D \land E, \Gamma_1 \fCenter A$
  \DisplayProof
  \quad\text{ यात }\quad
  \bottomAlignProof
  \Axiom$w:D \fCenter D$
  \Deduce$w:D, \Gamma_1 \fCenter A$
  \DisplayProof
  \]
  विगमन गृहीतक~$\delta_1$ ला लागू होते,
  कारण~$\delta$ मधील अनुमानांची संख्या ही
  $\delta_1$ मधील अनुमानांच्या संख्येपेक्षा~$1$
  ने अधिक आहे.
येथे $w$ ही संपूर्ण $\delta$ मध्ये न येणारी नवी खूण
घ्यायची. संदर्भ हे प्रत्येक नियमाच्या आधारसंदर्भांपासून
पुन्हा मोजायचे; इतर शाखेतले $x:D\land E$ किंवा
इतर गृहीतक काढायचे नाही. $\Gamma_1$ मध्ये त्या
इतर शाखांतून राहणारी सर्व गृहीतके समाविष्ट आहेत.

  विगमन गृहीतकानुसार,
  $D, \Gamma_1' \Sequent \Delta$ ची
  \Log{G2i}-सिद्धता~$\pi_1$ असते.
  \LeftR{\land} लावून~सिद्धता~$\pi$ मिळते:
  \[
  \AxiomC{}
  \RightLabel{$\pi_1$}
  \Deduce$D, \Gamma_1' \fCenter \Delta$
  \RightLabel{\LeftR{\land}}
  \UnaryInf$D \land E, \Gamma_1' \fCenter \Delta$
  \DisplayProof
  \]
  $A \ident \lfalse$ असेल, तर उत्तरपक्ष रिकामाच
  ठेवतो; उजवीकडील दुर्बलीकरण लागत नाही:
  \[
  \AxiomC{}
  \RightLabel{$\pi_1$}
  \Deduce$D, \Gamma_1' \fCenter $
  \RightLabel{\LeftR{\land}}
  \UnaryInf$D \land E, \Gamma_1' \fCenter $
  \DisplayProof
  \]
  \item $C \ident D \lor E$ असेल,
  तर ते~\Elim{\lor} चे प्रमुख आधारविधान असले
  पाहिजे आणि हे अनुमान मुख्य शाखेवरील शेवटचे
  अनुमान~$R$ असले पाहिजे. म्हणजे~$\delta$ चे
  रूप असे आहे:
  \[
  \Axiom$x:D \lor E \fCenter D \lor E$
  \AxiomC{}
  \RightLabel{$\delta_1$}
  \Deduce$y:D, \Gamma_1 \fCenter A$
  \AxiomC{}
  \RightLabel{$\delta_2$}
  \Deduce$z:E, \Gamma_2 \fCenter A$
  \RightLabel{\Elim{\lor}}
  \TrinaryInf$x:D \lor E, \Gamma_1, \Gamma_2 \fCenter A$
  \DisplayProof
  \]
  विगमन गृहीतक~सिद्धता $\delta_1$ आणि~$\delta_2$
  यांना लागू होते; त्यामुळे अनुक्रमे
  $D, \Gamma_1' \Sequent \Delta$ आणि
  $E, \Gamma_2' \Sequent \Delta$ यांच्या
  \Log{G2i}-सिद्धता $\pi_1$ व~$\pi_2$ मिळतात.
  त्यांच्यावर~$\LeftR{\lor}$ लावून~$\pi$ मिळते:
  \[
  \AxiomC{}
  \RightLabel{$\pi_1$}
  \Deduce$D, \Gamma_1' \fCenter \Delta$
  \AxiomC{}
  \RightLabel{$\pi_2$}
  \Deduce$E, \Gamma_2' \fCenter \Delta$
  \RightLabel{\LeftR{\lor}}
  \BinaryInf$D \lor E, \Gamma_1', \Gamma_2' \fCenter \Delta$
  \DisplayProof
  \]
  $\Elim{\lor}$ ने~$y: D$ किंवा~$z: E$
  निरस्त केले नसेल (म्हणजे ही गृहीतके गौण
  आधारविधानांच्या संदर्भात नसतील), तर ती आवश्यक
  \LeftR{\Weakening} ने जोडायची. $A\ident\lfalse$
  असेल, तर उत्तरपक्ष रिकामाच ठेवायचा; त्यात असत्य
  जोडायचे नाही.
  उदाहरणार्थ:
  \[
  \AxiomC{}
  \RightLabel{$\pi_1$}
  \Deduce$\Gamma_1' \fCenter $
  \RightLabel{\LeftR{\Weakening}}
  \UnaryInf$D, \Gamma_1' \fCenter $
  \AxiomC{}
  \RightLabel{$\pi_2$}
  \Deduce$\Gamma_2' \fCenter $
  \RightLabel{\LeftR{\Weakening}}
  \UnaryInf$E, \Gamma_2' \fCenter $
  \RightLabel{\LeftR{\lor}}
  \BinaryInf$D \lor E, \Gamma_1', \Gamma_2' \fCenter $
  \DisplayProof
  \]

  \item $C \ident D \lif E$ असेल,
  तर~$\delta$ चे रूप असे आहे:
  \[
  \Axiom$x:D \lif E \fCenter D \lif E$
  \AxiomC{}
  \RightLabel{$\delta_1$}
  \Deduce$\Gamma_1 \fCenter D$
  \RightLabel{\Elim{\lif}}
  \BinaryInf$x:D \lif E, \Gamma_1 \fCenter E$
  \RightLabel{$\delta_2$}
  \Deduce$x:D \lif E, \Gamma_1, \Gamma_2 \fCenter A$
  \DisplayProof
  \]
  येथे~$x:D \lif E, \Gamma_1 \Sequent E$
  समाविष्ट करणाऱ्या मुख्य शाखेत~\Intro{\lif},
  \Intro{\lnot} यांचे प्रमुख आधारविधान किंवा
  \Elim{\lor}, \Elim{\lexists} यांचे गौण
  आधारविधान नसते. त्यामुळे~$x:D \lif E, \Gamma_1$
  मधील संदर्भ-सूत्रे मुख्य शाखेवर बदलत
  नाहीत. म्हणून अंतिम क्रमवर्तीच्या संदर्भात
  $x:D \lif E, \Gamma_1$ असले पाहिजे.
  शिवाय,~\Elim{\lif} ने संपणाऱ्या उप-सिद्धता
  ऐवजी~$w:E \Sequent E$ ठेवून आणि मुख्य
  शाखेवरील प्रत्येक क्रमवर्तीच्या संदर्भात
  $x:D \lif E, \Gamma_1$ ऐवजी~$w:E$
  ठेवून,~सिद्धता च्या~$\delta_2$ या भागाचे
  योग्य सिद्धता~$\delta_2'$ मध्ये रूपांतर
  करता येते; म्हणजे:
  \[
  \bottomAlignProof
  \Axiom$x:D \lif E, \Gamma_1 \fCenter E$
  \RightLabel{$\delta_2$}
  \Deduce$x:D \lif E, \Gamma_1, \Gamma_2 \fCenter A$
  \DisplayProof
  \quad\text{ यात }\quad
  \bottomAlignProof
  \Axiom$w:E \fCenter E$
  \RightLabel{$\delta_2'$}
  \Deduce$w:E, \Gamma_2 \fCenter A$
  \DisplayProof
  \]
  विगमन गृहीतक~$\delta_1$ आणि~$\delta_2'$
  यांना लागू होते; कारण~$\delta$ मधील
  अनुमानांची संख्या ही~$\delta_1$ व~$\delta_2$
  मधील अनुमानांच्या संख्यांची बेरीज अधिक~$1$
  इतकी आहे. हा एक अधिकचा टप्पा मुख्य शाखेच्या सुरुवातीच्या
  \Elim{\lif} अनुमानासाठी आहे. (तेच अनुमान
  $\delta$ मधील शेवटचे अनुमान~$R$ असेल,
  तर~$\delta_2$ मध्ये~$0$ अनुमाने असतात.)
याही प्रसंगी $w$ ही $\delta$ मधील सर्व खुणांपासून
वेगळी नवी खूण असावी. मुख्य शाखेवरील संदर्भ नियमांनुसार
पुन्हा मोजताना इतर शाखांची गृहीतके जपायची;
$\Gamma_2$ मध्ये ती सर्व राहणारी गृहीतके येतात.
मूळ $\Gamma_1$ शी सामाईक असलेली गृहीतके नुसती
वजा करून त्या शाखांतून काढायची नाहीत.

  विगमन गृहीतकानुसार $\pi_1$ ही
  $\Gamma_1'\Sequent\Delta_D$ ची आणि $\pi_2$ ही
  $E,\Gamma_2'\Sequent\Delta$ ची \Log{G2i}-सिद्धता आहे;
  येथे $D\ident\lfalse$ असेल तर $\Delta_D=\emptyset$,
  अन्यथा $\Delta_D=\{D\}$ आहे. आवश्यक असल्यास
  पहिल्या सिद्धतेला उजवीकडील दुर्बलीकरण जोडून
  $\Gamma_1'\Sequent D$ ची $\pi_1'$ मिळते.
  \LeftR{\lif} लावून~सिद्धता~$\pi$ मिळते:
  \[
  \AxiomC{}
  \RightLabel{$\pi_1'$}
  \Deduce$\Gamma_1' \fCenter D$
  \AxiomC{}
  \RightLabel{$\pi_2$}
  \Deduce$E, \Gamma_2' \fCenter \Delta$
  \RightLabel{\LeftR{\lif}}
  \BinaryInf$D \lif E, \Gamma_1', \Gamma_2' \fCenter \Delta$
  \DisplayProof
  \]
  $D \ident \lfalse$ असेल, तर पहिल्या आधाराच्या
  रिकाम्या उत्तरपक्षात~$D$ जोडण्यासाठी एक
  \RightR{\Weakening} लावतो. $A \ident \lfalse$
  असेल, तर दुसऱ्या आधाराचा आणि अंतिम क्रमवर्तीचा
  उत्तरपक्ष रिकामाच ठेवतो. दोन्ही असत्य असताना नमुना असा:
  \[
  \AxiomC{}
  \RightLabel{$\pi_1$}
  \Deduce$\Gamma_1' \fCenter $
  \RightLabel{\RightR{\Weakening}}
  \UnaryInf$\Gamma_1' \fCenter D$
  \AxiomC{}
  \RightLabel{$\pi_2$}
  \Deduce$E, \Gamma_2' \fCenter $
  \RightLabel{\LeftR{\lif}}
  \BinaryInf$D \lif E, \Gamma_1', \Gamma_2' \fCenter $
  \DisplayProof
  \]

सामाईक खुणांची गृहीतके दोन आधारसंदर्भांत आली असल्यास
त्यांच्या खुणा काढल्यावर बहुसंचात जादा प्रती येऊ शकतात.
आवश्यक $\LeftR{\Contraction}$ लावून मूळ संचाच्या
खुणा काढून मिळणारा नेमका बहुसंच ठेवायचा.
\end{enumerate}
उरलेले प्रसंग याच प्रकारे हाताळता येतात;
ते सरावासाठी ठेवले आहेत.
\end{proof}

\begin{cor}
सामान्यरूप~\Log{N1i}-सिद्धता मध्ये येणारे प्रत्येक
सूत्र हे $\lfalse$ असते किंवा अंतिम सूत्राच्या
अथवा उघड्या गृहीतकाच्या एखाद्या उप-सूत्राचे
उदाहरण असते. सामान्यरूप~\Log{N2i}-सिद्धता साठी
अंतिम क्रमवर्तीतील सर्व सूत्रे घ्यायची.
येथे संख्यापकांमुळे उपसूत्रात मुक्त झालेल्या चलांच्या
जागी पदे बसवून मिळणारी उदाहरणेही समाविष्ट आहेत.
विधानीय प्रकरणात आदेशनाची ही अट लागत नाही.
\end{cor}

\begin{proof}
  \ref{pt:nor:tr:prop:normal-N2i-to-G2i} मधील अंतःप्रज्ञावादी
  रचना सामान्यरूप~\Log{N2i}-सिद्धता चे कट-विरहित
  \Log{G2i}-सिद्धता मध्ये रूपांतर करते.
  या रचनेत प्रत्येक असत्याव्यतिरिक्त सूत्र अनुरूप
  क्रमवर्तीच्या सिद्धतेत राहते. असत्य मात्र रिकाम्या
  उत्तरपक्षाने दर्शवले जाऊ शकते.
  कट-विरहित क्रमवर्ती नियम उलट तपासल्यास विधानीय
  आधारसूत्रे ही निष्कर्षातील सूत्रांची उपसूत्रे असतात;
  संख्यापक नियमांत त्यांच्या पद-आदेशनाची उदाहरणे येतात.
  त्यामुळे उपसूत्रांचा येथे सांगितलेला गुणधर्म मिळतो.
  \Log{N1i} चे~\Log{N2i} मध्ये आधीचे रूपांतर सूत्रे
  आणि उघडी गृहीतके जपते, म्हणून त्यालाही हा निष्कर्ष लागू होतो.
\end{proof}



\chapter{सिद्धता-शोध}\label{pt:ps:chap}










\section{प्रस्तावना}\label{pt:ps:int:sec}

\Log{G3c} किंवा~\Log{N1i} सारख्या सिद्धता
पद्धतींची स्वयंसिद्धके आणि अनुमाननियम ठरवतात की
क्रमवर्ती किंवा सूत्रे यांच्या कोणत्या वृक्षांना
योग्य सिद्धता म्हणायचे. असा क्रमवर्तींचा किंवा
सूत्रांचा वृक्ष दिला असेल (आवश्यक असल्यास
प्रत्येक टप्प्यावर कोणता नियम लागू झाला, किंवा
गृहीतकांच्या आणि त्यांना निरस्त करणाऱ्या अनुमाननियमांच्या
खुणा यांसारखी अतिरिक्त माहितीही असेल), तर
स्वयंसिद्धके आणि नियम यांच्या व्याख्यांवरून तो
वृक्ष योग्य सिद्धता आहे का हे तपासता येते.
पण दिलेल्या क्रमवर्तीची किंवा दिलेल्या गृहीतकांपासून
एखाद्या सूत्राची सिद्धता \emph{शोधणे}
हा वेगळा आणि बहुतेकदा अधिक कठीण प्रश्न आहे.
हाच \emph{सिद्धता-शोध} प्रश्न.

सिद्धता-शोधासाठी एक अगदी सोपी पद्धत आहे.
सूत्रे यांना सांत वर्णमालेतील चिन्हांचे अनुक्रम
समजू शकतो. क्रमवर्ती आणि क्रमवर्तींचे किंवा
सूत्रांचे वृक्षही सूत्रांचे अनुक्रम
मानता येतात (वृक्षाची शाखा दाखवण्यासाठी विशेष
कंस वापरून). चिन्हांचा एखादा अनुक्रम योग्य सिद्धता
दर्शवतो का, हे स्वयंसिद्धके आणि नियम यांमुळे
यांत्रिकरीत्या तपासता येते. परिणामी, सर्व \emph{शक्य}
योग्य सिद्धता यांत्रिकरीत्या निर्माण करता येतात:
सिद्धता दर्शवणाऱ्या वर्णमालेतील चिन्हांचे सर्व
अनुक्रम तयार करून प्रत्येक उमेदवार अनुक्रम योग्य
सिद्धता आहे का ते तपासा. ही सोपी सिद्धता-शोध
पद्धत अशी: दिलेल्या क्रमवर्तीची सिद्धता मिळेपर्यंत
सर्व योग्य सिद्धता निर्माण करा. (किंवा दिलेल्या
गृहीतकांपैकी काही मुक्त गृहीतके असलेल्या दिलेल्या
सूत्राची सिद्धता मिळेपर्यंत तसे करा.)

याहून चांगली पद्धत म्हणजे हाताने सिद्धता
शोधताना वापरायची सहज पद्धत: अंतिम क्रमवर्तीपासून
सुरुवात करा, एक सूत्र निवडा, आणि अनुमाननियम
``उलट्या'' दिशेने लावून एक किंवा अनेक आधारविधाने
निर्माण करा. त्यांची सिद्धता असतील, तर शेवटच्या
अनुमानासह त्यांतून दिलेल्या क्रमवर्तीची सिद्धता
बनते. नैसर्गिक निगमन पद्धतीत सिद्धता
शोधण्यासाठीही अशीच, पण अधिक गुंतागुंतीची, पद्धत
वापरली जाईल.

मात्र ही पद्धत हमखास यशस्वी नाही:
चुकीचा नियम किंवा सूत्रांचा चुकीचा क्रम
निवडल्यास मार्ग बंद होऊ शकतो. ती पूर्णपणे
निर्दिष्टही नाही: प्रत्येक टप्प्यावर निवडण्यासाठी
एकापेक्षा अधिक सूत्र किंवा उलट्या दिशेने
लावण्यासाठी एकापेक्षा अधिक नियम असू शकतात.
\emph{सिद्धता-शोध अल्गोरिदम} म्हणजे पूर्णपणे
निर्दिष्ट अशी प्रक्रिया की सिद्धता अस्तित्वात
असेल तर ती शोधून काढते (म्हणजे ती हमखास यशस्वी आहे).

काही सिद्धता पद्धतींसाठी सिद्धता-शोध
अल्गोरिदम इतरांपेक्षा सोपे असतात. क्रमवर्ती-पद्धतीत
सूत्राचा मुख्य संकारक आणि ते डावीकडे
की उजवीकडे आहे यावर उलट्या दिशेने कोणता नियम
लावायचा हे ठरते. उदाहरणार्थ,~$A \land B$
उजवीकडे असलेल्या~$\Gamma \Sequent \Delta, A \land B$
ची सिद्धता शोधायची असेल, तर~\RightR{\land}
उलटा लावून अनुरूप दोन आधारविधाने
$\Gamma \Sequent \Delta, A$ आणि
$\Gamma \Sequent \Delta, B$ निर्माण करावीत.
पण पद्धत~\Log{G1c} असेल, तर आकुंचनामुळे अधिक
पर्याय मिळतात आणि काम कठीण होते: उलट्या दिशेने
\RightR{\Contraction} लावून
$\Gamma \Sequent \Delta, A \land B, A \land B$
हे आधारविधानही तयार करता येईल. मग
$\Gamma \Sequent \Delta, A \land B, A \land B, A \land B$
हे, आणि असे पुढे. \Log{G2c} मध्ये अडचण अधिक आहे.
\RightR{\land} उलटा लावताना~$\Gamma_1$, $\Gamma_2$
हे~$\Gamma = \Gamma_1 \cup \Gamma_2$ पूर्ण करतील
असा अंदाज, आणि~$\Delta_1$, $\Delta_2$ हे
$\Delta = \Delta_1 \cup \Delta_2$ पूर्ण करतील
असा अंदाज बांधून,~$\Gamma_1 \Sequent \Delta_1, A$
आणि~$\Gamma_2 \Sequent \Delta_2, B$ ही
आधारविधाने वापरून पाहावी लागतात. चुकीच्या अंदाजाने
पुढे मार्ग बंद झाला, तर सुरुवातीला परत जाऊन
वेगळे~$\Gamma_1$, $\Gamma_2$, $\Delta_1$, $\Delta_2$
निवडावे लागतात. सूत्र हे~$A \lor B$
असेल, तर~\RightR{\lor} चे दोन पर्यायांपैकी एक
निवडून~$\Gamma \Sequent \Delta, A$ किंवा
$\Gamma \Sequent \Delta, B$ या दोनपैकी
एक आधारविधान निर्माण करावे लागते. चुकीच्या
निवडीवर पुन्हा मागे फिरावे लागते.

\Log{G3c} मध्ये या अडचणी नाहीत.
तिथे संरचनात्मक नियम नाहीत; मुख्य संकारक
आणि सूत्राची बाजू यांवर अनुमाननियम ठरतो.
म्हणून~\Log{G3c} हे~\Log{G1c} किंवा~\Log{G2c}
पेक्षा सिद्धता-शोधासाठी अधिक सोयीचे आहे.
यशस्वी शोधातून~\Log{G3c} मधील सिद्धता मिळते;
नंतर त्या सिद्धता चे~\Log{G1c} किंवा~\Log{G2c}
(किंवा~\Log{N1c}) मध्ये रूपांतर करता येते.
त्यामुळे~\Log{G3c} साठी सिद्धता-शोध अल्गोरिदम
आणि~\Log{G3c} सिद्धता चे इतर पद्धतींतील
सिद्धता मध्ये रूपांतर करणारे अल्गोरिदम मिळाले,
तर त्या इतर पद्धतींचा सिद्धता-शोध प्रश्नही सुटतो.

सिद्धता-शोध अल्गोरिदम हमखास यशस्वी आहे
हे कसे समजेल? त्याने सिद्धता शोधली नाही,
तर सिद्धता अस्तित्वातच नाही हे दाखवावे लागेल.
कारण अल्गोरिदम चुकीचा निवडला, तर सिद्धता
असूनही ती न सापडण्याचे प्रसंग येऊ शकतात.
अगदी सोप्या पद्धतीत ही अडचण नसते: ती सर्व
शक्य सिद्धता तपासते. पण सिद्धता-शोध
अल्गोरिदम कार्यक्षम असावा आणि शक्य तितके कमी
प्रयत्न वापरावेत अशी अपेक्षा असते.

पुढील बंद-पद शोध आणि त्याच्या पद-प्रतिमानाची सिद्धता
वाक्यांच्या क्रमवर्तींसाठी मांडली आहे; खालील पूर्तिचे
वर्णन या व्याप्तीत वाचायचे.
सिद्धता-शोध अल्गोरिदम हमखास यशस्वी आहे
हे दाखवण्याची मुख्य कल्पना अशी: शोध अयशस्वी
झाल्यास त्याच्या अपयशाच्या स्वरूपावरून त्या
क्रमवर्तीची सिद्धता अस्तित्वात असू शकत नाही हे दाखवा.
हे दाखवण्यासाठी ती वैध नाही असे दाखवू.
अभिजात प्रथम-क्रम तर्कशास्त्रात
$\Gamma \Sequent \Delta$ ही क्रमवर्ती वैध
असण्याचा अर्थ असा की प्रत्येक संरचना
मध्ये~$\Gamma$ मधील किमान एक सूत्र
असत्य असेल किंवा~$\Delta$ मधील किमान एक
सूत्र सत्य असेल. \Log{G3c} \emph{निर्दोष}
आहे; म्हणजे तो फक्त वैध क्रमवर्ती सिद्ध करतो.
हे सिद्धता वर विगमन करून सिद्ध होते:
स्वयंसिद्धके स्पष्टपणे वैध आहेत आणि नियमांची
आधारविधाने वैध असतील तर निष्कर्षही वैध असतो.
सिद्धता-शोध अल्गोरिदम हमखास यशस्वी आहे हे
दाखवण्यासाठी,~$\Gamma \Sequent \Delta$ ची
सिद्धता शोधण्यात अपयश आल्यास अशी
संरचना परिभाषित करता येते हे दाखवतो
की तिच्यात~$\Gamma$ मधील सर्व सूत्रे
सत्य आणि~$\Delta$ मधील सर्व सूत्रे
असत्य आहेत. यावरून~\Log{G3c} \emph{संपूर्ण}
आहे हेही सिद्ध होते:~$\Gamma \Sequent \Delta$
वैध असेल, तर तिची~\Log{G3c}-सिद्धता असते.
अयशस्वी सिद्धता-शोधामुळे~$\Gamma \Sequent \Delta$
अवैध ठरते; म्हणून वैध क्रमवर्तीसाठीचा प्रत्येक
सिद्धता-शोध यशस्वी ठरलाच पाहिजे.













\section{\Log{G3c} साठी सिद्धता-शोध अल्गोरिदम}\label{pt:ps:sal:sec}

\Log{G3c} साठी सिद्धता-शोध अल्गोरिदम मांडू.
हा एकमेव शक्य अल्गोरिदम नाही आणि सर्वांत कार्यक्षमही
नाही. पण तो बरोबर आहे:~$\Gamma \Sequent \Delta$
या क्रमवर्तीची सिद्धता असेल, तर अखेरीस तो
थांबतो आणि ती सिद्धता शोधतो. तो विश्वासार्हही
आहे: सिद्धता न सापडता तो थांबला, किंवा कधीच
समाप्त झाला नाही, तर~$\Gamma \Sequent \Delta$
मुळातच वैध नाही.
या बंद-पद शोधात सुरुवातीच्या क्रमवर्तीतील सर्व सूत्रे
वाक्ये, म्हणजे मुक्त चल नसलेली सूत्रे, असावीत.

अल्गोरिदमचा अत्यावश्यक गुणधर्म असा की
$\Gamma \Sequent \Delta$ मधील किंवा सिद्धता-शोधात
निर्माण झालेल्या क्रमवर्तीतील प्रत्येक सूत्र
न्यूनीकृत होते: उलट्या दिशेने लावलेल्या
अनुमाननियमात त्याला मुख्य सूत्र मानले जाते,
आणि गरजेइतक्या वेळा असे केले जाते.
या गुणधर्माला \emph{न्याय्यता} म्हणू.

अल्गोरिदम टप्प्याटप्प्याने चालतो.
प्रत्येक टप्प्याचे फलित क्रमवर्तींचा वृक्ष असतो.
पुढील टप्प्याचे फलित मिळवताना आधीच स्वयंसिद्धक
नसलेल्या प्रत्येक सर्वांत वरच्या क्रमवर्तीतील
सूत्राचे न्यूनीकरण करतो.

~$0$ व्या टप्प्यावर~$\Gamma \Sequent \Delta$
ने सुरुवात करतो. न्याय्यता सुनिश्चित करण्यासाठी
$\Gamma$ आणि~$\Delta$ मधील सूत्रांच्या
प्रत्येक आस्थितीला, म्हणजे वेगवेगळ्या सूत्रे
च्या आस्थितींना, वेगळा संख्यात्मक निर्देशांक
देतो. उदाहरणार्थ, सूत्रे डावीकडून उजवीकडे
मोजा. या संख्या ऊर्ध्वनिर्देशांक म्हणून लिहितो.
उदाहरणार्थ,~$A, B \Sequent C$ ची सिद्धता
शोधायची असेल, तर~$0$ व्या टप्प्याचे फलित
$A^1, B^2 \Sequent C^3$ असते. अशा निर्देशांकित
क्रमवर्तीतील ऊर्ध्वनिर्देशांकांपैकी सर्वांत मोठा
~$n$ हा तिचा \emph{कमाल निर्देशांक}
(उदाहरणात~$3$) आहे.

क्रमवर्तीत संख्यापक असतील, तर
उजवीकडील~$\lexists[x][A(x)]$ किंवा डावीकडील
$\lforall[x][A(x)]$ न्यूनीकृत करताना कोणती पदे
वापरायची तेही ठरवावे लागते. ~$\Gamma$ आणि~$\Delta$
मधील फलनचिन्हे वापरून~स्थिरांक
$\Obj c_0$, $\Obj c_1$, $\Obj c_2$, \dots
यांपासून बनणारी पदेच लागतात. या स्थिरांक
पासून बनणाऱ्या सर्व पदांचा
एक निश्चित क्रम दिला आहे असे मानू; म्हणजे,
उदाहरणार्थ,~``$\Gamma \Sequent \Delta$ मध्ये
न आलेला पहिला स्थिरांक'' असे म्हणता येईल.
अल्गोरिदमने बनवलेल्या वृक्षातील प्रत्येक
निर्देशांकित क्रमवर्तीशी स्थिरांकाचा एक संच
संबंधित असतो. ~$0$ व्या टप्प्यातील क्रमवर्तीशी
तिच्यात येणाऱ्या सर्व स्थिरांकाचा संच
(असतील तर) संबंधित असतो; त्यात एकही
स्थिरांक नसतील, तर~$\{\Obj c_0\}$ असतो.

~$n$ व्या टप्प्यावर संबंधित स्थिरांक
च्या संचांसह निर्देशांकित क्रमवर्तींचा वृक्ष
अल्गोरिदमने निर्माण केला आहे असे मानू.
~$n+1$ व्या टप्प्यावर प्रत्येक सर्वांत वरच्या
क्रमवर्तीसाठी पुढीलपैकी योग्य ती कृती करतो.

~$A$ हे सूत्र क्रमवर्तीच्या डाव्या
आणि उजव्या दोन्ही बाजूंना असेल, किंवा डावीकडे
$\lfalse$ असेल, तर काही करत नाही: अशा क्रमवर्तींच्या
~\Log{G3c} मधील सिद्धता असतात; या विशिष्ट
सर्वांत वरच्या क्रमवर्तीसाठी पुढे काही करायचे नसते.

क्रमवर्तीत अणुरूप नसलेले एकही सूत्र
नसेल, तर अल्गोरिदम अयशस्वी म्हणून थांबवा.
$\Gamma \Sequent \Delta$ ची सिद्धता नाही.

अन्यथा सर्वांत लहान निर्देशांक~$i$ असलेले
अणुरूप नसलेले सूत्र~$A$ निवडा. मग विचाराधीन
सर्वांत वरची क्रमवर्ती~$A^i, \Pi \Sequent \Lambda$
किंवा~$\Pi \Sequent \Lambda, A^i$ अशी असते.
त्या क्रमवर्तीचा कमाल निर्देशांक~$k$ आणि
तिच्याशी संबंधित स्थिरांकाचा संच~$C$ असू द्या.

~$A^i$ चे ``न्यूनीकरण'' करू:~$A$
चा मुख्य संकारक आणि~$A^i$ डावीकडे आहे
की उजवीकडे, यांवर ठरणाऱ्या अनुमाननियमानुसार
नव्या सर्वांत वरच्या क्रमवर्ती तयार करतो.

\begin{enumerate}
  \item $A \ident B \land C$ असून ते डावीकडे
  असेल, तर~$(B \land C)^i, \Pi \Sequent \Lambda$
  च्या वर एक नवी सर्वोच्च क्रमवर्ती जोडा:
  \[
  \Axiom$B^{k+1}, C^{k+2}, \Pi \fCenter \Lambda$
  \RightLabel{\LeftR{\land}}
  \UnaryInf$(B \land C)^i, \Pi \fCenter \Lambda$
  \DisplayProof
  \]
  नव्या क्रमवर्तीशी संबंधित स्थिरांकाचा संच~$C$ आहे.
  \item $A \ident B \land C$ असून ते उजवीकडे
  असेल, तर~$\Pi \Sequent \Lambda, (B \land C)^i$
  च्या वर दोन नव्या सर्वोच्च क्रमवर्ती जोडा:
  \[
  \Axiom$\Pi \fCenter \Lambda, B^{k+1}$
  \Axiom$\Pi \fCenter \Lambda, C^{k+1}$
  \RightLabel{\RightR{\land}}
  \BinaryInf$\Pi \fCenter \Lambda, (B \land C)^i$
  \DisplayProof
  \]
  प्रत्येक नव्या क्रमवर्तीशी संबंधित स्थिरांकाचा संच~$C$ आहे.

  \item $A \ident \lnot B$,
  $A \ident B \lor C$ आणि~$A \ident B \lif C$
  हे प्रसंगही याच प्रकारे हाताळता येतात;
  ते सरावासाठी ठेवले आहेत.

  \item $A \ident \lexists[x][B(x)]$ असून ते
  डावीकडे असेल, तर~$\lexists[x][B(x)]^i,
  \Pi \Sequent \Lambda$ च्या वर एक नवी
  सर्वोच्च क्रमवर्ती जोडा:
  \[
  \Axiom$B(c)^{k+1}, \Pi \fCenter \Lambda$
  \RightLabel{\LeftR{\lexists}}
  \UnaryInf$\lexists[x][B(x)]^i, \Pi \fCenter \Lambda$
  \DisplayProof
  \]
  येथे~$c$ हा~$C$ मध्ये नसलेला पहिला स्थिरांक
  आहे. नव्या क्रमवर्तीशी संबंधित स्थिरांक
  चा संच~$C \cup \{c\}$ आहे.
  \item $A \ident \lexists[x][B(x)]$ असून ते
  उजवीकडे असेल, तर~$\Pi \Sequent \Lambda,
  \lexists[x][B(x)]^i$ च्या वर एक नवी
  सर्वोच्च क्रमवर्ती जोडा:
  \[
  \Axiom$\Pi \fCenter \Lambda, \lexists[x][B(x)]^{k+1}, B(t)^{k+2}$
  \RightLabel{\RightR{\lexists}}
  \UnaryInf$\Pi \fCenter \Lambda, \lexists[x][B(x)]^i$
  \DisplayProof
  \]
  येथे~$t$ हे~$C$ मधीलच स्थिरांक असलेले,
  आणि या क्रमवर्तीखाली उजवीकडील
  $\lexists[x][B(x)]$ च्या आधीच्या न्यूनीकरणात
  न वापरलेले पहिले पद आहे. अशी सर्व पदे वापरून
  झाली असतील, तर~$C$ मधील पहिला स्थिरांक पुन्हा वापरा.
  $C$ सुरुवातीपासून अरिक्त आहे आणि पुढे त्यात फक्त भर
  पडते; म्हणून हा स्थिरांक उपलब्ध आहे आणि~$C$ बाहेरचा
  नवा स्थिरांक नकळत क्रमवर्तीत येत नाही.
  नव्या क्रमवर्तीशी संबंधित स्थिरांकाचा संच~$C$.
  \item $A \ident \lforall[x][B(x)]$
  हे प्रसंगही याच प्रकारे हाताळता येतात;
  ते सरावासाठी ठेवले आहेत.
\end{enumerate}

~$n+1$ व्या टप्प्यावर कोणत्याही सर्वांत वरच्या
क्रमवर्तीच्या वर नवी क्रमवर्ती जोडली गेली नसेल,
तर अल्गोरिदम यशस्वीपणे संपतो. अशा वेळी
~$n+1$ व्या टप्प्यातील वृक्षाचे सर्व निर्देशांक
काढल्यावर~$\Gamma \Sequent \Delta$ ची
सिद्धता मिळते. सर्वांत वरच्या क्रमवर्तींचे
रूप~$A, \Pi \Sequent \Lambda, A$ किंवा
$\lfalse, \Pi \Sequent \Lambda$ असे असते.
सर्व अनुमाने~\Log{G3c} मधील योग्य अनुमाने असतात.
विधानीय अनुमानांसाठी हे स्पष्ट आहे. प्रत्येक
टप्प्यावर क्रमवर्तीशी संबंधित स्थिरांकाच्या
संच~$C$ मध्ये त्या क्रमवर्तीतील सर्व स्थिरांक असतात.
म्हणून~\LeftR{\lexists} आणि~\RightR{\lforall}
यांनी न्यूनीकरण करताना आयगेन-चराची अट पूर्ण होते:
आयगेन-चर~$c$ हा निष्कर्षाशी संबंधित संच~$C$
मध्ये नसलेला स्थिरांक असल्याने तो
~$c$ निष्कर्षातही आलेला नव्हता. म्हणून:

\begin{prop}
~\Log{G3c} साठीचा सिद्धता-शोध अल्गोरिदम बरोबर आहे:
तो यशस्वीपणे संपला, तर सुरुवातीच्या
क्रमवर्ती~$\Gamma \Sequent \Delta$ ची
सिद्धता असते.
\end{prop}













\section{संपूर्णता}\label{pt:ps:cpl:sec}

आता सिद्धता-शोध अल्गोरिदम विश्वासार्ह आहे हे दाखवू:
तो सिद्धता न सापडता अयशस्वी म्हणून थांबला तरी किंवा कधीच समाप्त झाला नाही तरी,
सुरुवातीच्या क्रमवर्ती~$\Gamma \Sequent \Delta$ ची
सिद्धता नसते. यासाठी अशी संरचना~$\Struct M$
बांधू की प्रत्येक~$A \in \Gamma$ साठी~$\Sat{M}{A}$
आणि प्रत्येक~$A \in \Delta$ साठी~$\Sat/{M}{A}$.
म्हणजे~$\Gamma$ मधील सर्व सूत्रे सत्य,
तर~$\Delta$ मधील सर्व सूत्रे हे~$\Struct M$
मध्ये असत्य आहेत. त्यामुळे~$\Gamma \Sequent \Delta$
वैध नाही. \Log{G3c} निर्दोष असल्याने
$\Gamma \Sequent \Delta$ ची सिद्धता नाही.
येथे सुरुवातीच्या क्रमवर्तीतील सर्व सूत्रे वाक्ये आहेत,
म्हणजे त्यांत मुक्त चल नाहीत. पुढील बंद-पद अल्गोरिदम
आणि पद-प्रतिमानाची सिद्धता या व्याप्तीत वापरायची.

सूत्रांची जोडी~$\Theta, \Xi$ आणि
स्थिरांकाचा संच~$C$ यांपासून~$\Struct{M}$ बांधतो.
हे~$\Theta, \Xi$ आणि~$C$ आपल्याला अयशस्वी
(अयशस्वी म्हणून थांबलेल्या किंवा समाप्त न होणाऱ्या) सिद्धता-शोधाच्या
\emph{अपयश-शाखे}तून मिळतात. अपयश-शाखा म्हणजे
क्रमवर्तींचा~$\Pi_n \Sequent \Lambda_n$ अनुक्रम
आणि~स्थिरांकाचे संच~$C_n$; येथे
$\Pi_0 \Sequent \Lambda_0$ ही सुरुवातीची
अंतिम क्रमवर्ती~$\Gamma \Sequent \Delta$ असते,
आणि~$\Pi_{n+1} \Sequent \Lambda_{n+1}$
हे~$n+1$ व्या टप्प्यावर~$\Pi_n \Sequent \Lambda_n$
पासून निर्माण झालेले असे आधारविधान असते की
अल्गोरिदमला त्याची सिद्धता सापडत नाही.
प्रत्येक~$n$ साठी असे आधारविधान असलेच पाहिजे;
नसते तर अल्गोरिदमला सिद्धता सापडली असती.
~$C_n$ हा~$\Pi_n \Sequent \Lambda_n$
शी संबंधित स्थिरांकाचा संच घेऊ.

अल्गोरिदम फक्त अणुरूप सूत्रे असलेल्या
सर्वांत वरच्या क्रमवर्तीवर अयशस्वी म्हणून थांबला, तर त्या
क्रमवर्तीवर संपणारी शाखा ही अपयश-शाखा आहे.
अल्गोरिदम कधीच समाप्त झाला नाही, तर तो वाढत्या
मोठ्या वृक्षांची निर्मिती करत राहतो. हे वृक्ष
वाढून एक अनंत वृक्ष तयार होतो असे समजू शकतो
(अल्गोरिदम अनंत टप्पे चालल्यानंतर जो वृक्ष असेल तो).
हा अनंत वृक्ष सांत-शाखित असतो: प्रत्येक शिखराला
फक्त एक किंवा दोन बालशिखरे असतात---म्हणजे त्याच्या
तत्काळ वरच्या क्रमवर्ती. क्योनिग (K\H{o}nig)
यांच्या उपप्रमेयानुसार,
प्रत्येक अनंत सांत-शाखित वृक्षाला अनंत शाखा असते.
सिद्धता-शोध कायम चालू राहिल्यास अशी अनंत शाखा
अपयश-शाखा ठरते.

~$\Pi_n \Sequent \Lambda_n$ आणि~$C_n$
यांचा अनुक्रम अपयश-शाखा असेल, तर
$\Theta = \bigcup_n \Pi_n$ आणि~$\Xi = \bigcup_n \Lambda_n$
घेऊ (सूत्रे वरील निर्देशांक आणि त्यांच्या
अनेक आस्थिती काढून); तसेच~$C = \bigcup_n C_n$ घेऊ.

आता~$\Struct{M}$ परिभाषित करू.
\begin{enumerate}
\item $\Domain{M}$ हा~$C$ मधील स्थिरांक आणि
$\Gamma \Sequent \Delta$ मधील फलने
(असतील तर) वापरून, पण चले न वापरता,
बनवता येणाऱ्या पदांचा संच आहे.
\item $c \in C$ असेल, तर~$\Assign{c}{M} = c$.
\item $f$ हे~$\Gamma \Sequent \Delta$ मधील
$m$-स्थानी फलन असेल आणि~$t_1$,
\dots, $t_m \in \Domain{M}$ असतील, तर
$\Assign{f}{M}(t_1, \dots, t_m) = f(t_1, \dots, t_m)$.
\item $R$ हे~$\Gamma \Sequent \Delta$ मधील
$m$-स्थानी विधेय असेल, तर
$\Assign{R}{M} = \Setabs{\tuple{t_1, \dots, t_m} \in
\Domain{M}^m}{R(t_1, \dots, t_m) \in \Theta}$.
\end{enumerate}
~$\Struct{M}$ हे \emph{पद-प्रतिमान} आहे, कारण
त्याचे विश्व पदांचा संच आहे आणि स्थिरांक व
फलने यांचा अर्थ असा लावला आहे की
$\Value{t}{M} = t$ मिळते. ~$\Theta$ मधील अणुरूप
सूत्रे वापरून~$\Assign{R}{M}$ अशी परिभाषित
केली आहे की~$\Sat{M}{R(t_1, \dots, t_m)}$ तेव्हाच
आणि तेव्हाच जेव्हा~$R(t_1, \dots, t_m) \in \Theta$.

\begin{lem}\label{pt:ps:cpl:lem:truth}
प्रत्येक~$A \in \Theta \cup \Xi$ साठी:
\begin{enumerate}
  \item $A \in \Theta$ असेल, तर~$\Sat{M}{A}$.
  \item $A \in \Xi$ असेल, तर~$\Sat/{M}{A}$.
\end{enumerate}
\end{lem}

\begin{proof}
~$\depth{A}$ वर विगमन करू.

प्रथम~$A$ अणुरूप आहे असे मानू;
म्हणजे~$A \ident \lfalse$ किंवा
$A \ident R(t_1, \dots, t_m)$.

~$\Pi_n \Sequent \Lambda_n$ ही अपयश-शाखा
असल्याने कोणत्याही~$\Pi_n$ मध्ये~$\lfalse$
असू शकत नाही (नाहीतर~$\Pi_n \Sequent \Lambda_n$
स्वयंसिद्धक ठरेल). ~$\Struct{M}$ च्या व्याख्येनुसार
$\Sat{M}{R(t_1, \dots, t_m)}$ तेव्हा आणि केवळ तेव्हा
जेव्हा~$R(t_1, \dots, t_m) \in \Theta$.
दोन्ही मिळून: $A \in \Theta$ असेल, तर
$\Sat{M}{A}$.

~$A$ अणुरूप असून~$A \in \Pi_n$ असेल,
तर~$A \in \Pi_{n+1}$; आणि~$A \in \Lambda_n$
असेल, तर~$A \in \Lambda_{n+1}$.
(सिद्धता-शोध अल्गोरिदम पुढील क्रमवर्ती कशा
बनवतो यावरून हे मिळते.) म्हणून~$A \in \Theta \cap \Xi$
असेल, तर एखाद्या~$n$ साठी~$A \in \Pi_n \cap
\Lambda_n$. मग~$\Pi_n \Sequent \Lambda_n$
स्वयंसिद्धक ठरते; पण अपयश-शाखेत स्वयंसिद्धके
असू शकत नाहीत. म्हणून अणुरूप~$A$ साठी
$A \notin \Theta \cap \Xi$; परिणामी~$A \in \Xi$
असेल तर~$A \notin \Theta$. त्यामुळे~$\Struct{M}$
च्या व्याख्येनुसार~$A \in \Xi$ असताना
$\Sat/{M}{A}$.

विगमन-टप्प्यासाठी,~$A$ पेक्षा कमी
खोली असलेल्या सर्व सूत्रे साठी
(1) आणि~(2) खरे आहेत असे मानू. ~$A$ साठी
(1) व~(2) सिद्ध करण्यासाठी प्रसंग वेगळे पाहू.

येथील $A$ अणुरूप नसलेले आहे. या परिच्छेदात
“सर्वांत लहान निर्देशांक” म्हणजे अणुरूप नसलेल्या
आस्थितींमधील सर्वांत लहान निर्देशांक; अणुरूप आस्थिती
जपल्या जातात, पण न्यूनीकरणासाठी निवडल्या जात नाहीत.
प्रथम,~$\Pi_n \Sequent \Lambda_n$ मध्ये
$A^i$ आल्यास ते~$\Pi_{n+1} \Sequent \Lambda_{n+1}$
मध्येही~$A^i$ म्हणून येते, जोपर्यंत~$i$ हा~$\Pi_n \Sequent \Lambda_n$
मधील सर्वांत लहान निर्देशांक नसतो. प्रत्येक
टप्प्यावर जोडलेल्या सर्वांत वरच्या क्रमवर्तींत
सर्वांत लहान निर्देशांक वाढतो. म्हणून अखेरीस
$A^i$ असलेली अशी~$\Pi_n \Sequent \Lambda_n$
मिळते की~$i$ हा सर्वांत लहान निर्देशांक असतो.
~$n+1$ व्या टप्प्यावर अल्गोरिदम~$A^i$ चे
न्यूनीकरण करतो. दुसऱ्या शब्दांत,
$A \in \Theta$ असेल, तर एखाद्या~$n$ साठी
$\Pi_n = \Pi_n', A^i$ आणि~$i$ सर्वांत
लहान निर्देशांक असतो. तसेच~$A \in \Xi$
असेल, तर एखाद्या~$n$ साठी
$\Lambda_n = \Lambda_n', A^i$ आणि~$i$
सर्वांत लहान निर्देशांक असतो.

\begin{enumerate}
  \item $A \in \Theta$ आणि~$A \ident B \land C$
  असेल, तर एखाद्या~$n$ साठी
  $\Pi_n = \Pi_n', (B \land C)^i$.
  ~$\Pi_{n+1} \Sequent \Lambda_{n+1}$ हे
  \LeftR{\land} चे अनुरूप आधारविधान आहे;
  म्हणजे~$\Pi_{n+1} = \Pi_n', B^{k+1}, C^{k+2}$.
  त्यामुळे~$B, C \in \Theta$. विगमन गृहीतकानुसार
  $\Sat{M}{B}$ आणि~$\Sat{M}{C}$.
  ~$\Sat{M}{}$ च्या व्याख्येवरून
  $\Sat{M}{B \land C}$.

  \item $A \in \Xi$ आणि~$A \ident B \land C$
  असेल, तर एखाद्या~$n$ साठी
  $\Lambda_n = \Lambda_n', B \land C$.
  ~$\Pi_{n+1} \Sequent \Lambda_{n+1}$ हे
  निष्कर्ष~$\Pi_n \Sequent \Lambda'_n, B \land C$
  असलेल्या~\RightR{\land} च्या दोन आधारविधानांपैकी
  एक आहे; म्हणजे~$\Lambda_{n+1} = \Lambda_n', B^{k+1}$
  किंवा~$\Lambda_{n+1} = \Lambda_n', C^{k+1}$.
  त्यामुळे~$B \in \Xi$ किंवा~$C \in \Xi$.
  विगमन गृहीतकानुसार~$\Sat/{M}{B}$ किंवा
  $\Sat/{M}{C}$. ~$\Sat{M}{}$ च्या व्याख्येवरून
  $\Sat/{M}{B \land C}$.

  \item $A \ident \lnot B$,
  $A \ident B \lor C$ आणि~$A \ident B \lif C$
  हे प्रसंगही याच प्रकारे हाताळता येतात;
  ते सरावासाठी ठेवले आहेत.

  \item $A \in \Theta$ आणि
  $A \ident \lexists[x][B(x)]$ असेल,
  तर एखाद्या~$n$ साठी
  $\Pi_n = \Pi_n', \lexists[x][B(x)]^i$.
  ~$\Pi_{n+1} \Sequent \Lambda_{n+1}$ हे
  \LeftR{\lexists} चे अनुरूप आधारविधान आहे;
  म्हणजे एखाद्या~$c$ साठी
  $\Pi_{n+1} = \Pi_n', B(c)^{k+1}$.
  त्यामुळे~$B(c) \in \Theta$ आणि~$c \in C$.
  विगमन गृहीतकानुसार~$\Sat{M}{B(c)}$.
  ~$\Sat{M}{}$ च्या व्याख्येवरून
  $\Sat{M}{\lexists[x][B(x)]}$.

  \item $A \in \Xi$ आणि
  $A \ident \lexists[x][B(x)]$ असेल,
  तर एखाद्या~$n$ साठी
  $\Lambda_n = \Lambda_n', \lexists[x][B(x)]^i$.
  प्रत्येक वेळी~$\lexists[x][B(x)]$ चे
  न्यूनीकरण करताना~$\lexists[x][B(x)]$
  आधारविधानाच्या उजव्या
  बाजूस नव्या निर्देशांकासह राहते. म्हणून
  अपयश-शाखेत ते उजव्या बाजूस असेल, तर
  $\lexists[x][B(x)]$ चे अनंत वेळा
  न्यूनीकरण होते. प्रत्येक~$t \in \Domain{M}$
  हे अखेरीस अपयश-शाखेत~``$C_n$ मधीलच स्थिरांक
  असलेले आणि उजवीकडील~$\lexists[x][B(x)]$
  च्या आधीच्या कोणत्याही न्यूनीकरणात न वापरलेले
  पहिले पद'' ठरते. त्यामुळे प्रत्येक
  $t \in \Domain{M}$ साठी एखाद्या~$n$ ला
  $B(t) \in \Lambda_n$, आणि म्हणून~$B(t) \in \Xi$.
  विगमन गृहीतकानुसार प्रत्येक~$t \in \Domain{M}$
  साठी~$\Sat/{M}{B(t)}$. ~$\Sat{M}{}$
  च्या व्याख्येवरून~$\Sat/{M}{\lexists[x][B(x)]}$.

  \item $A \ident \lforall[x][B(x)]$
  हे प्रसंगही याच प्रकारे हाताळता येतात;
  ते सरावासाठी ठेवले आहेत.
\end{enumerate}
\end{proof}

\begin{cor}\label{pt:ps:cpl:cor:complete}
~\Log{G3c} साठीचा सिद्धता-शोध अल्गोरिदम संपूर्ण
आहे: तो अयशस्वी म्हणून थांबला किंवा कधीच समाप्त झाला नाही,
तर सुरुवातीची क्रमवर्ती~$\Gamma \Sequent \Delta$
वैध नसते आणि म्हणून तिची सिद्धता नसते.
\end{cor}

\begin{proof}
  अल्गोरिदम अयशस्वी म्हणून थांबला किंवा कधीच समाप्त झाला नाही,
  तर अशी अपयश-शाखा असते की तिच्यापासून
  $\Struct{M}$ परिभाषित करता येते.
  \ref{pt:ps:cpl:lem:truth} नुसार प्रत्येक~$A \in \Gamma$
  साठी~$\Sat{M}{A}$ आणि प्रत्येक~$A \in \Delta$
  साठी~$\Sat/{M}{A}$. (लक्षात घ्या की
  $\Pi_0 = \Gamma$ आणि~$\Lambda_0 = \Delta$;
  म्हणून~$\Gamma \subseteq \Theta$ आणि
  $\Delta \subseteq \Xi$.) त्यामुळे
  $\Gamma \Sequent \Delta$ वैध नाही.
  ~\Log{G3c} च्या निर्दोषतेवरून
  $\Gamma \Sequent \Delta$ ची सिद्धता नाही.
\end{proof}

\begin{cor}
वाक्यांच्या क्रमवर्तींसाठी~\Log{G3c} संपूर्ण आहे:
$\Gamma \Sequent \Delta$ वैध असेल,
तर~$\Log{G3c} \Proves \Gamma \Sequent \Delta$.
\end{cor}

\begin{proof}
  व्यत्यस्त विधान सिद्ध करू.
  $\Log{G3c} \Proves/ \Gamma \Sequent \Delta$
  असे मानू. मग शोधण्याजोगी सिद्धता नसल्याने
  सिद्धता-शोध अल्गोरिदम अयशस्वी म्हणून थांबला पाहिजे
  किंवा कायम चालू राहिला पाहिजे.
  \ref{pt:ps:cpl:cor:complete} नुसार
  $\Gamma \Sequent \Delta$ वैध नाही.
\end{proof}













\section{टॅब्लो}\label{pt:ps:tab:sec}

टॅब्लो पद्धती या क्रमवर्ती-कलनाशी जवळून
संबंधित सिद्धता पद्धती आहेत. हाताने किंवा
संगणकीय अंमलबजावणीत सिद्धता-शोधासाठी त्या विशेष
सोयीच्या आहेत. टॅब्लोत क्रमवर्तींचे वृक्ष न बांधता
सिद्धता ही सूत्रांचा वृक्ष असते.
पण नैसर्गिक निगमनापेक्षा वेगळे म्हणजे त्यांचे नियम
क्रमवर्ती-कलनाच्या नियमांशी जुळतात. टॅब्लो
सिद्धता म्हणजे मूलतः प्रतिमानाचा स्पष्ट शोध.
म्हणून त्यांना कधी \emph{चिन्हार्थविषयक} टॅब्लो,
तर कधी \emph{सत्यता-वृक्ष} म्हणतात.

\emph{चिन्हांकित} टॅब्लो हा चिन्हांकित
सूत्रांचा वृक्ष आहे. चिन्हांकित सूत्राचे रूप
$\sFmla{\True}{A}$ किंवा~$\sFmla{\False}{A}$
असे असू शकते. वृक्ष खालच्या दिशेने वाढतात.
सुरुवातीला आपल्याला काय सिद्ध करायचे ते
दर्शवणारी सूत्रे असतात. उदाहरणार्थ,
$A, B \Sequent C, D$ सिद्ध करण्यासाठी
वृक्षाची सुरुवात
$\sFmla{\True}{A}, \sFmla{\True}{B},
\sFmla{\False}{C}, \sFmla{\False}{D}$ अशी
होईल. ~$A$ आणि~$B$ सत्य, पण~$C$ आणि~$D$
असत्य ठरवणारी संरचना म्हणजे
$A, B \Sequent C, D$ अवैध असल्याचे
दाखवणारी संरचना.

टॅब्लोचे तळाचे पानशिखर शाखा-विस्तार
नियमांनी वाढवता येते. हे नियम त्याच शाखेत नवी
शिखरे जोडतात किंवा खाली दोन भावंड शिखरे जोडून
शाखा दोन नव्या शाखांत विभागतात. त्याच शाखेतील
त्या शिखराच्या वरच्या कोणत्याही चिन्हांकित
सूत्राला नियम लागू करता येतात. अभिजात
टॅब्लो पद्धत~\Log{Tc} साठीचे शाखा-विस्तार नियम
\ref{pt:ps:tab:tab:Tc} मध्ये दिले आहेत.








\begin{table}
\begin{center}
\begin{tabular}{@{}cc@{}}
\AxiomC{\sFmla{\True}{\lnot A}}
\RightLabel{\TRule{\True}{\lnot}}
\UnaryInfC{\sFmla{\False}{A}}
\DisplayProof
&
\AxiomC{\sFmla{\False}{\lnot A}}
\RightLabel{\TRule{\False}{\lnot}}
\UnaryInfC{\sFmla{\True}{A}}
\DisplayProof
\\[6ex]
\AxiomC{\sFmla{\True}{A \land B}}
\RightLabel{\TRule{\True}{\land}}
\UnaryInfC{\sFmla{\True}{A}}
\noLine
\UnaryInfC{\sFmla{\True}{B}}
\DisplayProof
&
\AxiomC{\sFmla{\False}{A \land B}}
\RightLabel{\TRule{\False}{\land}}
\UnaryInfC{$\sFmla{\False}{A} \quad \mid \quad \sFmla{\False}{B}$}
\DisplayProof
\\[6ex]
\AxiomC{\sFmla{\True}{A \lor B}}
\RightLabel{\TRule{\True}{\lor}}
\UnaryInfC{$\sFmla{\True}{A} \quad \mid \quad \sFmla{\True}{B}$}
\DisplayProof
&
\AxiomC{\sFmla{\False}{A \lor B}}
\RightLabel{\TRule{\False}{\lor}}
\UnaryInfC{\sFmla{\False}{A}}
\noLine
\UnaryInfC{\sFmla{\False}{B}}
\DisplayProof
\\[6ex]
\AxiomC{\sFmla{\True}{A \lif B}}
\RightLabel{\TRule{\True}{\lif}}
\UnaryInfC{$\sFmla{\False}{A} \quad \mid \quad \sFmla{\True}{B}$}
\DisplayProof
&
\AxiomC{\sFmla{\False}{A \lif B}}
\RightLabel{\TRule{\False}{\lif}}
\UnaryInfC{\sFmla{\True}{A}}
\noLine
\UnaryInfC{\sFmla{\False}{B}}
\DisplayProof
\\[6ex]
\AxiomC{\sFmla{\True}{\lforall[x][A(x)]}}
\RightLabel{\TRule{\True}{\forall}}
\UnaryInfC{\sFmla{\True}{A(t)}}
\DisplayProof
&
\AxiomC{\sFmla{\False}{\lforall[x][A(x)]}}
\RightLabel{\TRule{\False}{\lforall}}
\UnaryInfC{\sFmla{\False}{A(a)}}
\DisplayProof
\\[6ex]
\AxiomC{\sFmla{\True}{\lexists[x][A(x)]}}
\RightLabel{\TRule{\True}{\lexists}}
\UnaryInfC{\sFmla{\True}{A(a)}}
\DisplayProof
&
\AxiomC{\sFmla{\False}{\lexists[x][A(x)]}}
\RightLabel{\TRule{\False}{\lexists}}
\UnaryInfC{\sFmla{\False}{A(t)}}
\DisplayProof
\end{tabular}
\end{center}
\caption{\Log{Tc} चे नियम. $\TRule{\False}{\lforall}$ आणि
$\TRule{\True}{\lexists}$ मध्ये स्थिरांक~$a$ वरच्या
शाखेत आलेला नसावा; म्हणजे तो शाखेसाठी नवा असावा.}
\label{pt:ps:tab:tab:Tc}
\end{table}




नियम म्हणजे~\Log{G3c} चेच नियम उलटे
दाखवले आहेत. चिन्हे~$\True$ आणि~$\False$
अनुक्रमे क्रमवर्तीच्या डाव्या किंवा उजव्या बाजूला
सूत्र मुख्य आहे हे दर्शवतात.

टॅब्लोची शाखा~$\sFmla{\True}{A}$ आणि
$\sFmla{\False}{A}$ ही दोन्ही सूत्रे धरत असेल,
किंवा तिच्यावर~$\sFmla{\True}{\lfalse}$ असेल,
तर ती \emph{बंद} आहे. प्रत्येक शाखा बंद असेल,
तर संपूर्ण टॅब्लो बंद म्हणतो. उदाहरणार्थ,
$A \lor (B \land C) \Sequent (A \lor B) \land
(A \lor C)$ साठी हा बंद टॅब्लो आहे; म्हणजे
$\sFmla{\True}{A \lor (B \land C)}$ आणि
$\sFmla{\False}{(A \lor B) \land (A \lor C)}$
यांपासून सुरुवात होते:
\begin{oltableau}
  [\sFmla{\True}{\formula{A} \lor (\formula{B} \land \formula{C})},
    just=\TAss, checked
    [\sFmla{\False}{(\formula{A} \lor \formula{B}) \land
        (\formula{A} \lor \formula{C})}, just=\TAss, checked
      [\sFmla{\True}{\formula{A}}, just={\TRule{\True}{\lor}[1]}
        [\sFmla{\False}{\formula{A} \lor \formula{B}},
          just={\TRule{\False}{\land}[2]}, checked
          [\sFmla{\False}{\formula{A}},
            just={\TRule{\False}{\lor}[4]}
            [\sFmla{\False}{\formula{B}},
              just={\TRule{\False}{\lor}[4]}, close]
          ]
        ]
        [\sFmla{\False}{\formula{A} \lor \formula{C}},
          just={\TRule{\False}{\land}[2]}, checked
          [\sFmla{\False}{\formula{A}},
            just={\TRule{\False}{\lor}[4]}
            [\sFmla{\False}{\formula{C}},
              just={\TRule{\False}{\lor}[4]}, close]
          ]
        ]
      ]
      [\sFmla{\True}{\formula{B} \land \formula{C}},
        just={\TRule{\True}{\lor}[1]}, checked
        [\sFmla{\False}{\formula{A} \lor \formula{B}},
          just={\TRule{\False}{\land}[2]}, checked
          [\sFmla{\True}{\formula{B}},
            just={\TRule{\True}{\land}[3]}, move by=2
            [\sFmla{\True}{\formula{C}},
              just={\TRule{\True}{\land}[3]}
              [\sFmla{\False}{\formula{A}},
                just={\TRule{\False}{\lor}[4]}
                [\sFmla{\False}{\formula{B}},
                  just={\TRule{\False}{\lor}[4]},close
                ]
              ]
            ]
          ]
        ]
        [\sFmla{\False}{\formula{A} \lor \formula{C}},
          just={\TRule{\False}{\land}[2]}, checked
          [\sFmla{\True}{\formula{B}},
            just={\TRule{\True}{\land}[3]}, move by=2
              [\sFmla{\True}{\formula{C}},
                just={\TRule{\True}{\land}[3]}
                [\sFmla{\False}{\formula{A}},
                  just={\TRule{\False}{\lor}[4]}
                  [\sFmla{\False}{\formula{C}},
                  just={\TRule{\False}{\lor}[4]},close
                ]
              ]
            ]
          ]
        ]
      ]
    ]
  ]
\end{oltableau}

खूणा दर्शवतात की संबंधित नियम त्या
सूत्राच्या खालच्या सर्व शाखांवर लागू
केला आहे. ~$\otimes$ हे चिन्ह शाखा बंद असल्याचे
दर्शवते. नियमानंतरच्या ओळ-क्रमांकांवरून नियम
कोणत्या चिन्हांकित सूत्रे ला लागू केला
हे कळते (म्हणजे दर्शवलेल्या ओळीवरील त्याच
शाखेतील सूत्राला).

टॅब्लोमध्ये सिद्धता-शोध टप्प्याटप्प्याने करतो.
येथील बंद-पद शोधासाठी सुरुवातीची सूत्रे वाक्ये असावीत.
$0$ व्या टप्प्यावर ती वर लिहा. प्रत्येक टप्प्यावर
बंद शाखा पुढे वाढवू नका. प्रत्येक उघड्या शाखेत अद्याप
खूण न केलेले सर्वांत वरचे अणुरूप नसलेले सूत्र निवडा,
पण पुढील परिच्छेदातील पुन्हा वापरायचे संख्यापक त्यातून
वगळा. त्याला शाखा-विस्तार नियम लावा. संबंधित सर्व
उघड्या शाखांवर नियम लागू झाल्यावर त्याला खूण करा.
अणुरूप सूत्रांना विस्तार-नियम लागत नाही.
(वरचे विधानीय उदाहरण या पद्धतीने बनवले आहे.)

$\sFmla{\True}{\lforall}$ आणि
$\sFmla{\False}{\lexists}$ ही येथे संख्यापक-प्रकारांची
संक्षिप्त रूपे आहेत; संख्यापक एकटाच पूर्ण सूत्र असल्याचा
अर्थ नाही. पूर्ण रूपे~$\sFmla{\True}{\lforall[x][B(x)]}$
आणि~$\sFmla{\False}{\lexists[x][B(x)]}$ अशी आहेत.
या सूत्रांवर कायमची खूण करत नाही: सर्व योग्य पदांसाठी
अनुरूप नियम अखेरीस लावले गेले पाहिजेत.
प्रत्येक उघड्या शाखेबरोबर स्थिरांकांचा संच~$C$ ठेवा.
सुरुवातीला त्यात सुरुवातीच्या सूत्रांतील सर्व स्थिरांक
घ्या; ते नसतील, तर~$\{\Obj c_0\}$ घ्या.
$\TRule{\False}{\lforall}$ किंवा~$\TRule{\True}{\lexists}$
साठी आणलेला नवा स्थिरांक~$C$ मध्ये जोडा, तो सूत्रात
प्रत्यक्ष उरला नाही तरी; तो आधीच्या~$C$ बाहेरचा असावा.
शाखा फुटल्यावर दोन्ही शाखांना त्या संचाची प्रत द्या.
सामान्य सूत्र हाताळल्यानंतर, त्या टप्प्यावर शाखेतील
प्रत्येक पुन्हा वापरायचा संख्यापक एकदा हाताळा.
निश्चित पद-क्रमातून~$C$ मधील स्थिरांक आणि सुरुवातीची
फलनचिन्हे वापरून बनणारे पहिले पद~$t$ निवडा ज्यासाठी
अनुरूप~$\sFmla{\True}{B(t)}$ किंवा
$\sFmla{\False}{B(t)}$ शाखेत नाही, आणि ते उदाहरण जोडा.
असे पद नसेल, तर त्या टप्प्यावर या सूत्रासाठी काही
जोडू नका; पुढे~$C$ वाढल्यास पुन्हा तपासा.

या शोधातून बंद टॅब्लो मिळाला नाही, तर सांत उघडी
शाखा पूर्णपणे संतृप्त झाली असेल किंवा शोध अनंत,
सांत-शाखित वृक्ष तयार करत राहील. येथे सांत शाखा
संतृप्त असण्याचा अर्थ असा: पुन्हा वापरायचे संख्यापक
सोडून सर्व अणुरूप नसलेली सूत्रे हाताळली आहेत,
आणि प्रत्येक पुन्हा वापरायच्या संख्यापकासाठी~$C$ व
सुरुवातीच्या फलनचिन्हांपासून बनणाऱ्या प्रत्येक पदाचे
अनुरूप चिन्हांकित उदाहरण आधीच शाखेत आहे.
अनंत वृक्षाला अनंत शाखा असते. निश्चित न्याय्य पद-क्रमामुळे
तिच्यावरही प्रत्येक आवश्यक संख्यापक-उदाहरण अखेरीस येते.
\ref{pt:ps:cpl:sec} प्रमाणे~$\Struct{M}$ बांधता येते:
$\Theta = \Setabs{A}{\sFmla{\True}{A} \text{ शाखेवर}}$
आणि~$\Xi = \Setabs{A}{\sFmla{\False}{A} \text{ शाखेवर}}$
घ्या; स्थिरांकांसाठी शाखेबरोबर नोंदलेल्या संचांचे संघटन
घ्या. तो संच अरिक्त आहे. मग शाखेवर
$\sFmla{\True}{A}$ असेल तर~$\Sat{M}{A}$,
आणि~$\sFmla{\False}{A}$ असेल तर~$\Sat/{M}{A}$.
संतृप्त शाखेसाठी संख्यापक-उदाहरणे सर्व पदांसाठी उपलब्ध
असल्याने पद-प्रतिमानातील पूर्वीचे विगमन लागू होते.

या पद्धतीने तयार झालेले टॅब्लो प्रत्येक
शिखराला क्रमवर्ती देऊन~\Log{G3c} मधील सिद्धतांमध्ये
सहज रूपांतरित करता येतात. सुरुवातीची
सूत्रे ही~$\Gamma \Sequent \Delta$
क्रमवर्तीशी संबंधित असतील, तर त्या प्रत्येक
सूत्राला~$\Gamma \Sequent \Delta$ ही खूण द्या.
~$\sFmla{\True}{A}$ या चिन्हांकित सूत्र
वर शाखा-विस्तार नियम लावून शाखा वाढवली, तर
पानाला~$A, \Pi \Sequent \Lambda$ ही खूण द्या;
~$\sFmla{\False}{A}$ या चिन्हांकित सूत्र
वर नियम लावला,
तर पानाला~$\Pi \Sequent \Lambda, A$ ही खूण द्या.
टॅब्लोत~\TRule{\True}{\star} नियम वापरला,
तर क्रमवर्तीवर~\LeftR{\star} नियम लावा
(मुख्य सूत्र~$A$);~\TRule{\False}{\star}
वापरला असेल, तर~\RightR{\star} लावा.
नियमाची आधारविधाने ही टॅब्लोत नव्याने
जोडलेल्या अनुरूप शिखरांना दिलेल्या क्रमवर्ती आहेत.
~$\sFmla{\True}{A}$ आणि~$\sFmla{\False}{A}$
यांनी टॅब्लो बंद झाल्यावर अंतिम पानशिखराची
खूण~$A, \Pi \Sequent \Lambda, A$
या रूपातील क्रमवर्ती असेल. $\sFmla{\True}{\lfalse}$
ने शाखा बंद झाली असेल, तर पानाची खूण
$\lfalse,\Pi\Sequent\Lambda$ या स्वयंसिद्धकाचे रूप असेल.
टॅब्लोतील काही सलग
शिखरांना तीच क्रमवर्ती दिली जाईल; तिच्या
पुनरावृत्ती काढा. उदाहरणार्थ, वरच्या टॅब्लोचे
पुढील सिद्धता मध्ये रूपांतर होते:
\[\small
\Axiom$A \fCenter A, B$
\RightLabel{\RightR{\lor}}
\UnaryInf$A \fCenter A \lor B$
\Axiom$A \fCenter A, C$
\RightLabel{\RightR{\lor}}
\UnaryInf$A \fCenter A \lor C$
\RightLabel{\RightR{\land}}
\BinaryInf$A \fCenter (A \lor B) \land (A \lor C)$
\Axiom$B, C \fCenter A, B$
\RightLabel{\RightR{\lor}}
\UnaryInf$B, C \fCenter A \lor B$
\RightLabel{\LeftR{\land}}
\UnaryInf$B \land C \fCenter A \lor B$
\Axiom$B, C \fCenter A, C$
\RightLabel{\RightR{\lor}}
\UnaryInf$B, C \fCenter A \lor C$
\RightLabel{\LeftR{\land}}
\UnaryInf$B \land C \fCenter A \lor C$
\RightLabel{\RightR{\land}}
\BinaryInf$B \land C \fCenter (A \lor B) \land
(A \lor C)$
\RightLabel{\LeftR{\lor}}
\BinaryInf$A \lor (B \land C) \fCenter (A \lor B) \land
(A \lor C)$
\DisplayProof
\]



\chapter{विधाने म्हणजे प्रकार}\label{int:pty:chap}
\begin{editorial}
  करी--हॉवर्ड अनुरूपतेवरील या प्रकरणाचा हा \emph{अतिशय प्रायोगिक}
  मसुदा आहे. याला अधिक स्पष्टीकरण आणि प्रेरणा हवी; यात चुका
  व गाळलेले भागही असण्याची शक्यता आहे. सामान्यीकरणाची सिद्धता
  तपासून विस्तारित केली पाहिजे. गुणाकार-प्रकाराची उदाहरणे नाहीत.
  क्रमपरिवर्तन आणि सरलीकरण रूपांतरणांचा समावेश केलेला नाही.
  येथे मूलतः गृहीत धरलेल्या (प्रकारयुक्त) लॅम्डा कलनाविषयीही
  साहित्य आल्यावर हे प्रकरण अधिक समजेल. अत्यंत सावधगिरीने वापरा.
\end{editorial}









\section{परिचय}\label{int:pty:int:sec}

लॅम्डा कलन हे संगणनाचे एक साधे प्रतिमान आहे. ते चरांपासून
बनलेल्या पदांवर कार्य करते. $N$ आणि $M$ ही पदे असतील, तर
$(NM)$ हेही एक पद आहे (``$N$ चे~$M$ वर उपयोजन''). $N$ हे पद
असेल, तर $\lambd[x][N]$ हेही पद आहे. $N$ मध्ये~$x$ हे चर
असेल, तर $\lambd[x][N]$ मधील~$x$ ची संबंधित स्थाने
$\lambd[x]$ या संकारकाने बद्ध होतात; त्यामुळे ती मुक्त राहत
नाहीत. $(\lambd[x][N]M)$ या रूपातील पदाचे
$\Subst{N}{M}{x}$ मध्ये $\beta$-संकुचन होते असे म्हणतात
(म्हणजे~$N$ मधील~$x$ ची प्रत्येक मुक्त आस्थिती~$M$ ने
बदलल्यावर मिळणारे पद). $N$ मधील एखाद्या उपपदाचे
$\beta$-संकुचन करून~$N'$ मिळत असेल, तर~$N$ चे~$N'$ मध्ये
$\beta$-परिवर्तन होते, आणि आपण $N \redone N'$ असे लिहितो.
$\lambd$-कलनातील संगणन म्हणजे
$N \redone N_1 \redone N_2 \redone \dots$ असा क्रम.
हा क्रम~$N_k$ या अशा पदावर संपला की त्याच्या कोणत्याही
उपपदाचे $\beta$-संकुचन होऊ शकत नाही, तर त्या पदाला
\emph{सामान्य} किंवा~$N$ चे सामान्य रूप म्हणतात.
$\lambd$-कलनातील संगणनाचा विचार अशा क्रमांच्या रूपात करा.
सामान्य रूप मिळाल्यास संगणन थांबते, आणि ते सामान्य रूपच
संगणनाचे फल असते. हे साधे संगणन-प्रतिमान ट्यूरिंग-संपूर्ण
ठरते: प्रत्येक आंशिक पुनरावर्ती फलनाची गणना
$\lambd$-कलनातील एका पदाने करता येते.

हास्केल करी आणि विल्यम हॉवर्ड यांनी स्वतंत्रपणे नैसर्गिक
निगमन आणि $\lambd$-कलन यांतील जवळचे साम्य शोधले.
सहज अर्थाने $\lambd[x][N]$ हे पद एक फलन आहे: $x$ हा त्याचा
आर्ग्युमेंट, आणि~$x$ दिल्यावर मूल्य कसे मोजायचे हे~$N$
सांगते. म्हणून $(\lambd[x][N]M)$ म्हणजे
$\lambd[x][N]$ हे फलन~$M$ या आर्ग्युमेंटवर लावणे;
$\beta$-संकुचनाने त्याचे मूल्यांकन कसे करायचे ते कळते:
$N$ मध्ये~$x$ ऐवजी~$M$ ठेवायचे (आणि मिळालेल्या पदाचे
मूल्यांकन पुढे चालू ठेवायचे).
आता या फलनांना निश्चित प्रांत आणि मूल्यसंच आहेत असे समजा.
$\lambd[x][N]$ चा प्रांत~$A$ आणि मूल्यसंच~$B$ असेल, तर
$x$ ला फक्त~$A$ मधील आर्ग्युमेंट घेता येतील आणि~$N$ ने
$B$ मधील मूल्ये द्यावीत. त्यामुळे $\lambd[x][N]$ हे
$A \to B$ असे फलन आणि $M \in A$ असेल, तेव्हाच
$(\lambd[x][N]M)$ ला अर्थ असावा. ही माहिती
$\lambd$-कलनात घातली की \emph{प्रकारयुक्त} $\lambd$-कलन
मिळते. प्रकारयुक्त $\lambd$-कलनात प्रत्येक
$(\lambd[x][N]M)$ अभिव्यक्ती ग्राह्य नसते;
$\lambd[x][N]$ आणि~$M$ यांचे ``प्रकार'' जुळतात,
म्हणजे $\lambd[x][N]$ चा प्रकार $A \lif B$ आणि~$M$ चा
प्रकार~$A$ असतो, तेव्हाच ती ग्राह्य असते.
$x$ चा प्रकार~$A$ आणि~$N$ चा प्रकार~$B$ असेल,
तेव्हाच $\lambd[x][N]$ चा प्रकार $A \to B$ असतो.
सर्वसाधारणपणे प्रत्येक $(M_1M_2)$ अभिव्यक्ती ग्राह्य नसते.
$(M_1M_2)$ हे पद ग्राह्य होण्यासाठी~$M_1$ चा प्रकार $A \to B$
आणि~$M_2$ चा प्रकार~$A$ असावा;
अशा वेळी $(M_1M_2)$ चा प्रकार~$B$ असतो.

हे प्रकार-नियम आता नैसर्गिक निगमनातील
निष्पत्ती शी स्पष्टपणे जुळतात: प्रकारांना
सूत्रे, आणि निष्पत्ती ना $\lambd$-पदे
अनुरूप असतात. उदाहरणार्थ, $A \lif B$ ची
निष्पत्ती $\lambd[\typeof{x}{A}][N]$
या पदाशी जुळते; त्याचा फलन-प्रकार $A \lif B$ आहे.
नैसर्गिक निगमनाचे अनुमान-नियम प्रकारयुक्त लॅम्डा कलनातील
प्रकार-नियमांशी जुळतात; उदाहरणार्थ,
\begin{prooftree}
  \AxiomC{$\Discharge{A}{x}$}
  \DeduceC{$B$}
  \DischargeRule{\Intro\lif}{x}
  \UnaryInfC{$A \lif B$}
  \DisplayProof\bottomAlignProof
  \quad\text{याला अनुरूप आहे}\quad
  \Axiom$x: A \fCenter N: B$
  \UnaryInf$ \fCenter \lambd[\typeof{x}{A}][N] : A \lif B$
\end{prooftree}
उजवीकडील नियमाचा अर्थ असा की~$x$ चा प्रकार~$A$ आणि
$N$ चा प्रकार~$B$ असेल, तर
$\lambd[\typeof{x}{A}][N]$ चा प्रकार~$A \lif B$ असतो.

$\Elim\lif$ हा नियम संयुक्त पदांसाठीच्या प्रकार-नियमाशी
जुळतो; म्हणजे,
\[
  \AxiomC{$A \lif B$}
  \AxiomC{$A$}
  \RightLabel{\Elim\lif}
  \BinaryInfC{$B$}
  \DisplayProof
  \quad\text{याला अनुरूप आहे}\quad
  \Axiom$\fCenter M_1: A \lif B$
  \Axiom$\fCenter M_2: A$
  \BinaryInf$ \fCenter (M_1M_2) : B$
  \DisplayProof
\]
\Intro\lif{} या नियमानंतर लगेच \Elim\lif{} नियम
आला, तर निष्पत्ती सोपी करता येते:
\begin{gather*}
  \AxiomC{$\Discharge{A}{x}$}
  \DeduceC{$B$}
  \DischargeRule{\Intro\lif}{x}
  \UnaryInfC{$A \lif B$}
  \AxiomC{}
  \DeduceC{$A$}
  \RightLabel{\Elim\lif}
  \BinaryInfC{$B$}
  \DisplayProof
  \quad\redone\quad
  \AxiomC{}
  \DeduceC{$A$}
  \DeduceC{$B$}
  \DisplayProof
\end{gather*}
हे लॅम्डा पदांच्या $\beta$-संकुचनाशी जुळते:
\[
(\lambd[\typeof{x}{A}][N]M) \quad\redone\quad
\Subst{N}{M}{x}.
\]

अशीच अनुरूपता $\land$ च्या नियमांचा ``गुणाकार''
प्रकारांशी आणि $\lor$ च्या नियमांचा ``बेरीज'' प्रकारांशीही आहे.

साध्या प्रकारयुक्त लॅम्डा कलनातील पदे आणि नैसर्गिक निगमनातील
निष्पत्ती यांतील या अनुरूपतेला ``करी--हॉवर्ड'',
``सिद्धता म्हणजे कार्यक्रम'', किंवा ``विधाने म्हणजे प्रकार''
अनुरूपता म्हणतात. सूत्रे (विधाने) प्रकारांशी आणि
सिद्धता पदांशी जुळतात; त्याशिवाय इतर अनुरूपता पुढीलप्रमाणे:
\begin{center}
  \begin{tabular}{c c c}
    तर्कशास्त्र & कार्यक्रम \\
    \hline
    विधान/सूत्र & प्रकार \\
    सिद्धता & पद \\
    गृहीतक & चर \\
    निरसित गृहीतक & बद्ध चर \\
    अनिरसित गृहीतक & मुक्त चर \\
    $A \lif B$ & फलन-प्रकार \\
    $A \land B$ & गुणाकार-प्रकार \\
    $A \lor B$ & बेरीज-प्रकार \\
    $\lfalse$ & रिक्त प्रकार \\
    \hline
  \end{tabular}
\end{center}

करी--हॉवर्ड अनुरूपता ही स्वयंचलित सिद्धता-साहाय्यकांची आणि
कार्यक्रमांसाठीच्या प्रकार-तपासकांची एक आधारभूत कल्पना आहे.
कारण एखादे विधान सिद्ध करणारी सिद्धता तपासणे (जसे आपण वर केले)
म्हणजे एखाद्या कार्यक्रमाला (पदाला) जाहीर केलेला प्रकार आहे का
हे तपासणे होय.












\section{सिद्धता-पदे}\label{int:pty:ter:sec}

\Log{N1i} मध्ये गृहीतकांना निरसन-खुणा म्हणून चलांची
चिन्हे देतो (उदा., $A^x$). \Log{N2i} मध्ये प्रत्येक
$\Gamma \Sequent A$ क्रमवर्तीच्या डावीकडील संदर्भ~$\Gamma$
हा सुसंगत चल–सूत्र जोड्यांचा संच असतो. त्यामुळे
\Log{N2i} मधील सिद्धता च्या प्रत्येक क्रमवर्तीवर
आतापर्यंतच्या सिद्धता ने काय सिद्ध केले आहे
(म्हणजे~$A$) एवढीच नव्हे, तर त्या टप्प्यावर कोणती
गृहीतके खुली आहेत (म्हणजे~$\Gamma$) ही माहितीही असते.
$A$ \emph{कसे} सिद्ध झाले याची माहितीही प्रत्येक
क्रमवर्तीत घातली तर? त्यासाठी \emph{पद-चिन्हांकित}
\Log{tN2i} पद्धती परिभाषित करता येते; तिच्या
क्रमवर्तींचे रूप $\Gamma \Sequent N : A$ असते.
$\Gamma$ आणि~$A$ यांचे अर्थ आधीप्रमाणेच आहेत:
$x: B$ या रूपातील जोड्यांचा संच, आणि क्रमवर्तीवर
संपणाऱ्या उप-सिद्धता ने~$\Gamma$ मधील गृहीतकांवरून
निष्पन्न होते असे दाखवलेले सूत्र. ~$N$ हे त्या
क्रमवर्तीवर संपणाऱ्या सिद्धता चे सांकेतिकीकरण करणारे
विशिष्ट पद असेल.

प्रत्येक क्रमवर्तीला देण्यात येणारी पदे $x$, $y$,
\dots या चलांपासून आणि आतापर्यंत झालेल्या अनुमानांचे
सांकेतिकीकरण करणाऱ्या फलन-चिन्हांपासून बनतात.
प्रत्येक अनुमान-नियमासाठी स्वतंत्र फलन-चिन्ह लागते.
परिचय नियमांना अनुरूप फलन-चिन्हांना
``रचनाकार'' आणि निरसन नियमांना अनुरूप चिन्हांना
``विघटक'' म्हणतात. नियमाची जितकी आधारविधाने, तितके
प्रत्येक फलन-चिन्हाचे आर्ग्युमेंट असतात. उदाहरणार्थ,
\Intro{\land} आणि \Elim{\land}[1] साठी अनुक्रमे
$c^\land$ (दोन-स्थानी) आणि $d^\land_1$ (एक-स्थानी)
ही फलन-चिन्हे वापरता येतील. \Log{tN2i} मध्ये हे
नियम असे होतील:
\[
\Axiom$\Gamma_1 \fCenter N: A$
\Axiom$\Gamma_2 \fCenter M: B$
\RightLabel{\Intro{\land}}
\BinaryInf$\Gamma_1, \Gamma_2 \fCenter c^\land(N, M) : A \land B $
\DisplayProof
\quad
\Axiom$\Gamma \fCenter N: A \land B$
\RightLabel{\Elim{\land}[1]}
\UnaryInf$\Gamma \fCenter d^\land_1(N) : A$
\DisplayProof
\]
अनुमान-नियम एखादे गृहीतक निरसित करत असेल,
म्हणजे आधारविधानात असलेले खुणित सूत्र~$x:A$
निष्कर्षाच्या संदर्भात नसेल, तर ते सिद्धता-पदातही
दर्शवावे लागते. अशा नियमांशी संबंधित फलन-चिन्हे
त्या चलांना \emph{बद्ध} करतात. उदाहरणार्थ,
\Intro{\lif} नियम $x:A, \Gamma \Sequent B$ या
आधारविधानापासून $\Gamma \Sequent A \lif B$
या निष्कर्षापर्यंत जातो. रचनाकार~$c^{\lif}$ एक
आर्ग्युमेंट घेतो, पण चल~$x$ बद्ध होतो हेही त्याने
दर्शवले पाहिजे. खूण~$x$ असलेले गृहीतक निरसित
झाले आहे हे सिद्धता-पद अशा प्रकारे व्यक्त करते.
उदाहरणार्थ, आधारविधानाचे सिद्धता-पद~$N$ असेल,
तर निष्कर्षाचे सिद्धता-पद $c^{\lif}([x]N)$ आणि
नियम असा लिहिता येईल:
\[
\Axiom$x:A, \Gamma \fCenter N: B$
\RightLabel{\Intro{\lif}}
\UnaryInf$\Gamma \fCenter c^{\lif}([x]N) : A \lif B$
\DisplayProof
\]

सिद्धता पूर्णपणे पुन्हा उभारण्यासाठी आणखी
काही माहिती लागते. नियमाच्या आधारविधानांतील
सूत्रे वरून निष्कर्षातील सूत्र पूर्णपणे
ठरत नसेल, तर ती माहिती पदात जोडली पाहिजे.
\Intro{\lif} मध्ये असे होते; म्हणून आपण
$c^{\lif}([x:A]N)$ असे लिहू. \Intro{\lor}
आणि~\FalseInt मध्येही असेच होते; म्हणून ती
माहिती वाहून नेणारी फलन-चिन्हे वापरू, उदाहरणार्थ:
\[
\Axiom$\Gamma \fCenter N:A$
\RightLabel{\Intro{\lor}[1]}
\UnaryInf$\Gamma \fCenter c^{\lor{B}}_1(N):A \lor B$
\DisplayProof
\quad
\Axiom$\Gamma \fCenter N:\lfalse$
\RightLabel{\FalseInt}
\UnaryInf$\Gamma \fCenter d^\lfalse_{A}(N):A$
\DisplayProof
\]

या कल्पनांवरून क्रमवर्ती-शैलीतील नैसर्गिक निगमनाची
अशी पद्धती बांधता येते की प्रत्येक क्रमवर्तीचे रूप
$\Gamma \Sequent N : A$ असते आणि~$N$ हे
\emph{सिद्धता-पद} असते.

अशी सिद्धता-पदे आणि तथाकथित \emph{प्रकारयुक्त लॅम्डा कलन}
यातील पदे यांचा जवळचा संबंध आहे. तो लक्षात घेऊन वरची
सर्वसाधारण रचनाकार-विघटक चिन्हे न वापरता लॅम्डा
कलनाच्या साहित्यातील चिन्हे वापरू. उदाहरणार्थ,
$c^{\lif}([x:A]N)$ ऐवजी $\lambd[\typeof{x}{A}][N]$,
$d^{\lif}(N,M)$ ऐवजी $(NM)$,
$c^\land(N,M)$ ऐवजी $\pair{N}{M}$, आणि
$d^\land_i(N)$ ऐवजी~$\proj{i}{N}$ लिहू.

\begin{defn}
  सिद्धता-पदे पुढील नियमांनी विगमनाने तयार होतात:
  \begin{enumerate}
  \item प्रत्येक चर~$x$ हे सिद्धता-पद आहे (स्वयंसिद्धके).
  \item $x$ हे चर, $A$ हे सूत्र, आणि~$N$ हे
    सिद्धता-पद असेल, तर $\lambd[\typeof{x}{A}][N]$
    हेही सिद्धता-पद आहे~(\Intro\lif).
  \item $N$ आणि~$M$ ही सिद्धता-पदे असतील, तर $(NM)$
    हेही सिद्धता-पद आहे~(\Elim\lif).
  \item $N$ आणि~$M$ ही सिद्धता-पदे असतील, तर
    $\pair{N}{M}$ हे सिद्धता-पद आहे~(\Intro\land).
  \item $N$ हे सिद्धता-पद असेल, तर $\proj{i}{N}$
    हेही सिद्धता-पद आहे; येथे~$i$ हे~$1$ किंवा~$2$
    असते~(\Elim\land[i]).
  \item $N$ हे सिद्धता-पद आणि~$A$ हे सूत्र असेल,
    तर $\inj[A]{i}{N}$ हे सिद्धता-पद आहे; येथे~$i$
    हे~$1$ किंवा~$2$ असते~(\Intro\lor[i]).
  \item $M$, $N_1$, $N_2$ ही सिद्धता-पदे आणि~$x_1$,
    $x_2$ ही चले असतील, तर
    $\dcase{M}{x_1}{N_1}{x_2}{N_2}$ हे सिद्धता-पद
    आहे~(\Elim\lor).
  \item $N$ हे सिद्धता-पद आणि~$A$ हे सूत्र
    असेल, तर $\abort{A}{N}$ हे सिद्धता-पद आहे~(\FalseInt).
  \end{enumerate}
\end{defn}

प्रत्येक सिद्धता-पद हे एखाद्या सिद्धता चे
सिद्धता-पद असेलच असे नाही. उदाहरणार्थ,
$\pair{N}{M}$ हे पद~\Intro{\land} अनुमानावर
संपणाऱ्या सिद्धतेचे सांकेतिकीकरण करते; म्हणून तिचे
अंतिम सूत्र $A \land B$ हे संयोगसूत्र असते.
हे~\Elim{\lif} नियमाचे मुख्य आधारविधान होऊ शकत
नाही; त्यामुळे $(\pair{N}{M}P)$ हे योग्य सिद्धतेचे
सिद्धता-पद होऊ शकत नाही. \Log{tN2ip} च्या नियमांचे
पालन करणारी सिद्धता-पदेच \emph{योग्य} सिद्धता-पदे
आहेत. हे नियम \ref{int:pty:ter:tab:tN2ip} मध्ये दिले आहेत.








\begin{table}
\begin{tabular}{@{}c@{}c@{}}
\multicolumn{2}{l}{\textbf{स्वयंसिद्धके:}}\\
\multicolumn{2}{c}{$x:A \Sequent x:A$}\\[3ex]
\multicolumn{2}{l}{\textbf{अनुमान-नियम:}}\\[3ex]
\multicolumn{2}{c}{\Axiom$x: A, \Gamma \fCenter N:B$
\RightLabel{\Intro{\lif}}
\UnaryInf$\Gamma \fCenter \lambd[\typeof{x}{A}][N] : A \lif B $
\DisplayProof}
\\[4ex]
\multicolumn{2}{c}{\Axiom$\Gamma_1 \fCenter N:A \lif B$
\Axiom$\Gamma_2 \fCenter M:A$
\RightLabel{\Elim{\lif}}
\BinaryInf$\Gamma_1, \Gamma_2 \fCenter (NM) : B$
\DisplayProof}
\\[4ex]
\multicolumn{2}{c}{\Axiom$\Gamma_1 \fCenter N_1 : A$
\Axiom$\Gamma_2 \fCenter N_2 : B$
\RightLabel{\Intro{\land}}
\BinaryInf$\Gamma_1, \Gamma_2 \fCenter \pair{N_1}{N_2} : A \land B $
\DisplayProof}
\\[4ex]
\Axiom$\Gamma \fCenter N: A \land B$
\RightLabel{\Elim{\land}[1]}
\UnaryInf$\Gamma \fCenter \proj{1}{N} : A$
\DisplayProof
&
\Axiom$\Gamma \fCenter N:A \land B$
\RightLabel{\Elim{\land}[2]}
\UnaryInf$\Gamma \fCenter \proj{2}{N} : B$
\DisplayProof\\[3ex]
\\[4ex]
\Axiom$\Gamma \fCenter N: A$
\RightLabel{\Intro{\lor}[1]}
\UnaryInf$\Gamma \fCenter \inj[B]{1}{N} : A \lor B$
\DisplayProof
&
\Axiom$ \Gamma \fCenter N : B$
\RightLabel{\Intro{\lor}[2]}
\UnaryInf$\Gamma \fCenter \inj[A]{2}{N} : A \lor B$
\DisplayProof\\[3ex]
\multicolumn{2}{c}{
\begin{tabular}[b]{@{}l@{}}
\Axiom$\Gamma_1 \fCenter M : A \lor B$
\Axiom$x:A, \Gamma_2 \fCenter N_1 : C$
\Axiom$y:B, \Gamma_3 \fCenter N_2: C$
\RightLabel{\Elim{\lor}}
\TrinaryInf$\Gamma_1, \Gamma_2, \Gamma_3 \fCenter
\dcase{M}{x}{N_1}{y}{N_2} : C$
\DisplayProof
\end{tabular}
}
\\[4ex]
\multicolumn{2}{c}{\Axiom$\Gamma \fCenter N: \lfalse$
\RightLabel{\FalseInt}
\UnaryInf$\Gamma \fCenter \abort{A}{N} : A$
\DisplayProof}
\end{tabular}
\caption{विधानीय पद-चिन्हांकित अंतःप्रज्ञावादी नैसर्गिक निगमन
पद्धतीचे नियम \Log{tN2ip}.}
\label{int:pty:ter:tab:tN2ip}
\end{table}
















\section{निष्पत्तीचे सिद्धता-पदांत रूपांतर}\label{int:pty:pt:sec}

विधानीय नैसर्गिक निगमनातील सिद्धता चे सिद्धता-पदांत रूपांतर,
म्हणजे~\Log{N2i} मधील सिद्धता चे~\Log{tN2i}
मधील सिद्धता मध्ये भाषांतर, अगदी सरळ आहे.
निष्पत्तीमधील प्रत्येक क्रमवर्तीच्या उजव्या
बाजूच्या सूत्राच्या डावीकडे सिद्धता-पद लिहायचे;
म्हणजे $\Gamma \Sequent M : A$ या रूपातील
क्रमवर्ती मिळतात. मग संदर्भ~$\Gamma$ मध्ये~$M$ हे
$A$ ची \emph{साक्ष देते} असे म्हणू. कोणते सिद्धता-पद
लिहायचे हे पूर्णपणे \Log{tN2i} च्या नियमांनी ठरते.

\begin{prop}
विधानीय \Log{N2i} मधील $\Gamma \Sequent A$ ची प्रत्येक
सिद्धता~$\delta$ ही~\Log{tN2i} मधील
$\Gamma \Sequent N : A$ च्या सिद्धता~$\delta'$ शी
अनुरूप असते. सिद्धता-पद~$N$ हे~$\delta$ वरून
एकमेवपणे ठरते.
\end{prop}

\begin{proof}
~$\delta$ च्या उंचीवर विगमन करू. $\delta$ हे
$x:A \Sequent A$ हे स्वयंसिद्धक असेल, तर
$\delta'$ हे~$x:A \Sequent x:A$ हे स्वयंसिद्धक आहे.
$\delta$ एखाद्या अनुमान-नियमावर संपत असेल, तर
विगमन-गृहीतक प्रत्येक आधारविधानासाठी सिद्धता-पद आणि
ते पद असलेल्या क्रमवर्तीची \Log{tN2i}-सिद्धता देते.
मग \Log{tN2i} चा अनुरूप नियम लावतो;
\Log{tN2i} चा अनुमान-नियमच अनुरूप सिद्धता-पद निश्चित करतो.
\end{proof}
येथे दिलेल्या पद-नियमांची व्याप्ती विधानीय आहे;
संख्यापकांचे स्वतंत्र पद-संकारक परिभाषित केलेले नाहीत.

\begin{ex}
  \Log{N1i} मधील $(A \land B) \lif (B \land A)$
  ची पुढील अगदी साधी सिद्धता घ्या:
  \[
  \AxiomC{$\Discharge{A\land B}{x}$}
  \RightLabel{\Elim\land[2]}
  \UnaryInfC{$B$}
  \AxiomC{$\Discharge{A\land B}{x}$}
  \RightLabel{\Elim\land[1]}
  \UnaryInfC{$A$}
  \RightLabel{\Intro\land}
  \BinaryInfC{$B \land A$}
  \DischargeRule{\Intro\lif}{x}
  \UnaryInfC{$(A \land B) \lif (B \land A)$}
  \DisplayProof
  \]
  \Log{N2i} मध्ये ती अशी दिसते:
  \[
  \Axiom$x : A\land B \fCenter A\land B$
  \RightLabel{\Elim\land[2]}
  \UnaryInf$x : A\land B \fCenter B$
  \Axiom$x : A\land B \fCenter A\land B$
  \RightLabel{\Elim\land[1]}
  \UnaryInf$x : A\land B \fCenter A$
  \RightLabel{\Intro\land}
  \BinaryInf$x : A\land B \fCenter B \land A$
  \RightLabel{\Intro\lif}
  \UnaryInf$\fCenter (A \land B) \lif (B \land A)$
  \DisplayProof
  \]
  सिद्धता-पदे वरपासून खाली देऊ. स्वयंसिद्धकांची
  सिद्धता-पदे ही संदर्भातील~$x$ या चलासारखीच आहेत.
  मग $\Elim\land[i]$ शी संबंधित सिद्धता-पदे अनुक्रमे
  $\proj{2}{x}$ आणि~$\proj{1}{x}$ अशी ठरतात.
  \Intro{\land} लावताना ही दोन सिद्धता-पदे एकत्र करतो:
  $\pair{\proj{2}{x}}{\proj{1}{x}}$.
  शेवटी \Intro{\lif} नियमाने $\lambd$-अमूर्तीकरण
  केले जाते; त्यामुळे~$x$ बद्ध होतो आणि~$x$ ही खूण
  $A \land B$ या गृहीतकाला होती हे नोंदले जाते:
  $\lambd[\typeof{x}{A \land
  B}][\pair{\proj{2}{x}}{\proj{1}{x}}]$.
  मग \Log{tN2i} मधील सिद्धता अशी आहे:
  \[
  \Axiom$x : A\land B \fCenter x: A\land B$
  \RightLabel{\Elim\land[2]}
  \UnaryInf$x : A\land B \fCenter \proj{2}{x} : B$
  \Axiom$x : A\land B \fCenter x: A\land B$
  \RightLabel{\Elim\land[1]}
  \UnaryInf$x : A\land B \fCenter \proj{1}{x} : A$
  \RightLabel{\Intro\land}
  \BinaryInf$x : A\land B \fCenter \pair{\proj{2}{x}}{\proj{1}{x}} :
    B \land A$
  \RightLabel{\Intro\lif}
  \UnaryInf$\fCenter \lambd[\typeof{x}{A \land
  B}][\pair{\proj{2}{x}}{\proj{1}{x}}] : (A \land B) \lif (B \land A)$
  \DisplayProof
  \]
\end{ex}












\section{सिद्धता-पदांवरून निष्पत्ती पुन्हा मिळवणे}\label{int:pty:tp:sec}

आता उलट दिशा पाहू: सिद्धता-पद~$N$ वरून
\Log{N2i} मधील सिद्धता पुन्हा मिळवणे.
सर्वसाधारणपणे हे दिलेल्या सिद्धता-पदात मुक्त असलेल्या
प्रत्येक चल~$x$ साठी $x : A$ अशी जोडी धरणाऱ्या
संदर्भाच्या सापेक्षच करता येते. पण आधी मुक्त चले
नसलेले उदाहरण घेऊ, म्हणजे हे पद:
\[
\lambd[\typeof{x}{A \land
  B}][\pair{\proj{2}{x}}{\proj{1}{x}}]
\]
याचे रूप $\lambd[\typeof{x}{A \land B}][N_1]$ आहे.
अशा रूपाचे पद निर्माण करणारा \Log{tN2i} चा
एकमेव नियम~\Intro{\lif} असल्यामुळे~$N$ ने
सांकेतिक केलेली सिद्धता~\Intro{\lif} वर संपली
पाहिजे, म्हणजे तिचा शेवट असा आहे:
\[
  \Axiom$x : A\land B \fCenter \pair{\proj{2}{x}}{\proj{1}{x}} :
    C$
  \RightLabel{\Intro\lif}
  \UnaryInf$\fCenter \lambd[\typeof{x}{A \land
  B}][\pair{\proj{2}{x}}{\proj{1}{x}}] : (A \land B) \lif C$
  \DisplayProof
  \]
$C$ काय आहे हे अद्याप माहीत नाही; पण पूर्वांग
$A \land B$ आहे, आणि आधारविधानाच्या संदर्भात
$x:A \land B$ असलेच पाहिजे हे माहीत आहे.

आता आधारविधानाशी संबंधित सिद्धता-पद निश्चित झाले आहे:
$\pair{\proj{2}{x}}{\proj{1}{x}}$.
\emph{त्याने} सांकेतिक केलेली सिद्धता~\Intro{\land}
वर संपली पाहिजे. ~$C$ काय आहे हे अजून माहीत
नाही; पण तो~\Intro{\land} चा निष्कर्ष असल्यामुळे
संयोग असलाच पाहिजे, म्हणजे
$C \ident (C_1 \land C_2)$. पुनर्रचना पुढे चालू ठेवू:
\[
  \Axiom$x : A\land B \fCenter \proj{2}{x} : C_1$
  \Axiom$x : A\land B \fCenter \proj{1}{x} : C_2$
  \RightLabel{\Intro\land}
  \BinaryInf$x : A\land B \fCenter \pair{\proj{2}{x}}{\proj{1}{x}} :
    C_1 \land C_2$
  \RightLabel{\Intro\lif}
  \UnaryInf$\fCenter \lambd[\typeof{x}{A \land
  B}][\pair{\proj{2}{x}}{\proj{1}{x}}] : (A \land B) \lif (C_1 \land C_2)$
  \DisplayProof
  \]
$\pi_2(x)$ हे सिद्धता-पद सांगते की डावीकडील वरची
क्रमवर्ती~\Elim{\land}[2] ने मिळाली; म्हणजे तिचे
आधारविधान $D \land C_1$ या रूपाचे आहे. उजवीकडे
$C_2$ कडे नेणारे अनुमान~\Elim{\land}[1] असले
पाहिजे आणि त्याचे आधारविधान $C_2 \land E$ या
रूपाचे असले पाहिजे, हे ठरवता येते.
\[
  \Axiom$x : A\land B \fCenter x: D\land C_1$
  \RightLabel{\Elim\land[2]}
  \UnaryInf$x : A\land B \fCenter \proj{2}{x} : C_1$
  \Axiom$x : A\land B \fCenter x: C_2\land E$
  \RightLabel{\Elim\land[1]}
  \UnaryInf$x : A\land B \fCenter \proj{1}{x} : C_2$
  \RightLabel{\Intro\land}
  \BinaryInf$x : A\land B \fCenter \pair{\proj{2}{x}}{\proj{1}{x}} :
    C_1 \land C_2$
  \RightLabel{\Intro\lif}
  \UnaryInf$\fCenter \lambd[\typeof{x}{A \land
  B}][\pair{\proj{2}{x}}{\proj{1}{x}}] : (A \land B) \lif (C_1 \land C_2)$
  \DisplayProof
  \]
सर्वांत वरच्या क्रमवर्तींची सिद्धता-पदे आता चले
असल्यामुळे त्या स्वयंसिद्धकेच असल्या पाहिजेत.
त्यावरून $C_1$, $C_2$, $D$, आणि $E$ ठरतात:
$D \ident A$ आणि $C_1 \ident B$ नसतील,
तर डावी क्रमवर्ती स्वयंसिद्धक नाही; तसेच
$C_2 \ident A$ आणि $E \ident B$ नसतील,
तर उजवी क्रमवर्तीही स्वयंसिद्धक नाही.
अशा रीतीने संपूर्ण सिद्धता पुन्हा मिळाली:
\[
  \Axiom$x : A\land B \fCenter x: A\land B$
  \RightLabel{\Elim\land[2]}
  \UnaryInf$x : A\land B \fCenter \proj{2}{x} : B$
  \Axiom$x : A\land B \fCenter x: A\land B$
  \RightLabel{\Elim\land[1]}
  \UnaryInf$x : A\land B \fCenter \proj{1}{x} : A$
  \RightLabel{\Intro\land}
  \BinaryInf$x : A\land B \fCenter \pair{\proj{2}{x}}{\proj{1}{x}} :
    B \land A$
  \RightLabel{\Intro\lif}
  \UnaryInf$\fCenter \lambd[\typeof{x}{A \land
  B}][\pair{\proj{2}{x}}{\proj{1}{x}}] : (A \land B) \lif (B \land A)$
  \DisplayProof
  \]













\section{प्रकार}\label{int:pty:typ:sec}

$(MN)$ सारखे सिद्धता-पद स्वतःहून कोणत्या सिद्धतेचे
सिद्धता-पद आहे हे एकमेवपणे ठरवत नाही. कारण मुक्त चले,
म्हणजे स्वयंसिद्धकांची सिद्धता-पदे, संबंधित स्वयंसिद्धकांतील
सूत्रे सांगत नाहीत. पण ती सुसंगत संदर्भ~$\Gamma$ मध्ये
दिली आणि पद योग्य प्रकाराचे असेल, तर सिद्धता \emph{एकमेवपणे} ठरते.

\begin{defn}
सिद्धता-पद~$N$ साठी \emph{संदर्भ} म्हणजे प्रकार-खुणांचा
सुसंगत संच~$\Gamma$: कोणत्याही चलाला दोन वेगळे प्रकार
देता येत नाहीत, आणि~$N$ च्या प्रत्येक मुक्त चलासाठी
$\Gamma$ मध्ये $x:A$ अशी नेमकी एक जोडी असते.
\end{defn}

या व्याख्येनुसार~$\Gamma$ मध्ये $x:A$ ही जोडी
\emph{फक्त}~$N$ च्या मुक्त चलांसाठीच असावी अशी अट नाही.
उदाहरणार्थ, $\Gamma = \{x:A, y:B, z:C\}$ हा
$\pair{x}{y}$ साठी संदर्भ आहे.

\begin{defn}\label{int:pty:typ:defn:type}
  ~$\Gamma$ हा~$N$ साठी संदर्भ आहे असे माना.
  ~$\Gamma$ च्या सापेक्ष~$N$ चा \emph{प्रकार}
  पुढीलप्रमाणे विगमनाने परिभाषित करतो:
  \begin{enumerate}
  \item $x:A \in \Gamma$ असेल तेव्हाच~$x$ चा प्रकार~$A$ असतो.
  \item $x:A, \Gamma$ च्या सापेक्ष~$N$ चा प्रकार~$B$
    असेल, तर $\lambd[\typeof{x}{A}][N]$ चा प्रकार
    $A \lif B$ असतो.
  \item $N$ चा प्रकार $A \lif B$ आणि~$M$ चा प्रकार~$A$
    असेल, तर $(NM)$ चा प्रकार~$B$ असतो.
  \item $N$ चा प्रकार~$A$ आणि~$M$ चा प्रकार~$B$
    असेल, तर $\pair{N}{M}$ चा प्रकार $A \land B$ असतो.
  \item $N$ चा प्रकार~$A_1 \land A_2$ असेल, तर
    $\proj{i}{N}$ चा प्रकार~$A_i$ असतो.
  \item $N$ चा प्रकार~$A$ असेल, तर
    $\inj[B]{i}{N}$ चा प्रकार $i=1$ असताना
    $A \lor B$ आणि $i=2$ असताना $B \lor A$ असतो.
  \item $M$ चा प्रकार $A \lor B$, तसेच
    $x_1:A, \Gamma$ च्या सापेक्ष~$N_1$ चा प्रकार~$C$
    आणि $x_2:B, \Gamma$ च्या सापेक्ष~$N_2$ चा
    प्रकार~$C$ असेल, तर
    $\dcase{M}{x_1}{N_1}{x_2}{N_2}$ चा प्रकार~$C$ असतो.
  \item $N$ चा प्रकार~$\lfalse$ असेल, तर
    $\abort{A}{N}$ चा प्रकार~$A$ असतो.
  \end{enumerate}
\end{defn}
या नियमांत न बसणाऱ्या पदाला प्रकार नाही. अमूर्तीकरणाच्या
आणि प्रकरण-विभाजनाच्या बद्ध चलांना संदर्भातील चलांपासून वेगळी
नवी नावे देऊन वरील संदर्भ-वाढ करायची.

\begin{prop}
सिद्धता-पद~$N$ साठीच्या सुसंगत संदर्भ~$\Gamma$ च्या
सापेक्ष~$N$ ला जास्तीत जास्त एक प्रकार असतो. प्रकार
अस्तित्वात असेल, तर तो एकमेव असतो; प्रत्येक कच्च्या
सिद्धता-पदाला प्रकार मिळतोच असे नाही.
\end{prop}

\begin{proof}
  ~$N$ वर विगमनाने.
\end{proof}

संदर्भ~$\Gamma$ च्या सापेक्ष~$N$ चा प्रकार~$A$
असेल, तर $\Gamma \Sequent N : A$ असे लिहितो.
\ref{int:pty:typ:defn:type} मधील उपवाक्ये परिचित वाटली पाहिजेत:
त्यात एक लहानसा फरक वगळता \Log{tN2i} चेच नियम
आहेत—संदर्भ~$\Gamma$ सर्वत्र तोच राहतो. ही
व्याख्या \Log{tN3i} या कलनाच्या रूपात मांडता येते;
त्याचे नियम \ref{int:pty:typ:tab:tN3ip} मध्ये दिले आहेत.









\begin{table}
\begin{tabular}{@{}c@{}c@{}}
\multicolumn{2}{l}{\textbf{स्वयंसिद्धके:}}\\
\multicolumn{2}{c}{$\Gamma \Sequent x:A$ जेव्हा $x:A \in \Gamma$}\\[3ex]
\multicolumn{2}{l}{\textbf{अनुमान-नियम:}}\\[3ex]
\multicolumn{2}{c}{\Axiom$x: A, \Gamma \fCenter N:B$
\RightLabel{\Intro{\lif}}
\UnaryInf$\Gamma \fCenter \lambd[\typeof{x}{A}][N] : A \lif B $
\DisplayProof}
\\[4ex]
\multicolumn{2}{c}{\Axiom$\Gamma \fCenter N:A \lif B$
\Axiom$\Gamma \fCenter M:A$
\RightLabel{\Elim{\lif}}
\BinaryInf$\Gamma \fCenter (NM) : B$
\DisplayProof}
\\[4ex]
\multicolumn{2}{c}{\Axiom$\Gamma \fCenter N_1 : A$
\Axiom$\Gamma \fCenter N_2 : B$
\RightLabel{\Intro{\land}}
\BinaryInf$\Gamma \fCenter \pair{N_1}{N_2} : A \land B $
\DisplayProof}
\\[4ex]
\Axiom$\Gamma \fCenter N: A \land B$
\RightLabel{\Elim{\land}[1]}
\UnaryInf$\Gamma \fCenter \proj{1}{N} : A$
\DisplayProof
&
\Axiom$\Gamma \fCenter N:A \land B$
\RightLabel{\Elim{\land}[2]}
\UnaryInf$\Gamma \fCenter \proj{2}{N} : B$
\DisplayProof\\[3ex]
\\[4ex]
\Axiom$\Gamma \fCenter N: A$
\RightLabel{\Intro{\lor}[1]}
\UnaryInf$\Gamma \fCenter \inj[B]{1}{N} : A \lor B$
\DisplayProof
&
\Axiom$ \Gamma \fCenter N : A$
\RightLabel{\Intro{\lor}[2]}
\UnaryInf$\Gamma \fCenter \inj[B]{2}{N} : B \lor A$
\DisplayProof\\[3ex]
\multicolumn{2}{c}{
\begin{tabular}[b]{@{}l@{}}
\Axiom$\Gamma \fCenter M : A \lor B$
\Axiom$x:A, \Gamma \fCenter N_1 : C$
\Axiom$y:B, \Gamma \fCenter N_2: C$
\RightLabel{\Elim{\lor}}
\TrinaryInf$\Gamma \fCenter
\dcase{M}{x}{N_1}{y}{N_2} : C$
\DisplayProof
\end{tabular}
}
\\[4ex]
\multicolumn{2}{c}{\Axiom$\Gamma \fCenter N: \lfalse$
\RightLabel{\FalseInt}
\UnaryInf$\Gamma \fCenter \abort{A}{N} : A$
\DisplayProof}
\end{tabular}
\caption{प्रकार-निर्देशन पद्धती म्हणून वापरल्या जाणाऱ्या
\Log{tN3ip} या विधानीय पद-चिन्हांकित अंतःप्रज्ञावादी
नैसर्गिक निगमन पद्धतीचे नियम.}
\label{int:pty:typ:tab:tN3ip}
\end{table}




सिद्धता-पदे ही $\lambd$-कलनाच्या एका आवृत्तीतील
पदे आहेत: \emph{गुणाकार, बेरीज आणि रिक्त प्रकार
असलेले प्रकारयुक्त $\lambd$-कलन}. पारंपरिक
$\lambd$-कलनात फक्त चलांपासून बनलेली पदे,
\emph{$\lambd$-अमूर्तीकरणे}~$\lambd[x][N]$, आणि
उपयोजन $(MN)$ असते; त्यात प्रकारयुक्त
$\lambd$-अमूर्तीकरण~$\lambd[\typeof{x}{A}][N]$
यातील~$A$ हे प्रकार-चिन्ह नसते. ~$\lambd$-कलन
हे ट्यूरिंग-संपूर्ण संगणन-प्रतिमान आहे (म्हणजे प्रत्येक
आंशिक संगणनीय फलन त्यातील पदाने दर्शवता येते).
सिद्धता-पदे आणि प्रकार-नियम देतात ती~$\lambd$-कलनाची
आवृत्ती निर्बंधित संगणन-प्रतिमान आहे; पण तिच्या मूलभूत
कल्पना त्याच आहेत. हे निर्बंध प्रकार-निर्देशनातून येतात:
फक्त \emph{योग्य प्रकाराची} पदे, म्हणजे
संदर्भ~$\Gamma$ च्या सापेक्ष प्रकार असलेली
सिद्धता-पदे, ग्राह्य असतात. उदाहरणार्थ,
$\lambd[x][(xx)]$ हे पारंपरिक प्रकारविरहित
$\lambd$-कलनात ग्राह्य पद आहे; पण
$\lambd[\typeof{x}{A}][(xx)]$ हे कोणत्याही~$A$
साठी योग्य प्रकाराचे नाही. कारण $(xx)$ ला प्रकार
असण्यासाठी पहिल्या~$x$ चा प्रकार $A \lif B$
आणि दुसऱ्या~$x$ चा प्रकार~$A$ असावा लागेल.
पण~$A$ हे $A \lif B$ चे योग्य उप-सूत्र
असल्यामुळे हे अशक्य आहे.

एखाद्या प्रकाराचे पद ज्या वस्तूला नाव देते, तिचा
प्रकार म्हणून प्रकारांकडे सहजपणे पाहता येते. म्हणून
$A \land B$ या प्रकाराचे पद पहिला घटक~$A$
आणि दुसरा~$B$ या प्रकाराचा असलेली जोडी निवडते;
म्हणजे गुणाकार~$A \times B$ मधील घटक.
$A \lif B$ या प्रकाराचे पद प्रकार~$A$ च्या
वस्तूंपासून प्रकार~$B$ च्या वस्तूंकडे जाणारे फलन आहे.
$A \lor B$ या प्रकारचे पद हे प्रकार~$A$ किंवा
प्रकार~$B$ ची वस्तू असते; तसेच ती दोहोंपैकी कोणत्या
प्रकाराची आहे ही माहिती त्यात असते. हे~$A$ आणि~$B$
या दोन संचांच्या असंयुक्त संघासारखे आहे:
$\{\tuple{i,x} : i = 1 \text{ किंवा } 2, x \in A
\text{ जर } i = 1, x \in B \text{ जर } i=2\}$.
$\lfalse$ हा प्रकार रिक्त संच आहे; म्हणजे त्या प्रकाराचे
पद कोणत्याही वस्तूला नाव देऊ शकत नाही.
नैसर्गिक निगमन (\Log{tN3i} सह) सुसंगत असल्यामुळे
खुली गृहीतके नसताना~$\lfalse$ ची सिद्धता नसते.
म्हणून~$\lfalse$ प्रकाराचे बंद सिद्धता-पद
(मुक्त चले नसलेले पद) नसते.

प्रकारयुक्त $\lambd$-कलनाचे संकारक योग्य प्रकाराच्या
वस्तूंना सहजपणे नाव देणारी पदे तयार करतात, अर्थात
त्यांना लावलेली पदेही ठरावीक प्रकारच्या वस्तूंना नाव
देत असतील तर. उदाहरणार्थ, ``$N_1 : A$ आणि
$N_2 : B$ असतील, तर
$\pair{N_1}{N_2} : A \land B$'' याचा अर्थ:
$N_1$ ने प्रकार~$A$ च्या वस्तूला आणि~$N_2$ ने
प्रकार~$B$ च्या वस्तूला नाव दिले, तर
$\pair{N_1}{N_2}$ हे प्रकार $A \land B$ च्या
वस्तूला नाव देते.













\section{न्यूनीकरण}\label{int:pty:red:sec}

नैसर्गिक निगमनाच्या निष्पत्तीमध्ये
परिचय नियमापाठोपाठ निरसन नियम
आला, तर तो वळसा ``न्यूनीत'' करता येतो. उदाहरणार्थ,
$\Intro{\land}$ नंतर $\Elim{\land}$ आले, तर ते
``एकमेकांना रद्द'' करतात: $\Elim{\land}$ चा निष्कर्ष
नेहमीच \Intro{\land} ने तयार केलेल्या संयोगातील
एका घटकासारखा असतो, आणि $\Intro{\land}$ ची दोन्ही
आधारविधाने तेच संयोगघटक असतात. म्हणून एखाद्या
संयोगघटकाची निष्पत्ती हवी असेल, तर संयोग
प्रस्तुत करून त्याचे निरसन न करता तीच निष्पत्ती
थेट वापरता येते. आपल्याला असे मिळते:
\begin{gather*}
  \bottomAlignProof
  \AxiomC{}
  \RightLabel{$\delta_1$}
  \DeduceC{$A_1$}
  \AxiomC{}
  \RightLabel{$\delta_2$}
  \DeduceC{$A_2$}
  \RightLabel{$\Intro{\land}$}
  \BinaryInfC{$A_1 \land A_2$}
  \RightLabel{$\Elim{\land}_i$}
  \UnaryInfC{$A_i$}
  \DisplayProof\bottomAlignProof
  \quad
  \redone
  \quad
  \AxiomC{}
  \RightLabel{$\delta_i$}
  \DeduceC{$A_i$}
  \DisplayProof
\end{gather*}

या सोपीकरणातून संबंधित सिद्धता-पदांवरील तसाच
न्यूनीकरण-संबंध दिसतो. ~$N_1$ हे~$\delta_1$ चे
आणि~$N_2$ हे~$\delta_2$ चे सिद्धता-पद असेल,
तर डावीकडील सिद्धतेचे सिद्धता-पद एका पायरीत
उजवीकडील सिद्धतेच्या सिद्धता-पदात न्यूनीत होते:
\[
  \proj{i}{\pair{N_1}{N_2}} \redone N_i
\]
प्रकारयुक्त लॅम्डा कलनात हा गुणाकार-प्रकाराचा
\emph{बीटा-न्यूनीकरण} नियम आहे.

नैसर्गिक निगमनात \Intro{\land} नंतर
\Elim{\land} अशी अनुमानांची जोडी, म्हणजे न्यूनीकरण
लागू होणारी जोडी, हिला \emph{कट} म्हणतात. प्रकारयुक्त
लॅम्डा कलनात $\redone$ च्या डावीकडील पदाला
\emph{रेडेक्स}, तर उजवीकडील पदाला \emph{संकुचनफल}
म्हणतात. प्रकारविरहित लॅम्डा कलनात फक्त
$(\lambd[x][N])Q$ यालाच रेडेक्स मानतात; पण
प्रकारयुक्त लॅम्डा कलनात $\pair{N}{M}$,
$\proj{i}{N}$, $\inj[A]{i}{N}$,
$\dcase{N}{x_1}{M_1}{x_2}{M_2}$, आणि
$\abort{A}{N}$ यांसारखी पदेही विन्यासात येतात,
आणि त्यांना अनुरूप रेडेक्सही असतात.

त्याचप्रमाणे वियोगासाठीही न्यूनीकरण आहे:
\begin{gather*}
  \bottomAlignProof
  \AxiomC{}
  \RightLabel{$\delta$}
  \DeduceC{$A_i$}
  \RightLabel{$\Intro{\lor}$}
  \UnaryInfC{$A_1 \lor A_2$}
  \AxiomC{$\Discharge{A_1}{x_1}$}
  \RightLabel{$\delta_1$}
  \DeduceC{$C$}
  \AxiomC{$\Discharge{A_2}{x_2}$}
  \RightLabel{$\delta_2$}
  \DeduceC{$C$}
  \DischargeRule{$\Elim{\lor}$}{x_1,x_2}
  \TrinaryInfC{$C$}
  \DisplayProof\bottomAlignProof
  \quad
  \redone
  \quad
  \AxiomC{}
  \RightLabel{$\delta$}
  \DeduceC{$A_i$}
  \RightLabel{$\delta_i$}
  \DeduceC{$C$}
  \DisplayProof
\end{gather*}
याला अनुरूप सिद्धता-पदांवरील न्यूनीकरण असे आहे:
\begin{gather*}
  \dcase{\inj[A_{3-i}]{i}{M}}{x_1}{N_1}{x_2}{N_2} \redone \Subst{N_i}{M}{x_i}
\end{gather*}
येथे~$M$ हे~$\delta$ चे, आणि~$N_i$ हे~$\delta_i$
चे सिद्धता-पद आहे. हा बेरीज-प्रकारांसाठीचा
बीटा-न्यूनीकरण नियम आहे. येथे
$\Subst{M}{N_i}{x}$ म्हणजे~$M$ मधील~$x$ या
चराची सर्व मुक्त स्थाने~$N_i$ ने बदलणे.
हे मुक्त-चल प्रतिस्थापन बद्ध-चल ग्रहण होऊ न देता
करायचे; आवश्यक बद्ध चलांना आधी नवी नावे द्यायची.
विन्यासातील प्रत्येक संकारकालाच रेडेक्स म्हणत नाही:
मुख्य बीटा-रेडेक्स हे वरच्या गुणाकाराच्या आणि पुढच्या
बेरीज व फलनाच्या विशिष्ट रूपांचे असतात.

फलन-प्रकाराचे न्यूनीकरण म्हणजे $\Intro\lif$ नंतर
$\Elim\lif$ आल्याने झालेला वळसा काढणे:
\begin{gather*}
  \bottomAlignProof
  \AxiomC{$\Discharge{A}{x}$}
  \RightLabel{$\delta$}
  \DeduceC{$B$}
  \DischargeRule{$\Intro{\lif}$}{x}
  \UnaryInfC{$A \lif B$}
  \AxiomC{}
  \RightLabel{$\delta'$}
  \DeduceC{$A$}
  \RightLabel{$\Elim{\lif}$}
  \BinaryInfC{$B$}
  \DisplayProof\bottomAlignProof
  \quad
  \redone
  \quad
  \AxiomC{}
  \RightLabel{$\delta'$}
  \DeduceC{$A$}
  \RightLabel{$\delta$}
  \DeduceC{$B$}
  \DisplayProof
\end{gather*}
सिद्धता-पदांसाठी हे नेहमीचे बीटा-न्यूनीकरण आहे:
\begin{gather*}
  (\lambd[\typeof{x}{A}][N]M)
  \redone \Subst{N}{M}{x}
\end{gather*}

सिद्धता वरील न्यूनीकरण रूपांतरणांखेरीज
क्रमपरिवर्तन रूपांतरणांचाही विचार करता येतो.
ती \Elim{\lor} नंतर दुसरा निरसन नियम
असलेल्या (उप-)सिद्धता वर लागू होतात, उदाहरणार्थ:
\begin{multline*}
  \bottomAlignProof
  \AxiomC{}
  \RightLabel{$\delta$}
  \DeduceC{$A_i$}
  \RightLabel{$\Intro{\lor}$}
  \UnaryInfC{$A_1 \lor A_2$}
  \AxiomC{$\Discharge{A_1}{x_1}$}
  \RightLabel{$\delta_1$}
  \DeduceC{$C_1 \land C_2$}
  \AxiomC{$\Discharge{A_2}{x_2}$}
  \RightLabel{$\delta_2$}
  \DeduceC{$C_1 \land C_2$}
  \DischargeRule{$\Elim{\lor}$}{x_1,x_2}
  \TrinaryInfC{$C_1 \land C_2$}
  \RightLabel{\Elim{\land}[i]}
  \UnaryInfC{$C_i$}
  \DisplayProof\bottomAlignProof
  \quad
  \redone
  \\
  \AxiomC{}
  \RightLabel{$\delta$}
  \DeduceC{$A_i$}
  \RightLabel{$\Intro{\lor}$}
  \UnaryInfC{$A_1 \lor A_2$}
  \AxiomC{$\Discharge{A_1}{x_1}$}
  \RightLabel{$\delta_1$}
  \DeduceC{$C_1 \land C_2$}
  \RightLabel{\Elim{\land}[i]}
  \UnaryInfC{$C_i$}
  \AxiomC{$\Discharge{A_2}{x_2}$}
  \RightLabel{$\delta_2$}
  \DeduceC{$C_1 \land C_2$}
  \RightLabel{\Elim{\land}[i]}
  \UnaryInfC{$C_i$}
  \DischargeRule{$\Elim{\lor}$}{x_1,x_2}
  \TrinaryInfC{$C_i$}
  \DisplayProof
\end{multline*}
~$\delta$, $\delta_1$, आणि~$\delta_2$ यांची
सिद्धता-पदे अनुक्रमे~$M$, $N_1$, आणि~$N_2$
असतील, तर सिद्धता-पदांसाठी, म्हणजेच~$\lambd$
पदांसाठी, अनुरूप न्यूनीकरण नियम असा आहे:
\[
  \proj{i}{\dcase{M}{x_1}{N_1}{x_2}{N_2}} \redone
\dcase{M}{x_1}{\proj{i}{N_1}}{x_2}{\proj{i}{N_2}}
\]
इतर निरसनांसाठीही वियोग-निरसनाच्या निष्कर्षावरील
निरसन दोन्ही शाखांच्या निष्कर्षांवर हलवायचे. बाहेरील
आदान किंवा उपपद आत जाताना~$x_1,x_2$ यांनी त्याचे
मुक्त चल पकडू नयेत म्हणून या बद्ध चलांना आधी नवी
नावे द्यायची. असत्याच्या सरलीकरणांसाठी नैसर्गिक
निगमनाच्या सामान्यरूपीकरण प्रकरणातील अनुरूप
रूपांतरणे वापरायची.

अर्थात, वळसे केवळ सिद्धता च्या शेवटीच नव्हे,
तर सिद्धता च्या आतही काढता येतात: वळशावर संपणाऱ्या
उप-सिद्धता ला न्यूनीकरण रूपांतरण लावायचे.
तसेच $\lambd$-पदांच्या उपपदांवर $\beta$-न्यूनीकरण
लागू होते. $N$ हे~$M$ चे उपपद असेल आणि
$N \redone N'$ असेल, तर~$M$ मधील उपपद~$N$ हे~$N'$ ने
बदलून~$M'$ मिळवता येते. अशा वेळीही $M \redone M'$
असे लिहितो. $M = M_1 \redone M_2 \redone M_3
\redone \dots \redone M_k = M'$ असा क्रम असेल
(यात $k=1$, म्हणजे $M = M'$ ही बाबही येते),
तर $M \red M'$ असे लिहितो.

ज्या सिद्धता च्या कोणत्याही उप-सिद्धता वर
रूपांतरण लागू होत नाही, तिला \emph{सामान्य} म्हणतात.
तसेच ज्या $\lambd$-पदावर (त्याच्या कोणत्याही उपपदावर)
रूपांतरण लागू होत नाही, त्याला \emph{सामान्य} म्हणतात.
सामान्य $\lambd$-पदाला \emph{सामान्य रूप} असेही म्हणतात.














\section{प्रकार-संरक्षण}\label{int:pty:tpr:sec}

विचारात घेतलेली न्यूनीकरण आणि क्रमपरिवर्तन रूपांतरणे
\Log{tN2i} आणि \Log{tN3i} या पद्धतींसाठीही थेट
परिभाषित करता येतात. \Log{tN3i} पद्धत संदर्भ~$\Gamma$
च्या सापेक्ष $\lambd$-पद~$N$ चा प्रकार असे
सूत्र~$A$ म्हणून परिभाषित करते की
$\Gamma \Sequent N : A$ ची \Log{tN3i} मध्ये
सिद्धता आहे. $\lambd$-पदांसाठीची
पद-रूपांतरणे सिद्धता ची न्यूनीकरण
आणि क्रमपरिवर्तन रूपांतरणे तंतोतंत प्रतिबिंबित करतात.
त्याचा महत्त्वाचा परिणाम म्हणजे संदर्भ~$\Gamma$ च्या
सापेक्ष $\lambd$-पदाचा प्रकार अनुरूप रूपांतरणानंतरही
तोच राहतो. या गुणधर्माला \emph{प्रकार-संरक्षण} म्हणतात.

\begin{prop}
संदर्भ~$\Gamma$ च्या सापेक्ष $N$ चा प्रकार~$A$
असेल आणि $N \redone N'$ असेल, तर~$\Gamma$ च्याच
सापेक्ष $N'$ चा प्रकारही~$A$ असतो.
\end{prop}

\begin{proof}
  ~$N$ वर आणि $\Gamma \Sequent N : A$ च्या
  सिद्धता~$\delta$ च्या उंचीवर विगमन करून पुढील
  गोष्ट सहज सिद्ध करता येते:~$M$ हे~$N$ चे उपपद
  असेल, तर~$\delta$ मध्ये $\Gamma' \Sequent M : B$
  वर संपणारी उप-सिद्धता~$\delta_1$ असते.
  $M$ हा रेडेक्स असेल, तर~$\delta_1$ वर रूपांतरण
  लागू होते. ~$\delta_1$ ला न्यूनीकरण लावल्यावर
  $\Gamma' \Sequent M' : B$ ची सिद्धता~$\delta_1'$
  मिळते. \Log{tN3i} च्या रूपांतरणांची तपासणी केली,
  तर~$M'$ हे~$M$ चे संकुचनफल आहे. ~$\delta$ मध्ये
  $\delta_1$ च्या जागी~$\delta_1'$ ठेवल्यावर
  $\Gamma \Sequent N' : A$ ची सिद्धता~$\delta'$
  मिळते; येथे~$N$ मधील उपपद~$M$ च्या जागी~$M'$
  ठेवल्यावर~$N'$ मिळते.

  उदाहरणार्थ, \Intro{\land} नंतर
  \Elim{\land} यासाठीचे न्यूनीकरण रूपांतरण घ्या:
  \[
  \bottomAlignProof
  \AxiomC{}
  \Deduce$\Gamma \fCenter N_1 : B_1$
  \AxiomC{}
  \Deduce$\Gamma \fCenter N_2 : B_2$
  \RightLabel{\Intro{\land}}
  \BinaryInf$\Gamma \fCenter \pair{N_1}{N_2} : B_1 \land B_2$
  \RightLabel{\Elim{\land}[i]}
  \UnaryInf$\Gamma \fCenter \proj{i}{\pair{N_1}{N_2}} : B_i$
  \DisplayProof
  \quad\redone\quad
  \bottomAlignProof
  \AxiomC{}
  \Deduce$\Gamma \fCenter N_i : B_i$
  \DisplayProof
  \]
  उजवीकडील सिद्धता चे सिद्धता-पद, म्हणजे~$\delta_1'$
  हे~$N_i$ आहे. डावीकडील~($\delta_1$) सिद्धता
  च्या $\proj{i}{\pair{N_1}{N_2}}$ या सिद्धता-पदावर
  $\beta$-परिवर्तन लावल्यावर मिळणारे हेच फल आहे.

  किंवा \Intro{\lif} नंतर \Elim{\lif} येणारे
  अनुमान आणि त्यावर लागू होणारे न्यूनीकरण रूपांतरण घ्या:
  \begin{multline*}
  \bottomAlignProof
  \AxiomC{$\delta_2$}
  \Deduce$x: B_1, \Gamma \fCenter N : B_2$
  \RightLabel{\Intro{\lif}}
  \UnaryInf$\Gamma \fCenter \lambd[\typeof{x}{B_1}][N] : B_1 \lif B_2$
  \AxiomC{}
  \RightLabel{$\delta_3$}
  \Deduce$\Gamma \fCenter M : B_1$
  \RightLabel{\Elim{\lif}}
  \BinaryInf$\Gamma \fCenter (\lambd[\typeof{x}{B_1}][N] M) : B_2$
  \DisplayProof
  \quad\redone\\
  \bottomAlignProof
  \AxiomC{}
  \RightLabel{$\delta_3$}
  \Deduce$\Gamma \fCenter M : B_1$
  \RightLabel{$\Subst{\delta_2}{\delta_3}{x:B_1}$}
  \Deduce$\Gamma \fCenter \Subst{N}{M}{x} : B_2$
  \DisplayProof
  \end{multline*}
  पुन्हा उजवीकडील सिद्धता चे सिद्धता-पद हे
  डावीकडील सिद्धता च्या सिद्धता-पदाचे
  $\beta$-न्यूनीकरण केल्यावर मिळणारेच फल आहे.
  अर्थात, ~$\delta_2$ च्या उंचीवर विगमन करून पुढील
  गोष्ट काळजीपूर्वक तपासावी लागेल:~$\delta_2$
  मधील $\Gamma' \Sequent x:B_1$ या रूपातील सर्व
  स्वयंसिद्धकांच्या जागी $\Gamma \Sequent M:B_1$
  ची सिद्धता~$\delta_3$ ठेवल्यास खरोखर
  $\Gamma \Sequent \Subst{N}{M}{x} : B_2$
  ची सिद्धता~$\Subst{\delta_2}{\delta_3}{x:B_1}$
  मिळते.
\end{proof}
हे प्रतिस्थापन करताना बद्ध चलांना आधी सुरक्षित नवी
नावे द्यायची आणि संबंधित खुणांनी बांधलेली स्वयंसिद्धकेच
बदलायची. येथे संदर्भ वाढवता येणारी \Log{tN3ip}
पद्धत वापरतो. \Log{tN2ip} मध्ये वापरलेल्या गृहीतकांचा
अचूक संच नोंदतो; शाखा किंवा न वापरलेले आदान गळल्यास
तो संच कमी होऊ शकतो, म्हणून तिथे तोच अचूक संदर्भ
राहतो असा दावा नाही.

सिद्धता आणि $\lambd$-पदे यांतील निकटच्या
अनुरूपतेचा आणखी एक परिणाम आहे. संदर्भ~$\Gamma$
च्या सापेक्ष~$A$ या प्रकाराचे $\lambd$-पद~$N$
मध्ये रेडेक्स~$M$ असेल (म्हणजे ते सामान्य रूपात
नसेल), तर $\Gamma \Sequent N : A$ ची अनुरूप
\Log{tN3i} सिद्धता~$\delta$ यात
$\Gamma' \Sequent M: B$ वर संपणारी अशी
उप-सिद्धता असते की तिच्यावर रूपांतरण लागू होते.
संकुचनफलाने~$M$ बदलून $N \redone N'$ होत असेल,
तर~$\delta$ ला न्यूनीकरण लावून
$\Gamma \Sequent N' : A$ ची सिद्धता मिळते.
म्हणून सामान्य रूपात नसलेले प्रत्येक $\lambd$-पद
दुसऱ्या पदात अनुरूप रूपांतरणाने न्यूनीत होते. या गुणधर्माला
\emph{प्रगती} म्हणतात.
येथील ‘सामान्य’ याचा अर्थ घोषित रूपांतरण लागू
नसलेले पद असा आहे. खुले सामान्य पद हे मूल्याच्या
नेहमीच्या संगणकीय रूपात असतेच, किंवा प्रत्येक टप्प्यावर
फक्त मुख्य बीटा-पायरी लागू होतेच, असा दावा नाही.

प्रगती आणि प्रकार-संरक्षण मिळून $\lambd$-पदांचे
घोषित पद-न्यूनीकरण ``अडकत'' नाही याची खात्री देतात:
एखादे पद योग्य प्रकाराचे असेल आणि सामान्य रूपात
नसेल, तर ते दुसऱ्या योग्य प्रकाराच्या पदात न्यूनीत
होते (आणि त्या पदाचा प्रकारही तोच असतो).














\section{सामान्यीकरण}\label{int:pty:nor:sec}

या विभागात योग्य प्रकाराच्या सिद्धता-पदांना सामान्य रूप
मिळते हे दाखवू. यासाठी आधी सिद्ध केलेले नैसर्गिक
निगमनाचे सामान्यरूपीकरण प्रमेय आणि विधाने-हे-प्रकार
अनुरूपता वापरू. येथे पूर्ण सामान्यरूपीकरणासाठी
न्यूनीकरणासह आवश्यक क्रमपरिवर्तन आणि असत्याच्या
सरलीकरणाचे अनुरूप पद-रूपांतरणही धरतो; फक्त तीन
मुख्य $\beta$-नियम पुरेसे असल्याचा दावा नाही.

प्रथम पदांची गुंतागुंत मोजणारी काही फलने परिभाषित करू.
सूत्रांची \emph{लांबी}~$\len{A}$ पुढीलप्रमाणे:
\begin{align*}
  \len{p} &= 0\\
  \len{A \land B} &= \len{A} + \len{B} + 1\\
  \len{A \lor B} &= \len{A} + \len{B} + 1 \\
  \len{A \lif B} &= \len{A} + \len{B} + 1.
\end{align*}
मुख्य $\beta$-रेडेक्स~$M$ ची गुंतागुंत त्याच्या \emph{कट-कोटी}ने,
म्हणजे~$\cutrank{M}$ ने मोजतो:
\begin{align*}
  \cutrank{(\lambd[\typeof{x}{A}][\typeof{N}{B}])Q} &=
  \len{A} + \len{B} + 1 \\
  \cutrank{\ande{i}{\andi{\typeof{M}{A}}{\typeof{N}{B}}}} &=
  \len{A} + \len{B} + 1 \\
  \cutrank{\ore{\inj[A_{3-i}]{i}{\typeof{M}{A_i}}}
    {\typeof{x_1}{A_1}}{\typeof{N_1}{C}}
    {\typeof{x_2}{A_2}}{\typeof{N_2}{C}}} &=
  \len{A_1} + \len{A_2} + 1
\end{align*}
या तीन मुख्य बीटा-प्रकारांसाठी सिद्धता-पदाची
गुंतागुंत त्यातील सर्वाधिक गुंतागुंतीच्या रेडेक्सने मोजतो;
असे रेडेक्स नसतील, तर ती~$0$ असते. हा मापदंड
क्रमपरिवर्तन किंवा असत्याच्या सरलीकरणाचे मोजमाप नाही:
\[
  \maxrank{M} = \max \{\cutrank{N} | N \text{ हे उपपद आहे } M \text{ आणि रेडेक्स आहे}\}
\]
अणुरूप प्रकार आणि~$\lfalse$ यांची लांबी~$0$ धरतो.
वरील महत्तम मूल्य घेण्यासाठीचा संच रिकामा असेल, तर व्याख्येने
$\maxrank{M}=0$ घेतो.
बेरीज-प्रकाराच्या पदात~$\inj[A_{3-i}]{i}{M}$ वापरतो:
$M$ चा प्रकार~$A_i$ आहे आणि वरची प्रकार-खूण
निवडला न गेलेला दुसरा प्रकार सांगते. त्यामुळे पूर्ण
प्रकार~$A_1\lor A_2$ मिळतो, प्रकार-तक्त्याप्रमाणे.

\begin{lem}\label{int:pty:nor:lem:subst}
  $\Subst{M}{\typeof{N}{A}}{\typeof{x}{A}}$ हा मुख्य बीटा-रेडेक्स असेल
  आणि $M \not\ident x$ असेल, तर पुढीलपैकी एक शक्यता असते:
  \begin{enumerate}
  \item $M$ स्वतः रेडेक्स आहे; किंवा
  \item $M$ चे रूप $\ande{i}{x}$ आणि $N$ चे रूप
    $\andi{P_1}{P_2}$ आहे;
  \item $M$ चे रूप $\ore{x}{x_1}{P_1}{x_2}{P_2}$ आणि $N$ चे
    रूप $\inj{i}{Q}$ आहे;
  \item $M$ चे रूप $x Q$ आणि $N$ चे रूप
    $\lambd[x][P]$ आहे.
  \end{enumerate}
  पहिल्या बाबतीत $\cutrank{\Subst{M}{N}{x}} = \cutrank{M}$;
  इतर बाबतींत $\cutrank{\Subst{M}{N}{x}} = \len{A}$.
\end{lem}
\begin{proof}
  ~$M$ वर विगमनाने सिद्ध करू.
  \begin{enumerate}
  \item $M$ हे एकच चर $y$ असेल आणि $y \not\ident x$ असेल,
    तर $\Subst{y}{N}{x}$ हे $y$ आहे; म्हणून ते रेडेक्स नाही.
  \item $M$ चे रूप $\andi{N_1}{N_2}$, किंवा $\lambd[x][N]$,
    किंवा $\ori{i}{A}{N}$ असेल, तर
    $\Subst{M}{\typeof{N}{A}}{\typeof{x}{A}}$ चेही तेच
    रूप असते; म्हणून ते रेडेक्स नसते.
  \item $M$ चे रूप $\ande{i}{P}$ असेल, तर दोन बाबी पाहू.
    \begin{enumerate}
      \item $P$ चे रूप $\andi{P_1}{P_2}$ असेल, तर
        $M \ident \ande{i}{\andi{P_1}{P_2}}$ हा रेडेक्स आहे;
        आणि स्पष्टच,
        \[
        \Subst{M}{N}{x} \ident
        \ande{i}{\andi{\Subst{P_1}{N}{x}}{\Subst{P_2}{N}{x}}}
        \]
        हाही रेडेक्स आहे. दोघांच्या कट-कोटी समान आहेत.
      \item $P$ एकच चर असेल, तर प्रतिस्थापनाने रेडेक्स
        मिळण्यासाठी तो~$x$ असलाच पाहिजे, आणि~$N$ चे रूप
        $\andi{P_1}{P_2}$ असले पाहिजे. आता
        \[
        \Subst{M}{N}{x} \ident \Subst{\ande{i}{x}}{\andi{P_1}{P_2}}{x},
        \]
        विचारात घ्या. हे $\ande{i}{\andi{P_1}{P_2}}$ आहे.
        त्याची कट-कोटी $\cutrank{\ande{i}{\andi{P_1}{P_2}}}$
        म्हणजेच $\len{A}$ आहे.
    \end{enumerate}
  \end{enumerate}
  $\ore{N}{x_1}{N_1}{x_2}{N_2}$ आणि $P Q$ या बाबीही
  अशाच प्रकारच्या आहेत.
\end{proof}
येथील प्रतिस्थापन मुक्त चलांना बद्ध होऊ न देता करायचे;
आवश्यकतेनुसार बद्ध चलांना आधी नवी नावे द्यायची.
नवीन रेडेक्सच्या वर्गीकरणात बदललेल्या पदाच्या
मुळाशी असलेला रेडेक्स विचारात घेतला आहे, सर्व
आतील रेडेक्स नव्हेत. $\abort{B}{P}$ चे मूळ
तीन मुख्य बीटा-रेडेक्सपैकी कोणतेही नसते.

\begin{lem}
  $M$ या मुख्य बीटा-रेडेक्सचे $M'$ मध्ये संकुचन होत असेल आणि~$M$ च्या प्रत्येक
  योग्य रेडेक्स उपपद~$N$ साठी
  $\cutrank{M} > \cutrank{N}$ असेल, तर
  $\cutrank{M} > \maxrank{M'}$.
\end{lem}

\begin{proof}
  बाबींनुसार सिद्ध करू.
  \begin{enumerate}
  \item $M$ चे रूप $\ande{i}{\andi{M_1}{M_2}}$ असेल, तर
    $M'$ हे $M_i$ आहे. $M_i$ चे कोणतेही उपपद हे~$M$ चे
    योग्य उपपदही असल्यामुळे दावा सिद्ध होतो.
  \item $M$ चे रूप
    $(\lambd[\typeof{x}{A}][N])\typeof{Q}{A}$ असेल,
    तर $M'$ हे $\Subst{N}{\typeof{Q}{A}}{\typeof{x}{A}}$
    आहे. $M'$ मधील रेडेक्स घ्या. एकतर तो~$Q$ च्या बसवलेल्या प्रतीत आहे; अशा वेळी
    त्याची कट-कोटी गृहीतकामुळे मूळ~$M$ पेक्षा कमी आहे.
    अन्यथा~$N$ मध्ये त्याला
    अनुरूप व समान कट-कोटीचा रेडेक्स आहे; गृहीतकामुळे ती
    कट-कोटी $\cutrank{M}$ पेक्षा कमी आहे. अन्यथा त्याची
    कट-कोटी $\len{A}$ आहे; व्याख्येनुसार ती
    $\cutrank{(\lambd[x^{A}][N])Q}$ पेक्षा कमी आहे.
  \item $M$ चे रूप
    \[
    \ore{\inj[A_{3-i}]{i}{\typeof{N}{A_i}}}
        {\typeof{x_1}{A_1}}{\typeof{N_1}{C}}
        {\typeof{x_2}{A_2}}{\typeof{N_2}{C}},
    \]
    असेल, तर $M' \ident \Subst{N_i}{N}{\typeof{x_i}{A_i}}$.
    ~$M'$ मधील रेडेक्स घ्या. एकतर तो बसवलेल्या~$N$ च्या प्रतीत आहे; अशा वेळी
    गृहीतकामुळे त्याची कट-कोटी मूळ~$M$ पेक्षा कमी आहे.
    अन्यथा~$N_i$ मध्ये त्याला
    अनुरूप व समान कट-कोटीचा रेडेक्स आहे; गृहीतकामुळे ती
    कट-कोटी $\cutrank{M}$ पेक्षा कमी आहे. अन्यथा त्याची
    कट-कोटी $\len{A_i}$ आहे; व्याख्येनुसार ती
    $\cutrank{\ore{\inj[A_{3-i}]{i}{\typeof{N}{A_i}}}
      {\typeof{x_1}{A_1}}{\typeof{N_1}{C}}
      {\typeof{x_2}{A_2}}{\typeof{N_2}{C}}}$ पेक्षा कमी आहे.
  \end{enumerate}
\end{proof}

\begin{thm}
  प्रत्येक योग्य प्रकाराचे सिद्धता-पद पूर्ण सामान्यरूपीकरणाच्या
  अनुरूप पद-रूपांतरणांनी सामान्य रूपात न्यूनीत करता येते;
  प्रत्येक अंतःप्रज्ञावादी निष्पत्ती सामान्य निष्पत्तीमध्ये न्यूनीत करता येते.
\end{thm}

\begin{proof}
  नैसर्गिक निगमनाच्या सामान्यरूपीकरण प्रकरणातील
  \ref{pt:nor:thm:thm:normalization-N2i} वापरू.
  योग्य प्रकाराच्या~$M$ साठी त्याची प्रकार-सिद्धता घ्या.
  अनावश्यक संदर्भ-गृहीतके काढून आणि बद्ध चलांची नावे
  सुरक्षित ठेवून तिच्याशी अनुरूप \Log{N2i}-सिद्धता
  मिळते. तिचे सामान्यरूप करा.

  प्रत्येक रूपांतरणाला पुन्हा सिद्धता-पदे द्या. मुख्य
  रूपांतरणे म्हणजे येथे दिलेले फलन, गुणाकार आणि बेरीज
  प्रकारांचे बीटा-नियम. वियोग-निरसनापुढील निरसन आत
  हलवणे हे संबंधित क्रमपरिवर्तन आहे; त्यात बद्ध चल
  आत हलवलेल्या पदांच्या मुक्त चलांपासून नवे ठेवायचे.
  \FalseInt{} च्या सरलीकरणासाठीही त्या प्रकरणातील
  रूपांतरणाची अनुरूप पद-आवृत्ती घ्यायची. यामुळे मूळ
  पदापासून सामान्य सिद्धतेच्या पदापर्यंत सांत क्रम मिळतो.
  अंतिम पदात लागू होणारे रूपांतरण असते, तर अंतिम
  सिद्धता सामान्य नसती. आवश्यकतेनुसार उरलेला संदर्भ
  पुन्हा वाढवून \Log{tN3ip} मध्ये मूळ प्रकार जपता येतो.

  म्हणून दोन्ही दावे मिळतात. वरचा रेडेक्स-मापदंड
  एकटाच संपूर्ण सिद्धता देत नाही: एका उपपदाचे संकुचन
  बाहेरच्या संदर्भात नवा रेडेक्स उघडू शकते; विशेषतः
  वियोग-निरसनाच्या निष्कर्षाचा प्रकार त्याच्या वियोगाच्या
  प्रकारापेक्षा मोठा असू शकतो. म्हणून नुसती महत्तम
  कट-कोटीच्या रेडेक्सांची संख्या घटते असे गृहीत धरलेले नाही.
\end{proof}

प्रकारयुक्त $\lambd$-पदांचे सामान्यीकरण होते याचा अर्थ
अशा प्रत्येक पदाला \emph{सामान्य रूप मिळते}. $\lambd$-पदांचे
सामान्य रूपापर्यंत $\beta$-न्यूनीकरण करणे म्हणजे संगणन चालवणे,
आणि सामान्य रूप म्हणजे त्याचे मूल्य, असे मानल्यास वरच्या पूर्ण सामान्यरूपीकरणाच्या पद्धतीने प्रत्येक
योग्य प्रकाराचे पद अखेरीस सामान्य रूपात पोहोचते.
शिवाय प्रकार-संरक्षणामुळे त्या मूल्याचा प्रकार योग्य असतो.
उदाहरणार्थ,~$N$ चा प्रकार $A \lif B$ असेल
(म्हणजे ते प्रकार~$A$ च्या आर्ग्युमेंटवर चालणारे आणि
प्रकार~$B$ चे मूल्य देणारे फलन असेल) आणि प्रकार~$A$
च्या~$M$ या पदावर (निवेशावर) ते लावले, तर $(NM)$ चे
$N'$ या पदात (निर्गमात) न्यूनीकरण होते आणि त्या निर्गमाचा
प्रकार~$B$ असेल याची हमी मिळते.

मुख्य बीटा-न्यूनीकरणासाठी प्रकारयुक्त $\lambd$-कलनाचा
आणखी एक गुणधर्म खरा आहे: बीटा-पायऱ्या कोणत्याही क्रमाने
लावल्या तरी योग्य प्रकाराच्या पदातून अनंत क्रम मिळत नाही.
याला \emph{प्रबळ सामान्यीकरण} म्हणतात. हा अधिक बलवान
परिणाम येथे स्वतंत्रपणे सिद्ध केलेला नाही; वरचे प्रमेय
एका पूर्ण सामान्यरूपीकरण-क्रमाचे अस्तित्व देते.
खालील संगमशीलतेचा विचार मुख्य बीटा-न्यूनीकरणाचा आहे;
सर्व वाढीव क्रमपरिवर्तनांचे स्वतंत्र संगमशीलता-प्रमाणपत्र नाही.

संगणन चालवण्याच्या वेगवेगळ्या पद्धतींनी \emph{तेच}
मूल्य मिळते हे कसे कळते? आतापर्यंत (प्रकार-संरक्षणामुळे)
एवढेच माहीत आहे की मिळालेल्या सर्व मूल्यांचा प्रकार एकच
असतो. प्रकारयुक्त $\lambd$-कलनाला आणखी एक महत्त्वाचा
गुणधर्म आहे: ते \emph{संगमशील}, म्हणजे त्याला
\emph{चर्च–रॉसर गुणधर्म} आहे. $N \red N_1$ आणि
$N \red N_2$ असतील, तर एखादे पद~$N'$ असे असते की
$N_1 \red N'$ आणि $N_2 \red N'$. लक्षात ठेवा की
$N \red N'$ याचा अर्थ $\redone$-न्यूनीकरणाच्या
शून्य किंवा अधिक पायऱ्यांनी~$N'$ मिळते.
$N_1$ सामान्य रूपात असेल, तर $N_1 \red N'$ असणारे
एकमेव पद~$N'$ म्हणजे~$N_1$ स्वतःच (शून्य
$\redone$-न्यूनीकरण पायऱ्यांनी). त्यामुळे~$N_1$ आणि
$N_2$ दोन्ही सामान्य रूपात असतील, तर $N_1 = N' = N_2$.
दुसऱ्या शब्दांत, $\redone$ प्रबळ सामान्यीकारी असल्यामुळे
पद न्यूनीत करण्याची प्रत्येक पद्धत अखेरीस सामान्य रूप देते.
$\red$ ला चर्च–रॉसर गुणधर्म असल्यामुळे कोणतीही दोन
सामान्य रूपे समान असतात. म्हणून प्रकारयुक्त
$\lambd$-कलनातील संगणन नेहमी थांबते आणि ते कसेही
चालवले तरी मिळणारे मूल्य (सामान्य रूप) तेच असते.

संबंध संगमशील आहे हे सिद्ध करणे सहसा अवघड असते.
तथापि, पुढील दुर्बल अट सहज सिद्ध करता येते:

\begin{prop}[दुर्बल चर्च–रॉसर गुणधर्म]
$N \redone N_1$ आणि $N \redone N_2$ असतील, तर एखादे
पद~$N'$ असे असते की $N_1 \red N'$ आणि $N_2 \red N'$.
\end{prop}

\begin{proof}
  $N$ मधील रेडेक्स~$M_1$ किंवा~$M_2$ त्याच्या संकुचनफलाने
  बदलल्यास अनुक्रमे~$N_1$ किंवा~$N_2$ मिळतात.
  $M_1$ आणि~$M_2$ हे~$N$ चे उपपद असल्याने त्यांचे
  अंशतः छेदन होऊ शकत नाही. म्हणजे एक रेडेक्स दुसऱ्याचे
  उपपद असेल, किंवा ते~$N$ मधील वेगवेगळ्या न छेदणाऱ्या
  स्थानी असतील. दुसरी बाब घ्या:~$N$ चे रूप $N(M_1,M_2)$
  आहे. $M_1$ च्या जागी त्याचे संकुचनफल~$M_1'$ ठेवल्यास
  $N(M_1',M_2)$ हे~$N_1$ मिळते; $M_2$ च्या जागी त्याचे
  संकुचनफल~$M_2'$ ठेवल्यास $N(M_1,M_2')$ हे~$N_2$
  मिळते. दोघांचेही $N(M_1',M_2')$ मध्ये न्यूनीकरण होते.
  दुसऱ्या बाबतीत~$M_2$ हे~$M_1$ चे उपपद आहे असे माना
  (उलट बाब सममित आहे), आणि~$M_1$ कोणत्या प्रकारचा
  रेडेक्स आहे यानुसार बाबी वेगळ्या करा. उदाहरणार्थ,
  $M_1 \ident \proj{1}{\pair{O_1}{O_2}}$ असेल, तर त्याचे
  $O_1$ मध्ये न्यूनीकरण होते. $M_2$ हे~$O_1$ चे उपपद
  असेल, तर~$M_2$ चे परिवर्तन केल्यावर~$O_1'$ मिळते.
  म्हणून आधी~$M_1$ आणि नंतर~$M_2$ चे परिवर्तन केल्यावर
  $O_1'$ मिळते. पण आधी~$M_2$ चे परिवर्तन केल्यावर
  $\proj{1}{\pair{O_1'}{O_2}}$ मिळते, आणि त्याचेही
  $O_1'$ मध्ये न्यूनीकरण होते.
\end{proof}
दोन्ही निवडी तोच रेडेक्स असतील, तर निष्कर्ष तत्काळ
मिळतो. वरील गुणाकार उदाहरणात~$M_2$ हे~$O_2$ मध्ये
असेल, तर दोन्ही क्रम~$O_1$ येथे मिळतात. फलनाच्या
बीटा-रेडेक्समध्ये आतील पायरी देहात असेल, तर
प्रतिस्थापनानंतर ती अनुरूप स्थानी लावायची; ती आदानात
असेल, तर बसवलेल्या सर्व प्रतींत लावायची, किंवा आदान
वापरले नसेल, तर शून्य पायऱ्या घ्यायच्या. बेरीज-प्रकारात
निवडलेल्या शाखेसाठी हाच युक्तिवाद लागू होतो; टाकलेल्या
शाखेतील पायरीची गरज उरत नाही. सुरक्षित प्रतिस्थापन
आणि बद्ध चलांची सुसंगत नावे वापरल्याने सामायिक पद
मिळते. पदांची समानता येथे बद्ध-चल नावबदलापर्यंत आहे.

\begin{lem}[न्यूमनची प्रमेयिका]
$\redone$ प्रबळ सामान्यीकारी असेल आणि त्याला दुर्बल
चर्च–रॉसर गुणधर्म असेल, तर त्याचा शून्य किंवा अधिक पायऱ्यांचा संबंध~$\red$ संगमशील आहे.
\end{lem}

\begin{proof}
  $\redone$ प्रबळ सामान्यीकारी असल्यामुळे प्रत्येक पूर्ण
  न्यूनीकरण-क्रम सामान्य रूपात संपतो. सामान्य रूपे एकमेव असणे आणि
  $\red$ संगमशील असणे या सममूल्य अटी आहेत. म्हणून एखाद्या
  पद~$N$ ला~$N'$ आणि~$N''$ ही दोन वेगळी सामान्य रूपे
  आहेत असे माना; ती पुढील न्यूनीकरण-क्रमांनी मिळतात:
  \begin{align*}
  N & \redone N_1' \redone N_2' \redone \dots \redone N_k' = N' \text{ आणि}\\
  N & \redone N_1'' \redone N_2'' \redone \dots \redone N_l'' = N''
  \end{align*}
  (न्यूनीकरण-क्रमाची लांबी $>0$ असलीच पाहिजे, अन्यथा
  $N$ आधीच सामान्य रूपात असेल.)

  $N_1' = N_1''$ असेल, तर त्याला दोन वेगळी सामान्य रूपे
  ($N'$ आणि $N''$) आहेत. $N_1' \neq N_1''$ असेल, तर
  दुर्बल चर्च–रॉसर गुणधर्मामुळे एखादे~$N'''$ असे आहे की
  $N_1' \red N'''$ आणि $N_1'' \red N'''$.
  $N' \neq N''$ असल्यामुळे $N''' \neq N'$ किंवा
  $N''' \neq N''$ यांपैकी किमान एक अट खरी आहे.
  निर्णायक मुद्दा असा: प्रबळ सामान्यीकरणामुळे या सामायिक
  पदाला एक सामान्य रूप मिळते. ते मूळ दोन वेगळ्या सामान्य
  रूपांपैकी किमान एकापेक्षा वेगळे असते. त्यामुळे~$N_1'$
  किंवा~$N_1''$ यांपैकी एकाला दोन वेगळी
  सामान्य रूपे आहेत. आपण दाखवले की~$N$ ला दोन वेगळी
  सामान्य रूपे असतील, तर एखादे~$N_1$ असे आहे की
  $N \redone N_1$ आणि~$N_1$ लाही दोन वेगळी सामान्य रूपे
  आहेत. ही प्रक्रिया पुढे चालू ठेवून
  $N \redone N_1 \redone N_2 \redone \dots$ असा क्रम
  मिळेल; पण $\redone$ प्रबळ सामान्यीकारी असल्यामुळे
  हे अशक्य आहे.
\end{proof}



\chapter{क्रमवर्ती कलन}\label{pt:seq:chap}










\section{परिचय}\label{pt:seq:int:sec}

क्रमवर्ती कलन या सिद्धता पद्धती गेरहार्ड
गेन्ट्झेनने १९३० च्या दशकात मांडल्या. व्यावहारिक
सिद्ध करणे पद्धती म्हणून त्यांचा उद्देश तुलनेने कमी
होता. त्याऐवजी, सिद्धता ची संभाव्य रचना आणि
त्यांचे गुणधर्म यांविषयी काही निष्कर्ष सिद्ध करण्यासाठी
गेन्ट्झेनने त्यांचा वापर केला. नेमके असेच प्रश्न
सिद्धता-उपपत्तीच्या अभ्यासाचे विषय आहेत.

नावाप्रमाणेच क्रमवर्ती कलने सूत्रे वर
नव्हे तर \emph{क्रमवर्तींवर} काम करतात. क्रमवर्ती
म्हणजे सूत्रांच्या क्रमिकांची किंवा बहुसंचांची
जोडी. गेन्ट्झेनच्या मूळ पद्धतींमध्ये (\Log{LK} आणि
\Log{LJ}) क्रमिका वापरल्या होत्या; आता
सिद्धता-उपपत्तीचे अभ्यासक बहुसंच वापरणे पसंत करतात.
बहुसंच हा संचासारखा असतो, पण त्यात एखादा
घटक एकाहून अधिक वेळा येऊ शकतो.
संचाप्रमाणेच, मात्र, घटक चा क्रम महत्त्वाचा
नसतो. म्हणून $\{A, A, B\}$ आणि
$\{A, B, A\}$ ही एकाच बहुसंचाची रूपे आहेत.
पण हा बहुसंच $\{A,
B\}$ आणि $\{A, A, B, B\}$ यांहून वेगळा आहे,
कारण त्यांतील $A$ आणि $B$ यांच्या प्रतींच्या
संख्या वेगवेगळ्या आहेत.

सिद्धता-उपपत्तीमध्ये सूत्रांचे बहुसंच
परंपरेने मोठ्या ग्रीक अक्षरांनी लिहितात. म्हणून
$\Gamma$, $\Delta$, $\Pi$, $\Lambda$, किंवा
$\Theta$ लिहिल्यास आपण ज्या तर्कशास्त्रात काम
करतो त्यातील सूत्रांचे बहुसंच अभिप्रेत असतात.
तसेच परंपरेने, $\tuple{\Gamma, \Delta}$ असे लिहिण्याऐवजी
सिद्धता-उपपत्तीचे अभ्यासक दोन घटक क्रमवर्ती-चिन्हाने
वेगळे करतात; येथे आपण~$\Sequent$ वापरतो.

\begin{defn}[क्रमवर्ती]
\emph{क्रमवर्ती} म्हणजे पुढील रूपाची अभिव्यक्ती:
\[
\Gamma \Sequent \Delta
\]
येथे $\Gamma$ आणि $\Delta$ हे सूत्रांचे
सान्त (कदाचित रिकामे) बहुसंच असतात. $\Gamma$ ला
\emph{पूर्वांग}, तर $\Delta$ ला \emph{उत्तरांग} म्हणतात.
\end{defn}

$G$ हे~$\Gamma$ मधील सूत्रांचे संयोजन
मानू (रिकामे संयोजन~$\ltrue$); आणि $D$ हे~$\Delta$
मधील सूत्रांचे वियोजन मानू (रिकामे वियोजन
~$\lfalse$). मग क्रमवर्ती~$\Gamma \Sequent \Delta$
काही रचना~$\Struct{M}$ मध्ये सत्य असेल
तेव्हाच $G \lif D$ सत्य असते. अभिजात
तर्कशास्त्रात याचा अर्थ असा:~$\Gamma$ मधील
सूत्रांपैकी किमान एक असत्य असते किंवा~$\Delta$
मधील सूत्रांपैकी किमान एक सत्य असते.

$\Gamma$ हा सूत्रांचा बहुसंच असेल
तर~$\Gamma$ मध्ये~$A$ जोडल्यावर मिळणारा बहुसंच $\Gamma, A$
(किंवा $A, \Gamma$) असा लिहितो. $\Delta$ हाही
सूत्रांचा बहुसंच असेल तर $\Gamma, \Delta$
म्हणजे त्या दोन बहुसंचांचे बहुसंच-संयोजन: $\Gamma$
मध्ये $A$ ची \emph{प्रतींची संख्या}~n (म्हणजे~$A$
च्या $n$ प्रती) आणि $\Delta$ मध्ये ती~$m$ असेल,
तर $\Gamma, \Delta$ मध्ये $A$ ची प्रतींची संख्या
~$n+m$ असते. क्रमवर्तीच्या एका बाजूचा बहुसंच रिकामा
असेल तर बहुतेक वेळा ती जागा कोरी ठेवतो. उदाहरणार्थ,
$\Sequent A$ या क्रमवर्तीचे पूर्वांग रिकामे असून
उत्तरांगात~$A$ ची नेमकी~$1$ प्रत असते.

क्रमवर्ती कलनातील सिद्धता म्हणजे
क्रमवर्तींचा वृक्ष: त्याच्या सर्वांत वरच्या (किंवा
प्रारंभीच्या) क्रमवर्ती संबंधित पद्धतीतील स्वयंसिद्धके
असतात; आणि एखादी क्रमवर्ती एक किंवा दोन क्रमवर्तींच्या
खाली असेल, तर ती पद्धतीतील निगमन नियमाने त्यांच्यापासून
योग्य रीतीने मिळाली पाहिजे. प्रथम आपण अभिजात
तर्कशास्त्राच्या दोन भिन्न क्रमवर्ती पद्धती, $\Log{G1c}$
आणि~$\Log{G3c}$, पाहू. त्यांची स्वयंसिद्धके आणि नियम
तक्ते~\ref{pt:seq:int:tab:G1c} आणि~\ref{pt:seq:int:tab:G3c}
मध्ये दिले आहेत.








\begin{table}
\begin{tabular}{@{}c@{}c@{}}
\multicolumn{2}{@{}c@{}}{स्वयंसिद्धके}\\[2ex]
$A \Sequent A$ & $\lfalse \Sequent \quad$
\\[2ex]
\multicolumn{2}{@{}c@{}}{संरचनात्मक नियम}\\[2ex]
\Axiom$ \Gamma \fCenter \Delta $
\RightLabel{\LeftR{\Weakening}}
\UnaryInf$ A, \Gamma \fCenter \Delta$
\DisplayProof
&
\Axiom$ \Gamma \fCenter \Delta$
\RightLabel{\RightR{\Weakening}}
\UnaryInf$ \Gamma \fCenter \Delta, A$
\DisplayProof
\\[4ex]
\Axiom$ A, A, \Gamma \fCenter \Delta $
\RightLabel{\LeftR{\Contraction}}
\UnaryInf$ A, \Gamma \fCenter \Delta$
\DisplayProof
&
\Axiom$ \Gamma \fCenter \Delta, A, A$
\RightLabel{\RightR{\Contraction}}
\UnaryInf$ \Gamma \fCenter \Delta, A$
\DisplayProof
\\[4ex]
\multicolumn{2}{@{}c@{}}{तार्किक नियम}\\[2ex]
\Axiom$ \Gamma \fCenter \Delta, A $
\RightLabel{\LeftR{\lnot}}
\UnaryInf$ \lnot A, \Gamma \fCenter \Delta$
\DisplayProof
&
\Axiom$A, \Gamma \fCenter \Delta$
\RightLabel{\RightR{\lnot}}
\UnaryInf$ \Gamma \fCenter \Delta, \lnot A $
\DisplayProof
\\[4ex]\begin{tabular}[t]{@{}l@{}}
\Axiom$ A, \Gamma \fCenter \Delta$
\RightLabel{\LeftR{\land}}
\UnaryInf$ A \land B, \Gamma \fCenter \Delta$
\DisplayProof
\\[3ex]
\Axiom$B, \Gamma \fCenter \Delta$
\RightLabel{\LeftR{\land}}
\UnaryInf$A \land B, \Gamma \fCenter \Delta$
\DisplayProof\\[3ex]
\end{tabular}
&
\Axiom$\Gamma \fCenter \Delta, A$
\Axiom$ \Gamma \fCenter \Delta, B$
\RightLabel{\RightR{\land}}
\BinaryInf$ \Gamma \fCenter \Delta, A \land B $
\DisplayProof
\\[4ex]\Axiom$A, \Gamma \fCenter \Delta$
\Axiom$B, \Gamma \fCenter \Delta$
\RightLabel{\LeftR{\lor}}
\BinaryInf$A \lor B, \Gamma \fCenter \Delta$
\DisplayProof
&
\begin{tabular}[t]{@{}r@{}}
\Axiom$\Gamma \fCenter \Delta, A$
\RightLabel{\RightR{\lor}}
\UnaryInf$ \Gamma \fCenter \Delta, A \lor B$
\DisplayProof
\\[3ex]
\Axiom$ \Gamma \fCenter \Delta, B$
\RightLabel{\RightR{\lor}}
\UnaryInf$ \Gamma \fCenter \Delta, A \lor B$
\DisplayProof\\[3ex]
\end{tabular}
\\[4ex]
\Axiom$ \Gamma \fCenter \Delta, A$
\Axiom$ B, \Gamma \fCenter \Delta$
\RightLabel{\LeftR{\lif}}
\BinaryInf$ A \lif B, \Gamma \fCenter \Delta$
\DisplayProof
&
\Axiom$ A, \Gamma \fCenter \Delta, B$
\RightLabel{\RightR{\lif}}
\UnaryInf$ \Gamma \fCenter \Delta, A \lif B $
\DisplayProof
\\[4ex]
\Axiom$ A(t), \Gamma \fCenter \Delta$
\RightLabel{\LeftR{\lforall}}
\UnaryInf$ \lforall[x][A(x)],\Gamma \fCenter \Delta$
\DisplayProof
&
\Axiom$ \Gamma \fCenter \Delta, A(a) $
\RightLabel{\RightR{\lforall}}
\UnaryInf$ \Gamma \fCenter \Delta, \lforall[x][A(x)]$
\DisplayProof
\\[4ex]
\Axiom$ A(a), \Gamma \fCenter \Delta $
\RightLabel{\LeftR{\lexists}}
\UnaryInf$ \lexists[x][A(x)], \Gamma \fCenter \Delta$
\DisplayProof
&
\Axiom$ \Gamma \fCenter \Delta, A(t) $
\RightLabel{\RightR{\lexists}}
\UnaryInf$ \Gamma \fCenter \Delta, \lexists[x][A(x)]$
\DisplayProof
\end{tabular}
\caption{\Log{G1c} चे नियम. $\RightR{\forall}$ आणि
$\LeftR{\exists}$ मध्ये स्थिरांक~$a$ हा
निष्कर्ष-क्रमवर्तीमध्ये येऊ नये. $\Gamma$ आणि
$\Delta$ हे सूत्रांचे बहुसंच आहेत.}
\label{pt:seq:int:tab:G1c}
\end{table}










\begin{table}
\begin{tabular}{@{}c@{}c@{}}
\multicolumn{2}{@{}c@{}}{स्वयंसिद्धके}\\[2ex]
$C, \Gamma \Sequent \Delta, C$ & $\lfalse, \Gamma \Sequent \Delta$
\\[2ex]
\multicolumn{2}{@{}c@{}}{तार्किक नियम}\\[2ex]
\Axiom$ \Gamma \fCenter \Delta, A $
\RightLabel{\LeftR{\lnot}}
\UnaryInf$ \lnot A, \Gamma \fCenter \Delta$
\DisplayProof
&
\Axiom$A, \Gamma \fCenter \Delta$
\RightLabel{\RightR{\lnot}}
\UnaryInf$ \Gamma \fCenter \Delta, \lnot A $
\DisplayProof
\\[4ex]
\Axiom$ A, B, \Gamma \fCenter \Delta$
\RightLabel{\LeftR{\land}}
\UnaryInf$ A \land B, \Gamma \fCenter \Delta$
\DisplayProof
&
\Axiom$\Gamma \fCenter \Delta, A$
\Axiom$ \Gamma \fCenter \Delta, B$
\RightLabel{\RightR{\land}}
\BinaryInf$ \Gamma \fCenter \Delta, A \land B $
\DisplayProof
\\[4ex]\Axiom$A, \Gamma \fCenter \Delta$
\Axiom$B, \Gamma \fCenter \Delta$
\RightLabel{\LeftR{\lor}}
\BinaryInf$A \lor B, \Gamma \fCenter \Delta$
\DisplayProof
&
\Axiom$\Gamma \fCenter \Delta, A, B$
\RightLabel{\RightR{\lor}}
\UnaryInf$ \Gamma \fCenter \Delta, A \lor B$
\DisplayProof
\\[4ex]
\Axiom$ \Gamma \fCenter \Delta, A$
\Axiom$ B, \Gamma \fCenter \Delta$
\RightLabel{\LeftR{\lif}}
\BinaryInf$ A \lif B, \Gamma \fCenter \Delta$
\DisplayProof
&
\Axiom$ A, \Gamma \fCenter \Delta, B$
\RightLabel{\RightR{\lif}}
\UnaryInf$ \Gamma \fCenter \Delta, A \lif B $
\DisplayProof
\\[4ex]
\Axiom$ A(t), \lforall[x][A(x)], \Gamma \fCenter \Delta$
\RightLabel{\LeftR{\lforall}}
\UnaryInf$ \lforall[x][A(x)],\Gamma \fCenter \Delta$
\DisplayProof
&
\Axiom$ \Gamma \fCenter \Delta, A(a) $
\RightLabel{\RightR{\lforall}}
\UnaryInf$ \Gamma \fCenter \Delta, \lforall[x][A(x)]$
\DisplayProof
\\[4ex]
\Axiom$ A(a), \Gamma \fCenter \Delta $
\RightLabel{\LeftR{\lexists}}
\UnaryInf$ \lexists[x][A(x)], \Gamma \fCenter \Delta$
\DisplayProof
&
\Axiom$ \Gamma \fCenter \Delta, \lexists[x][A(x)], A(t) $
\RightLabel{\RightR{\lexists}}
\UnaryInf$ \Gamma \fCenter \Delta, \lexists[x][A(x)]$
\DisplayProof
\end{tabular}
\caption{\Log{G3c} चे नियम. $\RightR{\forall}$ आणि
$\LeftR{\exists}$ मध्ये स्थिरांक~$a$ हा
निष्कर्ष-क्रमवर्तीमध्ये येऊ नये. $\Gamma$ आणि
$\Delta$ हे सूत्रांचे बहुसंच आहेत.
स्वयंसिद्धकांमध्ये सूत्र~$C$ अण्वीय असले पाहिजे.}
\label{pt:seq:int:tab:G3c}
\end{table}




या पद्धतींमध्ये सूक्ष्म फरक आहेत आणि त्यामुळे
त्यांचे सिद्धता-उपपत्तीय गुणधर्म वेगवेगळे आहेत.
\Log{G1c} ची स्वयंसिद्धके $A
\Sequent A$ अशी असतात; त्यात इतर कोणतेही
सूत्रांचालत नाहीत. \Log{G3c} ची स्वयंसिद्धके
$C, \Gamma \Sequent \Delta, C$ या रूपाची
असतात; त्यात~$\Gamma$ आणि~$\Delta$ मधील इतर
``बाजूची'' सूत्रे मनमानी असू शकतात, पण
दोन्ही बाजूंना येणारे~$C$ हे सूत्र
\emph{अण्वीय} असले पाहिजे. \Log{G1c} मध्ये
$\LeftR{\land}$ आणि $\RightR{\lor}$ नियमांच्या
प्रत्येकी दोन आवृत्त्या आहेत, तर \Log{G3c} मध्ये
प्रत्येकी एकच आहे. तसेच \Log{G1c} मध्ये
``आकुंचन'' (\LeftR{\Contraction}, \RightR{\Contraction})
आणि ``दुर्बलीकरण'' (\LeftR{\Weakening}, \RightR{\Weakening})
हे अतिरिक्त ``संरचनात्मक नियम'' आहेत.
यांशिवाय असत्याचे स्वयंसिद्धकही आहे: \Log{G1c} मध्ये
$\lfalse\Sequent\quad$, आणि \Log{G3c} मध्ये
$\lfalse,\Gamma\Sequent\Delta$. त्यामुळे समान
अणुरूप सूत्र दोन्ही बाजूंना असण्याचा प्रसंग हाच
स्वयंसिद्धकांचा एकमेव प्रसंग नाही.

\Log{G1c} आणि \Log{G3c} या दोन्ही
पद्धती अभिजात तर्कशास्त्रासाठी निर्दोष व संपूर्ण
आहेत. म्हणजे या पद्धतींमध्ये सिद्धतायोग्य असलेली
प्रत्येक क्रमवर्ती प्रत्येक संरचनांमध्ये सत्य
असते; आणि प्रत्येक संरचनांमध्ये सत्य असलेल्या
प्रत्येक क्रमवर्तीची सिद्धता असते. म्हणून एखाद्या
क्रमवर्तीची सिद्धता \emph{आहे का}, या प्रश्नाचे
उत्तर तर्कशास्त्राच्या इतर पद्धतींनीही (उदा.,
प्रतिमान-उपपत्तीने) मिळवता येते. आणि दोन्ही पद्धती
नेमक्या प्रत्येक संरचनांमध्ये सत्य असलेल्या
क्रमवर्ती सिद्ध करतात; म्हणून त्या विशेषतः
त्याच क्रमवर्ती सिद्ध करतात.

पण सिद्धता-उपपत्तीला एखादी गोष्ट
\emph{सिद्धतायोग्य आहे का} एवढाच नव्हे तर ती
\emph{कशी} सिद्ध होते हाही प्रश्न महत्त्वाचा असतो.
सिद्धता-उपपत्तीचे अभ्यासक असे प्रश्न विचारतात:
``एखादा विशिष्ट नियम नेहमी टाळता येईल का?'',
``नव्या क्रमवर्ती सिद्धतायोग्य न होता हा नियम पद्धतीत
जोडता येईल का?'', किंवा ``एका पद्धतीतील सिद्धता
दुसऱ्या पद्धतीतील सिद्धता मध्ये रूपांतरित करता
येतील का?'' अशा प्रश्नाचे उत्तर ``होय'' असेल, तर
त्या तथ्याची सिद्धता बहुधा सिद्धता च्या
गुंतागुंतीवर किंवा, समतुल्यपणे, त्या सिद्धता च्या रचनेवर
विगमन करून दिली जाते. त्यातून केवळ तथ्याचीच
सिद्धता मिळत नाही (उदा., \Log{G1c} ने एखादी
क्रमवर्ती सिद्ध केली तर तीच क्रमवर्ती \Log{G3c}
नेही सिद्ध होते), तर एक प्रक्रिया मिळते (उदा.,
\Log{G1c} मधील सिद्धता ना \Log{G3c} मधील
सिद्धता मध्ये रूपांतरित करण्याची).

सिद्धता-उपपत्तीतील एक मध्यवर्ती निष्कर्ष म्हणजे
\emph{कट-निरसन प्रमेय}. त्यानुसार पुढील कट नियम
\[
\Axiom$\Gamma \fCenter \Delta, A$
\Axiom$A, \Pi \fCenter \Lambda$
\RightLabel{\Cut}
\BinaryInf$\Gamma, \Pi \fCenter \Delta, \Lambda$
\DisplayProof
\]
क्रमवर्ती पद्धतींमध्ये जोडला तरी कोणत्या क्रमवर्ती
सिद्धतायोग्य आहेत यांत बदल होत नाही. याहूनही,
त्या निष्कर्षाची सिद्धता अशी पद्धत देते की
\Log{G1c} (किंवा \Log{G3c}) चे नियम आणि~\Cut{}
वापरून केलेल्या कोणत्याही सिद्धता चे रूपांतर
फक्त \Log{G1c} (किंवा \Log{G3c}) चेच नियम
वापरणाऱ्या सिद्धतेत करता येते. दुसऱ्या शब्दांत, ही
प्रक्रिया~\Cut{} वापरणाऱ्या सिद्धता मधून
हा नियम ``निरसित'' करते; म्हणून हे नाव.













\section{नियम आणि सिद्धता}\label{pt:seq:rp:sec}

आधी काही व्याख्या देऊ आणि काही संज्ञांचा अर्थ निश्चित करू.

उदाहरणार्थ, $A$ आणि~$B$ असलेल्या क्रमवर्तींपासून
$A \land B$ असलेल्या क्रमवर्ती सिद्ध करणे करू देणारे \Log{G3c} मधील
दोन नियम पाहा:
\[
\Axiom$ A, B, \Gamma \fCenter \Delta$
\RightLabel{\LeftR{\land}}
\UnaryInf$ A \land B, \Gamma \fCenter \Delta$
\DisplayProof
\qquad
\Axiom$\Gamma \fCenter \Delta, A$
\Axiom$ \Gamma \fCenter \Delta, B$
\RightLabel{\RightR{\land}}
\BinaryInf$ \Gamma \fCenter \Delta, A \land B $
\DisplayProof
\]
रेषेच्या वरच्या क्रमवर्तींना \emph{आधारविधाने}, तर खालच्या
क्रमवर्तीला \emph{निष्कर्ष} म्हणतात. प्रत्येक नियमातील
$\Gamma$ आणि $\Delta$ यांतील सूत्रे ना \emph{संदर्भ}
किंवा \emph{बाजूची सूत्रे} म्हणतात. निष्कर्षातील
संदर्भाबाहेरचे सूत्र हे \emph{मुख्य} सूत्र असते;
आणि आधारविधानांतील संदर्भाबाहेरच्या सूत्रे ना
\emph{सक्रिय} (किंवा \emph{सहायक}) सूत्रे म्हणतात.

परिमाणकांचे नियम थोडे अधिक गुंतागुंतीचे आहेत. उदाहरणार्थ,
$\lforall$ साठीचे \Log{G1c} मधील नियम असे आहेत:
\[\Axiom$ A(t), \Gamma \fCenter \Delta$
\RightLabel{\LeftR{\lforall}}
\UnaryInf$ \lforall[x][A(x)],\Gamma \fCenter \Delta$
\DisplayProof
\qquad
\Axiom$ \Gamma \fCenter \Delta, A(a) $
\RightLabel{\RightR{\lforall}}
\UnaryInf$ \Gamma \fCenter \Delta, \lforall[x][A(x)]$
\DisplayProof
\]
\LeftR{\lforall} नियमातील सक्रिय सूत्र हे~$A(t)$ आहे;
येथे $t$ हे बंद पद, म्हणजे चर नसलेले पद, आहे. एखादे
निगमन योग्य असते—म्हणजे ते \LeftR{\lforall} या
रूपरेषात्मक नियमाशी जुळते—जर एखादे सूत्र~$A$,
चर~$x$ आणि बंद पद~$t$ असे असतील की निष्कर्षातील मुख्य
सूत्र $\lforall[x][A]$ असेल आणि आधारविधानातील
सक्रिय सूत्र हे~$\Subst{A}{t}{x}$ असेल.
\RightR{\lforall} नियमही याचप्रमाणे कार्य करतो; मात्र
कोणत्याही $t$ पदाऐवजी आधारविधानातील सक्रिय सूत्र
$\Subst{A}{a}{x}$ असे असले पाहिजे, जेथे $a$ हा
स्थिरांक आहे. या अचराला \RightR{\lforall}
निगमनाचा \emph{आयगेन चर} म्हणतात. खालच्या क्रमवर्तीमध्ये
$a$ येत नसेल—म्हणजेच $\Gamma$, $\Delta$ किंवा~$A$
यांत $a$ नसेल—तरच हे निगमन योग्य असते. ही
\emph{आयगेन चराची अट} होय.

\Log{G1c} आणि \Log{G3c} यांसारख्या क्रमवर्ती पद्धतींत अशा
रूपरेषात्मक रीतीने दिलेल्या नियमांचा संग्रह आणि स्वयंसिद्धके
म्हणून ग्राह्य अशा, त्याच रीतीने दिलेल्या, काही क्रमवर्ती असतात.
अशा पद्धतीतील सिद्धता म्हणजे स्वयंसिद्धकांपासून निगमन नियमांनी
विगमनाने तयार होणारे क्रमवर्तींचे सान्त वृक्ष. प्रत्येक पान
स्वयंसिद्धक असते; आणि इतर प्रत्येक क्रमवर्ती तिच्या लगेच
वरच्या क्रमवर्तींपासून पद्धतीतील एका निगमन नियमाने मिळते.

\begin{defn}\label{pt:seq:rp:defn:proof}
क्रमवर्ती पद्धतीतील \emph{सिद्धता} म्हणजे खालीलप्रमाणे
विगमनाने बनवलेले क्रमवर्तींचे सान्त वृक्ष:
\begin{enumerate}
  \item\label{pt:seq:rp:defn:prf:axiom} प्रत्येक स्वयंसिद्धक ही सिद्धता असते.
  \item\label{pt:seq:rp:defn:prf:one-prem} $\pi_1$ ही $S_1$ या क्रमवर्तीवर
  संपणारी सिद्धता असेल आणि $S_1$ पासून एका आधारविधानाच्या
  $R$ नियमाने $S$ ही क्रमवर्ती मिळत असेल, तर
  \[
  \AxiomC{}
  \RightLabel{$\pi_1$}
  \DeduceC{$S_1$}
  \RightLabel{$R$}
  \UnaryInfC{$S$}
  \DisplayProof
  \]
  ही सिद्धता असते.
  \item\label{pt:seq:rp:defn:prf:two-prem} $\pi_1$ आणि $\pi_2$ या
  अनुक्रमे $S_1$ आणि~$S_2$ या क्रमवर्तींवर संपणाऱ्या सिद्धता
  असतील, आणि $S_1$ व $S_2$ यांच्यापासून दोन आधारविधानांच्या $R$ नियमाने
  $S$ मिळत असेल, तर
  \[
  \AxiomC{}
  \RightLabel{$\pi_1$}
  \DeduceC{$S_1$}
  \AxiomC{}
  \RightLabel{$\pi_2$}
  \DeduceC{$S_2$}
  \RightLabel{$R$}
  \BinaryInfC{$S$}
  \DisplayProof
  \]
  ही सिद्धता असते.
\end{enumerate}
\end{defn}

\ref{pt:seq:rp:defn:proof} मध्ये सिद्धता ची व्याख्या क्रमवर्तींचे वृक्ष म्हणून
दिली आहे. हे वृक्ष ``वरच्या दिशेने'' वाढतात असे आपण मानतो.
सर्वांत वरची (पानांवरील) स्थाने स्वयंसिद्धके असतात; त्यांना
\emph{आरंभीच्या क्रमवर्ती} म्हणतात. सर्वांत खालची क्रमवर्ती
\emph{अंतिम क्रमवर्ती} असते. $\pi$ ही $S$ अंतिम क्रमवर्ती
असलेली सिद्धता असेल तर $\pi \Proves S$ असेही लिहितो.
$S$ ची सिद्धता असेल तर $\Proves S$ असे लिहितो; काही वेळा
ज्या पद्धतीत सिद्धता आहे तीही $\Proves$ च्या डावीकडे दाखवतो,
उदा., $\Log{G1c} \Proves A, B \Sequent A \land
B$. आरंभीच्या क्रमवर्ती वगळता सिद्धता मधील प्रत्येक
क्रमवर्ती एखाद्या $R$ नियमाने झालेल्या निगमनाचा
\emph{निष्कर्ष} असते. सिद्धता मधील एखाद्या क्रमवर्तीच्या
लगेच वर असलेल्या एक किंवा दोन क्रमवर्तींना (वृक्षातील तिच्या
पालक-स्थानांना) त्या निगमनाची \emph{आधारविधाने} म्हणतात.

एखादी सिद्धता विगमनाने बनवताना वापरलेल्या सिद्धता ना तिच्या
\emph{उप-सिद्धता} म्हणतात. एखाद्या निगमनाच्या आधारविधानांवर
संपणाऱ्या उप-सिद्धता ना त्या निगमनाच्या \emph{लगेचच्या
उप-सिद्धता} म्हणतात.

अनेकदा सिद्धता ची गुंतागुंत एखाद्या नैसर्गिक संख्येने
मोजावी लागते. सिद्धता मधील क्रमवर्तींची संख्या मोजता येईल;
पण अनेकदा अधिक उपयोगी मोजमाप म्हणजे सिद्धता ची उंची:
अंतिम क्रमवर्तीपासून आरंभीच्या क्रमवर्तीपर्यंतच्या सर्वाधिक
निगमन-पायऱ्यांची संख्या. तिची विगमनात्मक व्याख्याही देता येते.

\begin{defn}
सिद्धता~$\pi$ ची \emph{उंची}~$\pheight{\pi}$ पुढीलप्रमाणे
विगमनाने ठरते.
\begin{enumerate}
  \item $\pi$ मध्ये एकच स्वयंसिद्धक असेल, तर $\pheight{\pi} = 0$.
  \item $\pi$ चा शेवट एका आधारविधानाच्या निगमनाने होत असेल आणि
  तिची लगेचची उप-सिद्धता~$\pi_1$ असेल, तर $\pheight{\pi} = \pheight{\pi_1} + 1$.
  \item $\pi$ चा शेवट दोन आधारविधानांच्या निगमनाने होत असेल आणि
  तिच्या लगेचच्या उप-सिद्धता~$\pi_1$ आणि~$\pi_2$ असतील, तर $\pheight{\pi} =
  \max(\pheight{\pi_1}, \pheight{\pi_2}) + 1$.
\end{enumerate}
$S$ या क्रमवर्तीची $\le n$ उंचीची सिद्धता असेल, तर $\Proves[n] S$ असे लिहितो.
\end{defn}

\begin{ex}
\Log{G3c} मध्ये $C, \Gamma \Sequent \Delta, C$ या रूपाची
प्रत्येक क्रमवर्ती स्वयंसिद्धक असते; येथे $C$ अण्वीय असले पाहिजे.
$D$ आणि $E$ ही अण्वीय वाक्ये असोत. मग $D, E \Sequent D$
आणि $D, E \Sequent E$ ही $C, \Gamma \Sequent \Delta, C$
या स्वयंसिद्धकाची उदाहरणे आहेत. उदाहरणार्थ, $C \ident E$,
$\Gamma = \{D\}$ आणि $\Delta = \emptyset$ घेतल्यास
$C, \Gamma \Sequent \Delta, C$ पासून नेमकी $D, E
\Sequent E$ ही क्रमवर्ती मिळते. (लक्षात घ्या, आपण \emph{बहुसंच} वापरतो;
म्हणून $D, E \Sequent E$ आणि $E, D \Sequent E$ या एकाच
क्रमवर्ती आहेत.)

$D, E \Sequent D$ ही \Log{G3c} च्या स्वयंसिद्धकाचे उदाहरण
असल्यामुळे ती स्वतःच \Log{G3c} मधील सिद्धता असते. त्याचप्रमाणे
$D, E \Sequent E$ हीही \ref{pt:seq:rp:defn:proof}\ref{pt:seq:rp:defn:prf:axiom}
नुसार तशीच असते. त्यांना अनुक्रमे $\pi_1$ आणि~$\pi_2$ म्हणू.
\ref{pt:seq:rp:defn:prf:two-prem} नुसार, या दोन सिद्धता घेऊन
दोन आधारविधानांच्या $\RightR{\land}$ नियमाने नवी सिद्धता
($\pi_3$) तयार करता येते:
\[
\Axiom$D, E \fCenter D$
\Axiom$D, E \fCenter E$
\RightLabel{\RightR{\land}}
\BinaryInf$D, E  \fCenter D \land E $
\DisplayProof
\]
$\pi_3$ पासून पुढे सिद्धता~$\pi_4$ मिळते:
\[
\Axiom$D, E \fCenter D$
\Axiom$D, E \fCenter E$
\RightLabel{\RightR{\land}}
\insertBetweenHyps{\hspace{5ex}}
\BinaryInf$D, E  \fCenter D \land E $
\LeftSubproofLabel{$\pi_3$}
\RightLabel{\LeftR{\land}}
\UnaryInf$D \land E  \fCenter D \land E$
\RightSubproofLabel{$\pi_4$}
\DisplayProof
\]
येथे एका आधारविधानाचा $\LeftR{\land}$ नियम वापरला आहे
(\ref{pt:seq:rp:defn:proof}\ref{pt:seq:rp:defn:prf:one-prem} नुसार). $\pi_4$ मधील
शेवटच्या निगमनाची लगेचची उप-सिद्धता ही~$\pi_3$ आहे.
$\RightR{\land}$ निगमनाच्या लगेचच्या उप-सिद्धता $\pi_1$
आणि~$\pi_2$ आहेत. $\pi_4$ च्या उप-सिद्धता मध्ये $\pi_1$,
$\pi_2$, $\pi_3$ आणि स्वतः $\pi_4$ यांचा समावेश होतो.
$\pheight{\pi_1} = \pheight{\pi_2} = 0$,
$\pheight{\pi_3} = 1$ आणि $\pheight{\pi_4} = 2$.
म्हणून कोणत्याही अण्वीय $D$ आणि~$E$ साठी
$\Log{G3c} \Proves[2] D \land E  \fCenter D
\land E$ हे दाखवले आहे.
\end{ex}













\section{सिद्धतांची उदाहरणे}\label{pt:seq:ex:sec}

क्रमवर्तींच्या सिद्धता तयार करताना बहुतेक
वेळा अंतिम क्रमवर्तीपासून वरच्या दिशेने काम करणे
सोयीचे ठरते: त्यातील एखादे सूत्र सक्रिय
सूत्र म्हणून निवडायचे, त्या क्रमवर्तीच्या वर
संबंधित नियमाचे पूर्वपक्ष लिहायचे, आणि स्वयंसिद्धके
मिळेपर्यंत ही प्रक्रिया पुन्हा करायची. उदाहरणार्थ,
पुढील क्रमवर्ती सिद्ध करायची आहे असे माना:
\[(C \land D) \lif E \Sequent C \lif (D \lif E)\]
$C \lif (D \lif E)$ विचारात घ्या: ते $A \lif B$
या रूपाचे असून क्रमवर्तीच्या उजवीकडे आहे. पूर्ण
क्रमवर्तीची \Log{G1c} च्या $\RightR{\lif}$ नियमाशी
जुळवणी करू:
\[\Axiom$ A, \Gamma \fCenter \Delta, B$
\RightLabel{\RightR{\lif}}
\UnaryInf$ \Gamma \fCenter \Delta, A \lif B $
\DisplayProof\]
यावरून $\Gamma = \{(C \land D) \lif E\}$,
$\Delta =
\emptyset$, $A \ident C$ आणि $B \ident D \lif E$
असे घ्यावे लागते. मग सिद्धतेचे शेवटचे निगमन असे असू शकेल:
\[\Axiom$ C, (C \land D) \lif E \fCenter D \lif E$
\RightLabel{\RightR{\lif}}
\UnaryInf$(C \land D) \lif E \fCenter C \lif (D \lif E)$
\DisplayProof\]
हीच कल्पना पुन्हा लावून $D \lif E$ ला सक्रिय
सूत्र मानल्यास पुढील मिळते:
\[\Axiom$ D, C, (C \land D) \lif E \fCenter E$
\RightLabel{\RightR{\lif}}
\UnaryInf$C, (C \land D) \lif E \fCenter D \lif E$
\RightLabel{\RightR{\lif}}
\UnaryInf$(C \land D) \lif E \fCenter C \lif (D \lif E)$
\DisplayProof\]
आता $(C \land D) \lif E$ वर लक्ष केंद्रित करून
\Log{G1c} चा $\LeftR{\lif}$ नियम वापरावा लागेल:
\[\Axiom$ \Gamma \fCenter \Delta, A$
\Axiom$ B, \Gamma \fCenter \Delta$
\RightLabel{\LeftR{\lif}}
\BinaryInf$ A \lif B, \Gamma \fCenter \Delta$
\DisplayProof\]
त्यातून पुढील मिळते:
\[\Axiom$D, C \fCenter E, C \land D$
\Axiom$D, C, E \fCenter E$
\RightLabel{\LeftR{\lif}}
\BinaryInf$ D, C, (C \land D) \lif E \fCenter E$
\RightLabel{\RightR{\lif}}
\UnaryInf$C, (C \land D) \lif E \fCenter D \lif E$
\RightLabel{\RightR{\lif}}
\UnaryInf$(C \land D) \lif E \fCenter C \lif (D \lif E)$
\DisplayProof\]
वरच्या डाव्या क्रमवर्तीतील उजवीकडचे सूत्र
$C \land D$ प्रमुख मानून $\RightR{\land}$ नियम
लावल्यास पुढील मिळते:
\[
\Axiom$D, C \fCenter E, C$
\Axiom$D, C \fCenter E, D$
\RightLabel{\RightR{\land}}
\BinaryInf$D, C \fCenter E, C \land D$
\Axiom$D, C, E \fCenter E$
\RightLabel{\LeftR{\lif}}
\BinaryInf$ D, C, (C \land D) \lif E \fCenter E$
\RightLabel{\RightR{\lif}}
\UnaryInf$C, (C \land D) \lif E \fCenter D \lif E$
\RightLabel{\RightR{\lif}}
\UnaryInf$(C \land D) \lif E \fCenter C \lif (D \lif E)$
\DisplayProof
\]
आतापर्यंत सगळे ठीक आहे. आपण वापरलेले नियम
\Log{G1c} आणि~\Log{G3c} मध्ये सारखे असल्याने हे
सर्व \Log{G3c} मध्येही चालते. या सिद्धता च्या
अपूर्ण भागातील वरच्या सर्व क्रमवर्तींच्या डावीकडे
आणि उजवीकडे तेच सूत्र येते; पण त्या अद्याप
अगदी स्वयंसिद्धके झालेल्या नाहीत.

\Log{G1c} मध्ये वरच्या क्रमवर्तींच्या
सिद्धता $A \Sequent A$ या रूपाच्या स्वयंसिद्धकांपासून
द्याव्या लागतात. दुर्बलीकरण नियमांनी ते सहज शक्य
आहे. उदाहरणार्थ, सर्वांत डावीकडची वरची क्रमवर्ती
$D, C \Sequent
E, C$ पुढीलप्रमाणे सिद्ध करता येते:
\[
\Axiom$C \fCenter C$
\RightLabel{\LeftR{\Weakening}}
\UnaryInf$D, C \fCenter C$
\RightLabel{\RightR{\Weakening}}
\UnaryInf$D, C \fCenter E, C$
\]
एकत्रितपणे, \Log{G1c} मधील सिद्धता अशी दिसते:
\[
\Axiom$C \fCenter C$
\RightLabel{\LeftR{\Weakening}}
\UnaryInf$D, C \fCenter C$
\RightLabel{\RightR{\Weakening}}
\UnaryInf$D, C \fCenter E, C$
\Axiom$D \fCenter D$
\RightLabel{\LeftR{\Weakening}}
\UnaryInf$D, C \fCenter D$
\RightLabel{\RightR{\Weakening}}
\UnaryInf$D, C \fCenter E, D$
\RightLabel{\RightR{\land}}
\BinaryInf$D, C \fCenter E, C \land D$
\Axiom$E \fCenter E$
\RightLabel{\LeftR{\Weakening}}
\UnaryInf$C, E \fCenter E$
\RightLabel{\LeftR{\Weakening}}
\UnaryInf$D, C, E \fCenter E$
\RightLabel{\LeftR{\lif}}
\BinaryInf$ D, C, (C \land D) \lif E \fCenter E$
\RightLabel{\RightR{\lif}}
\UnaryInf$C, (C \land D) \lif E \fCenter D \lif E$
\RightLabel{\RightR{\lif}}
\UnaryInf$(C \land D) \lif E \fCenter C \lif (D \lif E)$
\DisplayProof
\]
या सिद्धतेची रचना (विशेषतः तिची उंची) सूत्रे
$C$, $D$, आणि~$E$ काय आहेत यावर अवलंबून नाही.
वरील सिद्धता मध्ये या चिन्हांच्या जागी कोणतीही
सूत्रे घातली तरी मिळणारी रचना~\Log{G1c}
मधील सिद्धता राहते. \Log{G1c} मधील ही
सिद्धता \emph{रूपरेषात्मक} आहे.

\Log{G3c} मध्ये स्वयंसिद्धके
$A, \Gamma \Sequent
\Delta, A$ या रूपाची क्रमवर्ती असतात; त्यात~$A$
अण्वीय असले पाहिजे. म्हणून $D, C \Sequent
E, C$ ही \Log{G3c} ची स्वयंसिद्धकासारखी दिसत
असली ($\Gamma = \{D\}$; $\Delta =
\{E\}$), तरी~$C$ प्रत्यक्षात अण्वीय असेल तेव्हाच
ती \emph{खरोखर} स्वयंसिद्धक असते. $C$, $D$, आणि $E$ ही अण्वीय सूत्रे
असतील, तर हा सिद्धता-वृक्ष \Log{G3c} मध्ये पूर्ण
सिद्धता असतो. ही पुरेशी अट आहे; इतर बाजूच्या सूत्रामुळे
एखादी वरची क्रमवर्ती स्वयंसिद्धक असण्याची शक्यताही आहे.

काटेकोरपणे पाहता, आपण अजून
\[\Log{G3c} \Proves (C \land D) \lif E \fCenter C \lif (D \lif
E),\] हे दाखवलेले नाही. तरी व्यवहारात एवढेच
करणे आपल्याला पुरेसे आहे (आणि शक्यही आहे).
कारण $A, \Gamma \Sequent \Delta, A$ या रूपाची
प्रत्येक क्रमवर्ती \Log{G3c} मध्ये सिद्धतायोग्य
असते, जरी~$A$ स्वतः अण्वीय नसले तरी.
हे~$A$ च्या खोली~$\depth{A}$ वर विगमन करून
दाखवता येते.

\begin{defn}\label{pt:seq:ex:defn:depth}
सूत्राची \emph{खोली}~$\depth{A}$ पुढीलप्रमाणे
परिभाषित करतो:
\begin{enumerate}
  \item $A$ अण्वीय असेल तर $\depth{A} = 0$.
  \item $A \ident \lnot B$, $A \ident \lforall[x][B]$,
  किंवा $A \ident \lexists[x][B]$ असेल तर
  $\depth{A} = \depth{B} + 1$.
  \item $A \ident (B \land C)$, $(B \lor C)$, किंवा $(B \lif
  C)$ असेल तर $\depth{A} = \max(\depth{B},\depth{C}) + 1$.
\end{enumerate}
\end{defn}

कोणत्याही पद~$t$ साठी $\depth{A(x)} =
\depth{A(t)}$ सिद्ध करणे कठीण नाही. आपल्याला
महत्त्वाचे हे की तार्किक नियमात सक्रिय विघटक सूत्रांपेक्षा प्रमुख
सूत्राची खोली नेहमी अधिक असते. संरचनात्मक
नियमांसाठी असा दावा नाही; तसेच संख्यापकाच्या काही
\Log{G3c} नियमांत जपलेली प्रमुख आस्थिती स्वतः त्याच
खोलीची असते. उदाहरणार्थ:
\begin{align*}
\depth{P(x)} = \depth{P(a)} & = 0,\\
\depth{\lforall[x][P(x)]} & = 1,\\
\depth{(\lforall[x][P(x)] \lif P(a))} & = \max(1,0) + 1 = 2.
\end{align*}
याचा वापर करून~$\depth{A}$ वर विगमनाने काही
निष्कर्ष सिद्ध करता येतात; पुढील निष्कर्ष त्यांपैकी एक.

\begin{prop}\label{pt:seq:ex:prop:G1c-axioms}
कोणत्याही~$A$ साठी, $A \Sequent A$ ची अशी
$\Log{G1c}$-सिद्धता आहे की तिच्यातील सर्व
स्वयंसिद्धके अण्वीय आहेत.
\end{prop}

\begin{proof}
  $\depth{A}$ वर विगमन करू. $\depth{A} = 0$
  असेल तर~$A$ अण्वीय आहे आणि $A \Sequent A$
  हीच अपेक्षित रूपाची \Log{G1c} मधील सिद्धता आहे.
  $\depth{A} > 0$ असेल तर~$A$ कमी खोलीच्या
  सूत्रांपासून बनते, उदा., $A \ident \lnot B$,
  $A \ident
  (B \land C)$, किंवा $A \ident \lforall[x][A(x)]$.
  विगमनात्मक गृहीतक असे: $\depth{D} < \depth{A}$
  असलेल्या प्रत्येक~$D$ साठी अण्वीय स्वयंसिद्धकांपासून
  $\Log{G1c} \Proves D \Sequent D$.

$A \ident \lnot B$ असेल, तर
  $\depth{B} < \depth{A}$; विगमनात्मक
  गृहीतकानुसार अण्वीय स्वयंसिद्धकांपासून
  $B \Sequent B$ ची \Log{G1c} मधील
  सिद्धता~$\pi_1$ आहे. तिचा पुढीलप्रमाणे
  $\lnot B \Sequent \lnot B$ च्या सिद्धता
  पर्यंत विस्तार करता येतो:
\[
  \AxiomC{}
  \RightLabel{$\pi_1$}
  \Deduce$B \fCenter B$
  \RightLabel{\LeftR{\lnot}}
  \UnaryInf$\lnot B, B \fCenter $
  \RightLabel{\RightR{\lnot}}
  \UnaryInf$\lnot B \fCenter \lnot B$
  \DisplayProof
\]

$A \ident (B \land C)$ असेल, तर
  $\depth{B} < \depth{A}$ आणि
  $\depth{C} < \depth{A}$. विगमनात्मक गृहीतकानुसार
  अण्वीय स्वयंसिद्धकांपासून $B
  \Sequent B$ आणि $C \Sequent C$ यांच्या
  \Log{G1c} मधील सिद्धता~$\pi_1$ आणि~$\pi_2$
  आहेत. त्यांचा पुढीलप्रमाणे
  $B \land C \Sequent B \land C$ च्या
  सिद्धता पर्यंत विस्तार करता येतो:
\[
  \AxiomC{}
  \RightLabel{$\pi_1$}
  \Deduce$B \fCenter B$
  \RightLabel{\LeftR{\Weakening}}
  \UnaryInf$B, C \fCenter B$
  \RightLabel{\LeftR{\land}}
  \UnaryInf$B \land C, C\fCenter B$
  \RightLabel{\LeftR{\land}}
  \UnaryInf$B \land C, B \land C \fCenter B$
  \AxiomC{}
  \RightLabel{$\pi_2$}
  \Deduce$C \fCenter C$
  \RightLabel{\LeftR{\Weakening}}
  \UnaryInf$B, C \fCenter C$
  \RightLabel{\LeftR{\land}}
  \UnaryInf$B \land C, C \fCenter C$
  \RightLabel{\LeftR{\land}}
  \UnaryInf$B \land C, B \land C \fCenter C$
  \RightLabel{\RightR{\land}}
  \BinaryInf$B \land C, B \land C \fCenter B \land C$
  \RightLabel{\LeftR{\Contraction}}
  \UnaryInf$B \land C \fCenter B \land C$
  \DisplayProof
  \]

आता $A \ident \lforall[x][B(x)]$ असे माना.
  $\lforall[x][B(x)]$ मध्ये न येणारा
  स्थिरांक~$c$ घ्या. मग $\depth{B(c)} =
  \depth{B(x)} < \depth{A}$; विगमनात्मक गृहीतकानुसार
  अण्वीय स्वयंसिद्धकांपासून $B(c)
  \Sequent B(c)$ ची \Log{G1c} मधील
  सिद्धता~$\pi_1$ आहे. तिचा पुढीलप्रमाणे
  $\lforall[x][B(x)] \Sequent \lforall[x][B(x)]$
  च्या सिद्धता पर्यंत विस्तार करता येतो:
\[
  \AxiomC{}
  \RightLabel{$\pi_1$}
  \Deduce$B(c) \fCenter B(c)$
  \RightLabel{\LeftR{\lforall}}
  \UnaryInf$\lforall[x][B(x)] \fCenter B(c)$
  \RightLabel{\RightR{\lforall}}
  \UnaryInf$\lforall[x][B(x)] \fCenter \lforall[x][B(x)]$
  \DisplayProof
\]

उरलेले प्रसंग प्रश्न म्हणून सोडले आहेत.
\end{proof}

\begin{prob}
  $A \ident (B \lor C)$, $A \ident (B \lif C)$
  आणि $A \ident \lexists[x][B(x)]$ या प्रसंगांचे
  काम करून \ref{pt:seq:ex:prop:G1c-axioms} ची सिद्धता पूर्ण करा.
\end{prob}

आपल्या कामासाठी $A \ident \lnot B$ या
प्रसंगात पुढील सिद्धता ही वापरता आली असती:
\[
\AxiomC{}
\RightLabel{$\pi_1$}
\Deduce$B \fCenter B$
\RightLabel{\RightR{\lnot}}
\UnaryInf$\fCenter B, \lnot B$
\RightLabel{\LeftR{\lnot}}
\UnaryInf$\lnot B \fCenter \lnot B$
\DisplayProof
\]
मात्र काही क्रमवर्ती पद्धतींमध्ये हे चालले नसते.
अंतःप्रज्ञावादी क्रमवर्ती कलन~\Log{G1i} हे
\Log{G1c} सारखेच आहे; फरक इतकाच की प्रत्येक
क्रमवर्तीच्या उजव्या बाजूस जास्तीत जास्त एकच
सूत्र ठेवता येते. $\Delta = \emptyset$
असेल, तर पहिली सिद्धता ही~\Log{G1i} मध्येही
सिद्धता असेल; परंतु दुसरी नसेल, कारण तिच्या
एका क्रमवर्तीच्या उजवीकडे $B$ आणि~$\lnot B$
दोन्ही आहेत.

$A \ident \lforall[x][B(x)]$ असताना
\Log{G1c} मध्येही नियमांचा दुसरा क्रम वापरता
आला नसता. पुढील रचना \emph{सिद्धता नाही}:
\[
  \AxiomC{}
  \RightLabel{$\pi_1$}
  \Deduce$B(c) \fCenter B(c)$
  \RightLabel{\RightR{\lforall}}
  \UnaryInf$B(c) \fCenter \lforall[x][B(x)]$
  \RightLabel{\LeftR{\lforall}}
  \UnaryInf$\lforall[x][B(x)] \fCenter \lforall[x][B(x)]$
  \DisplayProof
\]
कारण त्यात~$\RightR{\lforall}$ ची आयगेन चराची
अट मोडली जाते.
येथे~$x$ हे~$B(x)$ मध्ये खरोखर मुक्त आहे असे मानतो,
म्हणून~$c$ डावीकडील~$B(c)$ मध्ये येतो. संख्यापक
निष्फळ असेल आणि~$c$ त्या सूत्रात येतच नसेल, तर
दाखवलेल्या आयगेन-चर अटीचे उल्लंघन होत नाही.

\begin{prop}\label{pt:seq:ex:prop:G3c-axioms}
$A, \Gamma
\Sequent \Delta, A$ या रूपाची कोणतीही क्रमवर्ती
(~$A$ \emph{अण्वीय} असण्याची आवश्यकता नाही)
\Log{G3c} मध्ये सिद्धतायोग्य आहे.
\end{prop}

\begin{proof}
ही सिद्धता \ref{pt:seq:ex:prop:G1c-axioms} च्या सिद्धतेसारखीच
आहे, पण विगमनात्मक गृहीतकाबाबत काळजी घ्यावी
लागते. $n =
\depth{A} = 0$ हा प्रसंग पुन्हा सोपा आहे.
$n=0$ असेल, तर~$A$ अण्वीय आहे, आणि $A,
\Gamma \Sequent \Delta, A$ हे~\Log{G3c}
चे स्वयंसिद्धक आहे.

आता $n>0$ असे माना. पुन्हा प्रसंग वेगळे करू.
विगमनात्मक गृहीतक असे: $\depth{D} < n$ असलेल्या
सर्व~$D$ साठी आणि सर्व~$\Pi$ व~$\Lambda$
साठी $\Log{G3c} \Proves D, \Pi \Sequent \Lambda, D$.
$A \ident (B \land C)$ हा प्रसंग पाहू. विगमनात्मक
गृहीतक दोनदा वापरू: एकदा $D = B$,
$\Pi = C, \Gamma$, आणि $\Lambda
= \Delta$ घेऊन $B, C, \Gamma \Sequent
\Delta, B$ ची सिद्धता~$\pi_1$ मिळते; आणि
दुसऱ्यांदा $D = C$, $\Pi = B, \Gamma$, आणि
$\Lambda = \Delta$ घेऊन $B, C, \Gamma \Sequent
\Delta, C$ ची सिद्धता~$\pi_2$ मिळते.
$B \land C, \Gamma \Sequent
\Delta, B \land C$ ची सिद्धता पुढीलप्रमाणे:
\[
\AxiomC{}
\RightLabel{$\pi_1$}
\Deduce$B, C, \Gamma \fCenter \Delta, B$
\RightLabel{\LeftR{\land}}
\UnaryInf$B \land C, \Gamma \fCenter \Delta, B$
\AxiomC{}
\RightLabel{$\pi_2$}
\Deduce$B, C, \Gamma \fCenter \Delta, C$
\RightLabel{\LeftR{\land}}
\UnaryInf$B \land C, \Gamma \fCenter \Delta, C$
\RightLabel{\RightR{\land}}
\BinaryInf$B \land C, \Gamma \fCenter \Delta, B \land C$
\DisplayProof
\]

$A \ident \lforall[x][B(x)]$ साठी
~$\lforall[x][B(x)]$, $\Gamma$ किंवा~$\Delta$
मध्ये न येणारा स्थिरांक निवडा.
$D = B(c)$, $\Pi = \lforall[x][B(x)], \Gamma$
आणि $\Lambda = \Delta$ घेऊन विगमनात्मक गृहीतक
लावा. त्यातून $B(c),
\lforall[x][B(x)], \Gamma \Sequent \Delta, B(c)$
ची सिद्धता~$\pi_1$ मिळते; तिच्यापासून:
\[
\AxiomC{}
\RightLabel{$\pi_1$}
\Deduce$B(c), \lforall[x][B(x)], \Gamma \fCenter \Delta, B(c)$
\RightLabel{\LeftR{\lforall}}
\UnaryInf$\lforall[x][B(x)], \Gamma \fCenter \Delta, B(c)$
\RightLabel{\RightR{\lforall}}
\UnaryInf$\lforall[x][B(x)], \Gamma \fCenter \Delta, \lforall[x][B(x)]$
\DisplayProof
\]
आपण~$c$ ची केलेली निवड~$\RightR{\lforall}$
ची आयगेन चराची अट पूर्ण करते.
\end{proof}

\begin{prob}
  $A \ident (B \lif C)$ आणि $A \ident
  \lexists[x][B(x)]$ या प्रसंगांचे काम करून
  \ref{pt:seq:ex:prop:G3c-axioms} ची सिद्धता पुढे न्या.
\end{prob}













\section{नियमांचा अर्थ}\label{pt:seq:irl:sec}

क्रमवर्तीचा अर्थ आठवा: संरचना~$\Struct{M}$ मध्ये
क्रमवर्ती~$\Gamma \Sequent \Delta$ सत्य असते, जर
किमान एक $A \in \Gamma$ असत्य असेल किंवा किमान
एक~$A \in \Delta$ सत्य असेल. \Log{G3c} चे
तार्किक नियम त्यांच्या प्रमुख सूत्रांच्या
सत्यतेच्या अटी नियमबद्ध करतात असे वाचता येते.
\RightR{\lif} विचारात घ्या. त्याचे प्रमुख सूत्र
$A \lif B$ आहे. $A$ असत्य असेल किंवा~$B$
सत्य असेल तेव्हा ते सत्य असते. म्हणून क्रमवर्ती
$\ \Sequent A \lif B$ ही $A \Sequent B$
सत्य असेल तेव्हाच सत्य असते. या क्रमवर्तींच्या
पूर्वांगात आणि उत्तरांगात संदर्भातील सूत्रे
$\Gamma$, $\Delta$ जोडल्यास \RightR{\lif}
नियमाचे नेमके निष्कर्ष व पूर्वपक्ष मिळतात.
\LeftR{\lif} साठीही स्थिती अशीच आहे. $A
\lif B$ हे $A$~सत्य आणि $B$~असत्य असेल तेव्हा
असत्य असते. म्हणून $A \lif B
\Sequent \ $ ही क्रमवर्ती $\ \Sequent A$ आणि
$B \Sequent \ $ या दोन्ही सत्य असतील तेव्हाच
सत्य असते. \LeftR{\lif} चा निष्कर्ष त्याच्या
पूर्वपक्षांच्या संयुक्त सत्यतेशी समतुल्य आहे. या
रचनेमुळे \Log{G3c} चे नियम निर्दोष असतात (सर्व
पूर्वपक्ष सत्य असतील तर निष्कर्षही सत्य असतो).
त्यांच्यामुळे पूर्णताही मिळते.

खरे तर, कोणत्याही सत्यता-फलनीय संयोजकाच्या
सत्यतेच्या अटी अशा क्रमवर्तींच्या संचाने मांडता येतात.
कारण अभिजात तर्कशास्त्रातील प्रत्येक सत्यता-फलन
संयुक्त सामान्य रूपातील सूत्र ने दाखवता येते:
अण्वीय आणि निषेधित अण्वीय सूत्रे यांच्या
वियोजनांचे संयोजन. उदाहरणार्थ, $A
\oplus B$ हे $A$ आणि $B$ यांपैकी नेमके एकच
सत्य असेल तेव्हाच सत्य मानू; यालाच ``अनन्य किंवा''
म्हणतात. मग $A \oplus B$ हे $(A \lor B) \land
(\lnot A \lor \lnot B)$ सत्य असेल तेव्हाच सत्य
असते; आणि $(\lnot A \lor B)
\land (A \lor \lnot B)$ सत्य असेल तेव्हा ते असत्य
असते. म्हणून $\oplus$ साठी क्रमवर्ती-कलनाचे नियम
पुढीलप्रमाणे देता येतील:
\begin{align*}
&
\Axiom$\Gamma \fCenter \Delta, A, B$
\Axiom$A, B, \Gamma \fCenter \Delta$
\RightLabel{\RightR{\oplus}}
\BinaryInf$\Gamma \fCenter \Delta, A \oplus B$
\DisplayProof
\\
& \Axiom$A, \Gamma \fCenter \Delta, B$
\Axiom$B, \Gamma \fCenter \Delta, A$
\RightLabel{\LeftR{\oplus}}
\BinaryInf$A \oplus B, \Gamma \fCenter \Delta$
\DisplayProof
\end{align*}
हे नियमही निर्दोष व पूर्ण असतील. त्यांचे निष्कर्ष सर्व
पूर्वपक्ष सत्य असतील तेव्हाच सत्य असतात.

\begin{prob}
$A \downarrow B$ हे $A$ किंवा~$B$ यांपैकी
एकही सत्य नसल्यास आणि तरच सत्य आहे असे माना
(``NOR''). त्याचे \Log{G3c} मधील नियम कोणते असतील?
($A \downarrow B$ सत्य असेल तेव्हाच एकाच वेळी
सत्य असणाऱ्या दोन क्रमवर्ती शोधा. त्यातून
\RightR{\downarrow} नियम मिळेल. त्यानंतर $A
\downarrow B$ असत्य असेल तेव्हाच सत्य असणारी एक
क्रमवर्ती शोधा. त्यातून \LeftR{\downarrow} नियम मिळेल.)
\end{prob}

\Log{G3c} मध्ये प्रत्येक संयोजक आणि संख्यक
यांचे नेमके दोन नियम आहेत: एक डावीकडचा आणि एक
उजवीकडचा. येथे~$\lfalse$ हे स्वतंत्र तार्किक स्थिरांक
आहे; त्यासाठी डावीकडील स्वयंसिद्धक असते, उजवीकडील
स्वतंत्र नियम नसतो. पण \Log{G1c} मध्ये दोन
\LeftR{\land} आणि दोन \RightR{\lor} नियम आहेत.
याचे कारण असे की \Log{G1c} चे प्रतिमान असलेल्या
गेन्ट्झेनच्या मूळ~\Log{LK} पद्धतीला एक महत्त्वाच्या
अभिजातेतर तर्कशास्त्रासाठी, म्हणजे अंतःप्रज्ञावादी
तर्कशास्त्रासाठी,~\Log{LJ} हा प्रकार होता.
अंतःप्रज्ञावादी तर्कशास्त्रासाठी निर्दोष व पूर्ण
पद्धत मिळवण्याकरिता~\Log{LJ} मधील प्रत्येक
क्रमवर्तीच्या उत्तरांगात जास्तीत जास्त एकच
सूत्र ठेवण्याचे बंधन आहे. त्यामुळे
\Log{G3c} चा \RightR{\lor} नियम वापरता येत
नाही, कारण त्याच्या पूर्वपक्षाच्या उत्तरांगात $A$
आणि~$B$ ही दोन्ही येतात. म्हणून गेन्ट्झेनने
\Log{LJ} साठी \RightR{\lor} नियमांची जोडी
वापरली, आणि तीच~\Log{LK} साठीही वापरली.
तदनुरूप, अंतःप्रज्ञावादी प्रकार~\Log{G1i} हा
प्रत्येक उत्तरांगात जास्तीत जास्त एकच सूत्र
ठेवण्याचे बंधन घातलेली~\Log{G1c} पद्धत आहे.
(\Log{G3c} चाही अंतःप्रज्ञावादी प्रकार आहे;
मात्र त्यातील काही नियम \Log{G3c} मधील संबंधित
नियमांहून वेगळे आहेत. उदाहरणार्थ, त्यातही
\RightR{\lor} नियमांची जोडी वापरतात.)

\Log{G1c} चे संरचनात्मक नियम निर्दोष आहेत
हे सहज दिसते. उदाहरणार्थ, $\Gamma \Sequent \Delta$
च्या डावीकडे असत्य सूत्र किंवा उजवीकडे सत्य
सूत्र असेल, तर $\Gamma \Sequent
\Delta, A$ मध्येही तेच असते; $A$ सत्य आहे की
असत्य, याने फरक पडत नाही. म्हणून
\RightR{\Weakening} चा पूर्वपक्ष सत्य असेल तर
निष्कर्षही सत्य असतो. येथे उलट दिशेचा दावा मात्र
सत्य नाही: निष्कर्ष $A$ सत्य असल्यामुळे सत्य असू
शकतो; तेव्हा पूर्वपक्ष सत्य असेल याची खात्री नाही.

संरचनात्मक नियम मुळात का हवेत? उदाहरणार्थ,
$\ \Sequent A \lor \lnot A$ सिद्ध करण्यासाठी
आकुंचन (\RightR{\Contraction},
\LeftR{\Contraction}) लागते. $A \Sequent A$
पासून सुरुवात केल्यास, आधी~\RightR{\lnot} वापरून
$\Sequent A, \lnot A$ मिळते. मग
\RightR{\lor} च्या दोन आवृत्त्या क्रमाने वापरता येतात:
एकदा~$\lnot A$ पासून $A \lor \lnot
A$ मिळवण्यासाठी आणि एकदा~$A$ पासून
$A \lor \lnot A$ मिळवण्यासाठी. त्यातून $\ \Sequent A \lor \lnot
A, A \lor \lnot A$ मिळते. म्हणून \Log{G1c}
मध्ये $\Sequent A \lor \lnot A$ मिळवण्यासाठी
सिद्धता मधील त्याच सूत्राच्या दोन
उपस्थितींचे ``आकुंचन'' करता आले पाहिजे:
\[
\Axiom$A \fCenter A$
\RightLabel{\RightR{\lnot}}
\UnaryInf$\fCenter A, \lnot A$
\RightLabel{\RightR{\lor}}
\UnaryInf$\fCenter A, A \lor \lnot A$
\RightLabel{\RightR{\lor}}
\UnaryInf$\fCenter A \lor \lnot A, A \lor \lnot A$
\RightLabel{\RightR{\Contraction}}
\UnaryInf$\fCenter A \lor \lnot A$
\DisplayProof
\]
खरे तर, दुर्बलीकरण किंवा आकुंचन नसताना \Log{G1c}
मध्ये सिद्ध न करता येणाऱ्या वैध क्रमवर्ती आहेत.
उदाहरणार्थ,~$A$ आणि~$B$ ही वेगळी अणुरूप
सूत्रे घ्या. दुर्बलीकरण व आकुंचन वगळलेल्या तक्त्यातील
\Log{G1c} नियमांनी क्रमवर्ती
$A \Sequent B \lif A$ ची सिद्धता
\RightR{\lif} नेच संपली पाहिजे (कारण फक्त
$B \lif A$ प्रमुख असू शकते); त्यासाठी
$B, A \Sequent A$ हा पूर्वपक्ष हवा. तो
$A \Sequent A$ या स्वयंसिद्धकापासून मिळवायला
\LeftR{\Weakening} लागते:
\[
\Axiom$A \fCenter A$
\RightLabel{\LeftR{\Weakening}}
\UnaryInf$B, A \fCenter A$
\RightLabel{\RightR{\lif}}
\UnaryInf$A \fCenter B \lif A$
\DisplayProof
\]

\Log{G3c} मध्ये आकुंचन आणि दुर्बलीकरण
हे नियम अजिबात लागत नाहीत. दुर्बलीकरण आधीच
स्वयंसिद्धकांत अंतर्भूत आहे. दोन पूर्वपक्ष असलेल्या
नियमांमधील संदर्भ~$\Gamma$, $\Delta$ यांच्या दोन
प्रती एकत्र करून आकुंचन बहुतेक ठिकाणी टाळता येते.
\Log{G1c} आणि \Log{G3c} या दोन्हींमध्ये असे
``संदर्भ सामायिक करणारे'' नियम आहेत. एकच पूर्वपक्ष
असलेल्या नियमांत आकुंचन फक्त \RightR{\lor},
\LeftR{\land}, \RightR{\lexists} आणि~\LeftR{\lforall}
यांसाठी लागते. \Log{G3c} मधील या नियमांच्या
आवृत्त्या आकुंचन पूर्णपणे अनावश्यक करतात.
\Log{G2c} ही आणखी एक पद्धत आहे (
\ref{pt:seq:irl:tab:G2c} पाहा); तिच्या दोन-पूर्वपक्षी नियमांना
स्वतंत्र संदर्भ असतात. \Log{G2c} मध्ये आकुंचन
\Log{G1c} पेक्षाही अधिक महत्त्वाचे आहे.








\begin{table}
\begin{tabular}{@{}c@{\hskip-5ex}c@{}}
\multicolumn{2}{@{}c@{}}{स्वयंसिद्धके}\\[2ex]
$A \Sequent A$ & $\lfalse \Sequent \quad$
\\[2ex]
\multicolumn{2}{@{}c@{}}{संरचनात्मक नियम}\\[2ex]
\Axiom$ \Gamma \fCenter \Delta $
\RightLabel{\LeftR{\Weakening}}
\UnaryInf$ A, \Gamma \fCenter \Delta$
\DisplayProof
&
\Axiom$ \Gamma \fCenter \Delta$
\RightLabel{\RightR{\Weakening}}
\UnaryInf$ \Gamma \fCenter \Delta, A$
\DisplayProof
\\[4ex]
\Axiom$ A, A, \Gamma \fCenter \Delta $
\RightLabel{\LeftR{\Contraction}}
\UnaryInf$ A, \Gamma \fCenter \Delta$
\DisplayProof
&
\Axiom$ \Gamma \fCenter \Delta, A, A$
\RightLabel{\RightR{\Contraction}}
\UnaryInf$ \Gamma \fCenter \Delta, A$
\DisplayProof
\\[4ex]
\multicolumn{2}{@{}c@{}}{तार्किक नियम}\\[2ex]
\Axiom$ \Gamma \fCenter \Delta, A $
\RightLabel{\LeftR{\lnot}}
\UnaryInf$ \lnot A, \Gamma \fCenter \Delta$
\DisplayProof
&
\Axiom$A, \Gamma \fCenter \Delta$
\RightLabel{\RightR{\lnot}}
\UnaryInf$ \Gamma \fCenter \Delta, \lnot A $
\DisplayProof
\\[4ex]\begin{tabular}[t]{@{}l@{}}
\Axiom$ A, \Gamma \fCenter \Delta$
\RightLabel{\LeftR{\land}}
\UnaryInf$ A \land B, \Gamma \fCenter \Delta$
\DisplayProof
\\[3ex]
\Axiom$B, \Gamma \fCenter \Delta$
\RightLabel{\LeftR{\land}}
\UnaryInf$A \land B, \Gamma \fCenter \Delta$
\DisplayProof\\[3ex]
\end{tabular}
&
\Axiom$\Gamma_1 \fCenter \Delta_1, A$
\Axiom$ \Gamma_2 \fCenter \Delta_2, B$
\RightLabel{\RightR{\land}}
\BinaryInf$ \Gamma_1, \Gamma_2 \fCenter \Delta_1, \Delta_2, A \land B $
\DisplayProof
\\[4ex]\Axiom$A, \Gamma_1 \fCenter \Delta_1$
\Axiom$B, \Gamma_2 \fCenter \Delta_2$
\RightLabel{\LeftR{\lor}}
\BinaryInf$A \lor B, \Gamma_1, \Gamma_2 \fCenter \Delta_1, \Delta_2$
\DisplayProof
&
\begin{tabular}[t]{@{}r@{}}
\Axiom$\Gamma \fCenter \Delta, A$
\RightLabel{\RightR{\lor}}
\UnaryInf$ \Gamma \fCenter \Delta, A \lor B$
\DisplayProof
\\[3ex]
\Axiom$ \Gamma \fCenter \Delta, B$
\RightLabel{\RightR{\lor}}
\UnaryInf$ \Gamma \fCenter \Delta, A \lor B$
\DisplayProof\\[3ex]
\end{tabular}
\\[4ex]
\Axiom$ \Gamma_1 \fCenter \Delta_1, A$
\Axiom$ B, \Gamma_2 \fCenter \Delta_2$
\RightLabel{\LeftR{\lif}}
\BinaryInf$ A \lif B, \Gamma_1, \Gamma_2 \fCenter \Delta_1, \Delta_2$
\DisplayProof
&
\Axiom$ A, \Gamma \fCenter \Delta, B$
\RightLabel{\RightR{\lif}}
\UnaryInf$ \Gamma \fCenter \Delta, A \lif B $
\DisplayProof
\\[4ex]
\Axiom$ A(t), \Gamma \fCenter \Delta$
\RightLabel{\LeftR{\lforall}}
\UnaryInf$ \lforall[x][A(x)],\Gamma \fCenter \Delta$
\DisplayProof
&
\Axiom$ \Gamma \fCenter \Delta, A(a) $
\RightLabel{\RightR{\lforall}}
\UnaryInf$ \Gamma \fCenter \Delta, \lforall[x][A(x)]$
\DisplayProof
\\[4ex]
\Axiom$ A(a), \Gamma \fCenter \Delta $
\RightLabel{\LeftR{\lexists}}
\UnaryInf$ \lexists[x][A(x)], \Gamma \fCenter \Delta$
\DisplayProof
&
\Axiom$ \Gamma \fCenter \Delta, A(t) $
\RightLabel{\RightR{\lexists}}
\UnaryInf$ \Gamma \fCenter \Delta, \lexists[x][A(x)]$
\DisplayProof
\end{tabular}
\caption{\Log{G2c} चे नियम. $\RightR{\forall}$ आणि
$\LeftR{\exists}$ मध्ये स्थिरांक~$a$ हा
निष्कर्ष-क्रमवर्तीमध्ये येऊ नये. $\Gamma$ आणि
$\Delta$ हे सूत्रांचे बहुसंच आहेत.}
\label{pt:seq:irl:tab:G2c}
\end{table}
















\section{नियमित सिद्धता आणि प्रतिस्थापन}\label{pt:seq:qua:sec}

संख्यकांचे नियम काही गुंतागुंत निर्माण करतात;
कारण~\LeftR{\lexists} आणि~\RightR{\lforall} या
नियमांना ``आयगेन चराच्या अटी'' लागू होतात. या
बहुतेक अडचणी त्या निगमनांतील स्थिरांक~$a$
पुन्हा कधीही वापरला जाणार नाही याची खात्री करून
सोडवता येतात. म्हणून पुढील व्याख्या देऊ.

\begin{defn}
\Log{G1c} किंवा \Log{G3c} मधील सिद्धता~$\pi$
\emph{नियमित} असते, जर
\begin{enumerate}
  \item कोणताही आयगेन चर~$a$ एकाहून अधिक
  \LeftR{\lexists} किंवा~\RightR{\lforall} निगमनांचा
  आयगेन चर म्हणून वापरला जात नसेल; आणि
  \item $a$ हा \LeftR{\lexists} किंवा~\RightR{\lforall}
  निगमनाचा आयगेन चर असेल, तर तो~$\pi$ मध्ये
  फक्त त्या निगमनाच्या वरच येत असेल.
\end{enumerate}
\end{defn}

\begin{lem}\label{pt:seq:qua:lem:subst-regular} $\pi$ ही
नियमित असून तिची अंतिम क्रमवर्ती $\Gamma \Sequent
\Delta$ आहे असे माना. तसेच~$c$ हा~$\pi$ मध्ये
आयगेन चर म्हणून न वापरलेला स्थिरांक असो,
आणि~$t$ हे~$\pi$ चा कोणताही आयगेन चर नसलेले बंद पद असो.
येथील संख्यापक-नियमांत पद बंद असण्याची अट आहे;
त्या संकेतानुसार ही आदेशनाची मर्यादा आवश्यक आहे.
मग~$\pi$ मधील~$c$ ची प्रत्येक उपस्थिती~$t$ ने
बदलून मिळणारी~$\Subst{\pi}{t}{c}$ ही नियमित
सिद्धता असते. $\Subst{\pi}{t}{c}$ ची अंतिम क्रमवर्ती
$\Subst{\Gamma}{t}{c}
\Sequent \Subst{\Delta}{t}{c}$ असते; येथे
$\Subst{\Gamma}{t}{c} =
\Setabs{A(t)}{A(c) \in \Gamma}$.
\end{lem}

\begin{proof}
  $\pheight{\pi}$ वर विगमन करू. $\pheight{\pi} = 0$
  असेल तर ती फक्त एक स्वयंसिद्धक आहे. मग
  $\Subst{\pi}{t}{c}$ हेही स्वयंसिद्धक असते, कारण
  त्यातील कोणतेही सूत्र~$A(c)$ हे~$A(t)$
  होते.

$\pheight{\pi} > 0$ असेल तर~$\pi$
  मधील शेवटच्या निगमनानुसार प्रसंग वेगळे करू.
  संरचनात्मक नियमांचे (\Log{G1c} मध्ये) आणि
  विधानात्मक संयोजकांच्या तार्किक नियमांचे प्रसंग
  सोपे आहेत; ते प्रश्न म्हणून सोडतो. $\pi$ चा शेवट
  संख्यकाच्या निगमनाने होणारे प्रसंग पाहू.
  \begin{enumerate}
    \item शेवटचे निगमन $\RightR{\lforall}$ असेल, तर
    $\pi$ या रूपाची असते:
\[
    \AxiomC{}
    \RightLabel{$\pi'$}
    \Deduce$\Gamma(c) \fCenter \Delta'(c), A(a,c)$
    \RightLabel{\RightR{\lforall}}
    \UnaryInf$\Gamma(c) \fCenter \Delta'(c), \lforall[x][A(x,c)]$
    \DisplayProof
    \]
    विगमनात्मक गृहीतकानुसार $\Subst{\pi'}{t}{c}$
    ही $\Gamma(t) \fCenter \Delta'(t), A(a,t)$ ची
    नियमित सिद्धता आहे. $a$ हा~$\pi$ चा आयगेन
    चर असल्यामुळे~$a$ हे~$t$ मध्ये येत नाही.
    म्हणून $\Subst{\pi}{t}{c}$ मधील शेवटचे निगमन,
\[
    \AxiomC{}
    \RightLabel{$\Subst{\pi'}{t}{c}$}
    \Deduce$\Gamma(t) \fCenter \Delta'(t), A(a,t)$
    \RightLabel{\RightR{\lforall}}
    \UnaryInf$\Gamma(t) \fCenter \Delta'(t), \lforall[x][A(x,t)]$
    \DisplayProof
    \]
    योग्य आहे: त्याच्या निष्कर्षात~$a$ येत नाही.
    \item शेवटचे निगमन \Log{G3c} मधील
    $\LeftR{\lforall}$ असेल, तर $\pi$ चे रूप असे:
\[
    \AxiomC{}
    \RightLabel{$\pi'$}
    \Deduce$A(s(c),c), \lforall[x][A(x,c)], \Gamma(c) \fCenter \Delta'(c)$
    \RightLabel{\LeftR{\lforall}}
    \UnaryInf$\lforall[x][A(x,c)], \Gamma(c) \fCenter \Delta'(c)$
    \DisplayProof
    \]
    विगमनात्मक गृहीतकानुसार $\Subst{\pi'}{t}{c}$
    ही $A(s(t),t), \lforall[x][A(x,t)], \Gamma(t) \fCenter
    \Delta'(t)$ ची नियमित सिद्धता आहे.
    $\Subst{\pi}{t}{c}$ मधील शेवटचे निगमन,
\[
    \AxiomC{}
    \RightLabel{$\Subst{\pi'}{t}{c}$}
    \Deduce$A(s(t),t), \lforall[x][A(x,t)], \Gamma(t) \fCenter \Delta'(t)$
    \RightLabel{\LeftR{\lforall}}
    \UnaryInf$\lforall[x][A(x,t)], \Gamma(t) \fCenter \Delta'(t)$
    \DisplayProof
    \]
    योग्य आहे. \Log{G1c} मधील \LeftR{\lforall}
    साठीचा प्रसंगही असाच, पण काहीसा सोपा आहे.
    \item शेवटचे निगमन~\RightR{\lexists} किंवा
    \LeftR{\lexists} असेल, तर ते प्रश्न म्हणून सोडले आहे.
  \end{enumerate}
  प्रत्येक प्रसंगात आपण नियमित सिद्धता च्या
  शेवटी एक क्रमवर्ती जोडली (दोन पूर्वपक्ष असलेल्या
  नियमात दोन नियमित सिद्धता वापरल्या). नव्या
  अंतिम क्रमवर्तीमध्ये~$\pi$ च्या अंतिम
  क्रमवर्तीपेक्षा फक्त~$c$ च्या जागी पद~$t$ आलेले
  असते. तसेच $\pi$ आणि $\Subst{\pi}{t}{c}$
  यांचे आयगेन चर तेच राहतात; कारण~$t$ मध्ये
  $\pi$ चा कोणताही आयगेन चर नाही, नव्या अंतिम
  क्रमवर्तीमध्येही $\Subst{\pi}{t}{c}$ चा आयगेन चर
  येत नाही. त्यामुळे~$\Subst{\pi}{t}{c}$ नियमित असते.
\end{proof}

\begin{prop}
$\pi$ ही \Log{G1c} किंवा \Log{G3c} मधील
सिद्धता असेल, तर तिच्याशी तीच रचना
(विशेषतः तीच उंची) असलेली नियमित
सिद्धता~$\pi'$ मिळते, जी फक्त आयगेन चर
बदलल्यामुळे~$\pi$ पेक्षा वेगळी असते.
\end{prop}

\begin{proof}
फक्त या सिद्धतेपुरते,~$\pi$ मधील आयगेन चराचे
एखादे निगमन \emph{स्वच्छ} म्हणू, जर त्याचा आयगेन
चर इतर कोणत्याही निगमनाचा आयगेन चर नसेल आणि
त्या निगमनाच्या तत्काळ उप-सिद्धता व्यतिरिक्त
~$\pi$ मध्ये कुठेही येत नसेल; अन्यथा ते
\emph{अस्वच्छ} म्हणू. $\pi$ मधील अस्वच्छ आयगेन
चराच्या निगमनांची संख्या~$n$ असो. $n$ वर
विगमन करून~$\pi$ चे अपेक्षित नियमित
सिद्धता~$\pi'$ मध्ये रूपांतर करता येते हे
दाखवू. $n=0$ असेल तर~$\pi$ मधील आयगेन चराची
सर्व निगमने स्वच्छ आहेत; म्हणून~$\pi$ नियमित
असते आणि सिद्ध करण्यास काही उरत नाही.
अन्यथा, सर्वांत वरचे अस्वच्छ आयगेन-चर निगमन,
उदा.~\RightR{\lforall}, निवडा. त्यात संपणारी
~$\pi$ ची उप-सिद्धता~$\pi_1$ पुढील रूपाची आहे:
\[
    \AxiomC{}
    \RightLabel{$\pi_2$}
    \Deduce$\Gamma \fCenter \Delta, A(c)$
    \RightLabel{\RightR{\lforall}}
    \UnaryInf$\Gamma \fCenter \Delta, \lforall[x][A(x)]$
    \DisplayProof
    \]
$c$ हा आयगेन चर असल्यामुळे तो~$\Gamma$,
$\Delta$, किंवा $\lforall[x][A(x)]$ मध्ये
येत नाही. $\pi_2$ नियमित आहे, कारण तिच्यातील
कोणतेही आयगेन-चर निगमन अस्वच्छ नाही.
~$\pi$ मध्ये न येणारा स्थिरांक~$a$ घ्या.
\ref{pt:seq:qua:lem:subst-regular} नुसार
$\Subst{\pi_2}{a}{c}$ ही $\Gamma \fCenter \Delta,
A(a)$ ची नियमित सिद्धता आहे; तिच्यावर
~$\RightR{\lforall}$ लावता येतो. $\pi$ मधील
~$\pi_1$ च्या जागी पुढील सिद्धता~$\pi_1'$ ठेवू:
\[
    \AxiomC{}
    \RightLabel{$\Subst{\pi_2}{a}{c}$}
    \Deduce$\Gamma \fCenter \Delta, A(a)$
    \RightLabel{\RightR{\lforall}}
    \UnaryInf$\Gamma \fCenter \Delta, \lforall[x][A(x)]$
    \DisplayProof
    \]
त्यामुळे एक अस्वच्छ आयगेन-चर निगमन स्वच्छ बनते;
फरक फक्त आयगेन चर~$c$ च्या जागी~$a$ आल्याचा
आहे. परिणामी सिद्धतेतील अस्वच्छ आयगेन-चर
निगमनांची संख्या~$n-1$ पेक्षा जास्त राहत नाही,
म्हणून विगमनात्मक गृहीतकातून निष्कर्ष मिळतो.
\end{proof}

नुकताच सिद्ध केलेला निष्कर्ष आपण स्पष्टपणे
संदर्भ म्हणून वापरणार नाही; पण पुढील सर्व
सिद्धता नियमित आहेत असे मानू. त्या नियमित
नसल्या तर आयगेन चर बदलून त्यांना नेहमी नियमित
बनवता येते.













\section{अनुज्ञेय आणि व्युत्पाद्य नियम}\label{pt:seq:adm:sec}

सिद्धता-उपपत्तीतील अनेक निष्कर्ष एखादा विशिष्ट नियम वापरून
सिद्ध करता येणारी प्रत्येक गोष्ट तो नियम न वापरताही
सिद्ध करता येते का, यासंबंधी असतात किंवा अशा
निष्कर्षांचा वापर करतात. कट-निरसन प्रमेय हे याचे सर्वांत
प्रसिद्ध उदाहरण आहे. त्यानुसार,~\Cut{} असलेल्या क्रमवर्ती
पद्धतीचे नियम वापरून सिद्ध करता येणारी कोणतीही क्रमवर्ती
~\Cut न वापरता सिद्ध करता येते. हा गुणधर्म इतका
महत्त्वाचा आहे की त्याला स्वतंत्र नाव आहे.

\begin{defn}
एखाद्या क्रमवर्ती पद्धतीत नियम~$R$ \emph{अनुज्ञेय} असतो,
जर त्या पद्धतीत नियम~$R$ जोडल्यावर सिद्ध करता येणारी
प्रत्येक क्रमवर्ती~$R$ शिवायही सिद्ध करता येत असेल.
\end{defn}

\begin{ex}\label{pt:seq:adm:ex:land-deriv}
नियम~$\LeftR{\land}$ चे \Log{G1c} आणि \Log{G3c} मध्ये
वेगवेगळे रूप आहे. \Log{G1c} मध्ये या नियमाच्या दोन आवृत्त्या
आहेत: एकीत प्रमुख सूत्र~$A \land B$ मधील
डावा संयोज्य~$A$ पूर्वपक्षात असतो, तर दुसरीत उजवा
संयोज्य~$B$ असतो:
\[
\Axiom$ A, \Gamma \fCenter \Delta$
\RightLabel{$\LeftR{\land}_1$}
\UnaryInf$ A \land B, \Gamma \fCenter \Delta$
\DisplayProof
\qquad
\Axiom$B, \Gamma \fCenter \Delta$
\RightLabel{$\LeftR{\land}_2$}
\UnaryInf$A \land B, \Gamma \fCenter \Delta$
\DisplayProof
\]
याउलट, \Log{G3c} मध्ये एकच नियम आहे:
\[
\Axiom$ A, B, \Gamma \fCenter \Delta$
\RightLabel{$\LeftR{\land}_3$}
\UnaryInf$ A \land B, \Gamma \fCenter \Delta$
\DisplayProof
\]
आता ``\Log{G3c} चा नियम~$\LeftR{\land}_3$
हा \Log{G1c} मध्ये अनुज्ञेय आहे का?'' असा प्रश्न विचारता येईल.
उत्तर होय आहे:~$\LeftR{\land}_3$ वापरून केलेल्या कोणत्याही
निगमनाचे \Log{G1c} च्याच नियमांनी थेट अनुकरण करता येते.
त्यासाठी~$\LeftR{\land}_1$ आणि $\LeftR{\land}_2$
यांबरोबर संरचनात्मक नियम~$\LeftR{\Contraction}$ ही लागेल:
\[
\Axiom$A, B, \Gamma \fCenter \Delta$
\RightLabel{$\LeftR{\land}_1$}
\UnaryInf$A \land B, B, \Gamma \fCenter \Delta$
\RightLabel{$\LeftR{\land}_2$}
\UnaryInf$A \land B, A \land B, \Gamma \fCenter \Delta$
\RightLabel{\LeftR{\Contraction}}
\UnaryInf$A \land B, \Gamma \fCenter \Delta$
\DisplayProof
\]
\end{ex}

एका नियमाने केलेल्या निगमनाचे इतर नियमांनी
अनुकरण करणे हा नियम अनुज्ञेय असल्याचे दाखवण्याचा एक मार्ग
आहे. पण तो एकमेव मार्ग नाही; काही नियम अनुज्ञेय असूनही
अशा प्रकारे त्यांचे अनुकरण करता येत नाही. ज्या नियमांचे
असे अनुकरण करता येते त्यांना \emph{व्युत्पाद्य} म्हणतात.

\begin{defn}
एखादा नियम~$R$,
\[
\AxiomC{$S_1$}
\AxiomC{$\dots$}
\AxiomC{$S_n$}
\RightLabel{$R$}
\TrinaryInfC{$S$}
\DisplayProof
\]
पद्धतीचे नियम आणि स्वयंसिद्धके, तसेच $S_1$, \dots,~$S_n$
या रूपांच्या अतिरिक्त प्रारंभीच्या क्रमवर्ती वापरून
क्रमवर्ती~$S$ ची रूपरेषात्मक सिद्धता असेल, तर तो नियम
त्या पद्धतीत \emph{व्युत्पाद्य} म्हणतात.
\end{defn}

\ref{pt:seq:adm:ex:land-deriv} मधील सिद्धता दाखवते
की~$\LeftR{\land}_3$ हा नियम \Log{G1c} मध्ये व्युत्पाद्य
आहे. कारण ती~$\LeftR{\land}_3$ च्या पूर्वपक्षांना प्रारंभीच्या
क्रमवर्ती मानून, फक्त \Log{G1c} चे नियम वापरून,
$\LeftR{\land}_3$ च्या निष्कर्षाची रूपरेषात्मक सिद्धता
आहे. (त्या सिद्धता मध्ये \Log{G1c}
मधील कोणतेही स्वयंसिद्धक वापरलेले नाही.)

\begin{prop}\label{pt:seq:adm:prop:lor-G3-adm}
\Log{G3c} चे नियम~$\LeftR{\land}$ आणि $\RightR{\lor}$
हे \Log{G1c} मध्ये व्युत्पाद्य आहेत.
\end{prop}

\begin{proof}
  $\LeftR{\land}$ साठीची सिद्धता \ref{pt:seq:adm:ex:land-deriv} मध्ये
  दिली आहे. $\RightR{\lor}$ साठीची सिद्धता प्रश्न म्हणून सोडली आहे.
\end{proof}

\begin{prob}
\Log{G3c} चा नियम~$\RightR{\lor}$ हा \Log{G1c} मध्ये
व्युत्पाद्य आहे हे दाखवा.
\end{prob}

\begin{prob}
$\liff$ हा परिभाषित संकारक आहे असे माना:
$A \liff B$ हे $(A \lif B) \land (B \lif A)$
याचे संक्षिप्त रूप आहे. पुढील नियम
\begin{enumerate}
\item \Axiom$\Gamma, A, B \fCenter \Delta$
\Axiom$\Gamma \fCenter \Delta, A, B$
\RightLabel{\LeftR{\liff}}
\BinaryInf$A \liff B, \Gamma \fCenter \Delta$
\DisplayProof
\item
\Axiom$A, \Gamma \fCenter \Delta, B$
\Axiom$B, \Gamma \fCenter \Delta, A$
\RightLabel{\RightR{\liff}}
\BinaryInf$\Gamma \fCenter \Delta, A \liff B$
\DisplayProof
\end{enumerate}
(a)~\Log{G1c} आणि (b)~\Log{G3c} मध्ये व्युत्पाद्य
आहेत हे दाखवा.
\end{prob}

\begin{ex}\label{pt:seq:adm:ex:weak-adm}
\Log{G1c} चा नियम~$\RightR{\Weakening}$,
\[
\Axiom$ \Gamma \fCenter \Delta$
\RightLabel{\RightR{\Weakening}}
\UnaryInf$ \Gamma \fCenter \Delta, A$
\DisplayProof
\]
हा \Log{G3c} मध्ये अनुज्ञेय आहे. सिद्धतेची कल्पना सोपी
आहे: \Log{G3c} मध्ये पूर्वपक्ष~$\Gamma \Sequent \Delta$
याची सिद्धता~$\pi$ आहे असे माना. $\pi$ मधील
प्रत्येक क्रमवर्तीच्या निष्कर्षभागात सूत्र~$A$ जोडा.
स्वयंसिद्धके स्वयंसिद्धकेच राहतात; योग्य निगमनेही योग्यच
राहतात. नव्या सिद्धता मध्ये हे सूत्र सर्वत्र
जोडलेले असते; तिची अंतिम क्रमवर्ती
$\Gamma \Sequent \Delta, A$ असते.
यापूर्वी~$\pi$ चे आयगेन स्थिरांक~$A$ मध्ये येणाऱ्या
स्थिरांकांपासून वेगळे नवी नावे देऊन घ्यायचे. नावबदल
अंतिम क्रमवर्ती आणि उंची जपतो; त्यानंतरच~$A$ सर्व
क्रमवर्तींमध्ये जोडायचे, जेणेकरून संख्यापक-नियमांची
आयगेन-चर अट जपली जाते.

तथापि, नियम~$\RightR{\Weakening}$ हा
\Log{G3c} मध्ये व्युत्पाद्य नाही. समजा तो व्युत्पाद्य
आहे; म्हणजे \Log{G3c} ची स्वयंसिद्धके व नियम आणि
कदाचित अतिरिक्त प्रारंभीची क्रमवर्ती~$\Gamma \Sequent
\Delta$ वापरून $\Gamma \Sequent \Delta, A$ ची
सिद्धता मिळते. $\RightR{\Weakening}$ व्युत्पाद्य
असण्यासाठी आवश्यक असल्याप्रमाणे ही सिद्धता पूर्णतः
रूपरेषात्मक असेल, तर $\Gamma$ आणि~$\Delta$ काढून
टाकून तिच्यापासून~$\Sequent A$ ची सिद्धता
बनवता येईल. त्यामुळे अतिरिक्त प्रारंभीच्या क्रमवर्ती
$\Gamma \Sequent \Delta$ चे रूपांतर रिकाम्या
क्रमवर्ती~$\quad \Sequent \quad$ मध्ये होईल. रिकाम्या
क्रमवर्तीमध्ये सूत्रे अजिबात नसतात, तर \Log{G3c}
च्या प्रत्येक नियमाच्या पूर्वपक्षात किमान एक सक्रिय
सूत्र आवश्यक असते. म्हणून या रूपरेषात्मक सिद्धतेतील
कोणत्याही निगमनाचा पूर्वपक्ष रिकामी क्रमवर्ती असू शकत
नाही. परिणामी, अतिरिक्त प्रारंभीची क्रमवर्ती~$\Gamma
\Sequent \Delta$ प्रत्यक्ष न वापरल्यासच ती सिद्धता
रूपरेषात्मक असू शकते. दुसऱ्या शब्दांत,
$\Gamma \Sequent \Delta, A$ चे \Log{G3c} चीच
स्वयंसिद्धके व नियम वापरून रूपरेषात्मक सिद्धता मिळेल;
मग या रूपाची कोणतीही क्रमवर्ती \Log{G3c} मध्ये
सिद्धता असलेली ठरेल. पण हे अशक्य आहे: \Log{G3c}
निर्दोष असल्यामुळे प्रत्येक संरचनांमध्ये सत्य
असलेल्या क्रमवर्तीच ते सिद्ध करू शकते. उदाहरणार्थ,
$\Gamma = \Delta = \emptyset$ आणि $A \ident \lfalse$
घेतल्यास, \Log{G3c} ला $\quad \Sequent \lfalse$
ही क्रमवर्ती सिद्ध करावी लागेल; ते तसे करत नाही.
\end{ex}

आता \Log{G3c} मध्ये दुर्बलीकरण हा अनुज्ञेय
नियम आहे, याची नीट विगमनात्मक सिद्धता देऊ. खरे तर,
वरील कल्पनेप्रमाणे किंचित अधिक प्रबळ निष्कर्ष मिळतो:
$\pi$ मधील प्रत्येक क्रमवर्तीच्या उजवीकडे~$A$
जोडल्याने सिद्धता ची रचना बदलत नाही; विशेषतः तिची
उंची वाढत नाही. हे महत्त्वाचे असल्याने स्वतंत्र व्याख्या
देऊ.

\begin{defn}
एखादा नियम~$R$,
\[
\AxiomC{$S_1$}
\AxiomC{$\dots$}
\AxiomC{$S_n$}
\RightLabel{$R$}
\TrinaryInfC{$S$}
\DisplayProof
\]
पद्धतीत \emph{उंची जपून अनुज्ञेय} (किंवा
\emph{उंची जपून-अनुज्ञेय}) असतो, जर $i = 1$, \dots,~$n$
साठी $\Proves[h_i] S_i$ असताना $m = \max(h_1, \dots,
h_n)$ घेऊन $\Proves_m S$ मिळत असेल.
\end{defn}

नियम उंची जपून अनुज्ञेय असण्यासाठी पुढील अट
पुरेशी आहे: पूर्वपक्ष~$S_i$ यांच्या उंची~$h_i =
\pheight{\pi_i}$ असलेल्या सिद्धता~$\pi_i$ दिल्या
असता, निष्कर्षाची अशी सिद्धता~$\pi$ असावी की
तिची उंची~$\pheight{\pi}$ ही~$h_i$ यांच्या कमाल
मूल्यापेक्षा मोठी नाही. विशेषतः, एकच पूर्वपक्ष~$S_1$
असेल तर $\pheight{\pi} \le \pheight{\pi_1}$ पुरेसे
आहे. $\RightR{\Weakening}$ आणि $\LeftR{\Weakening}$
बाबतीत तर प्रत्यक्षात $\pheight{\pi} =
\pheight{\pi_1}$ असते; पण समानता आवश्यक नाही.

\begin{prop}\label{pt:seq:adm:prop:weak-G3c-adm}
\Log{G1c} चे नियम~$\RightR{\Weakening}$ आणि
$\LeftR{\Weakening}$ हे \Log{G3c} मध्ये
उंची जपून अनुज्ञेय आहेत.
\end{prop}

\begin{proof}
  $\pheight{\pi} = n$ उंची असलेली सिद्धता~$\pi$
  वापरून $\Log{G3c} \Proves \Gamma \Sequent \Delta$
  असेल, तर $\pheight{\pi'} = n$ असलेली
  $\Gamma \Sequent \Delta, A$ ची \Log{G3c}-सिद्धता
  $\pi'$ आहे, हे दाखवायचे आहे. $n$ वर विगमन करू.
  $n=0$ असेल तर~$\pi$ हे फक्त $\Gamma \Sequent
  \Delta$ स्वयंसिद्धक आहे (म्हणजे एखादे अण्वीय~$C$
  हे~$\Gamma$ आणि~$\Delta$ दोन्हीत येते, किंवा
  $\lfalse$ डावीकडे असते); म्हणून
  $\Gamma \Sequent \Delta, A$ हेही स्वयंसिद्धक आहे.
  $n = \pheight{\pi} >0$ असेल तर सिद्धता~$\pi$
  मध्ये शेवटचे एक निगमन असते. कोणता नियम वापरला यानुसार
  उदाहरणे वेगळी करतो. एकच उदाहरण देऊ: शेवटचे निगमन
  $\RightR{\land}$ होते असे समजा. मग $\Delta =
  \Delta', B \land C$; शेवटच्या निगमनातील तत्काळ
  उप-सिद्धता~$\pi_1$ आणि~$\pi_2$ यांच्या अंतिम
  क्रमवर्ती अनुक्रमे $\Gamma \Sequent \Delta', B$
  आणि $\Gamma \Sequent \Delta', C$ असतात. म्हणजे
  $\pi$ चे रूप असे असते:
\[
  \AxiomC{}
  \RightLabel{$\pi_1$}
  \Deduce$\Gamma \fCenter \Delta', B$
  \AxiomC{}
  \RightLabel{$\pi_2$}
  \Deduce$\Gamma \fCenter \Delta', C$
  \RightLabel{\RightR{\land}}
  \BinaryInf$\Gamma \fCenter \Delta', B \land C$
  \DisplayProof
  \]
  तसेच $n =
  \max(\pheight{\pi_1}, \pheight{\pi_2}) + 1$.
  म्हणून $\pheight{\pi_1} <
  n$ आणि $\pheight{\pi_2} < n$; त्यामुळे विगमनात्मक
  गृहीतक लागू होते: $\Gamma \Sequent
  \Delta', B, A$ आणि $\Gamma \Sequent \Delta', C, A$
  यांच्या अनुक्रमे सिद्धता~$\pi_1'$ आणि~$\pi_2'$
  मिळतात. $\RightR{\land}$ नियमाने $\Gamma
  \Sequent \Delta', B \land C, A$ ची सिद्धता~$\pi'$
  मिळते:
\[
  \AxiomC{}
  \RightLabel{$\pi_1'$}
  \Deduce$\Gamma \fCenter \Delta', B, A$
  \AxiomC{}
  \RightLabel{$\pi_2'$}
  \Deduce$\Gamma \fCenter \Delta', C, A$
  \RightLabel{\RightR{\land}}
  \BinaryInf$\Gamma \fCenter \Delta', B \land C, A$
  \DisplayProof
  \]
  आपल्याला $\pheight{\pi'} =
  \max(\pheight{\pi_1'}, \pheight{\pi_2'}) + 1 = \max(\pheight{\pi_1},
  \pheight{\pi_2}) + 1 = \pheight{\pi}$ मिळते.
\end{proof}

\ref{pt:seq:adm:prop:weak-G3c-adm} अत्यंत उपयोगी आहे.
त्याचा वारंवार वापर दाखवण्यासाठी एक संकेतन ठरवू:
$\pi$ ही~$\Gamma \Sequent \Delta$ ची सिद्धता
असेल, तर $\pi[\Pi \Sequent \Lambda]$ म्हणजे डावीकडे
$\Pi$ मधील सर्व सूत्रे आणि उजवीकडे~$\Lambda$
मधील सर्व सूत्रे जोडण्यासाठी दुर्बलीकरणाच्या
अनुज्ञेयतेचा केलेला वापर. परिणामी ती $\Gamma, \Pi
\Sequent \Delta, \Lambda$ ची सिद्धता असते.













\section{नियमांची व्युत्क्रमणीयता}\label{pt:seq:inv:sec}

\begin{defn}
एखादा नियम~$R$,
\[
\AxiomC{$S_1$}
\AxiomC{$\dots$}
\AxiomC{$S_n$}
\RightLabel{$R$}
\TrinaryInfC{$S$}
\DisplayProof
\]
पद्धतीत \emph{व्युत्क्रमणीय} असतो, जर $\Proves S$
असताना $i = 1$, \dots,~$n$ साठी $\Proves[h_i] S_i$
मिळत असेल. तो \emph{उंची जपून व्युत्क्रमणीय}
(किंवा \emph{उंची जपून-व्युत्क्रमणीय}) असतो, जर
$\Proves_h S$ असताना $i = 1$, \dots,~$n$
साठी $\Proves[h] S_i$ मिळत असेल.
\end{defn}

\begin{lem}\label{pt:seq:inv:lem:G3c-invert}
\Log{G3c} मधील $\lnot$, $\land$, $\lor$, आणि $\lif$
यांचे नियम उंची जपून व्युत्क्रमणीय आहेत;
म्हणजे:
\begin{enumerate}
  \item \RightR{\lnot}: $\Log{G3c} \Proves[n] \Gamma \Sequent
  \Delta, \lnot A$ असेल, तर $\Log{G3c} \Proves[n] A, \Gamma \Sequent
  \Delta$.
  \item \LeftR{\lnot}: $\Log{G3c} \Proves[n] \Gamma, \lnot A
  \Sequent \Delta$ असेल, तर $\Log{G3c} \Proves[n] \Gamma \Sequent \Delta,
  A$.
  \item \RightR{\land}: $\Log{G3c} \Proves[n] \Gamma \Sequent
  \Delta, A \land B$ असेल, तर $\Log{G3c} \Proves[n] \Gamma \Sequent
  \Delta, A$ आणि $\Log{G3c} \Proves[n] \Gamma \Sequent
  \Delta, B$.
  \item \LeftR{\land}: $\Log{G3c} \Proves[n] \Gamma, A \land B
  \Sequent \Delta$ असेल, तर $\Log{G3c} \Proves[n] A, B, \Gamma \Sequent
  \Delta$.
  \item \RightR{\lor}: $\Log{G3c} \Proves[n] \Gamma \Sequent
  \Delta, A \lor B$ असेल, तर $\Log{G3c} \Proves[n] \Gamma \Sequent
  \Delta, A, B$.
  \item \LeftR{\lor}: $\Log{G3c} \Proves[n] A \lor B, \Gamma \Sequent
  \Delta$ असेल, तर $\Log{G3c} \Proves[n] A, \Gamma \Sequent
  \Delta$ आणि $\Log{G3c} \Proves[n] B, \Gamma \Sequent
  \Delta$.
  \item \RightR{\lif}: $\Log{G3c} \Proves[n] \Gamma \Sequent
  \Delta, A \lif B$ असेल, तर $\Log{G3c} \Proves[n] A, \Gamma \Sequent
  \Delta, B$.
  \item \LeftR{\lif}: $\Log{G3c} \Proves[n] A \lif B, \Gamma
  \Sequent \Delta$ असेल, तर $\Log{G3c} \Proves[n] \Gamma \Sequent \Delta,
  A$ आणि $\Log{G3c} \Proves[n] B, \Gamma \Sequent \Delta$.
\end{enumerate}
\end{lem}

\begin{proof}
वरील प्रमेयिकेत दिलेल्या विधानीय नियम~$R$ साठी पुढील
गोष्ट दाखवायची आहे:~$R$ च्या निष्कर्ष~$S$ ची
सिद्धता~$\pi$ असेल, तर $R$ च्या पूर्वपक्ष~$S_i$
यांच्या सिद्धता~$\pi_i$ सुद्धा असतात, आणि
$\pheight{\pi_i}
\le \pheight{\pi}$. नियम~\LeftR{\land} विचारात घ्या:
अंतिम क्रमवर्ती $A \land B, \Gamma
\Sequent \Delta$ या रूपाची असलेली सिद्धता~$\pi$
आहे असे समजा. $A, B, \Gamma \Sequent \Delta$
ची अशी सिद्धता~$\pi'$ आहे की $\pheight{\pi'} \le
\pheight{\pi}$, हे दाखवायचे आहे. $n = \pheight{\pi}$
वर विगमन करून हे सिद्ध करू.

$n = 0$ असेल तर $\pi$ फक्त एकच स्वयंसिद्धक
असतो: $A \land B, \Gamma
\Sequent \Delta$. \Log{G3c} मधील स्वयंसिद्धके
$C, \Gamma
\Sequent \Delta, C$ या रूपाची असतात, आणि~$C$
अण्वीय असतो. म्हणजेच $C$ हा $A \land B$
असू शकत नाही. $A \land B, \Gamma \Sequent \Delta$
स्वयंसिद्धक असेल, तर एखादा अण्वीय~$C$ हा~$\Gamma$
आणि~$\Delta$ दोहोंचा घटक असला पाहिजे.
मग $A, B, \Gamma \Sequent \Delta$ हाही
स्वयंसिद्धक असतो आणि आपल्याला हवी असलेली
सिद्धता~$\pi'$ मिळाली (तिची उंची ही~$0$).
दुसरा स्वयंसिद्धक-प्रसंग म्हणजे~$\lfalse\in\Gamma$.
संयोगाची आस्थिती बदलल्यावरही तो डावीकडे राहतो;
म्हणून अपेक्षित क्रमवर्ती पुन्हा उंची~$0$ चे
असत्य-स्वयंसिद्धक आहे.

$n>0$ असेल, तर सिद्धतेचा शेवट निगमनाने
होतो. $n$ साठी दावा सिद्ध करायचा आहे; परंतु संदर्भ
वेगळा असला तरी उंची~$<n$ असलेल्या कोणत्याही
सिद्धता साठी दावा खरा आहे असे मानता येते. म्हणजे
कोणत्याही $k<n$ आणि कोणत्याही $\Pi$ व $\Lambda$
साठी $A \land B, \Pi \Sequent \Lambda$ ची
सिद्धता उंची~$k$ असेल, तर $A, B, \Pi
\Sequent \Lambda$ ची सिद्धता उंची~$\le k$
असलेली मिळते.

दोन प्रसंग आहेत: डावीकडील $A \land
B$ प्रमुख असतो किंवा नसतो. तो प्रमुख असेल तर
शेवटचे निगमन~\LeftR{\land} असते आणि तत्काळ
उप-सिद्धता~$\pi_1$ ची अंतिम क्रमवर्ती
$A, B, \Gamma \Sequent \Delta$ असते. या
प्रसंगात निष्कर्ष मिळतो, कारण $\pheight{\pi_1} <
\pheight{\pi}$.

$A \land B$ प्रमुख नसेल तर दुसरे
एखादे सूत्र प्रमुख असते आणि $A
\land B$ हा शेवटच्या निगमनाच्या संदर्भात असतो.
मग शेवटच्या नियम~$R$ च्या तत्काळ उप-सिद्धता
ना विगमनात्मक गृहीतक लावता येते. त्यातून~$\pi$
च्या तत्काळ उप-सिद्धता च्या उंचींपेक्षा मोठी
उंची नसलेली सिद्धता~$\pi_1'$ किंवा
दोन सिद्धता~$\pi_1'$ आणि~$\pi_2'$ मिळतात.
$\pi_1'$ आणि~$\pi_2'$ यांच्या अंतिम क्रमवर्तींमध्ये
डावीकडील संदर्भात $A \land B$ ऐवजी $A, B$
येतात. त्यांना नियम~$R$ लावून सिद्धता~$\pi'$
मिळते. उदाहरणार्थ,~$\pi$ चा शेवट~\LeftR{\lif}
ने होत असेल तर:
\[
\AxiomC{}
\RightLabel{$\pi_1$}
\Deduce$A \land B, \Gamma' \fCenter \Delta, C$
\AxiomC{}
\RightLabel{$\pi_2$}
\Deduce$A \land B, \Gamma', D \fCenter \Delta$
\RightLabel{\LeftR{\lif}}
\BinaryInf$A \land B, \Gamma', C \lif D \fCenter \Delta$
\DisplayProof
\]
येथे डावीकडील $C \lif D$ प्रमुख असून डावीकडील
संदर्भ $A \land B, \Gamma'$ आहे (म्हणजे
$\Gamma = \Gamma', C \lif D$). विगमनात्मक
गृहीतकाने अनुक्रमे $A, B, \Gamma' \Sequent \Delta, C$
आणि $A, B, \Gamma', D \Sequent \Delta$
यांच्या सिद्धता~$\pi_1'$
आणि~$\pi_2'$ मिळतात. या दोन क्रमवर्ती पुन्हा
\LeftR{\lif} चे पूर्वपक्ष बनून~$\pi'$ देतात:
\[
\AxiomC{}
\RightLabel{$\pi_1'$}
\Deduce$A, B, \Gamma' \fCenter \Delta, C$
\AxiomC{}
\RightLabel{$\pi_2'$}
\Deduce$A, B, \Gamma', D \fCenter \Delta$
\RightLabel{\LeftR{\lif}}
\BinaryInf$A, B, \Gamma', C \lif D \fCenter \Delta$
\DisplayProof
\]
आपल्याला पुढील मिळते:
\begin{align*}
\pheight{\pi} & = \max(\pheight{\pi_1}, \pheight{\pi_2}) + 1\\
\pheight{\pi'} & = \max(\pheight{\pi_1'}, \pheight{\pi_2'}) + 1
\intertext{व्याख्येनुसार; आणि}
\pheight{\pi_1'} & \le \pheight{\pi_1}\\
\pheight{\pi_2'} & \le \pheight{\pi_2}
\intertext{विगमन गृहीतकानुसार. म्हणून,}
\pheight{\pi'} & \le \pheight{\pi}.
\end{align*}
\end{proof}

\begin{prob}
\RightR{\lif} साठी \ref{pt:seq:inv:lem:G3c-invert} ची सिद्धता द्या.
\end{prob}

\begin{lem}\label{pt:seq:inv:lem:invert-quant}
\RightR{\lforall} आणि \LeftR{\lexists} हे नियम
पुढील अर्थाने उंची जपून व्युत्क्रमणीय आहेत:
\begin{enumerate}
  \item $\Log{G3c} \Proves[n] \Gamma \Sequent \Delta,
  \lforall[x][A(x)]$ असेल, तर $\Log{G3c} \Proves[n] \Gamma \Sequent
  \Delta, A(c)$; आणि
  \item $\Log{G3c} \Proves[n] \lexists[x][A(x)], \Gamma \Sequent
  \Delta$ असेल, तर $\Log{G3c} \Proves[n] A(c), \Gamma \Sequent \Delta$,
\end{enumerate}
येथे प्रत्येक प्रसंगात~$c$ हा अनुक्रमे $\Gamma$, $\Delta$
आणि $\lforall[x][A(x)]$ किंवा $\lexists[x][A(x)]$
यांत न येणारा कोणताही स्थिरांक आहे.
\end{lem}

\begin{proof}
  आपण (1) सिद्ध करू आणि (2) प्रश्न म्हणून सोडू.
  दिलेल्या~$c$ पासून सर्व अंतर्गत आयगेन स्थिरांक
  वेगळे होतील अशी सुरक्षित नियमित नावे आधी घ्या.
  युक्तिवाद \ref{pt:seq:inv:lem:G3c-invert} च्या सिद्धतेसारखाच
  आहे. विगमनात्मक टप्प्यात पुन्हा दोन प्रसंग असतात.
  $\lforall[x][A(x)]$ प्रमुख असण्याचा प्रसंग सोपा आहे:
  शेवटच्या $\RightR{\lforall}$ निगमनाचा पूर्वपक्ष
  अपेक्षित $\Gamma
  \Sequent \Delta, A(a)$ या रूपाचा असतो आणि
  त्यावर संपणाऱ्या तत्काळ उप-सिद्धता~$\pi_1$ ची
  $\pheight{\pi_1} < \pheight{\pi}$ असते.
  शिवाय, स्वचलाच्या अटीमुळे $a$ हा
  $\Gamma$, $\Delta$, किंवा $\lforall[x][A(x)]$
  यांत येत नाही. $\pi_1$ नियमित आहे असे मानता
  येत असल्याने \ref{pt:seq:qua:lem:subst-regular} नुसार
  $\Subst{\pi_1}{c}{a}$ ही $\Gamma \Sequent \Delta, A(c)$
  ची त्याच उंची ची सिद्धता आहे.
उंची~$0$ चा प्रसंगही जपला जातो: संख्यापकीय सूत्र
अणुरूप नसल्याने समान अणुरूप सूत्राचा किंवा डावीकडील
असत्याचा साक्षी उरलेल्या संदर्भात आहे; म्हणून त्याची
जागा~$A(c)$ ने बदलल्यावरही क्रमवर्ती स्वयंसिद्धक आहे.
नव्या नियमाच्या आयगेन स्थिरांकांनाही~$c$ पासून वेगळे
नावे देऊन उरलेल्या विधानीय प्रसंगांप्रमाणे नियम पुन्हा
लावायचा.

दुसऱ्या प्रसंगात $\lforall[x][A(x)]$ हे
  शेवटच्या निगमनात प्रमुख नसते; दुसरे सूत्र प्रमुख
  असते. पुन्हा एका उदाहरणाने प्रसंग पाहू. शेवटचे
  निगमन~\RightR{\land} आहे असे समजा.
  $\Delta$ मधील एखादे सूत्र~$B \land C$
  या रूपाचे असून प्रमुख आहे. सिद्धता~$\pi$
  अशी संपते:
\[
\AxiomC{}
\RightLabel{$\pi_1$}
\Deduce$\Gamma \fCenter \Delta', \lforall[x][A(x)], B$
\AxiomC{}
\RightLabel{$\pi_2$}
\Deduce$\Gamma \fCenter \Delta', \lforall[x][A(x)], C$
\RightLabel{\RightR{\land}}
\BinaryInf$\Gamma \fCenter \Delta', \lforall[x][A(x)], B \land C$
\DisplayProof
\]
येथे $\Delta = \Delta', B \land C$.
$\Delta$ मध्ये आणि पूर्ण निष्कर्षाच्या उरलेल्या
भागातही न येणारा स्थिरांक~$c$ घ्या;
तो~$\Delta', B$ किंवा $\Delta', C$ मध्येही
येत नाही हे उघड आहे. म्हणून विगमनात्मक गृहीतक
$\pi_1$ आणि~$\pi_2$ ना लागू होते; आपल्याकडे
सिद्धता~$\pi_1'$ आणि $\pi_2'$ आहेत, ज्यांसाठी
$\pheight{\pi_1'} \le \pheight{\pi_1}$ आणि
$\pheight{\pi_2'} \le
\pheight{\pi_2}$. अपेक्षित $\Gamma \Sequent
\Delta, A(c)$ ची सिद्धता~$\pi'$ अशी:
\[
\AxiomC{}
\RightLabel{$\pi_1'$}
\Deduce$\Gamma \fCenter \Delta', A(c), B$
\AxiomC{}
\RightLabel{$\pi_2'$}
\Deduce$\Gamma \fCenter \Delta', A(c), C$
\RightLabel{\RightR{\land}}
\BinaryInf$\Gamma \fCenter \Delta', A(c), B \land C$
\DisplayProof
\]
आणि $\pheight{\pi'} = \max(\pheight{\pi_1'}, \pheight{\pi_2'})+1 \le \pheight{\pi}$.
\end{proof}

\ref{pt:seq:inv:lem:G3c-invert} चा कट-निरसन प्रमेयाच्या
सिद्धतेत महत्त्वाचा उपयोग होतो. पण त्याने पुढील
निष्कर्षही दाखवता येतो:

\begin{prop}\label{pt:seq:inv:prop:G3c-cont-adm}
\Log{G1c} चे आकुंचन नियम~\LeftR{\Contraction} आणि
\RightR{\Contraction} हे \Log{G3c} मध्ये
उंची जपून अनुज्ञेय आहेत.
\end{prop}

\begin{proof}
\RightR{\Contraction} आणि \LeftR{\Contraction}
दोन्हींसाठी एकाच वेळी युक्तिवाद देऊ.
\Log{G3c} मध्ये $\pi$ ही $\Gamma \Sequent
\Delta, A, A$ ची किंवा $A, A, \Gamma \Sequent
\Delta$ ची सिद्धता आहे असे माना. अनुक्रमे
$\Gamma \Sequent \Delta, A$ किंवा $A,
\Gamma \Sequent \Delta$ ची अशी \Log{G3c}-सिद्धता
~$\pi'$ आहे की $\pheight{\pi'} \le
\pheight{\pi}$, हे दाखवायचे आहे. ते
$n = \pheight{\pi}$ वर विगमन करून करू.

$n=0$ असेल, तर~$\pi$ हे फक्त स्वयंसिद्धक
आहे. $\Gamma \Sequent \Delta, A,
A$ हे \Log{G3c} चे स्वयंसिद्धक असेल तर
$\Gamma \Sequent \Delta, A$ हेही स्वयंसिद्धक
असते: कारण एखादा अण्वीय~$C$ हा~$\Gamma$
आणि~$\Delta$ दोहोंचा घटक असतो, किंवा
$A$ अण्वीय असून~$\Gamma$ चा घटक
असतो. अंतिम क्रमवर्ती $A, A, \Gamma \Sequent
\Delta$ असेल, तरी हाच युक्तिवाद लागू होतो.
डावीकडील असत्यामुळे स्वयंसिद्धक असेल, तर आकुंचनानंतर
किमान एक असत्य आस्थिती राहते; म्हणून तो प्रसंगही
उंची~$0$ चा असतो.

$n>0$ असेल, तर~$\pi$ चा शेवट काही
नियम~$R$ ने होतो. या शेवटच्या निगमनात $A$
प्रमुख नसेल, तर \ref{pt:seq:inv:lem:G3c-invert} च्या
सिद्धतेप्रमाणे~$\pi$ च्या तत्काळ उप-सिद्धता
ना विगमनात्मक गृहीतक लावून आणि मिळालेल्या
सिद्धता(s) ना पुन्हा तोच नियम~$R$ लावून
सिद्धता~$\pi'$ मिळवता येते. विगमनात्मक गृहीतक
असे आहे: $\Pi \Sequent \Lambda, D,
D$ किंवा $D, D, \Pi \Sequent \Lambda$ यांची
उंची~$k<n$ असलेली सिद्धता असेल, तर
अनुक्रमे $\Pi \Sequent \Lambda, D$ किंवा $D, \Pi \Sequent
\Lambda$ यांची उंची~$\le k$ असलेली
सिद्धता आहे. (संदर्भ किंवा आकुंचित होणारे
सूत्र हे~$\Gamma$, $\Delta$, किंवा~$A$
यांहून वेगळे असले तरी विगमनात्मक गृहीतक लागू होते!)

उदाहरणार्थ, $R$ हा~$\RightR{\land}$
असून~$\pi$ चा शेवट असा आहे असे समजा:
\[
\AxiomC{}
\RightLabel{$\pi_1$}
\Deduce$\Gamma \fCenter \Delta', B, A, A$
\AxiomC{}
\RightLabel{$\pi_2$}
\Deduce$\Gamma \fCenter \Delta', C, A, A$
\RightLabel{\RightR{\land}}
\BinaryInf$\Gamma \fCenter \Delta', B \land C, A, A$
\DisplayProof
\]
येथे $\Delta = \Delta', B \land C$. विगमनात्मक
गृहीतकाने अनुक्रमे $\Gamma \fCenter \Delta', B, A$
आणि $\Gamma \fCenter \Delta', C, A$ यांच्या
सिद्धता~$\pi_1'$ आणि $\pi_2'$ मिळतात, आणि
$\pheight{\pi_1'} \le  \pheight{\pi_1}$ व
$\pheight{\pi_2'} \le
\pheight{\pi_2}$ असते. मग सिद्धता~$\pi'$ अशी:
\[
\AxiomC{}
\RightLabel{$\pi_1'$}
\Deduce$\Gamma \fCenter \Delta', B, A$
\AxiomC{}
\RightLabel{$\pi_2'$}
\Deduce$\Gamma \fCenter \Delta', C, A$
\RightLabel{\RightR{\land}}
\BinaryInf$\Gamma \fCenter \Delta', B \land C, A$
\DisplayProof
\]

शेवटच्या निगमनात~$A$ प्रमुख असेल, तर
~$\pi$ च्या तत्काळ उप-सिद्धता ना व्युत्क्रमण
पूर्वप्रमेय (\ref{pt:seq:inv:lem:G3c-invert}) लावतो, मग
विगमनात्मक गृहीतक लावतो, आणि शेवटी नियम~$R$
पुन्हा लावतो. उदाहरणार्थ, $\pi$ ही
$\Gamma \Sequent \Delta, A, A$ ची सिद्धता
असून $A
\ident (B \land C)$ हे शेवटच्या निगमनात
प्रमुख आहे असे समजा. मग~$\pi$ चे रूप असे:
\[
\AxiomC{}
\RightLabel{$\pi_1$}
\Deduce$\Gamma \fCenter \Delta, B \land C, B$
\AxiomC{}
\RightLabel{$\pi_2$}
\Deduce$\Gamma \fCenter \Delta, B \land C, C$
\RightLabel{\RightR{\land}}
\BinaryInf$\Gamma \fCenter \Delta, B \land C, B \land C$
\DisplayProof
\]
\ref{pt:seq:inv:lem:G3c-invert} हे~$\pi_1$ ला लावल्यास
$\Gamma \Sequent \Delta, B, B$ ची
सिद्धता~$\pi_1'$ मिळते. (त्यातून
$\Gamma \Sequent \Delta, C, B$ ची सिद्धता
सुद्धा मिळते, पण आपल्याला ती लागत नाही.)
तसेच~$\pi_2$ ला \ref{pt:seq:inv:lem:G3c-invert} लावल्यास
$\Gamma \Sequent \Delta, C, C$ ची
सिद्धता~$\pi_2'$ मिळते. \RightR{\land} चे
व्युत्क्रमण उंची जपते, म्हणून
$\pheight{\pi_1'} \le  \pheight{\pi_1} < \pheight{\pi}$
आणि $\pheight{\pi_2'} \le  \pheight{\pi_2} < \pheight{\pi}$.
त्यामुळे~$\pi_1'$ आणि $\pi_2'$ ना विगमनात्मक गृहीतक
लागू होते: अनुक्रमे $\Gamma \Sequent \Delta, B$
आणि $\Gamma \Sequent \Delta, C$ यांच्या
सिद्धता~$\pi_1''$ आणि $\pi_2''$ मिळतात, आणि
$\pheight{\pi_1''}
\le \pheight{\pi_1'} \le \pheight{\pi_1}$ तसेच
$\pheight{\pi_2''} \le
\pheight{\pi_2'} \le \pheight{\pi_2}$ असते.
मग सिद्धता~$\pi'$ अशी:
\[
\AxiomC{}
\RightLabel{$\pi_1''$}
\Deduce$\Gamma \fCenter \Delta, B$
\AxiomC{}
\RightLabel{$\pi_2''$}
\Deduce$\Gamma \fCenter \Delta, C$
\RightLabel{\RightR{\land}}
\BinaryInf$\Gamma \fCenter \Delta, B \land C$
\DisplayProof
\]
तिच्यासाठी $\pheight{\pi'} \le \pheight{\pi}$.


सूत्र प्रमुख असताना संख्यकांच्या
नियमांबाबत विशेष काळजी घ्यावी लागते.

$A \ident \lforall[x][B(x)]$ हे उजवीकडे
प्रमुख आहे असे समजा. मग~$\pi$ चा शेवट असा:
\[
\AxiomC{}
\RightLabel{$\pi_1$}
\Deduce$\Gamma \fCenter \Delta, \lforall[x][B(x)], B(a)$
\RightLabel{\RightR{\lforall}}
\UnaryInf$\Gamma \fCenter \Delta, \lforall[x][B(x)], \lforall[x][B(x)]$
\DisplayProof
\]
\ref{pt:seq:inv:lem:invert-quant} नुसार, $\Gamma \fCenter \Delta,
\lforall[x][B(x)], B(a)$ मध्ये न येणाऱ्या
कोणत्याही~$c$ साठी $\Gamma \fCenter \Delta, B(c),
B(a)$ ची उंची~$\le \pheight{\pi_1}$
असलेली सिद्धता~$\pi_1'$ मिळते. $\pi_1'$
नियमित करून तिचे अंतर्गत आयगेन स्थिरांक~$a$ पासून
वेगळे आहेत असे मानता येते; म्हणून
\ref{pt:seq:qua:lem:subst-regular} नुसार
सिद्धता~$\Subst{\pi_1'}{a}{c}$ ही $\Gamma \fCenter \Delta, B(a),
B(a)$ ची त्याचप्रमाणे उंची~$\le \pheight{\pi_1}$
असलेली सिद्धता आहे. आता विगमनात्मक गृहीतक
लागू होऊन $\Gamma
\fCenter \Delta, B(a)$ ची सिद्धता~$\pi_1''$
मिळते. \RightR{\lforall} नियम लावून~$\pi'$
मिळवू:
\[
\AxiomC{}
\RightLabel{$\pi_1''$}
\Deduce$\Gamma \fCenter \Delta, B(a)$
\RightLabel{\RightR{\lforall}}
\UnaryInf$\Gamma \fCenter \Delta, \lforall[x][B(x)]$
\DisplayProof
\]
आणि $\pheight{\pi'} \le \pheight{\pi}$.


आता $A \ident \lexists[x][B(x)]$ हे
उजवीकडे प्रमुख आहे असे समजा. मग~$\pi$ चे रूप असे:
\[
\AxiomC{}
\RightLabel{$\pi_1$}
\Deduce$\Gamma \fCenter \Delta, \lexists[x][B(x)], \lexists[x][B(x)], B(t)$
\RightLabel{\RightR{\lexists}}
\UnaryInf$\Gamma \fCenter \Delta, \lexists[x][B(x)], \lexists[x][B(x)]$
\DisplayProof
\]
विगमनात्मक गृहीतक~$\pi_1$ ला लागू होते; त्यामुळे
$\Gamma \Sequent \Delta, \lexists[x][B(x)], B(t)$
ची सिद्धता~$\pi_1'$ मिळते.
(येथे व्युत्क्रमण पूर्वप्रमेयाची गरज नाही!) मग
सिद्धता~$\pi'$ अशी:
\[
\AxiomC{}
\RightLabel{$\pi_1'$}
\Deduce$\Gamma \fCenter \Delta,  \lexists[x][B(x)], B(t)$
\RightLabel{\RightR{\lexists}}
\UnaryInf$\Gamma \fCenter \Delta, \lexists[x][B(x)]$
\DisplayProof
\]
तिच्यासाठी $\pheight{\pi'} \le \pheight{\pi}$.
\end{proof}

\begin{prob}
\LeftR{\Contraction} साठी \ref{pt:seq:inv:prop:G3c-cont-adm}
च्या सिद्धतेची रूपरेषा द्या. $A \ident (B \land C)$
आणि $A \ident (B
\lif C)$ हे प्रसंग वर्णा.
\end{prob}

\begin{prob}
उरलेल्या संख्यकांच्या प्रसंगांसाठी, म्हणजे $A \ident
\lforall[x][B(x)]$ आणि $A \ident \lexists[x][B(x)]$
असताना, \LeftR{\Contraction} साठी
\ref{pt:seq:inv:prop:G3c-cont-adm} च्या सिद्धतेची रूपरेषा द्या.
\end{prob}













\section{\Log{G1c} आणि \Log{G3c} यांतील रूपांतरे}\label{pt:seq:tr:sec}

\Log{G1c} आणि \Log{G3c} या दोन्ही पद्धती अभिजात तर्कशास्त्रासाठी
निर्दोष व संपूर्ण आहेत, असे आपण म्हटले होते. म्हणजे प्रत्येक
संरचनांमध्ये सत्य असलेल्या आणि फक्त त्याच क्रमवर्ती त्या
सिद्ध करतात. विशेषतः, दोन्ही त्याच क्रमवर्ती सिद्ध करतात. आता
निर्दोषता व संपूर्णतेच्या प्रमेयांचा वळसा न घालता, केवळ
सिद्धता-उपपत्तीय मार्गाने हे तथ्य सिद्ध करता येते. अशा
सिद्धतेत \Log{G1c} मधील सिद्धता चे \Log{G3c} मधील
सिद्धता मध्ये आणि उलट दिशेने \emph{रूपांतरे} मिळतात.

\begin{prop}
$\Log{G1c} \Proves S$ असेल, तर $\Log{G3c} \Proves S$.
\end{prop}

\begin{proof}
  $\pi$ ही \Log{G1c} मधील सिद्धता असेल, तर \Log{G3c} मध्ये
  $S$ ची सिद्धता~$\pi'$ आहे, हे दाखवू. यासाठी $n
  = \pheight{\pi}$ वर विगमन करू.

  $n = 0$ असेल, तर $\pi$ हे स्वयंसिद्धक असते. $\lfalse \Sequent
  \quad$ हे स्वयंसिद्धक \Log{G3c} मध्येही आहे (संदर्भ $\Gamma,
  \Delta$ रिकामे घ्या). \Log{G1c} मधील $A \Sequent A$ ही
  स्वयंसिद्धके $A$ अण्वीय असण्याची अट न घालता ग्राह्य असतात.
  पण अशा सर्व क्रमवर्ती \Log{G3c} मध्ये सिद्धतायोग्य आहेत, हे आपण
  \ref{pt:seq:ex:prop:G3c-axioms} मध्ये सिद्ध केले आहे.

  $n>0$ असेल, तर $\pi$ चा शेवट $R$ नियमाने होतो. विगमनाचे
  गृहीतक सांगते की त्या नियमाच्या आधारविधानांच्या \Log{G3c}
  मधील सिद्धता आहेत; म्हणजे लगेचच्या उप-सिद्धता ची
  \Log{G3c} मध्ये रूपांतरे आहेत. $R$ हा \Log{G3c} मधीलही नियम
  असेल, तर आधारविधानांच्या रूपांतरित सिद्धता वर तो लागू करू.
  $R$ हा \Log{G1c} मधील $\LeftR{\land}$ असेल,
  तर आधारविधानात नसलेले दुसरे संयोजक सूत्र डावीकडे
  दुर्बलीकरणाने जोडा आणि \Log{G3c} चा $\LeftR{\land}$
  नियम लावा. $\RightR{\lor}$ साठी नसलेले दुसरे वियोजक
  सूत्र उजवीकडे जोडा आणि \Log{G3c} चा $\RightR{\lor}$
  नियम लावा. \Log{G1c} च्या $\LeftR{\lforall}$ किंवा
  $\RightR{\lexists}$ साठी त्याच बाजूला मुख्य संख्यापकीय
  सूत्र दुर्बलीकरणाने आधी जोडा; मग ते जपणारा \Log{G3c}
  नियम लावा. दुर्बलीकरणाची अनुज्ञेयता
  \ref{pt:seq:adm:prop:weak-G3c-adm} देते. $R$ स्वतः
  दुर्बलीकरणाचा नियम असेल, तर तीच प्रमेयिका वापरा;
  आकुंचनाचा नियम असेल, तर
  \ref{pt:seq:inv:prop:G3c-cont-adm} वापरा.
\end{proof}

\begin{prop}
$\Log{G3c} \Proves S$ असेल, तर $\Log{G1c} \Proves S$.
\end{prop}

\begin{proof}
  पुन्हा $S$ ची सिद्धता~$\pi$ हिच्या उंची वर विगमन करू.
  \Log{G3c} मधील प्रत्येक स्वयंसिद्धक \Log{G1c} मध्ये एका
  स्वयंसिद्धकापासून \LeftR{\Weakening} आणि
  \RightR{\Weakening} नियम काही वेळा वापरून सिद्ध करता येते:
  \[
  \Axiom$C \fCenter C$
  \RightLabel{\LeftR{\Weakening}}
  \Deduce$C, \Gamma \fCenter C$
  \RightLabel{\RightR{\Weakening}}
  \Deduce$C, \Gamma \fCenter \Delta, C$
  \DisplayProof
\qquad
  \Axiom$\lfalse \fCenter$
  \RightLabel{\LeftR{\Weakening}}
  \Deduce$\lfalse, \Gamma \fCenter$
  \RightLabel{\RightR{\Weakening}}
  \Deduce$\lfalse, \Gamma \fCenter \Delta$
  \DisplayProof
  \]
  $\pi$ चा शेवट $R$ नियमाने होत असेल, तर विगमनाच्या गृहीतकाने
  आधारविधानांच्या \Log{G1c} मधील सिद्धता मिळतात;
  नियम दोन्ही पद्धतींत समान असेल, तर त्या सिद्धता वर
  तोच $R$ नियम लागू करता येतो. $\LeftR{\lforall}$ आणि
  $\RightR{\lexists}$ मध्ये \Log{G3c} च्या आधारविधानात
  मुख्य संख्यापकीय सूत्रही जपलेले असते. ते संदर्भात ठेवून
  \Log{G1c} चा नियम लावल्यास निष्कर्षात मुख्य सूत्राच्या
  दोन प्रती मिळतात; अनुक्रमे डाव्या किंवा उजव्या आकुंचनाने
  त्या एक करता येतात. उरलेले अपवाद म्हणजे \LeftR{\land}
  आणि \RightR{\lor}; त्यांच्या \Log{G1c} आणि
  \Log{G3c} मधील आवृत्त्या वेगळ्या आहेत. \Log{G3c} मधील
  आवृत्त्या \Log{G1c} मध्ये पुढीलप्रमाणे व्युत्पाद्य आहेत:
  \[
  \Axiom$A, B, \Gamma \fCenter \Delta$
  \RightLabel{\LeftR{\land}}
  \UnaryInf$A \land B, B, \Gamma \fCenter \Delta$
  \RightLabel{\LeftR{\land}}
  \UnaryInf$A \land B, A \land B, \Gamma \fCenter \Delta$
  \RightLabel{\LeftR{\Contraction}}
  \UnaryInf$A \land B, \Gamma \fCenter \Delta$
    \DisplayProof
\qquad
  \Axiom$\Gamma \fCenter \Delta, A, B$
  \RightLabel{\RightR{\lor}}
  \UnaryInf$\Gamma \fCenter \Delta, A \lor B, B$
  \RightLabel{\RightR{\lor}}
  \UnaryInf$\Gamma \fCenter \Delta, A \lor B, A \lor B$
  \RightLabel{\RightR{\Contraction}}
  \UnaryInf$\Gamma \fCenter \Delta, A \lor B$
    \DisplayProof
  \]
\end{proof}
\part{इतिहास}\label{his:part}
\chapter{चरित्रे}\label{his:bio:chap}










\section{गेऑर्ग कँटर}\label{his:bio:can:sec}



गेऑर्ग कँटर (\textsc{गे}-ऑर्ग \textsc{कान}-टोर) यांच्या एका जुन्या चरित्रात असा दावा केला होता की रशियातील सेंट पीटर्सबर्गकडे जाणाऱ्या जहाजावर त्यांचा जन्म झाला आणि ते तिथे सापडले; त्यांचे पालक कोण होते हे माहीत नव्हते. मात्र हा दावा खरा नाही. कँटर यांचा जन्म १८४५ मध्ये सेंट पीटर्सबर्ग येथे झाला होता.

कँटर यांनी १८६७ मध्ये बर्लिन विद्यापीठातून गणितातील डॉक्टरेट मिळवली. संच उपपत्तीतील त्यांच्या कार्यासाठी ते ओळखले जातात आणि संशोधनाची स्वतंत्र शाखा म्हणून संच उपपत्तीची स्थापना करण्याचे श्रेय त्यांना दिले जाते. निरनिराळ्या आकारमानांचे अनंत संच असतात हे सर्वप्रथम त्यांनी सिद्ध केले. त्यांच्या उपपत्तींबद्दल, विशेषतः अनंतांच्या उपपत्तीबद्दल, तत्कालीन गणितज्ञांत बरीच चर्चा झाली आणि त्यांचे कार्य वादग्रस्त ठरले.

कँटर यांच्या धार्मिक श्रद्धा आणि गणितीय कार्य यांची घट्ट सांगड होती. अनंतपार संख्यांची उपपत्ती आपल्याला थेट ईश्वराकडून कळली, असा दावाही त्यांनी केला होता. आयुष्याच्या पुढील काळात त्यांना मानसिक आजाराचा त्रास झाला. १८८४ पासून त्यांना रुग्णालयात दाखल करावे लागले; नंतरच्या काळात असे अधिक वारंवार घडले. या स्रोताच्या मांडणीनुसार, त्यांच्या कार्यावरील तीव्र टीका आणि गणितज्ञ लिओपोल्ड क्रोनेकर यांच्याशी झालेला मतभेद यांमुळे त्यांना नैराश्य आले आणि गणितातील रस कमी झाला. नैराश्याच्या काळात कँटर तत्त्वज्ञान आणि साहित्याकडे वळत. शेक्सपियरची नाटके फ्रान्सिस बेकन यांनी लिहिली होती, असा सिद्धांतही त्यांनी प्रकाशित केला.

६ जानेवारी १९१८ रोजी हाल येथील आरोग्यधामात कँटर यांचे निधन झाले.

\paragraph{अधिक वाचन.}
कँटर यांच्या सविस्तर चरित्रांसाठी Dauben (1990) आणि Grattan-Guinness (1971) पाहा. त्यांच्या अपरंपरागत मतांचे वर्णन बीबीसी रेडिओ ४ वरील \emph{A Brief History of Mathematics} (Sautoy 2014) या कार्यक्रमातही आहे. कँटर यांच्या उपपत्ती रॅपच्या रूपात ऐकायच्या असतील, तर Rose (2012) पाहा.













\section{ॲलोन्झो चर्च}\label{his:bio:chu:sec}



ॲलोन्झो चर्च यांचा जन्म १४ जून १९०३ रोजी वॉशिंग्टन, डीसी येथे झाला. लहानपणी हवेच्या बंदुकीशी संबंधित एका अपघातामुळे त्यांचा एक डोळा आंधळा झाला. कनेक्टिकट येथील महाविद्यालयपूर्व शाळेचे शिक्षण त्यांनी १९२० मध्ये पूर्ण केले आणि त्याच वर्षी प्रिन्स्टन विद्यापीठात शिक्षण सुरू केले. १९२७ मध्ये त्यांनी डॉक्टरेट पूर्ण केली. दोन वर्षे परदेशात राहिल्यानंतर चर्च प्रिन्स्टनला परतले. ते अत्यंत नम्र आणि काळजीपूर्वक काम करणारे म्हणून ओळखले जात. फळ्यावरचे त्यांचे लेखन निर्दोष असे. महत्त्वाचे कागद जतन करण्यासाठी ते त्यांवर ड्यूको सिमेंट या पारदर्शक गोंदाचा पातळ थर लावत. शैक्षणिक कामाव्यतिरिक्त त्यांना विज्ञानकथा मासिके वाचायला आवडत; कथांमध्ये काही चूक आढळली तर संपादकांना पत्र लिहायलाही ते कचरत नसत.

चर्च यांचे शैक्षणिक योगदान मोठे होते. त्यांचे विद्यार्थी स्टीफन क्लिनी आणि बार्कली रॉसर यांच्यासह त्यांनी प्रभावी गणनीयतेची एक उपपत्ती, म्हणजे लॅम्डा कलन, विकसित केले. हे काम ॲलन ट्यूरिंग यांनी ट्यूरिंग यंत्र विकसित केले त्या कामापासून स्वतंत्रपणे झाले. गणनीयतेच्या या दोन व्याख्या सममूल्य आहेत. त्यांतून आज \emph{चर्च–ट्यूरिंग प्रबंध} म्हणून ओळखले जाणारे प्रतिपादन पुढे येते: नैसर्गिक संख्यांवरील एखादे फलन प्रभावीपणे गणनीय असते जर आणि फक्त जर ते ट्यूरिंग यंत्राने (किंवा लॅम्डा कलनाने) गणनीय असेल. चर्च यांनी आज \emph{चर्च यांचे प्रमेय} म्हणून ओळखला जाणारा निष्कर्षही सिद्ध केला: प्रथम-क्रम सूत्रांच्या वैधतेची निर्णय समस्या सोडवता येत नाही.

चर्च यांनी वृद्धापकाळापर्यंत काम चालू ठेवले. १९६७ मध्ये त्यांनी प्रिन्स्टन सोडून यूसीएलए येथे प्रवेश केला. १९९० मध्ये निवृत्त होईपर्यंत ते तेथे प्राध्यापक होते. ११ ऑगस्ट १९९५ रोजी वयाच्या ९२व्या वर्षी त्यांचे निधन झाले.

\paragraph{अधिक वाचन.}
चर्च यांचे संक्षिप्त चरित्र Enderton (2019) येथे पाहा. लॅम्डा कलन, चर्च यांचा प्रबंध आणि निर्णय समस्या यांवरील चर्च यांचे मूळ लेख Church (1936; 1936) येथे आहेत. १९३०च्या दशकातील प्रिन्स्टनच्या गणिती समुदायाविषयी चर्च यांची मुलाखत Aspray (1984) यांनी नोंदवली आहे. चर्च यांनी १९३६ ते १९७९ या काळात \emph{Journal of Symbolic Logic} साठी अनेक पुस्तक परीक्षणे लिहिली. ती सर्व जॉन मॅकफार्लेन यांच्या संकेतस्थळावर संग्रहित आहेत (MacFarlane 2015).













\section{गरहार्ड गेंटझेन}\label{his:bio:gen:sec}



गरहार्ड गेंटझेन हे मुख्यतः संरचनात्मक सिद्धता उपपत्तीचे निर्माते म्हणून, आणि विशेषतः नैसर्गिक निगमन व क्रमवर्ती कलन या निष्पत्ती-पद्धती निर्माण केल्याबद्दल ओळखले जातात. त्यांचा जन्म २४ नोव्हेंबर १९०९ रोजी जर्मनीतील ग्राइफ्सवाल्ड येथे झाला. महाविद्यालयपूर्व शाळेत जाण्यापूर्वी गरहार्ड यांचे तीन वर्षे घरीच शिक्षण झाले. शाळेत गेल्यावर शिक्षणाच्या बाबतीत ते आपल्या बहुतेक सहाध्यायांच्या मागे होते. तरीही ते उत्तम विद्यार्थी होते आणि गणितातील त्यांची विशेष क्षमता दिसून येत होती. त्यांची आवड विविध विषयांत होती. उदाहरणार्थ, त्यांनी आपल्या आईसाठी कविता आणि शाळेच्या नाट्यगृहासाठी नाटकेही लिहिली.

गेंटझेन यांनी विद्यापीठीय शिक्षणाची सुरुवात ग्राइफ्सवाल्ड विद्यापीठात केली; त्यानंतर ते गॉटिंगेन (G\"{o}ttingen), म्यूनिख आणि बर्लिन येथे गेले. १९३३ मध्ये त्यांनी गॉटिंगेन (G\"{o}ttingen) विद्यापीठातून हेरमान वाइल यांच्या मार्गदर्शनाखाली डॉक्टरेट मिळवली. (त्यांच्या बहुतांश कामाचे मार्गदर्शन पॉल बर्नायस यांनी केले होते; मात्र नाझींनी त्यांना विद्यापीठातून काढून टाकले.) १९३४ मध्ये गेंटझेन डेव्हिड हिल्बर्ट यांचे सहाय्यक म्हणून काम करू लागले. क्रमवर्ती कलन आणि नैसर्गिक निगमन या निष्पत्ती-पद्धतींवरील त्यांचे \emph{Untersuchungen \"{u}ber das logische Schlie\ss en I--II
  [Investigations Into Logical Deduction I--II]} हे लेख १९३५ मध्ये प्रकाशित झाले. १९३६ मध्ये त्यांनी पेआनो स्वयंसिद्धकांची सुसंगती सिद्ध केली.

गेंटझेन यांचा नाझींशी असलेला संबंध गुंतागुंतीचा होता. त्यांचे मार्गदर्शक बर्नायस यांना जर्मनी सोडण्यास भाग पाडले जात असतानाच गेंटझेन नाझींची निमलष्करी संघटना असलेल्या एसएच्या विद्यापीठीय शाखेत सामील झाले. अनेक जर्मनांप्रमाणे ते नाझी पक्षाचेही सदस्य होते. युद्धकाळात त्यांनी हवाई गुप्तचर विभागात दूरसंचार अधिकारी म्हणून काम केले. मात्र १९४२ मध्ये मानसिक आरोग्य अचानक खालावल्याने त्यांना सेवेतून मुक्त करण्यात आले. गेंटझेन यांची निष्ठा नाझी पक्षाशी होती की शैक्षणिक प्रगती सुनिश्चित करण्यासाठी ते पक्षात सामील झाले, हे स्पष्ट नाही.

१९४३ मध्ये गेंटझेन यांना प्राग येथील जर्मन विद्यापीठाच्या गणित संस्थेत शैक्षणिक पद देऊ करण्यात आले; त्यांनी ते स्वीकारले. मात्र १९४५ मध्ये प्रागच्या नागरिकांनी जर्मन अधिपत्याविरुद्ध बंड केले. सोव्हिएत सैन्य शहरात दाखल झाले आणि विद्यापीठातील सर्व प्राध्यापकांना अटक केली, असे या स्रोतामध्ये म्हटले आहे. नाझी संघटनांतील सदस्यत्वामुळे गेंटझेन यांना सक्तीच्या श्रमांच्या छावणीत नेण्यात आले. ४ ऑगस्ट १९४५ रोजी वयाच्या ३५व्या वर्षी कारागृहातील त्यांच्या कोठडीत कुपोषणामुळे त्यांचे निधन झाले.

\paragraph{अधिक वाचन.}
गेंटझेन यांचे संपूर्ण चरित्र Menzler-Trott (2007) येथे पाहा. नाझी राजवटीखालील गणितज्ञांविषयीचे Segal (2014) यांचे पुस्तकही वाचनीय आहे; त्यात गेंटझेन यांच्या आयुष्याविषयी एक संक्षिप्त नोंद आहे. तार्किक निगमनावरील गेंटझेन यांचे मूळ जर्मन लेख (Gentzen 1935; Gentzen 1935) येथे उपलब्ध आहेत. त्यांच्या लेखांची इंग्रजी भाषांतरे Szabo (1969) यांनी एका खंडात एकत्र केली आहेत; त्यात त्यांचे संक्षिप्त चरित्रही आहे.














\section{कुर्ट गोडेल}\label{his:bio:god:sec}



कुर्ट गोडेल (\textsc{गर}-डल) यांचा जन्म २८ एप्रिल १९०६ रोजी ऑस्ट्रो-हंगेरी साम्राज्यातील ब्र्युन येथे झाला; ते ठिकाण आता झेक प्रजासत्ताकातील ब्र्नो आहे. लहानपणी ते जिज्ञासू आणि हुशार असल्याने घरचे त्यांना “Der kleine Herr Warum” (लहानसे 'का-महोदय') असे म्हणत. प्राथमिक शाळेपासूनच त्यांची शैक्षणिक कामगिरी उत्कृष्ट होती; सर्वोच्च गुणांपेक्षा कमी गुण त्यांना फक्त गणितात मिळाले. तब्येत बरी नसल्याने गोडेल अनेकदा शाळेत गैरहजर राहत आणि त्यांना शारीरिक शिक्षणातून सूट मिळाली होती. बालपणी त्यांना संधिवातजन्य तापाचे निदान झाले. वैद्यकीय तपासणीत त्याच्या उलट निष्कर्ष आला असतानाही, त्या आजाराचा आपल्या हृदयावर कायमचा परिणाम झाला आहे, असे त्यांना आयुष्यभर वाटत राहिले.

गोडेल यांनी १९२४ मध्ये व्हिएन्ना विद्यापीठात शिक्षण सुरू केले आणि १९२९ मध्ये डॉक्टरेटचा प्रबंध पूर्ण केला; डॉक्टरेटची पदवी १९३० मध्ये मिळाली. सुरुवातीला त्यांचा भौतिकशास्त्र शिकण्याचा विचार होता; पण काही अंशी तत्त्वज्ञ रुडॉल्फ कार्नाप यांच्या प्रभावामुळे त्यांची आवड गणिताकडे, विशेषतः तर्कशास्त्राकडे वळली. हान्स हान यांच्या मार्गदर्शनाखाली लिहिलेल्या प्रबंधात त्यांनी समानतेसह प्रथम-क्रम विधेय तर्कशास्त्राचे संपूर्णता प्रमेय सिद्ध केले (G\"odel 1929). या प्रबंधानंतर अवघ्या वर्षभरात त्यांनी आपले सर्वात प्रसिद्ध निष्कर्ष, पहिले आणि दुसरे अपूर्णता प्रमेय, मिळवले (हे G{\"{o}}del 1931 मध्ये प्रसिद्ध झाले). व्हिएन्नातील वास्तव्यात गोडेल 'व्हिएन्ना मंडळ' या विज्ञानाभिमुख तत्त्वज्ञांच्या समूहात सक्रिय होते. या समूहातील कार्नाप यांच्या कार्यावर गोडेल यांच्या निष्कर्षांचा विशेष प्रभाव पडला.

१९३८ मध्ये गोडेल यांनी अडेल निंबुर्स्की यांच्याशी विवाह केला. त्यांचे पालक या विवाहावर नाराज होते: अडेल त्यांच्यापेक्षा सहा वर्षांनी मोठ्या होत्या, त्यांचा आधी घटस्फोट झाला होता आणि त्या एका नाइटक्लबमध्ये नर्तकी म्हणून काम करत होत्या. मात्र सामाजिक दबावाचा गोडेल यांच्यावर परिणाम झाला नाही; त्यांच्या निधनापर्यंत त्यांचा विवाह आनंदी राहिला, असे स्रोत सांगतो.

१९३८ मध्ये नाझी जर्मनीने ऑस्ट्रियाचा विलय केल्यानंतर गोडेल आणि अडेल अमेरिकेत स्थलांतरित झाले. न्यू जर्सीतील प्रिन्स्टन येथील इन्स्टिट्यूट फॉर ॲडव्हान्स्ड स्टडी येथे गोडेल यांनी पद स्वीकारले. ते अंतर्मुख आणि काहीसे विलक्षण स्वभावाचे असले, तरी प्रिन्स्टनमधील त्यांचा काळ सहकार्यपूर्ण आणि फलदायी ठरला. त्यांनी संच उपपत्ती, तत्त्वज्ञान आणि भौतिकशास्त्र यांवर निबंध प्रकाशित केले. याच संस्थेतील सहकारी अल्बर्ट आइन्स्टाइन यांच्याशी त्यांची विशेष घनिष्ठ मैत्री झाली.

आयुष्याच्या उत्तरार्धात गोडेल यांचे मानसिक आरोग्य खालावले. १९७७ मध्ये त्यांच्या पत्नीला रुग्णालयात दाखल करावे लागल्याने त्या त्यांच्यासाठी जेवण बनवू शकत नव्हत्या. गोडेल यांना आयुष्यभर मानसिक आरोग्याच्या समस्या होत्या; त्यातून त्यांच्या मनात संशय बळावला. आपल्याला विष घातले जाईल या तीव्र भीतीने त्यांनी अन्न घेणे बंद केले. १४ जानेवारी १९७८ रोजी प्रिन्स्टन येथे उपासमारीने त्यांचे निधन झाले, असे या स्रोताचे वर्णन आहे.

\paragraph{अधिक वाचन.}
गोडेल यांचे सविस्तर चरित्र John Dawson (1997) येथे पाहा. त्यांच्या चरित्राविषयीचे आणखी लेख आणि तर्कशास्त्र व तत्त्वज्ञानातील योगदानावरील निबंध Wang (1990), Baaz (2011), Takeuti (2003) आणि Sigmund (2007) येथे मिळतील.

गोडेल यांचा डॉक्टरेट प्रबंध मूळ जर्मन भाषेत (G\"odel 1929) येथे उपलब्ध आहे. अपूर्णता प्रमेयांचा मूळ लेख (G{\"{o}}del 1931) येथे आहे. त्यांचे प्रकाशित व अप्रकाशित सर्व लेखन आणि निवडक पत्रव्यवहार यांची इंग्रजी आवृत्ती त्यांच्या \emph{Collected Papers} Feferman (1986; 1990) मध्ये उपलब्ध आहे.

गोडेल यांच्या अपूर्णता प्रमेयांच्या सविस्तर विवेचनासाठी Smith (2013) पाहा. या प्रमेयांवरील अनौपचारिक तात्त्विक चर्चेसाठी मार्क लिन्झेनमायर यांचे पॉडकास्ट (Linsenmayer 2014) पाहा.













\section{एमी नोएथर}\label{his:bio:noe:sec}



एमी नोएथर (\textsc{नर}-टर) यांचा जन्म २३ मार्च १८८२ रोजी जर्मनीतील एरलांगेन येथे एका उच्च मध्यमवर्गीय विद्वान कुटुंबात झाला. त्यांना “आधुनिक बीजगणिताची जननी” म्हटले जाते. स्त्रियांच्या शिक्षणातील मोठ्या अडथळ्यांवर मात करत त्यांनी गणित आणि भौतिकशास्त्र या दोन्ही क्षेत्रांत मूलगामी योगदान दिले. त्या काळी जर्मनीतील मुलींनी कलांचे शिक्षण घ्यावे अशी अपेक्षा होती आणि त्यांना महाविद्यालयपूर्व शाळांमध्ये प्रवेश नव्हता. तरीही नोएथर यांनी गॉटिंगेन आणि एरलांगेन विद्यापीठांतील वर्गांत श्रोती म्हणून उपस्थिती लावली; एरलांगेन येथे त्यांचे वडील गणिताचे प्राध्यापक होते. अखेरीस १९०४ मध्ये महिला विद्यार्थिनींना प्रवेश देण्याचे एरलांगेनचे धोरण बदलल्यावर त्या तेथे विद्यार्थिनी म्हणून दाखल झाल्या. १९०७ मध्ये त्यांनी गणितात डॉक्टरेट मिळवली.

पात्रता असूनही नोएथर यांना कारकिर्दीत मोठा विरोध सहन करावा लागला. १९०८ ते १९१५ या काळात त्यांनी एरलांगेन येथे विनावेतन अध्यापन केले. याच काळात तत्कालीन अग्रगण्य गणितज्ञ डेव्हिड हिल्बर्ट यांचे त्यांच्याकडे लक्ष गेले आणि त्यांनी नोएथर यांना गॉटिंगेनला बोलावले. मात्र स्त्रियांना प्राध्यापकपदे मिळण्यास बंदी होती. त्यामुळे त्यांना हिल्बर्ट यांच्या नावाखालीच व्याख्याने देता आली, तीही विनावेतन. याच काळात त्यांनी आज 'नोएथर यांचे प्रमेय' म्हणून ओळखला जाणारा निष्कर्ष सिद्ध केला; सैद्धांतिक भौतिकशास्त्रात आजही त्याचा उपयोग होतो. अखेरीस १९१९ मध्ये त्यांना अध्यापनाचा अधिकार मिळाला. विद्यापीठातील सहकाऱ्यांचा विरोध सुरूच असताना हिल्बर्ट यांनी “महाशयांनो, विद्याशाखेची सभा स्नानगृह नव्हे,” असे उत्तर दिल्याचे सांगितले जाते.

१९२०च्या दशकाच्या उत्तरार्धात नोएथर यांनी अमूर्त बीजगणितावर लक्ष केंद्रित केले आणि त्यांच्या योगदानाने या क्षेत्रात आमूलाग्र बदल घडवला. आपल्या सिद्धतांमध्ये त्या अनेकदा तथाकथित आरोही साखळी अट वापरत. या अटीनुसार विशिष्ट संचांची काटेकोरपणे वाढत जाणारी अनंत साखळी नसते. उदाहरणार्थ, आता नोएथरियन वलये म्हणून ओळखल्या जाणाऱ्या काही बीजगणितीय रचनांत आदर्शांचा $I_1
\subsetneq I_2 \subsetneq \dots$ असा अनंत अनुक्रम नसतो. ही अट कोणत्याही अंशतः क्रमावर व्यापक करता येते (बीजगणितात आदर्शांना उपसंच-संबंधाने लावलेला क्रम हे तिचे विशेष प्रकरण आहे). तिची द्वैती अवरोही साखळी अटही विचारात घेता येते: अंशतः क्रमातील प्रत्येक काटेकोरपणे \emph{उतरता} अनुक्रम अखेरीस थांबतो. एखादा अंशतः क्रम ही अवरोही साखळी अट पूर्ण करत असेल, तर $<$ या $\Nat$ वरील क्रमाने जसे विगमन करता येते, तसे त्या क्रमानेही विगमन करता येते. अशा क्रमांना \emph{सुस्थापित} किंवा \emph{नोएथरियन} म्हणतात; त्यानुसार मिळणाऱ्या सिद्धता-तत्त्वाला \emph{नोएथरियन विगमन} म्हणतात.

नोएथर ज्यू होत्या. १९३३ मध्ये नाझी सत्तेवर आल्यावर त्यांना पदावरून काढून टाकण्यात आले. सुदैवाने त्यांना अमेरिकेतील पेनसिल्व्हेनिया राज्यातील ब्रिन मॉर येथे तात्पुरत्या पदासाठी स्थलांतर करता आले. त्या काळात त्यांनी प्रिन्स्टनमध्येही व्याख्याने दिली; मात्र ते विद्यापीठ स्त्रियांना अनुकूल नाही, असे त्यांना वाटले (Dick 1981, पृ.~81). १९३५ मध्ये त्यांच्या गर्भाशयातील गाठ काढण्यासाठी शस्त्रक्रिया झाली. शस्त्रक्रियेनंतर झालेल्या संसर्गामुळे त्यांचे निधन झाले आणि ब्रिन मॉर येथे त्यांचे दफन झाले.

\paragraph{अधिक वाचन.}
नोएथर यांच्या चरित्रासाठी Dick (1981) पाहा. पेरिमीटर इन्स्टिट्यूट फॉर थिअरेटिकल फिजिक्सची त्यांच्या जीवनावरील आणि प्रभावावरील व्याख्याने ऑनलाइन उपलब्ध आहेत (Institute 2015). वाचनाऐवजी ऐकायचे असल्यास, \emph{Stuff You Missed in History Class} या मालिकेतील नोएथर यांच्या जीवनावरील आणि प्रभावावरील पॉडकास्ट (Frey 2015) पाहा. नोएथर यांचे संकलित लेखन मूळ जर्मन भाषेत (Jacobson 1983) येथे उपलब्ध आहे.













\section{रोझा पेटर}\label{his:bio:pet:sec}



रोझा पेटर यांचा जन्म १७ फेब्रुवारी १९०५ रोजी हंगेरीतील बुडापेस्ट येथे रोझा पोलित्झर या नावाने झाला. पुनरावर्ती फलनांवरील कार्यासाठी त्या विशेष ओळखल्या जातात; पुनरावर्तन उपपत्ती हे क्षेत्र निर्माण होण्यात हे कार्य अत्यावश्यक ठरले.

पेटर यांचे बालपण राजकीय उलथापालथीच्या काळात गेले; त्या किशोरवयीन असताना पहिले महायुद्ध सुरू होते. तरीही त्यांना बुडापेस्टमधील प्रतिष्ठित मारिया तेरेझिया मुलींच्या शाळेत शिकता आले आणि १९२२ मध्ये त्यांनी तेथील शिक्षण पूर्ण केले. त्यानंतर त्यांनी बुडापेस्टच्या पाझमान्य पेटर विद्यापीठात शिक्षण घेतले; पुढे त्याचे नाव लोरांड एटव्होश विद्यापीठ झाले. वडिलांच्या आग्रहामुळे त्यांनी सुरुवातीला रसायनशास्त्र घेतले, पण नंतर गणिताकडे वळल्या आणि १९२७ मध्ये पदवीधर झाल्या. माध्यमिक शाळेत गणित शिकवण्याची पात्रता असूनही त्या काळची आर्थिक परिस्थिती भीषण होती; महामंदीने जागतिक अर्थव्यवस्था प्रभावित झाली होती. म्हणून त्यांनी शिकवण्या आणि खाजगी गणित अध्यापन अशी प्रसंगोपात्त कामे केली. अखेरीस गणितातील पदव्युत्तर शिक्षणासाठी त्या पुन्हा विद्यापीठात आल्या. सुरुवातीला त्यांचा संख्या उपपत्तीवर काम करण्याचा विचार होता; पण आपले निष्कर्ष आधीच सिद्ध झाले आहेत हे कळल्यावर त्यांनी गणितच सोडून देण्याचा जवळजवळ निर्णय घेतला. त्यांना गोडेल यांच्या अपूर्णता प्रमेयांवर काम करण्यास प्रोत्साहन मिळाले आणि नकळत त्यांनी त्यांचे काही निष्कर्ष वेगळ्या पद्धतीने सिद्ध केले. त्यामुळे त्यांचा आत्मविश्वास परत आला. डेव्हिड हिल्बर्ट यांच्या पायाभूत कार्यक्रमातून प्रेरणा घेऊन पेटर यांनी पुनरावर्तन उपपत्तीवरील आपले पहिले लेख लिहिले. १९३५ मध्ये त्यांनी डॉक्टरेट मिळवली आणि १९३७ मध्ये त्या \emph{Journal of
  Symbolic Logic} च्या संपादक झाल्या.

पेटर यांचे आरंभीचे लेख पुनरावर्ती फलनांच्या उपपत्तीचे संस्थापक योगदान म्हणून व्यापकपणे मानले जातात. (P{\'e}ter 1935a) मध्ये त्यांनी पुनरावर्तनाच्या विविध प्रकारांमधील संबंध तपासले. (P\'eter 1935b) मध्ये त्यांनी पुनरावर्ती पद्धतीने परिभाषित केलेले एक विशिष्ट फलन आदिम पुनरावर्ती नाही, हे दाखवले. यामुळे विल्हेल्म ॲकरमन यांचा आधीचा निष्कर्ष सोपा झाला. पेटर यांचे हे सुलभ केलेले फलन आज अनेकदा ॲकरमन फलन म्हणतात; काही वेळा अधिक अचूकपणे ॲकरमन–पेटर फलनही म्हणतात. पुनरावर्ती फलनांच्या उपपत्तीवरील पहिले पुस्तक त्यांनी लिहिले (P\'eter 1951).

त्यांच्या कार्याचे महत्त्व आणि प्रभाव मोठा असूनही पेटर यांना १९४५ पर्यंत पूर्णवेळ अध्यापनाचे पद मिळाले नाही. दुसऱ्या महायुद्धात नाझींनी हंगेरीवर कब्जा केला असताना यहुदीविरोधी कायद्यांमुळे त्यांना शिकवण्याची परवानगी नव्हती. १९४४ मध्ये सरकारने बुडापेस्टमध्ये ज्यू लोकांसाठी बंदिस्त वस्ती निर्माण केली; ती शहराच्या इतर भागांपासून अलग ठेवली गेली होती आणि सशस्त्र रक्षक तिच्यावर पहारा देत होते. १९४५ मध्ये ही वस्ती मुक्त होईपर्यंत पेटर यांना तेथे राहण्यास भाग पाडले गेले. त्यानंतर त्यांनी बुडापेस्ट टीचर्स ट्रेनिंग कॉलेजमध्ये आणि १९५५ पासून एटव्होश लोरांड विद्यापीठात अध्यापन केले. ‘गणितातील अकादमिक डॉक्टर’ हा दर्जा मिळवणाऱ्या त्या पहिल्या हंगेरियन महिला गणितज्ञ आणि हंगेरियन विज्ञान अकादमीवर निवडून जाणाऱ्या पहिल्या महिला होत्या.

पेटर या गणित उत्कटतेने शिकवणाऱ्या अध्यापिका म्हणून ओळखल्या जात. केवळ व्याख्यान देण्याऐवजी विद्यार्थ्यांसोबत गणिती समस्यांचे स्वरूप आणि सौंदर्य शोधणे त्यांना अधिक आवडत असे. म्हणून विद्यार्थी त्यांना प्रेमाने “रोझा मावशी” म्हणत. १९७७ मध्ये वयाच्या ७१व्या वर्षी त्यांचे निधन झाले.

\paragraph{अधिक वाचन.}
अधिक चरित्रात्मक वाचनासाठी (O'Connor 2014) आणि (Andr\'asfai 1986) पाहा. पेटर यांची एक छोटी मुलाखत Tamassy (1994) यांनी घेतली होती. गणिताविषयी रंजक वाचनासाठी पेटर यांचे \emph{Playing With Infinity} हे पुस्तक (P\'eter 2010) पाहा.












\section{ज्युलिया रॉबिन्सन}\label{his:bio:rob:sec}



ज्युलिया बोमन रॉबिन्सन या अमेरिकन गणितज्ञ होत्या. त्या मुख्यतः निर्णय समस्यांवरील कार्यासाठी, आणि विशेषतः हिल्बर्ट यांच्या दहाव्या समस्येच्या समाधानातील योगदानासाठी ओळखल्या जातात. रॉबिन्सन यांचा जन्म ८ डिसेंबर १९१९ रोजी मिसुरी राज्यातील सेंट लुईस येथे झाला. लहानपणापासूनच संख्यांविषयी कुतूहल होते, अशी त्यांची आठवण आहे (Reid 1986, पृ.~4). वयाच्या नवव्या वर्षी त्यांना स्कार्लेट फीव्हर झाला आणि नंतर संधिवातजन्य तापाचे अनेक झटके आले. त्यामुळे त्यांना बराच वेळ अंथरुणात घालवावा लागला आणि शिक्षणात त्या मागे पडल्या. खाजगी शिक्षकांच्या मदतीने त्यांनी ही तूट भरून काढली; मात्र आजाराचे शारीरिक परिणाम त्यांच्या पुढील आयुष्यावरही राहिले.

बालपणीच्या अडचणी असूनही रॉबिन्सन यांनी गणित आणि विज्ञानातील अनेक पुरस्कारांसह माध्यमिक शिक्षण पूर्ण केले. विद्यापीठीय शिक्षणाची सुरुवात त्यांनी सॅन डिएगो स्टेट कॉलेजमध्ये केली आणि शेवटच्या वर्षी कॅलिफोर्निया विद्यापीठ, बर्कली येथे प्रवेश घेतला. तेथे गणितज्ञ राफाएल रॉबिन्सन यांचा त्यांच्यावर प्रभाव पडला. दोघांची मैत्री झाली आणि १९४१ मध्ये त्यांनी विवाह केला. विद्याशाखेतील सदस्याची पत्नी असल्याने ज्युलिया यांना बर्कलीच्या गणित विभागात अध्यापनाची परवानगी नव्हती. त्या गणिताचे वर्ग श्रोती म्हणून घेत राहिल्या, पण विद्यापीठ सोडून कुटुंब सुरू करण्याची त्यांची इच्छा होती. विवाहानंतर थोड्याच काळात त्यांना न्यूमोनिया झाला. बालपणीच्या संधिवातजन्य तापामुळे त्यांच्या हृदयावर बरीच व्रणयुक्त ऊतकांची वाढ झाली आहे, असे त्यांना सांगण्यात आले. या स्थितीची तीव्रता पाहून डॉक्टरांनी त्या चाळीस वर्षांपलीकडे जगणार नाहीत असा अंदाज व्यक्त केला आणि त्यांना मुले न होऊ देण्याचा सल्ला दिला (Reid 1986, पृ.~13).

रॉबिन्सन दीर्घकाळ नैराश्यात होत्या; अखेरीस त्यांनी गणिताचे शिक्षण सुरू ठेवायचे ठरवले. त्या बर्कलीला परतल्या आणि ॲल्फ्रेड टार्स्की यांच्या मार्गदर्शनाखाली १९४८ मध्ये डॉक्टरेट पूर्ण केली. वास्तविक संख्यांची प्रथम-क्रम उपपत्ती निर्णेय आहे, हे टार्स्की यांनी दाखवले होते; गोडेल यांच्या कार्यावरून नैसर्गिक संख्यांची प्रथम-क्रम उपपत्ती अनिर्णेय आहे, हे निष्पन्न होत होते. परिमेय संख्यांची प्रथम-क्रम उपपत्ती निर्णेय आहे की नाही, हा मोठा अनिर्णित प्रश्न होता. आपल्या प्रबंधात (1949) रॉबिन्सन यांनी ती अनिर्णेय आहे, हे सिद्ध केले.

निर्णय समस्यांत रस असल्याने रॉबिन्सन यांनी पुढे हिल्बर्ट यांच्या दहाव्या समस्येचे समाधान शोधण्याचा प्रयत्न केला. १९०० मध्ये डेव्हिड हिल्बर्ट यांनी मांडलेल्या २३ प्रसिद्ध गणिती समस्यांपैकी ही एक होती. पूर्णांक गुणांक असलेल्या $3x^2 - 2y +3 = 0$ यांसारख्या बहुपदी समीकरणाला पूर्णांकांमध्ये उत्तर आहे की नाही, हे मर्यादित वेळेत सांगणारा कलनविधी आहे का, असा दहाव्या समस्येचा प्रश्न आहे. अशा प्रश्नांना \emph{डायोफँटाइन समस्या} म्हणतात. काही सुरुवातीच्या यशांनंतर रॉबिन्सन यांनी याच समस्येवर काम करणारे मार्टिन डेव्हिस आणि हिलरी पुटनम यांच्यासोबत काम सुरू केले. अज्ञात राशी घातांकांतही येऊ शकतात अशा घातांकीय डायोफँटाइन समस्या अनिर्णेय आहेत, हे त्यांनी दाखवले. तसेच पुढे “J.R.” म्हणून ओळखल्या गेलेल्या एका विशिष्ट गृहीतकावरून हिल्बर्ट यांची दहावी समस्या अनिर्णेय आहे हे निष्पन्न होते, असेही त्यांनी सिद्ध केले (Davis 1961). १९६०च्या दशकातही रॉबिन्सन या समस्येवर काम करत राहिल्या. अखेरीस १९७० मध्ये तरुण रशियन गणितज्ञ युरी मातियासेविच यांनी J.R. गृहीतक सिद्ध केले. एकत्रित निष्कर्षाला आता मातियासेविच–रॉबिन्सन–डेव्हिस–पुटनम प्रमेय, संक्षेपात MRDP प्रमेय, म्हणतात. मातियासेविच आणि रॉबिन्सन यांची मैत्री झाली आणि त्यांनी अनेक लेखांवर एकत्र काम केले. मातियासेविच यांना लिहिलेल्या एका पत्रात रॉबिन्सन यांनी असे म्हटले होते: “प्रत्यक्षात, हजारो मैल दूर राहूनही एकत्र काम केल्यामुळे आपल्यापैकी कोणीही एकट्याने जितकी प्रगती केली असती त्याहून अधिक प्रगती आपण करत आहोत, याचा मला खूप आनंद आहे” (Matijasevich 1992, पृ.~45).

रॉबिन्सन या अमेरिकन मॅथेमॅटिकल सोसायटीच्या पहिल्या महिला अध्यक्ष होत्या. राष्ट्रीय विज्ञान अकादमीच्या गणित विभागात निवडून जाणाऱ्या त्या पहिल्या महिला होत्या. रक्ताच्या कर्करोगाचे निदान झाल्यानंतर ३० जुलै १९८५ रोजी वयाच्या ६५व्या वर्षी त्यांचे निधन झाले.

\paragraph{अधिक वाचन.}
रॉबिन्सन यांचे गणितीय लेख त्यांच्या \textit{Collected Works} (Robinson 1996) मध्ये उपलब्ध आहेत. त्यात राष्ट्रीय विज्ञान अकादमीसाठी लिहिलेल्या त्यांच्या चरित्रात्मक स्मृतिलेखाचे पुनर्मुद्रणही आहे (Feferman 1994). रॉबिन्सन यांची मोठी बहीण कॉन्स्टन्स रीड यांनी मुलाखतींवर आधारित “Autobiography of Julia” (Reid 1986) आणि त्यांचे सविस्तर चरित्र (Reid 1996) प्रकाशित केले. रॉबिन्सन आणि हिल्बर्ट यांच्या दहाव्या समस्येवरील लघुपटाचे दिग्दर्शन जॉर्ज चिचेरी यांनी केले (Csicsery 2016). मातियासेविच यांचे रॉबिन्सन यांच्यासोबतचे सहकार्य आणि त्यांच्या कार्यावरील रॉबिन्सन यांचा प्रभाव यांविषयीच्या संक्षिप्त स्मृतिलेखासाठी (Matijasevich 1992) पाहा.














\section{बर्ट्रंड रसेल}\label{his:bio:rus:sec}



बर्ट्रंड रसेल यांना आधुनिक विश्लेषणात्मक तत्त्वज्ञानाच्या संस्थापकांपैकी एक मानले जाते. १८ मे १८७२ रोजी जन्मलेले रसेल तत्त्वज्ञान आणि तर्कशास्त्रातील कामासाठी तर ओळखले जातच; याशिवाय त्यांनी अनेक विषयांवर सर्वसामान्य वाचकांसाठी पुस्तके लिहिली. आयुष्यभर ते सक्रिय राजकीय कार्यकर्तेही होते.

रसेल यांचा जन्म वेल्समधील मॉनमथशरच्या ट्रेलेच येथे झाला. त्यांचे आईवडील ब्रिटिश उमराव वर्गातील होते. ते स्वतंत्र विचारांचे होते आणि त्या काळी बोस्टनमधील क्रांतिकारी विचारांच्या लोकांशीही त्यांची मैत्री होती. दुर्दैवाने रसेल लहान असतानाच त्यांच्या आईवडिलांचे निधन झाले; म्हणून ते आजीआजोबांकडे राहायला गेले. तेथे त्यांचे धार्मिक संस्कारांत संगोपन झाले, जे त्यांच्या आईवडिलांना कोणत्याही परिस्थितीत टाळायचे होते. नैतिकतेच्या सर्व बाबतीत त्यांची आजी अतिशय कडक होती. किशोरवयात त्यांचे बहुतांश शिक्षण खाजगी शिक्षकांकडून घरीच झाले.

विश्लेषणात्मक तत्त्वज्ञानावर, विशेषतः तर्कशास्त्रावर, रसेल यांचा प्रभाव फार मोठा आहे. त्यांनी केंब्रिजच्या ट्रिनिटी कॉलेजमध्ये गणित आणि तत्त्वज्ञान शिकले; तेथे गणितज्ञ आणि तत्त्वज्ञ ॲल्फ्रेड नॉर्थ व्हाइटहेड यांचा त्यांच्यावर प्रभाव पडला. १९१० मध्ये रसेल आणि व्हाइटहेड यांनी \emph{Principia Mathematica} चा पहिला खंड प्रकाशित केला. त्यात गणिताचे तर्कशास्त्रात न्यूनीकरण करता येते, या भूमिकेचे त्यांनी समर्थन केले. पुढे रसेल यांनी शेकडो पुस्तके, निबंध आणि राजकीय पत्रके प्रकाशित केली. १९५० मध्ये त्यांना साहित्यातील नोबेल पारितोषिक मिळाले.

रसेल राजकारणात आणि सामाजिक चळवळींत खोलवर सहभागी होते. पहिल्या महायुद्धात शांततावादी कृती आणि निषेध यांमुळे त्यांना अटक झाली आणि सहा महिन्यांच्या तुरुंगवासाची शिक्षा झाली. तुरुंगात त्यांना वाचता आणि लिहिता येत असे; हा अनुभव “बराच सुखावह” वाटल्याचा त्यांचा दावा होता. पुढेही त्यांनी शांततेसाठी काम केले, मात्र सर्व परिस्थितींमध्ये युद्धाला विरोध करावा अशी त्यांची भूमिका नव्हती. १९६१ मध्ये अण्वस्त्रनिःशस्त्रीकरणाच्या सभेत सहभागी झाल्यामुळे त्यांना पुन्हा तुरुंगवास झाला. १९४८ मध्ये ते विमान अपघातातूनही बचावले. या स्रोताच्या कथनानुसार, तेव्हा धूम्रपानासाठी राखीव विभागात बसलेले लोकच बचावले; म्हणून धूम्रपानामुळेच आपला जीव वाचला, असा रसेल यांचा दावा होता. रसेल यांनी चार विवाह केले, पण विवाहबाह्य संबंध ठेवत असल्याची त्यांची ख्याती होती. २ फेब्रुवारी १९७० रोजी वयाच्या ९७व्या वर्षी वेल्समधील पेनरिनड्यूड्रेथ येथे त्यांचे निधन झाले.

\paragraph{अधिक वाचन.}
रसेल यांनी १८७२ ते १९६७ या कालखंडातील आपले जीवन सांगणारे तीन भागांचे आत्मचरित्र लिहिले (Russell 1967; Russell 1968; Russell 1969). मॅकमास्टर विद्यापीठातील बर्ट्रंड रसेल संशोधन केंद्रात त्यांच्या संग्रहित दस्तऐवजांचा ठेवा आहे. त्यांच्या संकलित लेखनाचे खंड, शोधता येणाऱ्या अनुक्रमणिका आणि संग्रहण प्रकल्पांविषयी केंद्राच्या संकेतस्थळावर Duncan (2015) माहिती आहे. रसेल यांचा \emph{On
  Denoting} (Russell 1905) हा लेख विसाव्या शतकातील विश्लेषणात्मक तत्त्वज्ञानातील अभिजात लेख मानला जातो.

रसेल यांच्यावरील स्टॅनफर्ड एन्सायक्लोपीडिया ऑफ फिलॉसॉफीमधील नोंदीत (Irvine 2015) इच्छा आणि राजकीय सिद्धांताविषयी रसेल बोलत असल्याच्या ध्वनिमुद्रिका आहेत. रसेल यांच्या अनेक चित्रित मुलाखती ऑनलाइन उपलब्ध आहेत. उदाहरणार्थ, धूम्रपान आणि विमान अपघातातील सहभागाविषयी त्यांना बोलताना पाहण्यासाठी Russell (n.d.) पाहा. त्यांच्या काही पुस्तकांची, त्यांत \emph{Introduction to Mathematical Philosophy} चा समावेश आहे, विनामूल्य ध्वनिपुस्तके LibriVox (n.d.) येथे उपलब्ध आहेत.













\section{ॲल्फ्रेड टार्स्की}\label{his:bio:tar:sec}



ॲल्फ्रेड टार्स्की यांचा जन्म १४ जानेवारी १९०१ रोजी पोलंडमधील वॉर्सा येथे झाला; तेव्हा ते शहर रशियन साम्राज्याचा भाग होते. टार्स्की यांचे वर्णन “नेपोलियनसारखे” असे केले जात असे: ते उत्साही, बोलके आणि तीव्र स्वभावाचे होते. त्यांची ऊर्जा व्याख्यानांतही दिसत असे. एकदा व्याख्यानादरम्यान सिगारेट टाकताना त्यांनी कचरापेटीला आग लावली आणि त्यानंतर त्या इमारतीत त्यांना पुन्हा व्याख्यान देण्यास बंदी घालण्यात आली.

लहानपणापासूनच टार्स्की यांना ज्ञानाची तीव्र ओढ होती. आयुष्याच्या पुढील काळात ते विद्यार्थ्यांना सांगत की तर्कशास्त्रात त्यांना 'B' गुण मिळाले म्हणून त्यांनी तो विषय शिकला; तोच त्यांचा 'B' गुण मिळालेला एकमेव विषय होता. मात्र त्यांच्या माध्यमिक शाळेच्या नोंदींनुसार तर्कशास्त्रासह सर्व विषयांत त्यांना 'A' गुण मिळाले होते. त्यांनी १९१८ ते १९२४ या काळात वॉर्सा विद्यापीठात शिक्षण घेतले. सुरुवातीला त्यांचा जीवशास्त्र शिकण्याचा विचार होता; पण विद्यापीठ हे वॉर्सा तर्कशास्त्र व तत्त्वज्ञान परंपरेचे केंद्र असल्याने त्यांना गणित, तत्त्वज्ञान आणि तर्कशास्त्रात रस निर्माण झाला. १९२४ मध्ये त्यांनी स्तानिस्वाव लेश्न्येव्स्की (Stanis\l{}aw Le\'{s}niewski) यांच्या मार्गदर्शनाखाली डॉक्टरेट मिळवली.

१९३९ मध्ये अमेरिकेत स्थलांतर करण्यापूर्वी टार्स्की यांनी वॉर्सा येथे माध्यमिक शाळेत शिकवत असताना आपले काही सर्वात महत्त्वाचे काम पूर्ण केले. तार्किक परिणाम आणि तार्किक सत्य यांवरील त्यांचे लेख याच काळात लिहिले गेले. १९३९ मध्ये ते व्याख्यानांच्या दौऱ्यासाठी अमेरिकेत गेले होते. त्याच भेटीदरम्यान जर्मनीने पोलंडवर आक्रमण केले; ज्यू वंशामुळे टार्स्की परत जाऊ शकले नाहीत. त्यांची पत्नी आणि मुले युद्ध संपेपर्यंत पोलंडमध्ये राहिली; नंतर त्यांनाही अमेरिकेत स्थलांतर करता आले. टार्स्की यांनी हार्वर्ड, सिटी कॉलेज ऑफ न्यू यॉर्क, प्रिन्स्टन येथील इन्स्टिट्यूट फॉर ॲडव्हान्स्ड स्टडी आणि अखेरीस कॅलिफोर्निया विद्यापीठ, बर्कली येथे अध्यापन केले. बर्कलीमध्ये त्यांनी तर्कशास्त्र आणि विज्ञानपद्धती यांचा बहुविषयक अभ्यास कार्यक्रम स्थापन केला. २६ ऑक्टोबर १९८३ रोजी वयाच्या ८२व्या वर्षी टार्स्की यांचे निधन झाले.

\paragraph{अधिक वाचन.}
टार्स्की यांच्या जीवनाविषयी अधिक माहितीसाठी \emph{Alfred Tarski: Life
  and Logic} हे चरित्र (Feferman 2004) पाहा. तार्किक परिणाम आणि सत्य यांवरील टार्स्की यांचे मूलभूत लेख इंग्रजीत (Corcoran 1983) येथे उपलब्ध आहेत. त्यांचे सर्व मूळ लेख चार खंडांच्या मालिकेत संकलित केले आहेत (Tarski 1981).














\section{ॲलन ट्यूरिंग}\label{his:bio:tur:sec}

ॲलन ट्यूरिंग यांचा जन्म २३ जून १९१२ रोजी लंडनमधील मायडा व्हेल येथे झाला. त्यांना सैद्धांतिक संगणकशास्त्राचे जनक मानले जाते. भौतिक विज्ञान आणि गणित यांविषयी त्यांना लहानपणापासूनच रस होता. पण ते शिकत असलेल्या शाळांत साहित्य आणि अभिजात विषयांवर भर असल्याने त्यांच्या आवडीला फार वाव मिळाला नाही. परिणामी, शाळेत त्यांची कामगिरी कमकुवत राहिली आणि अनेक शिक्षकांनी त्यांना तंबी दिली.



ट्यूरिंग यांनी केंब्रिजच्या किंग्ज कॉलेजमध्ये गणिताचे पदवीशिक्षण घेतले. गणनीय फलन ही संकल्पना नेमकी परिभाषित करण्यासाठी आणि निर्णय समस्या अनिर्णेय आहे हे सिद्ध करण्यासाठी त्यांनी १९३६ मध्ये पुढे ट्यूरिंग यंत्र म्हणून ओळखली गेलेली संकल्पना विकसित केली. ॲलोन्झो चर्च यांनी स्वतःच्या लॅम्डा कलनाच्या साहाय्याने हा निष्कर्ष आधीच सिद्ध केला होता. तरीही चर्च यांच्या निष्कर्षाचा उल्लेख करून ट्यूरिंग यांचा लेख प्रकाशित झाला. चर्च यांच्या निमंत्रणावरून ट्यूरिंग प्रिन्स्टनला गेले. तेथे ते १९३६ ते १९३८ या काळात राहिले आणि चर्च यांच्या मार्गदर्शनाखाली डॉक्टरेट मिळवली.

तर्कशास्त्रातील रस असूनही भौतिक विज्ञानातील ट्यूरिंग यांची जुनी आवड कायम होती. दुसऱ्या महायुद्धात ब्रिटनच्या ब्लेचली पार्क येथील संकेतभेदन विभागात काम करताना त्यांच्या व्यावहारिक कौशल्याचा उपयोग झाला. जर्मन नौदलाच्या संदेशांमध्ये वापरला जाणारा एनिग्मा संकेत भेदण्यात त्यांची मध्यवर्ती भूमिका होती. संख्याशास्त्र आणि संकेतलेखनातील त्यांचे कौशल्य, तसेच विद्युत् आणि यांत्रिक तंत्रांचा वापर, यांमुळे चमूला “बॉम्ब” (bombe) नावाचे संकेतभेदक यंत्र तयार करून हा संकेत भेदता आला. त्यांच्या कल्पनांनी जगातील पहिला प्रोग्राम करता येणारा इलेक्ट्रॉनिक संगणक, कोलॉसस, तयार करण्यासही मदत केली; ब्लेचली पार्क येथेच तो जर्मन लोरेन्झ संकेत भेदण्यासाठी वापरला गेला.

ट्यूरिंग समलैंगिक होते. तरीही १९४१ मध्ये त्यांनी ब्लेचली पार्कमधील सहकारी जोन क्लार्क यांना लग्नाची मागणी घातली. नंतर त्यांनी साखरपुडा मोडला आणि आपण समलैंगिक आहोत असे त्यांना सांगितले. ट्यूरिंग यांचे आयुष्यात अनेक पुरुष जोडीदार होते; त्या काळी ब्रिटनमध्ये पुरुषांमधील समलैंगिक संबंध कायद्याने गुन्हा होते. १९५२ मध्ये त्यावेळच्या जोडीदाराच्या एका मित्राने ट्यूरिंग यांच्या घरात चोरी केली. पोलिसांत तक्रार करताना ट्यूरिंग यांनी आपले समलैंगिक संबंध असल्याचे सांगितले; अशा संबंधांना लवकरच कायदेशीर मान्यता मिळेल, असा त्यांचा समज होता. मात्र तसे झाले नाही आणि त्यांच्यावर “गंभीर अश्लीलता” (gross indecency) या तत्कालीन गुन्ह्याचा आरोप ठेवण्यात आला. तुरुंगवासाऐवजी ट्यूरिंग यांनी कामेच्छा कमी करणारे संप्रेरक उपचार निवडले. ८ जून १९५४ रोजी सायनाइडचे अतिप्रमाण सेवन झाल्याने ते मृतावस्थेत आढळले; स्रोताच्या मते ती बहुधा आत्महत्या असावी. २०१३ मध्ये राणी एलिझाबेथ दुसरी यांनी त्यांना राजमाफी दिली.

\paragraph{अधिक वाचन.}
ॲलन ट्यूरिंग यांच्या सविस्तर चरित्रासाठी Hodges (2014) पाहा. त्यांच्या जीवनावर आणि कार्यावर आधारित \emph{Breaking the Code} हे नाटक १९९६ मध्ये डेरेक जेकोबी यांनी ट्यूरिंग यांची भूमिका केलेल्या दूरचित्रवाणी कार्यक्रमाच्या रूपात सादर झाले. अकादमी पुरस्कारासाठी नामांकन मिळालेला \emph{The Imitation Game} हा बेनेडिक्ट कंबरबॅच आणि कीरा नाइटली यांच्या भूमिका असलेला चित्रपटही ट्यूरिंग यांचे जीवन आणि ब्लेचली पार्कमधील काळ यांवर सैलपणे आधारलेला आहे (Tyldum 2014).

Radiolab (2012) येथे ट्यूरिंग यांच्या जीवन आणि कार्यावरील अनेक पॉडकास्ट आहेत. बीबीसी होरायझनचा \emph{The Strange Life and Death of
  Dr.~Turing} हा माहितीपट ऑनलाइन पाहता येतो (Sykes 1992). ट्यूरिंग यांच्या जन्मशताब्दीनिमित्त २०१२ मध्ये तयार केलेले कार्यरत लेगो ट्यूरिंग यंत्र दाखवणारा छोटा व्हिडिओ (Theelen 2012) येथे आहे.

ट्यूरिंग यंत्र आणि निर्णय समस्येवरील ट्यूरिंग यांचा मूळ लेख Turing (1937) हा आहे.














\section{एर्न्स्ट झर्मेलो}\label{his:bio:zer:sec}



एर्न्स्ट झर्मेलो यांचा जन्म २७ जुलै १८७१ रोजी जर्मनीतील बर्लिन येथे झाला. त्यांना पाच बहिणी होत्या; मात्र कुटुंबातील सदस्यांची तब्येत नाजूक असल्याने त्यांपैकी केवळ तीनच जणी प्रौढ वयापर्यंत जगल्या. झर्मेलो लहान असतानाच त्यांच्या आईवडिलांचेही निधन झाले आणि वयाच्या सतराव्या वर्षी ते व त्यांची भावंडे अनाथ झाली. झर्मेलो यांना कलांत, विशेषतः कवितेत, मोठा रस होता. ते तीक्ष्ण बुद्धीचे, चतुर आणि टीकात्मक स्वभावाचे म्हणून ओळखले जात. त्यांच्या सर्वाधिक प्रसिद्ध गणितीय कामगिरीत १९०४ मध्ये निवडीचे स्वयंसिद्धक मांडणे आणि १९०८ मध्ये संच उपपत्तीची स्वयंसिद्धकीय मांडणी करणे यांचा समावेश होतो.

विद्यापीठात झर्मेलो यांच्या आवडी विविध होत्या. त्यांनी भौतिकशास्त्र, गणित आणि तत्त्वज्ञान यांचे अभ्यासक्रम घेतले. हेरमान श्वार्त्स यांच्या मार्गदर्शनाखाली त्यांनी १८९४ मध्ये बर्लिन विद्यापीठात \emph{Investigations in
  the Calculus of Variations} हा प्रबंध पूर्ण केला. १८९७ मध्ये पुढील शिक्षणासाठी त्यांनी गॉटिंगेन विद्यापीठात जाण्याचे ठरवले; तेथे डेव्हिड हिल्बर्ट यांच्या गणिताच्या पायांवरील कार्याचा त्यांच्यावर मोठा प्रभाव पडला. १८९९ मध्ये ते प्राध्यापकपदासाठी पात्र झाले, पण वेतन मिळणारे प्राध्यापकपद अकरा वर्षांनीच मिळाले. त्यांच्या विचित्र वागण्यामुळे आणि “अस्वस्थ उतावळेपणा”मुळे हा विलंब झाला असावा, असा स्रोताचा अंदाज आहे.

अखेरीस १९१० मध्ये झर्मेलो यांना झुरिक विद्यापीठात वेतन मिळणारे प्राध्यापकपद मिळाले; मात्र क्षयरोगामुळे १९१६ मध्ये त्यांना निवृत्त व्हावे लागले. बरे झाल्यावर १९२६ मध्ये त्यांना फ्रायबुर्ग विद्यापीठात मानद प्राध्यापकपद देण्यात आले. या काळात त्यांनी गणिताच्या पायांवर काम केले. थोराल्फ स्कोलेम आणि कुर्ट गोडेल यांच्या कार्यावर ते नाराज झाले आणि त्यांनी आपल्या लेखांत त्यांच्या दृष्टिकोनांवर जाहीर टीका केली. १९३५ मध्ये हिटलरला अभिवादन करण्यास नकार दिल्यामुळे त्यांच्याविरुद्ध तक्रार झाल्यानंतर त्यांना मानद प्राध्यापकपदाचा राजीनामा देणे भाग पडले.

झर्मेलो यांच्या आयुष्याचा उत्तरार्ध एकाकीपणात गेला. १९३५ मध्ये राजीनामा द्यावा लागल्यानंतर त्यांनी गणिताचे काम सोडले. ते ग्रामीण भागात जाऊन साधेपणाने राहू लागले. १९४४ मध्ये त्यांनी विवाह केला. त्यांची दृष्टी मंदावत असल्याने ते पत्नीवर पूर्णपणे अवलंबून झाले. १९५१ पर्यंत त्यांची दृष्टी पूर्णपणे गेली होती. २१ मे १९५३ रोजी जर्मनीतील ग्युंटरस्टाल येथे त्यांचे निधन झाले.

\paragraph{अधिक वाचन.}
झर्मेलो यांच्या सविस्तर चरित्रासाठी Ebbinghaus (2015) पाहा. त्यांचे १९०४ आणि १९०८ मधील मूलभूत लेख मूळ जर्मन भाषेत (Zermelo 1904; Zermelo 1908) येथे उपलब्ध आहेत. भौतिकशास्त्रावरील लेखनासह त्यांचे संकलित लेख इंग्रजी भाषांतरात (Ebbinghaus 2010; Ebbinghaus 2013) येथे उपलब्ध आहेत.




\chapter{संच उपपत्तीचा इतिहास आणि दंतकथा}\label{his:set:chap}
\begin{editorial}
या प्रकरणात टिम बटन यांच्या Open Set Theory या ग्रंथातील ऐतिहासिक प्रस्तावना
समाविष्ट केली आहे.
\end{editorial}




\section{अतिसूक्ष्म राशी आणि अवकलन}\label{his:set:infinitesimals:sec}

न्यूटन आणि लाइब्निझ यांनी सतराव्या शतकाच्या अखेरीस एकमेकांपासून स्वतंत्रपणे कलनाचा शोध लावला. कलनाचा एक विशेष महत्त्वाचा उपयोग म्हणजे \emph{अवकलन}. ढोबळपणे सांगायचे तर, एखाद्या बिंदूवर फलनाच्या “बदलाचा दर” किंवा उतार निश्चित करणे हा अवकलनाचा उद्देश असतो.

ही कल्पना एका ठळक उदाहरणाने समजते. खाली काळ्या रेषेत दाखवलेले $f(x) = \nicefrac{x^2}{4} + \nicefrac{1}{2}$ हे फलन पाहा:
\begin{center}
	\begin{tikzpicture}[scale=1]
	\draw[->, gray] (-1,0) -- (4.25,0) node[right] {$x$};
	\draw[->, gray] (0,-0.25) -- (0,5.25) node[above] {$f(x)$};

	\foreach \x/\xtext in {1/1, 2/2, 3/3, 4/4}
	\draw[shift={(\x,0)}, gray] (0pt,2pt) -- (0pt,-2pt) node[below] {$\xtext$};

	\foreach \y/\ytext in {1/1, 2/2, 3/3, 4/4, 5/5}
	\draw[shift={(0,\y)}, gray] (2pt,0pt) -- (-2pt,0pt) node[left] {$\ytext$};

	\draw[oldiagcolorC] (0.5, 9/16) -- (3.5, 57/16) -- (3.5, 9/16)--cycle;
	\draw[oldiagcolorD] (.5, 9/16) -- (2.5, 33/16) -- (2.5, 9/16) -- cycle;
	\draw[oldiagcolorE] (.5, 9/16) -- (1.5, 17/16) -- (1.5, 9/16) -- cycle;
	\draw[thick] (-1,.75) parabola bend (0,0.5) (4,4.5);
	\end{tikzpicture}
\end{center}
आपल्याला $c = \nicefrac{1}{2}$ येथे या फलनाचा उतार शोधायचा आहे, असे समजा. प्रथम असा त्रिकोण काढू की त्याचा कर्ण त्या बिंदूवरील उताराचे सन्निकटन देतो; वरील लाल त्रिकोण हे एक उदाहरण आहे. आपल्या त्रिकोणाच्या पायाची लांबी $\beta$ असेल, तर त्याची उंची $f(\nicefrac{1}{2}+\beta) - f(\nicefrac{1}{2})$ असेल. म्हणून कर्णाचा उतार असा मिळतो:
\[
\frac{f(\nicefrac{1}{2}+\beta) - f(\nicefrac{1}{2})}{\beta}.
\]
म्हणून पायाची लांबी~$3$ असलेल्या \textcolor{oldiagcolorC}{लाल} त्रिकोणाचा उतार नेमका~$1$ आहे. त्याहून लहान, पायाची लांबी~$2$ असलेल्या \textcolor{oldiagcolorD}{निळा} त्रिकोणाचा कर्ण अधिक चांगले सन्निकटन देतो; त्याचा उतार $\nicefrac{3}{4}$ आहे. त्याहूनही लहान, पायाची लांबी~$1$ असलेला \textcolor{oldiagcolorE}{हिरवा} त्रिकोण आणखी चांगले सन्निकटन देतो; त्याचा उतार $\nicefrac{1}{2}$ आहे.

त्रिकोण जसजसे लहान होत जातात तसतसे आपल्याला अधिक चांगली सन्निकटने मिळतात. म्हणून आपण असे म्हणू शकू: \emph{अतिसूक्ष्म} पायाची लांबी असलेल्या त्रिकोणाचा कर्ण $c = \nicefrac{1}{2}$ येथील नेमका उतार देतो. अशा रीतीने $c$ या बिंदूवरील $f$ फलनाच्या (पहिल्या) अवकलजाचे सूत्र मिळेल:
\[
{f'}(c) = \frac{f(c+\beta) - f(c)}{\beta} \text{ येथे $\beta$ अतिसूक्ष्म आहे.}
\]
ढोबळपणे पाहता, न्यूटन आणि लाइब्निझ यांचे म्हणणे हेच होते.

मात्र असे म्हटल्यावर आपण त्यांना विचारले पाहिजे: \emph{अतिसूक्ष्म} म्हणजे नेमके काय? येथे एक गंभीर द्विधा निर्माण होते. $\beta = 0$ असेल, तर $f'$ ची व्याख्या नीट होत नाही, कारण त्यात $0$ ने भागाकार करावा लागतो. पण $\beta > 0$ असेल, तर आपल्याला नेमका उतार नव्हे, तर उताराचे \emph{सन्निकटन}च मिळते.

ही नंतरच्या काळातील कल्पना त्या काळावर लादलेली शंका नाही. न्यूटन यांच्या अनुयायांवर टीका करताना बर्कली असे म्हणतात:
\begin{quote}
चिन्हे कोणतीही गोष्ट किंवा काहीच नाही असे दर्शवू शकतात, हे मी मान्य करतो. त्यामुळे मूळ संकेतनातील $c + \beta$ या पदात $\beta$ हे वाढीचे प्रमाण किंवा काहीच नाही, यांपैकी काहीही दर्शवू शकले असते. पण तुम्ही त्याला यांपैकी जो अर्थ द्याल, त्या अर्थाशी सुसंगतपणे युक्तिवाद केला पाहिजे; दुहेरी अर्थाचा आधार घेऊन पुढे जाता कामा नये. तसे करणे म्हणजे उघड कुतर्क होय. (Berkeley 1734, \S{}XIII; आधीच्या मजकुराशी जुळण्यासाठी चल बदलले आहेत.)
\end{quote}
बर्कली यांच्या टीकेपासून अतिसूक्ष्म राशींवरील कलनाचा बचाव करताना, “अतिसूक्ष्म राशी” हा शब्दप्रयोग केवळ रूपकात्मक आहे, असे कोणी म्हणू शकेल. आपण \emph{खरोखर फार लहान} त्रिकोण घेतला, तर स्पर्शरेषेचे \emph{पुरेसे चांगले} सन्निकटन मिळेल, असे त्यांचे म्हणणे असेल. यावरही बर्कली यांचे उत्तर होते: अभियांत्रिकीसाठी ते पुरेसे असेलही, पण त्यामुळे गणिताच्या \emph{स्थानाला} धक्का बसतो; कारण
\begin{quote}
आपल्याला असे सांगितले जाते की \emph{in rebus mathematicis errores qu\`{a}m minimi
non sunt contemnendi}. [गणितात अगदी लहानात लहान चुकाही दुर्लक्षण्याजोग्या नसतात.] (Berkeley 1734, \S{}IX)
\end{quote}
तिरप्या अक्षरांतील हा उतारा न्यूटन यांच्या स्वतःच्या \emph{Quadrature of Curves} (1704) या ग्रंथातील विधानाशी जवळजवळ शब्दशः जुळतो.

बर्कली यांचे तात्त्विक आक्षेप अत्यंत नेमके होते. तरीसुद्धा कलन अत्यंत यशस्वी ठरले आणि गणितज्ञांनी चुका न करता त्याचा वापर चालू ठेवला.












\section{सीमांची काटेकोर व्याख्या}\label{his:set:limits:sec}

आजकाल आधीच्या समस्येवरचे प्रमाणित उत्तर म्हणजे अतिसूक्ष्म राशींचा आधार सोडणे. ते कसे करता येते ते पाहू.

आपण पाहिले की $\beta$ जसजसा लहान होतो तसतशी उताराची अधिक चांगली सन्निकटने मिळतात. खरे तर $\beta$ हा $0$ च्या कितीही जवळ जात असताना, $f'(c)$ या आधीच्या सूत्रातील भागाकाराचा मान आपल्याला हव्या असलेल्या उताराकडे झुकतो. त्यामुळे $\beta = 0$ \emph{असतानाची} स्थिती पाहण्याऐवजी, $\beta$ हा $0$ कडे जात असताना $f'(c)$ च्या त्या सूत्रातील भागाकाराचा \emph{कल} पाहणे पुरेसे आहे.

अशा प्रकारे मांडल्यास, पुढील स्वरूपाच्या विधानांचा अर्थ स्पष्ट करणे हे सर्वसाधारण आव्हान ठरते:
\begin{center}
$x$ हा $c$ कडे जात असताना $g(x)$ चे मूल्य $\ell$ कडे झुकते.
\end{center}
हे अधिक संक्षेपात असे लिहिता येते:
\[
\lim_{x \rightarrow c}g(x) = \ell.
\]
एकोणिसाव्या शतकात कॉशी यांच्या आधीच्या कार्यावर आधार घेत वायेरश्ट्रास यांनी या विधानाची पूर्णपणे काटेकोर व्याख्या दिली. $x$ हा $c$ च्या योग्य इतका जवळ घेतल्यास $g(x)$ हे $\ell$ च्या आपल्याला हवे तितके जवळ नेता येते, हीच त्यामागील कल्पना आहे. अधिक नेमकेपणे, $\lim_{x \rightarrow c} g(x) = \ell$ याचा अर्थ असा ठरवू:
\[
(\forall\epsilon > 0)(\exists \delta > 0)\forall x \left(0 < |x - c| < \delta \lif |g(x) - \ell| < \epsilon \right).
\]

येथील उभ्या रेषा निरपेक्ष मूल्य दर्शवतात. म्हणजे $x \geq 0$ असेल तेव्हा $|x| = x$, आणि $x < 0$ असेल तेव्हा $|x| =-x$; \emph{त्या} फलनाचा आलेख पुढीलप्रमाणे दाखवता येतो:
\begin{center}
	\begin{tikzpicture}[scale=1]
	\draw[->, gray] (-2.5,0) -- (2.5,0) node[right] {$x$};
	\draw[->, gray] (0,-.2) -- (0,2.5) node[above] {$|x|$};

	\foreach \x/\xtext in {-2/-2, -1/-1, 1/1, 2/2}
	\draw[shift={(\x,0)}] (0pt,2pt) -- (0pt,-2pt) node[below] {{$\xtext$}};

	\foreach \y/\ytext in {1/1, 2/2}
	\draw[shift={(0,\y)}] (2pt,0pt) -- (-2pt,0pt) node[left] {{$\ytext$}};

	\draw[thick] (-2.5, 2.5)--(0,0)--(2.5, 2.5);
	\end{tikzpicture}
\end{center}
म्हणून या व्याख्येचा ढोबळ अर्थ असा: $c$ पेक्षा वेगळी पण त्याच्या $\delta$ इतक्या जवळची आदाने निवडून (म्हणजे $0 < |x - c| < \delta$), आपली “त्रुटी” $\epsilon$ पेक्षा लहान करता येते (म्हणजे $|g(x) - \ell| < \epsilon$).

सीमेची संकल्पना व्याख्येने निश्चित केल्यावर आपण अतिसूक्ष्म राशी पूर्णपणे टाळू शकतो. त्यासाठी $c$ येथील $f$ चा उतार असा ठरवू:
\[
	{f}'(c) = \lim_{x \rightarrow 0}\left(\frac{f(c +x) - f(c)}{x}\right) \text{ जेथे सीमा अस्तित्वात आहे}.
\]
मात्र “जेथे सीमा अस्तित्वात आहे” ही अट व्याख्येत का हवी, हे लक्षात घेणे महत्त्वाचे आहे. साधे उदाहरण म्हणून आपण आत्ताच ज्याचा आलेख पाहिला तो $f(x) = |x|$ विचारात घ्या. येथे $f'(0)$ ची व्याख्या होत नाही: $0$ कडे “उजवीकडून” गेल्यास उतार नेहमी $1$ असतो, तर $0$ कडे “डावीकडून” गेल्यास उतार नेहमी $-1$ असतो; म्हणून सीमा अस्तित्वात नाही. त्यामुळे अशी सीमा अस्तित्वात असेल तर आणि तरच फलन~$f$ हे~$x$ येथे \emph{अवकलनीय} आहे, असेही म्हणता येईल.

अवकलनासाठी \emph{सीमेची} संकल्पना कशी वापरायची ते आपण पाहिले. त्याच संकल्पनेने \emph{सतत} फलनाची कल्पनाही व्याख्येने निश्चित करता येते. (१८१७ पर्यंत बोल्झानो यांना हे तत्त्वतः उमगले होते.) कॉशी–वायेरश्ट्रास यांची सातत्याची मांडणी अशी आहे. ढोबळपणे: फलनाच्या निर्गतात ठरावीक अचूकता हवी असेल, तर आदानात पुरेशी अचूकता मागून ती मिळेल, अशा वेळी फलन~$f$ त्या बिंदूवर सतत असते. अधिक नेमकेपणे: $x$ हा शून्याकडे जात असताना $f(c + x)$ आणि $f(c)$ यांतील फरकही $0$ कडे जात असेल, तर $f$ हे $c$ येथे सतत आहे. दुसऱ्या शब्दांत, $f(c) = \lim_{x \rightarrow c} f(x)$ असेल तर आणि तरच $f$ हे $c$ येथे \emph{सतत} आहे.

यापुढे गेलो तर आपला विषय असलेल्या संच उपपत्तीपासून दूर जाऊन आपण वास्तव विश्लेषणात शिरू. म्हणून आता थांबून मुख्य बोध सांगू. एकोणिसाव्या शतकात गणितज्ञांनी काटेकोरपणे व्याख्या केलेल्या \emph{सीमेच्या} संकल्पनेचा वापर करून अतिसूक्ष्म राशींशिवाय काम करायला शिकले.







\section{विलक्षण रचना}\label{his:set:pathology:sec}

मात्र \emph{सीमेच्या} व्याख्येमुळे काही खरोखर “विलक्षण” रचना शक्य असल्याचे दिसून आले.

१८३०च्या सुमारास बोल्झानो यांनी असे एक फलन शोधले जे \emph{सर्वत्र सतत} होते, पण \emph{कोठेही अवकलनीय नव्हते}. (दुर्दैवाने बोल्झानो यांनी हे कधी प्रकाशित केले नाही. वायेरश्ट्रास यांनी स्वतंत्रपणे हीच कल्पना शोधल्यामुळे गणितज्ञांना १८७२ मध्ये तिची प्रथम ओळख झाली.)\footnote{हा इतिहास \href{http://en.wikipedia.org/wiki/Weierstrass_function}{वायेरश्ट्रास फलनावरील विकिपीडिया लेखाच्या} अत्यंत तपशीलवार तळटीपांत नोंदवला आहे.} हे किमान आश्चर्यकारक तरी होतेच. $|x|$ सारखी सर्वत्र सतत, पण एखाद्या विशिष्ट बिंदूवर अवकलनीय नसलेली फलने सहज सापडतात. परंतु सर्वत्र सतत असूनही \emph{कोठेही} अवकलनीय नसलेले फलन अगदी वेगळेच आहे. असे फलन कसे काढाल याचा क्षणभर विचार करा. ते सतत असण्यासाठी कागदावरून लेखणी कधीही न उचलता त्याचा आलेख काढता आला पाहिजे; पण ते कोठेही अवकलनीय नसण्यासाठी लेखणीची दिशा सतत अचानक बदलावी लागेल.

यानंतर आणखी “विलक्षण” रचना समोर आल्या. ५ जानेवारी १८७४ रोजी कँटर यांनी डेडेकिंड यांना पत्र लिहून पुढील प्रश्न विचारला:
\begin{quote}
एखाद्या पृष्ठभागाची (उदाहरणार्थ सीमेसह चौरसाची) एखाद्या रेषेशी (उदाहरणार्थ टोकांसह सरळ रेषेशी) अशी एकास-एक संगती लावता येईल का की पृष्ठभागावरील प्रत्येक बिंदूला रेषेवरील एक बिंदू मिळेल आणि उलट रेषेवरील प्रत्येक बिंदूला पृष्ठभागावरील एक बिंदू मिळेल?

या प्रश्नाचे उत्तर शोधणे अजूनही मला फार कठीण वाटते—तरी येथेही \emph{नाही} असे म्हणण्याची इतकी तीव्र ओढ वाटते की त्यासाठी सिद्धतेची गरज जवळजवळ नसावी असे वाटू लागते.
[Gouv\^{e}a 2011 मधून उद्धृत]
\end{quote}
पण १८७७ मध्ये कँटर यांनी स्वतःचा अंदाज चुकीचा असल्याचे सिद्ध केले. प्रत्यक्षात रेषा आणि चौरस यांवरील बिंदूंची संख्या अगदी सारखी असते. २९ जून १८७७ रोजी त्यांनी डेडेकिंड यांना “\emph{je le vois, mais je ne le crois pas}” असे लिहिले; म्हणजे, “ते मला दिसते, पण त्यावर माझा विश्वास बसत नाही.” गणिताच्या “प्रचलित इतिहासकथनात” याचा अर्थ अनेकदा असा घेतला जातो की त्या काळच्या गणितज्ञांना हे नवे निष्कर्ष \emph{अक्षरशः अविश्वसनीय} वाटत होते. (हा पत्रव्यवहार Gouv\^{e}a (2011) येथे दिला आहे आणि \ref{his:set:mythology:sec} मध्ये आपण त्याकडे पुन्हा वळू. कँटर यांच्या सिद्धतेचा आराखडा \ref{his:set:cantorplane:sec} मध्ये दिला आहे.)

कँटर यांच्या निष्कर्षाने प्रेरित होऊन, एखाद्या रेषेचे \emph{अखंडपणे} प्रतलावर प्रतिचित्रण करून ते प्रतल भरून काढता येईल का, असा विचार पियानो यांनी सुरू केला. असे प्रतिचित्रण म्हणजे \emph{अवकाश भरणारा वक्र} असेल. 1890 मध्ये पियानो यांनी असाच वक्र रचला. हे अंतर्ज्ञानाला खरोखर विरुद्ध वाटते: युक्लिडने रेषेची व्याख्या “रुंदीविरहित लांबी” अशी केली होती (पुस्तक I, व्याख्या २); पण रेषा योग्य प्रकारे वळवली, तर तिची लांबी रुंदी व्यापू शकते, हे पियानो यांनी दाखवले. 1891 मध्ये हिल्बर्ट यांनी याहून किंचित अधिक सहज समजणारा अवकाश भरणारा वक्र आणि त्याची काही चित्रे दिली. हा वक्र टप्प्याटप्प्याने रचला जातो; रचनेचे पहिले सहा टप्पे येथे दिले आहेत:
\begin{center}
\begin{tikzpicture}
\begin{scope}[xshift=20pt, yshift=20pt]
	\draw [oldiagcolorC, l-system={Hilbert curve, step=40pt, angle=90, axiom=L, order=1}] lindenmayer system;
\end{scope}
\begin{scope}[xshift=110pt, yshift=10pt]
	\draw [oldiagcolorC, l-system={Hilbert curve, step=20pt, angle=90, axiom=L, order=2}] lindenmayer system;
\end{scope}
\begin{scope}[xshift=205pt, yshift=5pt]
	\draw [oldiagcolorC, l-system={Hilbert curve, step=10pt, angle=90, axiom=L, order=3}] lindenmayer system;
\end{scope}
\begin{scope}[xshift=2.5pt, yshift=-97.5pt]
	\draw [oldiagcolorC, l-system={Hilbert curve, step=5pt, angle=90, axiom=L, order=4}] lindenmayer system;
\end{scope}
\begin{scope}[xshift=101.25pt, yshift=-98.75pt]
	\draw [oldiagcolorC, l-system={Hilbert curve, step=2.5pt, angle=90, axiom=L, order=5}] lindenmayer system;
\end{scope}
\begin{scope}[xshift=200.625pt, yshift=-99.375pt]
	\draw [oldiagcolorC, l-system={Hilbert curve, step=1.25pt, angle=90, axiom=L, order=6}] lindenmayer system;
\end{scope}
\draw[gray] (0pt,0pt) rectangle (80pt, 80pt);
\draw[gray] (100pt,0pt) rectangle (180pt, 80pt);
\draw[gray] (200pt,0pt) rectangle (280pt, 80pt);
\draw[gray] (0pt,-100pt) rectangle (80pt, -20pt);
\draw[gray] (100pt,-100pt) rectangle (180pt, -20pt);
\draw[gray] (200pt,-100pt) rectangle (280pt, -20pt);
\end{tikzpicture}
\end{center}
सीमेत—जी संकल्पना आता काटेकोरपणे व्याख्यित झाली होती—संपूर्ण चौरस एकसंध \textcolor{oldiagcolorC}{लाल} रंगाने भरतो. तसेच, हिल्बर्ट यांचा वक्र सर्वत्र सतत आहे, पण कोठेही अवकलनीय नाही; अंतर्ज्ञानाने सांगायचे तर, अनंत टप्प्यांच्या सीमेत फलन प्रत्येक क्षणी अचानक दिशा बदलते. (हिल्बर्ट यांच्या रचनेचा अधिक तपशीलवार आराखडा \ref{his:set:hilbertcurve:sec} मध्ये पाहू.)

या “विलक्षण” भूमितीय रचनांना, योग्य की अयोग्य, भूमितीय अंतर्ज्ञानावर विसंबून राहण्याबाबत शंका घेण्याचे कारण मानले गेले. गणित करण्याच्या एका नव्या पद्धतीच्या \emph{प्रचाराचे} त्या जणू साधन बनल्या; या पद्धतीचा विकास पुढे संच उपपत्तीत झाला. या विषयाभोवती नंतर उभ्या राहिलेल्या दंतकथांमध्ये, हे निष्कर्ष पूर्णपणे काटेकोर आणि तितकेच धक्कादायक होते, असे वारंवार सांगितले गेले. त्यामुळे त्यांनी दोन कामे केली: भूमितीय अंतर्ज्ञानावर अवलंबून राहण्याविरुद्ध सावध केले आणि नव्या विचारपद्धतींची फलद्रूपता दाखवली.








\section{इतिहासापेक्षा दंतकथाच अधिक?}\label{his:set:mythology:sec}

या घटनांकडे शतकाहून अधिक काळानंतर मागे वळून पाहताना, त्यांच्याबद्दल प्रचलित झालेले निष्कर्ष विनातपासणी स्वीकारू नयेत. ते गणिती निष्कर्ष नक्कीच नवे, उत्सुकता वाढवणारे आणि आश्चर्यकारक होते. पण ते खरोखर किती धक्कादायक होते? आणि भूमितीय अंतर्ज्ञानावर विसंबू नये, हे त्यांनी खरोखर दाखवले का?

धक्क्याच्या प्रश्नावर Gouv\^{e}a (2011) असे निदर्शनास आणतात की कँटर यांनी डेडेकिंड यांना लिहिलेले प्रसिद्ध “\emph{je le vois, mais je ne le crois
pas}” हे वाक्य त्याच्या संदर्भातून बरेचसे वेगळे केले जाते. तो संदर्भ पुढे दिला आहे (Gouv\^{e}a यांच्याकडून उद्धृत):
\begin{quote}
या विषयाबद्दलच्या माझ्या उत्साहापोटी मी तुमच्या दयाळूपणाची आणि संयमाची इतकी मागणी करत असेन, तर मला क्षमा करा. मी तुम्हाला अलीकडे पाठवलेले विचार मलाही इतके अनपेक्षित, इतके नवे वाटतात की, आदरणीय मित्रा, त्यांची अचूकता तुम्ही निश्चित करेपर्यंत मला मनःशांती मिळू शकत नाही. तुम्ही माझ्याशी सहमत होत नाही तोपर्यंत मी एवढेच म्हणू शकतो:
\emph{je le vois, mais je ne le crois pas.}
\end{quote}
आपला निष्कर्ष “इतका अनपेक्षित, इतका नवा” आहे हे कँटर यांना माहीत होते. पण तो निष्कर्ष त्यांना खरोखर \emph{अविश्वसनीय} वाटत होता का, याबद्दल शंका आहे. Gouv\^{e}a यांनी दाखवल्याप्रमाणे, त्यांनी डेडेकिंड यांना केवळ स्वतः दिलेल्या सिद्धतेची तपासणी करायला सांगितले होते.

भूमितीय अंतर्ज्ञानाच्या प्रश्नावर पाहिले तर, पियानो यांनी अवकाश भरणारा वक्र प्रकाशित करताना कोणतीही चित्रे दिली नव्हती. पण हिल्बर्ट यांनी स्वतःचा वक्र प्रकाशित केला तेव्हा आपला हेतू त्यांनी स्पष्ट केला: वाचकांनी “पुढील भूमितीय अंतर्ज्ञानाचा आधार घेतल्यास” पियानो यांचा निष्कर्ष त्यांना नीट समजावा, म्हणून त्यांनी \ref{his:set:pathology:sec} मध्ये दिलेल्यांसारखी \emph{चित्रे} त्यात जोडली.

अधिक सर्वसाधारणपणे सांगायचे तर, प्रकाशित सिद्धतांमध्ये चित्रे वापरण्याची पद्धत काहीशी कमी झाली असली तरी गणितज्ञ सिद्धता करताना \emph{वारंवार} चित्रे वापरतात, ही गोष्ट नाकारता येत नाही. (ढोबळपणे सांगायचे तर, भूमितीय अंतर्ज्ञानावर कधी विसंबता येते हे चांगल्या गणितज्ञांना माहीत असते.)

थोडक्यात, अतिरंजित प्रचारावर विश्वास ठेवू नका; निदान तो विनातपासणी स्वीकारू नका. याबद्दल अधिक माहितीसाठी Giaquinto (2007) वाचता येईल.








\section{रेषा आणि प्रतल: कँटर यांचा निष्कर्ष}\label{his:set:cantorplane:sec}

श्रेडर--बर्नस्टाइन प्रमेयाच्या सिद्धतेशी संबंधित काही परिस्थितींचा संबंध \ref{his:set:pathology:sec} मध्ये चर्चिलेल्या इतिहासाशी आहे. आठवा: १८७७ मध्ये कँटर यांनी चौरसावर जितके बिंदू आहेत तितकेच त्याच्या एका बाजूवरही आहेत, हे सिद्ध केले. त्यांच्या सुरुवातीच्या प्रयत्नातील अंक एकाआड एक लिहिण्याची कल्पना येथे द्विमान रूपात मांडू. मूळ ऐतिहासिक प्रयत्न दशमान विस्तारांवर आधारलेला होता; खालील द्विमान मांडणी ही त्याचे आधुनिक रूपांतर आहे ((Gouv\^{e}a 2011, पृ.~201--202)).

$\unitline$ ही एकक रेषा, म्हणजे $[0,1]$ या बिंदूंचा संच, असू द्या. $\unitsquare$ हा एकक चौरस, म्हणजे $\unitline \times \unitline$ या बिंदूंचा संच, असू द्या. या संज्ञांत कँटर यांनी $\cardeq{\unitline}{\unitsquare}$ सिद्ध केले. त्यांनी डेडेकिंड यांना पत्रातून कळवलेल्या कल्पनेचे द्विमान रूपांतर पुढे दिले आहे.

\begin{thm}\label{his:set:cantorplane:thm:cantorplane}
$\cardeq{\unitline}{\unitsquare}$
\end{thm}

\begin{proof}[सिद्धता: पहिला भाग.]
$a,b\in\unitline$ निश्चित करा. द्विमान विस्तारासाठी पुढील
एकच पद्धत ठरवू: $0=0.000\ldots$; दोन निरूपणे असलेल्या
धन संख्येसाठी शेवटी अनंत $1$ असलेले निरूपण वापरू.
विशेषतः $\nicefrac12=0.0111\ldots$ आणि $1=0.111\ldots$.
इतर संख्येचे द्विमान निरूपण एकमेव असते. या निवडीने
प्रत्येक धन संख्येच्या विस्तारात अनंत $1$ येतात. असे लिहा:
\begin{align*}
a &= 0.a_1a_2a_3a_4\dots\\
b &= 0.b_1b_2b_3b_4\dots
\intertext{आता पुढीलप्रमाणे दिलेले $f\colon\unitsquare\to\unitline$ हे फलन विचारात घ्या:}
f(a,b)&=0.a_1b_1a_2b_2a_3b_3a_4b_4\dots
\end{align*}
परिणामी विस्तारात अनंत $1$ असतात, जोपर्यंत $a=b=0$
नाहीत; त्या प्रकरणात तो पूर्णपणे शून्यांचा असतो.
म्हणून परिणामी विस्तारही वर ठरवलेल्या पद्धतीतीलच आहे.
त्याच्या अंकांचे एकमेवपण वापरल्यास, $f(a,b)=f(c,d)$
असताना प्रत्येक $n\in\Nat$ साठी $a_n=c_n$ आणि
$b_n=d_n$ मिळते. म्हणून $a=c$ आणि $b=d$;
म्हणजे $f$ हे एकास-एक फलन आहे.
\end{proof}

मात्र हे फलन आच्छादक फलन नाही. उदाहरणार्थ,
द्विमान संख्या $z=0.1101010101\ldots=\nicefrac56$ घ्या.
तिचा द्विमान विस्तार सांत नाही, म्हणून त्या विस्ताराची दुसरी
निवड उपलब्ध नाही. $f(a,b)=z$ असते, तर तिचे विषम आणि
सम स्थानांवरील अंक वेगळे केल्याने अनुक्रमे
\begin{align*}
a&=0.100000\ldots=\nicefrac12,\\
b&=0.111111\ldots=1
\end{align*}
असे विस्तार मिळाले असते. पण $a$ चा हा विस्तार आपण
ठरवलेल्या पद्धतीला मान्य नाही: $\nicefrac12$ साठी
$0.011111\ldots$ वापरायचा आहे. त्यामुळे $z$ साठी
कोणतीही जोडी $\tuple{a,b}$ मिळत नाही.
उदाहरणार्थ, $0.101010\ldots$ हा मात्र
या पद्धतीत $f(1,0)$ आहे; म्हणून तो येथे प्रतिउदाहरण नाही.

अंक एकाआड एक लिहिण्याच्या ऐतिहासिक प्रयत्नात
निरूपणांची निवड महत्त्वाची आहे, हे डेडेकिंड यांनी
निदर्शनास आणले. वरच्या निश्चित द्विमान पद्धतीत $f$
हे एकास-एक फलन आहे, पण आच्छादक फलन नाही.
म्हणून या पहिल्या भागाने फक्त
$\cardle{\unitsquare}{\unitline}$ दाखवले आहे;
$\cardeq{\unitsquare}{\unitline}$ अजून सिद्ध झालेले नाही.

कँटर यांनी डेडेकिंड यांना जवळजवळ लगेच उत्तर लिहून सिद्धता साधारणपणे पुढीलप्रमाणे पूर्ण करता येईल असे सुचवले:

\begin{proof}[सिद्धता: पूर्ण रूप.]
आपण $\cardle{\unitsquare}{\unitline}$ दाखवले आहे. पण $\unitline$ वरून $\unitsquare$ मध्ये जाणारे एकास-एक फलन उघडच आहे: रेषा चौरसाच्या एका बाजूवर सरळ ठेवा. म्हणून $\cardle{\unitline}{\unitsquare}$ आणि $\cardle{\unitsquare}{\unitline}$. श्रेडर--बर्नस्टाइन प्रमेयानुसार (\ref{sfr:siz:sb:thm:schroder-bernstein}) $\cardeq{\unitline}{\unitsquare}$.
\end{proof}

मात्र कँटर यांना अर्थातच शेवटची ओळ या रूपात पूर्ण करता आली नसती, कारण श्रेडर--बर्नस्टाइन प्रमेय तेव्हा अजून सिद्ध झाले नव्हते. पुढे कँटर यांनी हा सर्वसाधारण अंदाज मांडला, तरी त्याची समाधानकारक सिद्धता १८९७ पर्यंत मिळाली नव्हती. (म्हणूनच १८७७ मध्ये पुढे कँटर यांनी \ref{his:set:cantorplane:thm:cantorplane} ची श्रेडर--बर्नस्टाइन प्रमेयाचा आधार न घेणारी वेगळी सिद्धता दिली.)









\section{परिशिष्ट: हिल्बर्ट यांचे अवकाश भरणारे वक्र}\label{his:set:hilbertcurve:sec}

रेषा आणि चौरस यांवर नेमके सारखेच बिंदू असतात, या कँटर यांच्या सिद्धतेमुळे (\ref{his:set:cantorplane:thm:cantorplane}) अवकाश भरणारा \emph{वक्र} असू शकेल का, हा प्रश्न पियानो यांना कसा पडला, याचा उल्लेख आपण \ref{his:set:pathology:sec} मध्ये केला होता. १८९० मध्ये पियानो यांना त्याचे होकारार्थी उत्तर मिळाले. या विभागात हिल्बर्ट यांचा अवकाश भरणारा वक्र (एका लहानशा बदलासह) कसा रचता येतो, हे अनौपचारिक रीतीने स्पष्ट करू.\footnote{अधिक काटेकोर स्पष्टीकरणासाठी (Rose 2010) पाहा. या बदलात खालील वक्रांचे लाल भाग समाविष्ट केले आहेत. त्यामुळे वक्र सतत आहे हे तपासणे थोडे सोपे होते.}

$h$ हे फलन $h_1$, $h_2$, $h_3$, \dots\@ या फलन-अनुक्रमाची सीमा म्हणून व्याख्यित करायचे आहे. प्रथम रचनेचे वर्णन करू. मग ती रचना अवकाश भरते हे दाखवू. त्यानंतर ती खरोखर वक्र आहे हे दाखवू.

$h$ चा लक्ष्यसंच एकक चौरस $\unitsquare$ घेऊ. $h$ चे पहिले सन्निकटन, म्हणजे $h_1$, असे दिसते:
\begin{center}
	\begin{tikzpicture}
	\draw[gray] (0pt, 0pt) rectangle (80pt, 80pt);
	\draw[step = 40pt, gray, very thin] (0pt, 0pt) grid (80pt, 80pt);
	\draw[oldiagcolorC, thick] (20pt, 0)--(20pt, 20pt);
	\draw[oldiagcolorC, thick] (60pt, 0)--(60pt, 20pt);
	\begin{scope}[xshift=20pt, yshift=20pt]
	\draw [black, thick, l-system={Hilbert curve, step=40pt, angle=90, axiom=L, order=1}]  lindenmayer system;
	\end{scope}
	\end{tikzpicture}
\end{center}
रचनेचा मागोवा घेण्यासाठी चौरसावर $2 \times 2$ जाळी आखली आहे. वक्राची सुरुवात खालच्या डाव्या चतुर्थांशात होते, तो वरच्या डाव्या, मग वरच्या उजव्या आणि शेवटी खालच्या उजव्या चतुर्थांशात जातो, असे मानता येते. रचनेचा दुसरा टप्पा, म्हणजे $h_2$, असा आहे:
\begin{center}
	\begin{tikzpicture}
	\draw[gray] (0pt, 0pt) rectangle (80pt, 80pt);
	\draw[step = 20pt, gray, very thin] (0pt, 0pt) grid (80pt, 80pt);
	\draw[oldiagcolorC, thick] (0pt, 10pt)--(10pt, 10pt);
	\draw[oldiagcolorC, thick] (70pt, 10pt)--(80pt, 10pt);
	\draw[thick, oldiagcolorE] (10pt, 30pt)--(10pt, 50pt);
	\draw[thick, oldiagcolorE] (30pt, 50pt)--(50pt, 50pt);
	\draw[thick, oldiagcolorE] (70pt, 50pt)--(70pt, 30pt); \begin{scope}[xshift=10pt, yshift=30pt]
	\draw [thick, rotate=270,black, l-system={Hilbert curve, step=20pt, angle=90, axiom=L, order=1}]  lindenmayer system;
	\end{scope}
	\begin{scope}[xshift=10pt, yshift=50pt]
	\draw [thick, black, l-system={Hilbert curve, step=20pt, angle=90, axiom=L, order=1}]  lindenmayer system;
	\end{scope}
	\begin{scope}[xshift=50pt, yshift=50pt]
	\draw [thick, black, l-system={Hilbert curve, step=20pt, angle=90, axiom=L, order=1}]  lindenmayer system;
	\end{scope}
	\begin{scope}[xshift=70pt, yshift=10pt]
	\draw [thick, rotate=90,black, l-system={Hilbert curve, step=20pt, angle=90, axiom=L, order=1}]  lindenmayer system;
	\end{scope}
	\end{tikzpicture}
\end{center}
$h_2$ कसा रचला ते वेगवेगळ्या रंगांमुळे स्पष्ट होते. प्रथम $h_1$ च्या \textcolor{oldiagcolorC}{लाल} नसलेल्या भागाच्या लहान प्रमाणातील प्रती चौरसाच्या खालच्या डाव्या, वरच्या डाव्या, वरच्या उजव्या आणि खालच्या उजव्या भागांत ठेवतो (त्या काळ्या रंगात दाखवल्या आहेत). मग या चार आकृत्या एकमेकांना (\textcolor{oldiagcolorE}{हिरवा} रेषांनी) जोडतो. शेवटी संपूर्ण आकृती चौरसाच्या सीमेला (\textcolor{oldiagcolorC}{लाल} रेषांनी) जोडतो.

आता $h_3$ पाहू. $h_1$ च्या लाल नसलेल्या भागाच्या एकमेकांना जोडलेल्या चार लहान प्रती वापरून $h_2$ जसा तयार केला, तसाच $h_2$ च्या लाल नसलेल्या भागाच्या चार लहान प्रतींमधून $h_3$ तयार होतो (प्रती काळ्या रंगात दाखवल्या आहेत). नंतर त्या (\textcolor{oldiagcolorE}{हिरवा} रेषांनी) जोडल्या जातात आणि शेवटी चौरसाच्या सीमेला (\textcolor{oldiagcolorC}{लाल} रेषांनी) जोडल्या जातात.
\begin{center}
	\begin{tikzpicture}
	\draw[gray] (0pt, 0pt) rectangle (80pt, 80pt);
	\draw[step = 10pt, gray, very thin] (0pt, 0pt) grid (80pt, 80pt);
	\draw[thick, oldiagcolorC](5pt, 0pt)--(5pt, 5pt);
	\draw[thick, oldiagcolorC] (75pt, 0pt)--(75pt, 5pt);
	\draw[thick, oldiagcolorE] (75pt, 45pt)--(75pt, 35pt);
	\draw[thick, oldiagcolorE] (5pt, 35pt)--(5pt, 45pt);
	\draw[thick, oldiagcolorE] (35pt, 45pt)--(45pt, 45pt);
	\draw[thick, oldiagcolorE] (75pt, 45pt)--(75pt, 35pt); \begin{scope}[xshift=5pt, yshift=35pt]
	\draw [rotate=270, thick, black, l-system={Hilbert curve, step=10pt, angle=90, axiom=L, order=2}]  lindenmayer system;
	\end{scope}
	\begin{scope}[xshift=5pt, yshift=45pt]
	\draw [thick, black, l-system={Hilbert curve, step=10pt, angle=90, axiom=L, order=2}]  lindenmayer system;
	\end{scope}
	\begin{scope}[xshift=45pt, yshift=45pt]
	\draw [thick, black, l-system={Hilbert curve, step=10pt, angle=90, axiom=L, order=2}]  lindenmayer system;
	\end{scope}
	\begin{scope}[xshift=75pt, yshift=5pt]
	\draw [thick, rotate=90,black, l-system={Hilbert curve, step=10pt, angle=90, axiom=L, order=2}]  lindenmayer system;
	\end{scope}
	\end{tikzpicture}
\end{center}
आता सर्व टप्प्यांतील प्राचल सुसंगतपणे ठरवू.
$h_n$ भेट देतो त्या क्रमाने बंद जाळी-घरे
$Q_{n,0},\ldots,Q_{n,4^n-1}$ म्हणा. प्रत्येक घरासाठी
एकक अंतरालाचा समान लांबीचा भाग
\[
I_{n,k}=\left[\frac{k}{4^n},\frac{k+1}{4^n}\right]
\qquad(0\leq k<4^n)
\]
वापरू. या भागाच्या पहिल्या अर्ध्यावर घराच्या प्रवेशबिंदूपासून
मध्यबिंदूपर्यंत आणि दुसऱ्या अर्ध्यावर मध्यबिंदूपासून
निर्गमबिंदूपर्यंत रेषीय रीतीने जाऊ. लागोपाठच्या घरांचा
निर्गम आणि प्रवेश त्यांच्या सामायिक बाजूचा मध्यबिंदू
असतो; संपूर्ण मार्गाचे पहिले आणि शेवटचे बिंदू आकृत्यांतील
लाल भागांच्या बाह्य टोकांवर असतात. त्यामुळे हेच रेखाखंड
मिळतात आणि प्रत्येक $h_n\colon\unitline\to\unitsquare$
सतत असतो.
आकृत्यांतील फिरवलेल्या लहान प्रतींचा क्रम पुनरावृत्तीने
वापरतो: $I_{n,k}$ चे चार समान उपभाग पुढील टप्प्यात
$Q_{n,k}$ च्या चार उपघरांना त्याच भेटीच्या क्रमाने
जोडले जातात. म्हणून प्रत्येक $m\geq n$ साठी
\[
h_m(I_{n,k})\subseteq Q_{n,k}.
\]
येथे बंद घरे वापरली आहेत; अंतरालांचे सामायिक टोक दोन्ही
संबंधित घरांच्या सामायिक सीमेतच जाते.
प्रत्येक $Q_{n,k}$ च्या कर्णाची लांबी
\[
D_n=\sqrt{\left(\nicefrac{1}{2^{n}}\right)^2+
\left(\nicefrac{1}{2^{n}}\right)^2}
=2^{(\frac{1}{2}-n)}
\]
असते. $x\in I_{n,k}$ घेतल्यास $h_m(x)$ आणि $h_n(x)$
हे दोन्ही $Q_{n,k}$ मध्ये आहेत. त्यामुळे
\[
\sup_{x\in\unitline}\|h_m(x)-h_n(x)\|\leq D_n
\qquad(m\geq n).
\]
$D_n\to0$ आणि $\Real^2$ पूर्ण असल्याने प्रत्येक
$x\in\unitline$ साठी $h_n(x)$ ची सीमा अस्तित्वात आहे.
म्हणून
\[
h(x)=\lim_{n\to\infty}h_n(x)
\]
असे व्याख्यित करू शकतो. शिवाय
$\|h(x)-h_n(x)\|\leq D_n$ सर्व $x$ साठी असल्याने
हे अभिसरण समसमान आहे; बंद घरात सीमा घेतल्याने
$h(I_{n,k})\subseteq Q_{n,k}$ सुद्धा मिळते.

आता $h$ अवकाश भरतो हे दाखवू. $p\in\unitsquare$ घ्या.
$p$ धारण करणारे पहिले जाळी-घर निवडा, मग प्रत्येक
टप्प्यात त्याचे $p$ धारण करणारे उपघर निवडा. या
घरांना जोडलेले $I_{n,k_n}$ हे एकमेकांत अंतर्भूत,
रिकामे नसलेले बंद अंतराल आहेत आणि त्यांच्या लांबी
$4^{-n}\to0$ होतात. त्यांच्या छेदात एक बिंदू $x$
असतो. प्रत्येक $n$ साठी $h(x)$ आणि $p$ हे दोन्ही
$Q_{n,k_n}$ मध्ये असल्याने $\|h(x)-p\|\leq D_n$.
$D_n\to0$ म्हणून $h(x)=p$. अशा प्रकारे
$h(\unitline)=\unitsquare$.


आता $h$ हा खरोखर \emph{वक्र} आहे हे दाखवायचे आहे. त्यासाठी वक्राची संकल्पना व्याख्यित करावी लागेल. आधुनिक व्याख्या १८८७ मध्ये जॉर्डन यांनी दिलेल्या व्याख्येवर आधारलेली आहे (म्हणजे पहिला अवकाश भरणारा वक्र येण्याच्या अवघ्या काही वर्षांपूर्वीची).

\begin{defn}
वक्र म्हणजे $\unitline$ वरून $\Real^2$ मध्ये जाणारे सतत फलन.
\end{defn}

ही कल्पना बरीच सहज समजते: वक्र एकक रेषेचे बिंदू
प्रतलातील बिंदूंमध्ये सतत रीतीने पाठवतो. आपले $h$ हे
$\unitline$ वरून $\Real^2$ मध्ये जाणारे फलन आहे.
म्हणून $h$ सतत आहे हेच दाखवायचे उरते.
\ref{his:set:limits:sec} मधील $\epsilon$/$\delta$ व्याख्येनुसार,
प्रत्येक $x_0\in\unitline$ आणि $\epsilon>0$ साठी
असा $\delta>0$ मिळायला हवा की
$x\in\unitline$ आणि $|x-x_0|<\delta$ असल्यास
$\|h(x)-h(x_0)\|<\epsilon$.
अंतरालाच्या टोकांवरही शेजार $\unitline$ च्या सापेक्ष
घ्यायचा आहे.

$x_0$ आणि $\epsilon>0$ निश्चित करा.
$D_n<\epsilon/3$ होईल असा $n$ घ्या.
$h_n$ हा सतत बहुभुजी मार्ग असल्याने असा $\delta>0$
मिळतो की $x\in\unitline$ आणि $|x-x_0|<\delta$
असताना $\|h_n(x)-h_n(x_0)\|<\epsilon/3$.

त्या $x$ साठी त्रिकोण-असमानतेने
\begin{align*}
\|h(x)-h(x_0)\|
&\leq\|h(x)-h_n(x)\|+\|h_n(x)-h_n(x_0)\|
+\|h_n(x_0)-h(x_0)\|\\
&\leq D_n+\|h_n(x)-h_n(x_0)\|+D_n
<\epsilon.
\end{align*}
म्हणून $h$ प्रत्येक $x_0$ येथे सतत आहे आणि व्याख्येनुसार
एक वक्र आहे. आधी सिद्ध केलेल्या
$h(\unitline)=\unitsquare$ मुळे तो अवकाश भरणारा वक्र आहे.





\part{संदर्भसामग्री}\label{ref:part}
\begin{editorial}
  या भागात पाठ्यपुस्तकाच्या परिशिष्टांत समाविष्ट करण्यासाठी वर्णमाला,
  संकेतपद्धती, व्याख्या, नियम इत्यादींच्या विविध याद्या एकत्र केल्या आहेत.
\end{editorial}
\chapter{ग्रीक वर्णमाला}\label{ref:grk:chap}








\begin{center}
    \begin{tabular}{lll@{\qquad}lll}
    अल्फा &	$\alpha$ & $A$ &
    न्यू & $\nu$ & $N$ \\
    बीटा & $\beta$ & $B$ &
    झाय & $\xi$ & $\Xi$ \\
    गॅमा &	$\gamma$ & $\Gamma$ &
    ओमिक्रॉन & $o$ & $O$ \\
    डेल्टा & $\delta$ & $\Delta$ &
    पाय & $\pi$ & $\Pi$ \\
    एप्सिलॉन & $\epsilon$ & $E$ &
    रो & $\rho$ & $P$ \\
    झीटा & $\zeta$ & $Z$ &
    सिग्मा &	$\sigma$ & $\Sigma$ \\
    एटा & $\eta$ & $H$ &
    टाऊ & $\tau$ & $T$ \\
    थीटा & $\theta$ & $\Theta$ &
    अप्सिलॉन & $\upsilon$ & $\Upsilon$ \\
    आयोटा & $\iota$ & $I$ &
    फाय & $\varphi$ & $\Phi$ \\
    कॅप्पा & $\kappa$ & $K$ &
    खाय & $\chi$ & $X$ \\
    लॅम्ब्डा & $\lambda$ & $\Lambda$ &
    साय & $\psi$ & $\Psi$ \\
    म्यू & $\mu$ & $M$ &
    ओमेगा & $\omega$ & $\Omega$
    \end{tabular}
\end{center}



\chapter{फ्राक्टूर वर्णमाला}\label{ref:frk:chap}








\begin{center}
    \begin{tabular}{llll@{\qquad}llll}
    A &	$\mathfrak{A}$ & a & $\mathfrak{a}$ &
    N &	$\mathfrak{N}$ & n & $\mathfrak{n}$ \\
    B &	$\mathfrak{B}$ & b & $\mathfrak{b}$ &
    O &	$\mathfrak{O}$ & o & $\mathfrak{o}$ \\
    C &	$\mathfrak{C}$ & c & $\mathfrak{c}$ &
    P &	$\mathfrak{P}$ & p & $\mathfrak{p}$ \\
    D &	$\mathfrak{D}$ & d & $\mathfrak{d}$ &
    Q &	$\mathfrak{Q}$ & q & $\mathfrak{q}$ \\
    E &	$\mathfrak{E}$ & e & $\mathfrak{e}$ &
    R &	$\mathfrak{R}$ & r & $\mathfrak{r}$ \\
    F &	$\mathfrak{F}$ & f & $\mathfrak{f}$ &
    S &	$\mathfrak{S}$ & s & $\mathfrak{s}$ \\
    G &	$\mathfrak{G}$ & g & $\mathfrak{g}$ &
    T &	$\mathfrak{T}$ & t & $\mathfrak{t}$ \\
    H &	$\mathfrak{H}$ & h & $\mathfrak{h}$ &
    U &	$\mathfrak{U}$ & u & $\mathfrak{u}$ \\
    I &	$\mathfrak{I}$ & i & $\mathfrak{i}$ &
    V &	$\mathfrak{V}$ & v & $\mathfrak{v}$ \\
    J &	$\mathfrak{J}$ & j & $\mathfrak{j}$ &
    W &	$\mathfrak{W}$ & w & $\mathfrak{w}$ \\
    K &	$\mathfrak{K}$ & k & $\mathfrak{k}$ &
    X &	$\mathfrak{X}$ & x & $\mathfrak{x}$ \\
    L &	$\mathfrak{L}$ & l & $\mathfrak{l}$ &
    Y &	$\mathfrak{Y}$ & y & $\mathfrak{y}$ \\
    M &	$\mathfrak{M}$ & m & $\mathfrak{m}$ &
    Z &	$\mathfrak{Z}$ & z & $\mathfrak{z}$
    \end{tabular}
    \end{center}



\part{पूरक आणि पर्यायी विभाग}\label{supp:part}
\chapter{मूळ संकलनमार्गाबाहेरील विभाग}\label{supp:chap}
या विभागांचे स्रोत-एककांमध्ये भाषांतर केलेले आहे. मुख्य संकलनमार्गात नसलेली सामग्री येथे विषयानुसार पूरक स्वरूपात दिली आहे.












\section{निष्पन्नता चे गुणधर्म}\label{fol:axd:prv-alt:sec}

\begin{prop}[एकस्वनिकता] जर $\Gamma \subseteq \Delta$ आणि $\Gamma \Proves \varphi$
असेल, तर $\Delta \Proves \varphi$ सुद्धा. \end{prop}

\begin{proof} कोणताही परिमित $\Gamma_0 \subseteq \Gamma$ हा~$\Delta$ चासुद्धा
परिमित उपसंच असतो. म्हणून $\Gamma_0 \Proves \varphi$ ची निष्पत्ती
$\Delta \Proves \varphi$ हेही दाखवते. \end{proof}

\begin{prop}\label{fol:axd:prv-alt:prop:provability} $\Gamma \Proves \varphi$ तेव्हा आणि केवळ
तेव्हाच, जेव्हा $\Gamma\cup \{\lnot\varphi \}$ विसंगत असतो. \end{prop}

\begin{prop} \begin{enumerate}
\item \label{fol:axd:prv-alt:prop:provability-contr} जर $\Gamma \Proves \varphi$ आणि $\Gamma
\cup \{ \varphi\} \Proves \lfalse$ असेल, तर $\Gamma$ विसंगत असतो.

\item \label{fol:axd:prv-alt:prop:provability-lnot} जर $\Gamma \cup \{\varphi\}
\Proves \lfalse$ असेल, तर $\Gamma \Proves \lnot \varphi$.

\item \label{fol:axd:prv-alt:prop:provability-exhaustive} जर $\Gamma \cup \{\varphi\}
\Proves \lfalse$ आणि $\Gamma \cup \{\lnot \varphi\} \Proves
\lfalse$ असेल, तर $\Gamma \Proves \lfalse$.

\item \label{fol:axd:prv-alt:prop:provability-lor-left} जर $\Gamma \cup \{\varphi\}
\Proves \lfalse$ आणि $\Gamma \cup \{\psi\} \Proves
\lfalse$ असेल, तर $\Gamma \cup \{\varphi \lor \psi\} \Proves \lfalse$.

\item \label{fol:axd:prv-alt:prop:provability-lor-right} जर $\Gamma
\Proves \varphi$ किंवा $\Gamma \Proves \psi$ असेल, तर $\Gamma
\Proves \varphi \lor \psi$.

\item \label{fol:axd:prv-alt:prop:provability-land-left} जर $\Gamma
\Proves \varphi \land \psi$ असेल, तर $\Gamma \Proves \varphi$ आणि
$\Gamma \Proves \psi$.

\item \label{fol:axd:prv-alt:prop:provability-land-right} जर $\Gamma
\Proves \varphi$ आणि $\Gamma \Proves \psi$ असेल, तर $\Gamma
\Proves \varphi \land \psi$.

\item \label{fol:axd:prv-alt:prop:provability-mp} जर $\Gamma \Proves
\varphi$ आणि $\Gamma \Proves \varphi \lif \psi$ असेल, तर $\Gamma
\Proves \psi$.

\item \label{fol:axd:prv-alt:prop:provability-lif} जर $\Gamma \Proves
\lnot \varphi$ किंवा $\Gamma \Proves \psi$ असेल, तर $\Gamma \Proves
\varphi \lif \psi$. \end{enumerate} \end{prop}

\begin{proof}
\begin{enumerate}
\item $\Gamma_0 \cup \{\varphi\} \Proves \lfalse$ असे गृहीत धरा.

\begin{derivation}
1. & $\Gamma_0 \cup \{\varphi\} \Proves \lfalse$ & गृहीत \\
2. & $\Gamma_1 \Proves \varphi$ & गृहीत\\
3. & $\Gamma_0 \cup \Gamma_1 \Proves \lfalse$ & कट नियम \\
\end{derivation}

कारण $\Gamma_0 \subseteq \Gamma$ आणि $\Gamma_1 \subseteq \Gamma$,
$\Gamma_0 \cup \Gamma_1 \subseteq \Gamma$; म्हणून $\Gamma \Proves \lfalse$.

\item $\Gamma_0 \cup\{\varphi\} \Proves \lfalse$ असे गृहीत धरा.

\begin{derivation}
1. & $\Gamma_0 \cup\{\varphi\} \Proves \lfalse$ & गृहीत \\
2. & $\Gamma_0 \Proves \varphi \lif \lfalse$ & निगमन प्रमेय\\
3. & $\Gamma_0 \Proves (\varphi \lif \lfalse) \lif \lnot \varphi$ & Ax0\\
4. & $\Gamma_0 \Proves \lnot \varphi$ & MP 2, 3 \\
\end{derivation}


\item $\Gamma_0 \cup \{ \varphi \} \Proves \lfalse$ आणि
$\Gamma_1 \cup \{\lnot \varphi \} \Proves \lfalse$ असे गृहीत धरा.


\begin{derivation}
1. & $\Gamma_0 \cup\{\varphi\} \Proves \lfalse$ & गृहीत \\
2. & $\Gamma_0 \Proves \varphi \lif \lfalse$ & निगमन प्रमेय \\
3. & $\Gamma_1 \cup \{\lnot \varphi \} \Proves \lfalse$ & गृहीत\\
4. & $\Gamma_0 \Proves (\varphi \lif \lfalse) \lif \lnot \varphi$ & Ax0\\
5. & $\Gamma_0 \Proves \lnot \varphi$ & MP 2, 4\\
6. & $\Gamma_0 \cup \Gamma_1 \Proves \lfalse$ & कट नियम 3, 5 \\
\end{derivation}

कारण $\Gamma_0 \subseteq \Gamma$ आणि $\Gamma_1 \subseteq \Gamma$,
$\Gamma_0 \cup \Gamma_1 \subseteq \Gamma$. म्हणून $\Gamma \Proves \lfalse$.


\item
सराव.


\item
$\Gamma_0 \Proves \varphi$ असे गृहीत धरा.

\begin{derivation}
1. & $\Gamma_0 \Proves \varphi$ & गृहीत \\
2. & $\Gamma_0 \Proves \varphi \lif (\varphi \lor \psi)$ & Ax0\\
3. & $\Gamma_0 \Proves \varphi \lor \psi$ & MP 1, 2 \\
\end{derivation}

म्हणून $\Gamma \Proves \varphi \lor \psi$. $\Gamma \Proves \psi$ असल्यास
सिद्धताही याचप्रमाणे मिळते.


\item
$\Gamma_0 \Proves \varphi \land \psi$ असे गृहीत धरा

\begin{derivation}
1. & $\Gamma_0 \Proves \varphi \land \psi$ & गृहीत \\
2. & $\Gamma_0 \Proves (\varphi\land \psi) \lif \varphi $ & Ax0\\
3. & $\Gamma_0 \Proves \varphi$ & MP 1, 2 \\
\end{derivation}


म्हणून $\Gamma \Proves \varphi$. याचप्रमाणे सुरू होणारे निष्पत्ती
$\Gamma \Proves \psi$ हे दाखवते.


\item सराव.


\item   $\Gamma_0 \Proves \varphi$ आणि $\Gamma_0 \Proves \varphi \lif \psi $ ही दोन्ही गृहीत धरा.

\begin{derivation}
1. & $\Gamma_0 \Proves \varphi$ & गृहीत \\
2. & $\Gamma_0 \Proves \varphi \lif \psi $ & गृहीत\\
3. & $\Gamma_0 \Proves \psi$ & MP 1, 2\\
\end{derivation}

कारण $\Gamma_0 \subseteq \Gamma$, यावरून $\Gamma
\Proves \psi$ मिळते.



\item प्रथम $\Gamma \Proves \lnot \varphi$ असे माना.

\begin{derivation}
1. & $\Gamma_0 \Proves \lnot \varphi$ & गृहीत \\
2. & $\Gamma_0 \Proves \lnot \varphi \lif (\varphi \lif \psi) $ & Ax0\\
3. & $\Gamma_0 \Proves \varphi \lif \psi$ & MP 1,
2\\ \end{derivation}

$\Gamma \Proves \psi$ यासाठीही सिद्धता याचप्रमाणे आहे:

\begin{derivation}
1. & $\Gamma_0 \Proves \psi$ & गृहीत \\
2. & $\Gamma_0 \Proves \psi \lif (\varphi \lif \psi) $ & Ax0\\
3. & $\Gamma_0 \Proves \varphi \lif \psi$ & MP 1, 2\\
\end{derivation}

\end{enumerate}
\end{proof}



\begin{thm}[व्यापकीकरण] \label{fol:axd:prv-alt:thm:Generalization} जर $\Gamma \Proves
\varphi$ असेल आणि $x$ हे $\Gamma$ मधील कोणत्याही सूत्रामध्ये मुक्त नसेल,
तर $\Gamma \Proves \lforall[x][\varphi]$.
\end{thm}

\begin{proof} $\mathsf{Thm}_\Gamma$ वर विगमन वापरा: जर $\varphi$ स्वयंसिद्धक
असेल, तर $\lforall[x][\varphi]$ हेसुद्धा स्वयंसिद्धक असते. जर $\varphi\in \Gamma$
असेल, तर $x$ हे $\varphi$ मध्ये मुक्त नसते. म्हणून $\varphi \lif \lforall[x][\varphi]$
हे स्वयंसिद्धक (\textbf{Ax3}) आहे आणि विध्यात्मक अनुमानाने $\Gamma
\Proves \lforall[x][\varphi]$. आता $\Gamma \Proves \psi$ आणि $\Gamma \Proves
\psi \lif \varphi$ यांमुळे विध्यात्मक अनुमानाने $\varphi$ मिळते असे माना. विगमन
गृहीतकानुसार $\Gamma \Proves \lforall[x][\psi]$ आणि $\Gamma \Proves
\lforall[x][( \psi \lif \varphi)]$. परंतु \textbf{Ax2} मुळे पुढीलही मिळते:
\[ \Gamma \Proves
\lforall[x][( \psi \lif \varphi)] \lif (\lforall[x][ \psi] \lif \lforall[x][\varphi]), \]
आणि विध्यात्मक अनुमान दोनदा वापरल्यावर अपेक्षेप्रमाणे $\Gamma \Proves
\lforall[x][\varphi]$ मिळते. \end{proof}

\begin{thm}\label{fol:axd:prv-alt:thm:WeakGen} (\emph{स्थिरांकांवरील दुर्बल व्यापकीकरण})
जर $\Gamma \Proves \varphi$ असेल आणि $c$ हा $\Gamma$ मध्ये न येणारा
स्थिरांक असेल, तर $\varphi$ मध्ये न येणारा असा चल $x$ असतो की
$\Gamma \Proves \lforall[x][\Subst{\varphi}{x}{c}]$; आणि सिद्धतेत $c$ येत नाही.
\end{thm}

\begin{proof} $\varphi_1,\ldots,\varphi_n$ ही $\Gamma$ पासून $\varphi$ ची सिद्धता
असू द्या; म्हणजे $\varphi=\varphi_n$. $\varphi_1,\ldots,\varphi_n$ यांमध्ये न येणारा चल $x$
निवडा आणि पुढील नवा अनुक्रम पाहा:
$\Subst{\varphi_1}{x}{c},\ldots,\Subst{\varphi_n}{x}{c}.$ हा अनुक्रम
$\Gamma$ पासून $\Subst{\varphi}{x}{c}$ ची सिद्धता असतो. खरेच, प्रत्येक
$i=1,\ldots,n$ साठी: \begin{itemize}
\item जर $\varphi_i$ स्वयंसिद्धक असेल, तर $\Subst{\varphi_i}{x}{c}$ हेसुद्धा स्वयंसिद्धक असते;
\item जर $\varphi_i \in \Gamma$ असेल, तर $c$ हे $\Gamma$ मध्ये नसल्यामुळे
$\Subst{\varphi_i}{x}{c} = \varphi_i \in \Gamma$;
\item जर $\varphi_i$ हे $\varphi_j$ आणि $\varphi_j \lif \varphi_i$ यांपासून मिळाले असेल,
तर $\Subst{\varphi_i}{x}{c}$ हे $\Subst{\varphi_j}{x}{c}$ आणि
$\Subst{(\varphi_j \lif \varphi_i)}{x}{c}
= \Subst{\varphi_j}{x}{c} \lif \Subst{\varphi_i}{x}{c}$ यांपासून विध्यात्मक
अनुमानाने मिळते. \end{itemize}
नव्या अनुक्रमात स्थिरांक $c$ उरलेला नाही, हे स्पष्ट आहे. आता
$\Subst{\varphi}{x}{c}$ च्या निष्पत्तीमध्ये येणाऱ्या $\Gamma$ मधील
सूत्रांचा संच $\Gamma'$ असू द्या. मग $\Gamma'$ मधील कोणत्याही
सूत्रामध्ये $x$ मुक्त नाही आणि $\Gamma' \Proves \Subst{\varphi}{x}{c}$.
व्यापकीकरणाने (\ref{fol:axd:prv-alt:thm:Generalization}) $\Gamma'
\Proves \lforall[x][\Subst{\varphi}{x}{c}]$; म्हणून एकस्वनिकतेने
अपेक्षेप्रमाणे $\Gamma \Proves
\lforall[x][\Subst{\varphi}{x}{c}]$. \end{proof}


\begin{explain}
\ref{fol:axd:prv-alt:thm:WeakGen} मधील $x$ हा $\varphi$ मध्ये (अजिबात) येत नाही, ही अट
कमकुवत करून $x$ हा $\varphi$ मध्ये मुक्त नाही आणि $\varphi$ मध्ये $c$ च्या जागी
बसवण्यासाठी आदेशनासाठी मुक्त आहे, इतक्याच अटी ठेवणे हे आपले उद्दिष्ट आहे.
बद्ध चल बदलल्यास हे करता येते; म्हणून प्रथम तो बदल सिद्ध करू.
\end{explain}

\begin{lem}
\label{fol:axd:prv-alt:lem:free-for}
जर $x$ हा $\varphi$ मध्ये $c$ च्या जागी बसवण्यासाठी मुक्त असेल आणि $y$
हा $\varphi$ मध्ये मुक्त नसेल, तर $x$ हा $\Subst{\varphi}{y}{c}$ मध्ये
$y$ च्या जागी बसवण्यासाठी आदेशनासाठी मुक्त असतो.
\end{lem}

\begin{proof} जर $x$ हा $\Subst{\varphi}{y}{c}$ मध्ये $y$ च्या जागी
बसवण्यासाठी आदेशनासाठी मुक्त नसेल, तर $\Subst{\varphi}{y}{c}$ मधील $y$ चा एखादा
मुक्त उद्भव $\lforall[x]$ या संख्यापकाच्या कार्यक्षेत्रात येतो. पण
गृहीतकानुसार $y$ हा $\varphi$ मध्ये मुक्त नाही; म्हणून असे सर्व उद्भव
$\Subst{}{y}{c}$ या आदेशनातूनच येतात. अशा रीतीने $c$ चा एखादा
उद्भव $\lforall[x]$ च्या कार्यक्षेत्रात येतो आणि $x$ हा $\varphi$ मध्ये
$c$ च्या जागी बसवण्यासाठी आदेशनासाठी मुक्त राहत नाही.
\end{proof}


\begin{lem}[बद्ध चराचा बदल]
\label{fol:axd:prv-alt:lem:ChangeBdVar}
जर $x$ आणि $y$ हे $\varphi$ मध्ये मुक्त नसतील आणि $\varphi$ मध्ये $c$ च्या जागी
बसवण्यासाठी दोन्ही आदेशनासाठी मुक्त असतील, तर
$\Proves \lforall[x][\Subst{\varphi}{x}{c}] \equiv
\lforall[y][\Subst{\varphi}{y}{c}]$ (आणि सिद्धतेत $c$ येत नाही).
\end{lem}

\begin{proof}
गृहीतकांनुसार $y$ हा $\varphi$ मध्ये $c$ च्या जागी बसवण्यासाठी आदेशनासाठी मुक्त
आहे आणि $x$ हा $\varphi$ मध्ये मुक्त नाही. म्हणून मागील
\ref{fol:axd:prv-alt:lem:free-for} वरून $y$ हा $\Subst{\varphi}{x}{c}$ मध्ये $x$ च्या
जागी बसवण्यासाठी आदेशनासाठी मुक्त आहे. त्यामुळे
$\lforall[x][\Subst{\varphi}{x}{c}] \lif \Subst{\varphi}{x}{c} \Subst{}{y}{x}$
हे स्वयंसिद्धक (\textbf{Ax1}) आहे. $x$ हा $\varphi$ मध्ये मुक्त नसल्यामुळे
$\Subst{\varphi}{x}{c} \Subst{}{y}{x} = \Subst{\varphi}{y}{c}$; म्हणून
\[
\Proves \lforall[x][\Subst{\varphi}{x}{c}] \lif \Subst{\varphi}{y}{c},
\]
आणि निगमन प्रमेयाने $\lforall[x][\Subst{\varphi}{x}{c}] \Proves
\Subst{\varphi}{y}{c}$. $y$ हा $\varphi$ मध्ये मुक्त नसल्यामुळे तो
$\lforall[x][\Subst{\varphi}{x}{c}]$ मध्येही मुक्त नाही. त्यामुळे
व्यापकीकरणाने $\lforall[x][\Subst{\varphi}{x}{c}] \Proves
\lforall[y][\Subst{\varphi}{y}{c}]$; पुन्हा निगमन प्रमेयाने
$\Proves \lforall[x][\Subst{\varphi}{x}{c}] \lif
\lforall[y][\Subst{\varphi}{y}{c}]$. उलट दिशेची
$\Proves \lforall[y][\Subst{\varphi}{y}{c}] \lif \lforall[x]
[\Subst{\varphi}{x}{c}]$ ही सिद्धता पूर्णपणे सममित आहे; त्यामुळे
निष्कर्ष Taut वापरून मिळतो.
\end{proof}

\begin{thm}[स्थिरांकांवरील प्रबळ व्यापकीकरण]
\label{fol:axd:prv-alt:thm:StrongGen}
जर $\Gamma \Proves \varphi$ असेल, स्थिरांक $c$ हा $\Gamma$ मध्ये येत
नसेल, आणि $x$ हा $\varphi$ मध्ये मुक्त नसेल पण $\varphi$ मध्ये $c$ च्या जागी
बसवण्यासाठी आदेशनासाठी मुक्त असेल, तर
$\Gamma \Proves \lforall[x][\Subst{\varphi}{x}{c}]$.
\end{thm}

\begin{proof}
$\Gamma \Proves \varphi$ आणि $c$ हा $\Gamma$ मध्ये नाही; म्हणून दुर्बल
व्यापकीकरणाने $\varphi$ मध्ये न येणारा चल $y$ असतो की
$\Gamma \Proves \lforall[y][\Subst{\varphi}{y}{c}]$. $y$ हा $\varphi$ मध्ये
(अजिबात) येत नसल्यामुळे तो $\varphi$ मध्ये मुक्तही नसतो आणि $\varphi$ मध्ये
$c$ च्या जागी बसवण्यासाठी आदेशनासाठी मुक्त असतो. गृहीतकानुसार $x$ हाही
$\varphi$ मध्ये मुक्त नाही आणि $\varphi$ मध्ये $c$ च्या जागी बसवण्यासाठी आदेशनासाठी मुक्त
आहे; म्हणून बद्ध चल बदलण्याच्या अटी पूर्ण होतात. त्यामुळे
$\Proves \lforall[x][\Subst{\varphi}{x}{c}] \equiv
\lforall[y][\Subst{\varphi}{y}{c}]$; म्हणून
$\Gamma \Proves \lforall[x][\Subst{\varphi}{x}{c}]$.
\end{proof}

\begin{thm}
\label{fol:axd:prv-alt:thm:provability-quantifiers}
\begin{enumerate}
\item जर $\Gamma \Proves \varphi(t)$ असेल, तर $\Gamma \Proves
  \lexists[x][\varphi(x)]$.
\item जर $\Gamma \Proves \lforall[x][\varphi(x)]$ असेल, तर $\Gamma
  \Proves \varphi(t)$.
\end{enumerate}
\end{thm}

\begin{proof}
\begin{enumerate}

\item सराव.
\item असे गृहीत धरा की  $\Gamma_0 \Proves \lforall[x][\varphi(x)]$.

\begin{derivation}
1. & $\Gamma_0 \Proves \lforall[x][\varphi(x)]$ & गृहीत \\
2. & $\Gamma_0 \Proves \lforall[x][\varphi(x)] \lif \Subst{\varphi}{t}{x}$ & Ax1 \\
3. & $\Gamma_0 \Proves \Subst{\varphi}{t}{x}$ & MP 1, 2 \\
\end{derivation}

\end{enumerate}
\end{proof}



\begin{quote}\small\textbf{स्रोतनोंद.} SOL6-A285: निषेधाच्या गुणधर्माची सिद्धता आता त्याच्या स्वतःच्या गृहीतकापासून not-A काढते. बद्ध चर बदलातील सार्विक संख्यापक फक्त अभिव्यंजनाच्या पूर्वांगावर आहे; मूळ संपूर्ण अभिव्यंजनावरील व्याप्ती योग्य आदेशन-स्वयंसिद्धक नव्हती.\end{quote}
\begin{quote}\small\textbf{स्रोतनोंद.} OLINC-321: मूळ हिशोब निषेधाच्या विधानाऐवजी आधीच्या समतुल्यतेची उलट दिशा सिद्ध करत होता. SOL6-A285 मध्ये A या अतिरिक्त गृहीतकापासून असत्य मिळते या आधारावर not-A काढला आहे.\end{quote}
\begin{quote}\small\textbf{स्रोतनोंद.} OLINC-326: पहिल्या गृहीतकातील आधारसंच पुढील हिशोबाप्रमाणे Gamma\_0 केला आहे.\end{quote}
\begin{quote}\small\textbf{स्रोतनोंद.} OLINC-322: A आणि B यांच्या संयुक्तातून A काढणाऱ्या हिशोबात स्रोताच्या तिसऱ्या ओळीत चुकीने A किंवा B आहे.\end{quote}
\begin{quote}\small\textbf{स्रोतनोंद.} OLINC-323: या हिशोबात Gamma\_1 परिभाषित नाही; वापरलेला परिमित आधारसंच Gamma\_0 आहे.\end{quote}
\begin{quote}\small\textbf{स्रोतनोंद.} OLINC-325: सिद्धता-अनुक्रमाच्या प्रत्येक टप्प्यातील स्वयंसिद्धक A\_i आहे; स्रोताने A म्हटले आहे.\end{quote}
\begin{quote}\small\textbf{स्रोतनोंद.} OLINC-324: पहिल्या आदेशनात स्रोताने c ऐवजी x लिहिला आहे; उरलेल्या विधानानुसार c केला आहे.\end{quote}











\section{वाक्य यांचे महत्तम सुसंगत संच}\label{fol:com:mcs:sec}

\begin{defn}[महत्तम सुसंगत संच]
\label{fol:com:mcs:mcs}
वाक्यांचा संच~$\Gamma$ \emph{महत्तम सुसंगत} तेव्हा आणि केवळ तेव्हाच असतो, जेव्हा
\begin{enumerate}
\item $\Gamma$ सुसंगत असतो, आणि
\item जर $\Gamma \subsetneq \Gamma'$ असेल, तर $\Gamma'$ विसंगत असतो.
\end{enumerate}
\end{defn}

\begin{explain}
वरील व्याख्येशी समतुल्य अशी पर्यायी व्याख्या ही: वाक्यांचा संच $\Gamma$
\emph{महत्तम सुसंगत} तेव्हा आणि केवळ तेव्हाच असतो, जेव्हा
\begin{enumerate}
\item $\Gamma$ सुसंगत असतो, आणि
\item जर $\Gamma \cup \{ \varphi \}$ सुसंगत असेल, तर $\varphi \in \Gamma$ असते.
\end{enumerate}
म्हणजेच~$\Gamma$ मध्ये आधीपासून नसलेली वाक्ये महत्तम सुसंगत
संच~$\Gamma$ मध्ये जोडली, तर मिळणारा मोठा संच विसंगत होतो.
\end{explain}

\begin{explain}
संपूर्णतेच्या सिद्धतेत महत्तम सुसंगत संच महत्त्वाचे आहेत. कारण प्रत्येक
सुसंगत वाक्यांचा संच~$\Gamma$ हा एखाद्या महत्तम सुसंगत
संचात~$\Gamma^*$ सामावतो, याची खात्री देता येते. शिवाय प्रत्येक
वाक्य~$\varphi$ साठी महत्तम सुसंगत संचात $\varphi$ किंवा त्याचा
निषेध $ \lnot \varphi$ यांपैकी एखादे वाक्य असते. विशेषतः अणुरूप वाक्ये साठी हे खरे
आहे. म्हणून स्थिरांक घालून योग्य रीतीने विस्तारलेल्या भाषेतील
महत्तम सुसंगत संचापासून आपण संरचना रचू शकतो: कोणती अणुरूप
वाक्ये~$\Gamma^*$ मध्ये आहेत यावरून विधेयांचे
अर्थनिर्वचन ठरवतो. मग या संरचनांमध्ये $\Gamma^*$ मधील
सर्व वाक्ये (आणि त्यामुळे $\Gamma$ मधीलसुद्धा) सत्य ठरतात,
हे दाखवता येते. या शेवटच्या विधानाच्या सिद्धतेसाठी
$\lnot \varphi \in \Gamma^*$ तेव्हा आणि केवळ तेव्हाच $\varphi \notin
\Gamma^*$, तसेच $(\varphi \lor \psi) \in \Gamma^*$ तेव्हा आणि केवळ तेव्हाच
$\varphi \in \Gamma^*$ किंवा $\psi \in \Gamma^*$, इत्यादी गुणधर्म लागतात.
\end{explain}

\begin{prop}
\label{fol:com:mcs:prop:mcs}
$\Gamma$ महत्तम सुसंगत आहे असे माना. मग:
\begin{enumerate}
\item \label{fol:com:mcs:prop:mcs-prov-in} जर $\Gamma \Proves \varphi$ असेल, तर $\varphi \in
  \Gamma$.

\item \label{fol:com:mcs:prop:mcs-either-or} कोणत्याही $\varphi$ साठी $\varphi \in
  \Gamma$ किंवा $\lnot \varphi \in \Gamma$.

\item \label{fol:com:mcs:prop:mcs-and} $(\varphi \land \psi) \in \Gamma$
  तेव्हा आणि केवळ तेव्हाच $\varphi \in \Gamma$ आणि $\psi \in \Gamma$ ही दोन्ही.

\item \label{fol:com:mcs:prop:mcs-or} $(\varphi \lor \psi) \in \Gamma$
  तेव्हा आणि केवळ तेव्हाच $\varphi \in \Gamma$ किंवा $\psi \in \Gamma$.

\item \label{fol:com:mcs:prop:mcs-if} $(\varphi \lif \psi) \in \Gamma$
  तेव्हा आणि केवळ तेव्हाच $\varphi \notin \Gamma$ किंवा $\psi \in \Gamma$.
\end{enumerate}
\end{prop}

\begin{proof}
पुढील सर्व भागांत $\Gamma$ महत्तम सुसंगत आहे असे मानू.
\begin{enumerate}
\item जर $\Gamma \Proves \varphi$ असेल, तर $\varphi \in \Gamma$.

$\Gamma \Proves \varphi$ असे माना. उलट $\varphi \notin \Gamma$ असेही
मानल्यास, $\Gamma$ महत्तम सुसंगत असल्यामुळे $\Gamma \cup \{\varphi\}$
विसंगत असतो; म्हणजे $\Gamma \cup \{\varphi\} \Proves \lfalse$.
\ref{fol:seq:prv:prop:provability-contr}
वरून $\Gamma$ विसंगत ठरतो. पण $\Gamma$ सुसंगत आहे या गृहीतकाला
हा व्याघात आहे. म्हणून $\varphi \notin \Gamma$ शक्य नाही; म्हणजे
$\varphi \in \Gamma$.

\item कोणत्याही $\varphi$ साठी $\varphi \in \Gamma$ किंवा $\lnot \varphi \in \Gamma$.

उलट, एखाद्या~$\varphi$ साठी $\varphi \notin \Gamma$ आणि
$\lnot \varphi \notin \Gamma$ असे दोन्ही आहेत असे माना. $\Gamma$ महत्तम
सुसंगत असल्यामुळे $\Gamma \cup \{\varphi\}$ आणि $\Gamma \cup \{\lnot \varphi\}$
दोन्ही विसंगत आहेत. म्हणून $\Gamma \cup \{\varphi\} \Proves \lfalse$
आणि $\Gamma \cup \{\lnot \varphi\} \Proves \lfalse$.
\ref{fol:seq:prv:prop:provability-exhaustive}
वरून $\Gamma$ विसंगत ठरतो, हा व्याघात. म्हणून असे वाक्य~$\varphi$
असू शकत नाही. प्रत्येक $\varphi$ साठी $\varphi \in \Gamma$ किंवा
$\lnot \varphi \in \Gamma$.

\item

$(\varphi \land \psi) \in \Gamma$ तेव्हा आणि केवळ तेव्हाच $\varphi \in \Gamma$
आणि $\psi \in \Gamma$ ही दोन्ही:

पुढील दिशेसाठी $(\varphi \land \psi) \in \Gamma$ असे माना. मग
$\Gamma \Proves \varphi \land \psi$. आता
\ref{fol:seq:ppr:prop:provability-land-left}
वरून $\Gamma \Proves \varphi$ आणि $\Gamma \Proves \psi$. तसेच
\ref{fol:com:mcs:prop:mcs-prov-in} वरून अपेक्षेप्रमाणे $\varphi \in \Gamma$
आणि $\psi \in \Gamma$.

उलट दिशेसाठी $\varphi \in \Gamma$ आणि $\psi \in \Gamma$ असे माना.
मग $\Gamma \Proves \varphi$ आणि $\Gamma \Proves \psi$. आता
\ref{fol:seq:ppr:prop:provability-land-right}
वरून $\Gamma \Proves \varphi \land \psi$. \ref{fol:com:mcs:prop:mcs-prov-in}
वरून $(\varphi \land \psi) \in \Gamma$.

\item

$(\varphi \lor \psi) \in \Gamma$ तेव्हा आणि केवळ तेव्हाच $\varphi \in \Gamma$
किंवा $\psi \in \Gamma$.

पुढील दिशेचा प्रतिव्यत्यास सिद्ध करण्यासाठी $\varphi \notin \Gamma$
आणि $\psi \notin \Gamma$ असे माना. आपल्याला $(\varphi \lor \psi)
\notin \Gamma$ दाखवायचे आहे. $\Gamma$ महत्तम सुसंगत असल्यामुळे
$\Gamma \cup \{\varphi\} \Proves \lfalse$ आणि
$\Gamma \cup \{\psi\} \Proves \lfalse$. आता
\ref{fol:axd:prv-alt:prop:provability-lor-left}
वरून $\Gamma \cup \{(\varphi \lor \psi)\}$ विसंगत ठरतो. म्हणून
अपेक्षेप्रमाणे $(\varphi \lor \psi) \notin \Gamma$.

उलट दिशेसाठी $\varphi \in \Gamma$ किंवा $\psi \in \Gamma$ असे
माना. मग $\Gamma \Proves \varphi$ किंवा $\Gamma \Proves \psi$. आता
\ref{fol:axd:prv-alt:prop:provability-lor-right}
वरून $\Gamma \Proves \varphi \lor \psi$. \ref{fol:com:mcs:prop:mcs-prov-in}
वरून अपेक्षेप्रमाणे $(\varphi \lor \psi) \in \Gamma$.

\item

$(\varphi \lif \psi) \in \Gamma$ तेव्हा आणि केवळ तेव्हाच
$\varphi \notin \Gamma$ किंवा $\psi \in \Gamma$:

पुढील दिशेसाठी $(\varphi \lif \psi) \in \Gamma$ असे माना.
उलट $\varphi \in \Gamma$ आणि $\psi \notin \Gamma$ असेही माना.
या गृहीतकांवरून $\Gamma \Proves \varphi \lif \psi$ आणि
$\Gamma \Proves \varphi$. आता
\ref{fol:axd:prv-alt:prop:provability-mp}
वरून $\Gamma \Proves \psi$. परंतु \ref{fol:com:mcs:prop:mcs-prov-in} वरून
$\psi \in \Gamma$, हा $\psi \notin \Gamma$ या गृहीतकाला व्याघात.

उलट दिशेसाठी प्रथम $\varphi \notin \Gamma$ हे प्रकरण पाहू.
\ref{fol:com:mcs:prop:mcs-either-or} वरून $\lnot \varphi \in \Gamma$ आणि म्हणून
$\Gamma \Proves \lnot \varphi$. आता
\ref{fol:seq:ppr:prop:provability-lif}
वरून $\Gamma \Proves \varphi \lif \psi$.
पुन्हा \ref{fol:com:mcs:prop:mcs-prov-in} वरून अपेक्षेप्रमाणे
$(\varphi \lif \psi) \in \Gamma$.

आता $\psi \in \Gamma$ हे प्रकरण पाहू. मग $\Gamma \Proves \psi$,
आणि
\ref{fol:seq:ppr:prop:provability-lif}
वरून $\Gamma \Proves \varphi \lif \psi$. \ref{fol:com:mcs:prop:mcs-prov-in}
वरून $(\varphi \lif \psi) \in \Gamma$.
\end{enumerate}
\end{proof}

\begin{prob}
\ref{fol:com:mcs:prop:mcs} ची सिद्धता पूर्ण करा.
\end{prob}












\section{परिचय}\label{fol:syn:int:sec}

प्रथम-क्रम तर्कशास्त्राची उपपत्ती आणि अधिसिद्धांत विकसित करण्यासाठी,
आपण प्रथम त्याच्या व्यंजकांची विन्यासमीमांसा आणि चिन्हार्थमीमांसा
परिभाषित केली पाहिजे. प्रथम-क्रम तर्कशास्त्राची व्यंजके म्हणजे पदे आणि
सूत्रे. पदे चले, स्थिरांक आणि फलने
यांपासून तयार होतात. सूत्रे मात्र विधेये आणि पदे
यांपासून तयार होतात (यांतून सर्वांत लहान, ``आण्विक'' सूत्रे
तयार होतात); मग तार्किक संयोजके आणि संख्यापक वापरून आण्विक
सूत्रांपासून अधिक जटिल सूत्रे तयार करता येतात. रचनानियम
मांडण्याचे अनेक निरनिराळे मार्ग आहेत; आपण त्यांपैकी केवळ एक शक्य
मार्ग देतो. इतर पद्धती वेगळी चिन्हे निवडतील, संयोजकांचे वेगळे संच
आदिम म्हणून निवडतील आणि कंसांचा वेगळ्या रीतीने वापर करतील (किंवा
तथाकथित पोलिश संकेतपद्धतीप्रमाणे कंस अजिबात वापरणार नाहीत). तरी सर्व
दृष्टिकोनांत सामाईक बाब ही की रचनानियम पदांचा आणि सूत्रांचा
संच \emph{विगमनाने} परिभाषित करतात. हे योग्य रीतीने केल्यास,
रचनानियमांनुसार प्रत्येक व्यंजक मूलतः एकाच रीतीने तयार होऊ शकते.
त्यामुळे तयार होणारी व्यंजके \emph{एकमेव रीतीने वाचनीय} असतात. परिणामी,
त्यांची ही विगमनाधारित व्याख्या आपल्याला त्याच पद्धतीने---विगमनाधारित
व्याख्येने---त्या व्यंजकांना अर्थ देता येतो, याची खात्री देते.

व्यंजकांना अर्थ देणे हे चिन्हार्थमीमांसेचे क्षेत्र आहे. चिन्हार्थमीमांसेतील
मध्यवर्ती संकल्पना म्हणजे संरचना मधील पूर्ति. संरचना
भाषेच्या मूलभूत घटकांना अर्थ देते: प्रांत म्हणजे वस्तूंचा अरिक्त
संच. संख्यापक या प्रांतातील घटकांवर फिरतात असा त्यांचा अर्थ लावला जातो,
स्थिरांक ना प्रांतातील घटक नेमून दिले जातात, फलने ना
प्रांत पासून त्याच प्रांताकडे जाणारी फलने नेमून दिली जातात आणि
विधेये ना प्रांत वरील संबंध नेमून दिले जातात. प्रांत
आणि मूलभूत शब्दसंग्रहातील चिन्हांना केलेल्या नेमणुका मिळून
संरचना बनते.
चले ही सूत्रे मध्ये येऊ शकतात; त्यांना चिन्हार्थ देण्यासाठी
आपल्याला प्रांत मधील घटक ही नेमून द्यावी लागतात---हेच
चर-मूल्यांकन होय. शेवटी, पूर्ति-संबंध या सर्व गोष्टी एकत्र आणतो.
सूत्राची संरचना~$\Struct{M}$ मध्ये चर-मूल्यांकन~$s$ च्या सापेक्ष पूर्ति होऊ शकते; हे $\Sat{M}{\varphi}[s]$ असे
लिहितात. हा संबंधदेखील~$\varphi$ च्या संरचनेवरील विगमनाने परिभाषित केला जातो.
उदाहरणार्थ, $\varphi \land \psi$ ची पूर्ति परिभाषित करताना तार्किक संयोजकांचे
सत्यता-कोष्टक वापरून ती $\varphi$ आणि~$\psi$ यांची पूर्ति होते (किंवा होत नाही)
यांच्या आधारे सांगितली जाते. त्यानंतर असे निष्पन्न होते की सूत्र~$\varphi$
हे वाक्य, म्हणजे त्यात मुक्त चरे नसतील, तर चराचे
मूल्यांकन असंबद्ध ठरते; त्यामुळे संरचनांमध्ये वाक्यांची
सरळ पूर्ति होते (किंवा होत नाही) असे आपण म्हणू शकतो.

वाक्यांसाठीच्या पूर्ति-संबंध~$\Sat{M}{\varphi}$ च्या आधारे आपण वैधता,
तार्किक निष्पन्नता आणि पूर्ततायोग्यता या मूलभूत चिन्हार्थविषयक संकल्पना
परिभाषित करू शकतो. एखाद्या वाक्याची प्रत्येक संरचनेत पूर्ति
होत असेल, तर ते वैध, $\Entails \varphi$, असते. वाक्यांच्या एखाद्या
संचातून ते तार्किकरीत्या निष्पन्न होते, $\Gamma \Entails \varphi$, जर~$\Gamma$ मधील
सर्व वाक्यांची पूर्ति करणारे प्रत्येक संरचना हे~$\varphi$ चीही
पूर्ति करत असेल. आणि वाक्यांच्या एखाद्या संचातील सर्व
वाक्यांची एकाच वेळी पूर्ति करणारे एखादे संरचना असेल, तर तो
संच पूर्ततायोग्य असतो. सूत्रे विगमनाधारित रीतीने परिभाषित केली आहेत
आणि पूर्ति हीही सूत्रांच्या संरचनेवरील विगमनाने परिभाषित केली जाते;
म्हणून आपल्या चिन्हार्थमीमांसेचे गुणधर्म सिद्ध करण्यासाठी आणि परिभाषित
चिन्हार्थविषयक संकल्पनांचा परस्परसंबंध लावण्यासाठी विगमन वापरता येते.












\section{$C$ मधील फलने~$\Th{Q}$ मध्ये निरूपणीय असतात}\label{inc:req:cre:sec}

$C$ मधील प्रत्येक फलन~$\Th{Q}$ मध्ये निरूपणीय आहे, हे दाखवायचे
आहे. अखेरीस $C$ मधील प्रत्येक $k$-आदानी फलन
$f(x_0,\dots,x_{k-1})$ ला त्याचे निरूपण करणारे सूत्र
$\varphi_f(x_0,\dots,x_{k-1},y)$ कसे द्यायचे, हे दाखवावे लागेल.

\begin{lem}
$n$ आणि $m$ या नैसर्गिक संख्या असू द्या. $n \neq m$ असेल,
तर $Q \Proves \num n \neq \num m$.
\end{lem}

\begin{proof}
$n$ वर विगमन वापरून प्रत्येक $m$ साठी $n \neq m$ असेल, तर
$Q \Proves \eq/[\num n][\num m]$, हे दाखवू.

आधार टप्प्यात $n = 0$. जर $m$ हा $0$ च्या बरोबर नसेल,
तर एखाद्या नैसर्गिक संख्या $k$ साठी $m = k +
1$. आपल्याकडे $\lforall[x][\eq/[0][x']]$ असे सांगणारे स्वयंसिद्धक
आहे. संख्यापकाच्या स्वयंसिद्धकाने $x$ च्या जागी
$\num k$ ठेवून $\eq/[0][\num k']$ हा निष्कर्ष मिळतो. पण
$\num k'$ म्हणजेच $\num m$.

विगमन टप्प्यात विधान $n$ साठी खरे आहे असे मानून $n+1$
पाहू. $m$ ही कोणतीही नैसर्गिक संख्या असू द्या. दोन शक्यता आहेत:
$m = 0$, किंवा एखाद्या $k$ साठी $m = k+1$. पहिली बाब
वरप्रमाणे हाताळता येते. दुसऱ्या बाबतीत $n+1 \neq
k+1$ असे समजा. मग $n \neq k$. $n$ साठीच्या विगमन गृहीतकाने
$\Th{Q} \Proves \eq/[\num n][\num k]$. आपल्याकडे
$\lforall[x][\lforall[y][\eq[x'][y'] \lif \eq[x][y]]]$
असे सांगणारे स्वयंसिद्धक आहे. संख्यापकाच्या स्वयंसिद्धकाने
$\eq[\num n'][\num k'] \lif \eq[\num n][\num
  k]$ मिळते. विधानीय तर्कशास्त्र वापरून $\Th{Q}$ मध्ये
$\eq/[\num n][\num k] \lif \eq/[\num n'][\num k']$
हे निष्पन्न होते. विध्यात्मक अनुमानाने
$\eq/[\num n'][\num k']$ मिळते; आणि $\num k'$ म्हणजे
$\num m$ असल्यामुळे हेच अपेक्षित होते.
\end{proof}

या पूर्वप्रमेयाचा आशय मर्यादित आहे: वेगवेगळे संख्यांक
वेगवेगळ्या वस्तू दर्शवतात, हे $\Th{Q}$ सिद्ध करू शकते. उदाहरणार्थ,
$\Th{Q}$ हे $0'' \neq 0'''$ सिद्ध करते. पण हे सर्वसाधारणपणे
दाखवण्यासाठी थोडी काळजी घ्यावी लागते. येथे वापरलेले विगमन
$\Th{Q}$ च्या \emph{बाहेर} आहे, हेही लक्षात घ्या.

शून्य, उत्तरवर्ती, बेरीज, गुणाकार, समतेचे जातिबोधक फल
आणि प्रक्षेपण यांचे निरूपण करता येईल. प्रत्येक बाबतीत योग्य
निरूपक सूत्र सहज दिसते. उदाहरणार्थ, शून्याचे निरूपण
\[
y = 0,
\]
उत्तरवर्तीचे निरूपण
\[
x_0' = y,
\]
आणि बेरीजेचे निरूपण
\[
x_0 + x_1 = y.
\]
खरे काम म्हणजे संबंधित विधाने $\Th{Q}$ सिद्ध करू शकते, हे
दाखवणे. उदाहरणार्थ, वरील सूत्र बेरीजेचे निरूपण करते असे
म्हणण्यासाठी नैसर्गिक संख्यांच्या प्रत्येक जोडी $m$ आणि $n$
साठी $\Th{Q}$ पुढील सिद्ध करते, हे दाखवावे लागते:
\[
\eq[\num n + \num m][\num {n+m}]
\]
आणि
\[
\lforall[y][(\eq[(\num n + \num m)][y] \lif \eq[y][\num {n+m}])].
\]

संयोजनाचे काय? $h$ पुढीलप्रमाणे परिभाषित आहे असे माना:
\[
h(x_0,\dots,x_{l-1}) = f(g_0(x_0,\dots,x_{l-1}), \dots,
g_{k-1}(x_0,\dots,x_{l-1})).
\]
येथे $f,g_0,\dots,g_{k-1}$ यांचे अनुक्रमे निरूपण करणारी
सूत्रे $\varphi_f, \varphi_{g_0}, \dots,
\varphi_{g_{k-1}}$ आपल्याला आधीच मिळाली आहेत. आता $h$ चे निरूपण
करणारे सूत्र $\varphi_h$ परिभाषित करता येते. त्यासाठी
$\varphi_h(x_0,\dots,x_{l-1},y)$ पुढील असे घेऊ:
\begin{multline*}
  \lexists[z_0\dots][\lexists[z_{k-1}][(\varphi_{g_0}(x_0,\dots,x_{l-1},z_0) \land
  \dots \land \varphi_{g_{k-1}}(x_0,\dots,x_{l-1},z_{k-1}) \land\\
  \varphi_f(z_0,\dots,z_{k-1},y))]].
\end{multline*}


शेवटी अपरिबद्ध शोध पाहू. $g(x,\vec z)$ नियमित आहे आणि
$\Th{Q}$ मध्ये, समजा $\varphi_g(x,\vec
z,y)$ या सूत्राने, निरूपणीय आहे असे माना. $f$ ची व्याख्या
$f(\vec z) = \mu x \; g(x,\vec z)$ अशी असू द्या. $f$ चे
निरूपण करणारे सूत्र $\varphi_f(\vec z,y)$ शोधायचे आहे. नैसर्गिक
निवड अशी:
\[
\varphi_f(\vec z,y) \ident \varphi_g(y,\vec z,0) \land \lforall[w][(w < y
\lif \lnot \varphi_g(w,\vec z,0))].
\]

आणखी काही पूर्वप्रमेये वापरून हे सूत्र योग्य आहे असे दाखवता येते; उदाहरणार्थ:

\begin{lem}
  प्रत्येक चल $x$ आणि प्रत्येक नैसर्गिक संख्या $n$ साठी $\Th{Q}$
  हे $\eq[(x'
  + \num n)][(x + \num n)']$ सिद्ध करते.
\end{lem}

पुन्हा लक्षात घ्या: यावरून $\Th{Q}$ हे
$\lforall[x][\lforall[y][(x' + y) = (x +y)']]$
सिद्ध करते असा सिद्धतायोग्यतेचा दावा चुकीचा आहे.

\begin{proof}
नेहमीप्रमाणे $n$ वर विगमन वापरू. आधार टप्प्यात $n =
0$ असेल, तर $\Th{Q}$ हे $\eq[(x'+0)][(x+0)']$
सिद्ध करते, हे दाखवायचे आहे. पुढील विधाने मिळतात:
\begin{eqnarray*}
& & \eq[(x' + 0)][x'] \quad \text{स्वयंसिद्धक 4 वरून} \\
& & \eq[(x + 0)][x] \quad \text{स्वयंसिद्धक 4 वरून} \\
& & \eq[(x+0)'][x'] \quad \text{ओळ 2 वरून} \\
& & \eq[(x' + 0)][(x + 0)'] \quad \text{ओळी 1 आणि 3 वरून}
\end{eqnarray*}
विगमन टप्प्यात $\Th{Q}$ मध्ये $\eq[(x' +
  \num n)][(x + \num n)']$ आधी सिद्ध झाले आहे असे
माना. $\num{n+1}$ म्हणजे $\num
n'$ असल्यामुळे $\Th{Q}$ हे
$\eq[(x' + \num n')][(x + \num
n')']$ सिद्ध करते, हे दाखवायचे आहे. पुढील समताशृंखला
$\Th{Q}$ मध्ये सिद्ध करता येते:
\begin{eqnarray*}
(x' + \num n') & \eq & (x' + \num n)' \quad \text{स्वयंसिद्धक 5} \\
& \eq & (x + \num n')' \quad \text{विगमन गृहीतकावरून}
\end{eqnarray*}
\end{proof}

\begin{lem}
\begin{enumerate}
\item $\Th{Q}$ हे $\lnot (x < \num 0)$ सिद्ध करते.
\item प्रत्येक नैसर्गिक संख्या $n$ साठी $\Th{Q}$ पुढील सिद्ध करते:
\[
x < \num {n+1} \lif (\eq[x][\num 0] \lor \dots \lor \eq[x][\num n]).
\]
\end{enumerate}
\end{lem}

\begin{proof}
१ आणि २ चा काही भाग अनौपचारिकपणे पाहू; म्हणजे औपचारिक
निष्पत्ती कसे रचावे यासाठी फक्त संकेत देऊ.

भाग १ साठी $<$ च्या व्याख्येनुसार $\Th{Q}$ मध्ये
$\lnot \lexists[y][\eq[(y'+x)][0]]$ सिद्ध करायचे आहे. प्रथम-क्रम
तर्कशास्त्राच्या स्वयंसिद्धकांनी आणि नियमांनी हे
$\lforall[y][\eq/[(y' +
    x)][0]]$ याच्याशी समतुल्य आहे. कल्पना अशी: $\eq[(y'+x)][0]$
असे माना. $x$ हा 0 असेल, तर $\eq[(y' + 0)][0]$.
पण $\Th{Q}$ च्या स्वयंसिद्धक 4 वरून
$\eq[(y' + 0)][y']$, आणि स्वयंसिद्धक 2 वरून
$\eq/[y'][0]$; हा व्याघात. म्हणून
$\lforall[y][\eq/[(y' +x)][0]]$. $x$ हा $0$ नसेल,
तर स्वयंसिद्धक 3 वरून असा $z$ असतो की
$\eq[x][z']$. मग $\eq[(y' + z')][0]$.
स्वयंसिद्धक~5 वरून $\eq[(y' + z)'][0]$, हेही
स्वयंसिद्धक~2 शी विसंगत आहे.

भाग २ साठी $n$ वर विगमन वापरा. प्रथम $n =0$
हा आधार टप्पा पाहू. त्या वेळी
$x < \num 1 \lif \eq[x][\num
  0]$ दाखवायचे आहे. $x < \num 1$ असे माना.
$<$ च्या परिभाषक स्वयंसिद्धकाने
$\lexists[y][\eq[(y'+x)][0']]$ मिळते. $y$ कडे हा
गुणधर्म आहे असे माना; म्हणजे $\eq[y'+x][0']$.

$\eq[x][0]$ दाखवायचे आहे. स्वयंसिद्धक 3 वरून
$x$ हा 0 नसेल, तर एखाद्या $z$ साठी तो $z'$ च्या
बरोबर असतो. मग $\eq[(y' + z')][0']$.
$\Th{Q}$ च्या स्वयंसिद्धक~5 वरून
$\eq[(y'+z)'][0']$. स्वयंसिद्धक 1 वरून
$\eq[(y' +
    z)][0]$. पण व्याख्येने याचा अर्थ
$z < 0$; हे भाग~१ शी विसंगत आहे.
\end{proof}

\begin{explain}
संगणनक्षम फलनांचा संच म्हणजेच $\Th{Q}$ मध्ये निरूपणीय
फलनांचा संच, असे लक्षणन आपण दाखवले आहे. प्रत्यक्षात हा
पुरावा अधिक व्यापक आहे. निरूपणीयतेच्या व्याख्येवरून
$\Th{Q}$ चा विस्तार असलेल्या (किंवा ज्यात $\Th{Q}$ चे
अर्थनिर्धारण करता येते अशा) कोणत्याही उपपत्तीत
संगणनक्षम फलनांचे निरूपण करता येते, हे सहज दिसते.
उलट दिशेसाठी पूर्वीच्या निरूपणीयतेच्या विभागातील अटी
आवश्यक आहेत: निष्पत्ती व्यवस्थेतील उपपत्ती सुसंगत
असावी, भिन्न संख्यांकांची असमता सिद्ध व्हावी आणि
सिद्धतांची तपासणी संगणनक्षम असावी. या अटींखाली
प्रत्येक निरूपणीय फलन संगणनक्षम असते. उदाहरणार्थ, संगणनक्षम फलनांचा संच
म्हणजे पेआनो अंकगणितात किंवा अगदी झर्मेलो--फ्रेंकेल संच
उपपत्तीत निरूपणीय फलनांचा संच, असे त्या उपपत्ती सुसंगत
आहेत या गृहीतकाखाली लक्षणन करता येते.
गोडेल यांनी नोंदवल्याप्रमाणे हे काहीसे आश्चर्यकारक आहे.
सिद्धतायोग्यतेचे प्रश्न कोणती उपपत्ती विचारात घेतो यावर
फार अवलंबून असतात, हे आपण पाहू: साधारणपणे स्वयंसिद्धके
जितकी बलवान, तितके अधिक सिद्ध करता येते. या अटी पूर्ण करणाऱ्या आणि संगणनक्षम फलनांचे निरूपण
करणाऱ्या स्वयंसिद्धकीय उपपत्त्यांत निरूपणीय फलने
नेमकी संगणनक्षम फलनेच असतात.

\end{explain}



\begin{quote}\small\textbf{स्रोतनोंद.} SOL6-A289: संयोजनातील सर्व मध्यवर्ती मूल्ये आणि अंतिम फलसूत्र यांना अस्तित्व-संख्यापकांची एकच व्याप्ती दिली आहे. निरूपणीय फलनांच्या उलट दाव्यास सुसंगतता, वेगळ्या संख्यांकांची सिद्ध असमता आणि संगणनक्षम सिद्धता-तपासणीच्या अटी आवश्यक आहेत. मूळ एक-दिशीय आधार टप्प्यात उलट दिशा गहाळ असल्याची पूर्वीची नोंद चुकीची होती; ती मागे घेतली आहे.\end{quote}
\begin{quote}\small\textbf{स्रोतनोंद.} OLINC-328: मूळ संयोजन-सूत्रातील संख्यापक आणि संयुक्त-विधानाचे कंस SOL6-A289 मध्ये दुरुस्त केले आहेत. मध्यवर्ती z चलांच्या अस्तित्वाची व्याप्ती सर्व g सूत्रे आणि अंतिम f सूत्र व्यापते.\end{quote}
\begin{quote}\small\textbf{स्रोतनोंद.} OLINC-327: SOL6-A289 मध्ये पूर्वीच्या OLINC-035 प्रमाणे सुसंगतता, भिन्न संख्यांकांची सिद्ध असमता आणि सिद्धतांची संगणनक्षम तपासणी या अटी उलट दाव्यात स्पष्ट घातल्या आहेत.\end{quote}









\section{फलनांचा संच~$C$}\label{inc:req:cee:sec}

$C$ हा पुढील फलने समाविष्ट करणारा सर्वात लहान फलनसंच असू द्या:
\begin{enumerate}
\item $0$,
\item उत्तरवर्ती फलन,
\item बेरीज,
\item गुणाकार,
\item प्रक्षेपण फलने, आणि
\item समतेचे जातिबोधक फल, $\Char{=}$;
\end{enumerate}
आणि जो पुढील क्रियांखाली बंद आहे:
\begin{enumerate}
\item संयोजन, आणि
\item नियमित फलनांवर लावलेला अपरिबद्ध शोध.
\end{enumerate}
ही शेवटची मर्यादा म्हणजे $\mu$ क्रिया फक्त तिचा परिणाम
सर्वत्र परिभाषित असेल तेव्हाच वापरता येते, एवढेच. याची
\emph{सामान्य पुनरावर्ती} फलनांच्या व्याख्येशी तुलना करा:
येथे बेरीज, गुणाकार आणि $\Char{=}$ समाविष्ट केले आहेत,
पण आदिम पुनरावर्तन वगळले आहे.

$C$ मधील प्रत्येक फलन पुनरावर्ती आहे, हे स्पष्ट आहे; कारण
बेरीज, गुणाकार आणि $\Char{=}$ ही फलने पुनरावर्ती आहेत.
याचा उलट दावा खरा आहे, हेही दाखवू. म्हणजे $C$ मधील
इतर क्रिया वापरून आदिम पुनरावर्तन करता येते.












\section{विधानसंच}\label{int:sem:pro:sec}

कधीकधी संबंधात्मक प्रतिमान~$\mModel{M}$ मधील ज्या
जगतांत सूत्र~$\varphi$ सत्य असते, त्या जगांच्या संचाला नाव देणे
उपयुक्त ठरते. त्याला $\varphi$ ने परिभाषित केलेला
\emph{विधानसंच} म्हणू आणि $\Prop{M}{\varphi}$ असे लिहू.

\begin{defn}
  $V$ या फलनाचा विस्तार $\Prop{M}{\varphi} \colon \Frm
  \to \Pow{W}$ अशा फलनात पुढीलप्रमाणे विगमनाने करू:
  \begin{enumerate}
  \item $\Prop{M}{\lfalse} = \emptyset$
  \item $\Prop{M}{p} = V(p)$
  \item $\Prop{M}{\lnot \psi} = {}$
    \[
      \Setabs{w \in W}{\text{$Rww'$ असेल अशा प्रत्येक $w'$ साठी $w'
          \notin \Prop{M}{\psi}$}}.
    \]
  \item $\Prop{M}{\psi \land \chi} = \Prop{M}{\psi} \cap \Prop{M}{\chi}$.
  \item $\Prop{M}{\psi \lor \chi} = \Prop{M}{\psi} \cup \Prop{M}{\chi}$.
  \item $\Prop{M}{\psi \lif \chi} = {}$
    \[
      \Setabs{w \in W}{\text{$Rww'$ असेल अशा प्रत्येक $w'$ साठी,
          $w' \in \Prop{M}{\psi}$ असेल तर $w' \in \Prop{M}{\chi}$}}.
    \]
  \end{enumerate}
\end{defn}


\begin{prop}
  \label{int:sem:pro:prop:propositions-true}
  $\mSat{M}{\varphi}[w]$ तेव्हा आणि केवळ तेव्हाच $w \in \Prop{M}{\varphi}$.
\end{prop}


\begin{proof}
  सराव.
\end{proof}

\begin{prob}
  \ref{int:sem:pro:prop:propositions-true} सिद्ध करा.
\end{prob}



\begin{quote}\small\textbf{स्रोतनोंद.} OLINC-331: या चिन्हात A हे बदलते सूत्र आहे; त्याला त्या सूत्राने परिभाषित केलेल्या जगांचा संच जोडणारे फलन असा आशय आहे. ही चिन्हवाचनाची स्पष्टता आहे, सिद्धतेतील उणीव किंवा निश्चित नवा स्रोतदोष नव्हे.\end{quote}
\begin{quote}\small\textbf{स्रोतनोंद.} OLINC-330: स्रोताने पूर्ति-संबंधाच्या जागी फक्त प्रतिमानाचे नाव छापणारा mModel मॅक्रो वापरला आहे; येथे mSat केला आहे.\end{quote}










\section{याद्या}\label{lam:ldf:lst:sec}

यादी $[M_1, M_2, \ldots, M_n]$ पुढीलप्रमाणे परिभाषित करता येते.
येथे $s$ आणि $z$ ही परस्पर भिन्न आणि कोणत्याही $M_i$ मध्ये
मुक्त नसलेली चरे ही पदे गुंफण्यापूर्वी निवडली आहेत:
\[
  [M_1, M_2, \ldots, M_n] \ident \lambd[sz][s M_1 (s M_2 (\ldots (s M_n z)))]
\]
अनौपचारिकपणे, यादीचे निरूपण करणारे पद हे त्या यादीचे
“संचायक” फलन असते. हे फलन दोन पदे स्वीकारते: पहिले
असे फलन असते, जे आतापर्यंत साचलेले मूल्य आणि नवा घटक
घेऊन नवे संचयित मूल्य देते; दुसरे म्हणजे संचयाचे आरंभीचे
मूल्य. ते एकेका घटकाचा संचय करते आणि शेवटी परिणाम देते.

\begin{digress}
  फलनीय प्रोग्रामलेखनाशी परिचय असेल, तर ही व्याख्या
  \emph{उजवी घडी} (right fold) यासारखी वाटेल:
  यादीतील घटकांवर उजवी घडी घालणारे फलन म्हणूनच
  यादी परिभाषित होते.
\end{digress}

या संकेतावर आधारलेली काही उपयुक्त फलने सहज परिभाषित
करता येतात:
\begin{align*}
  \fn{Sum} & \ident
  \lambd[l][l \, (\lambd[xa][\fn{Add} \, x \, a])\,  \num{0}]\\
  \fn{Len} & \ident
  \lambd[l][l \, (\lambd[xa][\fn{Add} \, a \, \num{1}]) \num{0}]
\end{align*}

$\fn{Sum}$ हे चर्च अंकांच्या यादीतील बेरजेची गणना
करते. ते यादीवर संचय करते: आरंभीचे मूल्य
$\num{0}$ असते आणि प्रत्येक घटकासाठी
$\fn{Add} \, x\, a$ हे नवे मूल्य देते
(येथे $x$ हा चालू घटक आणि $a$ हे संचयित मूल्य आहे).
परिणाम म्हणजे सर्व घटकांची बेरीज.

\begin{prob}
  $\fn{Len}$ काय करते? स्पष्ट करा.
\end{prob}

\begin{prob}
  यादीचा क्रम उलटवणारे फलन परिभाषित करा.
\end{prob}


\begin{quote}\small\textbf{स्रोतनोंद.} SOL6-A292: यादीतील पदांच्या मुक्त चरांना बाहेरच्या lambda ने पकडू नये म्हणून संचायक आणि आरंभीच्या मूल्याची चरे यादी गुंफण्यापूर्वी निवडली आहेत. ही अट सर्वसाधारण पदांसाठी आवश्यक आहे; सर्व बंद पदांसाठी ती आपोआप पूर्ण होते.\end{quote}









\section{परिवर्तन आणि न्यूनीकरण}\label{lam:int:con:sec}

लॅम्डा कलनात परिवर्तनाचे किंवा न्यूनीकरणाचे तीन प्रकार आहेत.
“रूपांतरण” एकदिश आहे की द्विदिश, यावर परिवर्तन आणि
न्यूनीकरण यांतील फरक अवलंबून असतो; हे आपण पुढे पाहू.

$\alpha$-परिवर्तनात चरांची नावे बदलली जातात.
$\lambda x. x$ आणि $\lambda y.y$ ही दोन्ही पदे एकच फलन
(ओळख फलन) दर्शवतात असे मानणे स्वाभाविक आहे.
ही कल्पना आपण पुढीलप्रमाणे औपचारिक करू:




\begin{quote}\small\textbf{स्रोतनोंद.} OLINC-332: गोठवलेला स्रोत येथेच थांबतो; आश्वासित औपचारिक व्याख्या किंवा उर्वरित प्रकार या विभागात दिलेले नाहीत.\end{quote}










\section{सिद्धता-उपपत्तीय संकल्पना}\label{mvl:seq:prf:sec}


\begin{explain}
जशा आपण अनेक महत्त्वाच्या चिन्हार्थविषयक संकल्पना (वैधता, तार्किक निष्पन्नता,
पूर्ततायोग्यता) परिभाषित केल्या आहेत, तशाच त्यांना अनुरूप
\emph{सिद्धता-उपपत्तीय संकल्पना} आता परिभाषित करू. संरचनांमध्ये
वाक्यांची पूर्ति होते की नाही, याचा आधार घेऊन या संकल्पना परिभाषित
केल्या जात नाहीत; त्याऐवजी ठरावीक क्रमवर्तींची निष्पन्नता किंवा
अनिष्पन्नता यांचा आधार घेतला जातो. या दोन प्रकारच्या संकल्पना
एकमेकांशी जुळतात, हा एक महत्त्वाचा शोध होता. त्या जुळतात, हाच
\emph{निर्दोषता प्रमेय} आणि \emph{संपूर्णता प्रमेय} यांचा आशय आहे.
\end{explain}

\begin{defn}[प्रमेये]
वाक्य~$\varphi$ हे \emph{प्रमेय} असते, जर $\Log{LK}$ मध्ये
$\quad \Sequent \varphi$ या क्रमवर्तीची निष्पत्ती असेल. $\varphi$ हे प्रमेय
असेल तर आपण $\Proves \varphi$ लिहितो आणि ते प्रमेय नसेल तर $\Proves/ \varphi$ लिहितो.
\end{defn}

\begin{defn}[निष्पन्नता]
वाक्य~$\varphi$ हे वाक्यांच्या संच~$\Gamma$ \emph{पासून
निष्पन्न करता येण्याजोगे} असते, म्हणजे $\Gamma \Proves \varphi$, तेव्हा आणि केवळ तेव्हाच,
जेव्हा एखादा सांत उपसंच~$\Gamma_0 \subseteq \Gamma$ आणि $\Gamma_0$ मधील
वाक्यांची एखादी क्रमिका $\Gamma_0'$ अशी असते की $\Log{LK}$
$\Gamma_0' \Sequent \varphi$ हे निष्पन्न करते. $\varphi$ हे $\Gamma$ पासून
निष्पन्न करता येण्याजोगे नसेल, तर आपण $\Gamma \Proves/ \varphi$ लिहितो.
\end{defn}

आकुंचन, दुर्बलीकरण आणि अदलाबदल नियमांमुळे $\Gamma_0'$ मधील
वाक्यांचा क्रम आणि त्यांची पुनरावृत्ती कितीदा होते हे महत्त्वाचे नसते:
$\Gamma_0' \Sequent \varphi$ ही क्रमवर्ती निष्पन्न करता येण्याजोगे असेल, तर
$\Gamma_0'$ मध्ये असलेली तीच वाक्ये असणाऱ्या कोणत्याही
$\Gamma_0''$ साठी $\Gamma_0'' \Sequent \varphi$ ही क्रमवर्तीसुद्धा निष्पन्न करता
येते. उदाहरणार्थ, $\Gamma_0 = \{\psi, \chi\}$ असेल, तर
$\Gamma_0' = \tuple{\psi, \psi, \chi}$ आणि $\Gamma_0'' =
\tuple{\chi, \chi, \psi}$ या दोन्ही क्रमिकांमध्ये $\Gamma_0$ मधीलच
वाक्ये आहेत. यांपैकी एक क्रमिका असणारी क्रमवर्ती निष्पन्न करता येण्याजोगे असेल,
तर दुसरी असणारी क्रमवर्तीसुद्धा निष्पन्न करता येते; उदा.:
\begin{prooftree}
  \AxiomC{}
  \Deduce$\psi, \psi, \chi \fCenter \varphi$
  \RightLabel{\LeftR{\Contraction}}
  \UnaryInf$\psi, \chi \fCenter \varphi$
  \RightLabel{\LeftR{\Exchange}}
  \UnaryInf$\chi, \psi \fCenter \varphi$
  \RightLabel{\LeftR{\Weakening}}
  \UnaryInf$\chi, \chi, \psi \fCenter \varphi$
\end{prooftree}
यापुढे $\Gamma_0$ हा वाक्यांचा सांत संच असेल, तर
$\Gamma_0 \Sequent \varphi$ याचा अर्थ असा कोणताही क्रमवर्ती घेऊ की ज्याच्या
पूर्वांगात $\Gamma_0$ मधील वाक्यांची क्रमिका आहे; तसेच आवश्यक असल्यास
आकुंचन, अदलाबदल आणि दुर्बलीकरणे गृहीत धरू.

\begin{defn}[सुसंगतता]
वाक्यांचा संच~$\Gamma$ हा \emph{विसंगत} असतो तेव्हा आणि केवळ तेव्हाच, जेव्हा
एखादा सांत उपसंच~$\Gamma_0 \subseteq \Gamma$ असा असतो की $\Log{LK}$
$\Gamma_0 \Sequent \quad$ हे निष्पन्न करते. $\Gamma$ विसंगत नसेल,
म्हणजे प्रत्येक सांत $\Gamma_0 \subseteq \Gamma$ साठी $\Log{LK}$
$\Gamma_0 \Sequent \quad$ हे निष्पन्न करत नसेल, तर त्याला
\emph{सुसंगत} म्हणतो.
\end{defn}

\begin{prop}[परावर्तिता]
\label{mvl:seq:prf:prop:reflexivity}
$\varphi \in \Gamma$ असेल, तर $\Gamma \Proves \varphi$.
\end{prop}

\begin{proof}
$\varphi \Sequent \varphi$ ही आरंभीची क्रमवर्ती निष्पन्न करता येण्याजोगे आहे आणि $\{\varphi\}
\subseteq \Gamma$.
\end{proof}

\begin{prop}[एकस्वनिकता]
\label{mvl:seq:prf:prop:monotonicity}
$\Gamma \subseteq \Delta$ आणि $\Gamma \Proves \varphi$ असेल, तर $\Delta
\Proves \varphi$.
\end{prop}

\begin{proof}
$\Gamma \Proves \varphi$ असे समजा; म्हणजे, एखादा सांत $\Gamma_0
\subseteq \Gamma$ असा आहे की $\Gamma_0 \Sequent \varphi$ ही क्रमवर्ती
निष्पन्न करता येण्याजोगे आहे. $\Gamma \subseteq \Delta$ असल्यामुळे $\Gamma_0$ हा
$\Delta$ चासुद्धा सांत उपसंच आहे. म्हणून $\Gamma_0 \Sequent \varphi$ ची
निष्पत्ती ही $\Delta \Proves \varphi$ हेसुद्धा दाखवते.
\end{proof}

\begin{prop}[संक्रमकता]
\label{mvl:seq:prf:prop:transitivity}
$\Gamma \Proves \varphi$ आणि $\{\varphi\} \cup \Delta \Proves
\psi$ असेल, तर $\Gamma \cup \Delta \Proves \psi$.
\end{prop}

\begin{proof}
$\Gamma \Proves \varphi$ असेल, तर एखादा सांत $\Gamma_0 \subseteq \Gamma$
आणि $\Gamma_0 \Sequent \varphi$ ची निष्पत्ती~$\pi_0$ असते.
$\{\varphi\} \cup \Delta \Proves \psi$ असेल, तर एखाद्या सांत उपसंचासाठी
$\Delta_0 \subseteq \Delta$, $\varphi, \Delta_0 \Sequent \psi$ ची
निष्पत्ती~$\pi_1$ असते. पुढील निष्पत्ती विचारात घ्या:
\begin{prooftree}
  \AxiomC{}
  \RightLabel{$\pi_0$}
  \Deduce$\Gamma_0 \fCenter \varphi$
  \AxiomC{}
  \RightLabel{$\pi_1$}
  \Deduce$\varphi, \Delta_0 \fCenter \psi$
  \RightLabel{\Cut}
  \BinaryInf$\Gamma_0, \Delta_0 \fCenter \psi$
\end{prooftree}
$\Gamma_0 \cup \Delta_0 \subseteq \Gamma \cup \Delta$ असल्यामुळे यावरून
$\Gamma \cup \Delta \Proves \psi$ हे दिसते.
\end{proof}

विशेषतः, याचा अर्थ असा की $\Gamma \Proves \varphi$ आणि $\varphi
\Proves \psi$ असेल, तर $\Gamma \Proves \psi$. तसेच $\varphi_1,
\dots, \varphi_n \Proves \psi$ असेल आणि प्रत्येक~$i$ साठी $\Gamma \Proves \varphi_i$
असेल, तर $\Gamma \Proves \psi$, हेसुद्धा यावरून मिळते.

\begin{prop}
\label{mvl:seq:prf:prop:incons}
$\Gamma$ विसंगत असतो तेव्हा आणि केवळ तेव्हाच, जेव्हा प्रत्येक
  वाक्य~$\varphi$ साठी $\Gamma \Proves {\varphi}$.
\end{prop}

\begin{proof}
सराव.
\end{proof}

\begin{prob}
\ref{fol:seq:ptn:prop:incons} सिद्ध करा.
\end{prob}

\begin{prop}[संहतता]
  \label{mvl:seq:prf:prop:proves-compact}
  \begin{enumerate}
  \item $\Gamma \Proves \varphi$ असेल, तर एखादा सांत उपसंच $\Gamma_0
    \subseteq \Gamma$ असा आहे की $\Gamma_0 \Proves \varphi$.
  \item $\Gamma$ चा प्रत्येक सांत उपसंच सुसंगत असेल, तर $\Gamma$ सुसंगत
    असतो.
  \end{enumerate}
\end{prop}

\begin{proof}
  \begin{enumerate}
    \item $\Gamma \Proves \varphi$ असेल, तर एखादा सांत उपसंच
      $\Gamma_0 \subseteq \Gamma$ असा आहे की $\Gamma_0
      \Sequent \varphi$ या क्रमवर्तीची निष्पत्ती आहे. परिणामी,
      $\Gamma_0 \Proves \varphi$.
    \item $\Gamma$ विसंगत असेल, तर एखादा सांत
      उपसंच~$\Gamma_0 \subseteq \Gamma$ असा आहे की $\Log{LK}$
      $\Gamma_0 \Sequent \quad$ हे निष्पन्न करते. पण मग $\Gamma_0$ हा
      $\Gamma$ चा विसंगत सांत उपसंच आहे.
  \end{enumerate}
\end{proof}



\begin{quote}\small\textbf{स्रोतनोंद.} OLINC-333: हा विभाग अभिजात LK आणि द्विस्थानी क्रमवर्ती वापरतो; आधीच्या बहुमूल्य विभागातील n-स्थानी कलनाची येथे व्याख्या केलेली नाही.\end{quote}










\section{अंकगणिताची अमानक प्रतिमाने}\label{mod:bas:nsa:sec}

\begin{defn}
$\Lang{L_N}$ ही अंकगणिताची भाषा असू द्या. तिच्यात
स्थिरांक $\Obj{0}$, द्विस्थानी विधेय $<$,
एकस्थानी फलन $\prime$, आणि द्विस्थानी फलने
$+$ व $\times$ आहेत.
\begin{enumerate}
\item अंकगणिताचे \emph{मानक प्रतिमान} म्हणजे $\Lang{L_N}$ साठीची
  संरचना $\Struct{N}$: तिचा प्रांत $\Nat = \{ 0, 1, 2, \dots\}$
  असतो; $\Obj{0}$ चा अर्थ $0$, $<$ चा अर्थ $\Nat$ वरील लहान-असण्याचा
  संबंध, तर $\prime$, $+$ आणि $\times$ यांचे अर्थ अनुक्रमे
  $\Nat$ वरील उत्तरवर्ती, बेरीज आणि गुणाकार असे असतात.
\item \emph{सत्य अंकगणित} म्हणजे $\Theory{N}$ ही उपपत्ती.
\end{enumerate}
\end{defn}

$\Lang{L_N}$ मध्ये काम करताना $\Obj{0}$ वर उत्तरवर्ती फलन
$n$ वेळा लावून मिळणाऱ्या $\Obj{0}^{\prime\dots\prime}$
या रूपातील पदाचे संक्षिप्त रूप~$\num{n}$ लिहितो.

\begin{defn}
$\Lang{L_N}$ साठीची संरचना~$\Struct{M}$ ही
$\Struct{N} \iso \Struct{M}$ असेल तेव्हा आणि केवळ तेव्हाच
\emph{मानक} असते.
\end{defn}

\begin{thm}\label{mod:bas:nsa:thm:non-std}
सत्य अंकगणिताची अमानक गणनीय प्रतिमाने आहेत.
\end{thm}

\begin{proof}
$\Lang{L_N}$ मध्ये नवा स्थिरांक~$c$ घालून तिचा विस्तार करा
आणि पुढील उपपत्ती विचारात घ्या:
\[
\Theory{N} \cup \Setabs{\num{n} < c}{n \in \Nat}.
\]
ही उपपत्ती सांत पूर्ततायोग्य आहे. त्यामुळे संहततेनुसार तिला
$\Struct{M}$ हे प्रतिमान आहे; अधोगामी लोव्हेनहाइम--स्कोलेम
प्रमेयानुसार ते गणनीय घेता येते. $\Domain{M}$ हा
$\Struct{M}$ चा प्रांत असताना, $\Lang{L_N}$ मधील अतार्किक
स्थिरांकांचे $\Struct{M}$ मधील अर्थनिर्धारण असे लिहू:
$\mathbf{z} = \Assign{\Obj{0}}{M} \in \Domain{M}$,
${\prec} = \Assign{<}{M} \subseteq \Domain{M}^2$,
$* = \Assign{\prime}{M} \colon \Domain{M} \to \Domain{M}$,
आणि $\oplus = \Assign{+}{M}, \otimes =
\Assign{\times}{M} \colon \Domain{M}^2 \to \Domain{M}$.
प्रत्येक $x \in \Domain{M}$ साठी, $x$ वर~$*$ लावल्याने
मिळणाऱ्या $\Domain{M}$ मधील घटकाला $x^*$ लिहू.


आता, $h$ हे $\Struct{N}$ आणि $\Struct{M}$ यांचे
समरूपण असेल, तर काही $n \in \Nat$ साठी
$h(n) = \Assign{c}{M}$ असते. म्हणून $\Struct{N}$ मध्ये
$s(x) = n$ अशी कोणतीही चलनियुक्ती $s$ घ्या. मग
$\Sat{N}{\eq[\num{n}][x]}[s]$; तसेच
\ref{mod:bas:iso:thm:isom} च्या सिद्धतेनुसार
$\Sat{M}{\eq[\num{n}][x]}[h \circ s]$.
त्यामुळे $\Assign{c}{M} = \mathbf{z}^{*\cdots *}$
(येथे $*$ ही क्रिया $n$ वेळा केली आहे).
पण गृहीतकानुसार $\Sat{M}{\num{n} < c}$ आहे आणि
$\prec$ हा संबंध अपरावर्ती आहे; म्हणून हे अशक्य आहे.
त्यामुळे $\Struct{M}$ अमानक आहे.
\end{proof}

\begin{prob}
संच $X$ वरील संबंध $R$ हा तेव्हा आणि केवळ तेव्हाच
\emph{सुस्थापित} असतो, जेव्हा $R$ मध्ये अनंत अवरोही साखळी
नसते; म्हणजे $\dots x_2Rx_1Rx_0$ अशी अट पूर्ण करणारे
$X$ मधील $x_0$, $x_1$, $x_2$,~\dots नसतात.
झेर्मेलो--फ्रेंकेल संच उपपत्ती~$ZF$ सुसंगत आहे असे मानून,
$ZF$ ची असुस्थापित प्रतिमाने आहेत हे दाखवा; म्हणजे
$\dots x_2 \in x_1 \in x_0$ असणारी प्रतिमाने
$\mathfrak{M}$ आहेत हे सिद्ध करा.
\end{prob}

\ref{mod:bas:nsa:thm:non-std} मधील अमानक प्रतिमान
$\Struct{M}$ हे मानक प्रतिमानाशी प्राथमिक सममूल्य असल्यामुळे
$\Struct{M}$ चे अनेक गुणधर्म काढता येतात. या विभागाचा
उर्वरित भाग अशा गुणधर्मांच्या अभ्यासासाठी आहे; त्यातून
$\Theory{N}$ च्या अमानक गणनीय प्रतिमानांचे
अचूक लक्षण मिळेल.

\begin{enumerate}
\item $\Domain{M}$ मधील कोणताही घटक स्वतःहून $\prec$-लहान
  नसतो: $\lforall[x][\lnot x < x]$ हे वाक्य $\Struct{N}$
  मध्ये सत्य असल्यामुळे $\Struct{M}$ मध्येही सत्य आहे.
\item याच रीतीने $\prec$ हा $\Domain{M}$ वरील
  \emph{रेषीय क्रम} आहे, म्हणजे तो $\Domain{M}$ वरील
  पूर्ण, अपरावर्ती आणि संक्रमक संबंध आहे.
\item $\mathbf{z}$ हा $\Domain{M}$ मधील $\prec$-लघुतम घटक आहे.
\item $\Domain{M}$ मधील प्रत्येक घटक त्याच्या $*$-उत्तरवर्तीपेक्षा
  $\prec$-लहान असतो; आणि $x^*$ हा $x$ पेक्षा मोठ्या
  $\Domain{M}$ मधील घटकांपैकी $\prec$-लघुतम असतो.
\item $\Struct{M}$ मध्ये $\Nat$ शी समरूप असा $\prec$ चा
  आरंभीचा खंड आहे: $\mathbf{z}, \mathbf{z}^*,
  \mathbf{z}^{**}, \dots$. त्याला $\Domain{M}$ चा
  \emph{मानक भाग} म्हणतो. $\Domain{M}$ मधील इतर प्रत्येक
  घटक \emph{अमानक} असतो. $\Domain{M}$ मध्ये अमानक
  घटक असलेच पाहिजेत;
  अन्यथा \ref{mod:bas:nsa:thm:non-std} च्या सिद्धतेतील $h$ हे
  समरूपण ठरेल. $n, m, \dots$ ही चले आपण
  $\Struct{M}$ च्या या मानक भागावर चालणारी मानतो.
\item प्रत्येक अमानक घटक प्रत्येक मानक घटकापेक्षा मोठा
  असतो; कारण प्रत्येक $n \in \Nat$ साठी
  \[
  \Sat{N}{\lforall[z][(\lnot(\eq[z][\Obj{0}] \lor
  \dots \lor \eq[z][\num{n}]) \lif \num{n} < z)]},
  \]
  म्हणून $z \in \Domain{M}$ सर्व मानक घटकांहून
  वेगळा असेल, तर तो त्या सर्वांपेक्षा \emph{मोठा} असतो.
\item $\mathbf{z}$ सोडून $\Domain{M}$ मधील प्रत्येक
  घटक हा एखाद्या एकमेव $\Domain{M}$ मधील घटकाचा
  $*$-उत्तरवर्ती असतो; त्या पूर्ववर्तीला $^*x$ लिहितो.
  $x = y^*$ असेल, तर $x$ आणि $y$ यांपैकी एक मानक
  असल्यास दोन्ही मानक असतात (आणि एक अमानक असल्यास
  दोन्ही अमानक असतात).
\item $\Domain{M}$ वर $\approx$ हा सममूल्यता संबंध
  परिभाषित करा: एखाद्या \emph{मानक} $n$ साठी
  $x \oplus n = y$ किंवा $y \oplus n =x$
  असेल तेव्हा आणि केवळ तेव्हाच $x \approx y$.
  दुसऱ्या शब्दांत, $x$ आणि $y$ यांच्यात सांत अंतर
  असेल तेव्हा आणि केवळ तेव्हाच $x \approx y$.
  $n$ आणि $m$ मानक असतील, तर $n \approx m$.
  $x$ चा \emph{खंड} म्हणजे सममूल्यता वर्ग
  $[x] = \Setabs{y}{x \approx y}$.
\item $x \prec y$ पण $x \not\approx y$ असे समजा.
  $\Sat{N}{\lforall[x][\lforall[y][(x < y \lif (x' < y \lor x' =
      y))]]}$ असल्यामुळे, $x^* \prec y$ किंवा
  $x^* = y$. दुसरी शक्यता $x \approx y$ सूचित
  करत असल्यामुळे अशक्य आहे; म्हणून
  $x^* \prec y$. त्याचप्रमाणे $x \prec y$ आणि
  $x \not\approx y$ असतील, तर $x \prec {^*y}$.
  म्हणून $x \prec y$ आणि $x \not\approx y$ असतील,
  तर प्रत्येक $w \approx x$ हा प्रत्येक $v \approx y$
  पेक्षा $\prec$-लहान असतो. त्यानुसार, प्रत्येक
  अमानक खंड $[x]$ पुढील द्विदिश अनंत साखळी बनवतो:
  \[
  \cdots \prec  {^{**}x} \prec {^*}x \prec x \prec x^* \prec x^{**}
  \prec \cdots
  \]
  पूर्णांकांच्या क्रमप्रकारासारखा असल्यामुळे तिला
  $Z$-साखळी म्हणतात.


\item $\prec$ हा क्रम खंडांवरही नेऊ शकतो: $x \prec y$
  आणि ते भिन्न खंडांत असतील, तर $x$ चा खंड
  $y$ च्या खंडापेक्षा लहान असतो. अमानक घटक असणारा
  खंड \emph{अमानक} असतो. मानक खंड हा लघुतम खंड आहे.

\item लघुतम अमानक खंड नाही: $y$ अमानक असेल,
  तर $x \prec y$ आणि $x \not\approx y$ अशी अट
  पूर्ण करणारा अमानक $x$ असतो. सिद्धता: मानक
  प्रतिमान $\Struct{N}$ मध्ये प्रत्येक संख्येला दोनने
  भाग जातो, कदाचित बाकी एक उरते; म्हणजे
  $\Sat{N}{\lforall[y][\lexists[x][(y = x + x \lor y = x + x +
        \Obj{0}')]]}$. प्राथमिक सममूल्यतेमुळे
  प्रत्येक $y \in \Domain{M}$ साठी असा
  $x \in \Domain{M}$ असतो की $x \oplus x = y$
  किंवा $x \oplus x \oplus \mathbf{z}^*= y$.
  $x$ मानक असेल, तर $y$ सुद्धा मानक असेल;
  म्हणून $x$ अमानक आहे. शिवाय $x$ आणि $y$
  वेगळ्या खंडांत येतात, म्हणजे $x \not\approx y$.
  तसे नसल्यास, एखाद्या मानक $n$ साठी
  $x \oplus n = y$ असेल. $y = x \oplus x$
  असेल, तर $x \oplus n = x \oplus x$; बेरीज
  रद्द करण्याच्या नियमानुसार $x = n$ येईल.
  हा नियम $\Struct{N}$ मध्ये आणि म्हणून
  $\Struct{M}$ मध्येही चालतो; पण मग $x$
  मानक ठरेल. $y = x \oplus x \oplus \mathbf{z}^*$
  या प्रसंगातही तसेच होते.

\item याच प्रकारच्या युक्तिवादाने सर्वांत मोठा खंडही नाही.
\item खंडांचा क्रम सघन आहे: $[x]$ हा $[y]$
  पेक्षा लहान असेल ($x \not\approx y$), तर
  त्यांच्या मध्ये दोन्हींपेक्षा वेगळा खंड $[z]$
  असतो. $x \prec y$ असे समजा. आधीप्रमाणे
  $x \oplus y$ ला दोनने भाग जातो (कदाचित
  बाकी उरते); म्हणून काही $u \in \Domain{M}$
  साठी $x \oplus y = u \oplus u$ किंवा
  $x \oplus y = u \oplus u \oplus \mathbf{z}^*$.
  $u$ हा $x$ आणि $y$ यांचा मध्य आहे आणि
  त्यांच्यामध्ये येतो. $x \oplus y = u \oplus u$
  असे समजा (दुसरा प्रसंगही तसाच): $u \approx x$
  असेल, तर एखाद्या मानक $n$ साठी
  \[
  x \oplus y = x \oplus n \oplus x \oplus n,
  \]
  म्हणून $y = x \oplus n \oplus n$; यातून
  $x \approx y$ येते, जे गृहीतकाविरुद्ध आहे.
  म्हणून $u \not\approx x$. त्याच युक्तिवादाने
  $u \not\approx y$.
\end{enumerate}
त्यामुळे अमानक खंडांचा क्रम परिमेय संख्यांप्रमाणे
आहे: तो अंत्यबिंदू नसलेला गणनीय रेषीय क्रम आहे.
यावरून सत्य अंकगणिताच्या कोणत्याही दोन
गणनीय अमानक प्रतिमानांचे,
$\Struct{M}_1$ आणि $\Struct{M_2}$ यांचे,
फक्त $<$ आणि $=$ असणाऱ्या भाषेतील
न्यूनीकृत प्रतिमान समरूप असतात.
खरोखर, समरूपण $h$ असे बनवता येते:
$\Struct{M_1}$ आणि $\Struct{M_2}$ यांचे
मानक भाग मानक प्रतिमान $\Struct{N}$ शी,
म्हणून एकमेकांशी, समरूप असतात. अमानक भाग
बनवणारे खंड परिमेयांप्रमाणे क्रमबद्ध असतात;
म्हणून \ref{mod:bas:dlo:thm:cantorQ} नुसार ते
समरूप असतात. प्रत्येक खंडातील कोणतेही
घटक जुळवून खंडांचे समरूपण खंडांच्या
\emph{आत} विस्तारता येते; मग
$\Struct{M_1}$ मधील $x$ च्या उत्तरवर्तीची
प्रतिमा ही $\Struct{M_2}$ मधील $x$ च्या
प्रतिमेचा उत्तरवर्ती ठरवता येते.
यावरून मात्र $\mathfrak{M}_1$ आणि
$\mathfrak{M}_2$ ही संपूर्ण अंकगणितीय
भाषेत समरूप आहेत असे \emph{निघत नाही}
(कारण समरूपता नेहमी चिन्हसंचाच्या संदर्भातच
असते). $\Domain{M_1}$ आणि $\Domain{M_2}$
वर बेरीज व गुणाकार परिभाषित करण्याचे
असमरूप मार्ग आहेत. एका संपूर्ण उपपत्तीच्या
गणनीय प्रतिमानांची संख्या नेमकी~2
असू शकत नाही, या वॉट यांच्या प्रसिद्ध
प्रमेयावरूनही हे निघते.

\begin{prob}
अंकगणिताच्या अमानक प्रतिमानात सर्वांत मोठा
खंड असू शकत नाही हे दाखवा.
\end{prob}

\begin{prob}
$\Lang{L}$ ही $<$ हा एकच विधेय
($\eq$ व्यतिरिक्त) असणारी प्रथम-क्रमाची
भाषा असू द्या आणि
$\Struct{N} = (\Nat, <)$ असू द्या.
$\Struct{N}$ चे सर्व सांत किंवा सहसांत
उपसंच परिभाषणीय आहेत. \emph{हेच एकमेव}
$\Struct{N}$ चे परिभाषणीय उपसंच आहेत हे दाखवा.

(सूचना: प्रथम, $\Obj{prc}(x,y)$ हे
$\Lang{L}$ मधील “$x$ हा $y$ चा लगतचा
पूर्ववर्ती आहे” याचे संक्षिप्त रूप असलेले
सूत्र असू द्या:
\[
x<y \land \lnot \lexists[z][(x<z \land z < y)].
\]
आता $\Struct{N}$ च्या कोणत्याही परिभाषणीय
उपसंचाला $\Lang{L}$ मधील सूत्र $\varphi(x)$
अनुरूप असते. अशा प्रत्येक $\varphi$ साठी
पुढील वाक्य $\theta$ विचारात घ्या:
\[
\lexists[x][\lforall[y][\lforall[z][((x<y \land x<z \land \Obj{prc}(y,z)
  \land \varphi(y)) \lif \varphi(z))]]].
\]
$\Sat{N}{\theta}$ असेल तेव्हा आणि केवळ
तेव्हाच $\varphi$ ने परिभाषित केलेला
$\Struct{N}$ चा उपसंच सांत किंवा
सहसांत असतो हे दाखवा.

आता $\Struct{M}$ हे $\Struct{N}$
शी प्राथमिक सममूल्य अमानक प्रतिमान असू द्या.
$a \in \Domain{M}$ अमानक असेल, तर
$b, c \in \Domain{M}$ हे $a$ पेक्षा
मोठे घ्या आणि $b$ हा~$c$ चा लगतचा
पूर्ववर्ती असू द्या. मग $\Domain{M}$
चे असे स्वसमरूपण $h$ आहे की $h(b)=c$
(का?). त्यामुळे $b$ ने $\varphi$ ची पूर्ति
केली, तर $c$ नेही ती केली पाहिजे (का?).
म्हणून $\theta$ हे $\Struct{M}$ मध्ये सत्य
आहे आणि त्यामुळे $\Struct{N}$ मध्येही.
पण याचा अर्थ असा की $\varphi$ ने परिभाषित
केलेला $\Struct{N}$ चा उपसंच सांत
किंवा सहसांत आहे.
\end{prob}



\begin{quote}\small\textbf{स्रोतनोंद.} SOL6-A293: अमानक प्रतिमानाच्या भाषेसाठी L\_N आणि संबंधाच्या प्रांतासाठी Domain(M) वापरले आहेत. वेगळ्या खंडांतील घटकांच्या क्रमात x चा उत्तरवर्तीही y पेक्षा लहान असतो, हा मिळालेला निष्कर्ष आता स्पष्ट आहे.\end{quote}
\begin{quote}\small\textbf{स्रोतनोंद.} OLINC-334: SOL6-A293 मध्ये आधी निश्चित केलेली अंकगणिताची भाषा L\_N आणि प्रतिमानाचा प्रांत Domain(M) स्पष्ट वापरला आहे; ऐतिहासिक विसंगत संकेत स्रोतामध्ये जतन आहेत.\end{quote}
\begin{quote}\small\textbf{स्रोतनोंद.} OLINC-336: गोठवलेल्या स्रोताचा 'प्रत्येक खंड' हा दावा मानक खंडासाठी चुकीचा आहे; द्विदिश अनंत साखळी फक्त अमानक खंडांना लागू होते.\end{quote}
\begin{quote}\small\textbf{स्रोतनोंद.} OLINC-337: SOL6-A293 मध्ये या प्रसंगाचा प्रत्यक्ष निष्कर्ष x चा उत्तरवर्ती y पेक्षा लहान आहे असा केला आहे; आधीचेच गृहीतक पुन्हा सांगणारे मूळ वाक्य दुरुस्त केले आहे.\end{quote}
\begin{quote}\small\textbf{स्रोतनोंद.} OLINC-338: भिन्न खंडांसाठी x आणि y असण्याची अट स्रोताच्या पहिल्या वाक्यात गाळलेली आहे; एकाच खंडातील x<y त्यांचा खंड स्वतःहून लहान करत नाही.\end{quote}
\begin{quote}\small\textbf{स्रोतनोंद.} OLINC-335: स्रोताच्या अर्धे करण्याच्या सूत्रात प्रत्येक y साठी 'प्रत्येक x' असे चुकीने आहे; अस्तित्वदर्शक x आवश्यक आहे आणि येथे दुरुस्त केला आहे.\end{quote}










\section{सामान्य मोडल तर्कशास्त्रे}\label{nml:syn:nor:sec}

मोडल तर्कशास्त्र म्हणजे योग्य अटी पूर्ण करणारा
सूत्रांचा संच असे मानता येते. काही मोडल
तर्कशास्त्रे विशेष महत्त्वाची आहेत: त्यांची अर्थलक्षणे
आणि उपयोग रंजक आहेत, आणि विशेषतः त्यांचा
सैद्धान्तिक अभ्यास चांगला झालेला आहे. यांनाच
सामान्य मोडल तर्कशास्त्रे म्हणतात.

\begin{defn}\label{nml:syn:nor:defn:modal-logic}
\emph{मोडल तर्कशास्त्र}~$\Log L$ म्हणजे मूलभूत
मोडल भाषेतील सूत्रांचा असा संच की
\begin{enumerate}
\item त्यात सर्व विधानात्मक सर्वत:सत्य सूत्रे असतात;
\item तो एकरूप आदेशनाखाली बंद असतो; आणि
\item तो मोडस पोनेन्सखाली बंद असतो; म्हणजे $\varphi$ आणि
  $(\varphi \lif \psi) \in \Log L$ असतील, तर
  $\psi \in \Log L$ सुद्धा असते.
\end{enumerate}
\end{defn}

\begin{ex}
फारसे रंजक नसले तरी ही व्याख्या पूर्ण करणारे,
म्हणून मोडल तर्कशास्त्रे असणारे, सूत्रे
चे काही संच असे आहेत:
\begin{enumerate}
\item सर्व सर्वत:सत्य सूत्रांचा संच;
\item सर्व सूत्रांचा संच (विसंगत मोडल तर्कशास्त्र).
\end{enumerate}
\end{ex}

\begin{defn}\label{defn:normal-modal-logic}
मोडल तर्कशास्त्र~$\Log L$ हे तेव्हा आणि केवळ
तेव्हाच \emph{सामान्य} असते, जेव्हा
\begin{enumerate}
\item त्यात पुढील दोन्ही असतात:
\begin{align}
\tag{K} \Box(p \lif q) & \lif (\Box p \lif \Box q) \\
\tag{Dual} \Diamond p & \liff \lnot \Box \lnot p
\end{align}
\item आणि ते अनिवार्यीकरणाखाली बंद असते; म्हणजे
  $\varphi \in \Log L$ असेल, तर
  $\Box \varphi \in \Log L$ सुद्धा असते.
\end{enumerate}
\end{defn}

मोडल तर्कशास्त्रे परिभाषित करण्याचे दोन प्रमुख
मार्ग आपण पाहू. पहिला सिद्धता-उपपत्तीय: विशिष्ट
निष्पत्ती-पद्धतींमध्ये निष्पन्न करता येणाऱ्या
सूत्रांचा संच. दुसरा प्रतिमान-उपपत्तीय:
प्रतिमानांच्या विशिष्ट वर्गांत वैध असणाऱ्या
सूत्रांचे संच. साध्या विधानात्मक तर्कशास्त्रात
(किंवा प्रथम-क्रम तर्कशास्त्रात) हाच प्रकार दिसतो.
तेथेही सूत्रांचा संच दोन प्रकारे ठरवता येतो:
सर्व सत्यतामूल्य-नियुक्त्यांत सत्य असलेल्या
सर्वत:सत्य सूत्रांचा संच म्हणून, किंवा एखाद्या
निष्पत्ती-पद्धतीतील सर्व निष्पन्न करता येण्याजोगे
सूत्रांचा संच म्हणून. सामान्य मोडल
तर्कशास्त्रांतही संबंधीय प्रतिमाने आणि स्वयंसिद्धकी
(हिल्बर्ट-प्रकारच्या) सिद्धता पद्धती वापरून हे
दुहेरी वर्णन शक्य होते.












\section{क्रमवर्ती-रूप नैसर्गिक निगमन}\label{int:pty:snd:sec}

$\Gamma$ ही मुक्त न केलेली गृहीतके आणि~$\varphi$ हा
निष्कर्ष असलेली नैसर्गिक निगमनाची निष्पत्ती असेल,
तर $\Gamma \Sequent \varphi$ असे लिहू;~$\Gamma$ रिकामे
असेल, तर $\Sequent \varphi$ असे लिहू.

$\Gamma \cup \{\varphi_1, \dots, \varphi_n\}$ साठी
$\Gamma, \varphi_1, \dots, \varphi_n$ आणि $\Gamma \cup \Delta$
साठी $\Gamma, \Delta$ असे लिहू.

$\Gamma \Sequent \varphi \land \psi$ असेल, म्हणजे~$\Gamma$
ही मुक्त न केलेली गृहीतके आणि $\varphi \land \psi$ हे
अंतिम सूत्र असलेली निष्पत्ती असेल,
तर तळाशी $\Elim{\land}$ लावून तीच
मुक्त न केलेली गृहीतके आणि~$\varphi$ हा निष्कर्ष
असलेली निष्पत्ती मिळते. दुसऱ्या शब्दांत,
$\Gamma \Sequent \varphi \land \psi$ असेल, तर
$\Gamma \Sequent \varphi$.
\begin{prooftree}
  \Axiom$\Gamma \fCenter \varphi \land \psi$
  \RightLabel{$\Elim{\land}$}
  \UnaryInf$\Gamma \fCenter \varphi$
  \DisplayProof\qquad\bottomAlignProof
  \Axiom$\Gamma \fCenter \varphi \land \psi$
  \RightLabel{$\Elim{\land}$}
  \UnaryInf$\Gamma \fCenter \psi$
\end{prooftree}
$\Elim{\land}$ ही खूण नैसर्गिक निगमनातील त्याच
नावाच्या नियमाशी असलेला संबंध सूचित करते.

त्याचप्रमाणे $\Gamma, \varphi \Sequent \psi$ असेल,
म्हणजे $\Gamma, \varphi$ ही मुक्त न केलेली गृहीतके
आणि~$\psi$ हे अंतिम सूत्र असलेली
निष्पत्ती असेल, तर $\Intro{\lif}$ नियम लावून
$\Gamma$ ही मुक्त न केलेली गृहीतके आणि
$\varphi \lif \psi$ हे अंतिम सूत्र असलेली
निष्पत्ती मिळते; म्हणजे
$\Gamma \Sequent \varphi \lif \psi$.
यात गृहीतकांचे मुक्त करणे अधिक स्पष्ट झाले आहे.
\begin{prooftree}
  \Axiom $\Gamma, \varphi \fCenter \psi$
  \RightLabel{$\Intro{\lif}$}
  \UnaryInf $\Gamma \fCenter \varphi \lif \psi$
\end{prooftree}

इतर नियमांवरूनही अशाच प्रकारे निष्कर्ष काढता
येतात. ते पुढीलप्रमाणे स्पष्ट केले आहेत:
\begin{gather*}
  \Axiom $\Gamma \fCenter \varphi$
  \Axiom $\Delta \fCenter \psi$
  \RightLabel{$\Intro{\land}$}
  \BinaryInf $\Gamma,\Delta \fCenter \varphi \land \psi$
  \DisplayProof\\
  \Axiom $\Gamma \fCenter \varphi \land \psi$
  \RightLabel{$\Elim{\land}_1$}
  \UnaryInf $\Gamma \fCenter \varphi$
  \DisplayProof
  \qquad
  \Axiom $\Gamma \fCenter \varphi \land \psi$
  \RightLabel{$\Elim{\land}_2$}
  \UnaryInf $\Gamma \fCenter \psi$
  \DisplayProof
  \\
  \Axiom $\Gamma \fCenter \varphi$
  \RightLabel{$\Intro{\lor}_1$}
  \UnaryInf $\Gamma \fCenter \varphi \lor \psi$
  \DisplayProof\qquad
  \Axiom $\Gamma \fCenter \psi$
  \RightLabel{$\Intro{\lor}_2$}
  \UnaryInf $\Gamma \fCenter \varphi \lor \psi$
  \DisplayProof\\
  \Axiom $\Gamma \fCenter \varphi \lor \psi$
  \Axiom $\Delta, \varphi \fCenter \chi$
  \Axiom $\Delta', \psi \fCenter \chi$
  \RightLabel{$\Elim{\lor}$}
  \TrinaryInf $\Gamma, \Delta, \Delta' \fCenter \chi$
  \DisplayProof
  \\
  \Axiom $\Gamma, \varphi \fCenter \psi$
  \RightLabel{$\Intro{\lif}$}
  \UnaryInf $\Gamma \fCenter \varphi \lif \psi$
  \DisplayProof
  \qquad
  \Axiom $\Delta \fCenter \varphi \lif \psi$
  \Axiom $\Gamma \fCenter \varphi$
  \RightLabel{$\Elim{\lif}$}
  \BinaryInf $\Gamma, \Delta \fCenter \psi$
  \DisplayProof
  \\
  \Axiom $\Gamma \fCenter \lfalse$
  \RightLabel{$\FalseInt$}
  \UnaryInf $\Gamma \fCenter \varphi$
  \DisplayProof
\end{gather*}

कोणतेही गृहीतक स्वतःच~$\varphi$ वरून~$\varphi$ ची
निष्पत्ती असते; म्हणजे $\varphi \Sequent \varphi$
हे नेहमीच असते.

\begin{prooftree}
  \AxiomC{$ $}
  \UnaryInfC{$\varphi \fCenter \varphi$}
\end{prooftree}

हे नियम एकत्रितपणे कोणत्या नैसर्गिक निगमनाच्या
निष्पत्त्या अस्तित्वात आहेत यांचे कलन म्हणून
घेता येतात. त्यांना नैसर्गिक निगमनाचे चिन्हांकनात्मक
रूपांतरही मानता येते; प्रत्येक टप्प्यावर कोणते
सूत्र निष्पन्न झाले आहे एवढेच नव्हे,
तर ते कोणत्या मुक्त न केलेली गृहीतकांवरून
निष्पन्न झाले आहे हेही त्यात नोंदते.

\begin{prooftree}
  \Axiom$\varphi \fCenter \varphi$
  \UnaryInf $\varphi \fCenter \varphi \lor (\varphi \lif \lfalse)$
  \Axiom $\psi \fCenter \psi$
  \BinaryInf$\varphi, \psi \fCenter \lfalse$
  \UnaryInf$\psi \fCenter \varphi \lif \lfalse$
  \UnaryInf$\psi \fCenter \varphi \lor (\varphi \lif \lfalse)$
  \Axiom$\psi \fCenter \psi$
  \BinaryInf$\psi \fCenter \lfalse$
  \UnaryInf$\fCenter \psi \lif \lfalse$
\end{prooftree}
येथे~$\psi$ हे $(\varphi \lor (\varphi \lif \lfalse))\lif \lfalse$
याचे संक्षिप्त रूप आहे.



\begin{quote}\small\textbf{स्रोतनोंद.} SOL6-A325: अंतिम उदाहरणातील संदर्भात चुकून आलेले अपूर्ण अभिव्यंजन-चिन्ह व कंस काढले; दिलेला B चा अर्थ आणि सिद्धतेचा निष्कर्ष जपला.\end{quote}











\section{विधानात्मक तर्कशास्त्राची संपूर्णता}\label{pl:prp:com:sec}

\begin{defn}
\label{pl:prp:com:def:MaxCon}
सूत्रांचा संच~$\Gamma$ \emph{महत्तम सुसंगत} असतो, जर तो
सुसंगत असेल आणि $\Gamma \subseteq \Delta$ असा कोणताही सुसंगत
संच~$\Delta$ घेतल्यास $\Gamma = \Delta$ असेल.
\end{defn}

\begin{lem}[सत्यता पूर्वप्रमेय]
  \label{pl:prp:com:lem:truth}
$\Gamma$ महत्तम सुसंगत असू द्या; मग:
\begin{enumerate}
\item \label{pl:prp:com:lem:truth-1} $\varphi \in \Gamma$ तेव्हा आणि केवळ तेव्हा $\lnot \varphi
  \notin \Gamma$;
\item \label{pl:prp:com:lem:truth-2} $\Gamma \Proves \varphi$ तेव्हा आणि केवळ तेव्हा $\varphi \in \Gamma$;
\item \label{pl:prp:com:lem:truth-3} $\varphi \lif \psi \in \Gamma$ तेव्हा आणि केवळ तेव्हा
  $\varphi \notin \Gamma$ किंवा $\psi \in \Gamma$.
\end{enumerate}
\end{lem}

\begin{proof}
\begin{enumerate}
\item $\varphi \in \Gamma$ असे माना. $\lnot \varphi \in \Gamma$ सुद्धा असेल,
  तर $\Gamma$ विसंगत ठरतो. $\varphi$ आणि $\lnot\varphi$ यांपैकी एकही
  $\Gamma$ मध्ये नसेल, तर \readerexternalref{pl:prp:axd:prop:phi} नुसार
  $\Gamma\cup\{\varphi\}$ आणि $\Gamma \cup \{\lnot \varphi\}$ यांपैकी
  एक संच सुसंगत असतो; म्हणून $\Gamma$ महत्तम नसतो. उलटी दिशाही
  याचप्रमाणे सिद्ध होते.
\item $\varphi \in \Gamma$ असेल, तर अर्थातच $\Gamma \Proves \varphi$.
  उलट, $\Gamma \Proves \varphi$ असे माना. $\varphi \notin \Gamma$ असेल,
  तर \ref{pl:prp:com:lem:truth-1} नुसार $\lnot \varphi \in \Gamma$. त्यामुळे
  $\Gamma \Proves \varphi$ आणि $\Gamma \Proves \lnot \varphi$, हा व्याघात.
\item सरावासाठी.
\end{enumerate}
\end{proof}

\ref{pl:prp:com:lem:truth}\ref{pl:prp:com:lem:truth-2} ची डावीकडून उजवीकडील
दिशा दाखवते की $\Gamma$ \emph{निगमनतः बंद} आहे: म्हणजे
कोणत्याही~$\varphi$ साठी $\Gamma \Proves \varphi$ असेल, तर $\varphi
\in \Gamma$.

\begin{prob}
  \ref{pl:prp:com:lem:truth} ची सिद्धता पूर्ण करा.
\end{prob}

\begin{thm}[संपूर्णता]
\label{pl:prp:com:thm:completeness}
$\Gamma$ सुसंगत असेल, तर $\Gamma$ पूर्ततायोग्य असतो.
\end{thm}

\begin{proof}
$\varphi_0$, $\varphi_1$,~\dots ही भाषेतील सर्व सूत्रांची
पूर्ण यादी असू द्या. पुढीलप्रमाणे सूत्रांच्या संचांची
$\Gamma_0$, $\Gamma_1$,~\dots ही वाढती क्रमिका
आवर्ती रीतीने परिभाषित करू:
\begin{align*}
\Gamma_0 & = \Gamma\\
\Gamma_{n+1} &  =
\begin{cases}
  \Gamma_n \cup \{\varphi_n\} & \text{जर $\Gamma_n \cup \{\varphi_n\}$ सुसंगत असेल;}\\
  \Gamma_n \cup \{\lnot \varphi_n\} & \text{अन्यथा}.
\end{cases}
\end{align*}
मग पुढील व्याख्या करू:
\[
\Gamma^* = \bigcup_{0\le n}\Gamma_n.
\]
आता क्रमाने पुढील तथ्ये सिद्ध करायची आहेत:
\begin{enumerate}
\item प्रत्येक $n$ साठी $\Gamma_n$ हा संच सुसंगत आहे
  ($n$ वर विगमनाने, \readerexternalref{pl:prp:axd:prop:phi} वापरून);
\item $\Gamma^*$ सुसंगत आहे;
\item $\Gamma^*$ महत्तम आहे.
\end{enumerate}
मग $p_i \in \Gamma^*$ असेल तेव्हाच $v(p_i) = \True$
असे ठरवून सत्य-मूल्यांकन~$v$ परिभाषित करू. $\varphi$ वर विगमन
केल्यास अधिक गुंतागुंतीच्या वाक्ये साठीही
$\Gamma^*$ मधील सदस्यत्व आणि $v$ नुसार सत्यता जुळतात,
हे दिसते: $\overline{v}(\varphi) = \True$ तेव्हा आणि केवळ तेव्हा
$\varphi \in \Gamma^*$. विशेषतः, $v$ हे $\Gamma^*$ ची पूर्तता
करते. $\Gamma \subseteq \Gamma^*$ असल्याने $\Gamma$
सुद्धा पूर्ततायोग्य आहे, हेच सिद्ध करायचे होते.
\end{proof}

\begin{cor}
$\Gamma \Entails \varphi$ असेल, तर $\Gamma \Proves \varphi$.
\end{cor}

\begin{proof}
$\Gamma\Proves/ \varphi$ असेल, तर \readerexternalref{pl:prp:axd:prop:prov-incons} नुसार
$\Gamma \cup \{\lnot \varphi\}$ सुसंगत आहे. म्हणून संपूर्णता
प्रमेयानुसार $\Gamma \cup \{\lnot \varphi\}$ पूर्ततायोग्य आहे;
त्यामुळे $\Gamma \Entails/ \varphi$.
\end{proof}

\begin{prop}[संहतता प्रमेय]
\label{pl:prp:com:prop:compactness}
$\Gamma$ पूर्ततायोग्य असेल तेव्हा आणि केवळ तेव्हा $\Gamma$ चा प्रत्येक
\emph{सान्त} उपसंच~$\Gamma_0$ पूर्ततायोग्य असेल.
\end{prop}

\begin{proof}
$\Gamma$ अपूर्ततायोग्य असेल तेव्हा आणि केवळ तेव्हा तो विसंगत असेल;
तेव्हा आणि केवळ तेव्हा $\Gamma$ चा एखादा सान्त उपसंच~$\Gamma_0$
विसंगत असेल; आणि तेव्हा आणि केवळ तेव्हा $\Gamma$ चा एखादा सान्त
उपसंच~$\Gamma_0$ अपूर्ततायोग्य असेल.
\end{proof}

\begin{cor}
$\Gamma \Entails \varphi$ असेल तेव्हा आणि केवळ तेव्हा $\Gamma$ च्या
एखाद्या सान्त उपसंच~$\Gamma_0$ साठी $\Gamma_0 \Entails \varphi$ असेल.
\end{cor}



\begin{quote}\small\textbf{स्रोतनोंद.} SOL6-A338: द्विदिशा दर्शवणारा वाक्प्रयोग आणि विगमनाची संज्ञा सुसंगत केली; संक्षिप्त सिद्धता आणि सराव जपले.\end{quote}











\section{विधानात्मक तर्कशास्त्राची निर्दोषता}\label{pl:prp:snd:sec}

\begin{thm}[निर्दोषता]
\label{pl:prp:snd:thm:soundness}
$\Gamma\Proves\varphi$ असेल, तर
$\Gamma \Entails \varphi$.
\end{thm}

\begin{proof}
प्रमेयांवर विगमन करू. $\varphi$ स्वयंसिद्धक असेल, तर $\Entails
\varphi$; म्हणून विशेषतः $\Gamma \Entails\varphi$. तसेच $\varphi \in \Gamma$
असेल, तर $\Gamma \Entails\varphi$. विगमनात्मक पायरीसाठी $\chi$
आणि $\chi\lif\varphi$ यांपासून \emph{विध्यात्मक अनुमानाने} $\varphi$
मिळाले असे माना. तसेच विगमनाच्या गृहीतकानुसार प्रमेय
$\chi\lif \varphi$ आणि $\chi$ यांसाठी खरे आहे असे माना.
\ref{pl:syn:sem:prop:semanticalfacts} मधील तिसऱ्या मुद्द्यावरून
अपेक्षित निष्कर्ष मिळतो.
\end{proof}

\begin{cor}
$\Gamma$ पूर्ततायोग्य असेल, तर $\Gamma$ सुसंगत असतो. म्हणून
विधानात्मक तर्कशास्त्र सुसंगत आहे.
\end{cor}



\begin{quote}\small\textbf{स्रोतनोंद.} SOL6-A339: निर्दोषतेच्या सिद्धतेतील विगमनाची संज्ञा सुसंगत केली; सर्व गणिती आधारप्रसंग व पायरी जपली.\end{quote}










\section{द्वितीय-क्रम तर्कशास्त्राची भाषा}\label{sol:syn:lan:sec}

प्रथम-क्रम तर्कशास्त्राप्रमाणेच द्वितीय-क्रम तर्कशास्त्रातील
अभिव्यक्ती मूलभूत शब्दसंग्रहापासून बनतात. त्यात \emph{चले},
\emph{स्थिरांक}, \emph{विधेये} आणि कधी कधी
\emph{फलने} असतात. यांच्याबरोबर तार्किक संकारके,
संख्यक आणि कंस व स्वल्पविरामांसारखी विरामचिन्हे वापरून
\emph{पदे} आणि \emph{सूत्रे} तयार होतात. फरक असा की
वस्तूंसाठीच्या चलांव्यतिरिक्त द्वितीय-क्रम तर्कशास्त्रात संबंध
आणि फलने यांसाठीही चले असतात आणि त्यांच्यावर संख्यापन
करता येते. म्हणून द्वितीय-क्रम तर्कशास्त्राची तार्किक चिन्हे
म्हणजे प्रथम-क्रम तर्कशास्त्राची चिन्हे, तसेच पुढील चिन्हे:

\begin{enumerate}
\item प्रत्येक स्थानसंख्या~$n$ साठी द्वितीय-क्रम संबंध चलांचा
  गणनीय अनंत संच: $\Obj V_0^n$, $\Obj V_1^n$, $\Obj V_2^n$, \dots
\item द्वितीय-क्रम फलन चलांचा गणनीय अनंत संच:
  $\Obj u_0^n$, $\Obj u_1^n$, $\Obj u_2^n$, \dots
\end{enumerate}

प्रथम-क्रम तर्कशास्त्रात एकरूपता~$\eq$ सहसा समाविष्ट
असते. प्रथम-क्रम तर्कशास्त्रात $\Lang{L}$ या भाषेतील
अतार्किक चिन्हे आपल्याला काही अर्थपूर्ण गोष्ट व्यक्त करता
यावी यासाठी अत्यावश्यक असतात. अतार्किक चिन्हे नसलेली
वाक्ये नक्कीच आहेत, पण केवळ~$\eq$ वापरून
फारसे अर्थपूर्ण काही सांगता येत नाही. द्वितीय-क्रम
तर्कशास्त्रात मात्र संबंधचल आणि फलनचल अमर्याद प्रमाणात
उपलब्ध असल्याने, अतार्किक चिन्हांचा वेगळा साठा नसतानाही
प्रथम-क्रम भाषेत सांगता येणारी कोणतीही गोष्ट सांगता येते.












\section{समरूपता}\label{sfr:fun:iso:sec}

\begin{explain}
\emph{समरूपता} म्हणजे संबंधित संचांची रचना जपणारे एकास-एक
आच्छादन. ही रचना संचांतील घटक मधील संबंधांवर अवलंबून
असते. पुढील दोन संच पाहा: $X=\{1,2,3\}$ आणि $Y=\{4,5,6\}$.
दोन्ही संचांवर उत्तरवर्ती, लहान असणे आणि मोठे असणे हे संबंध
आहेत. या दोन संचांमधील समरूपता म्हणजे या रचना जपणारे
एकास-एक आच्छादन. म्हणून $f \colon X
\to Y$ हे एकास-एक व आच्छादक फलन समरूपता असेल, जर
$i<j$ तेव्हा आणि केवळ तेव्हा $f(i)<f(j)$, $i>j$ तेव्हा आणि केवळ तेव्हा $f(i)>f(j)$, आणि $j$ हा $i$ चा उत्तरवर्ती असेल
तेव्हा आणि केवळ तेव्हा $f(j)$ हा $f(i)$ चा उत्तरवर्ती असेल.
\end{explain}

\begin{defn}[समरूपता]
$U$ ही जोडी $\langle X, R\rangle$ आणि $V$ ही जोडी $\langle
Y, S\rangle$ असो; येथे $X$ आणि $Y$ हे संच, तर $R$ आणि $S$ हे
अनुक्रमे $X$ आणि $Y$ वरील संबंध आहेत. $X$ पासून $Y$ कडे जाणारे
एकास-एक आच्छादन~$f$ हे $U$ पासून $V$ कडे जाणारी \emph{समरूपता}
असेल तेव्हा आणि केवळ तेव्हा ते संबंधात्मक रचना जपते; म्हणजे
$X$ मधील कोणत्याही $x_{1}$ आणि $x_{2}$ साठी
$\tuple{x_1,x_2} \in R$ तेव्हा आणि केवळ तेव्हा
$\tuple{f(x_1),f(x_2)} \in S$.
\end{defn}

\begin{ex}
$X=\{1,2,3\}$ आणि $Y=\{4,5,6\}$ हे दोन संच, तसेच
लहान असणे आणि मोठे असणे हे संबंध पाहा. $f(x) = 7-x$ असे
असलेले $f\colon X \to
Y$ हे फलन $\tuple{X,<}$ आणि $\tuple{Y,>}$ यांमधील
समरूपता आहे.
\end{ex}



\begin{quote}\small\textbf{स्रोतनोंद.} SOL6-A340: समरूपतेच्या व्याख्येतील द्विदिशा स्पष्ट केली; सर्व संबंध आणि उदाहरणे जपली.\end{quote}











\section{परिचय}\label{sfr:ind:int:sec}

\begin{explain}
विगमन ही तर्कशास्त्रातील गणितात वारंवार वापरली जाणारी
सिद्धतेची पद्धत आहे. नैसर्गिक संख्यांवर केलेले विगमन
तुम्हाला गणितातून परिचित असेल. प्रत्यक्षात अनेक प्रकारचे
गणिती व तार्किक वस्तुसंच विगमनाला अनुकूल असतात.

सर्वसाधारणपणे, विगमनाने सार्वत्रिक दावा सिद्ध करता येतो;
म्हणजे एखाद्या प्रकारच्या प्रत्येक वस्तूला एक गुणधर्म आहे हे
दाखवता येते. विशेषतः, “सार्वत्रिक परिचय” या सिद्धता-पद्धतीने
सिद्धता करणे फार गुंतागुंतीचे असेल तेव्हा विगमन उपयुक्त ठरते.
विशिष्ट रीतीने बनवलेल्या गणिती वस्तूंवरच विगमन चालते:
संचातील प्रत्येक घटक मूलभूत असेल किंवा मूलभूत घटकांपासून
बनलेला असेल, तर त्या संचावर विगमन करता येते. अधिक नेमकेपणाने:
\end{explain}

\begin{defn}
एखादा संच~$S$ हा फलन~$f$ च्या खाली \emph{बंद} असतो तेव्हा
आणि केवळ तेव्हा $x \in
S$ असताना $f(x) \in S$ असते.
\end{defn}

\begin{explain}
उदाहरणार्थ, नैसर्गिक संख्यांचा संच ($\mathbb{N}$) उत्तरवर्ती
फलन $f(x) = x+1$ च्या खाली बंद आहे: $x$ नैसर्गिक संख्या
असेल, तर $x+1$ सुद्धा नैसर्गिक संख्या असते. मात्र नैसर्गिक
संख्यांचा संच पूर्ववर्ती फलन $f(x) = x-1$ च्या खाली बंद
नाही; कारण $0$ नैसर्गिक संख्या आहे, पण $-1$ नाही.
\end{explain}

\begin{defn}
फलन~$f(x)$ च्या प्रांतातील कोणत्याही वस्तू~$a$ साठी $P(a)$
असल्यास $P(f(a))$ असेल (म्हणजे $a$ ला गुणधर्म~$P$ असेल,
तर $f(a)$ लाही असेल), तर गुणधर्म~$P$ फलन~$f$ च्या खाली
\emph{जपला जातो} असे म्हणतात.
\end{defn}

\begin{explain}
“सम संख्या असणे” हा गुणधर्म $f(x) = x+2$ या फलनाखाली
जपला जातो; कारण $x$ सम असेल, तर $x+2$ सुद्धा सम असतो.
मात्र हा गुणधर्म उत्तरवर्ती फलनाखाली जपला जात नाही;
कारण $x$ सम असेल, तर $x+1$ विषम असतो.

नैसर्गिक संख्यांवर विगमन चालते, कारण प्रत्येक नैसर्गिक
संख्येपर्यंत 0 पासून उत्तरवर्ती फलन सान्त वेळा लावून
पोहोचता येते. ज्या संचावर विगमन करता येते, त्या प्रत्येक
संचाचे हे वैशिष्ट्य असते.
\end{explain}

\begin{defn}[$\Nat$ वरील विगमन]
0 ला गुणधर्म~$P$ असेल आणि $P$ हा गुणधर्म उत्तरवर्ती
फलनाखाली जपला जात असेल, तर प्रत्येक नैसर्गिक संख्येला
गुणधर्म~$P$ असतो.
\end{defn}

\begin{ex} 0 पासून $n$ पर्यंतच्या नैसर्गिक
संख्यांची बेरीज $\frac{n(n+1)}{2}$ असते. \\

\textbf{विगमनाचा आधार.} 0 पासून 0 पर्यंतच्या नैसर्गिक
संख्यांची बेरीज $0=
\frac{0\cdot 1}{2}$ असते.\\

\textbf{विगमनाची पायरी.} गुणधर्म $k$ साठी खरा आहे असे
माना. तो $k+1$ साठीही खरा आहे हे दाखवू. म्हणजे $P(k)$
खरे आहे असे माना; अर्थात,
\[ 0 + 1 + \ldots + k = \frac{k(k+1)}{2} \]
$P(k+1)$ खरे आहे, म्हणजे
\[ 0 + 1 + \ldots + k + (k+1) = \frac{(k+1)(k+2)}{2} \]
हे $P(k)$ या गृहीतकाचा उपयोग करून दाखवू.

\begin{align*}
0 + 1 + \ldots + k &= \frac{k(k+1)}{2} \tag{गृहीतक} \\
0 + 1 + \ldots + k + (k+1) &= \frac{k(k+1)}{2} + (k+1)
\tag{दोन्ही बाजूंना $k+1$ मिळवा} \\
&= \frac{k(k+1) + 2(k+1)}{2} \\
&= \frac{k^2 + k + 2k +2}{2} \\
&=\frac{(k+1)(k+2)}{2}
\end{align*}
अशा रीतीने विगमनाची पायरी सिद्ध होते आणि दावा सिद्ध झाला.
\end{ex}

\begin{explain}
सूत्रे बाबतीत मूलभूत घटक अण्वीय सूत्रे असतात.
कमी गुंतागुंतीची सूत्रे संकारकांनी जोडून अधिक गुंतागुंतीची
सूत्रे मिळतात. प्रत्येक संकारकाला सूत्रे वरचे
असे एक फलन अनुरूप मानता येते की ते संबंधित संकारक वापरून
एका किंवा अधिक सूत्रांपासून नवे सूत्र देते. उदाहरणार्थ,
$\varphi$ आणि $\psi$ ही दोन सूत्रे जोडून संयुती
बनवणारे फलन $\mathcal E _\land (\varphi,
\psi) = \varphi \land \psi$ असे लिहिता येते. यांना
\emph{सूत्ररचना-फलने} म्हणू. सूत्रांचा संच या
फलनांच्या खाली बंद आहे: $\varphi$ आणि $\psi$ ही
सूत्रे असतील, तर $\lnot \varphi$,
$\varphi \land
\psi$ इत्यादीही तशीच असतात.
\end{explain}

\begin{defn}[सूत्रे वरील विगमन]
प्रत्येक अण्वीय सूत्राला गुणधर्म~$P$ असेल आणि
सूत्र-रचना-फलनांच्या खाली $P$ जपला जात असेल, तर प्रत्येक
सूत्राला गुणधर्म~$P$ असतो.
\end{defn}

\begin{explain}
या व्याख्येतून सूत्रे वरील विगमनात्मक सिद्धतेची
“कृती” सुचते. प्रत्येक सूत्राला गुणधर्म~$P$
आहे हे दाखवायचे असल्यास पुढील पायऱ्या पाळा:\\

\textbf{विगमनाचा आधार.} $\varphi$ हे अण्वीय
सूत्र असो. [\ldots] म्हणून $\varphi$ ला गुणधर्म~$P$ असतो.

\textbf{विगमनाची पायरी.} $\varphi$ आणि $\psi$ ही दोन्ही
गुणधर्म~$P$ असलेली सूत्रे असोत.

प्रकरण $\lnot$: [\ldots] म्हणून $\lnot \varphi$ ला गुणधर्म~$P$ असतो.

प्रकरण $\land$: [\ldots] म्हणून $\varphi \land \psi$ ला गुणधर्म~$P$ असतो.

प्रकरण $\lor$: [\ldots] म्हणून $\varphi \lor \psi$ ला गुणधर्म~$P$ असतो.

प्रकरण $\lif$: [\ldots] म्हणून $\varphi \lif \psi$ ला गुणधर्म~$P$ असतो.

प्रकरण $\liff$: [\ldots] म्हणून $\varphi \liff \psi$ ला गुणधर्म~$P$ असतो.

प्रकरण $\lforall[]$: [\ldots] म्हणून $\lforall[x] \varphi$ ला गुणधर्म~$P$ असतो.

प्रकरण $\lexists[]$: [\ldots] म्हणून $\lexists[x] \varphi$ ला गुणधर्म~$P$ असतो. \\

म्हणून प्रत्येक सूत्राला गुणधर्म~$P$ असतो.
\end{explain}



\begin{quote}\small\textbf{स्रोतनोंद.} SOL6-A341: सूत्ररचना नेहमी दोनच सूत्रांपासून होत नाही हे दुरुस्त केले; बेरीजच्या आधारप्रसंगाची मर्यादा आधीच्या विधानाशी जुळवली आणि परिचय-नियमाची संज्ञा सुसंगत केली.\end{quote}










\section{संचांबद्दलच्या सिद्धता}\label{sfr:set:prf:sec}

\begin{editorial}
हा विभाग \verb|content/method/proofs| ने अधिक्रमित केला आहे आणि
या प्रकरणातून पूर्वनिर्धारितपणे काढला आहे.
\end{editorial}

\begin{explain}
आतापर्यंत ओळख करून दिलेले संच आणि संकेतलेखन
संच निर्दिष्ट करण्यासाठी व त्यांच्यातील संबंध व्यक्त
करण्यासाठी सोयीची लघुरूपे देतात. अशा संबंधांबद्दलचे
दावेही अनेकदा सिद्ध करावे लागतात. गणिती सिद्धतांची
ओळख नसेल, तर हे तुमच्यासाठी नवीन असेल. म्हणून एक
सोपे उदाहरण टप्प्याटप्प्याने पाहू. कोणत्याही संच~$X$
आणि $Y$ साठी $X \cap
(X \cup Y) = X$ हे नेहमी खरे असते, हे सिद्ध करू.
अशी संच-समानता कशी सिद्ध करायची? संच केवळ त्यांच्या
घटक वरून ठरतात, म्हणजे त्यांत तेच घटक
असतील तेव्हाच ते समान असतात, हे आठवा. म्हणून येथे
(a) $X \cap (X \cup
Y)$ मधील प्रत्येक घटक हा $X$ चाही घटक
आहे आणि, उलट, (b) $X$ मधील प्रत्येक घटक
हा $X \cap (X \cup Y)$ चाही घटक आहे, हे
सिद्ध करायचे आहे. दुसऱ्या शब्दांत, (a) $X \cap (X \cup Y) \subseteq X$
आणि (b) $X \subseteq X \cap (X \cup Y)$ हे दोन्ही दाखवतो.

“$X \cap (X
\cup Y)$ मधील प्रत्येक घटक~$z$ हा $X$ चाही घटक
आहे” असा सर्वसाधारण दावा सिद्ध करताना प्रथम कोणताही
$z \in X \cap (X \cup Y)$ दिला आहे असे मानतो आणि या
गृहीतकापासून $z \in X$ सिद्ध करतो. ही पद्धत “सर्वसाधारण
सोपाधिक सिद्धता” म्हणून परिचित असेल. या सिद्धतेत
गृहीतक आणि निष्कर्ष यांतील व्याख्यांचाही उपयोग करावा
लागेल; येथे उदाहरणार्थ “$\cap$” आणि “$\cup$” यांच्या
व्याख्यांचा. म्हणून (a) साठीचा युक्तिवाद असा:


\begin{quote}
(a) प्रथम $X \cap (X \cup Y) \subseteq X$ दाखवायचे आहे;
म्हणजे $\subseteq$ च्या व्याख्येनुसार कोणत्याही~$z$ साठी
$z \in X \cap (X \cup Y)$ असेल, तर $z \in X$ हे दाखवायचे.
म्हणून $z \in X \cap (X \cup
Y)$ असे माना. $z$ हा दोन संचांच्या छेदाचा घटक
तेव्हा आणि केवळ तेव्हा असतो जेव्हा तो दोन्ही संचांचा
घटक असतो. म्हणून $z \in X$ आणि $z \in X \cup Y$
असे निष्पन्न होते. विशेषतः $z \in X$, हेच दाखवायचे होते.
\end{quote}

याने सिद्धतेचा पहिला भाग पूर्ण झाला. शेवटच्या
पायरीत आपण एक तथ्य वापरले: एखाद्या गृहीतकापासून
संयुती ($z \in X$ आणि $z \in X
\cup Y$) निष्पन्न झाली, तर तिचा प्रत्येक संयुतीघटक
त्याच गृहीतकापासून निष्पन्न होतो. हा नियम “संयुती-निरसन”
किंवा $\land$Elim म्हणून परिचित असेल. आता (b) सिद्ध करू:

\begin{quote}
(b) आता $X \subseteq X \cap (X \cup Y)$ सिद्ध करू; म्हणजे
$\subseteq$ च्या व्याख्येनुसार कोणत्याही~$z$ साठी
$z \in X$ असेल, तर $z \in X \cap
(X \cup Y)$ हे दाखवू. $z \in X$ असे माना. $z \in X
\cap (X \cup Y)$ दाखवण्यासाठी “$\cap$” च्या व्याख्येनुसार
(i) $z \in X$ आणि (ii) $z \in X \cup Y$ हे दाखवावे
लागेल. (i) हे आपले गृहीतकच आहे; म्हणून अधिक
सिद्ध करण्याची गरज नाही. (ii) साठी आठवा की $z$
हा संचांच्या संयोगाचा घटक तेव्हा आणि केवळ तेव्हा
असतो जेव्हा तो त्या संचांपैकी किमान एकाचा घटक
असतो. $z \in X$ आणि $X \cup Y$ हा $X$ व $Y$ यांचा
संयोग असल्याने येथे ही अट पूर्ण होते. म्हणून
$z \in X \cup Y$. आपण (i) $z \in X$ आणि
(ii) $z \in X \cup Y$ ही दोन्ही दाखवली; त्यामुळे
“$\cap$” च्या व्याख्येनुसार $z \in X \cap (X \cup Y)$.
\end{quote}

हा युक्तिवाद काहीसा लांबलचक होता; पण संच आणि
त्यांच्यातील संबंधांविषयी आपण कसा विचार करतो हे
तो दाखवतो. सहसा आपण इतक्या स्पष्टपणे सर्व काही
लिहीत नाही; विशेषतः प्रत्येक व्याख्या पुन्हा सांगत
नाही. अधिक प्रगत पाठ्यात याच निकालाची सिद्धता
खूप संक्षिप्त असेल. ती पुढीलप्रमाणे दिसू शकते.
\end{explain}

\begin{prop}[शोषण]
सर्व संच $X$, $Y$ साठी,
\[
X \cap (X \cup Y) = X
\]
\end{prop}

\begin{proof}
(a) $z \in X \cap (X \cup Y)$ असे माना. मग $z \in X$;
म्हणून $X \cap (X
  \cup Y) \subseteq X$.

(b) आता $z \in X$ असे माना. मग $z \in X \cup Y$ सुद्धा;
त्यामुळे $z \in X \cap (X \cup Y)$ सुद्धा. म्हणून
$X \subseteq X \cap (X \cup
  Y)$.
\end{proof}

\begin{prob}
$X \cup (X \cap Y) = X$ हे तपशीलवार सिद्ध करा.
मग त्याची संक्षिप्त सिद्धता द्या. (सूचना:
$X \cup (X \cap Y) \subseteq X$ या दिशेसाठी
प्रकरणांनुसार सिद्धता, म्हणजे $\lor$Elim, लागेल.)
\end{prob}



\begin{quote}\small\textbf{स्रोतनोंद.} SOL6-A342: छेद आणि संयोगातील सदस्यत्वासाठी द्विदिशेचे शब्दरूप स्पष्ट केले; दोन्ही सिद्धता आणि सराव जपले.\end{quote}

\section{अंतःप्रज्ञावादी कट-निर्मूलनातील पूरक प्रसंग}
$A \ident (B \lif C)$ असेल, तर
$\pi$ चा शेवट असा असतो:
\[
\AxiomC{}
\RightLabel{$\pi_1'$}
\Deduce$B, \Gamma \fCenter \Delta, C$
\RightLabel{\RightR{\lif}}
\UnaryInf$\Gamma \fCenter \Delta, B \lif C$
\LeftSubproofLabel{$\pi_1$}
\AxiomC{}
\RightLabel{$\pi_2'$}
\Deduce$B \lif C, \Pi \fCenter \Lambda, B$
\AxiomC{}
\RightLabel{$\pi_2''$}
\Deduce$C, \Pi \fCenter \Lambda$
\RightLabel{\LeftR{\lif}}
\BinaryInf$B \lif C, \Pi \fCenter \Lambda$
\RightSubproofLabel{$\pi_2$}
\RightLabel{\Cut}
\BinaryInf$\Gamma, \Pi \fCenter \Delta, \Lambda$
\DisplayProof
\]
\Cut ने संपणारी सिद्धता पुढीलप्रमाणे:
\[
\AxiomC{}
\RightLabel{$\pi_1'$}
\Deduce$B, \Gamma \fCenter \Delta, C$
\RightLabel{\RightR{\lif}}
\UnaryInf$\Gamma \fCenter \Delta, B \lif C$
\LeftSubproofLabel{$\pi_1$}
\AxiomC{}
\RightLabel{$\pi_2'$}
\Deduce$B \lif C, \Pi \fCenter \Lambda, B$
\RightLabel{\Cut}
\BinaryInf$\Gamma, \Pi \fCenter \Delta, \Lambda, B$
\DisplayProof
\]
तिची कटची उंची
$\pi$ च्या कट-उंचीपेक्षा कमी आहे.
विगमन गृहीतकाने
$\Gamma, \Pi \fCenter \Delta, \Lambda, B$
याची सिद्धता~$\pi_1''$ मिळते.
\Cut ने संपणारी दुसरी
सिद्धता अशी आहे:
\[
\AxiomC{}
\RightLabel{$\pi_1'$}
\Deduce$B, \Gamma \fCenter \Delta, C$
\AxiomC{}
\RightLabel{$\pi_2''$}
\Deduce$C, \Pi \fCenter \Lambda$
\RightLabel{\Cut}
\BinaryInf$B, \Gamma, \Pi \fCenter \Delta, \Lambda$
\DisplayProof
\]
तिची कट-कोटी आणि
कटची उंची या दोन्ही
$\pi$ च्या कट-कोटी आणि
कटच्या उंची पेक्षा कमी आहेत.
म्हणून विगमन गृहीतक आणि
दुर्बलीकरण वापरून
$B, \Gamma, \Pi, \Pi \fCenter \Delta, \Lambda,
\Lambda$ याची
~\Log{G3c} मधील सिद्धता~$\pi_2'''$
मिळते. आता पुढील
\Cut ने संपणाऱ्या
सिद्धता वर विगमन
गृहीतक लावा:
\[
\AxiomC{}
\RightLabel{$\pi_1''$}
\Deduce$\Gamma, \Pi \fCenter \Delta, \Lambda, B$
\AxiomC{}
\RightLabel{$\pi_2'''$}
\Deduce$B, \Gamma, \Pi, \Pi \fCenter \Delta, \Lambda,
\Lambda$
\RightLabel{\Cut}
\BinaryInf$\Gamma, \Gamma, \Pi, \Pi, \Pi \fCenter \Delta, \Delta, \Lambda, \Lambda, \Lambda$
\DisplayProof
\]
(हिची कट-कोटी
$< \cutr{\pi}$ आहे.)
त्यातून
$\Gamma,
\Gamma, \Pi, \Pi, \Pi \fCenter \Delta, \Delta, \Lambda, \Lambda, \Lambda$
याची सिद्धता मिळते.
मग \ref{pt:seq:inv:prop:G3c-cont-adm}
लावून
$\Gamma, \Pi \fCenter \Delta, \Lambda$
याची सिद्धता~$\pi'$ मिळते.

\section{पूरक rules-G1i नियमसारणी}








\begin{table}
\begin{tabular}{@{}c@{}c@{}}
\multicolumn{2}{@{}c@{}}{स्वयंसिद्धके}\\[2ex]
$A \Sequent A$ & $\lfalse \Sequent \quad$
\\[2ex]
\multicolumn{2}{@{}c@{}}{संरचनात्मक नियम}\\[2ex]
\Axiom$ \Gamma \fCenter \Delta $
\RightLabel{\LeftR{\Weakening}}
\UnaryInf$ A, \Gamma \fCenter \Delta$
\DisplayProof
&
\Axiom$ \Gamma \fCenter $
\RightLabel{\RightR{\Weakening}}
\UnaryInf$ \Gamma \fCenter A$
\DisplayProof
\\[4ex]
\Axiom$ A, A, \Gamma \fCenter \Delta $
\RightLabel{\LeftR{\Contraction}}
\UnaryInf$ A, \Gamma \fCenter \Delta$
\DisplayProof
&
\\[4ex]
\multicolumn{2}{@{}c@{}}{तार्किक नियम}\\[2ex]
\Axiom$ \Gamma \fCenter A $
\RightLabel{\LeftR{\lnot}}
\UnaryInf$ \lnot A, \Gamma \fCenter $
\DisplayProof
&
\Axiom$A, \Gamma \fCenter $
\RightLabel{\RightR{\lnot}}
\UnaryInf$ \Gamma \fCenter \lnot A $
\DisplayProof
\\[4ex]\begin{tabular}[t]{@{}l@{}}
\Axiom$ A, \Gamma \fCenter \Delta$
\RightLabel{\LeftR{\land}}
\UnaryInf$ A \land B, \Gamma \fCenter \Delta$
\DisplayProof
\\[3ex]
\Axiom$B, \Gamma \fCenter \Delta$
\RightLabel{\LeftR{\land}}
\UnaryInf$A \land B, \Gamma \fCenter \Delta$
\DisplayProof\\[3ex]
\end{tabular}
&
\Axiom$\Gamma \fCenter A$
\Axiom$ \Gamma \fCenter B$
\RightLabel{\RightR{\land}}
\BinaryInf$ \Gamma \fCenter A \land B $
\DisplayProof
\\[4ex]\Axiom$A, \Gamma \fCenter \Delta$
\Axiom$B, \Gamma \fCenter \Delta$
\RightLabel{\LeftR{\lor}}
\BinaryInf$A \lor B, \Gamma \fCenter \Delta$
\DisplayProof
&
\begin{tabular}[t]{@{}r@{}}
\Axiom$\Gamma \fCenter A$
\RightLabel{\RightR{\lor}}
\UnaryInf$ \Gamma \fCenter A \lor B$
\DisplayProof
\\[3ex]
\Axiom$ \Gamma \fCenter B$
\RightLabel{\RightR{\lor}}
\UnaryInf$ \Gamma \fCenter A \lor B$
\DisplayProof\\[3ex]
\end{tabular}
\\[4ex]
\Axiom$ \Gamma \fCenter A$
\Axiom$ B, \Gamma \fCenter \Delta$
\RightLabel{\LeftR{\lif}}
\BinaryInf$ A \lif B, \Gamma \fCenter \Delta$
\DisplayProof
&
\Axiom$ A, \Gamma \fCenter B$
\RightLabel{\RightR{\lif}}
\UnaryInf$ \Gamma \fCenter A \lif B $
\DisplayProof
\\[4ex]
\Axiom$ A(t), \Gamma \fCenter \Delta$
\RightLabel{\LeftR{\lforall}}
\UnaryInf$ \lforall[x][A(x)],\Gamma \fCenter \Delta$
\DisplayProof
&
\Axiom$ \Gamma \fCenter A(a) $
\RightLabel{\RightR{\lforall}}
\UnaryInf$ \Gamma \fCenter \lforall[x][A(x)]$
\DisplayProof
\\[4ex]
\Axiom$ A(a), \Gamma \fCenter \Delta $
\RightLabel{\LeftR{\lexists}}
\UnaryInf$ \lexists[x][A(x)], \Gamma \fCenter \Delta$
\DisplayProof
&
\Axiom$ \Gamma \fCenter A(t) $
\RightLabel{\RightR{\lexists}}
\UnaryInf$ \Gamma \fCenter \lexists[x][A(x)]$
\DisplayProof
\end{tabular}
\caption{अंतःप्रज्ञावादी तर्कशास्त्रासाठी \Log{G1i} चे नियम.
$\RightR{\forall}$ आणि $\LeftR{\exists}$ मध्ये स्थिरांक~$a$ हा
निष्कर्ष-क्रमवर्तीमध्ये येऊ नये. $\Gamma$ हा सूत्रांचा
बहुसंच असून $\Delta$ रिकामा असू शकतो किंवा त्यात एकच
सूत्र असू शकते. \RightR{\Contraction} नियम नाही. \Log{G1m}
म्हणजे~\RightR{\Weakening} आणि तक्त्यातील
$\lfalse \Sequent \quad$ हे असत्याचे स्वयंसिद्धक
वगळलेली \Log{G1i}.}
\label{pt:seq:supp:tab:G1i}
\end{table}



\section{पूरक rules-G3i नियमसारणी}








\begin{table}
\begin{tabular}{@{}c@{}c@{}}
\multicolumn{2}{@{}c@{}}{स्वयंसिद्धके}\\[2ex]
$C, \Gamma \Sequent C$ & $\lfalse, \Gamma \Sequent \Delta$
\\[2ex]
\multicolumn{2}{@{}c@{}}{तार्किक नियम}\\[2ex]
\Axiom$\lnot A, \Gamma \fCenter A $
\RightLabel{\LeftR{\lnot}}
\UnaryInf$\lnot A, \Gamma \fCenter \Delta$
\DisplayProof
&
\Axiom$A, \Gamma \fCenter$
\RightLabel{\RightR{\lnot}}
\UnaryInf$ \Gamma \fCenter \lnot A $
\DisplayProof
\\[4ex]
\Axiom$ A, B, \Gamma \fCenter \Delta$
\RightLabel{\LeftR{\land}}
\UnaryInf$ A \land B, \Gamma \fCenter \Delta$
\DisplayProof
&
\Axiom$\Gamma \fCenter A$
\Axiom$ \Gamma \fCenter B$
\RightLabel{\RightR{\land}}
\BinaryInf$ \Gamma \fCenter A \land B $
\DisplayProof
\\[4ex]\Axiom$A, \Gamma \fCenter \Delta$
\Axiom$B, \Gamma \fCenter \Delta$
\RightLabel{\LeftR{\lor}}
\BinaryInf$A \lor B, \Gamma \fCenter \Delta$
\DisplayProof
&
\begin{tabular}[t]{@{}r@{}}
\Axiom$\Gamma \fCenter A$
\RightLabel{\RightR{\lor}}
\UnaryInf$ \Gamma \fCenter A \lor B$
\DisplayProof
\\[3ex]
\Axiom$ \Gamma \fCenter B$
\RightLabel{\RightR{\lor}}
\UnaryInf$ \Gamma \fCenter A \lor B$
\DisplayProof\\[3ex]
\end{tabular}
\\[4ex]
\Axiom$A \lif B, \Gamma \fCenter A$
\Axiom$B, \Gamma \fCenter \Delta$
\RightLabel{\LeftR{\lif}}
\BinaryInf$A \lif B, \Gamma \fCenter \Delta$
\DisplayProof
&
\Axiom$A, \Gamma \fCenter B$
\RightLabel{\RightR{\lif}}
\UnaryInf$ \Gamma \fCenter A \lif B $
\DisplayProof
\\[4ex]
\Axiom$ A(t), \lforall[x][A(x)], \Gamma \fCenter \Delta$
\RightLabel{\LeftR{\lforall}}
\UnaryInf$ \lforall[x][A(x)],\Gamma \fCenter \Delta$
\DisplayProof
&
\Axiom$ \Gamma \fCenter A(a) $
\RightLabel{\RightR{\lforall}}
\UnaryInf$ \Gamma \fCenter \lforall[x][A(x)]$
\DisplayProof
\\[4ex]
\Axiom$ A(a), \Gamma \fCenter \Delta $
\RightLabel{\LeftR{\lexists}}
\UnaryInf$ \lexists[x][A(x)], \Gamma \fCenter \Delta$
\DisplayProof
&
\Axiom$ \Gamma \fCenter A(t) $
\RightLabel{\RightR{\lexists}}
\UnaryInf$ \Gamma \fCenter \lexists[x][A(x)]$
\DisplayProof
\end{tabular}
\caption{अंतःप्रज्ञावादी तर्कशास्त्रासाठी \Log{G3i} चे नियम.
$\RightR{\forall}$ आणि $\LeftR{\exists}$ मध्ये स्थिरांक~$a$ हा
निष्कर्ष-क्रमवर्तीमध्ये येऊ नये. $\Gamma$ हा सूत्रांचा
बहुसंच असून $\Delta$ रिकामा असू शकतो किंवा त्यात एकच
सूत्र असू शकते. स्वयंसिद्धकांत सूत्र~$C$ अण्वीय
असावे. \Log{G1m} म्हणजे \Log{G1i} मधून
\RightR{\Weakening} आणि $\lfalse \Sequent \quad$
हे स्वयंसिद्धक वगळून मिळणारी पद्धत.}
\label{pt:seq:supp:tab:G3i}
\end{table}



\section{पूरक rules-LK नियमसारणी}







\begin{table}
\begin{tabular}{@{}c@{}c@{}}
\multicolumn{2}{@{}c@{}}{स्वयंसिद्धके}\\[2ex]
$A \Sequent A$ & $\lfalse \Sequent \quad$
\\[2ex]
\multicolumn{2}{@{}c@{}}{संरचनात्मक नियम}\\[2ex]
\Axiom$ \Gamma \fCenter \Delta $
\RightLabel{\LeftR{\Weakening}}
\UnaryInf$ A, \Gamma \fCenter \Delta$
\DisplayProof
&
\Axiom$ \Gamma \fCenter \Delta$
\RightLabel{\RightR{\Weakening}}
\UnaryInf$ \Gamma \fCenter \Delta, A$
\DisplayProof
\\[4ex]
\Axiom$ A, A, \Gamma \fCenter \Delta $
\RightLabel{\LeftR{\Contraction}}
\UnaryInf$ A, \Gamma \fCenter \Delta$
\DisplayProof
&
\Axiom$ \Gamma \fCenter \Delta, A, A$
\RightLabel{\RightR{\Contraction}}
\UnaryInf$ \Gamma \fCenter \Delta, A$
\DisplayProof
\\[4ex]
\Axiom$ \Gamma, A, B, \Pi \fCenter \Delta $
\RightLabel{\LeftR{\Exchange}}
\UnaryInf$\Gamma, B, A, \Pi \fCenter \Delta$
\DisplayProof
&
\Axiom$ \Gamma \fCenter \Delta, A, B, \Lambda$
\RightLabel{\RightR{\Exchange}}
\UnaryInf$ \Gamma \fCenter \Delta, B, A, \Lambda$
\DisplayProof
\\[4ex]
\multicolumn{2}{@{}c@{}}{तार्किक नियम}\\[2ex]
\Axiom$ \Gamma \fCenter \Delta, A $
\RightLabel{\LeftR{\lnot}}
\UnaryInf$ \lnot A, \Gamma \fCenter \Delta$
\DisplayProof
&
\Axiom$A, \Gamma \fCenter \Delta$
\RightLabel{\RightR{\lnot}}
\UnaryInf$ \Gamma \fCenter \Delta, \lnot A $
\DisplayProof
\\[4ex]\begin{tabular}[t]{@{}l@{}}
\Axiom$ A, \Gamma \fCenter \Delta$
\RightLabel{\LeftR{\land}}
\UnaryInf$ A \land B, \Gamma \fCenter \Delta$
\DisplayProof
\\[3ex]
\Axiom$B, \Gamma \fCenter \Delta$
\RightLabel{\LeftR{\land}}
\UnaryInf$A \land B, \Gamma \fCenter \Delta$
\DisplayProof\\[3ex]
\end{tabular}
&
\Axiom$\Gamma \fCenter \Delta, A$
\Axiom$ \Gamma \fCenter \Delta, B$
\RightLabel{\RightR{\land}}
\BinaryInf$ \Gamma \fCenter \Delta, A \land B $
\DisplayProof
\\[4ex]\Axiom$A, \Gamma \fCenter \Delta$
\Axiom$B, \Gamma \fCenter \Delta$
\RightLabel{\LeftR{\lor}}
\BinaryInf$A \lor B, \Gamma \fCenter \Delta$
\DisplayProof
&
\begin{tabular}[t]{@{}r@{}}
\Axiom$\Gamma \fCenter \Delta, A$
\RightLabel{\RightR{\lor}}
\UnaryInf$ \Gamma \fCenter \Delta, A \lor B$
\DisplayProof
\\[3ex]
\Axiom$ \Gamma \fCenter \Delta, B$
\RightLabel{\RightR{\lor}}
\UnaryInf$ \Gamma \fCenter \Delta, A \lor B$
\DisplayProof\\[3ex]
\end{tabular}
\\[4ex]
\Axiom$ \Gamma \fCenter \Delta, A$
\Axiom$ B, \Gamma \fCenter \Delta$
\RightLabel{\LeftR{\lif}}
\BinaryInf$ A \lif B, \Gamma \fCenter \Delta$
\DisplayProof
&
\Axiom$ A, \Gamma \fCenter \Delta, B$
\RightLabel{\RightR{\lif}}
\UnaryInf$ \Gamma \fCenter \Delta, A \lif B $
\DisplayProof
\\[4ex]
\Axiom$ A(t), \Gamma \fCenter \Delta$
\RightLabel{\LeftR{\lforall}}
\UnaryInf$ \lforall[x][A(x)],\Gamma \fCenter \Delta$
\DisplayProof
&
\Axiom$ \Gamma \fCenter \Delta, A(a) $
\RightLabel{\RightR{\lforall}}
\UnaryInf$ \Gamma \fCenter \Delta, \lforall[x][A(x)]$
\DisplayProof
\\[4ex]
\Axiom$ A(a), \Gamma \fCenter \Delta $
\RightLabel{\LeftR{\lexists}}
\UnaryInf$ \lexists[x][A(x)], \Gamma \fCenter \Delta$
\DisplayProof
&
\Axiom$ \Gamma \fCenter \Delta, A(t) $
\RightLabel{\RightR{\lexists}}
\UnaryInf$ \Gamma \fCenter \Delta, \lexists[x][A(x)]$
\DisplayProof
\end{tabular}
\caption{\Log{LK} चे नियम. $\RightR{\forall}$ आणि
$\LeftR{\exists}$ मध्ये स्थिरांक~$a$ हा
निष्कर्ष-क्रमवर्तीमध्ये येऊ नये. $\Gamma$, $\Delta$, $\Pi$ आणि~$\Lambda$
या सूत्रांच्या क्रमिका आहेत.}
\label{pt:seq:supp:tab:LK}
\end{table}



\section{पूरक rules-mG3i नियमसारणी}








\begin{table}
\begin{tabular}{@{}c@{}c@{}}
\multicolumn{2}{@{}c@{}}{स्वयंसिद्धके}\\[2ex]
$C, \Gamma \Sequent \Delta, C$ & $\lfalse, \Gamma \Sequent \Delta$
\\[2ex]
\multicolumn{2}{@{}c@{}}{तार्किक नियम}\\[2ex]
\Axiom$\lnot A, \Gamma \fCenter \Delta, A $
\RightLabel{\LeftR{\lnot}}
\UnaryInf$\lnot A, \Gamma \fCenter \Delta$
\DisplayProof
&
\Axiom$A, \Gamma \fCenter$
\RightLabel{\RightR{\lnot}}
\UnaryInf$ \Gamma \fCenter \Delta, \lnot A $
\DisplayProof
\\[4ex]
\Axiom$ A, B, \Gamma \fCenter \Delta$
\RightLabel{\LeftR{\land}}
\UnaryInf$ A \land B, \Gamma \fCenter \Delta$
\DisplayProof
&
\Axiom$\Gamma \fCenter \Delta, A$
\Axiom$ \Gamma \fCenter \Delta, B$
\RightLabel{\RightR{\land}}
\BinaryInf$ \Gamma \fCenter \Delta, A \land B $
\DisplayProof
\\[4ex]\Axiom$A, \Gamma \fCenter \Delta$
\Axiom$B, \Gamma \fCenter \Delta$
\RightLabel{\LeftR{\lor}}
\BinaryInf$A \lor B, \Gamma \fCenter \Delta$
\DisplayProof
&
\Axiom$\Gamma \fCenter \Delta, A, B$
\RightLabel{\RightR{\lor}}
\UnaryInf$ \Gamma \fCenter \Delta, A \lor B$
\DisplayProof
\\[4ex]
\Axiom$A \lif B, \Gamma \fCenter \Delta, A$
\Axiom$ B, \Gamma \fCenter \Delta$
\RightLabel{\LeftR{\lif}}
\BinaryInf$A \lif B, \Gamma \fCenter \Delta$
\DisplayProof
&
\Axiom$ A, \Gamma \fCenter B$
\RightLabel{\RightR{\lif}}
\UnaryInf$ \Gamma \fCenter \Delta, A \lif B $
\DisplayProof
\\[4ex]
\Axiom$ A(t), \lforall[x][A(x)], \Gamma \fCenter \Delta$
\RightLabel{\LeftR{\lforall}}
\UnaryInf$ \lforall[x][A(x)],\Gamma \fCenter \Delta$
\DisplayProof
&
\Axiom$ \Gamma \fCenter A(a) $
\RightLabel{\RightR{\lforall}}
\UnaryInf$ \Gamma \fCenter \lforall[x][A(x)]$
\DisplayProof
\\[4ex]
\Axiom$ A(a), \Gamma \fCenter \Delta $
\RightLabel{\LeftR{\lexists}}
\UnaryInf$ \lexists[x][A(x)], \Gamma \fCenter \Delta$
\DisplayProof
&
\Axiom$ \Gamma \fCenter \Delta, \lexists[x][A(x)], A(t) $
\RightLabel{\RightR{\lexists}}
\UnaryInf$ \Gamma \fCenter \Delta, \lexists[x][A(x)]$
\DisplayProof
\end{tabular}
\caption{अंतःप्रज्ञावादी तर्कशास्त्रासाठी बहुनिष्कर्षी \Log{mG3i} चे नियम.
$\RightR{\forall}$ आणि $\LeftR{\exists}$ मध्ये स्थिरांक~$a$ हा
निष्कर्ष-क्रमवर्तीमध्ये येऊ नये. $\Gamma$ आणि $\Delta$ हे सूत्रांचे
बहुसंच आहेत. स्वयंसिद्धकांत सूत्र~$C$ अण्वीय असावे.}
\label{pt:seq:supp:tab:mG3i}
\end{table}



\section{मूळ संपूर्ण-सामग्री चालकाची संपादकीय नोंद}
खालील नोंद गोठवलेल्या मूळ संपूर्ण-सामग्री चालकाची आहे; प्रस्तुत ग्रंथाच्या संपादकीय स्थितीचा ती नवा दावा करत नाही.
\begin{editorial}
ही फाइल ओपन लॉजिक प्रकल्पातील सर्व सामग्री लोड करते. अशा
संपादकीय नोंदी दिसत असल्यास, ही सामग्री कशी सादर केली जाणार
याचा विचार न करता फाइल संकलित केली गेली आहे, हे त्यावरून
समजते. तुम्हाला हे वाचता येत असेल, तर या PDF वरून शिकवणे
किंवा अभ्यास करणे बहुधा \emph{योग्य नाही}.

अध्यापनासाठी किंवा स्व-अध्ययनासाठी अधिक योग्य मजकूर तयार
करता यावा यासाठी ओपन लॉजिक प्रकल्पात अनेक सोयी आहेत.
उदाहरणार्थ, नेहमीच्या मांडणीत सर्व तार्किक संयोजक आदिम
मानले जातात आणि पुराव्यांमध्ये प्रत्येक संयोजकाचे सर्व प्रकार
तपासले जातात. पण त्यांपैकी काही प्रकार सरावासाठी सोडणे अधिक
चांगले असते. ओपन लॉजिक प्रकल्पाचे कामही अजून सुरू आहे.
सहयोग आणि सुधारणा वाढाव्यात म्हणून मसुदा स्वरूपातील किंवा
सरावप्रश्न नसलेली सामग्रीसुद्धा येथे समाविष्ट आहे. अभ्यासक्रमासाठी
तयार केलेल्या PDF मध्ये असे विभाग वगळले जातात.

अध्यापन आणि अभ्यासासाठी अधिक योग्य PDF पाहण्यासाठी
\href{http://builds.openlogicproject.org/}{OLP संकेतस्थळावरील नमुना अभ्यासक्रम}
पहा. स्वतःचा मजकूर तयार करण्यासाठी
\href{https://github.com/OpenLogicProject/OpenLogic/tree/master/courses/sample}{नमुना चालक फाइल}
वापरून सुरुवात करता येईल; अधिक सजवलेल्या आणि प्रगत उदाहरणांसाठी
या प्रकल्पातून तयार झालेल्या पाठ्यपुस्तकांचे स्रोत पाहता येतील.
\end{editorial}
\section{मूळ विभाग-चालकांच्या संपादकीय नोंदी}\label{source-driver:notes}
खालील नोंदी मूळ विभाग आणि प्रकरण-चालकांमधील आहेत. त्या मूळ ग्रंथाच्या संकलनमार्गांविषयी सांगतात; प्रस्तुत आवृत्तीविषयी नवे संपादकीय दावे नाहीत.
\subsection{संचांचे आकारमान — मूळ संपादकीय नोंद}\label{source-driver:OLP-0027}
\begin{editorial}
या प्रकरणात प्रगणने, गणनीयता आणि अगणनीयता यांची चर्चा आहे.
काही विभागांच्या दोन आवृत्त्या आहेत: अधिक प्राथमिक आवृत्तीत प्रगणने
म्हणजे याद्या किंवा $\PosInt$ पासूनची आच्छादने मानली आहेत;
अधिक अमूर्त आवृत्तीत प्रगणनांची व्याख्या $\Nat$ बरोबरच्या
एकास-एक आच्छादनांनी केली आहे.
\end{editorial}
\begin{editorial}
पुढील \ref{sfr:siz:enm-alt:sec},
\ref{sfr:siz:nen-alt:sec} आणि \ref{sfr:siz:red-alt:sec}
हे अनुक्रमे \ref{sfr:siz:enm:sec},
\ref{sfr:siz:nen:sec} आणि \ref{sfr:siz:red:sec}
यांचे पर्यायी विभाग आहेत. टिम बटन यांनी त्यांच्या Open Set Theory
या ग्रंथात वापरण्यासाठी ते तयार केले. ते किंचित अधिक प्रगत आहेत आणि
संच सिद्धांताच्या संदर्भात अधिक योग्य अशी गणनीयतेची भिन्न व्याख्या
वापरतात: यादी देता येणे किंवा $\PosInt$ पासूनच्या आच्छादक फलनाची
व्याप्ती असणे याऐवजी $\Nat$ किंवा त्याच्या आरंभीच्या खंडाशी एकास-एक
आच्छादन असणे.
\end{editorial}
\subsection{प्रथम-क्रम तर्कशास्त्र — मूळ संपादकीय नोंद}\label{source-driver:OLP-0138}
\begin{editorial}
  या भागात प्रथम-क्रम तर्कशास्त्राची अधिउपपत्ती संपूर्णतेपर्यंत मांडली
  आहे. सध्या ती विधानीय तर्कशास्त्राच्या स्वतंत्र विवेचनावर अवलंबून
  नाही; प्रत्येक गोष्ट सिद्ध केली आहे. ``FOL'' हा टॅग false असल्यास
  स्रोत फायली संख्यापकांवरील सामग्री वगळतात (आणि
  ``संरचना'' ऐवजी ``सत्य-मूल्यांकन'', $\Struct{M}$ ऐवजी
  $\pAssign{v}$, इत्यादी प्रतिस्थापित करतात). प्रत्यक्षात, विधानीय
  तर्कशास्त्राच्या भागातील बहुतेक सामग्री म्हणजे ``FOL'' टॅग बंद
  ठेवलेली प्रथम-क्रम सामग्रीच आहे.

  विधानीय तर्कशास्त्राचा भाग समाविष्ट केल्यास त्यामुळे बरीच
  पुनरावृत्ती होते. मात्र, या भागाला विधानीय तर्कशास्त्रातील सामग्री
  विचारात घेता यावी (आणि आधीच समाविष्ट केलेली सामग्री वगळता यावी,
  तसेच विधानीय भागातील संबंधित ठिकाणांचे संदर्भ देऊन सिद्धता
  संक्षिप्त करता याव्यात), अशी सोय करण्याचे नियोजन आहे.
\end{editorial}
\subsection{प्रतिमान उपपत्ती — मूळ संपादकीय नोंद}\label{source-driver:OLP-0182}
\begin{editorial}
  प्रतिमान उपपत्तीवरील साहित्य अपूर्ण आणि प्रयोगात्मक आहे. सध्या ते केवळ
  आल्डो अँटोनेली यांच्या प्रतिमान उपपत्तीवरील टिपणांचे रूपांतर आहे; मात्र
  प्रथम-क्रम तर्कशास्त्राच्या भागात समाविष्ट असलेले विषय (उपपत्त्या, संपूर्णता,
  संहतता) त्यातून वगळले आहेत. पाठ्यपुस्तकासाठी उपयुक्त ठरण्याकरिता त्यात बरीच
  अधिक प्रस्तावना, प्रेरणा व स्पष्टीकरण तसेच सरावप्रश्नांची गरज आहे. अँडी अराना
  खास या भागावर काम करण्याचे नियोजन करीत आहेत (\gitissue{65}).
\end{editorial}
\subsection{संगणनक्षमता — मूळ संपादकीय नोंद}\label{source-driver:OLP-0208}
\begin{editorial}
  हा भाग जेरेमी अविगॅड यांच्या संगणनक्षमता सिद्धांतावरील टिपांवर आधारित
  आहे. अजूनपर्यंत फक्त पुनरावर्ती फलनांच्या प्रकरणात सराव आहेत; तसेच
  प्रेरणा, उदाहरणे, तपशील आणि सराव घालून सर्वच सामग्रीचा विस्तार करणे
  उपयुक्त ठरेल.
\end{editorial}
\subsection{मोडल टॅब्लो — मूळ संपादकीय नोंद}\label{source-driver:OLP-0460}
\begin{editorial}
  मोडल तर्कशास्त्रातील पूर्वचिन्हयुक्त टॅब्लोंवरील
  हे प्रकरण अजून मसुदा आहे. अधिक उदाहरणे,
  संपूर्णतेच्या सिद्धता, आणि बंद टॅब्लो शोधण्याचे
  प्रयत्न निष्फळ ठरल्यास त्यातून प्रतिप्रतिमाने
  कशी मिळवता येतात याची चर्चा आवश्यक आहे.
\end{editorial}
\subsection{मोडल क्रमवर्ती कलन — मूळ संपादकीय नोंद}\label{source-driver:OLP-0470}
\begin{editorial}
  मोडल तर्कशास्त्रातील क्रमवर्ती कलनांवरील हे प्रकरण
  अजून मसुदा आहे. अधिक उदाहरणे, निर्दोषतेचे आणि
  संपूर्णतेचे पुरावे आवश्यक आहेत.
\end{editorial}
\subsection{कालिक तर्कशास्त्रे — मूळ संपादकीय नोंद}\label{source-driver:OLP-0476}
\begin{editorial}
  या प्रकरणात कालिक तर्कशास्त्रांचा आढावा घेतला आहे.
\end{editorial}
\subsection{संच, संबंध, फलने — मूळ संपादकीय नोंद}\label{source-driver:OLP-0720}
\begin{editorial}
  या फाइलमध्ये अनौपचारिक संचसिद्धान्तावरील या भागाची
  पहिली चार प्रकरणे मूळ, संथ गतीने मांडलेल्या आवृत्तीत
  समाविष्ट आहेत. संपूर्ण भाग
  \verb|sets-functions-relations-complete.tex| मध्ये आहे.
  या भागातील सामग्री मूलभूत अनौपचारिक संचसिद्धान्ताची
  बऱ्यापैकी पूर्ण ओळख करून देते. विद्यार्थ्यांना ही
  पार्श्वभूमी आधीच आहे असे गृहीत धरता येत नसेल, तर
  अभ्यासक्रमाची सुरुवात या सामग्रीच्या उजळणीने,
  निदान संच, फलने आणि संबंध या भागाच्या उजळणीने,
  करणे बहुधा योग्य ठरेल. त्यामुळे पुढील कोणत्याही
  तार्किक विभागांसाठी आवश्यक असलेल्या गणिती
  संकेतलेखनाचे मूलभूत कौशल्य सर्व विद्यार्थ्यांना
  प्राप्त होईल.
\end{editorial}
\subsection{संचांचे आकारमान — मूळ संपादकीय नोंद}\label{source-driver:OLP-0722}
\begin{editorial}
या प्रकरणात प्रगणने, गणनीयता आणि अगणनीयता यांची चर्चा आहे.
काही विभागांच्या दोन आवृत्त्या आहेत: संचसिद्धान्तावर अधिक
तपशीलवार चर्चा होत असल्यास वापरण्यासाठी एक अधिक प्राथमिक
आणि एक अधिक प्रगत आवृत्ती. या फाइलमध्ये फक्त प्राथमिक
आवृत्त्या समाविष्ट आहेत. अधिक प्रगत आवृत्त्या
\verb|size-of-sets-complete| मध्ये आहेत.
\end{editorial}
\chapter*{स्रोतावरील संपादकीय नोंदी}
\markboth{स्रोतावरील संपादकीय नोंदी}{स्रोतावरील संपादकीय नोंदी}
\addcontentsline{toc}{chapter}{स्रोतावरील संपादकीय नोंदी}
\section*{भाषांतरित स्रोत-एककांतील दुरुस्ती-नोंदी}
खालील नोंदी भाषांतरित स्रोत-एककांतील स्पष्ट दुरुस्त्या आणि निवडी सांगतात. विभागाजवळ आधीच दिसणाऱ्या नोंदी येथे पुन्हा दिलेल्या नाहीत.
\begin{quote}\small\textbf{OLP-0054 — स्रोतनोंद.} MRINF-003 पात्रता: गोठवलेल्या इंग्रजीतील अंतःस्थ cardeq रचना A\ensuremath{\approx}B\ensuremath{\approx}C अशी वैध रीतीने विस्तारते; ती स्रोतदोष नाही. मराठीतील स्पष्ट B\ensuremath{\approx}C निष्कर्ष त्याच तुल्यबलता-साखळीतील, पुढील उपयोगासाठी केंद्रित केलेली सममूल्य मांडणी आहे.\end{quote}
\begin{quote}\small\textbf{OLP-0058 — स्रोतनोंद.} MRPL-001: गोठवलेल्या इंग्रजीतील पहिल्या संक्षेपात जुळणारा आरंभीचा कंस नसतानाही शेवटचा कंस आहे. मराठीत मानक नकार-विकल्पयोग सूत्र दिले असून इंग्रजी बाइट्स बदललेले नाहीत.\end{quote}
\begin{quote}\small\textbf{OLP-0119 — स्रोतनोंद.} OLAXD-001: गोठवलेल्या इंग्रजी सूत्रात सर्वांत बाहेरचा उजवा कंस सुटला आहे. मराठी सूत्रात तो कंस जोडला असून गोठवलेले इंग्रजी बाइट्स बदललेले नाहीत.\end{quote}
\begin{quote}\small\textbf{OLP-0120 — स्रोतनोंद.} OLAXD-002: गोठवलेल्या इंग्रजी सूत्रात सर्वांत बाहेरचा एक उजवा कंस सुटला आहे. मराठी सूत्रात तो कंस जोडला असून गोठवलेले इंग्रजी बाइट्स बदललेले नाहीत.\end{quote}
\begin{quote}\small\textbf{OLP-0120 — स्रोतनोंद.} OLAXD-003: गोठवलेल्या इंग्रजी निष्कर्षात आवश्यक A-सशर्त भाग सुटला आहे. मराठी निष्कर्षात तो पुनर्स्थापित केला असून गोठवलेले इंग्रजी बाइट्स बदललेले नाहीत.\end{quote}
\begin{quote}\small\textbf{OLP-0122 — स्रोतनोंद.} OLAXD-004: गोठवलेल्या इंग्रजी मजकुरात दुसरा संदर्भ चुकून पुन्हा ax:land1 असा आहे. दुसऱ्या संयुक्ताचे निगमन सिद्ध करण्यासाठी मराठीत तो ax:land2 असा दुरुस्त केला असून गोठवलेले इंग्रजी बाइट्स बदललेले नाहीत.\end{quote}
\begin{quote}\small\textbf{OLP-0122 — स्रोतनोंद.} OLAXD-005: गोठवलेल्या इंग्रजी मजकुरात पहिला संदर्भ ax:lnot1 असा आहे; पण दाखवलेली दोन्ही सूत्रे थेट ax:lnot2 ची उदाहरणे आहेत. मराठीत संदर्भ दुरुस्त केला असून गोठवलेले इंग्रजी बाइट्स बदललेले नाहीत.\end{quote}
\begin{quote}\small\textbf{OLP-0123 — स्रोतनोंद.} OLAXD-006: गोठवलेल्या इंग्रजी मजकुरात शेवटचे पाऊल पुन्हा निगमन प्रमेयाने मिळते असे म्हटले आहे. प्रत्यक्षात आधीच्या सशर्त निष्कर्षावर सत्य स्थिरांकाचे स्वयंसिद्धक आणि विध्यात्मक अनुमान लावल्यावर अपेक्षित निष्कर्ष मिळतो; मराठीत हे कारण स्पष्ट केले असून गोठवलेले इंग्रजी बाइट्स बदललेले नाहीत.\end{quote}
\begin{quote}\small\textbf{OLP-0124 — स्रोतनोंद.} OLAXD-007: गोठवलेल्या इंग्रजी मजकुरातील QR निष्कर्षाच्या संख्यापक-व्याप्तीत B(x) आधी आवश्यक उद्गारचिन्ह सुटले आहे. मराठीत ते {!}B(x) असे पुनर्स्थापित केले असून गोठवलेले इंग्रजी बाइट्स बदललेले नाहीत.\end{quote}
\begin{quote}\small\textbf{OLP-0124 — स्रोतनोंद.} OLAXD-008: गोठवलेल्या इंग्रजी मजकुरातील एका पूर्तता-विधानात B(c) आधी आवश्यक उद्गारचिन्ह सुटले आहे. मराठीत ते {!}B(c) असे पुनर्स्थापित केले असून गोठवलेले इंग्रजी बाइट्स बदललेले नाहीत.\end{quote}
\begin{quote}\small\textbf{OLP-0125 — स्रोतनोंद.} OLAXD-009: गोठवलेल्या इंग्रजी स्रोताच्या अध्याय-शीर्ष टिप्पणीमध्ये जुने axiomatic-proofs नाव आहे; मार्ग आणि olfileid यांच्याशी सुसंगत असे axiomatic-deduction नाव येथे वापरले आहे.\end{quote}
\begin{quote}\small\textbf{OLP-0130 — स्रोतनोंद.} OLCOM-001: गोठवलेल्या इंग्रजी स्रोताच्या व्याख्याविषयक टिपेत सार्वत्रिक सूत्रातील A\_n नंतर (x\_n) सुटले आहे. मराठीत ते पुनर्स्थापित केले असून गोठवलेले इंग्रजी बाइट्स बदललेले नाहीत.\end{quote}
\begin{quote}\small\textbf{OLP-0132 — स्रोतनोंद.} OLCOM-002: गोठवलेल्या इंग्रजी स्रोताच्या सर्ववाची विगमन-प्रसंगातील अंतिम सूत्रात B ऐवजी A आले आहे. मराठीत प्रसंगाशी जुळणारे B पुनर्स्थापित केले असून गोठवलेले इंग्रजी बाइट्स बदललेले नाहीत.\end{quote}
\begin{quote}\small\textbf{OLP-0132 — स्रोतनोंद.} OLCOM-003: गोठवलेल्या इंग्रजी स्रोताच्या सत्यता पूर्वप्रमेयातील संख्यकीकृत प्रसंगांत चार वेळा पदे म्हटले आहे; संदर्भित निष्पत्ती आणि संतृप्त-दृष्टांत गुणधर्म बंद पदांसाठी आहेत. मराठीत बंद पदे असे स्पष्ट केले असून गोठवलेले इंग्रजी बाइट्स बदललेले नाहीत.\end{quote}
\begin{quote}\small\textbf{OLP-0133 — स्रोतनोंद.} OLCOM-004: गोठवलेल्या इंग्रजी स्रोताच्या फलन-पदातील t\_\{i+1\} नंतर सलग दोन स्वल्पविराम आहेत. मराठीत एकच स्वल्पविराम ठेवला असून गोठवलेले इंग्रजी बाइट्स बदललेले नाहीत.\end{quote}
\begin{quote}\small\textbf{OLP-0133 — स्रोतनोंद.} OLCOM-005: गोठवलेल्या इंग्रजी स्रोताच्या प्रतिनिधी-स्वातंत्र्य उदाहरणात दुसऱ्या प्रतिनिधी t' विषयी बोलताना न-पूर्ति सूत्रात पुन्हा t आले आहे. मराठीत t' पुनर्स्थापित केले असून गोठवलेले इंग्रजी बाइट्स बदललेले नाहीत.\end{quote}
\begin{quote}\small\textbf{OLP-0136 — स्रोतनोंद.} OLCOM-006: गोठवलेल्या इंग्रजी स्रोताच्या notFOL शाखेत सत्यता पूर्वप्रमेयासाठी prop:ccs च्या जागी prop:fsat-ccs वापरण्याचा गंतव्य-संदर्भ FOL टॅगच्या आत गेल्यामुळे दिसत नाही. मराठीत समान संदर्भ notFOL शाखेलाही स्पष्टपणे लागू केला असून सर्व गोठवलेले संदर्भ जतन केले आहेत.\end{quote}
\begin{quote}\small\textbf{OLP-0140 — स्रोतनोंद.} OLFOL-001: गोठवलेल्या इंग्रजी मजकुरातील या तार्किक निष्पन्नता-सूत्रात आणि पुढील दोन पुनरुक्तींमध्ये सार्वत्रिक सूत्राचा बंद करणारा चौकोनी कंस चुकीच्या जागी आहे किंवा गहाळ आहे, तर निष्कर्षानंतर एक जादा कंस आहे. मराठी मजकुरात तिन्ही ठिकाणी समान, सुव्यवस्थित आधारविधाने आणि निष्कर्ष दिले आहेत.\end{quote}
\begin{quote}\small\textbf{OLP-0143 — स्रोतनोंद.} OLFOL-002: गोठवलेल्या इंग्रजी मजकुरात अनेक स्थाने असू शकणारी चिन्हे constants म्हटली आहेत. स्थिरचिन्हे एकच वस्तू निर्देशित करतात; अनेक स्थाने असू शकणारी चिन्हे predicates असल्यामुळे मराठी मजकुरात ते दुरुस्त केले आहे.\end{quote}
\begin{quote}\small\textbf{OLP-0143 — स्रोतनोंद.} OLFOL-003: याच उदाहरणात प्रांत \{0, 1, 2\} असा ठरवला आहे; गोठवलेल्या इंग्रजी मजकुरातील \{1, 2, 3\} ही यादी 0 वगळते आणि प्रांताबाहेरील 3 समाविष्ट करते. मराठी मजकुरात प्रांताशी जुळणारी \{0, 1, 2\} ही यादी वापरली आहे.\end{quote}
\begin{quote}\small\textbf{OLP-0146 — स्रोतनोंद.} OLFOL-004: गोठवलेल्या इंग्रजीतील पहिल्या सार्वत्रिक सूत्रात Atom चा दुसरा फलसाधक चुकीने सार्वत्रिक सूत्राबाहेर ठेवला आहे. पुढील सर्वसाधारण रूपाशी जुळण्यासाठी मराठीत P(v\_0) हा पूर्ण आण्विक सूत्रभाग संख्यापकाच्या व्याप्तीत ठेवला आहे.\end{quote}
\begin{quote}\small\textbf{OLP-0152 — स्रोतनोंद.} OLFOL-005: गोठवलेल्या इंग्रजीतील पहिल्या संक्षेपात जुळणारा आरंभीचा कंस नसतानाही शेवटचा कंस आहे. मराठीत जुळत नसलेला कंस वगळला असून इंग्रजी बाइट्स बदललेले नाहीत.\end{quote}
\begin{quote}\small\textbf{OLP-0156 — स्रोतनोंद.} OLFOL-006: गोठवलेल्या इंग्रजी व्याख्येत k-स्थानी फलनासाठी m\_0 ते m\_k असे k+1 निर्देशांक आणि तितकेच युक्तिवाद दिले आहेत. येथे k युक्तिवादांसाठी m\_1 ते m\_k ही सुसंगत मांडणी वापरली आहे.\end{quote}
\begin{quote}\small\textbf{OLP-0156 — स्रोतनोंद.} OLFOL-007: हा पुरावा सामान्य भाषा L साठी आहे; गोठवलेल्या इंग्रजीतील या परिच्छेदाच्या दोन ठिकाणी चुकून विधानीय भाषा L\_0 आली आहे. येथे प्रमेयाशी सुसंगत L वापरले आहे.\end{quote}
\begin{quote}\small\textbf{OLP-0156 — स्रोतनोंद.} OLFOL-008: गोठवलेल्या इंग्रजीत या एकाच रचना-बाबीत विन्यासात्मक अनन्यता-चिन्हाऐवजी पर्यायी सममूल्यता-चिन्ह आले आहे. व्याख्या आणि इतर सर्व बाबींशी सुसंगत विन्यासात्मक अनन्यता येथे वापरली आहे.\end{quote}
\begin{quote}\small\textbf{OLP-0161 — स्रोतनोंद.} OLFOL-010: गोठवलेल्या इंग्रजी मजकुरातील ``single-two place relation'' हे जोडचिन्ह चुकीचे आहे; व्याख्येला आवश्यक असलेला एकच द्विस्थानी संबंध मराठीत स्पष्ट केला आहे.\end{quote}
\begin{quote}\small\textbf{OLP-0162 — स्रोतनोंद.} OLFOL-011: गोठवलेल्या इंग्रजी प्रदर्शनात पहिल्या ओळीच्या शेवटी आणि पुढील ओळीच्या सुरुवातीला सलग दोन समानता-चिन्हे येतात. मराठी प्रदर्शनात पहिल्या ओळीच्या शेवटचे अनावश्यक चिन्ह वगळले आहे.\end{quote}
\begin{quote}\small\textbf{OLP-0163 — स्रोतनोंद.} OLFOL-013: गोठवलेल्या इंग्रजी मजकुराच्या अस्तित्ववाची शाखेत पूर्तिची आवश्यक अट वगळली आहे. समांतर वैश्विक शाखा आणि औपचारिक व्याख्या यांच्याशी जुळणारी अट मराठीत पुनःस्थापित केली आहे.\end{quote}
\begin{quote}\small\textbf{OLP-0163 — स्रोतनोंद.} OLFOL-012: गोठवलेल्या इंग्रजी मजकुरात संबंधाचे अर्थनिर्धारण चर-मूल्यांकनावर अवलंबून असल्यासारखा अनावश्यक निर्देश जोडला आहे. रचनेतील संबंध स्थिर असल्यामुळे तो निर्देश मराठीत काढला आहे.\end{quote}
\begin{quote}\small\textbf{OLP-0163 — स्रोतनोंद.} OLFOL-014: गोठवलेल्या इंग्रजी उदाहरणात नकारानंतर दोन ठिकाणी एक जादा उघडा कंस आणि एका सूत्रात जादा स्वल्पविराम आहे. समांतर सूत्रे व व्याख्या यांनुसार मराठी आवृत्तीत ते दुरुस्त केले आहेत.\end{quote}
\begin{quote}\small\textbf{OLP-0163 — स्रोतनोंद.} OLFOL-015: गोठवलेल्या इंग्रजी उदाहरणात m = 2, 3, 4 साठी उत्तरवर्ती R(a,x) असत्य म्हटले आहे; m = 2 साठी ते सत्य आहे. सशर्त सूत्राचा पूर्ववर्ती R(x,a) या तिन्ही मूल्यांसाठी असत्य असल्यामुळे मराठीत तोच वापरला आहे.\end{quote}
\begin{quote}\small\textbf{OLP-0163 — स्रोतनोंद.} OLFOL-016: गोठवलेल्या इंग्रजी मजकुरात दुसऱ्या रंजक प्रकरणातील चलाचे m हे नाव गाळले आहे. संदर्भ आणि लगेच पुढील गणना यांनुसार ते मराठीत पुनःस्थापित केले आहे.\end{quote}
\begin{quote}\small\textbf{OLP-0163 — स्रोतनोंद.} OLFOL-017: गोठवलेल्या इंग्रजी सारांशात बाह्य वैश्विक चलाला चुकून n म्हटले आहे. सर्व चार बाह्य मूल्यांचा संदर्भ m नेच असल्यामुळे मराठीत m वापरले आहे.\end{quote}
\begin{quote}\small\textbf{OLP-0163 — स्रोतनोंद.} OLFOL-018: गोठवलेल्या इंग्रजी मजकुरात m = 1 आणि m = 2 दोन्हीसाठी n = 4 हा प्रतिदृष्टांत दिला आहे; पण (2,4) हा R मध्ये आहे. मराठीत अनुक्रमे n = 4 आणि n = 1 हे योग्य प्रतिदृष्टांत दिले आहेत.\end{quote}
\begin{quote}\small\textbf{OLP-0164 — स्रोतनोंद.} OLFOL-019: गोठवलेल्या इंग्रजी सिद्धतेतील दुसरी k-सदस्यीय क्रमित सूची t\_i पासून सुरू होते आणि त्यामुळे i मुक्त राहतो. पहिल्या सूचीशी व घटकनिहाय समानतेशी जुळण्यासाठी मराठीत ती t\_1 पासून सुरू केली आहे.\}\end{quote}
\begin{quote}\small\textbf{OLP-0164 — स्रोतनोंद.} OLFOL-020: गोठवलेल्या इंग्रजी वैश्विक बाबतीत s\_1' आणि s\_2' ही दोन्ही अपरिभाषित s पासून बनवली आहेत. पुढील पर्याय व एकमत दाव्यांना आवश्यक असल्याप्रमाणे मराठीत ती अनुक्रमे s\_1 आणि s\_2 पासून बनवली आहेत.\}\}\end{quote}
\begin{quote}\small\textbf{OLP-0164 — स्रोतनोंद.} OLFOL-021: गोठवलेल्या इंग्रजी व्याख्येच्या सुरुवातीच्या नामपदबंधात Gamma सलग दोनदा आले आहे. अर्थ न बदलणारी दुसरी पुनरावृत्ती मराठीत काढली आहे.\end{quote}
\begin{quote}\small\textbf{OLP-0165 — स्रोतनोंद.} OLFOL-022: गोठवलेल्या इंग्रजी गणनेच्या पहिल्या ओळीच्या शेवटी आणि पुढील ओळीच्या सुरुवातीला सलग दोन समानताचिन्हे दिसतात. मराठी आवृत्तीत पहिल्या ओळीचे अनावश्यक चिन्ह काढले आहे.\end{quote}
\begin{quote}\small\textbf{OLP-0165 — स्रोतनोंद.} OLFOL-023: या प्रकरणातील सूत्र-आदेशनाची व्याख्या t' हे A मध्ये x साठी आदेशनासाठी मुक्त असतानाच लागू होते. गोठवलेल्या इंग्रजी विधानात ही आवश्यक पूर्वअट नसल्यामुळे मराठीत ती स्पष्टपणे जोडली आहे.\end{quote}
\begin{quote}\small\textbf{OLP-0171 — स्रोतनोंद.} OLFOL-024: गोठवलेल्या इंग्रजीतील दुसऱ्या सूत्राच्या शेवटच्या v\_2 पुढे वस्तुभाषेचे निर्देशचिन्ह चुकले आहे. समान परिच्छेदातील सर्व इतर चर-आस्थिती आणि अभिप्रेत सूत्राशी जुळण्यासाठी मराठीत वस्तुभाषेतील v\_2 केले असून गोठवलेले इंग्रजी बाइट्स बदललेले नाहीत.\end{quote}
\begin{quote}\small\textbf{OLP-0178 — स्रोतनोंद.} OLFOL-025: बाब (6) मध्ये x चा प्रकार tau दिलेला असूनही गोठवलेल्या इंग्रजीतील पुढील स्पष्टीकरणात तो sigma म्हटला आहे. फलनाचा प्रांत आणि tau ते sigma हा प्रकार जुळण्यासाठी मराठीत tau केले असून गोठवलेले इंग्रजी बाइट्स बदललेले नाहीत.\end{quote}
\begin{quote}\small\textbf{OLP-0187 — स्रोतनोंद.} OLFOL-026 — गोठवलेल्या इंग्रजी सूत्रात M-प्राइममधील पदमूल्य काढताना f चे अर्थनिर्धारण चुकून M मधील दिले आहे. व्याख्या आणि पुढील गणना यांनुसार येथे M-प्राइम असणे आवश्यक आहे; मराठी सूत्रात ती खूण दुरुस्त केली आहे.\end{quote}
\begin{quote}\small\textbf{OLP-0189 — स्रोतनोंद.} OLFOL-027 — गोठवलेल्या इंग्रजी व्याख्येत क्रमिकेची लांबी आणि पुनरावर्तनाचा प्रतांक या दोन्हींसाठी n वापरल्यामुळे I\_0 येथे चलरहित आण्विक सूत्रांपुरते चुकून मर्यादित होते. समान क्रमिकेची लांबी k अशी स्वतंत्र खूण घेऊन मराठी व्याख्या दुरुस्त केली आहे.\end{quote}
\begin{quote}\small\textbf{OLP-0190 — स्रोतनोंद.} OLFOL-028 — गोठवलेल्या इंग्रजी पुराव्यात p च्या प्रांतात a आधीच असण्याचे आणि p रिक्त असण्याचे प्रसंग वगळले आहेत; परंतु I मध्ये रिक्त नकाशा स्पष्टपणे समाविष्ट आहे. मराठी पुराव्यात सूत्रे न बदलता हे दोन प्रसंग जोडले आहेत.\end{quote}
\begin{quote}\small\textbf{OLP-0192 — स्रोतनोंद.} OLFOL-029 — गोठवलेल्या इंग्रजीतील या अस्तित्ववाचक पुराव्याच्या विधेयात साक्षीदार x हा पहिला युक्तिवाद चुकून वगळला आहे. त्याच परिच्छेदातील सुसंगतता-विधान आणि पुढील प्रत्येक मानक संख्यांकाचे सूत्र या दोन्हींनुसार मराठी सूत्रात x पुनःस्थापित केला आहे.\end{quote}
\begin{quote}\small\textbf{OLP-0193 — स्रोतनोंद.} OLFOL-030 — गोठवलेल्या इंग्रजीतील शेवटच्या तर्कात x हा s च्या प्रांतात नसतो असे म्हटले आहे; परंतु s चा प्रांत M असल्यामुळे x त्यात असतोच. आच्छादकता अपयशी होण्यासाठी x हा s च्या व्याप्तीत नसणे आवश्यक आहे, आणि मराठी मजकुरात ते दुरुस्त केले आहे.\end{quote}
\begin{quote}\small\textbf{OLP-0194 — स्रोतनोंद.} OLFOL-031 — गोठवलेल्या इंग्रजीतील पहिल्या सूत्रात c चे अर्थनिर्धारण M मध्ये घेतले आहे; पण c हे विस्तारित भाषेचे चिन्ह असल्यामुळे L\_A-रचना M मध्ये ते परिभाषितच नाही. मराठी सूत्रात ते योग्य विस्तारित रचना M\textasciicircum{}c मध्ये घेतले आहे.\end{quote}
\begin{quote}\small\textbf{OLP-0194 — स्रोतनोंद.} OLFOL-033 — गोठवलेल्या इंग्रजी पुराव्यात मनमाना सांत उपसंच घेतल्यानंतर लगेच सर्वांत मोठा k निवडला आहे; पण त्या उपसंचात c-विषयक असमता एकही नसेल, तर असा k नसतो. मराठी पुराव्यात तो रिकामा प्रसंग स्वतंत्रपणे हाताळला आहे.\end{quote}
\begin{quote}\small\textbf{OLP-0194 — स्रोतनोंद.} OLFOL-032 — विधानात गणनीय अमानक प्रतिमानाचा दावा आहे; गोठवलेल्या इंग्रजी पुराव्यात संहततेने प्रतिमान मिळवल्यानंतर त्याची गणनीयता सिद्ध केलेली नाही. उपलब्ध पूर्वप्रमेयानुसार अधोगामी लोव्हेनहाइम--स्कोलेम प्रमेय लावून ती पायरी मराठीत स्पष्ट केली आहे.\end{quote}
\begin{quote}\small\textbf{OLP-0195 — स्रोतनोंद.} OLFOL-034 — गोठवलेल्या इंग्रजीतील प्रकरण-विभागणीत K च्या प्रांताबाहेरील b लिहिले आहे; पण लगेचची गणना y = a हेच प्रकरण हाताळते. मराठी स्पष्टीकरणात b ऐवजी a केले आहे.\end{quote}
\begin{quote}\small\textbf{OLP-0195 — स्रोतनोंद.} OLFOL-035 — गोठवलेल्या इंग्रजीतील तिसऱ्या गणनेच्या शेवटी मुक्त y आहे; त्या ओळीची डावी बाजू आणि L ची बेरीज-सारणी दोन्ही येथे a आवश्यक ठरवतात. मराठी सूत्रात y ऐवजी a केले आहे.\end{quote}
\begin{quote}\small\textbf{OLP-0196 — स्रोतनोंद.} OLFOL-036 — गोठवलेल्या इंग्रजी विधानात प्रत्येक x ला पूर्ववर्ती असल्याचा दावा आहे, पण त्याच विधानानुसार शून्य हा लघुतम घटक असून Q\_2 नुसार तो कोणाचाही उत्तरवर्ती नाही. मराठीत पूर्ववर्तीचा दावा अशून्य x पुरता मर्यादित केला आहे; पुढील पुराव्यातील पुनरुक्तीही त्यानुसार दुरुस्त केली आहे.\end{quote}
\begin{quote}\small\textbf{OLP-0196 — स्रोतनोंद.} OLFOL-037 — गोठवलेल्या इंग्रजी पुराव्यात या एकाच गणनेत आधी न परिभाषित केलेले opulus चिन्ह वापरले आहे; विभागभर बेरीजेसाठी nsplus निश्चित केलेले असल्यामुळे मराठी सूत्रांत तेच चिन्ह पुनःस्थापित केले आहे.\end{quote}
\begin{quote}\small\textbf{OLP-0196 — स्रोतनोंद.} OLFOL-038 — गोठवलेल्या इंग्रजी निष्कर्षात कोणतेही M गणनीय आहे असे गृहीत न धरता अमानक खंड गणनीय असल्याचे म्हटले आहे; अमानक PA-प्रतिमाने अगणनीयही असू शकतात. पुढील वाक्याच्या गणनीय-प्रतिमान व्याप्तीशी जुळवून मराठी दावा गणनीय प्रतिमानापुरता स्पष्ट केला आहे.\end{quote}
\begin{quote}\small\textbf{OLP-0197 — स्रोतनोंद.} OLFOL-040 — गोठवलेल्या इंग्रजी संच-वर्णनात क्रमित जोडीतील x ऐवजी असंबद्ध n वर अट आहे. K च्या घोषित संबंधाशी आणि आधीच्या उदाहरणाशी जुळवून मराठी सूत्रात x बांधला आहे.\end{quote}
\begin{quote}\small\textbf{OLP-0197 — स्रोतनोंद.} OLFOL-039 — गोठवलेल्या इंग्रजी सूत्रात n>0 साठी g(n)=n+1 दिल्यामुळे 0 आणि 1 प्रतिमेत येत नाहीत आणि g द्विएकी राहत नाही. पुढील सर्व स्थानांतरित क्रियांशी सुसंगत असा अपेक्षित नकाशा g(n)=n-1 आहे.\end{quote}
\begin{quote}\small\textbf{OLP-0200 — स्रोतनोंद.} OLFOL-041 — गोठवलेल्या इंग्रजीतील या निष्पन्नतेच्या उजव्या बाजूस आधी परिभाषित केलेल्या H ऐवजी असंबद्ध ग्रीक डेल्टा छापले आहे. पुढील वाक्यातील प्रतिपरिवर्तनासाठी आणि संपूर्ण युक्तिवादासाठी येथे H चाच निषेध आवश्यक आहे; मराठी सूत्रात H पुनःस्थापित केला आहे.\end{quote}
\begin{quote}\small\textbf{OLP-0201 — स्रोतनोंद.} OLFOL-042 — गोठवलेल्या इंग्रजीतील या स्थानांतरण-सूत्रात h ला M'\_2 मधील P चे अर्थनिर्धारण दिले आहे; पण h चा प्रांत M\_1 आहे आणि पुढील घटकनिहाय अटही M'\_1 मधील P चे अर्थनिर्धारणच स्थानांतरित करते. मराठी सूत्रात M'\_1 पुनःस्थापित केले आहे.\end{quote}
\begin{quote}\small\textbf{OLP-0202 — स्रोतनोंद.} OLFOL-043 — गोठविलेल्या इंग्रजी विधानातील “if and only if” या पदबंधात “if” गाळले आहे; मराठी विधानात आशयाला आवश्यक असलेला द्विशर्त संबंध स्पष्ट केला आहे.\end{quote}
\begin{quote}\small\textbf{OLP-0202 — स्रोतनोंद.} OLFOL-044 — गोठविलेल्या इंग्रजी सूत्राच्या शेवटच्या अणुसूत्रात P'c\_1…c\_n अशी अपूर्ण कच्ची मांडणी आहे; येथे आधीपासून वापरलेले सुबद्ध अणुसूत्र पुनर्स्थापित केले आहे.\end{quote}
\begin{quote}\small\textbf{OLP-0205 — स्रोतनोंद.} OLFOL-045 — गोठवलेल्या इंग्रजीत M' ही L शी अनुरूप असल्याचे म्हटले आहे; पण पुनर्नामकरणात भाषा L शी नव्हे, तर रचना M शी अनुरूप रचना अपेक्षित आहे. मराठी वाक्यात M पुनःस्थापित केले आहे.\end{quote}
\begin{quote}\small\textbf{OLP-0206 — स्रोतनोंद.} OLFOL-046 — गोठवलेल्या इंग्रजीत या मोठ्या रचनेसाठी आधीच्या डाव्या उपरचनेचेच M चिन्ह पुन्हा वापरले आहे आणि N-तारांकित प्रांतातील तारका कंसाबाहेर चुकीच्या पातळीवर ठेवली आहे. मराठीत मोठ्या रचनेला A म्हटले असून N च्या प्रांताची तारका M च्या समांतर केली आहे.\end{quote}
\begin{quote}\small\textbf{OLP-0206 — स्रोतनोंद.} OLFOL-047 — सापेक्षीकरणाने मिळणारे D\_1 हे सर्वसाधारणपणे अमूर्त तर्कशास्त्रातील वाक्य असते, प्रथम-क्रम वाक्य नव्हे; D\_2 मात्र प्रथम-क्रम आहे. गोठवलेल्या इंग्रजीतील दोन्हींना first-order म्हणणारा दावा मराठीत दुरुस्त केला आहे.\end{quote}
\begin{quote}\small\textbf{OLP-0206 — स्रोतनोंद.} OLFOL-048 — गोठवलेल्या इंग्रजीत गणनीय मोठ्या प्रतिमानासाठी आणि त्यातील डाव्या उपरचनेसाठी M\_0 हेच चिन्ह दोनदा वापरले आहे. मराठीत मोठे प्रतिमान A\_0 आणि उपरचना M\_0, N\_0 अशी भिन्न चिन्हे वापरली आहेत.\end{quote}
\begin{quote}\small\textbf{OLP-0207 — स्रोतनोंद.} OLFOL-049 — गोठवलेल्या इंग्रजीतील प्रमेय-विधानात normal हे आवश्यक गृहीतक राहून गेले आहे; सिद्धता समरूपता, विस्तार आणि सापेक्षीकरण यांसह सामान्यतेचे गुणधर्म स्पष्टपणे वापरते. मराठी विधानात सामान्य हे गृहीतक पुनःस्थापित केले आहे.\end{quote}
\begin{quote}\small\textbf{OLP-0207 — स्रोतनोंद.} OLFOL-050 — गोठवलेल्या इंग्रजीत या पायरीतील उपअनुक्रमाला M\_n ऐवजी M चा उपअनुक्रम म्हटले आहे. मराठीत योग्य n-अधोलिखित पुनःस्थापित केले आहे.\end{quote}
\begin{quote}\small\textbf{OLP-0207 — स्रोतनोंद.} OLNOTE-001 — सर्व M\_n चा संयोग स्पष्टपणे घडवण्यासाठी स्थिरांच्या समान अर्थनिर्धारणांबाहेर परस्पर विभक्त प्रांत असलेल्या समरूपी प्रती निवडल्या आहेत. ही निवड पुरेशी आणि सोयीची आहे; संयोगासाठी ती अनिवार्य नाही.\end{quote}
\begin{quote}\small\textbf{OLP-0207 — स्रोतनोंद.} SOL6-A019 — स्थिरांचे समानतेचे आकृतिबंध भिन्न असतील, तर समरूपतेने सर्व रचनांतील स्थिरांना त्याच वस्तू नेमता येत नाहीत. प्रथम समान आण्विक सिद्धांताचा अनंत उपअनुक्रम निवडला आहे, त्यानंतर स्थिरांच्या अर्थनिर्धारणांचे एकीकरण केले आहे. वाढत्या मूळ निर्देशांकांचा नव्या निर्देशांकांपेक्षा लहान नसण्याचा गुणधर्म आवश्यक सममूल्यतेची खोली राखतो.\end{quote}
\begin{quote}\small\textbf{OLP-0207 — स्रोतनोंद.} OLFOL-051 — गोठवलेल्या इंग्रजीत या मोठ्या संकेतबद्ध रचनेसाठी पुन्हा M हेच चिन्ह वापरले आहे, जरी M आधीच सर्व M\_n च्या संयोगासाठी वापरलेले आहे. मराठीत मोठ्या रचनेसाठी A वापरून दोन वेगळ्या वस्तू स्पष्ट केल्या आहेत.\end{quote}
\begin{quote}\small\textbf{OLP-0207 — स्रोतनोंद.} OLFOL-052 — E हे अमूर्त तर्कशास्त्रातील वाक्य आहे; म्हणून M\_n आणि N\_n वरील दोन पूर्ति-चिन्हांना L अधोलिखित आवश्यक आहे. गोठवलेल्या इंग्रजीतील साधे models चिन्ह मराठीत models\_L असे दुरुस्त केले आहे.\end{quote}
\begin{quote}\small\textbf{OLP-0207 — स्रोतनोंद.} OLFOL-053 — गोठवलेल्या इंग्रजीतील संहततेचा नेहमीचा युक्तिवाद संक्षिप्त आहे. मराठी सिद्धतेत नवीन स्थिर प्रत्येक मानक घटकाच्या वर ठेवणारी वाक्ये स्पष्ट करून अमानक साक्षीचा निष्कर्ष उलगडला आहे.\end{quote}
\begin{quote}\small\textbf{OLP-0207 — स्रोतनोंद.} OLFOL-054 — गोठवलेल्या इंग्रजीत अंतिम मोठ्या प्रतिमानासाठी M\textasciicircum{}* या दोन विसंगत TeX रूपांचा वापर आहे आणि ते आधीच्या M\textasciicircum{}* वर्धित रचनेशीही धडकते. मराठीत संकेतबद्ध मोठ्या प्रतिमानाचे A\textasciicircum{}* हेच चिन्ह सातत्याने वापरले आहे.\end{quote}
\begin{quote}\small\textbf{OLP-0220 — स्रोतनोंद.} SOL6-A029 — रिकाम्या क्रमिकेसाठी 0 आणि 1 दोन्ही संकेतांक मान्य आहेत. पहिल्या concat पुनरावर्तनात दुसरी क्रमिका रिकामी असेल तर पहिला संकेतांक जतन होतो; परिबद्ध लघुतम शोध मात्र रिकाम्या परिणामासाठी 0 निवडतो. दोन पर्यायी संकेतांक-फलनांतील हा मर्यादित भेद येथे स्पष्ट केला आहे.\end{quote}
\begin{quote}\small\textbf{OLP-0226 — स्रोतनोंद.} OLCMP-012 — गोठवलेल्या इंग्रजीत काही e आंशिक पुनरावर्ती फलनाचा निर्देशांक नसू शकतो असे म्हटले आहे. पण आधीच दिलेल्या प्रमाण रूपातील T आणि U मुळे प्रत्येक नैसर्गिक e हा एका आंशिक पुनरावर्ती फलनाचा निर्देशांक आहे. मराठीत ही अपवादस्थिती काढून सिद्धतेतील तदनुरूप प्रकरणे दुरुस्त केली आहेत.\end{quote}
\begin{quote}\small\textbf{OLP-0232 — स्रोतनोंद.} OLCMP-013 — गोठवलेल्या इंग्रजी स्पष्टीकरणात कार्यक्रमाचा निर्देशांक दोन ठिकाणी x असा चुकून आला आहे; प्रमेयात आणि त्याच स्पष्टीकरणात तो e आहे. मराठीत दोन्ही ठिकाणी e ठेवला आहे. इंग्रजीतील ``It you think'' ही अक्षरदोषयुक्त सुरुवातही ``जर'' अशी वाचली आहे.\end{quote}
\begin{quote}\small\textbf{OLP-0233 — स्रोतनोंद.} OLCMP-014 — गोठवलेल्या इंग्रजीत शेवटचे g हे केवळ ``recursive'' म्हटले आहे. f आंशिक पुनरावर्ती असल्याने g सुद्धा आंशिक पुनरावर्तीच हमखास असते; ते सर्वत्र परिभाषित असणे आवश्यक नाही. मराठीत हा गुणधर्म स्पष्ट केला आहे.\end{quote}
\begin{quote}\small\textbf{OLP-0234 — स्रोतनोंद.} OLCMP-015 — गोठवलेल्या इंग्रजीतील पहिल्या वाक्यात ``total for the partial computable functions'' असे आले आहे. मागील प्रमेय आणि पुढील विरोध पाहता येथे अभिप्रेत शब्द ``universal'' आहे; मराठीत तोच अर्थ घेतला आहे.\end{quote}
\begin{quote}\small\textbf{OLP-0236 — स्रोतनोंद.} OLCMP-016 — गोठवलेल्या इंग्रजीत S \ensuremath{\in} S नंतर X \ensuremath{\notin} S असे आले आहे, पण X परिभाषित नाही. S च्या व्याख्येवरून आवश्यक विधान S \ensuremath{\notin} S आहे; मराठीत तोच निर्देशांक ठेवला आहे.\end{quote}
\begin{quote}\small\textbf{OLP-0239 — स्रोतनोंद.} OLCMP-017 — गोठवलेल्या इंग्रजीतील व्याप्तीच्या उलट दिशेच्या युक्तिवादात x अनिर्दिष्ट आहे. येथे साक्षीदार z चा पहिला घटक (z)\_0 हा योग्य आदान आहे; मराठीत तो स्पष्ट केला आहे.\end{quote}
\begin{quote}\small\textbf{OLP-0239 — स्रोतनोंद.} OLCMP-018 — गोठवलेल्या इंग्रजीत या दिशेची सिद्धता थेट S ही संगणनक्षम फलनाची व्याप्ती आहे असे धरते; पण व्याख्येनुसार रिकामा S देखील c.e. आहे आणि तो सर्वत्र परिभाषित फलनाची व्याप्ती असू शकत नाही. रिकाम्या संचासाठी कुठेही परिभाषित नसलेल्या फलनाचा स्वतंत्र आधार मराठीत दिला आहे.\end{quote}
\begin{quote}\small\textbf{OLP-0240 — स्रोतनोंद.} SOL6-A032 — गोठवलेल्या इंग्रजीतील len(z)=2 ही कसोटी अवैध prime-power संख्यांनाही लागू होऊ शकते; 42 मध्ये 2 आणि 3 आहेत, 5 नाही पण 7 आहे. परिभाषित length helper चा सर्वांत लहान साक्षीदार index 1 आहे, म्हणून length 2 येते. जोडीची आवश्यक-संपूर्ण कसोटी z=tuple((z)\_0,(z)\_1) ठेवून K\_0 चा अचूक domain मिळतो.\end{quote}
\begin{quote}\small\textbf{OLP-0241 — स्रोतनोंद.} OLCMP-019 — गोठवलेल्या इंग्रजीतील पहिल्या दोन सिद्धता दोन्ही c.e. संचांना सर्वत्र परिभाषित प्रगणक आहेत असे धरतात. c.e. च्या व्याख्येत रिकामा संचही आहे, पण तो अशा फलनाची व्याप्ती नसतो. रिकाम्या संचाची स्वतंत्र सोपी स्थिती येथे स्पष्ट केली आहे.\end{quote}
\begin{quote}\small\textbf{OLP-0242 — स्रोतनोंद.} OLCMP-020 — गोठवलेल्या इंग्रजी सिद्धतेत A चे परिभाषाक्षेत्र cfind\_d असे ठरवूनही अंतिम कसोटीत T(e,x,h(x)) आणि cfind\_e लिहिले आहे; ती कसोटी प्रत्यक्षात A च्या पूरकाची आहे. येथे d वापरले आहे. पुढील अनौपचारिक स्पष्टीकरणातील e,f ही जोडीही d,e अशी दुरुस्त केली आहे.\end{quote}
\begin{quote}\small\textbf{OLP-0243 — स्रोतनोंद.} OLCMP-021 — गोठवलेल्या इंग्रजीत K\_0 ची पहिली व्याख्या जोडी (e,x) वापरते, पण लगेचची समतुल्य व्याख्या (x,e) अशी उलट क्रमाने लिहिते. W\_e मध्ये x असल्याचा अर्थ पहिल्या व्याख्येतील (e,x) जोडीशी जुळतो; म्हणून येथे तीच क्रमरचना ठेवली आहे.\end{quote}
\begin{quote}\small\textbf{OLP-0244 — स्रोतनोंद.} OLCMP-022 — गोठवलेल्या इंग्रजी प्रश्नात अनेक-एक न्यूनीकरण f चे प्रकारलेखन f:A\ensuremath{\to}B असे आहे; मात्र आधीची व्याख्या f:N\ensuremath{\to}N अशी आहे आणि जातिबोधक फलांचे समीकरण संपूर्ण N वर तपासायचे आहे. म्हणून येथे आधीच्या व्याख्येशी सुसंगत प्रकारलेखन वापरले आहे.\end{quote}
\begin{quote}\small\textbf{OLP-0245 — स्रोतनोंद.} OLCMP-023 — गोठवलेल्या इंग्रजीतील शेवटचे सिद्धतावाक्य K चे K\_0 कडे न्यूनीकरण सांगते. त्यामुळे K\_0 च्या संपूर्णतेवरून K ची संपूर्णता सिद्ध होत नाही; आवश्यक दिशा K\_0 चे K कडे अशी आहे. पुढील स्वतंत्र प्रश्न मात्र मूळ K चे K\_0 कडे न्यूनीकरणच विचारतो, आणि तो तसाच ठेवला आहे.\end{quote}
\begin{quote}\small\textbf{OLP-0246 — स्रोतनोंद.} SOL6-A031 — अनेक-एक न्यूनीकरण संपूर्ण N वर अट पूर्ण करणारे total computable फलन असते. गोठवलेल्या इंग्रजीतील g(w) जोडीचे घटक वेगळे काढून थेट वापरते; मात्र प्रत्येक नैसर्गिक संख्या या prime-power coding मधील वैध जोडीचा संकेतांक नाही. अचूक पुनःसंकेतबद्धीकरणाची कसोटी आणि अवैध संकेतांकासाठी nowhere-defined फलनाचा निर्देशांक येथे जोडला आहे.\end{quote}
\begin{quote}\small\textbf{OLP-0250 — स्रोतनोंद.} OLCMP-024 — गोठवलेल्या इंग्रजी प्रमेयात आंशिक संगणनक्षम f नसल्याचे म्हटले आहे, पण सिद्धतेची पहिली ओळ फक्त सर्वत्र परिभाषित संगणनक्षम f घेते. पुढील g ची रचना f(e) अपरिभाषित असतानाही आंशिक संगणनक्षम राहते; म्हणून येथे आवश्यक सर्वसाधारण आंशिक f घेतले आहे.\end{quote}
\begin{quote}\small\textbf{OLP-0254 — स्रोतनोंद.} SOL6-A040 — गोठवलेल्या इंग्रजी परिचयात शीर्ष सर्वांत डावीकडील, म्हणजे end-marker च्या घरावर सुरू होते. पुढील औपचारिक initial configuration आणि सम यंत्राची सर्व उदाहरणे आदानाच्या पहिल्या घरावर सुरू होतात. मराठी परिचय त्या operative व्याख्येशी सुसंगत केला; गोठवलेले source bytes जतन आहेत.\end{quote}
\begin{quote}\small\textbf{OLP-0254 — स्रोतनोंद.} SOL6-A037 — गोठवलेल्या इंग्रजी इतिहासातील कोलॉससचे ट्यूरिंग यांना दिलेले बांधणी-श्रेय, SSEM ला सर्वसाधारण वापराचे म्हटलेला दावा आणि सर्व आधीची यंत्रे हाताने नव्याने जोडावी लागत हा दावा दुरुस्त केला आहे. Z3 साठी पालकांचा दिवाणखाना हे अनिश्चित ठिकाण वगळले; Zuse archive मध्ये तो तपशील Z1 शी जोडला आहे. प्रत्यक्ष पुन्हा पाहिलेल्या संस्थात्मक स्रोतांचे search extractions: \href{https://www.tnmoc.org/colossus}{राष्ट्रीय संगणन संग्रहालय}, \href{https://www.manchester.ac.uk/about/news/university-celebrates-babys-65th-birthday/}{मँचेस्टर विद्यापीठ} आणि \href{https://zuse.zib.de/z1}{झुसे संग्रह}. गणितीय व्याख्यांमध्ये बदल नाही.\end{quote}
\begin{quote}\small\textbf{OLP-0255 — स्रोतनोंद.} SOL6-A038 — गोठवलेल्या इंग्रजीत कोणत्याही यंत्राच्या कोणत्याही आदानासाठी स्थितिवर्णने पाहून प्रदान निश्चित करता येते असे म्हटले आहे. हा दावा थांबणाऱ्या संगणनांपुरता स्पष्ट केला; याच विभागातील विषम आदानावरचे सम यंत्र कायम चालते आणि त्याला अंतिम प्रदान नसते.\end{quote}
\begin{quote}\small\textbf{OLP-0255 — स्रोतनोंद.} OLTUR-001 — गोठवलेल्या इंग्रजीत आरंभीची अवस्था one म्हटली आहे; लगेच आधीचे चित्र आणि संक्रमण फलन शून्य अवस्था दाखवतात. मराठीत त्यांच्याशी सुसंगत शून्य अवस्था वापरली आहे.\end{quote}
\begin{quote}\small\textbf{OLP-0257 — स्रोतनोंद.} OLTUR-002 — या औपचारिक प्रारंभिक स्थितिवर्णनात शीर्ष आदानाच्या पहिल्या घरावर आहे; गोठवलेल्या इंग्रजी परिचयात ते फितीच्या डाव्या टोकाच्या घरावर असल्याचे म्हटले आहे. दोन्ही मूळ कथने गोठवलेल्या स्रोतांत जतन आहेत; मराठी परिचय आणि येथे दिलेली व्याख्या एकाच प्रारंभिक स्थानाशी सुसंगत केली आहेत.\end{quote}
\begin{quote}\small\textbf{OLP-0257 — स्रोतनोंद.} OLTUR-004 — गोठवलेली व्याख्या रिकाम्या आदानालाही परवानगी देते; पण दाखवलेल्या त्रयीत शीर्षाचा क्रमांक एक आणि फितीच्या क्रमिकेची लांबीही एक असल्याने आधीची काटेकोर लहानपणाची अट मोडते. मराठीत रिकामे घर समाविष्ट करण्याची मर्यादित स्पष्टता दिली आहे; सूत्र जतन आहे.\end{quote}
\begin{quote}\small\textbf{OLP-0257 — स्रोतनोंद.} OLTUR-003 — गोठवलेल्या इंग्रजीत या स्पष्टीकरणात आदान डाव्या टोकाच्या चिन्हाच्या डावीकडे सुरू होते असे शब्दशः येते; लगेच आधीचे सूत्र आणि परिच्छेद उजवी बाजू निश्चित करतात. मराठीत उजवीकडे असे स्पष्ट केले आहे.\end{quote}
\begin{quote}\small\textbf{OLP-0257 — स्रोतनोंद.} SOL6-A039 — सांत धावेतील अंतिम स्थितिवर्णनाला क्रमिकेत पुढचे स्थितिवर्णन नसते; अनंत धावेतील प्रत्येकाला असते. गोठवलेल्या व्याख्येतील प्रत्येक या शब्दाचा scope स्पष्ट केला. डावीकडील हालचालीच्या clause मधील अन्यथा फक्त डावीकडील हालचालीला लागू होते; उजवीकडील सूचनेला डाव्या टोकावर थांबवणारा पर्यायी successor मिळत नाही. रिकाम्या आदानासाठीचे आधीचे blank-padding स्पष्टीकरण कायम ठेवून मुख्य displayed initial formula फक्त nonempty input साठी म्हटले आहे.\end{quote}
\begin{quote}\small\textbf{OLP-0257 — स्रोतनोंद.} OLTUR-005 — आधीच्या विभागात डाव्या टोकाचे चिन्ह बदलून लिहिण्यास परवानगी आहे; गोठवलेल्या प्रदान-सूत्रात अंतिम फीत मात्र त्या चिन्हानेच सुरू होत असल्याचे गृहीत आहे. मराठीत या मर्यादेची स्पष्ट नोंद केली आहे.\end{quote}
\begin{quote}\small\textbf{OLP-0258 — स्रोतनोंद.} SOL6-A041 — गोठवलेल्या addition diagram मध्ये प्रारंभिक अवस्थेत रेघ वाचणारा loop दुसऱ्या अवस्थेकडे जातो. त्यामुळे मधला blank भरला जात नाही; उदाहरणार्थ तीन आणि दोन या आदानांचे output सलग पाच रेघा नसते. हा loop प्रारंभिक अवस्थेकडे परत नेला आहे; separator भरल्यावरच पुढच्या अवस्थेत जाते.\end{quote}
\begin{quote}\small\textbf{OLP-0258 — स्रोतनोंद.} SOL6-A041 — या mover exercise मधील गोठवलेला संदर्भ disciplined doubler कडे होता. आधीच्या परिच्छेदातील रचना मूळ सहा-अवस्थांच्या doubler ला mover जोडते; संदर्भ त्याच यंत्राकडे दुरुस्त केला.\end{quote}
\begin{quote}\small\textbf{OLP-0260 — स्रोतनोंद.} OLTUR-006 — गोठवलेल्या मजकुरात शीर्ष डाव्या टोकाचे चिन्ह सापडेपर्यंत डावीकडे नेण्याची पद्धत दिली आहे. पण आधीची यंत्र-व्याख्या ते चिन्ह पुसण्यास किंवा इतर घरांत लिहिण्यास मनाई करत नाही. त्यामुळे हा केवळ मर्यादित उदाहरणार्थ उपाय आहे; सर्वसाधारण शिस्तबद्ध रूपांतरासाठी अधिक अचूक चिन्ह-व्यवस्था आवश्यक आहे.\end{quote}
\begin{quote}\small\textbf{OLP-0260 — स्रोतनोंद.} SOL6-A042 — गोठवलेल्या disciplined addition diagram मध्येही प्रारंभिक stroke-loop चुकीने दुसऱ्या अवस्थेत जात होता. पहिला input block पार करताना प्रारंभिक अवस्थेतच राहणारा loop दिला; separator भरल्यावर पुढची अवस्था घेतली जाते. अंतिम head घर एकवर आणि designated h अवस्थेत पोहोचण्याची रचना कायम.\end{quote}
\begin{quote}\small\textbf{OLP-0261 — स्रोतनोंद.} OLTUR-007 — गोठवलेल्या इंग्रजीतील संक्रमण फलनाच्या पहिल्या शाखेत फक्त अवस्था Q मध्ये असणे ही अट आहे; त्यामुळे तिसऱ्या, अपरिभाषित-संक्रमण शाखेशी ती छेदते. येथे पहिली शाखा \ensuremath{\delta} परिभाषित असतानाच लागू होईल असे स्पष्ट केले आहे. सर्व चिन्हे आणि इतर शाखा जतन आहेत.\end{quote}
\begin{quote}\small\textbf{OLP-0261 — स्रोतनोंद.} SOL6-A043 — या विभागातील तिन्ही addition diagrams मध्ये पहिला stroke-loop चुकीने दुसऱ्या अवस्थेकडे होता; पहिला input block पार करताना प्रारंभिक अवस्थेत राहणारा loop केला. अंतिम combined machine रेघांचा एक दुप्पट सलग खंड ठेवते, पण त्याच्या डावीकडे blanks उरतात. म्हणून ते रेघांचा खंड दुप्पट करते एवढेच उदाहरण सांगते; unary numerical-output व्याख्येनुसार फलन संगणित करण्यासाठी आधीच्या mover प्रमाणे output फितीच्या आरंभी आणावे लागेल.\end{quote}
\begin{quote}\small\textbf{OLP-0262 — स्रोतनोंद.} SOL6-A044 — गोठवलेल्या boundary-simulation footnote मध्ये दुसऱ्या घरात end-marker लिहिण्याचे प्रकरण आहे; घर शून्यवरील marker पुसण्याला मात्र मूळ machine definition परवानगी देते. त्या स्थानाची ओळख जतन करणारी व्यवस्था लागते हे स्पष्ट केले. अनेक tapes ची संख्या प्रत्येक machine साठी सांत आहे; अनंत grid चे वेगळे उदाहरण कायम.\end{quote}
\begin{quote}\small\textbf{OLP-0265 — स्रोतनोंद.} SOL6-A045 — सांत संख्येच्या आदानांवरच परिभाषित असलेल्या फलनाचा दावा आंशिक संगणनक्षमतेचा आहे. त्या सांत सारणीतील मूल्ये दिली असता table lookup आणि बाहेरच्या inputs वर न थांबणे याने ते फलन संगणित करता येते; प्रत्येक नैसर्गिक आदानावर परिभाषित असल्याचा दावा नाही.\end{quote}
\begin{quote}\small\textbf{OLP-0267 — स्रोतनोंद.} SOL6-A048 — गोठवलेली यादी पहिल्या क्रमांकापासून लिहिलेली आहे, पण universal function चे निर्देशांक-परिभाषाक्षेत्र सर्व नैसर्गिक संख्या आहे आणि त्यात शून्य येते. येथे एकाच प्रभावी प्रगणनाला शून्यापासून निर्देशांक दिले; प्रत्येक नैसर्गिक निर्देशांकाला यंत्र मिळते. एकापासून निर्देशांक देणे स्वतः चुकीचे नाही, पण त्या पद्धतीत शून्याचा स्वतंत्र विस्तार सांगावा लागला असता.\end{quote}
\begin{quote}\small\textbf{OLP-0267 — स्रोतनोंद.} SOL6-A046 — गोठवलेल्या universal simulation sketch मध्ये zero-based tape-block numbering, रिकामे input padding, शून्य head वर left-stay आणि उजवीकडील tape extension स्पष्ट नव्हते. ते operative configuration नियमांशी जुळवले. Output cleanup मध्ये marker आणि trailing blanks वगळून शुद्ध unary output तपासणे आवश्यक आहे; प्रत्येक non-stroke block वर लगेच rejection केल्यास सर्व वैध initial markers आणि अनेक halting tapes चुकीने नाकारले जातात. अमान्य output असल्यास simulator थांबतो पण nonnumeric output ठेवतो, म्हणून halting equivalence आणि partial-function undefinedness दोन्ही जपली जातात. प्रत्यक्ष पूर्ण universal instruction table बांधल्याचा दावा नाही.\end{quote}
\begin{quote}\small\textbf{OLP-0268 — स्रोतनोंद.} SOL6-A049 — universal machine विभागात निश्चित केलेले तेच शून्यापासूनचे प्रभावी प्रगणन येथे वापरले आहे. गोठवलेल्या एकापासूनच्या यादीत शून्य निर्देशांकाचे यंत्र सांगितलेले नव्हते; या मांडणीत प्रत्येक नैसर्गिक निर्देशांकाला यंत्र मिळते.\end{quote}
\begin{quote}\small\textbf{OLP-0268 — स्रोतनोंद.} OLTUR-008 — गोठवलेल्या इंग्रजी पुराव्यात पहिले घर पाहत थांबणे म्हणजेच शिस्तबद्ध असणे असे कंसात म्हटले आहे. आधीच्या व्याख्येनुसार शिस्तबद्ध यंत्राला याशिवाय खास थांबण्याची अवस्था, टोकाचे चिन्ह राखणे आणि घर ० वरून डावीकडे न जाणे या अटीही आहेत. येथे आधीच्या रूपांतरानुसार पूर्ण शिस्तबद्ध यंत्र गृहीत धरले आहे; त्यातील पहिल्या घरावर थांबण्याचा गुणधर्म या जोडणीसाठी वापरला जातो.\end{quote}
\begin{quote}\small\textbf{OLP-0270 — स्रोतनोंद.} OLTUR-010 — गोठवलेल्या इंग्रजी मजकुरात A(x,y) बनवणाऱ्या अटींना sentences म्हटले आहे. प्रदर्शित अटींमध्ये x आणि y मुक्त आहेत; म्हणून त्या स्वतंत्र वाक्ये नसून उघडी सूत्रे आहेत. पुढील संक्रमण-स्वयंसिद्धकांत त्या सार्वत्रिक संख्यापकांच्या व्याप्तीत वापरल्या जातात.\end{quote}
\begin{quote}\small\textbf{OLP-0270 — स्रोतनोंद.} OLTUR-009 — गोठवलेल्या इंग्रजी डावीकडील संक्रमणाच्या पहिल्या फलितात A(x,y) आहे. पण त्या संक्रमणात चिन्ह लिहिले जाते x' या घरात; A च्या आधीच्या व्याख्येनुसार त्याने लिहिले जाणारे घर वगळले पाहिजे. येथे A(x',y) केले आहे. दुसऱ्या, घर शून्यवरील शाखेतील A(0,y) बदललेले नाही.\end{quote}
\begin{quote}\small\textbf{OLP-0271 — स्रोतनोंद.} OLTUR-013 — गोठवलेल्या इंग्रजी आढाव्यात एका निष्पन्नतेत T(M,w) ऐवजी आवश्यक {!}T(M,w) मधील {!} गाळलेले आहे; पुढील उदाहरणात संपूर्ण बांधणी M साठी असताना अचानक T यंत्र म्हटले आहे. मराठीत दोन्ही ठिकाणी सुसंगत {!}T(M,w) आणि M वापरले आहेत.\end{quote}
\begin{quote}\small\textbf{OLP-0271 — स्रोतनोंद.} SOL6-A051 — गोठवलेल्या इंग्रजी व्याख्येत आरंभीचे सर्वांत उजवे रिकामे नसलेले घर घेतले आहे. आदानाच्या शेवटी रिकामी चिन्हे असतील, तर ती मर्यादा आदानाच्या पूर्ण लांबीपेक्षा लहान होऊ शकते. दिलेली क्रम-स्वयंसिद्धके सर्व रचनांत लागोपाठच्या संख्यांकांमध्ये इतर घटक नसल्याची खात्री देत नाहीत; त्यामुळे लहान मर्यादेपलीकडील सार्वत्रिक रिकामेपण आरंभीच्या अटींवरून मिळत नाही. संपूर्ण आदानाची लांबी घेऊन मूळ विगमनाचा तर्क राखला आहे.\end{quote}
\begin{quote}\small\textbf{OLP-0271 — स्रोतनोंद.} OLTUR-012 — गोठवलेल्या इंग्रजी पुराव्यात आधी थांबणारी जोडी (q,\ensuremath{\sigma}) निवडल्यानंतर या निष्पन्नतेत अचानक (q',\ensuremath{\sigma}') येते, आणि एका S-विधेयापूर्वी Obj चिन्हही गाळले आहे. येथे आधी निवडलेलीच (q,\ensuremath{\sigma}) जोडी व भाषेतील तेच S-विधेय वापरले आहे; अस्तित्वविधानाचा निष्कर्ष बदलत नाही.\end{quote}
\begin{quote}\small\textbf{OLP-0271 — स्रोतनोंद.} OLTUR-011 — गोठवलेल्या इंग्रजी लेम्मात यंत्र n पायऱ्यांनंतर थांबलेले नसावे अशी अट आहे, पण पुढील पुरावा n व्या पायरीवर मिळणाऱ्या थांबणाऱ्या स्थितिवर्णनासाठीही तेच लेम्मा वापरतो. येथे चालक्रम n-पायऱ्यांच्या स्थितिवर्णनापर्यंत पोहोचणे ही अचूक अट घेतली आहे; ते स्थितिवर्णन थांबणारे असू शकते.\end{quote}
\begin{quote}\small\textbf{OLP-0271 — स्रोतनोंद.} OLTUR-014 — गोठवलेल्या इंग्रजी पडताळणीत डावीकडील संक्रमणाच्या पहिल्या शाखेत A(x,y) आणि त्याच्या उदाहरणात A(l,n) पुन्हा आले आहे. OLTUR-009 नुसार लिहिलेले घर x'=l'=m वगळले पाहिजे. म्हणून येथे A(x',y) आणि A(l',n) वापरले आहेत; घर शून्यवरील शाखा बदललेली नाही.\end{quote}
\begin{quote}\small\textbf{OLP-0271 — स्रोतनोंद.} OLTUR-015 — गोठवलेल्या इंग्रजी उजवीकडील व डावीकडील पुढील-स्थिती प्रदर्शनांत प्रत्येक घरासाठी जुने \ensuremath{\sigma}\_i चिन्हच पुन्हा दाखवले आहे, जरी आधीचे वाक्य लिहिलेल्या घरात \ensuremath{\sigma}' येते असे म्हणते आणि पुढील स्पष्टीकरण जुने त्या घराचे विधान वगळते. येथे वरची + खूण पुढील-स्थिती चिन्हांसाठी आहे: लिहिलेल्या घरात नवे चिन्ह आणि इतरत्र जुने चिन्ह. त्यामुळे प्रदर्शनाचे वाचन स्थितिवर्णनाच्या व्याख्येशी जुळते.\end{quote}
\begin{quote}\small\textbf{OLP-0272 — स्रोतनोंद.} SOL6-A052 — गोठवलेल्या इंग्रजी बांधणीत E सर्वांत डाव्या फीतघरावर थांबते म्हटले आहे, पण ते घर अंतखूणेचे आहे; पुढे चालवायच्या D चे प्रारंभिक शीर्षस्थान पहिले आदानघर आहे. येथे E ने तयार केलेले वाक्य त्या घरापासून ठेवून त्याच घरावर थांबण्याची अट घेतली आहे.\end{quote}
\begin{quote}\small\textbf{OLP-0272 — स्रोतनोंद.} OLTUR-016 — गोठवलेल्या इंग्रजी पुराव्यात आदान-वाक्य प्रथम B असे लिहिले आहे, पण त्याच परिच्छेदात पुढे {!}B सरकवले जाते आणि संपूर्ण प्रमेयात {!}B हेच वाक्य आहे. येथे प्रारंभीही {!}B लिहिले आहे; नकार आणि पूर्ततायोग्यतेची युक्ती अपरिवर्तित आहे.\end{quote}
\begin{quote}\small\textbf{OLP-0273 — स्रोतनोंद.} OLTUR-017 — गोठवलेल्या इंग्रजी सांत प्रतिरूपात वरची मर्यादा थांबण्याच्या पायऱ्या k आणि आदानाची लांबी यांतील मोठी संख्या घेतली आहे. यंत्र घर एकपासून सुरू होऊन k उजव्या हालचालींनंतर k+1 या घरात पोहोचू शकते. तसेच सर्वोच्च घटकावर n<n ठेवलेल्या सांत क्रमामुळे n हे शेवटचे आदान-चिन्ह असलेले घर नसावे; अन्यथा रिकाम्या-शेपटीची अट त्या घरालाच रिकामे म्हणेल. म्हणून मराठीत दोन्ही संख्यांपेक्षा एकाने मोठी मर्यादा घेतली आहे.\end{quote}
\begin{quote}\small\textbf{OLP-0273 — स्रोतनोंद.} OLTUR-018 — गोठवलेल्या इंग्रजी सरावातील एकमेव सूचना पुढील अवस्थेला अव्याख्यायित q म्हणते, जरी एकच अवस्था q0 दिली आहे; तसेच M'' च्या उत्तरवर्ती फलनात एका ठिकाणी M' लिहिले आहे. मराठीत दोन्ही ठिकाणी दिलेलीच एक-अवस्थेची यंत्रणा आणि M'' रचना वापरली आहे.\end{quote}
\begin{quote}\small\textbf{OLP-0273 — स्रोतनोंद.} OLTUR-019 — गोठवलेल्या इंग्रजी सांत-प्रतिरूप बांधणीत डावीकडे जाण्याच्या सकारात्मक-स्थान शाखेत स्थिरता-अट लिहिलेल्या घराऐवजी नव्या शीर्षस्थानाला वगळते; त्याच शाखेत नव्या वेळेसाठी B अटही गाळली आहे, जरी लगतचा मजकूर प्रत्येक संक्रमणात ती जोडल्याचे सांगतो. येथे आधीच्या OLTUR-009 नुसार लिहिलेले घर वगळले आहे आणि या शाखेतही नव्या वेळेची अट घातली आहे. शून्यावरील शाखा बदललेली नाही.\end{quote}
\begin{quote}\small\textbf{OLP-0273 — स्रोतनोंद.} OLTUR-020 — गोठवलेल्या इंग्रजी लेम्माच्या पुराव्यात आणि लगेचच्या सरावात थांबण्याचे वाक्य E असे लिहिले आहे; सरावात T' ऐवजी जुने T आले आहे. येथे दोन्ही ठिकाणी लेम्मातील संपूर्ण T' आणि {!}E संयुक्त विधानच तपासायचे आहे.\end{quote}
\begin{quote}\small\textbf{OLP-0273 — स्रोतनोंद.} OLTUR-021 — गोठवलेल्या इंग्रजी प्रतिपरिवर्तनात गृहीत प्रतिरूप अचानक जुन्या T आणि {!}E संयुक्त विधानाचे म्हटले आहे, जरी लेम्मा आणि पुढील निष्पन्नता T' वर अवलंबून आहेत. येथे लेम्मातीलच T' आणि {!}E गृहीत घेतले आहेत.\end{quote}
\begin{quote}\small\textbf{OLP-0273 — स्रोतनोंद.} OLTUR-022 — गोठवलेल्या इंग्रजी मजकुरात आरंभीच्या वेळेसाठीही B(0) निष्पन्न असल्याचा दावा आहे. पण B(y') ही अट संक्रमणांच्या निष्कर्षातच आहे; आरंभीच्या वेळेसाठी ती दिलेली नाही. येथे B(n) फक्त किमान एक पायरी झाल्यावर वापरले आहे. परस्परभिन्न संख्यांक मिळवण्यासाठी हेच पुरेसे आहे.\end{quote}
\begin{quote}\small\textbf{OLP-0273 — स्रोतनोंद.} SOL6-A053 — गोठवलेल्या इंग्रजी उपप्रमेयातील निष्पत्तीपद्धतीसाठी पुढील पुरावा निष्पत्तींचे प्रभावी प्रगणन वापरतो. ती आवश्यक अट येथे स्पष्ट केली आहे; केवळ सांत-वैध वाक्यांचा अप्रभावी स्वयंसिद्धकसंच घेतल्याने विरोधाभास मिळत नाही. सर्व सांत रचनांतील सत्यता हा सममूल्यतेचा पूर्ण उजवा भाग आहे.\end{quote}
\begin{quote}\small\textbf{OLP-0276 — स्रोतनोंद.} SOL6-A054 — गोठवलेल्या इंग्रजी परिच्छेदात फ्रेगेचे logicist उद्दिष्ट सर्व गणितासाठी सांगितले आहे. त्याच्या 1884 च्या Grundlagen मध्ये अंकगणित विश्लेषक, तर भूमिती अनुभवपूर्व संश्लेषक मानली आहे; येथे हा फरक जपला.\end{quote}
\begin{quote}\small\textbf{OLP-0276 — स्रोतनोंद.} SOL6-A054 — 16 जून 1902 चे रसेलचे फ्रेगेला पत्र McMaster Russell Archives मध्ये नोंदलेले आहे. येथे 1902 हा विरोधाभास कळवण्याचा वर्षसंदर्भ आहे; शोध लागल्याची तारीख म्हणून तो दिलेला नाही.\end{quote}
\begin{quote}\small\textbf{OLP-0276 — स्रोतनोंद.} SOL6-A054 — मूळ ऐतिहासिक सारांशातील सर्व प्रणालींबद्दलचा दावा येथे आवश्यक प्रभावीपणा व अंकगणितीय सामर्थ्य यांच्या अटींसह दिला आहे. केवळ सुसंगततेवरून स्वतःच्या सुसंगतता-विधानाच्या खंडनाची अशक्यता मिळत नाही.\end{quote}
\begin{quote}\small\textbf{OLP-0277 — स्रोतनोंद.} SOL6-A055 — मूळ परिच्छेदातील प्रतिनिधित्वाची अट पहिल्या अपूर्णता प्रमेयासाठी पुरेशी सामर्थ्याची अट आहे. दुसऱ्या प्रमेयासाठी अतिरिक्त अटी पुढील विभागात दिल्या आहेत. सुसंगतता व प्रभावी स्वयंसिद्धकीय मांडणी या इतर आवश्यक अटीही निष्कर्षात स्पष्ट केल्या.\end{quote}
\begin{quote}\small\textbf{OLP-0278 — स्रोतनोंद.} SOL6-A056 — गोठवलेल्या इंग्रजी आढाव्यात मूळ Gödel sentence आणि केवळ consistency गृहीत धरणाऱ्या व्यापक प्रमेयाची गल्लत होते. मूळ वाक्याच्या स्वतंत्रतेसाठी अतिरिक्त गृहीतक आणि Rosser च्या बदललेल्या रचनेचा फरक येथे स्पष्ट केला.\end{quote}
\begin{quote}\small\textbf{OLP-0283 — स्रोतनोंद.} SOL6-A063 — मर्यादा सिद्ध करताना गरज नसलेली पदे वगळून केवळ अंतिम पदाच्या उपपदांची रचनाक्रमिका निवडली आहे. मनमानी रचनाक्रमिकेत इतर मोठी पदे असू शकतात; त्यांच्यासाठी ही लांबीची मर्यादा दावा केलेली नाही.\end{quote}
\begin{quote}\small\textbf{OLP-0284 — स्रोतनोंद.} SOL6-A064 — मूळ आण्विक-सूत्र चाचणीतील argument-code tuple साठी z<x ही मर्यादा चुकीची आहे. एका वैध द्विस्थानी विधेयाला दोन स्थिरपदे दिल्यास हा tuple code अंतिम formula code पेक्षा मोठा असू शकतो. योग्य आदिम पुनरावर्ती मर्यादा आणि tuple length ही predicate arity इतकीच असण्याची अट येथे स्पष्ट केली.\end{quote}
\begin{quote}\small\textbf{OLP-0284 — स्रोतनोंद.} SOL6-A064 — मूळ sentence चाचणीत सूत्र असण्याची Frm(x) अट सुटली आहे. तिच्याशिवाय रिक्त अनुक्रमाच्या code वर दोन्ही bounded universal अटी रिक्तपणे खऱ्या ठरतात. येथे आवश्यक सूत्र-अट जोडली.\end{quote}
\begin{quote}\small\textbf{OLP-0286 — स्रोतनोंद.} SOL6-A065 — मूळ proof-code चाचण्या decoded fields वरच विसंबतात. येथे sequence codes, नेमका दोन-बाजूंचा sequent code, नेमकी initial/unary/binary proof-node रूपे आणि दाखवलेल्या दोन्ही नियमांचे syntax domains तपासले. लांबी दोन असली तरी extra prime factor असलेल्या 42 चा sequent code म्हणून स्वीकार किंवा initial node ला तिसरा अनावश्यक घटक जोडण्याचा स्वीकार यामुळे थांबतो.\end{quote}
\begin{quote}\small\textbf{OLP-0287 — स्रोतनोंद.} SOL6-A066 — मूळ OpenAssum च्या मार्गसंकेतांकासाठी पूर्ण subtree क्रमिकेचा code हा strict upper bound घेतला आहे. फक्त एक गृहीतक असलेल्या निष्पत्तीसाठी आवश्यक singleton मार्गाचा code त्या क्रमिकेच्या code इतकाच असतो; म्हणून मूळ कसोटी त्या गृहीतकाला उघडे मानत नाही. येथे वैध nonempty मार्गाला सर्वसाधारण sequence bound दिला. Exact node guard ने extra fields वगळले, आणि Assum ने प्रत्यक्ष subtree क्रमिकेच्या पूर्ण लांबीपर्यंत शोध घेतला.\end{quote}
\begin{quote}\small\textbf{OLP-0288 — स्रोतनोंद.} SOL6-A067 — decoded ओळी व लांबी तपासणे म्हणजे canonical क्रमिकासंकेतांक तपासणे नव्हे. एका axiom च्या singleton proof-code ला त्याच्या support बाहेरील अविभाज्य घटकाने गुणले तरी जुन्या length आणि line projections बदलत नाहीत; जुनी Deriv कसोटी ते अवैध संख्या स्वीकारते. येथे Deriv(d) आणि antecedent-list s साठी मागील Seq अट जोडली.\end{quote}
\begin{quote}\small\textbf{OLP-0292 — स्रोतनोंद.} SOL6-A070 — समशेषतेचा समान शेष हा अर्थ धन भाजकासाठी आहे. शून्य भाजकाचा totalized शेष शून्य घेतला तरी तो divisibility द्वारे दिलेल्या समशेषतेशी समतुल्य नसतो. म्हणून या व्याख्येत आणि सुनझीच्या प्रमेयात धन भाजक स्पष्ट केला. बीटा फलाच्या बांधणीत सर्व भाजक आधीपासून धन आहेत.\end{quote}
\begin{quote}\small\textbf{OLP-0326 — स्रोतनोंद.} SOL6-A091: पर्याय मूल्यांकन निश्चित x, X किंवा u वर बदलते; त्याचे कोणत्याही y वरील मूल्य डाव्या बाजूला दाखवले आहे. पूरक संचासाठी disjointness बरोबर पूर्ण आधारक्षेत्र झाकणे आवश्यक आहे. Conjunction encoding मधील X मुक्त नसण्याची अट capture टाळते.\end{quote}
\begin{quote}\small\textbf{OLP-0329 — स्रोतनोंद.} SOL6-A092: Inf डेडेकिंड अनंतता व्यक्त करते. येथे सर्वसाधारण अनंततेशी आणि Fin च्या सांततेशी सममूल्यत्वासाठी निवडीचे स्वयंसिद्धक मानणारा सामान्य संचसैद्धान्तिक संदर्भ घेतला आहे. Orbit ची पुनरावृत्ती काढल्यावर repetition-free प्रगणना मिळते.\end{quote}
\begin{quote}\small\textbf{OLP-0331 — स्रोतनोंद.} SOL6-A093: संगणनक्षम प्रगणनाचे विधान प्रभावी भाषांसाठी आणि प्रभावी सिद्धता प्रणालींसाठी आहे. द्वितीय क्रमाच्या मानक चिन्हार्थमीमांसेत निर्दोष व संपूर्ण प्रभावी प्रणाली नसते. प्रगणनीय प्रतिमानाचा लोव्हेनहाइम स्कोलेम दावा प्रगणनीय भाषेसाठी आहे.\end{quote}
\begin{quote}\small\textbf{OLP-0332 — स्रोतनोंद.} SOL6-A094: विगमन रूपबंधातील इतर मुक्त parameters सार्विकरीत्या बद्ध घेतले आहेत. Order proof मध्ये tail चा अंतर्भाव ही त्या सूत्राच्या पूर्ततेसाठी आवश्यक अट आहे; ती पुरेशी होण्यासाठी successor closure हवी. Tail स्वतः वरील सूत्र पूर्ण करतो, म्हणून ते पूर्ण करणाऱ्या सर्व संचांचा छेद योग्य order relation देतो.\end{quote}
\begin{quote}\small\textbf{OLP-0333 — स्रोतनोंद.} SOL6-A095: गैरस्वयंसिद्धकीकरणाचा निष्कर्ष प्रभावी सिद्धता प्रणालींविषयी आहे. Reduction मध्ये प्रथम क्रमाची अंकगणितीय वाक्ये वापरली असून नव्या चलांची निवड capture टाळणारी आहे.\end{quote}
\begin{quote}\small\textbf{OLP-0338 — स्रोतनोंद.} SOL6-A096: Injectivity फक्त तुलना केल्या जाणाऱ्या X वरील प्रतिबंधासाठी आहे. फलन पूर्ण प्रांतावर परिभाषित असले तरी X बाहेरील मूल्ये मनमानी असू शकतात; त्यांवर injectivity मागितलेली नाही.\end{quote}
\begin{quote}\small\textbf{OLP-0339 — स्रोतनोंद.} SOL6-A097: मूळ Inf, Count आणि Aleph-one सूत्रांचे आधी निदर्शित दोष येथे प्रत्यक्ष दुरुस्त केले आहेत. X-to-X closure, restricted injectivity, empty-set countability आणि strictly smaller cardinality च्या उपसंचांची अट आवश्यक आहेत.\end{quote}
\begin{quote}\small\textbf{OLP-0340 — स्रोतनोंद.} SOL6-A098: दुरुस्त cardinality definitions नंतर powerset coding आणि unique-code Cont वापरले आहेत. पूर्ण प्रांत continuum इतकाच असण्याची अट आता full-domain Y आणि Cont(Y) मागते; जुन्या अपुऱ्या identity-function range अटीऐवजी ही अट आहे. CH आणि NCH चे dependent explanations त्यानुसार दुरुस्त केले आहेत.\end{quote}
\begin{quote}\small\textbf{OLP-0344 — स्रोतनोंद.} SOL6-A102: फक्त कोणत्याही लॅम्डाच्या व्याप्तीत येणे हा बद्ध वापराचा निकष नाही. त्याच नावाचे संबंधित अमूर्तीकरण त्या वापराला बांधते. नामांतर करताना मुक्त वापर आणि इतर बंधन-संबंध जपले आहेत.\end{quote}
\begin{quote}\small\textbf{OLP-0345 — स्रोतनोंद.} SOL6-A103: आदेशनात चराचे फक्त मुक्त वापर बदलतात. आदेशित पदातील मुक्त चर पकडले जाऊ नयेत म्हणून आवश्यक बद्ध चरांचे आधी नामांतर केले जाते.\end{quote}
\begin{quote}\small\textbf{OLP-0347 — स्रोतनोंद.} SOL6-A105: या करीकरणाच्या न्यूनीकरणांत अमूर्त केलेली चरे परस्पर भिन्न आणि सर्व आदेशित पदांतील मुक्त चरांपासून वेगळी निवडली आहेत. आवश्यक नामांतर आधी केले जाते; अन्यथा displayed intermediate expression मुक्त चर पकडू शकते.\end{quote}
\begin{quote}\small\textbf{OLP-0350 — स्रोतनोंद.} SOL6-A111: क्लिनी प्रमाण रूपाने total primitive-recursive functions, total predicate चे लघुतमीकरण आणि search ला force करणारे शेवटचे उपयोजन पुरते आहे. या सिद्धतेत naive partial composition किंवा partial primitive-recursion construction गृहीत धरलेली नाही.\end{quote}
\begin{quote}\small\textbf{OLP-0351 — स्रोतनोंद.} SOL6-A108: आधीच्या recursive-function व्याख्येतील शून्य फलन एकस्थानी आहे. चर्चचा शून्य अंक आणि तो अंक नेहमी परत करणारे एकस्थानी फलन यांत अतिरिक्त अमूर्तीकरणाचा फरक आहे.\end{quote}
\begin{quote}\small\textbf{OLP-0352 — स्रोतनोंद.} SOL6-A109: सरळ nested application ची construction total numeric functions साठी आहे. आंशिक strict composition चे सामान्य विधान पुढील स्वतंत्र computability equivalence पूर्ण झाल्यावर मिळते; main proof येथे total closure पुरती वापरते.\end{quote}
\begin{quote}\small\textbf{OLP-0353 — स्रोतनोंद.} SOL6-A110: प्राचलांसहित प्रतिनिधी आणि प्राचले सामावलेली फलनमूल्ये यांची नावे वेगळी घेतली आहेत. मागील recursion value ला फक्त उरलेली प्राचले उपयोजली आहेत; recursion index पुन्हा उपयोजलेला नाही. Numeric total प्रकरण आणि arbitrary beta-recursion equations यांच्या सिद्धतांचे कार्य वेगळे स्पष्ट केले आहे.\end{quote}
\begin{quote}\small\textbf{OLP-0354 — स्रोतनोंद.} SOL6-A113: दिलेल्या पदातील मुक्त चर पकडले जाऊ नयेत म्हणून diagonal construction मधील अमूर्त चर आधी ताजे निवडले आहे. सर्व fixed-point expressions जपली आहेत.\end{quote}
\begin{quote}\small\textbf{OLP-0355 — स्रोतनोंद.} SOL6-A112: Displayed selector-search construction प्रथम total numeric test साठी सिद्ध केली आहे. Search थांबत नसल्यास त्याचे head computationही थांबत नाही. क्लिनी प्रमाण रूपातील अंतिम extractor आधी search force केल्याने general computability theorem पूर्ण होते; arbitrary partial test चे closure त्यानंतर घेतले आहे.\end{quote}
\begin{quote}\small\textbf{OLP-0362 — स्रोतनोंद.} SOL6-A120: एका inner substitution ची definedness मिळाल्यावर दुसऱ्याची definedness आपोआप मिळत नाही. सर्व आवश्यक नावे एका सांत forbidden संचाबाहेर निवडून, सर्वांत आतल्या अमूर्तीकरणांपासून नामांतर केले आहे. Binding-tree argument ने arbitrary alpha-equivalent प्रतिनिधींसाठी परिणामी alpha-class ची uniqueness सिद्ध केली आहे.\end{quote}
\begin{quote}\small\textbf{OLP-0363 — स्रोतनोंद.} SOL6-A121: नावांच्या यादीत एकाच नावाचे अनेक वापर असतील, तर पहिली जागा घ्यायची आहे. त्यामुळे सर्वांत जवळचे अमूर्तीकरण चराला बांधते. पुढील उलट रूपांतर फक्त प्रत्येक निर्देशांकाला उपलब्ध जागा असलेल्या पदांसाठी परिभाषित आहे.\end{quote}
\begin{quote}\small\textbf{OLP-0369 — स्रोतनोंद.} SOL6-A125: दोन आदेशनांच्या क्रमाची अदलाबदल करताना आवश्यक ताजेपणाच्या अटी आणि आदेशनाचे संयोजनसूत्र स्पष्ट केले आहे. आधी दुरुस्त केलेला अमूर्तीकरणाच्या प्रकरणातील न्यूनीत आदेशित प्रतिनिधी जपला आहे.\end{quote}
\begin{quote}\small\textbf{OLP-0371 — स्रोतनोंद.} SOL6-A126: इटा संकुचनाने फलनस्थानील अमूर्तीकरण नाहीसे होऊ शकते. त्यामुळे बीटा पुराव्यातील पहिली चार प्रकरणे तशीच लागू होत नाहीत. मुख्य triangle दावा आणि अमूर्तीकरणाच्या उपयोजनाचा सहाय्यक दावा यांचे एकत्र रचनात्मक विगमन देऊन ही उणीव पूर्ण केली आहे.\end{quote}
\begin{quote}\small\textbf{OLP-0375 — स्रोतनोंद.} SOL6-A128: मूळ e b हे पद शून्य घातांकासाठी identity मध्ये न्यूनीत होते; ते बीटा-सामान्य चर्च अंक एक नाही. बाहेरून चर्च अंकाचे दोन प्राचल लावल्याने सर्व घातांकांसाठी आवश्यक बीटा-सामान्य अंक मिळतो. शून्याच्या शून्य घाताचा येथे एक असा पुनरावर्ती संकेत आहे.\end{quote}
\begin{quote}\small\textbf{OLP-0394 — स्रोतनोंद.} SOL6-A140: दिलेल्या मोडल कोष्टकांत मधल्या U मूल्यावरील विरोधाभासाची संभाव्यता T होते; तिचे निषेधन F होते, मूळात दिलेले U नाही. मॅट्रिक्सचा चार-संयोजकीय भाषाखंड स्पष्ट केला. मूळ 1920 चा व्याख्यानवृत्तांत प्रत्यक्ष पाहून ऐतिहासिक वर्ष दुरुस्त केले; उद्धरणातील noon म्हणजे दुपारी बारा वाजता.\end{quote}
\begin{quote}\small\textbf{OLP-0395 — स्रोतनोंद.} SOL6-A141: दोन्ही क्लिनी मॅट्रिक्स येथे मानक L0 च्या असत्यता-स्थिर नसलेल्या चार-संयोजकीय खंडावर आहेत. सर्व चलांना U देणाऱ्या मूल्यनिर्णयनावरील सर्वतः सत्य सूत्र नसण्याचा पुरावा या व्याप्तीत लागू होतो.\end{quote}
\begin{quote}\small\textbf{OLP-0397 — स्रोतनोंद.} SOL6-A142: फक्त निर्दिष्ट मूल्ये बदलून लूकासेविच मॅट्रिक्सला R-Mingle म्हणणे अयोग्य होते. R-Mingle चे वेगळे सशर्त कोष्टक Kooi आणि Tamminga यांच्या primary paper मधील विभाग 4, पृ. 10 शी जुळवले; जोडलेल्या असत्यता-स्थिराचे मूल्य F स्पष्ट केले. LP आणि Hallden यांच्या classical tautologies ची समानता केवळ समान चार-संयोजकीय खंडात आहे.\end{quote}
\begin{quote}\small\textbf{OLP-0400 — स्रोतनोंद.} SOL6-A143: अनंतमूल्य आणि सांतमूल्य लूकासेविच मॅट्रिक्स येथे मानक L0 च्या चार-संयोजकीय खंडावर आहेत. याने मागील त्रिमूल्य मॅट्रिक्सशी भाषेची समानता जपली. सांत आधारसंचांसाठीचे प्रतिलोम हे आधीचे योग्य restriction जपले.\end{quote}
\begin{quote}\small\textbf{OLP-0404 — स्रोतनोंद.} SOL6-A144: प्रमेयाच्या क्रमवर्तीत फक्त निर्दिष्ट सत्यतामूल्यांच्या जागी सूत्र असते; इतर जागा रिकाम्या असतात. हा आवश्यक scope स्पष्ट केला. व्युत्पादन सांत वृक्ष आहे आणि त्यात तार्किक व संरचनात्मक नियम दोन्ही येतात. दाखवलेले Gamma उदाहरण सांत क्रमिकेसाठी; अनंत आधारसंचासाठी आधीच्या सांत-उपसंच व्याख्येचा वापर होतो.\end{quote}
\begin{quote}\small\textbf{OLP-0411 — स्रोतनोंद.} SOL6-A146: क्रमाने केलेल्या आदेशनाच्या दोन्ही मध्यवर्ती displays मध्ये संपूर्ण सूत्र दुसऱ्या आदेशनाच्या आत घेतले. पूर्वपक्षातील चलही बदलले पाहिजेत; केवळ उत्तरपक्ष बदलणे अपुरे आहे. अंतिम विस्तारित सूत्रे मूळातच योग्य होती.\end{quote}
\begin{quote}\small\textbf{OLP-0417 — स्रोतनोंद.} SOL6-A149: रूपबंधाच्या notation मध्ये केवळ विधानचल बदलतात; असत्यासारखे तार्किक स्थिरांक बदलत नाहीत. आदेशन-दृष्टांताच्या वैधतेच्या सिद्धतेतील नवीन valuation सर्व विधानचलांवर पूर्ण केली आणि तिचे truth sets प्रतिमानाच्या W मधील जगांपुरते स्पष्ट केले.\end{quote}
\begin{quote}\small\textbf{OLP-0421 — स्रोतनोंद.} SOL6-A150: एका असत्य दृष्टांताचे प्रदर्शन म्हणजे नेमका एकच दृष्टांत असत्य आहे असा दावा नाही. त्याच एक-जगी प्रतिमानात असत्य स्थिरांक घातल्यावर आणखी असत्य दृष्टांत मिळतो.\end{quote}
\begin{quote}\small\textbf{OLP-0423 — स्रोतनोंद.} SOL6-A151: D च्या proof मध्ये V हा model tuple मध्ये आहे; त्याची विशेष निवड करण्याची गरज नाही, असा independence चा अचूक अर्थ घेतला. B आणि 5 च्या proofs मध्ये वापरलेली पद्धत प्रतिपरिवर्तन आहे.\end{quote}
\begin{quote}\small\textbf{OLP-0426 — स्रोतनोंद.} SOL6-A153: मानक भाषांतराच्या प्रत्येक मोडल पायरीत नवीन बद्ध चल निवडण्याची अट स्पष्ट केली. Monadic हे द्वितीय-क्रम संख्यापकांपुरते बंधन आहे; व्यक्तिचलांवरील प्रथम-क्रम संख्यापक अनुमत आहेत. परावर्तितेच्या उदाहरणात व्यक्तिचलाचे मूल्य, प्रत्यक्ष जग आणि successor set वेगळे करून नेमणूक स्पष्ट केली.\end{quote}
\begin{quote}\small\textbf{OLP-0428 — स्रोतनोंद.} SOL6-A155: सामान्य मोडल तर्कशास्त्र मिळण्याचा दावा K आणि आवश्यक असल्यास Dual हे स्वयंसिद्धके घेतल्याच्या संदर्भात आहे. निर्दोषतेसाठी फक्त characteristic सूत्रे नव्हे, तर सर्व स्वयंसिद्धकांचे सर्व आदेशन-दृष्टांत प्रतिमानाच्या सर्व जगांत सत्य असणे आवश्यक आहे.\end{quote}
\begin{quote}\small\textbf{OLP-0431 — स्रोतनोंद.} SOL6-A156: दहा ओळींच्या K-सिद्धतेत ओळ 6 मधील implication चे बाह्य कोष्ठक उघडले आणि ओळी 8 व 9 मधील प्रत्येकी एक जादा बंद कोष्ठक काढले. ओळी 3 व 5 चे K-दृष्टांत आणि सर्व inference references आधीच योग्य होते; ते जपले.\end{quote}
\begin{quote}\small\textbf{OLP-0432 — स्रोतनोंद.} SOL6-A157: replacement-rule च्या संक्षिप्त नोंदीत “नवीन सूत्र for जुने सूत्र” असा क्रम वापरला आहे. म्हणून C(A) वरून C(B) मिळवताना label B for A आहे; हा क्रम खालील double-negation examples मधील p for not-not-p शी जुळतो.\end{quote}
\begin{quote}\small\textbf{OLP-0436 — स्रोतनोंद.} SOL6-A159: निर्दोषता विभागाच्या प्रस्तावनेतील axioms सत्य असण्याची अट पुढील प्रमेयाप्रमाणे सर्व आदेशन-दृष्टांत सर्व जगांत सत्य असण्याची आहे. मूळ शब्दांकनातील characteristic सूत्र आणि पूर्ण schema यांतील संदिग्धता दूर केली; प्रमेय व त्याची सिद्धता बदललेली नाही.\end{quote}
\begin{quote}\small\textbf{OLP-0438 — स्रोतनोंद.} SOL6-A160: प्रणालीतील derivability च्या व्याख्येचा संदर्भ त्याच्या योग्य विभागाकडे वळवला. गृहीतकांची संख्या शून्य असू शकते हे स्पष्ट केले; रिकाम्या premise set साठी कोणतीही गृहीतके न वापरणारी सिद्धता मूळ प्रणालीतील सिद्धताच असते.\end{quote}
\begin{quote}\small\textbf{OLP-0440 — स्रोतनोंद.} SOL6-A161: सापेक्ष विसंगतता आणि योग्य प्रतिमानांचा वर्ग. एखाद्या मजबूत प्रणालीच्या सापेक्ष विसंगत Gamma मनमानी Kripke प्रतिमानात satisfiable असू शकतो. उदाहरणार्थ KT च्या सापेक्ष Boxp आणि not-p विसंगत आहेत, पण nonreflexive model मध्ये एकत्र सत्य होतात. No-model दावा संबंधित system च्या सर्व axiom instances सर्व worlds मध्ये सत्य असलेल्या models पुरता केला. पुढील consistency definition आणि तीन facts बदलले नाहीत.\end{quote}
\begin{quote}\small\textbf{OLP-0442 — स्रोतनोंद.} SOL6-A162: सर्वसाधारण canonical construction मध्ये अनेक जगांची गरज. मूळ प्रेरक वाक्य कोणतेही A सत्य करण्यासाठी multiple worlds आवश्यक आहेत असे म्हणत होते. p सारखी सूत्रे single-world model मध्ये सत्य होतात; मात्र सर्व consistent formulas साठी single-world construction पुरेशी नाही. आवश्यकतेचा scope सर्वसाधारण रचनेवर केला. प्रत्यक्ष canonical world set सर्व complete consistent sets घेतो हा पुढील युक्तिवाद जपला.\end{quote}
\begin{quote}\small\textbf{OLP-0443 — स्रोतनोंद.} SOL6-A163: Complete consistent sets मधील disjunction आणि joint scope. Disjunction membership च्या iff proof मध्ये मूळाने फक्त forward direction दिली होती. दोन्ही classical introduction tautologies व deductive closure ने reverse direction जोडली. Biconditional च्या reverse proof मधील दोन्ही नसतील या Marathi wording मुळे not-both आणि both-not यांची गल्लत होत होती; दोन्ही statements एकत्र true असण्याच्या प्रत्येक शक्यतेचा निषेध स्पष्ट केला. आधीचे negation आणि biconditional-antecedent repairs जपले.\end{quote}
\begin{quote}\small\textbf{OLP-0444 — स्रोतनोंद.} SOL6-A164: Lindenbaum proof मधील रिकाम्या finite premise list चा निर्देशांक. Finite witness list शून्य-लांबीची असू शकते. मूळ proof मध्ये त्या list मधील indices चा largest घेतला होता, पण empty list ला largest नसतो. Max of zero plus the listed indices ही नेमकी convention दिली; k=0 साठी n=0 मिळतो. Increasing chain मधील finite witness containment आणि सर्व Delta\_n ची consistency proof जपली. पूर्वीची exhaustive enumeration साठी length-at-most-n दुरुस्ती योग्य.\end{quote}
\begin{quote}\small\textbf{OLP-0445 — स्रोतनोंद.} SOL6-A165: Canonical accessibility चे iff आणि शून्य गृहीतकांचा boxing lemma. Diamond branch मधील explanatory clause ने R खरे होण्यासाठी sufficient अट दिली होती; अपेक्षित implication साठी R ते condition ही दिशा आवश्यक आहे. औपचारिक पुढील definition प्रमाणे iff स्पष्ट केला. Boxing lemma च्या k=0 प्रकरणात RK नव्हे तर NEC थेट लागू केला; k>=1 साठी मूळ RK proof जपला. Modal equivalence च्या वाक्यात missing शी जोडले. पूर्वीचे k/n आणि system-subscript repairs योग्य.\end{quote}
\begin{quote}\small\textbf{OLP-0446 — स्रोतनोंद.} SOL6-A166: Canonical model चा nonempty domain आणि valuation ची world व्याप्ती. Normal logic च्या आधीच्या definition मध्ये inconsistent full logic वगळलेली नाही. अशा Sigma साठी complete consistent sets नसतात आणि W रिकामा होतो, जे basic model definition शी विसंगत आहे. Canonical model construction consistent normal systems साठी केली आणि Lindenbaum ने W nonempty दाखवले. Valuation truth set स्पष्टपणे Delta in W पुरता केला. Inconsistent systems साठी completeness चा explosion case स्वतंत्र हाताळता येतो; रिकामा set एक वैध Kripke model असल्याचा दावा नाही.\end{quote}
\begin{quote}\small\textbf{OLP-0447 — स्रोतनोंद.} SOL6-A167: Truth lemma मधील accessibility bridge आणि विगमन गृहीतक. Diamond च्या forward case मध्ये Box-inverse containment वरून Diamond-image containment मिळवण्यासाठी box-iff-diamond lemma आवश्यक आहे; जुन्या prop-diamond reference ने ती containment थेट मिळत नाही. योग्य lemma दाखवला. Reverse case मध्ये B चे membership आणि Diamond satisfaction यांच्या मध्ये induction hypothesis चे B satisfaction पाऊल जोडले. Theorem मध्ये consistent normal Sigma व Delta in W scope स्पष्ट केला; सर्व Boolean व Box cases आणि prior probAnd repair जपले.\end{quote}
\begin{quote}\small\textbf{OLP-0448 — स्रोतनोंद.} SOL6-A168: Determination theorem च्या hypotheses आणि canonical proof पद्धतीचा scope. Canonical model अस्तित्वासाठी consistent normal Sigma ही hypothesis determination theorem मध्ये स्पष्ट केली. K ची consistency पूर्वीच्या soundness आणि falsity च्या nonvalidity वरून मिळते, म्हणून K-completeness proof ला ही अट लागू होते. Arbitrary class completeness साठी canonical model membership नेहमी आवश्यक असे नाही; ही अट या canonical-model proof पद्धतीसाठी आवश्यक असल्याचे शब्दांकन केले. Determination iff आणि K corollary ची दोन्ही directions जपल्या.\end{quote}
\begin{quote}\small\textbf{OLP-0449 — स्रोतनोंद.} SOL6-A169: Canonical frame theorem ची consistency आणि weak-density witness. Canonical frame theorem आणि supplementary proposition सुसंगत normal systems पुरते स्पष्ट केले. Functionality proof मध्ये biconditional schema वरून D मिळतो हे नोंदले; partial functionality plus seriality देऊन exact one successor मिळतो. Weak-density proof मधील finite Diamond-premise list रिकामी मानल्यास त्याचे displayed nonempty conjunction steps अपूर्ण होतात. Box-inverse Delta1 हा consistent Delta2 चा subset असल्याने inconsistent witness ला किमान एक Diamond premise लागतो; m>=1 व n>=0 हे नेमके bounds दिले. पाच canonical relation proofs व all remaining displayed steps जपले.\end{quote}
\begin{quote}\small\textbf{OLP-0451 — स्रोतनोंद.} SOL6-A170: finite-model search मध्ये निश्चित frame व सूत्रातील वेगवेगळ्या चलांच्या valuations एवढाच count आहे. Frame-class test साठी दिलेली सांत schema-list आणि stopping bound साठी effective computation स्पष्ट केली. Universal-frame प्रेरक quotient केवळ लक्ष्य सूत्रातील सांत variables वर घेतला; सर्व भाषेतील variables वर agreement मागितल्यास finite quotient मिळेलच असे नाही. Formal filtration मध्ये सर्व चलांची valuation class-image set ने होते, पण selected formulas पुरते truth preservation असते. Box induction ची दुसरी दिशा पूर्ण केली.\end{quote}
\begin{quote}\small\textbf{OLP-0456 — स्रोतनोंद.} SOL6-A171: प्रणालीतील प्रत्येक सूत्र प्रतिमानाच्या प्रत्येक जगावर सत्य असणे आणि चौकट-गुणधर्म असणे हे स्वतंत्र आहेत. S5 च्या कोणत्याही प्रतिमानात लक्ष्य सूत्र सत्य असेल तर त्याचा निषेध सिद्ध होत नाही; संपूर्णतेवरून परावर्ती युक्लिडी प्रतिमान मिळते. त्यानंतर वैश्विक प्रतिमान आणि सांत निस्यंदनाचे मूळ पाऊल लागू होते.\end{quote}
\begin{quote}\small\textbf{OLP-0457 — स्रोतनोंद.} SOL6-A172: निर्णेयतेच्या प्रस्तावनेत दिलेली सांत आणि प्रभावी रूपबंध-यादी स्पष्ट केली. S5 च्या प्रतिमान-शोधात प्रत्येक आकारासाठी निश्चित जगांची नावे, वैश्विक संबंध आणि लक्ष्य सूत्रातील चलांचीच सांत मूल्यांकन-यादी वापरली जाते; उर्वरित चलांचे मूल्यांकन रिक्त ठेवता येते.\end{quote}
\begin{quote}\small\textbf{OLP-0458 — स्रोतनोंद.} SOL6-A173: अतिरिक्त अटीमुळे प्राप्यतेच्या जोड्या वाढत नाहीत; त्या कमी होतीलच असे नाही. दोन्ही निस्यंदने समान होऊ शकतात, म्हणून अधिक सूक्ष्मपणाच्या तुलनेत समतेची शक्यता स्पष्ट केली.\end{quote}
\begin{quote}\small\textbf{OLP-0459 — स्रोतनोंद.} SOL6-A174: रिक्त संचही मोडलदृष्ट्या बंद असतो; अनंततेचा दावा रिक्त नसलेल्या संचांसाठी आहे. संक्रमकतेच्या सिद्धतेत C1 ज्या modal सूत्रांवर परिमाणित आहे ती Gamma मध्ये असल्याची अट स्पष्ट घेतली; त्यातून उपसूत्र व मोडल बंदतेचा उपयोग वैध ठरतो.\end{quote}
\begin{quote}\small\textbf{OLP-0462 — स्रोतनोंद.} SOL6-A175: नव्या पूर्वचिन्हाच्या अटीचे प्रतिउदाहरण दाखवणाऱ्या केवळ Box-पर्यायात असत्य Box सूत्रावर लागू होणाऱ्या नियमाचे लेबल असत्य Box असे दुरुस्त केले. उदाहरणातील जाणीवपूर्वक शिथिल केलेली नवेपणाची अट बदललेली नाही.\end{quote}
\begin{quote}\small\textbf{OLP-0464 — स्रोतनोंद.} SOL6-A176: निर्दोषतेच्या सिद्धतेत मोडल प्रतिमान आणि मोडल पूर्तीची संज्ञा वापरली. नियमानुसार तयार झालेल्या शाखेची पूर्वचिन्हे आरंभीच्या रिक्त नसलेल्या खंडांखाली बंद असतात. म्हणून नवे पूर्वचिन्ह वाढवताना त्याचे आधीपासूनचे वंशज नसतात आणि फक्त पालकापासून नव्या पूर्वचिन्हापर्यंतची प्राप्यता तपासणे पुरेसे असते.\end{quote}
\begin{quote}\small\textbf{OLP-0465 — स्रोतनोंद.} SOL6-A179: अतिरिक्त नियमांच्या उदाहरणात 5 स्वयंसिद्धकाचे योग्य सूत्र आणि त्याचा बंद टॅब्लो दिला. मूळातील Box A वरून Box Diamond A हे वेगळे सूत्र होते. योग्य Diamond A गृहीतकासाठी सत्य Diamond नियमाने दुसरे नवे पूर्वचिन्ह आणले; शेवटचा असत्य Diamond नियम त्याच वापरलेल्या पूर्वचिन्हावर लागू होतो.\end{quote}
\begin{quote}\small\textbf{OLP-0468 — स्रोतनोंद.} SOL6-A180: संपूर्णतेच्या रचनेत बंद आणि संपूर्ण शाखा वेगळ्या केल्या. गृहीतकांचे समान पूर्वचिन्ह 1 स्पष्ट केले; रिक्त संचाचे प्रकरण वेगळे हाताळले. सांत रचनेचे थांबणे मोडल खोली आणि सांत संततीवरून स्पष्ट केले. असांत संचासाठी प्रत्येक सांत उपसंचाची पूर्ततायोग्यता, सुसंगततेची सांत सिद्धता, लिंडेनबाउमचे पूर्वप्रमेय आणि कॅनॉनिकल सत्यता पूर्वप्रमेय वापरून सर्वसाधारण निष्कर्ष जपला.\end{quote}
\begin{quote}\small\textbf{OLP-0469 — स्रोतनोंद.} SOL6-A181: निर्णय प्रक्रियेतील नवे पूर्वचिन्ह देणारा नियम प्रत्येक पूर्वचिन्हयुक्त चिन्हांकित सूत्रासाठी एकदा लावायचा असतो; वापरलेली पूर्वचिन्हे घेणाऱ्या नियमाची व्याप्ती त्याच्या तात्काळ विस्तारांपुरती आहे. थांबण्याचा सांत इनपुटवरील पुरावा स्पष्ट केला. मनमानी गृहीतकसंचासाठी संपूर्णतेचा निष्कर्ष जपून, प्रत्यक्ष सांत निर्णय प्रक्रियेचे विधान सांत गृहीतकसंचापुरते केले. उदाहरणांतील वापरलेली पूर्वचिन्हे 1.n या आकाराची आहेत हे स्पष्ट केले; सर्व सहा टॅब्लो आणि दोन्ही प्रतिमाने जपली.\end{quote}
\begin{quote}\small\textbf{OLP-0471 — स्रोतनोंद.} SOL6-A183: S5 साठी कट-मुक्त कलन ज्ञात नसल्याचा मूळ दावा ऐतिहासिक मसुद्यापुरता केला. प्रत्यक्ष वाचलेल्या Aghaei आणि Mohammadi यांच्या 2018 लेखाचा स्रोत जोडला; त्यातील G3S5 आणि पुढे दिलेला विशिष्ट LK-विस्तार यांची व्याप्ती वेगळी ठेवली.\end{quote}
\begin{quote}\small\textbf{OLP-0473 — स्रोतनोंद.} SOL6-A184: Dual च्या सिद्धतेतील पहिल्या नकार-पायरीत उजव्या बाजूचे सूत्र नकारासह डावीकडे जाते. त्यामुळे चुकीचे उजवा-नकार नियम-नाव डावा-नकार असे केले. अंतिम संयोग-नियमाच्या पायरीत द्विशर्ती दोन सशर्त सूत्रांच्या संयोगाचे संक्षिप्त रूप म्हणून घेतली आहे हे स्पष्ट केले; मूळ सिद्धतेतील सूत्रे जपली.\end{quote}
\begin{quote}\small\textbf{OLP-0474 — स्रोतनोंद.} SOL6-A185: मोडल क्रमवर्ती प्रणाल्यांच्या तक्त्यात मूळ नियमांची नावे primitive Box/Diamond निवडीनुसार केली; दोन्ही primitive असताना दोन्ही नियम दाखवले. Diamond-only कट-सिद्धतेतील डाव्या सूत्रांच्या अदलाबदलीचे चुकीचे उजवे नियम-नाव डावे केले. कटविना असंपूर्णतेचा दावा या विभागातील विशिष्ट कलनापुरता स्पष्ट केला.\end{quote}
\begin{quote}\small\textbf{OLP-0481 — स्रोतनोंद.} SOL6-A187: इतिहासात कालिक स्थान स्पष्ट होण्यासाठी प्रत्येक अवस्था-घटना स्वतंत्र मानली; प्रणालीची स्थिती पुन्हा आली तरी तिच्या वेगळ्या वेळच्या घटना वेगळ्या आहेत. Boolean उपसूत्रे त्याच बिंदू व इतिहासावर तपासायची हे स्पष्ट केले. शेवटच्या उदाहरणातील कधीच पुढे सत्य होत नाही हा सार्विक नकार दुरुस्त केला.\end{quote}
\begin{quote}\small\textbf{OLP-0484 — स्रोतनोंद.} SOL6-A188: मोडल-विरहिततेची अट सर्व कर्त्यांसाठी केली. समूहज्ञानाचा संयोग फक्त सीमित समूहासाठी साधा formula abbreviation आहे हे स्पष्ट केले; अनंत समूहाचे स्वतंत्र अर्थनिर्धारण व सामायिक ज्ञानाचा भाषा-विस्तार वेगळा दिला. वैयक्तिक आणि सामायिक ज्ञान नेहमी काटेकोरपणे वेगळे असतात हा अर्थ टाळला.\end{quote}
\begin{quote}\small\textbf{OLP-0485 — स्रोतनोंद.} SOL6-A192: प्रतिमानातील R हा कर्त्यानुसार निर्देशित संबंधांचा परिवार आहे हे स्पष्ट केले. दोन कर्त्यांचे संबंध सारखे असले तरी त्यांची नावे या परिवारात राखली जातात. V(p) हा W चा उपसंच आणि कर्त्याचे a हे गणिती चिन्ह स्पष्ट केले.\end{quote}
\begin{quote}\small\textbf{OLP-0486 — स्रोतनोंद.} SOL6-A189: संक्रमण-संवरणाच्या powers मध्ये मूळ zero-based indexing ऐवजी पूर्वीच्या व्याख्येशी जुळणारे R\textasciicircum{}1=R आणि n>0 union वापरले. Positive-path आणि reflexive-path conventions वेगळ्या ठेवल्या; मूळ common-knowledge semantics बदलली नाही.\end{quote}
\begin{quote}\small\textbf{OLP-0487 — स्रोतनोंद.} SOL6-A191: सकारात्मक व नकारात्मक आत्मनिरीक्षणाच्या प्रतिउदाहरणांना वेगळे पूर्वांग दिले. एखादी गोष्ट माहीत असणे हे नकारात्मक आत्मनिरीक्षणाचे पूर्वांग नाही; त्या तत्त्वासाठी माहीत नसणे आवश्यक आहे. तक्त्यातील मूळ सूत्रे जपली.\end{quote}
\begin{quote}\small\textbf{OLP-0488 — स्रोतनोंद.} SOL6-A190: द्विअनुकरणासाठी graphs समरूप असण्याची गरज नाही; forth/back उत्तरवर्ती अनुरूपता आवश्यक आहे हे स्पष्ट केले. चित्रात न दिलेल्या valuations आणि इतर कर्त्यांचे संबंध नेमके निवडल्यावरच दाखवलेला तीन-जोड्यांचा संबंध द्विअनुकरण आहे असे मांडले; मूळ graph जपला.\end{quote}
\begin{quote}\small\textbf{OLP-0490 — स्रोतनोंद.} SOL6-A193: p ची घोषणा आणि p असूनही b ला p माहीत नाही याची घोषणा समान प्रतिमान देतात हा मूळ दावा चुकीचा होता. पहिल्यात दोन जगे, दुसऱ्यात फक्त w1 उरते; actual arrows व valuations वरून हे स्पष्ट केले. असत्य घोषणा-पूर्वांगात उरलेल्या प्रतिमानात मूल्यांकन करायचे नाही हेही स्पष्ट केले; म्हणून रिकामा restriction undefined point निर्माण करत नाही.\end{quote}
\begin{quote}\small\textbf{OLP-0493 — स्रोतनोंद.} SOL6-A345: SOL6-A195 मधील गणिताने वैध पण ऐच्छिक अपरिमेयता-सिद्धता भाषांतरातून मागे घेतली. मूळाची संक्षिप्त मांडणी, निवडलेली मूल्ये आणि घाताची गणना जपली; विस्तृत सिद्धता अंमलात न आणलेली सुधारणा-सूचना म्हणून नोंदवली आहे.\end{quote}
\begin{quote}\small\textbf{OLP-0499 — स्रोतनोंद.} SOL6-A196: अंतःप्रज्ञावादी रचनाशील पडताळणी आणि साधारण epistemic knowledge operator यांतील फरक स्पष्ट केला. मूळ disjunction चे स्पष्टीकरण सामान्य knowledge operator साठी चुकीचे होते; मागील प्रत्यक्ष प्रतिमानातील counterexample आणि पुढील विभागातील योग्य local forcing clauses दिल्या.\end{quote}
\begin{quote}\small\textbf{OLP-0523 — स्रोतनोंद.} SOL6-A213: सर्वांत लहान पूर्वांग-गोलकाचा उल्लेख आकृतीपुरता स्पष्ट केला. दोन शेजारच्या दर्शवलेल्या थरांची जवळीक आणि मनमानी गोलकांतील जवळीक वेगळी आहे; सामान्य गोलक-अट सर्वात जवळचे जग नसण्यालाही लागू होते.\end{quote}
\begin{quote}\small\textbf{OLP-0524 — स्रोतनोंद.} SOL6-A214: रिकाम्या गोलकाची परवानगी व अनंत उपकुलांवरील closure स्पष्ट केली. सर्वांत लहान पूर्वांग-गोलक नसण्याचा पूर्ण प्रतिमान दिला; फक्त उतरत्या क्रमिकेवरून तो निष्कर्ष काढता येत नाही. पूर्वीचा लहान पूर्वांग-गोलकांसाठीचा source repair जपला.\end{quote}
\begin{quote}\small\textbf{OLP-0525 — स्रोतनोंद.} SOL6-A217: या आकृत्यांत सर्वाधिक जवळची पूर्वांग-जगे उपलब्ध आहेत, ही अट स्पष्ट केली. सर्वांत जवळचे जग नसलेल्या प्रतिमानांत सामान्य गोलक-अटच वापरायची; कठोर अभिव्यंजनाशी अवश्यतेची तुलना S5 च्या संदर्भात आहे.\end{quote}
\begin{quote}\small\textbf{OLP-0533 — स्रोतनोंद.} SOL6-A223: ही स्तराधारित संचरचनेची अनौपचारिक रूपरेषा आहे. रसेलचा तोच युक्तिवाद रोखणे आणि संपूर्ण औपचारिक उपपत्तीची सुसंगतता सिद्ध करणे वेगळे दावे आहेत; येथे पहिल्या दाव्याचे कारण दिले आहे.\end{quote}
\begin{quote}\small\textbf{OLP-0534 — स्रोतनोंद.} SOL6-A224: आरंभीच्या टप्प्याला आधीचे टप्पे नसतात. संचयी मांडणीमुळे रसेलचा अमर्याद संच तयार होत नाही हा युक्तिवाद आहे; कोणत्याही औपचारिक संच उपपत्तीच्या पूर्ण सुसंगततेची सिद्धता नाही.\end{quote}
\begin{quote}\small\textbf{OLP-0536 — स्रोतनोंद.} SOL6-A226: शेवटच्या चरणात सदस्यत्वाचा द्वितीय-क्रम संख्यापक असलेले संकल्पनारचनेचे उदाहरण वापरले आहे; पूर्ण रूपबंधापेक्षा या विशिष्ट उदाहरणाची उपलब्धता आवश्यक आहे. मर्यादित रूपबंधांची व्याप्ती वेगळी असू शकते; Sean Walsh, Fragments of Frege’s Grundgesetze and Gödel’s Constructible Universe, विभाग 2, व्याख्या 2.1--2.2, \url{https://arxiv.org/abs/1407.3861}, येथे ही मर्यादा स्पष्ट केली आहे.\end{quote}
\begin{quote}\small\textbf{OLP-0539 — स्रोतनोंद.} SOL6-A227: स्वसदस्यत्व-विरोधी युक्तिवाद प्रथम निर्माण-टप्पा आणि त्याच्या आधी तयार झालेले सदस्य वापरतो. सदस्यांचा शेवटचा टप्पा असावा किंवा पुढचा टप्पा लगतचा असावा अशी अतिरिक्त अट नाही.\end{quote}
\begin{quote}\small\textbf{OLP-0543 — स्रोतनोंद.} SOL6-A229: मूळ मनमाना अनंतता-साक्षी संचावर एकास-एकपणा थेट मानण्याऐवजी लघुतम संवरणावर तो सिद्ध केला. नियमिततेचे स्वयंसिद्धक वापरलेले नाही. इटाय बेन-याकोव यांच्या MATH 770: Foundations of Mathematics, विभाग 5.2.1, पृ.\textasciitilde{}86, मधील लघुतम संच व त्याच्या सदस्यांचे गुणधर्म या मार्गाशी जुळतात: \url{https://people.math.wisc.edu/~awmille1/old/m571-08/ben-yaacov.pdf}.\end{quote}
\begin{quote}\small\textbf{OLP-0550 — स्रोतनोंद.} SOL6-A230: सुक्रमणाचा संबंध दिलेल्या संचावर आहे. उपसंचात पूर्ववर्ती नसणे आणि त्यातील इतर सर्व सदस्यांपेक्षा लहान असणे यांतील भाषांतरातील भेद स्पष्ट केला. मूळाची संक्षिप्त सिद्धता जपली आहे; स्वतंत्र विस्तार केलेला नाही.\end{quote}
\begin{quote}\small\textbf{OLP-0551 — स्रोतनोंद.} SOL6-A231: द्विदिश अटी आणि विभक्तीकरणाचा मराठी शब्द दुरुस्त केला. मूळातील वैध प्रतिलोम-तर्क, प्रतिमेचा समतुल्य संकेत आणि संक्षिप्त तुलना-सिद्धता जपली आहेत; त्यांचे स्वतंत्र स्पष्टीकरण भाषांतरात घातलेले नाही.\end{quote}
\begin{quote}\small\textbf{OLP-0554 — स्रोतनोंद.} SOL6-A234: संच-फलनाच्या प्रतिमेत पदमूल्य वापरताना घटक फलनाच्या प्रांतात असावा ही अट स्पष्ट केली. मनमाना संच घेतल्यास प्रतिमा फलनाच्या आलेखावरून परिभाषित करून विभक्तीकरणाने मिळते; यासाठी प्रतिस्थापन लागत नाही.\end{quote}
\begin{quote}\small\textbf{OLP-0569 — स्रोतनोंद.} SOL6-A245: मूळाच्या प्रतिमा-तुलनेत निवडीची पुरेशी अट स्पष्ट केली. आच्छादक फलनापासून उलट दिशेचे एकास-एक फलन निवडीच्या स्वयंसिद्धकाखाली मिळते; प्रतिस्थापनाला स्वतः निवड आवश्यक आहे किंवा ही तुलना पूर्ण निवडीशी समतुल्य आहे असा दावा नाही.\end{quote}
\begin{quote}\small\textbf{OLP-0573 — स्रोतनोंद.} SOL6-A246: सर्व संक्रमक संचांत विभक्तीकरण सत्य असते हा मूळ दावा चुकीचा आहे. आवश्यक उपसंच-समावेशकत्वाची अट सिद्धतेच्या लघुतम प्रतिमानात आणि सरावाच्या विधानात जोडली. सरावाची सिद्धता दिलेली नाही. सांत स्वयंसिद्धकीकरणीयतेच्या पहिल्या निष्कर्षात सुसंगततेची अट आणि प्रतिमानांच्या अरिक्ततेचा आधार स्पष्ट केला.\end{quote}
\begin{quote}\small\textbf{OLP-0579 — स्रोतनोंद.} SOL6-A249: घातांकनाच्या येथे दिलेल्या दोन्ही व्याख्यांत आधार शून्येतर असण्याची अट स्पष्ट केली. शून्य आधारासाठी सीमा-सूत्रातील संयोगात शून्य घातांकाचा समावेश केल्यास चुकीचा परिणाम मिळतो; आधीच्या सरावप्रश्नात ही अट होती, पण व्याख्येत नव्हती. मूळ घातांकन सूत्रे व न सोडवलेला सराव जपले.\end{quote}
\begin{quote}\small\textbf{OLP-0597 — स्रोतनोंद.} SOL6-A263: निवड नसलेल्या उदाहरणात आधीच्या सुक्रमित संचांकांवर आधारित बेथ संकेताऐवजी वास्तव संख्यांचा नेमका दावा दिला. फेफरमन–लेव्ही यांच्या मूळ सारांशात, ZF सुसंगत असेल तर वास्तव संख्यांचा संच गणनीय संचांचे गणनीय युतीकरण असणेही सुसंगत आहे, असे म्हटले आहे: Notices of the American Mathematical Society 10 (1963), पृ. 593, सारांश 63T-389, https://www.ams.org/journals/notices/196310/196310FullIssue.pdf. येथे पूर्ण forcing सिद्धतेचे स्वतंत्र प्रमाणीकरण केलेले नाही.\end{quote}
\begin{quote}\small\textbf{OLP-0622 — स्रोतनोंद.} SOL6-A272: कँटर यांच्या उपचाराच्या आरंभाचे स्रोतामधील 1894 हे वर्ष 1884 केले आहे. ऑलिव्हर डाइझर यांच्या Einführung in die Mengenlehre मधील चरित्रकालक्रम, पृ. 218, मे 1884 मध्ये हाल येथील आरोग्यधामातील वास्तव्य नोंदवतो; लिओपोल्डिनाचे सदस्यचरित्रही 1884 पासून आजार आणि पुन्हा पुन्हा उपचार नोंदवते: https://www.leopoldina.org/en/members/member-list/detail/georg-cantor. येथे संपूर्ण वैद्यकीय नोंदींचे स्वतंत्र पुनरावलोकन केलेले नाही.\end{quote}
\begin{quote}\small\textbf{OLP-0624 — स्रोतनोंद.} SOL6-A273: गेंटझेन यांच्या तार्किक निगमनावरील I–II लेखांसाठी Springerची मूळ प्रकाशकनोंद 21 जुलै 1933 हा प्राप्तिदिनांक आणि 1935 हे प्रकाशनवर्ष देते. म्हणून 1934 शी जोडलेले निर्मितीवर्ष सांगण्याऐवजी निश्चित प्रकाशनवर्ष दिले आहे. मूळ लेख: https://doi.org/10.1007/BF01201353 आणि https://doi.org/10.1007/BF01201363. येथे संपूर्ण चरित्राचे स्वतंत्र ऐतिहासिक प्रमाणीकरण केलेले नाही.\end{quote}
\begin{quote}\small\textbf{OLP-0625 — स्रोतनोंद.} SOL6-A274: गोडेल यांचा प्रबंध 1929 मध्ये पूर्ण झाला आणि डॉक्टरेटची पदवी 1930 मध्ये मिळाली, हा भेद स्पष्ट केला आहे. प्रबंधानंतर वर्षभरात मिळालेल्या अपूर्णता निष्कर्षांचा निर्देश प्रबंधाशीच जोडला आहे; 1931 प्रकाशनवर्ष जपले आहे. व्हिएन्ना विद्यापीठाचे इतिहासवर्णन आणि IAS पदवीनोंद: https://geschichte.univie.ac.at/de/artikel/von-der-berechenbarkeit-zum-modernen-computer आणि https://www.ias.edu/scholars/godel.\end{quote}
\begin{quote}\small\textbf{OLP-0629 — स्रोतनोंद.} SOL6-A275: मूळातील सहा महिन्यांचा उल्लेख शिक्षेचा आहे; प्रत्यक्ष सुटका १४ सप्टेंबर १९१८ रोजी झाली. तसेच आयुष्यभर बिनशर्त शांततावाद हा दावा मागे घेतला: मॅकमास्टरच्या 30909 या पत्रनोंदीतील रसेल यांच्या स्वतःच्या विधानात काही परिस्थितींमध्ये बलप्रयोग समर्थनीय मानला आहे. स्रोत: https://russell-letters.mcmaster.ca/chronology आणि https://bracers.mcmaster.ca/30909.\end{quote}
\begin{quote}\small\textbf{OLP-0631 — स्रोतनोंद.} SOL6-A276: GCHQच्या नोंदीनुसार जोन क्लार्क आणि ट्यूरिंग यांचा साखरपुडा १९४१ मध्ये झाला. बॉम्ब यंत्र विद्युत् आणि यांत्रिक तंत्रावर चालत होते, असे The National Museum of Computing स्पष्ट करते; ते इलेक्ट्रॉनिक कोलॉससपेक्षा वेगळे होते. स्रोत: https://www.gchq.gov.uk/information/joan-clarke आणि https://www.tnmoc.org/bombe.\end{quote}
\begin{quote}\small\textbf{OLP-0632 — स्रोतनोंद.} SOL6-A277: झुरिकमधील १९१० चे पद वेतन मिळणारे होते; त्यामुळे त्यापूर्वीची गॉटिंगेनमधील १९०५ ची प्राध्यापक नेमणूक नाकारली जात नाही. फ्रायबुर्गची मानद नेमणूक १९२६ मध्ये झाली. हिटलरला अभिवादन करण्यास नकार दिल्यामुळे तक्रार झाल्यानंतर १९३५ मध्ये राजीनामा द्यावा लागला, अशी विद्यापीठाची नोंद आहे. स्रोत: https://uni-freiburg.de/math/geschichte-i/ आणि https://home.mathematik.uni-freiburg.de/mildenberger/freiburg2014/zermelo.html.\end{quote}
\begin{quote}\small\textbf{OLP-0635 — स्रोतनोंद.} OLINC-314: स्रोताच्या सीमा-व्याख्येत x=c वगळलेले नाही; मराठी सूत्रात 0<|x-c| ही आवश्यक अट घातली आहे.\end{quote}
\begin{quote}\small\textbf{OLP-0638 — स्रोतनोंद.} SOL6-A281: एकच द्विमान निरूपण-पद्धत वापरून एकास-एक फलनाच्या सिद्धतेतील अंक-एकमेवपणाची अट पूर्ण केली आहे. मूळ प्रतिमेबाहेरील संख्येचे उदाहरण [0,1] वर प्रतिमेतच मिळते; खरे प्रतिउदाहरण 5/6 दिले आहे. हा आधुनिक द्विमान युक्तिवाद मूळ दशमान पत्रातील शब्दशः सिद्धता म्हणून मांडलेला नाही.\end{quote}
\begin{quote}\small\textbf{OLP-0639 — स्रोतनोंद.} SOL6-A282: मूळ चित्रे जपली आहेत. समान प्राचल-अंतरालांवरील जाळी-घरांची सुसंगत विभागणी स्पष्ट केल्याने सीमा समसमान असल्याचे आणि प्रतिमा पूर्ण चौरस असल्याचे सिद्ध होते; सातत्याची अट प्रत्येक आदान बिंदूवर पूर्ण केली आहे.\end{quote}
\begin{quote}\small\textbf{OLP-0639 — स्रोतनोंद.} OLINC-317: स्रोताने x चा प्रांत चौरस दिला; येथे एकक रेषा हा योग्य प्रांत जपला आहे.\end{quote}
\begin{quote}\small\textbf{OLP-0639 — स्रोतनोंद.} OLINC-318: मूळ बिंदुनिहाय-सीमा युक्तिवादातील उणीव SOL6-A282 मध्ये दुरुस्त केली: समान प्राचल-अंतरालांवरील घर-अंतर्भाव, समसमान सीमा आणि अंतर्भूत बंद अंतरालांनी प्रतिमेतील प्रत्येक बिंदूची साक्ष मिळते.\end{quote}
\begin{quote}\small\textbf{OLP-0639 — स्रोतनोंद.} OLINC-319: प्रत्येक p साठी कुठलातरी खुला अंतराल देणे ही सातत्याची अट नव्हती. SOL6-A282 मध्ये प्रत्येक दिलेल्या x0 वरील epsilon-delta अट समसमान अभिसरणाने सिद्ध केली आहे.\end{quote}
\begin{quote}\small\textbf{OLP-0639 — स्रोतनोंद.} OLINC-320: मूळ शेवटच्या परिच्छेदातील असंबद्ध (a,b) उल्लेख आणि 'show to show' पुनरुक्ती नवीन सिद्धतेत आलेली नाहीत.\end{quote}
\begin{quote}\small\textbf{OLP-0655 — स्रोतनोंद.} SOL6-A294: उलटवण्यात प्रत्येक कोटीच्या कटांची संख्या वाढत नाही, ही सिद्ध सीमा वापरली आहे. अणु-कटातील विगमन कमी उंचीच्या तत्काळ उपसिद्धतेवर लावले आहे. कमी कोटीचे कट असतानाच्या दुर्बलीकरण व आकुंचनाचा आवश्यक आधार स्पष्ट केला आहे; सरावाचे प्रसंग सोडवलेले नाहीत.\end{quote}
\begin{quote}\small\textbf{OLP-0655 — स्रोतनोंद.} OLINC-339: स्रोताचे Delta=Delta,A हे चक्रीय समीकरण; आवश्यक विभागणी Delta=Delta',A अशी दुरुस्त केली आहे.\end{quote}
\begin{quote}\small\textbf{OLP-0655 — स्रोतनोंद.} OLINC-340: स्रोत येथे सर्वोच्च कटासाठीचे वेगळे उपप्रमेय दाखवतो; या परिच्छेदातील कोटी-कमीकरण मागील उपप्रमेयावर आधारित आहे.\end{quote}
\begin{quote}\small\textbf{OLP-0656 — स्रोतनोंद.} SOL6-A295: CutCS फक्त कट-सूत्र काढतो आणि उरलेला सामायिक संदर्भ जपतो. अस्तित्व-सूत्राचा संदर्भ पहिल्या वृक्षात जपला आहे; अभिव्यंजनाच्या वृक्षात आवश्यक दुर्बलीकरणाने दोन्ही आधारांचा संदर्भ जुळवला आहे आणि खालच्या कटचा योग्य निष्कर्ष दिला आहे.\end{quote}
\begin{quote}\small\textbf{OLP-0656 — स्रोतनोंद.} OLINC-341: स्रोत दुसऱ्या प्रसंगात pi\_2 लिहितो; उजवा आधार स्वयंसिद्धक असल्यास आवश्यक सिद्धता डाव्या आधाराची pi\_1 असते.\end{quote}
\begin{quote}\small\textbf{OLP-0656 — स्रोतनोंद.} OLINC-342: तिसऱ्या वृक्षात स्रोताचे Delta हे Delta' हवे; अन्यथा संदर्भातील B-and-C ची अतिरिक्त प्रत येते.\end{quote}
\begin{quote}\small\textbf{OLP-0656 — स्रोतनोंद.} OLINC-343: अंतिम वृक्षात Delta' आणि B-and-C च्या दोन प्रती हव्यात; स्रोताच्या Delta मुळे तीन प्रती येतात.\end{quote}
\begin{quote}\small\textbf{OLP-0658 — स्रोतनोंद.} SOL6-A296: चर-साक्षीचा प्रसंग भाषा-सीमेवरून पूर्ण केला आहे. सामान्य अंतर्वेशक सूत्ररूपात असतो; बंद संदर्भांतील उपयोगांसाठी मुक्त चरे सार्वत्रिकपणे बांधून वाक्यरूप घेतले आहे. OLINC-344 मधील ऐतिहासिक missing-case नोंद मूळ स्रोतासाठी कायम आहे, सध्याचा प्रसंग दुरुस्त आहे.\end{quote}
\begin{quote}\small\textbf{OLP-0658 — स्रोतनोंद.} SOL6-A303: बेथच्या परिभाष्यता परिच्छेदात उरलेल्या इंग्रजी predicate ऐवजी मराठी विधेय-संज्ञा वापरली आहे.\end{quote}
\begin{quote}\small\textbf{OLP-0659 — स्रोतनोंद.} SOL6-A297: सामायिक संदर्भ असलेल्या कटच्या वृक्षात CutCS नियमलेबल वापरले आहे.\end{quote}
\begin{quote}\small\textbf{OLP-0661 — स्रोतनोंद.} SOL6-A298: अग्र-संख्यापक सूत्राची matrix संख्यापकविरहित असल्याचे स्पष्ट केले आहे. आकुंचन-उलटवण्याने क्रममान वाढत नाही; अनुमाने काढल्यास ते कमी होऊ शकते.\end{quote}
\begin{quote}\small\textbf{OLP-0662 — स्रोतनोंद.} SOL6-A300: कलम-संयोजनापूर्वी बाहेरच्या सिद्धतेच्या मुक्ती-खुणा आत बसवलेल्या सिद्धतेच्या उघड्या गृहीतकांच्या खुणांपासून वेगळ्या ठेवायच्या; अन्यथा समान सूत्र असले तरी गृहीतक अनवधानाने मुक्त होऊ शकते.\end{quote}
\begin{quote}\small\textbf{OLP-0665 — स्रोतनोंद.} SOL6-A301: या प्रकरणाच्या बंद पदे वापरण्याच्या संकेताशी आदेशनाची अट जुळवली आहे. एका अस्वच्छ आयगेन-चर अनुमानाची दुरुस्ती केल्यावर अस्वच्छ अनुमानांची संख्या किमान एकाने घटते; नेमकी एकानेच घटते असा दावा नाही.\end{quote}
\begin{quote}\small\textbf{OLP-0671 — स्रोतनोंद.} SOL6-A302: अभिजात असत्यता नियमातून अंतिम सूत्र असत्य असल्यास आणखी कट वापरून रिकामा उत्तरपक्ष मिळवला आहे; कोणतेही गृहीतक मुक्त न होणाऱ्या प्रसंगाचाही अपेक्षित निष्कर्ष स्पष्ट केला आहे.\end{quote}
\begin{quote}\small\textbf{OLP-0672 — स्रोतनोंद.} SOL6-A307: उपसूत्राच्या गुणधर्माची परिचयातील सूचना पुढील अचूक निष्कर्षाशी जुळवली: उघडी गृहीतके, संख्यापकांची पद-उदाहरणे आणि आवश्यक असत्य समाविष्ट आहेत.\end{quote}
\begin{quote}\small\textbf{OLP-0675 — स्रोतनोंद.} SOL6-A304: क्रमपरिवर्तनापूर्वी बद्ध गृहीतक-खुणा सुरक्षित नवी नावे द्यायची; उपसिद्धतेच्या प्रत्येक प्रतीतील आयगेन स्थिरांक स्वतंत्र ठेवून नियमितता जपायची.\end{quote}
\begin{quote}\small\textbf{OLP-0676 — स्रोतनोंद.} SOL6-A305: न्यूनीकरणात सुरक्षित खुणा व आयगेन स्थिरांक वापरायचे. नव्या उघड्या गृहीतकांचा संच मूळच्या संचात समाविष्ट असतो; तो नेहमी तंतोतंत समान असेल किंवा खूण x पूर्णपणे नाहीशी होईल असा दावा नाही.\end{quote}
\begin{quote}\small\textbf{OLP-0676 — स्रोतनोंद.} SOL6-A309: नव्या परिच्छेदातील सिद्धता-नावे योग्य ग्रीक मुद्रण-संकेताने पुनर्स्थापित केली आहेत.\end{quote}
\begin{quote}\small\textbf{OLP-0678 — स्रोतनोंद.} SOL6-A306: मुख्य शाखेतील नवे गृहीतक नवी खूण वापरते. असत्य निष्कर्षासाठी रिकामा उत्तरपक्ष जपला आहे. उपसूत्राच्या गुणधर्मात उघडी गृहीतके, संख्यापकांची उदाहरणे आणि असत्य यांचा नेमका अपवाद स्पष्ट केला आहे.\end{quote}
\begin{quote}\small\textbf{OLP-0678 — स्रोतनोंद.} SOL6-A308: संकुचनाचा नियम योग्य मुद्रण-संकेताने दिला आहे. असत्य आधारासाठी मूळ रिकामा उत्तरपक्ष आणि दुर्बलीकरणाने मिळणारा आधार यांना स्वतंत्र सिद्धता-नावे दिली आहेत.\end{quote}
\begin{quote}\small\textbf{OLP-0678 — स्रोतनोंद.} SOL6-A310: नव्या अटींतील सूत्रे, सिद्धता-नावे आणि नियमांची नावे योग्य TeX संकेतांनी पुनर्स्थापित केली आहेत.\end{quote}
\begin{quote}\small\textbf{OLP-0679 — स्रोतनोंद.} SOL6-A311: सिद्धता न सापडता होणारे अपयशी थांबणे वेगळे केले आहे. न्याय्यतेत अणुरूप नसलेल्या आस्थितींचे निर्देशांक वापरले आहेत; वाक्यांच्या क्रमवर्तींसाठीची बंद-पद व्याप्ती स्पष्ट आहे.\end{quote}
\begin{quote}\small\textbf{OLP-0680 — स्रोतनोंद.} SOL6-A314: पुढील पद-प्रतिमानावर आधारित संपूर्णतेची चर्चा वाक्यांच्या क्रमवर्तींसाठी आहे, ही व्याप्ती परिचयात स्पष्ट केली आहे.\end{quote}
\begin{quote}\small\textbf{OLP-0683 — स्रोतनोंद.} SOL6-A312: सर्व योग्य पदे वापरून झाल्यावर संबंधित अरिक्त संचातील स्थिरांक पुन्हा वापरायचा; संचाबाहेरचा स्थिरांक आणून पुढील आयगेन-चर अट मोडायची नाही.\end{quote}
\begin{quote}\small\textbf{OLP-0684 — स्रोतनोंद.} SOL6-A316: सत्य-चिन्हांकित असत्यानेही शाखा बंद होते. पुन्हा वापरायच्या संख्यापकांसाठी अरिक्त स्थिरांक-संच, नवे उदाहरण नसतानाची कृती आणि संतृप्त उघड्या शाखेची अट स्पष्ट केली आहे.\end{quote}
\begin{quote}\small\textbf{OLP-0687 — स्रोतनोंद.} SOL6-A318: कच्च्या पद-विन्यासाऐवजी योग्य प्रकाराच्या पदांची व्याप्ती स्पष्ट केला. अपूर्ण रेडेक्स-मोजणी युक्तिवादाऐवजी आधीच्या नैसर्गिक निगमन सामान्यरूपीकरण प्रमेयाचा आणि सर्व अनुरूप रूपांतरणांचा वापर केला; बीटा-प्रबळ सामान्यीकरण हा स्वतंत्र अधिक बलवान परिणाम आहे.\end{quote}
\begin{quote}\small\textbf{OLP-0687 — स्रोतनोंद.} SOL6-A326: बेरीज-प्रकाराच्या उदाहरणांत तक्त्यातील injection चिन्ह वापरले आणि त्याची वरची खूण दुसरा प्रकार सांगते हे स्पष्ट केले.\end{quote}
\begin{quote}\small\textbf{OLP-0688 — स्रोतनोंद.} SOL6-A323: खुणित संदर्भ सुसंगत संच आहे आणि पदांची रचना विगमनाने आहे, हे आधीच्या नियमांशी जुळवले.\end{quote}
\begin{quote}\small\textbf{OLP-0689 — स्रोतनोंद.} SOL6-A324: येथे दिलेल्या पद-नियमांसाठी सिद्धतेची विधानीय व्याप्ती स्पष्ट केला; संख्यापकांचे पद-नियम दिले असल्याचा दावा नाही.\end{quote}
\begin{quote}\small\textbf{OLP-0691 — स्रोतनोंद.} SOL6-A321: मुख्य बीटा-नियम आणि पूर्ण सामान्यरूपीकरणातील क्रमपरिवर्तन व असत्य-सरलीकरण वेगळे केले; आत हलवलेल्या पदांचे मुक्त चल सुरक्षित ठेवले.\end{quote}
\begin{quote}\small\textbf{OLP-0693 — स्रोतनोंद.} SOL6-A320: संदर्भात खुणित प्रकार असल्यावर क्रमवर्ती स्वयंसिद्धक मिळते, अशी नियमाची दिशा दुरुस्त केली.\end{quote}
\begin{quote}\small\textbf{OLP-0696 — स्रोतनोंद.} SOL6-A322: प्रकार-संरक्षण संदर्भ वाढवता येणाऱ्या प्रकार-पद्धतीत आहे; अचूक वापरलेली गृहीतके घटू शकतात. क्रमपरिवर्तनाला बीटा-पायरी मानलेले नाही.\end{quote}
\begin{quote}\small\textbf{OLP-0696 — स्रोतनोंद.} SOL6-A327: पद-रूपांतरणे या शब्दाशी अनुरूप मराठी लिंग-सुसंगती दुरुस्त केली.\end{quote}
\begin{quote}\small\textbf{OLP-0697 — स्रोतनोंद.} SOL6-A319: प्रत्येक पदाला प्रकार असतो हा असत्य दावा काढला: योग्य प्रकार मिळाल्यास तो एकमेव असतो. सर्व संदर्भ-खुणांची सुसंगती आणि बद्ध चलांची सुरक्षित नावे स्पष्ट केली.\end{quote}
\begin{quote}\small\textbf{OLP-0698 — स्रोतनोंद.} SOL6-A328: दुर्बलीकरणापूर्वी आयगेन स्थिरांकांची सुरक्षित नावे घेतली; उंची-शून्य प्रसंगात डावीकडील असत्य-स्वयंसिद्धकही समाविष्ट केले.\end{quote}
\begin{quote}\small\textbf{OLP-0699 — स्रोतनोंद.} SOL6-A329: दोन-नियमांच्या विधानात असत्य स्थिरांकाचा अपवाद स्पष्ट केला; दुर्बलीकरणाची आवश्यकता दाखवणाऱ्या उदाहरणाला योग्य अणुरूप व नियम-संच scope दिला.\end{quote}
\begin{quote}\small\textbf{OLP-0700 — स्रोतनोंद.} SOL6-A330: दोन्ही क्रमवर्ती तक्त्यांतील असत्याचे स्वतंत्र स्वयंसिद्धक प्रस्तावनेतही स्पष्ट केले; बहुसंचाचे सर्व अर्थ जपले.\end{quote}
\begin{quote}\small\textbf{OLP-0701 — स्रोतनोंद.} SOL6-A331: व्युत्क्रमण व आकुंचनाच्या आधारप्रसंगांत असत्याचे स्वयंसिद्धक समाविष्ट केले; प्रतिस्थापनापूर्वी अंतर्गत आयगेन स्थिरांक सुरक्षित ठेवले. उंची आणि आधारविधानांची संख्या वेगळ्या अक्षरांनी लिहिली.\end{quote}
\begin{quote}\small\textbf{OLP-0702 — स्रोतनोंद.} SOL6-A332: संयोग-परिचयाच्या निष्कर्षात दोन पूर्वांग प्रती ठेवून पुढील आकुंचन खरे केले; खोलीच्या आणि स्वयंसिद्धकाच्या दाव्यांचा योग्य scope व निष्फळ संख्यापकाचा अपवाद स्पष्ट केला.\end{quote}
\begin{quote}\small\textbf{OLP-0703 — स्रोतनोंद.} SOL6-A313: क्रमवर्ती संख्यापक-नियमांच्या बंद पदे वापरण्याच्या संकेताशी आदेशनाची अट जुळवली आहे; आधीची नियमितीकरण आणि विगमनातील घट जपली आहे.\end{quote}
\begin{quote}\small\textbf{OLP-0705 — स्रोतनोंद.} SOL6-A333: लघुतम पद्धतीच्या व्याख्येत तक्त्यातील नेमके असत्याचे स्वयंसिद्धक आणि उजवे दुर्बलीकरण वगळले आहे.\end{quote}
\begin{quote}\small\textbf{OLP-0708 — स्रोतनोंद.} SOL6-A334: G1m ची व्याख्या G1i च्या तक्त्याशी जुळवली; उजवे दुर्बलीकरण आणि असत्याचे स्वयंसिद्धक दोन्ही वगळावे लागतात.\end{quote}
\begin{quote}\small\textbf{OLP-0710 — स्रोतनोंद.} SOL6-A335: तक्त्यात परिभाषित न केलेल्या mG1i आणि mG1m बद्दलचे मूळ विधान वगळले; दिलेले mG3i नियम जपले.\end{quote}
\begin{quote}\small\textbf{OLP-0711 — स्रोतनोंद.} SOL6-A336: सिद्धतेच्या व्याख्येतील विगमनाची संज्ञा सुसंगत केली आणि उघडे अवतरणचिन्ह बंद केले.\end{quote}
\begin{quote}\small\textbf{OLP-0713 — स्रोतनोंद.} SOL6-A337: G1c ते G3c भाषांतरातील उलट दिशेचा संदर्भ बदलला; दोन्ही दिशांत मुख्य संख्यापकीय सूत्र जपणाऱ्या नियमांचे आवश्यक दुर्बलीकरण किंवा आकुंचन स्पष्ट केले.\end{quote}


या नोंदी अनुवादकाने वेगळ्या जोडल्या आहेत; त्या मूळ मजकुराचा भाग नाहीत.
मूळ इंग्रजी स्रोताचे बाइट्स बदललेले नाहीत. मराठी TeX फाइलांतील नोंदवलेल्या स्रोतदुरुस्त्या या संपादकीय नोंदी आणि निर्णयनोंदवहीत स्पष्ट केल्या आहेत.
\begin{enumerate}
\item विभाग \ref{sfr:rel:set:sec} मधील $K=L\cup I$ आणि $H=G\cup I$ येथे
$I$ चे स्वतंत्र नामकरण राहिले आहे. आधीच्या कर्णाच्या चर्चेनुसार त्याचा अर्थ
$\Id{\Nat}$ हा एकरूपता संबंध असा घ्यावा.
\item विभाग \ref{sfr:rel:tre:sec} मधील शाखेच्या व्याख्येत $z\in X\setminus B$
असे छापले आहे. वृक्षाचा आधारसंच $A$ आहे; येथे $X$ ऐवजी $A$ अभिप्रेत दिसतो.
\item त्याच विभागातील आरंभीच्या खंडांनुसार बंद असणाऱ्या उपवृक्षाच्या उदाहरणात
$A$ अरिक्त असणे आवश्यक आहे: आधीच्या व्याख्येनुसार वृक्षाला मूळ असते,
म्हणून रिक्त संच त्या व्याख्येत बसत नाही.
\item फलन प्रकरणातील OLFUN-001 ते OLFUN-005 या पाच
गोठवलेल्या-स्रोत निष्कर्षांची दुरुस्ती संबंधित परिच्छेदांलगत स्वतंत्र
``स्रोतदुरुस्ती'' नोंदींमध्ये दिली आहे. मूळ इंग्रजी बाइट्स बदललेले नाहीत.
\item संचांच्या आकारमानाच्या प्रकरणातील OLSIZ-001 ते OLSIZ-010
या दहा सामायिक गोठवलेल्या-स्रोत निष्कर्षांची आणि MRSIZ-001 ते MRSIZ-004
या चार स्थानिक निरीक्षणांची नोंद संबंधित परिच्छेदांलगत दिली आहे. मूळ
इंग्रजी बाइट्स बदललेले नाहीत.
\item सामायिक OLSIZ-011 सूचना अचूक बाइट-पुनर्तपासणीनंतर चुकीची ठरून
मागे घेण्यात आली; तिच्यावर आधारित कोणतीही मराठी दुरुस्ती केलेली नाही.
\item अंकगणितीकरण प्रकरणातील MRARITH-001 ते MRARITH-012 या बारा
गोठवलेल्या-स्रोत निरीक्षणांतील दुरुस्त्या किंवा स्पष्ट खुलासे संबंधित
परिच्छेदांलगत स्वतंत्र ``स्रोतदुरुस्ती'' नोंदींमध्ये दिले आहेत. मूळ
इंग्रजी बाइट्स बदललेले नाहीत.
\item अनंत संच प्रकरणातील MRINF-001 ते MRINF-003 या तीन स्रोत-निरीक्षणांची
नोंद ठेवली आहे. MRINF-001 मधील व्याकरणदोषाचा अभिप्रेत अर्थ थेट दिला आहे;
MRINF-002 मधील अप्रकट आधारसंचाचे गृहीतक तज्ज्ञ-पुनरावलोकनासाठी राखले आहे;
MRINF-003 मधील स्रोतदोषाचा दावा पुनर्तपासणीनंतर चुकीचा ठरून मागे घेतला आहे.
अंतःस्थ \texttt{cardeq} रचना वैध आहे; मराठी निष्कर्ष सममूल्य केंद्रित
मांडणी म्हणून जतन केला आहे. मूळ इंग्रजी बाइट्स बदललेले नाहीत.
\item विधानीय तर्कशास्त्र प्रकरणातील MRPL-001 आणि MRPL-002 या दोन
गोठवलेल्या-स्रोत चिन्हदोषांच्या दुरुस्त्या संबंधित परिच्छेदांलगत स्वतंत्र
``स्रोतदुरुस्ती'' नोंदींमध्ये दिल्या आहेत. मूळ इंग्रजी बाइट्स बदललेले नाहीत.
\item \textbf{MRPRF-001}: गोठवलेल्या इंग्रजी टॅब्लो उदाहरणातील
शेवटच्या दोन पायऱ्या $\sFmla{\True}{\varphi \land \psi}$ चे घटक उलगडतात,
पण त्यांची नियम-खूण चुकून $\TRule{\True}{\lif}$ अशी आहे. या वाचकात ती
अर्थानुसार $\TRule{\True}{\land}$ केली आहे; संरेखित इंग्रजी आणि मराठी
स्रोत-बाइट्स बदललेले नाहीत.

\item \textbf{OLSEQ-001}: OLP-0075 मधील चार पायऱ्यांत पूर्वपक्षातील
सूत्रांची अदलाबदल असूनही नियम-खूण उजवी अदलाबदल अशी आहे. या वाचकात त्या
चार खुणा डावी अदलाबदल अशा केल्या आहेत; संरेखित स्रोतांत मूळ रूप जतन आहे.
\item \textbf{OLSEQ-002}: त्याच उदाहरणातील दोन कथनसूत्रांत दुसऱ्या
सूत्रापूर्वीचे नकारचिन्ह गाळले आहे. लगतचे ध्येय, सममितीचे स्पष्टीकरण आणि
सिद्धता-वृक्ष यांनुसार या वाचकात ती दोन नकारचिन्हे पुनःस्थापित केली आहेत.
\item \textbf{OLSEQ-003}: OLP-0072 च्या शेवटच्या परिच्छेदातील ``पदावर
निर्बंध नाही'' याचा अर्थ पद बंद असण्याव्यतिरिक्त आणखी निर्बंध नाही असा आहे;
मराठी मजकूर ही आधी स्पष्ट केलेली अट सांगतो.
\item \textbf{OLND-001}: OLP-0087 मध्येही संख्यापक-नियमांसाठी पद बंद
असण्याची उभी अट शेवटच्या स्पष्टीकरणात स्पष्ट केली आहे.
\item \textbf{OLND-002}: OLP-0087 मधील आयगेन चलाचा एकत्रित सारांश सर्व
आधारविधानांतून चल वगळल्यासारखा वाचतो; मराठी सारांश प्रत्येक नियमाची योग्य
आधारविधान, निष्कर्ष आणि अमुक्त गृहीतकांची अट वेगवेगळी राखतो.
\item \textbf{OLND-003}: OLP-0089 मध्ये आकृतीतील उजवी शाखा चुकून डावी
म्हटली आहे; वाचकात आकृती आणि पुढील परिच्छेदाशी सुसंगत ``उजवीकडील'' आहे.
\item \textbf{OLND-004}: OLP-0089 मधील एकच अनुमान कधी नकार-विलोपन तर
कधी असत्यता-प्रवेशन असे खूणांकित आहे. अनुमान अर्थतः समान असल्याने मूळची
पर्यायी खूण जतन केली आहे.
\item \textbf{OLND-005}: OLP-0090 मधील आयगेन-चल तपासणीत मुख्य
आधारविधानातील नकार गाळला आहे आणि नियमाची व्याप्ती अपुरी सांगितली आहे;
मराठी मजकूर योग्य नकार व तात्पुरत्या गृहीतकाचा अपवाद स्पष्ट करतो.
\item \textbf{OLND-006}: OLP-0090 मधील सार्वत्रिक विलोपनासाठी ``कोणतेही
पद'' म्हणजे उभ्या नियमानुसार कोणतेही बंद पद; मराठी स्पष्टीकरण ती अट राखते.
\item \textbf{OLND-007}: OLP-0095 च्या एका सरावातील नियम-खूण
\verb|\Elim{\forall}| ही प्रकरणात इतरत्र असलेल्या
\verb|\Elim{\lforall}| शी दृश्यतः समतुल्य आहे. संरेखित रूप जतन केले आहे.

\item \textbf{OLTAB-001--OLTAB-004}: टॅब्लो प्रकरणातील आधीच्या
उदाहरणांमध्ये विषय, बंद पदांची अट आणि एका ओळ-संदर्भाशी संबंधित स्रोतदोष
दुरुस्त किंवा स्पष्ट केले आहेत; संरेखित इंग्रजी आणि मराठी स्रोत-बाइट्स जतन आहेत.
\item \textbf{OLTAB-005--OLTAB-009}: टॅब्लो सिद्धतायोग्यता आणि सुसंगतता
विभागातील संच-लेखन, नियम-पूर्वपद, ओळ-संदर्भ आणि शाखा-वर्णनातील निश्चित
स्रोतदोषांचे मराठी रूपात मर्यादित निराकरण केले आहे; मूळ इंग्रजी बाइट्स जतन आहेत.
\item \textbf{OLTAB-010}: विधानीय परिणामांच्या आठ टॅब्लो नोड्समधील
चिन्हांकित-सूत्र मॅक्रोची गहाळ युक्ती-सिमा पुनर्स्थापित केली आहे; संरेखित
स्रोतांत गोठवलेले रूप जतन आहे.
\item \textbf{OLTAB-011}: टॅब्लो निर्दोषता सिद्धतेतील सार्वत्रिक नियमांच्या
पुनरावृत्त B/A रूपबंध-विसंगतीत निष्कर्ष आणि अर्थविषयक पायऱ्यांशी सुसंगत
A-रूप वापरले आहे; गोठवलेले रूप संरेखित स्रोतांत जतन आहे.
\item \textbf{OLTAB-012}: एकरूपता-टॅब्लोच्या सममिती आणि संक्रमकता
स्पष्टीकरणांत लगतच्या वृक्षांशी विसंगत असलेली दोन कथनसूत्रे दुरुस्त केली आहेत;
औपचारिक वृक्ष अपरिवर्तित आहेत.

\item \textbf{OLAXD-001--OLAXD-009}: स्वयंसिद्धकीय निष्पत्ती प्रकरणातील
कंस, अंतिम निष्कर्ष, स्वयंसिद्धक-संदर्भ, संख्यापक-नियम आणि जुनी
फाइल-माहिती यांतील नऊ निश्चित स्रोतदोष मराठी संरेखित स्रोतांत मर्यादितपणे
दुरुस्त केले आहेत; गोठवलेले इंग्रजी बाइट्स जतन आहेत.
\item \textbf{OLCOM-001--OLCOM-006}: संपूर्णता प्रकरणातील सूत्र-युक्तिवाद,
चल-सुसंगती, बंद-पद अट, विरामचिन्हे, प्रतिनिधी आणि टॅग-संदर्भ यांतील सहा
निश्चित स्रोतदोषांची दुरुस्ती किंवा स्पष्टता मराठी स्रोतांत नोंदवली आहे.
 \item \textbf{OLFOL-001--OLFOL-040}: प्रथम-क्रम तर्कशास्त्र आणि प्रतिमान
उपपत्तीच्या पुढील प्रकरणांतील एकोणचाळीस सूत्रात्मक, निर्देशांकविषयक,
अर्थविषयक आणि सिद्धताविषयक स्रोतदोष किंवा अपूर्ण पायऱ्या मर्यादितपणे
दुरुस्त किंवा स्पष्ट केल्या आहेत. OLFOL-009 मधील अतिरिक्त आकुंचित कंसाची
सूचना पुनर्तपासणीनंतर चुकीची ठरून मागे घेतली आहे; गोठवलेली रचना आधीपासून
संतुलित आहे. प्रत्येक नोंदीचा अचूक स्रोत-स्थान, पर्याय आणि
पुनरावलोकन-प्रश्न सोबतच्या तज्ज्ञ-पुनरावलोकन नोंदींत आहे.

\item \textbf{OLCMP-001--OLCMP-024}:\par संगणनक्षमता प्रकरणांतील औपचारिक
मांडणी, चिन्हांकने आणि सिद्धतापायऱ्यांबद्दलच्या स्रोत-निरीक्षणांची संरेखित
स्रोतांत स्वतंत्र नोंद आहे. मूळ इंग्रजी बाइट्स जतन आहेत; मराठी वाचक-मजकूर
नोंदवलेल्या मर्यादित स्पष्टता आणि दुरुस्त्या वापरतो.

\item \textbf{OLTUR-001--OLTUR-007}: ट्यूरिंग यंत्रांच्या उदाहरणांतील
आरंभीची अवस्था, फितीचे टोक, रिकामे आदान, प्रदान आणि यंत्र-संयोजन
यांबद्दलची सात स्रोत-निरीक्षणे संरेखित स्रोतांत नोंदवली आहेत. मूळ
इंग्रजी बाइट्स जतन आहेत; मराठी वाचक-मजकूर मर्यादित स्पष्टता वापरतो.

\item \textbf{OLTUR-008--OLTUR-022}: अनिर्णेयता आणि सांत प्रतिरूपांच्या
प्रकरणातील स्रोत-निरीक्षणे संबंधित मराठी मजकुरालगत दिली आहेत. विशेषतः
यंत्र-चालक्रमाचे निरूपण, शेवटच्या पायरीचे स्थितिवर्णन आणि सांत प्रतिरूपाची
मर्यादा यांतील दुरुस्त्या मूळ इंग्रजी मजकूर न बदलता नोंदवल्या आहेत.

\item \textbf{OLINC-001--OLINC-005}: अपूर्णतेच्या परिचयातील उद्दिष्टांची
क्रमवारी, प्रमेयांच्या गृहीतकांची व्याप्ती आणि संबंध-निरूपणातील उपपत्तीचा
निर्देश यांसंबंधीची निरीक्षणे नोंदवली आहेत. विकर्ण सिद्धतेतील दोन
वगळलेले निर्देशांक आणि विस्तारित उपपत्तींच्या सुसंगततेची अट मराठीत
स्पष्ट केली आहे; मूळ इंग्रजी मजकूर बदललेला नाही.

\item \textbf{OLINC-006--OLINC-043}: विन्यासमीमांसेच्या
अंकगणितीकरणात आणि $Q$ मधील निरूपणीयतेत मूळ इंग्रजीतील
सूत्रे, संख्यांकांचे निर्देश, सिद्धतापायऱ्या व गृहीतके
यांबद्दलची स्रोत-निरीक्षणे स्वतंत्र नोंदवली आहेत.
संबंधित मराठी विभागांत त्यांची मर्यादित दुरुस्ती
किंवा स्पष्टता दिली आहे; मूळ इंग्रजी बाइट्स बदललेले नाहीत.

\item \textbf{OLINC-044--OLINC-054}: उपपत्ती आणि संगणनक्षमता या
प्रकरणातील मूळ इंग्रजीतील सूत्रे, गृहीतकांची व्याप्ती आणि
प्रतिमानातील सत्य यांबद्दलची स्रोत-निरीक्षणे स्वतंत्र नोंदवली
आहेत. संबंधित मराठी विभागांत त्यांची मर्यादित दुरुस्ती
किंवा स्पष्टता दिली आहे; मूळ इंग्रजी बाइट्स बदललेले नाहीत.

\item \textbf{OLINC-055--OLINC-059}: अपूर्णता आणि सिद्धतायोग्यता
या प्रकरणातील मूळ इंग्रजीत वस्तुभाषेतील सिद्धतायोग्यता
विधेय, गोडेल-वाक्याचा संकेतनांक आणि सुसंगततेचे सूत्र
यांमध्ये आढळलेल्या विसंगती स्वतंत्र नोंदवल्या आहेत.
मराठी आवृत्तीत संबंधित सूत्रांची मर्यादित दुरुस्ती
केली आहे; मूळ इंग्रजी बाइट्स बदललेले नाहीत.

\item \textbf{OLINC-060--OLINC-061}: द्वितीय-क्रम चिन्हार्थमीमांसेच्या
मूळ इंग्रजीत संक्रमणाविषयी वापरलेले $R^*$ हे चिन्ह आधीच्या प्रकरणातील
व्याख्येशी विसंगत आहे; मराठी आवृत्तीत सकारात्मक संक्रमणासाठी $R^+$
ठेवले आहे. मर्यादित संचाच्या गणनेविषयीच्या सिद्धतेत शेवटच्या घटकाचे
फलनमूल्य स्पष्ट केले आहे. दोन्ही मुद्दे स्वतंत्र स्रोत-नोंदींत नोंदले आहेत.

\item \textbf{OLINC-062--OLINC-065}: द्वितीय-क्रम अंकगणितात
बेरीज परिभाषित करणाऱ्या सूत्रातील बांधलेला चल सुसंगत केला आहे.
स्वयंसिद्धकीकरणाच्या सिद्धतेतील तुटलेली गणित-सीमा दुरुस्त केली आहे.
संहततेच्या प्रमेयाला स्वतंत्र खूण दिली आहे आणि सांत उपसंचावरील
कमाल बंधाचा निर्देश त्या सांत उपसंचाकडेच केला आहे.
मूळ इंग्रजी बाइट्स बदललेले नाहीत.

\item \textbf{OLINC-066--OLINC-072}: द्वितीय-क्रम
संचसिद्धान्तावरील मूळ प्रकरण स्वतःच निष्कर्षांत दोष
असू शकतात अशी सूचना देते. उपसंचांच्या अनंतता,
गणनीयता आणि अलेफ-एक यांसाठीची काही मूळ सूत्रे
अपेक्षित गुणधर्म व्यक्त करत नव्हती. मराठी सूत्रांत
\textbf{SOL6-A096--SOL6-A098} या दुरुस्त्या केल्या आहेत:
एकास-एक अट संबंधित उपसंचापुरती घेतली आहे; अनंतता,
गणनीयता आणि अलेफ-एक यांच्या अटी दुरुस्त केल्या आहेत;
सांतत्याचे सांकेतीकरण आणि पूर्ण प्रांताचे आकारमान
त्यांनुसार दुरुस्त केले आहे. निवडीसह संचसिद्धान्ताचा
आधी स्पष्ट केलेला संदर्भ लागू आहे. मूळ इंग्रजी बाइट्स
आणि ऐतिहासिक दोषनोंदी जतन आहेत.

\item \textbf{OLINC-073--OLINC-081}: लॅम्डा कलनाच्या
परिचयात एका विभागाचा जुना प्रकरण-संकेत, करीकरणाच्या
अंतिम आदेशनातील चुकलेले पद, फलनाच्या स्थानसंख्येतील
विसंगती, आणि आदिम पुनरावर्तन व लघुतमीकरणाच्या
पुराव्यांतील स्थानिक त्रुटी नोंदवल्या आहेत. मराठी
मजकुरातील नेमक्या दुरुस्त्या व स्पष्टीकरणे स्रोत-नोंदीत
दर्शवली आहेत; मूळ इंग्रजी बाइट्स बदललेले नाहीत.

\item \textbf{OLINC-082--OLINC-102}: लॅम्डा पदांच्या
विन्यासमीमांसेतील एकमेव वाचनीयता, मुक्त चर, आदेशन,
अल्फा-परिवर्तन, सममूल्यतावर्ग आणि ईटा नियम यांमधील
गोठवलेल्या इंग्रजी स्रोताच्या स्थानिक त्रुटी स्वतंत्र
स्रोत-नोंदीत नोंदवल्या आहेत. गणिती अर्थ स्पष्ट करणाऱ्या
मराठी दुरुस्त्या तिथे नेमक्या ओळखता येतात; मूळ इंग्रजी
बाइट्स बदललेले नाहीत. अल्फा-परिवर्तनाचा सरावातील
दोन सारख्या प्रश्नजोड्या स्रोताप्रमाणे राखल्या आहेत.

\item \textbf{OLINC-103--OLINC-113}: चर्च–रॉसर प्रकरणातील
गोठवलेल्या इंग्रजी स्रोताच्या स्थानिक सूत्रचिन्ह, नियम,
चल आणि सिद्धता-अंतराच्या नोंदी स्वतंत्र स्रोत-नोंदीत
आहेत. दुरुस्त्या मराठी आवृत्तीत स्पष्ट केल्या आहेत;
बीटा–ईटा पूर्ण विकासातील आच्छादन-प्रकरणांची उणीव
\textbf{SOL6-A126} मध्ये दुरुस्त केली आहे: मुख्य
दावा आणि साहाय्यक दावा एकत्रित रचनात्मक विगमनाने
सिद्ध केले आहेत. मूळ इंग्रजी बाइट्स
बदललेले नाहीत.

\item \textbf{OLINC-114--OLINC-128}: लॅम्डा-परिभाष्यता
प्रकरणातील सूत्रे, आदानांची संख्या, न्यूनीकरणाच्या पायऱ्या,
पुनरावर्तन आणि स्थिरबिंदू संयोजकांच्या नोंदी स्वतंत्र
स्रोत-नोंदीत आहेत. जिथे गणिती अर्थ स्पष्ट आहे तिथे
मराठीतील स्थानिक दुरुस्त्या नोंदवल्या आहेत.
क्रमगुणिताच्या स्व-आदेशानंतरच्या सूत्रातील बाह्य
लॅम्डाच्या व्याप्तीचा दोष \textbf{SOL6-A132} मध्ये
दुरुस्त केला आहे: दोन्ही शाखा एकाच बाह्य शरीरात
आहेत; उरलेला Fac स्वसंदर्भ जपला आहे.
मूळ इंग्रजी बाइट्स बदललेले नाहीत.

\item \textbf{OLINC-129--OLINC-131}: बहुमूल्य तर्कशास्त्रातील
अभिजात तुलनेच्या गोठवलेल्या इंग्रजी प्रमेयाची व्याप्ती
अतिविस्तृत आहे आणि अंतिम सिद्धतेत पूर्तीऐवजी निष्पन्नतेची
चिन्हे चुकून वापरली आहेत. मराठी आवृत्तीत विधान सामाईक
चार-संयोजक खंडापुरते स्पष्ट केले आहे, सर्व विधानीय चलांसाठी
अभिजात मूल्यांची अट नमूद केली आहे आणि प्रत्युदाहरणातील
पूर्तीची चिन्हे दुरुस्त केली आहेत. मूळ इंग्रजी बाइट्स बदललेले नाहीत.

\item \textbf{OLINC-132--OLINC-137}: त्रिमूल्य तर्कशास्त्रांच्या
गोठवलेल्या इंग्रजी मजकुरात संधीच्या समानतेतील पुनरावृत्ती,
एका सरावातील अतिरिक्त कंस, विगमनाच्या आधारातील अटीविना
दिलेली समानता आणि संधीच्या पुराव्यात दुसऱ्या घटकाऐवजी
पहिलाच घटक पुन्हा लिहिलेला आहे. मराठी आवृत्तीत स्पष्ट
गणिती दुरुस्त्या नोंदवून केल्या आहेत. R-Mingle च्या
मॅट्रिक्समधील चुकीचा अभिव्यंजन-तक्ता आणि अतिरिक्त
असत्य-स्थिरांकाची अस्पष्ट व्याप्ती \textbf{SOL6-A142}
मध्ये दुरुस्त केली आहे. योग्य RM3 तक्ता आणि स्पष्ट
असत्य-स्थिरांक विस्तार दिला आहे. मूळ इंग्रजी
बाइट्स बदललेले नाहीत.

\item \textbf{OLINC-138--OLINC-142}: अनंतमूल्य तर्कशास्त्रांच्या
गोठवलेल्या इंग्रजी मजकुरात $V_\infty$ च्या संचवर्णनातील
शून्य हर, $V_m$ च्या संचवर्णनातील गणनाभेद आणि G\"odel
सत्यता-फलनातील चुकीची गणिती सीमाचिन्हे दुरुस्त केली आहेत.
अनंत गृहीतकांसाठी ससीम-मूल्य निष्पन्नतेवरून अनंतमूल्य
निष्पन्नता सर्वसाधारणपणे मिळत नाही; म्हणून केवळ ही
प्रतिदिशा ससीम गृहीतकांसाठी मांडली आहे. मूळ
इंग्रजी बाइट्स बदललेले नाहीत.

\item \textbf{OLINC-143--OLINC-145}: क्रमवर्ती कलनाच्या
गोठवलेल्या इंग्रजी मजकुरात डाव्या बाजूच्या सूत्रांची संख्या
क्रमवर्ती आणि संबंधित संधी यांत विसंगत आहे; एका
मूल्यनात मूल्यांकनाचा निर्देशांक गाळला आहे; आणि
$n$-बाजूंच्या क्रमवर्तीच्या प्रत्येक घटकाऐवजी फक्त
पहिल्या घटकाचा निर्देशांक लिहिला आहे. मराठी आवृत्तीत
या तीन गणिती नोंदी दुरुस्त केल्या आहेत. मूळ इंग्रजी
बाइट्स बदललेले नाहीत.

\item \textbf{OLINC-146--OLINC-153}: सामान्य मोडल
तर्कशास्त्रातील जगनिरपेक्ष अनिवार्यता, परिकर्मीची
कंसजोडी, द्विपक्षीय अभिव्यंजनाचा निवडक संकेत,
पूर्तीतील जगाचा निर्देशांक, आदेशन-दृष्टांतांची विगमन
सिद्धता आणि प्रतिमानाचे क्रमित त्रिक या ठिकाणी
गोठवलेल्या इंग्रजी स्रोताच्या मर्यादित दुरुस्त्या
मराठी मजकुरात नोंदल्या आहेत. मूळ इंग्रजी बाइट्स
बदललेले नाहीत.

\item \textbf{OLINC-154--OLINC-164}: चौकट-स्तरावरील
अनुरूपता, प्रतिउदाहरणांतील प्राप्यता-अटी, पाच रूपबंधांची
व्यत्यास सिद्धता, संहततेच्या साखळीचे निर्देशांक आणि
मोडल सूत्रांचे मानक भाषांतर यांत गोठवलेल्या इंग्रजी
स्रोताच्या मर्यादित दुरुस्त्या मराठी मजकुरात नोंदल्या
आहेत. मूळ इंग्रजी बाइट्स बदललेले नाहीत.

\item \textbf{OLINC-165--OLINC-174}: मोडल प्रणालींची
सिद्धता, व्युत्पन्न नियम, प्रतिमानांद्वारे वेगळेपणा
आणि सुसंगततेचे पुरावे यांतील गोठवलेल्या इंग्रजी
स्रोताच्या मर्यादित दुरुस्त्या मराठी मजकुरात नोंदल्या
आहेत. मूळ इंग्रजी बाइट्स बदललेले नाहीत.

\item \textbf{OLINC-175--OLINC-180}: संपूर्ण सुसंगत संचांचे
पुरावे, लिंडेनबाउमच्या प्रगणनातील लांबीची अट,
मोडल निष्पत्तींतील निर्देशांक आणि \(\Sigma\) ची खूण,
तसेच सत्यता पूर्वप्रमेयातील सराव-टॅग यांतील
गोठवलेल्या इंग्रजी स्रोताच्या मर्यादित दुरुस्त्या
मराठी मजकुरात नोंदल्या आहेत. मूळ इंग्रजी बाइट्स
बदललेले नाहीत.

\item \textbf{OLINC-181--OLINC-190}: निस्यंदनांच्या
भागात तुल्यतावर्गांचे आकार, मूल्यांकनांचे प्रांत,
भागाकार-जगांच्या खुणा, \(S5\) साठी प्रतिमानांचा
वर्ग, प्रमेयाशी जुळणारा सिद्धतेचा क्रम आणि
युक्लिडी चित्रांतील आवश्यक स्व-बाण यांबाबत
गोठवलेल्या इंग्रजी स्रोताच्या मर्यादित दुरुस्त्या
मराठी मजकुरात नोंदल्या आहेत. मूळ इंग्रजी बाइट्स
बदललेले नाहीत.

\item \textbf{OLINC-191--OLINC-208}: मोडल टॅब्लोच्या
नियमांतील चिन्हे व पूर्वचिन्हे, निर्दोषता आणि
संपूर्णतेच्या पुराव्यांतील काही सूत्रे, तसेच
प्रतिदृष्टांत प्रतिमानांच्या उदाहरणांतील
चुकीच्या खुणा मराठीत दुरुस्त केल्या आहेत.
\textbf{OLINC-204} मध्ये नोंदलेली सांत
\(\Gamma\) वरून सर्वसाधारण \(\Gamma\) कडे
जाणाऱ्या स्रोत-पुराव्यातील उणीव
\textbf{SOL6-A180} मध्ये दुरुस्त केली आहे.
असांत प्रकरणासाठी सांत उपसंचांची सुसंगतता,
आधीचे लिंडेनबाउम आणि सत्यता पूर्वप्रमेय
वापरले आहेत. मूळ इंग्रजी बाइट्स बदललेले नाहीत.

\item \textbf{OLINC-209}: कालिक भाषेतील सूत्रनिर्मितीच्या
नियमात मूळ इंग्रजीत साधा \(F\) आला आहे; त्याच
विभागातील संकारकांची यादी आणि सत्यतेची व्याख्या
\(\Ftemp\) वापरतात. मराठी नियमात \(\Ftemp\)
ठेवून ही एक-चिन्हीय विसंगती दुरुस्त केली आहे.
मूळ इंग्रजी बाइट्स बदललेले नाहीत.

\item \textbf{OLINC-210}: मूळ इंग्रजी प्रस्तावनेत
\emph{Knowledge and Belief} या ग्रंथाच्या लेखकाचे
पहिले नाव ``Jaako'' असे आहे. विद्यापीठाच्या
ग्रंथसूचीनुसार योग्य स्पेलिंग ``Jaakko'' आहे;
मराठीत ``याक्को हिंटिक्का'' वापरले आहे.
\item \textbf{OLINC-211}: द्विअनुकरणाच्या ``पुढे'' आणि
``मागे'' या दोन्ही अटींमध्ये मूळ इंग्रजीत कर्त्यांचा
संच \(A\) आला आहे. प्रकरणात परिभाषित कर्तृचिन्हांचा
संच \(G\) असल्याने मराठी अटींमध्ये \(G\) वापरले आहे.
\item \textbf{OLINC-212}: सार्वजनिक घोषणा भाषेतील
संयोजकांच्या यादीत मूळ इंग्रजीत द्विसशर्त राहिला
आहे; त्याच विभागातील सूत्रनिर्मितीच्या नियमात तो
असल्याने मराठी यादीत समाविष्ट केला आहे.
\item \textbf{OLINC-213}: घोषणेच्या रिक्तपणे सत्य
असण्याच्या उदाहरणात मूळ इंग्रजीत उत्तरसूत्राचे
चिन्हांकन राहिले आहे. मराठीत व्याख्येनुसार
\([\varphi] \psi\) हे रूप वापरले आहे.
मूळ इंग्रजी बाइट्स बदललेले नाहीत.

\item \textbf{OLINC-214}: BHK उदाहरणातील फलनाच्या
फलप्रदेशात मूळ इंग्रजीत एका ठिकाणी \(C\) आले
आहे; सूत्रानुसार मराठीत \(\chi\) वापरले आहे.
\item \textbf{OLINC-215}: विकल्पयोगाच्या डाव्या
समावेशनात मूळ इंग्रजीत \(M_1\) चे प्रतिमान
\(\langle 1,M_2\rangle\) दिले आहे. योग्य
\(\langle 1,M_1\rangle\) मराठीत वापरले आहे.
\item \textbf{OLINC-216}: संयोगाच्या BHK स्पष्टीकरणात
मूळ इंग्रजीत \(\varphi_1\land \varphi_1\) आले आहे;
जोडीचा दुसरा घटक \(\varphi_2\) ची रचना असल्याने
मराठीत \(\varphi_1\land \varphi_2\) वापरले आहे.
मूळ इंग्रजी बाइट्स बदललेले नाहीत.

\item \textbf{OLINC-217}: मूळ इंग्रजीत
स्थानिक निष्पन्नतेच्या दाव्याच्या सिद्धतेत
सर्व जगांतील सत्यतेचे अधिक बळकट गृहीतक
वापरले आहे; मराठीत आधी स्थानिक दावा
सिद्ध करून मग त्यावरून सर्व जगांतील
दावा मिळवला आहे.
\item \textbf{OLINC-218}: मूळ इंग्रजीत एका
ठिकाणी मूळ संच \(X\) लाच संस्थिति म्हटले
आहे. व्याख्येनुसार \(\mathcal O\) ही संस्थिति
आणि तिच्यासह \(X\) हा सांस्थितिक अवकाश आहे;
मराठीत हा भेद ठेवला आहे.
\item \textbf{OLINC-219}: मूळ इंग्रजीत
निष्पन्नता आणि सशर्त विधानाच्या दोन
अटींमध्ये \(\subset\) चिन्ह आहे.
या चिन्हाचा अर्थ संकेतपद्धतीनुसार बदलतो;
येथे समानताही चालते हे स्पष्ट करण्यासाठी
मराठीत \(\subseteq\) वापरले आहे.
उचित उपसंचच अभिप्रेत असल्याचा जुना
दोष-दावा मागे घेतला आहे.
मूळ इंग्रजी बाइट्स बदललेले नाहीत.

\item \textbf{OLINC-220}: स्वयंसिद्धकीय निर्दोषतेच्या
गृहीतक-प्रसंगात मूळ इंग्रजीतील पूर्ति-सूत्राला
जादा \(\Gamma\) घटक आहे; मराठीत \(A_n\)
ची पूर्ति योग्य रूपात लिहिली आहे.
\item \textbf{OLINC-221--222}: नैसर्गिक
निगमनाच्या निर्दोषता-सिद्धतेत संधियोगाच्या
निष्कर्षात चुकीची सूत्रे, तसेच विकल्पयोगाच्या
प्रसंगात जगाचे स्थान आणि दोन एकघटक संच
चुकले आहेत; मराठीत संदर्भाशी जुळणारी
रूपे वापरली आहेत.
\item \textbf{OLINC-223--224}: लिंडेनबाउमच्या
रचनेतील निवडीसाठी परिमित आधीच्या
निर्देशांकांचा युक्तिवाद स्पष्ट केला आहे;
कॅनॉनिकल प्रतिमानातील एकस्वनिकतेसाठी
विगमनाचा विस्तार-लांबी हा निर्देशांक
स्पष्ट केला आहे.
\item \textbf{OLINC-225}: सत्यता पूर्वप्रमेयाच्या
प्रस्तावनेत प्रारंभीच्या \(\Delta\) ऐवजी
त्या जगाचा \(\Delta(\sigma)\) अभाज्य संच
उल्लेखिला आहे.
\item \textbf{OLINC-227}: आधीच्या मोडल
संपूर्णता प्रकरणात मूल्यांकन \(V^\Sigma\)
पुढील गद्य कंस चुकून गणितात गेला होता;
मराठीत कंस गणिताबाहेर ठेवला आहे.
\item \textbf{OLINC-226}: निर्णेयतेच्या
स्रोतसिद्धतेत जगांचे केवळ विधानचरांच्या
सत्यसंचांवर आधारित भागाकार सशर्त
सूत्राची सत्यता जपत नाही. मूळ दोन
जगे असंबद्ध असताना नव्या उपसंच-क्रमाने
त्यांत प्राप्यता निर्माण होऊ शकते.
स्रोतदोषाची नोंद जपली आहे; सध्याच्या
मराठी विभागात योग्य उपसूत्र-गाळणी व
तिची सिद्धता दिली आहे. प्रमेय बदललेले नाही.
मूळ इंग्रजी बाइट्स बदललेले नाहीत.

\item \textbf{OLINC-228--230}: अंतःप्रज्ञावादी
टॅब्लो प्रकरणातील चुकून आलेला मोडल उल्लेख,
सत्याचा प्राप्य जगांमध्ये टिकाव, आणि
सशर्त सूत्रासाठीच्या शाखांवरील सत्य-असत्य
चिन्हे मराठीत संदर्भाशी जुळवली आहेत.
\item \textbf{OLINC-231}: उदाहरणातील
दोन संधियोग-नियमांच्या कारणनिर्देशांत
ओळ ४ ऐवजी त्या सूत्राची ओळ ७ दिली आहे.
\item \textbf{OLINC-232--235}: निर्दोषतेच्या
स्रोतसिद्धतेतील प्रतिप्रतिमानाचे चिन्ह,
दोन नियमांचे निष्कर्ष-संच आणि अंतिम
तार्किक निष्पन्नता यांतील चुका
दुरुस्त केल्या आहेत. \textbf{OLINC-233}
मधील मनमानी पूर्वचिन्ह-संचाविषयीची
व्याख्यात्मक उणीव दुरुस्त केली आहे:
सर्व वापरलेल्या पूर्वचिन्ह-विस्तारांसाठी
प्राप्यता आवश्यक आहे. नवे साक्षी-पूर्वचिन्ह
किंवा त्याचा विस्तार आधी वापरलेला नसावा,
या अटीने जुने अर्थनिर्धारण जपले जाते.
इतर नियमांच्या सिद्धता मूळप्रमाणे
स्वाध्याय आहेत. मूळ इंग्रजी बाइट्स
बदललेले नाहीत.

\item \textbf{OLINC-236}: कठोर अभिव्यंजनाच्या
एक अवैध निष्पन्नता-दाव्यात मूळ इंग्रजीत
वास्तविक अभिव्यंजनाचे निषेधन आले आहे.
मराठीत कठोर अभिव्यंजनाचे निषेधन वापरले आहे.
\item \textbf{OLINC-237}: कठोर अभिव्यंजन
आवश्यकतेने सत्य किंवा असत्य ठरते या दाव्याला
S5 ची व्याप्ती मराठीत स्पष्ट केली आहे.
\item \textbf{OLINC-238}: स्टालनाकर आणि
लुईस यांना एकच अद्वितीय सर्वाधिक जवळचे जग
गृहीत धरणारी मूळ सत्यता-अट दुरुस्त केली आहे.
स्टालनाकरची निवडलेल्या जगाची अट आणि
लुईसची गोलक-अट वेगळ्या दिल्या आहेत.
लुईसच्या विश्लेषणात समसमान जवळची अनेक जगे
किंवा सर्वात जवळचे जग नसणेही शक्य आहे.
आधीची स्रोतनोंद हा भेद जपते. मूळ इंग्रजी
बाइट्स बदललेले नाहीत.

\item \textbf{OLINC-239}: गोलक प्रतिमानाच्या
सत्यता-अटीतील ‘‘एखाद्या गोलकात
अट पूर्ण झाली, तर त्यातील
सर्व लहान गोलकांतही होते’’
हा मूळ इंग्रजी दावा अति-विस्तृत
आहे. पूर्वांग सत्य असलेले जग
धारण करणाऱ्या लहान गोलकांपुरताच
तो मराठीत मर्यादित केला आहे.
\item \textbf{OLINC-240}: बाह्य अवकाशातील
आगकाडीच्या उदाहरणात मूळ गद्य
दोनदा ‘‘पेटवणे’’ म्हणते; उद्धृत
अनुमान आणि अपेक्षित प्रयोग
यांच्यानुसार मराठीत ‘‘ओढणे’’
वापरले आहे.
\item \textbf{OLINC-241}: संक्रमणकतेच्या
प्रतिमानात \(q\lif r\) सर्व
संबंधित जगांत सत्य असूनही
मूळ सूत्रात पूर्तिचिन्हावर
निषेध आहे. मराठीत तो
निषेध काढला आहे.
\item \textbf{OLINC-242}: प्रतिपरिवर्तनाच्या
प्रतिमानात मूळ इंग्रजीत
एका जगाभोवतीची गोलक-प्रणाली
संपूर्ण फलन \(O\) म्हणून,
आणि घोषित \(M_1\) प्रतिमान
पुढे \(M\) म्हणून आले आहे;
मराठीत \(O_w\) व \(M_1\)
अशी सुसंगत नावे दिली आहेत.
मूळ इंग्रजी बाइट्स बदललेले नाहीत.

\item \textbf{OLINC-243}: स्वसमावेशकतेवरील चर्चेत
रसेल संचाचे घटक वर्णन करताना मूळ इंग्रजीत
चुकून दुहेरी निषेध आला आहे. सूत्रानुसार
स्वतःचे सदस्य नसलेले संच अभिप्रेत आहेत;
मराठी गद्यात तेच स्पष्ट केले आहे.
\item \textbf{OLINC-244}: दुष्टचक्र तत्त्व
नाकारण्यानंतर मर्यादित संचरचनेची गरज येते,
असे मूळ गद्य उलटे म्हणते. उद्धृत तत्त्व
आणि पुढील युक्तिवादानुसार ते स्वीकारल्यावर
मर्यादा येते, असे मराठीत दुरुस्त केले आहे.
\end{enumerate}


\paragraph{क्रम-समरूपणाच्या पुराव्यातील मुद्रित चूक.} मूळ इंग्रजी आवृत्तीत क्रम-समरूपणाच्या एका पुराव्यात $B_{b_2}$ ला फलनाचा प्रांत म्हटले आहे; प्रदर्शित $f:A_{a_2}\to B_{b_2}$ नुसार तो सहप्रांत आहे. मराठी मजकुरात हा शब्द दुरुस्त केला आहे.

\paragraph{पुनरावर्तनातील आधारप्रकरण.} मूळ इंग्रजीत रिकाम्या प्रांताच्या फलनासाठी पुनरावर्तनाच्या व्याख्येत प्रकरण दिलेले नाही, तरी पुराव्यात त्या प्रकरणाचे मूल्य वापरले आहे. मराठी मांडणीत ते प्रकरण स्पष्ट केले आहे.
\paragraph{पायाभूततेच्या पुराव्यातील बंधित चल.} मूळ सूत्रात आधी निश्चित केलेल्या $B$ ऐवजी अपरिभाषित $b$ आले आहे; मराठी सूत्रात $B$ वापरले आहे.
\paragraph{प्रतांकावरील पुराव्यातील अट.} मूळ पुराव्यात $x\in V_\alpha$ मानून लगेच $x\notin V_\alpha$ असे म्हटले आहे. मराठी मजकुरात हा निष्कर्ष प्रतांक आणि टप्पा समान असल्याच्या प्रकरणापुरता स्पष्ट केला आहे.

\paragraph{प्रतिबिंबनाच्या पुराव्यातील सूत्रदुरुस्त्या.} मूळ इंग्रजीत साक्षीदार-टप्प्याच्या अटीत एक अतिरिक्त बंद कंस आणि $\bigcup_{m<\omega}S_n$ मध्ये निर्देशांक-विसंगती आहे; मराठी सूत्रे अनुक्रमे संतुलित कंस आणि $S_m$ वापरतात. प्रतिस्थापनाच्या प्रतिबिंबनाधारित पुराव्यातील प्रतिमा-संचाच्या अटीचा हरवलेला बंद कंसही पुरवला आहे.
\paragraph{दुर्बल प्रतिबिंबनाचे नाव आणि सापेक्षीकरण.} मूळ मजकूर दुर्बल प्रतिबिंबनाची व्याख्या करण्याआधी चुकून प्रतिस्थापनाचे दुर्बल रूप म्हणतो; मराठी मांडणीत प्रतिबिंबन म्हटले आहे. पुढील परिशिष्टात $(\exists N\,\psi(N))^M$ उलगडताना संक्रमकतेच्या अटीवरील चुकीचे $N$ अधिलेख $M$ केले आहे.
\paragraph{विभक्तीकरणाविषयी मूळ पुराव्यातील त्रुटी.} सांत विस्ताराच्या पुराव्यात आणि त्यानंतरच्या प्रश्नात प्रत्येक संक्रमक संच विभक्तीकरणाचे प्रत्येक सापेक्षीकृत उदाहरण पूर्ण करतो, असा दावा आहे. तो खोटा आहे: $X=3=\{0,1,2\}$ हा संक्रमक संच घ्या. $A=2\in X$ आणि $x=1$ निवडणारी अट घेतल्यास अपेक्षित $\{1\}\notin X$. मराठी पुराव्यात आता लघुतम संक्रमक व उपसंच-समावेशक प्रतिमान वापरले आहे; प्रश्नाच्या विधानातही आवश्यक उपसंच-समावेशकत्वाची अट दिली आहे. प्रश्नाची सिद्धता वाचकांसाठी ठेवली आहे.

\paragraph{क्रमसूचक बेरजेच्या रचनेतील तीन दुरुस्त्या.} मूळ इंग्रजीत संबंधाच्या घटकजोड्या $D$ मधून घेतल्यासारख्या लिहिल्या आहेत; त्या $D\times D$ मधून असतात. पुढे अनुक्रमक बेरीज उलगडताना दोन आधीच खूण केलेल्या घटकसंचांमध्ये पुन्हा विभक्त बेरीज लिहिली आहे; मराठीत साधे संघटन वापरले आहे. शून्याशी बेरीज करताना रिकाम्या घटकाऐवजी एकघटकी संच लिहिलेला आहे; तो रिक्त संच केला आहे.
\paragraph{प्रतांकाच्या सरावातील चिन्ह.} मूळ दुसऱ्या सरावात अपेक्षित प्रतांक आणि वरची मर्यादा यांच्यामध्ये समानतेचे चिन्ह सुटले आहे; मराठीत ते पुरवले आहे. एका निष्क्रिय टिप्पणीत विसंगत गणित-सीमाचिन्हही जोडले आहे.
\paragraph{क्रमसूचक गुणाकाराचे सीमा प्रकरण.} मूळ सार्वत्रिक सूत्रातील कठोर लघुतम ऊर्ध्वबंध शून्य डाव्या घटकासाठी चुकीचा परिणाम देतो. मराठी सूत्रात आधीच्या गुणाकारांचे संघटन वापरले आहे; त्यामुळे शून्याचे सीमा क्रमसूचक संख्येशी गुणाकार शून्य राहतो.
\paragraph{क्रमसूचक घातांकनाचा फलनसंच आणि शून्य आधार.} मूळ रचनेत आधार आणि घातांक यांचा फलनसंचातील क्रम उलटा आहे. मराठी व्याख्या आणि सरावात घातांकावरून आधाराकडे जाणारी, केवळ सांत ठिकाणी अशून्य असणारी फलने घेतली आहेत. शून्य आधार आणि $\omega$ घातांकासाठी मात्र मूळ सीमा-सूत्रात शून्य घातांकाचे मूल्य $1$ संघटित होते, तर फलनसंच रिकामा असतो. त्यामुळे मुख्य दोन्ही व्याख्या आणि सराव शून्येतर आधारापुरते केले आहेत; सरावाची सिद्धता दिलेली नाही.

\paragraph{सांततेच्या विधानातील अर्थदुरुस्ती.} मूळ इंग्रजीत $\card{A}\notin\omega$ या अटीचा अर्थ चुकून “$A$ नैसर्गिक संख्या नाही” असा दिला आहे. सांत पण नैसर्गिक संख्या नसलेले संच असतात; म्हणून मराठीत “$A$ सांत नाही” असे म्हटले आहे.
\paragraph{अनंत उत्तरवर्ती संख्येच्या पुराव्यातील मधला टप्पा.} मूळ पुरावा उत्तरवर्ती क्रमसूचक संख्या अनंत असल्यास तिची पूर्ववर्ती अनंत असते, यासाठी सांत क्रमसूचक संख्यांच्या तुल्यबलतेवरील प्रतिज्ञेचा संदर्भ देतो. तो संदर्भ एकटाच पुरेसा नाही. पूर्ववर्ती सांत असती तर तिची उत्तरवर्तीही सांत असती, हा आवश्यक टप्पा मराठीत स्पष्ट केला आहे.

\paragraph{संचांक घातांकनाची संकेतरचना.} घातसंचाच्या संचांकावरील पुराव्यात आणि घातांकनाच्या मर्यादेत मूळ स्रोत साधी वरलिखित घात-संकेतरचना वापरतो. त्या ठिकाणी मराठीत संचांक घातांकनाची स्पष्ट संकेतरचना वापरली आहे.
\paragraph{कॅनॉनिकल क्रमातील सांत प्रकरण.} आद्यखंडाच्या आकारमानाचा पुरावा मूळ स्रोत फक्त अनंत कमाल निर्देशांकासाठी स्पष्ट करतो. कमाल निर्देशांक सांत असेल तर आद्यखंड सांत असतो, हे आवश्यक प्रकरण मराठीत जोडले आहे.
\paragraph{घातांकनाच्या पुराव्यातील फलनांचे प्रकार.} वियुक्त युतीवरील फलनाचे विभाजन दोन फलनांची क्रमित जोडी देते; त्यांच्या आलेखांचा कार्तीय गुणाकार देत नाही. तसेच पुनर्गटित फलनाचा प्रत्यक्ष प्रांत कार्तीय गुणाकार आहे, त्याचा संचांक-प्रतिनिधी नाही. मराठीत हे दोन्ही प्रकार दुरुस्त केले आहेत.
\paragraph{शून्य घातांकाचा अपवाद.} अनंत संचांकाचा सांत घात स्वतः त्या संचांकाइतका असतो, या विधानात शून्य घातांक वगळावा लागतो. शून्य घाताचे मूल्य एक असते. GCH अंतर्गत पुढील घात-मर्यादेतही अनंत आधार आणि शून्येतर घातांक या आवश्यक अटी मराठीत स्पष्ट केल्या आहेत.
\paragraph{अलेफ अनुक्रमाच्या पुराव्यातील टप्पे.} प्रत्येक अनंत संचांक अलेफ अनुक्रमात येतो, या पुराव्यात लघुतम अनंत संचांकाचे आधारप्रकरण आणि अनंत संचांकांपुरते विगमन आवश्यक आहे. हे आवश्यक टप्पे मराठीत स्पष्ट केले आहेत. एकमेवतेविषयीचा वैध मूळ संक्षिप्त तर्क स्वतंत्र विस्तार न करता जपला आहे.
\paragraph{स्थिरबिंदू अनुक्रमातील वाढ.} मूळ अंतिम रचनेत प्रारंभिक संचांक आधीच बेथ-स्थिरबिंदू असेल तर पुनरावर्तन स्थिर राहते. त्यामुळे दावा केलेली काटेकोर वाढ आणि नंतरचे एकास-एक फलन मिळत नाही. मराठीत सुरुवात प्रारंभिक संचांकाच्या उत्तरवर्ती संचांकापासून केली आहे. पुढील अनुक्रमाचे शून्यवे मूल्य शून्य आहे; ते बेथ-स्थिरबिंदू नाही, हा अपवादही स्पष्ट केला आहे.


\paragraph{टार्स्की--स्कॉटच्या आक्षेपातील रिक्त संच.} तुल्यबल संचांचे सर्व प्रतिनिधी एकत्र केल्यास प्रतांकाची मर्यादा राहत नाही, हा आक्षेप अरिक्त आधारसंचासाठी आहे. रिक्त आधारसंचाचा वर्ग मात्र एकघटकी संच आहे. मराठीत हा अपवाद स्पष्ट केला आहे.
\paragraph{हार्टोग्सच्या पुराव्यातील संच आणि निर्देशांक.} मूळ पुराव्यात सुक्रमित आधारसंच हा संबंधाचा उपसंच असल्यासारखा छापला आहे; तो मूळ आधारसंचाचा उपसंच असतो. घूर्णनाशी संबंध नसलेल्या याच पुराव्यात फलनाने क्रम नेण्याच्या सूत्रात निश्चित क्रमसूचक संख्या आणि अंतर्गत निर्देशांक यांना एकच नाव दिले आहे. मराठीत वेगळी, प्रांतात बंधित नावे वापरली आहेत.
\paragraph{निवडीच्या पुराव्यातील थांबण्याची व्याख्या.} मूळ पुढील-मूल्य अटीतील असमानतेची दिशा आधीच्या निवडी पुसून टाकते. ती दुरुस्त करून, एकास-एक आणि आच्छादक फलनाचा दावा फक्त थांबण्यापूर्वीच्या अंशासाठी केला आहे. रिक्त आधारसंचाचे प्रकरण आधी सोडवले आहे; न थांबण्याचा विरोधाभास प्रत्यक्ष आधारसंचावरील हार्टोग्स निष्कर्षाने मिळतो.
\paragraph{गणनीय निवडीच्या मोजणीतील निर्देशांक.} आधीच्या उपसंचांच्या युतीकरणात मूळ स्रोत चालू निर्देशांक पुन्हा वापरतो. मराठीत युतीकरणाचा योग्य चल निर्देशांक वापरला आहे आणि शून्यवे प्रकरण वेगळे केले आहे. गणनीय युतीकरणाच्या पुराव्यात गणनीय निवड पुरेशी आहे; या विभागात उलटी सममूल्यता सिद्ध केलेली नाही, हेही स्पष्ट केले आहे.
\paragraph{निवडसंच आणि प्रतिमा.} प्रत्येक सदस्याशी एकघटकी छेद असण्याच्या अटीत बाहेरील अतिरिक्त घटकांना बंदी नाही. निवडफलनाची प्रतिमा म्हणून ओळखण्याआधी ते अतिरिक्त घटक वगळले आहेत.
\paragraph{भूमितीय विधानांचे प्रांत.} बानाख--टार्स्कीचे घनगोल येथे धन त्रिज्येचे आणि त्रिमितीय आहेत; वितालीची वर्तुळे धन त्रिज्येची आहेत. पृष्ठभागाच्या विघटनातून पृष्ठभागाच्या प्रती मिळतात, घन अंतर्भागाच्या नव्हेत. अंतरालावरील स्पर्शरेषा-फलनाच्या सूत्रातील अतिरिक्त बंद कंस काढला आहे.
\paragraph{वितालीच्या घूर्णन गटाची दुरुस्ती.} परिमेय संख्येइतक्या रेडियनांची घूर्णने पूर्ण फेरीच्या मर्यादेत घेतल्यास व्यस्त आणि संयोजनासाठी संवरण मिळत नाही. म्हणून पूर्ण फेरीच्या परिमेय भागांची घूर्णने घेतली आहेत. शून्य घूर्णनाचा व्यस्त स्वतंत्रपणे हाताळला आहे. दोन अर्ध्यांना स्वतंत्र आधार म्हणण्याऐवजी जनक संच म्हटले आहे; प्रत्येक अर्ध्यावरील कोन-दुप्पट प्रतिचित्रण एकास-एक आणि आच्छादक असते. विभाजनातील अपरिभाषित पहिल्या-अर्ध्याचे नावही योग्य गटाच्या नावाने बदलले आहे.
\paragraph{त्रिमितीय रूपरेषेतील उरलेली पुरावा-उणीव.} मूळ मजकूर प्रत्येक मुक्त गटासाठी दुप्पट विघटन सांगतो; येथे प्रत्यक्ष वापरलेले दोन जनकांवरील मुक्त गटाचे प्रकरण म्हटले आहे. पण त्रिमितीय घूर्णने त्यांच्या अक्षांवरील काही बिंदू स्थिर ठेवतात. वर्तुळावरील मुक्त क्रियेचा पुरावा त्या बिंदूंना थेट लागू होत नाही. त्यांचा स्वतंत्र विचार या रूपरेषेत दिलेला नाही; म्हणून ही पूर्ण पुरावा नसल्याची उणीव स्पष्ट ठेवली आहे.
\paragraph{मापाच्या पुराव्यातील क्षेत्र आणि चल.} माप ऋणेतर असल्याची आवश्यक अट आणि अमापनीयता ही निश्चित मापाच्या क्षेत्राविषयी असल्याचा अर्थ स्पष्ट केला आहे. शेवटच्या पुराव्यात घूर्णनाचा चल चुकून निवडलेल्या बिंदूंच्या संचात घेतला होता; तो घूर्णन गटात घेतला आहे.


\paragraph{सिद्धता-प्रकरणातील आवृत्ती आणि विभागक्रम.} मूळ इंग्रजी आवृत्तीतील नैसर्गिक भाषेचा उल्लेख या मराठी आवृत्तीसाठी स्पष्ट केला आहे. सिद्धतेचे प्रकार ताबडतोब पुढच्या विभागात येतात असे दोन उल्लेख आहेत; प्रत्यक्ष स्रोतक्रमात व्याख्यांचा विभाग आधी असल्यामुळे पुढील विभागांचा सामान्य निर्देश केला आहे. व्याख्यांच्या विभागाचा मूळ भाग-निर्देशांकही पद्धती भागाशी जुळवला आहे.
\paragraph{अस्तित्ववाची विधानातील नव्या नावाचे कार्य.} निष्कर्षात नाव नसणे एवढेच पुरेसे नाही: नव्या वस्तूविषयी अतिरिक्त, मुक्त न केलेली गृहीतके वापरता कामा नयेत. तात्पुरते सदस्यत्व गृहीतक मुक्त केल्यावर मूळ अस्तित्ववाची विधान आणि इतर अनुमत गृहीतके यांचा आधार उरतो. मूळ स्पष्टीकरणात अरिक्त संचाऐवजी बंधित चलाचे नाव आले होते; संचाचे नाव वापरले आहे.
\paragraph{संचांच्या उदाहरणांतील मुद्रणदोष.} एका संयोग-नॉनसमावेशनाच्या स्पष्टीकरणात अपरिभाषित संचाचे नाव आले होते; योग्य संयोग वापरला आहे. शोषणाच्या विस्तृत सिद्धतेत पहिला समावेशन-निष्कर्ष दोनदा दिला होता; दुसरा निष्कर्ष उलट समावेशनाने बदलला आहे. दोन उदाहरणांतील न जुळणारे कंस दुरुस्त केले आहेत.


\paragraph{नैसर्गिक संख्यांवरील विगमनाचे निर्देशांक.} मूळ विवेचनाच्या एका पायरीत $k$ हा चल असताना त्याच्या जागी $1$ ठेवण्याऐवजी दुसऱ्या चलाचे नाव आले आहे; योग्य पायरीचल वापरला आहे. फाशांच्या उदाहरणात शून्य फाशांचे प्रकरण प्रस्तावनेत असले, तरी औपचारिक विगमन $1$ पासून सुरू होते. शून्य फाशांच्या रिकाम्या बेरजेचे प्रकरण स्पष्टपणे जोडले आहे.
\paragraph{प्रबल विगमनातील रिक्त पूर्वभाग.} $k=0$ असताना $l<0$ अशी कोणतीही नैसर्गिक संख्या नसते. या रिक्त पूर्वभागातील मूळ मजकुरात बंधित $l$ ऐवजी $P(0)$ छापले आहे; विगमनाच्या प्रत्यक्ष गृहीतकानुसार $P(l)$ वापरले आहे.
\paragraph{उचित पूर्वखंडाच्या पुराव्याची व्याप्ती.} मूळ व्याख्येत रिकामी चिन्हमाला उचित पूर्वखंड ठरते, पण तीत उघडणारे आणि बंद करणारे कंस शून्य असतात. म्हणून कंसांची काटेकोर असमानता तिच्यासाठी असत्य आहे. मूळ पुराव्यातील पाचही प्रकरणे अरिक्त पूर्वखंडांचीच आहेत; मराठीत ती आवश्यक अट व्याख्या, विधान आणि पुराव्यात स्पष्ट केली आहे.


\paragraph{इतिहासातील चरित्रविषयक दुरुस्त्या.} चर्च यांच्या प्रबंधाशी संबंधित गणनीयतेचे प्रतिपादन आणि प्रथम-क्रम निर्णय समस्येवरील त्यांचे प्रमेय यांतील भेद स्पष्ट केला आहे. पेटर यांच्या जन्मनावातील उलटलेला उल्लेख, ज्युलिया रॉबिन्सन यांच्या नॅशनल अकॅडमीतील निवडीचा गणित विभागाशी संबंध, ट्यूरिंगवरील चित्रपटातील कलाकारांची नावे आणि झर्मेलो यांच्या गॉटिंगेनमधील अभ्यासस्थानाची वर्तनी तपासून दुरुस्त केली आहे.


\paragraph{सीमेची व्याख्या आणि चित्रातील निर्देशांक.} मूळ मर्यादा-व्याख्येत $x=c$ वगळलेले नाही; तसे केल्याशिवाय सामान्य सीमेची अट मिळत नाही. मराठी सूत्रात $0<|x-c|$ ही अट स्पष्ट केली आहे. निरपेक्ष मूल्याच्या आलेखात $-1$ या बिंदूचे चिन्ह दुरुस्त केले आहे. तसेच आधीच्या अवकलज-सूत्रातील बदलत्या भागाकाराचा उल्लेख स्पष्ट केला आहे. \paragraph{द्विमान विस्तार.} \textbf{SOL6-A281} मध्ये द्विमान निरूपणांची एकच निवड ठरवून एकास-एक फलनाची अट पूर्ण केली आहे. मूळ missing-value उदाहरण या प्रांतावर प्रतिमेतच येते; त्याऐवजी प्रतिमेबाहेरचे $5/6$ उदाहरण दिले आहे. ऐतिहासिक दशमान प्रयत्न आणि त्याचे आधुनिक द्विमान रूपांतर वेगळे केले आहेत. \paragraph{हिल्बर्ट वक्राच्या सिद्धतेची दुरुस्ती.} प्रांत एकक रेषा ठेवून \textbf{SOL6-A282} मध्ये प्रत्येक स्तराचे समान बंद प्राचल-अंतराल आणि जाळी-घरांचे सुसंगत refinement दिले आहे. घरकर्णाच्या मर्यादेने समसमान अभिसरण मिळते. अंतर्भूत बंद अंतराल प्रत्येक चौरस-बिंदूची प्रतिमापूर्व साक्ष देतात; प्रत्येक आदान बिंदूवरील सातत्य epsilon-delta अटीने सिद्ध केले आहे. मूळ तीन आकृत्या तशाच जपल्या आहेत.

\chapter*{संदर्भ}
\markboth{संदर्भ}{संदर्भ}

\addcontentsline{toc}{chapter}{संदर्भ}

खालील संदर्भ मूळ स्रोताच्या ग्रंथसूचीवरून घेतले आहेत. लेखकांची नावे, प्रकाशनांची मूळ शीर्षके आणि ग्रंथसूचीतील तपशील ओळखीसाठी मूळ भाषेत ठेवले आहेत. ही नोंद दुव्यांच्या सध्याच्या उपलब्धतेची पडताळणी नाही.


\par\medskip\phantomsection\label{bib:Turing1937}
Alan M. Turing (1937). \emph{On computable numbers, with an application to the ``{E}ntscheidungsproblem''}. \emph{Proceedings of the London Mathematical Society, 2nd Series}, 42, 230--265.

\par\medskip\phantomsection\label{bib:Tarski1981}
Alfred Tarski (1981). \emph{The Collected Works of {A}lfred {T}arski}. Birkh\"{a}user. Basel.

\par\medskip\phantomsection\label{bib:Church1936}
Alonzo Church (1936). \emph{A note on the {E}ntscheidungsproblem}. \emph{The Journal of Symbolic Logic}, 1, 40--41.

\par\medskip\phantomsection\label{bib:Church1936a}
Alonzo Church (1936). \emph{An Unsolvable Problem of Elementary Number Theory}. \emph{American Journal of Mathematics}, 58, 345--363.

\par\medskip\phantomsection\label{bib:Andrasfai1986}
Andr\'asfai, B\'ela (1986). \emph{R\'ozsa ({R}osa) {P}\'eter}. \emph{Periodica Polytechnica Electrical Engineering}, 30(2-3), 139--145. \url{http://www.pp.bme.hu/ee/article/view/4651}.

\par\medskip\phantomsection\label{bib:Aspray1984}
Aspray, William (1984). \emph{The {P}rinceton Mathematics Community in the 1930s: {A}lonzo {C}hurch}. Interview. \url{http://www.princeton.edu/mudd/finding_aids/mathoral/pmc05.htm}.

\par\medskip\phantomsection\label{bib:Baaz2011}
Baaz, Matthias; Papadimitriou, Christos H.; Putnam, Hilary W.; Scott, Dana S.; Harper Jr., Charles L. (2011). \emph{Kurt {G}\"odel and the Foundations of Mathematics: Horizons of Truth}. Cambridge University Press. Cambridge.

\par\medskip\phantomsection\label{bib:BanachTarski1924}
Banach, Stefan; Tarski, Alfred (1924). \emph{Sur la d\'{e}composition des ensembles de points en parties respectivement congruentes}. \emph{Fundamenta Mathematicae}, 6, 244--77.

\par\medskip\phantomsection\label{bib:Benacerraf1965}
Benacerraf, Paul (1965). \emph{What numbers could not be}. \emph{The Philosophical Review}, 74(1), 47--73.

\par\medskip\phantomsection\label{bib:Berkeley1734}
Berkeley, George (1734). \emph{The Analyst; or, a Discourse Adressed to an Infidel Mathematician}.

\par\medskip\phantomsection\label{bib:Russell1905}
Bertrand Russell (1905). \emph{On Denoting}. \emph{Mind}, 14, 479--493.

\par\medskip\phantomsection\label{bib:Russell1967}
Bertrand Russell (1967). \emph{The Autobiography of Bertrand Russell}. Allen and Unwin. London.

\par\medskip\phantomsection\label{bib:Russell1968}
Bertrand Russell (1968). \emph{The Autobiography of Bertrand Russell}. Allen and Unwin. London.

\par\medskip\phantomsection\label{bib:Russell1969}
Bertrand Russell (1969). \emph{The Autobiography of Bertrand Russell}. Allen and Unwin. London.

\par\medskip\phantomsection\label{bib:RussellND}
Bertrand Russell (n.d.). \emph{{B}ertrand {R}ussell on Smoking}. Video Interview. \url{https://www.youtube.com/watch?v=80oLTiVW_lc}.

\par\medskip\phantomsection\label{bib:Boolos1971}
Boolos, George (1971). \emph{The Iterative Conception of Set}. \emph{The Journal of Philosophy}, 68(8), 215--31.

\par\medskip\phantomsection\label{bib:Boolos1989}
Boolos, George (1989). \emph{Iteration Again}. \emph{Philosophical Topics}, 17(2), 5--21.

\par\medskip\phantomsection\label{bib:Boolos2000}
Boolos, George (2000). \emph{Must we Believe in Set Theory?}. मध्ये: \emph{Between Logic and Intuition: Essays in Honor of {Charles} {Parsons}}. पृ.~257--68. Cambridge University Press. Cambridge.

\par\medskip\phantomsection\label{bib:Burali-Forti1897}
Burali-Forti, Cesare (1897). \emph{Una questione sui numeri transfiniti}. \emph{Rendiconti del Circolo Matematico di Palermo}, 11, 154--64.

\par\medskip\phantomsection\label{bib:ButtonLT1}
Button, Tim (2021). \emph{Level Theory, Part~1: Axiomatizing the Bare Idea of a Cumulative Hierarchy of Sets}. \emph{The Bulletin of Symbolic Logic}, 27(4), 436-460.

\par\medskip\phantomsection\label{bib:Cantor1878}
Cantor, Georg (1878). \emph{Ein {Beitrag} zur {Mannigfaltigkeitslehre}}. \emph{Journal für die reine und angewandte Mathematik}, 84, 242--58.

\par\medskip\phantomsection\label{bib:Cantor1883}
Cantor, Georg (1883). \emph{Grundlagen einer allgemeinen {Mannigfaltigkeitslehre}. Ein mathematisch-philosophischer {Versuch} in der {Lehre} des {Unendlichen}}. Teubner. Leipzig.

\par\medskip\phantomsection\label{bib:Cantor1892}
Cantor, Georg (1892). \emph{{\"Uber} eine elementare {Frage} der {Mannigfaltigkeitslehre}}. \emph{Jahresbericht der deutschen Mathematiker-Vereinigung}, 1, 75--8. Teubner. Leipzig.

\par\medskip\phantomsection\label{bib:Cohen1963}
Cohen, Paul J. (1963). \emph{The independence of the continuum hypothesis}. \emph{Proceedings of the National Academy of Sciences of the United States of America}, 24, 556--557.

\par\medskip\phantomsection\label{bib:Cohen1966}
Cohen, Paul J. (1966). \emph{Set Theory and the Continuum Hypothesis}. Benjamin. Reading, MA.

\par\medskip\phantomsection\label{bib:Reid1996}
Constance Reid (1996). \emph{Julia: A Life in Mathematics}. Cambridge University Press. Cambridge. \url{https://books.google.ca/books?id=lRtSzQyHf9UC&lpg=PP1&pg=PP1#v=onepage&q&f=false}.

\par\medskip\phantomsection\label{bib:Conway2006}
Conway, John (2006). \emph{The Power of Mathematics}. मध्ये: \emph{Power}. Cambridge University Press. Cambridge. \url{http://www.cs.toronto.edu/~mackay/conway.pdf}.

\par\medskip\phantomsection\label{bib:Csicsery2016}
Csicsery, George (2016). \emph{Zala Films: {J}ulia {R}obinson and {H}ilbert's Tenth Problem}. \url{http://www.zalafilms.com/films/juliarobinson.html}.

\par\medskip\phantomsection\label{bib:Velleman2019}
Daniel J. Velleman (2019). \emph{How to Prove It: A Structured Approach}. 3rd आवृत्ती. Cambridge University Press. Cambridge.

\par\medskip\phantomsection\label{bib:Solow2013}
Daniel Solow (2013). \emph{How to Read and Do Proofs}. Wiley. Hoboken, NJ.

\par\medskip\phantomsection\label{bib:Dauben1990}
Dauben, Joseph (1990). \emph{{G}eorg {C}antor: {H}is {M}athematics and {P}hilosophy of the {I}nfinite}. Princeton University Press. Princeton.

\par\medskip\phantomsection\label{bib:Dedekind1888}
Dedekind, Richard (1888). \emph{Was sind und was sollen die {Zahlen}?}. Vieweg. Braunschweig.

\par\medskip\phantomsection\label{bib:Dick1981}
Dick, Auguste (1981). \emph{Emmy {N}oether 1882--1935}. Birkh\"auser. Boston.

\par\medskip\phantomsection\label{bib:Duncan2015}
Duncan, Arlene (2015). \emph{The {B}ertrand {R}ussell {R}esearch {C}entre}. \url{http://russell.mcmaster.ca/}.

\par\medskip\phantomsection\label{bib:Ebbinghaus2015}
Ebbinghaus, Heinz-Dieter (2015). \emph{Ernst {Z}ermelo: {A}n {A}pproach to his {L}ife and {W}ork}. Springer-Verlag. Berlin.

\par\medskip\phantomsection\label{bib:Ebbinghaus2010}
Ebbinghaus, Heinz-Dieter; Fraser, Craig G.; Kanamori, Akihiro (2010). \emph{Ernst Zermelo. Collected Works}. Springer-Verlag. Berlin.

\par\medskip\phantomsection\label{bib:Ebbinghaus2013}
Ebbinghaus, Heinz-Dieter; Kanamori, Akihiro (2013). \emph{Ernst Zermelo: Collected Works}. Springer-Verlag. Berlin.

\par\medskip\phantomsection\label{bib:Steinhart2018}
Eric Steinhart (2018). \emph{More Precisely: The Math You Need to Do Philosophy}. 2nd आवृत्ती. Broadview. Peterborough, ON.

\par\medskip\phantomsection\label{bib:Zermelo1904}
Ernst Zermelo (1904). \emph{Beweis, da{\ss} jede {M}enge wohlgeordnet werden kann}. \emph{Mathe\-ma\-tische An\-nalen}, 59, 514--516. {E}nglish translation in \hyperref[bib:Ebbinghaus2010]{Ebbinghaus 2010, pp.~115--119}.

\par\medskip\phantomsection\label{bib:Cheng2004}
Eugenia Cheng (2004). \emph{How to Write Proofs: A Quick Quide}. \url{https://eugeniacheng.com/wp-content/uploads/2017/02/cheng-proofguide.pdf}.

\par\medskip\phantomsection\label{bib:Feferman2004}
Feferman, Anita; Feferman, Solomon (2004). \emph{Alfred Tarski: Life and Logic}. Cambridge University Press. Cambridge.

\par\medskip\phantomsection\label{bib:Feferman1994}
Feferman, Solomon (1994). \emph{{J}ulia {B}owman {R}obinson 1919--1985}. \emph{Biographical Memoirs of the National Academy of Sciences}, 63, 1--28. National Academies Press. \url{http://www.nasonline.org/publications/biographical-memoirs/memoir-pdfs/robinson-julia.pdf}.

\par\medskip\phantomsection\label{bib:Godel1986}
Feferman, Solomon; Dawson Jr., John W.; Kleene, Stephen C.; Moore, Gregory H.; Solovay, Robert M.; van Heijenoort, Jean (1986). \emph{Kurt G\"odel: Collected Works. Vol. 1: Publications 1929--1936}. Oxford University Press. Oxford.

\par\medskip\phantomsection\label{bib:Godel1990}
Feferman, Solomon; Dawson Jr., John W.; Kleene, Stephen C.; Moore, Gregory H.; Solovay, Robert M.; van Heijenoort, Jean (1990). \emph{Kurt G\"odel: Collected Works. Vol. 2: Publications 1938--1974}. Oxford University Press. Oxford.

\par\medskip\phantomsection\label{bib:Fraenkel1922}
Fraenkel, Abraham (1922). \emph{\"{U}ber den {Begriff} `definit' und die {Unabh\"angigkeit} des {Auswahlaxioms}}. \emph{Sitzungsberichte der Preussischen Akadademie der Wissenschaften, Physikalisch-mathematische Klasse}, 253--257.

\par\medskip\phantomsection\label{bib:Frege1884}
Frege, Gottlob (1884). \emph{Die {Grundlagen} der {Arithmetik}: Eine logisch mathematische {Untersuchung} \"{u}ber den {Begriff} der {Zahl}}. Wilhelm Koebner. Breslau. Translation in \hyperref[bib:Frege1953]{Frege 1953}.

\par\medskip\phantomsection\label{bib:Frege1953}
Frege, Gottlob (1953). \emph{Foundations of Arithmetic}. 2nd आवृत्ती. Basil Blackwell \& Mott. Oxford.

\par\medskip\phantomsection\label{bib:Frey2015}
Frey, Holly; Wilson, Tracy V. (2015). \emph{Stuff you Missed in History Class: {E}mmy {N}oether, Mathematics Trailblazer}. Podcast audio. \url{https://www.iheart.com/podcast/stuff-you-missed-in-history-cl-21124503/episode/emmy-noether-mathematics-trailblazer-30207491/}.

\par\medskip\phantomsection\label{bib:Godel1938}
G\"odel, Kurt (1938). \emph{The consistency of the axiom of choice and the generalized continuum hypothesis}. \emph{Proceedings of the National Academy of Sciences of the United States of America}, 50, 1143--8.

\par\medskip\phantomsection\label{bib:Gentzen1935a}
Gerhard Gentzen (1935). \emph{Untersuchungen \"uber das logische {S}chlie\ss en {I}}. \emph{Mathe\-mati\-sche Zeit\-schrift}, 39, 176--210. {E}nglish translation in \hyperref[bib:Gentzen1969]{Szabo 1969}, pp.~68--131. \url{https://doi.org/10.1007/BF01201353}.

\par\medskip\phantomsection\label{bib:Gentzen1935b}
Gerhard Gentzen (1935). \emph{Untersuchungen \"uber das logische {S}chlie\ss en {II}}. \emph{Mathe\-mati\-sche Zeit\-schrift}, 39, 176--210, 405--431. {E}nglish translation in \hyperref[bib:Gentzen1969]{Szabo 1969}, pp.~68--131. \url{https://doi.org/10.1007/BF01201363}.

\par\medskip\phantomsection\label{bib:Giaquinto2007}
Giaquinto, Marcus (2007). \emph{Visual Thinking in Mathematics}. Oxford University Press. Oxford.

\par\medskip\phantomsection\label{bib:Gouvea2011}
Gouv\^{e}a, Fernando Q. (2011). \emph{Was {C}antor Surprised?}. \emph{American Mathematical Monthly}, 118(3), 198-209.

\par\medskip\phantomsection\label{bib:Grattan-Guinness1971}
Grattan-Guinness, Ivor (1971). \emph{Towards a Biography of {G}eorg {C}antor}. \emph{Annals of Science}, 27(4), 345--391.

\par\medskip\phantomsection\label{bib:Wagon2016}
Grzegorz {Tomkowicz}; Stan {Wagon} (2016). \emph{The {Banach}-{Tarski} Paradox}. Cambridge University Press. Cambridge.

\par\medskip\phantomsection\label{bib:Hartogs1915}
Hartogs, Friedrich (1915). \emph{\"{U}ber das {P}roblem der {W}ohlordnung}. \emph{Mathematische Annalen}, 76, 438--43.

\par\medskip\phantomsection\label{bib:Hausdorff1914}
Hausdorff, Felix (1914). \emph{Bemerkung über den {Inhalt} von {Punktmengen}}. \emph{Mathematische Annalen}, 75, 428--34.

\par\medskip\phantomsection\label{bib:Heijenoort1967}
Heijenoort, Jean van (1967). \emph{From {F}rege to {G}\"{o}del: A Source Book in Mathematical Logic, 1879--1931}. Harvard University Press. Cambridge, MA.

\par\medskip\phantomsection\label{bib:EndertonND}
Herbert B. Enderton (2019). \emph{{A}lonzo {C}hurch: {L}ife and {W}ork}. मध्ये: \emph{The Collected Works of {A}lonzo {C}hurch}. MIT Press. Cambridge, MA.

\par\medskip\phantomsection\label{bib:Hilbert1891}
Hilbert, David (1891). \emph{\"{U}ber die stetige {Abbildung} einer {Linie} auf ein {Fl\"{a}chenst\"{u}ck}}. \emph{Mathematische Annalen}, 38(3), 459--460.

\par\medskip\phantomsection\label{bib:EwaldSieg2013}
Hilbert, David (2013). \emph{{D}avid {H}ilbert's Lectures on the Foundations of Arithmetic and Logic 1917--1933}. Springer. Heidelberg.

\par\medskip\phantomsection\label{bib:Hodges2014}
Hodges, Andrew (2014). \emph{Alan Turing: The Enigma}. Vintage. London.

\par\medskip\phantomsection\label{bib:Hume1740}
Hume, David (1740). \emph{A Treatise of Human Nature}. London.

\par\medskip\phantomsection\label{bib:Incurvati2020}
Incurvati, Luca (2020). \emph{Conceptions of Set and the Foundations of Mathematics}. Cambridge University Press. Cambridge.

\par\medskip\phantomsection\label{bib:Irvine2015}
Irvine, Andrew David (2015). \emph{Sound Clips of {B}ertrand {R}ussell Speaking}. \url{http://plato.stanford.edu/entries/russell/russell-soundclips.html}.

\par\medskip\phantomsection\label{bib:Noether1983}
Jacobson, Nathan (1983). \emph{Emmy {N}oether: {G}esammelte {A}bhandlungen---{C}ollected {P}apers}. Springer-Verlag. Berlin.

\par\medskip\phantomsection\label{bib:Tarski1983}
John Corcoran (1983). \emph{Logic, Semantics, Metamathematics}. 2nd आवृत्ती. Hackett. Indianapolis.

\par\medskip\phantomsection\label{bib:Dawson1997}
John Dawson, Jr. (1997). \emph{Logical Dilemmas: The Life and Work of Kurt {G}\"odel}. CRC Press. Boca Raton.

\par\medskip\phantomsection\label{bib:Robinson1949}
Julia Robinson (1949). \emph{Definability and Decision Problems in Arithmetic}. \emph{The Journal of Symbolic Logic}, 14(2), 98-114. \url{http://www.jstor.org/stable/2266510}.

\par\medskip\phantomsection\label{bib:KatzKatz2012}
Katz, Karin Usadi; Katz, Mikhail G. (2012). \emph{{S}tevin Numbers and Reality}. \emph{Foundations of Science}, 17(2), 109--23.

\par\medskip\phantomsection\label{bib:Kunen1980}
Kunen, Kenneth (1980). \emph{Set Theory: An Introduction to Independence Proofs}. North Holland. New York.

\par\medskip\phantomsection\label{bib:Godel1929}
Kurt G\"odel (1929). \emph{{\"U}ber die {V}ollst\"andigkeit des {L}ogikkalk\"uls [{O}n the Completeness of the Calculus of Logic]}. Reprinted and translated in \hyperref[bib:Godel1986]{Feferman 1986}, pp.~60--101.

\par\medskip\phantomsection\label{bib:Godel1931}
Kurt G{\"{o}}del (1931). \emph{\"Uber formal unentscheidbare {S}\"atze der \emph{{P}rincipia {M}athematica} und verwandter {S}ysteme {I} [{O}n Formally Undecidable Propositions of \emph{{P}rincipia {M}athematica} and Related Systems~{I}]}. \emph{Monatshefte f\"ur Mathematik und Physik}, 38, 173--198. Reprinted and translated in \hyperref[bib:Godel1986]{Feferman 1986}, pp.~144--195. \url{https://doi.org/10.1007/BF01700692}.

\par\medskip\phantomsection\label{bib:Levy1960}
L\'{e}vy, Azriel (1960). \emph{Axiom Schemata of Strong Infinity in Axiomatic Set Theory}. \emph{Pacific Journal of Mathematics}, 10(1), 223--38.

\par\medskip\phantomsection\label{bib:LibriVoxND}
LibriVox (n.d.). \emph{Bertrand {R}ussell}. Collection of public domain audiobooks. \url{https://librivox.org/author/1508?primary_key=1508&search_category=author&search_page=1&search_form=get_results}.

\par\medskip\phantomsection\label{bib:Linnebo2010}
Linnebo, \O{}ystein (2010). \emph{Predicative and Impredicative Definitions}. \emph{Internet Encyclopedia of Philosophy}. \url{http://www.iep.utm.edu/predicat/}.

\par\medskip\phantomsection\label{bib:Linsenmayer2014}
Linsenmayer, Mark (2014). \emph{The Partially Examined Life: {G}\"odel on Math}. Podcast audio. \url{http://www.partiallyexaminedlife.com/2014/06/16/ep95-godel/}.

\par\medskip\phantomsection\label{bib:MacFarlane2015}
MacFarlane, John (2015). \emph{Alonzo {C}hurch's {JSL} Reviews}. \url{http://johnmacfarlane.net/church.html}.

\par\medskip\phantomsection\label{bib:Magnus2021}
Magnus, P. D.; Button, Tim; Loftis, J. Robert; {Thomas-Bolduc}, Aaron; Trueman, Robert; Zach, Richard (2021). \emph{Forall {\emph{x}}: {{Calgary}}. {{An}} Introduction to Formal Logic}. F21 आवृत्ती. {Open Logic Project}. {Calgary}. \url{https://forallx.openlogicproject.org/}.

\par\medskip\phantomsection\label{bib:Sautoy2014}
Marcus du Sautoy (2014). \emph{{A} Brief History of Mathematics: {G}eorg {C}antor}. Audio Recording. \url{http://www.bbc.co.uk/programmes/b00ss1j0}.

\par\medskip\phantomsection\label{bib:DavisPutnamRobinson1961}
Martin Davis; Hilary Putnam; Julia Robinson (1961). \emph{The Decision Problem for Exponential {D}iophantine Equations}. \emph{Annals of Mathematics}, 74(3), 425-436. \url{http://www.jstor.org/stable/1970289}.

\par\medskip\phantomsection\label{bib:Matijasevich1992}
Matijasevich, Yuri (1992). \emph{My Collaboration with {J}ulia {R}obinson}. \emph{The Mathematical Intelligencer}, 14(4), 38--45.

\par\medskip\phantomsection\label{bib:Menzler-Trott2007}
Menzler-Trott, Eckart (2007). \emph{Logic's Lost Genius: The Life of {G}erhard {G}entzen}. American Mathematical Society. Providence.

\par\medskip\phantomsection\label{bib:Hutchings2003}
Michael Hutchings (2003). \emph{Introduction to Mathematical Arguments}. \url{https://math.berkeley.edu/~hutching/teach/proofs.pdf}.

\par\medskip\phantomsection\label{bib:Montague1961}
Montague, Richard (1961). \emph{Semantic closure and non-finite axiomatizability {I}}. मध्ये: \emph{Infinitistic Methods: Proceedings of the Symposium on Foundations of Mathematics (Warsaw 1959)}. पृ.~45--69. Pergamon. New York.

\par\medskip\phantomsection\label{bib:Montague1965}
Montague, Richard (1965). \emph{Set theory and higher-order logic}. मध्ये: \emph{Formal systems and recursive functions}. पृ.~131--48. North-Holland. Amsterdam. Proceedings of the Eight Logic Colloquium, July 1963.

\par\medskip\phantomsection\label{bib:Imitation2014}
Morten Tyldum (2014). \emph{The Imitation Game}. Motion picture.

\par\medskip\phantomsection\label{bib:OConnorRobertson:RN}
O'Connor, John J.; Robertson, Edmund F. (2005). \emph{The real numbers: {Stevin} to {Hilbert}}. \url{http://www-history.mcs.st-and.ac.uk/HistTopics/Real_numbers_2.html}.

\par\medskip\phantomsection\label{bib:Oconnor2014}
O'Connor, John J.; Robertson, Edmund F. (2014). \emph{R\'{o}zsa {P}\'{e}ter}. \url{http://www-groups.dcs.st-and.ac.uk/~history/Biographies/Peter.html}.

\par\medskip\phantomsection\label{bib:Peter1935b}
P\'eter, R\'ozsa (1935b). \emph{Konstruktion nichtrekursiver {F}unktionen}. \emph{Mathematische Annalen}, 111, 42--60. \url{https://doi.org/10.1007/BF01472200}.

\par\medskip\phantomsection\label{bib:Peter1951}
P\'eter, R\'ozsa (1951). \emph{Rekursive Funktionen}. Akademiai Kiado. Budapest. English translation in \hyperref[bib:Peter1967]{P\'eter 1967}.

\par\medskip\phantomsection\label{bib:Peter1967}
P\'eter, R\'ozsa (1967). \emph{Recursive Functions}. Academic Press. New York.

\par\medskip\phantomsection\label{bib:Peter2010}
P\'eter, R\'ozsa (2010). \emph{Playing with Infinity}. Dover. New York. \url{https://books.google.ca/books?id=6V3wNs4uv_4C&lpg=PP1&ots=BkQZaHcR99&lr&pg=PP1#v=onepage&q&f=false}.

\par\medskip\phantomsection\label{bib:Peano1890}
Peano, Giuseppe (1890). \emph{Sur une courbe, qui remplit toute une aire plane}. \emph{Mathematische Annalen}, 36(1), 157--60.

\par\medskip\phantomsection\label{bib:Maddy1988a}
Penelope Maddy (1988). \emph{Believing the Axioms~{I}}. \emph{The Journal of Symbolic Logic}, 53(2), 481--511.

\par\medskip\phantomsection\label{bib:Maddy1988b}
Penelope Maddy (1988). \emph{Believing the Axioms~{II}}. \emph{The Journal of Symbolic Logic}, 53(3), 736--64.

\par\medskip\phantomsection\label{bib:Perimeter2015}
Perimeter Institute (2015). \emph{Emmy {N}oether: Her Life, Work, and Influence}. Video Lecture. \url{https://www.youtube.com/watch?v=tNNyAyMRsgE}.

\par\medskip\phantomsection\label{bib:Potter2004}
Potter, Michael (2004). \emph{Set Theory and its Philosophy}. Oxford University Press. Oxford.

\par\medskip\phantomsection\label{bib:Peter1935a}
P{\'e}ter, R{\'o}zsa (1935a). \emph{{\"U}ber den {Z}usammenhang der verschiedenen {B}egriffe der rekursiven {F}unktion}. \emph{Mathematische Annalen}, 110, 612--632. \url{https://doi.org/10.1007/BF01448046}.

\par\medskip\phantomsection\label{bib:Radiolab2012}
Radiolab (2012). \emph{The {T}uring Problem}. Podcast audio. \url{http://www.radiolab.org/story/193037-turing-problem/}.

\par\medskip\phantomsection\label{bib:Ramsey1925}
Ramsey, Frank Plumpton (1925). \emph{The Foundations of Mathematics}. \emph{Proceedings of the London Mathematical Society}, 25, 338--384.

\par\medskip\phantomsection\label{bib:Reid1986}
Reid, Constance (1986). \emph{The Autobiography of {J}ulia {R}obinson}. \emph{The College Mathematics Journal}, 17, 3--21.

\par\medskip\phantomsection\label{bib:Hammack2013}
Richard Hammack (2013). \emph{Book of Proof}. Virginia Commonwealth University. Richmond, VA. \url{http://www.people.vcu.edu/~rhammack/BookOfProof/BookOfProof.pdf}.

\par\medskip\phantomsection\label{bib:Heck2012}
Richard Kimberly Heck (2012). \emph{Reading Frege's Grundgesetze}. Oxford University Press. Oxford.

\par\medskip\phantomsection\label{bib:Robinson1996}
Robinson, Julia (1996). \emph{The Collected Works of Julia Robinson}. American Mathematical Society. Providence.

\par\medskip\phantomsection\label{bib:Robinson1947}
Robinson, Raphael (1947). \emph{On the decomposition of spheres}. \emph{Fundamenta Mathematicae}, 34(1), 246--60.

\par\medskip\phantomsection\label{bib:Rose2012}
Rose, Daniel (2012). \emph{A Song about {G}eorg {C}antor}. Audio Recording. \url{https://www.youtube.com/watch?v=QUP5Z4Fb5k4}.

\par\medskip\phantomsection\label{bib:Rose2010}
Rose, Nicholas J (2010). \emph{Hilbert-Type Space-Filling Curves}. \url{https://web.archive.org/web/20151010184939/http://www4.ncsu.edu/~njrose/pdfFiles/HilbertCurve.pdf}.

\par\medskip\phantomsection\label{bib:Russell1919}
Russell, Bertrand (1919). \emph{Introduction to Mathematical Philosophy}. Allen \& Unwin. London.

\par\medskip\phantomsection\label{bib:Scott1974}
Scott, Dana (1974). \emph{Axiomatizing Set Theory}. मध्ये: \emph{Axiomatic Set Theory II}. पृ.~207--14. American Mathematical Society. Proceedings of the Symposium in Pure Mathematics of the American Mathematical Society, July--August 1967.

\par\medskip\phantomsection\label{bib:Segal2014}
Segal, Sanford L. (2014). \emph{Mathematicians under the Nazis}. Princeton University Press. Princeton.

\par\medskip\phantomsection\label{bib:Shoenfield:AST}
Shoenfield, Joseph R. (1977). \emph{Axioms of Set Theory}. मध्ये: \emph{Handbook of Mathematical Logic}. पृ.~321--44. North-Holland. London.

\par\medskip\phantomsection\label{bib:Sigmund2007}
Sigmund, Karl; Dawson, John; M\"uhlberger, Kurt; Enzensberger, Hans Magnus; Kennedy, Juliette (2007). \emph{Kurt {G}\"odel: {D}as {A}lbum--{T}he {A}lbum}. \emph{The Mathematical Intelligencer}, 29(3), 73--76. Springer.

\par\medskip\phantomsection\label{bib:Skolem1922}
Skolem, Thoralf (1922). \emph{Einige {Bemerkungen} zur axiomatischen {Begr\"undung} der {Mengenlehre}}. मध्ये: \emph{Wissenschaftliche Vortr\"age gehalten auf dem f\"unften Kongress der skandanivschen Mathematiker in Helsingfors vom 4.\ bis zum 7.\ Juli 1922}. पृ.~137--52. Akademiska Bokhandeln.

\par\medskip\phantomsection\label{bib:Smith2013}
Smith, Peter (2013). \emph{An Introduction to G\"odel's Theorems}. Cambridge University Press. Cambridge.

\par\medskip\phantomsection\label{bib:Smullyan1968}
Smullyan, Raymond M. (1968). \emph{First-Order Logic}. Springer. New York, NY. Corrected reprint, New York, NY: Dover, 1995.

\par\medskip\phantomsection\label{bib:FefermanLevy1963}
Solomon {Feferman}; Azriel {Levy} (1963). \emph{Independence results in set theory by {Cohen}'s method {II}}. \emph{Notices of the American Mathematical Society}, 10, 593.

\par\medskip\phantomsection\label{bib:Sykes1992}
Sykes, Christopher (1992). \emph{{BBC} {H}orizon: The Strange Life and Death of {D}r.~{T}uring}. \url{https://www.youtube.com/watch?v=gyusnGbBSHE}.

\par\medskip\phantomsection\label{bib:Gentzen1969}
Szabo, Manfred E. (1969). \emph{The Collected Papers of {G}erhard {G}entzen}. मध्ये: \emph{The Collected Papers of {G}erhard {G}entzen}. North-Holland. Amsterdam.

\par\medskip\phantomsection\label{bib:Takeuti2003}
Takeuti, Gaisi; Passell, Nicholas; Yasugi, Mariko (2003). \emph{Memoirs of a Proof Theorist: G\"odel and Other Logicians}. World Scientific. Singapore.

\par\medskip\phantomsection\label{bib:Tamassy1994}
Tamassy, Istvan (1994). \emph{Interview with {R}\'oza {P}\'eter}. \emph{Modern Logic}, 4(3), 277--280.

\par\medskip\phantomsection\label{bib:Sandstrum2019}
Ted Sandstrum (2019). \emph{Mathematical Reasoning: Writing and Proof}. Grand Valley State University. Allendale, MI. \url{https://scholarworks.gvsu.edu/books/7/}.

\par\medskip\phantomsection\label{bib:Theelen2012}
Theelen, Andre (2012). \emph{LEGO Turing Machine}. \url{https://www.youtube.com/watch?v=FTSAiF9AHN4}.

\par\medskip\phantomsection\label{bib:Vitali1905}
Vitali, Giuseppe (1905). \emph{Sul problema della misura dei gruppi di punti di una retta}. Gamberini e Parmeggiani. Bologna.

\par\medskip\phantomsection\label{bib:VonNeumann1925}
von Neumann, John (1925). \emph{Eine {Axiomatisierung} der {Mengenlehre}}. \emph{Journal f\"{u}r die reine und angewandte {Mathematik}}, 154, 219--40.

\par\medskip\phantomsection\label{bib:Wang1990}
Wang, Hao (1990). \emph{Reflections on Kurt G{\"{o}}del}. MIT Press. Cambridge.

\par\medskip\phantomsection\label{bib:Weston2003}
Weston, Tom (2003). \emph{The {Banach}-{Tarski} Paradox}. \url{http://people.math.umass.edu/~weston/oldpapers/banach.pdf}.

\par\medskip\phantomsection\label{bib:WhiteheadRussell1910}
Whitehead, Alfred North; Russell, Bertrand (1910). \emph{Principia Mathematica}. Cambridge University Press. Cambridge.

\par\medskip\phantomsection\label{bib:Zermelo1908}
Zermelo, Ernst (1908). \emph{Untersuchungen {\"u}ber die {G}rundlagen der {M}engenlehre~{I}}. \emph{Mathematische Annalen}, 65(2), 261--281. {E}nglish translation in \hyperref[bib:Ebbinghaus2010]{Ebbinghaus 2010, pp.~189-229}.

\par\medskip\phantomsection\label{bib:Zermelo1908Untersuchungen}
Zermelo, Ernst (1908). \emph{Untersuchungen \"{u}ber die {G}rundlagen der {M}engenlehre {I}}. \emph{Mathematische Annalen}, 65, 261--81.

\par\medskip\phantomsection\label{bib:Zuckerman1973}
Zuckerman, Martin M. (1973). \emph{Formation sequences for propositional formulas}. \emph{Notre Dame Journal of Formal Logic}, 14(1), 134--138.
\end{document}
